\chapter{Results} \label{chap:results}


In this chapter the results of the executed simulations are presented and discussed in terms of their relevance.
First, the simulations are verified by means of a comparison with the results of the flat \acrshort{tbl} simulations.
Next the results of case R100deg30 are analysed in terms of the optimisation capability.
The simulations of R100deg90 and R500deg18 are compared to analyse how different curvatures affect the \acrshort{tbl}.
Finally, a comparative analysis is performed on the optimised simulations of the three cases.


\section{flat section verification}\label{sec:flatTBLVerif}
The following section is employed as a reliability test of the governed results for the curved \acrshort{tbl} simulations.
As described in \autoref{sec:pressFluc}, we detected pressure fluctuations in the snapshots during the course of the simulations.
To ensure that these fluctuations leave the long-term statistics unaffected, we verified the statistics of the curved \acrshort{tbl} simulations in the straight inlet.
In order to undertake this task, we compare the statistics of case R100deg30 in optimised and base configurations with the ones from the flat \acrfull{tbl} setup and the reference data for flat \acrshort{tbl} simulations from \citet{schlatter_assessment_2010}, which has previously been introduced.
The statistics are taken at $\mathrm{Re}_{\theta} = 700$ in all simulations performed. This corresponds to a position of around $s_{Re_{\theta}= 700} = \frac{1}{2} s_p$.
Consequently, the impact of the curvature should be negligible at this position, and all setups should demonstrate analogous behaviour.
Following this, the results of this comparison are presented for case R100deg30. 
In order to circumvent the presentation of repetitive results, analogous analyses for the remaining curved configurations are not shown here.
%\\[-20pt]
\begin{figure}[h!]
\centering
\begin{subfigure}[t]{0.495\textwidth}
    \subcaption{}
    \centering
    %% Creator: Matplotlib, PGF backend
%%
%% To include the figure in your LaTeX document, write
%%   \input{<filename>.pgf}
%%
%% Make sure the required packages are loaded in your preamble
%%   \usepackage{pgf}
%%
%% Also ensure that all the required font packages are loaded; for instance,
%% the lmodern package is sometimes necessary when using math font.
%%   \usepackage{lmodern}
%%
%% Figures using additional raster images can only be included by \input if
%% they are in the same directory as the main LaTeX file. For loading figures
%% from other directories you can use the `import` package
%%   \usepackage{import}
%%
%% and then include the figures with
%%   \import{<path to file>}{<filename>.pgf}
%%
%% Matplotlib used the following preamble
%%   \def\mathdefault#1{#1}
%%   \everymath=\expandafter{\the\everymath\displaystyle}
%%   \IfFileExists{scrextend.sty}{
%%     \usepackage[fontsize=11.000000pt]{scrextend}
%%   }{
%%     \renewcommand{\normalsize}{\fontsize{11.000000}{13.200000}\selectfont}
%%     \normalsize
%%   }
%%   
%%   \makeatletter\@ifpackageloaded{underscore}{}{\usepackage[strings]{underscore}}\makeatother
%%
\begingroup%
\makeatletter%
\begin{pgfpicture}%
\pgfpathrectangle{\pgfpointorigin}{\pgfqpoint{3.118110in}{2.494488in}}%
\pgfusepath{use as bounding box, clip}%
\begin{pgfscope}%
\pgfsetbuttcap%
\pgfsetmiterjoin%
\definecolor{currentfill}{rgb}{1.000000,1.000000,1.000000}%
\pgfsetfillcolor{currentfill}%
\pgfsetlinewidth{0.000000pt}%
\definecolor{currentstroke}{rgb}{1.000000,1.000000,1.000000}%
\pgfsetstrokecolor{currentstroke}%
\pgfsetdash{}{0pt}%
\pgfpathmoveto{\pgfqpoint{0.000000in}{0.000000in}}%
\pgfpathlineto{\pgfqpoint{3.118110in}{0.000000in}}%
\pgfpathlineto{\pgfqpoint{3.118110in}{2.494488in}}%
\pgfpathlineto{\pgfqpoint{0.000000in}{2.494488in}}%
\pgfpathlineto{\pgfqpoint{0.000000in}{0.000000in}}%
\pgfpathclose%
\pgfusepath{fill}%
\end{pgfscope}%
\begin{pgfscope}%
\pgfsetbuttcap%
\pgfsetmiterjoin%
\definecolor{currentfill}{rgb}{1.000000,1.000000,1.000000}%
\pgfsetfillcolor{currentfill}%
\pgfsetlinewidth{0.000000pt}%
\definecolor{currentstroke}{rgb}{0.000000,0.000000,0.000000}%
\pgfsetstrokecolor{currentstroke}%
\pgfsetstrokeopacity{0.000000}%
\pgfsetdash{}{0pt}%
\pgfpathmoveto{\pgfqpoint{0.650000in}{0.420000in}}%
\pgfpathlineto{\pgfqpoint{3.038110in}{0.420000in}}%
\pgfpathlineto{\pgfqpoint{3.038110in}{2.414488in}}%
\pgfpathlineto{\pgfqpoint{0.650000in}{2.414488in}}%
\pgfpathlineto{\pgfqpoint{0.650000in}{0.420000in}}%
\pgfpathclose%
\pgfusepath{fill}%
\end{pgfscope}%
\begin{pgfscope}%
\pgfsetbuttcap%
\pgfsetroundjoin%
\definecolor{currentfill}{rgb}{0.000000,0.000000,0.000000}%
\pgfsetfillcolor{currentfill}%
\pgfsetlinewidth{0.501875pt}%
\definecolor{currentstroke}{rgb}{0.000000,0.000000,0.000000}%
\pgfsetstrokecolor{currentstroke}%
\pgfsetdash{}{0pt}%
\pgfsys@defobject{currentmarker}{\pgfqpoint{0.000000in}{0.000000in}}{\pgfqpoint{0.000000in}{0.041667in}}{%
\pgfpathmoveto{\pgfqpoint{0.000000in}{0.000000in}}%
\pgfpathlineto{\pgfqpoint{0.000000in}{0.041667in}}%
\pgfusepath{stroke,fill}%
}%
\begin{pgfscope}%
\pgfsys@transformshift{0.878501in}{0.420000in}%
\pgfsys@useobject{currentmarker}{}%
\end{pgfscope}%
\end{pgfscope}%
\begin{pgfscope}%
\pgfsetbuttcap%
\pgfsetroundjoin%
\definecolor{currentfill}{rgb}{0.000000,0.000000,0.000000}%
\pgfsetfillcolor{currentfill}%
\pgfsetlinewidth{0.501875pt}%
\definecolor{currentstroke}{rgb}{0.000000,0.000000,0.000000}%
\pgfsetstrokecolor{currentstroke}%
\pgfsetdash{}{0pt}%
\pgfsys@defobject{currentmarker}{\pgfqpoint{0.000000in}{-0.041667in}}{\pgfqpoint{0.000000in}{0.000000in}}{%
\pgfpathmoveto{\pgfqpoint{0.000000in}{0.000000in}}%
\pgfpathlineto{\pgfqpoint{0.000000in}{-0.041667in}}%
\pgfusepath{stroke,fill}%
}%
\begin{pgfscope}%
\pgfsys@transformshift{0.878501in}{2.414488in}%
\pgfsys@useobject{currentmarker}{}%
\end{pgfscope}%
\end{pgfscope}%
\begin{pgfscope}%
\definecolor{textcolor}{rgb}{0.000000,0.000000,0.000000}%
\pgfsetstrokecolor{textcolor}%
\pgfsetfillcolor{textcolor}%
\pgftext[x=0.878501in,y=0.371389in,,top]{\color{textcolor}{\rmfamily\fontsize{11.000000}{13.200000}\selectfont\catcode`\^=\active\def^{\ifmmode\sp\else\^{}\fi}\catcode`\%=\active\def%{\%}$\mathdefault{10^{0}}$}}%
\end{pgfscope}%
\begin{pgfscope}%
\pgfsetbuttcap%
\pgfsetroundjoin%
\definecolor{currentfill}{rgb}{0.000000,0.000000,0.000000}%
\pgfsetfillcolor{currentfill}%
\pgfsetlinewidth{0.501875pt}%
\definecolor{currentstroke}{rgb}{0.000000,0.000000,0.000000}%
\pgfsetstrokecolor{currentstroke}%
\pgfsetdash{}{0pt}%
\pgfsys@defobject{currentmarker}{\pgfqpoint{0.000000in}{0.000000in}}{\pgfqpoint{0.000000in}{0.041667in}}{%
\pgfpathmoveto{\pgfqpoint{0.000000in}{0.000000in}}%
\pgfpathlineto{\pgfqpoint{0.000000in}{0.041667in}}%
\pgfusepath{stroke,fill}%
}%
\begin{pgfscope}%
\pgfsys@transformshift{1.637564in}{0.420000in}%
\pgfsys@useobject{currentmarker}{}%
\end{pgfscope}%
\end{pgfscope}%
\begin{pgfscope}%
\pgfsetbuttcap%
\pgfsetroundjoin%
\definecolor{currentfill}{rgb}{0.000000,0.000000,0.000000}%
\pgfsetfillcolor{currentfill}%
\pgfsetlinewidth{0.501875pt}%
\definecolor{currentstroke}{rgb}{0.000000,0.000000,0.000000}%
\pgfsetstrokecolor{currentstroke}%
\pgfsetdash{}{0pt}%
\pgfsys@defobject{currentmarker}{\pgfqpoint{0.000000in}{-0.041667in}}{\pgfqpoint{0.000000in}{0.000000in}}{%
\pgfpathmoveto{\pgfqpoint{0.000000in}{0.000000in}}%
\pgfpathlineto{\pgfqpoint{0.000000in}{-0.041667in}}%
\pgfusepath{stroke,fill}%
}%
\begin{pgfscope}%
\pgfsys@transformshift{1.637564in}{2.414488in}%
\pgfsys@useobject{currentmarker}{}%
\end{pgfscope}%
\end{pgfscope}%
\begin{pgfscope}%
\definecolor{textcolor}{rgb}{0.000000,0.000000,0.000000}%
\pgfsetstrokecolor{textcolor}%
\pgfsetfillcolor{textcolor}%
\pgftext[x=1.637564in,y=0.371389in,,top]{\color{textcolor}{\rmfamily\fontsize{11.000000}{13.200000}\selectfont\catcode`\^=\active\def^{\ifmmode\sp\else\^{}\fi}\catcode`\%=\active\def%{\%}$\mathdefault{10^{1}}$}}%
\end{pgfscope}%
\begin{pgfscope}%
\pgfsetbuttcap%
\pgfsetroundjoin%
\definecolor{currentfill}{rgb}{0.000000,0.000000,0.000000}%
\pgfsetfillcolor{currentfill}%
\pgfsetlinewidth{0.501875pt}%
\definecolor{currentstroke}{rgb}{0.000000,0.000000,0.000000}%
\pgfsetstrokecolor{currentstroke}%
\pgfsetdash{}{0pt}%
\pgfsys@defobject{currentmarker}{\pgfqpoint{0.000000in}{0.000000in}}{\pgfqpoint{0.000000in}{0.041667in}}{%
\pgfpathmoveto{\pgfqpoint{0.000000in}{0.000000in}}%
\pgfpathlineto{\pgfqpoint{0.000000in}{0.041667in}}%
\pgfusepath{stroke,fill}%
}%
\begin{pgfscope}%
\pgfsys@transformshift{2.396627in}{0.420000in}%
\pgfsys@useobject{currentmarker}{}%
\end{pgfscope}%
\end{pgfscope}%
\begin{pgfscope}%
\pgfsetbuttcap%
\pgfsetroundjoin%
\definecolor{currentfill}{rgb}{0.000000,0.000000,0.000000}%
\pgfsetfillcolor{currentfill}%
\pgfsetlinewidth{0.501875pt}%
\definecolor{currentstroke}{rgb}{0.000000,0.000000,0.000000}%
\pgfsetstrokecolor{currentstroke}%
\pgfsetdash{}{0pt}%
\pgfsys@defobject{currentmarker}{\pgfqpoint{0.000000in}{-0.041667in}}{\pgfqpoint{0.000000in}{0.000000in}}{%
\pgfpathmoveto{\pgfqpoint{0.000000in}{0.000000in}}%
\pgfpathlineto{\pgfqpoint{0.000000in}{-0.041667in}}%
\pgfusepath{stroke,fill}%
}%
\begin{pgfscope}%
\pgfsys@transformshift{2.396627in}{2.414488in}%
\pgfsys@useobject{currentmarker}{}%
\end{pgfscope}%
\end{pgfscope}%
\begin{pgfscope}%
\definecolor{textcolor}{rgb}{0.000000,0.000000,0.000000}%
\pgfsetstrokecolor{textcolor}%
\pgfsetfillcolor{textcolor}%
\pgftext[x=2.396627in,y=0.371389in,,top]{\color{textcolor}{\rmfamily\fontsize{11.000000}{13.200000}\selectfont\catcode`\^=\active\def^{\ifmmode\sp\else\^{}\fi}\catcode`\%=\active\def%{\%}$\mathdefault{10^{2}}$}}%
\end{pgfscope}%
\begin{pgfscope}%
\pgfsetbuttcap%
\pgfsetroundjoin%
\definecolor{currentfill}{rgb}{0.000000,0.000000,0.000000}%
\pgfsetfillcolor{currentfill}%
\pgfsetlinewidth{0.401500pt}%
\definecolor{currentstroke}{rgb}{0.000000,0.000000,0.000000}%
\pgfsetstrokecolor{currentstroke}%
\pgfsetdash{}{0pt}%
\pgfsys@defobject{currentmarker}{\pgfqpoint{0.000000in}{0.000000in}}{\pgfqpoint{0.000000in}{0.020833in}}{%
\pgfpathmoveto{\pgfqpoint{0.000000in}{0.000000in}}%
\pgfpathlineto{\pgfqpoint{0.000000in}{0.020833in}}%
\pgfusepath{stroke,fill}%
}%
\begin{pgfscope}%
\pgfsys@transformshift{0.650000in}{0.420000in}%
\pgfsys@useobject{currentmarker}{}%
\end{pgfscope}%
\end{pgfscope}%
\begin{pgfscope}%
\pgfsetbuttcap%
\pgfsetroundjoin%
\definecolor{currentfill}{rgb}{0.000000,0.000000,0.000000}%
\pgfsetfillcolor{currentfill}%
\pgfsetlinewidth{0.401500pt}%
\definecolor{currentstroke}{rgb}{0.000000,0.000000,0.000000}%
\pgfsetstrokecolor{currentstroke}%
\pgfsetdash{}{0pt}%
\pgfsys@defobject{currentmarker}{\pgfqpoint{0.000000in}{-0.020833in}}{\pgfqpoint{0.000000in}{0.000000in}}{%
\pgfpathmoveto{\pgfqpoint{0.000000in}{0.000000in}}%
\pgfpathlineto{\pgfqpoint{0.000000in}{-0.020833in}}%
\pgfusepath{stroke,fill}%
}%
\begin{pgfscope}%
\pgfsys@transformshift{0.650000in}{2.414488in}%
\pgfsys@useobject{currentmarker}{}%
\end{pgfscope}%
\end{pgfscope}%
\begin{pgfscope}%
\pgfsetbuttcap%
\pgfsetroundjoin%
\definecolor{currentfill}{rgb}{0.000000,0.000000,0.000000}%
\pgfsetfillcolor{currentfill}%
\pgfsetlinewidth{0.401500pt}%
\definecolor{currentstroke}{rgb}{0.000000,0.000000,0.000000}%
\pgfsetstrokecolor{currentstroke}%
\pgfsetdash{}{0pt}%
\pgfsys@defobject{currentmarker}{\pgfqpoint{0.000000in}{0.000000in}}{\pgfqpoint{0.000000in}{0.020833in}}{%
\pgfpathmoveto{\pgfqpoint{0.000000in}{0.000000in}}%
\pgfpathlineto{\pgfqpoint{0.000000in}{0.020833in}}%
\pgfusepath{stroke,fill}%
}%
\begin{pgfscope}%
\pgfsys@transformshift{0.710104in}{0.420000in}%
\pgfsys@useobject{currentmarker}{}%
\end{pgfscope}%
\end{pgfscope}%
\begin{pgfscope}%
\pgfsetbuttcap%
\pgfsetroundjoin%
\definecolor{currentfill}{rgb}{0.000000,0.000000,0.000000}%
\pgfsetfillcolor{currentfill}%
\pgfsetlinewidth{0.401500pt}%
\definecolor{currentstroke}{rgb}{0.000000,0.000000,0.000000}%
\pgfsetstrokecolor{currentstroke}%
\pgfsetdash{}{0pt}%
\pgfsys@defobject{currentmarker}{\pgfqpoint{0.000000in}{-0.020833in}}{\pgfqpoint{0.000000in}{0.000000in}}{%
\pgfpathmoveto{\pgfqpoint{0.000000in}{0.000000in}}%
\pgfpathlineto{\pgfqpoint{0.000000in}{-0.020833in}}%
\pgfusepath{stroke,fill}%
}%
\begin{pgfscope}%
\pgfsys@transformshift{0.710104in}{2.414488in}%
\pgfsys@useobject{currentmarker}{}%
\end{pgfscope}%
\end{pgfscope}%
\begin{pgfscope}%
\pgfsetbuttcap%
\pgfsetroundjoin%
\definecolor{currentfill}{rgb}{0.000000,0.000000,0.000000}%
\pgfsetfillcolor{currentfill}%
\pgfsetlinewidth{0.401500pt}%
\definecolor{currentstroke}{rgb}{0.000000,0.000000,0.000000}%
\pgfsetstrokecolor{currentstroke}%
\pgfsetdash{}{0pt}%
\pgfsys@defobject{currentmarker}{\pgfqpoint{0.000000in}{0.000000in}}{\pgfqpoint{0.000000in}{0.020833in}}{%
\pgfpathmoveto{\pgfqpoint{0.000000in}{0.000000in}}%
\pgfpathlineto{\pgfqpoint{0.000000in}{0.020833in}}%
\pgfusepath{stroke,fill}%
}%
\begin{pgfscope}%
\pgfsys@transformshift{0.760920in}{0.420000in}%
\pgfsys@useobject{currentmarker}{}%
\end{pgfscope}%
\end{pgfscope}%
\begin{pgfscope}%
\pgfsetbuttcap%
\pgfsetroundjoin%
\definecolor{currentfill}{rgb}{0.000000,0.000000,0.000000}%
\pgfsetfillcolor{currentfill}%
\pgfsetlinewidth{0.401500pt}%
\definecolor{currentstroke}{rgb}{0.000000,0.000000,0.000000}%
\pgfsetstrokecolor{currentstroke}%
\pgfsetdash{}{0pt}%
\pgfsys@defobject{currentmarker}{\pgfqpoint{0.000000in}{-0.020833in}}{\pgfqpoint{0.000000in}{0.000000in}}{%
\pgfpathmoveto{\pgfqpoint{0.000000in}{0.000000in}}%
\pgfpathlineto{\pgfqpoint{0.000000in}{-0.020833in}}%
\pgfusepath{stroke,fill}%
}%
\begin{pgfscope}%
\pgfsys@transformshift{0.760920in}{2.414488in}%
\pgfsys@useobject{currentmarker}{}%
\end{pgfscope}%
\end{pgfscope}%
\begin{pgfscope}%
\pgfsetbuttcap%
\pgfsetroundjoin%
\definecolor{currentfill}{rgb}{0.000000,0.000000,0.000000}%
\pgfsetfillcolor{currentfill}%
\pgfsetlinewidth{0.401500pt}%
\definecolor{currentstroke}{rgb}{0.000000,0.000000,0.000000}%
\pgfsetstrokecolor{currentstroke}%
\pgfsetdash{}{0pt}%
\pgfsys@defobject{currentmarker}{\pgfqpoint{0.000000in}{0.000000in}}{\pgfqpoint{0.000000in}{0.020833in}}{%
\pgfpathmoveto{\pgfqpoint{0.000000in}{0.000000in}}%
\pgfpathlineto{\pgfqpoint{0.000000in}{0.020833in}}%
\pgfusepath{stroke,fill}%
}%
\begin{pgfscope}%
\pgfsys@transformshift{0.804940in}{0.420000in}%
\pgfsys@useobject{currentmarker}{}%
\end{pgfscope}%
\end{pgfscope}%
\begin{pgfscope}%
\pgfsetbuttcap%
\pgfsetroundjoin%
\definecolor{currentfill}{rgb}{0.000000,0.000000,0.000000}%
\pgfsetfillcolor{currentfill}%
\pgfsetlinewidth{0.401500pt}%
\definecolor{currentstroke}{rgb}{0.000000,0.000000,0.000000}%
\pgfsetstrokecolor{currentstroke}%
\pgfsetdash{}{0pt}%
\pgfsys@defobject{currentmarker}{\pgfqpoint{0.000000in}{-0.020833in}}{\pgfqpoint{0.000000in}{0.000000in}}{%
\pgfpathmoveto{\pgfqpoint{0.000000in}{0.000000in}}%
\pgfpathlineto{\pgfqpoint{0.000000in}{-0.020833in}}%
\pgfusepath{stroke,fill}%
}%
\begin{pgfscope}%
\pgfsys@transformshift{0.804940in}{2.414488in}%
\pgfsys@useobject{currentmarker}{}%
\end{pgfscope}%
\end{pgfscope}%
\begin{pgfscope}%
\pgfsetbuttcap%
\pgfsetroundjoin%
\definecolor{currentfill}{rgb}{0.000000,0.000000,0.000000}%
\pgfsetfillcolor{currentfill}%
\pgfsetlinewidth{0.401500pt}%
\definecolor{currentstroke}{rgb}{0.000000,0.000000,0.000000}%
\pgfsetstrokecolor{currentstroke}%
\pgfsetdash{}{0pt}%
\pgfsys@defobject{currentmarker}{\pgfqpoint{0.000000in}{0.000000in}}{\pgfqpoint{0.000000in}{0.020833in}}{%
\pgfpathmoveto{\pgfqpoint{0.000000in}{0.000000in}}%
\pgfpathlineto{\pgfqpoint{0.000000in}{0.020833in}}%
\pgfusepath{stroke,fill}%
}%
\begin{pgfscope}%
\pgfsys@transformshift{0.843768in}{0.420000in}%
\pgfsys@useobject{currentmarker}{}%
\end{pgfscope}%
\end{pgfscope}%
\begin{pgfscope}%
\pgfsetbuttcap%
\pgfsetroundjoin%
\definecolor{currentfill}{rgb}{0.000000,0.000000,0.000000}%
\pgfsetfillcolor{currentfill}%
\pgfsetlinewidth{0.401500pt}%
\definecolor{currentstroke}{rgb}{0.000000,0.000000,0.000000}%
\pgfsetstrokecolor{currentstroke}%
\pgfsetdash{}{0pt}%
\pgfsys@defobject{currentmarker}{\pgfqpoint{0.000000in}{-0.020833in}}{\pgfqpoint{0.000000in}{0.000000in}}{%
\pgfpathmoveto{\pgfqpoint{0.000000in}{0.000000in}}%
\pgfpathlineto{\pgfqpoint{0.000000in}{-0.020833in}}%
\pgfusepath{stroke,fill}%
}%
\begin{pgfscope}%
\pgfsys@transformshift{0.843768in}{2.414488in}%
\pgfsys@useobject{currentmarker}{}%
\end{pgfscope}%
\end{pgfscope}%
\begin{pgfscope}%
\pgfsetbuttcap%
\pgfsetroundjoin%
\definecolor{currentfill}{rgb}{0.000000,0.000000,0.000000}%
\pgfsetfillcolor{currentfill}%
\pgfsetlinewidth{0.401500pt}%
\definecolor{currentstroke}{rgb}{0.000000,0.000000,0.000000}%
\pgfsetstrokecolor{currentstroke}%
\pgfsetdash{}{0pt}%
\pgfsys@defobject{currentmarker}{\pgfqpoint{0.000000in}{0.000000in}}{\pgfqpoint{0.000000in}{0.020833in}}{%
\pgfpathmoveto{\pgfqpoint{0.000000in}{0.000000in}}%
\pgfpathlineto{\pgfqpoint{0.000000in}{0.020833in}}%
\pgfusepath{stroke,fill}%
}%
\begin{pgfscope}%
\pgfsys@transformshift{1.107002in}{0.420000in}%
\pgfsys@useobject{currentmarker}{}%
\end{pgfscope}%
\end{pgfscope}%
\begin{pgfscope}%
\pgfsetbuttcap%
\pgfsetroundjoin%
\definecolor{currentfill}{rgb}{0.000000,0.000000,0.000000}%
\pgfsetfillcolor{currentfill}%
\pgfsetlinewidth{0.401500pt}%
\definecolor{currentstroke}{rgb}{0.000000,0.000000,0.000000}%
\pgfsetstrokecolor{currentstroke}%
\pgfsetdash{}{0pt}%
\pgfsys@defobject{currentmarker}{\pgfqpoint{0.000000in}{-0.020833in}}{\pgfqpoint{0.000000in}{0.000000in}}{%
\pgfpathmoveto{\pgfqpoint{0.000000in}{0.000000in}}%
\pgfpathlineto{\pgfqpoint{0.000000in}{-0.020833in}}%
\pgfusepath{stroke,fill}%
}%
\begin{pgfscope}%
\pgfsys@transformshift{1.107002in}{2.414488in}%
\pgfsys@useobject{currentmarker}{}%
\end{pgfscope}%
\end{pgfscope}%
\begin{pgfscope}%
\pgfsetbuttcap%
\pgfsetroundjoin%
\definecolor{currentfill}{rgb}{0.000000,0.000000,0.000000}%
\pgfsetfillcolor{currentfill}%
\pgfsetlinewidth{0.401500pt}%
\definecolor{currentstroke}{rgb}{0.000000,0.000000,0.000000}%
\pgfsetstrokecolor{currentstroke}%
\pgfsetdash{}{0pt}%
\pgfsys@defobject{currentmarker}{\pgfqpoint{0.000000in}{0.000000in}}{\pgfqpoint{0.000000in}{0.020833in}}{%
\pgfpathmoveto{\pgfqpoint{0.000000in}{0.000000in}}%
\pgfpathlineto{\pgfqpoint{0.000000in}{0.020833in}}%
\pgfusepath{stroke,fill}%
}%
\begin{pgfscope}%
\pgfsys@transformshift{1.240666in}{0.420000in}%
\pgfsys@useobject{currentmarker}{}%
\end{pgfscope}%
\end{pgfscope}%
\begin{pgfscope}%
\pgfsetbuttcap%
\pgfsetroundjoin%
\definecolor{currentfill}{rgb}{0.000000,0.000000,0.000000}%
\pgfsetfillcolor{currentfill}%
\pgfsetlinewidth{0.401500pt}%
\definecolor{currentstroke}{rgb}{0.000000,0.000000,0.000000}%
\pgfsetstrokecolor{currentstroke}%
\pgfsetdash{}{0pt}%
\pgfsys@defobject{currentmarker}{\pgfqpoint{0.000000in}{-0.020833in}}{\pgfqpoint{0.000000in}{0.000000in}}{%
\pgfpathmoveto{\pgfqpoint{0.000000in}{0.000000in}}%
\pgfpathlineto{\pgfqpoint{0.000000in}{-0.020833in}}%
\pgfusepath{stroke,fill}%
}%
\begin{pgfscope}%
\pgfsys@transformshift{1.240666in}{2.414488in}%
\pgfsys@useobject{currentmarker}{}%
\end{pgfscope}%
\end{pgfscope}%
\begin{pgfscope}%
\pgfsetbuttcap%
\pgfsetroundjoin%
\definecolor{currentfill}{rgb}{0.000000,0.000000,0.000000}%
\pgfsetfillcolor{currentfill}%
\pgfsetlinewidth{0.401500pt}%
\definecolor{currentstroke}{rgb}{0.000000,0.000000,0.000000}%
\pgfsetstrokecolor{currentstroke}%
\pgfsetdash{}{0pt}%
\pgfsys@defobject{currentmarker}{\pgfqpoint{0.000000in}{0.000000in}}{\pgfqpoint{0.000000in}{0.020833in}}{%
\pgfpathmoveto{\pgfqpoint{0.000000in}{0.000000in}}%
\pgfpathlineto{\pgfqpoint{0.000000in}{0.020833in}}%
\pgfusepath{stroke,fill}%
}%
\begin{pgfscope}%
\pgfsys@transformshift{1.335502in}{0.420000in}%
\pgfsys@useobject{currentmarker}{}%
\end{pgfscope}%
\end{pgfscope}%
\begin{pgfscope}%
\pgfsetbuttcap%
\pgfsetroundjoin%
\definecolor{currentfill}{rgb}{0.000000,0.000000,0.000000}%
\pgfsetfillcolor{currentfill}%
\pgfsetlinewidth{0.401500pt}%
\definecolor{currentstroke}{rgb}{0.000000,0.000000,0.000000}%
\pgfsetstrokecolor{currentstroke}%
\pgfsetdash{}{0pt}%
\pgfsys@defobject{currentmarker}{\pgfqpoint{0.000000in}{-0.020833in}}{\pgfqpoint{0.000000in}{0.000000in}}{%
\pgfpathmoveto{\pgfqpoint{0.000000in}{0.000000in}}%
\pgfpathlineto{\pgfqpoint{0.000000in}{-0.020833in}}%
\pgfusepath{stroke,fill}%
}%
\begin{pgfscope}%
\pgfsys@transformshift{1.335502in}{2.414488in}%
\pgfsys@useobject{currentmarker}{}%
\end{pgfscope}%
\end{pgfscope}%
\begin{pgfscope}%
\pgfsetbuttcap%
\pgfsetroundjoin%
\definecolor{currentfill}{rgb}{0.000000,0.000000,0.000000}%
\pgfsetfillcolor{currentfill}%
\pgfsetlinewidth{0.401500pt}%
\definecolor{currentstroke}{rgb}{0.000000,0.000000,0.000000}%
\pgfsetstrokecolor{currentstroke}%
\pgfsetdash{}{0pt}%
\pgfsys@defobject{currentmarker}{\pgfqpoint{0.000000in}{0.000000in}}{\pgfqpoint{0.000000in}{0.020833in}}{%
\pgfpathmoveto{\pgfqpoint{0.000000in}{0.000000in}}%
\pgfpathlineto{\pgfqpoint{0.000000in}{0.020833in}}%
\pgfusepath{stroke,fill}%
}%
\begin{pgfscope}%
\pgfsys@transformshift{1.409063in}{0.420000in}%
\pgfsys@useobject{currentmarker}{}%
\end{pgfscope}%
\end{pgfscope}%
\begin{pgfscope}%
\pgfsetbuttcap%
\pgfsetroundjoin%
\definecolor{currentfill}{rgb}{0.000000,0.000000,0.000000}%
\pgfsetfillcolor{currentfill}%
\pgfsetlinewidth{0.401500pt}%
\definecolor{currentstroke}{rgb}{0.000000,0.000000,0.000000}%
\pgfsetstrokecolor{currentstroke}%
\pgfsetdash{}{0pt}%
\pgfsys@defobject{currentmarker}{\pgfqpoint{0.000000in}{-0.020833in}}{\pgfqpoint{0.000000in}{0.000000in}}{%
\pgfpathmoveto{\pgfqpoint{0.000000in}{0.000000in}}%
\pgfpathlineto{\pgfqpoint{0.000000in}{-0.020833in}}%
\pgfusepath{stroke,fill}%
}%
\begin{pgfscope}%
\pgfsys@transformshift{1.409063in}{2.414488in}%
\pgfsys@useobject{currentmarker}{}%
\end{pgfscope}%
\end{pgfscope}%
\begin{pgfscope}%
\pgfsetbuttcap%
\pgfsetroundjoin%
\definecolor{currentfill}{rgb}{0.000000,0.000000,0.000000}%
\pgfsetfillcolor{currentfill}%
\pgfsetlinewidth{0.401500pt}%
\definecolor{currentstroke}{rgb}{0.000000,0.000000,0.000000}%
\pgfsetstrokecolor{currentstroke}%
\pgfsetdash{}{0pt}%
\pgfsys@defobject{currentmarker}{\pgfqpoint{0.000000in}{0.000000in}}{\pgfqpoint{0.000000in}{0.020833in}}{%
\pgfpathmoveto{\pgfqpoint{0.000000in}{0.000000in}}%
\pgfpathlineto{\pgfqpoint{0.000000in}{0.020833in}}%
\pgfusepath{stroke,fill}%
}%
\begin{pgfscope}%
\pgfsys@transformshift{1.469167in}{0.420000in}%
\pgfsys@useobject{currentmarker}{}%
\end{pgfscope}%
\end{pgfscope}%
\begin{pgfscope}%
\pgfsetbuttcap%
\pgfsetroundjoin%
\definecolor{currentfill}{rgb}{0.000000,0.000000,0.000000}%
\pgfsetfillcolor{currentfill}%
\pgfsetlinewidth{0.401500pt}%
\definecolor{currentstroke}{rgb}{0.000000,0.000000,0.000000}%
\pgfsetstrokecolor{currentstroke}%
\pgfsetdash{}{0pt}%
\pgfsys@defobject{currentmarker}{\pgfqpoint{0.000000in}{-0.020833in}}{\pgfqpoint{0.000000in}{0.000000in}}{%
\pgfpathmoveto{\pgfqpoint{0.000000in}{0.000000in}}%
\pgfpathlineto{\pgfqpoint{0.000000in}{-0.020833in}}%
\pgfusepath{stroke,fill}%
}%
\begin{pgfscope}%
\pgfsys@transformshift{1.469167in}{2.414488in}%
\pgfsys@useobject{currentmarker}{}%
\end{pgfscope}%
\end{pgfscope}%
\begin{pgfscope}%
\pgfsetbuttcap%
\pgfsetroundjoin%
\definecolor{currentfill}{rgb}{0.000000,0.000000,0.000000}%
\pgfsetfillcolor{currentfill}%
\pgfsetlinewidth{0.401500pt}%
\definecolor{currentstroke}{rgb}{0.000000,0.000000,0.000000}%
\pgfsetstrokecolor{currentstroke}%
\pgfsetdash{}{0pt}%
\pgfsys@defobject{currentmarker}{\pgfqpoint{0.000000in}{0.000000in}}{\pgfqpoint{0.000000in}{0.020833in}}{%
\pgfpathmoveto{\pgfqpoint{0.000000in}{0.000000in}}%
\pgfpathlineto{\pgfqpoint{0.000000in}{0.020833in}}%
\pgfusepath{stroke,fill}%
}%
\begin{pgfscope}%
\pgfsys@transformshift{1.519984in}{0.420000in}%
\pgfsys@useobject{currentmarker}{}%
\end{pgfscope}%
\end{pgfscope}%
\begin{pgfscope}%
\pgfsetbuttcap%
\pgfsetroundjoin%
\definecolor{currentfill}{rgb}{0.000000,0.000000,0.000000}%
\pgfsetfillcolor{currentfill}%
\pgfsetlinewidth{0.401500pt}%
\definecolor{currentstroke}{rgb}{0.000000,0.000000,0.000000}%
\pgfsetstrokecolor{currentstroke}%
\pgfsetdash{}{0pt}%
\pgfsys@defobject{currentmarker}{\pgfqpoint{0.000000in}{-0.020833in}}{\pgfqpoint{0.000000in}{0.000000in}}{%
\pgfpathmoveto{\pgfqpoint{0.000000in}{0.000000in}}%
\pgfpathlineto{\pgfqpoint{0.000000in}{-0.020833in}}%
\pgfusepath{stroke,fill}%
}%
\begin{pgfscope}%
\pgfsys@transformshift{1.519984in}{2.414488in}%
\pgfsys@useobject{currentmarker}{}%
\end{pgfscope}%
\end{pgfscope}%
\begin{pgfscope}%
\pgfsetbuttcap%
\pgfsetroundjoin%
\definecolor{currentfill}{rgb}{0.000000,0.000000,0.000000}%
\pgfsetfillcolor{currentfill}%
\pgfsetlinewidth{0.401500pt}%
\definecolor{currentstroke}{rgb}{0.000000,0.000000,0.000000}%
\pgfsetstrokecolor{currentstroke}%
\pgfsetdash{}{0pt}%
\pgfsys@defobject{currentmarker}{\pgfqpoint{0.000000in}{0.000000in}}{\pgfqpoint{0.000000in}{0.020833in}}{%
\pgfpathmoveto{\pgfqpoint{0.000000in}{0.000000in}}%
\pgfpathlineto{\pgfqpoint{0.000000in}{0.020833in}}%
\pgfusepath{stroke,fill}%
}%
\begin{pgfscope}%
\pgfsys@transformshift{1.564003in}{0.420000in}%
\pgfsys@useobject{currentmarker}{}%
\end{pgfscope}%
\end{pgfscope}%
\begin{pgfscope}%
\pgfsetbuttcap%
\pgfsetroundjoin%
\definecolor{currentfill}{rgb}{0.000000,0.000000,0.000000}%
\pgfsetfillcolor{currentfill}%
\pgfsetlinewidth{0.401500pt}%
\definecolor{currentstroke}{rgb}{0.000000,0.000000,0.000000}%
\pgfsetstrokecolor{currentstroke}%
\pgfsetdash{}{0pt}%
\pgfsys@defobject{currentmarker}{\pgfqpoint{0.000000in}{-0.020833in}}{\pgfqpoint{0.000000in}{0.000000in}}{%
\pgfpathmoveto{\pgfqpoint{0.000000in}{0.000000in}}%
\pgfpathlineto{\pgfqpoint{0.000000in}{-0.020833in}}%
\pgfusepath{stroke,fill}%
}%
\begin{pgfscope}%
\pgfsys@transformshift{1.564003in}{2.414488in}%
\pgfsys@useobject{currentmarker}{}%
\end{pgfscope}%
\end{pgfscope}%
\begin{pgfscope}%
\pgfsetbuttcap%
\pgfsetroundjoin%
\definecolor{currentfill}{rgb}{0.000000,0.000000,0.000000}%
\pgfsetfillcolor{currentfill}%
\pgfsetlinewidth{0.401500pt}%
\definecolor{currentstroke}{rgb}{0.000000,0.000000,0.000000}%
\pgfsetstrokecolor{currentstroke}%
\pgfsetdash{}{0pt}%
\pgfsys@defobject{currentmarker}{\pgfqpoint{0.000000in}{0.000000in}}{\pgfqpoint{0.000000in}{0.020833in}}{%
\pgfpathmoveto{\pgfqpoint{0.000000in}{0.000000in}}%
\pgfpathlineto{\pgfqpoint{0.000000in}{0.020833in}}%
\pgfusepath{stroke,fill}%
}%
\begin{pgfscope}%
\pgfsys@transformshift{1.602831in}{0.420000in}%
\pgfsys@useobject{currentmarker}{}%
\end{pgfscope}%
\end{pgfscope}%
\begin{pgfscope}%
\pgfsetbuttcap%
\pgfsetroundjoin%
\definecolor{currentfill}{rgb}{0.000000,0.000000,0.000000}%
\pgfsetfillcolor{currentfill}%
\pgfsetlinewidth{0.401500pt}%
\definecolor{currentstroke}{rgb}{0.000000,0.000000,0.000000}%
\pgfsetstrokecolor{currentstroke}%
\pgfsetdash{}{0pt}%
\pgfsys@defobject{currentmarker}{\pgfqpoint{0.000000in}{-0.020833in}}{\pgfqpoint{0.000000in}{0.000000in}}{%
\pgfpathmoveto{\pgfqpoint{0.000000in}{0.000000in}}%
\pgfpathlineto{\pgfqpoint{0.000000in}{-0.020833in}}%
\pgfusepath{stroke,fill}%
}%
\begin{pgfscope}%
\pgfsys@transformshift{1.602831in}{2.414488in}%
\pgfsys@useobject{currentmarker}{}%
\end{pgfscope}%
\end{pgfscope}%
\begin{pgfscope}%
\pgfsetbuttcap%
\pgfsetroundjoin%
\definecolor{currentfill}{rgb}{0.000000,0.000000,0.000000}%
\pgfsetfillcolor{currentfill}%
\pgfsetlinewidth{0.401500pt}%
\definecolor{currentstroke}{rgb}{0.000000,0.000000,0.000000}%
\pgfsetstrokecolor{currentstroke}%
\pgfsetdash{}{0pt}%
\pgfsys@defobject{currentmarker}{\pgfqpoint{0.000000in}{0.000000in}}{\pgfqpoint{0.000000in}{0.020833in}}{%
\pgfpathmoveto{\pgfqpoint{0.000000in}{0.000000in}}%
\pgfpathlineto{\pgfqpoint{0.000000in}{0.020833in}}%
\pgfusepath{stroke,fill}%
}%
\begin{pgfscope}%
\pgfsys@transformshift{1.866065in}{0.420000in}%
\pgfsys@useobject{currentmarker}{}%
\end{pgfscope}%
\end{pgfscope}%
\begin{pgfscope}%
\pgfsetbuttcap%
\pgfsetroundjoin%
\definecolor{currentfill}{rgb}{0.000000,0.000000,0.000000}%
\pgfsetfillcolor{currentfill}%
\pgfsetlinewidth{0.401500pt}%
\definecolor{currentstroke}{rgb}{0.000000,0.000000,0.000000}%
\pgfsetstrokecolor{currentstroke}%
\pgfsetdash{}{0pt}%
\pgfsys@defobject{currentmarker}{\pgfqpoint{0.000000in}{-0.020833in}}{\pgfqpoint{0.000000in}{0.000000in}}{%
\pgfpathmoveto{\pgfqpoint{0.000000in}{0.000000in}}%
\pgfpathlineto{\pgfqpoint{0.000000in}{-0.020833in}}%
\pgfusepath{stroke,fill}%
}%
\begin{pgfscope}%
\pgfsys@transformshift{1.866065in}{2.414488in}%
\pgfsys@useobject{currentmarker}{}%
\end{pgfscope}%
\end{pgfscope}%
\begin{pgfscope}%
\pgfsetbuttcap%
\pgfsetroundjoin%
\definecolor{currentfill}{rgb}{0.000000,0.000000,0.000000}%
\pgfsetfillcolor{currentfill}%
\pgfsetlinewidth{0.401500pt}%
\definecolor{currentstroke}{rgb}{0.000000,0.000000,0.000000}%
\pgfsetstrokecolor{currentstroke}%
\pgfsetdash{}{0pt}%
\pgfsys@defobject{currentmarker}{\pgfqpoint{0.000000in}{0.000000in}}{\pgfqpoint{0.000000in}{0.020833in}}{%
\pgfpathmoveto{\pgfqpoint{0.000000in}{0.000000in}}%
\pgfpathlineto{\pgfqpoint{0.000000in}{0.020833in}}%
\pgfusepath{stroke,fill}%
}%
\begin{pgfscope}%
\pgfsys@transformshift{1.999729in}{0.420000in}%
\pgfsys@useobject{currentmarker}{}%
\end{pgfscope}%
\end{pgfscope}%
\begin{pgfscope}%
\pgfsetbuttcap%
\pgfsetroundjoin%
\definecolor{currentfill}{rgb}{0.000000,0.000000,0.000000}%
\pgfsetfillcolor{currentfill}%
\pgfsetlinewidth{0.401500pt}%
\definecolor{currentstroke}{rgb}{0.000000,0.000000,0.000000}%
\pgfsetstrokecolor{currentstroke}%
\pgfsetdash{}{0pt}%
\pgfsys@defobject{currentmarker}{\pgfqpoint{0.000000in}{-0.020833in}}{\pgfqpoint{0.000000in}{0.000000in}}{%
\pgfpathmoveto{\pgfqpoint{0.000000in}{0.000000in}}%
\pgfpathlineto{\pgfqpoint{0.000000in}{-0.020833in}}%
\pgfusepath{stroke,fill}%
}%
\begin{pgfscope}%
\pgfsys@transformshift{1.999729in}{2.414488in}%
\pgfsys@useobject{currentmarker}{}%
\end{pgfscope}%
\end{pgfscope}%
\begin{pgfscope}%
\pgfsetbuttcap%
\pgfsetroundjoin%
\definecolor{currentfill}{rgb}{0.000000,0.000000,0.000000}%
\pgfsetfillcolor{currentfill}%
\pgfsetlinewidth{0.401500pt}%
\definecolor{currentstroke}{rgb}{0.000000,0.000000,0.000000}%
\pgfsetstrokecolor{currentstroke}%
\pgfsetdash{}{0pt}%
\pgfsys@defobject{currentmarker}{\pgfqpoint{0.000000in}{0.000000in}}{\pgfqpoint{0.000000in}{0.020833in}}{%
\pgfpathmoveto{\pgfqpoint{0.000000in}{0.000000in}}%
\pgfpathlineto{\pgfqpoint{0.000000in}{0.020833in}}%
\pgfusepath{stroke,fill}%
}%
\begin{pgfscope}%
\pgfsys@transformshift{2.094566in}{0.420000in}%
\pgfsys@useobject{currentmarker}{}%
\end{pgfscope}%
\end{pgfscope}%
\begin{pgfscope}%
\pgfsetbuttcap%
\pgfsetroundjoin%
\definecolor{currentfill}{rgb}{0.000000,0.000000,0.000000}%
\pgfsetfillcolor{currentfill}%
\pgfsetlinewidth{0.401500pt}%
\definecolor{currentstroke}{rgb}{0.000000,0.000000,0.000000}%
\pgfsetstrokecolor{currentstroke}%
\pgfsetdash{}{0pt}%
\pgfsys@defobject{currentmarker}{\pgfqpoint{0.000000in}{-0.020833in}}{\pgfqpoint{0.000000in}{0.000000in}}{%
\pgfpathmoveto{\pgfqpoint{0.000000in}{0.000000in}}%
\pgfpathlineto{\pgfqpoint{0.000000in}{-0.020833in}}%
\pgfusepath{stroke,fill}%
}%
\begin{pgfscope}%
\pgfsys@transformshift{2.094566in}{2.414488in}%
\pgfsys@useobject{currentmarker}{}%
\end{pgfscope}%
\end{pgfscope}%
\begin{pgfscope}%
\pgfsetbuttcap%
\pgfsetroundjoin%
\definecolor{currentfill}{rgb}{0.000000,0.000000,0.000000}%
\pgfsetfillcolor{currentfill}%
\pgfsetlinewidth{0.401500pt}%
\definecolor{currentstroke}{rgb}{0.000000,0.000000,0.000000}%
\pgfsetstrokecolor{currentstroke}%
\pgfsetdash{}{0pt}%
\pgfsys@defobject{currentmarker}{\pgfqpoint{0.000000in}{0.000000in}}{\pgfqpoint{0.000000in}{0.020833in}}{%
\pgfpathmoveto{\pgfqpoint{0.000000in}{0.000000in}}%
\pgfpathlineto{\pgfqpoint{0.000000in}{0.020833in}}%
\pgfusepath{stroke,fill}%
}%
\begin{pgfscope}%
\pgfsys@transformshift{2.168126in}{0.420000in}%
\pgfsys@useobject{currentmarker}{}%
\end{pgfscope}%
\end{pgfscope}%
\begin{pgfscope}%
\pgfsetbuttcap%
\pgfsetroundjoin%
\definecolor{currentfill}{rgb}{0.000000,0.000000,0.000000}%
\pgfsetfillcolor{currentfill}%
\pgfsetlinewidth{0.401500pt}%
\definecolor{currentstroke}{rgb}{0.000000,0.000000,0.000000}%
\pgfsetstrokecolor{currentstroke}%
\pgfsetdash{}{0pt}%
\pgfsys@defobject{currentmarker}{\pgfqpoint{0.000000in}{-0.020833in}}{\pgfqpoint{0.000000in}{0.000000in}}{%
\pgfpathmoveto{\pgfqpoint{0.000000in}{0.000000in}}%
\pgfpathlineto{\pgfqpoint{0.000000in}{-0.020833in}}%
\pgfusepath{stroke,fill}%
}%
\begin{pgfscope}%
\pgfsys@transformshift{2.168126in}{2.414488in}%
\pgfsys@useobject{currentmarker}{}%
\end{pgfscope}%
\end{pgfscope}%
\begin{pgfscope}%
\pgfsetbuttcap%
\pgfsetroundjoin%
\definecolor{currentfill}{rgb}{0.000000,0.000000,0.000000}%
\pgfsetfillcolor{currentfill}%
\pgfsetlinewidth{0.401500pt}%
\definecolor{currentstroke}{rgb}{0.000000,0.000000,0.000000}%
\pgfsetstrokecolor{currentstroke}%
\pgfsetdash{}{0pt}%
\pgfsys@defobject{currentmarker}{\pgfqpoint{0.000000in}{0.000000in}}{\pgfqpoint{0.000000in}{0.020833in}}{%
\pgfpathmoveto{\pgfqpoint{0.000000in}{0.000000in}}%
\pgfpathlineto{\pgfqpoint{0.000000in}{0.020833in}}%
\pgfusepath{stroke,fill}%
}%
\begin{pgfscope}%
\pgfsys@transformshift{2.228230in}{0.420000in}%
\pgfsys@useobject{currentmarker}{}%
\end{pgfscope}%
\end{pgfscope}%
\begin{pgfscope}%
\pgfsetbuttcap%
\pgfsetroundjoin%
\definecolor{currentfill}{rgb}{0.000000,0.000000,0.000000}%
\pgfsetfillcolor{currentfill}%
\pgfsetlinewidth{0.401500pt}%
\definecolor{currentstroke}{rgb}{0.000000,0.000000,0.000000}%
\pgfsetstrokecolor{currentstroke}%
\pgfsetdash{}{0pt}%
\pgfsys@defobject{currentmarker}{\pgfqpoint{0.000000in}{-0.020833in}}{\pgfqpoint{0.000000in}{0.000000in}}{%
\pgfpathmoveto{\pgfqpoint{0.000000in}{0.000000in}}%
\pgfpathlineto{\pgfqpoint{0.000000in}{-0.020833in}}%
\pgfusepath{stroke,fill}%
}%
\begin{pgfscope}%
\pgfsys@transformshift{2.228230in}{2.414488in}%
\pgfsys@useobject{currentmarker}{}%
\end{pgfscope}%
\end{pgfscope}%
\begin{pgfscope}%
\pgfsetbuttcap%
\pgfsetroundjoin%
\definecolor{currentfill}{rgb}{0.000000,0.000000,0.000000}%
\pgfsetfillcolor{currentfill}%
\pgfsetlinewidth{0.401500pt}%
\definecolor{currentstroke}{rgb}{0.000000,0.000000,0.000000}%
\pgfsetstrokecolor{currentstroke}%
\pgfsetdash{}{0pt}%
\pgfsys@defobject{currentmarker}{\pgfqpoint{0.000000in}{0.000000in}}{\pgfqpoint{0.000000in}{0.020833in}}{%
\pgfpathmoveto{\pgfqpoint{0.000000in}{0.000000in}}%
\pgfpathlineto{\pgfqpoint{0.000000in}{0.020833in}}%
\pgfusepath{stroke,fill}%
}%
\begin{pgfscope}%
\pgfsys@transformshift{2.279047in}{0.420000in}%
\pgfsys@useobject{currentmarker}{}%
\end{pgfscope}%
\end{pgfscope}%
\begin{pgfscope}%
\pgfsetbuttcap%
\pgfsetroundjoin%
\definecolor{currentfill}{rgb}{0.000000,0.000000,0.000000}%
\pgfsetfillcolor{currentfill}%
\pgfsetlinewidth{0.401500pt}%
\definecolor{currentstroke}{rgb}{0.000000,0.000000,0.000000}%
\pgfsetstrokecolor{currentstroke}%
\pgfsetdash{}{0pt}%
\pgfsys@defobject{currentmarker}{\pgfqpoint{0.000000in}{-0.020833in}}{\pgfqpoint{0.000000in}{0.000000in}}{%
\pgfpathmoveto{\pgfqpoint{0.000000in}{0.000000in}}%
\pgfpathlineto{\pgfqpoint{0.000000in}{-0.020833in}}%
\pgfusepath{stroke,fill}%
}%
\begin{pgfscope}%
\pgfsys@transformshift{2.279047in}{2.414488in}%
\pgfsys@useobject{currentmarker}{}%
\end{pgfscope}%
\end{pgfscope}%
\begin{pgfscope}%
\pgfsetbuttcap%
\pgfsetroundjoin%
\definecolor{currentfill}{rgb}{0.000000,0.000000,0.000000}%
\pgfsetfillcolor{currentfill}%
\pgfsetlinewidth{0.401500pt}%
\definecolor{currentstroke}{rgb}{0.000000,0.000000,0.000000}%
\pgfsetstrokecolor{currentstroke}%
\pgfsetdash{}{0pt}%
\pgfsys@defobject{currentmarker}{\pgfqpoint{0.000000in}{0.000000in}}{\pgfqpoint{0.000000in}{0.020833in}}{%
\pgfpathmoveto{\pgfqpoint{0.000000in}{0.000000in}}%
\pgfpathlineto{\pgfqpoint{0.000000in}{0.020833in}}%
\pgfusepath{stroke,fill}%
}%
\begin{pgfscope}%
\pgfsys@transformshift{2.323066in}{0.420000in}%
\pgfsys@useobject{currentmarker}{}%
\end{pgfscope}%
\end{pgfscope}%
\begin{pgfscope}%
\pgfsetbuttcap%
\pgfsetroundjoin%
\definecolor{currentfill}{rgb}{0.000000,0.000000,0.000000}%
\pgfsetfillcolor{currentfill}%
\pgfsetlinewidth{0.401500pt}%
\definecolor{currentstroke}{rgb}{0.000000,0.000000,0.000000}%
\pgfsetstrokecolor{currentstroke}%
\pgfsetdash{}{0pt}%
\pgfsys@defobject{currentmarker}{\pgfqpoint{0.000000in}{-0.020833in}}{\pgfqpoint{0.000000in}{0.000000in}}{%
\pgfpathmoveto{\pgfqpoint{0.000000in}{0.000000in}}%
\pgfpathlineto{\pgfqpoint{0.000000in}{-0.020833in}}%
\pgfusepath{stroke,fill}%
}%
\begin{pgfscope}%
\pgfsys@transformshift{2.323066in}{2.414488in}%
\pgfsys@useobject{currentmarker}{}%
\end{pgfscope}%
\end{pgfscope}%
\begin{pgfscope}%
\pgfsetbuttcap%
\pgfsetroundjoin%
\definecolor{currentfill}{rgb}{0.000000,0.000000,0.000000}%
\pgfsetfillcolor{currentfill}%
\pgfsetlinewidth{0.401500pt}%
\definecolor{currentstroke}{rgb}{0.000000,0.000000,0.000000}%
\pgfsetstrokecolor{currentstroke}%
\pgfsetdash{}{0pt}%
\pgfsys@defobject{currentmarker}{\pgfqpoint{0.000000in}{0.000000in}}{\pgfqpoint{0.000000in}{0.020833in}}{%
\pgfpathmoveto{\pgfqpoint{0.000000in}{0.000000in}}%
\pgfpathlineto{\pgfqpoint{0.000000in}{0.020833in}}%
\pgfusepath{stroke,fill}%
}%
\begin{pgfscope}%
\pgfsys@transformshift{2.361894in}{0.420000in}%
\pgfsys@useobject{currentmarker}{}%
\end{pgfscope}%
\end{pgfscope}%
\begin{pgfscope}%
\pgfsetbuttcap%
\pgfsetroundjoin%
\definecolor{currentfill}{rgb}{0.000000,0.000000,0.000000}%
\pgfsetfillcolor{currentfill}%
\pgfsetlinewidth{0.401500pt}%
\definecolor{currentstroke}{rgb}{0.000000,0.000000,0.000000}%
\pgfsetstrokecolor{currentstroke}%
\pgfsetdash{}{0pt}%
\pgfsys@defobject{currentmarker}{\pgfqpoint{0.000000in}{-0.020833in}}{\pgfqpoint{0.000000in}{0.000000in}}{%
\pgfpathmoveto{\pgfqpoint{0.000000in}{0.000000in}}%
\pgfpathlineto{\pgfqpoint{0.000000in}{-0.020833in}}%
\pgfusepath{stroke,fill}%
}%
\begin{pgfscope}%
\pgfsys@transformshift{2.361894in}{2.414488in}%
\pgfsys@useobject{currentmarker}{}%
\end{pgfscope}%
\end{pgfscope}%
\begin{pgfscope}%
\pgfsetbuttcap%
\pgfsetroundjoin%
\definecolor{currentfill}{rgb}{0.000000,0.000000,0.000000}%
\pgfsetfillcolor{currentfill}%
\pgfsetlinewidth{0.401500pt}%
\definecolor{currentstroke}{rgb}{0.000000,0.000000,0.000000}%
\pgfsetstrokecolor{currentstroke}%
\pgfsetdash{}{0pt}%
\pgfsys@defobject{currentmarker}{\pgfqpoint{0.000000in}{0.000000in}}{\pgfqpoint{0.000000in}{0.020833in}}{%
\pgfpathmoveto{\pgfqpoint{0.000000in}{0.000000in}}%
\pgfpathlineto{\pgfqpoint{0.000000in}{0.020833in}}%
\pgfusepath{stroke,fill}%
}%
\begin{pgfscope}%
\pgfsys@transformshift{2.625128in}{0.420000in}%
\pgfsys@useobject{currentmarker}{}%
\end{pgfscope}%
\end{pgfscope}%
\begin{pgfscope}%
\pgfsetbuttcap%
\pgfsetroundjoin%
\definecolor{currentfill}{rgb}{0.000000,0.000000,0.000000}%
\pgfsetfillcolor{currentfill}%
\pgfsetlinewidth{0.401500pt}%
\definecolor{currentstroke}{rgb}{0.000000,0.000000,0.000000}%
\pgfsetstrokecolor{currentstroke}%
\pgfsetdash{}{0pt}%
\pgfsys@defobject{currentmarker}{\pgfqpoint{0.000000in}{-0.020833in}}{\pgfqpoint{0.000000in}{0.000000in}}{%
\pgfpathmoveto{\pgfqpoint{0.000000in}{0.000000in}}%
\pgfpathlineto{\pgfqpoint{0.000000in}{-0.020833in}}%
\pgfusepath{stroke,fill}%
}%
\begin{pgfscope}%
\pgfsys@transformshift{2.625128in}{2.414488in}%
\pgfsys@useobject{currentmarker}{}%
\end{pgfscope}%
\end{pgfscope}%
\begin{pgfscope}%
\pgfsetbuttcap%
\pgfsetroundjoin%
\definecolor{currentfill}{rgb}{0.000000,0.000000,0.000000}%
\pgfsetfillcolor{currentfill}%
\pgfsetlinewidth{0.401500pt}%
\definecolor{currentstroke}{rgb}{0.000000,0.000000,0.000000}%
\pgfsetstrokecolor{currentstroke}%
\pgfsetdash{}{0pt}%
\pgfsys@defobject{currentmarker}{\pgfqpoint{0.000000in}{0.000000in}}{\pgfqpoint{0.000000in}{0.020833in}}{%
\pgfpathmoveto{\pgfqpoint{0.000000in}{0.000000in}}%
\pgfpathlineto{\pgfqpoint{0.000000in}{0.020833in}}%
\pgfusepath{stroke,fill}%
}%
\begin{pgfscope}%
\pgfsys@transformshift{2.758792in}{0.420000in}%
\pgfsys@useobject{currentmarker}{}%
\end{pgfscope}%
\end{pgfscope}%
\begin{pgfscope}%
\pgfsetbuttcap%
\pgfsetroundjoin%
\definecolor{currentfill}{rgb}{0.000000,0.000000,0.000000}%
\pgfsetfillcolor{currentfill}%
\pgfsetlinewidth{0.401500pt}%
\definecolor{currentstroke}{rgb}{0.000000,0.000000,0.000000}%
\pgfsetstrokecolor{currentstroke}%
\pgfsetdash{}{0pt}%
\pgfsys@defobject{currentmarker}{\pgfqpoint{0.000000in}{-0.020833in}}{\pgfqpoint{0.000000in}{0.000000in}}{%
\pgfpathmoveto{\pgfqpoint{0.000000in}{0.000000in}}%
\pgfpathlineto{\pgfqpoint{0.000000in}{-0.020833in}}%
\pgfusepath{stroke,fill}%
}%
\begin{pgfscope}%
\pgfsys@transformshift{2.758792in}{2.414488in}%
\pgfsys@useobject{currentmarker}{}%
\end{pgfscope}%
\end{pgfscope}%
\begin{pgfscope}%
\pgfsetbuttcap%
\pgfsetroundjoin%
\definecolor{currentfill}{rgb}{0.000000,0.000000,0.000000}%
\pgfsetfillcolor{currentfill}%
\pgfsetlinewidth{0.401500pt}%
\definecolor{currentstroke}{rgb}{0.000000,0.000000,0.000000}%
\pgfsetstrokecolor{currentstroke}%
\pgfsetdash{}{0pt}%
\pgfsys@defobject{currentmarker}{\pgfqpoint{0.000000in}{0.000000in}}{\pgfqpoint{0.000000in}{0.020833in}}{%
\pgfpathmoveto{\pgfqpoint{0.000000in}{0.000000in}}%
\pgfpathlineto{\pgfqpoint{0.000000in}{0.020833in}}%
\pgfusepath{stroke,fill}%
}%
\begin{pgfscope}%
\pgfsys@transformshift{2.853629in}{0.420000in}%
\pgfsys@useobject{currentmarker}{}%
\end{pgfscope}%
\end{pgfscope}%
\begin{pgfscope}%
\pgfsetbuttcap%
\pgfsetroundjoin%
\definecolor{currentfill}{rgb}{0.000000,0.000000,0.000000}%
\pgfsetfillcolor{currentfill}%
\pgfsetlinewidth{0.401500pt}%
\definecolor{currentstroke}{rgb}{0.000000,0.000000,0.000000}%
\pgfsetstrokecolor{currentstroke}%
\pgfsetdash{}{0pt}%
\pgfsys@defobject{currentmarker}{\pgfqpoint{0.000000in}{-0.020833in}}{\pgfqpoint{0.000000in}{0.000000in}}{%
\pgfpathmoveto{\pgfqpoint{0.000000in}{0.000000in}}%
\pgfpathlineto{\pgfqpoint{0.000000in}{-0.020833in}}%
\pgfusepath{stroke,fill}%
}%
\begin{pgfscope}%
\pgfsys@transformshift{2.853629in}{2.414488in}%
\pgfsys@useobject{currentmarker}{}%
\end{pgfscope}%
\end{pgfscope}%
\begin{pgfscope}%
\pgfsetbuttcap%
\pgfsetroundjoin%
\definecolor{currentfill}{rgb}{0.000000,0.000000,0.000000}%
\pgfsetfillcolor{currentfill}%
\pgfsetlinewidth{0.401500pt}%
\definecolor{currentstroke}{rgb}{0.000000,0.000000,0.000000}%
\pgfsetstrokecolor{currentstroke}%
\pgfsetdash{}{0pt}%
\pgfsys@defobject{currentmarker}{\pgfqpoint{0.000000in}{0.000000in}}{\pgfqpoint{0.000000in}{0.020833in}}{%
\pgfpathmoveto{\pgfqpoint{0.000000in}{0.000000in}}%
\pgfpathlineto{\pgfqpoint{0.000000in}{0.020833in}}%
\pgfusepath{stroke,fill}%
}%
\begin{pgfscope}%
\pgfsys@transformshift{2.927190in}{0.420000in}%
\pgfsys@useobject{currentmarker}{}%
\end{pgfscope}%
\end{pgfscope}%
\begin{pgfscope}%
\pgfsetbuttcap%
\pgfsetroundjoin%
\definecolor{currentfill}{rgb}{0.000000,0.000000,0.000000}%
\pgfsetfillcolor{currentfill}%
\pgfsetlinewidth{0.401500pt}%
\definecolor{currentstroke}{rgb}{0.000000,0.000000,0.000000}%
\pgfsetstrokecolor{currentstroke}%
\pgfsetdash{}{0pt}%
\pgfsys@defobject{currentmarker}{\pgfqpoint{0.000000in}{-0.020833in}}{\pgfqpoint{0.000000in}{0.000000in}}{%
\pgfpathmoveto{\pgfqpoint{0.000000in}{0.000000in}}%
\pgfpathlineto{\pgfqpoint{0.000000in}{-0.020833in}}%
\pgfusepath{stroke,fill}%
}%
\begin{pgfscope}%
\pgfsys@transformshift{2.927190in}{2.414488in}%
\pgfsys@useobject{currentmarker}{}%
\end{pgfscope}%
\end{pgfscope}%
\begin{pgfscope}%
\pgfsetbuttcap%
\pgfsetroundjoin%
\definecolor{currentfill}{rgb}{0.000000,0.000000,0.000000}%
\pgfsetfillcolor{currentfill}%
\pgfsetlinewidth{0.401500pt}%
\definecolor{currentstroke}{rgb}{0.000000,0.000000,0.000000}%
\pgfsetstrokecolor{currentstroke}%
\pgfsetdash{}{0pt}%
\pgfsys@defobject{currentmarker}{\pgfqpoint{0.000000in}{0.000000in}}{\pgfqpoint{0.000000in}{0.020833in}}{%
\pgfpathmoveto{\pgfqpoint{0.000000in}{0.000000in}}%
\pgfpathlineto{\pgfqpoint{0.000000in}{0.020833in}}%
\pgfusepath{stroke,fill}%
}%
\begin{pgfscope}%
\pgfsys@transformshift{2.987293in}{0.420000in}%
\pgfsys@useobject{currentmarker}{}%
\end{pgfscope}%
\end{pgfscope}%
\begin{pgfscope}%
\pgfsetbuttcap%
\pgfsetroundjoin%
\definecolor{currentfill}{rgb}{0.000000,0.000000,0.000000}%
\pgfsetfillcolor{currentfill}%
\pgfsetlinewidth{0.401500pt}%
\definecolor{currentstroke}{rgb}{0.000000,0.000000,0.000000}%
\pgfsetstrokecolor{currentstroke}%
\pgfsetdash{}{0pt}%
\pgfsys@defobject{currentmarker}{\pgfqpoint{0.000000in}{-0.020833in}}{\pgfqpoint{0.000000in}{0.000000in}}{%
\pgfpathmoveto{\pgfqpoint{0.000000in}{0.000000in}}%
\pgfpathlineto{\pgfqpoint{0.000000in}{-0.020833in}}%
\pgfusepath{stroke,fill}%
}%
\begin{pgfscope}%
\pgfsys@transformshift{2.987293in}{2.414488in}%
\pgfsys@useobject{currentmarker}{}%
\end{pgfscope}%
\end{pgfscope}%
\begin{pgfscope}%
\pgfsetbuttcap%
\pgfsetroundjoin%
\definecolor{currentfill}{rgb}{0.000000,0.000000,0.000000}%
\pgfsetfillcolor{currentfill}%
\pgfsetlinewidth{0.401500pt}%
\definecolor{currentstroke}{rgb}{0.000000,0.000000,0.000000}%
\pgfsetstrokecolor{currentstroke}%
\pgfsetdash{}{0pt}%
\pgfsys@defobject{currentmarker}{\pgfqpoint{0.000000in}{0.000000in}}{\pgfqpoint{0.000000in}{0.020833in}}{%
\pgfpathmoveto{\pgfqpoint{0.000000in}{0.000000in}}%
\pgfpathlineto{\pgfqpoint{0.000000in}{0.020833in}}%
\pgfusepath{stroke,fill}%
}%
\begin{pgfscope}%
\pgfsys@transformshift{3.038110in}{0.420000in}%
\pgfsys@useobject{currentmarker}{}%
\end{pgfscope}%
\end{pgfscope}%
\begin{pgfscope}%
\pgfsetbuttcap%
\pgfsetroundjoin%
\definecolor{currentfill}{rgb}{0.000000,0.000000,0.000000}%
\pgfsetfillcolor{currentfill}%
\pgfsetlinewidth{0.401500pt}%
\definecolor{currentstroke}{rgb}{0.000000,0.000000,0.000000}%
\pgfsetstrokecolor{currentstroke}%
\pgfsetdash{}{0pt}%
\pgfsys@defobject{currentmarker}{\pgfqpoint{0.000000in}{-0.020833in}}{\pgfqpoint{0.000000in}{0.000000in}}{%
\pgfpathmoveto{\pgfqpoint{0.000000in}{0.000000in}}%
\pgfpathlineto{\pgfqpoint{0.000000in}{-0.020833in}}%
\pgfusepath{stroke,fill}%
}%
\begin{pgfscope}%
\pgfsys@transformshift{3.038110in}{2.414488in}%
\pgfsys@useobject{currentmarker}{}%
\end{pgfscope}%
\end{pgfscope}%
\begin{pgfscope}%
\definecolor{textcolor}{rgb}{0.000000,0.000000,0.000000}%
\pgfsetstrokecolor{textcolor}%
\pgfsetfillcolor{textcolor}%
\pgftext[x=1.844055in,y=0.208426in,,top]{\color{textcolor}{\rmfamily\fontsize{12.000000}{14.400000}\selectfont\catcode`\^=\active\def^{\ifmmode\sp\else\^{}\fi}\catcode`\%=\active\def%{\%}$y^+$}}%
\end{pgfscope}%
\begin{pgfscope}%
\pgfsetbuttcap%
\pgfsetroundjoin%
\definecolor{currentfill}{rgb}{0.000000,0.000000,0.000000}%
\pgfsetfillcolor{currentfill}%
\pgfsetlinewidth{0.501875pt}%
\definecolor{currentstroke}{rgb}{0.000000,0.000000,0.000000}%
\pgfsetstrokecolor{currentstroke}%
\pgfsetdash{}{0pt}%
\pgfsys@defobject{currentmarker}{\pgfqpoint{0.000000in}{0.000000in}}{\pgfqpoint{0.041667in}{0.000000in}}{%
\pgfpathmoveto{\pgfqpoint{0.000000in}{0.000000in}}%
\pgfpathlineto{\pgfqpoint{0.041667in}{0.000000in}}%
\pgfusepath{stroke,fill}%
}%
\begin{pgfscope}%
\pgfsys@transformshift{0.650000in}{0.420000in}%
\pgfsys@useobject{currentmarker}{}%
\end{pgfscope}%
\end{pgfscope}%
\begin{pgfscope}%
\pgfsetbuttcap%
\pgfsetroundjoin%
\definecolor{currentfill}{rgb}{0.000000,0.000000,0.000000}%
\pgfsetfillcolor{currentfill}%
\pgfsetlinewidth{0.501875pt}%
\definecolor{currentstroke}{rgb}{0.000000,0.000000,0.000000}%
\pgfsetstrokecolor{currentstroke}%
\pgfsetdash{}{0pt}%
\pgfsys@defobject{currentmarker}{\pgfqpoint{-0.041667in}{0.000000in}}{\pgfqpoint{-0.000000in}{0.000000in}}{%
\pgfpathmoveto{\pgfqpoint{-0.000000in}{0.000000in}}%
\pgfpathlineto{\pgfqpoint{-0.041667in}{0.000000in}}%
\pgfusepath{stroke,fill}%
}%
\begin{pgfscope}%
\pgfsys@transformshift{3.038110in}{0.420000in}%
\pgfsys@useobject{currentmarker}{}%
\end{pgfscope}%
\end{pgfscope}%
\begin{pgfscope}%
\definecolor{textcolor}{rgb}{0.000000,0.000000,0.000000}%
\pgfsetstrokecolor{textcolor}%
\pgfsetfillcolor{textcolor}%
\pgftext[x=0.525347in, y=0.367193in, left, base]{\color{textcolor}{\rmfamily\fontsize{11.000000}{13.200000}\selectfont\catcode`\^=\active\def^{\ifmmode\sp\else\^{}\fi}\catcode`\%=\active\def%{\%}0}}%
\end{pgfscope}%
\begin{pgfscope}%
\pgfsetbuttcap%
\pgfsetroundjoin%
\definecolor{currentfill}{rgb}{0.000000,0.000000,0.000000}%
\pgfsetfillcolor{currentfill}%
\pgfsetlinewidth{0.501875pt}%
\definecolor{currentstroke}{rgb}{0.000000,0.000000,0.000000}%
\pgfsetstrokecolor{currentstroke}%
\pgfsetdash{}{0pt}%
\pgfsys@defobject{currentmarker}{\pgfqpoint{0.000000in}{0.000000in}}{\pgfqpoint{0.041667in}{0.000000in}}{%
\pgfpathmoveto{\pgfqpoint{0.000000in}{0.000000in}}%
\pgfpathlineto{\pgfqpoint{0.041667in}{0.000000in}}%
\pgfusepath{stroke,fill}%
}%
\begin{pgfscope}%
\pgfsys@transformshift{0.650000in}{0.818898in}%
\pgfsys@useobject{currentmarker}{}%
\end{pgfscope}%
\end{pgfscope}%
\begin{pgfscope}%
\pgfsetbuttcap%
\pgfsetroundjoin%
\definecolor{currentfill}{rgb}{0.000000,0.000000,0.000000}%
\pgfsetfillcolor{currentfill}%
\pgfsetlinewidth{0.501875pt}%
\definecolor{currentstroke}{rgb}{0.000000,0.000000,0.000000}%
\pgfsetstrokecolor{currentstroke}%
\pgfsetdash{}{0pt}%
\pgfsys@defobject{currentmarker}{\pgfqpoint{-0.041667in}{0.000000in}}{\pgfqpoint{-0.000000in}{0.000000in}}{%
\pgfpathmoveto{\pgfqpoint{-0.000000in}{0.000000in}}%
\pgfpathlineto{\pgfqpoint{-0.041667in}{0.000000in}}%
\pgfusepath{stroke,fill}%
}%
\begin{pgfscope}%
\pgfsys@transformshift{3.038110in}{0.818898in}%
\pgfsys@useobject{currentmarker}{}%
\end{pgfscope}%
\end{pgfscope}%
\begin{pgfscope}%
\definecolor{textcolor}{rgb}{0.000000,0.000000,0.000000}%
\pgfsetstrokecolor{textcolor}%
\pgfsetfillcolor{textcolor}%
\pgftext[x=0.525347in, y=0.766091in, left, base]{\color{textcolor}{\rmfamily\fontsize{11.000000}{13.200000}\selectfont\catcode`\^=\active\def^{\ifmmode\sp\else\^{}\fi}\catcode`\%=\active\def%{\%}5}}%
\end{pgfscope}%
\begin{pgfscope}%
\pgfsetbuttcap%
\pgfsetroundjoin%
\definecolor{currentfill}{rgb}{0.000000,0.000000,0.000000}%
\pgfsetfillcolor{currentfill}%
\pgfsetlinewidth{0.501875pt}%
\definecolor{currentstroke}{rgb}{0.000000,0.000000,0.000000}%
\pgfsetstrokecolor{currentstroke}%
\pgfsetdash{}{0pt}%
\pgfsys@defobject{currentmarker}{\pgfqpoint{0.000000in}{0.000000in}}{\pgfqpoint{0.041667in}{0.000000in}}{%
\pgfpathmoveto{\pgfqpoint{0.000000in}{0.000000in}}%
\pgfpathlineto{\pgfqpoint{0.041667in}{0.000000in}}%
\pgfusepath{stroke,fill}%
}%
\begin{pgfscope}%
\pgfsys@transformshift{0.650000in}{1.217795in}%
\pgfsys@useobject{currentmarker}{}%
\end{pgfscope}%
\end{pgfscope}%
\begin{pgfscope}%
\pgfsetbuttcap%
\pgfsetroundjoin%
\definecolor{currentfill}{rgb}{0.000000,0.000000,0.000000}%
\pgfsetfillcolor{currentfill}%
\pgfsetlinewidth{0.501875pt}%
\definecolor{currentstroke}{rgb}{0.000000,0.000000,0.000000}%
\pgfsetstrokecolor{currentstroke}%
\pgfsetdash{}{0pt}%
\pgfsys@defobject{currentmarker}{\pgfqpoint{-0.041667in}{0.000000in}}{\pgfqpoint{-0.000000in}{0.000000in}}{%
\pgfpathmoveto{\pgfqpoint{-0.000000in}{0.000000in}}%
\pgfpathlineto{\pgfqpoint{-0.041667in}{0.000000in}}%
\pgfusepath{stroke,fill}%
}%
\begin{pgfscope}%
\pgfsys@transformshift{3.038110in}{1.217795in}%
\pgfsys@useobject{currentmarker}{}%
\end{pgfscope}%
\end{pgfscope}%
\begin{pgfscope}%
\definecolor{textcolor}{rgb}{0.000000,0.000000,0.000000}%
\pgfsetstrokecolor{textcolor}%
\pgfsetfillcolor{textcolor}%
\pgftext[x=0.449305in, y=1.164989in, left, base]{\color{textcolor}{\rmfamily\fontsize{11.000000}{13.200000}\selectfont\catcode`\^=\active\def^{\ifmmode\sp\else\^{}\fi}\catcode`\%=\active\def%{\%}10}}%
\end{pgfscope}%
\begin{pgfscope}%
\pgfsetbuttcap%
\pgfsetroundjoin%
\definecolor{currentfill}{rgb}{0.000000,0.000000,0.000000}%
\pgfsetfillcolor{currentfill}%
\pgfsetlinewidth{0.501875pt}%
\definecolor{currentstroke}{rgb}{0.000000,0.000000,0.000000}%
\pgfsetstrokecolor{currentstroke}%
\pgfsetdash{}{0pt}%
\pgfsys@defobject{currentmarker}{\pgfqpoint{0.000000in}{0.000000in}}{\pgfqpoint{0.041667in}{0.000000in}}{%
\pgfpathmoveto{\pgfqpoint{0.000000in}{0.000000in}}%
\pgfpathlineto{\pgfqpoint{0.041667in}{0.000000in}}%
\pgfusepath{stroke,fill}%
}%
\begin{pgfscope}%
\pgfsys@transformshift{0.650000in}{1.616693in}%
\pgfsys@useobject{currentmarker}{}%
\end{pgfscope}%
\end{pgfscope}%
\begin{pgfscope}%
\pgfsetbuttcap%
\pgfsetroundjoin%
\definecolor{currentfill}{rgb}{0.000000,0.000000,0.000000}%
\pgfsetfillcolor{currentfill}%
\pgfsetlinewidth{0.501875pt}%
\definecolor{currentstroke}{rgb}{0.000000,0.000000,0.000000}%
\pgfsetstrokecolor{currentstroke}%
\pgfsetdash{}{0pt}%
\pgfsys@defobject{currentmarker}{\pgfqpoint{-0.041667in}{0.000000in}}{\pgfqpoint{-0.000000in}{0.000000in}}{%
\pgfpathmoveto{\pgfqpoint{-0.000000in}{0.000000in}}%
\pgfpathlineto{\pgfqpoint{-0.041667in}{0.000000in}}%
\pgfusepath{stroke,fill}%
}%
\begin{pgfscope}%
\pgfsys@transformshift{3.038110in}{1.616693in}%
\pgfsys@useobject{currentmarker}{}%
\end{pgfscope}%
\end{pgfscope}%
\begin{pgfscope}%
\definecolor{textcolor}{rgb}{0.000000,0.000000,0.000000}%
\pgfsetstrokecolor{textcolor}%
\pgfsetfillcolor{textcolor}%
\pgftext[x=0.449305in, y=1.563886in, left, base]{\color{textcolor}{\rmfamily\fontsize{11.000000}{13.200000}\selectfont\catcode`\^=\active\def^{\ifmmode\sp\else\^{}\fi}\catcode`\%=\active\def%{\%}15}}%
\end{pgfscope}%
\begin{pgfscope}%
\pgfsetbuttcap%
\pgfsetroundjoin%
\definecolor{currentfill}{rgb}{0.000000,0.000000,0.000000}%
\pgfsetfillcolor{currentfill}%
\pgfsetlinewidth{0.501875pt}%
\definecolor{currentstroke}{rgb}{0.000000,0.000000,0.000000}%
\pgfsetstrokecolor{currentstroke}%
\pgfsetdash{}{0pt}%
\pgfsys@defobject{currentmarker}{\pgfqpoint{0.000000in}{0.000000in}}{\pgfqpoint{0.041667in}{0.000000in}}{%
\pgfpathmoveto{\pgfqpoint{0.000000in}{0.000000in}}%
\pgfpathlineto{\pgfqpoint{0.041667in}{0.000000in}}%
\pgfusepath{stroke,fill}%
}%
\begin{pgfscope}%
\pgfsys@transformshift{0.650000in}{2.015590in}%
\pgfsys@useobject{currentmarker}{}%
\end{pgfscope}%
\end{pgfscope}%
\begin{pgfscope}%
\pgfsetbuttcap%
\pgfsetroundjoin%
\definecolor{currentfill}{rgb}{0.000000,0.000000,0.000000}%
\pgfsetfillcolor{currentfill}%
\pgfsetlinewidth{0.501875pt}%
\definecolor{currentstroke}{rgb}{0.000000,0.000000,0.000000}%
\pgfsetstrokecolor{currentstroke}%
\pgfsetdash{}{0pt}%
\pgfsys@defobject{currentmarker}{\pgfqpoint{-0.041667in}{0.000000in}}{\pgfqpoint{-0.000000in}{0.000000in}}{%
\pgfpathmoveto{\pgfqpoint{-0.000000in}{0.000000in}}%
\pgfpathlineto{\pgfqpoint{-0.041667in}{0.000000in}}%
\pgfusepath{stroke,fill}%
}%
\begin{pgfscope}%
\pgfsys@transformshift{3.038110in}{2.015590in}%
\pgfsys@useobject{currentmarker}{}%
\end{pgfscope}%
\end{pgfscope}%
\begin{pgfscope}%
\definecolor{textcolor}{rgb}{0.000000,0.000000,0.000000}%
\pgfsetstrokecolor{textcolor}%
\pgfsetfillcolor{textcolor}%
\pgftext[x=0.449305in, y=1.962784in, left, base]{\color{textcolor}{\rmfamily\fontsize{11.000000}{13.200000}\selectfont\catcode`\^=\active\def^{\ifmmode\sp\else\^{}\fi}\catcode`\%=\active\def%{\%}20}}%
\end{pgfscope}%
\begin{pgfscope}%
\pgfsetbuttcap%
\pgfsetroundjoin%
\definecolor{currentfill}{rgb}{0.000000,0.000000,0.000000}%
\pgfsetfillcolor{currentfill}%
\pgfsetlinewidth{0.501875pt}%
\definecolor{currentstroke}{rgb}{0.000000,0.000000,0.000000}%
\pgfsetstrokecolor{currentstroke}%
\pgfsetdash{}{0pt}%
\pgfsys@defobject{currentmarker}{\pgfqpoint{0.000000in}{0.000000in}}{\pgfqpoint{0.041667in}{0.000000in}}{%
\pgfpathmoveto{\pgfqpoint{0.000000in}{0.000000in}}%
\pgfpathlineto{\pgfqpoint{0.041667in}{0.000000in}}%
\pgfusepath{stroke,fill}%
}%
\begin{pgfscope}%
\pgfsys@transformshift{0.650000in}{2.414488in}%
\pgfsys@useobject{currentmarker}{}%
\end{pgfscope}%
\end{pgfscope}%
\begin{pgfscope}%
\pgfsetbuttcap%
\pgfsetroundjoin%
\definecolor{currentfill}{rgb}{0.000000,0.000000,0.000000}%
\pgfsetfillcolor{currentfill}%
\pgfsetlinewidth{0.501875pt}%
\definecolor{currentstroke}{rgb}{0.000000,0.000000,0.000000}%
\pgfsetstrokecolor{currentstroke}%
\pgfsetdash{}{0pt}%
\pgfsys@defobject{currentmarker}{\pgfqpoint{-0.041667in}{0.000000in}}{\pgfqpoint{-0.000000in}{0.000000in}}{%
\pgfpathmoveto{\pgfqpoint{-0.000000in}{0.000000in}}%
\pgfpathlineto{\pgfqpoint{-0.041667in}{0.000000in}}%
\pgfusepath{stroke,fill}%
}%
\begin{pgfscope}%
\pgfsys@transformshift{3.038110in}{2.414488in}%
\pgfsys@useobject{currentmarker}{}%
\end{pgfscope}%
\end{pgfscope}%
\begin{pgfscope}%
\definecolor{textcolor}{rgb}{0.000000,0.000000,0.000000}%
\pgfsetstrokecolor{textcolor}%
\pgfsetfillcolor{textcolor}%
\pgftext[x=0.449305in, y=2.361681in, left, base]{\color{textcolor}{\rmfamily\fontsize{11.000000}{13.200000}\selectfont\catcode`\^=\active\def^{\ifmmode\sp\else\^{}\fi}\catcode`\%=\active\def%{\%}25}}%
\end{pgfscope}%
\begin{pgfscope}%
\pgfsetbuttcap%
\pgfsetroundjoin%
\definecolor{currentfill}{rgb}{0.000000,0.000000,0.000000}%
\pgfsetfillcolor{currentfill}%
\pgfsetlinewidth{0.401500pt}%
\definecolor{currentstroke}{rgb}{0.000000,0.000000,0.000000}%
\pgfsetstrokecolor{currentstroke}%
\pgfsetdash{}{0pt}%
\pgfsys@defobject{currentmarker}{\pgfqpoint{0.000000in}{0.000000in}}{\pgfqpoint{0.020833in}{0.000000in}}{%
\pgfpathmoveto{\pgfqpoint{0.000000in}{0.000000in}}%
\pgfpathlineto{\pgfqpoint{0.020833in}{0.000000in}}%
\pgfusepath{stroke,fill}%
}%
\begin{pgfscope}%
\pgfsys@transformshift{0.650000in}{0.499780in}%
\pgfsys@useobject{currentmarker}{}%
\end{pgfscope}%
\end{pgfscope}%
\begin{pgfscope}%
\pgfsetbuttcap%
\pgfsetroundjoin%
\definecolor{currentfill}{rgb}{0.000000,0.000000,0.000000}%
\pgfsetfillcolor{currentfill}%
\pgfsetlinewidth{0.401500pt}%
\definecolor{currentstroke}{rgb}{0.000000,0.000000,0.000000}%
\pgfsetstrokecolor{currentstroke}%
\pgfsetdash{}{0pt}%
\pgfsys@defobject{currentmarker}{\pgfqpoint{-0.020833in}{0.000000in}}{\pgfqpoint{-0.000000in}{0.000000in}}{%
\pgfpathmoveto{\pgfqpoint{-0.000000in}{0.000000in}}%
\pgfpathlineto{\pgfqpoint{-0.020833in}{0.000000in}}%
\pgfusepath{stroke,fill}%
}%
\begin{pgfscope}%
\pgfsys@transformshift{3.038110in}{0.499780in}%
\pgfsys@useobject{currentmarker}{}%
\end{pgfscope}%
\end{pgfscope}%
\begin{pgfscope}%
\pgfsetbuttcap%
\pgfsetroundjoin%
\definecolor{currentfill}{rgb}{0.000000,0.000000,0.000000}%
\pgfsetfillcolor{currentfill}%
\pgfsetlinewidth{0.401500pt}%
\definecolor{currentstroke}{rgb}{0.000000,0.000000,0.000000}%
\pgfsetstrokecolor{currentstroke}%
\pgfsetdash{}{0pt}%
\pgfsys@defobject{currentmarker}{\pgfqpoint{0.000000in}{0.000000in}}{\pgfqpoint{0.020833in}{0.000000in}}{%
\pgfpathmoveto{\pgfqpoint{0.000000in}{0.000000in}}%
\pgfpathlineto{\pgfqpoint{0.020833in}{0.000000in}}%
\pgfusepath{stroke,fill}%
}%
\begin{pgfscope}%
\pgfsys@transformshift{0.650000in}{0.579559in}%
\pgfsys@useobject{currentmarker}{}%
\end{pgfscope}%
\end{pgfscope}%
\begin{pgfscope}%
\pgfsetbuttcap%
\pgfsetroundjoin%
\definecolor{currentfill}{rgb}{0.000000,0.000000,0.000000}%
\pgfsetfillcolor{currentfill}%
\pgfsetlinewidth{0.401500pt}%
\definecolor{currentstroke}{rgb}{0.000000,0.000000,0.000000}%
\pgfsetstrokecolor{currentstroke}%
\pgfsetdash{}{0pt}%
\pgfsys@defobject{currentmarker}{\pgfqpoint{-0.020833in}{0.000000in}}{\pgfqpoint{-0.000000in}{0.000000in}}{%
\pgfpathmoveto{\pgfqpoint{-0.000000in}{0.000000in}}%
\pgfpathlineto{\pgfqpoint{-0.020833in}{0.000000in}}%
\pgfusepath{stroke,fill}%
}%
\begin{pgfscope}%
\pgfsys@transformshift{3.038110in}{0.579559in}%
\pgfsys@useobject{currentmarker}{}%
\end{pgfscope}%
\end{pgfscope}%
\begin{pgfscope}%
\pgfsetbuttcap%
\pgfsetroundjoin%
\definecolor{currentfill}{rgb}{0.000000,0.000000,0.000000}%
\pgfsetfillcolor{currentfill}%
\pgfsetlinewidth{0.401500pt}%
\definecolor{currentstroke}{rgb}{0.000000,0.000000,0.000000}%
\pgfsetstrokecolor{currentstroke}%
\pgfsetdash{}{0pt}%
\pgfsys@defobject{currentmarker}{\pgfqpoint{0.000000in}{0.000000in}}{\pgfqpoint{0.020833in}{0.000000in}}{%
\pgfpathmoveto{\pgfqpoint{0.000000in}{0.000000in}}%
\pgfpathlineto{\pgfqpoint{0.020833in}{0.000000in}}%
\pgfusepath{stroke,fill}%
}%
\begin{pgfscope}%
\pgfsys@transformshift{0.650000in}{0.659339in}%
\pgfsys@useobject{currentmarker}{}%
\end{pgfscope}%
\end{pgfscope}%
\begin{pgfscope}%
\pgfsetbuttcap%
\pgfsetroundjoin%
\definecolor{currentfill}{rgb}{0.000000,0.000000,0.000000}%
\pgfsetfillcolor{currentfill}%
\pgfsetlinewidth{0.401500pt}%
\definecolor{currentstroke}{rgb}{0.000000,0.000000,0.000000}%
\pgfsetstrokecolor{currentstroke}%
\pgfsetdash{}{0pt}%
\pgfsys@defobject{currentmarker}{\pgfqpoint{-0.020833in}{0.000000in}}{\pgfqpoint{-0.000000in}{0.000000in}}{%
\pgfpathmoveto{\pgfqpoint{-0.000000in}{0.000000in}}%
\pgfpathlineto{\pgfqpoint{-0.020833in}{0.000000in}}%
\pgfusepath{stroke,fill}%
}%
\begin{pgfscope}%
\pgfsys@transformshift{3.038110in}{0.659339in}%
\pgfsys@useobject{currentmarker}{}%
\end{pgfscope}%
\end{pgfscope}%
\begin{pgfscope}%
\pgfsetbuttcap%
\pgfsetroundjoin%
\definecolor{currentfill}{rgb}{0.000000,0.000000,0.000000}%
\pgfsetfillcolor{currentfill}%
\pgfsetlinewidth{0.401500pt}%
\definecolor{currentstroke}{rgb}{0.000000,0.000000,0.000000}%
\pgfsetstrokecolor{currentstroke}%
\pgfsetdash{}{0pt}%
\pgfsys@defobject{currentmarker}{\pgfqpoint{0.000000in}{0.000000in}}{\pgfqpoint{0.020833in}{0.000000in}}{%
\pgfpathmoveto{\pgfqpoint{0.000000in}{0.000000in}}%
\pgfpathlineto{\pgfqpoint{0.020833in}{0.000000in}}%
\pgfusepath{stroke,fill}%
}%
\begin{pgfscope}%
\pgfsys@transformshift{0.650000in}{0.739118in}%
\pgfsys@useobject{currentmarker}{}%
\end{pgfscope}%
\end{pgfscope}%
\begin{pgfscope}%
\pgfsetbuttcap%
\pgfsetroundjoin%
\definecolor{currentfill}{rgb}{0.000000,0.000000,0.000000}%
\pgfsetfillcolor{currentfill}%
\pgfsetlinewidth{0.401500pt}%
\definecolor{currentstroke}{rgb}{0.000000,0.000000,0.000000}%
\pgfsetstrokecolor{currentstroke}%
\pgfsetdash{}{0pt}%
\pgfsys@defobject{currentmarker}{\pgfqpoint{-0.020833in}{0.000000in}}{\pgfqpoint{-0.000000in}{0.000000in}}{%
\pgfpathmoveto{\pgfqpoint{-0.000000in}{0.000000in}}%
\pgfpathlineto{\pgfqpoint{-0.020833in}{0.000000in}}%
\pgfusepath{stroke,fill}%
}%
\begin{pgfscope}%
\pgfsys@transformshift{3.038110in}{0.739118in}%
\pgfsys@useobject{currentmarker}{}%
\end{pgfscope}%
\end{pgfscope}%
\begin{pgfscope}%
\pgfsetbuttcap%
\pgfsetroundjoin%
\definecolor{currentfill}{rgb}{0.000000,0.000000,0.000000}%
\pgfsetfillcolor{currentfill}%
\pgfsetlinewidth{0.401500pt}%
\definecolor{currentstroke}{rgb}{0.000000,0.000000,0.000000}%
\pgfsetstrokecolor{currentstroke}%
\pgfsetdash{}{0pt}%
\pgfsys@defobject{currentmarker}{\pgfqpoint{0.000000in}{0.000000in}}{\pgfqpoint{0.020833in}{0.000000in}}{%
\pgfpathmoveto{\pgfqpoint{0.000000in}{0.000000in}}%
\pgfpathlineto{\pgfqpoint{0.020833in}{0.000000in}}%
\pgfusepath{stroke,fill}%
}%
\begin{pgfscope}%
\pgfsys@transformshift{0.650000in}{0.898677in}%
\pgfsys@useobject{currentmarker}{}%
\end{pgfscope}%
\end{pgfscope}%
\begin{pgfscope}%
\pgfsetbuttcap%
\pgfsetroundjoin%
\definecolor{currentfill}{rgb}{0.000000,0.000000,0.000000}%
\pgfsetfillcolor{currentfill}%
\pgfsetlinewidth{0.401500pt}%
\definecolor{currentstroke}{rgb}{0.000000,0.000000,0.000000}%
\pgfsetstrokecolor{currentstroke}%
\pgfsetdash{}{0pt}%
\pgfsys@defobject{currentmarker}{\pgfqpoint{-0.020833in}{0.000000in}}{\pgfqpoint{-0.000000in}{0.000000in}}{%
\pgfpathmoveto{\pgfqpoint{-0.000000in}{0.000000in}}%
\pgfpathlineto{\pgfqpoint{-0.020833in}{0.000000in}}%
\pgfusepath{stroke,fill}%
}%
\begin{pgfscope}%
\pgfsys@transformshift{3.038110in}{0.898677in}%
\pgfsys@useobject{currentmarker}{}%
\end{pgfscope}%
\end{pgfscope}%
\begin{pgfscope}%
\pgfsetbuttcap%
\pgfsetroundjoin%
\definecolor{currentfill}{rgb}{0.000000,0.000000,0.000000}%
\pgfsetfillcolor{currentfill}%
\pgfsetlinewidth{0.401500pt}%
\definecolor{currentstroke}{rgb}{0.000000,0.000000,0.000000}%
\pgfsetstrokecolor{currentstroke}%
\pgfsetdash{}{0pt}%
\pgfsys@defobject{currentmarker}{\pgfqpoint{0.000000in}{0.000000in}}{\pgfqpoint{0.020833in}{0.000000in}}{%
\pgfpathmoveto{\pgfqpoint{0.000000in}{0.000000in}}%
\pgfpathlineto{\pgfqpoint{0.020833in}{0.000000in}}%
\pgfusepath{stroke,fill}%
}%
\begin{pgfscope}%
\pgfsys@transformshift{0.650000in}{0.978457in}%
\pgfsys@useobject{currentmarker}{}%
\end{pgfscope}%
\end{pgfscope}%
\begin{pgfscope}%
\pgfsetbuttcap%
\pgfsetroundjoin%
\definecolor{currentfill}{rgb}{0.000000,0.000000,0.000000}%
\pgfsetfillcolor{currentfill}%
\pgfsetlinewidth{0.401500pt}%
\definecolor{currentstroke}{rgb}{0.000000,0.000000,0.000000}%
\pgfsetstrokecolor{currentstroke}%
\pgfsetdash{}{0pt}%
\pgfsys@defobject{currentmarker}{\pgfqpoint{-0.020833in}{0.000000in}}{\pgfqpoint{-0.000000in}{0.000000in}}{%
\pgfpathmoveto{\pgfqpoint{-0.000000in}{0.000000in}}%
\pgfpathlineto{\pgfqpoint{-0.020833in}{0.000000in}}%
\pgfusepath{stroke,fill}%
}%
\begin{pgfscope}%
\pgfsys@transformshift{3.038110in}{0.978457in}%
\pgfsys@useobject{currentmarker}{}%
\end{pgfscope}%
\end{pgfscope}%
\begin{pgfscope}%
\pgfsetbuttcap%
\pgfsetroundjoin%
\definecolor{currentfill}{rgb}{0.000000,0.000000,0.000000}%
\pgfsetfillcolor{currentfill}%
\pgfsetlinewidth{0.401500pt}%
\definecolor{currentstroke}{rgb}{0.000000,0.000000,0.000000}%
\pgfsetstrokecolor{currentstroke}%
\pgfsetdash{}{0pt}%
\pgfsys@defobject{currentmarker}{\pgfqpoint{0.000000in}{0.000000in}}{\pgfqpoint{0.020833in}{0.000000in}}{%
\pgfpathmoveto{\pgfqpoint{0.000000in}{0.000000in}}%
\pgfpathlineto{\pgfqpoint{0.020833in}{0.000000in}}%
\pgfusepath{stroke,fill}%
}%
\begin{pgfscope}%
\pgfsys@transformshift{0.650000in}{1.058236in}%
\pgfsys@useobject{currentmarker}{}%
\end{pgfscope}%
\end{pgfscope}%
\begin{pgfscope}%
\pgfsetbuttcap%
\pgfsetroundjoin%
\definecolor{currentfill}{rgb}{0.000000,0.000000,0.000000}%
\pgfsetfillcolor{currentfill}%
\pgfsetlinewidth{0.401500pt}%
\definecolor{currentstroke}{rgb}{0.000000,0.000000,0.000000}%
\pgfsetstrokecolor{currentstroke}%
\pgfsetdash{}{0pt}%
\pgfsys@defobject{currentmarker}{\pgfqpoint{-0.020833in}{0.000000in}}{\pgfqpoint{-0.000000in}{0.000000in}}{%
\pgfpathmoveto{\pgfqpoint{-0.000000in}{0.000000in}}%
\pgfpathlineto{\pgfqpoint{-0.020833in}{0.000000in}}%
\pgfusepath{stroke,fill}%
}%
\begin{pgfscope}%
\pgfsys@transformshift{3.038110in}{1.058236in}%
\pgfsys@useobject{currentmarker}{}%
\end{pgfscope}%
\end{pgfscope}%
\begin{pgfscope}%
\pgfsetbuttcap%
\pgfsetroundjoin%
\definecolor{currentfill}{rgb}{0.000000,0.000000,0.000000}%
\pgfsetfillcolor{currentfill}%
\pgfsetlinewidth{0.401500pt}%
\definecolor{currentstroke}{rgb}{0.000000,0.000000,0.000000}%
\pgfsetstrokecolor{currentstroke}%
\pgfsetdash{}{0pt}%
\pgfsys@defobject{currentmarker}{\pgfqpoint{0.000000in}{0.000000in}}{\pgfqpoint{0.020833in}{0.000000in}}{%
\pgfpathmoveto{\pgfqpoint{0.000000in}{0.000000in}}%
\pgfpathlineto{\pgfqpoint{0.020833in}{0.000000in}}%
\pgfusepath{stroke,fill}%
}%
\begin{pgfscope}%
\pgfsys@transformshift{0.650000in}{1.138016in}%
\pgfsys@useobject{currentmarker}{}%
\end{pgfscope}%
\end{pgfscope}%
\begin{pgfscope}%
\pgfsetbuttcap%
\pgfsetroundjoin%
\definecolor{currentfill}{rgb}{0.000000,0.000000,0.000000}%
\pgfsetfillcolor{currentfill}%
\pgfsetlinewidth{0.401500pt}%
\definecolor{currentstroke}{rgb}{0.000000,0.000000,0.000000}%
\pgfsetstrokecolor{currentstroke}%
\pgfsetdash{}{0pt}%
\pgfsys@defobject{currentmarker}{\pgfqpoint{-0.020833in}{0.000000in}}{\pgfqpoint{-0.000000in}{0.000000in}}{%
\pgfpathmoveto{\pgfqpoint{-0.000000in}{0.000000in}}%
\pgfpathlineto{\pgfqpoint{-0.020833in}{0.000000in}}%
\pgfusepath{stroke,fill}%
}%
\begin{pgfscope}%
\pgfsys@transformshift{3.038110in}{1.138016in}%
\pgfsys@useobject{currentmarker}{}%
\end{pgfscope}%
\end{pgfscope}%
\begin{pgfscope}%
\pgfsetbuttcap%
\pgfsetroundjoin%
\definecolor{currentfill}{rgb}{0.000000,0.000000,0.000000}%
\pgfsetfillcolor{currentfill}%
\pgfsetlinewidth{0.401500pt}%
\definecolor{currentstroke}{rgb}{0.000000,0.000000,0.000000}%
\pgfsetstrokecolor{currentstroke}%
\pgfsetdash{}{0pt}%
\pgfsys@defobject{currentmarker}{\pgfqpoint{0.000000in}{0.000000in}}{\pgfqpoint{0.020833in}{0.000000in}}{%
\pgfpathmoveto{\pgfqpoint{0.000000in}{0.000000in}}%
\pgfpathlineto{\pgfqpoint{0.020833in}{0.000000in}}%
\pgfusepath{stroke,fill}%
}%
\begin{pgfscope}%
\pgfsys@transformshift{0.650000in}{1.297575in}%
\pgfsys@useobject{currentmarker}{}%
\end{pgfscope}%
\end{pgfscope}%
\begin{pgfscope}%
\pgfsetbuttcap%
\pgfsetroundjoin%
\definecolor{currentfill}{rgb}{0.000000,0.000000,0.000000}%
\pgfsetfillcolor{currentfill}%
\pgfsetlinewidth{0.401500pt}%
\definecolor{currentstroke}{rgb}{0.000000,0.000000,0.000000}%
\pgfsetstrokecolor{currentstroke}%
\pgfsetdash{}{0pt}%
\pgfsys@defobject{currentmarker}{\pgfqpoint{-0.020833in}{0.000000in}}{\pgfqpoint{-0.000000in}{0.000000in}}{%
\pgfpathmoveto{\pgfqpoint{-0.000000in}{0.000000in}}%
\pgfpathlineto{\pgfqpoint{-0.020833in}{0.000000in}}%
\pgfusepath{stroke,fill}%
}%
\begin{pgfscope}%
\pgfsys@transformshift{3.038110in}{1.297575in}%
\pgfsys@useobject{currentmarker}{}%
\end{pgfscope}%
\end{pgfscope}%
\begin{pgfscope}%
\pgfsetbuttcap%
\pgfsetroundjoin%
\definecolor{currentfill}{rgb}{0.000000,0.000000,0.000000}%
\pgfsetfillcolor{currentfill}%
\pgfsetlinewidth{0.401500pt}%
\definecolor{currentstroke}{rgb}{0.000000,0.000000,0.000000}%
\pgfsetstrokecolor{currentstroke}%
\pgfsetdash{}{0pt}%
\pgfsys@defobject{currentmarker}{\pgfqpoint{0.000000in}{0.000000in}}{\pgfqpoint{0.020833in}{0.000000in}}{%
\pgfpathmoveto{\pgfqpoint{0.000000in}{0.000000in}}%
\pgfpathlineto{\pgfqpoint{0.020833in}{0.000000in}}%
\pgfusepath{stroke,fill}%
}%
\begin{pgfscope}%
\pgfsys@transformshift{0.650000in}{1.377354in}%
\pgfsys@useobject{currentmarker}{}%
\end{pgfscope}%
\end{pgfscope}%
\begin{pgfscope}%
\pgfsetbuttcap%
\pgfsetroundjoin%
\definecolor{currentfill}{rgb}{0.000000,0.000000,0.000000}%
\pgfsetfillcolor{currentfill}%
\pgfsetlinewidth{0.401500pt}%
\definecolor{currentstroke}{rgb}{0.000000,0.000000,0.000000}%
\pgfsetstrokecolor{currentstroke}%
\pgfsetdash{}{0pt}%
\pgfsys@defobject{currentmarker}{\pgfqpoint{-0.020833in}{0.000000in}}{\pgfqpoint{-0.000000in}{0.000000in}}{%
\pgfpathmoveto{\pgfqpoint{-0.000000in}{0.000000in}}%
\pgfpathlineto{\pgfqpoint{-0.020833in}{0.000000in}}%
\pgfusepath{stroke,fill}%
}%
\begin{pgfscope}%
\pgfsys@transformshift{3.038110in}{1.377354in}%
\pgfsys@useobject{currentmarker}{}%
\end{pgfscope}%
\end{pgfscope}%
\begin{pgfscope}%
\pgfsetbuttcap%
\pgfsetroundjoin%
\definecolor{currentfill}{rgb}{0.000000,0.000000,0.000000}%
\pgfsetfillcolor{currentfill}%
\pgfsetlinewidth{0.401500pt}%
\definecolor{currentstroke}{rgb}{0.000000,0.000000,0.000000}%
\pgfsetstrokecolor{currentstroke}%
\pgfsetdash{}{0pt}%
\pgfsys@defobject{currentmarker}{\pgfqpoint{0.000000in}{0.000000in}}{\pgfqpoint{0.020833in}{0.000000in}}{%
\pgfpathmoveto{\pgfqpoint{0.000000in}{0.000000in}}%
\pgfpathlineto{\pgfqpoint{0.020833in}{0.000000in}}%
\pgfusepath{stroke,fill}%
}%
\begin{pgfscope}%
\pgfsys@transformshift{0.650000in}{1.457134in}%
\pgfsys@useobject{currentmarker}{}%
\end{pgfscope}%
\end{pgfscope}%
\begin{pgfscope}%
\pgfsetbuttcap%
\pgfsetroundjoin%
\definecolor{currentfill}{rgb}{0.000000,0.000000,0.000000}%
\pgfsetfillcolor{currentfill}%
\pgfsetlinewidth{0.401500pt}%
\definecolor{currentstroke}{rgb}{0.000000,0.000000,0.000000}%
\pgfsetstrokecolor{currentstroke}%
\pgfsetdash{}{0pt}%
\pgfsys@defobject{currentmarker}{\pgfqpoint{-0.020833in}{0.000000in}}{\pgfqpoint{-0.000000in}{0.000000in}}{%
\pgfpathmoveto{\pgfqpoint{-0.000000in}{0.000000in}}%
\pgfpathlineto{\pgfqpoint{-0.020833in}{0.000000in}}%
\pgfusepath{stroke,fill}%
}%
\begin{pgfscope}%
\pgfsys@transformshift{3.038110in}{1.457134in}%
\pgfsys@useobject{currentmarker}{}%
\end{pgfscope}%
\end{pgfscope}%
\begin{pgfscope}%
\pgfsetbuttcap%
\pgfsetroundjoin%
\definecolor{currentfill}{rgb}{0.000000,0.000000,0.000000}%
\pgfsetfillcolor{currentfill}%
\pgfsetlinewidth{0.401500pt}%
\definecolor{currentstroke}{rgb}{0.000000,0.000000,0.000000}%
\pgfsetstrokecolor{currentstroke}%
\pgfsetdash{}{0pt}%
\pgfsys@defobject{currentmarker}{\pgfqpoint{0.000000in}{0.000000in}}{\pgfqpoint{0.020833in}{0.000000in}}{%
\pgfpathmoveto{\pgfqpoint{0.000000in}{0.000000in}}%
\pgfpathlineto{\pgfqpoint{0.020833in}{0.000000in}}%
\pgfusepath{stroke,fill}%
}%
\begin{pgfscope}%
\pgfsys@transformshift{0.650000in}{1.536913in}%
\pgfsys@useobject{currentmarker}{}%
\end{pgfscope}%
\end{pgfscope}%
\begin{pgfscope}%
\pgfsetbuttcap%
\pgfsetroundjoin%
\definecolor{currentfill}{rgb}{0.000000,0.000000,0.000000}%
\pgfsetfillcolor{currentfill}%
\pgfsetlinewidth{0.401500pt}%
\definecolor{currentstroke}{rgb}{0.000000,0.000000,0.000000}%
\pgfsetstrokecolor{currentstroke}%
\pgfsetdash{}{0pt}%
\pgfsys@defobject{currentmarker}{\pgfqpoint{-0.020833in}{0.000000in}}{\pgfqpoint{-0.000000in}{0.000000in}}{%
\pgfpathmoveto{\pgfqpoint{-0.000000in}{0.000000in}}%
\pgfpathlineto{\pgfqpoint{-0.020833in}{0.000000in}}%
\pgfusepath{stroke,fill}%
}%
\begin{pgfscope}%
\pgfsys@transformshift{3.038110in}{1.536913in}%
\pgfsys@useobject{currentmarker}{}%
\end{pgfscope}%
\end{pgfscope}%
\begin{pgfscope}%
\pgfsetbuttcap%
\pgfsetroundjoin%
\definecolor{currentfill}{rgb}{0.000000,0.000000,0.000000}%
\pgfsetfillcolor{currentfill}%
\pgfsetlinewidth{0.401500pt}%
\definecolor{currentstroke}{rgb}{0.000000,0.000000,0.000000}%
\pgfsetstrokecolor{currentstroke}%
\pgfsetdash{}{0pt}%
\pgfsys@defobject{currentmarker}{\pgfqpoint{0.000000in}{0.000000in}}{\pgfqpoint{0.020833in}{0.000000in}}{%
\pgfpathmoveto{\pgfqpoint{0.000000in}{0.000000in}}%
\pgfpathlineto{\pgfqpoint{0.020833in}{0.000000in}}%
\pgfusepath{stroke,fill}%
}%
\begin{pgfscope}%
\pgfsys@transformshift{0.650000in}{1.696472in}%
\pgfsys@useobject{currentmarker}{}%
\end{pgfscope}%
\end{pgfscope}%
\begin{pgfscope}%
\pgfsetbuttcap%
\pgfsetroundjoin%
\definecolor{currentfill}{rgb}{0.000000,0.000000,0.000000}%
\pgfsetfillcolor{currentfill}%
\pgfsetlinewidth{0.401500pt}%
\definecolor{currentstroke}{rgb}{0.000000,0.000000,0.000000}%
\pgfsetstrokecolor{currentstroke}%
\pgfsetdash{}{0pt}%
\pgfsys@defobject{currentmarker}{\pgfqpoint{-0.020833in}{0.000000in}}{\pgfqpoint{-0.000000in}{0.000000in}}{%
\pgfpathmoveto{\pgfqpoint{-0.000000in}{0.000000in}}%
\pgfpathlineto{\pgfqpoint{-0.020833in}{0.000000in}}%
\pgfusepath{stroke,fill}%
}%
\begin{pgfscope}%
\pgfsys@transformshift{3.038110in}{1.696472in}%
\pgfsys@useobject{currentmarker}{}%
\end{pgfscope}%
\end{pgfscope}%
\begin{pgfscope}%
\pgfsetbuttcap%
\pgfsetroundjoin%
\definecolor{currentfill}{rgb}{0.000000,0.000000,0.000000}%
\pgfsetfillcolor{currentfill}%
\pgfsetlinewidth{0.401500pt}%
\definecolor{currentstroke}{rgb}{0.000000,0.000000,0.000000}%
\pgfsetstrokecolor{currentstroke}%
\pgfsetdash{}{0pt}%
\pgfsys@defobject{currentmarker}{\pgfqpoint{0.000000in}{0.000000in}}{\pgfqpoint{0.020833in}{0.000000in}}{%
\pgfpathmoveto{\pgfqpoint{0.000000in}{0.000000in}}%
\pgfpathlineto{\pgfqpoint{0.020833in}{0.000000in}}%
\pgfusepath{stroke,fill}%
}%
\begin{pgfscope}%
\pgfsys@transformshift{0.650000in}{1.776252in}%
\pgfsys@useobject{currentmarker}{}%
\end{pgfscope}%
\end{pgfscope}%
\begin{pgfscope}%
\pgfsetbuttcap%
\pgfsetroundjoin%
\definecolor{currentfill}{rgb}{0.000000,0.000000,0.000000}%
\pgfsetfillcolor{currentfill}%
\pgfsetlinewidth{0.401500pt}%
\definecolor{currentstroke}{rgb}{0.000000,0.000000,0.000000}%
\pgfsetstrokecolor{currentstroke}%
\pgfsetdash{}{0pt}%
\pgfsys@defobject{currentmarker}{\pgfqpoint{-0.020833in}{0.000000in}}{\pgfqpoint{-0.000000in}{0.000000in}}{%
\pgfpathmoveto{\pgfqpoint{-0.000000in}{0.000000in}}%
\pgfpathlineto{\pgfqpoint{-0.020833in}{0.000000in}}%
\pgfusepath{stroke,fill}%
}%
\begin{pgfscope}%
\pgfsys@transformshift{3.038110in}{1.776252in}%
\pgfsys@useobject{currentmarker}{}%
\end{pgfscope}%
\end{pgfscope}%
\begin{pgfscope}%
\pgfsetbuttcap%
\pgfsetroundjoin%
\definecolor{currentfill}{rgb}{0.000000,0.000000,0.000000}%
\pgfsetfillcolor{currentfill}%
\pgfsetlinewidth{0.401500pt}%
\definecolor{currentstroke}{rgb}{0.000000,0.000000,0.000000}%
\pgfsetstrokecolor{currentstroke}%
\pgfsetdash{}{0pt}%
\pgfsys@defobject{currentmarker}{\pgfqpoint{0.000000in}{0.000000in}}{\pgfqpoint{0.020833in}{0.000000in}}{%
\pgfpathmoveto{\pgfqpoint{0.000000in}{0.000000in}}%
\pgfpathlineto{\pgfqpoint{0.020833in}{0.000000in}}%
\pgfusepath{stroke,fill}%
}%
\begin{pgfscope}%
\pgfsys@transformshift{0.650000in}{1.856031in}%
\pgfsys@useobject{currentmarker}{}%
\end{pgfscope}%
\end{pgfscope}%
\begin{pgfscope}%
\pgfsetbuttcap%
\pgfsetroundjoin%
\definecolor{currentfill}{rgb}{0.000000,0.000000,0.000000}%
\pgfsetfillcolor{currentfill}%
\pgfsetlinewidth{0.401500pt}%
\definecolor{currentstroke}{rgb}{0.000000,0.000000,0.000000}%
\pgfsetstrokecolor{currentstroke}%
\pgfsetdash{}{0pt}%
\pgfsys@defobject{currentmarker}{\pgfqpoint{-0.020833in}{0.000000in}}{\pgfqpoint{-0.000000in}{0.000000in}}{%
\pgfpathmoveto{\pgfqpoint{-0.000000in}{0.000000in}}%
\pgfpathlineto{\pgfqpoint{-0.020833in}{0.000000in}}%
\pgfusepath{stroke,fill}%
}%
\begin{pgfscope}%
\pgfsys@transformshift{3.038110in}{1.856031in}%
\pgfsys@useobject{currentmarker}{}%
\end{pgfscope}%
\end{pgfscope}%
\begin{pgfscope}%
\pgfsetbuttcap%
\pgfsetroundjoin%
\definecolor{currentfill}{rgb}{0.000000,0.000000,0.000000}%
\pgfsetfillcolor{currentfill}%
\pgfsetlinewidth{0.401500pt}%
\definecolor{currentstroke}{rgb}{0.000000,0.000000,0.000000}%
\pgfsetstrokecolor{currentstroke}%
\pgfsetdash{}{0pt}%
\pgfsys@defobject{currentmarker}{\pgfqpoint{0.000000in}{0.000000in}}{\pgfqpoint{0.020833in}{0.000000in}}{%
\pgfpathmoveto{\pgfqpoint{0.000000in}{0.000000in}}%
\pgfpathlineto{\pgfqpoint{0.020833in}{0.000000in}}%
\pgfusepath{stroke,fill}%
}%
\begin{pgfscope}%
\pgfsys@transformshift{0.650000in}{1.935811in}%
\pgfsys@useobject{currentmarker}{}%
\end{pgfscope}%
\end{pgfscope}%
\begin{pgfscope}%
\pgfsetbuttcap%
\pgfsetroundjoin%
\definecolor{currentfill}{rgb}{0.000000,0.000000,0.000000}%
\pgfsetfillcolor{currentfill}%
\pgfsetlinewidth{0.401500pt}%
\definecolor{currentstroke}{rgb}{0.000000,0.000000,0.000000}%
\pgfsetstrokecolor{currentstroke}%
\pgfsetdash{}{0pt}%
\pgfsys@defobject{currentmarker}{\pgfqpoint{-0.020833in}{0.000000in}}{\pgfqpoint{-0.000000in}{0.000000in}}{%
\pgfpathmoveto{\pgfqpoint{-0.000000in}{0.000000in}}%
\pgfpathlineto{\pgfqpoint{-0.020833in}{0.000000in}}%
\pgfusepath{stroke,fill}%
}%
\begin{pgfscope}%
\pgfsys@transformshift{3.038110in}{1.935811in}%
\pgfsys@useobject{currentmarker}{}%
\end{pgfscope}%
\end{pgfscope}%
\begin{pgfscope}%
\pgfsetbuttcap%
\pgfsetroundjoin%
\definecolor{currentfill}{rgb}{0.000000,0.000000,0.000000}%
\pgfsetfillcolor{currentfill}%
\pgfsetlinewidth{0.401500pt}%
\definecolor{currentstroke}{rgb}{0.000000,0.000000,0.000000}%
\pgfsetstrokecolor{currentstroke}%
\pgfsetdash{}{0pt}%
\pgfsys@defobject{currentmarker}{\pgfqpoint{0.000000in}{0.000000in}}{\pgfqpoint{0.020833in}{0.000000in}}{%
\pgfpathmoveto{\pgfqpoint{0.000000in}{0.000000in}}%
\pgfpathlineto{\pgfqpoint{0.020833in}{0.000000in}}%
\pgfusepath{stroke,fill}%
}%
\begin{pgfscope}%
\pgfsys@transformshift{0.650000in}{2.095370in}%
\pgfsys@useobject{currentmarker}{}%
\end{pgfscope}%
\end{pgfscope}%
\begin{pgfscope}%
\pgfsetbuttcap%
\pgfsetroundjoin%
\definecolor{currentfill}{rgb}{0.000000,0.000000,0.000000}%
\pgfsetfillcolor{currentfill}%
\pgfsetlinewidth{0.401500pt}%
\definecolor{currentstroke}{rgb}{0.000000,0.000000,0.000000}%
\pgfsetstrokecolor{currentstroke}%
\pgfsetdash{}{0pt}%
\pgfsys@defobject{currentmarker}{\pgfqpoint{-0.020833in}{0.000000in}}{\pgfqpoint{-0.000000in}{0.000000in}}{%
\pgfpathmoveto{\pgfqpoint{-0.000000in}{0.000000in}}%
\pgfpathlineto{\pgfqpoint{-0.020833in}{0.000000in}}%
\pgfusepath{stroke,fill}%
}%
\begin{pgfscope}%
\pgfsys@transformshift{3.038110in}{2.095370in}%
\pgfsys@useobject{currentmarker}{}%
\end{pgfscope}%
\end{pgfscope}%
\begin{pgfscope}%
\pgfsetbuttcap%
\pgfsetroundjoin%
\definecolor{currentfill}{rgb}{0.000000,0.000000,0.000000}%
\pgfsetfillcolor{currentfill}%
\pgfsetlinewidth{0.401500pt}%
\definecolor{currentstroke}{rgb}{0.000000,0.000000,0.000000}%
\pgfsetstrokecolor{currentstroke}%
\pgfsetdash{}{0pt}%
\pgfsys@defobject{currentmarker}{\pgfqpoint{0.000000in}{0.000000in}}{\pgfqpoint{0.020833in}{0.000000in}}{%
\pgfpathmoveto{\pgfqpoint{0.000000in}{0.000000in}}%
\pgfpathlineto{\pgfqpoint{0.020833in}{0.000000in}}%
\pgfusepath{stroke,fill}%
}%
\begin{pgfscope}%
\pgfsys@transformshift{0.650000in}{2.175150in}%
\pgfsys@useobject{currentmarker}{}%
\end{pgfscope}%
\end{pgfscope}%
\begin{pgfscope}%
\pgfsetbuttcap%
\pgfsetroundjoin%
\definecolor{currentfill}{rgb}{0.000000,0.000000,0.000000}%
\pgfsetfillcolor{currentfill}%
\pgfsetlinewidth{0.401500pt}%
\definecolor{currentstroke}{rgb}{0.000000,0.000000,0.000000}%
\pgfsetstrokecolor{currentstroke}%
\pgfsetdash{}{0pt}%
\pgfsys@defobject{currentmarker}{\pgfqpoint{-0.020833in}{0.000000in}}{\pgfqpoint{-0.000000in}{0.000000in}}{%
\pgfpathmoveto{\pgfqpoint{-0.000000in}{0.000000in}}%
\pgfpathlineto{\pgfqpoint{-0.020833in}{0.000000in}}%
\pgfusepath{stroke,fill}%
}%
\begin{pgfscope}%
\pgfsys@transformshift{3.038110in}{2.175150in}%
\pgfsys@useobject{currentmarker}{}%
\end{pgfscope}%
\end{pgfscope}%
\begin{pgfscope}%
\pgfsetbuttcap%
\pgfsetroundjoin%
\definecolor{currentfill}{rgb}{0.000000,0.000000,0.000000}%
\pgfsetfillcolor{currentfill}%
\pgfsetlinewidth{0.401500pt}%
\definecolor{currentstroke}{rgb}{0.000000,0.000000,0.000000}%
\pgfsetstrokecolor{currentstroke}%
\pgfsetdash{}{0pt}%
\pgfsys@defobject{currentmarker}{\pgfqpoint{0.000000in}{0.000000in}}{\pgfqpoint{0.020833in}{0.000000in}}{%
\pgfpathmoveto{\pgfqpoint{0.000000in}{0.000000in}}%
\pgfpathlineto{\pgfqpoint{0.020833in}{0.000000in}}%
\pgfusepath{stroke,fill}%
}%
\begin{pgfscope}%
\pgfsys@transformshift{0.650000in}{2.254929in}%
\pgfsys@useobject{currentmarker}{}%
\end{pgfscope}%
\end{pgfscope}%
\begin{pgfscope}%
\pgfsetbuttcap%
\pgfsetroundjoin%
\definecolor{currentfill}{rgb}{0.000000,0.000000,0.000000}%
\pgfsetfillcolor{currentfill}%
\pgfsetlinewidth{0.401500pt}%
\definecolor{currentstroke}{rgb}{0.000000,0.000000,0.000000}%
\pgfsetstrokecolor{currentstroke}%
\pgfsetdash{}{0pt}%
\pgfsys@defobject{currentmarker}{\pgfqpoint{-0.020833in}{0.000000in}}{\pgfqpoint{-0.000000in}{0.000000in}}{%
\pgfpathmoveto{\pgfqpoint{-0.000000in}{0.000000in}}%
\pgfpathlineto{\pgfqpoint{-0.020833in}{0.000000in}}%
\pgfusepath{stroke,fill}%
}%
\begin{pgfscope}%
\pgfsys@transformshift{3.038110in}{2.254929in}%
\pgfsys@useobject{currentmarker}{}%
\end{pgfscope}%
\end{pgfscope}%
\begin{pgfscope}%
\pgfsetbuttcap%
\pgfsetroundjoin%
\definecolor{currentfill}{rgb}{0.000000,0.000000,0.000000}%
\pgfsetfillcolor{currentfill}%
\pgfsetlinewidth{0.401500pt}%
\definecolor{currentstroke}{rgb}{0.000000,0.000000,0.000000}%
\pgfsetstrokecolor{currentstroke}%
\pgfsetdash{}{0pt}%
\pgfsys@defobject{currentmarker}{\pgfqpoint{0.000000in}{0.000000in}}{\pgfqpoint{0.020833in}{0.000000in}}{%
\pgfpathmoveto{\pgfqpoint{0.000000in}{0.000000in}}%
\pgfpathlineto{\pgfqpoint{0.020833in}{0.000000in}}%
\pgfusepath{stroke,fill}%
}%
\begin{pgfscope}%
\pgfsys@transformshift{0.650000in}{2.334709in}%
\pgfsys@useobject{currentmarker}{}%
\end{pgfscope}%
\end{pgfscope}%
\begin{pgfscope}%
\pgfsetbuttcap%
\pgfsetroundjoin%
\definecolor{currentfill}{rgb}{0.000000,0.000000,0.000000}%
\pgfsetfillcolor{currentfill}%
\pgfsetlinewidth{0.401500pt}%
\definecolor{currentstroke}{rgb}{0.000000,0.000000,0.000000}%
\pgfsetstrokecolor{currentstroke}%
\pgfsetdash{}{0pt}%
\pgfsys@defobject{currentmarker}{\pgfqpoint{-0.020833in}{0.000000in}}{\pgfqpoint{-0.000000in}{0.000000in}}{%
\pgfpathmoveto{\pgfqpoint{-0.000000in}{0.000000in}}%
\pgfpathlineto{\pgfqpoint{-0.020833in}{0.000000in}}%
\pgfusepath{stroke,fill}%
}%
\begin{pgfscope}%
\pgfsys@transformshift{3.038110in}{2.334709in}%
\pgfsys@useobject{currentmarker}{}%
\end{pgfscope}%
\end{pgfscope}%
\begin{pgfscope}%
\definecolor{textcolor}{rgb}{0.000000,0.000000,0.000000}%
\pgfsetstrokecolor{textcolor}%
\pgfsetfillcolor{textcolor}%
\pgftext[x=0.421528in,y=1.417244in,,bottom,rotate=90.000000]{\color{textcolor}{\rmfamily\fontsize{12.000000}{14.400000}\selectfont\catcode`\^=\active\def^{\ifmmode\sp\else\^{}\fi}\catcode`\%=\active\def%{\%}$U^+$}}%
\end{pgfscope}%
\begin{pgfscope}%
\pgfpathrectangle{\pgfqpoint{0.650000in}{0.420000in}}{\pgfqpoint{2.388110in}{1.994488in}}%
\pgfusepath{clip}%
\pgfsetbuttcap%
\pgfsetroundjoin%
\pgfsetlinewidth{0.501875pt}%
\definecolor{currentstroke}{rgb}{0.000000,0.000000,0.000000}%
\pgfsetstrokecolor{currentstroke}%
\pgfsetdash{{1.850000pt}{0.800000pt}}{0.000000pt}%
\pgfpathmoveto{\pgfqpoint{0.645000in}{0.461627in}}%
\pgfpathlineto{\pgfqpoint{0.792011in}{0.481369in}}%
\pgfpathlineto{\pgfqpoint{0.925675in}{0.512053in}}%
\pgfpathlineto{\pgfqpoint{1.020511in}{0.542738in}}%
\pgfpathlineto{\pgfqpoint{1.094072in}{0.573422in}}%
\pgfpathlineto{\pgfqpoint{1.154176in}{0.604107in}}%
\pgfpathlineto{\pgfqpoint{1.204993in}{0.634791in}}%
\pgfpathlineto{\pgfqpoint{1.249012in}{0.665475in}}%
\pgfpathlineto{\pgfqpoint{1.287840in}{0.696160in}}%
\pgfpathlineto{\pgfqpoint{1.322573in}{0.726844in}}%
\pgfpathlineto{\pgfqpoint{1.353993in}{0.757529in}}%
\pgfpathlineto{\pgfqpoint{1.382677in}{0.788213in}}%
\pgfpathlineto{\pgfqpoint{1.409063in}{0.818898in}}%
\pgfpathlineto{\pgfqpoint{1.433493in}{0.849582in}}%
\pgfpathlineto{\pgfqpoint{1.456237in}{0.880266in}}%
\pgfpathlineto{\pgfqpoint{1.477513in}{0.910951in}}%
\pgfpathlineto{\pgfqpoint{1.497498in}{0.941635in}}%
\pgfpathlineto{\pgfqpoint{1.516341in}{0.972320in}}%
\pgfpathlineto{\pgfqpoint{1.534165in}{1.003004in}}%
\pgfpathlineto{\pgfqpoint{1.551074in}{1.033689in}}%
\pgfpathlineto{\pgfqpoint{1.567158in}{1.064373in}}%
\pgfpathlineto{\pgfqpoint{1.582493in}{1.095058in}}%
\pgfpathlineto{\pgfqpoint{1.597147in}{1.125742in}}%
\pgfpathlineto{\pgfqpoint{1.611177in}{1.156426in}}%
\pgfpathlineto{\pgfqpoint{1.624635in}{1.187111in}}%
\pgfpathlineto{\pgfqpoint{1.637564in}{1.217795in}}%
\pgfpathlineto{\pgfqpoint{1.650005in}{1.248480in}}%
\pgfpathlineto{\pgfqpoint{1.661994in}{1.279164in}}%
\pgfpathlineto{\pgfqpoint{1.673562in}{1.309849in}}%
\pgfpathlineto{\pgfqpoint{1.684738in}{1.340533in}}%
\pgfpathlineto{\pgfqpoint{1.695548in}{1.371217in}}%
\pgfpathlineto{\pgfqpoint{1.706014in}{1.401902in}}%
\pgfpathlineto{\pgfqpoint{1.716158in}{1.432586in}}%
\pgfpathlineto{\pgfqpoint{1.725999in}{1.463271in}}%
\pgfpathlineto{\pgfqpoint{1.735555in}{1.493955in}}%
\pgfpathlineto{\pgfqpoint{1.744842in}{1.524640in}}%
\pgfpathlineto{\pgfqpoint{1.753874in}{1.555324in}}%
\pgfpathlineto{\pgfqpoint{1.762665in}{1.586008in}}%
\pgfpathlineto{\pgfqpoint{1.771228in}{1.616693in}}%
\pgfusepath{stroke}%
\end{pgfscope}%
\begin{pgfscope}%
\pgfpathrectangle{\pgfqpoint{0.650000in}{0.420000in}}{\pgfqpoint{2.388110in}{1.994488in}}%
\pgfusepath{clip}%
\pgfsetbuttcap%
\pgfsetroundjoin%
\pgfsetlinewidth{0.501875pt}%
\definecolor{currentstroke}{rgb}{0.000000,0.000000,0.000000}%
\pgfsetstrokecolor{currentstroke}%
\pgfsetdash{{1.850000pt}{0.800000pt}}{0.000000pt}%
\pgfpathmoveto{\pgfqpoint{1.637564in}{1.282900in}}%
\pgfpathlineto{\pgfqpoint{3.043110in}{2.112542in}}%
\pgfusepath{stroke}%
\end{pgfscope}%
\begin{pgfscope}%
\pgfpathrectangle{\pgfqpoint{0.650000in}{0.420000in}}{\pgfqpoint{2.388110in}{1.994488in}}%
\pgfusepath{clip}%
\pgfsetrectcap%
\pgfsetroundjoin%
\pgfsetlinewidth{0.602250pt}%
\definecolor{currentstroke}{rgb}{0.000000,0.000000,0.000000}%
\pgfsetstrokecolor{currentstroke}%
\pgfsetdash{}{0pt}%
\pgfpathmoveto{\pgfqpoint{0.645000in}{0.460932in}}%
\pgfpathlineto{\pgfqpoint{0.742708in}{0.472840in}}%
\pgfpathlineto{\pgfqpoint{0.889821in}{0.502542in}}%
\pgfpathlineto{\pgfqpoint{1.010016in}{0.538782in}}%
\pgfpathlineto{\pgfqpoint{1.111637in}{0.581437in}}%
\pgfpathlineto{\pgfqpoint{1.199660in}{0.630235in}}%
\pgfpathlineto{\pgfqpoint{1.277299in}{0.684661in}}%
\pgfpathlineto{\pgfqpoint{1.346745in}{0.743875in}}%
\pgfpathlineto{\pgfqpoint{1.409562in}{0.806669in}}%
\pgfpathlineto{\pgfqpoint{1.466906in}{0.871526in}}%
\pgfpathlineto{\pgfqpoint{1.519654in}{0.936757in}}%
\pgfpathlineto{\pgfqpoint{1.568486in}{1.000709in}}%
\pgfpathlineto{\pgfqpoint{1.696400in}{1.172527in}}%
\pgfpathlineto{\pgfqpoint{1.734049in}{1.220921in}}%
\pgfpathlineto{\pgfqpoint{1.769658in}{1.264656in}}%
\pgfpathlineto{\pgfqpoint{1.803436in}{1.303962in}}%
\pgfpathlineto{\pgfqpoint{1.835562in}{1.339191in}}%
\pgfpathlineto{\pgfqpoint{1.866188in}{1.370751in}}%
\pgfpathlineto{\pgfqpoint{1.895450in}{1.399062in}}%
\pgfpathlineto{\pgfqpoint{1.923461in}{1.424533in}}%
\pgfpathlineto{\pgfqpoint{1.950325in}{1.447544in}}%
\pgfpathlineto{\pgfqpoint{1.976131in}{1.468444in}}%
\pgfpathlineto{\pgfqpoint{2.000959in}{1.487545in}}%
\pgfpathlineto{\pgfqpoint{2.047957in}{1.521427in}}%
\pgfpathlineto{\pgfqpoint{2.091803in}{1.551030in}}%
\pgfpathlineto{\pgfqpoint{2.171545in}{1.602537in}}%
\pgfpathlineto{\pgfqpoint{2.225539in}{1.637608in}}%
\pgfpathlineto{\pgfqpoint{2.275392in}{1.671425in}}%
\pgfpathlineto{\pgfqpoint{2.306620in}{1.693695in}}%
\pgfpathlineto{\pgfqpoint{2.336412in}{1.715934in}}%
\pgfpathlineto{\pgfqpoint{2.364891in}{1.738223in}}%
\pgfpathlineto{\pgfqpoint{2.392168in}{1.760564in}}%
\pgfpathlineto{\pgfqpoint{2.431033in}{1.794182in}}%
\pgfpathlineto{\pgfqpoint{2.467681in}{1.827844in}}%
\pgfpathlineto{\pgfqpoint{2.502345in}{1.861306in}}%
\pgfpathlineto{\pgfqpoint{2.545818in}{1.905111in}}%
\pgfpathlineto{\pgfqpoint{2.605912in}{1.965978in}}%
\pgfpathlineto{\pgfqpoint{2.624738in}{1.983805in}}%
\pgfpathlineto{\pgfqpoint{2.643017in}{1.999782in}}%
\pgfpathlineto{\pgfqpoint{2.660778in}{2.013523in}}%
\pgfpathlineto{\pgfqpoint{2.678050in}{2.024725in}}%
\pgfpathlineto{\pgfqpoint{2.694856in}{2.033312in}}%
\pgfpathlineto{\pgfqpoint{2.711220in}{2.039458in}}%
\pgfpathlineto{\pgfqpoint{2.727164in}{2.043535in}}%
\pgfpathlineto{\pgfqpoint{2.742707in}{2.046038in}}%
\pgfpathlineto{\pgfqpoint{2.757868in}{2.047455in}}%
\pgfpathlineto{\pgfqpoint{2.779929in}{2.048402in}}%
\pgfpathlineto{\pgfqpoint{2.828500in}{2.048835in}}%
\pgfpathlineto{\pgfqpoint{3.043110in}{2.049429in}}%
\pgfpathlineto{\pgfqpoint{3.043110in}{2.049429in}}%
\pgfusepath{stroke}%
\end{pgfscope}%
\begin{pgfscope}%
\pgfpathrectangle{\pgfqpoint{0.650000in}{0.420000in}}{\pgfqpoint{2.388110in}{1.994488in}}%
\pgfusepath{clip}%
\pgfsetrectcap%
\pgfsetroundjoin%
\pgfsetlinewidth{1.003750pt}%
\definecolor{currentstroke}{rgb}{0.870588,0.466667,0.682353}%
\pgfsetstrokecolor{currentstroke}%
\pgfsetdash{}{0pt}%
\pgfpathmoveto{\pgfqpoint{0.645000in}{0.459272in}}%
\pgfpathlineto{\pgfqpoint{0.686817in}{0.464539in}}%
\pgfpathlineto{\pgfqpoint{0.749359in}{0.473839in}}%
\pgfpathlineto{\pgfqpoint{0.805873in}{0.483901in}}%
\pgfpathlineto{\pgfqpoint{0.857743in}{0.494779in}}%
\pgfpathlineto{\pgfqpoint{0.905952in}{0.506538in}}%
\pgfpathlineto{\pgfqpoint{0.951207in}{0.519248in}}%
\pgfpathlineto{\pgfqpoint{0.994037in}{0.532980in}}%
\pgfpathlineto{\pgfqpoint{1.034846in}{0.547811in}}%
\pgfpathlineto{\pgfqpoint{1.073948in}{0.563823in}}%
\pgfpathlineto{\pgfqpoint{1.111595in}{0.581096in}}%
\pgfpathlineto{\pgfqpoint{1.147990in}{0.599714in}}%
\pgfpathlineto{\pgfqpoint{1.183300in}{0.619760in}}%
\pgfpathlineto{\pgfqpoint{1.217661in}{0.641312in}}%
\pgfpathlineto{\pgfqpoint{1.251191in}{0.664443in}}%
\pgfpathlineto{\pgfqpoint{1.283987in}{0.689213in}}%
\pgfpathlineto{\pgfqpoint{1.316132in}{0.715665in}}%
\pgfpathlineto{\pgfqpoint{1.347698in}{0.743822in}}%
\pgfpathlineto{\pgfqpoint{1.378746in}{0.773676in}}%
\pgfpathlineto{\pgfqpoint{1.409332in}{0.805184in}}%
\pgfpathlineto{\pgfqpoint{1.439502in}{0.838263in}}%
\pgfpathlineto{\pgfqpoint{1.469297in}{0.872784in}}%
\pgfpathlineto{\pgfqpoint{1.498753in}{0.908572in}}%
\pgfpathlineto{\pgfqpoint{1.527904in}{0.945405in}}%
\pgfpathlineto{\pgfqpoint{1.571118in}{1.002028in}}%
\pgfpathlineto{\pgfqpoint{1.655995in}{1.116399in}}%
\pgfpathlineto{\pgfqpoint{1.697791in}{1.172022in}}%
\pgfpathlineto{\pgfqpoint{1.739229in}{1.225357in}}%
\pgfpathlineto{\pgfqpoint{1.766679in}{1.259292in}}%
\pgfpathlineto{\pgfqpoint{1.794002in}{1.291756in}}%
\pgfpathlineto{\pgfqpoint{1.821209in}{1.322650in}}%
\pgfpathlineto{\pgfqpoint{1.848310in}{1.351921in}}%
\pgfpathlineto{\pgfqpoint{1.875312in}{1.379561in}}%
\pgfpathlineto{\pgfqpoint{1.902225in}{1.405596in}}%
\pgfpathlineto{\pgfqpoint{1.929055in}{1.430084in}}%
\pgfpathlineto{\pgfqpoint{1.955810in}{1.453116in}}%
\pgfpathlineto{\pgfqpoint{1.982494in}{1.474800in}}%
\pgfpathlineto{\pgfqpoint{2.009115in}{1.495253in}}%
\pgfpathlineto{\pgfqpoint{2.035676in}{1.514602in}}%
\pgfpathlineto{\pgfqpoint{2.075418in}{1.541868in}}%
\pgfpathlineto{\pgfqpoint{2.115051in}{1.567526in}}%
\pgfpathlineto{\pgfqpoint{2.180900in}{1.608496in}}%
\pgfpathlineto{\pgfqpoint{2.259636in}{1.657673in}}%
\pgfpathlineto{\pgfqpoint{2.311986in}{1.691740in}}%
\pgfpathlineto{\pgfqpoint{2.351187in}{1.718568in}}%
\pgfpathlineto{\pgfqpoint{2.390341in}{1.747057in}}%
\pgfpathlineto{\pgfqpoint{2.416421in}{1.767274in}}%
\pgfpathlineto{\pgfqpoint{2.442484in}{1.788639in}}%
\pgfpathlineto{\pgfqpoint{2.468532in}{1.811118in}}%
\pgfpathlineto{\pgfqpoint{2.507577in}{1.846624in}}%
\pgfpathlineto{\pgfqpoint{2.546569in}{1.884224in}}%
\pgfpathlineto{\pgfqpoint{2.584795in}{1.922960in}}%
\pgfpathlineto{\pgfqpoint{2.652123in}{1.992135in}}%
\pgfpathlineto{\pgfqpoint{2.671840in}{2.011094in}}%
\pgfpathlineto{\pgfqpoint{2.690445in}{2.027668in}}%
\pgfpathlineto{\pgfqpoint{2.708055in}{2.041721in}}%
\pgfpathlineto{\pgfqpoint{2.724772in}{2.053158in}}%
\pgfpathlineto{\pgfqpoint{2.740682in}{2.062015in}}%
\pgfpathlineto{\pgfqpoint{2.755859in}{2.068486in}}%
\pgfpathlineto{\pgfqpoint{2.770369in}{2.072932in}}%
\pgfpathlineto{\pgfqpoint{2.784266in}{2.075802in}}%
\pgfpathlineto{\pgfqpoint{2.797601in}{2.077522in}}%
\pgfpathlineto{\pgfqpoint{2.816644in}{2.078704in}}%
\pgfpathlineto{\pgfqpoint{2.846125in}{2.079094in}}%
\pgfpathlineto{\pgfqpoint{2.959536in}{2.078605in}}%
\pgfpathlineto{\pgfqpoint{3.043110in}{2.078280in}}%
\pgfpathlineto{\pgfqpoint{3.043110in}{2.078280in}}%
\pgfusepath{stroke}%
\end{pgfscope}%
\begin{pgfscope}%
\pgfpathrectangle{\pgfqpoint{0.650000in}{0.420000in}}{\pgfqpoint{2.388110in}{1.994488in}}%
\pgfusepath{clip}%
\pgfsetrectcap%
\pgfsetroundjoin%
\pgfsetlinewidth{1.003750pt}%
\definecolor{currentstroke}{rgb}{0.498039,0.737255,0.254902}%
\pgfsetstrokecolor{currentstroke}%
\pgfsetdash{}{0pt}%
\pgfpathmoveto{\pgfqpoint{0.645000in}{0.459276in}}%
\pgfpathlineto{\pgfqpoint{0.688321in}{0.464742in}}%
\pgfpathlineto{\pgfqpoint{0.750863in}{0.474085in}}%
\pgfpathlineto{\pgfqpoint{0.807378in}{0.484192in}}%
\pgfpathlineto{\pgfqpoint{0.859248in}{0.495119in}}%
\pgfpathlineto{\pgfqpoint{0.907456in}{0.506932in}}%
\pgfpathlineto{\pgfqpoint{0.952711in}{0.519699in}}%
\pgfpathlineto{\pgfqpoint{0.995541in}{0.533493in}}%
\pgfpathlineto{\pgfqpoint{1.036350in}{0.548391in}}%
\pgfpathlineto{\pgfqpoint{1.075452in}{0.564474in}}%
\pgfpathlineto{\pgfqpoint{1.113099in}{0.581823in}}%
\pgfpathlineto{\pgfqpoint{1.149495in}{0.600523in}}%
\pgfpathlineto{\pgfqpoint{1.184804in}{0.620655in}}%
\pgfpathlineto{\pgfqpoint{1.219166in}{0.642298in}}%
\pgfpathlineto{\pgfqpoint{1.252695in}{0.665524in}}%
\pgfpathlineto{\pgfqpoint{1.285491in}{0.690393in}}%
\pgfpathlineto{\pgfqpoint{1.317636in}{0.716946in}}%
\pgfpathlineto{\pgfqpoint{1.349202in}{0.745205in}}%
\pgfpathlineto{\pgfqpoint{1.380251in}{0.775161in}}%
\pgfpathlineto{\pgfqpoint{1.410836in}{0.806769in}}%
\pgfpathlineto{\pgfqpoint{1.441006in}{0.839943in}}%
\pgfpathlineto{\pgfqpoint{1.470801in}{0.874551in}}%
\pgfpathlineto{\pgfqpoint{1.500257in}{0.910414in}}%
\pgfpathlineto{\pgfqpoint{1.529408in}{0.947306in}}%
\pgfpathlineto{\pgfqpoint{1.572622in}{1.003976in}}%
\pgfpathlineto{\pgfqpoint{1.713145in}{1.191587in}}%
\pgfpathlineto{\pgfqpoint{1.754475in}{1.243634in}}%
\pgfpathlineto{\pgfqpoint{1.781860in}{1.276565in}}%
\pgfpathlineto{\pgfqpoint{1.809124in}{1.307946in}}%
\pgfpathlineto{\pgfqpoint{1.836277in}{1.337717in}}%
\pgfpathlineto{\pgfqpoint{1.863327in}{1.365857in}}%
\pgfpathlineto{\pgfqpoint{1.890284in}{1.392381in}}%
\pgfpathlineto{\pgfqpoint{1.917154in}{1.417328in}}%
\pgfpathlineto{\pgfqpoint{1.943946in}{1.440767in}}%
\pgfpathlineto{\pgfqpoint{1.970665in}{1.462793in}}%
\pgfpathlineto{\pgfqpoint{1.997317in}{1.483521in}}%
\pgfpathlineto{\pgfqpoint{2.023907in}{1.503074in}}%
\pgfpathlineto{\pgfqpoint{2.050441in}{1.521586in}}%
\pgfpathlineto{\pgfqpoint{2.090145in}{1.547725in}}%
\pgfpathlineto{\pgfqpoint{2.129745in}{1.572422in}}%
\pgfpathlineto{\pgfqpoint{2.274237in}{1.660758in}}%
\pgfpathlineto{\pgfqpoint{2.313490in}{1.686350in}}%
\pgfpathlineto{\pgfqpoint{2.352691in}{1.713259in}}%
\pgfpathlineto{\pgfqpoint{2.391846in}{1.741934in}}%
\pgfpathlineto{\pgfqpoint{2.417926in}{1.762244in}}%
\pgfpathlineto{\pgfqpoint{2.443989in}{1.783567in}}%
\pgfpathlineto{\pgfqpoint{2.483055in}{1.817384in}}%
\pgfpathlineto{\pgfqpoint{2.522089in}{1.853077in}}%
\pgfpathlineto{\pgfqpoint{2.560996in}{1.890312in}}%
\pgfpathlineto{\pgfqpoint{2.598477in}{1.928063in}}%
\pgfpathlineto{\pgfqpoint{2.653627in}{1.984509in}}%
\pgfpathlineto{\pgfqpoint{2.673345in}{2.003481in}}%
\pgfpathlineto{\pgfqpoint{2.691949in}{2.020025in}}%
\pgfpathlineto{\pgfqpoint{2.709559in}{2.034016in}}%
\pgfpathlineto{\pgfqpoint{2.726276in}{2.045399in}}%
\pgfpathlineto{\pgfqpoint{2.742186in}{2.054262in}}%
\pgfpathlineto{\pgfqpoint{2.757364in}{2.060821in}}%
\pgfpathlineto{\pgfqpoint{2.771873in}{2.065326in}}%
\pgfpathlineto{\pgfqpoint{2.785770in}{2.068144in}}%
\pgfpathlineto{\pgfqpoint{2.799106in}{2.069777in}}%
\pgfpathlineto{\pgfqpoint{2.818149in}{2.070944in}}%
\pgfpathlineto{\pgfqpoint{2.847629in}{2.071387in}}%
\pgfpathlineto{\pgfqpoint{2.952945in}{2.070985in}}%
\pgfpathlineto{\pgfqpoint{3.043110in}{2.070616in}}%
\pgfpathlineto{\pgfqpoint{3.043110in}{2.070616in}}%
\pgfusepath{stroke}%
\end{pgfscope}%
\begin{pgfscope}%
\pgfpathrectangle{\pgfqpoint{0.650000in}{0.420000in}}{\pgfqpoint{2.388110in}{1.994488in}}%
\pgfusepath{clip}%
\pgfsetrectcap%
\pgfsetroundjoin%
\pgfsetlinewidth{1.003750pt}%
\definecolor{currentstroke}{rgb}{0.400000,0.760784,0.643137}%
\pgfsetstrokecolor{currentstroke}%
\pgfsetdash{}{0pt}%
\pgfpathmoveto{\pgfqpoint{0.645000in}{0.461652in}}%
\pgfpathlineto{\pgfqpoint{0.795375in}{0.481942in}}%
\pgfpathlineto{\pgfqpoint{0.932330in}{0.513817in}}%
\pgfpathlineto{\pgfqpoint{1.030468in}{0.546267in}}%
\pgfpathlineto{\pgfqpoint{1.107341in}{0.579248in}}%
\pgfpathlineto{\pgfqpoint{1.170768in}{0.612694in}}%
\pgfpathlineto{\pgfqpoint{1.224919in}{0.646518in}}%
\pgfpathlineto{\pgfqpoint{1.272283in}{0.680614in}}%
\pgfpathlineto{\pgfqpoint{1.314467in}{0.714859in}}%
\pgfpathlineto{\pgfqpoint{1.352566in}{0.749116in}}%
\pgfpathlineto{\pgfqpoint{1.387363in}{0.783244in}}%
\pgfpathlineto{\pgfqpoint{1.419435in}{0.817098in}}%
\pgfpathlineto{\pgfqpoint{1.449220in}{0.850537in}}%
\pgfpathlineto{\pgfqpoint{1.477059in}{0.883428in}}%
\pgfpathlineto{\pgfqpoint{1.503223in}{0.915648in}}%
\pgfpathlineto{\pgfqpoint{1.551356in}{0.977659in}}%
\pgfpathlineto{\pgfqpoint{1.615321in}{1.063483in}}%
\pgfpathlineto{\pgfqpoint{1.689425in}{1.163026in}}%
\pgfpathlineto{\pgfqpoint{1.722877in}{1.206442in}}%
\pgfpathlineto{\pgfqpoint{1.754414in}{1.245893in}}%
\pgfpathlineto{\pgfqpoint{1.784307in}{1.281675in}}%
\pgfpathlineto{\pgfqpoint{1.812775in}{1.314119in}}%
\pgfpathlineto{\pgfqpoint{1.839994in}{1.343561in}}%
\pgfpathlineto{\pgfqpoint{1.866113in}{1.370323in}}%
\pgfpathlineto{\pgfqpoint{1.891255in}{1.394708in}}%
\pgfpathlineto{\pgfqpoint{1.915524in}{1.417000in}}%
\pgfpathlineto{\pgfqpoint{1.939010in}{1.437458in}}%
\pgfpathlineto{\pgfqpoint{1.972935in}{1.465229in}}%
\pgfpathlineto{\pgfqpoint{2.005484in}{1.490123in}}%
\pgfpathlineto{\pgfqpoint{2.036837in}{1.512734in}}%
\pgfpathlineto{\pgfqpoint{2.077027in}{1.540175in}}%
\pgfpathlineto{\pgfqpoint{2.125063in}{1.571411in}}%
\pgfpathlineto{\pgfqpoint{2.258177in}{1.656574in}}%
\pgfpathlineto{\pgfqpoint{2.307984in}{1.690688in}}%
\pgfpathlineto{\pgfqpoint{2.348355in}{1.719807in}}%
\pgfpathlineto{\pgfqpoint{2.380011in}{1.744003in}}%
\pgfpathlineto{\pgfqpoint{2.411163in}{1.769252in}}%
\pgfpathlineto{\pgfqpoint{2.441864in}{1.795464in}}%
\pgfpathlineto{\pgfqpoint{2.472160in}{1.822638in}}%
\pgfpathlineto{\pgfqpoint{2.509520in}{1.857866in}}%
\pgfpathlineto{\pgfqpoint{2.539043in}{1.887180in}}%
\pgfpathlineto{\pgfqpoint{2.582792in}{1.932392in}}%
\pgfpathlineto{\pgfqpoint{2.661582in}{2.014915in}}%
\pgfpathlineto{\pgfqpoint{2.682800in}{2.035040in}}%
\pgfpathlineto{\pgfqpoint{2.696888in}{2.047430in}}%
\pgfpathlineto{\pgfqpoint{2.710933in}{2.058480in}}%
\pgfpathlineto{\pgfqpoint{2.724937in}{2.067778in}}%
\pgfpathlineto{\pgfqpoint{2.738901in}{2.075215in}}%
\pgfpathlineto{\pgfqpoint{2.752828in}{2.080842in}}%
\pgfpathlineto{\pgfqpoint{2.766719in}{2.084760in}}%
\pgfpathlineto{\pgfqpoint{2.780575in}{2.087203in}}%
\pgfpathlineto{\pgfqpoint{2.794399in}{2.088607in}}%
\pgfpathlineto{\pgfqpoint{2.815076in}{2.089540in}}%
\pgfpathlineto{\pgfqpoint{2.856234in}{2.089721in}}%
\pgfpathlineto{\pgfqpoint{3.043110in}{2.088740in}}%
\pgfpathlineto{\pgfqpoint{3.043110in}{2.088740in}}%
\pgfusepath{stroke}%
\end{pgfscope}%
\begin{pgfscope}%
\pgfsetrectcap%
\pgfsetmiterjoin%
\pgfsetlinewidth{0.803000pt}%
\definecolor{currentstroke}{rgb}{0.000000,0.000000,0.000000}%
\pgfsetstrokecolor{currentstroke}%
\pgfsetdash{}{0pt}%
\pgfpathmoveto{\pgfqpoint{0.650000in}{0.420000in}}%
\pgfpathlineto{\pgfqpoint{0.650000in}{2.414488in}}%
\pgfusepath{stroke}%
\end{pgfscope}%
\begin{pgfscope}%
\pgfsetrectcap%
\pgfsetmiterjoin%
\pgfsetlinewidth{0.803000pt}%
\definecolor{currentstroke}{rgb}{0.000000,0.000000,0.000000}%
\pgfsetstrokecolor{currentstroke}%
\pgfsetdash{}{0pt}%
\pgfpathmoveto{\pgfqpoint{3.038110in}{0.420000in}}%
\pgfpathlineto{\pgfqpoint{3.038110in}{2.414488in}}%
\pgfusepath{stroke}%
\end{pgfscope}%
\begin{pgfscope}%
\pgfsetrectcap%
\pgfsetmiterjoin%
\pgfsetlinewidth{0.803000pt}%
\definecolor{currentstroke}{rgb}{0.000000,0.000000,0.000000}%
\pgfsetstrokecolor{currentstroke}%
\pgfsetdash{}{0pt}%
\pgfpathmoveto{\pgfqpoint{0.650000in}{0.420000in}}%
\pgfpathlineto{\pgfqpoint{3.038110in}{0.420000in}}%
\pgfusepath{stroke}%
\end{pgfscope}%
\begin{pgfscope}%
\pgfsetrectcap%
\pgfsetmiterjoin%
\pgfsetlinewidth{0.803000pt}%
\definecolor{currentstroke}{rgb}{0.000000,0.000000,0.000000}%
\pgfsetstrokecolor{currentstroke}%
\pgfsetdash{}{0pt}%
\pgfpathmoveto{\pgfqpoint{0.650000in}{2.414488in}}%
\pgfpathlineto{\pgfqpoint{3.038110in}{2.414488in}}%
\pgfusepath{stroke}%
\end{pgfscope}%
\begin{pgfscope}%
\pgfsetrectcap%
\pgfsetroundjoin%
\pgfsetlinewidth{0.602250pt}%
\definecolor{currentstroke}{rgb}{0.000000,0.000000,0.000000}%
\pgfsetstrokecolor{currentstroke}%
\pgfsetdash{}{0pt}%
\pgfpathmoveto{\pgfqpoint{0.750000in}{2.275599in}}%
\pgfpathlineto{\pgfqpoint{0.861111in}{2.275599in}}%
\pgfpathlineto{\pgfqpoint{0.972222in}{2.275599in}}%
\pgfusepath{stroke}%
\end{pgfscope}%
\begin{pgfscope}%
\definecolor{textcolor}{rgb}{0.000000,0.000000,0.000000}%
\pgfsetstrokecolor{textcolor}%
\pgfsetfillcolor{textcolor}%
\pgftext[x=1.061111in,y=2.236710in,left,base]{\color{textcolor}{\rmfamily\fontsize{8.000000}{9.600000}\selectfont\catcode`\^=\active\def^{\ifmmode\sp\else\^{}\fi}\catcode`\%=\active\def%{\%}ref}}%
\end{pgfscope}%
\begin{pgfscope}%
\pgfsetrectcap%
\pgfsetroundjoin%
\pgfsetlinewidth{1.003750pt}%
\definecolor{currentstroke}{rgb}{0.870588,0.466667,0.682353}%
\pgfsetstrokecolor{currentstroke}%
\pgfsetdash{}{0pt}%
\pgfpathmoveto{\pgfqpoint{0.750000in}{2.120661in}}%
\pgfpathlineto{\pgfqpoint{0.861111in}{2.120661in}}%
\pgfpathlineto{\pgfqpoint{0.972222in}{2.120661in}}%
\pgfusepath{stroke}%
\end{pgfscope}%
\begin{pgfscope}%
\definecolor{textcolor}{rgb}{0.000000,0.000000,0.000000}%
\pgfsetstrokecolor{textcolor}%
\pgfsetfillcolor{textcolor}%
\pgftext[x=1.061111in,y=2.081772in,left,base]{\color{textcolor}{\rmfamily\fontsize{8.000000}{9.600000}\selectfont\catcode`\^=\active\def^{\ifmmode\sp\else\^{}\fi}\catcode`\%=\active\def%{\%}base}}%
\end{pgfscope}%
\begin{pgfscope}%
\pgfsetrectcap%
\pgfsetroundjoin%
\pgfsetlinewidth{1.003750pt}%
\definecolor{currentstroke}{rgb}{0.498039,0.737255,0.254902}%
\pgfsetstrokecolor{currentstroke}%
\pgfsetdash{}{0pt}%
\pgfpathmoveto{\pgfqpoint{0.750000in}{1.965723in}}%
\pgfpathlineto{\pgfqpoint{0.861111in}{1.965723in}}%
\pgfpathlineto{\pgfqpoint{0.972222in}{1.965723in}}%
\pgfusepath{stroke}%
\end{pgfscope}%
\begin{pgfscope}%
\definecolor{textcolor}{rgb}{0.000000,0.000000,0.000000}%
\pgfsetstrokecolor{textcolor}%
\pgfsetfillcolor{textcolor}%
\pgftext[x=1.061111in,y=1.926834in,left,base]{\color{textcolor}{\rmfamily\fontsize{8.000000}{9.600000}\selectfont\catcode`\^=\active\def^{\ifmmode\sp\else\^{}\fi}\catcode`\%=\active\def%{\%}optimised}}%
\end{pgfscope}%
\begin{pgfscope}%
\pgfsetrectcap%
\pgfsetroundjoin%
\pgfsetlinewidth{1.003750pt}%
\definecolor{currentstroke}{rgb}{0.400000,0.760784,0.643137}%
\pgfsetstrokecolor{currentstroke}%
\pgfsetdash{}{0pt}%
\pgfpathmoveto{\pgfqpoint{0.750000in}{1.810784in}}%
\pgfpathlineto{\pgfqpoint{0.861111in}{1.810784in}}%
\pgfpathlineto{\pgfqpoint{0.972222in}{1.810784in}}%
\pgfusepath{stroke}%
\end{pgfscope}%
\begin{pgfscope}%
\definecolor{textcolor}{rgb}{0.000000,0.000000,0.000000}%
\pgfsetstrokecolor{textcolor}%
\pgfsetfillcolor{textcolor}%
\pgftext[x=1.061111in,y=1.771895in,left,base]{\color{textcolor}{\rmfamily\fontsize{8.000000}{9.600000}\selectfont\catcode`\^=\active\def^{\ifmmode\sp\else\^{}\fi}\catcode`\%=\active\def%{\%}planeTBL}}%
\end{pgfscope}%
\end{pgfpicture}%
\makeatother%
\endgroup%

    \label{fig:inflowSecVerifU}
\end{subfigure}
\hfill
\begin{subfigure}[t]{0.495\textwidth}
    \centering
    \subcaption{}
    %% Creator: Matplotlib, PGF backend
%%
%% To include the figure in your LaTeX document, write
%%   \input{<filename>.pgf}
%%
%% Make sure the required packages are loaded in your preamble
%%   \usepackage{pgf}
%%
%% Also ensure that all the required font packages are loaded; for instance,
%% the lmodern package is sometimes necessary when using math font.
%%   \usepackage{lmodern}
%%
%% Figures using additional raster images can only be included by \input if
%% they are in the same directory as the main LaTeX file. For loading figures
%% from other directories you can use the `import` package
%%   \usepackage{import}
%%
%% and then include the figures with
%%   \import{<path to file>}{<filename>.pgf}
%%
%% Matplotlib used the following preamble
%%   \def\mathdefault#1{#1}
%%   \everymath=\expandafter{\the\everymath\displaystyle}
%%   \IfFileExists{scrextend.sty}{
%%     \usepackage[fontsize=11.000000pt]{scrextend}
%%   }{
%%     \renewcommand{\normalsize}{\fontsize{11.000000}{13.200000}\selectfont}
%%     \normalsize
%%   }
%%   
%%   \makeatletter\@ifpackageloaded{underscore}{}{\usepackage[strings]{underscore}}\makeatother
%%
\begingroup%
\makeatletter%
\begin{pgfpicture}%
\pgfpathrectangle{\pgfpointorigin}{\pgfqpoint{3.118110in}{2.494488in}}%
\pgfusepath{use as bounding box, clip}%
\begin{pgfscope}%
\pgfsetbuttcap%
\pgfsetmiterjoin%
\definecolor{currentfill}{rgb}{1.000000,1.000000,1.000000}%
\pgfsetfillcolor{currentfill}%
\pgfsetlinewidth{0.000000pt}%
\definecolor{currentstroke}{rgb}{1.000000,1.000000,1.000000}%
\pgfsetstrokecolor{currentstroke}%
\pgfsetdash{}{0pt}%
\pgfpathmoveto{\pgfqpoint{0.000000in}{0.000000in}}%
\pgfpathlineto{\pgfqpoint{3.118110in}{0.000000in}}%
\pgfpathlineto{\pgfqpoint{3.118110in}{2.494488in}}%
\pgfpathlineto{\pgfqpoint{0.000000in}{2.494488in}}%
\pgfpathlineto{\pgfqpoint{0.000000in}{0.000000in}}%
\pgfpathclose%
\pgfusepath{fill}%
\end{pgfscope}%
\begin{pgfscope}%
\pgfsetbuttcap%
\pgfsetmiterjoin%
\definecolor{currentfill}{rgb}{1.000000,1.000000,1.000000}%
\pgfsetfillcolor{currentfill}%
\pgfsetlinewidth{0.000000pt}%
\definecolor{currentstroke}{rgb}{0.000000,0.000000,0.000000}%
\pgfsetstrokecolor{currentstroke}%
\pgfsetstrokeopacity{0.000000}%
\pgfsetdash{}{0pt}%
\pgfpathmoveto{\pgfqpoint{0.650000in}{0.420000in}}%
\pgfpathlineto{\pgfqpoint{3.038110in}{0.420000in}}%
\pgfpathlineto{\pgfqpoint{3.038110in}{2.414488in}}%
\pgfpathlineto{\pgfqpoint{0.650000in}{2.414488in}}%
\pgfpathlineto{\pgfqpoint{0.650000in}{0.420000in}}%
\pgfpathclose%
\pgfusepath{fill}%
\end{pgfscope}%
\begin{pgfscope}%
\pgfsetbuttcap%
\pgfsetroundjoin%
\definecolor{currentfill}{rgb}{0.000000,0.000000,0.000000}%
\pgfsetfillcolor{currentfill}%
\pgfsetlinewidth{0.501875pt}%
\definecolor{currentstroke}{rgb}{0.000000,0.000000,0.000000}%
\pgfsetstrokecolor{currentstroke}%
\pgfsetdash{}{0pt}%
\pgfsys@defobject{currentmarker}{\pgfqpoint{0.000000in}{0.000000in}}{\pgfqpoint{0.000000in}{0.041667in}}{%
\pgfpathmoveto{\pgfqpoint{0.000000in}{0.000000in}}%
\pgfpathlineto{\pgfqpoint{0.000000in}{0.041667in}}%
\pgfusepath{stroke,fill}%
}%
\begin{pgfscope}%
\pgfsys@transformshift{0.878501in}{0.420000in}%
\pgfsys@useobject{currentmarker}{}%
\end{pgfscope}%
\end{pgfscope}%
\begin{pgfscope}%
\pgfsetbuttcap%
\pgfsetroundjoin%
\definecolor{currentfill}{rgb}{0.000000,0.000000,0.000000}%
\pgfsetfillcolor{currentfill}%
\pgfsetlinewidth{0.501875pt}%
\definecolor{currentstroke}{rgb}{0.000000,0.000000,0.000000}%
\pgfsetstrokecolor{currentstroke}%
\pgfsetdash{}{0pt}%
\pgfsys@defobject{currentmarker}{\pgfqpoint{0.000000in}{-0.041667in}}{\pgfqpoint{0.000000in}{0.000000in}}{%
\pgfpathmoveto{\pgfqpoint{0.000000in}{0.000000in}}%
\pgfpathlineto{\pgfqpoint{0.000000in}{-0.041667in}}%
\pgfusepath{stroke,fill}%
}%
\begin{pgfscope}%
\pgfsys@transformshift{0.878501in}{2.414488in}%
\pgfsys@useobject{currentmarker}{}%
\end{pgfscope}%
\end{pgfscope}%
\begin{pgfscope}%
\definecolor{textcolor}{rgb}{0.000000,0.000000,0.000000}%
\pgfsetstrokecolor{textcolor}%
\pgfsetfillcolor{textcolor}%
\pgftext[x=0.878501in,y=0.371389in,,top]{\color{textcolor}{\rmfamily\fontsize{11.000000}{13.200000}\selectfont\catcode`\^=\active\def^{\ifmmode\sp\else\^{}\fi}\catcode`\%=\active\def%{\%}$\mathdefault{10^{0}}$}}%
\end{pgfscope}%
\begin{pgfscope}%
\pgfsetbuttcap%
\pgfsetroundjoin%
\definecolor{currentfill}{rgb}{0.000000,0.000000,0.000000}%
\pgfsetfillcolor{currentfill}%
\pgfsetlinewidth{0.501875pt}%
\definecolor{currentstroke}{rgb}{0.000000,0.000000,0.000000}%
\pgfsetstrokecolor{currentstroke}%
\pgfsetdash{}{0pt}%
\pgfsys@defobject{currentmarker}{\pgfqpoint{0.000000in}{0.000000in}}{\pgfqpoint{0.000000in}{0.041667in}}{%
\pgfpathmoveto{\pgfqpoint{0.000000in}{0.000000in}}%
\pgfpathlineto{\pgfqpoint{0.000000in}{0.041667in}}%
\pgfusepath{stroke,fill}%
}%
\begin{pgfscope}%
\pgfsys@transformshift{1.637564in}{0.420000in}%
\pgfsys@useobject{currentmarker}{}%
\end{pgfscope}%
\end{pgfscope}%
\begin{pgfscope}%
\pgfsetbuttcap%
\pgfsetroundjoin%
\definecolor{currentfill}{rgb}{0.000000,0.000000,0.000000}%
\pgfsetfillcolor{currentfill}%
\pgfsetlinewidth{0.501875pt}%
\definecolor{currentstroke}{rgb}{0.000000,0.000000,0.000000}%
\pgfsetstrokecolor{currentstroke}%
\pgfsetdash{}{0pt}%
\pgfsys@defobject{currentmarker}{\pgfqpoint{0.000000in}{-0.041667in}}{\pgfqpoint{0.000000in}{0.000000in}}{%
\pgfpathmoveto{\pgfqpoint{0.000000in}{0.000000in}}%
\pgfpathlineto{\pgfqpoint{0.000000in}{-0.041667in}}%
\pgfusepath{stroke,fill}%
}%
\begin{pgfscope}%
\pgfsys@transformshift{1.637564in}{2.414488in}%
\pgfsys@useobject{currentmarker}{}%
\end{pgfscope}%
\end{pgfscope}%
\begin{pgfscope}%
\definecolor{textcolor}{rgb}{0.000000,0.000000,0.000000}%
\pgfsetstrokecolor{textcolor}%
\pgfsetfillcolor{textcolor}%
\pgftext[x=1.637564in,y=0.371389in,,top]{\color{textcolor}{\rmfamily\fontsize{11.000000}{13.200000}\selectfont\catcode`\^=\active\def^{\ifmmode\sp\else\^{}\fi}\catcode`\%=\active\def%{\%}$\mathdefault{10^{1}}$}}%
\end{pgfscope}%
\begin{pgfscope}%
\pgfsetbuttcap%
\pgfsetroundjoin%
\definecolor{currentfill}{rgb}{0.000000,0.000000,0.000000}%
\pgfsetfillcolor{currentfill}%
\pgfsetlinewidth{0.501875pt}%
\definecolor{currentstroke}{rgb}{0.000000,0.000000,0.000000}%
\pgfsetstrokecolor{currentstroke}%
\pgfsetdash{}{0pt}%
\pgfsys@defobject{currentmarker}{\pgfqpoint{0.000000in}{0.000000in}}{\pgfqpoint{0.000000in}{0.041667in}}{%
\pgfpathmoveto{\pgfqpoint{0.000000in}{0.000000in}}%
\pgfpathlineto{\pgfqpoint{0.000000in}{0.041667in}}%
\pgfusepath{stroke,fill}%
}%
\begin{pgfscope}%
\pgfsys@transformshift{2.396627in}{0.420000in}%
\pgfsys@useobject{currentmarker}{}%
\end{pgfscope}%
\end{pgfscope}%
\begin{pgfscope}%
\pgfsetbuttcap%
\pgfsetroundjoin%
\definecolor{currentfill}{rgb}{0.000000,0.000000,0.000000}%
\pgfsetfillcolor{currentfill}%
\pgfsetlinewidth{0.501875pt}%
\definecolor{currentstroke}{rgb}{0.000000,0.000000,0.000000}%
\pgfsetstrokecolor{currentstroke}%
\pgfsetdash{}{0pt}%
\pgfsys@defobject{currentmarker}{\pgfqpoint{0.000000in}{-0.041667in}}{\pgfqpoint{0.000000in}{0.000000in}}{%
\pgfpathmoveto{\pgfqpoint{0.000000in}{0.000000in}}%
\pgfpathlineto{\pgfqpoint{0.000000in}{-0.041667in}}%
\pgfusepath{stroke,fill}%
}%
\begin{pgfscope}%
\pgfsys@transformshift{2.396627in}{2.414488in}%
\pgfsys@useobject{currentmarker}{}%
\end{pgfscope}%
\end{pgfscope}%
\begin{pgfscope}%
\definecolor{textcolor}{rgb}{0.000000,0.000000,0.000000}%
\pgfsetstrokecolor{textcolor}%
\pgfsetfillcolor{textcolor}%
\pgftext[x=2.396627in,y=0.371389in,,top]{\color{textcolor}{\rmfamily\fontsize{11.000000}{13.200000}\selectfont\catcode`\^=\active\def^{\ifmmode\sp\else\^{}\fi}\catcode`\%=\active\def%{\%}$\mathdefault{10^{2}}$}}%
\end{pgfscope}%
\begin{pgfscope}%
\pgfsetbuttcap%
\pgfsetroundjoin%
\definecolor{currentfill}{rgb}{0.000000,0.000000,0.000000}%
\pgfsetfillcolor{currentfill}%
\pgfsetlinewidth{0.401500pt}%
\definecolor{currentstroke}{rgb}{0.000000,0.000000,0.000000}%
\pgfsetstrokecolor{currentstroke}%
\pgfsetdash{}{0pt}%
\pgfsys@defobject{currentmarker}{\pgfqpoint{0.000000in}{0.000000in}}{\pgfqpoint{0.000000in}{0.020833in}}{%
\pgfpathmoveto{\pgfqpoint{0.000000in}{0.000000in}}%
\pgfpathlineto{\pgfqpoint{0.000000in}{0.020833in}}%
\pgfusepath{stroke,fill}%
}%
\begin{pgfscope}%
\pgfsys@transformshift{0.650000in}{0.420000in}%
\pgfsys@useobject{currentmarker}{}%
\end{pgfscope}%
\end{pgfscope}%
\begin{pgfscope}%
\pgfsetbuttcap%
\pgfsetroundjoin%
\definecolor{currentfill}{rgb}{0.000000,0.000000,0.000000}%
\pgfsetfillcolor{currentfill}%
\pgfsetlinewidth{0.401500pt}%
\definecolor{currentstroke}{rgb}{0.000000,0.000000,0.000000}%
\pgfsetstrokecolor{currentstroke}%
\pgfsetdash{}{0pt}%
\pgfsys@defobject{currentmarker}{\pgfqpoint{0.000000in}{-0.020833in}}{\pgfqpoint{0.000000in}{0.000000in}}{%
\pgfpathmoveto{\pgfqpoint{0.000000in}{0.000000in}}%
\pgfpathlineto{\pgfqpoint{0.000000in}{-0.020833in}}%
\pgfusepath{stroke,fill}%
}%
\begin{pgfscope}%
\pgfsys@transformshift{0.650000in}{2.414488in}%
\pgfsys@useobject{currentmarker}{}%
\end{pgfscope}%
\end{pgfscope}%
\begin{pgfscope}%
\pgfsetbuttcap%
\pgfsetroundjoin%
\definecolor{currentfill}{rgb}{0.000000,0.000000,0.000000}%
\pgfsetfillcolor{currentfill}%
\pgfsetlinewidth{0.401500pt}%
\definecolor{currentstroke}{rgb}{0.000000,0.000000,0.000000}%
\pgfsetstrokecolor{currentstroke}%
\pgfsetdash{}{0pt}%
\pgfsys@defobject{currentmarker}{\pgfqpoint{0.000000in}{0.000000in}}{\pgfqpoint{0.000000in}{0.020833in}}{%
\pgfpathmoveto{\pgfqpoint{0.000000in}{0.000000in}}%
\pgfpathlineto{\pgfqpoint{0.000000in}{0.020833in}}%
\pgfusepath{stroke,fill}%
}%
\begin{pgfscope}%
\pgfsys@transformshift{0.710104in}{0.420000in}%
\pgfsys@useobject{currentmarker}{}%
\end{pgfscope}%
\end{pgfscope}%
\begin{pgfscope}%
\pgfsetbuttcap%
\pgfsetroundjoin%
\definecolor{currentfill}{rgb}{0.000000,0.000000,0.000000}%
\pgfsetfillcolor{currentfill}%
\pgfsetlinewidth{0.401500pt}%
\definecolor{currentstroke}{rgb}{0.000000,0.000000,0.000000}%
\pgfsetstrokecolor{currentstroke}%
\pgfsetdash{}{0pt}%
\pgfsys@defobject{currentmarker}{\pgfqpoint{0.000000in}{-0.020833in}}{\pgfqpoint{0.000000in}{0.000000in}}{%
\pgfpathmoveto{\pgfqpoint{0.000000in}{0.000000in}}%
\pgfpathlineto{\pgfqpoint{0.000000in}{-0.020833in}}%
\pgfusepath{stroke,fill}%
}%
\begin{pgfscope}%
\pgfsys@transformshift{0.710104in}{2.414488in}%
\pgfsys@useobject{currentmarker}{}%
\end{pgfscope}%
\end{pgfscope}%
\begin{pgfscope}%
\pgfsetbuttcap%
\pgfsetroundjoin%
\definecolor{currentfill}{rgb}{0.000000,0.000000,0.000000}%
\pgfsetfillcolor{currentfill}%
\pgfsetlinewidth{0.401500pt}%
\definecolor{currentstroke}{rgb}{0.000000,0.000000,0.000000}%
\pgfsetstrokecolor{currentstroke}%
\pgfsetdash{}{0pt}%
\pgfsys@defobject{currentmarker}{\pgfqpoint{0.000000in}{0.000000in}}{\pgfqpoint{0.000000in}{0.020833in}}{%
\pgfpathmoveto{\pgfqpoint{0.000000in}{0.000000in}}%
\pgfpathlineto{\pgfqpoint{0.000000in}{0.020833in}}%
\pgfusepath{stroke,fill}%
}%
\begin{pgfscope}%
\pgfsys@transformshift{0.760920in}{0.420000in}%
\pgfsys@useobject{currentmarker}{}%
\end{pgfscope}%
\end{pgfscope}%
\begin{pgfscope}%
\pgfsetbuttcap%
\pgfsetroundjoin%
\definecolor{currentfill}{rgb}{0.000000,0.000000,0.000000}%
\pgfsetfillcolor{currentfill}%
\pgfsetlinewidth{0.401500pt}%
\definecolor{currentstroke}{rgb}{0.000000,0.000000,0.000000}%
\pgfsetstrokecolor{currentstroke}%
\pgfsetdash{}{0pt}%
\pgfsys@defobject{currentmarker}{\pgfqpoint{0.000000in}{-0.020833in}}{\pgfqpoint{0.000000in}{0.000000in}}{%
\pgfpathmoveto{\pgfqpoint{0.000000in}{0.000000in}}%
\pgfpathlineto{\pgfqpoint{0.000000in}{-0.020833in}}%
\pgfusepath{stroke,fill}%
}%
\begin{pgfscope}%
\pgfsys@transformshift{0.760920in}{2.414488in}%
\pgfsys@useobject{currentmarker}{}%
\end{pgfscope}%
\end{pgfscope}%
\begin{pgfscope}%
\pgfsetbuttcap%
\pgfsetroundjoin%
\definecolor{currentfill}{rgb}{0.000000,0.000000,0.000000}%
\pgfsetfillcolor{currentfill}%
\pgfsetlinewidth{0.401500pt}%
\definecolor{currentstroke}{rgb}{0.000000,0.000000,0.000000}%
\pgfsetstrokecolor{currentstroke}%
\pgfsetdash{}{0pt}%
\pgfsys@defobject{currentmarker}{\pgfqpoint{0.000000in}{0.000000in}}{\pgfqpoint{0.000000in}{0.020833in}}{%
\pgfpathmoveto{\pgfqpoint{0.000000in}{0.000000in}}%
\pgfpathlineto{\pgfqpoint{0.000000in}{0.020833in}}%
\pgfusepath{stroke,fill}%
}%
\begin{pgfscope}%
\pgfsys@transformshift{0.804940in}{0.420000in}%
\pgfsys@useobject{currentmarker}{}%
\end{pgfscope}%
\end{pgfscope}%
\begin{pgfscope}%
\pgfsetbuttcap%
\pgfsetroundjoin%
\definecolor{currentfill}{rgb}{0.000000,0.000000,0.000000}%
\pgfsetfillcolor{currentfill}%
\pgfsetlinewidth{0.401500pt}%
\definecolor{currentstroke}{rgb}{0.000000,0.000000,0.000000}%
\pgfsetstrokecolor{currentstroke}%
\pgfsetdash{}{0pt}%
\pgfsys@defobject{currentmarker}{\pgfqpoint{0.000000in}{-0.020833in}}{\pgfqpoint{0.000000in}{0.000000in}}{%
\pgfpathmoveto{\pgfqpoint{0.000000in}{0.000000in}}%
\pgfpathlineto{\pgfqpoint{0.000000in}{-0.020833in}}%
\pgfusepath{stroke,fill}%
}%
\begin{pgfscope}%
\pgfsys@transformshift{0.804940in}{2.414488in}%
\pgfsys@useobject{currentmarker}{}%
\end{pgfscope}%
\end{pgfscope}%
\begin{pgfscope}%
\pgfsetbuttcap%
\pgfsetroundjoin%
\definecolor{currentfill}{rgb}{0.000000,0.000000,0.000000}%
\pgfsetfillcolor{currentfill}%
\pgfsetlinewidth{0.401500pt}%
\definecolor{currentstroke}{rgb}{0.000000,0.000000,0.000000}%
\pgfsetstrokecolor{currentstroke}%
\pgfsetdash{}{0pt}%
\pgfsys@defobject{currentmarker}{\pgfqpoint{0.000000in}{0.000000in}}{\pgfqpoint{0.000000in}{0.020833in}}{%
\pgfpathmoveto{\pgfqpoint{0.000000in}{0.000000in}}%
\pgfpathlineto{\pgfqpoint{0.000000in}{0.020833in}}%
\pgfusepath{stroke,fill}%
}%
\begin{pgfscope}%
\pgfsys@transformshift{0.843768in}{0.420000in}%
\pgfsys@useobject{currentmarker}{}%
\end{pgfscope}%
\end{pgfscope}%
\begin{pgfscope}%
\pgfsetbuttcap%
\pgfsetroundjoin%
\definecolor{currentfill}{rgb}{0.000000,0.000000,0.000000}%
\pgfsetfillcolor{currentfill}%
\pgfsetlinewidth{0.401500pt}%
\definecolor{currentstroke}{rgb}{0.000000,0.000000,0.000000}%
\pgfsetstrokecolor{currentstroke}%
\pgfsetdash{}{0pt}%
\pgfsys@defobject{currentmarker}{\pgfqpoint{0.000000in}{-0.020833in}}{\pgfqpoint{0.000000in}{0.000000in}}{%
\pgfpathmoveto{\pgfqpoint{0.000000in}{0.000000in}}%
\pgfpathlineto{\pgfqpoint{0.000000in}{-0.020833in}}%
\pgfusepath{stroke,fill}%
}%
\begin{pgfscope}%
\pgfsys@transformshift{0.843768in}{2.414488in}%
\pgfsys@useobject{currentmarker}{}%
\end{pgfscope}%
\end{pgfscope}%
\begin{pgfscope}%
\pgfsetbuttcap%
\pgfsetroundjoin%
\definecolor{currentfill}{rgb}{0.000000,0.000000,0.000000}%
\pgfsetfillcolor{currentfill}%
\pgfsetlinewidth{0.401500pt}%
\definecolor{currentstroke}{rgb}{0.000000,0.000000,0.000000}%
\pgfsetstrokecolor{currentstroke}%
\pgfsetdash{}{0pt}%
\pgfsys@defobject{currentmarker}{\pgfqpoint{0.000000in}{0.000000in}}{\pgfqpoint{0.000000in}{0.020833in}}{%
\pgfpathmoveto{\pgfqpoint{0.000000in}{0.000000in}}%
\pgfpathlineto{\pgfqpoint{0.000000in}{0.020833in}}%
\pgfusepath{stroke,fill}%
}%
\begin{pgfscope}%
\pgfsys@transformshift{1.107002in}{0.420000in}%
\pgfsys@useobject{currentmarker}{}%
\end{pgfscope}%
\end{pgfscope}%
\begin{pgfscope}%
\pgfsetbuttcap%
\pgfsetroundjoin%
\definecolor{currentfill}{rgb}{0.000000,0.000000,0.000000}%
\pgfsetfillcolor{currentfill}%
\pgfsetlinewidth{0.401500pt}%
\definecolor{currentstroke}{rgb}{0.000000,0.000000,0.000000}%
\pgfsetstrokecolor{currentstroke}%
\pgfsetdash{}{0pt}%
\pgfsys@defobject{currentmarker}{\pgfqpoint{0.000000in}{-0.020833in}}{\pgfqpoint{0.000000in}{0.000000in}}{%
\pgfpathmoveto{\pgfqpoint{0.000000in}{0.000000in}}%
\pgfpathlineto{\pgfqpoint{0.000000in}{-0.020833in}}%
\pgfusepath{stroke,fill}%
}%
\begin{pgfscope}%
\pgfsys@transformshift{1.107002in}{2.414488in}%
\pgfsys@useobject{currentmarker}{}%
\end{pgfscope}%
\end{pgfscope}%
\begin{pgfscope}%
\pgfsetbuttcap%
\pgfsetroundjoin%
\definecolor{currentfill}{rgb}{0.000000,0.000000,0.000000}%
\pgfsetfillcolor{currentfill}%
\pgfsetlinewidth{0.401500pt}%
\definecolor{currentstroke}{rgb}{0.000000,0.000000,0.000000}%
\pgfsetstrokecolor{currentstroke}%
\pgfsetdash{}{0pt}%
\pgfsys@defobject{currentmarker}{\pgfqpoint{0.000000in}{0.000000in}}{\pgfqpoint{0.000000in}{0.020833in}}{%
\pgfpathmoveto{\pgfqpoint{0.000000in}{0.000000in}}%
\pgfpathlineto{\pgfqpoint{0.000000in}{0.020833in}}%
\pgfusepath{stroke,fill}%
}%
\begin{pgfscope}%
\pgfsys@transformshift{1.240666in}{0.420000in}%
\pgfsys@useobject{currentmarker}{}%
\end{pgfscope}%
\end{pgfscope}%
\begin{pgfscope}%
\pgfsetbuttcap%
\pgfsetroundjoin%
\definecolor{currentfill}{rgb}{0.000000,0.000000,0.000000}%
\pgfsetfillcolor{currentfill}%
\pgfsetlinewidth{0.401500pt}%
\definecolor{currentstroke}{rgb}{0.000000,0.000000,0.000000}%
\pgfsetstrokecolor{currentstroke}%
\pgfsetdash{}{0pt}%
\pgfsys@defobject{currentmarker}{\pgfqpoint{0.000000in}{-0.020833in}}{\pgfqpoint{0.000000in}{0.000000in}}{%
\pgfpathmoveto{\pgfqpoint{0.000000in}{0.000000in}}%
\pgfpathlineto{\pgfqpoint{0.000000in}{-0.020833in}}%
\pgfusepath{stroke,fill}%
}%
\begin{pgfscope}%
\pgfsys@transformshift{1.240666in}{2.414488in}%
\pgfsys@useobject{currentmarker}{}%
\end{pgfscope}%
\end{pgfscope}%
\begin{pgfscope}%
\pgfsetbuttcap%
\pgfsetroundjoin%
\definecolor{currentfill}{rgb}{0.000000,0.000000,0.000000}%
\pgfsetfillcolor{currentfill}%
\pgfsetlinewidth{0.401500pt}%
\definecolor{currentstroke}{rgb}{0.000000,0.000000,0.000000}%
\pgfsetstrokecolor{currentstroke}%
\pgfsetdash{}{0pt}%
\pgfsys@defobject{currentmarker}{\pgfqpoint{0.000000in}{0.000000in}}{\pgfqpoint{0.000000in}{0.020833in}}{%
\pgfpathmoveto{\pgfqpoint{0.000000in}{0.000000in}}%
\pgfpathlineto{\pgfqpoint{0.000000in}{0.020833in}}%
\pgfusepath{stroke,fill}%
}%
\begin{pgfscope}%
\pgfsys@transformshift{1.335502in}{0.420000in}%
\pgfsys@useobject{currentmarker}{}%
\end{pgfscope}%
\end{pgfscope}%
\begin{pgfscope}%
\pgfsetbuttcap%
\pgfsetroundjoin%
\definecolor{currentfill}{rgb}{0.000000,0.000000,0.000000}%
\pgfsetfillcolor{currentfill}%
\pgfsetlinewidth{0.401500pt}%
\definecolor{currentstroke}{rgb}{0.000000,0.000000,0.000000}%
\pgfsetstrokecolor{currentstroke}%
\pgfsetdash{}{0pt}%
\pgfsys@defobject{currentmarker}{\pgfqpoint{0.000000in}{-0.020833in}}{\pgfqpoint{0.000000in}{0.000000in}}{%
\pgfpathmoveto{\pgfqpoint{0.000000in}{0.000000in}}%
\pgfpathlineto{\pgfqpoint{0.000000in}{-0.020833in}}%
\pgfusepath{stroke,fill}%
}%
\begin{pgfscope}%
\pgfsys@transformshift{1.335502in}{2.414488in}%
\pgfsys@useobject{currentmarker}{}%
\end{pgfscope}%
\end{pgfscope}%
\begin{pgfscope}%
\pgfsetbuttcap%
\pgfsetroundjoin%
\definecolor{currentfill}{rgb}{0.000000,0.000000,0.000000}%
\pgfsetfillcolor{currentfill}%
\pgfsetlinewidth{0.401500pt}%
\definecolor{currentstroke}{rgb}{0.000000,0.000000,0.000000}%
\pgfsetstrokecolor{currentstroke}%
\pgfsetdash{}{0pt}%
\pgfsys@defobject{currentmarker}{\pgfqpoint{0.000000in}{0.000000in}}{\pgfqpoint{0.000000in}{0.020833in}}{%
\pgfpathmoveto{\pgfqpoint{0.000000in}{0.000000in}}%
\pgfpathlineto{\pgfqpoint{0.000000in}{0.020833in}}%
\pgfusepath{stroke,fill}%
}%
\begin{pgfscope}%
\pgfsys@transformshift{1.409063in}{0.420000in}%
\pgfsys@useobject{currentmarker}{}%
\end{pgfscope}%
\end{pgfscope}%
\begin{pgfscope}%
\pgfsetbuttcap%
\pgfsetroundjoin%
\definecolor{currentfill}{rgb}{0.000000,0.000000,0.000000}%
\pgfsetfillcolor{currentfill}%
\pgfsetlinewidth{0.401500pt}%
\definecolor{currentstroke}{rgb}{0.000000,0.000000,0.000000}%
\pgfsetstrokecolor{currentstroke}%
\pgfsetdash{}{0pt}%
\pgfsys@defobject{currentmarker}{\pgfqpoint{0.000000in}{-0.020833in}}{\pgfqpoint{0.000000in}{0.000000in}}{%
\pgfpathmoveto{\pgfqpoint{0.000000in}{0.000000in}}%
\pgfpathlineto{\pgfqpoint{0.000000in}{-0.020833in}}%
\pgfusepath{stroke,fill}%
}%
\begin{pgfscope}%
\pgfsys@transformshift{1.409063in}{2.414488in}%
\pgfsys@useobject{currentmarker}{}%
\end{pgfscope}%
\end{pgfscope}%
\begin{pgfscope}%
\pgfsetbuttcap%
\pgfsetroundjoin%
\definecolor{currentfill}{rgb}{0.000000,0.000000,0.000000}%
\pgfsetfillcolor{currentfill}%
\pgfsetlinewidth{0.401500pt}%
\definecolor{currentstroke}{rgb}{0.000000,0.000000,0.000000}%
\pgfsetstrokecolor{currentstroke}%
\pgfsetdash{}{0pt}%
\pgfsys@defobject{currentmarker}{\pgfqpoint{0.000000in}{0.000000in}}{\pgfqpoint{0.000000in}{0.020833in}}{%
\pgfpathmoveto{\pgfqpoint{0.000000in}{0.000000in}}%
\pgfpathlineto{\pgfqpoint{0.000000in}{0.020833in}}%
\pgfusepath{stroke,fill}%
}%
\begin{pgfscope}%
\pgfsys@transformshift{1.469167in}{0.420000in}%
\pgfsys@useobject{currentmarker}{}%
\end{pgfscope}%
\end{pgfscope}%
\begin{pgfscope}%
\pgfsetbuttcap%
\pgfsetroundjoin%
\definecolor{currentfill}{rgb}{0.000000,0.000000,0.000000}%
\pgfsetfillcolor{currentfill}%
\pgfsetlinewidth{0.401500pt}%
\definecolor{currentstroke}{rgb}{0.000000,0.000000,0.000000}%
\pgfsetstrokecolor{currentstroke}%
\pgfsetdash{}{0pt}%
\pgfsys@defobject{currentmarker}{\pgfqpoint{0.000000in}{-0.020833in}}{\pgfqpoint{0.000000in}{0.000000in}}{%
\pgfpathmoveto{\pgfqpoint{0.000000in}{0.000000in}}%
\pgfpathlineto{\pgfqpoint{0.000000in}{-0.020833in}}%
\pgfusepath{stroke,fill}%
}%
\begin{pgfscope}%
\pgfsys@transformshift{1.469167in}{2.414488in}%
\pgfsys@useobject{currentmarker}{}%
\end{pgfscope}%
\end{pgfscope}%
\begin{pgfscope}%
\pgfsetbuttcap%
\pgfsetroundjoin%
\definecolor{currentfill}{rgb}{0.000000,0.000000,0.000000}%
\pgfsetfillcolor{currentfill}%
\pgfsetlinewidth{0.401500pt}%
\definecolor{currentstroke}{rgb}{0.000000,0.000000,0.000000}%
\pgfsetstrokecolor{currentstroke}%
\pgfsetdash{}{0pt}%
\pgfsys@defobject{currentmarker}{\pgfqpoint{0.000000in}{0.000000in}}{\pgfqpoint{0.000000in}{0.020833in}}{%
\pgfpathmoveto{\pgfqpoint{0.000000in}{0.000000in}}%
\pgfpathlineto{\pgfqpoint{0.000000in}{0.020833in}}%
\pgfusepath{stroke,fill}%
}%
\begin{pgfscope}%
\pgfsys@transformshift{1.519984in}{0.420000in}%
\pgfsys@useobject{currentmarker}{}%
\end{pgfscope}%
\end{pgfscope}%
\begin{pgfscope}%
\pgfsetbuttcap%
\pgfsetroundjoin%
\definecolor{currentfill}{rgb}{0.000000,0.000000,0.000000}%
\pgfsetfillcolor{currentfill}%
\pgfsetlinewidth{0.401500pt}%
\definecolor{currentstroke}{rgb}{0.000000,0.000000,0.000000}%
\pgfsetstrokecolor{currentstroke}%
\pgfsetdash{}{0pt}%
\pgfsys@defobject{currentmarker}{\pgfqpoint{0.000000in}{-0.020833in}}{\pgfqpoint{0.000000in}{0.000000in}}{%
\pgfpathmoveto{\pgfqpoint{0.000000in}{0.000000in}}%
\pgfpathlineto{\pgfqpoint{0.000000in}{-0.020833in}}%
\pgfusepath{stroke,fill}%
}%
\begin{pgfscope}%
\pgfsys@transformshift{1.519984in}{2.414488in}%
\pgfsys@useobject{currentmarker}{}%
\end{pgfscope}%
\end{pgfscope}%
\begin{pgfscope}%
\pgfsetbuttcap%
\pgfsetroundjoin%
\definecolor{currentfill}{rgb}{0.000000,0.000000,0.000000}%
\pgfsetfillcolor{currentfill}%
\pgfsetlinewidth{0.401500pt}%
\definecolor{currentstroke}{rgb}{0.000000,0.000000,0.000000}%
\pgfsetstrokecolor{currentstroke}%
\pgfsetdash{}{0pt}%
\pgfsys@defobject{currentmarker}{\pgfqpoint{0.000000in}{0.000000in}}{\pgfqpoint{0.000000in}{0.020833in}}{%
\pgfpathmoveto{\pgfqpoint{0.000000in}{0.000000in}}%
\pgfpathlineto{\pgfqpoint{0.000000in}{0.020833in}}%
\pgfusepath{stroke,fill}%
}%
\begin{pgfscope}%
\pgfsys@transformshift{1.564003in}{0.420000in}%
\pgfsys@useobject{currentmarker}{}%
\end{pgfscope}%
\end{pgfscope}%
\begin{pgfscope}%
\pgfsetbuttcap%
\pgfsetroundjoin%
\definecolor{currentfill}{rgb}{0.000000,0.000000,0.000000}%
\pgfsetfillcolor{currentfill}%
\pgfsetlinewidth{0.401500pt}%
\definecolor{currentstroke}{rgb}{0.000000,0.000000,0.000000}%
\pgfsetstrokecolor{currentstroke}%
\pgfsetdash{}{0pt}%
\pgfsys@defobject{currentmarker}{\pgfqpoint{0.000000in}{-0.020833in}}{\pgfqpoint{0.000000in}{0.000000in}}{%
\pgfpathmoveto{\pgfqpoint{0.000000in}{0.000000in}}%
\pgfpathlineto{\pgfqpoint{0.000000in}{-0.020833in}}%
\pgfusepath{stroke,fill}%
}%
\begin{pgfscope}%
\pgfsys@transformshift{1.564003in}{2.414488in}%
\pgfsys@useobject{currentmarker}{}%
\end{pgfscope}%
\end{pgfscope}%
\begin{pgfscope}%
\pgfsetbuttcap%
\pgfsetroundjoin%
\definecolor{currentfill}{rgb}{0.000000,0.000000,0.000000}%
\pgfsetfillcolor{currentfill}%
\pgfsetlinewidth{0.401500pt}%
\definecolor{currentstroke}{rgb}{0.000000,0.000000,0.000000}%
\pgfsetstrokecolor{currentstroke}%
\pgfsetdash{}{0pt}%
\pgfsys@defobject{currentmarker}{\pgfqpoint{0.000000in}{0.000000in}}{\pgfqpoint{0.000000in}{0.020833in}}{%
\pgfpathmoveto{\pgfqpoint{0.000000in}{0.000000in}}%
\pgfpathlineto{\pgfqpoint{0.000000in}{0.020833in}}%
\pgfusepath{stroke,fill}%
}%
\begin{pgfscope}%
\pgfsys@transformshift{1.602831in}{0.420000in}%
\pgfsys@useobject{currentmarker}{}%
\end{pgfscope}%
\end{pgfscope}%
\begin{pgfscope}%
\pgfsetbuttcap%
\pgfsetroundjoin%
\definecolor{currentfill}{rgb}{0.000000,0.000000,0.000000}%
\pgfsetfillcolor{currentfill}%
\pgfsetlinewidth{0.401500pt}%
\definecolor{currentstroke}{rgb}{0.000000,0.000000,0.000000}%
\pgfsetstrokecolor{currentstroke}%
\pgfsetdash{}{0pt}%
\pgfsys@defobject{currentmarker}{\pgfqpoint{0.000000in}{-0.020833in}}{\pgfqpoint{0.000000in}{0.000000in}}{%
\pgfpathmoveto{\pgfqpoint{0.000000in}{0.000000in}}%
\pgfpathlineto{\pgfqpoint{0.000000in}{-0.020833in}}%
\pgfusepath{stroke,fill}%
}%
\begin{pgfscope}%
\pgfsys@transformshift{1.602831in}{2.414488in}%
\pgfsys@useobject{currentmarker}{}%
\end{pgfscope}%
\end{pgfscope}%
\begin{pgfscope}%
\pgfsetbuttcap%
\pgfsetroundjoin%
\definecolor{currentfill}{rgb}{0.000000,0.000000,0.000000}%
\pgfsetfillcolor{currentfill}%
\pgfsetlinewidth{0.401500pt}%
\definecolor{currentstroke}{rgb}{0.000000,0.000000,0.000000}%
\pgfsetstrokecolor{currentstroke}%
\pgfsetdash{}{0pt}%
\pgfsys@defobject{currentmarker}{\pgfqpoint{0.000000in}{0.000000in}}{\pgfqpoint{0.000000in}{0.020833in}}{%
\pgfpathmoveto{\pgfqpoint{0.000000in}{0.000000in}}%
\pgfpathlineto{\pgfqpoint{0.000000in}{0.020833in}}%
\pgfusepath{stroke,fill}%
}%
\begin{pgfscope}%
\pgfsys@transformshift{1.866065in}{0.420000in}%
\pgfsys@useobject{currentmarker}{}%
\end{pgfscope}%
\end{pgfscope}%
\begin{pgfscope}%
\pgfsetbuttcap%
\pgfsetroundjoin%
\definecolor{currentfill}{rgb}{0.000000,0.000000,0.000000}%
\pgfsetfillcolor{currentfill}%
\pgfsetlinewidth{0.401500pt}%
\definecolor{currentstroke}{rgb}{0.000000,0.000000,0.000000}%
\pgfsetstrokecolor{currentstroke}%
\pgfsetdash{}{0pt}%
\pgfsys@defobject{currentmarker}{\pgfqpoint{0.000000in}{-0.020833in}}{\pgfqpoint{0.000000in}{0.000000in}}{%
\pgfpathmoveto{\pgfqpoint{0.000000in}{0.000000in}}%
\pgfpathlineto{\pgfqpoint{0.000000in}{-0.020833in}}%
\pgfusepath{stroke,fill}%
}%
\begin{pgfscope}%
\pgfsys@transformshift{1.866065in}{2.414488in}%
\pgfsys@useobject{currentmarker}{}%
\end{pgfscope}%
\end{pgfscope}%
\begin{pgfscope}%
\pgfsetbuttcap%
\pgfsetroundjoin%
\definecolor{currentfill}{rgb}{0.000000,0.000000,0.000000}%
\pgfsetfillcolor{currentfill}%
\pgfsetlinewidth{0.401500pt}%
\definecolor{currentstroke}{rgb}{0.000000,0.000000,0.000000}%
\pgfsetstrokecolor{currentstroke}%
\pgfsetdash{}{0pt}%
\pgfsys@defobject{currentmarker}{\pgfqpoint{0.000000in}{0.000000in}}{\pgfqpoint{0.000000in}{0.020833in}}{%
\pgfpathmoveto{\pgfqpoint{0.000000in}{0.000000in}}%
\pgfpathlineto{\pgfqpoint{0.000000in}{0.020833in}}%
\pgfusepath{stroke,fill}%
}%
\begin{pgfscope}%
\pgfsys@transformshift{1.999729in}{0.420000in}%
\pgfsys@useobject{currentmarker}{}%
\end{pgfscope}%
\end{pgfscope}%
\begin{pgfscope}%
\pgfsetbuttcap%
\pgfsetroundjoin%
\definecolor{currentfill}{rgb}{0.000000,0.000000,0.000000}%
\pgfsetfillcolor{currentfill}%
\pgfsetlinewidth{0.401500pt}%
\definecolor{currentstroke}{rgb}{0.000000,0.000000,0.000000}%
\pgfsetstrokecolor{currentstroke}%
\pgfsetdash{}{0pt}%
\pgfsys@defobject{currentmarker}{\pgfqpoint{0.000000in}{-0.020833in}}{\pgfqpoint{0.000000in}{0.000000in}}{%
\pgfpathmoveto{\pgfqpoint{0.000000in}{0.000000in}}%
\pgfpathlineto{\pgfqpoint{0.000000in}{-0.020833in}}%
\pgfusepath{stroke,fill}%
}%
\begin{pgfscope}%
\pgfsys@transformshift{1.999729in}{2.414488in}%
\pgfsys@useobject{currentmarker}{}%
\end{pgfscope}%
\end{pgfscope}%
\begin{pgfscope}%
\pgfsetbuttcap%
\pgfsetroundjoin%
\definecolor{currentfill}{rgb}{0.000000,0.000000,0.000000}%
\pgfsetfillcolor{currentfill}%
\pgfsetlinewidth{0.401500pt}%
\definecolor{currentstroke}{rgb}{0.000000,0.000000,0.000000}%
\pgfsetstrokecolor{currentstroke}%
\pgfsetdash{}{0pt}%
\pgfsys@defobject{currentmarker}{\pgfqpoint{0.000000in}{0.000000in}}{\pgfqpoint{0.000000in}{0.020833in}}{%
\pgfpathmoveto{\pgfqpoint{0.000000in}{0.000000in}}%
\pgfpathlineto{\pgfqpoint{0.000000in}{0.020833in}}%
\pgfusepath{stroke,fill}%
}%
\begin{pgfscope}%
\pgfsys@transformshift{2.094566in}{0.420000in}%
\pgfsys@useobject{currentmarker}{}%
\end{pgfscope}%
\end{pgfscope}%
\begin{pgfscope}%
\pgfsetbuttcap%
\pgfsetroundjoin%
\definecolor{currentfill}{rgb}{0.000000,0.000000,0.000000}%
\pgfsetfillcolor{currentfill}%
\pgfsetlinewidth{0.401500pt}%
\definecolor{currentstroke}{rgb}{0.000000,0.000000,0.000000}%
\pgfsetstrokecolor{currentstroke}%
\pgfsetdash{}{0pt}%
\pgfsys@defobject{currentmarker}{\pgfqpoint{0.000000in}{-0.020833in}}{\pgfqpoint{0.000000in}{0.000000in}}{%
\pgfpathmoveto{\pgfqpoint{0.000000in}{0.000000in}}%
\pgfpathlineto{\pgfqpoint{0.000000in}{-0.020833in}}%
\pgfusepath{stroke,fill}%
}%
\begin{pgfscope}%
\pgfsys@transformshift{2.094566in}{2.414488in}%
\pgfsys@useobject{currentmarker}{}%
\end{pgfscope}%
\end{pgfscope}%
\begin{pgfscope}%
\pgfsetbuttcap%
\pgfsetroundjoin%
\definecolor{currentfill}{rgb}{0.000000,0.000000,0.000000}%
\pgfsetfillcolor{currentfill}%
\pgfsetlinewidth{0.401500pt}%
\definecolor{currentstroke}{rgb}{0.000000,0.000000,0.000000}%
\pgfsetstrokecolor{currentstroke}%
\pgfsetdash{}{0pt}%
\pgfsys@defobject{currentmarker}{\pgfqpoint{0.000000in}{0.000000in}}{\pgfqpoint{0.000000in}{0.020833in}}{%
\pgfpathmoveto{\pgfqpoint{0.000000in}{0.000000in}}%
\pgfpathlineto{\pgfqpoint{0.000000in}{0.020833in}}%
\pgfusepath{stroke,fill}%
}%
\begin{pgfscope}%
\pgfsys@transformshift{2.168126in}{0.420000in}%
\pgfsys@useobject{currentmarker}{}%
\end{pgfscope}%
\end{pgfscope}%
\begin{pgfscope}%
\pgfsetbuttcap%
\pgfsetroundjoin%
\definecolor{currentfill}{rgb}{0.000000,0.000000,0.000000}%
\pgfsetfillcolor{currentfill}%
\pgfsetlinewidth{0.401500pt}%
\definecolor{currentstroke}{rgb}{0.000000,0.000000,0.000000}%
\pgfsetstrokecolor{currentstroke}%
\pgfsetdash{}{0pt}%
\pgfsys@defobject{currentmarker}{\pgfqpoint{0.000000in}{-0.020833in}}{\pgfqpoint{0.000000in}{0.000000in}}{%
\pgfpathmoveto{\pgfqpoint{0.000000in}{0.000000in}}%
\pgfpathlineto{\pgfqpoint{0.000000in}{-0.020833in}}%
\pgfusepath{stroke,fill}%
}%
\begin{pgfscope}%
\pgfsys@transformshift{2.168126in}{2.414488in}%
\pgfsys@useobject{currentmarker}{}%
\end{pgfscope}%
\end{pgfscope}%
\begin{pgfscope}%
\pgfsetbuttcap%
\pgfsetroundjoin%
\definecolor{currentfill}{rgb}{0.000000,0.000000,0.000000}%
\pgfsetfillcolor{currentfill}%
\pgfsetlinewidth{0.401500pt}%
\definecolor{currentstroke}{rgb}{0.000000,0.000000,0.000000}%
\pgfsetstrokecolor{currentstroke}%
\pgfsetdash{}{0pt}%
\pgfsys@defobject{currentmarker}{\pgfqpoint{0.000000in}{0.000000in}}{\pgfqpoint{0.000000in}{0.020833in}}{%
\pgfpathmoveto{\pgfqpoint{0.000000in}{0.000000in}}%
\pgfpathlineto{\pgfqpoint{0.000000in}{0.020833in}}%
\pgfusepath{stroke,fill}%
}%
\begin{pgfscope}%
\pgfsys@transformshift{2.228230in}{0.420000in}%
\pgfsys@useobject{currentmarker}{}%
\end{pgfscope}%
\end{pgfscope}%
\begin{pgfscope}%
\pgfsetbuttcap%
\pgfsetroundjoin%
\definecolor{currentfill}{rgb}{0.000000,0.000000,0.000000}%
\pgfsetfillcolor{currentfill}%
\pgfsetlinewidth{0.401500pt}%
\definecolor{currentstroke}{rgb}{0.000000,0.000000,0.000000}%
\pgfsetstrokecolor{currentstroke}%
\pgfsetdash{}{0pt}%
\pgfsys@defobject{currentmarker}{\pgfqpoint{0.000000in}{-0.020833in}}{\pgfqpoint{0.000000in}{0.000000in}}{%
\pgfpathmoveto{\pgfqpoint{0.000000in}{0.000000in}}%
\pgfpathlineto{\pgfqpoint{0.000000in}{-0.020833in}}%
\pgfusepath{stroke,fill}%
}%
\begin{pgfscope}%
\pgfsys@transformshift{2.228230in}{2.414488in}%
\pgfsys@useobject{currentmarker}{}%
\end{pgfscope}%
\end{pgfscope}%
\begin{pgfscope}%
\pgfsetbuttcap%
\pgfsetroundjoin%
\definecolor{currentfill}{rgb}{0.000000,0.000000,0.000000}%
\pgfsetfillcolor{currentfill}%
\pgfsetlinewidth{0.401500pt}%
\definecolor{currentstroke}{rgb}{0.000000,0.000000,0.000000}%
\pgfsetstrokecolor{currentstroke}%
\pgfsetdash{}{0pt}%
\pgfsys@defobject{currentmarker}{\pgfqpoint{0.000000in}{0.000000in}}{\pgfqpoint{0.000000in}{0.020833in}}{%
\pgfpathmoveto{\pgfqpoint{0.000000in}{0.000000in}}%
\pgfpathlineto{\pgfqpoint{0.000000in}{0.020833in}}%
\pgfusepath{stroke,fill}%
}%
\begin{pgfscope}%
\pgfsys@transformshift{2.279047in}{0.420000in}%
\pgfsys@useobject{currentmarker}{}%
\end{pgfscope}%
\end{pgfscope}%
\begin{pgfscope}%
\pgfsetbuttcap%
\pgfsetroundjoin%
\definecolor{currentfill}{rgb}{0.000000,0.000000,0.000000}%
\pgfsetfillcolor{currentfill}%
\pgfsetlinewidth{0.401500pt}%
\definecolor{currentstroke}{rgb}{0.000000,0.000000,0.000000}%
\pgfsetstrokecolor{currentstroke}%
\pgfsetdash{}{0pt}%
\pgfsys@defobject{currentmarker}{\pgfqpoint{0.000000in}{-0.020833in}}{\pgfqpoint{0.000000in}{0.000000in}}{%
\pgfpathmoveto{\pgfqpoint{0.000000in}{0.000000in}}%
\pgfpathlineto{\pgfqpoint{0.000000in}{-0.020833in}}%
\pgfusepath{stroke,fill}%
}%
\begin{pgfscope}%
\pgfsys@transformshift{2.279047in}{2.414488in}%
\pgfsys@useobject{currentmarker}{}%
\end{pgfscope}%
\end{pgfscope}%
\begin{pgfscope}%
\pgfsetbuttcap%
\pgfsetroundjoin%
\definecolor{currentfill}{rgb}{0.000000,0.000000,0.000000}%
\pgfsetfillcolor{currentfill}%
\pgfsetlinewidth{0.401500pt}%
\definecolor{currentstroke}{rgb}{0.000000,0.000000,0.000000}%
\pgfsetstrokecolor{currentstroke}%
\pgfsetdash{}{0pt}%
\pgfsys@defobject{currentmarker}{\pgfqpoint{0.000000in}{0.000000in}}{\pgfqpoint{0.000000in}{0.020833in}}{%
\pgfpathmoveto{\pgfqpoint{0.000000in}{0.000000in}}%
\pgfpathlineto{\pgfqpoint{0.000000in}{0.020833in}}%
\pgfusepath{stroke,fill}%
}%
\begin{pgfscope}%
\pgfsys@transformshift{2.323066in}{0.420000in}%
\pgfsys@useobject{currentmarker}{}%
\end{pgfscope}%
\end{pgfscope}%
\begin{pgfscope}%
\pgfsetbuttcap%
\pgfsetroundjoin%
\definecolor{currentfill}{rgb}{0.000000,0.000000,0.000000}%
\pgfsetfillcolor{currentfill}%
\pgfsetlinewidth{0.401500pt}%
\definecolor{currentstroke}{rgb}{0.000000,0.000000,0.000000}%
\pgfsetstrokecolor{currentstroke}%
\pgfsetdash{}{0pt}%
\pgfsys@defobject{currentmarker}{\pgfqpoint{0.000000in}{-0.020833in}}{\pgfqpoint{0.000000in}{0.000000in}}{%
\pgfpathmoveto{\pgfqpoint{0.000000in}{0.000000in}}%
\pgfpathlineto{\pgfqpoint{0.000000in}{-0.020833in}}%
\pgfusepath{stroke,fill}%
}%
\begin{pgfscope}%
\pgfsys@transformshift{2.323066in}{2.414488in}%
\pgfsys@useobject{currentmarker}{}%
\end{pgfscope}%
\end{pgfscope}%
\begin{pgfscope}%
\pgfsetbuttcap%
\pgfsetroundjoin%
\definecolor{currentfill}{rgb}{0.000000,0.000000,0.000000}%
\pgfsetfillcolor{currentfill}%
\pgfsetlinewidth{0.401500pt}%
\definecolor{currentstroke}{rgb}{0.000000,0.000000,0.000000}%
\pgfsetstrokecolor{currentstroke}%
\pgfsetdash{}{0pt}%
\pgfsys@defobject{currentmarker}{\pgfqpoint{0.000000in}{0.000000in}}{\pgfqpoint{0.000000in}{0.020833in}}{%
\pgfpathmoveto{\pgfqpoint{0.000000in}{0.000000in}}%
\pgfpathlineto{\pgfqpoint{0.000000in}{0.020833in}}%
\pgfusepath{stroke,fill}%
}%
\begin{pgfscope}%
\pgfsys@transformshift{2.361894in}{0.420000in}%
\pgfsys@useobject{currentmarker}{}%
\end{pgfscope}%
\end{pgfscope}%
\begin{pgfscope}%
\pgfsetbuttcap%
\pgfsetroundjoin%
\definecolor{currentfill}{rgb}{0.000000,0.000000,0.000000}%
\pgfsetfillcolor{currentfill}%
\pgfsetlinewidth{0.401500pt}%
\definecolor{currentstroke}{rgb}{0.000000,0.000000,0.000000}%
\pgfsetstrokecolor{currentstroke}%
\pgfsetdash{}{0pt}%
\pgfsys@defobject{currentmarker}{\pgfqpoint{0.000000in}{-0.020833in}}{\pgfqpoint{0.000000in}{0.000000in}}{%
\pgfpathmoveto{\pgfqpoint{0.000000in}{0.000000in}}%
\pgfpathlineto{\pgfqpoint{0.000000in}{-0.020833in}}%
\pgfusepath{stroke,fill}%
}%
\begin{pgfscope}%
\pgfsys@transformshift{2.361894in}{2.414488in}%
\pgfsys@useobject{currentmarker}{}%
\end{pgfscope}%
\end{pgfscope}%
\begin{pgfscope}%
\pgfsetbuttcap%
\pgfsetroundjoin%
\definecolor{currentfill}{rgb}{0.000000,0.000000,0.000000}%
\pgfsetfillcolor{currentfill}%
\pgfsetlinewidth{0.401500pt}%
\definecolor{currentstroke}{rgb}{0.000000,0.000000,0.000000}%
\pgfsetstrokecolor{currentstroke}%
\pgfsetdash{}{0pt}%
\pgfsys@defobject{currentmarker}{\pgfqpoint{0.000000in}{0.000000in}}{\pgfqpoint{0.000000in}{0.020833in}}{%
\pgfpathmoveto{\pgfqpoint{0.000000in}{0.000000in}}%
\pgfpathlineto{\pgfqpoint{0.000000in}{0.020833in}}%
\pgfusepath{stroke,fill}%
}%
\begin{pgfscope}%
\pgfsys@transformshift{2.625128in}{0.420000in}%
\pgfsys@useobject{currentmarker}{}%
\end{pgfscope}%
\end{pgfscope}%
\begin{pgfscope}%
\pgfsetbuttcap%
\pgfsetroundjoin%
\definecolor{currentfill}{rgb}{0.000000,0.000000,0.000000}%
\pgfsetfillcolor{currentfill}%
\pgfsetlinewidth{0.401500pt}%
\definecolor{currentstroke}{rgb}{0.000000,0.000000,0.000000}%
\pgfsetstrokecolor{currentstroke}%
\pgfsetdash{}{0pt}%
\pgfsys@defobject{currentmarker}{\pgfqpoint{0.000000in}{-0.020833in}}{\pgfqpoint{0.000000in}{0.000000in}}{%
\pgfpathmoveto{\pgfqpoint{0.000000in}{0.000000in}}%
\pgfpathlineto{\pgfqpoint{0.000000in}{-0.020833in}}%
\pgfusepath{stroke,fill}%
}%
\begin{pgfscope}%
\pgfsys@transformshift{2.625128in}{2.414488in}%
\pgfsys@useobject{currentmarker}{}%
\end{pgfscope}%
\end{pgfscope}%
\begin{pgfscope}%
\pgfsetbuttcap%
\pgfsetroundjoin%
\definecolor{currentfill}{rgb}{0.000000,0.000000,0.000000}%
\pgfsetfillcolor{currentfill}%
\pgfsetlinewidth{0.401500pt}%
\definecolor{currentstroke}{rgb}{0.000000,0.000000,0.000000}%
\pgfsetstrokecolor{currentstroke}%
\pgfsetdash{}{0pt}%
\pgfsys@defobject{currentmarker}{\pgfqpoint{0.000000in}{0.000000in}}{\pgfqpoint{0.000000in}{0.020833in}}{%
\pgfpathmoveto{\pgfqpoint{0.000000in}{0.000000in}}%
\pgfpathlineto{\pgfqpoint{0.000000in}{0.020833in}}%
\pgfusepath{stroke,fill}%
}%
\begin{pgfscope}%
\pgfsys@transformshift{2.758792in}{0.420000in}%
\pgfsys@useobject{currentmarker}{}%
\end{pgfscope}%
\end{pgfscope}%
\begin{pgfscope}%
\pgfsetbuttcap%
\pgfsetroundjoin%
\definecolor{currentfill}{rgb}{0.000000,0.000000,0.000000}%
\pgfsetfillcolor{currentfill}%
\pgfsetlinewidth{0.401500pt}%
\definecolor{currentstroke}{rgb}{0.000000,0.000000,0.000000}%
\pgfsetstrokecolor{currentstroke}%
\pgfsetdash{}{0pt}%
\pgfsys@defobject{currentmarker}{\pgfqpoint{0.000000in}{-0.020833in}}{\pgfqpoint{0.000000in}{0.000000in}}{%
\pgfpathmoveto{\pgfqpoint{0.000000in}{0.000000in}}%
\pgfpathlineto{\pgfqpoint{0.000000in}{-0.020833in}}%
\pgfusepath{stroke,fill}%
}%
\begin{pgfscope}%
\pgfsys@transformshift{2.758792in}{2.414488in}%
\pgfsys@useobject{currentmarker}{}%
\end{pgfscope}%
\end{pgfscope}%
\begin{pgfscope}%
\pgfsetbuttcap%
\pgfsetroundjoin%
\definecolor{currentfill}{rgb}{0.000000,0.000000,0.000000}%
\pgfsetfillcolor{currentfill}%
\pgfsetlinewidth{0.401500pt}%
\definecolor{currentstroke}{rgb}{0.000000,0.000000,0.000000}%
\pgfsetstrokecolor{currentstroke}%
\pgfsetdash{}{0pt}%
\pgfsys@defobject{currentmarker}{\pgfqpoint{0.000000in}{0.000000in}}{\pgfqpoint{0.000000in}{0.020833in}}{%
\pgfpathmoveto{\pgfqpoint{0.000000in}{0.000000in}}%
\pgfpathlineto{\pgfqpoint{0.000000in}{0.020833in}}%
\pgfusepath{stroke,fill}%
}%
\begin{pgfscope}%
\pgfsys@transformshift{2.853629in}{0.420000in}%
\pgfsys@useobject{currentmarker}{}%
\end{pgfscope}%
\end{pgfscope}%
\begin{pgfscope}%
\pgfsetbuttcap%
\pgfsetroundjoin%
\definecolor{currentfill}{rgb}{0.000000,0.000000,0.000000}%
\pgfsetfillcolor{currentfill}%
\pgfsetlinewidth{0.401500pt}%
\definecolor{currentstroke}{rgb}{0.000000,0.000000,0.000000}%
\pgfsetstrokecolor{currentstroke}%
\pgfsetdash{}{0pt}%
\pgfsys@defobject{currentmarker}{\pgfqpoint{0.000000in}{-0.020833in}}{\pgfqpoint{0.000000in}{0.000000in}}{%
\pgfpathmoveto{\pgfqpoint{0.000000in}{0.000000in}}%
\pgfpathlineto{\pgfqpoint{0.000000in}{-0.020833in}}%
\pgfusepath{stroke,fill}%
}%
\begin{pgfscope}%
\pgfsys@transformshift{2.853629in}{2.414488in}%
\pgfsys@useobject{currentmarker}{}%
\end{pgfscope}%
\end{pgfscope}%
\begin{pgfscope}%
\pgfsetbuttcap%
\pgfsetroundjoin%
\definecolor{currentfill}{rgb}{0.000000,0.000000,0.000000}%
\pgfsetfillcolor{currentfill}%
\pgfsetlinewidth{0.401500pt}%
\definecolor{currentstroke}{rgb}{0.000000,0.000000,0.000000}%
\pgfsetstrokecolor{currentstroke}%
\pgfsetdash{}{0pt}%
\pgfsys@defobject{currentmarker}{\pgfqpoint{0.000000in}{0.000000in}}{\pgfqpoint{0.000000in}{0.020833in}}{%
\pgfpathmoveto{\pgfqpoint{0.000000in}{0.000000in}}%
\pgfpathlineto{\pgfqpoint{0.000000in}{0.020833in}}%
\pgfusepath{stroke,fill}%
}%
\begin{pgfscope}%
\pgfsys@transformshift{2.927190in}{0.420000in}%
\pgfsys@useobject{currentmarker}{}%
\end{pgfscope}%
\end{pgfscope}%
\begin{pgfscope}%
\pgfsetbuttcap%
\pgfsetroundjoin%
\definecolor{currentfill}{rgb}{0.000000,0.000000,0.000000}%
\pgfsetfillcolor{currentfill}%
\pgfsetlinewidth{0.401500pt}%
\definecolor{currentstroke}{rgb}{0.000000,0.000000,0.000000}%
\pgfsetstrokecolor{currentstroke}%
\pgfsetdash{}{0pt}%
\pgfsys@defobject{currentmarker}{\pgfqpoint{0.000000in}{-0.020833in}}{\pgfqpoint{0.000000in}{0.000000in}}{%
\pgfpathmoveto{\pgfqpoint{0.000000in}{0.000000in}}%
\pgfpathlineto{\pgfqpoint{0.000000in}{-0.020833in}}%
\pgfusepath{stroke,fill}%
}%
\begin{pgfscope}%
\pgfsys@transformshift{2.927190in}{2.414488in}%
\pgfsys@useobject{currentmarker}{}%
\end{pgfscope}%
\end{pgfscope}%
\begin{pgfscope}%
\pgfsetbuttcap%
\pgfsetroundjoin%
\definecolor{currentfill}{rgb}{0.000000,0.000000,0.000000}%
\pgfsetfillcolor{currentfill}%
\pgfsetlinewidth{0.401500pt}%
\definecolor{currentstroke}{rgb}{0.000000,0.000000,0.000000}%
\pgfsetstrokecolor{currentstroke}%
\pgfsetdash{}{0pt}%
\pgfsys@defobject{currentmarker}{\pgfqpoint{0.000000in}{0.000000in}}{\pgfqpoint{0.000000in}{0.020833in}}{%
\pgfpathmoveto{\pgfqpoint{0.000000in}{0.000000in}}%
\pgfpathlineto{\pgfqpoint{0.000000in}{0.020833in}}%
\pgfusepath{stroke,fill}%
}%
\begin{pgfscope}%
\pgfsys@transformshift{2.987293in}{0.420000in}%
\pgfsys@useobject{currentmarker}{}%
\end{pgfscope}%
\end{pgfscope}%
\begin{pgfscope}%
\pgfsetbuttcap%
\pgfsetroundjoin%
\definecolor{currentfill}{rgb}{0.000000,0.000000,0.000000}%
\pgfsetfillcolor{currentfill}%
\pgfsetlinewidth{0.401500pt}%
\definecolor{currentstroke}{rgb}{0.000000,0.000000,0.000000}%
\pgfsetstrokecolor{currentstroke}%
\pgfsetdash{}{0pt}%
\pgfsys@defobject{currentmarker}{\pgfqpoint{0.000000in}{-0.020833in}}{\pgfqpoint{0.000000in}{0.000000in}}{%
\pgfpathmoveto{\pgfqpoint{0.000000in}{0.000000in}}%
\pgfpathlineto{\pgfqpoint{0.000000in}{-0.020833in}}%
\pgfusepath{stroke,fill}%
}%
\begin{pgfscope}%
\pgfsys@transformshift{2.987293in}{2.414488in}%
\pgfsys@useobject{currentmarker}{}%
\end{pgfscope}%
\end{pgfscope}%
\begin{pgfscope}%
\pgfsetbuttcap%
\pgfsetroundjoin%
\definecolor{currentfill}{rgb}{0.000000,0.000000,0.000000}%
\pgfsetfillcolor{currentfill}%
\pgfsetlinewidth{0.401500pt}%
\definecolor{currentstroke}{rgb}{0.000000,0.000000,0.000000}%
\pgfsetstrokecolor{currentstroke}%
\pgfsetdash{}{0pt}%
\pgfsys@defobject{currentmarker}{\pgfqpoint{0.000000in}{0.000000in}}{\pgfqpoint{0.000000in}{0.020833in}}{%
\pgfpathmoveto{\pgfqpoint{0.000000in}{0.000000in}}%
\pgfpathlineto{\pgfqpoint{0.000000in}{0.020833in}}%
\pgfusepath{stroke,fill}%
}%
\begin{pgfscope}%
\pgfsys@transformshift{3.038110in}{0.420000in}%
\pgfsys@useobject{currentmarker}{}%
\end{pgfscope}%
\end{pgfscope}%
\begin{pgfscope}%
\pgfsetbuttcap%
\pgfsetroundjoin%
\definecolor{currentfill}{rgb}{0.000000,0.000000,0.000000}%
\pgfsetfillcolor{currentfill}%
\pgfsetlinewidth{0.401500pt}%
\definecolor{currentstroke}{rgb}{0.000000,0.000000,0.000000}%
\pgfsetstrokecolor{currentstroke}%
\pgfsetdash{}{0pt}%
\pgfsys@defobject{currentmarker}{\pgfqpoint{0.000000in}{-0.020833in}}{\pgfqpoint{0.000000in}{0.000000in}}{%
\pgfpathmoveto{\pgfqpoint{0.000000in}{0.000000in}}%
\pgfpathlineto{\pgfqpoint{0.000000in}{-0.020833in}}%
\pgfusepath{stroke,fill}%
}%
\begin{pgfscope}%
\pgfsys@transformshift{3.038110in}{2.414488in}%
\pgfsys@useobject{currentmarker}{}%
\end{pgfscope}%
\end{pgfscope}%
\begin{pgfscope}%
\definecolor{textcolor}{rgb}{0.000000,0.000000,0.000000}%
\pgfsetstrokecolor{textcolor}%
\pgfsetfillcolor{textcolor}%
\pgftext[x=1.844055in,y=0.208426in,,top]{\color{textcolor}{\rmfamily\fontsize{12.000000}{14.400000}\selectfont\catcode`\^=\active\def^{\ifmmode\sp\else\^{}\fi}\catcode`\%=\active\def%{\%}$y^{+}$}}%
\end{pgfscope}%
\begin{pgfscope}%
\pgfsetbuttcap%
\pgfsetroundjoin%
\definecolor{currentfill}{rgb}{0.000000,0.000000,0.000000}%
\pgfsetfillcolor{currentfill}%
\pgfsetlinewidth{0.501875pt}%
\definecolor{currentstroke}{rgb}{0.000000,0.000000,0.000000}%
\pgfsetstrokecolor{currentstroke}%
\pgfsetdash{}{0pt}%
\pgfsys@defobject{currentmarker}{\pgfqpoint{0.000000in}{0.000000in}}{\pgfqpoint{0.041667in}{0.000000in}}{%
\pgfpathmoveto{\pgfqpoint{0.000000in}{0.000000in}}%
\pgfpathlineto{\pgfqpoint{0.041667in}{0.000000in}}%
\pgfusepath{stroke,fill}%
}%
\begin{pgfscope}%
\pgfsys@transformshift{0.650000in}{0.496711in}%
\pgfsys@useobject{currentmarker}{}%
\end{pgfscope}%
\end{pgfscope}%
\begin{pgfscope}%
\pgfsetbuttcap%
\pgfsetroundjoin%
\definecolor{currentfill}{rgb}{0.000000,0.000000,0.000000}%
\pgfsetfillcolor{currentfill}%
\pgfsetlinewidth{0.501875pt}%
\definecolor{currentstroke}{rgb}{0.000000,0.000000,0.000000}%
\pgfsetstrokecolor{currentstroke}%
\pgfsetdash{}{0pt}%
\pgfsys@defobject{currentmarker}{\pgfqpoint{-0.041667in}{0.000000in}}{\pgfqpoint{-0.000000in}{0.000000in}}{%
\pgfpathmoveto{\pgfqpoint{-0.000000in}{0.000000in}}%
\pgfpathlineto{\pgfqpoint{-0.041667in}{0.000000in}}%
\pgfusepath{stroke,fill}%
}%
\begin{pgfscope}%
\pgfsys@transformshift{3.038110in}{0.496711in}%
\pgfsys@useobject{currentmarker}{}%
\end{pgfscope}%
\end{pgfscope}%
\begin{pgfscope}%
\definecolor{textcolor}{rgb}{0.000000,0.000000,0.000000}%
\pgfsetstrokecolor{textcolor}%
\pgfsetfillcolor{textcolor}%
\pgftext[x=0.288772in, y=0.443904in, left, base]{\color{textcolor}{\rmfamily\fontsize{11.000000}{13.200000}\selectfont\catcode`\^=\active\def^{\ifmmode\sp\else\^{}\fi}\catcode`\%=\active\def%{\%}\ensuremath{-}2.5}}%
\end{pgfscope}%
\begin{pgfscope}%
\pgfsetbuttcap%
\pgfsetroundjoin%
\definecolor{currentfill}{rgb}{0.000000,0.000000,0.000000}%
\pgfsetfillcolor{currentfill}%
\pgfsetlinewidth{0.501875pt}%
\definecolor{currentstroke}{rgb}{0.000000,0.000000,0.000000}%
\pgfsetstrokecolor{currentstroke}%
\pgfsetdash{}{0pt}%
\pgfsys@defobject{currentmarker}{\pgfqpoint{0.000000in}{0.000000in}}{\pgfqpoint{0.041667in}{0.000000in}}{%
\pgfpathmoveto{\pgfqpoint{0.000000in}{0.000000in}}%
\pgfpathlineto{\pgfqpoint{0.041667in}{0.000000in}}%
\pgfusepath{stroke,fill}%
}%
\begin{pgfscope}%
\pgfsys@transformshift{0.650000in}{0.880266in}%
\pgfsys@useobject{currentmarker}{}%
\end{pgfscope}%
\end{pgfscope}%
\begin{pgfscope}%
\pgfsetbuttcap%
\pgfsetroundjoin%
\definecolor{currentfill}{rgb}{0.000000,0.000000,0.000000}%
\pgfsetfillcolor{currentfill}%
\pgfsetlinewidth{0.501875pt}%
\definecolor{currentstroke}{rgb}{0.000000,0.000000,0.000000}%
\pgfsetstrokecolor{currentstroke}%
\pgfsetdash{}{0pt}%
\pgfsys@defobject{currentmarker}{\pgfqpoint{-0.041667in}{0.000000in}}{\pgfqpoint{-0.000000in}{0.000000in}}{%
\pgfpathmoveto{\pgfqpoint{-0.000000in}{0.000000in}}%
\pgfpathlineto{\pgfqpoint{-0.041667in}{0.000000in}}%
\pgfusepath{stroke,fill}%
}%
\begin{pgfscope}%
\pgfsys@transformshift{3.038110in}{0.880266in}%
\pgfsys@useobject{currentmarker}{}%
\end{pgfscope}%
\end{pgfscope}%
\begin{pgfscope}%
\definecolor{textcolor}{rgb}{0.000000,0.000000,0.000000}%
\pgfsetstrokecolor{textcolor}%
\pgfsetfillcolor{textcolor}%
\pgftext[x=0.407060in, y=0.827460in, left, base]{\color{textcolor}{\rmfamily\fontsize{11.000000}{13.200000}\selectfont\catcode`\^=\active\def^{\ifmmode\sp\else\^{}\fi}\catcode`\%=\active\def%{\%}0.0}}%
\end{pgfscope}%
\begin{pgfscope}%
\pgfsetbuttcap%
\pgfsetroundjoin%
\definecolor{currentfill}{rgb}{0.000000,0.000000,0.000000}%
\pgfsetfillcolor{currentfill}%
\pgfsetlinewidth{0.501875pt}%
\definecolor{currentstroke}{rgb}{0.000000,0.000000,0.000000}%
\pgfsetstrokecolor{currentstroke}%
\pgfsetdash{}{0pt}%
\pgfsys@defobject{currentmarker}{\pgfqpoint{0.000000in}{0.000000in}}{\pgfqpoint{0.041667in}{0.000000in}}{%
\pgfpathmoveto{\pgfqpoint{0.000000in}{0.000000in}}%
\pgfpathlineto{\pgfqpoint{0.041667in}{0.000000in}}%
\pgfusepath{stroke,fill}%
}%
\begin{pgfscope}%
\pgfsys@transformshift{0.650000in}{1.263822in}%
\pgfsys@useobject{currentmarker}{}%
\end{pgfscope}%
\end{pgfscope}%
\begin{pgfscope}%
\pgfsetbuttcap%
\pgfsetroundjoin%
\definecolor{currentfill}{rgb}{0.000000,0.000000,0.000000}%
\pgfsetfillcolor{currentfill}%
\pgfsetlinewidth{0.501875pt}%
\definecolor{currentstroke}{rgb}{0.000000,0.000000,0.000000}%
\pgfsetstrokecolor{currentstroke}%
\pgfsetdash{}{0pt}%
\pgfsys@defobject{currentmarker}{\pgfqpoint{-0.041667in}{0.000000in}}{\pgfqpoint{-0.000000in}{0.000000in}}{%
\pgfpathmoveto{\pgfqpoint{-0.000000in}{0.000000in}}%
\pgfpathlineto{\pgfqpoint{-0.041667in}{0.000000in}}%
\pgfusepath{stroke,fill}%
}%
\begin{pgfscope}%
\pgfsys@transformshift{3.038110in}{1.263822in}%
\pgfsys@useobject{currentmarker}{}%
\end{pgfscope}%
\end{pgfscope}%
\begin{pgfscope}%
\definecolor{textcolor}{rgb}{0.000000,0.000000,0.000000}%
\pgfsetstrokecolor{textcolor}%
\pgfsetfillcolor{textcolor}%
\pgftext[x=0.407060in, y=1.211015in, left, base]{\color{textcolor}{\rmfamily\fontsize{11.000000}{13.200000}\selectfont\catcode`\^=\active\def^{\ifmmode\sp\else\^{}\fi}\catcode`\%=\active\def%{\%}2.5}}%
\end{pgfscope}%
\begin{pgfscope}%
\pgfsetbuttcap%
\pgfsetroundjoin%
\definecolor{currentfill}{rgb}{0.000000,0.000000,0.000000}%
\pgfsetfillcolor{currentfill}%
\pgfsetlinewidth{0.501875pt}%
\definecolor{currentstroke}{rgb}{0.000000,0.000000,0.000000}%
\pgfsetstrokecolor{currentstroke}%
\pgfsetdash{}{0pt}%
\pgfsys@defobject{currentmarker}{\pgfqpoint{0.000000in}{0.000000in}}{\pgfqpoint{0.041667in}{0.000000in}}{%
\pgfpathmoveto{\pgfqpoint{0.000000in}{0.000000in}}%
\pgfpathlineto{\pgfqpoint{0.041667in}{0.000000in}}%
\pgfusepath{stroke,fill}%
}%
\begin{pgfscope}%
\pgfsys@transformshift{0.650000in}{1.647377in}%
\pgfsys@useobject{currentmarker}{}%
\end{pgfscope}%
\end{pgfscope}%
\begin{pgfscope}%
\pgfsetbuttcap%
\pgfsetroundjoin%
\definecolor{currentfill}{rgb}{0.000000,0.000000,0.000000}%
\pgfsetfillcolor{currentfill}%
\pgfsetlinewidth{0.501875pt}%
\definecolor{currentstroke}{rgb}{0.000000,0.000000,0.000000}%
\pgfsetstrokecolor{currentstroke}%
\pgfsetdash{}{0pt}%
\pgfsys@defobject{currentmarker}{\pgfqpoint{-0.041667in}{0.000000in}}{\pgfqpoint{-0.000000in}{0.000000in}}{%
\pgfpathmoveto{\pgfqpoint{-0.000000in}{0.000000in}}%
\pgfpathlineto{\pgfqpoint{-0.041667in}{0.000000in}}%
\pgfusepath{stroke,fill}%
}%
\begin{pgfscope}%
\pgfsys@transformshift{3.038110in}{1.647377in}%
\pgfsys@useobject{currentmarker}{}%
\end{pgfscope}%
\end{pgfscope}%
\begin{pgfscope}%
\definecolor{textcolor}{rgb}{0.000000,0.000000,0.000000}%
\pgfsetstrokecolor{textcolor}%
\pgfsetfillcolor{textcolor}%
\pgftext[x=0.407060in, y=1.594571in, left, base]{\color{textcolor}{\rmfamily\fontsize{11.000000}{13.200000}\selectfont\catcode`\^=\active\def^{\ifmmode\sp\else\^{}\fi}\catcode`\%=\active\def%{\%}5.0}}%
\end{pgfscope}%
\begin{pgfscope}%
\pgfsetbuttcap%
\pgfsetroundjoin%
\definecolor{currentfill}{rgb}{0.000000,0.000000,0.000000}%
\pgfsetfillcolor{currentfill}%
\pgfsetlinewidth{0.501875pt}%
\definecolor{currentstroke}{rgb}{0.000000,0.000000,0.000000}%
\pgfsetstrokecolor{currentstroke}%
\pgfsetdash{}{0pt}%
\pgfsys@defobject{currentmarker}{\pgfqpoint{0.000000in}{0.000000in}}{\pgfqpoint{0.041667in}{0.000000in}}{%
\pgfpathmoveto{\pgfqpoint{0.000000in}{0.000000in}}%
\pgfpathlineto{\pgfqpoint{0.041667in}{0.000000in}}%
\pgfusepath{stroke,fill}%
}%
\begin{pgfscope}%
\pgfsys@transformshift{0.650000in}{2.030933in}%
\pgfsys@useobject{currentmarker}{}%
\end{pgfscope}%
\end{pgfscope}%
\begin{pgfscope}%
\pgfsetbuttcap%
\pgfsetroundjoin%
\definecolor{currentfill}{rgb}{0.000000,0.000000,0.000000}%
\pgfsetfillcolor{currentfill}%
\pgfsetlinewidth{0.501875pt}%
\definecolor{currentstroke}{rgb}{0.000000,0.000000,0.000000}%
\pgfsetstrokecolor{currentstroke}%
\pgfsetdash{}{0pt}%
\pgfsys@defobject{currentmarker}{\pgfqpoint{-0.041667in}{0.000000in}}{\pgfqpoint{-0.000000in}{0.000000in}}{%
\pgfpathmoveto{\pgfqpoint{-0.000000in}{0.000000in}}%
\pgfpathlineto{\pgfqpoint{-0.041667in}{0.000000in}}%
\pgfusepath{stroke,fill}%
}%
\begin{pgfscope}%
\pgfsys@transformshift{3.038110in}{2.030933in}%
\pgfsys@useobject{currentmarker}{}%
\end{pgfscope}%
\end{pgfscope}%
\begin{pgfscope}%
\definecolor{textcolor}{rgb}{0.000000,0.000000,0.000000}%
\pgfsetstrokecolor{textcolor}%
\pgfsetfillcolor{textcolor}%
\pgftext[x=0.407060in, y=1.978126in, left, base]{\color{textcolor}{\rmfamily\fontsize{11.000000}{13.200000}\selectfont\catcode`\^=\active\def^{\ifmmode\sp\else\^{}\fi}\catcode`\%=\active\def%{\%}7.5}}%
\end{pgfscope}%
\begin{pgfscope}%
\pgfsetbuttcap%
\pgfsetroundjoin%
\definecolor{currentfill}{rgb}{0.000000,0.000000,0.000000}%
\pgfsetfillcolor{currentfill}%
\pgfsetlinewidth{0.501875pt}%
\definecolor{currentstroke}{rgb}{0.000000,0.000000,0.000000}%
\pgfsetstrokecolor{currentstroke}%
\pgfsetdash{}{0pt}%
\pgfsys@defobject{currentmarker}{\pgfqpoint{0.000000in}{0.000000in}}{\pgfqpoint{0.041667in}{0.000000in}}{%
\pgfpathmoveto{\pgfqpoint{0.000000in}{0.000000in}}%
\pgfpathlineto{\pgfqpoint{0.041667in}{0.000000in}}%
\pgfusepath{stroke,fill}%
}%
\begin{pgfscope}%
\pgfsys@transformshift{0.650000in}{2.414488in}%
\pgfsys@useobject{currentmarker}{}%
\end{pgfscope}%
\end{pgfscope}%
\begin{pgfscope}%
\pgfsetbuttcap%
\pgfsetroundjoin%
\definecolor{currentfill}{rgb}{0.000000,0.000000,0.000000}%
\pgfsetfillcolor{currentfill}%
\pgfsetlinewidth{0.501875pt}%
\definecolor{currentstroke}{rgb}{0.000000,0.000000,0.000000}%
\pgfsetstrokecolor{currentstroke}%
\pgfsetdash{}{0pt}%
\pgfsys@defobject{currentmarker}{\pgfqpoint{-0.041667in}{0.000000in}}{\pgfqpoint{-0.000000in}{0.000000in}}{%
\pgfpathmoveto{\pgfqpoint{-0.000000in}{0.000000in}}%
\pgfpathlineto{\pgfqpoint{-0.041667in}{0.000000in}}%
\pgfusepath{stroke,fill}%
}%
\begin{pgfscope}%
\pgfsys@transformshift{3.038110in}{2.414488in}%
\pgfsys@useobject{currentmarker}{}%
\end{pgfscope}%
\end{pgfscope}%
\begin{pgfscope}%
\definecolor{textcolor}{rgb}{0.000000,0.000000,0.000000}%
\pgfsetstrokecolor{textcolor}%
\pgfsetfillcolor{textcolor}%
\pgftext[x=0.331018in, y=2.361681in, left, base]{\color{textcolor}{\rmfamily\fontsize{11.000000}{13.200000}\selectfont\catcode`\^=\active\def^{\ifmmode\sp\else\^{}\fi}\catcode`\%=\active\def%{\%}10.0}}%
\end{pgfscope}%
\begin{pgfscope}%
\pgfsetbuttcap%
\pgfsetroundjoin%
\definecolor{currentfill}{rgb}{0.000000,0.000000,0.000000}%
\pgfsetfillcolor{currentfill}%
\pgfsetlinewidth{0.401500pt}%
\definecolor{currentstroke}{rgb}{0.000000,0.000000,0.000000}%
\pgfsetstrokecolor{currentstroke}%
\pgfsetdash{}{0pt}%
\pgfsys@defobject{currentmarker}{\pgfqpoint{0.000000in}{0.000000in}}{\pgfqpoint{0.020833in}{0.000000in}}{%
\pgfpathmoveto{\pgfqpoint{0.000000in}{0.000000in}}%
\pgfpathlineto{\pgfqpoint{0.020833in}{0.000000in}}%
\pgfusepath{stroke,fill}%
}%
\begin{pgfscope}%
\pgfsys@transformshift{0.650000in}{0.420000in}%
\pgfsys@useobject{currentmarker}{}%
\end{pgfscope}%
\end{pgfscope}%
\begin{pgfscope}%
\pgfsetbuttcap%
\pgfsetroundjoin%
\definecolor{currentfill}{rgb}{0.000000,0.000000,0.000000}%
\pgfsetfillcolor{currentfill}%
\pgfsetlinewidth{0.401500pt}%
\definecolor{currentstroke}{rgb}{0.000000,0.000000,0.000000}%
\pgfsetstrokecolor{currentstroke}%
\pgfsetdash{}{0pt}%
\pgfsys@defobject{currentmarker}{\pgfqpoint{-0.020833in}{0.000000in}}{\pgfqpoint{-0.000000in}{0.000000in}}{%
\pgfpathmoveto{\pgfqpoint{-0.000000in}{0.000000in}}%
\pgfpathlineto{\pgfqpoint{-0.020833in}{0.000000in}}%
\pgfusepath{stroke,fill}%
}%
\begin{pgfscope}%
\pgfsys@transformshift{3.038110in}{0.420000in}%
\pgfsys@useobject{currentmarker}{}%
\end{pgfscope}%
\end{pgfscope}%
\begin{pgfscope}%
\pgfsetbuttcap%
\pgfsetroundjoin%
\definecolor{currentfill}{rgb}{0.000000,0.000000,0.000000}%
\pgfsetfillcolor{currentfill}%
\pgfsetlinewidth{0.401500pt}%
\definecolor{currentstroke}{rgb}{0.000000,0.000000,0.000000}%
\pgfsetstrokecolor{currentstroke}%
\pgfsetdash{}{0pt}%
\pgfsys@defobject{currentmarker}{\pgfqpoint{0.000000in}{0.000000in}}{\pgfqpoint{0.020833in}{0.000000in}}{%
\pgfpathmoveto{\pgfqpoint{0.000000in}{0.000000in}}%
\pgfpathlineto{\pgfqpoint{0.020833in}{0.000000in}}%
\pgfusepath{stroke,fill}%
}%
\begin{pgfscope}%
\pgfsys@transformshift{0.650000in}{0.573422in}%
\pgfsys@useobject{currentmarker}{}%
\end{pgfscope}%
\end{pgfscope}%
\begin{pgfscope}%
\pgfsetbuttcap%
\pgfsetroundjoin%
\definecolor{currentfill}{rgb}{0.000000,0.000000,0.000000}%
\pgfsetfillcolor{currentfill}%
\pgfsetlinewidth{0.401500pt}%
\definecolor{currentstroke}{rgb}{0.000000,0.000000,0.000000}%
\pgfsetstrokecolor{currentstroke}%
\pgfsetdash{}{0pt}%
\pgfsys@defobject{currentmarker}{\pgfqpoint{-0.020833in}{0.000000in}}{\pgfqpoint{-0.000000in}{0.000000in}}{%
\pgfpathmoveto{\pgfqpoint{-0.000000in}{0.000000in}}%
\pgfpathlineto{\pgfqpoint{-0.020833in}{0.000000in}}%
\pgfusepath{stroke,fill}%
}%
\begin{pgfscope}%
\pgfsys@transformshift{3.038110in}{0.573422in}%
\pgfsys@useobject{currentmarker}{}%
\end{pgfscope}%
\end{pgfscope}%
\begin{pgfscope}%
\pgfsetbuttcap%
\pgfsetroundjoin%
\definecolor{currentfill}{rgb}{0.000000,0.000000,0.000000}%
\pgfsetfillcolor{currentfill}%
\pgfsetlinewidth{0.401500pt}%
\definecolor{currentstroke}{rgb}{0.000000,0.000000,0.000000}%
\pgfsetstrokecolor{currentstroke}%
\pgfsetdash{}{0pt}%
\pgfsys@defobject{currentmarker}{\pgfqpoint{0.000000in}{0.000000in}}{\pgfqpoint{0.020833in}{0.000000in}}{%
\pgfpathmoveto{\pgfqpoint{0.000000in}{0.000000in}}%
\pgfpathlineto{\pgfqpoint{0.020833in}{0.000000in}}%
\pgfusepath{stroke,fill}%
}%
\begin{pgfscope}%
\pgfsys@transformshift{0.650000in}{0.650133in}%
\pgfsys@useobject{currentmarker}{}%
\end{pgfscope}%
\end{pgfscope}%
\begin{pgfscope}%
\pgfsetbuttcap%
\pgfsetroundjoin%
\definecolor{currentfill}{rgb}{0.000000,0.000000,0.000000}%
\pgfsetfillcolor{currentfill}%
\pgfsetlinewidth{0.401500pt}%
\definecolor{currentstroke}{rgb}{0.000000,0.000000,0.000000}%
\pgfsetstrokecolor{currentstroke}%
\pgfsetdash{}{0pt}%
\pgfsys@defobject{currentmarker}{\pgfqpoint{-0.020833in}{0.000000in}}{\pgfqpoint{-0.000000in}{0.000000in}}{%
\pgfpathmoveto{\pgfqpoint{-0.000000in}{0.000000in}}%
\pgfpathlineto{\pgfqpoint{-0.020833in}{0.000000in}}%
\pgfusepath{stroke,fill}%
}%
\begin{pgfscope}%
\pgfsys@transformshift{3.038110in}{0.650133in}%
\pgfsys@useobject{currentmarker}{}%
\end{pgfscope}%
\end{pgfscope}%
\begin{pgfscope}%
\pgfsetbuttcap%
\pgfsetroundjoin%
\definecolor{currentfill}{rgb}{0.000000,0.000000,0.000000}%
\pgfsetfillcolor{currentfill}%
\pgfsetlinewidth{0.401500pt}%
\definecolor{currentstroke}{rgb}{0.000000,0.000000,0.000000}%
\pgfsetstrokecolor{currentstroke}%
\pgfsetdash{}{0pt}%
\pgfsys@defobject{currentmarker}{\pgfqpoint{0.000000in}{0.000000in}}{\pgfqpoint{0.020833in}{0.000000in}}{%
\pgfpathmoveto{\pgfqpoint{0.000000in}{0.000000in}}%
\pgfpathlineto{\pgfqpoint{0.020833in}{0.000000in}}%
\pgfusepath{stroke,fill}%
}%
\begin{pgfscope}%
\pgfsys@transformshift{0.650000in}{0.726844in}%
\pgfsys@useobject{currentmarker}{}%
\end{pgfscope}%
\end{pgfscope}%
\begin{pgfscope}%
\pgfsetbuttcap%
\pgfsetroundjoin%
\definecolor{currentfill}{rgb}{0.000000,0.000000,0.000000}%
\pgfsetfillcolor{currentfill}%
\pgfsetlinewidth{0.401500pt}%
\definecolor{currentstroke}{rgb}{0.000000,0.000000,0.000000}%
\pgfsetstrokecolor{currentstroke}%
\pgfsetdash{}{0pt}%
\pgfsys@defobject{currentmarker}{\pgfqpoint{-0.020833in}{0.000000in}}{\pgfqpoint{-0.000000in}{0.000000in}}{%
\pgfpathmoveto{\pgfqpoint{-0.000000in}{0.000000in}}%
\pgfpathlineto{\pgfqpoint{-0.020833in}{0.000000in}}%
\pgfusepath{stroke,fill}%
}%
\begin{pgfscope}%
\pgfsys@transformshift{3.038110in}{0.726844in}%
\pgfsys@useobject{currentmarker}{}%
\end{pgfscope}%
\end{pgfscope}%
\begin{pgfscope}%
\pgfsetbuttcap%
\pgfsetroundjoin%
\definecolor{currentfill}{rgb}{0.000000,0.000000,0.000000}%
\pgfsetfillcolor{currentfill}%
\pgfsetlinewidth{0.401500pt}%
\definecolor{currentstroke}{rgb}{0.000000,0.000000,0.000000}%
\pgfsetstrokecolor{currentstroke}%
\pgfsetdash{}{0pt}%
\pgfsys@defobject{currentmarker}{\pgfqpoint{0.000000in}{0.000000in}}{\pgfqpoint{0.020833in}{0.000000in}}{%
\pgfpathmoveto{\pgfqpoint{0.000000in}{0.000000in}}%
\pgfpathlineto{\pgfqpoint{0.020833in}{0.000000in}}%
\pgfusepath{stroke,fill}%
}%
\begin{pgfscope}%
\pgfsys@transformshift{0.650000in}{0.803555in}%
\pgfsys@useobject{currentmarker}{}%
\end{pgfscope}%
\end{pgfscope}%
\begin{pgfscope}%
\pgfsetbuttcap%
\pgfsetroundjoin%
\definecolor{currentfill}{rgb}{0.000000,0.000000,0.000000}%
\pgfsetfillcolor{currentfill}%
\pgfsetlinewidth{0.401500pt}%
\definecolor{currentstroke}{rgb}{0.000000,0.000000,0.000000}%
\pgfsetstrokecolor{currentstroke}%
\pgfsetdash{}{0pt}%
\pgfsys@defobject{currentmarker}{\pgfqpoint{-0.020833in}{0.000000in}}{\pgfqpoint{-0.000000in}{0.000000in}}{%
\pgfpathmoveto{\pgfqpoint{-0.000000in}{0.000000in}}%
\pgfpathlineto{\pgfqpoint{-0.020833in}{0.000000in}}%
\pgfusepath{stroke,fill}%
}%
\begin{pgfscope}%
\pgfsys@transformshift{3.038110in}{0.803555in}%
\pgfsys@useobject{currentmarker}{}%
\end{pgfscope}%
\end{pgfscope}%
\begin{pgfscope}%
\pgfsetbuttcap%
\pgfsetroundjoin%
\definecolor{currentfill}{rgb}{0.000000,0.000000,0.000000}%
\pgfsetfillcolor{currentfill}%
\pgfsetlinewidth{0.401500pt}%
\definecolor{currentstroke}{rgb}{0.000000,0.000000,0.000000}%
\pgfsetstrokecolor{currentstroke}%
\pgfsetdash{}{0pt}%
\pgfsys@defobject{currentmarker}{\pgfqpoint{0.000000in}{0.000000in}}{\pgfqpoint{0.020833in}{0.000000in}}{%
\pgfpathmoveto{\pgfqpoint{0.000000in}{0.000000in}}%
\pgfpathlineto{\pgfqpoint{0.020833in}{0.000000in}}%
\pgfusepath{stroke,fill}%
}%
\begin{pgfscope}%
\pgfsys@transformshift{0.650000in}{0.956978in}%
\pgfsys@useobject{currentmarker}{}%
\end{pgfscope}%
\end{pgfscope}%
\begin{pgfscope}%
\pgfsetbuttcap%
\pgfsetroundjoin%
\definecolor{currentfill}{rgb}{0.000000,0.000000,0.000000}%
\pgfsetfillcolor{currentfill}%
\pgfsetlinewidth{0.401500pt}%
\definecolor{currentstroke}{rgb}{0.000000,0.000000,0.000000}%
\pgfsetstrokecolor{currentstroke}%
\pgfsetdash{}{0pt}%
\pgfsys@defobject{currentmarker}{\pgfqpoint{-0.020833in}{0.000000in}}{\pgfqpoint{-0.000000in}{0.000000in}}{%
\pgfpathmoveto{\pgfqpoint{-0.000000in}{0.000000in}}%
\pgfpathlineto{\pgfqpoint{-0.020833in}{0.000000in}}%
\pgfusepath{stroke,fill}%
}%
\begin{pgfscope}%
\pgfsys@transformshift{3.038110in}{0.956978in}%
\pgfsys@useobject{currentmarker}{}%
\end{pgfscope}%
\end{pgfscope}%
\begin{pgfscope}%
\pgfsetbuttcap%
\pgfsetroundjoin%
\definecolor{currentfill}{rgb}{0.000000,0.000000,0.000000}%
\pgfsetfillcolor{currentfill}%
\pgfsetlinewidth{0.401500pt}%
\definecolor{currentstroke}{rgb}{0.000000,0.000000,0.000000}%
\pgfsetstrokecolor{currentstroke}%
\pgfsetdash{}{0pt}%
\pgfsys@defobject{currentmarker}{\pgfqpoint{0.000000in}{0.000000in}}{\pgfqpoint{0.020833in}{0.000000in}}{%
\pgfpathmoveto{\pgfqpoint{0.000000in}{0.000000in}}%
\pgfpathlineto{\pgfqpoint{0.020833in}{0.000000in}}%
\pgfusepath{stroke,fill}%
}%
\begin{pgfscope}%
\pgfsys@transformshift{0.650000in}{1.033689in}%
\pgfsys@useobject{currentmarker}{}%
\end{pgfscope}%
\end{pgfscope}%
\begin{pgfscope}%
\pgfsetbuttcap%
\pgfsetroundjoin%
\definecolor{currentfill}{rgb}{0.000000,0.000000,0.000000}%
\pgfsetfillcolor{currentfill}%
\pgfsetlinewidth{0.401500pt}%
\definecolor{currentstroke}{rgb}{0.000000,0.000000,0.000000}%
\pgfsetstrokecolor{currentstroke}%
\pgfsetdash{}{0pt}%
\pgfsys@defobject{currentmarker}{\pgfqpoint{-0.020833in}{0.000000in}}{\pgfqpoint{-0.000000in}{0.000000in}}{%
\pgfpathmoveto{\pgfqpoint{-0.000000in}{0.000000in}}%
\pgfpathlineto{\pgfqpoint{-0.020833in}{0.000000in}}%
\pgfusepath{stroke,fill}%
}%
\begin{pgfscope}%
\pgfsys@transformshift{3.038110in}{1.033689in}%
\pgfsys@useobject{currentmarker}{}%
\end{pgfscope}%
\end{pgfscope}%
\begin{pgfscope}%
\pgfsetbuttcap%
\pgfsetroundjoin%
\definecolor{currentfill}{rgb}{0.000000,0.000000,0.000000}%
\pgfsetfillcolor{currentfill}%
\pgfsetlinewidth{0.401500pt}%
\definecolor{currentstroke}{rgb}{0.000000,0.000000,0.000000}%
\pgfsetstrokecolor{currentstroke}%
\pgfsetdash{}{0pt}%
\pgfsys@defobject{currentmarker}{\pgfqpoint{0.000000in}{0.000000in}}{\pgfqpoint{0.020833in}{0.000000in}}{%
\pgfpathmoveto{\pgfqpoint{0.000000in}{0.000000in}}%
\pgfpathlineto{\pgfqpoint{0.020833in}{0.000000in}}%
\pgfusepath{stroke,fill}%
}%
\begin{pgfscope}%
\pgfsys@transformshift{0.650000in}{1.110400in}%
\pgfsys@useobject{currentmarker}{}%
\end{pgfscope}%
\end{pgfscope}%
\begin{pgfscope}%
\pgfsetbuttcap%
\pgfsetroundjoin%
\definecolor{currentfill}{rgb}{0.000000,0.000000,0.000000}%
\pgfsetfillcolor{currentfill}%
\pgfsetlinewidth{0.401500pt}%
\definecolor{currentstroke}{rgb}{0.000000,0.000000,0.000000}%
\pgfsetstrokecolor{currentstroke}%
\pgfsetdash{}{0pt}%
\pgfsys@defobject{currentmarker}{\pgfqpoint{-0.020833in}{0.000000in}}{\pgfqpoint{-0.000000in}{0.000000in}}{%
\pgfpathmoveto{\pgfqpoint{-0.000000in}{0.000000in}}%
\pgfpathlineto{\pgfqpoint{-0.020833in}{0.000000in}}%
\pgfusepath{stroke,fill}%
}%
\begin{pgfscope}%
\pgfsys@transformshift{3.038110in}{1.110400in}%
\pgfsys@useobject{currentmarker}{}%
\end{pgfscope}%
\end{pgfscope}%
\begin{pgfscope}%
\pgfsetbuttcap%
\pgfsetroundjoin%
\definecolor{currentfill}{rgb}{0.000000,0.000000,0.000000}%
\pgfsetfillcolor{currentfill}%
\pgfsetlinewidth{0.401500pt}%
\definecolor{currentstroke}{rgb}{0.000000,0.000000,0.000000}%
\pgfsetstrokecolor{currentstroke}%
\pgfsetdash{}{0pt}%
\pgfsys@defobject{currentmarker}{\pgfqpoint{0.000000in}{0.000000in}}{\pgfqpoint{0.020833in}{0.000000in}}{%
\pgfpathmoveto{\pgfqpoint{0.000000in}{0.000000in}}%
\pgfpathlineto{\pgfqpoint{0.020833in}{0.000000in}}%
\pgfusepath{stroke,fill}%
}%
\begin{pgfscope}%
\pgfsys@transformshift{0.650000in}{1.187111in}%
\pgfsys@useobject{currentmarker}{}%
\end{pgfscope}%
\end{pgfscope}%
\begin{pgfscope}%
\pgfsetbuttcap%
\pgfsetroundjoin%
\definecolor{currentfill}{rgb}{0.000000,0.000000,0.000000}%
\pgfsetfillcolor{currentfill}%
\pgfsetlinewidth{0.401500pt}%
\definecolor{currentstroke}{rgb}{0.000000,0.000000,0.000000}%
\pgfsetstrokecolor{currentstroke}%
\pgfsetdash{}{0pt}%
\pgfsys@defobject{currentmarker}{\pgfqpoint{-0.020833in}{0.000000in}}{\pgfqpoint{-0.000000in}{0.000000in}}{%
\pgfpathmoveto{\pgfqpoint{-0.000000in}{0.000000in}}%
\pgfpathlineto{\pgfqpoint{-0.020833in}{0.000000in}}%
\pgfusepath{stroke,fill}%
}%
\begin{pgfscope}%
\pgfsys@transformshift{3.038110in}{1.187111in}%
\pgfsys@useobject{currentmarker}{}%
\end{pgfscope}%
\end{pgfscope}%
\begin{pgfscope}%
\pgfsetbuttcap%
\pgfsetroundjoin%
\definecolor{currentfill}{rgb}{0.000000,0.000000,0.000000}%
\pgfsetfillcolor{currentfill}%
\pgfsetlinewidth{0.401500pt}%
\definecolor{currentstroke}{rgb}{0.000000,0.000000,0.000000}%
\pgfsetstrokecolor{currentstroke}%
\pgfsetdash{}{0pt}%
\pgfsys@defobject{currentmarker}{\pgfqpoint{0.000000in}{0.000000in}}{\pgfqpoint{0.020833in}{0.000000in}}{%
\pgfpathmoveto{\pgfqpoint{0.000000in}{0.000000in}}%
\pgfpathlineto{\pgfqpoint{0.020833in}{0.000000in}}%
\pgfusepath{stroke,fill}%
}%
\begin{pgfscope}%
\pgfsys@transformshift{0.650000in}{1.340533in}%
\pgfsys@useobject{currentmarker}{}%
\end{pgfscope}%
\end{pgfscope}%
\begin{pgfscope}%
\pgfsetbuttcap%
\pgfsetroundjoin%
\definecolor{currentfill}{rgb}{0.000000,0.000000,0.000000}%
\pgfsetfillcolor{currentfill}%
\pgfsetlinewidth{0.401500pt}%
\definecolor{currentstroke}{rgb}{0.000000,0.000000,0.000000}%
\pgfsetstrokecolor{currentstroke}%
\pgfsetdash{}{0pt}%
\pgfsys@defobject{currentmarker}{\pgfqpoint{-0.020833in}{0.000000in}}{\pgfqpoint{-0.000000in}{0.000000in}}{%
\pgfpathmoveto{\pgfqpoint{-0.000000in}{0.000000in}}%
\pgfpathlineto{\pgfqpoint{-0.020833in}{0.000000in}}%
\pgfusepath{stroke,fill}%
}%
\begin{pgfscope}%
\pgfsys@transformshift{3.038110in}{1.340533in}%
\pgfsys@useobject{currentmarker}{}%
\end{pgfscope}%
\end{pgfscope}%
\begin{pgfscope}%
\pgfsetbuttcap%
\pgfsetroundjoin%
\definecolor{currentfill}{rgb}{0.000000,0.000000,0.000000}%
\pgfsetfillcolor{currentfill}%
\pgfsetlinewidth{0.401500pt}%
\definecolor{currentstroke}{rgb}{0.000000,0.000000,0.000000}%
\pgfsetstrokecolor{currentstroke}%
\pgfsetdash{}{0pt}%
\pgfsys@defobject{currentmarker}{\pgfqpoint{0.000000in}{0.000000in}}{\pgfqpoint{0.020833in}{0.000000in}}{%
\pgfpathmoveto{\pgfqpoint{0.000000in}{0.000000in}}%
\pgfpathlineto{\pgfqpoint{0.020833in}{0.000000in}}%
\pgfusepath{stroke,fill}%
}%
\begin{pgfscope}%
\pgfsys@transformshift{0.650000in}{1.417244in}%
\pgfsys@useobject{currentmarker}{}%
\end{pgfscope}%
\end{pgfscope}%
\begin{pgfscope}%
\pgfsetbuttcap%
\pgfsetroundjoin%
\definecolor{currentfill}{rgb}{0.000000,0.000000,0.000000}%
\pgfsetfillcolor{currentfill}%
\pgfsetlinewidth{0.401500pt}%
\definecolor{currentstroke}{rgb}{0.000000,0.000000,0.000000}%
\pgfsetstrokecolor{currentstroke}%
\pgfsetdash{}{0pt}%
\pgfsys@defobject{currentmarker}{\pgfqpoint{-0.020833in}{0.000000in}}{\pgfqpoint{-0.000000in}{0.000000in}}{%
\pgfpathmoveto{\pgfqpoint{-0.000000in}{0.000000in}}%
\pgfpathlineto{\pgfqpoint{-0.020833in}{0.000000in}}%
\pgfusepath{stroke,fill}%
}%
\begin{pgfscope}%
\pgfsys@transformshift{3.038110in}{1.417244in}%
\pgfsys@useobject{currentmarker}{}%
\end{pgfscope}%
\end{pgfscope}%
\begin{pgfscope}%
\pgfsetbuttcap%
\pgfsetroundjoin%
\definecolor{currentfill}{rgb}{0.000000,0.000000,0.000000}%
\pgfsetfillcolor{currentfill}%
\pgfsetlinewidth{0.401500pt}%
\definecolor{currentstroke}{rgb}{0.000000,0.000000,0.000000}%
\pgfsetstrokecolor{currentstroke}%
\pgfsetdash{}{0pt}%
\pgfsys@defobject{currentmarker}{\pgfqpoint{0.000000in}{0.000000in}}{\pgfqpoint{0.020833in}{0.000000in}}{%
\pgfpathmoveto{\pgfqpoint{0.000000in}{0.000000in}}%
\pgfpathlineto{\pgfqpoint{0.020833in}{0.000000in}}%
\pgfusepath{stroke,fill}%
}%
\begin{pgfscope}%
\pgfsys@transformshift{0.650000in}{1.493955in}%
\pgfsys@useobject{currentmarker}{}%
\end{pgfscope}%
\end{pgfscope}%
\begin{pgfscope}%
\pgfsetbuttcap%
\pgfsetroundjoin%
\definecolor{currentfill}{rgb}{0.000000,0.000000,0.000000}%
\pgfsetfillcolor{currentfill}%
\pgfsetlinewidth{0.401500pt}%
\definecolor{currentstroke}{rgb}{0.000000,0.000000,0.000000}%
\pgfsetstrokecolor{currentstroke}%
\pgfsetdash{}{0pt}%
\pgfsys@defobject{currentmarker}{\pgfqpoint{-0.020833in}{0.000000in}}{\pgfqpoint{-0.000000in}{0.000000in}}{%
\pgfpathmoveto{\pgfqpoint{-0.000000in}{0.000000in}}%
\pgfpathlineto{\pgfqpoint{-0.020833in}{0.000000in}}%
\pgfusepath{stroke,fill}%
}%
\begin{pgfscope}%
\pgfsys@transformshift{3.038110in}{1.493955in}%
\pgfsys@useobject{currentmarker}{}%
\end{pgfscope}%
\end{pgfscope}%
\begin{pgfscope}%
\pgfsetbuttcap%
\pgfsetroundjoin%
\definecolor{currentfill}{rgb}{0.000000,0.000000,0.000000}%
\pgfsetfillcolor{currentfill}%
\pgfsetlinewidth{0.401500pt}%
\definecolor{currentstroke}{rgb}{0.000000,0.000000,0.000000}%
\pgfsetstrokecolor{currentstroke}%
\pgfsetdash{}{0pt}%
\pgfsys@defobject{currentmarker}{\pgfqpoint{0.000000in}{0.000000in}}{\pgfqpoint{0.020833in}{0.000000in}}{%
\pgfpathmoveto{\pgfqpoint{0.000000in}{0.000000in}}%
\pgfpathlineto{\pgfqpoint{0.020833in}{0.000000in}}%
\pgfusepath{stroke,fill}%
}%
\begin{pgfscope}%
\pgfsys@transformshift{0.650000in}{1.570666in}%
\pgfsys@useobject{currentmarker}{}%
\end{pgfscope}%
\end{pgfscope}%
\begin{pgfscope}%
\pgfsetbuttcap%
\pgfsetroundjoin%
\definecolor{currentfill}{rgb}{0.000000,0.000000,0.000000}%
\pgfsetfillcolor{currentfill}%
\pgfsetlinewidth{0.401500pt}%
\definecolor{currentstroke}{rgb}{0.000000,0.000000,0.000000}%
\pgfsetstrokecolor{currentstroke}%
\pgfsetdash{}{0pt}%
\pgfsys@defobject{currentmarker}{\pgfqpoint{-0.020833in}{0.000000in}}{\pgfqpoint{-0.000000in}{0.000000in}}{%
\pgfpathmoveto{\pgfqpoint{-0.000000in}{0.000000in}}%
\pgfpathlineto{\pgfqpoint{-0.020833in}{0.000000in}}%
\pgfusepath{stroke,fill}%
}%
\begin{pgfscope}%
\pgfsys@transformshift{3.038110in}{1.570666in}%
\pgfsys@useobject{currentmarker}{}%
\end{pgfscope}%
\end{pgfscope}%
\begin{pgfscope}%
\pgfsetbuttcap%
\pgfsetroundjoin%
\definecolor{currentfill}{rgb}{0.000000,0.000000,0.000000}%
\pgfsetfillcolor{currentfill}%
\pgfsetlinewidth{0.401500pt}%
\definecolor{currentstroke}{rgb}{0.000000,0.000000,0.000000}%
\pgfsetstrokecolor{currentstroke}%
\pgfsetdash{}{0pt}%
\pgfsys@defobject{currentmarker}{\pgfqpoint{0.000000in}{0.000000in}}{\pgfqpoint{0.020833in}{0.000000in}}{%
\pgfpathmoveto{\pgfqpoint{0.000000in}{0.000000in}}%
\pgfpathlineto{\pgfqpoint{0.020833in}{0.000000in}}%
\pgfusepath{stroke,fill}%
}%
\begin{pgfscope}%
\pgfsys@transformshift{0.650000in}{1.724088in}%
\pgfsys@useobject{currentmarker}{}%
\end{pgfscope}%
\end{pgfscope}%
\begin{pgfscope}%
\pgfsetbuttcap%
\pgfsetroundjoin%
\definecolor{currentfill}{rgb}{0.000000,0.000000,0.000000}%
\pgfsetfillcolor{currentfill}%
\pgfsetlinewidth{0.401500pt}%
\definecolor{currentstroke}{rgb}{0.000000,0.000000,0.000000}%
\pgfsetstrokecolor{currentstroke}%
\pgfsetdash{}{0pt}%
\pgfsys@defobject{currentmarker}{\pgfqpoint{-0.020833in}{0.000000in}}{\pgfqpoint{-0.000000in}{0.000000in}}{%
\pgfpathmoveto{\pgfqpoint{-0.000000in}{0.000000in}}%
\pgfpathlineto{\pgfqpoint{-0.020833in}{0.000000in}}%
\pgfusepath{stroke,fill}%
}%
\begin{pgfscope}%
\pgfsys@transformshift{3.038110in}{1.724088in}%
\pgfsys@useobject{currentmarker}{}%
\end{pgfscope}%
\end{pgfscope}%
\begin{pgfscope}%
\pgfsetbuttcap%
\pgfsetroundjoin%
\definecolor{currentfill}{rgb}{0.000000,0.000000,0.000000}%
\pgfsetfillcolor{currentfill}%
\pgfsetlinewidth{0.401500pt}%
\definecolor{currentstroke}{rgb}{0.000000,0.000000,0.000000}%
\pgfsetstrokecolor{currentstroke}%
\pgfsetdash{}{0pt}%
\pgfsys@defobject{currentmarker}{\pgfqpoint{0.000000in}{0.000000in}}{\pgfqpoint{0.020833in}{0.000000in}}{%
\pgfpathmoveto{\pgfqpoint{0.000000in}{0.000000in}}%
\pgfpathlineto{\pgfqpoint{0.020833in}{0.000000in}}%
\pgfusepath{stroke,fill}%
}%
\begin{pgfscope}%
\pgfsys@transformshift{0.650000in}{1.800799in}%
\pgfsys@useobject{currentmarker}{}%
\end{pgfscope}%
\end{pgfscope}%
\begin{pgfscope}%
\pgfsetbuttcap%
\pgfsetroundjoin%
\definecolor{currentfill}{rgb}{0.000000,0.000000,0.000000}%
\pgfsetfillcolor{currentfill}%
\pgfsetlinewidth{0.401500pt}%
\definecolor{currentstroke}{rgb}{0.000000,0.000000,0.000000}%
\pgfsetstrokecolor{currentstroke}%
\pgfsetdash{}{0pt}%
\pgfsys@defobject{currentmarker}{\pgfqpoint{-0.020833in}{0.000000in}}{\pgfqpoint{-0.000000in}{0.000000in}}{%
\pgfpathmoveto{\pgfqpoint{-0.000000in}{0.000000in}}%
\pgfpathlineto{\pgfqpoint{-0.020833in}{0.000000in}}%
\pgfusepath{stroke,fill}%
}%
\begin{pgfscope}%
\pgfsys@transformshift{3.038110in}{1.800799in}%
\pgfsys@useobject{currentmarker}{}%
\end{pgfscope}%
\end{pgfscope}%
\begin{pgfscope}%
\pgfsetbuttcap%
\pgfsetroundjoin%
\definecolor{currentfill}{rgb}{0.000000,0.000000,0.000000}%
\pgfsetfillcolor{currentfill}%
\pgfsetlinewidth{0.401500pt}%
\definecolor{currentstroke}{rgb}{0.000000,0.000000,0.000000}%
\pgfsetstrokecolor{currentstroke}%
\pgfsetdash{}{0pt}%
\pgfsys@defobject{currentmarker}{\pgfqpoint{0.000000in}{0.000000in}}{\pgfqpoint{0.020833in}{0.000000in}}{%
\pgfpathmoveto{\pgfqpoint{0.000000in}{0.000000in}}%
\pgfpathlineto{\pgfqpoint{0.020833in}{0.000000in}}%
\pgfusepath{stroke,fill}%
}%
\begin{pgfscope}%
\pgfsys@transformshift{0.650000in}{1.877511in}%
\pgfsys@useobject{currentmarker}{}%
\end{pgfscope}%
\end{pgfscope}%
\begin{pgfscope}%
\pgfsetbuttcap%
\pgfsetroundjoin%
\definecolor{currentfill}{rgb}{0.000000,0.000000,0.000000}%
\pgfsetfillcolor{currentfill}%
\pgfsetlinewidth{0.401500pt}%
\definecolor{currentstroke}{rgb}{0.000000,0.000000,0.000000}%
\pgfsetstrokecolor{currentstroke}%
\pgfsetdash{}{0pt}%
\pgfsys@defobject{currentmarker}{\pgfqpoint{-0.020833in}{0.000000in}}{\pgfqpoint{-0.000000in}{0.000000in}}{%
\pgfpathmoveto{\pgfqpoint{-0.000000in}{0.000000in}}%
\pgfpathlineto{\pgfqpoint{-0.020833in}{0.000000in}}%
\pgfusepath{stroke,fill}%
}%
\begin{pgfscope}%
\pgfsys@transformshift{3.038110in}{1.877511in}%
\pgfsys@useobject{currentmarker}{}%
\end{pgfscope}%
\end{pgfscope}%
\begin{pgfscope}%
\pgfsetbuttcap%
\pgfsetroundjoin%
\definecolor{currentfill}{rgb}{0.000000,0.000000,0.000000}%
\pgfsetfillcolor{currentfill}%
\pgfsetlinewidth{0.401500pt}%
\definecolor{currentstroke}{rgb}{0.000000,0.000000,0.000000}%
\pgfsetstrokecolor{currentstroke}%
\pgfsetdash{}{0pt}%
\pgfsys@defobject{currentmarker}{\pgfqpoint{0.000000in}{0.000000in}}{\pgfqpoint{0.020833in}{0.000000in}}{%
\pgfpathmoveto{\pgfqpoint{0.000000in}{0.000000in}}%
\pgfpathlineto{\pgfqpoint{0.020833in}{0.000000in}}%
\pgfusepath{stroke,fill}%
}%
\begin{pgfscope}%
\pgfsys@transformshift{0.650000in}{1.954222in}%
\pgfsys@useobject{currentmarker}{}%
\end{pgfscope}%
\end{pgfscope}%
\begin{pgfscope}%
\pgfsetbuttcap%
\pgfsetroundjoin%
\definecolor{currentfill}{rgb}{0.000000,0.000000,0.000000}%
\pgfsetfillcolor{currentfill}%
\pgfsetlinewidth{0.401500pt}%
\definecolor{currentstroke}{rgb}{0.000000,0.000000,0.000000}%
\pgfsetstrokecolor{currentstroke}%
\pgfsetdash{}{0pt}%
\pgfsys@defobject{currentmarker}{\pgfqpoint{-0.020833in}{0.000000in}}{\pgfqpoint{-0.000000in}{0.000000in}}{%
\pgfpathmoveto{\pgfqpoint{-0.000000in}{0.000000in}}%
\pgfpathlineto{\pgfqpoint{-0.020833in}{0.000000in}}%
\pgfusepath{stroke,fill}%
}%
\begin{pgfscope}%
\pgfsys@transformshift{3.038110in}{1.954222in}%
\pgfsys@useobject{currentmarker}{}%
\end{pgfscope}%
\end{pgfscope}%
\begin{pgfscope}%
\pgfsetbuttcap%
\pgfsetroundjoin%
\definecolor{currentfill}{rgb}{0.000000,0.000000,0.000000}%
\pgfsetfillcolor{currentfill}%
\pgfsetlinewidth{0.401500pt}%
\definecolor{currentstroke}{rgb}{0.000000,0.000000,0.000000}%
\pgfsetstrokecolor{currentstroke}%
\pgfsetdash{}{0pt}%
\pgfsys@defobject{currentmarker}{\pgfqpoint{0.000000in}{0.000000in}}{\pgfqpoint{0.020833in}{0.000000in}}{%
\pgfpathmoveto{\pgfqpoint{0.000000in}{0.000000in}}%
\pgfpathlineto{\pgfqpoint{0.020833in}{0.000000in}}%
\pgfusepath{stroke,fill}%
}%
\begin{pgfscope}%
\pgfsys@transformshift{0.650000in}{2.107644in}%
\pgfsys@useobject{currentmarker}{}%
\end{pgfscope}%
\end{pgfscope}%
\begin{pgfscope}%
\pgfsetbuttcap%
\pgfsetroundjoin%
\definecolor{currentfill}{rgb}{0.000000,0.000000,0.000000}%
\pgfsetfillcolor{currentfill}%
\pgfsetlinewidth{0.401500pt}%
\definecolor{currentstroke}{rgb}{0.000000,0.000000,0.000000}%
\pgfsetstrokecolor{currentstroke}%
\pgfsetdash{}{0pt}%
\pgfsys@defobject{currentmarker}{\pgfqpoint{-0.020833in}{0.000000in}}{\pgfqpoint{-0.000000in}{0.000000in}}{%
\pgfpathmoveto{\pgfqpoint{-0.000000in}{0.000000in}}%
\pgfpathlineto{\pgfqpoint{-0.020833in}{0.000000in}}%
\pgfusepath{stroke,fill}%
}%
\begin{pgfscope}%
\pgfsys@transformshift{3.038110in}{2.107644in}%
\pgfsys@useobject{currentmarker}{}%
\end{pgfscope}%
\end{pgfscope}%
\begin{pgfscope}%
\pgfsetbuttcap%
\pgfsetroundjoin%
\definecolor{currentfill}{rgb}{0.000000,0.000000,0.000000}%
\pgfsetfillcolor{currentfill}%
\pgfsetlinewidth{0.401500pt}%
\definecolor{currentstroke}{rgb}{0.000000,0.000000,0.000000}%
\pgfsetstrokecolor{currentstroke}%
\pgfsetdash{}{0pt}%
\pgfsys@defobject{currentmarker}{\pgfqpoint{0.000000in}{0.000000in}}{\pgfqpoint{0.020833in}{0.000000in}}{%
\pgfpathmoveto{\pgfqpoint{0.000000in}{0.000000in}}%
\pgfpathlineto{\pgfqpoint{0.020833in}{0.000000in}}%
\pgfusepath{stroke,fill}%
}%
\begin{pgfscope}%
\pgfsys@transformshift{0.650000in}{2.184355in}%
\pgfsys@useobject{currentmarker}{}%
\end{pgfscope}%
\end{pgfscope}%
\begin{pgfscope}%
\pgfsetbuttcap%
\pgfsetroundjoin%
\definecolor{currentfill}{rgb}{0.000000,0.000000,0.000000}%
\pgfsetfillcolor{currentfill}%
\pgfsetlinewidth{0.401500pt}%
\definecolor{currentstroke}{rgb}{0.000000,0.000000,0.000000}%
\pgfsetstrokecolor{currentstroke}%
\pgfsetdash{}{0pt}%
\pgfsys@defobject{currentmarker}{\pgfqpoint{-0.020833in}{0.000000in}}{\pgfqpoint{-0.000000in}{0.000000in}}{%
\pgfpathmoveto{\pgfqpoint{-0.000000in}{0.000000in}}%
\pgfpathlineto{\pgfqpoint{-0.020833in}{0.000000in}}%
\pgfusepath{stroke,fill}%
}%
\begin{pgfscope}%
\pgfsys@transformshift{3.038110in}{2.184355in}%
\pgfsys@useobject{currentmarker}{}%
\end{pgfscope}%
\end{pgfscope}%
\begin{pgfscope}%
\pgfsetbuttcap%
\pgfsetroundjoin%
\definecolor{currentfill}{rgb}{0.000000,0.000000,0.000000}%
\pgfsetfillcolor{currentfill}%
\pgfsetlinewidth{0.401500pt}%
\definecolor{currentstroke}{rgb}{0.000000,0.000000,0.000000}%
\pgfsetstrokecolor{currentstroke}%
\pgfsetdash{}{0pt}%
\pgfsys@defobject{currentmarker}{\pgfqpoint{0.000000in}{0.000000in}}{\pgfqpoint{0.020833in}{0.000000in}}{%
\pgfpathmoveto{\pgfqpoint{0.000000in}{0.000000in}}%
\pgfpathlineto{\pgfqpoint{0.020833in}{0.000000in}}%
\pgfusepath{stroke,fill}%
}%
\begin{pgfscope}%
\pgfsys@transformshift{0.650000in}{2.261066in}%
\pgfsys@useobject{currentmarker}{}%
\end{pgfscope}%
\end{pgfscope}%
\begin{pgfscope}%
\pgfsetbuttcap%
\pgfsetroundjoin%
\definecolor{currentfill}{rgb}{0.000000,0.000000,0.000000}%
\pgfsetfillcolor{currentfill}%
\pgfsetlinewidth{0.401500pt}%
\definecolor{currentstroke}{rgb}{0.000000,0.000000,0.000000}%
\pgfsetstrokecolor{currentstroke}%
\pgfsetdash{}{0pt}%
\pgfsys@defobject{currentmarker}{\pgfqpoint{-0.020833in}{0.000000in}}{\pgfqpoint{-0.000000in}{0.000000in}}{%
\pgfpathmoveto{\pgfqpoint{-0.000000in}{0.000000in}}%
\pgfpathlineto{\pgfqpoint{-0.020833in}{0.000000in}}%
\pgfusepath{stroke,fill}%
}%
\begin{pgfscope}%
\pgfsys@transformshift{3.038110in}{2.261066in}%
\pgfsys@useobject{currentmarker}{}%
\end{pgfscope}%
\end{pgfscope}%
\begin{pgfscope}%
\pgfsetbuttcap%
\pgfsetroundjoin%
\definecolor{currentfill}{rgb}{0.000000,0.000000,0.000000}%
\pgfsetfillcolor{currentfill}%
\pgfsetlinewidth{0.401500pt}%
\definecolor{currentstroke}{rgb}{0.000000,0.000000,0.000000}%
\pgfsetstrokecolor{currentstroke}%
\pgfsetdash{}{0pt}%
\pgfsys@defobject{currentmarker}{\pgfqpoint{0.000000in}{0.000000in}}{\pgfqpoint{0.020833in}{0.000000in}}{%
\pgfpathmoveto{\pgfqpoint{0.000000in}{0.000000in}}%
\pgfpathlineto{\pgfqpoint{0.020833in}{0.000000in}}%
\pgfusepath{stroke,fill}%
}%
\begin{pgfscope}%
\pgfsys@transformshift{0.650000in}{2.337777in}%
\pgfsys@useobject{currentmarker}{}%
\end{pgfscope}%
\end{pgfscope}%
\begin{pgfscope}%
\pgfsetbuttcap%
\pgfsetroundjoin%
\definecolor{currentfill}{rgb}{0.000000,0.000000,0.000000}%
\pgfsetfillcolor{currentfill}%
\pgfsetlinewidth{0.401500pt}%
\definecolor{currentstroke}{rgb}{0.000000,0.000000,0.000000}%
\pgfsetstrokecolor{currentstroke}%
\pgfsetdash{}{0pt}%
\pgfsys@defobject{currentmarker}{\pgfqpoint{-0.020833in}{0.000000in}}{\pgfqpoint{-0.000000in}{0.000000in}}{%
\pgfpathmoveto{\pgfqpoint{-0.000000in}{0.000000in}}%
\pgfpathlineto{\pgfqpoint{-0.020833in}{0.000000in}}%
\pgfusepath{stroke,fill}%
}%
\begin{pgfscope}%
\pgfsys@transformshift{3.038110in}{2.337777in}%
\pgfsys@useobject{currentmarker}{}%
\end{pgfscope}%
\end{pgfscope}%
\begin{pgfscope}%
\definecolor{textcolor}{rgb}{0.000000,0.000000,0.000000}%
\pgfsetstrokecolor{textcolor}%
\pgfsetfillcolor{textcolor}%
\pgftext[x=0.260995in,y=1.417244in,,bottom,rotate=90.000000]{\color{textcolor}{\rmfamily\fontsize{12.000000}{14.400000}\selectfont\catcode`\^=\active\def^{\ifmmode\sp\else\^{}\fi}\catcode`\%=\active\def%{\%}$\langle u_i u_j \rangle^+$}}%
\end{pgfscope}%
\begin{pgfscope}%
\pgfpathrectangle{\pgfqpoint{0.650000in}{0.420000in}}{\pgfqpoint{2.388110in}{1.994488in}}%
\pgfusepath{clip}%
\pgfsetrectcap%
\pgfsetroundjoin%
\pgfsetlinewidth{0.602250pt}%
\definecolor{currentstroke}{rgb}{0.000000,0.000000,0.000000}%
\pgfsetstrokecolor{currentstroke}%
\pgfsetdash{}{0pt}%
\pgfpathmoveto{\pgfqpoint{0.645000in}{0.887152in}}%
\pgfpathlineto{\pgfqpoint{0.742708in}{0.890883in}}%
\pgfpathlineto{\pgfqpoint{0.889821in}{0.906045in}}%
\pgfpathlineto{\pgfqpoint{1.010016in}{0.933343in}}%
\pgfpathlineto{\pgfqpoint{1.111637in}{0.977495in}}%
\pgfpathlineto{\pgfqpoint{1.199660in}{1.042920in}}%
\pgfpathlineto{\pgfqpoint{1.277299in}{1.132285in}}%
\pgfpathlineto{\pgfqpoint{1.346745in}{1.244735in}}%
\pgfpathlineto{\pgfqpoint{1.409562in}{1.374722in}}%
\pgfpathlineto{\pgfqpoint{1.466906in}{1.512424in}}%
\pgfpathlineto{\pgfqpoint{1.519654in}{1.645910in}}%
\pgfpathlineto{\pgfqpoint{1.568486in}{1.764084in}}%
\pgfpathlineto{\pgfqpoint{1.613944in}{1.859048in}}%
\pgfpathlineto{\pgfqpoint{1.656463in}{1.927050in}}%
\pgfpathlineto{\pgfqpoint{1.696400in}{1.968088in}}%
\pgfpathlineto{\pgfqpoint{1.734049in}{1.984787in}}%
\pgfpathlineto{\pgfqpoint{1.769658in}{1.981173in}}%
\pgfpathlineto{\pgfqpoint{1.803436in}{1.961698in}}%
\pgfpathlineto{\pgfqpoint{1.835562in}{1.930608in}}%
\pgfpathlineto{\pgfqpoint{1.866188in}{1.891631in}}%
\pgfpathlineto{\pgfqpoint{1.895450in}{1.847856in}}%
\pgfpathlineto{\pgfqpoint{1.923461in}{1.801718in}}%
\pgfpathlineto{\pgfqpoint{1.976131in}{1.709214in}}%
\pgfpathlineto{\pgfqpoint{2.024880in}{1.623414in}}%
\pgfpathlineto{\pgfqpoint{2.047957in}{1.584437in}}%
\pgfpathlineto{\pgfqpoint{2.070247in}{1.548422in}}%
\pgfpathlineto{\pgfqpoint{2.091803in}{1.515394in}}%
\pgfpathlineto{\pgfqpoint{2.112670in}{1.485223in}}%
\pgfpathlineto{\pgfqpoint{2.132891in}{1.457728in}}%
\pgfpathlineto{\pgfqpoint{2.152504in}{1.432736in}}%
\pgfpathlineto{\pgfqpoint{2.171545in}{1.410074in}}%
\pgfpathlineto{\pgfqpoint{2.190045in}{1.389501in}}%
\pgfpathlineto{\pgfqpoint{2.208034in}{1.370760in}}%
\pgfpathlineto{\pgfqpoint{2.225539in}{1.353600in}}%
\pgfpathlineto{\pgfqpoint{2.242585in}{1.337819in}}%
\pgfpathlineto{\pgfqpoint{2.275392in}{1.309811in}}%
\pgfpathlineto{\pgfqpoint{2.306620in}{1.285550in}}%
\pgfpathlineto{\pgfqpoint{2.350808in}{1.253322in}}%
\pgfpathlineto{\pgfqpoint{2.405385in}{1.213712in}}%
\pgfpathlineto{\pgfqpoint{2.443483in}{1.184142in}}%
\pgfpathlineto{\pgfqpoint{2.467681in}{1.164138in}}%
\pgfpathlineto{\pgfqpoint{2.490998in}{1.143749in}}%
\pgfpathlineto{\pgfqpoint{2.513495in}{1.122749in}}%
\pgfpathlineto{\pgfqpoint{2.535225in}{1.100830in}}%
\pgfpathlineto{\pgfqpoint{2.556238in}{1.077919in}}%
\pgfpathlineto{\pgfqpoint{2.576578in}{1.053848in}}%
\pgfpathlineto{\pgfqpoint{2.596285in}{1.028797in}}%
\pgfpathlineto{\pgfqpoint{2.624738in}{0.990229in}}%
\pgfpathlineto{\pgfqpoint{2.651960in}{0.953499in}}%
\pgfpathlineto{\pgfqpoint{2.669474in}{0.932317in}}%
\pgfpathlineto{\pgfqpoint{2.678050in}{0.923244in}}%
\pgfpathlineto{\pgfqpoint{2.686509in}{0.915280in}}%
\pgfpathlineto{\pgfqpoint{2.694856in}{0.908425in}}%
\pgfpathlineto{\pgfqpoint{2.703092in}{0.902656in}}%
\pgfpathlineto{\pgfqpoint{2.711220in}{0.897910in}}%
\pgfpathlineto{\pgfqpoint{2.719243in}{0.894088in}}%
\pgfpathlineto{\pgfqpoint{2.727164in}{0.891066in}}%
\pgfpathlineto{\pgfqpoint{2.742707in}{0.886906in}}%
\pgfpathlineto{\pgfqpoint{2.757868in}{0.884486in}}%
\pgfpathlineto{\pgfqpoint{2.779929in}{0.882600in}}%
\pgfpathlineto{\pgfqpoint{2.815015in}{0.881278in}}%
\pgfpathlineto{\pgfqpoint{2.873454in}{0.880570in}}%
\pgfpathlineto{\pgfqpoint{3.025161in}{0.880290in}}%
\pgfpathlineto{\pgfqpoint{3.043110in}{0.880284in}}%
\pgfpathlineto{\pgfqpoint{3.043110in}{0.880284in}}%
\pgfusepath{stroke}%
\end{pgfscope}%
\begin{pgfscope}%
\pgfpathrectangle{\pgfqpoint{0.650000in}{0.420000in}}{\pgfqpoint{2.388110in}{1.994488in}}%
\pgfusepath{clip}%
\pgfsetrectcap%
\pgfsetroundjoin%
\pgfsetlinewidth{0.602250pt}%
\definecolor{currentstroke}{rgb}{0.000000,0.000000,0.000000}%
\pgfsetstrokecolor{currentstroke}%
\pgfsetdash{}{0pt}%
\pgfpathmoveto{\pgfqpoint{0.645000in}{0.880268in}}%
\pgfpathlineto{\pgfqpoint{1.199660in}{0.880772in}}%
\pgfpathlineto{\pgfqpoint{1.346745in}{0.882328in}}%
\pgfpathlineto{\pgfqpoint{1.409562in}{0.883859in}}%
\pgfpathlineto{\pgfqpoint{1.466906in}{0.886073in}}%
\pgfpathlineto{\pgfqpoint{1.519654in}{0.889090in}}%
\pgfpathlineto{\pgfqpoint{1.568486in}{0.892997in}}%
\pgfpathlineto{\pgfqpoint{1.613944in}{0.897849in}}%
\pgfpathlineto{\pgfqpoint{1.656463in}{0.903658in}}%
\pgfpathlineto{\pgfqpoint{1.696400in}{0.910390in}}%
\pgfpathlineto{\pgfqpoint{1.734049in}{0.917961in}}%
\pgfpathlineto{\pgfqpoint{1.769658in}{0.926247in}}%
\pgfpathlineto{\pgfqpoint{1.803436in}{0.935086in}}%
\pgfpathlineto{\pgfqpoint{1.835562in}{0.944296in}}%
\pgfpathlineto{\pgfqpoint{1.895450in}{0.963061in}}%
\pgfpathlineto{\pgfqpoint{2.000959in}{0.997311in}}%
\pgfpathlineto{\pgfqpoint{2.047957in}{1.010901in}}%
\pgfpathlineto{\pgfqpoint{2.070247in}{1.016583in}}%
\pgfpathlineto{\pgfqpoint{2.091803in}{1.021498in}}%
\pgfpathlineto{\pgfqpoint{2.112670in}{1.025647in}}%
\pgfpathlineto{\pgfqpoint{2.132891in}{1.029048in}}%
\pgfpathlineto{\pgfqpoint{2.152504in}{1.031732in}}%
\pgfpathlineto{\pgfqpoint{2.171545in}{1.033742in}}%
\pgfpathlineto{\pgfqpoint{2.208034in}{1.035960in}}%
\pgfpathlineto{\pgfqpoint{2.242585in}{1.036093in}}%
\pgfpathlineto{\pgfqpoint{2.275392in}{1.034453in}}%
\pgfpathlineto{\pgfqpoint{2.306620in}{1.031398in}}%
\pgfpathlineto{\pgfqpoint{2.336412in}{1.027244in}}%
\pgfpathlineto{\pgfqpoint{2.364891in}{1.022216in}}%
\pgfpathlineto{\pgfqpoint{2.392168in}{1.016457in}}%
\pgfpathlineto{\pgfqpoint{2.418337in}{1.010111in}}%
\pgfpathlineto{\pgfqpoint{2.443483in}{1.003212in}}%
\pgfpathlineto{\pgfqpoint{2.479445in}{0.991921in}}%
\pgfpathlineto{\pgfqpoint{2.513495in}{0.979775in}}%
\pgfpathlineto{\pgfqpoint{2.545818in}{0.967006in}}%
\pgfpathlineto{\pgfqpoint{2.586508in}{0.949611in}}%
\pgfpathlineto{\pgfqpoint{2.678050in}{0.909493in}}%
\pgfpathlineto{\pgfqpoint{2.703092in}{0.900127in}}%
\pgfpathlineto{\pgfqpoint{2.719243in}{0.895023in}}%
\pgfpathlineto{\pgfqpoint{2.734984in}{0.890912in}}%
\pgfpathlineto{\pgfqpoint{2.750334in}{0.887766in}}%
\pgfpathlineto{\pgfqpoint{2.772663in}{0.884583in}}%
\pgfpathlineto{\pgfqpoint{2.794207in}{0.882756in}}%
\pgfpathlineto{\pgfqpoint{2.821795in}{0.881537in}}%
\pgfpathlineto{\pgfqpoint{2.873454in}{0.880692in}}%
\pgfpathlineto{\pgfqpoint{2.990820in}{0.880321in}}%
\pgfpathlineto{\pgfqpoint{3.043110in}{0.880291in}}%
\pgfpathlineto{\pgfqpoint{3.043110in}{0.880291in}}%
\pgfusepath{stroke}%
\end{pgfscope}%
\begin{pgfscope}%
\pgfpathrectangle{\pgfqpoint{0.650000in}{0.420000in}}{\pgfqpoint{2.388110in}{1.994488in}}%
\pgfusepath{clip}%
\pgfsetrectcap%
\pgfsetroundjoin%
\pgfsetlinewidth{0.602250pt}%
\definecolor{currentstroke}{rgb}{0.000000,0.000000,0.000000}%
\pgfsetstrokecolor{currentstroke}%
\pgfsetdash{}{0pt}%
\pgfpathmoveto{\pgfqpoint{0.645000in}{0.882674in}}%
\pgfpathlineto{\pgfqpoint{0.742708in}{0.883916in}}%
\pgfpathlineto{\pgfqpoint{0.889821in}{0.888421in}}%
\pgfpathlineto{\pgfqpoint{1.010016in}{0.895478in}}%
\pgfpathlineto{\pgfqpoint{1.111637in}{0.905217in}}%
\pgfpathlineto{\pgfqpoint{1.199660in}{0.917413in}}%
\pgfpathlineto{\pgfqpoint{1.277299in}{0.931557in}}%
\pgfpathlineto{\pgfqpoint{1.346745in}{0.947000in}}%
\pgfpathlineto{\pgfqpoint{1.409562in}{0.963115in}}%
\pgfpathlineto{\pgfqpoint{1.466906in}{0.979416in}}%
\pgfpathlineto{\pgfqpoint{1.519654in}{0.995585in}}%
\pgfpathlineto{\pgfqpoint{1.568486in}{1.011432in}}%
\pgfpathlineto{\pgfqpoint{1.656463in}{1.041583in}}%
\pgfpathlineto{\pgfqpoint{1.734049in}{1.068465in}}%
\pgfpathlineto{\pgfqpoint{1.769658in}{1.080186in}}%
\pgfpathlineto{\pgfqpoint{1.803436in}{1.090589in}}%
\pgfpathlineto{\pgfqpoint{1.835562in}{1.099639in}}%
\pgfpathlineto{\pgfqpoint{1.866188in}{1.107366in}}%
\pgfpathlineto{\pgfqpoint{1.895450in}{1.113845in}}%
\pgfpathlineto{\pgfqpoint{1.923461in}{1.119187in}}%
\pgfpathlineto{\pgfqpoint{1.950325in}{1.123524in}}%
\pgfpathlineto{\pgfqpoint{1.976131in}{1.126987in}}%
\pgfpathlineto{\pgfqpoint{2.024880in}{1.131739in}}%
\pgfpathlineto{\pgfqpoint{2.070247in}{1.134161in}}%
\pgfpathlineto{\pgfqpoint{2.112670in}{1.134686in}}%
\pgfpathlineto{\pgfqpoint{2.152504in}{1.133573in}}%
\pgfpathlineto{\pgfqpoint{2.190045in}{1.130917in}}%
\pgfpathlineto{\pgfqpoint{2.225539in}{1.126649in}}%
\pgfpathlineto{\pgfqpoint{2.259196in}{1.121158in}}%
\pgfpathlineto{\pgfqpoint{2.291194in}{1.114442in}}%
\pgfpathlineto{\pgfqpoint{2.321687in}{1.106666in}}%
\pgfpathlineto{\pgfqpoint{2.350808in}{1.098048in}}%
\pgfpathlineto{\pgfqpoint{2.378674in}{1.088644in}}%
\pgfpathlineto{\pgfqpoint{2.405385in}{1.078258in}}%
\pgfpathlineto{\pgfqpoint{2.431033in}{1.066771in}}%
\pgfpathlineto{\pgfqpoint{2.455696in}{1.054396in}}%
\pgfpathlineto{\pgfqpoint{2.479445in}{1.041195in}}%
\pgfpathlineto{\pgfqpoint{2.502345in}{1.027017in}}%
\pgfpathlineto{\pgfqpoint{2.524453in}{1.012048in}}%
\pgfpathlineto{\pgfqpoint{2.556238in}{0.988608in}}%
\pgfpathlineto{\pgfqpoint{2.624738in}{0.935785in}}%
\pgfpathlineto{\pgfqpoint{2.643017in}{0.923161in}}%
\pgfpathlineto{\pgfqpoint{2.660778in}{0.912234in}}%
\pgfpathlineto{\pgfqpoint{2.678050in}{0.903220in}}%
\pgfpathlineto{\pgfqpoint{2.694856in}{0.896152in}}%
\pgfpathlineto{\pgfqpoint{2.711220in}{0.890882in}}%
\pgfpathlineto{\pgfqpoint{2.727164in}{0.887116in}}%
\pgfpathlineto{\pgfqpoint{2.742707in}{0.884555in}}%
\pgfpathlineto{\pgfqpoint{2.765310in}{0.882346in}}%
\pgfpathlineto{\pgfqpoint{2.794207in}{0.881098in}}%
\pgfpathlineto{\pgfqpoint{2.841691in}{0.880506in}}%
\pgfpathlineto{\pgfqpoint{2.995856in}{0.880280in}}%
\pgfpathlineto{\pgfqpoint{3.043110in}{0.880273in}}%
\pgfpathlineto{\pgfqpoint{3.043110in}{0.880273in}}%
\pgfusepath{stroke}%
\end{pgfscope}%
\begin{pgfscope}%
\pgfpathrectangle{\pgfqpoint{0.650000in}{0.420000in}}{\pgfqpoint{2.388110in}{1.994488in}}%
\pgfusepath{clip}%
\pgfsetrectcap%
\pgfsetroundjoin%
\pgfsetlinewidth{0.602250pt}%
\definecolor{currentstroke}{rgb}{0.000000,0.000000,0.000000}%
\pgfsetstrokecolor{currentstroke}%
\pgfsetdash{}{0pt}%
\pgfpathmoveto{\pgfqpoint{0.645000in}{0.880239in}}%
\pgfpathlineto{\pgfqpoint{1.010016in}{0.879695in}}%
\pgfpathlineto{\pgfqpoint{1.111637in}{0.878818in}}%
\pgfpathlineto{\pgfqpoint{1.199660in}{0.877077in}}%
\pgfpathlineto{\pgfqpoint{1.277299in}{0.874018in}}%
\pgfpathlineto{\pgfqpoint{1.346745in}{0.869209in}}%
\pgfpathlineto{\pgfqpoint{1.409562in}{0.862382in}}%
\pgfpathlineto{\pgfqpoint{1.466906in}{0.853557in}}%
\pgfpathlineto{\pgfqpoint{1.519654in}{0.843073in}}%
\pgfpathlineto{\pgfqpoint{1.568486in}{0.831526in}}%
\pgfpathlineto{\pgfqpoint{1.656463in}{0.808009in}}%
\pgfpathlineto{\pgfqpoint{1.734049in}{0.787537in}}%
\pgfpathlineto{\pgfqpoint{1.769658in}{0.779131in}}%
\pgfpathlineto{\pgfqpoint{1.803436in}{0.771999in}}%
\pgfpathlineto{\pgfqpoint{1.835562in}{0.766061in}}%
\pgfpathlineto{\pgfqpoint{1.866188in}{0.761199in}}%
\pgfpathlineto{\pgfqpoint{1.895450in}{0.757280in}}%
\pgfpathlineto{\pgfqpoint{1.923461in}{0.754174in}}%
\pgfpathlineto{\pgfqpoint{1.976131in}{0.749934in}}%
\pgfpathlineto{\pgfqpoint{2.024880in}{0.747687in}}%
\pgfpathlineto{\pgfqpoint{2.070247in}{0.746873in}}%
\pgfpathlineto{\pgfqpoint{2.112670in}{0.747089in}}%
\pgfpathlineto{\pgfqpoint{2.171545in}{0.748808in}}%
\pgfpathlineto{\pgfqpoint{2.225539in}{0.751895in}}%
\pgfpathlineto{\pgfqpoint{2.275392in}{0.756121in}}%
\pgfpathlineto{\pgfqpoint{2.321687in}{0.761255in}}%
\pgfpathlineto{\pgfqpoint{2.364891in}{0.767485in}}%
\pgfpathlineto{\pgfqpoint{2.405385in}{0.774913in}}%
\pgfpathlineto{\pgfqpoint{2.443483in}{0.783500in}}%
\pgfpathlineto{\pgfqpoint{2.479445in}{0.793126in}}%
\pgfpathlineto{\pgfqpoint{2.513495in}{0.803783in}}%
\pgfpathlineto{\pgfqpoint{2.545818in}{0.815321in}}%
\pgfpathlineto{\pgfqpoint{2.576578in}{0.827483in}}%
\pgfpathlineto{\pgfqpoint{2.651960in}{0.858448in}}%
\pgfpathlineto{\pgfqpoint{2.669474in}{0.864557in}}%
\pgfpathlineto{\pgfqpoint{2.686509in}{0.869561in}}%
\pgfpathlineto{\pgfqpoint{2.703092in}{0.873415in}}%
\pgfpathlineto{\pgfqpoint{2.719243in}{0.876166in}}%
\pgfpathlineto{\pgfqpoint{2.742707in}{0.878601in}}%
\pgfpathlineto{\pgfqpoint{2.765310in}{0.879692in}}%
\pgfpathlineto{\pgfqpoint{2.808158in}{0.880215in}}%
\pgfpathlineto{\pgfqpoint{3.043110in}{0.880266in}}%
\pgfpathlineto{\pgfqpoint{3.043110in}{0.880266in}}%
\pgfusepath{stroke}%
\end{pgfscope}%
\begin{pgfscope}%
\pgfpathrectangle{\pgfqpoint{0.650000in}{0.420000in}}{\pgfqpoint{2.388110in}{1.994488in}}%
\pgfusepath{clip}%
\pgfsetrectcap%
\pgfsetroundjoin%
\pgfsetlinewidth{1.003750pt}%
\definecolor{currentstroke}{rgb}{0.945098,0.713725,0.854902}%
\pgfsetstrokecolor{currentstroke}%
\pgfsetdash{}{0pt}%
\pgfpathmoveto{\pgfqpoint{0.645000in}{0.887242in}}%
\pgfpathlineto{\pgfqpoint{0.686817in}{0.889071in}}%
\pgfpathlineto{\pgfqpoint{0.749359in}{0.892753in}}%
\pgfpathlineto{\pgfqpoint{0.805873in}{0.897477in}}%
\pgfpathlineto{\pgfqpoint{0.857743in}{0.903452in}}%
\pgfpathlineto{\pgfqpoint{0.905952in}{0.910922in}}%
\pgfpathlineto{\pgfqpoint{0.928912in}{0.915304in}}%
\pgfpathlineto{\pgfqpoint{0.951207in}{0.920167in}}%
\pgfpathlineto{\pgfqpoint{0.972898in}{0.925556in}}%
\pgfpathlineto{\pgfqpoint{0.994037in}{0.931513in}}%
\pgfpathlineto{\pgfqpoint{1.014673in}{0.938088in}}%
\pgfpathlineto{\pgfqpoint{1.034846in}{0.945331in}}%
\pgfpathlineto{\pgfqpoint{1.054593in}{0.953296in}}%
\pgfpathlineto{\pgfqpoint{1.073948in}{0.962039in}}%
\pgfpathlineto{\pgfqpoint{1.092940in}{0.971619in}}%
\pgfpathlineto{\pgfqpoint{1.111595in}{0.982100in}}%
\pgfpathlineto{\pgfqpoint{1.129938in}{0.993542in}}%
\pgfpathlineto{\pgfqpoint{1.147990in}{1.006013in}}%
\pgfpathlineto{\pgfqpoint{1.165771in}{1.019578in}}%
\pgfpathlineto{\pgfqpoint{1.183300in}{1.034302in}}%
\pgfpathlineto{\pgfqpoint{1.200591in}{1.050250in}}%
\pgfpathlineto{\pgfqpoint{1.217661in}{1.067484in}}%
\pgfpathlineto{\pgfqpoint{1.234524in}{1.086063in}}%
\pgfpathlineto{\pgfqpoint{1.251191in}{1.106038in}}%
\pgfpathlineto{\pgfqpoint{1.267675in}{1.127456in}}%
\pgfpathlineto{\pgfqpoint{1.283987in}{1.150353in}}%
\pgfpathlineto{\pgfqpoint{1.300136in}{1.174755in}}%
\pgfpathlineto{\pgfqpoint{1.316132in}{1.200671in}}%
\pgfpathlineto{\pgfqpoint{1.331983in}{1.228100in}}%
\pgfpathlineto{\pgfqpoint{1.347698in}{1.257017in}}%
\pgfpathlineto{\pgfqpoint{1.363283in}{1.287380in}}%
\pgfpathlineto{\pgfqpoint{1.378746in}{1.319124in}}%
\pgfpathlineto{\pgfqpoint{1.394094in}{1.352159in}}%
\pgfpathlineto{\pgfqpoint{1.424466in}{1.421626in}}%
\pgfpathlineto{\pgfqpoint{1.454444in}{1.494565in}}%
\pgfpathlineto{\pgfqpoint{1.498753in}{1.606878in}}%
\pgfpathlineto{\pgfqpoint{1.542373in}{1.716610in}}%
\pgfpathlineto{\pgfqpoint{1.571118in}{1.784652in}}%
\pgfpathlineto{\pgfqpoint{1.585398in}{1.816320in}}%
\pgfpathlineto{\pgfqpoint{1.599622in}{1.846081in}}%
\pgfpathlineto{\pgfqpoint{1.613791in}{1.873705in}}%
\pgfpathlineto{\pgfqpoint{1.627908in}{1.898987in}}%
\pgfpathlineto{\pgfqpoint{1.641975in}{1.921750in}}%
\pgfpathlineto{\pgfqpoint{1.655995in}{1.941843in}}%
\pgfpathlineto{\pgfqpoint{1.669970in}{1.959152in}}%
\pgfpathlineto{\pgfqpoint{1.683901in}{1.973585in}}%
\pgfpathlineto{\pgfqpoint{1.697791in}{1.985099in}}%
\pgfpathlineto{\pgfqpoint{1.711641in}{1.993666in}}%
\pgfpathlineto{\pgfqpoint{1.725453in}{1.999293in}}%
\pgfpathlineto{\pgfqpoint{1.739229in}{2.002023in}}%
\pgfpathlineto{\pgfqpoint{1.752971in}{2.001913in}}%
\pgfpathlineto{\pgfqpoint{1.766679in}{1.999049in}}%
\pgfpathlineto{\pgfqpoint{1.780356in}{1.993541in}}%
\pgfpathlineto{\pgfqpoint{1.794002in}{1.985519in}}%
\pgfpathlineto{\pgfqpoint{1.807620in}{1.975127in}}%
\pgfpathlineto{\pgfqpoint{1.821209in}{1.962512in}}%
\pgfpathlineto{\pgfqpoint{1.834772in}{1.947852in}}%
\pgfpathlineto{\pgfqpoint{1.848310in}{1.931327in}}%
\pgfpathlineto{\pgfqpoint{1.861823in}{1.913126in}}%
\pgfpathlineto{\pgfqpoint{1.875312in}{1.893463in}}%
\pgfpathlineto{\pgfqpoint{1.888779in}{1.872529in}}%
\pgfpathlineto{\pgfqpoint{1.915650in}{1.827710in}}%
\pgfpathlineto{\pgfqpoint{1.942442in}{1.780276in}}%
\pgfpathlineto{\pgfqpoint{2.022403in}{1.636301in}}%
\pgfpathlineto{\pgfqpoint{2.048936in}{1.591176in}}%
\pgfpathlineto{\pgfqpoint{2.075418in}{1.548734in}}%
\pgfpathlineto{\pgfqpoint{2.101851in}{1.509510in}}%
\pgfpathlineto{\pgfqpoint{2.128241in}{1.473673in}}%
\pgfpathlineto{\pgfqpoint{2.154589in}{1.441054in}}%
\pgfpathlineto{\pgfqpoint{2.180900in}{1.411347in}}%
\pgfpathlineto{\pgfqpoint{2.207177in}{1.384318in}}%
\pgfpathlineto{\pgfqpoint{2.233421in}{1.359882in}}%
\pgfpathlineto{\pgfqpoint{2.259636in}{1.337995in}}%
\pgfpathlineto{\pgfqpoint{2.285823in}{1.318444in}}%
\pgfpathlineto{\pgfqpoint{2.311986in}{1.300694in}}%
\pgfpathlineto{\pgfqpoint{2.338125in}{1.284105in}}%
\pgfpathlineto{\pgfqpoint{2.364243in}{1.268451in}}%
\pgfpathlineto{\pgfqpoint{2.442484in}{1.222833in}}%
\pgfpathlineto{\pgfqpoint{2.468532in}{1.205933in}}%
\pgfpathlineto{\pgfqpoint{2.494565in}{1.187233in}}%
\pgfpathlineto{\pgfqpoint{2.520585in}{1.166725in}}%
\pgfpathlineto{\pgfqpoint{2.546569in}{1.144761in}}%
\pgfpathlineto{\pgfqpoint{2.559491in}{1.133118in}}%
\pgfpathlineto{\pgfqpoint{2.572270in}{1.120843in}}%
\pgfpathlineto{\pgfqpoint{2.584795in}{1.107935in}}%
\pgfpathlineto{\pgfqpoint{2.608756in}{1.081002in}}%
\pgfpathlineto{\pgfqpoint{2.631151in}{1.053843in}}%
\pgfpathlineto{\pgfqpoint{2.652123in}{1.026557in}}%
\pgfpathlineto{\pgfqpoint{2.681274in}{0.985724in}}%
\pgfpathlineto{\pgfqpoint{2.699367in}{0.960423in}}%
\pgfpathlineto{\pgfqpoint{2.716519in}{0.938370in}}%
\pgfpathlineto{\pgfqpoint{2.724772in}{0.928854in}}%
\pgfpathlineto{\pgfqpoint{2.732823in}{0.920433in}}%
\pgfpathlineto{\pgfqpoint{2.740682in}{0.913085in}}%
\pgfpathlineto{\pgfqpoint{2.748358in}{0.906794in}}%
\pgfpathlineto{\pgfqpoint{2.755859in}{0.901491in}}%
\pgfpathlineto{\pgfqpoint{2.763194in}{0.897093in}}%
\pgfpathlineto{\pgfqpoint{2.770369in}{0.893545in}}%
\pgfpathlineto{\pgfqpoint{2.784266in}{0.888481in}}%
\pgfpathlineto{\pgfqpoint{2.797601in}{0.885515in}}%
\pgfpathlineto{\pgfqpoint{2.810418in}{0.883837in}}%
\pgfpathlineto{\pgfqpoint{2.834647in}{0.882174in}}%
\pgfpathlineto{\pgfqpoint{2.867947in}{0.881155in}}%
\pgfpathlineto{\pgfqpoint{2.925889in}{0.880558in}}%
\pgfpathlineto{\pgfqpoint{3.040648in}{0.880312in}}%
\pgfpathlineto{\pgfqpoint{3.043110in}{0.880319in}}%
\pgfpathlineto{\pgfqpoint{3.043110in}{0.880319in}}%
\pgfusepath{stroke}%
\end{pgfscope}%
\begin{pgfscope}%
\pgfpathrectangle{\pgfqpoint{0.650000in}{0.420000in}}{\pgfqpoint{2.388110in}{1.994488in}}%
\pgfusepath{clip}%
\pgfsetrectcap%
\pgfsetroundjoin%
\pgfsetlinewidth{1.003750pt}%
\definecolor{currentstroke}{rgb}{0.870588,0.466667,0.682353}%
\pgfsetstrokecolor{currentstroke}%
\pgfsetdash{}{0pt}%
\pgfpathmoveto{\pgfqpoint{0.645000in}{0.880268in}}%
\pgfpathlineto{\pgfqpoint{1.200591in}{0.880789in}}%
\pgfpathlineto{\pgfqpoint{1.316132in}{0.881849in}}%
\pgfpathlineto{\pgfqpoint{1.394094in}{0.883444in}}%
\pgfpathlineto{\pgfqpoint{1.454444in}{0.885551in}}%
\pgfpathlineto{\pgfqpoint{1.513365in}{0.888709in}}%
\pgfpathlineto{\pgfqpoint{1.556777in}{0.891972in}}%
\pgfpathlineto{\pgfqpoint{1.599622in}{0.896176in}}%
\pgfpathlineto{\pgfqpoint{1.641975in}{0.901476in}}%
\pgfpathlineto{\pgfqpoint{1.683901in}{0.908012in}}%
\pgfpathlineto{\pgfqpoint{1.725453in}{0.915878in}}%
\pgfpathlineto{\pgfqpoint{1.766679in}{0.925108in}}%
\pgfpathlineto{\pgfqpoint{1.807620in}{0.935639in}}%
\pgfpathlineto{\pgfqpoint{1.848310in}{0.947302in}}%
\pgfpathlineto{\pgfqpoint{1.902225in}{0.964109in}}%
\pgfpathlineto{\pgfqpoint{2.022403in}{1.002343in}}%
\pgfpathlineto{\pgfqpoint{2.062183in}{1.013475in}}%
\pgfpathlineto{\pgfqpoint{2.088640in}{1.020031in}}%
\pgfpathlineto{\pgfqpoint{2.115051in}{1.025764in}}%
\pgfpathlineto{\pgfqpoint{2.141420in}{1.030590in}}%
\pgfpathlineto{\pgfqpoint{2.167749in}{1.034451in}}%
\pgfpathlineto{\pgfqpoint{2.194043in}{1.037313in}}%
\pgfpathlineto{\pgfqpoint{2.220302in}{1.039171in}}%
\pgfpathlineto{\pgfqpoint{2.246532in}{1.040049in}}%
\pgfpathlineto{\pgfqpoint{2.272733in}{1.040004in}}%
\pgfpathlineto{\pgfqpoint{2.298907in}{1.039098in}}%
\pgfpathlineto{\pgfqpoint{2.325058in}{1.037383in}}%
\pgfpathlineto{\pgfqpoint{2.351187in}{1.034847in}}%
\pgfpathlineto{\pgfqpoint{2.377295in}{1.031402in}}%
\pgfpathlineto{\pgfqpoint{2.403383in}{1.026995in}}%
\pgfpathlineto{\pgfqpoint{2.429455in}{1.021694in}}%
\pgfpathlineto{\pgfqpoint{2.455510in}{1.015616in}}%
\pgfpathlineto{\pgfqpoint{2.481550in}{1.008731in}}%
\pgfpathlineto{\pgfqpoint{2.507577in}{1.000937in}}%
\pgfpathlineto{\pgfqpoint{2.533587in}{0.992281in}}%
\pgfpathlineto{\pgfqpoint{2.572270in}{0.978045in}}%
\pgfpathlineto{\pgfqpoint{2.608756in}{0.963175in}}%
\pgfpathlineto{\pgfqpoint{2.652123in}{0.943796in}}%
\pgfpathlineto{\pgfqpoint{2.724772in}{0.911142in}}%
\pgfpathlineto{\pgfqpoint{2.748358in}{0.901740in}}%
\pgfpathlineto{\pgfqpoint{2.770369in}{0.894495in}}%
\pgfpathlineto{\pgfqpoint{2.791001in}{0.889255in}}%
\pgfpathlineto{\pgfqpoint{2.810418in}{0.885757in}}%
\pgfpathlineto{\pgfqpoint{2.828755in}{0.883619in}}%
\pgfpathlineto{\pgfqpoint{2.851717in}{0.882100in}}%
\pgfpathlineto{\pgfqpoint{2.888414in}{0.881048in}}%
\pgfpathlineto{\pgfqpoint{2.959536in}{0.880466in}}%
\pgfpathlineto{\pgfqpoint{3.043110in}{0.880314in}}%
\pgfpathlineto{\pgfqpoint{3.043110in}{0.880314in}}%
\pgfusepath{stroke}%
\end{pgfscope}%
\begin{pgfscope}%
\pgfpathrectangle{\pgfqpoint{0.650000in}{0.420000in}}{\pgfqpoint{2.388110in}{1.994488in}}%
\pgfusepath{clip}%
\pgfsetrectcap%
\pgfsetroundjoin%
\pgfsetlinewidth{1.003750pt}%
\definecolor{currentstroke}{rgb}{0.772549,0.105882,0.490196}%
\pgfsetstrokecolor{currentstroke}%
\pgfsetdash{}{0pt}%
\pgfpathmoveto{\pgfqpoint{0.645000in}{0.882467in}}%
\pgfpathlineto{\pgfqpoint{0.749359in}{0.884192in}}%
\pgfpathlineto{\pgfqpoint{0.832320in}{0.886441in}}%
\pgfpathlineto{\pgfqpoint{0.905952in}{0.889410in}}%
\pgfpathlineto{\pgfqpoint{0.972898in}{0.893206in}}%
\pgfpathlineto{\pgfqpoint{1.034846in}{0.897927in}}%
\pgfpathlineto{\pgfqpoint{1.092940in}{0.903661in}}%
\pgfpathlineto{\pgfqpoint{1.147990in}{0.910472in}}%
\pgfpathlineto{\pgfqpoint{1.200591in}{0.918393in}}%
\pgfpathlineto{\pgfqpoint{1.251191in}{0.927418in}}%
\pgfpathlineto{\pgfqpoint{1.300136in}{0.937503in}}%
\pgfpathlineto{\pgfqpoint{1.347698in}{0.948569in}}%
\pgfpathlineto{\pgfqpoint{1.394094in}{0.960513in}}%
\pgfpathlineto{\pgfqpoint{1.439502in}{0.973224in}}%
\pgfpathlineto{\pgfqpoint{1.498753in}{0.991197in}}%
\pgfpathlineto{\pgfqpoint{1.556777in}{1.010169in}}%
\pgfpathlineto{\pgfqpoint{1.613791in}{1.029961in}}%
\pgfpathlineto{\pgfqpoint{1.794002in}{1.093861in}}%
\pgfpathlineto{\pgfqpoint{1.834772in}{1.106165in}}%
\pgfpathlineto{\pgfqpoint{1.861823in}{1.113341in}}%
\pgfpathlineto{\pgfqpoint{1.888779in}{1.119599in}}%
\pgfpathlineto{\pgfqpoint{1.915650in}{1.124901in}}%
\pgfpathlineto{\pgfqpoint{1.942442in}{1.129237in}}%
\pgfpathlineto{\pgfqpoint{1.969160in}{1.132626in}}%
\pgfpathlineto{\pgfqpoint{1.995812in}{1.135124in}}%
\pgfpathlineto{\pgfqpoint{2.035676in}{1.137462in}}%
\pgfpathlineto{\pgfqpoint{2.075418in}{1.138621in}}%
\pgfpathlineto{\pgfqpoint{2.128241in}{1.139019in}}%
\pgfpathlineto{\pgfqpoint{2.167749in}{1.138316in}}%
\pgfpathlineto{\pgfqpoint{2.207177in}{1.136292in}}%
\pgfpathlineto{\pgfqpoint{2.233421in}{1.134004in}}%
\pgfpathlineto{\pgfqpoint{2.259636in}{1.130914in}}%
\pgfpathlineto{\pgfqpoint{2.285823in}{1.126941in}}%
\pgfpathlineto{\pgfqpoint{2.311986in}{1.122050in}}%
\pgfpathlineto{\pgfqpoint{2.351187in}{1.113383in}}%
\pgfpathlineto{\pgfqpoint{2.390341in}{1.103219in}}%
\pgfpathlineto{\pgfqpoint{2.416421in}{1.095374in}}%
\pgfpathlineto{\pgfqpoint{2.442484in}{1.086299in}}%
\pgfpathlineto{\pgfqpoint{2.468532in}{1.075298in}}%
\pgfpathlineto{\pgfqpoint{2.494565in}{1.062081in}}%
\pgfpathlineto{\pgfqpoint{2.520585in}{1.047249in}}%
\pgfpathlineto{\pgfqpoint{2.546569in}{1.031235in}}%
\pgfpathlineto{\pgfqpoint{2.572270in}{1.013954in}}%
\pgfpathlineto{\pgfqpoint{2.596973in}{0.995856in}}%
\pgfpathlineto{\pgfqpoint{2.681274in}{0.931820in}}%
\pgfpathlineto{\pgfqpoint{2.699367in}{0.920143in}}%
\pgfpathlineto{\pgfqpoint{2.716519in}{0.910309in}}%
\pgfpathlineto{\pgfqpoint{2.732823in}{0.902049in}}%
\pgfpathlineto{\pgfqpoint{2.748358in}{0.895455in}}%
\pgfpathlineto{\pgfqpoint{2.763194in}{0.890523in}}%
\pgfpathlineto{\pgfqpoint{2.777391in}{0.887011in}}%
\pgfpathlineto{\pgfqpoint{2.791001in}{0.884615in}}%
\pgfpathlineto{\pgfqpoint{2.810418in}{0.882529in}}%
\pgfpathlineto{\pgfqpoint{2.834647in}{0.881239in}}%
\pgfpathlineto{\pgfqpoint{2.873184in}{0.880569in}}%
\pgfpathlineto{\pgfqpoint{2.986400in}{0.880290in}}%
\pgfpathlineto{\pgfqpoint{3.043110in}{0.880273in}}%
\pgfpathlineto{\pgfqpoint{3.043110in}{0.880273in}}%
\pgfusepath{stroke}%
\end{pgfscope}%
\begin{pgfscope}%
\pgfpathrectangle{\pgfqpoint{0.650000in}{0.420000in}}{\pgfqpoint{2.388110in}{1.994488in}}%
\pgfusepath{clip}%
\pgfsetrectcap%
\pgfsetroundjoin%
\pgfsetlinewidth{1.003750pt}%
\definecolor{currentstroke}{rgb}{0.556863,0.003922,0.321569}%
\pgfsetstrokecolor{currentstroke}%
\pgfsetdash{}{0pt}%
\pgfpathmoveto{\pgfqpoint{0.645000in}{0.880246in}}%
\pgfpathlineto{\pgfqpoint{1.014673in}{0.879666in}}%
\pgfpathlineto{\pgfqpoint{1.129938in}{0.878554in}}%
\pgfpathlineto{\pgfqpoint{1.217661in}{0.876534in}}%
\pgfpathlineto{\pgfqpoint{1.283987in}{0.873678in}}%
\pgfpathlineto{\pgfqpoint{1.331983in}{0.870494in}}%
\pgfpathlineto{\pgfqpoint{1.378746in}{0.866171in}}%
\pgfpathlineto{\pgfqpoint{1.424466in}{0.860521in}}%
\pgfpathlineto{\pgfqpoint{1.469297in}{0.853427in}}%
\pgfpathlineto{\pgfqpoint{1.513365in}{0.844896in}}%
\pgfpathlineto{\pgfqpoint{1.556777in}{0.835095in}}%
\pgfpathlineto{\pgfqpoint{1.613791in}{0.820637in}}%
\pgfpathlineto{\pgfqpoint{1.725453in}{0.791283in}}%
\pgfpathlineto{\pgfqpoint{1.766679in}{0.781563in}}%
\pgfpathlineto{\pgfqpoint{1.807620in}{0.773039in}}%
\pgfpathlineto{\pgfqpoint{1.848310in}{0.765843in}}%
\pgfpathlineto{\pgfqpoint{1.888779in}{0.759993in}}%
\pgfpathlineto{\pgfqpoint{1.929055in}{0.755426in}}%
\pgfpathlineto{\pgfqpoint{1.969160in}{0.752031in}}%
\pgfpathlineto{\pgfqpoint{2.009115in}{0.749645in}}%
\pgfpathlineto{\pgfqpoint{2.062183in}{0.747792in}}%
\pgfpathlineto{\pgfqpoint{2.115051in}{0.747310in}}%
\pgfpathlineto{\pgfqpoint{2.167749in}{0.748139in}}%
\pgfpathlineto{\pgfqpoint{2.220302in}{0.750214in}}%
\pgfpathlineto{\pgfqpoint{2.272733in}{0.753614in}}%
\pgfpathlineto{\pgfqpoint{2.325058in}{0.758374in}}%
\pgfpathlineto{\pgfqpoint{2.377295in}{0.764594in}}%
\pgfpathlineto{\pgfqpoint{2.429455in}{0.772109in}}%
\pgfpathlineto{\pgfqpoint{2.468532in}{0.778861in}}%
\pgfpathlineto{\pgfqpoint{2.494565in}{0.784341in}}%
\pgfpathlineto{\pgfqpoint{2.520585in}{0.790842in}}%
\pgfpathlineto{\pgfqpoint{2.546569in}{0.798326in}}%
\pgfpathlineto{\pgfqpoint{2.572270in}{0.806621in}}%
\pgfpathlineto{\pgfqpoint{2.608756in}{0.819872in}}%
\pgfpathlineto{\pgfqpoint{2.652123in}{0.837281in}}%
\pgfpathlineto{\pgfqpoint{2.699367in}{0.856606in}}%
\pgfpathlineto{\pgfqpoint{2.724772in}{0.865641in}}%
\pgfpathlineto{\pgfqpoint{2.740682in}{0.870270in}}%
\pgfpathlineto{\pgfqpoint{2.755859in}{0.873765in}}%
\pgfpathlineto{\pgfqpoint{2.777391in}{0.877227in}}%
\pgfpathlineto{\pgfqpoint{2.797601in}{0.879070in}}%
\pgfpathlineto{\pgfqpoint{2.822755in}{0.879973in}}%
\pgfpathlineto{\pgfqpoint{2.873184in}{0.880256in}}%
\pgfpathlineto{\pgfqpoint{3.043110in}{0.880266in}}%
\pgfpathlineto{\pgfqpoint{3.043110in}{0.880266in}}%
\pgfusepath{stroke}%
\end{pgfscope}%
\begin{pgfscope}%
\pgfpathrectangle{\pgfqpoint{0.650000in}{0.420000in}}{\pgfqpoint{2.388110in}{1.994488in}}%
\pgfusepath{clip}%
\pgfsetrectcap%
\pgfsetroundjoin%
\pgfsetlinewidth{1.003750pt}%
\definecolor{currentstroke}{rgb}{0.721569,0.882353,0.525490}%
\pgfsetstrokecolor{currentstroke}%
\pgfsetdash{}{0pt}%
\pgfpathmoveto{\pgfqpoint{0.645000in}{0.887261in}}%
\pgfpathlineto{\pgfqpoint{0.688321in}{0.889171in}}%
\pgfpathlineto{\pgfqpoint{0.750863in}{0.892896in}}%
\pgfpathlineto{\pgfqpoint{0.807378in}{0.897676in}}%
\pgfpathlineto{\pgfqpoint{0.859248in}{0.903723in}}%
\pgfpathlineto{\pgfqpoint{0.907456in}{0.911282in}}%
\pgfpathlineto{\pgfqpoint{0.930416in}{0.915717in}}%
\pgfpathlineto{\pgfqpoint{0.952711in}{0.920639in}}%
\pgfpathlineto{\pgfqpoint{0.974402in}{0.926092in}}%
\pgfpathlineto{\pgfqpoint{0.995541in}{0.932121in}}%
\pgfpathlineto{\pgfqpoint{1.016177in}{0.938775in}}%
\pgfpathlineto{\pgfqpoint{1.036350in}{0.946104in}}%
\pgfpathlineto{\pgfqpoint{1.056097in}{0.954163in}}%
\pgfpathlineto{\pgfqpoint{1.075452in}{0.963008in}}%
\pgfpathlineto{\pgfqpoint{1.094444in}{0.972701in}}%
\pgfpathlineto{\pgfqpoint{1.113099in}{0.983302in}}%
\pgfpathlineto{\pgfqpoint{1.131442in}{0.994875in}}%
\pgfpathlineto{\pgfqpoint{1.149495in}{1.007486in}}%
\pgfpathlineto{\pgfqpoint{1.167276in}{1.021200in}}%
\pgfpathlineto{\pgfqpoint{1.184804in}{1.036084in}}%
\pgfpathlineto{\pgfqpoint{1.202095in}{1.052201in}}%
\pgfpathlineto{\pgfqpoint{1.219166in}{1.069613in}}%
\pgfpathlineto{\pgfqpoint{1.236028in}{1.088379in}}%
\pgfpathlineto{\pgfqpoint{1.252695in}{1.108551in}}%
\pgfpathlineto{\pgfqpoint{1.269179in}{1.130173in}}%
\pgfpathlineto{\pgfqpoint{1.285491in}{1.153281in}}%
\pgfpathlineto{\pgfqpoint{1.301640in}{1.177899in}}%
\pgfpathlineto{\pgfqpoint{1.317636in}{1.204039in}}%
\pgfpathlineto{\pgfqpoint{1.333487in}{1.231695in}}%
\pgfpathlineto{\pgfqpoint{1.349202in}{1.260843in}}%
\pgfpathlineto{\pgfqpoint{1.364787in}{1.291440in}}%
\pgfpathlineto{\pgfqpoint{1.380251in}{1.323420in}}%
\pgfpathlineto{\pgfqpoint{1.395598in}{1.356692in}}%
\pgfpathlineto{\pgfqpoint{1.425971in}{1.426634in}}%
\pgfpathlineto{\pgfqpoint{1.455948in}{1.500043in}}%
\pgfpathlineto{\pgfqpoint{1.500257in}{1.613034in}}%
\pgfpathlineto{\pgfqpoint{1.543878in}{1.723375in}}%
\pgfpathlineto{\pgfqpoint{1.572622in}{1.791746in}}%
\pgfpathlineto{\pgfqpoint{1.586903in}{1.823546in}}%
\pgfpathlineto{\pgfqpoint{1.601126in}{1.853414in}}%
\pgfpathlineto{\pgfqpoint{1.615295in}{1.881117in}}%
\pgfpathlineto{\pgfqpoint{1.629412in}{1.906447in}}%
\pgfpathlineto{\pgfqpoint{1.643480in}{1.929225in}}%
\pgfpathlineto{\pgfqpoint{1.657499in}{1.949302in}}%
\pgfpathlineto{\pgfqpoint{1.671474in}{1.966563in}}%
\pgfpathlineto{\pgfqpoint{1.685405in}{1.980924in}}%
\pgfpathlineto{\pgfqpoint{1.699295in}{1.992340in}}%
\pgfpathlineto{\pgfqpoint{1.713145in}{2.000796in}}%
\pgfpathlineto{\pgfqpoint{1.726957in}{2.006309in}}%
\pgfpathlineto{\pgfqpoint{1.740734in}{2.008924in}}%
\pgfpathlineto{\pgfqpoint{1.754475in}{2.008719in}}%
\pgfpathlineto{\pgfqpoint{1.768184in}{2.005795in}}%
\pgfpathlineto{\pgfqpoint{1.781860in}{2.000277in}}%
\pgfpathlineto{\pgfqpoint{1.795507in}{1.992309in}}%
\pgfpathlineto{\pgfqpoint{1.809124in}{1.982045in}}%
\pgfpathlineto{\pgfqpoint{1.822714in}{1.969656in}}%
\pgfpathlineto{\pgfqpoint{1.836277in}{1.955319in}}%
\pgfpathlineto{\pgfqpoint{1.849814in}{1.939220in}}%
\pgfpathlineto{\pgfqpoint{1.863327in}{1.921552in}}%
\pgfpathlineto{\pgfqpoint{1.876817in}{1.902511in}}%
\pgfpathlineto{\pgfqpoint{1.890284in}{1.882289in}}%
\pgfpathlineto{\pgfqpoint{1.917154in}{1.839011in}}%
\pgfpathlineto{\pgfqpoint{1.943946in}{1.793051in}}%
\pgfpathlineto{\pgfqpoint{2.037181in}{1.628296in}}%
\pgfpathlineto{\pgfqpoint{2.063688in}{1.584535in}}%
\pgfpathlineto{\pgfqpoint{2.090145in}{1.543481in}}%
\pgfpathlineto{\pgfqpoint{2.116556in}{1.505372in}}%
\pgfpathlineto{\pgfqpoint{2.142924in}{1.470210in}}%
\pgfpathlineto{\pgfqpoint{2.169254in}{1.437873in}}%
\pgfpathlineto{\pgfqpoint{2.195547in}{1.408272in}}%
\pgfpathlineto{\pgfqpoint{2.221807in}{1.381420in}}%
\pgfpathlineto{\pgfqpoint{2.248036in}{1.357146in}}%
\pgfpathlineto{\pgfqpoint{2.274237in}{1.335290in}}%
\pgfpathlineto{\pgfqpoint{2.300412in}{1.315701in}}%
\pgfpathlineto{\pgfqpoint{2.326562in}{1.298143in}}%
\pgfpathlineto{\pgfqpoint{2.352691in}{1.282262in}}%
\pgfpathlineto{\pgfqpoint{2.404888in}{1.251553in}}%
\pgfpathlineto{\pgfqpoint{2.430959in}{1.234542in}}%
\pgfpathlineto{\pgfqpoint{2.470036in}{1.207462in}}%
\pgfpathlineto{\pgfqpoint{2.496070in}{1.188044in}}%
\pgfpathlineto{\pgfqpoint{2.522089in}{1.167251in}}%
\pgfpathlineto{\pgfqpoint{2.548073in}{1.145236in}}%
\pgfpathlineto{\pgfqpoint{2.573775in}{1.121368in}}%
\pgfpathlineto{\pgfqpoint{2.586300in}{1.108490in}}%
\pgfpathlineto{\pgfqpoint{2.610260in}{1.081430in}}%
\pgfpathlineto{\pgfqpoint{2.632655in}{1.053640in}}%
\pgfpathlineto{\pgfqpoint{2.653627in}{1.025490in}}%
\pgfpathlineto{\pgfqpoint{2.700872in}{0.960184in}}%
\pgfpathlineto{\pgfqpoint{2.718024in}{0.938849in}}%
\pgfpathlineto{\pgfqpoint{2.734327in}{0.920979in}}%
\pgfpathlineto{\pgfqpoint{2.742186in}{0.913447in}}%
\pgfpathlineto{\pgfqpoint{2.749862in}{0.906910in}}%
\pgfpathlineto{\pgfqpoint{2.757364in}{0.901428in}}%
\pgfpathlineto{\pgfqpoint{2.764698in}{0.896963in}}%
\pgfpathlineto{\pgfqpoint{2.771873in}{0.893440in}}%
\pgfpathlineto{\pgfqpoint{2.785770in}{0.888626in}}%
\pgfpathlineto{\pgfqpoint{2.799106in}{0.885756in}}%
\pgfpathlineto{\pgfqpoint{2.818149in}{0.883372in}}%
\pgfpathlineto{\pgfqpoint{2.841940in}{0.881952in}}%
\pgfpathlineto{\pgfqpoint{2.879843in}{0.880995in}}%
\pgfpathlineto{\pgfqpoint{2.944646in}{0.880479in}}%
\pgfpathlineto{\pgfqpoint{3.043110in}{0.880304in}}%
\pgfpathlineto{\pgfqpoint{3.043110in}{0.880304in}}%
\pgfusepath{stroke}%
\end{pgfscope}%
\begin{pgfscope}%
\pgfpathrectangle{\pgfqpoint{0.650000in}{0.420000in}}{\pgfqpoint{2.388110in}{1.994488in}}%
\pgfusepath{clip}%
\pgfsetrectcap%
\pgfsetroundjoin%
\pgfsetlinewidth{1.003750pt}%
\definecolor{currentstroke}{rgb}{0.498039,0.737255,0.254902}%
\pgfsetstrokecolor{currentstroke}%
\pgfsetdash{}{0pt}%
\pgfpathmoveto{\pgfqpoint{0.645000in}{0.880268in}}%
\pgfpathlineto{\pgfqpoint{1.202095in}{0.880789in}}%
\pgfpathlineto{\pgfqpoint{1.317636in}{0.881847in}}%
\pgfpathlineto{\pgfqpoint{1.395598in}{0.883441in}}%
\pgfpathlineto{\pgfqpoint{1.455948in}{0.885544in}}%
\pgfpathlineto{\pgfqpoint{1.514869in}{0.888696in}}%
\pgfpathlineto{\pgfqpoint{1.558281in}{0.891950in}}%
\pgfpathlineto{\pgfqpoint{1.601126in}{0.896137in}}%
\pgfpathlineto{\pgfqpoint{1.643480in}{0.901409in}}%
\pgfpathlineto{\pgfqpoint{1.685405in}{0.907898in}}%
\pgfpathlineto{\pgfqpoint{1.726957in}{0.915697in}}%
\pgfpathlineto{\pgfqpoint{1.768184in}{0.924835in}}%
\pgfpathlineto{\pgfqpoint{1.809124in}{0.935250in}}%
\pgfpathlineto{\pgfqpoint{1.849814in}{0.946784in}}%
\pgfpathlineto{\pgfqpoint{1.903729in}{0.963453in}}%
\pgfpathlineto{\pgfqpoint{2.023907in}{1.001815in}}%
\pgfpathlineto{\pgfqpoint{2.063688in}{1.013125in}}%
\pgfpathlineto{\pgfqpoint{2.103356in}{1.022928in}}%
\pgfpathlineto{\pgfqpoint{2.129745in}{1.028460in}}%
\pgfpathlineto{\pgfqpoint{2.156094in}{1.033078in}}%
\pgfpathlineto{\pgfqpoint{2.182405in}{1.036694in}}%
\pgfpathlineto{\pgfqpoint{2.208681in}{1.039247in}}%
\pgfpathlineto{\pgfqpoint{2.234925in}{1.040722in}}%
\pgfpathlineto{\pgfqpoint{2.261140in}{1.041129in}}%
\pgfpathlineto{\pgfqpoint{2.287328in}{1.040519in}}%
\pgfpathlineto{\pgfqpoint{2.313490in}{1.038971in}}%
\pgfpathlineto{\pgfqpoint{2.339629in}{1.036533in}}%
\pgfpathlineto{\pgfqpoint{2.365747in}{1.033191in}}%
\pgfpathlineto{\pgfqpoint{2.391846in}{1.028969in}}%
\pgfpathlineto{\pgfqpoint{2.430959in}{1.021358in}}%
\pgfpathlineto{\pgfqpoint{2.470036in}{1.012491in}}%
\pgfpathlineto{\pgfqpoint{2.496070in}{1.005698in}}%
\pgfpathlineto{\pgfqpoint{2.522089in}{0.997957in}}%
\pgfpathlineto{\pgfqpoint{2.548073in}{0.989097in}}%
\pgfpathlineto{\pgfqpoint{2.586300in}{0.974350in}}%
\pgfpathlineto{\pgfqpoint{2.621647in}{0.959498in}}%
\pgfpathlineto{\pgfqpoint{2.663633in}{0.940422in}}%
\pgfpathlineto{\pgfqpoint{2.726276in}{0.911768in}}%
\pgfpathlineto{\pgfqpoint{2.749862in}{0.902373in}}%
\pgfpathlineto{\pgfqpoint{2.771873in}{0.894945in}}%
\pgfpathlineto{\pgfqpoint{2.792505in}{0.889475in}}%
\pgfpathlineto{\pgfqpoint{2.811922in}{0.885886in}}%
\pgfpathlineto{\pgfqpoint{2.830259in}{0.883728in}}%
\pgfpathlineto{\pgfqpoint{2.853221in}{0.882188in}}%
\pgfpathlineto{\pgfqpoint{2.889918in}{0.881065in}}%
\pgfpathlineto{\pgfqpoint{2.957018in}{0.880472in}}%
\pgfpathlineto{\pgfqpoint{3.043110in}{0.880310in}}%
\pgfpathlineto{\pgfqpoint{3.043110in}{0.880310in}}%
\pgfusepath{stroke}%
\end{pgfscope}%
\begin{pgfscope}%
\pgfpathrectangle{\pgfqpoint{0.650000in}{0.420000in}}{\pgfqpoint{2.388110in}{1.994488in}}%
\pgfusepath{clip}%
\pgfsetrectcap%
\pgfsetroundjoin%
\pgfsetlinewidth{1.003750pt}%
\definecolor{currentstroke}{rgb}{0.301961,0.572549,0.129412}%
\pgfsetstrokecolor{currentstroke}%
\pgfsetdash{}{0pt}%
\pgfpathmoveto{\pgfqpoint{0.645000in}{0.882485in}}%
\pgfpathlineto{\pgfqpoint{0.750863in}{0.884257in}}%
\pgfpathlineto{\pgfqpoint{0.833824in}{0.886545in}}%
\pgfpathlineto{\pgfqpoint{0.907456in}{0.889567in}}%
\pgfpathlineto{\pgfqpoint{0.974402in}{0.893431in}}%
\pgfpathlineto{\pgfqpoint{1.036350in}{0.898241in}}%
\pgfpathlineto{\pgfqpoint{1.094444in}{0.904085in}}%
\pgfpathlineto{\pgfqpoint{1.149495in}{0.911033in}}%
\pgfpathlineto{\pgfqpoint{1.202095in}{0.919119in}}%
\pgfpathlineto{\pgfqpoint{1.252695in}{0.928340in}}%
\pgfpathlineto{\pgfqpoint{1.301640in}{0.938654in}}%
\pgfpathlineto{\pgfqpoint{1.349202in}{0.949976in}}%
\pgfpathlineto{\pgfqpoint{1.395598in}{0.962194in}}%
\pgfpathlineto{\pgfqpoint{1.441006in}{0.975176in}}%
\pgfpathlineto{\pgfqpoint{1.500257in}{0.993444in}}%
\pgfpathlineto{\pgfqpoint{1.558281in}{1.012547in}}%
\pgfpathlineto{\pgfqpoint{1.629412in}{1.037184in}}%
\pgfpathlineto{\pgfqpoint{1.754475in}{1.081014in}}%
\pgfpathlineto{\pgfqpoint{1.795507in}{1.094320in}}%
\pgfpathlineto{\pgfqpoint{1.836277in}{1.106287in}}%
\pgfpathlineto{\pgfqpoint{1.863327in}{1.113313in}}%
\pgfpathlineto{\pgfqpoint{1.890284in}{1.119490in}}%
\pgfpathlineto{\pgfqpoint{1.917154in}{1.124820in}}%
\pgfpathlineto{\pgfqpoint{1.957314in}{1.131355in}}%
\pgfpathlineto{\pgfqpoint{1.997317in}{1.136266in}}%
\pgfpathlineto{\pgfqpoint{2.023907in}{1.138591in}}%
\pgfpathlineto{\pgfqpoint{2.050441in}{1.140094in}}%
\pgfpathlineto{\pgfqpoint{2.090145in}{1.140881in}}%
\pgfpathlineto{\pgfqpoint{2.129745in}{1.140337in}}%
\pgfpathlineto{\pgfqpoint{2.182405in}{1.138197in}}%
\pgfpathlineto{\pgfqpoint{2.221807in}{1.135377in}}%
\pgfpathlineto{\pgfqpoint{2.261140in}{1.131105in}}%
\pgfpathlineto{\pgfqpoint{2.287328in}{1.127343in}}%
\pgfpathlineto{\pgfqpoint{2.313490in}{1.122603in}}%
\pgfpathlineto{\pgfqpoint{2.339629in}{1.116716in}}%
\pgfpathlineto{\pgfqpoint{2.365747in}{1.109741in}}%
\pgfpathlineto{\pgfqpoint{2.391846in}{1.101939in}}%
\pgfpathlineto{\pgfqpoint{2.417926in}{1.093199in}}%
\pgfpathlineto{\pgfqpoint{2.443989in}{1.083189in}}%
\pgfpathlineto{\pgfqpoint{2.470036in}{1.072067in}}%
\pgfpathlineto{\pgfqpoint{2.496070in}{1.059863in}}%
\pgfpathlineto{\pgfqpoint{2.522089in}{1.046356in}}%
\pgfpathlineto{\pgfqpoint{2.548073in}{1.031351in}}%
\pgfpathlineto{\pgfqpoint{2.573775in}{1.014817in}}%
\pgfpathlineto{\pgfqpoint{2.598477in}{0.997217in}}%
\pgfpathlineto{\pgfqpoint{2.691949in}{0.926835in}}%
\pgfpathlineto{\pgfqpoint{2.709559in}{0.915413in}}%
\pgfpathlineto{\pgfqpoint{2.726276in}{0.906005in}}%
\pgfpathlineto{\pgfqpoint{2.742186in}{0.898508in}}%
\pgfpathlineto{\pgfqpoint{2.757364in}{0.892771in}}%
\pgfpathlineto{\pgfqpoint{2.771873in}{0.888595in}}%
\pgfpathlineto{\pgfqpoint{2.785770in}{0.885699in}}%
\pgfpathlineto{\pgfqpoint{2.805576in}{0.883065in}}%
\pgfpathlineto{\pgfqpoint{2.824259in}{0.881730in}}%
\pgfpathlineto{\pgfqpoint{2.853221in}{0.880851in}}%
\pgfpathlineto{\pgfqpoint{2.918417in}{0.880382in}}%
\pgfpathlineto{\pgfqpoint{3.043110in}{0.880274in}}%
\pgfpathlineto{\pgfqpoint{3.043110in}{0.880274in}}%
\pgfusepath{stroke}%
\end{pgfscope}%
\begin{pgfscope}%
\pgfpathrectangle{\pgfqpoint{0.650000in}{0.420000in}}{\pgfqpoint{2.388110in}{1.994488in}}%
\pgfusepath{clip}%
\pgfsetrectcap%
\pgfsetroundjoin%
\pgfsetlinewidth{1.003750pt}%
\definecolor{currentstroke}{rgb}{0.152941,0.392157,0.098039}%
\pgfsetstrokecolor{currentstroke}%
\pgfsetdash{}{0pt}%
\pgfpathmoveto{\pgfqpoint{0.645000in}{0.880245in}}%
\pgfpathlineto{\pgfqpoint{1.016177in}{0.879657in}}%
\pgfpathlineto{\pgfqpoint{1.131442in}{0.878529in}}%
\pgfpathlineto{\pgfqpoint{1.202095in}{0.877008in}}%
\pgfpathlineto{\pgfqpoint{1.269179in}{0.874452in}}%
\pgfpathlineto{\pgfqpoint{1.317636in}{0.871571in}}%
\pgfpathlineto{\pgfqpoint{1.364787in}{0.867620in}}%
\pgfpathlineto{\pgfqpoint{1.410836in}{0.862394in}}%
\pgfpathlineto{\pgfqpoint{1.455948in}{0.855749in}}%
\pgfpathlineto{\pgfqpoint{1.500257in}{0.847646in}}%
\pgfpathlineto{\pgfqpoint{1.543878in}{0.838194in}}%
\pgfpathlineto{\pgfqpoint{1.586903in}{0.827667in}}%
\pgfpathlineto{\pgfqpoint{1.657499in}{0.808923in}}%
\pgfpathlineto{\pgfqpoint{1.726957in}{0.790766in}}%
\pgfpathlineto{\pgfqpoint{1.768184in}{0.781043in}}%
\pgfpathlineto{\pgfqpoint{1.809124in}{0.772561in}}%
\pgfpathlineto{\pgfqpoint{1.849814in}{0.765442in}}%
\pgfpathlineto{\pgfqpoint{1.890284in}{0.759681in}}%
\pgfpathlineto{\pgfqpoint{1.930560in}{0.755186in}}%
\pgfpathlineto{\pgfqpoint{1.970665in}{0.751832in}}%
\pgfpathlineto{\pgfqpoint{2.010619in}{0.749463in}}%
\pgfpathlineto{\pgfqpoint{2.063688in}{0.747578in}}%
\pgfpathlineto{\pgfqpoint{2.116556in}{0.746946in}}%
\pgfpathlineto{\pgfqpoint{2.169254in}{0.747490in}}%
\pgfpathlineto{\pgfqpoint{2.221807in}{0.749138in}}%
\pgfpathlineto{\pgfqpoint{2.274237in}{0.751888in}}%
\pgfpathlineto{\pgfqpoint{2.326562in}{0.755885in}}%
\pgfpathlineto{\pgfqpoint{2.365747in}{0.760013in}}%
\pgfpathlineto{\pgfqpoint{2.404888in}{0.765475in}}%
\pgfpathlineto{\pgfqpoint{2.443989in}{0.772347in}}%
\pgfpathlineto{\pgfqpoint{2.483055in}{0.780577in}}%
\pgfpathlineto{\pgfqpoint{2.522089in}{0.790171in}}%
\pgfpathlineto{\pgfqpoint{2.548073in}{0.797452in}}%
\pgfpathlineto{\pgfqpoint{2.573775in}{0.805627in}}%
\pgfpathlineto{\pgfqpoint{2.598477in}{0.814581in}}%
\pgfpathlineto{\pgfqpoint{2.632655in}{0.828494in}}%
\pgfpathlineto{\pgfqpoint{2.691949in}{0.853188in}}%
\pgfpathlineto{\pgfqpoint{2.718024in}{0.862841in}}%
\pgfpathlineto{\pgfqpoint{2.734327in}{0.867996in}}%
\pgfpathlineto{\pgfqpoint{2.749862in}{0.872118in}}%
\pgfpathlineto{\pgfqpoint{2.764698in}{0.875222in}}%
\pgfpathlineto{\pgfqpoint{2.778895in}{0.877351in}}%
\pgfpathlineto{\pgfqpoint{2.799106in}{0.879138in}}%
\pgfpathlineto{\pgfqpoint{2.824259in}{0.880038in}}%
\pgfpathlineto{\pgfqpoint{2.879843in}{0.880258in}}%
\pgfpathlineto{\pgfqpoint{3.043110in}{0.880267in}}%
\pgfpathlineto{\pgfqpoint{3.043110in}{0.880267in}}%
\pgfusepath{stroke}%
\end{pgfscope}%
\begin{pgfscope}%
\pgfpathrectangle{\pgfqpoint{0.650000in}{0.420000in}}{\pgfqpoint{2.388110in}{1.994488in}}%
\pgfusepath{clip}%
\pgfsetrectcap%
\pgfsetroundjoin%
\pgfsetlinewidth{1.003750pt}%
\definecolor{currentstroke}{rgb}{0.501961,0.803922,0.756863}%
\pgfsetstrokecolor{currentstroke}%
\pgfsetdash{}{0pt}%
\pgfpathmoveto{\pgfqpoint{0.645000in}{0.887749in}}%
\pgfpathlineto{\pgfqpoint{0.795272in}{0.894891in}}%
\pgfpathlineto{\pgfqpoint{0.932227in}{0.913622in}}%
\pgfpathlineto{\pgfqpoint{1.030365in}{0.940337in}}%
\pgfpathlineto{\pgfqpoint{1.107239in}{0.975162in}}%
\pgfpathlineto{\pgfqpoint{1.170666in}{1.017987in}}%
\pgfpathlineto{\pgfqpoint{1.224817in}{1.068395in}}%
\pgfpathlineto{\pgfqpoint{1.272181in}{1.125634in}}%
\pgfpathlineto{\pgfqpoint{1.314364in}{1.188625in}}%
\pgfpathlineto{\pgfqpoint{1.352464in}{1.256031in}}%
\pgfpathlineto{\pgfqpoint{1.387260in}{1.326339in}}%
\pgfpathlineto{\pgfqpoint{1.419332in}{1.397960in}}%
\pgfpathlineto{\pgfqpoint{1.449117in}{1.469336in}}%
\pgfpathlineto{\pgfqpoint{1.503121in}{1.605724in}}%
\pgfpathlineto{\pgfqpoint{1.527827in}{1.668410in}}%
\pgfpathlineto{\pgfqpoint{1.551254in}{1.726256in}}%
\pgfpathlineto{\pgfqpoint{1.573549in}{1.778681in}}%
\pgfpathlineto{\pgfqpoint{1.594835in}{1.825323in}}%
\pgfpathlineto{\pgfqpoint{1.615218in}{1.866005in}}%
\pgfpathlineto{\pgfqpoint{1.634786in}{1.900717in}}%
\pgfpathlineto{\pgfqpoint{1.653617in}{1.929597in}}%
\pgfpathlineto{\pgfqpoint{1.671776in}{1.952879in}}%
\pgfpathlineto{\pgfqpoint{1.689323in}{1.970874in}}%
\pgfpathlineto{\pgfqpoint{1.706307in}{1.983941in}}%
\pgfpathlineto{\pgfqpoint{1.722774in}{1.992465in}}%
\pgfpathlineto{\pgfqpoint{1.738764in}{1.996839in}}%
\pgfpathlineto{\pgfqpoint{1.754311in}{1.997460in}}%
\pgfpathlineto{\pgfqpoint{1.769449in}{1.994720in}}%
\pgfpathlineto{\pgfqpoint{1.784205in}{1.989002in}}%
\pgfpathlineto{\pgfqpoint{1.798605in}{1.980662in}}%
\pgfpathlineto{\pgfqpoint{1.812672in}{1.970056in}}%
\pgfpathlineto{\pgfqpoint{1.826428in}{1.957507in}}%
\pgfpathlineto{\pgfqpoint{1.839891in}{1.943321in}}%
\pgfpathlineto{\pgfqpoint{1.853080in}{1.927785in}}%
\pgfpathlineto{\pgfqpoint{1.866010in}{1.911161in}}%
\pgfpathlineto{\pgfqpoint{1.878696in}{1.893669in}}%
\pgfpathlineto{\pgfqpoint{1.903390in}{1.856868in}}%
\pgfpathlineto{\pgfqpoint{1.927257in}{1.818675in}}%
\pgfpathlineto{\pgfqpoint{1.961687in}{1.760769in}}%
\pgfpathlineto{\pgfqpoint{2.026407in}{1.650112in}}%
\pgfpathlineto{\pgfqpoint{2.057044in}{1.600261in}}%
\pgfpathlineto{\pgfqpoint{2.086715in}{1.554786in}}%
\pgfpathlineto{\pgfqpoint{2.106014in}{1.526935in}}%
\pgfpathlineto{\pgfqpoint{2.124961in}{1.501005in}}%
\pgfpathlineto{\pgfqpoint{2.143582in}{1.476857in}}%
\pgfpathlineto{\pgfqpoint{2.170951in}{1.443545in}}%
\pgfpathlineto{\pgfqpoint{2.197707in}{1.413084in}}%
\pgfpathlineto{\pgfqpoint{2.223909in}{1.385119in}}%
\pgfpathlineto{\pgfqpoint{2.241096in}{1.367951in}}%
\pgfpathlineto{\pgfqpoint{2.258075in}{1.352143in}}%
\pgfpathlineto{\pgfqpoint{2.274857in}{1.337783in}}%
\pgfpathlineto{\pgfqpoint{2.291456in}{1.324750in}}%
\pgfpathlineto{\pgfqpoint{2.324144in}{1.300996in}}%
\pgfpathlineto{\pgfqpoint{2.356217in}{1.278649in}}%
\pgfpathlineto{\pgfqpoint{2.387741in}{1.258494in}}%
\pgfpathlineto{\pgfqpoint{2.434126in}{1.229373in}}%
\pgfpathlineto{\pgfqpoint{2.456957in}{1.213593in}}%
\pgfpathlineto{\pgfqpoint{2.479572in}{1.196209in}}%
\pgfpathlineto{\pgfqpoint{2.538941in}{1.147997in}}%
\pgfpathlineto{\pgfqpoint{2.560891in}{1.128694in}}%
\pgfpathlineto{\pgfqpoint{2.575439in}{1.114429in}}%
\pgfpathlineto{\pgfqpoint{2.589923in}{1.098477in}}%
\pgfpathlineto{\pgfqpoint{2.611535in}{1.072200in}}%
\pgfpathlineto{\pgfqpoint{2.633019in}{1.044902in}}%
\pgfpathlineto{\pgfqpoint{2.654383in}{1.015961in}}%
\pgfpathlineto{\pgfqpoint{2.703813in}{0.946722in}}%
\pgfpathlineto{\pgfqpoint{2.717837in}{0.929578in}}%
\pgfpathlineto{\pgfqpoint{2.731821in}{0.914698in}}%
\pgfpathlineto{\pgfqpoint{2.738799in}{0.908372in}}%
\pgfpathlineto{\pgfqpoint{2.745767in}{0.902929in}}%
\pgfpathlineto{\pgfqpoint{2.752725in}{0.898354in}}%
\pgfpathlineto{\pgfqpoint{2.759675in}{0.894611in}}%
\pgfpathlineto{\pgfqpoint{2.766616in}{0.891632in}}%
\pgfpathlineto{\pgfqpoint{2.780473in}{0.887481in}}%
\pgfpathlineto{\pgfqpoint{2.794296in}{0.884913in}}%
\pgfpathlineto{\pgfqpoint{2.808088in}{0.883426in}}%
\pgfpathlineto{\pgfqpoint{2.835583in}{0.881898in}}%
\pgfpathlineto{\pgfqpoint{2.876622in}{0.880906in}}%
\pgfpathlineto{\pgfqpoint{2.991797in}{0.880354in}}%
\pgfpathlineto{\pgfqpoint{3.043110in}{0.880304in}}%
\pgfpathlineto{\pgfqpoint{3.043110in}{0.880304in}}%
\pgfusepath{stroke}%
\end{pgfscope}%
\begin{pgfscope}%
\pgfpathrectangle{\pgfqpoint{0.650000in}{0.420000in}}{\pgfqpoint{2.388110in}{1.994488in}}%
\pgfusepath{clip}%
\pgfsetrectcap%
\pgfsetroundjoin%
\pgfsetlinewidth{1.003750pt}%
\definecolor{currentstroke}{rgb}{0.400000,0.760784,0.643137}%
\pgfsetstrokecolor{currentstroke}%
\pgfsetdash{}{0pt}%
\pgfpathmoveto{\pgfqpoint{0.645000in}{0.880269in}}%
\pgfpathlineto{\pgfqpoint{1.170666in}{0.880659in}}%
\pgfpathlineto{\pgfqpoint{1.314364in}{0.881854in}}%
\pgfpathlineto{\pgfqpoint{1.387260in}{0.883327in}}%
\pgfpathlineto{\pgfqpoint{1.449117in}{0.885440in}}%
\pgfpathlineto{\pgfqpoint{1.503121in}{0.888245in}}%
\pgfpathlineto{\pgfqpoint{1.551254in}{0.891767in}}%
\pgfpathlineto{\pgfqpoint{1.594835in}{0.896006in}}%
\pgfpathlineto{\pgfqpoint{1.634786in}{0.900946in}}%
\pgfpathlineto{\pgfqpoint{1.671776in}{0.906550in}}%
\pgfpathlineto{\pgfqpoint{1.706307in}{0.912762in}}%
\pgfpathlineto{\pgfqpoint{1.738764in}{0.919504in}}%
\pgfpathlineto{\pgfqpoint{1.784205in}{0.930434in}}%
\pgfpathlineto{\pgfqpoint{1.826428in}{0.942056in}}%
\pgfpathlineto{\pgfqpoint{1.866010in}{0.954034in}}%
\pgfpathlineto{\pgfqpoint{1.927257in}{0.973919in}}%
\pgfpathlineto{\pgfqpoint{2.005382in}{0.999435in}}%
\pgfpathlineto{\pgfqpoint{2.046945in}{1.011823in}}%
\pgfpathlineto{\pgfqpoint{2.076925in}{1.019810in}}%
\pgfpathlineto{\pgfqpoint{2.106014in}{1.026633in}}%
\pgfpathlineto{\pgfqpoint{2.134311in}{1.032251in}}%
\pgfpathlineto{\pgfqpoint{2.161899in}{1.036642in}}%
\pgfpathlineto{\pgfqpoint{2.188853in}{1.039864in}}%
\pgfpathlineto{\pgfqpoint{2.215233in}{1.042050in}}%
\pgfpathlineto{\pgfqpoint{2.241096in}{1.043232in}}%
\pgfpathlineto{\pgfqpoint{2.266490in}{1.043344in}}%
\pgfpathlineto{\pgfqpoint{2.291456in}{1.042414in}}%
\pgfpathlineto{\pgfqpoint{2.316032in}{1.040566in}}%
\pgfpathlineto{\pgfqpoint{2.348252in}{1.036945in}}%
\pgfpathlineto{\pgfqpoint{2.387741in}{1.031095in}}%
\pgfpathlineto{\pgfqpoint{2.426464in}{1.024396in}}%
\pgfpathlineto{\pgfqpoint{2.449372in}{1.019559in}}%
\pgfpathlineto{\pgfqpoint{2.472057in}{1.013722in}}%
\pgfpathlineto{\pgfqpoint{2.501987in}{1.004767in}}%
\pgfpathlineto{\pgfqpoint{2.531588in}{0.994809in}}%
\pgfpathlineto{\pgfqpoint{2.560891in}{0.983675in}}%
\pgfpathlineto{\pgfqpoint{2.604346in}{0.965517in}}%
\pgfpathlineto{\pgfqpoint{2.640153in}{0.949341in}}%
\pgfpathlineto{\pgfqpoint{2.738799in}{0.903445in}}%
\pgfpathlineto{\pgfqpoint{2.759675in}{0.895928in}}%
\pgfpathlineto{\pgfqpoint{2.773549in}{0.891853in}}%
\pgfpathlineto{\pgfqpoint{2.787389in}{0.888598in}}%
\pgfpathlineto{\pgfqpoint{2.808088in}{0.885177in}}%
\pgfpathlineto{\pgfqpoint{2.828720in}{0.883108in}}%
\pgfpathlineto{\pgfqpoint{2.856132in}{0.881688in}}%
\pgfpathlineto{\pgfqpoint{2.897058in}{0.880826in}}%
\pgfpathlineto{\pgfqpoint{2.985061in}{0.880369in}}%
\pgfpathlineto{\pgfqpoint{3.043110in}{0.880303in}}%
\pgfpathlineto{\pgfqpoint{3.043110in}{0.880303in}}%
\pgfusepath{stroke}%
\end{pgfscope}%
\begin{pgfscope}%
\pgfpathrectangle{\pgfqpoint{0.650000in}{0.420000in}}{\pgfqpoint{2.388110in}{1.994488in}}%
\pgfusepath{clip}%
\pgfsetrectcap%
\pgfsetroundjoin%
\pgfsetlinewidth{1.003750pt}%
\definecolor{currentstroke}{rgb}{0.207843,0.592157,0.560784}%
\pgfsetstrokecolor{currentstroke}%
\pgfsetdash{}{0pt}%
\pgfpathmoveto{\pgfqpoint{0.645000in}{0.882974in}}%
\pgfpathlineto{\pgfqpoint{0.795272in}{0.885409in}}%
\pgfpathlineto{\pgfqpoint{0.932227in}{0.891002in}}%
\pgfpathlineto{\pgfqpoint{1.030365in}{0.897982in}}%
\pgfpathlineto{\pgfqpoint{1.107239in}{0.905970in}}%
\pgfpathlineto{\pgfqpoint{1.170666in}{0.914649in}}%
\pgfpathlineto{\pgfqpoint{1.224817in}{0.923761in}}%
\pgfpathlineto{\pgfqpoint{1.272181in}{0.933100in}}%
\pgfpathlineto{\pgfqpoint{1.314364in}{0.942508in}}%
\pgfpathlineto{\pgfqpoint{1.352464in}{0.951870in}}%
\pgfpathlineto{\pgfqpoint{1.419332in}{0.970167in}}%
\pgfpathlineto{\pgfqpoint{1.476957in}{0.987665in}}%
\pgfpathlineto{\pgfqpoint{1.527827in}{1.004293in}}%
\pgfpathlineto{\pgfqpoint{1.594835in}{1.027671in}}%
\pgfpathlineto{\pgfqpoint{1.769449in}{1.090721in}}%
\pgfpathlineto{\pgfqpoint{1.812672in}{1.104236in}}%
\pgfpathlineto{\pgfqpoint{1.839891in}{1.111784in}}%
\pgfpathlineto{\pgfqpoint{1.866010in}{1.118213in}}%
\pgfpathlineto{\pgfqpoint{1.891152in}{1.123586in}}%
\pgfpathlineto{\pgfqpoint{1.927257in}{1.129871in}}%
\pgfpathlineto{\pgfqpoint{1.961687in}{1.134442in}}%
\pgfpathlineto{\pgfqpoint{1.994673in}{1.137765in}}%
\pgfpathlineto{\pgfqpoint{2.046945in}{1.141600in}}%
\pgfpathlineto{\pgfqpoint{2.086715in}{1.143564in}}%
\pgfpathlineto{\pgfqpoint{2.115530in}{1.144103in}}%
\pgfpathlineto{\pgfqpoint{2.152777in}{1.143536in}}%
\pgfpathlineto{\pgfqpoint{2.188853in}{1.141684in}}%
\pgfpathlineto{\pgfqpoint{2.241096in}{1.137103in}}%
\pgfpathlineto{\pgfqpoint{2.283179in}{1.132998in}}%
\pgfpathlineto{\pgfqpoint{2.307881in}{1.129335in}}%
\pgfpathlineto{\pgfqpoint{2.332217in}{1.124383in}}%
\pgfpathlineto{\pgfqpoint{2.356217in}{1.118343in}}%
\pgfpathlineto{\pgfqpoint{2.379908in}{1.111449in}}%
\pgfpathlineto{\pgfqpoint{2.464519in}{1.083663in}}%
\pgfpathlineto{\pgfqpoint{2.479572in}{1.077465in}}%
\pgfpathlineto{\pgfqpoint{2.494537in}{1.069786in}}%
\pgfpathlineto{\pgfqpoint{2.516827in}{1.056314in}}%
\pgfpathlineto{\pgfqpoint{2.546275in}{1.036846in}}%
\pgfpathlineto{\pgfqpoint{2.568173in}{1.020810in}}%
\pgfpathlineto{\pgfqpoint{2.597142in}{0.997024in}}%
\pgfpathlineto{\pgfqpoint{2.625872in}{0.973329in}}%
\pgfpathlineto{\pgfqpoint{2.689747in}{0.924017in}}%
\pgfpathlineto{\pgfqpoint{2.703813in}{0.914691in}}%
\pgfpathlineto{\pgfqpoint{2.717837in}{0.906408in}}%
\pgfpathlineto{\pgfqpoint{2.731821in}{0.899589in}}%
\pgfpathlineto{\pgfqpoint{2.745767in}{0.894092in}}%
\pgfpathlineto{\pgfqpoint{2.759675in}{0.889711in}}%
\pgfpathlineto{\pgfqpoint{2.773549in}{0.886432in}}%
\pgfpathlineto{\pgfqpoint{2.787389in}{0.884155in}}%
\pgfpathlineto{\pgfqpoint{2.808088in}{0.882186in}}%
\pgfpathlineto{\pgfqpoint{2.835583in}{0.881036in}}%
\pgfpathlineto{\pgfqpoint{2.890252in}{0.880430in}}%
\pgfpathlineto{\pgfqpoint{3.043110in}{0.880270in}}%
\pgfpathlineto{\pgfqpoint{3.043110in}{0.880270in}}%
\pgfusepath{stroke}%
\end{pgfscope}%
\begin{pgfscope}%
\pgfpathrectangle{\pgfqpoint{0.650000in}{0.420000in}}{\pgfqpoint{2.388110in}{1.994488in}}%
\pgfusepath{clip}%
\pgfsetrectcap%
\pgfsetroundjoin%
\pgfsetlinewidth{1.003750pt}%
\definecolor{currentstroke}{rgb}{0.003922,0.400000,0.368627}%
\pgfsetstrokecolor{currentstroke}%
\pgfsetdash{}{0pt}%
\pgfpathmoveto{\pgfqpoint{0.645000in}{0.880233in}}%
\pgfpathlineto{\pgfqpoint{1.030365in}{0.879582in}}%
\pgfpathlineto{\pgfqpoint{1.170666in}{0.877806in}}%
\pgfpathlineto{\pgfqpoint{1.224817in}{0.876290in}}%
\pgfpathlineto{\pgfqpoint{1.272181in}{0.874279in}}%
\pgfpathlineto{\pgfqpoint{1.314364in}{0.871742in}}%
\pgfpathlineto{\pgfqpoint{1.352464in}{0.868672in}}%
\pgfpathlineto{\pgfqpoint{1.387260in}{0.865084in}}%
\pgfpathlineto{\pgfqpoint{1.419332in}{0.861015in}}%
\pgfpathlineto{\pgfqpoint{1.449117in}{0.856518in}}%
\pgfpathlineto{\pgfqpoint{1.476957in}{0.851660in}}%
\pgfpathlineto{\pgfqpoint{1.527827in}{0.841164in}}%
\pgfpathlineto{\pgfqpoint{1.573549in}{0.830142in}}%
\pgfpathlineto{\pgfqpoint{1.634786in}{0.813831in}}%
\pgfpathlineto{\pgfqpoint{1.722774in}{0.790180in}}%
\pgfpathlineto{\pgfqpoint{1.769449in}{0.778906in}}%
\pgfpathlineto{\pgfqpoint{1.812672in}{0.769799in}}%
\pgfpathlineto{\pgfqpoint{1.853080in}{0.762600in}}%
\pgfpathlineto{\pgfqpoint{1.891152in}{0.757000in}}%
\pgfpathlineto{\pgfqpoint{1.927257in}{0.752728in}}%
\pgfpathlineto{\pgfqpoint{1.972832in}{0.748710in}}%
\pgfpathlineto{\pgfqpoint{2.015958in}{0.746204in}}%
\pgfpathlineto{\pgfqpoint{2.057044in}{0.744883in}}%
\pgfpathlineto{\pgfqpoint{2.106014in}{0.744531in}}%
\pgfpathlineto{\pgfqpoint{2.152777in}{0.745281in}}%
\pgfpathlineto{\pgfqpoint{2.197707in}{0.747035in}}%
\pgfpathlineto{\pgfqpoint{2.241096in}{0.749780in}}%
\pgfpathlineto{\pgfqpoint{2.324144in}{0.756500in}}%
\pgfpathlineto{\pgfqpoint{2.364147in}{0.760709in}}%
\pgfpathlineto{\pgfqpoint{2.403317in}{0.766065in}}%
\pgfpathlineto{\pgfqpoint{2.434126in}{0.771216in}}%
\pgfpathlineto{\pgfqpoint{2.464519in}{0.777344in}}%
\pgfpathlineto{\pgfqpoint{2.501987in}{0.785979in}}%
\pgfpathlineto{\pgfqpoint{2.531588in}{0.793721in}}%
\pgfpathlineto{\pgfqpoint{2.568173in}{0.804254in}}%
\pgfpathlineto{\pgfqpoint{2.589923in}{0.811676in}}%
\pgfpathlineto{\pgfqpoint{2.618710in}{0.823296in}}%
\pgfpathlineto{\pgfqpoint{2.668564in}{0.844686in}}%
\pgfpathlineto{\pgfqpoint{2.703813in}{0.859793in}}%
\pgfpathlineto{\pgfqpoint{2.724834in}{0.867446in}}%
\pgfpathlineto{\pgfqpoint{2.738799in}{0.871587in}}%
\pgfpathlineto{\pgfqpoint{2.752725in}{0.874823in}}%
\pgfpathlineto{\pgfqpoint{2.766616in}{0.877080in}}%
\pgfpathlineto{\pgfqpoint{2.787389in}{0.878949in}}%
\pgfpathlineto{\pgfqpoint{2.814973in}{0.879919in}}%
\pgfpathlineto{\pgfqpoint{2.869798in}{0.880255in}}%
\pgfpathlineto{\pgfqpoint{3.043110in}{0.880268in}}%
\pgfpathlineto{\pgfqpoint{3.043110in}{0.880268in}}%
\pgfusepath{stroke}%
\end{pgfscope}%
\begin{pgfscope}%
\pgfsetrectcap%
\pgfsetmiterjoin%
\pgfsetlinewidth{0.803000pt}%
\definecolor{currentstroke}{rgb}{0.000000,0.000000,0.000000}%
\pgfsetstrokecolor{currentstroke}%
\pgfsetdash{}{0pt}%
\pgfpathmoveto{\pgfqpoint{0.650000in}{0.420000in}}%
\pgfpathlineto{\pgfqpoint{0.650000in}{2.414488in}}%
\pgfusepath{stroke}%
\end{pgfscope}%
\begin{pgfscope}%
\pgfsetrectcap%
\pgfsetmiterjoin%
\pgfsetlinewidth{0.803000pt}%
\definecolor{currentstroke}{rgb}{0.000000,0.000000,0.000000}%
\pgfsetstrokecolor{currentstroke}%
\pgfsetdash{}{0pt}%
\pgfpathmoveto{\pgfqpoint{3.038110in}{0.420000in}}%
\pgfpathlineto{\pgfqpoint{3.038110in}{2.414488in}}%
\pgfusepath{stroke}%
\end{pgfscope}%
\begin{pgfscope}%
\pgfsetrectcap%
\pgfsetmiterjoin%
\pgfsetlinewidth{0.803000pt}%
\definecolor{currentstroke}{rgb}{0.000000,0.000000,0.000000}%
\pgfsetstrokecolor{currentstroke}%
\pgfsetdash{}{0pt}%
\pgfpathmoveto{\pgfqpoint{0.650000in}{0.420000in}}%
\pgfpathlineto{\pgfqpoint{3.038110in}{0.420000in}}%
\pgfusepath{stroke}%
\end{pgfscope}%
\begin{pgfscope}%
\pgfsetrectcap%
\pgfsetmiterjoin%
\pgfsetlinewidth{0.803000pt}%
\definecolor{currentstroke}{rgb}{0.000000,0.000000,0.000000}%
\pgfsetstrokecolor{currentstroke}%
\pgfsetdash{}{0pt}%
\pgfpathmoveto{\pgfqpoint{0.650000in}{2.414488in}}%
\pgfpathlineto{\pgfqpoint{3.038110in}{2.414488in}}%
\pgfusepath{stroke}%
\end{pgfscope}%
\begin{pgfscope}%
\pgfsetrectcap%
\pgfsetroundjoin%
\pgfsetlinewidth{0.602250pt}%
\definecolor{currentstroke}{rgb}{0.000000,0.000000,0.000000}%
\pgfsetstrokecolor{currentstroke}%
\pgfsetdash{}{0pt}%
\pgfpathmoveto{\pgfqpoint{2.108975in}{2.275599in}}%
\pgfpathlineto{\pgfqpoint{2.220086in}{2.275599in}}%
\pgfpathlineto{\pgfqpoint{2.331197in}{2.275599in}}%
\pgfusepath{stroke}%
\end{pgfscope}%
\begin{pgfscope}%
\definecolor{textcolor}{rgb}{0.000000,0.000000,0.000000}%
\pgfsetstrokecolor{textcolor}%
\pgfsetfillcolor{textcolor}%
\pgftext[x=2.420086in,y=2.236710in,left,base]{\color{textcolor}{\rmfamily\fontsize{8.000000}{9.600000}\selectfont\catcode`\^=\active\def^{\ifmmode\sp\else\^{}\fi}\catcode`\%=\active\def%{\%}ref}}%
\end{pgfscope}%
\begin{pgfscope}%
\pgfsetrectcap%
\pgfsetroundjoin%
\pgfsetlinewidth{1.003750pt}%
\definecolor{currentstroke}{rgb}{0.945098,0.713725,0.854902}%
\pgfsetstrokecolor{currentstroke}%
\pgfsetdash{}{0pt}%
\pgfpathmoveto{\pgfqpoint{2.108975in}{2.120661in}}%
\pgfpathlineto{\pgfqpoint{2.220086in}{2.120661in}}%
\pgfpathlineto{\pgfqpoint{2.331197in}{2.120661in}}%
\pgfusepath{stroke}%
\end{pgfscope}%
\begin{pgfscope}%
\definecolor{textcolor}{rgb}{0.000000,0.000000,0.000000}%
\pgfsetstrokecolor{textcolor}%
\pgfsetfillcolor{textcolor}%
\pgftext[x=2.420086in,y=2.081772in,left,base]{\color{textcolor}{\rmfamily\fontsize{8.000000}{9.600000}\selectfont\catcode`\^=\active\def^{\ifmmode\sp\else\^{}\fi}\catcode`\%=\active\def%{\%}base}}%
\end{pgfscope}%
\begin{pgfscope}%
\pgfsetrectcap%
\pgfsetroundjoin%
\pgfsetlinewidth{1.003750pt}%
\definecolor{currentstroke}{rgb}{0.721569,0.882353,0.525490}%
\pgfsetstrokecolor{currentstroke}%
\pgfsetdash{}{0pt}%
\pgfpathmoveto{\pgfqpoint{2.108975in}{1.965723in}}%
\pgfpathlineto{\pgfqpoint{2.220086in}{1.965723in}}%
\pgfpathlineto{\pgfqpoint{2.331197in}{1.965723in}}%
\pgfusepath{stroke}%
\end{pgfscope}%
\begin{pgfscope}%
\definecolor{textcolor}{rgb}{0.000000,0.000000,0.000000}%
\pgfsetstrokecolor{textcolor}%
\pgfsetfillcolor{textcolor}%
\pgftext[x=2.420086in,y=1.926834in,left,base]{\color{textcolor}{\rmfamily\fontsize{8.000000}{9.600000}\selectfont\catcode`\^=\active\def^{\ifmmode\sp\else\^{}\fi}\catcode`\%=\active\def%{\%}optimised}}%
\end{pgfscope}%
\begin{pgfscope}%
\pgfsetrectcap%
\pgfsetroundjoin%
\pgfsetlinewidth{1.003750pt}%
\definecolor{currentstroke}{rgb}{0.501961,0.803922,0.756863}%
\pgfsetstrokecolor{currentstroke}%
\pgfsetdash{}{0pt}%
\pgfpathmoveto{\pgfqpoint{2.108975in}{1.810784in}}%
\pgfpathlineto{\pgfqpoint{2.220086in}{1.810784in}}%
\pgfpathlineto{\pgfqpoint{2.331197in}{1.810784in}}%
\pgfusepath{stroke}%
\end{pgfscope}%
\begin{pgfscope}%
\definecolor{textcolor}{rgb}{0.000000,0.000000,0.000000}%
\pgfsetstrokecolor{textcolor}%
\pgfsetfillcolor{textcolor}%
\pgftext[x=2.420086in,y=1.771895in,left,base]{\color{textcolor}{\rmfamily\fontsize{8.000000}{9.600000}\selectfont\catcode`\^=\active\def^{\ifmmode\sp\else\^{}\fi}\catcode`\%=\active\def%{\%}planeTBL}}%
\end{pgfscope}%
\end{pgfpicture}%
\makeatother%
\endgroup%

    \label{fig:inflowSecVerifFluct}
\end{subfigure}
\\[-30pt]
\begin{subfigure}[t]{0.495\textwidth}
    \centering
    \subcaption{}
    %% Creator: Matplotlib, PGF backend
%%
%% To include the figure in your LaTeX document, write
%%   \input{<filename>.pgf}
%%
%% Make sure the required packages are loaded in your preamble
%%   \usepackage{pgf}
%%
%% Also ensure that all the required font packages are loaded; for instance,
%% the lmodern package is sometimes necessary when using math font.
%%   \usepackage{lmodern}
%%
%% Figures using additional raster images can only be included by \input if
%% they are in the same directory as the main LaTeX file. For loading figures
%% from other directories you can use the `import` package
%%   \usepackage{import}
%%
%% and then include the figures with
%%   \import{<path to file>}{<filename>.pgf}
%%
%% Matplotlib used the following preamble
%%   \def\mathdefault#1{#1}
%%   \everymath=\expandafter{\the\everymath\displaystyle}
%%   \IfFileExists{scrextend.sty}{
%%     \usepackage[fontsize=11.000000pt]{scrextend}
%%   }{
%%     \renewcommand{\normalsize}{\fontsize{11.000000}{13.200000}\selectfont}
%%     \normalsize
%%   }
%%   
%%   \makeatletter\@ifpackageloaded{underscore}{}{\usepackage[strings]{underscore}}\makeatother
%%
\begingroup%
\makeatletter%
\begin{pgfpicture}%
\pgfpathrectangle{\pgfpointorigin}{\pgfqpoint{3.118110in}{2.494488in}}%
\pgfusepath{use as bounding box, clip}%
\begin{pgfscope}%
\pgfsetbuttcap%
\pgfsetmiterjoin%
\definecolor{currentfill}{rgb}{1.000000,1.000000,1.000000}%
\pgfsetfillcolor{currentfill}%
\pgfsetlinewidth{0.000000pt}%
\definecolor{currentstroke}{rgb}{1.000000,1.000000,1.000000}%
\pgfsetstrokecolor{currentstroke}%
\pgfsetdash{}{0pt}%
\pgfpathmoveto{\pgfqpoint{0.000000in}{0.000000in}}%
\pgfpathlineto{\pgfqpoint{3.118110in}{0.000000in}}%
\pgfpathlineto{\pgfqpoint{3.118110in}{2.494488in}}%
\pgfpathlineto{\pgfqpoint{0.000000in}{2.494488in}}%
\pgfpathlineto{\pgfqpoint{0.000000in}{0.000000in}}%
\pgfpathclose%
\pgfusepath{fill}%
\end{pgfscope}%
\begin{pgfscope}%
\pgfsetbuttcap%
\pgfsetmiterjoin%
\definecolor{currentfill}{rgb}{1.000000,1.000000,1.000000}%
\pgfsetfillcolor{currentfill}%
\pgfsetlinewidth{0.000000pt}%
\definecolor{currentstroke}{rgb}{0.000000,0.000000,0.000000}%
\pgfsetstrokecolor{currentstroke}%
\pgfsetstrokeopacity{0.000000}%
\pgfsetdash{}{0pt}%
\pgfpathmoveto{\pgfqpoint{0.650000in}{0.420000in}}%
\pgfpathlineto{\pgfqpoint{3.038110in}{0.420000in}}%
\pgfpathlineto{\pgfqpoint{3.038110in}{2.414488in}}%
\pgfpathlineto{\pgfqpoint{0.650000in}{2.414488in}}%
\pgfpathlineto{\pgfqpoint{0.650000in}{0.420000in}}%
\pgfpathclose%
\pgfusepath{fill}%
\end{pgfscope}%
\begin{pgfscope}%
\pgfsetbuttcap%
\pgfsetroundjoin%
\definecolor{currentfill}{rgb}{0.000000,0.000000,0.000000}%
\pgfsetfillcolor{currentfill}%
\pgfsetlinewidth{0.501875pt}%
\definecolor{currentstroke}{rgb}{0.000000,0.000000,0.000000}%
\pgfsetstrokecolor{currentstroke}%
\pgfsetdash{}{0pt}%
\pgfsys@defobject{currentmarker}{\pgfqpoint{0.000000in}{0.000000in}}{\pgfqpoint{0.000000in}{0.041667in}}{%
\pgfpathmoveto{\pgfqpoint{0.000000in}{0.000000in}}%
\pgfpathlineto{\pgfqpoint{0.000000in}{0.041667in}}%
\pgfusepath{stroke,fill}%
}%
\begin{pgfscope}%
\pgfsys@transformshift{0.878501in}{0.420000in}%
\pgfsys@useobject{currentmarker}{}%
\end{pgfscope}%
\end{pgfscope}%
\begin{pgfscope}%
\pgfsetbuttcap%
\pgfsetroundjoin%
\definecolor{currentfill}{rgb}{0.000000,0.000000,0.000000}%
\pgfsetfillcolor{currentfill}%
\pgfsetlinewidth{0.501875pt}%
\definecolor{currentstroke}{rgb}{0.000000,0.000000,0.000000}%
\pgfsetstrokecolor{currentstroke}%
\pgfsetdash{}{0pt}%
\pgfsys@defobject{currentmarker}{\pgfqpoint{0.000000in}{-0.041667in}}{\pgfqpoint{0.000000in}{0.000000in}}{%
\pgfpathmoveto{\pgfqpoint{0.000000in}{0.000000in}}%
\pgfpathlineto{\pgfqpoint{0.000000in}{-0.041667in}}%
\pgfusepath{stroke,fill}%
}%
\begin{pgfscope}%
\pgfsys@transformshift{0.878501in}{2.414488in}%
\pgfsys@useobject{currentmarker}{}%
\end{pgfscope}%
\end{pgfscope}%
\begin{pgfscope}%
\definecolor{textcolor}{rgb}{0.000000,0.000000,0.000000}%
\pgfsetstrokecolor{textcolor}%
\pgfsetfillcolor{textcolor}%
\pgftext[x=0.878501in,y=0.371389in,,top]{\color{textcolor}{\rmfamily\fontsize{11.000000}{13.200000}\selectfont\catcode`\^=\active\def^{\ifmmode\sp\else\^{}\fi}\catcode`\%=\active\def%{\%}$\mathdefault{10^{0}}$}}%
\end{pgfscope}%
\begin{pgfscope}%
\pgfsetbuttcap%
\pgfsetroundjoin%
\definecolor{currentfill}{rgb}{0.000000,0.000000,0.000000}%
\pgfsetfillcolor{currentfill}%
\pgfsetlinewidth{0.501875pt}%
\definecolor{currentstroke}{rgb}{0.000000,0.000000,0.000000}%
\pgfsetstrokecolor{currentstroke}%
\pgfsetdash{}{0pt}%
\pgfsys@defobject{currentmarker}{\pgfqpoint{0.000000in}{0.000000in}}{\pgfqpoint{0.000000in}{0.041667in}}{%
\pgfpathmoveto{\pgfqpoint{0.000000in}{0.000000in}}%
\pgfpathlineto{\pgfqpoint{0.000000in}{0.041667in}}%
\pgfusepath{stroke,fill}%
}%
\begin{pgfscope}%
\pgfsys@transformshift{1.637564in}{0.420000in}%
\pgfsys@useobject{currentmarker}{}%
\end{pgfscope}%
\end{pgfscope}%
\begin{pgfscope}%
\pgfsetbuttcap%
\pgfsetroundjoin%
\definecolor{currentfill}{rgb}{0.000000,0.000000,0.000000}%
\pgfsetfillcolor{currentfill}%
\pgfsetlinewidth{0.501875pt}%
\definecolor{currentstroke}{rgb}{0.000000,0.000000,0.000000}%
\pgfsetstrokecolor{currentstroke}%
\pgfsetdash{}{0pt}%
\pgfsys@defobject{currentmarker}{\pgfqpoint{0.000000in}{-0.041667in}}{\pgfqpoint{0.000000in}{0.000000in}}{%
\pgfpathmoveto{\pgfqpoint{0.000000in}{0.000000in}}%
\pgfpathlineto{\pgfqpoint{0.000000in}{-0.041667in}}%
\pgfusepath{stroke,fill}%
}%
\begin{pgfscope}%
\pgfsys@transformshift{1.637564in}{2.414488in}%
\pgfsys@useobject{currentmarker}{}%
\end{pgfscope}%
\end{pgfscope}%
\begin{pgfscope}%
\definecolor{textcolor}{rgb}{0.000000,0.000000,0.000000}%
\pgfsetstrokecolor{textcolor}%
\pgfsetfillcolor{textcolor}%
\pgftext[x=1.637564in,y=0.371389in,,top]{\color{textcolor}{\rmfamily\fontsize{11.000000}{13.200000}\selectfont\catcode`\^=\active\def^{\ifmmode\sp\else\^{}\fi}\catcode`\%=\active\def%{\%}$\mathdefault{10^{1}}$}}%
\end{pgfscope}%
\begin{pgfscope}%
\pgfsetbuttcap%
\pgfsetroundjoin%
\definecolor{currentfill}{rgb}{0.000000,0.000000,0.000000}%
\pgfsetfillcolor{currentfill}%
\pgfsetlinewidth{0.501875pt}%
\definecolor{currentstroke}{rgb}{0.000000,0.000000,0.000000}%
\pgfsetstrokecolor{currentstroke}%
\pgfsetdash{}{0pt}%
\pgfsys@defobject{currentmarker}{\pgfqpoint{0.000000in}{0.000000in}}{\pgfqpoint{0.000000in}{0.041667in}}{%
\pgfpathmoveto{\pgfqpoint{0.000000in}{0.000000in}}%
\pgfpathlineto{\pgfqpoint{0.000000in}{0.041667in}}%
\pgfusepath{stroke,fill}%
}%
\begin{pgfscope}%
\pgfsys@transformshift{2.396627in}{0.420000in}%
\pgfsys@useobject{currentmarker}{}%
\end{pgfscope}%
\end{pgfscope}%
\begin{pgfscope}%
\pgfsetbuttcap%
\pgfsetroundjoin%
\definecolor{currentfill}{rgb}{0.000000,0.000000,0.000000}%
\pgfsetfillcolor{currentfill}%
\pgfsetlinewidth{0.501875pt}%
\definecolor{currentstroke}{rgb}{0.000000,0.000000,0.000000}%
\pgfsetstrokecolor{currentstroke}%
\pgfsetdash{}{0pt}%
\pgfsys@defobject{currentmarker}{\pgfqpoint{0.000000in}{-0.041667in}}{\pgfqpoint{0.000000in}{0.000000in}}{%
\pgfpathmoveto{\pgfqpoint{0.000000in}{0.000000in}}%
\pgfpathlineto{\pgfqpoint{0.000000in}{-0.041667in}}%
\pgfusepath{stroke,fill}%
}%
\begin{pgfscope}%
\pgfsys@transformshift{2.396627in}{2.414488in}%
\pgfsys@useobject{currentmarker}{}%
\end{pgfscope}%
\end{pgfscope}%
\begin{pgfscope}%
\definecolor{textcolor}{rgb}{0.000000,0.000000,0.000000}%
\pgfsetstrokecolor{textcolor}%
\pgfsetfillcolor{textcolor}%
\pgftext[x=2.396627in,y=0.371389in,,top]{\color{textcolor}{\rmfamily\fontsize{11.000000}{13.200000}\selectfont\catcode`\^=\active\def^{\ifmmode\sp\else\^{}\fi}\catcode`\%=\active\def%{\%}$\mathdefault{10^{2}}$}}%
\end{pgfscope}%
\begin{pgfscope}%
\pgfsetbuttcap%
\pgfsetroundjoin%
\definecolor{currentfill}{rgb}{0.000000,0.000000,0.000000}%
\pgfsetfillcolor{currentfill}%
\pgfsetlinewidth{0.401500pt}%
\definecolor{currentstroke}{rgb}{0.000000,0.000000,0.000000}%
\pgfsetstrokecolor{currentstroke}%
\pgfsetdash{}{0pt}%
\pgfsys@defobject{currentmarker}{\pgfqpoint{0.000000in}{0.000000in}}{\pgfqpoint{0.000000in}{0.020833in}}{%
\pgfpathmoveto{\pgfqpoint{0.000000in}{0.000000in}}%
\pgfpathlineto{\pgfqpoint{0.000000in}{0.020833in}}%
\pgfusepath{stroke,fill}%
}%
\begin{pgfscope}%
\pgfsys@transformshift{0.650000in}{0.420000in}%
\pgfsys@useobject{currentmarker}{}%
\end{pgfscope}%
\end{pgfscope}%
\begin{pgfscope}%
\pgfsetbuttcap%
\pgfsetroundjoin%
\definecolor{currentfill}{rgb}{0.000000,0.000000,0.000000}%
\pgfsetfillcolor{currentfill}%
\pgfsetlinewidth{0.401500pt}%
\definecolor{currentstroke}{rgb}{0.000000,0.000000,0.000000}%
\pgfsetstrokecolor{currentstroke}%
\pgfsetdash{}{0pt}%
\pgfsys@defobject{currentmarker}{\pgfqpoint{0.000000in}{-0.020833in}}{\pgfqpoint{0.000000in}{0.000000in}}{%
\pgfpathmoveto{\pgfqpoint{0.000000in}{0.000000in}}%
\pgfpathlineto{\pgfqpoint{0.000000in}{-0.020833in}}%
\pgfusepath{stroke,fill}%
}%
\begin{pgfscope}%
\pgfsys@transformshift{0.650000in}{2.414488in}%
\pgfsys@useobject{currentmarker}{}%
\end{pgfscope}%
\end{pgfscope}%
\begin{pgfscope}%
\pgfsetbuttcap%
\pgfsetroundjoin%
\definecolor{currentfill}{rgb}{0.000000,0.000000,0.000000}%
\pgfsetfillcolor{currentfill}%
\pgfsetlinewidth{0.401500pt}%
\definecolor{currentstroke}{rgb}{0.000000,0.000000,0.000000}%
\pgfsetstrokecolor{currentstroke}%
\pgfsetdash{}{0pt}%
\pgfsys@defobject{currentmarker}{\pgfqpoint{0.000000in}{0.000000in}}{\pgfqpoint{0.000000in}{0.020833in}}{%
\pgfpathmoveto{\pgfqpoint{0.000000in}{0.000000in}}%
\pgfpathlineto{\pgfqpoint{0.000000in}{0.020833in}}%
\pgfusepath{stroke,fill}%
}%
\begin{pgfscope}%
\pgfsys@transformshift{0.710104in}{0.420000in}%
\pgfsys@useobject{currentmarker}{}%
\end{pgfscope}%
\end{pgfscope}%
\begin{pgfscope}%
\pgfsetbuttcap%
\pgfsetroundjoin%
\definecolor{currentfill}{rgb}{0.000000,0.000000,0.000000}%
\pgfsetfillcolor{currentfill}%
\pgfsetlinewidth{0.401500pt}%
\definecolor{currentstroke}{rgb}{0.000000,0.000000,0.000000}%
\pgfsetstrokecolor{currentstroke}%
\pgfsetdash{}{0pt}%
\pgfsys@defobject{currentmarker}{\pgfqpoint{0.000000in}{-0.020833in}}{\pgfqpoint{0.000000in}{0.000000in}}{%
\pgfpathmoveto{\pgfqpoint{0.000000in}{0.000000in}}%
\pgfpathlineto{\pgfqpoint{0.000000in}{-0.020833in}}%
\pgfusepath{stroke,fill}%
}%
\begin{pgfscope}%
\pgfsys@transformshift{0.710104in}{2.414488in}%
\pgfsys@useobject{currentmarker}{}%
\end{pgfscope}%
\end{pgfscope}%
\begin{pgfscope}%
\pgfsetbuttcap%
\pgfsetroundjoin%
\definecolor{currentfill}{rgb}{0.000000,0.000000,0.000000}%
\pgfsetfillcolor{currentfill}%
\pgfsetlinewidth{0.401500pt}%
\definecolor{currentstroke}{rgb}{0.000000,0.000000,0.000000}%
\pgfsetstrokecolor{currentstroke}%
\pgfsetdash{}{0pt}%
\pgfsys@defobject{currentmarker}{\pgfqpoint{0.000000in}{0.000000in}}{\pgfqpoint{0.000000in}{0.020833in}}{%
\pgfpathmoveto{\pgfqpoint{0.000000in}{0.000000in}}%
\pgfpathlineto{\pgfqpoint{0.000000in}{0.020833in}}%
\pgfusepath{stroke,fill}%
}%
\begin{pgfscope}%
\pgfsys@transformshift{0.760920in}{0.420000in}%
\pgfsys@useobject{currentmarker}{}%
\end{pgfscope}%
\end{pgfscope}%
\begin{pgfscope}%
\pgfsetbuttcap%
\pgfsetroundjoin%
\definecolor{currentfill}{rgb}{0.000000,0.000000,0.000000}%
\pgfsetfillcolor{currentfill}%
\pgfsetlinewidth{0.401500pt}%
\definecolor{currentstroke}{rgb}{0.000000,0.000000,0.000000}%
\pgfsetstrokecolor{currentstroke}%
\pgfsetdash{}{0pt}%
\pgfsys@defobject{currentmarker}{\pgfqpoint{0.000000in}{-0.020833in}}{\pgfqpoint{0.000000in}{0.000000in}}{%
\pgfpathmoveto{\pgfqpoint{0.000000in}{0.000000in}}%
\pgfpathlineto{\pgfqpoint{0.000000in}{-0.020833in}}%
\pgfusepath{stroke,fill}%
}%
\begin{pgfscope}%
\pgfsys@transformshift{0.760920in}{2.414488in}%
\pgfsys@useobject{currentmarker}{}%
\end{pgfscope}%
\end{pgfscope}%
\begin{pgfscope}%
\pgfsetbuttcap%
\pgfsetroundjoin%
\definecolor{currentfill}{rgb}{0.000000,0.000000,0.000000}%
\pgfsetfillcolor{currentfill}%
\pgfsetlinewidth{0.401500pt}%
\definecolor{currentstroke}{rgb}{0.000000,0.000000,0.000000}%
\pgfsetstrokecolor{currentstroke}%
\pgfsetdash{}{0pt}%
\pgfsys@defobject{currentmarker}{\pgfqpoint{0.000000in}{0.000000in}}{\pgfqpoint{0.000000in}{0.020833in}}{%
\pgfpathmoveto{\pgfqpoint{0.000000in}{0.000000in}}%
\pgfpathlineto{\pgfqpoint{0.000000in}{0.020833in}}%
\pgfusepath{stroke,fill}%
}%
\begin{pgfscope}%
\pgfsys@transformshift{0.804940in}{0.420000in}%
\pgfsys@useobject{currentmarker}{}%
\end{pgfscope}%
\end{pgfscope}%
\begin{pgfscope}%
\pgfsetbuttcap%
\pgfsetroundjoin%
\definecolor{currentfill}{rgb}{0.000000,0.000000,0.000000}%
\pgfsetfillcolor{currentfill}%
\pgfsetlinewidth{0.401500pt}%
\definecolor{currentstroke}{rgb}{0.000000,0.000000,0.000000}%
\pgfsetstrokecolor{currentstroke}%
\pgfsetdash{}{0pt}%
\pgfsys@defobject{currentmarker}{\pgfqpoint{0.000000in}{-0.020833in}}{\pgfqpoint{0.000000in}{0.000000in}}{%
\pgfpathmoveto{\pgfqpoint{0.000000in}{0.000000in}}%
\pgfpathlineto{\pgfqpoint{0.000000in}{-0.020833in}}%
\pgfusepath{stroke,fill}%
}%
\begin{pgfscope}%
\pgfsys@transformshift{0.804940in}{2.414488in}%
\pgfsys@useobject{currentmarker}{}%
\end{pgfscope}%
\end{pgfscope}%
\begin{pgfscope}%
\pgfsetbuttcap%
\pgfsetroundjoin%
\definecolor{currentfill}{rgb}{0.000000,0.000000,0.000000}%
\pgfsetfillcolor{currentfill}%
\pgfsetlinewidth{0.401500pt}%
\definecolor{currentstroke}{rgb}{0.000000,0.000000,0.000000}%
\pgfsetstrokecolor{currentstroke}%
\pgfsetdash{}{0pt}%
\pgfsys@defobject{currentmarker}{\pgfqpoint{0.000000in}{0.000000in}}{\pgfqpoint{0.000000in}{0.020833in}}{%
\pgfpathmoveto{\pgfqpoint{0.000000in}{0.000000in}}%
\pgfpathlineto{\pgfqpoint{0.000000in}{0.020833in}}%
\pgfusepath{stroke,fill}%
}%
\begin{pgfscope}%
\pgfsys@transformshift{0.843768in}{0.420000in}%
\pgfsys@useobject{currentmarker}{}%
\end{pgfscope}%
\end{pgfscope}%
\begin{pgfscope}%
\pgfsetbuttcap%
\pgfsetroundjoin%
\definecolor{currentfill}{rgb}{0.000000,0.000000,0.000000}%
\pgfsetfillcolor{currentfill}%
\pgfsetlinewidth{0.401500pt}%
\definecolor{currentstroke}{rgb}{0.000000,0.000000,0.000000}%
\pgfsetstrokecolor{currentstroke}%
\pgfsetdash{}{0pt}%
\pgfsys@defobject{currentmarker}{\pgfqpoint{0.000000in}{-0.020833in}}{\pgfqpoint{0.000000in}{0.000000in}}{%
\pgfpathmoveto{\pgfqpoint{0.000000in}{0.000000in}}%
\pgfpathlineto{\pgfqpoint{0.000000in}{-0.020833in}}%
\pgfusepath{stroke,fill}%
}%
\begin{pgfscope}%
\pgfsys@transformshift{0.843768in}{2.414488in}%
\pgfsys@useobject{currentmarker}{}%
\end{pgfscope}%
\end{pgfscope}%
\begin{pgfscope}%
\pgfsetbuttcap%
\pgfsetroundjoin%
\definecolor{currentfill}{rgb}{0.000000,0.000000,0.000000}%
\pgfsetfillcolor{currentfill}%
\pgfsetlinewidth{0.401500pt}%
\definecolor{currentstroke}{rgb}{0.000000,0.000000,0.000000}%
\pgfsetstrokecolor{currentstroke}%
\pgfsetdash{}{0pt}%
\pgfsys@defobject{currentmarker}{\pgfqpoint{0.000000in}{0.000000in}}{\pgfqpoint{0.000000in}{0.020833in}}{%
\pgfpathmoveto{\pgfqpoint{0.000000in}{0.000000in}}%
\pgfpathlineto{\pgfqpoint{0.000000in}{0.020833in}}%
\pgfusepath{stroke,fill}%
}%
\begin{pgfscope}%
\pgfsys@transformshift{1.107002in}{0.420000in}%
\pgfsys@useobject{currentmarker}{}%
\end{pgfscope}%
\end{pgfscope}%
\begin{pgfscope}%
\pgfsetbuttcap%
\pgfsetroundjoin%
\definecolor{currentfill}{rgb}{0.000000,0.000000,0.000000}%
\pgfsetfillcolor{currentfill}%
\pgfsetlinewidth{0.401500pt}%
\definecolor{currentstroke}{rgb}{0.000000,0.000000,0.000000}%
\pgfsetstrokecolor{currentstroke}%
\pgfsetdash{}{0pt}%
\pgfsys@defobject{currentmarker}{\pgfqpoint{0.000000in}{-0.020833in}}{\pgfqpoint{0.000000in}{0.000000in}}{%
\pgfpathmoveto{\pgfqpoint{0.000000in}{0.000000in}}%
\pgfpathlineto{\pgfqpoint{0.000000in}{-0.020833in}}%
\pgfusepath{stroke,fill}%
}%
\begin{pgfscope}%
\pgfsys@transformshift{1.107002in}{2.414488in}%
\pgfsys@useobject{currentmarker}{}%
\end{pgfscope}%
\end{pgfscope}%
\begin{pgfscope}%
\pgfsetbuttcap%
\pgfsetroundjoin%
\definecolor{currentfill}{rgb}{0.000000,0.000000,0.000000}%
\pgfsetfillcolor{currentfill}%
\pgfsetlinewidth{0.401500pt}%
\definecolor{currentstroke}{rgb}{0.000000,0.000000,0.000000}%
\pgfsetstrokecolor{currentstroke}%
\pgfsetdash{}{0pt}%
\pgfsys@defobject{currentmarker}{\pgfqpoint{0.000000in}{0.000000in}}{\pgfqpoint{0.000000in}{0.020833in}}{%
\pgfpathmoveto{\pgfqpoint{0.000000in}{0.000000in}}%
\pgfpathlineto{\pgfqpoint{0.000000in}{0.020833in}}%
\pgfusepath{stroke,fill}%
}%
\begin{pgfscope}%
\pgfsys@transformshift{1.240666in}{0.420000in}%
\pgfsys@useobject{currentmarker}{}%
\end{pgfscope}%
\end{pgfscope}%
\begin{pgfscope}%
\pgfsetbuttcap%
\pgfsetroundjoin%
\definecolor{currentfill}{rgb}{0.000000,0.000000,0.000000}%
\pgfsetfillcolor{currentfill}%
\pgfsetlinewidth{0.401500pt}%
\definecolor{currentstroke}{rgb}{0.000000,0.000000,0.000000}%
\pgfsetstrokecolor{currentstroke}%
\pgfsetdash{}{0pt}%
\pgfsys@defobject{currentmarker}{\pgfqpoint{0.000000in}{-0.020833in}}{\pgfqpoint{0.000000in}{0.000000in}}{%
\pgfpathmoveto{\pgfqpoint{0.000000in}{0.000000in}}%
\pgfpathlineto{\pgfqpoint{0.000000in}{-0.020833in}}%
\pgfusepath{stroke,fill}%
}%
\begin{pgfscope}%
\pgfsys@transformshift{1.240666in}{2.414488in}%
\pgfsys@useobject{currentmarker}{}%
\end{pgfscope}%
\end{pgfscope}%
\begin{pgfscope}%
\pgfsetbuttcap%
\pgfsetroundjoin%
\definecolor{currentfill}{rgb}{0.000000,0.000000,0.000000}%
\pgfsetfillcolor{currentfill}%
\pgfsetlinewidth{0.401500pt}%
\definecolor{currentstroke}{rgb}{0.000000,0.000000,0.000000}%
\pgfsetstrokecolor{currentstroke}%
\pgfsetdash{}{0pt}%
\pgfsys@defobject{currentmarker}{\pgfqpoint{0.000000in}{0.000000in}}{\pgfqpoint{0.000000in}{0.020833in}}{%
\pgfpathmoveto{\pgfqpoint{0.000000in}{0.000000in}}%
\pgfpathlineto{\pgfqpoint{0.000000in}{0.020833in}}%
\pgfusepath{stroke,fill}%
}%
\begin{pgfscope}%
\pgfsys@transformshift{1.335502in}{0.420000in}%
\pgfsys@useobject{currentmarker}{}%
\end{pgfscope}%
\end{pgfscope}%
\begin{pgfscope}%
\pgfsetbuttcap%
\pgfsetroundjoin%
\definecolor{currentfill}{rgb}{0.000000,0.000000,0.000000}%
\pgfsetfillcolor{currentfill}%
\pgfsetlinewidth{0.401500pt}%
\definecolor{currentstroke}{rgb}{0.000000,0.000000,0.000000}%
\pgfsetstrokecolor{currentstroke}%
\pgfsetdash{}{0pt}%
\pgfsys@defobject{currentmarker}{\pgfqpoint{0.000000in}{-0.020833in}}{\pgfqpoint{0.000000in}{0.000000in}}{%
\pgfpathmoveto{\pgfqpoint{0.000000in}{0.000000in}}%
\pgfpathlineto{\pgfqpoint{0.000000in}{-0.020833in}}%
\pgfusepath{stroke,fill}%
}%
\begin{pgfscope}%
\pgfsys@transformshift{1.335502in}{2.414488in}%
\pgfsys@useobject{currentmarker}{}%
\end{pgfscope}%
\end{pgfscope}%
\begin{pgfscope}%
\pgfsetbuttcap%
\pgfsetroundjoin%
\definecolor{currentfill}{rgb}{0.000000,0.000000,0.000000}%
\pgfsetfillcolor{currentfill}%
\pgfsetlinewidth{0.401500pt}%
\definecolor{currentstroke}{rgb}{0.000000,0.000000,0.000000}%
\pgfsetstrokecolor{currentstroke}%
\pgfsetdash{}{0pt}%
\pgfsys@defobject{currentmarker}{\pgfqpoint{0.000000in}{0.000000in}}{\pgfqpoint{0.000000in}{0.020833in}}{%
\pgfpathmoveto{\pgfqpoint{0.000000in}{0.000000in}}%
\pgfpathlineto{\pgfqpoint{0.000000in}{0.020833in}}%
\pgfusepath{stroke,fill}%
}%
\begin{pgfscope}%
\pgfsys@transformshift{1.409063in}{0.420000in}%
\pgfsys@useobject{currentmarker}{}%
\end{pgfscope}%
\end{pgfscope}%
\begin{pgfscope}%
\pgfsetbuttcap%
\pgfsetroundjoin%
\definecolor{currentfill}{rgb}{0.000000,0.000000,0.000000}%
\pgfsetfillcolor{currentfill}%
\pgfsetlinewidth{0.401500pt}%
\definecolor{currentstroke}{rgb}{0.000000,0.000000,0.000000}%
\pgfsetstrokecolor{currentstroke}%
\pgfsetdash{}{0pt}%
\pgfsys@defobject{currentmarker}{\pgfqpoint{0.000000in}{-0.020833in}}{\pgfqpoint{0.000000in}{0.000000in}}{%
\pgfpathmoveto{\pgfqpoint{0.000000in}{0.000000in}}%
\pgfpathlineto{\pgfqpoint{0.000000in}{-0.020833in}}%
\pgfusepath{stroke,fill}%
}%
\begin{pgfscope}%
\pgfsys@transformshift{1.409063in}{2.414488in}%
\pgfsys@useobject{currentmarker}{}%
\end{pgfscope}%
\end{pgfscope}%
\begin{pgfscope}%
\pgfsetbuttcap%
\pgfsetroundjoin%
\definecolor{currentfill}{rgb}{0.000000,0.000000,0.000000}%
\pgfsetfillcolor{currentfill}%
\pgfsetlinewidth{0.401500pt}%
\definecolor{currentstroke}{rgb}{0.000000,0.000000,0.000000}%
\pgfsetstrokecolor{currentstroke}%
\pgfsetdash{}{0pt}%
\pgfsys@defobject{currentmarker}{\pgfqpoint{0.000000in}{0.000000in}}{\pgfqpoint{0.000000in}{0.020833in}}{%
\pgfpathmoveto{\pgfqpoint{0.000000in}{0.000000in}}%
\pgfpathlineto{\pgfqpoint{0.000000in}{0.020833in}}%
\pgfusepath{stroke,fill}%
}%
\begin{pgfscope}%
\pgfsys@transformshift{1.469167in}{0.420000in}%
\pgfsys@useobject{currentmarker}{}%
\end{pgfscope}%
\end{pgfscope}%
\begin{pgfscope}%
\pgfsetbuttcap%
\pgfsetroundjoin%
\definecolor{currentfill}{rgb}{0.000000,0.000000,0.000000}%
\pgfsetfillcolor{currentfill}%
\pgfsetlinewidth{0.401500pt}%
\definecolor{currentstroke}{rgb}{0.000000,0.000000,0.000000}%
\pgfsetstrokecolor{currentstroke}%
\pgfsetdash{}{0pt}%
\pgfsys@defobject{currentmarker}{\pgfqpoint{0.000000in}{-0.020833in}}{\pgfqpoint{0.000000in}{0.000000in}}{%
\pgfpathmoveto{\pgfqpoint{0.000000in}{0.000000in}}%
\pgfpathlineto{\pgfqpoint{0.000000in}{-0.020833in}}%
\pgfusepath{stroke,fill}%
}%
\begin{pgfscope}%
\pgfsys@transformshift{1.469167in}{2.414488in}%
\pgfsys@useobject{currentmarker}{}%
\end{pgfscope}%
\end{pgfscope}%
\begin{pgfscope}%
\pgfsetbuttcap%
\pgfsetroundjoin%
\definecolor{currentfill}{rgb}{0.000000,0.000000,0.000000}%
\pgfsetfillcolor{currentfill}%
\pgfsetlinewidth{0.401500pt}%
\definecolor{currentstroke}{rgb}{0.000000,0.000000,0.000000}%
\pgfsetstrokecolor{currentstroke}%
\pgfsetdash{}{0pt}%
\pgfsys@defobject{currentmarker}{\pgfqpoint{0.000000in}{0.000000in}}{\pgfqpoint{0.000000in}{0.020833in}}{%
\pgfpathmoveto{\pgfqpoint{0.000000in}{0.000000in}}%
\pgfpathlineto{\pgfqpoint{0.000000in}{0.020833in}}%
\pgfusepath{stroke,fill}%
}%
\begin{pgfscope}%
\pgfsys@transformshift{1.519984in}{0.420000in}%
\pgfsys@useobject{currentmarker}{}%
\end{pgfscope}%
\end{pgfscope}%
\begin{pgfscope}%
\pgfsetbuttcap%
\pgfsetroundjoin%
\definecolor{currentfill}{rgb}{0.000000,0.000000,0.000000}%
\pgfsetfillcolor{currentfill}%
\pgfsetlinewidth{0.401500pt}%
\definecolor{currentstroke}{rgb}{0.000000,0.000000,0.000000}%
\pgfsetstrokecolor{currentstroke}%
\pgfsetdash{}{0pt}%
\pgfsys@defobject{currentmarker}{\pgfqpoint{0.000000in}{-0.020833in}}{\pgfqpoint{0.000000in}{0.000000in}}{%
\pgfpathmoveto{\pgfqpoint{0.000000in}{0.000000in}}%
\pgfpathlineto{\pgfqpoint{0.000000in}{-0.020833in}}%
\pgfusepath{stroke,fill}%
}%
\begin{pgfscope}%
\pgfsys@transformshift{1.519984in}{2.414488in}%
\pgfsys@useobject{currentmarker}{}%
\end{pgfscope}%
\end{pgfscope}%
\begin{pgfscope}%
\pgfsetbuttcap%
\pgfsetroundjoin%
\definecolor{currentfill}{rgb}{0.000000,0.000000,0.000000}%
\pgfsetfillcolor{currentfill}%
\pgfsetlinewidth{0.401500pt}%
\definecolor{currentstroke}{rgb}{0.000000,0.000000,0.000000}%
\pgfsetstrokecolor{currentstroke}%
\pgfsetdash{}{0pt}%
\pgfsys@defobject{currentmarker}{\pgfqpoint{0.000000in}{0.000000in}}{\pgfqpoint{0.000000in}{0.020833in}}{%
\pgfpathmoveto{\pgfqpoint{0.000000in}{0.000000in}}%
\pgfpathlineto{\pgfqpoint{0.000000in}{0.020833in}}%
\pgfusepath{stroke,fill}%
}%
\begin{pgfscope}%
\pgfsys@transformshift{1.564003in}{0.420000in}%
\pgfsys@useobject{currentmarker}{}%
\end{pgfscope}%
\end{pgfscope}%
\begin{pgfscope}%
\pgfsetbuttcap%
\pgfsetroundjoin%
\definecolor{currentfill}{rgb}{0.000000,0.000000,0.000000}%
\pgfsetfillcolor{currentfill}%
\pgfsetlinewidth{0.401500pt}%
\definecolor{currentstroke}{rgb}{0.000000,0.000000,0.000000}%
\pgfsetstrokecolor{currentstroke}%
\pgfsetdash{}{0pt}%
\pgfsys@defobject{currentmarker}{\pgfqpoint{0.000000in}{-0.020833in}}{\pgfqpoint{0.000000in}{0.000000in}}{%
\pgfpathmoveto{\pgfqpoint{0.000000in}{0.000000in}}%
\pgfpathlineto{\pgfqpoint{0.000000in}{-0.020833in}}%
\pgfusepath{stroke,fill}%
}%
\begin{pgfscope}%
\pgfsys@transformshift{1.564003in}{2.414488in}%
\pgfsys@useobject{currentmarker}{}%
\end{pgfscope}%
\end{pgfscope}%
\begin{pgfscope}%
\pgfsetbuttcap%
\pgfsetroundjoin%
\definecolor{currentfill}{rgb}{0.000000,0.000000,0.000000}%
\pgfsetfillcolor{currentfill}%
\pgfsetlinewidth{0.401500pt}%
\definecolor{currentstroke}{rgb}{0.000000,0.000000,0.000000}%
\pgfsetstrokecolor{currentstroke}%
\pgfsetdash{}{0pt}%
\pgfsys@defobject{currentmarker}{\pgfqpoint{0.000000in}{0.000000in}}{\pgfqpoint{0.000000in}{0.020833in}}{%
\pgfpathmoveto{\pgfqpoint{0.000000in}{0.000000in}}%
\pgfpathlineto{\pgfqpoint{0.000000in}{0.020833in}}%
\pgfusepath{stroke,fill}%
}%
\begin{pgfscope}%
\pgfsys@transformshift{1.602831in}{0.420000in}%
\pgfsys@useobject{currentmarker}{}%
\end{pgfscope}%
\end{pgfscope}%
\begin{pgfscope}%
\pgfsetbuttcap%
\pgfsetroundjoin%
\definecolor{currentfill}{rgb}{0.000000,0.000000,0.000000}%
\pgfsetfillcolor{currentfill}%
\pgfsetlinewidth{0.401500pt}%
\definecolor{currentstroke}{rgb}{0.000000,0.000000,0.000000}%
\pgfsetstrokecolor{currentstroke}%
\pgfsetdash{}{0pt}%
\pgfsys@defobject{currentmarker}{\pgfqpoint{0.000000in}{-0.020833in}}{\pgfqpoint{0.000000in}{0.000000in}}{%
\pgfpathmoveto{\pgfqpoint{0.000000in}{0.000000in}}%
\pgfpathlineto{\pgfqpoint{0.000000in}{-0.020833in}}%
\pgfusepath{stroke,fill}%
}%
\begin{pgfscope}%
\pgfsys@transformshift{1.602831in}{2.414488in}%
\pgfsys@useobject{currentmarker}{}%
\end{pgfscope}%
\end{pgfscope}%
\begin{pgfscope}%
\pgfsetbuttcap%
\pgfsetroundjoin%
\definecolor{currentfill}{rgb}{0.000000,0.000000,0.000000}%
\pgfsetfillcolor{currentfill}%
\pgfsetlinewidth{0.401500pt}%
\definecolor{currentstroke}{rgb}{0.000000,0.000000,0.000000}%
\pgfsetstrokecolor{currentstroke}%
\pgfsetdash{}{0pt}%
\pgfsys@defobject{currentmarker}{\pgfqpoint{0.000000in}{0.000000in}}{\pgfqpoint{0.000000in}{0.020833in}}{%
\pgfpathmoveto{\pgfqpoint{0.000000in}{0.000000in}}%
\pgfpathlineto{\pgfqpoint{0.000000in}{0.020833in}}%
\pgfusepath{stroke,fill}%
}%
\begin{pgfscope}%
\pgfsys@transformshift{1.866065in}{0.420000in}%
\pgfsys@useobject{currentmarker}{}%
\end{pgfscope}%
\end{pgfscope}%
\begin{pgfscope}%
\pgfsetbuttcap%
\pgfsetroundjoin%
\definecolor{currentfill}{rgb}{0.000000,0.000000,0.000000}%
\pgfsetfillcolor{currentfill}%
\pgfsetlinewidth{0.401500pt}%
\definecolor{currentstroke}{rgb}{0.000000,0.000000,0.000000}%
\pgfsetstrokecolor{currentstroke}%
\pgfsetdash{}{0pt}%
\pgfsys@defobject{currentmarker}{\pgfqpoint{0.000000in}{-0.020833in}}{\pgfqpoint{0.000000in}{0.000000in}}{%
\pgfpathmoveto{\pgfqpoint{0.000000in}{0.000000in}}%
\pgfpathlineto{\pgfqpoint{0.000000in}{-0.020833in}}%
\pgfusepath{stroke,fill}%
}%
\begin{pgfscope}%
\pgfsys@transformshift{1.866065in}{2.414488in}%
\pgfsys@useobject{currentmarker}{}%
\end{pgfscope}%
\end{pgfscope}%
\begin{pgfscope}%
\pgfsetbuttcap%
\pgfsetroundjoin%
\definecolor{currentfill}{rgb}{0.000000,0.000000,0.000000}%
\pgfsetfillcolor{currentfill}%
\pgfsetlinewidth{0.401500pt}%
\definecolor{currentstroke}{rgb}{0.000000,0.000000,0.000000}%
\pgfsetstrokecolor{currentstroke}%
\pgfsetdash{}{0pt}%
\pgfsys@defobject{currentmarker}{\pgfqpoint{0.000000in}{0.000000in}}{\pgfqpoint{0.000000in}{0.020833in}}{%
\pgfpathmoveto{\pgfqpoint{0.000000in}{0.000000in}}%
\pgfpathlineto{\pgfqpoint{0.000000in}{0.020833in}}%
\pgfusepath{stroke,fill}%
}%
\begin{pgfscope}%
\pgfsys@transformshift{1.999729in}{0.420000in}%
\pgfsys@useobject{currentmarker}{}%
\end{pgfscope}%
\end{pgfscope}%
\begin{pgfscope}%
\pgfsetbuttcap%
\pgfsetroundjoin%
\definecolor{currentfill}{rgb}{0.000000,0.000000,0.000000}%
\pgfsetfillcolor{currentfill}%
\pgfsetlinewidth{0.401500pt}%
\definecolor{currentstroke}{rgb}{0.000000,0.000000,0.000000}%
\pgfsetstrokecolor{currentstroke}%
\pgfsetdash{}{0pt}%
\pgfsys@defobject{currentmarker}{\pgfqpoint{0.000000in}{-0.020833in}}{\pgfqpoint{0.000000in}{0.000000in}}{%
\pgfpathmoveto{\pgfqpoint{0.000000in}{0.000000in}}%
\pgfpathlineto{\pgfqpoint{0.000000in}{-0.020833in}}%
\pgfusepath{stroke,fill}%
}%
\begin{pgfscope}%
\pgfsys@transformshift{1.999729in}{2.414488in}%
\pgfsys@useobject{currentmarker}{}%
\end{pgfscope}%
\end{pgfscope}%
\begin{pgfscope}%
\pgfsetbuttcap%
\pgfsetroundjoin%
\definecolor{currentfill}{rgb}{0.000000,0.000000,0.000000}%
\pgfsetfillcolor{currentfill}%
\pgfsetlinewidth{0.401500pt}%
\definecolor{currentstroke}{rgb}{0.000000,0.000000,0.000000}%
\pgfsetstrokecolor{currentstroke}%
\pgfsetdash{}{0pt}%
\pgfsys@defobject{currentmarker}{\pgfqpoint{0.000000in}{0.000000in}}{\pgfqpoint{0.000000in}{0.020833in}}{%
\pgfpathmoveto{\pgfqpoint{0.000000in}{0.000000in}}%
\pgfpathlineto{\pgfqpoint{0.000000in}{0.020833in}}%
\pgfusepath{stroke,fill}%
}%
\begin{pgfscope}%
\pgfsys@transformshift{2.094566in}{0.420000in}%
\pgfsys@useobject{currentmarker}{}%
\end{pgfscope}%
\end{pgfscope}%
\begin{pgfscope}%
\pgfsetbuttcap%
\pgfsetroundjoin%
\definecolor{currentfill}{rgb}{0.000000,0.000000,0.000000}%
\pgfsetfillcolor{currentfill}%
\pgfsetlinewidth{0.401500pt}%
\definecolor{currentstroke}{rgb}{0.000000,0.000000,0.000000}%
\pgfsetstrokecolor{currentstroke}%
\pgfsetdash{}{0pt}%
\pgfsys@defobject{currentmarker}{\pgfqpoint{0.000000in}{-0.020833in}}{\pgfqpoint{0.000000in}{0.000000in}}{%
\pgfpathmoveto{\pgfqpoint{0.000000in}{0.000000in}}%
\pgfpathlineto{\pgfqpoint{0.000000in}{-0.020833in}}%
\pgfusepath{stroke,fill}%
}%
\begin{pgfscope}%
\pgfsys@transformshift{2.094566in}{2.414488in}%
\pgfsys@useobject{currentmarker}{}%
\end{pgfscope}%
\end{pgfscope}%
\begin{pgfscope}%
\pgfsetbuttcap%
\pgfsetroundjoin%
\definecolor{currentfill}{rgb}{0.000000,0.000000,0.000000}%
\pgfsetfillcolor{currentfill}%
\pgfsetlinewidth{0.401500pt}%
\definecolor{currentstroke}{rgb}{0.000000,0.000000,0.000000}%
\pgfsetstrokecolor{currentstroke}%
\pgfsetdash{}{0pt}%
\pgfsys@defobject{currentmarker}{\pgfqpoint{0.000000in}{0.000000in}}{\pgfqpoint{0.000000in}{0.020833in}}{%
\pgfpathmoveto{\pgfqpoint{0.000000in}{0.000000in}}%
\pgfpathlineto{\pgfqpoint{0.000000in}{0.020833in}}%
\pgfusepath{stroke,fill}%
}%
\begin{pgfscope}%
\pgfsys@transformshift{2.168126in}{0.420000in}%
\pgfsys@useobject{currentmarker}{}%
\end{pgfscope}%
\end{pgfscope}%
\begin{pgfscope}%
\pgfsetbuttcap%
\pgfsetroundjoin%
\definecolor{currentfill}{rgb}{0.000000,0.000000,0.000000}%
\pgfsetfillcolor{currentfill}%
\pgfsetlinewidth{0.401500pt}%
\definecolor{currentstroke}{rgb}{0.000000,0.000000,0.000000}%
\pgfsetstrokecolor{currentstroke}%
\pgfsetdash{}{0pt}%
\pgfsys@defobject{currentmarker}{\pgfqpoint{0.000000in}{-0.020833in}}{\pgfqpoint{0.000000in}{0.000000in}}{%
\pgfpathmoveto{\pgfqpoint{0.000000in}{0.000000in}}%
\pgfpathlineto{\pgfqpoint{0.000000in}{-0.020833in}}%
\pgfusepath{stroke,fill}%
}%
\begin{pgfscope}%
\pgfsys@transformshift{2.168126in}{2.414488in}%
\pgfsys@useobject{currentmarker}{}%
\end{pgfscope}%
\end{pgfscope}%
\begin{pgfscope}%
\pgfsetbuttcap%
\pgfsetroundjoin%
\definecolor{currentfill}{rgb}{0.000000,0.000000,0.000000}%
\pgfsetfillcolor{currentfill}%
\pgfsetlinewidth{0.401500pt}%
\definecolor{currentstroke}{rgb}{0.000000,0.000000,0.000000}%
\pgfsetstrokecolor{currentstroke}%
\pgfsetdash{}{0pt}%
\pgfsys@defobject{currentmarker}{\pgfqpoint{0.000000in}{0.000000in}}{\pgfqpoint{0.000000in}{0.020833in}}{%
\pgfpathmoveto{\pgfqpoint{0.000000in}{0.000000in}}%
\pgfpathlineto{\pgfqpoint{0.000000in}{0.020833in}}%
\pgfusepath{stroke,fill}%
}%
\begin{pgfscope}%
\pgfsys@transformshift{2.228230in}{0.420000in}%
\pgfsys@useobject{currentmarker}{}%
\end{pgfscope}%
\end{pgfscope}%
\begin{pgfscope}%
\pgfsetbuttcap%
\pgfsetroundjoin%
\definecolor{currentfill}{rgb}{0.000000,0.000000,0.000000}%
\pgfsetfillcolor{currentfill}%
\pgfsetlinewidth{0.401500pt}%
\definecolor{currentstroke}{rgb}{0.000000,0.000000,0.000000}%
\pgfsetstrokecolor{currentstroke}%
\pgfsetdash{}{0pt}%
\pgfsys@defobject{currentmarker}{\pgfqpoint{0.000000in}{-0.020833in}}{\pgfqpoint{0.000000in}{0.000000in}}{%
\pgfpathmoveto{\pgfqpoint{0.000000in}{0.000000in}}%
\pgfpathlineto{\pgfqpoint{0.000000in}{-0.020833in}}%
\pgfusepath{stroke,fill}%
}%
\begin{pgfscope}%
\pgfsys@transformshift{2.228230in}{2.414488in}%
\pgfsys@useobject{currentmarker}{}%
\end{pgfscope}%
\end{pgfscope}%
\begin{pgfscope}%
\pgfsetbuttcap%
\pgfsetroundjoin%
\definecolor{currentfill}{rgb}{0.000000,0.000000,0.000000}%
\pgfsetfillcolor{currentfill}%
\pgfsetlinewidth{0.401500pt}%
\definecolor{currentstroke}{rgb}{0.000000,0.000000,0.000000}%
\pgfsetstrokecolor{currentstroke}%
\pgfsetdash{}{0pt}%
\pgfsys@defobject{currentmarker}{\pgfqpoint{0.000000in}{0.000000in}}{\pgfqpoint{0.000000in}{0.020833in}}{%
\pgfpathmoveto{\pgfqpoint{0.000000in}{0.000000in}}%
\pgfpathlineto{\pgfqpoint{0.000000in}{0.020833in}}%
\pgfusepath{stroke,fill}%
}%
\begin{pgfscope}%
\pgfsys@transformshift{2.279047in}{0.420000in}%
\pgfsys@useobject{currentmarker}{}%
\end{pgfscope}%
\end{pgfscope}%
\begin{pgfscope}%
\pgfsetbuttcap%
\pgfsetroundjoin%
\definecolor{currentfill}{rgb}{0.000000,0.000000,0.000000}%
\pgfsetfillcolor{currentfill}%
\pgfsetlinewidth{0.401500pt}%
\definecolor{currentstroke}{rgb}{0.000000,0.000000,0.000000}%
\pgfsetstrokecolor{currentstroke}%
\pgfsetdash{}{0pt}%
\pgfsys@defobject{currentmarker}{\pgfqpoint{0.000000in}{-0.020833in}}{\pgfqpoint{0.000000in}{0.000000in}}{%
\pgfpathmoveto{\pgfqpoint{0.000000in}{0.000000in}}%
\pgfpathlineto{\pgfqpoint{0.000000in}{-0.020833in}}%
\pgfusepath{stroke,fill}%
}%
\begin{pgfscope}%
\pgfsys@transformshift{2.279047in}{2.414488in}%
\pgfsys@useobject{currentmarker}{}%
\end{pgfscope}%
\end{pgfscope}%
\begin{pgfscope}%
\pgfsetbuttcap%
\pgfsetroundjoin%
\definecolor{currentfill}{rgb}{0.000000,0.000000,0.000000}%
\pgfsetfillcolor{currentfill}%
\pgfsetlinewidth{0.401500pt}%
\definecolor{currentstroke}{rgb}{0.000000,0.000000,0.000000}%
\pgfsetstrokecolor{currentstroke}%
\pgfsetdash{}{0pt}%
\pgfsys@defobject{currentmarker}{\pgfqpoint{0.000000in}{0.000000in}}{\pgfqpoint{0.000000in}{0.020833in}}{%
\pgfpathmoveto{\pgfqpoint{0.000000in}{0.000000in}}%
\pgfpathlineto{\pgfqpoint{0.000000in}{0.020833in}}%
\pgfusepath{stroke,fill}%
}%
\begin{pgfscope}%
\pgfsys@transformshift{2.323066in}{0.420000in}%
\pgfsys@useobject{currentmarker}{}%
\end{pgfscope}%
\end{pgfscope}%
\begin{pgfscope}%
\pgfsetbuttcap%
\pgfsetroundjoin%
\definecolor{currentfill}{rgb}{0.000000,0.000000,0.000000}%
\pgfsetfillcolor{currentfill}%
\pgfsetlinewidth{0.401500pt}%
\definecolor{currentstroke}{rgb}{0.000000,0.000000,0.000000}%
\pgfsetstrokecolor{currentstroke}%
\pgfsetdash{}{0pt}%
\pgfsys@defobject{currentmarker}{\pgfqpoint{0.000000in}{-0.020833in}}{\pgfqpoint{0.000000in}{0.000000in}}{%
\pgfpathmoveto{\pgfqpoint{0.000000in}{0.000000in}}%
\pgfpathlineto{\pgfqpoint{0.000000in}{-0.020833in}}%
\pgfusepath{stroke,fill}%
}%
\begin{pgfscope}%
\pgfsys@transformshift{2.323066in}{2.414488in}%
\pgfsys@useobject{currentmarker}{}%
\end{pgfscope}%
\end{pgfscope}%
\begin{pgfscope}%
\pgfsetbuttcap%
\pgfsetroundjoin%
\definecolor{currentfill}{rgb}{0.000000,0.000000,0.000000}%
\pgfsetfillcolor{currentfill}%
\pgfsetlinewidth{0.401500pt}%
\definecolor{currentstroke}{rgb}{0.000000,0.000000,0.000000}%
\pgfsetstrokecolor{currentstroke}%
\pgfsetdash{}{0pt}%
\pgfsys@defobject{currentmarker}{\pgfqpoint{0.000000in}{0.000000in}}{\pgfqpoint{0.000000in}{0.020833in}}{%
\pgfpathmoveto{\pgfqpoint{0.000000in}{0.000000in}}%
\pgfpathlineto{\pgfqpoint{0.000000in}{0.020833in}}%
\pgfusepath{stroke,fill}%
}%
\begin{pgfscope}%
\pgfsys@transformshift{2.361894in}{0.420000in}%
\pgfsys@useobject{currentmarker}{}%
\end{pgfscope}%
\end{pgfscope}%
\begin{pgfscope}%
\pgfsetbuttcap%
\pgfsetroundjoin%
\definecolor{currentfill}{rgb}{0.000000,0.000000,0.000000}%
\pgfsetfillcolor{currentfill}%
\pgfsetlinewidth{0.401500pt}%
\definecolor{currentstroke}{rgb}{0.000000,0.000000,0.000000}%
\pgfsetstrokecolor{currentstroke}%
\pgfsetdash{}{0pt}%
\pgfsys@defobject{currentmarker}{\pgfqpoint{0.000000in}{-0.020833in}}{\pgfqpoint{0.000000in}{0.000000in}}{%
\pgfpathmoveto{\pgfqpoint{0.000000in}{0.000000in}}%
\pgfpathlineto{\pgfqpoint{0.000000in}{-0.020833in}}%
\pgfusepath{stroke,fill}%
}%
\begin{pgfscope}%
\pgfsys@transformshift{2.361894in}{2.414488in}%
\pgfsys@useobject{currentmarker}{}%
\end{pgfscope}%
\end{pgfscope}%
\begin{pgfscope}%
\pgfsetbuttcap%
\pgfsetroundjoin%
\definecolor{currentfill}{rgb}{0.000000,0.000000,0.000000}%
\pgfsetfillcolor{currentfill}%
\pgfsetlinewidth{0.401500pt}%
\definecolor{currentstroke}{rgb}{0.000000,0.000000,0.000000}%
\pgfsetstrokecolor{currentstroke}%
\pgfsetdash{}{0pt}%
\pgfsys@defobject{currentmarker}{\pgfqpoint{0.000000in}{0.000000in}}{\pgfqpoint{0.000000in}{0.020833in}}{%
\pgfpathmoveto{\pgfqpoint{0.000000in}{0.000000in}}%
\pgfpathlineto{\pgfqpoint{0.000000in}{0.020833in}}%
\pgfusepath{stroke,fill}%
}%
\begin{pgfscope}%
\pgfsys@transformshift{2.625128in}{0.420000in}%
\pgfsys@useobject{currentmarker}{}%
\end{pgfscope}%
\end{pgfscope}%
\begin{pgfscope}%
\pgfsetbuttcap%
\pgfsetroundjoin%
\definecolor{currentfill}{rgb}{0.000000,0.000000,0.000000}%
\pgfsetfillcolor{currentfill}%
\pgfsetlinewidth{0.401500pt}%
\definecolor{currentstroke}{rgb}{0.000000,0.000000,0.000000}%
\pgfsetstrokecolor{currentstroke}%
\pgfsetdash{}{0pt}%
\pgfsys@defobject{currentmarker}{\pgfqpoint{0.000000in}{-0.020833in}}{\pgfqpoint{0.000000in}{0.000000in}}{%
\pgfpathmoveto{\pgfqpoint{0.000000in}{0.000000in}}%
\pgfpathlineto{\pgfqpoint{0.000000in}{-0.020833in}}%
\pgfusepath{stroke,fill}%
}%
\begin{pgfscope}%
\pgfsys@transformshift{2.625128in}{2.414488in}%
\pgfsys@useobject{currentmarker}{}%
\end{pgfscope}%
\end{pgfscope}%
\begin{pgfscope}%
\pgfsetbuttcap%
\pgfsetroundjoin%
\definecolor{currentfill}{rgb}{0.000000,0.000000,0.000000}%
\pgfsetfillcolor{currentfill}%
\pgfsetlinewidth{0.401500pt}%
\definecolor{currentstroke}{rgb}{0.000000,0.000000,0.000000}%
\pgfsetstrokecolor{currentstroke}%
\pgfsetdash{}{0pt}%
\pgfsys@defobject{currentmarker}{\pgfqpoint{0.000000in}{0.000000in}}{\pgfqpoint{0.000000in}{0.020833in}}{%
\pgfpathmoveto{\pgfqpoint{0.000000in}{0.000000in}}%
\pgfpathlineto{\pgfqpoint{0.000000in}{0.020833in}}%
\pgfusepath{stroke,fill}%
}%
\begin{pgfscope}%
\pgfsys@transformshift{2.758792in}{0.420000in}%
\pgfsys@useobject{currentmarker}{}%
\end{pgfscope}%
\end{pgfscope}%
\begin{pgfscope}%
\pgfsetbuttcap%
\pgfsetroundjoin%
\definecolor{currentfill}{rgb}{0.000000,0.000000,0.000000}%
\pgfsetfillcolor{currentfill}%
\pgfsetlinewidth{0.401500pt}%
\definecolor{currentstroke}{rgb}{0.000000,0.000000,0.000000}%
\pgfsetstrokecolor{currentstroke}%
\pgfsetdash{}{0pt}%
\pgfsys@defobject{currentmarker}{\pgfqpoint{0.000000in}{-0.020833in}}{\pgfqpoint{0.000000in}{0.000000in}}{%
\pgfpathmoveto{\pgfqpoint{0.000000in}{0.000000in}}%
\pgfpathlineto{\pgfqpoint{0.000000in}{-0.020833in}}%
\pgfusepath{stroke,fill}%
}%
\begin{pgfscope}%
\pgfsys@transformshift{2.758792in}{2.414488in}%
\pgfsys@useobject{currentmarker}{}%
\end{pgfscope}%
\end{pgfscope}%
\begin{pgfscope}%
\pgfsetbuttcap%
\pgfsetroundjoin%
\definecolor{currentfill}{rgb}{0.000000,0.000000,0.000000}%
\pgfsetfillcolor{currentfill}%
\pgfsetlinewidth{0.401500pt}%
\definecolor{currentstroke}{rgb}{0.000000,0.000000,0.000000}%
\pgfsetstrokecolor{currentstroke}%
\pgfsetdash{}{0pt}%
\pgfsys@defobject{currentmarker}{\pgfqpoint{0.000000in}{0.000000in}}{\pgfqpoint{0.000000in}{0.020833in}}{%
\pgfpathmoveto{\pgfqpoint{0.000000in}{0.000000in}}%
\pgfpathlineto{\pgfqpoint{0.000000in}{0.020833in}}%
\pgfusepath{stroke,fill}%
}%
\begin{pgfscope}%
\pgfsys@transformshift{2.853629in}{0.420000in}%
\pgfsys@useobject{currentmarker}{}%
\end{pgfscope}%
\end{pgfscope}%
\begin{pgfscope}%
\pgfsetbuttcap%
\pgfsetroundjoin%
\definecolor{currentfill}{rgb}{0.000000,0.000000,0.000000}%
\pgfsetfillcolor{currentfill}%
\pgfsetlinewidth{0.401500pt}%
\definecolor{currentstroke}{rgb}{0.000000,0.000000,0.000000}%
\pgfsetstrokecolor{currentstroke}%
\pgfsetdash{}{0pt}%
\pgfsys@defobject{currentmarker}{\pgfqpoint{0.000000in}{-0.020833in}}{\pgfqpoint{0.000000in}{0.000000in}}{%
\pgfpathmoveto{\pgfqpoint{0.000000in}{0.000000in}}%
\pgfpathlineto{\pgfqpoint{0.000000in}{-0.020833in}}%
\pgfusepath{stroke,fill}%
}%
\begin{pgfscope}%
\pgfsys@transformshift{2.853629in}{2.414488in}%
\pgfsys@useobject{currentmarker}{}%
\end{pgfscope}%
\end{pgfscope}%
\begin{pgfscope}%
\pgfsetbuttcap%
\pgfsetroundjoin%
\definecolor{currentfill}{rgb}{0.000000,0.000000,0.000000}%
\pgfsetfillcolor{currentfill}%
\pgfsetlinewidth{0.401500pt}%
\definecolor{currentstroke}{rgb}{0.000000,0.000000,0.000000}%
\pgfsetstrokecolor{currentstroke}%
\pgfsetdash{}{0pt}%
\pgfsys@defobject{currentmarker}{\pgfqpoint{0.000000in}{0.000000in}}{\pgfqpoint{0.000000in}{0.020833in}}{%
\pgfpathmoveto{\pgfqpoint{0.000000in}{0.000000in}}%
\pgfpathlineto{\pgfqpoint{0.000000in}{0.020833in}}%
\pgfusepath{stroke,fill}%
}%
\begin{pgfscope}%
\pgfsys@transformshift{2.927190in}{0.420000in}%
\pgfsys@useobject{currentmarker}{}%
\end{pgfscope}%
\end{pgfscope}%
\begin{pgfscope}%
\pgfsetbuttcap%
\pgfsetroundjoin%
\definecolor{currentfill}{rgb}{0.000000,0.000000,0.000000}%
\pgfsetfillcolor{currentfill}%
\pgfsetlinewidth{0.401500pt}%
\definecolor{currentstroke}{rgb}{0.000000,0.000000,0.000000}%
\pgfsetstrokecolor{currentstroke}%
\pgfsetdash{}{0pt}%
\pgfsys@defobject{currentmarker}{\pgfqpoint{0.000000in}{-0.020833in}}{\pgfqpoint{0.000000in}{0.000000in}}{%
\pgfpathmoveto{\pgfqpoint{0.000000in}{0.000000in}}%
\pgfpathlineto{\pgfqpoint{0.000000in}{-0.020833in}}%
\pgfusepath{stroke,fill}%
}%
\begin{pgfscope}%
\pgfsys@transformshift{2.927190in}{2.414488in}%
\pgfsys@useobject{currentmarker}{}%
\end{pgfscope}%
\end{pgfscope}%
\begin{pgfscope}%
\pgfsetbuttcap%
\pgfsetroundjoin%
\definecolor{currentfill}{rgb}{0.000000,0.000000,0.000000}%
\pgfsetfillcolor{currentfill}%
\pgfsetlinewidth{0.401500pt}%
\definecolor{currentstroke}{rgb}{0.000000,0.000000,0.000000}%
\pgfsetstrokecolor{currentstroke}%
\pgfsetdash{}{0pt}%
\pgfsys@defobject{currentmarker}{\pgfqpoint{0.000000in}{0.000000in}}{\pgfqpoint{0.000000in}{0.020833in}}{%
\pgfpathmoveto{\pgfqpoint{0.000000in}{0.000000in}}%
\pgfpathlineto{\pgfqpoint{0.000000in}{0.020833in}}%
\pgfusepath{stroke,fill}%
}%
\begin{pgfscope}%
\pgfsys@transformshift{2.987293in}{0.420000in}%
\pgfsys@useobject{currentmarker}{}%
\end{pgfscope}%
\end{pgfscope}%
\begin{pgfscope}%
\pgfsetbuttcap%
\pgfsetroundjoin%
\definecolor{currentfill}{rgb}{0.000000,0.000000,0.000000}%
\pgfsetfillcolor{currentfill}%
\pgfsetlinewidth{0.401500pt}%
\definecolor{currentstroke}{rgb}{0.000000,0.000000,0.000000}%
\pgfsetstrokecolor{currentstroke}%
\pgfsetdash{}{0pt}%
\pgfsys@defobject{currentmarker}{\pgfqpoint{0.000000in}{-0.020833in}}{\pgfqpoint{0.000000in}{0.000000in}}{%
\pgfpathmoveto{\pgfqpoint{0.000000in}{0.000000in}}%
\pgfpathlineto{\pgfqpoint{0.000000in}{-0.020833in}}%
\pgfusepath{stroke,fill}%
}%
\begin{pgfscope}%
\pgfsys@transformshift{2.987293in}{2.414488in}%
\pgfsys@useobject{currentmarker}{}%
\end{pgfscope}%
\end{pgfscope}%
\begin{pgfscope}%
\pgfsetbuttcap%
\pgfsetroundjoin%
\definecolor{currentfill}{rgb}{0.000000,0.000000,0.000000}%
\pgfsetfillcolor{currentfill}%
\pgfsetlinewidth{0.401500pt}%
\definecolor{currentstroke}{rgb}{0.000000,0.000000,0.000000}%
\pgfsetstrokecolor{currentstroke}%
\pgfsetdash{}{0pt}%
\pgfsys@defobject{currentmarker}{\pgfqpoint{0.000000in}{0.000000in}}{\pgfqpoint{0.000000in}{0.020833in}}{%
\pgfpathmoveto{\pgfqpoint{0.000000in}{0.000000in}}%
\pgfpathlineto{\pgfqpoint{0.000000in}{0.020833in}}%
\pgfusepath{stroke,fill}%
}%
\begin{pgfscope}%
\pgfsys@transformshift{3.038110in}{0.420000in}%
\pgfsys@useobject{currentmarker}{}%
\end{pgfscope}%
\end{pgfscope}%
\begin{pgfscope}%
\pgfsetbuttcap%
\pgfsetroundjoin%
\definecolor{currentfill}{rgb}{0.000000,0.000000,0.000000}%
\pgfsetfillcolor{currentfill}%
\pgfsetlinewidth{0.401500pt}%
\definecolor{currentstroke}{rgb}{0.000000,0.000000,0.000000}%
\pgfsetstrokecolor{currentstroke}%
\pgfsetdash{}{0pt}%
\pgfsys@defobject{currentmarker}{\pgfqpoint{0.000000in}{-0.020833in}}{\pgfqpoint{0.000000in}{0.000000in}}{%
\pgfpathmoveto{\pgfqpoint{0.000000in}{0.000000in}}%
\pgfpathlineto{\pgfqpoint{0.000000in}{-0.020833in}}%
\pgfusepath{stroke,fill}%
}%
\begin{pgfscope}%
\pgfsys@transformshift{3.038110in}{2.414488in}%
\pgfsys@useobject{currentmarker}{}%
\end{pgfscope}%
\end{pgfscope}%
\begin{pgfscope}%
\definecolor{textcolor}{rgb}{0.000000,0.000000,0.000000}%
\pgfsetstrokecolor{textcolor}%
\pgfsetfillcolor{textcolor}%
\pgftext[x=1.844055in,y=0.208426in,,top]{\color{textcolor}{\rmfamily\fontsize{12.000000}{14.400000}\selectfont\catcode`\^=\active\def^{\ifmmode\sp\else\^{}\fi}\catcode`\%=\active\def%{\%}$y^{+}$}}%
\end{pgfscope}%
\begin{pgfscope}%
\pgfsetbuttcap%
\pgfsetroundjoin%
\definecolor{currentfill}{rgb}{0.000000,0.000000,0.000000}%
\pgfsetfillcolor{currentfill}%
\pgfsetlinewidth{0.501875pt}%
\definecolor{currentstroke}{rgb}{0.000000,0.000000,0.000000}%
\pgfsetstrokecolor{currentstroke}%
\pgfsetdash{}{0pt}%
\pgfsys@defobject{currentmarker}{\pgfqpoint{0.000000in}{0.000000in}}{\pgfqpoint{0.041667in}{0.000000in}}{%
\pgfpathmoveto{\pgfqpoint{0.000000in}{0.000000in}}%
\pgfpathlineto{\pgfqpoint{0.041667in}{0.000000in}}%
\pgfusepath{stroke,fill}%
}%
\begin{pgfscope}%
\pgfsys@transformshift{0.650000in}{0.615427in}%
\pgfsys@useobject{currentmarker}{}%
\end{pgfscope}%
\end{pgfscope}%
\begin{pgfscope}%
\pgfsetbuttcap%
\pgfsetroundjoin%
\definecolor{currentfill}{rgb}{0.000000,0.000000,0.000000}%
\pgfsetfillcolor{currentfill}%
\pgfsetlinewidth{0.501875pt}%
\definecolor{currentstroke}{rgb}{0.000000,0.000000,0.000000}%
\pgfsetstrokecolor{currentstroke}%
\pgfsetdash{}{0pt}%
\pgfsys@defobject{currentmarker}{\pgfqpoint{-0.041667in}{0.000000in}}{\pgfqpoint{-0.000000in}{0.000000in}}{%
\pgfpathmoveto{\pgfqpoint{-0.000000in}{0.000000in}}%
\pgfpathlineto{\pgfqpoint{-0.041667in}{0.000000in}}%
\pgfusepath{stroke,fill}%
}%
\begin{pgfscope}%
\pgfsys@transformshift{3.038110in}{0.615427in}%
\pgfsys@useobject{currentmarker}{}%
\end{pgfscope}%
\end{pgfscope}%
\begin{pgfscope}%
\definecolor{textcolor}{rgb}{0.000000,0.000000,0.000000}%
\pgfsetstrokecolor{textcolor}%
\pgfsetfillcolor{textcolor}%
\pgftext[x=0.288772in, y=0.562620in, left, base]{\color{textcolor}{\rmfamily\fontsize{11.000000}{13.200000}\selectfont\catcode`\^=\active\def^{\ifmmode\sp\else\^{}\fi}\catcode`\%=\active\def%{\%}\ensuremath{-}0.2}}%
\end{pgfscope}%
\begin{pgfscope}%
\pgfsetbuttcap%
\pgfsetroundjoin%
\definecolor{currentfill}{rgb}{0.000000,0.000000,0.000000}%
\pgfsetfillcolor{currentfill}%
\pgfsetlinewidth{0.501875pt}%
\definecolor{currentstroke}{rgb}{0.000000,0.000000,0.000000}%
\pgfsetstrokecolor{currentstroke}%
\pgfsetdash{}{0pt}%
\pgfsys@defobject{currentmarker}{\pgfqpoint{0.000000in}{0.000000in}}{\pgfqpoint{0.041667in}{0.000000in}}{%
\pgfpathmoveto{\pgfqpoint{0.000000in}{0.000000in}}%
\pgfpathlineto{\pgfqpoint{0.041667in}{0.000000in}}%
\pgfusepath{stroke,fill}%
}%
\begin{pgfscope}%
\pgfsys@transformshift{0.650000in}{0.995958in}%
\pgfsys@useobject{currentmarker}{}%
\end{pgfscope}%
\end{pgfscope}%
\begin{pgfscope}%
\pgfsetbuttcap%
\pgfsetroundjoin%
\definecolor{currentfill}{rgb}{0.000000,0.000000,0.000000}%
\pgfsetfillcolor{currentfill}%
\pgfsetlinewidth{0.501875pt}%
\definecolor{currentstroke}{rgb}{0.000000,0.000000,0.000000}%
\pgfsetstrokecolor{currentstroke}%
\pgfsetdash{}{0pt}%
\pgfsys@defobject{currentmarker}{\pgfqpoint{-0.041667in}{0.000000in}}{\pgfqpoint{-0.000000in}{0.000000in}}{%
\pgfpathmoveto{\pgfqpoint{-0.000000in}{0.000000in}}%
\pgfpathlineto{\pgfqpoint{-0.041667in}{0.000000in}}%
\pgfusepath{stroke,fill}%
}%
\begin{pgfscope}%
\pgfsys@transformshift{3.038110in}{0.995958in}%
\pgfsys@useobject{currentmarker}{}%
\end{pgfscope}%
\end{pgfscope}%
\begin{pgfscope}%
\definecolor{textcolor}{rgb}{0.000000,0.000000,0.000000}%
\pgfsetstrokecolor{textcolor}%
\pgfsetfillcolor{textcolor}%
\pgftext[x=0.288772in, y=0.943151in, left, base]{\color{textcolor}{\rmfamily\fontsize{11.000000}{13.200000}\selectfont\catcode`\^=\active\def^{\ifmmode\sp\else\^{}\fi}\catcode`\%=\active\def%{\%}\ensuremath{-}0.1}}%
\end{pgfscope}%
\begin{pgfscope}%
\pgfsetbuttcap%
\pgfsetroundjoin%
\definecolor{currentfill}{rgb}{0.000000,0.000000,0.000000}%
\pgfsetfillcolor{currentfill}%
\pgfsetlinewidth{0.501875pt}%
\definecolor{currentstroke}{rgb}{0.000000,0.000000,0.000000}%
\pgfsetstrokecolor{currentstroke}%
\pgfsetdash{}{0pt}%
\pgfsys@defobject{currentmarker}{\pgfqpoint{0.000000in}{0.000000in}}{\pgfqpoint{0.041667in}{0.000000in}}{%
\pgfpathmoveto{\pgfqpoint{0.000000in}{0.000000in}}%
\pgfpathlineto{\pgfqpoint{0.041667in}{0.000000in}}%
\pgfusepath{stroke,fill}%
}%
\begin{pgfscope}%
\pgfsys@transformshift{0.650000in}{1.376489in}%
\pgfsys@useobject{currentmarker}{}%
\end{pgfscope}%
\end{pgfscope}%
\begin{pgfscope}%
\pgfsetbuttcap%
\pgfsetroundjoin%
\definecolor{currentfill}{rgb}{0.000000,0.000000,0.000000}%
\pgfsetfillcolor{currentfill}%
\pgfsetlinewidth{0.501875pt}%
\definecolor{currentstroke}{rgb}{0.000000,0.000000,0.000000}%
\pgfsetstrokecolor{currentstroke}%
\pgfsetdash{}{0pt}%
\pgfsys@defobject{currentmarker}{\pgfqpoint{-0.041667in}{0.000000in}}{\pgfqpoint{-0.000000in}{0.000000in}}{%
\pgfpathmoveto{\pgfqpoint{-0.000000in}{0.000000in}}%
\pgfpathlineto{\pgfqpoint{-0.041667in}{0.000000in}}%
\pgfusepath{stroke,fill}%
}%
\begin{pgfscope}%
\pgfsys@transformshift{3.038110in}{1.376489in}%
\pgfsys@useobject{currentmarker}{}%
\end{pgfscope}%
\end{pgfscope}%
\begin{pgfscope}%
\definecolor{textcolor}{rgb}{0.000000,0.000000,0.000000}%
\pgfsetstrokecolor{textcolor}%
\pgfsetfillcolor{textcolor}%
\pgftext[x=0.407060in, y=1.323682in, left, base]{\color{textcolor}{\rmfamily\fontsize{11.000000}{13.200000}\selectfont\catcode`\^=\active\def^{\ifmmode\sp\else\^{}\fi}\catcode`\%=\active\def%{\%}0.0}}%
\end{pgfscope}%
\begin{pgfscope}%
\pgfsetbuttcap%
\pgfsetroundjoin%
\definecolor{currentfill}{rgb}{0.000000,0.000000,0.000000}%
\pgfsetfillcolor{currentfill}%
\pgfsetlinewidth{0.501875pt}%
\definecolor{currentstroke}{rgb}{0.000000,0.000000,0.000000}%
\pgfsetstrokecolor{currentstroke}%
\pgfsetdash{}{0pt}%
\pgfsys@defobject{currentmarker}{\pgfqpoint{0.000000in}{0.000000in}}{\pgfqpoint{0.041667in}{0.000000in}}{%
\pgfpathmoveto{\pgfqpoint{0.000000in}{0.000000in}}%
\pgfpathlineto{\pgfqpoint{0.041667in}{0.000000in}}%
\pgfusepath{stroke,fill}%
}%
\begin{pgfscope}%
\pgfsys@transformshift{0.650000in}{1.757020in}%
\pgfsys@useobject{currentmarker}{}%
\end{pgfscope}%
\end{pgfscope}%
\begin{pgfscope}%
\pgfsetbuttcap%
\pgfsetroundjoin%
\definecolor{currentfill}{rgb}{0.000000,0.000000,0.000000}%
\pgfsetfillcolor{currentfill}%
\pgfsetlinewidth{0.501875pt}%
\definecolor{currentstroke}{rgb}{0.000000,0.000000,0.000000}%
\pgfsetstrokecolor{currentstroke}%
\pgfsetdash{}{0pt}%
\pgfsys@defobject{currentmarker}{\pgfqpoint{-0.041667in}{0.000000in}}{\pgfqpoint{-0.000000in}{0.000000in}}{%
\pgfpathmoveto{\pgfqpoint{-0.000000in}{0.000000in}}%
\pgfpathlineto{\pgfqpoint{-0.041667in}{0.000000in}}%
\pgfusepath{stroke,fill}%
}%
\begin{pgfscope}%
\pgfsys@transformshift{3.038110in}{1.757020in}%
\pgfsys@useobject{currentmarker}{}%
\end{pgfscope}%
\end{pgfscope}%
\begin{pgfscope}%
\definecolor{textcolor}{rgb}{0.000000,0.000000,0.000000}%
\pgfsetstrokecolor{textcolor}%
\pgfsetfillcolor{textcolor}%
\pgftext[x=0.407060in, y=1.704213in, left, base]{\color{textcolor}{\rmfamily\fontsize{11.000000}{13.200000}\selectfont\catcode`\^=\active\def^{\ifmmode\sp\else\^{}\fi}\catcode`\%=\active\def%{\%}0.1}}%
\end{pgfscope}%
\begin{pgfscope}%
\pgfsetbuttcap%
\pgfsetroundjoin%
\definecolor{currentfill}{rgb}{0.000000,0.000000,0.000000}%
\pgfsetfillcolor{currentfill}%
\pgfsetlinewidth{0.501875pt}%
\definecolor{currentstroke}{rgb}{0.000000,0.000000,0.000000}%
\pgfsetstrokecolor{currentstroke}%
\pgfsetdash{}{0pt}%
\pgfsys@defobject{currentmarker}{\pgfqpoint{0.000000in}{0.000000in}}{\pgfqpoint{0.041667in}{0.000000in}}{%
\pgfpathmoveto{\pgfqpoint{0.000000in}{0.000000in}}%
\pgfpathlineto{\pgfqpoint{0.041667in}{0.000000in}}%
\pgfusepath{stroke,fill}%
}%
\begin{pgfscope}%
\pgfsys@transformshift{0.650000in}{2.137551in}%
\pgfsys@useobject{currentmarker}{}%
\end{pgfscope}%
\end{pgfscope}%
\begin{pgfscope}%
\pgfsetbuttcap%
\pgfsetroundjoin%
\definecolor{currentfill}{rgb}{0.000000,0.000000,0.000000}%
\pgfsetfillcolor{currentfill}%
\pgfsetlinewidth{0.501875pt}%
\definecolor{currentstroke}{rgb}{0.000000,0.000000,0.000000}%
\pgfsetstrokecolor{currentstroke}%
\pgfsetdash{}{0pt}%
\pgfsys@defobject{currentmarker}{\pgfqpoint{-0.041667in}{0.000000in}}{\pgfqpoint{-0.000000in}{0.000000in}}{%
\pgfpathmoveto{\pgfqpoint{-0.000000in}{0.000000in}}%
\pgfpathlineto{\pgfqpoint{-0.041667in}{0.000000in}}%
\pgfusepath{stroke,fill}%
}%
\begin{pgfscope}%
\pgfsys@transformshift{3.038110in}{2.137551in}%
\pgfsys@useobject{currentmarker}{}%
\end{pgfscope}%
\end{pgfscope}%
\begin{pgfscope}%
\definecolor{textcolor}{rgb}{0.000000,0.000000,0.000000}%
\pgfsetstrokecolor{textcolor}%
\pgfsetfillcolor{textcolor}%
\pgftext[x=0.407060in, y=2.084744in, left, base]{\color{textcolor}{\rmfamily\fontsize{11.000000}{13.200000}\selectfont\catcode`\^=\active\def^{\ifmmode\sp\else\^{}\fi}\catcode`\%=\active\def%{\%}0.2}}%
\end{pgfscope}%
\begin{pgfscope}%
\pgfsetbuttcap%
\pgfsetroundjoin%
\definecolor{currentfill}{rgb}{0.000000,0.000000,0.000000}%
\pgfsetfillcolor{currentfill}%
\pgfsetlinewidth{0.401500pt}%
\definecolor{currentstroke}{rgb}{0.000000,0.000000,0.000000}%
\pgfsetstrokecolor{currentstroke}%
\pgfsetdash{}{0pt}%
\pgfsys@defobject{currentmarker}{\pgfqpoint{0.000000in}{0.000000in}}{\pgfqpoint{0.020833in}{0.000000in}}{%
\pgfpathmoveto{\pgfqpoint{0.000000in}{0.000000in}}%
\pgfpathlineto{\pgfqpoint{0.020833in}{0.000000in}}%
\pgfusepath{stroke,fill}%
}%
\begin{pgfscope}%
\pgfsys@transformshift{0.650000in}{0.463214in}%
\pgfsys@useobject{currentmarker}{}%
\end{pgfscope}%
\end{pgfscope}%
\begin{pgfscope}%
\pgfsetbuttcap%
\pgfsetroundjoin%
\definecolor{currentfill}{rgb}{0.000000,0.000000,0.000000}%
\pgfsetfillcolor{currentfill}%
\pgfsetlinewidth{0.401500pt}%
\definecolor{currentstroke}{rgb}{0.000000,0.000000,0.000000}%
\pgfsetstrokecolor{currentstroke}%
\pgfsetdash{}{0pt}%
\pgfsys@defobject{currentmarker}{\pgfqpoint{-0.020833in}{0.000000in}}{\pgfqpoint{-0.000000in}{0.000000in}}{%
\pgfpathmoveto{\pgfqpoint{-0.000000in}{0.000000in}}%
\pgfpathlineto{\pgfqpoint{-0.020833in}{0.000000in}}%
\pgfusepath{stroke,fill}%
}%
\begin{pgfscope}%
\pgfsys@transformshift{3.038110in}{0.463214in}%
\pgfsys@useobject{currentmarker}{}%
\end{pgfscope}%
\end{pgfscope}%
\begin{pgfscope}%
\pgfsetbuttcap%
\pgfsetroundjoin%
\definecolor{currentfill}{rgb}{0.000000,0.000000,0.000000}%
\pgfsetfillcolor{currentfill}%
\pgfsetlinewidth{0.401500pt}%
\definecolor{currentstroke}{rgb}{0.000000,0.000000,0.000000}%
\pgfsetstrokecolor{currentstroke}%
\pgfsetdash{}{0pt}%
\pgfsys@defobject{currentmarker}{\pgfqpoint{0.000000in}{0.000000in}}{\pgfqpoint{0.020833in}{0.000000in}}{%
\pgfpathmoveto{\pgfqpoint{0.000000in}{0.000000in}}%
\pgfpathlineto{\pgfqpoint{0.020833in}{0.000000in}}%
\pgfusepath{stroke,fill}%
}%
\begin{pgfscope}%
\pgfsys@transformshift{0.650000in}{0.539321in}%
\pgfsys@useobject{currentmarker}{}%
\end{pgfscope}%
\end{pgfscope}%
\begin{pgfscope}%
\pgfsetbuttcap%
\pgfsetroundjoin%
\definecolor{currentfill}{rgb}{0.000000,0.000000,0.000000}%
\pgfsetfillcolor{currentfill}%
\pgfsetlinewidth{0.401500pt}%
\definecolor{currentstroke}{rgb}{0.000000,0.000000,0.000000}%
\pgfsetstrokecolor{currentstroke}%
\pgfsetdash{}{0pt}%
\pgfsys@defobject{currentmarker}{\pgfqpoint{-0.020833in}{0.000000in}}{\pgfqpoint{-0.000000in}{0.000000in}}{%
\pgfpathmoveto{\pgfqpoint{-0.000000in}{0.000000in}}%
\pgfpathlineto{\pgfqpoint{-0.020833in}{0.000000in}}%
\pgfusepath{stroke,fill}%
}%
\begin{pgfscope}%
\pgfsys@transformshift{3.038110in}{0.539321in}%
\pgfsys@useobject{currentmarker}{}%
\end{pgfscope}%
\end{pgfscope}%
\begin{pgfscope}%
\pgfsetbuttcap%
\pgfsetroundjoin%
\definecolor{currentfill}{rgb}{0.000000,0.000000,0.000000}%
\pgfsetfillcolor{currentfill}%
\pgfsetlinewidth{0.401500pt}%
\definecolor{currentstroke}{rgb}{0.000000,0.000000,0.000000}%
\pgfsetstrokecolor{currentstroke}%
\pgfsetdash{}{0pt}%
\pgfsys@defobject{currentmarker}{\pgfqpoint{0.000000in}{0.000000in}}{\pgfqpoint{0.020833in}{0.000000in}}{%
\pgfpathmoveto{\pgfqpoint{0.000000in}{0.000000in}}%
\pgfpathlineto{\pgfqpoint{0.020833in}{0.000000in}}%
\pgfusepath{stroke,fill}%
}%
\begin{pgfscope}%
\pgfsys@transformshift{0.650000in}{0.691533in}%
\pgfsys@useobject{currentmarker}{}%
\end{pgfscope}%
\end{pgfscope}%
\begin{pgfscope}%
\pgfsetbuttcap%
\pgfsetroundjoin%
\definecolor{currentfill}{rgb}{0.000000,0.000000,0.000000}%
\pgfsetfillcolor{currentfill}%
\pgfsetlinewidth{0.401500pt}%
\definecolor{currentstroke}{rgb}{0.000000,0.000000,0.000000}%
\pgfsetstrokecolor{currentstroke}%
\pgfsetdash{}{0pt}%
\pgfsys@defobject{currentmarker}{\pgfqpoint{-0.020833in}{0.000000in}}{\pgfqpoint{-0.000000in}{0.000000in}}{%
\pgfpathmoveto{\pgfqpoint{-0.000000in}{0.000000in}}%
\pgfpathlineto{\pgfqpoint{-0.020833in}{0.000000in}}%
\pgfusepath{stroke,fill}%
}%
\begin{pgfscope}%
\pgfsys@transformshift{3.038110in}{0.691533in}%
\pgfsys@useobject{currentmarker}{}%
\end{pgfscope}%
\end{pgfscope}%
\begin{pgfscope}%
\pgfsetbuttcap%
\pgfsetroundjoin%
\definecolor{currentfill}{rgb}{0.000000,0.000000,0.000000}%
\pgfsetfillcolor{currentfill}%
\pgfsetlinewidth{0.401500pt}%
\definecolor{currentstroke}{rgb}{0.000000,0.000000,0.000000}%
\pgfsetstrokecolor{currentstroke}%
\pgfsetdash{}{0pt}%
\pgfsys@defobject{currentmarker}{\pgfqpoint{0.000000in}{0.000000in}}{\pgfqpoint{0.020833in}{0.000000in}}{%
\pgfpathmoveto{\pgfqpoint{0.000000in}{0.000000in}}%
\pgfpathlineto{\pgfqpoint{0.020833in}{0.000000in}}%
\pgfusepath{stroke,fill}%
}%
\begin{pgfscope}%
\pgfsys@transformshift{0.650000in}{0.767639in}%
\pgfsys@useobject{currentmarker}{}%
\end{pgfscope}%
\end{pgfscope}%
\begin{pgfscope}%
\pgfsetbuttcap%
\pgfsetroundjoin%
\definecolor{currentfill}{rgb}{0.000000,0.000000,0.000000}%
\pgfsetfillcolor{currentfill}%
\pgfsetlinewidth{0.401500pt}%
\definecolor{currentstroke}{rgb}{0.000000,0.000000,0.000000}%
\pgfsetstrokecolor{currentstroke}%
\pgfsetdash{}{0pt}%
\pgfsys@defobject{currentmarker}{\pgfqpoint{-0.020833in}{0.000000in}}{\pgfqpoint{-0.000000in}{0.000000in}}{%
\pgfpathmoveto{\pgfqpoint{-0.000000in}{0.000000in}}%
\pgfpathlineto{\pgfqpoint{-0.020833in}{0.000000in}}%
\pgfusepath{stroke,fill}%
}%
\begin{pgfscope}%
\pgfsys@transformshift{3.038110in}{0.767639in}%
\pgfsys@useobject{currentmarker}{}%
\end{pgfscope}%
\end{pgfscope}%
\begin{pgfscope}%
\pgfsetbuttcap%
\pgfsetroundjoin%
\definecolor{currentfill}{rgb}{0.000000,0.000000,0.000000}%
\pgfsetfillcolor{currentfill}%
\pgfsetlinewidth{0.401500pt}%
\definecolor{currentstroke}{rgb}{0.000000,0.000000,0.000000}%
\pgfsetstrokecolor{currentstroke}%
\pgfsetdash{}{0pt}%
\pgfsys@defobject{currentmarker}{\pgfqpoint{0.000000in}{0.000000in}}{\pgfqpoint{0.020833in}{0.000000in}}{%
\pgfpathmoveto{\pgfqpoint{0.000000in}{0.000000in}}%
\pgfpathlineto{\pgfqpoint{0.020833in}{0.000000in}}%
\pgfusepath{stroke,fill}%
}%
\begin{pgfscope}%
\pgfsys@transformshift{0.650000in}{0.843745in}%
\pgfsys@useobject{currentmarker}{}%
\end{pgfscope}%
\end{pgfscope}%
\begin{pgfscope}%
\pgfsetbuttcap%
\pgfsetroundjoin%
\definecolor{currentfill}{rgb}{0.000000,0.000000,0.000000}%
\pgfsetfillcolor{currentfill}%
\pgfsetlinewidth{0.401500pt}%
\definecolor{currentstroke}{rgb}{0.000000,0.000000,0.000000}%
\pgfsetstrokecolor{currentstroke}%
\pgfsetdash{}{0pt}%
\pgfsys@defobject{currentmarker}{\pgfqpoint{-0.020833in}{0.000000in}}{\pgfqpoint{-0.000000in}{0.000000in}}{%
\pgfpathmoveto{\pgfqpoint{-0.000000in}{0.000000in}}%
\pgfpathlineto{\pgfqpoint{-0.020833in}{0.000000in}}%
\pgfusepath{stroke,fill}%
}%
\begin{pgfscope}%
\pgfsys@transformshift{3.038110in}{0.843745in}%
\pgfsys@useobject{currentmarker}{}%
\end{pgfscope}%
\end{pgfscope}%
\begin{pgfscope}%
\pgfsetbuttcap%
\pgfsetroundjoin%
\definecolor{currentfill}{rgb}{0.000000,0.000000,0.000000}%
\pgfsetfillcolor{currentfill}%
\pgfsetlinewidth{0.401500pt}%
\definecolor{currentstroke}{rgb}{0.000000,0.000000,0.000000}%
\pgfsetstrokecolor{currentstroke}%
\pgfsetdash{}{0pt}%
\pgfsys@defobject{currentmarker}{\pgfqpoint{0.000000in}{0.000000in}}{\pgfqpoint{0.020833in}{0.000000in}}{%
\pgfpathmoveto{\pgfqpoint{0.000000in}{0.000000in}}%
\pgfpathlineto{\pgfqpoint{0.020833in}{0.000000in}}%
\pgfusepath{stroke,fill}%
}%
\begin{pgfscope}%
\pgfsys@transformshift{0.650000in}{0.919852in}%
\pgfsys@useobject{currentmarker}{}%
\end{pgfscope}%
\end{pgfscope}%
\begin{pgfscope}%
\pgfsetbuttcap%
\pgfsetroundjoin%
\definecolor{currentfill}{rgb}{0.000000,0.000000,0.000000}%
\pgfsetfillcolor{currentfill}%
\pgfsetlinewidth{0.401500pt}%
\definecolor{currentstroke}{rgb}{0.000000,0.000000,0.000000}%
\pgfsetstrokecolor{currentstroke}%
\pgfsetdash{}{0pt}%
\pgfsys@defobject{currentmarker}{\pgfqpoint{-0.020833in}{0.000000in}}{\pgfqpoint{-0.000000in}{0.000000in}}{%
\pgfpathmoveto{\pgfqpoint{-0.000000in}{0.000000in}}%
\pgfpathlineto{\pgfqpoint{-0.020833in}{0.000000in}}%
\pgfusepath{stroke,fill}%
}%
\begin{pgfscope}%
\pgfsys@transformshift{3.038110in}{0.919852in}%
\pgfsys@useobject{currentmarker}{}%
\end{pgfscope}%
\end{pgfscope}%
\begin{pgfscope}%
\pgfsetbuttcap%
\pgfsetroundjoin%
\definecolor{currentfill}{rgb}{0.000000,0.000000,0.000000}%
\pgfsetfillcolor{currentfill}%
\pgfsetlinewidth{0.401500pt}%
\definecolor{currentstroke}{rgb}{0.000000,0.000000,0.000000}%
\pgfsetstrokecolor{currentstroke}%
\pgfsetdash{}{0pt}%
\pgfsys@defobject{currentmarker}{\pgfqpoint{0.000000in}{0.000000in}}{\pgfqpoint{0.020833in}{0.000000in}}{%
\pgfpathmoveto{\pgfqpoint{0.000000in}{0.000000in}}%
\pgfpathlineto{\pgfqpoint{0.020833in}{0.000000in}}%
\pgfusepath{stroke,fill}%
}%
\begin{pgfscope}%
\pgfsys@transformshift{0.650000in}{1.072064in}%
\pgfsys@useobject{currentmarker}{}%
\end{pgfscope}%
\end{pgfscope}%
\begin{pgfscope}%
\pgfsetbuttcap%
\pgfsetroundjoin%
\definecolor{currentfill}{rgb}{0.000000,0.000000,0.000000}%
\pgfsetfillcolor{currentfill}%
\pgfsetlinewidth{0.401500pt}%
\definecolor{currentstroke}{rgb}{0.000000,0.000000,0.000000}%
\pgfsetstrokecolor{currentstroke}%
\pgfsetdash{}{0pt}%
\pgfsys@defobject{currentmarker}{\pgfqpoint{-0.020833in}{0.000000in}}{\pgfqpoint{-0.000000in}{0.000000in}}{%
\pgfpathmoveto{\pgfqpoint{-0.000000in}{0.000000in}}%
\pgfpathlineto{\pgfqpoint{-0.020833in}{0.000000in}}%
\pgfusepath{stroke,fill}%
}%
\begin{pgfscope}%
\pgfsys@transformshift{3.038110in}{1.072064in}%
\pgfsys@useobject{currentmarker}{}%
\end{pgfscope}%
\end{pgfscope}%
\begin{pgfscope}%
\pgfsetbuttcap%
\pgfsetroundjoin%
\definecolor{currentfill}{rgb}{0.000000,0.000000,0.000000}%
\pgfsetfillcolor{currentfill}%
\pgfsetlinewidth{0.401500pt}%
\definecolor{currentstroke}{rgb}{0.000000,0.000000,0.000000}%
\pgfsetstrokecolor{currentstroke}%
\pgfsetdash{}{0pt}%
\pgfsys@defobject{currentmarker}{\pgfqpoint{0.000000in}{0.000000in}}{\pgfqpoint{0.020833in}{0.000000in}}{%
\pgfpathmoveto{\pgfqpoint{0.000000in}{0.000000in}}%
\pgfpathlineto{\pgfqpoint{0.020833in}{0.000000in}}%
\pgfusepath{stroke,fill}%
}%
\begin{pgfscope}%
\pgfsys@transformshift{0.650000in}{1.148170in}%
\pgfsys@useobject{currentmarker}{}%
\end{pgfscope}%
\end{pgfscope}%
\begin{pgfscope}%
\pgfsetbuttcap%
\pgfsetroundjoin%
\definecolor{currentfill}{rgb}{0.000000,0.000000,0.000000}%
\pgfsetfillcolor{currentfill}%
\pgfsetlinewidth{0.401500pt}%
\definecolor{currentstroke}{rgb}{0.000000,0.000000,0.000000}%
\pgfsetstrokecolor{currentstroke}%
\pgfsetdash{}{0pt}%
\pgfsys@defobject{currentmarker}{\pgfqpoint{-0.020833in}{0.000000in}}{\pgfqpoint{-0.000000in}{0.000000in}}{%
\pgfpathmoveto{\pgfqpoint{-0.000000in}{0.000000in}}%
\pgfpathlineto{\pgfqpoint{-0.020833in}{0.000000in}}%
\pgfusepath{stroke,fill}%
}%
\begin{pgfscope}%
\pgfsys@transformshift{3.038110in}{1.148170in}%
\pgfsys@useobject{currentmarker}{}%
\end{pgfscope}%
\end{pgfscope}%
\begin{pgfscope}%
\pgfsetbuttcap%
\pgfsetroundjoin%
\definecolor{currentfill}{rgb}{0.000000,0.000000,0.000000}%
\pgfsetfillcolor{currentfill}%
\pgfsetlinewidth{0.401500pt}%
\definecolor{currentstroke}{rgb}{0.000000,0.000000,0.000000}%
\pgfsetstrokecolor{currentstroke}%
\pgfsetdash{}{0pt}%
\pgfsys@defobject{currentmarker}{\pgfqpoint{0.000000in}{0.000000in}}{\pgfqpoint{0.020833in}{0.000000in}}{%
\pgfpathmoveto{\pgfqpoint{0.000000in}{0.000000in}}%
\pgfpathlineto{\pgfqpoint{0.020833in}{0.000000in}}%
\pgfusepath{stroke,fill}%
}%
\begin{pgfscope}%
\pgfsys@transformshift{0.650000in}{1.224276in}%
\pgfsys@useobject{currentmarker}{}%
\end{pgfscope}%
\end{pgfscope}%
\begin{pgfscope}%
\pgfsetbuttcap%
\pgfsetroundjoin%
\definecolor{currentfill}{rgb}{0.000000,0.000000,0.000000}%
\pgfsetfillcolor{currentfill}%
\pgfsetlinewidth{0.401500pt}%
\definecolor{currentstroke}{rgb}{0.000000,0.000000,0.000000}%
\pgfsetstrokecolor{currentstroke}%
\pgfsetdash{}{0pt}%
\pgfsys@defobject{currentmarker}{\pgfqpoint{-0.020833in}{0.000000in}}{\pgfqpoint{-0.000000in}{0.000000in}}{%
\pgfpathmoveto{\pgfqpoint{-0.000000in}{0.000000in}}%
\pgfpathlineto{\pgfqpoint{-0.020833in}{0.000000in}}%
\pgfusepath{stroke,fill}%
}%
\begin{pgfscope}%
\pgfsys@transformshift{3.038110in}{1.224276in}%
\pgfsys@useobject{currentmarker}{}%
\end{pgfscope}%
\end{pgfscope}%
\begin{pgfscope}%
\pgfsetbuttcap%
\pgfsetroundjoin%
\definecolor{currentfill}{rgb}{0.000000,0.000000,0.000000}%
\pgfsetfillcolor{currentfill}%
\pgfsetlinewidth{0.401500pt}%
\definecolor{currentstroke}{rgb}{0.000000,0.000000,0.000000}%
\pgfsetstrokecolor{currentstroke}%
\pgfsetdash{}{0pt}%
\pgfsys@defobject{currentmarker}{\pgfqpoint{0.000000in}{0.000000in}}{\pgfqpoint{0.020833in}{0.000000in}}{%
\pgfpathmoveto{\pgfqpoint{0.000000in}{0.000000in}}%
\pgfpathlineto{\pgfqpoint{0.020833in}{0.000000in}}%
\pgfusepath{stroke,fill}%
}%
\begin{pgfscope}%
\pgfsys@transformshift{0.650000in}{1.300383in}%
\pgfsys@useobject{currentmarker}{}%
\end{pgfscope}%
\end{pgfscope}%
\begin{pgfscope}%
\pgfsetbuttcap%
\pgfsetroundjoin%
\definecolor{currentfill}{rgb}{0.000000,0.000000,0.000000}%
\pgfsetfillcolor{currentfill}%
\pgfsetlinewidth{0.401500pt}%
\definecolor{currentstroke}{rgb}{0.000000,0.000000,0.000000}%
\pgfsetstrokecolor{currentstroke}%
\pgfsetdash{}{0pt}%
\pgfsys@defobject{currentmarker}{\pgfqpoint{-0.020833in}{0.000000in}}{\pgfqpoint{-0.000000in}{0.000000in}}{%
\pgfpathmoveto{\pgfqpoint{-0.000000in}{0.000000in}}%
\pgfpathlineto{\pgfqpoint{-0.020833in}{0.000000in}}%
\pgfusepath{stroke,fill}%
}%
\begin{pgfscope}%
\pgfsys@transformshift{3.038110in}{1.300383in}%
\pgfsys@useobject{currentmarker}{}%
\end{pgfscope}%
\end{pgfscope}%
\begin{pgfscope}%
\pgfsetbuttcap%
\pgfsetroundjoin%
\definecolor{currentfill}{rgb}{0.000000,0.000000,0.000000}%
\pgfsetfillcolor{currentfill}%
\pgfsetlinewidth{0.401500pt}%
\definecolor{currentstroke}{rgb}{0.000000,0.000000,0.000000}%
\pgfsetstrokecolor{currentstroke}%
\pgfsetdash{}{0pt}%
\pgfsys@defobject{currentmarker}{\pgfqpoint{0.000000in}{0.000000in}}{\pgfqpoint{0.020833in}{0.000000in}}{%
\pgfpathmoveto{\pgfqpoint{0.000000in}{0.000000in}}%
\pgfpathlineto{\pgfqpoint{0.020833in}{0.000000in}}%
\pgfusepath{stroke,fill}%
}%
\begin{pgfscope}%
\pgfsys@transformshift{0.650000in}{1.452595in}%
\pgfsys@useobject{currentmarker}{}%
\end{pgfscope}%
\end{pgfscope}%
\begin{pgfscope}%
\pgfsetbuttcap%
\pgfsetroundjoin%
\definecolor{currentfill}{rgb}{0.000000,0.000000,0.000000}%
\pgfsetfillcolor{currentfill}%
\pgfsetlinewidth{0.401500pt}%
\definecolor{currentstroke}{rgb}{0.000000,0.000000,0.000000}%
\pgfsetstrokecolor{currentstroke}%
\pgfsetdash{}{0pt}%
\pgfsys@defobject{currentmarker}{\pgfqpoint{-0.020833in}{0.000000in}}{\pgfqpoint{-0.000000in}{0.000000in}}{%
\pgfpathmoveto{\pgfqpoint{-0.000000in}{0.000000in}}%
\pgfpathlineto{\pgfqpoint{-0.020833in}{0.000000in}}%
\pgfusepath{stroke,fill}%
}%
\begin{pgfscope}%
\pgfsys@transformshift{3.038110in}{1.452595in}%
\pgfsys@useobject{currentmarker}{}%
\end{pgfscope}%
\end{pgfscope}%
\begin{pgfscope}%
\pgfsetbuttcap%
\pgfsetroundjoin%
\definecolor{currentfill}{rgb}{0.000000,0.000000,0.000000}%
\pgfsetfillcolor{currentfill}%
\pgfsetlinewidth{0.401500pt}%
\definecolor{currentstroke}{rgb}{0.000000,0.000000,0.000000}%
\pgfsetstrokecolor{currentstroke}%
\pgfsetdash{}{0pt}%
\pgfsys@defobject{currentmarker}{\pgfqpoint{0.000000in}{0.000000in}}{\pgfqpoint{0.020833in}{0.000000in}}{%
\pgfpathmoveto{\pgfqpoint{0.000000in}{0.000000in}}%
\pgfpathlineto{\pgfqpoint{0.020833in}{0.000000in}}%
\pgfusepath{stroke,fill}%
}%
\begin{pgfscope}%
\pgfsys@transformshift{0.650000in}{1.528701in}%
\pgfsys@useobject{currentmarker}{}%
\end{pgfscope}%
\end{pgfscope}%
\begin{pgfscope}%
\pgfsetbuttcap%
\pgfsetroundjoin%
\definecolor{currentfill}{rgb}{0.000000,0.000000,0.000000}%
\pgfsetfillcolor{currentfill}%
\pgfsetlinewidth{0.401500pt}%
\definecolor{currentstroke}{rgb}{0.000000,0.000000,0.000000}%
\pgfsetstrokecolor{currentstroke}%
\pgfsetdash{}{0pt}%
\pgfsys@defobject{currentmarker}{\pgfqpoint{-0.020833in}{0.000000in}}{\pgfqpoint{-0.000000in}{0.000000in}}{%
\pgfpathmoveto{\pgfqpoint{-0.000000in}{0.000000in}}%
\pgfpathlineto{\pgfqpoint{-0.020833in}{0.000000in}}%
\pgfusepath{stroke,fill}%
}%
\begin{pgfscope}%
\pgfsys@transformshift{3.038110in}{1.528701in}%
\pgfsys@useobject{currentmarker}{}%
\end{pgfscope}%
\end{pgfscope}%
\begin{pgfscope}%
\pgfsetbuttcap%
\pgfsetroundjoin%
\definecolor{currentfill}{rgb}{0.000000,0.000000,0.000000}%
\pgfsetfillcolor{currentfill}%
\pgfsetlinewidth{0.401500pt}%
\definecolor{currentstroke}{rgb}{0.000000,0.000000,0.000000}%
\pgfsetstrokecolor{currentstroke}%
\pgfsetdash{}{0pt}%
\pgfsys@defobject{currentmarker}{\pgfqpoint{0.000000in}{0.000000in}}{\pgfqpoint{0.020833in}{0.000000in}}{%
\pgfpathmoveto{\pgfqpoint{0.000000in}{0.000000in}}%
\pgfpathlineto{\pgfqpoint{0.020833in}{0.000000in}}%
\pgfusepath{stroke,fill}%
}%
\begin{pgfscope}%
\pgfsys@transformshift{0.650000in}{1.604807in}%
\pgfsys@useobject{currentmarker}{}%
\end{pgfscope}%
\end{pgfscope}%
\begin{pgfscope}%
\pgfsetbuttcap%
\pgfsetroundjoin%
\definecolor{currentfill}{rgb}{0.000000,0.000000,0.000000}%
\pgfsetfillcolor{currentfill}%
\pgfsetlinewidth{0.401500pt}%
\definecolor{currentstroke}{rgb}{0.000000,0.000000,0.000000}%
\pgfsetstrokecolor{currentstroke}%
\pgfsetdash{}{0pt}%
\pgfsys@defobject{currentmarker}{\pgfqpoint{-0.020833in}{0.000000in}}{\pgfqpoint{-0.000000in}{0.000000in}}{%
\pgfpathmoveto{\pgfqpoint{-0.000000in}{0.000000in}}%
\pgfpathlineto{\pgfqpoint{-0.020833in}{0.000000in}}%
\pgfusepath{stroke,fill}%
}%
\begin{pgfscope}%
\pgfsys@transformshift{3.038110in}{1.604807in}%
\pgfsys@useobject{currentmarker}{}%
\end{pgfscope}%
\end{pgfscope}%
\begin{pgfscope}%
\pgfsetbuttcap%
\pgfsetroundjoin%
\definecolor{currentfill}{rgb}{0.000000,0.000000,0.000000}%
\pgfsetfillcolor{currentfill}%
\pgfsetlinewidth{0.401500pt}%
\definecolor{currentstroke}{rgb}{0.000000,0.000000,0.000000}%
\pgfsetstrokecolor{currentstroke}%
\pgfsetdash{}{0pt}%
\pgfsys@defobject{currentmarker}{\pgfqpoint{0.000000in}{0.000000in}}{\pgfqpoint{0.020833in}{0.000000in}}{%
\pgfpathmoveto{\pgfqpoint{0.000000in}{0.000000in}}%
\pgfpathlineto{\pgfqpoint{0.020833in}{0.000000in}}%
\pgfusepath{stroke,fill}%
}%
\begin{pgfscope}%
\pgfsys@transformshift{0.650000in}{1.680914in}%
\pgfsys@useobject{currentmarker}{}%
\end{pgfscope}%
\end{pgfscope}%
\begin{pgfscope}%
\pgfsetbuttcap%
\pgfsetroundjoin%
\definecolor{currentfill}{rgb}{0.000000,0.000000,0.000000}%
\pgfsetfillcolor{currentfill}%
\pgfsetlinewidth{0.401500pt}%
\definecolor{currentstroke}{rgb}{0.000000,0.000000,0.000000}%
\pgfsetstrokecolor{currentstroke}%
\pgfsetdash{}{0pt}%
\pgfsys@defobject{currentmarker}{\pgfqpoint{-0.020833in}{0.000000in}}{\pgfqpoint{-0.000000in}{0.000000in}}{%
\pgfpathmoveto{\pgfqpoint{-0.000000in}{0.000000in}}%
\pgfpathlineto{\pgfqpoint{-0.020833in}{0.000000in}}%
\pgfusepath{stroke,fill}%
}%
\begin{pgfscope}%
\pgfsys@transformshift{3.038110in}{1.680914in}%
\pgfsys@useobject{currentmarker}{}%
\end{pgfscope}%
\end{pgfscope}%
\begin{pgfscope}%
\pgfsetbuttcap%
\pgfsetroundjoin%
\definecolor{currentfill}{rgb}{0.000000,0.000000,0.000000}%
\pgfsetfillcolor{currentfill}%
\pgfsetlinewidth{0.401500pt}%
\definecolor{currentstroke}{rgb}{0.000000,0.000000,0.000000}%
\pgfsetstrokecolor{currentstroke}%
\pgfsetdash{}{0pt}%
\pgfsys@defobject{currentmarker}{\pgfqpoint{0.000000in}{0.000000in}}{\pgfqpoint{0.020833in}{0.000000in}}{%
\pgfpathmoveto{\pgfqpoint{0.000000in}{0.000000in}}%
\pgfpathlineto{\pgfqpoint{0.020833in}{0.000000in}}%
\pgfusepath{stroke,fill}%
}%
\begin{pgfscope}%
\pgfsys@transformshift{0.650000in}{1.833126in}%
\pgfsys@useobject{currentmarker}{}%
\end{pgfscope}%
\end{pgfscope}%
\begin{pgfscope}%
\pgfsetbuttcap%
\pgfsetroundjoin%
\definecolor{currentfill}{rgb}{0.000000,0.000000,0.000000}%
\pgfsetfillcolor{currentfill}%
\pgfsetlinewidth{0.401500pt}%
\definecolor{currentstroke}{rgb}{0.000000,0.000000,0.000000}%
\pgfsetstrokecolor{currentstroke}%
\pgfsetdash{}{0pt}%
\pgfsys@defobject{currentmarker}{\pgfqpoint{-0.020833in}{0.000000in}}{\pgfqpoint{-0.000000in}{0.000000in}}{%
\pgfpathmoveto{\pgfqpoint{-0.000000in}{0.000000in}}%
\pgfpathlineto{\pgfqpoint{-0.020833in}{0.000000in}}%
\pgfusepath{stroke,fill}%
}%
\begin{pgfscope}%
\pgfsys@transformshift{3.038110in}{1.833126in}%
\pgfsys@useobject{currentmarker}{}%
\end{pgfscope}%
\end{pgfscope}%
\begin{pgfscope}%
\pgfsetbuttcap%
\pgfsetroundjoin%
\definecolor{currentfill}{rgb}{0.000000,0.000000,0.000000}%
\pgfsetfillcolor{currentfill}%
\pgfsetlinewidth{0.401500pt}%
\definecolor{currentstroke}{rgb}{0.000000,0.000000,0.000000}%
\pgfsetstrokecolor{currentstroke}%
\pgfsetdash{}{0pt}%
\pgfsys@defobject{currentmarker}{\pgfqpoint{0.000000in}{0.000000in}}{\pgfqpoint{0.020833in}{0.000000in}}{%
\pgfpathmoveto{\pgfqpoint{0.000000in}{0.000000in}}%
\pgfpathlineto{\pgfqpoint{0.020833in}{0.000000in}}%
\pgfusepath{stroke,fill}%
}%
\begin{pgfscope}%
\pgfsys@transformshift{0.650000in}{1.909232in}%
\pgfsys@useobject{currentmarker}{}%
\end{pgfscope}%
\end{pgfscope}%
\begin{pgfscope}%
\pgfsetbuttcap%
\pgfsetroundjoin%
\definecolor{currentfill}{rgb}{0.000000,0.000000,0.000000}%
\pgfsetfillcolor{currentfill}%
\pgfsetlinewidth{0.401500pt}%
\definecolor{currentstroke}{rgb}{0.000000,0.000000,0.000000}%
\pgfsetstrokecolor{currentstroke}%
\pgfsetdash{}{0pt}%
\pgfsys@defobject{currentmarker}{\pgfqpoint{-0.020833in}{0.000000in}}{\pgfqpoint{-0.000000in}{0.000000in}}{%
\pgfpathmoveto{\pgfqpoint{-0.000000in}{0.000000in}}%
\pgfpathlineto{\pgfqpoint{-0.020833in}{0.000000in}}%
\pgfusepath{stroke,fill}%
}%
\begin{pgfscope}%
\pgfsys@transformshift{3.038110in}{1.909232in}%
\pgfsys@useobject{currentmarker}{}%
\end{pgfscope}%
\end{pgfscope}%
\begin{pgfscope}%
\pgfsetbuttcap%
\pgfsetroundjoin%
\definecolor{currentfill}{rgb}{0.000000,0.000000,0.000000}%
\pgfsetfillcolor{currentfill}%
\pgfsetlinewidth{0.401500pt}%
\definecolor{currentstroke}{rgb}{0.000000,0.000000,0.000000}%
\pgfsetstrokecolor{currentstroke}%
\pgfsetdash{}{0pt}%
\pgfsys@defobject{currentmarker}{\pgfqpoint{0.000000in}{0.000000in}}{\pgfqpoint{0.020833in}{0.000000in}}{%
\pgfpathmoveto{\pgfqpoint{0.000000in}{0.000000in}}%
\pgfpathlineto{\pgfqpoint{0.020833in}{0.000000in}}%
\pgfusepath{stroke,fill}%
}%
\begin{pgfscope}%
\pgfsys@transformshift{0.650000in}{1.985338in}%
\pgfsys@useobject{currentmarker}{}%
\end{pgfscope}%
\end{pgfscope}%
\begin{pgfscope}%
\pgfsetbuttcap%
\pgfsetroundjoin%
\definecolor{currentfill}{rgb}{0.000000,0.000000,0.000000}%
\pgfsetfillcolor{currentfill}%
\pgfsetlinewidth{0.401500pt}%
\definecolor{currentstroke}{rgb}{0.000000,0.000000,0.000000}%
\pgfsetstrokecolor{currentstroke}%
\pgfsetdash{}{0pt}%
\pgfsys@defobject{currentmarker}{\pgfqpoint{-0.020833in}{0.000000in}}{\pgfqpoint{-0.000000in}{0.000000in}}{%
\pgfpathmoveto{\pgfqpoint{-0.000000in}{0.000000in}}%
\pgfpathlineto{\pgfqpoint{-0.020833in}{0.000000in}}%
\pgfusepath{stroke,fill}%
}%
\begin{pgfscope}%
\pgfsys@transformshift{3.038110in}{1.985338in}%
\pgfsys@useobject{currentmarker}{}%
\end{pgfscope}%
\end{pgfscope}%
\begin{pgfscope}%
\pgfsetbuttcap%
\pgfsetroundjoin%
\definecolor{currentfill}{rgb}{0.000000,0.000000,0.000000}%
\pgfsetfillcolor{currentfill}%
\pgfsetlinewidth{0.401500pt}%
\definecolor{currentstroke}{rgb}{0.000000,0.000000,0.000000}%
\pgfsetstrokecolor{currentstroke}%
\pgfsetdash{}{0pt}%
\pgfsys@defobject{currentmarker}{\pgfqpoint{0.000000in}{0.000000in}}{\pgfqpoint{0.020833in}{0.000000in}}{%
\pgfpathmoveto{\pgfqpoint{0.000000in}{0.000000in}}%
\pgfpathlineto{\pgfqpoint{0.020833in}{0.000000in}}%
\pgfusepath{stroke,fill}%
}%
\begin{pgfscope}%
\pgfsys@transformshift{0.650000in}{2.061445in}%
\pgfsys@useobject{currentmarker}{}%
\end{pgfscope}%
\end{pgfscope}%
\begin{pgfscope}%
\pgfsetbuttcap%
\pgfsetroundjoin%
\definecolor{currentfill}{rgb}{0.000000,0.000000,0.000000}%
\pgfsetfillcolor{currentfill}%
\pgfsetlinewidth{0.401500pt}%
\definecolor{currentstroke}{rgb}{0.000000,0.000000,0.000000}%
\pgfsetstrokecolor{currentstroke}%
\pgfsetdash{}{0pt}%
\pgfsys@defobject{currentmarker}{\pgfqpoint{-0.020833in}{0.000000in}}{\pgfqpoint{-0.000000in}{0.000000in}}{%
\pgfpathmoveto{\pgfqpoint{-0.000000in}{0.000000in}}%
\pgfpathlineto{\pgfqpoint{-0.020833in}{0.000000in}}%
\pgfusepath{stroke,fill}%
}%
\begin{pgfscope}%
\pgfsys@transformshift{3.038110in}{2.061445in}%
\pgfsys@useobject{currentmarker}{}%
\end{pgfscope}%
\end{pgfscope}%
\begin{pgfscope}%
\pgfsetbuttcap%
\pgfsetroundjoin%
\definecolor{currentfill}{rgb}{0.000000,0.000000,0.000000}%
\pgfsetfillcolor{currentfill}%
\pgfsetlinewidth{0.401500pt}%
\definecolor{currentstroke}{rgb}{0.000000,0.000000,0.000000}%
\pgfsetstrokecolor{currentstroke}%
\pgfsetdash{}{0pt}%
\pgfsys@defobject{currentmarker}{\pgfqpoint{0.000000in}{0.000000in}}{\pgfqpoint{0.020833in}{0.000000in}}{%
\pgfpathmoveto{\pgfqpoint{0.000000in}{0.000000in}}%
\pgfpathlineto{\pgfqpoint{0.020833in}{0.000000in}}%
\pgfusepath{stroke,fill}%
}%
\begin{pgfscope}%
\pgfsys@transformshift{0.650000in}{2.213657in}%
\pgfsys@useobject{currentmarker}{}%
\end{pgfscope}%
\end{pgfscope}%
\begin{pgfscope}%
\pgfsetbuttcap%
\pgfsetroundjoin%
\definecolor{currentfill}{rgb}{0.000000,0.000000,0.000000}%
\pgfsetfillcolor{currentfill}%
\pgfsetlinewidth{0.401500pt}%
\definecolor{currentstroke}{rgb}{0.000000,0.000000,0.000000}%
\pgfsetstrokecolor{currentstroke}%
\pgfsetdash{}{0pt}%
\pgfsys@defobject{currentmarker}{\pgfqpoint{-0.020833in}{0.000000in}}{\pgfqpoint{-0.000000in}{0.000000in}}{%
\pgfpathmoveto{\pgfqpoint{-0.000000in}{0.000000in}}%
\pgfpathlineto{\pgfqpoint{-0.020833in}{0.000000in}}%
\pgfusepath{stroke,fill}%
}%
\begin{pgfscope}%
\pgfsys@transformshift{3.038110in}{2.213657in}%
\pgfsys@useobject{currentmarker}{}%
\end{pgfscope}%
\end{pgfscope}%
\begin{pgfscope}%
\pgfsetbuttcap%
\pgfsetroundjoin%
\definecolor{currentfill}{rgb}{0.000000,0.000000,0.000000}%
\pgfsetfillcolor{currentfill}%
\pgfsetlinewidth{0.401500pt}%
\definecolor{currentstroke}{rgb}{0.000000,0.000000,0.000000}%
\pgfsetstrokecolor{currentstroke}%
\pgfsetdash{}{0pt}%
\pgfsys@defobject{currentmarker}{\pgfqpoint{0.000000in}{0.000000in}}{\pgfqpoint{0.020833in}{0.000000in}}{%
\pgfpathmoveto{\pgfqpoint{0.000000in}{0.000000in}}%
\pgfpathlineto{\pgfqpoint{0.020833in}{0.000000in}}%
\pgfusepath{stroke,fill}%
}%
\begin{pgfscope}%
\pgfsys@transformshift{0.650000in}{2.289763in}%
\pgfsys@useobject{currentmarker}{}%
\end{pgfscope}%
\end{pgfscope}%
\begin{pgfscope}%
\pgfsetbuttcap%
\pgfsetroundjoin%
\definecolor{currentfill}{rgb}{0.000000,0.000000,0.000000}%
\pgfsetfillcolor{currentfill}%
\pgfsetlinewidth{0.401500pt}%
\definecolor{currentstroke}{rgb}{0.000000,0.000000,0.000000}%
\pgfsetstrokecolor{currentstroke}%
\pgfsetdash{}{0pt}%
\pgfsys@defobject{currentmarker}{\pgfqpoint{-0.020833in}{0.000000in}}{\pgfqpoint{-0.000000in}{0.000000in}}{%
\pgfpathmoveto{\pgfqpoint{-0.000000in}{0.000000in}}%
\pgfpathlineto{\pgfqpoint{-0.020833in}{0.000000in}}%
\pgfusepath{stroke,fill}%
}%
\begin{pgfscope}%
\pgfsys@transformshift{3.038110in}{2.289763in}%
\pgfsys@useobject{currentmarker}{}%
\end{pgfscope}%
\end{pgfscope}%
\begin{pgfscope}%
\pgfsetbuttcap%
\pgfsetroundjoin%
\definecolor{currentfill}{rgb}{0.000000,0.000000,0.000000}%
\pgfsetfillcolor{currentfill}%
\pgfsetlinewidth{0.401500pt}%
\definecolor{currentstroke}{rgb}{0.000000,0.000000,0.000000}%
\pgfsetstrokecolor{currentstroke}%
\pgfsetdash{}{0pt}%
\pgfsys@defobject{currentmarker}{\pgfqpoint{0.000000in}{0.000000in}}{\pgfqpoint{0.020833in}{0.000000in}}{%
\pgfpathmoveto{\pgfqpoint{0.000000in}{0.000000in}}%
\pgfpathlineto{\pgfqpoint{0.020833in}{0.000000in}}%
\pgfusepath{stroke,fill}%
}%
\begin{pgfscope}%
\pgfsys@transformshift{0.650000in}{2.365869in}%
\pgfsys@useobject{currentmarker}{}%
\end{pgfscope}%
\end{pgfscope}%
\begin{pgfscope}%
\pgfsetbuttcap%
\pgfsetroundjoin%
\definecolor{currentfill}{rgb}{0.000000,0.000000,0.000000}%
\pgfsetfillcolor{currentfill}%
\pgfsetlinewidth{0.401500pt}%
\definecolor{currentstroke}{rgb}{0.000000,0.000000,0.000000}%
\pgfsetstrokecolor{currentstroke}%
\pgfsetdash{}{0pt}%
\pgfsys@defobject{currentmarker}{\pgfqpoint{-0.020833in}{0.000000in}}{\pgfqpoint{-0.000000in}{0.000000in}}{%
\pgfpathmoveto{\pgfqpoint{-0.000000in}{0.000000in}}%
\pgfpathlineto{\pgfqpoint{-0.020833in}{0.000000in}}%
\pgfusepath{stroke,fill}%
}%
\begin{pgfscope}%
\pgfsys@transformshift{3.038110in}{2.365869in}%
\pgfsys@useobject{currentmarker}{}%
\end{pgfscope}%
\end{pgfscope}%
\begin{pgfscope}%
\definecolor{textcolor}{rgb}{0.000000,0.000000,0.000000}%
\pgfsetstrokecolor{textcolor}%
\pgfsetfillcolor{textcolor}%
\pgftext[x=0.260995in,y=1.417244in,,bottom,rotate=90.000000]{\color{textcolor}{\rmfamily\fontsize{12.000000}{14.400000}\selectfont\catcode`\^=\active\def^{\ifmmode\sp\else\^{}\fi}\catcode`\%=\active\def%{\%}TKE budget}}%
\end{pgfscope}%
\begin{pgfscope}%
\pgfpathrectangle{\pgfqpoint{0.650000in}{0.420000in}}{\pgfqpoint{2.388110in}{1.994488in}}%
\pgfusepath{clip}%
\pgfsetrectcap%
\pgfsetroundjoin%
\pgfsetlinewidth{0.602250pt}%
\definecolor{currentstroke}{rgb}{0.000000,0.000000,0.000000}%
\pgfsetstrokecolor{currentstroke}%
\pgfsetdash{}{0pt}%
\pgfpathmoveto{\pgfqpoint{0.645000in}{1.377182in}}%
\pgfpathlineto{\pgfqpoint{0.742708in}{1.377695in}}%
\pgfpathlineto{\pgfqpoint{0.889821in}{1.381173in}}%
\pgfpathlineto{\pgfqpoint{1.010016in}{1.390652in}}%
\pgfpathlineto{\pgfqpoint{1.111637in}{1.412173in}}%
\pgfpathlineto{\pgfqpoint{1.199660in}{1.454175in}}%
\pgfpathlineto{\pgfqpoint{1.277299in}{1.525569in}}%
\pgfpathlineto{\pgfqpoint{1.346745in}{1.631644in}}%
\pgfpathlineto{\pgfqpoint{1.409562in}{1.769287in}}%
\pgfpathlineto{\pgfqpoint{1.466906in}{1.924712in}}%
\pgfpathlineto{\pgfqpoint{1.519654in}{2.076258in}}%
\pgfpathlineto{\pgfqpoint{1.568486in}{2.201505in}}%
\pgfpathlineto{\pgfqpoint{1.613944in}{2.284710in}}%
\pgfpathlineto{\pgfqpoint{1.656463in}{2.320591in}}%
\pgfpathlineto{\pgfqpoint{1.696400in}{2.313425in}}%
\pgfpathlineto{\pgfqpoint{1.734049in}{2.273226in}}%
\pgfpathlineto{\pgfqpoint{1.769658in}{2.211626in}}%
\pgfpathlineto{\pgfqpoint{1.803436in}{2.139004in}}%
\pgfpathlineto{\pgfqpoint{1.835562in}{2.063229in}}%
\pgfpathlineto{\pgfqpoint{1.866188in}{1.989542in}}%
\pgfpathlineto{\pgfqpoint{1.895450in}{1.921031in}}%
\pgfpathlineto{\pgfqpoint{1.923461in}{1.859219in}}%
\pgfpathlineto{\pgfqpoint{1.950325in}{1.804650in}}%
\pgfpathlineto{\pgfqpoint{1.976131in}{1.757203in}}%
\pgfpathlineto{\pgfqpoint{2.000959in}{1.716403in}}%
\pgfpathlineto{\pgfqpoint{2.024880in}{1.681631in}}%
\pgfpathlineto{\pgfqpoint{2.047957in}{1.652173in}}%
\pgfpathlineto{\pgfqpoint{2.070247in}{1.627262in}}%
\pgfpathlineto{\pgfqpoint{2.091803in}{1.606271in}}%
\pgfpathlineto{\pgfqpoint{2.112670in}{1.588576in}}%
\pgfpathlineto{\pgfqpoint{2.132891in}{1.573585in}}%
\pgfpathlineto{\pgfqpoint{2.152504in}{1.560792in}}%
\pgfpathlineto{\pgfqpoint{2.171545in}{1.549787in}}%
\pgfpathlineto{\pgfqpoint{2.190045in}{1.540225in}}%
\pgfpathlineto{\pgfqpoint{2.208034in}{1.531872in}}%
\pgfpathlineto{\pgfqpoint{2.242585in}{1.517962in}}%
\pgfpathlineto{\pgfqpoint{2.275392in}{1.506509in}}%
\pgfpathlineto{\pgfqpoint{2.306620in}{1.496763in}}%
\pgfpathlineto{\pgfqpoint{2.431033in}{1.460200in}}%
\pgfpathlineto{\pgfqpoint{2.479445in}{1.444413in}}%
\pgfpathlineto{\pgfqpoint{2.524453in}{1.428707in}}%
\pgfpathlineto{\pgfqpoint{2.624738in}{1.392210in}}%
\pgfpathlineto{\pgfqpoint{2.643017in}{1.386895in}}%
\pgfpathlineto{\pgfqpoint{2.660778in}{1.382755in}}%
\pgfpathlineto{\pgfqpoint{2.678050in}{1.379857in}}%
\pgfpathlineto{\pgfqpoint{2.694856in}{1.378088in}}%
\pgfpathlineto{\pgfqpoint{2.719243in}{1.376888in}}%
\pgfpathlineto{\pgfqpoint{2.757868in}{1.376507in}}%
\pgfpathlineto{\pgfqpoint{3.043110in}{1.376489in}}%
\pgfpathlineto{\pgfqpoint{3.043110in}{1.376489in}}%
\pgfusepath{stroke}%
\end{pgfscope}%
\begin{pgfscope}%
\pgfpathrectangle{\pgfqpoint{0.650000in}{0.420000in}}{\pgfqpoint{2.388110in}{1.994488in}}%
\pgfusepath{clip}%
\pgfsetbuttcap%
\pgfsetroundjoin%
\pgfsetlinewidth{0.602250pt}%
\definecolor{currentstroke}{rgb}{0.000000,0.000000,0.000000}%
\pgfsetstrokecolor{currentstroke}%
\pgfsetdash{{3.000000pt}{0.600000pt}}{0.000000pt}%
\pgfpathmoveto{\pgfqpoint{0.645000in}{1.376490in}}%
\pgfpathlineto{\pgfqpoint{1.346745in}{1.376886in}}%
\pgfpathlineto{\pgfqpoint{1.613944in}{1.378291in}}%
\pgfpathlineto{\pgfqpoint{1.866188in}{1.380247in}}%
\pgfpathlineto{\pgfqpoint{2.070247in}{1.377739in}}%
\pgfpathlineto{\pgfqpoint{2.132891in}{1.376407in}}%
\pgfpathlineto{\pgfqpoint{2.208034in}{1.375330in}}%
\pgfpathlineto{\pgfqpoint{2.259196in}{1.375127in}}%
\pgfpathlineto{\pgfqpoint{2.306620in}{1.374628in}}%
\pgfpathlineto{\pgfqpoint{2.350808in}{1.373211in}}%
\pgfpathlineto{\pgfqpoint{2.378674in}{1.372370in}}%
\pgfpathlineto{\pgfqpoint{2.405385in}{1.371146in}}%
\pgfpathlineto{\pgfqpoint{2.502345in}{1.368030in}}%
\pgfpathlineto{\pgfqpoint{2.624738in}{1.364222in}}%
\pgfpathlineto{\pgfqpoint{2.651960in}{1.365082in}}%
\pgfpathlineto{\pgfqpoint{2.678050in}{1.367266in}}%
\pgfpathlineto{\pgfqpoint{2.742707in}{1.373881in}}%
\pgfpathlineto{\pgfqpoint{2.772663in}{1.375373in}}%
\pgfpathlineto{\pgfqpoint{2.815015in}{1.376159in}}%
\pgfpathlineto{\pgfqpoint{2.915251in}{1.376458in}}%
\pgfpathlineto{\pgfqpoint{3.043110in}{1.376486in}}%
\pgfpathlineto{\pgfqpoint{3.043110in}{1.376486in}}%
\pgfusepath{stroke}%
\end{pgfscope}%
\begin{pgfscope}%
\pgfpathrectangle{\pgfqpoint{0.650000in}{0.420000in}}{\pgfqpoint{2.388110in}{1.994488in}}%
\pgfusepath{clip}%
\pgfsetbuttcap%
\pgfsetroundjoin%
\pgfsetlinewidth{0.602250pt}%
\definecolor{currentstroke}{rgb}{0.000000,0.000000,0.000000}%
\pgfsetstrokecolor{currentstroke}%
\pgfsetdash{{1.800000pt}{0.600000pt}{0.600000pt}{0.600000pt}}{0.000000pt}%
\pgfpathmoveto{\pgfqpoint{0.645000in}{0.585841in}}%
\pgfpathlineto{\pgfqpoint{0.742708in}{0.599455in}}%
\pgfpathlineto{\pgfqpoint{0.889821in}{0.625382in}}%
\pgfpathlineto{\pgfqpoint{1.010016in}{0.649866in}}%
\pgfpathlineto{\pgfqpoint{1.111637in}{0.674818in}}%
\pgfpathlineto{\pgfqpoint{1.199660in}{0.703404in}}%
\pgfpathlineto{\pgfqpoint{1.277299in}{0.737212in}}%
\pgfpathlineto{\pgfqpoint{1.409562in}{0.805947in}}%
\pgfpathlineto{\pgfqpoint{1.466906in}{0.828228in}}%
\pgfpathlineto{\pgfqpoint{1.519654in}{0.838400in}}%
\pgfpathlineto{\pgfqpoint{1.568486in}{0.839586in}}%
\pgfpathlineto{\pgfqpoint{1.613944in}{0.837581in}}%
\pgfpathlineto{\pgfqpoint{1.656463in}{0.837654in}}%
\pgfpathlineto{\pgfqpoint{1.696400in}{0.842778in}}%
\pgfpathlineto{\pgfqpoint{1.734049in}{0.853617in}}%
\pgfpathlineto{\pgfqpoint{1.769658in}{0.869437in}}%
\pgfpathlineto{\pgfqpoint{1.803436in}{0.888998in}}%
\pgfpathlineto{\pgfqpoint{1.835562in}{0.911041in}}%
\pgfpathlineto{\pgfqpoint{1.866188in}{0.934545in}}%
\pgfpathlineto{\pgfqpoint{1.895450in}{0.958769in}}%
\pgfpathlineto{\pgfqpoint{1.950325in}{1.007028in}}%
\pgfpathlineto{\pgfqpoint{2.024880in}{1.073550in}}%
\pgfpathlineto{\pgfqpoint{2.070247in}{1.111705in}}%
\pgfpathlineto{\pgfqpoint{2.091803in}{1.128803in}}%
\pgfpathlineto{\pgfqpoint{2.112670in}{1.144536in}}%
\pgfpathlineto{\pgfqpoint{2.132891in}{1.158983in}}%
\pgfpathlineto{\pgfqpoint{2.171545in}{1.184565in}}%
\pgfpathlineto{\pgfqpoint{2.208034in}{1.206236in}}%
\pgfpathlineto{\pgfqpoint{2.242585in}{1.224709in}}%
\pgfpathlineto{\pgfqpoint{2.275392in}{1.240586in}}%
\pgfpathlineto{\pgfqpoint{2.306620in}{1.254291in}}%
\pgfpathlineto{\pgfqpoint{2.350808in}{1.271988in}}%
\pgfpathlineto{\pgfqpoint{2.392168in}{1.287333in}}%
\pgfpathlineto{\pgfqpoint{2.443483in}{1.304940in}}%
\pgfpathlineto{\pgfqpoint{2.513495in}{1.327416in}}%
\pgfpathlineto{\pgfqpoint{2.566489in}{1.343366in}}%
\pgfpathlineto{\pgfqpoint{2.605912in}{1.354164in}}%
\pgfpathlineto{\pgfqpoint{2.633944in}{1.360919in}}%
\pgfpathlineto{\pgfqpoint{2.660778in}{1.366340in}}%
\pgfpathlineto{\pgfqpoint{2.686509in}{1.370418in}}%
\pgfpathlineto{\pgfqpoint{2.711220in}{1.373184in}}%
\pgfpathlineto{\pgfqpoint{2.734984in}{1.374867in}}%
\pgfpathlineto{\pgfqpoint{2.765310in}{1.375949in}}%
\pgfpathlineto{\pgfqpoint{2.815015in}{1.376424in}}%
\pgfpathlineto{\pgfqpoint{3.043110in}{1.376488in}}%
\pgfpathlineto{\pgfqpoint{3.043110in}{1.376488in}}%
\pgfusepath{stroke}%
\end{pgfscope}%
\begin{pgfscope}%
\pgfpathrectangle{\pgfqpoint{0.650000in}{0.420000in}}{\pgfqpoint{2.388110in}{1.994488in}}%
\pgfusepath{clip}%
\pgfsetbuttcap%
\pgfsetroundjoin%
\pgfsetlinewidth{0.602250pt}%
\definecolor{currentstroke}{rgb}{0.000000,0.000000,0.000000}%
\pgfsetstrokecolor{currentstroke}%
\pgfsetdash{{0.600000pt}{0.600000pt}}{0.000000pt}%
\pgfpathmoveto{\pgfqpoint{0.645000in}{1.407335in}}%
\pgfpathlineto{\pgfqpoint{0.742708in}{1.414525in}}%
\pgfpathlineto{\pgfqpoint{1.010016in}{1.437806in}}%
\pgfpathlineto{\pgfqpoint{1.111637in}{1.444111in}}%
\pgfpathlineto{\pgfqpoint{1.199660in}{1.446359in}}%
\pgfpathlineto{\pgfqpoint{1.277299in}{1.445079in}}%
\pgfpathlineto{\pgfqpoint{1.346745in}{1.441169in}}%
\pgfpathlineto{\pgfqpoint{1.409562in}{1.435430in}}%
\pgfpathlineto{\pgfqpoint{1.466906in}{1.428433in}}%
\pgfpathlineto{\pgfqpoint{1.519654in}{1.420430in}}%
\pgfpathlineto{\pgfqpoint{1.568486in}{1.411619in}}%
\pgfpathlineto{\pgfqpoint{1.613944in}{1.402249in}}%
\pgfpathlineto{\pgfqpoint{1.696400in}{1.383980in}}%
\pgfpathlineto{\pgfqpoint{1.734049in}{1.376348in}}%
\pgfpathlineto{\pgfqpoint{1.769658in}{1.370373in}}%
\pgfpathlineto{\pgfqpoint{1.803436in}{1.366293in}}%
\pgfpathlineto{\pgfqpoint{1.835562in}{1.364058in}}%
\pgfpathlineto{\pgfqpoint{1.866188in}{1.363489in}}%
\pgfpathlineto{\pgfqpoint{1.895450in}{1.364238in}}%
\pgfpathlineto{\pgfqpoint{1.923461in}{1.365968in}}%
\pgfpathlineto{\pgfqpoint{2.000959in}{1.373139in}}%
\pgfpathlineto{\pgfqpoint{2.047957in}{1.377212in}}%
\pgfpathlineto{\pgfqpoint{2.091803in}{1.379997in}}%
\pgfpathlineto{\pgfqpoint{2.132891in}{1.381119in}}%
\pgfpathlineto{\pgfqpoint{2.190045in}{1.380701in}}%
\pgfpathlineto{\pgfqpoint{2.242585in}{1.379652in}}%
\pgfpathlineto{\pgfqpoint{2.364891in}{1.376337in}}%
\pgfpathlineto{\pgfqpoint{2.479445in}{1.373909in}}%
\pgfpathlineto{\pgfqpoint{2.576578in}{1.371822in}}%
\pgfpathlineto{\pgfqpoint{2.615396in}{1.371729in}}%
\pgfpathlineto{\pgfqpoint{2.660778in}{1.372791in}}%
\pgfpathlineto{\pgfqpoint{2.765310in}{1.376612in}}%
\pgfpathlineto{\pgfqpoint{2.841691in}{1.376619in}}%
\pgfpathlineto{\pgfqpoint{3.043110in}{1.376492in}}%
\pgfpathlineto{\pgfqpoint{3.043110in}{1.376492in}}%
\pgfusepath{stroke}%
\end{pgfscope}%
\begin{pgfscope}%
\pgfpathrectangle{\pgfqpoint{0.650000in}{0.420000in}}{\pgfqpoint{2.388110in}{1.994488in}}%
\pgfusepath{clip}%
\pgfsetbuttcap%
\pgfsetroundjoin%
\pgfsetlinewidth{0.602250pt}%
\definecolor{currentstroke}{rgb}{0.000000,0.000000,0.000000}%
\pgfsetstrokecolor{currentstroke}%
\pgfsetdash{{3.000000pt}{3.000000pt}}{0.000000pt}%
\pgfpathmoveto{\pgfqpoint{0.645000in}{1.377502in}}%
\pgfpathlineto{\pgfqpoint{0.742708in}{1.378245in}}%
\pgfpathlineto{\pgfqpoint{0.889821in}{1.383135in}}%
\pgfpathlineto{\pgfqpoint{1.010016in}{1.395667in}}%
\pgfpathlineto{\pgfqpoint{1.111637in}{1.421077in}}%
\pgfpathlineto{\pgfqpoint{1.199660in}{1.461645in}}%
\pgfpathlineto{\pgfqpoint{1.277299in}{1.510673in}}%
\pgfpathlineto{\pgfqpoint{1.346745in}{1.550041in}}%
\pgfpathlineto{\pgfqpoint{1.409562in}{1.557196in}}%
\pgfpathlineto{\pgfqpoint{1.466906in}{1.518465in}}%
\pgfpathlineto{\pgfqpoint{1.519654in}{1.438414in}}%
\pgfpathlineto{\pgfqpoint{1.568486in}{1.337638in}}%
\pgfpathlineto{\pgfqpoint{1.613944in}{1.241734in}}%
\pgfpathlineto{\pgfqpoint{1.656463in}{1.170003in}}%
\pgfpathlineto{\pgfqpoint{1.696400in}{1.130505in}}%
\pgfpathlineto{\pgfqpoint{1.734049in}{1.121409in}}%
\pgfpathlineto{\pgfqpoint{1.769658in}{1.135451in}}%
\pgfpathlineto{\pgfqpoint{1.803436in}{1.163997in}}%
\pgfpathlineto{\pgfqpoint{1.835562in}{1.199465in}}%
\pgfpathlineto{\pgfqpoint{1.895450in}{1.270585in}}%
\pgfpathlineto{\pgfqpoint{1.923461in}{1.300588in}}%
\pgfpathlineto{\pgfqpoint{1.950325in}{1.325468in}}%
\pgfpathlineto{\pgfqpoint{1.976131in}{1.345369in}}%
\pgfpathlineto{\pgfqpoint{2.000959in}{1.360773in}}%
\pgfpathlineto{\pgfqpoint{2.024880in}{1.372090in}}%
\pgfpathlineto{\pgfqpoint{2.047957in}{1.380016in}}%
\pgfpathlineto{\pgfqpoint{2.070247in}{1.385371in}}%
\pgfpathlineto{\pgfqpoint{2.091803in}{1.388942in}}%
\pgfpathlineto{\pgfqpoint{2.112670in}{1.390920in}}%
\pgfpathlineto{\pgfqpoint{2.132891in}{1.391551in}}%
\pgfpathlineto{\pgfqpoint{2.171545in}{1.390719in}}%
\pgfpathlineto{\pgfqpoint{2.208034in}{1.388445in}}%
\pgfpathlineto{\pgfqpoint{2.259196in}{1.383274in}}%
\pgfpathlineto{\pgfqpoint{2.291194in}{1.380042in}}%
\pgfpathlineto{\pgfqpoint{2.350808in}{1.376572in}}%
\pgfpathlineto{\pgfqpoint{2.378674in}{1.375773in}}%
\pgfpathlineto{\pgfqpoint{2.443483in}{1.376243in}}%
\pgfpathlineto{\pgfqpoint{2.479445in}{1.378754in}}%
\pgfpathlineto{\pgfqpoint{2.513495in}{1.381893in}}%
\pgfpathlineto{\pgfqpoint{2.556238in}{1.387117in}}%
\pgfpathlineto{\pgfqpoint{2.596285in}{1.392583in}}%
\pgfpathlineto{\pgfqpoint{2.615396in}{1.394406in}}%
\pgfpathlineto{\pgfqpoint{2.633944in}{1.395166in}}%
\pgfpathlineto{\pgfqpoint{2.651960in}{1.394881in}}%
\pgfpathlineto{\pgfqpoint{2.669474in}{1.393101in}}%
\pgfpathlineto{\pgfqpoint{2.694856in}{1.388596in}}%
\pgfpathlineto{\pgfqpoint{2.727164in}{1.382396in}}%
\pgfpathlineto{\pgfqpoint{2.750334in}{1.379256in}}%
\pgfpathlineto{\pgfqpoint{2.772663in}{1.377639in}}%
\pgfpathlineto{\pgfqpoint{2.801222in}{1.376796in}}%
\pgfpathlineto{\pgfqpoint{2.860950in}{1.376499in}}%
\pgfpathlineto{\pgfqpoint{3.043110in}{1.376489in}}%
\pgfpathlineto{\pgfqpoint{3.043110in}{1.376489in}}%
\pgfusepath{stroke}%
\end{pgfscope}%
\begin{pgfscope}%
\pgfpathrectangle{\pgfqpoint{0.650000in}{0.420000in}}{\pgfqpoint{2.388110in}{1.994488in}}%
\pgfusepath{clip}%
\pgfsetbuttcap%
\pgfsetroundjoin%
\pgfsetlinewidth{0.602250pt}%
\definecolor{currentstroke}{rgb}{0.000000,0.000000,0.000000}%
\pgfsetstrokecolor{currentstroke}%
\pgfsetdash{{1.800000pt}{3.000000pt}{0.600000pt}{3.000000pt}}{0.000000pt}%
\pgfpathmoveto{\pgfqpoint{0.645000in}{2.134600in}}%
\pgfpathlineto{\pgfqpoint{0.742708in}{2.112502in}}%
\pgfpathlineto{\pgfqpoint{0.889821in}{2.065153in}}%
\pgfpathlineto{\pgfqpoint{1.010016in}{2.008407in}}%
\pgfpathlineto{\pgfqpoint{1.111637in}{1.930186in}}%
\pgfpathlineto{\pgfqpoint{1.199660in}{1.816720in}}%
\pgfpathlineto{\pgfqpoint{1.277299in}{1.663668in}}%
\pgfpathlineto{\pgfqpoint{1.346745in}{1.485672in}}%
\pgfpathlineto{\pgfqpoint{1.409562in}{1.314008in}}%
\pgfpathlineto{\pgfqpoint{1.466906in}{1.181797in}}%
\pgfpathlineto{\pgfqpoint{1.519654in}{1.107828in}}%
\pgfpathlineto{\pgfqpoint{1.568486in}{1.090637in}}%
\pgfpathlineto{\pgfqpoint{1.613944in}{1.114332in}}%
\pgfpathlineto{\pgfqpoint{1.656463in}{1.159118in}}%
\pgfpathlineto{\pgfqpoint{1.696400in}{1.209054in}}%
\pgfpathlineto{\pgfqpoint{1.734049in}{1.254682in}}%
\pgfpathlineto{\pgfqpoint{1.769658in}{1.291990in}}%
\pgfpathlineto{\pgfqpoint{1.803436in}{1.320321in}}%
\pgfpathlineto{\pgfqpoint{1.835562in}{1.340753in}}%
\pgfpathlineto{\pgfqpoint{1.866188in}{1.354941in}}%
\pgfpathlineto{\pgfqpoint{1.895450in}{1.364394in}}%
\pgfpathlineto{\pgfqpoint{1.923461in}{1.370473in}}%
\pgfpathlineto{\pgfqpoint{1.950325in}{1.374245in}}%
\pgfpathlineto{\pgfqpoint{1.976131in}{1.376423in}}%
\pgfpathlineto{\pgfqpoint{2.024880in}{1.378059in}}%
\pgfpathlineto{\pgfqpoint{2.070247in}{1.378147in}}%
\pgfpathlineto{\pgfqpoint{2.171545in}{1.376952in}}%
\pgfpathlineto{\pgfqpoint{2.259196in}{1.376527in}}%
\pgfpathlineto{\pgfqpoint{2.378674in}{1.376297in}}%
\pgfpathlineto{\pgfqpoint{2.576578in}{1.376643in}}%
\pgfpathlineto{\pgfqpoint{2.772663in}{1.376675in}}%
\pgfpathlineto{\pgfqpoint{2.915251in}{1.376491in}}%
\pgfpathlineto{\pgfqpoint{3.043110in}{1.376489in}}%
\pgfpathlineto{\pgfqpoint{3.043110in}{1.376489in}}%
\pgfusepath{stroke}%
\end{pgfscope}%
\begin{pgfscope}%
\pgfpathrectangle{\pgfqpoint{0.650000in}{0.420000in}}{\pgfqpoint{2.388110in}{1.994488in}}%
\pgfusepath{clip}%
\pgfsetbuttcap%
\pgfsetroundjoin%
\pgfsetlinewidth{0.602250pt}%
\definecolor{currentstroke}{rgb}{0.000000,0.000000,0.000000}%
\pgfsetstrokecolor{currentstroke}%
\pgfsetdash{{0.600000pt}{3.000000pt}}{0.000000pt}%
\pgfpathmoveto{\pgfqpoint{0.645000in}{1.376506in}}%
\pgfpathlineto{\pgfqpoint{1.976131in}{1.376545in}}%
\pgfpathlineto{\pgfqpoint{3.043110in}{1.376488in}}%
\pgfpathlineto{\pgfqpoint{3.043110in}{1.376488in}}%
\pgfusepath{stroke}%
\end{pgfscope}%
\begin{pgfscope}%
\pgfpathrectangle{\pgfqpoint{0.650000in}{0.420000in}}{\pgfqpoint{2.388110in}{1.994488in}}%
\pgfusepath{clip}%
\pgfsetrectcap%
\pgfsetroundjoin%
\pgfsetlinewidth{1.003750pt}%
\definecolor{currentstroke}{rgb}{0.870588,0.466667,0.682353}%
\pgfsetstrokecolor{currentstroke}%
\pgfsetdash{}{0pt}%
\pgfpathmoveto{\pgfqpoint{0.645000in}{1.377008in}}%
\pgfpathlineto{\pgfqpoint{0.749359in}{1.377808in}}%
\pgfpathlineto{\pgfqpoint{0.832320in}{1.379289in}}%
\pgfpathlineto{\pgfqpoint{0.882256in}{1.380904in}}%
\pgfpathlineto{\pgfqpoint{0.928912in}{1.383252in}}%
\pgfpathlineto{\pgfqpoint{0.972898in}{1.386599in}}%
\pgfpathlineto{\pgfqpoint{1.014673in}{1.391290in}}%
\pgfpathlineto{\pgfqpoint{1.034846in}{1.394273in}}%
\pgfpathlineto{\pgfqpoint{1.054593in}{1.397765in}}%
\pgfpathlineto{\pgfqpoint{1.073948in}{1.401837in}}%
\pgfpathlineto{\pgfqpoint{1.092940in}{1.406571in}}%
\pgfpathlineto{\pgfqpoint{1.111595in}{1.412053in}}%
\pgfpathlineto{\pgfqpoint{1.129938in}{1.418380in}}%
\pgfpathlineto{\pgfqpoint{1.147990in}{1.425654in}}%
\pgfpathlineto{\pgfqpoint{1.165771in}{1.433988in}}%
\pgfpathlineto{\pgfqpoint{1.183300in}{1.443498in}}%
\pgfpathlineto{\pgfqpoint{1.200591in}{1.454309in}}%
\pgfpathlineto{\pgfqpoint{1.217661in}{1.466549in}}%
\pgfpathlineto{\pgfqpoint{1.234524in}{1.480346in}}%
\pgfpathlineto{\pgfqpoint{1.251191in}{1.495831in}}%
\pgfpathlineto{\pgfqpoint{1.267675in}{1.513129in}}%
\pgfpathlineto{\pgfqpoint{1.283987in}{1.532357in}}%
\pgfpathlineto{\pgfqpoint{1.300136in}{1.553620in}}%
\pgfpathlineto{\pgfqpoint{1.316132in}{1.577007in}}%
\pgfpathlineto{\pgfqpoint{1.331983in}{1.602580in}}%
\pgfpathlineto{\pgfqpoint{1.347698in}{1.630375in}}%
\pgfpathlineto{\pgfqpoint{1.363283in}{1.660390in}}%
\pgfpathlineto{\pgfqpoint{1.378746in}{1.692582in}}%
\pgfpathlineto{\pgfqpoint{1.394094in}{1.726858in}}%
\pgfpathlineto{\pgfqpoint{1.409332in}{1.763072in}}%
\pgfpathlineto{\pgfqpoint{1.424466in}{1.801024in}}%
\pgfpathlineto{\pgfqpoint{1.454444in}{1.881036in}}%
\pgfpathlineto{\pgfqpoint{1.527904in}{2.086280in}}%
\pgfpathlineto{\pgfqpoint{1.542373in}{2.124225in}}%
\pgfpathlineto{\pgfqpoint{1.556777in}{2.159857in}}%
\pgfpathlineto{\pgfqpoint{1.571118in}{2.192646in}}%
\pgfpathlineto{\pgfqpoint{1.585398in}{2.222099in}}%
\pgfpathlineto{\pgfqpoint{1.599622in}{2.247777in}}%
\pgfpathlineto{\pgfqpoint{1.613791in}{2.269308in}}%
\pgfpathlineto{\pgfqpoint{1.627908in}{2.286405in}}%
\pgfpathlineto{\pgfqpoint{1.641975in}{2.298868in}}%
\pgfpathlineto{\pgfqpoint{1.655995in}{2.306593in}}%
\pgfpathlineto{\pgfqpoint{1.669970in}{2.309567in}}%
\pgfpathlineto{\pgfqpoint{1.683901in}{2.307871in}}%
\pgfpathlineto{\pgfqpoint{1.697791in}{2.301668in}}%
\pgfpathlineto{\pgfqpoint{1.711641in}{2.291200in}}%
\pgfpathlineto{\pgfqpoint{1.725453in}{2.276773in}}%
\pgfpathlineto{\pgfqpoint{1.739229in}{2.258743in}}%
\pgfpathlineto{\pgfqpoint{1.752971in}{2.237507in}}%
\pgfpathlineto{\pgfqpoint{1.766679in}{2.213487in}}%
\pgfpathlineto{\pgfqpoint{1.780356in}{2.187116in}}%
\pgfpathlineto{\pgfqpoint{1.794002in}{2.158824in}}%
\pgfpathlineto{\pgfqpoint{1.821209in}{2.098147in}}%
\pgfpathlineto{\pgfqpoint{1.915650in}{1.879234in}}%
\pgfpathlineto{\pgfqpoint{1.942442in}{1.823067in}}%
\pgfpathlineto{\pgfqpoint{1.955810in}{1.796803in}}%
\pgfpathlineto{\pgfqpoint{1.969160in}{1.771822in}}%
\pgfpathlineto{\pgfqpoint{1.982494in}{1.748141in}}%
\pgfpathlineto{\pgfqpoint{1.995812in}{1.725766in}}%
\pgfpathlineto{\pgfqpoint{2.009115in}{1.704698in}}%
\pgfpathlineto{\pgfqpoint{2.022403in}{1.684935in}}%
\pgfpathlineto{\pgfqpoint{2.035676in}{1.666468in}}%
\pgfpathlineto{\pgfqpoint{2.048936in}{1.649278in}}%
\pgfpathlineto{\pgfqpoint{2.062183in}{1.633339in}}%
\pgfpathlineto{\pgfqpoint{2.075418in}{1.618624in}}%
\pgfpathlineto{\pgfqpoint{2.088640in}{1.605104in}}%
\pgfpathlineto{\pgfqpoint{2.101851in}{1.592740in}}%
\pgfpathlineto{\pgfqpoint{2.115051in}{1.581477in}}%
\pgfpathlineto{\pgfqpoint{2.128241in}{1.571246in}}%
\pgfpathlineto{\pgfqpoint{2.141420in}{1.561958in}}%
\pgfpathlineto{\pgfqpoint{2.154589in}{1.553511in}}%
\pgfpathlineto{\pgfqpoint{2.180900in}{1.538649in}}%
\pgfpathlineto{\pgfqpoint{2.207177in}{1.525739in}}%
\pgfpathlineto{\pgfqpoint{2.233421in}{1.514267in}}%
\pgfpathlineto{\pgfqpoint{2.259636in}{1.504055in}}%
\pgfpathlineto{\pgfqpoint{2.285823in}{1.494909in}}%
\pgfpathlineto{\pgfqpoint{2.311986in}{1.486684in}}%
\pgfpathlineto{\pgfqpoint{2.338125in}{1.479390in}}%
\pgfpathlineto{\pgfqpoint{2.364243in}{1.472953in}}%
\pgfpathlineto{\pgfqpoint{2.403383in}{1.464671in}}%
\pgfpathlineto{\pgfqpoint{2.455510in}{1.454266in}}%
\pgfpathlineto{\pgfqpoint{2.494565in}{1.444555in}}%
\pgfpathlineto{\pgfqpoint{2.546569in}{1.430274in}}%
\pgfpathlineto{\pgfqpoint{2.596973in}{1.414797in}}%
\pgfpathlineto{\pgfqpoint{2.662129in}{1.393976in}}%
\pgfpathlineto{\pgfqpoint{2.690445in}{1.386373in}}%
\pgfpathlineto{\pgfqpoint{2.708055in}{1.382643in}}%
\pgfpathlineto{\pgfqpoint{2.724772in}{1.379983in}}%
\pgfpathlineto{\pgfqpoint{2.748358in}{1.377714in}}%
\pgfpathlineto{\pgfqpoint{2.770369in}{1.376820in}}%
\pgfpathlineto{\pgfqpoint{2.816644in}{1.376494in}}%
\pgfpathlineto{\pgfqpoint{3.043110in}{1.376489in}}%
\pgfpathlineto{\pgfqpoint{3.043110in}{1.376489in}}%
\pgfusepath{stroke}%
\end{pgfscope}%
\begin{pgfscope}%
\pgfpathrectangle{\pgfqpoint{0.650000in}{0.420000in}}{\pgfqpoint{2.388110in}{1.994488in}}%
\pgfusepath{clip}%
\pgfsetbuttcap%
\pgfsetroundjoin%
\pgfsetlinewidth{1.003750pt}%
\definecolor{currentstroke}{rgb}{0.870588,0.466667,0.682353}%
\pgfsetstrokecolor{currentstroke}%
\pgfsetdash{{5.000000pt}{1.000000pt}}{0.000000pt}%
\pgfpathmoveto{\pgfqpoint{0.645000in}{1.376487in}}%
\pgfpathlineto{\pgfqpoint{1.378746in}{1.375959in}}%
\pgfpathlineto{\pgfqpoint{1.711641in}{1.375144in}}%
\pgfpathlineto{\pgfqpoint{1.834772in}{1.376136in}}%
\pgfpathlineto{\pgfqpoint{1.982494in}{1.378459in}}%
\pgfpathlineto{\pgfqpoint{2.022403in}{1.378932in}}%
\pgfpathlineto{\pgfqpoint{2.048936in}{1.378439in}}%
\pgfpathlineto{\pgfqpoint{2.115051in}{1.375356in}}%
\pgfpathlineto{\pgfqpoint{2.167749in}{1.372064in}}%
\pgfpathlineto{\pgfqpoint{2.233421in}{1.366111in}}%
\pgfpathlineto{\pgfqpoint{2.272733in}{1.361906in}}%
\pgfpathlineto{\pgfqpoint{2.325058in}{1.355746in}}%
\pgfpathlineto{\pgfqpoint{2.351187in}{1.353811in}}%
\pgfpathlineto{\pgfqpoint{2.403383in}{1.351310in}}%
\pgfpathlineto{\pgfqpoint{2.429455in}{1.351255in}}%
\pgfpathlineto{\pgfqpoint{2.520585in}{1.353993in}}%
\pgfpathlineto{\pgfqpoint{2.572270in}{1.355946in}}%
\pgfpathlineto{\pgfqpoint{2.681274in}{1.362024in}}%
\pgfpathlineto{\pgfqpoint{2.724772in}{1.366490in}}%
\pgfpathlineto{\pgfqpoint{2.797601in}{1.374227in}}%
\pgfpathlineto{\pgfqpoint{2.822755in}{1.375414in}}%
\pgfpathlineto{\pgfqpoint{2.867947in}{1.376287in}}%
\pgfpathlineto{\pgfqpoint{2.978945in}{1.376443in}}%
\pgfpathlineto{\pgfqpoint{3.043110in}{1.376492in}}%
\pgfpathlineto{\pgfqpoint{3.043110in}{1.376492in}}%
\pgfusepath{stroke}%
\end{pgfscope}%
\begin{pgfscope}%
\pgfpathrectangle{\pgfqpoint{0.650000in}{0.420000in}}{\pgfqpoint{2.388110in}{1.994488in}}%
\pgfusepath{clip}%
\pgfsetbuttcap%
\pgfsetroundjoin%
\pgfsetlinewidth{1.003750pt}%
\definecolor{currentstroke}{rgb}{0.870588,0.466667,0.682353}%
\pgfsetstrokecolor{currentstroke}%
\pgfsetdash{{3.000000pt}{1.000000pt}{1.000000pt}{1.000000pt}}{0.000000pt}%
\pgfpathmoveto{\pgfqpoint{0.645000in}{0.583214in}}%
\pgfpathlineto{\pgfqpoint{0.686817in}{0.589687in}}%
\pgfpathlineto{\pgfqpoint{0.749359in}{0.600316in}}%
\pgfpathlineto{\pgfqpoint{0.805873in}{0.610821in}}%
\pgfpathlineto{\pgfqpoint{0.882256in}{0.626342in}}%
\pgfpathlineto{\pgfqpoint{0.951207in}{0.641688in}}%
\pgfpathlineto{\pgfqpoint{1.014673in}{0.657140in}}%
\pgfpathlineto{\pgfqpoint{1.054593in}{0.667725in}}%
\pgfpathlineto{\pgfqpoint{1.092940in}{0.678751in}}%
\pgfpathlineto{\pgfqpoint{1.129938in}{0.690433in}}%
\pgfpathlineto{\pgfqpoint{1.165771in}{0.702996in}}%
\pgfpathlineto{\pgfqpoint{1.200591in}{0.716628in}}%
\pgfpathlineto{\pgfqpoint{1.234524in}{0.731418in}}%
\pgfpathlineto{\pgfqpoint{1.267675in}{0.747256in}}%
\pgfpathlineto{\pgfqpoint{1.316132in}{0.772243in}}%
\pgfpathlineto{\pgfqpoint{1.378746in}{0.805395in}}%
\pgfpathlineto{\pgfqpoint{1.409332in}{0.820154in}}%
\pgfpathlineto{\pgfqpoint{1.424466in}{0.826662in}}%
\pgfpathlineto{\pgfqpoint{1.439502in}{0.832459in}}%
\pgfpathlineto{\pgfqpoint{1.454444in}{0.837476in}}%
\pgfpathlineto{\pgfqpoint{1.469297in}{0.841669in}}%
\pgfpathlineto{\pgfqpoint{1.484065in}{0.845019in}}%
\pgfpathlineto{\pgfqpoint{1.498753in}{0.847516in}}%
\pgfpathlineto{\pgfqpoint{1.513365in}{0.849171in}}%
\pgfpathlineto{\pgfqpoint{1.527904in}{0.850021in}}%
\pgfpathlineto{\pgfqpoint{1.542373in}{0.850158in}}%
\pgfpathlineto{\pgfqpoint{1.571118in}{0.848932in}}%
\pgfpathlineto{\pgfqpoint{1.627908in}{0.845190in}}%
\pgfpathlineto{\pgfqpoint{1.655995in}{0.844927in}}%
\pgfpathlineto{\pgfqpoint{1.669970in}{0.845661in}}%
\pgfpathlineto{\pgfqpoint{1.683901in}{0.847092in}}%
\pgfpathlineto{\pgfqpoint{1.697791in}{0.849281in}}%
\pgfpathlineto{\pgfqpoint{1.711641in}{0.852259in}}%
\pgfpathlineto{\pgfqpoint{1.725453in}{0.856040in}}%
\pgfpathlineto{\pgfqpoint{1.739229in}{0.860617in}}%
\pgfpathlineto{\pgfqpoint{1.752971in}{0.865974in}}%
\pgfpathlineto{\pgfqpoint{1.766679in}{0.872085in}}%
\pgfpathlineto{\pgfqpoint{1.780356in}{0.878931in}}%
\pgfpathlineto{\pgfqpoint{1.794002in}{0.886487in}}%
\pgfpathlineto{\pgfqpoint{1.821209in}{0.903495in}}%
\pgfpathlineto{\pgfqpoint{1.848310in}{0.922429in}}%
\pgfpathlineto{\pgfqpoint{1.875312in}{0.942922in}}%
\pgfpathlineto{\pgfqpoint{1.902225in}{0.964882in}}%
\pgfpathlineto{\pgfqpoint{1.942442in}{0.999524in}}%
\pgfpathlineto{\pgfqpoint{1.995812in}{1.047133in}}%
\pgfpathlineto{\pgfqpoint{2.035676in}{1.082451in}}%
\pgfpathlineto{\pgfqpoint{2.062183in}{1.104963in}}%
\pgfpathlineto{\pgfqpoint{2.088640in}{1.126390in}}%
\pgfpathlineto{\pgfqpoint{2.115051in}{1.146664in}}%
\pgfpathlineto{\pgfqpoint{2.141420in}{1.165591in}}%
\pgfpathlineto{\pgfqpoint{2.167749in}{1.182914in}}%
\pgfpathlineto{\pgfqpoint{2.194043in}{1.198747in}}%
\pgfpathlineto{\pgfqpoint{2.220302in}{1.213317in}}%
\pgfpathlineto{\pgfqpoint{2.246532in}{1.226800in}}%
\pgfpathlineto{\pgfqpoint{2.272733in}{1.239318in}}%
\pgfpathlineto{\pgfqpoint{2.298907in}{1.250923in}}%
\pgfpathlineto{\pgfqpoint{2.325058in}{1.261553in}}%
\pgfpathlineto{\pgfqpoint{2.364243in}{1.275819in}}%
\pgfpathlineto{\pgfqpoint{2.403383in}{1.288579in}}%
\pgfpathlineto{\pgfqpoint{2.455510in}{1.303898in}}%
\pgfpathlineto{\pgfqpoint{2.559491in}{1.333172in}}%
\pgfpathlineto{\pgfqpoint{2.608756in}{1.346151in}}%
\pgfpathlineto{\pgfqpoint{2.652123in}{1.356437in}}%
\pgfpathlineto{\pgfqpoint{2.690445in}{1.364299in}}%
\pgfpathlineto{\pgfqpoint{2.716519in}{1.368670in}}%
\pgfpathlineto{\pgfqpoint{2.740682in}{1.371804in}}%
\pgfpathlineto{\pgfqpoint{2.770369in}{1.374414in}}%
\pgfpathlineto{\pgfqpoint{2.804072in}{1.375924in}}%
\pgfpathlineto{\pgfqpoint{2.846125in}{1.376423in}}%
\pgfpathlineto{\pgfqpoint{3.037507in}{1.376489in}}%
\pgfpathlineto{\pgfqpoint{3.043110in}{1.376489in}}%
\pgfpathlineto{\pgfqpoint{3.043110in}{1.376489in}}%
\pgfusepath{stroke}%
\end{pgfscope}%
\begin{pgfscope}%
\pgfpathrectangle{\pgfqpoint{0.650000in}{0.420000in}}{\pgfqpoint{2.388110in}{1.994488in}}%
\pgfusepath{clip}%
\pgfsetbuttcap%
\pgfsetroundjoin%
\pgfsetlinewidth{1.003750pt}%
\definecolor{currentstroke}{rgb}{0.870588,0.466667,0.682353}%
\pgfsetstrokecolor{currentstroke}%
\pgfsetdash{{1.000000pt}{1.000000pt}}{0.000000pt}%
\pgfpathmoveto{\pgfqpoint{0.645000in}{1.407607in}}%
\pgfpathlineto{\pgfqpoint{0.718942in}{1.413493in}}%
\pgfpathlineto{\pgfqpoint{0.857743in}{1.425931in}}%
\pgfpathlineto{\pgfqpoint{0.994037in}{1.438057in}}%
\pgfpathlineto{\pgfqpoint{1.054593in}{1.442491in}}%
\pgfpathlineto{\pgfqpoint{1.111595in}{1.445628in}}%
\pgfpathlineto{\pgfqpoint{1.165771in}{1.447391in}}%
\pgfpathlineto{\pgfqpoint{1.217661in}{1.447769in}}%
\pgfpathlineto{\pgfqpoint{1.267675in}{1.446758in}}%
\pgfpathlineto{\pgfqpoint{1.316132in}{1.444471in}}%
\pgfpathlineto{\pgfqpoint{1.363283in}{1.441042in}}%
\pgfpathlineto{\pgfqpoint{1.409332in}{1.436503in}}%
\pgfpathlineto{\pgfqpoint{1.454444in}{1.430860in}}%
\pgfpathlineto{\pgfqpoint{1.498753in}{1.424287in}}%
\pgfpathlineto{\pgfqpoint{1.542373in}{1.416835in}}%
\pgfpathlineto{\pgfqpoint{1.599622in}{1.405720in}}%
\pgfpathlineto{\pgfqpoint{1.739229in}{1.376992in}}%
\pgfpathlineto{\pgfqpoint{1.766679in}{1.372522in}}%
\pgfpathlineto{\pgfqpoint{1.794002in}{1.368925in}}%
\pgfpathlineto{\pgfqpoint{1.821209in}{1.366343in}}%
\pgfpathlineto{\pgfqpoint{1.848310in}{1.364835in}}%
\pgfpathlineto{\pgfqpoint{1.875312in}{1.364376in}}%
\pgfpathlineto{\pgfqpoint{1.902225in}{1.364949in}}%
\pgfpathlineto{\pgfqpoint{1.929055in}{1.366516in}}%
\pgfpathlineto{\pgfqpoint{1.969160in}{1.370106in}}%
\pgfpathlineto{\pgfqpoint{2.048936in}{1.377771in}}%
\pgfpathlineto{\pgfqpoint{2.088640in}{1.380152in}}%
\pgfpathlineto{\pgfqpoint{2.154589in}{1.382301in}}%
\pgfpathlineto{\pgfqpoint{2.194043in}{1.382644in}}%
\pgfpathlineto{\pgfqpoint{2.220302in}{1.381855in}}%
\pgfpathlineto{\pgfqpoint{2.325058in}{1.376640in}}%
\pgfpathlineto{\pgfqpoint{2.364243in}{1.376108in}}%
\pgfpathlineto{\pgfqpoint{2.429455in}{1.376307in}}%
\pgfpathlineto{\pgfqpoint{2.481550in}{1.374513in}}%
\pgfpathlineto{\pgfqpoint{2.584795in}{1.373885in}}%
\pgfpathlineto{\pgfqpoint{2.662129in}{1.371970in}}%
\pgfpathlineto{\pgfqpoint{2.690445in}{1.372391in}}%
\pgfpathlineto{\pgfqpoint{2.816644in}{1.376603in}}%
\pgfpathlineto{\pgfqpoint{2.916912in}{1.376489in}}%
\pgfpathlineto{\pgfqpoint{3.043110in}{1.376432in}}%
\pgfpathlineto{\pgfqpoint{3.043110in}{1.376432in}}%
\pgfusepath{stroke}%
\end{pgfscope}%
\begin{pgfscope}%
\pgfpathrectangle{\pgfqpoint{0.650000in}{0.420000in}}{\pgfqpoint{2.388110in}{1.994488in}}%
\pgfusepath{clip}%
\pgfsetbuttcap%
\pgfsetroundjoin%
\pgfsetlinewidth{1.003750pt}%
\definecolor{currentstroke}{rgb}{0.870588,0.466667,0.682353}%
\pgfsetstrokecolor{currentstroke}%
\pgfsetdash{{5.000000pt}{5.000000pt}}{0.000000pt}%
\pgfpathmoveto{\pgfqpoint{0.645000in}{1.377349in}}%
\pgfpathlineto{\pgfqpoint{0.749359in}{1.378607in}}%
\pgfpathlineto{\pgfqpoint{0.832320in}{1.380850in}}%
\pgfpathlineto{\pgfqpoint{0.882256in}{1.383236in}}%
\pgfpathlineto{\pgfqpoint{0.928912in}{1.386625in}}%
\pgfpathlineto{\pgfqpoint{0.972898in}{1.391327in}}%
\pgfpathlineto{\pgfqpoint{1.014673in}{1.397694in}}%
\pgfpathlineto{\pgfqpoint{1.034846in}{1.401620in}}%
\pgfpathlineto{\pgfqpoint{1.054593in}{1.406105in}}%
\pgfpathlineto{\pgfqpoint{1.073948in}{1.411194in}}%
\pgfpathlineto{\pgfqpoint{1.092940in}{1.416926in}}%
\pgfpathlineto{\pgfqpoint{1.111595in}{1.423335in}}%
\pgfpathlineto{\pgfqpoint{1.129938in}{1.430443in}}%
\pgfpathlineto{\pgfqpoint{1.147990in}{1.438261in}}%
\pgfpathlineto{\pgfqpoint{1.165771in}{1.446780in}}%
\pgfpathlineto{\pgfqpoint{1.183300in}{1.455969in}}%
\pgfpathlineto{\pgfqpoint{1.200591in}{1.465772in}}%
\pgfpathlineto{\pgfqpoint{1.234524in}{1.486841in}}%
\pgfpathlineto{\pgfqpoint{1.300136in}{1.530090in}}%
\pgfpathlineto{\pgfqpoint{1.316132in}{1.539702in}}%
\pgfpathlineto{\pgfqpoint{1.331983in}{1.548210in}}%
\pgfpathlineto{\pgfqpoint{1.347698in}{1.555252in}}%
\pgfpathlineto{\pgfqpoint{1.363283in}{1.560452in}}%
\pgfpathlineto{\pgfqpoint{1.378746in}{1.563438in}}%
\pgfpathlineto{\pgfqpoint{1.394094in}{1.563864in}}%
\pgfpathlineto{\pgfqpoint{1.409332in}{1.561427in}}%
\pgfpathlineto{\pgfqpoint{1.424466in}{1.555880in}}%
\pgfpathlineto{\pgfqpoint{1.439502in}{1.547041in}}%
\pgfpathlineto{\pgfqpoint{1.454444in}{1.534805in}}%
\pgfpathlineto{\pgfqpoint{1.469297in}{1.519165in}}%
\pgfpathlineto{\pgfqpoint{1.484065in}{1.500219in}}%
\pgfpathlineto{\pgfqpoint{1.498753in}{1.478179in}}%
\pgfpathlineto{\pgfqpoint{1.513365in}{1.453376in}}%
\pgfpathlineto{\pgfqpoint{1.527904in}{1.426238in}}%
\pgfpathlineto{\pgfqpoint{1.542373in}{1.397289in}}%
\pgfpathlineto{\pgfqpoint{1.571118in}{1.336391in}}%
\pgfpathlineto{\pgfqpoint{1.599622in}{1.275883in}}%
\pgfpathlineto{\pgfqpoint{1.613791in}{1.247432in}}%
\pgfpathlineto{\pgfqpoint{1.627908in}{1.221008in}}%
\pgfpathlineto{\pgfqpoint{1.641975in}{1.197138in}}%
\pgfpathlineto{\pgfqpoint{1.655995in}{1.176267in}}%
\pgfpathlineto{\pgfqpoint{1.669970in}{1.158750in}}%
\pgfpathlineto{\pgfqpoint{1.683901in}{1.144829in}}%
\pgfpathlineto{\pgfqpoint{1.697791in}{1.134629in}}%
\pgfpathlineto{\pgfqpoint{1.711641in}{1.128154in}}%
\pgfpathlineto{\pgfqpoint{1.725453in}{1.125299in}}%
\pgfpathlineto{\pgfqpoint{1.739229in}{1.125867in}}%
\pgfpathlineto{\pgfqpoint{1.752971in}{1.129580in}}%
\pgfpathlineto{\pgfqpoint{1.766679in}{1.136108in}}%
\pgfpathlineto{\pgfqpoint{1.780356in}{1.145093in}}%
\pgfpathlineto{\pgfqpoint{1.794002in}{1.156178in}}%
\pgfpathlineto{\pgfqpoint{1.807620in}{1.168981in}}%
\pgfpathlineto{\pgfqpoint{1.821209in}{1.183091in}}%
\pgfpathlineto{\pgfqpoint{1.848310in}{1.213610in}}%
\pgfpathlineto{\pgfqpoint{1.888779in}{1.260329in}}%
\pgfpathlineto{\pgfqpoint{1.915650in}{1.289182in}}%
\pgfpathlineto{\pgfqpoint{1.929055in}{1.302322in}}%
\pgfpathlineto{\pgfqpoint{1.942442in}{1.314452in}}%
\pgfpathlineto{\pgfqpoint{1.955810in}{1.325530in}}%
\pgfpathlineto{\pgfqpoint{1.969160in}{1.335529in}}%
\pgfpathlineto{\pgfqpoint{1.982494in}{1.344461in}}%
\pgfpathlineto{\pgfqpoint{1.995812in}{1.352380in}}%
\pgfpathlineto{\pgfqpoint{2.009115in}{1.359345in}}%
\pgfpathlineto{\pgfqpoint{2.022403in}{1.365387in}}%
\pgfpathlineto{\pgfqpoint{2.035676in}{1.370543in}}%
\pgfpathlineto{\pgfqpoint{2.048936in}{1.374896in}}%
\pgfpathlineto{\pgfqpoint{2.062183in}{1.378559in}}%
\pgfpathlineto{\pgfqpoint{2.075418in}{1.381608in}}%
\pgfpathlineto{\pgfqpoint{2.088640in}{1.384059in}}%
\pgfpathlineto{\pgfqpoint{2.101851in}{1.385905in}}%
\pgfpathlineto{\pgfqpoint{2.115051in}{1.387116in}}%
\pgfpathlineto{\pgfqpoint{2.128241in}{1.387660in}}%
\pgfpathlineto{\pgfqpoint{2.154589in}{1.386955in}}%
\pgfpathlineto{\pgfqpoint{2.233421in}{1.381546in}}%
\pgfpathlineto{\pgfqpoint{2.338125in}{1.380424in}}%
\pgfpathlineto{\pgfqpoint{2.377295in}{1.378690in}}%
\pgfpathlineto{\pgfqpoint{2.429455in}{1.375604in}}%
\pgfpathlineto{\pgfqpoint{2.455510in}{1.375234in}}%
\pgfpathlineto{\pgfqpoint{2.507577in}{1.376402in}}%
\pgfpathlineto{\pgfqpoint{2.533587in}{1.377435in}}%
\pgfpathlineto{\pgfqpoint{2.584795in}{1.381635in}}%
\pgfpathlineto{\pgfqpoint{2.608756in}{1.384190in}}%
\pgfpathlineto{\pgfqpoint{2.662129in}{1.392215in}}%
\pgfpathlineto{\pgfqpoint{2.681274in}{1.393093in}}%
\pgfpathlineto{\pgfqpoint{2.716519in}{1.392674in}}%
\pgfpathlineto{\pgfqpoint{2.732823in}{1.389789in}}%
\pgfpathlineto{\pgfqpoint{2.755859in}{1.385196in}}%
\pgfpathlineto{\pgfqpoint{2.777391in}{1.381804in}}%
\pgfpathlineto{\pgfqpoint{2.804072in}{1.378784in}}%
\pgfpathlineto{\pgfqpoint{2.828755in}{1.376683in}}%
\pgfpathlineto{\pgfqpoint{2.857216in}{1.376048in}}%
\pgfpathlineto{\pgfqpoint{2.867947in}{1.376279in}}%
\pgfpathlineto{\pgfqpoint{2.883415in}{1.376160in}}%
\pgfpathlineto{\pgfqpoint{2.898190in}{1.375880in}}%
\pgfpathlineto{\pgfqpoint{2.934629in}{1.376156in}}%
\pgfpathlineto{\pgfqpoint{2.943142in}{1.377115in}}%
\pgfpathlineto{\pgfqpoint{2.963511in}{1.376726in}}%
\pgfpathlineto{\pgfqpoint{2.982694in}{1.376010in}}%
\pgfpathlineto{\pgfqpoint{2.990065in}{1.376705in}}%
\pgfpathlineto{\pgfqpoint{3.007803in}{1.376054in}}%
\pgfpathlineto{\pgfqpoint{3.018004in}{1.376057in}}%
\pgfpathlineto{\pgfqpoint{3.031134in}{1.376893in}}%
\pgfpathlineto{\pgfqpoint{3.043110in}{1.375470in}}%
\pgfpathlineto{\pgfqpoint{3.043110in}{1.375470in}}%
\pgfusepath{stroke}%
\end{pgfscope}%
\begin{pgfscope}%
\pgfpathrectangle{\pgfqpoint{0.650000in}{0.420000in}}{\pgfqpoint{2.388110in}{1.994488in}}%
\pgfusepath{clip}%
\pgfsetbuttcap%
\pgfsetroundjoin%
\pgfsetlinewidth{1.003750pt}%
\definecolor{currentstroke}{rgb}{0.870588,0.466667,0.682353}%
\pgfsetstrokecolor{currentstroke}%
\pgfsetdash{{3.000000pt}{5.000000pt}{1.000000pt}{5.000000pt}}{0.000000pt}%
\pgfpathmoveto{\pgfqpoint{0.645000in}{2.137144in}}%
\pgfpathlineto{\pgfqpoint{0.652928in}{2.135309in}}%
\pgfpathlineto{\pgfqpoint{0.686817in}{2.126858in}}%
\pgfpathlineto{\pgfqpoint{0.749359in}{2.109593in}}%
\pgfpathlineto{\pgfqpoint{0.805873in}{2.091875in}}%
\pgfpathlineto{\pgfqpoint{0.857743in}{2.073438in}}%
\pgfpathlineto{\pgfqpoint{0.882256in}{2.063836in}}%
\pgfpathlineto{\pgfqpoint{0.905952in}{2.053901in}}%
\pgfpathlineto{\pgfqpoint{0.928912in}{2.043566in}}%
\pgfpathlineto{\pgfqpoint{0.951207in}{2.032753in}}%
\pgfpathlineto{\pgfqpoint{0.972898in}{2.021378in}}%
\pgfpathlineto{\pgfqpoint{0.994037in}{2.009344in}}%
\pgfpathlineto{\pgfqpoint{1.014673in}{1.996546in}}%
\pgfpathlineto{\pgfqpoint{1.034846in}{1.982867in}}%
\pgfpathlineto{\pgfqpoint{1.054593in}{1.968184in}}%
\pgfpathlineto{\pgfqpoint{1.073948in}{1.952364in}}%
\pgfpathlineto{\pgfqpoint{1.092940in}{1.935264in}}%
\pgfpathlineto{\pgfqpoint{1.111595in}{1.916739in}}%
\pgfpathlineto{\pgfqpoint{1.129938in}{1.896643in}}%
\pgfpathlineto{\pgfqpoint{1.147990in}{1.874827in}}%
\pgfpathlineto{\pgfqpoint{1.165771in}{1.851156in}}%
\pgfpathlineto{\pgfqpoint{1.183300in}{1.825510in}}%
\pgfpathlineto{\pgfqpoint{1.200591in}{1.797798in}}%
\pgfpathlineto{\pgfqpoint{1.217661in}{1.767972in}}%
\pgfpathlineto{\pgfqpoint{1.234524in}{1.736049in}}%
\pgfpathlineto{\pgfqpoint{1.251191in}{1.702123in}}%
\pgfpathlineto{\pgfqpoint{1.267675in}{1.666346in}}%
\pgfpathlineto{\pgfqpoint{1.283987in}{1.628867in}}%
\pgfpathlineto{\pgfqpoint{1.300136in}{1.589792in}}%
\pgfpathlineto{\pgfqpoint{1.331983in}{1.507731in}}%
\pgfpathlineto{\pgfqpoint{1.409332in}{1.301814in}}%
\pgfpathlineto{\pgfqpoint{1.424466in}{1.264719in}}%
\pgfpathlineto{\pgfqpoint{1.439502in}{1.230239in}}%
\pgfpathlineto{\pgfqpoint{1.454444in}{1.198883in}}%
\pgfpathlineto{\pgfqpoint{1.469297in}{1.171110in}}%
\pgfpathlineto{\pgfqpoint{1.484065in}{1.147287in}}%
\pgfpathlineto{\pgfqpoint{1.498753in}{1.127670in}}%
\pgfpathlineto{\pgfqpoint{1.513365in}{1.112401in}}%
\pgfpathlineto{\pgfqpoint{1.527904in}{1.101556in}}%
\pgfpathlineto{\pgfqpoint{1.542373in}{1.095135in}}%
\pgfpathlineto{\pgfqpoint{1.556777in}{1.092944in}}%
\pgfpathlineto{\pgfqpoint{1.571118in}{1.094621in}}%
\pgfpathlineto{\pgfqpoint{1.585398in}{1.099795in}}%
\pgfpathlineto{\pgfqpoint{1.599622in}{1.108097in}}%
\pgfpathlineto{\pgfqpoint{1.613791in}{1.119077in}}%
\pgfpathlineto{\pgfqpoint{1.627908in}{1.132245in}}%
\pgfpathlineto{\pgfqpoint{1.641975in}{1.147129in}}%
\pgfpathlineto{\pgfqpoint{1.669970in}{1.180190in}}%
\pgfpathlineto{\pgfqpoint{1.711641in}{1.231511in}}%
\pgfpathlineto{\pgfqpoint{1.739229in}{1.263402in}}%
\pgfpathlineto{\pgfqpoint{1.752971in}{1.278018in}}%
\pgfpathlineto{\pgfqpoint{1.766679in}{1.291585in}}%
\pgfpathlineto{\pgfqpoint{1.780356in}{1.304070in}}%
\pgfpathlineto{\pgfqpoint{1.794002in}{1.315449in}}%
\pgfpathlineto{\pgfqpoint{1.807620in}{1.325707in}}%
\pgfpathlineto{\pgfqpoint{1.821209in}{1.334899in}}%
\pgfpathlineto{\pgfqpoint{1.834772in}{1.343153in}}%
\pgfpathlineto{\pgfqpoint{1.848310in}{1.350556in}}%
\pgfpathlineto{\pgfqpoint{1.861823in}{1.357111in}}%
\pgfpathlineto{\pgfqpoint{1.875312in}{1.362805in}}%
\pgfpathlineto{\pgfqpoint{1.888779in}{1.367649in}}%
\pgfpathlineto{\pgfqpoint{1.902225in}{1.371697in}}%
\pgfpathlineto{\pgfqpoint{1.915650in}{1.375053in}}%
\pgfpathlineto{\pgfqpoint{1.942442in}{1.380032in}}%
\pgfpathlineto{\pgfqpoint{1.969160in}{1.382949in}}%
\pgfpathlineto{\pgfqpoint{1.995812in}{1.384111in}}%
\pgfpathlineto{\pgfqpoint{2.035676in}{1.384235in}}%
\pgfpathlineto{\pgfqpoint{2.075418in}{1.383547in}}%
\pgfpathlineto{\pgfqpoint{2.115051in}{1.381710in}}%
\pgfpathlineto{\pgfqpoint{2.167749in}{1.379280in}}%
\pgfpathlineto{\pgfqpoint{2.220302in}{1.378158in}}%
\pgfpathlineto{\pgfqpoint{2.468532in}{1.376032in}}%
\pgfpathlineto{\pgfqpoint{2.791001in}{1.376717in}}%
\pgfpathlineto{\pgfqpoint{2.912330in}{1.376494in}}%
\pgfpathlineto{\pgfqpoint{3.043110in}{1.376489in}}%
\pgfpathlineto{\pgfqpoint{3.043110in}{1.376489in}}%
\pgfusepath{stroke}%
\end{pgfscope}%
\begin{pgfscope}%
\pgfpathrectangle{\pgfqpoint{0.650000in}{0.420000in}}{\pgfqpoint{2.388110in}{1.994488in}}%
\pgfusepath{clip}%
\pgfsetbuttcap%
\pgfsetroundjoin%
\pgfsetlinewidth{1.003750pt}%
\definecolor{currentstroke}{rgb}{0.870588,0.466667,0.682353}%
\pgfsetstrokecolor{currentstroke}%
\pgfsetdash{{1.000000pt}{5.000000pt}}{0.000000pt}%
\pgfpathmoveto{\pgfqpoint{0.645000in}{1.376367in}}%
\pgfpathlineto{\pgfqpoint{1.217661in}{1.376148in}}%
\pgfpathlineto{\pgfqpoint{1.300136in}{1.376414in}}%
\pgfpathlineto{\pgfqpoint{1.627908in}{1.377093in}}%
\pgfpathlineto{\pgfqpoint{1.711641in}{1.377937in}}%
\pgfpathlineto{\pgfqpoint{1.834772in}{1.379694in}}%
\pgfpathlineto{\pgfqpoint{1.902225in}{1.380114in}}%
\pgfpathlineto{\pgfqpoint{1.955810in}{1.379561in}}%
\pgfpathlineto{\pgfqpoint{2.009115in}{1.378028in}}%
\pgfpathlineto{\pgfqpoint{2.048936in}{1.375881in}}%
\pgfpathlineto{\pgfqpoint{2.141420in}{1.368841in}}%
\pgfpathlineto{\pgfqpoint{2.180900in}{1.364533in}}%
\pgfpathlineto{\pgfqpoint{2.220302in}{1.359937in}}%
\pgfpathlineto{\pgfqpoint{2.311986in}{1.352543in}}%
\pgfpathlineto{\pgfqpoint{2.338125in}{1.351563in}}%
\pgfpathlineto{\pgfqpoint{2.390341in}{1.351692in}}%
\pgfpathlineto{\pgfqpoint{2.429455in}{1.352964in}}%
\pgfpathlineto{\pgfqpoint{2.468532in}{1.355041in}}%
\pgfpathlineto{\pgfqpoint{2.533587in}{1.360453in}}%
\pgfpathlineto{\pgfqpoint{2.572270in}{1.363947in}}%
\pgfpathlineto{\pgfqpoint{2.608756in}{1.366636in}}%
\pgfpathlineto{\pgfqpoint{2.652123in}{1.371028in}}%
\pgfpathlineto{\pgfqpoint{2.724772in}{1.375963in}}%
\pgfpathlineto{\pgfqpoint{2.748358in}{1.375171in}}%
\pgfpathlineto{\pgfqpoint{2.777391in}{1.376010in}}%
\pgfpathlineto{\pgfqpoint{2.804072in}{1.376527in}}%
\pgfpathlineto{\pgfqpoint{2.857216in}{1.375914in}}%
\pgfpathlineto{\pgfqpoint{2.867947in}{1.376202in}}%
\pgfpathlineto{\pgfqpoint{2.878339in}{1.375976in}}%
\pgfpathlineto{\pgfqpoint{2.893338in}{1.376090in}}%
\pgfpathlineto{\pgfqpoint{2.907684in}{1.375886in}}%
\pgfpathlineto{\pgfqpoint{2.916912in}{1.375834in}}%
\pgfpathlineto{\pgfqpoint{2.925889in}{1.375285in}}%
\pgfpathlineto{\pgfqpoint{2.934629in}{1.376205in}}%
\pgfpathlineto{\pgfqpoint{2.943142in}{1.377222in}}%
\pgfpathlineto{\pgfqpoint{2.963511in}{1.376723in}}%
\pgfpathlineto{\pgfqpoint{2.982694in}{1.375954in}}%
\pgfpathlineto{\pgfqpoint{2.990065in}{1.376702in}}%
\pgfpathlineto{\pgfqpoint{3.021336in}{1.376283in}}%
\pgfpathlineto{\pgfqpoint{3.031134in}{1.376940in}}%
\pgfpathlineto{\pgfqpoint{3.043110in}{1.375417in}}%
\pgfpathlineto{\pgfqpoint{3.043110in}{1.375417in}}%
\pgfusepath{stroke}%
\end{pgfscope}%
\begin{pgfscope}%
\pgfpathrectangle{\pgfqpoint{0.650000in}{0.420000in}}{\pgfqpoint{2.388110in}{1.994488in}}%
\pgfusepath{clip}%
\pgfsetrectcap%
\pgfsetroundjoin%
\pgfsetlinewidth{1.003750pt}%
\definecolor{currentstroke}{rgb}{0.498039,0.737255,0.254902}%
\pgfsetstrokecolor{currentstroke}%
\pgfsetdash{}{0pt}%
\pgfpathmoveto{\pgfqpoint{0.645000in}{1.377012in}}%
\pgfpathlineto{\pgfqpoint{0.750863in}{1.377832in}}%
\pgfpathlineto{\pgfqpoint{0.833824in}{1.379340in}}%
\pgfpathlineto{\pgfqpoint{0.883760in}{1.380984in}}%
\pgfpathlineto{\pgfqpoint{0.930416in}{1.383373in}}%
\pgfpathlineto{\pgfqpoint{0.974402in}{1.386776in}}%
\pgfpathlineto{\pgfqpoint{1.016177in}{1.391544in}}%
\pgfpathlineto{\pgfqpoint{1.036350in}{1.394574in}}%
\pgfpathlineto{\pgfqpoint{1.056097in}{1.398120in}}%
\pgfpathlineto{\pgfqpoint{1.075452in}{1.402255in}}%
\pgfpathlineto{\pgfqpoint{1.094444in}{1.407060in}}%
\pgfpathlineto{\pgfqpoint{1.113099in}{1.412622in}}%
\pgfpathlineto{\pgfqpoint{1.131442in}{1.419039in}}%
\pgfpathlineto{\pgfqpoint{1.149495in}{1.426414in}}%
\pgfpathlineto{\pgfqpoint{1.167276in}{1.434859in}}%
\pgfpathlineto{\pgfqpoint{1.184804in}{1.444492in}}%
\pgfpathlineto{\pgfqpoint{1.202095in}{1.455437in}}%
\pgfpathlineto{\pgfqpoint{1.219166in}{1.467822in}}%
\pgfpathlineto{\pgfqpoint{1.236028in}{1.481775in}}%
\pgfpathlineto{\pgfqpoint{1.252695in}{1.497427in}}%
\pgfpathlineto{\pgfqpoint{1.269179in}{1.514900in}}%
\pgfpathlineto{\pgfqpoint{1.285491in}{1.534311in}}%
\pgfpathlineto{\pgfqpoint{1.301640in}{1.555763in}}%
\pgfpathlineto{\pgfqpoint{1.317636in}{1.579342in}}%
\pgfpathlineto{\pgfqpoint{1.333487in}{1.605107in}}%
\pgfpathlineto{\pgfqpoint{1.349202in}{1.633091in}}%
\pgfpathlineto{\pgfqpoint{1.364787in}{1.663289in}}%
\pgfpathlineto{\pgfqpoint{1.380251in}{1.695652in}}%
\pgfpathlineto{\pgfqpoint{1.395598in}{1.730084in}}%
\pgfpathlineto{\pgfqpoint{1.410836in}{1.766435in}}%
\pgfpathlineto{\pgfqpoint{1.425971in}{1.804499in}}%
\pgfpathlineto{\pgfqpoint{1.455948in}{1.884645in}}%
\pgfpathlineto{\pgfqpoint{1.529408in}{2.089490in}}%
\pgfpathlineto{\pgfqpoint{1.543878in}{2.127197in}}%
\pgfpathlineto{\pgfqpoint{1.558281in}{2.162533in}}%
\pgfpathlineto{\pgfqpoint{1.572622in}{2.194963in}}%
\pgfpathlineto{\pgfqpoint{1.586903in}{2.223994in}}%
\pgfpathlineto{\pgfqpoint{1.601126in}{2.249184in}}%
\pgfpathlineto{\pgfqpoint{1.615295in}{2.270162in}}%
\pgfpathlineto{\pgfqpoint{1.629412in}{2.286642in}}%
\pgfpathlineto{\pgfqpoint{1.643480in}{2.298430in}}%
\pgfpathlineto{\pgfqpoint{1.657499in}{2.305427in}}%
\pgfpathlineto{\pgfqpoint{1.671474in}{2.307632in}}%
\pgfpathlineto{\pgfqpoint{1.685405in}{2.305136in}}%
\pgfpathlineto{\pgfqpoint{1.699295in}{2.298122in}}%
\pgfpathlineto{\pgfqpoint{1.713145in}{2.286849in}}%
\pgfpathlineto{\pgfqpoint{1.726957in}{2.271646in}}%
\pgfpathlineto{\pgfqpoint{1.740734in}{2.252890in}}%
\pgfpathlineto{\pgfqpoint{1.754475in}{2.230998in}}%
\pgfpathlineto{\pgfqpoint{1.768184in}{2.206408in}}%
\pgfpathlineto{\pgfqpoint{1.781860in}{2.179564in}}%
\pgfpathlineto{\pgfqpoint{1.795507in}{2.150903in}}%
\pgfpathlineto{\pgfqpoint{1.822714in}{2.089799in}}%
\pgfpathlineto{\pgfqpoint{1.903729in}{1.901019in}}%
\pgfpathlineto{\pgfqpoint{1.930560in}{1.842926in}}%
\pgfpathlineto{\pgfqpoint{1.957314in}{1.789397in}}%
\pgfpathlineto{\pgfqpoint{1.970665in}{1.764549in}}%
\pgfpathlineto{\pgfqpoint{1.983999in}{1.741023in}}%
\pgfpathlineto{\pgfqpoint{1.997317in}{1.718824in}}%
\pgfpathlineto{\pgfqpoint{2.010619in}{1.697947in}}%
\pgfpathlineto{\pgfqpoint{2.023907in}{1.678383in}}%
\pgfpathlineto{\pgfqpoint{2.037181in}{1.660116in}}%
\pgfpathlineto{\pgfqpoint{2.050441in}{1.643122in}}%
\pgfpathlineto{\pgfqpoint{2.063688in}{1.627369in}}%
\pgfpathlineto{\pgfqpoint{2.076922in}{1.612827in}}%
\pgfpathlineto{\pgfqpoint{2.090145in}{1.599463in}}%
\pgfpathlineto{\pgfqpoint{2.103356in}{1.587237in}}%
\pgfpathlineto{\pgfqpoint{2.116556in}{1.576103in}}%
\pgfpathlineto{\pgfqpoint{2.129745in}{1.566003in}}%
\pgfpathlineto{\pgfqpoint{2.142924in}{1.556870in}}%
\pgfpathlineto{\pgfqpoint{2.156094in}{1.548621in}}%
\pgfpathlineto{\pgfqpoint{2.169254in}{1.541162in}}%
\pgfpathlineto{\pgfqpoint{2.195547in}{1.528211in}}%
\pgfpathlineto{\pgfqpoint{2.221807in}{1.517298in}}%
\pgfpathlineto{\pgfqpoint{2.248036in}{1.507833in}}%
\pgfpathlineto{\pgfqpoint{2.274237in}{1.499360in}}%
\pgfpathlineto{\pgfqpoint{2.313490in}{1.488077in}}%
\pgfpathlineto{\pgfqpoint{2.352691in}{1.478455in}}%
\pgfpathlineto{\pgfqpoint{2.457014in}{1.454633in}}%
\pgfpathlineto{\pgfqpoint{2.496070in}{1.444101in}}%
\pgfpathlineto{\pgfqpoint{2.621647in}{1.407302in}}%
\pgfpathlineto{\pgfqpoint{2.663633in}{1.393831in}}%
\pgfpathlineto{\pgfqpoint{2.682778in}{1.388505in}}%
\pgfpathlineto{\pgfqpoint{2.700872in}{1.384286in}}%
\pgfpathlineto{\pgfqpoint{2.718024in}{1.381156in}}%
\pgfpathlineto{\pgfqpoint{2.734327in}{1.379034in}}%
\pgfpathlineto{\pgfqpoint{2.757364in}{1.377312in}}%
\pgfpathlineto{\pgfqpoint{2.785770in}{1.376596in}}%
\pgfpathlineto{\pgfqpoint{2.874688in}{1.376489in}}%
\pgfpathlineto{\pgfqpoint{3.043110in}{1.376489in}}%
\pgfpathlineto{\pgfqpoint{3.043110in}{1.376489in}}%
\pgfusepath{stroke}%
\end{pgfscope}%
\begin{pgfscope}%
\pgfpathrectangle{\pgfqpoint{0.650000in}{0.420000in}}{\pgfqpoint{2.388110in}{1.994488in}}%
\pgfusepath{clip}%
\pgfsetbuttcap%
\pgfsetroundjoin%
\pgfsetlinewidth{1.003750pt}%
\definecolor{currentstroke}{rgb}{0.498039,0.737255,0.254902}%
\pgfsetstrokecolor{currentstroke}%
\pgfsetdash{{5.000000pt}{1.000000pt}}{0.000000pt}%
\pgfpathmoveto{\pgfqpoint{0.645000in}{1.376489in}}%
\pgfpathlineto{\pgfqpoint{1.500257in}{1.376994in}}%
\pgfpathlineto{\pgfqpoint{1.629412in}{1.378074in}}%
\pgfpathlineto{\pgfqpoint{1.726957in}{1.380103in}}%
\pgfpathlineto{\pgfqpoint{1.822714in}{1.382266in}}%
\pgfpathlineto{\pgfqpoint{2.037181in}{1.382657in}}%
\pgfpathlineto{\pgfqpoint{2.063688in}{1.382126in}}%
\pgfpathlineto{\pgfqpoint{2.090145in}{1.380402in}}%
\pgfpathlineto{\pgfqpoint{2.116556in}{1.377690in}}%
\pgfpathlineto{\pgfqpoint{2.142924in}{1.373882in}}%
\pgfpathlineto{\pgfqpoint{2.221807in}{1.360904in}}%
\pgfpathlineto{\pgfqpoint{2.248036in}{1.358081in}}%
\pgfpathlineto{\pgfqpoint{2.326562in}{1.352229in}}%
\pgfpathlineto{\pgfqpoint{2.365747in}{1.350428in}}%
\pgfpathlineto{\pgfqpoint{2.430959in}{1.348959in}}%
\pgfpathlineto{\pgfqpoint{2.457014in}{1.349417in}}%
\pgfpathlineto{\pgfqpoint{2.548073in}{1.353812in}}%
\pgfpathlineto{\pgfqpoint{2.586300in}{1.355832in}}%
\pgfpathlineto{\pgfqpoint{2.632655in}{1.357446in}}%
\pgfpathlineto{\pgfqpoint{2.663633in}{1.360255in}}%
\pgfpathlineto{\pgfqpoint{2.718024in}{1.365742in}}%
\pgfpathlineto{\pgfqpoint{2.749862in}{1.368927in}}%
\pgfpathlineto{\pgfqpoint{2.785770in}{1.373218in}}%
\pgfpathlineto{\pgfqpoint{2.811922in}{1.375009in}}%
\pgfpathlineto{\pgfqpoint{2.841940in}{1.375968in}}%
\pgfpathlineto{\pgfqpoint{2.927394in}{1.376349in}}%
\pgfpathlineto{\pgfqpoint{3.043110in}{1.376448in}}%
\pgfpathlineto{\pgfqpoint{3.043110in}{1.376448in}}%
\pgfusepath{stroke}%
\end{pgfscope}%
\begin{pgfscope}%
\pgfpathrectangle{\pgfqpoint{0.650000in}{0.420000in}}{\pgfqpoint{2.388110in}{1.994488in}}%
\pgfusepath{clip}%
\pgfsetbuttcap%
\pgfsetroundjoin%
\pgfsetlinewidth{1.003750pt}%
\definecolor{currentstroke}{rgb}{0.498039,0.737255,0.254902}%
\pgfsetstrokecolor{currentstroke}%
\pgfsetdash{{3.000000pt}{1.000000pt}{1.000000pt}{1.000000pt}}{0.000000pt}%
\pgfpathmoveto{\pgfqpoint{0.645000in}{0.579866in}}%
\pgfpathlineto{\pgfqpoint{0.688321in}{0.586546in}}%
\pgfpathlineto{\pgfqpoint{0.750863in}{0.597167in}}%
\pgfpathlineto{\pgfqpoint{0.807378in}{0.607685in}}%
\pgfpathlineto{\pgfqpoint{0.883760in}{0.623276in}}%
\pgfpathlineto{\pgfqpoint{0.952711in}{0.638776in}}%
\pgfpathlineto{\pgfqpoint{1.016177in}{0.654489in}}%
\pgfpathlineto{\pgfqpoint{1.056097in}{0.665320in}}%
\pgfpathlineto{\pgfqpoint{1.094444in}{0.676657in}}%
\pgfpathlineto{\pgfqpoint{1.131442in}{0.688713in}}%
\pgfpathlineto{\pgfqpoint{1.167276in}{0.701705in}}%
\pgfpathlineto{\pgfqpoint{1.202095in}{0.715808in}}%
\pgfpathlineto{\pgfqpoint{1.236028in}{0.731087in}}%
\pgfpathlineto{\pgfqpoint{1.269179in}{0.747404in}}%
\pgfpathlineto{\pgfqpoint{1.317636in}{0.773043in}}%
\pgfpathlineto{\pgfqpoint{1.380251in}{0.806964in}}%
\pgfpathlineto{\pgfqpoint{1.410836in}{0.822120in}}%
\pgfpathlineto{\pgfqpoint{1.425971in}{0.828842in}}%
\pgfpathlineto{\pgfqpoint{1.441006in}{0.834865in}}%
\pgfpathlineto{\pgfqpoint{1.455948in}{0.840122in}}%
\pgfpathlineto{\pgfqpoint{1.470801in}{0.844573in}}%
\pgfpathlineto{\pgfqpoint{1.485569in}{0.848196in}}%
\pgfpathlineto{\pgfqpoint{1.500257in}{0.850981in}}%
\pgfpathlineto{\pgfqpoint{1.514869in}{0.852929in}}%
\pgfpathlineto{\pgfqpoint{1.529408in}{0.854071in}}%
\pgfpathlineto{\pgfqpoint{1.543878in}{0.854492in}}%
\pgfpathlineto{\pgfqpoint{1.572622in}{0.853791in}}%
\pgfpathlineto{\pgfqpoint{1.629412in}{0.850826in}}%
\pgfpathlineto{\pgfqpoint{1.657499in}{0.850776in}}%
\pgfpathlineto{\pgfqpoint{1.671474in}{0.851577in}}%
\pgfpathlineto{\pgfqpoint{1.685405in}{0.853054in}}%
\pgfpathlineto{\pgfqpoint{1.699295in}{0.855268in}}%
\pgfpathlineto{\pgfqpoint{1.713145in}{0.858258in}}%
\pgfpathlineto{\pgfqpoint{1.726957in}{0.862039in}}%
\pgfpathlineto{\pgfqpoint{1.740734in}{0.866611in}}%
\pgfpathlineto{\pgfqpoint{1.754475in}{0.871964in}}%
\pgfpathlineto{\pgfqpoint{1.768184in}{0.878083in}}%
\pgfpathlineto{\pgfqpoint{1.781860in}{0.884955in}}%
\pgfpathlineto{\pgfqpoint{1.795507in}{0.892565in}}%
\pgfpathlineto{\pgfqpoint{1.822714in}{0.909776in}}%
\pgfpathlineto{\pgfqpoint{1.849814in}{0.929013in}}%
\pgfpathlineto{\pgfqpoint{1.876817in}{0.949833in}}%
\pgfpathlineto{\pgfqpoint{1.903729in}{0.972095in}}%
\pgfpathlineto{\pgfqpoint{1.957314in}{1.018779in}}%
\pgfpathlineto{\pgfqpoint{2.023907in}{1.076725in}}%
\pgfpathlineto{\pgfqpoint{2.063688in}{1.109742in}}%
\pgfpathlineto{\pgfqpoint{2.090145in}{1.130450in}}%
\pgfpathlineto{\pgfqpoint{2.116556in}{1.149963in}}%
\pgfpathlineto{\pgfqpoint{2.142924in}{1.168165in}}%
\pgfpathlineto{\pgfqpoint{2.169254in}{1.184985in}}%
\pgfpathlineto{\pgfqpoint{2.195547in}{1.200599in}}%
\pgfpathlineto{\pgfqpoint{2.221807in}{1.215104in}}%
\pgfpathlineto{\pgfqpoint{2.248036in}{1.228460in}}%
\pgfpathlineto{\pgfqpoint{2.274237in}{1.240665in}}%
\pgfpathlineto{\pgfqpoint{2.300412in}{1.251851in}}%
\pgfpathlineto{\pgfqpoint{2.339629in}{1.266988in}}%
\pgfpathlineto{\pgfqpoint{2.378799in}{1.280810in}}%
\pgfpathlineto{\pgfqpoint{2.417926in}{1.293553in}}%
\pgfpathlineto{\pgfqpoint{2.470036in}{1.308756in}}%
\pgfpathlineto{\pgfqpoint{2.548073in}{1.330245in}}%
\pgfpathlineto{\pgfqpoint{2.621647in}{1.348962in}}%
\pgfpathlineto{\pgfqpoint{2.663633in}{1.358726in}}%
\pgfpathlineto{\pgfqpoint{2.691949in}{1.364356in}}%
\pgfpathlineto{\pgfqpoint{2.718024in}{1.368660in}}%
\pgfpathlineto{\pgfqpoint{2.742186in}{1.371754in}}%
\pgfpathlineto{\pgfqpoint{2.771873in}{1.374388in}}%
\pgfpathlineto{\pgfqpoint{2.799106in}{1.375730in}}%
\pgfpathlineto{\pgfqpoint{2.836151in}{1.376364in}}%
\pgfpathlineto{\pgfqpoint{2.952945in}{1.376488in}}%
\pgfpathlineto{\pgfqpoint{3.043110in}{1.376489in}}%
\pgfpathlineto{\pgfqpoint{3.043110in}{1.376489in}}%
\pgfusepath{stroke}%
\end{pgfscope}%
\begin{pgfscope}%
\pgfpathrectangle{\pgfqpoint{0.650000in}{0.420000in}}{\pgfqpoint{2.388110in}{1.994488in}}%
\pgfusepath{clip}%
\pgfsetbuttcap%
\pgfsetroundjoin%
\pgfsetlinewidth{1.003750pt}%
\definecolor{currentstroke}{rgb}{0.498039,0.737255,0.254902}%
\pgfsetstrokecolor{currentstroke}%
\pgfsetdash{{1.000000pt}{1.000000pt}}{0.000000pt}%
\pgfpathmoveto{\pgfqpoint{0.645000in}{1.407289in}}%
\pgfpathlineto{\pgfqpoint{0.720447in}{1.413254in}}%
\pgfpathlineto{\pgfqpoint{0.833824in}{1.423304in}}%
\pgfpathlineto{\pgfqpoint{0.995541in}{1.437830in}}%
\pgfpathlineto{\pgfqpoint{1.056097in}{1.442324in}}%
\pgfpathlineto{\pgfqpoint{1.113099in}{1.445537in}}%
\pgfpathlineto{\pgfqpoint{1.167276in}{1.447380in}}%
\pgfpathlineto{\pgfqpoint{1.219166in}{1.447819in}}%
\pgfpathlineto{\pgfqpoint{1.269179in}{1.446821in}}%
\pgfpathlineto{\pgfqpoint{1.317636in}{1.444509in}}%
\pgfpathlineto{\pgfqpoint{1.364787in}{1.441052in}}%
\pgfpathlineto{\pgfqpoint{1.410836in}{1.436501in}}%
\pgfpathlineto{\pgfqpoint{1.455948in}{1.430848in}}%
\pgfpathlineto{\pgfqpoint{1.500257in}{1.424201in}}%
\pgfpathlineto{\pgfqpoint{1.543878in}{1.416526in}}%
\pgfpathlineto{\pgfqpoint{1.586903in}{1.407931in}}%
\pgfpathlineto{\pgfqpoint{1.643480in}{1.395435in}}%
\pgfpathlineto{\pgfqpoint{1.713145in}{1.380040in}}%
\pgfpathlineto{\pgfqpoint{1.754475in}{1.372459in}}%
\pgfpathlineto{\pgfqpoint{1.781860in}{1.368526in}}%
\pgfpathlineto{\pgfqpoint{1.809124in}{1.365700in}}%
\pgfpathlineto{\pgfqpoint{1.836277in}{1.364056in}}%
\pgfpathlineto{\pgfqpoint{1.863327in}{1.363619in}}%
\pgfpathlineto{\pgfqpoint{1.890284in}{1.364317in}}%
\pgfpathlineto{\pgfqpoint{1.917154in}{1.366030in}}%
\pgfpathlineto{\pgfqpoint{1.957314in}{1.369946in}}%
\pgfpathlineto{\pgfqpoint{2.037181in}{1.378347in}}%
\pgfpathlineto{\pgfqpoint{2.076922in}{1.380722in}}%
\pgfpathlineto{\pgfqpoint{2.129745in}{1.382322in}}%
\pgfpathlineto{\pgfqpoint{2.169254in}{1.382690in}}%
\pgfpathlineto{\pgfqpoint{2.208681in}{1.381851in}}%
\pgfpathlineto{\pgfqpoint{2.339629in}{1.377202in}}%
\pgfpathlineto{\pgfqpoint{2.391846in}{1.374563in}}%
\pgfpathlineto{\pgfqpoint{2.430959in}{1.374244in}}%
\pgfpathlineto{\pgfqpoint{2.496070in}{1.374743in}}%
\pgfpathlineto{\pgfqpoint{2.632655in}{1.373042in}}%
\pgfpathlineto{\pgfqpoint{2.718024in}{1.373577in}}%
\pgfpathlineto{\pgfqpoint{2.778895in}{1.375762in}}%
\pgfpathlineto{\pgfqpoint{2.818149in}{1.376697in}}%
\pgfpathlineto{\pgfqpoint{2.874688in}{1.376622in}}%
\pgfpathlineto{\pgfqpoint{3.043110in}{1.376489in}}%
\pgfpathlineto{\pgfqpoint{3.043110in}{1.376489in}}%
\pgfusepath{stroke}%
\end{pgfscope}%
\begin{pgfscope}%
\pgfpathrectangle{\pgfqpoint{0.650000in}{0.420000in}}{\pgfqpoint{2.388110in}{1.994488in}}%
\pgfusepath{clip}%
\pgfsetbuttcap%
\pgfsetroundjoin%
\pgfsetlinewidth{1.003750pt}%
\definecolor{currentstroke}{rgb}{0.498039,0.737255,0.254902}%
\pgfsetstrokecolor{currentstroke}%
\pgfsetdash{{5.000000pt}{5.000000pt}}{0.000000pt}%
\pgfpathmoveto{\pgfqpoint{0.645000in}{1.377360in}}%
\pgfpathlineto{\pgfqpoint{0.750863in}{1.378659in}}%
\pgfpathlineto{\pgfqpoint{0.833824in}{1.380955in}}%
\pgfpathlineto{\pgfqpoint{0.883760in}{1.383395in}}%
\pgfpathlineto{\pgfqpoint{0.930416in}{1.386859in}}%
\pgfpathlineto{\pgfqpoint{0.974402in}{1.391657in}}%
\pgfpathlineto{\pgfqpoint{1.016177in}{1.398145in}}%
\pgfpathlineto{\pgfqpoint{1.036350in}{1.402140in}}%
\pgfpathlineto{\pgfqpoint{1.056097in}{1.406698in}}%
\pgfpathlineto{\pgfqpoint{1.075452in}{1.411862in}}%
\pgfpathlineto{\pgfqpoint{1.094444in}{1.417671in}}%
\pgfpathlineto{\pgfqpoint{1.113099in}{1.424155in}}%
\pgfpathlineto{\pgfqpoint{1.131442in}{1.431335in}}%
\pgfpathlineto{\pgfqpoint{1.149495in}{1.439213in}}%
\pgfpathlineto{\pgfqpoint{1.167276in}{1.447778in}}%
\pgfpathlineto{\pgfqpoint{1.184804in}{1.456992in}}%
\pgfpathlineto{\pgfqpoint{1.219166in}{1.477083in}}%
\pgfpathlineto{\pgfqpoint{1.252695in}{1.498589in}}%
\pgfpathlineto{\pgfqpoint{1.285491in}{1.520038in}}%
\pgfpathlineto{\pgfqpoint{1.301640in}{1.530112in}}%
\pgfpathlineto{\pgfqpoint{1.317636in}{1.539351in}}%
\pgfpathlineto{\pgfqpoint{1.333487in}{1.547424in}}%
\pgfpathlineto{\pgfqpoint{1.349202in}{1.553977in}}%
\pgfpathlineto{\pgfqpoint{1.364787in}{1.558641in}}%
\pgfpathlineto{\pgfqpoint{1.380251in}{1.561052in}}%
\pgfpathlineto{\pgfqpoint{1.395598in}{1.560877in}}%
\pgfpathlineto{\pgfqpoint{1.410836in}{1.557828in}}%
\pgfpathlineto{\pgfqpoint{1.425971in}{1.551675in}}%
\pgfpathlineto{\pgfqpoint{1.441006in}{1.542256in}}%
\pgfpathlineto{\pgfqpoint{1.455948in}{1.529486in}}%
\pgfpathlineto{\pgfqpoint{1.470801in}{1.513370in}}%
\pgfpathlineto{\pgfqpoint{1.485569in}{1.494017in}}%
\pgfpathlineto{\pgfqpoint{1.500257in}{1.471646in}}%
\pgfpathlineto{\pgfqpoint{1.514869in}{1.446590in}}%
\pgfpathlineto{\pgfqpoint{1.529408in}{1.419281in}}%
\pgfpathlineto{\pgfqpoint{1.543878in}{1.390236in}}%
\pgfpathlineto{\pgfqpoint{1.615295in}{1.240761in}}%
\pgfpathlineto{\pgfqpoint{1.629412in}{1.214529in}}%
\pgfpathlineto{\pgfqpoint{1.643480in}{1.190878in}}%
\pgfpathlineto{\pgfqpoint{1.657499in}{1.170245in}}%
\pgfpathlineto{\pgfqpoint{1.671474in}{1.152968in}}%
\pgfpathlineto{\pgfqpoint{1.685405in}{1.139272in}}%
\pgfpathlineto{\pgfqpoint{1.699295in}{1.129264in}}%
\pgfpathlineto{\pgfqpoint{1.713145in}{1.122936in}}%
\pgfpathlineto{\pgfqpoint{1.726957in}{1.120178in}}%
\pgfpathlineto{\pgfqpoint{1.740734in}{1.120795in}}%
\pgfpathlineto{\pgfqpoint{1.754475in}{1.124511in}}%
\pgfpathlineto{\pgfqpoint{1.768184in}{1.130988in}}%
\pgfpathlineto{\pgfqpoint{1.781860in}{1.139865in}}%
\pgfpathlineto{\pgfqpoint{1.795507in}{1.150775in}}%
\pgfpathlineto{\pgfqpoint{1.809124in}{1.163345in}}%
\pgfpathlineto{\pgfqpoint{1.822714in}{1.177186in}}%
\pgfpathlineto{\pgfqpoint{1.849814in}{1.207180in}}%
\pgfpathlineto{\pgfqpoint{1.903729in}{1.268296in}}%
\pgfpathlineto{\pgfqpoint{1.930560in}{1.296025in}}%
\pgfpathlineto{\pgfqpoint{1.943946in}{1.308574in}}%
\pgfpathlineto{\pgfqpoint{1.957314in}{1.320063in}}%
\pgfpathlineto{\pgfqpoint{1.970665in}{1.330386in}}%
\pgfpathlineto{\pgfqpoint{1.983999in}{1.339512in}}%
\pgfpathlineto{\pgfqpoint{1.997317in}{1.347522in}}%
\pgfpathlineto{\pgfqpoint{2.010619in}{1.354557in}}%
\pgfpathlineto{\pgfqpoint{2.023907in}{1.360741in}}%
\pgfpathlineto{\pgfqpoint{2.037181in}{1.366166in}}%
\pgfpathlineto{\pgfqpoint{2.050441in}{1.370907in}}%
\pgfpathlineto{\pgfqpoint{2.063688in}{1.375027in}}%
\pgfpathlineto{\pgfqpoint{2.090145in}{1.381538in}}%
\pgfpathlineto{\pgfqpoint{2.116556in}{1.386065in}}%
\pgfpathlineto{\pgfqpoint{2.142924in}{1.388995in}}%
\pgfpathlineto{\pgfqpoint{2.169254in}{1.390705in}}%
\pgfpathlineto{\pgfqpoint{2.195547in}{1.390995in}}%
\pgfpathlineto{\pgfqpoint{2.221807in}{1.390086in}}%
\pgfpathlineto{\pgfqpoint{2.248036in}{1.388239in}}%
\pgfpathlineto{\pgfqpoint{2.313490in}{1.381942in}}%
\pgfpathlineto{\pgfqpoint{2.470036in}{1.375627in}}%
\pgfpathlineto{\pgfqpoint{2.496070in}{1.376172in}}%
\pgfpathlineto{\pgfqpoint{2.535091in}{1.379326in}}%
\pgfpathlineto{\pgfqpoint{2.560996in}{1.381167in}}%
\pgfpathlineto{\pgfqpoint{2.586300in}{1.382477in}}%
\pgfpathlineto{\pgfqpoint{2.610260in}{1.385157in}}%
\pgfpathlineto{\pgfqpoint{2.632655in}{1.388478in}}%
\pgfpathlineto{\pgfqpoint{2.653627in}{1.390551in}}%
\pgfpathlineto{\pgfqpoint{2.682778in}{1.391679in}}%
\pgfpathlineto{\pgfqpoint{2.718024in}{1.391937in}}%
\pgfpathlineto{\pgfqpoint{2.734327in}{1.389358in}}%
\pgfpathlineto{\pgfqpoint{2.778895in}{1.381863in}}%
\pgfpathlineto{\pgfqpoint{2.805576in}{1.378645in}}%
\pgfpathlineto{\pgfqpoint{2.830259in}{1.376398in}}%
\pgfpathlineto{\pgfqpoint{2.853221in}{1.376095in}}%
\pgfpathlineto{\pgfqpoint{2.889918in}{1.376439in}}%
\pgfpathlineto{\pgfqpoint{2.904475in}{1.375858in}}%
\pgfpathlineto{\pgfqpoint{2.918417in}{1.376002in}}%
\pgfpathlineto{\pgfqpoint{2.927394in}{1.375612in}}%
\pgfpathlineto{\pgfqpoint{2.936133in}{1.376408in}}%
\pgfpathlineto{\pgfqpoint{2.944646in}{1.377149in}}%
\pgfpathlineto{\pgfqpoint{2.968942in}{1.376496in}}%
\pgfpathlineto{\pgfqpoint{2.980450in}{1.375690in}}%
\pgfpathlineto{\pgfqpoint{2.995194in}{1.376690in}}%
\pgfpathlineto{\pgfqpoint{3.019509in}{1.376092in}}%
\pgfpathlineto{\pgfqpoint{3.035840in}{1.376875in}}%
\pgfpathlineto{\pgfqpoint{3.043110in}{1.375667in}}%
\pgfpathlineto{\pgfqpoint{3.043110in}{1.375667in}}%
\pgfusepath{stroke}%
\end{pgfscope}%
\begin{pgfscope}%
\pgfpathrectangle{\pgfqpoint{0.650000in}{0.420000in}}{\pgfqpoint{2.388110in}{1.994488in}}%
\pgfusepath{clip}%
\pgfsetbuttcap%
\pgfsetroundjoin%
\pgfsetlinewidth{1.003750pt}%
\definecolor{currentstroke}{rgb}{0.498039,0.737255,0.254902}%
\pgfsetstrokecolor{currentstroke}%
\pgfsetdash{{3.000000pt}{5.000000pt}{1.000000pt}{5.000000pt}}{0.000000pt}%
\pgfpathmoveto{\pgfqpoint{0.645000in}{2.140748in}}%
\pgfpathlineto{\pgfqpoint{0.654433in}{2.138568in}}%
\pgfpathlineto{\pgfqpoint{0.688321in}{2.130129in}}%
\pgfpathlineto{\pgfqpoint{0.750863in}{2.112867in}}%
\pgfpathlineto{\pgfqpoint{0.807378in}{2.095110in}}%
\pgfpathlineto{\pgfqpoint{0.833824in}{2.085963in}}%
\pgfpathlineto{\pgfqpoint{0.859248in}{2.076578in}}%
\pgfpathlineto{\pgfqpoint{0.883760in}{2.066902in}}%
\pgfpathlineto{\pgfqpoint{0.907456in}{2.056873in}}%
\pgfpathlineto{\pgfqpoint{0.930416in}{2.046422in}}%
\pgfpathlineto{\pgfqpoint{0.952711in}{2.035472in}}%
\pgfpathlineto{\pgfqpoint{0.974402in}{2.023932in}}%
\pgfpathlineto{\pgfqpoint{0.995541in}{2.011707in}}%
\pgfpathlineto{\pgfqpoint{1.016177in}{1.998689in}}%
\pgfpathlineto{\pgfqpoint{1.036350in}{1.984763in}}%
\pgfpathlineto{\pgfqpoint{1.056097in}{1.969802in}}%
\pgfpathlineto{\pgfqpoint{1.075452in}{1.953675in}}%
\pgfpathlineto{\pgfqpoint{1.094444in}{1.936243in}}%
\pgfpathlineto{\pgfqpoint{1.113099in}{1.917361in}}%
\pgfpathlineto{\pgfqpoint{1.131442in}{1.896886in}}%
\pgfpathlineto{\pgfqpoint{1.149495in}{1.874676in}}%
\pgfpathlineto{\pgfqpoint{1.167276in}{1.850602in}}%
\pgfpathlineto{\pgfqpoint{1.184804in}{1.824550in}}%
\pgfpathlineto{\pgfqpoint{1.202095in}{1.796438in}}%
\pgfpathlineto{\pgfqpoint{1.219166in}{1.766228in}}%
\pgfpathlineto{\pgfqpoint{1.236028in}{1.733947in}}%
\pgfpathlineto{\pgfqpoint{1.252695in}{1.699699in}}%
\pgfpathlineto{\pgfqpoint{1.269179in}{1.663648in}}%
\pgfpathlineto{\pgfqpoint{1.285491in}{1.625959in}}%
\pgfpathlineto{\pgfqpoint{1.301640in}{1.586747in}}%
\pgfpathlineto{\pgfqpoint{1.333487in}{1.504622in}}%
\pgfpathlineto{\pgfqpoint{1.395598in}{1.338306in}}%
\pgfpathlineto{\pgfqpoint{1.410836in}{1.299322in}}%
\pgfpathlineto{\pgfqpoint{1.425971in}{1.262429in}}%
\pgfpathlineto{\pgfqpoint{1.441006in}{1.228144in}}%
\pgfpathlineto{\pgfqpoint{1.455948in}{1.196955in}}%
\pgfpathlineto{\pgfqpoint{1.470801in}{1.169299in}}%
\pgfpathlineto{\pgfqpoint{1.485569in}{1.145545in}}%
\pgfpathlineto{\pgfqpoint{1.500257in}{1.125964in}}%
\pgfpathlineto{\pgfqpoint{1.514869in}{1.110732in}}%
\pgfpathlineto{\pgfqpoint{1.529408in}{1.099930in}}%
\pgfpathlineto{\pgfqpoint{1.543878in}{1.093526in}}%
\pgfpathlineto{\pgfqpoint{1.558281in}{1.091318in}}%
\pgfpathlineto{\pgfqpoint{1.572622in}{1.092991in}}%
\pgfpathlineto{\pgfqpoint{1.586903in}{1.098210in}}%
\pgfpathlineto{\pgfqpoint{1.601126in}{1.106605in}}%
\pgfpathlineto{\pgfqpoint{1.615295in}{1.117731in}}%
\pgfpathlineto{\pgfqpoint{1.629412in}{1.131107in}}%
\pgfpathlineto{\pgfqpoint{1.643480in}{1.146269in}}%
\pgfpathlineto{\pgfqpoint{1.657499in}{1.162745in}}%
\pgfpathlineto{\pgfqpoint{1.699295in}{1.215323in}}%
\pgfpathlineto{\pgfqpoint{1.726957in}{1.249442in}}%
\pgfpathlineto{\pgfqpoint{1.740734in}{1.265434in}}%
\pgfpathlineto{\pgfqpoint{1.754475in}{1.280428in}}%
\pgfpathlineto{\pgfqpoint{1.768184in}{1.294304in}}%
\pgfpathlineto{\pgfqpoint{1.781860in}{1.307011in}}%
\pgfpathlineto{\pgfqpoint{1.795507in}{1.318505in}}%
\pgfpathlineto{\pgfqpoint{1.809124in}{1.328761in}}%
\pgfpathlineto{\pgfqpoint{1.822714in}{1.337830in}}%
\pgfpathlineto{\pgfqpoint{1.836277in}{1.345847in}}%
\pgfpathlineto{\pgfqpoint{1.849814in}{1.352918in}}%
\pgfpathlineto{\pgfqpoint{1.863327in}{1.359081in}}%
\pgfpathlineto{\pgfqpoint{1.876817in}{1.364348in}}%
\pgfpathlineto{\pgfqpoint{1.890284in}{1.368743in}}%
\pgfpathlineto{\pgfqpoint{1.903729in}{1.372320in}}%
\pgfpathlineto{\pgfqpoint{1.917154in}{1.375188in}}%
\pgfpathlineto{\pgfqpoint{1.943946in}{1.379287in}}%
\pgfpathlineto{\pgfqpoint{1.970665in}{1.381831in}}%
\pgfpathlineto{\pgfqpoint{1.997317in}{1.383217in}}%
\pgfpathlineto{\pgfqpoint{2.023907in}{1.383594in}}%
\pgfpathlineto{\pgfqpoint{2.063688in}{1.382875in}}%
\pgfpathlineto{\pgfqpoint{2.129745in}{1.380670in}}%
\pgfpathlineto{\pgfqpoint{2.195547in}{1.378755in}}%
\pgfpathlineto{\pgfqpoint{2.287328in}{1.377273in}}%
\pgfpathlineto{\pgfqpoint{2.404888in}{1.376126in}}%
\pgfpathlineto{\pgfqpoint{2.522089in}{1.376312in}}%
\pgfpathlineto{\pgfqpoint{3.043110in}{1.376489in}}%
\pgfpathlineto{\pgfqpoint{3.043110in}{1.376489in}}%
\pgfusepath{stroke}%
\end{pgfscope}%
\begin{pgfscope}%
\pgfpathrectangle{\pgfqpoint{0.650000in}{0.420000in}}{\pgfqpoint{2.388110in}{1.994488in}}%
\pgfusepath{clip}%
\pgfsetbuttcap%
\pgfsetroundjoin%
\pgfsetlinewidth{1.003750pt}%
\definecolor{currentstroke}{rgb}{0.498039,0.737255,0.254902}%
\pgfsetstrokecolor{currentstroke}%
\pgfsetdash{{1.000000pt}{5.000000pt}}{0.000000pt}%
\pgfpathmoveto{\pgfqpoint{0.645000in}{1.376319in}}%
\pgfpathlineto{\pgfqpoint{1.601126in}{1.377199in}}%
\pgfpathlineto{\pgfqpoint{1.768184in}{1.378952in}}%
\pgfpathlineto{\pgfqpoint{1.822714in}{1.379140in}}%
\pgfpathlineto{\pgfqpoint{1.903729in}{1.378868in}}%
\pgfpathlineto{\pgfqpoint{1.957314in}{1.379011in}}%
\pgfpathlineto{\pgfqpoint{1.997317in}{1.377929in}}%
\pgfpathlineto{\pgfqpoint{2.063688in}{1.374781in}}%
\pgfpathlineto{\pgfqpoint{2.116556in}{1.370582in}}%
\pgfpathlineto{\pgfqpoint{2.195547in}{1.363148in}}%
\pgfpathlineto{\pgfqpoint{2.234925in}{1.359561in}}%
\pgfpathlineto{\pgfqpoint{2.326562in}{1.352772in}}%
\pgfpathlineto{\pgfqpoint{2.365747in}{1.351945in}}%
\pgfpathlineto{\pgfqpoint{2.417926in}{1.352152in}}%
\pgfpathlineto{\pgfqpoint{2.443989in}{1.352843in}}%
\pgfpathlineto{\pgfqpoint{2.483055in}{1.355227in}}%
\pgfpathlineto{\pgfqpoint{2.522089in}{1.358346in}}%
\pgfpathlineto{\pgfqpoint{2.560996in}{1.362315in}}%
\pgfpathlineto{\pgfqpoint{2.610260in}{1.365787in}}%
\pgfpathlineto{\pgfqpoint{2.653627in}{1.370297in}}%
\pgfpathlineto{\pgfqpoint{2.726276in}{1.375570in}}%
\pgfpathlineto{\pgfqpoint{2.742186in}{1.375095in}}%
\pgfpathlineto{\pgfqpoint{2.785770in}{1.376220in}}%
\pgfpathlineto{\pgfqpoint{2.811922in}{1.376315in}}%
\pgfpathlineto{\pgfqpoint{2.841940in}{1.375953in}}%
\pgfpathlineto{\pgfqpoint{2.936133in}{1.376340in}}%
\pgfpathlineto{\pgfqpoint{2.944646in}{1.377145in}}%
\pgfpathlineto{\pgfqpoint{2.968942in}{1.376501in}}%
\pgfpathlineto{\pgfqpoint{2.980450in}{1.375683in}}%
\pgfpathlineto{\pgfqpoint{2.995194in}{1.376644in}}%
\pgfpathlineto{\pgfqpoint{3.019509in}{1.376098in}}%
\pgfpathlineto{\pgfqpoint{3.035840in}{1.376752in}}%
\pgfpathlineto{\pgfqpoint{3.043110in}{1.375626in}}%
\pgfpathlineto{\pgfqpoint{3.043110in}{1.375626in}}%
\pgfusepath{stroke}%
\end{pgfscope}%
\begin{pgfscope}%
\pgfpathrectangle{\pgfqpoint{0.650000in}{0.420000in}}{\pgfqpoint{2.388110in}{1.994488in}}%
\pgfusepath{clip}%
\pgfsetrectcap%
\pgfsetroundjoin%
\pgfsetlinewidth{1.003750pt}%
\definecolor{currentstroke}{rgb}{0.400000,0.760784,0.643137}%
\pgfsetstrokecolor{currentstroke}%
\pgfsetdash{}{0pt}%
\pgfpathmoveto{\pgfqpoint{0.645000in}{1.377316in}}%
\pgfpathlineto{\pgfqpoint{0.795272in}{1.378421in}}%
\pgfpathlineto{\pgfqpoint{0.932227in}{1.383347in}}%
\pgfpathlineto{\pgfqpoint{1.030365in}{1.393431in}}%
\pgfpathlineto{\pgfqpoint{1.107239in}{1.410621in}}%
\pgfpathlineto{\pgfqpoint{1.170666in}{1.436642in}}%
\pgfpathlineto{\pgfqpoint{1.224817in}{1.472738in}}%
\pgfpathlineto{\pgfqpoint{1.272181in}{1.519463in}}%
\pgfpathlineto{\pgfqpoint{1.314364in}{1.576537in}}%
\pgfpathlineto{\pgfqpoint{1.352464in}{1.642825in}}%
\pgfpathlineto{\pgfqpoint{1.387260in}{1.716424in}}%
\pgfpathlineto{\pgfqpoint{1.419332in}{1.794838in}}%
\pgfpathlineto{\pgfqpoint{1.449117in}{1.875219in}}%
\pgfpathlineto{\pgfqpoint{1.527827in}{2.100045in}}%
\pgfpathlineto{\pgfqpoint{1.551254in}{2.161661in}}%
\pgfpathlineto{\pgfqpoint{1.573549in}{2.213908in}}%
\pgfpathlineto{\pgfqpoint{1.594835in}{2.255980in}}%
\pgfpathlineto{\pgfqpoint{1.615218in}{2.287641in}}%
\pgfpathlineto{\pgfqpoint{1.634786in}{2.309090in}}%
\pgfpathlineto{\pgfqpoint{1.653617in}{2.320881in}}%
\pgfpathlineto{\pgfqpoint{1.671776in}{2.323830in}}%
\pgfpathlineto{\pgfqpoint{1.689323in}{2.318927in}}%
\pgfpathlineto{\pgfqpoint{1.706307in}{2.307241in}}%
\pgfpathlineto{\pgfqpoint{1.722774in}{2.289854in}}%
\pgfpathlineto{\pgfqpoint{1.738764in}{2.267824in}}%
\pgfpathlineto{\pgfqpoint{1.754311in}{2.242131in}}%
\pgfpathlineto{\pgfqpoint{1.769449in}{2.213677in}}%
\pgfpathlineto{\pgfqpoint{1.784205in}{2.183258in}}%
\pgfpathlineto{\pgfqpoint{1.798605in}{2.151561in}}%
\pgfpathlineto{\pgfqpoint{1.826428in}{2.086583in}}%
\pgfpathlineto{\pgfqpoint{1.891152in}{1.931858in}}%
\pgfpathlineto{\pgfqpoint{1.915421in}{1.877351in}}%
\pgfpathlineto{\pgfqpoint{1.938907in}{1.828167in}}%
\pgfpathlineto{\pgfqpoint{1.961687in}{1.784359in}}%
\pgfpathlineto{\pgfqpoint{1.983825in}{1.745692in}}%
\pgfpathlineto{\pgfqpoint{2.005382in}{1.711730in}}%
\pgfpathlineto{\pgfqpoint{2.026407in}{1.682028in}}%
\pgfpathlineto{\pgfqpoint{2.046945in}{1.656051in}}%
\pgfpathlineto{\pgfqpoint{2.067036in}{1.633385in}}%
\pgfpathlineto{\pgfqpoint{2.086715in}{1.613601in}}%
\pgfpathlineto{\pgfqpoint{2.106014in}{1.596287in}}%
\pgfpathlineto{\pgfqpoint{2.124961in}{1.581075in}}%
\pgfpathlineto{\pgfqpoint{2.143582in}{1.567690in}}%
\pgfpathlineto{\pgfqpoint{2.161899in}{1.555940in}}%
\pgfpathlineto{\pgfqpoint{2.179935in}{1.545738in}}%
\pgfpathlineto{\pgfqpoint{2.197707in}{1.536902in}}%
\pgfpathlineto{\pgfqpoint{2.215233in}{1.529287in}}%
\pgfpathlineto{\pgfqpoint{2.241096in}{1.519435in}}%
\pgfpathlineto{\pgfqpoint{2.299689in}{1.498413in}}%
\pgfpathlineto{\pgfqpoint{2.324144in}{1.491230in}}%
\pgfpathlineto{\pgfqpoint{2.348252in}{1.485045in}}%
\pgfpathlineto{\pgfqpoint{2.434126in}{1.464206in}}%
\pgfpathlineto{\pgfqpoint{2.472057in}{1.453651in}}%
\pgfpathlineto{\pgfqpoint{2.582689in}{1.421124in}}%
\pgfpathlineto{\pgfqpoint{2.604346in}{1.414304in}}%
\pgfpathlineto{\pgfqpoint{2.661480in}{1.394289in}}%
\pgfpathlineto{\pgfqpoint{2.689747in}{1.386483in}}%
\pgfpathlineto{\pgfqpoint{2.710830in}{1.381690in}}%
\pgfpathlineto{\pgfqpoint{2.724834in}{1.379429in}}%
\pgfpathlineto{\pgfqpoint{2.745767in}{1.377482in}}%
\pgfpathlineto{\pgfqpoint{2.773549in}{1.376621in}}%
\pgfpathlineto{\pgfqpoint{2.849289in}{1.376489in}}%
\pgfpathlineto{\pgfqpoint{3.043110in}{1.376489in}}%
\pgfpathlineto{\pgfqpoint{3.043110in}{1.376489in}}%
\pgfusepath{stroke}%
\end{pgfscope}%
\begin{pgfscope}%
\pgfpathrectangle{\pgfqpoint{0.650000in}{0.420000in}}{\pgfqpoint{2.388110in}{1.994488in}}%
\pgfusepath{clip}%
\pgfsetbuttcap%
\pgfsetroundjoin%
\pgfsetlinewidth{1.003750pt}%
\definecolor{currentstroke}{rgb}{0.400000,0.760784,0.643137}%
\pgfsetstrokecolor{currentstroke}%
\pgfsetdash{{5.000000pt}{1.000000pt}}{0.000000pt}%
\pgfpathmoveto{\pgfqpoint{0.645000in}{1.376489in}}%
\pgfpathlineto{\pgfqpoint{1.387260in}{1.376904in}}%
\pgfpathlineto{\pgfqpoint{1.573549in}{1.378421in}}%
\pgfpathlineto{\pgfqpoint{1.706307in}{1.379035in}}%
\pgfpathlineto{\pgfqpoint{1.994673in}{1.378438in}}%
\pgfpathlineto{\pgfqpoint{2.036734in}{1.376874in}}%
\pgfpathlineto{\pgfqpoint{2.086715in}{1.374641in}}%
\pgfpathlineto{\pgfqpoint{2.124961in}{1.374494in}}%
\pgfpathlineto{\pgfqpoint{2.161899in}{1.375378in}}%
\pgfpathlineto{\pgfqpoint{2.206500in}{1.376882in}}%
\pgfpathlineto{\pgfqpoint{2.266490in}{1.377900in}}%
\pgfpathlineto{\pgfqpoint{2.307881in}{1.378586in}}%
\pgfpathlineto{\pgfqpoint{2.332217in}{1.378189in}}%
\pgfpathlineto{\pgfqpoint{2.356217in}{1.377571in}}%
\pgfpathlineto{\pgfqpoint{2.418776in}{1.377905in}}%
\pgfpathlineto{\pgfqpoint{2.441762in}{1.377884in}}%
\pgfpathlineto{\pgfqpoint{2.516827in}{1.375562in}}%
\pgfpathlineto{\pgfqpoint{2.531588in}{1.375529in}}%
\pgfpathlineto{\pgfqpoint{2.546275in}{1.374289in}}%
\pgfpathlineto{\pgfqpoint{2.575439in}{1.370972in}}%
\pgfpathlineto{\pgfqpoint{2.597142in}{1.369069in}}%
\pgfpathlineto{\pgfqpoint{2.654383in}{1.366219in}}%
\pgfpathlineto{\pgfqpoint{2.675636in}{1.366183in}}%
\pgfpathlineto{\pgfqpoint{2.717837in}{1.366902in}}%
\pgfpathlineto{\pgfqpoint{2.731821in}{1.367882in}}%
\pgfpathlineto{\pgfqpoint{2.787389in}{1.374315in}}%
\pgfpathlineto{\pgfqpoint{2.814973in}{1.375582in}}%
\pgfpathlineto{\pgfqpoint{2.849289in}{1.376063in}}%
\pgfpathlineto{\pgfqpoint{2.910653in}{1.376466in}}%
\pgfpathlineto{\pgfqpoint{2.978320in}{1.376525in}}%
\pgfpathlineto{\pgfqpoint{3.043110in}{1.376432in}}%
\pgfpathlineto{\pgfqpoint{3.043110in}{1.376432in}}%
\pgfusepath{stroke}%
\end{pgfscope}%
\begin{pgfscope}%
\pgfpathrectangle{\pgfqpoint{0.650000in}{0.420000in}}{\pgfqpoint{2.388110in}{1.994488in}}%
\pgfusepath{clip}%
\pgfsetbuttcap%
\pgfsetroundjoin%
\pgfsetlinewidth{1.003750pt}%
\definecolor{currentstroke}{rgb}{0.400000,0.760784,0.643137}%
\pgfsetstrokecolor{currentstroke}%
\pgfsetdash{{3.000000pt}{1.000000pt}{1.000000pt}{1.000000pt}}{0.000000pt}%
\pgfpathmoveto{\pgfqpoint{0.645000in}{0.576010in}}%
\pgfpathlineto{\pgfqpoint{0.795272in}{0.600860in}}%
\pgfpathlineto{\pgfqpoint{0.932227in}{0.630211in}}%
\pgfpathlineto{\pgfqpoint{1.030365in}{0.654433in}}%
\pgfpathlineto{\pgfqpoint{1.107239in}{0.676395in}}%
\pgfpathlineto{\pgfqpoint{1.170666in}{0.697911in}}%
\pgfpathlineto{\pgfqpoint{1.224817in}{0.719793in}}%
\pgfpathlineto{\pgfqpoint{1.272181in}{0.742129in}}%
\pgfpathlineto{\pgfqpoint{1.314364in}{0.764190in}}%
\pgfpathlineto{\pgfqpoint{1.352464in}{0.784856in}}%
\pgfpathlineto{\pgfqpoint{1.387260in}{0.803051in}}%
\pgfpathlineto{\pgfqpoint{1.419332in}{0.818193in}}%
\pgfpathlineto{\pgfqpoint{1.449117in}{0.829906in}}%
\pgfpathlineto{\pgfqpoint{1.476957in}{0.837870in}}%
\pgfpathlineto{\pgfqpoint{1.503121in}{0.842565in}}%
\pgfpathlineto{\pgfqpoint{1.527827in}{0.844635in}}%
\pgfpathlineto{\pgfqpoint{1.551254in}{0.844776in}}%
\pgfpathlineto{\pgfqpoint{1.653617in}{0.840239in}}%
\pgfpathlineto{\pgfqpoint{1.671776in}{0.841065in}}%
\pgfpathlineto{\pgfqpoint{1.689323in}{0.843068in}}%
\pgfpathlineto{\pgfqpoint{1.706307in}{0.846302in}}%
\pgfpathlineto{\pgfqpoint{1.722774in}{0.850682in}}%
\pgfpathlineto{\pgfqpoint{1.738764in}{0.856199in}}%
\pgfpathlineto{\pgfqpoint{1.754311in}{0.862474in}}%
\pgfpathlineto{\pgfqpoint{1.769449in}{0.869592in}}%
\pgfpathlineto{\pgfqpoint{1.784205in}{0.877349in}}%
\pgfpathlineto{\pgfqpoint{1.798605in}{0.885678in}}%
\pgfpathlineto{\pgfqpoint{1.826428in}{0.903729in}}%
\pgfpathlineto{\pgfqpoint{1.853080in}{0.922916in}}%
\pgfpathlineto{\pgfqpoint{1.878696in}{0.943138in}}%
\pgfpathlineto{\pgfqpoint{1.903390in}{0.963650in}}%
\pgfpathlineto{\pgfqpoint{1.950381in}{1.004466in}}%
\pgfpathlineto{\pgfqpoint{2.005382in}{1.052905in}}%
\pgfpathlineto{\pgfqpoint{2.067036in}{1.104990in}}%
\pgfpathlineto{\pgfqpoint{2.096410in}{1.128271in}}%
\pgfpathlineto{\pgfqpoint{2.124961in}{1.149671in}}%
\pgfpathlineto{\pgfqpoint{2.152777in}{1.169096in}}%
\pgfpathlineto{\pgfqpoint{2.179935in}{1.186408in}}%
\pgfpathlineto{\pgfqpoint{2.206500in}{1.201608in}}%
\pgfpathlineto{\pgfqpoint{2.241096in}{1.219373in}}%
\pgfpathlineto{\pgfqpoint{2.291456in}{1.243303in}}%
\pgfpathlineto{\pgfqpoint{2.324144in}{1.257039in}}%
\pgfpathlineto{\pgfqpoint{2.348252in}{1.266255in}}%
\pgfpathlineto{\pgfqpoint{2.387741in}{1.279899in}}%
\pgfpathlineto{\pgfqpoint{2.426464in}{1.292569in}}%
\pgfpathlineto{\pgfqpoint{2.538941in}{1.325712in}}%
\pgfpathlineto{\pgfqpoint{2.560891in}{1.331935in}}%
\pgfpathlineto{\pgfqpoint{2.604346in}{1.343984in}}%
\pgfpathlineto{\pgfqpoint{2.625872in}{1.349622in}}%
\pgfpathlineto{\pgfqpoint{2.661480in}{1.358436in}}%
\pgfpathlineto{\pgfqpoint{2.682697in}{1.363059in}}%
\pgfpathlineto{\pgfqpoint{2.731821in}{1.371318in}}%
\pgfpathlineto{\pgfqpoint{2.752725in}{1.373484in}}%
\pgfpathlineto{\pgfqpoint{2.780473in}{1.375399in}}%
\pgfpathlineto{\pgfqpoint{2.821850in}{1.376309in}}%
\pgfpathlineto{\pgfqpoint{2.917442in}{1.376487in}}%
\pgfpathlineto{\pgfqpoint{3.043110in}{1.376489in}}%
\pgfpathlineto{\pgfqpoint{3.043110in}{1.376489in}}%
\pgfusepath{stroke}%
\end{pgfscope}%
\begin{pgfscope}%
\pgfpathrectangle{\pgfqpoint{0.650000in}{0.420000in}}{\pgfqpoint{2.388110in}{1.994488in}}%
\pgfusepath{clip}%
\pgfsetbuttcap%
\pgfsetroundjoin%
\pgfsetlinewidth{1.003750pt}%
\definecolor{currentstroke}{rgb}{0.400000,0.760784,0.643137}%
\pgfsetstrokecolor{currentstroke}%
\pgfsetdash{{1.000000pt}{1.000000pt}}{0.000000pt}%
\pgfpathmoveto{\pgfqpoint{0.645000in}{1.409755in}}%
\pgfpathlineto{\pgfqpoint{0.795272in}{1.422116in}}%
\pgfpathlineto{\pgfqpoint{0.932227in}{1.435157in}}%
\pgfpathlineto{\pgfqpoint{1.030365in}{1.443437in}}%
\pgfpathlineto{\pgfqpoint{1.107239in}{1.448169in}}%
\pgfpathlineto{\pgfqpoint{1.170666in}{1.449975in}}%
\pgfpathlineto{\pgfqpoint{1.224817in}{1.450173in}}%
\pgfpathlineto{\pgfqpoint{1.272181in}{1.448903in}}%
\pgfpathlineto{\pgfqpoint{1.314364in}{1.446626in}}%
\pgfpathlineto{\pgfqpoint{1.352464in}{1.443627in}}%
\pgfpathlineto{\pgfqpoint{1.387260in}{1.440277in}}%
\pgfpathlineto{\pgfqpoint{1.449117in}{1.432550in}}%
\pgfpathlineto{\pgfqpoint{1.503121in}{1.424165in}}%
\pgfpathlineto{\pgfqpoint{1.551254in}{1.415202in}}%
\pgfpathlineto{\pgfqpoint{1.594835in}{1.406007in}}%
\pgfpathlineto{\pgfqpoint{1.671776in}{1.388799in}}%
\pgfpathlineto{\pgfqpoint{1.722774in}{1.378066in}}%
\pgfpathlineto{\pgfqpoint{1.754311in}{1.372770in}}%
\pgfpathlineto{\pgfqpoint{1.784205in}{1.368833in}}%
\pgfpathlineto{\pgfqpoint{1.812672in}{1.366230in}}%
\pgfpathlineto{\pgfqpoint{1.839891in}{1.364712in}}%
\pgfpathlineto{\pgfqpoint{1.866010in}{1.364243in}}%
\pgfpathlineto{\pgfqpoint{1.903390in}{1.365004in}}%
\pgfpathlineto{\pgfqpoint{1.950381in}{1.367698in}}%
\pgfpathlineto{\pgfqpoint{2.015958in}{1.373152in}}%
\pgfpathlineto{\pgfqpoint{2.046945in}{1.376302in}}%
\pgfpathlineto{\pgfqpoint{2.086715in}{1.380383in}}%
\pgfpathlineto{\pgfqpoint{2.106014in}{1.381519in}}%
\pgfpathlineto{\pgfqpoint{2.124961in}{1.381877in}}%
\pgfpathlineto{\pgfqpoint{2.152777in}{1.381249in}}%
\pgfpathlineto{\pgfqpoint{2.232530in}{1.378552in}}%
\pgfpathlineto{\pgfqpoint{2.291456in}{1.377105in}}%
\pgfpathlineto{\pgfqpoint{2.332217in}{1.375779in}}%
\pgfpathlineto{\pgfqpoint{2.426464in}{1.376424in}}%
\pgfpathlineto{\pgfqpoint{2.464519in}{1.375047in}}%
\pgfpathlineto{\pgfqpoint{2.516827in}{1.373603in}}%
\pgfpathlineto{\pgfqpoint{2.582689in}{1.373917in}}%
\pgfpathlineto{\pgfqpoint{2.618710in}{1.372633in}}%
\pgfpathlineto{\pgfqpoint{2.689747in}{1.372048in}}%
\pgfpathlineto{\pgfqpoint{2.710830in}{1.372814in}}%
\pgfpathlineto{\pgfqpoint{2.780473in}{1.376268in}}%
\pgfpathlineto{\pgfqpoint{2.828720in}{1.376748in}}%
\pgfpathlineto{\pgfqpoint{3.043110in}{1.376494in}}%
\pgfpathlineto{\pgfqpoint{3.043110in}{1.376494in}}%
\pgfusepath{stroke}%
\end{pgfscope}%
\begin{pgfscope}%
\pgfpathrectangle{\pgfqpoint{0.650000in}{0.420000in}}{\pgfqpoint{2.388110in}{1.994488in}}%
\pgfusepath{clip}%
\pgfsetbuttcap%
\pgfsetroundjoin%
\pgfsetlinewidth{1.003750pt}%
\definecolor{currentstroke}{rgb}{0.400000,0.760784,0.643137}%
\pgfsetstrokecolor{currentstroke}%
\pgfsetdash{{5.000000pt}{5.000000pt}}{0.000000pt}%
\pgfpathmoveto{\pgfqpoint{0.645000in}{1.377707in}}%
\pgfpathlineto{\pgfqpoint{0.795272in}{1.379319in}}%
\pgfpathlineto{\pgfqpoint{0.932227in}{1.386226in}}%
\pgfpathlineto{\pgfqpoint{1.030365in}{1.399512in}}%
\pgfpathlineto{\pgfqpoint{1.107239in}{1.420092in}}%
\pgfpathlineto{\pgfqpoint{1.170666in}{1.447207in}}%
\pgfpathlineto{\pgfqpoint{1.224817in}{1.478473in}}%
\pgfpathlineto{\pgfqpoint{1.272181in}{1.509900in}}%
\pgfpathlineto{\pgfqpoint{1.314364in}{1.536970in}}%
\pgfpathlineto{\pgfqpoint{1.352464in}{1.555675in}}%
\pgfpathlineto{\pgfqpoint{1.387260in}{1.562950in}}%
\pgfpathlineto{\pgfqpoint{1.419332in}{1.557229in}}%
\pgfpathlineto{\pgfqpoint{1.449117in}{1.538736in}}%
\pgfpathlineto{\pgfqpoint{1.476957in}{1.508937in}}%
\pgfpathlineto{\pgfqpoint{1.503121in}{1.470031in}}%
\pgfpathlineto{\pgfqpoint{1.527827in}{1.425135in}}%
\pgfpathlineto{\pgfqpoint{1.551254in}{1.377352in}}%
\pgfpathlineto{\pgfqpoint{1.594835in}{1.284084in}}%
\pgfpathlineto{\pgfqpoint{1.615218in}{1.242945in}}%
\pgfpathlineto{\pgfqpoint{1.634786in}{1.207442in}}%
\pgfpathlineto{\pgfqpoint{1.653617in}{1.178280in}}%
\pgfpathlineto{\pgfqpoint{1.671776in}{1.155844in}}%
\pgfpathlineto{\pgfqpoint{1.689323in}{1.139962in}}%
\pgfpathlineto{\pgfqpoint{1.706307in}{1.130141in}}%
\pgfpathlineto{\pgfqpoint{1.722774in}{1.125837in}}%
\pgfpathlineto{\pgfqpoint{1.738764in}{1.126245in}}%
\pgfpathlineto{\pgfqpoint{1.754311in}{1.130718in}}%
\pgfpathlineto{\pgfqpoint{1.769449in}{1.138388in}}%
\pgfpathlineto{\pgfqpoint{1.784205in}{1.148594in}}%
\pgfpathlineto{\pgfqpoint{1.798605in}{1.160732in}}%
\pgfpathlineto{\pgfqpoint{1.812672in}{1.174252in}}%
\pgfpathlineto{\pgfqpoint{1.839891in}{1.203461in}}%
\pgfpathlineto{\pgfqpoint{1.903390in}{1.276358in}}%
\pgfpathlineto{\pgfqpoint{1.927257in}{1.301270in}}%
\pgfpathlineto{\pgfqpoint{1.950381in}{1.322786in}}%
\pgfpathlineto{\pgfqpoint{1.961687in}{1.332106in}}%
\pgfpathlineto{\pgfqpoint{1.972832in}{1.340469in}}%
\pgfpathlineto{\pgfqpoint{1.983825in}{1.347939in}}%
\pgfpathlineto{\pgfqpoint{2.005382in}{1.360648in}}%
\pgfpathlineto{\pgfqpoint{2.026407in}{1.370600in}}%
\pgfpathlineto{\pgfqpoint{2.046945in}{1.378371in}}%
\pgfpathlineto{\pgfqpoint{2.067036in}{1.383831in}}%
\pgfpathlineto{\pgfqpoint{2.086715in}{1.387530in}}%
\pgfpathlineto{\pgfqpoint{2.106014in}{1.390028in}}%
\pgfpathlineto{\pgfqpoint{2.143582in}{1.392692in}}%
\pgfpathlineto{\pgfqpoint{2.170951in}{1.393447in}}%
\pgfpathlineto{\pgfqpoint{2.188853in}{1.392847in}}%
\pgfpathlineto{\pgfqpoint{2.215233in}{1.390901in}}%
\pgfpathlineto{\pgfqpoint{2.258075in}{1.386782in}}%
\pgfpathlineto{\pgfqpoint{2.332217in}{1.382425in}}%
\pgfpathlineto{\pgfqpoint{2.356217in}{1.379275in}}%
\pgfpathlineto{\pgfqpoint{2.379908in}{1.376485in}}%
\pgfpathlineto{\pgfqpoint{2.411061in}{1.373028in}}%
\pgfpathlineto{\pgfqpoint{2.426464in}{1.372667in}}%
\pgfpathlineto{\pgfqpoint{2.441762in}{1.373171in}}%
\pgfpathlineto{\pgfqpoint{2.449372in}{1.374307in}}%
\pgfpathlineto{\pgfqpoint{2.456957in}{1.374345in}}%
\pgfpathlineto{\pgfqpoint{2.487066in}{1.376461in}}%
\pgfpathlineto{\pgfqpoint{2.516827in}{1.376852in}}%
\pgfpathlineto{\pgfqpoint{2.524217in}{1.375919in}}%
\pgfpathlineto{\pgfqpoint{2.546275in}{1.376863in}}%
\pgfpathlineto{\pgfqpoint{2.575439in}{1.380755in}}%
\pgfpathlineto{\pgfqpoint{2.582689in}{1.381098in}}%
\pgfpathlineto{\pgfqpoint{2.589923in}{1.382773in}}%
\pgfpathlineto{\pgfqpoint{2.604346in}{1.384354in}}%
\pgfpathlineto{\pgfqpoint{2.618710in}{1.386561in}}%
\pgfpathlineto{\pgfqpoint{2.633019in}{1.388413in}}%
\pgfpathlineto{\pgfqpoint{2.647275in}{1.390707in}}%
\pgfpathlineto{\pgfqpoint{2.668564in}{1.392771in}}%
\pgfpathlineto{\pgfqpoint{2.703813in}{1.393313in}}%
\pgfpathlineto{\pgfqpoint{2.731821in}{1.390178in}}%
\pgfpathlineto{\pgfqpoint{2.773549in}{1.381068in}}%
\pgfpathlineto{\pgfqpoint{2.801196in}{1.378233in}}%
\pgfpathlineto{\pgfqpoint{2.828720in}{1.377071in}}%
\pgfpathlineto{\pgfqpoint{2.856132in}{1.376722in}}%
\pgfpathlineto{\pgfqpoint{2.876622in}{1.376481in}}%
\pgfpathlineto{\pgfqpoint{2.971574in}{1.376630in}}%
\pgfpathlineto{\pgfqpoint{2.985061in}{1.376270in}}%
\pgfpathlineto{\pgfqpoint{2.998529in}{1.376659in}}%
\pgfpathlineto{\pgfqpoint{3.025418in}{1.376546in}}%
\pgfpathlineto{\pgfqpoint{3.043110in}{1.376608in}}%
\pgfpathlineto{\pgfqpoint{3.043110in}{1.376608in}}%
\pgfusepath{stroke}%
\end{pgfscope}%
\begin{pgfscope}%
\pgfpathrectangle{\pgfqpoint{0.650000in}{0.420000in}}{\pgfqpoint{2.388110in}{1.994488in}}%
\pgfusepath{clip}%
\pgfsetbuttcap%
\pgfsetroundjoin%
\pgfsetlinewidth{1.003750pt}%
\definecolor{currentstroke}{rgb}{0.400000,0.760784,0.643137}%
\pgfsetstrokecolor{currentstroke}%
\pgfsetdash{{3.000000pt}{5.000000pt}{1.000000pt}{5.000000pt}}{0.000000pt}%
\pgfpathmoveto{\pgfqpoint{0.645000in}{2.141663in}}%
\pgfpathlineto{\pgfqpoint{0.795272in}{2.101726in}}%
\pgfpathlineto{\pgfqpoint{0.932227in}{2.047509in}}%
\pgfpathlineto{\pgfqpoint{1.030365in}{1.991595in}}%
\pgfpathlineto{\pgfqpoint{1.107239in}{1.927078in}}%
\pgfpathlineto{\pgfqpoint{1.170666in}{1.850417in}}%
\pgfpathlineto{\pgfqpoint{1.224817in}{1.761126in}}%
\pgfpathlineto{\pgfqpoint{1.272181in}{1.661735in}}%
\pgfpathlineto{\pgfqpoint{1.314364in}{1.557708in}}%
\pgfpathlineto{\pgfqpoint{1.352464in}{1.455022in}}%
\pgfpathlineto{\pgfqpoint{1.387260in}{1.359086in}}%
\pgfpathlineto{\pgfqpoint{1.419332in}{1.275025in}}%
\pgfpathlineto{\pgfqpoint{1.449117in}{1.205005in}}%
\pgfpathlineto{\pgfqpoint{1.476957in}{1.151549in}}%
\pgfpathlineto{\pgfqpoint{1.503121in}{1.114032in}}%
\pgfpathlineto{\pgfqpoint{1.527827in}{1.091505in}}%
\pgfpathlineto{\pgfqpoint{1.551254in}{1.082007in}}%
\pgfpathlineto{\pgfqpoint{1.573549in}{1.082965in}}%
\pgfpathlineto{\pgfqpoint{1.594835in}{1.092972in}}%
\pgfpathlineto{\pgfqpoint{1.615218in}{1.106905in}}%
\pgfpathlineto{\pgfqpoint{1.634786in}{1.126478in}}%
\pgfpathlineto{\pgfqpoint{1.653617in}{1.148245in}}%
\pgfpathlineto{\pgfqpoint{1.689323in}{1.193647in}}%
\pgfpathlineto{\pgfqpoint{1.722774in}{1.235967in}}%
\pgfpathlineto{\pgfqpoint{1.754311in}{1.272063in}}%
\pgfpathlineto{\pgfqpoint{1.769449in}{1.287816in}}%
\pgfpathlineto{\pgfqpoint{1.784205in}{1.301984in}}%
\pgfpathlineto{\pgfqpoint{1.798605in}{1.314700in}}%
\pgfpathlineto{\pgfqpoint{1.812672in}{1.325771in}}%
\pgfpathlineto{\pgfqpoint{1.826428in}{1.335728in}}%
\pgfpathlineto{\pgfqpoint{1.839891in}{1.344329in}}%
\pgfpathlineto{\pgfqpoint{1.853080in}{1.351582in}}%
\pgfpathlineto{\pgfqpoint{1.866010in}{1.357701in}}%
\pgfpathlineto{\pgfqpoint{1.878696in}{1.362736in}}%
\pgfpathlineto{\pgfqpoint{1.891152in}{1.366938in}}%
\pgfpathlineto{\pgfqpoint{1.903390in}{1.370349in}}%
\pgfpathlineto{\pgfqpoint{1.927257in}{1.375578in}}%
\pgfpathlineto{\pgfqpoint{1.961687in}{1.379567in}}%
\pgfpathlineto{\pgfqpoint{1.994673in}{1.382148in}}%
\pgfpathlineto{\pgfqpoint{2.036734in}{1.383050in}}%
\pgfpathlineto{\pgfqpoint{2.197707in}{1.378064in}}%
\pgfpathlineto{\pgfqpoint{2.249611in}{1.378099in}}%
\pgfpathlineto{\pgfqpoint{2.332217in}{1.376292in}}%
\pgfpathlineto{\pgfqpoint{2.356217in}{1.376702in}}%
\pgfpathlineto{\pgfqpoint{2.403317in}{1.376557in}}%
\pgfpathlineto{\pgfqpoint{2.494537in}{1.376092in}}%
\pgfpathlineto{\pgfqpoint{2.509417in}{1.376384in}}%
\pgfpathlineto{\pgfqpoint{2.524217in}{1.376220in}}%
\pgfpathlineto{\pgfqpoint{2.560891in}{1.375995in}}%
\pgfpathlineto{\pgfqpoint{2.589923in}{1.376252in}}%
\pgfpathlineto{\pgfqpoint{2.625872in}{1.376617in}}%
\pgfpathlineto{\pgfqpoint{2.654383in}{1.376438in}}%
\pgfpathlineto{\pgfqpoint{2.696785in}{1.376854in}}%
\pgfpathlineto{\pgfqpoint{2.752725in}{1.376955in}}%
\pgfpathlineto{\pgfqpoint{2.890252in}{1.376494in}}%
\pgfpathlineto{\pgfqpoint{3.043110in}{1.376490in}}%
\pgfpathlineto{\pgfqpoint{3.043110in}{1.376490in}}%
\pgfusepath{stroke}%
\end{pgfscope}%
\begin{pgfscope}%
\pgfpathrectangle{\pgfqpoint{0.650000in}{0.420000in}}{\pgfqpoint{2.388110in}{1.994488in}}%
\pgfusepath{clip}%
\pgfsetbuttcap%
\pgfsetroundjoin%
\pgfsetlinewidth{1.003750pt}%
\definecolor{currentstroke}{rgb}{0.400000,0.760784,0.643137}%
\pgfsetstrokecolor{currentstroke}%
\pgfsetdash{{1.000000pt}{5.000000pt}}{0.000000pt}%
\pgfpathmoveto{\pgfqpoint{0.645000in}{1.376496in}}%
\pgfpathlineto{\pgfqpoint{1.573549in}{1.376826in}}%
\pgfpathlineto{\pgfqpoint{1.594835in}{1.377688in}}%
\pgfpathlineto{\pgfqpoint{1.615218in}{1.376775in}}%
\pgfpathlineto{\pgfqpoint{1.689323in}{1.377117in}}%
\pgfpathlineto{\pgfqpoint{1.826428in}{1.376515in}}%
\pgfpathlineto{\pgfqpoint{1.927257in}{1.375484in}}%
\pgfpathlineto{\pgfqpoint{1.972832in}{1.375581in}}%
\pgfpathlineto{\pgfqpoint{2.046945in}{1.375708in}}%
\pgfpathlineto{\pgfqpoint{2.086715in}{1.375712in}}%
\pgfpathlineto{\pgfqpoint{2.096410in}{1.376249in}}%
\pgfpathlineto{\pgfqpoint{2.115530in}{1.376090in}}%
\pgfpathlineto{\pgfqpoint{2.316032in}{1.379795in}}%
\pgfpathlineto{\pgfqpoint{2.332217in}{1.379536in}}%
\pgfpathlineto{\pgfqpoint{2.364147in}{1.379618in}}%
\pgfpathlineto{\pgfqpoint{2.494537in}{1.379681in}}%
\pgfpathlineto{\pgfqpoint{2.524217in}{1.379070in}}%
\pgfpathlineto{\pgfqpoint{2.538941in}{1.379494in}}%
\pgfpathlineto{\pgfqpoint{2.575439in}{1.378597in}}%
\pgfpathlineto{\pgfqpoint{2.582689in}{1.377162in}}%
\pgfpathlineto{\pgfqpoint{2.589923in}{1.378616in}}%
\pgfpathlineto{\pgfqpoint{2.604346in}{1.378392in}}%
\pgfpathlineto{\pgfqpoint{2.618710in}{1.378418in}}%
\pgfpathlineto{\pgfqpoint{2.633019in}{1.377664in}}%
\pgfpathlineto{\pgfqpoint{2.710830in}{1.376012in}}%
\pgfpathlineto{\pgfqpoint{2.717837in}{1.376837in}}%
\pgfpathlineto{\pgfqpoint{2.759675in}{1.376471in}}%
\pgfpathlineto{\pgfqpoint{2.801196in}{1.376533in}}%
\pgfpathlineto{\pgfqpoint{2.964825in}{1.376500in}}%
\pgfpathlineto{\pgfqpoint{3.011982in}{1.376591in}}%
\pgfpathlineto{\pgfqpoint{3.043110in}{1.376558in}}%
\pgfpathlineto{\pgfqpoint{3.043110in}{1.376558in}}%
\pgfusepath{stroke}%
\end{pgfscope}%
\begin{pgfscope}%
\pgfsetrectcap%
\pgfsetmiterjoin%
\pgfsetlinewidth{0.803000pt}%
\definecolor{currentstroke}{rgb}{0.000000,0.000000,0.000000}%
\pgfsetstrokecolor{currentstroke}%
\pgfsetdash{}{0pt}%
\pgfpathmoveto{\pgfqpoint{0.650000in}{0.420000in}}%
\pgfpathlineto{\pgfqpoint{0.650000in}{2.414488in}}%
\pgfusepath{stroke}%
\end{pgfscope}%
\begin{pgfscope}%
\pgfsetrectcap%
\pgfsetmiterjoin%
\pgfsetlinewidth{0.803000pt}%
\definecolor{currentstroke}{rgb}{0.000000,0.000000,0.000000}%
\pgfsetstrokecolor{currentstroke}%
\pgfsetdash{}{0pt}%
\pgfpathmoveto{\pgfqpoint{3.038110in}{0.420000in}}%
\pgfpathlineto{\pgfqpoint{3.038110in}{2.414488in}}%
\pgfusepath{stroke}%
\end{pgfscope}%
\begin{pgfscope}%
\pgfsetrectcap%
\pgfsetmiterjoin%
\pgfsetlinewidth{0.803000pt}%
\definecolor{currentstroke}{rgb}{0.000000,0.000000,0.000000}%
\pgfsetstrokecolor{currentstroke}%
\pgfsetdash{}{0pt}%
\pgfpathmoveto{\pgfqpoint{0.650000in}{0.420000in}}%
\pgfpathlineto{\pgfqpoint{3.038110in}{0.420000in}}%
\pgfusepath{stroke}%
\end{pgfscope}%
\begin{pgfscope}%
\pgfsetrectcap%
\pgfsetmiterjoin%
\pgfsetlinewidth{0.803000pt}%
\definecolor{currentstroke}{rgb}{0.000000,0.000000,0.000000}%
\pgfsetstrokecolor{currentstroke}%
\pgfsetdash{}{0pt}%
\pgfpathmoveto{\pgfqpoint{0.650000in}{2.414488in}}%
\pgfpathlineto{\pgfqpoint{3.038110in}{2.414488in}}%
\pgfusepath{stroke}%
\end{pgfscope}%
\begin{pgfscope}%
\pgfsetrectcap%
\pgfsetroundjoin%
\pgfsetlinewidth{0.602250pt}%
\definecolor{currentstroke}{rgb}{0.000000,0.000000,0.000000}%
\pgfsetstrokecolor{currentstroke}%
\pgfsetdash{}{0pt}%
\pgfpathmoveto{\pgfqpoint{2.108975in}{2.275599in}}%
\pgfpathlineto{\pgfqpoint{2.220086in}{2.275599in}}%
\pgfpathlineto{\pgfqpoint{2.331197in}{2.275599in}}%
\pgfusepath{stroke}%
\end{pgfscope}%
\begin{pgfscope}%
\definecolor{textcolor}{rgb}{0.000000,0.000000,0.000000}%
\pgfsetstrokecolor{textcolor}%
\pgfsetfillcolor{textcolor}%
\pgftext[x=2.420086in,y=2.236710in,left,base]{\color{textcolor}{\rmfamily\fontsize{8.000000}{9.600000}\selectfont\catcode`\^=\active\def^{\ifmmode\sp\else\^{}\fi}\catcode`\%=\active\def%{\%}ref}}%
\end{pgfscope}%
\begin{pgfscope}%
\pgfsetrectcap%
\pgfsetroundjoin%
\pgfsetlinewidth{1.003750pt}%
\definecolor{currentstroke}{rgb}{0.870588,0.466667,0.682353}%
\pgfsetstrokecolor{currentstroke}%
\pgfsetdash{}{0pt}%
\pgfpathmoveto{\pgfqpoint{2.108975in}{2.120661in}}%
\pgfpathlineto{\pgfqpoint{2.220086in}{2.120661in}}%
\pgfpathlineto{\pgfqpoint{2.331197in}{2.120661in}}%
\pgfusepath{stroke}%
\end{pgfscope}%
\begin{pgfscope}%
\definecolor{textcolor}{rgb}{0.000000,0.000000,0.000000}%
\pgfsetstrokecolor{textcolor}%
\pgfsetfillcolor{textcolor}%
\pgftext[x=2.420086in,y=2.081772in,left,base]{\color{textcolor}{\rmfamily\fontsize{8.000000}{9.600000}\selectfont\catcode`\^=\active\def^{\ifmmode\sp\else\^{}\fi}\catcode`\%=\active\def%{\%}base}}%
\end{pgfscope}%
\begin{pgfscope}%
\pgfsetrectcap%
\pgfsetroundjoin%
\pgfsetlinewidth{1.003750pt}%
\definecolor{currentstroke}{rgb}{0.498039,0.737255,0.254902}%
\pgfsetstrokecolor{currentstroke}%
\pgfsetdash{}{0pt}%
\pgfpathmoveto{\pgfqpoint{2.108975in}{1.965723in}}%
\pgfpathlineto{\pgfqpoint{2.220086in}{1.965723in}}%
\pgfpathlineto{\pgfqpoint{2.331197in}{1.965723in}}%
\pgfusepath{stroke}%
\end{pgfscope}%
\begin{pgfscope}%
\definecolor{textcolor}{rgb}{0.000000,0.000000,0.000000}%
\pgfsetstrokecolor{textcolor}%
\pgfsetfillcolor{textcolor}%
\pgftext[x=2.420086in,y=1.926834in,left,base]{\color{textcolor}{\rmfamily\fontsize{8.000000}{9.600000}\selectfont\catcode`\^=\active\def^{\ifmmode\sp\else\^{}\fi}\catcode`\%=\active\def%{\%}optimised}}%
\end{pgfscope}%
\begin{pgfscope}%
\pgfsetrectcap%
\pgfsetroundjoin%
\pgfsetlinewidth{1.003750pt}%
\definecolor{currentstroke}{rgb}{0.400000,0.760784,0.643137}%
\pgfsetstrokecolor{currentstroke}%
\pgfsetdash{}{0pt}%
\pgfpathmoveto{\pgfqpoint{2.108975in}{1.810784in}}%
\pgfpathlineto{\pgfqpoint{2.220086in}{1.810784in}}%
\pgfpathlineto{\pgfqpoint{2.331197in}{1.810784in}}%
\pgfusepath{stroke}%
\end{pgfscope}%
\begin{pgfscope}%
\definecolor{textcolor}{rgb}{0.000000,0.000000,0.000000}%
\pgfsetstrokecolor{textcolor}%
\pgfsetfillcolor{textcolor}%
\pgftext[x=2.420086in,y=1.771895in,left,base]{\color{textcolor}{\rmfamily\fontsize{8.000000}{9.600000}\selectfont\catcode`\^=\active\def^{\ifmmode\sp\else\^{}\fi}\catcode`\%=\active\def%{\%}planeTBL}}%
\end{pgfscope}%
\end{pgfpicture}%
\makeatother%
\endgroup%

    \label{fig:inflowSecVerifBudget}
\end{subfigure}
\hfill
\begin{subfigure}[t]{0.495\textwidth}
    \centering
    \subcaption{}
    %% Creator: Matplotlib, PGF backend
%%
%% To include the figure in your LaTeX document, write
%%   \input{<filename>.pgf}
%%
%% Make sure the required packages are loaded in your preamble
%%   \usepackage{pgf}
%%
%% Also ensure that all the required font packages are loaded; for instance,
%% the lmodern package is sometimes necessary when using math font.
%%   \usepackage{lmodern}
%%
%% Figures using additional raster images can only be included by \input if
%% they are in the same directory as the main LaTeX file. For loading figures
%% from other directories you can use the `import` package
%%   \usepackage{import}
%%
%% and then include the figures with
%%   \import{<path to file>}{<filename>.pgf}
%%
%% Matplotlib used the following preamble
%%   \def\mathdefault#1{#1}
%%   \everymath=\expandafter{\the\everymath\displaystyle}
%%   \IfFileExists{scrextend.sty}{
%%     \usepackage[fontsize=11.000000pt]{scrextend}
%%   }{
%%     \renewcommand{\normalsize}{\fontsize{11.000000}{13.200000}\selectfont}
%%     \normalsize
%%   }
%%   
%%   \makeatletter\@ifpackageloaded{underscore}{}{\usepackage[strings]{underscore}}\makeatother
%%
\begingroup%
\makeatletter%
\begin{pgfpicture}%
\pgfpathrectangle{\pgfpointorigin}{\pgfqpoint{3.118110in}{2.494488in}}%
\pgfusepath{use as bounding box, clip}%
\begin{pgfscope}%
\pgfsetbuttcap%
\pgfsetmiterjoin%
\definecolor{currentfill}{rgb}{1.000000,1.000000,1.000000}%
\pgfsetfillcolor{currentfill}%
\pgfsetlinewidth{0.000000pt}%
\definecolor{currentstroke}{rgb}{1.000000,1.000000,1.000000}%
\pgfsetstrokecolor{currentstroke}%
\pgfsetdash{}{0pt}%
\pgfpathmoveto{\pgfqpoint{0.000000in}{0.000000in}}%
\pgfpathlineto{\pgfqpoint{3.118110in}{0.000000in}}%
\pgfpathlineto{\pgfqpoint{3.118110in}{2.494488in}}%
\pgfpathlineto{\pgfqpoint{0.000000in}{2.494488in}}%
\pgfpathlineto{\pgfqpoint{0.000000in}{0.000000in}}%
\pgfpathclose%
\pgfusepath{fill}%
\end{pgfscope}%
\begin{pgfscope}%
\pgfsetbuttcap%
\pgfsetmiterjoin%
\definecolor{currentfill}{rgb}{1.000000,1.000000,1.000000}%
\pgfsetfillcolor{currentfill}%
\pgfsetlinewidth{0.000000pt}%
\definecolor{currentstroke}{rgb}{0.000000,0.000000,0.000000}%
\pgfsetstrokecolor{currentstroke}%
\pgfsetstrokeopacity{0.000000}%
\pgfsetdash{}{0pt}%
\pgfpathmoveto{\pgfqpoint{0.650000in}{0.420000in}}%
\pgfpathlineto{\pgfqpoint{3.038110in}{0.420000in}}%
\pgfpathlineto{\pgfqpoint{3.038110in}{2.414488in}}%
\pgfpathlineto{\pgfqpoint{0.650000in}{2.414488in}}%
\pgfpathlineto{\pgfqpoint{0.650000in}{0.420000in}}%
\pgfpathclose%
\pgfusepath{fill}%
\end{pgfscope}%
\begin{pgfscope}%
\pgfsetbuttcap%
\pgfsetroundjoin%
\definecolor{currentfill}{rgb}{0.000000,0.000000,0.000000}%
\pgfsetfillcolor{currentfill}%
\pgfsetlinewidth{0.501875pt}%
\definecolor{currentstroke}{rgb}{0.000000,0.000000,0.000000}%
\pgfsetstrokecolor{currentstroke}%
\pgfsetdash{}{0pt}%
\pgfsys@defobject{currentmarker}{\pgfqpoint{0.000000in}{0.000000in}}{\pgfqpoint{0.000000in}{0.041667in}}{%
\pgfpathmoveto{\pgfqpoint{0.000000in}{0.000000in}}%
\pgfpathlineto{\pgfqpoint{0.000000in}{0.041667in}}%
\pgfusepath{stroke,fill}%
}%
\begin{pgfscope}%
\pgfsys@transformshift{0.650000in}{0.420000in}%
\pgfsys@useobject{currentmarker}{}%
\end{pgfscope}%
\end{pgfscope}%
\begin{pgfscope}%
\pgfsetbuttcap%
\pgfsetroundjoin%
\definecolor{currentfill}{rgb}{0.000000,0.000000,0.000000}%
\pgfsetfillcolor{currentfill}%
\pgfsetlinewidth{0.501875pt}%
\definecolor{currentstroke}{rgb}{0.000000,0.000000,0.000000}%
\pgfsetstrokecolor{currentstroke}%
\pgfsetdash{}{0pt}%
\pgfsys@defobject{currentmarker}{\pgfqpoint{0.000000in}{-0.041667in}}{\pgfqpoint{0.000000in}{0.000000in}}{%
\pgfpathmoveto{\pgfqpoint{0.000000in}{0.000000in}}%
\pgfpathlineto{\pgfqpoint{0.000000in}{-0.041667in}}%
\pgfusepath{stroke,fill}%
}%
\begin{pgfscope}%
\pgfsys@transformshift{0.650000in}{2.414488in}%
\pgfsys@useobject{currentmarker}{}%
\end{pgfscope}%
\end{pgfscope}%
\begin{pgfscope}%
\definecolor{textcolor}{rgb}{0.000000,0.000000,0.000000}%
\pgfsetstrokecolor{textcolor}%
\pgfsetfillcolor{textcolor}%
\pgftext[x=0.650000in,y=0.371389in,,top]{\color{textcolor}{\rmfamily\fontsize{11.000000}{13.200000}\selectfont\catcode`\^=\active\def^{\ifmmode\sp\else\^{}\fi}\catcode`\%=\active\def%{\%}0}}%
\end{pgfscope}%
\begin{pgfscope}%
\pgfsetbuttcap%
\pgfsetroundjoin%
\definecolor{currentfill}{rgb}{0.000000,0.000000,0.000000}%
\pgfsetfillcolor{currentfill}%
\pgfsetlinewidth{0.501875pt}%
\definecolor{currentstroke}{rgb}{0.000000,0.000000,0.000000}%
\pgfsetstrokecolor{currentstroke}%
\pgfsetdash{}{0pt}%
\pgfsys@defobject{currentmarker}{\pgfqpoint{0.000000in}{0.000000in}}{\pgfqpoint{0.000000in}{0.041667in}}{%
\pgfpathmoveto{\pgfqpoint{0.000000in}{0.000000in}}%
\pgfpathlineto{\pgfqpoint{0.000000in}{0.041667in}}%
\pgfusepath{stroke,fill}%
}%
\begin{pgfscope}%
\pgfsys@transformshift{1.551174in}{0.420000in}%
\pgfsys@useobject{currentmarker}{}%
\end{pgfscope}%
\end{pgfscope}%
\begin{pgfscope}%
\pgfsetbuttcap%
\pgfsetroundjoin%
\definecolor{currentfill}{rgb}{0.000000,0.000000,0.000000}%
\pgfsetfillcolor{currentfill}%
\pgfsetlinewidth{0.501875pt}%
\definecolor{currentstroke}{rgb}{0.000000,0.000000,0.000000}%
\pgfsetstrokecolor{currentstroke}%
\pgfsetdash{}{0pt}%
\pgfsys@defobject{currentmarker}{\pgfqpoint{0.000000in}{-0.041667in}}{\pgfqpoint{0.000000in}{0.000000in}}{%
\pgfpathmoveto{\pgfqpoint{0.000000in}{0.000000in}}%
\pgfpathlineto{\pgfqpoint{0.000000in}{-0.041667in}}%
\pgfusepath{stroke,fill}%
}%
\begin{pgfscope}%
\pgfsys@transformshift{1.551174in}{2.414488in}%
\pgfsys@useobject{currentmarker}{}%
\end{pgfscope}%
\end{pgfscope}%
\begin{pgfscope}%
\definecolor{textcolor}{rgb}{0.000000,0.000000,0.000000}%
\pgfsetstrokecolor{textcolor}%
\pgfsetfillcolor{textcolor}%
\pgftext[x=1.551174in,y=0.371389in,,top]{\color{textcolor}{\rmfamily\fontsize{11.000000}{13.200000}\selectfont\catcode`\^=\active\def^{\ifmmode\sp\else\^{}\fi}\catcode`\%=\active\def%{\%}20}}%
\end{pgfscope}%
\begin{pgfscope}%
\pgfsetbuttcap%
\pgfsetroundjoin%
\definecolor{currentfill}{rgb}{0.000000,0.000000,0.000000}%
\pgfsetfillcolor{currentfill}%
\pgfsetlinewidth{0.501875pt}%
\definecolor{currentstroke}{rgb}{0.000000,0.000000,0.000000}%
\pgfsetstrokecolor{currentstroke}%
\pgfsetdash{}{0pt}%
\pgfsys@defobject{currentmarker}{\pgfqpoint{0.000000in}{0.000000in}}{\pgfqpoint{0.000000in}{0.041667in}}{%
\pgfpathmoveto{\pgfqpoint{0.000000in}{0.000000in}}%
\pgfpathlineto{\pgfqpoint{0.000000in}{0.041667in}}%
\pgfusepath{stroke,fill}%
}%
\begin{pgfscope}%
\pgfsys@transformshift{2.452347in}{0.420000in}%
\pgfsys@useobject{currentmarker}{}%
\end{pgfscope}%
\end{pgfscope}%
\begin{pgfscope}%
\pgfsetbuttcap%
\pgfsetroundjoin%
\definecolor{currentfill}{rgb}{0.000000,0.000000,0.000000}%
\pgfsetfillcolor{currentfill}%
\pgfsetlinewidth{0.501875pt}%
\definecolor{currentstroke}{rgb}{0.000000,0.000000,0.000000}%
\pgfsetstrokecolor{currentstroke}%
\pgfsetdash{}{0pt}%
\pgfsys@defobject{currentmarker}{\pgfqpoint{0.000000in}{-0.041667in}}{\pgfqpoint{0.000000in}{0.000000in}}{%
\pgfpathmoveto{\pgfqpoint{0.000000in}{0.000000in}}%
\pgfpathlineto{\pgfqpoint{0.000000in}{-0.041667in}}%
\pgfusepath{stroke,fill}%
}%
\begin{pgfscope}%
\pgfsys@transformshift{2.452347in}{2.414488in}%
\pgfsys@useobject{currentmarker}{}%
\end{pgfscope}%
\end{pgfscope}%
\begin{pgfscope}%
\definecolor{textcolor}{rgb}{0.000000,0.000000,0.000000}%
\pgfsetstrokecolor{textcolor}%
\pgfsetfillcolor{textcolor}%
\pgftext[x=2.452347in,y=0.371389in,,top]{\color{textcolor}{\rmfamily\fontsize{11.000000}{13.200000}\selectfont\catcode`\^=\active\def^{\ifmmode\sp\else\^{}\fi}\catcode`\%=\active\def%{\%}40}}%
\end{pgfscope}%
\begin{pgfscope}%
\pgfsetbuttcap%
\pgfsetroundjoin%
\definecolor{currentfill}{rgb}{0.000000,0.000000,0.000000}%
\pgfsetfillcolor{currentfill}%
\pgfsetlinewidth{0.401500pt}%
\definecolor{currentstroke}{rgb}{0.000000,0.000000,0.000000}%
\pgfsetstrokecolor{currentstroke}%
\pgfsetdash{}{0pt}%
\pgfsys@defobject{currentmarker}{\pgfqpoint{0.000000in}{0.000000in}}{\pgfqpoint{0.000000in}{0.020833in}}{%
\pgfpathmoveto{\pgfqpoint{0.000000in}{0.000000in}}%
\pgfpathlineto{\pgfqpoint{0.000000in}{0.020833in}}%
\pgfusepath{stroke,fill}%
}%
\begin{pgfscope}%
\pgfsys@transformshift{0.875293in}{0.420000in}%
\pgfsys@useobject{currentmarker}{}%
\end{pgfscope}%
\end{pgfscope}%
\begin{pgfscope}%
\pgfsetbuttcap%
\pgfsetroundjoin%
\definecolor{currentfill}{rgb}{0.000000,0.000000,0.000000}%
\pgfsetfillcolor{currentfill}%
\pgfsetlinewidth{0.401500pt}%
\definecolor{currentstroke}{rgb}{0.000000,0.000000,0.000000}%
\pgfsetstrokecolor{currentstroke}%
\pgfsetdash{}{0pt}%
\pgfsys@defobject{currentmarker}{\pgfqpoint{0.000000in}{-0.020833in}}{\pgfqpoint{0.000000in}{0.000000in}}{%
\pgfpathmoveto{\pgfqpoint{0.000000in}{0.000000in}}%
\pgfpathlineto{\pgfqpoint{0.000000in}{-0.020833in}}%
\pgfusepath{stroke,fill}%
}%
\begin{pgfscope}%
\pgfsys@transformshift{0.875293in}{2.414488in}%
\pgfsys@useobject{currentmarker}{}%
\end{pgfscope}%
\end{pgfscope}%
\begin{pgfscope}%
\pgfsetbuttcap%
\pgfsetroundjoin%
\definecolor{currentfill}{rgb}{0.000000,0.000000,0.000000}%
\pgfsetfillcolor{currentfill}%
\pgfsetlinewidth{0.401500pt}%
\definecolor{currentstroke}{rgb}{0.000000,0.000000,0.000000}%
\pgfsetstrokecolor{currentstroke}%
\pgfsetdash{}{0pt}%
\pgfsys@defobject{currentmarker}{\pgfqpoint{0.000000in}{0.000000in}}{\pgfqpoint{0.000000in}{0.020833in}}{%
\pgfpathmoveto{\pgfqpoint{0.000000in}{0.000000in}}%
\pgfpathlineto{\pgfqpoint{0.000000in}{0.020833in}}%
\pgfusepath{stroke,fill}%
}%
\begin{pgfscope}%
\pgfsys@transformshift{1.100587in}{0.420000in}%
\pgfsys@useobject{currentmarker}{}%
\end{pgfscope}%
\end{pgfscope}%
\begin{pgfscope}%
\pgfsetbuttcap%
\pgfsetroundjoin%
\definecolor{currentfill}{rgb}{0.000000,0.000000,0.000000}%
\pgfsetfillcolor{currentfill}%
\pgfsetlinewidth{0.401500pt}%
\definecolor{currentstroke}{rgb}{0.000000,0.000000,0.000000}%
\pgfsetstrokecolor{currentstroke}%
\pgfsetdash{}{0pt}%
\pgfsys@defobject{currentmarker}{\pgfqpoint{0.000000in}{-0.020833in}}{\pgfqpoint{0.000000in}{0.000000in}}{%
\pgfpathmoveto{\pgfqpoint{0.000000in}{0.000000in}}%
\pgfpathlineto{\pgfqpoint{0.000000in}{-0.020833in}}%
\pgfusepath{stroke,fill}%
}%
\begin{pgfscope}%
\pgfsys@transformshift{1.100587in}{2.414488in}%
\pgfsys@useobject{currentmarker}{}%
\end{pgfscope}%
\end{pgfscope}%
\begin{pgfscope}%
\pgfsetbuttcap%
\pgfsetroundjoin%
\definecolor{currentfill}{rgb}{0.000000,0.000000,0.000000}%
\pgfsetfillcolor{currentfill}%
\pgfsetlinewidth{0.401500pt}%
\definecolor{currentstroke}{rgb}{0.000000,0.000000,0.000000}%
\pgfsetstrokecolor{currentstroke}%
\pgfsetdash{}{0pt}%
\pgfsys@defobject{currentmarker}{\pgfqpoint{0.000000in}{0.000000in}}{\pgfqpoint{0.000000in}{0.020833in}}{%
\pgfpathmoveto{\pgfqpoint{0.000000in}{0.000000in}}%
\pgfpathlineto{\pgfqpoint{0.000000in}{0.020833in}}%
\pgfusepath{stroke,fill}%
}%
\begin{pgfscope}%
\pgfsys@transformshift{1.325880in}{0.420000in}%
\pgfsys@useobject{currentmarker}{}%
\end{pgfscope}%
\end{pgfscope}%
\begin{pgfscope}%
\pgfsetbuttcap%
\pgfsetroundjoin%
\definecolor{currentfill}{rgb}{0.000000,0.000000,0.000000}%
\pgfsetfillcolor{currentfill}%
\pgfsetlinewidth{0.401500pt}%
\definecolor{currentstroke}{rgb}{0.000000,0.000000,0.000000}%
\pgfsetstrokecolor{currentstroke}%
\pgfsetdash{}{0pt}%
\pgfsys@defobject{currentmarker}{\pgfqpoint{0.000000in}{-0.020833in}}{\pgfqpoint{0.000000in}{0.000000in}}{%
\pgfpathmoveto{\pgfqpoint{0.000000in}{0.000000in}}%
\pgfpathlineto{\pgfqpoint{0.000000in}{-0.020833in}}%
\pgfusepath{stroke,fill}%
}%
\begin{pgfscope}%
\pgfsys@transformshift{1.325880in}{2.414488in}%
\pgfsys@useobject{currentmarker}{}%
\end{pgfscope}%
\end{pgfscope}%
\begin{pgfscope}%
\pgfsetbuttcap%
\pgfsetroundjoin%
\definecolor{currentfill}{rgb}{0.000000,0.000000,0.000000}%
\pgfsetfillcolor{currentfill}%
\pgfsetlinewidth{0.401500pt}%
\definecolor{currentstroke}{rgb}{0.000000,0.000000,0.000000}%
\pgfsetstrokecolor{currentstroke}%
\pgfsetdash{}{0pt}%
\pgfsys@defobject{currentmarker}{\pgfqpoint{0.000000in}{0.000000in}}{\pgfqpoint{0.000000in}{0.020833in}}{%
\pgfpathmoveto{\pgfqpoint{0.000000in}{0.000000in}}%
\pgfpathlineto{\pgfqpoint{0.000000in}{0.020833in}}%
\pgfusepath{stroke,fill}%
}%
\begin{pgfscope}%
\pgfsys@transformshift{1.776467in}{0.420000in}%
\pgfsys@useobject{currentmarker}{}%
\end{pgfscope}%
\end{pgfscope}%
\begin{pgfscope}%
\pgfsetbuttcap%
\pgfsetroundjoin%
\definecolor{currentfill}{rgb}{0.000000,0.000000,0.000000}%
\pgfsetfillcolor{currentfill}%
\pgfsetlinewidth{0.401500pt}%
\definecolor{currentstroke}{rgb}{0.000000,0.000000,0.000000}%
\pgfsetstrokecolor{currentstroke}%
\pgfsetdash{}{0pt}%
\pgfsys@defobject{currentmarker}{\pgfqpoint{0.000000in}{-0.020833in}}{\pgfqpoint{0.000000in}{0.000000in}}{%
\pgfpathmoveto{\pgfqpoint{0.000000in}{0.000000in}}%
\pgfpathlineto{\pgfqpoint{0.000000in}{-0.020833in}}%
\pgfusepath{stroke,fill}%
}%
\begin{pgfscope}%
\pgfsys@transformshift{1.776467in}{2.414488in}%
\pgfsys@useobject{currentmarker}{}%
\end{pgfscope}%
\end{pgfscope}%
\begin{pgfscope}%
\pgfsetbuttcap%
\pgfsetroundjoin%
\definecolor{currentfill}{rgb}{0.000000,0.000000,0.000000}%
\pgfsetfillcolor{currentfill}%
\pgfsetlinewidth{0.401500pt}%
\definecolor{currentstroke}{rgb}{0.000000,0.000000,0.000000}%
\pgfsetstrokecolor{currentstroke}%
\pgfsetdash{}{0pt}%
\pgfsys@defobject{currentmarker}{\pgfqpoint{0.000000in}{0.000000in}}{\pgfqpoint{0.000000in}{0.020833in}}{%
\pgfpathmoveto{\pgfqpoint{0.000000in}{0.000000in}}%
\pgfpathlineto{\pgfqpoint{0.000000in}{0.020833in}}%
\pgfusepath{stroke,fill}%
}%
\begin{pgfscope}%
\pgfsys@transformshift{2.001760in}{0.420000in}%
\pgfsys@useobject{currentmarker}{}%
\end{pgfscope}%
\end{pgfscope}%
\begin{pgfscope}%
\pgfsetbuttcap%
\pgfsetroundjoin%
\definecolor{currentfill}{rgb}{0.000000,0.000000,0.000000}%
\pgfsetfillcolor{currentfill}%
\pgfsetlinewidth{0.401500pt}%
\definecolor{currentstroke}{rgb}{0.000000,0.000000,0.000000}%
\pgfsetstrokecolor{currentstroke}%
\pgfsetdash{}{0pt}%
\pgfsys@defobject{currentmarker}{\pgfqpoint{0.000000in}{-0.020833in}}{\pgfqpoint{0.000000in}{0.000000in}}{%
\pgfpathmoveto{\pgfqpoint{0.000000in}{0.000000in}}%
\pgfpathlineto{\pgfqpoint{0.000000in}{-0.020833in}}%
\pgfusepath{stroke,fill}%
}%
\begin{pgfscope}%
\pgfsys@transformshift{2.001760in}{2.414488in}%
\pgfsys@useobject{currentmarker}{}%
\end{pgfscope}%
\end{pgfscope}%
\begin{pgfscope}%
\pgfsetbuttcap%
\pgfsetroundjoin%
\definecolor{currentfill}{rgb}{0.000000,0.000000,0.000000}%
\pgfsetfillcolor{currentfill}%
\pgfsetlinewidth{0.401500pt}%
\definecolor{currentstroke}{rgb}{0.000000,0.000000,0.000000}%
\pgfsetstrokecolor{currentstroke}%
\pgfsetdash{}{0pt}%
\pgfsys@defobject{currentmarker}{\pgfqpoint{0.000000in}{0.000000in}}{\pgfqpoint{0.000000in}{0.020833in}}{%
\pgfpathmoveto{\pgfqpoint{0.000000in}{0.000000in}}%
\pgfpathlineto{\pgfqpoint{0.000000in}{0.020833in}}%
\pgfusepath{stroke,fill}%
}%
\begin{pgfscope}%
\pgfsys@transformshift{2.227054in}{0.420000in}%
\pgfsys@useobject{currentmarker}{}%
\end{pgfscope}%
\end{pgfscope}%
\begin{pgfscope}%
\pgfsetbuttcap%
\pgfsetroundjoin%
\definecolor{currentfill}{rgb}{0.000000,0.000000,0.000000}%
\pgfsetfillcolor{currentfill}%
\pgfsetlinewidth{0.401500pt}%
\definecolor{currentstroke}{rgb}{0.000000,0.000000,0.000000}%
\pgfsetstrokecolor{currentstroke}%
\pgfsetdash{}{0pt}%
\pgfsys@defobject{currentmarker}{\pgfqpoint{0.000000in}{-0.020833in}}{\pgfqpoint{0.000000in}{0.000000in}}{%
\pgfpathmoveto{\pgfqpoint{0.000000in}{0.000000in}}%
\pgfpathlineto{\pgfqpoint{0.000000in}{-0.020833in}}%
\pgfusepath{stroke,fill}%
}%
\begin{pgfscope}%
\pgfsys@transformshift{2.227054in}{2.414488in}%
\pgfsys@useobject{currentmarker}{}%
\end{pgfscope}%
\end{pgfscope}%
\begin{pgfscope}%
\pgfsetbuttcap%
\pgfsetroundjoin%
\definecolor{currentfill}{rgb}{0.000000,0.000000,0.000000}%
\pgfsetfillcolor{currentfill}%
\pgfsetlinewidth{0.401500pt}%
\definecolor{currentstroke}{rgb}{0.000000,0.000000,0.000000}%
\pgfsetstrokecolor{currentstroke}%
\pgfsetdash{}{0pt}%
\pgfsys@defobject{currentmarker}{\pgfqpoint{0.000000in}{0.000000in}}{\pgfqpoint{0.000000in}{0.020833in}}{%
\pgfpathmoveto{\pgfqpoint{0.000000in}{0.000000in}}%
\pgfpathlineto{\pgfqpoint{0.000000in}{0.020833in}}%
\pgfusepath{stroke,fill}%
}%
\begin{pgfscope}%
\pgfsys@transformshift{2.677641in}{0.420000in}%
\pgfsys@useobject{currentmarker}{}%
\end{pgfscope}%
\end{pgfscope}%
\begin{pgfscope}%
\pgfsetbuttcap%
\pgfsetroundjoin%
\definecolor{currentfill}{rgb}{0.000000,0.000000,0.000000}%
\pgfsetfillcolor{currentfill}%
\pgfsetlinewidth{0.401500pt}%
\definecolor{currentstroke}{rgb}{0.000000,0.000000,0.000000}%
\pgfsetstrokecolor{currentstroke}%
\pgfsetdash{}{0pt}%
\pgfsys@defobject{currentmarker}{\pgfqpoint{0.000000in}{-0.020833in}}{\pgfqpoint{0.000000in}{0.000000in}}{%
\pgfpathmoveto{\pgfqpoint{0.000000in}{0.000000in}}%
\pgfpathlineto{\pgfqpoint{0.000000in}{-0.020833in}}%
\pgfusepath{stroke,fill}%
}%
\begin{pgfscope}%
\pgfsys@transformshift{2.677641in}{2.414488in}%
\pgfsys@useobject{currentmarker}{}%
\end{pgfscope}%
\end{pgfscope}%
\begin{pgfscope}%
\pgfsetbuttcap%
\pgfsetroundjoin%
\definecolor{currentfill}{rgb}{0.000000,0.000000,0.000000}%
\pgfsetfillcolor{currentfill}%
\pgfsetlinewidth{0.401500pt}%
\definecolor{currentstroke}{rgb}{0.000000,0.000000,0.000000}%
\pgfsetstrokecolor{currentstroke}%
\pgfsetdash{}{0pt}%
\pgfsys@defobject{currentmarker}{\pgfqpoint{0.000000in}{0.000000in}}{\pgfqpoint{0.000000in}{0.020833in}}{%
\pgfpathmoveto{\pgfqpoint{0.000000in}{0.000000in}}%
\pgfpathlineto{\pgfqpoint{0.000000in}{0.020833in}}%
\pgfusepath{stroke,fill}%
}%
\begin{pgfscope}%
\pgfsys@transformshift{2.902934in}{0.420000in}%
\pgfsys@useobject{currentmarker}{}%
\end{pgfscope}%
\end{pgfscope}%
\begin{pgfscope}%
\pgfsetbuttcap%
\pgfsetroundjoin%
\definecolor{currentfill}{rgb}{0.000000,0.000000,0.000000}%
\pgfsetfillcolor{currentfill}%
\pgfsetlinewidth{0.401500pt}%
\definecolor{currentstroke}{rgb}{0.000000,0.000000,0.000000}%
\pgfsetstrokecolor{currentstroke}%
\pgfsetdash{}{0pt}%
\pgfsys@defobject{currentmarker}{\pgfqpoint{0.000000in}{-0.020833in}}{\pgfqpoint{0.000000in}{0.000000in}}{%
\pgfpathmoveto{\pgfqpoint{0.000000in}{0.000000in}}%
\pgfpathlineto{\pgfqpoint{0.000000in}{-0.020833in}}%
\pgfusepath{stroke,fill}%
}%
\begin{pgfscope}%
\pgfsys@transformshift{2.902934in}{2.414488in}%
\pgfsys@useobject{currentmarker}{}%
\end{pgfscope}%
\end{pgfscope}%
\begin{pgfscope}%
\definecolor{textcolor}{rgb}{0.000000,0.000000,0.000000}%
\pgfsetstrokecolor{textcolor}%
\pgfsetfillcolor{textcolor}%
\pgftext[x=1.844055in,y=0.208426in,,top]{\color{textcolor}{\rmfamily\fontsize{12.000000}{14.400000}\selectfont\catcode`\^=\active\def^{\ifmmode\sp\else\^{}\fi}\catcode`\%=\active\def%{\%}$s/\delta_{0}$}}%
\end{pgfscope}%
\begin{pgfscope}%
\pgfsetbuttcap%
\pgfsetroundjoin%
\definecolor{currentfill}{rgb}{0.000000,0.000000,0.000000}%
\pgfsetfillcolor{currentfill}%
\pgfsetlinewidth{0.501875pt}%
\definecolor{currentstroke}{rgb}{0.000000,0.000000,0.000000}%
\pgfsetstrokecolor{currentstroke}%
\pgfsetdash{}{0pt}%
\pgfsys@defobject{currentmarker}{\pgfqpoint{0.000000in}{0.000000in}}{\pgfqpoint{0.041667in}{0.000000in}}{%
\pgfpathmoveto{\pgfqpoint{0.000000in}{0.000000in}}%
\pgfpathlineto{\pgfqpoint{0.041667in}{0.000000in}}%
\pgfusepath{stroke,fill}%
}%
\begin{pgfscope}%
\pgfsys@transformshift{0.650000in}{0.420000in}%
\pgfsys@useobject{currentmarker}{}%
\end{pgfscope}%
\end{pgfscope}%
\begin{pgfscope}%
\pgfsetbuttcap%
\pgfsetroundjoin%
\definecolor{currentfill}{rgb}{0.000000,0.000000,0.000000}%
\pgfsetfillcolor{currentfill}%
\pgfsetlinewidth{0.501875pt}%
\definecolor{currentstroke}{rgb}{0.000000,0.000000,0.000000}%
\pgfsetstrokecolor{currentstroke}%
\pgfsetdash{}{0pt}%
\pgfsys@defobject{currentmarker}{\pgfqpoint{-0.041667in}{0.000000in}}{\pgfqpoint{-0.000000in}{0.000000in}}{%
\pgfpathmoveto{\pgfqpoint{-0.000000in}{0.000000in}}%
\pgfpathlineto{\pgfqpoint{-0.041667in}{0.000000in}}%
\pgfusepath{stroke,fill}%
}%
\begin{pgfscope}%
\pgfsys@transformshift{3.038110in}{0.420000in}%
\pgfsys@useobject{currentmarker}{}%
\end{pgfscope}%
\end{pgfscope}%
\begin{pgfscope}%
\definecolor{textcolor}{rgb}{0.000000,0.000000,0.000000}%
\pgfsetstrokecolor{textcolor}%
\pgfsetfillcolor{textcolor}%
\pgftext[x=0.288772in, y=0.367193in, left, base]{\color{textcolor}{\rmfamily\fontsize{11.000000}{13.200000}\selectfont\catcode`\^=\active\def^{\ifmmode\sp\else\^{}\fi}\catcode`\%=\active\def%{\%}\ensuremath{-}1.0}}%
\end{pgfscope}%
\begin{pgfscope}%
\pgfsetbuttcap%
\pgfsetroundjoin%
\definecolor{currentfill}{rgb}{0.000000,0.000000,0.000000}%
\pgfsetfillcolor{currentfill}%
\pgfsetlinewidth{0.501875pt}%
\definecolor{currentstroke}{rgb}{0.000000,0.000000,0.000000}%
\pgfsetstrokecolor{currentstroke}%
\pgfsetdash{}{0pt}%
\pgfsys@defobject{currentmarker}{\pgfqpoint{0.000000in}{0.000000in}}{\pgfqpoint{0.041667in}{0.000000in}}{%
\pgfpathmoveto{\pgfqpoint{0.000000in}{0.000000in}}%
\pgfpathlineto{\pgfqpoint{0.041667in}{0.000000in}}%
\pgfusepath{stroke,fill}%
}%
\begin{pgfscope}%
\pgfsys@transformshift{0.650000in}{0.918622in}%
\pgfsys@useobject{currentmarker}{}%
\end{pgfscope}%
\end{pgfscope}%
\begin{pgfscope}%
\pgfsetbuttcap%
\pgfsetroundjoin%
\definecolor{currentfill}{rgb}{0.000000,0.000000,0.000000}%
\pgfsetfillcolor{currentfill}%
\pgfsetlinewidth{0.501875pt}%
\definecolor{currentstroke}{rgb}{0.000000,0.000000,0.000000}%
\pgfsetstrokecolor{currentstroke}%
\pgfsetdash{}{0pt}%
\pgfsys@defobject{currentmarker}{\pgfqpoint{-0.041667in}{0.000000in}}{\pgfqpoint{-0.000000in}{0.000000in}}{%
\pgfpathmoveto{\pgfqpoint{-0.000000in}{0.000000in}}%
\pgfpathlineto{\pgfqpoint{-0.041667in}{0.000000in}}%
\pgfusepath{stroke,fill}%
}%
\begin{pgfscope}%
\pgfsys@transformshift{3.038110in}{0.918622in}%
\pgfsys@useobject{currentmarker}{}%
\end{pgfscope}%
\end{pgfscope}%
\begin{pgfscope}%
\definecolor{textcolor}{rgb}{0.000000,0.000000,0.000000}%
\pgfsetstrokecolor{textcolor}%
\pgfsetfillcolor{textcolor}%
\pgftext[x=0.288772in, y=0.865815in, left, base]{\color{textcolor}{\rmfamily\fontsize{11.000000}{13.200000}\selectfont\catcode`\^=\active\def^{\ifmmode\sp\else\^{}\fi}\catcode`\%=\active\def%{\%}\ensuremath{-}0.5}}%
\end{pgfscope}%
\begin{pgfscope}%
\pgfsetbuttcap%
\pgfsetroundjoin%
\definecolor{currentfill}{rgb}{0.000000,0.000000,0.000000}%
\pgfsetfillcolor{currentfill}%
\pgfsetlinewidth{0.501875pt}%
\definecolor{currentstroke}{rgb}{0.000000,0.000000,0.000000}%
\pgfsetstrokecolor{currentstroke}%
\pgfsetdash{}{0pt}%
\pgfsys@defobject{currentmarker}{\pgfqpoint{0.000000in}{0.000000in}}{\pgfqpoint{0.041667in}{0.000000in}}{%
\pgfpathmoveto{\pgfqpoint{0.000000in}{0.000000in}}%
\pgfpathlineto{\pgfqpoint{0.041667in}{0.000000in}}%
\pgfusepath{stroke,fill}%
}%
\begin{pgfscope}%
\pgfsys@transformshift{0.650000in}{1.417244in}%
\pgfsys@useobject{currentmarker}{}%
\end{pgfscope}%
\end{pgfscope}%
\begin{pgfscope}%
\pgfsetbuttcap%
\pgfsetroundjoin%
\definecolor{currentfill}{rgb}{0.000000,0.000000,0.000000}%
\pgfsetfillcolor{currentfill}%
\pgfsetlinewidth{0.501875pt}%
\definecolor{currentstroke}{rgb}{0.000000,0.000000,0.000000}%
\pgfsetstrokecolor{currentstroke}%
\pgfsetdash{}{0pt}%
\pgfsys@defobject{currentmarker}{\pgfqpoint{-0.041667in}{0.000000in}}{\pgfqpoint{-0.000000in}{0.000000in}}{%
\pgfpathmoveto{\pgfqpoint{-0.000000in}{0.000000in}}%
\pgfpathlineto{\pgfqpoint{-0.041667in}{0.000000in}}%
\pgfusepath{stroke,fill}%
}%
\begin{pgfscope}%
\pgfsys@transformshift{3.038110in}{1.417244in}%
\pgfsys@useobject{currentmarker}{}%
\end{pgfscope}%
\end{pgfscope}%
\begin{pgfscope}%
\definecolor{textcolor}{rgb}{0.000000,0.000000,0.000000}%
\pgfsetstrokecolor{textcolor}%
\pgfsetfillcolor{textcolor}%
\pgftext[x=0.407060in, y=1.364437in, left, base]{\color{textcolor}{\rmfamily\fontsize{11.000000}{13.200000}\selectfont\catcode`\^=\active\def^{\ifmmode\sp\else\^{}\fi}\catcode`\%=\active\def%{\%}0.0}}%
\end{pgfscope}%
\begin{pgfscope}%
\pgfsetbuttcap%
\pgfsetroundjoin%
\definecolor{currentfill}{rgb}{0.000000,0.000000,0.000000}%
\pgfsetfillcolor{currentfill}%
\pgfsetlinewidth{0.501875pt}%
\definecolor{currentstroke}{rgb}{0.000000,0.000000,0.000000}%
\pgfsetstrokecolor{currentstroke}%
\pgfsetdash{}{0pt}%
\pgfsys@defobject{currentmarker}{\pgfqpoint{0.000000in}{0.000000in}}{\pgfqpoint{0.041667in}{0.000000in}}{%
\pgfpathmoveto{\pgfqpoint{0.000000in}{0.000000in}}%
\pgfpathlineto{\pgfqpoint{0.041667in}{0.000000in}}%
\pgfusepath{stroke,fill}%
}%
\begin{pgfscope}%
\pgfsys@transformshift{0.650000in}{1.915866in}%
\pgfsys@useobject{currentmarker}{}%
\end{pgfscope}%
\end{pgfscope}%
\begin{pgfscope}%
\pgfsetbuttcap%
\pgfsetroundjoin%
\definecolor{currentfill}{rgb}{0.000000,0.000000,0.000000}%
\pgfsetfillcolor{currentfill}%
\pgfsetlinewidth{0.501875pt}%
\definecolor{currentstroke}{rgb}{0.000000,0.000000,0.000000}%
\pgfsetstrokecolor{currentstroke}%
\pgfsetdash{}{0pt}%
\pgfsys@defobject{currentmarker}{\pgfqpoint{-0.041667in}{0.000000in}}{\pgfqpoint{-0.000000in}{0.000000in}}{%
\pgfpathmoveto{\pgfqpoint{-0.000000in}{0.000000in}}%
\pgfpathlineto{\pgfqpoint{-0.041667in}{0.000000in}}%
\pgfusepath{stroke,fill}%
}%
\begin{pgfscope}%
\pgfsys@transformshift{3.038110in}{1.915866in}%
\pgfsys@useobject{currentmarker}{}%
\end{pgfscope}%
\end{pgfscope}%
\begin{pgfscope}%
\definecolor{textcolor}{rgb}{0.000000,0.000000,0.000000}%
\pgfsetstrokecolor{textcolor}%
\pgfsetfillcolor{textcolor}%
\pgftext[x=0.407060in, y=1.863059in, left, base]{\color{textcolor}{\rmfamily\fontsize{11.000000}{13.200000}\selectfont\catcode`\^=\active\def^{\ifmmode\sp\else\^{}\fi}\catcode`\%=\active\def%{\%}0.5}}%
\end{pgfscope}%
\begin{pgfscope}%
\pgfsetbuttcap%
\pgfsetroundjoin%
\definecolor{currentfill}{rgb}{0.000000,0.000000,0.000000}%
\pgfsetfillcolor{currentfill}%
\pgfsetlinewidth{0.501875pt}%
\definecolor{currentstroke}{rgb}{0.000000,0.000000,0.000000}%
\pgfsetstrokecolor{currentstroke}%
\pgfsetdash{}{0pt}%
\pgfsys@defobject{currentmarker}{\pgfqpoint{0.000000in}{0.000000in}}{\pgfqpoint{0.041667in}{0.000000in}}{%
\pgfpathmoveto{\pgfqpoint{0.000000in}{0.000000in}}%
\pgfpathlineto{\pgfqpoint{0.041667in}{0.000000in}}%
\pgfusepath{stroke,fill}%
}%
\begin{pgfscope}%
\pgfsys@transformshift{0.650000in}{2.414488in}%
\pgfsys@useobject{currentmarker}{}%
\end{pgfscope}%
\end{pgfscope}%
\begin{pgfscope}%
\pgfsetbuttcap%
\pgfsetroundjoin%
\definecolor{currentfill}{rgb}{0.000000,0.000000,0.000000}%
\pgfsetfillcolor{currentfill}%
\pgfsetlinewidth{0.501875pt}%
\definecolor{currentstroke}{rgb}{0.000000,0.000000,0.000000}%
\pgfsetstrokecolor{currentstroke}%
\pgfsetdash{}{0pt}%
\pgfsys@defobject{currentmarker}{\pgfqpoint{-0.041667in}{0.000000in}}{\pgfqpoint{-0.000000in}{0.000000in}}{%
\pgfpathmoveto{\pgfqpoint{-0.000000in}{0.000000in}}%
\pgfpathlineto{\pgfqpoint{-0.041667in}{0.000000in}}%
\pgfusepath{stroke,fill}%
}%
\begin{pgfscope}%
\pgfsys@transformshift{3.038110in}{2.414488in}%
\pgfsys@useobject{currentmarker}{}%
\end{pgfscope}%
\end{pgfscope}%
\begin{pgfscope}%
\definecolor{textcolor}{rgb}{0.000000,0.000000,0.000000}%
\pgfsetstrokecolor{textcolor}%
\pgfsetfillcolor{textcolor}%
\pgftext[x=0.407060in, y=2.361681in, left, base]{\color{textcolor}{\rmfamily\fontsize{11.000000}{13.200000}\selectfont\catcode`\^=\active\def^{\ifmmode\sp\else\^{}\fi}\catcode`\%=\active\def%{\%}1.0}}%
\end{pgfscope}%
\begin{pgfscope}%
\pgfsetbuttcap%
\pgfsetroundjoin%
\definecolor{currentfill}{rgb}{0.000000,0.000000,0.000000}%
\pgfsetfillcolor{currentfill}%
\pgfsetlinewidth{0.401500pt}%
\definecolor{currentstroke}{rgb}{0.000000,0.000000,0.000000}%
\pgfsetstrokecolor{currentstroke}%
\pgfsetdash{}{0pt}%
\pgfsys@defobject{currentmarker}{\pgfqpoint{0.000000in}{0.000000in}}{\pgfqpoint{0.020833in}{0.000000in}}{%
\pgfpathmoveto{\pgfqpoint{0.000000in}{0.000000in}}%
\pgfpathlineto{\pgfqpoint{0.020833in}{0.000000in}}%
\pgfusepath{stroke,fill}%
}%
\begin{pgfscope}%
\pgfsys@transformshift{0.650000in}{0.519724in}%
\pgfsys@useobject{currentmarker}{}%
\end{pgfscope}%
\end{pgfscope}%
\begin{pgfscope}%
\pgfsetbuttcap%
\pgfsetroundjoin%
\definecolor{currentfill}{rgb}{0.000000,0.000000,0.000000}%
\pgfsetfillcolor{currentfill}%
\pgfsetlinewidth{0.401500pt}%
\definecolor{currentstroke}{rgb}{0.000000,0.000000,0.000000}%
\pgfsetstrokecolor{currentstroke}%
\pgfsetdash{}{0pt}%
\pgfsys@defobject{currentmarker}{\pgfqpoint{-0.020833in}{0.000000in}}{\pgfqpoint{-0.000000in}{0.000000in}}{%
\pgfpathmoveto{\pgfqpoint{-0.000000in}{0.000000in}}%
\pgfpathlineto{\pgfqpoint{-0.020833in}{0.000000in}}%
\pgfusepath{stroke,fill}%
}%
\begin{pgfscope}%
\pgfsys@transformshift{3.038110in}{0.519724in}%
\pgfsys@useobject{currentmarker}{}%
\end{pgfscope}%
\end{pgfscope}%
\begin{pgfscope}%
\pgfsetbuttcap%
\pgfsetroundjoin%
\definecolor{currentfill}{rgb}{0.000000,0.000000,0.000000}%
\pgfsetfillcolor{currentfill}%
\pgfsetlinewidth{0.401500pt}%
\definecolor{currentstroke}{rgb}{0.000000,0.000000,0.000000}%
\pgfsetstrokecolor{currentstroke}%
\pgfsetdash{}{0pt}%
\pgfsys@defobject{currentmarker}{\pgfqpoint{0.000000in}{0.000000in}}{\pgfqpoint{0.020833in}{0.000000in}}{%
\pgfpathmoveto{\pgfqpoint{0.000000in}{0.000000in}}%
\pgfpathlineto{\pgfqpoint{0.020833in}{0.000000in}}%
\pgfusepath{stroke,fill}%
}%
\begin{pgfscope}%
\pgfsys@transformshift{0.650000in}{0.619449in}%
\pgfsys@useobject{currentmarker}{}%
\end{pgfscope}%
\end{pgfscope}%
\begin{pgfscope}%
\pgfsetbuttcap%
\pgfsetroundjoin%
\definecolor{currentfill}{rgb}{0.000000,0.000000,0.000000}%
\pgfsetfillcolor{currentfill}%
\pgfsetlinewidth{0.401500pt}%
\definecolor{currentstroke}{rgb}{0.000000,0.000000,0.000000}%
\pgfsetstrokecolor{currentstroke}%
\pgfsetdash{}{0pt}%
\pgfsys@defobject{currentmarker}{\pgfqpoint{-0.020833in}{0.000000in}}{\pgfqpoint{-0.000000in}{0.000000in}}{%
\pgfpathmoveto{\pgfqpoint{-0.000000in}{0.000000in}}%
\pgfpathlineto{\pgfqpoint{-0.020833in}{0.000000in}}%
\pgfusepath{stroke,fill}%
}%
\begin{pgfscope}%
\pgfsys@transformshift{3.038110in}{0.619449in}%
\pgfsys@useobject{currentmarker}{}%
\end{pgfscope}%
\end{pgfscope}%
\begin{pgfscope}%
\pgfsetbuttcap%
\pgfsetroundjoin%
\definecolor{currentfill}{rgb}{0.000000,0.000000,0.000000}%
\pgfsetfillcolor{currentfill}%
\pgfsetlinewidth{0.401500pt}%
\definecolor{currentstroke}{rgb}{0.000000,0.000000,0.000000}%
\pgfsetstrokecolor{currentstroke}%
\pgfsetdash{}{0pt}%
\pgfsys@defobject{currentmarker}{\pgfqpoint{0.000000in}{0.000000in}}{\pgfqpoint{0.020833in}{0.000000in}}{%
\pgfpathmoveto{\pgfqpoint{0.000000in}{0.000000in}}%
\pgfpathlineto{\pgfqpoint{0.020833in}{0.000000in}}%
\pgfusepath{stroke,fill}%
}%
\begin{pgfscope}%
\pgfsys@transformshift{0.650000in}{0.719173in}%
\pgfsys@useobject{currentmarker}{}%
\end{pgfscope}%
\end{pgfscope}%
\begin{pgfscope}%
\pgfsetbuttcap%
\pgfsetroundjoin%
\definecolor{currentfill}{rgb}{0.000000,0.000000,0.000000}%
\pgfsetfillcolor{currentfill}%
\pgfsetlinewidth{0.401500pt}%
\definecolor{currentstroke}{rgb}{0.000000,0.000000,0.000000}%
\pgfsetstrokecolor{currentstroke}%
\pgfsetdash{}{0pt}%
\pgfsys@defobject{currentmarker}{\pgfqpoint{-0.020833in}{0.000000in}}{\pgfqpoint{-0.000000in}{0.000000in}}{%
\pgfpathmoveto{\pgfqpoint{-0.000000in}{0.000000in}}%
\pgfpathlineto{\pgfqpoint{-0.020833in}{0.000000in}}%
\pgfusepath{stroke,fill}%
}%
\begin{pgfscope}%
\pgfsys@transformshift{3.038110in}{0.719173in}%
\pgfsys@useobject{currentmarker}{}%
\end{pgfscope}%
\end{pgfscope}%
\begin{pgfscope}%
\pgfsetbuttcap%
\pgfsetroundjoin%
\definecolor{currentfill}{rgb}{0.000000,0.000000,0.000000}%
\pgfsetfillcolor{currentfill}%
\pgfsetlinewidth{0.401500pt}%
\definecolor{currentstroke}{rgb}{0.000000,0.000000,0.000000}%
\pgfsetstrokecolor{currentstroke}%
\pgfsetdash{}{0pt}%
\pgfsys@defobject{currentmarker}{\pgfqpoint{0.000000in}{0.000000in}}{\pgfqpoint{0.020833in}{0.000000in}}{%
\pgfpathmoveto{\pgfqpoint{0.000000in}{0.000000in}}%
\pgfpathlineto{\pgfqpoint{0.020833in}{0.000000in}}%
\pgfusepath{stroke,fill}%
}%
\begin{pgfscope}%
\pgfsys@transformshift{0.650000in}{0.818898in}%
\pgfsys@useobject{currentmarker}{}%
\end{pgfscope}%
\end{pgfscope}%
\begin{pgfscope}%
\pgfsetbuttcap%
\pgfsetroundjoin%
\definecolor{currentfill}{rgb}{0.000000,0.000000,0.000000}%
\pgfsetfillcolor{currentfill}%
\pgfsetlinewidth{0.401500pt}%
\definecolor{currentstroke}{rgb}{0.000000,0.000000,0.000000}%
\pgfsetstrokecolor{currentstroke}%
\pgfsetdash{}{0pt}%
\pgfsys@defobject{currentmarker}{\pgfqpoint{-0.020833in}{0.000000in}}{\pgfqpoint{-0.000000in}{0.000000in}}{%
\pgfpathmoveto{\pgfqpoint{-0.000000in}{0.000000in}}%
\pgfpathlineto{\pgfqpoint{-0.020833in}{0.000000in}}%
\pgfusepath{stroke,fill}%
}%
\begin{pgfscope}%
\pgfsys@transformshift{3.038110in}{0.818898in}%
\pgfsys@useobject{currentmarker}{}%
\end{pgfscope}%
\end{pgfscope}%
\begin{pgfscope}%
\pgfsetbuttcap%
\pgfsetroundjoin%
\definecolor{currentfill}{rgb}{0.000000,0.000000,0.000000}%
\pgfsetfillcolor{currentfill}%
\pgfsetlinewidth{0.401500pt}%
\definecolor{currentstroke}{rgb}{0.000000,0.000000,0.000000}%
\pgfsetstrokecolor{currentstroke}%
\pgfsetdash{}{0pt}%
\pgfsys@defobject{currentmarker}{\pgfqpoint{0.000000in}{0.000000in}}{\pgfqpoint{0.020833in}{0.000000in}}{%
\pgfpathmoveto{\pgfqpoint{0.000000in}{0.000000in}}%
\pgfpathlineto{\pgfqpoint{0.020833in}{0.000000in}}%
\pgfusepath{stroke,fill}%
}%
\begin{pgfscope}%
\pgfsys@transformshift{0.650000in}{1.018346in}%
\pgfsys@useobject{currentmarker}{}%
\end{pgfscope}%
\end{pgfscope}%
\begin{pgfscope}%
\pgfsetbuttcap%
\pgfsetroundjoin%
\definecolor{currentfill}{rgb}{0.000000,0.000000,0.000000}%
\pgfsetfillcolor{currentfill}%
\pgfsetlinewidth{0.401500pt}%
\definecolor{currentstroke}{rgb}{0.000000,0.000000,0.000000}%
\pgfsetstrokecolor{currentstroke}%
\pgfsetdash{}{0pt}%
\pgfsys@defobject{currentmarker}{\pgfqpoint{-0.020833in}{0.000000in}}{\pgfqpoint{-0.000000in}{0.000000in}}{%
\pgfpathmoveto{\pgfqpoint{-0.000000in}{0.000000in}}%
\pgfpathlineto{\pgfqpoint{-0.020833in}{0.000000in}}%
\pgfusepath{stroke,fill}%
}%
\begin{pgfscope}%
\pgfsys@transformshift{3.038110in}{1.018346in}%
\pgfsys@useobject{currentmarker}{}%
\end{pgfscope}%
\end{pgfscope}%
\begin{pgfscope}%
\pgfsetbuttcap%
\pgfsetroundjoin%
\definecolor{currentfill}{rgb}{0.000000,0.000000,0.000000}%
\pgfsetfillcolor{currentfill}%
\pgfsetlinewidth{0.401500pt}%
\definecolor{currentstroke}{rgb}{0.000000,0.000000,0.000000}%
\pgfsetstrokecolor{currentstroke}%
\pgfsetdash{}{0pt}%
\pgfsys@defobject{currentmarker}{\pgfqpoint{0.000000in}{0.000000in}}{\pgfqpoint{0.020833in}{0.000000in}}{%
\pgfpathmoveto{\pgfqpoint{0.000000in}{0.000000in}}%
\pgfpathlineto{\pgfqpoint{0.020833in}{0.000000in}}%
\pgfusepath{stroke,fill}%
}%
\begin{pgfscope}%
\pgfsys@transformshift{0.650000in}{1.118071in}%
\pgfsys@useobject{currentmarker}{}%
\end{pgfscope}%
\end{pgfscope}%
\begin{pgfscope}%
\pgfsetbuttcap%
\pgfsetroundjoin%
\definecolor{currentfill}{rgb}{0.000000,0.000000,0.000000}%
\pgfsetfillcolor{currentfill}%
\pgfsetlinewidth{0.401500pt}%
\definecolor{currentstroke}{rgb}{0.000000,0.000000,0.000000}%
\pgfsetstrokecolor{currentstroke}%
\pgfsetdash{}{0pt}%
\pgfsys@defobject{currentmarker}{\pgfqpoint{-0.020833in}{0.000000in}}{\pgfqpoint{-0.000000in}{0.000000in}}{%
\pgfpathmoveto{\pgfqpoint{-0.000000in}{0.000000in}}%
\pgfpathlineto{\pgfqpoint{-0.020833in}{0.000000in}}%
\pgfusepath{stroke,fill}%
}%
\begin{pgfscope}%
\pgfsys@transformshift{3.038110in}{1.118071in}%
\pgfsys@useobject{currentmarker}{}%
\end{pgfscope}%
\end{pgfscope}%
\begin{pgfscope}%
\pgfsetbuttcap%
\pgfsetroundjoin%
\definecolor{currentfill}{rgb}{0.000000,0.000000,0.000000}%
\pgfsetfillcolor{currentfill}%
\pgfsetlinewidth{0.401500pt}%
\definecolor{currentstroke}{rgb}{0.000000,0.000000,0.000000}%
\pgfsetstrokecolor{currentstroke}%
\pgfsetdash{}{0pt}%
\pgfsys@defobject{currentmarker}{\pgfqpoint{0.000000in}{0.000000in}}{\pgfqpoint{0.020833in}{0.000000in}}{%
\pgfpathmoveto{\pgfqpoint{0.000000in}{0.000000in}}%
\pgfpathlineto{\pgfqpoint{0.020833in}{0.000000in}}%
\pgfusepath{stroke,fill}%
}%
\begin{pgfscope}%
\pgfsys@transformshift{0.650000in}{1.217795in}%
\pgfsys@useobject{currentmarker}{}%
\end{pgfscope}%
\end{pgfscope}%
\begin{pgfscope}%
\pgfsetbuttcap%
\pgfsetroundjoin%
\definecolor{currentfill}{rgb}{0.000000,0.000000,0.000000}%
\pgfsetfillcolor{currentfill}%
\pgfsetlinewidth{0.401500pt}%
\definecolor{currentstroke}{rgb}{0.000000,0.000000,0.000000}%
\pgfsetstrokecolor{currentstroke}%
\pgfsetdash{}{0pt}%
\pgfsys@defobject{currentmarker}{\pgfqpoint{-0.020833in}{0.000000in}}{\pgfqpoint{-0.000000in}{0.000000in}}{%
\pgfpathmoveto{\pgfqpoint{-0.000000in}{0.000000in}}%
\pgfpathlineto{\pgfqpoint{-0.020833in}{0.000000in}}%
\pgfusepath{stroke,fill}%
}%
\begin{pgfscope}%
\pgfsys@transformshift{3.038110in}{1.217795in}%
\pgfsys@useobject{currentmarker}{}%
\end{pgfscope}%
\end{pgfscope}%
\begin{pgfscope}%
\pgfsetbuttcap%
\pgfsetroundjoin%
\definecolor{currentfill}{rgb}{0.000000,0.000000,0.000000}%
\pgfsetfillcolor{currentfill}%
\pgfsetlinewidth{0.401500pt}%
\definecolor{currentstroke}{rgb}{0.000000,0.000000,0.000000}%
\pgfsetstrokecolor{currentstroke}%
\pgfsetdash{}{0pt}%
\pgfsys@defobject{currentmarker}{\pgfqpoint{0.000000in}{0.000000in}}{\pgfqpoint{0.020833in}{0.000000in}}{%
\pgfpathmoveto{\pgfqpoint{0.000000in}{0.000000in}}%
\pgfpathlineto{\pgfqpoint{0.020833in}{0.000000in}}%
\pgfusepath{stroke,fill}%
}%
\begin{pgfscope}%
\pgfsys@transformshift{0.650000in}{1.317520in}%
\pgfsys@useobject{currentmarker}{}%
\end{pgfscope}%
\end{pgfscope}%
\begin{pgfscope}%
\pgfsetbuttcap%
\pgfsetroundjoin%
\definecolor{currentfill}{rgb}{0.000000,0.000000,0.000000}%
\pgfsetfillcolor{currentfill}%
\pgfsetlinewidth{0.401500pt}%
\definecolor{currentstroke}{rgb}{0.000000,0.000000,0.000000}%
\pgfsetstrokecolor{currentstroke}%
\pgfsetdash{}{0pt}%
\pgfsys@defobject{currentmarker}{\pgfqpoint{-0.020833in}{0.000000in}}{\pgfqpoint{-0.000000in}{0.000000in}}{%
\pgfpathmoveto{\pgfqpoint{-0.000000in}{0.000000in}}%
\pgfpathlineto{\pgfqpoint{-0.020833in}{0.000000in}}%
\pgfusepath{stroke,fill}%
}%
\begin{pgfscope}%
\pgfsys@transformshift{3.038110in}{1.317520in}%
\pgfsys@useobject{currentmarker}{}%
\end{pgfscope}%
\end{pgfscope}%
\begin{pgfscope}%
\pgfsetbuttcap%
\pgfsetroundjoin%
\definecolor{currentfill}{rgb}{0.000000,0.000000,0.000000}%
\pgfsetfillcolor{currentfill}%
\pgfsetlinewidth{0.401500pt}%
\definecolor{currentstroke}{rgb}{0.000000,0.000000,0.000000}%
\pgfsetstrokecolor{currentstroke}%
\pgfsetdash{}{0pt}%
\pgfsys@defobject{currentmarker}{\pgfqpoint{0.000000in}{0.000000in}}{\pgfqpoint{0.020833in}{0.000000in}}{%
\pgfpathmoveto{\pgfqpoint{0.000000in}{0.000000in}}%
\pgfpathlineto{\pgfqpoint{0.020833in}{0.000000in}}%
\pgfusepath{stroke,fill}%
}%
\begin{pgfscope}%
\pgfsys@transformshift{0.650000in}{1.516968in}%
\pgfsys@useobject{currentmarker}{}%
\end{pgfscope}%
\end{pgfscope}%
\begin{pgfscope}%
\pgfsetbuttcap%
\pgfsetroundjoin%
\definecolor{currentfill}{rgb}{0.000000,0.000000,0.000000}%
\pgfsetfillcolor{currentfill}%
\pgfsetlinewidth{0.401500pt}%
\definecolor{currentstroke}{rgb}{0.000000,0.000000,0.000000}%
\pgfsetstrokecolor{currentstroke}%
\pgfsetdash{}{0pt}%
\pgfsys@defobject{currentmarker}{\pgfqpoint{-0.020833in}{0.000000in}}{\pgfqpoint{-0.000000in}{0.000000in}}{%
\pgfpathmoveto{\pgfqpoint{-0.000000in}{0.000000in}}%
\pgfpathlineto{\pgfqpoint{-0.020833in}{0.000000in}}%
\pgfusepath{stroke,fill}%
}%
\begin{pgfscope}%
\pgfsys@transformshift{3.038110in}{1.516968in}%
\pgfsys@useobject{currentmarker}{}%
\end{pgfscope}%
\end{pgfscope}%
\begin{pgfscope}%
\pgfsetbuttcap%
\pgfsetroundjoin%
\definecolor{currentfill}{rgb}{0.000000,0.000000,0.000000}%
\pgfsetfillcolor{currentfill}%
\pgfsetlinewidth{0.401500pt}%
\definecolor{currentstroke}{rgb}{0.000000,0.000000,0.000000}%
\pgfsetstrokecolor{currentstroke}%
\pgfsetdash{}{0pt}%
\pgfsys@defobject{currentmarker}{\pgfqpoint{0.000000in}{0.000000in}}{\pgfqpoint{0.020833in}{0.000000in}}{%
\pgfpathmoveto{\pgfqpoint{0.000000in}{0.000000in}}%
\pgfpathlineto{\pgfqpoint{0.020833in}{0.000000in}}%
\pgfusepath{stroke,fill}%
}%
\begin{pgfscope}%
\pgfsys@transformshift{0.650000in}{1.616693in}%
\pgfsys@useobject{currentmarker}{}%
\end{pgfscope}%
\end{pgfscope}%
\begin{pgfscope}%
\pgfsetbuttcap%
\pgfsetroundjoin%
\definecolor{currentfill}{rgb}{0.000000,0.000000,0.000000}%
\pgfsetfillcolor{currentfill}%
\pgfsetlinewidth{0.401500pt}%
\definecolor{currentstroke}{rgb}{0.000000,0.000000,0.000000}%
\pgfsetstrokecolor{currentstroke}%
\pgfsetdash{}{0pt}%
\pgfsys@defobject{currentmarker}{\pgfqpoint{-0.020833in}{0.000000in}}{\pgfqpoint{-0.000000in}{0.000000in}}{%
\pgfpathmoveto{\pgfqpoint{-0.000000in}{0.000000in}}%
\pgfpathlineto{\pgfqpoint{-0.020833in}{0.000000in}}%
\pgfusepath{stroke,fill}%
}%
\begin{pgfscope}%
\pgfsys@transformshift{3.038110in}{1.616693in}%
\pgfsys@useobject{currentmarker}{}%
\end{pgfscope}%
\end{pgfscope}%
\begin{pgfscope}%
\pgfsetbuttcap%
\pgfsetroundjoin%
\definecolor{currentfill}{rgb}{0.000000,0.000000,0.000000}%
\pgfsetfillcolor{currentfill}%
\pgfsetlinewidth{0.401500pt}%
\definecolor{currentstroke}{rgb}{0.000000,0.000000,0.000000}%
\pgfsetstrokecolor{currentstroke}%
\pgfsetdash{}{0pt}%
\pgfsys@defobject{currentmarker}{\pgfqpoint{0.000000in}{0.000000in}}{\pgfqpoint{0.020833in}{0.000000in}}{%
\pgfpathmoveto{\pgfqpoint{0.000000in}{0.000000in}}%
\pgfpathlineto{\pgfqpoint{0.020833in}{0.000000in}}%
\pgfusepath{stroke,fill}%
}%
\begin{pgfscope}%
\pgfsys@transformshift{0.650000in}{1.716417in}%
\pgfsys@useobject{currentmarker}{}%
\end{pgfscope}%
\end{pgfscope}%
\begin{pgfscope}%
\pgfsetbuttcap%
\pgfsetroundjoin%
\definecolor{currentfill}{rgb}{0.000000,0.000000,0.000000}%
\pgfsetfillcolor{currentfill}%
\pgfsetlinewidth{0.401500pt}%
\definecolor{currentstroke}{rgb}{0.000000,0.000000,0.000000}%
\pgfsetstrokecolor{currentstroke}%
\pgfsetdash{}{0pt}%
\pgfsys@defobject{currentmarker}{\pgfqpoint{-0.020833in}{0.000000in}}{\pgfqpoint{-0.000000in}{0.000000in}}{%
\pgfpathmoveto{\pgfqpoint{-0.000000in}{0.000000in}}%
\pgfpathlineto{\pgfqpoint{-0.020833in}{0.000000in}}%
\pgfusepath{stroke,fill}%
}%
\begin{pgfscope}%
\pgfsys@transformshift{3.038110in}{1.716417in}%
\pgfsys@useobject{currentmarker}{}%
\end{pgfscope}%
\end{pgfscope}%
\begin{pgfscope}%
\pgfsetbuttcap%
\pgfsetroundjoin%
\definecolor{currentfill}{rgb}{0.000000,0.000000,0.000000}%
\pgfsetfillcolor{currentfill}%
\pgfsetlinewidth{0.401500pt}%
\definecolor{currentstroke}{rgb}{0.000000,0.000000,0.000000}%
\pgfsetstrokecolor{currentstroke}%
\pgfsetdash{}{0pt}%
\pgfsys@defobject{currentmarker}{\pgfqpoint{0.000000in}{0.000000in}}{\pgfqpoint{0.020833in}{0.000000in}}{%
\pgfpathmoveto{\pgfqpoint{0.000000in}{0.000000in}}%
\pgfpathlineto{\pgfqpoint{0.020833in}{0.000000in}}%
\pgfusepath{stroke,fill}%
}%
\begin{pgfscope}%
\pgfsys@transformshift{0.650000in}{1.816142in}%
\pgfsys@useobject{currentmarker}{}%
\end{pgfscope}%
\end{pgfscope}%
\begin{pgfscope}%
\pgfsetbuttcap%
\pgfsetroundjoin%
\definecolor{currentfill}{rgb}{0.000000,0.000000,0.000000}%
\pgfsetfillcolor{currentfill}%
\pgfsetlinewidth{0.401500pt}%
\definecolor{currentstroke}{rgb}{0.000000,0.000000,0.000000}%
\pgfsetstrokecolor{currentstroke}%
\pgfsetdash{}{0pt}%
\pgfsys@defobject{currentmarker}{\pgfqpoint{-0.020833in}{0.000000in}}{\pgfqpoint{-0.000000in}{0.000000in}}{%
\pgfpathmoveto{\pgfqpoint{-0.000000in}{0.000000in}}%
\pgfpathlineto{\pgfqpoint{-0.020833in}{0.000000in}}%
\pgfusepath{stroke,fill}%
}%
\begin{pgfscope}%
\pgfsys@transformshift{3.038110in}{1.816142in}%
\pgfsys@useobject{currentmarker}{}%
\end{pgfscope}%
\end{pgfscope}%
\begin{pgfscope}%
\pgfsetbuttcap%
\pgfsetroundjoin%
\definecolor{currentfill}{rgb}{0.000000,0.000000,0.000000}%
\pgfsetfillcolor{currentfill}%
\pgfsetlinewidth{0.401500pt}%
\definecolor{currentstroke}{rgb}{0.000000,0.000000,0.000000}%
\pgfsetstrokecolor{currentstroke}%
\pgfsetdash{}{0pt}%
\pgfsys@defobject{currentmarker}{\pgfqpoint{0.000000in}{0.000000in}}{\pgfqpoint{0.020833in}{0.000000in}}{%
\pgfpathmoveto{\pgfqpoint{0.000000in}{0.000000in}}%
\pgfpathlineto{\pgfqpoint{0.020833in}{0.000000in}}%
\pgfusepath{stroke,fill}%
}%
\begin{pgfscope}%
\pgfsys@transformshift{0.650000in}{2.015590in}%
\pgfsys@useobject{currentmarker}{}%
\end{pgfscope}%
\end{pgfscope}%
\begin{pgfscope}%
\pgfsetbuttcap%
\pgfsetroundjoin%
\definecolor{currentfill}{rgb}{0.000000,0.000000,0.000000}%
\pgfsetfillcolor{currentfill}%
\pgfsetlinewidth{0.401500pt}%
\definecolor{currentstroke}{rgb}{0.000000,0.000000,0.000000}%
\pgfsetstrokecolor{currentstroke}%
\pgfsetdash{}{0pt}%
\pgfsys@defobject{currentmarker}{\pgfqpoint{-0.020833in}{0.000000in}}{\pgfqpoint{-0.000000in}{0.000000in}}{%
\pgfpathmoveto{\pgfqpoint{-0.000000in}{0.000000in}}%
\pgfpathlineto{\pgfqpoint{-0.020833in}{0.000000in}}%
\pgfusepath{stroke,fill}%
}%
\begin{pgfscope}%
\pgfsys@transformshift{3.038110in}{2.015590in}%
\pgfsys@useobject{currentmarker}{}%
\end{pgfscope}%
\end{pgfscope}%
\begin{pgfscope}%
\pgfsetbuttcap%
\pgfsetroundjoin%
\definecolor{currentfill}{rgb}{0.000000,0.000000,0.000000}%
\pgfsetfillcolor{currentfill}%
\pgfsetlinewidth{0.401500pt}%
\definecolor{currentstroke}{rgb}{0.000000,0.000000,0.000000}%
\pgfsetstrokecolor{currentstroke}%
\pgfsetdash{}{0pt}%
\pgfsys@defobject{currentmarker}{\pgfqpoint{0.000000in}{0.000000in}}{\pgfqpoint{0.020833in}{0.000000in}}{%
\pgfpathmoveto{\pgfqpoint{0.000000in}{0.000000in}}%
\pgfpathlineto{\pgfqpoint{0.020833in}{0.000000in}}%
\pgfusepath{stroke,fill}%
}%
\begin{pgfscope}%
\pgfsys@transformshift{0.650000in}{2.115315in}%
\pgfsys@useobject{currentmarker}{}%
\end{pgfscope}%
\end{pgfscope}%
\begin{pgfscope}%
\pgfsetbuttcap%
\pgfsetroundjoin%
\definecolor{currentfill}{rgb}{0.000000,0.000000,0.000000}%
\pgfsetfillcolor{currentfill}%
\pgfsetlinewidth{0.401500pt}%
\definecolor{currentstroke}{rgb}{0.000000,0.000000,0.000000}%
\pgfsetstrokecolor{currentstroke}%
\pgfsetdash{}{0pt}%
\pgfsys@defobject{currentmarker}{\pgfqpoint{-0.020833in}{0.000000in}}{\pgfqpoint{-0.000000in}{0.000000in}}{%
\pgfpathmoveto{\pgfqpoint{-0.000000in}{0.000000in}}%
\pgfpathlineto{\pgfqpoint{-0.020833in}{0.000000in}}%
\pgfusepath{stroke,fill}%
}%
\begin{pgfscope}%
\pgfsys@transformshift{3.038110in}{2.115315in}%
\pgfsys@useobject{currentmarker}{}%
\end{pgfscope}%
\end{pgfscope}%
\begin{pgfscope}%
\pgfsetbuttcap%
\pgfsetroundjoin%
\definecolor{currentfill}{rgb}{0.000000,0.000000,0.000000}%
\pgfsetfillcolor{currentfill}%
\pgfsetlinewidth{0.401500pt}%
\definecolor{currentstroke}{rgb}{0.000000,0.000000,0.000000}%
\pgfsetstrokecolor{currentstroke}%
\pgfsetdash{}{0pt}%
\pgfsys@defobject{currentmarker}{\pgfqpoint{0.000000in}{0.000000in}}{\pgfqpoint{0.020833in}{0.000000in}}{%
\pgfpathmoveto{\pgfqpoint{0.000000in}{0.000000in}}%
\pgfpathlineto{\pgfqpoint{0.020833in}{0.000000in}}%
\pgfusepath{stroke,fill}%
}%
\begin{pgfscope}%
\pgfsys@transformshift{0.650000in}{2.215039in}%
\pgfsys@useobject{currentmarker}{}%
\end{pgfscope}%
\end{pgfscope}%
\begin{pgfscope}%
\pgfsetbuttcap%
\pgfsetroundjoin%
\definecolor{currentfill}{rgb}{0.000000,0.000000,0.000000}%
\pgfsetfillcolor{currentfill}%
\pgfsetlinewidth{0.401500pt}%
\definecolor{currentstroke}{rgb}{0.000000,0.000000,0.000000}%
\pgfsetstrokecolor{currentstroke}%
\pgfsetdash{}{0pt}%
\pgfsys@defobject{currentmarker}{\pgfqpoint{-0.020833in}{0.000000in}}{\pgfqpoint{-0.000000in}{0.000000in}}{%
\pgfpathmoveto{\pgfqpoint{-0.000000in}{0.000000in}}%
\pgfpathlineto{\pgfqpoint{-0.020833in}{0.000000in}}%
\pgfusepath{stroke,fill}%
}%
\begin{pgfscope}%
\pgfsys@transformshift{3.038110in}{2.215039in}%
\pgfsys@useobject{currentmarker}{}%
\end{pgfscope}%
\end{pgfscope}%
\begin{pgfscope}%
\pgfsetbuttcap%
\pgfsetroundjoin%
\definecolor{currentfill}{rgb}{0.000000,0.000000,0.000000}%
\pgfsetfillcolor{currentfill}%
\pgfsetlinewidth{0.401500pt}%
\definecolor{currentstroke}{rgb}{0.000000,0.000000,0.000000}%
\pgfsetstrokecolor{currentstroke}%
\pgfsetdash{}{0pt}%
\pgfsys@defobject{currentmarker}{\pgfqpoint{0.000000in}{0.000000in}}{\pgfqpoint{0.020833in}{0.000000in}}{%
\pgfpathmoveto{\pgfqpoint{0.000000in}{0.000000in}}%
\pgfpathlineto{\pgfqpoint{0.020833in}{0.000000in}}%
\pgfusepath{stroke,fill}%
}%
\begin{pgfscope}%
\pgfsys@transformshift{0.650000in}{2.314764in}%
\pgfsys@useobject{currentmarker}{}%
\end{pgfscope}%
\end{pgfscope}%
\begin{pgfscope}%
\pgfsetbuttcap%
\pgfsetroundjoin%
\definecolor{currentfill}{rgb}{0.000000,0.000000,0.000000}%
\pgfsetfillcolor{currentfill}%
\pgfsetlinewidth{0.401500pt}%
\definecolor{currentstroke}{rgb}{0.000000,0.000000,0.000000}%
\pgfsetstrokecolor{currentstroke}%
\pgfsetdash{}{0pt}%
\pgfsys@defobject{currentmarker}{\pgfqpoint{-0.020833in}{0.000000in}}{\pgfqpoint{-0.000000in}{0.000000in}}{%
\pgfpathmoveto{\pgfqpoint{-0.000000in}{0.000000in}}%
\pgfpathlineto{\pgfqpoint{-0.020833in}{0.000000in}}%
\pgfusepath{stroke,fill}%
}%
\begin{pgfscope}%
\pgfsys@transformshift{3.038110in}{2.314764in}%
\pgfsys@useobject{currentmarker}{}%
\end{pgfscope}%
\end{pgfscope}%
\begin{pgfscope}%
\definecolor{textcolor}{rgb}{0.000000,0.000000,0.000000}%
\pgfsetstrokecolor{textcolor}%
\pgfsetfillcolor{textcolor}%
\pgftext[x=0.260995in,y=1.417244in,,bottom,rotate=90.000000]{\color{textcolor}{\rmfamily\fontsize{12.000000}{14.400000}\selectfont\catcode`\^=\active\def^{\ifmmode\sp\else\^{}\fi}\catcode`\%=\active\def%{\%}$\beta$}}%
\end{pgfscope}%
\begin{pgfscope}%
\pgfpathrectangle{\pgfqpoint{0.650000in}{0.420000in}}{\pgfqpoint{2.388110in}{1.994488in}}%
\pgfusepath{clip}%
\pgfsetrectcap%
\pgfsetroundjoin%
\pgfsetlinewidth{1.003750pt}%
\definecolor{currentstroke}{rgb}{0.870588,0.466667,0.682353}%
\pgfsetstrokecolor{currentstroke}%
\pgfsetdash{}{0pt}%
\pgfpathmoveto{\pgfqpoint{0.650000in}{1.276262in}}%
\pgfpathlineto{\pgfqpoint{0.653599in}{1.275946in}}%
\pgfpathlineto{\pgfqpoint{0.665735in}{1.275243in}}%
\pgfpathlineto{\pgfqpoint{0.672065in}{1.275751in}}%
\pgfpathlineto{\pgfqpoint{0.676411in}{1.277373in}}%
\pgfpathlineto{\pgfqpoint{0.678603in}{1.278940in}}%
\pgfpathlineto{\pgfqpoint{0.680803in}{1.281249in}}%
\pgfpathlineto{\pgfqpoint{0.683008in}{1.284522in}}%
\pgfpathlineto{\pgfqpoint{0.685217in}{1.288975in}}%
\pgfpathlineto{\pgfqpoint{0.689641in}{1.302177in}}%
\pgfpathlineto{\pgfqpoint{0.694069in}{1.321614in}}%
\pgfpathlineto{\pgfqpoint{0.700713in}{1.359695in}}%
\pgfpathlineto{\pgfqpoint{0.707357in}{1.397694in}}%
\pgfpathlineto{\pgfqpoint{0.711787in}{1.416408in}}%
\pgfpathlineto{\pgfqpoint{0.714002in}{1.422869in}}%
\pgfpathlineto{\pgfqpoint{0.716216in}{1.427305in}}%
\pgfpathlineto{\pgfqpoint{0.718431in}{1.429779in}}%
\pgfpathlineto{\pgfqpoint{0.720646in}{1.430482in}}%
\pgfpathlineto{\pgfqpoint{0.722861in}{1.429689in}}%
\pgfpathlineto{\pgfqpoint{0.725076in}{1.427658in}}%
\pgfpathlineto{\pgfqpoint{0.729505in}{1.421246in}}%
\pgfpathlineto{\pgfqpoint{0.736150in}{1.410538in}}%
\pgfpathlineto{\pgfqpoint{0.740580in}{1.405799in}}%
\pgfpathlineto{\pgfqpoint{0.742794in}{1.404714in}}%
\pgfpathlineto{\pgfqpoint{0.745009in}{1.404614in}}%
\pgfpathlineto{\pgfqpoint{0.747224in}{1.405540in}}%
\pgfpathlineto{\pgfqpoint{0.749439in}{1.407475in}}%
\pgfpathlineto{\pgfqpoint{0.753868in}{1.414058in}}%
\pgfpathlineto{\pgfqpoint{0.758298in}{1.423433in}}%
\pgfpathlineto{\pgfqpoint{0.773802in}{1.460480in}}%
\pgfpathlineto{\pgfqpoint{0.778231in}{1.468248in}}%
\pgfpathlineto{\pgfqpoint{0.782661in}{1.473877in}}%
\pgfpathlineto{\pgfqpoint{0.787091in}{1.477394in}}%
\pgfpathlineto{\pgfqpoint{0.791520in}{1.478981in}}%
\pgfpathlineto{\pgfqpoint{0.795950in}{1.478942in}}%
\pgfpathlineto{\pgfqpoint{0.800380in}{1.477610in}}%
\pgfpathlineto{\pgfqpoint{0.807024in}{1.473833in}}%
\pgfpathlineto{\pgfqpoint{0.813669in}{1.468651in}}%
\pgfpathlineto{\pgfqpoint{0.824743in}{1.458176in}}%
\pgfpathlineto{\pgfqpoint{0.846891in}{1.435944in}}%
\pgfpathlineto{\pgfqpoint{0.853535in}{1.430752in}}%
\pgfpathlineto{\pgfqpoint{0.862395in}{1.425402in}}%
\pgfpathlineto{\pgfqpoint{0.875684in}{1.419204in}}%
\pgfpathlineto{\pgfqpoint{0.906691in}{1.405743in}}%
\pgfpathlineto{\pgfqpoint{0.915551in}{1.403331in}}%
\pgfpathlineto{\pgfqpoint{0.926625in}{1.401645in}}%
\pgfpathlineto{\pgfqpoint{0.959847in}{1.397821in}}%
\pgfpathlineto{\pgfqpoint{0.979780in}{1.395702in}}%
\pgfpathlineto{\pgfqpoint{1.001929in}{1.394559in}}%
\pgfpathlineto{\pgfqpoint{1.017432in}{1.394733in}}%
\pgfpathlineto{\pgfqpoint{1.035151in}{1.396175in}}%
\pgfpathlineto{\pgfqpoint{1.052870in}{1.397468in}}%
\pgfpathlineto{\pgfqpoint{1.079447in}{1.397688in}}%
\pgfpathlineto{\pgfqpoint{1.106025in}{1.397391in}}%
\pgfpathlineto{\pgfqpoint{1.123744in}{1.397276in}}%
\pgfpathlineto{\pgfqpoint{1.134818in}{1.398475in}}%
\pgfpathlineto{\pgfqpoint{1.161396in}{1.402280in}}%
\pgfpathlineto{\pgfqpoint{1.172470in}{1.402190in}}%
\pgfpathlineto{\pgfqpoint{1.199048in}{1.401091in}}%
\pgfpathlineto{\pgfqpoint{1.218981in}{1.402627in}}%
\pgfpathlineto{\pgfqpoint{1.234485in}{1.403181in}}%
\pgfpathlineto{\pgfqpoint{1.245559in}{1.402450in}}%
\pgfpathlineto{\pgfqpoint{1.267707in}{1.400316in}}%
\pgfpathlineto{\pgfqpoint{1.278781in}{1.400682in}}%
\pgfpathlineto{\pgfqpoint{1.298715in}{1.401654in}}%
\pgfpathlineto{\pgfqpoint{1.312004in}{1.400994in}}%
\pgfpathlineto{\pgfqpoint{1.331937in}{1.399956in}}%
\pgfpathlineto{\pgfqpoint{1.396167in}{1.399980in}}%
\pgfpathlineto{\pgfqpoint{1.467041in}{1.400077in}}%
\pgfpathlineto{\pgfqpoint{1.491404in}{1.399164in}}%
\pgfpathlineto{\pgfqpoint{1.504693in}{1.400502in}}%
\pgfpathlineto{\pgfqpoint{1.517982in}{1.401719in}}%
\pgfpathlineto{\pgfqpoint{1.529056in}{1.401453in}}%
\pgfpathlineto{\pgfqpoint{1.546775in}{1.399501in}}%
\pgfpathlineto{\pgfqpoint{1.582212in}{1.395260in}}%
\pgfpathlineto{\pgfqpoint{1.591071in}{1.395578in}}%
\pgfpathlineto{\pgfqpoint{1.602146in}{1.397341in}}%
\pgfpathlineto{\pgfqpoint{1.622079in}{1.401000in}}%
\pgfpathlineto{\pgfqpoint{1.630938in}{1.401416in}}%
\pgfpathlineto{\pgfqpoint{1.644227in}{1.400749in}}%
\pgfpathlineto{\pgfqpoint{1.708457in}{1.395836in}}%
\pgfpathlineto{\pgfqpoint{1.721746in}{1.396582in}}%
\pgfpathlineto{\pgfqpoint{1.743894in}{1.398177in}}%
\pgfpathlineto{\pgfqpoint{1.757183in}{1.397779in}}%
\pgfpathlineto{\pgfqpoint{1.779331in}{1.396776in}}%
\pgfpathlineto{\pgfqpoint{1.792620in}{1.397657in}}%
\pgfpathlineto{\pgfqpoint{1.810339in}{1.399147in}}%
\pgfpathlineto{\pgfqpoint{1.819198in}{1.398795in}}%
\pgfpathlineto{\pgfqpoint{1.828057in}{1.397437in}}%
\pgfpathlineto{\pgfqpoint{1.841346in}{1.394082in}}%
\pgfpathlineto{\pgfqpoint{1.852421in}{1.391446in}}%
\pgfpathlineto{\pgfqpoint{1.859065in}{1.390708in}}%
\pgfpathlineto{\pgfqpoint{1.865709in}{1.390927in}}%
\pgfpathlineto{\pgfqpoint{1.874569in}{1.392539in}}%
\pgfpathlineto{\pgfqpoint{1.890073in}{1.395875in}}%
\pgfpathlineto{\pgfqpoint{1.898932in}{1.396237in}}%
\pgfpathlineto{\pgfqpoint{1.907791in}{1.395343in}}%
\pgfpathlineto{\pgfqpoint{1.925510in}{1.393079in}}%
\pgfpathlineto{\pgfqpoint{1.934369in}{1.393130in}}%
\pgfpathlineto{\pgfqpoint{1.945443in}{1.394424in}}%
\pgfpathlineto{\pgfqpoint{1.969806in}{1.397912in}}%
\pgfpathlineto{\pgfqpoint{1.989740in}{1.398963in}}%
\pgfpathlineto{\pgfqpoint{2.009673in}{1.399285in}}%
\pgfpathlineto{\pgfqpoint{2.018532in}{1.398370in}}%
\pgfpathlineto{\pgfqpoint{2.029606in}{1.395991in}}%
\pgfpathlineto{\pgfqpoint{2.042895in}{1.393235in}}%
\pgfpathlineto{\pgfqpoint{2.051755in}{1.392616in}}%
\pgfpathlineto{\pgfqpoint{2.078333in}{1.391852in}}%
\pgfpathlineto{\pgfqpoint{2.089407in}{1.391450in}}%
\pgfpathlineto{\pgfqpoint{2.100481in}{1.392320in}}%
\pgfpathlineto{\pgfqpoint{2.124844in}{1.394954in}}%
\pgfpathlineto{\pgfqpoint{2.138133in}{1.395114in}}%
\pgfpathlineto{\pgfqpoint{2.149207in}{1.393839in}}%
\pgfpathlineto{\pgfqpoint{2.173570in}{1.390028in}}%
\pgfpathlineto{\pgfqpoint{2.191288in}{1.388840in}}%
\pgfpathlineto{\pgfqpoint{2.202363in}{1.389070in}}%
\pgfpathlineto{\pgfqpoint{2.211222in}{1.390396in}}%
\pgfpathlineto{\pgfqpoint{2.235585in}{1.394950in}}%
\pgfpathlineto{\pgfqpoint{2.248874in}{1.395578in}}%
\pgfpathlineto{\pgfqpoint{2.264378in}{1.395277in}}%
\pgfpathlineto{\pgfqpoint{2.277666in}{1.393915in}}%
\pgfpathlineto{\pgfqpoint{2.310889in}{1.389646in}}%
\pgfpathlineto{\pgfqpoint{2.328607in}{1.389036in}}%
\pgfpathlineto{\pgfqpoint{2.350756in}{1.389127in}}%
\pgfpathlineto{\pgfqpoint{2.364045in}{1.390659in}}%
\pgfpathlineto{\pgfqpoint{2.379548in}{1.392684in}}%
\pgfpathlineto{\pgfqpoint{2.392837in}{1.393171in}}%
\pgfpathlineto{\pgfqpoint{2.408341in}{1.392737in}}%
\pgfpathlineto{\pgfqpoint{2.423845in}{1.390950in}}%
\pgfpathlineto{\pgfqpoint{2.437134in}{1.389618in}}%
\pgfpathlineto{\pgfqpoint{2.448208in}{1.389702in}}%
\pgfpathlineto{\pgfqpoint{2.488075in}{1.391548in}}%
\pgfpathlineto{\pgfqpoint{2.516867in}{1.389396in}}%
\pgfpathlineto{\pgfqpoint{2.527942in}{1.390729in}}%
\pgfpathlineto{\pgfqpoint{2.541230in}{1.392344in}}%
\pgfpathlineto{\pgfqpoint{2.547875in}{1.392088in}}%
\pgfpathlineto{\pgfqpoint{2.554519in}{1.390817in}}%
\pgfpathlineto{\pgfqpoint{2.576668in}{1.384955in}}%
\pgfpathlineto{\pgfqpoint{2.583312in}{1.385205in}}%
\pgfpathlineto{\pgfqpoint{2.592171in}{1.387142in}}%
\pgfpathlineto{\pgfqpoint{2.601031in}{1.389028in}}%
\pgfpathlineto{\pgfqpoint{2.607675in}{1.389217in}}%
\pgfpathlineto{\pgfqpoint{2.614320in}{1.388172in}}%
\pgfpathlineto{\pgfqpoint{2.629823in}{1.384823in}}%
\pgfpathlineto{\pgfqpoint{2.636468in}{1.384661in}}%
\pgfpathlineto{\pgfqpoint{2.647542in}{1.385966in}}%
\pgfpathlineto{\pgfqpoint{2.660831in}{1.387502in}}%
\pgfpathlineto{\pgfqpoint{2.694053in}{1.390076in}}%
\pgfpathlineto{\pgfqpoint{2.707342in}{1.391443in}}%
\pgfpathlineto{\pgfqpoint{2.716201in}{1.391099in}}%
\pgfpathlineto{\pgfqpoint{2.722846in}{1.389863in}}%
\pgfpathlineto{\pgfqpoint{2.731705in}{1.387107in}}%
\pgfpathlineto{\pgfqpoint{2.749424in}{1.381120in}}%
\pgfpathlineto{\pgfqpoint{2.756068in}{1.379933in}}%
\pgfpathlineto{\pgfqpoint{2.764927in}{1.379633in}}%
\pgfpathlineto{\pgfqpoint{2.780431in}{1.380848in}}%
\pgfpathlineto{\pgfqpoint{2.798150in}{1.381688in}}%
\pgfpathlineto{\pgfqpoint{2.809224in}{1.381282in}}%
\pgfpathlineto{\pgfqpoint{2.831372in}{1.379770in}}%
\pgfpathlineto{\pgfqpoint{2.840231in}{1.380829in}}%
\pgfpathlineto{\pgfqpoint{2.853521in}{1.382790in}}%
\pgfpathlineto{\pgfqpoint{2.860165in}{1.382859in}}%
\pgfpathlineto{\pgfqpoint{2.869024in}{1.381722in}}%
\pgfpathlineto{\pgfqpoint{2.888958in}{1.378342in}}%
\pgfpathlineto{\pgfqpoint{2.897817in}{1.378140in}}%
\pgfpathlineto{\pgfqpoint{2.922180in}{1.378259in}}%
\pgfpathlineto{\pgfqpoint{2.931039in}{1.376870in}}%
\pgfpathlineto{\pgfqpoint{2.939899in}{1.374439in}}%
\pgfpathlineto{\pgfqpoint{2.966476in}{1.366093in}}%
\pgfpathlineto{\pgfqpoint{2.973121in}{1.365616in}}%
\pgfpathlineto{\pgfqpoint{2.995269in}{1.365765in}}%
\pgfpathlineto{\pgfqpoint{3.017417in}{1.363059in}}%
\pgfpathlineto{\pgfqpoint{3.039566in}{1.363366in}}%
\pgfpathlineto{\pgfqpoint{3.039566in}{1.363366in}}%
\pgfusepath{stroke}%
\end{pgfscope}%
\begin{pgfscope}%
\pgfpathrectangle{\pgfqpoint{0.650000in}{0.420000in}}{\pgfqpoint{2.388110in}{1.994488in}}%
\pgfusepath{clip}%
\pgfsetrectcap%
\pgfsetroundjoin%
\pgfsetlinewidth{1.003750pt}%
\definecolor{currentstroke}{rgb}{0.498039,0.737255,0.254902}%
\pgfsetstrokecolor{currentstroke}%
\pgfsetdash{}{0pt}%
\pgfpathmoveto{\pgfqpoint{0.650000in}{1.276442in}}%
\pgfpathlineto{\pgfqpoint{0.658147in}{1.275628in}}%
\pgfpathlineto{\pgfqpoint{0.670016in}{1.275065in}}%
\pgfpathlineto{\pgfqpoint{0.674345in}{1.275825in}}%
\pgfpathlineto{\pgfqpoint{0.678739in}{1.278220in}}%
\pgfpathlineto{\pgfqpoint{0.680949in}{1.280446in}}%
\pgfpathlineto{\pgfqpoint{0.683165in}{1.283640in}}%
\pgfpathlineto{\pgfqpoint{0.685384in}{1.288025in}}%
\pgfpathlineto{\pgfqpoint{0.689829in}{1.301128in}}%
\pgfpathlineto{\pgfqpoint{0.694278in}{1.320520in}}%
\pgfpathlineto{\pgfqpoint{0.700953in}{1.358539in}}%
\pgfpathlineto{\pgfqpoint{0.707629in}{1.396357in}}%
\pgfpathlineto{\pgfqpoint{0.712080in}{1.414852in}}%
\pgfpathlineto{\pgfqpoint{0.714305in}{1.421188in}}%
\pgfpathlineto{\pgfqpoint{0.716530in}{1.425481in}}%
\pgfpathlineto{\pgfqpoint{0.718755in}{1.427821in}}%
\pgfpathlineto{\pgfqpoint{0.720981in}{1.428395in}}%
\pgfpathlineto{\pgfqpoint{0.723206in}{1.427477in}}%
\pgfpathlineto{\pgfqpoint{0.725431in}{1.425342in}}%
\pgfpathlineto{\pgfqpoint{0.729882in}{1.418801in}}%
\pgfpathlineto{\pgfqpoint{0.736558in}{1.408140in}}%
\pgfpathlineto{\pgfqpoint{0.741009in}{1.403623in}}%
\pgfpathlineto{\pgfqpoint{0.743234in}{1.402706in}}%
\pgfpathlineto{\pgfqpoint{0.745459in}{1.402814in}}%
\pgfpathlineto{\pgfqpoint{0.747685in}{1.403980in}}%
\pgfpathlineto{\pgfqpoint{0.749910in}{1.406180in}}%
\pgfpathlineto{\pgfqpoint{0.754361in}{1.413359in}}%
\pgfpathlineto{\pgfqpoint{0.758811in}{1.423344in}}%
\pgfpathlineto{\pgfqpoint{0.774388in}{1.461713in}}%
\pgfpathlineto{\pgfqpoint{0.778839in}{1.469534in}}%
\pgfpathlineto{\pgfqpoint{0.783290in}{1.475105in}}%
\pgfpathlineto{\pgfqpoint{0.787740in}{1.478470in}}%
\pgfpathlineto{\pgfqpoint{0.792191in}{1.479859in}}%
\pgfpathlineto{\pgfqpoint{0.796642in}{1.479580in}}%
\pgfpathlineto{\pgfqpoint{0.801092in}{1.478015in}}%
\pgfpathlineto{\pgfqpoint{0.807768in}{1.473945in}}%
\pgfpathlineto{\pgfqpoint{0.816669in}{1.466687in}}%
\pgfpathlineto{\pgfqpoint{0.834472in}{1.450065in}}%
\pgfpathlineto{\pgfqpoint{0.854500in}{1.431792in}}%
\pgfpathlineto{\pgfqpoint{0.865626in}{1.422938in}}%
\pgfpathlineto{\pgfqpoint{0.876753in}{1.415600in}}%
\pgfpathlineto{\pgfqpoint{0.885654in}{1.410952in}}%
\pgfpathlineto{\pgfqpoint{0.896781in}{1.406552in}}%
\pgfpathlineto{\pgfqpoint{0.910133in}{1.402778in}}%
\pgfpathlineto{\pgfqpoint{0.930160in}{1.398540in}}%
\pgfpathlineto{\pgfqpoint{0.943512in}{1.396586in}}%
\pgfpathlineto{\pgfqpoint{0.956864in}{1.395825in}}%
\pgfpathlineto{\pgfqpoint{0.999145in}{1.395435in}}%
\pgfpathlineto{\pgfqpoint{1.025849in}{1.395849in}}%
\pgfpathlineto{\pgfqpoint{1.045877in}{1.395628in}}%
\pgfpathlineto{\pgfqpoint{1.088158in}{1.397408in}}%
\pgfpathlineto{\pgfqpoint{1.112636in}{1.397336in}}%
\pgfpathlineto{\pgfqpoint{1.130439in}{1.398490in}}%
\pgfpathlineto{\pgfqpoint{1.177171in}{1.402575in}}%
\pgfpathlineto{\pgfqpoint{1.192748in}{1.402431in}}%
\pgfpathlineto{\pgfqpoint{1.208325in}{1.401006in}}%
\pgfpathlineto{\pgfqpoint{1.226128in}{1.399310in}}%
\pgfpathlineto{\pgfqpoint{1.241705in}{1.399045in}}%
\pgfpathlineto{\pgfqpoint{1.326267in}{1.398679in}}%
\pgfpathlineto{\pgfqpoint{1.397477in}{1.396730in}}%
\pgfpathlineto{\pgfqpoint{1.450884in}{1.393610in}}%
\pgfpathlineto{\pgfqpoint{1.493165in}{1.389680in}}%
\pgfpathlineto{\pgfqpoint{1.506517in}{1.389465in}}%
\pgfpathlineto{\pgfqpoint{1.519869in}{1.390370in}}%
\pgfpathlineto{\pgfqpoint{1.542122in}{1.392163in}}%
\pgfpathlineto{\pgfqpoint{1.553249in}{1.391815in}}%
\pgfpathlineto{\pgfqpoint{1.591079in}{1.389427in}}%
\pgfpathlineto{\pgfqpoint{1.606657in}{1.388738in}}%
\pgfpathlineto{\pgfqpoint{1.631135in}{1.386998in}}%
\pgfpathlineto{\pgfqpoint{1.644487in}{1.387636in}}%
\pgfpathlineto{\pgfqpoint{1.666740in}{1.388997in}}%
\pgfpathlineto{\pgfqpoint{1.717922in}{1.389808in}}%
\pgfpathlineto{\pgfqpoint{1.737950in}{1.391438in}}%
\pgfpathlineto{\pgfqpoint{1.757978in}{1.393405in}}%
\pgfpathlineto{\pgfqpoint{1.766879in}{1.393116in}}%
\pgfpathlineto{\pgfqpoint{1.791358in}{1.391340in}}%
\pgfpathlineto{\pgfqpoint{1.804710in}{1.392433in}}%
\pgfpathlineto{\pgfqpoint{1.818062in}{1.393279in}}%
\pgfpathlineto{\pgfqpoint{1.829188in}{1.392658in}}%
\pgfpathlineto{\pgfqpoint{1.860343in}{1.389750in}}%
\pgfpathlineto{\pgfqpoint{1.878145in}{1.389421in}}%
\pgfpathlineto{\pgfqpoint{1.891497in}{1.390303in}}%
\pgfpathlineto{\pgfqpoint{1.911525in}{1.391838in}}%
\pgfpathlineto{\pgfqpoint{1.922651in}{1.391576in}}%
\pgfpathlineto{\pgfqpoint{1.936003in}{1.390017in}}%
\pgfpathlineto{\pgfqpoint{1.947130in}{1.388866in}}%
\pgfpathlineto{\pgfqpoint{1.956031in}{1.388999in}}%
\pgfpathlineto{\pgfqpoint{1.969383in}{1.390709in}}%
\pgfpathlineto{\pgfqpoint{1.982735in}{1.392144in}}%
\pgfpathlineto{\pgfqpoint{2.022791in}{1.393793in}}%
\pgfpathlineto{\pgfqpoint{2.038368in}{1.393127in}}%
\pgfpathlineto{\pgfqpoint{2.080649in}{1.389936in}}%
\pgfpathlineto{\pgfqpoint{2.091775in}{1.389408in}}%
\pgfpathlineto{\pgfqpoint{2.102902in}{1.390245in}}%
\pgfpathlineto{\pgfqpoint{2.118479in}{1.391558in}}%
\pgfpathlineto{\pgfqpoint{2.127380in}{1.391166in}}%
\pgfpathlineto{\pgfqpoint{2.185239in}{1.384807in}}%
\pgfpathlineto{\pgfqpoint{2.223069in}{1.378705in}}%
\pgfpathlineto{\pgfqpoint{2.231970in}{1.378681in}}%
\pgfpathlineto{\pgfqpoint{2.243097in}{1.379878in}}%
\pgfpathlineto{\pgfqpoint{2.258674in}{1.381746in}}%
\pgfpathlineto{\pgfqpoint{2.269801in}{1.381861in}}%
\pgfpathlineto{\pgfqpoint{2.285378in}{1.381859in}}%
\pgfpathlineto{\pgfqpoint{2.307631in}{1.382689in}}%
\pgfpathlineto{\pgfqpoint{2.316532in}{1.381422in}}%
\pgfpathlineto{\pgfqpoint{2.334335in}{1.378380in}}%
\pgfpathlineto{\pgfqpoint{2.345461in}{1.377879in}}%
\pgfpathlineto{\pgfqpoint{2.414446in}{1.377995in}}%
\pgfpathlineto{\pgfqpoint{2.436699in}{1.379154in}}%
\pgfpathlineto{\pgfqpoint{2.458953in}{1.379348in}}%
\pgfpathlineto{\pgfqpoint{2.492332in}{1.382150in}}%
\pgfpathlineto{\pgfqpoint{2.512360in}{1.382683in}}%
\pgfpathlineto{\pgfqpoint{2.539064in}{1.384916in}}%
\pgfpathlineto{\pgfqpoint{2.561317in}{1.385316in}}%
\pgfpathlineto{\pgfqpoint{2.588021in}{1.387499in}}%
\pgfpathlineto{\pgfqpoint{2.599147in}{1.386481in}}%
\pgfpathlineto{\pgfqpoint{2.612499in}{1.385203in}}%
\pgfpathlineto{\pgfqpoint{2.621401in}{1.385657in}}%
\pgfpathlineto{\pgfqpoint{2.630302in}{1.387347in}}%
\pgfpathlineto{\pgfqpoint{2.645879in}{1.390577in}}%
\pgfpathlineto{\pgfqpoint{2.654780in}{1.391103in}}%
\pgfpathlineto{\pgfqpoint{2.672583in}{1.391577in}}%
\pgfpathlineto{\pgfqpoint{2.681484in}{1.393200in}}%
\pgfpathlineto{\pgfqpoint{2.703737in}{1.398081in}}%
\pgfpathlineto{\pgfqpoint{2.712639in}{1.398625in}}%
\pgfpathlineto{\pgfqpoint{2.743793in}{1.398299in}}%
\pgfpathlineto{\pgfqpoint{2.768271in}{1.398738in}}%
\pgfpathlineto{\pgfqpoint{2.788299in}{1.399700in}}%
\pgfpathlineto{\pgfqpoint{2.803876in}{1.398877in}}%
\pgfpathlineto{\pgfqpoint{2.819454in}{1.398279in}}%
\pgfpathlineto{\pgfqpoint{2.852833in}{1.398298in}}%
\pgfpathlineto{\pgfqpoint{2.881762in}{1.396021in}}%
\pgfpathlineto{\pgfqpoint{2.890664in}{1.396846in}}%
\pgfpathlineto{\pgfqpoint{2.910692in}{1.399811in}}%
\pgfpathlineto{\pgfqpoint{2.917368in}{1.399414in}}%
\pgfpathlineto{\pgfqpoint{2.924044in}{1.397935in}}%
\pgfpathlineto{\pgfqpoint{2.946297in}{1.391759in}}%
\pgfpathlineto{\pgfqpoint{2.952973in}{1.391496in}}%
\pgfpathlineto{\pgfqpoint{2.961874in}{1.392450in}}%
\pgfpathlineto{\pgfqpoint{2.984127in}{1.395807in}}%
\pgfpathlineto{\pgfqpoint{2.993028in}{1.395361in}}%
\pgfpathlineto{\pgfqpoint{3.001930in}{1.393571in}}%
\pgfpathlineto{\pgfqpoint{3.035309in}{1.384877in}}%
\pgfpathlineto{\pgfqpoint{3.039760in}{1.384344in}}%
\pgfpathlineto{\pgfqpoint{3.039760in}{1.384344in}}%
\pgfusepath{stroke}%
\end{pgfscope}%
\begin{pgfscope}%
\pgfpathrectangle{\pgfqpoint{0.650000in}{0.420000in}}{\pgfqpoint{2.388110in}{1.994488in}}%
\pgfusepath{clip}%
\pgfsetrectcap%
\pgfsetroundjoin%
\pgfsetlinewidth{1.003750pt}%
\definecolor{currentstroke}{rgb}{0.400000,0.760784,0.643137}%
\pgfsetstrokecolor{currentstroke}%
\pgfsetdash{}{0pt}%
\pgfpathmoveto{\pgfqpoint{0.650000in}{1.281252in}}%
\pgfpathlineto{\pgfqpoint{0.656972in}{1.281845in}}%
\pgfpathlineto{\pgfqpoint{0.663530in}{1.283423in}}%
\pgfpathlineto{\pgfqpoint{0.670834in}{1.286378in}}%
\pgfpathlineto{\pgfqpoint{0.678383in}{1.290415in}}%
\pgfpathlineto{\pgfqpoint{0.684086in}{1.294789in}}%
\pgfpathlineto{\pgfqpoint{0.687893in}{1.299420in}}%
\pgfpathlineto{\pgfqpoint{0.691702in}{1.306656in}}%
\pgfpathlineto{\pgfqpoint{0.695510in}{1.317617in}}%
\pgfpathlineto{\pgfqpoint{0.699319in}{1.332853in}}%
\pgfpathlineto{\pgfqpoint{0.705033in}{1.361794in}}%
\pgfpathlineto{\pgfqpoint{0.710746in}{1.390430in}}%
\pgfpathlineto{\pgfqpoint{0.714555in}{1.404338in}}%
\pgfpathlineto{\pgfqpoint{0.716459in}{1.408965in}}%
\pgfpathlineto{\pgfqpoint{0.718364in}{1.411894in}}%
\pgfpathlineto{\pgfqpoint{0.720268in}{1.413202in}}%
\pgfpathlineto{\pgfqpoint{0.722173in}{1.412958in}}%
\pgfpathlineto{\pgfqpoint{0.724077in}{1.411332in}}%
\pgfpathlineto{\pgfqpoint{0.727886in}{1.404802in}}%
\pgfpathlineto{\pgfqpoint{0.733599in}{1.390525in}}%
\pgfpathlineto{\pgfqpoint{0.739313in}{1.376458in}}%
\pgfpathlineto{\pgfqpoint{0.743122in}{1.369596in}}%
\pgfpathlineto{\pgfqpoint{0.746930in}{1.365715in}}%
\pgfpathlineto{\pgfqpoint{0.748835in}{1.364990in}}%
\pgfpathlineto{\pgfqpoint{0.750739in}{1.365068in}}%
\pgfpathlineto{\pgfqpoint{0.752644in}{1.365913in}}%
\pgfpathlineto{\pgfqpoint{0.756453in}{1.369735in}}%
\pgfpathlineto{\pgfqpoint{0.760262in}{1.375716in}}%
\pgfpathlineto{\pgfqpoint{0.765975in}{1.387472in}}%
\pgfpathlineto{\pgfqpoint{0.786924in}{1.434496in}}%
\pgfpathlineto{\pgfqpoint{0.792637in}{1.444599in}}%
\pgfpathlineto{\pgfqpoint{0.798351in}{1.452851in}}%
\pgfpathlineto{\pgfqpoint{0.804064in}{1.459018in}}%
\pgfpathlineto{\pgfqpoint{0.809777in}{1.463079in}}%
\pgfpathlineto{\pgfqpoint{0.813586in}{1.464647in}}%
\pgfpathlineto{\pgfqpoint{0.819299in}{1.465464in}}%
\pgfpathlineto{\pgfqpoint{0.825013in}{1.464842in}}%
\pgfpathlineto{\pgfqpoint{0.832631in}{1.462648in}}%
\pgfpathlineto{\pgfqpoint{0.889764in}{1.443181in}}%
\pgfpathlineto{\pgfqpoint{0.906904in}{1.440069in}}%
\pgfpathlineto{\pgfqpoint{0.920235in}{1.437028in}}%
\pgfpathlineto{\pgfqpoint{0.931662in}{1.433174in}}%
\pgfpathlineto{\pgfqpoint{0.958324in}{1.423263in}}%
\pgfpathlineto{\pgfqpoint{0.975464in}{1.418609in}}%
\pgfpathlineto{\pgfqpoint{0.988795in}{1.416104in}}%
\pgfpathlineto{\pgfqpoint{1.013553in}{1.412304in}}%
\pgfpathlineto{\pgfqpoint{1.034502in}{1.408066in}}%
\pgfpathlineto{\pgfqpoint{1.064973in}{1.404122in}}%
\pgfpathlineto{\pgfqpoint{1.076400in}{1.403711in}}%
\pgfpathlineto{\pgfqpoint{1.118298in}{1.403970in}}%
\pgfpathlineto{\pgfqpoint{1.133533in}{1.403224in}}%
\pgfpathlineto{\pgfqpoint{1.143055in}{1.403980in}}%
\pgfpathlineto{\pgfqpoint{1.171622in}{1.407916in}}%
\pgfpathlineto{\pgfqpoint{1.198284in}{1.408710in}}%
\pgfpathlineto{\pgfqpoint{1.211615in}{1.409344in}}%
\pgfpathlineto{\pgfqpoint{1.240182in}{1.409276in}}%
\pgfpathlineto{\pgfqpoint{1.266844in}{1.411725in}}%
\pgfpathlineto{\pgfqpoint{1.299220in}{1.413176in}}%
\pgfpathlineto{\pgfqpoint{1.316360in}{1.414163in}}%
\pgfpathlineto{\pgfqpoint{1.327787in}{1.413484in}}%
\pgfpathlineto{\pgfqpoint{1.341118in}{1.412706in}}%
\pgfpathlineto{\pgfqpoint{1.350640in}{1.413491in}}%
\pgfpathlineto{\pgfqpoint{1.369684in}{1.415756in}}%
\pgfpathlineto{\pgfqpoint{1.379207in}{1.415280in}}%
\pgfpathlineto{\pgfqpoint{1.400156in}{1.413242in}}%
\pgfpathlineto{\pgfqpoint{1.409678in}{1.413810in}}%
\pgfpathlineto{\pgfqpoint{1.440149in}{1.416789in}}%
\pgfpathlineto{\pgfqpoint{1.463002in}{1.417519in}}%
\pgfpathlineto{\pgfqpoint{1.476333in}{1.416900in}}%
\pgfpathlineto{\pgfqpoint{1.493473in}{1.415494in}}%
\pgfpathlineto{\pgfqpoint{1.501091in}{1.416168in}}%
\pgfpathlineto{\pgfqpoint{1.523945in}{1.419739in}}%
\pgfpathlineto{\pgfqpoint{1.533467in}{1.419334in}}%
\pgfpathlineto{\pgfqpoint{1.552511in}{1.418184in}}%
\pgfpathlineto{\pgfqpoint{1.575365in}{1.418441in}}%
\pgfpathlineto{\pgfqpoint{1.594409in}{1.418885in}}%
\pgfpathlineto{\pgfqpoint{1.603931in}{1.417693in}}%
\pgfpathlineto{\pgfqpoint{1.626785in}{1.413773in}}%
\pgfpathlineto{\pgfqpoint{1.634402in}{1.414151in}}%
\pgfpathlineto{\pgfqpoint{1.645829in}{1.416183in}}%
\pgfpathlineto{\pgfqpoint{1.657256in}{1.418130in}}%
\pgfpathlineto{\pgfqpoint{1.664874in}{1.418472in}}%
\pgfpathlineto{\pgfqpoint{1.674396in}{1.417550in}}%
\pgfpathlineto{\pgfqpoint{1.693440in}{1.414824in}}%
\pgfpathlineto{\pgfqpoint{1.704867in}{1.414799in}}%
\pgfpathlineto{\pgfqpoint{1.722007in}{1.415115in}}%
\pgfpathlineto{\pgfqpoint{1.744860in}{1.414437in}}%
\pgfpathlineto{\pgfqpoint{1.777236in}{1.417928in}}%
\pgfpathlineto{\pgfqpoint{1.803898in}{1.422363in}}%
\pgfpathlineto{\pgfqpoint{1.813420in}{1.422238in}}%
\pgfpathlineto{\pgfqpoint{1.853414in}{1.419291in}}%
\pgfpathlineto{\pgfqpoint{1.862936in}{1.416737in}}%
\pgfpathlineto{\pgfqpoint{1.874363in}{1.413804in}}%
\pgfpathlineto{\pgfqpoint{1.880076in}{1.413427in}}%
\pgfpathlineto{\pgfqpoint{1.885789in}{1.414000in}}%
\pgfpathlineto{\pgfqpoint{1.895311in}{1.416286in}}%
\pgfpathlineto{\pgfqpoint{1.906738in}{1.418954in}}%
\pgfpathlineto{\pgfqpoint{1.916260in}{1.420004in}}%
\pgfpathlineto{\pgfqpoint{1.925783in}{1.419852in}}%
\pgfpathlineto{\pgfqpoint{1.944827in}{1.419001in}}%
\pgfpathlineto{\pgfqpoint{1.971489in}{1.420026in}}%
\pgfpathlineto{\pgfqpoint{1.982916in}{1.418733in}}%
\pgfpathlineto{\pgfqpoint{2.001961in}{1.416006in}}%
\pgfpathlineto{\pgfqpoint{2.013387in}{1.416022in}}%
\pgfpathlineto{\pgfqpoint{2.030527in}{1.417286in}}%
\pgfpathlineto{\pgfqpoint{2.040049in}{1.419134in}}%
\pgfpathlineto{\pgfqpoint{2.060998in}{1.424140in}}%
\pgfpathlineto{\pgfqpoint{2.068616in}{1.424201in}}%
\pgfpathlineto{\pgfqpoint{2.076234in}{1.423012in}}%
\pgfpathlineto{\pgfqpoint{2.091469in}{1.420010in}}%
\pgfpathlineto{\pgfqpoint{2.100992in}{1.419539in}}%
\pgfpathlineto{\pgfqpoint{2.127654in}{1.419861in}}%
\pgfpathlineto{\pgfqpoint{2.135272in}{1.421375in}}%
\pgfpathlineto{\pgfqpoint{2.150507in}{1.424990in}}%
\pgfpathlineto{\pgfqpoint{2.156221in}{1.425183in}}%
\pgfpathlineto{\pgfqpoint{2.163838in}{1.424181in}}%
\pgfpathlineto{\pgfqpoint{2.180978in}{1.421193in}}%
\pgfpathlineto{\pgfqpoint{2.192405in}{1.420848in}}%
\pgfpathlineto{\pgfqpoint{2.211450in}{1.421222in}}%
\pgfpathlineto{\pgfqpoint{2.219067in}{1.422657in}}%
\pgfpathlineto{\pgfqpoint{2.234303in}{1.426265in}}%
\pgfpathlineto{\pgfqpoint{2.240016in}{1.426230in}}%
\pgfpathlineto{\pgfqpoint{2.245730in}{1.424967in}}%
\pgfpathlineto{\pgfqpoint{2.253347in}{1.421846in}}%
\pgfpathlineto{\pgfqpoint{2.264774in}{1.416994in}}%
\pgfpathlineto{\pgfqpoint{2.272392in}{1.415207in}}%
\pgfpathlineto{\pgfqpoint{2.281914in}{1.414338in}}%
\pgfpathlineto{\pgfqpoint{2.299054in}{1.413988in}}%
\pgfpathlineto{\pgfqpoint{2.308576in}{1.414904in}}%
\pgfpathlineto{\pgfqpoint{2.320003in}{1.417207in}}%
\pgfpathlineto{\pgfqpoint{2.333334in}{1.419752in}}%
\pgfpathlineto{\pgfqpoint{2.342856in}{1.420253in}}%
\pgfpathlineto{\pgfqpoint{2.352379in}{1.420921in}}%
\pgfpathlineto{\pgfqpoint{2.359996in}{1.422896in}}%
\pgfpathlineto{\pgfqpoint{2.375232in}{1.427989in}}%
\pgfpathlineto{\pgfqpoint{2.380945in}{1.428408in}}%
\pgfpathlineto{\pgfqpoint{2.386659in}{1.427786in}}%
\pgfpathlineto{\pgfqpoint{2.419034in}{1.421949in}}%
\pgfpathlineto{\pgfqpoint{2.438079in}{1.420702in}}%
\pgfpathlineto{\pgfqpoint{2.445696in}{1.421163in}}%
\pgfpathlineto{\pgfqpoint{2.455219in}{1.422996in}}%
\pgfpathlineto{\pgfqpoint{2.468550in}{1.425902in}}%
\pgfpathlineto{\pgfqpoint{2.474263in}{1.426141in}}%
\pgfpathlineto{\pgfqpoint{2.479976in}{1.425443in}}%
\pgfpathlineto{\pgfqpoint{2.500925in}{1.421095in}}%
\pgfpathlineto{\pgfqpoint{2.508543in}{1.421439in}}%
\pgfpathlineto{\pgfqpoint{2.521874in}{1.422318in}}%
\pgfpathlineto{\pgfqpoint{2.531396in}{1.421277in}}%
\pgfpathlineto{\pgfqpoint{2.540919in}{1.420219in}}%
\pgfpathlineto{\pgfqpoint{2.548537in}{1.420709in}}%
\pgfpathlineto{\pgfqpoint{2.558059in}{1.422759in}}%
\pgfpathlineto{\pgfqpoint{2.573294in}{1.425908in}}%
\pgfpathlineto{\pgfqpoint{2.582817in}{1.426800in}}%
\pgfpathlineto{\pgfqpoint{2.590434in}{1.426338in}}%
\pgfpathlineto{\pgfqpoint{2.611383in}{1.423427in}}%
\pgfpathlineto{\pgfqpoint{2.638046in}{1.423981in}}%
\pgfpathlineto{\pgfqpoint{2.649472in}{1.425709in}}%
\pgfpathlineto{\pgfqpoint{2.660899in}{1.427120in}}%
\pgfpathlineto{\pgfqpoint{2.687561in}{1.428765in}}%
\pgfpathlineto{\pgfqpoint{2.697083in}{1.429703in}}%
\pgfpathlineto{\pgfqpoint{2.704701in}{1.429166in}}%
\pgfpathlineto{\pgfqpoint{2.716128in}{1.426565in}}%
\pgfpathlineto{\pgfqpoint{2.723746in}{1.425517in}}%
\pgfpathlineto{\pgfqpoint{2.729459in}{1.425886in}}%
\pgfpathlineto{\pgfqpoint{2.748504in}{1.428997in}}%
\pgfpathlineto{\pgfqpoint{2.754217in}{1.428135in}}%
\pgfpathlineto{\pgfqpoint{2.761835in}{1.425668in}}%
\pgfpathlineto{\pgfqpoint{2.780879in}{1.418759in}}%
\pgfpathlineto{\pgfqpoint{2.786592in}{1.417764in}}%
\pgfpathlineto{\pgfqpoint{2.792306in}{1.417807in}}%
\pgfpathlineto{\pgfqpoint{2.798019in}{1.418962in}}%
\pgfpathlineto{\pgfqpoint{2.805637in}{1.421766in}}%
\pgfpathlineto{\pgfqpoint{2.822777in}{1.428509in}}%
\pgfpathlineto{\pgfqpoint{2.841821in}{1.434691in}}%
\pgfpathlineto{\pgfqpoint{2.847535in}{1.435350in}}%
\pgfpathlineto{\pgfqpoint{2.853248in}{1.434763in}}%
\pgfpathlineto{\pgfqpoint{2.860866in}{1.432369in}}%
\pgfpathlineto{\pgfqpoint{2.872293in}{1.428443in}}%
\pgfpathlineto{\pgfqpoint{2.879910in}{1.427305in}}%
\pgfpathlineto{\pgfqpoint{2.889432in}{1.427273in}}%
\pgfpathlineto{\pgfqpoint{2.906572in}{1.427382in}}%
\pgfpathlineto{\pgfqpoint{2.914190in}{1.426498in}}%
\pgfpathlineto{\pgfqpoint{2.921808in}{1.424529in}}%
\pgfpathlineto{\pgfqpoint{2.937044in}{1.419816in}}%
\pgfpathlineto{\pgfqpoint{2.942757in}{1.419383in}}%
\pgfpathlineto{\pgfqpoint{2.950375in}{1.420264in}}%
\pgfpathlineto{\pgfqpoint{2.961801in}{1.422045in}}%
\pgfpathlineto{\pgfqpoint{2.969419in}{1.421934in}}%
\pgfpathlineto{\pgfqpoint{2.982750in}{1.420961in}}%
\pgfpathlineto{\pgfqpoint{2.988464in}{1.421771in}}%
\pgfpathlineto{\pgfqpoint{2.994177in}{1.423742in}}%
\pgfpathlineto{\pgfqpoint{3.015126in}{1.433189in}}%
\pgfpathlineto{\pgfqpoint{3.020839in}{1.433992in}}%
\pgfpathlineto{\pgfqpoint{3.026553in}{1.433726in}}%
\pgfpathlineto{\pgfqpoint{3.034170in}{1.432140in}}%
\pgfpathlineto{\pgfqpoint{3.039884in}{1.430615in}}%
\pgfpathlineto{\pgfqpoint{3.039884in}{1.430615in}}%
\pgfusepath{stroke}%
\end{pgfscope}%
\begin{pgfscope}%
\pgfsetrectcap%
\pgfsetmiterjoin%
\pgfsetlinewidth{0.803000pt}%
\definecolor{currentstroke}{rgb}{0.000000,0.000000,0.000000}%
\pgfsetstrokecolor{currentstroke}%
\pgfsetdash{}{0pt}%
\pgfpathmoveto{\pgfqpoint{0.650000in}{0.420000in}}%
\pgfpathlineto{\pgfqpoint{0.650000in}{2.414488in}}%
\pgfusepath{stroke}%
\end{pgfscope}%
\begin{pgfscope}%
\pgfsetrectcap%
\pgfsetmiterjoin%
\pgfsetlinewidth{0.803000pt}%
\definecolor{currentstroke}{rgb}{0.000000,0.000000,0.000000}%
\pgfsetstrokecolor{currentstroke}%
\pgfsetdash{}{0pt}%
\pgfpathmoveto{\pgfqpoint{3.038110in}{0.420000in}}%
\pgfpathlineto{\pgfqpoint{3.038110in}{2.414488in}}%
\pgfusepath{stroke}%
\end{pgfscope}%
\begin{pgfscope}%
\pgfsetrectcap%
\pgfsetmiterjoin%
\pgfsetlinewidth{0.803000pt}%
\definecolor{currentstroke}{rgb}{0.000000,0.000000,0.000000}%
\pgfsetstrokecolor{currentstroke}%
\pgfsetdash{}{0pt}%
\pgfpathmoveto{\pgfqpoint{0.650000in}{0.420000in}}%
\pgfpathlineto{\pgfqpoint{3.038110in}{0.420000in}}%
\pgfusepath{stroke}%
\end{pgfscope}%
\begin{pgfscope}%
\pgfsetrectcap%
\pgfsetmiterjoin%
\pgfsetlinewidth{0.803000pt}%
\definecolor{currentstroke}{rgb}{0.000000,0.000000,0.000000}%
\pgfsetstrokecolor{currentstroke}%
\pgfsetdash{}{0pt}%
\pgfpathmoveto{\pgfqpoint{0.650000in}{2.414488in}}%
\pgfpathlineto{\pgfqpoint{3.038110in}{2.414488in}}%
\pgfusepath{stroke}%
\end{pgfscope}%
\begin{pgfscope}%
\pgfsetrectcap%
\pgfsetroundjoin%
\pgfsetlinewidth{1.003750pt}%
\definecolor{currentstroke}{rgb}{0.870588,0.466667,0.682353}%
\pgfsetstrokecolor{currentstroke}%
\pgfsetdash{}{0pt}%
\pgfpathmoveto{\pgfqpoint{2.108975in}{2.275599in}}%
\pgfpathlineto{\pgfqpoint{2.220086in}{2.275599in}}%
\pgfpathlineto{\pgfqpoint{2.331197in}{2.275599in}}%
\pgfusepath{stroke}%
\end{pgfscope}%
\begin{pgfscope}%
\definecolor{textcolor}{rgb}{0.000000,0.000000,0.000000}%
\pgfsetstrokecolor{textcolor}%
\pgfsetfillcolor{textcolor}%
\pgftext[x=2.420086in,y=2.236710in,left,base]{\color{textcolor}{\rmfamily\fontsize{8.000000}{9.600000}\selectfont\catcode`\^=\active\def^{\ifmmode\sp\else\^{}\fi}\catcode`\%=\active\def%{\%}base}}%
\end{pgfscope}%
\begin{pgfscope}%
\pgfsetrectcap%
\pgfsetroundjoin%
\pgfsetlinewidth{1.003750pt}%
\definecolor{currentstroke}{rgb}{0.498039,0.737255,0.254902}%
\pgfsetstrokecolor{currentstroke}%
\pgfsetdash{}{0pt}%
\pgfpathmoveto{\pgfqpoint{2.108975in}{2.120661in}}%
\pgfpathlineto{\pgfqpoint{2.220086in}{2.120661in}}%
\pgfpathlineto{\pgfqpoint{2.331197in}{2.120661in}}%
\pgfusepath{stroke}%
\end{pgfscope}%
\begin{pgfscope}%
\definecolor{textcolor}{rgb}{0.000000,0.000000,0.000000}%
\pgfsetstrokecolor{textcolor}%
\pgfsetfillcolor{textcolor}%
\pgftext[x=2.420086in,y=2.081772in,left,base]{\color{textcolor}{\rmfamily\fontsize{8.000000}{9.600000}\selectfont\catcode`\^=\active\def^{\ifmmode\sp\else\^{}\fi}\catcode`\%=\active\def%{\%}optimised}}%
\end{pgfscope}%
\begin{pgfscope}%
\pgfsetrectcap%
\pgfsetroundjoin%
\pgfsetlinewidth{1.003750pt}%
\definecolor{currentstroke}{rgb}{0.400000,0.760784,0.643137}%
\pgfsetstrokecolor{currentstroke}%
\pgfsetdash{}{0pt}%
\pgfpathmoveto{\pgfqpoint{2.108975in}{1.965723in}}%
\pgfpathlineto{\pgfqpoint{2.220086in}{1.965723in}}%
\pgfpathlineto{\pgfqpoint{2.331197in}{1.965723in}}%
\pgfusepath{stroke}%
\end{pgfscope}%
\begin{pgfscope}%
\definecolor{textcolor}{rgb}{0.000000,0.000000,0.000000}%
\pgfsetstrokecolor{textcolor}%
\pgfsetfillcolor{textcolor}%
\pgftext[x=2.420086in,y=1.926834in,left,base]{\color{textcolor}{\rmfamily\fontsize{8.000000}{9.600000}\selectfont\catcode`\^=\active\def^{\ifmmode\sp\else\^{}\fi}\catcode`\%=\active\def%{\%}planeTBL}}%
\end{pgfscope}%
\end{pgfpicture}%
\makeatother%
\endgroup%

    \label{fig:inflowSecVerifBeta}
\end{subfigure}
\\[-15pt]
\caption[First and second order statistics and Clauser parameter of flat \acrshort{tbl} and R100deg30 base and optimised simulations at $\mathrm{Re}_{\theta} = 700$ with reference data from \citet{schlatter_assessment_2010} at $\mathrm{Re}_{\theta} = 677$]{First and second order statistics of flat \acrshort{tbl} and R100deg30 base and optimised simulations at $\mathrm{Re}_{\theta} = 700$ with reference data taken from \citet{schlatter_assessment_2010} at $\mathrm{Re}_{\theta} = 677$: (a) streamwise velocity profile in viscous scales; (b) Reynolds stresses $\langle u_i' u_j' \rangle$; budget of \acrshort{tke} = $\frac{1}{2} \langle u_i' u_i' \rangle$; (c) budget of \acrshort{tke} = $\frac{1}{2} \langle u_i' u_i' \rangle$, line styles denote budget terms according to \autoref{tab:budgetLinestyles};
(d) Clauser parameter over non-dimensionalised streamwise coordinate.}
\label{fig:inflowSecVerif}
\end{figure}

The first- and second-order statistics of the compared simulations are depicted in \autoref{fig:inflowSecVerif}a-c.
The streamwise velocity profile in viscous scaling is illustrated in \autoref{fig:flatTBLverifU}.
It is evident that all profiles match each other. 
As already discussed in \autoref{subsec:flatTBLResults}, only the free-stream velocity of the reference data is slightly lower than that of the other profiles, since they are reported at a lower value of $\mathrm{Re}_{\theta}$.
Therefore, the boundary layer thickness is smaller, resulting in a faster achievement of the free-stream velocity.
In \autoref{fig:inflowSecVerifFluct}, the Reynolds stresses are visualised. 
We observe a satisfactory alignment for all components in the performed simulations.
Again, the reference data show a minor discrepancy in the outer region of the boundary layer, in the higher $y^+$ regions, respectively.
This phenomenon can also be implied by the lower Reynolds number.
The budget of the \acrshort{tke}, visualised in \autoref{fig:inflowSecVerifBudget}, gives a similar impression.
In this particular metric, the three profiles were found to be highly consistent with one another and the reference data.
The statistical analysis conducted thus far has revealed no statistically significant differences between the flat and the curved \acrshort{tbl} simulations, encompassing both the base and the optimised simulation setup of the R100deg30 case.
The final quantity under discussion is the Clauser parameter $\beta$, illustrated in \autoref{fig:inflowSecVerifBeta}.
$\beta$ is shown here only in the straight region of the domains, as this region is the primary focus of the present verification.
The streamwise coordinate is normalised by the reference boundary layer thickness $\delta_0$.
The trend of all simulations looks similar, indicating that the parameter $\beta$ 
remains close to zero, with a minor balancing region observed at the domain's beginning. 
As the direction of flow over the streamwise direction is observed, a slight increase in the distance between the flat setup and the curved setups is evident. This increase reaches a maximum distance of approximately $0.1$ at the termination of the straight section.
This trend could already show an impact of the curvature, even in the upstream region.
However, all profiles show quite smooth trends of $\beta$ indicating low to negligible impact of the pressure fluctuations on the mean Clauser parameter.
In summary, the comparisons obtained demonstrate that, in the straight region, the simulations generate analogous profiles which leads to two primary conclusions.
First, it generally verifies the curved simulation setup in enabling an undisturbed \acrshort{tbl} development in the straight section.
Secondly, the results suggest that the observed pressure fluctuations, as discussed in \autoref{sec:pressFluc}, show no visible manipulation of the statistics.



\section{Curved turbulent boundary layer results}\label{sec:curvedResults}

In the following section the results of the curved \acrshort{tbl} simulations are discussed.
Firstly, the results of each tested case are presented in isolation to directly compare the optimised configuration and the base configuration.
Then, the optimised configurations of all cases are compared to study the effect of different geometrical configurations on the optimisation and the development of the \acrshort{tbl}.


\subsection{Optimisation performance analysis for case R100deg30}\label{subsec:R100deg30}
First, the results of case R100deg30 are presented. 
This case is analysed using multiple quantities to ensure a proper evolution of the \acrshort{tbl} over the curvature.
We initially focus on the first- and second-order statistics extracted at a value of $\mathrm{Re}_{\theta} = 1200$. 
This position is located within the curvature, with a certain distance to both ends. 
This is done to mitigate interfering effects of the geometry changes.
Then, the development of the boundary layer thickness and Reynolds numbers in the streamwise direction is presented. 
Here, we use as reference data from flat \acrshort{tbl} simulations performed by \citet{jimenez_near-wall_2013} extracted at $\mathrm{Re}_{\theta} = 1100$.
Finally, we will discuss the pressure-oriented parameters to evaluate the optimisation capability. 

\begin{figure}[h!]
\centering
\begin{subfigure}[t]{0.495\textwidth}
    \subcaption{}
    \centering
    %% Creator: Matplotlib, PGF backend
%%
%% To include the figure in your LaTeX document, write
%%   \input{<filename>.pgf}
%%
%% Make sure the required packages are loaded in your preamble
%%   \usepackage{pgf}
%%
%% Also ensure that all the required font packages are loaded; for instance,
%% the lmodern package is sometimes necessary when using math font.
%%   \usepackage{lmodern}
%%
%% Figures using additional raster images can only be included by \input if
%% they are in the same directory as the main LaTeX file. For loading figures
%% from other directories you can use the `import` package
%%   \usepackage{import}
%%
%% and then include the figures with
%%   \import{<path to file>}{<filename>.pgf}
%%
%% Matplotlib used the following preamble
%%   \def\mathdefault#1{#1}
%%   \everymath=\expandafter{\the\everymath\displaystyle}
%%   \IfFileExists{scrextend.sty}{
%%     \usepackage[fontsize=11.000000pt]{scrextend}
%%   }{
%%     \renewcommand{\normalsize}{\fontsize{11.000000}{13.200000}\selectfont}
%%     \normalsize
%%   }
%%   
%%   \makeatletter\@ifpackageloaded{underscore}{}{\usepackage[strings]{underscore}}\makeatother
%%
\begingroup%
\makeatletter%
\begin{pgfpicture}%
\pgfpathrectangle{\pgfpointorigin}{\pgfqpoint{3.118110in}{2.494488in}}%
\pgfusepath{use as bounding box, clip}%
\begin{pgfscope}%
\pgfsetbuttcap%
\pgfsetmiterjoin%
\definecolor{currentfill}{rgb}{1.000000,1.000000,1.000000}%
\pgfsetfillcolor{currentfill}%
\pgfsetlinewidth{0.000000pt}%
\definecolor{currentstroke}{rgb}{1.000000,1.000000,1.000000}%
\pgfsetstrokecolor{currentstroke}%
\pgfsetdash{}{0pt}%
\pgfpathmoveto{\pgfqpoint{0.000000in}{0.000000in}}%
\pgfpathlineto{\pgfqpoint{3.118110in}{0.000000in}}%
\pgfpathlineto{\pgfqpoint{3.118110in}{2.494488in}}%
\pgfpathlineto{\pgfqpoint{0.000000in}{2.494488in}}%
\pgfpathlineto{\pgfqpoint{0.000000in}{0.000000in}}%
\pgfpathclose%
\pgfusepath{fill}%
\end{pgfscope}%
\begin{pgfscope}%
\pgfsetbuttcap%
\pgfsetmiterjoin%
\definecolor{currentfill}{rgb}{1.000000,1.000000,1.000000}%
\pgfsetfillcolor{currentfill}%
\pgfsetlinewidth{0.000000pt}%
\definecolor{currentstroke}{rgb}{0.000000,0.000000,0.000000}%
\pgfsetstrokecolor{currentstroke}%
\pgfsetstrokeopacity{0.000000}%
\pgfsetdash{}{0pt}%
\pgfpathmoveto{\pgfqpoint{0.650000in}{0.420000in}}%
\pgfpathlineto{\pgfqpoint{3.038110in}{0.420000in}}%
\pgfpathlineto{\pgfqpoint{3.038110in}{2.414488in}}%
\pgfpathlineto{\pgfqpoint{0.650000in}{2.414488in}}%
\pgfpathlineto{\pgfqpoint{0.650000in}{0.420000in}}%
\pgfpathclose%
\pgfusepath{fill}%
\end{pgfscope}%
\begin{pgfscope}%
\pgfsetbuttcap%
\pgfsetroundjoin%
\definecolor{currentfill}{rgb}{0.000000,0.000000,0.000000}%
\pgfsetfillcolor{currentfill}%
\pgfsetlinewidth{0.501875pt}%
\definecolor{currentstroke}{rgb}{0.000000,0.000000,0.000000}%
\pgfsetstrokecolor{currentstroke}%
\pgfsetdash{}{0pt}%
\pgfsys@defobject{currentmarker}{\pgfqpoint{0.000000in}{0.000000in}}{\pgfqpoint{0.000000in}{0.041667in}}{%
\pgfpathmoveto{\pgfqpoint{0.000000in}{0.000000in}}%
\pgfpathlineto{\pgfqpoint{0.000000in}{0.041667in}}%
\pgfusepath{stroke,fill}%
}%
\begin{pgfscope}%
\pgfsys@transformshift{0.862677in}{0.420000in}%
\pgfsys@useobject{currentmarker}{}%
\end{pgfscope}%
\end{pgfscope}%
\begin{pgfscope}%
\pgfsetbuttcap%
\pgfsetroundjoin%
\definecolor{currentfill}{rgb}{0.000000,0.000000,0.000000}%
\pgfsetfillcolor{currentfill}%
\pgfsetlinewidth{0.501875pt}%
\definecolor{currentstroke}{rgb}{0.000000,0.000000,0.000000}%
\pgfsetstrokecolor{currentstroke}%
\pgfsetdash{}{0pt}%
\pgfsys@defobject{currentmarker}{\pgfqpoint{0.000000in}{-0.041667in}}{\pgfqpoint{0.000000in}{0.000000in}}{%
\pgfpathmoveto{\pgfqpoint{0.000000in}{0.000000in}}%
\pgfpathlineto{\pgfqpoint{0.000000in}{-0.041667in}}%
\pgfusepath{stroke,fill}%
}%
\begin{pgfscope}%
\pgfsys@transformshift{0.862677in}{2.414488in}%
\pgfsys@useobject{currentmarker}{}%
\end{pgfscope}%
\end{pgfscope}%
\begin{pgfscope}%
\definecolor{textcolor}{rgb}{0.000000,0.000000,0.000000}%
\pgfsetstrokecolor{textcolor}%
\pgfsetfillcolor{textcolor}%
\pgftext[x=0.862677in,y=0.371389in,,top]{\color{textcolor}{\rmfamily\fontsize{11.000000}{13.200000}\selectfont\catcode`\^=\active\def^{\ifmmode\sp\else\^{}\fi}\catcode`\%=\active\def%{\%}$\mathdefault{10^{0}}$}}%
\end{pgfscope}%
\begin{pgfscope}%
\pgfsetbuttcap%
\pgfsetroundjoin%
\definecolor{currentfill}{rgb}{0.000000,0.000000,0.000000}%
\pgfsetfillcolor{currentfill}%
\pgfsetlinewidth{0.501875pt}%
\definecolor{currentstroke}{rgb}{0.000000,0.000000,0.000000}%
\pgfsetstrokecolor{currentstroke}%
\pgfsetdash{}{0pt}%
\pgfsys@defobject{currentmarker}{\pgfqpoint{0.000000in}{0.000000in}}{\pgfqpoint{0.000000in}{0.041667in}}{%
\pgfpathmoveto{\pgfqpoint{0.000000in}{0.000000in}}%
\pgfpathlineto{\pgfqpoint{0.000000in}{0.041667in}}%
\pgfusepath{stroke,fill}%
}%
\begin{pgfscope}%
\pgfsys@transformshift{1.569174in}{0.420000in}%
\pgfsys@useobject{currentmarker}{}%
\end{pgfscope}%
\end{pgfscope}%
\begin{pgfscope}%
\pgfsetbuttcap%
\pgfsetroundjoin%
\definecolor{currentfill}{rgb}{0.000000,0.000000,0.000000}%
\pgfsetfillcolor{currentfill}%
\pgfsetlinewidth{0.501875pt}%
\definecolor{currentstroke}{rgb}{0.000000,0.000000,0.000000}%
\pgfsetstrokecolor{currentstroke}%
\pgfsetdash{}{0pt}%
\pgfsys@defobject{currentmarker}{\pgfqpoint{0.000000in}{-0.041667in}}{\pgfqpoint{0.000000in}{0.000000in}}{%
\pgfpathmoveto{\pgfqpoint{0.000000in}{0.000000in}}%
\pgfpathlineto{\pgfqpoint{0.000000in}{-0.041667in}}%
\pgfusepath{stroke,fill}%
}%
\begin{pgfscope}%
\pgfsys@transformshift{1.569174in}{2.414488in}%
\pgfsys@useobject{currentmarker}{}%
\end{pgfscope}%
\end{pgfscope}%
\begin{pgfscope}%
\definecolor{textcolor}{rgb}{0.000000,0.000000,0.000000}%
\pgfsetstrokecolor{textcolor}%
\pgfsetfillcolor{textcolor}%
\pgftext[x=1.569174in,y=0.371389in,,top]{\color{textcolor}{\rmfamily\fontsize{11.000000}{13.200000}\selectfont\catcode`\^=\active\def^{\ifmmode\sp\else\^{}\fi}\catcode`\%=\active\def%{\%}$\mathdefault{10^{1}}$}}%
\end{pgfscope}%
\begin{pgfscope}%
\pgfsetbuttcap%
\pgfsetroundjoin%
\definecolor{currentfill}{rgb}{0.000000,0.000000,0.000000}%
\pgfsetfillcolor{currentfill}%
\pgfsetlinewidth{0.501875pt}%
\definecolor{currentstroke}{rgb}{0.000000,0.000000,0.000000}%
\pgfsetstrokecolor{currentstroke}%
\pgfsetdash{}{0pt}%
\pgfsys@defobject{currentmarker}{\pgfqpoint{0.000000in}{0.000000in}}{\pgfqpoint{0.000000in}{0.041667in}}{%
\pgfpathmoveto{\pgfqpoint{0.000000in}{0.000000in}}%
\pgfpathlineto{\pgfqpoint{0.000000in}{0.041667in}}%
\pgfusepath{stroke,fill}%
}%
\begin{pgfscope}%
\pgfsys@transformshift{2.275671in}{0.420000in}%
\pgfsys@useobject{currentmarker}{}%
\end{pgfscope}%
\end{pgfscope}%
\begin{pgfscope}%
\pgfsetbuttcap%
\pgfsetroundjoin%
\definecolor{currentfill}{rgb}{0.000000,0.000000,0.000000}%
\pgfsetfillcolor{currentfill}%
\pgfsetlinewidth{0.501875pt}%
\definecolor{currentstroke}{rgb}{0.000000,0.000000,0.000000}%
\pgfsetstrokecolor{currentstroke}%
\pgfsetdash{}{0pt}%
\pgfsys@defobject{currentmarker}{\pgfqpoint{0.000000in}{-0.041667in}}{\pgfqpoint{0.000000in}{0.000000in}}{%
\pgfpathmoveto{\pgfqpoint{0.000000in}{0.000000in}}%
\pgfpathlineto{\pgfqpoint{0.000000in}{-0.041667in}}%
\pgfusepath{stroke,fill}%
}%
\begin{pgfscope}%
\pgfsys@transformshift{2.275671in}{2.414488in}%
\pgfsys@useobject{currentmarker}{}%
\end{pgfscope}%
\end{pgfscope}%
\begin{pgfscope}%
\definecolor{textcolor}{rgb}{0.000000,0.000000,0.000000}%
\pgfsetstrokecolor{textcolor}%
\pgfsetfillcolor{textcolor}%
\pgftext[x=2.275671in,y=0.371389in,,top]{\color{textcolor}{\rmfamily\fontsize{11.000000}{13.200000}\selectfont\catcode`\^=\active\def^{\ifmmode\sp\else\^{}\fi}\catcode`\%=\active\def%{\%}$\mathdefault{10^{2}}$}}%
\end{pgfscope}%
\begin{pgfscope}%
\pgfsetbuttcap%
\pgfsetroundjoin%
\definecolor{currentfill}{rgb}{0.000000,0.000000,0.000000}%
\pgfsetfillcolor{currentfill}%
\pgfsetlinewidth{0.501875pt}%
\definecolor{currentstroke}{rgb}{0.000000,0.000000,0.000000}%
\pgfsetstrokecolor{currentstroke}%
\pgfsetdash{}{0pt}%
\pgfsys@defobject{currentmarker}{\pgfqpoint{0.000000in}{0.000000in}}{\pgfqpoint{0.000000in}{0.041667in}}{%
\pgfpathmoveto{\pgfqpoint{0.000000in}{0.000000in}}%
\pgfpathlineto{\pgfqpoint{0.000000in}{0.041667in}}%
\pgfusepath{stroke,fill}%
}%
\begin{pgfscope}%
\pgfsys@transformshift{2.982169in}{0.420000in}%
\pgfsys@useobject{currentmarker}{}%
\end{pgfscope}%
\end{pgfscope}%
\begin{pgfscope}%
\pgfsetbuttcap%
\pgfsetroundjoin%
\definecolor{currentfill}{rgb}{0.000000,0.000000,0.000000}%
\pgfsetfillcolor{currentfill}%
\pgfsetlinewidth{0.501875pt}%
\definecolor{currentstroke}{rgb}{0.000000,0.000000,0.000000}%
\pgfsetstrokecolor{currentstroke}%
\pgfsetdash{}{0pt}%
\pgfsys@defobject{currentmarker}{\pgfqpoint{0.000000in}{-0.041667in}}{\pgfqpoint{0.000000in}{0.000000in}}{%
\pgfpathmoveto{\pgfqpoint{0.000000in}{0.000000in}}%
\pgfpathlineto{\pgfqpoint{0.000000in}{-0.041667in}}%
\pgfusepath{stroke,fill}%
}%
\begin{pgfscope}%
\pgfsys@transformshift{2.982169in}{2.414488in}%
\pgfsys@useobject{currentmarker}{}%
\end{pgfscope}%
\end{pgfscope}%
\begin{pgfscope}%
\definecolor{textcolor}{rgb}{0.000000,0.000000,0.000000}%
\pgfsetstrokecolor{textcolor}%
\pgfsetfillcolor{textcolor}%
\pgftext[x=2.982169in,y=0.371389in,,top]{\color{textcolor}{\rmfamily\fontsize{11.000000}{13.200000}\selectfont\catcode`\^=\active\def^{\ifmmode\sp\else\^{}\fi}\catcode`\%=\active\def%{\%}$\mathdefault{10^{3}}$}}%
\end{pgfscope}%
\begin{pgfscope}%
\pgfsetbuttcap%
\pgfsetroundjoin%
\definecolor{currentfill}{rgb}{0.000000,0.000000,0.000000}%
\pgfsetfillcolor{currentfill}%
\pgfsetlinewidth{0.401500pt}%
\definecolor{currentstroke}{rgb}{0.000000,0.000000,0.000000}%
\pgfsetstrokecolor{currentstroke}%
\pgfsetdash{}{0pt}%
\pgfsys@defobject{currentmarker}{\pgfqpoint{0.000000in}{0.000000in}}{\pgfqpoint{0.000000in}{0.020833in}}{%
\pgfpathmoveto{\pgfqpoint{0.000000in}{0.000000in}}%
\pgfpathlineto{\pgfqpoint{0.000000in}{0.020833in}}%
\pgfusepath{stroke,fill}%
}%
\begin{pgfscope}%
\pgfsys@transformshift{0.650000in}{0.420000in}%
\pgfsys@useobject{currentmarker}{}%
\end{pgfscope}%
\end{pgfscope}%
\begin{pgfscope}%
\pgfsetbuttcap%
\pgfsetroundjoin%
\definecolor{currentfill}{rgb}{0.000000,0.000000,0.000000}%
\pgfsetfillcolor{currentfill}%
\pgfsetlinewidth{0.401500pt}%
\definecolor{currentstroke}{rgb}{0.000000,0.000000,0.000000}%
\pgfsetstrokecolor{currentstroke}%
\pgfsetdash{}{0pt}%
\pgfsys@defobject{currentmarker}{\pgfqpoint{0.000000in}{-0.020833in}}{\pgfqpoint{0.000000in}{0.000000in}}{%
\pgfpathmoveto{\pgfqpoint{0.000000in}{0.000000in}}%
\pgfpathlineto{\pgfqpoint{0.000000in}{-0.020833in}}%
\pgfusepath{stroke,fill}%
}%
\begin{pgfscope}%
\pgfsys@transformshift{0.650000in}{2.414488in}%
\pgfsys@useobject{currentmarker}{}%
\end{pgfscope}%
\end{pgfscope}%
\begin{pgfscope}%
\pgfsetbuttcap%
\pgfsetroundjoin%
\definecolor{currentfill}{rgb}{0.000000,0.000000,0.000000}%
\pgfsetfillcolor{currentfill}%
\pgfsetlinewidth{0.401500pt}%
\definecolor{currentstroke}{rgb}{0.000000,0.000000,0.000000}%
\pgfsetstrokecolor{currentstroke}%
\pgfsetdash{}{0pt}%
\pgfsys@defobject{currentmarker}{\pgfqpoint{0.000000in}{0.000000in}}{\pgfqpoint{0.000000in}{0.020833in}}{%
\pgfpathmoveto{\pgfqpoint{0.000000in}{0.000000in}}%
\pgfpathlineto{\pgfqpoint{0.000000in}{0.020833in}}%
\pgfusepath{stroke,fill}%
}%
\begin{pgfscope}%
\pgfsys@transformshift{0.705941in}{0.420000in}%
\pgfsys@useobject{currentmarker}{}%
\end{pgfscope}%
\end{pgfscope}%
\begin{pgfscope}%
\pgfsetbuttcap%
\pgfsetroundjoin%
\definecolor{currentfill}{rgb}{0.000000,0.000000,0.000000}%
\pgfsetfillcolor{currentfill}%
\pgfsetlinewidth{0.401500pt}%
\definecolor{currentstroke}{rgb}{0.000000,0.000000,0.000000}%
\pgfsetstrokecolor{currentstroke}%
\pgfsetdash{}{0pt}%
\pgfsys@defobject{currentmarker}{\pgfqpoint{0.000000in}{-0.020833in}}{\pgfqpoint{0.000000in}{0.000000in}}{%
\pgfpathmoveto{\pgfqpoint{0.000000in}{0.000000in}}%
\pgfpathlineto{\pgfqpoint{0.000000in}{-0.020833in}}%
\pgfusepath{stroke,fill}%
}%
\begin{pgfscope}%
\pgfsys@transformshift{0.705941in}{2.414488in}%
\pgfsys@useobject{currentmarker}{}%
\end{pgfscope}%
\end{pgfscope}%
\begin{pgfscope}%
\pgfsetbuttcap%
\pgfsetroundjoin%
\definecolor{currentfill}{rgb}{0.000000,0.000000,0.000000}%
\pgfsetfillcolor{currentfill}%
\pgfsetlinewidth{0.401500pt}%
\definecolor{currentstroke}{rgb}{0.000000,0.000000,0.000000}%
\pgfsetstrokecolor{currentstroke}%
\pgfsetdash{}{0pt}%
\pgfsys@defobject{currentmarker}{\pgfqpoint{0.000000in}{0.000000in}}{\pgfqpoint{0.000000in}{0.020833in}}{%
\pgfpathmoveto{\pgfqpoint{0.000000in}{0.000000in}}%
\pgfpathlineto{\pgfqpoint{0.000000in}{0.020833in}}%
\pgfusepath{stroke,fill}%
}%
\begin{pgfscope}%
\pgfsys@transformshift{0.753239in}{0.420000in}%
\pgfsys@useobject{currentmarker}{}%
\end{pgfscope}%
\end{pgfscope}%
\begin{pgfscope}%
\pgfsetbuttcap%
\pgfsetroundjoin%
\definecolor{currentfill}{rgb}{0.000000,0.000000,0.000000}%
\pgfsetfillcolor{currentfill}%
\pgfsetlinewidth{0.401500pt}%
\definecolor{currentstroke}{rgb}{0.000000,0.000000,0.000000}%
\pgfsetstrokecolor{currentstroke}%
\pgfsetdash{}{0pt}%
\pgfsys@defobject{currentmarker}{\pgfqpoint{0.000000in}{-0.020833in}}{\pgfqpoint{0.000000in}{0.000000in}}{%
\pgfpathmoveto{\pgfqpoint{0.000000in}{0.000000in}}%
\pgfpathlineto{\pgfqpoint{0.000000in}{-0.020833in}}%
\pgfusepath{stroke,fill}%
}%
\begin{pgfscope}%
\pgfsys@transformshift{0.753239in}{2.414488in}%
\pgfsys@useobject{currentmarker}{}%
\end{pgfscope}%
\end{pgfscope}%
\begin{pgfscope}%
\pgfsetbuttcap%
\pgfsetroundjoin%
\definecolor{currentfill}{rgb}{0.000000,0.000000,0.000000}%
\pgfsetfillcolor{currentfill}%
\pgfsetlinewidth{0.401500pt}%
\definecolor{currentstroke}{rgb}{0.000000,0.000000,0.000000}%
\pgfsetstrokecolor{currentstroke}%
\pgfsetdash{}{0pt}%
\pgfsys@defobject{currentmarker}{\pgfqpoint{0.000000in}{0.000000in}}{\pgfqpoint{0.000000in}{0.020833in}}{%
\pgfpathmoveto{\pgfqpoint{0.000000in}{0.000000in}}%
\pgfpathlineto{\pgfqpoint{0.000000in}{0.020833in}}%
\pgfusepath{stroke,fill}%
}%
\begin{pgfscope}%
\pgfsys@transformshift{0.794210in}{0.420000in}%
\pgfsys@useobject{currentmarker}{}%
\end{pgfscope}%
\end{pgfscope}%
\begin{pgfscope}%
\pgfsetbuttcap%
\pgfsetroundjoin%
\definecolor{currentfill}{rgb}{0.000000,0.000000,0.000000}%
\pgfsetfillcolor{currentfill}%
\pgfsetlinewidth{0.401500pt}%
\definecolor{currentstroke}{rgb}{0.000000,0.000000,0.000000}%
\pgfsetstrokecolor{currentstroke}%
\pgfsetdash{}{0pt}%
\pgfsys@defobject{currentmarker}{\pgfqpoint{0.000000in}{-0.020833in}}{\pgfqpoint{0.000000in}{0.000000in}}{%
\pgfpathmoveto{\pgfqpoint{0.000000in}{0.000000in}}%
\pgfpathlineto{\pgfqpoint{0.000000in}{-0.020833in}}%
\pgfusepath{stroke,fill}%
}%
\begin{pgfscope}%
\pgfsys@transformshift{0.794210in}{2.414488in}%
\pgfsys@useobject{currentmarker}{}%
\end{pgfscope}%
\end{pgfscope}%
\begin{pgfscope}%
\pgfsetbuttcap%
\pgfsetroundjoin%
\definecolor{currentfill}{rgb}{0.000000,0.000000,0.000000}%
\pgfsetfillcolor{currentfill}%
\pgfsetlinewidth{0.401500pt}%
\definecolor{currentstroke}{rgb}{0.000000,0.000000,0.000000}%
\pgfsetstrokecolor{currentstroke}%
\pgfsetdash{}{0pt}%
\pgfsys@defobject{currentmarker}{\pgfqpoint{0.000000in}{0.000000in}}{\pgfqpoint{0.000000in}{0.020833in}}{%
\pgfpathmoveto{\pgfqpoint{0.000000in}{0.000000in}}%
\pgfpathlineto{\pgfqpoint{0.000000in}{0.020833in}}%
\pgfusepath{stroke,fill}%
}%
\begin{pgfscope}%
\pgfsys@transformshift{0.830349in}{0.420000in}%
\pgfsys@useobject{currentmarker}{}%
\end{pgfscope}%
\end{pgfscope}%
\begin{pgfscope}%
\pgfsetbuttcap%
\pgfsetroundjoin%
\definecolor{currentfill}{rgb}{0.000000,0.000000,0.000000}%
\pgfsetfillcolor{currentfill}%
\pgfsetlinewidth{0.401500pt}%
\definecolor{currentstroke}{rgb}{0.000000,0.000000,0.000000}%
\pgfsetstrokecolor{currentstroke}%
\pgfsetdash{}{0pt}%
\pgfsys@defobject{currentmarker}{\pgfqpoint{0.000000in}{-0.020833in}}{\pgfqpoint{0.000000in}{0.000000in}}{%
\pgfpathmoveto{\pgfqpoint{0.000000in}{0.000000in}}%
\pgfpathlineto{\pgfqpoint{0.000000in}{-0.020833in}}%
\pgfusepath{stroke,fill}%
}%
\begin{pgfscope}%
\pgfsys@transformshift{0.830349in}{2.414488in}%
\pgfsys@useobject{currentmarker}{}%
\end{pgfscope}%
\end{pgfscope}%
\begin{pgfscope}%
\pgfsetbuttcap%
\pgfsetroundjoin%
\definecolor{currentfill}{rgb}{0.000000,0.000000,0.000000}%
\pgfsetfillcolor{currentfill}%
\pgfsetlinewidth{0.401500pt}%
\definecolor{currentstroke}{rgb}{0.000000,0.000000,0.000000}%
\pgfsetstrokecolor{currentstroke}%
\pgfsetdash{}{0pt}%
\pgfsys@defobject{currentmarker}{\pgfqpoint{0.000000in}{0.000000in}}{\pgfqpoint{0.000000in}{0.020833in}}{%
\pgfpathmoveto{\pgfqpoint{0.000000in}{0.000000in}}%
\pgfpathlineto{\pgfqpoint{0.000000in}{0.020833in}}%
\pgfusepath{stroke,fill}%
}%
\begin{pgfscope}%
\pgfsys@transformshift{1.075354in}{0.420000in}%
\pgfsys@useobject{currentmarker}{}%
\end{pgfscope}%
\end{pgfscope}%
\begin{pgfscope}%
\pgfsetbuttcap%
\pgfsetroundjoin%
\definecolor{currentfill}{rgb}{0.000000,0.000000,0.000000}%
\pgfsetfillcolor{currentfill}%
\pgfsetlinewidth{0.401500pt}%
\definecolor{currentstroke}{rgb}{0.000000,0.000000,0.000000}%
\pgfsetstrokecolor{currentstroke}%
\pgfsetdash{}{0pt}%
\pgfsys@defobject{currentmarker}{\pgfqpoint{0.000000in}{-0.020833in}}{\pgfqpoint{0.000000in}{0.000000in}}{%
\pgfpathmoveto{\pgfqpoint{0.000000in}{0.000000in}}%
\pgfpathlineto{\pgfqpoint{0.000000in}{-0.020833in}}%
\pgfusepath{stroke,fill}%
}%
\begin{pgfscope}%
\pgfsys@transformshift{1.075354in}{2.414488in}%
\pgfsys@useobject{currentmarker}{}%
\end{pgfscope}%
\end{pgfscope}%
\begin{pgfscope}%
\pgfsetbuttcap%
\pgfsetroundjoin%
\definecolor{currentfill}{rgb}{0.000000,0.000000,0.000000}%
\pgfsetfillcolor{currentfill}%
\pgfsetlinewidth{0.401500pt}%
\definecolor{currentstroke}{rgb}{0.000000,0.000000,0.000000}%
\pgfsetstrokecolor{currentstroke}%
\pgfsetdash{}{0pt}%
\pgfsys@defobject{currentmarker}{\pgfqpoint{0.000000in}{0.000000in}}{\pgfqpoint{0.000000in}{0.020833in}}{%
\pgfpathmoveto{\pgfqpoint{0.000000in}{0.000000in}}%
\pgfpathlineto{\pgfqpoint{0.000000in}{0.020833in}}%
\pgfusepath{stroke,fill}%
}%
\begin{pgfscope}%
\pgfsys@transformshift{1.199762in}{0.420000in}%
\pgfsys@useobject{currentmarker}{}%
\end{pgfscope}%
\end{pgfscope}%
\begin{pgfscope}%
\pgfsetbuttcap%
\pgfsetroundjoin%
\definecolor{currentfill}{rgb}{0.000000,0.000000,0.000000}%
\pgfsetfillcolor{currentfill}%
\pgfsetlinewidth{0.401500pt}%
\definecolor{currentstroke}{rgb}{0.000000,0.000000,0.000000}%
\pgfsetstrokecolor{currentstroke}%
\pgfsetdash{}{0pt}%
\pgfsys@defobject{currentmarker}{\pgfqpoint{0.000000in}{-0.020833in}}{\pgfqpoint{0.000000in}{0.000000in}}{%
\pgfpathmoveto{\pgfqpoint{0.000000in}{0.000000in}}%
\pgfpathlineto{\pgfqpoint{0.000000in}{-0.020833in}}%
\pgfusepath{stroke,fill}%
}%
\begin{pgfscope}%
\pgfsys@transformshift{1.199762in}{2.414488in}%
\pgfsys@useobject{currentmarker}{}%
\end{pgfscope}%
\end{pgfscope}%
\begin{pgfscope}%
\pgfsetbuttcap%
\pgfsetroundjoin%
\definecolor{currentfill}{rgb}{0.000000,0.000000,0.000000}%
\pgfsetfillcolor{currentfill}%
\pgfsetlinewidth{0.401500pt}%
\definecolor{currentstroke}{rgb}{0.000000,0.000000,0.000000}%
\pgfsetstrokecolor{currentstroke}%
\pgfsetdash{}{0pt}%
\pgfsys@defobject{currentmarker}{\pgfqpoint{0.000000in}{0.000000in}}{\pgfqpoint{0.000000in}{0.020833in}}{%
\pgfpathmoveto{\pgfqpoint{0.000000in}{0.000000in}}%
\pgfpathlineto{\pgfqpoint{0.000000in}{0.020833in}}%
\pgfusepath{stroke,fill}%
}%
\begin{pgfscope}%
\pgfsys@transformshift{1.288031in}{0.420000in}%
\pgfsys@useobject{currentmarker}{}%
\end{pgfscope}%
\end{pgfscope}%
\begin{pgfscope}%
\pgfsetbuttcap%
\pgfsetroundjoin%
\definecolor{currentfill}{rgb}{0.000000,0.000000,0.000000}%
\pgfsetfillcolor{currentfill}%
\pgfsetlinewidth{0.401500pt}%
\definecolor{currentstroke}{rgb}{0.000000,0.000000,0.000000}%
\pgfsetstrokecolor{currentstroke}%
\pgfsetdash{}{0pt}%
\pgfsys@defobject{currentmarker}{\pgfqpoint{0.000000in}{-0.020833in}}{\pgfqpoint{0.000000in}{0.000000in}}{%
\pgfpathmoveto{\pgfqpoint{0.000000in}{0.000000in}}%
\pgfpathlineto{\pgfqpoint{0.000000in}{-0.020833in}}%
\pgfusepath{stroke,fill}%
}%
\begin{pgfscope}%
\pgfsys@transformshift{1.288031in}{2.414488in}%
\pgfsys@useobject{currentmarker}{}%
\end{pgfscope}%
\end{pgfscope}%
\begin{pgfscope}%
\pgfsetbuttcap%
\pgfsetroundjoin%
\definecolor{currentfill}{rgb}{0.000000,0.000000,0.000000}%
\pgfsetfillcolor{currentfill}%
\pgfsetlinewidth{0.401500pt}%
\definecolor{currentstroke}{rgb}{0.000000,0.000000,0.000000}%
\pgfsetstrokecolor{currentstroke}%
\pgfsetdash{}{0pt}%
\pgfsys@defobject{currentmarker}{\pgfqpoint{0.000000in}{0.000000in}}{\pgfqpoint{0.000000in}{0.020833in}}{%
\pgfpathmoveto{\pgfqpoint{0.000000in}{0.000000in}}%
\pgfpathlineto{\pgfqpoint{0.000000in}{0.020833in}}%
\pgfusepath{stroke,fill}%
}%
\begin{pgfscope}%
\pgfsys@transformshift{1.356497in}{0.420000in}%
\pgfsys@useobject{currentmarker}{}%
\end{pgfscope}%
\end{pgfscope}%
\begin{pgfscope}%
\pgfsetbuttcap%
\pgfsetroundjoin%
\definecolor{currentfill}{rgb}{0.000000,0.000000,0.000000}%
\pgfsetfillcolor{currentfill}%
\pgfsetlinewidth{0.401500pt}%
\definecolor{currentstroke}{rgb}{0.000000,0.000000,0.000000}%
\pgfsetstrokecolor{currentstroke}%
\pgfsetdash{}{0pt}%
\pgfsys@defobject{currentmarker}{\pgfqpoint{0.000000in}{-0.020833in}}{\pgfqpoint{0.000000in}{0.000000in}}{%
\pgfpathmoveto{\pgfqpoint{0.000000in}{0.000000in}}%
\pgfpathlineto{\pgfqpoint{0.000000in}{-0.020833in}}%
\pgfusepath{stroke,fill}%
}%
\begin{pgfscope}%
\pgfsys@transformshift{1.356497in}{2.414488in}%
\pgfsys@useobject{currentmarker}{}%
\end{pgfscope}%
\end{pgfscope}%
\begin{pgfscope}%
\pgfsetbuttcap%
\pgfsetroundjoin%
\definecolor{currentfill}{rgb}{0.000000,0.000000,0.000000}%
\pgfsetfillcolor{currentfill}%
\pgfsetlinewidth{0.401500pt}%
\definecolor{currentstroke}{rgb}{0.000000,0.000000,0.000000}%
\pgfsetstrokecolor{currentstroke}%
\pgfsetdash{}{0pt}%
\pgfsys@defobject{currentmarker}{\pgfqpoint{0.000000in}{0.000000in}}{\pgfqpoint{0.000000in}{0.020833in}}{%
\pgfpathmoveto{\pgfqpoint{0.000000in}{0.000000in}}%
\pgfpathlineto{\pgfqpoint{0.000000in}{0.020833in}}%
\pgfusepath{stroke,fill}%
}%
\begin{pgfscope}%
\pgfsys@transformshift{1.412439in}{0.420000in}%
\pgfsys@useobject{currentmarker}{}%
\end{pgfscope}%
\end{pgfscope}%
\begin{pgfscope}%
\pgfsetbuttcap%
\pgfsetroundjoin%
\definecolor{currentfill}{rgb}{0.000000,0.000000,0.000000}%
\pgfsetfillcolor{currentfill}%
\pgfsetlinewidth{0.401500pt}%
\definecolor{currentstroke}{rgb}{0.000000,0.000000,0.000000}%
\pgfsetstrokecolor{currentstroke}%
\pgfsetdash{}{0pt}%
\pgfsys@defobject{currentmarker}{\pgfqpoint{0.000000in}{-0.020833in}}{\pgfqpoint{0.000000in}{0.000000in}}{%
\pgfpathmoveto{\pgfqpoint{0.000000in}{0.000000in}}%
\pgfpathlineto{\pgfqpoint{0.000000in}{-0.020833in}}%
\pgfusepath{stroke,fill}%
}%
\begin{pgfscope}%
\pgfsys@transformshift{1.412439in}{2.414488in}%
\pgfsys@useobject{currentmarker}{}%
\end{pgfscope}%
\end{pgfscope}%
\begin{pgfscope}%
\pgfsetbuttcap%
\pgfsetroundjoin%
\definecolor{currentfill}{rgb}{0.000000,0.000000,0.000000}%
\pgfsetfillcolor{currentfill}%
\pgfsetlinewidth{0.401500pt}%
\definecolor{currentstroke}{rgb}{0.000000,0.000000,0.000000}%
\pgfsetstrokecolor{currentstroke}%
\pgfsetdash{}{0pt}%
\pgfsys@defobject{currentmarker}{\pgfqpoint{0.000000in}{0.000000in}}{\pgfqpoint{0.000000in}{0.020833in}}{%
\pgfpathmoveto{\pgfqpoint{0.000000in}{0.000000in}}%
\pgfpathlineto{\pgfqpoint{0.000000in}{0.020833in}}%
\pgfusepath{stroke,fill}%
}%
\begin{pgfscope}%
\pgfsys@transformshift{1.459736in}{0.420000in}%
\pgfsys@useobject{currentmarker}{}%
\end{pgfscope}%
\end{pgfscope}%
\begin{pgfscope}%
\pgfsetbuttcap%
\pgfsetroundjoin%
\definecolor{currentfill}{rgb}{0.000000,0.000000,0.000000}%
\pgfsetfillcolor{currentfill}%
\pgfsetlinewidth{0.401500pt}%
\definecolor{currentstroke}{rgb}{0.000000,0.000000,0.000000}%
\pgfsetstrokecolor{currentstroke}%
\pgfsetdash{}{0pt}%
\pgfsys@defobject{currentmarker}{\pgfqpoint{0.000000in}{-0.020833in}}{\pgfqpoint{0.000000in}{0.000000in}}{%
\pgfpathmoveto{\pgfqpoint{0.000000in}{0.000000in}}%
\pgfpathlineto{\pgfqpoint{0.000000in}{-0.020833in}}%
\pgfusepath{stroke,fill}%
}%
\begin{pgfscope}%
\pgfsys@transformshift{1.459736in}{2.414488in}%
\pgfsys@useobject{currentmarker}{}%
\end{pgfscope}%
\end{pgfscope}%
\begin{pgfscope}%
\pgfsetbuttcap%
\pgfsetroundjoin%
\definecolor{currentfill}{rgb}{0.000000,0.000000,0.000000}%
\pgfsetfillcolor{currentfill}%
\pgfsetlinewidth{0.401500pt}%
\definecolor{currentstroke}{rgb}{0.000000,0.000000,0.000000}%
\pgfsetstrokecolor{currentstroke}%
\pgfsetdash{}{0pt}%
\pgfsys@defobject{currentmarker}{\pgfqpoint{0.000000in}{0.000000in}}{\pgfqpoint{0.000000in}{0.020833in}}{%
\pgfpathmoveto{\pgfqpoint{0.000000in}{0.000000in}}%
\pgfpathlineto{\pgfqpoint{0.000000in}{0.020833in}}%
\pgfusepath{stroke,fill}%
}%
\begin{pgfscope}%
\pgfsys@transformshift{1.500708in}{0.420000in}%
\pgfsys@useobject{currentmarker}{}%
\end{pgfscope}%
\end{pgfscope}%
\begin{pgfscope}%
\pgfsetbuttcap%
\pgfsetroundjoin%
\definecolor{currentfill}{rgb}{0.000000,0.000000,0.000000}%
\pgfsetfillcolor{currentfill}%
\pgfsetlinewidth{0.401500pt}%
\definecolor{currentstroke}{rgb}{0.000000,0.000000,0.000000}%
\pgfsetstrokecolor{currentstroke}%
\pgfsetdash{}{0pt}%
\pgfsys@defobject{currentmarker}{\pgfqpoint{0.000000in}{-0.020833in}}{\pgfqpoint{0.000000in}{0.000000in}}{%
\pgfpathmoveto{\pgfqpoint{0.000000in}{0.000000in}}%
\pgfpathlineto{\pgfqpoint{0.000000in}{-0.020833in}}%
\pgfusepath{stroke,fill}%
}%
\begin{pgfscope}%
\pgfsys@transformshift{1.500708in}{2.414488in}%
\pgfsys@useobject{currentmarker}{}%
\end{pgfscope}%
\end{pgfscope}%
\begin{pgfscope}%
\pgfsetbuttcap%
\pgfsetroundjoin%
\definecolor{currentfill}{rgb}{0.000000,0.000000,0.000000}%
\pgfsetfillcolor{currentfill}%
\pgfsetlinewidth{0.401500pt}%
\definecolor{currentstroke}{rgb}{0.000000,0.000000,0.000000}%
\pgfsetstrokecolor{currentstroke}%
\pgfsetdash{}{0pt}%
\pgfsys@defobject{currentmarker}{\pgfqpoint{0.000000in}{0.000000in}}{\pgfqpoint{0.000000in}{0.020833in}}{%
\pgfpathmoveto{\pgfqpoint{0.000000in}{0.000000in}}%
\pgfpathlineto{\pgfqpoint{0.000000in}{0.020833in}}%
\pgfusepath{stroke,fill}%
}%
\begin{pgfscope}%
\pgfsys@transformshift{1.536847in}{0.420000in}%
\pgfsys@useobject{currentmarker}{}%
\end{pgfscope}%
\end{pgfscope}%
\begin{pgfscope}%
\pgfsetbuttcap%
\pgfsetroundjoin%
\definecolor{currentfill}{rgb}{0.000000,0.000000,0.000000}%
\pgfsetfillcolor{currentfill}%
\pgfsetlinewidth{0.401500pt}%
\definecolor{currentstroke}{rgb}{0.000000,0.000000,0.000000}%
\pgfsetstrokecolor{currentstroke}%
\pgfsetdash{}{0pt}%
\pgfsys@defobject{currentmarker}{\pgfqpoint{0.000000in}{-0.020833in}}{\pgfqpoint{0.000000in}{0.000000in}}{%
\pgfpathmoveto{\pgfqpoint{0.000000in}{0.000000in}}%
\pgfpathlineto{\pgfqpoint{0.000000in}{-0.020833in}}%
\pgfusepath{stroke,fill}%
}%
\begin{pgfscope}%
\pgfsys@transformshift{1.536847in}{2.414488in}%
\pgfsys@useobject{currentmarker}{}%
\end{pgfscope}%
\end{pgfscope}%
\begin{pgfscope}%
\pgfsetbuttcap%
\pgfsetroundjoin%
\definecolor{currentfill}{rgb}{0.000000,0.000000,0.000000}%
\pgfsetfillcolor{currentfill}%
\pgfsetlinewidth{0.401500pt}%
\definecolor{currentstroke}{rgb}{0.000000,0.000000,0.000000}%
\pgfsetstrokecolor{currentstroke}%
\pgfsetdash{}{0pt}%
\pgfsys@defobject{currentmarker}{\pgfqpoint{0.000000in}{0.000000in}}{\pgfqpoint{0.000000in}{0.020833in}}{%
\pgfpathmoveto{\pgfqpoint{0.000000in}{0.000000in}}%
\pgfpathlineto{\pgfqpoint{0.000000in}{0.020833in}}%
\pgfusepath{stroke,fill}%
}%
\begin{pgfscope}%
\pgfsys@transformshift{1.781851in}{0.420000in}%
\pgfsys@useobject{currentmarker}{}%
\end{pgfscope}%
\end{pgfscope}%
\begin{pgfscope}%
\pgfsetbuttcap%
\pgfsetroundjoin%
\definecolor{currentfill}{rgb}{0.000000,0.000000,0.000000}%
\pgfsetfillcolor{currentfill}%
\pgfsetlinewidth{0.401500pt}%
\definecolor{currentstroke}{rgb}{0.000000,0.000000,0.000000}%
\pgfsetstrokecolor{currentstroke}%
\pgfsetdash{}{0pt}%
\pgfsys@defobject{currentmarker}{\pgfqpoint{0.000000in}{-0.020833in}}{\pgfqpoint{0.000000in}{0.000000in}}{%
\pgfpathmoveto{\pgfqpoint{0.000000in}{0.000000in}}%
\pgfpathlineto{\pgfqpoint{0.000000in}{-0.020833in}}%
\pgfusepath{stroke,fill}%
}%
\begin{pgfscope}%
\pgfsys@transformshift{1.781851in}{2.414488in}%
\pgfsys@useobject{currentmarker}{}%
\end{pgfscope}%
\end{pgfscope}%
\begin{pgfscope}%
\pgfsetbuttcap%
\pgfsetroundjoin%
\definecolor{currentfill}{rgb}{0.000000,0.000000,0.000000}%
\pgfsetfillcolor{currentfill}%
\pgfsetlinewidth{0.401500pt}%
\definecolor{currentstroke}{rgb}{0.000000,0.000000,0.000000}%
\pgfsetstrokecolor{currentstroke}%
\pgfsetdash{}{0pt}%
\pgfsys@defobject{currentmarker}{\pgfqpoint{0.000000in}{0.000000in}}{\pgfqpoint{0.000000in}{0.020833in}}{%
\pgfpathmoveto{\pgfqpoint{0.000000in}{0.000000in}}%
\pgfpathlineto{\pgfqpoint{0.000000in}{0.020833in}}%
\pgfusepath{stroke,fill}%
}%
\begin{pgfscope}%
\pgfsys@transformshift{1.906259in}{0.420000in}%
\pgfsys@useobject{currentmarker}{}%
\end{pgfscope}%
\end{pgfscope}%
\begin{pgfscope}%
\pgfsetbuttcap%
\pgfsetroundjoin%
\definecolor{currentfill}{rgb}{0.000000,0.000000,0.000000}%
\pgfsetfillcolor{currentfill}%
\pgfsetlinewidth{0.401500pt}%
\definecolor{currentstroke}{rgb}{0.000000,0.000000,0.000000}%
\pgfsetstrokecolor{currentstroke}%
\pgfsetdash{}{0pt}%
\pgfsys@defobject{currentmarker}{\pgfqpoint{0.000000in}{-0.020833in}}{\pgfqpoint{0.000000in}{0.000000in}}{%
\pgfpathmoveto{\pgfqpoint{0.000000in}{0.000000in}}%
\pgfpathlineto{\pgfqpoint{0.000000in}{-0.020833in}}%
\pgfusepath{stroke,fill}%
}%
\begin{pgfscope}%
\pgfsys@transformshift{1.906259in}{2.414488in}%
\pgfsys@useobject{currentmarker}{}%
\end{pgfscope}%
\end{pgfscope}%
\begin{pgfscope}%
\pgfsetbuttcap%
\pgfsetroundjoin%
\definecolor{currentfill}{rgb}{0.000000,0.000000,0.000000}%
\pgfsetfillcolor{currentfill}%
\pgfsetlinewidth{0.401500pt}%
\definecolor{currentstroke}{rgb}{0.000000,0.000000,0.000000}%
\pgfsetstrokecolor{currentstroke}%
\pgfsetdash{}{0pt}%
\pgfsys@defobject{currentmarker}{\pgfqpoint{0.000000in}{0.000000in}}{\pgfqpoint{0.000000in}{0.020833in}}{%
\pgfpathmoveto{\pgfqpoint{0.000000in}{0.000000in}}%
\pgfpathlineto{\pgfqpoint{0.000000in}{0.020833in}}%
\pgfusepath{stroke,fill}%
}%
\begin{pgfscope}%
\pgfsys@transformshift{1.994528in}{0.420000in}%
\pgfsys@useobject{currentmarker}{}%
\end{pgfscope}%
\end{pgfscope}%
\begin{pgfscope}%
\pgfsetbuttcap%
\pgfsetroundjoin%
\definecolor{currentfill}{rgb}{0.000000,0.000000,0.000000}%
\pgfsetfillcolor{currentfill}%
\pgfsetlinewidth{0.401500pt}%
\definecolor{currentstroke}{rgb}{0.000000,0.000000,0.000000}%
\pgfsetstrokecolor{currentstroke}%
\pgfsetdash{}{0pt}%
\pgfsys@defobject{currentmarker}{\pgfqpoint{0.000000in}{-0.020833in}}{\pgfqpoint{0.000000in}{0.000000in}}{%
\pgfpathmoveto{\pgfqpoint{0.000000in}{0.000000in}}%
\pgfpathlineto{\pgfqpoint{0.000000in}{-0.020833in}}%
\pgfusepath{stroke,fill}%
}%
\begin{pgfscope}%
\pgfsys@transformshift{1.994528in}{2.414488in}%
\pgfsys@useobject{currentmarker}{}%
\end{pgfscope}%
\end{pgfscope}%
\begin{pgfscope}%
\pgfsetbuttcap%
\pgfsetroundjoin%
\definecolor{currentfill}{rgb}{0.000000,0.000000,0.000000}%
\pgfsetfillcolor{currentfill}%
\pgfsetlinewidth{0.401500pt}%
\definecolor{currentstroke}{rgb}{0.000000,0.000000,0.000000}%
\pgfsetstrokecolor{currentstroke}%
\pgfsetdash{}{0pt}%
\pgfsys@defobject{currentmarker}{\pgfqpoint{0.000000in}{0.000000in}}{\pgfqpoint{0.000000in}{0.020833in}}{%
\pgfpathmoveto{\pgfqpoint{0.000000in}{0.000000in}}%
\pgfpathlineto{\pgfqpoint{0.000000in}{0.020833in}}%
\pgfusepath{stroke,fill}%
}%
\begin{pgfscope}%
\pgfsys@transformshift{2.062995in}{0.420000in}%
\pgfsys@useobject{currentmarker}{}%
\end{pgfscope}%
\end{pgfscope}%
\begin{pgfscope}%
\pgfsetbuttcap%
\pgfsetroundjoin%
\definecolor{currentfill}{rgb}{0.000000,0.000000,0.000000}%
\pgfsetfillcolor{currentfill}%
\pgfsetlinewidth{0.401500pt}%
\definecolor{currentstroke}{rgb}{0.000000,0.000000,0.000000}%
\pgfsetstrokecolor{currentstroke}%
\pgfsetdash{}{0pt}%
\pgfsys@defobject{currentmarker}{\pgfqpoint{0.000000in}{-0.020833in}}{\pgfqpoint{0.000000in}{0.000000in}}{%
\pgfpathmoveto{\pgfqpoint{0.000000in}{0.000000in}}%
\pgfpathlineto{\pgfqpoint{0.000000in}{-0.020833in}}%
\pgfusepath{stroke,fill}%
}%
\begin{pgfscope}%
\pgfsys@transformshift{2.062995in}{2.414488in}%
\pgfsys@useobject{currentmarker}{}%
\end{pgfscope}%
\end{pgfscope}%
\begin{pgfscope}%
\pgfsetbuttcap%
\pgfsetroundjoin%
\definecolor{currentfill}{rgb}{0.000000,0.000000,0.000000}%
\pgfsetfillcolor{currentfill}%
\pgfsetlinewidth{0.401500pt}%
\definecolor{currentstroke}{rgb}{0.000000,0.000000,0.000000}%
\pgfsetstrokecolor{currentstroke}%
\pgfsetdash{}{0pt}%
\pgfsys@defobject{currentmarker}{\pgfqpoint{0.000000in}{0.000000in}}{\pgfqpoint{0.000000in}{0.020833in}}{%
\pgfpathmoveto{\pgfqpoint{0.000000in}{0.000000in}}%
\pgfpathlineto{\pgfqpoint{0.000000in}{0.020833in}}%
\pgfusepath{stroke,fill}%
}%
\begin{pgfscope}%
\pgfsys@transformshift{2.118936in}{0.420000in}%
\pgfsys@useobject{currentmarker}{}%
\end{pgfscope}%
\end{pgfscope}%
\begin{pgfscope}%
\pgfsetbuttcap%
\pgfsetroundjoin%
\definecolor{currentfill}{rgb}{0.000000,0.000000,0.000000}%
\pgfsetfillcolor{currentfill}%
\pgfsetlinewidth{0.401500pt}%
\definecolor{currentstroke}{rgb}{0.000000,0.000000,0.000000}%
\pgfsetstrokecolor{currentstroke}%
\pgfsetdash{}{0pt}%
\pgfsys@defobject{currentmarker}{\pgfqpoint{0.000000in}{-0.020833in}}{\pgfqpoint{0.000000in}{0.000000in}}{%
\pgfpathmoveto{\pgfqpoint{0.000000in}{0.000000in}}%
\pgfpathlineto{\pgfqpoint{0.000000in}{-0.020833in}}%
\pgfusepath{stroke,fill}%
}%
\begin{pgfscope}%
\pgfsys@transformshift{2.118936in}{2.414488in}%
\pgfsys@useobject{currentmarker}{}%
\end{pgfscope}%
\end{pgfscope}%
\begin{pgfscope}%
\pgfsetbuttcap%
\pgfsetroundjoin%
\definecolor{currentfill}{rgb}{0.000000,0.000000,0.000000}%
\pgfsetfillcolor{currentfill}%
\pgfsetlinewidth{0.401500pt}%
\definecolor{currentstroke}{rgb}{0.000000,0.000000,0.000000}%
\pgfsetstrokecolor{currentstroke}%
\pgfsetdash{}{0pt}%
\pgfsys@defobject{currentmarker}{\pgfqpoint{0.000000in}{0.000000in}}{\pgfqpoint{0.000000in}{0.020833in}}{%
\pgfpathmoveto{\pgfqpoint{0.000000in}{0.000000in}}%
\pgfpathlineto{\pgfqpoint{0.000000in}{0.020833in}}%
\pgfusepath{stroke,fill}%
}%
\begin{pgfscope}%
\pgfsys@transformshift{2.166234in}{0.420000in}%
\pgfsys@useobject{currentmarker}{}%
\end{pgfscope}%
\end{pgfscope}%
\begin{pgfscope}%
\pgfsetbuttcap%
\pgfsetroundjoin%
\definecolor{currentfill}{rgb}{0.000000,0.000000,0.000000}%
\pgfsetfillcolor{currentfill}%
\pgfsetlinewidth{0.401500pt}%
\definecolor{currentstroke}{rgb}{0.000000,0.000000,0.000000}%
\pgfsetstrokecolor{currentstroke}%
\pgfsetdash{}{0pt}%
\pgfsys@defobject{currentmarker}{\pgfqpoint{0.000000in}{-0.020833in}}{\pgfqpoint{0.000000in}{0.000000in}}{%
\pgfpathmoveto{\pgfqpoint{0.000000in}{0.000000in}}%
\pgfpathlineto{\pgfqpoint{0.000000in}{-0.020833in}}%
\pgfusepath{stroke,fill}%
}%
\begin{pgfscope}%
\pgfsys@transformshift{2.166234in}{2.414488in}%
\pgfsys@useobject{currentmarker}{}%
\end{pgfscope}%
\end{pgfscope}%
\begin{pgfscope}%
\pgfsetbuttcap%
\pgfsetroundjoin%
\definecolor{currentfill}{rgb}{0.000000,0.000000,0.000000}%
\pgfsetfillcolor{currentfill}%
\pgfsetlinewidth{0.401500pt}%
\definecolor{currentstroke}{rgb}{0.000000,0.000000,0.000000}%
\pgfsetstrokecolor{currentstroke}%
\pgfsetdash{}{0pt}%
\pgfsys@defobject{currentmarker}{\pgfqpoint{0.000000in}{0.000000in}}{\pgfqpoint{0.000000in}{0.020833in}}{%
\pgfpathmoveto{\pgfqpoint{0.000000in}{0.000000in}}%
\pgfpathlineto{\pgfqpoint{0.000000in}{0.020833in}}%
\pgfusepath{stroke,fill}%
}%
\begin{pgfscope}%
\pgfsys@transformshift{2.207205in}{0.420000in}%
\pgfsys@useobject{currentmarker}{}%
\end{pgfscope}%
\end{pgfscope}%
\begin{pgfscope}%
\pgfsetbuttcap%
\pgfsetroundjoin%
\definecolor{currentfill}{rgb}{0.000000,0.000000,0.000000}%
\pgfsetfillcolor{currentfill}%
\pgfsetlinewidth{0.401500pt}%
\definecolor{currentstroke}{rgb}{0.000000,0.000000,0.000000}%
\pgfsetstrokecolor{currentstroke}%
\pgfsetdash{}{0pt}%
\pgfsys@defobject{currentmarker}{\pgfqpoint{0.000000in}{-0.020833in}}{\pgfqpoint{0.000000in}{0.000000in}}{%
\pgfpathmoveto{\pgfqpoint{0.000000in}{0.000000in}}%
\pgfpathlineto{\pgfqpoint{0.000000in}{-0.020833in}}%
\pgfusepath{stroke,fill}%
}%
\begin{pgfscope}%
\pgfsys@transformshift{2.207205in}{2.414488in}%
\pgfsys@useobject{currentmarker}{}%
\end{pgfscope}%
\end{pgfscope}%
\begin{pgfscope}%
\pgfsetbuttcap%
\pgfsetroundjoin%
\definecolor{currentfill}{rgb}{0.000000,0.000000,0.000000}%
\pgfsetfillcolor{currentfill}%
\pgfsetlinewidth{0.401500pt}%
\definecolor{currentstroke}{rgb}{0.000000,0.000000,0.000000}%
\pgfsetstrokecolor{currentstroke}%
\pgfsetdash{}{0pt}%
\pgfsys@defobject{currentmarker}{\pgfqpoint{0.000000in}{0.000000in}}{\pgfqpoint{0.000000in}{0.020833in}}{%
\pgfpathmoveto{\pgfqpoint{0.000000in}{0.000000in}}%
\pgfpathlineto{\pgfqpoint{0.000000in}{0.020833in}}%
\pgfusepath{stroke,fill}%
}%
\begin{pgfscope}%
\pgfsys@transformshift{2.243344in}{0.420000in}%
\pgfsys@useobject{currentmarker}{}%
\end{pgfscope}%
\end{pgfscope}%
\begin{pgfscope}%
\pgfsetbuttcap%
\pgfsetroundjoin%
\definecolor{currentfill}{rgb}{0.000000,0.000000,0.000000}%
\pgfsetfillcolor{currentfill}%
\pgfsetlinewidth{0.401500pt}%
\definecolor{currentstroke}{rgb}{0.000000,0.000000,0.000000}%
\pgfsetstrokecolor{currentstroke}%
\pgfsetdash{}{0pt}%
\pgfsys@defobject{currentmarker}{\pgfqpoint{0.000000in}{-0.020833in}}{\pgfqpoint{0.000000in}{0.000000in}}{%
\pgfpathmoveto{\pgfqpoint{0.000000in}{0.000000in}}%
\pgfpathlineto{\pgfqpoint{0.000000in}{-0.020833in}}%
\pgfusepath{stroke,fill}%
}%
\begin{pgfscope}%
\pgfsys@transformshift{2.243344in}{2.414488in}%
\pgfsys@useobject{currentmarker}{}%
\end{pgfscope}%
\end{pgfscope}%
\begin{pgfscope}%
\pgfsetbuttcap%
\pgfsetroundjoin%
\definecolor{currentfill}{rgb}{0.000000,0.000000,0.000000}%
\pgfsetfillcolor{currentfill}%
\pgfsetlinewidth{0.401500pt}%
\definecolor{currentstroke}{rgb}{0.000000,0.000000,0.000000}%
\pgfsetstrokecolor{currentstroke}%
\pgfsetdash{}{0pt}%
\pgfsys@defobject{currentmarker}{\pgfqpoint{0.000000in}{0.000000in}}{\pgfqpoint{0.000000in}{0.020833in}}{%
\pgfpathmoveto{\pgfqpoint{0.000000in}{0.000000in}}%
\pgfpathlineto{\pgfqpoint{0.000000in}{0.020833in}}%
\pgfusepath{stroke,fill}%
}%
\begin{pgfscope}%
\pgfsys@transformshift{2.488348in}{0.420000in}%
\pgfsys@useobject{currentmarker}{}%
\end{pgfscope}%
\end{pgfscope}%
\begin{pgfscope}%
\pgfsetbuttcap%
\pgfsetroundjoin%
\definecolor{currentfill}{rgb}{0.000000,0.000000,0.000000}%
\pgfsetfillcolor{currentfill}%
\pgfsetlinewidth{0.401500pt}%
\definecolor{currentstroke}{rgb}{0.000000,0.000000,0.000000}%
\pgfsetstrokecolor{currentstroke}%
\pgfsetdash{}{0pt}%
\pgfsys@defobject{currentmarker}{\pgfqpoint{0.000000in}{-0.020833in}}{\pgfqpoint{0.000000in}{0.000000in}}{%
\pgfpathmoveto{\pgfqpoint{0.000000in}{0.000000in}}%
\pgfpathlineto{\pgfqpoint{0.000000in}{-0.020833in}}%
\pgfusepath{stroke,fill}%
}%
\begin{pgfscope}%
\pgfsys@transformshift{2.488348in}{2.414488in}%
\pgfsys@useobject{currentmarker}{}%
\end{pgfscope}%
\end{pgfscope}%
\begin{pgfscope}%
\pgfsetbuttcap%
\pgfsetroundjoin%
\definecolor{currentfill}{rgb}{0.000000,0.000000,0.000000}%
\pgfsetfillcolor{currentfill}%
\pgfsetlinewidth{0.401500pt}%
\definecolor{currentstroke}{rgb}{0.000000,0.000000,0.000000}%
\pgfsetstrokecolor{currentstroke}%
\pgfsetdash{}{0pt}%
\pgfsys@defobject{currentmarker}{\pgfqpoint{0.000000in}{0.000000in}}{\pgfqpoint{0.000000in}{0.020833in}}{%
\pgfpathmoveto{\pgfqpoint{0.000000in}{0.000000in}}%
\pgfpathlineto{\pgfqpoint{0.000000in}{0.020833in}}%
\pgfusepath{stroke,fill}%
}%
\begin{pgfscope}%
\pgfsys@transformshift{2.612756in}{0.420000in}%
\pgfsys@useobject{currentmarker}{}%
\end{pgfscope}%
\end{pgfscope}%
\begin{pgfscope}%
\pgfsetbuttcap%
\pgfsetroundjoin%
\definecolor{currentfill}{rgb}{0.000000,0.000000,0.000000}%
\pgfsetfillcolor{currentfill}%
\pgfsetlinewidth{0.401500pt}%
\definecolor{currentstroke}{rgb}{0.000000,0.000000,0.000000}%
\pgfsetstrokecolor{currentstroke}%
\pgfsetdash{}{0pt}%
\pgfsys@defobject{currentmarker}{\pgfqpoint{0.000000in}{-0.020833in}}{\pgfqpoint{0.000000in}{0.000000in}}{%
\pgfpathmoveto{\pgfqpoint{0.000000in}{0.000000in}}%
\pgfpathlineto{\pgfqpoint{0.000000in}{-0.020833in}}%
\pgfusepath{stroke,fill}%
}%
\begin{pgfscope}%
\pgfsys@transformshift{2.612756in}{2.414488in}%
\pgfsys@useobject{currentmarker}{}%
\end{pgfscope}%
\end{pgfscope}%
\begin{pgfscope}%
\pgfsetbuttcap%
\pgfsetroundjoin%
\definecolor{currentfill}{rgb}{0.000000,0.000000,0.000000}%
\pgfsetfillcolor{currentfill}%
\pgfsetlinewidth{0.401500pt}%
\definecolor{currentstroke}{rgb}{0.000000,0.000000,0.000000}%
\pgfsetstrokecolor{currentstroke}%
\pgfsetdash{}{0pt}%
\pgfsys@defobject{currentmarker}{\pgfqpoint{0.000000in}{0.000000in}}{\pgfqpoint{0.000000in}{0.020833in}}{%
\pgfpathmoveto{\pgfqpoint{0.000000in}{0.000000in}}%
\pgfpathlineto{\pgfqpoint{0.000000in}{0.020833in}}%
\pgfusepath{stroke,fill}%
}%
\begin{pgfscope}%
\pgfsys@transformshift{2.701025in}{0.420000in}%
\pgfsys@useobject{currentmarker}{}%
\end{pgfscope}%
\end{pgfscope}%
\begin{pgfscope}%
\pgfsetbuttcap%
\pgfsetroundjoin%
\definecolor{currentfill}{rgb}{0.000000,0.000000,0.000000}%
\pgfsetfillcolor{currentfill}%
\pgfsetlinewidth{0.401500pt}%
\definecolor{currentstroke}{rgb}{0.000000,0.000000,0.000000}%
\pgfsetstrokecolor{currentstroke}%
\pgfsetdash{}{0pt}%
\pgfsys@defobject{currentmarker}{\pgfqpoint{0.000000in}{-0.020833in}}{\pgfqpoint{0.000000in}{0.000000in}}{%
\pgfpathmoveto{\pgfqpoint{0.000000in}{0.000000in}}%
\pgfpathlineto{\pgfqpoint{0.000000in}{-0.020833in}}%
\pgfusepath{stroke,fill}%
}%
\begin{pgfscope}%
\pgfsys@transformshift{2.701025in}{2.414488in}%
\pgfsys@useobject{currentmarker}{}%
\end{pgfscope}%
\end{pgfscope}%
\begin{pgfscope}%
\pgfsetbuttcap%
\pgfsetroundjoin%
\definecolor{currentfill}{rgb}{0.000000,0.000000,0.000000}%
\pgfsetfillcolor{currentfill}%
\pgfsetlinewidth{0.401500pt}%
\definecolor{currentstroke}{rgb}{0.000000,0.000000,0.000000}%
\pgfsetstrokecolor{currentstroke}%
\pgfsetdash{}{0pt}%
\pgfsys@defobject{currentmarker}{\pgfqpoint{0.000000in}{0.000000in}}{\pgfqpoint{0.000000in}{0.020833in}}{%
\pgfpathmoveto{\pgfqpoint{0.000000in}{0.000000in}}%
\pgfpathlineto{\pgfqpoint{0.000000in}{0.020833in}}%
\pgfusepath{stroke,fill}%
}%
\begin{pgfscope}%
\pgfsys@transformshift{2.769492in}{0.420000in}%
\pgfsys@useobject{currentmarker}{}%
\end{pgfscope}%
\end{pgfscope}%
\begin{pgfscope}%
\pgfsetbuttcap%
\pgfsetroundjoin%
\definecolor{currentfill}{rgb}{0.000000,0.000000,0.000000}%
\pgfsetfillcolor{currentfill}%
\pgfsetlinewidth{0.401500pt}%
\definecolor{currentstroke}{rgb}{0.000000,0.000000,0.000000}%
\pgfsetstrokecolor{currentstroke}%
\pgfsetdash{}{0pt}%
\pgfsys@defobject{currentmarker}{\pgfqpoint{0.000000in}{-0.020833in}}{\pgfqpoint{0.000000in}{0.000000in}}{%
\pgfpathmoveto{\pgfqpoint{0.000000in}{0.000000in}}%
\pgfpathlineto{\pgfqpoint{0.000000in}{-0.020833in}}%
\pgfusepath{stroke,fill}%
}%
\begin{pgfscope}%
\pgfsys@transformshift{2.769492in}{2.414488in}%
\pgfsys@useobject{currentmarker}{}%
\end{pgfscope}%
\end{pgfscope}%
\begin{pgfscope}%
\pgfsetbuttcap%
\pgfsetroundjoin%
\definecolor{currentfill}{rgb}{0.000000,0.000000,0.000000}%
\pgfsetfillcolor{currentfill}%
\pgfsetlinewidth{0.401500pt}%
\definecolor{currentstroke}{rgb}{0.000000,0.000000,0.000000}%
\pgfsetstrokecolor{currentstroke}%
\pgfsetdash{}{0pt}%
\pgfsys@defobject{currentmarker}{\pgfqpoint{0.000000in}{0.000000in}}{\pgfqpoint{0.000000in}{0.020833in}}{%
\pgfpathmoveto{\pgfqpoint{0.000000in}{0.000000in}}%
\pgfpathlineto{\pgfqpoint{0.000000in}{0.020833in}}%
\pgfusepath{stroke,fill}%
}%
\begin{pgfscope}%
\pgfsys@transformshift{2.825433in}{0.420000in}%
\pgfsys@useobject{currentmarker}{}%
\end{pgfscope}%
\end{pgfscope}%
\begin{pgfscope}%
\pgfsetbuttcap%
\pgfsetroundjoin%
\definecolor{currentfill}{rgb}{0.000000,0.000000,0.000000}%
\pgfsetfillcolor{currentfill}%
\pgfsetlinewidth{0.401500pt}%
\definecolor{currentstroke}{rgb}{0.000000,0.000000,0.000000}%
\pgfsetstrokecolor{currentstroke}%
\pgfsetdash{}{0pt}%
\pgfsys@defobject{currentmarker}{\pgfqpoint{0.000000in}{-0.020833in}}{\pgfqpoint{0.000000in}{0.000000in}}{%
\pgfpathmoveto{\pgfqpoint{0.000000in}{0.000000in}}%
\pgfpathlineto{\pgfqpoint{0.000000in}{-0.020833in}}%
\pgfusepath{stroke,fill}%
}%
\begin{pgfscope}%
\pgfsys@transformshift{2.825433in}{2.414488in}%
\pgfsys@useobject{currentmarker}{}%
\end{pgfscope}%
\end{pgfscope}%
\begin{pgfscope}%
\pgfsetbuttcap%
\pgfsetroundjoin%
\definecolor{currentfill}{rgb}{0.000000,0.000000,0.000000}%
\pgfsetfillcolor{currentfill}%
\pgfsetlinewidth{0.401500pt}%
\definecolor{currentstroke}{rgb}{0.000000,0.000000,0.000000}%
\pgfsetstrokecolor{currentstroke}%
\pgfsetdash{}{0pt}%
\pgfsys@defobject{currentmarker}{\pgfqpoint{0.000000in}{0.000000in}}{\pgfqpoint{0.000000in}{0.020833in}}{%
\pgfpathmoveto{\pgfqpoint{0.000000in}{0.000000in}}%
\pgfpathlineto{\pgfqpoint{0.000000in}{0.020833in}}%
\pgfusepath{stroke,fill}%
}%
\begin{pgfscope}%
\pgfsys@transformshift{2.872731in}{0.420000in}%
\pgfsys@useobject{currentmarker}{}%
\end{pgfscope}%
\end{pgfscope}%
\begin{pgfscope}%
\pgfsetbuttcap%
\pgfsetroundjoin%
\definecolor{currentfill}{rgb}{0.000000,0.000000,0.000000}%
\pgfsetfillcolor{currentfill}%
\pgfsetlinewidth{0.401500pt}%
\definecolor{currentstroke}{rgb}{0.000000,0.000000,0.000000}%
\pgfsetstrokecolor{currentstroke}%
\pgfsetdash{}{0pt}%
\pgfsys@defobject{currentmarker}{\pgfqpoint{0.000000in}{-0.020833in}}{\pgfqpoint{0.000000in}{0.000000in}}{%
\pgfpathmoveto{\pgfqpoint{0.000000in}{0.000000in}}%
\pgfpathlineto{\pgfqpoint{0.000000in}{-0.020833in}}%
\pgfusepath{stroke,fill}%
}%
\begin{pgfscope}%
\pgfsys@transformshift{2.872731in}{2.414488in}%
\pgfsys@useobject{currentmarker}{}%
\end{pgfscope}%
\end{pgfscope}%
\begin{pgfscope}%
\pgfsetbuttcap%
\pgfsetroundjoin%
\definecolor{currentfill}{rgb}{0.000000,0.000000,0.000000}%
\pgfsetfillcolor{currentfill}%
\pgfsetlinewidth{0.401500pt}%
\definecolor{currentstroke}{rgb}{0.000000,0.000000,0.000000}%
\pgfsetstrokecolor{currentstroke}%
\pgfsetdash{}{0pt}%
\pgfsys@defobject{currentmarker}{\pgfqpoint{0.000000in}{0.000000in}}{\pgfqpoint{0.000000in}{0.020833in}}{%
\pgfpathmoveto{\pgfqpoint{0.000000in}{0.000000in}}%
\pgfpathlineto{\pgfqpoint{0.000000in}{0.020833in}}%
\pgfusepath{stroke,fill}%
}%
\begin{pgfscope}%
\pgfsys@transformshift{2.913702in}{0.420000in}%
\pgfsys@useobject{currentmarker}{}%
\end{pgfscope}%
\end{pgfscope}%
\begin{pgfscope}%
\pgfsetbuttcap%
\pgfsetroundjoin%
\definecolor{currentfill}{rgb}{0.000000,0.000000,0.000000}%
\pgfsetfillcolor{currentfill}%
\pgfsetlinewidth{0.401500pt}%
\definecolor{currentstroke}{rgb}{0.000000,0.000000,0.000000}%
\pgfsetstrokecolor{currentstroke}%
\pgfsetdash{}{0pt}%
\pgfsys@defobject{currentmarker}{\pgfqpoint{0.000000in}{-0.020833in}}{\pgfqpoint{0.000000in}{0.000000in}}{%
\pgfpathmoveto{\pgfqpoint{0.000000in}{0.000000in}}%
\pgfpathlineto{\pgfqpoint{0.000000in}{-0.020833in}}%
\pgfusepath{stroke,fill}%
}%
\begin{pgfscope}%
\pgfsys@transformshift{2.913702in}{2.414488in}%
\pgfsys@useobject{currentmarker}{}%
\end{pgfscope}%
\end{pgfscope}%
\begin{pgfscope}%
\pgfsetbuttcap%
\pgfsetroundjoin%
\definecolor{currentfill}{rgb}{0.000000,0.000000,0.000000}%
\pgfsetfillcolor{currentfill}%
\pgfsetlinewidth{0.401500pt}%
\definecolor{currentstroke}{rgb}{0.000000,0.000000,0.000000}%
\pgfsetstrokecolor{currentstroke}%
\pgfsetdash{}{0pt}%
\pgfsys@defobject{currentmarker}{\pgfqpoint{0.000000in}{0.000000in}}{\pgfqpoint{0.000000in}{0.020833in}}{%
\pgfpathmoveto{\pgfqpoint{0.000000in}{0.000000in}}%
\pgfpathlineto{\pgfqpoint{0.000000in}{0.020833in}}%
\pgfusepath{stroke,fill}%
}%
\begin{pgfscope}%
\pgfsys@transformshift{2.949841in}{0.420000in}%
\pgfsys@useobject{currentmarker}{}%
\end{pgfscope}%
\end{pgfscope}%
\begin{pgfscope}%
\pgfsetbuttcap%
\pgfsetroundjoin%
\definecolor{currentfill}{rgb}{0.000000,0.000000,0.000000}%
\pgfsetfillcolor{currentfill}%
\pgfsetlinewidth{0.401500pt}%
\definecolor{currentstroke}{rgb}{0.000000,0.000000,0.000000}%
\pgfsetstrokecolor{currentstroke}%
\pgfsetdash{}{0pt}%
\pgfsys@defobject{currentmarker}{\pgfqpoint{0.000000in}{-0.020833in}}{\pgfqpoint{0.000000in}{0.000000in}}{%
\pgfpathmoveto{\pgfqpoint{0.000000in}{0.000000in}}%
\pgfpathlineto{\pgfqpoint{0.000000in}{-0.020833in}}%
\pgfusepath{stroke,fill}%
}%
\begin{pgfscope}%
\pgfsys@transformshift{2.949841in}{2.414488in}%
\pgfsys@useobject{currentmarker}{}%
\end{pgfscope}%
\end{pgfscope}%
\begin{pgfscope}%
\definecolor{textcolor}{rgb}{0.000000,0.000000,0.000000}%
\pgfsetstrokecolor{textcolor}%
\pgfsetfillcolor{textcolor}%
\pgftext[x=1.844055in,y=0.208426in,,top]{\color{textcolor}{\rmfamily\fontsize{12.000000}{14.400000}\selectfont\catcode`\^=\active\def^{\ifmmode\sp\else\^{}\fi}\catcode`\%=\active\def%{\%}$y^+$}}%
\end{pgfscope}%
\begin{pgfscope}%
\pgfsetbuttcap%
\pgfsetroundjoin%
\definecolor{currentfill}{rgb}{0.000000,0.000000,0.000000}%
\pgfsetfillcolor{currentfill}%
\pgfsetlinewidth{0.501875pt}%
\definecolor{currentstroke}{rgb}{0.000000,0.000000,0.000000}%
\pgfsetstrokecolor{currentstroke}%
\pgfsetdash{}{0pt}%
\pgfsys@defobject{currentmarker}{\pgfqpoint{0.000000in}{0.000000in}}{\pgfqpoint{0.041667in}{0.000000in}}{%
\pgfpathmoveto{\pgfqpoint{0.000000in}{0.000000in}}%
\pgfpathlineto{\pgfqpoint{0.041667in}{0.000000in}}%
\pgfusepath{stroke,fill}%
}%
\begin{pgfscope}%
\pgfsys@transformshift{0.650000in}{0.420000in}%
\pgfsys@useobject{currentmarker}{}%
\end{pgfscope}%
\end{pgfscope}%
\begin{pgfscope}%
\pgfsetbuttcap%
\pgfsetroundjoin%
\definecolor{currentfill}{rgb}{0.000000,0.000000,0.000000}%
\pgfsetfillcolor{currentfill}%
\pgfsetlinewidth{0.501875pt}%
\definecolor{currentstroke}{rgb}{0.000000,0.000000,0.000000}%
\pgfsetstrokecolor{currentstroke}%
\pgfsetdash{}{0pt}%
\pgfsys@defobject{currentmarker}{\pgfqpoint{-0.041667in}{0.000000in}}{\pgfqpoint{-0.000000in}{0.000000in}}{%
\pgfpathmoveto{\pgfqpoint{-0.000000in}{0.000000in}}%
\pgfpathlineto{\pgfqpoint{-0.041667in}{0.000000in}}%
\pgfusepath{stroke,fill}%
}%
\begin{pgfscope}%
\pgfsys@transformshift{3.038110in}{0.420000in}%
\pgfsys@useobject{currentmarker}{}%
\end{pgfscope}%
\end{pgfscope}%
\begin{pgfscope}%
\definecolor{textcolor}{rgb}{0.000000,0.000000,0.000000}%
\pgfsetstrokecolor{textcolor}%
\pgfsetfillcolor{textcolor}%
\pgftext[x=0.525347in, y=0.367193in, left, base]{\color{textcolor}{\rmfamily\fontsize{11.000000}{13.200000}\selectfont\catcode`\^=\active\def^{\ifmmode\sp\else\^{}\fi}\catcode`\%=\active\def%{\%}0}}%
\end{pgfscope}%
\begin{pgfscope}%
\pgfsetbuttcap%
\pgfsetroundjoin%
\definecolor{currentfill}{rgb}{0.000000,0.000000,0.000000}%
\pgfsetfillcolor{currentfill}%
\pgfsetlinewidth{0.501875pt}%
\definecolor{currentstroke}{rgb}{0.000000,0.000000,0.000000}%
\pgfsetstrokecolor{currentstroke}%
\pgfsetdash{}{0pt}%
\pgfsys@defobject{currentmarker}{\pgfqpoint{0.000000in}{0.000000in}}{\pgfqpoint{0.041667in}{0.000000in}}{%
\pgfpathmoveto{\pgfqpoint{0.000000in}{0.000000in}}%
\pgfpathlineto{\pgfqpoint{0.041667in}{0.000000in}}%
\pgfusepath{stroke,fill}%
}%
\begin{pgfscope}%
\pgfsys@transformshift{0.650000in}{0.818898in}%
\pgfsys@useobject{currentmarker}{}%
\end{pgfscope}%
\end{pgfscope}%
\begin{pgfscope}%
\pgfsetbuttcap%
\pgfsetroundjoin%
\definecolor{currentfill}{rgb}{0.000000,0.000000,0.000000}%
\pgfsetfillcolor{currentfill}%
\pgfsetlinewidth{0.501875pt}%
\definecolor{currentstroke}{rgb}{0.000000,0.000000,0.000000}%
\pgfsetstrokecolor{currentstroke}%
\pgfsetdash{}{0pt}%
\pgfsys@defobject{currentmarker}{\pgfqpoint{-0.041667in}{0.000000in}}{\pgfqpoint{-0.000000in}{0.000000in}}{%
\pgfpathmoveto{\pgfqpoint{-0.000000in}{0.000000in}}%
\pgfpathlineto{\pgfqpoint{-0.041667in}{0.000000in}}%
\pgfusepath{stroke,fill}%
}%
\begin{pgfscope}%
\pgfsys@transformshift{3.038110in}{0.818898in}%
\pgfsys@useobject{currentmarker}{}%
\end{pgfscope}%
\end{pgfscope}%
\begin{pgfscope}%
\definecolor{textcolor}{rgb}{0.000000,0.000000,0.000000}%
\pgfsetstrokecolor{textcolor}%
\pgfsetfillcolor{textcolor}%
\pgftext[x=0.525347in, y=0.766091in, left, base]{\color{textcolor}{\rmfamily\fontsize{11.000000}{13.200000}\selectfont\catcode`\^=\active\def^{\ifmmode\sp\else\^{}\fi}\catcode`\%=\active\def%{\%}5}}%
\end{pgfscope}%
\begin{pgfscope}%
\pgfsetbuttcap%
\pgfsetroundjoin%
\definecolor{currentfill}{rgb}{0.000000,0.000000,0.000000}%
\pgfsetfillcolor{currentfill}%
\pgfsetlinewidth{0.501875pt}%
\definecolor{currentstroke}{rgb}{0.000000,0.000000,0.000000}%
\pgfsetstrokecolor{currentstroke}%
\pgfsetdash{}{0pt}%
\pgfsys@defobject{currentmarker}{\pgfqpoint{0.000000in}{0.000000in}}{\pgfqpoint{0.041667in}{0.000000in}}{%
\pgfpathmoveto{\pgfqpoint{0.000000in}{0.000000in}}%
\pgfpathlineto{\pgfqpoint{0.041667in}{0.000000in}}%
\pgfusepath{stroke,fill}%
}%
\begin{pgfscope}%
\pgfsys@transformshift{0.650000in}{1.217795in}%
\pgfsys@useobject{currentmarker}{}%
\end{pgfscope}%
\end{pgfscope}%
\begin{pgfscope}%
\pgfsetbuttcap%
\pgfsetroundjoin%
\definecolor{currentfill}{rgb}{0.000000,0.000000,0.000000}%
\pgfsetfillcolor{currentfill}%
\pgfsetlinewidth{0.501875pt}%
\definecolor{currentstroke}{rgb}{0.000000,0.000000,0.000000}%
\pgfsetstrokecolor{currentstroke}%
\pgfsetdash{}{0pt}%
\pgfsys@defobject{currentmarker}{\pgfqpoint{-0.041667in}{0.000000in}}{\pgfqpoint{-0.000000in}{0.000000in}}{%
\pgfpathmoveto{\pgfqpoint{-0.000000in}{0.000000in}}%
\pgfpathlineto{\pgfqpoint{-0.041667in}{0.000000in}}%
\pgfusepath{stroke,fill}%
}%
\begin{pgfscope}%
\pgfsys@transformshift{3.038110in}{1.217795in}%
\pgfsys@useobject{currentmarker}{}%
\end{pgfscope}%
\end{pgfscope}%
\begin{pgfscope}%
\definecolor{textcolor}{rgb}{0.000000,0.000000,0.000000}%
\pgfsetstrokecolor{textcolor}%
\pgfsetfillcolor{textcolor}%
\pgftext[x=0.449305in, y=1.164989in, left, base]{\color{textcolor}{\rmfamily\fontsize{11.000000}{13.200000}\selectfont\catcode`\^=\active\def^{\ifmmode\sp\else\^{}\fi}\catcode`\%=\active\def%{\%}10}}%
\end{pgfscope}%
\begin{pgfscope}%
\pgfsetbuttcap%
\pgfsetroundjoin%
\definecolor{currentfill}{rgb}{0.000000,0.000000,0.000000}%
\pgfsetfillcolor{currentfill}%
\pgfsetlinewidth{0.501875pt}%
\definecolor{currentstroke}{rgb}{0.000000,0.000000,0.000000}%
\pgfsetstrokecolor{currentstroke}%
\pgfsetdash{}{0pt}%
\pgfsys@defobject{currentmarker}{\pgfqpoint{0.000000in}{0.000000in}}{\pgfqpoint{0.041667in}{0.000000in}}{%
\pgfpathmoveto{\pgfqpoint{0.000000in}{0.000000in}}%
\pgfpathlineto{\pgfqpoint{0.041667in}{0.000000in}}%
\pgfusepath{stroke,fill}%
}%
\begin{pgfscope}%
\pgfsys@transformshift{0.650000in}{1.616693in}%
\pgfsys@useobject{currentmarker}{}%
\end{pgfscope}%
\end{pgfscope}%
\begin{pgfscope}%
\pgfsetbuttcap%
\pgfsetroundjoin%
\definecolor{currentfill}{rgb}{0.000000,0.000000,0.000000}%
\pgfsetfillcolor{currentfill}%
\pgfsetlinewidth{0.501875pt}%
\definecolor{currentstroke}{rgb}{0.000000,0.000000,0.000000}%
\pgfsetstrokecolor{currentstroke}%
\pgfsetdash{}{0pt}%
\pgfsys@defobject{currentmarker}{\pgfqpoint{-0.041667in}{0.000000in}}{\pgfqpoint{-0.000000in}{0.000000in}}{%
\pgfpathmoveto{\pgfqpoint{-0.000000in}{0.000000in}}%
\pgfpathlineto{\pgfqpoint{-0.041667in}{0.000000in}}%
\pgfusepath{stroke,fill}%
}%
\begin{pgfscope}%
\pgfsys@transformshift{3.038110in}{1.616693in}%
\pgfsys@useobject{currentmarker}{}%
\end{pgfscope}%
\end{pgfscope}%
\begin{pgfscope}%
\definecolor{textcolor}{rgb}{0.000000,0.000000,0.000000}%
\pgfsetstrokecolor{textcolor}%
\pgfsetfillcolor{textcolor}%
\pgftext[x=0.449305in, y=1.563886in, left, base]{\color{textcolor}{\rmfamily\fontsize{11.000000}{13.200000}\selectfont\catcode`\^=\active\def^{\ifmmode\sp\else\^{}\fi}\catcode`\%=\active\def%{\%}15}}%
\end{pgfscope}%
\begin{pgfscope}%
\pgfsetbuttcap%
\pgfsetroundjoin%
\definecolor{currentfill}{rgb}{0.000000,0.000000,0.000000}%
\pgfsetfillcolor{currentfill}%
\pgfsetlinewidth{0.501875pt}%
\definecolor{currentstroke}{rgb}{0.000000,0.000000,0.000000}%
\pgfsetstrokecolor{currentstroke}%
\pgfsetdash{}{0pt}%
\pgfsys@defobject{currentmarker}{\pgfqpoint{0.000000in}{0.000000in}}{\pgfqpoint{0.041667in}{0.000000in}}{%
\pgfpathmoveto{\pgfqpoint{0.000000in}{0.000000in}}%
\pgfpathlineto{\pgfqpoint{0.041667in}{0.000000in}}%
\pgfusepath{stroke,fill}%
}%
\begin{pgfscope}%
\pgfsys@transformshift{0.650000in}{2.015590in}%
\pgfsys@useobject{currentmarker}{}%
\end{pgfscope}%
\end{pgfscope}%
\begin{pgfscope}%
\pgfsetbuttcap%
\pgfsetroundjoin%
\definecolor{currentfill}{rgb}{0.000000,0.000000,0.000000}%
\pgfsetfillcolor{currentfill}%
\pgfsetlinewidth{0.501875pt}%
\definecolor{currentstroke}{rgb}{0.000000,0.000000,0.000000}%
\pgfsetstrokecolor{currentstroke}%
\pgfsetdash{}{0pt}%
\pgfsys@defobject{currentmarker}{\pgfqpoint{-0.041667in}{0.000000in}}{\pgfqpoint{-0.000000in}{0.000000in}}{%
\pgfpathmoveto{\pgfqpoint{-0.000000in}{0.000000in}}%
\pgfpathlineto{\pgfqpoint{-0.041667in}{0.000000in}}%
\pgfusepath{stroke,fill}%
}%
\begin{pgfscope}%
\pgfsys@transformshift{3.038110in}{2.015590in}%
\pgfsys@useobject{currentmarker}{}%
\end{pgfscope}%
\end{pgfscope}%
\begin{pgfscope}%
\definecolor{textcolor}{rgb}{0.000000,0.000000,0.000000}%
\pgfsetstrokecolor{textcolor}%
\pgfsetfillcolor{textcolor}%
\pgftext[x=0.449305in, y=1.962784in, left, base]{\color{textcolor}{\rmfamily\fontsize{11.000000}{13.200000}\selectfont\catcode`\^=\active\def^{\ifmmode\sp\else\^{}\fi}\catcode`\%=\active\def%{\%}20}}%
\end{pgfscope}%
\begin{pgfscope}%
\pgfsetbuttcap%
\pgfsetroundjoin%
\definecolor{currentfill}{rgb}{0.000000,0.000000,0.000000}%
\pgfsetfillcolor{currentfill}%
\pgfsetlinewidth{0.501875pt}%
\definecolor{currentstroke}{rgb}{0.000000,0.000000,0.000000}%
\pgfsetstrokecolor{currentstroke}%
\pgfsetdash{}{0pt}%
\pgfsys@defobject{currentmarker}{\pgfqpoint{0.000000in}{0.000000in}}{\pgfqpoint{0.041667in}{0.000000in}}{%
\pgfpathmoveto{\pgfqpoint{0.000000in}{0.000000in}}%
\pgfpathlineto{\pgfqpoint{0.041667in}{0.000000in}}%
\pgfusepath{stroke,fill}%
}%
\begin{pgfscope}%
\pgfsys@transformshift{0.650000in}{2.414488in}%
\pgfsys@useobject{currentmarker}{}%
\end{pgfscope}%
\end{pgfscope}%
\begin{pgfscope}%
\pgfsetbuttcap%
\pgfsetroundjoin%
\definecolor{currentfill}{rgb}{0.000000,0.000000,0.000000}%
\pgfsetfillcolor{currentfill}%
\pgfsetlinewidth{0.501875pt}%
\definecolor{currentstroke}{rgb}{0.000000,0.000000,0.000000}%
\pgfsetstrokecolor{currentstroke}%
\pgfsetdash{}{0pt}%
\pgfsys@defobject{currentmarker}{\pgfqpoint{-0.041667in}{0.000000in}}{\pgfqpoint{-0.000000in}{0.000000in}}{%
\pgfpathmoveto{\pgfqpoint{-0.000000in}{0.000000in}}%
\pgfpathlineto{\pgfqpoint{-0.041667in}{0.000000in}}%
\pgfusepath{stroke,fill}%
}%
\begin{pgfscope}%
\pgfsys@transformshift{3.038110in}{2.414488in}%
\pgfsys@useobject{currentmarker}{}%
\end{pgfscope}%
\end{pgfscope}%
\begin{pgfscope}%
\definecolor{textcolor}{rgb}{0.000000,0.000000,0.000000}%
\pgfsetstrokecolor{textcolor}%
\pgfsetfillcolor{textcolor}%
\pgftext[x=0.449305in, y=2.361681in, left, base]{\color{textcolor}{\rmfamily\fontsize{11.000000}{13.200000}\selectfont\catcode`\^=\active\def^{\ifmmode\sp\else\^{}\fi}\catcode`\%=\active\def%{\%}25}}%
\end{pgfscope}%
\begin{pgfscope}%
\pgfsetbuttcap%
\pgfsetroundjoin%
\definecolor{currentfill}{rgb}{0.000000,0.000000,0.000000}%
\pgfsetfillcolor{currentfill}%
\pgfsetlinewidth{0.401500pt}%
\definecolor{currentstroke}{rgb}{0.000000,0.000000,0.000000}%
\pgfsetstrokecolor{currentstroke}%
\pgfsetdash{}{0pt}%
\pgfsys@defobject{currentmarker}{\pgfqpoint{0.000000in}{0.000000in}}{\pgfqpoint{0.020833in}{0.000000in}}{%
\pgfpathmoveto{\pgfqpoint{0.000000in}{0.000000in}}%
\pgfpathlineto{\pgfqpoint{0.020833in}{0.000000in}}%
\pgfusepath{stroke,fill}%
}%
\begin{pgfscope}%
\pgfsys@transformshift{0.650000in}{0.499780in}%
\pgfsys@useobject{currentmarker}{}%
\end{pgfscope}%
\end{pgfscope}%
\begin{pgfscope}%
\pgfsetbuttcap%
\pgfsetroundjoin%
\definecolor{currentfill}{rgb}{0.000000,0.000000,0.000000}%
\pgfsetfillcolor{currentfill}%
\pgfsetlinewidth{0.401500pt}%
\definecolor{currentstroke}{rgb}{0.000000,0.000000,0.000000}%
\pgfsetstrokecolor{currentstroke}%
\pgfsetdash{}{0pt}%
\pgfsys@defobject{currentmarker}{\pgfqpoint{-0.020833in}{0.000000in}}{\pgfqpoint{-0.000000in}{0.000000in}}{%
\pgfpathmoveto{\pgfqpoint{-0.000000in}{0.000000in}}%
\pgfpathlineto{\pgfqpoint{-0.020833in}{0.000000in}}%
\pgfusepath{stroke,fill}%
}%
\begin{pgfscope}%
\pgfsys@transformshift{3.038110in}{0.499780in}%
\pgfsys@useobject{currentmarker}{}%
\end{pgfscope}%
\end{pgfscope}%
\begin{pgfscope}%
\pgfsetbuttcap%
\pgfsetroundjoin%
\definecolor{currentfill}{rgb}{0.000000,0.000000,0.000000}%
\pgfsetfillcolor{currentfill}%
\pgfsetlinewidth{0.401500pt}%
\definecolor{currentstroke}{rgb}{0.000000,0.000000,0.000000}%
\pgfsetstrokecolor{currentstroke}%
\pgfsetdash{}{0pt}%
\pgfsys@defobject{currentmarker}{\pgfqpoint{0.000000in}{0.000000in}}{\pgfqpoint{0.020833in}{0.000000in}}{%
\pgfpathmoveto{\pgfqpoint{0.000000in}{0.000000in}}%
\pgfpathlineto{\pgfqpoint{0.020833in}{0.000000in}}%
\pgfusepath{stroke,fill}%
}%
\begin{pgfscope}%
\pgfsys@transformshift{0.650000in}{0.579559in}%
\pgfsys@useobject{currentmarker}{}%
\end{pgfscope}%
\end{pgfscope}%
\begin{pgfscope}%
\pgfsetbuttcap%
\pgfsetroundjoin%
\definecolor{currentfill}{rgb}{0.000000,0.000000,0.000000}%
\pgfsetfillcolor{currentfill}%
\pgfsetlinewidth{0.401500pt}%
\definecolor{currentstroke}{rgb}{0.000000,0.000000,0.000000}%
\pgfsetstrokecolor{currentstroke}%
\pgfsetdash{}{0pt}%
\pgfsys@defobject{currentmarker}{\pgfqpoint{-0.020833in}{0.000000in}}{\pgfqpoint{-0.000000in}{0.000000in}}{%
\pgfpathmoveto{\pgfqpoint{-0.000000in}{0.000000in}}%
\pgfpathlineto{\pgfqpoint{-0.020833in}{0.000000in}}%
\pgfusepath{stroke,fill}%
}%
\begin{pgfscope}%
\pgfsys@transformshift{3.038110in}{0.579559in}%
\pgfsys@useobject{currentmarker}{}%
\end{pgfscope}%
\end{pgfscope}%
\begin{pgfscope}%
\pgfsetbuttcap%
\pgfsetroundjoin%
\definecolor{currentfill}{rgb}{0.000000,0.000000,0.000000}%
\pgfsetfillcolor{currentfill}%
\pgfsetlinewidth{0.401500pt}%
\definecolor{currentstroke}{rgb}{0.000000,0.000000,0.000000}%
\pgfsetstrokecolor{currentstroke}%
\pgfsetdash{}{0pt}%
\pgfsys@defobject{currentmarker}{\pgfqpoint{0.000000in}{0.000000in}}{\pgfqpoint{0.020833in}{0.000000in}}{%
\pgfpathmoveto{\pgfqpoint{0.000000in}{0.000000in}}%
\pgfpathlineto{\pgfqpoint{0.020833in}{0.000000in}}%
\pgfusepath{stroke,fill}%
}%
\begin{pgfscope}%
\pgfsys@transformshift{0.650000in}{0.659339in}%
\pgfsys@useobject{currentmarker}{}%
\end{pgfscope}%
\end{pgfscope}%
\begin{pgfscope}%
\pgfsetbuttcap%
\pgfsetroundjoin%
\definecolor{currentfill}{rgb}{0.000000,0.000000,0.000000}%
\pgfsetfillcolor{currentfill}%
\pgfsetlinewidth{0.401500pt}%
\definecolor{currentstroke}{rgb}{0.000000,0.000000,0.000000}%
\pgfsetstrokecolor{currentstroke}%
\pgfsetdash{}{0pt}%
\pgfsys@defobject{currentmarker}{\pgfqpoint{-0.020833in}{0.000000in}}{\pgfqpoint{-0.000000in}{0.000000in}}{%
\pgfpathmoveto{\pgfqpoint{-0.000000in}{0.000000in}}%
\pgfpathlineto{\pgfqpoint{-0.020833in}{0.000000in}}%
\pgfusepath{stroke,fill}%
}%
\begin{pgfscope}%
\pgfsys@transformshift{3.038110in}{0.659339in}%
\pgfsys@useobject{currentmarker}{}%
\end{pgfscope}%
\end{pgfscope}%
\begin{pgfscope}%
\pgfsetbuttcap%
\pgfsetroundjoin%
\definecolor{currentfill}{rgb}{0.000000,0.000000,0.000000}%
\pgfsetfillcolor{currentfill}%
\pgfsetlinewidth{0.401500pt}%
\definecolor{currentstroke}{rgb}{0.000000,0.000000,0.000000}%
\pgfsetstrokecolor{currentstroke}%
\pgfsetdash{}{0pt}%
\pgfsys@defobject{currentmarker}{\pgfqpoint{0.000000in}{0.000000in}}{\pgfqpoint{0.020833in}{0.000000in}}{%
\pgfpathmoveto{\pgfqpoint{0.000000in}{0.000000in}}%
\pgfpathlineto{\pgfqpoint{0.020833in}{0.000000in}}%
\pgfusepath{stroke,fill}%
}%
\begin{pgfscope}%
\pgfsys@transformshift{0.650000in}{0.739118in}%
\pgfsys@useobject{currentmarker}{}%
\end{pgfscope}%
\end{pgfscope}%
\begin{pgfscope}%
\pgfsetbuttcap%
\pgfsetroundjoin%
\definecolor{currentfill}{rgb}{0.000000,0.000000,0.000000}%
\pgfsetfillcolor{currentfill}%
\pgfsetlinewidth{0.401500pt}%
\definecolor{currentstroke}{rgb}{0.000000,0.000000,0.000000}%
\pgfsetstrokecolor{currentstroke}%
\pgfsetdash{}{0pt}%
\pgfsys@defobject{currentmarker}{\pgfqpoint{-0.020833in}{0.000000in}}{\pgfqpoint{-0.000000in}{0.000000in}}{%
\pgfpathmoveto{\pgfqpoint{-0.000000in}{0.000000in}}%
\pgfpathlineto{\pgfqpoint{-0.020833in}{0.000000in}}%
\pgfusepath{stroke,fill}%
}%
\begin{pgfscope}%
\pgfsys@transformshift{3.038110in}{0.739118in}%
\pgfsys@useobject{currentmarker}{}%
\end{pgfscope}%
\end{pgfscope}%
\begin{pgfscope}%
\pgfsetbuttcap%
\pgfsetroundjoin%
\definecolor{currentfill}{rgb}{0.000000,0.000000,0.000000}%
\pgfsetfillcolor{currentfill}%
\pgfsetlinewidth{0.401500pt}%
\definecolor{currentstroke}{rgb}{0.000000,0.000000,0.000000}%
\pgfsetstrokecolor{currentstroke}%
\pgfsetdash{}{0pt}%
\pgfsys@defobject{currentmarker}{\pgfqpoint{0.000000in}{0.000000in}}{\pgfqpoint{0.020833in}{0.000000in}}{%
\pgfpathmoveto{\pgfqpoint{0.000000in}{0.000000in}}%
\pgfpathlineto{\pgfqpoint{0.020833in}{0.000000in}}%
\pgfusepath{stroke,fill}%
}%
\begin{pgfscope}%
\pgfsys@transformshift{0.650000in}{0.898677in}%
\pgfsys@useobject{currentmarker}{}%
\end{pgfscope}%
\end{pgfscope}%
\begin{pgfscope}%
\pgfsetbuttcap%
\pgfsetroundjoin%
\definecolor{currentfill}{rgb}{0.000000,0.000000,0.000000}%
\pgfsetfillcolor{currentfill}%
\pgfsetlinewidth{0.401500pt}%
\definecolor{currentstroke}{rgb}{0.000000,0.000000,0.000000}%
\pgfsetstrokecolor{currentstroke}%
\pgfsetdash{}{0pt}%
\pgfsys@defobject{currentmarker}{\pgfqpoint{-0.020833in}{0.000000in}}{\pgfqpoint{-0.000000in}{0.000000in}}{%
\pgfpathmoveto{\pgfqpoint{-0.000000in}{0.000000in}}%
\pgfpathlineto{\pgfqpoint{-0.020833in}{0.000000in}}%
\pgfusepath{stroke,fill}%
}%
\begin{pgfscope}%
\pgfsys@transformshift{3.038110in}{0.898677in}%
\pgfsys@useobject{currentmarker}{}%
\end{pgfscope}%
\end{pgfscope}%
\begin{pgfscope}%
\pgfsetbuttcap%
\pgfsetroundjoin%
\definecolor{currentfill}{rgb}{0.000000,0.000000,0.000000}%
\pgfsetfillcolor{currentfill}%
\pgfsetlinewidth{0.401500pt}%
\definecolor{currentstroke}{rgb}{0.000000,0.000000,0.000000}%
\pgfsetstrokecolor{currentstroke}%
\pgfsetdash{}{0pt}%
\pgfsys@defobject{currentmarker}{\pgfqpoint{0.000000in}{0.000000in}}{\pgfqpoint{0.020833in}{0.000000in}}{%
\pgfpathmoveto{\pgfqpoint{0.000000in}{0.000000in}}%
\pgfpathlineto{\pgfqpoint{0.020833in}{0.000000in}}%
\pgfusepath{stroke,fill}%
}%
\begin{pgfscope}%
\pgfsys@transformshift{0.650000in}{0.978457in}%
\pgfsys@useobject{currentmarker}{}%
\end{pgfscope}%
\end{pgfscope}%
\begin{pgfscope}%
\pgfsetbuttcap%
\pgfsetroundjoin%
\definecolor{currentfill}{rgb}{0.000000,0.000000,0.000000}%
\pgfsetfillcolor{currentfill}%
\pgfsetlinewidth{0.401500pt}%
\definecolor{currentstroke}{rgb}{0.000000,0.000000,0.000000}%
\pgfsetstrokecolor{currentstroke}%
\pgfsetdash{}{0pt}%
\pgfsys@defobject{currentmarker}{\pgfqpoint{-0.020833in}{0.000000in}}{\pgfqpoint{-0.000000in}{0.000000in}}{%
\pgfpathmoveto{\pgfqpoint{-0.000000in}{0.000000in}}%
\pgfpathlineto{\pgfqpoint{-0.020833in}{0.000000in}}%
\pgfusepath{stroke,fill}%
}%
\begin{pgfscope}%
\pgfsys@transformshift{3.038110in}{0.978457in}%
\pgfsys@useobject{currentmarker}{}%
\end{pgfscope}%
\end{pgfscope}%
\begin{pgfscope}%
\pgfsetbuttcap%
\pgfsetroundjoin%
\definecolor{currentfill}{rgb}{0.000000,0.000000,0.000000}%
\pgfsetfillcolor{currentfill}%
\pgfsetlinewidth{0.401500pt}%
\definecolor{currentstroke}{rgb}{0.000000,0.000000,0.000000}%
\pgfsetstrokecolor{currentstroke}%
\pgfsetdash{}{0pt}%
\pgfsys@defobject{currentmarker}{\pgfqpoint{0.000000in}{0.000000in}}{\pgfqpoint{0.020833in}{0.000000in}}{%
\pgfpathmoveto{\pgfqpoint{0.000000in}{0.000000in}}%
\pgfpathlineto{\pgfqpoint{0.020833in}{0.000000in}}%
\pgfusepath{stroke,fill}%
}%
\begin{pgfscope}%
\pgfsys@transformshift{0.650000in}{1.058236in}%
\pgfsys@useobject{currentmarker}{}%
\end{pgfscope}%
\end{pgfscope}%
\begin{pgfscope}%
\pgfsetbuttcap%
\pgfsetroundjoin%
\definecolor{currentfill}{rgb}{0.000000,0.000000,0.000000}%
\pgfsetfillcolor{currentfill}%
\pgfsetlinewidth{0.401500pt}%
\definecolor{currentstroke}{rgb}{0.000000,0.000000,0.000000}%
\pgfsetstrokecolor{currentstroke}%
\pgfsetdash{}{0pt}%
\pgfsys@defobject{currentmarker}{\pgfqpoint{-0.020833in}{0.000000in}}{\pgfqpoint{-0.000000in}{0.000000in}}{%
\pgfpathmoveto{\pgfqpoint{-0.000000in}{0.000000in}}%
\pgfpathlineto{\pgfqpoint{-0.020833in}{0.000000in}}%
\pgfusepath{stroke,fill}%
}%
\begin{pgfscope}%
\pgfsys@transformshift{3.038110in}{1.058236in}%
\pgfsys@useobject{currentmarker}{}%
\end{pgfscope}%
\end{pgfscope}%
\begin{pgfscope}%
\pgfsetbuttcap%
\pgfsetroundjoin%
\definecolor{currentfill}{rgb}{0.000000,0.000000,0.000000}%
\pgfsetfillcolor{currentfill}%
\pgfsetlinewidth{0.401500pt}%
\definecolor{currentstroke}{rgb}{0.000000,0.000000,0.000000}%
\pgfsetstrokecolor{currentstroke}%
\pgfsetdash{}{0pt}%
\pgfsys@defobject{currentmarker}{\pgfqpoint{0.000000in}{0.000000in}}{\pgfqpoint{0.020833in}{0.000000in}}{%
\pgfpathmoveto{\pgfqpoint{0.000000in}{0.000000in}}%
\pgfpathlineto{\pgfqpoint{0.020833in}{0.000000in}}%
\pgfusepath{stroke,fill}%
}%
\begin{pgfscope}%
\pgfsys@transformshift{0.650000in}{1.138016in}%
\pgfsys@useobject{currentmarker}{}%
\end{pgfscope}%
\end{pgfscope}%
\begin{pgfscope}%
\pgfsetbuttcap%
\pgfsetroundjoin%
\definecolor{currentfill}{rgb}{0.000000,0.000000,0.000000}%
\pgfsetfillcolor{currentfill}%
\pgfsetlinewidth{0.401500pt}%
\definecolor{currentstroke}{rgb}{0.000000,0.000000,0.000000}%
\pgfsetstrokecolor{currentstroke}%
\pgfsetdash{}{0pt}%
\pgfsys@defobject{currentmarker}{\pgfqpoint{-0.020833in}{0.000000in}}{\pgfqpoint{-0.000000in}{0.000000in}}{%
\pgfpathmoveto{\pgfqpoint{-0.000000in}{0.000000in}}%
\pgfpathlineto{\pgfqpoint{-0.020833in}{0.000000in}}%
\pgfusepath{stroke,fill}%
}%
\begin{pgfscope}%
\pgfsys@transformshift{3.038110in}{1.138016in}%
\pgfsys@useobject{currentmarker}{}%
\end{pgfscope}%
\end{pgfscope}%
\begin{pgfscope}%
\pgfsetbuttcap%
\pgfsetroundjoin%
\definecolor{currentfill}{rgb}{0.000000,0.000000,0.000000}%
\pgfsetfillcolor{currentfill}%
\pgfsetlinewidth{0.401500pt}%
\definecolor{currentstroke}{rgb}{0.000000,0.000000,0.000000}%
\pgfsetstrokecolor{currentstroke}%
\pgfsetdash{}{0pt}%
\pgfsys@defobject{currentmarker}{\pgfqpoint{0.000000in}{0.000000in}}{\pgfqpoint{0.020833in}{0.000000in}}{%
\pgfpathmoveto{\pgfqpoint{0.000000in}{0.000000in}}%
\pgfpathlineto{\pgfqpoint{0.020833in}{0.000000in}}%
\pgfusepath{stroke,fill}%
}%
\begin{pgfscope}%
\pgfsys@transformshift{0.650000in}{1.297575in}%
\pgfsys@useobject{currentmarker}{}%
\end{pgfscope}%
\end{pgfscope}%
\begin{pgfscope}%
\pgfsetbuttcap%
\pgfsetroundjoin%
\definecolor{currentfill}{rgb}{0.000000,0.000000,0.000000}%
\pgfsetfillcolor{currentfill}%
\pgfsetlinewidth{0.401500pt}%
\definecolor{currentstroke}{rgb}{0.000000,0.000000,0.000000}%
\pgfsetstrokecolor{currentstroke}%
\pgfsetdash{}{0pt}%
\pgfsys@defobject{currentmarker}{\pgfqpoint{-0.020833in}{0.000000in}}{\pgfqpoint{-0.000000in}{0.000000in}}{%
\pgfpathmoveto{\pgfqpoint{-0.000000in}{0.000000in}}%
\pgfpathlineto{\pgfqpoint{-0.020833in}{0.000000in}}%
\pgfusepath{stroke,fill}%
}%
\begin{pgfscope}%
\pgfsys@transformshift{3.038110in}{1.297575in}%
\pgfsys@useobject{currentmarker}{}%
\end{pgfscope}%
\end{pgfscope}%
\begin{pgfscope}%
\pgfsetbuttcap%
\pgfsetroundjoin%
\definecolor{currentfill}{rgb}{0.000000,0.000000,0.000000}%
\pgfsetfillcolor{currentfill}%
\pgfsetlinewidth{0.401500pt}%
\definecolor{currentstroke}{rgb}{0.000000,0.000000,0.000000}%
\pgfsetstrokecolor{currentstroke}%
\pgfsetdash{}{0pt}%
\pgfsys@defobject{currentmarker}{\pgfqpoint{0.000000in}{0.000000in}}{\pgfqpoint{0.020833in}{0.000000in}}{%
\pgfpathmoveto{\pgfqpoint{0.000000in}{0.000000in}}%
\pgfpathlineto{\pgfqpoint{0.020833in}{0.000000in}}%
\pgfusepath{stroke,fill}%
}%
\begin{pgfscope}%
\pgfsys@transformshift{0.650000in}{1.377354in}%
\pgfsys@useobject{currentmarker}{}%
\end{pgfscope}%
\end{pgfscope}%
\begin{pgfscope}%
\pgfsetbuttcap%
\pgfsetroundjoin%
\definecolor{currentfill}{rgb}{0.000000,0.000000,0.000000}%
\pgfsetfillcolor{currentfill}%
\pgfsetlinewidth{0.401500pt}%
\definecolor{currentstroke}{rgb}{0.000000,0.000000,0.000000}%
\pgfsetstrokecolor{currentstroke}%
\pgfsetdash{}{0pt}%
\pgfsys@defobject{currentmarker}{\pgfqpoint{-0.020833in}{0.000000in}}{\pgfqpoint{-0.000000in}{0.000000in}}{%
\pgfpathmoveto{\pgfqpoint{-0.000000in}{0.000000in}}%
\pgfpathlineto{\pgfqpoint{-0.020833in}{0.000000in}}%
\pgfusepath{stroke,fill}%
}%
\begin{pgfscope}%
\pgfsys@transformshift{3.038110in}{1.377354in}%
\pgfsys@useobject{currentmarker}{}%
\end{pgfscope}%
\end{pgfscope}%
\begin{pgfscope}%
\pgfsetbuttcap%
\pgfsetroundjoin%
\definecolor{currentfill}{rgb}{0.000000,0.000000,0.000000}%
\pgfsetfillcolor{currentfill}%
\pgfsetlinewidth{0.401500pt}%
\definecolor{currentstroke}{rgb}{0.000000,0.000000,0.000000}%
\pgfsetstrokecolor{currentstroke}%
\pgfsetdash{}{0pt}%
\pgfsys@defobject{currentmarker}{\pgfqpoint{0.000000in}{0.000000in}}{\pgfqpoint{0.020833in}{0.000000in}}{%
\pgfpathmoveto{\pgfqpoint{0.000000in}{0.000000in}}%
\pgfpathlineto{\pgfqpoint{0.020833in}{0.000000in}}%
\pgfusepath{stroke,fill}%
}%
\begin{pgfscope}%
\pgfsys@transformshift{0.650000in}{1.457134in}%
\pgfsys@useobject{currentmarker}{}%
\end{pgfscope}%
\end{pgfscope}%
\begin{pgfscope}%
\pgfsetbuttcap%
\pgfsetroundjoin%
\definecolor{currentfill}{rgb}{0.000000,0.000000,0.000000}%
\pgfsetfillcolor{currentfill}%
\pgfsetlinewidth{0.401500pt}%
\definecolor{currentstroke}{rgb}{0.000000,0.000000,0.000000}%
\pgfsetstrokecolor{currentstroke}%
\pgfsetdash{}{0pt}%
\pgfsys@defobject{currentmarker}{\pgfqpoint{-0.020833in}{0.000000in}}{\pgfqpoint{-0.000000in}{0.000000in}}{%
\pgfpathmoveto{\pgfqpoint{-0.000000in}{0.000000in}}%
\pgfpathlineto{\pgfqpoint{-0.020833in}{0.000000in}}%
\pgfusepath{stroke,fill}%
}%
\begin{pgfscope}%
\pgfsys@transformshift{3.038110in}{1.457134in}%
\pgfsys@useobject{currentmarker}{}%
\end{pgfscope}%
\end{pgfscope}%
\begin{pgfscope}%
\pgfsetbuttcap%
\pgfsetroundjoin%
\definecolor{currentfill}{rgb}{0.000000,0.000000,0.000000}%
\pgfsetfillcolor{currentfill}%
\pgfsetlinewidth{0.401500pt}%
\definecolor{currentstroke}{rgb}{0.000000,0.000000,0.000000}%
\pgfsetstrokecolor{currentstroke}%
\pgfsetdash{}{0pt}%
\pgfsys@defobject{currentmarker}{\pgfqpoint{0.000000in}{0.000000in}}{\pgfqpoint{0.020833in}{0.000000in}}{%
\pgfpathmoveto{\pgfqpoint{0.000000in}{0.000000in}}%
\pgfpathlineto{\pgfqpoint{0.020833in}{0.000000in}}%
\pgfusepath{stroke,fill}%
}%
\begin{pgfscope}%
\pgfsys@transformshift{0.650000in}{1.536913in}%
\pgfsys@useobject{currentmarker}{}%
\end{pgfscope}%
\end{pgfscope}%
\begin{pgfscope}%
\pgfsetbuttcap%
\pgfsetroundjoin%
\definecolor{currentfill}{rgb}{0.000000,0.000000,0.000000}%
\pgfsetfillcolor{currentfill}%
\pgfsetlinewidth{0.401500pt}%
\definecolor{currentstroke}{rgb}{0.000000,0.000000,0.000000}%
\pgfsetstrokecolor{currentstroke}%
\pgfsetdash{}{0pt}%
\pgfsys@defobject{currentmarker}{\pgfqpoint{-0.020833in}{0.000000in}}{\pgfqpoint{-0.000000in}{0.000000in}}{%
\pgfpathmoveto{\pgfqpoint{-0.000000in}{0.000000in}}%
\pgfpathlineto{\pgfqpoint{-0.020833in}{0.000000in}}%
\pgfusepath{stroke,fill}%
}%
\begin{pgfscope}%
\pgfsys@transformshift{3.038110in}{1.536913in}%
\pgfsys@useobject{currentmarker}{}%
\end{pgfscope}%
\end{pgfscope}%
\begin{pgfscope}%
\pgfsetbuttcap%
\pgfsetroundjoin%
\definecolor{currentfill}{rgb}{0.000000,0.000000,0.000000}%
\pgfsetfillcolor{currentfill}%
\pgfsetlinewidth{0.401500pt}%
\definecolor{currentstroke}{rgb}{0.000000,0.000000,0.000000}%
\pgfsetstrokecolor{currentstroke}%
\pgfsetdash{}{0pt}%
\pgfsys@defobject{currentmarker}{\pgfqpoint{0.000000in}{0.000000in}}{\pgfqpoint{0.020833in}{0.000000in}}{%
\pgfpathmoveto{\pgfqpoint{0.000000in}{0.000000in}}%
\pgfpathlineto{\pgfqpoint{0.020833in}{0.000000in}}%
\pgfusepath{stroke,fill}%
}%
\begin{pgfscope}%
\pgfsys@transformshift{0.650000in}{1.696472in}%
\pgfsys@useobject{currentmarker}{}%
\end{pgfscope}%
\end{pgfscope}%
\begin{pgfscope}%
\pgfsetbuttcap%
\pgfsetroundjoin%
\definecolor{currentfill}{rgb}{0.000000,0.000000,0.000000}%
\pgfsetfillcolor{currentfill}%
\pgfsetlinewidth{0.401500pt}%
\definecolor{currentstroke}{rgb}{0.000000,0.000000,0.000000}%
\pgfsetstrokecolor{currentstroke}%
\pgfsetdash{}{0pt}%
\pgfsys@defobject{currentmarker}{\pgfqpoint{-0.020833in}{0.000000in}}{\pgfqpoint{-0.000000in}{0.000000in}}{%
\pgfpathmoveto{\pgfqpoint{-0.000000in}{0.000000in}}%
\pgfpathlineto{\pgfqpoint{-0.020833in}{0.000000in}}%
\pgfusepath{stroke,fill}%
}%
\begin{pgfscope}%
\pgfsys@transformshift{3.038110in}{1.696472in}%
\pgfsys@useobject{currentmarker}{}%
\end{pgfscope}%
\end{pgfscope}%
\begin{pgfscope}%
\pgfsetbuttcap%
\pgfsetroundjoin%
\definecolor{currentfill}{rgb}{0.000000,0.000000,0.000000}%
\pgfsetfillcolor{currentfill}%
\pgfsetlinewidth{0.401500pt}%
\definecolor{currentstroke}{rgb}{0.000000,0.000000,0.000000}%
\pgfsetstrokecolor{currentstroke}%
\pgfsetdash{}{0pt}%
\pgfsys@defobject{currentmarker}{\pgfqpoint{0.000000in}{0.000000in}}{\pgfqpoint{0.020833in}{0.000000in}}{%
\pgfpathmoveto{\pgfqpoint{0.000000in}{0.000000in}}%
\pgfpathlineto{\pgfqpoint{0.020833in}{0.000000in}}%
\pgfusepath{stroke,fill}%
}%
\begin{pgfscope}%
\pgfsys@transformshift{0.650000in}{1.776252in}%
\pgfsys@useobject{currentmarker}{}%
\end{pgfscope}%
\end{pgfscope}%
\begin{pgfscope}%
\pgfsetbuttcap%
\pgfsetroundjoin%
\definecolor{currentfill}{rgb}{0.000000,0.000000,0.000000}%
\pgfsetfillcolor{currentfill}%
\pgfsetlinewidth{0.401500pt}%
\definecolor{currentstroke}{rgb}{0.000000,0.000000,0.000000}%
\pgfsetstrokecolor{currentstroke}%
\pgfsetdash{}{0pt}%
\pgfsys@defobject{currentmarker}{\pgfqpoint{-0.020833in}{0.000000in}}{\pgfqpoint{-0.000000in}{0.000000in}}{%
\pgfpathmoveto{\pgfqpoint{-0.000000in}{0.000000in}}%
\pgfpathlineto{\pgfqpoint{-0.020833in}{0.000000in}}%
\pgfusepath{stroke,fill}%
}%
\begin{pgfscope}%
\pgfsys@transformshift{3.038110in}{1.776252in}%
\pgfsys@useobject{currentmarker}{}%
\end{pgfscope}%
\end{pgfscope}%
\begin{pgfscope}%
\pgfsetbuttcap%
\pgfsetroundjoin%
\definecolor{currentfill}{rgb}{0.000000,0.000000,0.000000}%
\pgfsetfillcolor{currentfill}%
\pgfsetlinewidth{0.401500pt}%
\definecolor{currentstroke}{rgb}{0.000000,0.000000,0.000000}%
\pgfsetstrokecolor{currentstroke}%
\pgfsetdash{}{0pt}%
\pgfsys@defobject{currentmarker}{\pgfqpoint{0.000000in}{0.000000in}}{\pgfqpoint{0.020833in}{0.000000in}}{%
\pgfpathmoveto{\pgfqpoint{0.000000in}{0.000000in}}%
\pgfpathlineto{\pgfqpoint{0.020833in}{0.000000in}}%
\pgfusepath{stroke,fill}%
}%
\begin{pgfscope}%
\pgfsys@transformshift{0.650000in}{1.856031in}%
\pgfsys@useobject{currentmarker}{}%
\end{pgfscope}%
\end{pgfscope}%
\begin{pgfscope}%
\pgfsetbuttcap%
\pgfsetroundjoin%
\definecolor{currentfill}{rgb}{0.000000,0.000000,0.000000}%
\pgfsetfillcolor{currentfill}%
\pgfsetlinewidth{0.401500pt}%
\definecolor{currentstroke}{rgb}{0.000000,0.000000,0.000000}%
\pgfsetstrokecolor{currentstroke}%
\pgfsetdash{}{0pt}%
\pgfsys@defobject{currentmarker}{\pgfqpoint{-0.020833in}{0.000000in}}{\pgfqpoint{-0.000000in}{0.000000in}}{%
\pgfpathmoveto{\pgfqpoint{-0.000000in}{0.000000in}}%
\pgfpathlineto{\pgfqpoint{-0.020833in}{0.000000in}}%
\pgfusepath{stroke,fill}%
}%
\begin{pgfscope}%
\pgfsys@transformshift{3.038110in}{1.856031in}%
\pgfsys@useobject{currentmarker}{}%
\end{pgfscope}%
\end{pgfscope}%
\begin{pgfscope}%
\pgfsetbuttcap%
\pgfsetroundjoin%
\definecolor{currentfill}{rgb}{0.000000,0.000000,0.000000}%
\pgfsetfillcolor{currentfill}%
\pgfsetlinewidth{0.401500pt}%
\definecolor{currentstroke}{rgb}{0.000000,0.000000,0.000000}%
\pgfsetstrokecolor{currentstroke}%
\pgfsetdash{}{0pt}%
\pgfsys@defobject{currentmarker}{\pgfqpoint{0.000000in}{0.000000in}}{\pgfqpoint{0.020833in}{0.000000in}}{%
\pgfpathmoveto{\pgfqpoint{0.000000in}{0.000000in}}%
\pgfpathlineto{\pgfqpoint{0.020833in}{0.000000in}}%
\pgfusepath{stroke,fill}%
}%
\begin{pgfscope}%
\pgfsys@transformshift{0.650000in}{1.935811in}%
\pgfsys@useobject{currentmarker}{}%
\end{pgfscope}%
\end{pgfscope}%
\begin{pgfscope}%
\pgfsetbuttcap%
\pgfsetroundjoin%
\definecolor{currentfill}{rgb}{0.000000,0.000000,0.000000}%
\pgfsetfillcolor{currentfill}%
\pgfsetlinewidth{0.401500pt}%
\definecolor{currentstroke}{rgb}{0.000000,0.000000,0.000000}%
\pgfsetstrokecolor{currentstroke}%
\pgfsetdash{}{0pt}%
\pgfsys@defobject{currentmarker}{\pgfqpoint{-0.020833in}{0.000000in}}{\pgfqpoint{-0.000000in}{0.000000in}}{%
\pgfpathmoveto{\pgfqpoint{-0.000000in}{0.000000in}}%
\pgfpathlineto{\pgfqpoint{-0.020833in}{0.000000in}}%
\pgfusepath{stroke,fill}%
}%
\begin{pgfscope}%
\pgfsys@transformshift{3.038110in}{1.935811in}%
\pgfsys@useobject{currentmarker}{}%
\end{pgfscope}%
\end{pgfscope}%
\begin{pgfscope}%
\pgfsetbuttcap%
\pgfsetroundjoin%
\definecolor{currentfill}{rgb}{0.000000,0.000000,0.000000}%
\pgfsetfillcolor{currentfill}%
\pgfsetlinewidth{0.401500pt}%
\definecolor{currentstroke}{rgb}{0.000000,0.000000,0.000000}%
\pgfsetstrokecolor{currentstroke}%
\pgfsetdash{}{0pt}%
\pgfsys@defobject{currentmarker}{\pgfqpoint{0.000000in}{0.000000in}}{\pgfqpoint{0.020833in}{0.000000in}}{%
\pgfpathmoveto{\pgfqpoint{0.000000in}{0.000000in}}%
\pgfpathlineto{\pgfqpoint{0.020833in}{0.000000in}}%
\pgfusepath{stroke,fill}%
}%
\begin{pgfscope}%
\pgfsys@transformshift{0.650000in}{2.095370in}%
\pgfsys@useobject{currentmarker}{}%
\end{pgfscope}%
\end{pgfscope}%
\begin{pgfscope}%
\pgfsetbuttcap%
\pgfsetroundjoin%
\definecolor{currentfill}{rgb}{0.000000,0.000000,0.000000}%
\pgfsetfillcolor{currentfill}%
\pgfsetlinewidth{0.401500pt}%
\definecolor{currentstroke}{rgb}{0.000000,0.000000,0.000000}%
\pgfsetstrokecolor{currentstroke}%
\pgfsetdash{}{0pt}%
\pgfsys@defobject{currentmarker}{\pgfqpoint{-0.020833in}{0.000000in}}{\pgfqpoint{-0.000000in}{0.000000in}}{%
\pgfpathmoveto{\pgfqpoint{-0.000000in}{0.000000in}}%
\pgfpathlineto{\pgfqpoint{-0.020833in}{0.000000in}}%
\pgfusepath{stroke,fill}%
}%
\begin{pgfscope}%
\pgfsys@transformshift{3.038110in}{2.095370in}%
\pgfsys@useobject{currentmarker}{}%
\end{pgfscope}%
\end{pgfscope}%
\begin{pgfscope}%
\pgfsetbuttcap%
\pgfsetroundjoin%
\definecolor{currentfill}{rgb}{0.000000,0.000000,0.000000}%
\pgfsetfillcolor{currentfill}%
\pgfsetlinewidth{0.401500pt}%
\definecolor{currentstroke}{rgb}{0.000000,0.000000,0.000000}%
\pgfsetstrokecolor{currentstroke}%
\pgfsetdash{}{0pt}%
\pgfsys@defobject{currentmarker}{\pgfqpoint{0.000000in}{0.000000in}}{\pgfqpoint{0.020833in}{0.000000in}}{%
\pgfpathmoveto{\pgfqpoint{0.000000in}{0.000000in}}%
\pgfpathlineto{\pgfqpoint{0.020833in}{0.000000in}}%
\pgfusepath{stroke,fill}%
}%
\begin{pgfscope}%
\pgfsys@transformshift{0.650000in}{2.175150in}%
\pgfsys@useobject{currentmarker}{}%
\end{pgfscope}%
\end{pgfscope}%
\begin{pgfscope}%
\pgfsetbuttcap%
\pgfsetroundjoin%
\definecolor{currentfill}{rgb}{0.000000,0.000000,0.000000}%
\pgfsetfillcolor{currentfill}%
\pgfsetlinewidth{0.401500pt}%
\definecolor{currentstroke}{rgb}{0.000000,0.000000,0.000000}%
\pgfsetstrokecolor{currentstroke}%
\pgfsetdash{}{0pt}%
\pgfsys@defobject{currentmarker}{\pgfqpoint{-0.020833in}{0.000000in}}{\pgfqpoint{-0.000000in}{0.000000in}}{%
\pgfpathmoveto{\pgfqpoint{-0.000000in}{0.000000in}}%
\pgfpathlineto{\pgfqpoint{-0.020833in}{0.000000in}}%
\pgfusepath{stroke,fill}%
}%
\begin{pgfscope}%
\pgfsys@transformshift{3.038110in}{2.175150in}%
\pgfsys@useobject{currentmarker}{}%
\end{pgfscope}%
\end{pgfscope}%
\begin{pgfscope}%
\pgfsetbuttcap%
\pgfsetroundjoin%
\definecolor{currentfill}{rgb}{0.000000,0.000000,0.000000}%
\pgfsetfillcolor{currentfill}%
\pgfsetlinewidth{0.401500pt}%
\definecolor{currentstroke}{rgb}{0.000000,0.000000,0.000000}%
\pgfsetstrokecolor{currentstroke}%
\pgfsetdash{}{0pt}%
\pgfsys@defobject{currentmarker}{\pgfqpoint{0.000000in}{0.000000in}}{\pgfqpoint{0.020833in}{0.000000in}}{%
\pgfpathmoveto{\pgfqpoint{0.000000in}{0.000000in}}%
\pgfpathlineto{\pgfqpoint{0.020833in}{0.000000in}}%
\pgfusepath{stroke,fill}%
}%
\begin{pgfscope}%
\pgfsys@transformshift{0.650000in}{2.254929in}%
\pgfsys@useobject{currentmarker}{}%
\end{pgfscope}%
\end{pgfscope}%
\begin{pgfscope}%
\pgfsetbuttcap%
\pgfsetroundjoin%
\definecolor{currentfill}{rgb}{0.000000,0.000000,0.000000}%
\pgfsetfillcolor{currentfill}%
\pgfsetlinewidth{0.401500pt}%
\definecolor{currentstroke}{rgb}{0.000000,0.000000,0.000000}%
\pgfsetstrokecolor{currentstroke}%
\pgfsetdash{}{0pt}%
\pgfsys@defobject{currentmarker}{\pgfqpoint{-0.020833in}{0.000000in}}{\pgfqpoint{-0.000000in}{0.000000in}}{%
\pgfpathmoveto{\pgfqpoint{-0.000000in}{0.000000in}}%
\pgfpathlineto{\pgfqpoint{-0.020833in}{0.000000in}}%
\pgfusepath{stroke,fill}%
}%
\begin{pgfscope}%
\pgfsys@transformshift{3.038110in}{2.254929in}%
\pgfsys@useobject{currentmarker}{}%
\end{pgfscope}%
\end{pgfscope}%
\begin{pgfscope}%
\pgfsetbuttcap%
\pgfsetroundjoin%
\definecolor{currentfill}{rgb}{0.000000,0.000000,0.000000}%
\pgfsetfillcolor{currentfill}%
\pgfsetlinewidth{0.401500pt}%
\definecolor{currentstroke}{rgb}{0.000000,0.000000,0.000000}%
\pgfsetstrokecolor{currentstroke}%
\pgfsetdash{}{0pt}%
\pgfsys@defobject{currentmarker}{\pgfqpoint{0.000000in}{0.000000in}}{\pgfqpoint{0.020833in}{0.000000in}}{%
\pgfpathmoveto{\pgfqpoint{0.000000in}{0.000000in}}%
\pgfpathlineto{\pgfqpoint{0.020833in}{0.000000in}}%
\pgfusepath{stroke,fill}%
}%
\begin{pgfscope}%
\pgfsys@transformshift{0.650000in}{2.334709in}%
\pgfsys@useobject{currentmarker}{}%
\end{pgfscope}%
\end{pgfscope}%
\begin{pgfscope}%
\pgfsetbuttcap%
\pgfsetroundjoin%
\definecolor{currentfill}{rgb}{0.000000,0.000000,0.000000}%
\pgfsetfillcolor{currentfill}%
\pgfsetlinewidth{0.401500pt}%
\definecolor{currentstroke}{rgb}{0.000000,0.000000,0.000000}%
\pgfsetstrokecolor{currentstroke}%
\pgfsetdash{}{0pt}%
\pgfsys@defobject{currentmarker}{\pgfqpoint{-0.020833in}{0.000000in}}{\pgfqpoint{-0.000000in}{0.000000in}}{%
\pgfpathmoveto{\pgfqpoint{-0.000000in}{0.000000in}}%
\pgfpathlineto{\pgfqpoint{-0.020833in}{0.000000in}}%
\pgfusepath{stroke,fill}%
}%
\begin{pgfscope}%
\pgfsys@transformshift{3.038110in}{2.334709in}%
\pgfsys@useobject{currentmarker}{}%
\end{pgfscope}%
\end{pgfscope}%
\begin{pgfscope}%
\definecolor{textcolor}{rgb}{0.000000,0.000000,0.000000}%
\pgfsetstrokecolor{textcolor}%
\pgfsetfillcolor{textcolor}%
\pgftext[x=0.421528in,y=1.417244in,,bottom,rotate=90.000000]{\color{textcolor}{\rmfamily\fontsize{12.000000}{14.400000}\selectfont\catcode`\^=\active\def^{\ifmmode\sp\else\^{}\fi}\catcode`\%=\active\def%{\%}$U^+$}}%
\end{pgfscope}%
\begin{pgfscope}%
\pgfpathrectangle{\pgfqpoint{0.650000in}{0.420000in}}{\pgfqpoint{2.388110in}{1.994488in}}%
\pgfusepath{clip}%
\pgfsetbuttcap%
\pgfsetroundjoin%
\pgfsetlinewidth{0.501875pt}%
\definecolor{currentstroke}{rgb}{0.000000,0.000000,0.000000}%
\pgfsetstrokecolor{currentstroke}%
\pgfsetdash{{1.850000pt}{0.800000pt}}{0.000000pt}%
\pgfpathmoveto{\pgfqpoint{0.645000in}{0.461577in}}%
\pgfpathlineto{\pgfqpoint{0.782176in}{0.481369in}}%
\pgfpathlineto{\pgfqpoint{0.906584in}{0.512053in}}%
\pgfpathlineto{\pgfqpoint{0.994853in}{0.542738in}}%
\pgfpathlineto{\pgfqpoint{1.063320in}{0.573422in}}%
\pgfpathlineto{\pgfqpoint{1.119261in}{0.604107in}}%
\pgfpathlineto{\pgfqpoint{1.166559in}{0.634791in}}%
\pgfpathlineto{\pgfqpoint{1.207530in}{0.665475in}}%
\pgfpathlineto{\pgfqpoint{1.243669in}{0.696160in}}%
\pgfpathlineto{\pgfqpoint{1.275997in}{0.726844in}}%
\pgfpathlineto{\pgfqpoint{1.305240in}{0.757529in}}%
\pgfpathlineto{\pgfqpoint{1.331938in}{0.788213in}}%
\pgfpathlineto{\pgfqpoint{1.356497in}{0.818898in}}%
\pgfpathlineto{\pgfqpoint{1.379236in}{0.849582in}}%
\pgfpathlineto{\pgfqpoint{1.400405in}{0.880266in}}%
\pgfpathlineto{\pgfqpoint{1.420207in}{0.910951in}}%
\pgfpathlineto{\pgfqpoint{1.438808in}{0.941635in}}%
\pgfpathlineto{\pgfqpoint{1.456346in}{0.972320in}}%
\pgfpathlineto{\pgfqpoint{1.472935in}{1.003004in}}%
\pgfpathlineto{\pgfqpoint{1.488674in}{1.033689in}}%
\pgfpathlineto{\pgfqpoint{1.503644in}{1.064373in}}%
\pgfpathlineto{\pgfqpoint{1.517917in}{1.095058in}}%
\pgfpathlineto{\pgfqpoint{1.531556in}{1.125742in}}%
\pgfpathlineto{\pgfqpoint{1.544615in}{1.156426in}}%
\pgfpathlineto{\pgfqpoint{1.557140in}{1.187111in}}%
\pgfpathlineto{\pgfqpoint{1.569174in}{1.217795in}}%
\pgfpathlineto{\pgfqpoint{1.580754in}{1.248480in}}%
\pgfpathlineto{\pgfqpoint{1.591913in}{1.279164in}}%
\pgfpathlineto{\pgfqpoint{1.602680in}{1.309849in}}%
\pgfpathlineto{\pgfqpoint{1.613082in}{1.340533in}}%
\pgfpathlineto{\pgfqpoint{1.623142in}{1.371217in}}%
\pgfpathlineto{\pgfqpoint{1.632884in}{1.401902in}}%
\pgfpathlineto{\pgfqpoint{1.642325in}{1.432586in}}%
\pgfpathlineto{\pgfqpoint{1.651485in}{1.463271in}}%
\pgfpathlineto{\pgfqpoint{1.660379in}{1.493955in}}%
\pgfpathlineto{\pgfqpoint{1.669023in}{1.524640in}}%
\pgfpathlineto{\pgfqpoint{1.677430in}{1.555324in}}%
\pgfpathlineto{\pgfqpoint{1.685612in}{1.586008in}}%
\pgfpathlineto{\pgfqpoint{1.693582in}{1.616693in}}%
\pgfusepath{stroke}%
\end{pgfscope}%
\begin{pgfscope}%
\pgfpathrectangle{\pgfqpoint{0.650000in}{0.420000in}}{\pgfqpoint{2.388110in}{1.994488in}}%
\pgfusepath{clip}%
\pgfsetbuttcap%
\pgfsetroundjoin%
\pgfsetlinewidth{0.501875pt}%
\definecolor{currentstroke}{rgb}{0.000000,0.000000,0.000000}%
\pgfsetstrokecolor{currentstroke}%
\pgfsetdash{{1.850000pt}{0.800000pt}}{0.000000pt}%
\pgfpathmoveto{\pgfqpoint{1.569174in}{1.282900in}}%
\pgfpathlineto{\pgfqpoint{3.016941in}{2.201046in}}%
\pgfpathlineto{\pgfqpoint{3.043110in}{2.217641in}}%
\pgfusepath{stroke}%
\end{pgfscope}%
\begin{pgfscope}%
\pgfpathrectangle{\pgfqpoint{0.650000in}{0.420000in}}{\pgfqpoint{2.388110in}{1.994488in}}%
\pgfusepath{clip}%
\pgfsetrectcap%
\pgfsetroundjoin%
\pgfsetlinewidth{0.602250pt}%
\definecolor{currentstroke}{rgb}{0.000000,0.000000,0.000000}%
\pgfsetstrokecolor{currentstroke}%
\pgfsetdash{}{0pt}%
\pgfpathmoveto{\pgfqpoint{0.645000in}{0.460785in}}%
\pgfpathlineto{\pgfqpoint{0.718874in}{0.469925in}}%
\pgfpathlineto{\pgfqpoint{0.855798in}{0.497990in}}%
\pgfpathlineto{\pgfqpoint{0.967671in}{0.532241in}}%
\pgfpathlineto{\pgfqpoint{1.062253in}{0.572576in}}%
\pgfpathlineto{\pgfqpoint{1.144181in}{0.618769in}}%
\pgfpathlineto{\pgfqpoint{1.216443in}{0.670391in}}%
\pgfpathlineto{\pgfqpoint{1.281080in}{0.726732in}}%
\pgfpathlineto{\pgfqpoint{1.339547in}{0.786757in}}%
\pgfpathlineto{\pgfqpoint{1.392920in}{0.849132in}}%
\pgfpathlineto{\pgfqpoint{1.442015in}{0.912336in}}%
\pgfpathlineto{\pgfqpoint{1.487466in}{0.974827in}}%
\pgfpathlineto{\pgfqpoint{1.641563in}{1.194575in}}%
\pgfpathlineto{\pgfqpoint{1.674706in}{1.239082in}}%
\pgfpathlineto{\pgfqpoint{1.706145in}{1.279306in}}%
\pgfpathlineto{\pgfqpoint{1.736046in}{1.315523in}}%
\pgfpathlineto{\pgfqpoint{1.764552in}{1.348084in}}%
\pgfpathlineto{\pgfqpoint{1.791786in}{1.377366in}}%
\pgfpathlineto{\pgfqpoint{1.817858in}{1.403746in}}%
\pgfpathlineto{\pgfqpoint{1.842862in}{1.427578in}}%
\pgfpathlineto{\pgfqpoint{1.866881in}{1.449188in}}%
\pgfpathlineto{\pgfqpoint{1.889989in}{1.468862in}}%
\pgfpathlineto{\pgfqpoint{1.912253in}{1.486860in}}%
\pgfpathlineto{\pgfqpoint{1.954479in}{1.518746in}}%
\pgfpathlineto{\pgfqpoint{1.993964in}{1.546444in}}%
\pgfpathlineto{\pgfqpoint{2.048762in}{1.582659in}}%
\pgfpathlineto{\pgfqpoint{2.160125in}{1.654527in}}%
\pgfpathlineto{\pgfqpoint{2.202211in}{1.683113in}}%
\pgfpathlineto{\pgfqpoint{2.241547in}{1.711236in}}%
\pgfpathlineto{\pgfqpoint{2.278463in}{1.739234in}}%
\pgfpathlineto{\pgfqpoint{2.313235in}{1.767331in}}%
\pgfpathlineto{\pgfqpoint{2.346093in}{1.795583in}}%
\pgfpathlineto{\pgfqpoint{2.377230in}{1.824025in}}%
\pgfpathlineto{\pgfqpoint{2.406815in}{1.852671in}}%
\pgfpathlineto{\pgfqpoint{2.434988in}{1.881577in}}%
\pgfpathlineto{\pgfqpoint{2.461875in}{1.910583in}}%
\pgfpathlineto{\pgfqpoint{2.495908in}{1.949053in}}%
\pgfpathlineto{\pgfqpoint{2.528065in}{1.987180in}}%
\pgfpathlineto{\pgfqpoint{2.594481in}{2.066892in}}%
\pgfpathlineto{\pgfqpoint{2.615015in}{2.089214in}}%
\pgfpathlineto{\pgfqpoint{2.628304in}{2.102174in}}%
\pgfpathlineto{\pgfqpoint{2.641289in}{2.113307in}}%
\pgfpathlineto{\pgfqpoint{2.653982in}{2.122469in}}%
\pgfpathlineto{\pgfqpoint{2.666395in}{2.129626in}}%
\pgfpathlineto{\pgfqpoint{2.678539in}{2.134913in}}%
\pgfpathlineto{\pgfqpoint{2.690425in}{2.138594in}}%
\pgfpathlineto{\pgfqpoint{2.702063in}{2.140992in}}%
\pgfpathlineto{\pgfqpoint{2.719073in}{2.142934in}}%
\pgfpathlineto{\pgfqpoint{2.740966in}{2.143882in}}%
\pgfpathlineto{\pgfqpoint{2.792128in}{2.144292in}}%
\pgfpathlineto{\pgfqpoint{3.043110in}{2.145531in}}%
\pgfpathlineto{\pgfqpoint{3.043110in}{2.145531in}}%
\pgfusepath{stroke}%
\end{pgfscope}%
\begin{pgfscope}%
\pgfpathrectangle{\pgfqpoint{0.650000in}{0.420000in}}{\pgfqpoint{2.388110in}{1.994488in}}%
\pgfusepath{clip}%
\pgfsetrectcap%
\pgfsetroundjoin%
\pgfsetlinewidth{1.003750pt}%
\definecolor{currentstroke}{rgb}{0.870588,0.466667,0.682353}%
\pgfsetstrokecolor{currentstroke}%
\pgfsetdash{}{0pt}%
\pgfpathmoveto{\pgfqpoint{0.645000in}{0.459243in}}%
\pgfpathlineto{\pgfqpoint{0.686665in}{0.464919in}}%
\pgfpathlineto{\pgfqpoint{0.741886in}{0.473809in}}%
\pgfpathlineto{\pgfqpoint{0.792190in}{0.483425in}}%
\pgfpathlineto{\pgfqpoint{0.838669in}{0.493827in}}%
\pgfpathlineto{\pgfqpoint{0.882094in}{0.505078in}}%
\pgfpathlineto{\pgfqpoint{0.923033in}{0.517242in}}%
\pgfpathlineto{\pgfqpoint{0.961915in}{0.530387in}}%
\pgfpathlineto{\pgfqpoint{0.999071in}{0.544584in}}%
\pgfpathlineto{\pgfqpoint{1.034762in}{0.559908in}}%
\pgfpathlineto{\pgfqpoint{1.069199in}{0.576439in}}%
\pgfpathlineto{\pgfqpoint{1.102550in}{0.594264in}}%
\pgfpathlineto{\pgfqpoint{1.134959in}{0.613474in}}%
\pgfpathlineto{\pgfqpoint{1.166541in}{0.634153in}}%
\pgfpathlineto{\pgfqpoint{1.197397in}{0.656355in}}%
\pgfpathlineto{\pgfqpoint{1.227610in}{0.680132in}}%
\pgfpathlineto{\pgfqpoint{1.257252in}{0.705558in}}%
\pgfpathlineto{\pgfqpoint{1.286384in}{0.732638in}}%
\pgfpathlineto{\pgfqpoint{1.315062in}{0.761405in}}%
\pgfpathlineto{\pgfqpoint{1.343330in}{0.791854in}}%
\pgfpathlineto{\pgfqpoint{1.371232in}{0.823909in}}%
\pgfpathlineto{\pgfqpoint{1.398802in}{0.857466in}}%
\pgfpathlineto{\pgfqpoint{1.426073in}{0.892401in}}%
\pgfpathlineto{\pgfqpoint{1.453072in}{0.928508in}}%
\pgfpathlineto{\pgfqpoint{1.493118in}{0.984302in}}%
\pgfpathlineto{\pgfqpoint{1.545777in}{1.060356in}}%
\pgfpathlineto{\pgfqpoint{1.623474in}{1.173003in}}%
\pgfpathlineto{\pgfqpoint{1.661845in}{1.226575in}}%
\pgfpathlineto{\pgfqpoint{1.699950in}{1.277371in}}%
\pgfpathlineto{\pgfqpoint{1.725222in}{1.309447in}}%
\pgfpathlineto{\pgfqpoint{1.750399in}{1.340008in}}%
\pgfpathlineto{\pgfqpoint{1.775489in}{1.369026in}}%
\pgfpathlineto{\pgfqpoint{1.800499in}{1.396489in}}%
\pgfpathlineto{\pgfqpoint{1.825435in}{1.422425in}}%
\pgfpathlineto{\pgfqpoint{1.850304in}{1.446909in}}%
\pgfpathlineto{\pgfqpoint{1.875110in}{1.470019in}}%
\pgfpathlineto{\pgfqpoint{1.899859in}{1.491849in}}%
\pgfpathlineto{\pgfqpoint{1.924555in}{1.512507in}}%
\pgfpathlineto{\pgfqpoint{1.961510in}{1.541556in}}%
\pgfpathlineto{\pgfqpoint{1.998368in}{1.568705in}}%
\pgfpathlineto{\pgfqpoint{2.047381in}{1.602862in}}%
\pgfpathlineto{\pgfqpoint{2.181555in}{1.694660in}}%
\pgfpathlineto{\pgfqpoint{2.230174in}{1.729824in}}%
\pgfpathlineto{\pgfqpoint{2.278722in}{1.766313in}}%
\pgfpathlineto{\pgfqpoint{2.315094in}{1.795240in}}%
\pgfpathlineto{\pgfqpoint{2.351435in}{1.826043in}}%
\pgfpathlineto{\pgfqpoint{2.387726in}{1.858571in}}%
\pgfpathlineto{\pgfqpoint{2.434639in}{1.902571in}}%
\pgfpathlineto{\pgfqpoint{2.466451in}{1.933777in}}%
\pgfpathlineto{\pgfqpoint{2.495284in}{1.963627in}}%
\pgfpathlineto{\pgfqpoint{2.529943in}{2.001401in}}%
\pgfpathlineto{\pgfqpoint{2.561082in}{2.037005in}}%
\pgfpathlineto{\pgfqpoint{2.676861in}{2.172634in}}%
\pgfpathlineto{\pgfqpoint{2.728158in}{2.232195in}}%
\pgfpathlineto{\pgfqpoint{2.744872in}{2.250698in}}%
\pgfpathlineto{\pgfqpoint{2.756836in}{2.262723in}}%
\pgfpathlineto{\pgfqpoint{2.768351in}{2.273088in}}%
\pgfpathlineto{\pgfqpoint{2.779449in}{2.281887in}}%
\pgfpathlineto{\pgfqpoint{2.790160in}{2.288992in}}%
\pgfpathlineto{\pgfqpoint{2.800510in}{2.294166in}}%
\pgfpathlineto{\pgfqpoint{2.810521in}{2.297627in}}%
\pgfpathlineto{\pgfqpoint{2.820217in}{2.299863in}}%
\pgfpathlineto{\pgfqpoint{2.832685in}{2.301632in}}%
\pgfpathlineto{\pgfqpoint{2.847590in}{2.302509in}}%
\pgfpathlineto{\pgfqpoint{2.861804in}{2.302254in}}%
\pgfpathlineto{\pgfqpoint{2.883261in}{2.300596in}}%
\pgfpathlineto{\pgfqpoint{2.919846in}{2.296425in}}%
\pgfpathlineto{\pgfqpoint{2.972533in}{2.289098in}}%
\pgfpathlineto{\pgfqpoint{3.027359in}{2.280199in}}%
\pgfpathlineto{\pgfqpoint{3.043110in}{2.277328in}}%
\pgfpathlineto{\pgfqpoint{3.043110in}{2.277328in}}%
\pgfusepath{stroke}%
\end{pgfscope}%
\begin{pgfscope}%
\pgfpathrectangle{\pgfqpoint{0.650000in}{0.420000in}}{\pgfqpoint{2.388110in}{1.994488in}}%
\pgfusepath{clip}%
\pgfsetrectcap%
\pgfsetroundjoin%
\pgfsetlinewidth{1.003750pt}%
\definecolor{currentstroke}{rgb}{0.498039,0.737255,0.254902}%
\pgfsetstrokecolor{currentstroke}%
\pgfsetdash{}{0pt}%
\pgfpathmoveto{\pgfqpoint{0.645000in}{0.459241in}}%
\pgfpathlineto{\pgfqpoint{0.682092in}{0.464222in}}%
\pgfpathlineto{\pgfqpoint{0.737313in}{0.472911in}}%
\pgfpathlineto{\pgfqpoint{0.787618in}{0.482311in}}%
\pgfpathlineto{\pgfqpoint{0.834096in}{0.492479in}}%
\pgfpathlineto{\pgfqpoint{0.877521in}{0.503477in}}%
\pgfpathlineto{\pgfqpoint{0.918461in}{0.515367in}}%
\pgfpathlineto{\pgfqpoint{0.957343in}{0.528216in}}%
\pgfpathlineto{\pgfqpoint{0.994499in}{0.542094in}}%
\pgfpathlineto{\pgfqpoint{1.030190in}{0.557074in}}%
\pgfpathlineto{\pgfqpoint{1.064626in}{0.573236in}}%
\pgfpathlineto{\pgfqpoint{1.097978in}{0.590664in}}%
\pgfpathlineto{\pgfqpoint{1.130386in}{0.609446in}}%
\pgfpathlineto{\pgfqpoint{1.161969in}{0.629668in}}%
\pgfpathlineto{\pgfqpoint{1.192824in}{0.651387in}}%
\pgfpathlineto{\pgfqpoint{1.223037in}{0.674657in}}%
\pgfpathlineto{\pgfqpoint{1.252679in}{0.699549in}}%
\pgfpathlineto{\pgfqpoint{1.281811in}{0.726077in}}%
\pgfpathlineto{\pgfqpoint{1.310489in}{0.754270in}}%
\pgfpathlineto{\pgfqpoint{1.338758in}{0.784125in}}%
\pgfpathlineto{\pgfqpoint{1.366659in}{0.815573in}}%
\pgfpathlineto{\pgfqpoint{1.394229in}{0.848515in}}%
\pgfpathlineto{\pgfqpoint{1.421500in}{0.882829in}}%
\pgfpathlineto{\pgfqpoint{1.448499in}{0.918318in}}%
\pgfpathlineto{\pgfqpoint{1.488545in}{0.973211in}}%
\pgfpathlineto{\pgfqpoint{1.541204in}{1.048118in}}%
\pgfpathlineto{\pgfqpoint{1.618901in}{1.159194in}}%
\pgfpathlineto{\pgfqpoint{1.657272in}{1.212041in}}%
\pgfpathlineto{\pgfqpoint{1.695377in}{1.262109in}}%
\pgfpathlineto{\pgfqpoint{1.720650in}{1.293689in}}%
\pgfpathlineto{\pgfqpoint{1.745827in}{1.323738in}}%
\pgfpathlineto{\pgfqpoint{1.770917in}{1.352234in}}%
\pgfpathlineto{\pgfqpoint{1.795926in}{1.379182in}}%
\pgfpathlineto{\pgfqpoint{1.820863in}{1.404625in}}%
\pgfpathlineto{\pgfqpoint{1.845731in}{1.428651in}}%
\pgfpathlineto{\pgfqpoint{1.870537in}{1.451357in}}%
\pgfpathlineto{\pgfqpoint{1.895286in}{1.472861in}}%
\pgfpathlineto{\pgfqpoint{1.932312in}{1.503147in}}%
\pgfpathlineto{\pgfqpoint{1.969233in}{1.531470in}}%
\pgfpathlineto{\pgfqpoint{2.006062in}{1.558329in}}%
\pgfpathlineto{\pgfqpoint{2.055040in}{1.592575in}}%
\pgfpathlineto{\pgfqpoint{2.201302in}{1.693268in}}%
\pgfpathlineto{\pgfqpoint{2.237745in}{1.720319in}}%
\pgfpathlineto{\pgfqpoint{2.274150in}{1.748924in}}%
\pgfpathlineto{\pgfqpoint{2.310521in}{1.779069in}}%
\pgfpathlineto{\pgfqpoint{2.346862in}{1.810468in}}%
\pgfpathlineto{\pgfqpoint{2.383154in}{1.843393in}}%
\pgfpathlineto{\pgfqpoint{2.418733in}{1.877140in}}%
\pgfpathlineto{\pgfqpoint{2.451632in}{1.909815in}}%
\pgfpathlineto{\pgfqpoint{2.481398in}{1.941108in}}%
\pgfpathlineto{\pgfqpoint{2.508530in}{1.971059in}}%
\pgfpathlineto{\pgfqpoint{2.549016in}{2.017891in}}%
\pgfpathlineto{\pgfqpoint{2.591455in}{2.068845in}}%
\pgfpathlineto{\pgfqpoint{2.677323in}{2.175117in}}%
\pgfpathlineto{\pgfqpoint{2.719261in}{2.227290in}}%
\pgfpathlineto{\pgfqpoint{2.736206in}{2.246830in}}%
\pgfpathlineto{\pgfqpoint{2.748327in}{2.259625in}}%
\pgfpathlineto{\pgfqpoint{2.759988in}{2.270603in}}%
\pgfpathlineto{\pgfqpoint{2.771222in}{2.279759in}}%
\pgfpathlineto{\pgfqpoint{2.782059in}{2.287024in}}%
\pgfpathlineto{\pgfqpoint{2.792526in}{2.292444in}}%
\pgfpathlineto{\pgfqpoint{2.802648in}{2.296224in}}%
\pgfpathlineto{\pgfqpoint{2.812446in}{2.298639in}}%
\pgfpathlineto{\pgfqpoint{2.821942in}{2.299961in}}%
\pgfpathlineto{\pgfqpoint{2.834162in}{2.300532in}}%
\pgfpathlineto{\pgfqpoint{2.851625in}{2.299951in}}%
\pgfpathlineto{\pgfqpoint{2.876087in}{2.297743in}}%
\pgfpathlineto{\pgfqpoint{2.922095in}{2.292183in}}%
\pgfpathlineto{\pgfqpoint{2.973715in}{2.284731in}}%
\pgfpathlineto{\pgfqpoint{3.022787in}{2.276480in}}%
\pgfpathlineto{\pgfqpoint{3.043110in}{2.272684in}}%
\pgfpathlineto{\pgfqpoint{3.043110in}{2.272684in}}%
\pgfusepath{stroke}%
\end{pgfscope}%
\begin{pgfscope}%
\pgfsetrectcap%
\pgfsetmiterjoin%
\pgfsetlinewidth{0.803000pt}%
\definecolor{currentstroke}{rgb}{0.000000,0.000000,0.000000}%
\pgfsetstrokecolor{currentstroke}%
\pgfsetdash{}{0pt}%
\pgfpathmoveto{\pgfqpoint{0.650000in}{0.420000in}}%
\pgfpathlineto{\pgfqpoint{0.650000in}{2.414488in}}%
\pgfusepath{stroke}%
\end{pgfscope}%
\begin{pgfscope}%
\pgfsetrectcap%
\pgfsetmiterjoin%
\pgfsetlinewidth{0.803000pt}%
\definecolor{currentstroke}{rgb}{0.000000,0.000000,0.000000}%
\pgfsetstrokecolor{currentstroke}%
\pgfsetdash{}{0pt}%
\pgfpathmoveto{\pgfqpoint{3.038110in}{0.420000in}}%
\pgfpathlineto{\pgfqpoint{3.038110in}{2.414488in}}%
\pgfusepath{stroke}%
\end{pgfscope}%
\begin{pgfscope}%
\pgfsetrectcap%
\pgfsetmiterjoin%
\pgfsetlinewidth{0.803000pt}%
\definecolor{currentstroke}{rgb}{0.000000,0.000000,0.000000}%
\pgfsetstrokecolor{currentstroke}%
\pgfsetdash{}{0pt}%
\pgfpathmoveto{\pgfqpoint{0.650000in}{0.420000in}}%
\pgfpathlineto{\pgfqpoint{3.038110in}{0.420000in}}%
\pgfusepath{stroke}%
\end{pgfscope}%
\begin{pgfscope}%
\pgfsetrectcap%
\pgfsetmiterjoin%
\pgfsetlinewidth{0.803000pt}%
\definecolor{currentstroke}{rgb}{0.000000,0.000000,0.000000}%
\pgfsetstrokecolor{currentstroke}%
\pgfsetdash{}{0pt}%
\pgfpathmoveto{\pgfqpoint{0.650000in}{2.414488in}}%
\pgfpathlineto{\pgfqpoint{3.038110in}{2.414488in}}%
\pgfusepath{stroke}%
\end{pgfscope}%
\begin{pgfscope}%
\pgfsetrectcap%
\pgfsetroundjoin%
\pgfsetlinewidth{0.602250pt}%
\definecolor{currentstroke}{rgb}{0.000000,0.000000,0.000000}%
\pgfsetstrokecolor{currentstroke}%
\pgfsetdash{}{0pt}%
\pgfpathmoveto{\pgfqpoint{0.750000in}{2.275599in}}%
\pgfpathlineto{\pgfqpoint{0.861111in}{2.275599in}}%
\pgfpathlineto{\pgfqpoint{0.972222in}{2.275599in}}%
\pgfusepath{stroke}%
\end{pgfscope}%
\begin{pgfscope}%
\definecolor{textcolor}{rgb}{0.000000,0.000000,0.000000}%
\pgfsetstrokecolor{textcolor}%
\pgfsetfillcolor{textcolor}%
\pgftext[x=1.061111in,y=2.236710in,left,base]{\color{textcolor}{\rmfamily\fontsize{8.000000}{9.600000}\selectfont\catcode`\^=\active\def^{\ifmmode\sp\else\^{}\fi}\catcode`\%=\active\def%{\%}ref}}%
\end{pgfscope}%
\begin{pgfscope}%
\pgfsetrectcap%
\pgfsetroundjoin%
\pgfsetlinewidth{1.003750pt}%
\definecolor{currentstroke}{rgb}{0.870588,0.466667,0.682353}%
\pgfsetstrokecolor{currentstroke}%
\pgfsetdash{}{0pt}%
\pgfpathmoveto{\pgfqpoint{0.750000in}{2.120661in}}%
\pgfpathlineto{\pgfqpoint{0.861111in}{2.120661in}}%
\pgfpathlineto{\pgfqpoint{0.972222in}{2.120661in}}%
\pgfusepath{stroke}%
\end{pgfscope}%
\begin{pgfscope}%
\definecolor{textcolor}{rgb}{0.000000,0.000000,0.000000}%
\pgfsetstrokecolor{textcolor}%
\pgfsetfillcolor{textcolor}%
\pgftext[x=1.061111in,y=2.081772in,left,base]{\color{textcolor}{\rmfamily\fontsize{8.000000}{9.600000}\selectfont\catcode`\^=\active\def^{\ifmmode\sp\else\^{}\fi}\catcode`\%=\active\def%{\%}base}}%
\end{pgfscope}%
\begin{pgfscope}%
\pgfsetrectcap%
\pgfsetroundjoin%
\pgfsetlinewidth{1.003750pt}%
\definecolor{currentstroke}{rgb}{0.498039,0.737255,0.254902}%
\pgfsetstrokecolor{currentstroke}%
\pgfsetdash{}{0pt}%
\pgfpathmoveto{\pgfqpoint{0.750000in}{1.965723in}}%
\pgfpathlineto{\pgfqpoint{0.861111in}{1.965723in}}%
\pgfpathlineto{\pgfqpoint{0.972222in}{1.965723in}}%
\pgfusepath{stroke}%
\end{pgfscope}%
\begin{pgfscope}%
\definecolor{textcolor}{rgb}{0.000000,0.000000,0.000000}%
\pgfsetstrokecolor{textcolor}%
\pgfsetfillcolor{textcolor}%
\pgftext[x=1.061111in,y=1.926834in,left,base]{\color{textcolor}{\rmfamily\fontsize{8.000000}{9.600000}\selectfont\catcode`\^=\active\def^{\ifmmode\sp\else\^{}\fi}\catcode`\%=\active\def%{\%}optimised}}%
\end{pgfscope}%
\end{pgfpicture}%
\makeatother%
\endgroup%

    \label{fig:R100deg30Re1200U}
\end{subfigure}
\hfill
\begin{subfigure}[t]{0.495\textwidth}
    \centering
    \subcaption{}
    %% Creator: Matplotlib, PGF backend
%%
%% To include the figure in your LaTeX document, write
%%   \input{<filename>.pgf}
%%
%% Make sure the required packages are loaded in your preamble
%%   \usepackage{pgf}
%%
%% Also ensure that all the required font packages are loaded; for instance,
%% the lmodern package is sometimes necessary when using math font.
%%   \usepackage{lmodern}
%%
%% Figures using additional raster images can only be included by \input if
%% they are in the same directory as the main LaTeX file. For loading figures
%% from other directories you can use the `import` package
%%   \usepackage{import}
%%
%% and then include the figures with
%%   \import{<path to file>}{<filename>.pgf}
%%
%% Matplotlib used the following preamble
%%   \def\mathdefault#1{#1}
%%   \everymath=\expandafter{\the\everymath\displaystyle}
%%   \IfFileExists{scrextend.sty}{
%%     \usepackage[fontsize=11.000000pt]{scrextend}
%%   }{
%%     \renewcommand{\normalsize}{\fontsize{11.000000}{13.200000}\selectfont}
%%     \normalsize
%%   }
%%   
%%   \makeatletter\@ifpackageloaded{underscore}{}{\usepackage[strings]{underscore}}\makeatother
%%
\begingroup%
\makeatletter%
\begin{pgfpicture}%
\pgfpathrectangle{\pgfpointorigin}{\pgfqpoint{3.118110in}{2.494488in}}%
\pgfusepath{use as bounding box, clip}%
\begin{pgfscope}%
\pgfsetbuttcap%
\pgfsetmiterjoin%
\definecolor{currentfill}{rgb}{1.000000,1.000000,1.000000}%
\pgfsetfillcolor{currentfill}%
\pgfsetlinewidth{0.000000pt}%
\definecolor{currentstroke}{rgb}{1.000000,1.000000,1.000000}%
\pgfsetstrokecolor{currentstroke}%
\pgfsetdash{}{0pt}%
\pgfpathmoveto{\pgfqpoint{0.000000in}{0.000000in}}%
\pgfpathlineto{\pgfqpoint{3.118110in}{0.000000in}}%
\pgfpathlineto{\pgfqpoint{3.118110in}{2.494488in}}%
\pgfpathlineto{\pgfqpoint{0.000000in}{2.494488in}}%
\pgfpathlineto{\pgfqpoint{0.000000in}{0.000000in}}%
\pgfpathclose%
\pgfusepath{fill}%
\end{pgfscope}%
\begin{pgfscope}%
\pgfsetbuttcap%
\pgfsetmiterjoin%
\definecolor{currentfill}{rgb}{1.000000,1.000000,1.000000}%
\pgfsetfillcolor{currentfill}%
\pgfsetlinewidth{0.000000pt}%
\definecolor{currentstroke}{rgb}{0.000000,0.000000,0.000000}%
\pgfsetstrokecolor{currentstroke}%
\pgfsetstrokeopacity{0.000000}%
\pgfsetdash{}{0pt}%
\pgfpathmoveto{\pgfqpoint{0.650000in}{0.420000in}}%
\pgfpathlineto{\pgfqpoint{3.038110in}{0.420000in}}%
\pgfpathlineto{\pgfqpoint{3.038110in}{2.414488in}}%
\pgfpathlineto{\pgfqpoint{0.650000in}{2.414488in}}%
\pgfpathlineto{\pgfqpoint{0.650000in}{0.420000in}}%
\pgfpathclose%
\pgfusepath{fill}%
\end{pgfscope}%
\begin{pgfscope}%
\pgfsetbuttcap%
\pgfsetroundjoin%
\definecolor{currentfill}{rgb}{0.000000,0.000000,0.000000}%
\pgfsetfillcolor{currentfill}%
\pgfsetlinewidth{0.501875pt}%
\definecolor{currentstroke}{rgb}{0.000000,0.000000,0.000000}%
\pgfsetstrokecolor{currentstroke}%
\pgfsetdash{}{0pt}%
\pgfsys@defobject{currentmarker}{\pgfqpoint{0.000000in}{0.000000in}}{\pgfqpoint{0.000000in}{0.041667in}}{%
\pgfpathmoveto{\pgfqpoint{0.000000in}{0.000000in}}%
\pgfpathlineto{\pgfqpoint{0.000000in}{0.041667in}}%
\pgfusepath{stroke,fill}%
}%
\begin{pgfscope}%
\pgfsys@transformshift{0.862677in}{0.420000in}%
\pgfsys@useobject{currentmarker}{}%
\end{pgfscope}%
\end{pgfscope}%
\begin{pgfscope}%
\pgfsetbuttcap%
\pgfsetroundjoin%
\definecolor{currentfill}{rgb}{0.000000,0.000000,0.000000}%
\pgfsetfillcolor{currentfill}%
\pgfsetlinewidth{0.501875pt}%
\definecolor{currentstroke}{rgb}{0.000000,0.000000,0.000000}%
\pgfsetstrokecolor{currentstroke}%
\pgfsetdash{}{0pt}%
\pgfsys@defobject{currentmarker}{\pgfqpoint{0.000000in}{-0.041667in}}{\pgfqpoint{0.000000in}{0.000000in}}{%
\pgfpathmoveto{\pgfqpoint{0.000000in}{0.000000in}}%
\pgfpathlineto{\pgfqpoint{0.000000in}{-0.041667in}}%
\pgfusepath{stroke,fill}%
}%
\begin{pgfscope}%
\pgfsys@transformshift{0.862677in}{2.414488in}%
\pgfsys@useobject{currentmarker}{}%
\end{pgfscope}%
\end{pgfscope}%
\begin{pgfscope}%
\definecolor{textcolor}{rgb}{0.000000,0.000000,0.000000}%
\pgfsetstrokecolor{textcolor}%
\pgfsetfillcolor{textcolor}%
\pgftext[x=0.862677in,y=0.371389in,,top]{\color{textcolor}{\rmfamily\fontsize{11.000000}{13.200000}\selectfont\catcode`\^=\active\def^{\ifmmode\sp\else\^{}\fi}\catcode`\%=\active\def%{\%}$\mathdefault{10^{0}}$}}%
\end{pgfscope}%
\begin{pgfscope}%
\pgfsetbuttcap%
\pgfsetroundjoin%
\definecolor{currentfill}{rgb}{0.000000,0.000000,0.000000}%
\pgfsetfillcolor{currentfill}%
\pgfsetlinewidth{0.501875pt}%
\definecolor{currentstroke}{rgb}{0.000000,0.000000,0.000000}%
\pgfsetstrokecolor{currentstroke}%
\pgfsetdash{}{0pt}%
\pgfsys@defobject{currentmarker}{\pgfqpoint{0.000000in}{0.000000in}}{\pgfqpoint{0.000000in}{0.041667in}}{%
\pgfpathmoveto{\pgfqpoint{0.000000in}{0.000000in}}%
\pgfpathlineto{\pgfqpoint{0.000000in}{0.041667in}}%
\pgfusepath{stroke,fill}%
}%
\begin{pgfscope}%
\pgfsys@transformshift{1.569174in}{0.420000in}%
\pgfsys@useobject{currentmarker}{}%
\end{pgfscope}%
\end{pgfscope}%
\begin{pgfscope}%
\pgfsetbuttcap%
\pgfsetroundjoin%
\definecolor{currentfill}{rgb}{0.000000,0.000000,0.000000}%
\pgfsetfillcolor{currentfill}%
\pgfsetlinewidth{0.501875pt}%
\definecolor{currentstroke}{rgb}{0.000000,0.000000,0.000000}%
\pgfsetstrokecolor{currentstroke}%
\pgfsetdash{}{0pt}%
\pgfsys@defobject{currentmarker}{\pgfqpoint{0.000000in}{-0.041667in}}{\pgfqpoint{0.000000in}{0.000000in}}{%
\pgfpathmoveto{\pgfqpoint{0.000000in}{0.000000in}}%
\pgfpathlineto{\pgfqpoint{0.000000in}{-0.041667in}}%
\pgfusepath{stroke,fill}%
}%
\begin{pgfscope}%
\pgfsys@transformshift{1.569174in}{2.414488in}%
\pgfsys@useobject{currentmarker}{}%
\end{pgfscope}%
\end{pgfscope}%
\begin{pgfscope}%
\definecolor{textcolor}{rgb}{0.000000,0.000000,0.000000}%
\pgfsetstrokecolor{textcolor}%
\pgfsetfillcolor{textcolor}%
\pgftext[x=1.569174in,y=0.371389in,,top]{\color{textcolor}{\rmfamily\fontsize{11.000000}{13.200000}\selectfont\catcode`\^=\active\def^{\ifmmode\sp\else\^{}\fi}\catcode`\%=\active\def%{\%}$\mathdefault{10^{1}}$}}%
\end{pgfscope}%
\begin{pgfscope}%
\pgfsetbuttcap%
\pgfsetroundjoin%
\definecolor{currentfill}{rgb}{0.000000,0.000000,0.000000}%
\pgfsetfillcolor{currentfill}%
\pgfsetlinewidth{0.501875pt}%
\definecolor{currentstroke}{rgb}{0.000000,0.000000,0.000000}%
\pgfsetstrokecolor{currentstroke}%
\pgfsetdash{}{0pt}%
\pgfsys@defobject{currentmarker}{\pgfqpoint{0.000000in}{0.000000in}}{\pgfqpoint{0.000000in}{0.041667in}}{%
\pgfpathmoveto{\pgfqpoint{0.000000in}{0.000000in}}%
\pgfpathlineto{\pgfqpoint{0.000000in}{0.041667in}}%
\pgfusepath{stroke,fill}%
}%
\begin{pgfscope}%
\pgfsys@transformshift{2.275671in}{0.420000in}%
\pgfsys@useobject{currentmarker}{}%
\end{pgfscope}%
\end{pgfscope}%
\begin{pgfscope}%
\pgfsetbuttcap%
\pgfsetroundjoin%
\definecolor{currentfill}{rgb}{0.000000,0.000000,0.000000}%
\pgfsetfillcolor{currentfill}%
\pgfsetlinewidth{0.501875pt}%
\definecolor{currentstroke}{rgb}{0.000000,0.000000,0.000000}%
\pgfsetstrokecolor{currentstroke}%
\pgfsetdash{}{0pt}%
\pgfsys@defobject{currentmarker}{\pgfqpoint{0.000000in}{-0.041667in}}{\pgfqpoint{0.000000in}{0.000000in}}{%
\pgfpathmoveto{\pgfqpoint{0.000000in}{0.000000in}}%
\pgfpathlineto{\pgfqpoint{0.000000in}{-0.041667in}}%
\pgfusepath{stroke,fill}%
}%
\begin{pgfscope}%
\pgfsys@transformshift{2.275671in}{2.414488in}%
\pgfsys@useobject{currentmarker}{}%
\end{pgfscope}%
\end{pgfscope}%
\begin{pgfscope}%
\definecolor{textcolor}{rgb}{0.000000,0.000000,0.000000}%
\pgfsetstrokecolor{textcolor}%
\pgfsetfillcolor{textcolor}%
\pgftext[x=2.275671in,y=0.371389in,,top]{\color{textcolor}{\rmfamily\fontsize{11.000000}{13.200000}\selectfont\catcode`\^=\active\def^{\ifmmode\sp\else\^{}\fi}\catcode`\%=\active\def%{\%}$\mathdefault{10^{2}}$}}%
\end{pgfscope}%
\begin{pgfscope}%
\pgfsetbuttcap%
\pgfsetroundjoin%
\definecolor{currentfill}{rgb}{0.000000,0.000000,0.000000}%
\pgfsetfillcolor{currentfill}%
\pgfsetlinewidth{0.501875pt}%
\definecolor{currentstroke}{rgb}{0.000000,0.000000,0.000000}%
\pgfsetstrokecolor{currentstroke}%
\pgfsetdash{}{0pt}%
\pgfsys@defobject{currentmarker}{\pgfqpoint{0.000000in}{0.000000in}}{\pgfqpoint{0.000000in}{0.041667in}}{%
\pgfpathmoveto{\pgfqpoint{0.000000in}{0.000000in}}%
\pgfpathlineto{\pgfqpoint{0.000000in}{0.041667in}}%
\pgfusepath{stroke,fill}%
}%
\begin{pgfscope}%
\pgfsys@transformshift{2.982169in}{0.420000in}%
\pgfsys@useobject{currentmarker}{}%
\end{pgfscope}%
\end{pgfscope}%
\begin{pgfscope}%
\pgfsetbuttcap%
\pgfsetroundjoin%
\definecolor{currentfill}{rgb}{0.000000,0.000000,0.000000}%
\pgfsetfillcolor{currentfill}%
\pgfsetlinewidth{0.501875pt}%
\definecolor{currentstroke}{rgb}{0.000000,0.000000,0.000000}%
\pgfsetstrokecolor{currentstroke}%
\pgfsetdash{}{0pt}%
\pgfsys@defobject{currentmarker}{\pgfqpoint{0.000000in}{-0.041667in}}{\pgfqpoint{0.000000in}{0.000000in}}{%
\pgfpathmoveto{\pgfqpoint{0.000000in}{0.000000in}}%
\pgfpathlineto{\pgfqpoint{0.000000in}{-0.041667in}}%
\pgfusepath{stroke,fill}%
}%
\begin{pgfscope}%
\pgfsys@transformshift{2.982169in}{2.414488in}%
\pgfsys@useobject{currentmarker}{}%
\end{pgfscope}%
\end{pgfscope}%
\begin{pgfscope}%
\definecolor{textcolor}{rgb}{0.000000,0.000000,0.000000}%
\pgfsetstrokecolor{textcolor}%
\pgfsetfillcolor{textcolor}%
\pgftext[x=2.982169in,y=0.371389in,,top]{\color{textcolor}{\rmfamily\fontsize{11.000000}{13.200000}\selectfont\catcode`\^=\active\def^{\ifmmode\sp\else\^{}\fi}\catcode`\%=\active\def%{\%}$\mathdefault{10^{3}}$}}%
\end{pgfscope}%
\begin{pgfscope}%
\pgfsetbuttcap%
\pgfsetroundjoin%
\definecolor{currentfill}{rgb}{0.000000,0.000000,0.000000}%
\pgfsetfillcolor{currentfill}%
\pgfsetlinewidth{0.401500pt}%
\definecolor{currentstroke}{rgb}{0.000000,0.000000,0.000000}%
\pgfsetstrokecolor{currentstroke}%
\pgfsetdash{}{0pt}%
\pgfsys@defobject{currentmarker}{\pgfqpoint{0.000000in}{0.000000in}}{\pgfqpoint{0.000000in}{0.020833in}}{%
\pgfpathmoveto{\pgfqpoint{0.000000in}{0.000000in}}%
\pgfpathlineto{\pgfqpoint{0.000000in}{0.020833in}}%
\pgfusepath{stroke,fill}%
}%
\begin{pgfscope}%
\pgfsys@transformshift{0.650000in}{0.420000in}%
\pgfsys@useobject{currentmarker}{}%
\end{pgfscope}%
\end{pgfscope}%
\begin{pgfscope}%
\pgfsetbuttcap%
\pgfsetroundjoin%
\definecolor{currentfill}{rgb}{0.000000,0.000000,0.000000}%
\pgfsetfillcolor{currentfill}%
\pgfsetlinewidth{0.401500pt}%
\definecolor{currentstroke}{rgb}{0.000000,0.000000,0.000000}%
\pgfsetstrokecolor{currentstroke}%
\pgfsetdash{}{0pt}%
\pgfsys@defobject{currentmarker}{\pgfqpoint{0.000000in}{-0.020833in}}{\pgfqpoint{0.000000in}{0.000000in}}{%
\pgfpathmoveto{\pgfqpoint{0.000000in}{0.000000in}}%
\pgfpathlineto{\pgfqpoint{0.000000in}{-0.020833in}}%
\pgfusepath{stroke,fill}%
}%
\begin{pgfscope}%
\pgfsys@transformshift{0.650000in}{2.414488in}%
\pgfsys@useobject{currentmarker}{}%
\end{pgfscope}%
\end{pgfscope}%
\begin{pgfscope}%
\pgfsetbuttcap%
\pgfsetroundjoin%
\definecolor{currentfill}{rgb}{0.000000,0.000000,0.000000}%
\pgfsetfillcolor{currentfill}%
\pgfsetlinewidth{0.401500pt}%
\definecolor{currentstroke}{rgb}{0.000000,0.000000,0.000000}%
\pgfsetstrokecolor{currentstroke}%
\pgfsetdash{}{0pt}%
\pgfsys@defobject{currentmarker}{\pgfqpoint{0.000000in}{0.000000in}}{\pgfqpoint{0.000000in}{0.020833in}}{%
\pgfpathmoveto{\pgfqpoint{0.000000in}{0.000000in}}%
\pgfpathlineto{\pgfqpoint{0.000000in}{0.020833in}}%
\pgfusepath{stroke,fill}%
}%
\begin{pgfscope}%
\pgfsys@transformshift{0.705941in}{0.420000in}%
\pgfsys@useobject{currentmarker}{}%
\end{pgfscope}%
\end{pgfscope}%
\begin{pgfscope}%
\pgfsetbuttcap%
\pgfsetroundjoin%
\definecolor{currentfill}{rgb}{0.000000,0.000000,0.000000}%
\pgfsetfillcolor{currentfill}%
\pgfsetlinewidth{0.401500pt}%
\definecolor{currentstroke}{rgb}{0.000000,0.000000,0.000000}%
\pgfsetstrokecolor{currentstroke}%
\pgfsetdash{}{0pt}%
\pgfsys@defobject{currentmarker}{\pgfqpoint{0.000000in}{-0.020833in}}{\pgfqpoint{0.000000in}{0.000000in}}{%
\pgfpathmoveto{\pgfqpoint{0.000000in}{0.000000in}}%
\pgfpathlineto{\pgfqpoint{0.000000in}{-0.020833in}}%
\pgfusepath{stroke,fill}%
}%
\begin{pgfscope}%
\pgfsys@transformshift{0.705941in}{2.414488in}%
\pgfsys@useobject{currentmarker}{}%
\end{pgfscope}%
\end{pgfscope}%
\begin{pgfscope}%
\pgfsetbuttcap%
\pgfsetroundjoin%
\definecolor{currentfill}{rgb}{0.000000,0.000000,0.000000}%
\pgfsetfillcolor{currentfill}%
\pgfsetlinewidth{0.401500pt}%
\definecolor{currentstroke}{rgb}{0.000000,0.000000,0.000000}%
\pgfsetstrokecolor{currentstroke}%
\pgfsetdash{}{0pt}%
\pgfsys@defobject{currentmarker}{\pgfqpoint{0.000000in}{0.000000in}}{\pgfqpoint{0.000000in}{0.020833in}}{%
\pgfpathmoveto{\pgfqpoint{0.000000in}{0.000000in}}%
\pgfpathlineto{\pgfqpoint{0.000000in}{0.020833in}}%
\pgfusepath{stroke,fill}%
}%
\begin{pgfscope}%
\pgfsys@transformshift{0.753239in}{0.420000in}%
\pgfsys@useobject{currentmarker}{}%
\end{pgfscope}%
\end{pgfscope}%
\begin{pgfscope}%
\pgfsetbuttcap%
\pgfsetroundjoin%
\definecolor{currentfill}{rgb}{0.000000,0.000000,0.000000}%
\pgfsetfillcolor{currentfill}%
\pgfsetlinewidth{0.401500pt}%
\definecolor{currentstroke}{rgb}{0.000000,0.000000,0.000000}%
\pgfsetstrokecolor{currentstroke}%
\pgfsetdash{}{0pt}%
\pgfsys@defobject{currentmarker}{\pgfqpoint{0.000000in}{-0.020833in}}{\pgfqpoint{0.000000in}{0.000000in}}{%
\pgfpathmoveto{\pgfqpoint{0.000000in}{0.000000in}}%
\pgfpathlineto{\pgfqpoint{0.000000in}{-0.020833in}}%
\pgfusepath{stroke,fill}%
}%
\begin{pgfscope}%
\pgfsys@transformshift{0.753239in}{2.414488in}%
\pgfsys@useobject{currentmarker}{}%
\end{pgfscope}%
\end{pgfscope}%
\begin{pgfscope}%
\pgfsetbuttcap%
\pgfsetroundjoin%
\definecolor{currentfill}{rgb}{0.000000,0.000000,0.000000}%
\pgfsetfillcolor{currentfill}%
\pgfsetlinewidth{0.401500pt}%
\definecolor{currentstroke}{rgb}{0.000000,0.000000,0.000000}%
\pgfsetstrokecolor{currentstroke}%
\pgfsetdash{}{0pt}%
\pgfsys@defobject{currentmarker}{\pgfqpoint{0.000000in}{0.000000in}}{\pgfqpoint{0.000000in}{0.020833in}}{%
\pgfpathmoveto{\pgfqpoint{0.000000in}{0.000000in}}%
\pgfpathlineto{\pgfqpoint{0.000000in}{0.020833in}}%
\pgfusepath{stroke,fill}%
}%
\begin{pgfscope}%
\pgfsys@transformshift{0.794210in}{0.420000in}%
\pgfsys@useobject{currentmarker}{}%
\end{pgfscope}%
\end{pgfscope}%
\begin{pgfscope}%
\pgfsetbuttcap%
\pgfsetroundjoin%
\definecolor{currentfill}{rgb}{0.000000,0.000000,0.000000}%
\pgfsetfillcolor{currentfill}%
\pgfsetlinewidth{0.401500pt}%
\definecolor{currentstroke}{rgb}{0.000000,0.000000,0.000000}%
\pgfsetstrokecolor{currentstroke}%
\pgfsetdash{}{0pt}%
\pgfsys@defobject{currentmarker}{\pgfqpoint{0.000000in}{-0.020833in}}{\pgfqpoint{0.000000in}{0.000000in}}{%
\pgfpathmoveto{\pgfqpoint{0.000000in}{0.000000in}}%
\pgfpathlineto{\pgfqpoint{0.000000in}{-0.020833in}}%
\pgfusepath{stroke,fill}%
}%
\begin{pgfscope}%
\pgfsys@transformshift{0.794210in}{2.414488in}%
\pgfsys@useobject{currentmarker}{}%
\end{pgfscope}%
\end{pgfscope}%
\begin{pgfscope}%
\pgfsetbuttcap%
\pgfsetroundjoin%
\definecolor{currentfill}{rgb}{0.000000,0.000000,0.000000}%
\pgfsetfillcolor{currentfill}%
\pgfsetlinewidth{0.401500pt}%
\definecolor{currentstroke}{rgb}{0.000000,0.000000,0.000000}%
\pgfsetstrokecolor{currentstroke}%
\pgfsetdash{}{0pt}%
\pgfsys@defobject{currentmarker}{\pgfqpoint{0.000000in}{0.000000in}}{\pgfqpoint{0.000000in}{0.020833in}}{%
\pgfpathmoveto{\pgfqpoint{0.000000in}{0.000000in}}%
\pgfpathlineto{\pgfqpoint{0.000000in}{0.020833in}}%
\pgfusepath{stroke,fill}%
}%
\begin{pgfscope}%
\pgfsys@transformshift{0.830349in}{0.420000in}%
\pgfsys@useobject{currentmarker}{}%
\end{pgfscope}%
\end{pgfscope}%
\begin{pgfscope}%
\pgfsetbuttcap%
\pgfsetroundjoin%
\definecolor{currentfill}{rgb}{0.000000,0.000000,0.000000}%
\pgfsetfillcolor{currentfill}%
\pgfsetlinewidth{0.401500pt}%
\definecolor{currentstroke}{rgb}{0.000000,0.000000,0.000000}%
\pgfsetstrokecolor{currentstroke}%
\pgfsetdash{}{0pt}%
\pgfsys@defobject{currentmarker}{\pgfqpoint{0.000000in}{-0.020833in}}{\pgfqpoint{0.000000in}{0.000000in}}{%
\pgfpathmoveto{\pgfqpoint{0.000000in}{0.000000in}}%
\pgfpathlineto{\pgfqpoint{0.000000in}{-0.020833in}}%
\pgfusepath{stroke,fill}%
}%
\begin{pgfscope}%
\pgfsys@transformshift{0.830349in}{2.414488in}%
\pgfsys@useobject{currentmarker}{}%
\end{pgfscope}%
\end{pgfscope}%
\begin{pgfscope}%
\pgfsetbuttcap%
\pgfsetroundjoin%
\definecolor{currentfill}{rgb}{0.000000,0.000000,0.000000}%
\pgfsetfillcolor{currentfill}%
\pgfsetlinewidth{0.401500pt}%
\definecolor{currentstroke}{rgb}{0.000000,0.000000,0.000000}%
\pgfsetstrokecolor{currentstroke}%
\pgfsetdash{}{0pt}%
\pgfsys@defobject{currentmarker}{\pgfqpoint{0.000000in}{0.000000in}}{\pgfqpoint{0.000000in}{0.020833in}}{%
\pgfpathmoveto{\pgfqpoint{0.000000in}{0.000000in}}%
\pgfpathlineto{\pgfqpoint{0.000000in}{0.020833in}}%
\pgfusepath{stroke,fill}%
}%
\begin{pgfscope}%
\pgfsys@transformshift{1.075354in}{0.420000in}%
\pgfsys@useobject{currentmarker}{}%
\end{pgfscope}%
\end{pgfscope}%
\begin{pgfscope}%
\pgfsetbuttcap%
\pgfsetroundjoin%
\definecolor{currentfill}{rgb}{0.000000,0.000000,0.000000}%
\pgfsetfillcolor{currentfill}%
\pgfsetlinewidth{0.401500pt}%
\definecolor{currentstroke}{rgb}{0.000000,0.000000,0.000000}%
\pgfsetstrokecolor{currentstroke}%
\pgfsetdash{}{0pt}%
\pgfsys@defobject{currentmarker}{\pgfqpoint{0.000000in}{-0.020833in}}{\pgfqpoint{0.000000in}{0.000000in}}{%
\pgfpathmoveto{\pgfqpoint{0.000000in}{0.000000in}}%
\pgfpathlineto{\pgfqpoint{0.000000in}{-0.020833in}}%
\pgfusepath{stroke,fill}%
}%
\begin{pgfscope}%
\pgfsys@transformshift{1.075354in}{2.414488in}%
\pgfsys@useobject{currentmarker}{}%
\end{pgfscope}%
\end{pgfscope}%
\begin{pgfscope}%
\pgfsetbuttcap%
\pgfsetroundjoin%
\definecolor{currentfill}{rgb}{0.000000,0.000000,0.000000}%
\pgfsetfillcolor{currentfill}%
\pgfsetlinewidth{0.401500pt}%
\definecolor{currentstroke}{rgb}{0.000000,0.000000,0.000000}%
\pgfsetstrokecolor{currentstroke}%
\pgfsetdash{}{0pt}%
\pgfsys@defobject{currentmarker}{\pgfqpoint{0.000000in}{0.000000in}}{\pgfqpoint{0.000000in}{0.020833in}}{%
\pgfpathmoveto{\pgfqpoint{0.000000in}{0.000000in}}%
\pgfpathlineto{\pgfqpoint{0.000000in}{0.020833in}}%
\pgfusepath{stroke,fill}%
}%
\begin{pgfscope}%
\pgfsys@transformshift{1.199762in}{0.420000in}%
\pgfsys@useobject{currentmarker}{}%
\end{pgfscope}%
\end{pgfscope}%
\begin{pgfscope}%
\pgfsetbuttcap%
\pgfsetroundjoin%
\definecolor{currentfill}{rgb}{0.000000,0.000000,0.000000}%
\pgfsetfillcolor{currentfill}%
\pgfsetlinewidth{0.401500pt}%
\definecolor{currentstroke}{rgb}{0.000000,0.000000,0.000000}%
\pgfsetstrokecolor{currentstroke}%
\pgfsetdash{}{0pt}%
\pgfsys@defobject{currentmarker}{\pgfqpoint{0.000000in}{-0.020833in}}{\pgfqpoint{0.000000in}{0.000000in}}{%
\pgfpathmoveto{\pgfqpoint{0.000000in}{0.000000in}}%
\pgfpathlineto{\pgfqpoint{0.000000in}{-0.020833in}}%
\pgfusepath{stroke,fill}%
}%
\begin{pgfscope}%
\pgfsys@transformshift{1.199762in}{2.414488in}%
\pgfsys@useobject{currentmarker}{}%
\end{pgfscope}%
\end{pgfscope}%
\begin{pgfscope}%
\pgfsetbuttcap%
\pgfsetroundjoin%
\definecolor{currentfill}{rgb}{0.000000,0.000000,0.000000}%
\pgfsetfillcolor{currentfill}%
\pgfsetlinewidth{0.401500pt}%
\definecolor{currentstroke}{rgb}{0.000000,0.000000,0.000000}%
\pgfsetstrokecolor{currentstroke}%
\pgfsetdash{}{0pt}%
\pgfsys@defobject{currentmarker}{\pgfqpoint{0.000000in}{0.000000in}}{\pgfqpoint{0.000000in}{0.020833in}}{%
\pgfpathmoveto{\pgfqpoint{0.000000in}{0.000000in}}%
\pgfpathlineto{\pgfqpoint{0.000000in}{0.020833in}}%
\pgfusepath{stroke,fill}%
}%
\begin{pgfscope}%
\pgfsys@transformshift{1.288031in}{0.420000in}%
\pgfsys@useobject{currentmarker}{}%
\end{pgfscope}%
\end{pgfscope}%
\begin{pgfscope}%
\pgfsetbuttcap%
\pgfsetroundjoin%
\definecolor{currentfill}{rgb}{0.000000,0.000000,0.000000}%
\pgfsetfillcolor{currentfill}%
\pgfsetlinewidth{0.401500pt}%
\definecolor{currentstroke}{rgb}{0.000000,0.000000,0.000000}%
\pgfsetstrokecolor{currentstroke}%
\pgfsetdash{}{0pt}%
\pgfsys@defobject{currentmarker}{\pgfqpoint{0.000000in}{-0.020833in}}{\pgfqpoint{0.000000in}{0.000000in}}{%
\pgfpathmoveto{\pgfqpoint{0.000000in}{0.000000in}}%
\pgfpathlineto{\pgfqpoint{0.000000in}{-0.020833in}}%
\pgfusepath{stroke,fill}%
}%
\begin{pgfscope}%
\pgfsys@transformshift{1.288031in}{2.414488in}%
\pgfsys@useobject{currentmarker}{}%
\end{pgfscope}%
\end{pgfscope}%
\begin{pgfscope}%
\pgfsetbuttcap%
\pgfsetroundjoin%
\definecolor{currentfill}{rgb}{0.000000,0.000000,0.000000}%
\pgfsetfillcolor{currentfill}%
\pgfsetlinewidth{0.401500pt}%
\definecolor{currentstroke}{rgb}{0.000000,0.000000,0.000000}%
\pgfsetstrokecolor{currentstroke}%
\pgfsetdash{}{0pt}%
\pgfsys@defobject{currentmarker}{\pgfqpoint{0.000000in}{0.000000in}}{\pgfqpoint{0.000000in}{0.020833in}}{%
\pgfpathmoveto{\pgfqpoint{0.000000in}{0.000000in}}%
\pgfpathlineto{\pgfqpoint{0.000000in}{0.020833in}}%
\pgfusepath{stroke,fill}%
}%
\begin{pgfscope}%
\pgfsys@transformshift{1.356497in}{0.420000in}%
\pgfsys@useobject{currentmarker}{}%
\end{pgfscope}%
\end{pgfscope}%
\begin{pgfscope}%
\pgfsetbuttcap%
\pgfsetroundjoin%
\definecolor{currentfill}{rgb}{0.000000,0.000000,0.000000}%
\pgfsetfillcolor{currentfill}%
\pgfsetlinewidth{0.401500pt}%
\definecolor{currentstroke}{rgb}{0.000000,0.000000,0.000000}%
\pgfsetstrokecolor{currentstroke}%
\pgfsetdash{}{0pt}%
\pgfsys@defobject{currentmarker}{\pgfqpoint{0.000000in}{-0.020833in}}{\pgfqpoint{0.000000in}{0.000000in}}{%
\pgfpathmoveto{\pgfqpoint{0.000000in}{0.000000in}}%
\pgfpathlineto{\pgfqpoint{0.000000in}{-0.020833in}}%
\pgfusepath{stroke,fill}%
}%
\begin{pgfscope}%
\pgfsys@transformshift{1.356497in}{2.414488in}%
\pgfsys@useobject{currentmarker}{}%
\end{pgfscope}%
\end{pgfscope}%
\begin{pgfscope}%
\pgfsetbuttcap%
\pgfsetroundjoin%
\definecolor{currentfill}{rgb}{0.000000,0.000000,0.000000}%
\pgfsetfillcolor{currentfill}%
\pgfsetlinewidth{0.401500pt}%
\definecolor{currentstroke}{rgb}{0.000000,0.000000,0.000000}%
\pgfsetstrokecolor{currentstroke}%
\pgfsetdash{}{0pt}%
\pgfsys@defobject{currentmarker}{\pgfqpoint{0.000000in}{0.000000in}}{\pgfqpoint{0.000000in}{0.020833in}}{%
\pgfpathmoveto{\pgfqpoint{0.000000in}{0.000000in}}%
\pgfpathlineto{\pgfqpoint{0.000000in}{0.020833in}}%
\pgfusepath{stroke,fill}%
}%
\begin{pgfscope}%
\pgfsys@transformshift{1.412439in}{0.420000in}%
\pgfsys@useobject{currentmarker}{}%
\end{pgfscope}%
\end{pgfscope}%
\begin{pgfscope}%
\pgfsetbuttcap%
\pgfsetroundjoin%
\definecolor{currentfill}{rgb}{0.000000,0.000000,0.000000}%
\pgfsetfillcolor{currentfill}%
\pgfsetlinewidth{0.401500pt}%
\definecolor{currentstroke}{rgb}{0.000000,0.000000,0.000000}%
\pgfsetstrokecolor{currentstroke}%
\pgfsetdash{}{0pt}%
\pgfsys@defobject{currentmarker}{\pgfqpoint{0.000000in}{-0.020833in}}{\pgfqpoint{0.000000in}{0.000000in}}{%
\pgfpathmoveto{\pgfqpoint{0.000000in}{0.000000in}}%
\pgfpathlineto{\pgfqpoint{0.000000in}{-0.020833in}}%
\pgfusepath{stroke,fill}%
}%
\begin{pgfscope}%
\pgfsys@transformshift{1.412439in}{2.414488in}%
\pgfsys@useobject{currentmarker}{}%
\end{pgfscope}%
\end{pgfscope}%
\begin{pgfscope}%
\pgfsetbuttcap%
\pgfsetroundjoin%
\definecolor{currentfill}{rgb}{0.000000,0.000000,0.000000}%
\pgfsetfillcolor{currentfill}%
\pgfsetlinewidth{0.401500pt}%
\definecolor{currentstroke}{rgb}{0.000000,0.000000,0.000000}%
\pgfsetstrokecolor{currentstroke}%
\pgfsetdash{}{0pt}%
\pgfsys@defobject{currentmarker}{\pgfqpoint{0.000000in}{0.000000in}}{\pgfqpoint{0.000000in}{0.020833in}}{%
\pgfpathmoveto{\pgfqpoint{0.000000in}{0.000000in}}%
\pgfpathlineto{\pgfqpoint{0.000000in}{0.020833in}}%
\pgfusepath{stroke,fill}%
}%
\begin{pgfscope}%
\pgfsys@transformshift{1.459736in}{0.420000in}%
\pgfsys@useobject{currentmarker}{}%
\end{pgfscope}%
\end{pgfscope}%
\begin{pgfscope}%
\pgfsetbuttcap%
\pgfsetroundjoin%
\definecolor{currentfill}{rgb}{0.000000,0.000000,0.000000}%
\pgfsetfillcolor{currentfill}%
\pgfsetlinewidth{0.401500pt}%
\definecolor{currentstroke}{rgb}{0.000000,0.000000,0.000000}%
\pgfsetstrokecolor{currentstroke}%
\pgfsetdash{}{0pt}%
\pgfsys@defobject{currentmarker}{\pgfqpoint{0.000000in}{-0.020833in}}{\pgfqpoint{0.000000in}{0.000000in}}{%
\pgfpathmoveto{\pgfqpoint{0.000000in}{0.000000in}}%
\pgfpathlineto{\pgfqpoint{0.000000in}{-0.020833in}}%
\pgfusepath{stroke,fill}%
}%
\begin{pgfscope}%
\pgfsys@transformshift{1.459736in}{2.414488in}%
\pgfsys@useobject{currentmarker}{}%
\end{pgfscope}%
\end{pgfscope}%
\begin{pgfscope}%
\pgfsetbuttcap%
\pgfsetroundjoin%
\definecolor{currentfill}{rgb}{0.000000,0.000000,0.000000}%
\pgfsetfillcolor{currentfill}%
\pgfsetlinewidth{0.401500pt}%
\definecolor{currentstroke}{rgb}{0.000000,0.000000,0.000000}%
\pgfsetstrokecolor{currentstroke}%
\pgfsetdash{}{0pt}%
\pgfsys@defobject{currentmarker}{\pgfqpoint{0.000000in}{0.000000in}}{\pgfqpoint{0.000000in}{0.020833in}}{%
\pgfpathmoveto{\pgfqpoint{0.000000in}{0.000000in}}%
\pgfpathlineto{\pgfqpoint{0.000000in}{0.020833in}}%
\pgfusepath{stroke,fill}%
}%
\begin{pgfscope}%
\pgfsys@transformshift{1.500708in}{0.420000in}%
\pgfsys@useobject{currentmarker}{}%
\end{pgfscope}%
\end{pgfscope}%
\begin{pgfscope}%
\pgfsetbuttcap%
\pgfsetroundjoin%
\definecolor{currentfill}{rgb}{0.000000,0.000000,0.000000}%
\pgfsetfillcolor{currentfill}%
\pgfsetlinewidth{0.401500pt}%
\definecolor{currentstroke}{rgb}{0.000000,0.000000,0.000000}%
\pgfsetstrokecolor{currentstroke}%
\pgfsetdash{}{0pt}%
\pgfsys@defobject{currentmarker}{\pgfqpoint{0.000000in}{-0.020833in}}{\pgfqpoint{0.000000in}{0.000000in}}{%
\pgfpathmoveto{\pgfqpoint{0.000000in}{0.000000in}}%
\pgfpathlineto{\pgfqpoint{0.000000in}{-0.020833in}}%
\pgfusepath{stroke,fill}%
}%
\begin{pgfscope}%
\pgfsys@transformshift{1.500708in}{2.414488in}%
\pgfsys@useobject{currentmarker}{}%
\end{pgfscope}%
\end{pgfscope}%
\begin{pgfscope}%
\pgfsetbuttcap%
\pgfsetroundjoin%
\definecolor{currentfill}{rgb}{0.000000,0.000000,0.000000}%
\pgfsetfillcolor{currentfill}%
\pgfsetlinewidth{0.401500pt}%
\definecolor{currentstroke}{rgb}{0.000000,0.000000,0.000000}%
\pgfsetstrokecolor{currentstroke}%
\pgfsetdash{}{0pt}%
\pgfsys@defobject{currentmarker}{\pgfqpoint{0.000000in}{0.000000in}}{\pgfqpoint{0.000000in}{0.020833in}}{%
\pgfpathmoveto{\pgfqpoint{0.000000in}{0.000000in}}%
\pgfpathlineto{\pgfqpoint{0.000000in}{0.020833in}}%
\pgfusepath{stroke,fill}%
}%
\begin{pgfscope}%
\pgfsys@transformshift{1.536847in}{0.420000in}%
\pgfsys@useobject{currentmarker}{}%
\end{pgfscope}%
\end{pgfscope}%
\begin{pgfscope}%
\pgfsetbuttcap%
\pgfsetroundjoin%
\definecolor{currentfill}{rgb}{0.000000,0.000000,0.000000}%
\pgfsetfillcolor{currentfill}%
\pgfsetlinewidth{0.401500pt}%
\definecolor{currentstroke}{rgb}{0.000000,0.000000,0.000000}%
\pgfsetstrokecolor{currentstroke}%
\pgfsetdash{}{0pt}%
\pgfsys@defobject{currentmarker}{\pgfqpoint{0.000000in}{-0.020833in}}{\pgfqpoint{0.000000in}{0.000000in}}{%
\pgfpathmoveto{\pgfqpoint{0.000000in}{0.000000in}}%
\pgfpathlineto{\pgfqpoint{0.000000in}{-0.020833in}}%
\pgfusepath{stroke,fill}%
}%
\begin{pgfscope}%
\pgfsys@transformshift{1.536847in}{2.414488in}%
\pgfsys@useobject{currentmarker}{}%
\end{pgfscope}%
\end{pgfscope}%
\begin{pgfscope}%
\pgfsetbuttcap%
\pgfsetroundjoin%
\definecolor{currentfill}{rgb}{0.000000,0.000000,0.000000}%
\pgfsetfillcolor{currentfill}%
\pgfsetlinewidth{0.401500pt}%
\definecolor{currentstroke}{rgb}{0.000000,0.000000,0.000000}%
\pgfsetstrokecolor{currentstroke}%
\pgfsetdash{}{0pt}%
\pgfsys@defobject{currentmarker}{\pgfqpoint{0.000000in}{0.000000in}}{\pgfqpoint{0.000000in}{0.020833in}}{%
\pgfpathmoveto{\pgfqpoint{0.000000in}{0.000000in}}%
\pgfpathlineto{\pgfqpoint{0.000000in}{0.020833in}}%
\pgfusepath{stroke,fill}%
}%
\begin{pgfscope}%
\pgfsys@transformshift{1.781851in}{0.420000in}%
\pgfsys@useobject{currentmarker}{}%
\end{pgfscope}%
\end{pgfscope}%
\begin{pgfscope}%
\pgfsetbuttcap%
\pgfsetroundjoin%
\definecolor{currentfill}{rgb}{0.000000,0.000000,0.000000}%
\pgfsetfillcolor{currentfill}%
\pgfsetlinewidth{0.401500pt}%
\definecolor{currentstroke}{rgb}{0.000000,0.000000,0.000000}%
\pgfsetstrokecolor{currentstroke}%
\pgfsetdash{}{0pt}%
\pgfsys@defobject{currentmarker}{\pgfqpoint{0.000000in}{-0.020833in}}{\pgfqpoint{0.000000in}{0.000000in}}{%
\pgfpathmoveto{\pgfqpoint{0.000000in}{0.000000in}}%
\pgfpathlineto{\pgfqpoint{0.000000in}{-0.020833in}}%
\pgfusepath{stroke,fill}%
}%
\begin{pgfscope}%
\pgfsys@transformshift{1.781851in}{2.414488in}%
\pgfsys@useobject{currentmarker}{}%
\end{pgfscope}%
\end{pgfscope}%
\begin{pgfscope}%
\pgfsetbuttcap%
\pgfsetroundjoin%
\definecolor{currentfill}{rgb}{0.000000,0.000000,0.000000}%
\pgfsetfillcolor{currentfill}%
\pgfsetlinewidth{0.401500pt}%
\definecolor{currentstroke}{rgb}{0.000000,0.000000,0.000000}%
\pgfsetstrokecolor{currentstroke}%
\pgfsetdash{}{0pt}%
\pgfsys@defobject{currentmarker}{\pgfqpoint{0.000000in}{0.000000in}}{\pgfqpoint{0.000000in}{0.020833in}}{%
\pgfpathmoveto{\pgfqpoint{0.000000in}{0.000000in}}%
\pgfpathlineto{\pgfqpoint{0.000000in}{0.020833in}}%
\pgfusepath{stroke,fill}%
}%
\begin{pgfscope}%
\pgfsys@transformshift{1.906259in}{0.420000in}%
\pgfsys@useobject{currentmarker}{}%
\end{pgfscope}%
\end{pgfscope}%
\begin{pgfscope}%
\pgfsetbuttcap%
\pgfsetroundjoin%
\definecolor{currentfill}{rgb}{0.000000,0.000000,0.000000}%
\pgfsetfillcolor{currentfill}%
\pgfsetlinewidth{0.401500pt}%
\definecolor{currentstroke}{rgb}{0.000000,0.000000,0.000000}%
\pgfsetstrokecolor{currentstroke}%
\pgfsetdash{}{0pt}%
\pgfsys@defobject{currentmarker}{\pgfqpoint{0.000000in}{-0.020833in}}{\pgfqpoint{0.000000in}{0.000000in}}{%
\pgfpathmoveto{\pgfqpoint{0.000000in}{0.000000in}}%
\pgfpathlineto{\pgfqpoint{0.000000in}{-0.020833in}}%
\pgfusepath{stroke,fill}%
}%
\begin{pgfscope}%
\pgfsys@transformshift{1.906259in}{2.414488in}%
\pgfsys@useobject{currentmarker}{}%
\end{pgfscope}%
\end{pgfscope}%
\begin{pgfscope}%
\pgfsetbuttcap%
\pgfsetroundjoin%
\definecolor{currentfill}{rgb}{0.000000,0.000000,0.000000}%
\pgfsetfillcolor{currentfill}%
\pgfsetlinewidth{0.401500pt}%
\definecolor{currentstroke}{rgb}{0.000000,0.000000,0.000000}%
\pgfsetstrokecolor{currentstroke}%
\pgfsetdash{}{0pt}%
\pgfsys@defobject{currentmarker}{\pgfqpoint{0.000000in}{0.000000in}}{\pgfqpoint{0.000000in}{0.020833in}}{%
\pgfpathmoveto{\pgfqpoint{0.000000in}{0.000000in}}%
\pgfpathlineto{\pgfqpoint{0.000000in}{0.020833in}}%
\pgfusepath{stroke,fill}%
}%
\begin{pgfscope}%
\pgfsys@transformshift{1.994528in}{0.420000in}%
\pgfsys@useobject{currentmarker}{}%
\end{pgfscope}%
\end{pgfscope}%
\begin{pgfscope}%
\pgfsetbuttcap%
\pgfsetroundjoin%
\definecolor{currentfill}{rgb}{0.000000,0.000000,0.000000}%
\pgfsetfillcolor{currentfill}%
\pgfsetlinewidth{0.401500pt}%
\definecolor{currentstroke}{rgb}{0.000000,0.000000,0.000000}%
\pgfsetstrokecolor{currentstroke}%
\pgfsetdash{}{0pt}%
\pgfsys@defobject{currentmarker}{\pgfqpoint{0.000000in}{-0.020833in}}{\pgfqpoint{0.000000in}{0.000000in}}{%
\pgfpathmoveto{\pgfqpoint{0.000000in}{0.000000in}}%
\pgfpathlineto{\pgfqpoint{0.000000in}{-0.020833in}}%
\pgfusepath{stroke,fill}%
}%
\begin{pgfscope}%
\pgfsys@transformshift{1.994528in}{2.414488in}%
\pgfsys@useobject{currentmarker}{}%
\end{pgfscope}%
\end{pgfscope}%
\begin{pgfscope}%
\pgfsetbuttcap%
\pgfsetroundjoin%
\definecolor{currentfill}{rgb}{0.000000,0.000000,0.000000}%
\pgfsetfillcolor{currentfill}%
\pgfsetlinewidth{0.401500pt}%
\definecolor{currentstroke}{rgb}{0.000000,0.000000,0.000000}%
\pgfsetstrokecolor{currentstroke}%
\pgfsetdash{}{0pt}%
\pgfsys@defobject{currentmarker}{\pgfqpoint{0.000000in}{0.000000in}}{\pgfqpoint{0.000000in}{0.020833in}}{%
\pgfpathmoveto{\pgfqpoint{0.000000in}{0.000000in}}%
\pgfpathlineto{\pgfqpoint{0.000000in}{0.020833in}}%
\pgfusepath{stroke,fill}%
}%
\begin{pgfscope}%
\pgfsys@transformshift{2.062995in}{0.420000in}%
\pgfsys@useobject{currentmarker}{}%
\end{pgfscope}%
\end{pgfscope}%
\begin{pgfscope}%
\pgfsetbuttcap%
\pgfsetroundjoin%
\definecolor{currentfill}{rgb}{0.000000,0.000000,0.000000}%
\pgfsetfillcolor{currentfill}%
\pgfsetlinewidth{0.401500pt}%
\definecolor{currentstroke}{rgb}{0.000000,0.000000,0.000000}%
\pgfsetstrokecolor{currentstroke}%
\pgfsetdash{}{0pt}%
\pgfsys@defobject{currentmarker}{\pgfqpoint{0.000000in}{-0.020833in}}{\pgfqpoint{0.000000in}{0.000000in}}{%
\pgfpathmoveto{\pgfqpoint{0.000000in}{0.000000in}}%
\pgfpathlineto{\pgfqpoint{0.000000in}{-0.020833in}}%
\pgfusepath{stroke,fill}%
}%
\begin{pgfscope}%
\pgfsys@transformshift{2.062995in}{2.414488in}%
\pgfsys@useobject{currentmarker}{}%
\end{pgfscope}%
\end{pgfscope}%
\begin{pgfscope}%
\pgfsetbuttcap%
\pgfsetroundjoin%
\definecolor{currentfill}{rgb}{0.000000,0.000000,0.000000}%
\pgfsetfillcolor{currentfill}%
\pgfsetlinewidth{0.401500pt}%
\definecolor{currentstroke}{rgb}{0.000000,0.000000,0.000000}%
\pgfsetstrokecolor{currentstroke}%
\pgfsetdash{}{0pt}%
\pgfsys@defobject{currentmarker}{\pgfqpoint{0.000000in}{0.000000in}}{\pgfqpoint{0.000000in}{0.020833in}}{%
\pgfpathmoveto{\pgfqpoint{0.000000in}{0.000000in}}%
\pgfpathlineto{\pgfqpoint{0.000000in}{0.020833in}}%
\pgfusepath{stroke,fill}%
}%
\begin{pgfscope}%
\pgfsys@transformshift{2.118936in}{0.420000in}%
\pgfsys@useobject{currentmarker}{}%
\end{pgfscope}%
\end{pgfscope}%
\begin{pgfscope}%
\pgfsetbuttcap%
\pgfsetroundjoin%
\definecolor{currentfill}{rgb}{0.000000,0.000000,0.000000}%
\pgfsetfillcolor{currentfill}%
\pgfsetlinewidth{0.401500pt}%
\definecolor{currentstroke}{rgb}{0.000000,0.000000,0.000000}%
\pgfsetstrokecolor{currentstroke}%
\pgfsetdash{}{0pt}%
\pgfsys@defobject{currentmarker}{\pgfqpoint{0.000000in}{-0.020833in}}{\pgfqpoint{0.000000in}{0.000000in}}{%
\pgfpathmoveto{\pgfqpoint{0.000000in}{0.000000in}}%
\pgfpathlineto{\pgfqpoint{0.000000in}{-0.020833in}}%
\pgfusepath{stroke,fill}%
}%
\begin{pgfscope}%
\pgfsys@transformshift{2.118936in}{2.414488in}%
\pgfsys@useobject{currentmarker}{}%
\end{pgfscope}%
\end{pgfscope}%
\begin{pgfscope}%
\pgfsetbuttcap%
\pgfsetroundjoin%
\definecolor{currentfill}{rgb}{0.000000,0.000000,0.000000}%
\pgfsetfillcolor{currentfill}%
\pgfsetlinewidth{0.401500pt}%
\definecolor{currentstroke}{rgb}{0.000000,0.000000,0.000000}%
\pgfsetstrokecolor{currentstroke}%
\pgfsetdash{}{0pt}%
\pgfsys@defobject{currentmarker}{\pgfqpoint{0.000000in}{0.000000in}}{\pgfqpoint{0.000000in}{0.020833in}}{%
\pgfpathmoveto{\pgfqpoint{0.000000in}{0.000000in}}%
\pgfpathlineto{\pgfqpoint{0.000000in}{0.020833in}}%
\pgfusepath{stroke,fill}%
}%
\begin{pgfscope}%
\pgfsys@transformshift{2.166234in}{0.420000in}%
\pgfsys@useobject{currentmarker}{}%
\end{pgfscope}%
\end{pgfscope}%
\begin{pgfscope}%
\pgfsetbuttcap%
\pgfsetroundjoin%
\definecolor{currentfill}{rgb}{0.000000,0.000000,0.000000}%
\pgfsetfillcolor{currentfill}%
\pgfsetlinewidth{0.401500pt}%
\definecolor{currentstroke}{rgb}{0.000000,0.000000,0.000000}%
\pgfsetstrokecolor{currentstroke}%
\pgfsetdash{}{0pt}%
\pgfsys@defobject{currentmarker}{\pgfqpoint{0.000000in}{-0.020833in}}{\pgfqpoint{0.000000in}{0.000000in}}{%
\pgfpathmoveto{\pgfqpoint{0.000000in}{0.000000in}}%
\pgfpathlineto{\pgfqpoint{0.000000in}{-0.020833in}}%
\pgfusepath{stroke,fill}%
}%
\begin{pgfscope}%
\pgfsys@transformshift{2.166234in}{2.414488in}%
\pgfsys@useobject{currentmarker}{}%
\end{pgfscope}%
\end{pgfscope}%
\begin{pgfscope}%
\pgfsetbuttcap%
\pgfsetroundjoin%
\definecolor{currentfill}{rgb}{0.000000,0.000000,0.000000}%
\pgfsetfillcolor{currentfill}%
\pgfsetlinewidth{0.401500pt}%
\definecolor{currentstroke}{rgb}{0.000000,0.000000,0.000000}%
\pgfsetstrokecolor{currentstroke}%
\pgfsetdash{}{0pt}%
\pgfsys@defobject{currentmarker}{\pgfqpoint{0.000000in}{0.000000in}}{\pgfqpoint{0.000000in}{0.020833in}}{%
\pgfpathmoveto{\pgfqpoint{0.000000in}{0.000000in}}%
\pgfpathlineto{\pgfqpoint{0.000000in}{0.020833in}}%
\pgfusepath{stroke,fill}%
}%
\begin{pgfscope}%
\pgfsys@transformshift{2.207205in}{0.420000in}%
\pgfsys@useobject{currentmarker}{}%
\end{pgfscope}%
\end{pgfscope}%
\begin{pgfscope}%
\pgfsetbuttcap%
\pgfsetroundjoin%
\definecolor{currentfill}{rgb}{0.000000,0.000000,0.000000}%
\pgfsetfillcolor{currentfill}%
\pgfsetlinewidth{0.401500pt}%
\definecolor{currentstroke}{rgb}{0.000000,0.000000,0.000000}%
\pgfsetstrokecolor{currentstroke}%
\pgfsetdash{}{0pt}%
\pgfsys@defobject{currentmarker}{\pgfqpoint{0.000000in}{-0.020833in}}{\pgfqpoint{0.000000in}{0.000000in}}{%
\pgfpathmoveto{\pgfqpoint{0.000000in}{0.000000in}}%
\pgfpathlineto{\pgfqpoint{0.000000in}{-0.020833in}}%
\pgfusepath{stroke,fill}%
}%
\begin{pgfscope}%
\pgfsys@transformshift{2.207205in}{2.414488in}%
\pgfsys@useobject{currentmarker}{}%
\end{pgfscope}%
\end{pgfscope}%
\begin{pgfscope}%
\pgfsetbuttcap%
\pgfsetroundjoin%
\definecolor{currentfill}{rgb}{0.000000,0.000000,0.000000}%
\pgfsetfillcolor{currentfill}%
\pgfsetlinewidth{0.401500pt}%
\definecolor{currentstroke}{rgb}{0.000000,0.000000,0.000000}%
\pgfsetstrokecolor{currentstroke}%
\pgfsetdash{}{0pt}%
\pgfsys@defobject{currentmarker}{\pgfqpoint{0.000000in}{0.000000in}}{\pgfqpoint{0.000000in}{0.020833in}}{%
\pgfpathmoveto{\pgfqpoint{0.000000in}{0.000000in}}%
\pgfpathlineto{\pgfqpoint{0.000000in}{0.020833in}}%
\pgfusepath{stroke,fill}%
}%
\begin{pgfscope}%
\pgfsys@transformshift{2.243344in}{0.420000in}%
\pgfsys@useobject{currentmarker}{}%
\end{pgfscope}%
\end{pgfscope}%
\begin{pgfscope}%
\pgfsetbuttcap%
\pgfsetroundjoin%
\definecolor{currentfill}{rgb}{0.000000,0.000000,0.000000}%
\pgfsetfillcolor{currentfill}%
\pgfsetlinewidth{0.401500pt}%
\definecolor{currentstroke}{rgb}{0.000000,0.000000,0.000000}%
\pgfsetstrokecolor{currentstroke}%
\pgfsetdash{}{0pt}%
\pgfsys@defobject{currentmarker}{\pgfqpoint{0.000000in}{-0.020833in}}{\pgfqpoint{0.000000in}{0.000000in}}{%
\pgfpathmoveto{\pgfqpoint{0.000000in}{0.000000in}}%
\pgfpathlineto{\pgfqpoint{0.000000in}{-0.020833in}}%
\pgfusepath{stroke,fill}%
}%
\begin{pgfscope}%
\pgfsys@transformshift{2.243344in}{2.414488in}%
\pgfsys@useobject{currentmarker}{}%
\end{pgfscope}%
\end{pgfscope}%
\begin{pgfscope}%
\pgfsetbuttcap%
\pgfsetroundjoin%
\definecolor{currentfill}{rgb}{0.000000,0.000000,0.000000}%
\pgfsetfillcolor{currentfill}%
\pgfsetlinewidth{0.401500pt}%
\definecolor{currentstroke}{rgb}{0.000000,0.000000,0.000000}%
\pgfsetstrokecolor{currentstroke}%
\pgfsetdash{}{0pt}%
\pgfsys@defobject{currentmarker}{\pgfqpoint{0.000000in}{0.000000in}}{\pgfqpoint{0.000000in}{0.020833in}}{%
\pgfpathmoveto{\pgfqpoint{0.000000in}{0.000000in}}%
\pgfpathlineto{\pgfqpoint{0.000000in}{0.020833in}}%
\pgfusepath{stroke,fill}%
}%
\begin{pgfscope}%
\pgfsys@transformshift{2.488348in}{0.420000in}%
\pgfsys@useobject{currentmarker}{}%
\end{pgfscope}%
\end{pgfscope}%
\begin{pgfscope}%
\pgfsetbuttcap%
\pgfsetroundjoin%
\definecolor{currentfill}{rgb}{0.000000,0.000000,0.000000}%
\pgfsetfillcolor{currentfill}%
\pgfsetlinewidth{0.401500pt}%
\definecolor{currentstroke}{rgb}{0.000000,0.000000,0.000000}%
\pgfsetstrokecolor{currentstroke}%
\pgfsetdash{}{0pt}%
\pgfsys@defobject{currentmarker}{\pgfqpoint{0.000000in}{-0.020833in}}{\pgfqpoint{0.000000in}{0.000000in}}{%
\pgfpathmoveto{\pgfqpoint{0.000000in}{0.000000in}}%
\pgfpathlineto{\pgfqpoint{0.000000in}{-0.020833in}}%
\pgfusepath{stroke,fill}%
}%
\begin{pgfscope}%
\pgfsys@transformshift{2.488348in}{2.414488in}%
\pgfsys@useobject{currentmarker}{}%
\end{pgfscope}%
\end{pgfscope}%
\begin{pgfscope}%
\pgfsetbuttcap%
\pgfsetroundjoin%
\definecolor{currentfill}{rgb}{0.000000,0.000000,0.000000}%
\pgfsetfillcolor{currentfill}%
\pgfsetlinewidth{0.401500pt}%
\definecolor{currentstroke}{rgb}{0.000000,0.000000,0.000000}%
\pgfsetstrokecolor{currentstroke}%
\pgfsetdash{}{0pt}%
\pgfsys@defobject{currentmarker}{\pgfqpoint{0.000000in}{0.000000in}}{\pgfqpoint{0.000000in}{0.020833in}}{%
\pgfpathmoveto{\pgfqpoint{0.000000in}{0.000000in}}%
\pgfpathlineto{\pgfqpoint{0.000000in}{0.020833in}}%
\pgfusepath{stroke,fill}%
}%
\begin{pgfscope}%
\pgfsys@transformshift{2.612756in}{0.420000in}%
\pgfsys@useobject{currentmarker}{}%
\end{pgfscope}%
\end{pgfscope}%
\begin{pgfscope}%
\pgfsetbuttcap%
\pgfsetroundjoin%
\definecolor{currentfill}{rgb}{0.000000,0.000000,0.000000}%
\pgfsetfillcolor{currentfill}%
\pgfsetlinewidth{0.401500pt}%
\definecolor{currentstroke}{rgb}{0.000000,0.000000,0.000000}%
\pgfsetstrokecolor{currentstroke}%
\pgfsetdash{}{0pt}%
\pgfsys@defobject{currentmarker}{\pgfqpoint{0.000000in}{-0.020833in}}{\pgfqpoint{0.000000in}{0.000000in}}{%
\pgfpathmoveto{\pgfqpoint{0.000000in}{0.000000in}}%
\pgfpathlineto{\pgfqpoint{0.000000in}{-0.020833in}}%
\pgfusepath{stroke,fill}%
}%
\begin{pgfscope}%
\pgfsys@transformshift{2.612756in}{2.414488in}%
\pgfsys@useobject{currentmarker}{}%
\end{pgfscope}%
\end{pgfscope}%
\begin{pgfscope}%
\pgfsetbuttcap%
\pgfsetroundjoin%
\definecolor{currentfill}{rgb}{0.000000,0.000000,0.000000}%
\pgfsetfillcolor{currentfill}%
\pgfsetlinewidth{0.401500pt}%
\definecolor{currentstroke}{rgb}{0.000000,0.000000,0.000000}%
\pgfsetstrokecolor{currentstroke}%
\pgfsetdash{}{0pt}%
\pgfsys@defobject{currentmarker}{\pgfqpoint{0.000000in}{0.000000in}}{\pgfqpoint{0.000000in}{0.020833in}}{%
\pgfpathmoveto{\pgfqpoint{0.000000in}{0.000000in}}%
\pgfpathlineto{\pgfqpoint{0.000000in}{0.020833in}}%
\pgfusepath{stroke,fill}%
}%
\begin{pgfscope}%
\pgfsys@transformshift{2.701025in}{0.420000in}%
\pgfsys@useobject{currentmarker}{}%
\end{pgfscope}%
\end{pgfscope}%
\begin{pgfscope}%
\pgfsetbuttcap%
\pgfsetroundjoin%
\definecolor{currentfill}{rgb}{0.000000,0.000000,0.000000}%
\pgfsetfillcolor{currentfill}%
\pgfsetlinewidth{0.401500pt}%
\definecolor{currentstroke}{rgb}{0.000000,0.000000,0.000000}%
\pgfsetstrokecolor{currentstroke}%
\pgfsetdash{}{0pt}%
\pgfsys@defobject{currentmarker}{\pgfqpoint{0.000000in}{-0.020833in}}{\pgfqpoint{0.000000in}{0.000000in}}{%
\pgfpathmoveto{\pgfqpoint{0.000000in}{0.000000in}}%
\pgfpathlineto{\pgfqpoint{0.000000in}{-0.020833in}}%
\pgfusepath{stroke,fill}%
}%
\begin{pgfscope}%
\pgfsys@transformshift{2.701025in}{2.414488in}%
\pgfsys@useobject{currentmarker}{}%
\end{pgfscope}%
\end{pgfscope}%
\begin{pgfscope}%
\pgfsetbuttcap%
\pgfsetroundjoin%
\definecolor{currentfill}{rgb}{0.000000,0.000000,0.000000}%
\pgfsetfillcolor{currentfill}%
\pgfsetlinewidth{0.401500pt}%
\definecolor{currentstroke}{rgb}{0.000000,0.000000,0.000000}%
\pgfsetstrokecolor{currentstroke}%
\pgfsetdash{}{0pt}%
\pgfsys@defobject{currentmarker}{\pgfqpoint{0.000000in}{0.000000in}}{\pgfqpoint{0.000000in}{0.020833in}}{%
\pgfpathmoveto{\pgfqpoint{0.000000in}{0.000000in}}%
\pgfpathlineto{\pgfqpoint{0.000000in}{0.020833in}}%
\pgfusepath{stroke,fill}%
}%
\begin{pgfscope}%
\pgfsys@transformshift{2.769492in}{0.420000in}%
\pgfsys@useobject{currentmarker}{}%
\end{pgfscope}%
\end{pgfscope}%
\begin{pgfscope}%
\pgfsetbuttcap%
\pgfsetroundjoin%
\definecolor{currentfill}{rgb}{0.000000,0.000000,0.000000}%
\pgfsetfillcolor{currentfill}%
\pgfsetlinewidth{0.401500pt}%
\definecolor{currentstroke}{rgb}{0.000000,0.000000,0.000000}%
\pgfsetstrokecolor{currentstroke}%
\pgfsetdash{}{0pt}%
\pgfsys@defobject{currentmarker}{\pgfqpoint{0.000000in}{-0.020833in}}{\pgfqpoint{0.000000in}{0.000000in}}{%
\pgfpathmoveto{\pgfqpoint{0.000000in}{0.000000in}}%
\pgfpathlineto{\pgfqpoint{0.000000in}{-0.020833in}}%
\pgfusepath{stroke,fill}%
}%
\begin{pgfscope}%
\pgfsys@transformshift{2.769492in}{2.414488in}%
\pgfsys@useobject{currentmarker}{}%
\end{pgfscope}%
\end{pgfscope}%
\begin{pgfscope}%
\pgfsetbuttcap%
\pgfsetroundjoin%
\definecolor{currentfill}{rgb}{0.000000,0.000000,0.000000}%
\pgfsetfillcolor{currentfill}%
\pgfsetlinewidth{0.401500pt}%
\definecolor{currentstroke}{rgb}{0.000000,0.000000,0.000000}%
\pgfsetstrokecolor{currentstroke}%
\pgfsetdash{}{0pt}%
\pgfsys@defobject{currentmarker}{\pgfqpoint{0.000000in}{0.000000in}}{\pgfqpoint{0.000000in}{0.020833in}}{%
\pgfpathmoveto{\pgfqpoint{0.000000in}{0.000000in}}%
\pgfpathlineto{\pgfqpoint{0.000000in}{0.020833in}}%
\pgfusepath{stroke,fill}%
}%
\begin{pgfscope}%
\pgfsys@transformshift{2.825433in}{0.420000in}%
\pgfsys@useobject{currentmarker}{}%
\end{pgfscope}%
\end{pgfscope}%
\begin{pgfscope}%
\pgfsetbuttcap%
\pgfsetroundjoin%
\definecolor{currentfill}{rgb}{0.000000,0.000000,0.000000}%
\pgfsetfillcolor{currentfill}%
\pgfsetlinewidth{0.401500pt}%
\definecolor{currentstroke}{rgb}{0.000000,0.000000,0.000000}%
\pgfsetstrokecolor{currentstroke}%
\pgfsetdash{}{0pt}%
\pgfsys@defobject{currentmarker}{\pgfqpoint{0.000000in}{-0.020833in}}{\pgfqpoint{0.000000in}{0.000000in}}{%
\pgfpathmoveto{\pgfqpoint{0.000000in}{0.000000in}}%
\pgfpathlineto{\pgfqpoint{0.000000in}{-0.020833in}}%
\pgfusepath{stroke,fill}%
}%
\begin{pgfscope}%
\pgfsys@transformshift{2.825433in}{2.414488in}%
\pgfsys@useobject{currentmarker}{}%
\end{pgfscope}%
\end{pgfscope}%
\begin{pgfscope}%
\pgfsetbuttcap%
\pgfsetroundjoin%
\definecolor{currentfill}{rgb}{0.000000,0.000000,0.000000}%
\pgfsetfillcolor{currentfill}%
\pgfsetlinewidth{0.401500pt}%
\definecolor{currentstroke}{rgb}{0.000000,0.000000,0.000000}%
\pgfsetstrokecolor{currentstroke}%
\pgfsetdash{}{0pt}%
\pgfsys@defobject{currentmarker}{\pgfqpoint{0.000000in}{0.000000in}}{\pgfqpoint{0.000000in}{0.020833in}}{%
\pgfpathmoveto{\pgfqpoint{0.000000in}{0.000000in}}%
\pgfpathlineto{\pgfqpoint{0.000000in}{0.020833in}}%
\pgfusepath{stroke,fill}%
}%
\begin{pgfscope}%
\pgfsys@transformshift{2.872731in}{0.420000in}%
\pgfsys@useobject{currentmarker}{}%
\end{pgfscope}%
\end{pgfscope}%
\begin{pgfscope}%
\pgfsetbuttcap%
\pgfsetroundjoin%
\definecolor{currentfill}{rgb}{0.000000,0.000000,0.000000}%
\pgfsetfillcolor{currentfill}%
\pgfsetlinewidth{0.401500pt}%
\definecolor{currentstroke}{rgb}{0.000000,0.000000,0.000000}%
\pgfsetstrokecolor{currentstroke}%
\pgfsetdash{}{0pt}%
\pgfsys@defobject{currentmarker}{\pgfqpoint{0.000000in}{-0.020833in}}{\pgfqpoint{0.000000in}{0.000000in}}{%
\pgfpathmoveto{\pgfqpoint{0.000000in}{0.000000in}}%
\pgfpathlineto{\pgfqpoint{0.000000in}{-0.020833in}}%
\pgfusepath{stroke,fill}%
}%
\begin{pgfscope}%
\pgfsys@transformshift{2.872731in}{2.414488in}%
\pgfsys@useobject{currentmarker}{}%
\end{pgfscope}%
\end{pgfscope}%
\begin{pgfscope}%
\pgfsetbuttcap%
\pgfsetroundjoin%
\definecolor{currentfill}{rgb}{0.000000,0.000000,0.000000}%
\pgfsetfillcolor{currentfill}%
\pgfsetlinewidth{0.401500pt}%
\definecolor{currentstroke}{rgb}{0.000000,0.000000,0.000000}%
\pgfsetstrokecolor{currentstroke}%
\pgfsetdash{}{0pt}%
\pgfsys@defobject{currentmarker}{\pgfqpoint{0.000000in}{0.000000in}}{\pgfqpoint{0.000000in}{0.020833in}}{%
\pgfpathmoveto{\pgfqpoint{0.000000in}{0.000000in}}%
\pgfpathlineto{\pgfqpoint{0.000000in}{0.020833in}}%
\pgfusepath{stroke,fill}%
}%
\begin{pgfscope}%
\pgfsys@transformshift{2.913702in}{0.420000in}%
\pgfsys@useobject{currentmarker}{}%
\end{pgfscope}%
\end{pgfscope}%
\begin{pgfscope}%
\pgfsetbuttcap%
\pgfsetroundjoin%
\definecolor{currentfill}{rgb}{0.000000,0.000000,0.000000}%
\pgfsetfillcolor{currentfill}%
\pgfsetlinewidth{0.401500pt}%
\definecolor{currentstroke}{rgb}{0.000000,0.000000,0.000000}%
\pgfsetstrokecolor{currentstroke}%
\pgfsetdash{}{0pt}%
\pgfsys@defobject{currentmarker}{\pgfqpoint{0.000000in}{-0.020833in}}{\pgfqpoint{0.000000in}{0.000000in}}{%
\pgfpathmoveto{\pgfqpoint{0.000000in}{0.000000in}}%
\pgfpathlineto{\pgfqpoint{0.000000in}{-0.020833in}}%
\pgfusepath{stroke,fill}%
}%
\begin{pgfscope}%
\pgfsys@transformshift{2.913702in}{2.414488in}%
\pgfsys@useobject{currentmarker}{}%
\end{pgfscope}%
\end{pgfscope}%
\begin{pgfscope}%
\pgfsetbuttcap%
\pgfsetroundjoin%
\definecolor{currentfill}{rgb}{0.000000,0.000000,0.000000}%
\pgfsetfillcolor{currentfill}%
\pgfsetlinewidth{0.401500pt}%
\definecolor{currentstroke}{rgb}{0.000000,0.000000,0.000000}%
\pgfsetstrokecolor{currentstroke}%
\pgfsetdash{}{0pt}%
\pgfsys@defobject{currentmarker}{\pgfqpoint{0.000000in}{0.000000in}}{\pgfqpoint{0.000000in}{0.020833in}}{%
\pgfpathmoveto{\pgfqpoint{0.000000in}{0.000000in}}%
\pgfpathlineto{\pgfqpoint{0.000000in}{0.020833in}}%
\pgfusepath{stroke,fill}%
}%
\begin{pgfscope}%
\pgfsys@transformshift{2.949841in}{0.420000in}%
\pgfsys@useobject{currentmarker}{}%
\end{pgfscope}%
\end{pgfscope}%
\begin{pgfscope}%
\pgfsetbuttcap%
\pgfsetroundjoin%
\definecolor{currentfill}{rgb}{0.000000,0.000000,0.000000}%
\pgfsetfillcolor{currentfill}%
\pgfsetlinewidth{0.401500pt}%
\definecolor{currentstroke}{rgb}{0.000000,0.000000,0.000000}%
\pgfsetstrokecolor{currentstroke}%
\pgfsetdash{}{0pt}%
\pgfsys@defobject{currentmarker}{\pgfqpoint{0.000000in}{-0.020833in}}{\pgfqpoint{0.000000in}{0.000000in}}{%
\pgfpathmoveto{\pgfqpoint{0.000000in}{0.000000in}}%
\pgfpathlineto{\pgfqpoint{0.000000in}{-0.020833in}}%
\pgfusepath{stroke,fill}%
}%
\begin{pgfscope}%
\pgfsys@transformshift{2.949841in}{2.414488in}%
\pgfsys@useobject{currentmarker}{}%
\end{pgfscope}%
\end{pgfscope}%
\begin{pgfscope}%
\definecolor{textcolor}{rgb}{0.000000,0.000000,0.000000}%
\pgfsetstrokecolor{textcolor}%
\pgfsetfillcolor{textcolor}%
\pgftext[x=1.844055in,y=0.208426in,,top]{\color{textcolor}{\rmfamily\fontsize{12.000000}{14.400000}\selectfont\catcode`\^=\active\def^{\ifmmode\sp\else\^{}\fi}\catcode`\%=\active\def%{\%}$y^{+}$}}%
\end{pgfscope}%
\begin{pgfscope}%
\pgfsetbuttcap%
\pgfsetroundjoin%
\definecolor{currentfill}{rgb}{0.000000,0.000000,0.000000}%
\pgfsetfillcolor{currentfill}%
\pgfsetlinewidth{0.501875pt}%
\definecolor{currentstroke}{rgb}{0.000000,0.000000,0.000000}%
\pgfsetstrokecolor{currentstroke}%
\pgfsetdash{}{0pt}%
\pgfsys@defobject{currentmarker}{\pgfqpoint{0.000000in}{0.000000in}}{\pgfqpoint{0.041667in}{0.000000in}}{%
\pgfpathmoveto{\pgfqpoint{0.000000in}{0.000000in}}%
\pgfpathlineto{\pgfqpoint{0.041667in}{0.000000in}}%
\pgfusepath{stroke,fill}%
}%
\begin{pgfscope}%
\pgfsys@transformshift{0.650000in}{0.496711in}%
\pgfsys@useobject{currentmarker}{}%
\end{pgfscope}%
\end{pgfscope}%
\begin{pgfscope}%
\pgfsetbuttcap%
\pgfsetroundjoin%
\definecolor{currentfill}{rgb}{0.000000,0.000000,0.000000}%
\pgfsetfillcolor{currentfill}%
\pgfsetlinewidth{0.501875pt}%
\definecolor{currentstroke}{rgb}{0.000000,0.000000,0.000000}%
\pgfsetstrokecolor{currentstroke}%
\pgfsetdash{}{0pt}%
\pgfsys@defobject{currentmarker}{\pgfqpoint{-0.041667in}{0.000000in}}{\pgfqpoint{-0.000000in}{0.000000in}}{%
\pgfpathmoveto{\pgfqpoint{-0.000000in}{0.000000in}}%
\pgfpathlineto{\pgfqpoint{-0.041667in}{0.000000in}}%
\pgfusepath{stroke,fill}%
}%
\begin{pgfscope}%
\pgfsys@transformshift{3.038110in}{0.496711in}%
\pgfsys@useobject{currentmarker}{}%
\end{pgfscope}%
\end{pgfscope}%
\begin{pgfscope}%
\definecolor{textcolor}{rgb}{0.000000,0.000000,0.000000}%
\pgfsetstrokecolor{textcolor}%
\pgfsetfillcolor{textcolor}%
\pgftext[x=0.288772in, y=0.443904in, left, base]{\color{textcolor}{\rmfamily\fontsize{11.000000}{13.200000}\selectfont\catcode`\^=\active\def^{\ifmmode\sp\else\^{}\fi}\catcode`\%=\active\def%{\%}\ensuremath{-}2.5}}%
\end{pgfscope}%
\begin{pgfscope}%
\pgfsetbuttcap%
\pgfsetroundjoin%
\definecolor{currentfill}{rgb}{0.000000,0.000000,0.000000}%
\pgfsetfillcolor{currentfill}%
\pgfsetlinewidth{0.501875pt}%
\definecolor{currentstroke}{rgb}{0.000000,0.000000,0.000000}%
\pgfsetstrokecolor{currentstroke}%
\pgfsetdash{}{0pt}%
\pgfsys@defobject{currentmarker}{\pgfqpoint{0.000000in}{0.000000in}}{\pgfqpoint{0.041667in}{0.000000in}}{%
\pgfpathmoveto{\pgfqpoint{0.000000in}{0.000000in}}%
\pgfpathlineto{\pgfqpoint{0.041667in}{0.000000in}}%
\pgfusepath{stroke,fill}%
}%
\begin{pgfscope}%
\pgfsys@transformshift{0.650000in}{0.880266in}%
\pgfsys@useobject{currentmarker}{}%
\end{pgfscope}%
\end{pgfscope}%
\begin{pgfscope}%
\pgfsetbuttcap%
\pgfsetroundjoin%
\definecolor{currentfill}{rgb}{0.000000,0.000000,0.000000}%
\pgfsetfillcolor{currentfill}%
\pgfsetlinewidth{0.501875pt}%
\definecolor{currentstroke}{rgb}{0.000000,0.000000,0.000000}%
\pgfsetstrokecolor{currentstroke}%
\pgfsetdash{}{0pt}%
\pgfsys@defobject{currentmarker}{\pgfqpoint{-0.041667in}{0.000000in}}{\pgfqpoint{-0.000000in}{0.000000in}}{%
\pgfpathmoveto{\pgfqpoint{-0.000000in}{0.000000in}}%
\pgfpathlineto{\pgfqpoint{-0.041667in}{0.000000in}}%
\pgfusepath{stroke,fill}%
}%
\begin{pgfscope}%
\pgfsys@transformshift{3.038110in}{0.880266in}%
\pgfsys@useobject{currentmarker}{}%
\end{pgfscope}%
\end{pgfscope}%
\begin{pgfscope}%
\definecolor{textcolor}{rgb}{0.000000,0.000000,0.000000}%
\pgfsetstrokecolor{textcolor}%
\pgfsetfillcolor{textcolor}%
\pgftext[x=0.407060in, y=0.827460in, left, base]{\color{textcolor}{\rmfamily\fontsize{11.000000}{13.200000}\selectfont\catcode`\^=\active\def^{\ifmmode\sp\else\^{}\fi}\catcode`\%=\active\def%{\%}0.0}}%
\end{pgfscope}%
\begin{pgfscope}%
\pgfsetbuttcap%
\pgfsetroundjoin%
\definecolor{currentfill}{rgb}{0.000000,0.000000,0.000000}%
\pgfsetfillcolor{currentfill}%
\pgfsetlinewidth{0.501875pt}%
\definecolor{currentstroke}{rgb}{0.000000,0.000000,0.000000}%
\pgfsetstrokecolor{currentstroke}%
\pgfsetdash{}{0pt}%
\pgfsys@defobject{currentmarker}{\pgfqpoint{0.000000in}{0.000000in}}{\pgfqpoint{0.041667in}{0.000000in}}{%
\pgfpathmoveto{\pgfqpoint{0.000000in}{0.000000in}}%
\pgfpathlineto{\pgfqpoint{0.041667in}{0.000000in}}%
\pgfusepath{stroke,fill}%
}%
\begin{pgfscope}%
\pgfsys@transformshift{0.650000in}{1.263822in}%
\pgfsys@useobject{currentmarker}{}%
\end{pgfscope}%
\end{pgfscope}%
\begin{pgfscope}%
\pgfsetbuttcap%
\pgfsetroundjoin%
\definecolor{currentfill}{rgb}{0.000000,0.000000,0.000000}%
\pgfsetfillcolor{currentfill}%
\pgfsetlinewidth{0.501875pt}%
\definecolor{currentstroke}{rgb}{0.000000,0.000000,0.000000}%
\pgfsetstrokecolor{currentstroke}%
\pgfsetdash{}{0pt}%
\pgfsys@defobject{currentmarker}{\pgfqpoint{-0.041667in}{0.000000in}}{\pgfqpoint{-0.000000in}{0.000000in}}{%
\pgfpathmoveto{\pgfqpoint{-0.000000in}{0.000000in}}%
\pgfpathlineto{\pgfqpoint{-0.041667in}{0.000000in}}%
\pgfusepath{stroke,fill}%
}%
\begin{pgfscope}%
\pgfsys@transformshift{3.038110in}{1.263822in}%
\pgfsys@useobject{currentmarker}{}%
\end{pgfscope}%
\end{pgfscope}%
\begin{pgfscope}%
\definecolor{textcolor}{rgb}{0.000000,0.000000,0.000000}%
\pgfsetstrokecolor{textcolor}%
\pgfsetfillcolor{textcolor}%
\pgftext[x=0.407060in, y=1.211015in, left, base]{\color{textcolor}{\rmfamily\fontsize{11.000000}{13.200000}\selectfont\catcode`\^=\active\def^{\ifmmode\sp\else\^{}\fi}\catcode`\%=\active\def%{\%}2.5}}%
\end{pgfscope}%
\begin{pgfscope}%
\pgfsetbuttcap%
\pgfsetroundjoin%
\definecolor{currentfill}{rgb}{0.000000,0.000000,0.000000}%
\pgfsetfillcolor{currentfill}%
\pgfsetlinewidth{0.501875pt}%
\definecolor{currentstroke}{rgb}{0.000000,0.000000,0.000000}%
\pgfsetstrokecolor{currentstroke}%
\pgfsetdash{}{0pt}%
\pgfsys@defobject{currentmarker}{\pgfqpoint{0.000000in}{0.000000in}}{\pgfqpoint{0.041667in}{0.000000in}}{%
\pgfpathmoveto{\pgfqpoint{0.000000in}{0.000000in}}%
\pgfpathlineto{\pgfqpoint{0.041667in}{0.000000in}}%
\pgfusepath{stroke,fill}%
}%
\begin{pgfscope}%
\pgfsys@transformshift{0.650000in}{1.647377in}%
\pgfsys@useobject{currentmarker}{}%
\end{pgfscope}%
\end{pgfscope}%
\begin{pgfscope}%
\pgfsetbuttcap%
\pgfsetroundjoin%
\definecolor{currentfill}{rgb}{0.000000,0.000000,0.000000}%
\pgfsetfillcolor{currentfill}%
\pgfsetlinewidth{0.501875pt}%
\definecolor{currentstroke}{rgb}{0.000000,0.000000,0.000000}%
\pgfsetstrokecolor{currentstroke}%
\pgfsetdash{}{0pt}%
\pgfsys@defobject{currentmarker}{\pgfqpoint{-0.041667in}{0.000000in}}{\pgfqpoint{-0.000000in}{0.000000in}}{%
\pgfpathmoveto{\pgfqpoint{-0.000000in}{0.000000in}}%
\pgfpathlineto{\pgfqpoint{-0.041667in}{0.000000in}}%
\pgfusepath{stroke,fill}%
}%
\begin{pgfscope}%
\pgfsys@transformshift{3.038110in}{1.647377in}%
\pgfsys@useobject{currentmarker}{}%
\end{pgfscope}%
\end{pgfscope}%
\begin{pgfscope}%
\definecolor{textcolor}{rgb}{0.000000,0.000000,0.000000}%
\pgfsetstrokecolor{textcolor}%
\pgfsetfillcolor{textcolor}%
\pgftext[x=0.407060in, y=1.594571in, left, base]{\color{textcolor}{\rmfamily\fontsize{11.000000}{13.200000}\selectfont\catcode`\^=\active\def^{\ifmmode\sp\else\^{}\fi}\catcode`\%=\active\def%{\%}5.0}}%
\end{pgfscope}%
\begin{pgfscope}%
\pgfsetbuttcap%
\pgfsetroundjoin%
\definecolor{currentfill}{rgb}{0.000000,0.000000,0.000000}%
\pgfsetfillcolor{currentfill}%
\pgfsetlinewidth{0.501875pt}%
\definecolor{currentstroke}{rgb}{0.000000,0.000000,0.000000}%
\pgfsetstrokecolor{currentstroke}%
\pgfsetdash{}{0pt}%
\pgfsys@defobject{currentmarker}{\pgfqpoint{0.000000in}{0.000000in}}{\pgfqpoint{0.041667in}{0.000000in}}{%
\pgfpathmoveto{\pgfqpoint{0.000000in}{0.000000in}}%
\pgfpathlineto{\pgfqpoint{0.041667in}{0.000000in}}%
\pgfusepath{stroke,fill}%
}%
\begin{pgfscope}%
\pgfsys@transformshift{0.650000in}{2.030933in}%
\pgfsys@useobject{currentmarker}{}%
\end{pgfscope}%
\end{pgfscope}%
\begin{pgfscope}%
\pgfsetbuttcap%
\pgfsetroundjoin%
\definecolor{currentfill}{rgb}{0.000000,0.000000,0.000000}%
\pgfsetfillcolor{currentfill}%
\pgfsetlinewidth{0.501875pt}%
\definecolor{currentstroke}{rgb}{0.000000,0.000000,0.000000}%
\pgfsetstrokecolor{currentstroke}%
\pgfsetdash{}{0pt}%
\pgfsys@defobject{currentmarker}{\pgfqpoint{-0.041667in}{0.000000in}}{\pgfqpoint{-0.000000in}{0.000000in}}{%
\pgfpathmoveto{\pgfqpoint{-0.000000in}{0.000000in}}%
\pgfpathlineto{\pgfqpoint{-0.041667in}{0.000000in}}%
\pgfusepath{stroke,fill}%
}%
\begin{pgfscope}%
\pgfsys@transformshift{3.038110in}{2.030933in}%
\pgfsys@useobject{currentmarker}{}%
\end{pgfscope}%
\end{pgfscope}%
\begin{pgfscope}%
\definecolor{textcolor}{rgb}{0.000000,0.000000,0.000000}%
\pgfsetstrokecolor{textcolor}%
\pgfsetfillcolor{textcolor}%
\pgftext[x=0.407060in, y=1.978126in, left, base]{\color{textcolor}{\rmfamily\fontsize{11.000000}{13.200000}\selectfont\catcode`\^=\active\def^{\ifmmode\sp\else\^{}\fi}\catcode`\%=\active\def%{\%}7.5}}%
\end{pgfscope}%
\begin{pgfscope}%
\pgfsetbuttcap%
\pgfsetroundjoin%
\definecolor{currentfill}{rgb}{0.000000,0.000000,0.000000}%
\pgfsetfillcolor{currentfill}%
\pgfsetlinewidth{0.501875pt}%
\definecolor{currentstroke}{rgb}{0.000000,0.000000,0.000000}%
\pgfsetstrokecolor{currentstroke}%
\pgfsetdash{}{0pt}%
\pgfsys@defobject{currentmarker}{\pgfqpoint{0.000000in}{0.000000in}}{\pgfqpoint{0.041667in}{0.000000in}}{%
\pgfpathmoveto{\pgfqpoint{0.000000in}{0.000000in}}%
\pgfpathlineto{\pgfqpoint{0.041667in}{0.000000in}}%
\pgfusepath{stroke,fill}%
}%
\begin{pgfscope}%
\pgfsys@transformshift{0.650000in}{2.414488in}%
\pgfsys@useobject{currentmarker}{}%
\end{pgfscope}%
\end{pgfscope}%
\begin{pgfscope}%
\pgfsetbuttcap%
\pgfsetroundjoin%
\definecolor{currentfill}{rgb}{0.000000,0.000000,0.000000}%
\pgfsetfillcolor{currentfill}%
\pgfsetlinewidth{0.501875pt}%
\definecolor{currentstroke}{rgb}{0.000000,0.000000,0.000000}%
\pgfsetstrokecolor{currentstroke}%
\pgfsetdash{}{0pt}%
\pgfsys@defobject{currentmarker}{\pgfqpoint{-0.041667in}{0.000000in}}{\pgfqpoint{-0.000000in}{0.000000in}}{%
\pgfpathmoveto{\pgfqpoint{-0.000000in}{0.000000in}}%
\pgfpathlineto{\pgfqpoint{-0.041667in}{0.000000in}}%
\pgfusepath{stroke,fill}%
}%
\begin{pgfscope}%
\pgfsys@transformshift{3.038110in}{2.414488in}%
\pgfsys@useobject{currentmarker}{}%
\end{pgfscope}%
\end{pgfscope}%
\begin{pgfscope}%
\definecolor{textcolor}{rgb}{0.000000,0.000000,0.000000}%
\pgfsetstrokecolor{textcolor}%
\pgfsetfillcolor{textcolor}%
\pgftext[x=0.331018in, y=2.361681in, left, base]{\color{textcolor}{\rmfamily\fontsize{11.000000}{13.200000}\selectfont\catcode`\^=\active\def^{\ifmmode\sp\else\^{}\fi}\catcode`\%=\active\def%{\%}10.0}}%
\end{pgfscope}%
\begin{pgfscope}%
\pgfsetbuttcap%
\pgfsetroundjoin%
\definecolor{currentfill}{rgb}{0.000000,0.000000,0.000000}%
\pgfsetfillcolor{currentfill}%
\pgfsetlinewidth{0.401500pt}%
\definecolor{currentstroke}{rgb}{0.000000,0.000000,0.000000}%
\pgfsetstrokecolor{currentstroke}%
\pgfsetdash{}{0pt}%
\pgfsys@defobject{currentmarker}{\pgfqpoint{0.000000in}{0.000000in}}{\pgfqpoint{0.020833in}{0.000000in}}{%
\pgfpathmoveto{\pgfqpoint{0.000000in}{0.000000in}}%
\pgfpathlineto{\pgfqpoint{0.020833in}{0.000000in}}%
\pgfusepath{stroke,fill}%
}%
\begin{pgfscope}%
\pgfsys@transformshift{0.650000in}{0.420000in}%
\pgfsys@useobject{currentmarker}{}%
\end{pgfscope}%
\end{pgfscope}%
\begin{pgfscope}%
\pgfsetbuttcap%
\pgfsetroundjoin%
\definecolor{currentfill}{rgb}{0.000000,0.000000,0.000000}%
\pgfsetfillcolor{currentfill}%
\pgfsetlinewidth{0.401500pt}%
\definecolor{currentstroke}{rgb}{0.000000,0.000000,0.000000}%
\pgfsetstrokecolor{currentstroke}%
\pgfsetdash{}{0pt}%
\pgfsys@defobject{currentmarker}{\pgfqpoint{-0.020833in}{0.000000in}}{\pgfqpoint{-0.000000in}{0.000000in}}{%
\pgfpathmoveto{\pgfqpoint{-0.000000in}{0.000000in}}%
\pgfpathlineto{\pgfqpoint{-0.020833in}{0.000000in}}%
\pgfusepath{stroke,fill}%
}%
\begin{pgfscope}%
\pgfsys@transformshift{3.038110in}{0.420000in}%
\pgfsys@useobject{currentmarker}{}%
\end{pgfscope}%
\end{pgfscope}%
\begin{pgfscope}%
\pgfsetbuttcap%
\pgfsetroundjoin%
\definecolor{currentfill}{rgb}{0.000000,0.000000,0.000000}%
\pgfsetfillcolor{currentfill}%
\pgfsetlinewidth{0.401500pt}%
\definecolor{currentstroke}{rgb}{0.000000,0.000000,0.000000}%
\pgfsetstrokecolor{currentstroke}%
\pgfsetdash{}{0pt}%
\pgfsys@defobject{currentmarker}{\pgfqpoint{0.000000in}{0.000000in}}{\pgfqpoint{0.020833in}{0.000000in}}{%
\pgfpathmoveto{\pgfqpoint{0.000000in}{0.000000in}}%
\pgfpathlineto{\pgfqpoint{0.020833in}{0.000000in}}%
\pgfusepath{stroke,fill}%
}%
\begin{pgfscope}%
\pgfsys@transformshift{0.650000in}{0.573422in}%
\pgfsys@useobject{currentmarker}{}%
\end{pgfscope}%
\end{pgfscope}%
\begin{pgfscope}%
\pgfsetbuttcap%
\pgfsetroundjoin%
\definecolor{currentfill}{rgb}{0.000000,0.000000,0.000000}%
\pgfsetfillcolor{currentfill}%
\pgfsetlinewidth{0.401500pt}%
\definecolor{currentstroke}{rgb}{0.000000,0.000000,0.000000}%
\pgfsetstrokecolor{currentstroke}%
\pgfsetdash{}{0pt}%
\pgfsys@defobject{currentmarker}{\pgfqpoint{-0.020833in}{0.000000in}}{\pgfqpoint{-0.000000in}{0.000000in}}{%
\pgfpathmoveto{\pgfqpoint{-0.000000in}{0.000000in}}%
\pgfpathlineto{\pgfqpoint{-0.020833in}{0.000000in}}%
\pgfusepath{stroke,fill}%
}%
\begin{pgfscope}%
\pgfsys@transformshift{3.038110in}{0.573422in}%
\pgfsys@useobject{currentmarker}{}%
\end{pgfscope}%
\end{pgfscope}%
\begin{pgfscope}%
\pgfsetbuttcap%
\pgfsetroundjoin%
\definecolor{currentfill}{rgb}{0.000000,0.000000,0.000000}%
\pgfsetfillcolor{currentfill}%
\pgfsetlinewidth{0.401500pt}%
\definecolor{currentstroke}{rgb}{0.000000,0.000000,0.000000}%
\pgfsetstrokecolor{currentstroke}%
\pgfsetdash{}{0pt}%
\pgfsys@defobject{currentmarker}{\pgfqpoint{0.000000in}{0.000000in}}{\pgfqpoint{0.020833in}{0.000000in}}{%
\pgfpathmoveto{\pgfqpoint{0.000000in}{0.000000in}}%
\pgfpathlineto{\pgfqpoint{0.020833in}{0.000000in}}%
\pgfusepath{stroke,fill}%
}%
\begin{pgfscope}%
\pgfsys@transformshift{0.650000in}{0.650133in}%
\pgfsys@useobject{currentmarker}{}%
\end{pgfscope}%
\end{pgfscope}%
\begin{pgfscope}%
\pgfsetbuttcap%
\pgfsetroundjoin%
\definecolor{currentfill}{rgb}{0.000000,0.000000,0.000000}%
\pgfsetfillcolor{currentfill}%
\pgfsetlinewidth{0.401500pt}%
\definecolor{currentstroke}{rgb}{0.000000,0.000000,0.000000}%
\pgfsetstrokecolor{currentstroke}%
\pgfsetdash{}{0pt}%
\pgfsys@defobject{currentmarker}{\pgfqpoint{-0.020833in}{0.000000in}}{\pgfqpoint{-0.000000in}{0.000000in}}{%
\pgfpathmoveto{\pgfqpoint{-0.000000in}{0.000000in}}%
\pgfpathlineto{\pgfqpoint{-0.020833in}{0.000000in}}%
\pgfusepath{stroke,fill}%
}%
\begin{pgfscope}%
\pgfsys@transformshift{3.038110in}{0.650133in}%
\pgfsys@useobject{currentmarker}{}%
\end{pgfscope}%
\end{pgfscope}%
\begin{pgfscope}%
\pgfsetbuttcap%
\pgfsetroundjoin%
\definecolor{currentfill}{rgb}{0.000000,0.000000,0.000000}%
\pgfsetfillcolor{currentfill}%
\pgfsetlinewidth{0.401500pt}%
\definecolor{currentstroke}{rgb}{0.000000,0.000000,0.000000}%
\pgfsetstrokecolor{currentstroke}%
\pgfsetdash{}{0pt}%
\pgfsys@defobject{currentmarker}{\pgfqpoint{0.000000in}{0.000000in}}{\pgfqpoint{0.020833in}{0.000000in}}{%
\pgfpathmoveto{\pgfqpoint{0.000000in}{0.000000in}}%
\pgfpathlineto{\pgfqpoint{0.020833in}{0.000000in}}%
\pgfusepath{stroke,fill}%
}%
\begin{pgfscope}%
\pgfsys@transformshift{0.650000in}{0.726844in}%
\pgfsys@useobject{currentmarker}{}%
\end{pgfscope}%
\end{pgfscope}%
\begin{pgfscope}%
\pgfsetbuttcap%
\pgfsetroundjoin%
\definecolor{currentfill}{rgb}{0.000000,0.000000,0.000000}%
\pgfsetfillcolor{currentfill}%
\pgfsetlinewidth{0.401500pt}%
\definecolor{currentstroke}{rgb}{0.000000,0.000000,0.000000}%
\pgfsetstrokecolor{currentstroke}%
\pgfsetdash{}{0pt}%
\pgfsys@defobject{currentmarker}{\pgfqpoint{-0.020833in}{0.000000in}}{\pgfqpoint{-0.000000in}{0.000000in}}{%
\pgfpathmoveto{\pgfqpoint{-0.000000in}{0.000000in}}%
\pgfpathlineto{\pgfqpoint{-0.020833in}{0.000000in}}%
\pgfusepath{stroke,fill}%
}%
\begin{pgfscope}%
\pgfsys@transformshift{3.038110in}{0.726844in}%
\pgfsys@useobject{currentmarker}{}%
\end{pgfscope}%
\end{pgfscope}%
\begin{pgfscope}%
\pgfsetbuttcap%
\pgfsetroundjoin%
\definecolor{currentfill}{rgb}{0.000000,0.000000,0.000000}%
\pgfsetfillcolor{currentfill}%
\pgfsetlinewidth{0.401500pt}%
\definecolor{currentstroke}{rgb}{0.000000,0.000000,0.000000}%
\pgfsetstrokecolor{currentstroke}%
\pgfsetdash{}{0pt}%
\pgfsys@defobject{currentmarker}{\pgfqpoint{0.000000in}{0.000000in}}{\pgfqpoint{0.020833in}{0.000000in}}{%
\pgfpathmoveto{\pgfqpoint{0.000000in}{0.000000in}}%
\pgfpathlineto{\pgfqpoint{0.020833in}{0.000000in}}%
\pgfusepath{stroke,fill}%
}%
\begin{pgfscope}%
\pgfsys@transformshift{0.650000in}{0.803555in}%
\pgfsys@useobject{currentmarker}{}%
\end{pgfscope}%
\end{pgfscope}%
\begin{pgfscope}%
\pgfsetbuttcap%
\pgfsetroundjoin%
\definecolor{currentfill}{rgb}{0.000000,0.000000,0.000000}%
\pgfsetfillcolor{currentfill}%
\pgfsetlinewidth{0.401500pt}%
\definecolor{currentstroke}{rgb}{0.000000,0.000000,0.000000}%
\pgfsetstrokecolor{currentstroke}%
\pgfsetdash{}{0pt}%
\pgfsys@defobject{currentmarker}{\pgfqpoint{-0.020833in}{0.000000in}}{\pgfqpoint{-0.000000in}{0.000000in}}{%
\pgfpathmoveto{\pgfqpoint{-0.000000in}{0.000000in}}%
\pgfpathlineto{\pgfqpoint{-0.020833in}{0.000000in}}%
\pgfusepath{stroke,fill}%
}%
\begin{pgfscope}%
\pgfsys@transformshift{3.038110in}{0.803555in}%
\pgfsys@useobject{currentmarker}{}%
\end{pgfscope}%
\end{pgfscope}%
\begin{pgfscope}%
\pgfsetbuttcap%
\pgfsetroundjoin%
\definecolor{currentfill}{rgb}{0.000000,0.000000,0.000000}%
\pgfsetfillcolor{currentfill}%
\pgfsetlinewidth{0.401500pt}%
\definecolor{currentstroke}{rgb}{0.000000,0.000000,0.000000}%
\pgfsetstrokecolor{currentstroke}%
\pgfsetdash{}{0pt}%
\pgfsys@defobject{currentmarker}{\pgfqpoint{0.000000in}{0.000000in}}{\pgfqpoint{0.020833in}{0.000000in}}{%
\pgfpathmoveto{\pgfqpoint{0.000000in}{0.000000in}}%
\pgfpathlineto{\pgfqpoint{0.020833in}{0.000000in}}%
\pgfusepath{stroke,fill}%
}%
\begin{pgfscope}%
\pgfsys@transformshift{0.650000in}{0.956978in}%
\pgfsys@useobject{currentmarker}{}%
\end{pgfscope}%
\end{pgfscope}%
\begin{pgfscope}%
\pgfsetbuttcap%
\pgfsetroundjoin%
\definecolor{currentfill}{rgb}{0.000000,0.000000,0.000000}%
\pgfsetfillcolor{currentfill}%
\pgfsetlinewidth{0.401500pt}%
\definecolor{currentstroke}{rgb}{0.000000,0.000000,0.000000}%
\pgfsetstrokecolor{currentstroke}%
\pgfsetdash{}{0pt}%
\pgfsys@defobject{currentmarker}{\pgfqpoint{-0.020833in}{0.000000in}}{\pgfqpoint{-0.000000in}{0.000000in}}{%
\pgfpathmoveto{\pgfqpoint{-0.000000in}{0.000000in}}%
\pgfpathlineto{\pgfqpoint{-0.020833in}{0.000000in}}%
\pgfusepath{stroke,fill}%
}%
\begin{pgfscope}%
\pgfsys@transformshift{3.038110in}{0.956978in}%
\pgfsys@useobject{currentmarker}{}%
\end{pgfscope}%
\end{pgfscope}%
\begin{pgfscope}%
\pgfsetbuttcap%
\pgfsetroundjoin%
\definecolor{currentfill}{rgb}{0.000000,0.000000,0.000000}%
\pgfsetfillcolor{currentfill}%
\pgfsetlinewidth{0.401500pt}%
\definecolor{currentstroke}{rgb}{0.000000,0.000000,0.000000}%
\pgfsetstrokecolor{currentstroke}%
\pgfsetdash{}{0pt}%
\pgfsys@defobject{currentmarker}{\pgfqpoint{0.000000in}{0.000000in}}{\pgfqpoint{0.020833in}{0.000000in}}{%
\pgfpathmoveto{\pgfqpoint{0.000000in}{0.000000in}}%
\pgfpathlineto{\pgfqpoint{0.020833in}{0.000000in}}%
\pgfusepath{stroke,fill}%
}%
\begin{pgfscope}%
\pgfsys@transformshift{0.650000in}{1.033689in}%
\pgfsys@useobject{currentmarker}{}%
\end{pgfscope}%
\end{pgfscope}%
\begin{pgfscope}%
\pgfsetbuttcap%
\pgfsetroundjoin%
\definecolor{currentfill}{rgb}{0.000000,0.000000,0.000000}%
\pgfsetfillcolor{currentfill}%
\pgfsetlinewidth{0.401500pt}%
\definecolor{currentstroke}{rgb}{0.000000,0.000000,0.000000}%
\pgfsetstrokecolor{currentstroke}%
\pgfsetdash{}{0pt}%
\pgfsys@defobject{currentmarker}{\pgfqpoint{-0.020833in}{0.000000in}}{\pgfqpoint{-0.000000in}{0.000000in}}{%
\pgfpathmoveto{\pgfqpoint{-0.000000in}{0.000000in}}%
\pgfpathlineto{\pgfqpoint{-0.020833in}{0.000000in}}%
\pgfusepath{stroke,fill}%
}%
\begin{pgfscope}%
\pgfsys@transformshift{3.038110in}{1.033689in}%
\pgfsys@useobject{currentmarker}{}%
\end{pgfscope}%
\end{pgfscope}%
\begin{pgfscope}%
\pgfsetbuttcap%
\pgfsetroundjoin%
\definecolor{currentfill}{rgb}{0.000000,0.000000,0.000000}%
\pgfsetfillcolor{currentfill}%
\pgfsetlinewidth{0.401500pt}%
\definecolor{currentstroke}{rgb}{0.000000,0.000000,0.000000}%
\pgfsetstrokecolor{currentstroke}%
\pgfsetdash{}{0pt}%
\pgfsys@defobject{currentmarker}{\pgfqpoint{0.000000in}{0.000000in}}{\pgfqpoint{0.020833in}{0.000000in}}{%
\pgfpathmoveto{\pgfqpoint{0.000000in}{0.000000in}}%
\pgfpathlineto{\pgfqpoint{0.020833in}{0.000000in}}%
\pgfusepath{stroke,fill}%
}%
\begin{pgfscope}%
\pgfsys@transformshift{0.650000in}{1.110400in}%
\pgfsys@useobject{currentmarker}{}%
\end{pgfscope}%
\end{pgfscope}%
\begin{pgfscope}%
\pgfsetbuttcap%
\pgfsetroundjoin%
\definecolor{currentfill}{rgb}{0.000000,0.000000,0.000000}%
\pgfsetfillcolor{currentfill}%
\pgfsetlinewidth{0.401500pt}%
\definecolor{currentstroke}{rgb}{0.000000,0.000000,0.000000}%
\pgfsetstrokecolor{currentstroke}%
\pgfsetdash{}{0pt}%
\pgfsys@defobject{currentmarker}{\pgfqpoint{-0.020833in}{0.000000in}}{\pgfqpoint{-0.000000in}{0.000000in}}{%
\pgfpathmoveto{\pgfqpoint{-0.000000in}{0.000000in}}%
\pgfpathlineto{\pgfqpoint{-0.020833in}{0.000000in}}%
\pgfusepath{stroke,fill}%
}%
\begin{pgfscope}%
\pgfsys@transformshift{3.038110in}{1.110400in}%
\pgfsys@useobject{currentmarker}{}%
\end{pgfscope}%
\end{pgfscope}%
\begin{pgfscope}%
\pgfsetbuttcap%
\pgfsetroundjoin%
\definecolor{currentfill}{rgb}{0.000000,0.000000,0.000000}%
\pgfsetfillcolor{currentfill}%
\pgfsetlinewidth{0.401500pt}%
\definecolor{currentstroke}{rgb}{0.000000,0.000000,0.000000}%
\pgfsetstrokecolor{currentstroke}%
\pgfsetdash{}{0pt}%
\pgfsys@defobject{currentmarker}{\pgfqpoint{0.000000in}{0.000000in}}{\pgfqpoint{0.020833in}{0.000000in}}{%
\pgfpathmoveto{\pgfqpoint{0.000000in}{0.000000in}}%
\pgfpathlineto{\pgfqpoint{0.020833in}{0.000000in}}%
\pgfusepath{stroke,fill}%
}%
\begin{pgfscope}%
\pgfsys@transformshift{0.650000in}{1.187111in}%
\pgfsys@useobject{currentmarker}{}%
\end{pgfscope}%
\end{pgfscope}%
\begin{pgfscope}%
\pgfsetbuttcap%
\pgfsetroundjoin%
\definecolor{currentfill}{rgb}{0.000000,0.000000,0.000000}%
\pgfsetfillcolor{currentfill}%
\pgfsetlinewidth{0.401500pt}%
\definecolor{currentstroke}{rgb}{0.000000,0.000000,0.000000}%
\pgfsetstrokecolor{currentstroke}%
\pgfsetdash{}{0pt}%
\pgfsys@defobject{currentmarker}{\pgfqpoint{-0.020833in}{0.000000in}}{\pgfqpoint{-0.000000in}{0.000000in}}{%
\pgfpathmoveto{\pgfqpoint{-0.000000in}{0.000000in}}%
\pgfpathlineto{\pgfqpoint{-0.020833in}{0.000000in}}%
\pgfusepath{stroke,fill}%
}%
\begin{pgfscope}%
\pgfsys@transformshift{3.038110in}{1.187111in}%
\pgfsys@useobject{currentmarker}{}%
\end{pgfscope}%
\end{pgfscope}%
\begin{pgfscope}%
\pgfsetbuttcap%
\pgfsetroundjoin%
\definecolor{currentfill}{rgb}{0.000000,0.000000,0.000000}%
\pgfsetfillcolor{currentfill}%
\pgfsetlinewidth{0.401500pt}%
\definecolor{currentstroke}{rgb}{0.000000,0.000000,0.000000}%
\pgfsetstrokecolor{currentstroke}%
\pgfsetdash{}{0pt}%
\pgfsys@defobject{currentmarker}{\pgfqpoint{0.000000in}{0.000000in}}{\pgfqpoint{0.020833in}{0.000000in}}{%
\pgfpathmoveto{\pgfqpoint{0.000000in}{0.000000in}}%
\pgfpathlineto{\pgfqpoint{0.020833in}{0.000000in}}%
\pgfusepath{stroke,fill}%
}%
\begin{pgfscope}%
\pgfsys@transformshift{0.650000in}{1.340533in}%
\pgfsys@useobject{currentmarker}{}%
\end{pgfscope}%
\end{pgfscope}%
\begin{pgfscope}%
\pgfsetbuttcap%
\pgfsetroundjoin%
\definecolor{currentfill}{rgb}{0.000000,0.000000,0.000000}%
\pgfsetfillcolor{currentfill}%
\pgfsetlinewidth{0.401500pt}%
\definecolor{currentstroke}{rgb}{0.000000,0.000000,0.000000}%
\pgfsetstrokecolor{currentstroke}%
\pgfsetdash{}{0pt}%
\pgfsys@defobject{currentmarker}{\pgfqpoint{-0.020833in}{0.000000in}}{\pgfqpoint{-0.000000in}{0.000000in}}{%
\pgfpathmoveto{\pgfqpoint{-0.000000in}{0.000000in}}%
\pgfpathlineto{\pgfqpoint{-0.020833in}{0.000000in}}%
\pgfusepath{stroke,fill}%
}%
\begin{pgfscope}%
\pgfsys@transformshift{3.038110in}{1.340533in}%
\pgfsys@useobject{currentmarker}{}%
\end{pgfscope}%
\end{pgfscope}%
\begin{pgfscope}%
\pgfsetbuttcap%
\pgfsetroundjoin%
\definecolor{currentfill}{rgb}{0.000000,0.000000,0.000000}%
\pgfsetfillcolor{currentfill}%
\pgfsetlinewidth{0.401500pt}%
\definecolor{currentstroke}{rgb}{0.000000,0.000000,0.000000}%
\pgfsetstrokecolor{currentstroke}%
\pgfsetdash{}{0pt}%
\pgfsys@defobject{currentmarker}{\pgfqpoint{0.000000in}{0.000000in}}{\pgfqpoint{0.020833in}{0.000000in}}{%
\pgfpathmoveto{\pgfqpoint{0.000000in}{0.000000in}}%
\pgfpathlineto{\pgfqpoint{0.020833in}{0.000000in}}%
\pgfusepath{stroke,fill}%
}%
\begin{pgfscope}%
\pgfsys@transformshift{0.650000in}{1.417244in}%
\pgfsys@useobject{currentmarker}{}%
\end{pgfscope}%
\end{pgfscope}%
\begin{pgfscope}%
\pgfsetbuttcap%
\pgfsetroundjoin%
\definecolor{currentfill}{rgb}{0.000000,0.000000,0.000000}%
\pgfsetfillcolor{currentfill}%
\pgfsetlinewidth{0.401500pt}%
\definecolor{currentstroke}{rgb}{0.000000,0.000000,0.000000}%
\pgfsetstrokecolor{currentstroke}%
\pgfsetdash{}{0pt}%
\pgfsys@defobject{currentmarker}{\pgfqpoint{-0.020833in}{0.000000in}}{\pgfqpoint{-0.000000in}{0.000000in}}{%
\pgfpathmoveto{\pgfqpoint{-0.000000in}{0.000000in}}%
\pgfpathlineto{\pgfqpoint{-0.020833in}{0.000000in}}%
\pgfusepath{stroke,fill}%
}%
\begin{pgfscope}%
\pgfsys@transformshift{3.038110in}{1.417244in}%
\pgfsys@useobject{currentmarker}{}%
\end{pgfscope}%
\end{pgfscope}%
\begin{pgfscope}%
\pgfsetbuttcap%
\pgfsetroundjoin%
\definecolor{currentfill}{rgb}{0.000000,0.000000,0.000000}%
\pgfsetfillcolor{currentfill}%
\pgfsetlinewidth{0.401500pt}%
\definecolor{currentstroke}{rgb}{0.000000,0.000000,0.000000}%
\pgfsetstrokecolor{currentstroke}%
\pgfsetdash{}{0pt}%
\pgfsys@defobject{currentmarker}{\pgfqpoint{0.000000in}{0.000000in}}{\pgfqpoint{0.020833in}{0.000000in}}{%
\pgfpathmoveto{\pgfqpoint{0.000000in}{0.000000in}}%
\pgfpathlineto{\pgfqpoint{0.020833in}{0.000000in}}%
\pgfusepath{stroke,fill}%
}%
\begin{pgfscope}%
\pgfsys@transformshift{0.650000in}{1.493955in}%
\pgfsys@useobject{currentmarker}{}%
\end{pgfscope}%
\end{pgfscope}%
\begin{pgfscope}%
\pgfsetbuttcap%
\pgfsetroundjoin%
\definecolor{currentfill}{rgb}{0.000000,0.000000,0.000000}%
\pgfsetfillcolor{currentfill}%
\pgfsetlinewidth{0.401500pt}%
\definecolor{currentstroke}{rgb}{0.000000,0.000000,0.000000}%
\pgfsetstrokecolor{currentstroke}%
\pgfsetdash{}{0pt}%
\pgfsys@defobject{currentmarker}{\pgfqpoint{-0.020833in}{0.000000in}}{\pgfqpoint{-0.000000in}{0.000000in}}{%
\pgfpathmoveto{\pgfqpoint{-0.000000in}{0.000000in}}%
\pgfpathlineto{\pgfqpoint{-0.020833in}{0.000000in}}%
\pgfusepath{stroke,fill}%
}%
\begin{pgfscope}%
\pgfsys@transformshift{3.038110in}{1.493955in}%
\pgfsys@useobject{currentmarker}{}%
\end{pgfscope}%
\end{pgfscope}%
\begin{pgfscope}%
\pgfsetbuttcap%
\pgfsetroundjoin%
\definecolor{currentfill}{rgb}{0.000000,0.000000,0.000000}%
\pgfsetfillcolor{currentfill}%
\pgfsetlinewidth{0.401500pt}%
\definecolor{currentstroke}{rgb}{0.000000,0.000000,0.000000}%
\pgfsetstrokecolor{currentstroke}%
\pgfsetdash{}{0pt}%
\pgfsys@defobject{currentmarker}{\pgfqpoint{0.000000in}{0.000000in}}{\pgfqpoint{0.020833in}{0.000000in}}{%
\pgfpathmoveto{\pgfqpoint{0.000000in}{0.000000in}}%
\pgfpathlineto{\pgfqpoint{0.020833in}{0.000000in}}%
\pgfusepath{stroke,fill}%
}%
\begin{pgfscope}%
\pgfsys@transformshift{0.650000in}{1.570666in}%
\pgfsys@useobject{currentmarker}{}%
\end{pgfscope}%
\end{pgfscope}%
\begin{pgfscope}%
\pgfsetbuttcap%
\pgfsetroundjoin%
\definecolor{currentfill}{rgb}{0.000000,0.000000,0.000000}%
\pgfsetfillcolor{currentfill}%
\pgfsetlinewidth{0.401500pt}%
\definecolor{currentstroke}{rgb}{0.000000,0.000000,0.000000}%
\pgfsetstrokecolor{currentstroke}%
\pgfsetdash{}{0pt}%
\pgfsys@defobject{currentmarker}{\pgfqpoint{-0.020833in}{0.000000in}}{\pgfqpoint{-0.000000in}{0.000000in}}{%
\pgfpathmoveto{\pgfqpoint{-0.000000in}{0.000000in}}%
\pgfpathlineto{\pgfqpoint{-0.020833in}{0.000000in}}%
\pgfusepath{stroke,fill}%
}%
\begin{pgfscope}%
\pgfsys@transformshift{3.038110in}{1.570666in}%
\pgfsys@useobject{currentmarker}{}%
\end{pgfscope}%
\end{pgfscope}%
\begin{pgfscope}%
\pgfsetbuttcap%
\pgfsetroundjoin%
\definecolor{currentfill}{rgb}{0.000000,0.000000,0.000000}%
\pgfsetfillcolor{currentfill}%
\pgfsetlinewidth{0.401500pt}%
\definecolor{currentstroke}{rgb}{0.000000,0.000000,0.000000}%
\pgfsetstrokecolor{currentstroke}%
\pgfsetdash{}{0pt}%
\pgfsys@defobject{currentmarker}{\pgfqpoint{0.000000in}{0.000000in}}{\pgfqpoint{0.020833in}{0.000000in}}{%
\pgfpathmoveto{\pgfqpoint{0.000000in}{0.000000in}}%
\pgfpathlineto{\pgfqpoint{0.020833in}{0.000000in}}%
\pgfusepath{stroke,fill}%
}%
\begin{pgfscope}%
\pgfsys@transformshift{0.650000in}{1.724088in}%
\pgfsys@useobject{currentmarker}{}%
\end{pgfscope}%
\end{pgfscope}%
\begin{pgfscope}%
\pgfsetbuttcap%
\pgfsetroundjoin%
\definecolor{currentfill}{rgb}{0.000000,0.000000,0.000000}%
\pgfsetfillcolor{currentfill}%
\pgfsetlinewidth{0.401500pt}%
\definecolor{currentstroke}{rgb}{0.000000,0.000000,0.000000}%
\pgfsetstrokecolor{currentstroke}%
\pgfsetdash{}{0pt}%
\pgfsys@defobject{currentmarker}{\pgfqpoint{-0.020833in}{0.000000in}}{\pgfqpoint{-0.000000in}{0.000000in}}{%
\pgfpathmoveto{\pgfqpoint{-0.000000in}{0.000000in}}%
\pgfpathlineto{\pgfqpoint{-0.020833in}{0.000000in}}%
\pgfusepath{stroke,fill}%
}%
\begin{pgfscope}%
\pgfsys@transformshift{3.038110in}{1.724088in}%
\pgfsys@useobject{currentmarker}{}%
\end{pgfscope}%
\end{pgfscope}%
\begin{pgfscope}%
\pgfsetbuttcap%
\pgfsetroundjoin%
\definecolor{currentfill}{rgb}{0.000000,0.000000,0.000000}%
\pgfsetfillcolor{currentfill}%
\pgfsetlinewidth{0.401500pt}%
\definecolor{currentstroke}{rgb}{0.000000,0.000000,0.000000}%
\pgfsetstrokecolor{currentstroke}%
\pgfsetdash{}{0pt}%
\pgfsys@defobject{currentmarker}{\pgfqpoint{0.000000in}{0.000000in}}{\pgfqpoint{0.020833in}{0.000000in}}{%
\pgfpathmoveto{\pgfqpoint{0.000000in}{0.000000in}}%
\pgfpathlineto{\pgfqpoint{0.020833in}{0.000000in}}%
\pgfusepath{stroke,fill}%
}%
\begin{pgfscope}%
\pgfsys@transformshift{0.650000in}{1.800799in}%
\pgfsys@useobject{currentmarker}{}%
\end{pgfscope}%
\end{pgfscope}%
\begin{pgfscope}%
\pgfsetbuttcap%
\pgfsetroundjoin%
\definecolor{currentfill}{rgb}{0.000000,0.000000,0.000000}%
\pgfsetfillcolor{currentfill}%
\pgfsetlinewidth{0.401500pt}%
\definecolor{currentstroke}{rgb}{0.000000,0.000000,0.000000}%
\pgfsetstrokecolor{currentstroke}%
\pgfsetdash{}{0pt}%
\pgfsys@defobject{currentmarker}{\pgfqpoint{-0.020833in}{0.000000in}}{\pgfqpoint{-0.000000in}{0.000000in}}{%
\pgfpathmoveto{\pgfqpoint{-0.000000in}{0.000000in}}%
\pgfpathlineto{\pgfqpoint{-0.020833in}{0.000000in}}%
\pgfusepath{stroke,fill}%
}%
\begin{pgfscope}%
\pgfsys@transformshift{3.038110in}{1.800799in}%
\pgfsys@useobject{currentmarker}{}%
\end{pgfscope}%
\end{pgfscope}%
\begin{pgfscope}%
\pgfsetbuttcap%
\pgfsetroundjoin%
\definecolor{currentfill}{rgb}{0.000000,0.000000,0.000000}%
\pgfsetfillcolor{currentfill}%
\pgfsetlinewidth{0.401500pt}%
\definecolor{currentstroke}{rgb}{0.000000,0.000000,0.000000}%
\pgfsetstrokecolor{currentstroke}%
\pgfsetdash{}{0pt}%
\pgfsys@defobject{currentmarker}{\pgfqpoint{0.000000in}{0.000000in}}{\pgfqpoint{0.020833in}{0.000000in}}{%
\pgfpathmoveto{\pgfqpoint{0.000000in}{0.000000in}}%
\pgfpathlineto{\pgfqpoint{0.020833in}{0.000000in}}%
\pgfusepath{stroke,fill}%
}%
\begin{pgfscope}%
\pgfsys@transformshift{0.650000in}{1.877511in}%
\pgfsys@useobject{currentmarker}{}%
\end{pgfscope}%
\end{pgfscope}%
\begin{pgfscope}%
\pgfsetbuttcap%
\pgfsetroundjoin%
\definecolor{currentfill}{rgb}{0.000000,0.000000,0.000000}%
\pgfsetfillcolor{currentfill}%
\pgfsetlinewidth{0.401500pt}%
\definecolor{currentstroke}{rgb}{0.000000,0.000000,0.000000}%
\pgfsetstrokecolor{currentstroke}%
\pgfsetdash{}{0pt}%
\pgfsys@defobject{currentmarker}{\pgfqpoint{-0.020833in}{0.000000in}}{\pgfqpoint{-0.000000in}{0.000000in}}{%
\pgfpathmoveto{\pgfqpoint{-0.000000in}{0.000000in}}%
\pgfpathlineto{\pgfqpoint{-0.020833in}{0.000000in}}%
\pgfusepath{stroke,fill}%
}%
\begin{pgfscope}%
\pgfsys@transformshift{3.038110in}{1.877511in}%
\pgfsys@useobject{currentmarker}{}%
\end{pgfscope}%
\end{pgfscope}%
\begin{pgfscope}%
\pgfsetbuttcap%
\pgfsetroundjoin%
\definecolor{currentfill}{rgb}{0.000000,0.000000,0.000000}%
\pgfsetfillcolor{currentfill}%
\pgfsetlinewidth{0.401500pt}%
\definecolor{currentstroke}{rgb}{0.000000,0.000000,0.000000}%
\pgfsetstrokecolor{currentstroke}%
\pgfsetdash{}{0pt}%
\pgfsys@defobject{currentmarker}{\pgfqpoint{0.000000in}{0.000000in}}{\pgfqpoint{0.020833in}{0.000000in}}{%
\pgfpathmoveto{\pgfqpoint{0.000000in}{0.000000in}}%
\pgfpathlineto{\pgfqpoint{0.020833in}{0.000000in}}%
\pgfusepath{stroke,fill}%
}%
\begin{pgfscope}%
\pgfsys@transformshift{0.650000in}{1.954222in}%
\pgfsys@useobject{currentmarker}{}%
\end{pgfscope}%
\end{pgfscope}%
\begin{pgfscope}%
\pgfsetbuttcap%
\pgfsetroundjoin%
\definecolor{currentfill}{rgb}{0.000000,0.000000,0.000000}%
\pgfsetfillcolor{currentfill}%
\pgfsetlinewidth{0.401500pt}%
\definecolor{currentstroke}{rgb}{0.000000,0.000000,0.000000}%
\pgfsetstrokecolor{currentstroke}%
\pgfsetdash{}{0pt}%
\pgfsys@defobject{currentmarker}{\pgfqpoint{-0.020833in}{0.000000in}}{\pgfqpoint{-0.000000in}{0.000000in}}{%
\pgfpathmoveto{\pgfqpoint{-0.000000in}{0.000000in}}%
\pgfpathlineto{\pgfqpoint{-0.020833in}{0.000000in}}%
\pgfusepath{stroke,fill}%
}%
\begin{pgfscope}%
\pgfsys@transformshift{3.038110in}{1.954222in}%
\pgfsys@useobject{currentmarker}{}%
\end{pgfscope}%
\end{pgfscope}%
\begin{pgfscope}%
\pgfsetbuttcap%
\pgfsetroundjoin%
\definecolor{currentfill}{rgb}{0.000000,0.000000,0.000000}%
\pgfsetfillcolor{currentfill}%
\pgfsetlinewidth{0.401500pt}%
\definecolor{currentstroke}{rgb}{0.000000,0.000000,0.000000}%
\pgfsetstrokecolor{currentstroke}%
\pgfsetdash{}{0pt}%
\pgfsys@defobject{currentmarker}{\pgfqpoint{0.000000in}{0.000000in}}{\pgfqpoint{0.020833in}{0.000000in}}{%
\pgfpathmoveto{\pgfqpoint{0.000000in}{0.000000in}}%
\pgfpathlineto{\pgfqpoint{0.020833in}{0.000000in}}%
\pgfusepath{stroke,fill}%
}%
\begin{pgfscope}%
\pgfsys@transformshift{0.650000in}{2.107644in}%
\pgfsys@useobject{currentmarker}{}%
\end{pgfscope}%
\end{pgfscope}%
\begin{pgfscope}%
\pgfsetbuttcap%
\pgfsetroundjoin%
\definecolor{currentfill}{rgb}{0.000000,0.000000,0.000000}%
\pgfsetfillcolor{currentfill}%
\pgfsetlinewidth{0.401500pt}%
\definecolor{currentstroke}{rgb}{0.000000,0.000000,0.000000}%
\pgfsetstrokecolor{currentstroke}%
\pgfsetdash{}{0pt}%
\pgfsys@defobject{currentmarker}{\pgfqpoint{-0.020833in}{0.000000in}}{\pgfqpoint{-0.000000in}{0.000000in}}{%
\pgfpathmoveto{\pgfqpoint{-0.000000in}{0.000000in}}%
\pgfpathlineto{\pgfqpoint{-0.020833in}{0.000000in}}%
\pgfusepath{stroke,fill}%
}%
\begin{pgfscope}%
\pgfsys@transformshift{3.038110in}{2.107644in}%
\pgfsys@useobject{currentmarker}{}%
\end{pgfscope}%
\end{pgfscope}%
\begin{pgfscope}%
\pgfsetbuttcap%
\pgfsetroundjoin%
\definecolor{currentfill}{rgb}{0.000000,0.000000,0.000000}%
\pgfsetfillcolor{currentfill}%
\pgfsetlinewidth{0.401500pt}%
\definecolor{currentstroke}{rgb}{0.000000,0.000000,0.000000}%
\pgfsetstrokecolor{currentstroke}%
\pgfsetdash{}{0pt}%
\pgfsys@defobject{currentmarker}{\pgfqpoint{0.000000in}{0.000000in}}{\pgfqpoint{0.020833in}{0.000000in}}{%
\pgfpathmoveto{\pgfqpoint{0.000000in}{0.000000in}}%
\pgfpathlineto{\pgfqpoint{0.020833in}{0.000000in}}%
\pgfusepath{stroke,fill}%
}%
\begin{pgfscope}%
\pgfsys@transformshift{0.650000in}{2.184355in}%
\pgfsys@useobject{currentmarker}{}%
\end{pgfscope}%
\end{pgfscope}%
\begin{pgfscope}%
\pgfsetbuttcap%
\pgfsetroundjoin%
\definecolor{currentfill}{rgb}{0.000000,0.000000,0.000000}%
\pgfsetfillcolor{currentfill}%
\pgfsetlinewidth{0.401500pt}%
\definecolor{currentstroke}{rgb}{0.000000,0.000000,0.000000}%
\pgfsetstrokecolor{currentstroke}%
\pgfsetdash{}{0pt}%
\pgfsys@defobject{currentmarker}{\pgfqpoint{-0.020833in}{0.000000in}}{\pgfqpoint{-0.000000in}{0.000000in}}{%
\pgfpathmoveto{\pgfqpoint{-0.000000in}{0.000000in}}%
\pgfpathlineto{\pgfqpoint{-0.020833in}{0.000000in}}%
\pgfusepath{stroke,fill}%
}%
\begin{pgfscope}%
\pgfsys@transformshift{3.038110in}{2.184355in}%
\pgfsys@useobject{currentmarker}{}%
\end{pgfscope}%
\end{pgfscope}%
\begin{pgfscope}%
\pgfsetbuttcap%
\pgfsetroundjoin%
\definecolor{currentfill}{rgb}{0.000000,0.000000,0.000000}%
\pgfsetfillcolor{currentfill}%
\pgfsetlinewidth{0.401500pt}%
\definecolor{currentstroke}{rgb}{0.000000,0.000000,0.000000}%
\pgfsetstrokecolor{currentstroke}%
\pgfsetdash{}{0pt}%
\pgfsys@defobject{currentmarker}{\pgfqpoint{0.000000in}{0.000000in}}{\pgfqpoint{0.020833in}{0.000000in}}{%
\pgfpathmoveto{\pgfqpoint{0.000000in}{0.000000in}}%
\pgfpathlineto{\pgfqpoint{0.020833in}{0.000000in}}%
\pgfusepath{stroke,fill}%
}%
\begin{pgfscope}%
\pgfsys@transformshift{0.650000in}{2.261066in}%
\pgfsys@useobject{currentmarker}{}%
\end{pgfscope}%
\end{pgfscope}%
\begin{pgfscope}%
\pgfsetbuttcap%
\pgfsetroundjoin%
\definecolor{currentfill}{rgb}{0.000000,0.000000,0.000000}%
\pgfsetfillcolor{currentfill}%
\pgfsetlinewidth{0.401500pt}%
\definecolor{currentstroke}{rgb}{0.000000,0.000000,0.000000}%
\pgfsetstrokecolor{currentstroke}%
\pgfsetdash{}{0pt}%
\pgfsys@defobject{currentmarker}{\pgfqpoint{-0.020833in}{0.000000in}}{\pgfqpoint{-0.000000in}{0.000000in}}{%
\pgfpathmoveto{\pgfqpoint{-0.000000in}{0.000000in}}%
\pgfpathlineto{\pgfqpoint{-0.020833in}{0.000000in}}%
\pgfusepath{stroke,fill}%
}%
\begin{pgfscope}%
\pgfsys@transformshift{3.038110in}{2.261066in}%
\pgfsys@useobject{currentmarker}{}%
\end{pgfscope}%
\end{pgfscope}%
\begin{pgfscope}%
\pgfsetbuttcap%
\pgfsetroundjoin%
\definecolor{currentfill}{rgb}{0.000000,0.000000,0.000000}%
\pgfsetfillcolor{currentfill}%
\pgfsetlinewidth{0.401500pt}%
\definecolor{currentstroke}{rgb}{0.000000,0.000000,0.000000}%
\pgfsetstrokecolor{currentstroke}%
\pgfsetdash{}{0pt}%
\pgfsys@defobject{currentmarker}{\pgfqpoint{0.000000in}{0.000000in}}{\pgfqpoint{0.020833in}{0.000000in}}{%
\pgfpathmoveto{\pgfqpoint{0.000000in}{0.000000in}}%
\pgfpathlineto{\pgfqpoint{0.020833in}{0.000000in}}%
\pgfusepath{stroke,fill}%
}%
\begin{pgfscope}%
\pgfsys@transformshift{0.650000in}{2.337777in}%
\pgfsys@useobject{currentmarker}{}%
\end{pgfscope}%
\end{pgfscope}%
\begin{pgfscope}%
\pgfsetbuttcap%
\pgfsetroundjoin%
\definecolor{currentfill}{rgb}{0.000000,0.000000,0.000000}%
\pgfsetfillcolor{currentfill}%
\pgfsetlinewidth{0.401500pt}%
\definecolor{currentstroke}{rgb}{0.000000,0.000000,0.000000}%
\pgfsetstrokecolor{currentstroke}%
\pgfsetdash{}{0pt}%
\pgfsys@defobject{currentmarker}{\pgfqpoint{-0.020833in}{0.000000in}}{\pgfqpoint{-0.000000in}{0.000000in}}{%
\pgfpathmoveto{\pgfqpoint{-0.000000in}{0.000000in}}%
\pgfpathlineto{\pgfqpoint{-0.020833in}{0.000000in}}%
\pgfusepath{stroke,fill}%
}%
\begin{pgfscope}%
\pgfsys@transformshift{3.038110in}{2.337777in}%
\pgfsys@useobject{currentmarker}{}%
\end{pgfscope}%
\end{pgfscope}%
\begin{pgfscope}%
\definecolor{textcolor}{rgb}{0.000000,0.000000,0.000000}%
\pgfsetstrokecolor{textcolor}%
\pgfsetfillcolor{textcolor}%
\pgftext[x=0.260995in,y=1.417244in,,bottom,rotate=90.000000]{\color{textcolor}{\rmfamily\fontsize{12.000000}{14.400000}\selectfont\catcode`\^=\active\def^{\ifmmode\sp\else\^{}\fi}\catcode`\%=\active\def%{\%}$\langle u_i u_j \rangle^+$}}%
\end{pgfscope}%
\begin{pgfscope}%
\pgfpathrectangle{\pgfqpoint{0.650000in}{0.420000in}}{\pgfqpoint{2.388110in}{1.994488in}}%
\pgfusepath{clip}%
\pgfsetrectcap%
\pgfsetroundjoin%
\pgfsetlinewidth{0.602250pt}%
\definecolor{currentstroke}{rgb}{0.000000,0.000000,0.000000}%
\pgfsetstrokecolor{currentstroke}%
\pgfsetdash{}{0pt}%
\pgfpathmoveto{\pgfqpoint{0.645000in}{0.887260in}}%
\pgfpathlineto{\pgfqpoint{0.718874in}{0.890053in}}%
\pgfpathlineto{\pgfqpoint{0.855798in}{0.904023in}}%
\pgfpathlineto{\pgfqpoint{0.967671in}{0.929175in}}%
\pgfpathlineto{\pgfqpoint{1.062253in}{0.969906in}}%
\pgfpathlineto{\pgfqpoint{1.144181in}{1.030469in}}%
\pgfpathlineto{\pgfqpoint{1.216443in}{1.113748in}}%
\pgfpathlineto{\pgfqpoint{1.281080in}{1.219698in}}%
\pgfpathlineto{\pgfqpoint{1.339547in}{1.344142in}}%
\pgfpathlineto{\pgfqpoint{1.392920in}{1.478840in}}%
\pgfpathlineto{\pgfqpoint{1.487466in}{1.736213in}}%
\pgfpathlineto{\pgfqpoint{1.529776in}{1.839899in}}%
\pgfpathlineto{\pgfqpoint{1.569350in}{1.919350in}}%
\pgfpathlineto{\pgfqpoint{1.606521in}{1.973313in}}%
\pgfpathlineto{\pgfqpoint{1.641563in}{2.003248in}}%
\pgfpathlineto{\pgfqpoint{1.674706in}{2.012258in}}%
\pgfpathlineto{\pgfqpoint{1.706145in}{2.004164in}}%
\pgfpathlineto{\pgfqpoint{1.736046in}{1.982878in}}%
\pgfpathlineto{\pgfqpoint{1.764552in}{1.952034in}}%
\pgfpathlineto{\pgfqpoint{1.791786in}{1.914780in}}%
\pgfpathlineto{\pgfqpoint{1.817858in}{1.873690in}}%
\pgfpathlineto{\pgfqpoint{1.842862in}{1.830797in}}%
\pgfpathlineto{\pgfqpoint{1.889989in}{1.745477in}}%
\pgfpathlineto{\pgfqpoint{1.933732in}{1.666749in}}%
\pgfpathlineto{\pgfqpoint{1.954479in}{1.631006in}}%
\pgfpathlineto{\pgfqpoint{1.974542in}{1.597890in}}%
\pgfpathlineto{\pgfqpoint{1.993964in}{1.567372in}}%
\pgfpathlineto{\pgfqpoint{2.012785in}{1.539354in}}%
\pgfpathlineto{\pgfqpoint{2.031040in}{1.513729in}}%
\pgfpathlineto{\pgfqpoint{2.048762in}{1.490416in}}%
\pgfpathlineto{\pgfqpoint{2.065981in}{1.469247in}}%
\pgfpathlineto{\pgfqpoint{2.082724in}{1.449999in}}%
\pgfpathlineto{\pgfqpoint{2.099017in}{1.432513in}}%
\pgfpathlineto{\pgfqpoint{2.114883in}{1.416671in}}%
\pgfpathlineto{\pgfqpoint{2.130343in}{1.402341in}}%
\pgfpathlineto{\pgfqpoint{2.145418in}{1.389316in}}%
\pgfpathlineto{\pgfqpoint{2.174483in}{1.366327in}}%
\pgfpathlineto{\pgfqpoint{2.202211in}{1.346458in}}%
\pgfpathlineto{\pgfqpoint{2.228719in}{1.328810in}}%
\pgfpathlineto{\pgfqpoint{2.266409in}{1.304950in}}%
\pgfpathlineto{\pgfqpoint{2.313235in}{1.276736in}}%
\pgfpathlineto{\pgfqpoint{2.346093in}{1.256941in}}%
\pgfpathlineto{\pgfqpoint{2.377230in}{1.236648in}}%
\pgfpathlineto{\pgfqpoint{2.406815in}{1.215663in}}%
\pgfpathlineto{\pgfqpoint{2.425746in}{1.200971in}}%
\pgfpathlineto{\pgfqpoint{2.453049in}{1.178073in}}%
\pgfpathlineto{\pgfqpoint{2.479139in}{1.154172in}}%
\pgfpathlineto{\pgfqpoint{2.495908in}{1.137404in}}%
\pgfpathlineto{\pgfqpoint{2.512209in}{1.119866in}}%
\pgfpathlineto{\pgfqpoint{2.528065in}{1.101492in}}%
\pgfpathlineto{\pgfqpoint{2.543499in}{1.082238in}}%
\pgfpathlineto{\pgfqpoint{2.558532in}{1.061987in}}%
\pgfpathlineto{\pgfqpoint{2.580370in}{1.029992in}}%
\pgfpathlineto{\pgfqpoint{2.608252in}{0.985950in}}%
\pgfpathlineto{\pgfqpoint{2.628304in}{0.954209in}}%
\pgfpathlineto{\pgfqpoint{2.641289in}{0.935525in}}%
\pgfpathlineto{\pgfqpoint{2.653982in}{0.919728in}}%
\pgfpathlineto{\pgfqpoint{2.666395in}{0.907288in}}%
\pgfpathlineto{\pgfqpoint{2.672500in}{0.902310in}}%
\pgfpathlineto{\pgfqpoint{2.678539in}{0.898110in}}%
\pgfpathlineto{\pgfqpoint{2.690425in}{0.891847in}}%
\pgfpathlineto{\pgfqpoint{2.702063in}{0.887823in}}%
\pgfpathlineto{\pgfqpoint{2.713461in}{0.885317in}}%
\pgfpathlineto{\pgfqpoint{2.730129in}{0.883227in}}%
\pgfpathlineto{\pgfqpoint{2.751592in}{0.881877in}}%
\pgfpathlineto{\pgfqpoint{2.787225in}{0.880921in}}%
\pgfpathlineto{\pgfqpoint{2.860686in}{0.880403in}}%
\pgfpathlineto{\pgfqpoint{3.043110in}{0.880271in}}%
\pgfpathlineto{\pgfqpoint{3.043110in}{0.880271in}}%
\pgfusepath{stroke}%
\end{pgfscope}%
\begin{pgfscope}%
\pgfpathrectangle{\pgfqpoint{0.650000in}{0.420000in}}{\pgfqpoint{2.388110in}{1.994488in}}%
\pgfusepath{clip}%
\pgfsetrectcap%
\pgfsetroundjoin%
\pgfsetlinewidth{0.602250pt}%
\definecolor{currentstroke}{rgb}{0.000000,0.000000,0.000000}%
\pgfsetstrokecolor{currentstroke}%
\pgfsetdash{}{0pt}%
\pgfpathmoveto{\pgfqpoint{0.645000in}{0.880268in}}%
\pgfpathlineto{\pgfqpoint{1.144181in}{0.880722in}}%
\pgfpathlineto{\pgfqpoint{1.281080in}{0.882145in}}%
\pgfpathlineto{\pgfqpoint{1.339547in}{0.883559in}}%
\pgfpathlineto{\pgfqpoint{1.392920in}{0.885617in}}%
\pgfpathlineto{\pgfqpoint{1.442015in}{0.888438in}}%
\pgfpathlineto{\pgfqpoint{1.487466in}{0.892114in}}%
\pgfpathlineto{\pgfqpoint{1.529776in}{0.896703in}}%
\pgfpathlineto{\pgfqpoint{1.569350in}{0.902227in}}%
\pgfpathlineto{\pgfqpoint{1.606521in}{0.908667in}}%
\pgfpathlineto{\pgfqpoint{1.641563in}{0.915958in}}%
\pgfpathlineto{\pgfqpoint{1.674706in}{0.923998in}}%
\pgfpathlineto{\pgfqpoint{1.706145in}{0.932652in}}%
\pgfpathlineto{\pgfqpoint{1.736046in}{0.941761in}}%
\pgfpathlineto{\pgfqpoint{1.791786in}{0.960660in}}%
\pgfpathlineto{\pgfqpoint{1.866881in}{0.988348in}}%
\pgfpathlineto{\pgfqpoint{1.912253in}{1.004929in}}%
\pgfpathlineto{\pgfqpoint{1.954479in}{1.019326in}}%
\pgfpathlineto{\pgfqpoint{1.993964in}{1.031270in}}%
\pgfpathlineto{\pgfqpoint{2.031040in}{1.040693in}}%
\pgfpathlineto{\pgfqpoint{2.065981in}{1.047708in}}%
\pgfpathlineto{\pgfqpoint{2.099017in}{1.052554in}}%
\pgfpathlineto{\pgfqpoint{2.130343in}{1.055557in}}%
\pgfpathlineto{\pgfqpoint{2.160125in}{1.057075in}}%
\pgfpathlineto{\pgfqpoint{2.188507in}{1.057385in}}%
\pgfpathlineto{\pgfqpoint{2.215611in}{1.056724in}}%
\pgfpathlineto{\pgfqpoint{2.241547in}{1.055237in}}%
\pgfpathlineto{\pgfqpoint{2.278463in}{1.051757in}}%
\pgfpathlineto{\pgfqpoint{2.313235in}{1.046980in}}%
\pgfpathlineto{\pgfqpoint{2.346093in}{1.040999in}}%
\pgfpathlineto{\pgfqpoint{2.377230in}{1.033936in}}%
\pgfpathlineto{\pgfqpoint{2.406815in}{1.025935in}}%
\pgfpathlineto{\pgfqpoint{2.434988in}{1.017223in}}%
\pgfpathlineto{\pgfqpoint{2.461875in}{1.007642in}}%
\pgfpathlineto{\pgfqpoint{2.487584in}{0.997087in}}%
\pgfpathlineto{\pgfqpoint{2.512209in}{0.985755in}}%
\pgfpathlineto{\pgfqpoint{2.535833in}{0.973884in}}%
\pgfpathlineto{\pgfqpoint{2.573183in}{0.953521in}}%
\pgfpathlineto{\pgfqpoint{2.647671in}{0.912094in}}%
\pgfpathlineto{\pgfqpoint{2.666395in}{0.903079in}}%
\pgfpathlineto{\pgfqpoint{2.684514in}{0.895773in}}%
\pgfpathlineto{\pgfqpoint{2.702063in}{0.890299in}}%
\pgfpathlineto{\pgfqpoint{2.719073in}{0.886530in}}%
\pgfpathlineto{\pgfqpoint{2.735574in}{0.884122in}}%
\pgfpathlineto{\pgfqpoint{2.756829in}{0.882343in}}%
\pgfpathlineto{\pgfqpoint{2.787225in}{0.881200in}}%
\pgfpathlineto{\pgfqpoint{2.843265in}{0.880534in}}%
\pgfpathlineto{\pgfqpoint{2.994981in}{0.880281in}}%
\pgfpathlineto{\pgfqpoint{3.043110in}{0.880273in}}%
\pgfpathlineto{\pgfqpoint{3.043110in}{0.880273in}}%
\pgfusepath{stroke}%
\end{pgfscope}%
\begin{pgfscope}%
\pgfpathrectangle{\pgfqpoint{0.650000in}{0.420000in}}{\pgfqpoint{2.388110in}{1.994488in}}%
\pgfusepath{clip}%
\pgfsetrectcap%
\pgfsetroundjoin%
\pgfsetlinewidth{0.602250pt}%
\definecolor{currentstroke}{rgb}{0.000000,0.000000,0.000000}%
\pgfsetstrokecolor{currentstroke}%
\pgfsetdash{}{0pt}%
\pgfpathmoveto{\pgfqpoint{0.645000in}{0.882948in}}%
\pgfpathlineto{\pgfqpoint{0.718874in}{0.883976in}}%
\pgfpathlineto{\pgfqpoint{0.855798in}{0.888593in}}%
\pgfpathlineto{\pgfqpoint{0.967671in}{0.895884in}}%
\pgfpathlineto{\pgfqpoint{1.062253in}{0.906041in}}%
\pgfpathlineto{\pgfqpoint{1.144181in}{0.918894in}}%
\pgfpathlineto{\pgfqpoint{1.216443in}{0.933963in}}%
\pgfpathlineto{\pgfqpoint{1.281080in}{0.950589in}}%
\pgfpathlineto{\pgfqpoint{1.339547in}{0.968093in}}%
\pgfpathlineto{\pgfqpoint{1.392920in}{0.985906in}}%
\pgfpathlineto{\pgfqpoint{1.442015in}{1.003628in}}%
\pgfpathlineto{\pgfqpoint{1.487466in}{1.021010in}}%
\pgfpathlineto{\pgfqpoint{1.569350in}{1.054144in}}%
\pgfpathlineto{\pgfqpoint{1.674706in}{1.097606in}}%
\pgfpathlineto{\pgfqpoint{1.706145in}{1.109810in}}%
\pgfpathlineto{\pgfqpoint{1.736046in}{1.120682in}}%
\pgfpathlineto{\pgfqpoint{1.764552in}{1.130184in}}%
\pgfpathlineto{\pgfqpoint{1.791786in}{1.138330in}}%
\pgfpathlineto{\pgfqpoint{1.817858in}{1.145163in}}%
\pgfpathlineto{\pgfqpoint{1.842862in}{1.150757in}}%
\pgfpathlineto{\pgfqpoint{1.866881in}{1.155233in}}%
\pgfpathlineto{\pgfqpoint{1.889989in}{1.158738in}}%
\pgfpathlineto{\pgfqpoint{1.912253in}{1.161434in}}%
\pgfpathlineto{\pgfqpoint{1.954479in}{1.164907in}}%
\pgfpathlineto{\pgfqpoint{1.993964in}{1.166422in}}%
\pgfpathlineto{\pgfqpoint{2.031040in}{1.166489in}}%
\pgfpathlineto{\pgfqpoint{2.065981in}{1.165527in}}%
\pgfpathlineto{\pgfqpoint{2.099017in}{1.163666in}}%
\pgfpathlineto{\pgfqpoint{2.130343in}{1.160805in}}%
\pgfpathlineto{\pgfqpoint{2.160125in}{1.157004in}}%
\pgfpathlineto{\pgfqpoint{2.202211in}{1.150278in}}%
\pgfpathlineto{\pgfqpoint{2.241547in}{1.142380in}}%
\pgfpathlineto{\pgfqpoint{2.290280in}{1.131095in}}%
\pgfpathlineto{\pgfqpoint{2.324390in}{1.121909in}}%
\pgfpathlineto{\pgfqpoint{2.356654in}{1.111540in}}%
\pgfpathlineto{\pgfqpoint{2.387257in}{1.100176in}}%
\pgfpathlineto{\pgfqpoint{2.406815in}{1.091943in}}%
\pgfpathlineto{\pgfqpoint{2.425746in}{1.082807in}}%
\pgfpathlineto{\pgfqpoint{2.444088in}{1.072641in}}%
\pgfpathlineto{\pgfqpoint{2.461875in}{1.061503in}}%
\pgfpathlineto{\pgfqpoint{2.487584in}{1.043456in}}%
\pgfpathlineto{\pgfqpoint{2.504115in}{1.030859in}}%
\pgfpathlineto{\pgfqpoint{2.520191in}{1.017453in}}%
\pgfpathlineto{\pgfqpoint{2.543499in}{0.996454in}}%
\pgfpathlineto{\pgfqpoint{2.573183in}{0.968075in}}%
\pgfpathlineto{\pgfqpoint{2.601408in}{0.941608in}}%
\pgfpathlineto{\pgfqpoint{2.621698in}{0.924188in}}%
\pgfpathlineto{\pgfqpoint{2.634834in}{0.914278in}}%
\pgfpathlineto{\pgfqpoint{2.647671in}{0.905766in}}%
\pgfpathlineto{\pgfqpoint{2.660223in}{0.898821in}}%
\pgfpathlineto{\pgfqpoint{2.672500in}{0.893367in}}%
\pgfpathlineto{\pgfqpoint{2.684514in}{0.889274in}}%
\pgfpathlineto{\pgfqpoint{2.696274in}{0.886315in}}%
\pgfpathlineto{\pgfqpoint{2.713461in}{0.883503in}}%
\pgfpathlineto{\pgfqpoint{2.730129in}{0.881994in}}%
\pgfpathlineto{\pgfqpoint{2.756829in}{0.880934in}}%
\pgfpathlineto{\pgfqpoint{2.806572in}{0.880435in}}%
\pgfpathlineto{\pgfqpoint{2.994981in}{0.880270in}}%
\pgfpathlineto{\pgfqpoint{3.043110in}{0.880268in}}%
\pgfpathlineto{\pgfqpoint{3.043110in}{0.880268in}}%
\pgfusepath{stroke}%
\end{pgfscope}%
\begin{pgfscope}%
\pgfpathrectangle{\pgfqpoint{0.650000in}{0.420000in}}{\pgfqpoint{2.388110in}{1.994488in}}%
\pgfusepath{clip}%
\pgfsetrectcap%
\pgfsetroundjoin%
\pgfsetlinewidth{0.602250pt}%
\definecolor{currentstroke}{rgb}{0.000000,0.000000,0.000000}%
\pgfsetstrokecolor{currentstroke}%
\pgfsetdash{}{0pt}%
\pgfpathmoveto{\pgfqpoint{0.645000in}{0.880239in}}%
\pgfpathlineto{\pgfqpoint{0.967671in}{0.879770in}}%
\pgfpathlineto{\pgfqpoint{1.062253in}{0.879004in}}%
\pgfpathlineto{\pgfqpoint{1.144181in}{0.877475in}}%
\pgfpathlineto{\pgfqpoint{1.216443in}{0.874768in}}%
\pgfpathlineto{\pgfqpoint{1.281080in}{0.870470in}}%
\pgfpathlineto{\pgfqpoint{1.339547in}{0.864289in}}%
\pgfpathlineto{\pgfqpoint{1.392920in}{0.856176in}}%
\pgfpathlineto{\pgfqpoint{1.442015in}{0.846370in}}%
\pgfpathlineto{\pgfqpoint{1.487466in}{0.835360in}}%
\pgfpathlineto{\pgfqpoint{1.529776in}{0.823775in}}%
\pgfpathlineto{\pgfqpoint{1.641563in}{0.791298in}}%
\pgfpathlineto{\pgfqpoint{1.674706in}{0.782437in}}%
\pgfpathlineto{\pgfqpoint{1.706145in}{0.774774in}}%
\pgfpathlineto{\pgfqpoint{1.736046in}{0.768267in}}%
\pgfpathlineto{\pgfqpoint{1.764552in}{0.762821in}}%
\pgfpathlineto{\pgfqpoint{1.791786in}{0.758313in}}%
\pgfpathlineto{\pgfqpoint{1.842862in}{0.751607in}}%
\pgfpathlineto{\pgfqpoint{1.889989in}{0.747212in}}%
\pgfpathlineto{\pgfqpoint{1.933732in}{0.744424in}}%
\pgfpathlineto{\pgfqpoint{1.974542in}{0.742806in}}%
\pgfpathlineto{\pgfqpoint{2.031040in}{0.741951in}}%
\pgfpathlineto{\pgfqpoint{2.082724in}{0.742325in}}%
\pgfpathlineto{\pgfqpoint{2.145418in}{0.743922in}}%
\pgfpathlineto{\pgfqpoint{2.202211in}{0.746521in}}%
\pgfpathlineto{\pgfqpoint{2.254107in}{0.750205in}}%
\pgfpathlineto{\pgfqpoint{2.301868in}{0.754989in}}%
\pgfpathlineto{\pgfqpoint{2.335340in}{0.759318in}}%
\pgfpathlineto{\pgfqpoint{2.367031in}{0.764312in}}%
\pgfpathlineto{\pgfqpoint{2.397117in}{0.770057in}}%
\pgfpathlineto{\pgfqpoint{2.425746in}{0.776578in}}%
\pgfpathlineto{\pgfqpoint{2.453049in}{0.783894in}}%
\pgfpathlineto{\pgfqpoint{2.479139in}{0.791904in}}%
\pgfpathlineto{\pgfqpoint{2.504115in}{0.800669in}}%
\pgfpathlineto{\pgfqpoint{2.528065in}{0.810209in}}%
\pgfpathlineto{\pgfqpoint{2.551064in}{0.820336in}}%
\pgfpathlineto{\pgfqpoint{2.587469in}{0.837988in}}%
\pgfpathlineto{\pgfqpoint{2.628304in}{0.857672in}}%
\pgfpathlineto{\pgfqpoint{2.647671in}{0.865689in}}%
\pgfpathlineto{\pgfqpoint{2.660223in}{0.869999in}}%
\pgfpathlineto{\pgfqpoint{2.672500in}{0.873377in}}%
\pgfpathlineto{\pgfqpoint{2.690425in}{0.876853in}}%
\pgfpathlineto{\pgfqpoint{2.707791in}{0.878772in}}%
\pgfpathlineto{\pgfqpoint{2.730129in}{0.879850in}}%
\pgfpathlineto{\pgfqpoint{2.767152in}{0.880231in}}%
\pgfpathlineto{\pgfqpoint{3.043110in}{0.880266in}}%
\pgfpathlineto{\pgfqpoint{3.043110in}{0.880266in}}%
\pgfusepath{stroke}%
\end{pgfscope}%
\begin{pgfscope}%
\pgfpathrectangle{\pgfqpoint{0.650000in}{0.420000in}}{\pgfqpoint{2.388110in}{1.994488in}}%
\pgfusepath{clip}%
\pgfsetrectcap%
\pgfsetroundjoin%
\pgfsetlinewidth{1.003750pt}%
\definecolor{currentstroke}{rgb}{0.945098,0.713725,0.854902}%
\pgfsetstrokecolor{currentstroke}%
\pgfsetdash{}{0pt}%
\pgfpathmoveto{\pgfqpoint{0.645000in}{0.887406in}}%
\pgfpathlineto{\pgfqpoint{0.686665in}{0.889392in}}%
\pgfpathlineto{\pgfqpoint{0.741886in}{0.893016in}}%
\pgfpathlineto{\pgfqpoint{0.792190in}{0.897640in}}%
\pgfpathlineto{\pgfqpoint{0.838669in}{0.903467in}}%
\pgfpathlineto{\pgfqpoint{0.882094in}{0.910731in}}%
\pgfpathlineto{\pgfqpoint{0.902845in}{0.914984in}}%
\pgfpathlineto{\pgfqpoint{0.923033in}{0.919699in}}%
\pgfpathlineto{\pgfqpoint{0.942709in}{0.924915in}}%
\pgfpathlineto{\pgfqpoint{0.961915in}{0.930674in}}%
\pgfpathlineto{\pgfqpoint{0.980691in}{0.937017in}}%
\pgfpathlineto{\pgfqpoint{0.999071in}{0.943994in}}%
\pgfpathlineto{\pgfqpoint{1.017086in}{0.951652in}}%
\pgfpathlineto{\pgfqpoint{1.034762in}{0.960046in}}%
\pgfpathlineto{\pgfqpoint{1.052126in}{0.969232in}}%
\pgfpathlineto{\pgfqpoint{1.069199in}{0.979272in}}%
\pgfpathlineto{\pgfqpoint{1.086001in}{0.990229in}}%
\pgfpathlineto{\pgfqpoint{1.102550in}{1.002171in}}%
\pgfpathlineto{\pgfqpoint{1.118865in}{1.015166in}}%
\pgfpathlineto{\pgfqpoint{1.134959in}{1.029280in}}%
\pgfpathlineto{\pgfqpoint{1.150847in}{1.044575in}}%
\pgfpathlineto{\pgfqpoint{1.166541in}{1.061106in}}%
\pgfpathlineto{\pgfqpoint{1.182055in}{1.078918in}}%
\pgfpathlineto{\pgfqpoint{1.197397in}{1.098057in}}%
\pgfpathlineto{\pgfqpoint{1.212579in}{1.118583in}}%
\pgfpathlineto{\pgfqpoint{1.227610in}{1.140554in}}%
\pgfpathlineto{\pgfqpoint{1.242498in}{1.163997in}}%
\pgfpathlineto{\pgfqpoint{1.257252in}{1.188912in}}%
\pgfpathlineto{\pgfqpoint{1.271878in}{1.215308in}}%
\pgfpathlineto{\pgfqpoint{1.286384in}{1.243193in}}%
\pgfpathlineto{\pgfqpoint{1.300777in}{1.272551in}}%
\pgfpathlineto{\pgfqpoint{1.315062in}{1.303332in}}%
\pgfpathlineto{\pgfqpoint{1.329244in}{1.335460in}}%
\pgfpathlineto{\pgfqpoint{1.357325in}{1.403348in}}%
\pgfpathlineto{\pgfqpoint{1.385056in}{1.475234in}}%
\pgfpathlineto{\pgfqpoint{1.412473in}{1.549769in}}%
\pgfpathlineto{\pgfqpoint{1.479826in}{1.735513in}}%
\pgfpathlineto{\pgfqpoint{1.506356in}{1.804192in}}%
\pgfpathlineto{\pgfqpoint{1.519544in}{1.836312in}}%
\pgfpathlineto{\pgfqpoint{1.532684in}{1.866614in}}%
\pgfpathlineto{\pgfqpoint{1.545777in}{1.894890in}}%
\pgfpathlineto{\pgfqpoint{1.558826in}{1.920940in}}%
\pgfpathlineto{\pgfqpoint{1.571832in}{1.944602in}}%
\pgfpathlineto{\pgfqpoint{1.584799in}{1.965717in}}%
\pgfpathlineto{\pgfqpoint{1.597727in}{1.984167in}}%
\pgfpathlineto{\pgfqpoint{1.610618in}{1.999869in}}%
\pgfpathlineto{\pgfqpoint{1.623474in}{2.012782in}}%
\pgfpathlineto{\pgfqpoint{1.636296in}{2.022877in}}%
\pgfpathlineto{\pgfqpoint{1.649086in}{2.030169in}}%
\pgfpathlineto{\pgfqpoint{1.661845in}{2.034678in}}%
\pgfpathlineto{\pgfqpoint{1.674574in}{2.036469in}}%
\pgfpathlineto{\pgfqpoint{1.687276in}{2.035614in}}%
\pgfpathlineto{\pgfqpoint{1.699950in}{2.032227in}}%
\pgfpathlineto{\pgfqpoint{1.712599in}{2.026436in}}%
\pgfpathlineto{\pgfqpoint{1.725222in}{2.018376in}}%
\pgfpathlineto{\pgfqpoint{1.737822in}{2.008213in}}%
\pgfpathlineto{\pgfqpoint{1.750399in}{1.996106in}}%
\pgfpathlineto{\pgfqpoint{1.762955in}{1.982231in}}%
\pgfpathlineto{\pgfqpoint{1.775489in}{1.966766in}}%
\pgfpathlineto{\pgfqpoint{1.788004in}{1.949891in}}%
\pgfpathlineto{\pgfqpoint{1.800499in}{1.931775in}}%
\pgfpathlineto{\pgfqpoint{1.825435in}{1.892530in}}%
\pgfpathlineto{\pgfqpoint{1.850304in}{1.850384in}}%
\pgfpathlineto{\pgfqpoint{1.949203in}{1.677769in}}%
\pgfpathlineto{\pgfqpoint{1.973806in}{1.638359in}}%
\pgfpathlineto{\pgfqpoint{1.998368in}{1.601664in}}%
\pgfpathlineto{\pgfqpoint{2.022892in}{1.567742in}}%
\pgfpathlineto{\pgfqpoint{2.047381in}{1.536424in}}%
\pgfpathlineto{\pgfqpoint{2.071837in}{1.507744in}}%
\pgfpathlineto{\pgfqpoint{2.096264in}{1.481793in}}%
\pgfpathlineto{\pgfqpoint{2.120663in}{1.458286in}}%
\pgfpathlineto{\pgfqpoint{2.145037in}{1.436871in}}%
\pgfpathlineto{\pgfqpoint{2.169388in}{1.417285in}}%
\pgfpathlineto{\pgfqpoint{2.193717in}{1.399453in}}%
\pgfpathlineto{\pgfqpoint{2.218027in}{1.383450in}}%
\pgfpathlineto{\pgfqpoint{2.254457in}{1.361666in}}%
\pgfpathlineto{\pgfqpoint{2.302973in}{1.333193in}}%
\pgfpathlineto{\pgfqpoint{2.327210in}{1.317775in}}%
\pgfpathlineto{\pgfqpoint{2.434639in}{1.244930in}}%
\pgfpathlineto{\pgfqpoint{2.485971in}{1.207252in}}%
\pgfpathlineto{\pgfqpoint{2.513103in}{1.185753in}}%
\pgfpathlineto{\pgfqpoint{2.538029in}{1.163992in}}%
\pgfpathlineto{\pgfqpoint{2.568397in}{1.135344in}}%
\pgfpathlineto{\pgfqpoint{2.615231in}{1.090691in}}%
\pgfpathlineto{\pgfqpoint{2.633304in}{1.075678in}}%
\pgfpathlineto{\pgfqpoint{2.666538in}{1.049412in}}%
\pgfpathlineto{\pgfqpoint{2.681895in}{1.035174in}}%
\pgfpathlineto{\pgfqpoint{2.701245in}{1.015794in}}%
\pgfpathlineto{\pgfqpoint{2.714997in}{1.000377in}}%
\pgfpathlineto{\pgfqpoint{2.732423in}{0.978562in}}%
\pgfpathlineto{\pgfqpoint{2.748912in}{0.958149in}}%
\pgfpathlineto{\pgfqpoint{2.786631in}{0.913846in}}%
\pgfpathlineto{\pgfqpoint{2.793649in}{0.907752in}}%
\pgfpathlineto{\pgfqpoint{2.803883in}{0.900718in}}%
\pgfpathlineto{\pgfqpoint{2.813787in}{0.895648in}}%
\pgfpathlineto{\pgfqpoint{2.823382in}{0.891821in}}%
\pgfpathlineto{\pgfqpoint{2.835725in}{0.888185in}}%
\pgfpathlineto{\pgfqpoint{2.850486in}{0.885276in}}%
\pgfpathlineto{\pgfqpoint{2.867311in}{0.883301in}}%
\pgfpathlineto{\pgfqpoint{2.888398in}{0.882015in}}%
\pgfpathlineto{\pgfqpoint{2.928908in}{0.881031in}}%
\pgfpathlineto{\pgfqpoint{3.043110in}{0.880476in}}%
\pgfpathlineto{\pgfqpoint{3.043110in}{0.880476in}}%
\pgfusepath{stroke}%
\end{pgfscope}%
\begin{pgfscope}%
\pgfpathrectangle{\pgfqpoint{0.650000in}{0.420000in}}{\pgfqpoint{2.388110in}{1.994488in}}%
\pgfusepath{clip}%
\pgfsetrectcap%
\pgfsetroundjoin%
\pgfsetlinewidth{1.003750pt}%
\definecolor{currentstroke}{rgb}{0.870588,0.466667,0.682353}%
\pgfsetstrokecolor{currentstroke}%
\pgfsetdash{}{0pt}%
\pgfpathmoveto{\pgfqpoint{0.645000in}{0.880268in}}%
\pgfpathlineto{\pgfqpoint{1.150847in}{0.880784in}}%
\pgfpathlineto{\pgfqpoint{1.271878in}{0.882051in}}%
\pgfpathlineto{\pgfqpoint{1.343330in}{0.883783in}}%
\pgfpathlineto{\pgfqpoint{1.398802in}{0.886040in}}%
\pgfpathlineto{\pgfqpoint{1.453072in}{0.889389in}}%
\pgfpathlineto{\pgfqpoint{1.493118in}{0.892818in}}%
\pgfpathlineto{\pgfqpoint{1.532684in}{0.897207in}}%
\pgfpathlineto{\pgfqpoint{1.571832in}{0.902699in}}%
\pgfpathlineto{\pgfqpoint{1.610618in}{0.909428in}}%
\pgfpathlineto{\pgfqpoint{1.649086in}{0.917481in}}%
\pgfpathlineto{\pgfqpoint{1.687276in}{0.926900in}}%
\pgfpathlineto{\pgfqpoint{1.725222in}{0.937628in}}%
\pgfpathlineto{\pgfqpoint{1.762955in}{0.949531in}}%
\pgfpathlineto{\pgfqpoint{1.812976in}{0.966802in}}%
\pgfpathlineto{\pgfqpoint{1.949203in}{1.015382in}}%
\pgfpathlineto{\pgfqpoint{1.986092in}{1.026650in}}%
\pgfpathlineto{\pgfqpoint{2.022892in}{1.036309in}}%
\pgfpathlineto{\pgfqpoint{2.047381in}{1.041779in}}%
\pgfpathlineto{\pgfqpoint{2.071837in}{1.046433in}}%
\pgfpathlineto{\pgfqpoint{2.096264in}{1.050211in}}%
\pgfpathlineto{\pgfqpoint{2.120663in}{1.053084in}}%
\pgfpathlineto{\pgfqpoint{2.145037in}{1.055050in}}%
\pgfpathlineto{\pgfqpoint{2.169388in}{1.056139in}}%
\pgfpathlineto{\pgfqpoint{2.193717in}{1.056369in}}%
\pgfpathlineto{\pgfqpoint{2.218027in}{1.055766in}}%
\pgfpathlineto{\pgfqpoint{2.254457in}{1.053504in}}%
\pgfpathlineto{\pgfqpoint{2.290850in}{1.049753in}}%
\pgfpathlineto{\pgfqpoint{2.339324in}{1.043239in}}%
\pgfpathlineto{\pgfqpoint{2.399754in}{1.033977in}}%
\pgfpathlineto{\pgfqpoint{2.434639in}{1.027888in}}%
\pgfpathlineto{\pgfqpoint{2.466451in}{1.021326in}}%
\pgfpathlineto{\pgfqpoint{2.495284in}{1.014309in}}%
\pgfpathlineto{\pgfqpoint{2.529943in}{1.004498in}}%
\pgfpathlineto{\pgfqpoint{2.596028in}{0.984135in}}%
\pgfpathlineto{\pgfqpoint{2.627397in}{0.973205in}}%
\pgfpathlineto{\pgfqpoint{2.666538in}{0.958240in}}%
\pgfpathlineto{\pgfqpoint{2.696521in}{0.945594in}}%
\pgfpathlineto{\pgfqpoint{2.797099in}{0.901617in}}%
\pgfpathlineto{\pgfqpoint{2.810521in}{0.896670in}}%
\pgfpathlineto{\pgfqpoint{2.826514in}{0.892018in}}%
\pgfpathlineto{\pgfqpoint{2.844667in}{0.887959in}}%
\pgfpathlineto{\pgfqpoint{2.861804in}{0.885272in}}%
\pgfpathlineto{\pgfqpoint{2.880659in}{0.883495in}}%
\pgfpathlineto{\pgfqpoint{2.908129in}{0.882109in}}%
\pgfpathlineto{\pgfqpoint{2.948368in}{0.881339in}}%
\pgfpathlineto{\pgfqpoint{3.043110in}{0.880941in}}%
\pgfpathlineto{\pgfqpoint{3.043110in}{0.880941in}}%
\pgfusepath{stroke}%
\end{pgfscope}%
\begin{pgfscope}%
\pgfpathrectangle{\pgfqpoint{0.650000in}{0.420000in}}{\pgfqpoint{2.388110in}{1.994488in}}%
\pgfusepath{clip}%
\pgfsetrectcap%
\pgfsetroundjoin%
\pgfsetlinewidth{1.003750pt}%
\definecolor{currentstroke}{rgb}{0.772549,0.105882,0.490196}%
\pgfsetstrokecolor{currentstroke}%
\pgfsetdash{}{0pt}%
\pgfpathmoveto{\pgfqpoint{0.645000in}{0.882675in}}%
\pgfpathlineto{\pgfqpoint{0.741886in}{0.884562in}}%
\pgfpathlineto{\pgfqpoint{0.815854in}{0.886907in}}%
\pgfpathlineto{\pgfqpoint{0.882094in}{0.889986in}}%
\pgfpathlineto{\pgfqpoint{0.942709in}{0.893912in}}%
\pgfpathlineto{\pgfqpoint{0.999071in}{0.898787in}}%
\pgfpathlineto{\pgfqpoint{1.052126in}{0.904706in}}%
\pgfpathlineto{\pgfqpoint{1.102550in}{0.911742in}}%
\pgfpathlineto{\pgfqpoint{1.150847in}{0.919944in}}%
\pgfpathlineto{\pgfqpoint{1.197397in}{0.929313in}}%
\pgfpathlineto{\pgfqpoint{1.242498in}{0.939801in}}%
\pgfpathlineto{\pgfqpoint{1.286384in}{0.951331in}}%
\pgfpathlineto{\pgfqpoint{1.329244in}{0.963801in}}%
\pgfpathlineto{\pgfqpoint{1.371232in}{0.977096in}}%
\pgfpathlineto{\pgfqpoint{1.412473in}{0.991094in}}%
\pgfpathlineto{\pgfqpoint{1.466478in}{1.010698in}}%
\pgfpathlineto{\pgfqpoint{1.519544in}{1.031215in}}%
\pgfpathlineto{\pgfqpoint{1.571832in}{1.052517in}}%
\pgfpathlineto{\pgfqpoint{1.737822in}{1.121442in}}%
\pgfpathlineto{\pgfqpoint{1.762955in}{1.130443in}}%
\pgfpathlineto{\pgfqpoint{1.788004in}{1.138543in}}%
\pgfpathlineto{\pgfqpoint{1.812976in}{1.145678in}}%
\pgfpathlineto{\pgfqpoint{1.837878in}{1.151880in}}%
\pgfpathlineto{\pgfqpoint{1.875110in}{1.159662in}}%
\pgfpathlineto{\pgfqpoint{1.912213in}{1.165803in}}%
\pgfpathlineto{\pgfqpoint{1.949203in}{1.170252in}}%
\pgfpathlineto{\pgfqpoint{1.986092in}{1.173130in}}%
\pgfpathlineto{\pgfqpoint{2.010634in}{1.174178in}}%
\pgfpathlineto{\pgfqpoint{2.035140in}{1.174386in}}%
\pgfpathlineto{\pgfqpoint{2.059613in}{1.173650in}}%
\pgfpathlineto{\pgfqpoint{2.096264in}{1.171024in}}%
\pgfpathlineto{\pgfqpoint{2.157216in}{1.164936in}}%
\pgfpathlineto{\pgfqpoint{2.193717in}{1.160419in}}%
\pgfpathlineto{\pgfqpoint{2.218027in}{1.156588in}}%
\pgfpathlineto{\pgfqpoint{2.242318in}{1.151956in}}%
\pgfpathlineto{\pgfqpoint{2.278722in}{1.143590in}}%
\pgfpathlineto{\pgfqpoint{2.315094in}{1.134673in}}%
\pgfpathlineto{\pgfqpoint{2.363542in}{1.123619in}}%
\pgfpathlineto{\pgfqpoint{2.387726in}{1.116773in}}%
\pgfpathlineto{\pgfqpoint{2.434639in}{1.101636in}}%
\pgfpathlineto{\pgfqpoint{2.504323in}{1.077944in}}%
\pgfpathlineto{\pgfqpoint{2.521638in}{1.070517in}}%
\pgfpathlineto{\pgfqpoint{2.538029in}{1.062370in}}%
\pgfpathlineto{\pgfqpoint{2.561082in}{1.049209in}}%
\pgfpathlineto{\pgfqpoint{2.596028in}{1.028658in}}%
\pgfpathlineto{\pgfqpoint{2.621374in}{1.014056in}}%
\pgfpathlineto{\pgfqpoint{2.650370in}{0.995305in}}%
\pgfpathlineto{\pgfqpoint{2.686848in}{0.970239in}}%
\pgfpathlineto{\pgfqpoint{2.756836in}{0.918027in}}%
\pgfpathlineto{\pgfqpoint{2.768351in}{0.910136in}}%
\pgfpathlineto{\pgfqpoint{2.779449in}{0.903685in}}%
\pgfpathlineto{\pgfqpoint{2.793649in}{0.896824in}}%
\pgfpathlineto{\pgfqpoint{2.807220in}{0.891516in}}%
\pgfpathlineto{\pgfqpoint{2.820217in}{0.887739in}}%
\pgfpathlineto{\pgfqpoint{2.835725in}{0.884605in}}%
\pgfpathlineto{\pgfqpoint{2.850486in}{0.882601in}}%
\pgfpathlineto{\pgfqpoint{2.867311in}{0.881476in}}%
\pgfpathlineto{\pgfqpoint{2.903314in}{0.880549in}}%
\pgfpathlineto{\pgfqpoint{2.970590in}{0.880292in}}%
\pgfpathlineto{\pgfqpoint{3.043110in}{0.880269in}}%
\pgfpathlineto{\pgfqpoint{3.043110in}{0.880269in}}%
\pgfusepath{stroke}%
\end{pgfscope}%
\begin{pgfscope}%
\pgfpathrectangle{\pgfqpoint{0.650000in}{0.420000in}}{\pgfqpoint{2.388110in}{1.994488in}}%
\pgfusepath{clip}%
\pgfsetrectcap%
\pgfsetroundjoin%
\pgfsetlinewidth{1.003750pt}%
\definecolor{currentstroke}{rgb}{0.556863,0.003922,0.321569}%
\pgfsetstrokecolor{currentstroke}%
\pgfsetdash{}{0pt}%
\pgfpathmoveto{\pgfqpoint{0.645000in}{0.880245in}}%
\pgfpathlineto{\pgfqpoint{0.980691in}{0.879697in}}%
\pgfpathlineto{\pgfqpoint{1.086001in}{0.878670in}}%
\pgfpathlineto{\pgfqpoint{1.166541in}{0.876812in}}%
\pgfpathlineto{\pgfqpoint{1.227610in}{0.874186in}}%
\pgfpathlineto{\pgfqpoint{1.271878in}{0.871262in}}%
\pgfpathlineto{\pgfqpoint{1.315062in}{0.867282in}}%
\pgfpathlineto{\pgfqpoint{1.357325in}{0.862067in}}%
\pgfpathlineto{\pgfqpoint{1.398802in}{0.855449in}}%
\pgfpathlineto{\pgfqpoint{1.439605in}{0.847423in}}%
\pgfpathlineto{\pgfqpoint{1.479826in}{0.838086in}}%
\pgfpathlineto{\pgfqpoint{1.519544in}{0.827695in}}%
\pgfpathlineto{\pgfqpoint{1.584799in}{0.809096in}}%
\pgfpathlineto{\pgfqpoint{1.649086in}{0.790890in}}%
\pgfpathlineto{\pgfqpoint{1.687276in}{0.781011in}}%
\pgfpathlineto{\pgfqpoint{1.725222in}{0.772297in}}%
\pgfpathlineto{\pgfqpoint{1.762955in}{0.764844in}}%
\pgfpathlineto{\pgfqpoint{1.800499in}{0.758678in}}%
\pgfpathlineto{\pgfqpoint{1.837878in}{0.753786in}}%
\pgfpathlineto{\pgfqpoint{1.875110in}{0.750020in}}%
\pgfpathlineto{\pgfqpoint{1.912213in}{0.747263in}}%
\pgfpathlineto{\pgfqpoint{1.961510in}{0.744791in}}%
\pgfpathlineto{\pgfqpoint{2.022892in}{0.743096in}}%
\pgfpathlineto{\pgfqpoint{2.071837in}{0.742670in}}%
\pgfpathlineto{\pgfqpoint{2.120663in}{0.743474in}}%
\pgfpathlineto{\pgfqpoint{2.157216in}{0.745124in}}%
\pgfpathlineto{\pgfqpoint{2.205874in}{0.748494in}}%
\pgfpathlineto{\pgfqpoint{2.266591in}{0.753815in}}%
\pgfpathlineto{\pgfqpoint{2.315094in}{0.759300in}}%
\pgfpathlineto{\pgfqpoint{2.351435in}{0.764441in}}%
\pgfpathlineto{\pgfqpoint{2.411648in}{0.774289in}}%
\pgfpathlineto{\pgfqpoint{2.476366in}{0.786079in}}%
\pgfpathlineto{\pgfqpoint{2.513103in}{0.793674in}}%
\pgfpathlineto{\pgfqpoint{2.545907in}{0.801589in}}%
\pgfpathlineto{\pgfqpoint{2.582523in}{0.811932in}}%
\pgfpathlineto{\pgfqpoint{2.661243in}{0.834029in}}%
\pgfpathlineto{\pgfqpoint{2.701245in}{0.845428in}}%
\pgfpathlineto{\pgfqpoint{2.728158in}{0.854299in}}%
\pgfpathlineto{\pgfqpoint{2.760723in}{0.865790in}}%
\pgfpathlineto{\pgfqpoint{2.779449in}{0.871022in}}%
\pgfpathlineto{\pgfqpoint{2.793649in}{0.874132in}}%
\pgfpathlineto{\pgfqpoint{2.810521in}{0.876561in}}%
\pgfpathlineto{\pgfqpoint{2.832685in}{0.878456in}}%
\pgfpathlineto{\pgfqpoint{2.861804in}{0.879843in}}%
\pgfpathlineto{\pgfqpoint{2.903314in}{0.880187in}}%
\pgfpathlineto{\pgfqpoint{3.043110in}{0.880225in}}%
\pgfpathlineto{\pgfqpoint{3.043110in}{0.880225in}}%
\pgfusepath{stroke}%
\end{pgfscope}%
\begin{pgfscope}%
\pgfpathrectangle{\pgfqpoint{0.650000in}{0.420000in}}{\pgfqpoint{2.388110in}{1.994488in}}%
\pgfusepath{clip}%
\pgfsetrectcap%
\pgfsetroundjoin%
\pgfsetlinewidth{1.003750pt}%
\definecolor{currentstroke}{rgb}{0.721569,0.882353,0.525490}%
\pgfsetstrokecolor{currentstroke}%
\pgfsetdash{}{0pt}%
\pgfpathmoveto{\pgfqpoint{0.645000in}{0.887416in}}%
\pgfpathlineto{\pgfqpoint{0.682092in}{0.889164in}}%
\pgfpathlineto{\pgfqpoint{0.737313in}{0.892692in}}%
\pgfpathlineto{\pgfqpoint{0.787618in}{0.897187in}}%
\pgfpathlineto{\pgfqpoint{0.834096in}{0.902839in}}%
\pgfpathlineto{\pgfqpoint{0.877521in}{0.909870in}}%
\pgfpathlineto{\pgfqpoint{0.918461in}{0.918539in}}%
\pgfpathlineto{\pgfqpoint{0.938136in}{0.923577in}}%
\pgfpathlineto{\pgfqpoint{0.957343in}{0.929139in}}%
\pgfpathlineto{\pgfqpoint{0.976119in}{0.935268in}}%
\pgfpathlineto{\pgfqpoint{0.994499in}{0.942009in}}%
\pgfpathlineto{\pgfqpoint{1.012513in}{0.949413in}}%
\pgfpathlineto{\pgfqpoint{1.030190in}{0.957531in}}%
\pgfpathlineto{\pgfqpoint{1.047553in}{0.966418in}}%
\pgfpathlineto{\pgfqpoint{1.064626in}{0.976133in}}%
\pgfpathlineto{\pgfqpoint{1.081428in}{0.986736in}}%
\pgfpathlineto{\pgfqpoint{1.097978in}{0.998290in}}%
\pgfpathlineto{\pgfqpoint{1.114292in}{1.010859in}}%
\pgfpathlineto{\pgfqpoint{1.130386in}{1.024509in}}%
\pgfpathlineto{\pgfqpoint{1.146274in}{1.039301in}}%
\pgfpathlineto{\pgfqpoint{1.161969in}{1.055294in}}%
\pgfpathlineto{\pgfqpoint{1.177482in}{1.072540in}}%
\pgfpathlineto{\pgfqpoint{1.192824in}{1.091087in}}%
\pgfpathlineto{\pgfqpoint{1.208006in}{1.110993in}}%
\pgfpathlineto{\pgfqpoint{1.223037in}{1.132310in}}%
\pgfpathlineto{\pgfqpoint{1.237925in}{1.155068in}}%
\pgfpathlineto{\pgfqpoint{1.252679in}{1.179276in}}%
\pgfpathlineto{\pgfqpoint{1.267305in}{1.204945in}}%
\pgfpathlineto{\pgfqpoint{1.281811in}{1.232082in}}%
\pgfpathlineto{\pgfqpoint{1.296204in}{1.260674in}}%
\pgfpathlineto{\pgfqpoint{1.310489in}{1.290677in}}%
\pgfpathlineto{\pgfqpoint{1.324672in}{1.322022in}}%
\pgfpathlineto{\pgfqpoint{1.352752in}{1.388360in}}%
\pgfpathlineto{\pgfqpoint{1.380483in}{1.458734in}}%
\pgfpathlineto{\pgfqpoint{1.407900in}{1.531825in}}%
\pgfpathlineto{\pgfqpoint{1.475253in}{1.714505in}}%
\pgfpathlineto{\pgfqpoint{1.501784in}{1.782211in}}%
\pgfpathlineto{\pgfqpoint{1.514971in}{1.813901in}}%
\pgfpathlineto{\pgfqpoint{1.528111in}{1.843815in}}%
\pgfpathlineto{\pgfqpoint{1.541204in}{1.871746in}}%
\pgfpathlineto{\pgfqpoint{1.554253in}{1.897504in}}%
\pgfpathlineto{\pgfqpoint{1.567260in}{1.920924in}}%
\pgfpathlineto{\pgfqpoint{1.580226in}{1.941862in}}%
\pgfpathlineto{\pgfqpoint{1.593154in}{1.960207in}}%
\pgfpathlineto{\pgfqpoint{1.606045in}{1.975889in}}%
\pgfpathlineto{\pgfqpoint{1.618901in}{1.988866in}}%
\pgfpathlineto{\pgfqpoint{1.631723in}{1.999123in}}%
\pgfpathlineto{\pgfqpoint{1.644513in}{2.006669in}}%
\pgfpathlineto{\pgfqpoint{1.657272in}{2.011535in}}%
\pgfpathlineto{\pgfqpoint{1.670002in}{2.013770in}}%
\pgfpathlineto{\pgfqpoint{1.682703in}{2.013457in}}%
\pgfpathlineto{\pgfqpoint{1.695377in}{2.010681in}}%
\pgfpathlineto{\pgfqpoint{1.708026in}{2.005541in}}%
\pgfpathlineto{\pgfqpoint{1.720650in}{1.998162in}}%
\pgfpathlineto{\pgfqpoint{1.733250in}{1.988676in}}%
\pgfpathlineto{\pgfqpoint{1.745827in}{1.977222in}}%
\pgfpathlineto{\pgfqpoint{1.758382in}{1.963951in}}%
\pgfpathlineto{\pgfqpoint{1.770917in}{1.949026in}}%
\pgfpathlineto{\pgfqpoint{1.783431in}{1.932627in}}%
\pgfpathlineto{\pgfqpoint{1.795926in}{1.914939in}}%
\pgfpathlineto{\pgfqpoint{1.820863in}{1.876458in}}%
\pgfpathlineto{\pgfqpoint{1.845731in}{1.835141in}}%
\pgfpathlineto{\pgfqpoint{1.919983in}{1.708321in}}%
\pgfpathlineto{\pgfqpoint{1.944630in}{1.668613in}}%
\pgfpathlineto{\pgfqpoint{1.969233in}{1.631219in}}%
\pgfpathlineto{\pgfqpoint{1.993795in}{1.596392in}}%
\pgfpathlineto{\pgfqpoint{2.018319in}{1.564282in}}%
\pgfpathlineto{\pgfqpoint{2.042808in}{1.534892in}}%
\pgfpathlineto{\pgfqpoint{2.067264in}{1.508126in}}%
\pgfpathlineto{\pgfqpoint{2.091691in}{1.483862in}}%
\pgfpathlineto{\pgfqpoint{2.116091in}{1.461871in}}%
\pgfpathlineto{\pgfqpoint{2.140465in}{1.441638in}}%
\pgfpathlineto{\pgfqpoint{2.164815in}{1.422655in}}%
\pgfpathlineto{\pgfqpoint{2.201302in}{1.395942in}}%
\pgfpathlineto{\pgfqpoint{2.225602in}{1.379351in}}%
\pgfpathlineto{\pgfqpoint{2.249884in}{1.364128in}}%
\pgfpathlineto{\pgfqpoint{2.274150in}{1.350738in}}%
\pgfpathlineto{\pgfqpoint{2.298401in}{1.339057in}}%
\pgfpathlineto{\pgfqpoint{2.346862in}{1.317660in}}%
\pgfpathlineto{\pgfqpoint{2.383154in}{1.302048in}}%
\pgfpathlineto{\pgfqpoint{2.407075in}{1.290331in}}%
\pgfpathlineto{\pgfqpoint{2.418733in}{1.283711in}}%
\pgfpathlineto{\pgfqpoint{2.441034in}{1.269200in}}%
\pgfpathlineto{\pgfqpoint{2.461878in}{1.253413in}}%
\pgfpathlineto{\pgfqpoint{2.481398in}{1.236783in}}%
\pgfpathlineto{\pgfqpoint{2.508530in}{1.211400in}}%
\pgfpathlineto{\pgfqpoint{2.525371in}{1.195815in}}%
\pgfpathlineto{\pgfqpoint{2.563824in}{1.162056in}}%
\pgfpathlineto{\pgfqpoint{2.577950in}{1.147393in}}%
\pgfpathlineto{\pgfqpoint{2.591455in}{1.131865in}}%
\pgfpathlineto{\pgfqpoint{2.628731in}{1.087479in}}%
\pgfpathlineto{\pgfqpoint{2.645798in}{1.069415in}}%
\pgfpathlineto{\pgfqpoint{2.691948in}{1.022669in}}%
\pgfpathlineto{\pgfqpoint{2.710424in}{1.001605in}}%
\pgfpathlineto{\pgfqpoint{2.723586in}{0.984008in}}%
\pgfpathlineto{\pgfqpoint{2.756150in}{0.938082in}}%
\pgfpathlineto{\pgfqpoint{2.767523in}{0.923891in}}%
\pgfpathlineto{\pgfqpoint{2.778489in}{0.912171in}}%
\pgfpathlineto{\pgfqpoint{2.789076in}{0.903128in}}%
\pgfpathlineto{\pgfqpoint{2.799311in}{0.896417in}}%
\pgfpathlineto{\pgfqpoint{2.809215in}{0.891588in}}%
\pgfpathlineto{\pgfqpoint{2.818809in}{0.888338in}}%
\pgfpathlineto{\pgfqpoint{2.831152in}{0.885562in}}%
\pgfpathlineto{\pgfqpoint{2.845913in}{0.883537in}}%
\pgfpathlineto{\pgfqpoint{2.862738in}{0.882415in}}%
\pgfpathlineto{\pgfqpoint{2.901158in}{0.881252in}}%
\pgfpathlineto{\pgfqpoint{2.943796in}{0.880759in}}%
\pgfpathlineto{\pgfqpoint{3.035477in}{0.880485in}}%
\pgfpathlineto{\pgfqpoint{3.043110in}{0.880482in}}%
\pgfpathlineto{\pgfqpoint{3.043110in}{0.880482in}}%
\pgfusepath{stroke}%
\end{pgfscope}%
\begin{pgfscope}%
\pgfpathrectangle{\pgfqpoint{0.650000in}{0.420000in}}{\pgfqpoint{2.388110in}{1.994488in}}%
\pgfusepath{clip}%
\pgfsetrectcap%
\pgfsetroundjoin%
\pgfsetlinewidth{1.003750pt}%
\definecolor{currentstroke}{rgb}{0.498039,0.737255,0.254902}%
\pgfsetstrokecolor{currentstroke}%
\pgfsetdash{}{0pt}%
\pgfpathmoveto{\pgfqpoint{0.645000in}{0.880268in}}%
\pgfpathlineto{\pgfqpoint{1.146274in}{0.880766in}}%
\pgfpathlineto{\pgfqpoint{1.267305in}{0.881996in}}%
\pgfpathlineto{\pgfqpoint{1.338758in}{0.883685in}}%
\pgfpathlineto{\pgfqpoint{1.394229in}{0.885898in}}%
\pgfpathlineto{\pgfqpoint{1.448499in}{0.889194in}}%
\pgfpathlineto{\pgfqpoint{1.488545in}{0.892585in}}%
\pgfpathlineto{\pgfqpoint{1.528111in}{0.896937in}}%
\pgfpathlineto{\pgfqpoint{1.567260in}{0.902407in}}%
\pgfpathlineto{\pgfqpoint{1.606045in}{0.909140in}}%
\pgfpathlineto{\pgfqpoint{1.644513in}{0.917239in}}%
\pgfpathlineto{\pgfqpoint{1.682703in}{0.926757in}}%
\pgfpathlineto{\pgfqpoint{1.720650in}{0.937658in}}%
\pgfpathlineto{\pgfqpoint{1.758382in}{0.949811in}}%
\pgfpathlineto{\pgfqpoint{1.808403in}{0.967504in}}%
\pgfpathlineto{\pgfqpoint{1.932312in}{1.012735in}}%
\pgfpathlineto{\pgfqpoint{1.969233in}{1.024543in}}%
\pgfpathlineto{\pgfqpoint{2.006062in}{1.034971in}}%
\pgfpathlineto{\pgfqpoint{2.042808in}{1.043835in}}%
\pgfpathlineto{\pgfqpoint{2.067264in}{1.048783in}}%
\pgfpathlineto{\pgfqpoint{2.091691in}{1.052942in}}%
\pgfpathlineto{\pgfqpoint{2.116091in}{1.056333in}}%
\pgfpathlineto{\pgfqpoint{2.140465in}{1.058944in}}%
\pgfpathlineto{\pgfqpoint{2.164815in}{1.060688in}}%
\pgfpathlineto{\pgfqpoint{2.189145in}{1.061465in}}%
\pgfpathlineto{\pgfqpoint{2.213454in}{1.061302in}}%
\pgfpathlineto{\pgfqpoint{2.249884in}{1.059684in}}%
\pgfpathlineto{\pgfqpoint{2.286277in}{1.056666in}}%
\pgfpathlineto{\pgfqpoint{2.322638in}{1.052424in}}%
\pgfpathlineto{\pgfqpoint{2.358969in}{1.047311in}}%
\pgfpathlineto{\pgfqpoint{2.395181in}{1.041262in}}%
\pgfpathlineto{\pgfqpoint{2.441034in}{1.032471in}}%
\pgfpathlineto{\pgfqpoint{2.481398in}{1.023480in}}%
\pgfpathlineto{\pgfqpoint{2.517066in}{1.014220in}}%
\pgfpathlineto{\pgfqpoint{2.541335in}{1.006576in}}%
\pgfpathlineto{\pgfqpoint{2.570968in}{0.995462in}}%
\pgfpathlineto{\pgfqpoint{2.616802in}{0.977495in}}%
\pgfpathlineto{\pgfqpoint{2.701325in}{0.941131in}}%
\pgfpathlineto{\pgfqpoint{2.792526in}{0.899203in}}%
\pgfpathlineto{\pgfqpoint{2.809215in}{0.893356in}}%
\pgfpathlineto{\pgfqpoint{2.825043in}{0.889128in}}%
\pgfpathlineto{\pgfqpoint{2.840094in}{0.886342in}}%
\pgfpathlineto{\pgfqpoint{2.859997in}{0.883982in}}%
\pgfpathlineto{\pgfqpoint{2.881267in}{0.882588in}}%
\pgfpathlineto{\pgfqpoint{2.917564in}{0.881632in}}%
\pgfpathlineto{\pgfqpoint{2.986738in}{0.881072in}}%
\pgfpathlineto{\pgfqpoint{3.043110in}{0.880963in}}%
\pgfpathlineto{\pgfqpoint{3.043110in}{0.880963in}}%
\pgfusepath{stroke}%
\end{pgfscope}%
\begin{pgfscope}%
\pgfpathrectangle{\pgfqpoint{0.650000in}{0.420000in}}{\pgfqpoint{2.388110in}{1.994488in}}%
\pgfusepath{clip}%
\pgfsetrectcap%
\pgfsetroundjoin%
\pgfsetlinewidth{1.003750pt}%
\definecolor{currentstroke}{rgb}{0.301961,0.572549,0.129412}%
\pgfsetstrokecolor{currentstroke}%
\pgfsetdash{}{0pt}%
\pgfpathmoveto{\pgfqpoint{0.645000in}{0.882752in}}%
\pgfpathlineto{\pgfqpoint{0.737313in}{0.884569in}}%
\pgfpathlineto{\pgfqpoint{0.811281in}{0.886908in}}%
\pgfpathlineto{\pgfqpoint{0.877521in}{0.889982in}}%
\pgfpathlineto{\pgfqpoint{0.938136in}{0.893905in}}%
\pgfpathlineto{\pgfqpoint{0.994499in}{0.898785in}}%
\pgfpathlineto{\pgfqpoint{1.047553in}{0.904721in}}%
\pgfpathlineto{\pgfqpoint{1.097978in}{0.911797in}}%
\pgfpathlineto{\pgfqpoint{1.146274in}{0.920068in}}%
\pgfpathlineto{\pgfqpoint{1.192824in}{0.929549in}}%
\pgfpathlineto{\pgfqpoint{1.237925in}{0.940204in}}%
\pgfpathlineto{\pgfqpoint{1.281811in}{0.951962in}}%
\pgfpathlineto{\pgfqpoint{1.324672in}{0.964721in}}%
\pgfpathlineto{\pgfqpoint{1.366659in}{0.978354in}}%
\pgfpathlineto{\pgfqpoint{1.407900in}{0.992721in}}%
\pgfpathlineto{\pgfqpoint{1.461906in}{1.012815in}}%
\pgfpathlineto{\pgfqpoint{1.514971in}{1.033768in}}%
\pgfpathlineto{\pgfqpoint{1.580226in}{1.060857in}}%
\pgfpathlineto{\pgfqpoint{1.695377in}{1.109325in}}%
\pgfpathlineto{\pgfqpoint{1.733250in}{1.124089in}}%
\pgfpathlineto{\pgfqpoint{1.770917in}{1.137459in}}%
\pgfpathlineto{\pgfqpoint{1.795926in}{1.145405in}}%
\pgfpathlineto{\pgfqpoint{1.820863in}{1.152466in}}%
\pgfpathlineto{\pgfqpoint{1.845731in}{1.158603in}}%
\pgfpathlineto{\pgfqpoint{1.870537in}{1.163809in}}%
\pgfpathlineto{\pgfqpoint{1.895286in}{1.168112in}}%
\pgfpathlineto{\pgfqpoint{1.919983in}{1.171584in}}%
\pgfpathlineto{\pgfqpoint{1.944630in}{1.174293in}}%
\pgfpathlineto{\pgfqpoint{1.969233in}{1.176225in}}%
\pgfpathlineto{\pgfqpoint{2.006062in}{1.177666in}}%
\pgfpathlineto{\pgfqpoint{2.042808in}{1.177881in}}%
\pgfpathlineto{\pgfqpoint{2.079481in}{1.177061in}}%
\pgfpathlineto{\pgfqpoint{2.103894in}{1.175493in}}%
\pgfpathlineto{\pgfqpoint{2.140465in}{1.171456in}}%
\pgfpathlineto{\pgfqpoint{2.176983in}{1.167403in}}%
\pgfpathlineto{\pgfqpoint{2.237745in}{1.161396in}}%
\pgfpathlineto{\pgfqpoint{2.262019in}{1.157778in}}%
\pgfpathlineto{\pgfqpoint{2.286277in}{1.153106in}}%
\pgfpathlineto{\pgfqpoint{2.322638in}{1.144562in}}%
\pgfpathlineto{\pgfqpoint{2.358969in}{1.134411in}}%
\pgfpathlineto{\pgfqpoint{2.407075in}{1.119573in}}%
\pgfpathlineto{\pgfqpoint{2.451632in}{1.104760in}}%
\pgfpathlineto{\pgfqpoint{2.481398in}{1.093529in}}%
\pgfpathlineto{\pgfqpoint{2.499750in}{1.085895in}}%
\pgfpathlineto{\pgfqpoint{2.517066in}{1.077649in}}%
\pgfpathlineto{\pgfqpoint{2.549016in}{1.061945in}}%
\pgfpathlineto{\pgfqpoint{2.570968in}{1.051550in}}%
\pgfpathlineto{\pgfqpoint{2.584777in}{1.043716in}}%
\pgfpathlineto{\pgfqpoint{2.604390in}{1.030723in}}%
\pgfpathlineto{\pgfqpoint{2.628731in}{1.012736in}}%
\pgfpathlineto{\pgfqpoint{2.661965in}{0.988213in}}%
\pgfpathlineto{\pgfqpoint{2.677323in}{0.977161in}}%
\pgfpathlineto{\pgfqpoint{2.701325in}{0.957693in}}%
\pgfpathlineto{\pgfqpoint{2.723586in}{0.938539in}}%
\pgfpathlineto{\pgfqpoint{2.744340in}{0.920857in}}%
\pgfpathlineto{\pgfqpoint{2.756150in}{0.912258in}}%
\pgfpathlineto{\pgfqpoint{2.767523in}{0.905180in}}%
\pgfpathlineto{\pgfqpoint{2.782059in}{0.897651in}}%
\pgfpathlineto{\pgfqpoint{2.795937in}{0.891623in}}%
\pgfpathlineto{\pgfqpoint{2.805949in}{0.888389in}}%
\pgfpathlineto{\pgfqpoint{2.818809in}{0.885359in}}%
\pgfpathlineto{\pgfqpoint{2.834162in}{0.883112in}}%
\pgfpathlineto{\pgfqpoint{2.854441in}{0.881558in}}%
\pgfpathlineto{\pgfqpoint{2.881267in}{0.880694in}}%
\pgfpathlineto{\pgfqpoint{2.939576in}{0.880331in}}%
\pgfpathlineto{\pgfqpoint{3.043110in}{0.880268in}}%
\pgfpathlineto{\pgfqpoint{3.043110in}{0.880268in}}%
\pgfusepath{stroke}%
\end{pgfscope}%
\begin{pgfscope}%
\pgfpathrectangle{\pgfqpoint{0.650000in}{0.420000in}}{\pgfqpoint{2.388110in}{1.994488in}}%
\pgfusepath{clip}%
\pgfsetrectcap%
\pgfsetroundjoin%
\pgfsetlinewidth{1.003750pt}%
\definecolor{currentstroke}{rgb}{0.152941,0.392157,0.098039}%
\pgfsetstrokecolor{currentstroke}%
\pgfsetdash{}{0pt}%
\pgfpathmoveto{\pgfqpoint{0.645000in}{0.880246in}}%
\pgfpathlineto{\pgfqpoint{0.976119in}{0.879726in}}%
\pgfpathlineto{\pgfqpoint{1.097978in}{0.878479in}}%
\pgfpathlineto{\pgfqpoint{1.177482in}{0.876449in}}%
\pgfpathlineto{\pgfqpoint{1.237925in}{0.873616in}}%
\pgfpathlineto{\pgfqpoint{1.281811in}{0.870481in}}%
\pgfpathlineto{\pgfqpoint{1.324672in}{0.866261in}}%
\pgfpathlineto{\pgfqpoint{1.366659in}{0.860753in}}%
\pgfpathlineto{\pgfqpoint{1.407900in}{0.853867in}}%
\pgfpathlineto{\pgfqpoint{1.448499in}{0.845587in}}%
\pgfpathlineto{\pgfqpoint{1.488545in}{0.836060in}}%
\pgfpathlineto{\pgfqpoint{1.541204in}{0.821966in}}%
\pgfpathlineto{\pgfqpoint{1.670002in}{0.786311in}}%
\pgfpathlineto{\pgfqpoint{1.708026in}{0.777263in}}%
\pgfpathlineto{\pgfqpoint{1.745827in}{0.769451in}}%
\pgfpathlineto{\pgfqpoint{1.783431in}{0.762914in}}%
\pgfpathlineto{\pgfqpoint{1.820863in}{0.757642in}}%
\pgfpathlineto{\pgfqpoint{1.858142in}{0.753506in}}%
\pgfpathlineto{\pgfqpoint{1.895286in}{0.750352in}}%
\pgfpathlineto{\pgfqpoint{1.944630in}{0.747358in}}%
\pgfpathlineto{\pgfqpoint{1.993795in}{0.745518in}}%
\pgfpathlineto{\pgfqpoint{2.055040in}{0.744383in}}%
\pgfpathlineto{\pgfqpoint{2.103894in}{0.744381in}}%
\pgfpathlineto{\pgfqpoint{2.140465in}{0.745273in}}%
\pgfpathlineto{\pgfqpoint{2.176983in}{0.747098in}}%
\pgfpathlineto{\pgfqpoint{2.225602in}{0.750682in}}%
\pgfpathlineto{\pgfqpoint{2.286277in}{0.756267in}}%
\pgfpathlineto{\pgfqpoint{2.334751in}{0.761713in}}%
\pgfpathlineto{\pgfqpoint{2.383154in}{0.768318in}}%
\pgfpathlineto{\pgfqpoint{2.430067in}{0.775564in}}%
\pgfpathlineto{\pgfqpoint{2.471793in}{0.783532in}}%
\pgfpathlineto{\pgfqpoint{2.508530in}{0.791700in}}%
\pgfpathlineto{\pgfqpoint{2.577950in}{0.809273in}}%
\pgfpathlineto{\pgfqpoint{2.610658in}{0.818984in}}%
\pgfpathlineto{\pgfqpoint{2.651282in}{0.831248in}}%
\pgfpathlineto{\pgfqpoint{2.677323in}{0.838981in}}%
\pgfpathlineto{\pgfqpoint{2.705908in}{0.849179in}}%
\pgfpathlineto{\pgfqpoint{2.756150in}{0.867147in}}%
\pgfpathlineto{\pgfqpoint{2.774877in}{0.872214in}}%
\pgfpathlineto{\pgfqpoint{2.792526in}{0.875875in}}%
\pgfpathlineto{\pgfqpoint{2.809215in}{0.878192in}}%
\pgfpathlineto{\pgfqpoint{2.828112in}{0.879440in}}%
\pgfpathlineto{\pgfqpoint{2.857232in}{0.880011in}}%
\pgfpathlineto{\pgfqpoint{2.993939in}{0.880217in}}%
\pgfpathlineto{\pgfqpoint{3.043110in}{0.880207in}}%
\pgfpathlineto{\pgfqpoint{3.043110in}{0.880207in}}%
\pgfusepath{stroke}%
\end{pgfscope}%
\begin{pgfscope}%
\pgfsetrectcap%
\pgfsetmiterjoin%
\pgfsetlinewidth{0.803000pt}%
\definecolor{currentstroke}{rgb}{0.000000,0.000000,0.000000}%
\pgfsetstrokecolor{currentstroke}%
\pgfsetdash{}{0pt}%
\pgfpathmoveto{\pgfqpoint{0.650000in}{0.420000in}}%
\pgfpathlineto{\pgfqpoint{0.650000in}{2.414488in}}%
\pgfusepath{stroke}%
\end{pgfscope}%
\begin{pgfscope}%
\pgfsetrectcap%
\pgfsetmiterjoin%
\pgfsetlinewidth{0.803000pt}%
\definecolor{currentstroke}{rgb}{0.000000,0.000000,0.000000}%
\pgfsetstrokecolor{currentstroke}%
\pgfsetdash{}{0pt}%
\pgfpathmoveto{\pgfqpoint{3.038110in}{0.420000in}}%
\pgfpathlineto{\pgfqpoint{3.038110in}{2.414488in}}%
\pgfusepath{stroke}%
\end{pgfscope}%
\begin{pgfscope}%
\pgfsetrectcap%
\pgfsetmiterjoin%
\pgfsetlinewidth{0.803000pt}%
\definecolor{currentstroke}{rgb}{0.000000,0.000000,0.000000}%
\pgfsetstrokecolor{currentstroke}%
\pgfsetdash{}{0pt}%
\pgfpathmoveto{\pgfqpoint{0.650000in}{0.420000in}}%
\pgfpathlineto{\pgfqpoint{3.038110in}{0.420000in}}%
\pgfusepath{stroke}%
\end{pgfscope}%
\begin{pgfscope}%
\pgfsetrectcap%
\pgfsetmiterjoin%
\pgfsetlinewidth{0.803000pt}%
\definecolor{currentstroke}{rgb}{0.000000,0.000000,0.000000}%
\pgfsetstrokecolor{currentstroke}%
\pgfsetdash{}{0pt}%
\pgfpathmoveto{\pgfqpoint{0.650000in}{2.414488in}}%
\pgfpathlineto{\pgfqpoint{3.038110in}{2.414488in}}%
\pgfusepath{stroke}%
\end{pgfscope}%
\begin{pgfscope}%
\pgfsetrectcap%
\pgfsetroundjoin%
\pgfsetlinewidth{0.602250pt}%
\definecolor{currentstroke}{rgb}{0.000000,0.000000,0.000000}%
\pgfsetstrokecolor{currentstroke}%
\pgfsetdash{}{0pt}%
\pgfpathmoveto{\pgfqpoint{2.127879in}{2.275599in}}%
\pgfpathlineto{\pgfqpoint{2.238990in}{2.275599in}}%
\pgfpathlineto{\pgfqpoint{2.350101in}{2.275599in}}%
\pgfusepath{stroke}%
\end{pgfscope}%
\begin{pgfscope}%
\definecolor{textcolor}{rgb}{0.000000,0.000000,0.000000}%
\pgfsetstrokecolor{textcolor}%
\pgfsetfillcolor{textcolor}%
\pgftext[x=2.438990in,y=2.236710in,left,base]{\color{textcolor}{\rmfamily\fontsize{8.000000}{9.600000}\selectfont\catcode`\^=\active\def^{\ifmmode\sp\else\^{}\fi}\catcode`\%=\active\def%{\%}ref}}%
\end{pgfscope}%
\begin{pgfscope}%
\pgfsetrectcap%
\pgfsetroundjoin%
\pgfsetlinewidth{1.003750pt}%
\definecolor{currentstroke}{rgb}{0.945098,0.713725,0.854902}%
\pgfsetstrokecolor{currentstroke}%
\pgfsetdash{}{0pt}%
\pgfpathmoveto{\pgfqpoint{2.127879in}{2.120661in}}%
\pgfpathlineto{\pgfqpoint{2.238990in}{2.120661in}}%
\pgfpathlineto{\pgfqpoint{2.350101in}{2.120661in}}%
\pgfusepath{stroke}%
\end{pgfscope}%
\begin{pgfscope}%
\definecolor{textcolor}{rgb}{0.000000,0.000000,0.000000}%
\pgfsetstrokecolor{textcolor}%
\pgfsetfillcolor{textcolor}%
\pgftext[x=2.438990in,y=2.081772in,left,base]{\color{textcolor}{\rmfamily\fontsize{8.000000}{9.600000}\selectfont\catcode`\^=\active\def^{\ifmmode\sp\else\^{}\fi}\catcode`\%=\active\def%{\%}base}}%
\end{pgfscope}%
\begin{pgfscope}%
\pgfsetrectcap%
\pgfsetroundjoin%
\pgfsetlinewidth{1.003750pt}%
\definecolor{currentstroke}{rgb}{0.721569,0.882353,0.525490}%
\pgfsetstrokecolor{currentstroke}%
\pgfsetdash{}{0pt}%
\pgfpathmoveto{\pgfqpoint{2.127879in}{1.965723in}}%
\pgfpathlineto{\pgfqpoint{2.238990in}{1.965723in}}%
\pgfpathlineto{\pgfqpoint{2.350101in}{1.965723in}}%
\pgfusepath{stroke}%
\end{pgfscope}%
\begin{pgfscope}%
\definecolor{textcolor}{rgb}{0.000000,0.000000,0.000000}%
\pgfsetstrokecolor{textcolor}%
\pgfsetfillcolor{textcolor}%
\pgftext[x=2.438990in,y=1.926834in,left,base]{\color{textcolor}{\rmfamily\fontsize{8.000000}{9.600000}\selectfont\catcode`\^=\active\def^{\ifmmode\sp\else\^{}\fi}\catcode`\%=\active\def%{\%}optimised}}%
\end{pgfscope}%
\end{pgfpicture}%
\makeatother%
\endgroup%

    \label{fig:R100deg30Re1200Fluct}
\end{subfigure}
\\[-30pt]
\begin{subfigure}[t]{0.495\textwidth}
    \centering
    \subcaption{}
    %% Creator: Matplotlib, PGF backend
%%
%% To include the figure in your LaTeX document, write
%%   \input{<filename>.pgf}
%%
%% Make sure the required packages are loaded in your preamble
%%   \usepackage{pgf}
%%
%% Also ensure that all the required font packages are loaded; for instance,
%% the lmodern package is sometimes necessary when using math font.
%%   \usepackage{lmodern}
%%
%% Figures using additional raster images can only be included by \input if
%% they are in the same directory as the main LaTeX file. For loading figures
%% from other directories you can use the `import` package
%%   \usepackage{import}
%%
%% and then include the figures with
%%   \import{<path to file>}{<filename>.pgf}
%%
%% Matplotlib used the following preamble
%%   \def\mathdefault#1{#1}
%%   \everymath=\expandafter{\the\everymath\displaystyle}
%%   \IfFileExists{scrextend.sty}{
%%     \usepackage[fontsize=11.000000pt]{scrextend}
%%   }{
%%     \renewcommand{\normalsize}{\fontsize{11.000000}{13.200000}\selectfont}
%%     \normalsize
%%   }
%%   
%%   \makeatletter\@ifpackageloaded{underscore}{}{\usepackage[strings]{underscore}}\makeatother
%%
\begingroup%
\makeatletter%
\begin{pgfpicture}%
\pgfpathrectangle{\pgfpointorigin}{\pgfqpoint{3.118110in}{2.494488in}}%
\pgfusepath{use as bounding box, clip}%
\begin{pgfscope}%
\pgfsetbuttcap%
\pgfsetmiterjoin%
\definecolor{currentfill}{rgb}{1.000000,1.000000,1.000000}%
\pgfsetfillcolor{currentfill}%
\pgfsetlinewidth{0.000000pt}%
\definecolor{currentstroke}{rgb}{1.000000,1.000000,1.000000}%
\pgfsetstrokecolor{currentstroke}%
\pgfsetdash{}{0pt}%
\pgfpathmoveto{\pgfqpoint{0.000000in}{0.000000in}}%
\pgfpathlineto{\pgfqpoint{3.118110in}{0.000000in}}%
\pgfpathlineto{\pgfqpoint{3.118110in}{2.494488in}}%
\pgfpathlineto{\pgfqpoint{0.000000in}{2.494488in}}%
\pgfpathlineto{\pgfqpoint{0.000000in}{0.000000in}}%
\pgfpathclose%
\pgfusepath{fill}%
\end{pgfscope}%
\begin{pgfscope}%
\pgfsetbuttcap%
\pgfsetmiterjoin%
\definecolor{currentfill}{rgb}{1.000000,1.000000,1.000000}%
\pgfsetfillcolor{currentfill}%
\pgfsetlinewidth{0.000000pt}%
\definecolor{currentstroke}{rgb}{0.000000,0.000000,0.000000}%
\pgfsetstrokecolor{currentstroke}%
\pgfsetstrokeopacity{0.000000}%
\pgfsetdash{}{0pt}%
\pgfpathmoveto{\pgfqpoint{0.650000in}{0.420000in}}%
\pgfpathlineto{\pgfqpoint{3.038110in}{0.420000in}}%
\pgfpathlineto{\pgfqpoint{3.038110in}{2.414488in}}%
\pgfpathlineto{\pgfqpoint{0.650000in}{2.414488in}}%
\pgfpathlineto{\pgfqpoint{0.650000in}{0.420000in}}%
\pgfpathclose%
\pgfusepath{fill}%
\end{pgfscope}%
\begin{pgfscope}%
\pgfsetbuttcap%
\pgfsetroundjoin%
\definecolor{currentfill}{rgb}{0.000000,0.000000,0.000000}%
\pgfsetfillcolor{currentfill}%
\pgfsetlinewidth{0.501875pt}%
\definecolor{currentstroke}{rgb}{0.000000,0.000000,0.000000}%
\pgfsetstrokecolor{currentstroke}%
\pgfsetdash{}{0pt}%
\pgfsys@defobject{currentmarker}{\pgfqpoint{0.000000in}{0.000000in}}{\pgfqpoint{0.000000in}{0.041667in}}{%
\pgfpathmoveto{\pgfqpoint{0.000000in}{0.000000in}}%
\pgfpathlineto{\pgfqpoint{0.000000in}{0.041667in}}%
\pgfusepath{stroke,fill}%
}%
\begin{pgfscope}%
\pgfsys@transformshift{0.862677in}{0.420000in}%
\pgfsys@useobject{currentmarker}{}%
\end{pgfscope}%
\end{pgfscope}%
\begin{pgfscope}%
\pgfsetbuttcap%
\pgfsetroundjoin%
\definecolor{currentfill}{rgb}{0.000000,0.000000,0.000000}%
\pgfsetfillcolor{currentfill}%
\pgfsetlinewidth{0.501875pt}%
\definecolor{currentstroke}{rgb}{0.000000,0.000000,0.000000}%
\pgfsetstrokecolor{currentstroke}%
\pgfsetdash{}{0pt}%
\pgfsys@defobject{currentmarker}{\pgfqpoint{0.000000in}{-0.041667in}}{\pgfqpoint{0.000000in}{0.000000in}}{%
\pgfpathmoveto{\pgfqpoint{0.000000in}{0.000000in}}%
\pgfpathlineto{\pgfqpoint{0.000000in}{-0.041667in}}%
\pgfusepath{stroke,fill}%
}%
\begin{pgfscope}%
\pgfsys@transformshift{0.862677in}{2.414488in}%
\pgfsys@useobject{currentmarker}{}%
\end{pgfscope}%
\end{pgfscope}%
\begin{pgfscope}%
\definecolor{textcolor}{rgb}{0.000000,0.000000,0.000000}%
\pgfsetstrokecolor{textcolor}%
\pgfsetfillcolor{textcolor}%
\pgftext[x=0.862677in,y=0.371389in,,top]{\color{textcolor}{\rmfamily\fontsize{11.000000}{13.200000}\selectfont\catcode`\^=\active\def^{\ifmmode\sp\else\^{}\fi}\catcode`\%=\active\def%{\%}$\mathdefault{10^{0}}$}}%
\end{pgfscope}%
\begin{pgfscope}%
\pgfsetbuttcap%
\pgfsetroundjoin%
\definecolor{currentfill}{rgb}{0.000000,0.000000,0.000000}%
\pgfsetfillcolor{currentfill}%
\pgfsetlinewidth{0.501875pt}%
\definecolor{currentstroke}{rgb}{0.000000,0.000000,0.000000}%
\pgfsetstrokecolor{currentstroke}%
\pgfsetdash{}{0pt}%
\pgfsys@defobject{currentmarker}{\pgfqpoint{0.000000in}{0.000000in}}{\pgfqpoint{0.000000in}{0.041667in}}{%
\pgfpathmoveto{\pgfqpoint{0.000000in}{0.000000in}}%
\pgfpathlineto{\pgfqpoint{0.000000in}{0.041667in}}%
\pgfusepath{stroke,fill}%
}%
\begin{pgfscope}%
\pgfsys@transformshift{1.569174in}{0.420000in}%
\pgfsys@useobject{currentmarker}{}%
\end{pgfscope}%
\end{pgfscope}%
\begin{pgfscope}%
\pgfsetbuttcap%
\pgfsetroundjoin%
\definecolor{currentfill}{rgb}{0.000000,0.000000,0.000000}%
\pgfsetfillcolor{currentfill}%
\pgfsetlinewidth{0.501875pt}%
\definecolor{currentstroke}{rgb}{0.000000,0.000000,0.000000}%
\pgfsetstrokecolor{currentstroke}%
\pgfsetdash{}{0pt}%
\pgfsys@defobject{currentmarker}{\pgfqpoint{0.000000in}{-0.041667in}}{\pgfqpoint{0.000000in}{0.000000in}}{%
\pgfpathmoveto{\pgfqpoint{0.000000in}{0.000000in}}%
\pgfpathlineto{\pgfqpoint{0.000000in}{-0.041667in}}%
\pgfusepath{stroke,fill}%
}%
\begin{pgfscope}%
\pgfsys@transformshift{1.569174in}{2.414488in}%
\pgfsys@useobject{currentmarker}{}%
\end{pgfscope}%
\end{pgfscope}%
\begin{pgfscope}%
\definecolor{textcolor}{rgb}{0.000000,0.000000,0.000000}%
\pgfsetstrokecolor{textcolor}%
\pgfsetfillcolor{textcolor}%
\pgftext[x=1.569174in,y=0.371389in,,top]{\color{textcolor}{\rmfamily\fontsize{11.000000}{13.200000}\selectfont\catcode`\^=\active\def^{\ifmmode\sp\else\^{}\fi}\catcode`\%=\active\def%{\%}$\mathdefault{10^{1}}$}}%
\end{pgfscope}%
\begin{pgfscope}%
\pgfsetbuttcap%
\pgfsetroundjoin%
\definecolor{currentfill}{rgb}{0.000000,0.000000,0.000000}%
\pgfsetfillcolor{currentfill}%
\pgfsetlinewidth{0.501875pt}%
\definecolor{currentstroke}{rgb}{0.000000,0.000000,0.000000}%
\pgfsetstrokecolor{currentstroke}%
\pgfsetdash{}{0pt}%
\pgfsys@defobject{currentmarker}{\pgfqpoint{0.000000in}{0.000000in}}{\pgfqpoint{0.000000in}{0.041667in}}{%
\pgfpathmoveto{\pgfqpoint{0.000000in}{0.000000in}}%
\pgfpathlineto{\pgfqpoint{0.000000in}{0.041667in}}%
\pgfusepath{stroke,fill}%
}%
\begin{pgfscope}%
\pgfsys@transformshift{2.275671in}{0.420000in}%
\pgfsys@useobject{currentmarker}{}%
\end{pgfscope}%
\end{pgfscope}%
\begin{pgfscope}%
\pgfsetbuttcap%
\pgfsetroundjoin%
\definecolor{currentfill}{rgb}{0.000000,0.000000,0.000000}%
\pgfsetfillcolor{currentfill}%
\pgfsetlinewidth{0.501875pt}%
\definecolor{currentstroke}{rgb}{0.000000,0.000000,0.000000}%
\pgfsetstrokecolor{currentstroke}%
\pgfsetdash{}{0pt}%
\pgfsys@defobject{currentmarker}{\pgfqpoint{0.000000in}{-0.041667in}}{\pgfqpoint{0.000000in}{0.000000in}}{%
\pgfpathmoveto{\pgfqpoint{0.000000in}{0.000000in}}%
\pgfpathlineto{\pgfqpoint{0.000000in}{-0.041667in}}%
\pgfusepath{stroke,fill}%
}%
\begin{pgfscope}%
\pgfsys@transformshift{2.275671in}{2.414488in}%
\pgfsys@useobject{currentmarker}{}%
\end{pgfscope}%
\end{pgfscope}%
\begin{pgfscope}%
\definecolor{textcolor}{rgb}{0.000000,0.000000,0.000000}%
\pgfsetstrokecolor{textcolor}%
\pgfsetfillcolor{textcolor}%
\pgftext[x=2.275671in,y=0.371389in,,top]{\color{textcolor}{\rmfamily\fontsize{11.000000}{13.200000}\selectfont\catcode`\^=\active\def^{\ifmmode\sp\else\^{}\fi}\catcode`\%=\active\def%{\%}$\mathdefault{10^{2}}$}}%
\end{pgfscope}%
\begin{pgfscope}%
\pgfsetbuttcap%
\pgfsetroundjoin%
\definecolor{currentfill}{rgb}{0.000000,0.000000,0.000000}%
\pgfsetfillcolor{currentfill}%
\pgfsetlinewidth{0.501875pt}%
\definecolor{currentstroke}{rgb}{0.000000,0.000000,0.000000}%
\pgfsetstrokecolor{currentstroke}%
\pgfsetdash{}{0pt}%
\pgfsys@defobject{currentmarker}{\pgfqpoint{0.000000in}{0.000000in}}{\pgfqpoint{0.000000in}{0.041667in}}{%
\pgfpathmoveto{\pgfqpoint{0.000000in}{0.000000in}}%
\pgfpathlineto{\pgfqpoint{0.000000in}{0.041667in}}%
\pgfusepath{stroke,fill}%
}%
\begin{pgfscope}%
\pgfsys@transformshift{2.982169in}{0.420000in}%
\pgfsys@useobject{currentmarker}{}%
\end{pgfscope}%
\end{pgfscope}%
\begin{pgfscope}%
\pgfsetbuttcap%
\pgfsetroundjoin%
\definecolor{currentfill}{rgb}{0.000000,0.000000,0.000000}%
\pgfsetfillcolor{currentfill}%
\pgfsetlinewidth{0.501875pt}%
\definecolor{currentstroke}{rgb}{0.000000,0.000000,0.000000}%
\pgfsetstrokecolor{currentstroke}%
\pgfsetdash{}{0pt}%
\pgfsys@defobject{currentmarker}{\pgfqpoint{0.000000in}{-0.041667in}}{\pgfqpoint{0.000000in}{0.000000in}}{%
\pgfpathmoveto{\pgfqpoint{0.000000in}{0.000000in}}%
\pgfpathlineto{\pgfqpoint{0.000000in}{-0.041667in}}%
\pgfusepath{stroke,fill}%
}%
\begin{pgfscope}%
\pgfsys@transformshift{2.982169in}{2.414488in}%
\pgfsys@useobject{currentmarker}{}%
\end{pgfscope}%
\end{pgfscope}%
\begin{pgfscope}%
\definecolor{textcolor}{rgb}{0.000000,0.000000,0.000000}%
\pgfsetstrokecolor{textcolor}%
\pgfsetfillcolor{textcolor}%
\pgftext[x=2.982169in,y=0.371389in,,top]{\color{textcolor}{\rmfamily\fontsize{11.000000}{13.200000}\selectfont\catcode`\^=\active\def^{\ifmmode\sp\else\^{}\fi}\catcode`\%=\active\def%{\%}$\mathdefault{10^{3}}$}}%
\end{pgfscope}%
\begin{pgfscope}%
\pgfsetbuttcap%
\pgfsetroundjoin%
\definecolor{currentfill}{rgb}{0.000000,0.000000,0.000000}%
\pgfsetfillcolor{currentfill}%
\pgfsetlinewidth{0.401500pt}%
\definecolor{currentstroke}{rgb}{0.000000,0.000000,0.000000}%
\pgfsetstrokecolor{currentstroke}%
\pgfsetdash{}{0pt}%
\pgfsys@defobject{currentmarker}{\pgfqpoint{0.000000in}{0.000000in}}{\pgfqpoint{0.000000in}{0.020833in}}{%
\pgfpathmoveto{\pgfqpoint{0.000000in}{0.000000in}}%
\pgfpathlineto{\pgfqpoint{0.000000in}{0.020833in}}%
\pgfusepath{stroke,fill}%
}%
\begin{pgfscope}%
\pgfsys@transformshift{0.650000in}{0.420000in}%
\pgfsys@useobject{currentmarker}{}%
\end{pgfscope}%
\end{pgfscope}%
\begin{pgfscope}%
\pgfsetbuttcap%
\pgfsetroundjoin%
\definecolor{currentfill}{rgb}{0.000000,0.000000,0.000000}%
\pgfsetfillcolor{currentfill}%
\pgfsetlinewidth{0.401500pt}%
\definecolor{currentstroke}{rgb}{0.000000,0.000000,0.000000}%
\pgfsetstrokecolor{currentstroke}%
\pgfsetdash{}{0pt}%
\pgfsys@defobject{currentmarker}{\pgfqpoint{0.000000in}{-0.020833in}}{\pgfqpoint{0.000000in}{0.000000in}}{%
\pgfpathmoveto{\pgfqpoint{0.000000in}{0.000000in}}%
\pgfpathlineto{\pgfqpoint{0.000000in}{-0.020833in}}%
\pgfusepath{stroke,fill}%
}%
\begin{pgfscope}%
\pgfsys@transformshift{0.650000in}{2.414488in}%
\pgfsys@useobject{currentmarker}{}%
\end{pgfscope}%
\end{pgfscope}%
\begin{pgfscope}%
\pgfsetbuttcap%
\pgfsetroundjoin%
\definecolor{currentfill}{rgb}{0.000000,0.000000,0.000000}%
\pgfsetfillcolor{currentfill}%
\pgfsetlinewidth{0.401500pt}%
\definecolor{currentstroke}{rgb}{0.000000,0.000000,0.000000}%
\pgfsetstrokecolor{currentstroke}%
\pgfsetdash{}{0pt}%
\pgfsys@defobject{currentmarker}{\pgfqpoint{0.000000in}{0.000000in}}{\pgfqpoint{0.000000in}{0.020833in}}{%
\pgfpathmoveto{\pgfqpoint{0.000000in}{0.000000in}}%
\pgfpathlineto{\pgfqpoint{0.000000in}{0.020833in}}%
\pgfusepath{stroke,fill}%
}%
\begin{pgfscope}%
\pgfsys@transformshift{0.705941in}{0.420000in}%
\pgfsys@useobject{currentmarker}{}%
\end{pgfscope}%
\end{pgfscope}%
\begin{pgfscope}%
\pgfsetbuttcap%
\pgfsetroundjoin%
\definecolor{currentfill}{rgb}{0.000000,0.000000,0.000000}%
\pgfsetfillcolor{currentfill}%
\pgfsetlinewidth{0.401500pt}%
\definecolor{currentstroke}{rgb}{0.000000,0.000000,0.000000}%
\pgfsetstrokecolor{currentstroke}%
\pgfsetdash{}{0pt}%
\pgfsys@defobject{currentmarker}{\pgfqpoint{0.000000in}{-0.020833in}}{\pgfqpoint{0.000000in}{0.000000in}}{%
\pgfpathmoveto{\pgfqpoint{0.000000in}{0.000000in}}%
\pgfpathlineto{\pgfqpoint{0.000000in}{-0.020833in}}%
\pgfusepath{stroke,fill}%
}%
\begin{pgfscope}%
\pgfsys@transformshift{0.705941in}{2.414488in}%
\pgfsys@useobject{currentmarker}{}%
\end{pgfscope}%
\end{pgfscope}%
\begin{pgfscope}%
\pgfsetbuttcap%
\pgfsetroundjoin%
\definecolor{currentfill}{rgb}{0.000000,0.000000,0.000000}%
\pgfsetfillcolor{currentfill}%
\pgfsetlinewidth{0.401500pt}%
\definecolor{currentstroke}{rgb}{0.000000,0.000000,0.000000}%
\pgfsetstrokecolor{currentstroke}%
\pgfsetdash{}{0pt}%
\pgfsys@defobject{currentmarker}{\pgfqpoint{0.000000in}{0.000000in}}{\pgfqpoint{0.000000in}{0.020833in}}{%
\pgfpathmoveto{\pgfqpoint{0.000000in}{0.000000in}}%
\pgfpathlineto{\pgfqpoint{0.000000in}{0.020833in}}%
\pgfusepath{stroke,fill}%
}%
\begin{pgfscope}%
\pgfsys@transformshift{0.753239in}{0.420000in}%
\pgfsys@useobject{currentmarker}{}%
\end{pgfscope}%
\end{pgfscope}%
\begin{pgfscope}%
\pgfsetbuttcap%
\pgfsetroundjoin%
\definecolor{currentfill}{rgb}{0.000000,0.000000,0.000000}%
\pgfsetfillcolor{currentfill}%
\pgfsetlinewidth{0.401500pt}%
\definecolor{currentstroke}{rgb}{0.000000,0.000000,0.000000}%
\pgfsetstrokecolor{currentstroke}%
\pgfsetdash{}{0pt}%
\pgfsys@defobject{currentmarker}{\pgfqpoint{0.000000in}{-0.020833in}}{\pgfqpoint{0.000000in}{0.000000in}}{%
\pgfpathmoveto{\pgfqpoint{0.000000in}{0.000000in}}%
\pgfpathlineto{\pgfqpoint{0.000000in}{-0.020833in}}%
\pgfusepath{stroke,fill}%
}%
\begin{pgfscope}%
\pgfsys@transformshift{0.753239in}{2.414488in}%
\pgfsys@useobject{currentmarker}{}%
\end{pgfscope}%
\end{pgfscope}%
\begin{pgfscope}%
\pgfsetbuttcap%
\pgfsetroundjoin%
\definecolor{currentfill}{rgb}{0.000000,0.000000,0.000000}%
\pgfsetfillcolor{currentfill}%
\pgfsetlinewidth{0.401500pt}%
\definecolor{currentstroke}{rgb}{0.000000,0.000000,0.000000}%
\pgfsetstrokecolor{currentstroke}%
\pgfsetdash{}{0pt}%
\pgfsys@defobject{currentmarker}{\pgfqpoint{0.000000in}{0.000000in}}{\pgfqpoint{0.000000in}{0.020833in}}{%
\pgfpathmoveto{\pgfqpoint{0.000000in}{0.000000in}}%
\pgfpathlineto{\pgfqpoint{0.000000in}{0.020833in}}%
\pgfusepath{stroke,fill}%
}%
\begin{pgfscope}%
\pgfsys@transformshift{0.794210in}{0.420000in}%
\pgfsys@useobject{currentmarker}{}%
\end{pgfscope}%
\end{pgfscope}%
\begin{pgfscope}%
\pgfsetbuttcap%
\pgfsetroundjoin%
\definecolor{currentfill}{rgb}{0.000000,0.000000,0.000000}%
\pgfsetfillcolor{currentfill}%
\pgfsetlinewidth{0.401500pt}%
\definecolor{currentstroke}{rgb}{0.000000,0.000000,0.000000}%
\pgfsetstrokecolor{currentstroke}%
\pgfsetdash{}{0pt}%
\pgfsys@defobject{currentmarker}{\pgfqpoint{0.000000in}{-0.020833in}}{\pgfqpoint{0.000000in}{0.000000in}}{%
\pgfpathmoveto{\pgfqpoint{0.000000in}{0.000000in}}%
\pgfpathlineto{\pgfqpoint{0.000000in}{-0.020833in}}%
\pgfusepath{stroke,fill}%
}%
\begin{pgfscope}%
\pgfsys@transformshift{0.794210in}{2.414488in}%
\pgfsys@useobject{currentmarker}{}%
\end{pgfscope}%
\end{pgfscope}%
\begin{pgfscope}%
\pgfsetbuttcap%
\pgfsetroundjoin%
\definecolor{currentfill}{rgb}{0.000000,0.000000,0.000000}%
\pgfsetfillcolor{currentfill}%
\pgfsetlinewidth{0.401500pt}%
\definecolor{currentstroke}{rgb}{0.000000,0.000000,0.000000}%
\pgfsetstrokecolor{currentstroke}%
\pgfsetdash{}{0pt}%
\pgfsys@defobject{currentmarker}{\pgfqpoint{0.000000in}{0.000000in}}{\pgfqpoint{0.000000in}{0.020833in}}{%
\pgfpathmoveto{\pgfqpoint{0.000000in}{0.000000in}}%
\pgfpathlineto{\pgfqpoint{0.000000in}{0.020833in}}%
\pgfusepath{stroke,fill}%
}%
\begin{pgfscope}%
\pgfsys@transformshift{0.830349in}{0.420000in}%
\pgfsys@useobject{currentmarker}{}%
\end{pgfscope}%
\end{pgfscope}%
\begin{pgfscope}%
\pgfsetbuttcap%
\pgfsetroundjoin%
\definecolor{currentfill}{rgb}{0.000000,0.000000,0.000000}%
\pgfsetfillcolor{currentfill}%
\pgfsetlinewidth{0.401500pt}%
\definecolor{currentstroke}{rgb}{0.000000,0.000000,0.000000}%
\pgfsetstrokecolor{currentstroke}%
\pgfsetdash{}{0pt}%
\pgfsys@defobject{currentmarker}{\pgfqpoint{0.000000in}{-0.020833in}}{\pgfqpoint{0.000000in}{0.000000in}}{%
\pgfpathmoveto{\pgfqpoint{0.000000in}{0.000000in}}%
\pgfpathlineto{\pgfqpoint{0.000000in}{-0.020833in}}%
\pgfusepath{stroke,fill}%
}%
\begin{pgfscope}%
\pgfsys@transformshift{0.830349in}{2.414488in}%
\pgfsys@useobject{currentmarker}{}%
\end{pgfscope}%
\end{pgfscope}%
\begin{pgfscope}%
\pgfsetbuttcap%
\pgfsetroundjoin%
\definecolor{currentfill}{rgb}{0.000000,0.000000,0.000000}%
\pgfsetfillcolor{currentfill}%
\pgfsetlinewidth{0.401500pt}%
\definecolor{currentstroke}{rgb}{0.000000,0.000000,0.000000}%
\pgfsetstrokecolor{currentstroke}%
\pgfsetdash{}{0pt}%
\pgfsys@defobject{currentmarker}{\pgfqpoint{0.000000in}{0.000000in}}{\pgfqpoint{0.000000in}{0.020833in}}{%
\pgfpathmoveto{\pgfqpoint{0.000000in}{0.000000in}}%
\pgfpathlineto{\pgfqpoint{0.000000in}{0.020833in}}%
\pgfusepath{stroke,fill}%
}%
\begin{pgfscope}%
\pgfsys@transformshift{1.075354in}{0.420000in}%
\pgfsys@useobject{currentmarker}{}%
\end{pgfscope}%
\end{pgfscope}%
\begin{pgfscope}%
\pgfsetbuttcap%
\pgfsetroundjoin%
\definecolor{currentfill}{rgb}{0.000000,0.000000,0.000000}%
\pgfsetfillcolor{currentfill}%
\pgfsetlinewidth{0.401500pt}%
\definecolor{currentstroke}{rgb}{0.000000,0.000000,0.000000}%
\pgfsetstrokecolor{currentstroke}%
\pgfsetdash{}{0pt}%
\pgfsys@defobject{currentmarker}{\pgfqpoint{0.000000in}{-0.020833in}}{\pgfqpoint{0.000000in}{0.000000in}}{%
\pgfpathmoveto{\pgfqpoint{0.000000in}{0.000000in}}%
\pgfpathlineto{\pgfqpoint{0.000000in}{-0.020833in}}%
\pgfusepath{stroke,fill}%
}%
\begin{pgfscope}%
\pgfsys@transformshift{1.075354in}{2.414488in}%
\pgfsys@useobject{currentmarker}{}%
\end{pgfscope}%
\end{pgfscope}%
\begin{pgfscope}%
\pgfsetbuttcap%
\pgfsetroundjoin%
\definecolor{currentfill}{rgb}{0.000000,0.000000,0.000000}%
\pgfsetfillcolor{currentfill}%
\pgfsetlinewidth{0.401500pt}%
\definecolor{currentstroke}{rgb}{0.000000,0.000000,0.000000}%
\pgfsetstrokecolor{currentstroke}%
\pgfsetdash{}{0pt}%
\pgfsys@defobject{currentmarker}{\pgfqpoint{0.000000in}{0.000000in}}{\pgfqpoint{0.000000in}{0.020833in}}{%
\pgfpathmoveto{\pgfqpoint{0.000000in}{0.000000in}}%
\pgfpathlineto{\pgfqpoint{0.000000in}{0.020833in}}%
\pgfusepath{stroke,fill}%
}%
\begin{pgfscope}%
\pgfsys@transformshift{1.199762in}{0.420000in}%
\pgfsys@useobject{currentmarker}{}%
\end{pgfscope}%
\end{pgfscope}%
\begin{pgfscope}%
\pgfsetbuttcap%
\pgfsetroundjoin%
\definecolor{currentfill}{rgb}{0.000000,0.000000,0.000000}%
\pgfsetfillcolor{currentfill}%
\pgfsetlinewidth{0.401500pt}%
\definecolor{currentstroke}{rgb}{0.000000,0.000000,0.000000}%
\pgfsetstrokecolor{currentstroke}%
\pgfsetdash{}{0pt}%
\pgfsys@defobject{currentmarker}{\pgfqpoint{0.000000in}{-0.020833in}}{\pgfqpoint{0.000000in}{0.000000in}}{%
\pgfpathmoveto{\pgfqpoint{0.000000in}{0.000000in}}%
\pgfpathlineto{\pgfqpoint{0.000000in}{-0.020833in}}%
\pgfusepath{stroke,fill}%
}%
\begin{pgfscope}%
\pgfsys@transformshift{1.199762in}{2.414488in}%
\pgfsys@useobject{currentmarker}{}%
\end{pgfscope}%
\end{pgfscope}%
\begin{pgfscope}%
\pgfsetbuttcap%
\pgfsetroundjoin%
\definecolor{currentfill}{rgb}{0.000000,0.000000,0.000000}%
\pgfsetfillcolor{currentfill}%
\pgfsetlinewidth{0.401500pt}%
\definecolor{currentstroke}{rgb}{0.000000,0.000000,0.000000}%
\pgfsetstrokecolor{currentstroke}%
\pgfsetdash{}{0pt}%
\pgfsys@defobject{currentmarker}{\pgfqpoint{0.000000in}{0.000000in}}{\pgfqpoint{0.000000in}{0.020833in}}{%
\pgfpathmoveto{\pgfqpoint{0.000000in}{0.000000in}}%
\pgfpathlineto{\pgfqpoint{0.000000in}{0.020833in}}%
\pgfusepath{stroke,fill}%
}%
\begin{pgfscope}%
\pgfsys@transformshift{1.288031in}{0.420000in}%
\pgfsys@useobject{currentmarker}{}%
\end{pgfscope}%
\end{pgfscope}%
\begin{pgfscope}%
\pgfsetbuttcap%
\pgfsetroundjoin%
\definecolor{currentfill}{rgb}{0.000000,0.000000,0.000000}%
\pgfsetfillcolor{currentfill}%
\pgfsetlinewidth{0.401500pt}%
\definecolor{currentstroke}{rgb}{0.000000,0.000000,0.000000}%
\pgfsetstrokecolor{currentstroke}%
\pgfsetdash{}{0pt}%
\pgfsys@defobject{currentmarker}{\pgfqpoint{0.000000in}{-0.020833in}}{\pgfqpoint{0.000000in}{0.000000in}}{%
\pgfpathmoveto{\pgfqpoint{0.000000in}{0.000000in}}%
\pgfpathlineto{\pgfqpoint{0.000000in}{-0.020833in}}%
\pgfusepath{stroke,fill}%
}%
\begin{pgfscope}%
\pgfsys@transformshift{1.288031in}{2.414488in}%
\pgfsys@useobject{currentmarker}{}%
\end{pgfscope}%
\end{pgfscope}%
\begin{pgfscope}%
\pgfsetbuttcap%
\pgfsetroundjoin%
\definecolor{currentfill}{rgb}{0.000000,0.000000,0.000000}%
\pgfsetfillcolor{currentfill}%
\pgfsetlinewidth{0.401500pt}%
\definecolor{currentstroke}{rgb}{0.000000,0.000000,0.000000}%
\pgfsetstrokecolor{currentstroke}%
\pgfsetdash{}{0pt}%
\pgfsys@defobject{currentmarker}{\pgfqpoint{0.000000in}{0.000000in}}{\pgfqpoint{0.000000in}{0.020833in}}{%
\pgfpathmoveto{\pgfqpoint{0.000000in}{0.000000in}}%
\pgfpathlineto{\pgfqpoint{0.000000in}{0.020833in}}%
\pgfusepath{stroke,fill}%
}%
\begin{pgfscope}%
\pgfsys@transformshift{1.356497in}{0.420000in}%
\pgfsys@useobject{currentmarker}{}%
\end{pgfscope}%
\end{pgfscope}%
\begin{pgfscope}%
\pgfsetbuttcap%
\pgfsetroundjoin%
\definecolor{currentfill}{rgb}{0.000000,0.000000,0.000000}%
\pgfsetfillcolor{currentfill}%
\pgfsetlinewidth{0.401500pt}%
\definecolor{currentstroke}{rgb}{0.000000,0.000000,0.000000}%
\pgfsetstrokecolor{currentstroke}%
\pgfsetdash{}{0pt}%
\pgfsys@defobject{currentmarker}{\pgfqpoint{0.000000in}{-0.020833in}}{\pgfqpoint{0.000000in}{0.000000in}}{%
\pgfpathmoveto{\pgfqpoint{0.000000in}{0.000000in}}%
\pgfpathlineto{\pgfqpoint{0.000000in}{-0.020833in}}%
\pgfusepath{stroke,fill}%
}%
\begin{pgfscope}%
\pgfsys@transformshift{1.356497in}{2.414488in}%
\pgfsys@useobject{currentmarker}{}%
\end{pgfscope}%
\end{pgfscope}%
\begin{pgfscope}%
\pgfsetbuttcap%
\pgfsetroundjoin%
\definecolor{currentfill}{rgb}{0.000000,0.000000,0.000000}%
\pgfsetfillcolor{currentfill}%
\pgfsetlinewidth{0.401500pt}%
\definecolor{currentstroke}{rgb}{0.000000,0.000000,0.000000}%
\pgfsetstrokecolor{currentstroke}%
\pgfsetdash{}{0pt}%
\pgfsys@defobject{currentmarker}{\pgfqpoint{0.000000in}{0.000000in}}{\pgfqpoint{0.000000in}{0.020833in}}{%
\pgfpathmoveto{\pgfqpoint{0.000000in}{0.000000in}}%
\pgfpathlineto{\pgfqpoint{0.000000in}{0.020833in}}%
\pgfusepath{stroke,fill}%
}%
\begin{pgfscope}%
\pgfsys@transformshift{1.412439in}{0.420000in}%
\pgfsys@useobject{currentmarker}{}%
\end{pgfscope}%
\end{pgfscope}%
\begin{pgfscope}%
\pgfsetbuttcap%
\pgfsetroundjoin%
\definecolor{currentfill}{rgb}{0.000000,0.000000,0.000000}%
\pgfsetfillcolor{currentfill}%
\pgfsetlinewidth{0.401500pt}%
\definecolor{currentstroke}{rgb}{0.000000,0.000000,0.000000}%
\pgfsetstrokecolor{currentstroke}%
\pgfsetdash{}{0pt}%
\pgfsys@defobject{currentmarker}{\pgfqpoint{0.000000in}{-0.020833in}}{\pgfqpoint{0.000000in}{0.000000in}}{%
\pgfpathmoveto{\pgfqpoint{0.000000in}{0.000000in}}%
\pgfpathlineto{\pgfqpoint{0.000000in}{-0.020833in}}%
\pgfusepath{stroke,fill}%
}%
\begin{pgfscope}%
\pgfsys@transformshift{1.412439in}{2.414488in}%
\pgfsys@useobject{currentmarker}{}%
\end{pgfscope}%
\end{pgfscope}%
\begin{pgfscope}%
\pgfsetbuttcap%
\pgfsetroundjoin%
\definecolor{currentfill}{rgb}{0.000000,0.000000,0.000000}%
\pgfsetfillcolor{currentfill}%
\pgfsetlinewidth{0.401500pt}%
\definecolor{currentstroke}{rgb}{0.000000,0.000000,0.000000}%
\pgfsetstrokecolor{currentstroke}%
\pgfsetdash{}{0pt}%
\pgfsys@defobject{currentmarker}{\pgfqpoint{0.000000in}{0.000000in}}{\pgfqpoint{0.000000in}{0.020833in}}{%
\pgfpathmoveto{\pgfqpoint{0.000000in}{0.000000in}}%
\pgfpathlineto{\pgfqpoint{0.000000in}{0.020833in}}%
\pgfusepath{stroke,fill}%
}%
\begin{pgfscope}%
\pgfsys@transformshift{1.459736in}{0.420000in}%
\pgfsys@useobject{currentmarker}{}%
\end{pgfscope}%
\end{pgfscope}%
\begin{pgfscope}%
\pgfsetbuttcap%
\pgfsetroundjoin%
\definecolor{currentfill}{rgb}{0.000000,0.000000,0.000000}%
\pgfsetfillcolor{currentfill}%
\pgfsetlinewidth{0.401500pt}%
\definecolor{currentstroke}{rgb}{0.000000,0.000000,0.000000}%
\pgfsetstrokecolor{currentstroke}%
\pgfsetdash{}{0pt}%
\pgfsys@defobject{currentmarker}{\pgfqpoint{0.000000in}{-0.020833in}}{\pgfqpoint{0.000000in}{0.000000in}}{%
\pgfpathmoveto{\pgfqpoint{0.000000in}{0.000000in}}%
\pgfpathlineto{\pgfqpoint{0.000000in}{-0.020833in}}%
\pgfusepath{stroke,fill}%
}%
\begin{pgfscope}%
\pgfsys@transformshift{1.459736in}{2.414488in}%
\pgfsys@useobject{currentmarker}{}%
\end{pgfscope}%
\end{pgfscope}%
\begin{pgfscope}%
\pgfsetbuttcap%
\pgfsetroundjoin%
\definecolor{currentfill}{rgb}{0.000000,0.000000,0.000000}%
\pgfsetfillcolor{currentfill}%
\pgfsetlinewidth{0.401500pt}%
\definecolor{currentstroke}{rgb}{0.000000,0.000000,0.000000}%
\pgfsetstrokecolor{currentstroke}%
\pgfsetdash{}{0pt}%
\pgfsys@defobject{currentmarker}{\pgfqpoint{0.000000in}{0.000000in}}{\pgfqpoint{0.000000in}{0.020833in}}{%
\pgfpathmoveto{\pgfqpoint{0.000000in}{0.000000in}}%
\pgfpathlineto{\pgfqpoint{0.000000in}{0.020833in}}%
\pgfusepath{stroke,fill}%
}%
\begin{pgfscope}%
\pgfsys@transformshift{1.500708in}{0.420000in}%
\pgfsys@useobject{currentmarker}{}%
\end{pgfscope}%
\end{pgfscope}%
\begin{pgfscope}%
\pgfsetbuttcap%
\pgfsetroundjoin%
\definecolor{currentfill}{rgb}{0.000000,0.000000,0.000000}%
\pgfsetfillcolor{currentfill}%
\pgfsetlinewidth{0.401500pt}%
\definecolor{currentstroke}{rgb}{0.000000,0.000000,0.000000}%
\pgfsetstrokecolor{currentstroke}%
\pgfsetdash{}{0pt}%
\pgfsys@defobject{currentmarker}{\pgfqpoint{0.000000in}{-0.020833in}}{\pgfqpoint{0.000000in}{0.000000in}}{%
\pgfpathmoveto{\pgfqpoint{0.000000in}{0.000000in}}%
\pgfpathlineto{\pgfqpoint{0.000000in}{-0.020833in}}%
\pgfusepath{stroke,fill}%
}%
\begin{pgfscope}%
\pgfsys@transformshift{1.500708in}{2.414488in}%
\pgfsys@useobject{currentmarker}{}%
\end{pgfscope}%
\end{pgfscope}%
\begin{pgfscope}%
\pgfsetbuttcap%
\pgfsetroundjoin%
\definecolor{currentfill}{rgb}{0.000000,0.000000,0.000000}%
\pgfsetfillcolor{currentfill}%
\pgfsetlinewidth{0.401500pt}%
\definecolor{currentstroke}{rgb}{0.000000,0.000000,0.000000}%
\pgfsetstrokecolor{currentstroke}%
\pgfsetdash{}{0pt}%
\pgfsys@defobject{currentmarker}{\pgfqpoint{0.000000in}{0.000000in}}{\pgfqpoint{0.000000in}{0.020833in}}{%
\pgfpathmoveto{\pgfqpoint{0.000000in}{0.000000in}}%
\pgfpathlineto{\pgfqpoint{0.000000in}{0.020833in}}%
\pgfusepath{stroke,fill}%
}%
\begin{pgfscope}%
\pgfsys@transformshift{1.536847in}{0.420000in}%
\pgfsys@useobject{currentmarker}{}%
\end{pgfscope}%
\end{pgfscope}%
\begin{pgfscope}%
\pgfsetbuttcap%
\pgfsetroundjoin%
\definecolor{currentfill}{rgb}{0.000000,0.000000,0.000000}%
\pgfsetfillcolor{currentfill}%
\pgfsetlinewidth{0.401500pt}%
\definecolor{currentstroke}{rgb}{0.000000,0.000000,0.000000}%
\pgfsetstrokecolor{currentstroke}%
\pgfsetdash{}{0pt}%
\pgfsys@defobject{currentmarker}{\pgfqpoint{0.000000in}{-0.020833in}}{\pgfqpoint{0.000000in}{0.000000in}}{%
\pgfpathmoveto{\pgfqpoint{0.000000in}{0.000000in}}%
\pgfpathlineto{\pgfqpoint{0.000000in}{-0.020833in}}%
\pgfusepath{stroke,fill}%
}%
\begin{pgfscope}%
\pgfsys@transformshift{1.536847in}{2.414488in}%
\pgfsys@useobject{currentmarker}{}%
\end{pgfscope}%
\end{pgfscope}%
\begin{pgfscope}%
\pgfsetbuttcap%
\pgfsetroundjoin%
\definecolor{currentfill}{rgb}{0.000000,0.000000,0.000000}%
\pgfsetfillcolor{currentfill}%
\pgfsetlinewidth{0.401500pt}%
\definecolor{currentstroke}{rgb}{0.000000,0.000000,0.000000}%
\pgfsetstrokecolor{currentstroke}%
\pgfsetdash{}{0pt}%
\pgfsys@defobject{currentmarker}{\pgfqpoint{0.000000in}{0.000000in}}{\pgfqpoint{0.000000in}{0.020833in}}{%
\pgfpathmoveto{\pgfqpoint{0.000000in}{0.000000in}}%
\pgfpathlineto{\pgfqpoint{0.000000in}{0.020833in}}%
\pgfusepath{stroke,fill}%
}%
\begin{pgfscope}%
\pgfsys@transformshift{1.781851in}{0.420000in}%
\pgfsys@useobject{currentmarker}{}%
\end{pgfscope}%
\end{pgfscope}%
\begin{pgfscope}%
\pgfsetbuttcap%
\pgfsetroundjoin%
\definecolor{currentfill}{rgb}{0.000000,0.000000,0.000000}%
\pgfsetfillcolor{currentfill}%
\pgfsetlinewidth{0.401500pt}%
\definecolor{currentstroke}{rgb}{0.000000,0.000000,0.000000}%
\pgfsetstrokecolor{currentstroke}%
\pgfsetdash{}{0pt}%
\pgfsys@defobject{currentmarker}{\pgfqpoint{0.000000in}{-0.020833in}}{\pgfqpoint{0.000000in}{0.000000in}}{%
\pgfpathmoveto{\pgfqpoint{0.000000in}{0.000000in}}%
\pgfpathlineto{\pgfqpoint{0.000000in}{-0.020833in}}%
\pgfusepath{stroke,fill}%
}%
\begin{pgfscope}%
\pgfsys@transformshift{1.781851in}{2.414488in}%
\pgfsys@useobject{currentmarker}{}%
\end{pgfscope}%
\end{pgfscope}%
\begin{pgfscope}%
\pgfsetbuttcap%
\pgfsetroundjoin%
\definecolor{currentfill}{rgb}{0.000000,0.000000,0.000000}%
\pgfsetfillcolor{currentfill}%
\pgfsetlinewidth{0.401500pt}%
\definecolor{currentstroke}{rgb}{0.000000,0.000000,0.000000}%
\pgfsetstrokecolor{currentstroke}%
\pgfsetdash{}{0pt}%
\pgfsys@defobject{currentmarker}{\pgfqpoint{0.000000in}{0.000000in}}{\pgfqpoint{0.000000in}{0.020833in}}{%
\pgfpathmoveto{\pgfqpoint{0.000000in}{0.000000in}}%
\pgfpathlineto{\pgfqpoint{0.000000in}{0.020833in}}%
\pgfusepath{stroke,fill}%
}%
\begin{pgfscope}%
\pgfsys@transformshift{1.906259in}{0.420000in}%
\pgfsys@useobject{currentmarker}{}%
\end{pgfscope}%
\end{pgfscope}%
\begin{pgfscope}%
\pgfsetbuttcap%
\pgfsetroundjoin%
\definecolor{currentfill}{rgb}{0.000000,0.000000,0.000000}%
\pgfsetfillcolor{currentfill}%
\pgfsetlinewidth{0.401500pt}%
\definecolor{currentstroke}{rgb}{0.000000,0.000000,0.000000}%
\pgfsetstrokecolor{currentstroke}%
\pgfsetdash{}{0pt}%
\pgfsys@defobject{currentmarker}{\pgfqpoint{0.000000in}{-0.020833in}}{\pgfqpoint{0.000000in}{0.000000in}}{%
\pgfpathmoveto{\pgfqpoint{0.000000in}{0.000000in}}%
\pgfpathlineto{\pgfqpoint{0.000000in}{-0.020833in}}%
\pgfusepath{stroke,fill}%
}%
\begin{pgfscope}%
\pgfsys@transformshift{1.906259in}{2.414488in}%
\pgfsys@useobject{currentmarker}{}%
\end{pgfscope}%
\end{pgfscope}%
\begin{pgfscope}%
\pgfsetbuttcap%
\pgfsetroundjoin%
\definecolor{currentfill}{rgb}{0.000000,0.000000,0.000000}%
\pgfsetfillcolor{currentfill}%
\pgfsetlinewidth{0.401500pt}%
\definecolor{currentstroke}{rgb}{0.000000,0.000000,0.000000}%
\pgfsetstrokecolor{currentstroke}%
\pgfsetdash{}{0pt}%
\pgfsys@defobject{currentmarker}{\pgfqpoint{0.000000in}{0.000000in}}{\pgfqpoint{0.000000in}{0.020833in}}{%
\pgfpathmoveto{\pgfqpoint{0.000000in}{0.000000in}}%
\pgfpathlineto{\pgfqpoint{0.000000in}{0.020833in}}%
\pgfusepath{stroke,fill}%
}%
\begin{pgfscope}%
\pgfsys@transformshift{1.994528in}{0.420000in}%
\pgfsys@useobject{currentmarker}{}%
\end{pgfscope}%
\end{pgfscope}%
\begin{pgfscope}%
\pgfsetbuttcap%
\pgfsetroundjoin%
\definecolor{currentfill}{rgb}{0.000000,0.000000,0.000000}%
\pgfsetfillcolor{currentfill}%
\pgfsetlinewidth{0.401500pt}%
\definecolor{currentstroke}{rgb}{0.000000,0.000000,0.000000}%
\pgfsetstrokecolor{currentstroke}%
\pgfsetdash{}{0pt}%
\pgfsys@defobject{currentmarker}{\pgfqpoint{0.000000in}{-0.020833in}}{\pgfqpoint{0.000000in}{0.000000in}}{%
\pgfpathmoveto{\pgfqpoint{0.000000in}{0.000000in}}%
\pgfpathlineto{\pgfqpoint{0.000000in}{-0.020833in}}%
\pgfusepath{stroke,fill}%
}%
\begin{pgfscope}%
\pgfsys@transformshift{1.994528in}{2.414488in}%
\pgfsys@useobject{currentmarker}{}%
\end{pgfscope}%
\end{pgfscope}%
\begin{pgfscope}%
\pgfsetbuttcap%
\pgfsetroundjoin%
\definecolor{currentfill}{rgb}{0.000000,0.000000,0.000000}%
\pgfsetfillcolor{currentfill}%
\pgfsetlinewidth{0.401500pt}%
\definecolor{currentstroke}{rgb}{0.000000,0.000000,0.000000}%
\pgfsetstrokecolor{currentstroke}%
\pgfsetdash{}{0pt}%
\pgfsys@defobject{currentmarker}{\pgfqpoint{0.000000in}{0.000000in}}{\pgfqpoint{0.000000in}{0.020833in}}{%
\pgfpathmoveto{\pgfqpoint{0.000000in}{0.000000in}}%
\pgfpathlineto{\pgfqpoint{0.000000in}{0.020833in}}%
\pgfusepath{stroke,fill}%
}%
\begin{pgfscope}%
\pgfsys@transformshift{2.062995in}{0.420000in}%
\pgfsys@useobject{currentmarker}{}%
\end{pgfscope}%
\end{pgfscope}%
\begin{pgfscope}%
\pgfsetbuttcap%
\pgfsetroundjoin%
\definecolor{currentfill}{rgb}{0.000000,0.000000,0.000000}%
\pgfsetfillcolor{currentfill}%
\pgfsetlinewidth{0.401500pt}%
\definecolor{currentstroke}{rgb}{0.000000,0.000000,0.000000}%
\pgfsetstrokecolor{currentstroke}%
\pgfsetdash{}{0pt}%
\pgfsys@defobject{currentmarker}{\pgfqpoint{0.000000in}{-0.020833in}}{\pgfqpoint{0.000000in}{0.000000in}}{%
\pgfpathmoveto{\pgfqpoint{0.000000in}{0.000000in}}%
\pgfpathlineto{\pgfqpoint{0.000000in}{-0.020833in}}%
\pgfusepath{stroke,fill}%
}%
\begin{pgfscope}%
\pgfsys@transformshift{2.062995in}{2.414488in}%
\pgfsys@useobject{currentmarker}{}%
\end{pgfscope}%
\end{pgfscope}%
\begin{pgfscope}%
\pgfsetbuttcap%
\pgfsetroundjoin%
\definecolor{currentfill}{rgb}{0.000000,0.000000,0.000000}%
\pgfsetfillcolor{currentfill}%
\pgfsetlinewidth{0.401500pt}%
\definecolor{currentstroke}{rgb}{0.000000,0.000000,0.000000}%
\pgfsetstrokecolor{currentstroke}%
\pgfsetdash{}{0pt}%
\pgfsys@defobject{currentmarker}{\pgfqpoint{0.000000in}{0.000000in}}{\pgfqpoint{0.000000in}{0.020833in}}{%
\pgfpathmoveto{\pgfqpoint{0.000000in}{0.000000in}}%
\pgfpathlineto{\pgfqpoint{0.000000in}{0.020833in}}%
\pgfusepath{stroke,fill}%
}%
\begin{pgfscope}%
\pgfsys@transformshift{2.118936in}{0.420000in}%
\pgfsys@useobject{currentmarker}{}%
\end{pgfscope}%
\end{pgfscope}%
\begin{pgfscope}%
\pgfsetbuttcap%
\pgfsetroundjoin%
\definecolor{currentfill}{rgb}{0.000000,0.000000,0.000000}%
\pgfsetfillcolor{currentfill}%
\pgfsetlinewidth{0.401500pt}%
\definecolor{currentstroke}{rgb}{0.000000,0.000000,0.000000}%
\pgfsetstrokecolor{currentstroke}%
\pgfsetdash{}{0pt}%
\pgfsys@defobject{currentmarker}{\pgfqpoint{0.000000in}{-0.020833in}}{\pgfqpoint{0.000000in}{0.000000in}}{%
\pgfpathmoveto{\pgfqpoint{0.000000in}{0.000000in}}%
\pgfpathlineto{\pgfqpoint{0.000000in}{-0.020833in}}%
\pgfusepath{stroke,fill}%
}%
\begin{pgfscope}%
\pgfsys@transformshift{2.118936in}{2.414488in}%
\pgfsys@useobject{currentmarker}{}%
\end{pgfscope}%
\end{pgfscope}%
\begin{pgfscope}%
\pgfsetbuttcap%
\pgfsetroundjoin%
\definecolor{currentfill}{rgb}{0.000000,0.000000,0.000000}%
\pgfsetfillcolor{currentfill}%
\pgfsetlinewidth{0.401500pt}%
\definecolor{currentstroke}{rgb}{0.000000,0.000000,0.000000}%
\pgfsetstrokecolor{currentstroke}%
\pgfsetdash{}{0pt}%
\pgfsys@defobject{currentmarker}{\pgfqpoint{0.000000in}{0.000000in}}{\pgfqpoint{0.000000in}{0.020833in}}{%
\pgfpathmoveto{\pgfqpoint{0.000000in}{0.000000in}}%
\pgfpathlineto{\pgfqpoint{0.000000in}{0.020833in}}%
\pgfusepath{stroke,fill}%
}%
\begin{pgfscope}%
\pgfsys@transformshift{2.166234in}{0.420000in}%
\pgfsys@useobject{currentmarker}{}%
\end{pgfscope}%
\end{pgfscope}%
\begin{pgfscope}%
\pgfsetbuttcap%
\pgfsetroundjoin%
\definecolor{currentfill}{rgb}{0.000000,0.000000,0.000000}%
\pgfsetfillcolor{currentfill}%
\pgfsetlinewidth{0.401500pt}%
\definecolor{currentstroke}{rgb}{0.000000,0.000000,0.000000}%
\pgfsetstrokecolor{currentstroke}%
\pgfsetdash{}{0pt}%
\pgfsys@defobject{currentmarker}{\pgfqpoint{0.000000in}{-0.020833in}}{\pgfqpoint{0.000000in}{0.000000in}}{%
\pgfpathmoveto{\pgfqpoint{0.000000in}{0.000000in}}%
\pgfpathlineto{\pgfqpoint{0.000000in}{-0.020833in}}%
\pgfusepath{stroke,fill}%
}%
\begin{pgfscope}%
\pgfsys@transformshift{2.166234in}{2.414488in}%
\pgfsys@useobject{currentmarker}{}%
\end{pgfscope}%
\end{pgfscope}%
\begin{pgfscope}%
\pgfsetbuttcap%
\pgfsetroundjoin%
\definecolor{currentfill}{rgb}{0.000000,0.000000,0.000000}%
\pgfsetfillcolor{currentfill}%
\pgfsetlinewidth{0.401500pt}%
\definecolor{currentstroke}{rgb}{0.000000,0.000000,0.000000}%
\pgfsetstrokecolor{currentstroke}%
\pgfsetdash{}{0pt}%
\pgfsys@defobject{currentmarker}{\pgfqpoint{0.000000in}{0.000000in}}{\pgfqpoint{0.000000in}{0.020833in}}{%
\pgfpathmoveto{\pgfqpoint{0.000000in}{0.000000in}}%
\pgfpathlineto{\pgfqpoint{0.000000in}{0.020833in}}%
\pgfusepath{stroke,fill}%
}%
\begin{pgfscope}%
\pgfsys@transformshift{2.207205in}{0.420000in}%
\pgfsys@useobject{currentmarker}{}%
\end{pgfscope}%
\end{pgfscope}%
\begin{pgfscope}%
\pgfsetbuttcap%
\pgfsetroundjoin%
\definecolor{currentfill}{rgb}{0.000000,0.000000,0.000000}%
\pgfsetfillcolor{currentfill}%
\pgfsetlinewidth{0.401500pt}%
\definecolor{currentstroke}{rgb}{0.000000,0.000000,0.000000}%
\pgfsetstrokecolor{currentstroke}%
\pgfsetdash{}{0pt}%
\pgfsys@defobject{currentmarker}{\pgfqpoint{0.000000in}{-0.020833in}}{\pgfqpoint{0.000000in}{0.000000in}}{%
\pgfpathmoveto{\pgfqpoint{0.000000in}{0.000000in}}%
\pgfpathlineto{\pgfqpoint{0.000000in}{-0.020833in}}%
\pgfusepath{stroke,fill}%
}%
\begin{pgfscope}%
\pgfsys@transformshift{2.207205in}{2.414488in}%
\pgfsys@useobject{currentmarker}{}%
\end{pgfscope}%
\end{pgfscope}%
\begin{pgfscope}%
\pgfsetbuttcap%
\pgfsetroundjoin%
\definecolor{currentfill}{rgb}{0.000000,0.000000,0.000000}%
\pgfsetfillcolor{currentfill}%
\pgfsetlinewidth{0.401500pt}%
\definecolor{currentstroke}{rgb}{0.000000,0.000000,0.000000}%
\pgfsetstrokecolor{currentstroke}%
\pgfsetdash{}{0pt}%
\pgfsys@defobject{currentmarker}{\pgfqpoint{0.000000in}{0.000000in}}{\pgfqpoint{0.000000in}{0.020833in}}{%
\pgfpathmoveto{\pgfqpoint{0.000000in}{0.000000in}}%
\pgfpathlineto{\pgfqpoint{0.000000in}{0.020833in}}%
\pgfusepath{stroke,fill}%
}%
\begin{pgfscope}%
\pgfsys@transformshift{2.243344in}{0.420000in}%
\pgfsys@useobject{currentmarker}{}%
\end{pgfscope}%
\end{pgfscope}%
\begin{pgfscope}%
\pgfsetbuttcap%
\pgfsetroundjoin%
\definecolor{currentfill}{rgb}{0.000000,0.000000,0.000000}%
\pgfsetfillcolor{currentfill}%
\pgfsetlinewidth{0.401500pt}%
\definecolor{currentstroke}{rgb}{0.000000,0.000000,0.000000}%
\pgfsetstrokecolor{currentstroke}%
\pgfsetdash{}{0pt}%
\pgfsys@defobject{currentmarker}{\pgfqpoint{0.000000in}{-0.020833in}}{\pgfqpoint{0.000000in}{0.000000in}}{%
\pgfpathmoveto{\pgfqpoint{0.000000in}{0.000000in}}%
\pgfpathlineto{\pgfqpoint{0.000000in}{-0.020833in}}%
\pgfusepath{stroke,fill}%
}%
\begin{pgfscope}%
\pgfsys@transformshift{2.243344in}{2.414488in}%
\pgfsys@useobject{currentmarker}{}%
\end{pgfscope}%
\end{pgfscope}%
\begin{pgfscope}%
\pgfsetbuttcap%
\pgfsetroundjoin%
\definecolor{currentfill}{rgb}{0.000000,0.000000,0.000000}%
\pgfsetfillcolor{currentfill}%
\pgfsetlinewidth{0.401500pt}%
\definecolor{currentstroke}{rgb}{0.000000,0.000000,0.000000}%
\pgfsetstrokecolor{currentstroke}%
\pgfsetdash{}{0pt}%
\pgfsys@defobject{currentmarker}{\pgfqpoint{0.000000in}{0.000000in}}{\pgfqpoint{0.000000in}{0.020833in}}{%
\pgfpathmoveto{\pgfqpoint{0.000000in}{0.000000in}}%
\pgfpathlineto{\pgfqpoint{0.000000in}{0.020833in}}%
\pgfusepath{stroke,fill}%
}%
\begin{pgfscope}%
\pgfsys@transformshift{2.488348in}{0.420000in}%
\pgfsys@useobject{currentmarker}{}%
\end{pgfscope}%
\end{pgfscope}%
\begin{pgfscope}%
\pgfsetbuttcap%
\pgfsetroundjoin%
\definecolor{currentfill}{rgb}{0.000000,0.000000,0.000000}%
\pgfsetfillcolor{currentfill}%
\pgfsetlinewidth{0.401500pt}%
\definecolor{currentstroke}{rgb}{0.000000,0.000000,0.000000}%
\pgfsetstrokecolor{currentstroke}%
\pgfsetdash{}{0pt}%
\pgfsys@defobject{currentmarker}{\pgfqpoint{0.000000in}{-0.020833in}}{\pgfqpoint{0.000000in}{0.000000in}}{%
\pgfpathmoveto{\pgfqpoint{0.000000in}{0.000000in}}%
\pgfpathlineto{\pgfqpoint{0.000000in}{-0.020833in}}%
\pgfusepath{stroke,fill}%
}%
\begin{pgfscope}%
\pgfsys@transformshift{2.488348in}{2.414488in}%
\pgfsys@useobject{currentmarker}{}%
\end{pgfscope}%
\end{pgfscope}%
\begin{pgfscope}%
\pgfsetbuttcap%
\pgfsetroundjoin%
\definecolor{currentfill}{rgb}{0.000000,0.000000,0.000000}%
\pgfsetfillcolor{currentfill}%
\pgfsetlinewidth{0.401500pt}%
\definecolor{currentstroke}{rgb}{0.000000,0.000000,0.000000}%
\pgfsetstrokecolor{currentstroke}%
\pgfsetdash{}{0pt}%
\pgfsys@defobject{currentmarker}{\pgfqpoint{0.000000in}{0.000000in}}{\pgfqpoint{0.000000in}{0.020833in}}{%
\pgfpathmoveto{\pgfqpoint{0.000000in}{0.000000in}}%
\pgfpathlineto{\pgfqpoint{0.000000in}{0.020833in}}%
\pgfusepath{stroke,fill}%
}%
\begin{pgfscope}%
\pgfsys@transformshift{2.612756in}{0.420000in}%
\pgfsys@useobject{currentmarker}{}%
\end{pgfscope}%
\end{pgfscope}%
\begin{pgfscope}%
\pgfsetbuttcap%
\pgfsetroundjoin%
\definecolor{currentfill}{rgb}{0.000000,0.000000,0.000000}%
\pgfsetfillcolor{currentfill}%
\pgfsetlinewidth{0.401500pt}%
\definecolor{currentstroke}{rgb}{0.000000,0.000000,0.000000}%
\pgfsetstrokecolor{currentstroke}%
\pgfsetdash{}{0pt}%
\pgfsys@defobject{currentmarker}{\pgfqpoint{0.000000in}{-0.020833in}}{\pgfqpoint{0.000000in}{0.000000in}}{%
\pgfpathmoveto{\pgfqpoint{0.000000in}{0.000000in}}%
\pgfpathlineto{\pgfqpoint{0.000000in}{-0.020833in}}%
\pgfusepath{stroke,fill}%
}%
\begin{pgfscope}%
\pgfsys@transformshift{2.612756in}{2.414488in}%
\pgfsys@useobject{currentmarker}{}%
\end{pgfscope}%
\end{pgfscope}%
\begin{pgfscope}%
\pgfsetbuttcap%
\pgfsetroundjoin%
\definecolor{currentfill}{rgb}{0.000000,0.000000,0.000000}%
\pgfsetfillcolor{currentfill}%
\pgfsetlinewidth{0.401500pt}%
\definecolor{currentstroke}{rgb}{0.000000,0.000000,0.000000}%
\pgfsetstrokecolor{currentstroke}%
\pgfsetdash{}{0pt}%
\pgfsys@defobject{currentmarker}{\pgfqpoint{0.000000in}{0.000000in}}{\pgfqpoint{0.000000in}{0.020833in}}{%
\pgfpathmoveto{\pgfqpoint{0.000000in}{0.000000in}}%
\pgfpathlineto{\pgfqpoint{0.000000in}{0.020833in}}%
\pgfusepath{stroke,fill}%
}%
\begin{pgfscope}%
\pgfsys@transformshift{2.701025in}{0.420000in}%
\pgfsys@useobject{currentmarker}{}%
\end{pgfscope}%
\end{pgfscope}%
\begin{pgfscope}%
\pgfsetbuttcap%
\pgfsetroundjoin%
\definecolor{currentfill}{rgb}{0.000000,0.000000,0.000000}%
\pgfsetfillcolor{currentfill}%
\pgfsetlinewidth{0.401500pt}%
\definecolor{currentstroke}{rgb}{0.000000,0.000000,0.000000}%
\pgfsetstrokecolor{currentstroke}%
\pgfsetdash{}{0pt}%
\pgfsys@defobject{currentmarker}{\pgfqpoint{0.000000in}{-0.020833in}}{\pgfqpoint{0.000000in}{0.000000in}}{%
\pgfpathmoveto{\pgfqpoint{0.000000in}{0.000000in}}%
\pgfpathlineto{\pgfqpoint{0.000000in}{-0.020833in}}%
\pgfusepath{stroke,fill}%
}%
\begin{pgfscope}%
\pgfsys@transformshift{2.701025in}{2.414488in}%
\pgfsys@useobject{currentmarker}{}%
\end{pgfscope}%
\end{pgfscope}%
\begin{pgfscope}%
\pgfsetbuttcap%
\pgfsetroundjoin%
\definecolor{currentfill}{rgb}{0.000000,0.000000,0.000000}%
\pgfsetfillcolor{currentfill}%
\pgfsetlinewidth{0.401500pt}%
\definecolor{currentstroke}{rgb}{0.000000,0.000000,0.000000}%
\pgfsetstrokecolor{currentstroke}%
\pgfsetdash{}{0pt}%
\pgfsys@defobject{currentmarker}{\pgfqpoint{0.000000in}{0.000000in}}{\pgfqpoint{0.000000in}{0.020833in}}{%
\pgfpathmoveto{\pgfqpoint{0.000000in}{0.000000in}}%
\pgfpathlineto{\pgfqpoint{0.000000in}{0.020833in}}%
\pgfusepath{stroke,fill}%
}%
\begin{pgfscope}%
\pgfsys@transformshift{2.769492in}{0.420000in}%
\pgfsys@useobject{currentmarker}{}%
\end{pgfscope}%
\end{pgfscope}%
\begin{pgfscope}%
\pgfsetbuttcap%
\pgfsetroundjoin%
\definecolor{currentfill}{rgb}{0.000000,0.000000,0.000000}%
\pgfsetfillcolor{currentfill}%
\pgfsetlinewidth{0.401500pt}%
\definecolor{currentstroke}{rgb}{0.000000,0.000000,0.000000}%
\pgfsetstrokecolor{currentstroke}%
\pgfsetdash{}{0pt}%
\pgfsys@defobject{currentmarker}{\pgfqpoint{0.000000in}{-0.020833in}}{\pgfqpoint{0.000000in}{0.000000in}}{%
\pgfpathmoveto{\pgfqpoint{0.000000in}{0.000000in}}%
\pgfpathlineto{\pgfqpoint{0.000000in}{-0.020833in}}%
\pgfusepath{stroke,fill}%
}%
\begin{pgfscope}%
\pgfsys@transformshift{2.769492in}{2.414488in}%
\pgfsys@useobject{currentmarker}{}%
\end{pgfscope}%
\end{pgfscope}%
\begin{pgfscope}%
\pgfsetbuttcap%
\pgfsetroundjoin%
\definecolor{currentfill}{rgb}{0.000000,0.000000,0.000000}%
\pgfsetfillcolor{currentfill}%
\pgfsetlinewidth{0.401500pt}%
\definecolor{currentstroke}{rgb}{0.000000,0.000000,0.000000}%
\pgfsetstrokecolor{currentstroke}%
\pgfsetdash{}{0pt}%
\pgfsys@defobject{currentmarker}{\pgfqpoint{0.000000in}{0.000000in}}{\pgfqpoint{0.000000in}{0.020833in}}{%
\pgfpathmoveto{\pgfqpoint{0.000000in}{0.000000in}}%
\pgfpathlineto{\pgfqpoint{0.000000in}{0.020833in}}%
\pgfusepath{stroke,fill}%
}%
\begin{pgfscope}%
\pgfsys@transformshift{2.825433in}{0.420000in}%
\pgfsys@useobject{currentmarker}{}%
\end{pgfscope}%
\end{pgfscope}%
\begin{pgfscope}%
\pgfsetbuttcap%
\pgfsetroundjoin%
\definecolor{currentfill}{rgb}{0.000000,0.000000,0.000000}%
\pgfsetfillcolor{currentfill}%
\pgfsetlinewidth{0.401500pt}%
\definecolor{currentstroke}{rgb}{0.000000,0.000000,0.000000}%
\pgfsetstrokecolor{currentstroke}%
\pgfsetdash{}{0pt}%
\pgfsys@defobject{currentmarker}{\pgfqpoint{0.000000in}{-0.020833in}}{\pgfqpoint{0.000000in}{0.000000in}}{%
\pgfpathmoveto{\pgfqpoint{0.000000in}{0.000000in}}%
\pgfpathlineto{\pgfqpoint{0.000000in}{-0.020833in}}%
\pgfusepath{stroke,fill}%
}%
\begin{pgfscope}%
\pgfsys@transformshift{2.825433in}{2.414488in}%
\pgfsys@useobject{currentmarker}{}%
\end{pgfscope}%
\end{pgfscope}%
\begin{pgfscope}%
\pgfsetbuttcap%
\pgfsetroundjoin%
\definecolor{currentfill}{rgb}{0.000000,0.000000,0.000000}%
\pgfsetfillcolor{currentfill}%
\pgfsetlinewidth{0.401500pt}%
\definecolor{currentstroke}{rgb}{0.000000,0.000000,0.000000}%
\pgfsetstrokecolor{currentstroke}%
\pgfsetdash{}{0pt}%
\pgfsys@defobject{currentmarker}{\pgfqpoint{0.000000in}{0.000000in}}{\pgfqpoint{0.000000in}{0.020833in}}{%
\pgfpathmoveto{\pgfqpoint{0.000000in}{0.000000in}}%
\pgfpathlineto{\pgfqpoint{0.000000in}{0.020833in}}%
\pgfusepath{stroke,fill}%
}%
\begin{pgfscope}%
\pgfsys@transformshift{2.872731in}{0.420000in}%
\pgfsys@useobject{currentmarker}{}%
\end{pgfscope}%
\end{pgfscope}%
\begin{pgfscope}%
\pgfsetbuttcap%
\pgfsetroundjoin%
\definecolor{currentfill}{rgb}{0.000000,0.000000,0.000000}%
\pgfsetfillcolor{currentfill}%
\pgfsetlinewidth{0.401500pt}%
\definecolor{currentstroke}{rgb}{0.000000,0.000000,0.000000}%
\pgfsetstrokecolor{currentstroke}%
\pgfsetdash{}{0pt}%
\pgfsys@defobject{currentmarker}{\pgfqpoint{0.000000in}{-0.020833in}}{\pgfqpoint{0.000000in}{0.000000in}}{%
\pgfpathmoveto{\pgfqpoint{0.000000in}{0.000000in}}%
\pgfpathlineto{\pgfqpoint{0.000000in}{-0.020833in}}%
\pgfusepath{stroke,fill}%
}%
\begin{pgfscope}%
\pgfsys@transformshift{2.872731in}{2.414488in}%
\pgfsys@useobject{currentmarker}{}%
\end{pgfscope}%
\end{pgfscope}%
\begin{pgfscope}%
\pgfsetbuttcap%
\pgfsetroundjoin%
\definecolor{currentfill}{rgb}{0.000000,0.000000,0.000000}%
\pgfsetfillcolor{currentfill}%
\pgfsetlinewidth{0.401500pt}%
\definecolor{currentstroke}{rgb}{0.000000,0.000000,0.000000}%
\pgfsetstrokecolor{currentstroke}%
\pgfsetdash{}{0pt}%
\pgfsys@defobject{currentmarker}{\pgfqpoint{0.000000in}{0.000000in}}{\pgfqpoint{0.000000in}{0.020833in}}{%
\pgfpathmoveto{\pgfqpoint{0.000000in}{0.000000in}}%
\pgfpathlineto{\pgfqpoint{0.000000in}{0.020833in}}%
\pgfusepath{stroke,fill}%
}%
\begin{pgfscope}%
\pgfsys@transformshift{2.913702in}{0.420000in}%
\pgfsys@useobject{currentmarker}{}%
\end{pgfscope}%
\end{pgfscope}%
\begin{pgfscope}%
\pgfsetbuttcap%
\pgfsetroundjoin%
\definecolor{currentfill}{rgb}{0.000000,0.000000,0.000000}%
\pgfsetfillcolor{currentfill}%
\pgfsetlinewidth{0.401500pt}%
\definecolor{currentstroke}{rgb}{0.000000,0.000000,0.000000}%
\pgfsetstrokecolor{currentstroke}%
\pgfsetdash{}{0pt}%
\pgfsys@defobject{currentmarker}{\pgfqpoint{0.000000in}{-0.020833in}}{\pgfqpoint{0.000000in}{0.000000in}}{%
\pgfpathmoveto{\pgfqpoint{0.000000in}{0.000000in}}%
\pgfpathlineto{\pgfqpoint{0.000000in}{-0.020833in}}%
\pgfusepath{stroke,fill}%
}%
\begin{pgfscope}%
\pgfsys@transformshift{2.913702in}{2.414488in}%
\pgfsys@useobject{currentmarker}{}%
\end{pgfscope}%
\end{pgfscope}%
\begin{pgfscope}%
\pgfsetbuttcap%
\pgfsetroundjoin%
\definecolor{currentfill}{rgb}{0.000000,0.000000,0.000000}%
\pgfsetfillcolor{currentfill}%
\pgfsetlinewidth{0.401500pt}%
\definecolor{currentstroke}{rgb}{0.000000,0.000000,0.000000}%
\pgfsetstrokecolor{currentstroke}%
\pgfsetdash{}{0pt}%
\pgfsys@defobject{currentmarker}{\pgfqpoint{0.000000in}{0.000000in}}{\pgfqpoint{0.000000in}{0.020833in}}{%
\pgfpathmoveto{\pgfqpoint{0.000000in}{0.000000in}}%
\pgfpathlineto{\pgfqpoint{0.000000in}{0.020833in}}%
\pgfusepath{stroke,fill}%
}%
\begin{pgfscope}%
\pgfsys@transformshift{2.949841in}{0.420000in}%
\pgfsys@useobject{currentmarker}{}%
\end{pgfscope}%
\end{pgfscope}%
\begin{pgfscope}%
\pgfsetbuttcap%
\pgfsetroundjoin%
\definecolor{currentfill}{rgb}{0.000000,0.000000,0.000000}%
\pgfsetfillcolor{currentfill}%
\pgfsetlinewidth{0.401500pt}%
\definecolor{currentstroke}{rgb}{0.000000,0.000000,0.000000}%
\pgfsetstrokecolor{currentstroke}%
\pgfsetdash{}{0pt}%
\pgfsys@defobject{currentmarker}{\pgfqpoint{0.000000in}{-0.020833in}}{\pgfqpoint{0.000000in}{0.000000in}}{%
\pgfpathmoveto{\pgfqpoint{0.000000in}{0.000000in}}%
\pgfpathlineto{\pgfqpoint{0.000000in}{-0.020833in}}%
\pgfusepath{stroke,fill}%
}%
\begin{pgfscope}%
\pgfsys@transformshift{2.949841in}{2.414488in}%
\pgfsys@useobject{currentmarker}{}%
\end{pgfscope}%
\end{pgfscope}%
\begin{pgfscope}%
\definecolor{textcolor}{rgb}{0.000000,0.000000,0.000000}%
\pgfsetstrokecolor{textcolor}%
\pgfsetfillcolor{textcolor}%
\pgftext[x=1.844055in,y=0.208426in,,top]{\color{textcolor}{\rmfamily\fontsize{12.000000}{14.400000}\selectfont\catcode`\^=\active\def^{\ifmmode\sp\else\^{}\fi}\catcode`\%=\active\def%{\%}$y^{+}$}}%
\end{pgfscope}%
\begin{pgfscope}%
\pgfsetbuttcap%
\pgfsetroundjoin%
\definecolor{currentfill}{rgb}{0.000000,0.000000,0.000000}%
\pgfsetfillcolor{currentfill}%
\pgfsetlinewidth{0.501875pt}%
\definecolor{currentstroke}{rgb}{0.000000,0.000000,0.000000}%
\pgfsetstrokecolor{currentstroke}%
\pgfsetdash{}{0pt}%
\pgfsys@defobject{currentmarker}{\pgfqpoint{0.000000in}{0.000000in}}{\pgfqpoint{0.041667in}{0.000000in}}{%
\pgfpathmoveto{\pgfqpoint{0.000000in}{0.000000in}}%
\pgfpathlineto{\pgfqpoint{0.041667in}{0.000000in}}%
\pgfusepath{stroke,fill}%
}%
\begin{pgfscope}%
\pgfsys@transformshift{0.650000in}{0.648190in}%
\pgfsys@useobject{currentmarker}{}%
\end{pgfscope}%
\end{pgfscope}%
\begin{pgfscope}%
\pgfsetbuttcap%
\pgfsetroundjoin%
\definecolor{currentfill}{rgb}{0.000000,0.000000,0.000000}%
\pgfsetfillcolor{currentfill}%
\pgfsetlinewidth{0.501875pt}%
\definecolor{currentstroke}{rgb}{0.000000,0.000000,0.000000}%
\pgfsetstrokecolor{currentstroke}%
\pgfsetdash{}{0pt}%
\pgfsys@defobject{currentmarker}{\pgfqpoint{-0.041667in}{0.000000in}}{\pgfqpoint{-0.000000in}{0.000000in}}{%
\pgfpathmoveto{\pgfqpoint{-0.000000in}{0.000000in}}%
\pgfpathlineto{\pgfqpoint{-0.041667in}{0.000000in}}%
\pgfusepath{stroke,fill}%
}%
\begin{pgfscope}%
\pgfsys@transformshift{3.038110in}{0.648190in}%
\pgfsys@useobject{currentmarker}{}%
\end{pgfscope}%
\end{pgfscope}%
\begin{pgfscope}%
\definecolor{textcolor}{rgb}{0.000000,0.000000,0.000000}%
\pgfsetstrokecolor{textcolor}%
\pgfsetfillcolor{textcolor}%
\pgftext[x=0.288772in, y=0.595383in, left, base]{\color{textcolor}{\rmfamily\fontsize{11.000000}{13.200000}\selectfont\catcode`\^=\active\def^{\ifmmode\sp\else\^{}\fi}\catcode`\%=\active\def%{\%}\ensuremath{-}0.2}}%
\end{pgfscope}%
\begin{pgfscope}%
\pgfsetbuttcap%
\pgfsetroundjoin%
\definecolor{currentfill}{rgb}{0.000000,0.000000,0.000000}%
\pgfsetfillcolor{currentfill}%
\pgfsetlinewidth{0.501875pt}%
\definecolor{currentstroke}{rgb}{0.000000,0.000000,0.000000}%
\pgfsetstrokecolor{currentstroke}%
\pgfsetdash{}{0pt}%
\pgfsys@defobject{currentmarker}{\pgfqpoint{0.000000in}{0.000000in}}{\pgfqpoint{0.041667in}{0.000000in}}{%
\pgfpathmoveto{\pgfqpoint{0.000000in}{0.000000in}}%
\pgfpathlineto{\pgfqpoint{0.041667in}{0.000000in}}%
\pgfusepath{stroke,fill}%
}%
\begin{pgfscope}%
\pgfsys@transformshift{0.650000in}{1.021309in}%
\pgfsys@useobject{currentmarker}{}%
\end{pgfscope}%
\end{pgfscope}%
\begin{pgfscope}%
\pgfsetbuttcap%
\pgfsetroundjoin%
\definecolor{currentfill}{rgb}{0.000000,0.000000,0.000000}%
\pgfsetfillcolor{currentfill}%
\pgfsetlinewidth{0.501875pt}%
\definecolor{currentstroke}{rgb}{0.000000,0.000000,0.000000}%
\pgfsetstrokecolor{currentstroke}%
\pgfsetdash{}{0pt}%
\pgfsys@defobject{currentmarker}{\pgfqpoint{-0.041667in}{0.000000in}}{\pgfqpoint{-0.000000in}{0.000000in}}{%
\pgfpathmoveto{\pgfqpoint{-0.000000in}{0.000000in}}%
\pgfpathlineto{\pgfqpoint{-0.041667in}{0.000000in}}%
\pgfusepath{stroke,fill}%
}%
\begin{pgfscope}%
\pgfsys@transformshift{3.038110in}{1.021309in}%
\pgfsys@useobject{currentmarker}{}%
\end{pgfscope}%
\end{pgfscope}%
\begin{pgfscope}%
\definecolor{textcolor}{rgb}{0.000000,0.000000,0.000000}%
\pgfsetstrokecolor{textcolor}%
\pgfsetfillcolor{textcolor}%
\pgftext[x=0.288772in, y=0.968503in, left, base]{\color{textcolor}{\rmfamily\fontsize{11.000000}{13.200000}\selectfont\catcode`\^=\active\def^{\ifmmode\sp\else\^{}\fi}\catcode`\%=\active\def%{\%}\ensuremath{-}0.1}}%
\end{pgfscope}%
\begin{pgfscope}%
\pgfsetbuttcap%
\pgfsetroundjoin%
\definecolor{currentfill}{rgb}{0.000000,0.000000,0.000000}%
\pgfsetfillcolor{currentfill}%
\pgfsetlinewidth{0.501875pt}%
\definecolor{currentstroke}{rgb}{0.000000,0.000000,0.000000}%
\pgfsetstrokecolor{currentstroke}%
\pgfsetdash{}{0pt}%
\pgfsys@defobject{currentmarker}{\pgfqpoint{0.000000in}{0.000000in}}{\pgfqpoint{0.041667in}{0.000000in}}{%
\pgfpathmoveto{\pgfqpoint{0.000000in}{0.000000in}}%
\pgfpathlineto{\pgfqpoint{0.041667in}{0.000000in}}%
\pgfusepath{stroke,fill}%
}%
\begin{pgfscope}%
\pgfsys@transformshift{0.650000in}{1.394429in}%
\pgfsys@useobject{currentmarker}{}%
\end{pgfscope}%
\end{pgfscope}%
\begin{pgfscope}%
\pgfsetbuttcap%
\pgfsetroundjoin%
\definecolor{currentfill}{rgb}{0.000000,0.000000,0.000000}%
\pgfsetfillcolor{currentfill}%
\pgfsetlinewidth{0.501875pt}%
\definecolor{currentstroke}{rgb}{0.000000,0.000000,0.000000}%
\pgfsetstrokecolor{currentstroke}%
\pgfsetdash{}{0pt}%
\pgfsys@defobject{currentmarker}{\pgfqpoint{-0.041667in}{0.000000in}}{\pgfqpoint{-0.000000in}{0.000000in}}{%
\pgfpathmoveto{\pgfqpoint{-0.000000in}{0.000000in}}%
\pgfpathlineto{\pgfqpoint{-0.041667in}{0.000000in}}%
\pgfusepath{stroke,fill}%
}%
\begin{pgfscope}%
\pgfsys@transformshift{3.038110in}{1.394429in}%
\pgfsys@useobject{currentmarker}{}%
\end{pgfscope}%
\end{pgfscope}%
\begin{pgfscope}%
\definecolor{textcolor}{rgb}{0.000000,0.000000,0.000000}%
\pgfsetstrokecolor{textcolor}%
\pgfsetfillcolor{textcolor}%
\pgftext[x=0.407060in, y=1.341622in, left, base]{\color{textcolor}{\rmfamily\fontsize{11.000000}{13.200000}\selectfont\catcode`\^=\active\def^{\ifmmode\sp\else\^{}\fi}\catcode`\%=\active\def%{\%}0.0}}%
\end{pgfscope}%
\begin{pgfscope}%
\pgfsetbuttcap%
\pgfsetroundjoin%
\definecolor{currentfill}{rgb}{0.000000,0.000000,0.000000}%
\pgfsetfillcolor{currentfill}%
\pgfsetlinewidth{0.501875pt}%
\definecolor{currentstroke}{rgb}{0.000000,0.000000,0.000000}%
\pgfsetstrokecolor{currentstroke}%
\pgfsetdash{}{0pt}%
\pgfsys@defobject{currentmarker}{\pgfqpoint{0.000000in}{0.000000in}}{\pgfqpoint{0.041667in}{0.000000in}}{%
\pgfpathmoveto{\pgfqpoint{0.000000in}{0.000000in}}%
\pgfpathlineto{\pgfqpoint{0.041667in}{0.000000in}}%
\pgfusepath{stroke,fill}%
}%
\begin{pgfscope}%
\pgfsys@transformshift{0.650000in}{1.767549in}%
\pgfsys@useobject{currentmarker}{}%
\end{pgfscope}%
\end{pgfscope}%
\begin{pgfscope}%
\pgfsetbuttcap%
\pgfsetroundjoin%
\definecolor{currentfill}{rgb}{0.000000,0.000000,0.000000}%
\pgfsetfillcolor{currentfill}%
\pgfsetlinewidth{0.501875pt}%
\definecolor{currentstroke}{rgb}{0.000000,0.000000,0.000000}%
\pgfsetstrokecolor{currentstroke}%
\pgfsetdash{}{0pt}%
\pgfsys@defobject{currentmarker}{\pgfqpoint{-0.041667in}{0.000000in}}{\pgfqpoint{-0.000000in}{0.000000in}}{%
\pgfpathmoveto{\pgfqpoint{-0.000000in}{0.000000in}}%
\pgfpathlineto{\pgfqpoint{-0.041667in}{0.000000in}}%
\pgfusepath{stroke,fill}%
}%
\begin{pgfscope}%
\pgfsys@transformshift{3.038110in}{1.767549in}%
\pgfsys@useobject{currentmarker}{}%
\end{pgfscope}%
\end{pgfscope}%
\begin{pgfscope}%
\definecolor{textcolor}{rgb}{0.000000,0.000000,0.000000}%
\pgfsetstrokecolor{textcolor}%
\pgfsetfillcolor{textcolor}%
\pgftext[x=0.407060in, y=1.714742in, left, base]{\color{textcolor}{\rmfamily\fontsize{11.000000}{13.200000}\selectfont\catcode`\^=\active\def^{\ifmmode\sp\else\^{}\fi}\catcode`\%=\active\def%{\%}0.1}}%
\end{pgfscope}%
\begin{pgfscope}%
\pgfsetbuttcap%
\pgfsetroundjoin%
\definecolor{currentfill}{rgb}{0.000000,0.000000,0.000000}%
\pgfsetfillcolor{currentfill}%
\pgfsetlinewidth{0.501875pt}%
\definecolor{currentstroke}{rgb}{0.000000,0.000000,0.000000}%
\pgfsetstrokecolor{currentstroke}%
\pgfsetdash{}{0pt}%
\pgfsys@defobject{currentmarker}{\pgfqpoint{0.000000in}{0.000000in}}{\pgfqpoint{0.041667in}{0.000000in}}{%
\pgfpathmoveto{\pgfqpoint{0.000000in}{0.000000in}}%
\pgfpathlineto{\pgfqpoint{0.041667in}{0.000000in}}%
\pgfusepath{stroke,fill}%
}%
\begin{pgfscope}%
\pgfsys@transformshift{0.650000in}{2.140668in}%
\pgfsys@useobject{currentmarker}{}%
\end{pgfscope}%
\end{pgfscope}%
\begin{pgfscope}%
\pgfsetbuttcap%
\pgfsetroundjoin%
\definecolor{currentfill}{rgb}{0.000000,0.000000,0.000000}%
\pgfsetfillcolor{currentfill}%
\pgfsetlinewidth{0.501875pt}%
\definecolor{currentstroke}{rgb}{0.000000,0.000000,0.000000}%
\pgfsetstrokecolor{currentstroke}%
\pgfsetdash{}{0pt}%
\pgfsys@defobject{currentmarker}{\pgfqpoint{-0.041667in}{0.000000in}}{\pgfqpoint{-0.000000in}{0.000000in}}{%
\pgfpathmoveto{\pgfqpoint{-0.000000in}{0.000000in}}%
\pgfpathlineto{\pgfqpoint{-0.041667in}{0.000000in}}%
\pgfusepath{stroke,fill}%
}%
\begin{pgfscope}%
\pgfsys@transformshift{3.038110in}{2.140668in}%
\pgfsys@useobject{currentmarker}{}%
\end{pgfscope}%
\end{pgfscope}%
\begin{pgfscope}%
\definecolor{textcolor}{rgb}{0.000000,0.000000,0.000000}%
\pgfsetstrokecolor{textcolor}%
\pgfsetfillcolor{textcolor}%
\pgftext[x=0.407060in, y=2.087861in, left, base]{\color{textcolor}{\rmfamily\fontsize{11.000000}{13.200000}\selectfont\catcode`\^=\active\def^{\ifmmode\sp\else\^{}\fi}\catcode`\%=\active\def%{\%}0.2}}%
\end{pgfscope}%
\begin{pgfscope}%
\pgfsetbuttcap%
\pgfsetroundjoin%
\definecolor{currentfill}{rgb}{0.000000,0.000000,0.000000}%
\pgfsetfillcolor{currentfill}%
\pgfsetlinewidth{0.401500pt}%
\definecolor{currentstroke}{rgb}{0.000000,0.000000,0.000000}%
\pgfsetstrokecolor{currentstroke}%
\pgfsetdash{}{0pt}%
\pgfsys@defobject{currentmarker}{\pgfqpoint{0.000000in}{0.000000in}}{\pgfqpoint{0.020833in}{0.000000in}}{%
\pgfpathmoveto{\pgfqpoint{0.000000in}{0.000000in}}%
\pgfpathlineto{\pgfqpoint{0.020833in}{0.000000in}}%
\pgfusepath{stroke,fill}%
}%
\begin{pgfscope}%
\pgfsys@transformshift{0.650000in}{0.424318in}%
\pgfsys@useobject{currentmarker}{}%
\end{pgfscope}%
\end{pgfscope}%
\begin{pgfscope}%
\pgfsetbuttcap%
\pgfsetroundjoin%
\definecolor{currentfill}{rgb}{0.000000,0.000000,0.000000}%
\pgfsetfillcolor{currentfill}%
\pgfsetlinewidth{0.401500pt}%
\definecolor{currentstroke}{rgb}{0.000000,0.000000,0.000000}%
\pgfsetstrokecolor{currentstroke}%
\pgfsetdash{}{0pt}%
\pgfsys@defobject{currentmarker}{\pgfqpoint{-0.020833in}{0.000000in}}{\pgfqpoint{-0.000000in}{0.000000in}}{%
\pgfpathmoveto{\pgfqpoint{-0.000000in}{0.000000in}}%
\pgfpathlineto{\pgfqpoint{-0.020833in}{0.000000in}}%
\pgfusepath{stroke,fill}%
}%
\begin{pgfscope}%
\pgfsys@transformshift{3.038110in}{0.424318in}%
\pgfsys@useobject{currentmarker}{}%
\end{pgfscope}%
\end{pgfscope}%
\begin{pgfscope}%
\pgfsetbuttcap%
\pgfsetroundjoin%
\definecolor{currentfill}{rgb}{0.000000,0.000000,0.000000}%
\pgfsetfillcolor{currentfill}%
\pgfsetlinewidth{0.401500pt}%
\definecolor{currentstroke}{rgb}{0.000000,0.000000,0.000000}%
\pgfsetstrokecolor{currentstroke}%
\pgfsetdash{}{0pt}%
\pgfsys@defobject{currentmarker}{\pgfqpoint{0.000000in}{0.000000in}}{\pgfqpoint{0.020833in}{0.000000in}}{%
\pgfpathmoveto{\pgfqpoint{0.000000in}{0.000000in}}%
\pgfpathlineto{\pgfqpoint{0.020833in}{0.000000in}}%
\pgfusepath{stroke,fill}%
}%
\begin{pgfscope}%
\pgfsys@transformshift{0.650000in}{0.498942in}%
\pgfsys@useobject{currentmarker}{}%
\end{pgfscope}%
\end{pgfscope}%
\begin{pgfscope}%
\pgfsetbuttcap%
\pgfsetroundjoin%
\definecolor{currentfill}{rgb}{0.000000,0.000000,0.000000}%
\pgfsetfillcolor{currentfill}%
\pgfsetlinewidth{0.401500pt}%
\definecolor{currentstroke}{rgb}{0.000000,0.000000,0.000000}%
\pgfsetstrokecolor{currentstroke}%
\pgfsetdash{}{0pt}%
\pgfsys@defobject{currentmarker}{\pgfqpoint{-0.020833in}{0.000000in}}{\pgfqpoint{-0.000000in}{0.000000in}}{%
\pgfpathmoveto{\pgfqpoint{-0.000000in}{0.000000in}}%
\pgfpathlineto{\pgfqpoint{-0.020833in}{0.000000in}}%
\pgfusepath{stroke,fill}%
}%
\begin{pgfscope}%
\pgfsys@transformshift{3.038110in}{0.498942in}%
\pgfsys@useobject{currentmarker}{}%
\end{pgfscope}%
\end{pgfscope}%
\begin{pgfscope}%
\pgfsetbuttcap%
\pgfsetroundjoin%
\definecolor{currentfill}{rgb}{0.000000,0.000000,0.000000}%
\pgfsetfillcolor{currentfill}%
\pgfsetlinewidth{0.401500pt}%
\definecolor{currentstroke}{rgb}{0.000000,0.000000,0.000000}%
\pgfsetstrokecolor{currentstroke}%
\pgfsetdash{}{0pt}%
\pgfsys@defobject{currentmarker}{\pgfqpoint{0.000000in}{0.000000in}}{\pgfqpoint{0.020833in}{0.000000in}}{%
\pgfpathmoveto{\pgfqpoint{0.000000in}{0.000000in}}%
\pgfpathlineto{\pgfqpoint{0.020833in}{0.000000in}}%
\pgfusepath{stroke,fill}%
}%
\begin{pgfscope}%
\pgfsys@transformshift{0.650000in}{0.573566in}%
\pgfsys@useobject{currentmarker}{}%
\end{pgfscope}%
\end{pgfscope}%
\begin{pgfscope}%
\pgfsetbuttcap%
\pgfsetroundjoin%
\definecolor{currentfill}{rgb}{0.000000,0.000000,0.000000}%
\pgfsetfillcolor{currentfill}%
\pgfsetlinewidth{0.401500pt}%
\definecolor{currentstroke}{rgb}{0.000000,0.000000,0.000000}%
\pgfsetstrokecolor{currentstroke}%
\pgfsetdash{}{0pt}%
\pgfsys@defobject{currentmarker}{\pgfqpoint{-0.020833in}{0.000000in}}{\pgfqpoint{-0.000000in}{0.000000in}}{%
\pgfpathmoveto{\pgfqpoint{-0.000000in}{0.000000in}}%
\pgfpathlineto{\pgfqpoint{-0.020833in}{0.000000in}}%
\pgfusepath{stroke,fill}%
}%
\begin{pgfscope}%
\pgfsys@transformshift{3.038110in}{0.573566in}%
\pgfsys@useobject{currentmarker}{}%
\end{pgfscope}%
\end{pgfscope}%
\begin{pgfscope}%
\pgfsetbuttcap%
\pgfsetroundjoin%
\definecolor{currentfill}{rgb}{0.000000,0.000000,0.000000}%
\pgfsetfillcolor{currentfill}%
\pgfsetlinewidth{0.401500pt}%
\definecolor{currentstroke}{rgb}{0.000000,0.000000,0.000000}%
\pgfsetstrokecolor{currentstroke}%
\pgfsetdash{}{0pt}%
\pgfsys@defobject{currentmarker}{\pgfqpoint{0.000000in}{0.000000in}}{\pgfqpoint{0.020833in}{0.000000in}}{%
\pgfpathmoveto{\pgfqpoint{0.000000in}{0.000000in}}%
\pgfpathlineto{\pgfqpoint{0.020833in}{0.000000in}}%
\pgfusepath{stroke,fill}%
}%
\begin{pgfscope}%
\pgfsys@transformshift{0.650000in}{0.722814in}%
\pgfsys@useobject{currentmarker}{}%
\end{pgfscope}%
\end{pgfscope}%
\begin{pgfscope}%
\pgfsetbuttcap%
\pgfsetroundjoin%
\definecolor{currentfill}{rgb}{0.000000,0.000000,0.000000}%
\pgfsetfillcolor{currentfill}%
\pgfsetlinewidth{0.401500pt}%
\definecolor{currentstroke}{rgb}{0.000000,0.000000,0.000000}%
\pgfsetstrokecolor{currentstroke}%
\pgfsetdash{}{0pt}%
\pgfsys@defobject{currentmarker}{\pgfqpoint{-0.020833in}{0.000000in}}{\pgfqpoint{-0.000000in}{0.000000in}}{%
\pgfpathmoveto{\pgfqpoint{-0.000000in}{0.000000in}}%
\pgfpathlineto{\pgfqpoint{-0.020833in}{0.000000in}}%
\pgfusepath{stroke,fill}%
}%
\begin{pgfscope}%
\pgfsys@transformshift{3.038110in}{0.722814in}%
\pgfsys@useobject{currentmarker}{}%
\end{pgfscope}%
\end{pgfscope}%
\begin{pgfscope}%
\pgfsetbuttcap%
\pgfsetroundjoin%
\definecolor{currentfill}{rgb}{0.000000,0.000000,0.000000}%
\pgfsetfillcolor{currentfill}%
\pgfsetlinewidth{0.401500pt}%
\definecolor{currentstroke}{rgb}{0.000000,0.000000,0.000000}%
\pgfsetstrokecolor{currentstroke}%
\pgfsetdash{}{0pt}%
\pgfsys@defobject{currentmarker}{\pgfqpoint{0.000000in}{0.000000in}}{\pgfqpoint{0.020833in}{0.000000in}}{%
\pgfpathmoveto{\pgfqpoint{0.000000in}{0.000000in}}%
\pgfpathlineto{\pgfqpoint{0.020833in}{0.000000in}}%
\pgfusepath{stroke,fill}%
}%
\begin{pgfscope}%
\pgfsys@transformshift{0.650000in}{0.797438in}%
\pgfsys@useobject{currentmarker}{}%
\end{pgfscope}%
\end{pgfscope}%
\begin{pgfscope}%
\pgfsetbuttcap%
\pgfsetroundjoin%
\definecolor{currentfill}{rgb}{0.000000,0.000000,0.000000}%
\pgfsetfillcolor{currentfill}%
\pgfsetlinewidth{0.401500pt}%
\definecolor{currentstroke}{rgb}{0.000000,0.000000,0.000000}%
\pgfsetstrokecolor{currentstroke}%
\pgfsetdash{}{0pt}%
\pgfsys@defobject{currentmarker}{\pgfqpoint{-0.020833in}{0.000000in}}{\pgfqpoint{-0.000000in}{0.000000in}}{%
\pgfpathmoveto{\pgfqpoint{-0.000000in}{0.000000in}}%
\pgfpathlineto{\pgfqpoint{-0.020833in}{0.000000in}}%
\pgfusepath{stroke,fill}%
}%
\begin{pgfscope}%
\pgfsys@transformshift{3.038110in}{0.797438in}%
\pgfsys@useobject{currentmarker}{}%
\end{pgfscope}%
\end{pgfscope}%
\begin{pgfscope}%
\pgfsetbuttcap%
\pgfsetroundjoin%
\definecolor{currentfill}{rgb}{0.000000,0.000000,0.000000}%
\pgfsetfillcolor{currentfill}%
\pgfsetlinewidth{0.401500pt}%
\definecolor{currentstroke}{rgb}{0.000000,0.000000,0.000000}%
\pgfsetstrokecolor{currentstroke}%
\pgfsetdash{}{0pt}%
\pgfsys@defobject{currentmarker}{\pgfqpoint{0.000000in}{0.000000in}}{\pgfqpoint{0.020833in}{0.000000in}}{%
\pgfpathmoveto{\pgfqpoint{0.000000in}{0.000000in}}%
\pgfpathlineto{\pgfqpoint{0.020833in}{0.000000in}}%
\pgfusepath{stroke,fill}%
}%
\begin{pgfscope}%
\pgfsys@transformshift{0.650000in}{0.872061in}%
\pgfsys@useobject{currentmarker}{}%
\end{pgfscope}%
\end{pgfscope}%
\begin{pgfscope}%
\pgfsetbuttcap%
\pgfsetroundjoin%
\definecolor{currentfill}{rgb}{0.000000,0.000000,0.000000}%
\pgfsetfillcolor{currentfill}%
\pgfsetlinewidth{0.401500pt}%
\definecolor{currentstroke}{rgb}{0.000000,0.000000,0.000000}%
\pgfsetstrokecolor{currentstroke}%
\pgfsetdash{}{0pt}%
\pgfsys@defobject{currentmarker}{\pgfqpoint{-0.020833in}{0.000000in}}{\pgfqpoint{-0.000000in}{0.000000in}}{%
\pgfpathmoveto{\pgfqpoint{-0.000000in}{0.000000in}}%
\pgfpathlineto{\pgfqpoint{-0.020833in}{0.000000in}}%
\pgfusepath{stroke,fill}%
}%
\begin{pgfscope}%
\pgfsys@transformshift{3.038110in}{0.872061in}%
\pgfsys@useobject{currentmarker}{}%
\end{pgfscope}%
\end{pgfscope}%
\begin{pgfscope}%
\pgfsetbuttcap%
\pgfsetroundjoin%
\definecolor{currentfill}{rgb}{0.000000,0.000000,0.000000}%
\pgfsetfillcolor{currentfill}%
\pgfsetlinewidth{0.401500pt}%
\definecolor{currentstroke}{rgb}{0.000000,0.000000,0.000000}%
\pgfsetstrokecolor{currentstroke}%
\pgfsetdash{}{0pt}%
\pgfsys@defobject{currentmarker}{\pgfqpoint{0.000000in}{0.000000in}}{\pgfqpoint{0.020833in}{0.000000in}}{%
\pgfpathmoveto{\pgfqpoint{0.000000in}{0.000000in}}%
\pgfpathlineto{\pgfqpoint{0.020833in}{0.000000in}}%
\pgfusepath{stroke,fill}%
}%
\begin{pgfscope}%
\pgfsys@transformshift{0.650000in}{0.946685in}%
\pgfsys@useobject{currentmarker}{}%
\end{pgfscope}%
\end{pgfscope}%
\begin{pgfscope}%
\pgfsetbuttcap%
\pgfsetroundjoin%
\definecolor{currentfill}{rgb}{0.000000,0.000000,0.000000}%
\pgfsetfillcolor{currentfill}%
\pgfsetlinewidth{0.401500pt}%
\definecolor{currentstroke}{rgb}{0.000000,0.000000,0.000000}%
\pgfsetstrokecolor{currentstroke}%
\pgfsetdash{}{0pt}%
\pgfsys@defobject{currentmarker}{\pgfqpoint{-0.020833in}{0.000000in}}{\pgfqpoint{-0.000000in}{0.000000in}}{%
\pgfpathmoveto{\pgfqpoint{-0.000000in}{0.000000in}}%
\pgfpathlineto{\pgfqpoint{-0.020833in}{0.000000in}}%
\pgfusepath{stroke,fill}%
}%
\begin{pgfscope}%
\pgfsys@transformshift{3.038110in}{0.946685in}%
\pgfsys@useobject{currentmarker}{}%
\end{pgfscope}%
\end{pgfscope}%
\begin{pgfscope}%
\pgfsetbuttcap%
\pgfsetroundjoin%
\definecolor{currentfill}{rgb}{0.000000,0.000000,0.000000}%
\pgfsetfillcolor{currentfill}%
\pgfsetlinewidth{0.401500pt}%
\definecolor{currentstroke}{rgb}{0.000000,0.000000,0.000000}%
\pgfsetstrokecolor{currentstroke}%
\pgfsetdash{}{0pt}%
\pgfsys@defobject{currentmarker}{\pgfqpoint{0.000000in}{0.000000in}}{\pgfqpoint{0.020833in}{0.000000in}}{%
\pgfpathmoveto{\pgfqpoint{0.000000in}{0.000000in}}%
\pgfpathlineto{\pgfqpoint{0.020833in}{0.000000in}}%
\pgfusepath{stroke,fill}%
}%
\begin{pgfscope}%
\pgfsys@transformshift{0.650000in}{1.095933in}%
\pgfsys@useobject{currentmarker}{}%
\end{pgfscope}%
\end{pgfscope}%
\begin{pgfscope}%
\pgfsetbuttcap%
\pgfsetroundjoin%
\definecolor{currentfill}{rgb}{0.000000,0.000000,0.000000}%
\pgfsetfillcolor{currentfill}%
\pgfsetlinewidth{0.401500pt}%
\definecolor{currentstroke}{rgb}{0.000000,0.000000,0.000000}%
\pgfsetstrokecolor{currentstroke}%
\pgfsetdash{}{0pt}%
\pgfsys@defobject{currentmarker}{\pgfqpoint{-0.020833in}{0.000000in}}{\pgfqpoint{-0.000000in}{0.000000in}}{%
\pgfpathmoveto{\pgfqpoint{-0.000000in}{0.000000in}}%
\pgfpathlineto{\pgfqpoint{-0.020833in}{0.000000in}}%
\pgfusepath{stroke,fill}%
}%
\begin{pgfscope}%
\pgfsys@transformshift{3.038110in}{1.095933in}%
\pgfsys@useobject{currentmarker}{}%
\end{pgfscope}%
\end{pgfscope}%
\begin{pgfscope}%
\pgfsetbuttcap%
\pgfsetroundjoin%
\definecolor{currentfill}{rgb}{0.000000,0.000000,0.000000}%
\pgfsetfillcolor{currentfill}%
\pgfsetlinewidth{0.401500pt}%
\definecolor{currentstroke}{rgb}{0.000000,0.000000,0.000000}%
\pgfsetstrokecolor{currentstroke}%
\pgfsetdash{}{0pt}%
\pgfsys@defobject{currentmarker}{\pgfqpoint{0.000000in}{0.000000in}}{\pgfqpoint{0.020833in}{0.000000in}}{%
\pgfpathmoveto{\pgfqpoint{0.000000in}{0.000000in}}%
\pgfpathlineto{\pgfqpoint{0.020833in}{0.000000in}}%
\pgfusepath{stroke,fill}%
}%
\begin{pgfscope}%
\pgfsys@transformshift{0.650000in}{1.170557in}%
\pgfsys@useobject{currentmarker}{}%
\end{pgfscope}%
\end{pgfscope}%
\begin{pgfscope}%
\pgfsetbuttcap%
\pgfsetroundjoin%
\definecolor{currentfill}{rgb}{0.000000,0.000000,0.000000}%
\pgfsetfillcolor{currentfill}%
\pgfsetlinewidth{0.401500pt}%
\definecolor{currentstroke}{rgb}{0.000000,0.000000,0.000000}%
\pgfsetstrokecolor{currentstroke}%
\pgfsetdash{}{0pt}%
\pgfsys@defobject{currentmarker}{\pgfqpoint{-0.020833in}{0.000000in}}{\pgfqpoint{-0.000000in}{0.000000in}}{%
\pgfpathmoveto{\pgfqpoint{-0.000000in}{0.000000in}}%
\pgfpathlineto{\pgfqpoint{-0.020833in}{0.000000in}}%
\pgfusepath{stroke,fill}%
}%
\begin{pgfscope}%
\pgfsys@transformshift{3.038110in}{1.170557in}%
\pgfsys@useobject{currentmarker}{}%
\end{pgfscope}%
\end{pgfscope}%
\begin{pgfscope}%
\pgfsetbuttcap%
\pgfsetroundjoin%
\definecolor{currentfill}{rgb}{0.000000,0.000000,0.000000}%
\pgfsetfillcolor{currentfill}%
\pgfsetlinewidth{0.401500pt}%
\definecolor{currentstroke}{rgb}{0.000000,0.000000,0.000000}%
\pgfsetstrokecolor{currentstroke}%
\pgfsetdash{}{0pt}%
\pgfsys@defobject{currentmarker}{\pgfqpoint{0.000000in}{0.000000in}}{\pgfqpoint{0.020833in}{0.000000in}}{%
\pgfpathmoveto{\pgfqpoint{0.000000in}{0.000000in}}%
\pgfpathlineto{\pgfqpoint{0.020833in}{0.000000in}}%
\pgfusepath{stroke,fill}%
}%
\begin{pgfscope}%
\pgfsys@transformshift{0.650000in}{1.245181in}%
\pgfsys@useobject{currentmarker}{}%
\end{pgfscope}%
\end{pgfscope}%
\begin{pgfscope}%
\pgfsetbuttcap%
\pgfsetroundjoin%
\definecolor{currentfill}{rgb}{0.000000,0.000000,0.000000}%
\pgfsetfillcolor{currentfill}%
\pgfsetlinewidth{0.401500pt}%
\definecolor{currentstroke}{rgb}{0.000000,0.000000,0.000000}%
\pgfsetstrokecolor{currentstroke}%
\pgfsetdash{}{0pt}%
\pgfsys@defobject{currentmarker}{\pgfqpoint{-0.020833in}{0.000000in}}{\pgfqpoint{-0.000000in}{0.000000in}}{%
\pgfpathmoveto{\pgfqpoint{-0.000000in}{0.000000in}}%
\pgfpathlineto{\pgfqpoint{-0.020833in}{0.000000in}}%
\pgfusepath{stroke,fill}%
}%
\begin{pgfscope}%
\pgfsys@transformshift{3.038110in}{1.245181in}%
\pgfsys@useobject{currentmarker}{}%
\end{pgfscope}%
\end{pgfscope}%
\begin{pgfscope}%
\pgfsetbuttcap%
\pgfsetroundjoin%
\definecolor{currentfill}{rgb}{0.000000,0.000000,0.000000}%
\pgfsetfillcolor{currentfill}%
\pgfsetlinewidth{0.401500pt}%
\definecolor{currentstroke}{rgb}{0.000000,0.000000,0.000000}%
\pgfsetstrokecolor{currentstroke}%
\pgfsetdash{}{0pt}%
\pgfsys@defobject{currentmarker}{\pgfqpoint{0.000000in}{0.000000in}}{\pgfqpoint{0.020833in}{0.000000in}}{%
\pgfpathmoveto{\pgfqpoint{0.000000in}{0.000000in}}%
\pgfpathlineto{\pgfqpoint{0.020833in}{0.000000in}}%
\pgfusepath{stroke,fill}%
}%
\begin{pgfscope}%
\pgfsys@transformshift{0.650000in}{1.319805in}%
\pgfsys@useobject{currentmarker}{}%
\end{pgfscope}%
\end{pgfscope}%
\begin{pgfscope}%
\pgfsetbuttcap%
\pgfsetroundjoin%
\definecolor{currentfill}{rgb}{0.000000,0.000000,0.000000}%
\pgfsetfillcolor{currentfill}%
\pgfsetlinewidth{0.401500pt}%
\definecolor{currentstroke}{rgb}{0.000000,0.000000,0.000000}%
\pgfsetstrokecolor{currentstroke}%
\pgfsetdash{}{0pt}%
\pgfsys@defobject{currentmarker}{\pgfqpoint{-0.020833in}{0.000000in}}{\pgfqpoint{-0.000000in}{0.000000in}}{%
\pgfpathmoveto{\pgfqpoint{-0.000000in}{0.000000in}}%
\pgfpathlineto{\pgfqpoint{-0.020833in}{0.000000in}}%
\pgfusepath{stroke,fill}%
}%
\begin{pgfscope}%
\pgfsys@transformshift{3.038110in}{1.319805in}%
\pgfsys@useobject{currentmarker}{}%
\end{pgfscope}%
\end{pgfscope}%
\begin{pgfscope}%
\pgfsetbuttcap%
\pgfsetroundjoin%
\definecolor{currentfill}{rgb}{0.000000,0.000000,0.000000}%
\pgfsetfillcolor{currentfill}%
\pgfsetlinewidth{0.401500pt}%
\definecolor{currentstroke}{rgb}{0.000000,0.000000,0.000000}%
\pgfsetstrokecolor{currentstroke}%
\pgfsetdash{}{0pt}%
\pgfsys@defobject{currentmarker}{\pgfqpoint{0.000000in}{0.000000in}}{\pgfqpoint{0.020833in}{0.000000in}}{%
\pgfpathmoveto{\pgfqpoint{0.000000in}{0.000000in}}%
\pgfpathlineto{\pgfqpoint{0.020833in}{0.000000in}}%
\pgfusepath{stroke,fill}%
}%
\begin{pgfscope}%
\pgfsys@transformshift{0.650000in}{1.469053in}%
\pgfsys@useobject{currentmarker}{}%
\end{pgfscope}%
\end{pgfscope}%
\begin{pgfscope}%
\pgfsetbuttcap%
\pgfsetroundjoin%
\definecolor{currentfill}{rgb}{0.000000,0.000000,0.000000}%
\pgfsetfillcolor{currentfill}%
\pgfsetlinewidth{0.401500pt}%
\definecolor{currentstroke}{rgb}{0.000000,0.000000,0.000000}%
\pgfsetstrokecolor{currentstroke}%
\pgfsetdash{}{0pt}%
\pgfsys@defobject{currentmarker}{\pgfqpoint{-0.020833in}{0.000000in}}{\pgfqpoint{-0.000000in}{0.000000in}}{%
\pgfpathmoveto{\pgfqpoint{-0.000000in}{0.000000in}}%
\pgfpathlineto{\pgfqpoint{-0.020833in}{0.000000in}}%
\pgfusepath{stroke,fill}%
}%
\begin{pgfscope}%
\pgfsys@transformshift{3.038110in}{1.469053in}%
\pgfsys@useobject{currentmarker}{}%
\end{pgfscope}%
\end{pgfscope}%
\begin{pgfscope}%
\pgfsetbuttcap%
\pgfsetroundjoin%
\definecolor{currentfill}{rgb}{0.000000,0.000000,0.000000}%
\pgfsetfillcolor{currentfill}%
\pgfsetlinewidth{0.401500pt}%
\definecolor{currentstroke}{rgb}{0.000000,0.000000,0.000000}%
\pgfsetstrokecolor{currentstroke}%
\pgfsetdash{}{0pt}%
\pgfsys@defobject{currentmarker}{\pgfqpoint{0.000000in}{0.000000in}}{\pgfqpoint{0.020833in}{0.000000in}}{%
\pgfpathmoveto{\pgfqpoint{0.000000in}{0.000000in}}%
\pgfpathlineto{\pgfqpoint{0.020833in}{0.000000in}}%
\pgfusepath{stroke,fill}%
}%
\begin{pgfscope}%
\pgfsys@transformshift{0.650000in}{1.543677in}%
\pgfsys@useobject{currentmarker}{}%
\end{pgfscope}%
\end{pgfscope}%
\begin{pgfscope}%
\pgfsetbuttcap%
\pgfsetroundjoin%
\definecolor{currentfill}{rgb}{0.000000,0.000000,0.000000}%
\pgfsetfillcolor{currentfill}%
\pgfsetlinewidth{0.401500pt}%
\definecolor{currentstroke}{rgb}{0.000000,0.000000,0.000000}%
\pgfsetstrokecolor{currentstroke}%
\pgfsetdash{}{0pt}%
\pgfsys@defobject{currentmarker}{\pgfqpoint{-0.020833in}{0.000000in}}{\pgfqpoint{-0.000000in}{0.000000in}}{%
\pgfpathmoveto{\pgfqpoint{-0.000000in}{0.000000in}}%
\pgfpathlineto{\pgfqpoint{-0.020833in}{0.000000in}}%
\pgfusepath{stroke,fill}%
}%
\begin{pgfscope}%
\pgfsys@transformshift{3.038110in}{1.543677in}%
\pgfsys@useobject{currentmarker}{}%
\end{pgfscope}%
\end{pgfscope}%
\begin{pgfscope}%
\pgfsetbuttcap%
\pgfsetroundjoin%
\definecolor{currentfill}{rgb}{0.000000,0.000000,0.000000}%
\pgfsetfillcolor{currentfill}%
\pgfsetlinewidth{0.401500pt}%
\definecolor{currentstroke}{rgb}{0.000000,0.000000,0.000000}%
\pgfsetstrokecolor{currentstroke}%
\pgfsetdash{}{0pt}%
\pgfsys@defobject{currentmarker}{\pgfqpoint{0.000000in}{0.000000in}}{\pgfqpoint{0.020833in}{0.000000in}}{%
\pgfpathmoveto{\pgfqpoint{0.000000in}{0.000000in}}%
\pgfpathlineto{\pgfqpoint{0.020833in}{0.000000in}}%
\pgfusepath{stroke,fill}%
}%
\begin{pgfscope}%
\pgfsys@transformshift{0.650000in}{1.618301in}%
\pgfsys@useobject{currentmarker}{}%
\end{pgfscope}%
\end{pgfscope}%
\begin{pgfscope}%
\pgfsetbuttcap%
\pgfsetroundjoin%
\definecolor{currentfill}{rgb}{0.000000,0.000000,0.000000}%
\pgfsetfillcolor{currentfill}%
\pgfsetlinewidth{0.401500pt}%
\definecolor{currentstroke}{rgb}{0.000000,0.000000,0.000000}%
\pgfsetstrokecolor{currentstroke}%
\pgfsetdash{}{0pt}%
\pgfsys@defobject{currentmarker}{\pgfqpoint{-0.020833in}{0.000000in}}{\pgfqpoint{-0.000000in}{0.000000in}}{%
\pgfpathmoveto{\pgfqpoint{-0.000000in}{0.000000in}}%
\pgfpathlineto{\pgfqpoint{-0.020833in}{0.000000in}}%
\pgfusepath{stroke,fill}%
}%
\begin{pgfscope}%
\pgfsys@transformshift{3.038110in}{1.618301in}%
\pgfsys@useobject{currentmarker}{}%
\end{pgfscope}%
\end{pgfscope}%
\begin{pgfscope}%
\pgfsetbuttcap%
\pgfsetroundjoin%
\definecolor{currentfill}{rgb}{0.000000,0.000000,0.000000}%
\pgfsetfillcolor{currentfill}%
\pgfsetlinewidth{0.401500pt}%
\definecolor{currentstroke}{rgb}{0.000000,0.000000,0.000000}%
\pgfsetstrokecolor{currentstroke}%
\pgfsetdash{}{0pt}%
\pgfsys@defobject{currentmarker}{\pgfqpoint{0.000000in}{0.000000in}}{\pgfqpoint{0.020833in}{0.000000in}}{%
\pgfpathmoveto{\pgfqpoint{0.000000in}{0.000000in}}%
\pgfpathlineto{\pgfqpoint{0.020833in}{0.000000in}}%
\pgfusepath{stroke,fill}%
}%
\begin{pgfscope}%
\pgfsys@transformshift{0.650000in}{1.692925in}%
\pgfsys@useobject{currentmarker}{}%
\end{pgfscope}%
\end{pgfscope}%
\begin{pgfscope}%
\pgfsetbuttcap%
\pgfsetroundjoin%
\definecolor{currentfill}{rgb}{0.000000,0.000000,0.000000}%
\pgfsetfillcolor{currentfill}%
\pgfsetlinewidth{0.401500pt}%
\definecolor{currentstroke}{rgb}{0.000000,0.000000,0.000000}%
\pgfsetstrokecolor{currentstroke}%
\pgfsetdash{}{0pt}%
\pgfsys@defobject{currentmarker}{\pgfqpoint{-0.020833in}{0.000000in}}{\pgfqpoint{-0.000000in}{0.000000in}}{%
\pgfpathmoveto{\pgfqpoint{-0.000000in}{0.000000in}}%
\pgfpathlineto{\pgfqpoint{-0.020833in}{0.000000in}}%
\pgfusepath{stroke,fill}%
}%
\begin{pgfscope}%
\pgfsys@transformshift{3.038110in}{1.692925in}%
\pgfsys@useobject{currentmarker}{}%
\end{pgfscope}%
\end{pgfscope}%
\begin{pgfscope}%
\pgfsetbuttcap%
\pgfsetroundjoin%
\definecolor{currentfill}{rgb}{0.000000,0.000000,0.000000}%
\pgfsetfillcolor{currentfill}%
\pgfsetlinewidth{0.401500pt}%
\definecolor{currentstroke}{rgb}{0.000000,0.000000,0.000000}%
\pgfsetstrokecolor{currentstroke}%
\pgfsetdash{}{0pt}%
\pgfsys@defobject{currentmarker}{\pgfqpoint{0.000000in}{0.000000in}}{\pgfqpoint{0.020833in}{0.000000in}}{%
\pgfpathmoveto{\pgfqpoint{0.000000in}{0.000000in}}%
\pgfpathlineto{\pgfqpoint{0.020833in}{0.000000in}}%
\pgfusepath{stroke,fill}%
}%
\begin{pgfscope}%
\pgfsys@transformshift{0.650000in}{1.842172in}%
\pgfsys@useobject{currentmarker}{}%
\end{pgfscope}%
\end{pgfscope}%
\begin{pgfscope}%
\pgfsetbuttcap%
\pgfsetroundjoin%
\definecolor{currentfill}{rgb}{0.000000,0.000000,0.000000}%
\pgfsetfillcolor{currentfill}%
\pgfsetlinewidth{0.401500pt}%
\definecolor{currentstroke}{rgb}{0.000000,0.000000,0.000000}%
\pgfsetstrokecolor{currentstroke}%
\pgfsetdash{}{0pt}%
\pgfsys@defobject{currentmarker}{\pgfqpoint{-0.020833in}{0.000000in}}{\pgfqpoint{-0.000000in}{0.000000in}}{%
\pgfpathmoveto{\pgfqpoint{-0.000000in}{0.000000in}}%
\pgfpathlineto{\pgfqpoint{-0.020833in}{0.000000in}}%
\pgfusepath{stroke,fill}%
}%
\begin{pgfscope}%
\pgfsys@transformshift{3.038110in}{1.842172in}%
\pgfsys@useobject{currentmarker}{}%
\end{pgfscope}%
\end{pgfscope}%
\begin{pgfscope}%
\pgfsetbuttcap%
\pgfsetroundjoin%
\definecolor{currentfill}{rgb}{0.000000,0.000000,0.000000}%
\pgfsetfillcolor{currentfill}%
\pgfsetlinewidth{0.401500pt}%
\definecolor{currentstroke}{rgb}{0.000000,0.000000,0.000000}%
\pgfsetstrokecolor{currentstroke}%
\pgfsetdash{}{0pt}%
\pgfsys@defobject{currentmarker}{\pgfqpoint{0.000000in}{0.000000in}}{\pgfqpoint{0.020833in}{0.000000in}}{%
\pgfpathmoveto{\pgfqpoint{0.000000in}{0.000000in}}%
\pgfpathlineto{\pgfqpoint{0.020833in}{0.000000in}}%
\pgfusepath{stroke,fill}%
}%
\begin{pgfscope}%
\pgfsys@transformshift{0.650000in}{1.916796in}%
\pgfsys@useobject{currentmarker}{}%
\end{pgfscope}%
\end{pgfscope}%
\begin{pgfscope}%
\pgfsetbuttcap%
\pgfsetroundjoin%
\definecolor{currentfill}{rgb}{0.000000,0.000000,0.000000}%
\pgfsetfillcolor{currentfill}%
\pgfsetlinewidth{0.401500pt}%
\definecolor{currentstroke}{rgb}{0.000000,0.000000,0.000000}%
\pgfsetstrokecolor{currentstroke}%
\pgfsetdash{}{0pt}%
\pgfsys@defobject{currentmarker}{\pgfqpoint{-0.020833in}{0.000000in}}{\pgfqpoint{-0.000000in}{0.000000in}}{%
\pgfpathmoveto{\pgfqpoint{-0.000000in}{0.000000in}}%
\pgfpathlineto{\pgfqpoint{-0.020833in}{0.000000in}}%
\pgfusepath{stroke,fill}%
}%
\begin{pgfscope}%
\pgfsys@transformshift{3.038110in}{1.916796in}%
\pgfsys@useobject{currentmarker}{}%
\end{pgfscope}%
\end{pgfscope}%
\begin{pgfscope}%
\pgfsetbuttcap%
\pgfsetroundjoin%
\definecolor{currentfill}{rgb}{0.000000,0.000000,0.000000}%
\pgfsetfillcolor{currentfill}%
\pgfsetlinewidth{0.401500pt}%
\definecolor{currentstroke}{rgb}{0.000000,0.000000,0.000000}%
\pgfsetstrokecolor{currentstroke}%
\pgfsetdash{}{0pt}%
\pgfsys@defobject{currentmarker}{\pgfqpoint{0.000000in}{0.000000in}}{\pgfqpoint{0.020833in}{0.000000in}}{%
\pgfpathmoveto{\pgfqpoint{0.000000in}{0.000000in}}%
\pgfpathlineto{\pgfqpoint{0.020833in}{0.000000in}}%
\pgfusepath{stroke,fill}%
}%
\begin{pgfscope}%
\pgfsys@transformshift{0.650000in}{1.991420in}%
\pgfsys@useobject{currentmarker}{}%
\end{pgfscope}%
\end{pgfscope}%
\begin{pgfscope}%
\pgfsetbuttcap%
\pgfsetroundjoin%
\definecolor{currentfill}{rgb}{0.000000,0.000000,0.000000}%
\pgfsetfillcolor{currentfill}%
\pgfsetlinewidth{0.401500pt}%
\definecolor{currentstroke}{rgb}{0.000000,0.000000,0.000000}%
\pgfsetstrokecolor{currentstroke}%
\pgfsetdash{}{0pt}%
\pgfsys@defobject{currentmarker}{\pgfqpoint{-0.020833in}{0.000000in}}{\pgfqpoint{-0.000000in}{0.000000in}}{%
\pgfpathmoveto{\pgfqpoint{-0.000000in}{0.000000in}}%
\pgfpathlineto{\pgfqpoint{-0.020833in}{0.000000in}}%
\pgfusepath{stroke,fill}%
}%
\begin{pgfscope}%
\pgfsys@transformshift{3.038110in}{1.991420in}%
\pgfsys@useobject{currentmarker}{}%
\end{pgfscope}%
\end{pgfscope}%
\begin{pgfscope}%
\pgfsetbuttcap%
\pgfsetroundjoin%
\definecolor{currentfill}{rgb}{0.000000,0.000000,0.000000}%
\pgfsetfillcolor{currentfill}%
\pgfsetlinewidth{0.401500pt}%
\definecolor{currentstroke}{rgb}{0.000000,0.000000,0.000000}%
\pgfsetstrokecolor{currentstroke}%
\pgfsetdash{}{0pt}%
\pgfsys@defobject{currentmarker}{\pgfqpoint{0.000000in}{0.000000in}}{\pgfqpoint{0.020833in}{0.000000in}}{%
\pgfpathmoveto{\pgfqpoint{0.000000in}{0.000000in}}%
\pgfpathlineto{\pgfqpoint{0.020833in}{0.000000in}}%
\pgfusepath{stroke,fill}%
}%
\begin{pgfscope}%
\pgfsys@transformshift{0.650000in}{2.066044in}%
\pgfsys@useobject{currentmarker}{}%
\end{pgfscope}%
\end{pgfscope}%
\begin{pgfscope}%
\pgfsetbuttcap%
\pgfsetroundjoin%
\definecolor{currentfill}{rgb}{0.000000,0.000000,0.000000}%
\pgfsetfillcolor{currentfill}%
\pgfsetlinewidth{0.401500pt}%
\definecolor{currentstroke}{rgb}{0.000000,0.000000,0.000000}%
\pgfsetstrokecolor{currentstroke}%
\pgfsetdash{}{0pt}%
\pgfsys@defobject{currentmarker}{\pgfqpoint{-0.020833in}{0.000000in}}{\pgfqpoint{-0.000000in}{0.000000in}}{%
\pgfpathmoveto{\pgfqpoint{-0.000000in}{0.000000in}}%
\pgfpathlineto{\pgfqpoint{-0.020833in}{0.000000in}}%
\pgfusepath{stroke,fill}%
}%
\begin{pgfscope}%
\pgfsys@transformshift{3.038110in}{2.066044in}%
\pgfsys@useobject{currentmarker}{}%
\end{pgfscope}%
\end{pgfscope}%
\begin{pgfscope}%
\pgfsetbuttcap%
\pgfsetroundjoin%
\definecolor{currentfill}{rgb}{0.000000,0.000000,0.000000}%
\pgfsetfillcolor{currentfill}%
\pgfsetlinewidth{0.401500pt}%
\definecolor{currentstroke}{rgb}{0.000000,0.000000,0.000000}%
\pgfsetstrokecolor{currentstroke}%
\pgfsetdash{}{0pt}%
\pgfsys@defobject{currentmarker}{\pgfqpoint{0.000000in}{0.000000in}}{\pgfqpoint{0.020833in}{0.000000in}}{%
\pgfpathmoveto{\pgfqpoint{0.000000in}{0.000000in}}%
\pgfpathlineto{\pgfqpoint{0.020833in}{0.000000in}}%
\pgfusepath{stroke,fill}%
}%
\begin{pgfscope}%
\pgfsys@transformshift{0.650000in}{2.215292in}%
\pgfsys@useobject{currentmarker}{}%
\end{pgfscope}%
\end{pgfscope}%
\begin{pgfscope}%
\pgfsetbuttcap%
\pgfsetroundjoin%
\definecolor{currentfill}{rgb}{0.000000,0.000000,0.000000}%
\pgfsetfillcolor{currentfill}%
\pgfsetlinewidth{0.401500pt}%
\definecolor{currentstroke}{rgb}{0.000000,0.000000,0.000000}%
\pgfsetstrokecolor{currentstroke}%
\pgfsetdash{}{0pt}%
\pgfsys@defobject{currentmarker}{\pgfqpoint{-0.020833in}{0.000000in}}{\pgfqpoint{-0.000000in}{0.000000in}}{%
\pgfpathmoveto{\pgfqpoint{-0.000000in}{0.000000in}}%
\pgfpathlineto{\pgfqpoint{-0.020833in}{0.000000in}}%
\pgfusepath{stroke,fill}%
}%
\begin{pgfscope}%
\pgfsys@transformshift{3.038110in}{2.215292in}%
\pgfsys@useobject{currentmarker}{}%
\end{pgfscope}%
\end{pgfscope}%
\begin{pgfscope}%
\pgfsetbuttcap%
\pgfsetroundjoin%
\definecolor{currentfill}{rgb}{0.000000,0.000000,0.000000}%
\pgfsetfillcolor{currentfill}%
\pgfsetlinewidth{0.401500pt}%
\definecolor{currentstroke}{rgb}{0.000000,0.000000,0.000000}%
\pgfsetstrokecolor{currentstroke}%
\pgfsetdash{}{0pt}%
\pgfsys@defobject{currentmarker}{\pgfqpoint{0.000000in}{0.000000in}}{\pgfqpoint{0.020833in}{0.000000in}}{%
\pgfpathmoveto{\pgfqpoint{0.000000in}{0.000000in}}%
\pgfpathlineto{\pgfqpoint{0.020833in}{0.000000in}}%
\pgfusepath{stroke,fill}%
}%
\begin{pgfscope}%
\pgfsys@transformshift{0.650000in}{2.289916in}%
\pgfsys@useobject{currentmarker}{}%
\end{pgfscope}%
\end{pgfscope}%
\begin{pgfscope}%
\pgfsetbuttcap%
\pgfsetroundjoin%
\definecolor{currentfill}{rgb}{0.000000,0.000000,0.000000}%
\pgfsetfillcolor{currentfill}%
\pgfsetlinewidth{0.401500pt}%
\definecolor{currentstroke}{rgb}{0.000000,0.000000,0.000000}%
\pgfsetstrokecolor{currentstroke}%
\pgfsetdash{}{0pt}%
\pgfsys@defobject{currentmarker}{\pgfqpoint{-0.020833in}{0.000000in}}{\pgfqpoint{-0.000000in}{0.000000in}}{%
\pgfpathmoveto{\pgfqpoint{-0.000000in}{0.000000in}}%
\pgfpathlineto{\pgfqpoint{-0.020833in}{0.000000in}}%
\pgfusepath{stroke,fill}%
}%
\begin{pgfscope}%
\pgfsys@transformshift{3.038110in}{2.289916in}%
\pgfsys@useobject{currentmarker}{}%
\end{pgfscope}%
\end{pgfscope}%
\begin{pgfscope}%
\pgfsetbuttcap%
\pgfsetroundjoin%
\definecolor{currentfill}{rgb}{0.000000,0.000000,0.000000}%
\pgfsetfillcolor{currentfill}%
\pgfsetlinewidth{0.401500pt}%
\definecolor{currentstroke}{rgb}{0.000000,0.000000,0.000000}%
\pgfsetstrokecolor{currentstroke}%
\pgfsetdash{}{0pt}%
\pgfsys@defobject{currentmarker}{\pgfqpoint{0.000000in}{0.000000in}}{\pgfqpoint{0.020833in}{0.000000in}}{%
\pgfpathmoveto{\pgfqpoint{0.000000in}{0.000000in}}%
\pgfpathlineto{\pgfqpoint{0.020833in}{0.000000in}}%
\pgfusepath{stroke,fill}%
}%
\begin{pgfscope}%
\pgfsys@transformshift{0.650000in}{2.364540in}%
\pgfsys@useobject{currentmarker}{}%
\end{pgfscope}%
\end{pgfscope}%
\begin{pgfscope}%
\pgfsetbuttcap%
\pgfsetroundjoin%
\definecolor{currentfill}{rgb}{0.000000,0.000000,0.000000}%
\pgfsetfillcolor{currentfill}%
\pgfsetlinewidth{0.401500pt}%
\definecolor{currentstroke}{rgb}{0.000000,0.000000,0.000000}%
\pgfsetstrokecolor{currentstroke}%
\pgfsetdash{}{0pt}%
\pgfsys@defobject{currentmarker}{\pgfqpoint{-0.020833in}{0.000000in}}{\pgfqpoint{-0.000000in}{0.000000in}}{%
\pgfpathmoveto{\pgfqpoint{-0.000000in}{0.000000in}}%
\pgfpathlineto{\pgfqpoint{-0.020833in}{0.000000in}}%
\pgfusepath{stroke,fill}%
}%
\begin{pgfscope}%
\pgfsys@transformshift{3.038110in}{2.364540in}%
\pgfsys@useobject{currentmarker}{}%
\end{pgfscope}%
\end{pgfscope}%
\begin{pgfscope}%
\definecolor{textcolor}{rgb}{0.000000,0.000000,0.000000}%
\pgfsetstrokecolor{textcolor}%
\pgfsetfillcolor{textcolor}%
\pgftext[x=0.260995in,y=1.417244in,,bottom,rotate=90.000000]{\color{textcolor}{\rmfamily\fontsize{12.000000}{14.400000}\selectfont\catcode`\^=\active\def^{\ifmmode\sp\else\^{}\fi}\catcode`\%=\active\def%{\%}TKE budget}}%
\end{pgfscope}%
\begin{pgfscope}%
\pgfpathrectangle{\pgfqpoint{0.650000in}{0.420000in}}{\pgfqpoint{2.388110in}{1.994488in}}%
\pgfusepath{clip}%
\pgfsetrectcap%
\pgfsetroundjoin%
\pgfsetlinewidth{0.602250pt}%
\definecolor{currentstroke}{rgb}{0.000000,0.000000,0.000000}%
\pgfsetstrokecolor{currentstroke}%
\pgfsetdash{}{0pt}%
\pgfpathmoveto{\pgfqpoint{0.645000in}{1.395098in}}%
\pgfpathlineto{\pgfqpoint{0.718874in}{1.395451in}}%
\pgfpathlineto{\pgfqpoint{0.855798in}{1.398407in}}%
\pgfpathlineto{\pgfqpoint{0.967671in}{1.406487in}}%
\pgfpathlineto{\pgfqpoint{1.062253in}{1.424921in}}%
\pgfpathlineto{\pgfqpoint{1.144181in}{1.461176in}}%
\pgfpathlineto{\pgfqpoint{1.216443in}{1.523531in}}%
\pgfpathlineto{\pgfqpoint{1.281080in}{1.617774in}}%
\pgfpathlineto{\pgfqpoint{1.339547in}{1.742987in}}%
\pgfpathlineto{\pgfqpoint{1.392920in}{1.888909in}}%
\pgfpathlineto{\pgfqpoint{1.442015in}{2.037310in}}%
\pgfpathlineto{\pgfqpoint{1.487466in}{2.167386in}}%
\pgfpathlineto{\pgfqpoint{1.529776in}{2.262435in}}%
\pgfpathlineto{\pgfqpoint{1.569350in}{2.314284in}}%
\pgfpathlineto{\pgfqpoint{1.606521in}{2.323830in}}%
\pgfpathlineto{\pgfqpoint{1.641563in}{2.298451in}}%
\pgfpathlineto{\pgfqpoint{1.674706in}{2.248393in}}%
\pgfpathlineto{\pgfqpoint{1.706145in}{2.183740in}}%
\pgfpathlineto{\pgfqpoint{1.736046in}{2.112739in}}%
\pgfpathlineto{\pgfqpoint{1.791786in}{1.973026in}}%
\pgfpathlineto{\pgfqpoint{1.817858in}{1.910132in}}%
\pgfpathlineto{\pgfqpoint{1.842862in}{1.853504in}}%
\pgfpathlineto{\pgfqpoint{1.866881in}{1.803290in}}%
\pgfpathlineto{\pgfqpoint{1.889989in}{1.759255in}}%
\pgfpathlineto{\pgfqpoint{1.912253in}{1.721033in}}%
\pgfpathlineto{\pgfqpoint{1.933732in}{1.688214in}}%
\pgfpathlineto{\pgfqpoint{1.954479in}{1.660208in}}%
\pgfpathlineto{\pgfqpoint{1.974542in}{1.636356in}}%
\pgfpathlineto{\pgfqpoint{1.993964in}{1.616002in}}%
\pgfpathlineto{\pgfqpoint{2.012785in}{1.598576in}}%
\pgfpathlineto{\pgfqpoint{2.031040in}{1.583580in}}%
\pgfpathlineto{\pgfqpoint{2.048762in}{1.570631in}}%
\pgfpathlineto{\pgfqpoint{2.065981in}{1.559525in}}%
\pgfpathlineto{\pgfqpoint{2.082724in}{1.550096in}}%
\pgfpathlineto{\pgfqpoint{2.099017in}{1.541999in}}%
\pgfpathlineto{\pgfqpoint{2.114883in}{1.534887in}}%
\pgfpathlineto{\pgfqpoint{2.130343in}{1.528656in}}%
\pgfpathlineto{\pgfqpoint{2.160125in}{1.518286in}}%
\pgfpathlineto{\pgfqpoint{2.188507in}{1.509734in}}%
\pgfpathlineto{\pgfqpoint{2.228719in}{1.499162in}}%
\pgfpathlineto{\pgfqpoint{2.266409in}{1.490616in}}%
\pgfpathlineto{\pgfqpoint{2.434988in}{1.454728in}}%
\pgfpathlineto{\pgfqpoint{2.470571in}{1.445626in}}%
\pgfpathlineto{\pgfqpoint{2.520191in}{1.432174in}}%
\pgfpathlineto{\pgfqpoint{2.565904in}{1.418220in}}%
\pgfpathlineto{\pgfqpoint{2.601408in}{1.407402in}}%
\pgfpathlineto{\pgfqpoint{2.621698in}{1.402172in}}%
\pgfpathlineto{\pgfqpoint{2.641289in}{1.398339in}}%
\pgfpathlineto{\pgfqpoint{2.660223in}{1.396027in}}%
\pgfpathlineto{\pgfqpoint{2.678539in}{1.394942in}}%
\pgfpathlineto{\pgfqpoint{2.707791in}{1.394471in}}%
\pgfpathlineto{\pgfqpoint{2.885757in}{1.394429in}}%
\pgfpathlineto{\pgfqpoint{3.043110in}{1.394429in}}%
\pgfpathlineto{\pgfqpoint{3.043110in}{1.394429in}}%
\pgfusepath{stroke}%
\end{pgfscope}%
\begin{pgfscope}%
\pgfpathrectangle{\pgfqpoint{0.650000in}{0.420000in}}{\pgfqpoint{2.388110in}{1.994488in}}%
\pgfusepath{clip}%
\pgfsetbuttcap%
\pgfsetroundjoin%
\pgfsetlinewidth{0.602250pt}%
\definecolor{currentstroke}{rgb}{0.000000,0.000000,0.000000}%
\pgfsetstrokecolor{currentstroke}%
\pgfsetdash{{3.000000pt}{0.600000pt}}{0.000000pt}%
\pgfpathmoveto{\pgfqpoint{0.645000in}{1.394429in}}%
\pgfpathlineto{\pgfqpoint{1.569350in}{1.394940in}}%
\pgfpathlineto{\pgfqpoint{1.674706in}{1.395208in}}%
\pgfpathlineto{\pgfqpoint{1.889989in}{1.394214in}}%
\pgfpathlineto{\pgfqpoint{2.031040in}{1.395191in}}%
\pgfpathlineto{\pgfqpoint{2.065981in}{1.396303in}}%
\pgfpathlineto{\pgfqpoint{2.160125in}{1.395280in}}%
\pgfpathlineto{\pgfqpoint{2.215611in}{1.393736in}}%
\pgfpathlineto{\pgfqpoint{2.301868in}{1.393177in}}%
\pgfpathlineto{\pgfqpoint{2.406815in}{1.391195in}}%
\pgfpathlineto{\pgfqpoint{2.444088in}{1.389253in}}%
\pgfpathlineto{\pgfqpoint{2.528065in}{1.386913in}}%
\pgfpathlineto{\pgfqpoint{2.573183in}{1.385219in}}%
\pgfpathlineto{\pgfqpoint{2.615015in}{1.384830in}}%
\pgfpathlineto{\pgfqpoint{2.634834in}{1.385654in}}%
\pgfpathlineto{\pgfqpoint{2.653982in}{1.387294in}}%
\pgfpathlineto{\pgfqpoint{2.707791in}{1.392488in}}%
\pgfpathlineto{\pgfqpoint{2.740966in}{1.393788in}}%
\pgfpathlineto{\pgfqpoint{2.792128in}{1.394301in}}%
\pgfpathlineto{\pgfqpoint{2.994981in}{1.394427in}}%
\pgfpathlineto{\pgfqpoint{3.043110in}{1.394428in}}%
\pgfpathlineto{\pgfqpoint{3.043110in}{1.394428in}}%
\pgfusepath{stroke}%
\end{pgfscope}%
\begin{pgfscope}%
\pgfpathrectangle{\pgfqpoint{0.650000in}{0.420000in}}{\pgfqpoint{2.388110in}{1.994488in}}%
\pgfusepath{clip}%
\pgfsetbuttcap%
\pgfsetroundjoin%
\pgfsetlinewidth{0.602250pt}%
\definecolor{currentstroke}{rgb}{0.000000,0.000000,0.000000}%
\pgfsetstrokecolor{currentstroke}%
\pgfsetdash{{1.800000pt}{0.600000pt}{0.600000pt}{0.600000pt}}{0.000000pt}%
\pgfpathmoveto{\pgfqpoint{0.645000in}{0.576535in}}%
\pgfpathlineto{\pgfqpoint{0.718874in}{0.588516in}}%
\pgfpathlineto{\pgfqpoint{0.855798in}{0.617079in}}%
\pgfpathlineto{\pgfqpoint{0.967671in}{0.644279in}}%
\pgfpathlineto{\pgfqpoint{1.062253in}{0.671606in}}%
\pgfpathlineto{\pgfqpoint{1.144181in}{0.702023in}}%
\pgfpathlineto{\pgfqpoint{1.216443in}{0.737456in}}%
\pgfpathlineto{\pgfqpoint{1.281080in}{0.776093in}}%
\pgfpathlineto{\pgfqpoint{1.339547in}{0.812339in}}%
\pgfpathlineto{\pgfqpoint{1.392920in}{0.839956in}}%
\pgfpathlineto{\pgfqpoint{1.442015in}{0.855919in}}%
\pgfpathlineto{\pgfqpoint{1.487466in}{0.861901in}}%
\pgfpathlineto{\pgfqpoint{1.569350in}{0.863543in}}%
\pgfpathlineto{\pgfqpoint{1.606521in}{0.867978in}}%
\pgfpathlineto{\pgfqpoint{1.641563in}{0.877403in}}%
\pgfpathlineto{\pgfqpoint{1.674706in}{0.891553in}}%
\pgfpathlineto{\pgfqpoint{1.706145in}{0.909397in}}%
\pgfpathlineto{\pgfqpoint{1.736046in}{0.929781in}}%
\pgfpathlineto{\pgfqpoint{1.764552in}{0.951767in}}%
\pgfpathlineto{\pgfqpoint{1.791786in}{0.974646in}}%
\pgfpathlineto{\pgfqpoint{1.842862in}{1.020650in}}%
\pgfpathlineto{\pgfqpoint{1.933732in}{1.104384in}}%
\pgfpathlineto{\pgfqpoint{1.974542in}{1.139384in}}%
\pgfpathlineto{\pgfqpoint{1.993964in}{1.155031in}}%
\pgfpathlineto{\pgfqpoint{2.012785in}{1.169465in}}%
\pgfpathlineto{\pgfqpoint{2.048762in}{1.195040in}}%
\pgfpathlineto{\pgfqpoint{2.065981in}{1.206313in}}%
\pgfpathlineto{\pgfqpoint{2.099017in}{1.225920in}}%
\pgfpathlineto{\pgfqpoint{2.130343in}{1.242486in}}%
\pgfpathlineto{\pgfqpoint{2.160125in}{1.256666in}}%
\pgfpathlineto{\pgfqpoint{2.188507in}{1.268859in}}%
\pgfpathlineto{\pgfqpoint{2.215611in}{1.279457in}}%
\pgfpathlineto{\pgfqpoint{2.241547in}{1.288738in}}%
\pgfpathlineto{\pgfqpoint{2.278463in}{1.300701in}}%
\pgfpathlineto{\pgfqpoint{2.324390in}{1.314197in}}%
\pgfpathlineto{\pgfqpoint{2.377230in}{1.328186in}}%
\pgfpathlineto{\pgfqpoint{2.453049in}{1.346764in}}%
\pgfpathlineto{\pgfqpoint{2.565904in}{1.373628in}}%
\pgfpathlineto{\pgfqpoint{2.601408in}{1.381232in}}%
\pgfpathlineto{\pgfqpoint{2.628304in}{1.386096in}}%
\pgfpathlineto{\pgfqpoint{2.653982in}{1.389712in}}%
\pgfpathlineto{\pgfqpoint{2.678539in}{1.392116in}}%
\pgfpathlineto{\pgfqpoint{2.707791in}{1.393658in}}%
\pgfpathlineto{\pgfqpoint{2.746305in}{1.394305in}}%
\pgfpathlineto{\pgfqpoint{2.864951in}{1.394427in}}%
\pgfpathlineto{\pgfqpoint{3.043110in}{1.394429in}}%
\pgfpathlineto{\pgfqpoint{3.043110in}{1.394429in}}%
\pgfusepath{stroke}%
\end{pgfscope}%
\begin{pgfscope}%
\pgfpathrectangle{\pgfqpoint{0.650000in}{0.420000in}}{\pgfqpoint{2.388110in}{1.994488in}}%
\pgfusepath{clip}%
\pgfsetbuttcap%
\pgfsetroundjoin%
\pgfsetlinewidth{0.602250pt}%
\definecolor{currentstroke}{rgb}{0.000000,0.000000,0.000000}%
\pgfsetstrokecolor{currentstroke}%
\pgfsetdash{{0.600000pt}{0.600000pt}}{0.000000pt}%
\pgfpathmoveto{\pgfqpoint{0.645000in}{1.428556in}}%
\pgfpathlineto{\pgfqpoint{0.718874in}{1.434834in}}%
\pgfpathlineto{\pgfqpoint{0.967671in}{1.460793in}}%
\pgfpathlineto{\pgfqpoint{1.062253in}{1.468437in}}%
\pgfpathlineto{\pgfqpoint{1.144181in}{1.471802in}}%
\pgfpathlineto{\pgfqpoint{1.216443in}{1.471281in}}%
\pgfpathlineto{\pgfqpoint{1.281080in}{1.467748in}}%
\pgfpathlineto{\pgfqpoint{1.339547in}{1.462037in}}%
\pgfpathlineto{\pgfqpoint{1.392920in}{1.454837in}}%
\pgfpathlineto{\pgfqpoint{1.442015in}{1.446496in}}%
\pgfpathlineto{\pgfqpoint{1.487466in}{1.437302in}}%
\pgfpathlineto{\pgfqpoint{1.529776in}{1.427510in}}%
\pgfpathlineto{\pgfqpoint{1.641563in}{1.399745in}}%
\pgfpathlineto{\pgfqpoint{1.674706in}{1.392867in}}%
\pgfpathlineto{\pgfqpoint{1.706145in}{1.387826in}}%
\pgfpathlineto{\pgfqpoint{1.736046in}{1.384611in}}%
\pgfpathlineto{\pgfqpoint{1.764552in}{1.383087in}}%
\pgfpathlineto{\pgfqpoint{1.791786in}{1.382925in}}%
\pgfpathlineto{\pgfqpoint{1.817858in}{1.383859in}}%
\pgfpathlineto{\pgfqpoint{1.866881in}{1.387696in}}%
\pgfpathlineto{\pgfqpoint{1.933732in}{1.394267in}}%
\pgfpathlineto{\pgfqpoint{1.974542in}{1.397057in}}%
\pgfpathlineto{\pgfqpoint{2.012785in}{1.398006in}}%
\pgfpathlineto{\pgfqpoint{2.048762in}{1.398000in}}%
\pgfpathlineto{\pgfqpoint{2.202211in}{1.394831in}}%
\pgfpathlineto{\pgfqpoint{2.254107in}{1.393708in}}%
\pgfpathlineto{\pgfqpoint{2.335340in}{1.393982in}}%
\pgfpathlineto{\pgfqpoint{2.406815in}{1.393640in}}%
\pgfpathlineto{\pgfqpoint{2.487584in}{1.392616in}}%
\pgfpathlineto{\pgfqpoint{2.594481in}{1.390783in}}%
\pgfpathlineto{\pgfqpoint{2.634834in}{1.391418in}}%
\pgfpathlineto{\pgfqpoint{2.746305in}{1.394617in}}%
\pgfpathlineto{\pgfqpoint{3.043110in}{1.394430in}}%
\pgfpathlineto{\pgfqpoint{3.043110in}{1.394430in}}%
\pgfusepath{stroke}%
\end{pgfscope}%
\begin{pgfscope}%
\pgfpathrectangle{\pgfqpoint{0.650000in}{0.420000in}}{\pgfqpoint{2.388110in}{1.994488in}}%
\pgfusepath{clip}%
\pgfsetbuttcap%
\pgfsetroundjoin%
\pgfsetlinewidth{0.602250pt}%
\definecolor{currentstroke}{rgb}{0.000000,0.000000,0.000000}%
\pgfsetstrokecolor{currentstroke}%
\pgfsetdash{{3.000000pt}{3.000000pt}}{0.000000pt}%
\pgfpathmoveto{\pgfqpoint{0.645000in}{1.395460in}}%
\pgfpathlineto{\pgfqpoint{0.718874in}{1.396000in}}%
\pgfpathlineto{\pgfqpoint{0.855798in}{1.400399in}}%
\pgfpathlineto{\pgfqpoint{0.967671in}{1.411774in}}%
\pgfpathlineto{\pgfqpoint{1.062253in}{1.435184in}}%
\pgfpathlineto{\pgfqpoint{1.144181in}{1.473513in}}%
\pgfpathlineto{\pgfqpoint{1.281080in}{1.565130in}}%
\pgfpathlineto{\pgfqpoint{1.339547in}{1.581622in}}%
\pgfpathlineto{\pgfqpoint{1.392920in}{1.555855in}}%
\pgfpathlineto{\pgfqpoint{1.442015in}{1.487769in}}%
\pgfpathlineto{\pgfqpoint{1.487466in}{1.393216in}}%
\pgfpathlineto{\pgfqpoint{1.529776in}{1.295574in}}%
\pgfpathlineto{\pgfqpoint{1.569350in}{1.215298in}}%
\pgfpathlineto{\pgfqpoint{1.606521in}{1.163738in}}%
\pgfpathlineto{\pgfqpoint{1.641563in}{1.142444in}}%
\pgfpathlineto{\pgfqpoint{1.674706in}{1.146301in}}%
\pgfpathlineto{\pgfqpoint{1.706145in}{1.167439in}}%
\pgfpathlineto{\pgfqpoint{1.736046in}{1.198224in}}%
\pgfpathlineto{\pgfqpoint{1.764552in}{1.232661in}}%
\pgfpathlineto{\pgfqpoint{1.791786in}{1.266723in}}%
\pgfpathlineto{\pgfqpoint{1.817858in}{1.297701in}}%
\pgfpathlineto{\pgfqpoint{1.842862in}{1.324273in}}%
\pgfpathlineto{\pgfqpoint{1.866881in}{1.346376in}}%
\pgfpathlineto{\pgfqpoint{1.889989in}{1.364341in}}%
\pgfpathlineto{\pgfqpoint{1.912253in}{1.378480in}}%
\pgfpathlineto{\pgfqpoint{1.933732in}{1.389227in}}%
\pgfpathlineto{\pgfqpoint{1.954479in}{1.397171in}}%
\pgfpathlineto{\pgfqpoint{1.974542in}{1.403074in}}%
\pgfpathlineto{\pgfqpoint{1.993964in}{1.407420in}}%
\pgfpathlineto{\pgfqpoint{2.012785in}{1.410203in}}%
\pgfpathlineto{\pgfqpoint{2.031040in}{1.411449in}}%
\pgfpathlineto{\pgfqpoint{2.065981in}{1.411697in}}%
\pgfpathlineto{\pgfqpoint{2.099017in}{1.410841in}}%
\pgfpathlineto{\pgfqpoint{2.188507in}{1.403752in}}%
\pgfpathlineto{\pgfqpoint{2.228719in}{1.401014in}}%
\pgfpathlineto{\pgfqpoint{2.266409in}{1.397462in}}%
\pgfpathlineto{\pgfqpoint{2.313235in}{1.393470in}}%
\pgfpathlineto{\pgfqpoint{2.377230in}{1.390941in}}%
\pgfpathlineto{\pgfqpoint{2.406815in}{1.390884in}}%
\pgfpathlineto{\pgfqpoint{2.479139in}{1.394219in}}%
\pgfpathlineto{\pgfqpoint{2.535833in}{1.399944in}}%
\pgfpathlineto{\pgfqpoint{2.587469in}{1.406233in}}%
\pgfpathlineto{\pgfqpoint{2.615015in}{1.408402in}}%
\pgfpathlineto{\pgfqpoint{2.628304in}{1.408379in}}%
\pgfpathlineto{\pgfqpoint{2.641289in}{1.407592in}}%
\pgfpathlineto{\pgfqpoint{2.660223in}{1.404816in}}%
\pgfpathlineto{\pgfqpoint{2.696274in}{1.398475in}}%
\pgfpathlineto{\pgfqpoint{2.713461in}{1.396507in}}%
\pgfpathlineto{\pgfqpoint{2.735574in}{1.395156in}}%
\pgfpathlineto{\pgfqpoint{2.767152in}{1.394547in}}%
\pgfpathlineto{\pgfqpoint{2.873376in}{1.394429in}}%
\pgfpathlineto{\pgfqpoint{3.043110in}{1.394429in}}%
\pgfpathlineto{\pgfqpoint{3.043110in}{1.394429in}}%
\pgfusepath{stroke}%
\end{pgfscope}%
\begin{pgfscope}%
\pgfpathrectangle{\pgfqpoint{0.650000in}{0.420000in}}{\pgfqpoint{2.388110in}{1.994488in}}%
\pgfusepath{clip}%
\pgfsetbuttcap%
\pgfsetroundjoin%
\pgfsetlinewidth{0.602250pt}%
\definecolor{currentstroke}{rgb}{0.000000,0.000000,0.000000}%
\pgfsetstrokecolor{currentstroke}%
\pgfsetdash{{1.800000pt}{3.000000pt}{0.600000pt}{3.000000pt}}{0.000000pt}%
\pgfpathmoveto{\pgfqpoint{0.645000in}{2.176506in}}%
\pgfpathlineto{\pgfqpoint{0.718874in}{2.157324in}}%
\pgfpathlineto{\pgfqpoint{0.855798in}{2.107079in}}%
\pgfpathlineto{\pgfqpoint{0.967671in}{2.048791in}}%
\pgfpathlineto{\pgfqpoint{1.062253in}{1.971981in}}%
\pgfpathlineto{\pgfqpoint{1.144181in}{1.863638in}}%
\pgfpathlineto{\pgfqpoint{1.216443in}{1.717929in}}%
\pgfpathlineto{\pgfqpoint{1.281080in}{1.545502in}}%
\pgfpathlineto{\pgfqpoint{1.339547in}{1.373304in}}%
\pgfpathlineto{\pgfqpoint{1.392920in}{1.232733in}}%
\pgfpathlineto{\pgfqpoint{1.442015in}{1.144710in}}%
\pgfpathlineto{\pgfqpoint{1.487466in}{1.112257in}}%
\pgfpathlineto{\pgfqpoint{1.529776in}{1.123624in}}%
\pgfpathlineto{\pgfqpoint{1.569350in}{1.160894in}}%
\pgfpathlineto{\pgfqpoint{1.606521in}{1.207742in}}%
\pgfpathlineto{\pgfqpoint{1.641563in}{1.253280in}}%
\pgfpathlineto{\pgfqpoint{1.674706in}{1.292237in}}%
\pgfpathlineto{\pgfqpoint{1.706145in}{1.323144in}}%
\pgfpathlineto{\pgfqpoint{1.736046in}{1.346509in}}%
\pgfpathlineto{\pgfqpoint{1.764552in}{1.363439in}}%
\pgfpathlineto{\pgfqpoint{1.791786in}{1.375164in}}%
\pgfpathlineto{\pgfqpoint{1.817858in}{1.383127in}}%
\pgfpathlineto{\pgfqpoint{1.842862in}{1.388542in}}%
\pgfpathlineto{\pgfqpoint{1.866881in}{1.392102in}}%
\pgfpathlineto{\pgfqpoint{1.889989in}{1.394259in}}%
\pgfpathlineto{\pgfqpoint{1.912253in}{1.395449in}}%
\pgfpathlineto{\pgfqpoint{1.954479in}{1.396083in}}%
\pgfpathlineto{\pgfqpoint{2.215611in}{1.394429in}}%
\pgfpathlineto{\pgfqpoint{2.767152in}{1.394456in}}%
\pgfpathlineto{\pgfqpoint{3.043110in}{1.394429in}}%
\pgfpathlineto{\pgfqpoint{3.043110in}{1.394429in}}%
\pgfusepath{stroke}%
\end{pgfscope}%
\begin{pgfscope}%
\pgfpathrectangle{\pgfqpoint{0.650000in}{0.420000in}}{\pgfqpoint{2.388110in}{1.994488in}}%
\pgfusepath{clip}%
\pgfsetbuttcap%
\pgfsetroundjoin%
\pgfsetlinewidth{0.602250pt}%
\definecolor{currentstroke}{rgb}{0.000000,0.000000,0.000000}%
\pgfsetstrokecolor{currentstroke}%
\pgfsetdash{{0.600000pt}{3.000000pt}}{0.000000pt}%
\pgfpathmoveto{\pgfqpoint{0.645000in}{1.394438in}}%
\pgfpathlineto{\pgfqpoint{1.933732in}{1.394486in}}%
\pgfpathlineto{\pgfqpoint{2.543499in}{1.394400in}}%
\pgfpathlineto{\pgfqpoint{2.811301in}{1.394426in}}%
\pgfpathlineto{\pgfqpoint{3.043110in}{1.394428in}}%
\pgfpathlineto{\pgfqpoint{3.043110in}{1.394428in}}%
\pgfusepath{stroke}%
\end{pgfscope}%
\begin{pgfscope}%
\pgfpathrectangle{\pgfqpoint{0.650000in}{0.420000in}}{\pgfqpoint{2.388110in}{1.994488in}}%
\pgfusepath{clip}%
\pgfsetrectcap%
\pgfsetroundjoin%
\pgfsetlinewidth{1.003750pt}%
\definecolor{currentstroke}{rgb}{0.870588,0.466667,0.682353}%
\pgfsetstrokecolor{currentstroke}%
\pgfsetdash{}{0pt}%
\pgfpathmoveto{\pgfqpoint{0.645000in}{1.394950in}}%
\pgfpathlineto{\pgfqpoint{0.741886in}{1.395742in}}%
\pgfpathlineto{\pgfqpoint{0.815854in}{1.397142in}}%
\pgfpathlineto{\pgfqpoint{0.882094in}{1.399656in}}%
\pgfpathlineto{\pgfqpoint{0.923033in}{1.402273in}}%
\pgfpathlineto{\pgfqpoint{0.961915in}{1.405950in}}%
\pgfpathlineto{\pgfqpoint{0.999071in}{1.411031in}}%
\pgfpathlineto{\pgfqpoint{1.017086in}{1.414231in}}%
\pgfpathlineto{\pgfqpoint{1.034762in}{1.417955in}}%
\pgfpathlineto{\pgfqpoint{1.052126in}{1.422275in}}%
\pgfpathlineto{\pgfqpoint{1.069199in}{1.427274in}}%
\pgfpathlineto{\pgfqpoint{1.086001in}{1.433043in}}%
\pgfpathlineto{\pgfqpoint{1.102550in}{1.439683in}}%
\pgfpathlineto{\pgfqpoint{1.118865in}{1.447298in}}%
\pgfpathlineto{\pgfqpoint{1.134959in}{1.455997in}}%
\pgfpathlineto{\pgfqpoint{1.150847in}{1.465880in}}%
\pgfpathlineto{\pgfqpoint{1.166541in}{1.477038in}}%
\pgfpathlineto{\pgfqpoint{1.182055in}{1.489554in}}%
\pgfpathlineto{\pgfqpoint{1.197397in}{1.503536in}}%
\pgfpathlineto{\pgfqpoint{1.212579in}{1.519147in}}%
\pgfpathlineto{\pgfqpoint{1.227610in}{1.536594in}}%
\pgfpathlineto{\pgfqpoint{1.242498in}{1.555993in}}%
\pgfpathlineto{\pgfqpoint{1.257252in}{1.577328in}}%
\pgfpathlineto{\pgfqpoint{1.271878in}{1.600618in}}%
\pgfpathlineto{\pgfqpoint{1.286384in}{1.625999in}}%
\pgfpathlineto{\pgfqpoint{1.300777in}{1.653616in}}%
\pgfpathlineto{\pgfqpoint{1.315062in}{1.683491in}}%
\pgfpathlineto{\pgfqpoint{1.329244in}{1.715498in}}%
\pgfpathlineto{\pgfqpoint{1.343330in}{1.749434in}}%
\pgfpathlineto{\pgfqpoint{1.357325in}{1.785089in}}%
\pgfpathlineto{\pgfqpoint{1.371232in}{1.822286in}}%
\pgfpathlineto{\pgfqpoint{1.398802in}{1.900716in}}%
\pgfpathlineto{\pgfqpoint{1.453072in}{2.061732in}}%
\pgfpathlineto{\pgfqpoint{1.479826in}{2.136419in}}%
\pgfpathlineto{\pgfqpoint{1.493118in}{2.170975in}}%
\pgfpathlineto{\pgfqpoint{1.506356in}{2.203267in}}%
\pgfpathlineto{\pgfqpoint{1.519544in}{2.232652in}}%
\pgfpathlineto{\pgfqpoint{1.532684in}{2.258374in}}%
\pgfpathlineto{\pgfqpoint{1.545777in}{2.279985in}}%
\pgfpathlineto{\pgfqpoint{1.558826in}{2.297301in}}%
\pgfpathlineto{\pgfqpoint{1.571832in}{2.310031in}}%
\pgfpathlineto{\pgfqpoint{1.584799in}{2.317795in}}%
\pgfpathlineto{\pgfqpoint{1.597727in}{2.320696in}}%
\pgfpathlineto{\pgfqpoint{1.610618in}{2.319289in}}%
\pgfpathlineto{\pgfqpoint{1.623474in}{2.314061in}}%
\pgfpathlineto{\pgfqpoint{1.636296in}{2.305156in}}%
\pgfpathlineto{\pgfqpoint{1.649086in}{2.292532in}}%
\pgfpathlineto{\pgfqpoint{1.661845in}{2.276241in}}%
\pgfpathlineto{\pgfqpoint{1.674574in}{2.256570in}}%
\pgfpathlineto{\pgfqpoint{1.687276in}{2.233903in}}%
\pgfpathlineto{\pgfqpoint{1.699950in}{2.208662in}}%
\pgfpathlineto{\pgfqpoint{1.712599in}{2.181457in}}%
\pgfpathlineto{\pgfqpoint{1.737822in}{2.123565in}}%
\pgfpathlineto{\pgfqpoint{1.762955in}{2.062794in}}%
\pgfpathlineto{\pgfqpoint{1.812976in}{1.940236in}}%
\pgfpathlineto{\pgfqpoint{1.837878in}{1.883203in}}%
\pgfpathlineto{\pgfqpoint{1.862714in}{1.829505in}}%
\pgfpathlineto{\pgfqpoint{1.875110in}{1.804370in}}%
\pgfpathlineto{\pgfqpoint{1.887491in}{1.780563in}}%
\pgfpathlineto{\pgfqpoint{1.912213in}{1.736638in}}%
\pgfpathlineto{\pgfqpoint{1.924555in}{1.716379in}}%
\pgfpathlineto{\pgfqpoint{1.936885in}{1.697282in}}%
\pgfpathlineto{\pgfqpoint{1.949203in}{1.679441in}}%
\pgfpathlineto{\pgfqpoint{1.961510in}{1.662904in}}%
\pgfpathlineto{\pgfqpoint{1.973806in}{1.647596in}}%
\pgfpathlineto{\pgfqpoint{1.986092in}{1.633428in}}%
\pgfpathlineto{\pgfqpoint{1.998368in}{1.620361in}}%
\pgfpathlineto{\pgfqpoint{2.010634in}{1.608394in}}%
\pgfpathlineto{\pgfqpoint{2.022892in}{1.597518in}}%
\pgfpathlineto{\pgfqpoint{2.035140in}{1.587656in}}%
\pgfpathlineto{\pgfqpoint{2.047381in}{1.578690in}}%
\pgfpathlineto{\pgfqpoint{2.059613in}{1.570509in}}%
\pgfpathlineto{\pgfqpoint{2.084054in}{1.556163in}}%
\pgfpathlineto{\pgfqpoint{2.108467in}{1.543884in}}%
\pgfpathlineto{\pgfqpoint{2.132853in}{1.533253in}}%
\pgfpathlineto{\pgfqpoint{2.157216in}{1.523713in}}%
\pgfpathlineto{\pgfqpoint{2.218027in}{1.501940in}}%
\pgfpathlineto{\pgfqpoint{2.242318in}{1.493678in}}%
\pgfpathlineto{\pgfqpoint{2.266591in}{1.486634in}}%
\pgfpathlineto{\pgfqpoint{2.290850in}{1.480808in}}%
\pgfpathlineto{\pgfqpoint{2.375643in}{1.462287in}}%
\pgfpathlineto{\pgfqpoint{2.445607in}{1.445939in}}%
\pgfpathlineto{\pgfqpoint{2.476366in}{1.440673in}}%
\pgfpathlineto{\pgfqpoint{2.521638in}{1.432967in}}%
\pgfpathlineto{\pgfqpoint{2.561082in}{1.425802in}}%
\pgfpathlineto{\pgfqpoint{2.608963in}{1.416349in}}%
\pgfpathlineto{\pgfqpoint{2.644786in}{1.410640in}}%
\pgfpathlineto{\pgfqpoint{2.701245in}{1.403466in}}%
\pgfpathlineto{\pgfqpoint{2.764561in}{1.396487in}}%
\pgfpathlineto{\pgfqpoint{2.793649in}{1.394897in}}%
\pgfpathlineto{\pgfqpoint{2.823382in}{1.394455in}}%
\pgfpathlineto{\pgfqpoint{2.996727in}{1.394428in}}%
\pgfpathlineto{\pgfqpoint{3.043110in}{1.394428in}}%
\pgfpathlineto{\pgfqpoint{3.043110in}{1.394428in}}%
\pgfusepath{stroke}%
\end{pgfscope}%
\begin{pgfscope}%
\pgfpathrectangle{\pgfqpoint{0.650000in}{0.420000in}}{\pgfqpoint{2.388110in}{1.994488in}}%
\pgfusepath{clip}%
\pgfsetbuttcap%
\pgfsetroundjoin%
\pgfsetlinewidth{1.003750pt}%
\definecolor{currentstroke}{rgb}{0.870588,0.466667,0.682353}%
\pgfsetstrokecolor{currentstroke}%
\pgfsetdash{{5.000000pt}{1.000000pt}}{0.000000pt}%
\pgfpathmoveto{\pgfqpoint{0.645000in}{1.394429in}}%
\pgfpathlineto{\pgfqpoint{1.315062in}{1.393979in}}%
\pgfpathlineto{\pgfqpoint{1.493118in}{1.391955in}}%
\pgfpathlineto{\pgfqpoint{1.584799in}{1.389793in}}%
\pgfpathlineto{\pgfqpoint{1.674574in}{1.386731in}}%
\pgfpathlineto{\pgfqpoint{1.737822in}{1.384626in}}%
\pgfpathlineto{\pgfqpoint{1.775489in}{1.384518in}}%
\pgfpathlineto{\pgfqpoint{1.800499in}{1.385265in}}%
\pgfpathlineto{\pgfqpoint{1.850304in}{1.388396in}}%
\pgfpathlineto{\pgfqpoint{1.949203in}{1.395607in}}%
\pgfpathlineto{\pgfqpoint{2.071837in}{1.402312in}}%
\pgfpathlineto{\pgfqpoint{2.132853in}{1.408027in}}%
\pgfpathlineto{\pgfqpoint{2.157216in}{1.408816in}}%
\pgfpathlineto{\pgfqpoint{2.205874in}{1.408387in}}%
\pgfpathlineto{\pgfqpoint{2.254457in}{1.407544in}}%
\pgfpathlineto{\pgfqpoint{2.278722in}{1.407018in}}%
\pgfpathlineto{\pgfqpoint{2.302973in}{1.405155in}}%
\pgfpathlineto{\pgfqpoint{2.327210in}{1.402861in}}%
\pgfpathlineto{\pgfqpoint{2.351435in}{1.402026in}}%
\pgfpathlineto{\pgfqpoint{2.411648in}{1.402334in}}%
\pgfpathlineto{\pgfqpoint{2.538029in}{1.398490in}}%
\pgfpathlineto{\pgfqpoint{2.568397in}{1.396771in}}%
\pgfpathlineto{\pgfqpoint{2.627397in}{1.395745in}}%
\pgfpathlineto{\pgfqpoint{2.671743in}{1.395859in}}%
\pgfpathlineto{\pgfqpoint{2.740778in}{1.392257in}}%
\pgfpathlineto{\pgfqpoint{2.756836in}{1.392411in}}%
\pgfpathlineto{\pgfqpoint{2.775794in}{1.392688in}}%
\pgfpathlineto{\pgfqpoint{2.823382in}{1.393608in}}%
\pgfpathlineto{\pgfqpoint{2.864570in}{1.394037in}}%
\pgfpathlineto{\pgfqpoint{2.903314in}{1.394400in}}%
\pgfpathlineto{\pgfqpoint{2.942018in}{1.394407in}}%
\pgfpathlineto{\pgfqpoint{3.014121in}{1.394454in}}%
\pgfpathlineto{\pgfqpoint{3.036926in}{1.394465in}}%
\pgfpathlineto{\pgfqpoint{3.043110in}{1.394492in}}%
\pgfpathlineto{\pgfqpoint{3.043110in}{1.394492in}}%
\pgfusepath{stroke}%
\end{pgfscope}%
\begin{pgfscope}%
\pgfpathrectangle{\pgfqpoint{0.650000in}{0.420000in}}{\pgfqpoint{2.388110in}{1.994488in}}%
\pgfusepath{clip}%
\pgfsetbuttcap%
\pgfsetroundjoin%
\pgfsetlinewidth{1.003750pt}%
\definecolor{currentstroke}{rgb}{0.870588,0.466667,0.682353}%
\pgfsetstrokecolor{currentstroke}%
\pgfsetdash{{3.000000pt}{1.000000pt}{1.000000pt}{1.000000pt}}{0.000000pt}%
\pgfpathmoveto{\pgfqpoint{0.645000in}{0.545945in}}%
\pgfpathlineto{\pgfqpoint{0.656764in}{0.547370in}}%
\pgfpathlineto{\pgfqpoint{0.686665in}{0.552096in}}%
\pgfpathlineto{\pgfqpoint{0.714975in}{0.557721in}}%
\pgfpathlineto{\pgfqpoint{0.741886in}{0.564212in}}%
\pgfpathlineto{\pgfqpoint{0.767576in}{0.571497in}}%
\pgfpathlineto{\pgfqpoint{0.792190in}{0.579467in}}%
\pgfpathlineto{\pgfqpoint{0.815854in}{0.587977in}}%
\pgfpathlineto{\pgfqpoint{0.860723in}{0.605897in}}%
\pgfpathlineto{\pgfqpoint{0.923033in}{0.631931in}}%
\pgfpathlineto{\pgfqpoint{0.961915in}{0.646633in}}%
\pgfpathlineto{\pgfqpoint{0.999071in}{0.658661in}}%
\pgfpathlineto{\pgfqpoint{1.052126in}{0.674698in}}%
\pgfpathlineto{\pgfqpoint{1.069199in}{0.680898in}}%
\pgfpathlineto{\pgfqpoint{1.086001in}{0.688150in}}%
\pgfpathlineto{\pgfqpoint{1.102550in}{0.696655in}}%
\pgfpathlineto{\pgfqpoint{1.118865in}{0.706313in}}%
\pgfpathlineto{\pgfqpoint{1.150847in}{0.726438in}}%
\pgfpathlineto{\pgfqpoint{1.166541in}{0.734346in}}%
\pgfpathlineto{\pgfqpoint{1.182055in}{0.738928in}}%
\pgfpathlineto{\pgfqpoint{1.212579in}{0.741623in}}%
\pgfpathlineto{\pgfqpoint{1.227610in}{0.746565in}}%
\pgfpathlineto{\pgfqpoint{1.257252in}{0.762695in}}%
\pgfpathlineto{\pgfqpoint{1.271878in}{0.769086in}}%
\pgfpathlineto{\pgfqpoint{1.286384in}{0.776747in}}%
\pgfpathlineto{\pgfqpoint{1.300777in}{0.787569in}}%
\pgfpathlineto{\pgfqpoint{1.329244in}{0.813888in}}%
\pgfpathlineto{\pgfqpoint{1.343330in}{0.825325in}}%
\pgfpathlineto{\pgfqpoint{1.357325in}{0.834342in}}%
\pgfpathlineto{\pgfqpoint{1.371232in}{0.841520in}}%
\pgfpathlineto{\pgfqpoint{1.398802in}{0.854434in}}%
\pgfpathlineto{\pgfqpoint{1.412473in}{0.860087in}}%
\pgfpathlineto{\pgfqpoint{1.426073in}{0.863403in}}%
\pgfpathlineto{\pgfqpoint{1.439605in}{0.863796in}}%
\pgfpathlineto{\pgfqpoint{1.493118in}{0.858941in}}%
\pgfpathlineto{\pgfqpoint{1.532684in}{0.862059in}}%
\pgfpathlineto{\pgfqpoint{1.558826in}{0.861818in}}%
\pgfpathlineto{\pgfqpoint{1.584799in}{0.860513in}}%
\pgfpathlineto{\pgfqpoint{1.597727in}{0.859856in}}%
\pgfpathlineto{\pgfqpoint{1.610618in}{0.860501in}}%
\pgfpathlineto{\pgfqpoint{1.623474in}{0.862871in}}%
\pgfpathlineto{\pgfqpoint{1.636296in}{0.866603in}}%
\pgfpathlineto{\pgfqpoint{1.661845in}{0.876106in}}%
\pgfpathlineto{\pgfqpoint{1.687276in}{0.887429in}}%
\pgfpathlineto{\pgfqpoint{1.699950in}{0.893851in}}%
\pgfpathlineto{\pgfqpoint{1.712599in}{0.900953in}}%
\pgfpathlineto{\pgfqpoint{1.725222in}{0.908865in}}%
\pgfpathlineto{\pgfqpoint{1.750399in}{0.926780in}}%
\pgfpathlineto{\pgfqpoint{1.775489in}{0.946438in}}%
\pgfpathlineto{\pgfqpoint{1.800499in}{0.967674in}}%
\pgfpathlineto{\pgfqpoint{1.837878in}{1.001912in}}%
\pgfpathlineto{\pgfqpoint{1.875110in}{1.037415in}}%
\pgfpathlineto{\pgfqpoint{1.936885in}{1.096831in}}%
\pgfpathlineto{\pgfqpoint{1.961510in}{1.119133in}}%
\pgfpathlineto{\pgfqpoint{1.986092in}{1.140313in}}%
\pgfpathlineto{\pgfqpoint{2.010634in}{1.160316in}}%
\pgfpathlineto{\pgfqpoint{2.035140in}{1.178870in}}%
\pgfpathlineto{\pgfqpoint{2.059613in}{1.195918in}}%
\pgfpathlineto{\pgfqpoint{2.084054in}{1.211692in}}%
\pgfpathlineto{\pgfqpoint{2.108467in}{1.226207in}}%
\pgfpathlineto{\pgfqpoint{2.132853in}{1.239402in}}%
\pgfpathlineto{\pgfqpoint{2.157216in}{1.251526in}}%
\pgfpathlineto{\pgfqpoint{2.193717in}{1.268213in}}%
\pgfpathlineto{\pgfqpoint{2.230174in}{1.283332in}}%
\pgfpathlineto{\pgfqpoint{2.254457in}{1.292275in}}%
\pgfpathlineto{\pgfqpoint{2.290850in}{1.304139in}}%
\pgfpathlineto{\pgfqpoint{2.339324in}{1.318088in}}%
\pgfpathlineto{\pgfqpoint{2.387726in}{1.330748in}}%
\pgfpathlineto{\pgfqpoint{2.423306in}{1.339324in}}%
\pgfpathlineto{\pgfqpoint{2.456205in}{1.346220in}}%
\pgfpathlineto{\pgfqpoint{2.504323in}{1.354856in}}%
\pgfpathlineto{\pgfqpoint{2.627397in}{1.375045in}}%
\pgfpathlineto{\pgfqpoint{2.696521in}{1.383572in}}%
\pgfpathlineto{\pgfqpoint{2.736629in}{1.388020in}}%
\pgfpathlineto{\pgfqpoint{2.772095in}{1.391098in}}%
\pgfpathlineto{\pgfqpoint{2.813787in}{1.393414in}}%
\pgfpathlineto{\pgfqpoint{2.859014in}{1.394258in}}%
\pgfpathlineto{\pgfqpoint{2.937710in}{1.394427in}}%
\pgfpathlineto{\pgfqpoint{3.043110in}{1.394429in}}%
\pgfpathlineto{\pgfqpoint{3.043110in}{1.394429in}}%
\pgfusepath{stroke}%
\end{pgfscope}%
\begin{pgfscope}%
\pgfpathrectangle{\pgfqpoint{0.650000in}{0.420000in}}{\pgfqpoint{2.388110in}{1.994488in}}%
\pgfusepath{clip}%
\pgfsetbuttcap%
\pgfsetroundjoin%
\pgfsetlinewidth{1.003750pt}%
\definecolor{currentstroke}{rgb}{0.870588,0.466667,0.682353}%
\pgfsetstrokecolor{currentstroke}%
\pgfsetdash{{1.000000pt}{1.000000pt}}{0.000000pt}%
\pgfpathmoveto{\pgfqpoint{0.645000in}{1.428925in}}%
\pgfpathlineto{\pgfqpoint{0.714975in}{1.435768in}}%
\pgfpathlineto{\pgfqpoint{0.838669in}{1.449283in}}%
\pgfpathlineto{\pgfqpoint{0.961915in}{1.462661in}}%
\pgfpathlineto{\pgfqpoint{1.017086in}{1.467733in}}%
\pgfpathlineto{\pgfqpoint{1.069199in}{1.471480in}}%
\pgfpathlineto{\pgfqpoint{1.118865in}{1.473768in}}%
\pgfpathlineto{\pgfqpoint{1.150847in}{1.474441in}}%
\pgfpathlineto{\pgfqpoint{1.182055in}{1.474201in}}%
\pgfpathlineto{\pgfqpoint{1.242498in}{1.471478in}}%
\pgfpathlineto{\pgfqpoint{1.286384in}{1.468851in}}%
\pgfpathlineto{\pgfqpoint{1.315062in}{1.466211in}}%
\pgfpathlineto{\pgfqpoint{1.357325in}{1.460788in}}%
\pgfpathlineto{\pgfqpoint{1.412473in}{1.452490in}}%
\pgfpathlineto{\pgfqpoint{1.453072in}{1.444828in}}%
\pgfpathlineto{\pgfqpoint{1.506356in}{1.433042in}}%
\pgfpathlineto{\pgfqpoint{1.545777in}{1.423429in}}%
\pgfpathlineto{\pgfqpoint{1.661845in}{1.394032in}}%
\pgfpathlineto{\pgfqpoint{1.687276in}{1.389263in}}%
\pgfpathlineto{\pgfqpoint{1.712599in}{1.385609in}}%
\pgfpathlineto{\pgfqpoint{1.737822in}{1.383280in}}%
\pgfpathlineto{\pgfqpoint{1.762955in}{1.382285in}}%
\pgfpathlineto{\pgfqpoint{1.788004in}{1.382444in}}%
\pgfpathlineto{\pgfqpoint{1.825435in}{1.383956in}}%
\pgfpathlineto{\pgfqpoint{1.875110in}{1.388274in}}%
\pgfpathlineto{\pgfqpoint{1.912213in}{1.391463in}}%
\pgfpathlineto{\pgfqpoint{1.961510in}{1.394449in}}%
\pgfpathlineto{\pgfqpoint{2.010634in}{1.395608in}}%
\pgfpathlineto{\pgfqpoint{2.047381in}{1.395829in}}%
\pgfpathlineto{\pgfqpoint{2.181555in}{1.394118in}}%
\pgfpathlineto{\pgfqpoint{2.230174in}{1.392584in}}%
\pgfpathlineto{\pgfqpoint{2.315094in}{1.390581in}}%
\pgfpathlineto{\pgfqpoint{2.351435in}{1.392055in}}%
\pgfpathlineto{\pgfqpoint{2.387726in}{1.393486in}}%
\pgfpathlineto{\pgfqpoint{2.545907in}{1.394339in}}%
\pgfpathlineto{\pgfqpoint{2.568397in}{1.393551in}}%
\pgfpathlineto{\pgfqpoint{2.589350in}{1.394150in}}%
\pgfpathlineto{\pgfqpoint{2.608963in}{1.393113in}}%
\pgfpathlineto{\pgfqpoint{2.650370in}{1.393741in}}%
\pgfpathlineto{\pgfqpoint{2.671743in}{1.393463in}}%
\pgfpathlineto{\pgfqpoint{2.701245in}{1.392901in}}%
\pgfpathlineto{\pgfqpoint{2.803883in}{1.394014in}}%
\pgfpathlineto{\pgfqpoint{2.813787in}{1.394308in}}%
\pgfpathlineto{\pgfqpoint{2.829615in}{1.393912in}}%
\pgfpathlineto{\pgfqpoint{2.861804in}{1.395173in}}%
\pgfpathlineto{\pgfqpoint{2.875389in}{1.394244in}}%
\pgfpathlineto{\pgfqpoint{2.893451in}{1.394163in}}%
\pgfpathlineto{\pgfqpoint{2.910509in}{1.394671in}}%
\pgfpathlineto{\pgfqpoint{2.976382in}{1.394472in}}%
\pgfpathlineto{\pgfqpoint{2.985797in}{1.394429in}}%
\pgfpathlineto{\pgfqpoint{2.993127in}{1.395172in}}%
\pgfpathlineto{\pgfqpoint{2.996727in}{1.395083in}}%
\pgfpathlineto{\pgfqpoint{3.003804in}{1.394334in}}%
\pgfpathlineto{\pgfqpoint{3.012425in}{1.394682in}}%
\pgfpathlineto{\pgfqpoint{3.019153in}{1.394671in}}%
\pgfpathlineto{\pgfqpoint{3.025735in}{1.394480in}}%
\pgfpathlineto{\pgfqpoint{3.043110in}{1.394297in}}%
\pgfpathlineto{\pgfqpoint{3.043110in}{1.394297in}}%
\pgfusepath{stroke}%
\end{pgfscope}%
\begin{pgfscope}%
\pgfpathrectangle{\pgfqpoint{0.650000in}{0.420000in}}{\pgfqpoint{2.388110in}{1.994488in}}%
\pgfusepath{clip}%
\pgfsetbuttcap%
\pgfsetroundjoin%
\pgfsetlinewidth{1.003750pt}%
\definecolor{currentstroke}{rgb}{0.870588,0.466667,0.682353}%
\pgfsetstrokecolor{currentstroke}%
\pgfsetdash{{5.000000pt}{5.000000pt}}{0.000000pt}%
\pgfpathmoveto{\pgfqpoint{0.645000in}{1.395178in}}%
\pgfpathlineto{\pgfqpoint{0.714975in}{1.396007in}}%
\pgfpathlineto{\pgfqpoint{0.792190in}{1.398087in}}%
\pgfpathlineto{\pgfqpoint{0.838669in}{1.400302in}}%
\pgfpathlineto{\pgfqpoint{0.882094in}{1.403358in}}%
\pgfpathlineto{\pgfqpoint{0.923033in}{1.407468in}}%
\pgfpathlineto{\pgfqpoint{0.961915in}{1.412941in}}%
\pgfpathlineto{\pgfqpoint{0.980691in}{1.416314in}}%
\pgfpathlineto{\pgfqpoint{0.999071in}{1.420194in}}%
\pgfpathlineto{\pgfqpoint{1.017086in}{1.424645in}}%
\pgfpathlineto{\pgfqpoint{1.034762in}{1.429737in}}%
\pgfpathlineto{\pgfqpoint{1.052126in}{1.435526in}}%
\pgfpathlineto{\pgfqpoint{1.069199in}{1.442064in}}%
\pgfpathlineto{\pgfqpoint{1.086001in}{1.449379in}}%
\pgfpathlineto{\pgfqpoint{1.102550in}{1.457467in}}%
\pgfpathlineto{\pgfqpoint{1.118865in}{1.466290in}}%
\pgfpathlineto{\pgfqpoint{1.134959in}{1.475775in}}%
\pgfpathlineto{\pgfqpoint{1.166541in}{1.496256in}}%
\pgfpathlineto{\pgfqpoint{1.197397in}{1.518033in}}%
\pgfpathlineto{\pgfqpoint{1.227610in}{1.541041in}}%
\pgfpathlineto{\pgfqpoint{1.242498in}{1.552526in}}%
\pgfpathlineto{\pgfqpoint{1.257252in}{1.563092in}}%
\pgfpathlineto{\pgfqpoint{1.271878in}{1.572272in}}%
\pgfpathlineto{\pgfqpoint{1.286384in}{1.579954in}}%
\pgfpathlineto{\pgfqpoint{1.300777in}{1.586011in}}%
\pgfpathlineto{\pgfqpoint{1.315062in}{1.590198in}}%
\pgfpathlineto{\pgfqpoint{1.329244in}{1.592194in}}%
\pgfpathlineto{\pgfqpoint{1.343330in}{1.591713in}}%
\pgfpathlineto{\pgfqpoint{1.357325in}{1.588555in}}%
\pgfpathlineto{\pgfqpoint{1.371232in}{1.582458in}}%
\pgfpathlineto{\pgfqpoint{1.385056in}{1.573136in}}%
\pgfpathlineto{\pgfqpoint{1.398802in}{1.560345in}}%
\pgfpathlineto{\pgfqpoint{1.412473in}{1.544014in}}%
\pgfpathlineto{\pgfqpoint{1.426073in}{1.524445in}}%
\pgfpathlineto{\pgfqpoint{1.439605in}{1.502072in}}%
\pgfpathlineto{\pgfqpoint{1.453072in}{1.477201in}}%
\pgfpathlineto{\pgfqpoint{1.466478in}{1.449967in}}%
\pgfpathlineto{\pgfqpoint{1.479826in}{1.420643in}}%
\pgfpathlineto{\pgfqpoint{1.519544in}{1.328100in}}%
\pgfpathlineto{\pgfqpoint{1.532684in}{1.299726in}}%
\pgfpathlineto{\pgfqpoint{1.558826in}{1.247574in}}%
\pgfpathlineto{\pgfqpoint{1.571832in}{1.223275in}}%
\pgfpathlineto{\pgfqpoint{1.584799in}{1.202576in}}%
\pgfpathlineto{\pgfqpoint{1.597727in}{1.185561in}}%
\pgfpathlineto{\pgfqpoint{1.610618in}{1.170723in}}%
\pgfpathlineto{\pgfqpoint{1.623474in}{1.159571in}}%
\pgfpathlineto{\pgfqpoint{1.636296in}{1.154819in}}%
\pgfpathlineto{\pgfqpoint{1.649086in}{1.155546in}}%
\pgfpathlineto{\pgfqpoint{1.661845in}{1.158394in}}%
\pgfpathlineto{\pgfqpoint{1.674574in}{1.161859in}}%
\pgfpathlineto{\pgfqpoint{1.687276in}{1.167314in}}%
\pgfpathlineto{\pgfqpoint{1.699950in}{1.176046in}}%
\pgfpathlineto{\pgfqpoint{1.737822in}{1.208717in}}%
\pgfpathlineto{\pgfqpoint{1.750399in}{1.220203in}}%
\pgfpathlineto{\pgfqpoint{1.762955in}{1.233977in}}%
\pgfpathlineto{\pgfqpoint{1.788004in}{1.267643in}}%
\pgfpathlineto{\pgfqpoint{1.812976in}{1.294575in}}%
\pgfpathlineto{\pgfqpoint{1.825435in}{1.309568in}}%
\pgfpathlineto{\pgfqpoint{1.837878in}{1.322403in}}%
\pgfpathlineto{\pgfqpoint{1.875110in}{1.350628in}}%
\pgfpathlineto{\pgfqpoint{1.912213in}{1.375926in}}%
\pgfpathlineto{\pgfqpoint{1.924555in}{1.384572in}}%
\pgfpathlineto{\pgfqpoint{1.936885in}{1.389294in}}%
\pgfpathlineto{\pgfqpoint{1.949203in}{1.390881in}}%
\pgfpathlineto{\pgfqpoint{1.961510in}{1.393169in}}%
\pgfpathlineto{\pgfqpoint{1.986092in}{1.401765in}}%
\pgfpathlineto{\pgfqpoint{1.998368in}{1.403100in}}%
\pgfpathlineto{\pgfqpoint{2.010634in}{1.403002in}}%
\pgfpathlineto{\pgfqpoint{2.022892in}{1.404450in}}%
\pgfpathlineto{\pgfqpoint{2.047381in}{1.410635in}}%
\pgfpathlineto{\pgfqpoint{2.059613in}{1.411525in}}%
\pgfpathlineto{\pgfqpoint{2.084054in}{1.412197in}}%
\pgfpathlineto{\pgfqpoint{2.096264in}{1.413104in}}%
\pgfpathlineto{\pgfqpoint{2.108467in}{1.414616in}}%
\pgfpathlineto{\pgfqpoint{2.132853in}{1.419373in}}%
\pgfpathlineto{\pgfqpoint{2.145037in}{1.419751in}}%
\pgfpathlineto{\pgfqpoint{2.169388in}{1.415242in}}%
\pgfpathlineto{\pgfqpoint{2.193717in}{1.418028in}}%
\pgfpathlineto{\pgfqpoint{2.218027in}{1.417161in}}%
\pgfpathlineto{\pgfqpoint{2.230174in}{1.415682in}}%
\pgfpathlineto{\pgfqpoint{2.242318in}{1.411150in}}%
\pgfpathlineto{\pgfqpoint{2.254457in}{1.405621in}}%
\pgfpathlineto{\pgfqpoint{2.266591in}{1.403137in}}%
\pgfpathlineto{\pgfqpoint{2.278722in}{1.404195in}}%
\pgfpathlineto{\pgfqpoint{2.302973in}{1.408713in}}%
\pgfpathlineto{\pgfqpoint{2.315094in}{1.409110in}}%
\pgfpathlineto{\pgfqpoint{2.327210in}{1.406287in}}%
\pgfpathlineto{\pgfqpoint{2.351435in}{1.394721in}}%
\pgfpathlineto{\pgfqpoint{2.363542in}{1.391906in}}%
\pgfpathlineto{\pgfqpoint{2.375643in}{1.393067in}}%
\pgfpathlineto{\pgfqpoint{2.387726in}{1.396590in}}%
\pgfpathlineto{\pgfqpoint{2.399754in}{1.398655in}}%
\pgfpathlineto{\pgfqpoint{2.411648in}{1.396117in}}%
\pgfpathlineto{\pgfqpoint{2.423306in}{1.390564in}}%
\pgfpathlineto{\pgfqpoint{2.434639in}{1.387377in}}%
\pgfpathlineto{\pgfqpoint{2.445607in}{1.388978in}}%
\pgfpathlineto{\pgfqpoint{2.456205in}{1.391592in}}%
\pgfpathlineto{\pgfqpoint{2.466451in}{1.391178in}}%
\pgfpathlineto{\pgfqpoint{2.476366in}{1.388561in}}%
\pgfpathlineto{\pgfqpoint{2.485971in}{1.386904in}}%
\pgfpathlineto{\pgfqpoint{2.495284in}{1.387574in}}%
\pgfpathlineto{\pgfqpoint{2.504323in}{1.389708in}}%
\pgfpathlineto{\pgfqpoint{2.513103in}{1.392471in}}%
\pgfpathlineto{\pgfqpoint{2.521638in}{1.395963in}}%
\pgfpathlineto{\pgfqpoint{2.529943in}{1.400698in}}%
\pgfpathlineto{\pgfqpoint{2.538029in}{1.406140in}}%
\pgfpathlineto{\pgfqpoint{2.545907in}{1.408953in}}%
\pgfpathlineto{\pgfqpoint{2.553588in}{1.406865in}}%
\pgfpathlineto{\pgfqpoint{2.561082in}{1.403379in}}%
\pgfpathlineto{\pgfqpoint{2.568397in}{1.403065in}}%
\pgfpathlineto{\pgfqpoint{2.575541in}{1.405101in}}%
\pgfpathlineto{\pgfqpoint{2.582523in}{1.404918in}}%
\pgfpathlineto{\pgfqpoint{2.589350in}{1.401290in}}%
\pgfpathlineto{\pgfqpoint{2.596028in}{1.398512in}}%
\pgfpathlineto{\pgfqpoint{2.602563in}{1.399472in}}%
\pgfpathlineto{\pgfqpoint{2.608963in}{1.401475in}}%
\pgfpathlineto{\pgfqpoint{2.615231in}{1.401137in}}%
\pgfpathlineto{\pgfqpoint{2.621374in}{1.399020in}}%
\pgfpathlineto{\pgfqpoint{2.627397in}{1.397712in}}%
\pgfpathlineto{\pgfqpoint{2.639099in}{1.398652in}}%
\pgfpathlineto{\pgfqpoint{2.644786in}{1.399711in}}%
\pgfpathlineto{\pgfqpoint{2.655855in}{1.403540in}}%
\pgfpathlineto{\pgfqpoint{2.661243in}{1.404253in}}%
\pgfpathlineto{\pgfqpoint{2.666538in}{1.402791in}}%
\pgfpathlineto{\pgfqpoint{2.676861in}{1.396603in}}%
\pgfpathlineto{\pgfqpoint{2.681895in}{1.396504in}}%
\pgfpathlineto{\pgfqpoint{2.686848in}{1.398123in}}%
\pgfpathlineto{\pgfqpoint{2.691723in}{1.397811in}}%
\pgfpathlineto{\pgfqpoint{2.701245in}{1.390926in}}%
\pgfpathlineto{\pgfqpoint{2.705898in}{1.390802in}}%
\pgfpathlineto{\pgfqpoint{2.710481in}{1.394249in}}%
\pgfpathlineto{\pgfqpoint{2.714997in}{1.398857in}}%
\pgfpathlineto{\pgfqpoint{2.719447in}{1.401985in}}%
\pgfpathlineto{\pgfqpoint{2.723834in}{1.402452in}}%
\pgfpathlineto{\pgfqpoint{2.732423in}{1.401226in}}%
\pgfpathlineto{\pgfqpoint{2.736629in}{1.400851in}}%
\pgfpathlineto{\pgfqpoint{2.740778in}{1.398976in}}%
\pgfpathlineto{\pgfqpoint{2.752900in}{1.390736in}}%
\pgfpathlineto{\pgfqpoint{2.756836in}{1.389270in}}%
\pgfpathlineto{\pgfqpoint{2.760723in}{1.391733in}}%
\pgfpathlineto{\pgfqpoint{2.768351in}{1.403030in}}%
\pgfpathlineto{\pgfqpoint{2.772095in}{1.406065in}}%
\pgfpathlineto{\pgfqpoint{2.775794in}{1.404565in}}%
\pgfpathlineto{\pgfqpoint{2.783061in}{1.396029in}}%
\pgfpathlineto{\pgfqpoint{2.786631in}{1.394244in}}%
\pgfpathlineto{\pgfqpoint{2.793649in}{1.392643in}}%
\pgfpathlineto{\pgfqpoint{2.797099in}{1.391512in}}%
\pgfpathlineto{\pgfqpoint{2.800510in}{1.391051in}}%
\pgfpathlineto{\pgfqpoint{2.803883in}{1.392599in}}%
\pgfpathlineto{\pgfqpoint{2.810521in}{1.399912in}}%
\pgfpathlineto{\pgfqpoint{2.813787in}{1.399493in}}%
\pgfpathlineto{\pgfqpoint{2.820217in}{1.389869in}}%
\pgfpathlineto{\pgfqpoint{2.823382in}{1.388713in}}%
\pgfpathlineto{\pgfqpoint{2.826514in}{1.390672in}}%
\pgfpathlineto{\pgfqpoint{2.835725in}{1.398718in}}%
\pgfpathlineto{\pgfqpoint{2.838734in}{1.399848in}}%
\pgfpathlineto{\pgfqpoint{2.841715in}{1.399113in}}%
\pgfpathlineto{\pgfqpoint{2.844667in}{1.397645in}}%
\pgfpathlineto{\pgfqpoint{2.853355in}{1.394856in}}%
\pgfpathlineto{\pgfqpoint{2.856198in}{1.392028in}}%
\pgfpathlineto{\pgfqpoint{2.859014in}{1.391764in}}%
\pgfpathlineto{\pgfqpoint{2.861804in}{1.395125in}}%
\pgfpathlineto{\pgfqpoint{2.864570in}{1.397331in}}%
\pgfpathlineto{\pgfqpoint{2.867311in}{1.395287in}}%
\pgfpathlineto{\pgfqpoint{2.870027in}{1.392215in}}%
\pgfpathlineto{\pgfqpoint{2.872720in}{1.391771in}}%
\pgfpathlineto{\pgfqpoint{2.875389in}{1.393088in}}%
\pgfpathlineto{\pgfqpoint{2.878036in}{1.393268in}}%
\pgfpathlineto{\pgfqpoint{2.885840in}{1.389112in}}%
\pgfpathlineto{\pgfqpoint{2.890935in}{1.386786in}}%
\pgfpathlineto{\pgfqpoint{2.893451in}{1.387973in}}%
\pgfpathlineto{\pgfqpoint{2.898422in}{1.392612in}}%
\pgfpathlineto{\pgfqpoint{2.900878in}{1.393392in}}%
\pgfpathlineto{\pgfqpoint{2.903314in}{1.392472in}}%
\pgfpathlineto{\pgfqpoint{2.910509in}{1.384967in}}%
\pgfpathlineto{\pgfqpoint{2.912870in}{1.385746in}}%
\pgfpathlineto{\pgfqpoint{2.924410in}{1.397399in}}%
\pgfpathlineto{\pgfqpoint{2.926667in}{1.398728in}}%
\pgfpathlineto{\pgfqpoint{2.931132in}{1.396486in}}%
\pgfpathlineto{\pgfqpoint{2.933340in}{1.397677in}}%
\pgfpathlineto{\pgfqpoint{2.937710in}{1.404192in}}%
\pgfpathlineto{\pgfqpoint{2.939871in}{1.401900in}}%
\pgfpathlineto{\pgfqpoint{2.944149in}{1.390187in}}%
\pgfpathlineto{\pgfqpoint{2.946266in}{1.389986in}}%
\pgfpathlineto{\pgfqpoint{2.948368in}{1.391866in}}%
\pgfpathlineto{\pgfqpoint{2.950456in}{1.392171in}}%
\pgfpathlineto{\pgfqpoint{2.956636in}{1.391161in}}%
\pgfpathlineto{\pgfqpoint{2.960688in}{1.387499in}}%
\pgfpathlineto{\pgfqpoint{2.962695in}{1.389519in}}%
\pgfpathlineto{\pgfqpoint{2.964688in}{1.395879in}}%
\pgfpathlineto{\pgfqpoint{2.966668in}{1.404137in}}%
\pgfpathlineto{\pgfqpoint{2.968635in}{1.409948in}}%
\pgfpathlineto{\pgfqpoint{2.970590in}{1.410129in}}%
\pgfpathlineto{\pgfqpoint{2.972533in}{1.404798in}}%
\pgfpathlineto{\pgfqpoint{2.976382in}{1.390785in}}%
\pgfpathlineto{\pgfqpoint{2.978288in}{1.388121in}}%
\pgfpathlineto{\pgfqpoint{2.980183in}{1.388701in}}%
\pgfpathlineto{\pgfqpoint{2.982066in}{1.391212in}}%
\pgfpathlineto{\pgfqpoint{2.987646in}{1.401287in}}%
\pgfpathlineto{\pgfqpoint{2.991311in}{1.405549in}}%
\pgfpathlineto{\pgfqpoint{2.993127in}{1.409576in}}%
\pgfpathlineto{\pgfqpoint{2.994932in}{1.411839in}}%
\pgfpathlineto{\pgfqpoint{2.996727in}{1.409530in}}%
\pgfpathlineto{\pgfqpoint{3.000286in}{1.402986in}}%
\pgfpathlineto{\pgfqpoint{3.002050in}{1.399534in}}%
\pgfpathlineto{\pgfqpoint{3.005547in}{1.391094in}}%
\pgfpathlineto{\pgfqpoint{3.007281in}{1.391593in}}%
\pgfpathlineto{\pgfqpoint{3.009006in}{1.394092in}}%
\pgfpathlineto{\pgfqpoint{3.010720in}{1.394801in}}%
\pgfpathlineto{\pgfqpoint{3.012425in}{1.393899in}}%
\pgfpathlineto{\pgfqpoint{3.014121in}{1.395250in}}%
\pgfpathlineto{\pgfqpoint{3.015807in}{1.399113in}}%
\pgfpathlineto{\pgfqpoint{3.017485in}{1.401214in}}%
\pgfpathlineto{\pgfqpoint{3.019153in}{1.399935in}}%
\pgfpathlineto{\pgfqpoint{3.024103in}{1.392611in}}%
\pgfpathlineto{\pgfqpoint{3.025735in}{1.391372in}}%
\pgfpathlineto{\pgfqpoint{3.027359in}{1.392700in}}%
\pgfpathlineto{\pgfqpoint{3.030581in}{1.400572in}}%
\pgfpathlineto{\pgfqpoint{3.032180in}{1.401715in}}%
\pgfpathlineto{\pgfqpoint{3.038492in}{1.395647in}}%
\pgfpathlineto{\pgfqpoint{3.041600in}{1.387036in}}%
\pgfpathlineto{\pgfqpoint{3.043110in}{1.385462in}}%
\pgfpathlineto{\pgfqpoint{3.043110in}{1.385462in}}%
\pgfusepath{stroke}%
\end{pgfscope}%
\begin{pgfscope}%
\pgfpathrectangle{\pgfqpoint{0.650000in}{0.420000in}}{\pgfqpoint{2.388110in}{1.994488in}}%
\pgfusepath{clip}%
\pgfsetbuttcap%
\pgfsetroundjoin%
\pgfsetlinewidth{1.003750pt}%
\definecolor{currentstroke}{rgb}{0.870588,0.466667,0.682353}%
\pgfsetstrokecolor{currentstroke}%
\pgfsetdash{{3.000000pt}{5.000000pt}{1.000000pt}{5.000000pt}}{0.000000pt}%
\pgfpathmoveto{\pgfqpoint{0.645000in}{2.180999in}}%
\pgfpathlineto{\pgfqpoint{0.656764in}{2.181433in}}%
\pgfpathlineto{\pgfqpoint{0.686665in}{2.181703in}}%
\pgfpathlineto{\pgfqpoint{0.714975in}{2.180268in}}%
\pgfpathlineto{\pgfqpoint{0.741886in}{2.176459in}}%
\pgfpathlineto{\pgfqpoint{0.767576in}{2.169824in}}%
\pgfpathlineto{\pgfqpoint{0.792190in}{2.160162in}}%
\pgfpathlineto{\pgfqpoint{0.815854in}{2.147542in}}%
\pgfpathlineto{\pgfqpoint{0.838669in}{2.132296in}}%
\pgfpathlineto{\pgfqpoint{0.860723in}{2.115005in}}%
\pgfpathlineto{\pgfqpoint{0.882094in}{2.096455in}}%
\pgfpathlineto{\pgfqpoint{0.923033in}{2.059318in}}%
\pgfpathlineto{\pgfqpoint{0.942709in}{2.042602in}}%
\pgfpathlineto{\pgfqpoint{0.961915in}{2.028124in}}%
\pgfpathlineto{\pgfqpoint{0.980691in}{2.016249in}}%
\pgfpathlineto{\pgfqpoint{0.999071in}{2.006877in}}%
\pgfpathlineto{\pgfqpoint{1.052126in}{1.984338in}}%
\pgfpathlineto{\pgfqpoint{1.069199in}{1.972971in}}%
\pgfpathlineto{\pgfqpoint{1.086001in}{1.956232in}}%
\pgfpathlineto{\pgfqpoint{1.102550in}{1.932532in}}%
\pgfpathlineto{\pgfqpoint{1.118865in}{1.901401in}}%
\pgfpathlineto{\pgfqpoint{1.134959in}{1.864043in}}%
\pgfpathlineto{\pgfqpoint{1.166541in}{1.784109in}}%
\pgfpathlineto{\pgfqpoint{1.182055in}{1.748759in}}%
\pgfpathlineto{\pgfqpoint{1.197397in}{1.715977in}}%
\pgfpathlineto{\pgfqpoint{1.212579in}{1.679930in}}%
\pgfpathlineto{\pgfqpoint{1.242498in}{1.598581in}}%
\pgfpathlineto{\pgfqpoint{1.257252in}{1.568941in}}%
\pgfpathlineto{\pgfqpoint{1.271878in}{1.545879in}}%
\pgfpathlineto{\pgfqpoint{1.286384in}{1.518112in}}%
\pgfpathlineto{\pgfqpoint{1.300777in}{1.478993in}}%
\pgfpathlineto{\pgfqpoint{1.329244in}{1.381880in}}%
\pgfpathlineto{\pgfqpoint{1.343330in}{1.338693in}}%
\pgfpathlineto{\pgfqpoint{1.357325in}{1.303929in}}%
\pgfpathlineto{\pgfqpoint{1.385056in}{1.245490in}}%
\pgfpathlineto{\pgfqpoint{1.398802in}{1.212765in}}%
\pgfpathlineto{\pgfqpoint{1.412473in}{1.178071in}}%
\pgfpathlineto{\pgfqpoint{1.426073in}{1.146306in}}%
\pgfpathlineto{\pgfqpoint{1.439605in}{1.121737in}}%
\pgfpathlineto{\pgfqpoint{1.453072in}{1.107373in}}%
\pgfpathlineto{\pgfqpoint{1.466478in}{1.104855in}}%
\pgfpathlineto{\pgfqpoint{1.493118in}{1.117110in}}%
\pgfpathlineto{\pgfqpoint{1.519544in}{1.123239in}}%
\pgfpathlineto{\pgfqpoint{1.532684in}{1.130901in}}%
\pgfpathlineto{\pgfqpoint{1.571832in}{1.166838in}}%
\pgfpathlineto{\pgfqpoint{1.584799in}{1.180504in}}%
\pgfpathlineto{\pgfqpoint{1.597727in}{1.198161in}}%
\pgfpathlineto{\pgfqpoint{1.610618in}{1.218168in}}%
\pgfpathlineto{\pgfqpoint{1.623474in}{1.237076in}}%
\pgfpathlineto{\pgfqpoint{1.636296in}{1.253368in}}%
\pgfpathlineto{\pgfqpoint{1.661845in}{1.282976in}}%
\pgfpathlineto{\pgfqpoint{1.674574in}{1.296819in}}%
\pgfpathlineto{\pgfqpoint{1.699950in}{1.322467in}}%
\pgfpathlineto{\pgfqpoint{1.712599in}{1.334419in}}%
\pgfpathlineto{\pgfqpoint{1.725222in}{1.344510in}}%
\pgfpathlineto{\pgfqpoint{1.737822in}{1.352838in}}%
\pgfpathlineto{\pgfqpoint{1.762955in}{1.367315in}}%
\pgfpathlineto{\pgfqpoint{1.775489in}{1.373713in}}%
\pgfpathlineto{\pgfqpoint{1.788004in}{1.379432in}}%
\pgfpathlineto{\pgfqpoint{1.800499in}{1.384159in}}%
\pgfpathlineto{\pgfqpoint{1.812976in}{1.387759in}}%
\pgfpathlineto{\pgfqpoint{1.837878in}{1.393309in}}%
\pgfpathlineto{\pgfqpoint{1.862714in}{1.397831in}}%
\pgfpathlineto{\pgfqpoint{1.875110in}{1.399185in}}%
\pgfpathlineto{\pgfqpoint{1.899859in}{1.400090in}}%
\pgfpathlineto{\pgfqpoint{1.949203in}{1.400339in}}%
\pgfpathlineto{\pgfqpoint{1.986092in}{1.399459in}}%
\pgfpathlineto{\pgfqpoint{2.047381in}{1.396700in}}%
\pgfpathlineto{\pgfqpoint{2.120663in}{1.395538in}}%
\pgfpathlineto{\pgfqpoint{2.169388in}{1.394823in}}%
\pgfpathlineto{\pgfqpoint{2.445607in}{1.394295in}}%
\pgfpathlineto{\pgfqpoint{2.633304in}{1.394446in}}%
\pgfpathlineto{\pgfqpoint{3.043110in}{1.394429in}}%
\pgfpathlineto{\pgfqpoint{3.043110in}{1.394429in}}%
\pgfusepath{stroke}%
\end{pgfscope}%
\begin{pgfscope}%
\pgfpathrectangle{\pgfqpoint{0.650000in}{0.420000in}}{\pgfqpoint{2.388110in}{1.994488in}}%
\pgfusepath{clip}%
\pgfsetbuttcap%
\pgfsetroundjoin%
\pgfsetlinewidth{1.003750pt}%
\definecolor{currentstroke}{rgb}{0.870588,0.466667,0.682353}%
\pgfsetstrokecolor{currentstroke}%
\pgfsetdash{{1.000000pt}{5.000000pt}}{0.000000pt}%
\pgfpathmoveto{\pgfqpoint{0.645000in}{1.368282in}}%
\pgfpathlineto{\pgfqpoint{0.656764in}{1.371361in}}%
\pgfpathlineto{\pgfqpoint{0.714975in}{1.387487in}}%
\pgfpathlineto{\pgfqpoint{0.741886in}{1.393842in}}%
\pgfpathlineto{\pgfqpoint{0.767576in}{1.398323in}}%
\pgfpathlineto{\pgfqpoint{0.792190in}{1.400646in}}%
\pgfpathlineto{\pgfqpoint{0.815854in}{1.400757in}}%
\pgfpathlineto{\pgfqpoint{0.838669in}{1.398842in}}%
\pgfpathlineto{\pgfqpoint{0.860723in}{1.395318in}}%
\pgfpathlineto{\pgfqpoint{0.923033in}{1.381905in}}%
\pgfpathlineto{\pgfqpoint{0.942709in}{1.379188in}}%
\pgfpathlineto{\pgfqpoint{0.961915in}{1.378594in}}%
\pgfpathlineto{\pgfqpoint{0.980691in}{1.380575in}}%
\pgfpathlineto{\pgfqpoint{0.999071in}{1.385231in}}%
\pgfpathlineto{\pgfqpoint{1.017086in}{1.392223in}}%
\pgfpathlineto{\pgfqpoint{1.069199in}{1.416965in}}%
\pgfpathlineto{\pgfqpoint{1.086001in}{1.421485in}}%
\pgfpathlineto{\pgfqpoint{1.102550in}{1.421776in}}%
\pgfpathlineto{\pgfqpoint{1.118865in}{1.417331in}}%
\pgfpathlineto{\pgfqpoint{1.134959in}{1.408853in}}%
\pgfpathlineto{\pgfqpoint{1.166541in}{1.388428in}}%
\pgfpathlineto{\pgfqpoint{1.182055in}{1.380643in}}%
\pgfpathlineto{\pgfqpoint{1.197397in}{1.373673in}}%
\pgfpathlineto{\pgfqpoint{1.227610in}{1.356733in}}%
\pgfpathlineto{\pgfqpoint{1.242498in}{1.355518in}}%
\pgfpathlineto{\pgfqpoint{1.257252in}{1.364907in}}%
\pgfpathlineto{\pgfqpoint{1.271878in}{1.379807in}}%
\pgfpathlineto{\pgfqpoint{1.286384in}{1.391676in}}%
\pgfpathlineto{\pgfqpoint{1.300777in}{1.395769in}}%
\pgfpathlineto{\pgfqpoint{1.315062in}{1.393431in}}%
\pgfpathlineto{\pgfqpoint{1.329244in}{1.389733in}}%
\pgfpathlineto{\pgfqpoint{1.343330in}{1.389492in}}%
\pgfpathlineto{\pgfqpoint{1.357325in}{1.394195in}}%
\pgfpathlineto{\pgfqpoint{1.371232in}{1.401161in}}%
\pgfpathlineto{\pgfqpoint{1.385056in}{1.405594in}}%
\pgfpathlineto{\pgfqpoint{1.398802in}{1.404115in}}%
\pgfpathlineto{\pgfqpoint{1.412473in}{1.397027in}}%
\pgfpathlineto{\pgfqpoint{1.426073in}{1.387256in}}%
\pgfpathlineto{\pgfqpoint{1.439605in}{1.378264in}}%
\pgfpathlineto{\pgfqpoint{1.453072in}{1.374161in}}%
\pgfpathlineto{\pgfqpoint{1.466478in}{1.377778in}}%
\pgfpathlineto{\pgfqpoint{1.479826in}{1.386190in}}%
\pgfpathlineto{\pgfqpoint{1.493118in}{1.392761in}}%
\pgfpathlineto{\pgfqpoint{1.506356in}{1.394529in}}%
\pgfpathlineto{\pgfqpoint{1.519544in}{1.394557in}}%
\pgfpathlineto{\pgfqpoint{1.532684in}{1.396754in}}%
\pgfpathlineto{\pgfqpoint{1.545777in}{1.400109in}}%
\pgfpathlineto{\pgfqpoint{1.558826in}{1.400229in}}%
\pgfpathlineto{\pgfqpoint{1.584799in}{1.392251in}}%
\pgfpathlineto{\pgfqpoint{1.597727in}{1.391310in}}%
\pgfpathlineto{\pgfqpoint{1.623474in}{1.393071in}}%
\pgfpathlineto{\pgfqpoint{1.636296in}{1.395793in}}%
\pgfpathlineto{\pgfqpoint{1.649086in}{1.399868in}}%
\pgfpathlineto{\pgfqpoint{1.661845in}{1.402811in}}%
\pgfpathlineto{\pgfqpoint{1.674574in}{1.402889in}}%
\pgfpathlineto{\pgfqpoint{1.687276in}{1.401790in}}%
\pgfpathlineto{\pgfqpoint{1.712599in}{1.402479in}}%
\pgfpathlineto{\pgfqpoint{1.725222in}{1.401240in}}%
\pgfpathlineto{\pgfqpoint{1.750399in}{1.395697in}}%
\pgfpathlineto{\pgfqpoint{1.762955in}{1.395065in}}%
\pgfpathlineto{\pgfqpoint{1.788004in}{1.399686in}}%
\pgfpathlineto{\pgfqpoint{1.812976in}{1.398461in}}%
\pgfpathlineto{\pgfqpoint{1.825435in}{1.400365in}}%
\pgfpathlineto{\pgfqpoint{1.837878in}{1.401104in}}%
\pgfpathlineto{\pgfqpoint{1.862714in}{1.398102in}}%
\pgfpathlineto{\pgfqpoint{1.875110in}{1.398071in}}%
\pgfpathlineto{\pgfqpoint{1.899859in}{1.397107in}}%
\pgfpathlineto{\pgfqpoint{1.924555in}{1.401047in}}%
\pgfpathlineto{\pgfqpoint{1.936885in}{1.399788in}}%
\pgfpathlineto{\pgfqpoint{1.949203in}{1.396100in}}%
\pgfpathlineto{\pgfqpoint{1.961510in}{1.393884in}}%
\pgfpathlineto{\pgfqpoint{1.986092in}{1.395457in}}%
\pgfpathlineto{\pgfqpoint{1.998368in}{1.394121in}}%
\pgfpathlineto{\pgfqpoint{2.010634in}{1.391977in}}%
\pgfpathlineto{\pgfqpoint{2.022892in}{1.392105in}}%
\pgfpathlineto{\pgfqpoint{2.047381in}{1.397828in}}%
\pgfpathlineto{\pgfqpoint{2.084054in}{1.403007in}}%
\pgfpathlineto{\pgfqpoint{2.096264in}{1.405844in}}%
\pgfpathlineto{\pgfqpoint{2.108467in}{1.409484in}}%
\pgfpathlineto{\pgfqpoint{2.132853in}{1.418248in}}%
\pgfpathlineto{\pgfqpoint{2.145037in}{1.420201in}}%
\pgfpathlineto{\pgfqpoint{2.169388in}{1.417866in}}%
\pgfpathlineto{\pgfqpoint{2.193717in}{1.421736in}}%
\pgfpathlineto{\pgfqpoint{2.205874in}{1.421615in}}%
\pgfpathlineto{\pgfqpoint{2.218027in}{1.420904in}}%
\pgfpathlineto{\pgfqpoint{2.230174in}{1.419298in}}%
\pgfpathlineto{\pgfqpoint{2.242318in}{1.414763in}}%
\pgfpathlineto{\pgfqpoint{2.254457in}{1.409304in}}%
\pgfpathlineto{\pgfqpoint{2.266591in}{1.406883in}}%
\pgfpathlineto{\pgfqpoint{2.278722in}{1.407986in}}%
\pgfpathlineto{\pgfqpoint{2.302973in}{1.412379in}}%
\pgfpathlineto{\pgfqpoint{2.315094in}{1.412694in}}%
\pgfpathlineto{\pgfqpoint{2.327210in}{1.410042in}}%
\pgfpathlineto{\pgfqpoint{2.351435in}{1.400123in}}%
\pgfpathlineto{\pgfqpoint{2.363542in}{1.398505in}}%
\pgfpathlineto{\pgfqpoint{2.375643in}{1.400706in}}%
\pgfpathlineto{\pgfqpoint{2.387726in}{1.404885in}}%
\pgfpathlineto{\pgfqpoint{2.399754in}{1.407117in}}%
\pgfpathlineto{\pgfqpoint{2.411648in}{1.404320in}}%
\pgfpathlineto{\pgfqpoint{2.423306in}{1.398370in}}%
\pgfpathlineto{\pgfqpoint{2.434639in}{1.394948in}}%
\pgfpathlineto{\pgfqpoint{2.445607in}{1.396536in}}%
\pgfpathlineto{\pgfqpoint{2.456205in}{1.399247in}}%
\pgfpathlineto{\pgfqpoint{2.466451in}{1.398902in}}%
\pgfpathlineto{\pgfqpoint{2.476366in}{1.396210in}}%
\pgfpathlineto{\pgfqpoint{2.485971in}{1.394342in}}%
\pgfpathlineto{\pgfqpoint{2.495284in}{1.394793in}}%
\pgfpathlineto{\pgfqpoint{2.504323in}{1.396777in}}%
\pgfpathlineto{\pgfqpoint{2.513103in}{1.399386in}}%
\pgfpathlineto{\pgfqpoint{2.521638in}{1.402623in}}%
\pgfpathlineto{\pgfqpoint{2.529943in}{1.407004in}}%
\pgfpathlineto{\pgfqpoint{2.538029in}{1.411986in}}%
\pgfpathlineto{\pgfqpoint{2.545907in}{1.414080in}}%
\pgfpathlineto{\pgfqpoint{2.553588in}{1.410984in}}%
\pgfpathlineto{\pgfqpoint{2.561082in}{1.406542in}}%
\pgfpathlineto{\pgfqpoint{2.568397in}{1.405737in}}%
\pgfpathlineto{\pgfqpoint{2.575541in}{1.407753in}}%
\pgfpathlineto{\pgfqpoint{2.582523in}{1.407576in}}%
\pgfpathlineto{\pgfqpoint{2.596028in}{1.400213in}}%
\pgfpathlineto{\pgfqpoint{2.602563in}{1.400546in}}%
\pgfpathlineto{\pgfqpoint{2.608963in}{1.402018in}}%
\pgfpathlineto{\pgfqpoint{2.615231in}{1.401237in}}%
\pgfpathlineto{\pgfqpoint{2.621374in}{1.398793in}}%
\pgfpathlineto{\pgfqpoint{2.627397in}{1.397334in}}%
\pgfpathlineto{\pgfqpoint{2.633304in}{1.397563in}}%
\pgfpathlineto{\pgfqpoint{2.644786in}{1.399537in}}%
\pgfpathlineto{\pgfqpoint{2.655855in}{1.403359in}}%
\pgfpathlineto{\pgfqpoint{2.661243in}{1.404003in}}%
\pgfpathlineto{\pgfqpoint{2.666538in}{1.402398in}}%
\pgfpathlineto{\pgfqpoint{2.676861in}{1.395392in}}%
\pgfpathlineto{\pgfqpoint{2.681895in}{1.394714in}}%
\pgfpathlineto{\pgfqpoint{2.686848in}{1.395815in}}%
\pgfpathlineto{\pgfqpoint{2.691723in}{1.395108in}}%
\pgfpathlineto{\pgfqpoint{2.701245in}{1.387768in}}%
\pgfpathlineto{\pgfqpoint{2.705898in}{1.387496in}}%
\pgfpathlineto{\pgfqpoint{2.710481in}{1.390753in}}%
\pgfpathlineto{\pgfqpoint{2.714997in}{1.395097in}}%
\pgfpathlineto{\pgfqpoint{2.719447in}{1.397891in}}%
\pgfpathlineto{\pgfqpoint{2.723834in}{1.397989in}}%
\pgfpathlineto{\pgfqpoint{2.728158in}{1.396784in}}%
\pgfpathlineto{\pgfqpoint{2.736629in}{1.395783in}}%
\pgfpathlineto{\pgfqpoint{2.740778in}{1.393872in}}%
\pgfpathlineto{\pgfqpoint{2.752900in}{1.385664in}}%
\pgfpathlineto{\pgfqpoint{2.756836in}{1.384161in}}%
\pgfpathlineto{\pgfqpoint{2.760723in}{1.386670in}}%
\pgfpathlineto{\pgfqpoint{2.768351in}{1.398476in}}%
\pgfpathlineto{\pgfqpoint{2.772095in}{1.401816in}}%
\pgfpathlineto{\pgfqpoint{2.775794in}{1.400518in}}%
\pgfpathlineto{\pgfqpoint{2.783061in}{1.392201in}}%
\pgfpathlineto{\pgfqpoint{2.786631in}{1.390520in}}%
\pgfpathlineto{\pgfqpoint{2.800510in}{1.388130in}}%
\pgfpathlineto{\pgfqpoint{2.803883in}{1.390051in}}%
\pgfpathlineto{\pgfqpoint{2.810521in}{1.398014in}}%
\pgfpathlineto{\pgfqpoint{2.813787in}{1.397663in}}%
\pgfpathlineto{\pgfqpoint{2.820217in}{1.387801in}}%
\pgfpathlineto{\pgfqpoint{2.823382in}{1.386579in}}%
\pgfpathlineto{\pgfqpoint{2.826514in}{1.388596in}}%
\pgfpathlineto{\pgfqpoint{2.835725in}{1.397254in}}%
\pgfpathlineto{\pgfqpoint{2.838734in}{1.398609in}}%
\pgfpathlineto{\pgfqpoint{2.841715in}{1.398112in}}%
\pgfpathlineto{\pgfqpoint{2.844667in}{1.396892in}}%
\pgfpathlineto{\pgfqpoint{2.850486in}{1.396224in}}%
\pgfpathlineto{\pgfqpoint{2.853355in}{1.394668in}}%
\pgfpathlineto{\pgfqpoint{2.856198in}{1.392035in}}%
\pgfpathlineto{\pgfqpoint{2.859014in}{1.391948in}}%
\pgfpathlineto{\pgfqpoint{2.861804in}{1.395317in}}%
\pgfpathlineto{\pgfqpoint{2.864570in}{1.397335in}}%
\pgfpathlineto{\pgfqpoint{2.867311in}{1.395037in}}%
\pgfpathlineto{\pgfqpoint{2.870027in}{1.391790in}}%
\pgfpathlineto{\pgfqpoint{2.872720in}{1.391316in}}%
\pgfpathlineto{\pgfqpoint{2.875389in}{1.392716in}}%
\pgfpathlineto{\pgfqpoint{2.878036in}{1.393009in}}%
\pgfpathlineto{\pgfqpoint{2.888398in}{1.387491in}}%
\pgfpathlineto{\pgfqpoint{2.890935in}{1.386452in}}%
\pgfpathlineto{\pgfqpoint{2.893451in}{1.387631in}}%
\pgfpathlineto{\pgfqpoint{2.898422in}{1.392412in}}%
\pgfpathlineto{\pgfqpoint{2.900878in}{1.393293in}}%
\pgfpathlineto{\pgfqpoint{2.903314in}{1.392449in}}%
\pgfpathlineto{\pgfqpoint{2.910509in}{1.385040in}}%
\pgfpathlineto{\pgfqpoint{2.912870in}{1.385803in}}%
\pgfpathlineto{\pgfqpoint{2.924410in}{1.397344in}}%
\pgfpathlineto{\pgfqpoint{2.926667in}{1.398678in}}%
\pgfpathlineto{\pgfqpoint{2.931132in}{1.396751in}}%
\pgfpathlineto{\pgfqpoint{2.933340in}{1.398107in}}%
\pgfpathlineto{\pgfqpoint{2.935533in}{1.401980in}}%
\pgfpathlineto{\pgfqpoint{2.937710in}{1.404660in}}%
\pgfpathlineto{\pgfqpoint{2.939871in}{1.402258in}}%
\pgfpathlineto{\pgfqpoint{2.944149in}{1.390377in}}%
\pgfpathlineto{\pgfqpoint{2.946266in}{1.390165in}}%
\pgfpathlineto{\pgfqpoint{2.948368in}{1.392023in}}%
\pgfpathlineto{\pgfqpoint{2.950456in}{1.392282in}}%
\pgfpathlineto{\pgfqpoint{2.956636in}{1.391123in}}%
\pgfpathlineto{\pgfqpoint{2.960688in}{1.387380in}}%
\pgfpathlineto{\pgfqpoint{2.962695in}{1.389482in}}%
\pgfpathlineto{\pgfqpoint{2.964688in}{1.396036in}}%
\pgfpathlineto{\pgfqpoint{2.968635in}{1.410642in}}%
\pgfpathlineto{\pgfqpoint{2.970590in}{1.410892in}}%
\pgfpathlineto{\pgfqpoint{2.972533in}{1.405388in}}%
\pgfpathlineto{\pgfqpoint{2.976382in}{1.390830in}}%
\pgfpathlineto{\pgfqpoint{2.978288in}{1.388049in}}%
\pgfpathlineto{\pgfqpoint{2.980183in}{1.388625in}}%
\pgfpathlineto{\pgfqpoint{2.982066in}{1.391183in}}%
\pgfpathlineto{\pgfqpoint{2.987646in}{1.401402in}}%
\pgfpathlineto{\pgfqpoint{2.991311in}{1.406051in}}%
\pgfpathlineto{\pgfqpoint{2.993127in}{1.410390in}}%
\pgfpathlineto{\pgfqpoint{2.994932in}{1.412782in}}%
\pgfpathlineto{\pgfqpoint{2.996727in}{1.410282in}}%
\pgfpathlineto{\pgfqpoint{3.000286in}{1.403170in}}%
\pgfpathlineto{\pgfqpoint{3.002050in}{1.399558in}}%
\pgfpathlineto{\pgfqpoint{3.005547in}{1.390959in}}%
\pgfpathlineto{\pgfqpoint{3.007281in}{1.391553in}}%
\pgfpathlineto{\pgfqpoint{3.009006in}{1.394215in}}%
\pgfpathlineto{\pgfqpoint{3.010720in}{1.395045in}}%
\pgfpathlineto{\pgfqpoint{3.012425in}{1.394214in}}%
\pgfpathlineto{\pgfqpoint{3.014121in}{1.395655in}}%
\pgfpathlineto{\pgfqpoint{3.015807in}{1.399582in}}%
\pgfpathlineto{\pgfqpoint{3.017485in}{1.401597in}}%
\pgfpathlineto{\pgfqpoint{3.019153in}{1.400115in}}%
\pgfpathlineto{\pgfqpoint{3.024103in}{1.392634in}}%
\pgfpathlineto{\pgfqpoint{3.025735in}{1.391424in}}%
\pgfpathlineto{\pgfqpoint{3.027359in}{1.392757in}}%
\pgfpathlineto{\pgfqpoint{3.030581in}{1.400759in}}%
\pgfpathlineto{\pgfqpoint{3.032180in}{1.401995in}}%
\pgfpathlineto{\pgfqpoint{3.038492in}{1.395606in}}%
\pgfpathlineto{\pgfqpoint{3.041600in}{1.386879in}}%
\pgfpathlineto{\pgfqpoint{3.043110in}{1.385392in}}%
\pgfpathlineto{\pgfqpoint{3.043110in}{1.385392in}}%
\pgfusepath{stroke}%
\end{pgfscope}%
\begin{pgfscope}%
\pgfpathrectangle{\pgfqpoint{0.650000in}{0.420000in}}{\pgfqpoint{2.388110in}{1.994488in}}%
\pgfusepath{clip}%
\pgfsetrectcap%
\pgfsetroundjoin%
\pgfsetlinewidth{1.003750pt}%
\definecolor{currentstroke}{rgb}{0.498039,0.737255,0.254902}%
\pgfsetstrokecolor{currentstroke}%
\pgfsetdash{}{0pt}%
\pgfpathmoveto{\pgfqpoint{0.645000in}{1.394933in}}%
\pgfpathlineto{\pgfqpoint{0.763003in}{1.396018in}}%
\pgfpathlineto{\pgfqpoint{0.834096in}{1.397640in}}%
\pgfpathlineto{\pgfqpoint{0.877521in}{1.399371in}}%
\pgfpathlineto{\pgfqpoint{0.918461in}{1.401849in}}%
\pgfpathlineto{\pgfqpoint{0.957343in}{1.405334in}}%
\pgfpathlineto{\pgfqpoint{0.994499in}{1.410154in}}%
\pgfpathlineto{\pgfqpoint{1.030190in}{1.416725in}}%
\pgfpathlineto{\pgfqpoint{1.047553in}{1.420829in}}%
\pgfpathlineto{\pgfqpoint{1.064626in}{1.425579in}}%
\pgfpathlineto{\pgfqpoint{1.081428in}{1.431063in}}%
\pgfpathlineto{\pgfqpoint{1.097978in}{1.437374in}}%
\pgfpathlineto{\pgfqpoint{1.114292in}{1.444616in}}%
\pgfpathlineto{\pgfqpoint{1.130386in}{1.452894in}}%
\pgfpathlineto{\pgfqpoint{1.146274in}{1.462309in}}%
\pgfpathlineto{\pgfqpoint{1.161969in}{1.472959in}}%
\pgfpathlineto{\pgfqpoint{1.177482in}{1.484936in}}%
\pgfpathlineto{\pgfqpoint{1.192824in}{1.498345in}}%
\pgfpathlineto{\pgfqpoint{1.208006in}{1.513338in}}%
\pgfpathlineto{\pgfqpoint{1.223037in}{1.530090in}}%
\pgfpathlineto{\pgfqpoint{1.237925in}{1.548705in}}%
\pgfpathlineto{\pgfqpoint{1.252679in}{1.569199in}}%
\pgfpathlineto{\pgfqpoint{1.267305in}{1.591615in}}%
\pgfpathlineto{\pgfqpoint{1.281811in}{1.616097in}}%
\pgfpathlineto{\pgfqpoint{1.296204in}{1.642782in}}%
\pgfpathlineto{\pgfqpoint{1.310489in}{1.671690in}}%
\pgfpathlineto{\pgfqpoint{1.324672in}{1.702716in}}%
\pgfpathlineto{\pgfqpoint{1.338758in}{1.735681in}}%
\pgfpathlineto{\pgfqpoint{1.352752in}{1.770374in}}%
\pgfpathlineto{\pgfqpoint{1.366659in}{1.806626in}}%
\pgfpathlineto{\pgfqpoint{1.394229in}{1.883201in}}%
\pgfpathlineto{\pgfqpoint{1.461906in}{2.079371in}}%
\pgfpathlineto{\pgfqpoint{1.475253in}{2.115675in}}%
\pgfpathlineto{\pgfqpoint{1.488545in}{2.150086in}}%
\pgfpathlineto{\pgfqpoint{1.501784in}{2.182245in}}%
\pgfpathlineto{\pgfqpoint{1.514971in}{2.211550in}}%
\pgfpathlineto{\pgfqpoint{1.528111in}{2.237306in}}%
\pgfpathlineto{\pgfqpoint{1.541204in}{2.259077in}}%
\pgfpathlineto{\pgfqpoint{1.554253in}{2.276647in}}%
\pgfpathlineto{\pgfqpoint{1.567260in}{2.289728in}}%
\pgfpathlineto{\pgfqpoint{1.580226in}{2.298010in}}%
\pgfpathlineto{\pgfqpoint{1.593154in}{2.301549in}}%
\pgfpathlineto{\pgfqpoint{1.606045in}{2.300771in}}%
\pgfpathlineto{\pgfqpoint{1.618901in}{2.296064in}}%
\pgfpathlineto{\pgfqpoint{1.631723in}{2.287564in}}%
\pgfpathlineto{\pgfqpoint{1.644513in}{2.275299in}}%
\pgfpathlineto{\pgfqpoint{1.657272in}{2.259385in}}%
\pgfpathlineto{\pgfqpoint{1.670002in}{2.240144in}}%
\pgfpathlineto{\pgfqpoint{1.682703in}{2.217943in}}%
\pgfpathlineto{\pgfqpoint{1.695377in}{2.193170in}}%
\pgfpathlineto{\pgfqpoint{1.708026in}{2.166361in}}%
\pgfpathlineto{\pgfqpoint{1.733250in}{2.108917in}}%
\pgfpathlineto{\pgfqpoint{1.758382in}{2.048493in}}%
\pgfpathlineto{\pgfqpoint{1.795926in}{1.956948in}}%
\pgfpathlineto{\pgfqpoint{1.820863in}{1.899240in}}%
\pgfpathlineto{\pgfqpoint{1.845731in}{1.845242in}}%
\pgfpathlineto{\pgfqpoint{1.858142in}{1.819818in}}%
\pgfpathlineto{\pgfqpoint{1.870537in}{1.795714in}}%
\pgfpathlineto{\pgfqpoint{1.882919in}{1.773027in}}%
\pgfpathlineto{\pgfqpoint{1.895286in}{1.751689in}}%
\pgfpathlineto{\pgfqpoint{1.907641in}{1.731580in}}%
\pgfpathlineto{\pgfqpoint{1.919983in}{1.712617in}}%
\pgfpathlineto{\pgfqpoint{1.932312in}{1.694814in}}%
\pgfpathlineto{\pgfqpoint{1.944630in}{1.678221in}}%
\pgfpathlineto{\pgfqpoint{1.956937in}{1.662839in}}%
\pgfpathlineto{\pgfqpoint{1.969233in}{1.648584in}}%
\pgfpathlineto{\pgfqpoint{1.981519in}{1.635360in}}%
\pgfpathlineto{\pgfqpoint{1.993795in}{1.623097in}}%
\pgfpathlineto{\pgfqpoint{2.018319in}{1.601088in}}%
\pgfpathlineto{\pgfqpoint{2.042808in}{1.581769in}}%
\pgfpathlineto{\pgfqpoint{2.067264in}{1.565004in}}%
\pgfpathlineto{\pgfqpoint{2.091691in}{1.550578in}}%
\pgfpathlineto{\pgfqpoint{2.116091in}{1.538435in}}%
\pgfpathlineto{\pgfqpoint{2.140465in}{1.528353in}}%
\pgfpathlineto{\pgfqpoint{2.164815in}{1.519562in}}%
\pgfpathlineto{\pgfqpoint{2.201302in}{1.507822in}}%
\pgfpathlineto{\pgfqpoint{2.225602in}{1.500928in}}%
\pgfpathlineto{\pgfqpoint{2.262019in}{1.492152in}}%
\pgfpathlineto{\pgfqpoint{2.298401in}{1.483131in}}%
\pgfpathlineto{\pgfqpoint{2.346862in}{1.470658in}}%
\pgfpathlineto{\pgfqpoint{2.430067in}{1.452181in}}%
\pgfpathlineto{\pgfqpoint{2.651282in}{1.410361in}}%
\pgfpathlineto{\pgfqpoint{2.672288in}{1.407675in}}%
\pgfpathlineto{\pgfqpoint{2.696673in}{1.404364in}}%
\pgfpathlineto{\pgfqpoint{2.732056in}{1.399282in}}%
\pgfpathlineto{\pgfqpoint{2.759988in}{1.396304in}}%
\pgfpathlineto{\pgfqpoint{2.785587in}{1.394899in}}%
\pgfpathlineto{\pgfqpoint{2.815644in}{1.394447in}}%
\pgfpathlineto{\pgfqpoint{3.032353in}{1.394428in}}%
\pgfpathlineto{\pgfqpoint{3.043110in}{1.394428in}}%
\pgfpathlineto{\pgfqpoint{3.043110in}{1.394428in}}%
\pgfusepath{stroke}%
\end{pgfscope}%
\begin{pgfscope}%
\pgfpathrectangle{\pgfqpoint{0.650000in}{0.420000in}}{\pgfqpoint{2.388110in}{1.994488in}}%
\pgfusepath{clip}%
\pgfsetbuttcap%
\pgfsetroundjoin%
\pgfsetlinewidth{1.003750pt}%
\definecolor{currentstroke}{rgb}{0.498039,0.737255,0.254902}%
\pgfsetstrokecolor{currentstroke}%
\pgfsetdash{{5.000000pt}{1.000000pt}}{0.000000pt}%
\pgfpathmoveto{\pgfqpoint{0.645000in}{1.394456in}}%
\pgfpathlineto{\pgfqpoint{1.130386in}{1.395052in}}%
\pgfpathlineto{\pgfqpoint{1.281811in}{1.396291in}}%
\pgfpathlineto{\pgfqpoint{1.593154in}{1.399871in}}%
\pgfpathlineto{\pgfqpoint{1.695377in}{1.399340in}}%
\pgfpathlineto{\pgfqpoint{1.745827in}{1.398056in}}%
\pgfpathlineto{\pgfqpoint{1.808403in}{1.396500in}}%
\pgfpathlineto{\pgfqpoint{1.858142in}{1.396490in}}%
\pgfpathlineto{\pgfqpoint{1.907641in}{1.398145in}}%
\pgfpathlineto{\pgfqpoint{2.042808in}{1.402246in}}%
\pgfpathlineto{\pgfqpoint{2.103894in}{1.401606in}}%
\pgfpathlineto{\pgfqpoint{2.189145in}{1.400416in}}%
\pgfpathlineto{\pgfqpoint{2.249884in}{1.398087in}}%
\pgfpathlineto{\pgfqpoint{2.274150in}{1.398308in}}%
\pgfpathlineto{\pgfqpoint{2.322638in}{1.399371in}}%
\pgfpathlineto{\pgfqpoint{2.395181in}{1.399984in}}%
\pgfpathlineto{\pgfqpoint{2.461878in}{1.396455in}}%
\pgfpathlineto{\pgfqpoint{2.541335in}{1.395367in}}%
\pgfpathlineto{\pgfqpoint{2.604390in}{1.395282in}}%
\pgfpathlineto{\pgfqpoint{2.634526in}{1.395092in}}%
\pgfpathlineto{\pgfqpoint{2.661965in}{1.395049in}}%
\pgfpathlineto{\pgfqpoint{2.744340in}{1.391492in}}%
\pgfpathlineto{\pgfqpoint{2.789076in}{1.393089in}}%
\pgfpathlineto{\pgfqpoint{2.859997in}{1.394293in}}%
\pgfpathlineto{\pgfqpoint{2.937445in}{1.394474in}}%
\pgfpathlineto{\pgfqpoint{2.983073in}{1.394417in}}%
\pgfpathlineto{\pgfqpoint{3.043110in}{1.394453in}}%
\pgfpathlineto{\pgfqpoint{3.043110in}{1.394453in}}%
\pgfusepath{stroke}%
\end{pgfscope}%
\begin{pgfscope}%
\pgfpathrectangle{\pgfqpoint{0.650000in}{0.420000in}}{\pgfqpoint{2.388110in}{1.994488in}}%
\pgfusepath{clip}%
\pgfsetbuttcap%
\pgfsetroundjoin%
\pgfsetlinewidth{1.003750pt}%
\definecolor{currentstroke}{rgb}{0.498039,0.737255,0.254902}%
\pgfsetstrokecolor{currentstroke}%
\pgfsetdash{{3.000000pt}{1.000000pt}{1.000000pt}{1.000000pt}}{0.000000pt}%
\pgfpathmoveto{\pgfqpoint{0.645000in}{0.549674in}}%
\pgfpathlineto{\pgfqpoint{0.652191in}{0.550745in}}%
\pgfpathlineto{\pgfqpoint{0.682092in}{0.556223in}}%
\pgfpathlineto{\pgfqpoint{0.710402in}{0.562435in}}%
\pgfpathlineto{\pgfqpoint{0.737313in}{0.569316in}}%
\pgfpathlineto{\pgfqpoint{0.763003in}{0.576779in}}%
\pgfpathlineto{\pgfqpoint{0.787618in}{0.584711in}}%
\pgfpathlineto{\pgfqpoint{0.834096in}{0.601417in}}%
\pgfpathlineto{\pgfqpoint{0.918461in}{0.633678in}}%
\pgfpathlineto{\pgfqpoint{0.957343in}{0.647054in}}%
\pgfpathlineto{\pgfqpoint{0.994499in}{0.658270in}}%
\pgfpathlineto{\pgfqpoint{1.030190in}{0.668519in}}%
\pgfpathlineto{\pgfqpoint{1.047553in}{0.673971in}}%
\pgfpathlineto{\pgfqpoint{1.064626in}{0.680075in}}%
\pgfpathlineto{\pgfqpoint{1.081428in}{0.687108in}}%
\pgfpathlineto{\pgfqpoint{1.097978in}{0.695205in}}%
\pgfpathlineto{\pgfqpoint{1.114292in}{0.704262in}}%
\pgfpathlineto{\pgfqpoint{1.146274in}{0.723080in}}%
\pgfpathlineto{\pgfqpoint{1.161969in}{0.730840in}}%
\pgfpathlineto{\pgfqpoint{1.177482in}{0.736124in}}%
\pgfpathlineto{\pgfqpoint{1.208006in}{0.742332in}}%
\pgfpathlineto{\pgfqpoint{1.223037in}{0.748219in}}%
\pgfpathlineto{\pgfqpoint{1.252679in}{0.764201in}}%
\pgfpathlineto{\pgfqpoint{1.267305in}{0.771132in}}%
\pgfpathlineto{\pgfqpoint{1.281811in}{0.779370in}}%
\pgfpathlineto{\pgfqpoint{1.296204in}{0.790297in}}%
\pgfpathlineto{\pgfqpoint{1.324672in}{0.815931in}}%
\pgfpathlineto{\pgfqpoint{1.338758in}{0.827166in}}%
\pgfpathlineto{\pgfqpoint{1.352752in}{0.836228in}}%
\pgfpathlineto{\pgfqpoint{1.366659in}{0.843604in}}%
\pgfpathlineto{\pgfqpoint{1.394229in}{0.856907in}}%
\pgfpathlineto{\pgfqpoint{1.407900in}{0.862828in}}%
\pgfpathlineto{\pgfqpoint{1.421500in}{0.866678in}}%
\pgfpathlineto{\pgfqpoint{1.435032in}{0.867865in}}%
\pgfpathlineto{\pgfqpoint{1.461906in}{0.866383in}}%
\pgfpathlineto{\pgfqpoint{1.475253in}{0.865484in}}%
\pgfpathlineto{\pgfqpoint{1.488545in}{0.865306in}}%
\pgfpathlineto{\pgfqpoint{1.528111in}{0.867534in}}%
\pgfpathlineto{\pgfqpoint{1.567260in}{0.866102in}}%
\pgfpathlineto{\pgfqpoint{1.593154in}{0.864737in}}%
\pgfpathlineto{\pgfqpoint{1.606045in}{0.865361in}}%
\pgfpathlineto{\pgfqpoint{1.618901in}{0.867509in}}%
\pgfpathlineto{\pgfqpoint{1.631723in}{0.870945in}}%
\pgfpathlineto{\pgfqpoint{1.644513in}{0.875243in}}%
\pgfpathlineto{\pgfqpoint{1.670002in}{0.885571in}}%
\pgfpathlineto{\pgfqpoint{1.695377in}{0.898109in}}%
\pgfpathlineto{\pgfqpoint{1.708026in}{0.905347in}}%
\pgfpathlineto{\pgfqpoint{1.720650in}{0.913339in}}%
\pgfpathlineto{\pgfqpoint{1.745827in}{0.931191in}}%
\pgfpathlineto{\pgfqpoint{1.770917in}{0.950542in}}%
\pgfpathlineto{\pgfqpoint{1.795926in}{0.971183in}}%
\pgfpathlineto{\pgfqpoint{1.833305in}{1.004223in}}%
\pgfpathlineto{\pgfqpoint{1.870537in}{1.038704in}}%
\pgfpathlineto{\pgfqpoint{1.932312in}{1.096222in}}%
\pgfpathlineto{\pgfqpoint{1.969233in}{1.129054in}}%
\pgfpathlineto{\pgfqpoint{1.993795in}{1.149554in}}%
\pgfpathlineto{\pgfqpoint{2.018319in}{1.168901in}}%
\pgfpathlineto{\pgfqpoint{2.042808in}{1.187026in}}%
\pgfpathlineto{\pgfqpoint{2.067264in}{1.203731in}}%
\pgfpathlineto{\pgfqpoint{2.091691in}{1.218977in}}%
\pgfpathlineto{\pgfqpoint{2.116091in}{1.232715in}}%
\pgfpathlineto{\pgfqpoint{2.140465in}{1.244991in}}%
\pgfpathlineto{\pgfqpoint{2.164815in}{1.256143in}}%
\pgfpathlineto{\pgfqpoint{2.189145in}{1.266439in}}%
\pgfpathlineto{\pgfqpoint{2.213454in}{1.275903in}}%
\pgfpathlineto{\pgfqpoint{2.249884in}{1.288705in}}%
\pgfpathlineto{\pgfqpoint{2.298401in}{1.304338in}}%
\pgfpathlineto{\pgfqpoint{2.334751in}{1.314895in}}%
\pgfpathlineto{\pgfqpoint{2.383154in}{1.327262in}}%
\pgfpathlineto{\pgfqpoint{2.418733in}{1.335376in}}%
\pgfpathlineto{\pgfqpoint{2.481398in}{1.347635in}}%
\pgfpathlineto{\pgfqpoint{2.556509in}{1.361399in}}%
\pgfpathlineto{\pgfqpoint{2.604390in}{1.369501in}}%
\pgfpathlineto{\pgfqpoint{2.677323in}{1.381055in}}%
\pgfpathlineto{\pgfqpoint{2.752263in}{1.389834in}}%
\pgfpathlineto{\pgfqpoint{2.782059in}{1.392242in}}%
\pgfpathlineto{\pgfqpoint{2.815644in}{1.393779in}}%
\pgfpathlineto{\pgfqpoint{2.857232in}{1.394330in}}%
\pgfpathlineto{\pgfqpoint{2.983073in}{1.394429in}}%
\pgfpathlineto{\pgfqpoint{3.043110in}{1.394429in}}%
\pgfpathlineto{\pgfqpoint{3.043110in}{1.394429in}}%
\pgfusepath{stroke}%
\end{pgfscope}%
\begin{pgfscope}%
\pgfpathrectangle{\pgfqpoint{0.650000in}{0.420000in}}{\pgfqpoint{2.388110in}{1.994488in}}%
\pgfusepath{clip}%
\pgfsetbuttcap%
\pgfsetroundjoin%
\pgfsetlinewidth{1.003750pt}%
\definecolor{currentstroke}{rgb}{0.498039,0.737255,0.254902}%
\pgfsetstrokecolor{currentstroke}%
\pgfsetdash{{1.000000pt}{1.000000pt}}{0.000000pt}%
\pgfpathmoveto{\pgfqpoint{0.645000in}{1.429395in}}%
\pgfpathlineto{\pgfqpoint{0.710402in}{1.435695in}}%
\pgfpathlineto{\pgfqpoint{0.811281in}{1.446580in}}%
\pgfpathlineto{\pgfqpoint{0.957343in}{1.462656in}}%
\pgfpathlineto{\pgfqpoint{1.012513in}{1.467746in}}%
\pgfpathlineto{\pgfqpoint{1.064626in}{1.471543in}}%
\pgfpathlineto{\pgfqpoint{1.114292in}{1.473949in}}%
\pgfpathlineto{\pgfqpoint{1.146274in}{1.474658in}}%
\pgfpathlineto{\pgfqpoint{1.192824in}{1.474409in}}%
\pgfpathlineto{\pgfqpoint{1.223037in}{1.473399in}}%
\pgfpathlineto{\pgfqpoint{1.267305in}{1.470565in}}%
\pgfpathlineto{\pgfqpoint{1.310489in}{1.466797in}}%
\pgfpathlineto{\pgfqpoint{1.352752in}{1.461827in}}%
\pgfpathlineto{\pgfqpoint{1.394229in}{1.455817in}}%
\pgfpathlineto{\pgfqpoint{1.435032in}{1.448819in}}%
\pgfpathlineto{\pgfqpoint{1.475253in}{1.441028in}}%
\pgfpathlineto{\pgfqpoint{1.514971in}{1.432305in}}%
\pgfpathlineto{\pgfqpoint{1.580226in}{1.416284in}}%
\pgfpathlineto{\pgfqpoint{1.670002in}{1.393776in}}%
\pgfpathlineto{\pgfqpoint{1.695377in}{1.388602in}}%
\pgfpathlineto{\pgfqpoint{1.720650in}{1.384961in}}%
\pgfpathlineto{\pgfqpoint{1.745827in}{1.382718in}}%
\pgfpathlineto{\pgfqpoint{1.770917in}{1.381486in}}%
\pgfpathlineto{\pgfqpoint{1.795926in}{1.381472in}}%
\pgfpathlineto{\pgfqpoint{1.820863in}{1.382812in}}%
\pgfpathlineto{\pgfqpoint{1.858142in}{1.386284in}}%
\pgfpathlineto{\pgfqpoint{1.956937in}{1.395921in}}%
\pgfpathlineto{\pgfqpoint{1.993795in}{1.398351in}}%
\pgfpathlineto{\pgfqpoint{2.018319in}{1.398840in}}%
\pgfpathlineto{\pgfqpoint{2.042808in}{1.398102in}}%
\pgfpathlineto{\pgfqpoint{2.140465in}{1.392879in}}%
\pgfpathlineto{\pgfqpoint{2.237745in}{1.393729in}}%
\pgfpathlineto{\pgfqpoint{2.358969in}{1.393064in}}%
\pgfpathlineto{\pgfqpoint{2.430067in}{1.393926in}}%
\pgfpathlineto{\pgfqpoint{2.461878in}{1.393073in}}%
\pgfpathlineto{\pgfqpoint{2.490711in}{1.392527in}}%
\pgfpathlineto{\pgfqpoint{2.517066in}{1.393526in}}%
\pgfpathlineto{\pgfqpoint{2.549016in}{1.394566in}}%
\pgfpathlineto{\pgfqpoint{2.577950in}{1.394399in}}%
\pgfpathlineto{\pgfqpoint{2.604390in}{1.393424in}}%
\pgfpathlineto{\pgfqpoint{2.651282in}{1.393593in}}%
\pgfpathlineto{\pgfqpoint{2.677323in}{1.393521in}}%
\pgfpathlineto{\pgfqpoint{2.701325in}{1.393199in}}%
\pgfpathlineto{\pgfqpoint{2.732056in}{1.392903in}}%
\pgfpathlineto{\pgfqpoint{2.821942in}{1.394282in}}%
\pgfpathlineto{\pgfqpoint{2.843017in}{1.394442in}}%
\pgfpathlineto{\pgfqpoint{2.876087in}{1.394435in}}%
\pgfpathlineto{\pgfqpoint{3.043110in}{1.394401in}}%
\pgfpathlineto{\pgfqpoint{3.043110in}{1.394401in}}%
\pgfusepath{stroke}%
\end{pgfscope}%
\begin{pgfscope}%
\pgfpathrectangle{\pgfqpoint{0.650000in}{0.420000in}}{\pgfqpoint{2.388110in}{1.994488in}}%
\pgfusepath{clip}%
\pgfsetbuttcap%
\pgfsetroundjoin%
\pgfsetlinewidth{1.003750pt}%
\definecolor{currentstroke}{rgb}{0.498039,0.737255,0.254902}%
\pgfsetstrokecolor{currentstroke}%
\pgfsetdash{{5.000000pt}{5.000000pt}}{0.000000pt}%
\pgfpathmoveto{\pgfqpoint{0.645000in}{1.395312in}}%
\pgfpathlineto{\pgfqpoint{0.710402in}{1.396088in}}%
\pgfpathlineto{\pgfqpoint{0.787618in}{1.397976in}}%
\pgfpathlineto{\pgfqpoint{0.834096in}{1.399987in}}%
\pgfpathlineto{\pgfqpoint{0.877521in}{1.402818in}}%
\pgfpathlineto{\pgfqpoint{0.918461in}{1.406716in}}%
\pgfpathlineto{\pgfqpoint{0.957343in}{1.411994in}}%
\pgfpathlineto{\pgfqpoint{0.994499in}{1.419030in}}%
\pgfpathlineto{\pgfqpoint{1.012513in}{1.423341in}}%
\pgfpathlineto{\pgfqpoint{1.030190in}{1.428249in}}%
\pgfpathlineto{\pgfqpoint{1.047553in}{1.433804in}}%
\pgfpathlineto{\pgfqpoint{1.064626in}{1.440042in}}%
\pgfpathlineto{\pgfqpoint{1.081428in}{1.446989in}}%
\pgfpathlineto{\pgfqpoint{1.097978in}{1.454652in}}%
\pgfpathlineto{\pgfqpoint{1.114292in}{1.463015in}}%
\pgfpathlineto{\pgfqpoint{1.130386in}{1.472038in}}%
\pgfpathlineto{\pgfqpoint{1.146274in}{1.481665in}}%
\pgfpathlineto{\pgfqpoint{1.177482in}{1.502521in}}%
\pgfpathlineto{\pgfqpoint{1.208006in}{1.525179in}}%
\pgfpathlineto{\pgfqpoint{1.237925in}{1.547908in}}%
\pgfpathlineto{\pgfqpoint{1.252679in}{1.558262in}}%
\pgfpathlineto{\pgfqpoint{1.267305in}{1.567793in}}%
\pgfpathlineto{\pgfqpoint{1.281811in}{1.576487in}}%
\pgfpathlineto{\pgfqpoint{1.296204in}{1.584019in}}%
\pgfpathlineto{\pgfqpoint{1.310489in}{1.589811in}}%
\pgfpathlineto{\pgfqpoint{1.324672in}{1.593211in}}%
\pgfpathlineto{\pgfqpoint{1.338758in}{1.593713in}}%
\pgfpathlineto{\pgfqpoint{1.352752in}{1.591031in}}%
\pgfpathlineto{\pgfqpoint{1.366659in}{1.585114in}}%
\pgfpathlineto{\pgfqpoint{1.380483in}{1.576045in}}%
\pgfpathlineto{\pgfqpoint{1.394229in}{1.563999in}}%
\pgfpathlineto{\pgfqpoint{1.407900in}{1.549086in}}%
\pgfpathlineto{\pgfqpoint{1.421500in}{1.531271in}}%
\pgfpathlineto{\pgfqpoint{1.435032in}{1.510289in}}%
\pgfpathlineto{\pgfqpoint{1.448499in}{1.485912in}}%
\pgfpathlineto{\pgfqpoint{1.461906in}{1.458328in}}%
\pgfpathlineto{\pgfqpoint{1.488545in}{1.397250in}}%
\pgfpathlineto{\pgfqpoint{1.514971in}{1.334964in}}%
\pgfpathlineto{\pgfqpoint{1.541204in}{1.276858in}}%
\pgfpathlineto{\pgfqpoint{1.554253in}{1.249743in}}%
\pgfpathlineto{\pgfqpoint{1.567260in}{1.224089in}}%
\pgfpathlineto{\pgfqpoint{1.580226in}{1.201239in}}%
\pgfpathlineto{\pgfqpoint{1.593154in}{1.182786in}}%
\pgfpathlineto{\pgfqpoint{1.606045in}{1.169233in}}%
\pgfpathlineto{\pgfqpoint{1.618901in}{1.159953in}}%
\pgfpathlineto{\pgfqpoint{1.631723in}{1.153724in}}%
\pgfpathlineto{\pgfqpoint{1.644513in}{1.149528in}}%
\pgfpathlineto{\pgfqpoint{1.657272in}{1.148028in}}%
\pgfpathlineto{\pgfqpoint{1.670002in}{1.151734in}}%
\pgfpathlineto{\pgfqpoint{1.682703in}{1.162125in}}%
\pgfpathlineto{\pgfqpoint{1.708026in}{1.190437in}}%
\pgfpathlineto{\pgfqpoint{1.720650in}{1.201647in}}%
\pgfpathlineto{\pgfqpoint{1.745827in}{1.220800in}}%
\pgfpathlineto{\pgfqpoint{1.758382in}{1.231879in}}%
\pgfpathlineto{\pgfqpoint{1.770917in}{1.247393in}}%
\pgfpathlineto{\pgfqpoint{1.783431in}{1.265251in}}%
\pgfpathlineto{\pgfqpoint{1.795926in}{1.280767in}}%
\pgfpathlineto{\pgfqpoint{1.808403in}{1.293513in}}%
\pgfpathlineto{\pgfqpoint{1.820863in}{1.305152in}}%
\pgfpathlineto{\pgfqpoint{1.845731in}{1.325992in}}%
\pgfpathlineto{\pgfqpoint{1.858142in}{1.335544in}}%
\pgfpathlineto{\pgfqpoint{1.882919in}{1.352867in}}%
\pgfpathlineto{\pgfqpoint{1.895286in}{1.362546in}}%
\pgfpathlineto{\pgfqpoint{1.907641in}{1.373314in}}%
\pgfpathlineto{\pgfqpoint{1.919983in}{1.382707in}}%
\pgfpathlineto{\pgfqpoint{1.932312in}{1.388433in}}%
\pgfpathlineto{\pgfqpoint{1.956937in}{1.394388in}}%
\pgfpathlineto{\pgfqpoint{1.969233in}{1.398064in}}%
\pgfpathlineto{\pgfqpoint{1.981519in}{1.400904in}}%
\pgfpathlineto{\pgfqpoint{1.993795in}{1.401928in}}%
\pgfpathlineto{\pgfqpoint{2.006062in}{1.401734in}}%
\pgfpathlineto{\pgfqpoint{2.018319in}{1.400675in}}%
\pgfpathlineto{\pgfqpoint{2.030568in}{1.398994in}}%
\pgfpathlineto{\pgfqpoint{2.042808in}{1.398672in}}%
\pgfpathlineto{\pgfqpoint{2.067264in}{1.403211in}}%
\pgfpathlineto{\pgfqpoint{2.091691in}{1.405450in}}%
\pgfpathlineto{\pgfqpoint{2.103894in}{1.408438in}}%
\pgfpathlineto{\pgfqpoint{2.128281in}{1.416531in}}%
\pgfpathlineto{\pgfqpoint{2.140465in}{1.417956in}}%
\pgfpathlineto{\pgfqpoint{2.152643in}{1.415856in}}%
\pgfpathlineto{\pgfqpoint{2.176983in}{1.406355in}}%
\pgfpathlineto{\pgfqpoint{2.189145in}{1.403892in}}%
\pgfpathlineto{\pgfqpoint{2.201302in}{1.404170in}}%
\pgfpathlineto{\pgfqpoint{2.225602in}{1.406440in}}%
\pgfpathlineto{\pgfqpoint{2.237745in}{1.405887in}}%
\pgfpathlineto{\pgfqpoint{2.249884in}{1.403265in}}%
\pgfpathlineto{\pgfqpoint{2.262019in}{1.399326in}}%
\pgfpathlineto{\pgfqpoint{2.274150in}{1.397017in}}%
\pgfpathlineto{\pgfqpoint{2.286277in}{1.398076in}}%
\pgfpathlineto{\pgfqpoint{2.298401in}{1.400421in}}%
\pgfpathlineto{\pgfqpoint{2.310521in}{1.401042in}}%
\pgfpathlineto{\pgfqpoint{2.322638in}{1.398452in}}%
\pgfpathlineto{\pgfqpoint{2.334751in}{1.393213in}}%
\pgfpathlineto{\pgfqpoint{2.346862in}{1.388618in}}%
\pgfpathlineto{\pgfqpoint{2.358969in}{1.387571in}}%
\pgfpathlineto{\pgfqpoint{2.383154in}{1.391044in}}%
\pgfpathlineto{\pgfqpoint{2.395181in}{1.393676in}}%
\pgfpathlineto{\pgfqpoint{2.407075in}{1.398602in}}%
\pgfpathlineto{\pgfqpoint{2.418733in}{1.404100in}}%
\pgfpathlineto{\pgfqpoint{2.430067in}{1.405978in}}%
\pgfpathlineto{\pgfqpoint{2.441034in}{1.403392in}}%
\pgfpathlineto{\pgfqpoint{2.451632in}{1.400120in}}%
\pgfpathlineto{\pgfqpoint{2.461878in}{1.398965in}}%
\pgfpathlineto{\pgfqpoint{2.481398in}{1.398983in}}%
\pgfpathlineto{\pgfqpoint{2.490711in}{1.399577in}}%
\pgfpathlineto{\pgfqpoint{2.499750in}{1.400734in}}%
\pgfpathlineto{\pgfqpoint{2.517066in}{1.401438in}}%
\pgfpathlineto{\pgfqpoint{2.525371in}{1.403054in}}%
\pgfpathlineto{\pgfqpoint{2.533456in}{1.405941in}}%
\pgfpathlineto{\pgfqpoint{2.541335in}{1.407079in}}%
\pgfpathlineto{\pgfqpoint{2.549016in}{1.404868in}}%
\pgfpathlineto{\pgfqpoint{2.556509in}{1.400842in}}%
\pgfpathlineto{\pgfqpoint{2.570968in}{1.391835in}}%
\pgfpathlineto{\pgfqpoint{2.577950in}{1.389907in}}%
\pgfpathlineto{\pgfqpoint{2.584777in}{1.391134in}}%
\pgfpathlineto{\pgfqpoint{2.597990in}{1.396933in}}%
\pgfpathlineto{\pgfqpoint{2.604390in}{1.400638in}}%
\pgfpathlineto{\pgfqpoint{2.610658in}{1.403325in}}%
\pgfpathlineto{\pgfqpoint{2.616802in}{1.402930in}}%
\pgfpathlineto{\pgfqpoint{2.628731in}{1.397499in}}%
\pgfpathlineto{\pgfqpoint{2.634526in}{1.395827in}}%
\pgfpathlineto{\pgfqpoint{2.645798in}{1.394265in}}%
\pgfpathlineto{\pgfqpoint{2.656670in}{1.390614in}}%
\pgfpathlineto{\pgfqpoint{2.661965in}{1.391286in}}%
\pgfpathlineto{\pgfqpoint{2.672288in}{1.397283in}}%
\pgfpathlineto{\pgfqpoint{2.677323in}{1.398420in}}%
\pgfpathlineto{\pgfqpoint{2.687150in}{1.396399in}}%
\pgfpathlineto{\pgfqpoint{2.691948in}{1.396641in}}%
\pgfpathlineto{\pgfqpoint{2.701325in}{1.400990in}}%
\pgfpathlineto{\pgfqpoint{2.705908in}{1.401199in}}%
\pgfpathlineto{\pgfqpoint{2.714874in}{1.397680in}}%
\pgfpathlineto{\pgfqpoint{2.732056in}{1.397680in}}%
\pgfpathlineto{\pgfqpoint{2.740300in}{1.402812in}}%
\pgfpathlineto{\pgfqpoint{2.744340in}{1.402978in}}%
\pgfpathlineto{\pgfqpoint{2.752263in}{1.397565in}}%
\pgfpathlineto{\pgfqpoint{2.756150in}{1.396589in}}%
\pgfpathlineto{\pgfqpoint{2.759988in}{1.398347in}}%
\pgfpathlineto{\pgfqpoint{2.763778in}{1.401331in}}%
\pgfpathlineto{\pgfqpoint{2.767523in}{1.403043in}}%
\pgfpathlineto{\pgfqpoint{2.771222in}{1.403598in}}%
\pgfpathlineto{\pgfqpoint{2.774877in}{1.404908in}}%
\pgfpathlineto{\pgfqpoint{2.778489in}{1.407127in}}%
\pgfpathlineto{\pgfqpoint{2.782059in}{1.408241in}}%
\pgfpathlineto{\pgfqpoint{2.785587in}{1.406828in}}%
\pgfpathlineto{\pgfqpoint{2.789076in}{1.404257in}}%
\pgfpathlineto{\pgfqpoint{2.802648in}{1.398091in}}%
\pgfpathlineto{\pgfqpoint{2.805949in}{1.398166in}}%
\pgfpathlineto{\pgfqpoint{2.809215in}{1.399850in}}%
\pgfpathlineto{\pgfqpoint{2.812446in}{1.402545in}}%
\pgfpathlineto{\pgfqpoint{2.815644in}{1.406098in}}%
\pgfpathlineto{\pgfqpoint{2.818809in}{1.408782in}}%
\pgfpathlineto{\pgfqpoint{2.821942in}{1.407100in}}%
\pgfpathlineto{\pgfqpoint{2.825043in}{1.401576in}}%
\pgfpathlineto{\pgfqpoint{2.828112in}{1.397719in}}%
\pgfpathlineto{\pgfqpoint{2.834162in}{1.397591in}}%
\pgfpathlineto{\pgfqpoint{2.840094in}{1.394088in}}%
\pgfpathlineto{\pgfqpoint{2.843017in}{1.394387in}}%
\pgfpathlineto{\pgfqpoint{2.845913in}{1.396911in}}%
\pgfpathlineto{\pgfqpoint{2.848782in}{1.402154in}}%
\pgfpathlineto{\pgfqpoint{2.851625in}{1.409567in}}%
\pgfpathlineto{\pgfqpoint{2.854441in}{1.413984in}}%
\pgfpathlineto{\pgfqpoint{2.857232in}{1.411028in}}%
\pgfpathlineto{\pgfqpoint{2.859997in}{1.404184in}}%
\pgfpathlineto{\pgfqpoint{2.862738in}{1.399061in}}%
\pgfpathlineto{\pgfqpoint{2.865455in}{1.396781in}}%
\pgfpathlineto{\pgfqpoint{2.868147in}{1.396294in}}%
\pgfpathlineto{\pgfqpoint{2.873463in}{1.397646in}}%
\pgfpathlineto{\pgfqpoint{2.876087in}{1.397356in}}%
\pgfpathlineto{\pgfqpoint{2.881267in}{1.394547in}}%
\pgfpathlineto{\pgfqpoint{2.883825in}{1.394799in}}%
\pgfpathlineto{\pgfqpoint{2.888878in}{1.398828in}}%
\pgfpathlineto{\pgfqpoint{2.891374in}{1.398944in}}%
\pgfpathlineto{\pgfqpoint{2.896305in}{1.394849in}}%
\pgfpathlineto{\pgfqpoint{2.898741in}{1.394585in}}%
\pgfpathlineto{\pgfqpoint{2.901158in}{1.395428in}}%
\pgfpathlineto{\pgfqpoint{2.903556in}{1.395377in}}%
\pgfpathlineto{\pgfqpoint{2.910640in}{1.390887in}}%
\pgfpathlineto{\pgfqpoint{2.912965in}{1.391433in}}%
\pgfpathlineto{\pgfqpoint{2.917564in}{1.395329in}}%
\pgfpathlineto{\pgfqpoint{2.919838in}{1.397142in}}%
\pgfpathlineto{\pgfqpoint{2.922095in}{1.397219in}}%
\pgfpathlineto{\pgfqpoint{2.924335in}{1.395329in}}%
\pgfpathlineto{\pgfqpoint{2.928768in}{1.390060in}}%
\pgfpathlineto{\pgfqpoint{2.930960in}{1.390015in}}%
\pgfpathlineto{\pgfqpoint{2.933137in}{1.392812in}}%
\pgfpathlineto{\pgfqpoint{2.935299in}{1.396950in}}%
\pgfpathlineto{\pgfqpoint{2.937445in}{1.399533in}}%
\pgfpathlineto{\pgfqpoint{2.939576in}{1.398488in}}%
\pgfpathlineto{\pgfqpoint{2.943796in}{1.390858in}}%
\pgfpathlineto{\pgfqpoint{2.945884in}{1.389623in}}%
\pgfpathlineto{\pgfqpoint{2.947957in}{1.390648in}}%
\pgfpathlineto{\pgfqpoint{2.952064in}{1.396086in}}%
\pgfpathlineto{\pgfqpoint{2.954096in}{1.398928in}}%
\pgfpathlineto{\pgfqpoint{2.956116in}{1.399838in}}%
\pgfpathlineto{\pgfqpoint{2.958122in}{1.398172in}}%
\pgfpathlineto{\pgfqpoint{2.960115in}{1.395521in}}%
\pgfpathlineto{\pgfqpoint{2.962095in}{1.394739in}}%
\pgfpathlineto{\pgfqpoint{2.966018in}{1.399746in}}%
\pgfpathlineto{\pgfqpoint{2.967960in}{1.399123in}}%
\pgfpathlineto{\pgfqpoint{2.971809in}{1.392232in}}%
\pgfpathlineto{\pgfqpoint{2.973715in}{1.392519in}}%
\pgfpathlineto{\pgfqpoint{2.975610in}{1.394410in}}%
\pgfpathlineto{\pgfqpoint{2.977493in}{1.395159in}}%
\pgfpathlineto{\pgfqpoint{2.979364in}{1.393738in}}%
\pgfpathlineto{\pgfqpoint{2.981224in}{1.391137in}}%
\pgfpathlineto{\pgfqpoint{2.983073in}{1.389550in}}%
\pgfpathlineto{\pgfqpoint{2.984911in}{1.390916in}}%
\pgfpathlineto{\pgfqpoint{2.988554in}{1.399063in}}%
\pgfpathlineto{\pgfqpoint{2.990360in}{1.400346in}}%
\pgfpathlineto{\pgfqpoint{2.993939in}{1.397349in}}%
\pgfpathlineto{\pgfqpoint{2.997477in}{1.394837in}}%
\pgfpathlineto{\pgfqpoint{2.999231in}{1.394358in}}%
\pgfpathlineto{\pgfqpoint{3.000975in}{1.395215in}}%
\pgfpathlineto{\pgfqpoint{3.002709in}{1.396888in}}%
\pgfpathlineto{\pgfqpoint{3.004433in}{1.397470in}}%
\pgfpathlineto{\pgfqpoint{3.007853in}{1.394045in}}%
\pgfpathlineto{\pgfqpoint{3.009548in}{1.394216in}}%
\pgfpathlineto{\pgfqpoint{3.011235in}{1.395064in}}%
\pgfpathlineto{\pgfqpoint{3.014580in}{1.394108in}}%
\pgfpathlineto{\pgfqpoint{3.017889in}{1.396424in}}%
\pgfpathlineto{\pgfqpoint{3.021163in}{1.396151in}}%
\pgfpathlineto{\pgfqpoint{3.026009in}{1.398489in}}%
\pgfpathlineto{\pgfqpoint{3.027607in}{1.396063in}}%
\pgfpathlineto{\pgfqpoint{3.030779in}{1.389944in}}%
\pgfpathlineto{\pgfqpoint{3.032353in}{1.390590in}}%
\pgfpathlineto{\pgfqpoint{3.037027in}{1.398674in}}%
\pgfpathlineto{\pgfqpoint{3.038569in}{1.398812in}}%
\pgfpathlineto{\pgfqpoint{3.043110in}{1.394860in}}%
\pgfpathlineto{\pgfqpoint{3.043110in}{1.394860in}}%
\pgfusepath{stroke}%
\end{pgfscope}%
\begin{pgfscope}%
\pgfpathrectangle{\pgfqpoint{0.650000in}{0.420000in}}{\pgfqpoint{2.388110in}{1.994488in}}%
\pgfusepath{clip}%
\pgfsetbuttcap%
\pgfsetroundjoin%
\pgfsetlinewidth{1.003750pt}%
\definecolor{currentstroke}{rgb}{0.498039,0.737255,0.254902}%
\pgfsetstrokecolor{currentstroke}%
\pgfsetdash{{3.000000pt}{5.000000pt}{1.000000pt}{5.000000pt}}{0.000000pt}%
\pgfpathmoveto{\pgfqpoint{0.645000in}{2.188023in}}%
\pgfpathlineto{\pgfqpoint{0.652191in}{2.187439in}}%
\pgfpathlineto{\pgfqpoint{0.682092in}{2.183779in}}%
\pgfpathlineto{\pgfqpoint{0.710402in}{2.178606in}}%
\pgfpathlineto{\pgfqpoint{0.737313in}{2.171598in}}%
\pgfpathlineto{\pgfqpoint{0.763003in}{2.162580in}}%
\pgfpathlineto{\pgfqpoint{0.787618in}{2.151542in}}%
\pgfpathlineto{\pgfqpoint{0.811281in}{2.138643in}}%
\pgfpathlineto{\pgfqpoint{0.834096in}{2.124207in}}%
\pgfpathlineto{\pgfqpoint{0.856151in}{2.108696in}}%
\pgfpathlineto{\pgfqpoint{0.938136in}{2.047609in}}%
\pgfpathlineto{\pgfqpoint{0.957343in}{2.035183in}}%
\pgfpathlineto{\pgfqpoint{0.976119in}{2.024354in}}%
\pgfpathlineto{\pgfqpoint{1.030190in}{1.996624in}}%
\pgfpathlineto{\pgfqpoint{1.047553in}{1.985561in}}%
\pgfpathlineto{\pgfqpoint{1.064626in}{1.971294in}}%
\pgfpathlineto{\pgfqpoint{1.081428in}{1.952529in}}%
\pgfpathlineto{\pgfqpoint{1.097978in}{1.928482in}}%
\pgfpathlineto{\pgfqpoint{1.114292in}{1.899254in}}%
\pgfpathlineto{\pgfqpoint{1.130386in}{1.866069in}}%
\pgfpathlineto{\pgfqpoint{1.177482in}{1.762300in}}%
\pgfpathlineto{\pgfqpoint{1.192824in}{1.725764in}}%
\pgfpathlineto{\pgfqpoint{1.208006in}{1.684271in}}%
\pgfpathlineto{\pgfqpoint{1.223037in}{1.641305in}}%
\pgfpathlineto{\pgfqpoint{1.237925in}{1.604889in}}%
\pgfpathlineto{\pgfqpoint{1.252679in}{1.577969in}}%
\pgfpathlineto{\pgfqpoint{1.267305in}{1.554402in}}%
\pgfpathlineto{\pgfqpoint{1.281811in}{1.525240in}}%
\pgfpathlineto{\pgfqpoint{1.296204in}{1.486111in}}%
\pgfpathlineto{\pgfqpoint{1.338758in}{1.347141in}}%
\pgfpathlineto{\pgfqpoint{1.352752in}{1.310391in}}%
\pgfpathlineto{\pgfqpoint{1.394229in}{1.215176in}}%
\pgfpathlineto{\pgfqpoint{1.407900in}{1.181048in}}%
\pgfpathlineto{\pgfqpoint{1.421500in}{1.149688in}}%
\pgfpathlineto{\pgfqpoint{1.435032in}{1.125040in}}%
\pgfpathlineto{\pgfqpoint{1.448499in}{1.109902in}}%
\pgfpathlineto{\pgfqpoint{1.461906in}{1.105433in}}%
\pgfpathlineto{\pgfqpoint{1.475253in}{1.108576in}}%
\pgfpathlineto{\pgfqpoint{1.488545in}{1.113038in}}%
\pgfpathlineto{\pgfqpoint{1.501784in}{1.115602in}}%
\pgfpathlineto{\pgfqpoint{1.514971in}{1.118898in}}%
\pgfpathlineto{\pgfqpoint{1.528111in}{1.126675in}}%
\pgfpathlineto{\pgfqpoint{1.541204in}{1.138646in}}%
\pgfpathlineto{\pgfqpoint{1.567260in}{1.164568in}}%
\pgfpathlineto{\pgfqpoint{1.580226in}{1.179338in}}%
\pgfpathlineto{\pgfqpoint{1.593154in}{1.197347in}}%
\pgfpathlineto{\pgfqpoint{1.618901in}{1.235861in}}%
\pgfpathlineto{\pgfqpoint{1.631723in}{1.252408in}}%
\pgfpathlineto{\pgfqpoint{1.657272in}{1.282464in}}%
\pgfpathlineto{\pgfqpoint{1.670002in}{1.296247in}}%
\pgfpathlineto{\pgfqpoint{1.695377in}{1.321309in}}%
\pgfpathlineto{\pgfqpoint{1.708026in}{1.332871in}}%
\pgfpathlineto{\pgfqpoint{1.720650in}{1.342793in}}%
\pgfpathlineto{\pgfqpoint{1.733250in}{1.351111in}}%
\pgfpathlineto{\pgfqpoint{1.758382in}{1.365518in}}%
\pgfpathlineto{\pgfqpoint{1.783431in}{1.377790in}}%
\pgfpathlineto{\pgfqpoint{1.795926in}{1.382790in}}%
\pgfpathlineto{\pgfqpoint{1.808403in}{1.386796in}}%
\pgfpathlineto{\pgfqpoint{1.833305in}{1.393060in}}%
\pgfpathlineto{\pgfqpoint{1.858142in}{1.397675in}}%
\pgfpathlineto{\pgfqpoint{1.870537in}{1.399039in}}%
\pgfpathlineto{\pgfqpoint{1.895286in}{1.400327in}}%
\pgfpathlineto{\pgfqpoint{1.919983in}{1.400898in}}%
\pgfpathlineto{\pgfqpoint{1.944630in}{1.400376in}}%
\pgfpathlineto{\pgfqpoint{1.993795in}{1.398967in}}%
\pgfpathlineto{\pgfqpoint{2.140465in}{1.394992in}}%
\pgfpathlineto{\pgfqpoint{2.237745in}{1.394591in}}%
\pgfpathlineto{\pgfqpoint{2.517066in}{1.394447in}}%
\pgfpathlineto{\pgfqpoint{2.549016in}{1.394395in}}%
\pgfpathlineto{\pgfqpoint{2.591455in}{1.394302in}}%
\pgfpathlineto{\pgfqpoint{2.667170in}{1.394408in}}%
\pgfpathlineto{\pgfqpoint{2.719261in}{1.394362in}}%
\pgfpathlineto{\pgfqpoint{2.785587in}{1.394570in}}%
\pgfpathlineto{\pgfqpoint{3.043110in}{1.394429in}}%
\pgfpathlineto{\pgfqpoint{3.043110in}{1.394429in}}%
\pgfusepath{stroke}%
\end{pgfscope}%
\begin{pgfscope}%
\pgfpathrectangle{\pgfqpoint{0.650000in}{0.420000in}}{\pgfqpoint{2.388110in}{1.994488in}}%
\pgfusepath{clip}%
\pgfsetbuttcap%
\pgfsetroundjoin%
\pgfsetlinewidth{1.003750pt}%
\definecolor{currentstroke}{rgb}{0.498039,0.737255,0.254902}%
\pgfsetstrokecolor{currentstroke}%
\pgfsetdash{{1.000000pt}{5.000000pt}}{0.000000pt}%
\pgfpathmoveto{\pgfqpoint{0.645000in}{1.379648in}}%
\pgfpathlineto{\pgfqpoint{0.710402in}{1.390523in}}%
\pgfpathlineto{\pgfqpoint{0.737313in}{1.393982in}}%
\pgfpathlineto{\pgfqpoint{0.763003in}{1.396160in}}%
\pgfpathlineto{\pgfqpoint{0.787618in}{1.396961in}}%
\pgfpathlineto{\pgfqpoint{0.811281in}{1.396430in}}%
\pgfpathlineto{\pgfqpoint{0.834096in}{1.394768in}}%
\pgfpathlineto{\pgfqpoint{0.877521in}{1.389491in}}%
\pgfpathlineto{\pgfqpoint{0.898272in}{1.386824in}}%
\pgfpathlineto{\pgfqpoint{0.918461in}{1.384837in}}%
\pgfpathlineto{\pgfqpoint{0.938136in}{1.384004in}}%
\pgfpathlineto{\pgfqpoint{0.957343in}{1.384684in}}%
\pgfpathlineto{\pgfqpoint{0.976119in}{1.387043in}}%
\pgfpathlineto{\pgfqpoint{0.994499in}{1.390994in}}%
\pgfpathlineto{\pgfqpoint{1.012513in}{1.396160in}}%
\pgfpathlineto{\pgfqpoint{1.047553in}{1.407225in}}%
\pgfpathlineto{\pgfqpoint{1.064626in}{1.411212in}}%
\pgfpathlineto{\pgfqpoint{1.081428in}{1.412922in}}%
\pgfpathlineto{\pgfqpoint{1.097978in}{1.411805in}}%
\pgfpathlineto{\pgfqpoint{1.114292in}{1.407933in}}%
\pgfpathlineto{\pgfqpoint{1.130386in}{1.402142in}}%
\pgfpathlineto{\pgfqpoint{1.177482in}{1.383706in}}%
\pgfpathlineto{\pgfqpoint{1.192824in}{1.374663in}}%
\pgfpathlineto{\pgfqpoint{1.208006in}{1.362556in}}%
\pgfpathlineto{\pgfqpoint{1.223037in}{1.353319in}}%
\pgfpathlineto{\pgfqpoint{1.237925in}{1.354157in}}%
\pgfpathlineto{\pgfqpoint{1.252679in}{1.365055in}}%
\pgfpathlineto{\pgfqpoint{1.267305in}{1.379489in}}%
\pgfpathlineto{\pgfqpoint{1.281811in}{1.390787in}}%
\pgfpathlineto{\pgfqpoint{1.296204in}{1.395735in}}%
\pgfpathlineto{\pgfqpoint{1.310489in}{1.395198in}}%
\pgfpathlineto{\pgfqpoint{1.324672in}{1.392791in}}%
\pgfpathlineto{\pgfqpoint{1.338758in}{1.392248in}}%
\pgfpathlineto{\pgfqpoint{1.352752in}{1.395007in}}%
\pgfpathlineto{\pgfqpoint{1.366659in}{1.399467in}}%
\pgfpathlineto{\pgfqpoint{1.380483in}{1.402288in}}%
\pgfpathlineto{\pgfqpoint{1.394229in}{1.400974in}}%
\pgfpathlineto{\pgfqpoint{1.407900in}{1.395601in}}%
\pgfpathlineto{\pgfqpoint{1.435032in}{1.381153in}}%
\pgfpathlineto{\pgfqpoint{1.448499in}{1.377752in}}%
\pgfpathlineto{\pgfqpoint{1.461906in}{1.380131in}}%
\pgfpathlineto{\pgfqpoint{1.475253in}{1.386206in}}%
\pgfpathlineto{\pgfqpoint{1.488545in}{1.391132in}}%
\pgfpathlineto{\pgfqpoint{1.501784in}{1.392508in}}%
\pgfpathlineto{\pgfqpoint{1.514971in}{1.392349in}}%
\pgfpathlineto{\pgfqpoint{1.528111in}{1.393458in}}%
\pgfpathlineto{\pgfqpoint{1.541204in}{1.395406in}}%
\pgfpathlineto{\pgfqpoint{1.554253in}{1.395195in}}%
\pgfpathlineto{\pgfqpoint{1.580226in}{1.387884in}}%
\pgfpathlineto{\pgfqpoint{1.593154in}{1.387156in}}%
\pgfpathlineto{\pgfqpoint{1.606045in}{1.389831in}}%
\pgfpathlineto{\pgfqpoint{1.618901in}{1.393312in}}%
\pgfpathlineto{\pgfqpoint{1.631723in}{1.395199in}}%
\pgfpathlineto{\pgfqpoint{1.644513in}{1.395082in}}%
\pgfpathlineto{\pgfqpoint{1.657272in}{1.394220in}}%
\pgfpathlineto{\pgfqpoint{1.670002in}{1.394904in}}%
\pgfpathlineto{\pgfqpoint{1.682703in}{1.398988in}}%
\pgfpathlineto{\pgfqpoint{1.695377in}{1.404999in}}%
\pgfpathlineto{\pgfqpoint{1.708026in}{1.408568in}}%
\pgfpathlineto{\pgfqpoint{1.720650in}{1.407548in}}%
\pgfpathlineto{\pgfqpoint{1.733250in}{1.403472in}}%
\pgfpathlineto{\pgfqpoint{1.745827in}{1.398169in}}%
\pgfpathlineto{\pgfqpoint{1.758382in}{1.394116in}}%
\pgfpathlineto{\pgfqpoint{1.770917in}{1.394270in}}%
\pgfpathlineto{\pgfqpoint{1.783431in}{1.396949in}}%
\pgfpathlineto{\pgfqpoint{1.795926in}{1.397731in}}%
\pgfpathlineto{\pgfqpoint{1.858142in}{1.390751in}}%
\pgfpathlineto{\pgfqpoint{1.870537in}{1.389983in}}%
\pgfpathlineto{\pgfqpoint{1.882919in}{1.389906in}}%
\pgfpathlineto{\pgfqpoint{1.895286in}{1.391994in}}%
\pgfpathlineto{\pgfqpoint{1.919983in}{1.399846in}}%
\pgfpathlineto{\pgfqpoint{1.932312in}{1.400552in}}%
\pgfpathlineto{\pgfqpoint{1.944630in}{1.399214in}}%
\pgfpathlineto{\pgfqpoint{1.956937in}{1.398648in}}%
\pgfpathlineto{\pgfqpoint{1.981519in}{1.400696in}}%
\pgfpathlineto{\pgfqpoint{1.993795in}{1.400486in}}%
\pgfpathlineto{\pgfqpoint{2.006062in}{1.399356in}}%
\pgfpathlineto{\pgfqpoint{2.018319in}{1.397363in}}%
\pgfpathlineto{\pgfqpoint{2.030568in}{1.394583in}}%
\pgfpathlineto{\pgfqpoint{2.042808in}{1.393004in}}%
\pgfpathlineto{\pgfqpoint{2.079481in}{1.395147in}}%
\pgfpathlineto{\pgfqpoint{2.091691in}{1.395604in}}%
\pgfpathlineto{\pgfqpoint{2.103894in}{1.397990in}}%
\pgfpathlineto{\pgfqpoint{2.128281in}{1.406066in}}%
\pgfpathlineto{\pgfqpoint{2.140465in}{1.408333in}}%
\pgfpathlineto{\pgfqpoint{2.152643in}{1.407478in}}%
\pgfpathlineto{\pgfqpoint{2.176983in}{1.400325in}}%
\pgfpathlineto{\pgfqpoint{2.189145in}{1.398586in}}%
\pgfpathlineto{\pgfqpoint{2.201302in}{1.399383in}}%
\pgfpathlineto{\pgfqpoint{2.225602in}{1.402691in}}%
\pgfpathlineto{\pgfqpoint{2.237745in}{1.402821in}}%
\pgfpathlineto{\pgfqpoint{2.249884in}{1.401076in}}%
\pgfpathlineto{\pgfqpoint{2.262019in}{1.398193in}}%
\pgfpathlineto{\pgfqpoint{2.274150in}{1.396985in}}%
\pgfpathlineto{\pgfqpoint{2.286277in}{1.399028in}}%
\pgfpathlineto{\pgfqpoint{2.298401in}{1.402139in}}%
\pgfpathlineto{\pgfqpoint{2.310521in}{1.403300in}}%
\pgfpathlineto{\pgfqpoint{2.322638in}{1.401108in}}%
\pgfpathlineto{\pgfqpoint{2.334751in}{1.396284in}}%
\pgfpathlineto{\pgfqpoint{2.346862in}{1.392258in}}%
\pgfpathlineto{\pgfqpoint{2.358969in}{1.391971in}}%
\pgfpathlineto{\pgfqpoint{2.395181in}{1.399253in}}%
\pgfpathlineto{\pgfqpoint{2.418733in}{1.408471in}}%
\pgfpathlineto{\pgfqpoint{2.430067in}{1.409632in}}%
\pgfpathlineto{\pgfqpoint{2.441034in}{1.406313in}}%
\pgfpathlineto{\pgfqpoint{2.451632in}{1.402397in}}%
\pgfpathlineto{\pgfqpoint{2.461878in}{1.400759in}}%
\pgfpathlineto{\pgfqpoint{2.481398in}{1.400201in}}%
\pgfpathlineto{\pgfqpoint{2.490711in}{1.400777in}}%
\pgfpathlineto{\pgfqpoint{2.508530in}{1.402942in}}%
\pgfpathlineto{\pgfqpoint{2.517066in}{1.403284in}}%
\pgfpathlineto{\pgfqpoint{2.525371in}{1.404914in}}%
\pgfpathlineto{\pgfqpoint{2.533456in}{1.407745in}}%
\pgfpathlineto{\pgfqpoint{2.541335in}{1.408827in}}%
\pgfpathlineto{\pgfqpoint{2.549016in}{1.406560in}}%
\pgfpathlineto{\pgfqpoint{2.556509in}{1.402463in}}%
\pgfpathlineto{\pgfqpoint{2.570968in}{1.393180in}}%
\pgfpathlineto{\pgfqpoint{2.577950in}{1.390897in}}%
\pgfpathlineto{\pgfqpoint{2.584777in}{1.391619in}}%
\pgfpathlineto{\pgfqpoint{2.591455in}{1.393797in}}%
\pgfpathlineto{\pgfqpoint{2.597990in}{1.396661in}}%
\pgfpathlineto{\pgfqpoint{2.604390in}{1.400196in}}%
\pgfpathlineto{\pgfqpoint{2.610658in}{1.402756in}}%
\pgfpathlineto{\pgfqpoint{2.616802in}{1.402282in}}%
\pgfpathlineto{\pgfqpoint{2.628731in}{1.396846in}}%
\pgfpathlineto{\pgfqpoint{2.634526in}{1.395138in}}%
\pgfpathlineto{\pgfqpoint{2.645798in}{1.393357in}}%
\pgfpathlineto{\pgfqpoint{2.656670in}{1.389386in}}%
\pgfpathlineto{\pgfqpoint{2.661965in}{1.389881in}}%
\pgfpathlineto{\pgfqpoint{2.672288in}{1.395520in}}%
\pgfpathlineto{\pgfqpoint{2.677323in}{1.396502in}}%
\pgfpathlineto{\pgfqpoint{2.691948in}{1.394244in}}%
\pgfpathlineto{\pgfqpoint{2.701325in}{1.397646in}}%
\pgfpathlineto{\pgfqpoint{2.705908in}{1.397228in}}%
\pgfpathlineto{\pgfqpoint{2.714874in}{1.392461in}}%
\pgfpathlineto{\pgfqpoint{2.732056in}{1.391199in}}%
\pgfpathlineto{\pgfqpoint{2.736206in}{1.393532in}}%
\pgfpathlineto{\pgfqpoint{2.740300in}{1.396497in}}%
\pgfpathlineto{\pgfqpoint{2.744340in}{1.396879in}}%
\pgfpathlineto{\pgfqpoint{2.752263in}{1.392085in}}%
\pgfpathlineto{\pgfqpoint{2.756150in}{1.391399in}}%
\pgfpathlineto{\pgfqpoint{2.759988in}{1.393391in}}%
\pgfpathlineto{\pgfqpoint{2.763778in}{1.396587in}}%
\pgfpathlineto{\pgfqpoint{2.767523in}{1.398513in}}%
\pgfpathlineto{\pgfqpoint{2.771222in}{1.399269in}}%
\pgfpathlineto{\pgfqpoint{2.774877in}{1.400759in}}%
\pgfpathlineto{\pgfqpoint{2.778489in}{1.403151in}}%
\pgfpathlineto{\pgfqpoint{2.782059in}{1.404445in}}%
\pgfpathlineto{\pgfqpoint{2.785587in}{1.403237in}}%
\pgfpathlineto{\pgfqpoint{2.789076in}{1.400918in}}%
\pgfpathlineto{\pgfqpoint{2.802648in}{1.395855in}}%
\pgfpathlineto{\pgfqpoint{2.805949in}{1.396171in}}%
\pgfpathlineto{\pgfqpoint{2.809215in}{1.398078in}}%
\pgfpathlineto{\pgfqpoint{2.812446in}{1.400971in}}%
\pgfpathlineto{\pgfqpoint{2.818809in}{1.407536in}}%
\pgfpathlineto{\pgfqpoint{2.821942in}{1.406022in}}%
\pgfpathlineto{\pgfqpoint{2.825043in}{1.400691in}}%
\pgfpathlineto{\pgfqpoint{2.828112in}{1.397019in}}%
\pgfpathlineto{\pgfqpoint{2.834162in}{1.397151in}}%
\pgfpathlineto{\pgfqpoint{2.840094in}{1.393765in}}%
\pgfpathlineto{\pgfqpoint{2.843017in}{1.394058in}}%
\pgfpathlineto{\pgfqpoint{2.845913in}{1.396569in}}%
\pgfpathlineto{\pgfqpoint{2.848782in}{1.401805in}}%
\pgfpathlineto{\pgfqpoint{2.851625in}{1.409205in}}%
\pgfpathlineto{\pgfqpoint{2.854441in}{1.413611in}}%
\pgfpathlineto{\pgfqpoint{2.857232in}{1.410675in}}%
\pgfpathlineto{\pgfqpoint{2.859997in}{1.403876in}}%
\pgfpathlineto{\pgfqpoint{2.862738in}{1.398802in}}%
\pgfpathlineto{\pgfqpoint{2.865455in}{1.396567in}}%
\pgfpathlineto{\pgfqpoint{2.868147in}{1.396117in}}%
\pgfpathlineto{\pgfqpoint{2.873463in}{1.397509in}}%
\pgfpathlineto{\pgfqpoint{2.876087in}{1.397220in}}%
\pgfpathlineto{\pgfqpoint{2.881267in}{1.394387in}}%
\pgfpathlineto{\pgfqpoint{2.883825in}{1.394630in}}%
\pgfpathlineto{\pgfqpoint{2.888878in}{1.398729in}}%
\pgfpathlineto{\pgfqpoint{2.891374in}{1.398911in}}%
\pgfpathlineto{\pgfqpoint{2.896305in}{1.394874in}}%
\pgfpathlineto{\pgfqpoint{2.898741in}{1.394586in}}%
\pgfpathlineto{\pgfqpoint{2.901158in}{1.395395in}}%
\pgfpathlineto{\pgfqpoint{2.903556in}{1.395330in}}%
\pgfpathlineto{\pgfqpoint{2.910640in}{1.390906in}}%
\pgfpathlineto{\pgfqpoint{2.912965in}{1.391477in}}%
\pgfpathlineto{\pgfqpoint{2.917564in}{1.395386in}}%
\pgfpathlineto{\pgfqpoint{2.919838in}{1.397181in}}%
\pgfpathlineto{\pgfqpoint{2.922095in}{1.397239in}}%
\pgfpathlineto{\pgfqpoint{2.924335in}{1.395343in}}%
\pgfpathlineto{\pgfqpoint{2.928768in}{1.390089in}}%
\pgfpathlineto{\pgfqpoint{2.930960in}{1.390043in}}%
\pgfpathlineto{\pgfqpoint{2.933137in}{1.392834in}}%
\pgfpathlineto{\pgfqpoint{2.935299in}{1.396977in}}%
\pgfpathlineto{\pgfqpoint{2.937445in}{1.399579in}}%
\pgfpathlineto{\pgfqpoint{2.939576in}{1.398558in}}%
\pgfpathlineto{\pgfqpoint{2.943796in}{1.390941in}}%
\pgfpathlineto{\pgfqpoint{2.945884in}{1.389689in}}%
\pgfpathlineto{\pgfqpoint{2.947957in}{1.390698in}}%
\pgfpathlineto{\pgfqpoint{2.952064in}{1.396118in}}%
\pgfpathlineto{\pgfqpoint{2.954096in}{1.398943in}}%
\pgfpathlineto{\pgfqpoint{2.956116in}{1.399835in}}%
\pgfpathlineto{\pgfqpoint{2.958122in}{1.398160in}}%
\pgfpathlineto{\pgfqpoint{2.960115in}{1.395509in}}%
\pgfpathlineto{\pgfqpoint{2.962095in}{1.394720in}}%
\pgfpathlineto{\pgfqpoint{2.966018in}{1.399693in}}%
\pgfpathlineto{\pgfqpoint{2.967960in}{1.399071in}}%
\pgfpathlineto{\pgfqpoint{2.971809in}{1.392238in}}%
\pgfpathlineto{\pgfqpoint{2.973715in}{1.392556in}}%
\pgfpathlineto{\pgfqpoint{2.975610in}{1.394457in}}%
\pgfpathlineto{\pgfqpoint{2.977493in}{1.395193in}}%
\pgfpathlineto{\pgfqpoint{2.979364in}{1.393741in}}%
\pgfpathlineto{\pgfqpoint{2.981224in}{1.391110in}}%
\pgfpathlineto{\pgfqpoint{2.983073in}{1.389512in}}%
\pgfpathlineto{\pgfqpoint{2.984911in}{1.390894in}}%
\pgfpathlineto{\pgfqpoint{2.988554in}{1.399076in}}%
\pgfpathlineto{\pgfqpoint{2.990360in}{1.400346in}}%
\pgfpathlineto{\pgfqpoint{2.993939in}{1.397319in}}%
\pgfpathlineto{\pgfqpoint{2.997477in}{1.394802in}}%
\pgfpathlineto{\pgfqpoint{2.999231in}{1.394338in}}%
\pgfpathlineto{\pgfqpoint{3.000975in}{1.395223in}}%
\pgfpathlineto{\pgfqpoint{3.002709in}{1.396920in}}%
\pgfpathlineto{\pgfqpoint{3.004433in}{1.397503in}}%
\pgfpathlineto{\pgfqpoint{3.007853in}{1.394037in}}%
\pgfpathlineto{\pgfqpoint{3.009548in}{1.394208in}}%
\pgfpathlineto{\pgfqpoint{3.011235in}{1.395064in}}%
\pgfpathlineto{\pgfqpoint{3.014580in}{1.394068in}}%
\pgfpathlineto{\pgfqpoint{3.017889in}{1.396366in}}%
\pgfpathlineto{\pgfqpoint{3.021163in}{1.396133in}}%
\pgfpathlineto{\pgfqpoint{3.026009in}{1.398494in}}%
\pgfpathlineto{\pgfqpoint{3.027607in}{1.396058in}}%
\pgfpathlineto{\pgfqpoint{3.030779in}{1.389940in}}%
\pgfpathlineto{\pgfqpoint{3.032353in}{1.390612in}}%
\pgfpathlineto{\pgfqpoint{3.037027in}{1.398738in}}%
\pgfpathlineto{\pgfqpoint{3.038569in}{1.398867in}}%
\pgfpathlineto{\pgfqpoint{3.043110in}{1.394854in}}%
\pgfpathlineto{\pgfqpoint{3.043110in}{1.394854in}}%
\pgfusepath{stroke}%
\end{pgfscope}%
\begin{pgfscope}%
\pgfsetrectcap%
\pgfsetmiterjoin%
\pgfsetlinewidth{0.803000pt}%
\definecolor{currentstroke}{rgb}{0.000000,0.000000,0.000000}%
\pgfsetstrokecolor{currentstroke}%
\pgfsetdash{}{0pt}%
\pgfpathmoveto{\pgfqpoint{0.650000in}{0.420000in}}%
\pgfpathlineto{\pgfqpoint{0.650000in}{2.414488in}}%
\pgfusepath{stroke}%
\end{pgfscope}%
\begin{pgfscope}%
\pgfsetrectcap%
\pgfsetmiterjoin%
\pgfsetlinewidth{0.803000pt}%
\definecolor{currentstroke}{rgb}{0.000000,0.000000,0.000000}%
\pgfsetstrokecolor{currentstroke}%
\pgfsetdash{}{0pt}%
\pgfpathmoveto{\pgfqpoint{3.038110in}{0.420000in}}%
\pgfpathlineto{\pgfqpoint{3.038110in}{2.414488in}}%
\pgfusepath{stroke}%
\end{pgfscope}%
\begin{pgfscope}%
\pgfsetrectcap%
\pgfsetmiterjoin%
\pgfsetlinewidth{0.803000pt}%
\definecolor{currentstroke}{rgb}{0.000000,0.000000,0.000000}%
\pgfsetstrokecolor{currentstroke}%
\pgfsetdash{}{0pt}%
\pgfpathmoveto{\pgfqpoint{0.650000in}{0.420000in}}%
\pgfpathlineto{\pgfqpoint{3.038110in}{0.420000in}}%
\pgfusepath{stroke}%
\end{pgfscope}%
\begin{pgfscope}%
\pgfsetrectcap%
\pgfsetmiterjoin%
\pgfsetlinewidth{0.803000pt}%
\definecolor{currentstroke}{rgb}{0.000000,0.000000,0.000000}%
\pgfsetstrokecolor{currentstroke}%
\pgfsetdash{}{0pt}%
\pgfpathmoveto{\pgfqpoint{0.650000in}{2.414488in}}%
\pgfpathlineto{\pgfqpoint{3.038110in}{2.414488in}}%
\pgfusepath{stroke}%
\end{pgfscope}%
\begin{pgfscope}%
\pgfsetrectcap%
\pgfsetroundjoin%
\pgfsetlinewidth{0.602250pt}%
\definecolor{currentstroke}{rgb}{0.000000,0.000000,0.000000}%
\pgfsetstrokecolor{currentstroke}%
\pgfsetdash{}{0pt}%
\pgfpathmoveto{\pgfqpoint{2.127879in}{2.275599in}}%
\pgfpathlineto{\pgfqpoint{2.238990in}{2.275599in}}%
\pgfpathlineto{\pgfqpoint{2.350101in}{2.275599in}}%
\pgfusepath{stroke}%
\end{pgfscope}%
\begin{pgfscope}%
\definecolor{textcolor}{rgb}{0.000000,0.000000,0.000000}%
\pgfsetstrokecolor{textcolor}%
\pgfsetfillcolor{textcolor}%
\pgftext[x=2.438990in,y=2.236710in,left,base]{\color{textcolor}{\rmfamily\fontsize{8.000000}{9.600000}\selectfont\catcode`\^=\active\def^{\ifmmode\sp\else\^{}\fi}\catcode`\%=\active\def%{\%}ref}}%
\end{pgfscope}%
\begin{pgfscope}%
\pgfsetrectcap%
\pgfsetroundjoin%
\pgfsetlinewidth{1.003750pt}%
\definecolor{currentstroke}{rgb}{0.870588,0.466667,0.682353}%
\pgfsetstrokecolor{currentstroke}%
\pgfsetdash{}{0pt}%
\pgfpathmoveto{\pgfqpoint{2.127879in}{2.120661in}}%
\pgfpathlineto{\pgfqpoint{2.238990in}{2.120661in}}%
\pgfpathlineto{\pgfqpoint{2.350101in}{2.120661in}}%
\pgfusepath{stroke}%
\end{pgfscope}%
\begin{pgfscope}%
\definecolor{textcolor}{rgb}{0.000000,0.000000,0.000000}%
\pgfsetstrokecolor{textcolor}%
\pgfsetfillcolor{textcolor}%
\pgftext[x=2.438990in,y=2.081772in,left,base]{\color{textcolor}{\rmfamily\fontsize{8.000000}{9.600000}\selectfont\catcode`\^=\active\def^{\ifmmode\sp\else\^{}\fi}\catcode`\%=\active\def%{\%}base}}%
\end{pgfscope}%
\begin{pgfscope}%
\pgfsetrectcap%
\pgfsetroundjoin%
\pgfsetlinewidth{1.003750pt}%
\definecolor{currentstroke}{rgb}{0.498039,0.737255,0.254902}%
\pgfsetstrokecolor{currentstroke}%
\pgfsetdash{}{0pt}%
\pgfpathmoveto{\pgfqpoint{2.127879in}{1.965723in}}%
\pgfpathlineto{\pgfqpoint{2.238990in}{1.965723in}}%
\pgfpathlineto{\pgfqpoint{2.350101in}{1.965723in}}%
\pgfusepath{stroke}%
\end{pgfscope}%
\begin{pgfscope}%
\definecolor{textcolor}{rgb}{0.000000,0.000000,0.000000}%
\pgfsetstrokecolor{textcolor}%
\pgfsetfillcolor{textcolor}%
\pgftext[x=2.438990in,y=1.926834in,left,base]{\color{textcolor}{\rmfamily\fontsize{8.000000}{9.600000}\selectfont\catcode`\^=\active\def^{\ifmmode\sp\else\^{}\fi}\catcode`\%=\active\def%{\%}optimised}}%
\end{pgfscope}%
\end{pgfpicture}%
\makeatother%
\endgroup%

    \label{fig:R100deg30Re1200Budget}
\end{subfigure}
\hfill
\begin{subfigure}[t]{0.495\textwidth}
    \centering
    \subcaption{}
    %% Creator: Matplotlib, PGF backend
%%
%% To include the figure in your LaTeX document, write
%%   \input{<filename>.pgf}
%%
%% Make sure the required packages are loaded in your preamble
%%   \usepackage{pgf}
%%
%% Also ensure that all the required font packages are loaded; for instance,
%% the lmodern package is sometimes necessary when using math font.
%%   \usepackage{lmodern}
%%
%% Figures using additional raster images can only be included by \input if
%% they are in the same directory as the main LaTeX file. For loading figures
%% from other directories you can use the `import` package
%%   \usepackage{import}
%%
%% and then include the figures with
%%   \import{<path to file>}{<filename>.pgf}
%%
%% Matplotlib used the following preamble
%%   \def\mathdefault#1{#1}
%%   \everymath=\expandafter{\the\everymath\displaystyle}
%%   \IfFileExists{scrextend.sty}{
%%     \usepackage[fontsize=11.000000pt]{scrextend}
%%   }{
%%     \renewcommand{\normalsize}{\fontsize{11.000000}{13.200000}\selectfont}
%%     \normalsize
%%   }
%%   
%%   \makeatletter\@ifpackageloaded{underscore}{}{\usepackage[strings]{underscore}}\makeatother
%%
\begingroup%
\makeatletter%
\begin{pgfpicture}%
\pgfpathrectangle{\pgfpointorigin}{\pgfqpoint{3.118110in}{2.494488in}}%
\pgfusepath{use as bounding box, clip}%
\begin{pgfscope}%
\pgfsetbuttcap%
\pgfsetmiterjoin%
\definecolor{currentfill}{rgb}{1.000000,1.000000,1.000000}%
\pgfsetfillcolor{currentfill}%
\pgfsetlinewidth{0.000000pt}%
\definecolor{currentstroke}{rgb}{1.000000,1.000000,1.000000}%
\pgfsetstrokecolor{currentstroke}%
\pgfsetdash{}{0pt}%
\pgfpathmoveto{\pgfqpoint{0.000000in}{0.000000in}}%
\pgfpathlineto{\pgfqpoint{3.118110in}{0.000000in}}%
\pgfpathlineto{\pgfqpoint{3.118110in}{2.494488in}}%
\pgfpathlineto{\pgfqpoint{0.000000in}{2.494488in}}%
\pgfpathlineto{\pgfqpoint{0.000000in}{0.000000in}}%
\pgfpathclose%
\pgfusepath{fill}%
\end{pgfscope}%
\begin{pgfscope}%
\pgfsetbuttcap%
\pgfsetmiterjoin%
\definecolor{currentfill}{rgb}{1.000000,1.000000,1.000000}%
\pgfsetfillcolor{currentfill}%
\pgfsetlinewidth{0.000000pt}%
\definecolor{currentstroke}{rgb}{0.000000,0.000000,0.000000}%
\pgfsetstrokecolor{currentstroke}%
\pgfsetstrokeopacity{0.000000}%
\pgfsetdash{}{0pt}%
\pgfpathmoveto{\pgfqpoint{0.650000in}{0.420000in}}%
\pgfpathlineto{\pgfqpoint{3.038110in}{0.420000in}}%
\pgfpathlineto{\pgfqpoint{3.038110in}{2.414488in}}%
\pgfpathlineto{\pgfqpoint{0.650000in}{2.414488in}}%
\pgfpathlineto{\pgfqpoint{0.650000in}{0.420000in}}%
\pgfpathclose%
\pgfusepath{fill}%
\end{pgfscope}%
\begin{pgfscope}%
\pgfsetbuttcap%
\pgfsetroundjoin%
\definecolor{currentfill}{rgb}{0.000000,0.000000,0.000000}%
\pgfsetfillcolor{currentfill}%
\pgfsetlinewidth{0.501875pt}%
\definecolor{currentstroke}{rgb}{0.000000,0.000000,0.000000}%
\pgfsetstrokecolor{currentstroke}%
\pgfsetdash{}{0pt}%
\pgfsys@defobject{currentmarker}{\pgfqpoint{0.000000in}{0.000000in}}{\pgfqpoint{0.000000in}{0.041667in}}{%
\pgfpathmoveto{\pgfqpoint{0.000000in}{0.000000in}}%
\pgfpathlineto{\pgfqpoint{0.000000in}{0.041667in}}%
\pgfusepath{stroke,fill}%
}%
\begin{pgfscope}%
\pgfsys@transformshift{0.650000in}{0.420000in}%
\pgfsys@useobject{currentmarker}{}%
\end{pgfscope}%
\end{pgfscope}%
\begin{pgfscope}%
\pgfsetbuttcap%
\pgfsetroundjoin%
\definecolor{currentfill}{rgb}{0.000000,0.000000,0.000000}%
\pgfsetfillcolor{currentfill}%
\pgfsetlinewidth{0.501875pt}%
\definecolor{currentstroke}{rgb}{0.000000,0.000000,0.000000}%
\pgfsetstrokecolor{currentstroke}%
\pgfsetdash{}{0pt}%
\pgfsys@defobject{currentmarker}{\pgfqpoint{0.000000in}{-0.041667in}}{\pgfqpoint{0.000000in}{0.000000in}}{%
\pgfpathmoveto{\pgfqpoint{0.000000in}{0.000000in}}%
\pgfpathlineto{\pgfqpoint{0.000000in}{-0.041667in}}%
\pgfusepath{stroke,fill}%
}%
\begin{pgfscope}%
\pgfsys@transformshift{0.650000in}{2.414488in}%
\pgfsys@useobject{currentmarker}{}%
\end{pgfscope}%
\end{pgfscope}%
\begin{pgfscope}%
\definecolor{textcolor}{rgb}{0.000000,0.000000,0.000000}%
\pgfsetstrokecolor{textcolor}%
\pgfsetfillcolor{textcolor}%
\pgftext[x=0.650000in,y=0.371389in,,top]{\color{textcolor}{\rmfamily\fontsize{11.000000}{13.200000}\selectfont\catcode`\^=\active\def^{\ifmmode\sp\else\^{}\fi}\catcode`\%=\active\def%{\%}0}}%
\end{pgfscope}%
\begin{pgfscope}%
\pgfsetbuttcap%
\pgfsetroundjoin%
\definecolor{currentfill}{rgb}{0.000000,0.000000,0.000000}%
\pgfsetfillcolor{currentfill}%
\pgfsetlinewidth{0.501875pt}%
\definecolor{currentstroke}{rgb}{0.000000,0.000000,0.000000}%
\pgfsetstrokecolor{currentstroke}%
\pgfsetdash{}{0pt}%
\pgfsys@defobject{currentmarker}{\pgfqpoint{0.000000in}{0.000000in}}{\pgfqpoint{0.000000in}{0.041667in}}{%
\pgfpathmoveto{\pgfqpoint{0.000000in}{0.000000in}}%
\pgfpathlineto{\pgfqpoint{0.000000in}{0.041667in}}%
\pgfusepath{stroke,fill}%
}%
\begin{pgfscope}%
\pgfsys@transformshift{1.492631in}{0.420000in}%
\pgfsys@useobject{currentmarker}{}%
\end{pgfscope}%
\end{pgfscope}%
\begin{pgfscope}%
\pgfsetbuttcap%
\pgfsetroundjoin%
\definecolor{currentfill}{rgb}{0.000000,0.000000,0.000000}%
\pgfsetfillcolor{currentfill}%
\pgfsetlinewidth{0.501875pt}%
\definecolor{currentstroke}{rgb}{0.000000,0.000000,0.000000}%
\pgfsetstrokecolor{currentstroke}%
\pgfsetdash{}{0pt}%
\pgfsys@defobject{currentmarker}{\pgfqpoint{0.000000in}{-0.041667in}}{\pgfqpoint{0.000000in}{0.000000in}}{%
\pgfpathmoveto{\pgfqpoint{0.000000in}{0.000000in}}%
\pgfpathlineto{\pgfqpoint{0.000000in}{-0.041667in}}%
\pgfusepath{stroke,fill}%
}%
\begin{pgfscope}%
\pgfsys@transformshift{1.492631in}{2.414488in}%
\pgfsys@useobject{currentmarker}{}%
\end{pgfscope}%
\end{pgfscope}%
\begin{pgfscope}%
\definecolor{textcolor}{rgb}{0.000000,0.000000,0.000000}%
\pgfsetstrokecolor{textcolor}%
\pgfsetfillcolor{textcolor}%
\pgftext[x=1.492631in,y=0.371389in,,top]{\color{textcolor}{\rmfamily\fontsize{11.000000}{13.200000}\selectfont\catcode`\^=\active\def^{\ifmmode\sp\else\^{}\fi}\catcode`\%=\active\def%{\%}50}}%
\end{pgfscope}%
\begin{pgfscope}%
\pgfsetbuttcap%
\pgfsetroundjoin%
\definecolor{currentfill}{rgb}{0.000000,0.000000,0.000000}%
\pgfsetfillcolor{currentfill}%
\pgfsetlinewidth{0.501875pt}%
\definecolor{currentstroke}{rgb}{0.000000,0.000000,0.000000}%
\pgfsetstrokecolor{currentstroke}%
\pgfsetdash{}{0pt}%
\pgfsys@defobject{currentmarker}{\pgfqpoint{0.000000in}{0.000000in}}{\pgfqpoint{0.000000in}{0.041667in}}{%
\pgfpathmoveto{\pgfqpoint{0.000000in}{0.000000in}}%
\pgfpathlineto{\pgfqpoint{0.000000in}{0.041667in}}%
\pgfusepath{stroke,fill}%
}%
\begin{pgfscope}%
\pgfsys@transformshift{2.335263in}{0.420000in}%
\pgfsys@useobject{currentmarker}{}%
\end{pgfscope}%
\end{pgfscope}%
\begin{pgfscope}%
\pgfsetbuttcap%
\pgfsetroundjoin%
\definecolor{currentfill}{rgb}{0.000000,0.000000,0.000000}%
\pgfsetfillcolor{currentfill}%
\pgfsetlinewidth{0.501875pt}%
\definecolor{currentstroke}{rgb}{0.000000,0.000000,0.000000}%
\pgfsetstrokecolor{currentstroke}%
\pgfsetdash{}{0pt}%
\pgfsys@defobject{currentmarker}{\pgfqpoint{0.000000in}{-0.041667in}}{\pgfqpoint{0.000000in}{0.000000in}}{%
\pgfpathmoveto{\pgfqpoint{0.000000in}{0.000000in}}%
\pgfpathlineto{\pgfqpoint{0.000000in}{-0.041667in}}%
\pgfusepath{stroke,fill}%
}%
\begin{pgfscope}%
\pgfsys@transformshift{2.335263in}{2.414488in}%
\pgfsys@useobject{currentmarker}{}%
\end{pgfscope}%
\end{pgfscope}%
\begin{pgfscope}%
\definecolor{textcolor}{rgb}{0.000000,0.000000,0.000000}%
\pgfsetstrokecolor{textcolor}%
\pgfsetfillcolor{textcolor}%
\pgftext[x=2.335263in,y=0.371389in,,top]{\color{textcolor}{\rmfamily\fontsize{11.000000}{13.200000}\selectfont\catcode`\^=\active\def^{\ifmmode\sp\else\^{}\fi}\catcode`\%=\active\def%{\%}100}}%
\end{pgfscope}%
\begin{pgfscope}%
\pgfsetbuttcap%
\pgfsetroundjoin%
\definecolor{currentfill}{rgb}{0.000000,0.000000,0.000000}%
\pgfsetfillcolor{currentfill}%
\pgfsetlinewidth{0.401500pt}%
\definecolor{currentstroke}{rgb}{0.000000,0.000000,0.000000}%
\pgfsetstrokecolor{currentstroke}%
\pgfsetdash{}{0pt}%
\pgfsys@defobject{currentmarker}{\pgfqpoint{0.000000in}{0.000000in}}{\pgfqpoint{0.000000in}{0.020833in}}{%
\pgfpathmoveto{\pgfqpoint{0.000000in}{0.000000in}}%
\pgfpathlineto{\pgfqpoint{0.000000in}{0.020833in}}%
\pgfusepath{stroke,fill}%
}%
\begin{pgfscope}%
\pgfsys@transformshift{0.818526in}{0.420000in}%
\pgfsys@useobject{currentmarker}{}%
\end{pgfscope}%
\end{pgfscope}%
\begin{pgfscope}%
\pgfsetbuttcap%
\pgfsetroundjoin%
\definecolor{currentfill}{rgb}{0.000000,0.000000,0.000000}%
\pgfsetfillcolor{currentfill}%
\pgfsetlinewidth{0.401500pt}%
\definecolor{currentstroke}{rgb}{0.000000,0.000000,0.000000}%
\pgfsetstrokecolor{currentstroke}%
\pgfsetdash{}{0pt}%
\pgfsys@defobject{currentmarker}{\pgfqpoint{0.000000in}{-0.020833in}}{\pgfqpoint{0.000000in}{0.000000in}}{%
\pgfpathmoveto{\pgfqpoint{0.000000in}{0.000000in}}%
\pgfpathlineto{\pgfqpoint{0.000000in}{-0.020833in}}%
\pgfusepath{stroke,fill}%
}%
\begin{pgfscope}%
\pgfsys@transformshift{0.818526in}{2.414488in}%
\pgfsys@useobject{currentmarker}{}%
\end{pgfscope}%
\end{pgfscope}%
\begin{pgfscope}%
\pgfsetbuttcap%
\pgfsetroundjoin%
\definecolor{currentfill}{rgb}{0.000000,0.000000,0.000000}%
\pgfsetfillcolor{currentfill}%
\pgfsetlinewidth{0.401500pt}%
\definecolor{currentstroke}{rgb}{0.000000,0.000000,0.000000}%
\pgfsetstrokecolor{currentstroke}%
\pgfsetdash{}{0pt}%
\pgfsys@defobject{currentmarker}{\pgfqpoint{0.000000in}{0.000000in}}{\pgfqpoint{0.000000in}{0.020833in}}{%
\pgfpathmoveto{\pgfqpoint{0.000000in}{0.000000in}}%
\pgfpathlineto{\pgfqpoint{0.000000in}{0.020833in}}%
\pgfusepath{stroke,fill}%
}%
\begin{pgfscope}%
\pgfsys@transformshift{0.987053in}{0.420000in}%
\pgfsys@useobject{currentmarker}{}%
\end{pgfscope}%
\end{pgfscope}%
\begin{pgfscope}%
\pgfsetbuttcap%
\pgfsetroundjoin%
\definecolor{currentfill}{rgb}{0.000000,0.000000,0.000000}%
\pgfsetfillcolor{currentfill}%
\pgfsetlinewidth{0.401500pt}%
\definecolor{currentstroke}{rgb}{0.000000,0.000000,0.000000}%
\pgfsetstrokecolor{currentstroke}%
\pgfsetdash{}{0pt}%
\pgfsys@defobject{currentmarker}{\pgfqpoint{0.000000in}{-0.020833in}}{\pgfqpoint{0.000000in}{0.000000in}}{%
\pgfpathmoveto{\pgfqpoint{0.000000in}{0.000000in}}%
\pgfpathlineto{\pgfqpoint{0.000000in}{-0.020833in}}%
\pgfusepath{stroke,fill}%
}%
\begin{pgfscope}%
\pgfsys@transformshift{0.987053in}{2.414488in}%
\pgfsys@useobject{currentmarker}{}%
\end{pgfscope}%
\end{pgfscope}%
\begin{pgfscope}%
\pgfsetbuttcap%
\pgfsetroundjoin%
\definecolor{currentfill}{rgb}{0.000000,0.000000,0.000000}%
\pgfsetfillcolor{currentfill}%
\pgfsetlinewidth{0.401500pt}%
\definecolor{currentstroke}{rgb}{0.000000,0.000000,0.000000}%
\pgfsetstrokecolor{currentstroke}%
\pgfsetdash{}{0pt}%
\pgfsys@defobject{currentmarker}{\pgfqpoint{0.000000in}{0.000000in}}{\pgfqpoint{0.000000in}{0.020833in}}{%
\pgfpathmoveto{\pgfqpoint{0.000000in}{0.000000in}}%
\pgfpathlineto{\pgfqpoint{0.000000in}{0.020833in}}%
\pgfusepath{stroke,fill}%
}%
\begin{pgfscope}%
\pgfsys@transformshift{1.155579in}{0.420000in}%
\pgfsys@useobject{currentmarker}{}%
\end{pgfscope}%
\end{pgfscope}%
\begin{pgfscope}%
\pgfsetbuttcap%
\pgfsetroundjoin%
\definecolor{currentfill}{rgb}{0.000000,0.000000,0.000000}%
\pgfsetfillcolor{currentfill}%
\pgfsetlinewidth{0.401500pt}%
\definecolor{currentstroke}{rgb}{0.000000,0.000000,0.000000}%
\pgfsetstrokecolor{currentstroke}%
\pgfsetdash{}{0pt}%
\pgfsys@defobject{currentmarker}{\pgfqpoint{0.000000in}{-0.020833in}}{\pgfqpoint{0.000000in}{0.000000in}}{%
\pgfpathmoveto{\pgfqpoint{0.000000in}{0.000000in}}%
\pgfpathlineto{\pgfqpoint{0.000000in}{-0.020833in}}%
\pgfusepath{stroke,fill}%
}%
\begin{pgfscope}%
\pgfsys@transformshift{1.155579in}{2.414488in}%
\pgfsys@useobject{currentmarker}{}%
\end{pgfscope}%
\end{pgfscope}%
\begin{pgfscope}%
\pgfsetbuttcap%
\pgfsetroundjoin%
\definecolor{currentfill}{rgb}{0.000000,0.000000,0.000000}%
\pgfsetfillcolor{currentfill}%
\pgfsetlinewidth{0.401500pt}%
\definecolor{currentstroke}{rgb}{0.000000,0.000000,0.000000}%
\pgfsetstrokecolor{currentstroke}%
\pgfsetdash{}{0pt}%
\pgfsys@defobject{currentmarker}{\pgfqpoint{0.000000in}{0.000000in}}{\pgfqpoint{0.000000in}{0.020833in}}{%
\pgfpathmoveto{\pgfqpoint{0.000000in}{0.000000in}}%
\pgfpathlineto{\pgfqpoint{0.000000in}{0.020833in}}%
\pgfusepath{stroke,fill}%
}%
\begin{pgfscope}%
\pgfsys@transformshift{1.324105in}{0.420000in}%
\pgfsys@useobject{currentmarker}{}%
\end{pgfscope}%
\end{pgfscope}%
\begin{pgfscope}%
\pgfsetbuttcap%
\pgfsetroundjoin%
\definecolor{currentfill}{rgb}{0.000000,0.000000,0.000000}%
\pgfsetfillcolor{currentfill}%
\pgfsetlinewidth{0.401500pt}%
\definecolor{currentstroke}{rgb}{0.000000,0.000000,0.000000}%
\pgfsetstrokecolor{currentstroke}%
\pgfsetdash{}{0pt}%
\pgfsys@defobject{currentmarker}{\pgfqpoint{0.000000in}{-0.020833in}}{\pgfqpoint{0.000000in}{0.000000in}}{%
\pgfpathmoveto{\pgfqpoint{0.000000in}{0.000000in}}%
\pgfpathlineto{\pgfqpoint{0.000000in}{-0.020833in}}%
\pgfusepath{stroke,fill}%
}%
\begin{pgfscope}%
\pgfsys@transformshift{1.324105in}{2.414488in}%
\pgfsys@useobject{currentmarker}{}%
\end{pgfscope}%
\end{pgfscope}%
\begin{pgfscope}%
\pgfsetbuttcap%
\pgfsetroundjoin%
\definecolor{currentfill}{rgb}{0.000000,0.000000,0.000000}%
\pgfsetfillcolor{currentfill}%
\pgfsetlinewidth{0.401500pt}%
\definecolor{currentstroke}{rgb}{0.000000,0.000000,0.000000}%
\pgfsetstrokecolor{currentstroke}%
\pgfsetdash{}{0pt}%
\pgfsys@defobject{currentmarker}{\pgfqpoint{0.000000in}{0.000000in}}{\pgfqpoint{0.000000in}{0.020833in}}{%
\pgfpathmoveto{\pgfqpoint{0.000000in}{0.000000in}}%
\pgfpathlineto{\pgfqpoint{0.000000in}{0.020833in}}%
\pgfusepath{stroke,fill}%
}%
\begin{pgfscope}%
\pgfsys@transformshift{1.661158in}{0.420000in}%
\pgfsys@useobject{currentmarker}{}%
\end{pgfscope}%
\end{pgfscope}%
\begin{pgfscope}%
\pgfsetbuttcap%
\pgfsetroundjoin%
\definecolor{currentfill}{rgb}{0.000000,0.000000,0.000000}%
\pgfsetfillcolor{currentfill}%
\pgfsetlinewidth{0.401500pt}%
\definecolor{currentstroke}{rgb}{0.000000,0.000000,0.000000}%
\pgfsetstrokecolor{currentstroke}%
\pgfsetdash{}{0pt}%
\pgfsys@defobject{currentmarker}{\pgfqpoint{0.000000in}{-0.020833in}}{\pgfqpoint{0.000000in}{0.000000in}}{%
\pgfpathmoveto{\pgfqpoint{0.000000in}{0.000000in}}%
\pgfpathlineto{\pgfqpoint{0.000000in}{-0.020833in}}%
\pgfusepath{stroke,fill}%
}%
\begin{pgfscope}%
\pgfsys@transformshift{1.661158in}{2.414488in}%
\pgfsys@useobject{currentmarker}{}%
\end{pgfscope}%
\end{pgfscope}%
\begin{pgfscope}%
\pgfsetbuttcap%
\pgfsetroundjoin%
\definecolor{currentfill}{rgb}{0.000000,0.000000,0.000000}%
\pgfsetfillcolor{currentfill}%
\pgfsetlinewidth{0.401500pt}%
\definecolor{currentstroke}{rgb}{0.000000,0.000000,0.000000}%
\pgfsetstrokecolor{currentstroke}%
\pgfsetdash{}{0pt}%
\pgfsys@defobject{currentmarker}{\pgfqpoint{0.000000in}{0.000000in}}{\pgfqpoint{0.000000in}{0.020833in}}{%
\pgfpathmoveto{\pgfqpoint{0.000000in}{0.000000in}}%
\pgfpathlineto{\pgfqpoint{0.000000in}{0.020833in}}%
\pgfusepath{stroke,fill}%
}%
\begin{pgfscope}%
\pgfsys@transformshift{1.829684in}{0.420000in}%
\pgfsys@useobject{currentmarker}{}%
\end{pgfscope}%
\end{pgfscope}%
\begin{pgfscope}%
\pgfsetbuttcap%
\pgfsetroundjoin%
\definecolor{currentfill}{rgb}{0.000000,0.000000,0.000000}%
\pgfsetfillcolor{currentfill}%
\pgfsetlinewidth{0.401500pt}%
\definecolor{currentstroke}{rgb}{0.000000,0.000000,0.000000}%
\pgfsetstrokecolor{currentstroke}%
\pgfsetdash{}{0pt}%
\pgfsys@defobject{currentmarker}{\pgfqpoint{0.000000in}{-0.020833in}}{\pgfqpoint{0.000000in}{0.000000in}}{%
\pgfpathmoveto{\pgfqpoint{0.000000in}{0.000000in}}%
\pgfpathlineto{\pgfqpoint{0.000000in}{-0.020833in}}%
\pgfusepath{stroke,fill}%
}%
\begin{pgfscope}%
\pgfsys@transformshift{1.829684in}{2.414488in}%
\pgfsys@useobject{currentmarker}{}%
\end{pgfscope}%
\end{pgfscope}%
\begin{pgfscope}%
\pgfsetbuttcap%
\pgfsetroundjoin%
\definecolor{currentfill}{rgb}{0.000000,0.000000,0.000000}%
\pgfsetfillcolor{currentfill}%
\pgfsetlinewidth{0.401500pt}%
\definecolor{currentstroke}{rgb}{0.000000,0.000000,0.000000}%
\pgfsetstrokecolor{currentstroke}%
\pgfsetdash{}{0pt}%
\pgfsys@defobject{currentmarker}{\pgfqpoint{0.000000in}{0.000000in}}{\pgfqpoint{0.000000in}{0.020833in}}{%
\pgfpathmoveto{\pgfqpoint{0.000000in}{0.000000in}}%
\pgfpathlineto{\pgfqpoint{0.000000in}{0.020833in}}%
\pgfusepath{stroke,fill}%
}%
\begin{pgfscope}%
\pgfsys@transformshift{1.998210in}{0.420000in}%
\pgfsys@useobject{currentmarker}{}%
\end{pgfscope}%
\end{pgfscope}%
\begin{pgfscope}%
\pgfsetbuttcap%
\pgfsetroundjoin%
\definecolor{currentfill}{rgb}{0.000000,0.000000,0.000000}%
\pgfsetfillcolor{currentfill}%
\pgfsetlinewidth{0.401500pt}%
\definecolor{currentstroke}{rgb}{0.000000,0.000000,0.000000}%
\pgfsetstrokecolor{currentstroke}%
\pgfsetdash{}{0pt}%
\pgfsys@defobject{currentmarker}{\pgfqpoint{0.000000in}{-0.020833in}}{\pgfqpoint{0.000000in}{0.000000in}}{%
\pgfpathmoveto{\pgfqpoint{0.000000in}{0.000000in}}%
\pgfpathlineto{\pgfqpoint{0.000000in}{-0.020833in}}%
\pgfusepath{stroke,fill}%
}%
\begin{pgfscope}%
\pgfsys@transformshift{1.998210in}{2.414488in}%
\pgfsys@useobject{currentmarker}{}%
\end{pgfscope}%
\end{pgfscope}%
\begin{pgfscope}%
\pgfsetbuttcap%
\pgfsetroundjoin%
\definecolor{currentfill}{rgb}{0.000000,0.000000,0.000000}%
\pgfsetfillcolor{currentfill}%
\pgfsetlinewidth{0.401500pt}%
\definecolor{currentstroke}{rgb}{0.000000,0.000000,0.000000}%
\pgfsetstrokecolor{currentstroke}%
\pgfsetdash{}{0pt}%
\pgfsys@defobject{currentmarker}{\pgfqpoint{0.000000in}{0.000000in}}{\pgfqpoint{0.000000in}{0.020833in}}{%
\pgfpathmoveto{\pgfqpoint{0.000000in}{0.000000in}}%
\pgfpathlineto{\pgfqpoint{0.000000in}{0.020833in}}%
\pgfusepath{stroke,fill}%
}%
\begin{pgfscope}%
\pgfsys@transformshift{2.166737in}{0.420000in}%
\pgfsys@useobject{currentmarker}{}%
\end{pgfscope}%
\end{pgfscope}%
\begin{pgfscope}%
\pgfsetbuttcap%
\pgfsetroundjoin%
\definecolor{currentfill}{rgb}{0.000000,0.000000,0.000000}%
\pgfsetfillcolor{currentfill}%
\pgfsetlinewidth{0.401500pt}%
\definecolor{currentstroke}{rgb}{0.000000,0.000000,0.000000}%
\pgfsetstrokecolor{currentstroke}%
\pgfsetdash{}{0pt}%
\pgfsys@defobject{currentmarker}{\pgfqpoint{0.000000in}{-0.020833in}}{\pgfqpoint{0.000000in}{0.000000in}}{%
\pgfpathmoveto{\pgfqpoint{0.000000in}{0.000000in}}%
\pgfpathlineto{\pgfqpoint{0.000000in}{-0.020833in}}%
\pgfusepath{stroke,fill}%
}%
\begin{pgfscope}%
\pgfsys@transformshift{2.166737in}{2.414488in}%
\pgfsys@useobject{currentmarker}{}%
\end{pgfscope}%
\end{pgfscope}%
\begin{pgfscope}%
\pgfsetbuttcap%
\pgfsetroundjoin%
\definecolor{currentfill}{rgb}{0.000000,0.000000,0.000000}%
\pgfsetfillcolor{currentfill}%
\pgfsetlinewidth{0.401500pt}%
\definecolor{currentstroke}{rgb}{0.000000,0.000000,0.000000}%
\pgfsetstrokecolor{currentstroke}%
\pgfsetdash{}{0pt}%
\pgfsys@defobject{currentmarker}{\pgfqpoint{0.000000in}{0.000000in}}{\pgfqpoint{0.000000in}{0.020833in}}{%
\pgfpathmoveto{\pgfqpoint{0.000000in}{0.000000in}}%
\pgfpathlineto{\pgfqpoint{0.000000in}{0.020833in}}%
\pgfusepath{stroke,fill}%
}%
\begin{pgfscope}%
\pgfsys@transformshift{2.503789in}{0.420000in}%
\pgfsys@useobject{currentmarker}{}%
\end{pgfscope}%
\end{pgfscope}%
\begin{pgfscope}%
\pgfsetbuttcap%
\pgfsetroundjoin%
\definecolor{currentfill}{rgb}{0.000000,0.000000,0.000000}%
\pgfsetfillcolor{currentfill}%
\pgfsetlinewidth{0.401500pt}%
\definecolor{currentstroke}{rgb}{0.000000,0.000000,0.000000}%
\pgfsetstrokecolor{currentstroke}%
\pgfsetdash{}{0pt}%
\pgfsys@defobject{currentmarker}{\pgfqpoint{0.000000in}{-0.020833in}}{\pgfqpoint{0.000000in}{0.000000in}}{%
\pgfpathmoveto{\pgfqpoint{0.000000in}{0.000000in}}%
\pgfpathlineto{\pgfqpoint{0.000000in}{-0.020833in}}%
\pgfusepath{stroke,fill}%
}%
\begin{pgfscope}%
\pgfsys@transformshift{2.503789in}{2.414488in}%
\pgfsys@useobject{currentmarker}{}%
\end{pgfscope}%
\end{pgfscope}%
\begin{pgfscope}%
\pgfsetbuttcap%
\pgfsetroundjoin%
\definecolor{currentfill}{rgb}{0.000000,0.000000,0.000000}%
\pgfsetfillcolor{currentfill}%
\pgfsetlinewidth{0.401500pt}%
\definecolor{currentstroke}{rgb}{0.000000,0.000000,0.000000}%
\pgfsetstrokecolor{currentstroke}%
\pgfsetdash{}{0pt}%
\pgfsys@defobject{currentmarker}{\pgfqpoint{0.000000in}{0.000000in}}{\pgfqpoint{0.000000in}{0.020833in}}{%
\pgfpathmoveto{\pgfqpoint{0.000000in}{0.000000in}}%
\pgfpathlineto{\pgfqpoint{0.000000in}{0.020833in}}%
\pgfusepath{stroke,fill}%
}%
\begin{pgfscope}%
\pgfsys@transformshift{2.672315in}{0.420000in}%
\pgfsys@useobject{currentmarker}{}%
\end{pgfscope}%
\end{pgfscope}%
\begin{pgfscope}%
\pgfsetbuttcap%
\pgfsetroundjoin%
\definecolor{currentfill}{rgb}{0.000000,0.000000,0.000000}%
\pgfsetfillcolor{currentfill}%
\pgfsetlinewidth{0.401500pt}%
\definecolor{currentstroke}{rgb}{0.000000,0.000000,0.000000}%
\pgfsetstrokecolor{currentstroke}%
\pgfsetdash{}{0pt}%
\pgfsys@defobject{currentmarker}{\pgfqpoint{0.000000in}{-0.020833in}}{\pgfqpoint{0.000000in}{0.000000in}}{%
\pgfpathmoveto{\pgfqpoint{0.000000in}{0.000000in}}%
\pgfpathlineto{\pgfqpoint{0.000000in}{-0.020833in}}%
\pgfusepath{stroke,fill}%
}%
\begin{pgfscope}%
\pgfsys@transformshift{2.672315in}{2.414488in}%
\pgfsys@useobject{currentmarker}{}%
\end{pgfscope}%
\end{pgfscope}%
\begin{pgfscope}%
\pgfsetbuttcap%
\pgfsetroundjoin%
\definecolor{currentfill}{rgb}{0.000000,0.000000,0.000000}%
\pgfsetfillcolor{currentfill}%
\pgfsetlinewidth{0.401500pt}%
\definecolor{currentstroke}{rgb}{0.000000,0.000000,0.000000}%
\pgfsetstrokecolor{currentstroke}%
\pgfsetdash{}{0pt}%
\pgfsys@defobject{currentmarker}{\pgfqpoint{0.000000in}{0.000000in}}{\pgfqpoint{0.000000in}{0.020833in}}{%
\pgfpathmoveto{\pgfqpoint{0.000000in}{0.000000in}}%
\pgfpathlineto{\pgfqpoint{0.000000in}{0.020833in}}%
\pgfusepath{stroke,fill}%
}%
\begin{pgfscope}%
\pgfsys@transformshift{2.840842in}{0.420000in}%
\pgfsys@useobject{currentmarker}{}%
\end{pgfscope}%
\end{pgfscope}%
\begin{pgfscope}%
\pgfsetbuttcap%
\pgfsetroundjoin%
\definecolor{currentfill}{rgb}{0.000000,0.000000,0.000000}%
\pgfsetfillcolor{currentfill}%
\pgfsetlinewidth{0.401500pt}%
\definecolor{currentstroke}{rgb}{0.000000,0.000000,0.000000}%
\pgfsetstrokecolor{currentstroke}%
\pgfsetdash{}{0pt}%
\pgfsys@defobject{currentmarker}{\pgfqpoint{0.000000in}{-0.020833in}}{\pgfqpoint{0.000000in}{0.000000in}}{%
\pgfpathmoveto{\pgfqpoint{0.000000in}{0.000000in}}%
\pgfpathlineto{\pgfqpoint{0.000000in}{-0.020833in}}%
\pgfusepath{stroke,fill}%
}%
\begin{pgfscope}%
\pgfsys@transformshift{2.840842in}{2.414488in}%
\pgfsys@useobject{currentmarker}{}%
\end{pgfscope}%
\end{pgfscope}%
\begin{pgfscope}%
\pgfsetbuttcap%
\pgfsetroundjoin%
\definecolor{currentfill}{rgb}{0.000000,0.000000,0.000000}%
\pgfsetfillcolor{currentfill}%
\pgfsetlinewidth{0.401500pt}%
\definecolor{currentstroke}{rgb}{0.000000,0.000000,0.000000}%
\pgfsetstrokecolor{currentstroke}%
\pgfsetdash{}{0pt}%
\pgfsys@defobject{currentmarker}{\pgfqpoint{0.000000in}{0.000000in}}{\pgfqpoint{0.000000in}{0.020833in}}{%
\pgfpathmoveto{\pgfqpoint{0.000000in}{0.000000in}}%
\pgfpathlineto{\pgfqpoint{0.000000in}{0.020833in}}%
\pgfusepath{stroke,fill}%
}%
\begin{pgfscope}%
\pgfsys@transformshift{3.009368in}{0.420000in}%
\pgfsys@useobject{currentmarker}{}%
\end{pgfscope}%
\end{pgfscope}%
\begin{pgfscope}%
\pgfsetbuttcap%
\pgfsetroundjoin%
\definecolor{currentfill}{rgb}{0.000000,0.000000,0.000000}%
\pgfsetfillcolor{currentfill}%
\pgfsetlinewidth{0.401500pt}%
\definecolor{currentstroke}{rgb}{0.000000,0.000000,0.000000}%
\pgfsetstrokecolor{currentstroke}%
\pgfsetdash{}{0pt}%
\pgfsys@defobject{currentmarker}{\pgfqpoint{0.000000in}{-0.020833in}}{\pgfqpoint{0.000000in}{0.000000in}}{%
\pgfpathmoveto{\pgfqpoint{0.000000in}{0.000000in}}%
\pgfpathlineto{\pgfqpoint{0.000000in}{-0.020833in}}%
\pgfusepath{stroke,fill}%
}%
\begin{pgfscope}%
\pgfsys@transformshift{3.009368in}{2.414488in}%
\pgfsys@useobject{currentmarker}{}%
\end{pgfscope}%
\end{pgfscope}%
\begin{pgfscope}%
\definecolor{textcolor}{rgb}{0.000000,0.000000,0.000000}%
\pgfsetstrokecolor{textcolor}%
\pgfsetfillcolor{textcolor}%
\pgftext[x=1.844055in,y=0.208426in,,top]{\color{textcolor}{\rmfamily\fontsize{12.000000}{14.400000}\selectfont\catcode`\^=\active\def^{\ifmmode\sp\else\^{}\fi}\catcode`\%=\active\def%{\%}$s/\delta_{0}$}}%
\end{pgfscope}%
\begin{pgfscope}%
\pgfsetbuttcap%
\pgfsetroundjoin%
\definecolor{currentfill}{rgb}{0.000000,0.000000,0.000000}%
\pgfsetfillcolor{currentfill}%
\pgfsetlinewidth{0.501875pt}%
\definecolor{currentstroke}{rgb}{0.000000,0.000000,0.000000}%
\pgfsetstrokecolor{currentstroke}%
\pgfsetdash{}{0pt}%
\pgfsys@defobject{currentmarker}{\pgfqpoint{0.000000in}{0.000000in}}{\pgfqpoint{0.041667in}{0.000000in}}{%
\pgfpathmoveto{\pgfqpoint{0.000000in}{0.000000in}}%
\pgfpathlineto{\pgfqpoint{0.041667in}{0.000000in}}%
\pgfusepath{stroke,fill}%
}%
\begin{pgfscope}%
\pgfsys@transformshift{0.650000in}{0.490588in}%
\pgfsys@useobject{currentmarker}{}%
\end{pgfscope}%
\end{pgfscope}%
\begin{pgfscope}%
\pgfsetbuttcap%
\pgfsetroundjoin%
\definecolor{currentfill}{rgb}{0.000000,0.000000,0.000000}%
\pgfsetfillcolor{currentfill}%
\pgfsetlinewidth{0.501875pt}%
\definecolor{currentstroke}{rgb}{0.000000,0.000000,0.000000}%
\pgfsetstrokecolor{currentstroke}%
\pgfsetdash{}{0pt}%
\pgfsys@defobject{currentmarker}{\pgfqpoint{-0.041667in}{0.000000in}}{\pgfqpoint{-0.000000in}{0.000000in}}{%
\pgfpathmoveto{\pgfqpoint{-0.000000in}{0.000000in}}%
\pgfpathlineto{\pgfqpoint{-0.041667in}{0.000000in}}%
\pgfusepath{stroke,fill}%
}%
\begin{pgfscope}%
\pgfsys@transformshift{3.038110in}{0.490588in}%
\pgfsys@useobject{currentmarker}{}%
\end{pgfscope}%
\end{pgfscope}%
\begin{pgfscope}%
\definecolor{textcolor}{rgb}{0.000000,0.000000,0.000000}%
\pgfsetstrokecolor{textcolor}%
\pgfsetfillcolor{textcolor}%
\pgftext[x=0.407060in, y=0.437782in, left, base]{\color{textcolor}{\rmfamily\fontsize{11.000000}{13.200000}\selectfont\catcode`\^=\active\def^{\ifmmode\sp\else\^{}\fi}\catcode`\%=\active\def%{\%}0.0}}%
\end{pgfscope}%
\begin{pgfscope}%
\pgfsetbuttcap%
\pgfsetroundjoin%
\definecolor{currentfill}{rgb}{0.000000,0.000000,0.000000}%
\pgfsetfillcolor{currentfill}%
\pgfsetlinewidth{0.501875pt}%
\definecolor{currentstroke}{rgb}{0.000000,0.000000,0.000000}%
\pgfsetstrokecolor{currentstroke}%
\pgfsetdash{}{0pt}%
\pgfsys@defobject{currentmarker}{\pgfqpoint{0.000000in}{0.000000in}}{\pgfqpoint{0.041667in}{0.000000in}}{%
\pgfpathmoveto{\pgfqpoint{0.000000in}{0.000000in}}%
\pgfpathlineto{\pgfqpoint{0.041667in}{0.000000in}}%
\pgfusepath{stroke,fill}%
}%
\begin{pgfscope}%
\pgfsys@transformshift{0.650000in}{0.905684in}%
\pgfsys@useobject{currentmarker}{}%
\end{pgfscope}%
\end{pgfscope}%
\begin{pgfscope}%
\pgfsetbuttcap%
\pgfsetroundjoin%
\definecolor{currentfill}{rgb}{0.000000,0.000000,0.000000}%
\pgfsetfillcolor{currentfill}%
\pgfsetlinewidth{0.501875pt}%
\definecolor{currentstroke}{rgb}{0.000000,0.000000,0.000000}%
\pgfsetstrokecolor{currentstroke}%
\pgfsetdash{}{0pt}%
\pgfsys@defobject{currentmarker}{\pgfqpoint{-0.041667in}{0.000000in}}{\pgfqpoint{-0.000000in}{0.000000in}}{%
\pgfpathmoveto{\pgfqpoint{-0.000000in}{0.000000in}}%
\pgfpathlineto{\pgfqpoint{-0.041667in}{0.000000in}}%
\pgfusepath{stroke,fill}%
}%
\begin{pgfscope}%
\pgfsys@transformshift{3.038110in}{0.905684in}%
\pgfsys@useobject{currentmarker}{}%
\end{pgfscope}%
\end{pgfscope}%
\begin{pgfscope}%
\definecolor{textcolor}{rgb}{0.000000,0.000000,0.000000}%
\pgfsetstrokecolor{textcolor}%
\pgfsetfillcolor{textcolor}%
\pgftext[x=0.407060in, y=0.852878in, left, base]{\color{textcolor}{\rmfamily\fontsize{11.000000}{13.200000}\selectfont\catcode`\^=\active\def^{\ifmmode\sp\else\^{}\fi}\catcode`\%=\active\def%{\%}0.5}}%
\end{pgfscope}%
\begin{pgfscope}%
\pgfsetbuttcap%
\pgfsetroundjoin%
\definecolor{currentfill}{rgb}{0.000000,0.000000,0.000000}%
\pgfsetfillcolor{currentfill}%
\pgfsetlinewidth{0.501875pt}%
\definecolor{currentstroke}{rgb}{0.000000,0.000000,0.000000}%
\pgfsetstrokecolor{currentstroke}%
\pgfsetdash{}{0pt}%
\pgfsys@defobject{currentmarker}{\pgfqpoint{0.000000in}{0.000000in}}{\pgfqpoint{0.041667in}{0.000000in}}{%
\pgfpathmoveto{\pgfqpoint{0.000000in}{0.000000in}}%
\pgfpathlineto{\pgfqpoint{0.041667in}{0.000000in}}%
\pgfusepath{stroke,fill}%
}%
\begin{pgfscope}%
\pgfsys@transformshift{0.650000in}{1.320781in}%
\pgfsys@useobject{currentmarker}{}%
\end{pgfscope}%
\end{pgfscope}%
\begin{pgfscope}%
\pgfsetbuttcap%
\pgfsetroundjoin%
\definecolor{currentfill}{rgb}{0.000000,0.000000,0.000000}%
\pgfsetfillcolor{currentfill}%
\pgfsetlinewidth{0.501875pt}%
\definecolor{currentstroke}{rgb}{0.000000,0.000000,0.000000}%
\pgfsetstrokecolor{currentstroke}%
\pgfsetdash{}{0pt}%
\pgfsys@defobject{currentmarker}{\pgfqpoint{-0.041667in}{0.000000in}}{\pgfqpoint{-0.000000in}{0.000000in}}{%
\pgfpathmoveto{\pgfqpoint{-0.000000in}{0.000000in}}%
\pgfpathlineto{\pgfqpoint{-0.041667in}{0.000000in}}%
\pgfusepath{stroke,fill}%
}%
\begin{pgfscope}%
\pgfsys@transformshift{3.038110in}{1.320781in}%
\pgfsys@useobject{currentmarker}{}%
\end{pgfscope}%
\end{pgfscope}%
\begin{pgfscope}%
\definecolor{textcolor}{rgb}{0.000000,0.000000,0.000000}%
\pgfsetstrokecolor{textcolor}%
\pgfsetfillcolor{textcolor}%
\pgftext[x=0.407060in, y=1.267974in, left, base]{\color{textcolor}{\rmfamily\fontsize{11.000000}{13.200000}\selectfont\catcode`\^=\active\def^{\ifmmode\sp\else\^{}\fi}\catcode`\%=\active\def%{\%}1.0}}%
\end{pgfscope}%
\begin{pgfscope}%
\pgfsetbuttcap%
\pgfsetroundjoin%
\definecolor{currentfill}{rgb}{0.000000,0.000000,0.000000}%
\pgfsetfillcolor{currentfill}%
\pgfsetlinewidth{0.501875pt}%
\definecolor{currentstroke}{rgb}{0.000000,0.000000,0.000000}%
\pgfsetstrokecolor{currentstroke}%
\pgfsetdash{}{0pt}%
\pgfsys@defobject{currentmarker}{\pgfqpoint{0.000000in}{0.000000in}}{\pgfqpoint{0.041667in}{0.000000in}}{%
\pgfpathmoveto{\pgfqpoint{0.000000in}{0.000000in}}%
\pgfpathlineto{\pgfqpoint{0.041667in}{0.000000in}}%
\pgfusepath{stroke,fill}%
}%
\begin{pgfscope}%
\pgfsys@transformshift{0.650000in}{1.735877in}%
\pgfsys@useobject{currentmarker}{}%
\end{pgfscope}%
\end{pgfscope}%
\begin{pgfscope}%
\pgfsetbuttcap%
\pgfsetroundjoin%
\definecolor{currentfill}{rgb}{0.000000,0.000000,0.000000}%
\pgfsetfillcolor{currentfill}%
\pgfsetlinewidth{0.501875pt}%
\definecolor{currentstroke}{rgb}{0.000000,0.000000,0.000000}%
\pgfsetstrokecolor{currentstroke}%
\pgfsetdash{}{0pt}%
\pgfsys@defobject{currentmarker}{\pgfqpoint{-0.041667in}{0.000000in}}{\pgfqpoint{-0.000000in}{0.000000in}}{%
\pgfpathmoveto{\pgfqpoint{-0.000000in}{0.000000in}}%
\pgfpathlineto{\pgfqpoint{-0.041667in}{0.000000in}}%
\pgfusepath{stroke,fill}%
}%
\begin{pgfscope}%
\pgfsys@transformshift{3.038110in}{1.735877in}%
\pgfsys@useobject{currentmarker}{}%
\end{pgfscope}%
\end{pgfscope}%
\begin{pgfscope}%
\definecolor{textcolor}{rgb}{0.000000,0.000000,0.000000}%
\pgfsetstrokecolor{textcolor}%
\pgfsetfillcolor{textcolor}%
\pgftext[x=0.407060in, y=1.683070in, left, base]{\color{textcolor}{\rmfamily\fontsize{11.000000}{13.200000}\selectfont\catcode`\^=\active\def^{\ifmmode\sp\else\^{}\fi}\catcode`\%=\active\def%{\%}1.5}}%
\end{pgfscope}%
\begin{pgfscope}%
\pgfsetbuttcap%
\pgfsetroundjoin%
\definecolor{currentfill}{rgb}{0.000000,0.000000,0.000000}%
\pgfsetfillcolor{currentfill}%
\pgfsetlinewidth{0.501875pt}%
\definecolor{currentstroke}{rgb}{0.000000,0.000000,0.000000}%
\pgfsetstrokecolor{currentstroke}%
\pgfsetdash{}{0pt}%
\pgfsys@defobject{currentmarker}{\pgfqpoint{0.000000in}{0.000000in}}{\pgfqpoint{0.041667in}{0.000000in}}{%
\pgfpathmoveto{\pgfqpoint{0.000000in}{0.000000in}}%
\pgfpathlineto{\pgfqpoint{0.041667in}{0.000000in}}%
\pgfusepath{stroke,fill}%
}%
\begin{pgfscope}%
\pgfsys@transformshift{0.650000in}{2.150973in}%
\pgfsys@useobject{currentmarker}{}%
\end{pgfscope}%
\end{pgfscope}%
\begin{pgfscope}%
\pgfsetbuttcap%
\pgfsetroundjoin%
\definecolor{currentfill}{rgb}{0.000000,0.000000,0.000000}%
\pgfsetfillcolor{currentfill}%
\pgfsetlinewidth{0.501875pt}%
\definecolor{currentstroke}{rgb}{0.000000,0.000000,0.000000}%
\pgfsetstrokecolor{currentstroke}%
\pgfsetdash{}{0pt}%
\pgfsys@defobject{currentmarker}{\pgfqpoint{-0.041667in}{0.000000in}}{\pgfqpoint{-0.000000in}{0.000000in}}{%
\pgfpathmoveto{\pgfqpoint{-0.000000in}{0.000000in}}%
\pgfpathlineto{\pgfqpoint{-0.041667in}{0.000000in}}%
\pgfusepath{stroke,fill}%
}%
\begin{pgfscope}%
\pgfsys@transformshift{3.038110in}{2.150973in}%
\pgfsys@useobject{currentmarker}{}%
\end{pgfscope}%
\end{pgfscope}%
\begin{pgfscope}%
\definecolor{textcolor}{rgb}{0.000000,0.000000,0.000000}%
\pgfsetstrokecolor{textcolor}%
\pgfsetfillcolor{textcolor}%
\pgftext[x=0.407060in, y=2.098166in, left, base]{\color{textcolor}{\rmfamily\fontsize{11.000000}{13.200000}\selectfont\catcode`\^=\active\def^{\ifmmode\sp\else\^{}\fi}\catcode`\%=\active\def%{\%}2.0}}%
\end{pgfscope}%
\begin{pgfscope}%
\pgfsetbuttcap%
\pgfsetroundjoin%
\definecolor{currentfill}{rgb}{0.000000,0.000000,0.000000}%
\pgfsetfillcolor{currentfill}%
\pgfsetlinewidth{0.401500pt}%
\definecolor{currentstroke}{rgb}{0.000000,0.000000,0.000000}%
\pgfsetstrokecolor{currentstroke}%
\pgfsetdash{}{0pt}%
\pgfsys@defobject{currentmarker}{\pgfqpoint{0.000000in}{0.000000in}}{\pgfqpoint{0.020833in}{0.000000in}}{%
\pgfpathmoveto{\pgfqpoint{0.000000in}{0.000000in}}%
\pgfpathlineto{\pgfqpoint{0.020833in}{0.000000in}}%
\pgfusepath{stroke,fill}%
}%
\begin{pgfscope}%
\pgfsys@transformshift{0.650000in}{0.573607in}%
\pgfsys@useobject{currentmarker}{}%
\end{pgfscope}%
\end{pgfscope}%
\begin{pgfscope}%
\pgfsetbuttcap%
\pgfsetroundjoin%
\definecolor{currentfill}{rgb}{0.000000,0.000000,0.000000}%
\pgfsetfillcolor{currentfill}%
\pgfsetlinewidth{0.401500pt}%
\definecolor{currentstroke}{rgb}{0.000000,0.000000,0.000000}%
\pgfsetstrokecolor{currentstroke}%
\pgfsetdash{}{0pt}%
\pgfsys@defobject{currentmarker}{\pgfqpoint{-0.020833in}{0.000000in}}{\pgfqpoint{-0.000000in}{0.000000in}}{%
\pgfpathmoveto{\pgfqpoint{-0.000000in}{0.000000in}}%
\pgfpathlineto{\pgfqpoint{-0.020833in}{0.000000in}}%
\pgfusepath{stroke,fill}%
}%
\begin{pgfscope}%
\pgfsys@transformshift{3.038110in}{0.573607in}%
\pgfsys@useobject{currentmarker}{}%
\end{pgfscope}%
\end{pgfscope}%
\begin{pgfscope}%
\pgfsetbuttcap%
\pgfsetroundjoin%
\definecolor{currentfill}{rgb}{0.000000,0.000000,0.000000}%
\pgfsetfillcolor{currentfill}%
\pgfsetlinewidth{0.401500pt}%
\definecolor{currentstroke}{rgb}{0.000000,0.000000,0.000000}%
\pgfsetstrokecolor{currentstroke}%
\pgfsetdash{}{0pt}%
\pgfsys@defobject{currentmarker}{\pgfqpoint{0.000000in}{0.000000in}}{\pgfqpoint{0.020833in}{0.000000in}}{%
\pgfpathmoveto{\pgfqpoint{0.000000in}{0.000000in}}%
\pgfpathlineto{\pgfqpoint{0.020833in}{0.000000in}}%
\pgfusepath{stroke,fill}%
}%
\begin{pgfscope}%
\pgfsys@transformshift{0.650000in}{0.656627in}%
\pgfsys@useobject{currentmarker}{}%
\end{pgfscope}%
\end{pgfscope}%
\begin{pgfscope}%
\pgfsetbuttcap%
\pgfsetroundjoin%
\definecolor{currentfill}{rgb}{0.000000,0.000000,0.000000}%
\pgfsetfillcolor{currentfill}%
\pgfsetlinewidth{0.401500pt}%
\definecolor{currentstroke}{rgb}{0.000000,0.000000,0.000000}%
\pgfsetstrokecolor{currentstroke}%
\pgfsetdash{}{0pt}%
\pgfsys@defobject{currentmarker}{\pgfqpoint{-0.020833in}{0.000000in}}{\pgfqpoint{-0.000000in}{0.000000in}}{%
\pgfpathmoveto{\pgfqpoint{-0.000000in}{0.000000in}}%
\pgfpathlineto{\pgfqpoint{-0.020833in}{0.000000in}}%
\pgfusepath{stroke,fill}%
}%
\begin{pgfscope}%
\pgfsys@transformshift{3.038110in}{0.656627in}%
\pgfsys@useobject{currentmarker}{}%
\end{pgfscope}%
\end{pgfscope}%
\begin{pgfscope}%
\pgfsetbuttcap%
\pgfsetroundjoin%
\definecolor{currentfill}{rgb}{0.000000,0.000000,0.000000}%
\pgfsetfillcolor{currentfill}%
\pgfsetlinewidth{0.401500pt}%
\definecolor{currentstroke}{rgb}{0.000000,0.000000,0.000000}%
\pgfsetstrokecolor{currentstroke}%
\pgfsetdash{}{0pt}%
\pgfsys@defobject{currentmarker}{\pgfqpoint{0.000000in}{0.000000in}}{\pgfqpoint{0.020833in}{0.000000in}}{%
\pgfpathmoveto{\pgfqpoint{0.000000in}{0.000000in}}%
\pgfpathlineto{\pgfqpoint{0.020833in}{0.000000in}}%
\pgfusepath{stroke,fill}%
}%
\begin{pgfscope}%
\pgfsys@transformshift{0.650000in}{0.739646in}%
\pgfsys@useobject{currentmarker}{}%
\end{pgfscope}%
\end{pgfscope}%
\begin{pgfscope}%
\pgfsetbuttcap%
\pgfsetroundjoin%
\definecolor{currentfill}{rgb}{0.000000,0.000000,0.000000}%
\pgfsetfillcolor{currentfill}%
\pgfsetlinewidth{0.401500pt}%
\definecolor{currentstroke}{rgb}{0.000000,0.000000,0.000000}%
\pgfsetstrokecolor{currentstroke}%
\pgfsetdash{}{0pt}%
\pgfsys@defobject{currentmarker}{\pgfqpoint{-0.020833in}{0.000000in}}{\pgfqpoint{-0.000000in}{0.000000in}}{%
\pgfpathmoveto{\pgfqpoint{-0.000000in}{0.000000in}}%
\pgfpathlineto{\pgfqpoint{-0.020833in}{0.000000in}}%
\pgfusepath{stroke,fill}%
}%
\begin{pgfscope}%
\pgfsys@transformshift{3.038110in}{0.739646in}%
\pgfsys@useobject{currentmarker}{}%
\end{pgfscope}%
\end{pgfscope}%
\begin{pgfscope}%
\pgfsetbuttcap%
\pgfsetroundjoin%
\definecolor{currentfill}{rgb}{0.000000,0.000000,0.000000}%
\pgfsetfillcolor{currentfill}%
\pgfsetlinewidth{0.401500pt}%
\definecolor{currentstroke}{rgb}{0.000000,0.000000,0.000000}%
\pgfsetstrokecolor{currentstroke}%
\pgfsetdash{}{0pt}%
\pgfsys@defobject{currentmarker}{\pgfqpoint{0.000000in}{0.000000in}}{\pgfqpoint{0.020833in}{0.000000in}}{%
\pgfpathmoveto{\pgfqpoint{0.000000in}{0.000000in}}%
\pgfpathlineto{\pgfqpoint{0.020833in}{0.000000in}}%
\pgfusepath{stroke,fill}%
}%
\begin{pgfscope}%
\pgfsys@transformshift{0.650000in}{0.822665in}%
\pgfsys@useobject{currentmarker}{}%
\end{pgfscope}%
\end{pgfscope}%
\begin{pgfscope}%
\pgfsetbuttcap%
\pgfsetroundjoin%
\definecolor{currentfill}{rgb}{0.000000,0.000000,0.000000}%
\pgfsetfillcolor{currentfill}%
\pgfsetlinewidth{0.401500pt}%
\definecolor{currentstroke}{rgb}{0.000000,0.000000,0.000000}%
\pgfsetstrokecolor{currentstroke}%
\pgfsetdash{}{0pt}%
\pgfsys@defobject{currentmarker}{\pgfqpoint{-0.020833in}{0.000000in}}{\pgfqpoint{-0.000000in}{0.000000in}}{%
\pgfpathmoveto{\pgfqpoint{-0.000000in}{0.000000in}}%
\pgfpathlineto{\pgfqpoint{-0.020833in}{0.000000in}}%
\pgfusepath{stroke,fill}%
}%
\begin{pgfscope}%
\pgfsys@transformshift{3.038110in}{0.822665in}%
\pgfsys@useobject{currentmarker}{}%
\end{pgfscope}%
\end{pgfscope}%
\begin{pgfscope}%
\pgfsetbuttcap%
\pgfsetroundjoin%
\definecolor{currentfill}{rgb}{0.000000,0.000000,0.000000}%
\pgfsetfillcolor{currentfill}%
\pgfsetlinewidth{0.401500pt}%
\definecolor{currentstroke}{rgb}{0.000000,0.000000,0.000000}%
\pgfsetstrokecolor{currentstroke}%
\pgfsetdash{}{0pt}%
\pgfsys@defobject{currentmarker}{\pgfqpoint{0.000000in}{0.000000in}}{\pgfqpoint{0.020833in}{0.000000in}}{%
\pgfpathmoveto{\pgfqpoint{0.000000in}{0.000000in}}%
\pgfpathlineto{\pgfqpoint{0.020833in}{0.000000in}}%
\pgfusepath{stroke,fill}%
}%
\begin{pgfscope}%
\pgfsys@transformshift{0.650000in}{0.988704in}%
\pgfsys@useobject{currentmarker}{}%
\end{pgfscope}%
\end{pgfscope}%
\begin{pgfscope}%
\pgfsetbuttcap%
\pgfsetroundjoin%
\definecolor{currentfill}{rgb}{0.000000,0.000000,0.000000}%
\pgfsetfillcolor{currentfill}%
\pgfsetlinewidth{0.401500pt}%
\definecolor{currentstroke}{rgb}{0.000000,0.000000,0.000000}%
\pgfsetstrokecolor{currentstroke}%
\pgfsetdash{}{0pt}%
\pgfsys@defobject{currentmarker}{\pgfqpoint{-0.020833in}{0.000000in}}{\pgfqpoint{-0.000000in}{0.000000in}}{%
\pgfpathmoveto{\pgfqpoint{-0.000000in}{0.000000in}}%
\pgfpathlineto{\pgfqpoint{-0.020833in}{0.000000in}}%
\pgfusepath{stroke,fill}%
}%
\begin{pgfscope}%
\pgfsys@transformshift{3.038110in}{0.988704in}%
\pgfsys@useobject{currentmarker}{}%
\end{pgfscope}%
\end{pgfscope}%
\begin{pgfscope}%
\pgfsetbuttcap%
\pgfsetroundjoin%
\definecolor{currentfill}{rgb}{0.000000,0.000000,0.000000}%
\pgfsetfillcolor{currentfill}%
\pgfsetlinewidth{0.401500pt}%
\definecolor{currentstroke}{rgb}{0.000000,0.000000,0.000000}%
\pgfsetstrokecolor{currentstroke}%
\pgfsetdash{}{0pt}%
\pgfsys@defobject{currentmarker}{\pgfqpoint{0.000000in}{0.000000in}}{\pgfqpoint{0.020833in}{0.000000in}}{%
\pgfpathmoveto{\pgfqpoint{0.000000in}{0.000000in}}%
\pgfpathlineto{\pgfqpoint{0.020833in}{0.000000in}}%
\pgfusepath{stroke,fill}%
}%
\begin{pgfscope}%
\pgfsys@transformshift{0.650000in}{1.071723in}%
\pgfsys@useobject{currentmarker}{}%
\end{pgfscope}%
\end{pgfscope}%
\begin{pgfscope}%
\pgfsetbuttcap%
\pgfsetroundjoin%
\definecolor{currentfill}{rgb}{0.000000,0.000000,0.000000}%
\pgfsetfillcolor{currentfill}%
\pgfsetlinewidth{0.401500pt}%
\definecolor{currentstroke}{rgb}{0.000000,0.000000,0.000000}%
\pgfsetstrokecolor{currentstroke}%
\pgfsetdash{}{0pt}%
\pgfsys@defobject{currentmarker}{\pgfqpoint{-0.020833in}{0.000000in}}{\pgfqpoint{-0.000000in}{0.000000in}}{%
\pgfpathmoveto{\pgfqpoint{-0.000000in}{0.000000in}}%
\pgfpathlineto{\pgfqpoint{-0.020833in}{0.000000in}}%
\pgfusepath{stroke,fill}%
}%
\begin{pgfscope}%
\pgfsys@transformshift{3.038110in}{1.071723in}%
\pgfsys@useobject{currentmarker}{}%
\end{pgfscope}%
\end{pgfscope}%
\begin{pgfscope}%
\pgfsetbuttcap%
\pgfsetroundjoin%
\definecolor{currentfill}{rgb}{0.000000,0.000000,0.000000}%
\pgfsetfillcolor{currentfill}%
\pgfsetlinewidth{0.401500pt}%
\definecolor{currentstroke}{rgb}{0.000000,0.000000,0.000000}%
\pgfsetstrokecolor{currentstroke}%
\pgfsetdash{}{0pt}%
\pgfsys@defobject{currentmarker}{\pgfqpoint{0.000000in}{0.000000in}}{\pgfqpoint{0.020833in}{0.000000in}}{%
\pgfpathmoveto{\pgfqpoint{0.000000in}{0.000000in}}%
\pgfpathlineto{\pgfqpoint{0.020833in}{0.000000in}}%
\pgfusepath{stroke,fill}%
}%
\begin{pgfscope}%
\pgfsys@transformshift{0.650000in}{1.154742in}%
\pgfsys@useobject{currentmarker}{}%
\end{pgfscope}%
\end{pgfscope}%
\begin{pgfscope}%
\pgfsetbuttcap%
\pgfsetroundjoin%
\definecolor{currentfill}{rgb}{0.000000,0.000000,0.000000}%
\pgfsetfillcolor{currentfill}%
\pgfsetlinewidth{0.401500pt}%
\definecolor{currentstroke}{rgb}{0.000000,0.000000,0.000000}%
\pgfsetstrokecolor{currentstroke}%
\pgfsetdash{}{0pt}%
\pgfsys@defobject{currentmarker}{\pgfqpoint{-0.020833in}{0.000000in}}{\pgfqpoint{-0.000000in}{0.000000in}}{%
\pgfpathmoveto{\pgfqpoint{-0.000000in}{0.000000in}}%
\pgfpathlineto{\pgfqpoint{-0.020833in}{0.000000in}}%
\pgfusepath{stroke,fill}%
}%
\begin{pgfscope}%
\pgfsys@transformshift{3.038110in}{1.154742in}%
\pgfsys@useobject{currentmarker}{}%
\end{pgfscope}%
\end{pgfscope}%
\begin{pgfscope}%
\pgfsetbuttcap%
\pgfsetroundjoin%
\definecolor{currentfill}{rgb}{0.000000,0.000000,0.000000}%
\pgfsetfillcolor{currentfill}%
\pgfsetlinewidth{0.401500pt}%
\definecolor{currentstroke}{rgb}{0.000000,0.000000,0.000000}%
\pgfsetstrokecolor{currentstroke}%
\pgfsetdash{}{0pt}%
\pgfsys@defobject{currentmarker}{\pgfqpoint{0.000000in}{0.000000in}}{\pgfqpoint{0.020833in}{0.000000in}}{%
\pgfpathmoveto{\pgfqpoint{0.000000in}{0.000000in}}%
\pgfpathlineto{\pgfqpoint{0.020833in}{0.000000in}}%
\pgfusepath{stroke,fill}%
}%
\begin{pgfscope}%
\pgfsys@transformshift{0.650000in}{1.237761in}%
\pgfsys@useobject{currentmarker}{}%
\end{pgfscope}%
\end{pgfscope}%
\begin{pgfscope}%
\pgfsetbuttcap%
\pgfsetroundjoin%
\definecolor{currentfill}{rgb}{0.000000,0.000000,0.000000}%
\pgfsetfillcolor{currentfill}%
\pgfsetlinewidth{0.401500pt}%
\definecolor{currentstroke}{rgb}{0.000000,0.000000,0.000000}%
\pgfsetstrokecolor{currentstroke}%
\pgfsetdash{}{0pt}%
\pgfsys@defobject{currentmarker}{\pgfqpoint{-0.020833in}{0.000000in}}{\pgfqpoint{-0.000000in}{0.000000in}}{%
\pgfpathmoveto{\pgfqpoint{-0.000000in}{0.000000in}}%
\pgfpathlineto{\pgfqpoint{-0.020833in}{0.000000in}}%
\pgfusepath{stroke,fill}%
}%
\begin{pgfscope}%
\pgfsys@transformshift{3.038110in}{1.237761in}%
\pgfsys@useobject{currentmarker}{}%
\end{pgfscope}%
\end{pgfscope}%
\begin{pgfscope}%
\pgfsetbuttcap%
\pgfsetroundjoin%
\definecolor{currentfill}{rgb}{0.000000,0.000000,0.000000}%
\pgfsetfillcolor{currentfill}%
\pgfsetlinewidth{0.401500pt}%
\definecolor{currentstroke}{rgb}{0.000000,0.000000,0.000000}%
\pgfsetstrokecolor{currentstroke}%
\pgfsetdash{}{0pt}%
\pgfsys@defobject{currentmarker}{\pgfqpoint{0.000000in}{0.000000in}}{\pgfqpoint{0.020833in}{0.000000in}}{%
\pgfpathmoveto{\pgfqpoint{0.000000in}{0.000000in}}%
\pgfpathlineto{\pgfqpoint{0.020833in}{0.000000in}}%
\pgfusepath{stroke,fill}%
}%
\begin{pgfscope}%
\pgfsys@transformshift{0.650000in}{1.403800in}%
\pgfsys@useobject{currentmarker}{}%
\end{pgfscope}%
\end{pgfscope}%
\begin{pgfscope}%
\pgfsetbuttcap%
\pgfsetroundjoin%
\definecolor{currentfill}{rgb}{0.000000,0.000000,0.000000}%
\pgfsetfillcolor{currentfill}%
\pgfsetlinewidth{0.401500pt}%
\definecolor{currentstroke}{rgb}{0.000000,0.000000,0.000000}%
\pgfsetstrokecolor{currentstroke}%
\pgfsetdash{}{0pt}%
\pgfsys@defobject{currentmarker}{\pgfqpoint{-0.020833in}{0.000000in}}{\pgfqpoint{-0.000000in}{0.000000in}}{%
\pgfpathmoveto{\pgfqpoint{-0.000000in}{0.000000in}}%
\pgfpathlineto{\pgfqpoint{-0.020833in}{0.000000in}}%
\pgfusepath{stroke,fill}%
}%
\begin{pgfscope}%
\pgfsys@transformshift{3.038110in}{1.403800in}%
\pgfsys@useobject{currentmarker}{}%
\end{pgfscope}%
\end{pgfscope}%
\begin{pgfscope}%
\pgfsetbuttcap%
\pgfsetroundjoin%
\definecolor{currentfill}{rgb}{0.000000,0.000000,0.000000}%
\pgfsetfillcolor{currentfill}%
\pgfsetlinewidth{0.401500pt}%
\definecolor{currentstroke}{rgb}{0.000000,0.000000,0.000000}%
\pgfsetstrokecolor{currentstroke}%
\pgfsetdash{}{0pt}%
\pgfsys@defobject{currentmarker}{\pgfqpoint{0.000000in}{0.000000in}}{\pgfqpoint{0.020833in}{0.000000in}}{%
\pgfpathmoveto{\pgfqpoint{0.000000in}{0.000000in}}%
\pgfpathlineto{\pgfqpoint{0.020833in}{0.000000in}}%
\pgfusepath{stroke,fill}%
}%
\begin{pgfscope}%
\pgfsys@transformshift{0.650000in}{1.486819in}%
\pgfsys@useobject{currentmarker}{}%
\end{pgfscope}%
\end{pgfscope}%
\begin{pgfscope}%
\pgfsetbuttcap%
\pgfsetroundjoin%
\definecolor{currentfill}{rgb}{0.000000,0.000000,0.000000}%
\pgfsetfillcolor{currentfill}%
\pgfsetlinewidth{0.401500pt}%
\definecolor{currentstroke}{rgb}{0.000000,0.000000,0.000000}%
\pgfsetstrokecolor{currentstroke}%
\pgfsetdash{}{0pt}%
\pgfsys@defobject{currentmarker}{\pgfqpoint{-0.020833in}{0.000000in}}{\pgfqpoint{-0.000000in}{0.000000in}}{%
\pgfpathmoveto{\pgfqpoint{-0.000000in}{0.000000in}}%
\pgfpathlineto{\pgfqpoint{-0.020833in}{0.000000in}}%
\pgfusepath{stroke,fill}%
}%
\begin{pgfscope}%
\pgfsys@transformshift{3.038110in}{1.486819in}%
\pgfsys@useobject{currentmarker}{}%
\end{pgfscope}%
\end{pgfscope}%
\begin{pgfscope}%
\pgfsetbuttcap%
\pgfsetroundjoin%
\definecolor{currentfill}{rgb}{0.000000,0.000000,0.000000}%
\pgfsetfillcolor{currentfill}%
\pgfsetlinewidth{0.401500pt}%
\definecolor{currentstroke}{rgb}{0.000000,0.000000,0.000000}%
\pgfsetstrokecolor{currentstroke}%
\pgfsetdash{}{0pt}%
\pgfsys@defobject{currentmarker}{\pgfqpoint{0.000000in}{0.000000in}}{\pgfqpoint{0.020833in}{0.000000in}}{%
\pgfpathmoveto{\pgfqpoint{0.000000in}{0.000000in}}%
\pgfpathlineto{\pgfqpoint{0.020833in}{0.000000in}}%
\pgfusepath{stroke,fill}%
}%
\begin{pgfscope}%
\pgfsys@transformshift{0.650000in}{1.569838in}%
\pgfsys@useobject{currentmarker}{}%
\end{pgfscope}%
\end{pgfscope}%
\begin{pgfscope}%
\pgfsetbuttcap%
\pgfsetroundjoin%
\definecolor{currentfill}{rgb}{0.000000,0.000000,0.000000}%
\pgfsetfillcolor{currentfill}%
\pgfsetlinewidth{0.401500pt}%
\definecolor{currentstroke}{rgb}{0.000000,0.000000,0.000000}%
\pgfsetstrokecolor{currentstroke}%
\pgfsetdash{}{0pt}%
\pgfsys@defobject{currentmarker}{\pgfqpoint{-0.020833in}{0.000000in}}{\pgfqpoint{-0.000000in}{0.000000in}}{%
\pgfpathmoveto{\pgfqpoint{-0.000000in}{0.000000in}}%
\pgfpathlineto{\pgfqpoint{-0.020833in}{0.000000in}}%
\pgfusepath{stroke,fill}%
}%
\begin{pgfscope}%
\pgfsys@transformshift{3.038110in}{1.569838in}%
\pgfsys@useobject{currentmarker}{}%
\end{pgfscope}%
\end{pgfscope}%
\begin{pgfscope}%
\pgfsetbuttcap%
\pgfsetroundjoin%
\definecolor{currentfill}{rgb}{0.000000,0.000000,0.000000}%
\pgfsetfillcolor{currentfill}%
\pgfsetlinewidth{0.401500pt}%
\definecolor{currentstroke}{rgb}{0.000000,0.000000,0.000000}%
\pgfsetstrokecolor{currentstroke}%
\pgfsetdash{}{0pt}%
\pgfsys@defobject{currentmarker}{\pgfqpoint{0.000000in}{0.000000in}}{\pgfqpoint{0.020833in}{0.000000in}}{%
\pgfpathmoveto{\pgfqpoint{0.000000in}{0.000000in}}%
\pgfpathlineto{\pgfqpoint{0.020833in}{0.000000in}}%
\pgfusepath{stroke,fill}%
}%
\begin{pgfscope}%
\pgfsys@transformshift{0.650000in}{1.652858in}%
\pgfsys@useobject{currentmarker}{}%
\end{pgfscope}%
\end{pgfscope}%
\begin{pgfscope}%
\pgfsetbuttcap%
\pgfsetroundjoin%
\definecolor{currentfill}{rgb}{0.000000,0.000000,0.000000}%
\pgfsetfillcolor{currentfill}%
\pgfsetlinewidth{0.401500pt}%
\definecolor{currentstroke}{rgb}{0.000000,0.000000,0.000000}%
\pgfsetstrokecolor{currentstroke}%
\pgfsetdash{}{0pt}%
\pgfsys@defobject{currentmarker}{\pgfqpoint{-0.020833in}{0.000000in}}{\pgfqpoint{-0.000000in}{0.000000in}}{%
\pgfpathmoveto{\pgfqpoint{-0.000000in}{0.000000in}}%
\pgfpathlineto{\pgfqpoint{-0.020833in}{0.000000in}}%
\pgfusepath{stroke,fill}%
}%
\begin{pgfscope}%
\pgfsys@transformshift{3.038110in}{1.652858in}%
\pgfsys@useobject{currentmarker}{}%
\end{pgfscope}%
\end{pgfscope}%
\begin{pgfscope}%
\pgfsetbuttcap%
\pgfsetroundjoin%
\definecolor{currentfill}{rgb}{0.000000,0.000000,0.000000}%
\pgfsetfillcolor{currentfill}%
\pgfsetlinewidth{0.401500pt}%
\definecolor{currentstroke}{rgb}{0.000000,0.000000,0.000000}%
\pgfsetstrokecolor{currentstroke}%
\pgfsetdash{}{0pt}%
\pgfsys@defobject{currentmarker}{\pgfqpoint{0.000000in}{0.000000in}}{\pgfqpoint{0.020833in}{0.000000in}}{%
\pgfpathmoveto{\pgfqpoint{0.000000in}{0.000000in}}%
\pgfpathlineto{\pgfqpoint{0.020833in}{0.000000in}}%
\pgfusepath{stroke,fill}%
}%
\begin{pgfscope}%
\pgfsys@transformshift{0.650000in}{1.818896in}%
\pgfsys@useobject{currentmarker}{}%
\end{pgfscope}%
\end{pgfscope}%
\begin{pgfscope}%
\pgfsetbuttcap%
\pgfsetroundjoin%
\definecolor{currentfill}{rgb}{0.000000,0.000000,0.000000}%
\pgfsetfillcolor{currentfill}%
\pgfsetlinewidth{0.401500pt}%
\definecolor{currentstroke}{rgb}{0.000000,0.000000,0.000000}%
\pgfsetstrokecolor{currentstroke}%
\pgfsetdash{}{0pt}%
\pgfsys@defobject{currentmarker}{\pgfqpoint{-0.020833in}{0.000000in}}{\pgfqpoint{-0.000000in}{0.000000in}}{%
\pgfpathmoveto{\pgfqpoint{-0.000000in}{0.000000in}}%
\pgfpathlineto{\pgfqpoint{-0.020833in}{0.000000in}}%
\pgfusepath{stroke,fill}%
}%
\begin{pgfscope}%
\pgfsys@transformshift{3.038110in}{1.818896in}%
\pgfsys@useobject{currentmarker}{}%
\end{pgfscope}%
\end{pgfscope}%
\begin{pgfscope}%
\pgfsetbuttcap%
\pgfsetroundjoin%
\definecolor{currentfill}{rgb}{0.000000,0.000000,0.000000}%
\pgfsetfillcolor{currentfill}%
\pgfsetlinewidth{0.401500pt}%
\definecolor{currentstroke}{rgb}{0.000000,0.000000,0.000000}%
\pgfsetstrokecolor{currentstroke}%
\pgfsetdash{}{0pt}%
\pgfsys@defobject{currentmarker}{\pgfqpoint{0.000000in}{0.000000in}}{\pgfqpoint{0.020833in}{0.000000in}}{%
\pgfpathmoveto{\pgfqpoint{0.000000in}{0.000000in}}%
\pgfpathlineto{\pgfqpoint{0.020833in}{0.000000in}}%
\pgfusepath{stroke,fill}%
}%
\begin{pgfscope}%
\pgfsys@transformshift{0.650000in}{1.901915in}%
\pgfsys@useobject{currentmarker}{}%
\end{pgfscope}%
\end{pgfscope}%
\begin{pgfscope}%
\pgfsetbuttcap%
\pgfsetroundjoin%
\definecolor{currentfill}{rgb}{0.000000,0.000000,0.000000}%
\pgfsetfillcolor{currentfill}%
\pgfsetlinewidth{0.401500pt}%
\definecolor{currentstroke}{rgb}{0.000000,0.000000,0.000000}%
\pgfsetstrokecolor{currentstroke}%
\pgfsetdash{}{0pt}%
\pgfsys@defobject{currentmarker}{\pgfqpoint{-0.020833in}{0.000000in}}{\pgfqpoint{-0.000000in}{0.000000in}}{%
\pgfpathmoveto{\pgfqpoint{-0.000000in}{0.000000in}}%
\pgfpathlineto{\pgfqpoint{-0.020833in}{0.000000in}}%
\pgfusepath{stroke,fill}%
}%
\begin{pgfscope}%
\pgfsys@transformshift{3.038110in}{1.901915in}%
\pgfsys@useobject{currentmarker}{}%
\end{pgfscope}%
\end{pgfscope}%
\begin{pgfscope}%
\pgfsetbuttcap%
\pgfsetroundjoin%
\definecolor{currentfill}{rgb}{0.000000,0.000000,0.000000}%
\pgfsetfillcolor{currentfill}%
\pgfsetlinewidth{0.401500pt}%
\definecolor{currentstroke}{rgb}{0.000000,0.000000,0.000000}%
\pgfsetstrokecolor{currentstroke}%
\pgfsetdash{}{0pt}%
\pgfsys@defobject{currentmarker}{\pgfqpoint{0.000000in}{0.000000in}}{\pgfqpoint{0.020833in}{0.000000in}}{%
\pgfpathmoveto{\pgfqpoint{0.000000in}{0.000000in}}%
\pgfpathlineto{\pgfqpoint{0.020833in}{0.000000in}}%
\pgfusepath{stroke,fill}%
}%
\begin{pgfscope}%
\pgfsys@transformshift{0.650000in}{1.984935in}%
\pgfsys@useobject{currentmarker}{}%
\end{pgfscope}%
\end{pgfscope}%
\begin{pgfscope}%
\pgfsetbuttcap%
\pgfsetroundjoin%
\definecolor{currentfill}{rgb}{0.000000,0.000000,0.000000}%
\pgfsetfillcolor{currentfill}%
\pgfsetlinewidth{0.401500pt}%
\definecolor{currentstroke}{rgb}{0.000000,0.000000,0.000000}%
\pgfsetstrokecolor{currentstroke}%
\pgfsetdash{}{0pt}%
\pgfsys@defobject{currentmarker}{\pgfqpoint{-0.020833in}{0.000000in}}{\pgfqpoint{-0.000000in}{0.000000in}}{%
\pgfpathmoveto{\pgfqpoint{-0.000000in}{0.000000in}}%
\pgfpathlineto{\pgfqpoint{-0.020833in}{0.000000in}}%
\pgfusepath{stroke,fill}%
}%
\begin{pgfscope}%
\pgfsys@transformshift{3.038110in}{1.984935in}%
\pgfsys@useobject{currentmarker}{}%
\end{pgfscope}%
\end{pgfscope}%
\begin{pgfscope}%
\pgfsetbuttcap%
\pgfsetroundjoin%
\definecolor{currentfill}{rgb}{0.000000,0.000000,0.000000}%
\pgfsetfillcolor{currentfill}%
\pgfsetlinewidth{0.401500pt}%
\definecolor{currentstroke}{rgb}{0.000000,0.000000,0.000000}%
\pgfsetstrokecolor{currentstroke}%
\pgfsetdash{}{0pt}%
\pgfsys@defobject{currentmarker}{\pgfqpoint{0.000000in}{0.000000in}}{\pgfqpoint{0.020833in}{0.000000in}}{%
\pgfpathmoveto{\pgfqpoint{0.000000in}{0.000000in}}%
\pgfpathlineto{\pgfqpoint{0.020833in}{0.000000in}}%
\pgfusepath{stroke,fill}%
}%
\begin{pgfscope}%
\pgfsys@transformshift{0.650000in}{2.067954in}%
\pgfsys@useobject{currentmarker}{}%
\end{pgfscope}%
\end{pgfscope}%
\begin{pgfscope}%
\pgfsetbuttcap%
\pgfsetroundjoin%
\definecolor{currentfill}{rgb}{0.000000,0.000000,0.000000}%
\pgfsetfillcolor{currentfill}%
\pgfsetlinewidth{0.401500pt}%
\definecolor{currentstroke}{rgb}{0.000000,0.000000,0.000000}%
\pgfsetstrokecolor{currentstroke}%
\pgfsetdash{}{0pt}%
\pgfsys@defobject{currentmarker}{\pgfqpoint{-0.020833in}{0.000000in}}{\pgfqpoint{-0.000000in}{0.000000in}}{%
\pgfpathmoveto{\pgfqpoint{-0.000000in}{0.000000in}}%
\pgfpathlineto{\pgfqpoint{-0.020833in}{0.000000in}}%
\pgfusepath{stroke,fill}%
}%
\begin{pgfscope}%
\pgfsys@transformshift{3.038110in}{2.067954in}%
\pgfsys@useobject{currentmarker}{}%
\end{pgfscope}%
\end{pgfscope}%
\begin{pgfscope}%
\pgfsetbuttcap%
\pgfsetroundjoin%
\definecolor{currentfill}{rgb}{0.000000,0.000000,0.000000}%
\pgfsetfillcolor{currentfill}%
\pgfsetlinewidth{0.401500pt}%
\definecolor{currentstroke}{rgb}{0.000000,0.000000,0.000000}%
\pgfsetstrokecolor{currentstroke}%
\pgfsetdash{}{0pt}%
\pgfsys@defobject{currentmarker}{\pgfqpoint{0.000000in}{0.000000in}}{\pgfqpoint{0.020833in}{0.000000in}}{%
\pgfpathmoveto{\pgfqpoint{0.000000in}{0.000000in}}%
\pgfpathlineto{\pgfqpoint{0.020833in}{0.000000in}}%
\pgfusepath{stroke,fill}%
}%
\begin{pgfscope}%
\pgfsys@transformshift{0.650000in}{2.233992in}%
\pgfsys@useobject{currentmarker}{}%
\end{pgfscope}%
\end{pgfscope}%
\begin{pgfscope}%
\pgfsetbuttcap%
\pgfsetroundjoin%
\definecolor{currentfill}{rgb}{0.000000,0.000000,0.000000}%
\pgfsetfillcolor{currentfill}%
\pgfsetlinewidth{0.401500pt}%
\definecolor{currentstroke}{rgb}{0.000000,0.000000,0.000000}%
\pgfsetstrokecolor{currentstroke}%
\pgfsetdash{}{0pt}%
\pgfsys@defobject{currentmarker}{\pgfqpoint{-0.020833in}{0.000000in}}{\pgfqpoint{-0.000000in}{0.000000in}}{%
\pgfpathmoveto{\pgfqpoint{-0.000000in}{0.000000in}}%
\pgfpathlineto{\pgfqpoint{-0.020833in}{0.000000in}}%
\pgfusepath{stroke,fill}%
}%
\begin{pgfscope}%
\pgfsys@transformshift{3.038110in}{2.233992in}%
\pgfsys@useobject{currentmarker}{}%
\end{pgfscope}%
\end{pgfscope}%
\begin{pgfscope}%
\pgfsetbuttcap%
\pgfsetroundjoin%
\definecolor{currentfill}{rgb}{0.000000,0.000000,0.000000}%
\pgfsetfillcolor{currentfill}%
\pgfsetlinewidth{0.401500pt}%
\definecolor{currentstroke}{rgb}{0.000000,0.000000,0.000000}%
\pgfsetstrokecolor{currentstroke}%
\pgfsetdash{}{0pt}%
\pgfsys@defobject{currentmarker}{\pgfqpoint{0.000000in}{0.000000in}}{\pgfqpoint{0.020833in}{0.000000in}}{%
\pgfpathmoveto{\pgfqpoint{0.000000in}{0.000000in}}%
\pgfpathlineto{\pgfqpoint{0.020833in}{0.000000in}}%
\pgfusepath{stroke,fill}%
}%
\begin{pgfscope}%
\pgfsys@transformshift{0.650000in}{2.317012in}%
\pgfsys@useobject{currentmarker}{}%
\end{pgfscope}%
\end{pgfscope}%
\begin{pgfscope}%
\pgfsetbuttcap%
\pgfsetroundjoin%
\definecolor{currentfill}{rgb}{0.000000,0.000000,0.000000}%
\pgfsetfillcolor{currentfill}%
\pgfsetlinewidth{0.401500pt}%
\definecolor{currentstroke}{rgb}{0.000000,0.000000,0.000000}%
\pgfsetstrokecolor{currentstroke}%
\pgfsetdash{}{0pt}%
\pgfsys@defobject{currentmarker}{\pgfqpoint{-0.020833in}{0.000000in}}{\pgfqpoint{-0.000000in}{0.000000in}}{%
\pgfpathmoveto{\pgfqpoint{-0.000000in}{0.000000in}}%
\pgfpathlineto{\pgfqpoint{-0.020833in}{0.000000in}}%
\pgfusepath{stroke,fill}%
}%
\begin{pgfscope}%
\pgfsys@transformshift{3.038110in}{2.317012in}%
\pgfsys@useobject{currentmarker}{}%
\end{pgfscope}%
\end{pgfscope}%
\begin{pgfscope}%
\pgfsetbuttcap%
\pgfsetroundjoin%
\definecolor{currentfill}{rgb}{0.000000,0.000000,0.000000}%
\pgfsetfillcolor{currentfill}%
\pgfsetlinewidth{0.401500pt}%
\definecolor{currentstroke}{rgb}{0.000000,0.000000,0.000000}%
\pgfsetstrokecolor{currentstroke}%
\pgfsetdash{}{0pt}%
\pgfsys@defobject{currentmarker}{\pgfqpoint{0.000000in}{0.000000in}}{\pgfqpoint{0.020833in}{0.000000in}}{%
\pgfpathmoveto{\pgfqpoint{0.000000in}{0.000000in}}%
\pgfpathlineto{\pgfqpoint{0.020833in}{0.000000in}}%
\pgfusepath{stroke,fill}%
}%
\begin{pgfscope}%
\pgfsys@transformshift{0.650000in}{2.400031in}%
\pgfsys@useobject{currentmarker}{}%
\end{pgfscope}%
\end{pgfscope}%
\begin{pgfscope}%
\pgfsetbuttcap%
\pgfsetroundjoin%
\definecolor{currentfill}{rgb}{0.000000,0.000000,0.000000}%
\pgfsetfillcolor{currentfill}%
\pgfsetlinewidth{0.401500pt}%
\definecolor{currentstroke}{rgb}{0.000000,0.000000,0.000000}%
\pgfsetstrokecolor{currentstroke}%
\pgfsetdash{}{0pt}%
\pgfsys@defobject{currentmarker}{\pgfqpoint{-0.020833in}{0.000000in}}{\pgfqpoint{-0.000000in}{0.000000in}}{%
\pgfpathmoveto{\pgfqpoint{-0.000000in}{0.000000in}}%
\pgfpathlineto{\pgfqpoint{-0.020833in}{0.000000in}}%
\pgfusepath{stroke,fill}%
}%
\begin{pgfscope}%
\pgfsys@transformshift{3.038110in}{2.400031in}%
\pgfsys@useobject{currentmarker}{}%
\end{pgfscope}%
\end{pgfscope}%
\begin{pgfscope}%
\definecolor{textcolor}{rgb}{0.000000,0.000000,0.000000}%
\pgfsetstrokecolor{textcolor}%
\pgfsetfillcolor{textcolor}%
\pgftext[x=0.379282in,y=1.417244in,,bottom,rotate=90.000000]{\color{textcolor}{\rmfamily\fontsize{12.000000}{14.400000}\selectfont\catcode`\^=\active\def^{\ifmmode\sp\else\^{}\fi}\catcode`\%=\active\def%{\%}$\delta_{99}/\delta_{0},\ \theta/\delta_{0}$}}%
\end{pgfscope}%
\begin{pgfscope}%
\pgfpathrectangle{\pgfqpoint{0.650000in}{0.420000in}}{\pgfqpoint{2.388110in}{1.994488in}}%
\pgfusepath{clip}%
\pgfsetrectcap%
\pgfsetroundjoin%
\pgfsetlinewidth{1.003750pt}%
\definecolor{currentstroke}{rgb}{0.772549,0.105882,0.490196}%
\pgfsetstrokecolor{currentstroke}%
\pgfsetdash{}{0pt}%
\pgfpathmoveto{\pgfqpoint{0.650000in}{0.627724in}}%
\pgfpathlineto{\pgfqpoint{0.670595in}{0.634596in}}%
\pgfpathlineto{\pgfqpoint{0.689838in}{0.642922in}}%
\pgfpathlineto{\pgfqpoint{0.697677in}{0.648494in}}%
\pgfpathlineto{\pgfqpoint{0.711931in}{0.661251in}}%
\pgfpathlineto{\pgfqpoint{0.724760in}{0.675737in}}%
\pgfpathlineto{\pgfqpoint{0.739727in}{0.696242in}}%
\pgfpathlineto{\pgfqpoint{0.750417in}{0.712743in}}%
\pgfpathlineto{\pgfqpoint{0.771085in}{0.745472in}}%
\pgfpathlineto{\pgfqpoint{0.786765in}{0.769356in}}%
\pgfpathlineto{\pgfqpoint{0.828814in}{0.826891in}}%
\pgfpathlineto{\pgfqpoint{0.843068in}{0.844282in}}%
\pgfpathlineto{\pgfqpoint{0.851620in}{0.854365in}}%
\pgfpathlineto{\pgfqpoint{0.877278in}{0.883446in}}%
\pgfpathlineto{\pgfqpoint{0.896521in}{0.903444in}}%
\pgfpathlineto{\pgfqpoint{0.947835in}{0.954163in}}%
\pgfpathlineto{\pgfqpoint{1.026232in}{1.025992in}}%
\pgfpathlineto{\pgfqpoint{1.036210in}{1.034804in}}%
\pgfpathlineto{\pgfqpoint{1.061867in}{1.057293in}}%
\pgfpathlineto{\pgfqpoint{1.070419in}{1.064978in}}%
\pgfpathlineto{\pgfqpoint{1.103203in}{1.092828in}}%
\pgfpathlineto{\pgfqpoint{1.196567in}{1.169653in}}%
\pgfpathlineto{\pgfqpoint{1.207970in}{1.178828in}}%
\pgfpathlineto{\pgfqpoint{1.395410in}{1.323887in}}%
\pgfpathlineto{\pgfqpoint{1.425343in}{1.346082in}}%
\pgfpathlineto{\pgfqpoint{1.455989in}{1.369274in}}%
\pgfpathlineto{\pgfqpoint{1.469531in}{1.379483in}}%
\pgfpathlineto{\pgfqpoint{1.508017in}{1.405879in}}%
\pgfpathlineto{\pgfqpoint{1.529397in}{1.418267in}}%
\pgfpathlineto{\pgfqpoint{1.546502in}{1.426533in}}%
\pgfpathlineto{\pgfqpoint{1.554342in}{1.429335in}}%
\pgfpathlineto{\pgfqpoint{1.562182in}{1.431122in}}%
\pgfpathlineto{\pgfqpoint{1.569309in}{1.431670in}}%
\pgfpathlineto{\pgfqpoint{1.574298in}{1.431296in}}%
\pgfpathlineto{\pgfqpoint{1.577861in}{1.430397in}}%
\pgfpathlineto{\pgfqpoint{1.583563in}{1.428366in}}%
\pgfpathlineto{\pgfqpoint{1.589264in}{1.425086in}}%
\pgfpathlineto{\pgfqpoint{1.594966in}{1.420429in}}%
\pgfpathlineto{\pgfqpoint{1.599955in}{1.414973in}}%
\pgfpathlineto{\pgfqpoint{1.606369in}{1.406014in}}%
\pgfpathlineto{\pgfqpoint{1.624187in}{1.380018in}}%
\pgfpathlineto{\pgfqpoint{1.629888in}{1.374908in}}%
\pgfpathlineto{\pgfqpoint{1.635590in}{1.371404in}}%
\pgfpathlineto{\pgfqpoint{1.641291in}{1.369313in}}%
\pgfpathlineto{\pgfqpoint{1.650556in}{1.367465in}}%
\pgfpathlineto{\pgfqpoint{1.657683in}{1.368180in}}%
\pgfpathlineto{\pgfqpoint{1.666236in}{1.370302in}}%
\pgfpathlineto{\pgfqpoint{1.685479in}{1.376850in}}%
\pgfpathlineto{\pgfqpoint{1.699020in}{1.383316in}}%
\pgfpathlineto{\pgfqpoint{1.704009in}{1.385986in}}%
\pgfpathlineto{\pgfqpoint{1.704722in}{1.385539in}}%
\pgfpathlineto{\pgfqpoint{1.766014in}{1.419781in}}%
\pgfpathlineto{\pgfqpoint{1.773141in}{1.423921in}}%
\pgfpathlineto{\pgfqpoint{1.773854in}{1.423450in}}%
\pgfpathlineto{\pgfqpoint{1.795235in}{1.435697in}}%
\pgfpathlineto{\pgfqpoint{1.795948in}{1.435198in}}%
\pgfpathlineto{\pgfqpoint{1.818042in}{1.447308in}}%
\pgfpathlineto{\pgfqpoint{1.818754in}{1.446747in}}%
\pgfpathlineto{\pgfqpoint{1.843699in}{1.459157in}}%
\pgfpathlineto{\pgfqpoint{1.844412in}{1.458546in}}%
\pgfpathlineto{\pgfqpoint{1.870069in}{1.471005in}}%
\pgfpathlineto{\pgfqpoint{1.870782in}{1.470380in}}%
\pgfpathlineto{\pgfqpoint{1.894301in}{1.482391in}}%
\pgfpathlineto{\pgfqpoint{1.897864in}{1.484310in}}%
\pgfpathlineto{\pgfqpoint{1.898577in}{1.483726in}}%
\pgfpathlineto{\pgfqpoint{1.919245in}{1.494188in}}%
\pgfpathlineto{\pgfqpoint{1.928510in}{1.498536in}}%
\pgfpathlineto{\pgfqpoint{1.929223in}{1.497871in}}%
\pgfpathlineto{\pgfqpoint{1.995504in}{1.527508in}}%
\pgfpathlineto{\pgfqpoint{1.996217in}{1.526797in}}%
\pgfpathlineto{\pgfqpoint{2.018311in}{1.535917in}}%
\pgfpathlineto{\pgfqpoint{2.028289in}{1.539985in}}%
\pgfpathlineto{\pgfqpoint{2.029001in}{1.539232in}}%
\pgfpathlineto{\pgfqpoint{2.061786in}{1.553276in}}%
\pgfpathlineto{\pgfqpoint{2.062498in}{1.552524in}}%
\pgfpathlineto{\pgfqpoint{2.093145in}{1.566464in}}%
\pgfpathlineto{\pgfqpoint{2.095995in}{1.567802in}}%
\pgfpathlineto{\pgfqpoint{2.096708in}{1.567077in}}%
\pgfpathlineto{\pgfqpoint{2.128780in}{1.581884in}}%
\pgfpathlineto{\pgfqpoint{2.129492in}{1.581120in}}%
\pgfpathlineto{\pgfqpoint{2.150874in}{1.590294in}}%
\pgfpathlineto{\pgfqpoint{2.158713in}{1.593607in}}%
\pgfpathlineto{\pgfqpoint{2.159426in}{1.592787in}}%
\pgfpathlineto{\pgfqpoint{2.187934in}{1.604731in}}%
\pgfpathlineto{\pgfqpoint{2.188647in}{1.605005in}}%
\pgfpathlineto{\pgfqpoint{2.189360in}{1.604156in}}%
\pgfpathlineto{\pgfqpoint{2.224281in}{1.617390in}}%
\pgfpathlineto{\pgfqpoint{2.224994in}{1.616528in}}%
\pgfpathlineto{\pgfqpoint{2.259916in}{1.629332in}}%
\pgfpathlineto{\pgfqpoint{2.260629in}{1.628428in}}%
\pgfpathlineto{\pgfqpoint{2.282723in}{1.636685in}}%
\pgfpathlineto{\pgfqpoint{2.297689in}{1.642362in}}%
\pgfpathlineto{\pgfqpoint{2.298402in}{1.641460in}}%
\pgfpathlineto{\pgfqpoint{2.336175in}{1.655685in}}%
\pgfpathlineto{\pgfqpoint{2.336888in}{1.654758in}}%
\pgfpathlineto{\pgfqpoint{2.372524in}{1.668452in}}%
\pgfpathlineto{\pgfqpoint{2.373236in}{1.667486in}}%
\pgfpathlineto{\pgfqpoint{2.408872in}{1.681152in}}%
\pgfpathlineto{\pgfqpoint{2.409584in}{1.680155in}}%
\pgfpathlineto{\pgfqpoint{2.442368in}{1.692923in}}%
\pgfpathlineto{\pgfqpoint{2.443081in}{1.691925in}}%
\pgfpathlineto{\pgfqpoint{2.480142in}{1.707859in}}%
\pgfpathlineto{\pgfqpoint{2.480854in}{1.706897in}}%
\pgfpathlineto{\pgfqpoint{2.491545in}{1.712723in}}%
\pgfpathlineto{\pgfqpoint{2.503661in}{1.720599in}}%
\pgfpathlineto{\pgfqpoint{2.510788in}{1.725993in}}%
\pgfpathlineto{\pgfqpoint{2.511501in}{1.725258in}}%
\pgfpathlineto{\pgfqpoint{2.520053in}{1.732950in}}%
\pgfpathlineto{\pgfqpoint{2.527893in}{1.741792in}}%
\pgfpathlineto{\pgfqpoint{2.531456in}{1.746400in}}%
\pgfpathlineto{\pgfqpoint{2.532169in}{1.746057in}}%
\pgfpathlineto{\pgfqpoint{2.539296in}{1.757283in}}%
\pgfpathlineto{\pgfqpoint{2.544285in}{1.766946in}}%
\pgfpathlineto{\pgfqpoint{2.552125in}{1.784903in}}%
\pgfpathlineto{\pgfqpoint{2.557827in}{1.802667in}}%
\pgfpathlineto{\pgfqpoint{2.570655in}{1.845844in}}%
\pgfpathlineto{\pgfqpoint{2.572080in}{1.849412in}}%
\pgfpathlineto{\pgfqpoint{2.578495in}{1.867462in}}%
\pgfpathlineto{\pgfqpoint{2.583484in}{1.879284in}}%
\pgfpathlineto{\pgfqpoint{2.584196in}{1.879970in}}%
\pgfpathlineto{\pgfqpoint{2.591323in}{1.893698in}}%
\pgfpathlineto{\pgfqpoint{2.598451in}{1.904794in}}%
\pgfpathlineto{\pgfqpoint{2.600589in}{1.907732in}}%
\pgfpathlineto{\pgfqpoint{2.601302in}{1.907792in}}%
\pgfpathlineto{\pgfqpoint{2.609854in}{1.917958in}}%
\pgfpathlineto{\pgfqpoint{2.617694in}{1.925740in}}%
\pgfpathlineto{\pgfqpoint{2.621970in}{1.929399in}}%
\pgfpathlineto{\pgfqpoint{2.622683in}{1.929016in}}%
\pgfpathlineto{\pgfqpoint{2.631949in}{1.935999in}}%
\pgfpathlineto{\pgfqpoint{2.643352in}{1.943070in}}%
\pgfpathlineto{\pgfqpoint{2.649054in}{1.946227in}}%
\pgfpathlineto{\pgfqpoint{2.649766in}{1.945496in}}%
\pgfpathlineto{\pgfqpoint{2.668297in}{1.954578in}}%
\pgfpathlineto{\pgfqpoint{2.679700in}{1.959279in}}%
\pgfpathlineto{\pgfqpoint{2.680413in}{1.958279in}}%
\pgfpathlineto{\pgfqpoint{2.712485in}{1.971246in}}%
\pgfpathlineto{\pgfqpoint{2.713198in}{1.970088in}}%
\pgfpathlineto{\pgfqpoint{2.745271in}{1.982954in}}%
\pgfpathlineto{\pgfqpoint{2.745984in}{1.981624in}}%
\pgfpathlineto{\pgfqpoint{2.761663in}{1.989024in}}%
\pgfpathlineto{\pgfqpoint{2.772354in}{1.994453in}}%
\pgfpathlineto{\pgfqpoint{2.773067in}{1.992909in}}%
\pgfpathlineto{\pgfqpoint{2.795874in}{2.004038in}}%
\pgfpathlineto{\pgfqpoint{2.799437in}{2.005607in}}%
\pgfpathlineto{\pgfqpoint{2.800150in}{2.003662in}}%
\pgfpathlineto{\pgfqpoint{2.819393in}{2.012693in}}%
\pgfpathlineto{\pgfqpoint{2.826520in}{2.016341in}}%
\pgfpathlineto{\pgfqpoint{2.827233in}{2.014037in}}%
\pgfpathlineto{\pgfqpoint{2.851465in}{2.026105in}}%
\pgfpathlineto{\pgfqpoint{2.852178in}{2.023376in}}%
\pgfpathlineto{\pgfqpoint{2.876408in}{2.034008in}}%
\pgfpathlineto{\pgfqpoint{2.879259in}{2.035233in}}%
\pgfpathlineto{\pgfqpoint{2.879971in}{2.031715in}}%
\pgfpathlineto{\pgfqpoint{2.905627in}{2.042896in}}%
\pgfpathlineto{\pgfqpoint{2.906340in}{2.038401in}}%
\pgfpathlineto{\pgfqpoint{2.932708in}{2.048993in}}%
\pgfpathlineto{\pgfqpoint{2.933421in}{2.043192in}}%
\pgfpathlineto{\pgfqpoint{2.951238in}{2.050482in}}%
\pgfpathlineto{\pgfqpoint{2.959077in}{2.054224in}}%
\pgfpathlineto{\pgfqpoint{2.959789in}{2.047183in}}%
\pgfpathlineto{\pgfqpoint{2.981882in}{2.057972in}}%
\pgfpathlineto{\pgfqpoint{2.982595in}{2.049062in}}%
\pgfpathlineto{\pgfqpoint{3.001124in}{2.058428in}}%
\pgfpathlineto{\pgfqpoint{3.004688in}{2.060293in}}%
\pgfpathlineto{\pgfqpoint{3.005400in}{2.049329in}}%
\pgfpathlineto{\pgfqpoint{3.026066in}{2.059808in}}%
\pgfpathlineto{\pgfqpoint{3.026778in}{2.046351in}}%
\pgfpathlineto{\pgfqpoint{3.038110in}{2.051470in}}%
\pgfpathlineto{\pgfqpoint{3.038110in}{2.051470in}}%
\pgfusepath{stroke}%
\end{pgfscope}%
\begin{pgfscope}%
\pgfpathrectangle{\pgfqpoint{0.650000in}{0.420000in}}{\pgfqpoint{2.388110in}{1.994488in}}%
\pgfusepath{clip}%
\pgfsetrectcap%
\pgfsetroundjoin%
\pgfsetlinewidth{1.003750pt}%
\definecolor{currentstroke}{rgb}{0.945098,0.713725,0.854902}%
\pgfsetstrokecolor{currentstroke}%
\pgfsetdash{}{0pt}%
\pgfpathmoveto{\pgfqpoint{0.650000in}{0.510818in}}%
\pgfpathlineto{\pgfqpoint{0.675584in}{0.511931in}}%
\pgfpathlineto{\pgfqpoint{0.699103in}{0.514229in}}%
\pgfpathlineto{\pgfqpoint{0.710506in}{0.515668in}}%
\pgfpathlineto{\pgfqpoint{0.781063in}{0.525194in}}%
\pgfpathlineto{\pgfqpoint{0.889394in}{0.538623in}}%
\pgfpathlineto{\pgfqpoint{0.987033in}{0.549880in}}%
\pgfpathlineto{\pgfqpoint{1.081822in}{0.560090in}}%
\pgfpathlineto{\pgfqpoint{1.228638in}{0.574936in}}%
\pgfpathlineto{\pgfqpoint{1.274251in}{0.579392in}}%
\pgfpathlineto{\pgfqpoint{1.457415in}{0.596428in}}%
\pgfpathlineto{\pgfqpoint{1.526547in}{0.602094in}}%
\pgfpathlineto{\pgfqpoint{1.553629in}{0.603369in}}%
\pgfpathlineto{\pgfqpoint{1.572159in}{0.603246in}}%
\pgfpathlineto{\pgfqpoint{1.584275in}{0.602163in}}%
\pgfpathlineto{\pgfqpoint{1.594966in}{0.600262in}}%
\pgfpathlineto{\pgfqpoint{1.604231in}{0.597532in}}%
\pgfpathlineto{\pgfqpoint{1.615634in}{0.592864in}}%
\pgfpathlineto{\pgfqpoint{1.626325in}{0.588867in}}%
\pgfpathlineto{\pgfqpoint{1.636302in}{0.586431in}}%
\pgfpathlineto{\pgfqpoint{1.651269in}{0.584186in}}%
\pgfpathlineto{\pgfqpoint{1.666236in}{0.583776in}}%
\pgfpathlineto{\pgfqpoint{1.695457in}{0.584437in}}%
\pgfpathlineto{\pgfqpoint{1.726103in}{0.586246in}}%
\pgfpathlineto{\pgfqpoint{1.727528in}{0.585859in}}%
\pgfpathlineto{\pgfqpoint{1.773141in}{0.589276in}}%
\pgfpathlineto{\pgfqpoint{1.774567in}{0.588892in}}%
\pgfpathlineto{\pgfqpoint{1.795948in}{0.590213in}}%
\pgfpathlineto{\pgfqpoint{1.843699in}{0.593784in}}%
\pgfpathlineto{\pgfqpoint{1.845124in}{0.593397in}}%
\pgfpathlineto{\pgfqpoint{1.894301in}{0.597025in}}%
\pgfpathlineto{\pgfqpoint{1.928510in}{0.599346in}}%
\pgfpathlineto{\pgfqpoint{1.929936in}{0.598950in}}%
\pgfpathlineto{\pgfqpoint{1.995504in}{0.603672in}}%
\pgfpathlineto{\pgfqpoint{1.996930in}{0.603269in}}%
\pgfpathlineto{\pgfqpoint{2.061786in}{0.607744in}}%
\pgfpathlineto{\pgfqpoint{2.063211in}{0.607339in}}%
\pgfpathlineto{\pgfqpoint{2.128780in}{0.611895in}}%
\pgfpathlineto{\pgfqpoint{2.130205in}{0.611483in}}%
\pgfpathlineto{\pgfqpoint{2.188647in}{0.615350in}}%
\pgfpathlineto{\pgfqpoint{2.190073in}{0.614928in}}%
\pgfpathlineto{\pgfqpoint{2.224281in}{0.617430in}}%
\pgfpathlineto{\pgfqpoint{2.225706in}{0.616998in}}%
\pgfpathlineto{\pgfqpoint{2.297689in}{0.621504in}}%
\pgfpathlineto{\pgfqpoint{2.299115in}{0.621063in}}%
\pgfpathlineto{\pgfqpoint{2.336175in}{0.623640in}}%
\pgfpathlineto{\pgfqpoint{2.337601in}{0.623192in}}%
\pgfpathlineto{\pgfqpoint{2.408872in}{0.627611in}}%
\pgfpathlineto{\pgfqpoint{2.410297in}{0.627146in}}%
\pgfpathlineto{\pgfqpoint{2.500810in}{0.633960in}}%
\pgfpathlineto{\pgfqpoint{2.524329in}{0.637842in}}%
\pgfpathlineto{\pgfqpoint{2.540009in}{0.642459in}}%
\pgfpathlineto{\pgfqpoint{2.551412in}{0.647935in}}%
\pgfpathlineto{\pgfqpoint{2.562103in}{0.655577in}}%
\pgfpathlineto{\pgfqpoint{2.570655in}{0.662062in}}%
\pgfpathlineto{\pgfqpoint{2.572080in}{0.662705in}}%
\pgfpathlineto{\pgfqpoint{2.579920in}{0.667115in}}%
\pgfpathlineto{\pgfqpoint{2.586335in}{0.669634in}}%
\pgfpathlineto{\pgfqpoint{2.596313in}{0.672874in}}%
\pgfpathlineto{\pgfqpoint{2.606291in}{0.675004in}}%
\pgfpathlineto{\pgfqpoint{2.621970in}{0.677661in}}%
\pgfpathlineto{\pgfqpoint{2.624109in}{0.677702in}}%
\pgfpathlineto{\pgfqpoint{2.646203in}{0.680053in}}%
\pgfpathlineto{\pgfqpoint{2.674711in}{0.682006in}}%
\pgfpathlineto{\pgfqpoint{2.712485in}{0.684166in}}%
\pgfpathlineto{\pgfqpoint{2.713911in}{0.683895in}}%
\pgfpathlineto{\pgfqpoint{2.765227in}{0.686566in}}%
\pgfpathlineto{\pgfqpoint{2.799437in}{0.688122in}}%
\pgfpathlineto{\pgfqpoint{2.800863in}{0.687639in}}%
\pgfpathlineto{\pgfqpoint{2.826520in}{0.689191in}}%
\pgfpathlineto{\pgfqpoint{2.827946in}{0.688602in}}%
\pgfpathlineto{\pgfqpoint{2.851465in}{0.690022in}}%
\pgfpathlineto{\pgfqpoint{2.852890in}{0.689301in}}%
\pgfpathlineto{\pgfqpoint{2.879259in}{0.690827in}}%
\pgfpathlineto{\pgfqpoint{2.879971in}{0.689876in}}%
\pgfpathlineto{\pgfqpoint{2.905627in}{0.691189in}}%
\pgfpathlineto{\pgfqpoint{2.906340in}{0.689988in}}%
\pgfpathlineto{\pgfqpoint{2.932708in}{0.691338in}}%
\pgfpathlineto{\pgfqpoint{2.933421in}{0.689809in}}%
\pgfpathlineto{\pgfqpoint{2.959077in}{0.691055in}}%
\pgfpathlineto{\pgfqpoint{2.959789in}{0.689100in}}%
\pgfpathlineto{\pgfqpoint{2.981882in}{0.690288in}}%
\pgfpathlineto{\pgfqpoint{2.982595in}{0.687794in}}%
\pgfpathlineto{\pgfqpoint{3.004688in}{0.689129in}}%
\pgfpathlineto{\pgfqpoint{3.005400in}{0.685941in}}%
\pgfpathlineto{\pgfqpoint{3.026066in}{0.687353in}}%
\pgfpathlineto{\pgfqpoint{3.026778in}{0.683254in}}%
\pgfpathlineto{\pgfqpoint{3.038110in}{0.684030in}}%
\pgfpathlineto{\pgfqpoint{3.038110in}{0.684030in}}%
\pgfusepath{stroke}%
\end{pgfscope}%
\begin{pgfscope}%
\pgfpathrectangle{\pgfqpoint{0.650000in}{0.420000in}}{\pgfqpoint{2.388110in}{1.994488in}}%
\pgfusepath{clip}%
\pgfsetrectcap%
\pgfsetroundjoin%
\pgfsetlinewidth{1.003750pt}%
\definecolor{currentstroke}{rgb}{0.301961,0.572549,0.129412}%
\pgfsetstrokecolor{currentstroke}%
\pgfsetdash{}{0pt}%
\pgfpathmoveto{\pgfqpoint{0.650000in}{0.626645in}}%
\pgfpathlineto{\pgfqpoint{0.670433in}{0.633464in}}%
\pgfpathlineto{\pgfqpoint{0.689524in}{0.641725in}}%
\pgfpathlineto{\pgfqpoint{0.697302in}{0.647233in}}%
\pgfpathlineto{\pgfqpoint{0.711444in}{0.659924in}}%
\pgfpathlineto{\pgfqpoint{0.723465in}{0.673633in}}%
\pgfpathlineto{\pgfqpoint{0.740435in}{0.697690in}}%
\pgfpathlineto{\pgfqpoint{0.778618in}{0.756345in}}%
\pgfpathlineto{\pgfqpoint{0.819629in}{0.812763in}}%
\pgfpathlineto{\pgfqpoint{0.845085in}{0.843984in}}%
\pgfpathlineto{\pgfqpoint{0.850034in}{0.850120in}}%
\pgfpathlineto{\pgfqpoint{0.865590in}{0.868194in}}%
\pgfpathlineto{\pgfqpoint{0.904480in}{0.909934in}}%
\pgfpathlineto{\pgfqpoint{0.940542in}{0.945548in}}%
\pgfpathlineto{\pgfqpoint{0.950441in}{0.954903in}}%
\pgfpathlineto{\pgfqpoint{0.957512in}{0.961794in}}%
\pgfpathlineto{\pgfqpoint{0.978725in}{0.981406in}}%
\pgfpathlineto{\pgfqpoint{0.998524in}{0.999493in}}%
\pgfpathlineto{\pgfqpoint{1.010544in}{1.010361in}}%
\pgfpathlineto{\pgfqpoint{1.019736in}{1.018965in}}%
\pgfpathlineto{\pgfqpoint{1.070647in}{1.064291in}}%
\pgfpathlineto{\pgfqpoint{1.168933in}{1.146855in}}%
\pgfpathlineto{\pgfqpoint{1.192267in}{1.166628in}}%
\pgfpathlineto{\pgfqpoint{1.204995in}{1.177653in}}%
\pgfpathlineto{\pgfqpoint{1.258734in}{1.220193in}}%
\pgfpathlineto{\pgfqpoint{1.269340in}{1.228231in}}%
\pgfpathlineto{\pgfqpoint{1.328736in}{1.271700in}}%
\pgfpathlineto{\pgfqpoint{1.352777in}{1.286841in}}%
\pgfpathlineto{\pgfqpoint{1.367626in}{1.294713in}}%
\pgfpathlineto{\pgfqpoint{1.373990in}{1.297514in}}%
\pgfpathlineto{\pgfqpoint{1.386010in}{1.302109in}}%
\pgfpathlineto{\pgfqpoint{1.401566in}{1.307746in}}%
\pgfpathlineto{\pgfqpoint{1.411466in}{1.313222in}}%
\pgfpathlineto{\pgfqpoint{1.419951in}{1.319525in}}%
\pgfpathlineto{\pgfqpoint{1.434093in}{1.332705in}}%
\pgfpathlineto{\pgfqpoint{1.455305in}{1.356810in}}%
\pgfpathlineto{\pgfqpoint{1.466619in}{1.370917in}}%
\pgfpathlineto{\pgfqpoint{1.501266in}{1.413101in}}%
\pgfpathlineto{\pgfqpoint{1.513287in}{1.423507in}}%
\pgfpathlineto{\pgfqpoint{1.520358in}{1.427662in}}%
\pgfpathlineto{\pgfqpoint{1.527429in}{1.430469in}}%
\pgfpathlineto{\pgfqpoint{1.533793in}{1.431973in}}%
\pgfpathlineto{\pgfqpoint{1.540156in}{1.432426in}}%
\pgfpathlineto{\pgfqpoint{1.547227in}{1.431813in}}%
\pgfpathlineto{\pgfqpoint{1.553591in}{1.430157in}}%
\pgfpathlineto{\pgfqpoint{1.562783in}{1.426576in}}%
\pgfpathlineto{\pgfqpoint{1.570561in}{1.422361in}}%
\pgfpathlineto{\pgfqpoint{1.577632in}{1.417364in}}%
\pgfpathlineto{\pgfqpoint{1.589653in}{1.407100in}}%
\pgfpathlineto{\pgfqpoint{1.592481in}{1.403962in}}%
\pgfpathlineto{\pgfqpoint{1.617937in}{1.378464in}}%
\pgfpathlineto{\pgfqpoint{1.625007in}{1.373805in}}%
\pgfpathlineto{\pgfqpoint{1.632078in}{1.370756in}}%
\pgfpathlineto{\pgfqpoint{1.641978in}{1.367983in}}%
\pgfpathlineto{\pgfqpoint{1.649756in}{1.367858in}}%
\pgfpathlineto{\pgfqpoint{1.658241in}{1.368858in}}%
\pgfpathlineto{\pgfqpoint{1.662483in}{1.369722in}}%
\pgfpathlineto{\pgfqpoint{1.663191in}{1.369106in}}%
\pgfpathlineto{\pgfqpoint{1.673797in}{1.372179in}}%
\pgfpathlineto{\pgfqpoint{1.697131in}{1.380693in}}%
\pgfpathlineto{\pgfqpoint{1.702081in}{1.382875in}}%
\pgfpathlineto{\pgfqpoint{1.702788in}{1.382306in}}%
\pgfpathlineto{\pgfqpoint{1.716930in}{1.389303in}}%
\pgfpathlineto{\pgfqpoint{1.725415in}{1.394175in}}%
\pgfpathlineto{\pgfqpoint{1.726122in}{1.393677in}}%
\pgfpathlineto{\pgfqpoint{1.742385in}{1.404138in}}%
\pgfpathlineto{\pgfqpoint{1.748042in}{1.408052in}}%
\pgfpathlineto{\pgfqpoint{1.748749in}{1.407610in}}%
\pgfpathlineto{\pgfqpoint{1.762891in}{1.418190in}}%
\pgfpathlineto{\pgfqpoint{1.772083in}{1.425836in}}%
\pgfpathlineto{\pgfqpoint{1.772791in}{1.425534in}}%
\pgfpathlineto{\pgfqpoint{1.789054in}{1.440081in}}%
\pgfpathlineto{\pgfqpoint{1.798246in}{1.448445in}}%
\pgfpathlineto{\pgfqpoint{1.798953in}{1.448246in}}%
\pgfpathlineto{\pgfqpoint{1.815216in}{1.462339in}}%
\pgfpathlineto{\pgfqpoint{1.824409in}{1.469313in}}%
\pgfpathlineto{\pgfqpoint{1.825116in}{1.469055in}}%
\pgfpathlineto{\pgfqpoint{1.840672in}{1.479336in}}%
\pgfpathlineto{\pgfqpoint{1.848450in}{1.483482in}}%
\pgfpathlineto{\pgfqpoint{1.849157in}{1.483051in}}%
\pgfpathlineto{\pgfqpoint{1.858349in}{1.486685in}}%
\pgfpathlineto{\pgfqpoint{1.868248in}{1.489342in}}%
\pgfpathlineto{\pgfqpoint{1.873905in}{1.490556in}}%
\pgfpathlineto{\pgfqpoint{1.874612in}{1.489842in}}%
\pgfpathlineto{\pgfqpoint{1.900067in}{1.494742in}}%
\pgfpathlineto{\pgfqpoint{1.900775in}{1.493961in}}%
\pgfpathlineto{\pgfqpoint{1.913502in}{1.497464in}}%
\pgfpathlineto{\pgfqpoint{1.924109in}{1.501448in}}%
\pgfpathlineto{\pgfqpoint{1.925523in}{1.502069in}}%
\pgfpathlineto{\pgfqpoint{1.926230in}{1.501370in}}%
\pgfpathlineto{\pgfqpoint{1.938251in}{1.507380in}}%
\pgfpathlineto{\pgfqpoint{1.950271in}{1.514285in}}%
\pgfpathlineto{\pgfqpoint{1.950978in}{1.513683in}}%
\pgfpathlineto{\pgfqpoint{1.978555in}{1.530502in}}%
\pgfpathlineto{\pgfqpoint{1.979262in}{1.529961in}}%
\pgfpathlineto{\pgfqpoint{2.008960in}{1.548762in}}%
\pgfpathlineto{\pgfqpoint{2.009667in}{1.548278in}}%
\pgfpathlineto{\pgfqpoint{2.038658in}{1.566372in}}%
\pgfpathlineto{\pgfqpoint{2.039365in}{1.565911in}}%
\pgfpathlineto{\pgfqpoint{2.067649in}{1.582541in}}%
\pgfpathlineto{\pgfqpoint{2.068356in}{1.582038in}}%
\pgfpathlineto{\pgfqpoint{2.097347in}{1.596929in}}%
\pgfpathlineto{\pgfqpoint{2.098054in}{1.596377in}}%
\pgfpathlineto{\pgfqpoint{2.125631in}{1.608746in}}%
\pgfpathlineto{\pgfqpoint{2.126338in}{1.608162in}}%
\pgfpathlineto{\pgfqpoint{2.151086in}{1.619132in}}%
\pgfpathlineto{\pgfqpoint{2.151793in}{1.618502in}}%
\pgfpathlineto{\pgfqpoint{2.169471in}{1.625499in}}%
\pgfpathlineto{\pgfqpoint{2.175835in}{1.628022in}}%
\pgfpathlineto{\pgfqpoint{2.176542in}{1.627320in}}%
\pgfpathlineto{\pgfqpoint{2.200583in}{1.637129in}}%
\pgfpathlineto{\pgfqpoint{2.201290in}{1.636412in}}%
\pgfpathlineto{\pgfqpoint{2.228866in}{1.647835in}}%
\pgfpathlineto{\pgfqpoint{2.232401in}{1.649448in}}%
\pgfpathlineto{\pgfqpoint{2.233108in}{1.648752in}}%
\pgfpathlineto{\pgfqpoint{2.247958in}{1.656378in}}%
\pgfpathlineto{\pgfqpoint{2.273413in}{1.670756in}}%
\pgfpathlineto{\pgfqpoint{2.289676in}{1.681639in}}%
\pgfpathlineto{\pgfqpoint{2.290383in}{1.681089in}}%
\pgfpathlineto{\pgfqpoint{2.305232in}{1.690028in}}%
\pgfpathlineto{\pgfqpoint{2.315839in}{1.695144in}}%
\pgfpathlineto{\pgfqpoint{2.321496in}{1.697544in}}%
\pgfpathlineto{\pgfqpoint{2.322203in}{1.696862in}}%
\pgfpathlineto{\pgfqpoint{2.354022in}{1.709618in}}%
\pgfpathlineto{\pgfqpoint{2.354729in}{1.708901in}}%
\pgfpathlineto{\pgfqpoint{2.370993in}{1.716073in}}%
\pgfpathlineto{\pgfqpoint{2.383720in}{1.722395in}}%
\pgfpathlineto{\pgfqpoint{2.384427in}{1.721777in}}%
\pgfpathlineto{\pgfqpoint{2.395034in}{1.728144in}}%
\pgfpathlineto{\pgfqpoint{2.414833in}{1.741494in}}%
\pgfpathlineto{\pgfqpoint{2.429681in}{1.754368in}}%
\pgfpathlineto{\pgfqpoint{2.430389in}{1.755010in}}%
\pgfpathlineto{\pgfqpoint{2.431096in}{1.754694in}}%
\pgfpathlineto{\pgfqpoint{2.448773in}{1.771601in}}%
\pgfpathlineto{\pgfqpoint{2.452309in}{1.775162in}}%
\pgfpathlineto{\pgfqpoint{2.453016in}{1.774984in}}%
\pgfpathlineto{\pgfqpoint{2.466451in}{1.789443in}}%
\pgfpathlineto{\pgfqpoint{2.483421in}{1.809069in}}%
\pgfpathlineto{\pgfqpoint{2.508169in}{1.842289in}}%
\pgfpathlineto{\pgfqpoint{2.522311in}{1.865114in}}%
\pgfpathlineto{\pgfqpoint{2.537161in}{1.893290in}}%
\pgfpathlineto{\pgfqpoint{2.564737in}{1.950730in}}%
\pgfpathlineto{\pgfqpoint{2.572515in}{1.965032in}}%
\pgfpathlineto{\pgfqpoint{2.580293in}{1.978105in}}%
\pgfpathlineto{\pgfqpoint{2.592314in}{1.997503in}}%
\pgfpathlineto{\pgfqpoint{2.602921in}{2.012393in}}%
\pgfpathlineto{\pgfqpoint{2.613528in}{2.025655in}}%
\pgfpathlineto{\pgfqpoint{2.625548in}{2.038782in}}%
\pgfpathlineto{\pgfqpoint{2.640398in}{2.052844in}}%
\pgfpathlineto{\pgfqpoint{2.655954in}{2.066442in}}%
\pgfpathlineto{\pgfqpoint{2.691310in}{2.095790in}}%
\pgfpathlineto{\pgfqpoint{2.713230in}{2.111462in}}%
\pgfpathlineto{\pgfqpoint{2.764849in}{2.145497in}}%
\pgfpathlineto{\pgfqpoint{2.803740in}{2.169865in}}%
\pgfpathlineto{\pgfqpoint{2.814347in}{2.176219in}}%
\pgfpathlineto{\pgfqpoint{2.815762in}{2.177602in}}%
\pgfpathlineto{\pgfqpoint{2.828489in}{2.186245in}}%
\pgfpathlineto{\pgfqpoint{2.832731in}{2.189271in}}%
\pgfpathlineto{\pgfqpoint{2.834145in}{2.191563in}}%
\pgfpathlineto{\pgfqpoint{2.850408in}{2.202537in}}%
\pgfpathlineto{\pgfqpoint{2.851822in}{2.203493in}}%
\pgfpathlineto{\pgfqpoint{2.852529in}{2.208450in}}%
\pgfpathlineto{\pgfqpoint{2.909800in}{2.248168in}}%
\pgfpathlineto{\pgfqpoint{2.954345in}{2.276984in}}%
\pgfpathlineto{\pgfqpoint{2.965658in}{2.285511in}}%
\pgfpathlineto{\pgfqpoint{2.979799in}{2.297147in}}%
\pgfpathlineto{\pgfqpoint{2.991112in}{2.308687in}}%
\pgfpathlineto{\pgfqpoint{2.999597in}{2.317413in}}%
\pgfpathlineto{\pgfqpoint{3.004546in}{2.320795in}}%
\pgfpathlineto{\pgfqpoint{3.008784in}{2.322587in}}%
\pgfpathlineto{\pgfqpoint{3.014299in}{2.323601in}}%
\pgfpathlineto{\pgfqpoint{3.019322in}{2.323825in}}%
\pgfpathlineto{\pgfqpoint{3.019322in}{2.323825in}}%
\pgfusepath{stroke}%
\end{pgfscope}%
\begin{pgfscope}%
\pgfpathrectangle{\pgfqpoint{0.650000in}{0.420000in}}{\pgfqpoint{2.388110in}{1.994488in}}%
\pgfusepath{clip}%
\pgfsetrectcap%
\pgfsetroundjoin%
\pgfsetlinewidth{1.003750pt}%
\definecolor{currentstroke}{rgb}{0.721569,0.882353,0.525490}%
\pgfsetstrokecolor{currentstroke}%
\pgfsetdash{}{0pt}%
\pgfpathmoveto{\pgfqpoint{0.650000in}{0.510659in}}%
\pgfpathlineto{\pgfqpoint{0.675382in}{0.511763in}}%
\pgfpathlineto{\pgfqpoint{0.698717in}{0.514049in}}%
\pgfpathlineto{\pgfqpoint{0.710030in}{0.515472in}}%
\pgfpathlineto{\pgfqpoint{0.783568in}{0.525387in}}%
\pgfpathlineto{\pgfqpoint{0.930643in}{0.543242in}}%
\pgfpathlineto{\pgfqpoint{1.028929in}{0.554219in}}%
\pgfpathlineto{\pgfqpoint{1.180246in}{0.570043in}}%
\pgfpathlineto{\pgfqpoint{1.284896in}{0.580376in}}%
\pgfpathlineto{\pgfqpoint{1.351363in}{0.586040in}}%
\pgfpathlineto{\pgfqpoint{1.381061in}{0.587520in}}%
\pgfpathlineto{\pgfqpoint{1.419244in}{0.589243in}}%
\pgfpathlineto{\pgfqpoint{1.443992in}{0.592400in}}%
\pgfpathlineto{\pgfqpoint{1.512580in}{0.602137in}}%
\pgfpathlineto{\pgfqpoint{1.530964in}{0.603035in}}%
\pgfpathlineto{\pgfqpoint{1.547934in}{0.602778in}}%
\pgfpathlineto{\pgfqpoint{1.564905in}{0.601409in}}%
\pgfpathlineto{\pgfqpoint{1.580461in}{0.599127in}}%
\pgfpathlineto{\pgfqpoint{1.593188in}{0.596217in}}%
\pgfpathlineto{\pgfqpoint{1.608744in}{0.591641in}}%
\pgfpathlineto{\pgfqpoint{1.622886in}{0.587856in}}%
\pgfpathlineto{\pgfqpoint{1.635614in}{0.585671in}}%
\pgfpathlineto{\pgfqpoint{1.655413in}{0.583683in}}%
\pgfpathlineto{\pgfqpoint{1.689353in}{0.583023in}}%
\pgfpathlineto{\pgfqpoint{1.735314in}{0.585567in}}%
\pgfpathlineto{\pgfqpoint{1.769255in}{0.590118in}}%
\pgfpathlineto{\pgfqpoint{1.798246in}{0.595224in}}%
\pgfpathlineto{\pgfqpoint{1.799660in}{0.595037in}}%
\pgfpathlineto{\pgfqpoint{1.823701in}{0.599175in}}%
\pgfpathlineto{\pgfqpoint{1.834308in}{0.600136in}}%
\pgfpathlineto{\pgfqpoint{1.848450in}{0.601453in}}%
\pgfpathlineto{\pgfqpoint{1.849864in}{0.601139in}}%
\pgfpathlineto{\pgfqpoint{1.870370in}{0.601957in}}%
\pgfpathlineto{\pgfqpoint{1.885926in}{0.601696in}}%
\pgfpathlineto{\pgfqpoint{1.932594in}{0.602683in}}%
\pgfpathlineto{\pgfqpoint{1.971484in}{0.606649in}}%
\pgfpathlineto{\pgfqpoint{2.008960in}{0.611167in}}%
\pgfpathlineto{\pgfqpoint{2.010374in}{0.610871in}}%
\pgfpathlineto{\pgfqpoint{2.038658in}{0.614435in}}%
\pgfpathlineto{\pgfqpoint{2.040072in}{0.614149in}}%
\pgfpathlineto{\pgfqpoint{2.067649in}{0.617198in}}%
\pgfpathlineto{\pgfqpoint{2.069063in}{0.616905in}}%
\pgfpathlineto{\pgfqpoint{2.097347in}{0.619639in}}%
\pgfpathlineto{\pgfqpoint{2.098761in}{0.619333in}}%
\pgfpathlineto{\pgfqpoint{2.127045in}{0.621336in}}%
\pgfpathlineto{\pgfqpoint{2.151086in}{0.623170in}}%
\pgfpathlineto{\pgfqpoint{2.152501in}{0.622816in}}%
\pgfpathlineto{\pgfqpoint{2.175835in}{0.624421in}}%
\pgfpathlineto{\pgfqpoint{2.177249in}{0.624033in}}%
\pgfpathlineto{\pgfqpoint{2.250786in}{0.629253in}}%
\pgfpathlineto{\pgfqpoint{2.308061in}{0.635582in}}%
\pgfpathlineto{\pgfqpoint{2.321496in}{0.636850in}}%
\pgfpathlineto{\pgfqpoint{2.322910in}{0.636530in}}%
\pgfpathlineto{\pgfqpoint{2.354022in}{0.638649in}}%
\pgfpathlineto{\pgfqpoint{2.355437in}{0.638291in}}%
\pgfpathlineto{\pgfqpoint{2.395034in}{0.641417in}}%
\pgfpathlineto{\pgfqpoint{2.428267in}{0.646205in}}%
\pgfpathlineto{\pgfqpoint{2.477057in}{0.655817in}}%
\pgfpathlineto{\pgfqpoint{2.509583in}{0.664722in}}%
\pgfpathlineto{\pgfqpoint{2.532918in}{0.673220in}}%
\pgfpathlineto{\pgfqpoint{2.579586in}{0.693341in}}%
\pgfpathlineto{\pgfqpoint{2.582415in}{0.694171in}}%
\pgfpathlineto{\pgfqpoint{2.597264in}{0.698688in}}%
\pgfpathlineto{\pgfqpoint{2.614235in}{0.702997in}}%
\pgfpathlineto{\pgfqpoint{2.634034in}{0.707056in}}%
\pgfpathlineto{\pgfqpoint{2.659490in}{0.711442in}}%
\pgfpathlineto{\pgfqpoint{2.692017in}{0.715948in}}%
\pgfpathlineto{\pgfqpoint{2.725251in}{0.719359in}}%
\pgfpathlineto{\pgfqpoint{2.784649in}{0.724335in}}%
\pgfpathlineto{\pgfqpoint{2.851822in}{0.729989in}}%
\pgfpathlineto{\pgfqpoint{2.853236in}{0.731033in}}%
\pgfpathlineto{\pgfqpoint{2.916871in}{0.736514in}}%
\pgfpathlineto{\pgfqpoint{2.986163in}{0.742523in}}%
\pgfpathlineto{\pgfqpoint{3.007373in}{0.744768in}}%
\pgfpathlineto{\pgfqpoint{3.019322in}{0.745019in}}%
\pgfpathlineto{\pgfqpoint{3.019322in}{0.745019in}}%
\pgfusepath{stroke}%
\end{pgfscope}%
\begin{pgfscope}%
\pgfsetrectcap%
\pgfsetmiterjoin%
\pgfsetlinewidth{0.803000pt}%
\definecolor{currentstroke}{rgb}{0.000000,0.000000,0.000000}%
\pgfsetstrokecolor{currentstroke}%
\pgfsetdash{}{0pt}%
\pgfpathmoveto{\pgfqpoint{0.650000in}{0.420000in}}%
\pgfpathlineto{\pgfqpoint{0.650000in}{2.414488in}}%
\pgfusepath{stroke}%
\end{pgfscope}%
\begin{pgfscope}%
\pgfsetrectcap%
\pgfsetmiterjoin%
\pgfsetlinewidth{0.803000pt}%
\definecolor{currentstroke}{rgb}{0.000000,0.000000,0.000000}%
\pgfsetstrokecolor{currentstroke}%
\pgfsetdash{}{0pt}%
\pgfpathmoveto{\pgfqpoint{3.038110in}{0.420000in}}%
\pgfpathlineto{\pgfqpoint{3.038110in}{2.414488in}}%
\pgfusepath{stroke}%
\end{pgfscope}%
\begin{pgfscope}%
\pgfsetrectcap%
\pgfsetmiterjoin%
\pgfsetlinewidth{0.803000pt}%
\definecolor{currentstroke}{rgb}{0.000000,0.000000,0.000000}%
\pgfsetstrokecolor{currentstroke}%
\pgfsetdash{}{0pt}%
\pgfpathmoveto{\pgfqpoint{0.650000in}{0.420000in}}%
\pgfpathlineto{\pgfqpoint{3.038110in}{0.420000in}}%
\pgfusepath{stroke}%
\end{pgfscope}%
\begin{pgfscope}%
\pgfsetrectcap%
\pgfsetmiterjoin%
\pgfsetlinewidth{0.803000pt}%
\definecolor{currentstroke}{rgb}{0.000000,0.000000,0.000000}%
\pgfsetstrokecolor{currentstroke}%
\pgfsetdash{}{0pt}%
\pgfpathmoveto{\pgfqpoint{0.650000in}{2.414488in}}%
\pgfpathlineto{\pgfqpoint{3.038110in}{2.414488in}}%
\pgfusepath{stroke}%
\end{pgfscope}%
\begin{pgfscope}%
\pgfsetrectcap%
\pgfsetroundjoin%
\pgfsetlinewidth{1.003750pt}%
\definecolor{currentstroke}{rgb}{0.945098,0.713725,0.854902}%
\pgfsetstrokecolor{currentstroke}%
\pgfsetdash{}{0pt}%
\pgfpathmoveto{\pgfqpoint{0.750000in}{2.275599in}}%
\pgfpathlineto{\pgfqpoint{0.861111in}{2.275599in}}%
\pgfpathlineto{\pgfqpoint{0.972222in}{2.275599in}}%
\pgfusepath{stroke}%
\end{pgfscope}%
\begin{pgfscope}%
\definecolor{textcolor}{rgb}{0.000000,0.000000,0.000000}%
\pgfsetstrokecolor{textcolor}%
\pgfsetfillcolor{textcolor}%
\pgftext[x=1.061111in,y=2.236710in,left,base]{\color{textcolor}{\rmfamily\fontsize{8.000000}{9.600000}\selectfont\catcode`\^=\active\def^{\ifmmode\sp\else\^{}\fi}\catcode`\%=\active\def%{\%}base}}%
\end{pgfscope}%
\begin{pgfscope}%
\pgfsetrectcap%
\pgfsetroundjoin%
\pgfsetlinewidth{1.003750pt}%
\definecolor{currentstroke}{rgb}{0.721569,0.882353,0.525490}%
\pgfsetstrokecolor{currentstroke}%
\pgfsetdash{}{0pt}%
\pgfpathmoveto{\pgfqpoint{0.750000in}{2.120661in}}%
\pgfpathlineto{\pgfqpoint{0.861111in}{2.120661in}}%
\pgfpathlineto{\pgfqpoint{0.972222in}{2.120661in}}%
\pgfusepath{stroke}%
\end{pgfscope}%
\begin{pgfscope}%
\definecolor{textcolor}{rgb}{0.000000,0.000000,0.000000}%
\pgfsetstrokecolor{textcolor}%
\pgfsetfillcolor{textcolor}%
\pgftext[x=1.061111in,y=2.081772in,left,base]{\color{textcolor}{\rmfamily\fontsize{8.000000}{9.600000}\selectfont\catcode`\^=\active\def^{\ifmmode\sp\else\^{}\fi}\catcode`\%=\active\def%{\%}optimised}}%
\end{pgfscope}%
\end{pgfpicture}%
\makeatother%
\endgroup%

    \label{fig:R100deg30Re1200Thickness}
\end{subfigure}
\\[-15pt]
\caption[First and second order statistics of R100deg30 base and optimised simulations at $\mathrm{Re}_{\theta} = 1200$ with reference data taken from \citet{jimenez_near-wall_2013} at $\mathrm{Re}_{\theta} = 1100$, and streamwise boundary layer thickness and momentum thickness development normalised with reference boundary layer thickness $\delta_0$]{First and second order statistics of R100deg30 base and optimised simulations at $\mathrm{Re}_{\theta} = 1200$ with reference data taken from \citet{jimenez_near-wall_2013} at $\mathrm{Re}_{\theta} = 1100$: (a) streamwise velocity profile in viscous scales; (b) Reynolds stresses $\langle u_i' u_j' \rangle$; (c) budget of \acrshort{tke} = $\frac{1}{2} \langle u_i' u_i' \rangle$, line styles denote budget terms according to \autoref{tab:budgetLinestyles}; (d) streamwise boundary layer thickness and momentum thickness development normalised with reference boundary layer thickness $\delta_0$.}
\label{fig:R100deg30Re1200Stats}
\end{figure}

In \autoref{fig:R100deg30Re1200U} the velocity profile in viscous scales is depicted.
It can be seen that the base and optimised configuration show no difference in their velocity profiles.
Both also match the reference profile.
As already discussed in the previous shown velocity profiles, the only difference at the upper end of the boundary layer occurs because of the different extraction positions, respectively varying values of $\mathrm{Re}_{\theta}$. 
In addition, it can be seen that the profiles follow the analytical solutions of the \benennung{linear law} and the \benennung{log law}.
This indicates a correct transformation of the velocity components from fix $xyz$-coordinate system into streamwise and wall-normal components. 
Since the components in \name{Neko} are defined in a fixed coordinate system, all quantities are transformed in the post-processing.
The correct implementation can be justified by the velocity profiles since in transformed directions the qualitative shape in wall-normal direction should not change in the curvature.
In \autoref{fig:R100deg30Re1200Fluct} the Reynolds stress components are depicted.
Similar to the velocity profiles, the differences between the optimised and the base configuration are negligible.
Also compared to the reference data only small variations in the outer region of the \acrshort{tbl} are visible.
As a last local statistics we focus on the budgets of the \acrshort{tke} in \autoref{fig:R100deg30Re1200Budget}.
It can be seen that the trends of the terms show no significant differences to the ones taken at $\mathrm{Re}_{\theta} = 700$ (see \autoref{sec:flatTBLVerif}).
The trends of all separate terms look quite similar compared to the ones in the straight section. 
In addition, both configurations show similar behaviour.
The analysed statistics allow some first suggestions about the \acrshort{tbl} behaviour in the curved section:
First, all profiles show a high qualitative accordance to the profiles in the straight part of this domain and respectively with the ones from the flat \acrshort{tbl}.
This implies a small impact of the curvature to the local statistics at this explicit position.
\citet{muck_effect_1985} already made similar findings for mild convex curvatures.
However, this suggestion is very sensible to the domain geometry and the relative position of the extracted profile and thus cannot be generalised \citep{patel_longitudinal_1997}.


To analyse the global evolution of the \acrshort{tbl}, the normalised boundary layer and momentum thicknesses are displayed in \autoref{fig:R100deg30Re1200Thickness}.
Their growth is clearly impacted by the curvature.
Up to $s/\delta_0 \approx 55$ both quantities constantly grow with increasing streamwise length.
Then, a sudden drop of both is remarkable. 
This position coincides with the start of the curvature in the domain which implies that the geometry change impacts the boundary layer growth.
After a small period in streamwise direction, the profiles recover and increase again. 
However, both quantities show a reduced slope such that their growth is decelerated in the curvature compared to the straight section.
This phenomenon is known and was already detected by \citet{gillis_turbulent_1983}.
In the curvature the growth rate stays constant until the end of the curvature is reached.
A sudden change of both quantities is remarkable at $s/\delta_0 \approx 110$ which marks the end of the curvature.
Reversely to the start of the curvature, here a sudden increase of the growth rate can be detected in all quantities.
Similar, this region is quite short in streamwise direction until the slope levels at a comparable niveau as in the curvature.
This observation forces the statement that the curvature effects the growth of a boundary layer.
Besides this general behaviour for both configurations, some local differences between them can be detected.
In the straight part the configurations perfectly match each other. 
Right before the curvature starts, a small negative bump is visible for the optimised configuration.
The following drop of the thickness measures is similar in the configurations.
After the recovery of the drop, the profiles start to differ.
It can be seen that for both measures $\delta_{99}$ and $\theta$ the optimised profile develops with a higher growth rate than the base profile.
Additionally, in the optimised configuration the change in growth rate due to the end of the curvature is less compared to the based profile.
However, the optimised configuration appears to be affected earlier by the end of the curvature.

A similar observation was made in their results of \citet{appelbaum_systematic_2025}.
The authors used a different approach, as we described in \autoref{sec:curvedTheory}, to manipulate the pressure field in a concave, curved domain.
Therefore, they observed converse adaptations of the boundary layer growth. 
This finding was reflected in the \acrshort{tbl}, which showed an increased growth rate at the start of the curvature. 
Furthermore, the authors detected varying slopes for the growth rate depending on the pressure configuration.
These effects can be compared to the behaviour of the R100deg30 curved \acrshort{tbl} under study at the transition at the end of the curvature to the straight outlet where a similar effect can be observed.
This observation can be taken as an indication that the sensitivity of the boundary layer growth to the configuration may be driven more by the change of curvature itself rather than the specific direction of the transition. 
To strengthen this assumption, a detailed analysis within one geometrical setup would be required.

\begin{figure}[h!]
\centering\begin{subfigure}[t]{0.495\textwidth}
    \centering
    \subcaption{}
    %% Creator: Matplotlib, PGF backend
%%
%% To include the figure in your LaTeX document, write
%%   \input{<filename>.pgf}
%%
%% Make sure the required packages are loaded in your preamble
%%   \usepackage{pgf}
%%
%% Also ensure that all the required font packages are loaded; for instance,
%% the lmodern package is sometimes necessary when using math font.
%%   \usepackage{lmodern}
%%
%% Figures using additional raster images can only be included by \input if
%% they are in the same directory as the main LaTeX file. For loading figures
%% from other directories you can use the `import` package
%%   \usepackage{import}
%%
%% and then include the figures with
%%   \import{<path to file>}{<filename>.pgf}
%%
%% Matplotlib used the following preamble
%%   \def\mathdefault#1{#1}
%%   \everymath=\expandafter{\the\everymath\displaystyle}
%%   \IfFileExists{scrextend.sty}{
%%     \usepackage[fontsize=11.000000pt]{scrextend}
%%   }{
%%     \renewcommand{\normalsize}{\fontsize{11.000000}{13.200000}\selectfont}
%%     \normalsize
%%   }
%%   
%%   \makeatletter\@ifpackageloaded{underscore}{}{\usepackage[strings]{underscore}}\makeatother
%%
\begingroup%
\makeatletter%
\begin{pgfpicture}%
\pgfpathrectangle{\pgfpointorigin}{\pgfqpoint{3.118110in}{2.494488in}}%
\pgfusepath{use as bounding box, clip}%
\begin{pgfscope}%
\pgfsetbuttcap%
\pgfsetmiterjoin%
\definecolor{currentfill}{rgb}{1.000000,1.000000,1.000000}%
\pgfsetfillcolor{currentfill}%
\pgfsetlinewidth{0.000000pt}%
\definecolor{currentstroke}{rgb}{1.000000,1.000000,1.000000}%
\pgfsetstrokecolor{currentstroke}%
\pgfsetdash{}{0pt}%
\pgfpathmoveto{\pgfqpoint{0.000000in}{0.000000in}}%
\pgfpathlineto{\pgfqpoint{3.118110in}{0.000000in}}%
\pgfpathlineto{\pgfqpoint{3.118110in}{2.494488in}}%
\pgfpathlineto{\pgfqpoint{0.000000in}{2.494488in}}%
\pgfpathlineto{\pgfqpoint{0.000000in}{0.000000in}}%
\pgfpathclose%
\pgfusepath{fill}%
\end{pgfscope}%
\begin{pgfscope}%
\pgfsetbuttcap%
\pgfsetmiterjoin%
\definecolor{currentfill}{rgb}{1.000000,1.000000,1.000000}%
\pgfsetfillcolor{currentfill}%
\pgfsetlinewidth{0.000000pt}%
\definecolor{currentstroke}{rgb}{0.000000,0.000000,0.000000}%
\pgfsetstrokecolor{currentstroke}%
\pgfsetstrokeopacity{0.000000}%
\pgfsetdash{}{0pt}%
\pgfpathmoveto{\pgfqpoint{0.650000in}{0.420000in}}%
\pgfpathlineto{\pgfqpoint{3.038110in}{0.420000in}}%
\pgfpathlineto{\pgfqpoint{3.038110in}{2.414488in}}%
\pgfpathlineto{\pgfqpoint{0.650000in}{2.414488in}}%
\pgfpathlineto{\pgfqpoint{0.650000in}{0.420000in}}%
\pgfpathclose%
\pgfusepath{fill}%
\end{pgfscope}%
\begin{pgfscope}%
\pgfsetbuttcap%
\pgfsetroundjoin%
\definecolor{currentfill}{rgb}{0.000000,0.000000,0.000000}%
\pgfsetfillcolor{currentfill}%
\pgfsetlinewidth{0.501875pt}%
\definecolor{currentstroke}{rgb}{0.000000,0.000000,0.000000}%
\pgfsetstrokecolor{currentstroke}%
\pgfsetdash{}{0pt}%
\pgfsys@defobject{currentmarker}{\pgfqpoint{0.000000in}{0.000000in}}{\pgfqpoint{0.000000in}{0.041667in}}{%
\pgfpathmoveto{\pgfqpoint{0.000000in}{0.000000in}}%
\pgfpathlineto{\pgfqpoint{0.000000in}{0.041667in}}%
\pgfusepath{stroke,fill}%
}%
\begin{pgfscope}%
\pgfsys@transformshift{0.650000in}{0.420000in}%
\pgfsys@useobject{currentmarker}{}%
\end{pgfscope}%
\end{pgfscope}%
\begin{pgfscope}%
\pgfsetbuttcap%
\pgfsetroundjoin%
\definecolor{currentfill}{rgb}{0.000000,0.000000,0.000000}%
\pgfsetfillcolor{currentfill}%
\pgfsetlinewidth{0.501875pt}%
\definecolor{currentstroke}{rgb}{0.000000,0.000000,0.000000}%
\pgfsetstrokecolor{currentstroke}%
\pgfsetdash{}{0pt}%
\pgfsys@defobject{currentmarker}{\pgfqpoint{0.000000in}{-0.041667in}}{\pgfqpoint{0.000000in}{0.000000in}}{%
\pgfpathmoveto{\pgfqpoint{0.000000in}{0.000000in}}%
\pgfpathlineto{\pgfqpoint{0.000000in}{-0.041667in}}%
\pgfusepath{stroke,fill}%
}%
\begin{pgfscope}%
\pgfsys@transformshift{0.650000in}{2.414488in}%
\pgfsys@useobject{currentmarker}{}%
\end{pgfscope}%
\end{pgfscope}%
\begin{pgfscope}%
\definecolor{textcolor}{rgb}{0.000000,0.000000,0.000000}%
\pgfsetstrokecolor{textcolor}%
\pgfsetfillcolor{textcolor}%
\pgftext[x=0.650000in,y=0.371389in,,top]{\color{textcolor}{\rmfamily\fontsize{11.000000}{13.200000}\selectfont\catcode`\^=\active\def^{\ifmmode\sp\else\^{}\fi}\catcode`\%=\active\def%{\%}0}}%
\end{pgfscope}%
\begin{pgfscope}%
\pgfsetbuttcap%
\pgfsetroundjoin%
\definecolor{currentfill}{rgb}{0.000000,0.000000,0.000000}%
\pgfsetfillcolor{currentfill}%
\pgfsetlinewidth{0.501875pt}%
\definecolor{currentstroke}{rgb}{0.000000,0.000000,0.000000}%
\pgfsetstrokecolor{currentstroke}%
\pgfsetdash{}{0pt}%
\pgfsys@defobject{currentmarker}{\pgfqpoint{0.000000in}{0.000000in}}{\pgfqpoint{0.000000in}{0.041667in}}{%
\pgfpathmoveto{\pgfqpoint{0.000000in}{0.000000in}}%
\pgfpathlineto{\pgfqpoint{0.000000in}{0.041667in}}%
\pgfusepath{stroke,fill}%
}%
\begin{pgfscope}%
\pgfsys@transformshift{1.492631in}{0.420000in}%
\pgfsys@useobject{currentmarker}{}%
\end{pgfscope}%
\end{pgfscope}%
\begin{pgfscope}%
\pgfsetbuttcap%
\pgfsetroundjoin%
\definecolor{currentfill}{rgb}{0.000000,0.000000,0.000000}%
\pgfsetfillcolor{currentfill}%
\pgfsetlinewidth{0.501875pt}%
\definecolor{currentstroke}{rgb}{0.000000,0.000000,0.000000}%
\pgfsetstrokecolor{currentstroke}%
\pgfsetdash{}{0pt}%
\pgfsys@defobject{currentmarker}{\pgfqpoint{0.000000in}{-0.041667in}}{\pgfqpoint{0.000000in}{0.000000in}}{%
\pgfpathmoveto{\pgfqpoint{0.000000in}{0.000000in}}%
\pgfpathlineto{\pgfqpoint{0.000000in}{-0.041667in}}%
\pgfusepath{stroke,fill}%
}%
\begin{pgfscope}%
\pgfsys@transformshift{1.492631in}{2.414488in}%
\pgfsys@useobject{currentmarker}{}%
\end{pgfscope}%
\end{pgfscope}%
\begin{pgfscope}%
\definecolor{textcolor}{rgb}{0.000000,0.000000,0.000000}%
\pgfsetstrokecolor{textcolor}%
\pgfsetfillcolor{textcolor}%
\pgftext[x=1.492631in,y=0.371389in,,top]{\color{textcolor}{\rmfamily\fontsize{11.000000}{13.200000}\selectfont\catcode`\^=\active\def^{\ifmmode\sp\else\^{}\fi}\catcode`\%=\active\def%{\%}50}}%
\end{pgfscope}%
\begin{pgfscope}%
\pgfsetbuttcap%
\pgfsetroundjoin%
\definecolor{currentfill}{rgb}{0.000000,0.000000,0.000000}%
\pgfsetfillcolor{currentfill}%
\pgfsetlinewidth{0.501875pt}%
\definecolor{currentstroke}{rgb}{0.000000,0.000000,0.000000}%
\pgfsetstrokecolor{currentstroke}%
\pgfsetdash{}{0pt}%
\pgfsys@defobject{currentmarker}{\pgfqpoint{0.000000in}{0.000000in}}{\pgfqpoint{0.000000in}{0.041667in}}{%
\pgfpathmoveto{\pgfqpoint{0.000000in}{0.000000in}}%
\pgfpathlineto{\pgfqpoint{0.000000in}{0.041667in}}%
\pgfusepath{stroke,fill}%
}%
\begin{pgfscope}%
\pgfsys@transformshift{2.335263in}{0.420000in}%
\pgfsys@useobject{currentmarker}{}%
\end{pgfscope}%
\end{pgfscope}%
\begin{pgfscope}%
\pgfsetbuttcap%
\pgfsetroundjoin%
\definecolor{currentfill}{rgb}{0.000000,0.000000,0.000000}%
\pgfsetfillcolor{currentfill}%
\pgfsetlinewidth{0.501875pt}%
\definecolor{currentstroke}{rgb}{0.000000,0.000000,0.000000}%
\pgfsetstrokecolor{currentstroke}%
\pgfsetdash{}{0pt}%
\pgfsys@defobject{currentmarker}{\pgfqpoint{0.000000in}{-0.041667in}}{\pgfqpoint{0.000000in}{0.000000in}}{%
\pgfpathmoveto{\pgfqpoint{0.000000in}{0.000000in}}%
\pgfpathlineto{\pgfqpoint{0.000000in}{-0.041667in}}%
\pgfusepath{stroke,fill}%
}%
\begin{pgfscope}%
\pgfsys@transformshift{2.335263in}{2.414488in}%
\pgfsys@useobject{currentmarker}{}%
\end{pgfscope}%
\end{pgfscope}%
\begin{pgfscope}%
\definecolor{textcolor}{rgb}{0.000000,0.000000,0.000000}%
\pgfsetstrokecolor{textcolor}%
\pgfsetfillcolor{textcolor}%
\pgftext[x=2.335263in,y=0.371389in,,top]{\color{textcolor}{\rmfamily\fontsize{11.000000}{13.200000}\selectfont\catcode`\^=\active\def^{\ifmmode\sp\else\^{}\fi}\catcode`\%=\active\def%{\%}100}}%
\end{pgfscope}%
\begin{pgfscope}%
\pgfsetbuttcap%
\pgfsetroundjoin%
\definecolor{currentfill}{rgb}{0.000000,0.000000,0.000000}%
\pgfsetfillcolor{currentfill}%
\pgfsetlinewidth{0.401500pt}%
\definecolor{currentstroke}{rgb}{0.000000,0.000000,0.000000}%
\pgfsetstrokecolor{currentstroke}%
\pgfsetdash{}{0pt}%
\pgfsys@defobject{currentmarker}{\pgfqpoint{0.000000in}{0.000000in}}{\pgfqpoint{0.000000in}{0.020833in}}{%
\pgfpathmoveto{\pgfqpoint{0.000000in}{0.000000in}}%
\pgfpathlineto{\pgfqpoint{0.000000in}{0.020833in}}%
\pgfusepath{stroke,fill}%
}%
\begin{pgfscope}%
\pgfsys@transformshift{0.818526in}{0.420000in}%
\pgfsys@useobject{currentmarker}{}%
\end{pgfscope}%
\end{pgfscope}%
\begin{pgfscope}%
\pgfsetbuttcap%
\pgfsetroundjoin%
\definecolor{currentfill}{rgb}{0.000000,0.000000,0.000000}%
\pgfsetfillcolor{currentfill}%
\pgfsetlinewidth{0.401500pt}%
\definecolor{currentstroke}{rgb}{0.000000,0.000000,0.000000}%
\pgfsetstrokecolor{currentstroke}%
\pgfsetdash{}{0pt}%
\pgfsys@defobject{currentmarker}{\pgfqpoint{0.000000in}{-0.020833in}}{\pgfqpoint{0.000000in}{0.000000in}}{%
\pgfpathmoveto{\pgfqpoint{0.000000in}{0.000000in}}%
\pgfpathlineto{\pgfqpoint{0.000000in}{-0.020833in}}%
\pgfusepath{stroke,fill}%
}%
\begin{pgfscope}%
\pgfsys@transformshift{0.818526in}{2.414488in}%
\pgfsys@useobject{currentmarker}{}%
\end{pgfscope}%
\end{pgfscope}%
\begin{pgfscope}%
\pgfsetbuttcap%
\pgfsetroundjoin%
\definecolor{currentfill}{rgb}{0.000000,0.000000,0.000000}%
\pgfsetfillcolor{currentfill}%
\pgfsetlinewidth{0.401500pt}%
\definecolor{currentstroke}{rgb}{0.000000,0.000000,0.000000}%
\pgfsetstrokecolor{currentstroke}%
\pgfsetdash{}{0pt}%
\pgfsys@defobject{currentmarker}{\pgfqpoint{0.000000in}{0.000000in}}{\pgfqpoint{0.000000in}{0.020833in}}{%
\pgfpathmoveto{\pgfqpoint{0.000000in}{0.000000in}}%
\pgfpathlineto{\pgfqpoint{0.000000in}{0.020833in}}%
\pgfusepath{stroke,fill}%
}%
\begin{pgfscope}%
\pgfsys@transformshift{0.987053in}{0.420000in}%
\pgfsys@useobject{currentmarker}{}%
\end{pgfscope}%
\end{pgfscope}%
\begin{pgfscope}%
\pgfsetbuttcap%
\pgfsetroundjoin%
\definecolor{currentfill}{rgb}{0.000000,0.000000,0.000000}%
\pgfsetfillcolor{currentfill}%
\pgfsetlinewidth{0.401500pt}%
\definecolor{currentstroke}{rgb}{0.000000,0.000000,0.000000}%
\pgfsetstrokecolor{currentstroke}%
\pgfsetdash{}{0pt}%
\pgfsys@defobject{currentmarker}{\pgfqpoint{0.000000in}{-0.020833in}}{\pgfqpoint{0.000000in}{0.000000in}}{%
\pgfpathmoveto{\pgfqpoint{0.000000in}{0.000000in}}%
\pgfpathlineto{\pgfqpoint{0.000000in}{-0.020833in}}%
\pgfusepath{stroke,fill}%
}%
\begin{pgfscope}%
\pgfsys@transformshift{0.987053in}{2.414488in}%
\pgfsys@useobject{currentmarker}{}%
\end{pgfscope}%
\end{pgfscope}%
\begin{pgfscope}%
\pgfsetbuttcap%
\pgfsetroundjoin%
\definecolor{currentfill}{rgb}{0.000000,0.000000,0.000000}%
\pgfsetfillcolor{currentfill}%
\pgfsetlinewidth{0.401500pt}%
\definecolor{currentstroke}{rgb}{0.000000,0.000000,0.000000}%
\pgfsetstrokecolor{currentstroke}%
\pgfsetdash{}{0pt}%
\pgfsys@defobject{currentmarker}{\pgfqpoint{0.000000in}{0.000000in}}{\pgfqpoint{0.000000in}{0.020833in}}{%
\pgfpathmoveto{\pgfqpoint{0.000000in}{0.000000in}}%
\pgfpathlineto{\pgfqpoint{0.000000in}{0.020833in}}%
\pgfusepath{stroke,fill}%
}%
\begin{pgfscope}%
\pgfsys@transformshift{1.155579in}{0.420000in}%
\pgfsys@useobject{currentmarker}{}%
\end{pgfscope}%
\end{pgfscope}%
\begin{pgfscope}%
\pgfsetbuttcap%
\pgfsetroundjoin%
\definecolor{currentfill}{rgb}{0.000000,0.000000,0.000000}%
\pgfsetfillcolor{currentfill}%
\pgfsetlinewidth{0.401500pt}%
\definecolor{currentstroke}{rgb}{0.000000,0.000000,0.000000}%
\pgfsetstrokecolor{currentstroke}%
\pgfsetdash{}{0pt}%
\pgfsys@defobject{currentmarker}{\pgfqpoint{0.000000in}{-0.020833in}}{\pgfqpoint{0.000000in}{0.000000in}}{%
\pgfpathmoveto{\pgfqpoint{0.000000in}{0.000000in}}%
\pgfpathlineto{\pgfqpoint{0.000000in}{-0.020833in}}%
\pgfusepath{stroke,fill}%
}%
\begin{pgfscope}%
\pgfsys@transformshift{1.155579in}{2.414488in}%
\pgfsys@useobject{currentmarker}{}%
\end{pgfscope}%
\end{pgfscope}%
\begin{pgfscope}%
\pgfsetbuttcap%
\pgfsetroundjoin%
\definecolor{currentfill}{rgb}{0.000000,0.000000,0.000000}%
\pgfsetfillcolor{currentfill}%
\pgfsetlinewidth{0.401500pt}%
\definecolor{currentstroke}{rgb}{0.000000,0.000000,0.000000}%
\pgfsetstrokecolor{currentstroke}%
\pgfsetdash{}{0pt}%
\pgfsys@defobject{currentmarker}{\pgfqpoint{0.000000in}{0.000000in}}{\pgfqpoint{0.000000in}{0.020833in}}{%
\pgfpathmoveto{\pgfqpoint{0.000000in}{0.000000in}}%
\pgfpathlineto{\pgfqpoint{0.000000in}{0.020833in}}%
\pgfusepath{stroke,fill}%
}%
\begin{pgfscope}%
\pgfsys@transformshift{1.324105in}{0.420000in}%
\pgfsys@useobject{currentmarker}{}%
\end{pgfscope}%
\end{pgfscope}%
\begin{pgfscope}%
\pgfsetbuttcap%
\pgfsetroundjoin%
\definecolor{currentfill}{rgb}{0.000000,0.000000,0.000000}%
\pgfsetfillcolor{currentfill}%
\pgfsetlinewidth{0.401500pt}%
\definecolor{currentstroke}{rgb}{0.000000,0.000000,0.000000}%
\pgfsetstrokecolor{currentstroke}%
\pgfsetdash{}{0pt}%
\pgfsys@defobject{currentmarker}{\pgfqpoint{0.000000in}{-0.020833in}}{\pgfqpoint{0.000000in}{0.000000in}}{%
\pgfpathmoveto{\pgfqpoint{0.000000in}{0.000000in}}%
\pgfpathlineto{\pgfqpoint{0.000000in}{-0.020833in}}%
\pgfusepath{stroke,fill}%
}%
\begin{pgfscope}%
\pgfsys@transformshift{1.324105in}{2.414488in}%
\pgfsys@useobject{currentmarker}{}%
\end{pgfscope}%
\end{pgfscope}%
\begin{pgfscope}%
\pgfsetbuttcap%
\pgfsetroundjoin%
\definecolor{currentfill}{rgb}{0.000000,0.000000,0.000000}%
\pgfsetfillcolor{currentfill}%
\pgfsetlinewidth{0.401500pt}%
\definecolor{currentstroke}{rgb}{0.000000,0.000000,0.000000}%
\pgfsetstrokecolor{currentstroke}%
\pgfsetdash{}{0pt}%
\pgfsys@defobject{currentmarker}{\pgfqpoint{0.000000in}{0.000000in}}{\pgfqpoint{0.000000in}{0.020833in}}{%
\pgfpathmoveto{\pgfqpoint{0.000000in}{0.000000in}}%
\pgfpathlineto{\pgfqpoint{0.000000in}{0.020833in}}%
\pgfusepath{stroke,fill}%
}%
\begin{pgfscope}%
\pgfsys@transformshift{1.661158in}{0.420000in}%
\pgfsys@useobject{currentmarker}{}%
\end{pgfscope}%
\end{pgfscope}%
\begin{pgfscope}%
\pgfsetbuttcap%
\pgfsetroundjoin%
\definecolor{currentfill}{rgb}{0.000000,0.000000,0.000000}%
\pgfsetfillcolor{currentfill}%
\pgfsetlinewidth{0.401500pt}%
\definecolor{currentstroke}{rgb}{0.000000,0.000000,0.000000}%
\pgfsetstrokecolor{currentstroke}%
\pgfsetdash{}{0pt}%
\pgfsys@defobject{currentmarker}{\pgfqpoint{0.000000in}{-0.020833in}}{\pgfqpoint{0.000000in}{0.000000in}}{%
\pgfpathmoveto{\pgfqpoint{0.000000in}{0.000000in}}%
\pgfpathlineto{\pgfqpoint{0.000000in}{-0.020833in}}%
\pgfusepath{stroke,fill}%
}%
\begin{pgfscope}%
\pgfsys@transformshift{1.661158in}{2.414488in}%
\pgfsys@useobject{currentmarker}{}%
\end{pgfscope}%
\end{pgfscope}%
\begin{pgfscope}%
\pgfsetbuttcap%
\pgfsetroundjoin%
\definecolor{currentfill}{rgb}{0.000000,0.000000,0.000000}%
\pgfsetfillcolor{currentfill}%
\pgfsetlinewidth{0.401500pt}%
\definecolor{currentstroke}{rgb}{0.000000,0.000000,0.000000}%
\pgfsetstrokecolor{currentstroke}%
\pgfsetdash{}{0pt}%
\pgfsys@defobject{currentmarker}{\pgfqpoint{0.000000in}{0.000000in}}{\pgfqpoint{0.000000in}{0.020833in}}{%
\pgfpathmoveto{\pgfqpoint{0.000000in}{0.000000in}}%
\pgfpathlineto{\pgfqpoint{0.000000in}{0.020833in}}%
\pgfusepath{stroke,fill}%
}%
\begin{pgfscope}%
\pgfsys@transformshift{1.829684in}{0.420000in}%
\pgfsys@useobject{currentmarker}{}%
\end{pgfscope}%
\end{pgfscope}%
\begin{pgfscope}%
\pgfsetbuttcap%
\pgfsetroundjoin%
\definecolor{currentfill}{rgb}{0.000000,0.000000,0.000000}%
\pgfsetfillcolor{currentfill}%
\pgfsetlinewidth{0.401500pt}%
\definecolor{currentstroke}{rgb}{0.000000,0.000000,0.000000}%
\pgfsetstrokecolor{currentstroke}%
\pgfsetdash{}{0pt}%
\pgfsys@defobject{currentmarker}{\pgfqpoint{0.000000in}{-0.020833in}}{\pgfqpoint{0.000000in}{0.000000in}}{%
\pgfpathmoveto{\pgfqpoint{0.000000in}{0.000000in}}%
\pgfpathlineto{\pgfqpoint{0.000000in}{-0.020833in}}%
\pgfusepath{stroke,fill}%
}%
\begin{pgfscope}%
\pgfsys@transformshift{1.829684in}{2.414488in}%
\pgfsys@useobject{currentmarker}{}%
\end{pgfscope}%
\end{pgfscope}%
\begin{pgfscope}%
\pgfsetbuttcap%
\pgfsetroundjoin%
\definecolor{currentfill}{rgb}{0.000000,0.000000,0.000000}%
\pgfsetfillcolor{currentfill}%
\pgfsetlinewidth{0.401500pt}%
\definecolor{currentstroke}{rgb}{0.000000,0.000000,0.000000}%
\pgfsetstrokecolor{currentstroke}%
\pgfsetdash{}{0pt}%
\pgfsys@defobject{currentmarker}{\pgfqpoint{0.000000in}{0.000000in}}{\pgfqpoint{0.000000in}{0.020833in}}{%
\pgfpathmoveto{\pgfqpoint{0.000000in}{0.000000in}}%
\pgfpathlineto{\pgfqpoint{0.000000in}{0.020833in}}%
\pgfusepath{stroke,fill}%
}%
\begin{pgfscope}%
\pgfsys@transformshift{1.998210in}{0.420000in}%
\pgfsys@useobject{currentmarker}{}%
\end{pgfscope}%
\end{pgfscope}%
\begin{pgfscope}%
\pgfsetbuttcap%
\pgfsetroundjoin%
\definecolor{currentfill}{rgb}{0.000000,0.000000,0.000000}%
\pgfsetfillcolor{currentfill}%
\pgfsetlinewidth{0.401500pt}%
\definecolor{currentstroke}{rgb}{0.000000,0.000000,0.000000}%
\pgfsetstrokecolor{currentstroke}%
\pgfsetdash{}{0pt}%
\pgfsys@defobject{currentmarker}{\pgfqpoint{0.000000in}{-0.020833in}}{\pgfqpoint{0.000000in}{0.000000in}}{%
\pgfpathmoveto{\pgfqpoint{0.000000in}{0.000000in}}%
\pgfpathlineto{\pgfqpoint{0.000000in}{-0.020833in}}%
\pgfusepath{stroke,fill}%
}%
\begin{pgfscope}%
\pgfsys@transformshift{1.998210in}{2.414488in}%
\pgfsys@useobject{currentmarker}{}%
\end{pgfscope}%
\end{pgfscope}%
\begin{pgfscope}%
\pgfsetbuttcap%
\pgfsetroundjoin%
\definecolor{currentfill}{rgb}{0.000000,0.000000,0.000000}%
\pgfsetfillcolor{currentfill}%
\pgfsetlinewidth{0.401500pt}%
\definecolor{currentstroke}{rgb}{0.000000,0.000000,0.000000}%
\pgfsetstrokecolor{currentstroke}%
\pgfsetdash{}{0pt}%
\pgfsys@defobject{currentmarker}{\pgfqpoint{0.000000in}{0.000000in}}{\pgfqpoint{0.000000in}{0.020833in}}{%
\pgfpathmoveto{\pgfqpoint{0.000000in}{0.000000in}}%
\pgfpathlineto{\pgfqpoint{0.000000in}{0.020833in}}%
\pgfusepath{stroke,fill}%
}%
\begin{pgfscope}%
\pgfsys@transformshift{2.166737in}{0.420000in}%
\pgfsys@useobject{currentmarker}{}%
\end{pgfscope}%
\end{pgfscope}%
\begin{pgfscope}%
\pgfsetbuttcap%
\pgfsetroundjoin%
\definecolor{currentfill}{rgb}{0.000000,0.000000,0.000000}%
\pgfsetfillcolor{currentfill}%
\pgfsetlinewidth{0.401500pt}%
\definecolor{currentstroke}{rgb}{0.000000,0.000000,0.000000}%
\pgfsetstrokecolor{currentstroke}%
\pgfsetdash{}{0pt}%
\pgfsys@defobject{currentmarker}{\pgfqpoint{0.000000in}{-0.020833in}}{\pgfqpoint{0.000000in}{0.000000in}}{%
\pgfpathmoveto{\pgfqpoint{0.000000in}{0.000000in}}%
\pgfpathlineto{\pgfqpoint{0.000000in}{-0.020833in}}%
\pgfusepath{stroke,fill}%
}%
\begin{pgfscope}%
\pgfsys@transformshift{2.166737in}{2.414488in}%
\pgfsys@useobject{currentmarker}{}%
\end{pgfscope}%
\end{pgfscope}%
\begin{pgfscope}%
\pgfsetbuttcap%
\pgfsetroundjoin%
\definecolor{currentfill}{rgb}{0.000000,0.000000,0.000000}%
\pgfsetfillcolor{currentfill}%
\pgfsetlinewidth{0.401500pt}%
\definecolor{currentstroke}{rgb}{0.000000,0.000000,0.000000}%
\pgfsetstrokecolor{currentstroke}%
\pgfsetdash{}{0pt}%
\pgfsys@defobject{currentmarker}{\pgfqpoint{0.000000in}{0.000000in}}{\pgfqpoint{0.000000in}{0.020833in}}{%
\pgfpathmoveto{\pgfqpoint{0.000000in}{0.000000in}}%
\pgfpathlineto{\pgfqpoint{0.000000in}{0.020833in}}%
\pgfusepath{stroke,fill}%
}%
\begin{pgfscope}%
\pgfsys@transformshift{2.503789in}{0.420000in}%
\pgfsys@useobject{currentmarker}{}%
\end{pgfscope}%
\end{pgfscope}%
\begin{pgfscope}%
\pgfsetbuttcap%
\pgfsetroundjoin%
\definecolor{currentfill}{rgb}{0.000000,0.000000,0.000000}%
\pgfsetfillcolor{currentfill}%
\pgfsetlinewidth{0.401500pt}%
\definecolor{currentstroke}{rgb}{0.000000,0.000000,0.000000}%
\pgfsetstrokecolor{currentstroke}%
\pgfsetdash{}{0pt}%
\pgfsys@defobject{currentmarker}{\pgfqpoint{0.000000in}{-0.020833in}}{\pgfqpoint{0.000000in}{0.000000in}}{%
\pgfpathmoveto{\pgfqpoint{0.000000in}{0.000000in}}%
\pgfpathlineto{\pgfqpoint{0.000000in}{-0.020833in}}%
\pgfusepath{stroke,fill}%
}%
\begin{pgfscope}%
\pgfsys@transformshift{2.503789in}{2.414488in}%
\pgfsys@useobject{currentmarker}{}%
\end{pgfscope}%
\end{pgfscope}%
\begin{pgfscope}%
\pgfsetbuttcap%
\pgfsetroundjoin%
\definecolor{currentfill}{rgb}{0.000000,0.000000,0.000000}%
\pgfsetfillcolor{currentfill}%
\pgfsetlinewidth{0.401500pt}%
\definecolor{currentstroke}{rgb}{0.000000,0.000000,0.000000}%
\pgfsetstrokecolor{currentstroke}%
\pgfsetdash{}{0pt}%
\pgfsys@defobject{currentmarker}{\pgfqpoint{0.000000in}{0.000000in}}{\pgfqpoint{0.000000in}{0.020833in}}{%
\pgfpathmoveto{\pgfqpoint{0.000000in}{0.000000in}}%
\pgfpathlineto{\pgfqpoint{0.000000in}{0.020833in}}%
\pgfusepath{stroke,fill}%
}%
\begin{pgfscope}%
\pgfsys@transformshift{2.672315in}{0.420000in}%
\pgfsys@useobject{currentmarker}{}%
\end{pgfscope}%
\end{pgfscope}%
\begin{pgfscope}%
\pgfsetbuttcap%
\pgfsetroundjoin%
\definecolor{currentfill}{rgb}{0.000000,0.000000,0.000000}%
\pgfsetfillcolor{currentfill}%
\pgfsetlinewidth{0.401500pt}%
\definecolor{currentstroke}{rgb}{0.000000,0.000000,0.000000}%
\pgfsetstrokecolor{currentstroke}%
\pgfsetdash{}{0pt}%
\pgfsys@defobject{currentmarker}{\pgfqpoint{0.000000in}{-0.020833in}}{\pgfqpoint{0.000000in}{0.000000in}}{%
\pgfpathmoveto{\pgfqpoint{0.000000in}{0.000000in}}%
\pgfpathlineto{\pgfqpoint{0.000000in}{-0.020833in}}%
\pgfusepath{stroke,fill}%
}%
\begin{pgfscope}%
\pgfsys@transformshift{2.672315in}{2.414488in}%
\pgfsys@useobject{currentmarker}{}%
\end{pgfscope}%
\end{pgfscope}%
\begin{pgfscope}%
\pgfsetbuttcap%
\pgfsetroundjoin%
\definecolor{currentfill}{rgb}{0.000000,0.000000,0.000000}%
\pgfsetfillcolor{currentfill}%
\pgfsetlinewidth{0.401500pt}%
\definecolor{currentstroke}{rgb}{0.000000,0.000000,0.000000}%
\pgfsetstrokecolor{currentstroke}%
\pgfsetdash{}{0pt}%
\pgfsys@defobject{currentmarker}{\pgfqpoint{0.000000in}{0.000000in}}{\pgfqpoint{0.000000in}{0.020833in}}{%
\pgfpathmoveto{\pgfqpoint{0.000000in}{0.000000in}}%
\pgfpathlineto{\pgfqpoint{0.000000in}{0.020833in}}%
\pgfusepath{stroke,fill}%
}%
\begin{pgfscope}%
\pgfsys@transformshift{2.840842in}{0.420000in}%
\pgfsys@useobject{currentmarker}{}%
\end{pgfscope}%
\end{pgfscope}%
\begin{pgfscope}%
\pgfsetbuttcap%
\pgfsetroundjoin%
\definecolor{currentfill}{rgb}{0.000000,0.000000,0.000000}%
\pgfsetfillcolor{currentfill}%
\pgfsetlinewidth{0.401500pt}%
\definecolor{currentstroke}{rgb}{0.000000,0.000000,0.000000}%
\pgfsetstrokecolor{currentstroke}%
\pgfsetdash{}{0pt}%
\pgfsys@defobject{currentmarker}{\pgfqpoint{0.000000in}{-0.020833in}}{\pgfqpoint{0.000000in}{0.000000in}}{%
\pgfpathmoveto{\pgfqpoint{0.000000in}{0.000000in}}%
\pgfpathlineto{\pgfqpoint{0.000000in}{-0.020833in}}%
\pgfusepath{stroke,fill}%
}%
\begin{pgfscope}%
\pgfsys@transformshift{2.840842in}{2.414488in}%
\pgfsys@useobject{currentmarker}{}%
\end{pgfscope}%
\end{pgfscope}%
\begin{pgfscope}%
\pgfsetbuttcap%
\pgfsetroundjoin%
\definecolor{currentfill}{rgb}{0.000000,0.000000,0.000000}%
\pgfsetfillcolor{currentfill}%
\pgfsetlinewidth{0.401500pt}%
\definecolor{currentstroke}{rgb}{0.000000,0.000000,0.000000}%
\pgfsetstrokecolor{currentstroke}%
\pgfsetdash{}{0pt}%
\pgfsys@defobject{currentmarker}{\pgfqpoint{0.000000in}{0.000000in}}{\pgfqpoint{0.000000in}{0.020833in}}{%
\pgfpathmoveto{\pgfqpoint{0.000000in}{0.000000in}}%
\pgfpathlineto{\pgfqpoint{0.000000in}{0.020833in}}%
\pgfusepath{stroke,fill}%
}%
\begin{pgfscope}%
\pgfsys@transformshift{3.009368in}{0.420000in}%
\pgfsys@useobject{currentmarker}{}%
\end{pgfscope}%
\end{pgfscope}%
\begin{pgfscope}%
\pgfsetbuttcap%
\pgfsetroundjoin%
\definecolor{currentfill}{rgb}{0.000000,0.000000,0.000000}%
\pgfsetfillcolor{currentfill}%
\pgfsetlinewidth{0.401500pt}%
\definecolor{currentstroke}{rgb}{0.000000,0.000000,0.000000}%
\pgfsetstrokecolor{currentstroke}%
\pgfsetdash{}{0pt}%
\pgfsys@defobject{currentmarker}{\pgfqpoint{0.000000in}{-0.020833in}}{\pgfqpoint{0.000000in}{0.000000in}}{%
\pgfpathmoveto{\pgfqpoint{0.000000in}{0.000000in}}%
\pgfpathlineto{\pgfqpoint{0.000000in}{-0.020833in}}%
\pgfusepath{stroke,fill}%
}%
\begin{pgfscope}%
\pgfsys@transformshift{3.009368in}{2.414488in}%
\pgfsys@useobject{currentmarker}{}%
\end{pgfscope}%
\end{pgfscope}%
\begin{pgfscope}%
\definecolor{textcolor}{rgb}{0.000000,0.000000,0.000000}%
\pgfsetstrokecolor{textcolor}%
\pgfsetfillcolor{textcolor}%
\pgftext[x=1.844055in,y=0.208426in,,top]{\color{textcolor}{\rmfamily\fontsize{12.000000}{14.400000}\selectfont\catcode`\^=\active\def^{\ifmmode\sp\else\^{}\fi}\catcode`\%=\active\def%{\%}$s/\delta_{0}$}}%
\end{pgfscope}%
\begin{pgfscope}%
\pgfsetbuttcap%
\pgfsetroundjoin%
\definecolor{currentfill}{rgb}{0.000000,0.000000,0.000000}%
\pgfsetfillcolor{currentfill}%
\pgfsetlinewidth{0.501875pt}%
\definecolor{currentstroke}{rgb}{0.000000,0.000000,0.000000}%
\pgfsetstrokecolor{currentstroke}%
\pgfsetdash{}{0pt}%
\pgfsys@defobject{currentmarker}{\pgfqpoint{0.000000in}{0.000000in}}{\pgfqpoint{0.041667in}{0.000000in}}{%
\pgfpathmoveto{\pgfqpoint{0.000000in}{0.000000in}}%
\pgfpathlineto{\pgfqpoint{0.041667in}{0.000000in}}%
\pgfusepath{stroke,fill}%
}%
\begin{pgfscope}%
\pgfsys@transformshift{0.650000in}{0.420000in}%
\pgfsys@useobject{currentmarker}{}%
\end{pgfscope}%
\end{pgfscope}%
\begin{pgfscope}%
\pgfsetbuttcap%
\pgfsetroundjoin%
\definecolor{currentfill}{rgb}{0.000000,0.000000,0.000000}%
\pgfsetfillcolor{currentfill}%
\pgfsetlinewidth{0.501875pt}%
\definecolor{currentstroke}{rgb}{0.000000,0.000000,0.000000}%
\pgfsetstrokecolor{currentstroke}%
\pgfsetdash{}{0pt}%
\pgfsys@defobject{currentmarker}{\pgfqpoint{-0.041667in}{0.000000in}}{\pgfqpoint{-0.000000in}{0.000000in}}{%
\pgfpathmoveto{\pgfqpoint{-0.000000in}{0.000000in}}%
\pgfpathlineto{\pgfqpoint{-0.041667in}{0.000000in}}%
\pgfusepath{stroke,fill}%
}%
\begin{pgfscope}%
\pgfsys@transformshift{3.038110in}{0.420000in}%
\pgfsys@useobject{currentmarker}{}%
\end{pgfscope}%
\end{pgfscope}%
\begin{pgfscope}%
\definecolor{textcolor}{rgb}{0.000000,0.000000,0.000000}%
\pgfsetstrokecolor{textcolor}%
\pgfsetfillcolor{textcolor}%
\pgftext[x=0.407060in, y=0.367193in, left, base]{\color{textcolor}{\rmfamily\fontsize{11.000000}{13.200000}\selectfont\catcode`\^=\active\def^{\ifmmode\sp\else\^{}\fi}\catcode`\%=\active\def%{\%}\ensuremath{-}2}}%
\end{pgfscope}%
\begin{pgfscope}%
\pgfsetbuttcap%
\pgfsetroundjoin%
\definecolor{currentfill}{rgb}{0.000000,0.000000,0.000000}%
\pgfsetfillcolor{currentfill}%
\pgfsetlinewidth{0.501875pt}%
\definecolor{currentstroke}{rgb}{0.000000,0.000000,0.000000}%
\pgfsetstrokecolor{currentstroke}%
\pgfsetdash{}{0pt}%
\pgfsys@defobject{currentmarker}{\pgfqpoint{0.000000in}{0.000000in}}{\pgfqpoint{0.041667in}{0.000000in}}{%
\pgfpathmoveto{\pgfqpoint{0.000000in}{0.000000in}}%
\pgfpathlineto{\pgfqpoint{0.041667in}{0.000000in}}%
\pgfusepath{stroke,fill}%
}%
\begin{pgfscope}%
\pgfsys@transformshift{0.650000in}{0.918622in}%
\pgfsys@useobject{currentmarker}{}%
\end{pgfscope}%
\end{pgfscope}%
\begin{pgfscope}%
\pgfsetbuttcap%
\pgfsetroundjoin%
\definecolor{currentfill}{rgb}{0.000000,0.000000,0.000000}%
\pgfsetfillcolor{currentfill}%
\pgfsetlinewidth{0.501875pt}%
\definecolor{currentstroke}{rgb}{0.000000,0.000000,0.000000}%
\pgfsetstrokecolor{currentstroke}%
\pgfsetdash{}{0pt}%
\pgfsys@defobject{currentmarker}{\pgfqpoint{-0.041667in}{0.000000in}}{\pgfqpoint{-0.000000in}{0.000000in}}{%
\pgfpathmoveto{\pgfqpoint{-0.000000in}{0.000000in}}%
\pgfpathlineto{\pgfqpoint{-0.041667in}{0.000000in}}%
\pgfusepath{stroke,fill}%
}%
\begin{pgfscope}%
\pgfsys@transformshift{3.038110in}{0.918622in}%
\pgfsys@useobject{currentmarker}{}%
\end{pgfscope}%
\end{pgfscope}%
\begin{pgfscope}%
\definecolor{textcolor}{rgb}{0.000000,0.000000,0.000000}%
\pgfsetstrokecolor{textcolor}%
\pgfsetfillcolor{textcolor}%
\pgftext[x=0.407060in, y=0.865815in, left, base]{\color{textcolor}{\rmfamily\fontsize{11.000000}{13.200000}\selectfont\catcode`\^=\active\def^{\ifmmode\sp\else\^{}\fi}\catcode`\%=\active\def%{\%}\ensuremath{-}1}}%
\end{pgfscope}%
\begin{pgfscope}%
\pgfsetbuttcap%
\pgfsetroundjoin%
\definecolor{currentfill}{rgb}{0.000000,0.000000,0.000000}%
\pgfsetfillcolor{currentfill}%
\pgfsetlinewidth{0.501875pt}%
\definecolor{currentstroke}{rgb}{0.000000,0.000000,0.000000}%
\pgfsetstrokecolor{currentstroke}%
\pgfsetdash{}{0pt}%
\pgfsys@defobject{currentmarker}{\pgfqpoint{0.000000in}{0.000000in}}{\pgfqpoint{0.041667in}{0.000000in}}{%
\pgfpathmoveto{\pgfqpoint{0.000000in}{0.000000in}}%
\pgfpathlineto{\pgfqpoint{0.041667in}{0.000000in}}%
\pgfusepath{stroke,fill}%
}%
\begin{pgfscope}%
\pgfsys@transformshift{0.650000in}{1.417244in}%
\pgfsys@useobject{currentmarker}{}%
\end{pgfscope}%
\end{pgfscope}%
\begin{pgfscope}%
\pgfsetbuttcap%
\pgfsetroundjoin%
\definecolor{currentfill}{rgb}{0.000000,0.000000,0.000000}%
\pgfsetfillcolor{currentfill}%
\pgfsetlinewidth{0.501875pt}%
\definecolor{currentstroke}{rgb}{0.000000,0.000000,0.000000}%
\pgfsetstrokecolor{currentstroke}%
\pgfsetdash{}{0pt}%
\pgfsys@defobject{currentmarker}{\pgfqpoint{-0.041667in}{0.000000in}}{\pgfqpoint{-0.000000in}{0.000000in}}{%
\pgfpathmoveto{\pgfqpoint{-0.000000in}{0.000000in}}%
\pgfpathlineto{\pgfqpoint{-0.041667in}{0.000000in}}%
\pgfusepath{stroke,fill}%
}%
\begin{pgfscope}%
\pgfsys@transformshift{3.038110in}{1.417244in}%
\pgfsys@useobject{currentmarker}{}%
\end{pgfscope}%
\end{pgfscope}%
\begin{pgfscope}%
\definecolor{textcolor}{rgb}{0.000000,0.000000,0.000000}%
\pgfsetstrokecolor{textcolor}%
\pgfsetfillcolor{textcolor}%
\pgftext[x=0.525347in, y=1.364437in, left, base]{\color{textcolor}{\rmfamily\fontsize{11.000000}{13.200000}\selectfont\catcode`\^=\active\def^{\ifmmode\sp\else\^{}\fi}\catcode`\%=\active\def%{\%}0}}%
\end{pgfscope}%
\begin{pgfscope}%
\pgfsetbuttcap%
\pgfsetroundjoin%
\definecolor{currentfill}{rgb}{0.000000,0.000000,0.000000}%
\pgfsetfillcolor{currentfill}%
\pgfsetlinewidth{0.501875pt}%
\definecolor{currentstroke}{rgb}{0.000000,0.000000,0.000000}%
\pgfsetstrokecolor{currentstroke}%
\pgfsetdash{}{0pt}%
\pgfsys@defobject{currentmarker}{\pgfqpoint{0.000000in}{0.000000in}}{\pgfqpoint{0.041667in}{0.000000in}}{%
\pgfpathmoveto{\pgfqpoint{0.000000in}{0.000000in}}%
\pgfpathlineto{\pgfqpoint{0.041667in}{0.000000in}}%
\pgfusepath{stroke,fill}%
}%
\begin{pgfscope}%
\pgfsys@transformshift{0.650000in}{1.915866in}%
\pgfsys@useobject{currentmarker}{}%
\end{pgfscope}%
\end{pgfscope}%
\begin{pgfscope}%
\pgfsetbuttcap%
\pgfsetroundjoin%
\definecolor{currentfill}{rgb}{0.000000,0.000000,0.000000}%
\pgfsetfillcolor{currentfill}%
\pgfsetlinewidth{0.501875pt}%
\definecolor{currentstroke}{rgb}{0.000000,0.000000,0.000000}%
\pgfsetstrokecolor{currentstroke}%
\pgfsetdash{}{0pt}%
\pgfsys@defobject{currentmarker}{\pgfqpoint{-0.041667in}{0.000000in}}{\pgfqpoint{-0.000000in}{0.000000in}}{%
\pgfpathmoveto{\pgfqpoint{-0.000000in}{0.000000in}}%
\pgfpathlineto{\pgfqpoint{-0.041667in}{0.000000in}}%
\pgfusepath{stroke,fill}%
}%
\begin{pgfscope}%
\pgfsys@transformshift{3.038110in}{1.915866in}%
\pgfsys@useobject{currentmarker}{}%
\end{pgfscope}%
\end{pgfscope}%
\begin{pgfscope}%
\definecolor{textcolor}{rgb}{0.000000,0.000000,0.000000}%
\pgfsetstrokecolor{textcolor}%
\pgfsetfillcolor{textcolor}%
\pgftext[x=0.525347in, y=1.863059in, left, base]{\color{textcolor}{\rmfamily\fontsize{11.000000}{13.200000}\selectfont\catcode`\^=\active\def^{\ifmmode\sp\else\^{}\fi}\catcode`\%=\active\def%{\%}1}}%
\end{pgfscope}%
\begin{pgfscope}%
\pgfsetbuttcap%
\pgfsetroundjoin%
\definecolor{currentfill}{rgb}{0.000000,0.000000,0.000000}%
\pgfsetfillcolor{currentfill}%
\pgfsetlinewidth{0.501875pt}%
\definecolor{currentstroke}{rgb}{0.000000,0.000000,0.000000}%
\pgfsetstrokecolor{currentstroke}%
\pgfsetdash{}{0pt}%
\pgfsys@defobject{currentmarker}{\pgfqpoint{0.000000in}{0.000000in}}{\pgfqpoint{0.041667in}{0.000000in}}{%
\pgfpathmoveto{\pgfqpoint{0.000000in}{0.000000in}}%
\pgfpathlineto{\pgfqpoint{0.041667in}{0.000000in}}%
\pgfusepath{stroke,fill}%
}%
\begin{pgfscope}%
\pgfsys@transformshift{0.650000in}{2.414488in}%
\pgfsys@useobject{currentmarker}{}%
\end{pgfscope}%
\end{pgfscope}%
\begin{pgfscope}%
\pgfsetbuttcap%
\pgfsetroundjoin%
\definecolor{currentfill}{rgb}{0.000000,0.000000,0.000000}%
\pgfsetfillcolor{currentfill}%
\pgfsetlinewidth{0.501875pt}%
\definecolor{currentstroke}{rgb}{0.000000,0.000000,0.000000}%
\pgfsetstrokecolor{currentstroke}%
\pgfsetdash{}{0pt}%
\pgfsys@defobject{currentmarker}{\pgfqpoint{-0.041667in}{0.000000in}}{\pgfqpoint{-0.000000in}{0.000000in}}{%
\pgfpathmoveto{\pgfqpoint{-0.000000in}{0.000000in}}%
\pgfpathlineto{\pgfqpoint{-0.041667in}{0.000000in}}%
\pgfusepath{stroke,fill}%
}%
\begin{pgfscope}%
\pgfsys@transformshift{3.038110in}{2.414488in}%
\pgfsys@useobject{currentmarker}{}%
\end{pgfscope}%
\end{pgfscope}%
\begin{pgfscope}%
\definecolor{textcolor}{rgb}{0.000000,0.000000,0.000000}%
\pgfsetstrokecolor{textcolor}%
\pgfsetfillcolor{textcolor}%
\pgftext[x=0.525347in, y=2.361681in, left, base]{\color{textcolor}{\rmfamily\fontsize{11.000000}{13.200000}\selectfont\catcode`\^=\active\def^{\ifmmode\sp\else\^{}\fi}\catcode`\%=\active\def%{\%}2}}%
\end{pgfscope}%
\begin{pgfscope}%
\pgfsetbuttcap%
\pgfsetroundjoin%
\definecolor{currentfill}{rgb}{0.000000,0.000000,0.000000}%
\pgfsetfillcolor{currentfill}%
\pgfsetlinewidth{0.401500pt}%
\definecolor{currentstroke}{rgb}{0.000000,0.000000,0.000000}%
\pgfsetstrokecolor{currentstroke}%
\pgfsetdash{}{0pt}%
\pgfsys@defobject{currentmarker}{\pgfqpoint{0.000000in}{0.000000in}}{\pgfqpoint{0.020833in}{0.000000in}}{%
\pgfpathmoveto{\pgfqpoint{0.000000in}{0.000000in}}%
\pgfpathlineto{\pgfqpoint{0.020833in}{0.000000in}}%
\pgfusepath{stroke,fill}%
}%
\begin{pgfscope}%
\pgfsys@transformshift{0.650000in}{0.519724in}%
\pgfsys@useobject{currentmarker}{}%
\end{pgfscope}%
\end{pgfscope}%
\begin{pgfscope}%
\pgfsetbuttcap%
\pgfsetroundjoin%
\definecolor{currentfill}{rgb}{0.000000,0.000000,0.000000}%
\pgfsetfillcolor{currentfill}%
\pgfsetlinewidth{0.401500pt}%
\definecolor{currentstroke}{rgb}{0.000000,0.000000,0.000000}%
\pgfsetstrokecolor{currentstroke}%
\pgfsetdash{}{0pt}%
\pgfsys@defobject{currentmarker}{\pgfqpoint{-0.020833in}{0.000000in}}{\pgfqpoint{-0.000000in}{0.000000in}}{%
\pgfpathmoveto{\pgfqpoint{-0.000000in}{0.000000in}}%
\pgfpathlineto{\pgfqpoint{-0.020833in}{0.000000in}}%
\pgfusepath{stroke,fill}%
}%
\begin{pgfscope}%
\pgfsys@transformshift{3.038110in}{0.519724in}%
\pgfsys@useobject{currentmarker}{}%
\end{pgfscope}%
\end{pgfscope}%
\begin{pgfscope}%
\pgfsetbuttcap%
\pgfsetroundjoin%
\definecolor{currentfill}{rgb}{0.000000,0.000000,0.000000}%
\pgfsetfillcolor{currentfill}%
\pgfsetlinewidth{0.401500pt}%
\definecolor{currentstroke}{rgb}{0.000000,0.000000,0.000000}%
\pgfsetstrokecolor{currentstroke}%
\pgfsetdash{}{0pt}%
\pgfsys@defobject{currentmarker}{\pgfqpoint{0.000000in}{0.000000in}}{\pgfqpoint{0.020833in}{0.000000in}}{%
\pgfpathmoveto{\pgfqpoint{0.000000in}{0.000000in}}%
\pgfpathlineto{\pgfqpoint{0.020833in}{0.000000in}}%
\pgfusepath{stroke,fill}%
}%
\begin{pgfscope}%
\pgfsys@transformshift{0.650000in}{0.619449in}%
\pgfsys@useobject{currentmarker}{}%
\end{pgfscope}%
\end{pgfscope}%
\begin{pgfscope}%
\pgfsetbuttcap%
\pgfsetroundjoin%
\definecolor{currentfill}{rgb}{0.000000,0.000000,0.000000}%
\pgfsetfillcolor{currentfill}%
\pgfsetlinewidth{0.401500pt}%
\definecolor{currentstroke}{rgb}{0.000000,0.000000,0.000000}%
\pgfsetstrokecolor{currentstroke}%
\pgfsetdash{}{0pt}%
\pgfsys@defobject{currentmarker}{\pgfqpoint{-0.020833in}{0.000000in}}{\pgfqpoint{-0.000000in}{0.000000in}}{%
\pgfpathmoveto{\pgfqpoint{-0.000000in}{0.000000in}}%
\pgfpathlineto{\pgfqpoint{-0.020833in}{0.000000in}}%
\pgfusepath{stroke,fill}%
}%
\begin{pgfscope}%
\pgfsys@transformshift{3.038110in}{0.619449in}%
\pgfsys@useobject{currentmarker}{}%
\end{pgfscope}%
\end{pgfscope}%
\begin{pgfscope}%
\pgfsetbuttcap%
\pgfsetroundjoin%
\definecolor{currentfill}{rgb}{0.000000,0.000000,0.000000}%
\pgfsetfillcolor{currentfill}%
\pgfsetlinewidth{0.401500pt}%
\definecolor{currentstroke}{rgb}{0.000000,0.000000,0.000000}%
\pgfsetstrokecolor{currentstroke}%
\pgfsetdash{}{0pt}%
\pgfsys@defobject{currentmarker}{\pgfqpoint{0.000000in}{0.000000in}}{\pgfqpoint{0.020833in}{0.000000in}}{%
\pgfpathmoveto{\pgfqpoint{0.000000in}{0.000000in}}%
\pgfpathlineto{\pgfqpoint{0.020833in}{0.000000in}}%
\pgfusepath{stroke,fill}%
}%
\begin{pgfscope}%
\pgfsys@transformshift{0.650000in}{0.719173in}%
\pgfsys@useobject{currentmarker}{}%
\end{pgfscope}%
\end{pgfscope}%
\begin{pgfscope}%
\pgfsetbuttcap%
\pgfsetroundjoin%
\definecolor{currentfill}{rgb}{0.000000,0.000000,0.000000}%
\pgfsetfillcolor{currentfill}%
\pgfsetlinewidth{0.401500pt}%
\definecolor{currentstroke}{rgb}{0.000000,0.000000,0.000000}%
\pgfsetstrokecolor{currentstroke}%
\pgfsetdash{}{0pt}%
\pgfsys@defobject{currentmarker}{\pgfqpoint{-0.020833in}{0.000000in}}{\pgfqpoint{-0.000000in}{0.000000in}}{%
\pgfpathmoveto{\pgfqpoint{-0.000000in}{0.000000in}}%
\pgfpathlineto{\pgfqpoint{-0.020833in}{0.000000in}}%
\pgfusepath{stroke,fill}%
}%
\begin{pgfscope}%
\pgfsys@transformshift{3.038110in}{0.719173in}%
\pgfsys@useobject{currentmarker}{}%
\end{pgfscope}%
\end{pgfscope}%
\begin{pgfscope}%
\pgfsetbuttcap%
\pgfsetroundjoin%
\definecolor{currentfill}{rgb}{0.000000,0.000000,0.000000}%
\pgfsetfillcolor{currentfill}%
\pgfsetlinewidth{0.401500pt}%
\definecolor{currentstroke}{rgb}{0.000000,0.000000,0.000000}%
\pgfsetstrokecolor{currentstroke}%
\pgfsetdash{}{0pt}%
\pgfsys@defobject{currentmarker}{\pgfqpoint{0.000000in}{0.000000in}}{\pgfqpoint{0.020833in}{0.000000in}}{%
\pgfpathmoveto{\pgfqpoint{0.000000in}{0.000000in}}%
\pgfpathlineto{\pgfqpoint{0.020833in}{0.000000in}}%
\pgfusepath{stroke,fill}%
}%
\begin{pgfscope}%
\pgfsys@transformshift{0.650000in}{0.818898in}%
\pgfsys@useobject{currentmarker}{}%
\end{pgfscope}%
\end{pgfscope}%
\begin{pgfscope}%
\pgfsetbuttcap%
\pgfsetroundjoin%
\definecolor{currentfill}{rgb}{0.000000,0.000000,0.000000}%
\pgfsetfillcolor{currentfill}%
\pgfsetlinewidth{0.401500pt}%
\definecolor{currentstroke}{rgb}{0.000000,0.000000,0.000000}%
\pgfsetstrokecolor{currentstroke}%
\pgfsetdash{}{0pt}%
\pgfsys@defobject{currentmarker}{\pgfqpoint{-0.020833in}{0.000000in}}{\pgfqpoint{-0.000000in}{0.000000in}}{%
\pgfpathmoveto{\pgfqpoint{-0.000000in}{0.000000in}}%
\pgfpathlineto{\pgfqpoint{-0.020833in}{0.000000in}}%
\pgfusepath{stroke,fill}%
}%
\begin{pgfscope}%
\pgfsys@transformshift{3.038110in}{0.818898in}%
\pgfsys@useobject{currentmarker}{}%
\end{pgfscope}%
\end{pgfscope}%
\begin{pgfscope}%
\pgfsetbuttcap%
\pgfsetroundjoin%
\definecolor{currentfill}{rgb}{0.000000,0.000000,0.000000}%
\pgfsetfillcolor{currentfill}%
\pgfsetlinewidth{0.401500pt}%
\definecolor{currentstroke}{rgb}{0.000000,0.000000,0.000000}%
\pgfsetstrokecolor{currentstroke}%
\pgfsetdash{}{0pt}%
\pgfsys@defobject{currentmarker}{\pgfqpoint{0.000000in}{0.000000in}}{\pgfqpoint{0.020833in}{0.000000in}}{%
\pgfpathmoveto{\pgfqpoint{0.000000in}{0.000000in}}%
\pgfpathlineto{\pgfqpoint{0.020833in}{0.000000in}}%
\pgfusepath{stroke,fill}%
}%
\begin{pgfscope}%
\pgfsys@transformshift{0.650000in}{1.018346in}%
\pgfsys@useobject{currentmarker}{}%
\end{pgfscope}%
\end{pgfscope}%
\begin{pgfscope}%
\pgfsetbuttcap%
\pgfsetroundjoin%
\definecolor{currentfill}{rgb}{0.000000,0.000000,0.000000}%
\pgfsetfillcolor{currentfill}%
\pgfsetlinewidth{0.401500pt}%
\definecolor{currentstroke}{rgb}{0.000000,0.000000,0.000000}%
\pgfsetstrokecolor{currentstroke}%
\pgfsetdash{}{0pt}%
\pgfsys@defobject{currentmarker}{\pgfqpoint{-0.020833in}{0.000000in}}{\pgfqpoint{-0.000000in}{0.000000in}}{%
\pgfpathmoveto{\pgfqpoint{-0.000000in}{0.000000in}}%
\pgfpathlineto{\pgfqpoint{-0.020833in}{0.000000in}}%
\pgfusepath{stroke,fill}%
}%
\begin{pgfscope}%
\pgfsys@transformshift{3.038110in}{1.018346in}%
\pgfsys@useobject{currentmarker}{}%
\end{pgfscope}%
\end{pgfscope}%
\begin{pgfscope}%
\pgfsetbuttcap%
\pgfsetroundjoin%
\definecolor{currentfill}{rgb}{0.000000,0.000000,0.000000}%
\pgfsetfillcolor{currentfill}%
\pgfsetlinewidth{0.401500pt}%
\definecolor{currentstroke}{rgb}{0.000000,0.000000,0.000000}%
\pgfsetstrokecolor{currentstroke}%
\pgfsetdash{}{0pt}%
\pgfsys@defobject{currentmarker}{\pgfqpoint{0.000000in}{0.000000in}}{\pgfqpoint{0.020833in}{0.000000in}}{%
\pgfpathmoveto{\pgfqpoint{0.000000in}{0.000000in}}%
\pgfpathlineto{\pgfqpoint{0.020833in}{0.000000in}}%
\pgfusepath{stroke,fill}%
}%
\begin{pgfscope}%
\pgfsys@transformshift{0.650000in}{1.118071in}%
\pgfsys@useobject{currentmarker}{}%
\end{pgfscope}%
\end{pgfscope}%
\begin{pgfscope}%
\pgfsetbuttcap%
\pgfsetroundjoin%
\definecolor{currentfill}{rgb}{0.000000,0.000000,0.000000}%
\pgfsetfillcolor{currentfill}%
\pgfsetlinewidth{0.401500pt}%
\definecolor{currentstroke}{rgb}{0.000000,0.000000,0.000000}%
\pgfsetstrokecolor{currentstroke}%
\pgfsetdash{}{0pt}%
\pgfsys@defobject{currentmarker}{\pgfqpoint{-0.020833in}{0.000000in}}{\pgfqpoint{-0.000000in}{0.000000in}}{%
\pgfpathmoveto{\pgfqpoint{-0.000000in}{0.000000in}}%
\pgfpathlineto{\pgfqpoint{-0.020833in}{0.000000in}}%
\pgfusepath{stroke,fill}%
}%
\begin{pgfscope}%
\pgfsys@transformshift{3.038110in}{1.118071in}%
\pgfsys@useobject{currentmarker}{}%
\end{pgfscope}%
\end{pgfscope}%
\begin{pgfscope}%
\pgfsetbuttcap%
\pgfsetroundjoin%
\definecolor{currentfill}{rgb}{0.000000,0.000000,0.000000}%
\pgfsetfillcolor{currentfill}%
\pgfsetlinewidth{0.401500pt}%
\definecolor{currentstroke}{rgb}{0.000000,0.000000,0.000000}%
\pgfsetstrokecolor{currentstroke}%
\pgfsetdash{}{0pt}%
\pgfsys@defobject{currentmarker}{\pgfqpoint{0.000000in}{0.000000in}}{\pgfqpoint{0.020833in}{0.000000in}}{%
\pgfpathmoveto{\pgfqpoint{0.000000in}{0.000000in}}%
\pgfpathlineto{\pgfqpoint{0.020833in}{0.000000in}}%
\pgfusepath{stroke,fill}%
}%
\begin{pgfscope}%
\pgfsys@transformshift{0.650000in}{1.217795in}%
\pgfsys@useobject{currentmarker}{}%
\end{pgfscope}%
\end{pgfscope}%
\begin{pgfscope}%
\pgfsetbuttcap%
\pgfsetroundjoin%
\definecolor{currentfill}{rgb}{0.000000,0.000000,0.000000}%
\pgfsetfillcolor{currentfill}%
\pgfsetlinewidth{0.401500pt}%
\definecolor{currentstroke}{rgb}{0.000000,0.000000,0.000000}%
\pgfsetstrokecolor{currentstroke}%
\pgfsetdash{}{0pt}%
\pgfsys@defobject{currentmarker}{\pgfqpoint{-0.020833in}{0.000000in}}{\pgfqpoint{-0.000000in}{0.000000in}}{%
\pgfpathmoveto{\pgfqpoint{-0.000000in}{0.000000in}}%
\pgfpathlineto{\pgfqpoint{-0.020833in}{0.000000in}}%
\pgfusepath{stroke,fill}%
}%
\begin{pgfscope}%
\pgfsys@transformshift{3.038110in}{1.217795in}%
\pgfsys@useobject{currentmarker}{}%
\end{pgfscope}%
\end{pgfscope}%
\begin{pgfscope}%
\pgfsetbuttcap%
\pgfsetroundjoin%
\definecolor{currentfill}{rgb}{0.000000,0.000000,0.000000}%
\pgfsetfillcolor{currentfill}%
\pgfsetlinewidth{0.401500pt}%
\definecolor{currentstroke}{rgb}{0.000000,0.000000,0.000000}%
\pgfsetstrokecolor{currentstroke}%
\pgfsetdash{}{0pt}%
\pgfsys@defobject{currentmarker}{\pgfqpoint{0.000000in}{0.000000in}}{\pgfqpoint{0.020833in}{0.000000in}}{%
\pgfpathmoveto{\pgfqpoint{0.000000in}{0.000000in}}%
\pgfpathlineto{\pgfqpoint{0.020833in}{0.000000in}}%
\pgfusepath{stroke,fill}%
}%
\begin{pgfscope}%
\pgfsys@transformshift{0.650000in}{1.317520in}%
\pgfsys@useobject{currentmarker}{}%
\end{pgfscope}%
\end{pgfscope}%
\begin{pgfscope}%
\pgfsetbuttcap%
\pgfsetroundjoin%
\definecolor{currentfill}{rgb}{0.000000,0.000000,0.000000}%
\pgfsetfillcolor{currentfill}%
\pgfsetlinewidth{0.401500pt}%
\definecolor{currentstroke}{rgb}{0.000000,0.000000,0.000000}%
\pgfsetstrokecolor{currentstroke}%
\pgfsetdash{}{0pt}%
\pgfsys@defobject{currentmarker}{\pgfqpoint{-0.020833in}{0.000000in}}{\pgfqpoint{-0.000000in}{0.000000in}}{%
\pgfpathmoveto{\pgfqpoint{-0.000000in}{0.000000in}}%
\pgfpathlineto{\pgfqpoint{-0.020833in}{0.000000in}}%
\pgfusepath{stroke,fill}%
}%
\begin{pgfscope}%
\pgfsys@transformshift{3.038110in}{1.317520in}%
\pgfsys@useobject{currentmarker}{}%
\end{pgfscope}%
\end{pgfscope}%
\begin{pgfscope}%
\pgfsetbuttcap%
\pgfsetroundjoin%
\definecolor{currentfill}{rgb}{0.000000,0.000000,0.000000}%
\pgfsetfillcolor{currentfill}%
\pgfsetlinewidth{0.401500pt}%
\definecolor{currentstroke}{rgb}{0.000000,0.000000,0.000000}%
\pgfsetstrokecolor{currentstroke}%
\pgfsetdash{}{0pt}%
\pgfsys@defobject{currentmarker}{\pgfqpoint{0.000000in}{0.000000in}}{\pgfqpoint{0.020833in}{0.000000in}}{%
\pgfpathmoveto{\pgfqpoint{0.000000in}{0.000000in}}%
\pgfpathlineto{\pgfqpoint{0.020833in}{0.000000in}}%
\pgfusepath{stroke,fill}%
}%
\begin{pgfscope}%
\pgfsys@transformshift{0.650000in}{1.516968in}%
\pgfsys@useobject{currentmarker}{}%
\end{pgfscope}%
\end{pgfscope}%
\begin{pgfscope}%
\pgfsetbuttcap%
\pgfsetroundjoin%
\definecolor{currentfill}{rgb}{0.000000,0.000000,0.000000}%
\pgfsetfillcolor{currentfill}%
\pgfsetlinewidth{0.401500pt}%
\definecolor{currentstroke}{rgb}{0.000000,0.000000,0.000000}%
\pgfsetstrokecolor{currentstroke}%
\pgfsetdash{}{0pt}%
\pgfsys@defobject{currentmarker}{\pgfqpoint{-0.020833in}{0.000000in}}{\pgfqpoint{-0.000000in}{0.000000in}}{%
\pgfpathmoveto{\pgfqpoint{-0.000000in}{0.000000in}}%
\pgfpathlineto{\pgfqpoint{-0.020833in}{0.000000in}}%
\pgfusepath{stroke,fill}%
}%
\begin{pgfscope}%
\pgfsys@transformshift{3.038110in}{1.516968in}%
\pgfsys@useobject{currentmarker}{}%
\end{pgfscope}%
\end{pgfscope}%
\begin{pgfscope}%
\pgfsetbuttcap%
\pgfsetroundjoin%
\definecolor{currentfill}{rgb}{0.000000,0.000000,0.000000}%
\pgfsetfillcolor{currentfill}%
\pgfsetlinewidth{0.401500pt}%
\definecolor{currentstroke}{rgb}{0.000000,0.000000,0.000000}%
\pgfsetstrokecolor{currentstroke}%
\pgfsetdash{}{0pt}%
\pgfsys@defobject{currentmarker}{\pgfqpoint{0.000000in}{0.000000in}}{\pgfqpoint{0.020833in}{0.000000in}}{%
\pgfpathmoveto{\pgfqpoint{0.000000in}{0.000000in}}%
\pgfpathlineto{\pgfqpoint{0.020833in}{0.000000in}}%
\pgfusepath{stroke,fill}%
}%
\begin{pgfscope}%
\pgfsys@transformshift{0.650000in}{1.616693in}%
\pgfsys@useobject{currentmarker}{}%
\end{pgfscope}%
\end{pgfscope}%
\begin{pgfscope}%
\pgfsetbuttcap%
\pgfsetroundjoin%
\definecolor{currentfill}{rgb}{0.000000,0.000000,0.000000}%
\pgfsetfillcolor{currentfill}%
\pgfsetlinewidth{0.401500pt}%
\definecolor{currentstroke}{rgb}{0.000000,0.000000,0.000000}%
\pgfsetstrokecolor{currentstroke}%
\pgfsetdash{}{0pt}%
\pgfsys@defobject{currentmarker}{\pgfqpoint{-0.020833in}{0.000000in}}{\pgfqpoint{-0.000000in}{0.000000in}}{%
\pgfpathmoveto{\pgfqpoint{-0.000000in}{0.000000in}}%
\pgfpathlineto{\pgfqpoint{-0.020833in}{0.000000in}}%
\pgfusepath{stroke,fill}%
}%
\begin{pgfscope}%
\pgfsys@transformshift{3.038110in}{1.616693in}%
\pgfsys@useobject{currentmarker}{}%
\end{pgfscope}%
\end{pgfscope}%
\begin{pgfscope}%
\pgfsetbuttcap%
\pgfsetroundjoin%
\definecolor{currentfill}{rgb}{0.000000,0.000000,0.000000}%
\pgfsetfillcolor{currentfill}%
\pgfsetlinewidth{0.401500pt}%
\definecolor{currentstroke}{rgb}{0.000000,0.000000,0.000000}%
\pgfsetstrokecolor{currentstroke}%
\pgfsetdash{}{0pt}%
\pgfsys@defobject{currentmarker}{\pgfqpoint{0.000000in}{0.000000in}}{\pgfqpoint{0.020833in}{0.000000in}}{%
\pgfpathmoveto{\pgfqpoint{0.000000in}{0.000000in}}%
\pgfpathlineto{\pgfqpoint{0.020833in}{0.000000in}}%
\pgfusepath{stroke,fill}%
}%
\begin{pgfscope}%
\pgfsys@transformshift{0.650000in}{1.716417in}%
\pgfsys@useobject{currentmarker}{}%
\end{pgfscope}%
\end{pgfscope}%
\begin{pgfscope}%
\pgfsetbuttcap%
\pgfsetroundjoin%
\definecolor{currentfill}{rgb}{0.000000,0.000000,0.000000}%
\pgfsetfillcolor{currentfill}%
\pgfsetlinewidth{0.401500pt}%
\definecolor{currentstroke}{rgb}{0.000000,0.000000,0.000000}%
\pgfsetstrokecolor{currentstroke}%
\pgfsetdash{}{0pt}%
\pgfsys@defobject{currentmarker}{\pgfqpoint{-0.020833in}{0.000000in}}{\pgfqpoint{-0.000000in}{0.000000in}}{%
\pgfpathmoveto{\pgfqpoint{-0.000000in}{0.000000in}}%
\pgfpathlineto{\pgfqpoint{-0.020833in}{0.000000in}}%
\pgfusepath{stroke,fill}%
}%
\begin{pgfscope}%
\pgfsys@transformshift{3.038110in}{1.716417in}%
\pgfsys@useobject{currentmarker}{}%
\end{pgfscope}%
\end{pgfscope}%
\begin{pgfscope}%
\pgfsetbuttcap%
\pgfsetroundjoin%
\definecolor{currentfill}{rgb}{0.000000,0.000000,0.000000}%
\pgfsetfillcolor{currentfill}%
\pgfsetlinewidth{0.401500pt}%
\definecolor{currentstroke}{rgb}{0.000000,0.000000,0.000000}%
\pgfsetstrokecolor{currentstroke}%
\pgfsetdash{}{0pt}%
\pgfsys@defobject{currentmarker}{\pgfqpoint{0.000000in}{0.000000in}}{\pgfqpoint{0.020833in}{0.000000in}}{%
\pgfpathmoveto{\pgfqpoint{0.000000in}{0.000000in}}%
\pgfpathlineto{\pgfqpoint{0.020833in}{0.000000in}}%
\pgfusepath{stroke,fill}%
}%
\begin{pgfscope}%
\pgfsys@transformshift{0.650000in}{1.816142in}%
\pgfsys@useobject{currentmarker}{}%
\end{pgfscope}%
\end{pgfscope}%
\begin{pgfscope}%
\pgfsetbuttcap%
\pgfsetroundjoin%
\definecolor{currentfill}{rgb}{0.000000,0.000000,0.000000}%
\pgfsetfillcolor{currentfill}%
\pgfsetlinewidth{0.401500pt}%
\definecolor{currentstroke}{rgb}{0.000000,0.000000,0.000000}%
\pgfsetstrokecolor{currentstroke}%
\pgfsetdash{}{0pt}%
\pgfsys@defobject{currentmarker}{\pgfqpoint{-0.020833in}{0.000000in}}{\pgfqpoint{-0.000000in}{0.000000in}}{%
\pgfpathmoveto{\pgfqpoint{-0.000000in}{0.000000in}}%
\pgfpathlineto{\pgfqpoint{-0.020833in}{0.000000in}}%
\pgfusepath{stroke,fill}%
}%
\begin{pgfscope}%
\pgfsys@transformshift{3.038110in}{1.816142in}%
\pgfsys@useobject{currentmarker}{}%
\end{pgfscope}%
\end{pgfscope}%
\begin{pgfscope}%
\pgfsetbuttcap%
\pgfsetroundjoin%
\definecolor{currentfill}{rgb}{0.000000,0.000000,0.000000}%
\pgfsetfillcolor{currentfill}%
\pgfsetlinewidth{0.401500pt}%
\definecolor{currentstroke}{rgb}{0.000000,0.000000,0.000000}%
\pgfsetstrokecolor{currentstroke}%
\pgfsetdash{}{0pt}%
\pgfsys@defobject{currentmarker}{\pgfqpoint{0.000000in}{0.000000in}}{\pgfqpoint{0.020833in}{0.000000in}}{%
\pgfpathmoveto{\pgfqpoint{0.000000in}{0.000000in}}%
\pgfpathlineto{\pgfqpoint{0.020833in}{0.000000in}}%
\pgfusepath{stroke,fill}%
}%
\begin{pgfscope}%
\pgfsys@transformshift{0.650000in}{2.015590in}%
\pgfsys@useobject{currentmarker}{}%
\end{pgfscope}%
\end{pgfscope}%
\begin{pgfscope}%
\pgfsetbuttcap%
\pgfsetroundjoin%
\definecolor{currentfill}{rgb}{0.000000,0.000000,0.000000}%
\pgfsetfillcolor{currentfill}%
\pgfsetlinewidth{0.401500pt}%
\definecolor{currentstroke}{rgb}{0.000000,0.000000,0.000000}%
\pgfsetstrokecolor{currentstroke}%
\pgfsetdash{}{0pt}%
\pgfsys@defobject{currentmarker}{\pgfqpoint{-0.020833in}{0.000000in}}{\pgfqpoint{-0.000000in}{0.000000in}}{%
\pgfpathmoveto{\pgfqpoint{-0.000000in}{0.000000in}}%
\pgfpathlineto{\pgfqpoint{-0.020833in}{0.000000in}}%
\pgfusepath{stroke,fill}%
}%
\begin{pgfscope}%
\pgfsys@transformshift{3.038110in}{2.015590in}%
\pgfsys@useobject{currentmarker}{}%
\end{pgfscope}%
\end{pgfscope}%
\begin{pgfscope}%
\pgfsetbuttcap%
\pgfsetroundjoin%
\definecolor{currentfill}{rgb}{0.000000,0.000000,0.000000}%
\pgfsetfillcolor{currentfill}%
\pgfsetlinewidth{0.401500pt}%
\definecolor{currentstroke}{rgb}{0.000000,0.000000,0.000000}%
\pgfsetstrokecolor{currentstroke}%
\pgfsetdash{}{0pt}%
\pgfsys@defobject{currentmarker}{\pgfqpoint{0.000000in}{0.000000in}}{\pgfqpoint{0.020833in}{0.000000in}}{%
\pgfpathmoveto{\pgfqpoint{0.000000in}{0.000000in}}%
\pgfpathlineto{\pgfqpoint{0.020833in}{0.000000in}}%
\pgfusepath{stroke,fill}%
}%
\begin{pgfscope}%
\pgfsys@transformshift{0.650000in}{2.115315in}%
\pgfsys@useobject{currentmarker}{}%
\end{pgfscope}%
\end{pgfscope}%
\begin{pgfscope}%
\pgfsetbuttcap%
\pgfsetroundjoin%
\definecolor{currentfill}{rgb}{0.000000,0.000000,0.000000}%
\pgfsetfillcolor{currentfill}%
\pgfsetlinewidth{0.401500pt}%
\definecolor{currentstroke}{rgb}{0.000000,0.000000,0.000000}%
\pgfsetstrokecolor{currentstroke}%
\pgfsetdash{}{0pt}%
\pgfsys@defobject{currentmarker}{\pgfqpoint{-0.020833in}{0.000000in}}{\pgfqpoint{-0.000000in}{0.000000in}}{%
\pgfpathmoveto{\pgfqpoint{-0.000000in}{0.000000in}}%
\pgfpathlineto{\pgfqpoint{-0.020833in}{0.000000in}}%
\pgfusepath{stroke,fill}%
}%
\begin{pgfscope}%
\pgfsys@transformshift{3.038110in}{2.115315in}%
\pgfsys@useobject{currentmarker}{}%
\end{pgfscope}%
\end{pgfscope}%
\begin{pgfscope}%
\pgfsetbuttcap%
\pgfsetroundjoin%
\definecolor{currentfill}{rgb}{0.000000,0.000000,0.000000}%
\pgfsetfillcolor{currentfill}%
\pgfsetlinewidth{0.401500pt}%
\definecolor{currentstroke}{rgb}{0.000000,0.000000,0.000000}%
\pgfsetstrokecolor{currentstroke}%
\pgfsetdash{}{0pt}%
\pgfsys@defobject{currentmarker}{\pgfqpoint{0.000000in}{0.000000in}}{\pgfqpoint{0.020833in}{0.000000in}}{%
\pgfpathmoveto{\pgfqpoint{0.000000in}{0.000000in}}%
\pgfpathlineto{\pgfqpoint{0.020833in}{0.000000in}}%
\pgfusepath{stroke,fill}%
}%
\begin{pgfscope}%
\pgfsys@transformshift{0.650000in}{2.215039in}%
\pgfsys@useobject{currentmarker}{}%
\end{pgfscope}%
\end{pgfscope}%
\begin{pgfscope}%
\pgfsetbuttcap%
\pgfsetroundjoin%
\definecolor{currentfill}{rgb}{0.000000,0.000000,0.000000}%
\pgfsetfillcolor{currentfill}%
\pgfsetlinewidth{0.401500pt}%
\definecolor{currentstroke}{rgb}{0.000000,0.000000,0.000000}%
\pgfsetstrokecolor{currentstroke}%
\pgfsetdash{}{0pt}%
\pgfsys@defobject{currentmarker}{\pgfqpoint{-0.020833in}{0.000000in}}{\pgfqpoint{-0.000000in}{0.000000in}}{%
\pgfpathmoveto{\pgfqpoint{-0.000000in}{0.000000in}}%
\pgfpathlineto{\pgfqpoint{-0.020833in}{0.000000in}}%
\pgfusepath{stroke,fill}%
}%
\begin{pgfscope}%
\pgfsys@transformshift{3.038110in}{2.215039in}%
\pgfsys@useobject{currentmarker}{}%
\end{pgfscope}%
\end{pgfscope}%
\begin{pgfscope}%
\pgfsetbuttcap%
\pgfsetroundjoin%
\definecolor{currentfill}{rgb}{0.000000,0.000000,0.000000}%
\pgfsetfillcolor{currentfill}%
\pgfsetlinewidth{0.401500pt}%
\definecolor{currentstroke}{rgb}{0.000000,0.000000,0.000000}%
\pgfsetstrokecolor{currentstroke}%
\pgfsetdash{}{0pt}%
\pgfsys@defobject{currentmarker}{\pgfqpoint{0.000000in}{0.000000in}}{\pgfqpoint{0.020833in}{0.000000in}}{%
\pgfpathmoveto{\pgfqpoint{0.000000in}{0.000000in}}%
\pgfpathlineto{\pgfqpoint{0.020833in}{0.000000in}}%
\pgfusepath{stroke,fill}%
}%
\begin{pgfscope}%
\pgfsys@transformshift{0.650000in}{2.314764in}%
\pgfsys@useobject{currentmarker}{}%
\end{pgfscope}%
\end{pgfscope}%
\begin{pgfscope}%
\pgfsetbuttcap%
\pgfsetroundjoin%
\definecolor{currentfill}{rgb}{0.000000,0.000000,0.000000}%
\pgfsetfillcolor{currentfill}%
\pgfsetlinewidth{0.401500pt}%
\definecolor{currentstroke}{rgb}{0.000000,0.000000,0.000000}%
\pgfsetstrokecolor{currentstroke}%
\pgfsetdash{}{0pt}%
\pgfsys@defobject{currentmarker}{\pgfqpoint{-0.020833in}{0.000000in}}{\pgfqpoint{-0.000000in}{0.000000in}}{%
\pgfpathmoveto{\pgfqpoint{-0.000000in}{0.000000in}}%
\pgfpathlineto{\pgfqpoint{-0.020833in}{0.000000in}}%
\pgfusepath{stroke,fill}%
}%
\begin{pgfscope}%
\pgfsys@transformshift{3.038110in}{2.314764in}%
\pgfsys@useobject{currentmarker}{}%
\end{pgfscope}%
\end{pgfscope}%
\begin{pgfscope}%
\definecolor{textcolor}{rgb}{0.000000,0.000000,0.000000}%
\pgfsetstrokecolor{textcolor}%
\pgfsetfillcolor{textcolor}%
\pgftext[x=0.379282in,y=1.417244in,,bottom,rotate=90.000000]{\color{textcolor}{\rmfamily\fontsize{12.000000}{14.400000}\selectfont\catcode`\^=\active\def^{\ifmmode\sp\else\^{}\fi}\catcode`\%=\active\def%{\%}$\beta$}}%
\end{pgfscope}%
\begin{pgfscope}%
\pgfpathrectangle{\pgfqpoint{0.650000in}{0.420000in}}{\pgfqpoint{2.388110in}{1.994488in}}%
\pgfusepath{clip}%
\pgfsetrectcap%
\pgfsetroundjoin%
\pgfsetlinewidth{1.003750pt}%
\definecolor{currentstroke}{rgb}{0.870588,0.466667,0.682353}%
\pgfsetstrokecolor{currentstroke}%
\pgfsetdash{}{0pt}%
\pgfpathmoveto{\pgfqpoint{0.650000in}{1.348741in}}%
\pgfpathlineto{\pgfqpoint{0.659204in}{1.350044in}}%
\pgfpathlineto{\pgfqpoint{0.662044in}{1.351434in}}%
\pgfpathlineto{\pgfqpoint{0.663468in}{1.352877in}}%
\pgfpathlineto{\pgfqpoint{0.664893in}{1.355528in}}%
\pgfpathlineto{\pgfqpoint{0.666319in}{1.360159in}}%
\pgfpathlineto{\pgfqpoint{0.668457in}{1.371971in}}%
\pgfpathlineto{\pgfqpoint{0.674871in}{1.417574in}}%
\pgfpathlineto{\pgfqpoint{0.676296in}{1.420503in}}%
\pgfpathlineto{\pgfqpoint{0.677009in}{1.420577in}}%
\pgfpathlineto{\pgfqpoint{0.678435in}{1.418421in}}%
\pgfpathlineto{\pgfqpoint{0.680573in}{1.411489in}}%
\pgfpathlineto{\pgfqpoint{0.684136in}{1.399089in}}%
\pgfpathlineto{\pgfqpoint{0.686274in}{1.395310in}}%
\pgfpathlineto{\pgfqpoint{0.687700in}{1.394977in}}%
\pgfpathlineto{\pgfqpoint{0.689125in}{1.396309in}}%
\pgfpathlineto{\pgfqpoint{0.691263in}{1.400847in}}%
\pgfpathlineto{\pgfqpoint{0.695539in}{1.414568in}}%
\pgfpathlineto{\pgfqpoint{0.700528in}{1.429456in}}%
\pgfpathlineto{\pgfqpoint{0.703379in}{1.435121in}}%
\pgfpathlineto{\pgfqpoint{0.706230in}{1.438530in}}%
\pgfpathlineto{\pgfqpoint{0.709081in}{1.440005in}}%
\pgfpathlineto{\pgfqpoint{0.711931in}{1.440040in}}%
\pgfpathlineto{\pgfqpoint{0.715495in}{1.438676in}}%
\pgfpathlineto{\pgfqpoint{0.719771in}{1.435702in}}%
\pgfpathlineto{\pgfqpoint{0.724760in}{1.430870in}}%
\pgfpathlineto{\pgfqpoint{0.731887in}{1.424030in}}%
\pgfpathlineto{\pgfqpoint{0.736163in}{1.421605in}}%
\pgfpathlineto{\pgfqpoint{0.741865in}{1.419895in}}%
\pgfpathlineto{\pgfqpoint{0.753268in}{1.416755in}}%
\pgfpathlineto{\pgfqpoint{0.767522in}{1.410474in}}%
\pgfpathlineto{\pgfqpoint{0.776074in}{1.407363in}}%
\pgfpathlineto{\pgfqpoint{0.783201in}{1.406146in}}%
\pgfpathlineto{\pgfqpoint{0.814560in}{1.403383in}}%
\pgfpathlineto{\pgfqpoint{0.855897in}{1.403226in}}%
\pgfpathlineto{\pgfqpoint{0.861598in}{1.403585in}}%
\pgfpathlineto{\pgfqpoint{0.868725in}{1.404033in}}%
\pgfpathlineto{\pgfqpoint{0.895808in}{1.402188in}}%
\pgfpathlineto{\pgfqpoint{0.907211in}{1.401826in}}%
\pgfpathlineto{\pgfqpoint{0.916476in}{1.402773in}}%
\pgfpathlineto{\pgfqpoint{0.923603in}{1.402800in}}%
\pgfpathlineto{\pgfqpoint{0.937144in}{1.401792in}}%
\pgfpathlineto{\pgfqpoint{0.944984in}{1.401280in}}%
\pgfpathlineto{\pgfqpoint{0.958525in}{1.401366in}}%
\pgfpathlineto{\pgfqpoint{0.970641in}{1.400076in}}%
\pgfpathlineto{\pgfqpoint{0.977768in}{1.400343in}}%
\pgfpathlineto{\pgfqpoint{0.992735in}{1.401444in}}%
\pgfpathlineto{\pgfqpoint{1.006276in}{1.399998in}}%
\pgfpathlineto{\pgfqpoint{1.015541in}{1.398969in}}%
\pgfpathlineto{\pgfqpoint{1.026944in}{1.398647in}}%
\pgfpathlineto{\pgfqpoint{1.039773in}{1.399660in}}%
\pgfpathlineto{\pgfqpoint{1.049751in}{1.399310in}}%
\pgfpathlineto{\pgfqpoint{1.059016in}{1.398746in}}%
\pgfpathlineto{\pgfqpoint{1.071132in}{1.399159in}}%
\pgfpathlineto{\pgfqpoint{1.101065in}{1.397943in}}%
\pgfpathlineto{\pgfqpoint{1.106767in}{1.397586in}}%
\pgfpathlineto{\pgfqpoint{1.118883in}{1.395417in}}%
\pgfpathlineto{\pgfqpoint{1.127435in}{1.396322in}}%
\pgfpathlineto{\pgfqpoint{1.141689in}{1.398104in}}%
\pgfpathlineto{\pgfqpoint{1.153092in}{1.398315in}}%
\pgfpathlineto{\pgfqpoint{1.156656in}{1.397183in}}%
\pgfpathlineto{\pgfqpoint{1.165208in}{1.393687in}}%
\pgfpathlineto{\pgfqpoint{1.168772in}{1.393951in}}%
\pgfpathlineto{\pgfqpoint{1.180888in}{1.396334in}}%
\pgfpathlineto{\pgfqpoint{1.188015in}{1.395716in}}%
\pgfpathlineto{\pgfqpoint{1.198705in}{1.394317in}}%
\pgfpathlineto{\pgfqpoint{1.204407in}{1.395584in}}%
\pgfpathlineto{\pgfqpoint{1.207970in}{1.395884in}}%
\pgfpathlineto{\pgfqpoint{1.212246in}{1.394914in}}%
\pgfpathlineto{\pgfqpoint{1.216523in}{1.394118in}}%
\pgfpathlineto{\pgfqpoint{1.220799in}{1.394541in}}%
\pgfpathlineto{\pgfqpoint{1.228638in}{1.395313in}}%
\pgfpathlineto{\pgfqpoint{1.242892in}{1.395159in}}%
\pgfpathlineto{\pgfqpoint{1.250732in}{1.394309in}}%
\pgfpathlineto{\pgfqpoint{1.258572in}{1.394003in}}%
\pgfpathlineto{\pgfqpoint{1.269975in}{1.391609in}}%
\pgfpathlineto{\pgfqpoint{1.292781in}{1.394516in}}%
\pgfpathlineto{\pgfqpoint{1.305610in}{1.392607in}}%
\pgfpathlineto{\pgfqpoint{1.314875in}{1.392880in}}%
\pgfpathlineto{\pgfqpoint{1.323427in}{1.393056in}}%
\pgfpathlineto{\pgfqpoint{1.336969in}{1.391724in}}%
\pgfpathlineto{\pgfqpoint{1.342670in}{1.391921in}}%
\pgfpathlineto{\pgfqpoint{1.349085in}{1.390831in}}%
\pgfpathlineto{\pgfqpoint{1.356212in}{1.389667in}}%
\pgfpathlineto{\pgfqpoint{1.376167in}{1.390707in}}%
\pgfpathlineto{\pgfqpoint{1.388996in}{1.392833in}}%
\pgfpathlineto{\pgfqpoint{1.396835in}{1.392147in}}%
\pgfpathlineto{\pgfqpoint{1.405388in}{1.390664in}}%
\pgfpathlineto{\pgfqpoint{1.410377in}{1.389858in}}%
\pgfpathlineto{\pgfqpoint{1.413940in}{1.390495in}}%
\pgfpathlineto{\pgfqpoint{1.423205in}{1.393331in}}%
\pgfpathlineto{\pgfqpoint{1.426769in}{1.392523in}}%
\pgfpathlineto{\pgfqpoint{1.434608in}{1.390135in}}%
\pgfpathlineto{\pgfqpoint{1.440310in}{1.390047in}}%
\pgfpathlineto{\pgfqpoint{1.449575in}{1.390042in}}%
\pgfpathlineto{\pgfqpoint{1.475232in}{1.387950in}}%
\pgfpathlineto{\pgfqpoint{1.481647in}{1.387988in}}%
\pgfpathlineto{\pgfqpoint{1.492337in}{1.386511in}}%
\pgfpathlineto{\pgfqpoint{1.498751in}{1.386589in}}%
\pgfpathlineto{\pgfqpoint{1.504453in}{1.386720in}}%
\pgfpathlineto{\pgfqpoint{1.510867in}{1.386900in}}%
\pgfpathlineto{\pgfqpoint{1.515143in}{1.385733in}}%
\pgfpathlineto{\pgfqpoint{1.545077in}{1.373943in}}%
\pgfpathlineto{\pgfqpoint{1.550066in}{1.370348in}}%
\pgfpathlineto{\pgfqpoint{1.555767in}{1.364486in}}%
\pgfpathlineto{\pgfqpoint{1.571447in}{1.347139in}}%
\pgfpathlineto{\pgfqpoint{1.575723in}{1.339025in}}%
\pgfpathlineto{\pgfqpoint{1.582850in}{1.322106in}}%
\pgfpathlineto{\pgfqpoint{1.589264in}{1.304503in}}%
\pgfpathlineto{\pgfqpoint{1.593540in}{1.289338in}}%
\pgfpathlineto{\pgfqpoint{1.597817in}{1.269973in}}%
\pgfpathlineto{\pgfqpoint{1.602806in}{1.241475in}}%
\pgfpathlineto{\pgfqpoint{1.609933in}{1.198673in}}%
\pgfpathlineto{\pgfqpoint{1.612071in}{1.192263in}}%
\pgfpathlineto{\pgfqpoint{1.613496in}{1.191331in}}%
\pgfpathlineto{\pgfqpoint{1.614921in}{1.193248in}}%
\pgfpathlineto{\pgfqpoint{1.616347in}{1.197801in}}%
\pgfpathlineto{\pgfqpoint{1.619198in}{1.212784in}}%
\pgfpathlineto{\pgfqpoint{1.630601in}{1.282750in}}%
\pgfpathlineto{\pgfqpoint{1.635590in}{1.304214in}}%
\pgfpathlineto{\pgfqpoint{1.641291in}{1.323199in}}%
\pgfpathlineto{\pgfqpoint{1.647705in}{1.340724in}}%
\pgfpathlineto{\pgfqpoint{1.651982in}{1.349490in}}%
\pgfpathlineto{\pgfqpoint{1.656971in}{1.357004in}}%
\pgfpathlineto{\pgfqpoint{1.661960in}{1.362549in}}%
\pgfpathlineto{\pgfqpoint{1.668374in}{1.367750in}}%
\pgfpathlineto{\pgfqpoint{1.675501in}{1.373662in}}%
\pgfpathlineto{\pgfqpoint{1.684053in}{1.381219in}}%
\pgfpathlineto{\pgfqpoint{1.687617in}{1.382723in}}%
\pgfpathlineto{\pgfqpoint{1.694744in}{1.384002in}}%
\pgfpathlineto{\pgfqpoint{1.701158in}{1.385722in}}%
\pgfpathlineto{\pgfqpoint{1.710423in}{1.388176in}}%
\pgfpathlineto{\pgfqpoint{1.716125in}{1.388547in}}%
\pgfpathlineto{\pgfqpoint{1.724677in}{1.388688in}}%
\pgfpathlineto{\pgfqpoint{1.736081in}{1.390603in}}%
\pgfpathlineto{\pgfqpoint{1.749622in}{1.388503in}}%
\pgfpathlineto{\pgfqpoint{1.753898in}{1.389644in}}%
\pgfpathlineto{\pgfqpoint{1.758887in}{1.390629in}}%
\pgfpathlineto{\pgfqpoint{1.763876in}{1.390001in}}%
\pgfpathlineto{\pgfqpoint{1.768865in}{1.389776in}}%
\pgfpathlineto{\pgfqpoint{1.778843in}{1.390397in}}%
\pgfpathlineto{\pgfqpoint{1.785257in}{1.388448in}}%
\pgfpathlineto{\pgfqpoint{1.790246in}{1.387544in}}%
\pgfpathlineto{\pgfqpoint{1.795235in}{1.387548in}}%
\pgfpathlineto{\pgfqpoint{1.802362in}{1.390104in}}%
\pgfpathlineto{\pgfqpoint{1.807351in}{1.391735in}}%
\pgfpathlineto{\pgfqpoint{1.812340in}{1.391936in}}%
\pgfpathlineto{\pgfqpoint{1.834434in}{1.390294in}}%
\pgfpathlineto{\pgfqpoint{1.838710in}{1.389477in}}%
\pgfpathlineto{\pgfqpoint{1.841561in}{1.390035in}}%
\pgfpathlineto{\pgfqpoint{1.847975in}{1.391950in}}%
\pgfpathlineto{\pgfqpoint{1.850826in}{1.391183in}}%
\pgfpathlineto{\pgfqpoint{1.856528in}{1.387727in}}%
\pgfpathlineto{\pgfqpoint{1.860804in}{1.385884in}}%
\pgfpathlineto{\pgfqpoint{1.864367in}{1.385588in}}%
\pgfpathlineto{\pgfqpoint{1.880047in}{1.387381in}}%
\pgfpathlineto{\pgfqpoint{1.883610in}{1.388372in}}%
\pgfpathlineto{\pgfqpoint{1.886461in}{1.387911in}}%
\pgfpathlineto{\pgfqpoint{1.892875in}{1.385663in}}%
\pgfpathlineto{\pgfqpoint{1.895726in}{1.386369in}}%
\pgfpathlineto{\pgfqpoint{1.900715in}{1.388197in}}%
\pgfpathlineto{\pgfqpoint{1.903566in}{1.387944in}}%
\pgfpathlineto{\pgfqpoint{1.909980in}{1.386964in}}%
\pgfpathlineto{\pgfqpoint{1.921383in}{1.387614in}}%
\pgfpathlineto{\pgfqpoint{1.928510in}{1.386599in}}%
\pgfpathlineto{\pgfqpoint{1.935637in}{1.388252in}}%
\pgfpathlineto{\pgfqpoint{1.944190in}{1.392165in}}%
\pgfpathlineto{\pgfqpoint{1.947041in}{1.391693in}}%
\pgfpathlineto{\pgfqpoint{1.957018in}{1.388888in}}%
\pgfpathlineto{\pgfqpoint{1.963433in}{1.388867in}}%
\pgfpathlineto{\pgfqpoint{1.969847in}{1.388335in}}%
\pgfpathlineto{\pgfqpoint{1.976261in}{1.387710in}}%
\pgfpathlineto{\pgfqpoint{1.985527in}{1.387445in}}%
\pgfpathlineto{\pgfqpoint{1.991228in}{1.386271in}}%
\pgfpathlineto{\pgfqpoint{2.003344in}{1.381848in}}%
\pgfpathlineto{\pgfqpoint{2.012609in}{1.382559in}}%
\pgfpathlineto{\pgfqpoint{2.023300in}{1.384409in}}%
\pgfpathlineto{\pgfqpoint{2.035416in}{1.383825in}}%
\pgfpathlineto{\pgfqpoint{2.042543in}{1.385247in}}%
\pgfpathlineto{\pgfqpoint{2.049670in}{1.385699in}}%
\pgfpathlineto{\pgfqpoint{2.054659in}{1.386684in}}%
\pgfpathlineto{\pgfqpoint{2.057509in}{1.386018in}}%
\pgfpathlineto{\pgfqpoint{2.063924in}{1.384122in}}%
\pgfpathlineto{\pgfqpoint{2.068913in}{1.384095in}}%
\pgfpathlineto{\pgfqpoint{2.073902in}{1.385345in}}%
\pgfpathlineto{\pgfqpoint{2.081741in}{1.387520in}}%
\pgfpathlineto{\pgfqpoint{2.086730in}{1.387303in}}%
\pgfpathlineto{\pgfqpoint{2.093145in}{1.387268in}}%
\pgfpathlineto{\pgfqpoint{2.101697in}{1.388201in}}%
\pgfpathlineto{\pgfqpoint{2.104548in}{1.386412in}}%
\pgfpathlineto{\pgfqpoint{2.108824in}{1.383821in}}%
\pgfpathlineto{\pgfqpoint{2.111675in}{1.383754in}}%
\pgfpathlineto{\pgfqpoint{2.119515in}{1.385134in}}%
\pgfpathlineto{\pgfqpoint{2.123791in}{1.383798in}}%
\pgfpathlineto{\pgfqpoint{2.127354in}{1.382990in}}%
\pgfpathlineto{\pgfqpoint{2.130205in}{1.383525in}}%
\pgfpathlineto{\pgfqpoint{2.134481in}{1.385649in}}%
\pgfpathlineto{\pgfqpoint{2.138045in}{1.387151in}}%
\pgfpathlineto{\pgfqpoint{2.140896in}{1.387225in}}%
\pgfpathlineto{\pgfqpoint{2.145172in}{1.385892in}}%
\pgfpathlineto{\pgfqpoint{2.150874in}{1.384408in}}%
\pgfpathlineto{\pgfqpoint{2.161564in}{1.382205in}}%
\pgfpathlineto{\pgfqpoint{2.165840in}{1.381626in}}%
\pgfpathlineto{\pgfqpoint{2.170829in}{1.382277in}}%
\pgfpathlineto{\pgfqpoint{2.180807in}{1.384436in}}%
\pgfpathlineto{\pgfqpoint{2.185083in}{1.383611in}}%
\pgfpathlineto{\pgfqpoint{2.192923in}{1.381514in}}%
\pgfpathlineto{\pgfqpoint{2.197199in}{1.381905in}}%
\pgfpathlineto{\pgfqpoint{2.202901in}{1.382378in}}%
\pgfpathlineto{\pgfqpoint{2.210740in}{1.382465in}}%
\pgfpathlineto{\pgfqpoint{2.217154in}{1.382993in}}%
\pgfpathlineto{\pgfqpoint{2.220718in}{1.381851in}}%
\pgfpathlineto{\pgfqpoint{2.228557in}{1.378482in}}%
\pgfpathlineto{\pgfqpoint{2.231408in}{1.378721in}}%
\pgfpathlineto{\pgfqpoint{2.237110in}{1.380974in}}%
\pgfpathlineto{\pgfqpoint{2.242099in}{1.382303in}}%
\pgfpathlineto{\pgfqpoint{2.252789in}{1.384576in}}%
\pgfpathlineto{\pgfqpoint{2.256352in}{1.385177in}}%
\pgfpathlineto{\pgfqpoint{2.259203in}{1.384331in}}%
\pgfpathlineto{\pgfqpoint{2.268468in}{1.379600in}}%
\pgfpathlineto{\pgfqpoint{2.271319in}{1.379881in}}%
\pgfpathlineto{\pgfqpoint{2.275596in}{1.381900in}}%
\pgfpathlineto{\pgfqpoint{2.278446in}{1.382879in}}%
\pgfpathlineto{\pgfqpoint{2.280584in}{1.382476in}}%
\pgfpathlineto{\pgfqpoint{2.283435in}{1.380460in}}%
\pgfpathlineto{\pgfqpoint{2.286999in}{1.378011in}}%
\pgfpathlineto{\pgfqpoint{2.289137in}{1.377811in}}%
\pgfpathlineto{\pgfqpoint{2.291275in}{1.378721in}}%
\pgfpathlineto{\pgfqpoint{2.294838in}{1.381688in}}%
\pgfpathlineto{\pgfqpoint{2.299115in}{1.385232in}}%
\pgfpathlineto{\pgfqpoint{2.301253in}{1.385916in}}%
\pgfpathlineto{\pgfqpoint{2.303391in}{1.385646in}}%
\pgfpathlineto{\pgfqpoint{2.306242in}{1.384166in}}%
\pgfpathlineto{\pgfqpoint{2.317645in}{1.376840in}}%
\pgfpathlineto{\pgfqpoint{2.320496in}{1.377265in}}%
\pgfpathlineto{\pgfqpoint{2.323347in}{1.378980in}}%
\pgfpathlineto{\pgfqpoint{2.329761in}{1.383740in}}%
\pgfpathlineto{\pgfqpoint{2.331899in}{1.383592in}}%
\pgfpathlineto{\pgfqpoint{2.334750in}{1.381874in}}%
\pgfpathlineto{\pgfqpoint{2.339026in}{1.379319in}}%
\pgfpathlineto{\pgfqpoint{2.343302in}{1.378496in}}%
\pgfpathlineto{\pgfqpoint{2.348291in}{1.377190in}}%
\pgfpathlineto{\pgfqpoint{2.351855in}{1.376539in}}%
\pgfpathlineto{\pgfqpoint{2.356131in}{1.377044in}}%
\pgfpathlineto{\pgfqpoint{2.363971in}{1.379142in}}%
\pgfpathlineto{\pgfqpoint{2.376087in}{1.383715in}}%
\pgfpathlineto{\pgfqpoint{2.378938in}{1.382594in}}%
\pgfpathlineto{\pgfqpoint{2.383927in}{1.380014in}}%
\pgfpathlineto{\pgfqpoint{2.386778in}{1.379927in}}%
\pgfpathlineto{\pgfqpoint{2.391766in}{1.380208in}}%
\pgfpathlineto{\pgfqpoint{2.396755in}{1.380233in}}%
\pgfpathlineto{\pgfqpoint{2.403170in}{1.381557in}}%
\pgfpathlineto{\pgfqpoint{2.406021in}{1.380093in}}%
\pgfpathlineto{\pgfqpoint{2.410297in}{1.377731in}}%
\pgfpathlineto{\pgfqpoint{2.413148in}{1.377399in}}%
\pgfpathlineto{\pgfqpoint{2.416711in}{1.378219in}}%
\pgfpathlineto{\pgfqpoint{2.420275in}{1.380152in}}%
\pgfpathlineto{\pgfqpoint{2.425976in}{1.383850in}}%
\pgfpathlineto{\pgfqpoint{2.428827in}{1.384101in}}%
\pgfpathlineto{\pgfqpoint{2.433103in}{1.383995in}}%
\pgfpathlineto{\pgfqpoint{2.435241in}{1.385195in}}%
\pgfpathlineto{\pgfqpoint{2.438092in}{1.388615in}}%
\pgfpathlineto{\pgfqpoint{2.442368in}{1.394044in}}%
\pgfpathlineto{\pgfqpoint{2.443794in}{1.394773in}}%
\pgfpathlineto{\pgfqpoint{2.445219in}{1.394622in}}%
\pgfpathlineto{\pgfqpoint{2.447357in}{1.393158in}}%
\pgfpathlineto{\pgfqpoint{2.455909in}{1.385374in}}%
\pgfpathlineto{\pgfqpoint{2.459473in}{1.384185in}}%
\pgfpathlineto{\pgfqpoint{2.463037in}{1.384177in}}%
\pgfpathlineto{\pgfqpoint{2.467313in}{1.385517in}}%
\pgfpathlineto{\pgfqpoint{2.484418in}{1.393435in}}%
\pgfpathlineto{\pgfqpoint{2.489407in}{1.398485in}}%
\pgfpathlineto{\pgfqpoint{2.496534in}{1.407445in}}%
\pgfpathlineto{\pgfqpoint{2.505086in}{1.418818in}}%
\pgfpathlineto{\pgfqpoint{2.512213in}{1.423678in}}%
\pgfpathlineto{\pgfqpoint{2.515064in}{1.428936in}}%
\pgfpathlineto{\pgfqpoint{2.520053in}{1.441340in}}%
\pgfpathlineto{\pgfqpoint{2.525042in}{1.456416in}}%
\pgfpathlineto{\pgfqpoint{2.534307in}{1.489356in}}%
\pgfpathlineto{\pgfqpoint{2.538583in}{1.508206in}}%
\pgfpathlineto{\pgfqpoint{2.542860in}{1.532882in}}%
\pgfpathlineto{\pgfqpoint{2.547136in}{1.564919in}}%
\pgfpathlineto{\pgfqpoint{2.549987in}{1.591248in}}%
\pgfpathlineto{\pgfqpoint{2.554263in}{1.642002in}}%
\pgfpathlineto{\pgfqpoint{2.559965in}{1.712915in}}%
\pgfpathlineto{\pgfqpoint{2.562103in}{1.727880in}}%
\pgfpathlineto{\pgfqpoint{2.563528in}{1.731461in}}%
\pgfpathlineto{\pgfqpoint{2.564241in}{1.731265in}}%
\pgfpathlineto{\pgfqpoint{2.565666in}{1.727133in}}%
\pgfpathlineto{\pgfqpoint{2.567805in}{1.713186in}}%
\pgfpathlineto{\pgfqpoint{2.571368in}{1.677941in}}%
\pgfpathlineto{\pgfqpoint{2.578495in}{1.607642in}}%
\pgfpathlineto{\pgfqpoint{2.586335in}{1.544848in}}%
\pgfpathlineto{\pgfqpoint{2.591323in}{1.511801in}}%
\pgfpathlineto{\pgfqpoint{2.596313in}{1.487203in}}%
\pgfpathlineto{\pgfqpoint{2.604153in}{1.455005in}}%
\pgfpathlineto{\pgfqpoint{2.607716in}{1.444794in}}%
\pgfpathlineto{\pgfqpoint{2.610567in}{1.440111in}}%
\pgfpathlineto{\pgfqpoint{2.617694in}{1.432169in}}%
\pgfpathlineto{\pgfqpoint{2.620545in}{1.425327in}}%
\pgfpathlineto{\pgfqpoint{2.630523in}{1.398011in}}%
\pgfpathlineto{\pgfqpoint{2.633374in}{1.392777in}}%
\pgfpathlineto{\pgfqpoint{2.635512in}{1.390546in}}%
\pgfpathlineto{\pgfqpoint{2.637650in}{1.389899in}}%
\pgfpathlineto{\pgfqpoint{2.640501in}{1.390633in}}%
\pgfpathlineto{\pgfqpoint{2.644065in}{1.391466in}}%
\pgfpathlineto{\pgfqpoint{2.646203in}{1.390846in}}%
\pgfpathlineto{\pgfqpoint{2.649054in}{1.388645in}}%
\pgfpathlineto{\pgfqpoint{2.654042in}{1.384430in}}%
\pgfpathlineto{\pgfqpoint{2.656893in}{1.383626in}}%
\pgfpathlineto{\pgfqpoint{2.662595in}{1.382926in}}%
\pgfpathlineto{\pgfqpoint{2.665446in}{1.381015in}}%
\pgfpathlineto{\pgfqpoint{2.674711in}{1.373345in}}%
\pgfpathlineto{\pgfqpoint{2.679700in}{1.371056in}}%
\pgfpathlineto{\pgfqpoint{2.682551in}{1.370682in}}%
\pgfpathlineto{\pgfqpoint{2.684689in}{1.371520in}}%
\pgfpathlineto{\pgfqpoint{2.686827in}{1.373712in}}%
\pgfpathlineto{\pgfqpoint{2.693955in}{1.383086in}}%
\pgfpathlineto{\pgfqpoint{2.696093in}{1.383514in}}%
\pgfpathlineto{\pgfqpoint{2.698944in}{1.382569in}}%
\pgfpathlineto{\pgfqpoint{2.708209in}{1.377533in}}%
\pgfpathlineto{\pgfqpoint{2.711773in}{1.374984in}}%
\pgfpathlineto{\pgfqpoint{2.713911in}{1.374753in}}%
\pgfpathlineto{\pgfqpoint{2.720325in}{1.375725in}}%
\pgfpathlineto{\pgfqpoint{2.723889in}{1.375042in}}%
\pgfpathlineto{\pgfqpoint{2.725314in}{1.375714in}}%
\pgfpathlineto{\pgfqpoint{2.727453in}{1.378250in}}%
\pgfpathlineto{\pgfqpoint{2.731729in}{1.384107in}}%
\pgfpathlineto{\pgfqpoint{2.733155in}{1.384179in}}%
\pgfpathlineto{\pgfqpoint{2.734580in}{1.382663in}}%
\pgfpathlineto{\pgfqpoint{2.736718in}{1.377667in}}%
\pgfpathlineto{\pgfqpoint{2.740994in}{1.365502in}}%
\pgfpathlineto{\pgfqpoint{2.742420in}{1.363669in}}%
\pgfpathlineto{\pgfqpoint{2.743845in}{1.363617in}}%
\pgfpathlineto{\pgfqpoint{2.745271in}{1.365076in}}%
\pgfpathlineto{\pgfqpoint{2.748834in}{1.370391in}}%
\pgfpathlineto{\pgfqpoint{2.750260in}{1.371241in}}%
\pgfpathlineto{\pgfqpoint{2.751685in}{1.370962in}}%
\pgfpathlineto{\pgfqpoint{2.753823in}{1.368908in}}%
\pgfpathlineto{\pgfqpoint{2.758100in}{1.364352in}}%
\pgfpathlineto{\pgfqpoint{2.759525in}{1.364017in}}%
\pgfpathlineto{\pgfqpoint{2.760951in}{1.364690in}}%
\pgfpathlineto{\pgfqpoint{2.762376in}{1.366517in}}%
\pgfpathlineto{\pgfqpoint{2.764514in}{1.371473in}}%
\pgfpathlineto{\pgfqpoint{2.768791in}{1.386237in}}%
\pgfpathlineto{\pgfqpoint{2.770929in}{1.392164in}}%
\pgfpathlineto{\pgfqpoint{2.772354in}{1.394045in}}%
\pgfpathlineto{\pgfqpoint{2.773067in}{1.394244in}}%
\pgfpathlineto{\pgfqpoint{2.774492in}{1.392906in}}%
\pgfpathlineto{\pgfqpoint{2.775918in}{1.389570in}}%
\pgfpathlineto{\pgfqpoint{2.778768in}{1.379058in}}%
\pgfpathlineto{\pgfqpoint{2.783045in}{1.363261in}}%
\pgfpathlineto{\pgfqpoint{2.785183in}{1.358556in}}%
\pgfpathlineto{\pgfqpoint{2.787321in}{1.356619in}}%
\pgfpathlineto{\pgfqpoint{2.788746in}{1.356663in}}%
\pgfpathlineto{\pgfqpoint{2.790885in}{1.358094in}}%
\pgfpathlineto{\pgfqpoint{2.799437in}{1.365969in}}%
\pgfpathlineto{\pgfqpoint{2.803001in}{1.373051in}}%
\pgfpathlineto{\pgfqpoint{2.806565in}{1.380047in}}%
\pgfpathlineto{\pgfqpoint{2.807990in}{1.381336in}}%
\pgfpathlineto{\pgfqpoint{2.809415in}{1.381341in}}%
\pgfpathlineto{\pgfqpoint{2.810841in}{1.380099in}}%
\pgfpathlineto{\pgfqpoint{2.812979in}{1.376370in}}%
\pgfpathlineto{\pgfqpoint{2.818681in}{1.364767in}}%
\pgfpathlineto{\pgfqpoint{2.820106in}{1.363547in}}%
\pgfpathlineto{\pgfqpoint{2.821531in}{1.363488in}}%
\pgfpathlineto{\pgfqpoint{2.822957in}{1.364513in}}%
\pgfpathlineto{\pgfqpoint{2.825808in}{1.368718in}}%
\pgfpathlineto{\pgfqpoint{2.829371in}{1.373587in}}%
\pgfpathlineto{\pgfqpoint{2.831510in}{1.374306in}}%
\pgfpathlineto{\pgfqpoint{2.833648in}{1.373578in}}%
\pgfpathlineto{\pgfqpoint{2.837211in}{1.371961in}}%
\pgfpathlineto{\pgfqpoint{2.839350in}{1.371979in}}%
\pgfpathlineto{\pgfqpoint{2.850040in}{1.375256in}}%
\pgfpathlineto{\pgfqpoint{2.852890in}{1.377349in}}%
\pgfpathlineto{\pgfqpoint{2.858592in}{1.381806in}}%
\pgfpathlineto{\pgfqpoint{2.860729in}{1.381967in}}%
\pgfpathlineto{\pgfqpoint{2.862867in}{1.380882in}}%
\pgfpathlineto{\pgfqpoint{2.868569in}{1.376598in}}%
\pgfpathlineto{\pgfqpoint{2.871419in}{1.376370in}}%
\pgfpathlineto{\pgfqpoint{2.874270in}{1.376071in}}%
\pgfpathlineto{\pgfqpoint{2.876408in}{1.374578in}}%
\pgfpathlineto{\pgfqpoint{2.879259in}{1.370660in}}%
\pgfpathlineto{\pgfqpoint{2.883535in}{1.365370in}}%
\pgfpathlineto{\pgfqpoint{2.884960in}{1.364941in}}%
\pgfpathlineto{\pgfqpoint{2.886385in}{1.365387in}}%
\pgfpathlineto{\pgfqpoint{2.888523in}{1.367365in}}%
\pgfpathlineto{\pgfqpoint{2.894225in}{1.373584in}}%
\pgfpathlineto{\pgfqpoint{2.896363in}{1.373922in}}%
\pgfpathlineto{\pgfqpoint{2.898501in}{1.372867in}}%
\pgfpathlineto{\pgfqpoint{2.905627in}{1.367334in}}%
\pgfpathlineto{\pgfqpoint{2.908478in}{1.368782in}}%
\pgfpathlineto{\pgfqpoint{2.911328in}{1.372437in}}%
\pgfpathlineto{\pgfqpoint{2.915605in}{1.378150in}}%
\pgfpathlineto{\pgfqpoint{2.918455in}{1.380020in}}%
\pgfpathlineto{\pgfqpoint{2.921306in}{1.380139in}}%
\pgfpathlineto{\pgfqpoint{2.928432in}{1.378159in}}%
\pgfpathlineto{\pgfqpoint{2.934846in}{1.375150in}}%
\pgfpathlineto{\pgfqpoint{2.937697in}{1.375676in}}%
\pgfpathlineto{\pgfqpoint{2.941973in}{1.376901in}}%
\pgfpathlineto{\pgfqpoint{2.944111in}{1.376476in}}%
\pgfpathlineto{\pgfqpoint{2.946249in}{1.374645in}}%
\pgfpathlineto{\pgfqpoint{2.952663in}{1.366218in}}%
\pgfpathlineto{\pgfqpoint{2.954088in}{1.366310in}}%
\pgfpathlineto{\pgfqpoint{2.956226in}{1.368179in}}%
\pgfpathlineto{\pgfqpoint{2.961928in}{1.375010in}}%
\pgfpathlineto{\pgfqpoint{2.964066in}{1.375255in}}%
\pgfpathlineto{\pgfqpoint{2.966203in}{1.373948in}}%
\pgfpathlineto{\pgfqpoint{2.969767in}{1.371569in}}%
\pgfpathlineto{\pgfqpoint{2.971192in}{1.371807in}}%
\pgfpathlineto{\pgfqpoint{2.973330in}{1.373707in}}%
\pgfpathlineto{\pgfqpoint{2.982595in}{1.384619in}}%
\pgfpathlineto{\pgfqpoint{2.984021in}{1.384629in}}%
\pgfpathlineto{\pgfqpoint{2.986159in}{1.383121in}}%
\pgfpathlineto{\pgfqpoint{2.990435in}{1.379005in}}%
\pgfpathlineto{\pgfqpoint{2.993285in}{1.378003in}}%
\pgfpathlineto{\pgfqpoint{2.996849in}{1.377499in}}%
\pgfpathlineto{\pgfqpoint{2.998274in}{1.378004in}}%
\pgfpathlineto{\pgfqpoint{3.000412in}{1.380407in}}%
\pgfpathlineto{\pgfqpoint{3.007538in}{1.391847in}}%
\pgfpathlineto{\pgfqpoint{3.010389in}{1.393027in}}%
\pgfpathlineto{\pgfqpoint{3.012527in}{1.392583in}}%
\pgfpathlineto{\pgfqpoint{3.014665in}{1.390546in}}%
\pgfpathlineto{\pgfqpoint{3.020366in}{1.383335in}}%
\pgfpathlineto{\pgfqpoint{3.023930in}{1.380313in}}%
\pgfpathlineto{\pgfqpoint{3.027489in}{1.374693in}}%
\pgfpathlineto{\pgfqpoint{3.031010in}{1.368746in}}%
\pgfpathlineto{\pgfqpoint{3.034327in}{1.363810in}}%
\pgfpathlineto{\pgfqpoint{3.036519in}{1.358369in}}%
\pgfpathlineto{\pgfqpoint{3.038110in}{1.353079in}}%
\pgfpathlineto{\pgfqpoint{3.038110in}{1.353079in}}%
\pgfusepath{stroke}%
\end{pgfscope}%
\begin{pgfscope}%
\pgfpathrectangle{\pgfqpoint{0.650000in}{0.420000in}}{\pgfqpoint{2.388110in}{1.994488in}}%
\pgfusepath{clip}%
\pgfsetrectcap%
\pgfsetroundjoin%
\pgfsetlinewidth{1.003750pt}%
\definecolor{currentstroke}{rgb}{0.498039,0.737255,0.254902}%
\pgfsetstrokecolor{currentstroke}%
\pgfsetdash{}{0pt}%
\pgfpathmoveto{\pgfqpoint{0.650000in}{1.348774in}}%
\pgfpathlineto{\pgfqpoint{0.655023in}{1.348933in}}%
\pgfpathlineto{\pgfqpoint{0.659834in}{1.350361in}}%
\pgfpathlineto{\pgfqpoint{0.662656in}{1.352092in}}%
\pgfpathlineto{\pgfqpoint{0.664069in}{1.354039in}}%
\pgfpathlineto{\pgfqpoint{0.665483in}{1.357574in}}%
\pgfpathlineto{\pgfqpoint{0.667604in}{1.367319in}}%
\pgfpathlineto{\pgfqpoint{0.670433in}{1.388762in}}%
\pgfpathlineto{\pgfqpoint{0.673968in}{1.414492in}}%
\pgfpathlineto{\pgfqpoint{0.675382in}{1.419380in}}%
\pgfpathlineto{\pgfqpoint{0.676797in}{1.420455in}}%
\pgfpathlineto{\pgfqpoint{0.678211in}{1.418330in}}%
\pgfpathlineto{\pgfqpoint{0.680332in}{1.411467in}}%
\pgfpathlineto{\pgfqpoint{0.683868in}{1.399229in}}%
\pgfpathlineto{\pgfqpoint{0.685989in}{1.395580in}}%
\pgfpathlineto{\pgfqpoint{0.687403in}{1.395337in}}%
\pgfpathlineto{\pgfqpoint{0.688817in}{1.396757in}}%
\pgfpathlineto{\pgfqpoint{0.690938in}{1.401394in}}%
\pgfpathlineto{\pgfqpoint{0.695181in}{1.415026in}}%
\pgfpathlineto{\pgfqpoint{0.699424in}{1.427569in}}%
\pgfpathlineto{\pgfqpoint{0.702959in}{1.434486in}}%
\pgfpathlineto{\pgfqpoint{0.705787in}{1.437734in}}%
\pgfpathlineto{\pgfqpoint{0.708616in}{1.439294in}}%
\pgfpathlineto{\pgfqpoint{0.711444in}{1.439455in}}%
\pgfpathlineto{\pgfqpoint{0.714273in}{1.438447in}}%
\pgfpathlineto{\pgfqpoint{0.718515in}{1.435523in}}%
\pgfpathlineto{\pgfqpoint{0.733364in}{1.423852in}}%
\pgfpathlineto{\pgfqpoint{0.739021in}{1.421634in}}%
\pgfpathlineto{\pgfqpoint{0.747506in}{1.418524in}}%
\pgfpathlineto{\pgfqpoint{0.772961in}{1.407790in}}%
\pgfpathlineto{\pgfqpoint{0.792053in}{1.403995in}}%
\pgfpathlineto{\pgfqpoint{0.804073in}{1.403976in}}%
\pgfpathlineto{\pgfqpoint{0.833771in}{1.403282in}}%
\pgfpathlineto{\pgfqpoint{0.841549in}{1.402031in}}%
\pgfpathlineto{\pgfqpoint{0.846499in}{1.402491in}}%
\pgfpathlineto{\pgfqpoint{0.858519in}{1.404121in}}%
\pgfpathlineto{\pgfqpoint{0.871954in}{1.403818in}}%
\pgfpathlineto{\pgfqpoint{0.881146in}{1.403217in}}%
\pgfpathlineto{\pgfqpoint{0.895995in}{1.403405in}}%
\pgfpathlineto{\pgfqpoint{0.906602in}{1.402423in}}%
\pgfpathlineto{\pgfqpoint{0.923572in}{1.402573in}}%
\pgfpathlineto{\pgfqpoint{0.931350in}{1.402441in}}%
\pgfpathlineto{\pgfqpoint{0.939835in}{1.402439in}}%
\pgfpathlineto{\pgfqpoint{0.956805in}{1.400709in}}%
\pgfpathlineto{\pgfqpoint{0.963169in}{1.401007in}}%
\pgfpathlineto{\pgfqpoint{0.974483in}{1.402195in}}%
\pgfpathlineto{\pgfqpoint{0.980846in}{1.401180in}}%
\pgfpathlineto{\pgfqpoint{0.992867in}{1.399076in}}%
\pgfpathlineto{\pgfqpoint{0.998524in}{1.399236in}}%
\pgfpathlineto{\pgfqpoint{1.009837in}{1.400372in}}%
\pgfpathlineto{\pgfqpoint{1.026807in}{1.398938in}}%
\pgfpathlineto{\pgfqpoint{1.033878in}{1.399018in}}%
\pgfpathlineto{\pgfqpoint{1.043778in}{1.399501in}}%
\pgfpathlineto{\pgfqpoint{1.060041in}{1.398635in}}%
\pgfpathlineto{\pgfqpoint{1.069233in}{1.399931in}}%
\pgfpathlineto{\pgfqpoint{1.076304in}{1.399488in}}%
\pgfpathlineto{\pgfqpoint{1.087617in}{1.398092in}}%
\pgfpathlineto{\pgfqpoint{1.094688in}{1.397486in}}%
\pgfpathlineto{\pgfqpoint{1.101052in}{1.397077in}}%
\pgfpathlineto{\pgfqpoint{1.118022in}{1.397750in}}%
\pgfpathlineto{\pgfqpoint{1.137821in}{1.397497in}}%
\pgfpathlineto{\pgfqpoint{1.144185in}{1.397408in}}%
\pgfpathlineto{\pgfqpoint{1.149134in}{1.396106in}}%
\pgfpathlineto{\pgfqpoint{1.155498in}{1.394380in}}%
\pgfpathlineto{\pgfqpoint{1.160448in}{1.394431in}}%
\pgfpathlineto{\pgfqpoint{1.170347in}{1.394633in}}%
\pgfpathlineto{\pgfqpoint{1.176004in}{1.395267in}}%
\pgfpathlineto{\pgfqpoint{1.189439in}{1.398077in}}%
\pgfpathlineto{\pgfqpoint{1.193681in}{1.397133in}}%
\pgfpathlineto{\pgfqpoint{1.202873in}{1.394476in}}%
\pgfpathlineto{\pgfqpoint{1.209944in}{1.394515in}}%
\pgfpathlineto{\pgfqpoint{1.217722in}{1.394012in}}%
\pgfpathlineto{\pgfqpoint{1.231157in}{1.393010in}}%
\pgfpathlineto{\pgfqpoint{1.236107in}{1.393518in}}%
\pgfpathlineto{\pgfqpoint{1.242471in}{1.394280in}}%
\pgfpathlineto{\pgfqpoint{1.249541in}{1.393653in}}%
\pgfpathlineto{\pgfqpoint{1.257319in}{1.392710in}}%
\pgfpathlineto{\pgfqpoint{1.262269in}{1.393832in}}%
\pgfpathlineto{\pgfqpoint{1.266512in}{1.394561in}}%
\pgfpathlineto{\pgfqpoint{1.270754in}{1.394030in}}%
\pgfpathlineto{\pgfqpoint{1.277825in}{1.392931in}}%
\pgfpathlineto{\pgfqpoint{1.295502in}{1.392946in}}%
\pgfpathlineto{\pgfqpoint{1.306109in}{1.393235in}}%
\pgfpathlineto{\pgfqpoint{1.312473in}{1.391876in}}%
\pgfpathlineto{\pgfqpoint{1.319544in}{1.390419in}}%
\pgfpathlineto{\pgfqpoint{1.328736in}{1.389572in}}%
\pgfpathlineto{\pgfqpoint{1.338635in}{1.385651in}}%
\pgfpathlineto{\pgfqpoint{1.347827in}{1.385236in}}%
\pgfpathlineto{\pgfqpoint{1.357727in}{1.383001in}}%
\pgfpathlineto{\pgfqpoint{1.363383in}{1.381989in}}%
\pgfpathlineto{\pgfqpoint{1.369040in}{1.379677in}}%
\pgfpathlineto{\pgfqpoint{1.383182in}{1.372383in}}%
\pgfpathlineto{\pgfqpoint{1.389546in}{1.368043in}}%
\pgfpathlineto{\pgfqpoint{1.405809in}{1.354096in}}%
\pgfpathlineto{\pgfqpoint{1.412880in}{1.348791in}}%
\pgfpathlineto{\pgfqpoint{1.441871in}{1.329839in}}%
\pgfpathlineto{\pgfqpoint{1.446113in}{1.328811in}}%
\pgfpathlineto{\pgfqpoint{1.450356in}{1.329047in}}%
\pgfpathlineto{\pgfqpoint{1.454598in}{1.329100in}}%
\pgfpathlineto{\pgfqpoint{1.458134in}{1.327845in}}%
\pgfpathlineto{\pgfqpoint{1.463790in}{1.325695in}}%
\pgfpathlineto{\pgfqpoint{1.468033in}{1.325766in}}%
\pgfpathlineto{\pgfqpoint{1.477225in}{1.327554in}}%
\pgfpathlineto{\pgfqpoint{1.482175in}{1.329714in}}%
\pgfpathlineto{\pgfqpoint{1.487832in}{1.333517in}}%
\pgfpathlineto{\pgfqpoint{1.493488in}{1.338362in}}%
\pgfpathlineto{\pgfqpoint{1.501266in}{1.346592in}}%
\pgfpathlineto{\pgfqpoint{1.509751in}{1.354869in}}%
\pgfpathlineto{\pgfqpoint{1.515408in}{1.359170in}}%
\pgfpathlineto{\pgfqpoint{1.527429in}{1.366558in}}%
\pgfpathlineto{\pgfqpoint{1.536621in}{1.373943in}}%
\pgfpathlineto{\pgfqpoint{1.544399in}{1.379134in}}%
\pgfpathlineto{\pgfqpoint{1.548642in}{1.380423in}}%
\pgfpathlineto{\pgfqpoint{1.553591in}{1.380664in}}%
\pgfpathlineto{\pgfqpoint{1.558541in}{1.379896in}}%
\pgfpathlineto{\pgfqpoint{1.562076in}{1.378177in}}%
\pgfpathlineto{\pgfqpoint{1.565612in}{1.375136in}}%
\pgfpathlineto{\pgfqpoint{1.569854in}{1.370044in}}%
\pgfpathlineto{\pgfqpoint{1.574804in}{1.362204in}}%
\pgfpathlineto{\pgfqpoint{1.579754in}{1.351929in}}%
\pgfpathlineto{\pgfqpoint{1.583996in}{1.340928in}}%
\pgfpathlineto{\pgfqpoint{1.588239in}{1.326572in}}%
\pgfpathlineto{\pgfqpoint{1.591774in}{1.310745in}}%
\pgfpathlineto{\pgfqpoint{1.596017in}{1.285773in}}%
\pgfpathlineto{\pgfqpoint{1.602381in}{1.245734in}}%
\pgfpathlineto{\pgfqpoint{1.604502in}{1.238899in}}%
\pgfpathlineto{\pgfqpoint{1.605916in}{1.237876in}}%
\pgfpathlineto{\pgfqpoint{1.607330in}{1.239950in}}%
\pgfpathlineto{\pgfqpoint{1.608744in}{1.244892in}}%
\pgfpathlineto{\pgfqpoint{1.611573in}{1.261354in}}%
\pgfpathlineto{\pgfqpoint{1.623593in}{1.342705in}}%
\pgfpathlineto{\pgfqpoint{1.628543in}{1.365610in}}%
\pgfpathlineto{\pgfqpoint{1.634200in}{1.386105in}}%
\pgfpathlineto{\pgfqpoint{1.641978in}{1.409947in}}%
\pgfpathlineto{\pgfqpoint{1.646220in}{1.419775in}}%
\pgfpathlineto{\pgfqpoint{1.651170in}{1.428099in}}%
\pgfpathlineto{\pgfqpoint{1.661069in}{1.441600in}}%
\pgfpathlineto{\pgfqpoint{1.665312in}{1.445970in}}%
\pgfpathlineto{\pgfqpoint{1.668847in}{1.448299in}}%
\pgfpathlineto{\pgfqpoint{1.678747in}{1.452372in}}%
\pgfpathlineto{\pgfqpoint{1.681575in}{1.452775in}}%
\pgfpathlineto{\pgfqpoint{1.684403in}{1.452023in}}%
\pgfpathlineto{\pgfqpoint{1.692181in}{1.449666in}}%
\pgfpathlineto{\pgfqpoint{1.700667in}{1.448662in}}%
\pgfpathlineto{\pgfqpoint{1.706323in}{1.445864in}}%
\pgfpathlineto{\pgfqpoint{1.720465in}{1.438960in}}%
\pgfpathlineto{\pgfqpoint{1.728950in}{1.431795in}}%
\pgfpathlineto{\pgfqpoint{1.739557in}{1.423039in}}%
\pgfpathlineto{\pgfqpoint{1.749456in}{1.414858in}}%
\pgfpathlineto{\pgfqpoint{1.760770in}{1.403726in}}%
\pgfpathlineto{\pgfqpoint{1.765012in}{1.401415in}}%
\pgfpathlineto{\pgfqpoint{1.769962in}{1.400369in}}%
\pgfpathlineto{\pgfqpoint{1.776326in}{1.398870in}}%
\pgfpathlineto{\pgfqpoint{1.780569in}{1.398155in}}%
\pgfpathlineto{\pgfqpoint{1.784811in}{1.398744in}}%
\pgfpathlineto{\pgfqpoint{1.791882in}{1.401332in}}%
\pgfpathlineto{\pgfqpoint{1.796125in}{1.402242in}}%
\pgfpathlineto{\pgfqpoint{1.799660in}{1.401749in}}%
\pgfpathlineto{\pgfqpoint{1.808852in}{1.399264in}}%
\pgfpathlineto{\pgfqpoint{1.813095in}{1.400045in}}%
\pgfpathlineto{\pgfqpoint{1.816630in}{1.401778in}}%
\pgfpathlineto{\pgfqpoint{1.820166in}{1.404801in}}%
\pgfpathlineto{\pgfqpoint{1.844914in}{1.430349in}}%
\pgfpathlineto{\pgfqpoint{1.851985in}{1.436513in}}%
\pgfpathlineto{\pgfqpoint{1.858349in}{1.440947in}}%
\pgfpathlineto{\pgfqpoint{1.862592in}{1.442445in}}%
\pgfpathlineto{\pgfqpoint{1.866834in}{1.444049in}}%
\pgfpathlineto{\pgfqpoint{1.871077in}{1.447291in}}%
\pgfpathlineto{\pgfqpoint{1.874612in}{1.449767in}}%
\pgfpathlineto{\pgfqpoint{1.877441in}{1.450727in}}%
\pgfpathlineto{\pgfqpoint{1.882390in}{1.450896in}}%
\pgfpathlineto{\pgfqpoint{1.888754in}{1.451508in}}%
\pgfpathlineto{\pgfqpoint{1.897239in}{1.451979in}}%
\pgfpathlineto{\pgfqpoint{1.909967in}{1.452577in}}%
\pgfpathlineto{\pgfqpoint{1.913502in}{1.451140in}}%
\pgfpathlineto{\pgfqpoint{1.923402in}{1.446084in}}%
\pgfpathlineto{\pgfqpoint{1.930473in}{1.442698in}}%
\pgfpathlineto{\pgfqpoint{1.934008in}{1.439925in}}%
\pgfpathlineto{\pgfqpoint{1.943200in}{1.431412in}}%
\pgfpathlineto{\pgfqpoint{1.948150in}{1.429629in}}%
\pgfpathlineto{\pgfqpoint{1.952393in}{1.427656in}}%
\pgfpathlineto{\pgfqpoint{1.968656in}{1.417039in}}%
\pgfpathlineto{\pgfqpoint{1.975020in}{1.412846in}}%
\pgfpathlineto{\pgfqpoint{1.984212in}{1.408700in}}%
\pgfpathlineto{\pgfqpoint{1.992697in}{1.401210in}}%
\pgfpathlineto{\pgfqpoint{1.997647in}{1.399353in}}%
\pgfpathlineto{\pgfqpoint{2.004011in}{1.398033in}}%
\pgfpathlineto{\pgfqpoint{2.011789in}{1.397534in}}%
\pgfpathlineto{\pgfqpoint{2.016031in}{1.398870in}}%
\pgfpathlineto{\pgfqpoint{2.020274in}{1.399964in}}%
\pgfpathlineto{\pgfqpoint{2.025931in}{1.399955in}}%
\pgfpathlineto{\pgfqpoint{2.030880in}{1.398980in}}%
\pgfpathlineto{\pgfqpoint{2.037244in}{1.397201in}}%
\pgfpathlineto{\pgfqpoint{2.040072in}{1.397912in}}%
\pgfpathlineto{\pgfqpoint{2.044315in}{1.400353in}}%
\pgfpathlineto{\pgfqpoint{2.052800in}{1.405514in}}%
\pgfpathlineto{\pgfqpoint{2.055629in}{1.406028in}}%
\pgfpathlineto{\pgfqpoint{2.062699in}{1.406481in}}%
\pgfpathlineto{\pgfqpoint{2.065528in}{1.408437in}}%
\pgfpathlineto{\pgfqpoint{2.074720in}{1.416804in}}%
\pgfpathlineto{\pgfqpoint{2.078256in}{1.417510in}}%
\pgfpathlineto{\pgfqpoint{2.083205in}{1.418041in}}%
\pgfpathlineto{\pgfqpoint{2.086034in}{1.419580in}}%
\pgfpathlineto{\pgfqpoint{2.091690in}{1.423342in}}%
\pgfpathlineto{\pgfqpoint{2.095933in}{1.424288in}}%
\pgfpathlineto{\pgfqpoint{2.099468in}{1.425633in}}%
\pgfpathlineto{\pgfqpoint{2.103004in}{1.428569in}}%
\pgfpathlineto{\pgfqpoint{2.108661in}{1.433459in}}%
\pgfpathlineto{\pgfqpoint{2.113610in}{1.435917in}}%
\pgfpathlineto{\pgfqpoint{2.118560in}{1.438651in}}%
\pgfpathlineto{\pgfqpoint{2.125631in}{1.443248in}}%
\pgfpathlineto{\pgfqpoint{2.128459in}{1.443212in}}%
\pgfpathlineto{\pgfqpoint{2.131995in}{1.441437in}}%
\pgfpathlineto{\pgfqpoint{2.135530in}{1.439885in}}%
\pgfpathlineto{\pgfqpoint{2.137652in}{1.439960in}}%
\pgfpathlineto{\pgfqpoint{2.140480in}{1.441321in}}%
\pgfpathlineto{\pgfqpoint{2.150379in}{1.447594in}}%
\pgfpathlineto{\pgfqpoint{2.154622in}{1.449326in}}%
\pgfpathlineto{\pgfqpoint{2.157450in}{1.449325in}}%
\pgfpathlineto{\pgfqpoint{2.160279in}{1.448055in}}%
\pgfpathlineto{\pgfqpoint{2.165936in}{1.445172in}}%
\pgfpathlineto{\pgfqpoint{2.170178in}{1.444748in}}%
\pgfpathlineto{\pgfqpoint{2.175835in}{1.444206in}}%
\pgfpathlineto{\pgfqpoint{2.180077in}{1.443457in}}%
\pgfpathlineto{\pgfqpoint{2.183613in}{1.443783in}}%
\pgfpathlineto{\pgfqpoint{2.189976in}{1.444895in}}%
\pgfpathlineto{\pgfqpoint{2.194219in}{1.444002in}}%
\pgfpathlineto{\pgfqpoint{2.198461in}{1.443273in}}%
\pgfpathlineto{\pgfqpoint{2.206239in}{1.443023in}}%
\pgfpathlineto{\pgfqpoint{2.209775in}{1.441203in}}%
\pgfpathlineto{\pgfqpoint{2.226745in}{1.430733in}}%
\pgfpathlineto{\pgfqpoint{2.236644in}{1.425454in}}%
\pgfpathlineto{\pgfqpoint{2.242301in}{1.421273in}}%
\pgfpathlineto{\pgfqpoint{2.245836in}{1.417152in}}%
\pgfpathlineto{\pgfqpoint{2.251493in}{1.410155in}}%
\pgfpathlineto{\pgfqpoint{2.254321in}{1.408392in}}%
\pgfpathlineto{\pgfqpoint{2.258564in}{1.407429in}}%
\pgfpathlineto{\pgfqpoint{2.264221in}{1.406035in}}%
\pgfpathlineto{\pgfqpoint{2.274827in}{1.401911in}}%
\pgfpathlineto{\pgfqpoint{2.282605in}{1.402711in}}%
\pgfpathlineto{\pgfqpoint{2.287555in}{1.404389in}}%
\pgfpathlineto{\pgfqpoint{2.291797in}{1.405464in}}%
\pgfpathlineto{\pgfqpoint{2.295333in}{1.405238in}}%
\pgfpathlineto{\pgfqpoint{2.299575in}{1.403598in}}%
\pgfpathlineto{\pgfqpoint{2.304525in}{1.401643in}}%
\pgfpathlineto{\pgfqpoint{2.307354in}{1.401884in}}%
\pgfpathlineto{\pgfqpoint{2.309475in}{1.403222in}}%
\pgfpathlineto{\pgfqpoint{2.313010in}{1.407301in}}%
\pgfpathlineto{\pgfqpoint{2.315839in}{1.410030in}}%
\pgfpathlineto{\pgfqpoint{2.317960in}{1.410367in}}%
\pgfpathlineto{\pgfqpoint{2.320081in}{1.409246in}}%
\pgfpathlineto{\pgfqpoint{2.323617in}{1.406921in}}%
\pgfpathlineto{\pgfqpoint{2.325738in}{1.406975in}}%
\pgfpathlineto{\pgfqpoint{2.327859in}{1.408389in}}%
\pgfpathlineto{\pgfqpoint{2.332809in}{1.412265in}}%
\pgfpathlineto{\pgfqpoint{2.335638in}{1.412649in}}%
\pgfpathlineto{\pgfqpoint{2.348365in}{1.411949in}}%
\pgfpathlineto{\pgfqpoint{2.352608in}{1.413680in}}%
\pgfpathlineto{\pgfqpoint{2.358972in}{1.416242in}}%
\pgfpathlineto{\pgfqpoint{2.364629in}{1.417334in}}%
\pgfpathlineto{\pgfqpoint{2.375235in}{1.417348in}}%
\pgfpathlineto{\pgfqpoint{2.380892in}{1.416784in}}%
\pgfpathlineto{\pgfqpoint{2.383013in}{1.418059in}}%
\pgfpathlineto{\pgfqpoint{2.385842in}{1.421405in}}%
\pgfpathlineto{\pgfqpoint{2.390084in}{1.426552in}}%
\pgfpathlineto{\pgfqpoint{2.394327in}{1.429613in}}%
\pgfpathlineto{\pgfqpoint{2.398570in}{1.431772in}}%
\pgfpathlineto{\pgfqpoint{2.401398in}{1.432160in}}%
\pgfpathlineto{\pgfqpoint{2.404933in}{1.432377in}}%
\pgfpathlineto{\pgfqpoint{2.407055in}{1.433504in}}%
\pgfpathlineto{\pgfqpoint{2.414125in}{1.439104in}}%
\pgfpathlineto{\pgfqpoint{2.416954in}{1.438917in}}%
\pgfpathlineto{\pgfqpoint{2.420489in}{1.438363in}}%
\pgfpathlineto{\pgfqpoint{2.422610in}{1.439334in}}%
\pgfpathlineto{\pgfqpoint{2.424732in}{1.441808in}}%
\pgfpathlineto{\pgfqpoint{2.428267in}{1.448486in}}%
\pgfpathlineto{\pgfqpoint{2.436045in}{1.463850in}}%
\pgfpathlineto{\pgfqpoint{2.438874in}{1.467328in}}%
\pgfpathlineto{\pgfqpoint{2.441702in}{1.468928in}}%
\pgfpathlineto{\pgfqpoint{2.445945in}{1.470655in}}%
\pgfpathlineto{\pgfqpoint{2.448066in}{1.473118in}}%
\pgfpathlineto{\pgfqpoint{2.450894in}{1.478836in}}%
\pgfpathlineto{\pgfqpoint{2.472815in}{1.535037in}}%
\pgfpathlineto{\pgfqpoint{2.478471in}{1.552242in}}%
\pgfpathlineto{\pgfqpoint{2.482714in}{1.569376in}}%
\pgfpathlineto{\pgfqpoint{2.489078in}{1.596234in}}%
\pgfpathlineto{\pgfqpoint{2.496149in}{1.617470in}}%
\pgfpathlineto{\pgfqpoint{2.498977in}{1.631587in}}%
\pgfpathlineto{\pgfqpoint{2.504634in}{1.668142in}}%
\pgfpathlineto{\pgfqpoint{2.518069in}{1.760823in}}%
\pgfpathlineto{\pgfqpoint{2.523726in}{1.810851in}}%
\pgfpathlineto{\pgfqpoint{2.529383in}{1.872041in}}%
\pgfpathlineto{\pgfqpoint{2.534332in}{1.937926in}}%
\pgfpathlineto{\pgfqpoint{2.539989in}{2.030626in}}%
\pgfpathlineto{\pgfqpoint{2.544939in}{2.107000in}}%
\pgfpathlineto{\pgfqpoint{2.547060in}{2.124305in}}%
\pgfpathlineto{\pgfqpoint{2.548474in}{2.127282in}}%
\pgfpathlineto{\pgfqpoint{2.549181in}{2.126399in}}%
\pgfpathlineto{\pgfqpoint{2.550595in}{2.120502in}}%
\pgfpathlineto{\pgfqpoint{2.552717in}{2.102383in}}%
\pgfpathlineto{\pgfqpoint{2.556252in}{2.058481in}}%
\pgfpathlineto{\pgfqpoint{2.562616in}{1.980143in}}%
\pgfpathlineto{\pgfqpoint{2.568980in}{1.919319in}}%
\pgfpathlineto{\pgfqpoint{2.579586in}{1.824991in}}%
\pgfpathlineto{\pgfqpoint{2.585951in}{1.780727in}}%
\pgfpathlineto{\pgfqpoint{2.591607in}{1.747261in}}%
\pgfpathlineto{\pgfqpoint{2.602214in}{1.691673in}}%
\pgfpathlineto{\pgfqpoint{2.608578in}{1.656490in}}%
\pgfpathlineto{\pgfqpoint{2.617063in}{1.619901in}}%
\pgfpathlineto{\pgfqpoint{2.619892in}{1.610811in}}%
\pgfpathlineto{\pgfqpoint{2.622013in}{1.606705in}}%
\pgfpathlineto{\pgfqpoint{2.624134in}{1.604639in}}%
\pgfpathlineto{\pgfqpoint{2.628377in}{1.601689in}}%
\pgfpathlineto{\pgfqpoint{2.630498in}{1.598591in}}%
\pgfpathlineto{\pgfqpoint{2.633327in}{1.591869in}}%
\pgfpathlineto{\pgfqpoint{2.644640in}{1.557963in}}%
\pgfpathlineto{\pgfqpoint{2.653126in}{1.530889in}}%
\pgfpathlineto{\pgfqpoint{2.656661in}{1.524311in}}%
\pgfpathlineto{\pgfqpoint{2.660904in}{1.517050in}}%
\pgfpathlineto{\pgfqpoint{2.665146in}{1.506899in}}%
\pgfpathlineto{\pgfqpoint{2.667975in}{1.501056in}}%
\pgfpathlineto{\pgfqpoint{2.670096in}{1.498913in}}%
\pgfpathlineto{\pgfqpoint{2.672217in}{1.498631in}}%
\pgfpathlineto{\pgfqpoint{2.675046in}{1.498818in}}%
\pgfpathlineto{\pgfqpoint{2.676460in}{1.498039in}}%
\pgfpathlineto{\pgfqpoint{2.678582in}{1.495052in}}%
\pgfpathlineto{\pgfqpoint{2.681410in}{1.488309in}}%
\pgfpathlineto{\pgfqpoint{2.690603in}{1.461659in}}%
\pgfpathlineto{\pgfqpoint{2.696966in}{1.439998in}}%
\pgfpathlineto{\pgfqpoint{2.699088in}{1.436138in}}%
\pgfpathlineto{\pgfqpoint{2.701916in}{1.433605in}}%
\pgfpathlineto{\pgfqpoint{2.705452in}{1.432121in}}%
\pgfpathlineto{\pgfqpoint{2.707573in}{1.432452in}}%
\pgfpathlineto{\pgfqpoint{2.709695in}{1.434372in}}%
\pgfpathlineto{\pgfqpoint{2.714645in}{1.440195in}}%
\pgfpathlineto{\pgfqpoint{2.716059in}{1.439881in}}%
\pgfpathlineto{\pgfqpoint{2.717473in}{1.438030in}}%
\pgfpathlineto{\pgfqpoint{2.719595in}{1.432639in}}%
\pgfpathlineto{\pgfqpoint{2.724544in}{1.418027in}}%
\pgfpathlineto{\pgfqpoint{2.725958in}{1.416499in}}%
\pgfpathlineto{\pgfqpoint{2.727373in}{1.416716in}}%
\pgfpathlineto{\pgfqpoint{2.729494in}{1.419310in}}%
\pgfpathlineto{\pgfqpoint{2.732322in}{1.422731in}}%
\pgfpathlineto{\pgfqpoint{2.733737in}{1.422878in}}%
\pgfpathlineto{\pgfqpoint{2.735151in}{1.421682in}}%
\pgfpathlineto{\pgfqpoint{2.737272in}{1.417878in}}%
\pgfpathlineto{\pgfqpoint{2.742929in}{1.405968in}}%
\pgfpathlineto{\pgfqpoint{2.744343in}{1.404624in}}%
\pgfpathlineto{\pgfqpoint{2.745758in}{1.404471in}}%
\pgfpathlineto{\pgfqpoint{2.747172in}{1.405598in}}%
\pgfpathlineto{\pgfqpoint{2.749293in}{1.409457in}}%
\pgfpathlineto{\pgfqpoint{2.754950in}{1.422038in}}%
\pgfpathlineto{\pgfqpoint{2.756364in}{1.422714in}}%
\pgfpathlineto{\pgfqpoint{2.757778in}{1.421729in}}%
\pgfpathlineto{\pgfqpoint{2.759192in}{1.419148in}}%
\pgfpathlineto{\pgfqpoint{2.762021in}{1.410437in}}%
\pgfpathlineto{\pgfqpoint{2.768385in}{1.388161in}}%
\pgfpathlineto{\pgfqpoint{2.770506in}{1.383805in}}%
\pgfpathlineto{\pgfqpoint{2.772628in}{1.381677in}}%
\pgfpathlineto{\pgfqpoint{2.774749in}{1.381113in}}%
\pgfpathlineto{\pgfqpoint{2.779699in}{1.381345in}}%
\pgfpathlineto{\pgfqpoint{2.781820in}{1.382179in}}%
\pgfpathlineto{\pgfqpoint{2.784649in}{1.384920in}}%
\pgfpathlineto{\pgfqpoint{2.788184in}{1.388425in}}%
\pgfpathlineto{\pgfqpoint{2.789598in}{1.388430in}}%
\pgfpathlineto{\pgfqpoint{2.791013in}{1.387171in}}%
\pgfpathlineto{\pgfqpoint{2.793134in}{1.382998in}}%
\pgfpathlineto{\pgfqpoint{2.799498in}{1.367135in}}%
\pgfpathlineto{\pgfqpoint{2.801619in}{1.365208in}}%
\pgfpathlineto{\pgfqpoint{2.803033in}{1.365225in}}%
\pgfpathlineto{\pgfqpoint{2.805155in}{1.366785in}}%
\pgfpathlineto{\pgfqpoint{2.809397in}{1.370810in}}%
\pgfpathlineto{\pgfqpoint{2.810812in}{1.370816in}}%
\pgfpathlineto{\pgfqpoint{2.812226in}{1.369618in}}%
\pgfpathlineto{\pgfqpoint{2.814347in}{1.365590in}}%
\pgfpathlineto{\pgfqpoint{2.820004in}{1.352123in}}%
\pgfpathlineto{\pgfqpoint{2.821419in}{1.351103in}}%
\pgfpathlineto{\pgfqpoint{2.822832in}{1.351491in}}%
\pgfpathlineto{\pgfqpoint{2.824954in}{1.354078in}}%
\pgfpathlineto{\pgfqpoint{2.830610in}{1.362553in}}%
\pgfpathlineto{\pgfqpoint{2.832731in}{1.363919in}}%
\pgfpathlineto{\pgfqpoint{2.836267in}{1.364438in}}%
\pgfpathlineto{\pgfqpoint{2.839802in}{1.363973in}}%
\pgfpathlineto{\pgfqpoint{2.842630in}{1.362564in}}%
\pgfpathlineto{\pgfqpoint{2.846873in}{1.358621in}}%
\pgfpathlineto{\pgfqpoint{2.848994in}{1.357144in}}%
\pgfpathlineto{\pgfqpoint{2.851115in}{1.356954in}}%
\pgfpathlineto{\pgfqpoint{2.857478in}{1.358463in}}%
\pgfpathlineto{\pgfqpoint{2.859600in}{1.356661in}}%
\pgfpathlineto{\pgfqpoint{2.865963in}{1.348820in}}%
\pgfpathlineto{\pgfqpoint{2.867377in}{1.348639in}}%
\pgfpathlineto{\pgfqpoint{2.869498in}{1.349795in}}%
\pgfpathlineto{\pgfqpoint{2.875862in}{1.355290in}}%
\pgfpathlineto{\pgfqpoint{2.877276in}{1.355189in}}%
\pgfpathlineto{\pgfqpoint{2.879397in}{1.353599in}}%
\pgfpathlineto{\pgfqpoint{2.882225in}{1.349479in}}%
\pgfpathlineto{\pgfqpoint{2.885760in}{1.344254in}}%
\pgfpathlineto{\pgfqpoint{2.887882in}{1.343078in}}%
\pgfpathlineto{\pgfqpoint{2.889296in}{1.343667in}}%
\pgfpathlineto{\pgfqpoint{2.891417in}{1.346491in}}%
\pgfpathlineto{\pgfqpoint{2.897073in}{1.355924in}}%
\pgfpathlineto{\pgfqpoint{2.898488in}{1.356424in}}%
\pgfpathlineto{\pgfqpoint{2.899902in}{1.355839in}}%
\pgfpathlineto{\pgfqpoint{2.902023in}{1.353215in}}%
\pgfpathlineto{\pgfqpoint{2.908386in}{1.343173in}}%
\pgfpathlineto{\pgfqpoint{2.911922in}{1.340402in}}%
\pgfpathlineto{\pgfqpoint{2.914750in}{1.339241in}}%
\pgfpathlineto{\pgfqpoint{2.916871in}{1.339495in}}%
\pgfpathlineto{\pgfqpoint{2.919699in}{1.341375in}}%
\pgfpathlineto{\pgfqpoint{2.923942in}{1.344512in}}%
\pgfpathlineto{\pgfqpoint{2.926063in}{1.344890in}}%
\pgfpathlineto{\pgfqpoint{2.928184in}{1.343946in}}%
\pgfpathlineto{\pgfqpoint{2.934548in}{1.338695in}}%
\pgfpathlineto{\pgfqpoint{2.935961in}{1.339123in}}%
\pgfpathlineto{\pgfqpoint{2.938083in}{1.341198in}}%
\pgfpathlineto{\pgfqpoint{2.945153in}{1.350210in}}%
\pgfpathlineto{\pgfqpoint{2.947274in}{1.350753in}}%
\pgfpathlineto{\pgfqpoint{2.951517in}{1.350597in}}%
\pgfpathlineto{\pgfqpoint{2.953638in}{1.352346in}}%
\pgfpathlineto{\pgfqpoint{2.962123in}{1.362538in}}%
\pgfpathlineto{\pgfqpoint{2.964951in}{1.363737in}}%
\pgfpathlineto{\pgfqpoint{2.967072in}{1.363757in}}%
\pgfpathlineto{\pgfqpoint{2.970608in}{1.363627in}}%
\pgfpathlineto{\pgfqpoint{2.975557in}{1.364756in}}%
\pgfpathlineto{\pgfqpoint{2.977678in}{1.363576in}}%
\pgfpathlineto{\pgfqpoint{2.980506in}{1.361881in}}%
\pgfpathlineto{\pgfqpoint{2.981921in}{1.362190in}}%
\pgfpathlineto{\pgfqpoint{2.984042in}{1.364456in}}%
\pgfpathlineto{\pgfqpoint{2.989698in}{1.372335in}}%
\pgfpathlineto{\pgfqpoint{2.991819in}{1.373632in}}%
\pgfpathlineto{\pgfqpoint{2.993234in}{1.373575in}}%
\pgfpathlineto{\pgfqpoint{2.994648in}{1.372518in}}%
\pgfpathlineto{\pgfqpoint{2.996769in}{1.368973in}}%
\pgfpathlineto{\pgfqpoint{3.005253in}{1.349274in}}%
\pgfpathlineto{\pgfqpoint{3.008079in}{1.338517in}}%
\pgfpathlineto{\pgfqpoint{3.019322in}{1.285076in}}%
\pgfpathlineto{\pgfqpoint{3.019322in}{1.285076in}}%
\pgfusepath{stroke}%
\end{pgfscope}%
\begin{pgfscope}%
\pgfsetrectcap%
\pgfsetmiterjoin%
\pgfsetlinewidth{0.803000pt}%
\definecolor{currentstroke}{rgb}{0.000000,0.000000,0.000000}%
\pgfsetstrokecolor{currentstroke}%
\pgfsetdash{}{0pt}%
\pgfpathmoveto{\pgfqpoint{0.650000in}{0.420000in}}%
\pgfpathlineto{\pgfqpoint{0.650000in}{2.414488in}}%
\pgfusepath{stroke}%
\end{pgfscope}%
\begin{pgfscope}%
\pgfsetrectcap%
\pgfsetmiterjoin%
\pgfsetlinewidth{0.803000pt}%
\definecolor{currentstroke}{rgb}{0.000000,0.000000,0.000000}%
\pgfsetstrokecolor{currentstroke}%
\pgfsetdash{}{0pt}%
\pgfpathmoveto{\pgfqpoint{3.038110in}{0.420000in}}%
\pgfpathlineto{\pgfqpoint{3.038110in}{2.414488in}}%
\pgfusepath{stroke}%
\end{pgfscope}%
\begin{pgfscope}%
\pgfsetrectcap%
\pgfsetmiterjoin%
\pgfsetlinewidth{0.803000pt}%
\definecolor{currentstroke}{rgb}{0.000000,0.000000,0.000000}%
\pgfsetstrokecolor{currentstroke}%
\pgfsetdash{}{0pt}%
\pgfpathmoveto{\pgfqpoint{0.650000in}{0.420000in}}%
\pgfpathlineto{\pgfqpoint{3.038110in}{0.420000in}}%
\pgfusepath{stroke}%
\end{pgfscope}%
\begin{pgfscope}%
\pgfsetrectcap%
\pgfsetmiterjoin%
\pgfsetlinewidth{0.803000pt}%
\definecolor{currentstroke}{rgb}{0.000000,0.000000,0.000000}%
\pgfsetstrokecolor{currentstroke}%
\pgfsetdash{}{0pt}%
\pgfpathmoveto{\pgfqpoint{0.650000in}{2.414488in}}%
\pgfpathlineto{\pgfqpoint{3.038110in}{2.414488in}}%
\pgfusepath{stroke}%
\end{pgfscope}%
\begin{pgfscope}%
\pgfsetrectcap%
\pgfsetroundjoin%
\pgfsetlinewidth{1.003750pt}%
\definecolor{currentstroke}{rgb}{0.870588,0.466667,0.682353}%
\pgfsetstrokecolor{currentstroke}%
\pgfsetdash{}{0pt}%
\pgfpathmoveto{\pgfqpoint{0.750000in}{2.275599in}}%
\pgfpathlineto{\pgfqpoint{0.861111in}{2.275599in}}%
\pgfpathlineto{\pgfqpoint{0.972222in}{2.275599in}}%
\pgfusepath{stroke}%
\end{pgfscope}%
\begin{pgfscope}%
\definecolor{textcolor}{rgb}{0.000000,0.000000,0.000000}%
\pgfsetstrokecolor{textcolor}%
\pgfsetfillcolor{textcolor}%
\pgftext[x=1.061111in,y=2.236710in,left,base]{\color{textcolor}{\rmfamily\fontsize{8.000000}{9.600000}\selectfont\catcode`\^=\active\def^{\ifmmode\sp\else\^{}\fi}\catcode`\%=\active\def%{\%}base}}%
\end{pgfscope}%
\begin{pgfscope}%
\pgfsetrectcap%
\pgfsetroundjoin%
\pgfsetlinewidth{1.003750pt}%
\definecolor{currentstroke}{rgb}{0.498039,0.737255,0.254902}%
\pgfsetstrokecolor{currentstroke}%
\pgfsetdash{}{0pt}%
\pgfpathmoveto{\pgfqpoint{0.750000in}{2.120661in}}%
\pgfpathlineto{\pgfqpoint{0.861111in}{2.120661in}}%
\pgfpathlineto{\pgfqpoint{0.972222in}{2.120661in}}%
\pgfusepath{stroke}%
\end{pgfscope}%
\begin{pgfscope}%
\definecolor{textcolor}{rgb}{0.000000,0.000000,0.000000}%
\pgfsetstrokecolor{textcolor}%
\pgfsetfillcolor{textcolor}%
\pgftext[x=1.061111in,y=2.081772in,left,base]{\color{textcolor}{\rmfamily\fontsize{8.000000}{9.600000}\selectfont\catcode`\^=\active\def^{\ifmmode\sp\else\^{}\fi}\catcode`\%=\active\def%{\%}optimised}}%
\end{pgfscope}%
\end{pgfpicture}%
\makeatother%
\endgroup%

    \label{fig:R100deg30Re1200Beta}
\end{subfigure}
\hfill
\begin{subfigure}[t]{0.495\textwidth}
    \subcaption{}
    \centering
    %% Creator: Matplotlib, PGF backend
%%
%% To include the figure in your LaTeX document, write
%%   \input{<filename>.pgf}
%%
%% Make sure the required packages are loaded in your preamble
%%   \usepackage{pgf}
%%
%% Also ensure that all the required font packages are loaded; for instance,
%% the lmodern package is sometimes necessary when using math font.
%%   \usepackage{lmodern}
%%
%% Figures using additional raster images can only be included by \input if
%% they are in the same directory as the main LaTeX file. For loading figures
%% from other directories you can use the `import` package
%%   \usepackage{import}
%%
%% and then include the figures with
%%   \import{<path to file>}{<filename>.pgf}
%%
%% Matplotlib used the following preamble
%%   \def\mathdefault#1{#1}
%%   \everymath=\expandafter{\the\everymath\displaystyle}
%%   \IfFileExists{scrextend.sty}{
%%     \usepackage[fontsize=11.000000pt]{scrextend}
%%   }{
%%     \renewcommand{\normalsize}{\fontsize{11.000000}{13.200000}\selectfont}
%%     \normalsize
%%   }
%%   
%%   \makeatletter\@ifpackageloaded{underscore}{}{\usepackage[strings]{underscore}}\makeatother
%%
\begingroup%
\makeatletter%
\begin{pgfpicture}%
\pgfpathrectangle{\pgfpointorigin}{\pgfqpoint{3.118110in}{2.494488in}}%
\pgfusepath{use as bounding box, clip}%
\begin{pgfscope}%
\pgfsetbuttcap%
\pgfsetmiterjoin%
\definecolor{currentfill}{rgb}{1.000000,1.000000,1.000000}%
\pgfsetfillcolor{currentfill}%
\pgfsetlinewidth{0.000000pt}%
\definecolor{currentstroke}{rgb}{1.000000,1.000000,1.000000}%
\pgfsetstrokecolor{currentstroke}%
\pgfsetdash{}{0pt}%
\pgfpathmoveto{\pgfqpoint{0.000000in}{0.000000in}}%
\pgfpathlineto{\pgfqpoint{3.118110in}{0.000000in}}%
\pgfpathlineto{\pgfqpoint{3.118110in}{2.494488in}}%
\pgfpathlineto{\pgfqpoint{0.000000in}{2.494488in}}%
\pgfpathlineto{\pgfqpoint{0.000000in}{0.000000in}}%
\pgfpathclose%
\pgfusepath{fill}%
\end{pgfscope}%
\begin{pgfscope}%
\pgfsetbuttcap%
\pgfsetmiterjoin%
\definecolor{currentfill}{rgb}{1.000000,1.000000,1.000000}%
\pgfsetfillcolor{currentfill}%
\pgfsetlinewidth{0.000000pt}%
\definecolor{currentstroke}{rgb}{0.000000,0.000000,0.000000}%
\pgfsetstrokecolor{currentstroke}%
\pgfsetstrokeopacity{0.000000}%
\pgfsetdash{}{0pt}%
\pgfpathmoveto{\pgfqpoint{0.650000in}{0.420000in}}%
\pgfpathlineto{\pgfqpoint{3.038110in}{0.420000in}}%
\pgfpathlineto{\pgfqpoint{3.038110in}{2.414488in}}%
\pgfpathlineto{\pgfqpoint{0.650000in}{2.414488in}}%
\pgfpathlineto{\pgfqpoint{0.650000in}{0.420000in}}%
\pgfpathclose%
\pgfusepath{fill}%
\end{pgfscope}%
\begin{pgfscope}%
\pgfsetbuttcap%
\pgfsetroundjoin%
\definecolor{currentfill}{rgb}{0.000000,0.000000,0.000000}%
\pgfsetfillcolor{currentfill}%
\pgfsetlinewidth{0.501875pt}%
\definecolor{currentstroke}{rgb}{0.000000,0.000000,0.000000}%
\pgfsetstrokecolor{currentstroke}%
\pgfsetdash{}{0pt}%
\pgfsys@defobject{currentmarker}{\pgfqpoint{0.000000in}{0.000000in}}{\pgfqpoint{0.000000in}{0.041667in}}{%
\pgfpathmoveto{\pgfqpoint{0.000000in}{0.000000in}}%
\pgfpathlineto{\pgfqpoint{0.000000in}{0.041667in}}%
\pgfusepath{stroke,fill}%
}%
\begin{pgfscope}%
\pgfsys@transformshift{0.650000in}{0.420000in}%
\pgfsys@useobject{currentmarker}{}%
\end{pgfscope}%
\end{pgfscope}%
\begin{pgfscope}%
\pgfsetbuttcap%
\pgfsetroundjoin%
\definecolor{currentfill}{rgb}{0.000000,0.000000,0.000000}%
\pgfsetfillcolor{currentfill}%
\pgfsetlinewidth{0.501875pt}%
\definecolor{currentstroke}{rgb}{0.000000,0.000000,0.000000}%
\pgfsetstrokecolor{currentstroke}%
\pgfsetdash{}{0pt}%
\pgfsys@defobject{currentmarker}{\pgfqpoint{0.000000in}{-0.041667in}}{\pgfqpoint{0.000000in}{0.000000in}}{%
\pgfpathmoveto{\pgfqpoint{0.000000in}{0.000000in}}%
\pgfpathlineto{\pgfqpoint{0.000000in}{-0.041667in}}%
\pgfusepath{stroke,fill}%
}%
\begin{pgfscope}%
\pgfsys@transformshift{0.650000in}{2.414488in}%
\pgfsys@useobject{currentmarker}{}%
\end{pgfscope}%
\end{pgfscope}%
\begin{pgfscope}%
\definecolor{textcolor}{rgb}{0.000000,0.000000,0.000000}%
\pgfsetstrokecolor{textcolor}%
\pgfsetfillcolor{textcolor}%
\pgftext[x=0.650000in,y=0.371389in,,top]{\color{textcolor}{\rmfamily\fontsize{11.000000}{13.200000}\selectfont\catcode`\^=\active\def^{\ifmmode\sp\else\^{}\fi}\catcode`\%=\active\def%{\%}0}}%
\end{pgfscope}%
\begin{pgfscope}%
\pgfsetbuttcap%
\pgfsetroundjoin%
\definecolor{currentfill}{rgb}{0.000000,0.000000,0.000000}%
\pgfsetfillcolor{currentfill}%
\pgfsetlinewidth{0.501875pt}%
\definecolor{currentstroke}{rgb}{0.000000,0.000000,0.000000}%
\pgfsetstrokecolor{currentstroke}%
\pgfsetdash{}{0pt}%
\pgfsys@defobject{currentmarker}{\pgfqpoint{0.000000in}{0.000000in}}{\pgfqpoint{0.000000in}{0.041667in}}{%
\pgfpathmoveto{\pgfqpoint{0.000000in}{0.000000in}}%
\pgfpathlineto{\pgfqpoint{0.000000in}{0.041667in}}%
\pgfusepath{stroke,fill}%
}%
\begin{pgfscope}%
\pgfsys@transformshift{1.492631in}{0.420000in}%
\pgfsys@useobject{currentmarker}{}%
\end{pgfscope}%
\end{pgfscope}%
\begin{pgfscope}%
\pgfsetbuttcap%
\pgfsetroundjoin%
\definecolor{currentfill}{rgb}{0.000000,0.000000,0.000000}%
\pgfsetfillcolor{currentfill}%
\pgfsetlinewidth{0.501875pt}%
\definecolor{currentstroke}{rgb}{0.000000,0.000000,0.000000}%
\pgfsetstrokecolor{currentstroke}%
\pgfsetdash{}{0pt}%
\pgfsys@defobject{currentmarker}{\pgfqpoint{0.000000in}{-0.041667in}}{\pgfqpoint{0.000000in}{0.000000in}}{%
\pgfpathmoveto{\pgfqpoint{0.000000in}{0.000000in}}%
\pgfpathlineto{\pgfqpoint{0.000000in}{-0.041667in}}%
\pgfusepath{stroke,fill}%
}%
\begin{pgfscope}%
\pgfsys@transformshift{1.492631in}{2.414488in}%
\pgfsys@useobject{currentmarker}{}%
\end{pgfscope}%
\end{pgfscope}%
\begin{pgfscope}%
\definecolor{textcolor}{rgb}{0.000000,0.000000,0.000000}%
\pgfsetstrokecolor{textcolor}%
\pgfsetfillcolor{textcolor}%
\pgftext[x=1.492631in,y=0.371389in,,top]{\color{textcolor}{\rmfamily\fontsize{11.000000}{13.200000}\selectfont\catcode`\^=\active\def^{\ifmmode\sp\else\^{}\fi}\catcode`\%=\active\def%{\%}50}}%
\end{pgfscope}%
\begin{pgfscope}%
\pgfsetbuttcap%
\pgfsetroundjoin%
\definecolor{currentfill}{rgb}{0.000000,0.000000,0.000000}%
\pgfsetfillcolor{currentfill}%
\pgfsetlinewidth{0.501875pt}%
\definecolor{currentstroke}{rgb}{0.000000,0.000000,0.000000}%
\pgfsetstrokecolor{currentstroke}%
\pgfsetdash{}{0pt}%
\pgfsys@defobject{currentmarker}{\pgfqpoint{0.000000in}{0.000000in}}{\pgfqpoint{0.000000in}{0.041667in}}{%
\pgfpathmoveto{\pgfqpoint{0.000000in}{0.000000in}}%
\pgfpathlineto{\pgfqpoint{0.000000in}{0.041667in}}%
\pgfusepath{stroke,fill}%
}%
\begin{pgfscope}%
\pgfsys@transformshift{2.335263in}{0.420000in}%
\pgfsys@useobject{currentmarker}{}%
\end{pgfscope}%
\end{pgfscope}%
\begin{pgfscope}%
\pgfsetbuttcap%
\pgfsetroundjoin%
\definecolor{currentfill}{rgb}{0.000000,0.000000,0.000000}%
\pgfsetfillcolor{currentfill}%
\pgfsetlinewidth{0.501875pt}%
\definecolor{currentstroke}{rgb}{0.000000,0.000000,0.000000}%
\pgfsetstrokecolor{currentstroke}%
\pgfsetdash{}{0pt}%
\pgfsys@defobject{currentmarker}{\pgfqpoint{0.000000in}{-0.041667in}}{\pgfqpoint{0.000000in}{0.000000in}}{%
\pgfpathmoveto{\pgfqpoint{0.000000in}{0.000000in}}%
\pgfpathlineto{\pgfqpoint{0.000000in}{-0.041667in}}%
\pgfusepath{stroke,fill}%
}%
\begin{pgfscope}%
\pgfsys@transformshift{2.335263in}{2.414488in}%
\pgfsys@useobject{currentmarker}{}%
\end{pgfscope}%
\end{pgfscope}%
\begin{pgfscope}%
\definecolor{textcolor}{rgb}{0.000000,0.000000,0.000000}%
\pgfsetstrokecolor{textcolor}%
\pgfsetfillcolor{textcolor}%
\pgftext[x=2.335263in,y=0.371389in,,top]{\color{textcolor}{\rmfamily\fontsize{11.000000}{13.200000}\selectfont\catcode`\^=\active\def^{\ifmmode\sp\else\^{}\fi}\catcode`\%=\active\def%{\%}100}}%
\end{pgfscope}%
\begin{pgfscope}%
\pgfsetbuttcap%
\pgfsetroundjoin%
\definecolor{currentfill}{rgb}{0.000000,0.000000,0.000000}%
\pgfsetfillcolor{currentfill}%
\pgfsetlinewidth{0.401500pt}%
\definecolor{currentstroke}{rgb}{0.000000,0.000000,0.000000}%
\pgfsetstrokecolor{currentstroke}%
\pgfsetdash{}{0pt}%
\pgfsys@defobject{currentmarker}{\pgfqpoint{0.000000in}{0.000000in}}{\pgfqpoint{0.000000in}{0.020833in}}{%
\pgfpathmoveto{\pgfqpoint{0.000000in}{0.000000in}}%
\pgfpathlineto{\pgfqpoint{0.000000in}{0.020833in}}%
\pgfusepath{stroke,fill}%
}%
\begin{pgfscope}%
\pgfsys@transformshift{0.818526in}{0.420000in}%
\pgfsys@useobject{currentmarker}{}%
\end{pgfscope}%
\end{pgfscope}%
\begin{pgfscope}%
\pgfsetbuttcap%
\pgfsetroundjoin%
\definecolor{currentfill}{rgb}{0.000000,0.000000,0.000000}%
\pgfsetfillcolor{currentfill}%
\pgfsetlinewidth{0.401500pt}%
\definecolor{currentstroke}{rgb}{0.000000,0.000000,0.000000}%
\pgfsetstrokecolor{currentstroke}%
\pgfsetdash{}{0pt}%
\pgfsys@defobject{currentmarker}{\pgfqpoint{0.000000in}{-0.020833in}}{\pgfqpoint{0.000000in}{0.000000in}}{%
\pgfpathmoveto{\pgfqpoint{0.000000in}{0.000000in}}%
\pgfpathlineto{\pgfqpoint{0.000000in}{-0.020833in}}%
\pgfusepath{stroke,fill}%
}%
\begin{pgfscope}%
\pgfsys@transformshift{0.818526in}{2.414488in}%
\pgfsys@useobject{currentmarker}{}%
\end{pgfscope}%
\end{pgfscope}%
\begin{pgfscope}%
\pgfsetbuttcap%
\pgfsetroundjoin%
\definecolor{currentfill}{rgb}{0.000000,0.000000,0.000000}%
\pgfsetfillcolor{currentfill}%
\pgfsetlinewidth{0.401500pt}%
\definecolor{currentstroke}{rgb}{0.000000,0.000000,0.000000}%
\pgfsetstrokecolor{currentstroke}%
\pgfsetdash{}{0pt}%
\pgfsys@defobject{currentmarker}{\pgfqpoint{0.000000in}{0.000000in}}{\pgfqpoint{0.000000in}{0.020833in}}{%
\pgfpathmoveto{\pgfqpoint{0.000000in}{0.000000in}}%
\pgfpathlineto{\pgfqpoint{0.000000in}{0.020833in}}%
\pgfusepath{stroke,fill}%
}%
\begin{pgfscope}%
\pgfsys@transformshift{0.987053in}{0.420000in}%
\pgfsys@useobject{currentmarker}{}%
\end{pgfscope}%
\end{pgfscope}%
\begin{pgfscope}%
\pgfsetbuttcap%
\pgfsetroundjoin%
\definecolor{currentfill}{rgb}{0.000000,0.000000,0.000000}%
\pgfsetfillcolor{currentfill}%
\pgfsetlinewidth{0.401500pt}%
\definecolor{currentstroke}{rgb}{0.000000,0.000000,0.000000}%
\pgfsetstrokecolor{currentstroke}%
\pgfsetdash{}{0pt}%
\pgfsys@defobject{currentmarker}{\pgfqpoint{0.000000in}{-0.020833in}}{\pgfqpoint{0.000000in}{0.000000in}}{%
\pgfpathmoveto{\pgfqpoint{0.000000in}{0.000000in}}%
\pgfpathlineto{\pgfqpoint{0.000000in}{-0.020833in}}%
\pgfusepath{stroke,fill}%
}%
\begin{pgfscope}%
\pgfsys@transformshift{0.987053in}{2.414488in}%
\pgfsys@useobject{currentmarker}{}%
\end{pgfscope}%
\end{pgfscope}%
\begin{pgfscope}%
\pgfsetbuttcap%
\pgfsetroundjoin%
\definecolor{currentfill}{rgb}{0.000000,0.000000,0.000000}%
\pgfsetfillcolor{currentfill}%
\pgfsetlinewidth{0.401500pt}%
\definecolor{currentstroke}{rgb}{0.000000,0.000000,0.000000}%
\pgfsetstrokecolor{currentstroke}%
\pgfsetdash{}{0pt}%
\pgfsys@defobject{currentmarker}{\pgfqpoint{0.000000in}{0.000000in}}{\pgfqpoint{0.000000in}{0.020833in}}{%
\pgfpathmoveto{\pgfqpoint{0.000000in}{0.000000in}}%
\pgfpathlineto{\pgfqpoint{0.000000in}{0.020833in}}%
\pgfusepath{stroke,fill}%
}%
\begin{pgfscope}%
\pgfsys@transformshift{1.155579in}{0.420000in}%
\pgfsys@useobject{currentmarker}{}%
\end{pgfscope}%
\end{pgfscope}%
\begin{pgfscope}%
\pgfsetbuttcap%
\pgfsetroundjoin%
\definecolor{currentfill}{rgb}{0.000000,0.000000,0.000000}%
\pgfsetfillcolor{currentfill}%
\pgfsetlinewidth{0.401500pt}%
\definecolor{currentstroke}{rgb}{0.000000,0.000000,0.000000}%
\pgfsetstrokecolor{currentstroke}%
\pgfsetdash{}{0pt}%
\pgfsys@defobject{currentmarker}{\pgfqpoint{0.000000in}{-0.020833in}}{\pgfqpoint{0.000000in}{0.000000in}}{%
\pgfpathmoveto{\pgfqpoint{0.000000in}{0.000000in}}%
\pgfpathlineto{\pgfqpoint{0.000000in}{-0.020833in}}%
\pgfusepath{stroke,fill}%
}%
\begin{pgfscope}%
\pgfsys@transformshift{1.155579in}{2.414488in}%
\pgfsys@useobject{currentmarker}{}%
\end{pgfscope}%
\end{pgfscope}%
\begin{pgfscope}%
\pgfsetbuttcap%
\pgfsetroundjoin%
\definecolor{currentfill}{rgb}{0.000000,0.000000,0.000000}%
\pgfsetfillcolor{currentfill}%
\pgfsetlinewidth{0.401500pt}%
\definecolor{currentstroke}{rgb}{0.000000,0.000000,0.000000}%
\pgfsetstrokecolor{currentstroke}%
\pgfsetdash{}{0pt}%
\pgfsys@defobject{currentmarker}{\pgfqpoint{0.000000in}{0.000000in}}{\pgfqpoint{0.000000in}{0.020833in}}{%
\pgfpathmoveto{\pgfqpoint{0.000000in}{0.000000in}}%
\pgfpathlineto{\pgfqpoint{0.000000in}{0.020833in}}%
\pgfusepath{stroke,fill}%
}%
\begin{pgfscope}%
\pgfsys@transformshift{1.324105in}{0.420000in}%
\pgfsys@useobject{currentmarker}{}%
\end{pgfscope}%
\end{pgfscope}%
\begin{pgfscope}%
\pgfsetbuttcap%
\pgfsetroundjoin%
\definecolor{currentfill}{rgb}{0.000000,0.000000,0.000000}%
\pgfsetfillcolor{currentfill}%
\pgfsetlinewidth{0.401500pt}%
\definecolor{currentstroke}{rgb}{0.000000,0.000000,0.000000}%
\pgfsetstrokecolor{currentstroke}%
\pgfsetdash{}{0pt}%
\pgfsys@defobject{currentmarker}{\pgfqpoint{0.000000in}{-0.020833in}}{\pgfqpoint{0.000000in}{0.000000in}}{%
\pgfpathmoveto{\pgfqpoint{0.000000in}{0.000000in}}%
\pgfpathlineto{\pgfqpoint{0.000000in}{-0.020833in}}%
\pgfusepath{stroke,fill}%
}%
\begin{pgfscope}%
\pgfsys@transformshift{1.324105in}{2.414488in}%
\pgfsys@useobject{currentmarker}{}%
\end{pgfscope}%
\end{pgfscope}%
\begin{pgfscope}%
\pgfsetbuttcap%
\pgfsetroundjoin%
\definecolor{currentfill}{rgb}{0.000000,0.000000,0.000000}%
\pgfsetfillcolor{currentfill}%
\pgfsetlinewidth{0.401500pt}%
\definecolor{currentstroke}{rgb}{0.000000,0.000000,0.000000}%
\pgfsetstrokecolor{currentstroke}%
\pgfsetdash{}{0pt}%
\pgfsys@defobject{currentmarker}{\pgfqpoint{0.000000in}{0.000000in}}{\pgfqpoint{0.000000in}{0.020833in}}{%
\pgfpathmoveto{\pgfqpoint{0.000000in}{0.000000in}}%
\pgfpathlineto{\pgfqpoint{0.000000in}{0.020833in}}%
\pgfusepath{stroke,fill}%
}%
\begin{pgfscope}%
\pgfsys@transformshift{1.661158in}{0.420000in}%
\pgfsys@useobject{currentmarker}{}%
\end{pgfscope}%
\end{pgfscope}%
\begin{pgfscope}%
\pgfsetbuttcap%
\pgfsetroundjoin%
\definecolor{currentfill}{rgb}{0.000000,0.000000,0.000000}%
\pgfsetfillcolor{currentfill}%
\pgfsetlinewidth{0.401500pt}%
\definecolor{currentstroke}{rgb}{0.000000,0.000000,0.000000}%
\pgfsetstrokecolor{currentstroke}%
\pgfsetdash{}{0pt}%
\pgfsys@defobject{currentmarker}{\pgfqpoint{0.000000in}{-0.020833in}}{\pgfqpoint{0.000000in}{0.000000in}}{%
\pgfpathmoveto{\pgfqpoint{0.000000in}{0.000000in}}%
\pgfpathlineto{\pgfqpoint{0.000000in}{-0.020833in}}%
\pgfusepath{stroke,fill}%
}%
\begin{pgfscope}%
\pgfsys@transformshift{1.661158in}{2.414488in}%
\pgfsys@useobject{currentmarker}{}%
\end{pgfscope}%
\end{pgfscope}%
\begin{pgfscope}%
\pgfsetbuttcap%
\pgfsetroundjoin%
\definecolor{currentfill}{rgb}{0.000000,0.000000,0.000000}%
\pgfsetfillcolor{currentfill}%
\pgfsetlinewidth{0.401500pt}%
\definecolor{currentstroke}{rgb}{0.000000,0.000000,0.000000}%
\pgfsetstrokecolor{currentstroke}%
\pgfsetdash{}{0pt}%
\pgfsys@defobject{currentmarker}{\pgfqpoint{0.000000in}{0.000000in}}{\pgfqpoint{0.000000in}{0.020833in}}{%
\pgfpathmoveto{\pgfqpoint{0.000000in}{0.000000in}}%
\pgfpathlineto{\pgfqpoint{0.000000in}{0.020833in}}%
\pgfusepath{stroke,fill}%
}%
\begin{pgfscope}%
\pgfsys@transformshift{1.829684in}{0.420000in}%
\pgfsys@useobject{currentmarker}{}%
\end{pgfscope}%
\end{pgfscope}%
\begin{pgfscope}%
\pgfsetbuttcap%
\pgfsetroundjoin%
\definecolor{currentfill}{rgb}{0.000000,0.000000,0.000000}%
\pgfsetfillcolor{currentfill}%
\pgfsetlinewidth{0.401500pt}%
\definecolor{currentstroke}{rgb}{0.000000,0.000000,0.000000}%
\pgfsetstrokecolor{currentstroke}%
\pgfsetdash{}{0pt}%
\pgfsys@defobject{currentmarker}{\pgfqpoint{0.000000in}{-0.020833in}}{\pgfqpoint{0.000000in}{0.000000in}}{%
\pgfpathmoveto{\pgfqpoint{0.000000in}{0.000000in}}%
\pgfpathlineto{\pgfqpoint{0.000000in}{-0.020833in}}%
\pgfusepath{stroke,fill}%
}%
\begin{pgfscope}%
\pgfsys@transformshift{1.829684in}{2.414488in}%
\pgfsys@useobject{currentmarker}{}%
\end{pgfscope}%
\end{pgfscope}%
\begin{pgfscope}%
\pgfsetbuttcap%
\pgfsetroundjoin%
\definecolor{currentfill}{rgb}{0.000000,0.000000,0.000000}%
\pgfsetfillcolor{currentfill}%
\pgfsetlinewidth{0.401500pt}%
\definecolor{currentstroke}{rgb}{0.000000,0.000000,0.000000}%
\pgfsetstrokecolor{currentstroke}%
\pgfsetdash{}{0pt}%
\pgfsys@defobject{currentmarker}{\pgfqpoint{0.000000in}{0.000000in}}{\pgfqpoint{0.000000in}{0.020833in}}{%
\pgfpathmoveto{\pgfqpoint{0.000000in}{0.000000in}}%
\pgfpathlineto{\pgfqpoint{0.000000in}{0.020833in}}%
\pgfusepath{stroke,fill}%
}%
\begin{pgfscope}%
\pgfsys@transformshift{1.998210in}{0.420000in}%
\pgfsys@useobject{currentmarker}{}%
\end{pgfscope}%
\end{pgfscope}%
\begin{pgfscope}%
\pgfsetbuttcap%
\pgfsetroundjoin%
\definecolor{currentfill}{rgb}{0.000000,0.000000,0.000000}%
\pgfsetfillcolor{currentfill}%
\pgfsetlinewidth{0.401500pt}%
\definecolor{currentstroke}{rgb}{0.000000,0.000000,0.000000}%
\pgfsetstrokecolor{currentstroke}%
\pgfsetdash{}{0pt}%
\pgfsys@defobject{currentmarker}{\pgfqpoint{0.000000in}{-0.020833in}}{\pgfqpoint{0.000000in}{0.000000in}}{%
\pgfpathmoveto{\pgfqpoint{0.000000in}{0.000000in}}%
\pgfpathlineto{\pgfqpoint{0.000000in}{-0.020833in}}%
\pgfusepath{stroke,fill}%
}%
\begin{pgfscope}%
\pgfsys@transformshift{1.998210in}{2.414488in}%
\pgfsys@useobject{currentmarker}{}%
\end{pgfscope}%
\end{pgfscope}%
\begin{pgfscope}%
\pgfsetbuttcap%
\pgfsetroundjoin%
\definecolor{currentfill}{rgb}{0.000000,0.000000,0.000000}%
\pgfsetfillcolor{currentfill}%
\pgfsetlinewidth{0.401500pt}%
\definecolor{currentstroke}{rgb}{0.000000,0.000000,0.000000}%
\pgfsetstrokecolor{currentstroke}%
\pgfsetdash{}{0pt}%
\pgfsys@defobject{currentmarker}{\pgfqpoint{0.000000in}{0.000000in}}{\pgfqpoint{0.000000in}{0.020833in}}{%
\pgfpathmoveto{\pgfqpoint{0.000000in}{0.000000in}}%
\pgfpathlineto{\pgfqpoint{0.000000in}{0.020833in}}%
\pgfusepath{stroke,fill}%
}%
\begin{pgfscope}%
\pgfsys@transformshift{2.166737in}{0.420000in}%
\pgfsys@useobject{currentmarker}{}%
\end{pgfscope}%
\end{pgfscope}%
\begin{pgfscope}%
\pgfsetbuttcap%
\pgfsetroundjoin%
\definecolor{currentfill}{rgb}{0.000000,0.000000,0.000000}%
\pgfsetfillcolor{currentfill}%
\pgfsetlinewidth{0.401500pt}%
\definecolor{currentstroke}{rgb}{0.000000,0.000000,0.000000}%
\pgfsetstrokecolor{currentstroke}%
\pgfsetdash{}{0pt}%
\pgfsys@defobject{currentmarker}{\pgfqpoint{0.000000in}{-0.020833in}}{\pgfqpoint{0.000000in}{0.000000in}}{%
\pgfpathmoveto{\pgfqpoint{0.000000in}{0.000000in}}%
\pgfpathlineto{\pgfqpoint{0.000000in}{-0.020833in}}%
\pgfusepath{stroke,fill}%
}%
\begin{pgfscope}%
\pgfsys@transformshift{2.166737in}{2.414488in}%
\pgfsys@useobject{currentmarker}{}%
\end{pgfscope}%
\end{pgfscope}%
\begin{pgfscope}%
\pgfsetbuttcap%
\pgfsetroundjoin%
\definecolor{currentfill}{rgb}{0.000000,0.000000,0.000000}%
\pgfsetfillcolor{currentfill}%
\pgfsetlinewidth{0.401500pt}%
\definecolor{currentstroke}{rgb}{0.000000,0.000000,0.000000}%
\pgfsetstrokecolor{currentstroke}%
\pgfsetdash{}{0pt}%
\pgfsys@defobject{currentmarker}{\pgfqpoint{0.000000in}{0.000000in}}{\pgfqpoint{0.000000in}{0.020833in}}{%
\pgfpathmoveto{\pgfqpoint{0.000000in}{0.000000in}}%
\pgfpathlineto{\pgfqpoint{0.000000in}{0.020833in}}%
\pgfusepath{stroke,fill}%
}%
\begin{pgfscope}%
\pgfsys@transformshift{2.503789in}{0.420000in}%
\pgfsys@useobject{currentmarker}{}%
\end{pgfscope}%
\end{pgfscope}%
\begin{pgfscope}%
\pgfsetbuttcap%
\pgfsetroundjoin%
\definecolor{currentfill}{rgb}{0.000000,0.000000,0.000000}%
\pgfsetfillcolor{currentfill}%
\pgfsetlinewidth{0.401500pt}%
\definecolor{currentstroke}{rgb}{0.000000,0.000000,0.000000}%
\pgfsetstrokecolor{currentstroke}%
\pgfsetdash{}{0pt}%
\pgfsys@defobject{currentmarker}{\pgfqpoint{0.000000in}{-0.020833in}}{\pgfqpoint{0.000000in}{0.000000in}}{%
\pgfpathmoveto{\pgfqpoint{0.000000in}{0.000000in}}%
\pgfpathlineto{\pgfqpoint{0.000000in}{-0.020833in}}%
\pgfusepath{stroke,fill}%
}%
\begin{pgfscope}%
\pgfsys@transformshift{2.503789in}{2.414488in}%
\pgfsys@useobject{currentmarker}{}%
\end{pgfscope}%
\end{pgfscope}%
\begin{pgfscope}%
\pgfsetbuttcap%
\pgfsetroundjoin%
\definecolor{currentfill}{rgb}{0.000000,0.000000,0.000000}%
\pgfsetfillcolor{currentfill}%
\pgfsetlinewidth{0.401500pt}%
\definecolor{currentstroke}{rgb}{0.000000,0.000000,0.000000}%
\pgfsetstrokecolor{currentstroke}%
\pgfsetdash{}{0pt}%
\pgfsys@defobject{currentmarker}{\pgfqpoint{0.000000in}{0.000000in}}{\pgfqpoint{0.000000in}{0.020833in}}{%
\pgfpathmoveto{\pgfqpoint{0.000000in}{0.000000in}}%
\pgfpathlineto{\pgfqpoint{0.000000in}{0.020833in}}%
\pgfusepath{stroke,fill}%
}%
\begin{pgfscope}%
\pgfsys@transformshift{2.672315in}{0.420000in}%
\pgfsys@useobject{currentmarker}{}%
\end{pgfscope}%
\end{pgfscope}%
\begin{pgfscope}%
\pgfsetbuttcap%
\pgfsetroundjoin%
\definecolor{currentfill}{rgb}{0.000000,0.000000,0.000000}%
\pgfsetfillcolor{currentfill}%
\pgfsetlinewidth{0.401500pt}%
\definecolor{currentstroke}{rgb}{0.000000,0.000000,0.000000}%
\pgfsetstrokecolor{currentstroke}%
\pgfsetdash{}{0pt}%
\pgfsys@defobject{currentmarker}{\pgfqpoint{0.000000in}{-0.020833in}}{\pgfqpoint{0.000000in}{0.000000in}}{%
\pgfpathmoveto{\pgfqpoint{0.000000in}{0.000000in}}%
\pgfpathlineto{\pgfqpoint{0.000000in}{-0.020833in}}%
\pgfusepath{stroke,fill}%
}%
\begin{pgfscope}%
\pgfsys@transformshift{2.672315in}{2.414488in}%
\pgfsys@useobject{currentmarker}{}%
\end{pgfscope}%
\end{pgfscope}%
\begin{pgfscope}%
\pgfsetbuttcap%
\pgfsetroundjoin%
\definecolor{currentfill}{rgb}{0.000000,0.000000,0.000000}%
\pgfsetfillcolor{currentfill}%
\pgfsetlinewidth{0.401500pt}%
\definecolor{currentstroke}{rgb}{0.000000,0.000000,0.000000}%
\pgfsetstrokecolor{currentstroke}%
\pgfsetdash{}{0pt}%
\pgfsys@defobject{currentmarker}{\pgfqpoint{0.000000in}{0.000000in}}{\pgfqpoint{0.000000in}{0.020833in}}{%
\pgfpathmoveto{\pgfqpoint{0.000000in}{0.000000in}}%
\pgfpathlineto{\pgfqpoint{0.000000in}{0.020833in}}%
\pgfusepath{stroke,fill}%
}%
\begin{pgfscope}%
\pgfsys@transformshift{2.840842in}{0.420000in}%
\pgfsys@useobject{currentmarker}{}%
\end{pgfscope}%
\end{pgfscope}%
\begin{pgfscope}%
\pgfsetbuttcap%
\pgfsetroundjoin%
\definecolor{currentfill}{rgb}{0.000000,0.000000,0.000000}%
\pgfsetfillcolor{currentfill}%
\pgfsetlinewidth{0.401500pt}%
\definecolor{currentstroke}{rgb}{0.000000,0.000000,0.000000}%
\pgfsetstrokecolor{currentstroke}%
\pgfsetdash{}{0pt}%
\pgfsys@defobject{currentmarker}{\pgfqpoint{0.000000in}{-0.020833in}}{\pgfqpoint{0.000000in}{0.000000in}}{%
\pgfpathmoveto{\pgfqpoint{0.000000in}{0.000000in}}%
\pgfpathlineto{\pgfqpoint{0.000000in}{-0.020833in}}%
\pgfusepath{stroke,fill}%
}%
\begin{pgfscope}%
\pgfsys@transformshift{2.840842in}{2.414488in}%
\pgfsys@useobject{currentmarker}{}%
\end{pgfscope}%
\end{pgfscope}%
\begin{pgfscope}%
\pgfsetbuttcap%
\pgfsetroundjoin%
\definecolor{currentfill}{rgb}{0.000000,0.000000,0.000000}%
\pgfsetfillcolor{currentfill}%
\pgfsetlinewidth{0.401500pt}%
\definecolor{currentstroke}{rgb}{0.000000,0.000000,0.000000}%
\pgfsetstrokecolor{currentstroke}%
\pgfsetdash{}{0pt}%
\pgfsys@defobject{currentmarker}{\pgfqpoint{0.000000in}{0.000000in}}{\pgfqpoint{0.000000in}{0.020833in}}{%
\pgfpathmoveto{\pgfqpoint{0.000000in}{0.000000in}}%
\pgfpathlineto{\pgfqpoint{0.000000in}{0.020833in}}%
\pgfusepath{stroke,fill}%
}%
\begin{pgfscope}%
\pgfsys@transformshift{3.009368in}{0.420000in}%
\pgfsys@useobject{currentmarker}{}%
\end{pgfscope}%
\end{pgfscope}%
\begin{pgfscope}%
\pgfsetbuttcap%
\pgfsetroundjoin%
\definecolor{currentfill}{rgb}{0.000000,0.000000,0.000000}%
\pgfsetfillcolor{currentfill}%
\pgfsetlinewidth{0.401500pt}%
\definecolor{currentstroke}{rgb}{0.000000,0.000000,0.000000}%
\pgfsetstrokecolor{currentstroke}%
\pgfsetdash{}{0pt}%
\pgfsys@defobject{currentmarker}{\pgfqpoint{0.000000in}{-0.020833in}}{\pgfqpoint{0.000000in}{0.000000in}}{%
\pgfpathmoveto{\pgfqpoint{0.000000in}{0.000000in}}%
\pgfpathlineto{\pgfqpoint{0.000000in}{-0.020833in}}%
\pgfusepath{stroke,fill}%
}%
\begin{pgfscope}%
\pgfsys@transformshift{3.009368in}{2.414488in}%
\pgfsys@useobject{currentmarker}{}%
\end{pgfscope}%
\end{pgfscope}%
\begin{pgfscope}%
\definecolor{textcolor}{rgb}{0.000000,0.000000,0.000000}%
\pgfsetstrokecolor{textcolor}%
\pgfsetfillcolor{textcolor}%
\pgftext[x=1.844055in,y=0.208426in,,top]{\color{textcolor}{\rmfamily\fontsize{12.000000}{14.400000}\selectfont\catcode`\^=\active\def^{\ifmmode\sp\else\^{}\fi}\catcode`\%=\active\def%{\%}$s/\delta_{0}$}}%
\end{pgfscope}%
\begin{pgfscope}%
\pgfsetbuttcap%
\pgfsetroundjoin%
\definecolor{currentfill}{rgb}{0.000000,0.000000,0.000000}%
\pgfsetfillcolor{currentfill}%
\pgfsetlinewidth{0.501875pt}%
\definecolor{currentstroke}{rgb}{0.000000,0.000000,0.000000}%
\pgfsetstrokecolor{currentstroke}%
\pgfsetdash{}{0pt}%
\pgfsys@defobject{currentmarker}{\pgfqpoint{0.000000in}{0.000000in}}{\pgfqpoint{0.041667in}{0.000000in}}{%
\pgfpathmoveto{\pgfqpoint{0.000000in}{0.000000in}}%
\pgfpathlineto{\pgfqpoint{0.041667in}{0.000000in}}%
\pgfusepath{stroke,fill}%
}%
\begin{pgfscope}%
\pgfsys@transformshift{0.650000in}{0.420000in}%
\pgfsys@useobject{currentmarker}{}%
\end{pgfscope}%
\end{pgfscope}%
\begin{pgfscope}%
\pgfsetbuttcap%
\pgfsetroundjoin%
\definecolor{currentfill}{rgb}{0.000000,0.000000,0.000000}%
\pgfsetfillcolor{currentfill}%
\pgfsetlinewidth{0.501875pt}%
\definecolor{currentstroke}{rgb}{0.000000,0.000000,0.000000}%
\pgfsetstrokecolor{currentstroke}%
\pgfsetdash{}{0pt}%
\pgfsys@defobject{currentmarker}{\pgfqpoint{-0.041667in}{0.000000in}}{\pgfqpoint{-0.000000in}{0.000000in}}{%
\pgfpathmoveto{\pgfqpoint{-0.000000in}{0.000000in}}%
\pgfpathlineto{\pgfqpoint{-0.041667in}{0.000000in}}%
\pgfusepath{stroke,fill}%
}%
\begin{pgfscope}%
\pgfsys@transformshift{3.038110in}{0.420000in}%
\pgfsys@useobject{currentmarker}{}%
\end{pgfscope}%
\end{pgfscope}%
\begin{pgfscope}%
\definecolor{textcolor}{rgb}{0.000000,0.000000,0.000000}%
\pgfsetstrokecolor{textcolor}%
\pgfsetfillcolor{textcolor}%
\pgftext[x=0.212731in, y=0.367193in, left, base]{\color{textcolor}{\rmfamily\fontsize{11.000000}{13.200000}\selectfont\catcode`\^=\active\def^{\ifmmode\sp\else\^{}\fi}\catcode`\%=\active\def%{\%}\ensuremath{-}0.05}}%
\end{pgfscope}%
\begin{pgfscope}%
\pgfsetbuttcap%
\pgfsetroundjoin%
\definecolor{currentfill}{rgb}{0.000000,0.000000,0.000000}%
\pgfsetfillcolor{currentfill}%
\pgfsetlinewidth{0.501875pt}%
\definecolor{currentstroke}{rgb}{0.000000,0.000000,0.000000}%
\pgfsetstrokecolor{currentstroke}%
\pgfsetdash{}{0pt}%
\pgfsys@defobject{currentmarker}{\pgfqpoint{0.000000in}{0.000000in}}{\pgfqpoint{0.041667in}{0.000000in}}{%
\pgfpathmoveto{\pgfqpoint{0.000000in}{0.000000in}}%
\pgfpathlineto{\pgfqpoint{0.041667in}{0.000000in}}%
\pgfusepath{stroke,fill}%
}%
\begin{pgfscope}%
\pgfsys@transformshift{0.650000in}{1.084829in}%
\pgfsys@useobject{currentmarker}{}%
\end{pgfscope}%
\end{pgfscope}%
\begin{pgfscope}%
\pgfsetbuttcap%
\pgfsetroundjoin%
\definecolor{currentfill}{rgb}{0.000000,0.000000,0.000000}%
\pgfsetfillcolor{currentfill}%
\pgfsetlinewidth{0.501875pt}%
\definecolor{currentstroke}{rgb}{0.000000,0.000000,0.000000}%
\pgfsetstrokecolor{currentstroke}%
\pgfsetdash{}{0pt}%
\pgfsys@defobject{currentmarker}{\pgfqpoint{-0.041667in}{0.000000in}}{\pgfqpoint{-0.000000in}{0.000000in}}{%
\pgfpathmoveto{\pgfqpoint{-0.000000in}{0.000000in}}%
\pgfpathlineto{\pgfqpoint{-0.041667in}{0.000000in}}%
\pgfusepath{stroke,fill}%
}%
\begin{pgfscope}%
\pgfsys@transformshift{3.038110in}{1.084829in}%
\pgfsys@useobject{currentmarker}{}%
\end{pgfscope}%
\end{pgfscope}%
\begin{pgfscope}%
\definecolor{textcolor}{rgb}{0.000000,0.000000,0.000000}%
\pgfsetstrokecolor{textcolor}%
\pgfsetfillcolor{textcolor}%
\pgftext[x=0.331018in, y=1.032023in, left, base]{\color{textcolor}{\rmfamily\fontsize{11.000000}{13.200000}\selectfont\catcode`\^=\active\def^{\ifmmode\sp\else\^{}\fi}\catcode`\%=\active\def%{\%}0.00}}%
\end{pgfscope}%
\begin{pgfscope}%
\pgfsetbuttcap%
\pgfsetroundjoin%
\definecolor{currentfill}{rgb}{0.000000,0.000000,0.000000}%
\pgfsetfillcolor{currentfill}%
\pgfsetlinewidth{0.501875pt}%
\definecolor{currentstroke}{rgb}{0.000000,0.000000,0.000000}%
\pgfsetstrokecolor{currentstroke}%
\pgfsetdash{}{0pt}%
\pgfsys@defobject{currentmarker}{\pgfqpoint{0.000000in}{0.000000in}}{\pgfqpoint{0.041667in}{0.000000in}}{%
\pgfpathmoveto{\pgfqpoint{0.000000in}{0.000000in}}%
\pgfpathlineto{\pgfqpoint{0.041667in}{0.000000in}}%
\pgfusepath{stroke,fill}%
}%
\begin{pgfscope}%
\pgfsys@transformshift{0.650000in}{1.749659in}%
\pgfsys@useobject{currentmarker}{}%
\end{pgfscope}%
\end{pgfscope}%
\begin{pgfscope}%
\pgfsetbuttcap%
\pgfsetroundjoin%
\definecolor{currentfill}{rgb}{0.000000,0.000000,0.000000}%
\pgfsetfillcolor{currentfill}%
\pgfsetlinewidth{0.501875pt}%
\definecolor{currentstroke}{rgb}{0.000000,0.000000,0.000000}%
\pgfsetstrokecolor{currentstroke}%
\pgfsetdash{}{0pt}%
\pgfsys@defobject{currentmarker}{\pgfqpoint{-0.041667in}{0.000000in}}{\pgfqpoint{-0.000000in}{0.000000in}}{%
\pgfpathmoveto{\pgfqpoint{-0.000000in}{0.000000in}}%
\pgfpathlineto{\pgfqpoint{-0.041667in}{0.000000in}}%
\pgfusepath{stroke,fill}%
}%
\begin{pgfscope}%
\pgfsys@transformshift{3.038110in}{1.749659in}%
\pgfsys@useobject{currentmarker}{}%
\end{pgfscope}%
\end{pgfscope}%
\begin{pgfscope}%
\definecolor{textcolor}{rgb}{0.000000,0.000000,0.000000}%
\pgfsetstrokecolor{textcolor}%
\pgfsetfillcolor{textcolor}%
\pgftext[x=0.331018in, y=1.696852in, left, base]{\color{textcolor}{\rmfamily\fontsize{11.000000}{13.200000}\selectfont\catcode`\^=\active\def^{\ifmmode\sp\else\^{}\fi}\catcode`\%=\active\def%{\%}0.05}}%
\end{pgfscope}%
\begin{pgfscope}%
\pgfsetbuttcap%
\pgfsetroundjoin%
\definecolor{currentfill}{rgb}{0.000000,0.000000,0.000000}%
\pgfsetfillcolor{currentfill}%
\pgfsetlinewidth{0.501875pt}%
\definecolor{currentstroke}{rgb}{0.000000,0.000000,0.000000}%
\pgfsetstrokecolor{currentstroke}%
\pgfsetdash{}{0pt}%
\pgfsys@defobject{currentmarker}{\pgfqpoint{0.000000in}{0.000000in}}{\pgfqpoint{0.041667in}{0.000000in}}{%
\pgfpathmoveto{\pgfqpoint{0.000000in}{0.000000in}}%
\pgfpathlineto{\pgfqpoint{0.041667in}{0.000000in}}%
\pgfusepath{stroke,fill}%
}%
\begin{pgfscope}%
\pgfsys@transformshift{0.650000in}{2.414488in}%
\pgfsys@useobject{currentmarker}{}%
\end{pgfscope}%
\end{pgfscope}%
\begin{pgfscope}%
\pgfsetbuttcap%
\pgfsetroundjoin%
\definecolor{currentfill}{rgb}{0.000000,0.000000,0.000000}%
\pgfsetfillcolor{currentfill}%
\pgfsetlinewidth{0.501875pt}%
\definecolor{currentstroke}{rgb}{0.000000,0.000000,0.000000}%
\pgfsetstrokecolor{currentstroke}%
\pgfsetdash{}{0pt}%
\pgfsys@defobject{currentmarker}{\pgfqpoint{-0.041667in}{0.000000in}}{\pgfqpoint{-0.000000in}{0.000000in}}{%
\pgfpathmoveto{\pgfqpoint{-0.000000in}{0.000000in}}%
\pgfpathlineto{\pgfqpoint{-0.041667in}{0.000000in}}%
\pgfusepath{stroke,fill}%
}%
\begin{pgfscope}%
\pgfsys@transformshift{3.038110in}{2.414488in}%
\pgfsys@useobject{currentmarker}{}%
\end{pgfscope}%
\end{pgfscope}%
\begin{pgfscope}%
\definecolor{textcolor}{rgb}{0.000000,0.000000,0.000000}%
\pgfsetstrokecolor{textcolor}%
\pgfsetfillcolor{textcolor}%
\pgftext[x=0.331018in, y=2.361681in, left, base]{\color{textcolor}{\rmfamily\fontsize{11.000000}{13.200000}\selectfont\catcode`\^=\active\def^{\ifmmode\sp\else\^{}\fi}\catcode`\%=\active\def%{\%}0.10}}%
\end{pgfscope}%
\begin{pgfscope}%
\pgfsetbuttcap%
\pgfsetroundjoin%
\definecolor{currentfill}{rgb}{0.000000,0.000000,0.000000}%
\pgfsetfillcolor{currentfill}%
\pgfsetlinewidth{0.401500pt}%
\definecolor{currentstroke}{rgb}{0.000000,0.000000,0.000000}%
\pgfsetstrokecolor{currentstroke}%
\pgfsetdash{}{0pt}%
\pgfsys@defobject{currentmarker}{\pgfqpoint{0.000000in}{0.000000in}}{\pgfqpoint{0.020833in}{0.000000in}}{%
\pgfpathmoveto{\pgfqpoint{0.000000in}{0.000000in}}%
\pgfpathlineto{\pgfqpoint{0.020833in}{0.000000in}}%
\pgfusepath{stroke,fill}%
}%
\begin{pgfscope}%
\pgfsys@transformshift{0.650000in}{0.552966in}%
\pgfsys@useobject{currentmarker}{}%
\end{pgfscope}%
\end{pgfscope}%
\begin{pgfscope}%
\pgfsetbuttcap%
\pgfsetroundjoin%
\definecolor{currentfill}{rgb}{0.000000,0.000000,0.000000}%
\pgfsetfillcolor{currentfill}%
\pgfsetlinewidth{0.401500pt}%
\definecolor{currentstroke}{rgb}{0.000000,0.000000,0.000000}%
\pgfsetstrokecolor{currentstroke}%
\pgfsetdash{}{0pt}%
\pgfsys@defobject{currentmarker}{\pgfqpoint{-0.020833in}{0.000000in}}{\pgfqpoint{-0.000000in}{0.000000in}}{%
\pgfpathmoveto{\pgfqpoint{-0.000000in}{0.000000in}}%
\pgfpathlineto{\pgfqpoint{-0.020833in}{0.000000in}}%
\pgfusepath{stroke,fill}%
}%
\begin{pgfscope}%
\pgfsys@transformshift{3.038110in}{0.552966in}%
\pgfsys@useobject{currentmarker}{}%
\end{pgfscope}%
\end{pgfscope}%
\begin{pgfscope}%
\pgfsetbuttcap%
\pgfsetroundjoin%
\definecolor{currentfill}{rgb}{0.000000,0.000000,0.000000}%
\pgfsetfillcolor{currentfill}%
\pgfsetlinewidth{0.401500pt}%
\definecolor{currentstroke}{rgb}{0.000000,0.000000,0.000000}%
\pgfsetstrokecolor{currentstroke}%
\pgfsetdash{}{0pt}%
\pgfsys@defobject{currentmarker}{\pgfqpoint{0.000000in}{0.000000in}}{\pgfqpoint{0.020833in}{0.000000in}}{%
\pgfpathmoveto{\pgfqpoint{0.000000in}{0.000000in}}%
\pgfpathlineto{\pgfqpoint{0.020833in}{0.000000in}}%
\pgfusepath{stroke,fill}%
}%
\begin{pgfscope}%
\pgfsys@transformshift{0.650000in}{0.685932in}%
\pgfsys@useobject{currentmarker}{}%
\end{pgfscope}%
\end{pgfscope}%
\begin{pgfscope}%
\pgfsetbuttcap%
\pgfsetroundjoin%
\definecolor{currentfill}{rgb}{0.000000,0.000000,0.000000}%
\pgfsetfillcolor{currentfill}%
\pgfsetlinewidth{0.401500pt}%
\definecolor{currentstroke}{rgb}{0.000000,0.000000,0.000000}%
\pgfsetstrokecolor{currentstroke}%
\pgfsetdash{}{0pt}%
\pgfsys@defobject{currentmarker}{\pgfqpoint{-0.020833in}{0.000000in}}{\pgfqpoint{-0.000000in}{0.000000in}}{%
\pgfpathmoveto{\pgfqpoint{-0.000000in}{0.000000in}}%
\pgfpathlineto{\pgfqpoint{-0.020833in}{0.000000in}}%
\pgfusepath{stroke,fill}%
}%
\begin{pgfscope}%
\pgfsys@transformshift{3.038110in}{0.685932in}%
\pgfsys@useobject{currentmarker}{}%
\end{pgfscope}%
\end{pgfscope}%
\begin{pgfscope}%
\pgfsetbuttcap%
\pgfsetroundjoin%
\definecolor{currentfill}{rgb}{0.000000,0.000000,0.000000}%
\pgfsetfillcolor{currentfill}%
\pgfsetlinewidth{0.401500pt}%
\definecolor{currentstroke}{rgb}{0.000000,0.000000,0.000000}%
\pgfsetstrokecolor{currentstroke}%
\pgfsetdash{}{0pt}%
\pgfsys@defobject{currentmarker}{\pgfqpoint{0.000000in}{0.000000in}}{\pgfqpoint{0.020833in}{0.000000in}}{%
\pgfpathmoveto{\pgfqpoint{0.000000in}{0.000000in}}%
\pgfpathlineto{\pgfqpoint{0.020833in}{0.000000in}}%
\pgfusepath{stroke,fill}%
}%
\begin{pgfscope}%
\pgfsys@transformshift{0.650000in}{0.818898in}%
\pgfsys@useobject{currentmarker}{}%
\end{pgfscope}%
\end{pgfscope}%
\begin{pgfscope}%
\pgfsetbuttcap%
\pgfsetroundjoin%
\definecolor{currentfill}{rgb}{0.000000,0.000000,0.000000}%
\pgfsetfillcolor{currentfill}%
\pgfsetlinewidth{0.401500pt}%
\definecolor{currentstroke}{rgb}{0.000000,0.000000,0.000000}%
\pgfsetstrokecolor{currentstroke}%
\pgfsetdash{}{0pt}%
\pgfsys@defobject{currentmarker}{\pgfqpoint{-0.020833in}{0.000000in}}{\pgfqpoint{-0.000000in}{0.000000in}}{%
\pgfpathmoveto{\pgfqpoint{-0.000000in}{0.000000in}}%
\pgfpathlineto{\pgfqpoint{-0.020833in}{0.000000in}}%
\pgfusepath{stroke,fill}%
}%
\begin{pgfscope}%
\pgfsys@transformshift{3.038110in}{0.818898in}%
\pgfsys@useobject{currentmarker}{}%
\end{pgfscope}%
\end{pgfscope}%
\begin{pgfscope}%
\pgfsetbuttcap%
\pgfsetroundjoin%
\definecolor{currentfill}{rgb}{0.000000,0.000000,0.000000}%
\pgfsetfillcolor{currentfill}%
\pgfsetlinewidth{0.401500pt}%
\definecolor{currentstroke}{rgb}{0.000000,0.000000,0.000000}%
\pgfsetstrokecolor{currentstroke}%
\pgfsetdash{}{0pt}%
\pgfsys@defobject{currentmarker}{\pgfqpoint{0.000000in}{0.000000in}}{\pgfqpoint{0.020833in}{0.000000in}}{%
\pgfpathmoveto{\pgfqpoint{0.000000in}{0.000000in}}%
\pgfpathlineto{\pgfqpoint{0.020833in}{0.000000in}}%
\pgfusepath{stroke,fill}%
}%
\begin{pgfscope}%
\pgfsys@transformshift{0.650000in}{0.951863in}%
\pgfsys@useobject{currentmarker}{}%
\end{pgfscope}%
\end{pgfscope}%
\begin{pgfscope}%
\pgfsetbuttcap%
\pgfsetroundjoin%
\definecolor{currentfill}{rgb}{0.000000,0.000000,0.000000}%
\pgfsetfillcolor{currentfill}%
\pgfsetlinewidth{0.401500pt}%
\definecolor{currentstroke}{rgb}{0.000000,0.000000,0.000000}%
\pgfsetstrokecolor{currentstroke}%
\pgfsetdash{}{0pt}%
\pgfsys@defobject{currentmarker}{\pgfqpoint{-0.020833in}{0.000000in}}{\pgfqpoint{-0.000000in}{0.000000in}}{%
\pgfpathmoveto{\pgfqpoint{-0.000000in}{0.000000in}}%
\pgfpathlineto{\pgfqpoint{-0.020833in}{0.000000in}}%
\pgfusepath{stroke,fill}%
}%
\begin{pgfscope}%
\pgfsys@transformshift{3.038110in}{0.951863in}%
\pgfsys@useobject{currentmarker}{}%
\end{pgfscope}%
\end{pgfscope}%
\begin{pgfscope}%
\pgfsetbuttcap%
\pgfsetroundjoin%
\definecolor{currentfill}{rgb}{0.000000,0.000000,0.000000}%
\pgfsetfillcolor{currentfill}%
\pgfsetlinewidth{0.401500pt}%
\definecolor{currentstroke}{rgb}{0.000000,0.000000,0.000000}%
\pgfsetstrokecolor{currentstroke}%
\pgfsetdash{}{0pt}%
\pgfsys@defobject{currentmarker}{\pgfqpoint{0.000000in}{0.000000in}}{\pgfqpoint{0.020833in}{0.000000in}}{%
\pgfpathmoveto{\pgfqpoint{0.000000in}{0.000000in}}%
\pgfpathlineto{\pgfqpoint{0.020833in}{0.000000in}}%
\pgfusepath{stroke,fill}%
}%
\begin{pgfscope}%
\pgfsys@transformshift{0.650000in}{1.217795in}%
\pgfsys@useobject{currentmarker}{}%
\end{pgfscope}%
\end{pgfscope}%
\begin{pgfscope}%
\pgfsetbuttcap%
\pgfsetroundjoin%
\definecolor{currentfill}{rgb}{0.000000,0.000000,0.000000}%
\pgfsetfillcolor{currentfill}%
\pgfsetlinewidth{0.401500pt}%
\definecolor{currentstroke}{rgb}{0.000000,0.000000,0.000000}%
\pgfsetstrokecolor{currentstroke}%
\pgfsetdash{}{0pt}%
\pgfsys@defobject{currentmarker}{\pgfqpoint{-0.020833in}{0.000000in}}{\pgfqpoint{-0.000000in}{0.000000in}}{%
\pgfpathmoveto{\pgfqpoint{-0.000000in}{0.000000in}}%
\pgfpathlineto{\pgfqpoint{-0.020833in}{0.000000in}}%
\pgfusepath{stroke,fill}%
}%
\begin{pgfscope}%
\pgfsys@transformshift{3.038110in}{1.217795in}%
\pgfsys@useobject{currentmarker}{}%
\end{pgfscope}%
\end{pgfscope}%
\begin{pgfscope}%
\pgfsetbuttcap%
\pgfsetroundjoin%
\definecolor{currentfill}{rgb}{0.000000,0.000000,0.000000}%
\pgfsetfillcolor{currentfill}%
\pgfsetlinewidth{0.401500pt}%
\definecolor{currentstroke}{rgb}{0.000000,0.000000,0.000000}%
\pgfsetstrokecolor{currentstroke}%
\pgfsetdash{}{0pt}%
\pgfsys@defobject{currentmarker}{\pgfqpoint{0.000000in}{0.000000in}}{\pgfqpoint{0.020833in}{0.000000in}}{%
\pgfpathmoveto{\pgfqpoint{0.000000in}{0.000000in}}%
\pgfpathlineto{\pgfqpoint{0.020833in}{0.000000in}}%
\pgfusepath{stroke,fill}%
}%
\begin{pgfscope}%
\pgfsys@transformshift{0.650000in}{1.350761in}%
\pgfsys@useobject{currentmarker}{}%
\end{pgfscope}%
\end{pgfscope}%
\begin{pgfscope}%
\pgfsetbuttcap%
\pgfsetroundjoin%
\definecolor{currentfill}{rgb}{0.000000,0.000000,0.000000}%
\pgfsetfillcolor{currentfill}%
\pgfsetlinewidth{0.401500pt}%
\definecolor{currentstroke}{rgb}{0.000000,0.000000,0.000000}%
\pgfsetstrokecolor{currentstroke}%
\pgfsetdash{}{0pt}%
\pgfsys@defobject{currentmarker}{\pgfqpoint{-0.020833in}{0.000000in}}{\pgfqpoint{-0.000000in}{0.000000in}}{%
\pgfpathmoveto{\pgfqpoint{-0.000000in}{0.000000in}}%
\pgfpathlineto{\pgfqpoint{-0.020833in}{0.000000in}}%
\pgfusepath{stroke,fill}%
}%
\begin{pgfscope}%
\pgfsys@transformshift{3.038110in}{1.350761in}%
\pgfsys@useobject{currentmarker}{}%
\end{pgfscope}%
\end{pgfscope}%
\begin{pgfscope}%
\pgfsetbuttcap%
\pgfsetroundjoin%
\definecolor{currentfill}{rgb}{0.000000,0.000000,0.000000}%
\pgfsetfillcolor{currentfill}%
\pgfsetlinewidth{0.401500pt}%
\definecolor{currentstroke}{rgb}{0.000000,0.000000,0.000000}%
\pgfsetstrokecolor{currentstroke}%
\pgfsetdash{}{0pt}%
\pgfsys@defobject{currentmarker}{\pgfqpoint{0.000000in}{0.000000in}}{\pgfqpoint{0.020833in}{0.000000in}}{%
\pgfpathmoveto{\pgfqpoint{0.000000in}{0.000000in}}%
\pgfpathlineto{\pgfqpoint{0.020833in}{0.000000in}}%
\pgfusepath{stroke,fill}%
}%
\begin{pgfscope}%
\pgfsys@transformshift{0.650000in}{1.483727in}%
\pgfsys@useobject{currentmarker}{}%
\end{pgfscope}%
\end{pgfscope}%
\begin{pgfscope}%
\pgfsetbuttcap%
\pgfsetroundjoin%
\definecolor{currentfill}{rgb}{0.000000,0.000000,0.000000}%
\pgfsetfillcolor{currentfill}%
\pgfsetlinewidth{0.401500pt}%
\definecolor{currentstroke}{rgb}{0.000000,0.000000,0.000000}%
\pgfsetstrokecolor{currentstroke}%
\pgfsetdash{}{0pt}%
\pgfsys@defobject{currentmarker}{\pgfqpoint{-0.020833in}{0.000000in}}{\pgfqpoint{-0.000000in}{0.000000in}}{%
\pgfpathmoveto{\pgfqpoint{-0.000000in}{0.000000in}}%
\pgfpathlineto{\pgfqpoint{-0.020833in}{0.000000in}}%
\pgfusepath{stroke,fill}%
}%
\begin{pgfscope}%
\pgfsys@transformshift{3.038110in}{1.483727in}%
\pgfsys@useobject{currentmarker}{}%
\end{pgfscope}%
\end{pgfscope}%
\begin{pgfscope}%
\pgfsetbuttcap%
\pgfsetroundjoin%
\definecolor{currentfill}{rgb}{0.000000,0.000000,0.000000}%
\pgfsetfillcolor{currentfill}%
\pgfsetlinewidth{0.401500pt}%
\definecolor{currentstroke}{rgb}{0.000000,0.000000,0.000000}%
\pgfsetstrokecolor{currentstroke}%
\pgfsetdash{}{0pt}%
\pgfsys@defobject{currentmarker}{\pgfqpoint{0.000000in}{0.000000in}}{\pgfqpoint{0.020833in}{0.000000in}}{%
\pgfpathmoveto{\pgfqpoint{0.000000in}{0.000000in}}%
\pgfpathlineto{\pgfqpoint{0.020833in}{0.000000in}}%
\pgfusepath{stroke,fill}%
}%
\begin{pgfscope}%
\pgfsys@transformshift{0.650000in}{1.616693in}%
\pgfsys@useobject{currentmarker}{}%
\end{pgfscope}%
\end{pgfscope}%
\begin{pgfscope}%
\pgfsetbuttcap%
\pgfsetroundjoin%
\definecolor{currentfill}{rgb}{0.000000,0.000000,0.000000}%
\pgfsetfillcolor{currentfill}%
\pgfsetlinewidth{0.401500pt}%
\definecolor{currentstroke}{rgb}{0.000000,0.000000,0.000000}%
\pgfsetstrokecolor{currentstroke}%
\pgfsetdash{}{0pt}%
\pgfsys@defobject{currentmarker}{\pgfqpoint{-0.020833in}{0.000000in}}{\pgfqpoint{-0.000000in}{0.000000in}}{%
\pgfpathmoveto{\pgfqpoint{-0.000000in}{0.000000in}}%
\pgfpathlineto{\pgfqpoint{-0.020833in}{0.000000in}}%
\pgfusepath{stroke,fill}%
}%
\begin{pgfscope}%
\pgfsys@transformshift{3.038110in}{1.616693in}%
\pgfsys@useobject{currentmarker}{}%
\end{pgfscope}%
\end{pgfscope}%
\begin{pgfscope}%
\pgfsetbuttcap%
\pgfsetroundjoin%
\definecolor{currentfill}{rgb}{0.000000,0.000000,0.000000}%
\pgfsetfillcolor{currentfill}%
\pgfsetlinewidth{0.401500pt}%
\definecolor{currentstroke}{rgb}{0.000000,0.000000,0.000000}%
\pgfsetstrokecolor{currentstroke}%
\pgfsetdash{}{0pt}%
\pgfsys@defobject{currentmarker}{\pgfqpoint{0.000000in}{0.000000in}}{\pgfqpoint{0.020833in}{0.000000in}}{%
\pgfpathmoveto{\pgfqpoint{0.000000in}{0.000000in}}%
\pgfpathlineto{\pgfqpoint{0.020833in}{0.000000in}}%
\pgfusepath{stroke,fill}%
}%
\begin{pgfscope}%
\pgfsys@transformshift{0.650000in}{1.882625in}%
\pgfsys@useobject{currentmarker}{}%
\end{pgfscope}%
\end{pgfscope}%
\begin{pgfscope}%
\pgfsetbuttcap%
\pgfsetroundjoin%
\definecolor{currentfill}{rgb}{0.000000,0.000000,0.000000}%
\pgfsetfillcolor{currentfill}%
\pgfsetlinewidth{0.401500pt}%
\definecolor{currentstroke}{rgb}{0.000000,0.000000,0.000000}%
\pgfsetstrokecolor{currentstroke}%
\pgfsetdash{}{0pt}%
\pgfsys@defobject{currentmarker}{\pgfqpoint{-0.020833in}{0.000000in}}{\pgfqpoint{-0.000000in}{0.000000in}}{%
\pgfpathmoveto{\pgfqpoint{-0.000000in}{0.000000in}}%
\pgfpathlineto{\pgfqpoint{-0.020833in}{0.000000in}}%
\pgfusepath{stroke,fill}%
}%
\begin{pgfscope}%
\pgfsys@transformshift{3.038110in}{1.882625in}%
\pgfsys@useobject{currentmarker}{}%
\end{pgfscope}%
\end{pgfscope}%
\begin{pgfscope}%
\pgfsetbuttcap%
\pgfsetroundjoin%
\definecolor{currentfill}{rgb}{0.000000,0.000000,0.000000}%
\pgfsetfillcolor{currentfill}%
\pgfsetlinewidth{0.401500pt}%
\definecolor{currentstroke}{rgb}{0.000000,0.000000,0.000000}%
\pgfsetstrokecolor{currentstroke}%
\pgfsetdash{}{0pt}%
\pgfsys@defobject{currentmarker}{\pgfqpoint{0.000000in}{0.000000in}}{\pgfqpoint{0.020833in}{0.000000in}}{%
\pgfpathmoveto{\pgfqpoint{0.000000in}{0.000000in}}%
\pgfpathlineto{\pgfqpoint{0.020833in}{0.000000in}}%
\pgfusepath{stroke,fill}%
}%
\begin{pgfscope}%
\pgfsys@transformshift{0.650000in}{2.015590in}%
\pgfsys@useobject{currentmarker}{}%
\end{pgfscope}%
\end{pgfscope}%
\begin{pgfscope}%
\pgfsetbuttcap%
\pgfsetroundjoin%
\definecolor{currentfill}{rgb}{0.000000,0.000000,0.000000}%
\pgfsetfillcolor{currentfill}%
\pgfsetlinewidth{0.401500pt}%
\definecolor{currentstroke}{rgb}{0.000000,0.000000,0.000000}%
\pgfsetstrokecolor{currentstroke}%
\pgfsetdash{}{0pt}%
\pgfsys@defobject{currentmarker}{\pgfqpoint{-0.020833in}{0.000000in}}{\pgfqpoint{-0.000000in}{0.000000in}}{%
\pgfpathmoveto{\pgfqpoint{-0.000000in}{0.000000in}}%
\pgfpathlineto{\pgfqpoint{-0.020833in}{0.000000in}}%
\pgfusepath{stroke,fill}%
}%
\begin{pgfscope}%
\pgfsys@transformshift{3.038110in}{2.015590in}%
\pgfsys@useobject{currentmarker}{}%
\end{pgfscope}%
\end{pgfscope}%
\begin{pgfscope}%
\pgfsetbuttcap%
\pgfsetroundjoin%
\definecolor{currentfill}{rgb}{0.000000,0.000000,0.000000}%
\pgfsetfillcolor{currentfill}%
\pgfsetlinewidth{0.401500pt}%
\definecolor{currentstroke}{rgb}{0.000000,0.000000,0.000000}%
\pgfsetstrokecolor{currentstroke}%
\pgfsetdash{}{0pt}%
\pgfsys@defobject{currentmarker}{\pgfqpoint{0.000000in}{0.000000in}}{\pgfqpoint{0.020833in}{0.000000in}}{%
\pgfpathmoveto{\pgfqpoint{0.000000in}{0.000000in}}%
\pgfpathlineto{\pgfqpoint{0.020833in}{0.000000in}}%
\pgfusepath{stroke,fill}%
}%
\begin{pgfscope}%
\pgfsys@transformshift{0.650000in}{2.148556in}%
\pgfsys@useobject{currentmarker}{}%
\end{pgfscope}%
\end{pgfscope}%
\begin{pgfscope}%
\pgfsetbuttcap%
\pgfsetroundjoin%
\definecolor{currentfill}{rgb}{0.000000,0.000000,0.000000}%
\pgfsetfillcolor{currentfill}%
\pgfsetlinewidth{0.401500pt}%
\definecolor{currentstroke}{rgb}{0.000000,0.000000,0.000000}%
\pgfsetstrokecolor{currentstroke}%
\pgfsetdash{}{0pt}%
\pgfsys@defobject{currentmarker}{\pgfqpoint{-0.020833in}{0.000000in}}{\pgfqpoint{-0.000000in}{0.000000in}}{%
\pgfpathmoveto{\pgfqpoint{-0.000000in}{0.000000in}}%
\pgfpathlineto{\pgfqpoint{-0.020833in}{0.000000in}}%
\pgfusepath{stroke,fill}%
}%
\begin{pgfscope}%
\pgfsys@transformshift{3.038110in}{2.148556in}%
\pgfsys@useobject{currentmarker}{}%
\end{pgfscope}%
\end{pgfscope}%
\begin{pgfscope}%
\pgfsetbuttcap%
\pgfsetroundjoin%
\definecolor{currentfill}{rgb}{0.000000,0.000000,0.000000}%
\pgfsetfillcolor{currentfill}%
\pgfsetlinewidth{0.401500pt}%
\definecolor{currentstroke}{rgb}{0.000000,0.000000,0.000000}%
\pgfsetstrokecolor{currentstroke}%
\pgfsetdash{}{0pt}%
\pgfsys@defobject{currentmarker}{\pgfqpoint{0.000000in}{0.000000in}}{\pgfqpoint{0.020833in}{0.000000in}}{%
\pgfpathmoveto{\pgfqpoint{0.000000in}{0.000000in}}%
\pgfpathlineto{\pgfqpoint{0.020833in}{0.000000in}}%
\pgfusepath{stroke,fill}%
}%
\begin{pgfscope}%
\pgfsys@transformshift{0.650000in}{2.281522in}%
\pgfsys@useobject{currentmarker}{}%
\end{pgfscope}%
\end{pgfscope}%
\begin{pgfscope}%
\pgfsetbuttcap%
\pgfsetroundjoin%
\definecolor{currentfill}{rgb}{0.000000,0.000000,0.000000}%
\pgfsetfillcolor{currentfill}%
\pgfsetlinewidth{0.401500pt}%
\definecolor{currentstroke}{rgb}{0.000000,0.000000,0.000000}%
\pgfsetstrokecolor{currentstroke}%
\pgfsetdash{}{0pt}%
\pgfsys@defobject{currentmarker}{\pgfqpoint{-0.020833in}{0.000000in}}{\pgfqpoint{-0.000000in}{0.000000in}}{%
\pgfpathmoveto{\pgfqpoint{-0.000000in}{0.000000in}}%
\pgfpathlineto{\pgfqpoint{-0.020833in}{0.000000in}}%
\pgfusepath{stroke,fill}%
}%
\begin{pgfscope}%
\pgfsys@transformshift{3.038110in}{2.281522in}%
\pgfsys@useobject{currentmarker}{}%
\end{pgfscope}%
\end{pgfscope}%
\begin{pgfscope}%
\definecolor{textcolor}{rgb}{0.000000,0.000000,0.000000}%
\pgfsetstrokecolor{textcolor}%
\pgfsetfillcolor{textcolor}%
\pgftext[x=0.184953in,y=1.417244in,,bottom,rotate=90.000000]{\color{textcolor}{\rmfamily\fontsize{12.000000}{14.400000}\selectfont\catcode`\^=\active\def^{\ifmmode\sp\else\^{}\fi}\catcode`\%=\active\def%{\%}$p(s)-\langle p_{c}\rangle$}}%
\end{pgfscope}%
\begin{pgfscope}%
\pgfpathrectangle{\pgfqpoint{0.650000in}{0.420000in}}{\pgfqpoint{2.388110in}{1.994488in}}%
\pgfusepath{clip}%
\pgfsetrectcap%
\pgfsetroundjoin%
\pgfsetlinewidth{1.003750pt}%
\definecolor{currentstroke}{rgb}{0.870588,0.466667,0.682353}%
\pgfsetstrokecolor{currentstroke}%
\pgfsetdash{}{0pt}%
\pgfpathmoveto{\pgfqpoint{0.650000in}{2.047441in}}%
\pgfpathlineto{\pgfqpoint{0.666319in}{2.017279in}}%
\pgfpathlineto{\pgfqpoint{0.669882in}{2.012404in}}%
\pgfpathlineto{\pgfqpoint{0.672733in}{2.010352in}}%
\pgfpathlineto{\pgfqpoint{0.675584in}{2.010042in}}%
\pgfpathlineto{\pgfqpoint{0.680573in}{2.010198in}}%
\pgfpathlineto{\pgfqpoint{0.683423in}{2.008782in}}%
\pgfpathlineto{\pgfqpoint{0.688412in}{2.004379in}}%
\pgfpathlineto{\pgfqpoint{0.692689in}{2.001331in}}%
\pgfpathlineto{\pgfqpoint{0.695539in}{2.000495in}}%
\pgfpathlineto{\pgfqpoint{0.698390in}{2.000721in}}%
\pgfpathlineto{\pgfqpoint{0.701954in}{2.002315in}}%
\pgfpathlineto{\pgfqpoint{0.706230in}{2.005619in}}%
\pgfpathlineto{\pgfqpoint{0.714782in}{2.014160in}}%
\pgfpathlineto{\pgfqpoint{0.723335in}{2.022122in}}%
\pgfpathlineto{\pgfqpoint{0.729036in}{2.025837in}}%
\pgfpathlineto{\pgfqpoint{0.734738in}{2.028072in}}%
\pgfpathlineto{\pgfqpoint{0.742577in}{2.029754in}}%
\pgfpathlineto{\pgfqpoint{0.751130in}{2.030545in}}%
\pgfpathlineto{\pgfqpoint{0.757544in}{2.030141in}}%
\pgfpathlineto{\pgfqpoint{0.763958in}{2.028515in}}%
\pgfpathlineto{\pgfqpoint{0.770373in}{2.025649in}}%
\pgfpathlineto{\pgfqpoint{0.778212in}{2.020717in}}%
\pgfpathlineto{\pgfqpoint{0.792466in}{2.010076in}}%
\pgfpathlineto{\pgfqpoint{0.811709in}{1.994462in}}%
\pgfpathlineto{\pgfqpoint{0.832378in}{1.977680in}}%
\pgfpathlineto{\pgfqpoint{0.848057in}{1.965915in}}%
\pgfpathlineto{\pgfqpoint{0.892957in}{1.935239in}}%
\pgfpathlineto{\pgfqpoint{0.922890in}{1.915433in}}%
\pgfpathlineto{\pgfqpoint{0.972067in}{1.884817in}}%
\pgfpathlineto{\pgfqpoint{0.989171in}{1.874403in}}%
\pgfpathlineto{\pgfqpoint{1.014829in}{1.859232in}}%
\pgfpathlineto{\pgfqpoint{1.044762in}{1.841167in}}%
\pgfpathlineto{\pgfqpoint{1.118883in}{1.798484in}}%
\pgfpathlineto{\pgfqpoint{1.135987in}{1.788471in}}%
\pgfpathlineto{\pgfqpoint{1.174473in}{1.767072in}}%
\pgfpathlineto{\pgfqpoint{1.232202in}{1.735055in}}%
\pgfpathlineto{\pgfqpoint{1.268550in}{1.715220in}}%
\pgfpathlineto{\pgfqpoint{1.285654in}{1.705649in}}%
\pgfpathlineto{\pgfqpoint{1.332693in}{1.680630in}}%
\pgfpathlineto{\pgfqpoint{1.389708in}{1.649661in}}%
\pgfpathlineto{\pgfqpoint{1.413940in}{1.637206in}}%
\pgfpathlineto{\pgfqpoint{1.431758in}{1.628510in}}%
\pgfpathlineto{\pgfqpoint{1.470956in}{1.607698in}}%
\pgfpathlineto{\pgfqpoint{1.524409in}{1.577663in}}%
\pgfpathlineto{\pgfqpoint{1.538663in}{1.568084in}}%
\pgfpathlineto{\pgfqpoint{1.550778in}{1.558744in}}%
\pgfpathlineto{\pgfqpoint{1.559331in}{1.550856in}}%
\pgfpathlineto{\pgfqpoint{1.568596in}{1.540810in}}%
\pgfpathlineto{\pgfqpoint{1.576436in}{1.530692in}}%
\pgfpathlineto{\pgfqpoint{1.582850in}{1.520602in}}%
\pgfpathlineto{\pgfqpoint{1.589264in}{1.508503in}}%
\pgfpathlineto{\pgfqpoint{1.594966in}{1.495588in}}%
\pgfpathlineto{\pgfqpoint{1.600667in}{1.479785in}}%
\pgfpathlineto{\pgfqpoint{1.605656in}{1.462744in}}%
\pgfpathlineto{\pgfqpoint{1.611358in}{1.439107in}}%
\pgfpathlineto{\pgfqpoint{1.623474in}{1.387319in}}%
\pgfpathlineto{\pgfqpoint{1.629175in}{1.368584in}}%
\pgfpathlineto{\pgfqpoint{1.634877in}{1.353371in}}%
\pgfpathlineto{\pgfqpoint{1.641291in}{1.339404in}}%
\pgfpathlineto{\pgfqpoint{1.647705in}{1.328037in}}%
\pgfpathlineto{\pgfqpoint{1.654833in}{1.317725in}}%
\pgfpathlineto{\pgfqpoint{1.663385in}{1.307422in}}%
\pgfpathlineto{\pgfqpoint{1.673363in}{1.297207in}}%
\pgfpathlineto{\pgfqpoint{1.681915in}{1.289920in}}%
\pgfpathlineto{\pgfqpoint{1.694744in}{1.280832in}}%
\pgfpathlineto{\pgfqpoint{1.709711in}{1.271393in}}%
\pgfpathlineto{\pgfqpoint{1.734655in}{1.257399in}}%
\pgfpathlineto{\pgfqpoint{1.797373in}{1.224061in}}%
\pgfpathlineto{\pgfqpoint{1.810202in}{1.217707in}}%
\pgfpathlineto{\pgfqpoint{1.857953in}{1.194996in}}%
\pgfpathlineto{\pgfqpoint{1.907129in}{1.169097in}}%
\pgfpathlineto{\pgfqpoint{1.941339in}{1.152264in}}%
\pgfpathlineto{\pgfqpoint{1.966996in}{1.140835in}}%
\pgfpathlineto{\pgfqpoint{1.996930in}{1.126403in}}%
\pgfpathlineto{\pgfqpoint{2.039692in}{1.103881in}}%
\pgfpathlineto{\pgfqpoint{2.083167in}{1.083255in}}%
\pgfpathlineto{\pgfqpoint{2.121653in}{1.065880in}}%
\pgfpathlineto{\pgfqpoint{2.306242in}{0.978779in}}%
\pgfpathlineto{\pgfqpoint{2.318358in}{0.972954in}}%
\pgfpathlineto{\pgfqpoint{2.331186in}{0.966947in}}%
\pgfpathlineto{\pgfqpoint{2.345441in}{0.960310in}}%
\pgfpathlineto{\pgfqpoint{2.375374in}{0.946252in}}%
\pgfpathlineto{\pgfqpoint{2.392479in}{0.938658in}}%
\pgfpathlineto{\pgfqpoint{2.433816in}{0.920425in}}%
\pgfpathlineto{\pgfqpoint{2.443081in}{0.917358in}}%
\pgfpathlineto{\pgfqpoint{2.455197in}{0.913537in}}%
\pgfpathlineto{\pgfqpoint{2.481567in}{0.904046in}}%
\pgfpathlineto{\pgfqpoint{2.492258in}{0.901416in}}%
\pgfpathlineto{\pgfqpoint{2.500810in}{0.900469in}}%
\pgfpathlineto{\pgfqpoint{2.510788in}{0.900636in}}%
\pgfpathlineto{\pgfqpoint{2.518628in}{0.901645in}}%
\pgfpathlineto{\pgfqpoint{2.525042in}{0.903732in}}%
\pgfpathlineto{\pgfqpoint{2.530744in}{0.906800in}}%
\pgfpathlineto{\pgfqpoint{2.536445in}{0.911094in}}%
\pgfpathlineto{\pgfqpoint{2.542147in}{0.916836in}}%
\pgfpathlineto{\pgfqpoint{2.547136in}{0.923464in}}%
\pgfpathlineto{\pgfqpoint{2.552125in}{0.932089in}}%
\pgfpathlineto{\pgfqpoint{2.557114in}{0.943356in}}%
\pgfpathlineto{\pgfqpoint{2.563528in}{0.961526in}}%
\pgfpathlineto{\pgfqpoint{2.571368in}{0.983241in}}%
\pgfpathlineto{\pgfqpoint{2.576357in}{0.994055in}}%
\pgfpathlineto{\pgfqpoint{2.582058in}{1.003670in}}%
\pgfpathlineto{\pgfqpoint{2.587047in}{1.010093in}}%
\pgfpathlineto{\pgfqpoint{2.592036in}{1.014897in}}%
\pgfpathlineto{\pgfqpoint{2.597738in}{1.018866in}}%
\pgfpathlineto{\pgfqpoint{2.603440in}{1.021540in}}%
\pgfpathlineto{\pgfqpoint{2.609854in}{1.023318in}}%
\pgfpathlineto{\pgfqpoint{2.620545in}{1.025003in}}%
\pgfpathlineto{\pgfqpoint{2.626960in}{1.024926in}}%
\pgfpathlineto{\pgfqpoint{2.634799in}{1.023572in}}%
\pgfpathlineto{\pgfqpoint{2.660457in}{1.017041in}}%
\pgfpathlineto{\pgfqpoint{2.674711in}{1.012282in}}%
\pgfpathlineto{\pgfqpoint{2.698944in}{1.003549in}}%
\pgfpathlineto{\pgfqpoint{2.714624in}{0.998046in}}%
\pgfpathlineto{\pgfqpoint{2.733867in}{0.991366in}}%
\pgfpathlineto{\pgfqpoint{2.740994in}{0.988706in}}%
\pgfpathlineto{\pgfqpoint{2.760951in}{0.980138in}}%
\pgfpathlineto{\pgfqpoint{2.768791in}{0.977288in}}%
\pgfpathlineto{\pgfqpoint{2.783045in}{0.973196in}}%
\pgfpathlineto{\pgfqpoint{2.806565in}{0.962581in}}%
\pgfpathlineto{\pgfqpoint{2.818681in}{0.958200in}}%
\pgfpathlineto{\pgfqpoint{2.832935in}{0.952293in}}%
\pgfpathlineto{\pgfqpoint{2.869994in}{0.939137in}}%
\pgfpathlineto{\pgfqpoint{2.884960in}{0.933151in}}%
\pgfpathlineto{\pgfqpoint{2.897075in}{0.928280in}}%
\pgfpathlineto{\pgfqpoint{3.000412in}{0.890537in}}%
\pgfpathlineto{\pgfqpoint{3.009676in}{0.888124in}}%
\pgfpathlineto{\pgfqpoint{3.022504in}{0.884732in}}%
\pgfpathlineto{\pgfqpoint{3.032379in}{0.881055in}}%
\pgfpathlineto{\pgfqpoint{3.038110in}{0.878274in}}%
\pgfpathlineto{\pgfqpoint{3.038110in}{0.878274in}}%
\pgfusepath{stroke}%
\end{pgfscope}%
\begin{pgfscope}%
\pgfpathrectangle{\pgfqpoint{0.650000in}{0.420000in}}{\pgfqpoint{2.388110in}{1.994488in}}%
\pgfusepath{clip}%
\pgfsetrectcap%
\pgfsetroundjoin%
\pgfsetlinewidth{1.003750pt}%
\definecolor{currentstroke}{rgb}{0.498039,0.737255,0.254902}%
\pgfsetstrokecolor{currentstroke}%
\pgfsetdash{}{0pt}%
\pgfpathmoveto{\pgfqpoint{0.650000in}{1.807482in}}%
\pgfpathlineto{\pgfqpoint{0.666190in}{1.777335in}}%
\pgfpathlineto{\pgfqpoint{0.669726in}{1.772457in}}%
\pgfpathlineto{\pgfqpoint{0.672554in}{1.770395in}}%
\pgfpathlineto{\pgfqpoint{0.675382in}{1.770069in}}%
\pgfpathlineto{\pgfqpoint{0.680332in}{1.770205in}}%
\pgfpathlineto{\pgfqpoint{0.683160in}{1.768794in}}%
\pgfpathlineto{\pgfqpoint{0.688817in}{1.763831in}}%
\pgfpathlineto{\pgfqpoint{0.692353in}{1.761476in}}%
\pgfpathlineto{\pgfqpoint{0.695181in}{1.760699in}}%
\pgfpathlineto{\pgfqpoint{0.698009in}{1.760959in}}%
\pgfpathlineto{\pgfqpoint{0.701545in}{1.762532in}}%
\pgfpathlineto{\pgfqpoint{0.705787in}{1.765721in}}%
\pgfpathlineto{\pgfqpoint{0.713565in}{1.773296in}}%
\pgfpathlineto{\pgfqpoint{0.722051in}{1.781106in}}%
\pgfpathlineto{\pgfqpoint{0.728414in}{1.785374in}}%
\pgfpathlineto{\pgfqpoint{0.734778in}{1.788185in}}%
\pgfpathlineto{\pgfqpoint{0.742556in}{1.790329in}}%
\pgfpathlineto{\pgfqpoint{0.749627in}{1.791169in}}%
\pgfpathlineto{\pgfqpoint{0.755991in}{1.790738in}}%
\pgfpathlineto{\pgfqpoint{0.762355in}{1.789100in}}%
\pgfpathlineto{\pgfqpoint{0.768719in}{1.786194in}}%
\pgfpathlineto{\pgfqpoint{0.776497in}{1.781282in}}%
\pgfpathlineto{\pgfqpoint{0.787103in}{1.773386in}}%
\pgfpathlineto{\pgfqpoint{0.830943in}{1.737591in}}%
\pgfpathlineto{\pgfqpoint{0.865590in}{1.711729in}}%
\pgfpathlineto{\pgfqpoint{0.924986in}{1.673501in}}%
\pgfpathlineto{\pgfqpoint{1.012666in}{1.619090in}}%
\pgfpathlineto{\pgfqpoint{1.074183in}{1.582798in}}%
\pgfpathlineto{\pgfqpoint{1.094688in}{1.571045in}}%
\pgfpathlineto{\pgfqpoint{1.127214in}{1.552370in}}%
\pgfpathlineto{\pgfqpoint{1.270047in}{1.471359in}}%
\pgfpathlineto{\pgfqpoint{1.314594in}{1.446854in}}%
\pgfpathlineto{\pgfqpoint{1.335100in}{1.434455in}}%
\pgfpathlineto{\pgfqpoint{1.361969in}{1.415630in}}%
\pgfpathlineto{\pgfqpoint{1.374697in}{1.405543in}}%
\pgfpathlineto{\pgfqpoint{1.386010in}{1.395148in}}%
\pgfpathlineto{\pgfqpoint{1.395910in}{1.384545in}}%
\pgfpathlineto{\pgfqpoint{1.405809in}{1.372208in}}%
\pgfpathlineto{\pgfqpoint{1.417122in}{1.356153in}}%
\pgfpathlineto{\pgfqpoint{1.429850in}{1.336174in}}%
\pgfpathlineto{\pgfqpoint{1.443285in}{1.312865in}}%
\pgfpathlineto{\pgfqpoint{1.487832in}{1.233912in}}%
\pgfpathlineto{\pgfqpoint{1.498438in}{1.217846in}}%
\pgfpathlineto{\pgfqpoint{1.507630in}{1.205869in}}%
\pgfpathlineto{\pgfqpoint{1.518237in}{1.193948in}}%
\pgfpathlineto{\pgfqpoint{1.529550in}{1.182745in}}%
\pgfpathlineto{\pgfqpoint{1.539449in}{1.174515in}}%
\pgfpathlineto{\pgfqpoint{1.551470in}{1.166070in}}%
\pgfpathlineto{\pgfqpoint{1.567026in}{1.155014in}}%
\pgfpathlineto{\pgfqpoint{1.574097in}{1.148643in}}%
\pgfpathlineto{\pgfqpoint{1.580461in}{1.141454in}}%
\pgfpathlineto{\pgfqpoint{1.586117in}{1.133431in}}%
\pgfpathlineto{\pgfqpoint{1.591067in}{1.124612in}}%
\pgfpathlineto{\pgfqpoint{1.596017in}{1.113300in}}%
\pgfpathlineto{\pgfqpoint{1.600966in}{1.098635in}}%
\pgfpathlineto{\pgfqpoint{1.609452in}{1.068174in}}%
\pgfpathlineto{\pgfqpoint{1.615108in}{1.050459in}}%
\pgfpathlineto{\pgfqpoint{1.619351in}{1.040355in}}%
\pgfpathlineto{\pgfqpoint{1.623593in}{1.032707in}}%
\pgfpathlineto{\pgfqpoint{1.627836in}{1.027047in}}%
\pgfpathlineto{\pgfqpoint{1.632078in}{1.022978in}}%
\pgfpathlineto{\pgfqpoint{1.636321in}{1.020237in}}%
\pgfpathlineto{\pgfqpoint{1.640563in}{1.018694in}}%
\pgfpathlineto{\pgfqpoint{1.645513in}{1.018211in}}%
\pgfpathlineto{\pgfqpoint{1.651170in}{1.018942in}}%
\pgfpathlineto{\pgfqpoint{1.657534in}{1.020943in}}%
\pgfpathlineto{\pgfqpoint{1.664605in}{1.024437in}}%
\pgfpathlineto{\pgfqpoint{1.675211in}{1.031186in}}%
\pgfpathlineto{\pgfqpoint{1.704909in}{1.050884in}}%
\pgfpathlineto{\pgfqpoint{1.716930in}{1.057263in}}%
\pgfpathlineto{\pgfqpoint{1.726122in}{1.061017in}}%
\pgfpathlineto{\pgfqpoint{1.734607in}{1.063220in}}%
\pgfpathlineto{\pgfqpoint{1.743092in}{1.064263in}}%
\pgfpathlineto{\pgfqpoint{1.751578in}{1.064185in}}%
\pgfpathlineto{\pgfqpoint{1.759356in}{1.062899in}}%
\pgfpathlineto{\pgfqpoint{1.770669in}{1.059663in}}%
\pgfpathlineto{\pgfqpoint{1.818752in}{1.045210in}}%
\pgfpathlineto{\pgfqpoint{1.826530in}{1.043972in}}%
\pgfpathlineto{\pgfqpoint{1.835015in}{1.043866in}}%
\pgfpathlineto{\pgfqpoint{1.843500in}{1.044888in}}%
\pgfpathlineto{\pgfqpoint{1.851985in}{1.047016in}}%
\pgfpathlineto{\pgfqpoint{1.862592in}{1.050875in}}%
\pgfpathlineto{\pgfqpoint{1.874612in}{1.056259in}}%
\pgfpathlineto{\pgfqpoint{1.909260in}{1.074677in}}%
\pgfpathlineto{\pgfqpoint{1.921987in}{1.080981in}}%
\pgfpathlineto{\pgfqpoint{1.934715in}{1.086045in}}%
\pgfpathlineto{\pgfqpoint{1.943200in}{1.088287in}}%
\pgfpathlineto{\pgfqpoint{1.956635in}{1.090534in}}%
\pgfpathlineto{\pgfqpoint{1.967242in}{1.091182in}}%
\pgfpathlineto{\pgfqpoint{1.977848in}{1.090763in}}%
\pgfpathlineto{\pgfqpoint{1.989162in}{1.089387in}}%
\pgfpathlineto{\pgfqpoint{2.001889in}{1.086345in}}%
\pgfpathlineto{\pgfqpoint{2.045729in}{1.075294in}}%
\pgfpathlineto{\pgfqpoint{2.058457in}{1.073169in}}%
\pgfpathlineto{\pgfqpoint{2.071185in}{1.071749in}}%
\pgfpathlineto{\pgfqpoint{2.086741in}{1.071786in}}%
\pgfpathlineto{\pgfqpoint{2.100176in}{1.072852in}}%
\pgfpathlineto{\pgfqpoint{2.110075in}{1.074496in}}%
\pgfpathlineto{\pgfqpoint{2.123510in}{1.077955in}}%
\pgfpathlineto{\pgfqpoint{2.159572in}{1.089814in}}%
\pgfpathlineto{\pgfqpoint{2.175835in}{1.095245in}}%
\pgfpathlineto{\pgfqpoint{2.214724in}{1.106912in}}%
\pgfpathlineto{\pgfqpoint{2.227452in}{1.109342in}}%
\pgfpathlineto{\pgfqpoint{2.240887in}{1.110790in}}%
\pgfpathlineto{\pgfqpoint{2.250079in}{1.110807in}}%
\pgfpathlineto{\pgfqpoint{2.264928in}{1.109271in}}%
\pgfpathlineto{\pgfqpoint{2.284726in}{1.106161in}}%
\pgfpathlineto{\pgfqpoint{2.317960in}{1.101745in}}%
\pgfpathlineto{\pgfqpoint{2.358972in}{1.099292in}}%
\pgfpathlineto{\pgfqpoint{2.389377in}{1.099567in}}%
\pgfpathlineto{\pgfqpoint{2.402812in}{1.101299in}}%
\pgfpathlineto{\pgfqpoint{2.417661in}{1.104045in}}%
\pgfpathlineto{\pgfqpoint{2.430389in}{1.107157in}}%
\pgfpathlineto{\pgfqpoint{2.438166in}{1.110397in}}%
\pgfpathlineto{\pgfqpoint{2.452309in}{1.117689in}}%
\pgfpathlineto{\pgfqpoint{2.460087in}{1.123079in}}%
\pgfpathlineto{\pgfqpoint{2.468572in}{1.130383in}}%
\pgfpathlineto{\pgfqpoint{2.477057in}{1.139260in}}%
\pgfpathlineto{\pgfqpoint{2.484835in}{1.149073in}}%
\pgfpathlineto{\pgfqpoint{2.492613in}{1.160875in}}%
\pgfpathlineto{\pgfqpoint{2.501805in}{1.176778in}}%
\pgfpathlineto{\pgfqpoint{2.508877in}{1.191433in}}%
\pgfpathlineto{\pgfqpoint{2.516655in}{1.210171in}}%
\pgfpathlineto{\pgfqpoint{2.524433in}{1.231904in}}%
\pgfpathlineto{\pgfqpoint{2.531504in}{1.254858in}}%
\pgfpathlineto{\pgfqpoint{2.538575in}{1.281716in}}%
\pgfpathlineto{\pgfqpoint{2.545646in}{1.313316in}}%
\pgfpathlineto{\pgfqpoint{2.559080in}{1.374764in}}%
\pgfpathlineto{\pgfqpoint{2.566858in}{1.403791in}}%
\pgfpathlineto{\pgfqpoint{2.574636in}{1.428755in}}%
\pgfpathlineto{\pgfqpoint{2.582415in}{1.449969in}}%
\pgfpathlineto{\pgfqpoint{2.590193in}{1.468184in}}%
\pgfpathlineto{\pgfqpoint{2.598678in}{1.485279in}}%
\pgfpathlineto{\pgfqpoint{2.606456in}{1.498644in}}%
\pgfpathlineto{\pgfqpoint{2.614942in}{1.510856in}}%
\pgfpathlineto{\pgfqpoint{2.623427in}{1.521057in}}%
\pgfpathlineto{\pgfqpoint{2.637569in}{1.536175in}}%
\pgfpathlineto{\pgfqpoint{2.646762in}{1.544294in}}%
\pgfpathlineto{\pgfqpoint{2.655247in}{1.550332in}}%
\pgfpathlineto{\pgfqpoint{2.666561in}{1.556959in}}%
\pgfpathlineto{\pgfqpoint{2.683532in}{1.564985in}}%
\pgfpathlineto{\pgfqpoint{2.692017in}{1.567736in}}%
\pgfpathlineto{\pgfqpoint{2.699795in}{1.568984in}}%
\pgfpathlineto{\pgfqpoint{2.726666in}{1.571309in}}%
\pgfpathlineto{\pgfqpoint{2.766264in}{1.570593in}}%
\pgfpathlineto{\pgfqpoint{2.776870in}{1.568603in}}%
\pgfpathlineto{\pgfqpoint{2.805155in}{1.562252in}}%
\pgfpathlineto{\pgfqpoint{2.824954in}{1.555910in}}%
\pgfpathlineto{\pgfqpoint{2.841216in}{1.550663in}}%
\pgfpathlineto{\pgfqpoint{2.861721in}{1.543443in}}%
\pgfpathlineto{\pgfqpoint{2.912629in}{1.523490in}}%
\pgfpathlineto{\pgfqpoint{2.940204in}{1.511496in}}%
\pgfpathlineto{\pgfqpoint{2.958588in}{1.504556in}}%
\pgfpathlineto{\pgfqpoint{2.988284in}{1.495181in}}%
\pgfpathlineto{\pgfqpoint{3.005253in}{1.490010in}}%
\pgfpathlineto{\pgfqpoint{3.012962in}{1.486345in}}%
\pgfpathlineto{\pgfqpoint{3.019322in}{1.482175in}}%
\pgfpathlineto{\pgfqpoint{3.019322in}{1.482175in}}%
\pgfusepath{stroke}%
\end{pgfscope}%
\begin{pgfscope}%
\pgfpathrectangle{\pgfqpoint{0.650000in}{0.420000in}}{\pgfqpoint{2.388110in}{1.994488in}}%
\pgfusepath{clip}%
\pgfsetrectcap%
\pgfsetroundjoin%
\pgfsetlinewidth{0.803000pt}%
\definecolor{currentstroke}{rgb}{0.500000,0.500000,0.500000}%
\pgfsetstrokecolor{currentstroke}%
\pgfsetdash{}{0pt}%
\pgfpathmoveto{\pgfqpoint{0.650000in}{1.084829in}}%
\pgfpathlineto{\pgfqpoint{3.038110in}{1.084829in}}%
\pgfusepath{stroke}%
\end{pgfscope}%
\begin{pgfscope}%
\pgfsetrectcap%
\pgfsetmiterjoin%
\pgfsetlinewidth{0.803000pt}%
\definecolor{currentstroke}{rgb}{0.000000,0.000000,0.000000}%
\pgfsetstrokecolor{currentstroke}%
\pgfsetdash{}{0pt}%
\pgfpathmoveto{\pgfqpoint{0.650000in}{0.420000in}}%
\pgfpathlineto{\pgfqpoint{0.650000in}{2.414488in}}%
\pgfusepath{stroke}%
\end{pgfscope}%
\begin{pgfscope}%
\pgfsetrectcap%
\pgfsetmiterjoin%
\pgfsetlinewidth{0.803000pt}%
\definecolor{currentstroke}{rgb}{0.000000,0.000000,0.000000}%
\pgfsetstrokecolor{currentstroke}%
\pgfsetdash{}{0pt}%
\pgfpathmoveto{\pgfqpoint{3.038110in}{0.420000in}}%
\pgfpathlineto{\pgfqpoint{3.038110in}{2.414488in}}%
\pgfusepath{stroke}%
\end{pgfscope}%
\begin{pgfscope}%
\pgfsetrectcap%
\pgfsetmiterjoin%
\pgfsetlinewidth{0.803000pt}%
\definecolor{currentstroke}{rgb}{0.000000,0.000000,0.000000}%
\pgfsetstrokecolor{currentstroke}%
\pgfsetdash{}{0pt}%
\pgfpathmoveto{\pgfqpoint{0.650000in}{0.420000in}}%
\pgfpathlineto{\pgfqpoint{3.038110in}{0.420000in}}%
\pgfusepath{stroke}%
\end{pgfscope}%
\begin{pgfscope}%
\pgfsetrectcap%
\pgfsetmiterjoin%
\pgfsetlinewidth{0.803000pt}%
\definecolor{currentstroke}{rgb}{0.000000,0.000000,0.000000}%
\pgfsetstrokecolor{currentstroke}%
\pgfsetdash{}{0pt}%
\pgfpathmoveto{\pgfqpoint{0.650000in}{2.414488in}}%
\pgfpathlineto{\pgfqpoint{3.038110in}{2.414488in}}%
\pgfusepath{stroke}%
\end{pgfscope}%
\begin{pgfscope}%
\pgfsetrectcap%
\pgfsetroundjoin%
\pgfsetlinewidth{1.003750pt}%
\definecolor{currentstroke}{rgb}{0.870588,0.466667,0.682353}%
\pgfsetstrokecolor{currentstroke}%
\pgfsetdash{}{0pt}%
\pgfpathmoveto{\pgfqpoint{2.127879in}{2.275599in}}%
\pgfpathlineto{\pgfqpoint{2.238990in}{2.275599in}}%
\pgfpathlineto{\pgfqpoint{2.350101in}{2.275599in}}%
\pgfusepath{stroke}%
\end{pgfscope}%
\begin{pgfscope}%
\definecolor{textcolor}{rgb}{0.000000,0.000000,0.000000}%
\pgfsetstrokecolor{textcolor}%
\pgfsetfillcolor{textcolor}%
\pgftext[x=2.438990in,y=2.236710in,left,base]{\color{textcolor}{\rmfamily\fontsize{8.000000}{9.600000}\selectfont\catcode`\^=\active\def^{\ifmmode\sp\else\^{}\fi}\catcode`\%=\active\def%{\%}base}}%
\end{pgfscope}%
\begin{pgfscope}%
\pgfsetrectcap%
\pgfsetroundjoin%
\pgfsetlinewidth{1.003750pt}%
\definecolor{currentstroke}{rgb}{0.498039,0.737255,0.254902}%
\pgfsetstrokecolor{currentstroke}%
\pgfsetdash{}{0pt}%
\pgfpathmoveto{\pgfqpoint{2.127879in}{2.120661in}}%
\pgfpathlineto{\pgfqpoint{2.238990in}{2.120661in}}%
\pgfpathlineto{\pgfqpoint{2.350101in}{2.120661in}}%
\pgfusepath{stroke}%
\end{pgfscope}%
\begin{pgfscope}%
\definecolor{textcolor}{rgb}{0.000000,0.000000,0.000000}%
\pgfsetstrokecolor{textcolor}%
\pgfsetfillcolor{textcolor}%
\pgftext[x=2.438990in,y=2.081772in,left,base]{\color{textcolor}{\rmfamily\fontsize{8.000000}{9.600000}\selectfont\catcode`\^=\active\def^{\ifmmode\sp\else\^{}\fi}\catcode`\%=\active\def%{\%}optimised}}%
\end{pgfscope}%
\end{pgfpicture}%
\makeatother%
\endgroup%

    \label{fig:R100deg30Re1200p}
\end{subfigure}
\\[-20pt]
\begin{subfigure}[t]{0.495\textwidth}
    \centering
    \subcaption{}
    %% Creator: Matplotlib, PGF backend
%%
%% To include the figure in your LaTeX document, write
%%   \input{<filename>.pgf}
%%
%% Make sure the required packages are loaded in your preamble
%%   \usepackage{pgf}
%%
%% Also ensure that all the required font packages are loaded; for instance,
%% the lmodern package is sometimes necessary when using math font.
%%   \usepackage{lmodern}
%%
%% Figures using additional raster images can only be included by \input if
%% they are in the same directory as the main LaTeX file. For loading figures
%% from other directories you can use the `import` package
%%   \usepackage{import}
%%
%% and then include the figures with
%%   \import{<path to file>}{<filename>.pgf}
%%
%% Matplotlib used the following preamble
%%   \def\mathdefault#1{#1}
%%   \everymath=\expandafter{\the\everymath\displaystyle}
%%   \IfFileExists{scrextend.sty}{
%%     \usepackage[fontsize=11.000000pt]{scrextend}
%%   }{
%%     \renewcommand{\normalsize}{\fontsize{11.000000}{13.200000}\selectfont}
%%     \normalsize
%%   }
%%   
%%   \makeatletter\@ifpackageloaded{underscore}{}{\usepackage[strings]{underscore}}\makeatother
%%
\begingroup%
\makeatletter%
\begin{pgfpicture}%
\pgfpathrectangle{\pgfpointorigin}{\pgfqpoint{3.118110in}{2.494488in}}%
\pgfusepath{use as bounding box, clip}%
\begin{pgfscope}%
\pgfsetbuttcap%
\pgfsetmiterjoin%
\definecolor{currentfill}{rgb}{1.000000,1.000000,1.000000}%
\pgfsetfillcolor{currentfill}%
\pgfsetlinewidth{0.000000pt}%
\definecolor{currentstroke}{rgb}{1.000000,1.000000,1.000000}%
\pgfsetstrokecolor{currentstroke}%
\pgfsetdash{}{0pt}%
\pgfpathmoveto{\pgfqpoint{0.000000in}{0.000000in}}%
\pgfpathlineto{\pgfqpoint{3.118110in}{0.000000in}}%
\pgfpathlineto{\pgfqpoint{3.118110in}{2.494488in}}%
\pgfpathlineto{\pgfqpoint{0.000000in}{2.494488in}}%
\pgfpathlineto{\pgfqpoint{0.000000in}{0.000000in}}%
\pgfpathclose%
\pgfusepath{fill}%
\end{pgfscope}%
\begin{pgfscope}%
\pgfsetbuttcap%
\pgfsetmiterjoin%
\definecolor{currentfill}{rgb}{1.000000,1.000000,1.000000}%
\pgfsetfillcolor{currentfill}%
\pgfsetlinewidth{0.000000pt}%
\definecolor{currentstroke}{rgb}{0.000000,0.000000,0.000000}%
\pgfsetstrokecolor{currentstroke}%
\pgfsetstrokeopacity{0.000000}%
\pgfsetdash{}{0pt}%
\pgfpathmoveto{\pgfqpoint{0.650000in}{0.420000in}}%
\pgfpathlineto{\pgfqpoint{3.038110in}{0.420000in}}%
\pgfpathlineto{\pgfqpoint{3.038110in}{2.414488in}}%
\pgfpathlineto{\pgfqpoint{0.650000in}{2.414488in}}%
\pgfpathlineto{\pgfqpoint{0.650000in}{0.420000in}}%
\pgfpathclose%
\pgfusepath{fill}%
\end{pgfscope}%
\begin{pgfscope}%
\pgfsetbuttcap%
\pgfsetroundjoin%
\definecolor{currentfill}{rgb}{0.000000,0.000000,0.000000}%
\pgfsetfillcolor{currentfill}%
\pgfsetlinewidth{0.501875pt}%
\definecolor{currentstroke}{rgb}{0.000000,0.000000,0.000000}%
\pgfsetstrokecolor{currentstroke}%
\pgfsetdash{}{0pt}%
\pgfsys@defobject{currentmarker}{\pgfqpoint{0.000000in}{0.000000in}}{\pgfqpoint{0.000000in}{0.041667in}}{%
\pgfpathmoveto{\pgfqpoint{0.000000in}{0.000000in}}%
\pgfpathlineto{\pgfqpoint{0.000000in}{0.041667in}}%
\pgfusepath{stroke,fill}%
}%
\begin{pgfscope}%
\pgfsys@transformshift{0.650000in}{0.420000in}%
\pgfsys@useobject{currentmarker}{}%
\end{pgfscope}%
\end{pgfscope}%
\begin{pgfscope}%
\pgfsetbuttcap%
\pgfsetroundjoin%
\definecolor{currentfill}{rgb}{0.000000,0.000000,0.000000}%
\pgfsetfillcolor{currentfill}%
\pgfsetlinewidth{0.501875pt}%
\definecolor{currentstroke}{rgb}{0.000000,0.000000,0.000000}%
\pgfsetstrokecolor{currentstroke}%
\pgfsetdash{}{0pt}%
\pgfsys@defobject{currentmarker}{\pgfqpoint{0.000000in}{-0.041667in}}{\pgfqpoint{0.000000in}{0.000000in}}{%
\pgfpathmoveto{\pgfqpoint{0.000000in}{0.000000in}}%
\pgfpathlineto{\pgfqpoint{0.000000in}{-0.041667in}}%
\pgfusepath{stroke,fill}%
}%
\begin{pgfscope}%
\pgfsys@transformshift{0.650000in}{2.414488in}%
\pgfsys@useobject{currentmarker}{}%
\end{pgfscope}%
\end{pgfscope}%
\begin{pgfscope}%
\definecolor{textcolor}{rgb}{0.000000,0.000000,0.000000}%
\pgfsetstrokecolor{textcolor}%
\pgfsetfillcolor{textcolor}%
\pgftext[x=0.650000in,y=0.371389in,,top]{\color{textcolor}{\rmfamily\fontsize{11.000000}{13.200000}\selectfont\catcode`\^=\active\def^{\ifmmode\sp\else\^{}\fi}\catcode`\%=\active\def%{\%}0}}%
\end{pgfscope}%
\begin{pgfscope}%
\pgfsetbuttcap%
\pgfsetroundjoin%
\definecolor{currentfill}{rgb}{0.000000,0.000000,0.000000}%
\pgfsetfillcolor{currentfill}%
\pgfsetlinewidth{0.501875pt}%
\definecolor{currentstroke}{rgb}{0.000000,0.000000,0.000000}%
\pgfsetstrokecolor{currentstroke}%
\pgfsetdash{}{0pt}%
\pgfsys@defobject{currentmarker}{\pgfqpoint{0.000000in}{0.000000in}}{\pgfqpoint{0.000000in}{0.041667in}}{%
\pgfpathmoveto{\pgfqpoint{0.000000in}{0.000000in}}%
\pgfpathlineto{\pgfqpoint{0.000000in}{0.041667in}}%
\pgfusepath{stroke,fill}%
}%
\begin{pgfscope}%
\pgfsys@transformshift{1.492631in}{0.420000in}%
\pgfsys@useobject{currentmarker}{}%
\end{pgfscope}%
\end{pgfscope}%
\begin{pgfscope}%
\pgfsetbuttcap%
\pgfsetroundjoin%
\definecolor{currentfill}{rgb}{0.000000,0.000000,0.000000}%
\pgfsetfillcolor{currentfill}%
\pgfsetlinewidth{0.501875pt}%
\definecolor{currentstroke}{rgb}{0.000000,0.000000,0.000000}%
\pgfsetstrokecolor{currentstroke}%
\pgfsetdash{}{0pt}%
\pgfsys@defobject{currentmarker}{\pgfqpoint{0.000000in}{-0.041667in}}{\pgfqpoint{0.000000in}{0.000000in}}{%
\pgfpathmoveto{\pgfqpoint{0.000000in}{0.000000in}}%
\pgfpathlineto{\pgfqpoint{0.000000in}{-0.041667in}}%
\pgfusepath{stroke,fill}%
}%
\begin{pgfscope}%
\pgfsys@transformshift{1.492631in}{2.414488in}%
\pgfsys@useobject{currentmarker}{}%
\end{pgfscope}%
\end{pgfscope}%
\begin{pgfscope}%
\definecolor{textcolor}{rgb}{0.000000,0.000000,0.000000}%
\pgfsetstrokecolor{textcolor}%
\pgfsetfillcolor{textcolor}%
\pgftext[x=1.492631in,y=0.371389in,,top]{\color{textcolor}{\rmfamily\fontsize{11.000000}{13.200000}\selectfont\catcode`\^=\active\def^{\ifmmode\sp\else\^{}\fi}\catcode`\%=\active\def%{\%}50}}%
\end{pgfscope}%
\begin{pgfscope}%
\pgfsetbuttcap%
\pgfsetroundjoin%
\definecolor{currentfill}{rgb}{0.000000,0.000000,0.000000}%
\pgfsetfillcolor{currentfill}%
\pgfsetlinewidth{0.501875pt}%
\definecolor{currentstroke}{rgb}{0.000000,0.000000,0.000000}%
\pgfsetstrokecolor{currentstroke}%
\pgfsetdash{}{0pt}%
\pgfsys@defobject{currentmarker}{\pgfqpoint{0.000000in}{0.000000in}}{\pgfqpoint{0.000000in}{0.041667in}}{%
\pgfpathmoveto{\pgfqpoint{0.000000in}{0.000000in}}%
\pgfpathlineto{\pgfqpoint{0.000000in}{0.041667in}}%
\pgfusepath{stroke,fill}%
}%
\begin{pgfscope}%
\pgfsys@transformshift{2.335263in}{0.420000in}%
\pgfsys@useobject{currentmarker}{}%
\end{pgfscope}%
\end{pgfscope}%
\begin{pgfscope}%
\pgfsetbuttcap%
\pgfsetroundjoin%
\definecolor{currentfill}{rgb}{0.000000,0.000000,0.000000}%
\pgfsetfillcolor{currentfill}%
\pgfsetlinewidth{0.501875pt}%
\definecolor{currentstroke}{rgb}{0.000000,0.000000,0.000000}%
\pgfsetstrokecolor{currentstroke}%
\pgfsetdash{}{0pt}%
\pgfsys@defobject{currentmarker}{\pgfqpoint{0.000000in}{-0.041667in}}{\pgfqpoint{0.000000in}{0.000000in}}{%
\pgfpathmoveto{\pgfqpoint{0.000000in}{0.000000in}}%
\pgfpathlineto{\pgfqpoint{0.000000in}{-0.041667in}}%
\pgfusepath{stroke,fill}%
}%
\begin{pgfscope}%
\pgfsys@transformshift{2.335263in}{2.414488in}%
\pgfsys@useobject{currentmarker}{}%
\end{pgfscope}%
\end{pgfscope}%
\begin{pgfscope}%
\definecolor{textcolor}{rgb}{0.000000,0.000000,0.000000}%
\pgfsetstrokecolor{textcolor}%
\pgfsetfillcolor{textcolor}%
\pgftext[x=2.335263in,y=0.371389in,,top]{\color{textcolor}{\rmfamily\fontsize{11.000000}{13.200000}\selectfont\catcode`\^=\active\def^{\ifmmode\sp\else\^{}\fi}\catcode`\%=\active\def%{\%}100}}%
\end{pgfscope}%
\begin{pgfscope}%
\pgfsetbuttcap%
\pgfsetroundjoin%
\definecolor{currentfill}{rgb}{0.000000,0.000000,0.000000}%
\pgfsetfillcolor{currentfill}%
\pgfsetlinewidth{0.401500pt}%
\definecolor{currentstroke}{rgb}{0.000000,0.000000,0.000000}%
\pgfsetstrokecolor{currentstroke}%
\pgfsetdash{}{0pt}%
\pgfsys@defobject{currentmarker}{\pgfqpoint{0.000000in}{0.000000in}}{\pgfqpoint{0.000000in}{0.020833in}}{%
\pgfpathmoveto{\pgfqpoint{0.000000in}{0.000000in}}%
\pgfpathlineto{\pgfqpoint{0.000000in}{0.020833in}}%
\pgfusepath{stroke,fill}%
}%
\begin{pgfscope}%
\pgfsys@transformshift{0.818526in}{0.420000in}%
\pgfsys@useobject{currentmarker}{}%
\end{pgfscope}%
\end{pgfscope}%
\begin{pgfscope}%
\pgfsetbuttcap%
\pgfsetroundjoin%
\definecolor{currentfill}{rgb}{0.000000,0.000000,0.000000}%
\pgfsetfillcolor{currentfill}%
\pgfsetlinewidth{0.401500pt}%
\definecolor{currentstroke}{rgb}{0.000000,0.000000,0.000000}%
\pgfsetstrokecolor{currentstroke}%
\pgfsetdash{}{0pt}%
\pgfsys@defobject{currentmarker}{\pgfqpoint{0.000000in}{-0.020833in}}{\pgfqpoint{0.000000in}{0.000000in}}{%
\pgfpathmoveto{\pgfqpoint{0.000000in}{0.000000in}}%
\pgfpathlineto{\pgfqpoint{0.000000in}{-0.020833in}}%
\pgfusepath{stroke,fill}%
}%
\begin{pgfscope}%
\pgfsys@transformshift{0.818526in}{2.414488in}%
\pgfsys@useobject{currentmarker}{}%
\end{pgfscope}%
\end{pgfscope}%
\begin{pgfscope}%
\pgfsetbuttcap%
\pgfsetroundjoin%
\definecolor{currentfill}{rgb}{0.000000,0.000000,0.000000}%
\pgfsetfillcolor{currentfill}%
\pgfsetlinewidth{0.401500pt}%
\definecolor{currentstroke}{rgb}{0.000000,0.000000,0.000000}%
\pgfsetstrokecolor{currentstroke}%
\pgfsetdash{}{0pt}%
\pgfsys@defobject{currentmarker}{\pgfqpoint{0.000000in}{0.000000in}}{\pgfqpoint{0.000000in}{0.020833in}}{%
\pgfpathmoveto{\pgfqpoint{0.000000in}{0.000000in}}%
\pgfpathlineto{\pgfqpoint{0.000000in}{0.020833in}}%
\pgfusepath{stroke,fill}%
}%
\begin{pgfscope}%
\pgfsys@transformshift{0.987053in}{0.420000in}%
\pgfsys@useobject{currentmarker}{}%
\end{pgfscope}%
\end{pgfscope}%
\begin{pgfscope}%
\pgfsetbuttcap%
\pgfsetroundjoin%
\definecolor{currentfill}{rgb}{0.000000,0.000000,0.000000}%
\pgfsetfillcolor{currentfill}%
\pgfsetlinewidth{0.401500pt}%
\definecolor{currentstroke}{rgb}{0.000000,0.000000,0.000000}%
\pgfsetstrokecolor{currentstroke}%
\pgfsetdash{}{0pt}%
\pgfsys@defobject{currentmarker}{\pgfqpoint{0.000000in}{-0.020833in}}{\pgfqpoint{0.000000in}{0.000000in}}{%
\pgfpathmoveto{\pgfqpoint{0.000000in}{0.000000in}}%
\pgfpathlineto{\pgfqpoint{0.000000in}{-0.020833in}}%
\pgfusepath{stroke,fill}%
}%
\begin{pgfscope}%
\pgfsys@transformshift{0.987053in}{2.414488in}%
\pgfsys@useobject{currentmarker}{}%
\end{pgfscope}%
\end{pgfscope}%
\begin{pgfscope}%
\pgfsetbuttcap%
\pgfsetroundjoin%
\definecolor{currentfill}{rgb}{0.000000,0.000000,0.000000}%
\pgfsetfillcolor{currentfill}%
\pgfsetlinewidth{0.401500pt}%
\definecolor{currentstroke}{rgb}{0.000000,0.000000,0.000000}%
\pgfsetstrokecolor{currentstroke}%
\pgfsetdash{}{0pt}%
\pgfsys@defobject{currentmarker}{\pgfqpoint{0.000000in}{0.000000in}}{\pgfqpoint{0.000000in}{0.020833in}}{%
\pgfpathmoveto{\pgfqpoint{0.000000in}{0.000000in}}%
\pgfpathlineto{\pgfqpoint{0.000000in}{0.020833in}}%
\pgfusepath{stroke,fill}%
}%
\begin{pgfscope}%
\pgfsys@transformshift{1.155579in}{0.420000in}%
\pgfsys@useobject{currentmarker}{}%
\end{pgfscope}%
\end{pgfscope}%
\begin{pgfscope}%
\pgfsetbuttcap%
\pgfsetroundjoin%
\definecolor{currentfill}{rgb}{0.000000,0.000000,0.000000}%
\pgfsetfillcolor{currentfill}%
\pgfsetlinewidth{0.401500pt}%
\definecolor{currentstroke}{rgb}{0.000000,0.000000,0.000000}%
\pgfsetstrokecolor{currentstroke}%
\pgfsetdash{}{0pt}%
\pgfsys@defobject{currentmarker}{\pgfqpoint{0.000000in}{-0.020833in}}{\pgfqpoint{0.000000in}{0.000000in}}{%
\pgfpathmoveto{\pgfqpoint{0.000000in}{0.000000in}}%
\pgfpathlineto{\pgfqpoint{0.000000in}{-0.020833in}}%
\pgfusepath{stroke,fill}%
}%
\begin{pgfscope}%
\pgfsys@transformshift{1.155579in}{2.414488in}%
\pgfsys@useobject{currentmarker}{}%
\end{pgfscope}%
\end{pgfscope}%
\begin{pgfscope}%
\pgfsetbuttcap%
\pgfsetroundjoin%
\definecolor{currentfill}{rgb}{0.000000,0.000000,0.000000}%
\pgfsetfillcolor{currentfill}%
\pgfsetlinewidth{0.401500pt}%
\definecolor{currentstroke}{rgb}{0.000000,0.000000,0.000000}%
\pgfsetstrokecolor{currentstroke}%
\pgfsetdash{}{0pt}%
\pgfsys@defobject{currentmarker}{\pgfqpoint{0.000000in}{0.000000in}}{\pgfqpoint{0.000000in}{0.020833in}}{%
\pgfpathmoveto{\pgfqpoint{0.000000in}{0.000000in}}%
\pgfpathlineto{\pgfqpoint{0.000000in}{0.020833in}}%
\pgfusepath{stroke,fill}%
}%
\begin{pgfscope}%
\pgfsys@transformshift{1.324105in}{0.420000in}%
\pgfsys@useobject{currentmarker}{}%
\end{pgfscope}%
\end{pgfscope}%
\begin{pgfscope}%
\pgfsetbuttcap%
\pgfsetroundjoin%
\definecolor{currentfill}{rgb}{0.000000,0.000000,0.000000}%
\pgfsetfillcolor{currentfill}%
\pgfsetlinewidth{0.401500pt}%
\definecolor{currentstroke}{rgb}{0.000000,0.000000,0.000000}%
\pgfsetstrokecolor{currentstroke}%
\pgfsetdash{}{0pt}%
\pgfsys@defobject{currentmarker}{\pgfqpoint{0.000000in}{-0.020833in}}{\pgfqpoint{0.000000in}{0.000000in}}{%
\pgfpathmoveto{\pgfqpoint{0.000000in}{0.000000in}}%
\pgfpathlineto{\pgfqpoint{0.000000in}{-0.020833in}}%
\pgfusepath{stroke,fill}%
}%
\begin{pgfscope}%
\pgfsys@transformshift{1.324105in}{2.414488in}%
\pgfsys@useobject{currentmarker}{}%
\end{pgfscope}%
\end{pgfscope}%
\begin{pgfscope}%
\pgfsetbuttcap%
\pgfsetroundjoin%
\definecolor{currentfill}{rgb}{0.000000,0.000000,0.000000}%
\pgfsetfillcolor{currentfill}%
\pgfsetlinewidth{0.401500pt}%
\definecolor{currentstroke}{rgb}{0.000000,0.000000,0.000000}%
\pgfsetstrokecolor{currentstroke}%
\pgfsetdash{}{0pt}%
\pgfsys@defobject{currentmarker}{\pgfqpoint{0.000000in}{0.000000in}}{\pgfqpoint{0.000000in}{0.020833in}}{%
\pgfpathmoveto{\pgfqpoint{0.000000in}{0.000000in}}%
\pgfpathlineto{\pgfqpoint{0.000000in}{0.020833in}}%
\pgfusepath{stroke,fill}%
}%
\begin{pgfscope}%
\pgfsys@transformshift{1.661158in}{0.420000in}%
\pgfsys@useobject{currentmarker}{}%
\end{pgfscope}%
\end{pgfscope}%
\begin{pgfscope}%
\pgfsetbuttcap%
\pgfsetroundjoin%
\definecolor{currentfill}{rgb}{0.000000,0.000000,0.000000}%
\pgfsetfillcolor{currentfill}%
\pgfsetlinewidth{0.401500pt}%
\definecolor{currentstroke}{rgb}{0.000000,0.000000,0.000000}%
\pgfsetstrokecolor{currentstroke}%
\pgfsetdash{}{0pt}%
\pgfsys@defobject{currentmarker}{\pgfqpoint{0.000000in}{-0.020833in}}{\pgfqpoint{0.000000in}{0.000000in}}{%
\pgfpathmoveto{\pgfqpoint{0.000000in}{0.000000in}}%
\pgfpathlineto{\pgfqpoint{0.000000in}{-0.020833in}}%
\pgfusepath{stroke,fill}%
}%
\begin{pgfscope}%
\pgfsys@transformshift{1.661158in}{2.414488in}%
\pgfsys@useobject{currentmarker}{}%
\end{pgfscope}%
\end{pgfscope}%
\begin{pgfscope}%
\pgfsetbuttcap%
\pgfsetroundjoin%
\definecolor{currentfill}{rgb}{0.000000,0.000000,0.000000}%
\pgfsetfillcolor{currentfill}%
\pgfsetlinewidth{0.401500pt}%
\definecolor{currentstroke}{rgb}{0.000000,0.000000,0.000000}%
\pgfsetstrokecolor{currentstroke}%
\pgfsetdash{}{0pt}%
\pgfsys@defobject{currentmarker}{\pgfqpoint{0.000000in}{0.000000in}}{\pgfqpoint{0.000000in}{0.020833in}}{%
\pgfpathmoveto{\pgfqpoint{0.000000in}{0.000000in}}%
\pgfpathlineto{\pgfqpoint{0.000000in}{0.020833in}}%
\pgfusepath{stroke,fill}%
}%
\begin{pgfscope}%
\pgfsys@transformshift{1.829684in}{0.420000in}%
\pgfsys@useobject{currentmarker}{}%
\end{pgfscope}%
\end{pgfscope}%
\begin{pgfscope}%
\pgfsetbuttcap%
\pgfsetroundjoin%
\definecolor{currentfill}{rgb}{0.000000,0.000000,0.000000}%
\pgfsetfillcolor{currentfill}%
\pgfsetlinewidth{0.401500pt}%
\definecolor{currentstroke}{rgb}{0.000000,0.000000,0.000000}%
\pgfsetstrokecolor{currentstroke}%
\pgfsetdash{}{0pt}%
\pgfsys@defobject{currentmarker}{\pgfqpoint{0.000000in}{-0.020833in}}{\pgfqpoint{0.000000in}{0.000000in}}{%
\pgfpathmoveto{\pgfqpoint{0.000000in}{0.000000in}}%
\pgfpathlineto{\pgfqpoint{0.000000in}{-0.020833in}}%
\pgfusepath{stroke,fill}%
}%
\begin{pgfscope}%
\pgfsys@transformshift{1.829684in}{2.414488in}%
\pgfsys@useobject{currentmarker}{}%
\end{pgfscope}%
\end{pgfscope}%
\begin{pgfscope}%
\pgfsetbuttcap%
\pgfsetroundjoin%
\definecolor{currentfill}{rgb}{0.000000,0.000000,0.000000}%
\pgfsetfillcolor{currentfill}%
\pgfsetlinewidth{0.401500pt}%
\definecolor{currentstroke}{rgb}{0.000000,0.000000,0.000000}%
\pgfsetstrokecolor{currentstroke}%
\pgfsetdash{}{0pt}%
\pgfsys@defobject{currentmarker}{\pgfqpoint{0.000000in}{0.000000in}}{\pgfqpoint{0.000000in}{0.020833in}}{%
\pgfpathmoveto{\pgfqpoint{0.000000in}{0.000000in}}%
\pgfpathlineto{\pgfqpoint{0.000000in}{0.020833in}}%
\pgfusepath{stroke,fill}%
}%
\begin{pgfscope}%
\pgfsys@transformshift{1.998210in}{0.420000in}%
\pgfsys@useobject{currentmarker}{}%
\end{pgfscope}%
\end{pgfscope}%
\begin{pgfscope}%
\pgfsetbuttcap%
\pgfsetroundjoin%
\definecolor{currentfill}{rgb}{0.000000,0.000000,0.000000}%
\pgfsetfillcolor{currentfill}%
\pgfsetlinewidth{0.401500pt}%
\definecolor{currentstroke}{rgb}{0.000000,0.000000,0.000000}%
\pgfsetstrokecolor{currentstroke}%
\pgfsetdash{}{0pt}%
\pgfsys@defobject{currentmarker}{\pgfqpoint{0.000000in}{-0.020833in}}{\pgfqpoint{0.000000in}{0.000000in}}{%
\pgfpathmoveto{\pgfqpoint{0.000000in}{0.000000in}}%
\pgfpathlineto{\pgfqpoint{0.000000in}{-0.020833in}}%
\pgfusepath{stroke,fill}%
}%
\begin{pgfscope}%
\pgfsys@transformshift{1.998210in}{2.414488in}%
\pgfsys@useobject{currentmarker}{}%
\end{pgfscope}%
\end{pgfscope}%
\begin{pgfscope}%
\pgfsetbuttcap%
\pgfsetroundjoin%
\definecolor{currentfill}{rgb}{0.000000,0.000000,0.000000}%
\pgfsetfillcolor{currentfill}%
\pgfsetlinewidth{0.401500pt}%
\definecolor{currentstroke}{rgb}{0.000000,0.000000,0.000000}%
\pgfsetstrokecolor{currentstroke}%
\pgfsetdash{}{0pt}%
\pgfsys@defobject{currentmarker}{\pgfqpoint{0.000000in}{0.000000in}}{\pgfqpoint{0.000000in}{0.020833in}}{%
\pgfpathmoveto{\pgfqpoint{0.000000in}{0.000000in}}%
\pgfpathlineto{\pgfqpoint{0.000000in}{0.020833in}}%
\pgfusepath{stroke,fill}%
}%
\begin{pgfscope}%
\pgfsys@transformshift{2.166737in}{0.420000in}%
\pgfsys@useobject{currentmarker}{}%
\end{pgfscope}%
\end{pgfscope}%
\begin{pgfscope}%
\pgfsetbuttcap%
\pgfsetroundjoin%
\definecolor{currentfill}{rgb}{0.000000,0.000000,0.000000}%
\pgfsetfillcolor{currentfill}%
\pgfsetlinewidth{0.401500pt}%
\definecolor{currentstroke}{rgb}{0.000000,0.000000,0.000000}%
\pgfsetstrokecolor{currentstroke}%
\pgfsetdash{}{0pt}%
\pgfsys@defobject{currentmarker}{\pgfqpoint{0.000000in}{-0.020833in}}{\pgfqpoint{0.000000in}{0.000000in}}{%
\pgfpathmoveto{\pgfqpoint{0.000000in}{0.000000in}}%
\pgfpathlineto{\pgfqpoint{0.000000in}{-0.020833in}}%
\pgfusepath{stroke,fill}%
}%
\begin{pgfscope}%
\pgfsys@transformshift{2.166737in}{2.414488in}%
\pgfsys@useobject{currentmarker}{}%
\end{pgfscope}%
\end{pgfscope}%
\begin{pgfscope}%
\pgfsetbuttcap%
\pgfsetroundjoin%
\definecolor{currentfill}{rgb}{0.000000,0.000000,0.000000}%
\pgfsetfillcolor{currentfill}%
\pgfsetlinewidth{0.401500pt}%
\definecolor{currentstroke}{rgb}{0.000000,0.000000,0.000000}%
\pgfsetstrokecolor{currentstroke}%
\pgfsetdash{}{0pt}%
\pgfsys@defobject{currentmarker}{\pgfqpoint{0.000000in}{0.000000in}}{\pgfqpoint{0.000000in}{0.020833in}}{%
\pgfpathmoveto{\pgfqpoint{0.000000in}{0.000000in}}%
\pgfpathlineto{\pgfqpoint{0.000000in}{0.020833in}}%
\pgfusepath{stroke,fill}%
}%
\begin{pgfscope}%
\pgfsys@transformshift{2.503789in}{0.420000in}%
\pgfsys@useobject{currentmarker}{}%
\end{pgfscope}%
\end{pgfscope}%
\begin{pgfscope}%
\pgfsetbuttcap%
\pgfsetroundjoin%
\definecolor{currentfill}{rgb}{0.000000,0.000000,0.000000}%
\pgfsetfillcolor{currentfill}%
\pgfsetlinewidth{0.401500pt}%
\definecolor{currentstroke}{rgb}{0.000000,0.000000,0.000000}%
\pgfsetstrokecolor{currentstroke}%
\pgfsetdash{}{0pt}%
\pgfsys@defobject{currentmarker}{\pgfqpoint{0.000000in}{-0.020833in}}{\pgfqpoint{0.000000in}{0.000000in}}{%
\pgfpathmoveto{\pgfqpoint{0.000000in}{0.000000in}}%
\pgfpathlineto{\pgfqpoint{0.000000in}{-0.020833in}}%
\pgfusepath{stroke,fill}%
}%
\begin{pgfscope}%
\pgfsys@transformshift{2.503789in}{2.414488in}%
\pgfsys@useobject{currentmarker}{}%
\end{pgfscope}%
\end{pgfscope}%
\begin{pgfscope}%
\pgfsetbuttcap%
\pgfsetroundjoin%
\definecolor{currentfill}{rgb}{0.000000,0.000000,0.000000}%
\pgfsetfillcolor{currentfill}%
\pgfsetlinewidth{0.401500pt}%
\definecolor{currentstroke}{rgb}{0.000000,0.000000,0.000000}%
\pgfsetstrokecolor{currentstroke}%
\pgfsetdash{}{0pt}%
\pgfsys@defobject{currentmarker}{\pgfqpoint{0.000000in}{0.000000in}}{\pgfqpoint{0.000000in}{0.020833in}}{%
\pgfpathmoveto{\pgfqpoint{0.000000in}{0.000000in}}%
\pgfpathlineto{\pgfqpoint{0.000000in}{0.020833in}}%
\pgfusepath{stroke,fill}%
}%
\begin{pgfscope}%
\pgfsys@transformshift{2.672315in}{0.420000in}%
\pgfsys@useobject{currentmarker}{}%
\end{pgfscope}%
\end{pgfscope}%
\begin{pgfscope}%
\pgfsetbuttcap%
\pgfsetroundjoin%
\definecolor{currentfill}{rgb}{0.000000,0.000000,0.000000}%
\pgfsetfillcolor{currentfill}%
\pgfsetlinewidth{0.401500pt}%
\definecolor{currentstroke}{rgb}{0.000000,0.000000,0.000000}%
\pgfsetstrokecolor{currentstroke}%
\pgfsetdash{}{0pt}%
\pgfsys@defobject{currentmarker}{\pgfqpoint{0.000000in}{-0.020833in}}{\pgfqpoint{0.000000in}{0.000000in}}{%
\pgfpathmoveto{\pgfqpoint{0.000000in}{0.000000in}}%
\pgfpathlineto{\pgfqpoint{0.000000in}{-0.020833in}}%
\pgfusepath{stroke,fill}%
}%
\begin{pgfscope}%
\pgfsys@transformshift{2.672315in}{2.414488in}%
\pgfsys@useobject{currentmarker}{}%
\end{pgfscope}%
\end{pgfscope}%
\begin{pgfscope}%
\pgfsetbuttcap%
\pgfsetroundjoin%
\definecolor{currentfill}{rgb}{0.000000,0.000000,0.000000}%
\pgfsetfillcolor{currentfill}%
\pgfsetlinewidth{0.401500pt}%
\definecolor{currentstroke}{rgb}{0.000000,0.000000,0.000000}%
\pgfsetstrokecolor{currentstroke}%
\pgfsetdash{}{0pt}%
\pgfsys@defobject{currentmarker}{\pgfqpoint{0.000000in}{0.000000in}}{\pgfqpoint{0.000000in}{0.020833in}}{%
\pgfpathmoveto{\pgfqpoint{0.000000in}{0.000000in}}%
\pgfpathlineto{\pgfqpoint{0.000000in}{0.020833in}}%
\pgfusepath{stroke,fill}%
}%
\begin{pgfscope}%
\pgfsys@transformshift{2.840842in}{0.420000in}%
\pgfsys@useobject{currentmarker}{}%
\end{pgfscope}%
\end{pgfscope}%
\begin{pgfscope}%
\pgfsetbuttcap%
\pgfsetroundjoin%
\definecolor{currentfill}{rgb}{0.000000,0.000000,0.000000}%
\pgfsetfillcolor{currentfill}%
\pgfsetlinewidth{0.401500pt}%
\definecolor{currentstroke}{rgb}{0.000000,0.000000,0.000000}%
\pgfsetstrokecolor{currentstroke}%
\pgfsetdash{}{0pt}%
\pgfsys@defobject{currentmarker}{\pgfqpoint{0.000000in}{-0.020833in}}{\pgfqpoint{0.000000in}{0.000000in}}{%
\pgfpathmoveto{\pgfqpoint{0.000000in}{0.000000in}}%
\pgfpathlineto{\pgfqpoint{0.000000in}{-0.020833in}}%
\pgfusepath{stroke,fill}%
}%
\begin{pgfscope}%
\pgfsys@transformshift{2.840842in}{2.414488in}%
\pgfsys@useobject{currentmarker}{}%
\end{pgfscope}%
\end{pgfscope}%
\begin{pgfscope}%
\pgfsetbuttcap%
\pgfsetroundjoin%
\definecolor{currentfill}{rgb}{0.000000,0.000000,0.000000}%
\pgfsetfillcolor{currentfill}%
\pgfsetlinewidth{0.401500pt}%
\definecolor{currentstroke}{rgb}{0.000000,0.000000,0.000000}%
\pgfsetstrokecolor{currentstroke}%
\pgfsetdash{}{0pt}%
\pgfsys@defobject{currentmarker}{\pgfqpoint{0.000000in}{0.000000in}}{\pgfqpoint{0.000000in}{0.020833in}}{%
\pgfpathmoveto{\pgfqpoint{0.000000in}{0.000000in}}%
\pgfpathlineto{\pgfqpoint{0.000000in}{0.020833in}}%
\pgfusepath{stroke,fill}%
}%
\begin{pgfscope}%
\pgfsys@transformshift{3.009368in}{0.420000in}%
\pgfsys@useobject{currentmarker}{}%
\end{pgfscope}%
\end{pgfscope}%
\begin{pgfscope}%
\pgfsetbuttcap%
\pgfsetroundjoin%
\definecolor{currentfill}{rgb}{0.000000,0.000000,0.000000}%
\pgfsetfillcolor{currentfill}%
\pgfsetlinewidth{0.401500pt}%
\definecolor{currentstroke}{rgb}{0.000000,0.000000,0.000000}%
\pgfsetstrokecolor{currentstroke}%
\pgfsetdash{}{0pt}%
\pgfsys@defobject{currentmarker}{\pgfqpoint{0.000000in}{-0.020833in}}{\pgfqpoint{0.000000in}{0.000000in}}{%
\pgfpathmoveto{\pgfqpoint{0.000000in}{0.000000in}}%
\pgfpathlineto{\pgfqpoint{0.000000in}{-0.020833in}}%
\pgfusepath{stroke,fill}%
}%
\begin{pgfscope}%
\pgfsys@transformshift{3.009368in}{2.414488in}%
\pgfsys@useobject{currentmarker}{}%
\end{pgfscope}%
\end{pgfscope}%
\begin{pgfscope}%
\definecolor{textcolor}{rgb}{0.000000,0.000000,0.000000}%
\pgfsetstrokecolor{textcolor}%
\pgfsetfillcolor{textcolor}%
\pgftext[x=1.844055in,y=0.208426in,,top]{\color{textcolor}{\rmfamily\fontsize{12.000000}{14.400000}\selectfont\catcode`\^=\active\def^{\ifmmode\sp\else\^{}\fi}\catcode`\%=\active\def%{\%}$s/\delta_{0}$}}%
\end{pgfscope}%
\begin{pgfscope}%
\pgfsetbuttcap%
\pgfsetroundjoin%
\definecolor{currentfill}{rgb}{0.000000,0.000000,0.000000}%
\pgfsetfillcolor{currentfill}%
\pgfsetlinewidth{0.501875pt}%
\definecolor{currentstroke}{rgb}{0.000000,0.000000,0.000000}%
\pgfsetstrokecolor{currentstroke}%
\pgfsetdash{}{0pt}%
\pgfsys@defobject{currentmarker}{\pgfqpoint{0.000000in}{0.000000in}}{\pgfqpoint{0.041667in}{0.000000in}}{%
\pgfpathmoveto{\pgfqpoint{0.000000in}{0.000000in}}%
\pgfpathlineto{\pgfqpoint{0.041667in}{0.000000in}}%
\pgfusepath{stroke,fill}%
}%
\begin{pgfscope}%
\pgfsys@transformshift{0.650000in}{0.420000in}%
\pgfsys@useobject{currentmarker}{}%
\end{pgfscope}%
\end{pgfscope}%
\begin{pgfscope}%
\pgfsetbuttcap%
\pgfsetroundjoin%
\definecolor{currentfill}{rgb}{0.000000,0.000000,0.000000}%
\pgfsetfillcolor{currentfill}%
\pgfsetlinewidth{0.501875pt}%
\definecolor{currentstroke}{rgb}{0.000000,0.000000,0.000000}%
\pgfsetstrokecolor{currentstroke}%
\pgfsetdash{}{0pt}%
\pgfsys@defobject{currentmarker}{\pgfqpoint{-0.041667in}{0.000000in}}{\pgfqpoint{-0.000000in}{0.000000in}}{%
\pgfpathmoveto{\pgfqpoint{-0.000000in}{0.000000in}}%
\pgfpathlineto{\pgfqpoint{-0.041667in}{0.000000in}}%
\pgfusepath{stroke,fill}%
}%
\begin{pgfscope}%
\pgfsys@transformshift{3.038110in}{0.420000in}%
\pgfsys@useobject{currentmarker}{}%
\end{pgfscope}%
\end{pgfscope}%
\begin{pgfscope}%
\definecolor{textcolor}{rgb}{0.000000,0.000000,0.000000}%
\pgfsetstrokecolor{textcolor}%
\pgfsetfillcolor{textcolor}%
\pgftext[x=0.331018in, y=0.367193in, left, base]{\color{textcolor}{\rmfamily\fontsize{11.000000}{13.200000}\selectfont\catcode`\^=\active\def^{\ifmmode\sp\else\^{}\fi}\catcode`\%=\active\def%{\%}0.90}}%
\end{pgfscope}%
\begin{pgfscope}%
\pgfsetbuttcap%
\pgfsetroundjoin%
\definecolor{currentfill}{rgb}{0.000000,0.000000,0.000000}%
\pgfsetfillcolor{currentfill}%
\pgfsetlinewidth{0.501875pt}%
\definecolor{currentstroke}{rgb}{0.000000,0.000000,0.000000}%
\pgfsetstrokecolor{currentstroke}%
\pgfsetdash{}{0pt}%
\pgfsys@defobject{currentmarker}{\pgfqpoint{0.000000in}{0.000000in}}{\pgfqpoint{0.041667in}{0.000000in}}{%
\pgfpathmoveto{\pgfqpoint{0.000000in}{0.000000in}}%
\pgfpathlineto{\pgfqpoint{0.041667in}{0.000000in}}%
\pgfusepath{stroke,fill}%
}%
\begin{pgfscope}%
\pgfsys@transformshift{0.650000in}{0.918622in}%
\pgfsys@useobject{currentmarker}{}%
\end{pgfscope}%
\end{pgfscope}%
\begin{pgfscope}%
\pgfsetbuttcap%
\pgfsetroundjoin%
\definecolor{currentfill}{rgb}{0.000000,0.000000,0.000000}%
\pgfsetfillcolor{currentfill}%
\pgfsetlinewidth{0.501875pt}%
\definecolor{currentstroke}{rgb}{0.000000,0.000000,0.000000}%
\pgfsetstrokecolor{currentstroke}%
\pgfsetdash{}{0pt}%
\pgfsys@defobject{currentmarker}{\pgfqpoint{-0.041667in}{0.000000in}}{\pgfqpoint{-0.000000in}{0.000000in}}{%
\pgfpathmoveto{\pgfqpoint{-0.000000in}{0.000000in}}%
\pgfpathlineto{\pgfqpoint{-0.041667in}{0.000000in}}%
\pgfusepath{stroke,fill}%
}%
\begin{pgfscope}%
\pgfsys@transformshift{3.038110in}{0.918622in}%
\pgfsys@useobject{currentmarker}{}%
\end{pgfscope}%
\end{pgfscope}%
\begin{pgfscope}%
\definecolor{textcolor}{rgb}{0.000000,0.000000,0.000000}%
\pgfsetstrokecolor{textcolor}%
\pgfsetfillcolor{textcolor}%
\pgftext[x=0.331018in, y=0.865815in, left, base]{\color{textcolor}{\rmfamily\fontsize{11.000000}{13.200000}\selectfont\catcode`\^=\active\def^{\ifmmode\sp\else\^{}\fi}\catcode`\%=\active\def%{\%}0.95}}%
\end{pgfscope}%
\begin{pgfscope}%
\pgfsetbuttcap%
\pgfsetroundjoin%
\definecolor{currentfill}{rgb}{0.000000,0.000000,0.000000}%
\pgfsetfillcolor{currentfill}%
\pgfsetlinewidth{0.501875pt}%
\definecolor{currentstroke}{rgb}{0.000000,0.000000,0.000000}%
\pgfsetstrokecolor{currentstroke}%
\pgfsetdash{}{0pt}%
\pgfsys@defobject{currentmarker}{\pgfqpoint{0.000000in}{0.000000in}}{\pgfqpoint{0.041667in}{0.000000in}}{%
\pgfpathmoveto{\pgfqpoint{0.000000in}{0.000000in}}%
\pgfpathlineto{\pgfqpoint{0.041667in}{0.000000in}}%
\pgfusepath{stroke,fill}%
}%
\begin{pgfscope}%
\pgfsys@transformshift{0.650000in}{1.417244in}%
\pgfsys@useobject{currentmarker}{}%
\end{pgfscope}%
\end{pgfscope}%
\begin{pgfscope}%
\pgfsetbuttcap%
\pgfsetroundjoin%
\definecolor{currentfill}{rgb}{0.000000,0.000000,0.000000}%
\pgfsetfillcolor{currentfill}%
\pgfsetlinewidth{0.501875pt}%
\definecolor{currentstroke}{rgb}{0.000000,0.000000,0.000000}%
\pgfsetstrokecolor{currentstroke}%
\pgfsetdash{}{0pt}%
\pgfsys@defobject{currentmarker}{\pgfqpoint{-0.041667in}{0.000000in}}{\pgfqpoint{-0.000000in}{0.000000in}}{%
\pgfpathmoveto{\pgfqpoint{-0.000000in}{0.000000in}}%
\pgfpathlineto{\pgfqpoint{-0.041667in}{0.000000in}}%
\pgfusepath{stroke,fill}%
}%
\begin{pgfscope}%
\pgfsys@transformshift{3.038110in}{1.417244in}%
\pgfsys@useobject{currentmarker}{}%
\end{pgfscope}%
\end{pgfscope}%
\begin{pgfscope}%
\definecolor{textcolor}{rgb}{0.000000,0.000000,0.000000}%
\pgfsetstrokecolor{textcolor}%
\pgfsetfillcolor{textcolor}%
\pgftext[x=0.331018in, y=1.364437in, left, base]{\color{textcolor}{\rmfamily\fontsize{11.000000}{13.200000}\selectfont\catcode`\^=\active\def^{\ifmmode\sp\else\^{}\fi}\catcode`\%=\active\def%{\%}1.00}}%
\end{pgfscope}%
\begin{pgfscope}%
\pgfsetbuttcap%
\pgfsetroundjoin%
\definecolor{currentfill}{rgb}{0.000000,0.000000,0.000000}%
\pgfsetfillcolor{currentfill}%
\pgfsetlinewidth{0.501875pt}%
\definecolor{currentstroke}{rgb}{0.000000,0.000000,0.000000}%
\pgfsetstrokecolor{currentstroke}%
\pgfsetdash{}{0pt}%
\pgfsys@defobject{currentmarker}{\pgfqpoint{0.000000in}{0.000000in}}{\pgfqpoint{0.041667in}{0.000000in}}{%
\pgfpathmoveto{\pgfqpoint{0.000000in}{0.000000in}}%
\pgfpathlineto{\pgfqpoint{0.041667in}{0.000000in}}%
\pgfusepath{stroke,fill}%
}%
\begin{pgfscope}%
\pgfsys@transformshift{0.650000in}{1.915866in}%
\pgfsys@useobject{currentmarker}{}%
\end{pgfscope}%
\end{pgfscope}%
\begin{pgfscope}%
\pgfsetbuttcap%
\pgfsetroundjoin%
\definecolor{currentfill}{rgb}{0.000000,0.000000,0.000000}%
\pgfsetfillcolor{currentfill}%
\pgfsetlinewidth{0.501875pt}%
\definecolor{currentstroke}{rgb}{0.000000,0.000000,0.000000}%
\pgfsetstrokecolor{currentstroke}%
\pgfsetdash{}{0pt}%
\pgfsys@defobject{currentmarker}{\pgfqpoint{-0.041667in}{0.000000in}}{\pgfqpoint{-0.000000in}{0.000000in}}{%
\pgfpathmoveto{\pgfqpoint{-0.000000in}{0.000000in}}%
\pgfpathlineto{\pgfqpoint{-0.041667in}{0.000000in}}%
\pgfusepath{stroke,fill}%
}%
\begin{pgfscope}%
\pgfsys@transformshift{3.038110in}{1.915866in}%
\pgfsys@useobject{currentmarker}{}%
\end{pgfscope}%
\end{pgfscope}%
\begin{pgfscope}%
\definecolor{textcolor}{rgb}{0.000000,0.000000,0.000000}%
\pgfsetstrokecolor{textcolor}%
\pgfsetfillcolor{textcolor}%
\pgftext[x=0.331018in, y=1.863059in, left, base]{\color{textcolor}{\rmfamily\fontsize{11.000000}{13.200000}\selectfont\catcode`\^=\active\def^{\ifmmode\sp\else\^{}\fi}\catcode`\%=\active\def%{\%}1.05}}%
\end{pgfscope}%
\begin{pgfscope}%
\pgfsetbuttcap%
\pgfsetroundjoin%
\definecolor{currentfill}{rgb}{0.000000,0.000000,0.000000}%
\pgfsetfillcolor{currentfill}%
\pgfsetlinewidth{0.501875pt}%
\definecolor{currentstroke}{rgb}{0.000000,0.000000,0.000000}%
\pgfsetstrokecolor{currentstroke}%
\pgfsetdash{}{0pt}%
\pgfsys@defobject{currentmarker}{\pgfqpoint{0.000000in}{0.000000in}}{\pgfqpoint{0.041667in}{0.000000in}}{%
\pgfpathmoveto{\pgfqpoint{0.000000in}{0.000000in}}%
\pgfpathlineto{\pgfqpoint{0.041667in}{0.000000in}}%
\pgfusepath{stroke,fill}%
}%
\begin{pgfscope}%
\pgfsys@transformshift{0.650000in}{2.414488in}%
\pgfsys@useobject{currentmarker}{}%
\end{pgfscope}%
\end{pgfscope}%
\begin{pgfscope}%
\pgfsetbuttcap%
\pgfsetroundjoin%
\definecolor{currentfill}{rgb}{0.000000,0.000000,0.000000}%
\pgfsetfillcolor{currentfill}%
\pgfsetlinewidth{0.501875pt}%
\definecolor{currentstroke}{rgb}{0.000000,0.000000,0.000000}%
\pgfsetstrokecolor{currentstroke}%
\pgfsetdash{}{0pt}%
\pgfsys@defobject{currentmarker}{\pgfqpoint{-0.041667in}{0.000000in}}{\pgfqpoint{-0.000000in}{0.000000in}}{%
\pgfpathmoveto{\pgfqpoint{-0.000000in}{0.000000in}}%
\pgfpathlineto{\pgfqpoint{-0.041667in}{0.000000in}}%
\pgfusepath{stroke,fill}%
}%
\begin{pgfscope}%
\pgfsys@transformshift{3.038110in}{2.414488in}%
\pgfsys@useobject{currentmarker}{}%
\end{pgfscope}%
\end{pgfscope}%
\begin{pgfscope}%
\definecolor{textcolor}{rgb}{0.000000,0.000000,0.000000}%
\pgfsetstrokecolor{textcolor}%
\pgfsetfillcolor{textcolor}%
\pgftext[x=0.331018in, y=2.361681in, left, base]{\color{textcolor}{\rmfamily\fontsize{11.000000}{13.200000}\selectfont\catcode`\^=\active\def^{\ifmmode\sp\else\^{}\fi}\catcode`\%=\active\def%{\%}1.10}}%
\end{pgfscope}%
\begin{pgfscope}%
\pgfsetbuttcap%
\pgfsetroundjoin%
\definecolor{currentfill}{rgb}{0.000000,0.000000,0.000000}%
\pgfsetfillcolor{currentfill}%
\pgfsetlinewidth{0.401500pt}%
\definecolor{currentstroke}{rgb}{0.000000,0.000000,0.000000}%
\pgfsetstrokecolor{currentstroke}%
\pgfsetdash{}{0pt}%
\pgfsys@defobject{currentmarker}{\pgfqpoint{0.000000in}{0.000000in}}{\pgfqpoint{0.020833in}{0.000000in}}{%
\pgfpathmoveto{\pgfqpoint{0.000000in}{0.000000in}}%
\pgfpathlineto{\pgfqpoint{0.020833in}{0.000000in}}%
\pgfusepath{stroke,fill}%
}%
\begin{pgfscope}%
\pgfsys@transformshift{0.650000in}{0.519724in}%
\pgfsys@useobject{currentmarker}{}%
\end{pgfscope}%
\end{pgfscope}%
\begin{pgfscope}%
\pgfsetbuttcap%
\pgfsetroundjoin%
\definecolor{currentfill}{rgb}{0.000000,0.000000,0.000000}%
\pgfsetfillcolor{currentfill}%
\pgfsetlinewidth{0.401500pt}%
\definecolor{currentstroke}{rgb}{0.000000,0.000000,0.000000}%
\pgfsetstrokecolor{currentstroke}%
\pgfsetdash{}{0pt}%
\pgfsys@defobject{currentmarker}{\pgfqpoint{-0.020833in}{0.000000in}}{\pgfqpoint{-0.000000in}{0.000000in}}{%
\pgfpathmoveto{\pgfqpoint{-0.000000in}{0.000000in}}%
\pgfpathlineto{\pgfqpoint{-0.020833in}{0.000000in}}%
\pgfusepath{stroke,fill}%
}%
\begin{pgfscope}%
\pgfsys@transformshift{3.038110in}{0.519724in}%
\pgfsys@useobject{currentmarker}{}%
\end{pgfscope}%
\end{pgfscope}%
\begin{pgfscope}%
\pgfsetbuttcap%
\pgfsetroundjoin%
\definecolor{currentfill}{rgb}{0.000000,0.000000,0.000000}%
\pgfsetfillcolor{currentfill}%
\pgfsetlinewidth{0.401500pt}%
\definecolor{currentstroke}{rgb}{0.000000,0.000000,0.000000}%
\pgfsetstrokecolor{currentstroke}%
\pgfsetdash{}{0pt}%
\pgfsys@defobject{currentmarker}{\pgfqpoint{0.000000in}{0.000000in}}{\pgfqpoint{0.020833in}{0.000000in}}{%
\pgfpathmoveto{\pgfqpoint{0.000000in}{0.000000in}}%
\pgfpathlineto{\pgfqpoint{0.020833in}{0.000000in}}%
\pgfusepath{stroke,fill}%
}%
\begin{pgfscope}%
\pgfsys@transformshift{0.650000in}{0.619449in}%
\pgfsys@useobject{currentmarker}{}%
\end{pgfscope}%
\end{pgfscope}%
\begin{pgfscope}%
\pgfsetbuttcap%
\pgfsetroundjoin%
\definecolor{currentfill}{rgb}{0.000000,0.000000,0.000000}%
\pgfsetfillcolor{currentfill}%
\pgfsetlinewidth{0.401500pt}%
\definecolor{currentstroke}{rgb}{0.000000,0.000000,0.000000}%
\pgfsetstrokecolor{currentstroke}%
\pgfsetdash{}{0pt}%
\pgfsys@defobject{currentmarker}{\pgfqpoint{-0.020833in}{0.000000in}}{\pgfqpoint{-0.000000in}{0.000000in}}{%
\pgfpathmoveto{\pgfqpoint{-0.000000in}{0.000000in}}%
\pgfpathlineto{\pgfqpoint{-0.020833in}{0.000000in}}%
\pgfusepath{stroke,fill}%
}%
\begin{pgfscope}%
\pgfsys@transformshift{3.038110in}{0.619449in}%
\pgfsys@useobject{currentmarker}{}%
\end{pgfscope}%
\end{pgfscope}%
\begin{pgfscope}%
\pgfsetbuttcap%
\pgfsetroundjoin%
\definecolor{currentfill}{rgb}{0.000000,0.000000,0.000000}%
\pgfsetfillcolor{currentfill}%
\pgfsetlinewidth{0.401500pt}%
\definecolor{currentstroke}{rgb}{0.000000,0.000000,0.000000}%
\pgfsetstrokecolor{currentstroke}%
\pgfsetdash{}{0pt}%
\pgfsys@defobject{currentmarker}{\pgfqpoint{0.000000in}{0.000000in}}{\pgfqpoint{0.020833in}{0.000000in}}{%
\pgfpathmoveto{\pgfqpoint{0.000000in}{0.000000in}}%
\pgfpathlineto{\pgfqpoint{0.020833in}{0.000000in}}%
\pgfusepath{stroke,fill}%
}%
\begin{pgfscope}%
\pgfsys@transformshift{0.650000in}{0.719173in}%
\pgfsys@useobject{currentmarker}{}%
\end{pgfscope}%
\end{pgfscope}%
\begin{pgfscope}%
\pgfsetbuttcap%
\pgfsetroundjoin%
\definecolor{currentfill}{rgb}{0.000000,0.000000,0.000000}%
\pgfsetfillcolor{currentfill}%
\pgfsetlinewidth{0.401500pt}%
\definecolor{currentstroke}{rgb}{0.000000,0.000000,0.000000}%
\pgfsetstrokecolor{currentstroke}%
\pgfsetdash{}{0pt}%
\pgfsys@defobject{currentmarker}{\pgfqpoint{-0.020833in}{0.000000in}}{\pgfqpoint{-0.000000in}{0.000000in}}{%
\pgfpathmoveto{\pgfqpoint{-0.000000in}{0.000000in}}%
\pgfpathlineto{\pgfqpoint{-0.020833in}{0.000000in}}%
\pgfusepath{stroke,fill}%
}%
\begin{pgfscope}%
\pgfsys@transformshift{3.038110in}{0.719173in}%
\pgfsys@useobject{currentmarker}{}%
\end{pgfscope}%
\end{pgfscope}%
\begin{pgfscope}%
\pgfsetbuttcap%
\pgfsetroundjoin%
\definecolor{currentfill}{rgb}{0.000000,0.000000,0.000000}%
\pgfsetfillcolor{currentfill}%
\pgfsetlinewidth{0.401500pt}%
\definecolor{currentstroke}{rgb}{0.000000,0.000000,0.000000}%
\pgfsetstrokecolor{currentstroke}%
\pgfsetdash{}{0pt}%
\pgfsys@defobject{currentmarker}{\pgfqpoint{0.000000in}{0.000000in}}{\pgfqpoint{0.020833in}{0.000000in}}{%
\pgfpathmoveto{\pgfqpoint{0.000000in}{0.000000in}}%
\pgfpathlineto{\pgfqpoint{0.020833in}{0.000000in}}%
\pgfusepath{stroke,fill}%
}%
\begin{pgfscope}%
\pgfsys@transformshift{0.650000in}{0.818898in}%
\pgfsys@useobject{currentmarker}{}%
\end{pgfscope}%
\end{pgfscope}%
\begin{pgfscope}%
\pgfsetbuttcap%
\pgfsetroundjoin%
\definecolor{currentfill}{rgb}{0.000000,0.000000,0.000000}%
\pgfsetfillcolor{currentfill}%
\pgfsetlinewidth{0.401500pt}%
\definecolor{currentstroke}{rgb}{0.000000,0.000000,0.000000}%
\pgfsetstrokecolor{currentstroke}%
\pgfsetdash{}{0pt}%
\pgfsys@defobject{currentmarker}{\pgfqpoint{-0.020833in}{0.000000in}}{\pgfqpoint{-0.000000in}{0.000000in}}{%
\pgfpathmoveto{\pgfqpoint{-0.000000in}{0.000000in}}%
\pgfpathlineto{\pgfqpoint{-0.020833in}{0.000000in}}%
\pgfusepath{stroke,fill}%
}%
\begin{pgfscope}%
\pgfsys@transformshift{3.038110in}{0.818898in}%
\pgfsys@useobject{currentmarker}{}%
\end{pgfscope}%
\end{pgfscope}%
\begin{pgfscope}%
\pgfsetbuttcap%
\pgfsetroundjoin%
\definecolor{currentfill}{rgb}{0.000000,0.000000,0.000000}%
\pgfsetfillcolor{currentfill}%
\pgfsetlinewidth{0.401500pt}%
\definecolor{currentstroke}{rgb}{0.000000,0.000000,0.000000}%
\pgfsetstrokecolor{currentstroke}%
\pgfsetdash{}{0pt}%
\pgfsys@defobject{currentmarker}{\pgfqpoint{0.000000in}{0.000000in}}{\pgfqpoint{0.020833in}{0.000000in}}{%
\pgfpathmoveto{\pgfqpoint{0.000000in}{0.000000in}}%
\pgfpathlineto{\pgfqpoint{0.020833in}{0.000000in}}%
\pgfusepath{stroke,fill}%
}%
\begin{pgfscope}%
\pgfsys@transformshift{0.650000in}{1.018346in}%
\pgfsys@useobject{currentmarker}{}%
\end{pgfscope}%
\end{pgfscope}%
\begin{pgfscope}%
\pgfsetbuttcap%
\pgfsetroundjoin%
\definecolor{currentfill}{rgb}{0.000000,0.000000,0.000000}%
\pgfsetfillcolor{currentfill}%
\pgfsetlinewidth{0.401500pt}%
\definecolor{currentstroke}{rgb}{0.000000,0.000000,0.000000}%
\pgfsetstrokecolor{currentstroke}%
\pgfsetdash{}{0pt}%
\pgfsys@defobject{currentmarker}{\pgfqpoint{-0.020833in}{0.000000in}}{\pgfqpoint{-0.000000in}{0.000000in}}{%
\pgfpathmoveto{\pgfqpoint{-0.000000in}{0.000000in}}%
\pgfpathlineto{\pgfqpoint{-0.020833in}{0.000000in}}%
\pgfusepath{stroke,fill}%
}%
\begin{pgfscope}%
\pgfsys@transformshift{3.038110in}{1.018346in}%
\pgfsys@useobject{currentmarker}{}%
\end{pgfscope}%
\end{pgfscope}%
\begin{pgfscope}%
\pgfsetbuttcap%
\pgfsetroundjoin%
\definecolor{currentfill}{rgb}{0.000000,0.000000,0.000000}%
\pgfsetfillcolor{currentfill}%
\pgfsetlinewidth{0.401500pt}%
\definecolor{currentstroke}{rgb}{0.000000,0.000000,0.000000}%
\pgfsetstrokecolor{currentstroke}%
\pgfsetdash{}{0pt}%
\pgfsys@defobject{currentmarker}{\pgfqpoint{0.000000in}{0.000000in}}{\pgfqpoint{0.020833in}{0.000000in}}{%
\pgfpathmoveto{\pgfqpoint{0.000000in}{0.000000in}}%
\pgfpathlineto{\pgfqpoint{0.020833in}{0.000000in}}%
\pgfusepath{stroke,fill}%
}%
\begin{pgfscope}%
\pgfsys@transformshift{0.650000in}{1.118071in}%
\pgfsys@useobject{currentmarker}{}%
\end{pgfscope}%
\end{pgfscope}%
\begin{pgfscope}%
\pgfsetbuttcap%
\pgfsetroundjoin%
\definecolor{currentfill}{rgb}{0.000000,0.000000,0.000000}%
\pgfsetfillcolor{currentfill}%
\pgfsetlinewidth{0.401500pt}%
\definecolor{currentstroke}{rgb}{0.000000,0.000000,0.000000}%
\pgfsetstrokecolor{currentstroke}%
\pgfsetdash{}{0pt}%
\pgfsys@defobject{currentmarker}{\pgfqpoint{-0.020833in}{0.000000in}}{\pgfqpoint{-0.000000in}{0.000000in}}{%
\pgfpathmoveto{\pgfqpoint{-0.000000in}{0.000000in}}%
\pgfpathlineto{\pgfqpoint{-0.020833in}{0.000000in}}%
\pgfusepath{stroke,fill}%
}%
\begin{pgfscope}%
\pgfsys@transformshift{3.038110in}{1.118071in}%
\pgfsys@useobject{currentmarker}{}%
\end{pgfscope}%
\end{pgfscope}%
\begin{pgfscope}%
\pgfsetbuttcap%
\pgfsetroundjoin%
\definecolor{currentfill}{rgb}{0.000000,0.000000,0.000000}%
\pgfsetfillcolor{currentfill}%
\pgfsetlinewidth{0.401500pt}%
\definecolor{currentstroke}{rgb}{0.000000,0.000000,0.000000}%
\pgfsetstrokecolor{currentstroke}%
\pgfsetdash{}{0pt}%
\pgfsys@defobject{currentmarker}{\pgfqpoint{0.000000in}{0.000000in}}{\pgfqpoint{0.020833in}{0.000000in}}{%
\pgfpathmoveto{\pgfqpoint{0.000000in}{0.000000in}}%
\pgfpathlineto{\pgfqpoint{0.020833in}{0.000000in}}%
\pgfusepath{stroke,fill}%
}%
\begin{pgfscope}%
\pgfsys@transformshift{0.650000in}{1.217795in}%
\pgfsys@useobject{currentmarker}{}%
\end{pgfscope}%
\end{pgfscope}%
\begin{pgfscope}%
\pgfsetbuttcap%
\pgfsetroundjoin%
\definecolor{currentfill}{rgb}{0.000000,0.000000,0.000000}%
\pgfsetfillcolor{currentfill}%
\pgfsetlinewidth{0.401500pt}%
\definecolor{currentstroke}{rgb}{0.000000,0.000000,0.000000}%
\pgfsetstrokecolor{currentstroke}%
\pgfsetdash{}{0pt}%
\pgfsys@defobject{currentmarker}{\pgfqpoint{-0.020833in}{0.000000in}}{\pgfqpoint{-0.000000in}{0.000000in}}{%
\pgfpathmoveto{\pgfqpoint{-0.000000in}{0.000000in}}%
\pgfpathlineto{\pgfqpoint{-0.020833in}{0.000000in}}%
\pgfusepath{stroke,fill}%
}%
\begin{pgfscope}%
\pgfsys@transformshift{3.038110in}{1.217795in}%
\pgfsys@useobject{currentmarker}{}%
\end{pgfscope}%
\end{pgfscope}%
\begin{pgfscope}%
\pgfsetbuttcap%
\pgfsetroundjoin%
\definecolor{currentfill}{rgb}{0.000000,0.000000,0.000000}%
\pgfsetfillcolor{currentfill}%
\pgfsetlinewidth{0.401500pt}%
\definecolor{currentstroke}{rgb}{0.000000,0.000000,0.000000}%
\pgfsetstrokecolor{currentstroke}%
\pgfsetdash{}{0pt}%
\pgfsys@defobject{currentmarker}{\pgfqpoint{0.000000in}{0.000000in}}{\pgfqpoint{0.020833in}{0.000000in}}{%
\pgfpathmoveto{\pgfqpoint{0.000000in}{0.000000in}}%
\pgfpathlineto{\pgfqpoint{0.020833in}{0.000000in}}%
\pgfusepath{stroke,fill}%
}%
\begin{pgfscope}%
\pgfsys@transformshift{0.650000in}{1.317520in}%
\pgfsys@useobject{currentmarker}{}%
\end{pgfscope}%
\end{pgfscope}%
\begin{pgfscope}%
\pgfsetbuttcap%
\pgfsetroundjoin%
\definecolor{currentfill}{rgb}{0.000000,0.000000,0.000000}%
\pgfsetfillcolor{currentfill}%
\pgfsetlinewidth{0.401500pt}%
\definecolor{currentstroke}{rgb}{0.000000,0.000000,0.000000}%
\pgfsetstrokecolor{currentstroke}%
\pgfsetdash{}{0pt}%
\pgfsys@defobject{currentmarker}{\pgfqpoint{-0.020833in}{0.000000in}}{\pgfqpoint{-0.000000in}{0.000000in}}{%
\pgfpathmoveto{\pgfqpoint{-0.000000in}{0.000000in}}%
\pgfpathlineto{\pgfqpoint{-0.020833in}{0.000000in}}%
\pgfusepath{stroke,fill}%
}%
\begin{pgfscope}%
\pgfsys@transformshift{3.038110in}{1.317520in}%
\pgfsys@useobject{currentmarker}{}%
\end{pgfscope}%
\end{pgfscope}%
\begin{pgfscope}%
\pgfsetbuttcap%
\pgfsetroundjoin%
\definecolor{currentfill}{rgb}{0.000000,0.000000,0.000000}%
\pgfsetfillcolor{currentfill}%
\pgfsetlinewidth{0.401500pt}%
\definecolor{currentstroke}{rgb}{0.000000,0.000000,0.000000}%
\pgfsetstrokecolor{currentstroke}%
\pgfsetdash{}{0pt}%
\pgfsys@defobject{currentmarker}{\pgfqpoint{0.000000in}{0.000000in}}{\pgfqpoint{0.020833in}{0.000000in}}{%
\pgfpathmoveto{\pgfqpoint{0.000000in}{0.000000in}}%
\pgfpathlineto{\pgfqpoint{0.020833in}{0.000000in}}%
\pgfusepath{stroke,fill}%
}%
\begin{pgfscope}%
\pgfsys@transformshift{0.650000in}{1.516968in}%
\pgfsys@useobject{currentmarker}{}%
\end{pgfscope}%
\end{pgfscope}%
\begin{pgfscope}%
\pgfsetbuttcap%
\pgfsetroundjoin%
\definecolor{currentfill}{rgb}{0.000000,0.000000,0.000000}%
\pgfsetfillcolor{currentfill}%
\pgfsetlinewidth{0.401500pt}%
\definecolor{currentstroke}{rgb}{0.000000,0.000000,0.000000}%
\pgfsetstrokecolor{currentstroke}%
\pgfsetdash{}{0pt}%
\pgfsys@defobject{currentmarker}{\pgfqpoint{-0.020833in}{0.000000in}}{\pgfqpoint{-0.000000in}{0.000000in}}{%
\pgfpathmoveto{\pgfqpoint{-0.000000in}{0.000000in}}%
\pgfpathlineto{\pgfqpoint{-0.020833in}{0.000000in}}%
\pgfusepath{stroke,fill}%
}%
\begin{pgfscope}%
\pgfsys@transformshift{3.038110in}{1.516968in}%
\pgfsys@useobject{currentmarker}{}%
\end{pgfscope}%
\end{pgfscope}%
\begin{pgfscope}%
\pgfsetbuttcap%
\pgfsetroundjoin%
\definecolor{currentfill}{rgb}{0.000000,0.000000,0.000000}%
\pgfsetfillcolor{currentfill}%
\pgfsetlinewidth{0.401500pt}%
\definecolor{currentstroke}{rgb}{0.000000,0.000000,0.000000}%
\pgfsetstrokecolor{currentstroke}%
\pgfsetdash{}{0pt}%
\pgfsys@defobject{currentmarker}{\pgfqpoint{0.000000in}{0.000000in}}{\pgfqpoint{0.020833in}{0.000000in}}{%
\pgfpathmoveto{\pgfqpoint{0.000000in}{0.000000in}}%
\pgfpathlineto{\pgfqpoint{0.020833in}{0.000000in}}%
\pgfusepath{stroke,fill}%
}%
\begin{pgfscope}%
\pgfsys@transformshift{0.650000in}{1.616693in}%
\pgfsys@useobject{currentmarker}{}%
\end{pgfscope}%
\end{pgfscope}%
\begin{pgfscope}%
\pgfsetbuttcap%
\pgfsetroundjoin%
\definecolor{currentfill}{rgb}{0.000000,0.000000,0.000000}%
\pgfsetfillcolor{currentfill}%
\pgfsetlinewidth{0.401500pt}%
\definecolor{currentstroke}{rgb}{0.000000,0.000000,0.000000}%
\pgfsetstrokecolor{currentstroke}%
\pgfsetdash{}{0pt}%
\pgfsys@defobject{currentmarker}{\pgfqpoint{-0.020833in}{0.000000in}}{\pgfqpoint{-0.000000in}{0.000000in}}{%
\pgfpathmoveto{\pgfqpoint{-0.000000in}{0.000000in}}%
\pgfpathlineto{\pgfqpoint{-0.020833in}{0.000000in}}%
\pgfusepath{stroke,fill}%
}%
\begin{pgfscope}%
\pgfsys@transformshift{3.038110in}{1.616693in}%
\pgfsys@useobject{currentmarker}{}%
\end{pgfscope}%
\end{pgfscope}%
\begin{pgfscope}%
\pgfsetbuttcap%
\pgfsetroundjoin%
\definecolor{currentfill}{rgb}{0.000000,0.000000,0.000000}%
\pgfsetfillcolor{currentfill}%
\pgfsetlinewidth{0.401500pt}%
\definecolor{currentstroke}{rgb}{0.000000,0.000000,0.000000}%
\pgfsetstrokecolor{currentstroke}%
\pgfsetdash{}{0pt}%
\pgfsys@defobject{currentmarker}{\pgfqpoint{0.000000in}{0.000000in}}{\pgfqpoint{0.020833in}{0.000000in}}{%
\pgfpathmoveto{\pgfqpoint{0.000000in}{0.000000in}}%
\pgfpathlineto{\pgfqpoint{0.020833in}{0.000000in}}%
\pgfusepath{stroke,fill}%
}%
\begin{pgfscope}%
\pgfsys@transformshift{0.650000in}{1.716417in}%
\pgfsys@useobject{currentmarker}{}%
\end{pgfscope}%
\end{pgfscope}%
\begin{pgfscope}%
\pgfsetbuttcap%
\pgfsetroundjoin%
\definecolor{currentfill}{rgb}{0.000000,0.000000,0.000000}%
\pgfsetfillcolor{currentfill}%
\pgfsetlinewidth{0.401500pt}%
\definecolor{currentstroke}{rgb}{0.000000,0.000000,0.000000}%
\pgfsetstrokecolor{currentstroke}%
\pgfsetdash{}{0pt}%
\pgfsys@defobject{currentmarker}{\pgfqpoint{-0.020833in}{0.000000in}}{\pgfqpoint{-0.000000in}{0.000000in}}{%
\pgfpathmoveto{\pgfqpoint{-0.000000in}{0.000000in}}%
\pgfpathlineto{\pgfqpoint{-0.020833in}{0.000000in}}%
\pgfusepath{stroke,fill}%
}%
\begin{pgfscope}%
\pgfsys@transformshift{3.038110in}{1.716417in}%
\pgfsys@useobject{currentmarker}{}%
\end{pgfscope}%
\end{pgfscope}%
\begin{pgfscope}%
\pgfsetbuttcap%
\pgfsetroundjoin%
\definecolor{currentfill}{rgb}{0.000000,0.000000,0.000000}%
\pgfsetfillcolor{currentfill}%
\pgfsetlinewidth{0.401500pt}%
\definecolor{currentstroke}{rgb}{0.000000,0.000000,0.000000}%
\pgfsetstrokecolor{currentstroke}%
\pgfsetdash{}{0pt}%
\pgfsys@defobject{currentmarker}{\pgfqpoint{0.000000in}{0.000000in}}{\pgfqpoint{0.020833in}{0.000000in}}{%
\pgfpathmoveto{\pgfqpoint{0.000000in}{0.000000in}}%
\pgfpathlineto{\pgfqpoint{0.020833in}{0.000000in}}%
\pgfusepath{stroke,fill}%
}%
\begin{pgfscope}%
\pgfsys@transformshift{0.650000in}{1.816142in}%
\pgfsys@useobject{currentmarker}{}%
\end{pgfscope}%
\end{pgfscope}%
\begin{pgfscope}%
\pgfsetbuttcap%
\pgfsetroundjoin%
\definecolor{currentfill}{rgb}{0.000000,0.000000,0.000000}%
\pgfsetfillcolor{currentfill}%
\pgfsetlinewidth{0.401500pt}%
\definecolor{currentstroke}{rgb}{0.000000,0.000000,0.000000}%
\pgfsetstrokecolor{currentstroke}%
\pgfsetdash{}{0pt}%
\pgfsys@defobject{currentmarker}{\pgfqpoint{-0.020833in}{0.000000in}}{\pgfqpoint{-0.000000in}{0.000000in}}{%
\pgfpathmoveto{\pgfqpoint{-0.000000in}{0.000000in}}%
\pgfpathlineto{\pgfqpoint{-0.020833in}{0.000000in}}%
\pgfusepath{stroke,fill}%
}%
\begin{pgfscope}%
\pgfsys@transformshift{3.038110in}{1.816142in}%
\pgfsys@useobject{currentmarker}{}%
\end{pgfscope}%
\end{pgfscope}%
\begin{pgfscope}%
\pgfsetbuttcap%
\pgfsetroundjoin%
\definecolor{currentfill}{rgb}{0.000000,0.000000,0.000000}%
\pgfsetfillcolor{currentfill}%
\pgfsetlinewidth{0.401500pt}%
\definecolor{currentstroke}{rgb}{0.000000,0.000000,0.000000}%
\pgfsetstrokecolor{currentstroke}%
\pgfsetdash{}{0pt}%
\pgfsys@defobject{currentmarker}{\pgfqpoint{0.000000in}{0.000000in}}{\pgfqpoint{0.020833in}{0.000000in}}{%
\pgfpathmoveto{\pgfqpoint{0.000000in}{0.000000in}}%
\pgfpathlineto{\pgfqpoint{0.020833in}{0.000000in}}%
\pgfusepath{stroke,fill}%
}%
\begin{pgfscope}%
\pgfsys@transformshift{0.650000in}{2.015590in}%
\pgfsys@useobject{currentmarker}{}%
\end{pgfscope}%
\end{pgfscope}%
\begin{pgfscope}%
\pgfsetbuttcap%
\pgfsetroundjoin%
\definecolor{currentfill}{rgb}{0.000000,0.000000,0.000000}%
\pgfsetfillcolor{currentfill}%
\pgfsetlinewidth{0.401500pt}%
\definecolor{currentstroke}{rgb}{0.000000,0.000000,0.000000}%
\pgfsetstrokecolor{currentstroke}%
\pgfsetdash{}{0pt}%
\pgfsys@defobject{currentmarker}{\pgfqpoint{-0.020833in}{0.000000in}}{\pgfqpoint{-0.000000in}{0.000000in}}{%
\pgfpathmoveto{\pgfqpoint{-0.000000in}{0.000000in}}%
\pgfpathlineto{\pgfqpoint{-0.020833in}{0.000000in}}%
\pgfusepath{stroke,fill}%
}%
\begin{pgfscope}%
\pgfsys@transformshift{3.038110in}{2.015590in}%
\pgfsys@useobject{currentmarker}{}%
\end{pgfscope}%
\end{pgfscope}%
\begin{pgfscope}%
\pgfsetbuttcap%
\pgfsetroundjoin%
\definecolor{currentfill}{rgb}{0.000000,0.000000,0.000000}%
\pgfsetfillcolor{currentfill}%
\pgfsetlinewidth{0.401500pt}%
\definecolor{currentstroke}{rgb}{0.000000,0.000000,0.000000}%
\pgfsetstrokecolor{currentstroke}%
\pgfsetdash{}{0pt}%
\pgfsys@defobject{currentmarker}{\pgfqpoint{0.000000in}{0.000000in}}{\pgfqpoint{0.020833in}{0.000000in}}{%
\pgfpathmoveto{\pgfqpoint{0.000000in}{0.000000in}}%
\pgfpathlineto{\pgfqpoint{0.020833in}{0.000000in}}%
\pgfusepath{stroke,fill}%
}%
\begin{pgfscope}%
\pgfsys@transformshift{0.650000in}{2.115315in}%
\pgfsys@useobject{currentmarker}{}%
\end{pgfscope}%
\end{pgfscope}%
\begin{pgfscope}%
\pgfsetbuttcap%
\pgfsetroundjoin%
\definecolor{currentfill}{rgb}{0.000000,0.000000,0.000000}%
\pgfsetfillcolor{currentfill}%
\pgfsetlinewidth{0.401500pt}%
\definecolor{currentstroke}{rgb}{0.000000,0.000000,0.000000}%
\pgfsetstrokecolor{currentstroke}%
\pgfsetdash{}{0pt}%
\pgfsys@defobject{currentmarker}{\pgfqpoint{-0.020833in}{0.000000in}}{\pgfqpoint{-0.000000in}{0.000000in}}{%
\pgfpathmoveto{\pgfqpoint{-0.000000in}{0.000000in}}%
\pgfpathlineto{\pgfqpoint{-0.020833in}{0.000000in}}%
\pgfusepath{stroke,fill}%
}%
\begin{pgfscope}%
\pgfsys@transformshift{3.038110in}{2.115315in}%
\pgfsys@useobject{currentmarker}{}%
\end{pgfscope}%
\end{pgfscope}%
\begin{pgfscope}%
\pgfsetbuttcap%
\pgfsetroundjoin%
\definecolor{currentfill}{rgb}{0.000000,0.000000,0.000000}%
\pgfsetfillcolor{currentfill}%
\pgfsetlinewidth{0.401500pt}%
\definecolor{currentstroke}{rgb}{0.000000,0.000000,0.000000}%
\pgfsetstrokecolor{currentstroke}%
\pgfsetdash{}{0pt}%
\pgfsys@defobject{currentmarker}{\pgfqpoint{0.000000in}{0.000000in}}{\pgfqpoint{0.020833in}{0.000000in}}{%
\pgfpathmoveto{\pgfqpoint{0.000000in}{0.000000in}}%
\pgfpathlineto{\pgfqpoint{0.020833in}{0.000000in}}%
\pgfusepath{stroke,fill}%
}%
\begin{pgfscope}%
\pgfsys@transformshift{0.650000in}{2.215039in}%
\pgfsys@useobject{currentmarker}{}%
\end{pgfscope}%
\end{pgfscope}%
\begin{pgfscope}%
\pgfsetbuttcap%
\pgfsetroundjoin%
\definecolor{currentfill}{rgb}{0.000000,0.000000,0.000000}%
\pgfsetfillcolor{currentfill}%
\pgfsetlinewidth{0.401500pt}%
\definecolor{currentstroke}{rgb}{0.000000,0.000000,0.000000}%
\pgfsetstrokecolor{currentstroke}%
\pgfsetdash{}{0pt}%
\pgfsys@defobject{currentmarker}{\pgfqpoint{-0.020833in}{0.000000in}}{\pgfqpoint{-0.000000in}{0.000000in}}{%
\pgfpathmoveto{\pgfqpoint{-0.000000in}{0.000000in}}%
\pgfpathlineto{\pgfqpoint{-0.020833in}{0.000000in}}%
\pgfusepath{stroke,fill}%
}%
\begin{pgfscope}%
\pgfsys@transformshift{3.038110in}{2.215039in}%
\pgfsys@useobject{currentmarker}{}%
\end{pgfscope}%
\end{pgfscope}%
\begin{pgfscope}%
\pgfsetbuttcap%
\pgfsetroundjoin%
\definecolor{currentfill}{rgb}{0.000000,0.000000,0.000000}%
\pgfsetfillcolor{currentfill}%
\pgfsetlinewidth{0.401500pt}%
\definecolor{currentstroke}{rgb}{0.000000,0.000000,0.000000}%
\pgfsetstrokecolor{currentstroke}%
\pgfsetdash{}{0pt}%
\pgfsys@defobject{currentmarker}{\pgfqpoint{0.000000in}{0.000000in}}{\pgfqpoint{0.020833in}{0.000000in}}{%
\pgfpathmoveto{\pgfqpoint{0.000000in}{0.000000in}}%
\pgfpathlineto{\pgfqpoint{0.020833in}{0.000000in}}%
\pgfusepath{stroke,fill}%
}%
\begin{pgfscope}%
\pgfsys@transformshift{0.650000in}{2.314764in}%
\pgfsys@useobject{currentmarker}{}%
\end{pgfscope}%
\end{pgfscope}%
\begin{pgfscope}%
\pgfsetbuttcap%
\pgfsetroundjoin%
\definecolor{currentfill}{rgb}{0.000000,0.000000,0.000000}%
\pgfsetfillcolor{currentfill}%
\pgfsetlinewidth{0.401500pt}%
\definecolor{currentstroke}{rgb}{0.000000,0.000000,0.000000}%
\pgfsetstrokecolor{currentstroke}%
\pgfsetdash{}{0pt}%
\pgfsys@defobject{currentmarker}{\pgfqpoint{-0.020833in}{0.000000in}}{\pgfqpoint{-0.000000in}{0.000000in}}{%
\pgfpathmoveto{\pgfqpoint{-0.000000in}{0.000000in}}%
\pgfpathlineto{\pgfqpoint{-0.020833in}{0.000000in}}%
\pgfusepath{stroke,fill}%
}%
\begin{pgfscope}%
\pgfsys@transformshift{3.038110in}{2.314764in}%
\pgfsys@useobject{currentmarker}{}%
\end{pgfscope}%
\end{pgfscope}%
\begin{pgfscope}%
\definecolor{textcolor}{rgb}{0.000000,0.000000,0.000000}%
\pgfsetstrokecolor{textcolor}%
\pgfsetfillcolor{textcolor}%
\pgftext[x=0.303240in,y=1.417244in,,bottom,rotate=90.000000]{\color{textcolor}{\rmfamily\fontsize{12.000000}{14.400000}\selectfont\catcode`\^=\active\def^{\ifmmode\sp\else\^{}\fi}\catcode`\%=\active\def%{\%}$U/U_{0}$}}%
\end{pgfscope}%
\begin{pgfscope}%
\pgfpathrectangle{\pgfqpoint{0.650000in}{0.420000in}}{\pgfqpoint{2.388110in}{1.994488in}}%
\pgfusepath{clip}%
\pgfsetrectcap%
\pgfsetroundjoin%
\pgfsetlinewidth{1.003750pt}%
\definecolor{currentstroke}{rgb}{0.870588,0.466667,0.682353}%
\pgfsetstrokecolor{currentstroke}%
\pgfsetdash{}{0pt}%
\pgfpathmoveto{\pgfqpoint{0.650000in}{1.152338in}}%
\pgfpathlineto{\pgfqpoint{0.664181in}{1.154482in}}%
\pgfpathlineto{\pgfqpoint{0.679860in}{1.156698in}}%
\pgfpathlineto{\pgfqpoint{0.688412in}{1.157498in}}%
\pgfpathlineto{\pgfqpoint{0.699103in}{1.157815in}}%
\pgfpathlineto{\pgfqpoint{0.708368in}{1.157298in}}%
\pgfpathlineto{\pgfqpoint{0.738301in}{1.155833in}}%
\pgfpathlineto{\pgfqpoint{0.751843in}{1.156845in}}%
\pgfpathlineto{\pgfqpoint{0.763958in}{1.158935in}}%
\pgfpathlineto{\pgfqpoint{0.776787in}{1.162328in}}%
\pgfpathlineto{\pgfqpoint{0.791754in}{1.167526in}}%
\pgfpathlineto{\pgfqpoint{0.810997in}{1.175493in}}%
\pgfpathlineto{\pgfqpoint{0.845206in}{1.191255in}}%
\pgfpathlineto{\pgfqpoint{0.928592in}{1.230179in}}%
\pgfpathlineto{\pgfqpoint{0.987033in}{1.256243in}}%
\pgfpathlineto{\pgfqpoint{1.039060in}{1.278536in}}%
\pgfpathlineto{\pgfqpoint{1.098927in}{1.303407in}}%
\pgfpathlineto{\pgfqpoint{1.184451in}{1.337800in}}%
\pgfpathlineto{\pgfqpoint{1.256434in}{1.366073in}}%
\pgfpathlineto{\pgfqpoint{1.336969in}{1.396991in}}%
\pgfpathlineto{\pgfqpoint{1.350510in}{1.402097in}}%
\pgfpathlineto{\pgfqpoint{1.415366in}{1.426045in}}%
\pgfpathlineto{\pgfqpoint{1.490912in}{1.452929in}}%
\pgfpathlineto{\pgfqpoint{1.519420in}{1.461979in}}%
\pgfpathlineto{\pgfqpoint{1.535812in}{1.466231in}}%
\pgfpathlineto{\pgfqpoint{1.549353in}{1.468808in}}%
\pgfpathlineto{\pgfqpoint{1.554342in}{1.469304in}}%
\pgfpathlineto{\pgfqpoint{1.564320in}{1.469723in}}%
\pgfpathlineto{\pgfqpoint{1.572872in}{1.468934in}}%
\pgfpathlineto{\pgfqpoint{1.577861in}{1.467575in}}%
\pgfpathlineto{\pgfqpoint{1.584275in}{1.465157in}}%
\pgfpathlineto{\pgfqpoint{1.589977in}{1.461927in}}%
\pgfpathlineto{\pgfqpoint{1.595679in}{1.457451in}}%
\pgfpathlineto{\pgfqpoint{1.601380in}{1.451493in}}%
\pgfpathlineto{\pgfqpoint{1.607082in}{1.443928in}}%
\pgfpathlineto{\pgfqpoint{1.625612in}{1.417457in}}%
\pgfpathlineto{\pgfqpoint{1.631313in}{1.412105in}}%
\pgfpathlineto{\pgfqpoint{1.637015in}{1.408196in}}%
\pgfpathlineto{\pgfqpoint{1.643429in}{1.405211in}}%
\pgfpathlineto{\pgfqpoint{1.646993in}{1.404084in}}%
\pgfpathlineto{\pgfqpoint{1.647705in}{1.402491in}}%
\pgfpathlineto{\pgfqpoint{1.654833in}{1.401223in}}%
\pgfpathlineto{\pgfqpoint{1.663385in}{1.400954in}}%
\pgfpathlineto{\pgfqpoint{1.673363in}{1.401874in}}%
\pgfpathlineto{\pgfqpoint{1.677639in}{1.402567in}}%
\pgfpathlineto{\pgfqpoint{1.678352in}{1.401159in}}%
\pgfpathlineto{\pgfqpoint{1.691893in}{1.404200in}}%
\pgfpathlineto{\pgfqpoint{1.704009in}{1.407608in}}%
\pgfpathlineto{\pgfqpoint{1.704722in}{1.406248in}}%
\pgfpathlineto{\pgfqpoint{1.726103in}{1.413203in}}%
\pgfpathlineto{\pgfqpoint{1.726816in}{1.411861in}}%
\pgfpathlineto{\pgfqpoint{1.750335in}{1.420139in}}%
\pgfpathlineto{\pgfqpoint{1.751048in}{1.418799in}}%
\pgfpathlineto{\pgfqpoint{1.773141in}{1.426860in}}%
\pgfpathlineto{\pgfqpoint{1.773854in}{1.425517in}}%
\pgfpathlineto{\pgfqpoint{1.795235in}{1.433226in}}%
\pgfpathlineto{\pgfqpoint{1.795948in}{1.431866in}}%
\pgfpathlineto{\pgfqpoint{1.818042in}{1.439646in}}%
\pgfpathlineto{\pgfqpoint{1.818754in}{1.438277in}}%
\pgfpathlineto{\pgfqpoint{1.843699in}{1.447016in}}%
\pgfpathlineto{\pgfqpoint{1.844412in}{1.445652in}}%
\pgfpathlineto{\pgfqpoint{1.870069in}{1.454877in}}%
\pgfpathlineto{\pgfqpoint{1.870782in}{1.453512in}}%
\pgfpathlineto{\pgfqpoint{1.897864in}{1.463319in}}%
\pgfpathlineto{\pgfqpoint{1.898577in}{1.461944in}}%
\pgfpathlineto{\pgfqpoint{1.928510in}{1.472547in}}%
\pgfpathlineto{\pgfqpoint{1.929223in}{1.471153in}}%
\pgfpathlineto{\pgfqpoint{1.962720in}{1.482493in}}%
\pgfpathlineto{\pgfqpoint{1.963433in}{1.481080in}}%
\pgfpathlineto{\pgfqpoint{1.995504in}{1.491922in}}%
\pgfpathlineto{\pgfqpoint{1.996217in}{1.490504in}}%
\pgfpathlineto{\pgfqpoint{2.028289in}{1.501492in}}%
\pgfpathlineto{\pgfqpoint{2.029001in}{1.500065in}}%
\pgfpathlineto{\pgfqpoint{2.061786in}{1.511184in}}%
\pgfpathlineto{\pgfqpoint{2.062498in}{1.509741in}}%
\pgfpathlineto{\pgfqpoint{2.095995in}{1.520838in}}%
\pgfpathlineto{\pgfqpoint{2.096708in}{1.519383in}}%
\pgfpathlineto{\pgfqpoint{2.128780in}{1.529814in}}%
\pgfpathlineto{\pgfqpoint{2.129492in}{1.528339in}}%
\pgfpathlineto{\pgfqpoint{2.158713in}{1.537778in}}%
\pgfpathlineto{\pgfqpoint{2.159426in}{1.536286in}}%
\pgfpathlineto{\pgfqpoint{2.188647in}{1.545863in}}%
\pgfpathlineto{\pgfqpoint{2.189360in}{1.544357in}}%
\pgfpathlineto{\pgfqpoint{2.224281in}{1.555845in}}%
\pgfpathlineto{\pgfqpoint{2.224994in}{1.554319in}}%
\pgfpathlineto{\pgfqpoint{2.259916in}{1.565474in}}%
\pgfpathlineto{\pgfqpoint{2.260629in}{1.563920in}}%
\pgfpathlineto{\pgfqpoint{2.297689in}{1.575506in}}%
\pgfpathlineto{\pgfqpoint{2.298402in}{1.573929in}}%
\pgfpathlineto{\pgfqpoint{2.336175in}{1.585779in}}%
\pgfpathlineto{\pgfqpoint{2.336888in}{1.584177in}}%
\pgfpathlineto{\pgfqpoint{2.372524in}{1.595365in}}%
\pgfpathlineto{\pgfqpoint{2.373236in}{1.593727in}}%
\pgfpathlineto{\pgfqpoint{2.408872in}{1.604627in}}%
\pgfpathlineto{\pgfqpoint{2.409584in}{1.602951in}}%
\pgfpathlineto{\pgfqpoint{2.442368in}{1.613205in}}%
\pgfpathlineto{\pgfqpoint{2.443081in}{1.611511in}}%
\pgfpathlineto{\pgfqpoint{2.470164in}{1.620843in}}%
\pgfpathlineto{\pgfqpoint{2.480142in}{1.624744in}}%
\pgfpathlineto{\pgfqpoint{2.480854in}{1.623094in}}%
\pgfpathlineto{\pgfqpoint{2.494396in}{1.629473in}}%
\pgfpathlineto{\pgfqpoint{2.505086in}{1.635699in}}%
\pgfpathlineto{\pgfqpoint{2.510788in}{1.639677in}}%
\pgfpathlineto{\pgfqpoint{2.511501in}{1.638262in}}%
\pgfpathlineto{\pgfqpoint{2.520053in}{1.645551in}}%
\pgfpathlineto{\pgfqpoint{2.527180in}{1.653078in}}%
\pgfpathlineto{\pgfqpoint{2.531456in}{1.658451in}}%
\pgfpathlineto{\pgfqpoint{2.532169in}{1.657501in}}%
\pgfpathlineto{\pgfqpoint{2.538583in}{1.667325in}}%
\pgfpathlineto{\pgfqpoint{2.544285in}{1.678053in}}%
\pgfpathlineto{\pgfqpoint{2.547136in}{1.684300in}}%
\pgfpathlineto{\pgfqpoint{2.547849in}{1.684179in}}%
\pgfpathlineto{\pgfqpoint{2.553550in}{1.699340in}}%
\pgfpathlineto{\pgfqpoint{2.563528in}{1.730419in}}%
\pgfpathlineto{\pgfqpoint{2.570655in}{1.752737in}}%
\pgfpathlineto{\pgfqpoint{2.571368in}{1.753665in}}%
\pgfpathlineto{\pgfqpoint{2.577069in}{1.768027in}}%
\pgfpathlineto{\pgfqpoint{2.582771in}{1.779685in}}%
\pgfpathlineto{\pgfqpoint{2.585622in}{1.783779in}}%
\pgfpathlineto{\pgfqpoint{2.592036in}{1.793332in}}%
\pgfpathlineto{\pgfqpoint{2.599164in}{1.801822in}}%
\pgfpathlineto{\pgfqpoint{2.600589in}{1.803310in}}%
\pgfpathlineto{\pgfqpoint{2.601302in}{1.803160in}}%
\pgfpathlineto{\pgfqpoint{2.609142in}{1.810183in}}%
\pgfpathlineto{\pgfqpoint{2.617694in}{1.816332in}}%
\pgfpathlineto{\pgfqpoint{2.621970in}{1.818962in}}%
\pgfpathlineto{\pgfqpoint{2.622683in}{1.818466in}}%
\pgfpathlineto{\pgfqpoint{2.632661in}{1.823614in}}%
\pgfpathlineto{\pgfqpoint{2.645490in}{1.828865in}}%
\pgfpathlineto{\pgfqpoint{2.649054in}{1.830144in}}%
\pgfpathlineto{\pgfqpoint{2.649766in}{1.829385in}}%
\pgfpathlineto{\pgfqpoint{2.669010in}{1.835584in}}%
\pgfpathlineto{\pgfqpoint{2.679700in}{1.838520in}}%
\pgfpathlineto{\pgfqpoint{2.680413in}{1.837560in}}%
\pgfpathlineto{\pgfqpoint{2.712485in}{1.845716in}}%
\pgfpathlineto{\pgfqpoint{2.713198in}{1.844560in}}%
\pgfpathlineto{\pgfqpoint{2.745271in}{1.852039in}}%
\pgfpathlineto{\pgfqpoint{2.745984in}{1.850665in}}%
\pgfpathlineto{\pgfqpoint{2.772354in}{1.856879in}}%
\pgfpathlineto{\pgfqpoint{2.773067in}{1.855260in}}%
\pgfpathlineto{\pgfqpoint{2.799437in}{1.861368in}}%
\pgfpathlineto{\pgfqpoint{2.800150in}{1.859441in}}%
\pgfpathlineto{\pgfqpoint{2.826520in}{1.865698in}}%
\pgfpathlineto{\pgfqpoint{2.827233in}{1.863371in}}%
\pgfpathlineto{\pgfqpoint{2.851465in}{1.868912in}}%
\pgfpathlineto{\pgfqpoint{2.852178in}{1.866086in}}%
\pgfpathlineto{\pgfqpoint{2.879259in}{1.872017in}}%
\pgfpathlineto{\pgfqpoint{2.879971in}{1.868494in}}%
\pgfpathlineto{\pgfqpoint{2.905627in}{1.873482in}}%
\pgfpathlineto{\pgfqpoint{2.906340in}{1.869029in}}%
\pgfpathlineto{\pgfqpoint{2.932708in}{1.874192in}}%
\pgfpathlineto{\pgfqpoint{2.933421in}{1.868521in}}%
\pgfpathlineto{\pgfqpoint{2.959077in}{1.873408in}}%
\pgfpathlineto{\pgfqpoint{2.959789in}{1.866159in}}%
\pgfpathlineto{\pgfqpoint{2.981882in}{1.870862in}}%
\pgfpathlineto{\pgfqpoint{2.982595in}{1.861635in}}%
\pgfpathlineto{\pgfqpoint{3.004688in}{1.867157in}}%
\pgfpathlineto{\pgfqpoint{3.005400in}{1.855398in}}%
\pgfpathlineto{\pgfqpoint{3.022504in}{1.860497in}}%
\pgfpathlineto{\pgfqpoint{3.026066in}{1.861707in}}%
\pgfpathlineto{\pgfqpoint{3.026778in}{1.846651in}}%
\pgfpathlineto{\pgfqpoint{3.038110in}{1.850265in}}%
\pgfpathlineto{\pgfqpoint{3.038110in}{1.850265in}}%
\pgfusepath{stroke}%
\end{pgfscope}%
\begin{pgfscope}%
\pgfpathrectangle{\pgfqpoint{0.650000in}{0.420000in}}{\pgfqpoint{2.388110in}{1.994488in}}%
\pgfusepath{clip}%
\pgfsetrectcap%
\pgfsetroundjoin%
\pgfsetlinewidth{1.003750pt}%
\definecolor{currentstroke}{rgb}{0.498039,0.737255,0.254902}%
\pgfsetstrokecolor{currentstroke}%
\pgfsetdash{}{0pt}%
\pgfpathmoveto{\pgfqpoint{0.650000in}{1.134301in}}%
\pgfpathlineto{\pgfqpoint{0.664069in}{1.136447in}}%
\pgfpathlineto{\pgfqpoint{0.679625in}{1.138661in}}%
\pgfpathlineto{\pgfqpoint{0.688110in}{1.139461in}}%
\pgfpathlineto{\pgfqpoint{0.698717in}{1.139783in}}%
\pgfpathlineto{\pgfqpoint{0.707909in}{1.139287in}}%
\pgfpathlineto{\pgfqpoint{0.738314in}{1.137848in}}%
\pgfpathlineto{\pgfqpoint{0.751041in}{1.138841in}}%
\pgfpathlineto{\pgfqpoint{0.763062in}{1.141000in}}%
\pgfpathlineto{\pgfqpoint{0.775790in}{1.144482in}}%
\pgfpathlineto{\pgfqpoint{0.791346in}{1.150012in}}%
\pgfpathlineto{\pgfqpoint{0.813973in}{1.159544in}}%
\pgfpathlineto{\pgfqpoint{0.852863in}{1.177682in}}%
\pgfpathlineto{\pgfqpoint{0.905895in}{1.202368in}}%
\pgfpathlineto{\pgfqpoint{0.951856in}{1.223085in}}%
\pgfpathlineto{\pgfqpoint{1.007716in}{1.247706in}}%
\pgfpathlineto{\pgfqpoint{1.016201in}{1.251355in}}%
\pgfpathlineto{\pgfqpoint{1.067112in}{1.272879in}}%
\pgfpathlineto{\pgfqpoint{1.148427in}{1.306282in}}%
\pgfpathlineto{\pgfqpoint{1.251663in}{1.348022in}}%
\pgfpathlineto{\pgfqpoint{1.296210in}{1.364819in}}%
\pgfpathlineto{\pgfqpoint{1.333685in}{1.378017in}}%
\pgfpathlineto{\pgfqpoint{1.353484in}{1.383967in}}%
\pgfpathlineto{\pgfqpoint{1.369747in}{1.387778in}}%
\pgfpathlineto{\pgfqpoint{1.395910in}{1.393895in}}%
\pgfpathlineto{\pgfqpoint{1.401566in}{1.396734in}}%
\pgfpathlineto{\pgfqpoint{1.407223in}{1.400968in}}%
\pgfpathlineto{\pgfqpoint{1.414294in}{1.407947in}}%
\pgfpathlineto{\pgfqpoint{1.421365in}{1.417244in}}%
\pgfpathlineto{\pgfqpoint{1.429143in}{1.429481in}}%
\pgfpathlineto{\pgfqpoint{1.439042in}{1.447293in}}%
\pgfpathlineto{\pgfqpoint{1.451770in}{1.473130in}}%
\pgfpathlineto{\pgfqpoint{1.482882in}{1.537231in}}%
\pgfpathlineto{\pgfqpoint{1.492781in}{1.554516in}}%
\pgfpathlineto{\pgfqpoint{1.501266in}{1.566996in}}%
\pgfpathlineto{\pgfqpoint{1.509044in}{1.576135in}}%
\pgfpathlineto{\pgfqpoint{1.515408in}{1.581530in}}%
\pgfpathlineto{\pgfqpoint{1.521065in}{1.584834in}}%
\pgfpathlineto{\pgfqpoint{1.526015in}{1.586554in}}%
\pgfpathlineto{\pgfqpoint{1.530964in}{1.587182in}}%
\pgfpathlineto{\pgfqpoint{1.535914in}{1.586733in}}%
\pgfpathlineto{\pgfqpoint{1.540864in}{1.585232in}}%
\pgfpathlineto{\pgfqpoint{1.546520in}{1.582289in}}%
\pgfpathlineto{\pgfqpoint{1.552177in}{1.578042in}}%
\pgfpathlineto{\pgfqpoint{1.558541in}{1.572112in}}%
\pgfpathlineto{\pgfqpoint{1.566319in}{1.563246in}}%
\pgfpathlineto{\pgfqpoint{1.575511in}{1.550942in}}%
\pgfpathlineto{\pgfqpoint{1.587532in}{1.532798in}}%
\pgfpathlineto{\pgfqpoint{1.598138in}{1.514990in}}%
\pgfpathlineto{\pgfqpoint{1.638442in}{1.446668in}}%
\pgfpathlineto{\pgfqpoint{1.639856in}{1.443114in}}%
\pgfpathlineto{\pgfqpoint{1.656120in}{1.417967in}}%
\pgfpathlineto{\pgfqpoint{1.662483in}{1.408779in}}%
\pgfpathlineto{\pgfqpoint{1.663898in}{1.405278in}}%
\pgfpathlineto{\pgfqpoint{1.676625in}{1.388365in}}%
\pgfpathlineto{\pgfqpoint{1.682282in}{1.381530in}}%
\pgfpathlineto{\pgfqpoint{1.683696in}{1.378244in}}%
\pgfpathlineto{\pgfqpoint{1.694303in}{1.366910in}}%
\pgfpathlineto{\pgfqpoint{1.702081in}{1.359855in}}%
\pgfpathlineto{\pgfqpoint{1.703495in}{1.356963in}}%
\pgfpathlineto{\pgfqpoint{1.711980in}{1.350893in}}%
\pgfpathlineto{\pgfqpoint{1.719051in}{1.347106in}}%
\pgfpathlineto{\pgfqpoint{1.725415in}{1.344769in}}%
\pgfpathlineto{\pgfqpoint{1.726122in}{1.342790in}}%
\pgfpathlineto{\pgfqpoint{1.732486in}{1.341652in}}%
\pgfpathlineto{\pgfqpoint{1.738850in}{1.341640in}}%
\pgfpathlineto{\pgfqpoint{1.745214in}{1.342766in}}%
\pgfpathlineto{\pgfqpoint{1.748042in}{1.343625in}}%
\pgfpathlineto{\pgfqpoint{1.748749in}{1.342107in}}%
\pgfpathlineto{\pgfqpoint{1.755113in}{1.344955in}}%
\pgfpathlineto{\pgfqpoint{1.762184in}{1.349298in}}%
\pgfpathlineto{\pgfqpoint{1.770669in}{1.355907in}}%
\pgfpathlineto{\pgfqpoint{1.772083in}{1.357134in}}%
\pgfpathlineto{\pgfqpoint{1.772791in}{1.356095in}}%
\pgfpathlineto{\pgfqpoint{1.784811in}{1.367664in}}%
\pgfpathlineto{\pgfqpoint{1.798246in}{1.381185in}}%
\pgfpathlineto{\pgfqpoint{1.798953in}{1.380390in}}%
\pgfpathlineto{\pgfqpoint{1.809560in}{1.389956in}}%
\pgfpathlineto{\pgfqpoint{1.817338in}{1.395578in}}%
\pgfpathlineto{\pgfqpoint{1.824409in}{1.399304in}}%
\pgfpathlineto{\pgfqpoint{1.825116in}{1.398260in}}%
\pgfpathlineto{\pgfqpoint{1.831479in}{1.400163in}}%
\pgfpathlineto{\pgfqpoint{1.837136in}{1.400710in}}%
\pgfpathlineto{\pgfqpoint{1.842793in}{1.400155in}}%
\pgfpathlineto{\pgfqpoint{1.848450in}{1.398516in}}%
\pgfpathlineto{\pgfqpoint{1.849864in}{1.396614in}}%
\pgfpathlineto{\pgfqpoint{1.856228in}{1.393255in}}%
\pgfpathlineto{\pgfqpoint{1.863299in}{1.388210in}}%
\pgfpathlineto{\pgfqpoint{1.871784in}{1.380713in}}%
\pgfpathlineto{\pgfqpoint{1.873905in}{1.378652in}}%
\pgfpathlineto{\pgfqpoint{1.875319in}{1.375793in}}%
\pgfpathlineto{\pgfqpoint{1.900067in}{1.350037in}}%
\pgfpathlineto{\pgfqpoint{1.901482in}{1.347041in}}%
\pgfpathlineto{\pgfqpoint{1.911381in}{1.338219in}}%
\pgfpathlineto{\pgfqpoint{1.919866in}{1.332096in}}%
\pgfpathlineto{\pgfqpoint{1.925523in}{1.328863in}}%
\pgfpathlineto{\pgfqpoint{1.926230in}{1.326791in}}%
\pgfpathlineto{\pgfqpoint{1.934008in}{1.323615in}}%
\pgfpathlineto{\pgfqpoint{1.941079in}{1.321895in}}%
\pgfpathlineto{\pgfqpoint{1.948857in}{1.321212in}}%
\pgfpathlineto{\pgfqpoint{1.950271in}{1.321219in}}%
\pgfpathlineto{\pgfqpoint{1.950978in}{1.319511in}}%
\pgfpathlineto{\pgfqpoint{1.958756in}{1.320303in}}%
\pgfpathlineto{\pgfqpoint{1.967242in}{1.322279in}}%
\pgfpathlineto{\pgfqpoint{1.977141in}{1.325740in}}%
\pgfpathlineto{\pgfqpoint{1.978555in}{1.326315in}}%
\pgfpathlineto{\pgfqpoint{1.979262in}{1.324953in}}%
\pgfpathlineto{\pgfqpoint{1.991990in}{1.330927in}}%
\pgfpathlineto{\pgfqpoint{2.008960in}{1.339905in}}%
\pgfpathlineto{\pgfqpoint{2.009667in}{1.338730in}}%
\pgfpathlineto{\pgfqpoint{2.032294in}{1.350354in}}%
\pgfpathlineto{\pgfqpoint{2.038658in}{1.353208in}}%
\pgfpathlineto{\pgfqpoint{2.039365in}{1.352032in}}%
\pgfpathlineto{\pgfqpoint{2.052093in}{1.356826in}}%
\pgfpathlineto{\pgfqpoint{2.063407in}{1.359944in}}%
\pgfpathlineto{\pgfqpoint{2.067649in}{1.360812in}}%
\pgfpathlineto{\pgfqpoint{2.068356in}{1.359507in}}%
\pgfpathlineto{\pgfqpoint{2.078963in}{1.360834in}}%
\pgfpathlineto{\pgfqpoint{2.089569in}{1.361064in}}%
\pgfpathlineto{\pgfqpoint{2.097347in}{1.360523in}}%
\pgfpathlineto{\pgfqpoint{2.098054in}{1.359021in}}%
\pgfpathlineto{\pgfqpoint{2.109368in}{1.357194in}}%
\pgfpathlineto{\pgfqpoint{2.121388in}{1.354082in}}%
\pgfpathlineto{\pgfqpoint{2.125631in}{1.352717in}}%
\pgfpathlineto{\pgfqpoint{2.126338in}{1.351022in}}%
\pgfpathlineto{\pgfqpoint{2.140480in}{1.345520in}}%
\pgfpathlineto{\pgfqpoint{2.151086in}{1.340684in}}%
\pgfpathlineto{\pgfqpoint{2.152501in}{1.338496in}}%
\pgfpathlineto{\pgfqpoint{2.175835in}{1.326806in}}%
\pgfpathlineto{\pgfqpoint{2.177249in}{1.324521in}}%
\pgfpathlineto{\pgfqpoint{2.195633in}{1.315914in}}%
\pgfpathlineto{\pgfqpoint{2.200583in}{1.313931in}}%
\pgfpathlineto{\pgfqpoint{2.201290in}{1.312027in}}%
\pgfpathlineto{\pgfqpoint{2.211896in}{1.308579in}}%
\pgfpathlineto{\pgfqpoint{2.221795in}{1.306516in}}%
\pgfpathlineto{\pgfqpoint{2.230280in}{1.305829in}}%
\pgfpathlineto{\pgfqpoint{2.232401in}{1.305820in}}%
\pgfpathlineto{\pgfqpoint{2.233108in}{1.304165in}}%
\pgfpathlineto{\pgfqpoint{2.241594in}{1.304883in}}%
\pgfpathlineto{\pgfqpoint{2.250786in}{1.306796in}}%
\pgfpathlineto{\pgfqpoint{2.260685in}{1.310029in}}%
\pgfpathlineto{\pgfqpoint{2.263514in}{1.311140in}}%
\pgfpathlineto{\pgfqpoint{2.264221in}{1.309821in}}%
\pgfpathlineto{\pgfqpoint{2.281191in}{1.317423in}}%
\pgfpathlineto{\pgfqpoint{2.289676in}{1.321237in}}%
\pgfpathlineto{\pgfqpoint{2.290383in}{1.320025in}}%
\pgfpathlineto{\pgfqpoint{2.303111in}{1.324917in}}%
\pgfpathlineto{\pgfqpoint{2.313718in}{1.327776in}}%
\pgfpathlineto{\pgfqpoint{2.321496in}{1.329077in}}%
\pgfpathlineto{\pgfqpoint{2.322203in}{1.327701in}}%
\pgfpathlineto{\pgfqpoint{2.332809in}{1.328306in}}%
\pgfpathlineto{\pgfqpoint{2.346951in}{1.327825in}}%
\pgfpathlineto{\pgfqpoint{2.354022in}{1.327327in}}%
\pgfpathlineto{\pgfqpoint{2.354729in}{1.325774in}}%
\pgfpathlineto{\pgfqpoint{2.373114in}{1.324965in}}%
\pgfpathlineto{\pgfqpoint{2.383720in}{1.325443in}}%
\pgfpathlineto{\pgfqpoint{2.384427in}{1.323975in}}%
\pgfpathlineto{\pgfqpoint{2.396448in}{1.325727in}}%
\pgfpathlineto{\pgfqpoint{2.407762in}{1.328506in}}%
\pgfpathlineto{\pgfqpoint{2.408469in}{1.328712in}}%
\pgfpathlineto{\pgfqpoint{2.409176in}{1.327432in}}%
\pgfpathlineto{\pgfqpoint{2.425439in}{1.332965in}}%
\pgfpathlineto{\pgfqpoint{2.430389in}{1.334802in}}%
\pgfpathlineto{\pgfqpoint{2.431096in}{1.333661in}}%
\pgfpathlineto{\pgfqpoint{2.452309in}{1.341081in}}%
\pgfpathlineto{\pgfqpoint{2.453016in}{1.339994in}}%
\pgfpathlineto{\pgfqpoint{2.464329in}{1.342741in}}%
\pgfpathlineto{\pgfqpoint{2.472108in}{1.343952in}}%
\pgfpathlineto{\pgfqpoint{2.472815in}{1.342840in}}%
\pgfpathlineto{\pgfqpoint{2.482007in}{1.343292in}}%
\pgfpathlineto{\pgfqpoint{2.489078in}{1.342859in}}%
\pgfpathlineto{\pgfqpoint{2.489785in}{1.341680in}}%
\pgfpathlineto{\pgfqpoint{2.498977in}{1.340062in}}%
\pgfpathlineto{\pgfqpoint{2.503927in}{1.338696in}}%
\pgfpathlineto{\pgfqpoint{2.505341in}{1.337225in}}%
\pgfpathlineto{\pgfqpoint{2.515240in}{1.333347in}}%
\pgfpathlineto{\pgfqpoint{2.516655in}{1.332693in}}%
\pgfpathlineto{\pgfqpoint{2.518069in}{1.331069in}}%
\pgfpathlineto{\pgfqpoint{2.527968in}{1.325688in}}%
\pgfpathlineto{\pgfqpoint{2.529383in}{1.323954in}}%
\pgfpathlineto{\pgfqpoint{2.537867in}{1.318542in}}%
\pgfpathlineto{\pgfqpoint{2.539282in}{1.316796in}}%
\pgfpathlineto{\pgfqpoint{2.547767in}{1.311249in}}%
\pgfpathlineto{\pgfqpoint{2.549181in}{1.309599in}}%
\pgfpathlineto{\pgfqpoint{2.567565in}{1.297531in}}%
\pgfpathlineto{\pgfqpoint{2.568980in}{1.296044in}}%
\pgfpathlineto{\pgfqpoint{2.596557in}{1.279916in}}%
\pgfpathlineto{\pgfqpoint{2.612821in}{1.273179in}}%
\pgfpathlineto{\pgfqpoint{2.630498in}{1.267849in}}%
\pgfpathlineto{\pgfqpoint{2.643226in}{1.265568in}}%
\pgfpathlineto{\pgfqpoint{2.663025in}{1.263298in}}%
\pgfpathlineto{\pgfqpoint{2.680703in}{1.262728in}}%
\pgfpathlineto{\pgfqpoint{2.703330in}{1.263753in}}%
\pgfpathlineto{\pgfqpoint{2.729494in}{1.266244in}}%
\pgfpathlineto{\pgfqpoint{2.757778in}{1.270401in}}%
\pgfpathlineto{\pgfqpoint{2.779699in}{1.274606in}}%
\pgfpathlineto{\pgfqpoint{2.795962in}{1.278829in}}%
\pgfpathlineto{\pgfqpoint{2.821419in}{1.286513in}}%
\pgfpathlineto{\pgfqpoint{2.832731in}{1.290268in}}%
\pgfpathlineto{\pgfqpoint{2.834145in}{1.291758in}}%
\pgfpathlineto{\pgfqpoint{2.851822in}{1.298175in}}%
\pgfpathlineto{\pgfqpoint{2.852529in}{1.301880in}}%
\pgfpathlineto{\pgfqpoint{2.976264in}{1.356040in}}%
\pgfpathlineto{\pgfqpoint{2.985456in}{1.361449in}}%
\pgfpathlineto{\pgfqpoint{3.001718in}{1.371391in}}%
\pgfpathlineto{\pgfqpoint{3.007373in}{1.373177in}}%
\pgfpathlineto{\pgfqpoint{3.014945in}{1.374230in}}%
\pgfpathlineto{\pgfqpoint{3.019322in}{1.374574in}}%
\pgfpathlineto{\pgfqpoint{3.019322in}{1.374574in}}%
\pgfusepath{stroke}%
\end{pgfscope}%
\begin{pgfscope}%
\pgfpathrectangle{\pgfqpoint{0.650000in}{0.420000in}}{\pgfqpoint{2.388110in}{1.994488in}}%
\pgfusepath{clip}%
\pgfsetrectcap%
\pgfsetroundjoin%
\pgfsetlinewidth{0.803000pt}%
\definecolor{currentstroke}{rgb}{0.500000,0.500000,0.500000}%
\pgfsetstrokecolor{currentstroke}%
\pgfsetdash{}{0pt}%
\pgfpathmoveto{\pgfqpoint{0.650000in}{1.417244in}}%
\pgfpathlineto{\pgfqpoint{3.038110in}{1.417244in}}%
\pgfusepath{stroke}%
\end{pgfscope}%
\begin{pgfscope}%
\pgfsetrectcap%
\pgfsetmiterjoin%
\pgfsetlinewidth{0.803000pt}%
\definecolor{currentstroke}{rgb}{0.000000,0.000000,0.000000}%
\pgfsetstrokecolor{currentstroke}%
\pgfsetdash{}{0pt}%
\pgfpathmoveto{\pgfqpoint{0.650000in}{0.420000in}}%
\pgfpathlineto{\pgfqpoint{0.650000in}{2.414488in}}%
\pgfusepath{stroke}%
\end{pgfscope}%
\begin{pgfscope}%
\pgfsetrectcap%
\pgfsetmiterjoin%
\pgfsetlinewidth{0.803000pt}%
\definecolor{currentstroke}{rgb}{0.000000,0.000000,0.000000}%
\pgfsetstrokecolor{currentstroke}%
\pgfsetdash{}{0pt}%
\pgfpathmoveto{\pgfqpoint{3.038110in}{0.420000in}}%
\pgfpathlineto{\pgfqpoint{3.038110in}{2.414488in}}%
\pgfusepath{stroke}%
\end{pgfscope}%
\begin{pgfscope}%
\pgfsetrectcap%
\pgfsetmiterjoin%
\pgfsetlinewidth{0.803000pt}%
\definecolor{currentstroke}{rgb}{0.000000,0.000000,0.000000}%
\pgfsetstrokecolor{currentstroke}%
\pgfsetdash{}{0pt}%
\pgfpathmoveto{\pgfqpoint{0.650000in}{0.420000in}}%
\pgfpathlineto{\pgfqpoint{3.038110in}{0.420000in}}%
\pgfusepath{stroke}%
\end{pgfscope}%
\begin{pgfscope}%
\pgfsetrectcap%
\pgfsetmiterjoin%
\pgfsetlinewidth{0.803000pt}%
\definecolor{currentstroke}{rgb}{0.000000,0.000000,0.000000}%
\pgfsetstrokecolor{currentstroke}%
\pgfsetdash{}{0pt}%
\pgfpathmoveto{\pgfqpoint{0.650000in}{2.414488in}}%
\pgfpathlineto{\pgfqpoint{3.038110in}{2.414488in}}%
\pgfusepath{stroke}%
\end{pgfscope}%
\begin{pgfscope}%
\pgfsetrectcap%
\pgfsetroundjoin%
\pgfsetlinewidth{1.003750pt}%
\definecolor{currentstroke}{rgb}{0.870588,0.466667,0.682353}%
\pgfsetstrokecolor{currentstroke}%
\pgfsetdash{}{0pt}%
\pgfpathmoveto{\pgfqpoint{2.127879in}{2.275599in}}%
\pgfpathlineto{\pgfqpoint{2.238990in}{2.275599in}}%
\pgfpathlineto{\pgfqpoint{2.350101in}{2.275599in}}%
\pgfusepath{stroke}%
\end{pgfscope}%
\begin{pgfscope}%
\definecolor{textcolor}{rgb}{0.000000,0.000000,0.000000}%
\pgfsetstrokecolor{textcolor}%
\pgfsetfillcolor{textcolor}%
\pgftext[x=2.438990in,y=2.236710in,left,base]{\color{textcolor}{\rmfamily\fontsize{8.000000}{9.600000}\selectfont\catcode`\^=\active\def^{\ifmmode\sp\else\^{}\fi}\catcode`\%=\active\def%{\%}base}}%
\end{pgfscope}%
\begin{pgfscope}%
\pgfsetrectcap%
\pgfsetroundjoin%
\pgfsetlinewidth{1.003750pt}%
\definecolor{currentstroke}{rgb}{0.498039,0.737255,0.254902}%
\pgfsetstrokecolor{currentstroke}%
\pgfsetdash{}{0pt}%
\pgfpathmoveto{\pgfqpoint{2.127879in}{2.120661in}}%
\pgfpathlineto{\pgfqpoint{2.238990in}{2.120661in}}%
\pgfpathlineto{\pgfqpoint{2.350101in}{2.120661in}}%
\pgfusepath{stroke}%
\end{pgfscope}%
\begin{pgfscope}%
\definecolor{textcolor}{rgb}{0.000000,0.000000,0.000000}%
\pgfsetstrokecolor{textcolor}%
\pgfsetfillcolor{textcolor}%
\pgftext[x=2.438990in,y=2.081772in,left,base]{\color{textcolor}{\rmfamily\fontsize{8.000000}{9.600000}\selectfont\catcode`\^=\active\def^{\ifmmode\sp\else\^{}\fi}\catcode`\%=\active\def%{\%}optimised}}%
\end{pgfscope}%
\end{pgfpicture}%
\makeatother%
\endgroup%

    \label{fig:R100deg30Re1200Uinf}
\end{subfigure}
\\[-15pt]
\caption{Pressure analysis quantities: (a) Clauser parameter $\beta$; (b) pressure distribution, visualised as difference between local pressure and mean pressure in the curved section; (c) normalised free-stream velocity; all quantities plotted against normalised streamwise coordinate}
\label{fig:R100deg30Re1200Prelated}
\end{figure}

In \autoref{fig:R100deg30Re1200Prelated} the pressure-related quantities are visualised.
Firstly, we focus on \autoref{fig:R100deg30Re1200Beta}, in which the Clauser parameter is displayed over the normalised streamwise coordinate.
It is evident that a significant disparity exists between the base configuration and the optimised configuration.
In the flat segment, both configurations demonstrate a constant value of $\beta$ at zero, which has been already described in detail in \autoref{sec:flatTBLVerif}.
This segment corresponds to the straight inlet of the domain.
At $s/\delta_0 \approx 55$ a sharp negative spike in both configurations can be detected.
This spike corresponds to the start of the curvature in the geometry. 
A similar spike in positive direction is visible at $s/\delta_0 \approx 120$, which marks the end of the curve.
Here, it has to be noticed that the peak value for the optimised configuration is approximately two times the value of the base configuration.
Between the two spikes, the two configurations demonstrate significant disparities in the evolution of the Clauser parameter.
In the base configuration, the Clauser parameter remains constant over the entire curvature, exhibiting a slight decrease at a constant rate. 
In contrast, the optimised configuration gives rise to a fluctuating Clauser parameter in the curvature.
over the curved section, three positive bumps with an amplitude of $\Delta\beta \approx 0.2$ can be perceived.
A similar, albeit negative, bump is positioned right before the curved segment starts.

At first glance, this behaviour appears to be counterintuitive. The Clauser parameter represents the optimised quantity, and thus an improvement of $\beta$ in the optimised configuration relative to the base one is anticipated.
To understand this behaviour in a boarder context, we first need to take two more plots into account.
In \autoref{fig:R100deg30Re1200p}, the wall pressure distribution within the domain is depicted.
The streamwise direction is once more normalised with the reference boundary layer thickness. The pressure is visualised as a difference between the local pressure at the wall $p_w$ and the mean pressure at the wall within the curved segment of the domain $\langle p_c \rangle$.
In the first segment, where the domain is straight, both configurations exhibit a constantly decreasing pressure distribution with a similar gradient.
They remain parallel to each other in the straight section, but with the optimised configuration exhibits a lower absolute level.
Once the curvature starts, there is a substantial divergence in the behaviour of the two configurations.
In the base configuration, a steep drop of the pressure level is observed at the start of the curvature. 
Subsequent to this decline, the pressure recovers such that it decreases further with a similar gradient to that observed prior to the drop. 
A similar behaviour in the opposite direction can be observed at the end of the curvature.
A positive pressure jump can be detected in this area.
Afterwards, the pressure decreases with a constant gradient again.
Throughout the entire domain a pressure decrease of $\Delta(p_w-\langle p_c \rangle) \approx 0.09$ is visible.
The optimised configuration shows a different trend in the curved segment.
Once the curvature starts, the pressure trend changes from decreasing to a constant level with a minor positive trend.
Minor fluctuations in the curve can be noticed throughout the entire curved segment.
However, the value remains constant at $\approx 0$ in the curvature. 
At the end of the curved section, a steep rise of the pressure up to $0.04$ is observed.
Near the end of the domain, again, a slightly negative trend adjusts.


For a comprehensive analysis, we also have a look at the free-stream velocity as illustrated in \autoref{fig:R100deg30Re1200Uinf}.
The local free-stream velocity $U(s)$ relative to the free-stream velocity at the reference position $U_0$ is plotted over the normalised streamwise coordinate.
In both configurations, a constant positive trend is evident in the straight part.
As previously outlined, the two configurations diverge once the curvature starts.
The base configuration shows a negative drop at the start of the curvature, followed by an increase with the same gradient as observed before.
At the end of the curvature, a converse behaviour is observed.
Overall, the base configuration shows a positive trend with an absolute difference of $\Delta U/U_0 \approx 0.08$.
In the optimised configuration, conversely, a relatively positive spike occurs immediately prior to the start of the curvature.
Thereafter, the curve ceases to increase in the curved section. 
Instead, a minor negative trend is evident, but with a lower absolute gradient compared to the base configuration.
In addition, fluctuations are observable in the curved segment, similar as described in \autoref{fig:R100deg30Re1200p}.
As the curvature transitions into the straight outlet, a decrease in velocity becomes evident, which recovers until the end of the domain is reached.


The following discussion will address the interpretation of the presented plots. 
In \autoref{fig:R100deg30Re1200Beta} we have observed that the optimised configuration behaves contrary to the expectations.


Instead of an improved trend in terms of constancy at zero and smoothness, it is evident that higher fluctuations occur in comparison to the base configuration.
A closer look at the positions of the bumps suggests indicates a correlation with the positioning of the design points in the optimisation setup.
As described in \autoref{subsec:optSetup}, the optimisation is based on a finite number of design points whose position is adjusted with the objective of reducing the \acrfull{rms} of $\beta$.
The positioning of a \benennung{design point} will modify the size of the cross-sectional area, which impacts the surrounding area as well.
In accordance with the principle of mass conservation, this will result in a modification to the local velocity, consequently leading to an alteration in the momentum balance and, by extension, the local pressure gradient.
This behaviour is enforced if the design points are distributed sparsely over the arc.
This is due to the fact that the cross-sections located between the design points depend only on the spline shape selected by the optimiser.
This may lead to relatively high changes in the cross-sectional area through an unfavourable course of the spline.
A dense positioning of design points would prevent the aforementoined behaviour.



In the context of the case under discussion, this suggests a distribution of design points that is too sparse for maintaining a continuous control of the pressure-gradient results in observable fluctuations in the Clauser parameter.
\citet{morita_applying_2022} increased the number of design points in cases where they anticipated a higher level of necessary shape complexity at the upper wall to accomplish the desired $\beta$-distribution.
The setup studied here shows more extreme curvature characteristics in comparison to the ones already tested with the \acrshort{bo}-\acrshort{gpr}-optimisation. 
Consequently, no reference values are available concerning the adequate number of design points.
To investigate the exact impact of the number of design points on the fluctuations in the Clauser parameter, further investigations will be required.
However, \autoref{fig:R100deg30Re1200p} demonstrates that the optimisation does have an effect on the pressure field.

For a more detailed discussion of the effects in the pressure field, a split-up into \benennung{global} and \benennung{local} optimisation capabilities is useful.
In this context, we use the term \benennung{global} when analysing the general trend of a quantity over the entire curved section.
The \benennung{local} behaviour is defined as the short-range deviations of a quantity from this general trend.
Compared to the base configuration, a shift towards a more uniform pressure distribution on a global level is evident.
Subtracting the mean pressure in the curvature results in an approximately zero pressure level. 
It can thus be concluded that the global pressure level remains constant over the entire curvature.
In contrast, fluctuations are visible on the local level.
Their constitution reveals a striking similarity to the fluctuations observed in the $\beta$ distribution.
In the base configuration, the global pressure trend shows a decrease in the curved section, analogous to in the straight inlet.
Since the pressure declines with a constant slope, this constitution results in a constant pressure gradient.
In conclusion, the optimisation algorithm's capability to enforce a constant pressure differs significantly between the local and the global level when evaluated.
At the local level, the optimisation has been found to be ineffective in enforcing a constant pressure level.
In this instance, the base configuration shows a substantially more pronounced trend of smoothness.
Evaluating the global trend of the pressure, an opposite trend can be observed.
Regarding the global level, the optimisation capability is considerably more advanced.
The findings indicate that the optimisation process leads to a substantial modification of the pressure field, as evidenced by the almost constant pressure level in the curved domain.


\autoref{fig:R100deg30Re1200Uinf} is considered to give a further supporting argument.
The constantly increasing free-stream velocity in the base configuration is consistent with its decreasing pressure.
The decline in pressure provokes a negative pressure-gradient, already introduced as \acrshort{fpg}.
This \acrshort{fpg} accelerates the free-stream, a phenomenon that can be observed in \autoref{fig:R100deg30Re1200Uinf}.
In contrast, the free-stream velocity in the optimised configuration shows a more constant trend on a global level.
This behaviour corresponds to the observed pressure distribution which tends to zero in the curved section.
Therefore, the global pressure gradient is smaller, resulting in a more uniform free-stream velocity.
Nevertheless, at the local level, the free-stream profile of the optimised configuration also demonstrates fluctuations.
Given the observation that all three of the discussed plots exhibit fluctuations of a similar constitution, it can be suggested that these fluctuations occur for the same reason.
To examine this in detail, an analysis of their positions in relation to the positions of the design points will be necessary, but it is out of the scope of this thesis.
In conclusion, the findings indicate that while the optimisation created a $\beta$-distribution with fluctuations within the curved section, the pressure distribution is nonetheless shifted towards the desired, more uniform state on a global level.



\subsection{Curvature comparison for the cases R100deg90 and R500deg18}

The following section compares the cases R100deg90 and R500deg18.
First, the statistics extracted at $\mathrm{Re}_{\theta}=2000$ are analysed and compared with the flat-plate \acrshort{tbl} simulation data from \citet{schlatter_assessment_2010}.
Only the profiles from the optimised simulations are considered here, since the differences between the optimised and the base configuration are negligible in both cases.
The plots comparing the base and optimised configurations within one case are added in the attachment.

\begin{figure}[h!]
\centering
\begin{subfigure}[t]{0.495\textwidth}
    \subcaption{}
    \centering
    %% Creator: Matplotlib, PGF backend
%%
%% To include the figure in your LaTeX document, write
%%   \input{<filename>.pgf}
%%
%% Make sure the required packages are loaded in your preamble
%%   \usepackage{pgf}
%%
%% Also ensure that all the required font packages are loaded; for instance,
%% the lmodern package is sometimes necessary when using math font.
%%   \usepackage{lmodern}
%%
%% Figures using additional raster images can only be included by \input if
%% they are in the same directory as the main LaTeX file. For loading figures
%% from other directories you can use the `import` package
%%   \usepackage{import}
%%
%% and then include the figures with
%%   \import{<path to file>}{<filename>.pgf}
%%
%% Matplotlib used the following preamble
%%   \def\mathdefault#1{#1}
%%   \everymath=\expandafter{\the\everymath\displaystyle}
%%   \IfFileExists{scrextend.sty}{
%%     \usepackage[fontsize=11.000000pt]{scrextend}
%%   }{
%%     \renewcommand{\normalsize}{\fontsize{11.000000}{13.200000}\selectfont}
%%     \normalsize
%%   }
%%   
%%   \makeatletter\@ifpackageloaded{underscore}{}{\usepackage[strings]{underscore}}\makeatother
%%
\begingroup%
\makeatletter%
\begin{pgfpicture}%
\pgfpathrectangle{\pgfpointorigin}{\pgfqpoint{3.118110in}{2.494488in}}%
\pgfusepath{use as bounding box, clip}%
\begin{pgfscope}%
\pgfsetbuttcap%
\pgfsetmiterjoin%
\definecolor{currentfill}{rgb}{1.000000,1.000000,1.000000}%
\pgfsetfillcolor{currentfill}%
\pgfsetlinewidth{0.000000pt}%
\definecolor{currentstroke}{rgb}{1.000000,1.000000,1.000000}%
\pgfsetstrokecolor{currentstroke}%
\pgfsetdash{}{0pt}%
\pgfpathmoveto{\pgfqpoint{0.000000in}{0.000000in}}%
\pgfpathlineto{\pgfqpoint{3.118110in}{0.000000in}}%
\pgfpathlineto{\pgfqpoint{3.118110in}{2.494488in}}%
\pgfpathlineto{\pgfqpoint{0.000000in}{2.494488in}}%
\pgfpathlineto{\pgfqpoint{0.000000in}{0.000000in}}%
\pgfpathclose%
\pgfusepath{fill}%
\end{pgfscope}%
\begin{pgfscope}%
\pgfsetbuttcap%
\pgfsetmiterjoin%
\definecolor{currentfill}{rgb}{1.000000,1.000000,1.000000}%
\pgfsetfillcolor{currentfill}%
\pgfsetlinewidth{0.000000pt}%
\definecolor{currentstroke}{rgb}{0.000000,0.000000,0.000000}%
\pgfsetstrokecolor{currentstroke}%
\pgfsetstrokeopacity{0.000000}%
\pgfsetdash{}{0pt}%
\pgfpathmoveto{\pgfqpoint{0.650000in}{0.420000in}}%
\pgfpathlineto{\pgfqpoint{3.038110in}{0.420000in}}%
\pgfpathlineto{\pgfqpoint{3.038110in}{2.414488in}}%
\pgfpathlineto{\pgfqpoint{0.650000in}{2.414488in}}%
\pgfpathlineto{\pgfqpoint{0.650000in}{0.420000in}}%
\pgfpathclose%
\pgfusepath{fill}%
\end{pgfscope}%
\begin{pgfscope}%
\pgfsetbuttcap%
\pgfsetroundjoin%
\definecolor{currentfill}{rgb}{0.000000,0.000000,0.000000}%
\pgfsetfillcolor{currentfill}%
\pgfsetlinewidth{0.501875pt}%
\definecolor{currentstroke}{rgb}{0.000000,0.000000,0.000000}%
\pgfsetstrokecolor{currentstroke}%
\pgfsetdash{}{0pt}%
\pgfsys@defobject{currentmarker}{\pgfqpoint{0.000000in}{0.000000in}}{\pgfqpoint{0.000000in}{0.041667in}}{%
\pgfpathmoveto{\pgfqpoint{0.000000in}{0.000000in}}%
\pgfpathlineto{\pgfqpoint{0.000000in}{0.041667in}}%
\pgfusepath{stroke,fill}%
}%
\begin{pgfscope}%
\pgfsys@transformshift{0.849578in}{0.420000in}%
\pgfsys@useobject{currentmarker}{}%
\end{pgfscope}%
\end{pgfscope}%
\begin{pgfscope}%
\pgfsetbuttcap%
\pgfsetroundjoin%
\definecolor{currentfill}{rgb}{0.000000,0.000000,0.000000}%
\pgfsetfillcolor{currentfill}%
\pgfsetlinewidth{0.501875pt}%
\definecolor{currentstroke}{rgb}{0.000000,0.000000,0.000000}%
\pgfsetstrokecolor{currentstroke}%
\pgfsetdash{}{0pt}%
\pgfsys@defobject{currentmarker}{\pgfqpoint{0.000000in}{-0.041667in}}{\pgfqpoint{0.000000in}{0.000000in}}{%
\pgfpathmoveto{\pgfqpoint{0.000000in}{0.000000in}}%
\pgfpathlineto{\pgfqpoint{0.000000in}{-0.041667in}}%
\pgfusepath{stroke,fill}%
}%
\begin{pgfscope}%
\pgfsys@transformshift{0.849578in}{2.414488in}%
\pgfsys@useobject{currentmarker}{}%
\end{pgfscope}%
\end{pgfscope}%
\begin{pgfscope}%
\definecolor{textcolor}{rgb}{0.000000,0.000000,0.000000}%
\pgfsetstrokecolor{textcolor}%
\pgfsetfillcolor{textcolor}%
\pgftext[x=0.849578in,y=0.371389in,,top]{\color{textcolor}{\rmfamily\fontsize{11.000000}{13.200000}\selectfont\catcode`\^=\active\def^{\ifmmode\sp\else\^{}\fi}\catcode`\%=\active\def%{\%}$\mathdefault{10^{0}}$}}%
\end{pgfscope}%
\begin{pgfscope}%
\pgfsetbuttcap%
\pgfsetroundjoin%
\definecolor{currentfill}{rgb}{0.000000,0.000000,0.000000}%
\pgfsetfillcolor{currentfill}%
\pgfsetlinewidth{0.501875pt}%
\definecolor{currentstroke}{rgb}{0.000000,0.000000,0.000000}%
\pgfsetstrokecolor{currentstroke}%
\pgfsetdash{}{0pt}%
\pgfsys@defobject{currentmarker}{\pgfqpoint{0.000000in}{0.000000in}}{\pgfqpoint{0.000000in}{0.041667in}}{%
\pgfpathmoveto{\pgfqpoint{0.000000in}{0.000000in}}%
\pgfpathlineto{\pgfqpoint{0.000000in}{0.041667in}}%
\pgfusepath{stroke,fill}%
}%
\begin{pgfscope}%
\pgfsys@transformshift{1.512563in}{0.420000in}%
\pgfsys@useobject{currentmarker}{}%
\end{pgfscope}%
\end{pgfscope}%
\begin{pgfscope}%
\pgfsetbuttcap%
\pgfsetroundjoin%
\definecolor{currentfill}{rgb}{0.000000,0.000000,0.000000}%
\pgfsetfillcolor{currentfill}%
\pgfsetlinewidth{0.501875pt}%
\definecolor{currentstroke}{rgb}{0.000000,0.000000,0.000000}%
\pgfsetstrokecolor{currentstroke}%
\pgfsetdash{}{0pt}%
\pgfsys@defobject{currentmarker}{\pgfqpoint{0.000000in}{-0.041667in}}{\pgfqpoint{0.000000in}{0.000000in}}{%
\pgfpathmoveto{\pgfqpoint{0.000000in}{0.000000in}}%
\pgfpathlineto{\pgfqpoint{0.000000in}{-0.041667in}}%
\pgfusepath{stroke,fill}%
}%
\begin{pgfscope}%
\pgfsys@transformshift{1.512563in}{2.414488in}%
\pgfsys@useobject{currentmarker}{}%
\end{pgfscope}%
\end{pgfscope}%
\begin{pgfscope}%
\definecolor{textcolor}{rgb}{0.000000,0.000000,0.000000}%
\pgfsetstrokecolor{textcolor}%
\pgfsetfillcolor{textcolor}%
\pgftext[x=1.512563in,y=0.371389in,,top]{\color{textcolor}{\rmfamily\fontsize{11.000000}{13.200000}\selectfont\catcode`\^=\active\def^{\ifmmode\sp\else\^{}\fi}\catcode`\%=\active\def%{\%}$\mathdefault{10^{1}}$}}%
\end{pgfscope}%
\begin{pgfscope}%
\pgfsetbuttcap%
\pgfsetroundjoin%
\definecolor{currentfill}{rgb}{0.000000,0.000000,0.000000}%
\pgfsetfillcolor{currentfill}%
\pgfsetlinewidth{0.501875pt}%
\definecolor{currentstroke}{rgb}{0.000000,0.000000,0.000000}%
\pgfsetstrokecolor{currentstroke}%
\pgfsetdash{}{0pt}%
\pgfsys@defobject{currentmarker}{\pgfqpoint{0.000000in}{0.000000in}}{\pgfqpoint{0.000000in}{0.041667in}}{%
\pgfpathmoveto{\pgfqpoint{0.000000in}{0.000000in}}%
\pgfpathlineto{\pgfqpoint{0.000000in}{0.041667in}}%
\pgfusepath{stroke,fill}%
}%
\begin{pgfscope}%
\pgfsys@transformshift{2.175547in}{0.420000in}%
\pgfsys@useobject{currentmarker}{}%
\end{pgfscope}%
\end{pgfscope}%
\begin{pgfscope}%
\pgfsetbuttcap%
\pgfsetroundjoin%
\definecolor{currentfill}{rgb}{0.000000,0.000000,0.000000}%
\pgfsetfillcolor{currentfill}%
\pgfsetlinewidth{0.501875pt}%
\definecolor{currentstroke}{rgb}{0.000000,0.000000,0.000000}%
\pgfsetstrokecolor{currentstroke}%
\pgfsetdash{}{0pt}%
\pgfsys@defobject{currentmarker}{\pgfqpoint{0.000000in}{-0.041667in}}{\pgfqpoint{0.000000in}{0.000000in}}{%
\pgfpathmoveto{\pgfqpoint{0.000000in}{0.000000in}}%
\pgfpathlineto{\pgfqpoint{0.000000in}{-0.041667in}}%
\pgfusepath{stroke,fill}%
}%
\begin{pgfscope}%
\pgfsys@transformshift{2.175547in}{2.414488in}%
\pgfsys@useobject{currentmarker}{}%
\end{pgfscope}%
\end{pgfscope}%
\begin{pgfscope}%
\definecolor{textcolor}{rgb}{0.000000,0.000000,0.000000}%
\pgfsetstrokecolor{textcolor}%
\pgfsetfillcolor{textcolor}%
\pgftext[x=2.175547in,y=0.371389in,,top]{\color{textcolor}{\rmfamily\fontsize{11.000000}{13.200000}\selectfont\catcode`\^=\active\def^{\ifmmode\sp\else\^{}\fi}\catcode`\%=\active\def%{\%}$\mathdefault{10^{2}}$}}%
\end{pgfscope}%
\begin{pgfscope}%
\pgfsetbuttcap%
\pgfsetroundjoin%
\definecolor{currentfill}{rgb}{0.000000,0.000000,0.000000}%
\pgfsetfillcolor{currentfill}%
\pgfsetlinewidth{0.501875pt}%
\definecolor{currentstroke}{rgb}{0.000000,0.000000,0.000000}%
\pgfsetstrokecolor{currentstroke}%
\pgfsetdash{}{0pt}%
\pgfsys@defobject{currentmarker}{\pgfqpoint{0.000000in}{0.000000in}}{\pgfqpoint{0.000000in}{0.041667in}}{%
\pgfpathmoveto{\pgfqpoint{0.000000in}{0.000000in}}%
\pgfpathlineto{\pgfqpoint{0.000000in}{0.041667in}}%
\pgfusepath{stroke,fill}%
}%
\begin{pgfscope}%
\pgfsys@transformshift{2.838532in}{0.420000in}%
\pgfsys@useobject{currentmarker}{}%
\end{pgfscope}%
\end{pgfscope}%
\begin{pgfscope}%
\pgfsetbuttcap%
\pgfsetroundjoin%
\definecolor{currentfill}{rgb}{0.000000,0.000000,0.000000}%
\pgfsetfillcolor{currentfill}%
\pgfsetlinewidth{0.501875pt}%
\definecolor{currentstroke}{rgb}{0.000000,0.000000,0.000000}%
\pgfsetstrokecolor{currentstroke}%
\pgfsetdash{}{0pt}%
\pgfsys@defobject{currentmarker}{\pgfqpoint{0.000000in}{-0.041667in}}{\pgfqpoint{0.000000in}{0.000000in}}{%
\pgfpathmoveto{\pgfqpoint{0.000000in}{0.000000in}}%
\pgfpathlineto{\pgfqpoint{0.000000in}{-0.041667in}}%
\pgfusepath{stroke,fill}%
}%
\begin{pgfscope}%
\pgfsys@transformshift{2.838532in}{2.414488in}%
\pgfsys@useobject{currentmarker}{}%
\end{pgfscope}%
\end{pgfscope}%
\begin{pgfscope}%
\definecolor{textcolor}{rgb}{0.000000,0.000000,0.000000}%
\pgfsetstrokecolor{textcolor}%
\pgfsetfillcolor{textcolor}%
\pgftext[x=2.838532in,y=0.371389in,,top]{\color{textcolor}{\rmfamily\fontsize{11.000000}{13.200000}\selectfont\catcode`\^=\active\def^{\ifmmode\sp\else\^{}\fi}\catcode`\%=\active\def%{\%}$\mathdefault{10^{3}}$}}%
\end{pgfscope}%
\begin{pgfscope}%
\pgfsetbuttcap%
\pgfsetroundjoin%
\definecolor{currentfill}{rgb}{0.000000,0.000000,0.000000}%
\pgfsetfillcolor{currentfill}%
\pgfsetlinewidth{0.401500pt}%
\definecolor{currentstroke}{rgb}{0.000000,0.000000,0.000000}%
\pgfsetstrokecolor{currentstroke}%
\pgfsetdash{}{0pt}%
\pgfsys@defobject{currentmarker}{\pgfqpoint{0.000000in}{0.000000in}}{\pgfqpoint{0.000000in}{0.020833in}}{%
\pgfpathmoveto{\pgfqpoint{0.000000in}{0.000000in}}%
\pgfpathlineto{\pgfqpoint{0.000000in}{0.020833in}}%
\pgfusepath{stroke,fill}%
}%
\begin{pgfscope}%
\pgfsys@transformshift{0.650000in}{0.420000in}%
\pgfsys@useobject{currentmarker}{}%
\end{pgfscope}%
\end{pgfscope}%
\begin{pgfscope}%
\pgfsetbuttcap%
\pgfsetroundjoin%
\definecolor{currentfill}{rgb}{0.000000,0.000000,0.000000}%
\pgfsetfillcolor{currentfill}%
\pgfsetlinewidth{0.401500pt}%
\definecolor{currentstroke}{rgb}{0.000000,0.000000,0.000000}%
\pgfsetstrokecolor{currentstroke}%
\pgfsetdash{}{0pt}%
\pgfsys@defobject{currentmarker}{\pgfqpoint{0.000000in}{-0.020833in}}{\pgfqpoint{0.000000in}{0.000000in}}{%
\pgfpathmoveto{\pgfqpoint{0.000000in}{0.000000in}}%
\pgfpathlineto{\pgfqpoint{0.000000in}{-0.020833in}}%
\pgfusepath{stroke,fill}%
}%
\begin{pgfscope}%
\pgfsys@transformshift{0.650000in}{2.414488in}%
\pgfsys@useobject{currentmarker}{}%
\end{pgfscope}%
\end{pgfscope}%
\begin{pgfscope}%
\pgfsetbuttcap%
\pgfsetroundjoin%
\definecolor{currentfill}{rgb}{0.000000,0.000000,0.000000}%
\pgfsetfillcolor{currentfill}%
\pgfsetlinewidth{0.401500pt}%
\definecolor{currentstroke}{rgb}{0.000000,0.000000,0.000000}%
\pgfsetstrokecolor{currentstroke}%
\pgfsetdash{}{0pt}%
\pgfsys@defobject{currentmarker}{\pgfqpoint{0.000000in}{0.000000in}}{\pgfqpoint{0.000000in}{0.020833in}}{%
\pgfpathmoveto{\pgfqpoint{0.000000in}{0.000000in}}%
\pgfpathlineto{\pgfqpoint{0.000000in}{0.020833in}}%
\pgfusepath{stroke,fill}%
}%
\begin{pgfscope}%
\pgfsys@transformshift{0.702496in}{0.420000in}%
\pgfsys@useobject{currentmarker}{}%
\end{pgfscope}%
\end{pgfscope}%
\begin{pgfscope}%
\pgfsetbuttcap%
\pgfsetroundjoin%
\definecolor{currentfill}{rgb}{0.000000,0.000000,0.000000}%
\pgfsetfillcolor{currentfill}%
\pgfsetlinewidth{0.401500pt}%
\definecolor{currentstroke}{rgb}{0.000000,0.000000,0.000000}%
\pgfsetstrokecolor{currentstroke}%
\pgfsetdash{}{0pt}%
\pgfsys@defobject{currentmarker}{\pgfqpoint{0.000000in}{-0.020833in}}{\pgfqpoint{0.000000in}{0.000000in}}{%
\pgfpathmoveto{\pgfqpoint{0.000000in}{0.000000in}}%
\pgfpathlineto{\pgfqpoint{0.000000in}{-0.020833in}}%
\pgfusepath{stroke,fill}%
}%
\begin{pgfscope}%
\pgfsys@transformshift{0.702496in}{2.414488in}%
\pgfsys@useobject{currentmarker}{}%
\end{pgfscope}%
\end{pgfscope}%
\begin{pgfscope}%
\pgfsetbuttcap%
\pgfsetroundjoin%
\definecolor{currentfill}{rgb}{0.000000,0.000000,0.000000}%
\pgfsetfillcolor{currentfill}%
\pgfsetlinewidth{0.401500pt}%
\definecolor{currentstroke}{rgb}{0.000000,0.000000,0.000000}%
\pgfsetstrokecolor{currentstroke}%
\pgfsetdash{}{0pt}%
\pgfsys@defobject{currentmarker}{\pgfqpoint{0.000000in}{0.000000in}}{\pgfqpoint{0.000000in}{0.020833in}}{%
\pgfpathmoveto{\pgfqpoint{0.000000in}{0.000000in}}%
\pgfpathlineto{\pgfqpoint{0.000000in}{0.020833in}}%
\pgfusepath{stroke,fill}%
}%
\begin{pgfscope}%
\pgfsys@transformshift{0.746881in}{0.420000in}%
\pgfsys@useobject{currentmarker}{}%
\end{pgfscope}%
\end{pgfscope}%
\begin{pgfscope}%
\pgfsetbuttcap%
\pgfsetroundjoin%
\definecolor{currentfill}{rgb}{0.000000,0.000000,0.000000}%
\pgfsetfillcolor{currentfill}%
\pgfsetlinewidth{0.401500pt}%
\definecolor{currentstroke}{rgb}{0.000000,0.000000,0.000000}%
\pgfsetstrokecolor{currentstroke}%
\pgfsetdash{}{0pt}%
\pgfsys@defobject{currentmarker}{\pgfqpoint{0.000000in}{-0.020833in}}{\pgfqpoint{0.000000in}{0.000000in}}{%
\pgfpathmoveto{\pgfqpoint{0.000000in}{0.000000in}}%
\pgfpathlineto{\pgfqpoint{0.000000in}{-0.020833in}}%
\pgfusepath{stroke,fill}%
}%
\begin{pgfscope}%
\pgfsys@transformshift{0.746881in}{2.414488in}%
\pgfsys@useobject{currentmarker}{}%
\end{pgfscope}%
\end{pgfscope}%
\begin{pgfscope}%
\pgfsetbuttcap%
\pgfsetroundjoin%
\definecolor{currentfill}{rgb}{0.000000,0.000000,0.000000}%
\pgfsetfillcolor{currentfill}%
\pgfsetlinewidth{0.401500pt}%
\definecolor{currentstroke}{rgb}{0.000000,0.000000,0.000000}%
\pgfsetstrokecolor{currentstroke}%
\pgfsetdash{}{0pt}%
\pgfsys@defobject{currentmarker}{\pgfqpoint{0.000000in}{0.000000in}}{\pgfqpoint{0.000000in}{0.020833in}}{%
\pgfpathmoveto{\pgfqpoint{0.000000in}{0.000000in}}%
\pgfpathlineto{\pgfqpoint{0.000000in}{0.020833in}}%
\pgfusepath{stroke,fill}%
}%
\begin{pgfscope}%
\pgfsys@transformshift{0.785328in}{0.420000in}%
\pgfsys@useobject{currentmarker}{}%
\end{pgfscope}%
\end{pgfscope}%
\begin{pgfscope}%
\pgfsetbuttcap%
\pgfsetroundjoin%
\definecolor{currentfill}{rgb}{0.000000,0.000000,0.000000}%
\pgfsetfillcolor{currentfill}%
\pgfsetlinewidth{0.401500pt}%
\definecolor{currentstroke}{rgb}{0.000000,0.000000,0.000000}%
\pgfsetstrokecolor{currentstroke}%
\pgfsetdash{}{0pt}%
\pgfsys@defobject{currentmarker}{\pgfqpoint{0.000000in}{-0.020833in}}{\pgfqpoint{0.000000in}{0.000000in}}{%
\pgfpathmoveto{\pgfqpoint{0.000000in}{0.000000in}}%
\pgfpathlineto{\pgfqpoint{0.000000in}{-0.020833in}}%
\pgfusepath{stroke,fill}%
}%
\begin{pgfscope}%
\pgfsys@transformshift{0.785328in}{2.414488in}%
\pgfsys@useobject{currentmarker}{}%
\end{pgfscope}%
\end{pgfscope}%
\begin{pgfscope}%
\pgfsetbuttcap%
\pgfsetroundjoin%
\definecolor{currentfill}{rgb}{0.000000,0.000000,0.000000}%
\pgfsetfillcolor{currentfill}%
\pgfsetlinewidth{0.401500pt}%
\definecolor{currentstroke}{rgb}{0.000000,0.000000,0.000000}%
\pgfsetstrokecolor{currentstroke}%
\pgfsetdash{}{0pt}%
\pgfsys@defobject{currentmarker}{\pgfqpoint{0.000000in}{0.000000in}}{\pgfqpoint{0.000000in}{0.020833in}}{%
\pgfpathmoveto{\pgfqpoint{0.000000in}{0.000000in}}%
\pgfpathlineto{\pgfqpoint{0.000000in}{0.020833in}}%
\pgfusepath{stroke,fill}%
}%
\begin{pgfscope}%
\pgfsys@transformshift{0.819242in}{0.420000in}%
\pgfsys@useobject{currentmarker}{}%
\end{pgfscope}%
\end{pgfscope}%
\begin{pgfscope}%
\pgfsetbuttcap%
\pgfsetroundjoin%
\definecolor{currentfill}{rgb}{0.000000,0.000000,0.000000}%
\pgfsetfillcolor{currentfill}%
\pgfsetlinewidth{0.401500pt}%
\definecolor{currentstroke}{rgb}{0.000000,0.000000,0.000000}%
\pgfsetstrokecolor{currentstroke}%
\pgfsetdash{}{0pt}%
\pgfsys@defobject{currentmarker}{\pgfqpoint{0.000000in}{-0.020833in}}{\pgfqpoint{0.000000in}{0.000000in}}{%
\pgfpathmoveto{\pgfqpoint{0.000000in}{0.000000in}}%
\pgfpathlineto{\pgfqpoint{0.000000in}{-0.020833in}}%
\pgfusepath{stroke,fill}%
}%
\begin{pgfscope}%
\pgfsys@transformshift{0.819242in}{2.414488in}%
\pgfsys@useobject{currentmarker}{}%
\end{pgfscope}%
\end{pgfscope}%
\begin{pgfscope}%
\pgfsetbuttcap%
\pgfsetroundjoin%
\definecolor{currentfill}{rgb}{0.000000,0.000000,0.000000}%
\pgfsetfillcolor{currentfill}%
\pgfsetlinewidth{0.401500pt}%
\definecolor{currentstroke}{rgb}{0.000000,0.000000,0.000000}%
\pgfsetstrokecolor{currentstroke}%
\pgfsetdash{}{0pt}%
\pgfsys@defobject{currentmarker}{\pgfqpoint{0.000000in}{0.000000in}}{\pgfqpoint{0.000000in}{0.020833in}}{%
\pgfpathmoveto{\pgfqpoint{0.000000in}{0.000000in}}%
\pgfpathlineto{\pgfqpoint{0.000000in}{0.020833in}}%
\pgfusepath{stroke,fill}%
}%
\begin{pgfscope}%
\pgfsys@transformshift{1.049156in}{0.420000in}%
\pgfsys@useobject{currentmarker}{}%
\end{pgfscope}%
\end{pgfscope}%
\begin{pgfscope}%
\pgfsetbuttcap%
\pgfsetroundjoin%
\definecolor{currentfill}{rgb}{0.000000,0.000000,0.000000}%
\pgfsetfillcolor{currentfill}%
\pgfsetlinewidth{0.401500pt}%
\definecolor{currentstroke}{rgb}{0.000000,0.000000,0.000000}%
\pgfsetstrokecolor{currentstroke}%
\pgfsetdash{}{0pt}%
\pgfsys@defobject{currentmarker}{\pgfqpoint{0.000000in}{-0.020833in}}{\pgfqpoint{0.000000in}{0.000000in}}{%
\pgfpathmoveto{\pgfqpoint{0.000000in}{0.000000in}}%
\pgfpathlineto{\pgfqpoint{0.000000in}{-0.020833in}}%
\pgfusepath{stroke,fill}%
}%
\begin{pgfscope}%
\pgfsys@transformshift{1.049156in}{2.414488in}%
\pgfsys@useobject{currentmarker}{}%
\end{pgfscope}%
\end{pgfscope}%
\begin{pgfscope}%
\pgfsetbuttcap%
\pgfsetroundjoin%
\definecolor{currentfill}{rgb}{0.000000,0.000000,0.000000}%
\pgfsetfillcolor{currentfill}%
\pgfsetlinewidth{0.401500pt}%
\definecolor{currentstroke}{rgb}{0.000000,0.000000,0.000000}%
\pgfsetstrokecolor{currentstroke}%
\pgfsetdash{}{0pt}%
\pgfsys@defobject{currentmarker}{\pgfqpoint{0.000000in}{0.000000in}}{\pgfqpoint{0.000000in}{0.020833in}}{%
\pgfpathmoveto{\pgfqpoint{0.000000in}{0.000000in}}%
\pgfpathlineto{\pgfqpoint{0.000000in}{0.020833in}}%
\pgfusepath{stroke,fill}%
}%
\begin{pgfscope}%
\pgfsys@transformshift{1.165902in}{0.420000in}%
\pgfsys@useobject{currentmarker}{}%
\end{pgfscope}%
\end{pgfscope}%
\begin{pgfscope}%
\pgfsetbuttcap%
\pgfsetroundjoin%
\definecolor{currentfill}{rgb}{0.000000,0.000000,0.000000}%
\pgfsetfillcolor{currentfill}%
\pgfsetlinewidth{0.401500pt}%
\definecolor{currentstroke}{rgb}{0.000000,0.000000,0.000000}%
\pgfsetstrokecolor{currentstroke}%
\pgfsetdash{}{0pt}%
\pgfsys@defobject{currentmarker}{\pgfqpoint{0.000000in}{-0.020833in}}{\pgfqpoint{0.000000in}{0.000000in}}{%
\pgfpathmoveto{\pgfqpoint{0.000000in}{0.000000in}}%
\pgfpathlineto{\pgfqpoint{0.000000in}{-0.020833in}}%
\pgfusepath{stroke,fill}%
}%
\begin{pgfscope}%
\pgfsys@transformshift{1.165902in}{2.414488in}%
\pgfsys@useobject{currentmarker}{}%
\end{pgfscope}%
\end{pgfscope}%
\begin{pgfscope}%
\pgfsetbuttcap%
\pgfsetroundjoin%
\definecolor{currentfill}{rgb}{0.000000,0.000000,0.000000}%
\pgfsetfillcolor{currentfill}%
\pgfsetlinewidth{0.401500pt}%
\definecolor{currentstroke}{rgb}{0.000000,0.000000,0.000000}%
\pgfsetstrokecolor{currentstroke}%
\pgfsetdash{}{0pt}%
\pgfsys@defobject{currentmarker}{\pgfqpoint{0.000000in}{0.000000in}}{\pgfqpoint{0.000000in}{0.020833in}}{%
\pgfpathmoveto{\pgfqpoint{0.000000in}{0.000000in}}%
\pgfpathlineto{\pgfqpoint{0.000000in}{0.020833in}}%
\pgfusepath{stroke,fill}%
}%
\begin{pgfscope}%
\pgfsys@transformshift{1.248735in}{0.420000in}%
\pgfsys@useobject{currentmarker}{}%
\end{pgfscope}%
\end{pgfscope}%
\begin{pgfscope}%
\pgfsetbuttcap%
\pgfsetroundjoin%
\definecolor{currentfill}{rgb}{0.000000,0.000000,0.000000}%
\pgfsetfillcolor{currentfill}%
\pgfsetlinewidth{0.401500pt}%
\definecolor{currentstroke}{rgb}{0.000000,0.000000,0.000000}%
\pgfsetstrokecolor{currentstroke}%
\pgfsetdash{}{0pt}%
\pgfsys@defobject{currentmarker}{\pgfqpoint{0.000000in}{-0.020833in}}{\pgfqpoint{0.000000in}{0.000000in}}{%
\pgfpathmoveto{\pgfqpoint{0.000000in}{0.000000in}}%
\pgfpathlineto{\pgfqpoint{0.000000in}{-0.020833in}}%
\pgfusepath{stroke,fill}%
}%
\begin{pgfscope}%
\pgfsys@transformshift{1.248735in}{2.414488in}%
\pgfsys@useobject{currentmarker}{}%
\end{pgfscope}%
\end{pgfscope}%
\begin{pgfscope}%
\pgfsetbuttcap%
\pgfsetroundjoin%
\definecolor{currentfill}{rgb}{0.000000,0.000000,0.000000}%
\pgfsetfillcolor{currentfill}%
\pgfsetlinewidth{0.401500pt}%
\definecolor{currentstroke}{rgb}{0.000000,0.000000,0.000000}%
\pgfsetstrokecolor{currentstroke}%
\pgfsetdash{}{0pt}%
\pgfsys@defobject{currentmarker}{\pgfqpoint{0.000000in}{0.000000in}}{\pgfqpoint{0.000000in}{0.020833in}}{%
\pgfpathmoveto{\pgfqpoint{0.000000in}{0.000000in}}%
\pgfpathlineto{\pgfqpoint{0.000000in}{0.020833in}}%
\pgfusepath{stroke,fill}%
}%
\begin{pgfscope}%
\pgfsys@transformshift{1.312985in}{0.420000in}%
\pgfsys@useobject{currentmarker}{}%
\end{pgfscope}%
\end{pgfscope}%
\begin{pgfscope}%
\pgfsetbuttcap%
\pgfsetroundjoin%
\definecolor{currentfill}{rgb}{0.000000,0.000000,0.000000}%
\pgfsetfillcolor{currentfill}%
\pgfsetlinewidth{0.401500pt}%
\definecolor{currentstroke}{rgb}{0.000000,0.000000,0.000000}%
\pgfsetstrokecolor{currentstroke}%
\pgfsetdash{}{0pt}%
\pgfsys@defobject{currentmarker}{\pgfqpoint{0.000000in}{-0.020833in}}{\pgfqpoint{0.000000in}{0.000000in}}{%
\pgfpathmoveto{\pgfqpoint{0.000000in}{0.000000in}}%
\pgfpathlineto{\pgfqpoint{0.000000in}{-0.020833in}}%
\pgfusepath{stroke,fill}%
}%
\begin{pgfscope}%
\pgfsys@transformshift{1.312985in}{2.414488in}%
\pgfsys@useobject{currentmarker}{}%
\end{pgfscope}%
\end{pgfscope}%
\begin{pgfscope}%
\pgfsetbuttcap%
\pgfsetroundjoin%
\definecolor{currentfill}{rgb}{0.000000,0.000000,0.000000}%
\pgfsetfillcolor{currentfill}%
\pgfsetlinewidth{0.401500pt}%
\definecolor{currentstroke}{rgb}{0.000000,0.000000,0.000000}%
\pgfsetstrokecolor{currentstroke}%
\pgfsetdash{}{0pt}%
\pgfsys@defobject{currentmarker}{\pgfqpoint{0.000000in}{0.000000in}}{\pgfqpoint{0.000000in}{0.020833in}}{%
\pgfpathmoveto{\pgfqpoint{0.000000in}{0.000000in}}%
\pgfpathlineto{\pgfqpoint{0.000000in}{0.020833in}}%
\pgfusepath{stroke,fill}%
}%
\begin{pgfscope}%
\pgfsys@transformshift{1.365480in}{0.420000in}%
\pgfsys@useobject{currentmarker}{}%
\end{pgfscope}%
\end{pgfscope}%
\begin{pgfscope}%
\pgfsetbuttcap%
\pgfsetroundjoin%
\definecolor{currentfill}{rgb}{0.000000,0.000000,0.000000}%
\pgfsetfillcolor{currentfill}%
\pgfsetlinewidth{0.401500pt}%
\definecolor{currentstroke}{rgb}{0.000000,0.000000,0.000000}%
\pgfsetstrokecolor{currentstroke}%
\pgfsetdash{}{0pt}%
\pgfsys@defobject{currentmarker}{\pgfqpoint{0.000000in}{-0.020833in}}{\pgfqpoint{0.000000in}{0.000000in}}{%
\pgfpathmoveto{\pgfqpoint{0.000000in}{0.000000in}}%
\pgfpathlineto{\pgfqpoint{0.000000in}{-0.020833in}}%
\pgfusepath{stroke,fill}%
}%
\begin{pgfscope}%
\pgfsys@transformshift{1.365480in}{2.414488in}%
\pgfsys@useobject{currentmarker}{}%
\end{pgfscope}%
\end{pgfscope}%
\begin{pgfscope}%
\pgfsetbuttcap%
\pgfsetroundjoin%
\definecolor{currentfill}{rgb}{0.000000,0.000000,0.000000}%
\pgfsetfillcolor{currentfill}%
\pgfsetlinewidth{0.401500pt}%
\definecolor{currentstroke}{rgb}{0.000000,0.000000,0.000000}%
\pgfsetstrokecolor{currentstroke}%
\pgfsetdash{}{0pt}%
\pgfsys@defobject{currentmarker}{\pgfqpoint{0.000000in}{0.000000in}}{\pgfqpoint{0.000000in}{0.020833in}}{%
\pgfpathmoveto{\pgfqpoint{0.000000in}{0.000000in}}%
\pgfpathlineto{\pgfqpoint{0.000000in}{0.020833in}}%
\pgfusepath{stroke,fill}%
}%
\begin{pgfscope}%
\pgfsys@transformshift{1.409865in}{0.420000in}%
\pgfsys@useobject{currentmarker}{}%
\end{pgfscope}%
\end{pgfscope}%
\begin{pgfscope}%
\pgfsetbuttcap%
\pgfsetroundjoin%
\definecolor{currentfill}{rgb}{0.000000,0.000000,0.000000}%
\pgfsetfillcolor{currentfill}%
\pgfsetlinewidth{0.401500pt}%
\definecolor{currentstroke}{rgb}{0.000000,0.000000,0.000000}%
\pgfsetstrokecolor{currentstroke}%
\pgfsetdash{}{0pt}%
\pgfsys@defobject{currentmarker}{\pgfqpoint{0.000000in}{-0.020833in}}{\pgfqpoint{0.000000in}{0.000000in}}{%
\pgfpathmoveto{\pgfqpoint{0.000000in}{0.000000in}}%
\pgfpathlineto{\pgfqpoint{0.000000in}{-0.020833in}}%
\pgfusepath{stroke,fill}%
}%
\begin{pgfscope}%
\pgfsys@transformshift{1.409865in}{2.414488in}%
\pgfsys@useobject{currentmarker}{}%
\end{pgfscope}%
\end{pgfscope}%
\begin{pgfscope}%
\pgfsetbuttcap%
\pgfsetroundjoin%
\definecolor{currentfill}{rgb}{0.000000,0.000000,0.000000}%
\pgfsetfillcolor{currentfill}%
\pgfsetlinewidth{0.401500pt}%
\definecolor{currentstroke}{rgb}{0.000000,0.000000,0.000000}%
\pgfsetstrokecolor{currentstroke}%
\pgfsetdash{}{0pt}%
\pgfsys@defobject{currentmarker}{\pgfqpoint{0.000000in}{0.000000in}}{\pgfqpoint{0.000000in}{0.020833in}}{%
\pgfpathmoveto{\pgfqpoint{0.000000in}{0.000000in}}%
\pgfpathlineto{\pgfqpoint{0.000000in}{0.020833in}}%
\pgfusepath{stroke,fill}%
}%
\begin{pgfscope}%
\pgfsys@transformshift{1.448313in}{0.420000in}%
\pgfsys@useobject{currentmarker}{}%
\end{pgfscope}%
\end{pgfscope}%
\begin{pgfscope}%
\pgfsetbuttcap%
\pgfsetroundjoin%
\definecolor{currentfill}{rgb}{0.000000,0.000000,0.000000}%
\pgfsetfillcolor{currentfill}%
\pgfsetlinewidth{0.401500pt}%
\definecolor{currentstroke}{rgb}{0.000000,0.000000,0.000000}%
\pgfsetstrokecolor{currentstroke}%
\pgfsetdash{}{0pt}%
\pgfsys@defobject{currentmarker}{\pgfqpoint{0.000000in}{-0.020833in}}{\pgfqpoint{0.000000in}{0.000000in}}{%
\pgfpathmoveto{\pgfqpoint{0.000000in}{0.000000in}}%
\pgfpathlineto{\pgfqpoint{0.000000in}{-0.020833in}}%
\pgfusepath{stroke,fill}%
}%
\begin{pgfscope}%
\pgfsys@transformshift{1.448313in}{2.414488in}%
\pgfsys@useobject{currentmarker}{}%
\end{pgfscope}%
\end{pgfscope}%
\begin{pgfscope}%
\pgfsetbuttcap%
\pgfsetroundjoin%
\definecolor{currentfill}{rgb}{0.000000,0.000000,0.000000}%
\pgfsetfillcolor{currentfill}%
\pgfsetlinewidth{0.401500pt}%
\definecolor{currentstroke}{rgb}{0.000000,0.000000,0.000000}%
\pgfsetstrokecolor{currentstroke}%
\pgfsetdash{}{0pt}%
\pgfsys@defobject{currentmarker}{\pgfqpoint{0.000000in}{0.000000in}}{\pgfqpoint{0.000000in}{0.020833in}}{%
\pgfpathmoveto{\pgfqpoint{0.000000in}{0.000000in}}%
\pgfpathlineto{\pgfqpoint{0.000000in}{0.020833in}}%
\pgfusepath{stroke,fill}%
}%
\begin{pgfscope}%
\pgfsys@transformshift{1.482226in}{0.420000in}%
\pgfsys@useobject{currentmarker}{}%
\end{pgfscope}%
\end{pgfscope}%
\begin{pgfscope}%
\pgfsetbuttcap%
\pgfsetroundjoin%
\definecolor{currentfill}{rgb}{0.000000,0.000000,0.000000}%
\pgfsetfillcolor{currentfill}%
\pgfsetlinewidth{0.401500pt}%
\definecolor{currentstroke}{rgb}{0.000000,0.000000,0.000000}%
\pgfsetstrokecolor{currentstroke}%
\pgfsetdash{}{0pt}%
\pgfsys@defobject{currentmarker}{\pgfqpoint{0.000000in}{-0.020833in}}{\pgfqpoint{0.000000in}{0.000000in}}{%
\pgfpathmoveto{\pgfqpoint{0.000000in}{0.000000in}}%
\pgfpathlineto{\pgfqpoint{0.000000in}{-0.020833in}}%
\pgfusepath{stroke,fill}%
}%
\begin{pgfscope}%
\pgfsys@transformshift{1.482226in}{2.414488in}%
\pgfsys@useobject{currentmarker}{}%
\end{pgfscope}%
\end{pgfscope}%
\begin{pgfscope}%
\pgfsetbuttcap%
\pgfsetroundjoin%
\definecolor{currentfill}{rgb}{0.000000,0.000000,0.000000}%
\pgfsetfillcolor{currentfill}%
\pgfsetlinewidth{0.401500pt}%
\definecolor{currentstroke}{rgb}{0.000000,0.000000,0.000000}%
\pgfsetstrokecolor{currentstroke}%
\pgfsetdash{}{0pt}%
\pgfsys@defobject{currentmarker}{\pgfqpoint{0.000000in}{0.000000in}}{\pgfqpoint{0.000000in}{0.020833in}}{%
\pgfpathmoveto{\pgfqpoint{0.000000in}{0.000000in}}%
\pgfpathlineto{\pgfqpoint{0.000000in}{0.020833in}}%
\pgfusepath{stroke,fill}%
}%
\begin{pgfscope}%
\pgfsys@transformshift{1.712141in}{0.420000in}%
\pgfsys@useobject{currentmarker}{}%
\end{pgfscope}%
\end{pgfscope}%
\begin{pgfscope}%
\pgfsetbuttcap%
\pgfsetroundjoin%
\definecolor{currentfill}{rgb}{0.000000,0.000000,0.000000}%
\pgfsetfillcolor{currentfill}%
\pgfsetlinewidth{0.401500pt}%
\definecolor{currentstroke}{rgb}{0.000000,0.000000,0.000000}%
\pgfsetstrokecolor{currentstroke}%
\pgfsetdash{}{0pt}%
\pgfsys@defobject{currentmarker}{\pgfqpoint{0.000000in}{-0.020833in}}{\pgfqpoint{0.000000in}{0.000000in}}{%
\pgfpathmoveto{\pgfqpoint{0.000000in}{0.000000in}}%
\pgfpathlineto{\pgfqpoint{0.000000in}{-0.020833in}}%
\pgfusepath{stroke,fill}%
}%
\begin{pgfscope}%
\pgfsys@transformshift{1.712141in}{2.414488in}%
\pgfsys@useobject{currentmarker}{}%
\end{pgfscope}%
\end{pgfscope}%
\begin{pgfscope}%
\pgfsetbuttcap%
\pgfsetroundjoin%
\definecolor{currentfill}{rgb}{0.000000,0.000000,0.000000}%
\pgfsetfillcolor{currentfill}%
\pgfsetlinewidth{0.401500pt}%
\definecolor{currentstroke}{rgb}{0.000000,0.000000,0.000000}%
\pgfsetstrokecolor{currentstroke}%
\pgfsetdash{}{0pt}%
\pgfsys@defobject{currentmarker}{\pgfqpoint{0.000000in}{0.000000in}}{\pgfqpoint{0.000000in}{0.020833in}}{%
\pgfpathmoveto{\pgfqpoint{0.000000in}{0.000000in}}%
\pgfpathlineto{\pgfqpoint{0.000000in}{0.020833in}}%
\pgfusepath{stroke,fill}%
}%
\begin{pgfscope}%
\pgfsys@transformshift{1.828887in}{0.420000in}%
\pgfsys@useobject{currentmarker}{}%
\end{pgfscope}%
\end{pgfscope}%
\begin{pgfscope}%
\pgfsetbuttcap%
\pgfsetroundjoin%
\definecolor{currentfill}{rgb}{0.000000,0.000000,0.000000}%
\pgfsetfillcolor{currentfill}%
\pgfsetlinewidth{0.401500pt}%
\definecolor{currentstroke}{rgb}{0.000000,0.000000,0.000000}%
\pgfsetstrokecolor{currentstroke}%
\pgfsetdash{}{0pt}%
\pgfsys@defobject{currentmarker}{\pgfqpoint{0.000000in}{-0.020833in}}{\pgfqpoint{0.000000in}{0.000000in}}{%
\pgfpathmoveto{\pgfqpoint{0.000000in}{0.000000in}}%
\pgfpathlineto{\pgfqpoint{0.000000in}{-0.020833in}}%
\pgfusepath{stroke,fill}%
}%
\begin{pgfscope}%
\pgfsys@transformshift{1.828887in}{2.414488in}%
\pgfsys@useobject{currentmarker}{}%
\end{pgfscope}%
\end{pgfscope}%
\begin{pgfscope}%
\pgfsetbuttcap%
\pgfsetroundjoin%
\definecolor{currentfill}{rgb}{0.000000,0.000000,0.000000}%
\pgfsetfillcolor{currentfill}%
\pgfsetlinewidth{0.401500pt}%
\definecolor{currentstroke}{rgb}{0.000000,0.000000,0.000000}%
\pgfsetstrokecolor{currentstroke}%
\pgfsetdash{}{0pt}%
\pgfsys@defobject{currentmarker}{\pgfqpoint{0.000000in}{0.000000in}}{\pgfqpoint{0.000000in}{0.020833in}}{%
\pgfpathmoveto{\pgfqpoint{0.000000in}{0.000000in}}%
\pgfpathlineto{\pgfqpoint{0.000000in}{0.020833in}}%
\pgfusepath{stroke,fill}%
}%
\begin{pgfscope}%
\pgfsys@transformshift{1.911719in}{0.420000in}%
\pgfsys@useobject{currentmarker}{}%
\end{pgfscope}%
\end{pgfscope}%
\begin{pgfscope}%
\pgfsetbuttcap%
\pgfsetroundjoin%
\definecolor{currentfill}{rgb}{0.000000,0.000000,0.000000}%
\pgfsetfillcolor{currentfill}%
\pgfsetlinewidth{0.401500pt}%
\definecolor{currentstroke}{rgb}{0.000000,0.000000,0.000000}%
\pgfsetstrokecolor{currentstroke}%
\pgfsetdash{}{0pt}%
\pgfsys@defobject{currentmarker}{\pgfqpoint{0.000000in}{-0.020833in}}{\pgfqpoint{0.000000in}{0.000000in}}{%
\pgfpathmoveto{\pgfqpoint{0.000000in}{0.000000in}}%
\pgfpathlineto{\pgfqpoint{0.000000in}{-0.020833in}}%
\pgfusepath{stroke,fill}%
}%
\begin{pgfscope}%
\pgfsys@transformshift{1.911719in}{2.414488in}%
\pgfsys@useobject{currentmarker}{}%
\end{pgfscope}%
\end{pgfscope}%
\begin{pgfscope}%
\pgfsetbuttcap%
\pgfsetroundjoin%
\definecolor{currentfill}{rgb}{0.000000,0.000000,0.000000}%
\pgfsetfillcolor{currentfill}%
\pgfsetlinewidth{0.401500pt}%
\definecolor{currentstroke}{rgb}{0.000000,0.000000,0.000000}%
\pgfsetstrokecolor{currentstroke}%
\pgfsetdash{}{0pt}%
\pgfsys@defobject{currentmarker}{\pgfqpoint{0.000000in}{0.000000in}}{\pgfqpoint{0.000000in}{0.020833in}}{%
\pgfpathmoveto{\pgfqpoint{0.000000in}{0.000000in}}%
\pgfpathlineto{\pgfqpoint{0.000000in}{0.020833in}}%
\pgfusepath{stroke,fill}%
}%
\begin{pgfscope}%
\pgfsys@transformshift{1.975969in}{0.420000in}%
\pgfsys@useobject{currentmarker}{}%
\end{pgfscope}%
\end{pgfscope}%
\begin{pgfscope}%
\pgfsetbuttcap%
\pgfsetroundjoin%
\definecolor{currentfill}{rgb}{0.000000,0.000000,0.000000}%
\pgfsetfillcolor{currentfill}%
\pgfsetlinewidth{0.401500pt}%
\definecolor{currentstroke}{rgb}{0.000000,0.000000,0.000000}%
\pgfsetstrokecolor{currentstroke}%
\pgfsetdash{}{0pt}%
\pgfsys@defobject{currentmarker}{\pgfqpoint{0.000000in}{-0.020833in}}{\pgfqpoint{0.000000in}{0.000000in}}{%
\pgfpathmoveto{\pgfqpoint{0.000000in}{0.000000in}}%
\pgfpathlineto{\pgfqpoint{0.000000in}{-0.020833in}}%
\pgfusepath{stroke,fill}%
}%
\begin{pgfscope}%
\pgfsys@transformshift{1.975969in}{2.414488in}%
\pgfsys@useobject{currentmarker}{}%
\end{pgfscope}%
\end{pgfscope}%
\begin{pgfscope}%
\pgfsetbuttcap%
\pgfsetroundjoin%
\definecolor{currentfill}{rgb}{0.000000,0.000000,0.000000}%
\pgfsetfillcolor{currentfill}%
\pgfsetlinewidth{0.401500pt}%
\definecolor{currentstroke}{rgb}{0.000000,0.000000,0.000000}%
\pgfsetstrokecolor{currentstroke}%
\pgfsetdash{}{0pt}%
\pgfsys@defobject{currentmarker}{\pgfqpoint{0.000000in}{0.000000in}}{\pgfqpoint{0.000000in}{0.020833in}}{%
\pgfpathmoveto{\pgfqpoint{0.000000in}{0.000000in}}%
\pgfpathlineto{\pgfqpoint{0.000000in}{0.020833in}}%
\pgfusepath{stroke,fill}%
}%
\begin{pgfscope}%
\pgfsys@transformshift{2.028465in}{0.420000in}%
\pgfsys@useobject{currentmarker}{}%
\end{pgfscope}%
\end{pgfscope}%
\begin{pgfscope}%
\pgfsetbuttcap%
\pgfsetroundjoin%
\definecolor{currentfill}{rgb}{0.000000,0.000000,0.000000}%
\pgfsetfillcolor{currentfill}%
\pgfsetlinewidth{0.401500pt}%
\definecolor{currentstroke}{rgb}{0.000000,0.000000,0.000000}%
\pgfsetstrokecolor{currentstroke}%
\pgfsetdash{}{0pt}%
\pgfsys@defobject{currentmarker}{\pgfqpoint{0.000000in}{-0.020833in}}{\pgfqpoint{0.000000in}{0.000000in}}{%
\pgfpathmoveto{\pgfqpoint{0.000000in}{0.000000in}}%
\pgfpathlineto{\pgfqpoint{0.000000in}{-0.020833in}}%
\pgfusepath{stroke,fill}%
}%
\begin{pgfscope}%
\pgfsys@transformshift{2.028465in}{2.414488in}%
\pgfsys@useobject{currentmarker}{}%
\end{pgfscope}%
\end{pgfscope}%
\begin{pgfscope}%
\pgfsetbuttcap%
\pgfsetroundjoin%
\definecolor{currentfill}{rgb}{0.000000,0.000000,0.000000}%
\pgfsetfillcolor{currentfill}%
\pgfsetlinewidth{0.401500pt}%
\definecolor{currentstroke}{rgb}{0.000000,0.000000,0.000000}%
\pgfsetstrokecolor{currentstroke}%
\pgfsetdash{}{0pt}%
\pgfsys@defobject{currentmarker}{\pgfqpoint{0.000000in}{0.000000in}}{\pgfqpoint{0.000000in}{0.020833in}}{%
\pgfpathmoveto{\pgfqpoint{0.000000in}{0.000000in}}%
\pgfpathlineto{\pgfqpoint{0.000000in}{0.020833in}}%
\pgfusepath{stroke,fill}%
}%
\begin{pgfscope}%
\pgfsys@transformshift{2.072850in}{0.420000in}%
\pgfsys@useobject{currentmarker}{}%
\end{pgfscope}%
\end{pgfscope}%
\begin{pgfscope}%
\pgfsetbuttcap%
\pgfsetroundjoin%
\definecolor{currentfill}{rgb}{0.000000,0.000000,0.000000}%
\pgfsetfillcolor{currentfill}%
\pgfsetlinewidth{0.401500pt}%
\definecolor{currentstroke}{rgb}{0.000000,0.000000,0.000000}%
\pgfsetstrokecolor{currentstroke}%
\pgfsetdash{}{0pt}%
\pgfsys@defobject{currentmarker}{\pgfqpoint{0.000000in}{-0.020833in}}{\pgfqpoint{0.000000in}{0.000000in}}{%
\pgfpathmoveto{\pgfqpoint{0.000000in}{0.000000in}}%
\pgfpathlineto{\pgfqpoint{0.000000in}{-0.020833in}}%
\pgfusepath{stroke,fill}%
}%
\begin{pgfscope}%
\pgfsys@transformshift{2.072850in}{2.414488in}%
\pgfsys@useobject{currentmarker}{}%
\end{pgfscope}%
\end{pgfscope}%
\begin{pgfscope}%
\pgfsetbuttcap%
\pgfsetroundjoin%
\definecolor{currentfill}{rgb}{0.000000,0.000000,0.000000}%
\pgfsetfillcolor{currentfill}%
\pgfsetlinewidth{0.401500pt}%
\definecolor{currentstroke}{rgb}{0.000000,0.000000,0.000000}%
\pgfsetstrokecolor{currentstroke}%
\pgfsetdash{}{0pt}%
\pgfsys@defobject{currentmarker}{\pgfqpoint{0.000000in}{0.000000in}}{\pgfqpoint{0.000000in}{0.020833in}}{%
\pgfpathmoveto{\pgfqpoint{0.000000in}{0.000000in}}%
\pgfpathlineto{\pgfqpoint{0.000000in}{0.020833in}}%
\pgfusepath{stroke,fill}%
}%
\begin{pgfscope}%
\pgfsys@transformshift{2.111298in}{0.420000in}%
\pgfsys@useobject{currentmarker}{}%
\end{pgfscope}%
\end{pgfscope}%
\begin{pgfscope}%
\pgfsetbuttcap%
\pgfsetroundjoin%
\definecolor{currentfill}{rgb}{0.000000,0.000000,0.000000}%
\pgfsetfillcolor{currentfill}%
\pgfsetlinewidth{0.401500pt}%
\definecolor{currentstroke}{rgb}{0.000000,0.000000,0.000000}%
\pgfsetstrokecolor{currentstroke}%
\pgfsetdash{}{0pt}%
\pgfsys@defobject{currentmarker}{\pgfqpoint{0.000000in}{-0.020833in}}{\pgfqpoint{0.000000in}{0.000000in}}{%
\pgfpathmoveto{\pgfqpoint{0.000000in}{0.000000in}}%
\pgfpathlineto{\pgfqpoint{0.000000in}{-0.020833in}}%
\pgfusepath{stroke,fill}%
}%
\begin{pgfscope}%
\pgfsys@transformshift{2.111298in}{2.414488in}%
\pgfsys@useobject{currentmarker}{}%
\end{pgfscope}%
\end{pgfscope}%
\begin{pgfscope}%
\pgfsetbuttcap%
\pgfsetroundjoin%
\definecolor{currentfill}{rgb}{0.000000,0.000000,0.000000}%
\pgfsetfillcolor{currentfill}%
\pgfsetlinewidth{0.401500pt}%
\definecolor{currentstroke}{rgb}{0.000000,0.000000,0.000000}%
\pgfsetstrokecolor{currentstroke}%
\pgfsetdash{}{0pt}%
\pgfsys@defobject{currentmarker}{\pgfqpoint{0.000000in}{0.000000in}}{\pgfqpoint{0.000000in}{0.020833in}}{%
\pgfpathmoveto{\pgfqpoint{0.000000in}{0.000000in}}%
\pgfpathlineto{\pgfqpoint{0.000000in}{0.020833in}}%
\pgfusepath{stroke,fill}%
}%
\begin{pgfscope}%
\pgfsys@transformshift{2.145211in}{0.420000in}%
\pgfsys@useobject{currentmarker}{}%
\end{pgfscope}%
\end{pgfscope}%
\begin{pgfscope}%
\pgfsetbuttcap%
\pgfsetroundjoin%
\definecolor{currentfill}{rgb}{0.000000,0.000000,0.000000}%
\pgfsetfillcolor{currentfill}%
\pgfsetlinewidth{0.401500pt}%
\definecolor{currentstroke}{rgb}{0.000000,0.000000,0.000000}%
\pgfsetstrokecolor{currentstroke}%
\pgfsetdash{}{0pt}%
\pgfsys@defobject{currentmarker}{\pgfqpoint{0.000000in}{-0.020833in}}{\pgfqpoint{0.000000in}{0.000000in}}{%
\pgfpathmoveto{\pgfqpoint{0.000000in}{0.000000in}}%
\pgfpathlineto{\pgfqpoint{0.000000in}{-0.020833in}}%
\pgfusepath{stroke,fill}%
}%
\begin{pgfscope}%
\pgfsys@transformshift{2.145211in}{2.414488in}%
\pgfsys@useobject{currentmarker}{}%
\end{pgfscope}%
\end{pgfscope}%
\begin{pgfscope}%
\pgfsetbuttcap%
\pgfsetroundjoin%
\definecolor{currentfill}{rgb}{0.000000,0.000000,0.000000}%
\pgfsetfillcolor{currentfill}%
\pgfsetlinewidth{0.401500pt}%
\definecolor{currentstroke}{rgb}{0.000000,0.000000,0.000000}%
\pgfsetstrokecolor{currentstroke}%
\pgfsetdash{}{0pt}%
\pgfsys@defobject{currentmarker}{\pgfqpoint{0.000000in}{0.000000in}}{\pgfqpoint{0.000000in}{0.020833in}}{%
\pgfpathmoveto{\pgfqpoint{0.000000in}{0.000000in}}%
\pgfpathlineto{\pgfqpoint{0.000000in}{0.020833in}}%
\pgfusepath{stroke,fill}%
}%
\begin{pgfscope}%
\pgfsys@transformshift{2.375126in}{0.420000in}%
\pgfsys@useobject{currentmarker}{}%
\end{pgfscope}%
\end{pgfscope}%
\begin{pgfscope}%
\pgfsetbuttcap%
\pgfsetroundjoin%
\definecolor{currentfill}{rgb}{0.000000,0.000000,0.000000}%
\pgfsetfillcolor{currentfill}%
\pgfsetlinewidth{0.401500pt}%
\definecolor{currentstroke}{rgb}{0.000000,0.000000,0.000000}%
\pgfsetstrokecolor{currentstroke}%
\pgfsetdash{}{0pt}%
\pgfsys@defobject{currentmarker}{\pgfqpoint{0.000000in}{-0.020833in}}{\pgfqpoint{0.000000in}{0.000000in}}{%
\pgfpathmoveto{\pgfqpoint{0.000000in}{0.000000in}}%
\pgfpathlineto{\pgfqpoint{0.000000in}{-0.020833in}}%
\pgfusepath{stroke,fill}%
}%
\begin{pgfscope}%
\pgfsys@transformshift{2.375126in}{2.414488in}%
\pgfsys@useobject{currentmarker}{}%
\end{pgfscope}%
\end{pgfscope}%
\begin{pgfscope}%
\pgfsetbuttcap%
\pgfsetroundjoin%
\definecolor{currentfill}{rgb}{0.000000,0.000000,0.000000}%
\pgfsetfillcolor{currentfill}%
\pgfsetlinewidth{0.401500pt}%
\definecolor{currentstroke}{rgb}{0.000000,0.000000,0.000000}%
\pgfsetstrokecolor{currentstroke}%
\pgfsetdash{}{0pt}%
\pgfsys@defobject{currentmarker}{\pgfqpoint{0.000000in}{0.000000in}}{\pgfqpoint{0.000000in}{0.020833in}}{%
\pgfpathmoveto{\pgfqpoint{0.000000in}{0.000000in}}%
\pgfpathlineto{\pgfqpoint{0.000000in}{0.020833in}}%
\pgfusepath{stroke,fill}%
}%
\begin{pgfscope}%
\pgfsys@transformshift{2.491871in}{0.420000in}%
\pgfsys@useobject{currentmarker}{}%
\end{pgfscope}%
\end{pgfscope}%
\begin{pgfscope}%
\pgfsetbuttcap%
\pgfsetroundjoin%
\definecolor{currentfill}{rgb}{0.000000,0.000000,0.000000}%
\pgfsetfillcolor{currentfill}%
\pgfsetlinewidth{0.401500pt}%
\definecolor{currentstroke}{rgb}{0.000000,0.000000,0.000000}%
\pgfsetstrokecolor{currentstroke}%
\pgfsetdash{}{0pt}%
\pgfsys@defobject{currentmarker}{\pgfqpoint{0.000000in}{-0.020833in}}{\pgfqpoint{0.000000in}{0.000000in}}{%
\pgfpathmoveto{\pgfqpoint{0.000000in}{0.000000in}}%
\pgfpathlineto{\pgfqpoint{0.000000in}{-0.020833in}}%
\pgfusepath{stroke,fill}%
}%
\begin{pgfscope}%
\pgfsys@transformshift{2.491871in}{2.414488in}%
\pgfsys@useobject{currentmarker}{}%
\end{pgfscope}%
\end{pgfscope}%
\begin{pgfscope}%
\pgfsetbuttcap%
\pgfsetroundjoin%
\definecolor{currentfill}{rgb}{0.000000,0.000000,0.000000}%
\pgfsetfillcolor{currentfill}%
\pgfsetlinewidth{0.401500pt}%
\definecolor{currentstroke}{rgb}{0.000000,0.000000,0.000000}%
\pgfsetstrokecolor{currentstroke}%
\pgfsetdash{}{0pt}%
\pgfsys@defobject{currentmarker}{\pgfqpoint{0.000000in}{0.000000in}}{\pgfqpoint{0.000000in}{0.020833in}}{%
\pgfpathmoveto{\pgfqpoint{0.000000in}{0.000000in}}%
\pgfpathlineto{\pgfqpoint{0.000000in}{0.020833in}}%
\pgfusepath{stroke,fill}%
}%
\begin{pgfscope}%
\pgfsys@transformshift{2.574704in}{0.420000in}%
\pgfsys@useobject{currentmarker}{}%
\end{pgfscope}%
\end{pgfscope}%
\begin{pgfscope}%
\pgfsetbuttcap%
\pgfsetroundjoin%
\definecolor{currentfill}{rgb}{0.000000,0.000000,0.000000}%
\pgfsetfillcolor{currentfill}%
\pgfsetlinewidth{0.401500pt}%
\definecolor{currentstroke}{rgb}{0.000000,0.000000,0.000000}%
\pgfsetstrokecolor{currentstroke}%
\pgfsetdash{}{0pt}%
\pgfsys@defobject{currentmarker}{\pgfqpoint{0.000000in}{-0.020833in}}{\pgfqpoint{0.000000in}{0.000000in}}{%
\pgfpathmoveto{\pgfqpoint{0.000000in}{0.000000in}}%
\pgfpathlineto{\pgfqpoint{0.000000in}{-0.020833in}}%
\pgfusepath{stroke,fill}%
}%
\begin{pgfscope}%
\pgfsys@transformshift{2.574704in}{2.414488in}%
\pgfsys@useobject{currentmarker}{}%
\end{pgfscope}%
\end{pgfscope}%
\begin{pgfscope}%
\pgfsetbuttcap%
\pgfsetroundjoin%
\definecolor{currentfill}{rgb}{0.000000,0.000000,0.000000}%
\pgfsetfillcolor{currentfill}%
\pgfsetlinewidth{0.401500pt}%
\definecolor{currentstroke}{rgb}{0.000000,0.000000,0.000000}%
\pgfsetstrokecolor{currentstroke}%
\pgfsetdash{}{0pt}%
\pgfsys@defobject{currentmarker}{\pgfqpoint{0.000000in}{0.000000in}}{\pgfqpoint{0.000000in}{0.020833in}}{%
\pgfpathmoveto{\pgfqpoint{0.000000in}{0.000000in}}%
\pgfpathlineto{\pgfqpoint{0.000000in}{0.020833in}}%
\pgfusepath{stroke,fill}%
}%
\begin{pgfscope}%
\pgfsys@transformshift{2.638954in}{0.420000in}%
\pgfsys@useobject{currentmarker}{}%
\end{pgfscope}%
\end{pgfscope}%
\begin{pgfscope}%
\pgfsetbuttcap%
\pgfsetroundjoin%
\definecolor{currentfill}{rgb}{0.000000,0.000000,0.000000}%
\pgfsetfillcolor{currentfill}%
\pgfsetlinewidth{0.401500pt}%
\definecolor{currentstroke}{rgb}{0.000000,0.000000,0.000000}%
\pgfsetstrokecolor{currentstroke}%
\pgfsetdash{}{0pt}%
\pgfsys@defobject{currentmarker}{\pgfqpoint{0.000000in}{-0.020833in}}{\pgfqpoint{0.000000in}{0.000000in}}{%
\pgfpathmoveto{\pgfqpoint{0.000000in}{0.000000in}}%
\pgfpathlineto{\pgfqpoint{0.000000in}{-0.020833in}}%
\pgfusepath{stroke,fill}%
}%
\begin{pgfscope}%
\pgfsys@transformshift{2.638954in}{2.414488in}%
\pgfsys@useobject{currentmarker}{}%
\end{pgfscope}%
\end{pgfscope}%
\begin{pgfscope}%
\pgfsetbuttcap%
\pgfsetroundjoin%
\definecolor{currentfill}{rgb}{0.000000,0.000000,0.000000}%
\pgfsetfillcolor{currentfill}%
\pgfsetlinewidth{0.401500pt}%
\definecolor{currentstroke}{rgb}{0.000000,0.000000,0.000000}%
\pgfsetstrokecolor{currentstroke}%
\pgfsetdash{}{0pt}%
\pgfsys@defobject{currentmarker}{\pgfqpoint{0.000000in}{0.000000in}}{\pgfqpoint{0.000000in}{0.020833in}}{%
\pgfpathmoveto{\pgfqpoint{0.000000in}{0.000000in}}%
\pgfpathlineto{\pgfqpoint{0.000000in}{0.020833in}}%
\pgfusepath{stroke,fill}%
}%
\begin{pgfscope}%
\pgfsys@transformshift{2.691450in}{0.420000in}%
\pgfsys@useobject{currentmarker}{}%
\end{pgfscope}%
\end{pgfscope}%
\begin{pgfscope}%
\pgfsetbuttcap%
\pgfsetroundjoin%
\definecolor{currentfill}{rgb}{0.000000,0.000000,0.000000}%
\pgfsetfillcolor{currentfill}%
\pgfsetlinewidth{0.401500pt}%
\definecolor{currentstroke}{rgb}{0.000000,0.000000,0.000000}%
\pgfsetstrokecolor{currentstroke}%
\pgfsetdash{}{0pt}%
\pgfsys@defobject{currentmarker}{\pgfqpoint{0.000000in}{-0.020833in}}{\pgfqpoint{0.000000in}{0.000000in}}{%
\pgfpathmoveto{\pgfqpoint{0.000000in}{0.000000in}}%
\pgfpathlineto{\pgfqpoint{0.000000in}{-0.020833in}}%
\pgfusepath{stroke,fill}%
}%
\begin{pgfscope}%
\pgfsys@transformshift{2.691450in}{2.414488in}%
\pgfsys@useobject{currentmarker}{}%
\end{pgfscope}%
\end{pgfscope}%
\begin{pgfscope}%
\pgfsetbuttcap%
\pgfsetroundjoin%
\definecolor{currentfill}{rgb}{0.000000,0.000000,0.000000}%
\pgfsetfillcolor{currentfill}%
\pgfsetlinewidth{0.401500pt}%
\definecolor{currentstroke}{rgb}{0.000000,0.000000,0.000000}%
\pgfsetstrokecolor{currentstroke}%
\pgfsetdash{}{0pt}%
\pgfsys@defobject{currentmarker}{\pgfqpoint{0.000000in}{0.000000in}}{\pgfqpoint{0.000000in}{0.020833in}}{%
\pgfpathmoveto{\pgfqpoint{0.000000in}{0.000000in}}%
\pgfpathlineto{\pgfqpoint{0.000000in}{0.020833in}}%
\pgfusepath{stroke,fill}%
}%
\begin{pgfscope}%
\pgfsys@transformshift{2.735834in}{0.420000in}%
\pgfsys@useobject{currentmarker}{}%
\end{pgfscope}%
\end{pgfscope}%
\begin{pgfscope}%
\pgfsetbuttcap%
\pgfsetroundjoin%
\definecolor{currentfill}{rgb}{0.000000,0.000000,0.000000}%
\pgfsetfillcolor{currentfill}%
\pgfsetlinewidth{0.401500pt}%
\definecolor{currentstroke}{rgb}{0.000000,0.000000,0.000000}%
\pgfsetstrokecolor{currentstroke}%
\pgfsetdash{}{0pt}%
\pgfsys@defobject{currentmarker}{\pgfqpoint{0.000000in}{-0.020833in}}{\pgfqpoint{0.000000in}{0.000000in}}{%
\pgfpathmoveto{\pgfqpoint{0.000000in}{0.000000in}}%
\pgfpathlineto{\pgfqpoint{0.000000in}{-0.020833in}}%
\pgfusepath{stroke,fill}%
}%
\begin{pgfscope}%
\pgfsys@transformshift{2.735834in}{2.414488in}%
\pgfsys@useobject{currentmarker}{}%
\end{pgfscope}%
\end{pgfscope}%
\begin{pgfscope}%
\pgfsetbuttcap%
\pgfsetroundjoin%
\definecolor{currentfill}{rgb}{0.000000,0.000000,0.000000}%
\pgfsetfillcolor{currentfill}%
\pgfsetlinewidth{0.401500pt}%
\definecolor{currentstroke}{rgb}{0.000000,0.000000,0.000000}%
\pgfsetstrokecolor{currentstroke}%
\pgfsetdash{}{0pt}%
\pgfsys@defobject{currentmarker}{\pgfqpoint{0.000000in}{0.000000in}}{\pgfqpoint{0.000000in}{0.020833in}}{%
\pgfpathmoveto{\pgfqpoint{0.000000in}{0.000000in}}%
\pgfpathlineto{\pgfqpoint{0.000000in}{0.020833in}}%
\pgfusepath{stroke,fill}%
}%
\begin{pgfscope}%
\pgfsys@transformshift{2.774282in}{0.420000in}%
\pgfsys@useobject{currentmarker}{}%
\end{pgfscope}%
\end{pgfscope}%
\begin{pgfscope}%
\pgfsetbuttcap%
\pgfsetroundjoin%
\definecolor{currentfill}{rgb}{0.000000,0.000000,0.000000}%
\pgfsetfillcolor{currentfill}%
\pgfsetlinewidth{0.401500pt}%
\definecolor{currentstroke}{rgb}{0.000000,0.000000,0.000000}%
\pgfsetstrokecolor{currentstroke}%
\pgfsetdash{}{0pt}%
\pgfsys@defobject{currentmarker}{\pgfqpoint{0.000000in}{-0.020833in}}{\pgfqpoint{0.000000in}{0.000000in}}{%
\pgfpathmoveto{\pgfqpoint{0.000000in}{0.000000in}}%
\pgfpathlineto{\pgfqpoint{0.000000in}{-0.020833in}}%
\pgfusepath{stroke,fill}%
}%
\begin{pgfscope}%
\pgfsys@transformshift{2.774282in}{2.414488in}%
\pgfsys@useobject{currentmarker}{}%
\end{pgfscope}%
\end{pgfscope}%
\begin{pgfscope}%
\pgfsetbuttcap%
\pgfsetroundjoin%
\definecolor{currentfill}{rgb}{0.000000,0.000000,0.000000}%
\pgfsetfillcolor{currentfill}%
\pgfsetlinewidth{0.401500pt}%
\definecolor{currentstroke}{rgb}{0.000000,0.000000,0.000000}%
\pgfsetstrokecolor{currentstroke}%
\pgfsetdash{}{0pt}%
\pgfsys@defobject{currentmarker}{\pgfqpoint{0.000000in}{0.000000in}}{\pgfqpoint{0.000000in}{0.020833in}}{%
\pgfpathmoveto{\pgfqpoint{0.000000in}{0.000000in}}%
\pgfpathlineto{\pgfqpoint{0.000000in}{0.020833in}}%
\pgfusepath{stroke,fill}%
}%
\begin{pgfscope}%
\pgfsys@transformshift{2.808195in}{0.420000in}%
\pgfsys@useobject{currentmarker}{}%
\end{pgfscope}%
\end{pgfscope}%
\begin{pgfscope}%
\pgfsetbuttcap%
\pgfsetroundjoin%
\definecolor{currentfill}{rgb}{0.000000,0.000000,0.000000}%
\pgfsetfillcolor{currentfill}%
\pgfsetlinewidth{0.401500pt}%
\definecolor{currentstroke}{rgb}{0.000000,0.000000,0.000000}%
\pgfsetstrokecolor{currentstroke}%
\pgfsetdash{}{0pt}%
\pgfsys@defobject{currentmarker}{\pgfqpoint{0.000000in}{-0.020833in}}{\pgfqpoint{0.000000in}{0.000000in}}{%
\pgfpathmoveto{\pgfqpoint{0.000000in}{0.000000in}}%
\pgfpathlineto{\pgfqpoint{0.000000in}{-0.020833in}}%
\pgfusepath{stroke,fill}%
}%
\begin{pgfscope}%
\pgfsys@transformshift{2.808195in}{2.414488in}%
\pgfsys@useobject{currentmarker}{}%
\end{pgfscope}%
\end{pgfscope}%
\begin{pgfscope}%
\pgfsetbuttcap%
\pgfsetroundjoin%
\definecolor{currentfill}{rgb}{0.000000,0.000000,0.000000}%
\pgfsetfillcolor{currentfill}%
\pgfsetlinewidth{0.401500pt}%
\definecolor{currentstroke}{rgb}{0.000000,0.000000,0.000000}%
\pgfsetstrokecolor{currentstroke}%
\pgfsetdash{}{0pt}%
\pgfsys@defobject{currentmarker}{\pgfqpoint{0.000000in}{0.000000in}}{\pgfqpoint{0.000000in}{0.020833in}}{%
\pgfpathmoveto{\pgfqpoint{0.000000in}{0.000000in}}%
\pgfpathlineto{\pgfqpoint{0.000000in}{0.020833in}}%
\pgfusepath{stroke,fill}%
}%
\begin{pgfscope}%
\pgfsys@transformshift{3.038110in}{0.420000in}%
\pgfsys@useobject{currentmarker}{}%
\end{pgfscope}%
\end{pgfscope}%
\begin{pgfscope}%
\pgfsetbuttcap%
\pgfsetroundjoin%
\definecolor{currentfill}{rgb}{0.000000,0.000000,0.000000}%
\pgfsetfillcolor{currentfill}%
\pgfsetlinewidth{0.401500pt}%
\definecolor{currentstroke}{rgb}{0.000000,0.000000,0.000000}%
\pgfsetstrokecolor{currentstroke}%
\pgfsetdash{}{0pt}%
\pgfsys@defobject{currentmarker}{\pgfqpoint{0.000000in}{-0.020833in}}{\pgfqpoint{0.000000in}{0.000000in}}{%
\pgfpathmoveto{\pgfqpoint{0.000000in}{0.000000in}}%
\pgfpathlineto{\pgfqpoint{0.000000in}{-0.020833in}}%
\pgfusepath{stroke,fill}%
}%
\begin{pgfscope}%
\pgfsys@transformshift{3.038110in}{2.414488in}%
\pgfsys@useobject{currentmarker}{}%
\end{pgfscope}%
\end{pgfscope}%
\begin{pgfscope}%
\definecolor{textcolor}{rgb}{0.000000,0.000000,0.000000}%
\pgfsetstrokecolor{textcolor}%
\pgfsetfillcolor{textcolor}%
\pgftext[x=1.844055in,y=0.208426in,,top]{\color{textcolor}{\rmfamily\fontsize{12.000000}{14.400000}\selectfont\catcode`\^=\active\def^{\ifmmode\sp\else\^{}\fi}\catcode`\%=\active\def%{\%}$y^+$}}%
\end{pgfscope}%
\begin{pgfscope}%
\pgfsetbuttcap%
\pgfsetroundjoin%
\definecolor{currentfill}{rgb}{0.000000,0.000000,0.000000}%
\pgfsetfillcolor{currentfill}%
\pgfsetlinewidth{0.501875pt}%
\definecolor{currentstroke}{rgb}{0.000000,0.000000,0.000000}%
\pgfsetstrokecolor{currentstroke}%
\pgfsetdash{}{0pt}%
\pgfsys@defobject{currentmarker}{\pgfqpoint{0.000000in}{0.000000in}}{\pgfqpoint{0.041667in}{0.000000in}}{%
\pgfpathmoveto{\pgfqpoint{0.000000in}{0.000000in}}%
\pgfpathlineto{\pgfqpoint{0.041667in}{0.000000in}}%
\pgfusepath{stroke,fill}%
}%
\begin{pgfscope}%
\pgfsys@transformshift{0.650000in}{0.420000in}%
\pgfsys@useobject{currentmarker}{}%
\end{pgfscope}%
\end{pgfscope}%
\begin{pgfscope}%
\pgfsetbuttcap%
\pgfsetroundjoin%
\definecolor{currentfill}{rgb}{0.000000,0.000000,0.000000}%
\pgfsetfillcolor{currentfill}%
\pgfsetlinewidth{0.501875pt}%
\definecolor{currentstroke}{rgb}{0.000000,0.000000,0.000000}%
\pgfsetstrokecolor{currentstroke}%
\pgfsetdash{}{0pt}%
\pgfsys@defobject{currentmarker}{\pgfqpoint{-0.041667in}{0.000000in}}{\pgfqpoint{-0.000000in}{0.000000in}}{%
\pgfpathmoveto{\pgfqpoint{-0.000000in}{0.000000in}}%
\pgfpathlineto{\pgfqpoint{-0.041667in}{0.000000in}}%
\pgfusepath{stroke,fill}%
}%
\begin{pgfscope}%
\pgfsys@transformshift{3.038110in}{0.420000in}%
\pgfsys@useobject{currentmarker}{}%
\end{pgfscope}%
\end{pgfscope}%
\begin{pgfscope}%
\definecolor{textcolor}{rgb}{0.000000,0.000000,0.000000}%
\pgfsetstrokecolor{textcolor}%
\pgfsetfillcolor{textcolor}%
\pgftext[x=0.525347in, y=0.367193in, left, base]{\color{textcolor}{\rmfamily\fontsize{11.000000}{13.200000}\selectfont\catcode`\^=\active\def^{\ifmmode\sp\else\^{}\fi}\catcode`\%=\active\def%{\%}0}}%
\end{pgfscope}%
\begin{pgfscope}%
\pgfsetbuttcap%
\pgfsetroundjoin%
\definecolor{currentfill}{rgb}{0.000000,0.000000,0.000000}%
\pgfsetfillcolor{currentfill}%
\pgfsetlinewidth{0.501875pt}%
\definecolor{currentstroke}{rgb}{0.000000,0.000000,0.000000}%
\pgfsetstrokecolor{currentstroke}%
\pgfsetdash{}{0pt}%
\pgfsys@defobject{currentmarker}{\pgfqpoint{0.000000in}{0.000000in}}{\pgfqpoint{0.041667in}{0.000000in}}{%
\pgfpathmoveto{\pgfqpoint{0.000000in}{0.000000in}}%
\pgfpathlineto{\pgfqpoint{0.041667in}{0.000000in}}%
\pgfusepath{stroke,fill}%
}%
\begin{pgfscope}%
\pgfsys@transformshift{0.650000in}{0.752415in}%
\pgfsys@useobject{currentmarker}{}%
\end{pgfscope}%
\end{pgfscope}%
\begin{pgfscope}%
\pgfsetbuttcap%
\pgfsetroundjoin%
\definecolor{currentfill}{rgb}{0.000000,0.000000,0.000000}%
\pgfsetfillcolor{currentfill}%
\pgfsetlinewidth{0.501875pt}%
\definecolor{currentstroke}{rgb}{0.000000,0.000000,0.000000}%
\pgfsetstrokecolor{currentstroke}%
\pgfsetdash{}{0pt}%
\pgfsys@defobject{currentmarker}{\pgfqpoint{-0.041667in}{0.000000in}}{\pgfqpoint{-0.000000in}{0.000000in}}{%
\pgfpathmoveto{\pgfqpoint{-0.000000in}{0.000000in}}%
\pgfpathlineto{\pgfqpoint{-0.041667in}{0.000000in}}%
\pgfusepath{stroke,fill}%
}%
\begin{pgfscope}%
\pgfsys@transformshift{3.038110in}{0.752415in}%
\pgfsys@useobject{currentmarker}{}%
\end{pgfscope}%
\end{pgfscope}%
\begin{pgfscope}%
\definecolor{textcolor}{rgb}{0.000000,0.000000,0.000000}%
\pgfsetstrokecolor{textcolor}%
\pgfsetfillcolor{textcolor}%
\pgftext[x=0.525347in, y=0.699608in, left, base]{\color{textcolor}{\rmfamily\fontsize{11.000000}{13.200000}\selectfont\catcode`\^=\active\def^{\ifmmode\sp\else\^{}\fi}\catcode`\%=\active\def%{\%}5}}%
\end{pgfscope}%
\begin{pgfscope}%
\pgfsetbuttcap%
\pgfsetroundjoin%
\definecolor{currentfill}{rgb}{0.000000,0.000000,0.000000}%
\pgfsetfillcolor{currentfill}%
\pgfsetlinewidth{0.501875pt}%
\definecolor{currentstroke}{rgb}{0.000000,0.000000,0.000000}%
\pgfsetstrokecolor{currentstroke}%
\pgfsetdash{}{0pt}%
\pgfsys@defobject{currentmarker}{\pgfqpoint{0.000000in}{0.000000in}}{\pgfqpoint{0.041667in}{0.000000in}}{%
\pgfpathmoveto{\pgfqpoint{0.000000in}{0.000000in}}%
\pgfpathlineto{\pgfqpoint{0.041667in}{0.000000in}}%
\pgfusepath{stroke,fill}%
}%
\begin{pgfscope}%
\pgfsys@transformshift{0.650000in}{1.084829in}%
\pgfsys@useobject{currentmarker}{}%
\end{pgfscope}%
\end{pgfscope}%
\begin{pgfscope}%
\pgfsetbuttcap%
\pgfsetroundjoin%
\definecolor{currentfill}{rgb}{0.000000,0.000000,0.000000}%
\pgfsetfillcolor{currentfill}%
\pgfsetlinewidth{0.501875pt}%
\definecolor{currentstroke}{rgb}{0.000000,0.000000,0.000000}%
\pgfsetstrokecolor{currentstroke}%
\pgfsetdash{}{0pt}%
\pgfsys@defobject{currentmarker}{\pgfqpoint{-0.041667in}{0.000000in}}{\pgfqpoint{-0.000000in}{0.000000in}}{%
\pgfpathmoveto{\pgfqpoint{-0.000000in}{0.000000in}}%
\pgfpathlineto{\pgfqpoint{-0.041667in}{0.000000in}}%
\pgfusepath{stroke,fill}%
}%
\begin{pgfscope}%
\pgfsys@transformshift{3.038110in}{1.084829in}%
\pgfsys@useobject{currentmarker}{}%
\end{pgfscope}%
\end{pgfscope}%
\begin{pgfscope}%
\definecolor{textcolor}{rgb}{0.000000,0.000000,0.000000}%
\pgfsetstrokecolor{textcolor}%
\pgfsetfillcolor{textcolor}%
\pgftext[x=0.449305in, y=1.032023in, left, base]{\color{textcolor}{\rmfamily\fontsize{11.000000}{13.200000}\selectfont\catcode`\^=\active\def^{\ifmmode\sp\else\^{}\fi}\catcode`\%=\active\def%{\%}10}}%
\end{pgfscope}%
\begin{pgfscope}%
\pgfsetbuttcap%
\pgfsetroundjoin%
\definecolor{currentfill}{rgb}{0.000000,0.000000,0.000000}%
\pgfsetfillcolor{currentfill}%
\pgfsetlinewidth{0.501875pt}%
\definecolor{currentstroke}{rgb}{0.000000,0.000000,0.000000}%
\pgfsetstrokecolor{currentstroke}%
\pgfsetdash{}{0pt}%
\pgfsys@defobject{currentmarker}{\pgfqpoint{0.000000in}{0.000000in}}{\pgfqpoint{0.041667in}{0.000000in}}{%
\pgfpathmoveto{\pgfqpoint{0.000000in}{0.000000in}}%
\pgfpathlineto{\pgfqpoint{0.041667in}{0.000000in}}%
\pgfusepath{stroke,fill}%
}%
\begin{pgfscope}%
\pgfsys@transformshift{0.650000in}{1.417244in}%
\pgfsys@useobject{currentmarker}{}%
\end{pgfscope}%
\end{pgfscope}%
\begin{pgfscope}%
\pgfsetbuttcap%
\pgfsetroundjoin%
\definecolor{currentfill}{rgb}{0.000000,0.000000,0.000000}%
\pgfsetfillcolor{currentfill}%
\pgfsetlinewidth{0.501875pt}%
\definecolor{currentstroke}{rgb}{0.000000,0.000000,0.000000}%
\pgfsetstrokecolor{currentstroke}%
\pgfsetdash{}{0pt}%
\pgfsys@defobject{currentmarker}{\pgfqpoint{-0.041667in}{0.000000in}}{\pgfqpoint{-0.000000in}{0.000000in}}{%
\pgfpathmoveto{\pgfqpoint{-0.000000in}{0.000000in}}%
\pgfpathlineto{\pgfqpoint{-0.041667in}{0.000000in}}%
\pgfusepath{stroke,fill}%
}%
\begin{pgfscope}%
\pgfsys@transformshift{3.038110in}{1.417244in}%
\pgfsys@useobject{currentmarker}{}%
\end{pgfscope}%
\end{pgfscope}%
\begin{pgfscope}%
\definecolor{textcolor}{rgb}{0.000000,0.000000,0.000000}%
\pgfsetstrokecolor{textcolor}%
\pgfsetfillcolor{textcolor}%
\pgftext[x=0.449305in, y=1.364437in, left, base]{\color{textcolor}{\rmfamily\fontsize{11.000000}{13.200000}\selectfont\catcode`\^=\active\def^{\ifmmode\sp\else\^{}\fi}\catcode`\%=\active\def%{\%}15}}%
\end{pgfscope}%
\begin{pgfscope}%
\pgfsetbuttcap%
\pgfsetroundjoin%
\definecolor{currentfill}{rgb}{0.000000,0.000000,0.000000}%
\pgfsetfillcolor{currentfill}%
\pgfsetlinewidth{0.501875pt}%
\definecolor{currentstroke}{rgb}{0.000000,0.000000,0.000000}%
\pgfsetstrokecolor{currentstroke}%
\pgfsetdash{}{0pt}%
\pgfsys@defobject{currentmarker}{\pgfqpoint{0.000000in}{0.000000in}}{\pgfqpoint{0.041667in}{0.000000in}}{%
\pgfpathmoveto{\pgfqpoint{0.000000in}{0.000000in}}%
\pgfpathlineto{\pgfqpoint{0.041667in}{0.000000in}}%
\pgfusepath{stroke,fill}%
}%
\begin{pgfscope}%
\pgfsys@transformshift{0.650000in}{1.749659in}%
\pgfsys@useobject{currentmarker}{}%
\end{pgfscope}%
\end{pgfscope}%
\begin{pgfscope}%
\pgfsetbuttcap%
\pgfsetroundjoin%
\definecolor{currentfill}{rgb}{0.000000,0.000000,0.000000}%
\pgfsetfillcolor{currentfill}%
\pgfsetlinewidth{0.501875pt}%
\definecolor{currentstroke}{rgb}{0.000000,0.000000,0.000000}%
\pgfsetstrokecolor{currentstroke}%
\pgfsetdash{}{0pt}%
\pgfsys@defobject{currentmarker}{\pgfqpoint{-0.041667in}{0.000000in}}{\pgfqpoint{-0.000000in}{0.000000in}}{%
\pgfpathmoveto{\pgfqpoint{-0.000000in}{0.000000in}}%
\pgfpathlineto{\pgfqpoint{-0.041667in}{0.000000in}}%
\pgfusepath{stroke,fill}%
}%
\begin{pgfscope}%
\pgfsys@transformshift{3.038110in}{1.749659in}%
\pgfsys@useobject{currentmarker}{}%
\end{pgfscope}%
\end{pgfscope}%
\begin{pgfscope}%
\definecolor{textcolor}{rgb}{0.000000,0.000000,0.000000}%
\pgfsetstrokecolor{textcolor}%
\pgfsetfillcolor{textcolor}%
\pgftext[x=0.449305in, y=1.696852in, left, base]{\color{textcolor}{\rmfamily\fontsize{11.000000}{13.200000}\selectfont\catcode`\^=\active\def^{\ifmmode\sp\else\^{}\fi}\catcode`\%=\active\def%{\%}20}}%
\end{pgfscope}%
\begin{pgfscope}%
\pgfsetbuttcap%
\pgfsetroundjoin%
\definecolor{currentfill}{rgb}{0.000000,0.000000,0.000000}%
\pgfsetfillcolor{currentfill}%
\pgfsetlinewidth{0.501875pt}%
\definecolor{currentstroke}{rgb}{0.000000,0.000000,0.000000}%
\pgfsetstrokecolor{currentstroke}%
\pgfsetdash{}{0pt}%
\pgfsys@defobject{currentmarker}{\pgfqpoint{0.000000in}{0.000000in}}{\pgfqpoint{0.041667in}{0.000000in}}{%
\pgfpathmoveto{\pgfqpoint{0.000000in}{0.000000in}}%
\pgfpathlineto{\pgfqpoint{0.041667in}{0.000000in}}%
\pgfusepath{stroke,fill}%
}%
\begin{pgfscope}%
\pgfsys@transformshift{0.650000in}{2.082073in}%
\pgfsys@useobject{currentmarker}{}%
\end{pgfscope}%
\end{pgfscope}%
\begin{pgfscope}%
\pgfsetbuttcap%
\pgfsetroundjoin%
\definecolor{currentfill}{rgb}{0.000000,0.000000,0.000000}%
\pgfsetfillcolor{currentfill}%
\pgfsetlinewidth{0.501875pt}%
\definecolor{currentstroke}{rgb}{0.000000,0.000000,0.000000}%
\pgfsetstrokecolor{currentstroke}%
\pgfsetdash{}{0pt}%
\pgfsys@defobject{currentmarker}{\pgfqpoint{-0.041667in}{0.000000in}}{\pgfqpoint{-0.000000in}{0.000000in}}{%
\pgfpathmoveto{\pgfqpoint{-0.000000in}{0.000000in}}%
\pgfpathlineto{\pgfqpoint{-0.041667in}{0.000000in}}%
\pgfusepath{stroke,fill}%
}%
\begin{pgfscope}%
\pgfsys@transformshift{3.038110in}{2.082073in}%
\pgfsys@useobject{currentmarker}{}%
\end{pgfscope}%
\end{pgfscope}%
\begin{pgfscope}%
\definecolor{textcolor}{rgb}{0.000000,0.000000,0.000000}%
\pgfsetstrokecolor{textcolor}%
\pgfsetfillcolor{textcolor}%
\pgftext[x=0.449305in, y=2.029267in, left, base]{\color{textcolor}{\rmfamily\fontsize{11.000000}{13.200000}\selectfont\catcode`\^=\active\def^{\ifmmode\sp\else\^{}\fi}\catcode`\%=\active\def%{\%}25}}%
\end{pgfscope}%
\begin{pgfscope}%
\pgfsetbuttcap%
\pgfsetroundjoin%
\definecolor{currentfill}{rgb}{0.000000,0.000000,0.000000}%
\pgfsetfillcolor{currentfill}%
\pgfsetlinewidth{0.501875pt}%
\definecolor{currentstroke}{rgb}{0.000000,0.000000,0.000000}%
\pgfsetstrokecolor{currentstroke}%
\pgfsetdash{}{0pt}%
\pgfsys@defobject{currentmarker}{\pgfqpoint{0.000000in}{0.000000in}}{\pgfqpoint{0.041667in}{0.000000in}}{%
\pgfpathmoveto{\pgfqpoint{0.000000in}{0.000000in}}%
\pgfpathlineto{\pgfqpoint{0.041667in}{0.000000in}}%
\pgfusepath{stroke,fill}%
}%
\begin{pgfscope}%
\pgfsys@transformshift{0.650000in}{2.414488in}%
\pgfsys@useobject{currentmarker}{}%
\end{pgfscope}%
\end{pgfscope}%
\begin{pgfscope}%
\pgfsetbuttcap%
\pgfsetroundjoin%
\definecolor{currentfill}{rgb}{0.000000,0.000000,0.000000}%
\pgfsetfillcolor{currentfill}%
\pgfsetlinewidth{0.501875pt}%
\definecolor{currentstroke}{rgb}{0.000000,0.000000,0.000000}%
\pgfsetstrokecolor{currentstroke}%
\pgfsetdash{}{0pt}%
\pgfsys@defobject{currentmarker}{\pgfqpoint{-0.041667in}{0.000000in}}{\pgfqpoint{-0.000000in}{0.000000in}}{%
\pgfpathmoveto{\pgfqpoint{-0.000000in}{0.000000in}}%
\pgfpathlineto{\pgfqpoint{-0.041667in}{0.000000in}}%
\pgfusepath{stroke,fill}%
}%
\begin{pgfscope}%
\pgfsys@transformshift{3.038110in}{2.414488in}%
\pgfsys@useobject{currentmarker}{}%
\end{pgfscope}%
\end{pgfscope}%
\begin{pgfscope}%
\definecolor{textcolor}{rgb}{0.000000,0.000000,0.000000}%
\pgfsetstrokecolor{textcolor}%
\pgfsetfillcolor{textcolor}%
\pgftext[x=0.449305in, y=2.361681in, left, base]{\color{textcolor}{\rmfamily\fontsize{11.000000}{13.200000}\selectfont\catcode`\^=\active\def^{\ifmmode\sp\else\^{}\fi}\catcode`\%=\active\def%{\%}30}}%
\end{pgfscope}%
\begin{pgfscope}%
\pgfsetbuttcap%
\pgfsetroundjoin%
\definecolor{currentfill}{rgb}{0.000000,0.000000,0.000000}%
\pgfsetfillcolor{currentfill}%
\pgfsetlinewidth{0.401500pt}%
\definecolor{currentstroke}{rgb}{0.000000,0.000000,0.000000}%
\pgfsetstrokecolor{currentstroke}%
\pgfsetdash{}{0pt}%
\pgfsys@defobject{currentmarker}{\pgfqpoint{0.000000in}{0.000000in}}{\pgfqpoint{0.020833in}{0.000000in}}{%
\pgfpathmoveto{\pgfqpoint{0.000000in}{0.000000in}}%
\pgfpathlineto{\pgfqpoint{0.020833in}{0.000000in}}%
\pgfusepath{stroke,fill}%
}%
\begin{pgfscope}%
\pgfsys@transformshift{0.650000in}{0.486483in}%
\pgfsys@useobject{currentmarker}{}%
\end{pgfscope}%
\end{pgfscope}%
\begin{pgfscope}%
\pgfsetbuttcap%
\pgfsetroundjoin%
\definecolor{currentfill}{rgb}{0.000000,0.000000,0.000000}%
\pgfsetfillcolor{currentfill}%
\pgfsetlinewidth{0.401500pt}%
\definecolor{currentstroke}{rgb}{0.000000,0.000000,0.000000}%
\pgfsetstrokecolor{currentstroke}%
\pgfsetdash{}{0pt}%
\pgfsys@defobject{currentmarker}{\pgfqpoint{-0.020833in}{0.000000in}}{\pgfqpoint{-0.000000in}{0.000000in}}{%
\pgfpathmoveto{\pgfqpoint{-0.000000in}{0.000000in}}%
\pgfpathlineto{\pgfqpoint{-0.020833in}{0.000000in}}%
\pgfusepath{stroke,fill}%
}%
\begin{pgfscope}%
\pgfsys@transformshift{3.038110in}{0.486483in}%
\pgfsys@useobject{currentmarker}{}%
\end{pgfscope}%
\end{pgfscope}%
\begin{pgfscope}%
\pgfsetbuttcap%
\pgfsetroundjoin%
\definecolor{currentfill}{rgb}{0.000000,0.000000,0.000000}%
\pgfsetfillcolor{currentfill}%
\pgfsetlinewidth{0.401500pt}%
\definecolor{currentstroke}{rgb}{0.000000,0.000000,0.000000}%
\pgfsetstrokecolor{currentstroke}%
\pgfsetdash{}{0pt}%
\pgfsys@defobject{currentmarker}{\pgfqpoint{0.000000in}{0.000000in}}{\pgfqpoint{0.020833in}{0.000000in}}{%
\pgfpathmoveto{\pgfqpoint{0.000000in}{0.000000in}}%
\pgfpathlineto{\pgfqpoint{0.020833in}{0.000000in}}%
\pgfusepath{stroke,fill}%
}%
\begin{pgfscope}%
\pgfsys@transformshift{0.650000in}{0.552966in}%
\pgfsys@useobject{currentmarker}{}%
\end{pgfscope}%
\end{pgfscope}%
\begin{pgfscope}%
\pgfsetbuttcap%
\pgfsetroundjoin%
\definecolor{currentfill}{rgb}{0.000000,0.000000,0.000000}%
\pgfsetfillcolor{currentfill}%
\pgfsetlinewidth{0.401500pt}%
\definecolor{currentstroke}{rgb}{0.000000,0.000000,0.000000}%
\pgfsetstrokecolor{currentstroke}%
\pgfsetdash{}{0pt}%
\pgfsys@defobject{currentmarker}{\pgfqpoint{-0.020833in}{0.000000in}}{\pgfqpoint{-0.000000in}{0.000000in}}{%
\pgfpathmoveto{\pgfqpoint{-0.000000in}{0.000000in}}%
\pgfpathlineto{\pgfqpoint{-0.020833in}{0.000000in}}%
\pgfusepath{stroke,fill}%
}%
\begin{pgfscope}%
\pgfsys@transformshift{3.038110in}{0.552966in}%
\pgfsys@useobject{currentmarker}{}%
\end{pgfscope}%
\end{pgfscope}%
\begin{pgfscope}%
\pgfsetbuttcap%
\pgfsetroundjoin%
\definecolor{currentfill}{rgb}{0.000000,0.000000,0.000000}%
\pgfsetfillcolor{currentfill}%
\pgfsetlinewidth{0.401500pt}%
\definecolor{currentstroke}{rgb}{0.000000,0.000000,0.000000}%
\pgfsetstrokecolor{currentstroke}%
\pgfsetdash{}{0pt}%
\pgfsys@defobject{currentmarker}{\pgfqpoint{0.000000in}{0.000000in}}{\pgfqpoint{0.020833in}{0.000000in}}{%
\pgfpathmoveto{\pgfqpoint{0.000000in}{0.000000in}}%
\pgfpathlineto{\pgfqpoint{0.020833in}{0.000000in}}%
\pgfusepath{stroke,fill}%
}%
\begin{pgfscope}%
\pgfsys@transformshift{0.650000in}{0.619449in}%
\pgfsys@useobject{currentmarker}{}%
\end{pgfscope}%
\end{pgfscope}%
\begin{pgfscope}%
\pgfsetbuttcap%
\pgfsetroundjoin%
\definecolor{currentfill}{rgb}{0.000000,0.000000,0.000000}%
\pgfsetfillcolor{currentfill}%
\pgfsetlinewidth{0.401500pt}%
\definecolor{currentstroke}{rgb}{0.000000,0.000000,0.000000}%
\pgfsetstrokecolor{currentstroke}%
\pgfsetdash{}{0pt}%
\pgfsys@defobject{currentmarker}{\pgfqpoint{-0.020833in}{0.000000in}}{\pgfqpoint{-0.000000in}{0.000000in}}{%
\pgfpathmoveto{\pgfqpoint{-0.000000in}{0.000000in}}%
\pgfpathlineto{\pgfqpoint{-0.020833in}{0.000000in}}%
\pgfusepath{stroke,fill}%
}%
\begin{pgfscope}%
\pgfsys@transformshift{3.038110in}{0.619449in}%
\pgfsys@useobject{currentmarker}{}%
\end{pgfscope}%
\end{pgfscope}%
\begin{pgfscope}%
\pgfsetbuttcap%
\pgfsetroundjoin%
\definecolor{currentfill}{rgb}{0.000000,0.000000,0.000000}%
\pgfsetfillcolor{currentfill}%
\pgfsetlinewidth{0.401500pt}%
\definecolor{currentstroke}{rgb}{0.000000,0.000000,0.000000}%
\pgfsetstrokecolor{currentstroke}%
\pgfsetdash{}{0pt}%
\pgfsys@defobject{currentmarker}{\pgfqpoint{0.000000in}{0.000000in}}{\pgfqpoint{0.020833in}{0.000000in}}{%
\pgfpathmoveto{\pgfqpoint{0.000000in}{0.000000in}}%
\pgfpathlineto{\pgfqpoint{0.020833in}{0.000000in}}%
\pgfusepath{stroke,fill}%
}%
\begin{pgfscope}%
\pgfsys@transformshift{0.650000in}{0.685932in}%
\pgfsys@useobject{currentmarker}{}%
\end{pgfscope}%
\end{pgfscope}%
\begin{pgfscope}%
\pgfsetbuttcap%
\pgfsetroundjoin%
\definecolor{currentfill}{rgb}{0.000000,0.000000,0.000000}%
\pgfsetfillcolor{currentfill}%
\pgfsetlinewidth{0.401500pt}%
\definecolor{currentstroke}{rgb}{0.000000,0.000000,0.000000}%
\pgfsetstrokecolor{currentstroke}%
\pgfsetdash{}{0pt}%
\pgfsys@defobject{currentmarker}{\pgfqpoint{-0.020833in}{0.000000in}}{\pgfqpoint{-0.000000in}{0.000000in}}{%
\pgfpathmoveto{\pgfqpoint{-0.000000in}{0.000000in}}%
\pgfpathlineto{\pgfqpoint{-0.020833in}{0.000000in}}%
\pgfusepath{stroke,fill}%
}%
\begin{pgfscope}%
\pgfsys@transformshift{3.038110in}{0.685932in}%
\pgfsys@useobject{currentmarker}{}%
\end{pgfscope}%
\end{pgfscope}%
\begin{pgfscope}%
\pgfsetbuttcap%
\pgfsetroundjoin%
\definecolor{currentfill}{rgb}{0.000000,0.000000,0.000000}%
\pgfsetfillcolor{currentfill}%
\pgfsetlinewidth{0.401500pt}%
\definecolor{currentstroke}{rgb}{0.000000,0.000000,0.000000}%
\pgfsetstrokecolor{currentstroke}%
\pgfsetdash{}{0pt}%
\pgfsys@defobject{currentmarker}{\pgfqpoint{0.000000in}{0.000000in}}{\pgfqpoint{0.020833in}{0.000000in}}{%
\pgfpathmoveto{\pgfqpoint{0.000000in}{0.000000in}}%
\pgfpathlineto{\pgfqpoint{0.020833in}{0.000000in}}%
\pgfusepath{stroke,fill}%
}%
\begin{pgfscope}%
\pgfsys@transformshift{0.650000in}{0.818898in}%
\pgfsys@useobject{currentmarker}{}%
\end{pgfscope}%
\end{pgfscope}%
\begin{pgfscope}%
\pgfsetbuttcap%
\pgfsetroundjoin%
\definecolor{currentfill}{rgb}{0.000000,0.000000,0.000000}%
\pgfsetfillcolor{currentfill}%
\pgfsetlinewidth{0.401500pt}%
\definecolor{currentstroke}{rgb}{0.000000,0.000000,0.000000}%
\pgfsetstrokecolor{currentstroke}%
\pgfsetdash{}{0pt}%
\pgfsys@defobject{currentmarker}{\pgfqpoint{-0.020833in}{0.000000in}}{\pgfqpoint{-0.000000in}{0.000000in}}{%
\pgfpathmoveto{\pgfqpoint{-0.000000in}{0.000000in}}%
\pgfpathlineto{\pgfqpoint{-0.020833in}{0.000000in}}%
\pgfusepath{stroke,fill}%
}%
\begin{pgfscope}%
\pgfsys@transformshift{3.038110in}{0.818898in}%
\pgfsys@useobject{currentmarker}{}%
\end{pgfscope}%
\end{pgfscope}%
\begin{pgfscope}%
\pgfsetbuttcap%
\pgfsetroundjoin%
\definecolor{currentfill}{rgb}{0.000000,0.000000,0.000000}%
\pgfsetfillcolor{currentfill}%
\pgfsetlinewidth{0.401500pt}%
\definecolor{currentstroke}{rgb}{0.000000,0.000000,0.000000}%
\pgfsetstrokecolor{currentstroke}%
\pgfsetdash{}{0pt}%
\pgfsys@defobject{currentmarker}{\pgfqpoint{0.000000in}{0.000000in}}{\pgfqpoint{0.020833in}{0.000000in}}{%
\pgfpathmoveto{\pgfqpoint{0.000000in}{0.000000in}}%
\pgfpathlineto{\pgfqpoint{0.020833in}{0.000000in}}%
\pgfusepath{stroke,fill}%
}%
\begin{pgfscope}%
\pgfsys@transformshift{0.650000in}{0.885381in}%
\pgfsys@useobject{currentmarker}{}%
\end{pgfscope}%
\end{pgfscope}%
\begin{pgfscope}%
\pgfsetbuttcap%
\pgfsetroundjoin%
\definecolor{currentfill}{rgb}{0.000000,0.000000,0.000000}%
\pgfsetfillcolor{currentfill}%
\pgfsetlinewidth{0.401500pt}%
\definecolor{currentstroke}{rgb}{0.000000,0.000000,0.000000}%
\pgfsetstrokecolor{currentstroke}%
\pgfsetdash{}{0pt}%
\pgfsys@defobject{currentmarker}{\pgfqpoint{-0.020833in}{0.000000in}}{\pgfqpoint{-0.000000in}{0.000000in}}{%
\pgfpathmoveto{\pgfqpoint{-0.000000in}{0.000000in}}%
\pgfpathlineto{\pgfqpoint{-0.020833in}{0.000000in}}%
\pgfusepath{stroke,fill}%
}%
\begin{pgfscope}%
\pgfsys@transformshift{3.038110in}{0.885381in}%
\pgfsys@useobject{currentmarker}{}%
\end{pgfscope}%
\end{pgfscope}%
\begin{pgfscope}%
\pgfsetbuttcap%
\pgfsetroundjoin%
\definecolor{currentfill}{rgb}{0.000000,0.000000,0.000000}%
\pgfsetfillcolor{currentfill}%
\pgfsetlinewidth{0.401500pt}%
\definecolor{currentstroke}{rgb}{0.000000,0.000000,0.000000}%
\pgfsetstrokecolor{currentstroke}%
\pgfsetdash{}{0pt}%
\pgfsys@defobject{currentmarker}{\pgfqpoint{0.000000in}{0.000000in}}{\pgfqpoint{0.020833in}{0.000000in}}{%
\pgfpathmoveto{\pgfqpoint{0.000000in}{0.000000in}}%
\pgfpathlineto{\pgfqpoint{0.020833in}{0.000000in}}%
\pgfusepath{stroke,fill}%
}%
\begin{pgfscope}%
\pgfsys@transformshift{0.650000in}{0.951863in}%
\pgfsys@useobject{currentmarker}{}%
\end{pgfscope}%
\end{pgfscope}%
\begin{pgfscope}%
\pgfsetbuttcap%
\pgfsetroundjoin%
\definecolor{currentfill}{rgb}{0.000000,0.000000,0.000000}%
\pgfsetfillcolor{currentfill}%
\pgfsetlinewidth{0.401500pt}%
\definecolor{currentstroke}{rgb}{0.000000,0.000000,0.000000}%
\pgfsetstrokecolor{currentstroke}%
\pgfsetdash{}{0pt}%
\pgfsys@defobject{currentmarker}{\pgfqpoint{-0.020833in}{0.000000in}}{\pgfqpoint{-0.000000in}{0.000000in}}{%
\pgfpathmoveto{\pgfqpoint{-0.000000in}{0.000000in}}%
\pgfpathlineto{\pgfqpoint{-0.020833in}{0.000000in}}%
\pgfusepath{stroke,fill}%
}%
\begin{pgfscope}%
\pgfsys@transformshift{3.038110in}{0.951863in}%
\pgfsys@useobject{currentmarker}{}%
\end{pgfscope}%
\end{pgfscope}%
\begin{pgfscope}%
\pgfsetbuttcap%
\pgfsetroundjoin%
\definecolor{currentfill}{rgb}{0.000000,0.000000,0.000000}%
\pgfsetfillcolor{currentfill}%
\pgfsetlinewidth{0.401500pt}%
\definecolor{currentstroke}{rgb}{0.000000,0.000000,0.000000}%
\pgfsetstrokecolor{currentstroke}%
\pgfsetdash{}{0pt}%
\pgfsys@defobject{currentmarker}{\pgfqpoint{0.000000in}{0.000000in}}{\pgfqpoint{0.020833in}{0.000000in}}{%
\pgfpathmoveto{\pgfqpoint{0.000000in}{0.000000in}}%
\pgfpathlineto{\pgfqpoint{0.020833in}{0.000000in}}%
\pgfusepath{stroke,fill}%
}%
\begin{pgfscope}%
\pgfsys@transformshift{0.650000in}{1.018346in}%
\pgfsys@useobject{currentmarker}{}%
\end{pgfscope}%
\end{pgfscope}%
\begin{pgfscope}%
\pgfsetbuttcap%
\pgfsetroundjoin%
\definecolor{currentfill}{rgb}{0.000000,0.000000,0.000000}%
\pgfsetfillcolor{currentfill}%
\pgfsetlinewidth{0.401500pt}%
\definecolor{currentstroke}{rgb}{0.000000,0.000000,0.000000}%
\pgfsetstrokecolor{currentstroke}%
\pgfsetdash{}{0pt}%
\pgfsys@defobject{currentmarker}{\pgfqpoint{-0.020833in}{0.000000in}}{\pgfqpoint{-0.000000in}{0.000000in}}{%
\pgfpathmoveto{\pgfqpoint{-0.000000in}{0.000000in}}%
\pgfpathlineto{\pgfqpoint{-0.020833in}{0.000000in}}%
\pgfusepath{stroke,fill}%
}%
\begin{pgfscope}%
\pgfsys@transformshift{3.038110in}{1.018346in}%
\pgfsys@useobject{currentmarker}{}%
\end{pgfscope}%
\end{pgfscope}%
\begin{pgfscope}%
\pgfsetbuttcap%
\pgfsetroundjoin%
\definecolor{currentfill}{rgb}{0.000000,0.000000,0.000000}%
\pgfsetfillcolor{currentfill}%
\pgfsetlinewidth{0.401500pt}%
\definecolor{currentstroke}{rgb}{0.000000,0.000000,0.000000}%
\pgfsetstrokecolor{currentstroke}%
\pgfsetdash{}{0pt}%
\pgfsys@defobject{currentmarker}{\pgfqpoint{0.000000in}{0.000000in}}{\pgfqpoint{0.020833in}{0.000000in}}{%
\pgfpathmoveto{\pgfqpoint{0.000000in}{0.000000in}}%
\pgfpathlineto{\pgfqpoint{0.020833in}{0.000000in}}%
\pgfusepath{stroke,fill}%
}%
\begin{pgfscope}%
\pgfsys@transformshift{0.650000in}{1.151312in}%
\pgfsys@useobject{currentmarker}{}%
\end{pgfscope}%
\end{pgfscope}%
\begin{pgfscope}%
\pgfsetbuttcap%
\pgfsetroundjoin%
\definecolor{currentfill}{rgb}{0.000000,0.000000,0.000000}%
\pgfsetfillcolor{currentfill}%
\pgfsetlinewidth{0.401500pt}%
\definecolor{currentstroke}{rgb}{0.000000,0.000000,0.000000}%
\pgfsetstrokecolor{currentstroke}%
\pgfsetdash{}{0pt}%
\pgfsys@defobject{currentmarker}{\pgfqpoint{-0.020833in}{0.000000in}}{\pgfqpoint{-0.000000in}{0.000000in}}{%
\pgfpathmoveto{\pgfqpoint{-0.000000in}{0.000000in}}%
\pgfpathlineto{\pgfqpoint{-0.020833in}{0.000000in}}%
\pgfusepath{stroke,fill}%
}%
\begin{pgfscope}%
\pgfsys@transformshift{3.038110in}{1.151312in}%
\pgfsys@useobject{currentmarker}{}%
\end{pgfscope}%
\end{pgfscope}%
\begin{pgfscope}%
\pgfsetbuttcap%
\pgfsetroundjoin%
\definecolor{currentfill}{rgb}{0.000000,0.000000,0.000000}%
\pgfsetfillcolor{currentfill}%
\pgfsetlinewidth{0.401500pt}%
\definecolor{currentstroke}{rgb}{0.000000,0.000000,0.000000}%
\pgfsetstrokecolor{currentstroke}%
\pgfsetdash{}{0pt}%
\pgfsys@defobject{currentmarker}{\pgfqpoint{0.000000in}{0.000000in}}{\pgfqpoint{0.020833in}{0.000000in}}{%
\pgfpathmoveto{\pgfqpoint{0.000000in}{0.000000in}}%
\pgfpathlineto{\pgfqpoint{0.020833in}{0.000000in}}%
\pgfusepath{stroke,fill}%
}%
\begin{pgfscope}%
\pgfsys@transformshift{0.650000in}{1.217795in}%
\pgfsys@useobject{currentmarker}{}%
\end{pgfscope}%
\end{pgfscope}%
\begin{pgfscope}%
\pgfsetbuttcap%
\pgfsetroundjoin%
\definecolor{currentfill}{rgb}{0.000000,0.000000,0.000000}%
\pgfsetfillcolor{currentfill}%
\pgfsetlinewidth{0.401500pt}%
\definecolor{currentstroke}{rgb}{0.000000,0.000000,0.000000}%
\pgfsetstrokecolor{currentstroke}%
\pgfsetdash{}{0pt}%
\pgfsys@defobject{currentmarker}{\pgfqpoint{-0.020833in}{0.000000in}}{\pgfqpoint{-0.000000in}{0.000000in}}{%
\pgfpathmoveto{\pgfqpoint{-0.000000in}{0.000000in}}%
\pgfpathlineto{\pgfqpoint{-0.020833in}{0.000000in}}%
\pgfusepath{stroke,fill}%
}%
\begin{pgfscope}%
\pgfsys@transformshift{3.038110in}{1.217795in}%
\pgfsys@useobject{currentmarker}{}%
\end{pgfscope}%
\end{pgfscope}%
\begin{pgfscope}%
\pgfsetbuttcap%
\pgfsetroundjoin%
\definecolor{currentfill}{rgb}{0.000000,0.000000,0.000000}%
\pgfsetfillcolor{currentfill}%
\pgfsetlinewidth{0.401500pt}%
\definecolor{currentstroke}{rgb}{0.000000,0.000000,0.000000}%
\pgfsetstrokecolor{currentstroke}%
\pgfsetdash{}{0pt}%
\pgfsys@defobject{currentmarker}{\pgfqpoint{0.000000in}{0.000000in}}{\pgfqpoint{0.020833in}{0.000000in}}{%
\pgfpathmoveto{\pgfqpoint{0.000000in}{0.000000in}}%
\pgfpathlineto{\pgfqpoint{0.020833in}{0.000000in}}%
\pgfusepath{stroke,fill}%
}%
\begin{pgfscope}%
\pgfsys@transformshift{0.650000in}{1.284278in}%
\pgfsys@useobject{currentmarker}{}%
\end{pgfscope}%
\end{pgfscope}%
\begin{pgfscope}%
\pgfsetbuttcap%
\pgfsetroundjoin%
\definecolor{currentfill}{rgb}{0.000000,0.000000,0.000000}%
\pgfsetfillcolor{currentfill}%
\pgfsetlinewidth{0.401500pt}%
\definecolor{currentstroke}{rgb}{0.000000,0.000000,0.000000}%
\pgfsetstrokecolor{currentstroke}%
\pgfsetdash{}{0pt}%
\pgfsys@defobject{currentmarker}{\pgfqpoint{-0.020833in}{0.000000in}}{\pgfqpoint{-0.000000in}{0.000000in}}{%
\pgfpathmoveto{\pgfqpoint{-0.000000in}{0.000000in}}%
\pgfpathlineto{\pgfqpoint{-0.020833in}{0.000000in}}%
\pgfusepath{stroke,fill}%
}%
\begin{pgfscope}%
\pgfsys@transformshift{3.038110in}{1.284278in}%
\pgfsys@useobject{currentmarker}{}%
\end{pgfscope}%
\end{pgfscope}%
\begin{pgfscope}%
\pgfsetbuttcap%
\pgfsetroundjoin%
\definecolor{currentfill}{rgb}{0.000000,0.000000,0.000000}%
\pgfsetfillcolor{currentfill}%
\pgfsetlinewidth{0.401500pt}%
\definecolor{currentstroke}{rgb}{0.000000,0.000000,0.000000}%
\pgfsetstrokecolor{currentstroke}%
\pgfsetdash{}{0pt}%
\pgfsys@defobject{currentmarker}{\pgfqpoint{0.000000in}{0.000000in}}{\pgfqpoint{0.020833in}{0.000000in}}{%
\pgfpathmoveto{\pgfqpoint{0.000000in}{0.000000in}}%
\pgfpathlineto{\pgfqpoint{0.020833in}{0.000000in}}%
\pgfusepath{stroke,fill}%
}%
\begin{pgfscope}%
\pgfsys@transformshift{0.650000in}{1.350761in}%
\pgfsys@useobject{currentmarker}{}%
\end{pgfscope}%
\end{pgfscope}%
\begin{pgfscope}%
\pgfsetbuttcap%
\pgfsetroundjoin%
\definecolor{currentfill}{rgb}{0.000000,0.000000,0.000000}%
\pgfsetfillcolor{currentfill}%
\pgfsetlinewidth{0.401500pt}%
\definecolor{currentstroke}{rgb}{0.000000,0.000000,0.000000}%
\pgfsetstrokecolor{currentstroke}%
\pgfsetdash{}{0pt}%
\pgfsys@defobject{currentmarker}{\pgfqpoint{-0.020833in}{0.000000in}}{\pgfqpoint{-0.000000in}{0.000000in}}{%
\pgfpathmoveto{\pgfqpoint{-0.000000in}{0.000000in}}%
\pgfpathlineto{\pgfqpoint{-0.020833in}{0.000000in}}%
\pgfusepath{stroke,fill}%
}%
\begin{pgfscope}%
\pgfsys@transformshift{3.038110in}{1.350761in}%
\pgfsys@useobject{currentmarker}{}%
\end{pgfscope}%
\end{pgfscope}%
\begin{pgfscope}%
\pgfsetbuttcap%
\pgfsetroundjoin%
\definecolor{currentfill}{rgb}{0.000000,0.000000,0.000000}%
\pgfsetfillcolor{currentfill}%
\pgfsetlinewidth{0.401500pt}%
\definecolor{currentstroke}{rgb}{0.000000,0.000000,0.000000}%
\pgfsetstrokecolor{currentstroke}%
\pgfsetdash{}{0pt}%
\pgfsys@defobject{currentmarker}{\pgfqpoint{0.000000in}{0.000000in}}{\pgfqpoint{0.020833in}{0.000000in}}{%
\pgfpathmoveto{\pgfqpoint{0.000000in}{0.000000in}}%
\pgfpathlineto{\pgfqpoint{0.020833in}{0.000000in}}%
\pgfusepath{stroke,fill}%
}%
\begin{pgfscope}%
\pgfsys@transformshift{0.650000in}{1.483727in}%
\pgfsys@useobject{currentmarker}{}%
\end{pgfscope}%
\end{pgfscope}%
\begin{pgfscope}%
\pgfsetbuttcap%
\pgfsetroundjoin%
\definecolor{currentfill}{rgb}{0.000000,0.000000,0.000000}%
\pgfsetfillcolor{currentfill}%
\pgfsetlinewidth{0.401500pt}%
\definecolor{currentstroke}{rgb}{0.000000,0.000000,0.000000}%
\pgfsetstrokecolor{currentstroke}%
\pgfsetdash{}{0pt}%
\pgfsys@defobject{currentmarker}{\pgfqpoint{-0.020833in}{0.000000in}}{\pgfqpoint{-0.000000in}{0.000000in}}{%
\pgfpathmoveto{\pgfqpoint{-0.000000in}{0.000000in}}%
\pgfpathlineto{\pgfqpoint{-0.020833in}{0.000000in}}%
\pgfusepath{stroke,fill}%
}%
\begin{pgfscope}%
\pgfsys@transformshift{3.038110in}{1.483727in}%
\pgfsys@useobject{currentmarker}{}%
\end{pgfscope}%
\end{pgfscope}%
\begin{pgfscope}%
\pgfsetbuttcap%
\pgfsetroundjoin%
\definecolor{currentfill}{rgb}{0.000000,0.000000,0.000000}%
\pgfsetfillcolor{currentfill}%
\pgfsetlinewidth{0.401500pt}%
\definecolor{currentstroke}{rgb}{0.000000,0.000000,0.000000}%
\pgfsetstrokecolor{currentstroke}%
\pgfsetdash{}{0pt}%
\pgfsys@defobject{currentmarker}{\pgfqpoint{0.000000in}{0.000000in}}{\pgfqpoint{0.020833in}{0.000000in}}{%
\pgfpathmoveto{\pgfqpoint{0.000000in}{0.000000in}}%
\pgfpathlineto{\pgfqpoint{0.020833in}{0.000000in}}%
\pgfusepath{stroke,fill}%
}%
\begin{pgfscope}%
\pgfsys@transformshift{0.650000in}{1.550210in}%
\pgfsys@useobject{currentmarker}{}%
\end{pgfscope}%
\end{pgfscope}%
\begin{pgfscope}%
\pgfsetbuttcap%
\pgfsetroundjoin%
\definecolor{currentfill}{rgb}{0.000000,0.000000,0.000000}%
\pgfsetfillcolor{currentfill}%
\pgfsetlinewidth{0.401500pt}%
\definecolor{currentstroke}{rgb}{0.000000,0.000000,0.000000}%
\pgfsetstrokecolor{currentstroke}%
\pgfsetdash{}{0pt}%
\pgfsys@defobject{currentmarker}{\pgfqpoint{-0.020833in}{0.000000in}}{\pgfqpoint{-0.000000in}{0.000000in}}{%
\pgfpathmoveto{\pgfqpoint{-0.000000in}{0.000000in}}%
\pgfpathlineto{\pgfqpoint{-0.020833in}{0.000000in}}%
\pgfusepath{stroke,fill}%
}%
\begin{pgfscope}%
\pgfsys@transformshift{3.038110in}{1.550210in}%
\pgfsys@useobject{currentmarker}{}%
\end{pgfscope}%
\end{pgfscope}%
\begin{pgfscope}%
\pgfsetbuttcap%
\pgfsetroundjoin%
\definecolor{currentfill}{rgb}{0.000000,0.000000,0.000000}%
\pgfsetfillcolor{currentfill}%
\pgfsetlinewidth{0.401500pt}%
\definecolor{currentstroke}{rgb}{0.000000,0.000000,0.000000}%
\pgfsetstrokecolor{currentstroke}%
\pgfsetdash{}{0pt}%
\pgfsys@defobject{currentmarker}{\pgfqpoint{0.000000in}{0.000000in}}{\pgfqpoint{0.020833in}{0.000000in}}{%
\pgfpathmoveto{\pgfqpoint{0.000000in}{0.000000in}}%
\pgfpathlineto{\pgfqpoint{0.020833in}{0.000000in}}%
\pgfusepath{stroke,fill}%
}%
\begin{pgfscope}%
\pgfsys@transformshift{0.650000in}{1.616693in}%
\pgfsys@useobject{currentmarker}{}%
\end{pgfscope}%
\end{pgfscope}%
\begin{pgfscope}%
\pgfsetbuttcap%
\pgfsetroundjoin%
\definecolor{currentfill}{rgb}{0.000000,0.000000,0.000000}%
\pgfsetfillcolor{currentfill}%
\pgfsetlinewidth{0.401500pt}%
\definecolor{currentstroke}{rgb}{0.000000,0.000000,0.000000}%
\pgfsetstrokecolor{currentstroke}%
\pgfsetdash{}{0pt}%
\pgfsys@defobject{currentmarker}{\pgfqpoint{-0.020833in}{0.000000in}}{\pgfqpoint{-0.000000in}{0.000000in}}{%
\pgfpathmoveto{\pgfqpoint{-0.000000in}{0.000000in}}%
\pgfpathlineto{\pgfqpoint{-0.020833in}{0.000000in}}%
\pgfusepath{stroke,fill}%
}%
\begin{pgfscope}%
\pgfsys@transformshift{3.038110in}{1.616693in}%
\pgfsys@useobject{currentmarker}{}%
\end{pgfscope}%
\end{pgfscope}%
\begin{pgfscope}%
\pgfsetbuttcap%
\pgfsetroundjoin%
\definecolor{currentfill}{rgb}{0.000000,0.000000,0.000000}%
\pgfsetfillcolor{currentfill}%
\pgfsetlinewidth{0.401500pt}%
\definecolor{currentstroke}{rgb}{0.000000,0.000000,0.000000}%
\pgfsetstrokecolor{currentstroke}%
\pgfsetdash{}{0pt}%
\pgfsys@defobject{currentmarker}{\pgfqpoint{0.000000in}{0.000000in}}{\pgfqpoint{0.020833in}{0.000000in}}{%
\pgfpathmoveto{\pgfqpoint{0.000000in}{0.000000in}}%
\pgfpathlineto{\pgfqpoint{0.020833in}{0.000000in}}%
\pgfusepath{stroke,fill}%
}%
\begin{pgfscope}%
\pgfsys@transformshift{0.650000in}{1.683176in}%
\pgfsys@useobject{currentmarker}{}%
\end{pgfscope}%
\end{pgfscope}%
\begin{pgfscope}%
\pgfsetbuttcap%
\pgfsetroundjoin%
\definecolor{currentfill}{rgb}{0.000000,0.000000,0.000000}%
\pgfsetfillcolor{currentfill}%
\pgfsetlinewidth{0.401500pt}%
\definecolor{currentstroke}{rgb}{0.000000,0.000000,0.000000}%
\pgfsetstrokecolor{currentstroke}%
\pgfsetdash{}{0pt}%
\pgfsys@defobject{currentmarker}{\pgfqpoint{-0.020833in}{0.000000in}}{\pgfqpoint{-0.000000in}{0.000000in}}{%
\pgfpathmoveto{\pgfqpoint{-0.000000in}{0.000000in}}%
\pgfpathlineto{\pgfqpoint{-0.020833in}{0.000000in}}%
\pgfusepath{stroke,fill}%
}%
\begin{pgfscope}%
\pgfsys@transformshift{3.038110in}{1.683176in}%
\pgfsys@useobject{currentmarker}{}%
\end{pgfscope}%
\end{pgfscope}%
\begin{pgfscope}%
\pgfsetbuttcap%
\pgfsetroundjoin%
\definecolor{currentfill}{rgb}{0.000000,0.000000,0.000000}%
\pgfsetfillcolor{currentfill}%
\pgfsetlinewidth{0.401500pt}%
\definecolor{currentstroke}{rgb}{0.000000,0.000000,0.000000}%
\pgfsetstrokecolor{currentstroke}%
\pgfsetdash{}{0pt}%
\pgfsys@defobject{currentmarker}{\pgfqpoint{0.000000in}{0.000000in}}{\pgfqpoint{0.020833in}{0.000000in}}{%
\pgfpathmoveto{\pgfqpoint{0.000000in}{0.000000in}}%
\pgfpathlineto{\pgfqpoint{0.020833in}{0.000000in}}%
\pgfusepath{stroke,fill}%
}%
\begin{pgfscope}%
\pgfsys@transformshift{0.650000in}{1.816142in}%
\pgfsys@useobject{currentmarker}{}%
\end{pgfscope}%
\end{pgfscope}%
\begin{pgfscope}%
\pgfsetbuttcap%
\pgfsetroundjoin%
\definecolor{currentfill}{rgb}{0.000000,0.000000,0.000000}%
\pgfsetfillcolor{currentfill}%
\pgfsetlinewidth{0.401500pt}%
\definecolor{currentstroke}{rgb}{0.000000,0.000000,0.000000}%
\pgfsetstrokecolor{currentstroke}%
\pgfsetdash{}{0pt}%
\pgfsys@defobject{currentmarker}{\pgfqpoint{-0.020833in}{0.000000in}}{\pgfqpoint{-0.000000in}{0.000000in}}{%
\pgfpathmoveto{\pgfqpoint{-0.000000in}{0.000000in}}%
\pgfpathlineto{\pgfqpoint{-0.020833in}{0.000000in}}%
\pgfusepath{stroke,fill}%
}%
\begin{pgfscope}%
\pgfsys@transformshift{3.038110in}{1.816142in}%
\pgfsys@useobject{currentmarker}{}%
\end{pgfscope}%
\end{pgfscope}%
\begin{pgfscope}%
\pgfsetbuttcap%
\pgfsetroundjoin%
\definecolor{currentfill}{rgb}{0.000000,0.000000,0.000000}%
\pgfsetfillcolor{currentfill}%
\pgfsetlinewidth{0.401500pt}%
\definecolor{currentstroke}{rgb}{0.000000,0.000000,0.000000}%
\pgfsetstrokecolor{currentstroke}%
\pgfsetdash{}{0pt}%
\pgfsys@defobject{currentmarker}{\pgfqpoint{0.000000in}{0.000000in}}{\pgfqpoint{0.020833in}{0.000000in}}{%
\pgfpathmoveto{\pgfqpoint{0.000000in}{0.000000in}}%
\pgfpathlineto{\pgfqpoint{0.020833in}{0.000000in}}%
\pgfusepath{stroke,fill}%
}%
\begin{pgfscope}%
\pgfsys@transformshift{0.650000in}{1.882625in}%
\pgfsys@useobject{currentmarker}{}%
\end{pgfscope}%
\end{pgfscope}%
\begin{pgfscope}%
\pgfsetbuttcap%
\pgfsetroundjoin%
\definecolor{currentfill}{rgb}{0.000000,0.000000,0.000000}%
\pgfsetfillcolor{currentfill}%
\pgfsetlinewidth{0.401500pt}%
\definecolor{currentstroke}{rgb}{0.000000,0.000000,0.000000}%
\pgfsetstrokecolor{currentstroke}%
\pgfsetdash{}{0pt}%
\pgfsys@defobject{currentmarker}{\pgfqpoint{-0.020833in}{0.000000in}}{\pgfqpoint{-0.000000in}{0.000000in}}{%
\pgfpathmoveto{\pgfqpoint{-0.000000in}{0.000000in}}%
\pgfpathlineto{\pgfqpoint{-0.020833in}{0.000000in}}%
\pgfusepath{stroke,fill}%
}%
\begin{pgfscope}%
\pgfsys@transformshift{3.038110in}{1.882625in}%
\pgfsys@useobject{currentmarker}{}%
\end{pgfscope}%
\end{pgfscope}%
\begin{pgfscope}%
\pgfsetbuttcap%
\pgfsetroundjoin%
\definecolor{currentfill}{rgb}{0.000000,0.000000,0.000000}%
\pgfsetfillcolor{currentfill}%
\pgfsetlinewidth{0.401500pt}%
\definecolor{currentstroke}{rgb}{0.000000,0.000000,0.000000}%
\pgfsetstrokecolor{currentstroke}%
\pgfsetdash{}{0pt}%
\pgfsys@defobject{currentmarker}{\pgfqpoint{0.000000in}{0.000000in}}{\pgfqpoint{0.020833in}{0.000000in}}{%
\pgfpathmoveto{\pgfqpoint{0.000000in}{0.000000in}}%
\pgfpathlineto{\pgfqpoint{0.020833in}{0.000000in}}%
\pgfusepath{stroke,fill}%
}%
\begin{pgfscope}%
\pgfsys@transformshift{0.650000in}{1.949108in}%
\pgfsys@useobject{currentmarker}{}%
\end{pgfscope}%
\end{pgfscope}%
\begin{pgfscope}%
\pgfsetbuttcap%
\pgfsetroundjoin%
\definecolor{currentfill}{rgb}{0.000000,0.000000,0.000000}%
\pgfsetfillcolor{currentfill}%
\pgfsetlinewidth{0.401500pt}%
\definecolor{currentstroke}{rgb}{0.000000,0.000000,0.000000}%
\pgfsetstrokecolor{currentstroke}%
\pgfsetdash{}{0pt}%
\pgfsys@defobject{currentmarker}{\pgfqpoint{-0.020833in}{0.000000in}}{\pgfqpoint{-0.000000in}{0.000000in}}{%
\pgfpathmoveto{\pgfqpoint{-0.000000in}{0.000000in}}%
\pgfpathlineto{\pgfqpoint{-0.020833in}{0.000000in}}%
\pgfusepath{stroke,fill}%
}%
\begin{pgfscope}%
\pgfsys@transformshift{3.038110in}{1.949108in}%
\pgfsys@useobject{currentmarker}{}%
\end{pgfscope}%
\end{pgfscope}%
\begin{pgfscope}%
\pgfsetbuttcap%
\pgfsetroundjoin%
\definecolor{currentfill}{rgb}{0.000000,0.000000,0.000000}%
\pgfsetfillcolor{currentfill}%
\pgfsetlinewidth{0.401500pt}%
\definecolor{currentstroke}{rgb}{0.000000,0.000000,0.000000}%
\pgfsetstrokecolor{currentstroke}%
\pgfsetdash{}{0pt}%
\pgfsys@defobject{currentmarker}{\pgfqpoint{0.000000in}{0.000000in}}{\pgfqpoint{0.020833in}{0.000000in}}{%
\pgfpathmoveto{\pgfqpoint{0.000000in}{0.000000in}}%
\pgfpathlineto{\pgfqpoint{0.020833in}{0.000000in}}%
\pgfusepath{stroke,fill}%
}%
\begin{pgfscope}%
\pgfsys@transformshift{0.650000in}{2.015590in}%
\pgfsys@useobject{currentmarker}{}%
\end{pgfscope}%
\end{pgfscope}%
\begin{pgfscope}%
\pgfsetbuttcap%
\pgfsetroundjoin%
\definecolor{currentfill}{rgb}{0.000000,0.000000,0.000000}%
\pgfsetfillcolor{currentfill}%
\pgfsetlinewidth{0.401500pt}%
\definecolor{currentstroke}{rgb}{0.000000,0.000000,0.000000}%
\pgfsetstrokecolor{currentstroke}%
\pgfsetdash{}{0pt}%
\pgfsys@defobject{currentmarker}{\pgfqpoint{-0.020833in}{0.000000in}}{\pgfqpoint{-0.000000in}{0.000000in}}{%
\pgfpathmoveto{\pgfqpoint{-0.000000in}{0.000000in}}%
\pgfpathlineto{\pgfqpoint{-0.020833in}{0.000000in}}%
\pgfusepath{stroke,fill}%
}%
\begin{pgfscope}%
\pgfsys@transformshift{3.038110in}{2.015590in}%
\pgfsys@useobject{currentmarker}{}%
\end{pgfscope}%
\end{pgfscope}%
\begin{pgfscope}%
\pgfsetbuttcap%
\pgfsetroundjoin%
\definecolor{currentfill}{rgb}{0.000000,0.000000,0.000000}%
\pgfsetfillcolor{currentfill}%
\pgfsetlinewidth{0.401500pt}%
\definecolor{currentstroke}{rgb}{0.000000,0.000000,0.000000}%
\pgfsetstrokecolor{currentstroke}%
\pgfsetdash{}{0pt}%
\pgfsys@defobject{currentmarker}{\pgfqpoint{0.000000in}{0.000000in}}{\pgfqpoint{0.020833in}{0.000000in}}{%
\pgfpathmoveto{\pgfqpoint{0.000000in}{0.000000in}}%
\pgfpathlineto{\pgfqpoint{0.020833in}{0.000000in}}%
\pgfusepath{stroke,fill}%
}%
\begin{pgfscope}%
\pgfsys@transformshift{0.650000in}{2.148556in}%
\pgfsys@useobject{currentmarker}{}%
\end{pgfscope}%
\end{pgfscope}%
\begin{pgfscope}%
\pgfsetbuttcap%
\pgfsetroundjoin%
\definecolor{currentfill}{rgb}{0.000000,0.000000,0.000000}%
\pgfsetfillcolor{currentfill}%
\pgfsetlinewidth{0.401500pt}%
\definecolor{currentstroke}{rgb}{0.000000,0.000000,0.000000}%
\pgfsetstrokecolor{currentstroke}%
\pgfsetdash{}{0pt}%
\pgfsys@defobject{currentmarker}{\pgfqpoint{-0.020833in}{0.000000in}}{\pgfqpoint{-0.000000in}{0.000000in}}{%
\pgfpathmoveto{\pgfqpoint{-0.000000in}{0.000000in}}%
\pgfpathlineto{\pgfqpoint{-0.020833in}{0.000000in}}%
\pgfusepath{stroke,fill}%
}%
\begin{pgfscope}%
\pgfsys@transformshift{3.038110in}{2.148556in}%
\pgfsys@useobject{currentmarker}{}%
\end{pgfscope}%
\end{pgfscope}%
\begin{pgfscope}%
\pgfsetbuttcap%
\pgfsetroundjoin%
\definecolor{currentfill}{rgb}{0.000000,0.000000,0.000000}%
\pgfsetfillcolor{currentfill}%
\pgfsetlinewidth{0.401500pt}%
\definecolor{currentstroke}{rgb}{0.000000,0.000000,0.000000}%
\pgfsetstrokecolor{currentstroke}%
\pgfsetdash{}{0pt}%
\pgfsys@defobject{currentmarker}{\pgfqpoint{0.000000in}{0.000000in}}{\pgfqpoint{0.020833in}{0.000000in}}{%
\pgfpathmoveto{\pgfqpoint{0.000000in}{0.000000in}}%
\pgfpathlineto{\pgfqpoint{0.020833in}{0.000000in}}%
\pgfusepath{stroke,fill}%
}%
\begin{pgfscope}%
\pgfsys@transformshift{0.650000in}{2.215039in}%
\pgfsys@useobject{currentmarker}{}%
\end{pgfscope}%
\end{pgfscope}%
\begin{pgfscope}%
\pgfsetbuttcap%
\pgfsetroundjoin%
\definecolor{currentfill}{rgb}{0.000000,0.000000,0.000000}%
\pgfsetfillcolor{currentfill}%
\pgfsetlinewidth{0.401500pt}%
\definecolor{currentstroke}{rgb}{0.000000,0.000000,0.000000}%
\pgfsetstrokecolor{currentstroke}%
\pgfsetdash{}{0pt}%
\pgfsys@defobject{currentmarker}{\pgfqpoint{-0.020833in}{0.000000in}}{\pgfqpoint{-0.000000in}{0.000000in}}{%
\pgfpathmoveto{\pgfqpoint{-0.000000in}{0.000000in}}%
\pgfpathlineto{\pgfqpoint{-0.020833in}{0.000000in}}%
\pgfusepath{stroke,fill}%
}%
\begin{pgfscope}%
\pgfsys@transformshift{3.038110in}{2.215039in}%
\pgfsys@useobject{currentmarker}{}%
\end{pgfscope}%
\end{pgfscope}%
\begin{pgfscope}%
\pgfsetbuttcap%
\pgfsetroundjoin%
\definecolor{currentfill}{rgb}{0.000000,0.000000,0.000000}%
\pgfsetfillcolor{currentfill}%
\pgfsetlinewidth{0.401500pt}%
\definecolor{currentstroke}{rgb}{0.000000,0.000000,0.000000}%
\pgfsetstrokecolor{currentstroke}%
\pgfsetdash{}{0pt}%
\pgfsys@defobject{currentmarker}{\pgfqpoint{0.000000in}{0.000000in}}{\pgfqpoint{0.020833in}{0.000000in}}{%
\pgfpathmoveto{\pgfqpoint{0.000000in}{0.000000in}}%
\pgfpathlineto{\pgfqpoint{0.020833in}{0.000000in}}%
\pgfusepath{stroke,fill}%
}%
\begin{pgfscope}%
\pgfsys@transformshift{0.650000in}{2.281522in}%
\pgfsys@useobject{currentmarker}{}%
\end{pgfscope}%
\end{pgfscope}%
\begin{pgfscope}%
\pgfsetbuttcap%
\pgfsetroundjoin%
\definecolor{currentfill}{rgb}{0.000000,0.000000,0.000000}%
\pgfsetfillcolor{currentfill}%
\pgfsetlinewidth{0.401500pt}%
\definecolor{currentstroke}{rgb}{0.000000,0.000000,0.000000}%
\pgfsetstrokecolor{currentstroke}%
\pgfsetdash{}{0pt}%
\pgfsys@defobject{currentmarker}{\pgfqpoint{-0.020833in}{0.000000in}}{\pgfqpoint{-0.000000in}{0.000000in}}{%
\pgfpathmoveto{\pgfqpoint{-0.000000in}{0.000000in}}%
\pgfpathlineto{\pgfqpoint{-0.020833in}{0.000000in}}%
\pgfusepath{stroke,fill}%
}%
\begin{pgfscope}%
\pgfsys@transformshift{3.038110in}{2.281522in}%
\pgfsys@useobject{currentmarker}{}%
\end{pgfscope}%
\end{pgfscope}%
\begin{pgfscope}%
\pgfsetbuttcap%
\pgfsetroundjoin%
\definecolor{currentfill}{rgb}{0.000000,0.000000,0.000000}%
\pgfsetfillcolor{currentfill}%
\pgfsetlinewidth{0.401500pt}%
\definecolor{currentstroke}{rgb}{0.000000,0.000000,0.000000}%
\pgfsetstrokecolor{currentstroke}%
\pgfsetdash{}{0pt}%
\pgfsys@defobject{currentmarker}{\pgfqpoint{0.000000in}{0.000000in}}{\pgfqpoint{0.020833in}{0.000000in}}{%
\pgfpathmoveto{\pgfqpoint{0.000000in}{0.000000in}}%
\pgfpathlineto{\pgfqpoint{0.020833in}{0.000000in}}%
\pgfusepath{stroke,fill}%
}%
\begin{pgfscope}%
\pgfsys@transformshift{0.650000in}{2.348005in}%
\pgfsys@useobject{currentmarker}{}%
\end{pgfscope}%
\end{pgfscope}%
\begin{pgfscope}%
\pgfsetbuttcap%
\pgfsetroundjoin%
\definecolor{currentfill}{rgb}{0.000000,0.000000,0.000000}%
\pgfsetfillcolor{currentfill}%
\pgfsetlinewidth{0.401500pt}%
\definecolor{currentstroke}{rgb}{0.000000,0.000000,0.000000}%
\pgfsetstrokecolor{currentstroke}%
\pgfsetdash{}{0pt}%
\pgfsys@defobject{currentmarker}{\pgfqpoint{-0.020833in}{0.000000in}}{\pgfqpoint{-0.000000in}{0.000000in}}{%
\pgfpathmoveto{\pgfqpoint{-0.000000in}{0.000000in}}%
\pgfpathlineto{\pgfqpoint{-0.020833in}{0.000000in}}%
\pgfusepath{stroke,fill}%
}%
\begin{pgfscope}%
\pgfsys@transformshift{3.038110in}{2.348005in}%
\pgfsys@useobject{currentmarker}{}%
\end{pgfscope}%
\end{pgfscope}%
\begin{pgfscope}%
\definecolor{textcolor}{rgb}{0.000000,0.000000,0.000000}%
\pgfsetstrokecolor{textcolor}%
\pgfsetfillcolor{textcolor}%
\pgftext[x=0.421528in,y=1.417244in,,bottom,rotate=90.000000]{\color{textcolor}{\rmfamily\fontsize{12.000000}{14.400000}\selectfont\catcode`\^=\active\def^{\ifmmode\sp\else\^{}\fi}\catcode`\%=\active\def%{\%}$U^+$}}%
\end{pgfscope}%
\begin{pgfscope}%
\pgfpathrectangle{\pgfqpoint{0.650000in}{0.420000in}}{\pgfqpoint{2.388110in}{1.994488in}}%
\pgfusepath{clip}%
\pgfsetbuttcap%
\pgfsetroundjoin%
\pgfsetlinewidth{0.501875pt}%
\definecolor{currentstroke}{rgb}{0.000000,0.000000,0.000000}%
\pgfsetstrokecolor{currentstroke}%
\pgfsetdash{{1.850000pt}{0.800000pt}}{0.000000pt}%
\pgfpathmoveto{\pgfqpoint{0.645000in}{0.454608in}}%
\pgfpathlineto{\pgfqpoint{0.774036in}{0.471141in}}%
\pgfpathlineto{\pgfqpoint{0.890781in}{0.496711in}}%
\pgfpathlineto{\pgfqpoint{0.973614in}{0.522281in}}%
\pgfpathlineto{\pgfqpoint{1.037864in}{0.547852in}}%
\pgfpathlineto{\pgfqpoint{1.090360in}{0.573422in}}%
\pgfpathlineto{\pgfqpoint{1.134744in}{0.598993in}}%
\pgfpathlineto{\pgfqpoint{1.173192in}{0.624563in}}%
\pgfpathlineto{\pgfqpoint{1.207105in}{0.650133in}}%
\pgfpathlineto{\pgfqpoint{1.237442in}{0.675704in}}%
\pgfpathlineto{\pgfqpoint{1.264885in}{0.701274in}}%
\pgfpathlineto{\pgfqpoint{1.289938in}{0.726844in}}%
\pgfpathlineto{\pgfqpoint{1.312985in}{0.752415in}}%
\pgfpathlineto{\pgfqpoint{1.334323in}{0.777985in}}%
\pgfpathlineto{\pgfqpoint{1.354188in}{0.803555in}}%
\pgfpathlineto{\pgfqpoint{1.372770in}{0.829126in}}%
\pgfpathlineto{\pgfqpoint{1.390226in}{0.854696in}}%
\pgfpathlineto{\pgfqpoint{1.406684in}{0.880266in}}%
\pgfpathlineto{\pgfqpoint{1.422251in}{0.905837in}}%
\pgfpathlineto{\pgfqpoint{1.437020in}{0.931407in}}%
\pgfpathlineto{\pgfqpoint{1.451068in}{0.956978in}}%
\pgfpathlineto{\pgfqpoint{1.464463in}{0.982548in}}%
\pgfpathlineto{\pgfqpoint{1.477262in}{1.008118in}}%
\pgfpathlineto{\pgfqpoint{1.489516in}{1.033689in}}%
\pgfpathlineto{\pgfqpoint{1.501270in}{1.059259in}}%
\pgfpathlineto{\pgfqpoint{1.512563in}{1.084829in}}%
\pgfpathlineto{\pgfqpoint{1.523429in}{1.110400in}}%
\pgfpathlineto{\pgfqpoint{1.533901in}{1.135970in}}%
\pgfpathlineto{\pgfqpoint{1.544005in}{1.161540in}}%
\pgfpathlineto{\pgfqpoint{1.553766in}{1.187111in}}%
\pgfpathlineto{\pgfqpoint{1.563207in}{1.212681in}}%
\pgfpathlineto{\pgfqpoint{1.572349in}{1.238252in}}%
\pgfpathlineto{\pgfqpoint{1.581209in}{1.263822in}}%
\pgfpathlineto{\pgfqpoint{1.589804in}{1.289392in}}%
\pgfpathlineto{\pgfqpoint{1.598151in}{1.314963in}}%
\pgfpathlineto{\pgfqpoint{1.606262in}{1.340533in}}%
\pgfpathlineto{\pgfqpoint{1.614151in}{1.366103in}}%
\pgfpathlineto{\pgfqpoint{1.621829in}{1.391674in}}%
\pgfpathlineto{\pgfqpoint{1.629309in}{1.417244in}}%
\pgfusepath{stroke}%
\end{pgfscope}%
\begin{pgfscope}%
\pgfpathrectangle{\pgfqpoint{0.650000in}{0.420000in}}{\pgfqpoint{2.388110in}{1.994488in}}%
\pgfusepath{clip}%
\pgfsetbuttcap%
\pgfsetroundjoin%
\pgfsetlinewidth{0.501875pt}%
\definecolor{currentstroke}{rgb}{0.000000,0.000000,0.000000}%
\pgfsetstrokecolor{currentstroke}%
\pgfsetdash{{1.850000pt}{0.800000pt}}{0.000000pt}%
\pgfpathmoveto{\pgfqpoint{1.512563in}{1.139084in}}%
\pgfpathlineto{\pgfqpoint{2.871163in}{1.904205in}}%
\pgfpathlineto{\pgfqpoint{3.043110in}{2.001040in}}%
\pgfusepath{stroke}%
\end{pgfscope}%
\begin{pgfscope}%
\pgfpathrectangle{\pgfqpoint{0.650000in}{0.420000in}}{\pgfqpoint{2.388110in}{1.994488in}}%
\pgfusepath{clip}%
\pgfsetrectcap%
\pgfsetroundjoin%
\pgfsetlinewidth{0.602250pt}%
\definecolor{currentstroke}{rgb}{0.000000,0.000000,0.000000}%
\pgfsetstrokecolor{currentstroke}%
\pgfsetdash{}{0pt}%
\pgfpathmoveto{\pgfqpoint{0.645000in}{0.453619in}}%
\pgfpathlineto{\pgfqpoint{0.687793in}{0.457902in}}%
\pgfpathlineto{\pgfqpoint{0.816285in}{0.479211in}}%
\pgfpathlineto{\pgfqpoint{0.921267in}{0.505226in}}%
\pgfpathlineto{\pgfqpoint{1.010024in}{0.535887in}}%
\pgfpathlineto{\pgfqpoint{1.086906in}{0.571061in}}%
\pgfpathlineto{\pgfqpoint{1.154718in}{0.610492in}}%
\pgfpathlineto{\pgfqpoint{1.215374in}{0.653749in}}%
\pgfpathlineto{\pgfqpoint{1.270240in}{0.700182in}}%
\pgfpathlineto{\pgfqpoint{1.320326in}{0.748927in}}%
\pgfpathlineto{\pgfqpoint{1.366397in}{0.798949in}}%
\pgfpathlineto{\pgfqpoint{1.409048in}{0.849139in}}%
\pgfpathlineto{\pgfqpoint{1.448752in}{0.898422in}}%
\pgfpathlineto{\pgfqpoint{1.553655in}{1.032606in}}%
\pgfpathlineto{\pgfqpoint{1.584756in}{1.071186in}}%
\pgfpathlineto{\pgfqpoint{1.614259in}{1.106442in}}%
\pgfpathlineto{\pgfqpoint{1.642318in}{1.138473in}}%
\pgfpathlineto{\pgfqpoint{1.669069in}{1.167470in}}%
\pgfpathlineto{\pgfqpoint{1.694626in}{1.193678in}}%
\pgfpathlineto{\pgfqpoint{1.719092in}{1.217365in}}%
\pgfpathlineto{\pgfqpoint{1.742555in}{1.238800in}}%
\pgfpathlineto{\pgfqpoint{1.765095in}{1.258239in}}%
\pgfpathlineto{\pgfqpoint{1.786780in}{1.275920in}}%
\pgfpathlineto{\pgfqpoint{1.807673in}{1.292057in}}%
\pgfpathlineto{\pgfqpoint{1.847299in}{1.320453in}}%
\pgfpathlineto{\pgfqpoint{1.884352in}{1.344734in}}%
\pgfpathlineto{\pgfqpoint{1.919144in}{1.365923in}}%
\pgfpathlineto{\pgfqpoint{1.967645in}{1.393587in}}%
\pgfpathlineto{\pgfqpoint{2.040279in}{1.432931in}}%
\pgfpathlineto{\pgfqpoint{2.128472in}{1.480772in}}%
\pgfpathlineto{\pgfqpoint{2.183959in}{1.512327in}}%
\pgfpathlineto{\pgfqpoint{2.224704in}{1.536705in}}%
\pgfpathlineto{\pgfqpoint{2.262674in}{1.560649in}}%
\pgfpathlineto{\pgfqpoint{2.298213in}{1.584370in}}%
\pgfpathlineto{\pgfqpoint{2.331604in}{1.608009in}}%
\pgfpathlineto{\pgfqpoint{2.363083in}{1.631703in}}%
\pgfpathlineto{\pgfqpoint{2.392847in}{1.655421in}}%
\pgfpathlineto{\pgfqpoint{2.421066in}{1.679147in}}%
\pgfpathlineto{\pgfqpoint{2.454383in}{1.708750in}}%
\pgfpathlineto{\pgfqpoint{2.479623in}{1.732562in}}%
\pgfpathlineto{\pgfqpoint{2.509576in}{1.762435in}}%
\pgfpathlineto{\pgfqpoint{2.543407in}{1.798153in}}%
\pgfpathlineto{\pgfqpoint{2.575176in}{1.833280in}}%
\pgfpathlineto{\pgfqpoint{2.619426in}{1.884091in}}%
\pgfpathlineto{\pgfqpoint{2.660107in}{1.930660in}}%
\pgfpathlineto{\pgfqpoint{2.677170in}{1.948664in}}%
\pgfpathlineto{\pgfqpoint{2.689587in}{1.960437in}}%
\pgfpathlineto{\pgfqpoint{2.701695in}{1.970539in}}%
\pgfpathlineto{\pgfqpoint{2.713505in}{1.978853in}}%
\pgfpathlineto{\pgfqpoint{2.725030in}{1.985308in}}%
\pgfpathlineto{\pgfqpoint{2.736282in}{1.990005in}}%
\pgfpathlineto{\pgfqpoint{2.747270in}{1.993200in}}%
\pgfpathlineto{\pgfqpoint{2.758004in}{1.995230in}}%
\pgfpathlineto{\pgfqpoint{2.771937in}{1.996700in}}%
\pgfpathlineto{\pgfqpoint{2.792068in}{1.997552in}}%
\pgfpathlineto{\pgfqpoint{2.841653in}{1.998048in}}%
\pgfpathlineto{\pgfqpoint{3.043110in}{1.999362in}}%
\pgfpathlineto{\pgfqpoint{3.043110in}{1.999362in}}%
\pgfusepath{stroke}%
\end{pgfscope}%
\begin{pgfscope}%
\pgfpathrectangle{\pgfqpoint{0.650000in}{0.420000in}}{\pgfqpoint{2.388110in}{1.994488in}}%
\pgfusepath{clip}%
\pgfsetrectcap%
\pgfsetroundjoin%
\pgfsetlinewidth{1.003750pt}%
\definecolor{currentstroke}{rgb}{0.454902,0.678431,0.819608}%
\pgfsetstrokecolor{currentstroke}%
\pgfsetdash{}{0pt}%
\pgfpathmoveto{\pgfqpoint{0.645000in}{0.452999in}}%
\pgfpathlineto{\pgfqpoint{0.679586in}{0.456724in}}%
\pgfpathlineto{\pgfqpoint{0.744946in}{0.465721in}}%
\pgfpathlineto{\pgfqpoint{0.802042in}{0.475441in}}%
\pgfpathlineto{\pgfqpoint{0.853159in}{0.485959in}}%
\pgfpathlineto{\pgfqpoint{0.899759in}{0.497350in}}%
\pgfpathlineto{\pgfqpoint{0.942836in}{0.509682in}}%
\pgfpathlineto{\pgfqpoint{0.983096in}{0.523017in}}%
\pgfpathlineto{\pgfqpoint{1.021057in}{0.537406in}}%
\pgfpathlineto{\pgfqpoint{1.057113in}{0.552895in}}%
\pgfpathlineto{\pgfqpoint{1.091568in}{0.569527in}}%
\pgfpathlineto{\pgfqpoint{1.124662in}{0.587353in}}%
\pgfpathlineto{\pgfqpoint{1.156588in}{0.606444in}}%
\pgfpathlineto{\pgfqpoint{1.187505in}{0.626899in}}%
\pgfpathlineto{\pgfqpoint{1.217543in}{0.648826in}}%
\pgfpathlineto{\pgfqpoint{1.246809in}{0.672278in}}%
\pgfpathlineto{\pgfqpoint{1.275397in}{0.697055in}}%
\pgfpathlineto{\pgfqpoint{1.303382in}{0.723013in}}%
\pgfpathlineto{\pgfqpoint{1.330833in}{0.750283in}}%
\pgfpathlineto{\pgfqpoint{1.357807in}{0.778528in}}%
\pgfpathlineto{\pgfqpoint{1.384354in}{0.807963in}}%
\pgfpathlineto{\pgfqpoint{1.410517in}{0.838335in}}%
\pgfpathlineto{\pgfqpoint{1.461840in}{0.900548in}}%
\pgfpathlineto{\pgfqpoint{1.561290in}{1.024436in}}%
\pgfpathlineto{\pgfqpoint{1.609778in}{1.082227in}}%
\pgfpathlineto{\pgfqpoint{1.633774in}{1.109543in}}%
\pgfpathlineto{\pgfqpoint{1.657624in}{1.135537in}}%
\pgfpathlineto{\pgfqpoint{1.681340in}{1.160308in}}%
\pgfpathlineto{\pgfqpoint{1.704933in}{1.183792in}}%
\pgfpathlineto{\pgfqpoint{1.728413in}{1.206102in}}%
\pgfpathlineto{\pgfqpoint{1.751791in}{1.227083in}}%
\pgfpathlineto{\pgfqpoint{1.775073in}{1.246951in}}%
\pgfpathlineto{\pgfqpoint{1.798269in}{1.265727in}}%
\pgfpathlineto{\pgfqpoint{1.821384in}{1.283464in}}%
\pgfpathlineto{\pgfqpoint{1.844426in}{1.300254in}}%
\pgfpathlineto{\pgfqpoint{1.867400in}{1.316207in}}%
\pgfpathlineto{\pgfqpoint{1.913165in}{1.345897in}}%
\pgfpathlineto{\pgfqpoint{1.958719in}{1.373450in}}%
\pgfpathlineto{\pgfqpoint{2.004092in}{1.399589in}}%
\pgfpathlineto{\pgfqpoint{2.116902in}{1.463350in}}%
\pgfpathlineto{\pgfqpoint{2.161830in}{1.490465in}}%
\pgfpathlineto{\pgfqpoint{2.205858in}{1.519153in}}%
\pgfpathlineto{\pgfqpoint{2.244728in}{1.545786in}}%
\pgfpathlineto{\pgfqpoint{2.278969in}{1.570982in}}%
\pgfpathlineto{\pgfqpoint{2.309568in}{1.595184in}}%
\pgfpathlineto{\pgfqpoint{2.337224in}{1.618619in}}%
\pgfpathlineto{\pgfqpoint{2.362456in}{1.641642in}}%
\pgfpathlineto{\pgfqpoint{2.385654in}{1.663778in}}%
\pgfpathlineto{\pgfqpoint{2.417283in}{1.695799in}}%
\pgfpathlineto{\pgfqpoint{2.454683in}{1.735760in}}%
\pgfpathlineto{\pgfqpoint{2.487779in}{1.772658in}}%
\pgfpathlineto{\pgfqpoint{2.510322in}{1.799353in}}%
\pgfpathlineto{\pgfqpoint{2.563013in}{1.864913in}}%
\pgfpathlineto{\pgfqpoint{2.586135in}{1.896342in}}%
\pgfpathlineto{\pgfqpoint{2.612648in}{1.934567in}}%
\pgfpathlineto{\pgfqpoint{2.650565in}{1.990290in}}%
\pgfpathlineto{\pgfqpoint{2.721107in}{2.088360in}}%
\pgfpathlineto{\pgfqpoint{2.734685in}{2.104998in}}%
\pgfpathlineto{\pgfqpoint{2.744465in}{2.115387in}}%
\pgfpathlineto{\pgfqpoint{2.753923in}{2.123482in}}%
\pgfpathlineto{\pgfqpoint{2.763080in}{2.129619in}}%
\pgfpathlineto{\pgfqpoint{2.771955in}{2.134160in}}%
\pgfpathlineto{\pgfqpoint{2.780564in}{2.136946in}}%
\pgfpathlineto{\pgfqpoint{2.788924in}{2.138542in}}%
\pgfpathlineto{\pgfqpoint{2.797047in}{2.139383in}}%
\pgfpathlineto{\pgfqpoint{2.807534in}{2.139371in}}%
\pgfpathlineto{\pgfqpoint{2.822581in}{2.137784in}}%
\pgfpathlineto{\pgfqpoint{2.843773in}{2.134792in}}%
\pgfpathlineto{\pgfqpoint{2.871867in}{2.129815in}}%
\pgfpathlineto{\pgfqpoint{2.897463in}{2.124823in}}%
\pgfpathlineto{\pgfqpoint{2.967370in}{2.108706in}}%
\pgfpathlineto{\pgfqpoint{3.027200in}{2.092049in}}%
\pgfpathlineto{\pgfqpoint{3.030779in}{2.092173in}}%
\pgfpathlineto{\pgfqpoint{3.033141in}{2.093098in}}%
\pgfpathlineto{\pgfqpoint{3.035484in}{2.094968in}}%
\pgfpathlineto{\pgfqpoint{3.037807in}{2.097975in}}%
\pgfpathlineto{\pgfqpoint{3.041258in}{2.104626in}}%
\pgfpathlineto{\pgfqpoint{3.043110in}{2.108999in}}%
\pgfpathlineto{\pgfqpoint{3.043110in}{2.108999in}}%
\pgfusepath{stroke}%
\end{pgfscope}%
\begin{pgfscope}%
\pgfpathrectangle{\pgfqpoint{0.650000in}{0.420000in}}{\pgfqpoint{2.388110in}{1.994488in}}%
\pgfusepath{clip}%
\pgfsetrectcap%
\pgfsetroundjoin%
\pgfsetlinewidth{1.003750pt}%
\definecolor{currentstroke}{rgb}{0.501961,0.450980,0.674510}%
\pgfsetstrokecolor{currentstroke}%
\pgfsetdash{}{0pt}%
\pgfpathmoveto{\pgfqpoint{0.645000in}{0.452668in}}%
\pgfpathlineto{\pgfqpoint{0.693779in}{0.458615in}}%
\pgfpathlineto{\pgfqpoint{0.743169in}{0.465786in}}%
\pgfpathlineto{\pgfqpoint{0.788495in}{0.473540in}}%
\pgfpathlineto{\pgfqpoint{0.830619in}{0.481925in}}%
\pgfpathlineto{\pgfqpoint{0.870161in}{0.490992in}}%
\pgfpathlineto{\pgfqpoint{0.907582in}{0.500794in}}%
\pgfpathlineto{\pgfqpoint{0.943235in}{0.511390in}}%
\pgfpathlineto{\pgfqpoint{0.977396in}{0.522839in}}%
\pgfpathlineto{\pgfqpoint{1.010285in}{0.535204in}}%
\pgfpathlineto{\pgfqpoint{1.042080in}{0.548552in}}%
\pgfpathlineto{\pgfqpoint{1.072926in}{0.562947in}}%
\pgfpathlineto{\pgfqpoint{1.102943in}{0.578456in}}%
\pgfpathlineto{\pgfqpoint{1.132233in}{0.595142in}}%
\pgfpathlineto{\pgfqpoint{1.160881in}{0.613064in}}%
\pgfpathlineto{\pgfqpoint{1.188960in}{0.632274in}}%
\pgfpathlineto{\pgfqpoint{1.216534in}{0.652816in}}%
\pgfpathlineto{\pgfqpoint{1.243655in}{0.674719in}}%
\pgfpathlineto{\pgfqpoint{1.270372in}{0.697998in}}%
\pgfpathlineto{\pgfqpoint{1.296725in}{0.722647in}}%
\pgfpathlineto{\pgfqpoint{1.322750in}{0.748632in}}%
\pgfpathlineto{\pgfqpoint{1.348480in}{0.775881in}}%
\pgfpathlineto{\pgfqpoint{1.373943in}{0.804274in}}%
\pgfpathlineto{\pgfqpoint{1.399163in}{0.833654in}}%
\pgfpathlineto{\pgfqpoint{1.436586in}{0.879176in}}%
\pgfpathlineto{\pgfqpoint{1.485826in}{0.941477in}}%
\pgfpathlineto{\pgfqpoint{1.558528in}{1.034220in}}%
\pgfpathlineto{\pgfqpoint{1.594450in}{1.078527in}}%
\pgfpathlineto{\pgfqpoint{1.630133in}{1.120560in}}%
\pgfpathlineto{\pgfqpoint{1.653804in}{1.147108in}}%
\pgfpathlineto{\pgfqpoint{1.677389in}{1.172371in}}%
\pgfpathlineto{\pgfqpoint{1.700896in}{1.196302in}}%
\pgfpathlineto{\pgfqpoint{1.724330in}{1.218901in}}%
\pgfpathlineto{\pgfqpoint{1.747699in}{1.240190in}}%
\pgfpathlineto{\pgfqpoint{1.771006in}{1.260224in}}%
\pgfpathlineto{\pgfqpoint{1.794257in}{1.279077in}}%
\pgfpathlineto{\pgfqpoint{1.817457in}{1.296840in}}%
\pgfpathlineto{\pgfqpoint{1.852169in}{1.321659in}}%
\pgfpathlineto{\pgfqpoint{1.886786in}{1.344597in}}%
\pgfpathlineto{\pgfqpoint{1.921320in}{1.366011in}}%
\pgfpathlineto{\pgfqpoint{1.967251in}{1.392836in}}%
\pgfpathlineto{\pgfqpoint{2.024510in}{1.424701in}}%
\pgfpathlineto{\pgfqpoint{2.138624in}{1.487510in}}%
\pgfpathlineto{\pgfqpoint{2.184153in}{1.513912in}}%
\pgfpathlineto{\pgfqpoint{2.229630in}{1.541409in}}%
\pgfpathlineto{\pgfqpoint{2.263686in}{1.562983in}}%
\pgfpathlineto{\pgfqpoint{2.297074in}{1.585270in}}%
\pgfpathlineto{\pgfqpoint{2.327947in}{1.607301in}}%
\pgfpathlineto{\pgfqpoint{2.355880in}{1.628472in}}%
\pgfpathlineto{\pgfqpoint{2.389351in}{1.655284in}}%
\pgfpathlineto{\pgfqpoint{2.419333in}{1.680609in}}%
\pgfpathlineto{\pgfqpoint{2.452892in}{1.710235in}}%
\pgfpathlineto{\pgfqpoint{2.477180in}{1.732958in}}%
\pgfpathlineto{\pgfqpoint{2.504914in}{1.760757in}}%
\pgfpathlineto{\pgfqpoint{2.539738in}{1.797986in}}%
\pgfpathlineto{\pgfqpoint{2.587154in}{1.849231in}}%
\pgfpathlineto{\pgfqpoint{2.613711in}{1.879701in}}%
\pgfpathlineto{\pgfqpoint{2.651047in}{1.924507in}}%
\pgfpathlineto{\pgfqpoint{2.681241in}{1.959880in}}%
\pgfpathlineto{\pgfqpoint{2.700638in}{1.980936in}}%
\pgfpathlineto{\pgfqpoint{2.716283in}{1.996090in}}%
\pgfpathlineto{\pgfqpoint{2.731122in}{2.008651in}}%
\pgfpathlineto{\pgfqpoint{2.742929in}{2.017101in}}%
\pgfpathlineto{\pgfqpoint{2.752038in}{2.022428in}}%
\pgfpathlineto{\pgfqpoint{2.763034in}{2.027472in}}%
\pgfpathlineto{\pgfqpoint{2.773625in}{2.031039in}}%
\pgfpathlineto{\pgfqpoint{2.785840in}{2.033778in}}%
\pgfpathlineto{\pgfqpoint{2.797558in}{2.035176in}}%
\pgfpathlineto{\pgfqpoint{2.814287in}{2.035859in}}%
\pgfpathlineto{\pgfqpoint{2.838519in}{2.035500in}}%
\pgfpathlineto{\pgfqpoint{2.927735in}{2.032196in}}%
\pgfpathlineto{\pgfqpoint{3.014363in}{2.027868in}}%
\pgfpathlineto{\pgfqpoint{3.043110in}{2.026151in}}%
\pgfpathlineto{\pgfqpoint{3.043110in}{2.026151in}}%
\pgfusepath{stroke}%
\end{pgfscope}%
\begin{pgfscope}%
\pgfsetrectcap%
\pgfsetmiterjoin%
\pgfsetlinewidth{0.803000pt}%
\definecolor{currentstroke}{rgb}{0.000000,0.000000,0.000000}%
\pgfsetstrokecolor{currentstroke}%
\pgfsetdash{}{0pt}%
\pgfpathmoveto{\pgfqpoint{0.650000in}{0.420000in}}%
\pgfpathlineto{\pgfqpoint{0.650000in}{2.414488in}}%
\pgfusepath{stroke}%
\end{pgfscope}%
\begin{pgfscope}%
\pgfsetrectcap%
\pgfsetmiterjoin%
\pgfsetlinewidth{0.803000pt}%
\definecolor{currentstroke}{rgb}{0.000000,0.000000,0.000000}%
\pgfsetstrokecolor{currentstroke}%
\pgfsetdash{}{0pt}%
\pgfpathmoveto{\pgfqpoint{3.038110in}{0.420000in}}%
\pgfpathlineto{\pgfqpoint{3.038110in}{2.414488in}}%
\pgfusepath{stroke}%
\end{pgfscope}%
\begin{pgfscope}%
\pgfsetrectcap%
\pgfsetmiterjoin%
\pgfsetlinewidth{0.803000pt}%
\definecolor{currentstroke}{rgb}{0.000000,0.000000,0.000000}%
\pgfsetstrokecolor{currentstroke}%
\pgfsetdash{}{0pt}%
\pgfpathmoveto{\pgfqpoint{0.650000in}{0.420000in}}%
\pgfpathlineto{\pgfqpoint{3.038110in}{0.420000in}}%
\pgfusepath{stroke}%
\end{pgfscope}%
\begin{pgfscope}%
\pgfsetrectcap%
\pgfsetmiterjoin%
\pgfsetlinewidth{0.803000pt}%
\definecolor{currentstroke}{rgb}{0.000000,0.000000,0.000000}%
\pgfsetstrokecolor{currentstroke}%
\pgfsetdash{}{0pt}%
\pgfpathmoveto{\pgfqpoint{0.650000in}{2.414488in}}%
\pgfpathlineto{\pgfqpoint{3.038110in}{2.414488in}}%
\pgfusepath{stroke}%
\end{pgfscope}%
\begin{pgfscope}%
\pgfsetrectcap%
\pgfsetroundjoin%
\pgfsetlinewidth{0.602250pt}%
\definecolor{currentstroke}{rgb}{0.000000,0.000000,0.000000}%
\pgfsetstrokecolor{currentstroke}%
\pgfsetdash{}{0pt}%
\pgfpathmoveto{\pgfqpoint{0.750000in}{2.275599in}}%
\pgfpathlineto{\pgfqpoint{0.861111in}{2.275599in}}%
\pgfpathlineto{\pgfqpoint{0.972222in}{2.275599in}}%
\pgfusepath{stroke}%
\end{pgfscope}%
\begin{pgfscope}%
\definecolor{textcolor}{rgb}{0.000000,0.000000,0.000000}%
\pgfsetstrokecolor{textcolor}%
\pgfsetfillcolor{textcolor}%
\pgftext[x=1.061111in,y=2.236710in,left,base]{\color{textcolor}{\rmfamily\fontsize{8.000000}{9.600000}\selectfont\catcode`\^=\active\def^{\ifmmode\sp\else\^{}\fi}\catcode`\%=\active\def%{\%}ref}}%
\end{pgfscope}%
\begin{pgfscope}%
\pgfsetrectcap%
\pgfsetroundjoin%
\pgfsetlinewidth{1.003750pt}%
\definecolor{currentstroke}{rgb}{0.454902,0.678431,0.819608}%
\pgfsetstrokecolor{currentstroke}%
\pgfsetdash{}{0pt}%
\pgfpathmoveto{\pgfqpoint{0.750000in}{2.120661in}}%
\pgfpathlineto{\pgfqpoint{0.861111in}{2.120661in}}%
\pgfpathlineto{\pgfqpoint{0.972222in}{2.120661in}}%
\pgfusepath{stroke}%
\end{pgfscope}%
\begin{pgfscope}%
\definecolor{textcolor}{rgb}{0.000000,0.000000,0.000000}%
\pgfsetstrokecolor{textcolor}%
\pgfsetfillcolor{textcolor}%
\pgftext[x=1.061111in,y=2.081772in,left,base]{\color{textcolor}{\rmfamily\fontsize{8.000000}{9.600000}\selectfont\catcode`\^=\active\def^{\ifmmode\sp\else\^{}\fi}\catcode`\%=\active\def%{\%}R100deg90}}%
\end{pgfscope}%
\begin{pgfscope}%
\pgfsetrectcap%
\pgfsetroundjoin%
\pgfsetlinewidth{1.003750pt}%
\definecolor{currentstroke}{rgb}{0.501961,0.450980,0.674510}%
\pgfsetstrokecolor{currentstroke}%
\pgfsetdash{}{0pt}%
\pgfpathmoveto{\pgfqpoint{0.750000in}{1.965723in}}%
\pgfpathlineto{\pgfqpoint{0.861111in}{1.965723in}}%
\pgfpathlineto{\pgfqpoint{0.972222in}{1.965723in}}%
\pgfusepath{stroke}%
\end{pgfscope}%
\begin{pgfscope}%
\definecolor{textcolor}{rgb}{0.000000,0.000000,0.000000}%
\pgfsetstrokecolor{textcolor}%
\pgfsetfillcolor{textcolor}%
\pgftext[x=1.061111in,y=1.926834in,left,base]{\color{textcolor}{\rmfamily\fontsize{8.000000}{9.600000}\selectfont\catcode`\^=\active\def^{\ifmmode\sp\else\^{}\fi}\catcode`\%=\active\def%{\%}R500deg18}}%
\end{pgfscope}%
\end{pgfpicture}%
\makeatother%
\endgroup%

    \label{fig:R500deg18R100deg90_U}
\end{subfigure}
\hfill
\begin{subfigure}[t]{0.495\textwidth}
    \centering
    \subcaption{}
    %% Creator: Matplotlib, PGF backend
%%
%% To include the figure in your LaTeX document, write
%%   \input{<filename>.pgf}
%%
%% Make sure the required packages are loaded in your preamble
%%   \usepackage{pgf}
%%
%% Also ensure that all the required font packages are loaded; for instance,
%% the lmodern package is sometimes necessary when using math font.
%%   \usepackage{lmodern}
%%
%% Figures using additional raster images can only be included by \input if
%% they are in the same directory as the main LaTeX file. For loading figures
%% from other directories you can use the `import` package
%%   \usepackage{import}
%%
%% and then include the figures with
%%   \import{<path to file>}{<filename>.pgf}
%%
%% Matplotlib used the following preamble
%%   \def\mathdefault#1{#1}
%%   \everymath=\expandafter{\the\everymath\displaystyle}
%%   \IfFileExists{scrextend.sty}{
%%     \usepackage[fontsize=11.000000pt]{scrextend}
%%   }{
%%     \renewcommand{\normalsize}{\fontsize{11.000000}{13.200000}\selectfont}
%%     \normalsize
%%   }
%%   
%%   \makeatletter\@ifpackageloaded{underscore}{}{\usepackage[strings]{underscore}}\makeatother
%%
\begingroup%
\makeatletter%
\begin{pgfpicture}%
\pgfpathrectangle{\pgfpointorigin}{\pgfqpoint{3.118110in}{2.494488in}}%
\pgfusepath{use as bounding box, clip}%
\begin{pgfscope}%
\pgfsetbuttcap%
\pgfsetmiterjoin%
\definecolor{currentfill}{rgb}{1.000000,1.000000,1.000000}%
\pgfsetfillcolor{currentfill}%
\pgfsetlinewidth{0.000000pt}%
\definecolor{currentstroke}{rgb}{1.000000,1.000000,1.000000}%
\pgfsetstrokecolor{currentstroke}%
\pgfsetdash{}{0pt}%
\pgfpathmoveto{\pgfqpoint{0.000000in}{0.000000in}}%
\pgfpathlineto{\pgfqpoint{3.118110in}{0.000000in}}%
\pgfpathlineto{\pgfqpoint{3.118110in}{2.494488in}}%
\pgfpathlineto{\pgfqpoint{0.000000in}{2.494488in}}%
\pgfpathlineto{\pgfqpoint{0.000000in}{0.000000in}}%
\pgfpathclose%
\pgfusepath{fill}%
\end{pgfscope}%
\begin{pgfscope}%
\pgfsetbuttcap%
\pgfsetmiterjoin%
\definecolor{currentfill}{rgb}{1.000000,1.000000,1.000000}%
\pgfsetfillcolor{currentfill}%
\pgfsetlinewidth{0.000000pt}%
\definecolor{currentstroke}{rgb}{0.000000,0.000000,0.000000}%
\pgfsetstrokecolor{currentstroke}%
\pgfsetstrokeopacity{0.000000}%
\pgfsetdash{}{0pt}%
\pgfpathmoveto{\pgfqpoint{0.650000in}{0.420000in}}%
\pgfpathlineto{\pgfqpoint{3.038110in}{0.420000in}}%
\pgfpathlineto{\pgfqpoint{3.038110in}{2.414488in}}%
\pgfpathlineto{\pgfqpoint{0.650000in}{2.414488in}}%
\pgfpathlineto{\pgfqpoint{0.650000in}{0.420000in}}%
\pgfpathclose%
\pgfusepath{fill}%
\end{pgfscope}%
\begin{pgfscope}%
\pgfsetbuttcap%
\pgfsetroundjoin%
\definecolor{currentfill}{rgb}{0.000000,0.000000,0.000000}%
\pgfsetfillcolor{currentfill}%
\pgfsetlinewidth{0.501875pt}%
\definecolor{currentstroke}{rgb}{0.000000,0.000000,0.000000}%
\pgfsetstrokecolor{currentstroke}%
\pgfsetdash{}{0pt}%
\pgfsys@defobject{currentmarker}{\pgfqpoint{0.000000in}{0.000000in}}{\pgfqpoint{0.000000in}{0.041667in}}{%
\pgfpathmoveto{\pgfqpoint{0.000000in}{0.000000in}}%
\pgfpathlineto{\pgfqpoint{0.000000in}{0.041667in}}%
\pgfusepath{stroke,fill}%
}%
\begin{pgfscope}%
\pgfsys@transformshift{0.849578in}{0.420000in}%
\pgfsys@useobject{currentmarker}{}%
\end{pgfscope}%
\end{pgfscope}%
\begin{pgfscope}%
\pgfsetbuttcap%
\pgfsetroundjoin%
\definecolor{currentfill}{rgb}{0.000000,0.000000,0.000000}%
\pgfsetfillcolor{currentfill}%
\pgfsetlinewidth{0.501875pt}%
\definecolor{currentstroke}{rgb}{0.000000,0.000000,0.000000}%
\pgfsetstrokecolor{currentstroke}%
\pgfsetdash{}{0pt}%
\pgfsys@defobject{currentmarker}{\pgfqpoint{0.000000in}{-0.041667in}}{\pgfqpoint{0.000000in}{0.000000in}}{%
\pgfpathmoveto{\pgfqpoint{0.000000in}{0.000000in}}%
\pgfpathlineto{\pgfqpoint{0.000000in}{-0.041667in}}%
\pgfusepath{stroke,fill}%
}%
\begin{pgfscope}%
\pgfsys@transformshift{0.849578in}{2.414488in}%
\pgfsys@useobject{currentmarker}{}%
\end{pgfscope}%
\end{pgfscope}%
\begin{pgfscope}%
\definecolor{textcolor}{rgb}{0.000000,0.000000,0.000000}%
\pgfsetstrokecolor{textcolor}%
\pgfsetfillcolor{textcolor}%
\pgftext[x=0.849578in,y=0.371389in,,top]{\color{textcolor}{\rmfamily\fontsize{11.000000}{13.200000}\selectfont\catcode`\^=\active\def^{\ifmmode\sp\else\^{}\fi}\catcode`\%=\active\def%{\%}$\mathdefault{10^{0}}$}}%
\end{pgfscope}%
\begin{pgfscope}%
\pgfsetbuttcap%
\pgfsetroundjoin%
\definecolor{currentfill}{rgb}{0.000000,0.000000,0.000000}%
\pgfsetfillcolor{currentfill}%
\pgfsetlinewidth{0.501875pt}%
\definecolor{currentstroke}{rgb}{0.000000,0.000000,0.000000}%
\pgfsetstrokecolor{currentstroke}%
\pgfsetdash{}{0pt}%
\pgfsys@defobject{currentmarker}{\pgfqpoint{0.000000in}{0.000000in}}{\pgfqpoint{0.000000in}{0.041667in}}{%
\pgfpathmoveto{\pgfqpoint{0.000000in}{0.000000in}}%
\pgfpathlineto{\pgfqpoint{0.000000in}{0.041667in}}%
\pgfusepath{stroke,fill}%
}%
\begin{pgfscope}%
\pgfsys@transformshift{1.512563in}{0.420000in}%
\pgfsys@useobject{currentmarker}{}%
\end{pgfscope}%
\end{pgfscope}%
\begin{pgfscope}%
\pgfsetbuttcap%
\pgfsetroundjoin%
\definecolor{currentfill}{rgb}{0.000000,0.000000,0.000000}%
\pgfsetfillcolor{currentfill}%
\pgfsetlinewidth{0.501875pt}%
\definecolor{currentstroke}{rgb}{0.000000,0.000000,0.000000}%
\pgfsetstrokecolor{currentstroke}%
\pgfsetdash{}{0pt}%
\pgfsys@defobject{currentmarker}{\pgfqpoint{0.000000in}{-0.041667in}}{\pgfqpoint{0.000000in}{0.000000in}}{%
\pgfpathmoveto{\pgfqpoint{0.000000in}{0.000000in}}%
\pgfpathlineto{\pgfqpoint{0.000000in}{-0.041667in}}%
\pgfusepath{stroke,fill}%
}%
\begin{pgfscope}%
\pgfsys@transformshift{1.512563in}{2.414488in}%
\pgfsys@useobject{currentmarker}{}%
\end{pgfscope}%
\end{pgfscope}%
\begin{pgfscope}%
\definecolor{textcolor}{rgb}{0.000000,0.000000,0.000000}%
\pgfsetstrokecolor{textcolor}%
\pgfsetfillcolor{textcolor}%
\pgftext[x=1.512563in,y=0.371389in,,top]{\color{textcolor}{\rmfamily\fontsize{11.000000}{13.200000}\selectfont\catcode`\^=\active\def^{\ifmmode\sp\else\^{}\fi}\catcode`\%=\active\def%{\%}$\mathdefault{10^{1}}$}}%
\end{pgfscope}%
\begin{pgfscope}%
\pgfsetbuttcap%
\pgfsetroundjoin%
\definecolor{currentfill}{rgb}{0.000000,0.000000,0.000000}%
\pgfsetfillcolor{currentfill}%
\pgfsetlinewidth{0.501875pt}%
\definecolor{currentstroke}{rgb}{0.000000,0.000000,0.000000}%
\pgfsetstrokecolor{currentstroke}%
\pgfsetdash{}{0pt}%
\pgfsys@defobject{currentmarker}{\pgfqpoint{0.000000in}{0.000000in}}{\pgfqpoint{0.000000in}{0.041667in}}{%
\pgfpathmoveto{\pgfqpoint{0.000000in}{0.000000in}}%
\pgfpathlineto{\pgfqpoint{0.000000in}{0.041667in}}%
\pgfusepath{stroke,fill}%
}%
\begin{pgfscope}%
\pgfsys@transformshift{2.175547in}{0.420000in}%
\pgfsys@useobject{currentmarker}{}%
\end{pgfscope}%
\end{pgfscope}%
\begin{pgfscope}%
\pgfsetbuttcap%
\pgfsetroundjoin%
\definecolor{currentfill}{rgb}{0.000000,0.000000,0.000000}%
\pgfsetfillcolor{currentfill}%
\pgfsetlinewidth{0.501875pt}%
\definecolor{currentstroke}{rgb}{0.000000,0.000000,0.000000}%
\pgfsetstrokecolor{currentstroke}%
\pgfsetdash{}{0pt}%
\pgfsys@defobject{currentmarker}{\pgfqpoint{0.000000in}{-0.041667in}}{\pgfqpoint{0.000000in}{0.000000in}}{%
\pgfpathmoveto{\pgfqpoint{0.000000in}{0.000000in}}%
\pgfpathlineto{\pgfqpoint{0.000000in}{-0.041667in}}%
\pgfusepath{stroke,fill}%
}%
\begin{pgfscope}%
\pgfsys@transformshift{2.175547in}{2.414488in}%
\pgfsys@useobject{currentmarker}{}%
\end{pgfscope}%
\end{pgfscope}%
\begin{pgfscope}%
\definecolor{textcolor}{rgb}{0.000000,0.000000,0.000000}%
\pgfsetstrokecolor{textcolor}%
\pgfsetfillcolor{textcolor}%
\pgftext[x=2.175547in,y=0.371389in,,top]{\color{textcolor}{\rmfamily\fontsize{11.000000}{13.200000}\selectfont\catcode`\^=\active\def^{\ifmmode\sp\else\^{}\fi}\catcode`\%=\active\def%{\%}$\mathdefault{10^{2}}$}}%
\end{pgfscope}%
\begin{pgfscope}%
\pgfsetbuttcap%
\pgfsetroundjoin%
\definecolor{currentfill}{rgb}{0.000000,0.000000,0.000000}%
\pgfsetfillcolor{currentfill}%
\pgfsetlinewidth{0.501875pt}%
\definecolor{currentstroke}{rgb}{0.000000,0.000000,0.000000}%
\pgfsetstrokecolor{currentstroke}%
\pgfsetdash{}{0pt}%
\pgfsys@defobject{currentmarker}{\pgfqpoint{0.000000in}{0.000000in}}{\pgfqpoint{0.000000in}{0.041667in}}{%
\pgfpathmoveto{\pgfqpoint{0.000000in}{0.000000in}}%
\pgfpathlineto{\pgfqpoint{0.000000in}{0.041667in}}%
\pgfusepath{stroke,fill}%
}%
\begin{pgfscope}%
\pgfsys@transformshift{2.838532in}{0.420000in}%
\pgfsys@useobject{currentmarker}{}%
\end{pgfscope}%
\end{pgfscope}%
\begin{pgfscope}%
\pgfsetbuttcap%
\pgfsetroundjoin%
\definecolor{currentfill}{rgb}{0.000000,0.000000,0.000000}%
\pgfsetfillcolor{currentfill}%
\pgfsetlinewidth{0.501875pt}%
\definecolor{currentstroke}{rgb}{0.000000,0.000000,0.000000}%
\pgfsetstrokecolor{currentstroke}%
\pgfsetdash{}{0pt}%
\pgfsys@defobject{currentmarker}{\pgfqpoint{0.000000in}{-0.041667in}}{\pgfqpoint{0.000000in}{0.000000in}}{%
\pgfpathmoveto{\pgfqpoint{0.000000in}{0.000000in}}%
\pgfpathlineto{\pgfqpoint{0.000000in}{-0.041667in}}%
\pgfusepath{stroke,fill}%
}%
\begin{pgfscope}%
\pgfsys@transformshift{2.838532in}{2.414488in}%
\pgfsys@useobject{currentmarker}{}%
\end{pgfscope}%
\end{pgfscope}%
\begin{pgfscope}%
\definecolor{textcolor}{rgb}{0.000000,0.000000,0.000000}%
\pgfsetstrokecolor{textcolor}%
\pgfsetfillcolor{textcolor}%
\pgftext[x=2.838532in,y=0.371389in,,top]{\color{textcolor}{\rmfamily\fontsize{11.000000}{13.200000}\selectfont\catcode`\^=\active\def^{\ifmmode\sp\else\^{}\fi}\catcode`\%=\active\def%{\%}$\mathdefault{10^{3}}$}}%
\end{pgfscope}%
\begin{pgfscope}%
\pgfsetbuttcap%
\pgfsetroundjoin%
\definecolor{currentfill}{rgb}{0.000000,0.000000,0.000000}%
\pgfsetfillcolor{currentfill}%
\pgfsetlinewidth{0.401500pt}%
\definecolor{currentstroke}{rgb}{0.000000,0.000000,0.000000}%
\pgfsetstrokecolor{currentstroke}%
\pgfsetdash{}{0pt}%
\pgfsys@defobject{currentmarker}{\pgfqpoint{0.000000in}{0.000000in}}{\pgfqpoint{0.000000in}{0.020833in}}{%
\pgfpathmoveto{\pgfqpoint{0.000000in}{0.000000in}}%
\pgfpathlineto{\pgfqpoint{0.000000in}{0.020833in}}%
\pgfusepath{stroke,fill}%
}%
\begin{pgfscope}%
\pgfsys@transformshift{0.650000in}{0.420000in}%
\pgfsys@useobject{currentmarker}{}%
\end{pgfscope}%
\end{pgfscope}%
\begin{pgfscope}%
\pgfsetbuttcap%
\pgfsetroundjoin%
\definecolor{currentfill}{rgb}{0.000000,0.000000,0.000000}%
\pgfsetfillcolor{currentfill}%
\pgfsetlinewidth{0.401500pt}%
\definecolor{currentstroke}{rgb}{0.000000,0.000000,0.000000}%
\pgfsetstrokecolor{currentstroke}%
\pgfsetdash{}{0pt}%
\pgfsys@defobject{currentmarker}{\pgfqpoint{0.000000in}{-0.020833in}}{\pgfqpoint{0.000000in}{0.000000in}}{%
\pgfpathmoveto{\pgfqpoint{0.000000in}{0.000000in}}%
\pgfpathlineto{\pgfqpoint{0.000000in}{-0.020833in}}%
\pgfusepath{stroke,fill}%
}%
\begin{pgfscope}%
\pgfsys@transformshift{0.650000in}{2.414488in}%
\pgfsys@useobject{currentmarker}{}%
\end{pgfscope}%
\end{pgfscope}%
\begin{pgfscope}%
\pgfsetbuttcap%
\pgfsetroundjoin%
\definecolor{currentfill}{rgb}{0.000000,0.000000,0.000000}%
\pgfsetfillcolor{currentfill}%
\pgfsetlinewidth{0.401500pt}%
\definecolor{currentstroke}{rgb}{0.000000,0.000000,0.000000}%
\pgfsetstrokecolor{currentstroke}%
\pgfsetdash{}{0pt}%
\pgfsys@defobject{currentmarker}{\pgfqpoint{0.000000in}{0.000000in}}{\pgfqpoint{0.000000in}{0.020833in}}{%
\pgfpathmoveto{\pgfqpoint{0.000000in}{0.000000in}}%
\pgfpathlineto{\pgfqpoint{0.000000in}{0.020833in}}%
\pgfusepath{stroke,fill}%
}%
\begin{pgfscope}%
\pgfsys@transformshift{0.702496in}{0.420000in}%
\pgfsys@useobject{currentmarker}{}%
\end{pgfscope}%
\end{pgfscope}%
\begin{pgfscope}%
\pgfsetbuttcap%
\pgfsetroundjoin%
\definecolor{currentfill}{rgb}{0.000000,0.000000,0.000000}%
\pgfsetfillcolor{currentfill}%
\pgfsetlinewidth{0.401500pt}%
\definecolor{currentstroke}{rgb}{0.000000,0.000000,0.000000}%
\pgfsetstrokecolor{currentstroke}%
\pgfsetdash{}{0pt}%
\pgfsys@defobject{currentmarker}{\pgfqpoint{0.000000in}{-0.020833in}}{\pgfqpoint{0.000000in}{0.000000in}}{%
\pgfpathmoveto{\pgfqpoint{0.000000in}{0.000000in}}%
\pgfpathlineto{\pgfqpoint{0.000000in}{-0.020833in}}%
\pgfusepath{stroke,fill}%
}%
\begin{pgfscope}%
\pgfsys@transformshift{0.702496in}{2.414488in}%
\pgfsys@useobject{currentmarker}{}%
\end{pgfscope}%
\end{pgfscope}%
\begin{pgfscope}%
\pgfsetbuttcap%
\pgfsetroundjoin%
\definecolor{currentfill}{rgb}{0.000000,0.000000,0.000000}%
\pgfsetfillcolor{currentfill}%
\pgfsetlinewidth{0.401500pt}%
\definecolor{currentstroke}{rgb}{0.000000,0.000000,0.000000}%
\pgfsetstrokecolor{currentstroke}%
\pgfsetdash{}{0pt}%
\pgfsys@defobject{currentmarker}{\pgfqpoint{0.000000in}{0.000000in}}{\pgfqpoint{0.000000in}{0.020833in}}{%
\pgfpathmoveto{\pgfqpoint{0.000000in}{0.000000in}}%
\pgfpathlineto{\pgfqpoint{0.000000in}{0.020833in}}%
\pgfusepath{stroke,fill}%
}%
\begin{pgfscope}%
\pgfsys@transformshift{0.746881in}{0.420000in}%
\pgfsys@useobject{currentmarker}{}%
\end{pgfscope}%
\end{pgfscope}%
\begin{pgfscope}%
\pgfsetbuttcap%
\pgfsetroundjoin%
\definecolor{currentfill}{rgb}{0.000000,0.000000,0.000000}%
\pgfsetfillcolor{currentfill}%
\pgfsetlinewidth{0.401500pt}%
\definecolor{currentstroke}{rgb}{0.000000,0.000000,0.000000}%
\pgfsetstrokecolor{currentstroke}%
\pgfsetdash{}{0pt}%
\pgfsys@defobject{currentmarker}{\pgfqpoint{0.000000in}{-0.020833in}}{\pgfqpoint{0.000000in}{0.000000in}}{%
\pgfpathmoveto{\pgfqpoint{0.000000in}{0.000000in}}%
\pgfpathlineto{\pgfqpoint{0.000000in}{-0.020833in}}%
\pgfusepath{stroke,fill}%
}%
\begin{pgfscope}%
\pgfsys@transformshift{0.746881in}{2.414488in}%
\pgfsys@useobject{currentmarker}{}%
\end{pgfscope}%
\end{pgfscope}%
\begin{pgfscope}%
\pgfsetbuttcap%
\pgfsetroundjoin%
\definecolor{currentfill}{rgb}{0.000000,0.000000,0.000000}%
\pgfsetfillcolor{currentfill}%
\pgfsetlinewidth{0.401500pt}%
\definecolor{currentstroke}{rgb}{0.000000,0.000000,0.000000}%
\pgfsetstrokecolor{currentstroke}%
\pgfsetdash{}{0pt}%
\pgfsys@defobject{currentmarker}{\pgfqpoint{0.000000in}{0.000000in}}{\pgfqpoint{0.000000in}{0.020833in}}{%
\pgfpathmoveto{\pgfqpoint{0.000000in}{0.000000in}}%
\pgfpathlineto{\pgfqpoint{0.000000in}{0.020833in}}%
\pgfusepath{stroke,fill}%
}%
\begin{pgfscope}%
\pgfsys@transformshift{0.785328in}{0.420000in}%
\pgfsys@useobject{currentmarker}{}%
\end{pgfscope}%
\end{pgfscope}%
\begin{pgfscope}%
\pgfsetbuttcap%
\pgfsetroundjoin%
\definecolor{currentfill}{rgb}{0.000000,0.000000,0.000000}%
\pgfsetfillcolor{currentfill}%
\pgfsetlinewidth{0.401500pt}%
\definecolor{currentstroke}{rgb}{0.000000,0.000000,0.000000}%
\pgfsetstrokecolor{currentstroke}%
\pgfsetdash{}{0pt}%
\pgfsys@defobject{currentmarker}{\pgfqpoint{0.000000in}{-0.020833in}}{\pgfqpoint{0.000000in}{0.000000in}}{%
\pgfpathmoveto{\pgfqpoint{0.000000in}{0.000000in}}%
\pgfpathlineto{\pgfqpoint{0.000000in}{-0.020833in}}%
\pgfusepath{stroke,fill}%
}%
\begin{pgfscope}%
\pgfsys@transformshift{0.785328in}{2.414488in}%
\pgfsys@useobject{currentmarker}{}%
\end{pgfscope}%
\end{pgfscope}%
\begin{pgfscope}%
\pgfsetbuttcap%
\pgfsetroundjoin%
\definecolor{currentfill}{rgb}{0.000000,0.000000,0.000000}%
\pgfsetfillcolor{currentfill}%
\pgfsetlinewidth{0.401500pt}%
\definecolor{currentstroke}{rgb}{0.000000,0.000000,0.000000}%
\pgfsetstrokecolor{currentstroke}%
\pgfsetdash{}{0pt}%
\pgfsys@defobject{currentmarker}{\pgfqpoint{0.000000in}{0.000000in}}{\pgfqpoint{0.000000in}{0.020833in}}{%
\pgfpathmoveto{\pgfqpoint{0.000000in}{0.000000in}}%
\pgfpathlineto{\pgfqpoint{0.000000in}{0.020833in}}%
\pgfusepath{stroke,fill}%
}%
\begin{pgfscope}%
\pgfsys@transformshift{0.819242in}{0.420000in}%
\pgfsys@useobject{currentmarker}{}%
\end{pgfscope}%
\end{pgfscope}%
\begin{pgfscope}%
\pgfsetbuttcap%
\pgfsetroundjoin%
\definecolor{currentfill}{rgb}{0.000000,0.000000,0.000000}%
\pgfsetfillcolor{currentfill}%
\pgfsetlinewidth{0.401500pt}%
\definecolor{currentstroke}{rgb}{0.000000,0.000000,0.000000}%
\pgfsetstrokecolor{currentstroke}%
\pgfsetdash{}{0pt}%
\pgfsys@defobject{currentmarker}{\pgfqpoint{0.000000in}{-0.020833in}}{\pgfqpoint{0.000000in}{0.000000in}}{%
\pgfpathmoveto{\pgfqpoint{0.000000in}{0.000000in}}%
\pgfpathlineto{\pgfqpoint{0.000000in}{-0.020833in}}%
\pgfusepath{stroke,fill}%
}%
\begin{pgfscope}%
\pgfsys@transformshift{0.819242in}{2.414488in}%
\pgfsys@useobject{currentmarker}{}%
\end{pgfscope}%
\end{pgfscope}%
\begin{pgfscope}%
\pgfsetbuttcap%
\pgfsetroundjoin%
\definecolor{currentfill}{rgb}{0.000000,0.000000,0.000000}%
\pgfsetfillcolor{currentfill}%
\pgfsetlinewidth{0.401500pt}%
\definecolor{currentstroke}{rgb}{0.000000,0.000000,0.000000}%
\pgfsetstrokecolor{currentstroke}%
\pgfsetdash{}{0pt}%
\pgfsys@defobject{currentmarker}{\pgfqpoint{0.000000in}{0.000000in}}{\pgfqpoint{0.000000in}{0.020833in}}{%
\pgfpathmoveto{\pgfqpoint{0.000000in}{0.000000in}}%
\pgfpathlineto{\pgfqpoint{0.000000in}{0.020833in}}%
\pgfusepath{stroke,fill}%
}%
\begin{pgfscope}%
\pgfsys@transformshift{1.049156in}{0.420000in}%
\pgfsys@useobject{currentmarker}{}%
\end{pgfscope}%
\end{pgfscope}%
\begin{pgfscope}%
\pgfsetbuttcap%
\pgfsetroundjoin%
\definecolor{currentfill}{rgb}{0.000000,0.000000,0.000000}%
\pgfsetfillcolor{currentfill}%
\pgfsetlinewidth{0.401500pt}%
\definecolor{currentstroke}{rgb}{0.000000,0.000000,0.000000}%
\pgfsetstrokecolor{currentstroke}%
\pgfsetdash{}{0pt}%
\pgfsys@defobject{currentmarker}{\pgfqpoint{0.000000in}{-0.020833in}}{\pgfqpoint{0.000000in}{0.000000in}}{%
\pgfpathmoveto{\pgfqpoint{0.000000in}{0.000000in}}%
\pgfpathlineto{\pgfqpoint{0.000000in}{-0.020833in}}%
\pgfusepath{stroke,fill}%
}%
\begin{pgfscope}%
\pgfsys@transformshift{1.049156in}{2.414488in}%
\pgfsys@useobject{currentmarker}{}%
\end{pgfscope}%
\end{pgfscope}%
\begin{pgfscope}%
\pgfsetbuttcap%
\pgfsetroundjoin%
\definecolor{currentfill}{rgb}{0.000000,0.000000,0.000000}%
\pgfsetfillcolor{currentfill}%
\pgfsetlinewidth{0.401500pt}%
\definecolor{currentstroke}{rgb}{0.000000,0.000000,0.000000}%
\pgfsetstrokecolor{currentstroke}%
\pgfsetdash{}{0pt}%
\pgfsys@defobject{currentmarker}{\pgfqpoint{0.000000in}{0.000000in}}{\pgfqpoint{0.000000in}{0.020833in}}{%
\pgfpathmoveto{\pgfqpoint{0.000000in}{0.000000in}}%
\pgfpathlineto{\pgfqpoint{0.000000in}{0.020833in}}%
\pgfusepath{stroke,fill}%
}%
\begin{pgfscope}%
\pgfsys@transformshift{1.165902in}{0.420000in}%
\pgfsys@useobject{currentmarker}{}%
\end{pgfscope}%
\end{pgfscope}%
\begin{pgfscope}%
\pgfsetbuttcap%
\pgfsetroundjoin%
\definecolor{currentfill}{rgb}{0.000000,0.000000,0.000000}%
\pgfsetfillcolor{currentfill}%
\pgfsetlinewidth{0.401500pt}%
\definecolor{currentstroke}{rgb}{0.000000,0.000000,0.000000}%
\pgfsetstrokecolor{currentstroke}%
\pgfsetdash{}{0pt}%
\pgfsys@defobject{currentmarker}{\pgfqpoint{0.000000in}{-0.020833in}}{\pgfqpoint{0.000000in}{0.000000in}}{%
\pgfpathmoveto{\pgfqpoint{0.000000in}{0.000000in}}%
\pgfpathlineto{\pgfqpoint{0.000000in}{-0.020833in}}%
\pgfusepath{stroke,fill}%
}%
\begin{pgfscope}%
\pgfsys@transformshift{1.165902in}{2.414488in}%
\pgfsys@useobject{currentmarker}{}%
\end{pgfscope}%
\end{pgfscope}%
\begin{pgfscope}%
\pgfsetbuttcap%
\pgfsetroundjoin%
\definecolor{currentfill}{rgb}{0.000000,0.000000,0.000000}%
\pgfsetfillcolor{currentfill}%
\pgfsetlinewidth{0.401500pt}%
\definecolor{currentstroke}{rgb}{0.000000,0.000000,0.000000}%
\pgfsetstrokecolor{currentstroke}%
\pgfsetdash{}{0pt}%
\pgfsys@defobject{currentmarker}{\pgfqpoint{0.000000in}{0.000000in}}{\pgfqpoint{0.000000in}{0.020833in}}{%
\pgfpathmoveto{\pgfqpoint{0.000000in}{0.000000in}}%
\pgfpathlineto{\pgfqpoint{0.000000in}{0.020833in}}%
\pgfusepath{stroke,fill}%
}%
\begin{pgfscope}%
\pgfsys@transformshift{1.248735in}{0.420000in}%
\pgfsys@useobject{currentmarker}{}%
\end{pgfscope}%
\end{pgfscope}%
\begin{pgfscope}%
\pgfsetbuttcap%
\pgfsetroundjoin%
\definecolor{currentfill}{rgb}{0.000000,0.000000,0.000000}%
\pgfsetfillcolor{currentfill}%
\pgfsetlinewidth{0.401500pt}%
\definecolor{currentstroke}{rgb}{0.000000,0.000000,0.000000}%
\pgfsetstrokecolor{currentstroke}%
\pgfsetdash{}{0pt}%
\pgfsys@defobject{currentmarker}{\pgfqpoint{0.000000in}{-0.020833in}}{\pgfqpoint{0.000000in}{0.000000in}}{%
\pgfpathmoveto{\pgfqpoint{0.000000in}{0.000000in}}%
\pgfpathlineto{\pgfqpoint{0.000000in}{-0.020833in}}%
\pgfusepath{stroke,fill}%
}%
\begin{pgfscope}%
\pgfsys@transformshift{1.248735in}{2.414488in}%
\pgfsys@useobject{currentmarker}{}%
\end{pgfscope}%
\end{pgfscope}%
\begin{pgfscope}%
\pgfsetbuttcap%
\pgfsetroundjoin%
\definecolor{currentfill}{rgb}{0.000000,0.000000,0.000000}%
\pgfsetfillcolor{currentfill}%
\pgfsetlinewidth{0.401500pt}%
\definecolor{currentstroke}{rgb}{0.000000,0.000000,0.000000}%
\pgfsetstrokecolor{currentstroke}%
\pgfsetdash{}{0pt}%
\pgfsys@defobject{currentmarker}{\pgfqpoint{0.000000in}{0.000000in}}{\pgfqpoint{0.000000in}{0.020833in}}{%
\pgfpathmoveto{\pgfqpoint{0.000000in}{0.000000in}}%
\pgfpathlineto{\pgfqpoint{0.000000in}{0.020833in}}%
\pgfusepath{stroke,fill}%
}%
\begin{pgfscope}%
\pgfsys@transformshift{1.312985in}{0.420000in}%
\pgfsys@useobject{currentmarker}{}%
\end{pgfscope}%
\end{pgfscope}%
\begin{pgfscope}%
\pgfsetbuttcap%
\pgfsetroundjoin%
\definecolor{currentfill}{rgb}{0.000000,0.000000,0.000000}%
\pgfsetfillcolor{currentfill}%
\pgfsetlinewidth{0.401500pt}%
\definecolor{currentstroke}{rgb}{0.000000,0.000000,0.000000}%
\pgfsetstrokecolor{currentstroke}%
\pgfsetdash{}{0pt}%
\pgfsys@defobject{currentmarker}{\pgfqpoint{0.000000in}{-0.020833in}}{\pgfqpoint{0.000000in}{0.000000in}}{%
\pgfpathmoveto{\pgfqpoint{0.000000in}{0.000000in}}%
\pgfpathlineto{\pgfqpoint{0.000000in}{-0.020833in}}%
\pgfusepath{stroke,fill}%
}%
\begin{pgfscope}%
\pgfsys@transformshift{1.312985in}{2.414488in}%
\pgfsys@useobject{currentmarker}{}%
\end{pgfscope}%
\end{pgfscope}%
\begin{pgfscope}%
\pgfsetbuttcap%
\pgfsetroundjoin%
\definecolor{currentfill}{rgb}{0.000000,0.000000,0.000000}%
\pgfsetfillcolor{currentfill}%
\pgfsetlinewidth{0.401500pt}%
\definecolor{currentstroke}{rgb}{0.000000,0.000000,0.000000}%
\pgfsetstrokecolor{currentstroke}%
\pgfsetdash{}{0pt}%
\pgfsys@defobject{currentmarker}{\pgfqpoint{0.000000in}{0.000000in}}{\pgfqpoint{0.000000in}{0.020833in}}{%
\pgfpathmoveto{\pgfqpoint{0.000000in}{0.000000in}}%
\pgfpathlineto{\pgfqpoint{0.000000in}{0.020833in}}%
\pgfusepath{stroke,fill}%
}%
\begin{pgfscope}%
\pgfsys@transformshift{1.365480in}{0.420000in}%
\pgfsys@useobject{currentmarker}{}%
\end{pgfscope}%
\end{pgfscope}%
\begin{pgfscope}%
\pgfsetbuttcap%
\pgfsetroundjoin%
\definecolor{currentfill}{rgb}{0.000000,0.000000,0.000000}%
\pgfsetfillcolor{currentfill}%
\pgfsetlinewidth{0.401500pt}%
\definecolor{currentstroke}{rgb}{0.000000,0.000000,0.000000}%
\pgfsetstrokecolor{currentstroke}%
\pgfsetdash{}{0pt}%
\pgfsys@defobject{currentmarker}{\pgfqpoint{0.000000in}{-0.020833in}}{\pgfqpoint{0.000000in}{0.000000in}}{%
\pgfpathmoveto{\pgfqpoint{0.000000in}{0.000000in}}%
\pgfpathlineto{\pgfqpoint{0.000000in}{-0.020833in}}%
\pgfusepath{stroke,fill}%
}%
\begin{pgfscope}%
\pgfsys@transformshift{1.365480in}{2.414488in}%
\pgfsys@useobject{currentmarker}{}%
\end{pgfscope}%
\end{pgfscope}%
\begin{pgfscope}%
\pgfsetbuttcap%
\pgfsetroundjoin%
\definecolor{currentfill}{rgb}{0.000000,0.000000,0.000000}%
\pgfsetfillcolor{currentfill}%
\pgfsetlinewidth{0.401500pt}%
\definecolor{currentstroke}{rgb}{0.000000,0.000000,0.000000}%
\pgfsetstrokecolor{currentstroke}%
\pgfsetdash{}{0pt}%
\pgfsys@defobject{currentmarker}{\pgfqpoint{0.000000in}{0.000000in}}{\pgfqpoint{0.000000in}{0.020833in}}{%
\pgfpathmoveto{\pgfqpoint{0.000000in}{0.000000in}}%
\pgfpathlineto{\pgfqpoint{0.000000in}{0.020833in}}%
\pgfusepath{stroke,fill}%
}%
\begin{pgfscope}%
\pgfsys@transformshift{1.409865in}{0.420000in}%
\pgfsys@useobject{currentmarker}{}%
\end{pgfscope}%
\end{pgfscope}%
\begin{pgfscope}%
\pgfsetbuttcap%
\pgfsetroundjoin%
\definecolor{currentfill}{rgb}{0.000000,0.000000,0.000000}%
\pgfsetfillcolor{currentfill}%
\pgfsetlinewidth{0.401500pt}%
\definecolor{currentstroke}{rgb}{0.000000,0.000000,0.000000}%
\pgfsetstrokecolor{currentstroke}%
\pgfsetdash{}{0pt}%
\pgfsys@defobject{currentmarker}{\pgfqpoint{0.000000in}{-0.020833in}}{\pgfqpoint{0.000000in}{0.000000in}}{%
\pgfpathmoveto{\pgfqpoint{0.000000in}{0.000000in}}%
\pgfpathlineto{\pgfqpoint{0.000000in}{-0.020833in}}%
\pgfusepath{stroke,fill}%
}%
\begin{pgfscope}%
\pgfsys@transformshift{1.409865in}{2.414488in}%
\pgfsys@useobject{currentmarker}{}%
\end{pgfscope}%
\end{pgfscope}%
\begin{pgfscope}%
\pgfsetbuttcap%
\pgfsetroundjoin%
\definecolor{currentfill}{rgb}{0.000000,0.000000,0.000000}%
\pgfsetfillcolor{currentfill}%
\pgfsetlinewidth{0.401500pt}%
\definecolor{currentstroke}{rgb}{0.000000,0.000000,0.000000}%
\pgfsetstrokecolor{currentstroke}%
\pgfsetdash{}{0pt}%
\pgfsys@defobject{currentmarker}{\pgfqpoint{0.000000in}{0.000000in}}{\pgfqpoint{0.000000in}{0.020833in}}{%
\pgfpathmoveto{\pgfqpoint{0.000000in}{0.000000in}}%
\pgfpathlineto{\pgfqpoint{0.000000in}{0.020833in}}%
\pgfusepath{stroke,fill}%
}%
\begin{pgfscope}%
\pgfsys@transformshift{1.448313in}{0.420000in}%
\pgfsys@useobject{currentmarker}{}%
\end{pgfscope}%
\end{pgfscope}%
\begin{pgfscope}%
\pgfsetbuttcap%
\pgfsetroundjoin%
\definecolor{currentfill}{rgb}{0.000000,0.000000,0.000000}%
\pgfsetfillcolor{currentfill}%
\pgfsetlinewidth{0.401500pt}%
\definecolor{currentstroke}{rgb}{0.000000,0.000000,0.000000}%
\pgfsetstrokecolor{currentstroke}%
\pgfsetdash{}{0pt}%
\pgfsys@defobject{currentmarker}{\pgfqpoint{0.000000in}{-0.020833in}}{\pgfqpoint{0.000000in}{0.000000in}}{%
\pgfpathmoveto{\pgfqpoint{0.000000in}{0.000000in}}%
\pgfpathlineto{\pgfqpoint{0.000000in}{-0.020833in}}%
\pgfusepath{stroke,fill}%
}%
\begin{pgfscope}%
\pgfsys@transformshift{1.448313in}{2.414488in}%
\pgfsys@useobject{currentmarker}{}%
\end{pgfscope}%
\end{pgfscope}%
\begin{pgfscope}%
\pgfsetbuttcap%
\pgfsetroundjoin%
\definecolor{currentfill}{rgb}{0.000000,0.000000,0.000000}%
\pgfsetfillcolor{currentfill}%
\pgfsetlinewidth{0.401500pt}%
\definecolor{currentstroke}{rgb}{0.000000,0.000000,0.000000}%
\pgfsetstrokecolor{currentstroke}%
\pgfsetdash{}{0pt}%
\pgfsys@defobject{currentmarker}{\pgfqpoint{0.000000in}{0.000000in}}{\pgfqpoint{0.000000in}{0.020833in}}{%
\pgfpathmoveto{\pgfqpoint{0.000000in}{0.000000in}}%
\pgfpathlineto{\pgfqpoint{0.000000in}{0.020833in}}%
\pgfusepath{stroke,fill}%
}%
\begin{pgfscope}%
\pgfsys@transformshift{1.482226in}{0.420000in}%
\pgfsys@useobject{currentmarker}{}%
\end{pgfscope}%
\end{pgfscope}%
\begin{pgfscope}%
\pgfsetbuttcap%
\pgfsetroundjoin%
\definecolor{currentfill}{rgb}{0.000000,0.000000,0.000000}%
\pgfsetfillcolor{currentfill}%
\pgfsetlinewidth{0.401500pt}%
\definecolor{currentstroke}{rgb}{0.000000,0.000000,0.000000}%
\pgfsetstrokecolor{currentstroke}%
\pgfsetdash{}{0pt}%
\pgfsys@defobject{currentmarker}{\pgfqpoint{0.000000in}{-0.020833in}}{\pgfqpoint{0.000000in}{0.000000in}}{%
\pgfpathmoveto{\pgfqpoint{0.000000in}{0.000000in}}%
\pgfpathlineto{\pgfqpoint{0.000000in}{-0.020833in}}%
\pgfusepath{stroke,fill}%
}%
\begin{pgfscope}%
\pgfsys@transformshift{1.482226in}{2.414488in}%
\pgfsys@useobject{currentmarker}{}%
\end{pgfscope}%
\end{pgfscope}%
\begin{pgfscope}%
\pgfsetbuttcap%
\pgfsetroundjoin%
\definecolor{currentfill}{rgb}{0.000000,0.000000,0.000000}%
\pgfsetfillcolor{currentfill}%
\pgfsetlinewidth{0.401500pt}%
\definecolor{currentstroke}{rgb}{0.000000,0.000000,0.000000}%
\pgfsetstrokecolor{currentstroke}%
\pgfsetdash{}{0pt}%
\pgfsys@defobject{currentmarker}{\pgfqpoint{0.000000in}{0.000000in}}{\pgfqpoint{0.000000in}{0.020833in}}{%
\pgfpathmoveto{\pgfqpoint{0.000000in}{0.000000in}}%
\pgfpathlineto{\pgfqpoint{0.000000in}{0.020833in}}%
\pgfusepath{stroke,fill}%
}%
\begin{pgfscope}%
\pgfsys@transformshift{1.712141in}{0.420000in}%
\pgfsys@useobject{currentmarker}{}%
\end{pgfscope}%
\end{pgfscope}%
\begin{pgfscope}%
\pgfsetbuttcap%
\pgfsetroundjoin%
\definecolor{currentfill}{rgb}{0.000000,0.000000,0.000000}%
\pgfsetfillcolor{currentfill}%
\pgfsetlinewidth{0.401500pt}%
\definecolor{currentstroke}{rgb}{0.000000,0.000000,0.000000}%
\pgfsetstrokecolor{currentstroke}%
\pgfsetdash{}{0pt}%
\pgfsys@defobject{currentmarker}{\pgfqpoint{0.000000in}{-0.020833in}}{\pgfqpoint{0.000000in}{0.000000in}}{%
\pgfpathmoveto{\pgfqpoint{0.000000in}{0.000000in}}%
\pgfpathlineto{\pgfqpoint{0.000000in}{-0.020833in}}%
\pgfusepath{stroke,fill}%
}%
\begin{pgfscope}%
\pgfsys@transformshift{1.712141in}{2.414488in}%
\pgfsys@useobject{currentmarker}{}%
\end{pgfscope}%
\end{pgfscope}%
\begin{pgfscope}%
\pgfsetbuttcap%
\pgfsetroundjoin%
\definecolor{currentfill}{rgb}{0.000000,0.000000,0.000000}%
\pgfsetfillcolor{currentfill}%
\pgfsetlinewidth{0.401500pt}%
\definecolor{currentstroke}{rgb}{0.000000,0.000000,0.000000}%
\pgfsetstrokecolor{currentstroke}%
\pgfsetdash{}{0pt}%
\pgfsys@defobject{currentmarker}{\pgfqpoint{0.000000in}{0.000000in}}{\pgfqpoint{0.000000in}{0.020833in}}{%
\pgfpathmoveto{\pgfqpoint{0.000000in}{0.000000in}}%
\pgfpathlineto{\pgfqpoint{0.000000in}{0.020833in}}%
\pgfusepath{stroke,fill}%
}%
\begin{pgfscope}%
\pgfsys@transformshift{1.828887in}{0.420000in}%
\pgfsys@useobject{currentmarker}{}%
\end{pgfscope}%
\end{pgfscope}%
\begin{pgfscope}%
\pgfsetbuttcap%
\pgfsetroundjoin%
\definecolor{currentfill}{rgb}{0.000000,0.000000,0.000000}%
\pgfsetfillcolor{currentfill}%
\pgfsetlinewidth{0.401500pt}%
\definecolor{currentstroke}{rgb}{0.000000,0.000000,0.000000}%
\pgfsetstrokecolor{currentstroke}%
\pgfsetdash{}{0pt}%
\pgfsys@defobject{currentmarker}{\pgfqpoint{0.000000in}{-0.020833in}}{\pgfqpoint{0.000000in}{0.000000in}}{%
\pgfpathmoveto{\pgfqpoint{0.000000in}{0.000000in}}%
\pgfpathlineto{\pgfqpoint{0.000000in}{-0.020833in}}%
\pgfusepath{stroke,fill}%
}%
\begin{pgfscope}%
\pgfsys@transformshift{1.828887in}{2.414488in}%
\pgfsys@useobject{currentmarker}{}%
\end{pgfscope}%
\end{pgfscope}%
\begin{pgfscope}%
\pgfsetbuttcap%
\pgfsetroundjoin%
\definecolor{currentfill}{rgb}{0.000000,0.000000,0.000000}%
\pgfsetfillcolor{currentfill}%
\pgfsetlinewidth{0.401500pt}%
\definecolor{currentstroke}{rgb}{0.000000,0.000000,0.000000}%
\pgfsetstrokecolor{currentstroke}%
\pgfsetdash{}{0pt}%
\pgfsys@defobject{currentmarker}{\pgfqpoint{0.000000in}{0.000000in}}{\pgfqpoint{0.000000in}{0.020833in}}{%
\pgfpathmoveto{\pgfqpoint{0.000000in}{0.000000in}}%
\pgfpathlineto{\pgfqpoint{0.000000in}{0.020833in}}%
\pgfusepath{stroke,fill}%
}%
\begin{pgfscope}%
\pgfsys@transformshift{1.911719in}{0.420000in}%
\pgfsys@useobject{currentmarker}{}%
\end{pgfscope}%
\end{pgfscope}%
\begin{pgfscope}%
\pgfsetbuttcap%
\pgfsetroundjoin%
\definecolor{currentfill}{rgb}{0.000000,0.000000,0.000000}%
\pgfsetfillcolor{currentfill}%
\pgfsetlinewidth{0.401500pt}%
\definecolor{currentstroke}{rgb}{0.000000,0.000000,0.000000}%
\pgfsetstrokecolor{currentstroke}%
\pgfsetdash{}{0pt}%
\pgfsys@defobject{currentmarker}{\pgfqpoint{0.000000in}{-0.020833in}}{\pgfqpoint{0.000000in}{0.000000in}}{%
\pgfpathmoveto{\pgfqpoint{0.000000in}{0.000000in}}%
\pgfpathlineto{\pgfqpoint{0.000000in}{-0.020833in}}%
\pgfusepath{stroke,fill}%
}%
\begin{pgfscope}%
\pgfsys@transformshift{1.911719in}{2.414488in}%
\pgfsys@useobject{currentmarker}{}%
\end{pgfscope}%
\end{pgfscope}%
\begin{pgfscope}%
\pgfsetbuttcap%
\pgfsetroundjoin%
\definecolor{currentfill}{rgb}{0.000000,0.000000,0.000000}%
\pgfsetfillcolor{currentfill}%
\pgfsetlinewidth{0.401500pt}%
\definecolor{currentstroke}{rgb}{0.000000,0.000000,0.000000}%
\pgfsetstrokecolor{currentstroke}%
\pgfsetdash{}{0pt}%
\pgfsys@defobject{currentmarker}{\pgfqpoint{0.000000in}{0.000000in}}{\pgfqpoint{0.000000in}{0.020833in}}{%
\pgfpathmoveto{\pgfqpoint{0.000000in}{0.000000in}}%
\pgfpathlineto{\pgfqpoint{0.000000in}{0.020833in}}%
\pgfusepath{stroke,fill}%
}%
\begin{pgfscope}%
\pgfsys@transformshift{1.975969in}{0.420000in}%
\pgfsys@useobject{currentmarker}{}%
\end{pgfscope}%
\end{pgfscope}%
\begin{pgfscope}%
\pgfsetbuttcap%
\pgfsetroundjoin%
\definecolor{currentfill}{rgb}{0.000000,0.000000,0.000000}%
\pgfsetfillcolor{currentfill}%
\pgfsetlinewidth{0.401500pt}%
\definecolor{currentstroke}{rgb}{0.000000,0.000000,0.000000}%
\pgfsetstrokecolor{currentstroke}%
\pgfsetdash{}{0pt}%
\pgfsys@defobject{currentmarker}{\pgfqpoint{0.000000in}{-0.020833in}}{\pgfqpoint{0.000000in}{0.000000in}}{%
\pgfpathmoveto{\pgfqpoint{0.000000in}{0.000000in}}%
\pgfpathlineto{\pgfqpoint{0.000000in}{-0.020833in}}%
\pgfusepath{stroke,fill}%
}%
\begin{pgfscope}%
\pgfsys@transformshift{1.975969in}{2.414488in}%
\pgfsys@useobject{currentmarker}{}%
\end{pgfscope}%
\end{pgfscope}%
\begin{pgfscope}%
\pgfsetbuttcap%
\pgfsetroundjoin%
\definecolor{currentfill}{rgb}{0.000000,0.000000,0.000000}%
\pgfsetfillcolor{currentfill}%
\pgfsetlinewidth{0.401500pt}%
\definecolor{currentstroke}{rgb}{0.000000,0.000000,0.000000}%
\pgfsetstrokecolor{currentstroke}%
\pgfsetdash{}{0pt}%
\pgfsys@defobject{currentmarker}{\pgfqpoint{0.000000in}{0.000000in}}{\pgfqpoint{0.000000in}{0.020833in}}{%
\pgfpathmoveto{\pgfqpoint{0.000000in}{0.000000in}}%
\pgfpathlineto{\pgfqpoint{0.000000in}{0.020833in}}%
\pgfusepath{stroke,fill}%
}%
\begin{pgfscope}%
\pgfsys@transformshift{2.028465in}{0.420000in}%
\pgfsys@useobject{currentmarker}{}%
\end{pgfscope}%
\end{pgfscope}%
\begin{pgfscope}%
\pgfsetbuttcap%
\pgfsetroundjoin%
\definecolor{currentfill}{rgb}{0.000000,0.000000,0.000000}%
\pgfsetfillcolor{currentfill}%
\pgfsetlinewidth{0.401500pt}%
\definecolor{currentstroke}{rgb}{0.000000,0.000000,0.000000}%
\pgfsetstrokecolor{currentstroke}%
\pgfsetdash{}{0pt}%
\pgfsys@defobject{currentmarker}{\pgfqpoint{0.000000in}{-0.020833in}}{\pgfqpoint{0.000000in}{0.000000in}}{%
\pgfpathmoveto{\pgfqpoint{0.000000in}{0.000000in}}%
\pgfpathlineto{\pgfqpoint{0.000000in}{-0.020833in}}%
\pgfusepath{stroke,fill}%
}%
\begin{pgfscope}%
\pgfsys@transformshift{2.028465in}{2.414488in}%
\pgfsys@useobject{currentmarker}{}%
\end{pgfscope}%
\end{pgfscope}%
\begin{pgfscope}%
\pgfsetbuttcap%
\pgfsetroundjoin%
\definecolor{currentfill}{rgb}{0.000000,0.000000,0.000000}%
\pgfsetfillcolor{currentfill}%
\pgfsetlinewidth{0.401500pt}%
\definecolor{currentstroke}{rgb}{0.000000,0.000000,0.000000}%
\pgfsetstrokecolor{currentstroke}%
\pgfsetdash{}{0pt}%
\pgfsys@defobject{currentmarker}{\pgfqpoint{0.000000in}{0.000000in}}{\pgfqpoint{0.000000in}{0.020833in}}{%
\pgfpathmoveto{\pgfqpoint{0.000000in}{0.000000in}}%
\pgfpathlineto{\pgfqpoint{0.000000in}{0.020833in}}%
\pgfusepath{stroke,fill}%
}%
\begin{pgfscope}%
\pgfsys@transformshift{2.072850in}{0.420000in}%
\pgfsys@useobject{currentmarker}{}%
\end{pgfscope}%
\end{pgfscope}%
\begin{pgfscope}%
\pgfsetbuttcap%
\pgfsetroundjoin%
\definecolor{currentfill}{rgb}{0.000000,0.000000,0.000000}%
\pgfsetfillcolor{currentfill}%
\pgfsetlinewidth{0.401500pt}%
\definecolor{currentstroke}{rgb}{0.000000,0.000000,0.000000}%
\pgfsetstrokecolor{currentstroke}%
\pgfsetdash{}{0pt}%
\pgfsys@defobject{currentmarker}{\pgfqpoint{0.000000in}{-0.020833in}}{\pgfqpoint{0.000000in}{0.000000in}}{%
\pgfpathmoveto{\pgfqpoint{0.000000in}{0.000000in}}%
\pgfpathlineto{\pgfqpoint{0.000000in}{-0.020833in}}%
\pgfusepath{stroke,fill}%
}%
\begin{pgfscope}%
\pgfsys@transformshift{2.072850in}{2.414488in}%
\pgfsys@useobject{currentmarker}{}%
\end{pgfscope}%
\end{pgfscope}%
\begin{pgfscope}%
\pgfsetbuttcap%
\pgfsetroundjoin%
\definecolor{currentfill}{rgb}{0.000000,0.000000,0.000000}%
\pgfsetfillcolor{currentfill}%
\pgfsetlinewidth{0.401500pt}%
\definecolor{currentstroke}{rgb}{0.000000,0.000000,0.000000}%
\pgfsetstrokecolor{currentstroke}%
\pgfsetdash{}{0pt}%
\pgfsys@defobject{currentmarker}{\pgfqpoint{0.000000in}{0.000000in}}{\pgfqpoint{0.000000in}{0.020833in}}{%
\pgfpathmoveto{\pgfqpoint{0.000000in}{0.000000in}}%
\pgfpathlineto{\pgfqpoint{0.000000in}{0.020833in}}%
\pgfusepath{stroke,fill}%
}%
\begin{pgfscope}%
\pgfsys@transformshift{2.111298in}{0.420000in}%
\pgfsys@useobject{currentmarker}{}%
\end{pgfscope}%
\end{pgfscope}%
\begin{pgfscope}%
\pgfsetbuttcap%
\pgfsetroundjoin%
\definecolor{currentfill}{rgb}{0.000000,0.000000,0.000000}%
\pgfsetfillcolor{currentfill}%
\pgfsetlinewidth{0.401500pt}%
\definecolor{currentstroke}{rgb}{0.000000,0.000000,0.000000}%
\pgfsetstrokecolor{currentstroke}%
\pgfsetdash{}{0pt}%
\pgfsys@defobject{currentmarker}{\pgfqpoint{0.000000in}{-0.020833in}}{\pgfqpoint{0.000000in}{0.000000in}}{%
\pgfpathmoveto{\pgfqpoint{0.000000in}{0.000000in}}%
\pgfpathlineto{\pgfqpoint{0.000000in}{-0.020833in}}%
\pgfusepath{stroke,fill}%
}%
\begin{pgfscope}%
\pgfsys@transformshift{2.111298in}{2.414488in}%
\pgfsys@useobject{currentmarker}{}%
\end{pgfscope}%
\end{pgfscope}%
\begin{pgfscope}%
\pgfsetbuttcap%
\pgfsetroundjoin%
\definecolor{currentfill}{rgb}{0.000000,0.000000,0.000000}%
\pgfsetfillcolor{currentfill}%
\pgfsetlinewidth{0.401500pt}%
\definecolor{currentstroke}{rgb}{0.000000,0.000000,0.000000}%
\pgfsetstrokecolor{currentstroke}%
\pgfsetdash{}{0pt}%
\pgfsys@defobject{currentmarker}{\pgfqpoint{0.000000in}{0.000000in}}{\pgfqpoint{0.000000in}{0.020833in}}{%
\pgfpathmoveto{\pgfqpoint{0.000000in}{0.000000in}}%
\pgfpathlineto{\pgfqpoint{0.000000in}{0.020833in}}%
\pgfusepath{stroke,fill}%
}%
\begin{pgfscope}%
\pgfsys@transformshift{2.145211in}{0.420000in}%
\pgfsys@useobject{currentmarker}{}%
\end{pgfscope}%
\end{pgfscope}%
\begin{pgfscope}%
\pgfsetbuttcap%
\pgfsetroundjoin%
\definecolor{currentfill}{rgb}{0.000000,0.000000,0.000000}%
\pgfsetfillcolor{currentfill}%
\pgfsetlinewidth{0.401500pt}%
\definecolor{currentstroke}{rgb}{0.000000,0.000000,0.000000}%
\pgfsetstrokecolor{currentstroke}%
\pgfsetdash{}{0pt}%
\pgfsys@defobject{currentmarker}{\pgfqpoint{0.000000in}{-0.020833in}}{\pgfqpoint{0.000000in}{0.000000in}}{%
\pgfpathmoveto{\pgfqpoint{0.000000in}{0.000000in}}%
\pgfpathlineto{\pgfqpoint{0.000000in}{-0.020833in}}%
\pgfusepath{stroke,fill}%
}%
\begin{pgfscope}%
\pgfsys@transformshift{2.145211in}{2.414488in}%
\pgfsys@useobject{currentmarker}{}%
\end{pgfscope}%
\end{pgfscope}%
\begin{pgfscope}%
\pgfsetbuttcap%
\pgfsetroundjoin%
\definecolor{currentfill}{rgb}{0.000000,0.000000,0.000000}%
\pgfsetfillcolor{currentfill}%
\pgfsetlinewidth{0.401500pt}%
\definecolor{currentstroke}{rgb}{0.000000,0.000000,0.000000}%
\pgfsetstrokecolor{currentstroke}%
\pgfsetdash{}{0pt}%
\pgfsys@defobject{currentmarker}{\pgfqpoint{0.000000in}{0.000000in}}{\pgfqpoint{0.000000in}{0.020833in}}{%
\pgfpathmoveto{\pgfqpoint{0.000000in}{0.000000in}}%
\pgfpathlineto{\pgfqpoint{0.000000in}{0.020833in}}%
\pgfusepath{stroke,fill}%
}%
\begin{pgfscope}%
\pgfsys@transformshift{2.375126in}{0.420000in}%
\pgfsys@useobject{currentmarker}{}%
\end{pgfscope}%
\end{pgfscope}%
\begin{pgfscope}%
\pgfsetbuttcap%
\pgfsetroundjoin%
\definecolor{currentfill}{rgb}{0.000000,0.000000,0.000000}%
\pgfsetfillcolor{currentfill}%
\pgfsetlinewidth{0.401500pt}%
\definecolor{currentstroke}{rgb}{0.000000,0.000000,0.000000}%
\pgfsetstrokecolor{currentstroke}%
\pgfsetdash{}{0pt}%
\pgfsys@defobject{currentmarker}{\pgfqpoint{0.000000in}{-0.020833in}}{\pgfqpoint{0.000000in}{0.000000in}}{%
\pgfpathmoveto{\pgfqpoint{0.000000in}{0.000000in}}%
\pgfpathlineto{\pgfqpoint{0.000000in}{-0.020833in}}%
\pgfusepath{stroke,fill}%
}%
\begin{pgfscope}%
\pgfsys@transformshift{2.375126in}{2.414488in}%
\pgfsys@useobject{currentmarker}{}%
\end{pgfscope}%
\end{pgfscope}%
\begin{pgfscope}%
\pgfsetbuttcap%
\pgfsetroundjoin%
\definecolor{currentfill}{rgb}{0.000000,0.000000,0.000000}%
\pgfsetfillcolor{currentfill}%
\pgfsetlinewidth{0.401500pt}%
\definecolor{currentstroke}{rgb}{0.000000,0.000000,0.000000}%
\pgfsetstrokecolor{currentstroke}%
\pgfsetdash{}{0pt}%
\pgfsys@defobject{currentmarker}{\pgfqpoint{0.000000in}{0.000000in}}{\pgfqpoint{0.000000in}{0.020833in}}{%
\pgfpathmoveto{\pgfqpoint{0.000000in}{0.000000in}}%
\pgfpathlineto{\pgfqpoint{0.000000in}{0.020833in}}%
\pgfusepath{stroke,fill}%
}%
\begin{pgfscope}%
\pgfsys@transformshift{2.491871in}{0.420000in}%
\pgfsys@useobject{currentmarker}{}%
\end{pgfscope}%
\end{pgfscope}%
\begin{pgfscope}%
\pgfsetbuttcap%
\pgfsetroundjoin%
\definecolor{currentfill}{rgb}{0.000000,0.000000,0.000000}%
\pgfsetfillcolor{currentfill}%
\pgfsetlinewidth{0.401500pt}%
\definecolor{currentstroke}{rgb}{0.000000,0.000000,0.000000}%
\pgfsetstrokecolor{currentstroke}%
\pgfsetdash{}{0pt}%
\pgfsys@defobject{currentmarker}{\pgfqpoint{0.000000in}{-0.020833in}}{\pgfqpoint{0.000000in}{0.000000in}}{%
\pgfpathmoveto{\pgfqpoint{0.000000in}{0.000000in}}%
\pgfpathlineto{\pgfqpoint{0.000000in}{-0.020833in}}%
\pgfusepath{stroke,fill}%
}%
\begin{pgfscope}%
\pgfsys@transformshift{2.491871in}{2.414488in}%
\pgfsys@useobject{currentmarker}{}%
\end{pgfscope}%
\end{pgfscope}%
\begin{pgfscope}%
\pgfsetbuttcap%
\pgfsetroundjoin%
\definecolor{currentfill}{rgb}{0.000000,0.000000,0.000000}%
\pgfsetfillcolor{currentfill}%
\pgfsetlinewidth{0.401500pt}%
\definecolor{currentstroke}{rgb}{0.000000,0.000000,0.000000}%
\pgfsetstrokecolor{currentstroke}%
\pgfsetdash{}{0pt}%
\pgfsys@defobject{currentmarker}{\pgfqpoint{0.000000in}{0.000000in}}{\pgfqpoint{0.000000in}{0.020833in}}{%
\pgfpathmoveto{\pgfqpoint{0.000000in}{0.000000in}}%
\pgfpathlineto{\pgfqpoint{0.000000in}{0.020833in}}%
\pgfusepath{stroke,fill}%
}%
\begin{pgfscope}%
\pgfsys@transformshift{2.574704in}{0.420000in}%
\pgfsys@useobject{currentmarker}{}%
\end{pgfscope}%
\end{pgfscope}%
\begin{pgfscope}%
\pgfsetbuttcap%
\pgfsetroundjoin%
\definecolor{currentfill}{rgb}{0.000000,0.000000,0.000000}%
\pgfsetfillcolor{currentfill}%
\pgfsetlinewidth{0.401500pt}%
\definecolor{currentstroke}{rgb}{0.000000,0.000000,0.000000}%
\pgfsetstrokecolor{currentstroke}%
\pgfsetdash{}{0pt}%
\pgfsys@defobject{currentmarker}{\pgfqpoint{0.000000in}{-0.020833in}}{\pgfqpoint{0.000000in}{0.000000in}}{%
\pgfpathmoveto{\pgfqpoint{0.000000in}{0.000000in}}%
\pgfpathlineto{\pgfqpoint{0.000000in}{-0.020833in}}%
\pgfusepath{stroke,fill}%
}%
\begin{pgfscope}%
\pgfsys@transformshift{2.574704in}{2.414488in}%
\pgfsys@useobject{currentmarker}{}%
\end{pgfscope}%
\end{pgfscope}%
\begin{pgfscope}%
\pgfsetbuttcap%
\pgfsetroundjoin%
\definecolor{currentfill}{rgb}{0.000000,0.000000,0.000000}%
\pgfsetfillcolor{currentfill}%
\pgfsetlinewidth{0.401500pt}%
\definecolor{currentstroke}{rgb}{0.000000,0.000000,0.000000}%
\pgfsetstrokecolor{currentstroke}%
\pgfsetdash{}{0pt}%
\pgfsys@defobject{currentmarker}{\pgfqpoint{0.000000in}{0.000000in}}{\pgfqpoint{0.000000in}{0.020833in}}{%
\pgfpathmoveto{\pgfqpoint{0.000000in}{0.000000in}}%
\pgfpathlineto{\pgfqpoint{0.000000in}{0.020833in}}%
\pgfusepath{stroke,fill}%
}%
\begin{pgfscope}%
\pgfsys@transformshift{2.638954in}{0.420000in}%
\pgfsys@useobject{currentmarker}{}%
\end{pgfscope}%
\end{pgfscope}%
\begin{pgfscope}%
\pgfsetbuttcap%
\pgfsetroundjoin%
\definecolor{currentfill}{rgb}{0.000000,0.000000,0.000000}%
\pgfsetfillcolor{currentfill}%
\pgfsetlinewidth{0.401500pt}%
\definecolor{currentstroke}{rgb}{0.000000,0.000000,0.000000}%
\pgfsetstrokecolor{currentstroke}%
\pgfsetdash{}{0pt}%
\pgfsys@defobject{currentmarker}{\pgfqpoint{0.000000in}{-0.020833in}}{\pgfqpoint{0.000000in}{0.000000in}}{%
\pgfpathmoveto{\pgfqpoint{0.000000in}{0.000000in}}%
\pgfpathlineto{\pgfqpoint{0.000000in}{-0.020833in}}%
\pgfusepath{stroke,fill}%
}%
\begin{pgfscope}%
\pgfsys@transformshift{2.638954in}{2.414488in}%
\pgfsys@useobject{currentmarker}{}%
\end{pgfscope}%
\end{pgfscope}%
\begin{pgfscope}%
\pgfsetbuttcap%
\pgfsetroundjoin%
\definecolor{currentfill}{rgb}{0.000000,0.000000,0.000000}%
\pgfsetfillcolor{currentfill}%
\pgfsetlinewidth{0.401500pt}%
\definecolor{currentstroke}{rgb}{0.000000,0.000000,0.000000}%
\pgfsetstrokecolor{currentstroke}%
\pgfsetdash{}{0pt}%
\pgfsys@defobject{currentmarker}{\pgfqpoint{0.000000in}{0.000000in}}{\pgfqpoint{0.000000in}{0.020833in}}{%
\pgfpathmoveto{\pgfqpoint{0.000000in}{0.000000in}}%
\pgfpathlineto{\pgfqpoint{0.000000in}{0.020833in}}%
\pgfusepath{stroke,fill}%
}%
\begin{pgfscope}%
\pgfsys@transformshift{2.691450in}{0.420000in}%
\pgfsys@useobject{currentmarker}{}%
\end{pgfscope}%
\end{pgfscope}%
\begin{pgfscope}%
\pgfsetbuttcap%
\pgfsetroundjoin%
\definecolor{currentfill}{rgb}{0.000000,0.000000,0.000000}%
\pgfsetfillcolor{currentfill}%
\pgfsetlinewidth{0.401500pt}%
\definecolor{currentstroke}{rgb}{0.000000,0.000000,0.000000}%
\pgfsetstrokecolor{currentstroke}%
\pgfsetdash{}{0pt}%
\pgfsys@defobject{currentmarker}{\pgfqpoint{0.000000in}{-0.020833in}}{\pgfqpoint{0.000000in}{0.000000in}}{%
\pgfpathmoveto{\pgfqpoint{0.000000in}{0.000000in}}%
\pgfpathlineto{\pgfqpoint{0.000000in}{-0.020833in}}%
\pgfusepath{stroke,fill}%
}%
\begin{pgfscope}%
\pgfsys@transformshift{2.691450in}{2.414488in}%
\pgfsys@useobject{currentmarker}{}%
\end{pgfscope}%
\end{pgfscope}%
\begin{pgfscope}%
\pgfsetbuttcap%
\pgfsetroundjoin%
\definecolor{currentfill}{rgb}{0.000000,0.000000,0.000000}%
\pgfsetfillcolor{currentfill}%
\pgfsetlinewidth{0.401500pt}%
\definecolor{currentstroke}{rgb}{0.000000,0.000000,0.000000}%
\pgfsetstrokecolor{currentstroke}%
\pgfsetdash{}{0pt}%
\pgfsys@defobject{currentmarker}{\pgfqpoint{0.000000in}{0.000000in}}{\pgfqpoint{0.000000in}{0.020833in}}{%
\pgfpathmoveto{\pgfqpoint{0.000000in}{0.000000in}}%
\pgfpathlineto{\pgfqpoint{0.000000in}{0.020833in}}%
\pgfusepath{stroke,fill}%
}%
\begin{pgfscope}%
\pgfsys@transformshift{2.735834in}{0.420000in}%
\pgfsys@useobject{currentmarker}{}%
\end{pgfscope}%
\end{pgfscope}%
\begin{pgfscope}%
\pgfsetbuttcap%
\pgfsetroundjoin%
\definecolor{currentfill}{rgb}{0.000000,0.000000,0.000000}%
\pgfsetfillcolor{currentfill}%
\pgfsetlinewidth{0.401500pt}%
\definecolor{currentstroke}{rgb}{0.000000,0.000000,0.000000}%
\pgfsetstrokecolor{currentstroke}%
\pgfsetdash{}{0pt}%
\pgfsys@defobject{currentmarker}{\pgfqpoint{0.000000in}{-0.020833in}}{\pgfqpoint{0.000000in}{0.000000in}}{%
\pgfpathmoveto{\pgfqpoint{0.000000in}{0.000000in}}%
\pgfpathlineto{\pgfqpoint{0.000000in}{-0.020833in}}%
\pgfusepath{stroke,fill}%
}%
\begin{pgfscope}%
\pgfsys@transformshift{2.735834in}{2.414488in}%
\pgfsys@useobject{currentmarker}{}%
\end{pgfscope}%
\end{pgfscope}%
\begin{pgfscope}%
\pgfsetbuttcap%
\pgfsetroundjoin%
\definecolor{currentfill}{rgb}{0.000000,0.000000,0.000000}%
\pgfsetfillcolor{currentfill}%
\pgfsetlinewidth{0.401500pt}%
\definecolor{currentstroke}{rgb}{0.000000,0.000000,0.000000}%
\pgfsetstrokecolor{currentstroke}%
\pgfsetdash{}{0pt}%
\pgfsys@defobject{currentmarker}{\pgfqpoint{0.000000in}{0.000000in}}{\pgfqpoint{0.000000in}{0.020833in}}{%
\pgfpathmoveto{\pgfqpoint{0.000000in}{0.000000in}}%
\pgfpathlineto{\pgfqpoint{0.000000in}{0.020833in}}%
\pgfusepath{stroke,fill}%
}%
\begin{pgfscope}%
\pgfsys@transformshift{2.774282in}{0.420000in}%
\pgfsys@useobject{currentmarker}{}%
\end{pgfscope}%
\end{pgfscope}%
\begin{pgfscope}%
\pgfsetbuttcap%
\pgfsetroundjoin%
\definecolor{currentfill}{rgb}{0.000000,0.000000,0.000000}%
\pgfsetfillcolor{currentfill}%
\pgfsetlinewidth{0.401500pt}%
\definecolor{currentstroke}{rgb}{0.000000,0.000000,0.000000}%
\pgfsetstrokecolor{currentstroke}%
\pgfsetdash{}{0pt}%
\pgfsys@defobject{currentmarker}{\pgfqpoint{0.000000in}{-0.020833in}}{\pgfqpoint{0.000000in}{0.000000in}}{%
\pgfpathmoveto{\pgfqpoint{0.000000in}{0.000000in}}%
\pgfpathlineto{\pgfqpoint{0.000000in}{-0.020833in}}%
\pgfusepath{stroke,fill}%
}%
\begin{pgfscope}%
\pgfsys@transformshift{2.774282in}{2.414488in}%
\pgfsys@useobject{currentmarker}{}%
\end{pgfscope}%
\end{pgfscope}%
\begin{pgfscope}%
\pgfsetbuttcap%
\pgfsetroundjoin%
\definecolor{currentfill}{rgb}{0.000000,0.000000,0.000000}%
\pgfsetfillcolor{currentfill}%
\pgfsetlinewidth{0.401500pt}%
\definecolor{currentstroke}{rgb}{0.000000,0.000000,0.000000}%
\pgfsetstrokecolor{currentstroke}%
\pgfsetdash{}{0pt}%
\pgfsys@defobject{currentmarker}{\pgfqpoint{0.000000in}{0.000000in}}{\pgfqpoint{0.000000in}{0.020833in}}{%
\pgfpathmoveto{\pgfqpoint{0.000000in}{0.000000in}}%
\pgfpathlineto{\pgfqpoint{0.000000in}{0.020833in}}%
\pgfusepath{stroke,fill}%
}%
\begin{pgfscope}%
\pgfsys@transformshift{2.808195in}{0.420000in}%
\pgfsys@useobject{currentmarker}{}%
\end{pgfscope}%
\end{pgfscope}%
\begin{pgfscope}%
\pgfsetbuttcap%
\pgfsetroundjoin%
\definecolor{currentfill}{rgb}{0.000000,0.000000,0.000000}%
\pgfsetfillcolor{currentfill}%
\pgfsetlinewidth{0.401500pt}%
\definecolor{currentstroke}{rgb}{0.000000,0.000000,0.000000}%
\pgfsetstrokecolor{currentstroke}%
\pgfsetdash{}{0pt}%
\pgfsys@defobject{currentmarker}{\pgfqpoint{0.000000in}{-0.020833in}}{\pgfqpoint{0.000000in}{0.000000in}}{%
\pgfpathmoveto{\pgfqpoint{0.000000in}{0.000000in}}%
\pgfpathlineto{\pgfqpoint{0.000000in}{-0.020833in}}%
\pgfusepath{stroke,fill}%
}%
\begin{pgfscope}%
\pgfsys@transformshift{2.808195in}{2.414488in}%
\pgfsys@useobject{currentmarker}{}%
\end{pgfscope}%
\end{pgfscope}%
\begin{pgfscope}%
\pgfsetbuttcap%
\pgfsetroundjoin%
\definecolor{currentfill}{rgb}{0.000000,0.000000,0.000000}%
\pgfsetfillcolor{currentfill}%
\pgfsetlinewidth{0.401500pt}%
\definecolor{currentstroke}{rgb}{0.000000,0.000000,0.000000}%
\pgfsetstrokecolor{currentstroke}%
\pgfsetdash{}{0pt}%
\pgfsys@defobject{currentmarker}{\pgfqpoint{0.000000in}{0.000000in}}{\pgfqpoint{0.000000in}{0.020833in}}{%
\pgfpathmoveto{\pgfqpoint{0.000000in}{0.000000in}}%
\pgfpathlineto{\pgfqpoint{0.000000in}{0.020833in}}%
\pgfusepath{stroke,fill}%
}%
\begin{pgfscope}%
\pgfsys@transformshift{3.038110in}{0.420000in}%
\pgfsys@useobject{currentmarker}{}%
\end{pgfscope}%
\end{pgfscope}%
\begin{pgfscope}%
\pgfsetbuttcap%
\pgfsetroundjoin%
\definecolor{currentfill}{rgb}{0.000000,0.000000,0.000000}%
\pgfsetfillcolor{currentfill}%
\pgfsetlinewidth{0.401500pt}%
\definecolor{currentstroke}{rgb}{0.000000,0.000000,0.000000}%
\pgfsetstrokecolor{currentstroke}%
\pgfsetdash{}{0pt}%
\pgfsys@defobject{currentmarker}{\pgfqpoint{0.000000in}{-0.020833in}}{\pgfqpoint{0.000000in}{0.000000in}}{%
\pgfpathmoveto{\pgfqpoint{0.000000in}{0.000000in}}%
\pgfpathlineto{\pgfqpoint{0.000000in}{-0.020833in}}%
\pgfusepath{stroke,fill}%
}%
\begin{pgfscope}%
\pgfsys@transformshift{3.038110in}{2.414488in}%
\pgfsys@useobject{currentmarker}{}%
\end{pgfscope}%
\end{pgfscope}%
\begin{pgfscope}%
\definecolor{textcolor}{rgb}{0.000000,0.000000,0.000000}%
\pgfsetstrokecolor{textcolor}%
\pgfsetfillcolor{textcolor}%
\pgftext[x=1.844055in,y=0.208426in,,top]{\color{textcolor}{\rmfamily\fontsize{12.000000}{14.400000}\selectfont\catcode`\^=\active\def^{\ifmmode\sp\else\^{}\fi}\catcode`\%=\active\def%{\%}$y^{+}$}}%
\end{pgfscope}%
\begin{pgfscope}%
\pgfsetbuttcap%
\pgfsetroundjoin%
\definecolor{currentfill}{rgb}{0.000000,0.000000,0.000000}%
\pgfsetfillcolor{currentfill}%
\pgfsetlinewidth{0.501875pt}%
\definecolor{currentstroke}{rgb}{0.000000,0.000000,0.000000}%
\pgfsetstrokecolor{currentstroke}%
\pgfsetdash{}{0pt}%
\pgfsys@defobject{currentmarker}{\pgfqpoint{0.000000in}{0.000000in}}{\pgfqpoint{0.041667in}{0.000000in}}{%
\pgfpathmoveto{\pgfqpoint{0.000000in}{0.000000in}}%
\pgfpathlineto{\pgfqpoint{0.041667in}{0.000000in}}%
\pgfusepath{stroke,fill}%
}%
\begin{pgfscope}%
\pgfsys@transformshift{0.650000in}{0.496711in}%
\pgfsys@useobject{currentmarker}{}%
\end{pgfscope}%
\end{pgfscope}%
\begin{pgfscope}%
\pgfsetbuttcap%
\pgfsetroundjoin%
\definecolor{currentfill}{rgb}{0.000000,0.000000,0.000000}%
\pgfsetfillcolor{currentfill}%
\pgfsetlinewidth{0.501875pt}%
\definecolor{currentstroke}{rgb}{0.000000,0.000000,0.000000}%
\pgfsetstrokecolor{currentstroke}%
\pgfsetdash{}{0pt}%
\pgfsys@defobject{currentmarker}{\pgfqpoint{-0.041667in}{0.000000in}}{\pgfqpoint{-0.000000in}{0.000000in}}{%
\pgfpathmoveto{\pgfqpoint{-0.000000in}{0.000000in}}%
\pgfpathlineto{\pgfqpoint{-0.041667in}{0.000000in}}%
\pgfusepath{stroke,fill}%
}%
\begin{pgfscope}%
\pgfsys@transformshift{3.038110in}{0.496711in}%
\pgfsys@useobject{currentmarker}{}%
\end{pgfscope}%
\end{pgfscope}%
\begin{pgfscope}%
\definecolor{textcolor}{rgb}{0.000000,0.000000,0.000000}%
\pgfsetstrokecolor{textcolor}%
\pgfsetfillcolor{textcolor}%
\pgftext[x=0.288772in, y=0.443904in, left, base]{\color{textcolor}{\rmfamily\fontsize{11.000000}{13.200000}\selectfont\catcode`\^=\active\def^{\ifmmode\sp\else\^{}\fi}\catcode`\%=\active\def%{\%}\ensuremath{-}2.5}}%
\end{pgfscope}%
\begin{pgfscope}%
\pgfsetbuttcap%
\pgfsetroundjoin%
\definecolor{currentfill}{rgb}{0.000000,0.000000,0.000000}%
\pgfsetfillcolor{currentfill}%
\pgfsetlinewidth{0.501875pt}%
\definecolor{currentstroke}{rgb}{0.000000,0.000000,0.000000}%
\pgfsetstrokecolor{currentstroke}%
\pgfsetdash{}{0pt}%
\pgfsys@defobject{currentmarker}{\pgfqpoint{0.000000in}{0.000000in}}{\pgfqpoint{0.041667in}{0.000000in}}{%
\pgfpathmoveto{\pgfqpoint{0.000000in}{0.000000in}}%
\pgfpathlineto{\pgfqpoint{0.041667in}{0.000000in}}%
\pgfusepath{stroke,fill}%
}%
\begin{pgfscope}%
\pgfsys@transformshift{0.650000in}{0.880266in}%
\pgfsys@useobject{currentmarker}{}%
\end{pgfscope}%
\end{pgfscope}%
\begin{pgfscope}%
\pgfsetbuttcap%
\pgfsetroundjoin%
\definecolor{currentfill}{rgb}{0.000000,0.000000,0.000000}%
\pgfsetfillcolor{currentfill}%
\pgfsetlinewidth{0.501875pt}%
\definecolor{currentstroke}{rgb}{0.000000,0.000000,0.000000}%
\pgfsetstrokecolor{currentstroke}%
\pgfsetdash{}{0pt}%
\pgfsys@defobject{currentmarker}{\pgfqpoint{-0.041667in}{0.000000in}}{\pgfqpoint{-0.000000in}{0.000000in}}{%
\pgfpathmoveto{\pgfqpoint{-0.000000in}{0.000000in}}%
\pgfpathlineto{\pgfqpoint{-0.041667in}{0.000000in}}%
\pgfusepath{stroke,fill}%
}%
\begin{pgfscope}%
\pgfsys@transformshift{3.038110in}{0.880266in}%
\pgfsys@useobject{currentmarker}{}%
\end{pgfscope}%
\end{pgfscope}%
\begin{pgfscope}%
\definecolor{textcolor}{rgb}{0.000000,0.000000,0.000000}%
\pgfsetstrokecolor{textcolor}%
\pgfsetfillcolor{textcolor}%
\pgftext[x=0.407060in, y=0.827460in, left, base]{\color{textcolor}{\rmfamily\fontsize{11.000000}{13.200000}\selectfont\catcode`\^=\active\def^{\ifmmode\sp\else\^{}\fi}\catcode`\%=\active\def%{\%}0.0}}%
\end{pgfscope}%
\begin{pgfscope}%
\pgfsetbuttcap%
\pgfsetroundjoin%
\definecolor{currentfill}{rgb}{0.000000,0.000000,0.000000}%
\pgfsetfillcolor{currentfill}%
\pgfsetlinewidth{0.501875pt}%
\definecolor{currentstroke}{rgb}{0.000000,0.000000,0.000000}%
\pgfsetstrokecolor{currentstroke}%
\pgfsetdash{}{0pt}%
\pgfsys@defobject{currentmarker}{\pgfqpoint{0.000000in}{0.000000in}}{\pgfqpoint{0.041667in}{0.000000in}}{%
\pgfpathmoveto{\pgfqpoint{0.000000in}{0.000000in}}%
\pgfpathlineto{\pgfqpoint{0.041667in}{0.000000in}}%
\pgfusepath{stroke,fill}%
}%
\begin{pgfscope}%
\pgfsys@transformshift{0.650000in}{1.263822in}%
\pgfsys@useobject{currentmarker}{}%
\end{pgfscope}%
\end{pgfscope}%
\begin{pgfscope}%
\pgfsetbuttcap%
\pgfsetroundjoin%
\definecolor{currentfill}{rgb}{0.000000,0.000000,0.000000}%
\pgfsetfillcolor{currentfill}%
\pgfsetlinewidth{0.501875pt}%
\definecolor{currentstroke}{rgb}{0.000000,0.000000,0.000000}%
\pgfsetstrokecolor{currentstroke}%
\pgfsetdash{}{0pt}%
\pgfsys@defobject{currentmarker}{\pgfqpoint{-0.041667in}{0.000000in}}{\pgfqpoint{-0.000000in}{0.000000in}}{%
\pgfpathmoveto{\pgfqpoint{-0.000000in}{0.000000in}}%
\pgfpathlineto{\pgfqpoint{-0.041667in}{0.000000in}}%
\pgfusepath{stroke,fill}%
}%
\begin{pgfscope}%
\pgfsys@transformshift{3.038110in}{1.263822in}%
\pgfsys@useobject{currentmarker}{}%
\end{pgfscope}%
\end{pgfscope}%
\begin{pgfscope}%
\definecolor{textcolor}{rgb}{0.000000,0.000000,0.000000}%
\pgfsetstrokecolor{textcolor}%
\pgfsetfillcolor{textcolor}%
\pgftext[x=0.407060in, y=1.211015in, left, base]{\color{textcolor}{\rmfamily\fontsize{11.000000}{13.200000}\selectfont\catcode`\^=\active\def^{\ifmmode\sp\else\^{}\fi}\catcode`\%=\active\def%{\%}2.5}}%
\end{pgfscope}%
\begin{pgfscope}%
\pgfsetbuttcap%
\pgfsetroundjoin%
\definecolor{currentfill}{rgb}{0.000000,0.000000,0.000000}%
\pgfsetfillcolor{currentfill}%
\pgfsetlinewidth{0.501875pt}%
\definecolor{currentstroke}{rgb}{0.000000,0.000000,0.000000}%
\pgfsetstrokecolor{currentstroke}%
\pgfsetdash{}{0pt}%
\pgfsys@defobject{currentmarker}{\pgfqpoint{0.000000in}{0.000000in}}{\pgfqpoint{0.041667in}{0.000000in}}{%
\pgfpathmoveto{\pgfqpoint{0.000000in}{0.000000in}}%
\pgfpathlineto{\pgfqpoint{0.041667in}{0.000000in}}%
\pgfusepath{stroke,fill}%
}%
\begin{pgfscope}%
\pgfsys@transformshift{0.650000in}{1.647377in}%
\pgfsys@useobject{currentmarker}{}%
\end{pgfscope}%
\end{pgfscope}%
\begin{pgfscope}%
\pgfsetbuttcap%
\pgfsetroundjoin%
\definecolor{currentfill}{rgb}{0.000000,0.000000,0.000000}%
\pgfsetfillcolor{currentfill}%
\pgfsetlinewidth{0.501875pt}%
\definecolor{currentstroke}{rgb}{0.000000,0.000000,0.000000}%
\pgfsetstrokecolor{currentstroke}%
\pgfsetdash{}{0pt}%
\pgfsys@defobject{currentmarker}{\pgfqpoint{-0.041667in}{0.000000in}}{\pgfqpoint{-0.000000in}{0.000000in}}{%
\pgfpathmoveto{\pgfqpoint{-0.000000in}{0.000000in}}%
\pgfpathlineto{\pgfqpoint{-0.041667in}{0.000000in}}%
\pgfusepath{stroke,fill}%
}%
\begin{pgfscope}%
\pgfsys@transformshift{3.038110in}{1.647377in}%
\pgfsys@useobject{currentmarker}{}%
\end{pgfscope}%
\end{pgfscope}%
\begin{pgfscope}%
\definecolor{textcolor}{rgb}{0.000000,0.000000,0.000000}%
\pgfsetstrokecolor{textcolor}%
\pgfsetfillcolor{textcolor}%
\pgftext[x=0.407060in, y=1.594571in, left, base]{\color{textcolor}{\rmfamily\fontsize{11.000000}{13.200000}\selectfont\catcode`\^=\active\def^{\ifmmode\sp\else\^{}\fi}\catcode`\%=\active\def%{\%}5.0}}%
\end{pgfscope}%
\begin{pgfscope}%
\pgfsetbuttcap%
\pgfsetroundjoin%
\definecolor{currentfill}{rgb}{0.000000,0.000000,0.000000}%
\pgfsetfillcolor{currentfill}%
\pgfsetlinewidth{0.501875pt}%
\definecolor{currentstroke}{rgb}{0.000000,0.000000,0.000000}%
\pgfsetstrokecolor{currentstroke}%
\pgfsetdash{}{0pt}%
\pgfsys@defobject{currentmarker}{\pgfqpoint{0.000000in}{0.000000in}}{\pgfqpoint{0.041667in}{0.000000in}}{%
\pgfpathmoveto{\pgfqpoint{0.000000in}{0.000000in}}%
\pgfpathlineto{\pgfqpoint{0.041667in}{0.000000in}}%
\pgfusepath{stroke,fill}%
}%
\begin{pgfscope}%
\pgfsys@transformshift{0.650000in}{2.030933in}%
\pgfsys@useobject{currentmarker}{}%
\end{pgfscope}%
\end{pgfscope}%
\begin{pgfscope}%
\pgfsetbuttcap%
\pgfsetroundjoin%
\definecolor{currentfill}{rgb}{0.000000,0.000000,0.000000}%
\pgfsetfillcolor{currentfill}%
\pgfsetlinewidth{0.501875pt}%
\definecolor{currentstroke}{rgb}{0.000000,0.000000,0.000000}%
\pgfsetstrokecolor{currentstroke}%
\pgfsetdash{}{0pt}%
\pgfsys@defobject{currentmarker}{\pgfqpoint{-0.041667in}{0.000000in}}{\pgfqpoint{-0.000000in}{0.000000in}}{%
\pgfpathmoveto{\pgfqpoint{-0.000000in}{0.000000in}}%
\pgfpathlineto{\pgfqpoint{-0.041667in}{0.000000in}}%
\pgfusepath{stroke,fill}%
}%
\begin{pgfscope}%
\pgfsys@transformshift{3.038110in}{2.030933in}%
\pgfsys@useobject{currentmarker}{}%
\end{pgfscope}%
\end{pgfscope}%
\begin{pgfscope}%
\definecolor{textcolor}{rgb}{0.000000,0.000000,0.000000}%
\pgfsetstrokecolor{textcolor}%
\pgfsetfillcolor{textcolor}%
\pgftext[x=0.407060in, y=1.978126in, left, base]{\color{textcolor}{\rmfamily\fontsize{11.000000}{13.200000}\selectfont\catcode`\^=\active\def^{\ifmmode\sp\else\^{}\fi}\catcode`\%=\active\def%{\%}7.5}}%
\end{pgfscope}%
\begin{pgfscope}%
\pgfsetbuttcap%
\pgfsetroundjoin%
\definecolor{currentfill}{rgb}{0.000000,0.000000,0.000000}%
\pgfsetfillcolor{currentfill}%
\pgfsetlinewidth{0.501875pt}%
\definecolor{currentstroke}{rgb}{0.000000,0.000000,0.000000}%
\pgfsetstrokecolor{currentstroke}%
\pgfsetdash{}{0pt}%
\pgfsys@defobject{currentmarker}{\pgfqpoint{0.000000in}{0.000000in}}{\pgfqpoint{0.041667in}{0.000000in}}{%
\pgfpathmoveto{\pgfqpoint{0.000000in}{0.000000in}}%
\pgfpathlineto{\pgfqpoint{0.041667in}{0.000000in}}%
\pgfusepath{stroke,fill}%
}%
\begin{pgfscope}%
\pgfsys@transformshift{0.650000in}{2.414488in}%
\pgfsys@useobject{currentmarker}{}%
\end{pgfscope}%
\end{pgfscope}%
\begin{pgfscope}%
\pgfsetbuttcap%
\pgfsetroundjoin%
\definecolor{currentfill}{rgb}{0.000000,0.000000,0.000000}%
\pgfsetfillcolor{currentfill}%
\pgfsetlinewidth{0.501875pt}%
\definecolor{currentstroke}{rgb}{0.000000,0.000000,0.000000}%
\pgfsetstrokecolor{currentstroke}%
\pgfsetdash{}{0pt}%
\pgfsys@defobject{currentmarker}{\pgfqpoint{-0.041667in}{0.000000in}}{\pgfqpoint{-0.000000in}{0.000000in}}{%
\pgfpathmoveto{\pgfqpoint{-0.000000in}{0.000000in}}%
\pgfpathlineto{\pgfqpoint{-0.041667in}{0.000000in}}%
\pgfusepath{stroke,fill}%
}%
\begin{pgfscope}%
\pgfsys@transformshift{3.038110in}{2.414488in}%
\pgfsys@useobject{currentmarker}{}%
\end{pgfscope}%
\end{pgfscope}%
\begin{pgfscope}%
\definecolor{textcolor}{rgb}{0.000000,0.000000,0.000000}%
\pgfsetstrokecolor{textcolor}%
\pgfsetfillcolor{textcolor}%
\pgftext[x=0.331018in, y=2.361681in, left, base]{\color{textcolor}{\rmfamily\fontsize{11.000000}{13.200000}\selectfont\catcode`\^=\active\def^{\ifmmode\sp\else\^{}\fi}\catcode`\%=\active\def%{\%}10.0}}%
\end{pgfscope}%
\begin{pgfscope}%
\pgfsetbuttcap%
\pgfsetroundjoin%
\definecolor{currentfill}{rgb}{0.000000,0.000000,0.000000}%
\pgfsetfillcolor{currentfill}%
\pgfsetlinewidth{0.401500pt}%
\definecolor{currentstroke}{rgb}{0.000000,0.000000,0.000000}%
\pgfsetstrokecolor{currentstroke}%
\pgfsetdash{}{0pt}%
\pgfsys@defobject{currentmarker}{\pgfqpoint{0.000000in}{0.000000in}}{\pgfqpoint{0.020833in}{0.000000in}}{%
\pgfpathmoveto{\pgfqpoint{0.000000in}{0.000000in}}%
\pgfpathlineto{\pgfqpoint{0.020833in}{0.000000in}}%
\pgfusepath{stroke,fill}%
}%
\begin{pgfscope}%
\pgfsys@transformshift{0.650000in}{0.420000in}%
\pgfsys@useobject{currentmarker}{}%
\end{pgfscope}%
\end{pgfscope}%
\begin{pgfscope}%
\pgfsetbuttcap%
\pgfsetroundjoin%
\definecolor{currentfill}{rgb}{0.000000,0.000000,0.000000}%
\pgfsetfillcolor{currentfill}%
\pgfsetlinewidth{0.401500pt}%
\definecolor{currentstroke}{rgb}{0.000000,0.000000,0.000000}%
\pgfsetstrokecolor{currentstroke}%
\pgfsetdash{}{0pt}%
\pgfsys@defobject{currentmarker}{\pgfqpoint{-0.020833in}{0.000000in}}{\pgfqpoint{-0.000000in}{0.000000in}}{%
\pgfpathmoveto{\pgfqpoint{-0.000000in}{0.000000in}}%
\pgfpathlineto{\pgfqpoint{-0.020833in}{0.000000in}}%
\pgfusepath{stroke,fill}%
}%
\begin{pgfscope}%
\pgfsys@transformshift{3.038110in}{0.420000in}%
\pgfsys@useobject{currentmarker}{}%
\end{pgfscope}%
\end{pgfscope}%
\begin{pgfscope}%
\pgfsetbuttcap%
\pgfsetroundjoin%
\definecolor{currentfill}{rgb}{0.000000,0.000000,0.000000}%
\pgfsetfillcolor{currentfill}%
\pgfsetlinewidth{0.401500pt}%
\definecolor{currentstroke}{rgb}{0.000000,0.000000,0.000000}%
\pgfsetstrokecolor{currentstroke}%
\pgfsetdash{}{0pt}%
\pgfsys@defobject{currentmarker}{\pgfqpoint{0.000000in}{0.000000in}}{\pgfqpoint{0.020833in}{0.000000in}}{%
\pgfpathmoveto{\pgfqpoint{0.000000in}{0.000000in}}%
\pgfpathlineto{\pgfqpoint{0.020833in}{0.000000in}}%
\pgfusepath{stroke,fill}%
}%
\begin{pgfscope}%
\pgfsys@transformshift{0.650000in}{0.573422in}%
\pgfsys@useobject{currentmarker}{}%
\end{pgfscope}%
\end{pgfscope}%
\begin{pgfscope}%
\pgfsetbuttcap%
\pgfsetroundjoin%
\definecolor{currentfill}{rgb}{0.000000,0.000000,0.000000}%
\pgfsetfillcolor{currentfill}%
\pgfsetlinewidth{0.401500pt}%
\definecolor{currentstroke}{rgb}{0.000000,0.000000,0.000000}%
\pgfsetstrokecolor{currentstroke}%
\pgfsetdash{}{0pt}%
\pgfsys@defobject{currentmarker}{\pgfqpoint{-0.020833in}{0.000000in}}{\pgfqpoint{-0.000000in}{0.000000in}}{%
\pgfpathmoveto{\pgfqpoint{-0.000000in}{0.000000in}}%
\pgfpathlineto{\pgfqpoint{-0.020833in}{0.000000in}}%
\pgfusepath{stroke,fill}%
}%
\begin{pgfscope}%
\pgfsys@transformshift{3.038110in}{0.573422in}%
\pgfsys@useobject{currentmarker}{}%
\end{pgfscope}%
\end{pgfscope}%
\begin{pgfscope}%
\pgfsetbuttcap%
\pgfsetroundjoin%
\definecolor{currentfill}{rgb}{0.000000,0.000000,0.000000}%
\pgfsetfillcolor{currentfill}%
\pgfsetlinewidth{0.401500pt}%
\definecolor{currentstroke}{rgb}{0.000000,0.000000,0.000000}%
\pgfsetstrokecolor{currentstroke}%
\pgfsetdash{}{0pt}%
\pgfsys@defobject{currentmarker}{\pgfqpoint{0.000000in}{0.000000in}}{\pgfqpoint{0.020833in}{0.000000in}}{%
\pgfpathmoveto{\pgfqpoint{0.000000in}{0.000000in}}%
\pgfpathlineto{\pgfqpoint{0.020833in}{0.000000in}}%
\pgfusepath{stroke,fill}%
}%
\begin{pgfscope}%
\pgfsys@transformshift{0.650000in}{0.650133in}%
\pgfsys@useobject{currentmarker}{}%
\end{pgfscope}%
\end{pgfscope}%
\begin{pgfscope}%
\pgfsetbuttcap%
\pgfsetroundjoin%
\definecolor{currentfill}{rgb}{0.000000,0.000000,0.000000}%
\pgfsetfillcolor{currentfill}%
\pgfsetlinewidth{0.401500pt}%
\definecolor{currentstroke}{rgb}{0.000000,0.000000,0.000000}%
\pgfsetstrokecolor{currentstroke}%
\pgfsetdash{}{0pt}%
\pgfsys@defobject{currentmarker}{\pgfqpoint{-0.020833in}{0.000000in}}{\pgfqpoint{-0.000000in}{0.000000in}}{%
\pgfpathmoveto{\pgfqpoint{-0.000000in}{0.000000in}}%
\pgfpathlineto{\pgfqpoint{-0.020833in}{0.000000in}}%
\pgfusepath{stroke,fill}%
}%
\begin{pgfscope}%
\pgfsys@transformshift{3.038110in}{0.650133in}%
\pgfsys@useobject{currentmarker}{}%
\end{pgfscope}%
\end{pgfscope}%
\begin{pgfscope}%
\pgfsetbuttcap%
\pgfsetroundjoin%
\definecolor{currentfill}{rgb}{0.000000,0.000000,0.000000}%
\pgfsetfillcolor{currentfill}%
\pgfsetlinewidth{0.401500pt}%
\definecolor{currentstroke}{rgb}{0.000000,0.000000,0.000000}%
\pgfsetstrokecolor{currentstroke}%
\pgfsetdash{}{0pt}%
\pgfsys@defobject{currentmarker}{\pgfqpoint{0.000000in}{0.000000in}}{\pgfqpoint{0.020833in}{0.000000in}}{%
\pgfpathmoveto{\pgfqpoint{0.000000in}{0.000000in}}%
\pgfpathlineto{\pgfqpoint{0.020833in}{0.000000in}}%
\pgfusepath{stroke,fill}%
}%
\begin{pgfscope}%
\pgfsys@transformshift{0.650000in}{0.726844in}%
\pgfsys@useobject{currentmarker}{}%
\end{pgfscope}%
\end{pgfscope}%
\begin{pgfscope}%
\pgfsetbuttcap%
\pgfsetroundjoin%
\definecolor{currentfill}{rgb}{0.000000,0.000000,0.000000}%
\pgfsetfillcolor{currentfill}%
\pgfsetlinewidth{0.401500pt}%
\definecolor{currentstroke}{rgb}{0.000000,0.000000,0.000000}%
\pgfsetstrokecolor{currentstroke}%
\pgfsetdash{}{0pt}%
\pgfsys@defobject{currentmarker}{\pgfqpoint{-0.020833in}{0.000000in}}{\pgfqpoint{-0.000000in}{0.000000in}}{%
\pgfpathmoveto{\pgfqpoint{-0.000000in}{0.000000in}}%
\pgfpathlineto{\pgfqpoint{-0.020833in}{0.000000in}}%
\pgfusepath{stroke,fill}%
}%
\begin{pgfscope}%
\pgfsys@transformshift{3.038110in}{0.726844in}%
\pgfsys@useobject{currentmarker}{}%
\end{pgfscope}%
\end{pgfscope}%
\begin{pgfscope}%
\pgfsetbuttcap%
\pgfsetroundjoin%
\definecolor{currentfill}{rgb}{0.000000,0.000000,0.000000}%
\pgfsetfillcolor{currentfill}%
\pgfsetlinewidth{0.401500pt}%
\definecolor{currentstroke}{rgb}{0.000000,0.000000,0.000000}%
\pgfsetstrokecolor{currentstroke}%
\pgfsetdash{}{0pt}%
\pgfsys@defobject{currentmarker}{\pgfqpoint{0.000000in}{0.000000in}}{\pgfqpoint{0.020833in}{0.000000in}}{%
\pgfpathmoveto{\pgfqpoint{0.000000in}{0.000000in}}%
\pgfpathlineto{\pgfqpoint{0.020833in}{0.000000in}}%
\pgfusepath{stroke,fill}%
}%
\begin{pgfscope}%
\pgfsys@transformshift{0.650000in}{0.803555in}%
\pgfsys@useobject{currentmarker}{}%
\end{pgfscope}%
\end{pgfscope}%
\begin{pgfscope}%
\pgfsetbuttcap%
\pgfsetroundjoin%
\definecolor{currentfill}{rgb}{0.000000,0.000000,0.000000}%
\pgfsetfillcolor{currentfill}%
\pgfsetlinewidth{0.401500pt}%
\definecolor{currentstroke}{rgb}{0.000000,0.000000,0.000000}%
\pgfsetstrokecolor{currentstroke}%
\pgfsetdash{}{0pt}%
\pgfsys@defobject{currentmarker}{\pgfqpoint{-0.020833in}{0.000000in}}{\pgfqpoint{-0.000000in}{0.000000in}}{%
\pgfpathmoveto{\pgfqpoint{-0.000000in}{0.000000in}}%
\pgfpathlineto{\pgfqpoint{-0.020833in}{0.000000in}}%
\pgfusepath{stroke,fill}%
}%
\begin{pgfscope}%
\pgfsys@transformshift{3.038110in}{0.803555in}%
\pgfsys@useobject{currentmarker}{}%
\end{pgfscope}%
\end{pgfscope}%
\begin{pgfscope}%
\pgfsetbuttcap%
\pgfsetroundjoin%
\definecolor{currentfill}{rgb}{0.000000,0.000000,0.000000}%
\pgfsetfillcolor{currentfill}%
\pgfsetlinewidth{0.401500pt}%
\definecolor{currentstroke}{rgb}{0.000000,0.000000,0.000000}%
\pgfsetstrokecolor{currentstroke}%
\pgfsetdash{}{0pt}%
\pgfsys@defobject{currentmarker}{\pgfqpoint{0.000000in}{0.000000in}}{\pgfqpoint{0.020833in}{0.000000in}}{%
\pgfpathmoveto{\pgfqpoint{0.000000in}{0.000000in}}%
\pgfpathlineto{\pgfqpoint{0.020833in}{0.000000in}}%
\pgfusepath{stroke,fill}%
}%
\begin{pgfscope}%
\pgfsys@transformshift{0.650000in}{0.956978in}%
\pgfsys@useobject{currentmarker}{}%
\end{pgfscope}%
\end{pgfscope}%
\begin{pgfscope}%
\pgfsetbuttcap%
\pgfsetroundjoin%
\definecolor{currentfill}{rgb}{0.000000,0.000000,0.000000}%
\pgfsetfillcolor{currentfill}%
\pgfsetlinewidth{0.401500pt}%
\definecolor{currentstroke}{rgb}{0.000000,0.000000,0.000000}%
\pgfsetstrokecolor{currentstroke}%
\pgfsetdash{}{0pt}%
\pgfsys@defobject{currentmarker}{\pgfqpoint{-0.020833in}{0.000000in}}{\pgfqpoint{-0.000000in}{0.000000in}}{%
\pgfpathmoveto{\pgfqpoint{-0.000000in}{0.000000in}}%
\pgfpathlineto{\pgfqpoint{-0.020833in}{0.000000in}}%
\pgfusepath{stroke,fill}%
}%
\begin{pgfscope}%
\pgfsys@transformshift{3.038110in}{0.956978in}%
\pgfsys@useobject{currentmarker}{}%
\end{pgfscope}%
\end{pgfscope}%
\begin{pgfscope}%
\pgfsetbuttcap%
\pgfsetroundjoin%
\definecolor{currentfill}{rgb}{0.000000,0.000000,0.000000}%
\pgfsetfillcolor{currentfill}%
\pgfsetlinewidth{0.401500pt}%
\definecolor{currentstroke}{rgb}{0.000000,0.000000,0.000000}%
\pgfsetstrokecolor{currentstroke}%
\pgfsetdash{}{0pt}%
\pgfsys@defobject{currentmarker}{\pgfqpoint{0.000000in}{0.000000in}}{\pgfqpoint{0.020833in}{0.000000in}}{%
\pgfpathmoveto{\pgfqpoint{0.000000in}{0.000000in}}%
\pgfpathlineto{\pgfqpoint{0.020833in}{0.000000in}}%
\pgfusepath{stroke,fill}%
}%
\begin{pgfscope}%
\pgfsys@transformshift{0.650000in}{1.033689in}%
\pgfsys@useobject{currentmarker}{}%
\end{pgfscope}%
\end{pgfscope}%
\begin{pgfscope}%
\pgfsetbuttcap%
\pgfsetroundjoin%
\definecolor{currentfill}{rgb}{0.000000,0.000000,0.000000}%
\pgfsetfillcolor{currentfill}%
\pgfsetlinewidth{0.401500pt}%
\definecolor{currentstroke}{rgb}{0.000000,0.000000,0.000000}%
\pgfsetstrokecolor{currentstroke}%
\pgfsetdash{}{0pt}%
\pgfsys@defobject{currentmarker}{\pgfqpoint{-0.020833in}{0.000000in}}{\pgfqpoint{-0.000000in}{0.000000in}}{%
\pgfpathmoveto{\pgfqpoint{-0.000000in}{0.000000in}}%
\pgfpathlineto{\pgfqpoint{-0.020833in}{0.000000in}}%
\pgfusepath{stroke,fill}%
}%
\begin{pgfscope}%
\pgfsys@transformshift{3.038110in}{1.033689in}%
\pgfsys@useobject{currentmarker}{}%
\end{pgfscope}%
\end{pgfscope}%
\begin{pgfscope}%
\pgfsetbuttcap%
\pgfsetroundjoin%
\definecolor{currentfill}{rgb}{0.000000,0.000000,0.000000}%
\pgfsetfillcolor{currentfill}%
\pgfsetlinewidth{0.401500pt}%
\definecolor{currentstroke}{rgb}{0.000000,0.000000,0.000000}%
\pgfsetstrokecolor{currentstroke}%
\pgfsetdash{}{0pt}%
\pgfsys@defobject{currentmarker}{\pgfqpoint{0.000000in}{0.000000in}}{\pgfqpoint{0.020833in}{0.000000in}}{%
\pgfpathmoveto{\pgfqpoint{0.000000in}{0.000000in}}%
\pgfpathlineto{\pgfqpoint{0.020833in}{0.000000in}}%
\pgfusepath{stroke,fill}%
}%
\begin{pgfscope}%
\pgfsys@transformshift{0.650000in}{1.110400in}%
\pgfsys@useobject{currentmarker}{}%
\end{pgfscope}%
\end{pgfscope}%
\begin{pgfscope}%
\pgfsetbuttcap%
\pgfsetroundjoin%
\definecolor{currentfill}{rgb}{0.000000,0.000000,0.000000}%
\pgfsetfillcolor{currentfill}%
\pgfsetlinewidth{0.401500pt}%
\definecolor{currentstroke}{rgb}{0.000000,0.000000,0.000000}%
\pgfsetstrokecolor{currentstroke}%
\pgfsetdash{}{0pt}%
\pgfsys@defobject{currentmarker}{\pgfqpoint{-0.020833in}{0.000000in}}{\pgfqpoint{-0.000000in}{0.000000in}}{%
\pgfpathmoveto{\pgfqpoint{-0.000000in}{0.000000in}}%
\pgfpathlineto{\pgfqpoint{-0.020833in}{0.000000in}}%
\pgfusepath{stroke,fill}%
}%
\begin{pgfscope}%
\pgfsys@transformshift{3.038110in}{1.110400in}%
\pgfsys@useobject{currentmarker}{}%
\end{pgfscope}%
\end{pgfscope}%
\begin{pgfscope}%
\pgfsetbuttcap%
\pgfsetroundjoin%
\definecolor{currentfill}{rgb}{0.000000,0.000000,0.000000}%
\pgfsetfillcolor{currentfill}%
\pgfsetlinewidth{0.401500pt}%
\definecolor{currentstroke}{rgb}{0.000000,0.000000,0.000000}%
\pgfsetstrokecolor{currentstroke}%
\pgfsetdash{}{0pt}%
\pgfsys@defobject{currentmarker}{\pgfqpoint{0.000000in}{0.000000in}}{\pgfqpoint{0.020833in}{0.000000in}}{%
\pgfpathmoveto{\pgfqpoint{0.000000in}{0.000000in}}%
\pgfpathlineto{\pgfqpoint{0.020833in}{0.000000in}}%
\pgfusepath{stroke,fill}%
}%
\begin{pgfscope}%
\pgfsys@transformshift{0.650000in}{1.187111in}%
\pgfsys@useobject{currentmarker}{}%
\end{pgfscope}%
\end{pgfscope}%
\begin{pgfscope}%
\pgfsetbuttcap%
\pgfsetroundjoin%
\definecolor{currentfill}{rgb}{0.000000,0.000000,0.000000}%
\pgfsetfillcolor{currentfill}%
\pgfsetlinewidth{0.401500pt}%
\definecolor{currentstroke}{rgb}{0.000000,0.000000,0.000000}%
\pgfsetstrokecolor{currentstroke}%
\pgfsetdash{}{0pt}%
\pgfsys@defobject{currentmarker}{\pgfqpoint{-0.020833in}{0.000000in}}{\pgfqpoint{-0.000000in}{0.000000in}}{%
\pgfpathmoveto{\pgfqpoint{-0.000000in}{0.000000in}}%
\pgfpathlineto{\pgfqpoint{-0.020833in}{0.000000in}}%
\pgfusepath{stroke,fill}%
}%
\begin{pgfscope}%
\pgfsys@transformshift{3.038110in}{1.187111in}%
\pgfsys@useobject{currentmarker}{}%
\end{pgfscope}%
\end{pgfscope}%
\begin{pgfscope}%
\pgfsetbuttcap%
\pgfsetroundjoin%
\definecolor{currentfill}{rgb}{0.000000,0.000000,0.000000}%
\pgfsetfillcolor{currentfill}%
\pgfsetlinewidth{0.401500pt}%
\definecolor{currentstroke}{rgb}{0.000000,0.000000,0.000000}%
\pgfsetstrokecolor{currentstroke}%
\pgfsetdash{}{0pt}%
\pgfsys@defobject{currentmarker}{\pgfqpoint{0.000000in}{0.000000in}}{\pgfqpoint{0.020833in}{0.000000in}}{%
\pgfpathmoveto{\pgfqpoint{0.000000in}{0.000000in}}%
\pgfpathlineto{\pgfqpoint{0.020833in}{0.000000in}}%
\pgfusepath{stroke,fill}%
}%
\begin{pgfscope}%
\pgfsys@transformshift{0.650000in}{1.340533in}%
\pgfsys@useobject{currentmarker}{}%
\end{pgfscope}%
\end{pgfscope}%
\begin{pgfscope}%
\pgfsetbuttcap%
\pgfsetroundjoin%
\definecolor{currentfill}{rgb}{0.000000,0.000000,0.000000}%
\pgfsetfillcolor{currentfill}%
\pgfsetlinewidth{0.401500pt}%
\definecolor{currentstroke}{rgb}{0.000000,0.000000,0.000000}%
\pgfsetstrokecolor{currentstroke}%
\pgfsetdash{}{0pt}%
\pgfsys@defobject{currentmarker}{\pgfqpoint{-0.020833in}{0.000000in}}{\pgfqpoint{-0.000000in}{0.000000in}}{%
\pgfpathmoveto{\pgfqpoint{-0.000000in}{0.000000in}}%
\pgfpathlineto{\pgfqpoint{-0.020833in}{0.000000in}}%
\pgfusepath{stroke,fill}%
}%
\begin{pgfscope}%
\pgfsys@transformshift{3.038110in}{1.340533in}%
\pgfsys@useobject{currentmarker}{}%
\end{pgfscope}%
\end{pgfscope}%
\begin{pgfscope}%
\pgfsetbuttcap%
\pgfsetroundjoin%
\definecolor{currentfill}{rgb}{0.000000,0.000000,0.000000}%
\pgfsetfillcolor{currentfill}%
\pgfsetlinewidth{0.401500pt}%
\definecolor{currentstroke}{rgb}{0.000000,0.000000,0.000000}%
\pgfsetstrokecolor{currentstroke}%
\pgfsetdash{}{0pt}%
\pgfsys@defobject{currentmarker}{\pgfqpoint{0.000000in}{0.000000in}}{\pgfqpoint{0.020833in}{0.000000in}}{%
\pgfpathmoveto{\pgfqpoint{0.000000in}{0.000000in}}%
\pgfpathlineto{\pgfqpoint{0.020833in}{0.000000in}}%
\pgfusepath{stroke,fill}%
}%
\begin{pgfscope}%
\pgfsys@transformshift{0.650000in}{1.417244in}%
\pgfsys@useobject{currentmarker}{}%
\end{pgfscope}%
\end{pgfscope}%
\begin{pgfscope}%
\pgfsetbuttcap%
\pgfsetroundjoin%
\definecolor{currentfill}{rgb}{0.000000,0.000000,0.000000}%
\pgfsetfillcolor{currentfill}%
\pgfsetlinewidth{0.401500pt}%
\definecolor{currentstroke}{rgb}{0.000000,0.000000,0.000000}%
\pgfsetstrokecolor{currentstroke}%
\pgfsetdash{}{0pt}%
\pgfsys@defobject{currentmarker}{\pgfqpoint{-0.020833in}{0.000000in}}{\pgfqpoint{-0.000000in}{0.000000in}}{%
\pgfpathmoveto{\pgfqpoint{-0.000000in}{0.000000in}}%
\pgfpathlineto{\pgfqpoint{-0.020833in}{0.000000in}}%
\pgfusepath{stroke,fill}%
}%
\begin{pgfscope}%
\pgfsys@transformshift{3.038110in}{1.417244in}%
\pgfsys@useobject{currentmarker}{}%
\end{pgfscope}%
\end{pgfscope}%
\begin{pgfscope}%
\pgfsetbuttcap%
\pgfsetroundjoin%
\definecolor{currentfill}{rgb}{0.000000,0.000000,0.000000}%
\pgfsetfillcolor{currentfill}%
\pgfsetlinewidth{0.401500pt}%
\definecolor{currentstroke}{rgb}{0.000000,0.000000,0.000000}%
\pgfsetstrokecolor{currentstroke}%
\pgfsetdash{}{0pt}%
\pgfsys@defobject{currentmarker}{\pgfqpoint{0.000000in}{0.000000in}}{\pgfqpoint{0.020833in}{0.000000in}}{%
\pgfpathmoveto{\pgfqpoint{0.000000in}{0.000000in}}%
\pgfpathlineto{\pgfqpoint{0.020833in}{0.000000in}}%
\pgfusepath{stroke,fill}%
}%
\begin{pgfscope}%
\pgfsys@transformshift{0.650000in}{1.493955in}%
\pgfsys@useobject{currentmarker}{}%
\end{pgfscope}%
\end{pgfscope}%
\begin{pgfscope}%
\pgfsetbuttcap%
\pgfsetroundjoin%
\definecolor{currentfill}{rgb}{0.000000,0.000000,0.000000}%
\pgfsetfillcolor{currentfill}%
\pgfsetlinewidth{0.401500pt}%
\definecolor{currentstroke}{rgb}{0.000000,0.000000,0.000000}%
\pgfsetstrokecolor{currentstroke}%
\pgfsetdash{}{0pt}%
\pgfsys@defobject{currentmarker}{\pgfqpoint{-0.020833in}{0.000000in}}{\pgfqpoint{-0.000000in}{0.000000in}}{%
\pgfpathmoveto{\pgfqpoint{-0.000000in}{0.000000in}}%
\pgfpathlineto{\pgfqpoint{-0.020833in}{0.000000in}}%
\pgfusepath{stroke,fill}%
}%
\begin{pgfscope}%
\pgfsys@transformshift{3.038110in}{1.493955in}%
\pgfsys@useobject{currentmarker}{}%
\end{pgfscope}%
\end{pgfscope}%
\begin{pgfscope}%
\pgfsetbuttcap%
\pgfsetroundjoin%
\definecolor{currentfill}{rgb}{0.000000,0.000000,0.000000}%
\pgfsetfillcolor{currentfill}%
\pgfsetlinewidth{0.401500pt}%
\definecolor{currentstroke}{rgb}{0.000000,0.000000,0.000000}%
\pgfsetstrokecolor{currentstroke}%
\pgfsetdash{}{0pt}%
\pgfsys@defobject{currentmarker}{\pgfqpoint{0.000000in}{0.000000in}}{\pgfqpoint{0.020833in}{0.000000in}}{%
\pgfpathmoveto{\pgfqpoint{0.000000in}{0.000000in}}%
\pgfpathlineto{\pgfqpoint{0.020833in}{0.000000in}}%
\pgfusepath{stroke,fill}%
}%
\begin{pgfscope}%
\pgfsys@transformshift{0.650000in}{1.570666in}%
\pgfsys@useobject{currentmarker}{}%
\end{pgfscope}%
\end{pgfscope}%
\begin{pgfscope}%
\pgfsetbuttcap%
\pgfsetroundjoin%
\definecolor{currentfill}{rgb}{0.000000,0.000000,0.000000}%
\pgfsetfillcolor{currentfill}%
\pgfsetlinewidth{0.401500pt}%
\definecolor{currentstroke}{rgb}{0.000000,0.000000,0.000000}%
\pgfsetstrokecolor{currentstroke}%
\pgfsetdash{}{0pt}%
\pgfsys@defobject{currentmarker}{\pgfqpoint{-0.020833in}{0.000000in}}{\pgfqpoint{-0.000000in}{0.000000in}}{%
\pgfpathmoveto{\pgfqpoint{-0.000000in}{0.000000in}}%
\pgfpathlineto{\pgfqpoint{-0.020833in}{0.000000in}}%
\pgfusepath{stroke,fill}%
}%
\begin{pgfscope}%
\pgfsys@transformshift{3.038110in}{1.570666in}%
\pgfsys@useobject{currentmarker}{}%
\end{pgfscope}%
\end{pgfscope}%
\begin{pgfscope}%
\pgfsetbuttcap%
\pgfsetroundjoin%
\definecolor{currentfill}{rgb}{0.000000,0.000000,0.000000}%
\pgfsetfillcolor{currentfill}%
\pgfsetlinewidth{0.401500pt}%
\definecolor{currentstroke}{rgb}{0.000000,0.000000,0.000000}%
\pgfsetstrokecolor{currentstroke}%
\pgfsetdash{}{0pt}%
\pgfsys@defobject{currentmarker}{\pgfqpoint{0.000000in}{0.000000in}}{\pgfqpoint{0.020833in}{0.000000in}}{%
\pgfpathmoveto{\pgfqpoint{0.000000in}{0.000000in}}%
\pgfpathlineto{\pgfqpoint{0.020833in}{0.000000in}}%
\pgfusepath{stroke,fill}%
}%
\begin{pgfscope}%
\pgfsys@transformshift{0.650000in}{1.724088in}%
\pgfsys@useobject{currentmarker}{}%
\end{pgfscope}%
\end{pgfscope}%
\begin{pgfscope}%
\pgfsetbuttcap%
\pgfsetroundjoin%
\definecolor{currentfill}{rgb}{0.000000,0.000000,0.000000}%
\pgfsetfillcolor{currentfill}%
\pgfsetlinewidth{0.401500pt}%
\definecolor{currentstroke}{rgb}{0.000000,0.000000,0.000000}%
\pgfsetstrokecolor{currentstroke}%
\pgfsetdash{}{0pt}%
\pgfsys@defobject{currentmarker}{\pgfqpoint{-0.020833in}{0.000000in}}{\pgfqpoint{-0.000000in}{0.000000in}}{%
\pgfpathmoveto{\pgfqpoint{-0.000000in}{0.000000in}}%
\pgfpathlineto{\pgfqpoint{-0.020833in}{0.000000in}}%
\pgfusepath{stroke,fill}%
}%
\begin{pgfscope}%
\pgfsys@transformshift{3.038110in}{1.724088in}%
\pgfsys@useobject{currentmarker}{}%
\end{pgfscope}%
\end{pgfscope}%
\begin{pgfscope}%
\pgfsetbuttcap%
\pgfsetroundjoin%
\definecolor{currentfill}{rgb}{0.000000,0.000000,0.000000}%
\pgfsetfillcolor{currentfill}%
\pgfsetlinewidth{0.401500pt}%
\definecolor{currentstroke}{rgb}{0.000000,0.000000,0.000000}%
\pgfsetstrokecolor{currentstroke}%
\pgfsetdash{}{0pt}%
\pgfsys@defobject{currentmarker}{\pgfqpoint{0.000000in}{0.000000in}}{\pgfqpoint{0.020833in}{0.000000in}}{%
\pgfpathmoveto{\pgfqpoint{0.000000in}{0.000000in}}%
\pgfpathlineto{\pgfqpoint{0.020833in}{0.000000in}}%
\pgfusepath{stroke,fill}%
}%
\begin{pgfscope}%
\pgfsys@transformshift{0.650000in}{1.800799in}%
\pgfsys@useobject{currentmarker}{}%
\end{pgfscope}%
\end{pgfscope}%
\begin{pgfscope}%
\pgfsetbuttcap%
\pgfsetroundjoin%
\definecolor{currentfill}{rgb}{0.000000,0.000000,0.000000}%
\pgfsetfillcolor{currentfill}%
\pgfsetlinewidth{0.401500pt}%
\definecolor{currentstroke}{rgb}{0.000000,0.000000,0.000000}%
\pgfsetstrokecolor{currentstroke}%
\pgfsetdash{}{0pt}%
\pgfsys@defobject{currentmarker}{\pgfqpoint{-0.020833in}{0.000000in}}{\pgfqpoint{-0.000000in}{0.000000in}}{%
\pgfpathmoveto{\pgfqpoint{-0.000000in}{0.000000in}}%
\pgfpathlineto{\pgfqpoint{-0.020833in}{0.000000in}}%
\pgfusepath{stroke,fill}%
}%
\begin{pgfscope}%
\pgfsys@transformshift{3.038110in}{1.800799in}%
\pgfsys@useobject{currentmarker}{}%
\end{pgfscope}%
\end{pgfscope}%
\begin{pgfscope}%
\pgfsetbuttcap%
\pgfsetroundjoin%
\definecolor{currentfill}{rgb}{0.000000,0.000000,0.000000}%
\pgfsetfillcolor{currentfill}%
\pgfsetlinewidth{0.401500pt}%
\definecolor{currentstroke}{rgb}{0.000000,0.000000,0.000000}%
\pgfsetstrokecolor{currentstroke}%
\pgfsetdash{}{0pt}%
\pgfsys@defobject{currentmarker}{\pgfqpoint{0.000000in}{0.000000in}}{\pgfqpoint{0.020833in}{0.000000in}}{%
\pgfpathmoveto{\pgfqpoint{0.000000in}{0.000000in}}%
\pgfpathlineto{\pgfqpoint{0.020833in}{0.000000in}}%
\pgfusepath{stroke,fill}%
}%
\begin{pgfscope}%
\pgfsys@transformshift{0.650000in}{1.877511in}%
\pgfsys@useobject{currentmarker}{}%
\end{pgfscope}%
\end{pgfscope}%
\begin{pgfscope}%
\pgfsetbuttcap%
\pgfsetroundjoin%
\definecolor{currentfill}{rgb}{0.000000,0.000000,0.000000}%
\pgfsetfillcolor{currentfill}%
\pgfsetlinewidth{0.401500pt}%
\definecolor{currentstroke}{rgb}{0.000000,0.000000,0.000000}%
\pgfsetstrokecolor{currentstroke}%
\pgfsetdash{}{0pt}%
\pgfsys@defobject{currentmarker}{\pgfqpoint{-0.020833in}{0.000000in}}{\pgfqpoint{-0.000000in}{0.000000in}}{%
\pgfpathmoveto{\pgfqpoint{-0.000000in}{0.000000in}}%
\pgfpathlineto{\pgfqpoint{-0.020833in}{0.000000in}}%
\pgfusepath{stroke,fill}%
}%
\begin{pgfscope}%
\pgfsys@transformshift{3.038110in}{1.877511in}%
\pgfsys@useobject{currentmarker}{}%
\end{pgfscope}%
\end{pgfscope}%
\begin{pgfscope}%
\pgfsetbuttcap%
\pgfsetroundjoin%
\definecolor{currentfill}{rgb}{0.000000,0.000000,0.000000}%
\pgfsetfillcolor{currentfill}%
\pgfsetlinewidth{0.401500pt}%
\definecolor{currentstroke}{rgb}{0.000000,0.000000,0.000000}%
\pgfsetstrokecolor{currentstroke}%
\pgfsetdash{}{0pt}%
\pgfsys@defobject{currentmarker}{\pgfqpoint{0.000000in}{0.000000in}}{\pgfqpoint{0.020833in}{0.000000in}}{%
\pgfpathmoveto{\pgfqpoint{0.000000in}{0.000000in}}%
\pgfpathlineto{\pgfqpoint{0.020833in}{0.000000in}}%
\pgfusepath{stroke,fill}%
}%
\begin{pgfscope}%
\pgfsys@transformshift{0.650000in}{1.954222in}%
\pgfsys@useobject{currentmarker}{}%
\end{pgfscope}%
\end{pgfscope}%
\begin{pgfscope}%
\pgfsetbuttcap%
\pgfsetroundjoin%
\definecolor{currentfill}{rgb}{0.000000,0.000000,0.000000}%
\pgfsetfillcolor{currentfill}%
\pgfsetlinewidth{0.401500pt}%
\definecolor{currentstroke}{rgb}{0.000000,0.000000,0.000000}%
\pgfsetstrokecolor{currentstroke}%
\pgfsetdash{}{0pt}%
\pgfsys@defobject{currentmarker}{\pgfqpoint{-0.020833in}{0.000000in}}{\pgfqpoint{-0.000000in}{0.000000in}}{%
\pgfpathmoveto{\pgfqpoint{-0.000000in}{0.000000in}}%
\pgfpathlineto{\pgfqpoint{-0.020833in}{0.000000in}}%
\pgfusepath{stroke,fill}%
}%
\begin{pgfscope}%
\pgfsys@transformshift{3.038110in}{1.954222in}%
\pgfsys@useobject{currentmarker}{}%
\end{pgfscope}%
\end{pgfscope}%
\begin{pgfscope}%
\pgfsetbuttcap%
\pgfsetroundjoin%
\definecolor{currentfill}{rgb}{0.000000,0.000000,0.000000}%
\pgfsetfillcolor{currentfill}%
\pgfsetlinewidth{0.401500pt}%
\definecolor{currentstroke}{rgb}{0.000000,0.000000,0.000000}%
\pgfsetstrokecolor{currentstroke}%
\pgfsetdash{}{0pt}%
\pgfsys@defobject{currentmarker}{\pgfqpoint{0.000000in}{0.000000in}}{\pgfqpoint{0.020833in}{0.000000in}}{%
\pgfpathmoveto{\pgfqpoint{0.000000in}{0.000000in}}%
\pgfpathlineto{\pgfqpoint{0.020833in}{0.000000in}}%
\pgfusepath{stroke,fill}%
}%
\begin{pgfscope}%
\pgfsys@transformshift{0.650000in}{2.107644in}%
\pgfsys@useobject{currentmarker}{}%
\end{pgfscope}%
\end{pgfscope}%
\begin{pgfscope}%
\pgfsetbuttcap%
\pgfsetroundjoin%
\definecolor{currentfill}{rgb}{0.000000,0.000000,0.000000}%
\pgfsetfillcolor{currentfill}%
\pgfsetlinewidth{0.401500pt}%
\definecolor{currentstroke}{rgb}{0.000000,0.000000,0.000000}%
\pgfsetstrokecolor{currentstroke}%
\pgfsetdash{}{0pt}%
\pgfsys@defobject{currentmarker}{\pgfqpoint{-0.020833in}{0.000000in}}{\pgfqpoint{-0.000000in}{0.000000in}}{%
\pgfpathmoveto{\pgfqpoint{-0.000000in}{0.000000in}}%
\pgfpathlineto{\pgfqpoint{-0.020833in}{0.000000in}}%
\pgfusepath{stroke,fill}%
}%
\begin{pgfscope}%
\pgfsys@transformshift{3.038110in}{2.107644in}%
\pgfsys@useobject{currentmarker}{}%
\end{pgfscope}%
\end{pgfscope}%
\begin{pgfscope}%
\pgfsetbuttcap%
\pgfsetroundjoin%
\definecolor{currentfill}{rgb}{0.000000,0.000000,0.000000}%
\pgfsetfillcolor{currentfill}%
\pgfsetlinewidth{0.401500pt}%
\definecolor{currentstroke}{rgb}{0.000000,0.000000,0.000000}%
\pgfsetstrokecolor{currentstroke}%
\pgfsetdash{}{0pt}%
\pgfsys@defobject{currentmarker}{\pgfqpoint{0.000000in}{0.000000in}}{\pgfqpoint{0.020833in}{0.000000in}}{%
\pgfpathmoveto{\pgfqpoint{0.000000in}{0.000000in}}%
\pgfpathlineto{\pgfqpoint{0.020833in}{0.000000in}}%
\pgfusepath{stroke,fill}%
}%
\begin{pgfscope}%
\pgfsys@transformshift{0.650000in}{2.184355in}%
\pgfsys@useobject{currentmarker}{}%
\end{pgfscope}%
\end{pgfscope}%
\begin{pgfscope}%
\pgfsetbuttcap%
\pgfsetroundjoin%
\definecolor{currentfill}{rgb}{0.000000,0.000000,0.000000}%
\pgfsetfillcolor{currentfill}%
\pgfsetlinewidth{0.401500pt}%
\definecolor{currentstroke}{rgb}{0.000000,0.000000,0.000000}%
\pgfsetstrokecolor{currentstroke}%
\pgfsetdash{}{0pt}%
\pgfsys@defobject{currentmarker}{\pgfqpoint{-0.020833in}{0.000000in}}{\pgfqpoint{-0.000000in}{0.000000in}}{%
\pgfpathmoveto{\pgfqpoint{-0.000000in}{0.000000in}}%
\pgfpathlineto{\pgfqpoint{-0.020833in}{0.000000in}}%
\pgfusepath{stroke,fill}%
}%
\begin{pgfscope}%
\pgfsys@transformshift{3.038110in}{2.184355in}%
\pgfsys@useobject{currentmarker}{}%
\end{pgfscope}%
\end{pgfscope}%
\begin{pgfscope}%
\pgfsetbuttcap%
\pgfsetroundjoin%
\definecolor{currentfill}{rgb}{0.000000,0.000000,0.000000}%
\pgfsetfillcolor{currentfill}%
\pgfsetlinewidth{0.401500pt}%
\definecolor{currentstroke}{rgb}{0.000000,0.000000,0.000000}%
\pgfsetstrokecolor{currentstroke}%
\pgfsetdash{}{0pt}%
\pgfsys@defobject{currentmarker}{\pgfqpoint{0.000000in}{0.000000in}}{\pgfqpoint{0.020833in}{0.000000in}}{%
\pgfpathmoveto{\pgfqpoint{0.000000in}{0.000000in}}%
\pgfpathlineto{\pgfqpoint{0.020833in}{0.000000in}}%
\pgfusepath{stroke,fill}%
}%
\begin{pgfscope}%
\pgfsys@transformshift{0.650000in}{2.261066in}%
\pgfsys@useobject{currentmarker}{}%
\end{pgfscope}%
\end{pgfscope}%
\begin{pgfscope}%
\pgfsetbuttcap%
\pgfsetroundjoin%
\definecolor{currentfill}{rgb}{0.000000,0.000000,0.000000}%
\pgfsetfillcolor{currentfill}%
\pgfsetlinewidth{0.401500pt}%
\definecolor{currentstroke}{rgb}{0.000000,0.000000,0.000000}%
\pgfsetstrokecolor{currentstroke}%
\pgfsetdash{}{0pt}%
\pgfsys@defobject{currentmarker}{\pgfqpoint{-0.020833in}{0.000000in}}{\pgfqpoint{-0.000000in}{0.000000in}}{%
\pgfpathmoveto{\pgfqpoint{-0.000000in}{0.000000in}}%
\pgfpathlineto{\pgfqpoint{-0.020833in}{0.000000in}}%
\pgfusepath{stroke,fill}%
}%
\begin{pgfscope}%
\pgfsys@transformshift{3.038110in}{2.261066in}%
\pgfsys@useobject{currentmarker}{}%
\end{pgfscope}%
\end{pgfscope}%
\begin{pgfscope}%
\pgfsetbuttcap%
\pgfsetroundjoin%
\definecolor{currentfill}{rgb}{0.000000,0.000000,0.000000}%
\pgfsetfillcolor{currentfill}%
\pgfsetlinewidth{0.401500pt}%
\definecolor{currentstroke}{rgb}{0.000000,0.000000,0.000000}%
\pgfsetstrokecolor{currentstroke}%
\pgfsetdash{}{0pt}%
\pgfsys@defobject{currentmarker}{\pgfqpoint{0.000000in}{0.000000in}}{\pgfqpoint{0.020833in}{0.000000in}}{%
\pgfpathmoveto{\pgfqpoint{0.000000in}{0.000000in}}%
\pgfpathlineto{\pgfqpoint{0.020833in}{0.000000in}}%
\pgfusepath{stroke,fill}%
}%
\begin{pgfscope}%
\pgfsys@transformshift{0.650000in}{2.337777in}%
\pgfsys@useobject{currentmarker}{}%
\end{pgfscope}%
\end{pgfscope}%
\begin{pgfscope}%
\pgfsetbuttcap%
\pgfsetroundjoin%
\definecolor{currentfill}{rgb}{0.000000,0.000000,0.000000}%
\pgfsetfillcolor{currentfill}%
\pgfsetlinewidth{0.401500pt}%
\definecolor{currentstroke}{rgb}{0.000000,0.000000,0.000000}%
\pgfsetstrokecolor{currentstroke}%
\pgfsetdash{}{0pt}%
\pgfsys@defobject{currentmarker}{\pgfqpoint{-0.020833in}{0.000000in}}{\pgfqpoint{-0.000000in}{0.000000in}}{%
\pgfpathmoveto{\pgfqpoint{-0.000000in}{0.000000in}}%
\pgfpathlineto{\pgfqpoint{-0.020833in}{0.000000in}}%
\pgfusepath{stroke,fill}%
}%
\begin{pgfscope}%
\pgfsys@transformshift{3.038110in}{2.337777in}%
\pgfsys@useobject{currentmarker}{}%
\end{pgfscope}%
\end{pgfscope}%
\begin{pgfscope}%
\definecolor{textcolor}{rgb}{0.000000,0.000000,0.000000}%
\pgfsetstrokecolor{textcolor}%
\pgfsetfillcolor{textcolor}%
\pgftext[x=0.260995in,y=1.417244in,,bottom,rotate=90.000000]{\color{textcolor}{\rmfamily\fontsize{12.000000}{14.400000}\selectfont\catcode`\^=\active\def^{\ifmmode\sp\else\^{}\fi}\catcode`\%=\active\def%{\%}$\langle u_i u_j \rangle^+$}}%
\end{pgfscope}%
\begin{pgfscope}%
\pgfpathrectangle{\pgfqpoint{0.650000in}{0.420000in}}{\pgfqpoint{2.388110in}{1.994488in}}%
\pgfusepath{clip}%
\pgfsetrectcap%
\pgfsetroundjoin%
\pgfsetlinewidth{0.602250pt}%
\definecolor{currentstroke}{rgb}{0.000000,0.000000,0.000000}%
\pgfsetstrokecolor{currentstroke}%
\pgfsetdash{}{0pt}%
\pgfpathmoveto{\pgfqpoint{0.645000in}{0.887245in}}%
\pgfpathlineto{\pgfqpoint{0.687793in}{0.888737in}}%
\pgfpathlineto{\pgfqpoint{0.816285in}{0.900826in}}%
\pgfpathlineto{\pgfqpoint{0.921267in}{0.922599in}}%
\pgfpathlineto{\pgfqpoint{1.010024in}{0.957929in}}%
\pgfpathlineto{\pgfqpoint{1.086906in}{1.010733in}}%
\pgfpathlineto{\pgfqpoint{1.154718in}{1.084090in}}%
\pgfpathlineto{\pgfqpoint{1.215374in}{1.179013in}}%
\pgfpathlineto{\pgfqpoint{1.270240in}{1.293315in}}%
\pgfpathlineto{\pgfqpoint{1.320326in}{1.421246in}}%
\pgfpathlineto{\pgfqpoint{1.366397in}{1.554286in}}%
\pgfpathlineto{\pgfqpoint{1.409048in}{1.682905in}}%
\pgfpathlineto{\pgfqpoint{1.448752in}{1.798534in}}%
\pgfpathlineto{\pgfqpoint{1.485889in}{1.895020in}}%
\pgfpathlineto{\pgfqpoint{1.520771in}{1.969183in}}%
\pgfpathlineto{\pgfqpoint{1.553655in}{2.020564in}}%
\pgfpathlineto{\pgfqpoint{1.584756in}{2.050714in}}%
\pgfpathlineto{\pgfqpoint{1.614259in}{2.062380in}}%
\pgfpathlineto{\pgfqpoint{1.642318in}{2.058816in}}%
\pgfpathlineto{\pgfqpoint{1.669069in}{2.043305in}}%
\pgfpathlineto{\pgfqpoint{1.694626in}{2.018872in}}%
\pgfpathlineto{\pgfqpoint{1.719092in}{1.988148in}}%
\pgfpathlineto{\pgfqpoint{1.742555in}{1.953333in}}%
\pgfpathlineto{\pgfqpoint{1.765095in}{1.916211in}}%
\pgfpathlineto{\pgfqpoint{1.786780in}{1.878197in}}%
\pgfpathlineto{\pgfqpoint{1.866126in}{1.734971in}}%
\pgfpathlineto{\pgfqpoint{1.884352in}{1.703735in}}%
\pgfpathlineto{\pgfqpoint{1.902013in}{1.674690in}}%
\pgfpathlineto{\pgfqpoint{1.919144in}{1.647843in}}%
\pgfpathlineto{\pgfqpoint{1.935774in}{1.623145in}}%
\pgfpathlineto{\pgfqpoint{1.951933in}{1.600514in}}%
\pgfpathlineto{\pgfqpoint{1.967645in}{1.579839in}}%
\pgfpathlineto{\pgfqpoint{1.982934in}{1.560983in}}%
\pgfpathlineto{\pgfqpoint{1.997823in}{1.543800in}}%
\pgfpathlineto{\pgfqpoint{2.012331in}{1.528159in}}%
\pgfpathlineto{\pgfqpoint{2.026477in}{1.513933in}}%
\pgfpathlineto{\pgfqpoint{2.040279in}{1.500980in}}%
\pgfpathlineto{\pgfqpoint{2.066912in}{1.478315in}}%
\pgfpathlineto{\pgfqpoint{2.092348in}{1.459231in}}%
\pgfpathlineto{\pgfqpoint{2.116686in}{1.443088in}}%
\pgfpathlineto{\pgfqpoint{2.140016in}{1.429249in}}%
\pgfpathlineto{\pgfqpoint{2.173292in}{1.411240in}}%
\pgfpathlineto{\pgfqpoint{2.214793in}{1.390568in}}%
\pgfpathlineto{\pgfqpoint{2.323444in}{1.337454in}}%
\pgfpathlineto{\pgfqpoint{2.355381in}{1.320428in}}%
\pgfpathlineto{\pgfqpoint{2.378168in}{1.307437in}}%
\pgfpathlineto{\pgfqpoint{2.407140in}{1.289398in}}%
\pgfpathlineto{\pgfqpoint{2.427896in}{1.275287in}}%
\pgfpathlineto{\pgfqpoint{2.447883in}{1.260395in}}%
\pgfpathlineto{\pgfqpoint{2.473424in}{1.239818in}}%
\pgfpathlineto{\pgfqpoint{2.497797in}{1.218908in}}%
\pgfpathlineto{\pgfqpoint{2.532371in}{1.187396in}}%
\pgfpathlineto{\pgfqpoint{2.548839in}{1.170967in}}%
\pgfpathlineto{\pgfqpoint{2.564802in}{1.153373in}}%
\pgfpathlineto{\pgfqpoint{2.580286in}{1.134845in}}%
\pgfpathlineto{\pgfqpoint{2.595318in}{1.115279in}}%
\pgfpathlineto{\pgfqpoint{2.614695in}{1.087678in}}%
\pgfpathlineto{\pgfqpoint{2.628756in}{1.066094in}}%
\pgfpathlineto{\pgfqpoint{2.642434in}{1.043330in}}%
\pgfpathlineto{\pgfqpoint{2.664429in}{1.004280in}}%
\pgfpathlineto{\pgfqpoint{2.693657in}{0.952195in}}%
\pgfpathlineto{\pgfqpoint{2.705664in}{0.932876in}}%
\pgfpathlineto{\pgfqpoint{2.713505in}{0.921686in}}%
\pgfpathlineto{\pgfqpoint{2.721220in}{0.912190in}}%
\pgfpathlineto{\pgfqpoint{2.728811in}{0.904429in}}%
\pgfpathlineto{\pgfqpoint{2.736282in}{0.898280in}}%
\pgfpathlineto{\pgfqpoint{2.743636in}{0.893578in}}%
\pgfpathlineto{\pgfqpoint{2.750875in}{0.890059in}}%
\pgfpathlineto{\pgfqpoint{2.761527in}{0.886488in}}%
\pgfpathlineto{\pgfqpoint{2.771937in}{0.884343in}}%
\pgfpathlineto{\pgfqpoint{2.785457in}{0.882738in}}%
\pgfpathlineto{\pgfqpoint{2.805002in}{0.881584in}}%
\pgfpathlineto{\pgfqpoint{2.841653in}{0.880757in}}%
\pgfpathlineto{\pgfqpoint{2.923547in}{0.880349in}}%
\pgfpathlineto{\pgfqpoint{3.043110in}{0.880275in}}%
\pgfpathlineto{\pgfqpoint{3.043110in}{0.880275in}}%
\pgfusepath{stroke}%
\end{pgfscope}%
\begin{pgfscope}%
\pgfpathrectangle{\pgfqpoint{0.650000in}{0.420000in}}{\pgfqpoint{2.388110in}{1.994488in}}%
\pgfusepath{clip}%
\pgfsetrectcap%
\pgfsetroundjoin%
\pgfsetlinewidth{0.602250pt}%
\definecolor{currentstroke}{rgb}{0.000000,0.000000,0.000000}%
\pgfsetstrokecolor{currentstroke}%
\pgfsetdash{}{0pt}%
\pgfpathmoveto{\pgfqpoint{0.645000in}{0.880268in}}%
\pgfpathlineto{\pgfqpoint{1.086906in}{0.880626in}}%
\pgfpathlineto{\pgfqpoint{1.215374in}{0.881779in}}%
\pgfpathlineto{\pgfqpoint{1.270240in}{0.882942in}}%
\pgfpathlineto{\pgfqpoint{1.320326in}{0.884658in}}%
\pgfpathlineto{\pgfqpoint{1.366397in}{0.887037in}}%
\pgfpathlineto{\pgfqpoint{1.409048in}{0.890170in}}%
\pgfpathlineto{\pgfqpoint{1.448752in}{0.894125in}}%
\pgfpathlineto{\pgfqpoint{1.485889in}{0.898938in}}%
\pgfpathlineto{\pgfqpoint{1.520771in}{0.904611in}}%
\pgfpathlineto{\pgfqpoint{1.553655in}{0.911114in}}%
\pgfpathlineto{\pgfqpoint{1.584756in}{0.918381in}}%
\pgfpathlineto{\pgfqpoint{1.614259in}{0.926313in}}%
\pgfpathlineto{\pgfqpoint{1.642318in}{0.934788in}}%
\pgfpathlineto{\pgfqpoint{1.669069in}{0.943670in}}%
\pgfpathlineto{\pgfqpoint{1.719092in}{0.962071in}}%
\pgfpathlineto{\pgfqpoint{1.786780in}{0.989260in}}%
\pgfpathlineto{\pgfqpoint{1.847299in}{1.013548in}}%
\pgfpathlineto{\pgfqpoint{1.884352in}{1.027372in}}%
\pgfpathlineto{\pgfqpoint{1.919144in}{1.039131in}}%
\pgfpathlineto{\pgfqpoint{1.951933in}{1.048885in}}%
\pgfpathlineto{\pgfqpoint{1.982934in}{1.056800in}}%
\pgfpathlineto{\pgfqpoint{2.012331in}{1.063095in}}%
\pgfpathlineto{\pgfqpoint{2.040279in}{1.067988in}}%
\pgfpathlineto{\pgfqpoint{2.066912in}{1.071696in}}%
\pgfpathlineto{\pgfqpoint{2.104648in}{1.075487in}}%
\pgfpathlineto{\pgfqpoint{2.140016in}{1.077616in}}%
\pgfpathlineto{\pgfqpoint{2.173292in}{1.078505in}}%
\pgfpathlineto{\pgfqpoint{2.214793in}{1.078286in}}%
\pgfpathlineto{\pgfqpoint{2.262674in}{1.076553in}}%
\pgfpathlineto{\pgfqpoint{2.298213in}{1.074267in}}%
\pgfpathlineto{\pgfqpoint{2.331604in}{1.071189in}}%
\pgfpathlineto{\pgfqpoint{2.363083in}{1.067247in}}%
\pgfpathlineto{\pgfqpoint{2.392847in}{1.062548in}}%
\pgfpathlineto{\pgfqpoint{2.421066in}{1.057033in}}%
\pgfpathlineto{\pgfqpoint{2.447883in}{1.050752in}}%
\pgfpathlineto{\pgfqpoint{2.473424in}{1.043706in}}%
\pgfpathlineto{\pgfqpoint{2.497797in}{1.035708in}}%
\pgfpathlineto{\pgfqpoint{2.521097in}{1.026747in}}%
\pgfpathlineto{\pgfqpoint{2.543407in}{1.016831in}}%
\pgfpathlineto{\pgfqpoint{2.564802in}{1.006219in}}%
\pgfpathlineto{\pgfqpoint{2.585346in}{0.994853in}}%
\pgfpathlineto{\pgfqpoint{2.609919in}{0.979740in}}%
\pgfpathlineto{\pgfqpoint{2.637916in}{0.961269in}}%
\pgfpathlineto{\pgfqpoint{2.717378in}{0.907840in}}%
\pgfpathlineto{\pgfqpoint{2.732561in}{0.899677in}}%
\pgfpathlineto{\pgfqpoint{2.747270in}{0.893322in}}%
\pgfpathlineto{\pgfqpoint{2.761527in}{0.888731in}}%
\pgfpathlineto{\pgfqpoint{2.775355in}{0.885641in}}%
\pgfpathlineto{\pgfqpoint{2.792068in}{0.883316in}}%
\pgfpathlineto{\pgfqpoint{2.811330in}{0.881902in}}%
\pgfpathlineto{\pgfqpoint{2.841653in}{0.880964in}}%
\pgfpathlineto{\pgfqpoint{2.903990in}{0.880434in}}%
\pgfpathlineto{\pgfqpoint{3.043110in}{0.880279in}}%
\pgfpathlineto{\pgfqpoint{3.043110in}{0.880279in}}%
\pgfusepath{stroke}%
\end{pgfscope}%
\begin{pgfscope}%
\pgfpathrectangle{\pgfqpoint{0.650000in}{0.420000in}}{\pgfqpoint{2.388110in}{1.994488in}}%
\pgfusepath{clip}%
\pgfsetrectcap%
\pgfsetroundjoin%
\pgfsetlinewidth{0.602250pt}%
\definecolor{currentstroke}{rgb}{0.000000,0.000000,0.000000}%
\pgfsetstrokecolor{currentstroke}%
\pgfsetdash{}{0pt}%
\pgfpathmoveto{\pgfqpoint{0.645000in}{0.883160in}}%
\pgfpathlineto{\pgfqpoint{0.687793in}{0.883759in}}%
\pgfpathlineto{\pgfqpoint{0.816285in}{0.888163in}}%
\pgfpathlineto{\pgfqpoint{0.921267in}{0.895204in}}%
\pgfpathlineto{\pgfqpoint{1.010024in}{0.905162in}}%
\pgfpathlineto{\pgfqpoint{1.086906in}{0.917981in}}%
\pgfpathlineto{\pgfqpoint{1.154718in}{0.933299in}}%
\pgfpathlineto{\pgfqpoint{1.215374in}{0.950539in}}%
\pgfpathlineto{\pgfqpoint{1.270240in}{0.969041in}}%
\pgfpathlineto{\pgfqpoint{1.320326in}{0.988190in}}%
\pgfpathlineto{\pgfqpoint{1.366397in}{1.007505in}}%
\pgfpathlineto{\pgfqpoint{1.409048in}{1.026653in}}%
\pgfpathlineto{\pgfqpoint{1.448752in}{1.045418in}}%
\pgfpathlineto{\pgfqpoint{1.520771in}{1.081201in}}%
\pgfpathlineto{\pgfqpoint{1.614259in}{1.128288in}}%
\pgfpathlineto{\pgfqpoint{1.642318in}{1.141649in}}%
\pgfpathlineto{\pgfqpoint{1.669069in}{1.153670in}}%
\pgfpathlineto{\pgfqpoint{1.694626in}{1.164306in}}%
\pgfpathlineto{\pgfqpoint{1.719092in}{1.173558in}}%
\pgfpathlineto{\pgfqpoint{1.742555in}{1.181468in}}%
\pgfpathlineto{\pgfqpoint{1.765095in}{1.188115in}}%
\pgfpathlineto{\pgfqpoint{1.786780in}{1.193606in}}%
\pgfpathlineto{\pgfqpoint{1.807673in}{1.198063in}}%
\pgfpathlineto{\pgfqpoint{1.827829in}{1.201604in}}%
\pgfpathlineto{\pgfqpoint{1.847299in}{1.204346in}}%
\pgfpathlineto{\pgfqpoint{1.884352in}{1.207843in}}%
\pgfpathlineto{\pgfqpoint{1.919144in}{1.209235in}}%
\pgfpathlineto{\pgfqpoint{1.951933in}{1.209035in}}%
\pgfpathlineto{\pgfqpoint{1.982934in}{1.207614in}}%
\pgfpathlineto{\pgfqpoint{2.012331in}{1.205230in}}%
\pgfpathlineto{\pgfqpoint{2.053752in}{1.200550in}}%
\pgfpathlineto{\pgfqpoint{2.092348in}{1.195228in}}%
\pgfpathlineto{\pgfqpoint{2.140016in}{1.187421in}}%
\pgfpathlineto{\pgfqpoint{2.289540in}{1.160418in}}%
\pgfpathlineto{\pgfqpoint{2.355381in}{1.147227in}}%
\pgfpathlineto{\pgfqpoint{2.385557in}{1.139867in}}%
\pgfpathlineto{\pgfqpoint{2.414148in}{1.131664in}}%
\pgfpathlineto{\pgfqpoint{2.434641in}{1.124756in}}%
\pgfpathlineto{\pgfqpoint{2.460806in}{1.114707in}}%
\pgfpathlineto{\pgfqpoint{2.479623in}{1.106473in}}%
\pgfpathlineto{\pgfqpoint{2.497797in}{1.097415in}}%
\pgfpathlineto{\pgfqpoint{2.515368in}{1.087244in}}%
\pgfpathlineto{\pgfqpoint{2.532371in}{1.076266in}}%
\pgfpathlineto{\pgfqpoint{2.548839in}{1.064369in}}%
\pgfpathlineto{\pgfqpoint{2.564802in}{1.051417in}}%
\pgfpathlineto{\pgfqpoint{2.580286in}{1.037639in}}%
\pgfpathlineto{\pgfqpoint{2.600232in}{1.018214in}}%
\pgfpathlineto{\pgfqpoint{2.619426in}{0.998051in}}%
\pgfpathlineto{\pgfqpoint{2.646912in}{0.967574in}}%
\pgfpathlineto{\pgfqpoint{2.668712in}{0.943513in}}%
\pgfpathlineto{\pgfqpoint{2.685483in}{0.926506in}}%
\pgfpathlineto{\pgfqpoint{2.697692in}{0.915485in}}%
\pgfpathlineto{\pgfqpoint{2.709601in}{0.906219in}}%
\pgfpathlineto{\pgfqpoint{2.721220in}{0.898773in}}%
\pgfpathlineto{\pgfqpoint{2.732561in}{0.892993in}}%
\pgfpathlineto{\pgfqpoint{2.743636in}{0.888765in}}%
\pgfpathlineto{\pgfqpoint{2.754453in}{0.885795in}}%
\pgfpathlineto{\pgfqpoint{2.768494in}{0.883312in}}%
\pgfpathlineto{\pgfqpoint{2.785457in}{0.881723in}}%
\pgfpathlineto{\pgfqpoint{2.808177in}{0.880844in}}%
\pgfpathlineto{\pgfqpoint{2.856039in}{0.880403in}}%
\pgfpathlineto{\pgfqpoint{3.043110in}{0.880268in}}%
\pgfpathlineto{\pgfqpoint{3.043110in}{0.880268in}}%
\pgfusepath{stroke}%
\end{pgfscope}%
\begin{pgfscope}%
\pgfpathrectangle{\pgfqpoint{0.650000in}{0.420000in}}{\pgfqpoint{2.388110in}{1.994488in}}%
\pgfusepath{clip}%
\pgfsetrectcap%
\pgfsetroundjoin%
\pgfsetlinewidth{0.602250pt}%
\definecolor{currentstroke}{rgb}{0.000000,0.000000,0.000000}%
\pgfsetstrokecolor{currentstroke}%
\pgfsetdash{}{0pt}%
\pgfpathmoveto{\pgfqpoint{0.645000in}{0.880241in}}%
\pgfpathlineto{\pgfqpoint{0.921267in}{0.879887in}}%
\pgfpathlineto{\pgfqpoint{1.010024in}{0.879299in}}%
\pgfpathlineto{\pgfqpoint{1.086906in}{0.878114in}}%
\pgfpathlineto{\pgfqpoint{1.154718in}{0.875992in}}%
\pgfpathlineto{\pgfqpoint{1.215374in}{0.872564in}}%
\pgfpathlineto{\pgfqpoint{1.270240in}{0.867528in}}%
\pgfpathlineto{\pgfqpoint{1.320326in}{0.860744in}}%
\pgfpathlineto{\pgfqpoint{1.366397in}{0.852295in}}%
\pgfpathlineto{\pgfqpoint{1.409048in}{0.842490in}}%
\pgfpathlineto{\pgfqpoint{1.448752in}{0.831804in}}%
\pgfpathlineto{\pgfqpoint{1.520771in}{0.809922in}}%
\pgfpathlineto{\pgfqpoint{1.584756in}{0.790242in}}%
\pgfpathlineto{\pgfqpoint{1.614259in}{0.781844in}}%
\pgfpathlineto{\pgfqpoint{1.642318in}{0.774499in}}%
\pgfpathlineto{\pgfqpoint{1.669069in}{0.768171in}}%
\pgfpathlineto{\pgfqpoint{1.694626in}{0.762779in}}%
\pgfpathlineto{\pgfqpoint{1.742555in}{0.754386in}}%
\pgfpathlineto{\pgfqpoint{1.786780in}{0.748489in}}%
\pgfpathlineto{\pgfqpoint{1.827829in}{0.744389in}}%
\pgfpathlineto{\pgfqpoint{1.866126in}{0.741569in}}%
\pgfpathlineto{\pgfqpoint{1.919144in}{0.738968in}}%
\pgfpathlineto{\pgfqpoint{1.982934in}{0.737381in}}%
\pgfpathlineto{\pgfqpoint{2.053752in}{0.737012in}}%
\pgfpathlineto{\pgfqpoint{2.128472in}{0.737755in}}%
\pgfpathlineto{\pgfqpoint{2.204702in}{0.739568in}}%
\pgfpathlineto{\pgfqpoint{2.262674in}{0.741973in}}%
\pgfpathlineto{\pgfqpoint{2.306753in}{0.744770in}}%
\pgfpathlineto{\pgfqpoint{2.347570in}{0.748525in}}%
\pgfpathlineto{\pgfqpoint{2.385557in}{0.753219in}}%
\pgfpathlineto{\pgfqpoint{2.421066in}{0.758787in}}%
\pgfpathlineto{\pgfqpoint{2.454383in}{0.765173in}}%
\pgfpathlineto{\pgfqpoint{2.485750in}{0.772219in}}%
\pgfpathlineto{\pgfqpoint{2.515368in}{0.779984in}}%
\pgfpathlineto{\pgfqpoint{2.543407in}{0.788452in}}%
\pgfpathlineto{\pgfqpoint{2.570015in}{0.797761in}}%
\pgfpathlineto{\pgfqpoint{2.590356in}{0.805965in}}%
\pgfpathlineto{\pgfqpoint{2.614695in}{0.817097in}}%
\pgfpathlineto{\pgfqpoint{2.637916in}{0.828904in}}%
\pgfpathlineto{\pgfqpoint{2.701695in}{0.862661in}}%
\pgfpathlineto{\pgfqpoint{2.717378in}{0.869163in}}%
\pgfpathlineto{\pgfqpoint{2.728811in}{0.872897in}}%
\pgfpathlineto{\pgfqpoint{2.743636in}{0.876364in}}%
\pgfpathlineto{\pgfqpoint{2.758004in}{0.878424in}}%
\pgfpathlineto{\pgfqpoint{2.775355in}{0.879633in}}%
\pgfpathlineto{\pgfqpoint{2.805002in}{0.880200in}}%
\pgfpathlineto{\pgfqpoint{2.930625in}{0.880265in}}%
\pgfpathlineto{\pgfqpoint{3.043110in}{0.880266in}}%
\pgfpathlineto{\pgfqpoint{3.043110in}{0.880266in}}%
\pgfusepath{stroke}%
\end{pgfscope}%
\begin{pgfscope}%
\pgfpathrectangle{\pgfqpoint{0.650000in}{0.420000in}}{\pgfqpoint{2.388110in}{1.994488in}}%
\pgfusepath{clip}%
\pgfsetrectcap%
\pgfsetroundjoin%
\pgfsetlinewidth{1.003750pt}%
\definecolor{currentstroke}{rgb}{0.670588,0.850980,0.913725}%
\pgfsetstrokecolor{currentstroke}%
\pgfsetdash{}{0pt}%
\pgfpathmoveto{\pgfqpoint{0.645000in}{0.886882in}}%
\pgfpathlineto{\pgfqpoint{0.679586in}{0.888317in}}%
\pgfpathlineto{\pgfqpoint{0.744946in}{0.892674in}}%
\pgfpathlineto{\pgfqpoint{0.802042in}{0.898418in}}%
\pgfpathlineto{\pgfqpoint{0.853159in}{0.905845in}}%
\pgfpathlineto{\pgfqpoint{0.899759in}{0.915307in}}%
\pgfpathlineto{\pgfqpoint{0.942836in}{0.927203in}}%
\pgfpathlineto{\pgfqpoint{0.983096in}{0.941973in}}%
\pgfpathlineto{\pgfqpoint{1.021057in}{0.960080in}}%
\pgfpathlineto{\pgfqpoint{1.057113in}{0.982006in}}%
\pgfpathlineto{\pgfqpoint{1.091568in}{1.008245in}}%
\pgfpathlineto{\pgfqpoint{1.124662in}{1.039309in}}%
\pgfpathlineto{\pgfqpoint{1.156588in}{1.075730in}}%
\pgfpathlineto{\pgfqpoint{1.187505in}{1.118056in}}%
\pgfpathlineto{\pgfqpoint{1.217543in}{1.166772in}}%
\pgfpathlineto{\pgfqpoint{1.246809in}{1.222030in}}%
\pgfpathlineto{\pgfqpoint{1.275397in}{1.283123in}}%
\pgfpathlineto{\pgfqpoint{1.303382in}{1.349360in}}%
\pgfpathlineto{\pgfqpoint{1.330833in}{1.419873in}}%
\pgfpathlineto{\pgfqpoint{1.357807in}{1.493324in}}%
\pgfpathlineto{\pgfqpoint{1.436334in}{1.714946in}}%
\pgfpathlineto{\pgfqpoint{1.461840in}{1.781911in}}%
\pgfpathlineto{\pgfqpoint{1.487063in}{1.842010in}}%
\pgfpathlineto{\pgfqpoint{1.512030in}{1.893385in}}%
\pgfpathlineto{\pgfqpoint{1.536766in}{1.934517in}}%
\pgfpathlineto{\pgfqpoint{1.561290in}{1.965469in}}%
\pgfpathlineto{\pgfqpoint{1.585622in}{1.985407in}}%
\pgfpathlineto{\pgfqpoint{1.609778in}{1.995061in}}%
\pgfpathlineto{\pgfqpoint{1.633774in}{1.994565in}}%
\pgfpathlineto{\pgfqpoint{1.657624in}{1.985039in}}%
\pgfpathlineto{\pgfqpoint{1.681340in}{1.967675in}}%
\pgfpathlineto{\pgfqpoint{1.704933in}{1.943561in}}%
\pgfpathlineto{\pgfqpoint{1.728413in}{1.913874in}}%
\pgfpathlineto{\pgfqpoint{1.751791in}{1.880101in}}%
\pgfpathlineto{\pgfqpoint{1.775073in}{1.843175in}}%
\pgfpathlineto{\pgfqpoint{1.867400in}{1.689251in}}%
\pgfpathlineto{\pgfqpoint{1.890311in}{1.653437in}}%
\pgfpathlineto{\pgfqpoint{1.913165in}{1.619819in}}%
\pgfpathlineto{\pgfqpoint{1.935966in}{1.589265in}}%
\pgfpathlineto{\pgfqpoint{1.958719in}{1.562116in}}%
\pgfpathlineto{\pgfqpoint{1.981426in}{1.538124in}}%
\pgfpathlineto{\pgfqpoint{2.004092in}{1.516560in}}%
\pgfpathlineto{\pgfqpoint{2.049312in}{1.476531in}}%
\pgfpathlineto{\pgfqpoint{2.071871in}{1.457217in}}%
\pgfpathlineto{\pgfqpoint{2.094401in}{1.439185in}}%
\pgfpathlineto{\pgfqpoint{2.116902in}{1.423683in}}%
\pgfpathlineto{\pgfqpoint{2.139378in}{1.410721in}}%
\pgfpathlineto{\pgfqpoint{2.161830in}{1.400143in}}%
\pgfpathlineto{\pgfqpoint{2.205858in}{1.382375in}}%
\pgfpathlineto{\pgfqpoint{2.225949in}{1.374575in}}%
\pgfpathlineto{\pgfqpoint{2.278969in}{1.356652in}}%
\pgfpathlineto{\pgfqpoint{2.323728in}{1.337087in}}%
\pgfpathlineto{\pgfqpoint{2.337224in}{1.330081in}}%
\pgfpathlineto{\pgfqpoint{2.350117in}{1.322005in}}%
\pgfpathlineto{\pgfqpoint{2.385654in}{1.295909in}}%
\pgfpathlineto{\pgfqpoint{2.396587in}{1.289820in}}%
\pgfpathlineto{\pgfqpoint{2.445778in}{1.265544in}}%
\pgfpathlineto{\pgfqpoint{2.454683in}{1.260785in}}%
\pgfpathlineto{\pgfqpoint{2.463320in}{1.255176in}}%
\pgfpathlineto{\pgfqpoint{2.479855in}{1.242762in}}%
\pgfpathlineto{\pgfqpoint{2.495491in}{1.229813in}}%
\pgfpathlineto{\pgfqpoint{2.524426in}{1.203156in}}%
\pgfpathlineto{\pgfqpoint{2.563013in}{1.175559in}}%
\pgfpathlineto{\pgfqpoint{2.580526in}{1.158439in}}%
\pgfpathlineto{\pgfqpoint{2.591637in}{1.145801in}}%
\pgfpathlineto{\pgfqpoint{2.602334in}{1.131006in}}%
\pgfpathlineto{\pgfqpoint{2.612648in}{1.117092in}}%
\pgfpathlineto{\pgfqpoint{2.646090in}{1.075137in}}%
\pgfpathlineto{\pgfqpoint{2.667802in}{1.049935in}}%
\pgfpathlineto{\pgfqpoint{2.676048in}{1.037135in}}%
\pgfpathlineto{\pgfqpoint{2.706857in}{0.986674in}}%
\pgfpathlineto{\pgfqpoint{2.727976in}{0.954223in}}%
\pgfpathlineto{\pgfqpoint{2.737982in}{0.938112in}}%
\pgfpathlineto{\pgfqpoint{2.747652in}{0.924617in}}%
\pgfpathlineto{\pgfqpoint{2.753923in}{0.917114in}}%
\pgfpathlineto{\pgfqpoint{2.763080in}{0.908252in}}%
\pgfpathlineto{\pgfqpoint{2.771955in}{0.900686in}}%
\pgfpathlineto{\pgfqpoint{2.780564in}{0.895096in}}%
\pgfpathlineto{\pgfqpoint{2.788924in}{0.891113in}}%
\pgfpathlineto{\pgfqpoint{2.802338in}{0.886542in}}%
\pgfpathlineto{\pgfqpoint{2.812638in}{0.884866in}}%
\pgfpathlineto{\pgfqpoint{2.836881in}{0.882672in}}%
\pgfpathlineto{\pgfqpoint{2.887886in}{0.881349in}}%
\pgfpathlineto{\pgfqpoint{2.924419in}{0.880966in}}%
\pgfpathlineto{\pgfqpoint{3.012421in}{0.880681in}}%
\pgfpathlineto{\pgfqpoint{3.043110in}{0.880578in}}%
\pgfpathlineto{\pgfqpoint{3.043110in}{0.880578in}}%
\pgfusepath{stroke}%
\end{pgfscope}%
\begin{pgfscope}%
\pgfpathrectangle{\pgfqpoint{0.650000in}{0.420000in}}{\pgfqpoint{2.388110in}{1.994488in}}%
\pgfusepath{clip}%
\pgfsetrectcap%
\pgfsetroundjoin%
\pgfsetlinewidth{1.003750pt}%
\definecolor{currentstroke}{rgb}{0.454902,0.678431,0.819608}%
\pgfsetstrokecolor{currentstroke}%
\pgfsetdash{}{0pt}%
\pgfpathmoveto{\pgfqpoint{0.645000in}{0.880269in}}%
\pgfpathlineto{\pgfqpoint{1.124662in}{0.880814in}}%
\pgfpathlineto{\pgfqpoint{1.217543in}{0.881777in}}%
\pgfpathlineto{\pgfqpoint{1.303382in}{0.883905in}}%
\pgfpathlineto{\pgfqpoint{1.357807in}{0.886355in}}%
\pgfpathlineto{\pgfqpoint{1.410517in}{0.890016in}}%
\pgfpathlineto{\pgfqpoint{1.461840in}{0.895238in}}%
\pgfpathlineto{\pgfqpoint{1.512030in}{0.902427in}}%
\pgfpathlineto{\pgfqpoint{1.536766in}{0.906814in}}%
\pgfpathlineto{\pgfqpoint{1.561290in}{0.911827in}}%
\pgfpathlineto{\pgfqpoint{1.585622in}{0.917437in}}%
\pgfpathlineto{\pgfqpoint{1.609778in}{0.923697in}}%
\pgfpathlineto{\pgfqpoint{1.633774in}{0.930577in}}%
\pgfpathlineto{\pgfqpoint{1.681340in}{0.946021in}}%
\pgfpathlineto{\pgfqpoint{1.728413in}{0.963253in}}%
\pgfpathlineto{\pgfqpoint{1.867400in}{1.016152in}}%
\pgfpathlineto{\pgfqpoint{1.913165in}{1.031047in}}%
\pgfpathlineto{\pgfqpoint{1.935966in}{1.037524in}}%
\pgfpathlineto{\pgfqpoint{1.958719in}{1.043277in}}%
\pgfpathlineto{\pgfqpoint{1.981426in}{1.048221in}}%
\pgfpathlineto{\pgfqpoint{2.004092in}{1.052290in}}%
\pgfpathlineto{\pgfqpoint{2.026719in}{1.055617in}}%
\pgfpathlineto{\pgfqpoint{2.049312in}{1.058364in}}%
\pgfpathlineto{\pgfqpoint{2.094401in}{1.062228in}}%
\pgfpathlineto{\pgfqpoint{2.184260in}{1.067687in}}%
\pgfpathlineto{\pgfqpoint{2.225949in}{1.068637in}}%
\pgfpathlineto{\pgfqpoint{2.244728in}{1.068094in}}%
\pgfpathlineto{\pgfqpoint{2.278969in}{1.065582in}}%
\pgfpathlineto{\pgfqpoint{2.309568in}{1.062131in}}%
\pgfpathlineto{\pgfqpoint{2.337224in}{1.058241in}}%
\pgfpathlineto{\pgfqpoint{2.385654in}{1.054138in}}%
\pgfpathlineto{\pgfqpoint{2.417283in}{1.049700in}}%
\pgfpathlineto{\pgfqpoint{2.436590in}{1.046147in}}%
\pgfpathlineto{\pgfqpoint{2.463320in}{1.039355in}}%
\pgfpathlineto{\pgfqpoint{2.524426in}{1.022351in}}%
\pgfpathlineto{\pgfqpoint{2.563013in}{1.008635in}}%
\pgfpathlineto{\pgfqpoint{2.580526in}{1.000051in}}%
\pgfpathlineto{\pgfqpoint{2.597035in}{0.992807in}}%
\pgfpathlineto{\pgfqpoint{2.607537in}{0.988333in}}%
\pgfpathlineto{\pgfqpoint{2.622606in}{0.980431in}}%
\pgfpathlineto{\pgfqpoint{2.691864in}{0.942150in}}%
\pgfpathlineto{\pgfqpoint{2.710486in}{0.930525in}}%
\pgfpathlineto{\pgfqpoint{2.727976in}{0.919585in}}%
\pgfpathlineto{\pgfqpoint{2.741242in}{0.912818in}}%
\pgfpathlineto{\pgfqpoint{2.774853in}{0.897276in}}%
\pgfpathlineto{\pgfqpoint{2.791657in}{0.891541in}}%
\pgfpathlineto{\pgfqpoint{2.804948in}{0.888156in}}%
\pgfpathlineto{\pgfqpoint{2.817652in}{0.885968in}}%
\pgfpathlineto{\pgfqpoint{2.836881in}{0.884097in}}%
\pgfpathlineto{\pgfqpoint{2.867720in}{0.882687in}}%
\pgfpathlineto{\pgfqpoint{2.912153in}{0.881928in}}%
\pgfpathlineto{\pgfqpoint{3.043110in}{0.881571in}}%
\pgfpathlineto{\pgfqpoint{3.043110in}{0.881571in}}%
\pgfusepath{stroke}%
\end{pgfscope}%
\begin{pgfscope}%
\pgfpathrectangle{\pgfqpoint{0.650000in}{0.420000in}}{\pgfqpoint{2.388110in}{1.994488in}}%
\pgfusepath{clip}%
\pgfsetrectcap%
\pgfsetroundjoin%
\pgfsetlinewidth{1.003750pt}%
\definecolor{currentstroke}{rgb}{0.270588,0.458824,0.705882}%
\pgfsetstrokecolor{currentstroke}%
\pgfsetdash{}{0pt}%
\pgfpathmoveto{\pgfqpoint{0.645000in}{0.882996in}}%
\pgfpathlineto{\pgfqpoint{0.679586in}{0.883560in}}%
\pgfpathlineto{\pgfqpoint{0.744946in}{0.885208in}}%
\pgfpathlineto{\pgfqpoint{0.802042in}{0.887285in}}%
\pgfpathlineto{\pgfqpoint{0.853159in}{0.889835in}}%
\pgfpathlineto{\pgfqpoint{0.899759in}{0.892902in}}%
\pgfpathlineto{\pgfqpoint{0.942836in}{0.896524in}}%
\pgfpathlineto{\pgfqpoint{0.983096in}{0.900733in}}%
\pgfpathlineto{\pgfqpoint{1.021057in}{0.905551in}}%
\pgfpathlineto{\pgfqpoint{1.057113in}{0.910993in}}%
\pgfpathlineto{\pgfqpoint{1.091568in}{0.917062in}}%
\pgfpathlineto{\pgfqpoint{1.124662in}{0.923755in}}%
\pgfpathlineto{\pgfqpoint{1.156588in}{0.931067in}}%
\pgfpathlineto{\pgfqpoint{1.187505in}{0.938990in}}%
\pgfpathlineto{\pgfqpoint{1.217543in}{0.947518in}}%
\pgfpathlineto{\pgfqpoint{1.246809in}{0.956621in}}%
\pgfpathlineto{\pgfqpoint{1.275397in}{0.966187in}}%
\pgfpathlineto{\pgfqpoint{1.303382in}{0.976140in}}%
\pgfpathlineto{\pgfqpoint{1.357807in}{0.997143in}}%
\pgfpathlineto{\pgfqpoint{1.410517in}{1.019562in}}%
\pgfpathlineto{\pgfqpoint{1.461840in}{1.043047in}}%
\pgfpathlineto{\pgfqpoint{1.536766in}{1.079789in}}%
\pgfpathlineto{\pgfqpoint{1.633774in}{1.128352in}}%
\pgfpathlineto{\pgfqpoint{1.681340in}{1.150374in}}%
\pgfpathlineto{\pgfqpoint{1.704933in}{1.160385in}}%
\pgfpathlineto{\pgfqpoint{1.728413in}{1.169632in}}%
\pgfpathlineto{\pgfqpoint{1.751791in}{1.177933in}}%
\pgfpathlineto{\pgfqpoint{1.775073in}{1.185248in}}%
\pgfpathlineto{\pgfqpoint{1.798269in}{1.191403in}}%
\pgfpathlineto{\pgfqpoint{1.821384in}{1.196213in}}%
\pgfpathlineto{\pgfqpoint{1.844426in}{1.199688in}}%
\pgfpathlineto{\pgfqpoint{1.867400in}{1.201913in}}%
\pgfpathlineto{\pgfqpoint{1.890311in}{1.202957in}}%
\pgfpathlineto{\pgfqpoint{1.913165in}{1.202893in}}%
\pgfpathlineto{\pgfqpoint{1.935966in}{1.201858in}}%
\pgfpathlineto{\pgfqpoint{2.026719in}{1.194302in}}%
\pgfpathlineto{\pgfqpoint{2.049312in}{1.193363in}}%
\pgfpathlineto{\pgfqpoint{2.094401in}{1.192391in}}%
\pgfpathlineto{\pgfqpoint{2.116902in}{1.190034in}}%
\pgfpathlineto{\pgfqpoint{2.139378in}{1.185473in}}%
\pgfpathlineto{\pgfqpoint{2.205858in}{1.169637in}}%
\pgfpathlineto{\pgfqpoint{2.244728in}{1.161685in}}%
\pgfpathlineto{\pgfqpoint{2.262357in}{1.160014in}}%
\pgfpathlineto{\pgfqpoint{2.294675in}{1.159569in}}%
\pgfpathlineto{\pgfqpoint{2.309568in}{1.160055in}}%
\pgfpathlineto{\pgfqpoint{2.323728in}{1.161286in}}%
\pgfpathlineto{\pgfqpoint{2.337224in}{1.161794in}}%
\pgfpathlineto{\pgfqpoint{2.350117in}{1.161116in}}%
\pgfpathlineto{\pgfqpoint{2.362456in}{1.159201in}}%
\pgfpathlineto{\pgfqpoint{2.374288in}{1.156600in}}%
\pgfpathlineto{\pgfqpoint{2.385654in}{1.153264in}}%
\pgfpathlineto{\pgfqpoint{2.396587in}{1.148564in}}%
\pgfpathlineto{\pgfqpoint{2.417283in}{1.138103in}}%
\pgfpathlineto{\pgfqpoint{2.427098in}{1.133591in}}%
\pgfpathlineto{\pgfqpoint{2.436590in}{1.130116in}}%
\pgfpathlineto{\pgfqpoint{2.495491in}{1.113263in}}%
\pgfpathlineto{\pgfqpoint{2.544365in}{1.095219in}}%
\pgfpathlineto{\pgfqpoint{2.550716in}{1.092118in}}%
\pgfpathlineto{\pgfqpoint{2.556930in}{1.088389in}}%
\pgfpathlineto{\pgfqpoint{2.568970in}{1.079169in}}%
\pgfpathlineto{\pgfqpoint{2.580526in}{1.069851in}}%
\pgfpathlineto{\pgfqpoint{2.591637in}{1.058874in}}%
\pgfpathlineto{\pgfqpoint{2.607537in}{1.042851in}}%
\pgfpathlineto{\pgfqpoint{2.636925in}{1.013684in}}%
\pgfpathlineto{\pgfqpoint{2.646090in}{1.006366in}}%
\pgfpathlineto{\pgfqpoint{2.654972in}{0.997598in}}%
\pgfpathlineto{\pgfqpoint{2.676048in}{0.974896in}}%
\pgfpathlineto{\pgfqpoint{2.691864in}{0.956101in}}%
\pgfpathlineto{\pgfqpoint{2.703181in}{0.944974in}}%
\pgfpathlineto{\pgfqpoint{2.710486in}{0.938028in}}%
\pgfpathlineto{\pgfqpoint{2.724562in}{0.923804in}}%
\pgfpathlineto{\pgfqpoint{2.744465in}{0.906666in}}%
\pgfpathlineto{\pgfqpoint{2.753923in}{0.900605in}}%
\pgfpathlineto{\pgfqpoint{2.771955in}{0.890779in}}%
\pgfpathlineto{\pgfqpoint{2.780564in}{0.887686in}}%
\pgfpathlineto{\pgfqpoint{2.794365in}{0.884478in}}%
\pgfpathlineto{\pgfqpoint{2.807534in}{0.882518in}}%
\pgfpathlineto{\pgfqpoint{2.829820in}{0.881035in}}%
\pgfpathlineto{\pgfqpoint{2.854905in}{0.880496in}}%
\pgfpathlineto{\pgfqpoint{2.939459in}{0.880277in}}%
\pgfpathlineto{\pgfqpoint{3.043110in}{0.880267in}}%
\pgfpathlineto{\pgfqpoint{3.043110in}{0.880267in}}%
\pgfusepath{stroke}%
\end{pgfscope}%
\begin{pgfscope}%
\pgfpathrectangle{\pgfqpoint{0.650000in}{0.420000in}}{\pgfqpoint{2.388110in}{1.994488in}}%
\pgfusepath{clip}%
\pgfsetrectcap%
\pgfsetroundjoin%
\pgfsetlinewidth{1.003750pt}%
\definecolor{currentstroke}{rgb}{0.192157,0.211765,0.584314}%
\pgfsetstrokecolor{currentstroke}%
\pgfsetdash{}{0pt}%
\pgfpathmoveto{\pgfqpoint{0.645000in}{0.880246in}}%
\pgfpathlineto{\pgfqpoint{0.942836in}{0.879819in}}%
\pgfpathlineto{\pgfqpoint{1.057113in}{0.878778in}}%
\pgfpathlineto{\pgfqpoint{1.124662in}{0.877288in}}%
\pgfpathlineto{\pgfqpoint{1.187505in}{0.874725in}}%
\pgfpathlineto{\pgfqpoint{1.217543in}{0.872906in}}%
\pgfpathlineto{\pgfqpoint{1.275397in}{0.867673in}}%
\pgfpathlineto{\pgfqpoint{1.303382in}{0.864241in}}%
\pgfpathlineto{\pgfqpoint{1.330833in}{0.860189in}}%
\pgfpathlineto{\pgfqpoint{1.357807in}{0.855521in}}%
\pgfpathlineto{\pgfqpoint{1.384354in}{0.850192in}}%
\pgfpathlineto{\pgfqpoint{1.410517in}{0.844240in}}%
\pgfpathlineto{\pgfqpoint{1.461840in}{0.830882in}}%
\pgfpathlineto{\pgfqpoint{1.585622in}{0.794982in}}%
\pgfpathlineto{\pgfqpoint{1.633774in}{0.782515in}}%
\pgfpathlineto{\pgfqpoint{1.657624in}{0.776963in}}%
\pgfpathlineto{\pgfqpoint{1.704933in}{0.767690in}}%
\pgfpathlineto{\pgfqpoint{1.751791in}{0.760411in}}%
\pgfpathlineto{\pgfqpoint{1.798269in}{0.754977in}}%
\pgfpathlineto{\pgfqpoint{1.844426in}{0.750980in}}%
\pgfpathlineto{\pgfqpoint{1.890311in}{0.748094in}}%
\pgfpathlineto{\pgfqpoint{1.935966in}{0.746351in}}%
\pgfpathlineto{\pgfqpoint{2.004092in}{0.745102in}}%
\pgfpathlineto{\pgfqpoint{2.071871in}{0.745478in}}%
\pgfpathlineto{\pgfqpoint{2.139378in}{0.746758in}}%
\pgfpathlineto{\pgfqpoint{2.225949in}{0.749987in}}%
\pgfpathlineto{\pgfqpoint{2.262357in}{0.752052in}}%
\pgfpathlineto{\pgfqpoint{2.278969in}{0.753625in}}%
\pgfpathlineto{\pgfqpoint{2.294675in}{0.755691in}}%
\pgfpathlineto{\pgfqpoint{2.323728in}{0.761213in}}%
\pgfpathlineto{\pgfqpoint{2.350117in}{0.766783in}}%
\pgfpathlineto{\pgfqpoint{2.374288in}{0.770628in}}%
\pgfpathlineto{\pgfqpoint{2.396587in}{0.772911in}}%
\pgfpathlineto{\pgfqpoint{2.436590in}{0.776489in}}%
\pgfpathlineto{\pgfqpoint{2.454683in}{0.779226in}}%
\pgfpathlineto{\pgfqpoint{2.503002in}{0.789350in}}%
\pgfpathlineto{\pgfqpoint{2.544365in}{0.796351in}}%
\pgfpathlineto{\pgfqpoint{2.563013in}{0.801175in}}%
\pgfpathlineto{\pgfqpoint{2.586135in}{0.808762in}}%
\pgfpathlineto{\pgfqpoint{2.607537in}{0.816753in}}%
\pgfpathlineto{\pgfqpoint{2.641544in}{0.830567in}}%
\pgfpathlineto{\pgfqpoint{2.724562in}{0.859663in}}%
\pgfpathlineto{\pgfqpoint{2.747652in}{0.869417in}}%
\pgfpathlineto{\pgfqpoint{2.763080in}{0.873556in}}%
\pgfpathlineto{\pgfqpoint{2.777723in}{0.876738in}}%
\pgfpathlineto{\pgfqpoint{2.794365in}{0.878807in}}%
\pgfpathlineto{\pgfqpoint{2.810097in}{0.879873in}}%
\pgfpathlineto{\pgfqpoint{2.846034in}{0.880246in}}%
\pgfpathlineto{\pgfqpoint{3.043110in}{0.880159in}}%
\pgfpathlineto{\pgfqpoint{3.043110in}{0.880159in}}%
\pgfusepath{stroke}%
\end{pgfscope}%
\begin{pgfscope}%
\pgfpathrectangle{\pgfqpoint{0.650000in}{0.420000in}}{\pgfqpoint{2.388110in}{1.994488in}}%
\pgfusepath{clip}%
\pgfsetrectcap%
\pgfsetroundjoin%
\pgfsetlinewidth{1.003750pt}%
\definecolor{currentstroke}{rgb}{0.698039,0.670588,0.823529}%
\pgfsetstrokecolor{currentstroke}%
\pgfsetdash{}{0pt}%
\pgfpathmoveto{\pgfqpoint{0.645000in}{0.887533in}}%
\pgfpathlineto{\pgfqpoint{0.693779in}{0.890060in}}%
\pgfpathlineto{\pgfqpoint{0.743169in}{0.893671in}}%
\pgfpathlineto{\pgfqpoint{0.788495in}{0.898270in}}%
\pgfpathlineto{\pgfqpoint{0.830619in}{0.904057in}}%
\pgfpathlineto{\pgfqpoint{0.870161in}{0.911267in}}%
\pgfpathlineto{\pgfqpoint{0.907582in}{0.920166in}}%
\pgfpathlineto{\pgfqpoint{0.925610in}{0.925341in}}%
\pgfpathlineto{\pgfqpoint{0.943235in}{0.931053in}}%
\pgfpathlineto{\pgfqpoint{0.960487in}{0.937346in}}%
\pgfpathlineto{\pgfqpoint{0.977396in}{0.944265in}}%
\pgfpathlineto{\pgfqpoint{0.993988in}{0.951857in}}%
\pgfpathlineto{\pgfqpoint{1.010285in}{0.960171in}}%
\pgfpathlineto{\pgfqpoint{1.026310in}{0.969259in}}%
\pgfpathlineto{\pgfqpoint{1.042080in}{0.979174in}}%
\pgfpathlineto{\pgfqpoint{1.057613in}{0.989971in}}%
\pgfpathlineto{\pgfqpoint{1.072926in}{1.001706in}}%
\pgfpathlineto{\pgfqpoint{1.088031in}{1.014437in}}%
\pgfpathlineto{\pgfqpoint{1.102943in}{1.028221in}}%
\pgfpathlineto{\pgfqpoint{1.117673in}{1.043118in}}%
\pgfpathlineto{\pgfqpoint{1.132233in}{1.059186in}}%
\pgfpathlineto{\pgfqpoint{1.146632in}{1.076483in}}%
\pgfpathlineto{\pgfqpoint{1.160881in}{1.095063in}}%
\pgfpathlineto{\pgfqpoint{1.174988in}{1.114979in}}%
\pgfpathlineto{\pgfqpoint{1.188960in}{1.136279in}}%
\pgfpathlineto{\pgfqpoint{1.202807in}{1.159005in}}%
\pgfpathlineto{\pgfqpoint{1.216534in}{1.183189in}}%
\pgfpathlineto{\pgfqpoint{1.230148in}{1.208852in}}%
\pgfpathlineto{\pgfqpoint{1.243655in}{1.236001in}}%
\pgfpathlineto{\pgfqpoint{1.257062in}{1.264623in}}%
\pgfpathlineto{\pgfqpoint{1.270372in}{1.294684in}}%
\pgfpathlineto{\pgfqpoint{1.283592in}{1.326122in}}%
\pgfpathlineto{\pgfqpoint{1.309777in}{1.392733in}}%
\pgfpathlineto{\pgfqpoint{1.335650in}{1.463372in}}%
\pgfpathlineto{\pgfqpoint{1.361243in}{1.536667in}}%
\pgfpathlineto{\pgfqpoint{1.436586in}{1.756895in}}%
\pgfpathlineto{\pgfqpoint{1.461293in}{1.823961in}}%
\pgfpathlineto{\pgfqpoint{1.473581in}{1.855207in}}%
\pgfpathlineto{\pgfqpoint{1.485826in}{1.884606in}}%
\pgfpathlineto{\pgfqpoint{1.498032in}{1.911946in}}%
\pgfpathlineto{\pgfqpoint{1.510201in}{1.937054in}}%
\pgfpathlineto{\pgfqpoint{1.522333in}{1.959796in}}%
\pgfpathlineto{\pgfqpoint{1.534430in}{1.980062in}}%
\pgfpathlineto{\pgfqpoint{1.546495in}{1.997769in}}%
\pgfpathlineto{\pgfqpoint{1.558528in}{2.012854in}}%
\pgfpathlineto{\pgfqpoint{1.570530in}{2.025264in}}%
\pgfpathlineto{\pgfqpoint{1.582504in}{2.034977in}}%
\pgfpathlineto{\pgfqpoint{1.594450in}{2.042009in}}%
\pgfpathlineto{\pgfqpoint{1.606369in}{2.046414in}}%
\pgfpathlineto{\pgfqpoint{1.618263in}{2.048268in}}%
\pgfpathlineto{\pgfqpoint{1.630133in}{2.047658in}}%
\pgfpathlineto{\pgfqpoint{1.641980in}{2.044682in}}%
\pgfpathlineto{\pgfqpoint{1.653804in}{2.039443in}}%
\pgfpathlineto{\pgfqpoint{1.665607in}{2.032064in}}%
\pgfpathlineto{\pgfqpoint{1.677389in}{2.022688in}}%
\pgfpathlineto{\pgfqpoint{1.689152in}{2.011469in}}%
\pgfpathlineto{\pgfqpoint{1.700896in}{1.998552in}}%
\pgfpathlineto{\pgfqpoint{1.712622in}{1.984091in}}%
\pgfpathlineto{\pgfqpoint{1.724330in}{1.968239in}}%
\pgfpathlineto{\pgfqpoint{1.736022in}{1.951183in}}%
\pgfpathlineto{\pgfqpoint{1.759360in}{1.914114in}}%
\pgfpathlineto{\pgfqpoint{1.782638in}{1.874286in}}%
\pgfpathlineto{\pgfqpoint{1.863718in}{1.731396in}}%
\pgfpathlineto{\pgfqpoint{1.886786in}{1.693124in}}%
\pgfpathlineto{\pgfqpoint{1.909817in}{1.656851in}}%
\pgfpathlineto{\pgfqpoint{1.932814in}{1.622861in}}%
\pgfpathlineto{\pgfqpoint{1.955780in}{1.591414in}}%
\pgfpathlineto{\pgfqpoint{1.978716in}{1.562731in}}%
\pgfpathlineto{\pgfqpoint{2.001625in}{1.536956in}}%
\pgfpathlineto{\pgfqpoint{2.024510in}{1.513997in}}%
\pgfpathlineto{\pgfqpoint{2.047372in}{1.493702in}}%
\pgfpathlineto{\pgfqpoint{2.070213in}{1.475962in}}%
\pgfpathlineto{\pgfqpoint{2.093034in}{1.460555in}}%
\pgfpathlineto{\pgfqpoint{2.115837in}{1.447243in}}%
\pgfpathlineto{\pgfqpoint{2.138624in}{1.435749in}}%
\pgfpathlineto{\pgfqpoint{2.172776in}{1.420761in}}%
\pgfpathlineto{\pgfqpoint{2.195527in}{1.410279in}}%
\pgfpathlineto{\pgfqpoint{2.218265in}{1.397918in}}%
\pgfpathlineto{\pgfqpoint{2.240991in}{1.384050in}}%
\pgfpathlineto{\pgfqpoint{2.274973in}{1.361893in}}%
\pgfpathlineto{\pgfqpoint{2.307710in}{1.340565in}}%
\pgfpathlineto{\pgfqpoint{2.327947in}{1.329397in}}%
\pgfpathlineto{\pgfqpoint{2.346867in}{1.321075in}}%
\pgfpathlineto{\pgfqpoint{2.397144in}{1.301256in}}%
\pgfpathlineto{\pgfqpoint{2.412125in}{1.293456in}}%
\pgfpathlineto{\pgfqpoint{2.433230in}{1.280751in}}%
\pgfpathlineto{\pgfqpoint{2.446486in}{1.271893in}}%
\pgfpathlineto{\pgfqpoint{2.459159in}{1.261918in}}%
\pgfpathlineto{\pgfqpoint{2.471297in}{1.250483in}}%
\pgfpathlineto{\pgfqpoint{2.494139in}{1.228105in}}%
\pgfpathlineto{\pgfqpoint{2.525326in}{1.200917in}}%
\pgfpathlineto{\pgfqpoint{2.544386in}{1.184787in}}%
\pgfpathlineto{\pgfqpoint{2.574977in}{1.160907in}}%
\pgfpathlineto{\pgfqpoint{2.583152in}{1.153032in}}%
\pgfpathlineto{\pgfqpoint{2.594995in}{1.139809in}}%
\pgfpathlineto{\pgfqpoint{2.620869in}{1.108479in}}%
\pgfpathlineto{\pgfqpoint{2.641335in}{1.081696in}}%
\pgfpathlineto{\pgfqpoint{2.654213in}{1.063292in}}%
\pgfpathlineto{\pgfqpoint{2.669541in}{1.038562in}}%
\pgfpathlineto{\pgfqpoint{2.695228in}{0.994877in}}%
\pgfpathlineto{\pgfqpoint{2.711162in}{0.967857in}}%
\pgfpathlineto{\pgfqpoint{2.726260in}{0.945027in}}%
\pgfpathlineto{\pgfqpoint{2.738264in}{0.928702in}}%
\pgfpathlineto{\pgfqpoint{2.749788in}{0.915241in}}%
\pgfpathlineto{\pgfqpoint{2.758686in}{0.906438in}}%
\pgfpathlineto{\pgfqpoint{2.767317in}{0.899405in}}%
\pgfpathlineto{\pgfqpoint{2.775697in}{0.894308in}}%
\pgfpathlineto{\pgfqpoint{2.783840in}{0.890465in}}%
\pgfpathlineto{\pgfqpoint{2.793705in}{0.887269in}}%
\pgfpathlineto{\pgfqpoint{2.805113in}{0.884947in}}%
\pgfpathlineto{\pgfqpoint{2.823177in}{0.882786in}}%
\pgfpathlineto{\pgfqpoint{2.843457in}{0.881780in}}%
\pgfpathlineto{\pgfqpoint{2.878736in}{0.880969in}}%
\pgfpathlineto{\pgfqpoint{2.995747in}{0.880443in}}%
\pgfpathlineto{\pgfqpoint{3.043110in}{0.880407in}}%
\pgfpathlineto{\pgfqpoint{3.043110in}{0.880407in}}%
\pgfusepath{stroke}%
\end{pgfscope}%
\begin{pgfscope}%
\pgfpathrectangle{\pgfqpoint{0.650000in}{0.420000in}}{\pgfqpoint{2.388110in}{1.994488in}}%
\pgfusepath{clip}%
\pgfsetrectcap%
\pgfsetroundjoin%
\pgfsetlinewidth{1.003750pt}%
\definecolor{currentstroke}{rgb}{0.501961,0.450980,0.674510}%
\pgfsetstrokecolor{currentstroke}%
\pgfsetdash{}{0pt}%
\pgfpathmoveto{\pgfqpoint{0.645000in}{0.880268in}}%
\pgfpathlineto{\pgfqpoint{1.117673in}{0.880789in}}%
\pgfpathlineto{\pgfqpoint{1.230148in}{0.882048in}}%
\pgfpathlineto{\pgfqpoint{1.296725in}{0.883767in}}%
\pgfpathlineto{\pgfqpoint{1.348480in}{0.886006in}}%
\pgfpathlineto{\pgfqpoint{1.399163in}{0.889327in}}%
\pgfpathlineto{\pgfqpoint{1.436586in}{0.892731in}}%
\pgfpathlineto{\pgfqpoint{1.473581in}{0.897091in}}%
\pgfpathlineto{\pgfqpoint{1.510201in}{0.902562in}}%
\pgfpathlineto{\pgfqpoint{1.546495in}{0.909272in}}%
\pgfpathlineto{\pgfqpoint{1.582504in}{0.917327in}}%
\pgfpathlineto{\pgfqpoint{1.618263in}{0.926762in}}%
\pgfpathlineto{\pgfqpoint{1.653804in}{0.937537in}}%
\pgfpathlineto{\pgfqpoint{1.689152in}{0.949516in}}%
\pgfpathlineto{\pgfqpoint{1.724330in}{0.962491in}}%
\pgfpathlineto{\pgfqpoint{1.782638in}{0.985471in}}%
\pgfpathlineto{\pgfqpoint{1.852169in}{1.013008in}}%
\pgfpathlineto{\pgfqpoint{1.886786in}{1.025659in}}%
\pgfpathlineto{\pgfqpoint{1.921320in}{1.036958in}}%
\pgfpathlineto{\pgfqpoint{1.955780in}{1.046662in}}%
\pgfpathlineto{\pgfqpoint{1.990174in}{1.054774in}}%
\pgfpathlineto{\pgfqpoint{2.024510in}{1.061281in}}%
\pgfpathlineto{\pgfqpoint{2.058795in}{1.066281in}}%
\pgfpathlineto{\pgfqpoint{2.093034in}{1.070127in}}%
\pgfpathlineto{\pgfqpoint{2.138624in}{1.073823in}}%
\pgfpathlineto{\pgfqpoint{2.172776in}{1.075597in}}%
\pgfpathlineto{\pgfqpoint{2.206897in}{1.076342in}}%
\pgfpathlineto{\pgfqpoint{2.240991in}{1.075990in}}%
\pgfpathlineto{\pgfqpoint{2.274973in}{1.074690in}}%
\pgfpathlineto{\pgfqpoint{2.318002in}{1.071721in}}%
\pgfpathlineto{\pgfqpoint{2.355880in}{1.067878in}}%
\pgfpathlineto{\pgfqpoint{2.373102in}{1.065330in}}%
\pgfpathlineto{\pgfqpoint{2.397144in}{1.060617in}}%
\pgfpathlineto{\pgfqpoint{2.439934in}{1.051033in}}%
\pgfpathlineto{\pgfqpoint{2.471297in}{1.043052in}}%
\pgfpathlineto{\pgfqpoint{2.499577in}{1.034906in}}%
\pgfpathlineto{\pgfqpoint{2.548960in}{1.018331in}}%
\pgfpathlineto{\pgfqpoint{2.562262in}{1.012618in}}%
\pgfpathlineto{\pgfqpoint{2.594995in}{0.996225in}}%
\pgfpathlineto{\pgfqpoint{2.634673in}{0.975594in}}%
\pgfpathlineto{\pgfqpoint{2.651047in}{0.965039in}}%
\pgfpathlineto{\pgfqpoint{2.672511in}{0.951551in}}%
\pgfpathlineto{\pgfqpoint{2.711162in}{0.928101in}}%
\pgfpathlineto{\pgfqpoint{2.749788in}{0.906603in}}%
\pgfpathlineto{\pgfqpoint{2.769435in}{0.897335in}}%
\pgfpathlineto{\pgfqpoint{2.783840in}{0.891559in}}%
\pgfpathlineto{\pgfqpoint{2.795638in}{0.888108in}}%
\pgfpathlineto{\pgfqpoint{2.808818in}{0.885450in}}%
\pgfpathlineto{\pgfqpoint{2.824923in}{0.883397in}}%
\pgfpathlineto{\pgfqpoint{2.846702in}{0.881923in}}%
\pgfpathlineto{\pgfqpoint{2.881609in}{0.880944in}}%
\pgfpathlineto{\pgfqpoint{2.951114in}{0.880455in}}%
\pgfpathlineto{\pgfqpoint{3.043110in}{0.880326in}}%
\pgfpathlineto{\pgfqpoint{3.043110in}{0.880326in}}%
\pgfusepath{stroke}%
\end{pgfscope}%
\begin{pgfscope}%
\pgfpathrectangle{\pgfqpoint{0.650000in}{0.420000in}}{\pgfqpoint{2.388110in}{1.994488in}}%
\pgfusepath{clip}%
\pgfsetrectcap%
\pgfsetroundjoin%
\pgfsetlinewidth{1.003750pt}%
\definecolor{currentstroke}{rgb}{0.329412,0.152941,0.533333}%
\pgfsetstrokecolor{currentstroke}%
\pgfsetdash{}{0pt}%
\pgfpathmoveto{\pgfqpoint{0.645000in}{0.882988in}}%
\pgfpathlineto{\pgfqpoint{0.743169in}{0.885317in}}%
\pgfpathlineto{\pgfqpoint{0.809915in}{0.887914in}}%
\pgfpathlineto{\pgfqpoint{0.870161in}{0.891306in}}%
\pgfpathlineto{\pgfqpoint{0.925610in}{0.895617in}}%
\pgfpathlineto{\pgfqpoint{0.977396in}{0.900971in}}%
\pgfpathlineto{\pgfqpoint{1.026310in}{0.907477in}}%
\pgfpathlineto{\pgfqpoint{1.072926in}{0.915226in}}%
\pgfpathlineto{\pgfqpoint{1.117673in}{0.924280in}}%
\pgfpathlineto{\pgfqpoint{1.160881in}{0.934660in}}%
\pgfpathlineto{\pgfqpoint{1.202807in}{0.946343in}}%
\pgfpathlineto{\pgfqpoint{1.243655in}{0.959267in}}%
\pgfpathlineto{\pgfqpoint{1.283592in}{0.973330in}}%
\pgfpathlineto{\pgfqpoint{1.322750in}{0.988407in}}%
\pgfpathlineto{\pgfqpoint{1.361243in}{1.004347in}}%
\pgfpathlineto{\pgfqpoint{1.411689in}{1.026657in}}%
\pgfpathlineto{\pgfqpoint{1.461293in}{1.049843in}}%
\pgfpathlineto{\pgfqpoint{1.546495in}{1.091305in}}%
\pgfpathlineto{\pgfqpoint{1.618263in}{1.125793in}}%
\pgfpathlineto{\pgfqpoint{1.653804in}{1.141821in}}%
\pgfpathlineto{\pgfqpoint{1.689152in}{1.156513in}}%
\pgfpathlineto{\pgfqpoint{1.724330in}{1.169530in}}%
\pgfpathlineto{\pgfqpoint{1.747699in}{1.177153in}}%
\pgfpathlineto{\pgfqpoint{1.771006in}{1.183856in}}%
\pgfpathlineto{\pgfqpoint{1.794257in}{1.189570in}}%
\pgfpathlineto{\pgfqpoint{1.817457in}{1.194262in}}%
\pgfpathlineto{\pgfqpoint{1.840609in}{1.197940in}}%
\pgfpathlineto{\pgfqpoint{1.863718in}{1.200600in}}%
\pgfpathlineto{\pgfqpoint{1.886786in}{1.202187in}}%
\pgfpathlineto{\pgfqpoint{1.909817in}{1.202702in}}%
\pgfpathlineto{\pgfqpoint{1.932814in}{1.202262in}}%
\pgfpathlineto{\pgfqpoint{1.967251in}{1.200296in}}%
\pgfpathlineto{\pgfqpoint{2.013070in}{1.196258in}}%
\pgfpathlineto{\pgfqpoint{2.070213in}{1.189891in}}%
\pgfpathlineto{\pgfqpoint{2.104438in}{1.185215in}}%
\pgfpathlineto{\pgfqpoint{2.161396in}{1.176648in}}%
\pgfpathlineto{\pgfqpoint{2.229630in}{1.169852in}}%
\pgfpathlineto{\pgfqpoint{2.252347in}{1.165896in}}%
\pgfpathlineto{\pgfqpoint{2.297074in}{1.156186in}}%
\pgfpathlineto{\pgfqpoint{2.337562in}{1.147587in}}%
\pgfpathlineto{\pgfqpoint{2.373102in}{1.141956in}}%
\pgfpathlineto{\pgfqpoint{2.397144in}{1.137507in}}%
\pgfpathlineto{\pgfqpoint{2.412125in}{1.133708in}}%
\pgfpathlineto{\pgfqpoint{2.426365in}{1.129006in}}%
\pgfpathlineto{\pgfqpoint{2.446486in}{1.120870in}}%
\pgfpathlineto{\pgfqpoint{2.477180in}{1.107219in}}%
\pgfpathlineto{\pgfqpoint{2.520357in}{1.086408in}}%
\pgfpathlineto{\pgfqpoint{2.548960in}{1.071000in}}%
\pgfpathlineto{\pgfqpoint{2.579093in}{1.051324in}}%
\pgfpathlineto{\pgfqpoint{2.591101in}{1.042617in}}%
\pgfpathlineto{\pgfqpoint{2.606370in}{1.030151in}}%
\pgfpathlineto{\pgfqpoint{2.627854in}{1.010802in}}%
\pgfpathlineto{\pgfqpoint{2.644609in}{0.994862in}}%
\pgfpathlineto{\pgfqpoint{2.663507in}{0.974499in}}%
\pgfpathlineto{\pgfqpoint{2.681241in}{0.955564in}}%
\pgfpathlineto{\pgfqpoint{2.703305in}{0.934118in}}%
\pgfpathlineto{\pgfqpoint{2.716283in}{0.922614in}}%
\pgfpathlineto{\pgfqpoint{2.728702in}{0.913009in}}%
\pgfpathlineto{\pgfqpoint{2.742929in}{0.903478in}}%
\pgfpathlineto{\pgfqpoint{2.754271in}{0.896999in}}%
\pgfpathlineto{\pgfqpoint{2.765183in}{0.892036in}}%
\pgfpathlineto{\pgfqpoint{2.775697in}{0.888498in}}%
\pgfpathlineto{\pgfqpoint{2.789799in}{0.885075in}}%
\pgfpathlineto{\pgfqpoint{2.803243in}{0.882748in}}%
\pgfpathlineto{\pgfqpoint{2.817876in}{0.881527in}}%
\pgfpathlineto{\pgfqpoint{2.846702in}{0.880625in}}%
\pgfpathlineto{\pgfqpoint{2.902288in}{0.880309in}}%
\pgfpathlineto{\pgfqpoint{3.043110in}{0.880267in}}%
\pgfpathlineto{\pgfqpoint{3.043110in}{0.880267in}}%
\pgfusepath{stroke}%
\end{pgfscope}%
\begin{pgfscope}%
\pgfpathrectangle{\pgfqpoint{0.650000in}{0.420000in}}{\pgfqpoint{2.388110in}{1.994488in}}%
\pgfusepath{clip}%
\pgfsetrectcap%
\pgfsetroundjoin%
\pgfsetlinewidth{1.003750pt}%
\definecolor{currentstroke}{rgb}{0.176471,0.000000,0.294118}%
\pgfsetstrokecolor{currentstroke}%
\pgfsetdash{}{0pt}%
\pgfpathmoveto{\pgfqpoint{0.645000in}{0.880244in}}%
\pgfpathlineto{\pgfqpoint{0.960487in}{0.879692in}}%
\pgfpathlineto{\pgfqpoint{1.072926in}{0.878401in}}%
\pgfpathlineto{\pgfqpoint{1.146632in}{0.876334in}}%
\pgfpathlineto{\pgfqpoint{1.202807in}{0.873474in}}%
\pgfpathlineto{\pgfqpoint{1.243655in}{0.870334in}}%
\pgfpathlineto{\pgfqpoint{1.283592in}{0.866102in}}%
\pgfpathlineto{\pgfqpoint{1.322750in}{0.860600in}}%
\pgfpathlineto{\pgfqpoint{1.361243in}{0.853741in}}%
\pgfpathlineto{\pgfqpoint{1.399163in}{0.845515in}}%
\pgfpathlineto{\pgfqpoint{1.436586in}{0.836041in}}%
\pgfpathlineto{\pgfqpoint{1.485826in}{0.821957in}}%
\pgfpathlineto{\pgfqpoint{1.606369in}{0.785887in}}%
\pgfpathlineto{\pgfqpoint{1.641980in}{0.776538in}}%
\pgfpathlineto{\pgfqpoint{1.677389in}{0.768335in}}%
\pgfpathlineto{\pgfqpoint{1.712622in}{0.761350in}}%
\pgfpathlineto{\pgfqpoint{1.747699in}{0.755552in}}%
\pgfpathlineto{\pgfqpoint{1.782638in}{0.750878in}}%
\pgfpathlineto{\pgfqpoint{1.817457in}{0.747157in}}%
\pgfpathlineto{\pgfqpoint{1.863718in}{0.743482in}}%
\pgfpathlineto{\pgfqpoint{1.909817in}{0.740946in}}%
\pgfpathlineto{\pgfqpoint{1.967251in}{0.738906in}}%
\pgfpathlineto{\pgfqpoint{2.024510in}{0.737936in}}%
\pgfpathlineto{\pgfqpoint{2.093034in}{0.737959in}}%
\pgfpathlineto{\pgfqpoint{2.172776in}{0.738977in}}%
\pgfpathlineto{\pgfqpoint{2.206897in}{0.740347in}}%
\pgfpathlineto{\pgfqpoint{2.240991in}{0.742784in}}%
\pgfpathlineto{\pgfqpoint{2.364620in}{0.754130in}}%
\pgfpathlineto{\pgfqpoint{2.397144in}{0.758768in}}%
\pgfpathlineto{\pgfqpoint{2.471297in}{0.771052in}}%
\pgfpathlineto{\pgfqpoint{2.557896in}{0.792455in}}%
\pgfpathlineto{\pgfqpoint{2.602628in}{0.806839in}}%
\pgfpathlineto{\pgfqpoint{2.624383in}{0.815300in}}%
\pgfpathlineto{\pgfqpoint{2.654213in}{0.828551in}}%
\pgfpathlineto{\pgfqpoint{2.711162in}{0.854247in}}%
\pgfpathlineto{\pgfqpoint{2.738264in}{0.865501in}}%
\pgfpathlineto{\pgfqpoint{2.754271in}{0.870924in}}%
\pgfpathlineto{\pgfqpoint{2.769435in}{0.874643in}}%
\pgfpathlineto{\pgfqpoint{2.785840in}{0.877516in}}%
\pgfpathlineto{\pgfqpoint{2.803243in}{0.879295in}}%
\pgfpathlineto{\pgfqpoint{2.824923in}{0.880173in}}%
\pgfpathlineto{\pgfqpoint{2.877289in}{0.880270in}}%
\pgfpathlineto{\pgfqpoint{3.043110in}{0.880267in}}%
\pgfpathlineto{\pgfqpoint{3.043110in}{0.880267in}}%
\pgfusepath{stroke}%
\end{pgfscope}%
\begin{pgfscope}%
\pgfsetrectcap%
\pgfsetmiterjoin%
\pgfsetlinewidth{0.803000pt}%
\definecolor{currentstroke}{rgb}{0.000000,0.000000,0.000000}%
\pgfsetstrokecolor{currentstroke}%
\pgfsetdash{}{0pt}%
\pgfpathmoveto{\pgfqpoint{0.650000in}{0.420000in}}%
\pgfpathlineto{\pgfqpoint{0.650000in}{2.414488in}}%
\pgfusepath{stroke}%
\end{pgfscope}%
\begin{pgfscope}%
\pgfsetrectcap%
\pgfsetmiterjoin%
\pgfsetlinewidth{0.803000pt}%
\definecolor{currentstroke}{rgb}{0.000000,0.000000,0.000000}%
\pgfsetstrokecolor{currentstroke}%
\pgfsetdash{}{0pt}%
\pgfpathmoveto{\pgfqpoint{3.038110in}{0.420000in}}%
\pgfpathlineto{\pgfqpoint{3.038110in}{2.414488in}}%
\pgfusepath{stroke}%
\end{pgfscope}%
\begin{pgfscope}%
\pgfsetrectcap%
\pgfsetmiterjoin%
\pgfsetlinewidth{0.803000pt}%
\definecolor{currentstroke}{rgb}{0.000000,0.000000,0.000000}%
\pgfsetstrokecolor{currentstroke}%
\pgfsetdash{}{0pt}%
\pgfpathmoveto{\pgfqpoint{0.650000in}{0.420000in}}%
\pgfpathlineto{\pgfqpoint{3.038110in}{0.420000in}}%
\pgfusepath{stroke}%
\end{pgfscope}%
\begin{pgfscope}%
\pgfsetrectcap%
\pgfsetmiterjoin%
\pgfsetlinewidth{0.803000pt}%
\definecolor{currentstroke}{rgb}{0.000000,0.000000,0.000000}%
\pgfsetstrokecolor{currentstroke}%
\pgfsetdash{}{0pt}%
\pgfpathmoveto{\pgfqpoint{0.650000in}{2.414488in}}%
\pgfpathlineto{\pgfqpoint{3.038110in}{2.414488in}}%
\pgfusepath{stroke}%
\end{pgfscope}%
\begin{pgfscope}%
\pgfsetrectcap%
\pgfsetroundjoin%
\pgfsetlinewidth{0.602250pt}%
\definecolor{currentstroke}{rgb}{0.000000,0.000000,0.000000}%
\pgfsetstrokecolor{currentstroke}%
\pgfsetdash{}{0pt}%
\pgfpathmoveto{\pgfqpoint{2.067925in}{2.275599in}}%
\pgfpathlineto{\pgfqpoint{2.179036in}{2.275599in}}%
\pgfpathlineto{\pgfqpoint{2.290147in}{2.275599in}}%
\pgfusepath{stroke}%
\end{pgfscope}%
\begin{pgfscope}%
\definecolor{textcolor}{rgb}{0.000000,0.000000,0.000000}%
\pgfsetstrokecolor{textcolor}%
\pgfsetfillcolor{textcolor}%
\pgftext[x=2.379036in,y=2.236710in,left,base]{\color{textcolor}{\rmfamily\fontsize{8.000000}{9.600000}\selectfont\catcode`\^=\active\def^{\ifmmode\sp\else\^{}\fi}\catcode`\%=\active\def%{\%}ref}}%
\end{pgfscope}%
\begin{pgfscope}%
\pgfsetrectcap%
\pgfsetroundjoin%
\pgfsetlinewidth{1.003750pt}%
\definecolor{currentstroke}{rgb}{0.670588,0.850980,0.913725}%
\pgfsetstrokecolor{currentstroke}%
\pgfsetdash{}{0pt}%
\pgfpathmoveto{\pgfqpoint{2.067925in}{2.120661in}}%
\pgfpathlineto{\pgfqpoint{2.179036in}{2.120661in}}%
\pgfpathlineto{\pgfqpoint{2.290147in}{2.120661in}}%
\pgfusepath{stroke}%
\end{pgfscope}%
\begin{pgfscope}%
\definecolor{textcolor}{rgb}{0.000000,0.000000,0.000000}%
\pgfsetstrokecolor{textcolor}%
\pgfsetfillcolor{textcolor}%
\pgftext[x=2.379036in,y=2.081772in,left,base]{\color{textcolor}{\rmfamily\fontsize{8.000000}{9.600000}\selectfont\catcode`\^=\active\def^{\ifmmode\sp\else\^{}\fi}\catcode`\%=\active\def%{\%}R100deg90}}%
\end{pgfscope}%
\begin{pgfscope}%
\pgfsetrectcap%
\pgfsetroundjoin%
\pgfsetlinewidth{1.003750pt}%
\definecolor{currentstroke}{rgb}{0.698039,0.670588,0.823529}%
\pgfsetstrokecolor{currentstroke}%
\pgfsetdash{}{0pt}%
\pgfpathmoveto{\pgfqpoint{2.067925in}{1.965723in}}%
\pgfpathlineto{\pgfqpoint{2.179036in}{1.965723in}}%
\pgfpathlineto{\pgfqpoint{2.290147in}{1.965723in}}%
\pgfusepath{stroke}%
\end{pgfscope}%
\begin{pgfscope}%
\definecolor{textcolor}{rgb}{0.000000,0.000000,0.000000}%
\pgfsetstrokecolor{textcolor}%
\pgfsetfillcolor{textcolor}%
\pgftext[x=2.379036in,y=1.926834in,left,base]{\color{textcolor}{\rmfamily\fontsize{8.000000}{9.600000}\selectfont\catcode`\^=\active\def^{\ifmmode\sp\else\^{}\fi}\catcode`\%=\active\def%{\%}R500deg18}}%
\end{pgfscope}%
\end{pgfpicture}%
\makeatother%
\endgroup%

    \label{fig:R500deg18R100deg90_fluct}
\end{subfigure}
\\[-20pt]
\begin{subfigure}[t]{0.495\textwidth}
    \centering
    \subcaption{}
    %% Creator: Matplotlib, PGF backend
%%
%% To include the figure in your LaTeX document, write
%%   \input{<filename>.pgf}
%%
%% Make sure the required packages are loaded in your preamble
%%   \usepackage{pgf}
%%
%% Also ensure that all the required font packages are loaded; for instance,
%% the lmodern package is sometimes necessary when using math font.
%%   \usepackage{lmodern}
%%
%% Figures using additional raster images can only be included by \input if
%% they are in the same directory as the main LaTeX file. For loading figures
%% from other directories you can use the `import` package
%%   \usepackage{import}
%%
%% and then include the figures with
%%   \import{<path to file>}{<filename>.pgf}
%%
%% Matplotlib used the following preamble
%%   \def\mathdefault#1{#1}
%%   \everymath=\expandafter{\the\everymath\displaystyle}
%%   \IfFileExists{scrextend.sty}{
%%     \usepackage[fontsize=11.000000pt]{scrextend}
%%   }{
%%     \renewcommand{\normalsize}{\fontsize{11.000000}{13.200000}\selectfont}
%%     \normalsize
%%   }
%%   
%%   \makeatletter\@ifpackageloaded{underscore}{}{\usepackage[strings]{underscore}}\makeatother
%%
\begingroup%
\makeatletter%
\begin{pgfpicture}%
\pgfpathrectangle{\pgfpointorigin}{\pgfqpoint{3.118110in}{2.494488in}}%
\pgfusepath{use as bounding box, clip}%
\begin{pgfscope}%
\pgfsetbuttcap%
\pgfsetmiterjoin%
\definecolor{currentfill}{rgb}{1.000000,1.000000,1.000000}%
\pgfsetfillcolor{currentfill}%
\pgfsetlinewidth{0.000000pt}%
\definecolor{currentstroke}{rgb}{1.000000,1.000000,1.000000}%
\pgfsetstrokecolor{currentstroke}%
\pgfsetdash{}{0pt}%
\pgfpathmoveto{\pgfqpoint{0.000000in}{0.000000in}}%
\pgfpathlineto{\pgfqpoint{3.118110in}{0.000000in}}%
\pgfpathlineto{\pgfqpoint{3.118110in}{2.494488in}}%
\pgfpathlineto{\pgfqpoint{0.000000in}{2.494488in}}%
\pgfpathlineto{\pgfqpoint{0.000000in}{0.000000in}}%
\pgfpathclose%
\pgfusepath{fill}%
\end{pgfscope}%
\begin{pgfscope}%
\pgfsetbuttcap%
\pgfsetmiterjoin%
\definecolor{currentfill}{rgb}{1.000000,1.000000,1.000000}%
\pgfsetfillcolor{currentfill}%
\pgfsetlinewidth{0.000000pt}%
\definecolor{currentstroke}{rgb}{0.000000,0.000000,0.000000}%
\pgfsetstrokecolor{currentstroke}%
\pgfsetstrokeopacity{0.000000}%
\pgfsetdash{}{0pt}%
\pgfpathmoveto{\pgfqpoint{0.650000in}{0.420000in}}%
\pgfpathlineto{\pgfqpoint{3.038110in}{0.420000in}}%
\pgfpathlineto{\pgfqpoint{3.038110in}{2.414488in}}%
\pgfpathlineto{\pgfqpoint{0.650000in}{2.414488in}}%
\pgfpathlineto{\pgfqpoint{0.650000in}{0.420000in}}%
\pgfpathclose%
\pgfusepath{fill}%
\end{pgfscope}%
\begin{pgfscope}%
\pgfsetbuttcap%
\pgfsetroundjoin%
\definecolor{currentfill}{rgb}{0.000000,0.000000,0.000000}%
\pgfsetfillcolor{currentfill}%
\pgfsetlinewidth{0.501875pt}%
\definecolor{currentstroke}{rgb}{0.000000,0.000000,0.000000}%
\pgfsetstrokecolor{currentstroke}%
\pgfsetdash{}{0pt}%
\pgfsys@defobject{currentmarker}{\pgfqpoint{0.000000in}{0.000000in}}{\pgfqpoint{0.000000in}{0.041667in}}{%
\pgfpathmoveto{\pgfqpoint{0.000000in}{0.000000in}}%
\pgfpathlineto{\pgfqpoint{0.000000in}{0.041667in}}%
\pgfusepath{stroke,fill}%
}%
\begin{pgfscope}%
\pgfsys@transformshift{0.849578in}{0.420000in}%
\pgfsys@useobject{currentmarker}{}%
\end{pgfscope}%
\end{pgfscope}%
\begin{pgfscope}%
\pgfsetbuttcap%
\pgfsetroundjoin%
\definecolor{currentfill}{rgb}{0.000000,0.000000,0.000000}%
\pgfsetfillcolor{currentfill}%
\pgfsetlinewidth{0.501875pt}%
\definecolor{currentstroke}{rgb}{0.000000,0.000000,0.000000}%
\pgfsetstrokecolor{currentstroke}%
\pgfsetdash{}{0pt}%
\pgfsys@defobject{currentmarker}{\pgfqpoint{0.000000in}{-0.041667in}}{\pgfqpoint{0.000000in}{0.000000in}}{%
\pgfpathmoveto{\pgfqpoint{0.000000in}{0.000000in}}%
\pgfpathlineto{\pgfqpoint{0.000000in}{-0.041667in}}%
\pgfusepath{stroke,fill}%
}%
\begin{pgfscope}%
\pgfsys@transformshift{0.849578in}{2.414488in}%
\pgfsys@useobject{currentmarker}{}%
\end{pgfscope}%
\end{pgfscope}%
\begin{pgfscope}%
\definecolor{textcolor}{rgb}{0.000000,0.000000,0.000000}%
\pgfsetstrokecolor{textcolor}%
\pgfsetfillcolor{textcolor}%
\pgftext[x=0.849578in,y=0.371389in,,top]{\color{textcolor}{\rmfamily\fontsize{11.000000}{13.200000}\selectfont\catcode`\^=\active\def^{\ifmmode\sp\else\^{}\fi}\catcode`\%=\active\def%{\%}$\mathdefault{10^{0}}$}}%
\end{pgfscope}%
\begin{pgfscope}%
\pgfsetbuttcap%
\pgfsetroundjoin%
\definecolor{currentfill}{rgb}{0.000000,0.000000,0.000000}%
\pgfsetfillcolor{currentfill}%
\pgfsetlinewidth{0.501875pt}%
\definecolor{currentstroke}{rgb}{0.000000,0.000000,0.000000}%
\pgfsetstrokecolor{currentstroke}%
\pgfsetdash{}{0pt}%
\pgfsys@defobject{currentmarker}{\pgfqpoint{0.000000in}{0.000000in}}{\pgfqpoint{0.000000in}{0.041667in}}{%
\pgfpathmoveto{\pgfqpoint{0.000000in}{0.000000in}}%
\pgfpathlineto{\pgfqpoint{0.000000in}{0.041667in}}%
\pgfusepath{stroke,fill}%
}%
\begin{pgfscope}%
\pgfsys@transformshift{1.512563in}{0.420000in}%
\pgfsys@useobject{currentmarker}{}%
\end{pgfscope}%
\end{pgfscope}%
\begin{pgfscope}%
\pgfsetbuttcap%
\pgfsetroundjoin%
\definecolor{currentfill}{rgb}{0.000000,0.000000,0.000000}%
\pgfsetfillcolor{currentfill}%
\pgfsetlinewidth{0.501875pt}%
\definecolor{currentstroke}{rgb}{0.000000,0.000000,0.000000}%
\pgfsetstrokecolor{currentstroke}%
\pgfsetdash{}{0pt}%
\pgfsys@defobject{currentmarker}{\pgfqpoint{0.000000in}{-0.041667in}}{\pgfqpoint{0.000000in}{0.000000in}}{%
\pgfpathmoveto{\pgfqpoint{0.000000in}{0.000000in}}%
\pgfpathlineto{\pgfqpoint{0.000000in}{-0.041667in}}%
\pgfusepath{stroke,fill}%
}%
\begin{pgfscope}%
\pgfsys@transformshift{1.512563in}{2.414488in}%
\pgfsys@useobject{currentmarker}{}%
\end{pgfscope}%
\end{pgfscope}%
\begin{pgfscope}%
\definecolor{textcolor}{rgb}{0.000000,0.000000,0.000000}%
\pgfsetstrokecolor{textcolor}%
\pgfsetfillcolor{textcolor}%
\pgftext[x=1.512563in,y=0.371389in,,top]{\color{textcolor}{\rmfamily\fontsize{11.000000}{13.200000}\selectfont\catcode`\^=\active\def^{\ifmmode\sp\else\^{}\fi}\catcode`\%=\active\def%{\%}$\mathdefault{10^{1}}$}}%
\end{pgfscope}%
\begin{pgfscope}%
\pgfsetbuttcap%
\pgfsetroundjoin%
\definecolor{currentfill}{rgb}{0.000000,0.000000,0.000000}%
\pgfsetfillcolor{currentfill}%
\pgfsetlinewidth{0.501875pt}%
\definecolor{currentstroke}{rgb}{0.000000,0.000000,0.000000}%
\pgfsetstrokecolor{currentstroke}%
\pgfsetdash{}{0pt}%
\pgfsys@defobject{currentmarker}{\pgfqpoint{0.000000in}{0.000000in}}{\pgfqpoint{0.000000in}{0.041667in}}{%
\pgfpathmoveto{\pgfqpoint{0.000000in}{0.000000in}}%
\pgfpathlineto{\pgfqpoint{0.000000in}{0.041667in}}%
\pgfusepath{stroke,fill}%
}%
\begin{pgfscope}%
\pgfsys@transformshift{2.175547in}{0.420000in}%
\pgfsys@useobject{currentmarker}{}%
\end{pgfscope}%
\end{pgfscope}%
\begin{pgfscope}%
\pgfsetbuttcap%
\pgfsetroundjoin%
\definecolor{currentfill}{rgb}{0.000000,0.000000,0.000000}%
\pgfsetfillcolor{currentfill}%
\pgfsetlinewidth{0.501875pt}%
\definecolor{currentstroke}{rgb}{0.000000,0.000000,0.000000}%
\pgfsetstrokecolor{currentstroke}%
\pgfsetdash{}{0pt}%
\pgfsys@defobject{currentmarker}{\pgfqpoint{0.000000in}{-0.041667in}}{\pgfqpoint{0.000000in}{0.000000in}}{%
\pgfpathmoveto{\pgfqpoint{0.000000in}{0.000000in}}%
\pgfpathlineto{\pgfqpoint{0.000000in}{-0.041667in}}%
\pgfusepath{stroke,fill}%
}%
\begin{pgfscope}%
\pgfsys@transformshift{2.175547in}{2.414488in}%
\pgfsys@useobject{currentmarker}{}%
\end{pgfscope}%
\end{pgfscope}%
\begin{pgfscope}%
\definecolor{textcolor}{rgb}{0.000000,0.000000,0.000000}%
\pgfsetstrokecolor{textcolor}%
\pgfsetfillcolor{textcolor}%
\pgftext[x=2.175547in,y=0.371389in,,top]{\color{textcolor}{\rmfamily\fontsize{11.000000}{13.200000}\selectfont\catcode`\^=\active\def^{\ifmmode\sp\else\^{}\fi}\catcode`\%=\active\def%{\%}$\mathdefault{10^{2}}$}}%
\end{pgfscope}%
\begin{pgfscope}%
\pgfsetbuttcap%
\pgfsetroundjoin%
\definecolor{currentfill}{rgb}{0.000000,0.000000,0.000000}%
\pgfsetfillcolor{currentfill}%
\pgfsetlinewidth{0.501875pt}%
\definecolor{currentstroke}{rgb}{0.000000,0.000000,0.000000}%
\pgfsetstrokecolor{currentstroke}%
\pgfsetdash{}{0pt}%
\pgfsys@defobject{currentmarker}{\pgfqpoint{0.000000in}{0.000000in}}{\pgfqpoint{0.000000in}{0.041667in}}{%
\pgfpathmoveto{\pgfqpoint{0.000000in}{0.000000in}}%
\pgfpathlineto{\pgfqpoint{0.000000in}{0.041667in}}%
\pgfusepath{stroke,fill}%
}%
\begin{pgfscope}%
\pgfsys@transformshift{2.838532in}{0.420000in}%
\pgfsys@useobject{currentmarker}{}%
\end{pgfscope}%
\end{pgfscope}%
\begin{pgfscope}%
\pgfsetbuttcap%
\pgfsetroundjoin%
\definecolor{currentfill}{rgb}{0.000000,0.000000,0.000000}%
\pgfsetfillcolor{currentfill}%
\pgfsetlinewidth{0.501875pt}%
\definecolor{currentstroke}{rgb}{0.000000,0.000000,0.000000}%
\pgfsetstrokecolor{currentstroke}%
\pgfsetdash{}{0pt}%
\pgfsys@defobject{currentmarker}{\pgfqpoint{0.000000in}{-0.041667in}}{\pgfqpoint{0.000000in}{0.000000in}}{%
\pgfpathmoveto{\pgfqpoint{0.000000in}{0.000000in}}%
\pgfpathlineto{\pgfqpoint{0.000000in}{-0.041667in}}%
\pgfusepath{stroke,fill}%
}%
\begin{pgfscope}%
\pgfsys@transformshift{2.838532in}{2.414488in}%
\pgfsys@useobject{currentmarker}{}%
\end{pgfscope}%
\end{pgfscope}%
\begin{pgfscope}%
\definecolor{textcolor}{rgb}{0.000000,0.000000,0.000000}%
\pgfsetstrokecolor{textcolor}%
\pgfsetfillcolor{textcolor}%
\pgftext[x=2.838532in,y=0.371389in,,top]{\color{textcolor}{\rmfamily\fontsize{11.000000}{13.200000}\selectfont\catcode`\^=\active\def^{\ifmmode\sp\else\^{}\fi}\catcode`\%=\active\def%{\%}$\mathdefault{10^{3}}$}}%
\end{pgfscope}%
\begin{pgfscope}%
\pgfsetbuttcap%
\pgfsetroundjoin%
\definecolor{currentfill}{rgb}{0.000000,0.000000,0.000000}%
\pgfsetfillcolor{currentfill}%
\pgfsetlinewidth{0.401500pt}%
\definecolor{currentstroke}{rgb}{0.000000,0.000000,0.000000}%
\pgfsetstrokecolor{currentstroke}%
\pgfsetdash{}{0pt}%
\pgfsys@defobject{currentmarker}{\pgfqpoint{0.000000in}{0.000000in}}{\pgfqpoint{0.000000in}{0.020833in}}{%
\pgfpathmoveto{\pgfqpoint{0.000000in}{0.000000in}}%
\pgfpathlineto{\pgfqpoint{0.000000in}{0.020833in}}%
\pgfusepath{stroke,fill}%
}%
\begin{pgfscope}%
\pgfsys@transformshift{0.650000in}{0.420000in}%
\pgfsys@useobject{currentmarker}{}%
\end{pgfscope}%
\end{pgfscope}%
\begin{pgfscope}%
\pgfsetbuttcap%
\pgfsetroundjoin%
\definecolor{currentfill}{rgb}{0.000000,0.000000,0.000000}%
\pgfsetfillcolor{currentfill}%
\pgfsetlinewidth{0.401500pt}%
\definecolor{currentstroke}{rgb}{0.000000,0.000000,0.000000}%
\pgfsetstrokecolor{currentstroke}%
\pgfsetdash{}{0pt}%
\pgfsys@defobject{currentmarker}{\pgfqpoint{0.000000in}{-0.020833in}}{\pgfqpoint{0.000000in}{0.000000in}}{%
\pgfpathmoveto{\pgfqpoint{0.000000in}{0.000000in}}%
\pgfpathlineto{\pgfqpoint{0.000000in}{-0.020833in}}%
\pgfusepath{stroke,fill}%
}%
\begin{pgfscope}%
\pgfsys@transformshift{0.650000in}{2.414488in}%
\pgfsys@useobject{currentmarker}{}%
\end{pgfscope}%
\end{pgfscope}%
\begin{pgfscope}%
\pgfsetbuttcap%
\pgfsetroundjoin%
\definecolor{currentfill}{rgb}{0.000000,0.000000,0.000000}%
\pgfsetfillcolor{currentfill}%
\pgfsetlinewidth{0.401500pt}%
\definecolor{currentstroke}{rgb}{0.000000,0.000000,0.000000}%
\pgfsetstrokecolor{currentstroke}%
\pgfsetdash{}{0pt}%
\pgfsys@defobject{currentmarker}{\pgfqpoint{0.000000in}{0.000000in}}{\pgfqpoint{0.000000in}{0.020833in}}{%
\pgfpathmoveto{\pgfqpoint{0.000000in}{0.000000in}}%
\pgfpathlineto{\pgfqpoint{0.000000in}{0.020833in}}%
\pgfusepath{stroke,fill}%
}%
\begin{pgfscope}%
\pgfsys@transformshift{0.702496in}{0.420000in}%
\pgfsys@useobject{currentmarker}{}%
\end{pgfscope}%
\end{pgfscope}%
\begin{pgfscope}%
\pgfsetbuttcap%
\pgfsetroundjoin%
\definecolor{currentfill}{rgb}{0.000000,0.000000,0.000000}%
\pgfsetfillcolor{currentfill}%
\pgfsetlinewidth{0.401500pt}%
\definecolor{currentstroke}{rgb}{0.000000,0.000000,0.000000}%
\pgfsetstrokecolor{currentstroke}%
\pgfsetdash{}{0pt}%
\pgfsys@defobject{currentmarker}{\pgfqpoint{0.000000in}{-0.020833in}}{\pgfqpoint{0.000000in}{0.000000in}}{%
\pgfpathmoveto{\pgfqpoint{0.000000in}{0.000000in}}%
\pgfpathlineto{\pgfqpoint{0.000000in}{-0.020833in}}%
\pgfusepath{stroke,fill}%
}%
\begin{pgfscope}%
\pgfsys@transformshift{0.702496in}{2.414488in}%
\pgfsys@useobject{currentmarker}{}%
\end{pgfscope}%
\end{pgfscope}%
\begin{pgfscope}%
\pgfsetbuttcap%
\pgfsetroundjoin%
\definecolor{currentfill}{rgb}{0.000000,0.000000,0.000000}%
\pgfsetfillcolor{currentfill}%
\pgfsetlinewidth{0.401500pt}%
\definecolor{currentstroke}{rgb}{0.000000,0.000000,0.000000}%
\pgfsetstrokecolor{currentstroke}%
\pgfsetdash{}{0pt}%
\pgfsys@defobject{currentmarker}{\pgfqpoint{0.000000in}{0.000000in}}{\pgfqpoint{0.000000in}{0.020833in}}{%
\pgfpathmoveto{\pgfqpoint{0.000000in}{0.000000in}}%
\pgfpathlineto{\pgfqpoint{0.000000in}{0.020833in}}%
\pgfusepath{stroke,fill}%
}%
\begin{pgfscope}%
\pgfsys@transformshift{0.746881in}{0.420000in}%
\pgfsys@useobject{currentmarker}{}%
\end{pgfscope}%
\end{pgfscope}%
\begin{pgfscope}%
\pgfsetbuttcap%
\pgfsetroundjoin%
\definecolor{currentfill}{rgb}{0.000000,0.000000,0.000000}%
\pgfsetfillcolor{currentfill}%
\pgfsetlinewidth{0.401500pt}%
\definecolor{currentstroke}{rgb}{0.000000,0.000000,0.000000}%
\pgfsetstrokecolor{currentstroke}%
\pgfsetdash{}{0pt}%
\pgfsys@defobject{currentmarker}{\pgfqpoint{0.000000in}{-0.020833in}}{\pgfqpoint{0.000000in}{0.000000in}}{%
\pgfpathmoveto{\pgfqpoint{0.000000in}{0.000000in}}%
\pgfpathlineto{\pgfqpoint{0.000000in}{-0.020833in}}%
\pgfusepath{stroke,fill}%
}%
\begin{pgfscope}%
\pgfsys@transformshift{0.746881in}{2.414488in}%
\pgfsys@useobject{currentmarker}{}%
\end{pgfscope}%
\end{pgfscope}%
\begin{pgfscope}%
\pgfsetbuttcap%
\pgfsetroundjoin%
\definecolor{currentfill}{rgb}{0.000000,0.000000,0.000000}%
\pgfsetfillcolor{currentfill}%
\pgfsetlinewidth{0.401500pt}%
\definecolor{currentstroke}{rgb}{0.000000,0.000000,0.000000}%
\pgfsetstrokecolor{currentstroke}%
\pgfsetdash{}{0pt}%
\pgfsys@defobject{currentmarker}{\pgfqpoint{0.000000in}{0.000000in}}{\pgfqpoint{0.000000in}{0.020833in}}{%
\pgfpathmoveto{\pgfqpoint{0.000000in}{0.000000in}}%
\pgfpathlineto{\pgfqpoint{0.000000in}{0.020833in}}%
\pgfusepath{stroke,fill}%
}%
\begin{pgfscope}%
\pgfsys@transformshift{0.785328in}{0.420000in}%
\pgfsys@useobject{currentmarker}{}%
\end{pgfscope}%
\end{pgfscope}%
\begin{pgfscope}%
\pgfsetbuttcap%
\pgfsetroundjoin%
\definecolor{currentfill}{rgb}{0.000000,0.000000,0.000000}%
\pgfsetfillcolor{currentfill}%
\pgfsetlinewidth{0.401500pt}%
\definecolor{currentstroke}{rgb}{0.000000,0.000000,0.000000}%
\pgfsetstrokecolor{currentstroke}%
\pgfsetdash{}{0pt}%
\pgfsys@defobject{currentmarker}{\pgfqpoint{0.000000in}{-0.020833in}}{\pgfqpoint{0.000000in}{0.000000in}}{%
\pgfpathmoveto{\pgfqpoint{0.000000in}{0.000000in}}%
\pgfpathlineto{\pgfqpoint{0.000000in}{-0.020833in}}%
\pgfusepath{stroke,fill}%
}%
\begin{pgfscope}%
\pgfsys@transformshift{0.785328in}{2.414488in}%
\pgfsys@useobject{currentmarker}{}%
\end{pgfscope}%
\end{pgfscope}%
\begin{pgfscope}%
\pgfsetbuttcap%
\pgfsetroundjoin%
\definecolor{currentfill}{rgb}{0.000000,0.000000,0.000000}%
\pgfsetfillcolor{currentfill}%
\pgfsetlinewidth{0.401500pt}%
\definecolor{currentstroke}{rgb}{0.000000,0.000000,0.000000}%
\pgfsetstrokecolor{currentstroke}%
\pgfsetdash{}{0pt}%
\pgfsys@defobject{currentmarker}{\pgfqpoint{0.000000in}{0.000000in}}{\pgfqpoint{0.000000in}{0.020833in}}{%
\pgfpathmoveto{\pgfqpoint{0.000000in}{0.000000in}}%
\pgfpathlineto{\pgfqpoint{0.000000in}{0.020833in}}%
\pgfusepath{stroke,fill}%
}%
\begin{pgfscope}%
\pgfsys@transformshift{0.819242in}{0.420000in}%
\pgfsys@useobject{currentmarker}{}%
\end{pgfscope}%
\end{pgfscope}%
\begin{pgfscope}%
\pgfsetbuttcap%
\pgfsetroundjoin%
\definecolor{currentfill}{rgb}{0.000000,0.000000,0.000000}%
\pgfsetfillcolor{currentfill}%
\pgfsetlinewidth{0.401500pt}%
\definecolor{currentstroke}{rgb}{0.000000,0.000000,0.000000}%
\pgfsetstrokecolor{currentstroke}%
\pgfsetdash{}{0pt}%
\pgfsys@defobject{currentmarker}{\pgfqpoint{0.000000in}{-0.020833in}}{\pgfqpoint{0.000000in}{0.000000in}}{%
\pgfpathmoveto{\pgfqpoint{0.000000in}{0.000000in}}%
\pgfpathlineto{\pgfqpoint{0.000000in}{-0.020833in}}%
\pgfusepath{stroke,fill}%
}%
\begin{pgfscope}%
\pgfsys@transformshift{0.819242in}{2.414488in}%
\pgfsys@useobject{currentmarker}{}%
\end{pgfscope}%
\end{pgfscope}%
\begin{pgfscope}%
\pgfsetbuttcap%
\pgfsetroundjoin%
\definecolor{currentfill}{rgb}{0.000000,0.000000,0.000000}%
\pgfsetfillcolor{currentfill}%
\pgfsetlinewidth{0.401500pt}%
\definecolor{currentstroke}{rgb}{0.000000,0.000000,0.000000}%
\pgfsetstrokecolor{currentstroke}%
\pgfsetdash{}{0pt}%
\pgfsys@defobject{currentmarker}{\pgfqpoint{0.000000in}{0.000000in}}{\pgfqpoint{0.000000in}{0.020833in}}{%
\pgfpathmoveto{\pgfqpoint{0.000000in}{0.000000in}}%
\pgfpathlineto{\pgfqpoint{0.000000in}{0.020833in}}%
\pgfusepath{stroke,fill}%
}%
\begin{pgfscope}%
\pgfsys@transformshift{1.049156in}{0.420000in}%
\pgfsys@useobject{currentmarker}{}%
\end{pgfscope}%
\end{pgfscope}%
\begin{pgfscope}%
\pgfsetbuttcap%
\pgfsetroundjoin%
\definecolor{currentfill}{rgb}{0.000000,0.000000,0.000000}%
\pgfsetfillcolor{currentfill}%
\pgfsetlinewidth{0.401500pt}%
\definecolor{currentstroke}{rgb}{0.000000,0.000000,0.000000}%
\pgfsetstrokecolor{currentstroke}%
\pgfsetdash{}{0pt}%
\pgfsys@defobject{currentmarker}{\pgfqpoint{0.000000in}{-0.020833in}}{\pgfqpoint{0.000000in}{0.000000in}}{%
\pgfpathmoveto{\pgfqpoint{0.000000in}{0.000000in}}%
\pgfpathlineto{\pgfqpoint{0.000000in}{-0.020833in}}%
\pgfusepath{stroke,fill}%
}%
\begin{pgfscope}%
\pgfsys@transformshift{1.049156in}{2.414488in}%
\pgfsys@useobject{currentmarker}{}%
\end{pgfscope}%
\end{pgfscope}%
\begin{pgfscope}%
\pgfsetbuttcap%
\pgfsetroundjoin%
\definecolor{currentfill}{rgb}{0.000000,0.000000,0.000000}%
\pgfsetfillcolor{currentfill}%
\pgfsetlinewidth{0.401500pt}%
\definecolor{currentstroke}{rgb}{0.000000,0.000000,0.000000}%
\pgfsetstrokecolor{currentstroke}%
\pgfsetdash{}{0pt}%
\pgfsys@defobject{currentmarker}{\pgfqpoint{0.000000in}{0.000000in}}{\pgfqpoint{0.000000in}{0.020833in}}{%
\pgfpathmoveto{\pgfqpoint{0.000000in}{0.000000in}}%
\pgfpathlineto{\pgfqpoint{0.000000in}{0.020833in}}%
\pgfusepath{stroke,fill}%
}%
\begin{pgfscope}%
\pgfsys@transformshift{1.165902in}{0.420000in}%
\pgfsys@useobject{currentmarker}{}%
\end{pgfscope}%
\end{pgfscope}%
\begin{pgfscope}%
\pgfsetbuttcap%
\pgfsetroundjoin%
\definecolor{currentfill}{rgb}{0.000000,0.000000,0.000000}%
\pgfsetfillcolor{currentfill}%
\pgfsetlinewidth{0.401500pt}%
\definecolor{currentstroke}{rgb}{0.000000,0.000000,0.000000}%
\pgfsetstrokecolor{currentstroke}%
\pgfsetdash{}{0pt}%
\pgfsys@defobject{currentmarker}{\pgfqpoint{0.000000in}{-0.020833in}}{\pgfqpoint{0.000000in}{0.000000in}}{%
\pgfpathmoveto{\pgfqpoint{0.000000in}{0.000000in}}%
\pgfpathlineto{\pgfqpoint{0.000000in}{-0.020833in}}%
\pgfusepath{stroke,fill}%
}%
\begin{pgfscope}%
\pgfsys@transformshift{1.165902in}{2.414488in}%
\pgfsys@useobject{currentmarker}{}%
\end{pgfscope}%
\end{pgfscope}%
\begin{pgfscope}%
\pgfsetbuttcap%
\pgfsetroundjoin%
\definecolor{currentfill}{rgb}{0.000000,0.000000,0.000000}%
\pgfsetfillcolor{currentfill}%
\pgfsetlinewidth{0.401500pt}%
\definecolor{currentstroke}{rgb}{0.000000,0.000000,0.000000}%
\pgfsetstrokecolor{currentstroke}%
\pgfsetdash{}{0pt}%
\pgfsys@defobject{currentmarker}{\pgfqpoint{0.000000in}{0.000000in}}{\pgfqpoint{0.000000in}{0.020833in}}{%
\pgfpathmoveto{\pgfqpoint{0.000000in}{0.000000in}}%
\pgfpathlineto{\pgfqpoint{0.000000in}{0.020833in}}%
\pgfusepath{stroke,fill}%
}%
\begin{pgfscope}%
\pgfsys@transformshift{1.248735in}{0.420000in}%
\pgfsys@useobject{currentmarker}{}%
\end{pgfscope}%
\end{pgfscope}%
\begin{pgfscope}%
\pgfsetbuttcap%
\pgfsetroundjoin%
\definecolor{currentfill}{rgb}{0.000000,0.000000,0.000000}%
\pgfsetfillcolor{currentfill}%
\pgfsetlinewidth{0.401500pt}%
\definecolor{currentstroke}{rgb}{0.000000,0.000000,0.000000}%
\pgfsetstrokecolor{currentstroke}%
\pgfsetdash{}{0pt}%
\pgfsys@defobject{currentmarker}{\pgfqpoint{0.000000in}{-0.020833in}}{\pgfqpoint{0.000000in}{0.000000in}}{%
\pgfpathmoveto{\pgfqpoint{0.000000in}{0.000000in}}%
\pgfpathlineto{\pgfqpoint{0.000000in}{-0.020833in}}%
\pgfusepath{stroke,fill}%
}%
\begin{pgfscope}%
\pgfsys@transformshift{1.248735in}{2.414488in}%
\pgfsys@useobject{currentmarker}{}%
\end{pgfscope}%
\end{pgfscope}%
\begin{pgfscope}%
\pgfsetbuttcap%
\pgfsetroundjoin%
\definecolor{currentfill}{rgb}{0.000000,0.000000,0.000000}%
\pgfsetfillcolor{currentfill}%
\pgfsetlinewidth{0.401500pt}%
\definecolor{currentstroke}{rgb}{0.000000,0.000000,0.000000}%
\pgfsetstrokecolor{currentstroke}%
\pgfsetdash{}{0pt}%
\pgfsys@defobject{currentmarker}{\pgfqpoint{0.000000in}{0.000000in}}{\pgfqpoint{0.000000in}{0.020833in}}{%
\pgfpathmoveto{\pgfqpoint{0.000000in}{0.000000in}}%
\pgfpathlineto{\pgfqpoint{0.000000in}{0.020833in}}%
\pgfusepath{stroke,fill}%
}%
\begin{pgfscope}%
\pgfsys@transformshift{1.312985in}{0.420000in}%
\pgfsys@useobject{currentmarker}{}%
\end{pgfscope}%
\end{pgfscope}%
\begin{pgfscope}%
\pgfsetbuttcap%
\pgfsetroundjoin%
\definecolor{currentfill}{rgb}{0.000000,0.000000,0.000000}%
\pgfsetfillcolor{currentfill}%
\pgfsetlinewidth{0.401500pt}%
\definecolor{currentstroke}{rgb}{0.000000,0.000000,0.000000}%
\pgfsetstrokecolor{currentstroke}%
\pgfsetdash{}{0pt}%
\pgfsys@defobject{currentmarker}{\pgfqpoint{0.000000in}{-0.020833in}}{\pgfqpoint{0.000000in}{0.000000in}}{%
\pgfpathmoveto{\pgfqpoint{0.000000in}{0.000000in}}%
\pgfpathlineto{\pgfqpoint{0.000000in}{-0.020833in}}%
\pgfusepath{stroke,fill}%
}%
\begin{pgfscope}%
\pgfsys@transformshift{1.312985in}{2.414488in}%
\pgfsys@useobject{currentmarker}{}%
\end{pgfscope}%
\end{pgfscope}%
\begin{pgfscope}%
\pgfsetbuttcap%
\pgfsetroundjoin%
\definecolor{currentfill}{rgb}{0.000000,0.000000,0.000000}%
\pgfsetfillcolor{currentfill}%
\pgfsetlinewidth{0.401500pt}%
\definecolor{currentstroke}{rgb}{0.000000,0.000000,0.000000}%
\pgfsetstrokecolor{currentstroke}%
\pgfsetdash{}{0pt}%
\pgfsys@defobject{currentmarker}{\pgfqpoint{0.000000in}{0.000000in}}{\pgfqpoint{0.000000in}{0.020833in}}{%
\pgfpathmoveto{\pgfqpoint{0.000000in}{0.000000in}}%
\pgfpathlineto{\pgfqpoint{0.000000in}{0.020833in}}%
\pgfusepath{stroke,fill}%
}%
\begin{pgfscope}%
\pgfsys@transformshift{1.365480in}{0.420000in}%
\pgfsys@useobject{currentmarker}{}%
\end{pgfscope}%
\end{pgfscope}%
\begin{pgfscope}%
\pgfsetbuttcap%
\pgfsetroundjoin%
\definecolor{currentfill}{rgb}{0.000000,0.000000,0.000000}%
\pgfsetfillcolor{currentfill}%
\pgfsetlinewidth{0.401500pt}%
\definecolor{currentstroke}{rgb}{0.000000,0.000000,0.000000}%
\pgfsetstrokecolor{currentstroke}%
\pgfsetdash{}{0pt}%
\pgfsys@defobject{currentmarker}{\pgfqpoint{0.000000in}{-0.020833in}}{\pgfqpoint{0.000000in}{0.000000in}}{%
\pgfpathmoveto{\pgfqpoint{0.000000in}{0.000000in}}%
\pgfpathlineto{\pgfqpoint{0.000000in}{-0.020833in}}%
\pgfusepath{stroke,fill}%
}%
\begin{pgfscope}%
\pgfsys@transformshift{1.365480in}{2.414488in}%
\pgfsys@useobject{currentmarker}{}%
\end{pgfscope}%
\end{pgfscope}%
\begin{pgfscope}%
\pgfsetbuttcap%
\pgfsetroundjoin%
\definecolor{currentfill}{rgb}{0.000000,0.000000,0.000000}%
\pgfsetfillcolor{currentfill}%
\pgfsetlinewidth{0.401500pt}%
\definecolor{currentstroke}{rgb}{0.000000,0.000000,0.000000}%
\pgfsetstrokecolor{currentstroke}%
\pgfsetdash{}{0pt}%
\pgfsys@defobject{currentmarker}{\pgfqpoint{0.000000in}{0.000000in}}{\pgfqpoint{0.000000in}{0.020833in}}{%
\pgfpathmoveto{\pgfqpoint{0.000000in}{0.000000in}}%
\pgfpathlineto{\pgfqpoint{0.000000in}{0.020833in}}%
\pgfusepath{stroke,fill}%
}%
\begin{pgfscope}%
\pgfsys@transformshift{1.409865in}{0.420000in}%
\pgfsys@useobject{currentmarker}{}%
\end{pgfscope}%
\end{pgfscope}%
\begin{pgfscope}%
\pgfsetbuttcap%
\pgfsetroundjoin%
\definecolor{currentfill}{rgb}{0.000000,0.000000,0.000000}%
\pgfsetfillcolor{currentfill}%
\pgfsetlinewidth{0.401500pt}%
\definecolor{currentstroke}{rgb}{0.000000,0.000000,0.000000}%
\pgfsetstrokecolor{currentstroke}%
\pgfsetdash{}{0pt}%
\pgfsys@defobject{currentmarker}{\pgfqpoint{0.000000in}{-0.020833in}}{\pgfqpoint{0.000000in}{0.000000in}}{%
\pgfpathmoveto{\pgfqpoint{0.000000in}{0.000000in}}%
\pgfpathlineto{\pgfqpoint{0.000000in}{-0.020833in}}%
\pgfusepath{stroke,fill}%
}%
\begin{pgfscope}%
\pgfsys@transformshift{1.409865in}{2.414488in}%
\pgfsys@useobject{currentmarker}{}%
\end{pgfscope}%
\end{pgfscope}%
\begin{pgfscope}%
\pgfsetbuttcap%
\pgfsetroundjoin%
\definecolor{currentfill}{rgb}{0.000000,0.000000,0.000000}%
\pgfsetfillcolor{currentfill}%
\pgfsetlinewidth{0.401500pt}%
\definecolor{currentstroke}{rgb}{0.000000,0.000000,0.000000}%
\pgfsetstrokecolor{currentstroke}%
\pgfsetdash{}{0pt}%
\pgfsys@defobject{currentmarker}{\pgfqpoint{0.000000in}{0.000000in}}{\pgfqpoint{0.000000in}{0.020833in}}{%
\pgfpathmoveto{\pgfqpoint{0.000000in}{0.000000in}}%
\pgfpathlineto{\pgfqpoint{0.000000in}{0.020833in}}%
\pgfusepath{stroke,fill}%
}%
\begin{pgfscope}%
\pgfsys@transformshift{1.448313in}{0.420000in}%
\pgfsys@useobject{currentmarker}{}%
\end{pgfscope}%
\end{pgfscope}%
\begin{pgfscope}%
\pgfsetbuttcap%
\pgfsetroundjoin%
\definecolor{currentfill}{rgb}{0.000000,0.000000,0.000000}%
\pgfsetfillcolor{currentfill}%
\pgfsetlinewidth{0.401500pt}%
\definecolor{currentstroke}{rgb}{0.000000,0.000000,0.000000}%
\pgfsetstrokecolor{currentstroke}%
\pgfsetdash{}{0pt}%
\pgfsys@defobject{currentmarker}{\pgfqpoint{0.000000in}{-0.020833in}}{\pgfqpoint{0.000000in}{0.000000in}}{%
\pgfpathmoveto{\pgfqpoint{0.000000in}{0.000000in}}%
\pgfpathlineto{\pgfqpoint{0.000000in}{-0.020833in}}%
\pgfusepath{stroke,fill}%
}%
\begin{pgfscope}%
\pgfsys@transformshift{1.448313in}{2.414488in}%
\pgfsys@useobject{currentmarker}{}%
\end{pgfscope}%
\end{pgfscope}%
\begin{pgfscope}%
\pgfsetbuttcap%
\pgfsetroundjoin%
\definecolor{currentfill}{rgb}{0.000000,0.000000,0.000000}%
\pgfsetfillcolor{currentfill}%
\pgfsetlinewidth{0.401500pt}%
\definecolor{currentstroke}{rgb}{0.000000,0.000000,0.000000}%
\pgfsetstrokecolor{currentstroke}%
\pgfsetdash{}{0pt}%
\pgfsys@defobject{currentmarker}{\pgfqpoint{0.000000in}{0.000000in}}{\pgfqpoint{0.000000in}{0.020833in}}{%
\pgfpathmoveto{\pgfqpoint{0.000000in}{0.000000in}}%
\pgfpathlineto{\pgfqpoint{0.000000in}{0.020833in}}%
\pgfusepath{stroke,fill}%
}%
\begin{pgfscope}%
\pgfsys@transformshift{1.482226in}{0.420000in}%
\pgfsys@useobject{currentmarker}{}%
\end{pgfscope}%
\end{pgfscope}%
\begin{pgfscope}%
\pgfsetbuttcap%
\pgfsetroundjoin%
\definecolor{currentfill}{rgb}{0.000000,0.000000,0.000000}%
\pgfsetfillcolor{currentfill}%
\pgfsetlinewidth{0.401500pt}%
\definecolor{currentstroke}{rgb}{0.000000,0.000000,0.000000}%
\pgfsetstrokecolor{currentstroke}%
\pgfsetdash{}{0pt}%
\pgfsys@defobject{currentmarker}{\pgfqpoint{0.000000in}{-0.020833in}}{\pgfqpoint{0.000000in}{0.000000in}}{%
\pgfpathmoveto{\pgfqpoint{0.000000in}{0.000000in}}%
\pgfpathlineto{\pgfqpoint{0.000000in}{-0.020833in}}%
\pgfusepath{stroke,fill}%
}%
\begin{pgfscope}%
\pgfsys@transformshift{1.482226in}{2.414488in}%
\pgfsys@useobject{currentmarker}{}%
\end{pgfscope}%
\end{pgfscope}%
\begin{pgfscope}%
\pgfsetbuttcap%
\pgfsetroundjoin%
\definecolor{currentfill}{rgb}{0.000000,0.000000,0.000000}%
\pgfsetfillcolor{currentfill}%
\pgfsetlinewidth{0.401500pt}%
\definecolor{currentstroke}{rgb}{0.000000,0.000000,0.000000}%
\pgfsetstrokecolor{currentstroke}%
\pgfsetdash{}{0pt}%
\pgfsys@defobject{currentmarker}{\pgfqpoint{0.000000in}{0.000000in}}{\pgfqpoint{0.000000in}{0.020833in}}{%
\pgfpathmoveto{\pgfqpoint{0.000000in}{0.000000in}}%
\pgfpathlineto{\pgfqpoint{0.000000in}{0.020833in}}%
\pgfusepath{stroke,fill}%
}%
\begin{pgfscope}%
\pgfsys@transformshift{1.712141in}{0.420000in}%
\pgfsys@useobject{currentmarker}{}%
\end{pgfscope}%
\end{pgfscope}%
\begin{pgfscope}%
\pgfsetbuttcap%
\pgfsetroundjoin%
\definecolor{currentfill}{rgb}{0.000000,0.000000,0.000000}%
\pgfsetfillcolor{currentfill}%
\pgfsetlinewidth{0.401500pt}%
\definecolor{currentstroke}{rgb}{0.000000,0.000000,0.000000}%
\pgfsetstrokecolor{currentstroke}%
\pgfsetdash{}{0pt}%
\pgfsys@defobject{currentmarker}{\pgfqpoint{0.000000in}{-0.020833in}}{\pgfqpoint{0.000000in}{0.000000in}}{%
\pgfpathmoveto{\pgfqpoint{0.000000in}{0.000000in}}%
\pgfpathlineto{\pgfqpoint{0.000000in}{-0.020833in}}%
\pgfusepath{stroke,fill}%
}%
\begin{pgfscope}%
\pgfsys@transformshift{1.712141in}{2.414488in}%
\pgfsys@useobject{currentmarker}{}%
\end{pgfscope}%
\end{pgfscope}%
\begin{pgfscope}%
\pgfsetbuttcap%
\pgfsetroundjoin%
\definecolor{currentfill}{rgb}{0.000000,0.000000,0.000000}%
\pgfsetfillcolor{currentfill}%
\pgfsetlinewidth{0.401500pt}%
\definecolor{currentstroke}{rgb}{0.000000,0.000000,0.000000}%
\pgfsetstrokecolor{currentstroke}%
\pgfsetdash{}{0pt}%
\pgfsys@defobject{currentmarker}{\pgfqpoint{0.000000in}{0.000000in}}{\pgfqpoint{0.000000in}{0.020833in}}{%
\pgfpathmoveto{\pgfqpoint{0.000000in}{0.000000in}}%
\pgfpathlineto{\pgfqpoint{0.000000in}{0.020833in}}%
\pgfusepath{stroke,fill}%
}%
\begin{pgfscope}%
\pgfsys@transformshift{1.828887in}{0.420000in}%
\pgfsys@useobject{currentmarker}{}%
\end{pgfscope}%
\end{pgfscope}%
\begin{pgfscope}%
\pgfsetbuttcap%
\pgfsetroundjoin%
\definecolor{currentfill}{rgb}{0.000000,0.000000,0.000000}%
\pgfsetfillcolor{currentfill}%
\pgfsetlinewidth{0.401500pt}%
\definecolor{currentstroke}{rgb}{0.000000,0.000000,0.000000}%
\pgfsetstrokecolor{currentstroke}%
\pgfsetdash{}{0pt}%
\pgfsys@defobject{currentmarker}{\pgfqpoint{0.000000in}{-0.020833in}}{\pgfqpoint{0.000000in}{0.000000in}}{%
\pgfpathmoveto{\pgfqpoint{0.000000in}{0.000000in}}%
\pgfpathlineto{\pgfqpoint{0.000000in}{-0.020833in}}%
\pgfusepath{stroke,fill}%
}%
\begin{pgfscope}%
\pgfsys@transformshift{1.828887in}{2.414488in}%
\pgfsys@useobject{currentmarker}{}%
\end{pgfscope}%
\end{pgfscope}%
\begin{pgfscope}%
\pgfsetbuttcap%
\pgfsetroundjoin%
\definecolor{currentfill}{rgb}{0.000000,0.000000,0.000000}%
\pgfsetfillcolor{currentfill}%
\pgfsetlinewidth{0.401500pt}%
\definecolor{currentstroke}{rgb}{0.000000,0.000000,0.000000}%
\pgfsetstrokecolor{currentstroke}%
\pgfsetdash{}{0pt}%
\pgfsys@defobject{currentmarker}{\pgfqpoint{0.000000in}{0.000000in}}{\pgfqpoint{0.000000in}{0.020833in}}{%
\pgfpathmoveto{\pgfqpoint{0.000000in}{0.000000in}}%
\pgfpathlineto{\pgfqpoint{0.000000in}{0.020833in}}%
\pgfusepath{stroke,fill}%
}%
\begin{pgfscope}%
\pgfsys@transformshift{1.911719in}{0.420000in}%
\pgfsys@useobject{currentmarker}{}%
\end{pgfscope}%
\end{pgfscope}%
\begin{pgfscope}%
\pgfsetbuttcap%
\pgfsetroundjoin%
\definecolor{currentfill}{rgb}{0.000000,0.000000,0.000000}%
\pgfsetfillcolor{currentfill}%
\pgfsetlinewidth{0.401500pt}%
\definecolor{currentstroke}{rgb}{0.000000,0.000000,0.000000}%
\pgfsetstrokecolor{currentstroke}%
\pgfsetdash{}{0pt}%
\pgfsys@defobject{currentmarker}{\pgfqpoint{0.000000in}{-0.020833in}}{\pgfqpoint{0.000000in}{0.000000in}}{%
\pgfpathmoveto{\pgfqpoint{0.000000in}{0.000000in}}%
\pgfpathlineto{\pgfqpoint{0.000000in}{-0.020833in}}%
\pgfusepath{stroke,fill}%
}%
\begin{pgfscope}%
\pgfsys@transformshift{1.911719in}{2.414488in}%
\pgfsys@useobject{currentmarker}{}%
\end{pgfscope}%
\end{pgfscope}%
\begin{pgfscope}%
\pgfsetbuttcap%
\pgfsetroundjoin%
\definecolor{currentfill}{rgb}{0.000000,0.000000,0.000000}%
\pgfsetfillcolor{currentfill}%
\pgfsetlinewidth{0.401500pt}%
\definecolor{currentstroke}{rgb}{0.000000,0.000000,0.000000}%
\pgfsetstrokecolor{currentstroke}%
\pgfsetdash{}{0pt}%
\pgfsys@defobject{currentmarker}{\pgfqpoint{0.000000in}{0.000000in}}{\pgfqpoint{0.000000in}{0.020833in}}{%
\pgfpathmoveto{\pgfqpoint{0.000000in}{0.000000in}}%
\pgfpathlineto{\pgfqpoint{0.000000in}{0.020833in}}%
\pgfusepath{stroke,fill}%
}%
\begin{pgfscope}%
\pgfsys@transformshift{1.975969in}{0.420000in}%
\pgfsys@useobject{currentmarker}{}%
\end{pgfscope}%
\end{pgfscope}%
\begin{pgfscope}%
\pgfsetbuttcap%
\pgfsetroundjoin%
\definecolor{currentfill}{rgb}{0.000000,0.000000,0.000000}%
\pgfsetfillcolor{currentfill}%
\pgfsetlinewidth{0.401500pt}%
\definecolor{currentstroke}{rgb}{0.000000,0.000000,0.000000}%
\pgfsetstrokecolor{currentstroke}%
\pgfsetdash{}{0pt}%
\pgfsys@defobject{currentmarker}{\pgfqpoint{0.000000in}{-0.020833in}}{\pgfqpoint{0.000000in}{0.000000in}}{%
\pgfpathmoveto{\pgfqpoint{0.000000in}{0.000000in}}%
\pgfpathlineto{\pgfqpoint{0.000000in}{-0.020833in}}%
\pgfusepath{stroke,fill}%
}%
\begin{pgfscope}%
\pgfsys@transformshift{1.975969in}{2.414488in}%
\pgfsys@useobject{currentmarker}{}%
\end{pgfscope}%
\end{pgfscope}%
\begin{pgfscope}%
\pgfsetbuttcap%
\pgfsetroundjoin%
\definecolor{currentfill}{rgb}{0.000000,0.000000,0.000000}%
\pgfsetfillcolor{currentfill}%
\pgfsetlinewidth{0.401500pt}%
\definecolor{currentstroke}{rgb}{0.000000,0.000000,0.000000}%
\pgfsetstrokecolor{currentstroke}%
\pgfsetdash{}{0pt}%
\pgfsys@defobject{currentmarker}{\pgfqpoint{0.000000in}{0.000000in}}{\pgfqpoint{0.000000in}{0.020833in}}{%
\pgfpathmoveto{\pgfqpoint{0.000000in}{0.000000in}}%
\pgfpathlineto{\pgfqpoint{0.000000in}{0.020833in}}%
\pgfusepath{stroke,fill}%
}%
\begin{pgfscope}%
\pgfsys@transformshift{2.028465in}{0.420000in}%
\pgfsys@useobject{currentmarker}{}%
\end{pgfscope}%
\end{pgfscope}%
\begin{pgfscope}%
\pgfsetbuttcap%
\pgfsetroundjoin%
\definecolor{currentfill}{rgb}{0.000000,0.000000,0.000000}%
\pgfsetfillcolor{currentfill}%
\pgfsetlinewidth{0.401500pt}%
\definecolor{currentstroke}{rgb}{0.000000,0.000000,0.000000}%
\pgfsetstrokecolor{currentstroke}%
\pgfsetdash{}{0pt}%
\pgfsys@defobject{currentmarker}{\pgfqpoint{0.000000in}{-0.020833in}}{\pgfqpoint{0.000000in}{0.000000in}}{%
\pgfpathmoveto{\pgfqpoint{0.000000in}{0.000000in}}%
\pgfpathlineto{\pgfqpoint{0.000000in}{-0.020833in}}%
\pgfusepath{stroke,fill}%
}%
\begin{pgfscope}%
\pgfsys@transformshift{2.028465in}{2.414488in}%
\pgfsys@useobject{currentmarker}{}%
\end{pgfscope}%
\end{pgfscope}%
\begin{pgfscope}%
\pgfsetbuttcap%
\pgfsetroundjoin%
\definecolor{currentfill}{rgb}{0.000000,0.000000,0.000000}%
\pgfsetfillcolor{currentfill}%
\pgfsetlinewidth{0.401500pt}%
\definecolor{currentstroke}{rgb}{0.000000,0.000000,0.000000}%
\pgfsetstrokecolor{currentstroke}%
\pgfsetdash{}{0pt}%
\pgfsys@defobject{currentmarker}{\pgfqpoint{0.000000in}{0.000000in}}{\pgfqpoint{0.000000in}{0.020833in}}{%
\pgfpathmoveto{\pgfqpoint{0.000000in}{0.000000in}}%
\pgfpathlineto{\pgfqpoint{0.000000in}{0.020833in}}%
\pgfusepath{stroke,fill}%
}%
\begin{pgfscope}%
\pgfsys@transformshift{2.072850in}{0.420000in}%
\pgfsys@useobject{currentmarker}{}%
\end{pgfscope}%
\end{pgfscope}%
\begin{pgfscope}%
\pgfsetbuttcap%
\pgfsetroundjoin%
\definecolor{currentfill}{rgb}{0.000000,0.000000,0.000000}%
\pgfsetfillcolor{currentfill}%
\pgfsetlinewidth{0.401500pt}%
\definecolor{currentstroke}{rgb}{0.000000,0.000000,0.000000}%
\pgfsetstrokecolor{currentstroke}%
\pgfsetdash{}{0pt}%
\pgfsys@defobject{currentmarker}{\pgfqpoint{0.000000in}{-0.020833in}}{\pgfqpoint{0.000000in}{0.000000in}}{%
\pgfpathmoveto{\pgfqpoint{0.000000in}{0.000000in}}%
\pgfpathlineto{\pgfqpoint{0.000000in}{-0.020833in}}%
\pgfusepath{stroke,fill}%
}%
\begin{pgfscope}%
\pgfsys@transformshift{2.072850in}{2.414488in}%
\pgfsys@useobject{currentmarker}{}%
\end{pgfscope}%
\end{pgfscope}%
\begin{pgfscope}%
\pgfsetbuttcap%
\pgfsetroundjoin%
\definecolor{currentfill}{rgb}{0.000000,0.000000,0.000000}%
\pgfsetfillcolor{currentfill}%
\pgfsetlinewidth{0.401500pt}%
\definecolor{currentstroke}{rgb}{0.000000,0.000000,0.000000}%
\pgfsetstrokecolor{currentstroke}%
\pgfsetdash{}{0pt}%
\pgfsys@defobject{currentmarker}{\pgfqpoint{0.000000in}{0.000000in}}{\pgfqpoint{0.000000in}{0.020833in}}{%
\pgfpathmoveto{\pgfqpoint{0.000000in}{0.000000in}}%
\pgfpathlineto{\pgfqpoint{0.000000in}{0.020833in}}%
\pgfusepath{stroke,fill}%
}%
\begin{pgfscope}%
\pgfsys@transformshift{2.111298in}{0.420000in}%
\pgfsys@useobject{currentmarker}{}%
\end{pgfscope}%
\end{pgfscope}%
\begin{pgfscope}%
\pgfsetbuttcap%
\pgfsetroundjoin%
\definecolor{currentfill}{rgb}{0.000000,0.000000,0.000000}%
\pgfsetfillcolor{currentfill}%
\pgfsetlinewidth{0.401500pt}%
\definecolor{currentstroke}{rgb}{0.000000,0.000000,0.000000}%
\pgfsetstrokecolor{currentstroke}%
\pgfsetdash{}{0pt}%
\pgfsys@defobject{currentmarker}{\pgfqpoint{0.000000in}{-0.020833in}}{\pgfqpoint{0.000000in}{0.000000in}}{%
\pgfpathmoveto{\pgfqpoint{0.000000in}{0.000000in}}%
\pgfpathlineto{\pgfqpoint{0.000000in}{-0.020833in}}%
\pgfusepath{stroke,fill}%
}%
\begin{pgfscope}%
\pgfsys@transformshift{2.111298in}{2.414488in}%
\pgfsys@useobject{currentmarker}{}%
\end{pgfscope}%
\end{pgfscope}%
\begin{pgfscope}%
\pgfsetbuttcap%
\pgfsetroundjoin%
\definecolor{currentfill}{rgb}{0.000000,0.000000,0.000000}%
\pgfsetfillcolor{currentfill}%
\pgfsetlinewidth{0.401500pt}%
\definecolor{currentstroke}{rgb}{0.000000,0.000000,0.000000}%
\pgfsetstrokecolor{currentstroke}%
\pgfsetdash{}{0pt}%
\pgfsys@defobject{currentmarker}{\pgfqpoint{0.000000in}{0.000000in}}{\pgfqpoint{0.000000in}{0.020833in}}{%
\pgfpathmoveto{\pgfqpoint{0.000000in}{0.000000in}}%
\pgfpathlineto{\pgfqpoint{0.000000in}{0.020833in}}%
\pgfusepath{stroke,fill}%
}%
\begin{pgfscope}%
\pgfsys@transformshift{2.145211in}{0.420000in}%
\pgfsys@useobject{currentmarker}{}%
\end{pgfscope}%
\end{pgfscope}%
\begin{pgfscope}%
\pgfsetbuttcap%
\pgfsetroundjoin%
\definecolor{currentfill}{rgb}{0.000000,0.000000,0.000000}%
\pgfsetfillcolor{currentfill}%
\pgfsetlinewidth{0.401500pt}%
\definecolor{currentstroke}{rgb}{0.000000,0.000000,0.000000}%
\pgfsetstrokecolor{currentstroke}%
\pgfsetdash{}{0pt}%
\pgfsys@defobject{currentmarker}{\pgfqpoint{0.000000in}{-0.020833in}}{\pgfqpoint{0.000000in}{0.000000in}}{%
\pgfpathmoveto{\pgfqpoint{0.000000in}{0.000000in}}%
\pgfpathlineto{\pgfqpoint{0.000000in}{-0.020833in}}%
\pgfusepath{stroke,fill}%
}%
\begin{pgfscope}%
\pgfsys@transformshift{2.145211in}{2.414488in}%
\pgfsys@useobject{currentmarker}{}%
\end{pgfscope}%
\end{pgfscope}%
\begin{pgfscope}%
\pgfsetbuttcap%
\pgfsetroundjoin%
\definecolor{currentfill}{rgb}{0.000000,0.000000,0.000000}%
\pgfsetfillcolor{currentfill}%
\pgfsetlinewidth{0.401500pt}%
\definecolor{currentstroke}{rgb}{0.000000,0.000000,0.000000}%
\pgfsetstrokecolor{currentstroke}%
\pgfsetdash{}{0pt}%
\pgfsys@defobject{currentmarker}{\pgfqpoint{0.000000in}{0.000000in}}{\pgfqpoint{0.000000in}{0.020833in}}{%
\pgfpathmoveto{\pgfqpoint{0.000000in}{0.000000in}}%
\pgfpathlineto{\pgfqpoint{0.000000in}{0.020833in}}%
\pgfusepath{stroke,fill}%
}%
\begin{pgfscope}%
\pgfsys@transformshift{2.375126in}{0.420000in}%
\pgfsys@useobject{currentmarker}{}%
\end{pgfscope}%
\end{pgfscope}%
\begin{pgfscope}%
\pgfsetbuttcap%
\pgfsetroundjoin%
\definecolor{currentfill}{rgb}{0.000000,0.000000,0.000000}%
\pgfsetfillcolor{currentfill}%
\pgfsetlinewidth{0.401500pt}%
\definecolor{currentstroke}{rgb}{0.000000,0.000000,0.000000}%
\pgfsetstrokecolor{currentstroke}%
\pgfsetdash{}{0pt}%
\pgfsys@defobject{currentmarker}{\pgfqpoint{0.000000in}{-0.020833in}}{\pgfqpoint{0.000000in}{0.000000in}}{%
\pgfpathmoveto{\pgfqpoint{0.000000in}{0.000000in}}%
\pgfpathlineto{\pgfqpoint{0.000000in}{-0.020833in}}%
\pgfusepath{stroke,fill}%
}%
\begin{pgfscope}%
\pgfsys@transformshift{2.375126in}{2.414488in}%
\pgfsys@useobject{currentmarker}{}%
\end{pgfscope}%
\end{pgfscope}%
\begin{pgfscope}%
\pgfsetbuttcap%
\pgfsetroundjoin%
\definecolor{currentfill}{rgb}{0.000000,0.000000,0.000000}%
\pgfsetfillcolor{currentfill}%
\pgfsetlinewidth{0.401500pt}%
\definecolor{currentstroke}{rgb}{0.000000,0.000000,0.000000}%
\pgfsetstrokecolor{currentstroke}%
\pgfsetdash{}{0pt}%
\pgfsys@defobject{currentmarker}{\pgfqpoint{0.000000in}{0.000000in}}{\pgfqpoint{0.000000in}{0.020833in}}{%
\pgfpathmoveto{\pgfqpoint{0.000000in}{0.000000in}}%
\pgfpathlineto{\pgfqpoint{0.000000in}{0.020833in}}%
\pgfusepath{stroke,fill}%
}%
\begin{pgfscope}%
\pgfsys@transformshift{2.491871in}{0.420000in}%
\pgfsys@useobject{currentmarker}{}%
\end{pgfscope}%
\end{pgfscope}%
\begin{pgfscope}%
\pgfsetbuttcap%
\pgfsetroundjoin%
\definecolor{currentfill}{rgb}{0.000000,0.000000,0.000000}%
\pgfsetfillcolor{currentfill}%
\pgfsetlinewidth{0.401500pt}%
\definecolor{currentstroke}{rgb}{0.000000,0.000000,0.000000}%
\pgfsetstrokecolor{currentstroke}%
\pgfsetdash{}{0pt}%
\pgfsys@defobject{currentmarker}{\pgfqpoint{0.000000in}{-0.020833in}}{\pgfqpoint{0.000000in}{0.000000in}}{%
\pgfpathmoveto{\pgfqpoint{0.000000in}{0.000000in}}%
\pgfpathlineto{\pgfqpoint{0.000000in}{-0.020833in}}%
\pgfusepath{stroke,fill}%
}%
\begin{pgfscope}%
\pgfsys@transformshift{2.491871in}{2.414488in}%
\pgfsys@useobject{currentmarker}{}%
\end{pgfscope}%
\end{pgfscope}%
\begin{pgfscope}%
\pgfsetbuttcap%
\pgfsetroundjoin%
\definecolor{currentfill}{rgb}{0.000000,0.000000,0.000000}%
\pgfsetfillcolor{currentfill}%
\pgfsetlinewidth{0.401500pt}%
\definecolor{currentstroke}{rgb}{0.000000,0.000000,0.000000}%
\pgfsetstrokecolor{currentstroke}%
\pgfsetdash{}{0pt}%
\pgfsys@defobject{currentmarker}{\pgfqpoint{0.000000in}{0.000000in}}{\pgfqpoint{0.000000in}{0.020833in}}{%
\pgfpathmoveto{\pgfqpoint{0.000000in}{0.000000in}}%
\pgfpathlineto{\pgfqpoint{0.000000in}{0.020833in}}%
\pgfusepath{stroke,fill}%
}%
\begin{pgfscope}%
\pgfsys@transformshift{2.574704in}{0.420000in}%
\pgfsys@useobject{currentmarker}{}%
\end{pgfscope}%
\end{pgfscope}%
\begin{pgfscope}%
\pgfsetbuttcap%
\pgfsetroundjoin%
\definecolor{currentfill}{rgb}{0.000000,0.000000,0.000000}%
\pgfsetfillcolor{currentfill}%
\pgfsetlinewidth{0.401500pt}%
\definecolor{currentstroke}{rgb}{0.000000,0.000000,0.000000}%
\pgfsetstrokecolor{currentstroke}%
\pgfsetdash{}{0pt}%
\pgfsys@defobject{currentmarker}{\pgfqpoint{0.000000in}{-0.020833in}}{\pgfqpoint{0.000000in}{0.000000in}}{%
\pgfpathmoveto{\pgfqpoint{0.000000in}{0.000000in}}%
\pgfpathlineto{\pgfqpoint{0.000000in}{-0.020833in}}%
\pgfusepath{stroke,fill}%
}%
\begin{pgfscope}%
\pgfsys@transformshift{2.574704in}{2.414488in}%
\pgfsys@useobject{currentmarker}{}%
\end{pgfscope}%
\end{pgfscope}%
\begin{pgfscope}%
\pgfsetbuttcap%
\pgfsetroundjoin%
\definecolor{currentfill}{rgb}{0.000000,0.000000,0.000000}%
\pgfsetfillcolor{currentfill}%
\pgfsetlinewidth{0.401500pt}%
\definecolor{currentstroke}{rgb}{0.000000,0.000000,0.000000}%
\pgfsetstrokecolor{currentstroke}%
\pgfsetdash{}{0pt}%
\pgfsys@defobject{currentmarker}{\pgfqpoint{0.000000in}{0.000000in}}{\pgfqpoint{0.000000in}{0.020833in}}{%
\pgfpathmoveto{\pgfqpoint{0.000000in}{0.000000in}}%
\pgfpathlineto{\pgfqpoint{0.000000in}{0.020833in}}%
\pgfusepath{stroke,fill}%
}%
\begin{pgfscope}%
\pgfsys@transformshift{2.638954in}{0.420000in}%
\pgfsys@useobject{currentmarker}{}%
\end{pgfscope}%
\end{pgfscope}%
\begin{pgfscope}%
\pgfsetbuttcap%
\pgfsetroundjoin%
\definecolor{currentfill}{rgb}{0.000000,0.000000,0.000000}%
\pgfsetfillcolor{currentfill}%
\pgfsetlinewidth{0.401500pt}%
\definecolor{currentstroke}{rgb}{0.000000,0.000000,0.000000}%
\pgfsetstrokecolor{currentstroke}%
\pgfsetdash{}{0pt}%
\pgfsys@defobject{currentmarker}{\pgfqpoint{0.000000in}{-0.020833in}}{\pgfqpoint{0.000000in}{0.000000in}}{%
\pgfpathmoveto{\pgfqpoint{0.000000in}{0.000000in}}%
\pgfpathlineto{\pgfqpoint{0.000000in}{-0.020833in}}%
\pgfusepath{stroke,fill}%
}%
\begin{pgfscope}%
\pgfsys@transformshift{2.638954in}{2.414488in}%
\pgfsys@useobject{currentmarker}{}%
\end{pgfscope}%
\end{pgfscope}%
\begin{pgfscope}%
\pgfsetbuttcap%
\pgfsetroundjoin%
\definecolor{currentfill}{rgb}{0.000000,0.000000,0.000000}%
\pgfsetfillcolor{currentfill}%
\pgfsetlinewidth{0.401500pt}%
\definecolor{currentstroke}{rgb}{0.000000,0.000000,0.000000}%
\pgfsetstrokecolor{currentstroke}%
\pgfsetdash{}{0pt}%
\pgfsys@defobject{currentmarker}{\pgfqpoint{0.000000in}{0.000000in}}{\pgfqpoint{0.000000in}{0.020833in}}{%
\pgfpathmoveto{\pgfqpoint{0.000000in}{0.000000in}}%
\pgfpathlineto{\pgfqpoint{0.000000in}{0.020833in}}%
\pgfusepath{stroke,fill}%
}%
\begin{pgfscope}%
\pgfsys@transformshift{2.691450in}{0.420000in}%
\pgfsys@useobject{currentmarker}{}%
\end{pgfscope}%
\end{pgfscope}%
\begin{pgfscope}%
\pgfsetbuttcap%
\pgfsetroundjoin%
\definecolor{currentfill}{rgb}{0.000000,0.000000,0.000000}%
\pgfsetfillcolor{currentfill}%
\pgfsetlinewidth{0.401500pt}%
\definecolor{currentstroke}{rgb}{0.000000,0.000000,0.000000}%
\pgfsetstrokecolor{currentstroke}%
\pgfsetdash{}{0pt}%
\pgfsys@defobject{currentmarker}{\pgfqpoint{0.000000in}{-0.020833in}}{\pgfqpoint{0.000000in}{0.000000in}}{%
\pgfpathmoveto{\pgfqpoint{0.000000in}{0.000000in}}%
\pgfpathlineto{\pgfqpoint{0.000000in}{-0.020833in}}%
\pgfusepath{stroke,fill}%
}%
\begin{pgfscope}%
\pgfsys@transformshift{2.691450in}{2.414488in}%
\pgfsys@useobject{currentmarker}{}%
\end{pgfscope}%
\end{pgfscope}%
\begin{pgfscope}%
\pgfsetbuttcap%
\pgfsetroundjoin%
\definecolor{currentfill}{rgb}{0.000000,0.000000,0.000000}%
\pgfsetfillcolor{currentfill}%
\pgfsetlinewidth{0.401500pt}%
\definecolor{currentstroke}{rgb}{0.000000,0.000000,0.000000}%
\pgfsetstrokecolor{currentstroke}%
\pgfsetdash{}{0pt}%
\pgfsys@defobject{currentmarker}{\pgfqpoint{0.000000in}{0.000000in}}{\pgfqpoint{0.000000in}{0.020833in}}{%
\pgfpathmoveto{\pgfqpoint{0.000000in}{0.000000in}}%
\pgfpathlineto{\pgfqpoint{0.000000in}{0.020833in}}%
\pgfusepath{stroke,fill}%
}%
\begin{pgfscope}%
\pgfsys@transformshift{2.735834in}{0.420000in}%
\pgfsys@useobject{currentmarker}{}%
\end{pgfscope}%
\end{pgfscope}%
\begin{pgfscope}%
\pgfsetbuttcap%
\pgfsetroundjoin%
\definecolor{currentfill}{rgb}{0.000000,0.000000,0.000000}%
\pgfsetfillcolor{currentfill}%
\pgfsetlinewidth{0.401500pt}%
\definecolor{currentstroke}{rgb}{0.000000,0.000000,0.000000}%
\pgfsetstrokecolor{currentstroke}%
\pgfsetdash{}{0pt}%
\pgfsys@defobject{currentmarker}{\pgfqpoint{0.000000in}{-0.020833in}}{\pgfqpoint{0.000000in}{0.000000in}}{%
\pgfpathmoveto{\pgfqpoint{0.000000in}{0.000000in}}%
\pgfpathlineto{\pgfqpoint{0.000000in}{-0.020833in}}%
\pgfusepath{stroke,fill}%
}%
\begin{pgfscope}%
\pgfsys@transformshift{2.735834in}{2.414488in}%
\pgfsys@useobject{currentmarker}{}%
\end{pgfscope}%
\end{pgfscope}%
\begin{pgfscope}%
\pgfsetbuttcap%
\pgfsetroundjoin%
\definecolor{currentfill}{rgb}{0.000000,0.000000,0.000000}%
\pgfsetfillcolor{currentfill}%
\pgfsetlinewidth{0.401500pt}%
\definecolor{currentstroke}{rgb}{0.000000,0.000000,0.000000}%
\pgfsetstrokecolor{currentstroke}%
\pgfsetdash{}{0pt}%
\pgfsys@defobject{currentmarker}{\pgfqpoint{0.000000in}{0.000000in}}{\pgfqpoint{0.000000in}{0.020833in}}{%
\pgfpathmoveto{\pgfqpoint{0.000000in}{0.000000in}}%
\pgfpathlineto{\pgfqpoint{0.000000in}{0.020833in}}%
\pgfusepath{stroke,fill}%
}%
\begin{pgfscope}%
\pgfsys@transformshift{2.774282in}{0.420000in}%
\pgfsys@useobject{currentmarker}{}%
\end{pgfscope}%
\end{pgfscope}%
\begin{pgfscope}%
\pgfsetbuttcap%
\pgfsetroundjoin%
\definecolor{currentfill}{rgb}{0.000000,0.000000,0.000000}%
\pgfsetfillcolor{currentfill}%
\pgfsetlinewidth{0.401500pt}%
\definecolor{currentstroke}{rgb}{0.000000,0.000000,0.000000}%
\pgfsetstrokecolor{currentstroke}%
\pgfsetdash{}{0pt}%
\pgfsys@defobject{currentmarker}{\pgfqpoint{0.000000in}{-0.020833in}}{\pgfqpoint{0.000000in}{0.000000in}}{%
\pgfpathmoveto{\pgfqpoint{0.000000in}{0.000000in}}%
\pgfpathlineto{\pgfqpoint{0.000000in}{-0.020833in}}%
\pgfusepath{stroke,fill}%
}%
\begin{pgfscope}%
\pgfsys@transformshift{2.774282in}{2.414488in}%
\pgfsys@useobject{currentmarker}{}%
\end{pgfscope}%
\end{pgfscope}%
\begin{pgfscope}%
\pgfsetbuttcap%
\pgfsetroundjoin%
\definecolor{currentfill}{rgb}{0.000000,0.000000,0.000000}%
\pgfsetfillcolor{currentfill}%
\pgfsetlinewidth{0.401500pt}%
\definecolor{currentstroke}{rgb}{0.000000,0.000000,0.000000}%
\pgfsetstrokecolor{currentstroke}%
\pgfsetdash{}{0pt}%
\pgfsys@defobject{currentmarker}{\pgfqpoint{0.000000in}{0.000000in}}{\pgfqpoint{0.000000in}{0.020833in}}{%
\pgfpathmoveto{\pgfqpoint{0.000000in}{0.000000in}}%
\pgfpathlineto{\pgfqpoint{0.000000in}{0.020833in}}%
\pgfusepath{stroke,fill}%
}%
\begin{pgfscope}%
\pgfsys@transformshift{2.808195in}{0.420000in}%
\pgfsys@useobject{currentmarker}{}%
\end{pgfscope}%
\end{pgfscope}%
\begin{pgfscope}%
\pgfsetbuttcap%
\pgfsetroundjoin%
\definecolor{currentfill}{rgb}{0.000000,0.000000,0.000000}%
\pgfsetfillcolor{currentfill}%
\pgfsetlinewidth{0.401500pt}%
\definecolor{currentstroke}{rgb}{0.000000,0.000000,0.000000}%
\pgfsetstrokecolor{currentstroke}%
\pgfsetdash{}{0pt}%
\pgfsys@defobject{currentmarker}{\pgfqpoint{0.000000in}{-0.020833in}}{\pgfqpoint{0.000000in}{0.000000in}}{%
\pgfpathmoveto{\pgfqpoint{0.000000in}{0.000000in}}%
\pgfpathlineto{\pgfqpoint{0.000000in}{-0.020833in}}%
\pgfusepath{stroke,fill}%
}%
\begin{pgfscope}%
\pgfsys@transformshift{2.808195in}{2.414488in}%
\pgfsys@useobject{currentmarker}{}%
\end{pgfscope}%
\end{pgfscope}%
\begin{pgfscope}%
\pgfsetbuttcap%
\pgfsetroundjoin%
\definecolor{currentfill}{rgb}{0.000000,0.000000,0.000000}%
\pgfsetfillcolor{currentfill}%
\pgfsetlinewidth{0.401500pt}%
\definecolor{currentstroke}{rgb}{0.000000,0.000000,0.000000}%
\pgfsetstrokecolor{currentstroke}%
\pgfsetdash{}{0pt}%
\pgfsys@defobject{currentmarker}{\pgfqpoint{0.000000in}{0.000000in}}{\pgfqpoint{0.000000in}{0.020833in}}{%
\pgfpathmoveto{\pgfqpoint{0.000000in}{0.000000in}}%
\pgfpathlineto{\pgfqpoint{0.000000in}{0.020833in}}%
\pgfusepath{stroke,fill}%
}%
\begin{pgfscope}%
\pgfsys@transformshift{3.038110in}{0.420000in}%
\pgfsys@useobject{currentmarker}{}%
\end{pgfscope}%
\end{pgfscope}%
\begin{pgfscope}%
\pgfsetbuttcap%
\pgfsetroundjoin%
\definecolor{currentfill}{rgb}{0.000000,0.000000,0.000000}%
\pgfsetfillcolor{currentfill}%
\pgfsetlinewidth{0.401500pt}%
\definecolor{currentstroke}{rgb}{0.000000,0.000000,0.000000}%
\pgfsetstrokecolor{currentstroke}%
\pgfsetdash{}{0pt}%
\pgfsys@defobject{currentmarker}{\pgfqpoint{0.000000in}{-0.020833in}}{\pgfqpoint{0.000000in}{0.000000in}}{%
\pgfpathmoveto{\pgfqpoint{0.000000in}{0.000000in}}%
\pgfpathlineto{\pgfqpoint{0.000000in}{-0.020833in}}%
\pgfusepath{stroke,fill}%
}%
\begin{pgfscope}%
\pgfsys@transformshift{3.038110in}{2.414488in}%
\pgfsys@useobject{currentmarker}{}%
\end{pgfscope}%
\end{pgfscope}%
\begin{pgfscope}%
\definecolor{textcolor}{rgb}{0.000000,0.000000,0.000000}%
\pgfsetstrokecolor{textcolor}%
\pgfsetfillcolor{textcolor}%
\pgftext[x=1.844055in,y=0.208426in,,top]{\color{textcolor}{\rmfamily\fontsize{12.000000}{14.400000}\selectfont\catcode`\^=\active\def^{\ifmmode\sp\else\^{}\fi}\catcode`\%=\active\def%{\%}$y^{+}$}}%
\end{pgfscope}%
\begin{pgfscope}%
\pgfsetbuttcap%
\pgfsetroundjoin%
\definecolor{currentfill}{rgb}{0.000000,0.000000,0.000000}%
\pgfsetfillcolor{currentfill}%
\pgfsetlinewidth{0.501875pt}%
\definecolor{currentstroke}{rgb}{0.000000,0.000000,0.000000}%
\pgfsetstrokecolor{currentstroke}%
\pgfsetdash{}{0pt}%
\pgfsys@defobject{currentmarker}{\pgfqpoint{0.000000in}{0.000000in}}{\pgfqpoint{0.041667in}{0.000000in}}{%
\pgfpathmoveto{\pgfqpoint{0.000000in}{0.000000in}}%
\pgfpathlineto{\pgfqpoint{0.041667in}{0.000000in}}%
\pgfusepath{stroke,fill}%
}%
\begin{pgfscope}%
\pgfsys@transformshift{0.650000in}{0.699299in}%
\pgfsys@useobject{currentmarker}{}%
\end{pgfscope}%
\end{pgfscope}%
\begin{pgfscope}%
\pgfsetbuttcap%
\pgfsetroundjoin%
\definecolor{currentfill}{rgb}{0.000000,0.000000,0.000000}%
\pgfsetfillcolor{currentfill}%
\pgfsetlinewidth{0.501875pt}%
\definecolor{currentstroke}{rgb}{0.000000,0.000000,0.000000}%
\pgfsetstrokecolor{currentstroke}%
\pgfsetdash{}{0pt}%
\pgfsys@defobject{currentmarker}{\pgfqpoint{-0.041667in}{0.000000in}}{\pgfqpoint{-0.000000in}{0.000000in}}{%
\pgfpathmoveto{\pgfqpoint{-0.000000in}{0.000000in}}%
\pgfpathlineto{\pgfqpoint{-0.041667in}{0.000000in}}%
\pgfusepath{stroke,fill}%
}%
\begin{pgfscope}%
\pgfsys@transformshift{3.038110in}{0.699299in}%
\pgfsys@useobject{currentmarker}{}%
\end{pgfscope}%
\end{pgfscope}%
\begin{pgfscope}%
\definecolor{textcolor}{rgb}{0.000000,0.000000,0.000000}%
\pgfsetstrokecolor{textcolor}%
\pgfsetfillcolor{textcolor}%
\pgftext[x=0.288772in, y=0.646492in, left, base]{\color{textcolor}{\rmfamily\fontsize{11.000000}{13.200000}\selectfont\catcode`\^=\active\def^{\ifmmode\sp\else\^{}\fi}\catcode`\%=\active\def%{\%}\ensuremath{-}0.2}}%
\end{pgfscope}%
\begin{pgfscope}%
\pgfsetbuttcap%
\pgfsetroundjoin%
\definecolor{currentfill}{rgb}{0.000000,0.000000,0.000000}%
\pgfsetfillcolor{currentfill}%
\pgfsetlinewidth{0.501875pt}%
\definecolor{currentstroke}{rgb}{0.000000,0.000000,0.000000}%
\pgfsetstrokecolor{currentstroke}%
\pgfsetdash{}{0pt}%
\pgfsys@defobject{currentmarker}{\pgfqpoint{0.000000in}{0.000000in}}{\pgfqpoint{0.041667in}{0.000000in}}{%
\pgfpathmoveto{\pgfqpoint{0.000000in}{0.000000in}}%
\pgfpathlineto{\pgfqpoint{0.041667in}{0.000000in}}%
\pgfusepath{stroke,fill}%
}%
\begin{pgfscope}%
\pgfsys@transformshift{0.650000in}{1.058137in}%
\pgfsys@useobject{currentmarker}{}%
\end{pgfscope}%
\end{pgfscope}%
\begin{pgfscope}%
\pgfsetbuttcap%
\pgfsetroundjoin%
\definecolor{currentfill}{rgb}{0.000000,0.000000,0.000000}%
\pgfsetfillcolor{currentfill}%
\pgfsetlinewidth{0.501875pt}%
\definecolor{currentstroke}{rgb}{0.000000,0.000000,0.000000}%
\pgfsetstrokecolor{currentstroke}%
\pgfsetdash{}{0pt}%
\pgfsys@defobject{currentmarker}{\pgfqpoint{-0.041667in}{0.000000in}}{\pgfqpoint{-0.000000in}{0.000000in}}{%
\pgfpathmoveto{\pgfqpoint{-0.000000in}{0.000000in}}%
\pgfpathlineto{\pgfqpoint{-0.041667in}{0.000000in}}%
\pgfusepath{stroke,fill}%
}%
\begin{pgfscope}%
\pgfsys@transformshift{3.038110in}{1.058137in}%
\pgfsys@useobject{currentmarker}{}%
\end{pgfscope}%
\end{pgfscope}%
\begin{pgfscope}%
\definecolor{textcolor}{rgb}{0.000000,0.000000,0.000000}%
\pgfsetstrokecolor{textcolor}%
\pgfsetfillcolor{textcolor}%
\pgftext[x=0.288772in, y=1.005330in, left, base]{\color{textcolor}{\rmfamily\fontsize{11.000000}{13.200000}\selectfont\catcode`\^=\active\def^{\ifmmode\sp\else\^{}\fi}\catcode`\%=\active\def%{\%}\ensuremath{-}0.1}}%
\end{pgfscope}%
\begin{pgfscope}%
\pgfsetbuttcap%
\pgfsetroundjoin%
\definecolor{currentfill}{rgb}{0.000000,0.000000,0.000000}%
\pgfsetfillcolor{currentfill}%
\pgfsetlinewidth{0.501875pt}%
\definecolor{currentstroke}{rgb}{0.000000,0.000000,0.000000}%
\pgfsetstrokecolor{currentstroke}%
\pgfsetdash{}{0pt}%
\pgfsys@defobject{currentmarker}{\pgfqpoint{0.000000in}{0.000000in}}{\pgfqpoint{0.041667in}{0.000000in}}{%
\pgfpathmoveto{\pgfqpoint{0.000000in}{0.000000in}}%
\pgfpathlineto{\pgfqpoint{0.041667in}{0.000000in}}%
\pgfusepath{stroke,fill}%
}%
\begin{pgfscope}%
\pgfsys@transformshift{0.650000in}{1.416975in}%
\pgfsys@useobject{currentmarker}{}%
\end{pgfscope}%
\end{pgfscope}%
\begin{pgfscope}%
\pgfsetbuttcap%
\pgfsetroundjoin%
\definecolor{currentfill}{rgb}{0.000000,0.000000,0.000000}%
\pgfsetfillcolor{currentfill}%
\pgfsetlinewidth{0.501875pt}%
\definecolor{currentstroke}{rgb}{0.000000,0.000000,0.000000}%
\pgfsetstrokecolor{currentstroke}%
\pgfsetdash{}{0pt}%
\pgfsys@defobject{currentmarker}{\pgfqpoint{-0.041667in}{0.000000in}}{\pgfqpoint{-0.000000in}{0.000000in}}{%
\pgfpathmoveto{\pgfqpoint{-0.000000in}{0.000000in}}%
\pgfpathlineto{\pgfqpoint{-0.041667in}{0.000000in}}%
\pgfusepath{stroke,fill}%
}%
\begin{pgfscope}%
\pgfsys@transformshift{3.038110in}{1.416975in}%
\pgfsys@useobject{currentmarker}{}%
\end{pgfscope}%
\end{pgfscope}%
\begin{pgfscope}%
\definecolor{textcolor}{rgb}{0.000000,0.000000,0.000000}%
\pgfsetstrokecolor{textcolor}%
\pgfsetfillcolor{textcolor}%
\pgftext[x=0.407060in, y=1.364169in, left, base]{\color{textcolor}{\rmfamily\fontsize{11.000000}{13.200000}\selectfont\catcode`\^=\active\def^{\ifmmode\sp\else\^{}\fi}\catcode`\%=\active\def%{\%}0.0}}%
\end{pgfscope}%
\begin{pgfscope}%
\pgfsetbuttcap%
\pgfsetroundjoin%
\definecolor{currentfill}{rgb}{0.000000,0.000000,0.000000}%
\pgfsetfillcolor{currentfill}%
\pgfsetlinewidth{0.501875pt}%
\definecolor{currentstroke}{rgb}{0.000000,0.000000,0.000000}%
\pgfsetstrokecolor{currentstroke}%
\pgfsetdash{}{0pt}%
\pgfsys@defobject{currentmarker}{\pgfqpoint{0.000000in}{0.000000in}}{\pgfqpoint{0.041667in}{0.000000in}}{%
\pgfpathmoveto{\pgfqpoint{0.000000in}{0.000000in}}%
\pgfpathlineto{\pgfqpoint{0.041667in}{0.000000in}}%
\pgfusepath{stroke,fill}%
}%
\begin{pgfscope}%
\pgfsys@transformshift{0.650000in}{1.775813in}%
\pgfsys@useobject{currentmarker}{}%
\end{pgfscope}%
\end{pgfscope}%
\begin{pgfscope}%
\pgfsetbuttcap%
\pgfsetroundjoin%
\definecolor{currentfill}{rgb}{0.000000,0.000000,0.000000}%
\pgfsetfillcolor{currentfill}%
\pgfsetlinewidth{0.501875pt}%
\definecolor{currentstroke}{rgb}{0.000000,0.000000,0.000000}%
\pgfsetstrokecolor{currentstroke}%
\pgfsetdash{}{0pt}%
\pgfsys@defobject{currentmarker}{\pgfqpoint{-0.041667in}{0.000000in}}{\pgfqpoint{-0.000000in}{0.000000in}}{%
\pgfpathmoveto{\pgfqpoint{-0.000000in}{0.000000in}}%
\pgfpathlineto{\pgfqpoint{-0.041667in}{0.000000in}}%
\pgfusepath{stroke,fill}%
}%
\begin{pgfscope}%
\pgfsys@transformshift{3.038110in}{1.775813in}%
\pgfsys@useobject{currentmarker}{}%
\end{pgfscope}%
\end{pgfscope}%
\begin{pgfscope}%
\definecolor{textcolor}{rgb}{0.000000,0.000000,0.000000}%
\pgfsetstrokecolor{textcolor}%
\pgfsetfillcolor{textcolor}%
\pgftext[x=0.407060in, y=1.723007in, left, base]{\color{textcolor}{\rmfamily\fontsize{11.000000}{13.200000}\selectfont\catcode`\^=\active\def^{\ifmmode\sp\else\^{}\fi}\catcode`\%=\active\def%{\%}0.1}}%
\end{pgfscope}%
\begin{pgfscope}%
\pgfsetbuttcap%
\pgfsetroundjoin%
\definecolor{currentfill}{rgb}{0.000000,0.000000,0.000000}%
\pgfsetfillcolor{currentfill}%
\pgfsetlinewidth{0.501875pt}%
\definecolor{currentstroke}{rgb}{0.000000,0.000000,0.000000}%
\pgfsetstrokecolor{currentstroke}%
\pgfsetdash{}{0pt}%
\pgfsys@defobject{currentmarker}{\pgfqpoint{0.000000in}{0.000000in}}{\pgfqpoint{0.041667in}{0.000000in}}{%
\pgfpathmoveto{\pgfqpoint{0.000000in}{0.000000in}}%
\pgfpathlineto{\pgfqpoint{0.041667in}{0.000000in}}%
\pgfusepath{stroke,fill}%
}%
\begin{pgfscope}%
\pgfsys@transformshift{0.650000in}{2.134651in}%
\pgfsys@useobject{currentmarker}{}%
\end{pgfscope}%
\end{pgfscope}%
\begin{pgfscope}%
\pgfsetbuttcap%
\pgfsetroundjoin%
\definecolor{currentfill}{rgb}{0.000000,0.000000,0.000000}%
\pgfsetfillcolor{currentfill}%
\pgfsetlinewidth{0.501875pt}%
\definecolor{currentstroke}{rgb}{0.000000,0.000000,0.000000}%
\pgfsetstrokecolor{currentstroke}%
\pgfsetdash{}{0pt}%
\pgfsys@defobject{currentmarker}{\pgfqpoint{-0.041667in}{0.000000in}}{\pgfqpoint{-0.000000in}{0.000000in}}{%
\pgfpathmoveto{\pgfqpoint{-0.000000in}{0.000000in}}%
\pgfpathlineto{\pgfqpoint{-0.041667in}{0.000000in}}%
\pgfusepath{stroke,fill}%
}%
\begin{pgfscope}%
\pgfsys@transformshift{3.038110in}{2.134651in}%
\pgfsys@useobject{currentmarker}{}%
\end{pgfscope}%
\end{pgfscope}%
\begin{pgfscope}%
\definecolor{textcolor}{rgb}{0.000000,0.000000,0.000000}%
\pgfsetstrokecolor{textcolor}%
\pgfsetfillcolor{textcolor}%
\pgftext[x=0.407060in, y=2.081845in, left, base]{\color{textcolor}{\rmfamily\fontsize{11.000000}{13.200000}\selectfont\catcode`\^=\active\def^{\ifmmode\sp\else\^{}\fi}\catcode`\%=\active\def%{\%}0.2}}%
\end{pgfscope}%
\begin{pgfscope}%
\pgfsetbuttcap%
\pgfsetroundjoin%
\definecolor{currentfill}{rgb}{0.000000,0.000000,0.000000}%
\pgfsetfillcolor{currentfill}%
\pgfsetlinewidth{0.401500pt}%
\definecolor{currentstroke}{rgb}{0.000000,0.000000,0.000000}%
\pgfsetstrokecolor{currentstroke}%
\pgfsetdash{}{0pt}%
\pgfsys@defobject{currentmarker}{\pgfqpoint{0.000000in}{0.000000in}}{\pgfqpoint{0.020833in}{0.000000in}}{%
\pgfpathmoveto{\pgfqpoint{0.000000in}{0.000000in}}%
\pgfpathlineto{\pgfqpoint{0.020833in}{0.000000in}}%
\pgfusepath{stroke,fill}%
}%
\begin{pgfscope}%
\pgfsys@transformshift{0.650000in}{0.483996in}%
\pgfsys@useobject{currentmarker}{}%
\end{pgfscope}%
\end{pgfscope}%
\begin{pgfscope}%
\pgfsetbuttcap%
\pgfsetroundjoin%
\definecolor{currentfill}{rgb}{0.000000,0.000000,0.000000}%
\pgfsetfillcolor{currentfill}%
\pgfsetlinewidth{0.401500pt}%
\definecolor{currentstroke}{rgb}{0.000000,0.000000,0.000000}%
\pgfsetstrokecolor{currentstroke}%
\pgfsetdash{}{0pt}%
\pgfsys@defobject{currentmarker}{\pgfqpoint{-0.020833in}{0.000000in}}{\pgfqpoint{-0.000000in}{0.000000in}}{%
\pgfpathmoveto{\pgfqpoint{-0.000000in}{0.000000in}}%
\pgfpathlineto{\pgfqpoint{-0.020833in}{0.000000in}}%
\pgfusepath{stroke,fill}%
}%
\begin{pgfscope}%
\pgfsys@transformshift{3.038110in}{0.483996in}%
\pgfsys@useobject{currentmarker}{}%
\end{pgfscope}%
\end{pgfscope}%
\begin{pgfscope}%
\pgfsetbuttcap%
\pgfsetroundjoin%
\definecolor{currentfill}{rgb}{0.000000,0.000000,0.000000}%
\pgfsetfillcolor{currentfill}%
\pgfsetlinewidth{0.401500pt}%
\definecolor{currentstroke}{rgb}{0.000000,0.000000,0.000000}%
\pgfsetstrokecolor{currentstroke}%
\pgfsetdash{}{0pt}%
\pgfsys@defobject{currentmarker}{\pgfqpoint{0.000000in}{0.000000in}}{\pgfqpoint{0.020833in}{0.000000in}}{%
\pgfpathmoveto{\pgfqpoint{0.000000in}{0.000000in}}%
\pgfpathlineto{\pgfqpoint{0.020833in}{0.000000in}}%
\pgfusepath{stroke,fill}%
}%
\begin{pgfscope}%
\pgfsys@transformshift{0.650000in}{0.555764in}%
\pgfsys@useobject{currentmarker}{}%
\end{pgfscope}%
\end{pgfscope}%
\begin{pgfscope}%
\pgfsetbuttcap%
\pgfsetroundjoin%
\definecolor{currentfill}{rgb}{0.000000,0.000000,0.000000}%
\pgfsetfillcolor{currentfill}%
\pgfsetlinewidth{0.401500pt}%
\definecolor{currentstroke}{rgb}{0.000000,0.000000,0.000000}%
\pgfsetstrokecolor{currentstroke}%
\pgfsetdash{}{0pt}%
\pgfsys@defobject{currentmarker}{\pgfqpoint{-0.020833in}{0.000000in}}{\pgfqpoint{-0.000000in}{0.000000in}}{%
\pgfpathmoveto{\pgfqpoint{-0.000000in}{0.000000in}}%
\pgfpathlineto{\pgfqpoint{-0.020833in}{0.000000in}}%
\pgfusepath{stroke,fill}%
}%
\begin{pgfscope}%
\pgfsys@transformshift{3.038110in}{0.555764in}%
\pgfsys@useobject{currentmarker}{}%
\end{pgfscope}%
\end{pgfscope}%
\begin{pgfscope}%
\pgfsetbuttcap%
\pgfsetroundjoin%
\definecolor{currentfill}{rgb}{0.000000,0.000000,0.000000}%
\pgfsetfillcolor{currentfill}%
\pgfsetlinewidth{0.401500pt}%
\definecolor{currentstroke}{rgb}{0.000000,0.000000,0.000000}%
\pgfsetstrokecolor{currentstroke}%
\pgfsetdash{}{0pt}%
\pgfsys@defobject{currentmarker}{\pgfqpoint{0.000000in}{0.000000in}}{\pgfqpoint{0.020833in}{0.000000in}}{%
\pgfpathmoveto{\pgfqpoint{0.000000in}{0.000000in}}%
\pgfpathlineto{\pgfqpoint{0.020833in}{0.000000in}}%
\pgfusepath{stroke,fill}%
}%
\begin{pgfscope}%
\pgfsys@transformshift{0.650000in}{0.627531in}%
\pgfsys@useobject{currentmarker}{}%
\end{pgfscope}%
\end{pgfscope}%
\begin{pgfscope}%
\pgfsetbuttcap%
\pgfsetroundjoin%
\definecolor{currentfill}{rgb}{0.000000,0.000000,0.000000}%
\pgfsetfillcolor{currentfill}%
\pgfsetlinewidth{0.401500pt}%
\definecolor{currentstroke}{rgb}{0.000000,0.000000,0.000000}%
\pgfsetstrokecolor{currentstroke}%
\pgfsetdash{}{0pt}%
\pgfsys@defobject{currentmarker}{\pgfqpoint{-0.020833in}{0.000000in}}{\pgfqpoint{-0.000000in}{0.000000in}}{%
\pgfpathmoveto{\pgfqpoint{-0.000000in}{0.000000in}}%
\pgfpathlineto{\pgfqpoint{-0.020833in}{0.000000in}}%
\pgfusepath{stroke,fill}%
}%
\begin{pgfscope}%
\pgfsys@transformshift{3.038110in}{0.627531in}%
\pgfsys@useobject{currentmarker}{}%
\end{pgfscope}%
\end{pgfscope}%
\begin{pgfscope}%
\pgfsetbuttcap%
\pgfsetroundjoin%
\definecolor{currentfill}{rgb}{0.000000,0.000000,0.000000}%
\pgfsetfillcolor{currentfill}%
\pgfsetlinewidth{0.401500pt}%
\definecolor{currentstroke}{rgb}{0.000000,0.000000,0.000000}%
\pgfsetstrokecolor{currentstroke}%
\pgfsetdash{}{0pt}%
\pgfsys@defobject{currentmarker}{\pgfqpoint{0.000000in}{0.000000in}}{\pgfqpoint{0.020833in}{0.000000in}}{%
\pgfpathmoveto{\pgfqpoint{0.000000in}{0.000000in}}%
\pgfpathlineto{\pgfqpoint{0.020833in}{0.000000in}}%
\pgfusepath{stroke,fill}%
}%
\begin{pgfscope}%
\pgfsys@transformshift{0.650000in}{0.771067in}%
\pgfsys@useobject{currentmarker}{}%
\end{pgfscope}%
\end{pgfscope}%
\begin{pgfscope}%
\pgfsetbuttcap%
\pgfsetroundjoin%
\definecolor{currentfill}{rgb}{0.000000,0.000000,0.000000}%
\pgfsetfillcolor{currentfill}%
\pgfsetlinewidth{0.401500pt}%
\definecolor{currentstroke}{rgb}{0.000000,0.000000,0.000000}%
\pgfsetstrokecolor{currentstroke}%
\pgfsetdash{}{0pt}%
\pgfsys@defobject{currentmarker}{\pgfqpoint{-0.020833in}{0.000000in}}{\pgfqpoint{-0.000000in}{0.000000in}}{%
\pgfpathmoveto{\pgfqpoint{-0.000000in}{0.000000in}}%
\pgfpathlineto{\pgfqpoint{-0.020833in}{0.000000in}}%
\pgfusepath{stroke,fill}%
}%
\begin{pgfscope}%
\pgfsys@transformshift{3.038110in}{0.771067in}%
\pgfsys@useobject{currentmarker}{}%
\end{pgfscope}%
\end{pgfscope}%
\begin{pgfscope}%
\pgfsetbuttcap%
\pgfsetroundjoin%
\definecolor{currentfill}{rgb}{0.000000,0.000000,0.000000}%
\pgfsetfillcolor{currentfill}%
\pgfsetlinewidth{0.401500pt}%
\definecolor{currentstroke}{rgb}{0.000000,0.000000,0.000000}%
\pgfsetstrokecolor{currentstroke}%
\pgfsetdash{}{0pt}%
\pgfsys@defobject{currentmarker}{\pgfqpoint{0.000000in}{0.000000in}}{\pgfqpoint{0.020833in}{0.000000in}}{%
\pgfpathmoveto{\pgfqpoint{0.000000in}{0.000000in}}%
\pgfpathlineto{\pgfqpoint{0.020833in}{0.000000in}}%
\pgfusepath{stroke,fill}%
}%
\begin{pgfscope}%
\pgfsys@transformshift{0.650000in}{0.842834in}%
\pgfsys@useobject{currentmarker}{}%
\end{pgfscope}%
\end{pgfscope}%
\begin{pgfscope}%
\pgfsetbuttcap%
\pgfsetroundjoin%
\definecolor{currentfill}{rgb}{0.000000,0.000000,0.000000}%
\pgfsetfillcolor{currentfill}%
\pgfsetlinewidth{0.401500pt}%
\definecolor{currentstroke}{rgb}{0.000000,0.000000,0.000000}%
\pgfsetstrokecolor{currentstroke}%
\pgfsetdash{}{0pt}%
\pgfsys@defobject{currentmarker}{\pgfqpoint{-0.020833in}{0.000000in}}{\pgfqpoint{-0.000000in}{0.000000in}}{%
\pgfpathmoveto{\pgfqpoint{-0.000000in}{0.000000in}}%
\pgfpathlineto{\pgfqpoint{-0.020833in}{0.000000in}}%
\pgfusepath{stroke,fill}%
}%
\begin{pgfscope}%
\pgfsys@transformshift{3.038110in}{0.842834in}%
\pgfsys@useobject{currentmarker}{}%
\end{pgfscope}%
\end{pgfscope}%
\begin{pgfscope}%
\pgfsetbuttcap%
\pgfsetroundjoin%
\definecolor{currentfill}{rgb}{0.000000,0.000000,0.000000}%
\pgfsetfillcolor{currentfill}%
\pgfsetlinewidth{0.401500pt}%
\definecolor{currentstroke}{rgb}{0.000000,0.000000,0.000000}%
\pgfsetstrokecolor{currentstroke}%
\pgfsetdash{}{0pt}%
\pgfsys@defobject{currentmarker}{\pgfqpoint{0.000000in}{0.000000in}}{\pgfqpoint{0.020833in}{0.000000in}}{%
\pgfpathmoveto{\pgfqpoint{0.000000in}{0.000000in}}%
\pgfpathlineto{\pgfqpoint{0.020833in}{0.000000in}}%
\pgfusepath{stroke,fill}%
}%
\begin{pgfscope}%
\pgfsys@transformshift{0.650000in}{0.914602in}%
\pgfsys@useobject{currentmarker}{}%
\end{pgfscope}%
\end{pgfscope}%
\begin{pgfscope}%
\pgfsetbuttcap%
\pgfsetroundjoin%
\definecolor{currentfill}{rgb}{0.000000,0.000000,0.000000}%
\pgfsetfillcolor{currentfill}%
\pgfsetlinewidth{0.401500pt}%
\definecolor{currentstroke}{rgb}{0.000000,0.000000,0.000000}%
\pgfsetstrokecolor{currentstroke}%
\pgfsetdash{}{0pt}%
\pgfsys@defobject{currentmarker}{\pgfqpoint{-0.020833in}{0.000000in}}{\pgfqpoint{-0.000000in}{0.000000in}}{%
\pgfpathmoveto{\pgfqpoint{-0.000000in}{0.000000in}}%
\pgfpathlineto{\pgfqpoint{-0.020833in}{0.000000in}}%
\pgfusepath{stroke,fill}%
}%
\begin{pgfscope}%
\pgfsys@transformshift{3.038110in}{0.914602in}%
\pgfsys@useobject{currentmarker}{}%
\end{pgfscope}%
\end{pgfscope}%
\begin{pgfscope}%
\pgfsetbuttcap%
\pgfsetroundjoin%
\definecolor{currentfill}{rgb}{0.000000,0.000000,0.000000}%
\pgfsetfillcolor{currentfill}%
\pgfsetlinewidth{0.401500pt}%
\definecolor{currentstroke}{rgb}{0.000000,0.000000,0.000000}%
\pgfsetstrokecolor{currentstroke}%
\pgfsetdash{}{0pt}%
\pgfsys@defobject{currentmarker}{\pgfqpoint{0.000000in}{0.000000in}}{\pgfqpoint{0.020833in}{0.000000in}}{%
\pgfpathmoveto{\pgfqpoint{0.000000in}{0.000000in}}%
\pgfpathlineto{\pgfqpoint{0.020833in}{0.000000in}}%
\pgfusepath{stroke,fill}%
}%
\begin{pgfscope}%
\pgfsys@transformshift{0.650000in}{0.986370in}%
\pgfsys@useobject{currentmarker}{}%
\end{pgfscope}%
\end{pgfscope}%
\begin{pgfscope}%
\pgfsetbuttcap%
\pgfsetroundjoin%
\definecolor{currentfill}{rgb}{0.000000,0.000000,0.000000}%
\pgfsetfillcolor{currentfill}%
\pgfsetlinewidth{0.401500pt}%
\definecolor{currentstroke}{rgb}{0.000000,0.000000,0.000000}%
\pgfsetstrokecolor{currentstroke}%
\pgfsetdash{}{0pt}%
\pgfsys@defobject{currentmarker}{\pgfqpoint{-0.020833in}{0.000000in}}{\pgfqpoint{-0.000000in}{0.000000in}}{%
\pgfpathmoveto{\pgfqpoint{-0.000000in}{0.000000in}}%
\pgfpathlineto{\pgfqpoint{-0.020833in}{0.000000in}}%
\pgfusepath{stroke,fill}%
}%
\begin{pgfscope}%
\pgfsys@transformshift{3.038110in}{0.986370in}%
\pgfsys@useobject{currentmarker}{}%
\end{pgfscope}%
\end{pgfscope}%
\begin{pgfscope}%
\pgfsetbuttcap%
\pgfsetroundjoin%
\definecolor{currentfill}{rgb}{0.000000,0.000000,0.000000}%
\pgfsetfillcolor{currentfill}%
\pgfsetlinewidth{0.401500pt}%
\definecolor{currentstroke}{rgb}{0.000000,0.000000,0.000000}%
\pgfsetstrokecolor{currentstroke}%
\pgfsetdash{}{0pt}%
\pgfsys@defobject{currentmarker}{\pgfqpoint{0.000000in}{0.000000in}}{\pgfqpoint{0.020833in}{0.000000in}}{%
\pgfpathmoveto{\pgfqpoint{0.000000in}{0.000000in}}%
\pgfpathlineto{\pgfqpoint{0.020833in}{0.000000in}}%
\pgfusepath{stroke,fill}%
}%
\begin{pgfscope}%
\pgfsys@transformshift{0.650000in}{1.129905in}%
\pgfsys@useobject{currentmarker}{}%
\end{pgfscope}%
\end{pgfscope}%
\begin{pgfscope}%
\pgfsetbuttcap%
\pgfsetroundjoin%
\definecolor{currentfill}{rgb}{0.000000,0.000000,0.000000}%
\pgfsetfillcolor{currentfill}%
\pgfsetlinewidth{0.401500pt}%
\definecolor{currentstroke}{rgb}{0.000000,0.000000,0.000000}%
\pgfsetstrokecolor{currentstroke}%
\pgfsetdash{}{0pt}%
\pgfsys@defobject{currentmarker}{\pgfqpoint{-0.020833in}{0.000000in}}{\pgfqpoint{-0.000000in}{0.000000in}}{%
\pgfpathmoveto{\pgfqpoint{-0.000000in}{0.000000in}}%
\pgfpathlineto{\pgfqpoint{-0.020833in}{0.000000in}}%
\pgfusepath{stroke,fill}%
}%
\begin{pgfscope}%
\pgfsys@transformshift{3.038110in}{1.129905in}%
\pgfsys@useobject{currentmarker}{}%
\end{pgfscope}%
\end{pgfscope}%
\begin{pgfscope}%
\pgfsetbuttcap%
\pgfsetroundjoin%
\definecolor{currentfill}{rgb}{0.000000,0.000000,0.000000}%
\pgfsetfillcolor{currentfill}%
\pgfsetlinewidth{0.401500pt}%
\definecolor{currentstroke}{rgb}{0.000000,0.000000,0.000000}%
\pgfsetstrokecolor{currentstroke}%
\pgfsetdash{}{0pt}%
\pgfsys@defobject{currentmarker}{\pgfqpoint{0.000000in}{0.000000in}}{\pgfqpoint{0.020833in}{0.000000in}}{%
\pgfpathmoveto{\pgfqpoint{0.000000in}{0.000000in}}%
\pgfpathlineto{\pgfqpoint{0.020833in}{0.000000in}}%
\pgfusepath{stroke,fill}%
}%
\begin{pgfscope}%
\pgfsys@transformshift{0.650000in}{1.201672in}%
\pgfsys@useobject{currentmarker}{}%
\end{pgfscope}%
\end{pgfscope}%
\begin{pgfscope}%
\pgfsetbuttcap%
\pgfsetroundjoin%
\definecolor{currentfill}{rgb}{0.000000,0.000000,0.000000}%
\pgfsetfillcolor{currentfill}%
\pgfsetlinewidth{0.401500pt}%
\definecolor{currentstroke}{rgb}{0.000000,0.000000,0.000000}%
\pgfsetstrokecolor{currentstroke}%
\pgfsetdash{}{0pt}%
\pgfsys@defobject{currentmarker}{\pgfqpoint{-0.020833in}{0.000000in}}{\pgfqpoint{-0.000000in}{0.000000in}}{%
\pgfpathmoveto{\pgfqpoint{-0.000000in}{0.000000in}}%
\pgfpathlineto{\pgfqpoint{-0.020833in}{0.000000in}}%
\pgfusepath{stroke,fill}%
}%
\begin{pgfscope}%
\pgfsys@transformshift{3.038110in}{1.201672in}%
\pgfsys@useobject{currentmarker}{}%
\end{pgfscope}%
\end{pgfscope}%
\begin{pgfscope}%
\pgfsetbuttcap%
\pgfsetroundjoin%
\definecolor{currentfill}{rgb}{0.000000,0.000000,0.000000}%
\pgfsetfillcolor{currentfill}%
\pgfsetlinewidth{0.401500pt}%
\definecolor{currentstroke}{rgb}{0.000000,0.000000,0.000000}%
\pgfsetstrokecolor{currentstroke}%
\pgfsetdash{}{0pt}%
\pgfsys@defobject{currentmarker}{\pgfqpoint{0.000000in}{0.000000in}}{\pgfqpoint{0.020833in}{0.000000in}}{%
\pgfpathmoveto{\pgfqpoint{0.000000in}{0.000000in}}%
\pgfpathlineto{\pgfqpoint{0.020833in}{0.000000in}}%
\pgfusepath{stroke,fill}%
}%
\begin{pgfscope}%
\pgfsys@transformshift{0.650000in}{1.273440in}%
\pgfsys@useobject{currentmarker}{}%
\end{pgfscope}%
\end{pgfscope}%
\begin{pgfscope}%
\pgfsetbuttcap%
\pgfsetroundjoin%
\definecolor{currentfill}{rgb}{0.000000,0.000000,0.000000}%
\pgfsetfillcolor{currentfill}%
\pgfsetlinewidth{0.401500pt}%
\definecolor{currentstroke}{rgb}{0.000000,0.000000,0.000000}%
\pgfsetstrokecolor{currentstroke}%
\pgfsetdash{}{0pt}%
\pgfsys@defobject{currentmarker}{\pgfqpoint{-0.020833in}{0.000000in}}{\pgfqpoint{-0.000000in}{0.000000in}}{%
\pgfpathmoveto{\pgfqpoint{-0.000000in}{0.000000in}}%
\pgfpathlineto{\pgfqpoint{-0.020833in}{0.000000in}}%
\pgfusepath{stroke,fill}%
}%
\begin{pgfscope}%
\pgfsys@transformshift{3.038110in}{1.273440in}%
\pgfsys@useobject{currentmarker}{}%
\end{pgfscope}%
\end{pgfscope}%
\begin{pgfscope}%
\pgfsetbuttcap%
\pgfsetroundjoin%
\definecolor{currentfill}{rgb}{0.000000,0.000000,0.000000}%
\pgfsetfillcolor{currentfill}%
\pgfsetlinewidth{0.401500pt}%
\definecolor{currentstroke}{rgb}{0.000000,0.000000,0.000000}%
\pgfsetstrokecolor{currentstroke}%
\pgfsetdash{}{0pt}%
\pgfsys@defobject{currentmarker}{\pgfqpoint{0.000000in}{0.000000in}}{\pgfqpoint{0.020833in}{0.000000in}}{%
\pgfpathmoveto{\pgfqpoint{0.000000in}{0.000000in}}%
\pgfpathlineto{\pgfqpoint{0.020833in}{0.000000in}}%
\pgfusepath{stroke,fill}%
}%
\begin{pgfscope}%
\pgfsys@transformshift{0.650000in}{1.345208in}%
\pgfsys@useobject{currentmarker}{}%
\end{pgfscope}%
\end{pgfscope}%
\begin{pgfscope}%
\pgfsetbuttcap%
\pgfsetroundjoin%
\definecolor{currentfill}{rgb}{0.000000,0.000000,0.000000}%
\pgfsetfillcolor{currentfill}%
\pgfsetlinewidth{0.401500pt}%
\definecolor{currentstroke}{rgb}{0.000000,0.000000,0.000000}%
\pgfsetstrokecolor{currentstroke}%
\pgfsetdash{}{0pt}%
\pgfsys@defobject{currentmarker}{\pgfqpoint{-0.020833in}{0.000000in}}{\pgfqpoint{-0.000000in}{0.000000in}}{%
\pgfpathmoveto{\pgfqpoint{-0.000000in}{0.000000in}}%
\pgfpathlineto{\pgfqpoint{-0.020833in}{0.000000in}}%
\pgfusepath{stroke,fill}%
}%
\begin{pgfscope}%
\pgfsys@transformshift{3.038110in}{1.345208in}%
\pgfsys@useobject{currentmarker}{}%
\end{pgfscope}%
\end{pgfscope}%
\begin{pgfscope}%
\pgfsetbuttcap%
\pgfsetroundjoin%
\definecolor{currentfill}{rgb}{0.000000,0.000000,0.000000}%
\pgfsetfillcolor{currentfill}%
\pgfsetlinewidth{0.401500pt}%
\definecolor{currentstroke}{rgb}{0.000000,0.000000,0.000000}%
\pgfsetstrokecolor{currentstroke}%
\pgfsetdash{}{0pt}%
\pgfsys@defobject{currentmarker}{\pgfqpoint{0.000000in}{0.000000in}}{\pgfqpoint{0.020833in}{0.000000in}}{%
\pgfpathmoveto{\pgfqpoint{0.000000in}{0.000000in}}%
\pgfpathlineto{\pgfqpoint{0.020833in}{0.000000in}}%
\pgfusepath{stroke,fill}%
}%
\begin{pgfscope}%
\pgfsys@transformshift{0.650000in}{1.488743in}%
\pgfsys@useobject{currentmarker}{}%
\end{pgfscope}%
\end{pgfscope}%
\begin{pgfscope}%
\pgfsetbuttcap%
\pgfsetroundjoin%
\definecolor{currentfill}{rgb}{0.000000,0.000000,0.000000}%
\pgfsetfillcolor{currentfill}%
\pgfsetlinewidth{0.401500pt}%
\definecolor{currentstroke}{rgb}{0.000000,0.000000,0.000000}%
\pgfsetstrokecolor{currentstroke}%
\pgfsetdash{}{0pt}%
\pgfsys@defobject{currentmarker}{\pgfqpoint{-0.020833in}{0.000000in}}{\pgfqpoint{-0.000000in}{0.000000in}}{%
\pgfpathmoveto{\pgfqpoint{-0.000000in}{0.000000in}}%
\pgfpathlineto{\pgfqpoint{-0.020833in}{0.000000in}}%
\pgfusepath{stroke,fill}%
}%
\begin{pgfscope}%
\pgfsys@transformshift{3.038110in}{1.488743in}%
\pgfsys@useobject{currentmarker}{}%
\end{pgfscope}%
\end{pgfscope}%
\begin{pgfscope}%
\pgfsetbuttcap%
\pgfsetroundjoin%
\definecolor{currentfill}{rgb}{0.000000,0.000000,0.000000}%
\pgfsetfillcolor{currentfill}%
\pgfsetlinewidth{0.401500pt}%
\definecolor{currentstroke}{rgb}{0.000000,0.000000,0.000000}%
\pgfsetstrokecolor{currentstroke}%
\pgfsetdash{}{0pt}%
\pgfsys@defobject{currentmarker}{\pgfqpoint{0.000000in}{0.000000in}}{\pgfqpoint{0.020833in}{0.000000in}}{%
\pgfpathmoveto{\pgfqpoint{0.000000in}{0.000000in}}%
\pgfpathlineto{\pgfqpoint{0.020833in}{0.000000in}}%
\pgfusepath{stroke,fill}%
}%
\begin{pgfscope}%
\pgfsys@transformshift{0.650000in}{1.560510in}%
\pgfsys@useobject{currentmarker}{}%
\end{pgfscope}%
\end{pgfscope}%
\begin{pgfscope}%
\pgfsetbuttcap%
\pgfsetroundjoin%
\definecolor{currentfill}{rgb}{0.000000,0.000000,0.000000}%
\pgfsetfillcolor{currentfill}%
\pgfsetlinewidth{0.401500pt}%
\definecolor{currentstroke}{rgb}{0.000000,0.000000,0.000000}%
\pgfsetstrokecolor{currentstroke}%
\pgfsetdash{}{0pt}%
\pgfsys@defobject{currentmarker}{\pgfqpoint{-0.020833in}{0.000000in}}{\pgfqpoint{-0.000000in}{0.000000in}}{%
\pgfpathmoveto{\pgfqpoint{-0.000000in}{0.000000in}}%
\pgfpathlineto{\pgfqpoint{-0.020833in}{0.000000in}}%
\pgfusepath{stroke,fill}%
}%
\begin{pgfscope}%
\pgfsys@transformshift{3.038110in}{1.560510in}%
\pgfsys@useobject{currentmarker}{}%
\end{pgfscope}%
\end{pgfscope}%
\begin{pgfscope}%
\pgfsetbuttcap%
\pgfsetroundjoin%
\definecolor{currentfill}{rgb}{0.000000,0.000000,0.000000}%
\pgfsetfillcolor{currentfill}%
\pgfsetlinewidth{0.401500pt}%
\definecolor{currentstroke}{rgb}{0.000000,0.000000,0.000000}%
\pgfsetstrokecolor{currentstroke}%
\pgfsetdash{}{0pt}%
\pgfsys@defobject{currentmarker}{\pgfqpoint{0.000000in}{0.000000in}}{\pgfqpoint{0.020833in}{0.000000in}}{%
\pgfpathmoveto{\pgfqpoint{0.000000in}{0.000000in}}%
\pgfpathlineto{\pgfqpoint{0.020833in}{0.000000in}}%
\pgfusepath{stroke,fill}%
}%
\begin{pgfscope}%
\pgfsys@transformshift{0.650000in}{1.632278in}%
\pgfsys@useobject{currentmarker}{}%
\end{pgfscope}%
\end{pgfscope}%
\begin{pgfscope}%
\pgfsetbuttcap%
\pgfsetroundjoin%
\definecolor{currentfill}{rgb}{0.000000,0.000000,0.000000}%
\pgfsetfillcolor{currentfill}%
\pgfsetlinewidth{0.401500pt}%
\definecolor{currentstroke}{rgb}{0.000000,0.000000,0.000000}%
\pgfsetstrokecolor{currentstroke}%
\pgfsetdash{}{0pt}%
\pgfsys@defobject{currentmarker}{\pgfqpoint{-0.020833in}{0.000000in}}{\pgfqpoint{-0.000000in}{0.000000in}}{%
\pgfpathmoveto{\pgfqpoint{-0.000000in}{0.000000in}}%
\pgfpathlineto{\pgfqpoint{-0.020833in}{0.000000in}}%
\pgfusepath{stroke,fill}%
}%
\begin{pgfscope}%
\pgfsys@transformshift{3.038110in}{1.632278in}%
\pgfsys@useobject{currentmarker}{}%
\end{pgfscope}%
\end{pgfscope}%
\begin{pgfscope}%
\pgfsetbuttcap%
\pgfsetroundjoin%
\definecolor{currentfill}{rgb}{0.000000,0.000000,0.000000}%
\pgfsetfillcolor{currentfill}%
\pgfsetlinewidth{0.401500pt}%
\definecolor{currentstroke}{rgb}{0.000000,0.000000,0.000000}%
\pgfsetstrokecolor{currentstroke}%
\pgfsetdash{}{0pt}%
\pgfsys@defobject{currentmarker}{\pgfqpoint{0.000000in}{0.000000in}}{\pgfqpoint{0.020833in}{0.000000in}}{%
\pgfpathmoveto{\pgfqpoint{0.000000in}{0.000000in}}%
\pgfpathlineto{\pgfqpoint{0.020833in}{0.000000in}}%
\pgfusepath{stroke,fill}%
}%
\begin{pgfscope}%
\pgfsys@transformshift{0.650000in}{1.704046in}%
\pgfsys@useobject{currentmarker}{}%
\end{pgfscope}%
\end{pgfscope}%
\begin{pgfscope}%
\pgfsetbuttcap%
\pgfsetroundjoin%
\definecolor{currentfill}{rgb}{0.000000,0.000000,0.000000}%
\pgfsetfillcolor{currentfill}%
\pgfsetlinewidth{0.401500pt}%
\definecolor{currentstroke}{rgb}{0.000000,0.000000,0.000000}%
\pgfsetstrokecolor{currentstroke}%
\pgfsetdash{}{0pt}%
\pgfsys@defobject{currentmarker}{\pgfqpoint{-0.020833in}{0.000000in}}{\pgfqpoint{-0.000000in}{0.000000in}}{%
\pgfpathmoveto{\pgfqpoint{-0.000000in}{0.000000in}}%
\pgfpathlineto{\pgfqpoint{-0.020833in}{0.000000in}}%
\pgfusepath{stroke,fill}%
}%
\begin{pgfscope}%
\pgfsys@transformshift{3.038110in}{1.704046in}%
\pgfsys@useobject{currentmarker}{}%
\end{pgfscope}%
\end{pgfscope}%
\begin{pgfscope}%
\pgfsetbuttcap%
\pgfsetroundjoin%
\definecolor{currentfill}{rgb}{0.000000,0.000000,0.000000}%
\pgfsetfillcolor{currentfill}%
\pgfsetlinewidth{0.401500pt}%
\definecolor{currentstroke}{rgb}{0.000000,0.000000,0.000000}%
\pgfsetstrokecolor{currentstroke}%
\pgfsetdash{}{0pt}%
\pgfsys@defobject{currentmarker}{\pgfqpoint{0.000000in}{0.000000in}}{\pgfqpoint{0.020833in}{0.000000in}}{%
\pgfpathmoveto{\pgfqpoint{0.000000in}{0.000000in}}%
\pgfpathlineto{\pgfqpoint{0.020833in}{0.000000in}}%
\pgfusepath{stroke,fill}%
}%
\begin{pgfscope}%
\pgfsys@transformshift{0.650000in}{1.847581in}%
\pgfsys@useobject{currentmarker}{}%
\end{pgfscope}%
\end{pgfscope}%
\begin{pgfscope}%
\pgfsetbuttcap%
\pgfsetroundjoin%
\definecolor{currentfill}{rgb}{0.000000,0.000000,0.000000}%
\pgfsetfillcolor{currentfill}%
\pgfsetlinewidth{0.401500pt}%
\definecolor{currentstroke}{rgb}{0.000000,0.000000,0.000000}%
\pgfsetstrokecolor{currentstroke}%
\pgfsetdash{}{0pt}%
\pgfsys@defobject{currentmarker}{\pgfqpoint{-0.020833in}{0.000000in}}{\pgfqpoint{-0.000000in}{0.000000in}}{%
\pgfpathmoveto{\pgfqpoint{-0.000000in}{0.000000in}}%
\pgfpathlineto{\pgfqpoint{-0.020833in}{0.000000in}}%
\pgfusepath{stroke,fill}%
}%
\begin{pgfscope}%
\pgfsys@transformshift{3.038110in}{1.847581in}%
\pgfsys@useobject{currentmarker}{}%
\end{pgfscope}%
\end{pgfscope}%
\begin{pgfscope}%
\pgfsetbuttcap%
\pgfsetroundjoin%
\definecolor{currentfill}{rgb}{0.000000,0.000000,0.000000}%
\pgfsetfillcolor{currentfill}%
\pgfsetlinewidth{0.401500pt}%
\definecolor{currentstroke}{rgb}{0.000000,0.000000,0.000000}%
\pgfsetstrokecolor{currentstroke}%
\pgfsetdash{}{0pt}%
\pgfsys@defobject{currentmarker}{\pgfqpoint{0.000000in}{0.000000in}}{\pgfqpoint{0.020833in}{0.000000in}}{%
\pgfpathmoveto{\pgfqpoint{0.000000in}{0.000000in}}%
\pgfpathlineto{\pgfqpoint{0.020833in}{0.000000in}}%
\pgfusepath{stroke,fill}%
}%
\begin{pgfscope}%
\pgfsys@transformshift{0.650000in}{1.919349in}%
\pgfsys@useobject{currentmarker}{}%
\end{pgfscope}%
\end{pgfscope}%
\begin{pgfscope}%
\pgfsetbuttcap%
\pgfsetroundjoin%
\definecolor{currentfill}{rgb}{0.000000,0.000000,0.000000}%
\pgfsetfillcolor{currentfill}%
\pgfsetlinewidth{0.401500pt}%
\definecolor{currentstroke}{rgb}{0.000000,0.000000,0.000000}%
\pgfsetstrokecolor{currentstroke}%
\pgfsetdash{}{0pt}%
\pgfsys@defobject{currentmarker}{\pgfqpoint{-0.020833in}{0.000000in}}{\pgfqpoint{-0.000000in}{0.000000in}}{%
\pgfpathmoveto{\pgfqpoint{-0.000000in}{0.000000in}}%
\pgfpathlineto{\pgfqpoint{-0.020833in}{0.000000in}}%
\pgfusepath{stroke,fill}%
}%
\begin{pgfscope}%
\pgfsys@transformshift{3.038110in}{1.919349in}%
\pgfsys@useobject{currentmarker}{}%
\end{pgfscope}%
\end{pgfscope}%
\begin{pgfscope}%
\pgfsetbuttcap%
\pgfsetroundjoin%
\definecolor{currentfill}{rgb}{0.000000,0.000000,0.000000}%
\pgfsetfillcolor{currentfill}%
\pgfsetlinewidth{0.401500pt}%
\definecolor{currentstroke}{rgb}{0.000000,0.000000,0.000000}%
\pgfsetstrokecolor{currentstroke}%
\pgfsetdash{}{0pt}%
\pgfsys@defobject{currentmarker}{\pgfqpoint{0.000000in}{0.000000in}}{\pgfqpoint{0.020833in}{0.000000in}}{%
\pgfpathmoveto{\pgfqpoint{0.000000in}{0.000000in}}%
\pgfpathlineto{\pgfqpoint{0.020833in}{0.000000in}}%
\pgfusepath{stroke,fill}%
}%
\begin{pgfscope}%
\pgfsys@transformshift{0.650000in}{1.991116in}%
\pgfsys@useobject{currentmarker}{}%
\end{pgfscope}%
\end{pgfscope}%
\begin{pgfscope}%
\pgfsetbuttcap%
\pgfsetroundjoin%
\definecolor{currentfill}{rgb}{0.000000,0.000000,0.000000}%
\pgfsetfillcolor{currentfill}%
\pgfsetlinewidth{0.401500pt}%
\definecolor{currentstroke}{rgb}{0.000000,0.000000,0.000000}%
\pgfsetstrokecolor{currentstroke}%
\pgfsetdash{}{0pt}%
\pgfsys@defobject{currentmarker}{\pgfqpoint{-0.020833in}{0.000000in}}{\pgfqpoint{-0.000000in}{0.000000in}}{%
\pgfpathmoveto{\pgfqpoint{-0.000000in}{0.000000in}}%
\pgfpathlineto{\pgfqpoint{-0.020833in}{0.000000in}}%
\pgfusepath{stroke,fill}%
}%
\begin{pgfscope}%
\pgfsys@transformshift{3.038110in}{1.991116in}%
\pgfsys@useobject{currentmarker}{}%
\end{pgfscope}%
\end{pgfscope}%
\begin{pgfscope}%
\pgfsetbuttcap%
\pgfsetroundjoin%
\definecolor{currentfill}{rgb}{0.000000,0.000000,0.000000}%
\pgfsetfillcolor{currentfill}%
\pgfsetlinewidth{0.401500pt}%
\definecolor{currentstroke}{rgb}{0.000000,0.000000,0.000000}%
\pgfsetstrokecolor{currentstroke}%
\pgfsetdash{}{0pt}%
\pgfsys@defobject{currentmarker}{\pgfqpoint{0.000000in}{0.000000in}}{\pgfqpoint{0.020833in}{0.000000in}}{%
\pgfpathmoveto{\pgfqpoint{0.000000in}{0.000000in}}%
\pgfpathlineto{\pgfqpoint{0.020833in}{0.000000in}}%
\pgfusepath{stroke,fill}%
}%
\begin{pgfscope}%
\pgfsys@transformshift{0.650000in}{2.062884in}%
\pgfsys@useobject{currentmarker}{}%
\end{pgfscope}%
\end{pgfscope}%
\begin{pgfscope}%
\pgfsetbuttcap%
\pgfsetroundjoin%
\definecolor{currentfill}{rgb}{0.000000,0.000000,0.000000}%
\pgfsetfillcolor{currentfill}%
\pgfsetlinewidth{0.401500pt}%
\definecolor{currentstroke}{rgb}{0.000000,0.000000,0.000000}%
\pgfsetstrokecolor{currentstroke}%
\pgfsetdash{}{0pt}%
\pgfsys@defobject{currentmarker}{\pgfqpoint{-0.020833in}{0.000000in}}{\pgfqpoint{-0.000000in}{0.000000in}}{%
\pgfpathmoveto{\pgfqpoint{-0.000000in}{0.000000in}}%
\pgfpathlineto{\pgfqpoint{-0.020833in}{0.000000in}}%
\pgfusepath{stroke,fill}%
}%
\begin{pgfscope}%
\pgfsys@transformshift{3.038110in}{2.062884in}%
\pgfsys@useobject{currentmarker}{}%
\end{pgfscope}%
\end{pgfscope}%
\begin{pgfscope}%
\pgfsetbuttcap%
\pgfsetroundjoin%
\definecolor{currentfill}{rgb}{0.000000,0.000000,0.000000}%
\pgfsetfillcolor{currentfill}%
\pgfsetlinewidth{0.401500pt}%
\definecolor{currentstroke}{rgb}{0.000000,0.000000,0.000000}%
\pgfsetstrokecolor{currentstroke}%
\pgfsetdash{}{0pt}%
\pgfsys@defobject{currentmarker}{\pgfqpoint{0.000000in}{0.000000in}}{\pgfqpoint{0.020833in}{0.000000in}}{%
\pgfpathmoveto{\pgfqpoint{0.000000in}{0.000000in}}%
\pgfpathlineto{\pgfqpoint{0.020833in}{0.000000in}}%
\pgfusepath{stroke,fill}%
}%
\begin{pgfscope}%
\pgfsys@transformshift{0.650000in}{2.206419in}%
\pgfsys@useobject{currentmarker}{}%
\end{pgfscope}%
\end{pgfscope}%
\begin{pgfscope}%
\pgfsetbuttcap%
\pgfsetroundjoin%
\definecolor{currentfill}{rgb}{0.000000,0.000000,0.000000}%
\pgfsetfillcolor{currentfill}%
\pgfsetlinewidth{0.401500pt}%
\definecolor{currentstroke}{rgb}{0.000000,0.000000,0.000000}%
\pgfsetstrokecolor{currentstroke}%
\pgfsetdash{}{0pt}%
\pgfsys@defobject{currentmarker}{\pgfqpoint{-0.020833in}{0.000000in}}{\pgfqpoint{-0.000000in}{0.000000in}}{%
\pgfpathmoveto{\pgfqpoint{-0.000000in}{0.000000in}}%
\pgfpathlineto{\pgfqpoint{-0.020833in}{0.000000in}}%
\pgfusepath{stroke,fill}%
}%
\begin{pgfscope}%
\pgfsys@transformshift{3.038110in}{2.206419in}%
\pgfsys@useobject{currentmarker}{}%
\end{pgfscope}%
\end{pgfscope}%
\begin{pgfscope}%
\pgfsetbuttcap%
\pgfsetroundjoin%
\definecolor{currentfill}{rgb}{0.000000,0.000000,0.000000}%
\pgfsetfillcolor{currentfill}%
\pgfsetlinewidth{0.401500pt}%
\definecolor{currentstroke}{rgb}{0.000000,0.000000,0.000000}%
\pgfsetstrokecolor{currentstroke}%
\pgfsetdash{}{0pt}%
\pgfsys@defobject{currentmarker}{\pgfqpoint{0.000000in}{0.000000in}}{\pgfqpoint{0.020833in}{0.000000in}}{%
\pgfpathmoveto{\pgfqpoint{0.000000in}{0.000000in}}%
\pgfpathlineto{\pgfqpoint{0.020833in}{0.000000in}}%
\pgfusepath{stroke,fill}%
}%
\begin{pgfscope}%
\pgfsys@transformshift{0.650000in}{2.278187in}%
\pgfsys@useobject{currentmarker}{}%
\end{pgfscope}%
\end{pgfscope}%
\begin{pgfscope}%
\pgfsetbuttcap%
\pgfsetroundjoin%
\definecolor{currentfill}{rgb}{0.000000,0.000000,0.000000}%
\pgfsetfillcolor{currentfill}%
\pgfsetlinewidth{0.401500pt}%
\definecolor{currentstroke}{rgb}{0.000000,0.000000,0.000000}%
\pgfsetstrokecolor{currentstroke}%
\pgfsetdash{}{0pt}%
\pgfsys@defobject{currentmarker}{\pgfqpoint{-0.020833in}{0.000000in}}{\pgfqpoint{-0.000000in}{0.000000in}}{%
\pgfpathmoveto{\pgfqpoint{-0.000000in}{0.000000in}}%
\pgfpathlineto{\pgfqpoint{-0.020833in}{0.000000in}}%
\pgfusepath{stroke,fill}%
}%
\begin{pgfscope}%
\pgfsys@transformshift{3.038110in}{2.278187in}%
\pgfsys@useobject{currentmarker}{}%
\end{pgfscope}%
\end{pgfscope}%
\begin{pgfscope}%
\pgfsetbuttcap%
\pgfsetroundjoin%
\definecolor{currentfill}{rgb}{0.000000,0.000000,0.000000}%
\pgfsetfillcolor{currentfill}%
\pgfsetlinewidth{0.401500pt}%
\definecolor{currentstroke}{rgb}{0.000000,0.000000,0.000000}%
\pgfsetstrokecolor{currentstroke}%
\pgfsetdash{}{0pt}%
\pgfsys@defobject{currentmarker}{\pgfqpoint{0.000000in}{0.000000in}}{\pgfqpoint{0.020833in}{0.000000in}}{%
\pgfpathmoveto{\pgfqpoint{0.000000in}{0.000000in}}%
\pgfpathlineto{\pgfqpoint{0.020833in}{0.000000in}}%
\pgfusepath{stroke,fill}%
}%
\begin{pgfscope}%
\pgfsys@transformshift{0.650000in}{2.349954in}%
\pgfsys@useobject{currentmarker}{}%
\end{pgfscope}%
\end{pgfscope}%
\begin{pgfscope}%
\pgfsetbuttcap%
\pgfsetroundjoin%
\definecolor{currentfill}{rgb}{0.000000,0.000000,0.000000}%
\pgfsetfillcolor{currentfill}%
\pgfsetlinewidth{0.401500pt}%
\definecolor{currentstroke}{rgb}{0.000000,0.000000,0.000000}%
\pgfsetstrokecolor{currentstroke}%
\pgfsetdash{}{0pt}%
\pgfsys@defobject{currentmarker}{\pgfqpoint{-0.020833in}{0.000000in}}{\pgfqpoint{-0.000000in}{0.000000in}}{%
\pgfpathmoveto{\pgfqpoint{-0.000000in}{0.000000in}}%
\pgfpathlineto{\pgfqpoint{-0.020833in}{0.000000in}}%
\pgfusepath{stroke,fill}%
}%
\begin{pgfscope}%
\pgfsys@transformshift{3.038110in}{2.349954in}%
\pgfsys@useobject{currentmarker}{}%
\end{pgfscope}%
\end{pgfscope}%
\begin{pgfscope}%
\definecolor{textcolor}{rgb}{0.000000,0.000000,0.000000}%
\pgfsetstrokecolor{textcolor}%
\pgfsetfillcolor{textcolor}%
\pgftext[x=0.260995in,y=1.417244in,,bottom,rotate=90.000000]{\color{textcolor}{\rmfamily\fontsize{12.000000}{14.400000}\selectfont\catcode`\^=\active\def^{\ifmmode\sp\else\^{}\fi}\catcode`\%=\active\def%{\%}TKE budget}}%
\end{pgfscope}%
\begin{pgfscope}%
\pgfpathrectangle{\pgfqpoint{0.650000in}{0.420000in}}{\pgfqpoint{2.388110in}{1.994488in}}%
\pgfusepath{clip}%
\pgfsetrectcap%
\pgfsetroundjoin%
\pgfsetlinewidth{0.602250pt}%
\definecolor{currentstroke}{rgb}{0.000000,0.000000,0.000000}%
\pgfsetstrokecolor{currentstroke}%
\pgfsetdash{}{0pt}%
\pgfpathmoveto{\pgfqpoint{0.645000in}{1.417563in}}%
\pgfpathlineto{\pgfqpoint{0.687793in}{1.417723in}}%
\pgfpathlineto{\pgfqpoint{0.816285in}{1.419889in}}%
\pgfpathlineto{\pgfqpoint{0.921267in}{1.425832in}}%
\pgfpathlineto{\pgfqpoint{1.010024in}{1.439484in}}%
\pgfpathlineto{\pgfqpoint{1.086906in}{1.466650in}}%
\pgfpathlineto{\pgfqpoint{1.154718in}{1.514254in}}%
\pgfpathlineto{\pgfqpoint{1.215374in}{1.588226in}}%
\pgfpathlineto{\pgfqpoint{1.270240in}{1.690385in}}%
\pgfpathlineto{\pgfqpoint{1.320326in}{1.815767in}}%
\pgfpathlineto{\pgfqpoint{1.366397in}{1.952243in}}%
\pgfpathlineto{\pgfqpoint{1.409048in}{2.083237in}}%
\pgfpathlineto{\pgfqpoint{1.448752in}{2.192484in}}%
\pgfpathlineto{\pgfqpoint{1.485889in}{2.268499in}}%
\pgfpathlineto{\pgfqpoint{1.520771in}{2.306838in}}%
\pgfpathlineto{\pgfqpoint{1.553655in}{2.309673in}}%
\pgfpathlineto{\pgfqpoint{1.584756in}{2.283606in}}%
\pgfpathlineto{\pgfqpoint{1.614259in}{2.237044in}}%
\pgfpathlineto{\pgfqpoint{1.642318in}{2.178127in}}%
\pgfpathlineto{\pgfqpoint{1.669069in}{2.113570in}}%
\pgfpathlineto{\pgfqpoint{1.742555in}{1.927078in}}%
\pgfpathlineto{\pgfqpoint{1.765095in}{1.873964in}}%
\pgfpathlineto{\pgfqpoint{1.786780in}{1.826423in}}%
\pgfpathlineto{\pgfqpoint{1.807673in}{1.784339in}}%
\pgfpathlineto{\pgfqpoint{1.827829in}{1.747382in}}%
\pgfpathlineto{\pgfqpoint{1.847299in}{1.715111in}}%
\pgfpathlineto{\pgfqpoint{1.866126in}{1.687051in}}%
\pgfpathlineto{\pgfqpoint{1.884352in}{1.662727in}}%
\pgfpathlineto{\pgfqpoint{1.902013in}{1.641677in}}%
\pgfpathlineto{\pgfqpoint{1.919144in}{1.623480in}}%
\pgfpathlineto{\pgfqpoint{1.935774in}{1.607753in}}%
\pgfpathlineto{\pgfqpoint{1.951933in}{1.594142in}}%
\pgfpathlineto{\pgfqpoint{1.967645in}{1.582327in}}%
\pgfpathlineto{\pgfqpoint{1.982934in}{1.572036in}}%
\pgfpathlineto{\pgfqpoint{1.997823in}{1.563036in}}%
\pgfpathlineto{\pgfqpoint{2.012331in}{1.555112in}}%
\pgfpathlineto{\pgfqpoint{2.040279in}{1.541802in}}%
\pgfpathlineto{\pgfqpoint{2.066912in}{1.531113in}}%
\pgfpathlineto{\pgfqpoint{2.092348in}{1.522385in}}%
\pgfpathlineto{\pgfqpoint{2.128472in}{1.511796in}}%
\pgfpathlineto{\pgfqpoint{2.162417in}{1.503276in}}%
\pgfpathlineto{\pgfqpoint{2.194427in}{1.496425in}}%
\pgfpathlineto{\pgfqpoint{2.234442in}{1.489016in}}%
\pgfpathlineto{\pgfqpoint{2.289540in}{1.480475in}}%
\pgfpathlineto{\pgfqpoint{2.347570in}{1.472703in}}%
\pgfpathlineto{\pgfqpoint{2.521097in}{1.448940in}}%
\pgfpathlineto{\pgfqpoint{2.595318in}{1.436814in}}%
\pgfpathlineto{\pgfqpoint{2.693657in}{1.420097in}}%
\pgfpathlineto{\pgfqpoint{2.713505in}{1.418237in}}%
\pgfpathlineto{\pgfqpoint{2.736282in}{1.417250in}}%
\pgfpathlineto{\pgfqpoint{2.778748in}{1.416978in}}%
\pgfpathlineto{\pgfqpoint{3.043110in}{1.416975in}}%
\pgfpathlineto{\pgfqpoint{3.043110in}{1.416975in}}%
\pgfusepath{stroke}%
\end{pgfscope}%
\begin{pgfscope}%
\pgfpathrectangle{\pgfqpoint{0.650000in}{0.420000in}}{\pgfqpoint{2.388110in}{1.994488in}}%
\pgfusepath{clip}%
\pgfsetbuttcap%
\pgfsetroundjoin%
\pgfsetlinewidth{0.602250pt}%
\definecolor{currentstroke}{rgb}{0.000000,0.000000,0.000000}%
\pgfsetstrokecolor{currentstroke}%
\pgfsetdash{{3.000000pt}{0.600000pt}}{0.000000pt}%
\pgfpathmoveto{\pgfqpoint{0.645000in}{1.416975in}}%
\pgfpathlineto{\pgfqpoint{1.409048in}{1.417419in}}%
\pgfpathlineto{\pgfqpoint{1.765095in}{1.418343in}}%
\pgfpathlineto{\pgfqpoint{2.104648in}{1.417453in}}%
\pgfpathlineto{\pgfqpoint{2.537918in}{1.415221in}}%
\pgfpathlineto{\pgfqpoint{2.689587in}{1.412287in}}%
\pgfpathlineto{\pgfqpoint{2.717378in}{1.413389in}}%
\pgfpathlineto{\pgfqpoint{2.765024in}{1.416066in}}%
\pgfpathlineto{\pgfqpoint{2.801804in}{1.416766in}}%
\pgfpathlineto{\pgfqpoint{2.896390in}{1.416963in}}%
\pgfpathlineto{\pgfqpoint{3.043110in}{1.416975in}}%
\pgfpathlineto{\pgfqpoint{3.043110in}{1.416975in}}%
\pgfusepath{stroke}%
\end{pgfscope}%
\begin{pgfscope}%
\pgfpathrectangle{\pgfqpoint{0.650000in}{0.420000in}}{\pgfqpoint{2.388110in}{1.994488in}}%
\pgfusepath{clip}%
\pgfsetbuttcap%
\pgfsetroundjoin%
\pgfsetlinewidth{0.602250pt}%
\definecolor{currentstroke}{rgb}{0.000000,0.000000,0.000000}%
\pgfsetstrokecolor{currentstroke}%
\pgfsetdash{{1.800000pt}{0.600000pt}{0.600000pt}{0.600000pt}}{0.000000pt}%
\pgfpathmoveto{\pgfqpoint{0.645000in}{0.581761in}}%
\pgfpathlineto{\pgfqpoint{0.687793in}{0.589371in}}%
\pgfpathlineto{\pgfqpoint{0.816285in}{0.619397in}}%
\pgfpathlineto{\pgfqpoint{0.921267in}{0.648353in}}%
\pgfpathlineto{\pgfqpoint{1.010024in}{0.677080in}}%
\pgfpathlineto{\pgfqpoint{1.086906in}{0.708023in}}%
\pgfpathlineto{\pgfqpoint{1.154718in}{0.743353in}}%
\pgfpathlineto{\pgfqpoint{1.215374in}{0.782559in}}%
\pgfpathlineto{\pgfqpoint{1.270240in}{0.821690in}}%
\pgfpathlineto{\pgfqpoint{1.320326in}{0.855140in}}%
\pgfpathlineto{\pgfqpoint{1.366397in}{0.878784in}}%
\pgfpathlineto{\pgfqpoint{1.409048in}{0.892072in}}%
\pgfpathlineto{\pgfqpoint{1.448752in}{0.897824in}}%
\pgfpathlineto{\pgfqpoint{1.485889in}{0.900382in}}%
\pgfpathlineto{\pgfqpoint{1.520771in}{0.903661in}}%
\pgfpathlineto{\pgfqpoint{1.553655in}{0.910079in}}%
\pgfpathlineto{\pgfqpoint{1.584756in}{0.920492in}}%
\pgfpathlineto{\pgfqpoint{1.614259in}{0.934662in}}%
\pgfpathlineto{\pgfqpoint{1.642318in}{0.951821in}}%
\pgfpathlineto{\pgfqpoint{1.669069in}{0.971075in}}%
\pgfpathlineto{\pgfqpoint{1.694626in}{0.991629in}}%
\pgfpathlineto{\pgfqpoint{1.719092in}{1.012836in}}%
\pgfpathlineto{\pgfqpoint{1.765095in}{1.055362in}}%
\pgfpathlineto{\pgfqpoint{1.847299in}{1.133370in}}%
\pgfpathlineto{\pgfqpoint{1.884352in}{1.166425in}}%
\pgfpathlineto{\pgfqpoint{1.919144in}{1.195031in}}%
\pgfpathlineto{\pgfqpoint{1.951933in}{1.219431in}}%
\pgfpathlineto{\pgfqpoint{1.982934in}{1.240146in}}%
\pgfpathlineto{\pgfqpoint{2.012331in}{1.257671in}}%
\pgfpathlineto{\pgfqpoint{2.040279in}{1.272579in}}%
\pgfpathlineto{\pgfqpoint{2.066912in}{1.285343in}}%
\pgfpathlineto{\pgfqpoint{2.092348in}{1.296360in}}%
\pgfpathlineto{\pgfqpoint{2.116686in}{1.305963in}}%
\pgfpathlineto{\pgfqpoint{2.151329in}{1.318146in}}%
\pgfpathlineto{\pgfqpoint{2.183959in}{1.328178in}}%
\pgfpathlineto{\pgfqpoint{2.214793in}{1.336617in}}%
\pgfpathlineto{\pgfqpoint{2.253422in}{1.345894in}}%
\pgfpathlineto{\pgfqpoint{2.289540in}{1.353448in}}%
\pgfpathlineto{\pgfqpoint{2.339645in}{1.362619in}}%
\pgfpathlineto{\pgfqpoint{2.400041in}{1.372139in}}%
\pgfpathlineto{\pgfqpoint{2.485750in}{1.384270in}}%
\pgfpathlineto{\pgfqpoint{2.590356in}{1.399156in}}%
\pgfpathlineto{\pgfqpoint{2.668712in}{1.410008in}}%
\pgfpathlineto{\pgfqpoint{2.701695in}{1.413515in}}%
\pgfpathlineto{\pgfqpoint{2.732561in}{1.415638in}}%
\pgfpathlineto{\pgfqpoint{2.765024in}{1.416651in}}%
\pgfpathlineto{\pgfqpoint{2.820653in}{1.416963in}}%
\pgfpathlineto{\pgfqpoint{3.043110in}{1.416974in}}%
\pgfpathlineto{\pgfqpoint{3.043110in}{1.416974in}}%
\pgfusepath{stroke}%
\end{pgfscope}%
\begin{pgfscope}%
\pgfpathrectangle{\pgfqpoint{0.650000in}{0.420000in}}{\pgfqpoint{2.388110in}{1.994488in}}%
\pgfusepath{clip}%
\pgfsetbuttcap%
\pgfsetroundjoin%
\pgfsetlinewidth{0.602250pt}%
\definecolor{currentstroke}{rgb}{0.000000,0.000000,0.000000}%
\pgfsetstrokecolor{currentstroke}%
\pgfsetdash{{0.600000pt}{0.600000pt}}{0.000000pt}%
\pgfpathmoveto{\pgfqpoint{0.645000in}{1.453645in}}%
\pgfpathlineto{\pgfqpoint{0.687793in}{1.457605in}}%
\pgfpathlineto{\pgfqpoint{0.921267in}{1.485531in}}%
\pgfpathlineto{\pgfqpoint{1.010024in}{1.494697in}}%
\pgfpathlineto{\pgfqpoint{1.086906in}{1.499652in}}%
\pgfpathlineto{\pgfqpoint{1.154718in}{1.500583in}}%
\pgfpathlineto{\pgfqpoint{1.215374in}{1.498170in}}%
\pgfpathlineto{\pgfqpoint{1.270240in}{1.493237in}}%
\pgfpathlineto{\pgfqpoint{1.320326in}{1.486512in}}%
\pgfpathlineto{\pgfqpoint{1.366397in}{1.478485in}}%
\pgfpathlineto{\pgfqpoint{1.409048in}{1.469488in}}%
\pgfpathlineto{\pgfqpoint{1.448752in}{1.459792in}}%
\pgfpathlineto{\pgfqpoint{1.520771in}{1.439935in}}%
\pgfpathlineto{\pgfqpoint{1.553655in}{1.430839in}}%
\pgfpathlineto{\pgfqpoint{1.584756in}{1.423001in}}%
\pgfpathlineto{\pgfqpoint{1.614259in}{1.416805in}}%
\pgfpathlineto{\pgfqpoint{1.642318in}{1.412399in}}%
\pgfpathlineto{\pgfqpoint{1.669069in}{1.409744in}}%
\pgfpathlineto{\pgfqpoint{1.694626in}{1.408608in}}%
\pgfpathlineto{\pgfqpoint{1.719092in}{1.408688in}}%
\pgfpathlineto{\pgfqpoint{1.742555in}{1.409624in}}%
\pgfpathlineto{\pgfqpoint{1.786780in}{1.412782in}}%
\pgfpathlineto{\pgfqpoint{1.847299in}{1.417288in}}%
\pgfpathlineto{\pgfqpoint{1.884352in}{1.419054in}}%
\pgfpathlineto{\pgfqpoint{1.935774in}{1.420165in}}%
\pgfpathlineto{\pgfqpoint{1.982934in}{1.419944in}}%
\pgfpathlineto{\pgfqpoint{2.079773in}{1.417268in}}%
\pgfpathlineto{\pgfqpoint{2.162417in}{1.415195in}}%
\pgfpathlineto{\pgfqpoint{2.244013in}{1.415084in}}%
\pgfpathlineto{\pgfqpoint{2.497797in}{1.417187in}}%
\pgfpathlineto{\pgfqpoint{2.580286in}{1.416109in}}%
\pgfpathlineto{\pgfqpoint{2.672959in}{1.415189in}}%
\pgfpathlineto{\pgfqpoint{2.713505in}{1.415611in}}%
\pgfpathlineto{\pgfqpoint{2.778748in}{1.416981in}}%
\pgfpathlineto{\pgfqpoint{2.878062in}{1.416994in}}%
\pgfpathlineto{\pgfqpoint{3.043110in}{1.416976in}}%
\pgfpathlineto{\pgfqpoint{3.043110in}{1.416976in}}%
\pgfusepath{stroke}%
\end{pgfscope}%
\begin{pgfscope}%
\pgfpathrectangle{\pgfqpoint{0.650000in}{0.420000in}}{\pgfqpoint{2.388110in}{1.994488in}}%
\pgfusepath{clip}%
\pgfsetbuttcap%
\pgfsetroundjoin%
\pgfsetlinewidth{0.602250pt}%
\definecolor{currentstroke}{rgb}{0.000000,0.000000,0.000000}%
\pgfsetstrokecolor{currentstroke}%
\pgfsetdash{{3.000000pt}{3.000000pt}}{0.000000pt}%
\pgfpathmoveto{\pgfqpoint{0.645000in}{1.417973in}}%
\pgfpathlineto{\pgfqpoint{0.687793in}{1.418243in}}%
\pgfpathlineto{\pgfqpoint{0.816285in}{1.421817in}}%
\pgfpathlineto{\pgfqpoint{0.921267in}{1.431178in}}%
\pgfpathlineto{\pgfqpoint{1.010024in}{1.450867in}}%
\pgfpathlineto{\pgfqpoint{1.086906in}{1.484383in}}%
\pgfpathlineto{\pgfqpoint{1.154718in}{1.529800in}}%
\pgfpathlineto{\pgfqpoint{1.215374in}{1.576310in}}%
\pgfpathlineto{\pgfqpoint{1.270240in}{1.606114in}}%
\pgfpathlineto{\pgfqpoint{1.320326in}{1.602430in}}%
\pgfpathlineto{\pgfqpoint{1.366397in}{1.558622in}}%
\pgfpathlineto{\pgfqpoint{1.409048in}{1.481937in}}%
\pgfpathlineto{\pgfqpoint{1.448752in}{1.389763in}}%
\pgfpathlineto{\pgfqpoint{1.485889in}{1.301712in}}%
\pgfpathlineto{\pgfqpoint{1.520771in}{1.232639in}}%
\pgfpathlineto{\pgfqpoint{1.553655in}{1.189549in}}%
\pgfpathlineto{\pgfqpoint{1.584756in}{1.172271in}}%
\pgfpathlineto{\pgfqpoint{1.614259in}{1.176091in}}%
\pgfpathlineto{\pgfqpoint{1.642318in}{1.194549in}}%
\pgfpathlineto{\pgfqpoint{1.669069in}{1.221422in}}%
\pgfpathlineto{\pgfqpoint{1.694626in}{1.251746in}}%
\pgfpathlineto{\pgfqpoint{1.719092in}{1.282045in}}%
\pgfpathlineto{\pgfqpoint{1.742555in}{1.310183in}}%
\pgfpathlineto{\pgfqpoint{1.765095in}{1.335060in}}%
\pgfpathlineto{\pgfqpoint{1.786780in}{1.356341in}}%
\pgfpathlineto{\pgfqpoint{1.807673in}{1.374149in}}%
\pgfpathlineto{\pgfqpoint{1.827829in}{1.388815in}}%
\pgfpathlineto{\pgfqpoint{1.847299in}{1.400697in}}%
\pgfpathlineto{\pgfqpoint{1.866126in}{1.410175in}}%
\pgfpathlineto{\pgfqpoint{1.884352in}{1.417585in}}%
\pgfpathlineto{\pgfqpoint{1.902013in}{1.423229in}}%
\pgfpathlineto{\pgfqpoint{1.919144in}{1.427362in}}%
\pgfpathlineto{\pgfqpoint{1.935774in}{1.430282in}}%
\pgfpathlineto{\pgfqpoint{1.951933in}{1.432280in}}%
\pgfpathlineto{\pgfqpoint{1.982934in}{1.434368in}}%
\pgfpathlineto{\pgfqpoint{2.012331in}{1.434716in}}%
\pgfpathlineto{\pgfqpoint{2.053752in}{1.433640in}}%
\pgfpathlineto{\pgfqpoint{2.116686in}{1.430166in}}%
\pgfpathlineto{\pgfqpoint{2.173292in}{1.426670in}}%
\pgfpathlineto{\pgfqpoint{2.271775in}{1.419276in}}%
\pgfpathlineto{\pgfqpoint{2.323444in}{1.416342in}}%
\pgfpathlineto{\pgfqpoint{2.363083in}{1.415126in}}%
\pgfpathlineto{\pgfqpoint{2.434641in}{1.414064in}}%
\pgfpathlineto{\pgfqpoint{2.521097in}{1.414506in}}%
\pgfpathlineto{\pgfqpoint{2.575176in}{1.417425in}}%
\pgfpathlineto{\pgfqpoint{2.664429in}{1.423491in}}%
\pgfpathlineto{\pgfqpoint{2.685483in}{1.424262in}}%
\pgfpathlineto{\pgfqpoint{2.705664in}{1.423841in}}%
\pgfpathlineto{\pgfqpoint{2.721220in}{1.422469in}}%
\pgfpathlineto{\pgfqpoint{2.761527in}{1.418453in}}%
\pgfpathlineto{\pgfqpoint{2.785457in}{1.417414in}}%
\pgfpathlineto{\pgfqpoint{2.826759in}{1.417003in}}%
\pgfpathlineto{\pgfqpoint{3.043110in}{1.416975in}}%
\pgfpathlineto{\pgfqpoint{3.043110in}{1.416975in}}%
\pgfusepath{stroke}%
\end{pgfscope}%
\begin{pgfscope}%
\pgfpathrectangle{\pgfqpoint{0.650000in}{0.420000in}}{\pgfqpoint{2.388110in}{1.994488in}}%
\pgfusepath{clip}%
\pgfsetbuttcap%
\pgfsetroundjoin%
\pgfsetlinewidth{0.602250pt}%
\definecolor{currentstroke}{rgb}{0.000000,0.000000,0.000000}%
\pgfsetstrokecolor{currentstroke}%
\pgfsetdash{{1.800000pt}{3.000000pt}{0.600000pt}{3.000000pt}}{0.000000pt}%
\pgfpathmoveto{\pgfqpoint{0.645000in}{2.213930in}}%
\pgfpathlineto{\pgfqpoint{0.687793in}{2.201915in}}%
\pgfpathlineto{\pgfqpoint{0.816285in}{2.151054in}}%
\pgfpathlineto{\pgfqpoint{0.921267in}{2.093956in}}%
\pgfpathlineto{\pgfqpoint{1.010024in}{2.022719in}}%
\pgfpathlineto{\pgfqpoint{1.086906in}{1.926125in}}%
\pgfpathlineto{\pgfqpoint{1.154718in}{1.796824in}}%
\pgfpathlineto{\pgfqpoint{1.215374in}{1.639519in}}%
\pgfpathlineto{\pgfqpoint{1.270240in}{1.473308in}}%
\pgfpathlineto{\pgfqpoint{1.320326in}{1.324815in}}%
\pgfpathlineto{\pgfqpoint{1.366397in}{1.216430in}}%
\pgfpathlineto{\pgfqpoint{1.409048in}{1.157703in}}%
\pgfpathlineto{\pgfqpoint{1.448752in}{1.144419in}}%
\pgfpathlineto{\pgfqpoint{1.485889in}{1.163743in}}%
\pgfpathlineto{\pgfqpoint{1.520771in}{1.200867in}}%
\pgfpathlineto{\pgfqpoint{1.584756in}{1.284266in}}%
\pgfpathlineto{\pgfqpoint{1.614259in}{1.318932in}}%
\pgfpathlineto{\pgfqpoint{1.642318in}{1.346576in}}%
\pgfpathlineto{\pgfqpoint{1.669069in}{1.367641in}}%
\pgfpathlineto{\pgfqpoint{1.694626in}{1.383196in}}%
\pgfpathlineto{\pgfqpoint{1.719092in}{1.394436in}}%
\pgfpathlineto{\pgfqpoint{1.742555in}{1.402448in}}%
\pgfpathlineto{\pgfqpoint{1.765095in}{1.408092in}}%
\pgfpathlineto{\pgfqpoint{1.786780in}{1.411996in}}%
\pgfpathlineto{\pgfqpoint{1.807673in}{1.414586in}}%
\pgfpathlineto{\pgfqpoint{1.827829in}{1.416210in}}%
\pgfpathlineto{\pgfqpoint{1.866126in}{1.417688in}}%
\pgfpathlineto{\pgfqpoint{1.919144in}{1.418023in}}%
\pgfpathlineto{\pgfqpoint{2.693657in}{1.417092in}}%
\pgfpathlineto{\pgfqpoint{2.942126in}{1.416975in}}%
\pgfpathlineto{\pgfqpoint{3.043110in}{1.416975in}}%
\pgfpathlineto{\pgfqpoint{3.043110in}{1.416975in}}%
\pgfusepath{stroke}%
\end{pgfscope}%
\begin{pgfscope}%
\pgfpathrectangle{\pgfqpoint{0.650000in}{0.420000in}}{\pgfqpoint{2.388110in}{1.994488in}}%
\pgfusepath{clip}%
\pgfsetbuttcap%
\pgfsetroundjoin%
\pgfsetlinewidth{0.602250pt}%
\definecolor{currentstroke}{rgb}{0.000000,0.000000,0.000000}%
\pgfsetstrokecolor{currentstroke}%
\pgfsetdash{{0.600000pt}{3.000000pt}}{0.000000pt}%
\pgfpathmoveto{\pgfqpoint{0.645000in}{1.416972in}}%
\pgfpathlineto{\pgfqpoint{1.951933in}{1.417050in}}%
\pgfpathlineto{\pgfqpoint{3.043110in}{1.416975in}}%
\pgfpathlineto{\pgfqpoint{3.043110in}{1.416975in}}%
\pgfusepath{stroke}%
\end{pgfscope}%
\begin{pgfscope}%
\pgfpathrectangle{\pgfqpoint{0.650000in}{0.420000in}}{\pgfqpoint{2.388110in}{1.994488in}}%
\pgfusepath{clip}%
\pgfsetrectcap%
\pgfsetroundjoin%
\pgfsetlinewidth{1.003750pt}%
\definecolor{currentstroke}{rgb}{0.454902,0.678431,0.819608}%
\pgfsetstrokecolor{currentstroke}%
\pgfsetdash{}{0pt}%
\pgfpathmoveto{\pgfqpoint{0.645000in}{1.417432in}}%
\pgfpathlineto{\pgfqpoint{0.744946in}{1.418191in}}%
\pgfpathlineto{\pgfqpoint{0.802042in}{1.419201in}}%
\pgfpathlineto{\pgfqpoint{0.853159in}{1.420813in}}%
\pgfpathlineto{\pgfqpoint{0.899759in}{1.423300in}}%
\pgfpathlineto{\pgfqpoint{0.942836in}{1.427024in}}%
\pgfpathlineto{\pgfqpoint{0.983096in}{1.432438in}}%
\pgfpathlineto{\pgfqpoint{1.021057in}{1.440082in}}%
\pgfpathlineto{\pgfqpoint{1.057113in}{1.450579in}}%
\pgfpathlineto{\pgfqpoint{1.091568in}{1.464638in}}%
\pgfpathlineto{\pgfqpoint{1.124662in}{1.483133in}}%
\pgfpathlineto{\pgfqpoint{1.156588in}{1.507235in}}%
\pgfpathlineto{\pgfqpoint{1.187505in}{1.538451in}}%
\pgfpathlineto{\pgfqpoint{1.217543in}{1.578139in}}%
\pgfpathlineto{\pgfqpoint{1.246809in}{1.625761in}}%
\pgfpathlineto{\pgfqpoint{1.275397in}{1.676827in}}%
\pgfpathlineto{\pgfqpoint{1.303382in}{1.743917in}}%
\pgfpathlineto{\pgfqpoint{1.330833in}{1.804507in}}%
\pgfpathlineto{\pgfqpoint{1.357807in}{1.875909in}}%
\pgfpathlineto{\pgfqpoint{1.384354in}{1.957260in}}%
\pgfpathlineto{\pgfqpoint{1.410517in}{2.034189in}}%
\pgfpathlineto{\pgfqpoint{1.436334in}{2.098314in}}%
\pgfpathlineto{\pgfqpoint{1.461840in}{2.159602in}}%
\pgfpathlineto{\pgfqpoint{1.487063in}{2.208509in}}%
\pgfpathlineto{\pgfqpoint{1.512030in}{2.236913in}}%
\pgfpathlineto{\pgfqpoint{1.536766in}{2.229038in}}%
\pgfpathlineto{\pgfqpoint{1.561290in}{2.237059in}}%
\pgfpathlineto{\pgfqpoint{1.585622in}{2.215671in}}%
\pgfpathlineto{\pgfqpoint{1.609778in}{2.188203in}}%
\pgfpathlineto{\pgfqpoint{1.633774in}{2.129366in}}%
\pgfpathlineto{\pgfqpoint{1.657624in}{2.081908in}}%
\pgfpathlineto{\pgfqpoint{1.681340in}{2.040726in}}%
\pgfpathlineto{\pgfqpoint{1.704933in}{1.984582in}}%
\pgfpathlineto{\pgfqpoint{1.728413in}{1.926301in}}%
\pgfpathlineto{\pgfqpoint{1.751791in}{1.875538in}}%
\pgfpathlineto{\pgfqpoint{1.775073in}{1.830748in}}%
\pgfpathlineto{\pgfqpoint{1.798269in}{1.789266in}}%
\pgfpathlineto{\pgfqpoint{1.821384in}{1.743488in}}%
\pgfpathlineto{\pgfqpoint{1.867400in}{1.676503in}}%
\pgfpathlineto{\pgfqpoint{1.890311in}{1.647124in}}%
\pgfpathlineto{\pgfqpoint{1.913165in}{1.622040in}}%
\pgfpathlineto{\pgfqpoint{1.935966in}{1.603455in}}%
\pgfpathlineto{\pgfqpoint{1.958719in}{1.586885in}}%
\pgfpathlineto{\pgfqpoint{2.004092in}{1.556096in}}%
\pgfpathlineto{\pgfqpoint{2.026719in}{1.544000in}}%
\pgfpathlineto{\pgfqpoint{2.049312in}{1.533774in}}%
\pgfpathlineto{\pgfqpoint{2.071871in}{1.524279in}}%
\pgfpathlineto{\pgfqpoint{2.094401in}{1.517495in}}%
\pgfpathlineto{\pgfqpoint{2.139378in}{1.507027in}}%
\pgfpathlineto{\pgfqpoint{2.161830in}{1.504000in}}%
\pgfpathlineto{\pgfqpoint{2.184260in}{1.501784in}}%
\pgfpathlineto{\pgfqpoint{2.205858in}{1.495097in}}%
\pgfpathlineto{\pgfqpoint{2.225949in}{1.486323in}}%
\pgfpathlineto{\pgfqpoint{2.244728in}{1.488864in}}%
\pgfpathlineto{\pgfqpoint{2.262357in}{1.492117in}}%
\pgfpathlineto{\pgfqpoint{2.278969in}{1.483016in}}%
\pgfpathlineto{\pgfqpoint{2.294675in}{1.485785in}}%
\pgfpathlineto{\pgfqpoint{2.309568in}{1.479431in}}%
\pgfpathlineto{\pgfqpoint{2.337224in}{1.475705in}}%
\pgfpathlineto{\pgfqpoint{2.362456in}{1.471250in}}%
\pgfpathlineto{\pgfqpoint{2.374288in}{1.468544in}}%
\pgfpathlineto{\pgfqpoint{2.385654in}{1.468338in}}%
\pgfpathlineto{\pgfqpoint{2.396587in}{1.466861in}}%
\pgfpathlineto{\pgfqpoint{2.407121in}{1.464271in}}%
\pgfpathlineto{\pgfqpoint{2.427098in}{1.463786in}}%
\pgfpathlineto{\pgfqpoint{2.445778in}{1.459288in}}%
\pgfpathlineto{\pgfqpoint{2.454683in}{1.457455in}}%
\pgfpathlineto{\pgfqpoint{2.463320in}{1.457403in}}%
\pgfpathlineto{\pgfqpoint{2.479855in}{1.453556in}}%
\pgfpathlineto{\pgfqpoint{2.503002in}{1.451394in}}%
\pgfpathlineto{\pgfqpoint{2.517460in}{1.449935in}}%
\pgfpathlineto{\pgfqpoint{2.524426in}{1.449474in}}%
\pgfpathlineto{\pgfqpoint{2.531227in}{1.447153in}}%
\pgfpathlineto{\pgfqpoint{2.537871in}{1.446248in}}%
\pgfpathlineto{\pgfqpoint{2.544365in}{1.444750in}}%
\pgfpathlineto{\pgfqpoint{2.563013in}{1.443752in}}%
\pgfpathlineto{\pgfqpoint{2.568970in}{1.442105in}}%
\pgfpathlineto{\pgfqpoint{2.597035in}{1.439139in}}%
\pgfpathlineto{\pgfqpoint{2.607537in}{1.436820in}}%
\pgfpathlineto{\pgfqpoint{2.612648in}{1.436978in}}%
\pgfpathlineto{\pgfqpoint{2.622606in}{1.434402in}}%
\pgfpathlineto{\pgfqpoint{2.627458in}{1.434183in}}%
\pgfpathlineto{\pgfqpoint{2.636925in}{1.431244in}}%
\pgfpathlineto{\pgfqpoint{2.641544in}{1.431109in}}%
\pgfpathlineto{\pgfqpoint{2.646090in}{1.429688in}}%
\pgfpathlineto{\pgfqpoint{2.663589in}{1.427099in}}%
\pgfpathlineto{\pgfqpoint{2.667802in}{1.427877in}}%
\pgfpathlineto{\pgfqpoint{2.676048in}{1.425645in}}%
\pgfpathlineto{\pgfqpoint{2.684065in}{1.425561in}}%
\pgfpathlineto{\pgfqpoint{2.695686in}{1.423410in}}%
\pgfpathlineto{\pgfqpoint{2.703181in}{1.422985in}}%
\pgfpathlineto{\pgfqpoint{2.710486in}{1.421946in}}%
\pgfpathlineto{\pgfqpoint{2.734685in}{1.419428in}}%
\pgfpathlineto{\pgfqpoint{2.747652in}{1.418170in}}%
\pgfpathlineto{\pgfqpoint{2.757007in}{1.417596in}}%
\pgfpathlineto{\pgfqpoint{2.760060in}{1.417362in}}%
\pgfpathlineto{\pgfqpoint{2.763080in}{1.418483in}}%
\pgfpathlineto{\pgfqpoint{2.766068in}{1.417165in}}%
\pgfpathlineto{\pgfqpoint{2.769027in}{1.417638in}}%
\pgfpathlineto{\pgfqpoint{2.771955in}{1.416938in}}%
\pgfpathlineto{\pgfqpoint{2.780564in}{1.416987in}}%
\pgfpathlineto{\pgfqpoint{2.802338in}{1.417015in}}%
\pgfpathlineto{\pgfqpoint{2.825014in}{1.416955in}}%
\pgfpathlineto{\pgfqpoint{2.854905in}{1.416981in}}%
\pgfpathlineto{\pgfqpoint{2.949066in}{1.416974in}}%
\pgfpathlineto{\pgfqpoint{3.043110in}{1.416988in}}%
\pgfpathlineto{\pgfqpoint{3.043110in}{1.416988in}}%
\pgfusepath{stroke}%
\end{pgfscope}%
\begin{pgfscope}%
\pgfpathrectangle{\pgfqpoint{0.650000in}{0.420000in}}{\pgfqpoint{2.388110in}{1.994488in}}%
\pgfusepath{clip}%
\pgfsetbuttcap%
\pgfsetroundjoin%
\pgfsetlinewidth{1.003750pt}%
\definecolor{currentstroke}{rgb}{0.454902,0.678431,0.819608}%
\pgfsetstrokecolor{currentstroke}%
\pgfsetdash{{5.000000pt}{1.000000pt}}{0.000000pt}%
\pgfpathmoveto{\pgfqpoint{0.645000in}{1.417019in}}%
\pgfpathlineto{\pgfqpoint{1.187505in}{1.417673in}}%
\pgfpathlineto{\pgfqpoint{1.357807in}{1.419157in}}%
\pgfpathlineto{\pgfqpoint{1.436334in}{1.421197in}}%
\pgfpathlineto{\pgfqpoint{1.633774in}{1.425688in}}%
\pgfpathlineto{\pgfqpoint{1.657624in}{1.426791in}}%
\pgfpathlineto{\pgfqpoint{1.704933in}{1.427752in}}%
\pgfpathlineto{\pgfqpoint{1.775073in}{1.425634in}}%
\pgfpathlineto{\pgfqpoint{1.798269in}{1.424425in}}%
\pgfpathlineto{\pgfqpoint{1.821384in}{1.424484in}}%
\pgfpathlineto{\pgfqpoint{1.844426in}{1.422958in}}%
\pgfpathlineto{\pgfqpoint{1.935966in}{1.419984in}}%
\pgfpathlineto{\pgfqpoint{1.958719in}{1.418046in}}%
\pgfpathlineto{\pgfqpoint{2.004092in}{1.412535in}}%
\pgfpathlineto{\pgfqpoint{2.026719in}{1.411538in}}%
\pgfpathlineto{\pgfqpoint{2.049312in}{1.411868in}}%
\pgfpathlineto{\pgfqpoint{2.071871in}{1.410944in}}%
\pgfpathlineto{\pgfqpoint{2.094401in}{1.411173in}}%
\pgfpathlineto{\pgfqpoint{2.116902in}{1.412172in}}%
\pgfpathlineto{\pgfqpoint{2.139378in}{1.411125in}}%
\pgfpathlineto{\pgfqpoint{2.161830in}{1.411291in}}%
\pgfpathlineto{\pgfqpoint{2.184260in}{1.412348in}}%
\pgfpathlineto{\pgfqpoint{2.205858in}{1.412417in}}%
\pgfpathlineto{\pgfqpoint{2.244728in}{1.416759in}}%
\pgfpathlineto{\pgfqpoint{2.262357in}{1.417790in}}%
\pgfpathlineto{\pgfqpoint{2.278969in}{1.415684in}}%
\pgfpathlineto{\pgfqpoint{2.294675in}{1.415780in}}%
\pgfpathlineto{\pgfqpoint{2.309568in}{1.412766in}}%
\pgfpathlineto{\pgfqpoint{2.362456in}{1.407389in}}%
\pgfpathlineto{\pgfqpoint{2.374288in}{1.408125in}}%
\pgfpathlineto{\pgfqpoint{2.385654in}{1.409780in}}%
\pgfpathlineto{\pgfqpoint{2.407121in}{1.409245in}}%
\pgfpathlineto{\pgfqpoint{2.417283in}{1.410351in}}%
\pgfpathlineto{\pgfqpoint{2.427098in}{1.410447in}}%
\pgfpathlineto{\pgfqpoint{2.436590in}{1.411754in}}%
\pgfpathlineto{\pgfqpoint{2.445778in}{1.411634in}}%
\pgfpathlineto{\pgfqpoint{2.454683in}{1.410646in}}%
\pgfpathlineto{\pgfqpoint{2.463320in}{1.411060in}}%
\pgfpathlineto{\pgfqpoint{2.471706in}{1.410837in}}%
\pgfpathlineto{\pgfqpoint{2.479855in}{1.412314in}}%
\pgfpathlineto{\pgfqpoint{2.487779in}{1.411146in}}%
\pgfpathlineto{\pgfqpoint{2.495491in}{1.412192in}}%
\pgfpathlineto{\pgfqpoint{2.503002in}{1.414095in}}%
\pgfpathlineto{\pgfqpoint{2.510322in}{1.414037in}}%
\pgfpathlineto{\pgfqpoint{2.524426in}{1.416842in}}%
\pgfpathlineto{\pgfqpoint{2.531227in}{1.416244in}}%
\pgfpathlineto{\pgfqpoint{2.537871in}{1.414876in}}%
\pgfpathlineto{\pgfqpoint{2.568970in}{1.411463in}}%
\pgfpathlineto{\pgfqpoint{2.586135in}{1.412829in}}%
\pgfpathlineto{\pgfqpoint{2.591637in}{1.414074in}}%
\pgfpathlineto{\pgfqpoint{2.597035in}{1.414182in}}%
\pgfpathlineto{\pgfqpoint{2.607537in}{1.413088in}}%
\pgfpathlineto{\pgfqpoint{2.627458in}{1.413275in}}%
\pgfpathlineto{\pgfqpoint{2.632231in}{1.414065in}}%
\pgfpathlineto{\pgfqpoint{2.636925in}{1.413394in}}%
\pgfpathlineto{\pgfqpoint{2.646090in}{1.414282in}}%
\pgfpathlineto{\pgfqpoint{2.659313in}{1.416192in}}%
\pgfpathlineto{\pgfqpoint{2.671955in}{1.415406in}}%
\pgfpathlineto{\pgfqpoint{2.676048in}{1.413392in}}%
\pgfpathlineto{\pgfqpoint{2.684065in}{1.414455in}}%
\pgfpathlineto{\pgfqpoint{2.687991in}{1.414029in}}%
\pgfpathlineto{\pgfqpoint{2.691864in}{1.414872in}}%
\pgfpathlineto{\pgfqpoint{2.695686in}{1.414382in}}%
\pgfpathlineto{\pgfqpoint{2.699458in}{1.415477in}}%
\pgfpathlineto{\pgfqpoint{2.703181in}{1.415609in}}%
\pgfpathlineto{\pgfqpoint{2.706857in}{1.416972in}}%
\pgfpathlineto{\pgfqpoint{2.717610in}{1.415912in}}%
\pgfpathlineto{\pgfqpoint{2.724562in}{1.414599in}}%
\pgfpathlineto{\pgfqpoint{2.744465in}{1.414540in}}%
\pgfpathlineto{\pgfqpoint{2.747652in}{1.415518in}}%
\pgfpathlineto{\pgfqpoint{2.750804in}{1.417825in}}%
\pgfpathlineto{\pgfqpoint{2.753923in}{1.414956in}}%
\pgfpathlineto{\pgfqpoint{2.757007in}{1.417475in}}%
\pgfpathlineto{\pgfqpoint{2.760060in}{1.415835in}}%
\pgfpathlineto{\pgfqpoint{2.763080in}{1.417208in}}%
\pgfpathlineto{\pgfqpoint{2.766068in}{1.415609in}}%
\pgfpathlineto{\pgfqpoint{2.769027in}{1.416591in}}%
\pgfpathlineto{\pgfqpoint{2.774853in}{1.417032in}}%
\pgfpathlineto{\pgfqpoint{2.777723in}{1.415671in}}%
\pgfpathlineto{\pgfqpoint{2.780564in}{1.415880in}}%
\pgfpathlineto{\pgfqpoint{2.783378in}{1.415491in}}%
\pgfpathlineto{\pgfqpoint{2.786164in}{1.416110in}}%
\pgfpathlineto{\pgfqpoint{2.788924in}{1.415702in}}%
\pgfpathlineto{\pgfqpoint{2.791657in}{1.416474in}}%
\pgfpathlineto{\pgfqpoint{2.794365in}{1.415920in}}%
\pgfpathlineto{\pgfqpoint{2.799705in}{1.416385in}}%
\pgfpathlineto{\pgfqpoint{2.802338in}{1.415725in}}%
\pgfpathlineto{\pgfqpoint{2.804948in}{1.416726in}}%
\pgfpathlineto{\pgfqpoint{2.807534in}{1.416255in}}%
\pgfpathlineto{\pgfqpoint{2.810097in}{1.417309in}}%
\pgfpathlineto{\pgfqpoint{2.812638in}{1.417621in}}%
\pgfpathlineto{\pgfqpoint{2.815156in}{1.416919in}}%
\pgfpathlineto{\pgfqpoint{2.825014in}{1.416406in}}%
\pgfpathlineto{\pgfqpoint{2.827427in}{1.417020in}}%
\pgfpathlineto{\pgfqpoint{2.829820in}{1.416662in}}%
\pgfpathlineto{\pgfqpoint{2.832193in}{1.415569in}}%
\pgfpathlineto{\pgfqpoint{2.834546in}{1.416461in}}%
\pgfpathlineto{\pgfqpoint{2.843773in}{1.416821in}}%
\pgfpathlineto{\pgfqpoint{2.848278in}{1.416674in}}%
\pgfpathlineto{\pgfqpoint{2.850504in}{1.417278in}}%
\pgfpathlineto{\pgfqpoint{2.852713in}{1.416206in}}%
\pgfpathlineto{\pgfqpoint{2.859240in}{1.417081in}}%
\pgfpathlineto{\pgfqpoint{2.863511in}{1.417313in}}%
\pgfpathlineto{\pgfqpoint{2.908550in}{1.416995in}}%
\pgfpathlineto{\pgfqpoint{2.912153in}{1.417100in}}%
\pgfpathlineto{\pgfqpoint{2.913937in}{1.416621in}}%
\pgfpathlineto{\pgfqpoint{2.915711in}{1.417266in}}%
\pgfpathlineto{\pgfqpoint{2.917474in}{1.416655in}}%
\pgfpathlineto{\pgfqpoint{2.920967in}{1.417040in}}%
\pgfpathlineto{\pgfqpoint{2.939459in}{1.416916in}}%
\pgfpathlineto{\pgfqpoint{2.952199in}{1.416988in}}%
\pgfpathlineto{\pgfqpoint{2.971769in}{1.417050in}}%
\pgfpathlineto{\pgfqpoint{2.974665in}{1.416947in}}%
\pgfpathlineto{\pgfqpoint{3.007321in}{1.417064in}}%
\pgfpathlineto{\pgfqpoint{3.008605in}{1.416646in}}%
\pgfpathlineto{\pgfqpoint{3.012421in}{1.417112in}}%
\pgfpathlineto{\pgfqpoint{3.043110in}{1.417011in}}%
\pgfpathlineto{\pgfqpoint{3.043110in}{1.417011in}}%
\pgfusepath{stroke}%
\end{pgfscope}%
\begin{pgfscope}%
\pgfpathrectangle{\pgfqpoint{0.650000in}{0.420000in}}{\pgfqpoint{2.388110in}{1.994488in}}%
\pgfusepath{clip}%
\pgfsetbuttcap%
\pgfsetroundjoin%
\pgfsetlinewidth{1.003750pt}%
\definecolor{currentstroke}{rgb}{0.454902,0.678431,0.819608}%
\pgfsetstrokecolor{currentstroke}%
\pgfsetdash{{3.000000pt}{1.000000pt}{1.000000pt}{1.000000pt}}{0.000000pt}%
\pgfpathmoveto{\pgfqpoint{0.645000in}{0.545076in}}%
\pgfpathlineto{\pgfqpoint{0.679586in}{0.533824in}}%
\pgfpathlineto{\pgfqpoint{0.744946in}{0.549621in}}%
\pgfpathlineto{\pgfqpoint{0.802042in}{0.576152in}}%
\pgfpathlineto{\pgfqpoint{0.853159in}{0.607339in}}%
\pgfpathlineto{\pgfqpoint{0.899759in}{0.640036in}}%
\pgfpathlineto{\pgfqpoint{0.942836in}{0.672523in}}%
\pgfpathlineto{\pgfqpoint{1.057113in}{0.760980in}}%
\pgfpathlineto{\pgfqpoint{1.091568in}{0.786355in}}%
\pgfpathlineto{\pgfqpoint{1.124662in}{0.809483in}}%
\pgfpathlineto{\pgfqpoint{1.156588in}{0.830361in}}%
\pgfpathlineto{\pgfqpoint{1.187505in}{0.849030in}}%
\pgfpathlineto{\pgfqpoint{1.217543in}{0.865556in}}%
\pgfpathlineto{\pgfqpoint{1.246809in}{0.880030in}}%
\pgfpathlineto{\pgfqpoint{1.275397in}{0.892558in}}%
\pgfpathlineto{\pgfqpoint{1.303382in}{0.903261in}}%
\pgfpathlineto{\pgfqpoint{1.330833in}{0.912273in}}%
\pgfpathlineto{\pgfqpoint{1.357807in}{0.919739in}}%
\pgfpathlineto{\pgfqpoint{1.384354in}{0.925817in}}%
\pgfpathlineto{\pgfqpoint{1.410517in}{0.930675in}}%
\pgfpathlineto{\pgfqpoint{1.436334in}{0.934496in}}%
\pgfpathlineto{\pgfqpoint{1.461840in}{0.937472in}}%
\pgfpathlineto{\pgfqpoint{1.512030in}{0.941729in}}%
\pgfpathlineto{\pgfqpoint{1.536766in}{0.943465in}}%
\pgfpathlineto{\pgfqpoint{1.561290in}{0.942275in}}%
\pgfpathlineto{\pgfqpoint{1.585622in}{0.947990in}}%
\pgfpathlineto{\pgfqpoint{1.609778in}{0.960632in}}%
\pgfpathlineto{\pgfqpoint{1.633774in}{0.957945in}}%
\pgfpathlineto{\pgfqpoint{1.657624in}{0.972978in}}%
\pgfpathlineto{\pgfqpoint{1.681340in}{0.996134in}}%
\pgfpathlineto{\pgfqpoint{1.728413in}{1.030187in}}%
\pgfpathlineto{\pgfqpoint{1.751791in}{1.049551in}}%
\pgfpathlineto{\pgfqpoint{1.775073in}{1.072083in}}%
\pgfpathlineto{\pgfqpoint{1.798269in}{1.093631in}}%
\pgfpathlineto{\pgfqpoint{1.821384in}{1.113157in}}%
\pgfpathlineto{\pgfqpoint{1.844426in}{1.135000in}}%
\pgfpathlineto{\pgfqpoint{1.867400in}{1.155901in}}%
\pgfpathlineto{\pgfqpoint{1.913165in}{1.194095in}}%
\pgfpathlineto{\pgfqpoint{1.935966in}{1.211908in}}%
\pgfpathlineto{\pgfqpoint{1.958719in}{1.228096in}}%
\pgfpathlineto{\pgfqpoint{1.981426in}{1.243180in}}%
\pgfpathlineto{\pgfqpoint{2.004092in}{1.257047in}}%
\pgfpathlineto{\pgfqpoint{2.026719in}{1.270032in}}%
\pgfpathlineto{\pgfqpoint{2.049312in}{1.281725in}}%
\pgfpathlineto{\pgfqpoint{2.071871in}{1.292291in}}%
\pgfpathlineto{\pgfqpoint{2.094401in}{1.301317in}}%
\pgfpathlineto{\pgfqpoint{2.116902in}{1.309445in}}%
\pgfpathlineto{\pgfqpoint{2.139378in}{1.316860in}}%
\pgfpathlineto{\pgfqpoint{2.184260in}{1.329728in}}%
\pgfpathlineto{\pgfqpoint{2.205858in}{1.335280in}}%
\pgfpathlineto{\pgfqpoint{2.244728in}{1.343378in}}%
\pgfpathlineto{\pgfqpoint{2.337224in}{1.359656in}}%
\pgfpathlineto{\pgfqpoint{2.495491in}{1.383930in}}%
\pgfpathlineto{\pgfqpoint{2.531227in}{1.388678in}}%
\pgfpathlineto{\pgfqpoint{2.550716in}{1.391121in}}%
\pgfpathlineto{\pgfqpoint{2.568970in}{1.392901in}}%
\pgfpathlineto{\pgfqpoint{2.591637in}{1.395965in}}%
\pgfpathlineto{\pgfqpoint{2.622606in}{1.400169in}}%
\pgfpathlineto{\pgfqpoint{2.636925in}{1.402368in}}%
\pgfpathlineto{\pgfqpoint{2.641544in}{1.402458in}}%
\pgfpathlineto{\pgfqpoint{2.646090in}{1.403746in}}%
\pgfpathlineto{\pgfqpoint{2.734685in}{1.413853in}}%
\pgfpathlineto{\pgfqpoint{2.760060in}{1.415560in}}%
\pgfpathlineto{\pgfqpoint{2.763080in}{1.415149in}}%
\pgfpathlineto{\pgfqpoint{2.769027in}{1.415899in}}%
\pgfpathlineto{\pgfqpoint{2.807534in}{1.416811in}}%
\pgfpathlineto{\pgfqpoint{2.895573in}{1.416974in}}%
\pgfpathlineto{\pgfqpoint{3.043110in}{1.416975in}}%
\pgfpathlineto{\pgfqpoint{3.043110in}{1.416975in}}%
\pgfusepath{stroke}%
\end{pgfscope}%
\begin{pgfscope}%
\pgfpathrectangle{\pgfqpoint{0.650000in}{0.420000in}}{\pgfqpoint{2.388110in}{1.994488in}}%
\pgfusepath{clip}%
\pgfsetbuttcap%
\pgfsetroundjoin%
\pgfsetlinewidth{1.003750pt}%
\definecolor{currentstroke}{rgb}{0.454902,0.678431,0.819608}%
\pgfsetstrokecolor{currentstroke}%
\pgfsetdash{{1.000000pt}{1.000000pt}}{0.000000pt}%
\pgfpathmoveto{\pgfqpoint{0.645000in}{1.463431in}}%
\pgfpathlineto{\pgfqpoint{0.744946in}{1.472640in}}%
\pgfpathlineto{\pgfqpoint{0.853159in}{1.484245in}}%
\pgfpathlineto{\pgfqpoint{0.942836in}{1.494331in}}%
\pgfpathlineto{\pgfqpoint{0.983096in}{1.497686in}}%
\pgfpathlineto{\pgfqpoint{1.021057in}{1.499302in}}%
\pgfpathlineto{\pgfqpoint{1.057113in}{1.499208in}}%
\pgfpathlineto{\pgfqpoint{1.124662in}{1.497238in}}%
\pgfpathlineto{\pgfqpoint{1.156588in}{1.497721in}}%
\pgfpathlineto{\pgfqpoint{1.187505in}{1.499027in}}%
\pgfpathlineto{\pgfqpoint{1.217543in}{1.497766in}}%
\pgfpathlineto{\pgfqpoint{1.246809in}{1.489779in}}%
\pgfpathlineto{\pgfqpoint{1.275397in}{1.488434in}}%
\pgfpathlineto{\pgfqpoint{1.303382in}{1.473005in}}%
\pgfpathlineto{\pgfqpoint{1.330833in}{1.492914in}}%
\pgfpathlineto{\pgfqpoint{1.384354in}{1.472768in}}%
\pgfpathlineto{\pgfqpoint{1.410517in}{1.466947in}}%
\pgfpathlineto{\pgfqpoint{1.461840in}{1.457219in}}%
\pgfpathlineto{\pgfqpoint{1.487063in}{1.449591in}}%
\pgfpathlineto{\pgfqpoint{1.512030in}{1.447288in}}%
\pgfpathlineto{\pgfqpoint{1.536766in}{1.435641in}}%
\pgfpathlineto{\pgfqpoint{1.561290in}{1.432859in}}%
\pgfpathlineto{\pgfqpoint{1.585622in}{1.420264in}}%
\pgfpathlineto{\pgfqpoint{1.609778in}{1.425192in}}%
\pgfpathlineto{\pgfqpoint{1.633774in}{1.478666in}}%
\pgfpathlineto{\pgfqpoint{1.657624in}{1.413242in}}%
\pgfpathlineto{\pgfqpoint{1.681340in}{1.412714in}}%
\pgfpathlineto{\pgfqpoint{1.704933in}{1.398737in}}%
\pgfpathlineto{\pgfqpoint{1.728413in}{1.376904in}}%
\pgfpathlineto{\pgfqpoint{1.751791in}{1.414800in}}%
\pgfpathlineto{\pgfqpoint{1.775073in}{1.412974in}}%
\pgfpathlineto{\pgfqpoint{1.798269in}{1.421666in}}%
\pgfpathlineto{\pgfqpoint{1.821384in}{1.404890in}}%
\pgfpathlineto{\pgfqpoint{1.844426in}{1.408976in}}%
\pgfpathlineto{\pgfqpoint{1.867400in}{1.419367in}}%
\pgfpathlineto{\pgfqpoint{1.913165in}{1.422378in}}%
\pgfpathlineto{\pgfqpoint{1.935966in}{1.416475in}}%
\pgfpathlineto{\pgfqpoint{1.958719in}{1.419307in}}%
\pgfpathlineto{\pgfqpoint{1.981426in}{1.388962in}}%
\pgfpathlineto{\pgfqpoint{2.004092in}{1.416392in}}%
\pgfpathlineto{\pgfqpoint{2.026719in}{1.437874in}}%
\pgfpathlineto{\pgfqpoint{2.049312in}{1.419921in}}%
\pgfpathlineto{\pgfqpoint{2.071871in}{1.420911in}}%
\pgfpathlineto{\pgfqpoint{2.094401in}{1.416448in}}%
\pgfpathlineto{\pgfqpoint{2.116902in}{1.532004in}}%
\pgfpathlineto{\pgfqpoint{2.139378in}{1.415822in}}%
\pgfpathlineto{\pgfqpoint{2.161830in}{1.415896in}}%
\pgfpathlineto{\pgfqpoint{2.184260in}{1.418984in}}%
\pgfpathlineto{\pgfqpoint{2.205858in}{1.421220in}}%
\pgfpathlineto{\pgfqpoint{2.225949in}{1.487228in}}%
\pgfpathlineto{\pgfqpoint{2.244728in}{1.424857in}}%
\pgfpathlineto{\pgfqpoint{2.262357in}{1.336250in}}%
\pgfpathlineto{\pgfqpoint{2.278969in}{1.436260in}}%
\pgfpathlineto{\pgfqpoint{2.294675in}{1.464332in}}%
\pgfpathlineto{\pgfqpoint{2.309568in}{1.404093in}}%
\pgfpathlineto{\pgfqpoint{2.323728in}{1.415375in}}%
\pgfpathlineto{\pgfqpoint{2.337224in}{1.415954in}}%
\pgfpathlineto{\pgfqpoint{2.350117in}{1.413885in}}%
\pgfpathlineto{\pgfqpoint{2.362456in}{1.410514in}}%
\pgfpathlineto{\pgfqpoint{2.374288in}{1.410086in}}%
\pgfpathlineto{\pgfqpoint{2.385654in}{1.401157in}}%
\pgfpathlineto{\pgfqpoint{2.396587in}{1.410534in}}%
\pgfpathlineto{\pgfqpoint{2.407121in}{1.416448in}}%
\pgfpathlineto{\pgfqpoint{2.417283in}{1.434473in}}%
\pgfpathlineto{\pgfqpoint{2.427098in}{1.406211in}}%
\pgfpathlineto{\pgfqpoint{2.436590in}{1.415513in}}%
\pgfpathlineto{\pgfqpoint{2.445778in}{1.418719in}}%
\pgfpathlineto{\pgfqpoint{2.454683in}{1.417986in}}%
\pgfpathlineto{\pgfqpoint{2.463320in}{1.406004in}}%
\pgfpathlineto{\pgfqpoint{2.471706in}{1.410137in}}%
\pgfpathlineto{\pgfqpoint{2.479855in}{1.428922in}}%
\pgfpathlineto{\pgfqpoint{2.487779in}{1.413309in}}%
\pgfpathlineto{\pgfqpoint{2.495491in}{1.410052in}}%
\pgfpathlineto{\pgfqpoint{2.503002in}{1.419171in}}%
\pgfpathlineto{\pgfqpoint{2.510322in}{1.415318in}}%
\pgfpathlineto{\pgfqpoint{2.524426in}{1.416711in}}%
\pgfpathlineto{\pgfqpoint{2.531227in}{1.420560in}}%
\pgfpathlineto{\pgfqpoint{2.537871in}{1.417019in}}%
\pgfpathlineto{\pgfqpoint{2.544365in}{1.419854in}}%
\pgfpathlineto{\pgfqpoint{2.550716in}{1.420125in}}%
\pgfpathlineto{\pgfqpoint{2.556930in}{1.415022in}}%
\pgfpathlineto{\pgfqpoint{2.563013in}{1.416423in}}%
\pgfpathlineto{\pgfqpoint{2.568970in}{1.413302in}}%
\pgfpathlineto{\pgfqpoint{2.574806in}{1.416228in}}%
\pgfpathlineto{\pgfqpoint{2.580526in}{1.415136in}}%
\pgfpathlineto{\pgfqpoint{2.586135in}{1.417164in}}%
\pgfpathlineto{\pgfqpoint{2.591637in}{1.412838in}}%
\pgfpathlineto{\pgfqpoint{2.597035in}{1.419205in}}%
\pgfpathlineto{\pgfqpoint{2.602334in}{1.412957in}}%
\pgfpathlineto{\pgfqpoint{2.607537in}{1.423579in}}%
\pgfpathlineto{\pgfqpoint{2.612648in}{1.409224in}}%
\pgfpathlineto{\pgfqpoint{2.617670in}{1.429011in}}%
\pgfpathlineto{\pgfqpoint{2.622606in}{1.412725in}}%
\pgfpathlineto{\pgfqpoint{2.627458in}{1.420872in}}%
\pgfpathlineto{\pgfqpoint{2.632231in}{1.416469in}}%
\pgfpathlineto{\pgfqpoint{2.636925in}{1.420932in}}%
\pgfpathlineto{\pgfqpoint{2.641544in}{1.436437in}}%
\pgfpathlineto{\pgfqpoint{2.646090in}{1.420576in}}%
\pgfpathlineto{\pgfqpoint{2.650565in}{1.418529in}}%
\pgfpathlineto{\pgfqpoint{2.654972in}{1.418848in}}%
\pgfpathlineto{\pgfqpoint{2.659313in}{1.414846in}}%
\pgfpathlineto{\pgfqpoint{2.663589in}{1.419556in}}%
\pgfpathlineto{\pgfqpoint{2.667802in}{1.371568in}}%
\pgfpathlineto{\pgfqpoint{2.671955in}{1.413206in}}%
\pgfpathlineto{\pgfqpoint{2.676048in}{1.420916in}}%
\pgfpathlineto{\pgfqpoint{2.680085in}{1.416405in}}%
\pgfpathlineto{\pgfqpoint{2.687991in}{1.415505in}}%
\pgfpathlineto{\pgfqpoint{2.691864in}{1.419086in}}%
\pgfpathlineto{\pgfqpoint{2.695686in}{1.418931in}}%
\pgfpathlineto{\pgfqpoint{2.699458in}{1.414854in}}%
\pgfpathlineto{\pgfqpoint{2.703181in}{1.417205in}}%
\pgfpathlineto{\pgfqpoint{2.706857in}{1.414852in}}%
\pgfpathlineto{\pgfqpoint{2.710486in}{1.415833in}}%
\pgfpathlineto{\pgfqpoint{2.714070in}{1.414957in}}%
\pgfpathlineto{\pgfqpoint{2.717610in}{1.414690in}}%
\pgfpathlineto{\pgfqpoint{2.721107in}{1.421584in}}%
\pgfpathlineto{\pgfqpoint{2.724562in}{1.413000in}}%
\pgfpathlineto{\pgfqpoint{2.727976in}{1.417007in}}%
\pgfpathlineto{\pgfqpoint{2.731350in}{1.414158in}}%
\pgfpathlineto{\pgfqpoint{2.734685in}{1.413644in}}%
\pgfpathlineto{\pgfqpoint{2.737982in}{1.419472in}}%
\pgfpathlineto{\pgfqpoint{2.741242in}{1.414866in}}%
\pgfpathlineto{\pgfqpoint{2.744465in}{1.414174in}}%
\pgfpathlineto{\pgfqpoint{2.747652in}{1.416009in}}%
\pgfpathlineto{\pgfqpoint{2.750804in}{1.410732in}}%
\pgfpathlineto{\pgfqpoint{2.757007in}{1.429056in}}%
\pgfpathlineto{\pgfqpoint{2.760060in}{1.411056in}}%
\pgfpathlineto{\pgfqpoint{2.763080in}{1.284228in}}%
\pgfpathlineto{\pgfqpoint{2.766068in}{1.412893in}}%
\pgfpathlineto{\pgfqpoint{2.769027in}{1.380763in}}%
\pgfpathlineto{\pgfqpoint{2.774853in}{1.428843in}}%
\pgfpathlineto{\pgfqpoint{2.777723in}{1.427267in}}%
\pgfpathlineto{\pgfqpoint{2.780564in}{1.429088in}}%
\pgfpathlineto{\pgfqpoint{2.783378in}{1.402655in}}%
\pgfpathlineto{\pgfqpoint{2.786164in}{1.427217in}}%
\pgfpathlineto{\pgfqpoint{2.788924in}{1.419238in}}%
\pgfpathlineto{\pgfqpoint{2.791657in}{1.439071in}}%
\pgfpathlineto{\pgfqpoint{2.794365in}{1.400216in}}%
\pgfpathlineto{\pgfqpoint{2.797047in}{1.416542in}}%
\pgfpathlineto{\pgfqpoint{2.799705in}{1.395449in}}%
\pgfpathlineto{\pgfqpoint{2.802338in}{1.388527in}}%
\pgfpathlineto{\pgfqpoint{2.804948in}{1.437786in}}%
\pgfpathlineto{\pgfqpoint{2.807534in}{1.422041in}}%
\pgfpathlineto{\pgfqpoint{2.810097in}{1.412225in}}%
\pgfpathlineto{\pgfqpoint{2.812638in}{1.406395in}}%
\pgfpathlineto{\pgfqpoint{2.815156in}{1.429042in}}%
\pgfpathlineto{\pgfqpoint{2.817652in}{1.430562in}}%
\pgfpathlineto{\pgfqpoint{2.820127in}{1.420126in}}%
\pgfpathlineto{\pgfqpoint{2.822581in}{1.406400in}}%
\pgfpathlineto{\pgfqpoint{2.825014in}{1.428954in}}%
\pgfpathlineto{\pgfqpoint{2.827427in}{1.407840in}}%
\pgfpathlineto{\pgfqpoint{2.829820in}{1.410213in}}%
\pgfpathlineto{\pgfqpoint{2.832193in}{1.420005in}}%
\pgfpathlineto{\pgfqpoint{2.836881in}{1.401899in}}%
\pgfpathlineto{\pgfqpoint{2.839197in}{1.414386in}}%
\pgfpathlineto{\pgfqpoint{2.841494in}{1.411160in}}%
\pgfpathlineto{\pgfqpoint{2.843773in}{1.417841in}}%
\pgfpathlineto{\pgfqpoint{2.846034in}{1.426495in}}%
\pgfpathlineto{\pgfqpoint{2.848278in}{1.415470in}}%
\pgfpathlineto{\pgfqpoint{2.850504in}{1.419120in}}%
\pgfpathlineto{\pgfqpoint{2.852713in}{1.425946in}}%
\pgfpathlineto{\pgfqpoint{2.854905in}{1.407829in}}%
\pgfpathlineto{\pgfqpoint{2.857081in}{1.424269in}}%
\pgfpathlineto{\pgfqpoint{2.859240in}{1.413658in}}%
\pgfpathlineto{\pgfqpoint{2.861384in}{1.406635in}}%
\pgfpathlineto{\pgfqpoint{2.863511in}{1.413300in}}%
\pgfpathlineto{\pgfqpoint{2.865623in}{1.413732in}}%
\pgfpathlineto{\pgfqpoint{2.867720in}{1.417675in}}%
\pgfpathlineto{\pgfqpoint{2.869801in}{1.416478in}}%
\pgfpathlineto{\pgfqpoint{2.871867in}{1.409975in}}%
\pgfpathlineto{\pgfqpoint{2.873919in}{1.423448in}}%
\pgfpathlineto{\pgfqpoint{2.875956in}{1.417405in}}%
\pgfpathlineto{\pgfqpoint{2.877979in}{1.415253in}}%
\pgfpathlineto{\pgfqpoint{2.879988in}{1.416741in}}%
\pgfpathlineto{\pgfqpoint{2.881983in}{1.415684in}}%
\pgfpathlineto{\pgfqpoint{2.883964in}{1.415950in}}%
\pgfpathlineto{\pgfqpoint{2.885932in}{1.424748in}}%
\pgfpathlineto{\pgfqpoint{2.887886in}{1.413684in}}%
\pgfpathlineto{\pgfqpoint{2.889827in}{1.408729in}}%
\pgfpathlineto{\pgfqpoint{2.891755in}{1.427880in}}%
\pgfpathlineto{\pgfqpoint{2.893670in}{1.419781in}}%
\pgfpathlineto{\pgfqpoint{2.895573in}{1.414263in}}%
\pgfpathlineto{\pgfqpoint{2.897463in}{1.417411in}}%
\pgfpathlineto{\pgfqpoint{2.899341in}{1.414504in}}%
\pgfpathlineto{\pgfqpoint{2.901206in}{1.416038in}}%
\pgfpathlineto{\pgfqpoint{2.903060in}{1.416458in}}%
\pgfpathlineto{\pgfqpoint{2.904901in}{1.418571in}}%
\pgfpathlineto{\pgfqpoint{2.906731in}{1.408332in}}%
\pgfpathlineto{\pgfqpoint{2.908550in}{1.409360in}}%
\pgfpathlineto{\pgfqpoint{2.910357in}{1.417065in}}%
\pgfpathlineto{\pgfqpoint{2.912153in}{1.417574in}}%
\pgfpathlineto{\pgfqpoint{2.913937in}{1.422985in}}%
\pgfpathlineto{\pgfqpoint{2.915711in}{1.410003in}}%
\pgfpathlineto{\pgfqpoint{2.917474in}{1.431557in}}%
\pgfpathlineto{\pgfqpoint{2.919226in}{1.411561in}}%
\pgfpathlineto{\pgfqpoint{2.920967in}{1.419127in}}%
\pgfpathlineto{\pgfqpoint{2.922698in}{1.415372in}}%
\pgfpathlineto{\pgfqpoint{2.924419in}{1.415790in}}%
\pgfpathlineto{\pgfqpoint{2.926130in}{1.418825in}}%
\pgfpathlineto{\pgfqpoint{2.927830in}{1.414741in}}%
\pgfpathlineto{\pgfqpoint{2.929520in}{1.420020in}}%
\pgfpathlineto{\pgfqpoint{2.931201in}{1.410752in}}%
\pgfpathlineto{\pgfqpoint{2.932871in}{1.417248in}}%
\pgfpathlineto{\pgfqpoint{2.934533in}{1.417288in}}%
\pgfpathlineto{\pgfqpoint{2.936184in}{1.424880in}}%
\pgfpathlineto{\pgfqpoint{2.939459in}{1.419088in}}%
\pgfpathlineto{\pgfqpoint{2.941083in}{1.422330in}}%
\pgfpathlineto{\pgfqpoint{2.944303in}{1.415395in}}%
\pgfpathlineto{\pgfqpoint{2.945899in}{1.415709in}}%
\pgfpathlineto{\pgfqpoint{2.947487in}{1.417221in}}%
\pgfpathlineto{\pgfqpoint{2.949066in}{1.417717in}}%
\pgfpathlineto{\pgfqpoint{2.950637in}{1.417351in}}%
\pgfpathlineto{\pgfqpoint{2.952199in}{1.414999in}}%
\pgfpathlineto{\pgfqpoint{2.953752in}{1.417413in}}%
\pgfpathlineto{\pgfqpoint{2.955298in}{1.418459in}}%
\pgfpathlineto{\pgfqpoint{2.956835in}{1.418487in}}%
\pgfpathlineto{\pgfqpoint{2.958363in}{1.422978in}}%
\pgfpathlineto{\pgfqpoint{2.959884in}{1.413431in}}%
\pgfpathlineto{\pgfqpoint{2.961397in}{1.418071in}}%
\pgfpathlineto{\pgfqpoint{2.962902in}{1.418773in}}%
\pgfpathlineto{\pgfqpoint{2.964399in}{1.415442in}}%
\pgfpathlineto{\pgfqpoint{2.965888in}{1.418144in}}%
\pgfpathlineto{\pgfqpoint{2.967370in}{1.415716in}}%
\pgfpathlineto{\pgfqpoint{2.968844in}{1.416469in}}%
\pgfpathlineto{\pgfqpoint{2.970310in}{1.416214in}}%
\pgfpathlineto{\pgfqpoint{2.971769in}{1.418162in}}%
\pgfpathlineto{\pgfqpoint{2.973221in}{1.417168in}}%
\pgfpathlineto{\pgfqpoint{2.974665in}{1.417151in}}%
\pgfpathlineto{\pgfqpoint{2.976103in}{1.416174in}}%
\pgfpathlineto{\pgfqpoint{2.977533in}{1.423362in}}%
\pgfpathlineto{\pgfqpoint{2.978956in}{1.421098in}}%
\pgfpathlineto{\pgfqpoint{2.981781in}{1.414092in}}%
\pgfpathlineto{\pgfqpoint{2.983183in}{1.416290in}}%
\pgfpathlineto{\pgfqpoint{2.984578in}{1.412257in}}%
\pgfpathlineto{\pgfqpoint{2.985967in}{1.416090in}}%
\pgfpathlineto{\pgfqpoint{2.987349in}{1.421798in}}%
\pgfpathlineto{\pgfqpoint{2.988725in}{1.416906in}}%
\pgfpathlineto{\pgfqpoint{2.990094in}{1.416377in}}%
\pgfpathlineto{\pgfqpoint{2.991456in}{1.414062in}}%
\pgfpathlineto{\pgfqpoint{2.992812in}{1.420353in}}%
\pgfpathlineto{\pgfqpoint{2.994162in}{1.416284in}}%
\pgfpathlineto{\pgfqpoint{2.995505in}{1.416697in}}%
\pgfpathlineto{\pgfqpoint{2.996842in}{1.411934in}}%
\pgfpathlineto{\pgfqpoint{2.998173in}{1.432501in}}%
\pgfpathlineto{\pgfqpoint{2.999498in}{1.419136in}}%
\pgfpathlineto{\pgfqpoint{3.000817in}{1.414259in}}%
\pgfpathlineto{\pgfqpoint{3.002129in}{1.417473in}}%
\pgfpathlineto{\pgfqpoint{3.003436in}{1.412798in}}%
\pgfpathlineto{\pgfqpoint{3.004737in}{1.419081in}}%
\pgfpathlineto{\pgfqpoint{3.006032in}{1.411652in}}%
\pgfpathlineto{\pgfqpoint{3.007321in}{1.409587in}}%
\pgfpathlineto{\pgfqpoint{3.008605in}{1.411450in}}%
\pgfpathlineto{\pgfqpoint{3.009882in}{1.422337in}}%
\pgfpathlineto{\pgfqpoint{3.011155in}{1.416993in}}%
\pgfpathlineto{\pgfqpoint{3.012421in}{1.417469in}}%
\pgfpathlineto{\pgfqpoint{3.013682in}{1.415720in}}%
\pgfpathlineto{\pgfqpoint{3.016188in}{1.416393in}}%
\pgfpathlineto{\pgfqpoint{3.018672in}{1.419922in}}%
\pgfpathlineto{\pgfqpoint{3.022358in}{1.414485in}}%
\pgfpathlineto{\pgfqpoint{3.024789in}{1.417558in}}%
\pgfpathlineto{\pgfqpoint{3.025997in}{1.416758in}}%
\pgfpathlineto{\pgfqpoint{3.031963in}{1.417091in}}%
\pgfpathlineto{\pgfqpoint{3.033141in}{1.416236in}}%
\pgfpathlineto{\pgfqpoint{3.035484in}{1.417120in}}%
\pgfpathlineto{\pgfqpoint{3.037807in}{1.417366in}}%
\pgfpathlineto{\pgfqpoint{3.038962in}{1.416571in}}%
\pgfpathlineto{\pgfqpoint{3.040112in}{1.417697in}}%
\pgfpathlineto{\pgfqpoint{3.041258in}{1.414223in}}%
\pgfpathlineto{\pgfqpoint{3.042399in}{1.417809in}}%
\pgfpathlineto{\pgfqpoint{3.043110in}{1.416630in}}%
\pgfpathlineto{\pgfqpoint{3.043110in}{1.416630in}}%
\pgfusepath{stroke}%
\end{pgfscope}%
\begin{pgfscope}%
\pgfpathrectangle{\pgfqpoint{0.650000in}{0.420000in}}{\pgfqpoint{2.388110in}{1.994488in}}%
\pgfusepath{clip}%
\pgfsetbuttcap%
\pgfsetroundjoin%
\pgfsetlinewidth{1.003750pt}%
\definecolor{currentstroke}{rgb}{0.454902,0.678431,0.819608}%
\pgfsetstrokecolor{currentstroke}%
\pgfsetdash{{5.000000pt}{5.000000pt}}{0.000000pt}%
\pgfpathmoveto{\pgfqpoint{0.645000in}{1.417439in}}%
\pgfpathlineto{\pgfqpoint{0.679586in}{1.417574in}}%
\pgfpathlineto{\pgfqpoint{0.744946in}{1.419072in}}%
\pgfpathlineto{\pgfqpoint{0.802042in}{1.419654in}}%
\pgfpathlineto{\pgfqpoint{0.853159in}{1.422702in}}%
\pgfpathlineto{\pgfqpoint{0.899759in}{1.426553in}}%
\pgfpathlineto{\pgfqpoint{0.942836in}{1.432127in}}%
\pgfpathlineto{\pgfqpoint{0.983096in}{1.442081in}}%
\pgfpathlineto{\pgfqpoint{1.021057in}{1.452708in}}%
\pgfpathlineto{\pgfqpoint{1.057113in}{1.467340in}}%
\pgfpathlineto{\pgfqpoint{1.091568in}{1.482855in}}%
\pgfpathlineto{\pgfqpoint{1.124662in}{1.501192in}}%
\pgfpathlineto{\pgfqpoint{1.156588in}{1.523006in}}%
\pgfpathlineto{\pgfqpoint{1.187505in}{1.546137in}}%
\pgfpathlineto{\pgfqpoint{1.217543in}{1.569921in}}%
\pgfpathlineto{\pgfqpoint{1.246809in}{1.589136in}}%
\pgfpathlineto{\pgfqpoint{1.275397in}{1.598413in}}%
\pgfpathlineto{\pgfqpoint{1.303382in}{1.600727in}}%
\pgfpathlineto{\pgfqpoint{1.330833in}{1.590803in}}%
\pgfpathlineto{\pgfqpoint{1.357807in}{1.568957in}}%
\pgfpathlineto{\pgfqpoint{1.384354in}{1.534600in}}%
\pgfpathlineto{\pgfqpoint{1.410517in}{1.489339in}}%
\pgfpathlineto{\pgfqpoint{1.436334in}{1.436114in}}%
\pgfpathlineto{\pgfqpoint{1.487063in}{1.323943in}}%
\pgfpathlineto{\pgfqpoint{1.512030in}{1.275877in}}%
\pgfpathlineto{\pgfqpoint{1.536766in}{1.238768in}}%
\pgfpathlineto{\pgfqpoint{1.561290in}{1.213284in}}%
\pgfpathlineto{\pgfqpoint{1.585622in}{1.201032in}}%
\pgfpathlineto{\pgfqpoint{1.609778in}{1.201196in}}%
\pgfpathlineto{\pgfqpoint{1.633774in}{1.211845in}}%
\pgfpathlineto{\pgfqpoint{1.657624in}{1.230187in}}%
\pgfpathlineto{\pgfqpoint{1.681340in}{1.253263in}}%
\pgfpathlineto{\pgfqpoint{1.751791in}{1.329628in}}%
\pgfpathlineto{\pgfqpoint{1.775073in}{1.351934in}}%
\pgfpathlineto{\pgfqpoint{1.798269in}{1.371111in}}%
\pgfpathlineto{\pgfqpoint{1.821384in}{1.386895in}}%
\pgfpathlineto{\pgfqpoint{1.844426in}{1.399815in}}%
\pgfpathlineto{\pgfqpoint{1.867400in}{1.410194in}}%
\pgfpathlineto{\pgfqpoint{1.890311in}{1.418687in}}%
\pgfpathlineto{\pgfqpoint{1.913165in}{1.425515in}}%
\pgfpathlineto{\pgfqpoint{1.935966in}{1.430960in}}%
\pgfpathlineto{\pgfqpoint{1.958719in}{1.435314in}}%
\pgfpathlineto{\pgfqpoint{1.981426in}{1.438383in}}%
\pgfpathlineto{\pgfqpoint{2.004092in}{1.439453in}}%
\pgfpathlineto{\pgfqpoint{2.026719in}{1.438368in}}%
\pgfpathlineto{\pgfqpoint{2.071871in}{1.434064in}}%
\pgfpathlineto{\pgfqpoint{2.116902in}{1.431943in}}%
\pgfpathlineto{\pgfqpoint{2.139378in}{1.429035in}}%
\pgfpathlineto{\pgfqpoint{2.184260in}{1.419857in}}%
\pgfpathlineto{\pgfqpoint{2.205858in}{1.416335in}}%
\pgfpathlineto{\pgfqpoint{2.225949in}{1.413652in}}%
\pgfpathlineto{\pgfqpoint{2.244728in}{1.411722in}}%
\pgfpathlineto{\pgfqpoint{2.262357in}{1.411525in}}%
\pgfpathlineto{\pgfqpoint{2.278969in}{1.412862in}}%
\pgfpathlineto{\pgfqpoint{2.337224in}{1.420979in}}%
\pgfpathlineto{\pgfqpoint{2.350117in}{1.421653in}}%
\pgfpathlineto{\pgfqpoint{2.374288in}{1.420444in}}%
\pgfpathlineto{\pgfqpoint{2.417283in}{1.416912in}}%
\pgfpathlineto{\pgfqpoint{2.436590in}{1.417212in}}%
\pgfpathlineto{\pgfqpoint{2.463320in}{1.419156in}}%
\pgfpathlineto{\pgfqpoint{2.479855in}{1.418650in}}%
\pgfpathlineto{\pgfqpoint{2.487779in}{1.417941in}}%
\pgfpathlineto{\pgfqpoint{2.503002in}{1.414529in}}%
\pgfpathlineto{\pgfqpoint{2.517460in}{1.411118in}}%
\pgfpathlineto{\pgfqpoint{2.524426in}{1.410292in}}%
\pgfpathlineto{\pgfqpoint{2.531227in}{1.410460in}}%
\pgfpathlineto{\pgfqpoint{2.537871in}{1.411366in}}%
\pgfpathlineto{\pgfqpoint{2.550716in}{1.414573in}}%
\pgfpathlineto{\pgfqpoint{2.563013in}{1.418173in}}%
\pgfpathlineto{\pgfqpoint{2.568970in}{1.419178in}}%
\pgfpathlineto{\pgfqpoint{2.580526in}{1.419306in}}%
\pgfpathlineto{\pgfqpoint{2.607537in}{1.417992in}}%
\pgfpathlineto{\pgfqpoint{2.636925in}{1.417551in}}%
\pgfpathlineto{\pgfqpoint{2.659313in}{1.417293in}}%
\pgfpathlineto{\pgfqpoint{2.671955in}{1.418398in}}%
\pgfpathlineto{\pgfqpoint{2.684065in}{1.419600in}}%
\pgfpathlineto{\pgfqpoint{2.695686in}{1.419443in}}%
\pgfpathlineto{\pgfqpoint{2.706857in}{1.419166in}}%
\pgfpathlineto{\pgfqpoint{2.717610in}{1.420497in}}%
\pgfpathlineto{\pgfqpoint{2.727976in}{1.422001in}}%
\pgfpathlineto{\pgfqpoint{2.737982in}{1.421924in}}%
\pgfpathlineto{\pgfqpoint{2.786164in}{1.418385in}}%
\pgfpathlineto{\pgfqpoint{2.804948in}{1.417569in}}%
\pgfpathlineto{\pgfqpoint{2.820127in}{1.417185in}}%
\pgfpathlineto{\pgfqpoint{3.043110in}{1.417007in}}%
\pgfpathlineto{\pgfqpoint{3.043110in}{1.417007in}}%
\pgfusepath{stroke}%
\end{pgfscope}%
\begin{pgfscope}%
\pgfpathrectangle{\pgfqpoint{0.650000in}{0.420000in}}{\pgfqpoint{2.388110in}{1.994488in}}%
\pgfusepath{clip}%
\pgfsetbuttcap%
\pgfsetroundjoin%
\pgfsetlinewidth{1.003750pt}%
\definecolor{currentstroke}{rgb}{0.454902,0.678431,0.819608}%
\pgfsetstrokecolor{currentstroke}%
\pgfsetdash{{3.000000pt}{5.000000pt}{1.000000pt}{5.000000pt}}{0.000000pt}%
\pgfpathmoveto{\pgfqpoint{0.645000in}{2.137125in}}%
\pgfpathlineto{\pgfqpoint{0.679586in}{2.132680in}}%
\pgfpathlineto{\pgfqpoint{0.744946in}{2.121685in}}%
\pgfpathlineto{\pgfqpoint{0.802042in}{2.107234in}}%
\pgfpathlineto{\pgfqpoint{0.853159in}{2.087657in}}%
\pgfpathlineto{\pgfqpoint{0.899759in}{2.062260in}}%
\pgfpathlineto{\pgfqpoint{0.942836in}{2.033803in}}%
\pgfpathlineto{\pgfqpoint{0.983096in}{1.999594in}}%
\pgfpathlineto{\pgfqpoint{1.021057in}{1.959996in}}%
\pgfpathlineto{\pgfqpoint{1.057113in}{1.916026in}}%
\pgfpathlineto{\pgfqpoint{1.091568in}{1.868253in}}%
\pgfpathlineto{\pgfqpoint{1.124662in}{1.815047in}}%
\pgfpathlineto{\pgfqpoint{1.156588in}{1.754693in}}%
\pgfpathlineto{\pgfqpoint{1.187505in}{1.687034in}}%
\pgfpathlineto{\pgfqpoint{1.217543in}{1.613236in}}%
\pgfpathlineto{\pgfqpoint{1.303382in}{1.388374in}}%
\pgfpathlineto{\pgfqpoint{1.330833in}{1.323794in}}%
\pgfpathlineto{\pgfqpoint{1.357807in}{1.269892in}}%
\pgfpathlineto{\pgfqpoint{1.384354in}{1.228468in}}%
\pgfpathlineto{\pgfqpoint{1.410517in}{1.199912in}}%
\pgfpathlineto{\pgfqpoint{1.436334in}{1.184100in}}%
\pgfpathlineto{\pgfqpoint{1.461840in}{1.180504in}}%
\pgfpathlineto{\pgfqpoint{1.487063in}{1.188014in}}%
\pgfpathlineto{\pgfqpoint{1.512030in}{1.204458in}}%
\pgfpathlineto{\pgfqpoint{1.536766in}{1.227355in}}%
\pgfpathlineto{\pgfqpoint{1.561290in}{1.253676in}}%
\pgfpathlineto{\pgfqpoint{1.609778in}{1.307328in}}%
\pgfpathlineto{\pgfqpoint{1.633774in}{1.331321in}}%
\pgfpathlineto{\pgfqpoint{1.657624in}{1.352065in}}%
\pgfpathlineto{\pgfqpoint{1.681340in}{1.369336in}}%
\pgfpathlineto{\pgfqpoint{1.704933in}{1.383358in}}%
\pgfpathlineto{\pgfqpoint{1.728413in}{1.394313in}}%
\pgfpathlineto{\pgfqpoint{1.751791in}{1.402661in}}%
\pgfpathlineto{\pgfqpoint{1.775073in}{1.408784in}}%
\pgfpathlineto{\pgfqpoint{1.798269in}{1.413137in}}%
\pgfpathlineto{\pgfqpoint{1.821384in}{1.416128in}}%
\pgfpathlineto{\pgfqpoint{1.844426in}{1.418086in}}%
\pgfpathlineto{\pgfqpoint{1.890311in}{1.419924in}}%
\pgfpathlineto{\pgfqpoint{1.935966in}{1.420060in}}%
\pgfpathlineto{\pgfqpoint{2.184260in}{1.417337in}}%
\pgfpathlineto{\pgfqpoint{2.487779in}{1.416952in}}%
\pgfpathlineto{\pgfqpoint{2.663589in}{1.416936in}}%
\pgfpathlineto{\pgfqpoint{3.043110in}{1.416975in}}%
\pgfpathlineto{\pgfqpoint{3.043110in}{1.416975in}}%
\pgfusepath{stroke}%
\end{pgfscope}%
\begin{pgfscope}%
\pgfpathrectangle{\pgfqpoint{0.650000in}{0.420000in}}{\pgfqpoint{2.388110in}{1.994488in}}%
\pgfusepath{clip}%
\pgfsetbuttcap%
\pgfsetroundjoin%
\pgfsetlinewidth{1.003750pt}%
\definecolor{currentstroke}{rgb}{0.454902,0.678431,0.819608}%
\pgfsetstrokecolor{currentstroke}%
\pgfsetdash{{1.000000pt}{5.000000pt}}{0.000000pt}%
\pgfpathmoveto{\pgfqpoint{0.645000in}{1.312645in}}%
\pgfpathlineto{\pgfqpoint{0.679586in}{1.300295in}}%
\pgfpathlineto{\pgfqpoint{0.744946in}{1.313368in}}%
\pgfpathlineto{\pgfqpoint{0.802042in}{1.332912in}}%
\pgfpathlineto{\pgfqpoint{0.853159in}{1.354905in}}%
\pgfpathlineto{\pgfqpoint{0.983096in}{1.407827in}}%
\pgfpathlineto{\pgfqpoint{1.021057in}{1.417801in}}%
\pgfpathlineto{\pgfqpoint{1.057113in}{1.426563in}}%
\pgfpathlineto{\pgfqpoint{1.124662in}{1.438764in}}%
\pgfpathlineto{\pgfqpoint{1.187505in}{1.452475in}}%
\pgfpathlineto{\pgfqpoint{1.217543in}{1.457676in}}%
\pgfpathlineto{\pgfqpoint{1.246809in}{1.454803in}}%
\pgfpathlineto{\pgfqpoint{1.275397in}{1.450348in}}%
\pgfpathlineto{\pgfqpoint{1.303382in}{1.443078in}}%
\pgfpathlineto{\pgfqpoint{1.330833in}{1.458123in}}%
\pgfpathlineto{\pgfqpoint{1.357807in}{1.451269in}}%
\pgfpathlineto{\pgfqpoint{1.410517in}{1.456891in}}%
\pgfpathlineto{\pgfqpoint{1.436334in}{1.451123in}}%
\pgfpathlineto{\pgfqpoint{1.461840in}{1.450942in}}%
\pgfpathlineto{\pgfqpoint{1.487063in}{1.447184in}}%
\pgfpathlineto{\pgfqpoint{1.512030in}{1.444141in}}%
\pgfpathlineto{\pgfqpoint{1.536766in}{1.412790in}}%
\pgfpathlineto{\pgfqpoint{1.561290in}{1.418283in}}%
\pgfpathlineto{\pgfqpoint{1.585622in}{1.405965in}}%
\pgfpathlineto{\pgfqpoint{1.609778in}{1.422918in}}%
\pgfpathlineto{\pgfqpoint{1.633774in}{1.449954in}}%
\pgfpathlineto{\pgfqpoint{1.657624in}{1.392295in}}%
\pgfpathlineto{\pgfqpoint{1.681340in}{1.414538in}}%
\pgfpathlineto{\pgfqpoint{1.704933in}{1.401371in}}%
\pgfpathlineto{\pgfqpoint{1.728413in}{1.374819in}}%
\pgfpathlineto{\pgfqpoint{1.751791in}{1.414281in}}%
\pgfpathlineto{\pgfqpoint{1.775073in}{1.417282in}}%
\pgfpathlineto{\pgfqpoint{1.798269in}{1.428360in}}%
\pgfpathlineto{\pgfqpoint{1.821384in}{1.404166in}}%
\pgfpathlineto{\pgfqpoint{1.844426in}{1.410276in}}%
\pgfpathlineto{\pgfqpoint{1.867400in}{1.418468in}}%
\pgfpathlineto{\pgfqpoint{1.890311in}{1.418154in}}%
\pgfpathlineto{\pgfqpoint{1.913165in}{1.419237in}}%
\pgfpathlineto{\pgfqpoint{1.935966in}{1.417968in}}%
\pgfpathlineto{\pgfqpoint{1.958719in}{1.422545in}}%
\pgfpathlineto{\pgfqpoint{1.981426in}{1.392361in}}%
\pgfpathlineto{\pgfqpoint{2.004092in}{1.415577in}}%
\pgfpathlineto{\pgfqpoint{2.026719in}{1.435443in}}%
\pgfpathlineto{\pgfqpoint{2.049312in}{1.416683in}}%
\pgfpathlineto{\pgfqpoint{2.071871in}{1.415502in}}%
\pgfpathlineto{\pgfqpoint{2.094401in}{1.412311in}}%
\pgfpathlineto{\pgfqpoint{2.116902in}{1.530265in}}%
\pgfpathlineto{\pgfqpoint{2.139378in}{1.412474in}}%
\pgfpathlineto{\pgfqpoint{2.161830in}{1.411672in}}%
\pgfpathlineto{\pgfqpoint{2.184260in}{1.415161in}}%
\pgfpathlineto{\pgfqpoint{2.205858in}{1.412743in}}%
\pgfpathlineto{\pgfqpoint{2.225949in}{1.473843in}}%
\pgfpathlineto{\pgfqpoint{2.244728in}{1.417863in}}%
\pgfpathlineto{\pgfqpoint{2.262357in}{1.336468in}}%
\pgfpathlineto{\pgfqpoint{2.278969in}{1.429553in}}%
\pgfpathlineto{\pgfqpoint{2.294675in}{1.465113in}}%
\pgfpathlineto{\pgfqpoint{2.309568in}{1.400709in}}%
\pgfpathlineto{\pgfqpoint{2.323728in}{1.413372in}}%
\pgfpathlineto{\pgfqpoint{2.337224in}{1.414869in}}%
\pgfpathlineto{\pgfqpoint{2.350117in}{1.411439in}}%
\pgfpathlineto{\pgfqpoint{2.362456in}{1.405898in}}%
\pgfpathlineto{\pgfqpoint{2.374288in}{1.404749in}}%
\pgfpathlineto{\pgfqpoint{2.385654in}{1.397989in}}%
\pgfpathlineto{\pgfqpoint{2.396587in}{1.406433in}}%
\pgfpathlineto{\pgfqpoint{2.407121in}{1.410401in}}%
\pgfpathlineto{\pgfqpoint{2.417283in}{1.430034in}}%
\pgfpathlineto{\pgfqpoint{2.427098in}{1.403026in}}%
\pgfpathlineto{\pgfqpoint{2.436590in}{1.412887in}}%
\pgfpathlineto{\pgfqpoint{2.445778in}{1.416108in}}%
\pgfpathlineto{\pgfqpoint{2.454683in}{1.414704in}}%
\pgfpathlineto{\pgfqpoint{2.463320in}{1.404887in}}%
\pgfpathlineto{\pgfqpoint{2.471706in}{1.407773in}}%
\pgfpathlineto{\pgfqpoint{2.479855in}{1.426968in}}%
\pgfpathlineto{\pgfqpoint{2.487779in}{1.410262in}}%
\pgfpathlineto{\pgfqpoint{2.495491in}{1.406963in}}%
\pgfpathlineto{\pgfqpoint{2.503002in}{1.416145in}}%
\pgfpathlineto{\pgfqpoint{2.510322in}{1.410440in}}%
\pgfpathlineto{\pgfqpoint{2.517460in}{1.411408in}}%
\pgfpathlineto{\pgfqpoint{2.531227in}{1.415211in}}%
\pgfpathlineto{\pgfqpoint{2.537871in}{1.411247in}}%
\pgfpathlineto{\pgfqpoint{2.544365in}{1.414132in}}%
\pgfpathlineto{\pgfqpoint{2.550716in}{1.415774in}}%
\pgfpathlineto{\pgfqpoint{2.556930in}{1.411744in}}%
\pgfpathlineto{\pgfqpoint{2.563013in}{1.414478in}}%
\pgfpathlineto{\pgfqpoint{2.568970in}{1.410933in}}%
\pgfpathlineto{\pgfqpoint{2.574806in}{1.414879in}}%
\pgfpathlineto{\pgfqpoint{2.580526in}{1.414870in}}%
\pgfpathlineto{\pgfqpoint{2.586135in}{1.416368in}}%
\pgfpathlineto{\pgfqpoint{2.591637in}{1.413330in}}%
\pgfpathlineto{\pgfqpoint{2.597035in}{1.419332in}}%
\pgfpathlineto{\pgfqpoint{2.602334in}{1.412213in}}%
\pgfpathlineto{\pgfqpoint{2.607537in}{1.421521in}}%
\pgfpathlineto{\pgfqpoint{2.612648in}{1.408150in}}%
\pgfpathlineto{\pgfqpoint{2.617670in}{1.427857in}}%
\pgfpathlineto{\pgfqpoint{2.622606in}{1.411074in}}%
\pgfpathlineto{\pgfqpoint{2.627458in}{1.419461in}}%
\pgfpathlineto{\pgfqpoint{2.632231in}{1.414817in}}%
\pgfpathlineto{\pgfqpoint{2.636925in}{1.417645in}}%
\pgfpathlineto{\pgfqpoint{2.641544in}{1.433139in}}%
\pgfpathlineto{\pgfqpoint{2.646090in}{1.417713in}}%
\pgfpathlineto{\pgfqpoint{2.650565in}{1.416259in}}%
\pgfpathlineto{\pgfqpoint{2.654972in}{1.417394in}}%
\pgfpathlineto{\pgfqpoint{2.659313in}{1.413778in}}%
\pgfpathlineto{\pgfqpoint{2.663589in}{1.418066in}}%
\pgfpathlineto{\pgfqpoint{2.667802in}{1.371146in}}%
\pgfpathlineto{\pgfqpoint{2.671955in}{1.412407in}}%
\pgfpathlineto{\pgfqpoint{2.676048in}{1.418275in}}%
\pgfpathlineto{\pgfqpoint{2.680085in}{1.415414in}}%
\pgfpathlineto{\pgfqpoint{2.684065in}{1.415651in}}%
\pgfpathlineto{\pgfqpoint{2.687991in}{1.414660in}}%
\pgfpathlineto{\pgfqpoint{2.691864in}{1.418921in}}%
\pgfpathlineto{\pgfqpoint{2.695686in}{1.417878in}}%
\pgfpathlineto{\pgfqpoint{2.699458in}{1.415242in}}%
\pgfpathlineto{\pgfqpoint{2.703181in}{1.417601in}}%
\pgfpathlineto{\pgfqpoint{2.706857in}{1.416545in}}%
\pgfpathlineto{\pgfqpoint{2.710486in}{1.417468in}}%
\pgfpathlineto{\pgfqpoint{2.714070in}{1.416358in}}%
\pgfpathlineto{\pgfqpoint{2.717610in}{1.416708in}}%
\pgfpathlineto{\pgfqpoint{2.721107in}{1.422752in}}%
\pgfpathlineto{\pgfqpoint{2.724562in}{1.414920in}}%
\pgfpathlineto{\pgfqpoint{2.727976in}{1.419035in}}%
\pgfpathlineto{\pgfqpoint{2.731350in}{1.415928in}}%
\pgfpathlineto{\pgfqpoint{2.734685in}{1.415309in}}%
\pgfpathlineto{\pgfqpoint{2.737982in}{1.421410in}}%
\pgfpathlineto{\pgfqpoint{2.741242in}{1.416397in}}%
\pgfpathlineto{\pgfqpoint{2.744465in}{1.415286in}}%
\pgfpathlineto{\pgfqpoint{2.747652in}{1.417688in}}%
\pgfpathlineto{\pgfqpoint{2.750804in}{1.414624in}}%
\pgfpathlineto{\pgfqpoint{2.753923in}{1.420341in}}%
\pgfpathlineto{\pgfqpoint{2.757007in}{1.431848in}}%
\pgfpathlineto{\pgfqpoint{2.760060in}{1.411923in}}%
\pgfpathlineto{\pgfqpoint{2.763080in}{1.287021in}}%
\pgfpathlineto{\pgfqpoint{2.766068in}{1.413227in}}%
\pgfpathlineto{\pgfqpoint{2.769027in}{1.382485in}}%
\pgfpathlineto{\pgfqpoint{2.774853in}{1.430183in}}%
\pgfpathlineto{\pgfqpoint{2.777723in}{1.427037in}}%
\pgfpathlineto{\pgfqpoint{2.780564in}{1.428985in}}%
\pgfpathlineto{\pgfqpoint{2.783378in}{1.402186in}}%
\pgfpathlineto{\pgfqpoint{2.786164in}{1.427287in}}%
\pgfpathlineto{\pgfqpoint{2.788924in}{1.418998in}}%
\pgfpathlineto{\pgfqpoint{2.791657in}{1.439537in}}%
\pgfpathlineto{\pgfqpoint{2.794365in}{1.400104in}}%
\pgfpathlineto{\pgfqpoint{2.797047in}{1.416759in}}%
\pgfpathlineto{\pgfqpoint{2.799705in}{1.395624in}}%
\pgfpathlineto{\pgfqpoint{2.802338in}{1.387881in}}%
\pgfpathlineto{\pgfqpoint{2.804948in}{1.437947in}}%
\pgfpathlineto{\pgfqpoint{2.807534in}{1.421623in}}%
\pgfpathlineto{\pgfqpoint{2.810097in}{1.412794in}}%
\pgfpathlineto{\pgfqpoint{2.812638in}{1.407247in}}%
\pgfpathlineto{\pgfqpoint{2.815156in}{1.429026in}}%
\pgfpathlineto{\pgfqpoint{2.817652in}{1.430534in}}%
\pgfpathlineto{\pgfqpoint{2.820127in}{1.420080in}}%
\pgfpathlineto{\pgfqpoint{2.822581in}{1.406533in}}%
\pgfpathlineto{\pgfqpoint{2.825014in}{1.428518in}}%
\pgfpathlineto{\pgfqpoint{2.827427in}{1.408092in}}%
\pgfpathlineto{\pgfqpoint{2.829820in}{1.409998in}}%
\pgfpathlineto{\pgfqpoint{2.832193in}{1.418736in}}%
\pgfpathlineto{\pgfqpoint{2.836881in}{1.401613in}}%
\pgfpathlineto{\pgfqpoint{2.839197in}{1.414195in}}%
\pgfpathlineto{\pgfqpoint{2.841494in}{1.411717in}}%
\pgfpathlineto{\pgfqpoint{2.843773in}{1.417703in}}%
\pgfpathlineto{\pgfqpoint{2.846034in}{1.426283in}}%
\pgfpathlineto{\pgfqpoint{2.848278in}{1.415210in}}%
\pgfpathlineto{\pgfqpoint{2.850504in}{1.419389in}}%
\pgfpathlineto{\pgfqpoint{2.852713in}{1.425200in}}%
\pgfpathlineto{\pgfqpoint{2.854905in}{1.407396in}}%
\pgfpathlineto{\pgfqpoint{2.857081in}{1.424355in}}%
\pgfpathlineto{\pgfqpoint{2.859240in}{1.413751in}}%
\pgfpathlineto{\pgfqpoint{2.861384in}{1.406579in}}%
\pgfpathlineto{\pgfqpoint{2.863511in}{1.413694in}}%
\pgfpathlineto{\pgfqpoint{2.865623in}{1.414026in}}%
\pgfpathlineto{\pgfqpoint{2.867720in}{1.417581in}}%
\pgfpathlineto{\pgfqpoint{2.869801in}{1.416764in}}%
\pgfpathlineto{\pgfqpoint{2.871867in}{1.410239in}}%
\pgfpathlineto{\pgfqpoint{2.873919in}{1.423223in}}%
\pgfpathlineto{\pgfqpoint{2.875956in}{1.417357in}}%
\pgfpathlineto{\pgfqpoint{2.877979in}{1.415201in}}%
\pgfpathlineto{\pgfqpoint{2.879988in}{1.416696in}}%
\pgfpathlineto{\pgfqpoint{2.881983in}{1.415701in}}%
\pgfpathlineto{\pgfqpoint{2.883964in}{1.415950in}}%
\pgfpathlineto{\pgfqpoint{2.885932in}{1.425166in}}%
\pgfpathlineto{\pgfqpoint{2.887886in}{1.413922in}}%
\pgfpathlineto{\pgfqpoint{2.889827in}{1.408745in}}%
\pgfpathlineto{\pgfqpoint{2.891755in}{1.427633in}}%
\pgfpathlineto{\pgfqpoint{2.895573in}{1.414061in}}%
\pgfpathlineto{\pgfqpoint{2.897463in}{1.417528in}}%
\pgfpathlineto{\pgfqpoint{2.899341in}{1.414529in}}%
\pgfpathlineto{\pgfqpoint{2.901206in}{1.415933in}}%
\pgfpathlineto{\pgfqpoint{2.903060in}{1.416458in}}%
\pgfpathlineto{\pgfqpoint{2.904901in}{1.418351in}}%
\pgfpathlineto{\pgfqpoint{2.906731in}{1.407983in}}%
\pgfpathlineto{\pgfqpoint{2.908550in}{1.409431in}}%
\pgfpathlineto{\pgfqpoint{2.910357in}{1.417180in}}%
\pgfpathlineto{\pgfqpoint{2.912153in}{1.417667in}}%
\pgfpathlineto{\pgfqpoint{2.913937in}{1.422662in}}%
\pgfpathlineto{\pgfqpoint{2.915711in}{1.410340in}}%
\pgfpathlineto{\pgfqpoint{2.917474in}{1.431256in}}%
\pgfpathlineto{\pgfqpoint{2.919226in}{1.411584in}}%
\pgfpathlineto{\pgfqpoint{2.920967in}{1.419233in}}%
\pgfpathlineto{\pgfqpoint{2.922698in}{1.415507in}}%
\pgfpathlineto{\pgfqpoint{2.924419in}{1.415746in}}%
\pgfpathlineto{\pgfqpoint{2.926130in}{1.418698in}}%
\pgfpathlineto{\pgfqpoint{2.927830in}{1.414664in}}%
\pgfpathlineto{\pgfqpoint{2.929520in}{1.419849in}}%
\pgfpathlineto{\pgfqpoint{2.931201in}{1.410704in}}%
\pgfpathlineto{\pgfqpoint{2.932871in}{1.417323in}}%
\pgfpathlineto{\pgfqpoint{2.934533in}{1.417218in}}%
\pgfpathlineto{\pgfqpoint{2.936184in}{1.425010in}}%
\pgfpathlineto{\pgfqpoint{2.939459in}{1.419062in}}%
\pgfpathlineto{\pgfqpoint{2.941083in}{1.422243in}}%
\pgfpathlineto{\pgfqpoint{2.944303in}{1.415343in}}%
\pgfpathlineto{\pgfqpoint{2.945899in}{1.415622in}}%
\pgfpathlineto{\pgfqpoint{2.947487in}{1.417200in}}%
\pgfpathlineto{\pgfqpoint{2.949066in}{1.417715in}}%
\pgfpathlineto{\pgfqpoint{2.950637in}{1.417339in}}%
\pgfpathlineto{\pgfqpoint{2.952199in}{1.415012in}}%
\pgfpathlineto{\pgfqpoint{2.953752in}{1.417515in}}%
\pgfpathlineto{\pgfqpoint{2.955298in}{1.418499in}}%
\pgfpathlineto{\pgfqpoint{2.956835in}{1.418484in}}%
\pgfpathlineto{\pgfqpoint{2.958363in}{1.423269in}}%
\pgfpathlineto{\pgfqpoint{2.959884in}{1.413796in}}%
\pgfpathlineto{\pgfqpoint{2.961397in}{1.417976in}}%
\pgfpathlineto{\pgfqpoint{2.962902in}{1.418800in}}%
\pgfpathlineto{\pgfqpoint{2.964399in}{1.415395in}}%
\pgfpathlineto{\pgfqpoint{2.965888in}{1.417989in}}%
\pgfpathlineto{\pgfqpoint{2.967370in}{1.415652in}}%
\pgfpathlineto{\pgfqpoint{2.968844in}{1.416651in}}%
\pgfpathlineto{\pgfqpoint{2.970310in}{1.416391in}}%
\pgfpathlineto{\pgfqpoint{2.971769in}{1.418276in}}%
\pgfpathlineto{\pgfqpoint{2.973221in}{1.417005in}}%
\pgfpathlineto{\pgfqpoint{2.974665in}{1.417125in}}%
\pgfpathlineto{\pgfqpoint{2.976103in}{1.416162in}}%
\pgfpathlineto{\pgfqpoint{2.977533in}{1.423563in}}%
\pgfpathlineto{\pgfqpoint{2.978956in}{1.421209in}}%
\pgfpathlineto{\pgfqpoint{2.981781in}{1.414110in}}%
\pgfpathlineto{\pgfqpoint{2.983183in}{1.416258in}}%
\pgfpathlineto{\pgfqpoint{2.984578in}{1.412276in}}%
\pgfpathlineto{\pgfqpoint{2.985967in}{1.415993in}}%
\pgfpathlineto{\pgfqpoint{2.987349in}{1.421791in}}%
\pgfpathlineto{\pgfqpoint{2.988725in}{1.416947in}}%
\pgfpathlineto{\pgfqpoint{2.990094in}{1.416417in}}%
\pgfpathlineto{\pgfqpoint{2.991456in}{1.414024in}}%
\pgfpathlineto{\pgfqpoint{2.992812in}{1.420320in}}%
\pgfpathlineto{\pgfqpoint{2.994162in}{1.416202in}}%
\pgfpathlineto{\pgfqpoint{2.995505in}{1.416876in}}%
\pgfpathlineto{\pgfqpoint{2.996842in}{1.411837in}}%
\pgfpathlineto{\pgfqpoint{2.998173in}{1.432781in}}%
\pgfpathlineto{\pgfqpoint{2.999498in}{1.419088in}}%
\pgfpathlineto{\pgfqpoint{3.000817in}{1.414263in}}%
\pgfpathlineto{\pgfqpoint{3.002129in}{1.417557in}}%
\pgfpathlineto{\pgfqpoint{3.003436in}{1.412896in}}%
\pgfpathlineto{\pgfqpoint{3.004737in}{1.419563in}}%
\pgfpathlineto{\pgfqpoint{3.006032in}{1.411730in}}%
\pgfpathlineto{\pgfqpoint{3.007321in}{1.409701in}}%
\pgfpathlineto{\pgfqpoint{3.008605in}{1.411119in}}%
\pgfpathlineto{\pgfqpoint{3.009882in}{1.422242in}}%
\pgfpathlineto{\pgfqpoint{3.011155in}{1.417054in}}%
\pgfpathlineto{\pgfqpoint{3.012421in}{1.417596in}}%
\pgfpathlineto{\pgfqpoint{3.013682in}{1.415812in}}%
\pgfpathlineto{\pgfqpoint{3.016188in}{1.416203in}}%
\pgfpathlineto{\pgfqpoint{3.018672in}{1.419654in}}%
\pgfpathlineto{\pgfqpoint{3.022358in}{1.414857in}}%
\pgfpathlineto{\pgfqpoint{3.024789in}{1.417584in}}%
\pgfpathlineto{\pgfqpoint{3.025997in}{1.416891in}}%
\pgfpathlineto{\pgfqpoint{3.036648in}{1.417348in}}%
\pgfpathlineto{\pgfqpoint{3.037807in}{1.417444in}}%
\pgfpathlineto{\pgfqpoint{3.038962in}{1.416587in}}%
\pgfpathlineto{\pgfqpoint{3.040112in}{1.417822in}}%
\pgfpathlineto{\pgfqpoint{3.041258in}{1.414336in}}%
\pgfpathlineto{\pgfqpoint{3.042399in}{1.417790in}}%
\pgfpathlineto{\pgfqpoint{3.043110in}{1.416709in}}%
\pgfpathlineto{\pgfqpoint{3.043110in}{1.416709in}}%
\pgfusepath{stroke}%
\end{pgfscope}%
\begin{pgfscope}%
\pgfpathrectangle{\pgfqpoint{0.650000in}{0.420000in}}{\pgfqpoint{2.388110in}{1.994488in}}%
\pgfusepath{clip}%
\pgfsetrectcap%
\pgfsetroundjoin%
\pgfsetlinewidth{1.003750pt}%
\definecolor{currentstroke}{rgb}{0.501961,0.450980,0.674510}%
\pgfsetstrokecolor{currentstroke}%
\pgfsetdash{}{0pt}%
\pgfpathmoveto{\pgfqpoint{0.645000in}{1.417492in}}%
\pgfpathlineto{\pgfqpoint{0.766279in}{1.418672in}}%
\pgfpathlineto{\pgfqpoint{0.830619in}{1.420309in}}%
\pgfpathlineto{\pgfqpoint{0.870161in}{1.422050in}}%
\pgfpathlineto{\pgfqpoint{0.907582in}{1.424535in}}%
\pgfpathlineto{\pgfqpoint{0.943235in}{1.428018in}}%
\pgfpathlineto{\pgfqpoint{0.977396in}{1.432822in}}%
\pgfpathlineto{\pgfqpoint{1.010285in}{1.439351in}}%
\pgfpathlineto{\pgfqpoint{1.026310in}{1.443413in}}%
\pgfpathlineto{\pgfqpoint{1.042080in}{1.448101in}}%
\pgfpathlineto{\pgfqpoint{1.057613in}{1.453491in}}%
\pgfpathlineto{\pgfqpoint{1.072926in}{1.459668in}}%
\pgfpathlineto{\pgfqpoint{1.088031in}{1.466722in}}%
\pgfpathlineto{\pgfqpoint{1.102943in}{1.474751in}}%
\pgfpathlineto{\pgfqpoint{1.117673in}{1.483856in}}%
\pgfpathlineto{\pgfqpoint{1.132233in}{1.494144in}}%
\pgfpathlineto{\pgfqpoint{1.146632in}{1.505726in}}%
\pgfpathlineto{\pgfqpoint{1.160881in}{1.518714in}}%
\pgfpathlineto{\pgfqpoint{1.174988in}{1.533217in}}%
\pgfpathlineto{\pgfqpoint{1.188960in}{1.549346in}}%
\pgfpathlineto{\pgfqpoint{1.202807in}{1.567203in}}%
\pgfpathlineto{\pgfqpoint{1.216534in}{1.586878in}}%
\pgfpathlineto{\pgfqpoint{1.230148in}{1.608450in}}%
\pgfpathlineto{\pgfqpoint{1.243655in}{1.631976in}}%
\pgfpathlineto{\pgfqpoint{1.257062in}{1.657492in}}%
\pgfpathlineto{\pgfqpoint{1.270372in}{1.684998in}}%
\pgfpathlineto{\pgfqpoint{1.283592in}{1.714463in}}%
\pgfpathlineto{\pgfqpoint{1.296725in}{1.745810in}}%
\pgfpathlineto{\pgfqpoint{1.309777in}{1.778914in}}%
\pgfpathlineto{\pgfqpoint{1.335650in}{1.849631in}}%
\pgfpathlineto{\pgfqpoint{1.361243in}{1.924506in}}%
\pgfpathlineto{\pgfqpoint{1.411689in}{2.074935in}}%
\pgfpathlineto{\pgfqpoint{1.436586in}{2.143811in}}%
\pgfpathlineto{\pgfqpoint{1.448963in}{2.175092in}}%
\pgfpathlineto{\pgfqpoint{1.461293in}{2.203632in}}%
\pgfpathlineto{\pgfqpoint{1.473581in}{2.228954in}}%
\pgfpathlineto{\pgfqpoint{1.485826in}{2.250647in}}%
\pgfpathlineto{\pgfqpoint{1.498032in}{2.268379in}}%
\pgfpathlineto{\pgfqpoint{1.510201in}{2.281931in}}%
\pgfpathlineto{\pgfqpoint{1.522333in}{2.291206in}}%
\pgfpathlineto{\pgfqpoint{1.534430in}{2.296226in}}%
\pgfpathlineto{\pgfqpoint{1.546495in}{2.297102in}}%
\pgfpathlineto{\pgfqpoint{1.558528in}{2.293949in}}%
\pgfpathlineto{\pgfqpoint{1.570530in}{2.286826in}}%
\pgfpathlineto{\pgfqpoint{1.582504in}{2.275805in}}%
\pgfpathlineto{\pgfqpoint{1.594450in}{2.261115in}}%
\pgfpathlineto{\pgfqpoint{1.606369in}{2.243211in}}%
\pgfpathlineto{\pgfqpoint{1.618263in}{2.222607in}}%
\pgfpathlineto{\pgfqpoint{1.630133in}{2.199700in}}%
\pgfpathlineto{\pgfqpoint{1.641980in}{2.174763in}}%
\pgfpathlineto{\pgfqpoint{1.653804in}{2.148055in}}%
\pgfpathlineto{\pgfqpoint{1.677389in}{2.090821in}}%
\pgfpathlineto{\pgfqpoint{1.747699in}{1.915080in}}%
\pgfpathlineto{\pgfqpoint{1.771006in}{1.861408in}}%
\pgfpathlineto{\pgfqpoint{1.794257in}{1.812067in}}%
\pgfpathlineto{\pgfqpoint{1.817457in}{1.767596in}}%
\pgfpathlineto{\pgfqpoint{1.829039in}{1.747194in}}%
\pgfpathlineto{\pgfqpoint{1.840609in}{1.727996in}}%
\pgfpathlineto{\pgfqpoint{1.852169in}{1.709958in}}%
\pgfpathlineto{\pgfqpoint{1.863718in}{1.693034in}}%
\pgfpathlineto{\pgfqpoint{1.875257in}{1.677175in}}%
\pgfpathlineto{\pgfqpoint{1.886786in}{1.662373in}}%
\pgfpathlineto{\pgfqpoint{1.898306in}{1.648643in}}%
\pgfpathlineto{\pgfqpoint{1.909817in}{1.635961in}}%
\pgfpathlineto{\pgfqpoint{1.921320in}{1.624244in}}%
\pgfpathlineto{\pgfqpoint{1.932814in}{1.613408in}}%
\pgfpathlineto{\pgfqpoint{1.944301in}{1.603385in}}%
\pgfpathlineto{\pgfqpoint{1.967251in}{1.585521in}}%
\pgfpathlineto{\pgfqpoint{1.990174in}{1.570191in}}%
\pgfpathlineto{\pgfqpoint{2.013070in}{1.556932in}}%
\pgfpathlineto{\pgfqpoint{2.035944in}{1.545453in}}%
\pgfpathlineto{\pgfqpoint{2.058795in}{1.535446in}}%
\pgfpathlineto{\pgfqpoint{2.081626in}{1.526688in}}%
\pgfpathlineto{\pgfqpoint{2.104438in}{1.519089in}}%
\pgfpathlineto{\pgfqpoint{2.127233in}{1.512704in}}%
\pgfpathlineto{\pgfqpoint{2.172776in}{1.501891in}}%
\pgfpathlineto{\pgfqpoint{2.218265in}{1.491801in}}%
\pgfpathlineto{\pgfqpoint{2.263686in}{1.482884in}}%
\pgfpathlineto{\pgfqpoint{2.286134in}{1.479224in}}%
\pgfpathlineto{\pgfqpoint{2.327947in}{1.474159in}}%
\pgfpathlineto{\pgfqpoint{2.381341in}{1.467133in}}%
\pgfpathlineto{\pgfqpoint{2.459159in}{1.455923in}}%
\pgfpathlineto{\pgfqpoint{2.488596in}{1.452784in}}%
\pgfpathlineto{\pgfqpoint{2.510155in}{1.450325in}}%
\pgfpathlineto{\pgfqpoint{2.525326in}{1.448319in}}%
\pgfpathlineto{\pgfqpoint{2.562262in}{1.441585in}}%
\pgfpathlineto{\pgfqpoint{2.647846in}{1.429902in}}%
\pgfpathlineto{\pgfqpoint{2.684093in}{1.424391in}}%
\pgfpathlineto{\pgfqpoint{2.718810in}{1.420047in}}%
\pgfpathlineto{\pgfqpoint{2.747520in}{1.417881in}}%
\pgfpathlineto{\pgfqpoint{2.773625in}{1.417129in}}%
\pgfpathlineto{\pgfqpoint{2.835180in}{1.416975in}}%
\pgfpathlineto{\pgfqpoint{3.043110in}{1.416975in}}%
\pgfpathlineto{\pgfqpoint{3.043110in}{1.416975in}}%
\pgfusepath{stroke}%
\end{pgfscope}%
\begin{pgfscope}%
\pgfpathrectangle{\pgfqpoint{0.650000in}{0.420000in}}{\pgfqpoint{2.388110in}{1.994488in}}%
\pgfusepath{clip}%
\pgfsetbuttcap%
\pgfsetroundjoin%
\pgfsetlinewidth{1.003750pt}%
\definecolor{currentstroke}{rgb}{0.501961,0.450980,0.674510}%
\pgfsetstrokecolor{currentstroke}%
\pgfsetdash{{5.000000pt}{1.000000pt}}{0.000000pt}%
\pgfpathmoveto{\pgfqpoint{0.645000in}{1.416973in}}%
\pgfpathlineto{\pgfqpoint{1.202807in}{1.416392in}}%
\pgfpathlineto{\pgfqpoint{1.361243in}{1.414848in}}%
\pgfpathlineto{\pgfqpoint{1.558528in}{1.412905in}}%
\pgfpathlineto{\pgfqpoint{1.700896in}{1.412304in}}%
\pgfpathlineto{\pgfqpoint{1.759360in}{1.411377in}}%
\pgfpathlineto{\pgfqpoint{1.817457in}{1.411919in}}%
\pgfpathlineto{\pgfqpoint{1.921320in}{1.414040in}}%
\pgfpathlineto{\pgfqpoint{1.978716in}{1.412061in}}%
\pgfpathlineto{\pgfqpoint{2.001625in}{1.414236in}}%
\pgfpathlineto{\pgfqpoint{2.035944in}{1.417809in}}%
\pgfpathlineto{\pgfqpoint{2.081626in}{1.420673in}}%
\pgfpathlineto{\pgfqpoint{2.104438in}{1.420314in}}%
\pgfpathlineto{\pgfqpoint{2.161396in}{1.418028in}}%
\pgfpathlineto{\pgfqpoint{2.184153in}{1.418254in}}%
\pgfpathlineto{\pgfqpoint{2.218265in}{1.420095in}}%
\pgfpathlineto{\pgfqpoint{2.286134in}{1.426044in}}%
\pgfpathlineto{\pgfqpoint{2.307710in}{1.426915in}}%
\pgfpathlineto{\pgfqpoint{2.318002in}{1.426672in}}%
\pgfpathlineto{\pgfqpoint{2.327947in}{1.425683in}}%
\pgfpathlineto{\pgfqpoint{2.355880in}{1.421159in}}%
\pgfpathlineto{\pgfqpoint{2.373102in}{1.420289in}}%
\pgfpathlineto{\pgfqpoint{2.389351in}{1.420770in}}%
\pgfpathlineto{\pgfqpoint{2.412125in}{1.422143in}}%
\pgfpathlineto{\pgfqpoint{2.433230in}{1.422252in}}%
\pgfpathlineto{\pgfqpoint{2.459159in}{1.423882in}}%
\pgfpathlineto{\pgfqpoint{2.471297in}{1.422698in}}%
\pgfpathlineto{\pgfqpoint{2.488596in}{1.420510in}}%
\pgfpathlineto{\pgfqpoint{2.499577in}{1.420673in}}%
\pgfpathlineto{\pgfqpoint{2.515301in}{1.421594in}}%
\pgfpathlineto{\pgfqpoint{2.539738in}{1.421547in}}%
\pgfpathlineto{\pgfqpoint{2.557896in}{1.420368in}}%
\pgfpathlineto{\pgfqpoint{2.574977in}{1.421666in}}%
\pgfpathlineto{\pgfqpoint{2.587154in}{1.419923in}}%
\pgfpathlineto{\pgfqpoint{2.598836in}{1.418736in}}%
\pgfpathlineto{\pgfqpoint{2.624383in}{1.418145in}}%
\pgfpathlineto{\pgfqpoint{2.634673in}{1.417335in}}%
\pgfpathlineto{\pgfqpoint{2.647846in}{1.416723in}}%
\pgfpathlineto{\pgfqpoint{2.660442in}{1.415604in}}%
\pgfpathlineto{\pgfqpoint{2.721315in}{1.414656in}}%
\pgfpathlineto{\pgfqpoint{2.767317in}{1.415188in}}%
\pgfpathlineto{\pgfqpoint{2.805113in}{1.416337in}}%
\pgfpathlineto{\pgfqpoint{2.835180in}{1.416789in}}%
\pgfpathlineto{\pgfqpoint{2.853085in}{1.416880in}}%
\pgfpathlineto{\pgfqpoint{2.934950in}{1.416953in}}%
\pgfpathlineto{\pgfqpoint{2.989920in}{1.416974in}}%
\pgfpathlineto{\pgfqpoint{3.043110in}{1.416967in}}%
\pgfpathlineto{\pgfqpoint{3.043110in}{1.416967in}}%
\pgfusepath{stroke}%
\end{pgfscope}%
\begin{pgfscope}%
\pgfpathrectangle{\pgfqpoint{0.650000in}{0.420000in}}{\pgfqpoint{2.388110in}{1.994488in}}%
\pgfusepath{clip}%
\pgfsetbuttcap%
\pgfsetroundjoin%
\pgfsetlinewidth{1.003750pt}%
\definecolor{currentstroke}{rgb}{0.501961,0.450980,0.674510}%
\pgfsetstrokecolor{currentstroke}%
\pgfsetdash{{3.000000pt}{1.000000pt}{1.000000pt}{1.000000pt}}{0.000000pt}%
\pgfpathmoveto{\pgfqpoint{0.645000in}{0.593631in}}%
\pgfpathlineto{\pgfqpoint{0.667196in}{0.601047in}}%
\pgfpathlineto{\pgfqpoint{0.743169in}{0.615679in}}%
\pgfpathlineto{\pgfqpoint{0.788495in}{0.625558in}}%
\pgfpathlineto{\pgfqpoint{0.830619in}{0.636187in}}%
\pgfpathlineto{\pgfqpoint{0.870161in}{0.647748in}}%
\pgfpathlineto{\pgfqpoint{0.907582in}{0.660265in}}%
\pgfpathlineto{\pgfqpoint{0.943235in}{0.673682in}}%
\pgfpathlineto{\pgfqpoint{0.977396in}{0.687900in}}%
\pgfpathlineto{\pgfqpoint{1.010285in}{0.702799in}}%
\pgfpathlineto{\pgfqpoint{1.042080in}{0.718244in}}%
\pgfpathlineto{\pgfqpoint{1.088031in}{0.742126in}}%
\pgfpathlineto{\pgfqpoint{1.132233in}{0.766434in}}%
\pgfpathlineto{\pgfqpoint{1.283592in}{0.851134in}}%
\pgfpathlineto{\pgfqpoint{1.322750in}{0.870778in}}%
\pgfpathlineto{\pgfqpoint{1.348480in}{0.882535in}}%
\pgfpathlineto{\pgfqpoint{1.373943in}{0.893035in}}%
\pgfpathlineto{\pgfqpoint{1.399163in}{0.902115in}}%
\pgfpathlineto{\pgfqpoint{1.424163in}{0.909611in}}%
\pgfpathlineto{\pgfqpoint{1.448963in}{0.915355in}}%
\pgfpathlineto{\pgfqpoint{1.473581in}{0.919177in}}%
\pgfpathlineto{\pgfqpoint{1.498032in}{0.920907in}}%
\pgfpathlineto{\pgfqpoint{1.510201in}{0.920931in}}%
\pgfpathlineto{\pgfqpoint{1.522333in}{0.920367in}}%
\pgfpathlineto{\pgfqpoint{1.534430in}{0.919191in}}%
\pgfpathlineto{\pgfqpoint{1.546495in}{0.914181in}}%
\pgfpathlineto{\pgfqpoint{1.558528in}{0.916699in}}%
\pgfpathlineto{\pgfqpoint{1.582504in}{0.923480in}}%
\pgfpathlineto{\pgfqpoint{1.594450in}{0.927591in}}%
\pgfpathlineto{\pgfqpoint{1.606369in}{0.932418in}}%
\pgfpathlineto{\pgfqpoint{1.618263in}{0.938133in}}%
\pgfpathlineto{\pgfqpoint{1.630133in}{0.944697in}}%
\pgfpathlineto{\pgfqpoint{1.653804in}{0.959627in}}%
\pgfpathlineto{\pgfqpoint{1.677389in}{0.976381in}}%
\pgfpathlineto{\pgfqpoint{1.700896in}{0.995085in}}%
\pgfpathlineto{\pgfqpoint{1.724330in}{1.015190in}}%
\pgfpathlineto{\pgfqpoint{1.759360in}{1.047008in}}%
\pgfpathlineto{\pgfqpoint{1.817457in}{1.102016in}}%
\pgfpathlineto{\pgfqpoint{1.863718in}{1.145330in}}%
\pgfpathlineto{\pgfqpoint{1.898306in}{1.175853in}}%
\pgfpathlineto{\pgfqpoint{1.921320in}{1.194931in}}%
\pgfpathlineto{\pgfqpoint{1.944301in}{1.212761in}}%
\pgfpathlineto{\pgfqpoint{1.967251in}{1.229316in}}%
\pgfpathlineto{\pgfqpoint{1.990174in}{1.244678in}}%
\pgfpathlineto{\pgfqpoint{2.013070in}{1.258848in}}%
\pgfpathlineto{\pgfqpoint{2.035944in}{1.271739in}}%
\pgfpathlineto{\pgfqpoint{2.058795in}{1.283390in}}%
\pgfpathlineto{\pgfqpoint{2.081626in}{1.293834in}}%
\pgfpathlineto{\pgfqpoint{2.104438in}{1.303075in}}%
\pgfpathlineto{\pgfqpoint{2.138624in}{1.315147in}}%
\pgfpathlineto{\pgfqpoint{2.172776in}{1.326034in}}%
\pgfpathlineto{\pgfqpoint{2.206897in}{1.335681in}}%
\pgfpathlineto{\pgfqpoint{2.252347in}{1.346705in}}%
\pgfpathlineto{\pgfqpoint{2.286134in}{1.353953in}}%
\pgfpathlineto{\pgfqpoint{2.327947in}{1.361747in}}%
\pgfpathlineto{\pgfqpoint{2.364620in}{1.367596in}}%
\pgfpathlineto{\pgfqpoint{2.439934in}{1.378204in}}%
\pgfpathlineto{\pgfqpoint{2.494139in}{1.385348in}}%
\pgfpathlineto{\pgfqpoint{2.548960in}{1.392362in}}%
\pgfpathlineto{\pgfqpoint{2.602628in}{1.399656in}}%
\pgfpathlineto{\pgfqpoint{2.666540in}{1.407687in}}%
\pgfpathlineto{\pgfqpoint{2.703305in}{1.411692in}}%
\pgfpathlineto{\pgfqpoint{2.745234in}{1.414964in}}%
\pgfpathlineto{\pgfqpoint{2.783840in}{1.416473in}}%
\pgfpathlineto{\pgfqpoint{2.833496in}{1.416947in}}%
\pgfpathlineto{\pgfqpoint{3.043110in}{1.416975in}}%
\pgfpathlineto{\pgfqpoint{3.043110in}{1.416975in}}%
\pgfusepath{stroke}%
\end{pgfscope}%
\begin{pgfscope}%
\pgfpathrectangle{\pgfqpoint{0.650000in}{0.420000in}}{\pgfqpoint{2.388110in}{1.994488in}}%
\pgfusepath{clip}%
\pgfsetbuttcap%
\pgfsetroundjoin%
\pgfsetlinewidth{1.003750pt}%
\definecolor{currentstroke}{rgb}{0.501961,0.450980,0.674510}%
\pgfsetstrokecolor{currentstroke}%
\pgfsetdash{{1.000000pt}{1.000000pt}}{0.000000pt}%
\pgfpathmoveto{\pgfqpoint{0.645000in}{1.453611in}}%
\pgfpathlineto{\pgfqpoint{0.719048in}{1.461712in}}%
\pgfpathlineto{\pgfqpoint{0.809915in}{1.472983in}}%
\pgfpathlineto{\pgfqpoint{0.977396in}{1.494435in}}%
\pgfpathlineto{\pgfqpoint{1.026310in}{1.499645in}}%
\pgfpathlineto{\pgfqpoint{1.072926in}{1.503610in}}%
\pgfpathlineto{\pgfqpoint{1.117673in}{1.506236in}}%
\pgfpathlineto{\pgfqpoint{1.160881in}{1.507470in}}%
\pgfpathlineto{\pgfqpoint{1.202807in}{1.507286in}}%
\pgfpathlineto{\pgfqpoint{1.243655in}{1.505666in}}%
\pgfpathlineto{\pgfqpoint{1.283592in}{1.502548in}}%
\pgfpathlineto{\pgfqpoint{1.309777in}{1.499439in}}%
\pgfpathlineto{\pgfqpoint{1.335650in}{1.494876in}}%
\pgfpathlineto{\pgfqpoint{1.348480in}{1.491678in}}%
\pgfpathlineto{\pgfqpoint{1.386581in}{1.480331in}}%
\pgfpathlineto{\pgfqpoint{1.411689in}{1.473808in}}%
\pgfpathlineto{\pgfqpoint{1.448963in}{1.462100in}}%
\pgfpathlineto{\pgfqpoint{1.546495in}{1.436692in}}%
\pgfpathlineto{\pgfqpoint{1.558528in}{1.432481in}}%
\pgfpathlineto{\pgfqpoint{1.570530in}{1.427225in}}%
\pgfpathlineto{\pgfqpoint{1.594450in}{1.414074in}}%
\pgfpathlineto{\pgfqpoint{1.606369in}{1.407414in}}%
\pgfpathlineto{\pgfqpoint{1.618263in}{1.401470in}}%
\pgfpathlineto{\pgfqpoint{1.630133in}{1.396768in}}%
\pgfpathlineto{\pgfqpoint{1.641980in}{1.393816in}}%
\pgfpathlineto{\pgfqpoint{1.653804in}{1.392867in}}%
\pgfpathlineto{\pgfqpoint{1.665607in}{1.393881in}}%
\pgfpathlineto{\pgfqpoint{1.677389in}{1.396554in}}%
\pgfpathlineto{\pgfqpoint{1.700896in}{1.404038in}}%
\pgfpathlineto{\pgfqpoint{1.712622in}{1.406758in}}%
\pgfpathlineto{\pgfqpoint{1.724330in}{1.407353in}}%
\pgfpathlineto{\pgfqpoint{1.736022in}{1.405147in}}%
\pgfpathlineto{\pgfqpoint{1.747699in}{1.400107in}}%
\pgfpathlineto{\pgfqpoint{1.759360in}{1.392864in}}%
\pgfpathlineto{\pgfqpoint{1.771006in}{1.384667in}}%
\pgfpathlineto{\pgfqpoint{1.782638in}{1.377303in}}%
\pgfpathlineto{\pgfqpoint{1.794257in}{1.372763in}}%
\pgfpathlineto{\pgfqpoint{1.805863in}{1.372694in}}%
\pgfpathlineto{\pgfqpoint{1.817457in}{1.377925in}}%
\pgfpathlineto{\pgfqpoint{1.829039in}{1.388239in}}%
\pgfpathlineto{\pgfqpoint{1.840609in}{1.402311in}}%
\pgfpathlineto{\pgfqpoint{1.852169in}{1.417882in}}%
\pgfpathlineto{\pgfqpoint{1.863718in}{1.432449in}}%
\pgfpathlineto{\pgfqpoint{1.875257in}{1.444284in}}%
\pgfpathlineto{\pgfqpoint{1.886786in}{1.453095in}}%
\pgfpathlineto{\pgfqpoint{1.898306in}{1.459542in}}%
\pgfpathlineto{\pgfqpoint{1.909817in}{1.463994in}}%
\pgfpathlineto{\pgfqpoint{1.921320in}{1.466023in}}%
\pgfpathlineto{\pgfqpoint{1.932814in}{1.464689in}}%
\pgfpathlineto{\pgfqpoint{1.944301in}{1.458855in}}%
\pgfpathlineto{\pgfqpoint{1.955780in}{1.447781in}}%
\pgfpathlineto{\pgfqpoint{1.967251in}{1.432261in}}%
\pgfpathlineto{\pgfqpoint{1.978716in}{1.415448in}}%
\pgfpathlineto{\pgfqpoint{1.990174in}{1.402329in}}%
\pgfpathlineto{\pgfqpoint{2.001625in}{1.397572in}}%
\pgfpathlineto{\pgfqpoint{2.013070in}{1.402969in}}%
\pgfpathlineto{\pgfqpoint{2.024510in}{1.416367in}}%
\pgfpathlineto{\pgfqpoint{2.035944in}{1.432577in}}%
\pgfpathlineto{\pgfqpoint{2.047372in}{1.445234in}}%
\pgfpathlineto{\pgfqpoint{2.058795in}{1.448995in}}%
\pgfpathlineto{\pgfqpoint{2.070213in}{1.442319in}}%
\pgfpathlineto{\pgfqpoint{2.093034in}{1.416536in}}%
\pgfpathlineto{\pgfqpoint{2.104438in}{1.409163in}}%
\pgfpathlineto{\pgfqpoint{2.115837in}{1.406703in}}%
\pgfpathlineto{\pgfqpoint{2.138624in}{1.405927in}}%
\pgfpathlineto{\pgfqpoint{2.150012in}{1.406001in}}%
\pgfpathlineto{\pgfqpoint{2.172776in}{1.407918in}}%
\pgfpathlineto{\pgfqpoint{2.184153in}{1.409610in}}%
\pgfpathlineto{\pgfqpoint{2.195527in}{1.411901in}}%
\pgfpathlineto{\pgfqpoint{2.206897in}{1.413223in}}%
\pgfpathlineto{\pgfqpoint{2.218265in}{1.411310in}}%
\pgfpathlineto{\pgfqpoint{2.229630in}{1.406630in}}%
\pgfpathlineto{\pgfqpoint{2.240991in}{1.402581in}}%
\pgfpathlineto{\pgfqpoint{2.252347in}{1.401113in}}%
\pgfpathlineto{\pgfqpoint{2.263686in}{1.400955in}}%
\pgfpathlineto{\pgfqpoint{2.274973in}{1.401566in}}%
\pgfpathlineto{\pgfqpoint{2.286134in}{1.404917in}}%
\pgfpathlineto{\pgfqpoint{2.297074in}{1.410955in}}%
\pgfpathlineto{\pgfqpoint{2.307710in}{1.415543in}}%
\pgfpathlineto{\pgfqpoint{2.318002in}{1.415637in}}%
\pgfpathlineto{\pgfqpoint{2.327947in}{1.412998in}}%
\pgfpathlineto{\pgfqpoint{2.337562in}{1.411099in}}%
\pgfpathlineto{\pgfqpoint{2.346867in}{1.411337in}}%
\pgfpathlineto{\pgfqpoint{2.355880in}{1.413440in}}%
\pgfpathlineto{\pgfqpoint{2.364620in}{1.417054in}}%
\pgfpathlineto{\pgfqpoint{2.389351in}{1.430948in}}%
\pgfpathlineto{\pgfqpoint{2.397144in}{1.431742in}}%
\pgfpathlineto{\pgfqpoint{2.404732in}{1.426823in}}%
\pgfpathlineto{\pgfqpoint{2.419333in}{1.407423in}}%
\pgfpathlineto{\pgfqpoint{2.426365in}{1.403591in}}%
\pgfpathlineto{\pgfqpoint{2.433230in}{1.407598in}}%
\pgfpathlineto{\pgfqpoint{2.439934in}{1.415519in}}%
\pgfpathlineto{\pgfqpoint{2.446486in}{1.420410in}}%
\pgfpathlineto{\pgfqpoint{2.452892in}{1.418716in}}%
\pgfpathlineto{\pgfqpoint{2.459159in}{1.413049in}}%
\pgfpathlineto{\pgfqpoint{2.465292in}{1.408605in}}%
\pgfpathlineto{\pgfqpoint{2.471297in}{1.408325in}}%
\pgfpathlineto{\pgfqpoint{2.477180in}{1.411532in}}%
\pgfpathlineto{\pgfqpoint{2.488596in}{1.420377in}}%
\pgfpathlineto{\pgfqpoint{2.494139in}{1.423961in}}%
\pgfpathlineto{\pgfqpoint{2.499577in}{1.425912in}}%
\pgfpathlineto{\pgfqpoint{2.504914in}{1.425573in}}%
\pgfpathlineto{\pgfqpoint{2.515301in}{1.422147in}}%
\pgfpathlineto{\pgfqpoint{2.530211in}{1.419031in}}%
\pgfpathlineto{\pgfqpoint{2.535014in}{1.418993in}}%
\pgfpathlineto{\pgfqpoint{2.539738in}{1.420407in}}%
\pgfpathlineto{\pgfqpoint{2.544386in}{1.422705in}}%
\pgfpathlineto{\pgfqpoint{2.548960in}{1.424053in}}%
\pgfpathlineto{\pgfqpoint{2.553463in}{1.422976in}}%
\pgfpathlineto{\pgfqpoint{2.566563in}{1.414407in}}%
\pgfpathlineto{\pgfqpoint{2.570801in}{1.413500in}}%
\pgfpathlineto{\pgfqpoint{2.574977in}{1.414521in}}%
\pgfpathlineto{\pgfqpoint{2.583152in}{1.418817in}}%
\pgfpathlineto{\pgfqpoint{2.587154in}{1.418264in}}%
\pgfpathlineto{\pgfqpoint{2.591101in}{1.416170in}}%
\pgfpathlineto{\pgfqpoint{2.594995in}{1.415112in}}%
\pgfpathlineto{\pgfqpoint{2.602628in}{1.417783in}}%
\pgfpathlineto{\pgfqpoint{2.606370in}{1.417003in}}%
\pgfpathlineto{\pgfqpoint{2.613711in}{1.409280in}}%
\pgfpathlineto{\pgfqpoint{2.617312in}{1.407879in}}%
\pgfpathlineto{\pgfqpoint{2.620869in}{1.410899in}}%
\pgfpathlineto{\pgfqpoint{2.627854in}{1.424137in}}%
\pgfpathlineto{\pgfqpoint{2.631284in}{1.427108in}}%
\pgfpathlineto{\pgfqpoint{2.634673in}{1.424609in}}%
\pgfpathlineto{\pgfqpoint{2.638023in}{1.418841in}}%
\pgfpathlineto{\pgfqpoint{2.641335in}{1.414251in}}%
\pgfpathlineto{\pgfqpoint{2.644609in}{1.413564in}}%
\pgfpathlineto{\pgfqpoint{2.647846in}{1.415550in}}%
\pgfpathlineto{\pgfqpoint{2.651047in}{1.416854in}}%
\pgfpathlineto{\pgfqpoint{2.654213in}{1.415448in}}%
\pgfpathlineto{\pgfqpoint{2.660442in}{1.408774in}}%
\pgfpathlineto{\pgfqpoint{2.663507in}{1.408094in}}%
\pgfpathlineto{\pgfqpoint{2.666540in}{1.410873in}}%
\pgfpathlineto{\pgfqpoint{2.672511in}{1.420326in}}%
\pgfpathlineto{\pgfqpoint{2.675450in}{1.421714in}}%
\pgfpathlineto{\pgfqpoint{2.678360in}{1.419831in}}%
\pgfpathlineto{\pgfqpoint{2.681241in}{1.416700in}}%
\pgfpathlineto{\pgfqpoint{2.684093in}{1.414747in}}%
\pgfpathlineto{\pgfqpoint{2.686918in}{1.414797in}}%
\pgfpathlineto{\pgfqpoint{2.692484in}{1.416356in}}%
\pgfpathlineto{\pgfqpoint{2.695228in}{1.415762in}}%
\pgfpathlineto{\pgfqpoint{2.703305in}{1.410395in}}%
\pgfpathlineto{\pgfqpoint{2.705948in}{1.410034in}}%
\pgfpathlineto{\pgfqpoint{2.708567in}{1.411058in}}%
\pgfpathlineto{\pgfqpoint{2.716283in}{1.416045in}}%
\pgfpathlineto{\pgfqpoint{2.718810in}{1.416995in}}%
\pgfpathlineto{\pgfqpoint{2.721315in}{1.417113in}}%
\pgfpathlineto{\pgfqpoint{2.726260in}{1.416242in}}%
\pgfpathlineto{\pgfqpoint{2.731122in}{1.417226in}}%
\pgfpathlineto{\pgfqpoint{2.738264in}{1.420119in}}%
\pgfpathlineto{\pgfqpoint{2.740606in}{1.420302in}}%
\pgfpathlineto{\pgfqpoint{2.742929in}{1.419545in}}%
\pgfpathlineto{\pgfqpoint{2.749788in}{1.415624in}}%
\pgfpathlineto{\pgfqpoint{2.754271in}{1.415463in}}%
\pgfpathlineto{\pgfqpoint{2.758686in}{1.415639in}}%
\pgfpathlineto{\pgfqpoint{2.767317in}{1.414178in}}%
\pgfpathlineto{\pgfqpoint{2.777754in}{1.416643in}}%
\pgfpathlineto{\pgfqpoint{2.785840in}{1.415758in}}%
\pgfpathlineto{\pgfqpoint{2.793705in}{1.417286in}}%
\pgfpathlineto{\pgfqpoint{2.817876in}{1.417226in}}%
\pgfpathlineto{\pgfqpoint{2.826658in}{1.416611in}}%
\pgfpathlineto{\pgfqpoint{2.835180in}{1.417494in}}%
\pgfpathlineto{\pgfqpoint{2.863925in}{1.417133in}}%
\pgfpathlineto{\pgfqpoint{2.871426in}{1.416975in}}%
\pgfpathlineto{\pgfqpoint{2.883035in}{1.417132in}}%
\pgfpathlineto{\pgfqpoint{2.892822in}{1.416920in}}%
\pgfpathlineto{\pgfqpoint{2.908863in}{1.417029in}}%
\pgfpathlineto{\pgfqpoint{2.946587in}{1.417158in}}%
\pgfpathlineto{\pgfqpoint{2.952234in}{1.417167in}}%
\pgfpathlineto{\pgfqpoint{3.043110in}{1.417119in}}%
\pgfpathlineto{\pgfqpoint{3.043110in}{1.417119in}}%
\pgfusepath{stroke}%
\end{pgfscope}%
\begin{pgfscope}%
\pgfpathrectangle{\pgfqpoint{0.650000in}{0.420000in}}{\pgfqpoint{2.388110in}{1.994488in}}%
\pgfusepath{clip}%
\pgfsetbuttcap%
\pgfsetroundjoin%
\pgfsetlinewidth{1.003750pt}%
\definecolor{currentstroke}{rgb}{0.501961,0.450980,0.674510}%
\pgfsetstrokecolor{currentstroke}%
\pgfsetdash{{5.000000pt}{5.000000pt}}{0.000000pt}%
\pgfpathmoveto{\pgfqpoint{0.645000in}{1.419221in}}%
\pgfpathlineto{\pgfqpoint{0.766279in}{1.420984in}}%
\pgfpathlineto{\pgfqpoint{0.830619in}{1.423259in}}%
\pgfpathlineto{\pgfqpoint{0.870161in}{1.425790in}}%
\pgfpathlineto{\pgfqpoint{0.907582in}{1.429505in}}%
\pgfpathlineto{\pgfqpoint{0.943235in}{1.434759in}}%
\pgfpathlineto{\pgfqpoint{0.960487in}{1.438080in}}%
\pgfpathlineto{\pgfqpoint{0.977396in}{1.441924in}}%
\pgfpathlineto{\pgfqpoint{0.993988in}{1.446334in}}%
\pgfpathlineto{\pgfqpoint{1.010285in}{1.451349in}}%
\pgfpathlineto{\pgfqpoint{1.026310in}{1.457005in}}%
\pgfpathlineto{\pgfqpoint{1.042080in}{1.463326in}}%
\pgfpathlineto{\pgfqpoint{1.057613in}{1.470329in}}%
\pgfpathlineto{\pgfqpoint{1.072926in}{1.478018in}}%
\pgfpathlineto{\pgfqpoint{1.088031in}{1.486381in}}%
\pgfpathlineto{\pgfqpoint{1.102943in}{1.495389in}}%
\pgfpathlineto{\pgfqpoint{1.132233in}{1.515114in}}%
\pgfpathlineto{\pgfqpoint{1.160881in}{1.536480in}}%
\pgfpathlineto{\pgfqpoint{1.202807in}{1.568881in}}%
\pgfpathlineto{\pgfqpoint{1.216534in}{1.578886in}}%
\pgfpathlineto{\pgfqpoint{1.230148in}{1.588030in}}%
\pgfpathlineto{\pgfqpoint{1.243655in}{1.595996in}}%
\pgfpathlineto{\pgfqpoint{1.257062in}{1.602441in}}%
\pgfpathlineto{\pgfqpoint{1.270372in}{1.607022in}}%
\pgfpathlineto{\pgfqpoint{1.283592in}{1.609397in}}%
\pgfpathlineto{\pgfqpoint{1.296725in}{1.609253in}}%
\pgfpathlineto{\pgfqpoint{1.309777in}{1.606315in}}%
\pgfpathlineto{\pgfqpoint{1.322750in}{1.600385in}}%
\pgfpathlineto{\pgfqpoint{1.335650in}{1.591372in}}%
\pgfpathlineto{\pgfqpoint{1.348480in}{1.579307in}}%
\pgfpathlineto{\pgfqpoint{1.361243in}{1.564348in}}%
\pgfpathlineto{\pgfqpoint{1.373943in}{1.546693in}}%
\pgfpathlineto{\pgfqpoint{1.386581in}{1.526407in}}%
\pgfpathlineto{\pgfqpoint{1.399163in}{1.503313in}}%
\pgfpathlineto{\pgfqpoint{1.411689in}{1.477383in}}%
\pgfpathlineto{\pgfqpoint{1.436586in}{1.420037in}}%
\pgfpathlineto{\pgfqpoint{1.473581in}{1.333503in}}%
\pgfpathlineto{\pgfqpoint{1.498032in}{1.280588in}}%
\pgfpathlineto{\pgfqpoint{1.510201in}{1.256483in}}%
\pgfpathlineto{\pgfqpoint{1.522333in}{1.234426in}}%
\pgfpathlineto{\pgfqpoint{1.534430in}{1.214867in}}%
\pgfpathlineto{\pgfqpoint{1.546495in}{1.198419in}}%
\pgfpathlineto{\pgfqpoint{1.558528in}{1.185868in}}%
\pgfpathlineto{\pgfqpoint{1.570530in}{1.177873in}}%
\pgfpathlineto{\pgfqpoint{1.582504in}{1.174191in}}%
\pgfpathlineto{\pgfqpoint{1.594450in}{1.173625in}}%
\pgfpathlineto{\pgfqpoint{1.606369in}{1.175399in}}%
\pgfpathlineto{\pgfqpoint{1.618263in}{1.180063in}}%
\pgfpathlineto{\pgfqpoint{1.630133in}{1.187957in}}%
\pgfpathlineto{\pgfqpoint{1.653804in}{1.208515in}}%
\pgfpathlineto{\pgfqpoint{1.665607in}{1.219450in}}%
\pgfpathlineto{\pgfqpoint{1.677389in}{1.231966in}}%
\pgfpathlineto{\pgfqpoint{1.689152in}{1.247365in}}%
\pgfpathlineto{\pgfqpoint{1.712622in}{1.282992in}}%
\pgfpathlineto{\pgfqpoint{1.724330in}{1.297505in}}%
\pgfpathlineto{\pgfqpoint{1.736022in}{1.308361in}}%
\pgfpathlineto{\pgfqpoint{1.759360in}{1.327734in}}%
\pgfpathlineto{\pgfqpoint{1.782638in}{1.347435in}}%
\pgfpathlineto{\pgfqpoint{1.794257in}{1.356526in}}%
\pgfpathlineto{\pgfqpoint{1.817457in}{1.372389in}}%
\pgfpathlineto{\pgfqpoint{1.852169in}{1.395397in}}%
\pgfpathlineto{\pgfqpoint{1.863718in}{1.401684in}}%
\pgfpathlineto{\pgfqpoint{1.875257in}{1.404695in}}%
\pgfpathlineto{\pgfqpoint{1.886786in}{1.405785in}}%
\pgfpathlineto{\pgfqpoint{1.898306in}{1.408886in}}%
\pgfpathlineto{\pgfqpoint{1.921320in}{1.420268in}}%
\pgfpathlineto{\pgfqpoint{1.932814in}{1.424597in}}%
\pgfpathlineto{\pgfqpoint{1.944301in}{1.428074in}}%
\pgfpathlineto{\pgfqpoint{1.955780in}{1.430373in}}%
\pgfpathlineto{\pgfqpoint{1.967251in}{1.431027in}}%
\pgfpathlineto{\pgfqpoint{1.978716in}{1.430500in}}%
\pgfpathlineto{\pgfqpoint{2.001625in}{1.428660in}}%
\pgfpathlineto{\pgfqpoint{2.013070in}{1.429915in}}%
\pgfpathlineto{\pgfqpoint{2.024510in}{1.433080in}}%
\pgfpathlineto{\pgfqpoint{2.047372in}{1.440374in}}%
\pgfpathlineto{\pgfqpoint{2.058795in}{1.441425in}}%
\pgfpathlineto{\pgfqpoint{2.070213in}{1.439496in}}%
\pgfpathlineto{\pgfqpoint{2.104438in}{1.428850in}}%
\pgfpathlineto{\pgfqpoint{2.115837in}{1.426947in}}%
\pgfpathlineto{\pgfqpoint{2.127233in}{1.426795in}}%
\pgfpathlineto{\pgfqpoint{2.161396in}{1.428922in}}%
\pgfpathlineto{\pgfqpoint{2.206897in}{1.429500in}}%
\pgfpathlineto{\pgfqpoint{2.218265in}{1.427106in}}%
\pgfpathlineto{\pgfqpoint{2.229630in}{1.422348in}}%
\pgfpathlineto{\pgfqpoint{2.240991in}{1.418702in}}%
\pgfpathlineto{\pgfqpoint{2.252347in}{1.418192in}}%
\pgfpathlineto{\pgfqpoint{2.274973in}{1.419680in}}%
\pgfpathlineto{\pgfqpoint{2.297074in}{1.423553in}}%
\pgfpathlineto{\pgfqpoint{2.307710in}{1.422023in}}%
\pgfpathlineto{\pgfqpoint{2.318002in}{1.417887in}}%
\pgfpathlineto{\pgfqpoint{2.327947in}{1.415200in}}%
\pgfpathlineto{\pgfqpoint{2.337562in}{1.415501in}}%
\pgfpathlineto{\pgfqpoint{2.364620in}{1.421430in}}%
\pgfpathlineto{\pgfqpoint{2.373102in}{1.420766in}}%
\pgfpathlineto{\pgfqpoint{2.389351in}{1.418051in}}%
\pgfpathlineto{\pgfqpoint{2.397144in}{1.418373in}}%
\pgfpathlineto{\pgfqpoint{2.412125in}{1.421379in}}%
\pgfpathlineto{\pgfqpoint{2.433230in}{1.422266in}}%
\pgfpathlineto{\pgfqpoint{2.439934in}{1.420807in}}%
\pgfpathlineto{\pgfqpoint{2.452892in}{1.412040in}}%
\pgfpathlineto{\pgfqpoint{2.459159in}{1.409415in}}%
\pgfpathlineto{\pgfqpoint{2.465292in}{1.409708in}}%
\pgfpathlineto{\pgfqpoint{2.477180in}{1.414981in}}%
\pgfpathlineto{\pgfqpoint{2.482944in}{1.416026in}}%
\pgfpathlineto{\pgfqpoint{2.488596in}{1.414471in}}%
\pgfpathlineto{\pgfqpoint{2.504914in}{1.406421in}}%
\pgfpathlineto{\pgfqpoint{2.510155in}{1.405511in}}%
\pgfpathlineto{\pgfqpoint{2.515301in}{1.407037in}}%
\pgfpathlineto{\pgfqpoint{2.520357in}{1.409146in}}%
\pgfpathlineto{\pgfqpoint{2.525326in}{1.409160in}}%
\pgfpathlineto{\pgfqpoint{2.530211in}{1.407922in}}%
\pgfpathlineto{\pgfqpoint{2.535014in}{1.408398in}}%
\pgfpathlineto{\pgfqpoint{2.544386in}{1.413429in}}%
\pgfpathlineto{\pgfqpoint{2.548960in}{1.413794in}}%
\pgfpathlineto{\pgfqpoint{2.557896in}{1.412828in}}%
\pgfpathlineto{\pgfqpoint{2.562262in}{1.413748in}}%
\pgfpathlineto{\pgfqpoint{2.566563in}{1.416493in}}%
\pgfpathlineto{\pgfqpoint{2.570801in}{1.421921in}}%
\pgfpathlineto{\pgfqpoint{2.574977in}{1.428444in}}%
\pgfpathlineto{\pgfqpoint{2.579093in}{1.431490in}}%
\pgfpathlineto{\pgfqpoint{2.583152in}{1.428536in}}%
\pgfpathlineto{\pgfqpoint{2.591101in}{1.417266in}}%
\pgfpathlineto{\pgfqpoint{2.594995in}{1.414263in}}%
\pgfpathlineto{\pgfqpoint{2.598836in}{1.413228in}}%
\pgfpathlineto{\pgfqpoint{2.602628in}{1.414136in}}%
\pgfpathlineto{\pgfqpoint{2.610063in}{1.417648in}}%
\pgfpathlineto{\pgfqpoint{2.613711in}{1.418782in}}%
\pgfpathlineto{\pgfqpoint{2.617312in}{1.419175in}}%
\pgfpathlineto{\pgfqpoint{2.620869in}{1.418012in}}%
\pgfpathlineto{\pgfqpoint{2.624383in}{1.415788in}}%
\pgfpathlineto{\pgfqpoint{2.627854in}{1.414878in}}%
\pgfpathlineto{\pgfqpoint{2.634673in}{1.418683in}}%
\pgfpathlineto{\pgfqpoint{2.638023in}{1.418759in}}%
\pgfpathlineto{\pgfqpoint{2.644609in}{1.414866in}}%
\pgfpathlineto{\pgfqpoint{2.647846in}{1.414231in}}%
\pgfpathlineto{\pgfqpoint{2.657345in}{1.415025in}}%
\pgfpathlineto{\pgfqpoint{2.660442in}{1.416408in}}%
\pgfpathlineto{\pgfqpoint{2.666540in}{1.419827in}}%
\pgfpathlineto{\pgfqpoint{2.672511in}{1.422038in}}%
\pgfpathlineto{\pgfqpoint{2.675450in}{1.423759in}}%
\pgfpathlineto{\pgfqpoint{2.678360in}{1.424006in}}%
\pgfpathlineto{\pgfqpoint{2.681241in}{1.421222in}}%
\pgfpathlineto{\pgfqpoint{2.684093in}{1.416951in}}%
\pgfpathlineto{\pgfqpoint{2.686918in}{1.414675in}}%
\pgfpathlineto{\pgfqpoint{2.689714in}{1.415865in}}%
\pgfpathlineto{\pgfqpoint{2.695228in}{1.422521in}}%
\pgfpathlineto{\pgfqpoint{2.697946in}{1.424650in}}%
\pgfpathlineto{\pgfqpoint{2.700638in}{1.424655in}}%
\pgfpathlineto{\pgfqpoint{2.703305in}{1.422688in}}%
\pgfpathlineto{\pgfqpoint{2.708567in}{1.418015in}}%
\pgfpathlineto{\pgfqpoint{2.711162in}{1.417929in}}%
\pgfpathlineto{\pgfqpoint{2.713734in}{1.419870in}}%
\pgfpathlineto{\pgfqpoint{2.718810in}{1.425136in}}%
\pgfpathlineto{\pgfqpoint{2.721315in}{1.426188in}}%
\pgfpathlineto{\pgfqpoint{2.726260in}{1.426927in}}%
\pgfpathlineto{\pgfqpoint{2.728702in}{1.426224in}}%
\pgfpathlineto{\pgfqpoint{2.733523in}{1.421660in}}%
\pgfpathlineto{\pgfqpoint{2.735903in}{1.421765in}}%
\pgfpathlineto{\pgfqpoint{2.740606in}{1.426634in}}%
\pgfpathlineto{\pgfqpoint{2.742929in}{1.426502in}}%
\pgfpathlineto{\pgfqpoint{2.747520in}{1.422505in}}%
\pgfpathlineto{\pgfqpoint{2.749788in}{1.423336in}}%
\pgfpathlineto{\pgfqpoint{2.752038in}{1.425138in}}%
\pgfpathlineto{\pgfqpoint{2.754271in}{1.425158in}}%
\pgfpathlineto{\pgfqpoint{2.763034in}{1.417813in}}%
\pgfpathlineto{\pgfqpoint{2.767317in}{1.409432in}}%
\pgfpathlineto{\pgfqpoint{2.769435in}{1.408615in}}%
\pgfpathlineto{\pgfqpoint{2.775697in}{1.418832in}}%
\pgfpathlineto{\pgfqpoint{2.779797in}{1.420585in}}%
\pgfpathlineto{\pgfqpoint{2.781826in}{1.420036in}}%
\pgfpathlineto{\pgfqpoint{2.783840in}{1.418852in}}%
\pgfpathlineto{\pgfqpoint{2.785840in}{1.418584in}}%
\pgfpathlineto{\pgfqpoint{2.793705in}{1.421682in}}%
\pgfpathlineto{\pgfqpoint{2.795638in}{1.420327in}}%
\pgfpathlineto{\pgfqpoint{2.805113in}{1.407715in}}%
\pgfpathlineto{\pgfqpoint{2.806972in}{1.406488in}}%
\pgfpathlineto{\pgfqpoint{2.808818in}{1.408193in}}%
\pgfpathlineto{\pgfqpoint{2.810653in}{1.411033in}}%
\pgfpathlineto{\pgfqpoint{2.812476in}{1.410982in}}%
\pgfpathlineto{\pgfqpoint{2.816087in}{1.404409in}}%
\pgfpathlineto{\pgfqpoint{2.817876in}{1.406798in}}%
\pgfpathlineto{\pgfqpoint{2.821421in}{1.416924in}}%
\pgfpathlineto{\pgfqpoint{2.823177in}{1.417080in}}%
\pgfpathlineto{\pgfqpoint{2.826658in}{1.410812in}}%
\pgfpathlineto{\pgfqpoint{2.828383in}{1.409435in}}%
\pgfpathlineto{\pgfqpoint{2.830097in}{1.409697in}}%
\pgfpathlineto{\pgfqpoint{2.836854in}{1.416959in}}%
\pgfpathlineto{\pgfqpoint{2.840175in}{1.425199in}}%
\pgfpathlineto{\pgfqpoint{2.841820in}{1.426140in}}%
\pgfpathlineto{\pgfqpoint{2.843457in}{1.423203in}}%
\pgfpathlineto{\pgfqpoint{2.845084in}{1.419127in}}%
\pgfpathlineto{\pgfqpoint{2.846702in}{1.417379in}}%
\pgfpathlineto{\pgfqpoint{2.851502in}{1.420488in}}%
\pgfpathlineto{\pgfqpoint{2.853085in}{1.420159in}}%
\pgfpathlineto{\pgfqpoint{2.854659in}{1.420453in}}%
\pgfpathlineto{\pgfqpoint{2.857781in}{1.422428in}}%
\pgfpathlineto{\pgfqpoint{2.859329in}{1.422326in}}%
\pgfpathlineto{\pgfqpoint{2.863925in}{1.416224in}}%
\pgfpathlineto{\pgfqpoint{2.866949in}{1.420601in}}%
\pgfpathlineto{\pgfqpoint{2.868449in}{1.420277in}}%
\pgfpathlineto{\pgfqpoint{2.871426in}{1.412361in}}%
\pgfpathlineto{\pgfqpoint{2.872903in}{1.411050in}}%
\pgfpathlineto{\pgfqpoint{2.874372in}{1.412875in}}%
\pgfpathlineto{\pgfqpoint{2.877289in}{1.418139in}}%
\pgfpathlineto{\pgfqpoint{2.878736in}{1.419307in}}%
\pgfpathlineto{\pgfqpoint{2.881609in}{1.419889in}}%
\pgfpathlineto{\pgfqpoint{2.884454in}{1.419631in}}%
\pgfpathlineto{\pgfqpoint{2.887270in}{1.416995in}}%
\pgfpathlineto{\pgfqpoint{2.891444in}{1.412036in}}%
\pgfpathlineto{\pgfqpoint{2.892822in}{1.411510in}}%
\pgfpathlineto{\pgfqpoint{2.894194in}{1.412910in}}%
\pgfpathlineto{\pgfqpoint{2.896917in}{1.417147in}}%
\pgfpathlineto{\pgfqpoint{2.898269in}{1.417904in}}%
\pgfpathlineto{\pgfqpoint{2.902288in}{1.417702in}}%
\pgfpathlineto{\pgfqpoint{2.906251in}{1.419115in}}%
\pgfpathlineto{\pgfqpoint{2.908863in}{1.422248in}}%
\pgfpathlineto{\pgfqpoint{2.910161in}{1.422173in}}%
\pgfpathlineto{\pgfqpoint{2.914018in}{1.418024in}}%
\pgfpathlineto{\pgfqpoint{2.916561in}{1.414800in}}%
\pgfpathlineto{\pgfqpoint{2.917824in}{1.414861in}}%
\pgfpathlineto{\pgfqpoint{2.921581in}{1.417003in}}%
\pgfpathlineto{\pgfqpoint{2.924058in}{1.416941in}}%
\pgfpathlineto{\pgfqpoint{2.926514in}{1.418401in}}%
\pgfpathlineto{\pgfqpoint{2.927735in}{1.418125in}}%
\pgfpathlineto{\pgfqpoint{2.930160in}{1.416859in}}%
\pgfpathlineto{\pgfqpoint{2.932565in}{1.418010in}}%
\pgfpathlineto{\pgfqpoint{2.933760in}{1.417793in}}%
\pgfpathlineto{\pgfqpoint{2.937315in}{1.413911in}}%
\pgfpathlineto{\pgfqpoint{2.939661in}{1.414804in}}%
\pgfpathlineto{\pgfqpoint{2.945445in}{1.417161in}}%
\pgfpathlineto{\pgfqpoint{2.946587in}{1.417407in}}%
\pgfpathlineto{\pgfqpoint{2.949989in}{1.420081in}}%
\pgfpathlineto{\pgfqpoint{2.951114in}{1.419942in}}%
\pgfpathlineto{\pgfqpoint{2.953351in}{1.418177in}}%
\pgfpathlineto{\pgfqpoint{2.956674in}{1.418417in}}%
\pgfpathlineto{\pgfqpoint{2.959959in}{1.414497in}}%
\pgfpathlineto{\pgfqpoint{2.962128in}{1.416290in}}%
\pgfpathlineto{\pgfqpoint{2.964281in}{1.417456in}}%
\pgfpathlineto{\pgfqpoint{2.966418in}{1.417895in}}%
\pgfpathlineto{\pgfqpoint{2.967481in}{1.417201in}}%
\pgfpathlineto{\pgfqpoint{2.969595in}{1.414669in}}%
\pgfpathlineto{\pgfqpoint{2.972736in}{1.414394in}}%
\pgfpathlineto{\pgfqpoint{2.973776in}{1.413566in}}%
\pgfpathlineto{\pgfqpoint{2.977898in}{1.413526in}}%
\pgfpathlineto{\pgfqpoint{2.978919in}{1.415124in}}%
\pgfpathlineto{\pgfqpoint{2.980951in}{1.419979in}}%
\pgfpathlineto{\pgfqpoint{2.981961in}{1.419702in}}%
\pgfpathlineto{\pgfqpoint{2.983971in}{1.414834in}}%
\pgfpathlineto{\pgfqpoint{2.984971in}{1.414280in}}%
\pgfpathlineto{\pgfqpoint{2.985968in}{1.416095in}}%
\pgfpathlineto{\pgfqpoint{2.987950in}{1.421580in}}%
\pgfpathlineto{\pgfqpoint{2.988937in}{1.421925in}}%
\pgfpathlineto{\pgfqpoint{2.992848in}{1.416758in}}%
\pgfpathlineto{\pgfqpoint{2.995747in}{1.415677in}}%
\pgfpathlineto{\pgfqpoint{2.997664in}{1.416500in}}%
\pgfpathlineto{\pgfqpoint{2.999568in}{1.417847in}}%
\pgfpathlineto{\pgfqpoint{3.001460in}{1.417611in}}%
\pgfpathlineto{\pgfqpoint{3.006135in}{1.414634in}}%
\pgfpathlineto{\pgfqpoint{3.009821in}{1.417018in}}%
\pgfpathlineto{\pgfqpoint{3.014363in}{1.420436in}}%
\pgfpathlineto{\pgfqpoint{3.015263in}{1.419879in}}%
\pgfpathlineto{\pgfqpoint{3.018835in}{1.415401in}}%
\pgfpathlineto{\pgfqpoint{3.020604in}{1.416929in}}%
\pgfpathlineto{\pgfqpoint{3.022363in}{1.418223in}}%
\pgfpathlineto{\pgfqpoint{3.024981in}{1.417642in}}%
\pgfpathlineto{\pgfqpoint{3.027575in}{1.415958in}}%
\pgfpathlineto{\pgfqpoint{3.028435in}{1.416255in}}%
\pgfpathlineto{\pgfqpoint{3.031848in}{1.419974in}}%
\pgfpathlineto{\pgfqpoint{3.033540in}{1.417433in}}%
\pgfpathlineto{\pgfqpoint{3.035221in}{1.415703in}}%
\pgfpathlineto{\pgfqpoint{3.039383in}{1.417410in}}%
\pgfpathlineto{\pgfqpoint{3.043110in}{1.417080in}}%
\pgfpathlineto{\pgfqpoint{3.043110in}{1.417080in}}%
\pgfusepath{stroke}%
\end{pgfscope}%
\begin{pgfscope}%
\pgfpathrectangle{\pgfqpoint{0.650000in}{0.420000in}}{\pgfqpoint{2.388110in}{1.994488in}}%
\pgfusepath{clip}%
\pgfsetbuttcap%
\pgfsetroundjoin%
\pgfsetlinewidth{1.003750pt}%
\definecolor{currentstroke}{rgb}{0.501961,0.450980,0.674510}%
\pgfsetstrokecolor{currentstroke}%
\pgfsetdash{{3.000000pt}{5.000000pt}{1.000000pt}{5.000000pt}}{0.000000pt}%
\pgfpathmoveto{\pgfqpoint{0.645000in}{2.186555in}}%
\pgfpathlineto{\pgfqpoint{0.667196in}{2.185828in}}%
\pgfpathlineto{\pgfqpoint{0.693779in}{2.183917in}}%
\pgfpathlineto{\pgfqpoint{0.719048in}{2.180966in}}%
\pgfpathlineto{\pgfqpoint{0.743169in}{2.176932in}}%
\pgfpathlineto{\pgfqpoint{0.766279in}{2.171778in}}%
\pgfpathlineto{\pgfqpoint{0.788495in}{2.165469in}}%
\pgfpathlineto{\pgfqpoint{0.809915in}{2.157978in}}%
\pgfpathlineto{\pgfqpoint{0.830619in}{2.149282in}}%
\pgfpathlineto{\pgfqpoint{0.850681in}{2.139365in}}%
\pgfpathlineto{\pgfqpoint{0.870161in}{2.128219in}}%
\pgfpathlineto{\pgfqpoint{0.889112in}{2.115842in}}%
\pgfpathlineto{\pgfqpoint{0.907582in}{2.102242in}}%
\pgfpathlineto{\pgfqpoint{0.925610in}{2.087433in}}%
\pgfpathlineto{\pgfqpoint{0.943235in}{2.071439in}}%
\pgfpathlineto{\pgfqpoint{0.960487in}{2.054292in}}%
\pgfpathlineto{\pgfqpoint{0.977396in}{2.036031in}}%
\pgfpathlineto{\pgfqpoint{0.993988in}{2.016699in}}%
\pgfpathlineto{\pgfqpoint{1.010285in}{1.996350in}}%
\pgfpathlineto{\pgfqpoint{1.026310in}{1.975036in}}%
\pgfpathlineto{\pgfqpoint{1.042080in}{1.952811in}}%
\pgfpathlineto{\pgfqpoint{1.057613in}{1.929728in}}%
\pgfpathlineto{\pgfqpoint{1.072926in}{1.905830in}}%
\pgfpathlineto{\pgfqpoint{1.088031in}{1.881156in}}%
\pgfpathlineto{\pgfqpoint{1.117673in}{1.829535in}}%
\pgfpathlineto{\pgfqpoint{1.146632in}{1.774748in}}%
\pgfpathlineto{\pgfqpoint{1.160881in}{1.746006in}}%
\pgfpathlineto{\pgfqpoint{1.174988in}{1.716208in}}%
\pgfpathlineto{\pgfqpoint{1.188960in}{1.685195in}}%
\pgfpathlineto{\pgfqpoint{1.202807in}{1.652778in}}%
\pgfpathlineto{\pgfqpoint{1.216534in}{1.618754in}}%
\pgfpathlineto{\pgfqpoint{1.230148in}{1.582926in}}%
\pgfpathlineto{\pgfqpoint{1.243655in}{1.545146in}}%
\pgfpathlineto{\pgfqpoint{1.257062in}{1.505363in}}%
\pgfpathlineto{\pgfqpoint{1.283592in}{1.420620in}}%
\pgfpathlineto{\pgfqpoint{1.309777in}{1.334305in}}%
\pgfpathlineto{\pgfqpoint{1.322750in}{1.295098in}}%
\pgfpathlineto{\pgfqpoint{1.335650in}{1.262928in}}%
\pgfpathlineto{\pgfqpoint{1.348480in}{1.241786in}}%
\pgfpathlineto{\pgfqpoint{1.361243in}{1.232662in}}%
\pgfpathlineto{\pgfqpoint{1.373943in}{1.228368in}}%
\pgfpathlineto{\pgfqpoint{1.386581in}{1.216408in}}%
\pgfpathlineto{\pgfqpoint{1.411689in}{1.171204in}}%
\pgfpathlineto{\pgfqpoint{1.424163in}{1.158171in}}%
\pgfpathlineto{\pgfqpoint{1.436586in}{1.153663in}}%
\pgfpathlineto{\pgfqpoint{1.448963in}{1.153514in}}%
\pgfpathlineto{\pgfqpoint{1.461293in}{1.156039in}}%
\pgfpathlineto{\pgfqpoint{1.473581in}{1.161387in}}%
\pgfpathlineto{\pgfqpoint{1.485826in}{1.169906in}}%
\pgfpathlineto{\pgfqpoint{1.498032in}{1.181522in}}%
\pgfpathlineto{\pgfqpoint{1.510201in}{1.195673in}}%
\pgfpathlineto{\pgfqpoint{1.606369in}{1.319751in}}%
\pgfpathlineto{\pgfqpoint{1.618263in}{1.332111in}}%
\pgfpathlineto{\pgfqpoint{1.630133in}{1.343139in}}%
\pgfpathlineto{\pgfqpoint{1.653804in}{1.362674in}}%
\pgfpathlineto{\pgfqpoint{1.665607in}{1.371476in}}%
\pgfpathlineto{\pgfqpoint{1.677389in}{1.379455in}}%
\pgfpathlineto{\pgfqpoint{1.689152in}{1.386322in}}%
\pgfpathlineto{\pgfqpoint{1.700896in}{1.392052in}}%
\pgfpathlineto{\pgfqpoint{1.724330in}{1.401636in}}%
\pgfpathlineto{\pgfqpoint{1.747699in}{1.409743in}}%
\pgfpathlineto{\pgfqpoint{1.771006in}{1.415366in}}%
\pgfpathlineto{\pgfqpoint{1.794257in}{1.419039in}}%
\pgfpathlineto{\pgfqpoint{1.817457in}{1.421065in}}%
\pgfpathlineto{\pgfqpoint{1.840609in}{1.421730in}}%
\pgfpathlineto{\pgfqpoint{1.886786in}{1.421211in}}%
\pgfpathlineto{\pgfqpoint{2.138624in}{1.417626in}}%
\pgfpathlineto{\pgfqpoint{2.172776in}{1.417229in}}%
\pgfpathlineto{\pgfqpoint{2.206897in}{1.416546in}}%
\pgfpathlineto{\pgfqpoint{2.346867in}{1.416815in}}%
\pgfpathlineto{\pgfqpoint{2.412125in}{1.416686in}}%
\pgfpathlineto{\pgfqpoint{2.666540in}{1.416928in}}%
\pgfpathlineto{\pgfqpoint{2.752038in}{1.417046in}}%
\pgfpathlineto{\pgfqpoint{2.961045in}{1.416977in}}%
\pgfpathlineto{\pgfqpoint{3.043110in}{1.416977in}}%
\pgfpathlineto{\pgfqpoint{3.043110in}{1.416977in}}%
\pgfusepath{stroke}%
\end{pgfscope}%
\begin{pgfscope}%
\pgfpathrectangle{\pgfqpoint{0.650000in}{0.420000in}}{\pgfqpoint{2.388110in}{1.994488in}}%
\pgfusepath{clip}%
\pgfsetbuttcap%
\pgfsetroundjoin%
\pgfsetlinewidth{1.003750pt}%
\definecolor{currentstroke}{rgb}{0.501961,0.450980,0.674510}%
\pgfsetstrokecolor{currentstroke}%
\pgfsetdash{{1.000000pt}{5.000000pt}}{0.000000pt}%
\pgfpathmoveto{\pgfqpoint{0.645000in}{1.402606in}}%
\pgfpathlineto{\pgfqpoint{0.667196in}{1.411934in}}%
\pgfpathlineto{\pgfqpoint{0.693779in}{1.418432in}}%
\pgfpathlineto{\pgfqpoint{0.719048in}{1.423770in}}%
\pgfpathlineto{\pgfqpoint{0.743169in}{1.428089in}}%
\pgfpathlineto{\pgfqpoint{0.766279in}{1.431487in}}%
\pgfpathlineto{\pgfqpoint{0.788495in}{1.434035in}}%
\pgfpathlineto{\pgfqpoint{0.809915in}{1.435790in}}%
\pgfpathlineto{\pgfqpoint{0.830619in}{1.436799in}}%
\pgfpathlineto{\pgfqpoint{0.850681in}{1.437111in}}%
\pgfpathlineto{\pgfqpoint{0.870161in}{1.436776in}}%
\pgfpathlineto{\pgfqpoint{0.907582in}{1.434409in}}%
\pgfpathlineto{\pgfqpoint{0.943235in}{1.430276in}}%
\pgfpathlineto{\pgfqpoint{0.993988in}{1.422459in}}%
\pgfpathlineto{\pgfqpoint{1.026310in}{1.417546in}}%
\pgfpathlineto{\pgfqpoint{1.057613in}{1.414061in}}%
\pgfpathlineto{\pgfqpoint{1.072926in}{1.413154in}}%
\pgfpathlineto{\pgfqpoint{1.088031in}{1.412933in}}%
\pgfpathlineto{\pgfqpoint{1.102943in}{1.413467in}}%
\pgfpathlineto{\pgfqpoint{1.117673in}{1.414784in}}%
\pgfpathlineto{\pgfqpoint{1.132233in}{1.416871in}}%
\pgfpathlineto{\pgfqpoint{1.146632in}{1.419659in}}%
\pgfpathlineto{\pgfqpoint{1.174988in}{1.426755in}}%
\pgfpathlineto{\pgfqpoint{1.202807in}{1.434207in}}%
\pgfpathlineto{\pgfqpoint{1.216534in}{1.437193in}}%
\pgfpathlineto{\pgfqpoint{1.230148in}{1.439120in}}%
\pgfpathlineto{\pgfqpoint{1.243655in}{1.439561in}}%
\pgfpathlineto{\pgfqpoint{1.257062in}{1.438139in}}%
\pgfpathlineto{\pgfqpoint{1.270372in}{1.434619in}}%
\pgfpathlineto{\pgfqpoint{1.283592in}{1.429027in}}%
\pgfpathlineto{\pgfqpoint{1.322750in}{1.407705in}}%
\pgfpathlineto{\pgfqpoint{1.335650in}{1.405896in}}%
\pgfpathlineto{\pgfqpoint{1.348480in}{1.412150in}}%
\pgfpathlineto{\pgfqpoint{1.361243in}{1.427272in}}%
\pgfpathlineto{\pgfqpoint{1.373943in}{1.444362in}}%
\pgfpathlineto{\pgfqpoint{1.386581in}{1.451241in}}%
\pgfpathlineto{\pgfqpoint{1.399163in}{1.444128in}}%
\pgfpathlineto{\pgfqpoint{1.411689in}{1.432679in}}%
\pgfpathlineto{\pgfqpoint{1.424163in}{1.426177in}}%
\pgfpathlineto{\pgfqpoint{1.436586in}{1.424874in}}%
\pgfpathlineto{\pgfqpoint{1.461293in}{1.426384in}}%
\pgfpathlineto{\pgfqpoint{1.510201in}{1.429810in}}%
\pgfpathlineto{\pgfqpoint{1.522333in}{1.428961in}}%
\pgfpathlineto{\pgfqpoint{1.534430in}{1.426184in}}%
\pgfpathlineto{\pgfqpoint{1.546495in}{1.417917in}}%
\pgfpathlineto{\pgfqpoint{1.570530in}{1.413904in}}%
\pgfpathlineto{\pgfqpoint{1.582504in}{1.412259in}}%
\pgfpathlineto{\pgfqpoint{1.594450in}{1.409827in}}%
\pgfpathlineto{\pgfqpoint{1.618263in}{1.402054in}}%
\pgfpathlineto{\pgfqpoint{1.630133in}{1.399882in}}%
\pgfpathlineto{\pgfqpoint{1.641980in}{1.399312in}}%
\pgfpathlineto{\pgfqpoint{1.653804in}{1.399400in}}%
\pgfpathlineto{\pgfqpoint{1.665607in}{1.400211in}}%
\pgfpathlineto{\pgfqpoint{1.677389in}{1.402899in}}%
\pgfpathlineto{\pgfqpoint{1.689152in}{1.408300in}}%
\pgfpathlineto{\pgfqpoint{1.700896in}{1.415325in}}%
\pgfpathlineto{\pgfqpoint{1.712622in}{1.420679in}}%
\pgfpathlineto{\pgfqpoint{1.724330in}{1.420809in}}%
\pgfpathlineto{\pgfqpoint{1.736022in}{1.414993in}}%
\pgfpathlineto{\pgfqpoint{1.747699in}{1.405533in}}%
\pgfpathlineto{\pgfqpoint{1.782638in}{1.373728in}}%
\pgfpathlineto{\pgfqpoint{1.794257in}{1.366914in}}%
\pgfpathlineto{\pgfqpoint{1.805863in}{1.364609in}}%
\pgfpathlineto{\pgfqpoint{1.817457in}{1.368033in}}%
\pgfpathlineto{\pgfqpoint{1.829039in}{1.377440in}}%
\pgfpathlineto{\pgfqpoint{1.840609in}{1.391400in}}%
\pgfpathlineto{\pgfqpoint{1.852169in}{1.407603in}}%
\pgfpathlineto{\pgfqpoint{1.863718in}{1.422369in}}%
\pgfpathlineto{\pgfqpoint{1.875257in}{1.431971in}}%
\pgfpathlineto{\pgfqpoint{1.898306in}{1.443271in}}%
\pgfpathlineto{\pgfqpoint{1.909817in}{1.450311in}}%
\pgfpathlineto{\pgfqpoint{1.921320in}{1.455304in}}%
\pgfpathlineto{\pgfqpoint{1.932814in}{1.455970in}}%
\pgfpathlineto{\pgfqpoint{1.944301in}{1.451603in}}%
\pgfpathlineto{\pgfqpoint{1.955780in}{1.441212in}}%
\pgfpathlineto{\pgfqpoint{1.967251in}{1.425355in}}%
\pgfpathlineto{\pgfqpoint{1.978716in}{1.407914in}}%
\pgfpathlineto{\pgfqpoint{1.990174in}{1.394462in}}%
\pgfpathlineto{\pgfqpoint{2.001625in}{1.390410in}}%
\pgfpathlineto{\pgfqpoint{2.013070in}{1.398666in}}%
\pgfpathlineto{\pgfqpoint{2.024510in}{1.416804in}}%
\pgfpathlineto{\pgfqpoint{2.035944in}{1.438403in}}%
\pgfpathlineto{\pgfqpoint{2.047372in}{1.455906in}}%
\pgfpathlineto{\pgfqpoint{2.058795in}{1.462154in}}%
\pgfpathlineto{\pgfqpoint{2.070213in}{1.454943in}}%
\pgfpathlineto{\pgfqpoint{2.093034in}{1.423682in}}%
\pgfpathlineto{\pgfqpoint{2.104438in}{1.413224in}}%
\pgfpathlineto{\pgfqpoint{2.115837in}{1.409122in}}%
\pgfpathlineto{\pgfqpoint{2.127233in}{1.408949in}}%
\pgfpathlineto{\pgfqpoint{2.150012in}{1.411485in}}%
\pgfpathlineto{\pgfqpoint{2.172776in}{1.415381in}}%
\pgfpathlineto{\pgfqpoint{2.184153in}{1.417801in}}%
\pgfpathlineto{\pgfqpoint{2.195527in}{1.421243in}}%
\pgfpathlineto{\pgfqpoint{2.206897in}{1.423598in}}%
\pgfpathlineto{\pgfqpoint{2.218265in}{1.420614in}}%
\pgfpathlineto{\pgfqpoint{2.229630in}{1.412780in}}%
\pgfpathlineto{\pgfqpoint{2.240991in}{1.406848in}}%
\pgfpathlineto{\pgfqpoint{2.252347in}{1.406518in}}%
\pgfpathlineto{\pgfqpoint{2.263686in}{1.408446in}}%
\pgfpathlineto{\pgfqpoint{2.274973in}{1.411160in}}%
\pgfpathlineto{\pgfqpoint{2.286134in}{1.417820in}}%
\pgfpathlineto{\pgfqpoint{2.297074in}{1.426978in}}%
\pgfpathlineto{\pgfqpoint{2.307710in}{1.431140in}}%
\pgfpathlineto{\pgfqpoint{2.318002in}{1.427611in}}%
\pgfpathlineto{\pgfqpoint{2.327947in}{1.421881in}}%
\pgfpathlineto{\pgfqpoint{2.337562in}{1.419088in}}%
\pgfpathlineto{\pgfqpoint{2.346867in}{1.419981in}}%
\pgfpathlineto{\pgfqpoint{2.355880in}{1.423497in}}%
\pgfpathlineto{\pgfqpoint{2.373102in}{1.431663in}}%
\pgfpathlineto{\pgfqpoint{2.389351in}{1.438602in}}%
\pgfpathlineto{\pgfqpoint{2.397144in}{1.440402in}}%
\pgfpathlineto{\pgfqpoint{2.404732in}{1.437686in}}%
\pgfpathlineto{\pgfqpoint{2.412125in}{1.429548in}}%
\pgfpathlineto{\pgfqpoint{2.419333in}{1.419955in}}%
\pgfpathlineto{\pgfqpoint{2.426365in}{1.416086in}}%
\pgfpathlineto{\pgfqpoint{2.433230in}{1.420591in}}%
\pgfpathlineto{\pgfqpoint{2.439934in}{1.427442in}}%
\pgfpathlineto{\pgfqpoint{2.446486in}{1.428718in}}%
\pgfpathlineto{\pgfqpoint{2.452892in}{1.422746in}}%
\pgfpathlineto{\pgfqpoint{2.459159in}{1.414510in}}%
\pgfpathlineto{\pgfqpoint{2.465292in}{1.410185in}}%
\pgfpathlineto{\pgfqpoint{2.471297in}{1.412012in}}%
\pgfpathlineto{\pgfqpoint{2.482944in}{1.422391in}}%
\pgfpathlineto{\pgfqpoint{2.488596in}{1.424709in}}%
\pgfpathlineto{\pgfqpoint{2.494139in}{1.425357in}}%
\pgfpathlineto{\pgfqpoint{2.499577in}{1.424844in}}%
\pgfpathlineto{\pgfqpoint{2.510155in}{1.420497in}}%
\pgfpathlineto{\pgfqpoint{2.515301in}{1.420650in}}%
\pgfpathlineto{\pgfqpoint{2.520357in}{1.421652in}}%
\pgfpathlineto{\pgfqpoint{2.525326in}{1.420466in}}%
\pgfpathlineto{\pgfqpoint{2.530211in}{1.418060in}}%
\pgfpathlineto{\pgfqpoint{2.535014in}{1.418086in}}%
\pgfpathlineto{\pgfqpoint{2.544386in}{1.425385in}}%
\pgfpathlineto{\pgfqpoint{2.548960in}{1.426283in}}%
\pgfpathlineto{\pgfqpoint{2.553463in}{1.423959in}}%
\pgfpathlineto{\pgfqpoint{2.557896in}{1.420728in}}%
\pgfpathlineto{\pgfqpoint{2.562262in}{1.418873in}}%
\pgfpathlineto{\pgfqpoint{2.566563in}{1.419873in}}%
\pgfpathlineto{\pgfqpoint{2.570801in}{1.424776in}}%
\pgfpathlineto{\pgfqpoint{2.574977in}{1.432285in}}%
\pgfpathlineto{\pgfqpoint{2.579093in}{1.437370in}}%
\pgfpathlineto{\pgfqpoint{2.583152in}{1.435632in}}%
\pgfpathlineto{\pgfqpoint{2.591101in}{1.420710in}}%
\pgfpathlineto{\pgfqpoint{2.594995in}{1.416347in}}%
\pgfpathlineto{\pgfqpoint{2.598836in}{1.416287in}}%
\pgfpathlineto{\pgfqpoint{2.602628in}{1.418703in}}%
\pgfpathlineto{\pgfqpoint{2.606370in}{1.419846in}}%
\pgfpathlineto{\pgfqpoint{2.610063in}{1.417858in}}%
\pgfpathlineto{\pgfqpoint{2.613711in}{1.414838in}}%
\pgfpathlineto{\pgfqpoint{2.617312in}{1.413715in}}%
\pgfpathlineto{\pgfqpoint{2.620869in}{1.415320in}}%
\pgfpathlineto{\pgfqpoint{2.624383in}{1.419217in}}%
\pgfpathlineto{\pgfqpoint{2.627854in}{1.424593in}}%
\pgfpathlineto{\pgfqpoint{2.631284in}{1.428856in}}%
\pgfpathlineto{\pgfqpoint{2.634673in}{1.428268in}}%
\pgfpathlineto{\pgfqpoint{2.641335in}{1.415571in}}%
\pgfpathlineto{\pgfqpoint{2.644609in}{1.412662in}}%
\pgfpathlineto{\pgfqpoint{2.651047in}{1.414773in}}%
\pgfpathlineto{\pgfqpoint{2.654213in}{1.413148in}}%
\pgfpathlineto{\pgfqpoint{2.657345in}{1.409741in}}%
\pgfpathlineto{\pgfqpoint{2.660442in}{1.407687in}}%
\pgfpathlineto{\pgfqpoint{2.663507in}{1.408822in}}%
\pgfpathlineto{\pgfqpoint{2.666540in}{1.412791in}}%
\pgfpathlineto{\pgfqpoint{2.672511in}{1.423882in}}%
\pgfpathlineto{\pgfqpoint{2.675450in}{1.426751in}}%
\pgfpathlineto{\pgfqpoint{2.678360in}{1.424948in}}%
\pgfpathlineto{\pgfqpoint{2.684093in}{1.412657in}}%
\pgfpathlineto{\pgfqpoint{2.686918in}{1.410458in}}%
\pgfpathlineto{\pgfqpoint{2.689714in}{1.412712in}}%
\pgfpathlineto{\pgfqpoint{2.692484in}{1.416488in}}%
\pgfpathlineto{\pgfqpoint{2.695228in}{1.419123in}}%
\pgfpathlineto{\pgfqpoint{2.697946in}{1.419374in}}%
\pgfpathlineto{\pgfqpoint{2.700638in}{1.417012in}}%
\pgfpathlineto{\pgfqpoint{2.705948in}{1.409908in}}%
\pgfpathlineto{\pgfqpoint{2.708567in}{1.408919in}}%
\pgfpathlineto{\pgfqpoint{2.711162in}{1.410614in}}%
\pgfpathlineto{\pgfqpoint{2.718810in}{1.422121in}}%
\pgfpathlineto{\pgfqpoint{2.721315in}{1.423207in}}%
\pgfpathlineto{\pgfqpoint{2.728702in}{1.422662in}}%
\pgfpathlineto{\pgfqpoint{2.733523in}{1.419917in}}%
\pgfpathlineto{\pgfqpoint{2.735903in}{1.421070in}}%
\pgfpathlineto{\pgfqpoint{2.740606in}{1.426838in}}%
\pgfpathlineto{\pgfqpoint{2.742929in}{1.425859in}}%
\pgfpathlineto{\pgfqpoint{2.747520in}{1.418937in}}%
\pgfpathlineto{\pgfqpoint{2.749788in}{1.418860in}}%
\pgfpathlineto{\pgfqpoint{2.752038in}{1.420459in}}%
\pgfpathlineto{\pgfqpoint{2.754271in}{1.420748in}}%
\pgfpathlineto{\pgfqpoint{2.760868in}{1.415551in}}%
\pgfpathlineto{\pgfqpoint{2.763034in}{1.412761in}}%
\pgfpathlineto{\pgfqpoint{2.767317in}{1.404176in}}%
\pgfpathlineto{\pgfqpoint{2.769435in}{1.404047in}}%
\pgfpathlineto{\pgfqpoint{2.775697in}{1.416545in}}%
\pgfpathlineto{\pgfqpoint{2.777754in}{1.417728in}}%
\pgfpathlineto{\pgfqpoint{2.779797in}{1.418023in}}%
\pgfpathlineto{\pgfqpoint{2.785840in}{1.415670in}}%
\pgfpathlineto{\pgfqpoint{2.793705in}{1.420790in}}%
\pgfpathlineto{\pgfqpoint{2.795638in}{1.419536in}}%
\pgfpathlineto{\pgfqpoint{2.805113in}{1.407006in}}%
\pgfpathlineto{\pgfqpoint{2.806972in}{1.405672in}}%
\pgfpathlineto{\pgfqpoint{2.808818in}{1.407453in}}%
\pgfpathlineto{\pgfqpoint{2.810653in}{1.410549in}}%
\pgfpathlineto{\pgfqpoint{2.812476in}{1.410805in}}%
\pgfpathlineto{\pgfqpoint{2.816087in}{1.404430in}}%
\pgfpathlineto{\pgfqpoint{2.817876in}{1.406666in}}%
\pgfpathlineto{\pgfqpoint{2.821421in}{1.416439in}}%
\pgfpathlineto{\pgfqpoint{2.823177in}{1.416462in}}%
\pgfpathlineto{\pgfqpoint{2.826658in}{1.410212in}}%
\pgfpathlineto{\pgfqpoint{2.828383in}{1.409134in}}%
\pgfpathlineto{\pgfqpoint{2.830097in}{1.409796in}}%
\pgfpathlineto{\pgfqpoint{2.836854in}{1.417283in}}%
\pgfpathlineto{\pgfqpoint{2.840175in}{1.425714in}}%
\pgfpathlineto{\pgfqpoint{2.841820in}{1.426635in}}%
\pgfpathlineto{\pgfqpoint{2.843457in}{1.423600in}}%
\pgfpathlineto{\pgfqpoint{2.845084in}{1.419382in}}%
\pgfpathlineto{\pgfqpoint{2.846702in}{1.417498in}}%
\pgfpathlineto{\pgfqpoint{2.851502in}{1.420590in}}%
\pgfpathlineto{\pgfqpoint{2.854659in}{1.420734in}}%
\pgfpathlineto{\pgfqpoint{2.857781in}{1.422675in}}%
\pgfpathlineto{\pgfqpoint{2.859329in}{1.422557in}}%
\pgfpathlineto{\pgfqpoint{2.863925in}{1.416317in}}%
\pgfpathlineto{\pgfqpoint{2.866949in}{1.420883in}}%
\pgfpathlineto{\pgfqpoint{2.868449in}{1.420538in}}%
\pgfpathlineto{\pgfqpoint{2.871426in}{1.412333in}}%
\pgfpathlineto{\pgfqpoint{2.872903in}{1.410978in}}%
\pgfpathlineto{\pgfqpoint{2.874372in}{1.412917in}}%
\pgfpathlineto{\pgfqpoint{2.877289in}{1.418420in}}%
\pgfpathlineto{\pgfqpoint{2.878736in}{1.419536in}}%
\pgfpathlineto{\pgfqpoint{2.881609in}{1.419992in}}%
\pgfpathlineto{\pgfqpoint{2.884454in}{1.419781in}}%
\pgfpathlineto{\pgfqpoint{2.887270in}{1.417180in}}%
\pgfpathlineto{\pgfqpoint{2.892822in}{1.411427in}}%
\pgfpathlineto{\pgfqpoint{2.894194in}{1.412822in}}%
\pgfpathlineto{\pgfqpoint{2.896917in}{1.417155in}}%
\pgfpathlineto{\pgfqpoint{2.898269in}{1.417901in}}%
\pgfpathlineto{\pgfqpoint{2.902288in}{1.417607in}}%
\pgfpathlineto{\pgfqpoint{2.906251in}{1.419141in}}%
\pgfpathlineto{\pgfqpoint{2.908863in}{1.422276in}}%
\pgfpathlineto{\pgfqpoint{2.910161in}{1.422248in}}%
\pgfpathlineto{\pgfqpoint{2.912738in}{1.419862in}}%
\pgfpathlineto{\pgfqpoint{2.916561in}{1.414927in}}%
\pgfpathlineto{\pgfqpoint{2.917824in}{1.414955in}}%
\pgfpathlineto{\pgfqpoint{2.921581in}{1.417041in}}%
\pgfpathlineto{\pgfqpoint{2.924058in}{1.417045in}}%
\pgfpathlineto{\pgfqpoint{2.926514in}{1.418476in}}%
\pgfpathlineto{\pgfqpoint{2.927735in}{1.418177in}}%
\pgfpathlineto{\pgfqpoint{2.930160in}{1.416880in}}%
\pgfpathlineto{\pgfqpoint{2.932565in}{1.418062in}}%
\pgfpathlineto{\pgfqpoint{2.933760in}{1.417858in}}%
\pgfpathlineto{\pgfqpoint{2.937315in}{1.413903in}}%
\pgfpathlineto{\pgfqpoint{2.939661in}{1.414804in}}%
\pgfpathlineto{\pgfqpoint{2.945445in}{1.417360in}}%
\pgfpathlineto{\pgfqpoint{2.947726in}{1.418269in}}%
\pgfpathlineto{\pgfqpoint{2.949989in}{1.420134in}}%
\pgfpathlineto{\pgfqpoint{2.951114in}{1.420085in}}%
\pgfpathlineto{\pgfqpoint{2.953351in}{1.418371in}}%
\pgfpathlineto{\pgfqpoint{2.956674in}{1.418558in}}%
\pgfpathlineto{\pgfqpoint{2.958868in}{1.415094in}}%
\pgfpathlineto{\pgfqpoint{2.959959in}{1.414523in}}%
\pgfpathlineto{\pgfqpoint{2.966418in}{1.418031in}}%
\pgfpathlineto{\pgfqpoint{2.967481in}{1.417322in}}%
\pgfpathlineto{\pgfqpoint{2.969595in}{1.414652in}}%
\pgfpathlineto{\pgfqpoint{2.972736in}{1.414361in}}%
\pgfpathlineto{\pgfqpoint{2.973776in}{1.413578in}}%
\pgfpathlineto{\pgfqpoint{2.977898in}{1.413798in}}%
\pgfpathlineto{\pgfqpoint{2.978919in}{1.415391in}}%
\pgfpathlineto{\pgfqpoint{2.980951in}{1.420129in}}%
\pgfpathlineto{\pgfqpoint{2.981961in}{1.419786in}}%
\pgfpathlineto{\pgfqpoint{2.983971in}{1.414968in}}%
\pgfpathlineto{\pgfqpoint{2.984971in}{1.414515in}}%
\pgfpathlineto{\pgfqpoint{2.985968in}{1.416423in}}%
\pgfpathlineto{\pgfqpoint{2.987950in}{1.421934in}}%
\pgfpathlineto{\pgfqpoint{2.988937in}{1.422166in}}%
\pgfpathlineto{\pgfqpoint{2.992848in}{1.416926in}}%
\pgfpathlineto{\pgfqpoint{2.995747in}{1.415734in}}%
\pgfpathlineto{\pgfqpoint{2.997664in}{1.416574in}}%
\pgfpathlineto{\pgfqpoint{2.999568in}{1.418024in}}%
\pgfpathlineto{\pgfqpoint{3.000515in}{1.418175in}}%
\pgfpathlineto{\pgfqpoint{3.003339in}{1.416377in}}%
\pgfpathlineto{\pgfqpoint{3.005206in}{1.414724in}}%
\pgfpathlineto{\pgfqpoint{3.006135in}{1.414613in}}%
\pgfpathlineto{\pgfqpoint{3.008904in}{1.416501in}}%
\pgfpathlineto{\pgfqpoint{3.014363in}{1.420551in}}%
\pgfpathlineto{\pgfqpoint{3.015263in}{1.419997in}}%
\pgfpathlineto{\pgfqpoint{3.018835in}{1.415490in}}%
\pgfpathlineto{\pgfqpoint{3.020604in}{1.416923in}}%
\pgfpathlineto{\pgfqpoint{3.022363in}{1.418212in}}%
\pgfpathlineto{\pgfqpoint{3.024981in}{1.417709in}}%
\pgfpathlineto{\pgfqpoint{3.028435in}{1.416345in}}%
\pgfpathlineto{\pgfqpoint{3.031848in}{1.420106in}}%
\pgfpathlineto{\pgfqpoint{3.033540in}{1.417514in}}%
\pgfpathlineto{\pgfqpoint{3.035221in}{1.415844in}}%
\pgfpathlineto{\pgfqpoint{3.038555in}{1.417418in}}%
\pgfpathlineto{\pgfqpoint{3.041031in}{1.417208in}}%
\pgfpathlineto{\pgfqpoint{3.043110in}{1.417218in}}%
\pgfpathlineto{\pgfqpoint{3.043110in}{1.417218in}}%
\pgfusepath{stroke}%
\end{pgfscope}%
\begin{pgfscope}%
\pgfsetrectcap%
\pgfsetmiterjoin%
\pgfsetlinewidth{0.803000pt}%
\definecolor{currentstroke}{rgb}{0.000000,0.000000,0.000000}%
\pgfsetstrokecolor{currentstroke}%
\pgfsetdash{}{0pt}%
\pgfpathmoveto{\pgfqpoint{0.650000in}{0.420000in}}%
\pgfpathlineto{\pgfqpoint{0.650000in}{2.414488in}}%
\pgfusepath{stroke}%
\end{pgfscope}%
\begin{pgfscope}%
\pgfsetrectcap%
\pgfsetmiterjoin%
\pgfsetlinewidth{0.803000pt}%
\definecolor{currentstroke}{rgb}{0.000000,0.000000,0.000000}%
\pgfsetstrokecolor{currentstroke}%
\pgfsetdash{}{0pt}%
\pgfpathmoveto{\pgfqpoint{3.038110in}{0.420000in}}%
\pgfpathlineto{\pgfqpoint{3.038110in}{2.414488in}}%
\pgfusepath{stroke}%
\end{pgfscope}%
\begin{pgfscope}%
\pgfsetrectcap%
\pgfsetmiterjoin%
\pgfsetlinewidth{0.803000pt}%
\definecolor{currentstroke}{rgb}{0.000000,0.000000,0.000000}%
\pgfsetstrokecolor{currentstroke}%
\pgfsetdash{}{0pt}%
\pgfpathmoveto{\pgfqpoint{0.650000in}{0.420000in}}%
\pgfpathlineto{\pgfqpoint{3.038110in}{0.420000in}}%
\pgfusepath{stroke}%
\end{pgfscope}%
\begin{pgfscope}%
\pgfsetrectcap%
\pgfsetmiterjoin%
\pgfsetlinewidth{0.803000pt}%
\definecolor{currentstroke}{rgb}{0.000000,0.000000,0.000000}%
\pgfsetstrokecolor{currentstroke}%
\pgfsetdash{}{0pt}%
\pgfpathmoveto{\pgfqpoint{0.650000in}{2.414488in}}%
\pgfpathlineto{\pgfqpoint{3.038110in}{2.414488in}}%
\pgfusepath{stroke}%
\end{pgfscope}%
\begin{pgfscope}%
\pgfsetrectcap%
\pgfsetroundjoin%
\pgfsetlinewidth{0.602250pt}%
\definecolor{currentstroke}{rgb}{0.000000,0.000000,0.000000}%
\pgfsetstrokecolor{currentstroke}%
\pgfsetdash{}{0pt}%
\pgfpathmoveto{\pgfqpoint{2.127879in}{2.275599in}}%
\pgfpathlineto{\pgfqpoint{2.238990in}{2.275599in}}%
\pgfpathlineto{\pgfqpoint{2.350101in}{2.275599in}}%
\pgfusepath{stroke}%
\end{pgfscope}%
\begin{pgfscope}%
\definecolor{textcolor}{rgb}{0.000000,0.000000,0.000000}%
\pgfsetstrokecolor{textcolor}%
\pgfsetfillcolor{textcolor}%
\pgftext[x=2.438990in,y=2.236710in,left,base]{\color{textcolor}{\rmfamily\fontsize{8.000000}{9.600000}\selectfont\catcode`\^=\active\def^{\ifmmode\sp\else\^{}\fi}\catcode`\%=\active\def%{\%}ref}}%
\end{pgfscope}%
\begin{pgfscope}%
\pgfsetrectcap%
\pgfsetroundjoin%
\pgfsetlinewidth{1.003750pt}%
\definecolor{currentstroke}{rgb}{0.454902,0.678431,0.819608}%
\pgfsetstrokecolor{currentstroke}%
\pgfsetdash{}{0pt}%
\pgfpathmoveto{\pgfqpoint{2.127879in}{2.120661in}}%
\pgfpathlineto{\pgfqpoint{2.238990in}{2.120661in}}%
\pgfpathlineto{\pgfqpoint{2.350101in}{2.120661in}}%
\pgfusepath{stroke}%
\end{pgfscope}%
\begin{pgfscope}%
\definecolor{textcolor}{rgb}{0.000000,0.000000,0.000000}%
\pgfsetstrokecolor{textcolor}%
\pgfsetfillcolor{textcolor}%
\pgftext[x=2.438990in,y=2.081772in,left,base]{\color{textcolor}{\rmfamily\fontsize{8.000000}{9.600000}\selectfont\catcode`\^=\active\def^{\ifmmode\sp\else\^{}\fi}\catcode`\%=\active\def%{\%}base}}%
\end{pgfscope}%
\begin{pgfscope}%
\pgfsetrectcap%
\pgfsetroundjoin%
\pgfsetlinewidth{1.003750pt}%
\definecolor{currentstroke}{rgb}{0.501961,0.450980,0.674510}%
\pgfsetstrokecolor{currentstroke}%
\pgfsetdash{}{0pt}%
\pgfpathmoveto{\pgfqpoint{2.127879in}{1.965723in}}%
\pgfpathlineto{\pgfqpoint{2.238990in}{1.965723in}}%
\pgfpathlineto{\pgfqpoint{2.350101in}{1.965723in}}%
\pgfusepath{stroke}%
\end{pgfscope}%
\begin{pgfscope}%
\definecolor{textcolor}{rgb}{0.000000,0.000000,0.000000}%
\pgfsetstrokecolor{textcolor}%
\pgfsetfillcolor{textcolor}%
\pgftext[x=2.438990in,y=1.926834in,left,base]{\color{textcolor}{\rmfamily\fontsize{8.000000}{9.600000}\selectfont\catcode`\^=\active\def^{\ifmmode\sp\else\^{}\fi}\catcode`\%=\active\def%{\%}optimised}}%
\end{pgfscope}%
\end{pgfpicture}%
\makeatother%
\endgroup%

    \label{fig:R500deg18R100deg90_bud}
\end{subfigure}
\hfill
\begin{subfigure}[t]{0.495\textwidth}
    \centering
    \subcaption{}
    %% Creator: Matplotlib, PGF backend
%%
%% To include the figure in your LaTeX document, write
%%   \input{<filename>.pgf}
%%
%% Make sure the required packages are loaded in your preamble
%%   \usepackage{pgf}
%%
%% Also ensure that all the required font packages are loaded; for instance,
%% the lmodern package is sometimes necessary when using math font.
%%   \usepackage{lmodern}
%%
%% Figures using additional raster images can only be included by \input if
%% they are in the same directory as the main LaTeX file. For loading figures
%% from other directories you can use the `import` package
%%   \usepackage{import}
%%
%% and then include the figures with
%%   \import{<path to file>}{<filename>.pgf}
%%
%% Matplotlib used the following preamble
%%   \def\mathdefault#1{#1}
%%   \everymath=\expandafter{\the\everymath\displaystyle}
%%   \IfFileExists{scrextend.sty}{
%%     \usepackage[fontsize=11.000000pt]{scrextend}
%%   }{
%%     \renewcommand{\normalsize}{\fontsize{11.000000}{13.200000}\selectfont}
%%     \normalsize
%%   }
%%   
%%   \makeatletter\@ifpackageloaded{underscore}{}{\usepackage[strings]{underscore}}\makeatother
%%
\begingroup%
\makeatletter%
\begin{pgfpicture}%
\pgfpathrectangle{\pgfpointorigin}{\pgfqpoint{3.118110in}{2.494488in}}%
\pgfusepath{use as bounding box, clip}%
\begin{pgfscope}%
\pgfsetbuttcap%
\pgfsetmiterjoin%
\definecolor{currentfill}{rgb}{1.000000,1.000000,1.000000}%
\pgfsetfillcolor{currentfill}%
\pgfsetlinewidth{0.000000pt}%
\definecolor{currentstroke}{rgb}{1.000000,1.000000,1.000000}%
\pgfsetstrokecolor{currentstroke}%
\pgfsetdash{}{0pt}%
\pgfpathmoveto{\pgfqpoint{0.000000in}{0.000000in}}%
\pgfpathlineto{\pgfqpoint{3.118110in}{0.000000in}}%
\pgfpathlineto{\pgfqpoint{3.118110in}{2.494488in}}%
\pgfpathlineto{\pgfqpoint{0.000000in}{2.494488in}}%
\pgfpathlineto{\pgfqpoint{0.000000in}{0.000000in}}%
\pgfpathclose%
\pgfusepath{fill}%
\end{pgfscope}%
\begin{pgfscope}%
\pgfsetbuttcap%
\pgfsetmiterjoin%
\definecolor{currentfill}{rgb}{1.000000,1.000000,1.000000}%
\pgfsetfillcolor{currentfill}%
\pgfsetlinewidth{0.000000pt}%
\definecolor{currentstroke}{rgb}{0.000000,0.000000,0.000000}%
\pgfsetstrokecolor{currentstroke}%
\pgfsetstrokeopacity{0.000000}%
\pgfsetdash{}{0pt}%
\pgfpathmoveto{\pgfqpoint{0.650000in}{0.420000in}}%
\pgfpathlineto{\pgfqpoint{3.038110in}{0.420000in}}%
\pgfpathlineto{\pgfqpoint{3.038110in}{2.414488in}}%
\pgfpathlineto{\pgfqpoint{0.650000in}{2.414488in}}%
\pgfpathlineto{\pgfqpoint{0.650000in}{0.420000in}}%
\pgfpathclose%
\pgfusepath{fill}%
\end{pgfscope}%
\begin{pgfscope}%
\pgfsetbuttcap%
\pgfsetroundjoin%
\definecolor{currentfill}{rgb}{0.000000,0.000000,0.000000}%
\pgfsetfillcolor{currentfill}%
\pgfsetlinewidth{0.501875pt}%
\definecolor{currentstroke}{rgb}{0.000000,0.000000,0.000000}%
\pgfsetstrokecolor{currentstroke}%
\pgfsetdash{}{0pt}%
\pgfsys@defobject{currentmarker}{\pgfqpoint{0.000000in}{0.000000in}}{\pgfqpoint{0.000000in}{0.041667in}}{%
\pgfpathmoveto{\pgfqpoint{0.000000in}{0.000000in}}%
\pgfpathlineto{\pgfqpoint{0.000000in}{0.041667in}}%
\pgfusepath{stroke,fill}%
}%
\begin{pgfscope}%
\pgfsys@transformshift{0.650000in}{0.420000in}%
\pgfsys@useobject{currentmarker}{}%
\end{pgfscope}%
\end{pgfscope}%
\begin{pgfscope}%
\pgfsetbuttcap%
\pgfsetroundjoin%
\definecolor{currentfill}{rgb}{0.000000,0.000000,0.000000}%
\pgfsetfillcolor{currentfill}%
\pgfsetlinewidth{0.501875pt}%
\definecolor{currentstroke}{rgb}{0.000000,0.000000,0.000000}%
\pgfsetstrokecolor{currentstroke}%
\pgfsetdash{}{0pt}%
\pgfsys@defobject{currentmarker}{\pgfqpoint{0.000000in}{-0.041667in}}{\pgfqpoint{0.000000in}{0.000000in}}{%
\pgfpathmoveto{\pgfqpoint{0.000000in}{0.000000in}}%
\pgfpathlineto{\pgfqpoint{0.000000in}{-0.041667in}}%
\pgfusepath{stroke,fill}%
}%
\begin{pgfscope}%
\pgfsys@transformshift{0.650000in}{2.414488in}%
\pgfsys@useobject{currentmarker}{}%
\end{pgfscope}%
\end{pgfscope}%
\begin{pgfscope}%
\definecolor{textcolor}{rgb}{0.000000,0.000000,0.000000}%
\pgfsetstrokecolor{textcolor}%
\pgfsetfillcolor{textcolor}%
\pgftext[x=0.650000in,y=0.371389in,,top]{\color{textcolor}{\rmfamily\fontsize{11.000000}{13.200000}\selectfont\catcode`\^=\active\def^{\ifmmode\sp\else\^{}\fi}\catcode`\%=\active\def%{\%}0}}%
\end{pgfscope}%
\begin{pgfscope}%
\pgfsetbuttcap%
\pgfsetroundjoin%
\definecolor{currentfill}{rgb}{0.000000,0.000000,0.000000}%
\pgfsetfillcolor{currentfill}%
\pgfsetlinewidth{0.501875pt}%
\definecolor{currentstroke}{rgb}{0.000000,0.000000,0.000000}%
\pgfsetstrokecolor{currentstroke}%
\pgfsetdash{}{0pt}%
\pgfsys@defobject{currentmarker}{\pgfqpoint{0.000000in}{0.000000in}}{\pgfqpoint{0.000000in}{0.041667in}}{%
\pgfpathmoveto{\pgfqpoint{0.000000in}{0.000000in}}%
\pgfpathlineto{\pgfqpoint{0.000000in}{0.041667in}}%
\pgfusepath{stroke,fill}%
}%
\begin{pgfscope}%
\pgfsys@transformshift{1.493052in}{0.420000in}%
\pgfsys@useobject{currentmarker}{}%
\end{pgfscope}%
\end{pgfscope}%
\begin{pgfscope}%
\pgfsetbuttcap%
\pgfsetroundjoin%
\definecolor{currentfill}{rgb}{0.000000,0.000000,0.000000}%
\pgfsetfillcolor{currentfill}%
\pgfsetlinewidth{0.501875pt}%
\definecolor{currentstroke}{rgb}{0.000000,0.000000,0.000000}%
\pgfsetstrokecolor{currentstroke}%
\pgfsetdash{}{0pt}%
\pgfsys@defobject{currentmarker}{\pgfqpoint{0.000000in}{-0.041667in}}{\pgfqpoint{0.000000in}{0.000000in}}{%
\pgfpathmoveto{\pgfqpoint{0.000000in}{0.000000in}}%
\pgfpathlineto{\pgfqpoint{0.000000in}{-0.041667in}}%
\pgfusepath{stroke,fill}%
}%
\begin{pgfscope}%
\pgfsys@transformshift{1.493052in}{2.414488in}%
\pgfsys@useobject{currentmarker}{}%
\end{pgfscope}%
\end{pgfscope}%
\begin{pgfscope}%
\definecolor{textcolor}{rgb}{0.000000,0.000000,0.000000}%
\pgfsetstrokecolor{textcolor}%
\pgfsetfillcolor{textcolor}%
\pgftext[x=1.493052in,y=0.371389in,,top]{\color{textcolor}{\rmfamily\fontsize{11.000000}{13.200000}\selectfont\catcode`\^=\active\def^{\ifmmode\sp\else\^{}\fi}\catcode`\%=\active\def%{\%}100}}%
\end{pgfscope}%
\begin{pgfscope}%
\pgfsetbuttcap%
\pgfsetroundjoin%
\definecolor{currentfill}{rgb}{0.000000,0.000000,0.000000}%
\pgfsetfillcolor{currentfill}%
\pgfsetlinewidth{0.501875pt}%
\definecolor{currentstroke}{rgb}{0.000000,0.000000,0.000000}%
\pgfsetstrokecolor{currentstroke}%
\pgfsetdash{}{0pt}%
\pgfsys@defobject{currentmarker}{\pgfqpoint{0.000000in}{0.000000in}}{\pgfqpoint{0.000000in}{0.041667in}}{%
\pgfpathmoveto{\pgfqpoint{0.000000in}{0.000000in}}%
\pgfpathlineto{\pgfqpoint{0.000000in}{0.041667in}}%
\pgfusepath{stroke,fill}%
}%
\begin{pgfscope}%
\pgfsys@transformshift{2.336103in}{0.420000in}%
\pgfsys@useobject{currentmarker}{}%
\end{pgfscope}%
\end{pgfscope}%
\begin{pgfscope}%
\pgfsetbuttcap%
\pgfsetroundjoin%
\definecolor{currentfill}{rgb}{0.000000,0.000000,0.000000}%
\pgfsetfillcolor{currentfill}%
\pgfsetlinewidth{0.501875pt}%
\definecolor{currentstroke}{rgb}{0.000000,0.000000,0.000000}%
\pgfsetstrokecolor{currentstroke}%
\pgfsetdash{}{0pt}%
\pgfsys@defobject{currentmarker}{\pgfqpoint{0.000000in}{-0.041667in}}{\pgfqpoint{0.000000in}{0.000000in}}{%
\pgfpathmoveto{\pgfqpoint{0.000000in}{0.000000in}}%
\pgfpathlineto{\pgfqpoint{0.000000in}{-0.041667in}}%
\pgfusepath{stroke,fill}%
}%
\begin{pgfscope}%
\pgfsys@transformshift{2.336103in}{2.414488in}%
\pgfsys@useobject{currentmarker}{}%
\end{pgfscope}%
\end{pgfscope}%
\begin{pgfscope}%
\definecolor{textcolor}{rgb}{0.000000,0.000000,0.000000}%
\pgfsetstrokecolor{textcolor}%
\pgfsetfillcolor{textcolor}%
\pgftext[x=2.336103in,y=0.371389in,,top]{\color{textcolor}{\rmfamily\fontsize{11.000000}{13.200000}\selectfont\catcode`\^=\active\def^{\ifmmode\sp\else\^{}\fi}\catcode`\%=\active\def%{\%}200}}%
\end{pgfscope}%
\begin{pgfscope}%
\pgfsetbuttcap%
\pgfsetroundjoin%
\definecolor{currentfill}{rgb}{0.000000,0.000000,0.000000}%
\pgfsetfillcolor{currentfill}%
\pgfsetlinewidth{0.401500pt}%
\definecolor{currentstroke}{rgb}{0.000000,0.000000,0.000000}%
\pgfsetstrokecolor{currentstroke}%
\pgfsetdash{}{0pt}%
\pgfsys@defobject{currentmarker}{\pgfqpoint{0.000000in}{0.000000in}}{\pgfqpoint{0.000000in}{0.020833in}}{%
\pgfpathmoveto{\pgfqpoint{0.000000in}{0.000000in}}%
\pgfpathlineto{\pgfqpoint{0.000000in}{0.020833in}}%
\pgfusepath{stroke,fill}%
}%
\begin{pgfscope}%
\pgfsys@transformshift{0.818610in}{0.420000in}%
\pgfsys@useobject{currentmarker}{}%
\end{pgfscope}%
\end{pgfscope}%
\begin{pgfscope}%
\pgfsetbuttcap%
\pgfsetroundjoin%
\definecolor{currentfill}{rgb}{0.000000,0.000000,0.000000}%
\pgfsetfillcolor{currentfill}%
\pgfsetlinewidth{0.401500pt}%
\definecolor{currentstroke}{rgb}{0.000000,0.000000,0.000000}%
\pgfsetstrokecolor{currentstroke}%
\pgfsetdash{}{0pt}%
\pgfsys@defobject{currentmarker}{\pgfqpoint{0.000000in}{-0.020833in}}{\pgfqpoint{0.000000in}{0.000000in}}{%
\pgfpathmoveto{\pgfqpoint{0.000000in}{0.000000in}}%
\pgfpathlineto{\pgfqpoint{0.000000in}{-0.020833in}}%
\pgfusepath{stroke,fill}%
}%
\begin{pgfscope}%
\pgfsys@transformshift{0.818610in}{2.414488in}%
\pgfsys@useobject{currentmarker}{}%
\end{pgfscope}%
\end{pgfscope}%
\begin{pgfscope}%
\pgfsetbuttcap%
\pgfsetroundjoin%
\definecolor{currentfill}{rgb}{0.000000,0.000000,0.000000}%
\pgfsetfillcolor{currentfill}%
\pgfsetlinewidth{0.401500pt}%
\definecolor{currentstroke}{rgb}{0.000000,0.000000,0.000000}%
\pgfsetstrokecolor{currentstroke}%
\pgfsetdash{}{0pt}%
\pgfsys@defobject{currentmarker}{\pgfqpoint{0.000000in}{0.000000in}}{\pgfqpoint{0.000000in}{0.020833in}}{%
\pgfpathmoveto{\pgfqpoint{0.000000in}{0.000000in}}%
\pgfpathlineto{\pgfqpoint{0.000000in}{0.020833in}}%
\pgfusepath{stroke,fill}%
}%
\begin{pgfscope}%
\pgfsys@transformshift{0.987221in}{0.420000in}%
\pgfsys@useobject{currentmarker}{}%
\end{pgfscope}%
\end{pgfscope}%
\begin{pgfscope}%
\pgfsetbuttcap%
\pgfsetroundjoin%
\definecolor{currentfill}{rgb}{0.000000,0.000000,0.000000}%
\pgfsetfillcolor{currentfill}%
\pgfsetlinewidth{0.401500pt}%
\definecolor{currentstroke}{rgb}{0.000000,0.000000,0.000000}%
\pgfsetstrokecolor{currentstroke}%
\pgfsetdash{}{0pt}%
\pgfsys@defobject{currentmarker}{\pgfqpoint{0.000000in}{-0.020833in}}{\pgfqpoint{0.000000in}{0.000000in}}{%
\pgfpathmoveto{\pgfqpoint{0.000000in}{0.000000in}}%
\pgfpathlineto{\pgfqpoint{0.000000in}{-0.020833in}}%
\pgfusepath{stroke,fill}%
}%
\begin{pgfscope}%
\pgfsys@transformshift{0.987221in}{2.414488in}%
\pgfsys@useobject{currentmarker}{}%
\end{pgfscope}%
\end{pgfscope}%
\begin{pgfscope}%
\pgfsetbuttcap%
\pgfsetroundjoin%
\definecolor{currentfill}{rgb}{0.000000,0.000000,0.000000}%
\pgfsetfillcolor{currentfill}%
\pgfsetlinewidth{0.401500pt}%
\definecolor{currentstroke}{rgb}{0.000000,0.000000,0.000000}%
\pgfsetstrokecolor{currentstroke}%
\pgfsetdash{}{0pt}%
\pgfsys@defobject{currentmarker}{\pgfqpoint{0.000000in}{0.000000in}}{\pgfqpoint{0.000000in}{0.020833in}}{%
\pgfpathmoveto{\pgfqpoint{0.000000in}{0.000000in}}%
\pgfpathlineto{\pgfqpoint{0.000000in}{0.020833in}}%
\pgfusepath{stroke,fill}%
}%
\begin{pgfscope}%
\pgfsys@transformshift{1.155831in}{0.420000in}%
\pgfsys@useobject{currentmarker}{}%
\end{pgfscope}%
\end{pgfscope}%
\begin{pgfscope}%
\pgfsetbuttcap%
\pgfsetroundjoin%
\definecolor{currentfill}{rgb}{0.000000,0.000000,0.000000}%
\pgfsetfillcolor{currentfill}%
\pgfsetlinewidth{0.401500pt}%
\definecolor{currentstroke}{rgb}{0.000000,0.000000,0.000000}%
\pgfsetstrokecolor{currentstroke}%
\pgfsetdash{}{0pt}%
\pgfsys@defobject{currentmarker}{\pgfqpoint{0.000000in}{-0.020833in}}{\pgfqpoint{0.000000in}{0.000000in}}{%
\pgfpathmoveto{\pgfqpoint{0.000000in}{0.000000in}}%
\pgfpathlineto{\pgfqpoint{0.000000in}{-0.020833in}}%
\pgfusepath{stroke,fill}%
}%
\begin{pgfscope}%
\pgfsys@transformshift{1.155831in}{2.414488in}%
\pgfsys@useobject{currentmarker}{}%
\end{pgfscope}%
\end{pgfscope}%
\begin{pgfscope}%
\pgfsetbuttcap%
\pgfsetroundjoin%
\definecolor{currentfill}{rgb}{0.000000,0.000000,0.000000}%
\pgfsetfillcolor{currentfill}%
\pgfsetlinewidth{0.401500pt}%
\definecolor{currentstroke}{rgb}{0.000000,0.000000,0.000000}%
\pgfsetstrokecolor{currentstroke}%
\pgfsetdash{}{0pt}%
\pgfsys@defobject{currentmarker}{\pgfqpoint{0.000000in}{0.000000in}}{\pgfqpoint{0.000000in}{0.020833in}}{%
\pgfpathmoveto{\pgfqpoint{0.000000in}{0.000000in}}%
\pgfpathlineto{\pgfqpoint{0.000000in}{0.020833in}}%
\pgfusepath{stroke,fill}%
}%
\begin{pgfscope}%
\pgfsys@transformshift{1.324441in}{0.420000in}%
\pgfsys@useobject{currentmarker}{}%
\end{pgfscope}%
\end{pgfscope}%
\begin{pgfscope}%
\pgfsetbuttcap%
\pgfsetroundjoin%
\definecolor{currentfill}{rgb}{0.000000,0.000000,0.000000}%
\pgfsetfillcolor{currentfill}%
\pgfsetlinewidth{0.401500pt}%
\definecolor{currentstroke}{rgb}{0.000000,0.000000,0.000000}%
\pgfsetstrokecolor{currentstroke}%
\pgfsetdash{}{0pt}%
\pgfsys@defobject{currentmarker}{\pgfqpoint{0.000000in}{-0.020833in}}{\pgfqpoint{0.000000in}{0.000000in}}{%
\pgfpathmoveto{\pgfqpoint{0.000000in}{0.000000in}}%
\pgfpathlineto{\pgfqpoint{0.000000in}{-0.020833in}}%
\pgfusepath{stroke,fill}%
}%
\begin{pgfscope}%
\pgfsys@transformshift{1.324441in}{2.414488in}%
\pgfsys@useobject{currentmarker}{}%
\end{pgfscope}%
\end{pgfscope}%
\begin{pgfscope}%
\pgfsetbuttcap%
\pgfsetroundjoin%
\definecolor{currentfill}{rgb}{0.000000,0.000000,0.000000}%
\pgfsetfillcolor{currentfill}%
\pgfsetlinewidth{0.401500pt}%
\definecolor{currentstroke}{rgb}{0.000000,0.000000,0.000000}%
\pgfsetstrokecolor{currentstroke}%
\pgfsetdash{}{0pt}%
\pgfsys@defobject{currentmarker}{\pgfqpoint{0.000000in}{0.000000in}}{\pgfqpoint{0.000000in}{0.020833in}}{%
\pgfpathmoveto{\pgfqpoint{0.000000in}{0.000000in}}%
\pgfpathlineto{\pgfqpoint{0.000000in}{0.020833in}}%
\pgfusepath{stroke,fill}%
}%
\begin{pgfscope}%
\pgfsys@transformshift{1.661662in}{0.420000in}%
\pgfsys@useobject{currentmarker}{}%
\end{pgfscope}%
\end{pgfscope}%
\begin{pgfscope}%
\pgfsetbuttcap%
\pgfsetroundjoin%
\definecolor{currentfill}{rgb}{0.000000,0.000000,0.000000}%
\pgfsetfillcolor{currentfill}%
\pgfsetlinewidth{0.401500pt}%
\definecolor{currentstroke}{rgb}{0.000000,0.000000,0.000000}%
\pgfsetstrokecolor{currentstroke}%
\pgfsetdash{}{0pt}%
\pgfsys@defobject{currentmarker}{\pgfqpoint{0.000000in}{-0.020833in}}{\pgfqpoint{0.000000in}{0.000000in}}{%
\pgfpathmoveto{\pgfqpoint{0.000000in}{0.000000in}}%
\pgfpathlineto{\pgfqpoint{0.000000in}{-0.020833in}}%
\pgfusepath{stroke,fill}%
}%
\begin{pgfscope}%
\pgfsys@transformshift{1.661662in}{2.414488in}%
\pgfsys@useobject{currentmarker}{}%
\end{pgfscope}%
\end{pgfscope}%
\begin{pgfscope}%
\pgfsetbuttcap%
\pgfsetroundjoin%
\definecolor{currentfill}{rgb}{0.000000,0.000000,0.000000}%
\pgfsetfillcolor{currentfill}%
\pgfsetlinewidth{0.401500pt}%
\definecolor{currentstroke}{rgb}{0.000000,0.000000,0.000000}%
\pgfsetstrokecolor{currentstroke}%
\pgfsetdash{}{0pt}%
\pgfsys@defobject{currentmarker}{\pgfqpoint{0.000000in}{0.000000in}}{\pgfqpoint{0.000000in}{0.020833in}}{%
\pgfpathmoveto{\pgfqpoint{0.000000in}{0.000000in}}%
\pgfpathlineto{\pgfqpoint{0.000000in}{0.020833in}}%
\pgfusepath{stroke,fill}%
}%
\begin{pgfscope}%
\pgfsys@transformshift{1.830272in}{0.420000in}%
\pgfsys@useobject{currentmarker}{}%
\end{pgfscope}%
\end{pgfscope}%
\begin{pgfscope}%
\pgfsetbuttcap%
\pgfsetroundjoin%
\definecolor{currentfill}{rgb}{0.000000,0.000000,0.000000}%
\pgfsetfillcolor{currentfill}%
\pgfsetlinewidth{0.401500pt}%
\definecolor{currentstroke}{rgb}{0.000000,0.000000,0.000000}%
\pgfsetstrokecolor{currentstroke}%
\pgfsetdash{}{0pt}%
\pgfsys@defobject{currentmarker}{\pgfqpoint{0.000000in}{-0.020833in}}{\pgfqpoint{0.000000in}{0.000000in}}{%
\pgfpathmoveto{\pgfqpoint{0.000000in}{0.000000in}}%
\pgfpathlineto{\pgfqpoint{0.000000in}{-0.020833in}}%
\pgfusepath{stroke,fill}%
}%
\begin{pgfscope}%
\pgfsys@transformshift{1.830272in}{2.414488in}%
\pgfsys@useobject{currentmarker}{}%
\end{pgfscope}%
\end{pgfscope}%
\begin{pgfscope}%
\pgfsetbuttcap%
\pgfsetroundjoin%
\definecolor{currentfill}{rgb}{0.000000,0.000000,0.000000}%
\pgfsetfillcolor{currentfill}%
\pgfsetlinewidth{0.401500pt}%
\definecolor{currentstroke}{rgb}{0.000000,0.000000,0.000000}%
\pgfsetstrokecolor{currentstroke}%
\pgfsetdash{}{0pt}%
\pgfsys@defobject{currentmarker}{\pgfqpoint{0.000000in}{0.000000in}}{\pgfqpoint{0.000000in}{0.020833in}}{%
\pgfpathmoveto{\pgfqpoint{0.000000in}{0.000000in}}%
\pgfpathlineto{\pgfqpoint{0.000000in}{0.020833in}}%
\pgfusepath{stroke,fill}%
}%
\begin{pgfscope}%
\pgfsys@transformshift{1.998883in}{0.420000in}%
\pgfsys@useobject{currentmarker}{}%
\end{pgfscope}%
\end{pgfscope}%
\begin{pgfscope}%
\pgfsetbuttcap%
\pgfsetroundjoin%
\definecolor{currentfill}{rgb}{0.000000,0.000000,0.000000}%
\pgfsetfillcolor{currentfill}%
\pgfsetlinewidth{0.401500pt}%
\definecolor{currentstroke}{rgb}{0.000000,0.000000,0.000000}%
\pgfsetstrokecolor{currentstroke}%
\pgfsetdash{}{0pt}%
\pgfsys@defobject{currentmarker}{\pgfqpoint{0.000000in}{-0.020833in}}{\pgfqpoint{0.000000in}{0.000000in}}{%
\pgfpathmoveto{\pgfqpoint{0.000000in}{0.000000in}}%
\pgfpathlineto{\pgfqpoint{0.000000in}{-0.020833in}}%
\pgfusepath{stroke,fill}%
}%
\begin{pgfscope}%
\pgfsys@transformshift{1.998883in}{2.414488in}%
\pgfsys@useobject{currentmarker}{}%
\end{pgfscope}%
\end{pgfscope}%
\begin{pgfscope}%
\pgfsetbuttcap%
\pgfsetroundjoin%
\definecolor{currentfill}{rgb}{0.000000,0.000000,0.000000}%
\pgfsetfillcolor{currentfill}%
\pgfsetlinewidth{0.401500pt}%
\definecolor{currentstroke}{rgb}{0.000000,0.000000,0.000000}%
\pgfsetstrokecolor{currentstroke}%
\pgfsetdash{}{0pt}%
\pgfsys@defobject{currentmarker}{\pgfqpoint{0.000000in}{0.000000in}}{\pgfqpoint{0.000000in}{0.020833in}}{%
\pgfpathmoveto{\pgfqpoint{0.000000in}{0.000000in}}%
\pgfpathlineto{\pgfqpoint{0.000000in}{0.020833in}}%
\pgfusepath{stroke,fill}%
}%
\begin{pgfscope}%
\pgfsys@transformshift{2.167493in}{0.420000in}%
\pgfsys@useobject{currentmarker}{}%
\end{pgfscope}%
\end{pgfscope}%
\begin{pgfscope}%
\pgfsetbuttcap%
\pgfsetroundjoin%
\definecolor{currentfill}{rgb}{0.000000,0.000000,0.000000}%
\pgfsetfillcolor{currentfill}%
\pgfsetlinewidth{0.401500pt}%
\definecolor{currentstroke}{rgb}{0.000000,0.000000,0.000000}%
\pgfsetstrokecolor{currentstroke}%
\pgfsetdash{}{0pt}%
\pgfsys@defobject{currentmarker}{\pgfqpoint{0.000000in}{-0.020833in}}{\pgfqpoint{0.000000in}{0.000000in}}{%
\pgfpathmoveto{\pgfqpoint{0.000000in}{0.000000in}}%
\pgfpathlineto{\pgfqpoint{0.000000in}{-0.020833in}}%
\pgfusepath{stroke,fill}%
}%
\begin{pgfscope}%
\pgfsys@transformshift{2.167493in}{2.414488in}%
\pgfsys@useobject{currentmarker}{}%
\end{pgfscope}%
\end{pgfscope}%
\begin{pgfscope}%
\pgfsetbuttcap%
\pgfsetroundjoin%
\definecolor{currentfill}{rgb}{0.000000,0.000000,0.000000}%
\pgfsetfillcolor{currentfill}%
\pgfsetlinewidth{0.401500pt}%
\definecolor{currentstroke}{rgb}{0.000000,0.000000,0.000000}%
\pgfsetstrokecolor{currentstroke}%
\pgfsetdash{}{0pt}%
\pgfsys@defobject{currentmarker}{\pgfqpoint{0.000000in}{0.000000in}}{\pgfqpoint{0.000000in}{0.020833in}}{%
\pgfpathmoveto{\pgfqpoint{0.000000in}{0.000000in}}%
\pgfpathlineto{\pgfqpoint{0.000000in}{0.020833in}}%
\pgfusepath{stroke,fill}%
}%
\begin{pgfscope}%
\pgfsys@transformshift{2.504714in}{0.420000in}%
\pgfsys@useobject{currentmarker}{}%
\end{pgfscope}%
\end{pgfscope}%
\begin{pgfscope}%
\pgfsetbuttcap%
\pgfsetroundjoin%
\definecolor{currentfill}{rgb}{0.000000,0.000000,0.000000}%
\pgfsetfillcolor{currentfill}%
\pgfsetlinewidth{0.401500pt}%
\definecolor{currentstroke}{rgb}{0.000000,0.000000,0.000000}%
\pgfsetstrokecolor{currentstroke}%
\pgfsetdash{}{0pt}%
\pgfsys@defobject{currentmarker}{\pgfqpoint{0.000000in}{-0.020833in}}{\pgfqpoint{0.000000in}{0.000000in}}{%
\pgfpathmoveto{\pgfqpoint{0.000000in}{0.000000in}}%
\pgfpathlineto{\pgfqpoint{0.000000in}{-0.020833in}}%
\pgfusepath{stroke,fill}%
}%
\begin{pgfscope}%
\pgfsys@transformshift{2.504714in}{2.414488in}%
\pgfsys@useobject{currentmarker}{}%
\end{pgfscope}%
\end{pgfscope}%
\begin{pgfscope}%
\pgfsetbuttcap%
\pgfsetroundjoin%
\definecolor{currentfill}{rgb}{0.000000,0.000000,0.000000}%
\pgfsetfillcolor{currentfill}%
\pgfsetlinewidth{0.401500pt}%
\definecolor{currentstroke}{rgb}{0.000000,0.000000,0.000000}%
\pgfsetstrokecolor{currentstroke}%
\pgfsetdash{}{0pt}%
\pgfsys@defobject{currentmarker}{\pgfqpoint{0.000000in}{0.000000in}}{\pgfqpoint{0.000000in}{0.020833in}}{%
\pgfpathmoveto{\pgfqpoint{0.000000in}{0.000000in}}%
\pgfpathlineto{\pgfqpoint{0.000000in}{0.020833in}}%
\pgfusepath{stroke,fill}%
}%
\begin{pgfscope}%
\pgfsys@transformshift{2.673324in}{0.420000in}%
\pgfsys@useobject{currentmarker}{}%
\end{pgfscope}%
\end{pgfscope}%
\begin{pgfscope}%
\pgfsetbuttcap%
\pgfsetroundjoin%
\definecolor{currentfill}{rgb}{0.000000,0.000000,0.000000}%
\pgfsetfillcolor{currentfill}%
\pgfsetlinewidth{0.401500pt}%
\definecolor{currentstroke}{rgb}{0.000000,0.000000,0.000000}%
\pgfsetstrokecolor{currentstroke}%
\pgfsetdash{}{0pt}%
\pgfsys@defobject{currentmarker}{\pgfqpoint{0.000000in}{-0.020833in}}{\pgfqpoint{0.000000in}{0.000000in}}{%
\pgfpathmoveto{\pgfqpoint{0.000000in}{0.000000in}}%
\pgfpathlineto{\pgfqpoint{0.000000in}{-0.020833in}}%
\pgfusepath{stroke,fill}%
}%
\begin{pgfscope}%
\pgfsys@transformshift{2.673324in}{2.414488in}%
\pgfsys@useobject{currentmarker}{}%
\end{pgfscope}%
\end{pgfscope}%
\begin{pgfscope}%
\pgfsetbuttcap%
\pgfsetroundjoin%
\definecolor{currentfill}{rgb}{0.000000,0.000000,0.000000}%
\pgfsetfillcolor{currentfill}%
\pgfsetlinewidth{0.401500pt}%
\definecolor{currentstroke}{rgb}{0.000000,0.000000,0.000000}%
\pgfsetstrokecolor{currentstroke}%
\pgfsetdash{}{0pt}%
\pgfsys@defobject{currentmarker}{\pgfqpoint{0.000000in}{0.000000in}}{\pgfqpoint{0.000000in}{0.020833in}}{%
\pgfpathmoveto{\pgfqpoint{0.000000in}{0.000000in}}%
\pgfpathlineto{\pgfqpoint{0.000000in}{0.020833in}}%
\pgfusepath{stroke,fill}%
}%
\begin{pgfscope}%
\pgfsys@transformshift{2.841934in}{0.420000in}%
\pgfsys@useobject{currentmarker}{}%
\end{pgfscope}%
\end{pgfscope}%
\begin{pgfscope}%
\pgfsetbuttcap%
\pgfsetroundjoin%
\definecolor{currentfill}{rgb}{0.000000,0.000000,0.000000}%
\pgfsetfillcolor{currentfill}%
\pgfsetlinewidth{0.401500pt}%
\definecolor{currentstroke}{rgb}{0.000000,0.000000,0.000000}%
\pgfsetstrokecolor{currentstroke}%
\pgfsetdash{}{0pt}%
\pgfsys@defobject{currentmarker}{\pgfqpoint{0.000000in}{-0.020833in}}{\pgfqpoint{0.000000in}{0.000000in}}{%
\pgfpathmoveto{\pgfqpoint{0.000000in}{0.000000in}}%
\pgfpathlineto{\pgfqpoint{0.000000in}{-0.020833in}}%
\pgfusepath{stroke,fill}%
}%
\begin{pgfscope}%
\pgfsys@transformshift{2.841934in}{2.414488in}%
\pgfsys@useobject{currentmarker}{}%
\end{pgfscope}%
\end{pgfscope}%
\begin{pgfscope}%
\pgfsetbuttcap%
\pgfsetroundjoin%
\definecolor{currentfill}{rgb}{0.000000,0.000000,0.000000}%
\pgfsetfillcolor{currentfill}%
\pgfsetlinewidth{0.401500pt}%
\definecolor{currentstroke}{rgb}{0.000000,0.000000,0.000000}%
\pgfsetstrokecolor{currentstroke}%
\pgfsetdash{}{0pt}%
\pgfsys@defobject{currentmarker}{\pgfqpoint{0.000000in}{0.000000in}}{\pgfqpoint{0.000000in}{0.020833in}}{%
\pgfpathmoveto{\pgfqpoint{0.000000in}{0.000000in}}%
\pgfpathlineto{\pgfqpoint{0.000000in}{0.020833in}}%
\pgfusepath{stroke,fill}%
}%
\begin{pgfscope}%
\pgfsys@transformshift{3.010545in}{0.420000in}%
\pgfsys@useobject{currentmarker}{}%
\end{pgfscope}%
\end{pgfscope}%
\begin{pgfscope}%
\pgfsetbuttcap%
\pgfsetroundjoin%
\definecolor{currentfill}{rgb}{0.000000,0.000000,0.000000}%
\pgfsetfillcolor{currentfill}%
\pgfsetlinewidth{0.401500pt}%
\definecolor{currentstroke}{rgb}{0.000000,0.000000,0.000000}%
\pgfsetstrokecolor{currentstroke}%
\pgfsetdash{}{0pt}%
\pgfsys@defobject{currentmarker}{\pgfqpoint{0.000000in}{-0.020833in}}{\pgfqpoint{0.000000in}{0.000000in}}{%
\pgfpathmoveto{\pgfqpoint{0.000000in}{0.000000in}}%
\pgfpathlineto{\pgfqpoint{0.000000in}{-0.020833in}}%
\pgfusepath{stroke,fill}%
}%
\begin{pgfscope}%
\pgfsys@transformshift{3.010545in}{2.414488in}%
\pgfsys@useobject{currentmarker}{}%
\end{pgfscope}%
\end{pgfscope}%
\begin{pgfscope}%
\definecolor{textcolor}{rgb}{0.000000,0.000000,0.000000}%
\pgfsetstrokecolor{textcolor}%
\pgfsetfillcolor{textcolor}%
\pgftext[x=1.844055in,y=0.208426in,,top]{\color{textcolor}{\rmfamily\fontsize{12.000000}{14.400000}\selectfont\catcode`\^=\active\def^{\ifmmode\sp\else\^{}\fi}\catcode`\%=\active\def%{\%}$s/\delta_{0}$}}%
\end{pgfscope}%
\begin{pgfscope}%
\pgfsetbuttcap%
\pgfsetroundjoin%
\definecolor{currentfill}{rgb}{0.000000,0.000000,0.000000}%
\pgfsetfillcolor{currentfill}%
\pgfsetlinewidth{0.501875pt}%
\definecolor{currentstroke}{rgb}{0.000000,0.000000,0.000000}%
\pgfsetstrokecolor{currentstroke}%
\pgfsetdash{}{0pt}%
\pgfsys@defobject{currentmarker}{\pgfqpoint{0.000000in}{0.000000in}}{\pgfqpoint{0.041667in}{0.000000in}}{%
\pgfpathmoveto{\pgfqpoint{0.000000in}{0.000000in}}%
\pgfpathlineto{\pgfqpoint{0.041667in}{0.000000in}}%
\pgfusepath{stroke,fill}%
}%
\begin{pgfscope}%
\pgfsys@transformshift{0.650000in}{0.430501in}%
\pgfsys@useobject{currentmarker}{}%
\end{pgfscope}%
\end{pgfscope}%
\begin{pgfscope}%
\pgfsetbuttcap%
\pgfsetroundjoin%
\definecolor{currentfill}{rgb}{0.000000,0.000000,0.000000}%
\pgfsetfillcolor{currentfill}%
\pgfsetlinewidth{0.501875pt}%
\definecolor{currentstroke}{rgb}{0.000000,0.000000,0.000000}%
\pgfsetstrokecolor{currentstroke}%
\pgfsetdash{}{0pt}%
\pgfsys@defobject{currentmarker}{\pgfqpoint{-0.041667in}{0.000000in}}{\pgfqpoint{-0.000000in}{0.000000in}}{%
\pgfpathmoveto{\pgfqpoint{-0.000000in}{0.000000in}}%
\pgfpathlineto{\pgfqpoint{-0.041667in}{0.000000in}}%
\pgfusepath{stroke,fill}%
}%
\begin{pgfscope}%
\pgfsys@transformshift{3.038110in}{0.430501in}%
\pgfsys@useobject{currentmarker}{}%
\end{pgfscope}%
\end{pgfscope}%
\begin{pgfscope}%
\definecolor{textcolor}{rgb}{0.000000,0.000000,0.000000}%
\pgfsetstrokecolor{textcolor}%
\pgfsetfillcolor{textcolor}%
\pgftext[x=0.525347in, y=0.377694in, left, base]{\color{textcolor}{\rmfamily\fontsize{11.000000}{13.200000}\selectfont\catcode`\^=\active\def^{\ifmmode\sp\else\^{}\fi}\catcode`\%=\active\def%{\%}0}}%
\end{pgfscope}%
\begin{pgfscope}%
\pgfsetbuttcap%
\pgfsetroundjoin%
\definecolor{currentfill}{rgb}{0.000000,0.000000,0.000000}%
\pgfsetfillcolor{currentfill}%
\pgfsetlinewidth{0.501875pt}%
\definecolor{currentstroke}{rgb}{0.000000,0.000000,0.000000}%
\pgfsetstrokecolor{currentstroke}%
\pgfsetdash{}{0pt}%
\pgfsys@defobject{currentmarker}{\pgfqpoint{0.000000in}{0.000000in}}{\pgfqpoint{0.041667in}{0.000000in}}{%
\pgfpathmoveto{\pgfqpoint{0.000000in}{0.000000in}}%
\pgfpathlineto{\pgfqpoint{0.041667in}{0.000000in}}%
\pgfusepath{stroke,fill}%
}%
\begin{pgfscope}%
\pgfsys@transformshift{0.650000in}{0.923817in}%
\pgfsys@useobject{currentmarker}{}%
\end{pgfscope}%
\end{pgfscope}%
\begin{pgfscope}%
\pgfsetbuttcap%
\pgfsetroundjoin%
\definecolor{currentfill}{rgb}{0.000000,0.000000,0.000000}%
\pgfsetfillcolor{currentfill}%
\pgfsetlinewidth{0.501875pt}%
\definecolor{currentstroke}{rgb}{0.000000,0.000000,0.000000}%
\pgfsetstrokecolor{currentstroke}%
\pgfsetdash{}{0pt}%
\pgfsys@defobject{currentmarker}{\pgfqpoint{-0.041667in}{0.000000in}}{\pgfqpoint{-0.000000in}{0.000000in}}{%
\pgfpathmoveto{\pgfqpoint{-0.000000in}{0.000000in}}%
\pgfpathlineto{\pgfqpoint{-0.041667in}{0.000000in}}%
\pgfusepath{stroke,fill}%
}%
\begin{pgfscope}%
\pgfsys@transformshift{3.038110in}{0.923817in}%
\pgfsys@useobject{currentmarker}{}%
\end{pgfscope}%
\end{pgfscope}%
\begin{pgfscope}%
\definecolor{textcolor}{rgb}{0.000000,0.000000,0.000000}%
\pgfsetstrokecolor{textcolor}%
\pgfsetfillcolor{textcolor}%
\pgftext[x=0.525347in, y=0.871010in, left, base]{\color{textcolor}{\rmfamily\fontsize{11.000000}{13.200000}\selectfont\catcode`\^=\active\def^{\ifmmode\sp\else\^{}\fi}\catcode`\%=\active\def%{\%}1}}%
\end{pgfscope}%
\begin{pgfscope}%
\pgfsetbuttcap%
\pgfsetroundjoin%
\definecolor{currentfill}{rgb}{0.000000,0.000000,0.000000}%
\pgfsetfillcolor{currentfill}%
\pgfsetlinewidth{0.501875pt}%
\definecolor{currentstroke}{rgb}{0.000000,0.000000,0.000000}%
\pgfsetstrokecolor{currentstroke}%
\pgfsetdash{}{0pt}%
\pgfsys@defobject{currentmarker}{\pgfqpoint{0.000000in}{0.000000in}}{\pgfqpoint{0.041667in}{0.000000in}}{%
\pgfpathmoveto{\pgfqpoint{0.000000in}{0.000000in}}%
\pgfpathlineto{\pgfqpoint{0.041667in}{0.000000in}}%
\pgfusepath{stroke,fill}%
}%
\begin{pgfscope}%
\pgfsys@transformshift{0.650000in}{1.417133in}%
\pgfsys@useobject{currentmarker}{}%
\end{pgfscope}%
\end{pgfscope}%
\begin{pgfscope}%
\pgfsetbuttcap%
\pgfsetroundjoin%
\definecolor{currentfill}{rgb}{0.000000,0.000000,0.000000}%
\pgfsetfillcolor{currentfill}%
\pgfsetlinewidth{0.501875pt}%
\definecolor{currentstroke}{rgb}{0.000000,0.000000,0.000000}%
\pgfsetstrokecolor{currentstroke}%
\pgfsetdash{}{0pt}%
\pgfsys@defobject{currentmarker}{\pgfqpoint{-0.041667in}{0.000000in}}{\pgfqpoint{-0.000000in}{0.000000in}}{%
\pgfpathmoveto{\pgfqpoint{-0.000000in}{0.000000in}}%
\pgfpathlineto{\pgfqpoint{-0.041667in}{0.000000in}}%
\pgfusepath{stroke,fill}%
}%
\begin{pgfscope}%
\pgfsys@transformshift{3.038110in}{1.417133in}%
\pgfsys@useobject{currentmarker}{}%
\end{pgfscope}%
\end{pgfscope}%
\begin{pgfscope}%
\definecolor{textcolor}{rgb}{0.000000,0.000000,0.000000}%
\pgfsetstrokecolor{textcolor}%
\pgfsetfillcolor{textcolor}%
\pgftext[x=0.525347in, y=1.364326in, left, base]{\color{textcolor}{\rmfamily\fontsize{11.000000}{13.200000}\selectfont\catcode`\^=\active\def^{\ifmmode\sp\else\^{}\fi}\catcode`\%=\active\def%{\%}2}}%
\end{pgfscope}%
\begin{pgfscope}%
\pgfsetbuttcap%
\pgfsetroundjoin%
\definecolor{currentfill}{rgb}{0.000000,0.000000,0.000000}%
\pgfsetfillcolor{currentfill}%
\pgfsetlinewidth{0.501875pt}%
\definecolor{currentstroke}{rgb}{0.000000,0.000000,0.000000}%
\pgfsetstrokecolor{currentstroke}%
\pgfsetdash{}{0pt}%
\pgfsys@defobject{currentmarker}{\pgfqpoint{0.000000in}{0.000000in}}{\pgfqpoint{0.041667in}{0.000000in}}{%
\pgfpathmoveto{\pgfqpoint{0.000000in}{0.000000in}}%
\pgfpathlineto{\pgfqpoint{0.041667in}{0.000000in}}%
\pgfusepath{stroke,fill}%
}%
\begin{pgfscope}%
\pgfsys@transformshift{0.650000in}{1.910449in}%
\pgfsys@useobject{currentmarker}{}%
\end{pgfscope}%
\end{pgfscope}%
\begin{pgfscope}%
\pgfsetbuttcap%
\pgfsetroundjoin%
\definecolor{currentfill}{rgb}{0.000000,0.000000,0.000000}%
\pgfsetfillcolor{currentfill}%
\pgfsetlinewidth{0.501875pt}%
\definecolor{currentstroke}{rgb}{0.000000,0.000000,0.000000}%
\pgfsetstrokecolor{currentstroke}%
\pgfsetdash{}{0pt}%
\pgfsys@defobject{currentmarker}{\pgfqpoint{-0.041667in}{0.000000in}}{\pgfqpoint{-0.000000in}{0.000000in}}{%
\pgfpathmoveto{\pgfqpoint{-0.000000in}{0.000000in}}%
\pgfpathlineto{\pgfqpoint{-0.041667in}{0.000000in}}%
\pgfusepath{stroke,fill}%
}%
\begin{pgfscope}%
\pgfsys@transformshift{3.038110in}{1.910449in}%
\pgfsys@useobject{currentmarker}{}%
\end{pgfscope}%
\end{pgfscope}%
\begin{pgfscope}%
\definecolor{textcolor}{rgb}{0.000000,0.000000,0.000000}%
\pgfsetstrokecolor{textcolor}%
\pgfsetfillcolor{textcolor}%
\pgftext[x=0.525347in, y=1.857642in, left, base]{\color{textcolor}{\rmfamily\fontsize{11.000000}{13.200000}\selectfont\catcode`\^=\active\def^{\ifmmode\sp\else\^{}\fi}\catcode`\%=\active\def%{\%}3}}%
\end{pgfscope}%
\begin{pgfscope}%
\pgfsetbuttcap%
\pgfsetroundjoin%
\definecolor{currentfill}{rgb}{0.000000,0.000000,0.000000}%
\pgfsetfillcolor{currentfill}%
\pgfsetlinewidth{0.501875pt}%
\definecolor{currentstroke}{rgb}{0.000000,0.000000,0.000000}%
\pgfsetstrokecolor{currentstroke}%
\pgfsetdash{}{0pt}%
\pgfsys@defobject{currentmarker}{\pgfqpoint{0.000000in}{0.000000in}}{\pgfqpoint{0.041667in}{0.000000in}}{%
\pgfpathmoveto{\pgfqpoint{0.000000in}{0.000000in}}%
\pgfpathlineto{\pgfqpoint{0.041667in}{0.000000in}}%
\pgfusepath{stroke,fill}%
}%
\begin{pgfscope}%
\pgfsys@transformshift{0.650000in}{2.403765in}%
\pgfsys@useobject{currentmarker}{}%
\end{pgfscope}%
\end{pgfscope}%
\begin{pgfscope}%
\pgfsetbuttcap%
\pgfsetroundjoin%
\definecolor{currentfill}{rgb}{0.000000,0.000000,0.000000}%
\pgfsetfillcolor{currentfill}%
\pgfsetlinewidth{0.501875pt}%
\definecolor{currentstroke}{rgb}{0.000000,0.000000,0.000000}%
\pgfsetstrokecolor{currentstroke}%
\pgfsetdash{}{0pt}%
\pgfsys@defobject{currentmarker}{\pgfqpoint{-0.041667in}{0.000000in}}{\pgfqpoint{-0.000000in}{0.000000in}}{%
\pgfpathmoveto{\pgfqpoint{-0.000000in}{0.000000in}}%
\pgfpathlineto{\pgfqpoint{-0.041667in}{0.000000in}}%
\pgfusepath{stroke,fill}%
}%
\begin{pgfscope}%
\pgfsys@transformshift{3.038110in}{2.403765in}%
\pgfsys@useobject{currentmarker}{}%
\end{pgfscope}%
\end{pgfscope}%
\begin{pgfscope}%
\definecolor{textcolor}{rgb}{0.000000,0.000000,0.000000}%
\pgfsetstrokecolor{textcolor}%
\pgfsetfillcolor{textcolor}%
\pgftext[x=0.525347in, y=2.350958in, left, base]{\color{textcolor}{\rmfamily\fontsize{11.000000}{13.200000}\selectfont\catcode`\^=\active\def^{\ifmmode\sp\else\^{}\fi}\catcode`\%=\active\def%{\%}4}}%
\end{pgfscope}%
\begin{pgfscope}%
\pgfsetbuttcap%
\pgfsetroundjoin%
\definecolor{currentfill}{rgb}{0.000000,0.000000,0.000000}%
\pgfsetfillcolor{currentfill}%
\pgfsetlinewidth{0.401500pt}%
\definecolor{currentstroke}{rgb}{0.000000,0.000000,0.000000}%
\pgfsetstrokecolor{currentstroke}%
\pgfsetdash{}{0pt}%
\pgfsys@defobject{currentmarker}{\pgfqpoint{0.000000in}{0.000000in}}{\pgfqpoint{0.020833in}{0.000000in}}{%
\pgfpathmoveto{\pgfqpoint{0.000000in}{0.000000in}}%
\pgfpathlineto{\pgfqpoint{0.020833in}{0.000000in}}%
\pgfusepath{stroke,fill}%
}%
\begin{pgfscope}%
\pgfsys@transformshift{0.650000in}{0.529164in}%
\pgfsys@useobject{currentmarker}{}%
\end{pgfscope}%
\end{pgfscope}%
\begin{pgfscope}%
\pgfsetbuttcap%
\pgfsetroundjoin%
\definecolor{currentfill}{rgb}{0.000000,0.000000,0.000000}%
\pgfsetfillcolor{currentfill}%
\pgfsetlinewidth{0.401500pt}%
\definecolor{currentstroke}{rgb}{0.000000,0.000000,0.000000}%
\pgfsetstrokecolor{currentstroke}%
\pgfsetdash{}{0pt}%
\pgfsys@defobject{currentmarker}{\pgfqpoint{-0.020833in}{0.000000in}}{\pgfqpoint{-0.000000in}{0.000000in}}{%
\pgfpathmoveto{\pgfqpoint{-0.000000in}{0.000000in}}%
\pgfpathlineto{\pgfqpoint{-0.020833in}{0.000000in}}%
\pgfusepath{stroke,fill}%
}%
\begin{pgfscope}%
\pgfsys@transformshift{3.038110in}{0.529164in}%
\pgfsys@useobject{currentmarker}{}%
\end{pgfscope}%
\end{pgfscope}%
\begin{pgfscope}%
\pgfsetbuttcap%
\pgfsetroundjoin%
\definecolor{currentfill}{rgb}{0.000000,0.000000,0.000000}%
\pgfsetfillcolor{currentfill}%
\pgfsetlinewidth{0.401500pt}%
\definecolor{currentstroke}{rgb}{0.000000,0.000000,0.000000}%
\pgfsetstrokecolor{currentstroke}%
\pgfsetdash{}{0pt}%
\pgfsys@defobject{currentmarker}{\pgfqpoint{0.000000in}{0.000000in}}{\pgfqpoint{0.020833in}{0.000000in}}{%
\pgfpathmoveto{\pgfqpoint{0.000000in}{0.000000in}}%
\pgfpathlineto{\pgfqpoint{0.020833in}{0.000000in}}%
\pgfusepath{stroke,fill}%
}%
\begin{pgfscope}%
\pgfsys@transformshift{0.650000in}{0.627827in}%
\pgfsys@useobject{currentmarker}{}%
\end{pgfscope}%
\end{pgfscope}%
\begin{pgfscope}%
\pgfsetbuttcap%
\pgfsetroundjoin%
\definecolor{currentfill}{rgb}{0.000000,0.000000,0.000000}%
\pgfsetfillcolor{currentfill}%
\pgfsetlinewidth{0.401500pt}%
\definecolor{currentstroke}{rgb}{0.000000,0.000000,0.000000}%
\pgfsetstrokecolor{currentstroke}%
\pgfsetdash{}{0pt}%
\pgfsys@defobject{currentmarker}{\pgfqpoint{-0.020833in}{0.000000in}}{\pgfqpoint{-0.000000in}{0.000000in}}{%
\pgfpathmoveto{\pgfqpoint{-0.000000in}{0.000000in}}%
\pgfpathlineto{\pgfqpoint{-0.020833in}{0.000000in}}%
\pgfusepath{stroke,fill}%
}%
\begin{pgfscope}%
\pgfsys@transformshift{3.038110in}{0.627827in}%
\pgfsys@useobject{currentmarker}{}%
\end{pgfscope}%
\end{pgfscope}%
\begin{pgfscope}%
\pgfsetbuttcap%
\pgfsetroundjoin%
\definecolor{currentfill}{rgb}{0.000000,0.000000,0.000000}%
\pgfsetfillcolor{currentfill}%
\pgfsetlinewidth{0.401500pt}%
\definecolor{currentstroke}{rgb}{0.000000,0.000000,0.000000}%
\pgfsetstrokecolor{currentstroke}%
\pgfsetdash{}{0pt}%
\pgfsys@defobject{currentmarker}{\pgfqpoint{0.000000in}{0.000000in}}{\pgfqpoint{0.020833in}{0.000000in}}{%
\pgfpathmoveto{\pgfqpoint{0.000000in}{0.000000in}}%
\pgfpathlineto{\pgfqpoint{0.020833in}{0.000000in}}%
\pgfusepath{stroke,fill}%
}%
\begin{pgfscope}%
\pgfsys@transformshift{0.650000in}{0.726491in}%
\pgfsys@useobject{currentmarker}{}%
\end{pgfscope}%
\end{pgfscope}%
\begin{pgfscope}%
\pgfsetbuttcap%
\pgfsetroundjoin%
\definecolor{currentfill}{rgb}{0.000000,0.000000,0.000000}%
\pgfsetfillcolor{currentfill}%
\pgfsetlinewidth{0.401500pt}%
\definecolor{currentstroke}{rgb}{0.000000,0.000000,0.000000}%
\pgfsetstrokecolor{currentstroke}%
\pgfsetdash{}{0pt}%
\pgfsys@defobject{currentmarker}{\pgfqpoint{-0.020833in}{0.000000in}}{\pgfqpoint{-0.000000in}{0.000000in}}{%
\pgfpathmoveto{\pgfqpoint{-0.000000in}{0.000000in}}%
\pgfpathlineto{\pgfqpoint{-0.020833in}{0.000000in}}%
\pgfusepath{stroke,fill}%
}%
\begin{pgfscope}%
\pgfsys@transformshift{3.038110in}{0.726491in}%
\pgfsys@useobject{currentmarker}{}%
\end{pgfscope}%
\end{pgfscope}%
\begin{pgfscope}%
\pgfsetbuttcap%
\pgfsetroundjoin%
\definecolor{currentfill}{rgb}{0.000000,0.000000,0.000000}%
\pgfsetfillcolor{currentfill}%
\pgfsetlinewidth{0.401500pt}%
\definecolor{currentstroke}{rgb}{0.000000,0.000000,0.000000}%
\pgfsetstrokecolor{currentstroke}%
\pgfsetdash{}{0pt}%
\pgfsys@defobject{currentmarker}{\pgfqpoint{0.000000in}{0.000000in}}{\pgfqpoint{0.020833in}{0.000000in}}{%
\pgfpathmoveto{\pgfqpoint{0.000000in}{0.000000in}}%
\pgfpathlineto{\pgfqpoint{0.020833in}{0.000000in}}%
\pgfusepath{stroke,fill}%
}%
\begin{pgfscope}%
\pgfsys@transformshift{0.650000in}{0.825154in}%
\pgfsys@useobject{currentmarker}{}%
\end{pgfscope}%
\end{pgfscope}%
\begin{pgfscope}%
\pgfsetbuttcap%
\pgfsetroundjoin%
\definecolor{currentfill}{rgb}{0.000000,0.000000,0.000000}%
\pgfsetfillcolor{currentfill}%
\pgfsetlinewidth{0.401500pt}%
\definecolor{currentstroke}{rgb}{0.000000,0.000000,0.000000}%
\pgfsetstrokecolor{currentstroke}%
\pgfsetdash{}{0pt}%
\pgfsys@defobject{currentmarker}{\pgfqpoint{-0.020833in}{0.000000in}}{\pgfqpoint{-0.000000in}{0.000000in}}{%
\pgfpathmoveto{\pgfqpoint{-0.000000in}{0.000000in}}%
\pgfpathlineto{\pgfqpoint{-0.020833in}{0.000000in}}%
\pgfusepath{stroke,fill}%
}%
\begin{pgfscope}%
\pgfsys@transformshift{3.038110in}{0.825154in}%
\pgfsys@useobject{currentmarker}{}%
\end{pgfscope}%
\end{pgfscope}%
\begin{pgfscope}%
\pgfsetbuttcap%
\pgfsetroundjoin%
\definecolor{currentfill}{rgb}{0.000000,0.000000,0.000000}%
\pgfsetfillcolor{currentfill}%
\pgfsetlinewidth{0.401500pt}%
\definecolor{currentstroke}{rgb}{0.000000,0.000000,0.000000}%
\pgfsetstrokecolor{currentstroke}%
\pgfsetdash{}{0pt}%
\pgfsys@defobject{currentmarker}{\pgfqpoint{0.000000in}{0.000000in}}{\pgfqpoint{0.020833in}{0.000000in}}{%
\pgfpathmoveto{\pgfqpoint{0.000000in}{0.000000in}}%
\pgfpathlineto{\pgfqpoint{0.020833in}{0.000000in}}%
\pgfusepath{stroke,fill}%
}%
\begin{pgfscope}%
\pgfsys@transformshift{0.650000in}{1.022480in}%
\pgfsys@useobject{currentmarker}{}%
\end{pgfscope}%
\end{pgfscope}%
\begin{pgfscope}%
\pgfsetbuttcap%
\pgfsetroundjoin%
\definecolor{currentfill}{rgb}{0.000000,0.000000,0.000000}%
\pgfsetfillcolor{currentfill}%
\pgfsetlinewidth{0.401500pt}%
\definecolor{currentstroke}{rgb}{0.000000,0.000000,0.000000}%
\pgfsetstrokecolor{currentstroke}%
\pgfsetdash{}{0pt}%
\pgfsys@defobject{currentmarker}{\pgfqpoint{-0.020833in}{0.000000in}}{\pgfqpoint{-0.000000in}{0.000000in}}{%
\pgfpathmoveto{\pgfqpoint{-0.000000in}{0.000000in}}%
\pgfpathlineto{\pgfqpoint{-0.020833in}{0.000000in}}%
\pgfusepath{stroke,fill}%
}%
\begin{pgfscope}%
\pgfsys@transformshift{3.038110in}{1.022480in}%
\pgfsys@useobject{currentmarker}{}%
\end{pgfscope}%
\end{pgfscope}%
\begin{pgfscope}%
\pgfsetbuttcap%
\pgfsetroundjoin%
\definecolor{currentfill}{rgb}{0.000000,0.000000,0.000000}%
\pgfsetfillcolor{currentfill}%
\pgfsetlinewidth{0.401500pt}%
\definecolor{currentstroke}{rgb}{0.000000,0.000000,0.000000}%
\pgfsetstrokecolor{currentstroke}%
\pgfsetdash{}{0pt}%
\pgfsys@defobject{currentmarker}{\pgfqpoint{0.000000in}{0.000000in}}{\pgfqpoint{0.020833in}{0.000000in}}{%
\pgfpathmoveto{\pgfqpoint{0.000000in}{0.000000in}}%
\pgfpathlineto{\pgfqpoint{0.020833in}{0.000000in}}%
\pgfusepath{stroke,fill}%
}%
\begin{pgfscope}%
\pgfsys@transformshift{0.650000in}{1.121143in}%
\pgfsys@useobject{currentmarker}{}%
\end{pgfscope}%
\end{pgfscope}%
\begin{pgfscope}%
\pgfsetbuttcap%
\pgfsetroundjoin%
\definecolor{currentfill}{rgb}{0.000000,0.000000,0.000000}%
\pgfsetfillcolor{currentfill}%
\pgfsetlinewidth{0.401500pt}%
\definecolor{currentstroke}{rgb}{0.000000,0.000000,0.000000}%
\pgfsetstrokecolor{currentstroke}%
\pgfsetdash{}{0pt}%
\pgfsys@defobject{currentmarker}{\pgfqpoint{-0.020833in}{0.000000in}}{\pgfqpoint{-0.000000in}{0.000000in}}{%
\pgfpathmoveto{\pgfqpoint{-0.000000in}{0.000000in}}%
\pgfpathlineto{\pgfqpoint{-0.020833in}{0.000000in}}%
\pgfusepath{stroke,fill}%
}%
\begin{pgfscope}%
\pgfsys@transformshift{3.038110in}{1.121143in}%
\pgfsys@useobject{currentmarker}{}%
\end{pgfscope}%
\end{pgfscope}%
\begin{pgfscope}%
\pgfsetbuttcap%
\pgfsetroundjoin%
\definecolor{currentfill}{rgb}{0.000000,0.000000,0.000000}%
\pgfsetfillcolor{currentfill}%
\pgfsetlinewidth{0.401500pt}%
\definecolor{currentstroke}{rgb}{0.000000,0.000000,0.000000}%
\pgfsetstrokecolor{currentstroke}%
\pgfsetdash{}{0pt}%
\pgfsys@defobject{currentmarker}{\pgfqpoint{0.000000in}{0.000000in}}{\pgfqpoint{0.020833in}{0.000000in}}{%
\pgfpathmoveto{\pgfqpoint{0.000000in}{0.000000in}}%
\pgfpathlineto{\pgfqpoint{0.020833in}{0.000000in}}%
\pgfusepath{stroke,fill}%
}%
\begin{pgfscope}%
\pgfsys@transformshift{0.650000in}{1.219807in}%
\pgfsys@useobject{currentmarker}{}%
\end{pgfscope}%
\end{pgfscope}%
\begin{pgfscope}%
\pgfsetbuttcap%
\pgfsetroundjoin%
\definecolor{currentfill}{rgb}{0.000000,0.000000,0.000000}%
\pgfsetfillcolor{currentfill}%
\pgfsetlinewidth{0.401500pt}%
\definecolor{currentstroke}{rgb}{0.000000,0.000000,0.000000}%
\pgfsetstrokecolor{currentstroke}%
\pgfsetdash{}{0pt}%
\pgfsys@defobject{currentmarker}{\pgfqpoint{-0.020833in}{0.000000in}}{\pgfqpoint{-0.000000in}{0.000000in}}{%
\pgfpathmoveto{\pgfqpoint{-0.000000in}{0.000000in}}%
\pgfpathlineto{\pgfqpoint{-0.020833in}{0.000000in}}%
\pgfusepath{stroke,fill}%
}%
\begin{pgfscope}%
\pgfsys@transformshift{3.038110in}{1.219807in}%
\pgfsys@useobject{currentmarker}{}%
\end{pgfscope}%
\end{pgfscope}%
\begin{pgfscope}%
\pgfsetbuttcap%
\pgfsetroundjoin%
\definecolor{currentfill}{rgb}{0.000000,0.000000,0.000000}%
\pgfsetfillcolor{currentfill}%
\pgfsetlinewidth{0.401500pt}%
\definecolor{currentstroke}{rgb}{0.000000,0.000000,0.000000}%
\pgfsetstrokecolor{currentstroke}%
\pgfsetdash{}{0pt}%
\pgfsys@defobject{currentmarker}{\pgfqpoint{0.000000in}{0.000000in}}{\pgfqpoint{0.020833in}{0.000000in}}{%
\pgfpathmoveto{\pgfqpoint{0.000000in}{0.000000in}}%
\pgfpathlineto{\pgfqpoint{0.020833in}{0.000000in}}%
\pgfusepath{stroke,fill}%
}%
\begin{pgfscope}%
\pgfsys@transformshift{0.650000in}{1.318470in}%
\pgfsys@useobject{currentmarker}{}%
\end{pgfscope}%
\end{pgfscope}%
\begin{pgfscope}%
\pgfsetbuttcap%
\pgfsetroundjoin%
\definecolor{currentfill}{rgb}{0.000000,0.000000,0.000000}%
\pgfsetfillcolor{currentfill}%
\pgfsetlinewidth{0.401500pt}%
\definecolor{currentstroke}{rgb}{0.000000,0.000000,0.000000}%
\pgfsetstrokecolor{currentstroke}%
\pgfsetdash{}{0pt}%
\pgfsys@defobject{currentmarker}{\pgfqpoint{-0.020833in}{0.000000in}}{\pgfqpoint{-0.000000in}{0.000000in}}{%
\pgfpathmoveto{\pgfqpoint{-0.000000in}{0.000000in}}%
\pgfpathlineto{\pgfqpoint{-0.020833in}{0.000000in}}%
\pgfusepath{stroke,fill}%
}%
\begin{pgfscope}%
\pgfsys@transformshift{3.038110in}{1.318470in}%
\pgfsys@useobject{currentmarker}{}%
\end{pgfscope}%
\end{pgfscope}%
\begin{pgfscope}%
\pgfsetbuttcap%
\pgfsetroundjoin%
\definecolor{currentfill}{rgb}{0.000000,0.000000,0.000000}%
\pgfsetfillcolor{currentfill}%
\pgfsetlinewidth{0.401500pt}%
\definecolor{currentstroke}{rgb}{0.000000,0.000000,0.000000}%
\pgfsetstrokecolor{currentstroke}%
\pgfsetdash{}{0pt}%
\pgfsys@defobject{currentmarker}{\pgfqpoint{0.000000in}{0.000000in}}{\pgfqpoint{0.020833in}{0.000000in}}{%
\pgfpathmoveto{\pgfqpoint{0.000000in}{0.000000in}}%
\pgfpathlineto{\pgfqpoint{0.020833in}{0.000000in}}%
\pgfusepath{stroke,fill}%
}%
\begin{pgfscope}%
\pgfsys@transformshift{0.650000in}{1.515796in}%
\pgfsys@useobject{currentmarker}{}%
\end{pgfscope}%
\end{pgfscope}%
\begin{pgfscope}%
\pgfsetbuttcap%
\pgfsetroundjoin%
\definecolor{currentfill}{rgb}{0.000000,0.000000,0.000000}%
\pgfsetfillcolor{currentfill}%
\pgfsetlinewidth{0.401500pt}%
\definecolor{currentstroke}{rgb}{0.000000,0.000000,0.000000}%
\pgfsetstrokecolor{currentstroke}%
\pgfsetdash{}{0pt}%
\pgfsys@defobject{currentmarker}{\pgfqpoint{-0.020833in}{0.000000in}}{\pgfqpoint{-0.000000in}{0.000000in}}{%
\pgfpathmoveto{\pgfqpoint{-0.000000in}{0.000000in}}%
\pgfpathlineto{\pgfqpoint{-0.020833in}{0.000000in}}%
\pgfusepath{stroke,fill}%
}%
\begin{pgfscope}%
\pgfsys@transformshift{3.038110in}{1.515796in}%
\pgfsys@useobject{currentmarker}{}%
\end{pgfscope}%
\end{pgfscope}%
\begin{pgfscope}%
\pgfsetbuttcap%
\pgfsetroundjoin%
\definecolor{currentfill}{rgb}{0.000000,0.000000,0.000000}%
\pgfsetfillcolor{currentfill}%
\pgfsetlinewidth{0.401500pt}%
\definecolor{currentstroke}{rgb}{0.000000,0.000000,0.000000}%
\pgfsetstrokecolor{currentstroke}%
\pgfsetdash{}{0pt}%
\pgfsys@defobject{currentmarker}{\pgfqpoint{0.000000in}{0.000000in}}{\pgfqpoint{0.020833in}{0.000000in}}{%
\pgfpathmoveto{\pgfqpoint{0.000000in}{0.000000in}}%
\pgfpathlineto{\pgfqpoint{0.020833in}{0.000000in}}%
\pgfusepath{stroke,fill}%
}%
\begin{pgfscope}%
\pgfsys@transformshift{0.650000in}{1.614459in}%
\pgfsys@useobject{currentmarker}{}%
\end{pgfscope}%
\end{pgfscope}%
\begin{pgfscope}%
\pgfsetbuttcap%
\pgfsetroundjoin%
\definecolor{currentfill}{rgb}{0.000000,0.000000,0.000000}%
\pgfsetfillcolor{currentfill}%
\pgfsetlinewidth{0.401500pt}%
\definecolor{currentstroke}{rgb}{0.000000,0.000000,0.000000}%
\pgfsetstrokecolor{currentstroke}%
\pgfsetdash{}{0pt}%
\pgfsys@defobject{currentmarker}{\pgfqpoint{-0.020833in}{0.000000in}}{\pgfqpoint{-0.000000in}{0.000000in}}{%
\pgfpathmoveto{\pgfqpoint{-0.000000in}{0.000000in}}%
\pgfpathlineto{\pgfqpoint{-0.020833in}{0.000000in}}%
\pgfusepath{stroke,fill}%
}%
\begin{pgfscope}%
\pgfsys@transformshift{3.038110in}{1.614459in}%
\pgfsys@useobject{currentmarker}{}%
\end{pgfscope}%
\end{pgfscope}%
\begin{pgfscope}%
\pgfsetbuttcap%
\pgfsetroundjoin%
\definecolor{currentfill}{rgb}{0.000000,0.000000,0.000000}%
\pgfsetfillcolor{currentfill}%
\pgfsetlinewidth{0.401500pt}%
\definecolor{currentstroke}{rgb}{0.000000,0.000000,0.000000}%
\pgfsetstrokecolor{currentstroke}%
\pgfsetdash{}{0pt}%
\pgfsys@defobject{currentmarker}{\pgfqpoint{0.000000in}{0.000000in}}{\pgfqpoint{0.020833in}{0.000000in}}{%
\pgfpathmoveto{\pgfqpoint{0.000000in}{0.000000in}}%
\pgfpathlineto{\pgfqpoint{0.020833in}{0.000000in}}%
\pgfusepath{stroke,fill}%
}%
\begin{pgfscope}%
\pgfsys@transformshift{0.650000in}{1.713122in}%
\pgfsys@useobject{currentmarker}{}%
\end{pgfscope}%
\end{pgfscope}%
\begin{pgfscope}%
\pgfsetbuttcap%
\pgfsetroundjoin%
\definecolor{currentfill}{rgb}{0.000000,0.000000,0.000000}%
\pgfsetfillcolor{currentfill}%
\pgfsetlinewidth{0.401500pt}%
\definecolor{currentstroke}{rgb}{0.000000,0.000000,0.000000}%
\pgfsetstrokecolor{currentstroke}%
\pgfsetdash{}{0pt}%
\pgfsys@defobject{currentmarker}{\pgfqpoint{-0.020833in}{0.000000in}}{\pgfqpoint{-0.000000in}{0.000000in}}{%
\pgfpathmoveto{\pgfqpoint{-0.000000in}{0.000000in}}%
\pgfpathlineto{\pgfqpoint{-0.020833in}{0.000000in}}%
\pgfusepath{stroke,fill}%
}%
\begin{pgfscope}%
\pgfsys@transformshift{3.038110in}{1.713122in}%
\pgfsys@useobject{currentmarker}{}%
\end{pgfscope}%
\end{pgfscope}%
\begin{pgfscope}%
\pgfsetbuttcap%
\pgfsetroundjoin%
\definecolor{currentfill}{rgb}{0.000000,0.000000,0.000000}%
\pgfsetfillcolor{currentfill}%
\pgfsetlinewidth{0.401500pt}%
\definecolor{currentstroke}{rgb}{0.000000,0.000000,0.000000}%
\pgfsetstrokecolor{currentstroke}%
\pgfsetdash{}{0pt}%
\pgfsys@defobject{currentmarker}{\pgfqpoint{0.000000in}{0.000000in}}{\pgfqpoint{0.020833in}{0.000000in}}{%
\pgfpathmoveto{\pgfqpoint{0.000000in}{0.000000in}}%
\pgfpathlineto{\pgfqpoint{0.020833in}{0.000000in}}%
\pgfusepath{stroke,fill}%
}%
\begin{pgfscope}%
\pgfsys@transformshift{0.650000in}{1.811786in}%
\pgfsys@useobject{currentmarker}{}%
\end{pgfscope}%
\end{pgfscope}%
\begin{pgfscope}%
\pgfsetbuttcap%
\pgfsetroundjoin%
\definecolor{currentfill}{rgb}{0.000000,0.000000,0.000000}%
\pgfsetfillcolor{currentfill}%
\pgfsetlinewidth{0.401500pt}%
\definecolor{currentstroke}{rgb}{0.000000,0.000000,0.000000}%
\pgfsetstrokecolor{currentstroke}%
\pgfsetdash{}{0pt}%
\pgfsys@defobject{currentmarker}{\pgfqpoint{-0.020833in}{0.000000in}}{\pgfqpoint{-0.000000in}{0.000000in}}{%
\pgfpathmoveto{\pgfqpoint{-0.000000in}{0.000000in}}%
\pgfpathlineto{\pgfqpoint{-0.020833in}{0.000000in}}%
\pgfusepath{stroke,fill}%
}%
\begin{pgfscope}%
\pgfsys@transformshift{3.038110in}{1.811786in}%
\pgfsys@useobject{currentmarker}{}%
\end{pgfscope}%
\end{pgfscope}%
\begin{pgfscope}%
\pgfsetbuttcap%
\pgfsetroundjoin%
\definecolor{currentfill}{rgb}{0.000000,0.000000,0.000000}%
\pgfsetfillcolor{currentfill}%
\pgfsetlinewidth{0.401500pt}%
\definecolor{currentstroke}{rgb}{0.000000,0.000000,0.000000}%
\pgfsetstrokecolor{currentstroke}%
\pgfsetdash{}{0pt}%
\pgfsys@defobject{currentmarker}{\pgfqpoint{0.000000in}{0.000000in}}{\pgfqpoint{0.020833in}{0.000000in}}{%
\pgfpathmoveto{\pgfqpoint{0.000000in}{0.000000in}}%
\pgfpathlineto{\pgfqpoint{0.020833in}{0.000000in}}%
\pgfusepath{stroke,fill}%
}%
\begin{pgfscope}%
\pgfsys@transformshift{0.650000in}{2.009112in}%
\pgfsys@useobject{currentmarker}{}%
\end{pgfscope}%
\end{pgfscope}%
\begin{pgfscope}%
\pgfsetbuttcap%
\pgfsetroundjoin%
\definecolor{currentfill}{rgb}{0.000000,0.000000,0.000000}%
\pgfsetfillcolor{currentfill}%
\pgfsetlinewidth{0.401500pt}%
\definecolor{currentstroke}{rgb}{0.000000,0.000000,0.000000}%
\pgfsetstrokecolor{currentstroke}%
\pgfsetdash{}{0pt}%
\pgfsys@defobject{currentmarker}{\pgfqpoint{-0.020833in}{0.000000in}}{\pgfqpoint{-0.000000in}{0.000000in}}{%
\pgfpathmoveto{\pgfqpoint{-0.000000in}{0.000000in}}%
\pgfpathlineto{\pgfqpoint{-0.020833in}{0.000000in}}%
\pgfusepath{stroke,fill}%
}%
\begin{pgfscope}%
\pgfsys@transformshift{3.038110in}{2.009112in}%
\pgfsys@useobject{currentmarker}{}%
\end{pgfscope}%
\end{pgfscope}%
\begin{pgfscope}%
\pgfsetbuttcap%
\pgfsetroundjoin%
\definecolor{currentfill}{rgb}{0.000000,0.000000,0.000000}%
\pgfsetfillcolor{currentfill}%
\pgfsetlinewidth{0.401500pt}%
\definecolor{currentstroke}{rgb}{0.000000,0.000000,0.000000}%
\pgfsetstrokecolor{currentstroke}%
\pgfsetdash{}{0pt}%
\pgfsys@defobject{currentmarker}{\pgfqpoint{0.000000in}{0.000000in}}{\pgfqpoint{0.020833in}{0.000000in}}{%
\pgfpathmoveto{\pgfqpoint{0.000000in}{0.000000in}}%
\pgfpathlineto{\pgfqpoint{0.020833in}{0.000000in}}%
\pgfusepath{stroke,fill}%
}%
\begin{pgfscope}%
\pgfsys@transformshift{0.650000in}{2.107775in}%
\pgfsys@useobject{currentmarker}{}%
\end{pgfscope}%
\end{pgfscope}%
\begin{pgfscope}%
\pgfsetbuttcap%
\pgfsetroundjoin%
\definecolor{currentfill}{rgb}{0.000000,0.000000,0.000000}%
\pgfsetfillcolor{currentfill}%
\pgfsetlinewidth{0.401500pt}%
\definecolor{currentstroke}{rgb}{0.000000,0.000000,0.000000}%
\pgfsetstrokecolor{currentstroke}%
\pgfsetdash{}{0pt}%
\pgfsys@defobject{currentmarker}{\pgfqpoint{-0.020833in}{0.000000in}}{\pgfqpoint{-0.000000in}{0.000000in}}{%
\pgfpathmoveto{\pgfqpoint{-0.000000in}{0.000000in}}%
\pgfpathlineto{\pgfqpoint{-0.020833in}{0.000000in}}%
\pgfusepath{stroke,fill}%
}%
\begin{pgfscope}%
\pgfsys@transformshift{3.038110in}{2.107775in}%
\pgfsys@useobject{currentmarker}{}%
\end{pgfscope}%
\end{pgfscope}%
\begin{pgfscope}%
\pgfsetbuttcap%
\pgfsetroundjoin%
\definecolor{currentfill}{rgb}{0.000000,0.000000,0.000000}%
\pgfsetfillcolor{currentfill}%
\pgfsetlinewidth{0.401500pt}%
\definecolor{currentstroke}{rgb}{0.000000,0.000000,0.000000}%
\pgfsetstrokecolor{currentstroke}%
\pgfsetdash{}{0pt}%
\pgfsys@defobject{currentmarker}{\pgfqpoint{0.000000in}{0.000000in}}{\pgfqpoint{0.020833in}{0.000000in}}{%
\pgfpathmoveto{\pgfqpoint{0.000000in}{0.000000in}}%
\pgfpathlineto{\pgfqpoint{0.020833in}{0.000000in}}%
\pgfusepath{stroke,fill}%
}%
\begin{pgfscope}%
\pgfsys@transformshift{0.650000in}{2.206438in}%
\pgfsys@useobject{currentmarker}{}%
\end{pgfscope}%
\end{pgfscope}%
\begin{pgfscope}%
\pgfsetbuttcap%
\pgfsetroundjoin%
\definecolor{currentfill}{rgb}{0.000000,0.000000,0.000000}%
\pgfsetfillcolor{currentfill}%
\pgfsetlinewidth{0.401500pt}%
\definecolor{currentstroke}{rgb}{0.000000,0.000000,0.000000}%
\pgfsetstrokecolor{currentstroke}%
\pgfsetdash{}{0pt}%
\pgfsys@defobject{currentmarker}{\pgfqpoint{-0.020833in}{0.000000in}}{\pgfqpoint{-0.000000in}{0.000000in}}{%
\pgfpathmoveto{\pgfqpoint{-0.000000in}{0.000000in}}%
\pgfpathlineto{\pgfqpoint{-0.020833in}{0.000000in}}%
\pgfusepath{stroke,fill}%
}%
\begin{pgfscope}%
\pgfsys@transformshift{3.038110in}{2.206438in}%
\pgfsys@useobject{currentmarker}{}%
\end{pgfscope}%
\end{pgfscope}%
\begin{pgfscope}%
\pgfsetbuttcap%
\pgfsetroundjoin%
\definecolor{currentfill}{rgb}{0.000000,0.000000,0.000000}%
\pgfsetfillcolor{currentfill}%
\pgfsetlinewidth{0.401500pt}%
\definecolor{currentstroke}{rgb}{0.000000,0.000000,0.000000}%
\pgfsetstrokecolor{currentstroke}%
\pgfsetdash{}{0pt}%
\pgfsys@defobject{currentmarker}{\pgfqpoint{0.000000in}{0.000000in}}{\pgfqpoint{0.020833in}{0.000000in}}{%
\pgfpathmoveto{\pgfqpoint{0.000000in}{0.000000in}}%
\pgfpathlineto{\pgfqpoint{0.020833in}{0.000000in}}%
\pgfusepath{stroke,fill}%
}%
\begin{pgfscope}%
\pgfsys@transformshift{0.650000in}{2.305102in}%
\pgfsys@useobject{currentmarker}{}%
\end{pgfscope}%
\end{pgfscope}%
\begin{pgfscope}%
\pgfsetbuttcap%
\pgfsetroundjoin%
\definecolor{currentfill}{rgb}{0.000000,0.000000,0.000000}%
\pgfsetfillcolor{currentfill}%
\pgfsetlinewidth{0.401500pt}%
\definecolor{currentstroke}{rgb}{0.000000,0.000000,0.000000}%
\pgfsetstrokecolor{currentstroke}%
\pgfsetdash{}{0pt}%
\pgfsys@defobject{currentmarker}{\pgfqpoint{-0.020833in}{0.000000in}}{\pgfqpoint{-0.000000in}{0.000000in}}{%
\pgfpathmoveto{\pgfqpoint{-0.000000in}{0.000000in}}%
\pgfpathlineto{\pgfqpoint{-0.020833in}{0.000000in}}%
\pgfusepath{stroke,fill}%
}%
\begin{pgfscope}%
\pgfsys@transformshift{3.038110in}{2.305102in}%
\pgfsys@useobject{currentmarker}{}%
\end{pgfscope}%
\end{pgfscope}%
\begin{pgfscope}%
\definecolor{textcolor}{rgb}{0.000000,0.000000,0.000000}%
\pgfsetstrokecolor{textcolor}%
\pgfsetfillcolor{textcolor}%
\pgftext[x=0.497569in,y=1.417244in,,bottom,rotate=90.000000]{\color{textcolor}{\rmfamily\fontsize{12.000000}{14.400000}\selectfont\catcode`\^=\active\def^{\ifmmode\sp\else\^{}\fi}\catcode`\%=\active\def%{\%}$\delta_{99}/\delta_{0}$}}%
\end{pgfscope}%
\begin{pgfscope}%
\pgfpathrectangle{\pgfqpoint{0.650000in}{0.420000in}}{\pgfqpoint{2.388110in}{1.994488in}}%
\pgfusepath{clip}%
\pgfsetrectcap%
\pgfsetroundjoin%
\pgfsetlinewidth{1.003750pt}%
\definecolor{currentstroke}{rgb}{0.796078,0.094118,0.113725}%
\pgfsetstrokecolor{currentstroke}%
\pgfsetdash{}{0pt}%
\pgfpathmoveto{\pgfqpoint{0.650000in}{0.510659in}}%
\pgfpathlineto{\pgfqpoint{0.655376in}{0.512836in}}%
\pgfpathlineto{\pgfqpoint{0.665412in}{0.516490in}}%
\pgfpathlineto{\pgfqpoint{0.667562in}{0.517781in}}%
\pgfpathlineto{\pgfqpoint{0.674372in}{0.523612in}}%
\pgfpathlineto{\pgfqpoint{0.676881in}{0.527324in}}%
\pgfpathlineto{\pgfqpoint{0.686200in}{0.540883in}}%
\pgfpathlineto{\pgfqpoint{0.689784in}{0.547158in}}%
\pgfpathlineto{\pgfqpoint{0.691576in}{0.550483in}}%
\pgfpathlineto{\pgfqpoint{0.694801in}{0.556481in}}%
\pgfpathlineto{\pgfqpoint{0.696952in}{0.560646in}}%
\pgfpathlineto{\pgfqpoint{0.700536in}{0.567144in}}%
\pgfpathlineto{\pgfqpoint{0.702328in}{0.570594in}}%
\pgfpathlineto{\pgfqpoint{0.705912in}{0.577239in}}%
\pgfpathlineto{\pgfqpoint{0.709138in}{0.583132in}}%
\pgfpathlineto{\pgfqpoint{0.713080in}{0.589759in}}%
\pgfpathlineto{\pgfqpoint{0.715589in}{0.594231in}}%
\pgfpathlineto{\pgfqpoint{0.719532in}{0.600711in}}%
\pgfpathlineto{\pgfqpoint{0.723116in}{0.606613in}}%
\pgfpathlineto{\pgfqpoint{0.727775in}{0.613525in}}%
\pgfpathlineto{\pgfqpoint{0.729926in}{0.617043in}}%
\pgfpathlineto{\pgfqpoint{0.734585in}{0.624039in}}%
\pgfpathlineto{\pgfqpoint{0.738886in}{0.630494in}}%
\pgfpathlineto{\pgfqpoint{0.748205in}{0.643517in}}%
\pgfpathlineto{\pgfqpoint{0.761108in}{0.660015in}}%
\pgfpathlineto{\pgfqpoint{0.764333in}{0.664360in}}%
\pgfpathlineto{\pgfqpoint{0.770426in}{0.671773in}}%
\pgfpathlineto{\pgfqpoint{0.773652in}{0.676082in}}%
\pgfpathlineto{\pgfqpoint{0.780103in}{0.683779in}}%
\pgfpathlineto{\pgfqpoint{0.782971in}{0.687593in}}%
\pgfpathlineto{\pgfqpoint{0.789422in}{0.695003in}}%
\pgfpathlineto{\pgfqpoint{0.793006in}{0.699438in}}%
\pgfpathlineto{\pgfqpoint{0.804834in}{0.713042in}}%
\pgfpathlineto{\pgfqpoint{0.829564in}{0.740091in}}%
\pgfpathlineto{\pgfqpoint{0.835299in}{0.746239in}}%
\pgfpathlineto{\pgfqpoint{0.851069in}{0.762680in}}%
\pgfpathlineto{\pgfqpoint{0.871140in}{0.782843in}}%
\pgfpathlineto{\pgfqpoint{0.903397in}{0.815083in}}%
\pgfpathlineto{\pgfqpoint{0.914866in}{0.826613in}}%
\pgfpathlineto{\pgfqpoint{0.930278in}{0.841691in}}%
\pgfpathlineto{\pgfqpoint{0.974363in}{0.883479in}}%
\pgfpathlineto{\pgfqpoint{0.984398in}{0.892917in}}%
\pgfpathlineto{\pgfqpoint{0.991925in}{0.900252in}}%
\pgfpathlineto{\pgfqpoint{1.042102in}{0.945841in}}%
\pgfpathlineto{\pgfqpoint{1.060740in}{0.961578in}}%
\pgfpathlineto{\pgfqpoint{1.082961in}{0.978012in}}%
\pgfpathlineto{\pgfqpoint{1.090846in}{0.982611in}}%
\pgfpathlineto{\pgfqpoint{1.097656in}{0.985477in}}%
\pgfpathlineto{\pgfqpoint{1.103391in}{0.986708in}}%
\pgfpathlineto{\pgfqpoint{1.109125in}{0.986935in}}%
\pgfpathlineto{\pgfqpoint{1.114143in}{0.985999in}}%
\pgfpathlineto{\pgfqpoint{1.118444in}{0.984221in}}%
\pgfpathlineto{\pgfqpoint{1.121670in}{0.981855in}}%
\pgfpathlineto{\pgfqpoint{1.125612in}{0.978040in}}%
\pgfpathlineto{\pgfqpoint{1.130630in}{0.971472in}}%
\pgfpathlineto{\pgfqpoint{1.137082in}{0.961005in}}%
\pgfpathlineto{\pgfqpoint{1.142816in}{0.952937in}}%
\pgfpathlineto{\pgfqpoint{1.146759in}{0.949113in}}%
\pgfpathlineto{\pgfqpoint{1.151418in}{0.946174in}}%
\pgfpathlineto{\pgfqpoint{1.156078in}{0.944517in}}%
\pgfpathlineto{\pgfqpoint{1.162887in}{0.943516in}}%
\pgfpathlineto{\pgfqpoint{1.168622in}{0.944604in}}%
\pgfpathlineto{\pgfqpoint{1.176149in}{0.947272in}}%
\pgfpathlineto{\pgfqpoint{1.190127in}{0.953706in}}%
\pgfpathlineto{\pgfqpoint{1.201955in}{0.959497in}}%
\pgfpathlineto{\pgfqpoint{1.240305in}{0.980284in}}%
\pgfpathlineto{\pgfqpoint{1.255717in}{0.988929in}}%
\pgfpathlineto{\pgfqpoint{1.266111in}{0.994857in}}%
\pgfpathlineto{\pgfqpoint{1.266828in}{0.994617in}}%
\pgfpathlineto{\pgfqpoint{1.285465in}{1.004434in}}%
\pgfpathlineto{\pgfqpoint{1.286182in}{1.004164in}}%
\pgfpathlineto{\pgfqpoint{1.302311in}{1.013510in}}%
\pgfpathlineto{\pgfqpoint{1.304103in}{1.014543in}}%
\pgfpathlineto{\pgfqpoint{1.304820in}{1.014291in}}%
\pgfpathlineto{\pgfqpoint{1.329550in}{1.027690in}}%
\pgfpathlineto{\pgfqpoint{1.344245in}{1.035600in}}%
\pgfpathlineto{\pgfqpoint{1.344962in}{1.035297in}}%
\pgfpathlineto{\pgfqpoint{1.364675in}{1.045669in}}%
\pgfpathlineto{\pgfqpoint{1.365392in}{1.045392in}}%
\pgfpathlineto{\pgfqpoint{1.384746in}{1.056589in}}%
\pgfpathlineto{\pgfqpoint{1.385463in}{1.056279in}}%
\pgfpathlineto{\pgfqpoint{1.408760in}{1.068146in}}%
\pgfpathlineto{\pgfqpoint{1.409118in}{1.068313in}}%
\pgfpathlineto{\pgfqpoint{1.409835in}{1.067933in}}%
\pgfpathlineto{\pgfqpoint{1.479009in}{1.099690in}}%
\pgfpathlineto{\pgfqpoint{1.479726in}{1.099314in}}%
\pgfpathlineto{\pgfqpoint{1.505173in}{1.111063in}}%
\pgfpathlineto{\pgfqpoint{1.505890in}{1.110631in}}%
\pgfpathlineto{\pgfqpoint{1.523094in}{1.118368in}}%
\pgfpathlineto{\pgfqpoint{1.530621in}{1.122071in}}%
\pgfpathlineto{\pgfqpoint{1.531337in}{1.121672in}}%
\pgfpathlineto{\pgfqpoint{1.555351in}{1.132604in}}%
\pgfpathlineto{\pgfqpoint{1.556068in}{1.132202in}}%
\pgfpathlineto{\pgfqpoint{1.579365in}{1.143469in}}%
\pgfpathlineto{\pgfqpoint{1.580082in}{1.143064in}}%
\pgfpathlineto{\pgfqpoint{1.615923in}{1.160068in}}%
\pgfpathlineto{\pgfqpoint{1.625242in}{1.164102in}}%
\pgfpathlineto{\pgfqpoint{1.625959in}{1.163618in}}%
\pgfpathlineto{\pgfqpoint{1.653198in}{1.174811in}}%
\pgfpathlineto{\pgfqpoint{1.653915in}{1.174315in}}%
\pgfpathlineto{\pgfqpoint{1.671477in}{1.182353in}}%
\pgfpathlineto{\pgfqpoint{1.679721in}{1.185809in}}%
\pgfpathlineto{\pgfqpoint{1.680438in}{1.185276in}}%
\pgfpathlineto{\pgfqpoint{1.696566in}{1.192057in}}%
\pgfpathlineto{\pgfqpoint{1.706602in}{1.196793in}}%
\pgfpathlineto{\pgfqpoint{1.707319in}{1.196347in}}%
\pgfpathlineto{\pgfqpoint{1.730616in}{1.207329in}}%
\pgfpathlineto{\pgfqpoint{1.731333in}{1.206871in}}%
\pgfpathlineto{\pgfqpoint{1.754988in}{1.217757in}}%
\pgfpathlineto{\pgfqpoint{1.755705in}{1.218079in}}%
\pgfpathlineto{\pgfqpoint{1.756422in}{1.217611in}}%
\pgfpathlineto{\pgfqpoint{1.783662in}{1.229975in}}%
\pgfpathlineto{\pgfqpoint{1.784378in}{1.229469in}}%
\pgfpathlineto{\pgfqpoint{1.811976in}{1.242011in}}%
\pgfpathlineto{\pgfqpoint{1.814127in}{1.243004in}}%
\pgfpathlineto{\pgfqpoint{1.814843in}{1.242509in}}%
\pgfpathlineto{\pgfqpoint{1.838499in}{1.253323in}}%
\pgfpathlineto{\pgfqpoint{1.839215in}{1.252780in}}%
\pgfpathlineto{\pgfqpoint{1.863946in}{1.262932in}}%
\pgfpathlineto{\pgfqpoint{1.866455in}{1.264011in}}%
\pgfpathlineto{\pgfqpoint{1.867172in}{1.263472in}}%
\pgfpathlineto{\pgfqpoint{1.888676in}{1.271944in}}%
\pgfpathlineto{\pgfqpoint{1.895486in}{1.274616in}}%
\pgfpathlineto{\pgfqpoint{1.896203in}{1.274031in}}%
\pgfpathlineto{\pgfqpoint{1.923442in}{1.284643in}}%
\pgfpathlineto{\pgfqpoint{1.924159in}{1.284068in}}%
\pgfpathlineto{\pgfqpoint{1.952832in}{1.296243in}}%
\pgfpathlineto{\pgfqpoint{1.953549in}{1.295652in}}%
\pgfpathlineto{\pgfqpoint{1.972903in}{1.302463in}}%
\pgfpathlineto{\pgfqpoint{1.989031in}{1.308110in}}%
\pgfpathlineto{\pgfqpoint{1.989748in}{1.307488in}}%
\pgfpathlineto{\pgfqpoint{2.007310in}{1.314556in}}%
\pgfpathlineto{\pgfqpoint{2.019496in}{1.319308in}}%
\pgfpathlineto{\pgfqpoint{2.020213in}{1.318703in}}%
\pgfpathlineto{\pgfqpoint{2.051036in}{1.331190in}}%
\pgfpathlineto{\pgfqpoint{2.051753in}{1.330591in}}%
\pgfpathlineto{\pgfqpoint{2.080784in}{1.342474in}}%
\pgfpathlineto{\pgfqpoint{2.081501in}{1.341868in}}%
\pgfpathlineto{\pgfqpoint{2.108023in}{1.352186in}}%
\pgfpathlineto{\pgfqpoint{2.108740in}{1.351546in}}%
\pgfpathlineto{\pgfqpoint{2.124152in}{1.356979in}}%
\pgfpathlineto{\pgfqpoint{2.142431in}{1.364346in}}%
\pgfpathlineto{\pgfqpoint{2.153900in}{1.369185in}}%
\pgfpathlineto{\pgfqpoint{2.157126in}{1.370499in}}%
\pgfpathlineto{\pgfqpoint{2.157842in}{1.369838in}}%
\pgfpathlineto{\pgfqpoint{2.186515in}{1.381610in}}%
\pgfpathlineto{\pgfqpoint{2.186873in}{1.380808in}}%
\pgfpathlineto{\pgfqpoint{2.187231in}{1.380973in}}%
\pgfpathlineto{\pgfqpoint{2.212677in}{1.392665in}}%
\pgfpathlineto{\pgfqpoint{2.213394in}{1.391986in}}%
\pgfpathlineto{\pgfqpoint{2.236689in}{1.401447in}}%
\pgfpathlineto{\pgfqpoint{2.246365in}{1.405368in}}%
\pgfpathlineto{\pgfqpoint{2.246724in}{1.404475in}}%
\pgfpathlineto{\pgfqpoint{2.247082in}{1.404622in}}%
\pgfpathlineto{\pgfqpoint{2.277187in}{1.417045in}}%
\pgfpathlineto{\pgfqpoint{2.277545in}{1.416140in}}%
\pgfpathlineto{\pgfqpoint{2.277903in}{1.416291in}}%
\pgfpathlineto{\pgfqpoint{2.304066in}{1.427042in}}%
\pgfpathlineto{\pgfqpoint{2.304424in}{1.426143in}}%
\pgfpathlineto{\pgfqpoint{2.304782in}{1.426287in}}%
\pgfpathlineto{\pgfqpoint{2.321627in}{1.433586in}}%
\pgfpathlineto{\pgfqpoint{2.332378in}{1.438332in}}%
\pgfpathlineto{\pgfqpoint{2.332737in}{1.437403in}}%
\pgfpathlineto{\pgfqpoint{2.333095in}{1.437547in}}%
\pgfpathlineto{\pgfqpoint{2.354957in}{1.445736in}}%
\pgfpathlineto{\pgfqpoint{2.361049in}{1.447825in}}%
\pgfpathlineto{\pgfqpoint{2.361766in}{1.446956in}}%
\pgfpathlineto{\pgfqpoint{2.389003in}{1.455768in}}%
\pgfpathlineto{\pgfqpoint{2.395454in}{1.457507in}}%
\pgfpathlineto{\pgfqpoint{2.396171in}{1.456548in}}%
\pgfpathlineto{\pgfqpoint{2.407640in}{1.460195in}}%
\pgfpathlineto{\pgfqpoint{2.421617in}{1.465695in}}%
\pgfpathlineto{\pgfqpoint{2.430218in}{1.469373in}}%
\pgfpathlineto{\pgfqpoint{2.430576in}{1.468314in}}%
\pgfpathlineto{\pgfqpoint{2.430935in}{1.468453in}}%
\pgfpathlineto{\pgfqpoint{2.449212in}{1.474966in}}%
\pgfpathlineto{\pgfqpoint{2.462115in}{1.479085in}}%
\pgfpathlineto{\pgfqpoint{2.462832in}{1.478074in}}%
\pgfpathlineto{\pgfqpoint{2.482902in}{1.484956in}}%
\pgfpathlineto{\pgfqpoint{2.488994in}{1.487555in}}%
\pgfpathlineto{\pgfqpoint{2.489352in}{1.486417in}}%
\pgfpathlineto{\pgfqpoint{2.489711in}{1.486577in}}%
\pgfpathlineto{\pgfqpoint{2.513007in}{1.497393in}}%
\pgfpathlineto{\pgfqpoint{2.513365in}{1.496296in}}%
\pgfpathlineto{\pgfqpoint{2.513724in}{1.496465in}}%
\pgfpathlineto{\pgfqpoint{2.520533in}{1.500236in}}%
\pgfpathlineto{\pgfqpoint{2.528059in}{1.505691in}}%
\pgfpathlineto{\pgfqpoint{2.533794in}{1.511184in}}%
\pgfpathlineto{\pgfqpoint{2.534152in}{1.511577in}}%
\pgfpathlineto{\pgfqpoint{2.534511in}{1.510567in}}%
\pgfpathlineto{\pgfqpoint{2.534869in}{1.510971in}}%
\pgfpathlineto{\pgfqpoint{2.541320in}{1.519057in}}%
\pgfpathlineto{\pgfqpoint{2.546338in}{1.526809in}}%
\pgfpathlineto{\pgfqpoint{2.553506in}{1.540297in}}%
\pgfpathlineto{\pgfqpoint{2.557806in}{1.552442in}}%
\pgfpathlineto{\pgfqpoint{2.563182in}{1.571335in}}%
\pgfpathlineto{\pgfqpoint{2.570351in}{1.605505in}}%
\pgfpathlineto{\pgfqpoint{2.575011in}{1.626208in}}%
\pgfpathlineto{\pgfqpoint{2.580029in}{1.645284in}}%
\pgfpathlineto{\pgfqpoint{2.583613in}{1.655398in}}%
\pgfpathlineto{\pgfqpoint{2.588990in}{1.668925in}}%
\pgfpathlineto{\pgfqpoint{2.593291in}{1.677804in}}%
\pgfpathlineto{\pgfqpoint{2.593649in}{1.677355in}}%
\pgfpathlineto{\pgfqpoint{2.599743in}{1.687593in}}%
\pgfpathlineto{\pgfqpoint{2.605478in}{1.695116in}}%
\pgfpathlineto{\pgfqpoint{2.607270in}{1.697221in}}%
\pgfpathlineto{\pgfqpoint{2.607628in}{1.696305in}}%
\pgfpathlineto{\pgfqpoint{2.607987in}{1.696706in}}%
\pgfpathlineto{\pgfqpoint{2.614797in}{1.703396in}}%
\pgfpathlineto{\pgfqpoint{2.623041in}{1.710016in}}%
\pgfpathlineto{\pgfqpoint{2.625192in}{1.711427in}}%
\pgfpathlineto{\pgfqpoint{2.625550in}{1.710136in}}%
\pgfpathlineto{\pgfqpoint{2.625908in}{1.710360in}}%
\pgfpathlineto{\pgfqpoint{2.636303in}{1.716235in}}%
\pgfpathlineto{\pgfqpoint{2.646697in}{1.721649in}}%
\pgfpathlineto{\pgfqpoint{2.647056in}{1.719975in}}%
\pgfpathlineto{\pgfqpoint{2.647414in}{1.720137in}}%
\pgfpathlineto{\pgfqpoint{2.661752in}{1.725990in}}%
\pgfpathlineto{\pgfqpoint{2.670354in}{1.729784in}}%
\pgfpathlineto{\pgfqpoint{2.672146in}{1.730685in}}%
\pgfpathlineto{\pgfqpoint{2.672505in}{1.728692in}}%
\pgfpathlineto{\pgfqpoint{2.672863in}{1.728855in}}%
\pgfpathlineto{\pgfqpoint{2.688993in}{1.735558in}}%
\pgfpathlineto{\pgfqpoint{2.694728in}{1.737859in}}%
\pgfpathlineto{\pgfqpoint{2.695086in}{1.735440in}}%
\pgfpathlineto{\pgfqpoint{2.695444in}{1.735580in}}%
\pgfpathlineto{\pgfqpoint{2.717309in}{1.744462in}}%
\pgfpathlineto{\pgfqpoint{2.717667in}{1.741448in}}%
\pgfpathlineto{\pgfqpoint{2.718026in}{1.741601in}}%
\pgfpathlineto{\pgfqpoint{2.736306in}{1.748984in}}%
\pgfpathlineto{\pgfqpoint{2.741682in}{1.751051in}}%
\pgfpathlineto{\pgfqpoint{2.742041in}{1.747507in}}%
\pgfpathlineto{\pgfqpoint{2.742399in}{1.747651in}}%
\pgfpathlineto{\pgfqpoint{2.757095in}{1.754265in}}%
\pgfpathlineto{\pgfqpoint{2.764981in}{1.757686in}}%
\pgfpathlineto{\pgfqpoint{2.765339in}{1.753276in}}%
\pgfpathlineto{\pgfqpoint{2.765697in}{1.753427in}}%
\pgfpathlineto{\pgfqpoint{2.776451in}{1.758612in}}%
\pgfpathlineto{\pgfqpoint{2.783261in}{1.762184in}}%
\pgfpathlineto{\pgfqpoint{2.783619in}{1.756868in}}%
\pgfpathlineto{\pgfqpoint{2.783978in}{1.757043in}}%
\pgfpathlineto{\pgfqpoint{2.801541in}{1.765975in}}%
\pgfpathlineto{\pgfqpoint{2.801899in}{1.759516in}}%
\pgfpathlineto{\pgfqpoint{2.802258in}{1.759711in}}%
\pgfpathlineto{\pgfqpoint{2.808351in}{1.762369in}}%
\pgfpathlineto{\pgfqpoint{2.809785in}{1.762729in}}%
\pgfpathlineto{\pgfqpoint{2.809785in}{1.762729in}}%
\pgfusepath{stroke}%
\end{pgfscope}%
\begin{pgfscope}%
\pgfpathrectangle{\pgfqpoint{0.650000in}{0.420000in}}{\pgfqpoint{2.388110in}{1.994488in}}%
\pgfusepath{clip}%
\pgfsetrectcap%
\pgfsetroundjoin%
\pgfsetlinewidth{1.003750pt}%
\definecolor{currentstroke}{rgb}{0.454902,0.678431,0.819608}%
\pgfsetstrokecolor{currentstroke}%
\pgfsetdash{}{0pt}%
\pgfpathmoveto{\pgfqpoint{0.650000in}{0.510659in}}%
\pgfpathlineto{\pgfqpoint{0.655376in}{0.512835in}}%
\pgfpathlineto{\pgfqpoint{0.665412in}{0.516485in}}%
\pgfpathlineto{\pgfqpoint{0.667562in}{0.517778in}}%
\pgfpathlineto{\pgfqpoint{0.674372in}{0.523669in}}%
\pgfpathlineto{\pgfqpoint{0.681540in}{0.533959in}}%
\pgfpathlineto{\pgfqpoint{0.702328in}{0.570521in}}%
\pgfpathlineto{\pgfqpoint{0.705912in}{0.577095in}}%
\pgfpathlineto{\pgfqpoint{0.709138in}{0.582869in}}%
\pgfpathlineto{\pgfqpoint{0.713080in}{0.589499in}}%
\pgfpathlineto{\pgfqpoint{0.715589in}{0.593812in}}%
\pgfpathlineto{\pgfqpoint{0.719532in}{0.600369in}}%
\pgfpathlineto{\pgfqpoint{0.723474in}{0.606847in}}%
\pgfpathlineto{\pgfqpoint{0.728134in}{0.613771in}}%
\pgfpathlineto{\pgfqpoint{0.729926in}{0.616798in}}%
\pgfpathlineto{\pgfqpoint{0.734585in}{0.623918in}}%
\pgfpathlineto{\pgfqpoint{0.738528in}{0.629956in}}%
\pgfpathlineto{\pgfqpoint{0.743545in}{0.636867in}}%
\pgfpathlineto{\pgfqpoint{0.746054in}{0.640649in}}%
\pgfpathlineto{\pgfqpoint{0.751431in}{0.647896in}}%
\pgfpathlineto{\pgfqpoint{0.755373in}{0.653348in}}%
\pgfpathlineto{\pgfqpoint{0.760749in}{0.659792in}}%
\pgfpathlineto{\pgfqpoint{0.764333in}{0.664430in}}%
\pgfpathlineto{\pgfqpoint{0.770785in}{0.672048in}}%
\pgfpathlineto{\pgfqpoint{0.773294in}{0.675394in}}%
\pgfpathlineto{\pgfqpoint{0.779387in}{0.682656in}}%
\pgfpathlineto{\pgfqpoint{0.784046in}{0.688355in}}%
\pgfpathlineto{\pgfqpoint{0.799816in}{0.706905in}}%
\pgfpathlineto{\pgfqpoint{0.803400in}{0.711416in}}%
\pgfpathlineto{\pgfqpoint{0.817378in}{0.727622in}}%
\pgfpathlineto{\pgfqpoint{0.834582in}{0.746811in}}%
\pgfpathlineto{\pgfqpoint{0.850711in}{0.763735in}}%
\pgfpathlineto{\pgfqpoint{0.856445in}{0.770156in}}%
\pgfpathlineto{\pgfqpoint{0.869706in}{0.784336in}}%
\pgfpathlineto{\pgfqpoint{0.885835in}{0.801156in}}%
\pgfpathlineto{\pgfqpoint{0.895154in}{0.811644in}}%
\pgfpathlineto{\pgfqpoint{0.926694in}{0.844623in}}%
\pgfpathlineto{\pgfqpoint{0.934221in}{0.851779in}}%
\pgfpathlineto{\pgfqpoint{0.947840in}{0.863413in}}%
\pgfpathlineto{\pgfqpoint{0.966119in}{0.878268in}}%
\pgfpathlineto{\pgfqpoint{0.978305in}{0.885615in}}%
\pgfpathlineto{\pgfqpoint{0.986907in}{0.889873in}}%
\pgfpathlineto{\pgfqpoint{0.997301in}{0.893770in}}%
\pgfpathlineto{\pgfqpoint{1.003394in}{0.894955in}}%
\pgfpathlineto{\pgfqpoint{1.008412in}{0.894869in}}%
\pgfpathlineto{\pgfqpoint{1.014146in}{0.893647in}}%
\pgfpathlineto{\pgfqpoint{1.021315in}{0.891407in}}%
\pgfpathlineto{\pgfqpoint{1.028483in}{0.888767in}}%
\pgfpathlineto{\pgfqpoint{1.032067in}{0.888997in}}%
\pgfpathlineto{\pgfqpoint{1.034934in}{0.889738in}}%
\pgfpathlineto{\pgfqpoint{1.056797in}{0.891686in}}%
\pgfpathlineto{\pgfqpoint{1.087621in}{0.893190in}}%
\pgfpathlineto{\pgfqpoint{1.133856in}{0.895492in}}%
\pgfpathlineto{\pgfqpoint{1.155719in}{0.897756in}}%
\pgfpathlineto{\pgfqpoint{1.159303in}{0.900411in}}%
\pgfpathlineto{\pgfqpoint{1.165038in}{0.906312in}}%
\pgfpathlineto{\pgfqpoint{1.180808in}{0.925123in}}%
\pgfpathlineto{\pgfqpoint{1.191202in}{0.937770in}}%
\pgfpathlineto{\pgfqpoint{1.202313in}{0.950649in}}%
\pgfpathlineto{\pgfqpoint{1.212349in}{0.959706in}}%
\pgfpathlineto{\pgfqpoint{1.213065in}{0.959766in}}%
\pgfpathlineto{\pgfqpoint{1.218083in}{0.963229in}}%
\pgfpathlineto{\pgfqpoint{1.218800in}{0.963282in}}%
\pgfpathlineto{\pgfqpoint{1.224893in}{0.967630in}}%
\pgfpathlineto{\pgfqpoint{1.225610in}{0.967649in}}%
\pgfpathlineto{\pgfqpoint{1.229553in}{0.969833in}}%
\pgfpathlineto{\pgfqpoint{1.237438in}{0.974167in}}%
\pgfpathlineto{\pgfqpoint{1.251416in}{0.980740in}}%
\pgfpathlineto{\pgfqpoint{1.252491in}{0.981748in}}%
\pgfpathlineto{\pgfqpoint{1.261810in}{0.986860in}}%
\pgfpathlineto{\pgfqpoint{1.262885in}{0.987855in}}%
\pgfpathlineto{\pgfqpoint{1.276147in}{0.992749in}}%
\pgfpathlineto{\pgfqpoint{1.284748in}{0.995450in}}%
\pgfpathlineto{\pgfqpoint{1.287616in}{0.996050in}}%
\pgfpathlineto{\pgfqpoint{1.288333in}{0.995561in}}%
\pgfpathlineto{\pgfqpoint{1.317364in}{1.001337in}}%
\pgfpathlineto{\pgfqpoint{1.322024in}{1.003056in}}%
\pgfpathlineto{\pgfqpoint{1.322740in}{1.002591in}}%
\pgfpathlineto{\pgfqpoint{1.336719in}{1.008756in}}%
\pgfpathlineto{\pgfqpoint{1.351055in}{1.017075in}}%
\pgfpathlineto{\pgfqpoint{1.359299in}{1.023969in}}%
\pgfpathlineto{\pgfqpoint{1.373635in}{1.037433in}}%
\pgfpathlineto{\pgfqpoint{1.389764in}{1.055262in}}%
\pgfpathlineto{\pgfqpoint{1.396215in}{1.060493in}}%
\pgfpathlineto{\pgfqpoint{1.401950in}{1.063717in}}%
\pgfpathlineto{\pgfqpoint{1.405534in}{1.065386in}}%
\pgfpathlineto{\pgfqpoint{1.413061in}{1.069047in}}%
\pgfpathlineto{\pgfqpoint{1.415928in}{1.069714in}}%
\pgfpathlineto{\pgfqpoint{1.424889in}{1.072788in}}%
\pgfpathlineto{\pgfqpoint{1.445677in}{1.075453in}}%
\pgfpathlineto{\pgfqpoint{1.450336in}{1.075857in}}%
\pgfpathlineto{\pgfqpoint{1.451053in}{1.075305in}}%
\pgfpathlineto{\pgfqpoint{1.458938in}{1.076639in}}%
\pgfpathlineto{\pgfqpoint{1.483668in}{1.083572in}}%
\pgfpathlineto{\pgfqpoint{1.503381in}{1.093105in}}%
\pgfpathlineto{\pgfqpoint{1.528829in}{1.109197in}}%
\pgfpathlineto{\pgfqpoint{1.538147in}{1.115115in}}%
\pgfpathlineto{\pgfqpoint{1.538864in}{1.114912in}}%
\pgfpathlineto{\pgfqpoint{1.562878in}{1.130303in}}%
\pgfpathlineto{\pgfqpoint{1.563595in}{1.130158in}}%
\pgfpathlineto{\pgfqpoint{1.585458in}{1.144033in}}%
\pgfpathlineto{\pgfqpoint{1.586175in}{1.143780in}}%
\pgfpathlineto{\pgfqpoint{1.592985in}{1.146539in}}%
\pgfpathlineto{\pgfqpoint{1.611264in}{1.152489in}}%
\pgfpathlineto{\pgfqpoint{1.611981in}{1.152028in}}%
\pgfpathlineto{\pgfqpoint{1.632410in}{1.157531in}}%
\pgfpathlineto{\pgfqpoint{1.633127in}{1.157004in}}%
\pgfpathlineto{\pgfqpoint{1.647822in}{1.160118in}}%
\pgfpathlineto{\pgfqpoint{1.661084in}{1.163060in}}%
\pgfpathlineto{\pgfqpoint{1.662517in}{1.163498in}}%
\pgfpathlineto{\pgfqpoint{1.663234in}{1.162978in}}%
\pgfpathlineto{\pgfqpoint{1.672911in}{1.166831in}}%
\pgfpathlineto{\pgfqpoint{1.704452in}{1.182158in}}%
\pgfpathlineto{\pgfqpoint{1.705527in}{1.182728in}}%
\pgfpathlineto{\pgfqpoint{1.706244in}{1.182425in}}%
\pgfpathlineto{\pgfqpoint{1.729541in}{1.195109in}}%
\pgfpathlineto{\pgfqpoint{1.746028in}{1.204059in}}%
\pgfpathlineto{\pgfqpoint{1.764307in}{1.212054in}}%
\pgfpathlineto{\pgfqpoint{1.766457in}{1.212889in}}%
\pgfpathlineto{\pgfqpoint{1.767174in}{1.212549in}}%
\pgfpathlineto{\pgfqpoint{1.781511in}{1.217211in}}%
\pgfpathlineto{\pgfqpoint{1.797639in}{1.221937in}}%
\pgfpathlineto{\pgfqpoint{1.798356in}{1.221462in}}%
\pgfpathlineto{\pgfqpoint{1.825237in}{1.228883in}}%
\pgfpathlineto{\pgfqpoint{1.827030in}{1.229620in}}%
\pgfpathlineto{\pgfqpoint{1.827746in}{1.229257in}}%
\pgfpathlineto{\pgfqpoint{1.836707in}{1.233887in}}%
\pgfpathlineto{\pgfqpoint{1.845667in}{1.238422in}}%
\pgfpathlineto{\pgfqpoint{1.846384in}{1.238102in}}%
\pgfpathlineto{\pgfqpoint{1.867530in}{1.249878in}}%
\pgfpathlineto{\pgfqpoint{1.885092in}{1.261331in}}%
\pgfpathlineto{\pgfqpoint{1.894411in}{1.269508in}}%
\pgfpathlineto{\pgfqpoint{1.911257in}{1.283817in}}%
\pgfpathlineto{\pgfqpoint{1.911973in}{1.283921in}}%
\pgfpathlineto{\pgfqpoint{1.920217in}{1.290466in}}%
\pgfpathlineto{\pgfqpoint{1.927385in}{1.295711in}}%
\pgfpathlineto{\pgfqpoint{1.932044in}{1.297931in}}%
\pgfpathlineto{\pgfqpoint{1.932761in}{1.297723in}}%
\pgfpathlineto{\pgfqpoint{1.940646in}{1.299819in}}%
\pgfpathlineto{\pgfqpoint{1.945664in}{1.299771in}}%
\pgfpathlineto{\pgfqpoint{1.950682in}{1.298536in}}%
\pgfpathlineto{\pgfqpoint{1.957850in}{1.295411in}}%
\pgfpathlineto{\pgfqpoint{1.962868in}{1.292869in}}%
\pgfpathlineto{\pgfqpoint{1.963943in}{1.291613in}}%
\pgfpathlineto{\pgfqpoint{1.976487in}{1.285596in}}%
\pgfpathlineto{\pgfqpoint{1.981505in}{1.284331in}}%
\pgfpathlineto{\pgfqpoint{1.987956in}{1.284088in}}%
\pgfpathlineto{\pgfqpoint{1.989390in}{1.284134in}}%
\pgfpathlineto{\pgfqpoint{1.990107in}{1.283267in}}%
\pgfpathlineto{\pgfqpoint{1.998708in}{1.284193in}}%
\pgfpathlineto{\pgfqpoint{2.003009in}{1.285654in}}%
\pgfpathlineto{\pgfqpoint{2.008027in}{1.288713in}}%
\pgfpathlineto{\pgfqpoint{2.014478in}{1.294117in}}%
\pgfpathlineto{\pgfqpoint{2.048886in}{1.324591in}}%
\pgfpathlineto{\pgfqpoint{2.057130in}{1.329859in}}%
\pgfpathlineto{\pgfqpoint{2.066090in}{1.334456in}}%
\pgfpathlineto{\pgfqpoint{2.075767in}{1.337945in}}%
\pgfpathlineto{\pgfqpoint{2.084010in}{1.340185in}}%
\pgfpathlineto{\pgfqpoint{2.084727in}{1.339912in}}%
\pgfpathlineto{\pgfqpoint{2.092254in}{1.340818in}}%
\pgfpathlineto{\pgfqpoint{2.104439in}{1.342877in}}%
\pgfpathlineto{\pgfqpoint{2.113041in}{1.345342in}}%
\pgfpathlineto{\pgfqpoint{2.120926in}{1.347456in}}%
\pgfpathlineto{\pgfqpoint{2.121643in}{1.347166in}}%
\pgfpathlineto{\pgfqpoint{2.143506in}{1.352251in}}%
\pgfpathlineto{\pgfqpoint{2.144223in}{1.351858in}}%
\pgfpathlineto{\pgfqpoint{2.157126in}{1.354881in}}%
\pgfpathlineto{\pgfqpoint{2.167878in}{1.357432in}}%
\pgfpathlineto{\pgfqpoint{2.178988in}{1.359481in}}%
\pgfpathlineto{\pgfqpoint{2.179705in}{1.358971in}}%
\pgfpathlineto{\pgfqpoint{2.190815in}{1.362334in}}%
\pgfpathlineto{\pgfqpoint{2.203359in}{1.367199in}}%
\pgfpathlineto{\pgfqpoint{2.203717in}{1.367352in}}%
\pgfpathlineto{\pgfqpoint{2.204434in}{1.366971in}}%
\pgfpathlineto{\pgfqpoint{2.214469in}{1.372059in}}%
\pgfpathlineto{\pgfqpoint{2.230955in}{1.380899in}}%
\pgfpathlineto{\pgfqpoint{2.231672in}{1.380581in}}%
\pgfpathlineto{\pgfqpoint{2.262851in}{1.394243in}}%
\pgfpathlineto{\pgfqpoint{2.278979in}{1.401327in}}%
\pgfpathlineto{\pgfqpoint{2.306216in}{1.415615in}}%
\pgfpathlineto{\pgfqpoint{2.319118in}{1.421455in}}%
\pgfpathlineto{\pgfqpoint{2.335245in}{1.428473in}}%
\pgfpathlineto{\pgfqpoint{2.345639in}{1.433966in}}%
\pgfpathlineto{\pgfqpoint{2.353881in}{1.437004in}}%
\pgfpathlineto{\pgfqpoint{2.369292in}{1.441356in}}%
\pgfpathlineto{\pgfqpoint{2.392229in}{1.448493in}}%
\pgfpathlineto{\pgfqpoint{2.407640in}{1.456938in}}%
\pgfpathlineto{\pgfqpoint{2.415166in}{1.462152in}}%
\pgfpathlineto{\pgfqpoint{2.423050in}{1.468897in}}%
\pgfpathlineto{\pgfqpoint{2.434519in}{1.480004in}}%
\pgfpathlineto{\pgfqpoint{2.449929in}{1.497620in}}%
\pgfpathlineto{\pgfqpoint{2.468207in}{1.523350in}}%
\pgfpathlineto{\pgfqpoint{2.481826in}{1.545469in}}%
\pgfpathlineto{\pgfqpoint{2.491861in}{1.562598in}}%
\pgfpathlineto{\pgfqpoint{2.507989in}{1.590314in}}%
\pgfpathlineto{\pgfqpoint{2.523400in}{1.616621in}}%
\pgfpathlineto{\pgfqpoint{2.528418in}{1.624047in}}%
\pgfpathlineto{\pgfqpoint{2.544904in}{1.645381in}}%
\pgfpathlineto{\pgfqpoint{2.545621in}{1.645958in}}%
\pgfpathlineto{\pgfqpoint{2.552789in}{1.656399in}}%
\pgfpathlineto{\pgfqpoint{2.562466in}{1.673841in}}%
\pgfpathlineto{\pgfqpoint{2.570710in}{1.693386in}}%
\pgfpathlineto{\pgfqpoint{2.583613in}{1.732841in}}%
\pgfpathlineto{\pgfqpoint{2.595083in}{1.770185in}}%
\pgfpathlineto{\pgfqpoint{2.602969in}{1.794277in}}%
\pgfpathlineto{\pgfqpoint{2.612288in}{1.820350in}}%
\pgfpathlineto{\pgfqpoint{2.620532in}{1.840769in}}%
\pgfpathlineto{\pgfqpoint{2.628776in}{1.858109in}}%
\pgfpathlineto{\pgfqpoint{2.636303in}{1.872433in}}%
\pgfpathlineto{\pgfqpoint{2.647773in}{1.893610in}}%
\pgfpathlineto{\pgfqpoint{2.653866in}{1.902615in}}%
\pgfpathlineto{\pgfqpoint{2.663544in}{1.914705in}}%
\pgfpathlineto{\pgfqpoint{2.684692in}{1.935702in}}%
\pgfpathlineto{\pgfqpoint{2.695803in}{1.946189in}}%
\pgfpathlineto{\pgfqpoint{2.718026in}{1.970073in}}%
\pgfpathlineto{\pgfqpoint{2.724119in}{1.977871in}}%
\pgfpathlineto{\pgfqpoint{2.733797in}{1.991821in}}%
\pgfpathlineto{\pgfqpoint{2.742399in}{2.004433in}}%
\pgfpathlineto{\pgfqpoint{2.752794in}{2.016393in}}%
\pgfpathlineto{\pgfqpoint{2.768923in}{2.034994in}}%
\pgfpathlineto{\pgfqpoint{2.792939in}{2.070882in}}%
\pgfpathlineto{\pgfqpoint{2.797598in}{2.080137in}}%
\pgfpathlineto{\pgfqpoint{2.806917in}{2.100221in}}%
\pgfpathlineto{\pgfqpoint{2.809785in}{2.103207in}}%
\pgfpathlineto{\pgfqpoint{2.809785in}{2.103207in}}%
\pgfusepath{stroke}%
\end{pgfscope}%
\begin{pgfscope}%
\pgfpathrectangle{\pgfqpoint{0.650000in}{0.420000in}}{\pgfqpoint{2.388110in}{1.994488in}}%
\pgfusepath{clip}%
\pgfsetrectcap%
\pgfsetroundjoin%
\pgfsetlinewidth{1.003750pt}%
\definecolor{currentstroke}{rgb}{0.878431,0.509804,0.078431}%
\pgfsetstrokecolor{currentstroke}%
\pgfsetdash{}{0pt}%
\pgfpathmoveto{\pgfqpoint{0.650000in}{0.511804in}}%
\pgfpathlineto{\pgfqpoint{0.662804in}{0.516912in}}%
\pgfpathlineto{\pgfqpoint{0.667776in}{0.519991in}}%
\pgfpathlineto{\pgfqpoint{0.672335in}{0.525235in}}%
\pgfpathlineto{\pgfqpoint{0.676893in}{0.532038in}}%
\pgfpathlineto{\pgfqpoint{0.686010in}{0.547907in}}%
\pgfpathlineto{\pgfqpoint{0.692226in}{0.559789in}}%
\pgfpathlineto{\pgfqpoint{0.705901in}{0.584208in}}%
\pgfpathlineto{\pgfqpoint{0.711288in}{0.593468in}}%
\pgfpathlineto{\pgfqpoint{0.720404in}{0.608303in}}%
\pgfpathlineto{\pgfqpoint{0.737395in}{0.633479in}}%
\pgfpathlineto{\pgfqpoint{0.744854in}{0.643785in}}%
\pgfpathlineto{\pgfqpoint{0.752313in}{0.653713in}}%
\pgfpathlineto{\pgfqpoint{0.765988in}{0.670857in}}%
\pgfpathlineto{\pgfqpoint{0.783393in}{0.691376in}}%
\pgfpathlineto{\pgfqpoint{0.793752in}{0.703110in}}%
\pgfpathlineto{\pgfqpoint{0.806184in}{0.716521in}}%
\pgfpathlineto{\pgfqpoint{0.814472in}{0.725402in}}%
\pgfpathlineto{\pgfqpoint{0.822760in}{0.734431in}}%
\pgfpathlineto{\pgfqpoint{0.834363in}{0.746915in}}%
\pgfpathlineto{\pgfqpoint{0.957853in}{0.873330in}}%
\pgfpathlineto{\pgfqpoint{1.038660in}{0.949863in}}%
\pgfpathlineto{\pgfqpoint{1.072226in}{0.979304in}}%
\pgfpathlineto{\pgfqpoint{1.091288in}{0.993097in}}%
\pgfpathlineto{\pgfqpoint{1.109936in}{1.004440in}}%
\pgfpathlineto{\pgfqpoint{1.116981in}{1.007581in}}%
\pgfpathlineto{\pgfqpoint{1.122368in}{1.008903in}}%
\pgfpathlineto{\pgfqpoint{1.131070in}{1.009720in}}%
\pgfpathlineto{\pgfqpoint{1.138115in}{1.010931in}}%
\pgfpathlineto{\pgfqpoint{1.146403in}{1.013736in}}%
\pgfpathlineto{\pgfqpoint{1.153862in}{1.017691in}}%
\pgfpathlineto{\pgfqpoint{1.172509in}{1.029851in}}%
\pgfpathlineto{\pgfqpoint{1.187842in}{1.041328in}}%
\pgfpathlineto{\pgfqpoint{1.205661in}{1.055549in}}%
\pgfpathlineto{\pgfqpoint{1.248758in}{1.087912in}}%
\pgfpathlineto{\pgfqpoint{1.306359in}{1.130988in}}%
\pgfpathlineto{\pgfqpoint{1.319620in}{1.140649in}}%
\pgfpathlineto{\pgfqpoint{1.370591in}{1.176924in}}%
\pgfpathlineto{\pgfqpoint{1.402913in}{1.198823in}}%
\pgfpathlineto{\pgfqpoint{1.429435in}{1.216460in}}%
\pgfpathlineto{\pgfqpoint{1.447254in}{1.228556in}}%
\pgfpathlineto{\pgfqpoint{1.448911in}{1.229579in}}%
\pgfpathlineto{\pgfqpoint{1.469217in}{1.244989in}}%
\pgfpathlineto{\pgfqpoint{1.480405in}{1.253627in}}%
\pgfpathlineto{\pgfqpoint{1.482063in}{1.254632in}}%
\pgfpathlineto{\pgfqpoint{1.503612in}{1.269093in}}%
\pgfpathlineto{\pgfqpoint{1.536349in}{1.289297in}}%
\pgfpathlineto{\pgfqpoint{1.570744in}{1.311703in}}%
\pgfpathlineto{\pgfqpoint{1.592292in}{1.326809in}}%
\pgfpathlineto{\pgfqpoint{1.605553in}{1.337117in}}%
\pgfpathlineto{\pgfqpoint{1.616742in}{1.344437in}}%
\pgfpathlineto{\pgfqpoint{1.642434in}{1.360078in}}%
\pgfpathlineto{\pgfqpoint{1.652380in}{1.364992in}}%
\pgfpathlineto{\pgfqpoint{1.668541in}{1.373694in}}%
\pgfpathlineto{\pgfqpoint{1.680559in}{1.381673in}}%
\pgfpathlineto{\pgfqpoint{1.695891in}{1.391604in}}%
\pgfpathlineto{\pgfqpoint{1.710810in}{1.400685in}}%
\pgfpathlineto{\pgfqpoint{1.712053in}{1.401300in}}%
\pgfpathlineto{\pgfqpoint{1.725728in}{1.409987in}}%
\pgfpathlineto{\pgfqpoint{1.741889in}{1.420190in}}%
\pgfpathlineto{\pgfqpoint{1.754735in}{1.429017in}}%
\pgfpathlineto{\pgfqpoint{1.782500in}{1.446142in}}%
\pgfpathlineto{\pgfqpoint{1.794103in}{1.453340in}}%
\pgfpathlineto{\pgfqpoint{1.827669in}{1.477662in}}%
\pgfpathlineto{\pgfqpoint{1.877811in}{1.516223in}}%
\pgfpathlineto{\pgfqpoint{1.891072in}{1.523837in}}%
\pgfpathlineto{\pgfqpoint{1.910134in}{1.535684in}}%
\pgfpathlineto{\pgfqpoint{1.923394in}{1.544211in}}%
\pgfpathlineto{\pgfqpoint{1.934169in}{1.549340in}}%
\pgfpathlineto{\pgfqpoint{1.964834in}{1.564305in}}%
\pgfpathlineto{\pgfqpoint{1.968564in}{1.565247in}}%
\pgfpathlineto{\pgfqpoint{1.971879in}{1.567869in}}%
\pgfpathlineto{\pgfqpoint{1.976437in}{1.571711in}}%
\pgfpathlineto{\pgfqpoint{1.985968in}{1.577590in}}%
\pgfpathlineto{\pgfqpoint{2.036525in}{1.608074in}}%
\pgfpathlineto{\pgfqpoint{2.049371in}{1.616461in}}%
\pgfpathlineto{\pgfqpoint{2.063046in}{1.625525in}}%
\pgfpathlineto{\pgfqpoint{2.072163in}{1.631605in}}%
\pgfpathlineto{\pgfqpoint{2.097855in}{1.647845in}}%
\pgfpathlineto{\pgfqpoint{2.107386in}{1.652591in}}%
\pgfpathlineto{\pgfqpoint{2.138051in}{1.667916in}}%
\pgfpathlineto{\pgfqpoint{2.141781in}{1.668902in}}%
\pgfpathlineto{\pgfqpoint{2.148411in}{1.669465in}}%
\pgfpathlineto{\pgfqpoint{2.156285in}{1.675501in}}%
\pgfpathlineto{\pgfqpoint{2.157114in}{1.675572in}}%
\pgfpathlineto{\pgfqpoint{2.165402in}{1.678887in}}%
\pgfpathlineto{\pgfqpoint{2.180320in}{1.683975in}}%
\pgfpathlineto{\pgfqpoint{2.193580in}{1.690514in}}%
\pgfpathlineto{\pgfqpoint{2.208499in}{1.698199in}}%
\pgfpathlineto{\pgfqpoint{2.225074in}{1.706724in}}%
\pgfpathlineto{\pgfqpoint{2.247037in}{1.720346in}}%
\pgfpathlineto{\pgfqpoint{2.247866in}{1.720525in}}%
\pgfpathlineto{\pgfqpoint{2.259055in}{1.726692in}}%
\pgfpathlineto{\pgfqpoint{2.272316in}{1.734127in}}%
\pgfpathlineto{\pgfqpoint{2.279360in}{1.737767in}}%
\pgfpathlineto{\pgfqpoint{2.298837in}{1.749359in}}%
\pgfpathlineto{\pgfqpoint{2.311683in}{1.756833in}}%
\pgfpathlineto{\pgfqpoint{2.317899in}{1.759989in}}%
\pgfpathlineto{\pgfqpoint{2.318728in}{1.759994in}}%
\pgfpathlineto{\pgfqpoint{2.333646in}{1.765708in}}%
\pgfpathlineto{\pgfqpoint{2.334475in}{1.765682in}}%
\pgfpathlineto{\pgfqpoint{2.344835in}{1.770477in}}%
\pgfpathlineto{\pgfqpoint{2.359339in}{1.778392in}}%
\pgfpathlineto{\pgfqpoint{2.374257in}{1.786821in}}%
\pgfpathlineto{\pgfqpoint{2.375086in}{1.786908in}}%
\pgfpathlineto{\pgfqpoint{2.385446in}{1.791574in}}%
\pgfpathlineto{\pgfqpoint{2.399121in}{1.797996in}}%
\pgfpathlineto{\pgfqpoint{2.423570in}{1.811985in}}%
\pgfpathlineto{\pgfqpoint{2.424399in}{1.812126in}}%
\pgfpathlineto{\pgfqpoint{2.438488in}{1.819913in}}%
\pgfpathlineto{\pgfqpoint{2.439317in}{1.820027in}}%
\pgfpathlineto{\pgfqpoint{2.450506in}{1.827220in}}%
\pgfpathlineto{\pgfqpoint{2.463767in}{1.835694in}}%
\pgfpathlineto{\pgfqpoint{2.464595in}{1.835883in}}%
\pgfpathlineto{\pgfqpoint{2.478685in}{1.844626in}}%
\pgfpathlineto{\pgfqpoint{2.479513in}{1.844801in}}%
\pgfpathlineto{\pgfqpoint{2.511008in}{1.866357in}}%
\pgfpathlineto{\pgfqpoint{2.519710in}{1.875323in}}%
\pgfpathlineto{\pgfqpoint{2.530484in}{1.889782in}}%
\pgfpathlineto{\pgfqpoint{2.539601in}{1.904599in}}%
\pgfpathlineto{\pgfqpoint{2.546231in}{1.917840in}}%
\pgfpathlineto{\pgfqpoint{2.560320in}{1.948419in}}%
\pgfpathlineto{\pgfqpoint{2.567364in}{1.961230in}}%
\pgfpathlineto{\pgfqpoint{2.574408in}{1.972259in}}%
\pgfpathlineto{\pgfqpoint{2.581038in}{1.981025in}}%
\pgfpathlineto{\pgfqpoint{2.582281in}{1.982224in}}%
\pgfpathlineto{\pgfqpoint{2.589740in}{1.990000in}}%
\pgfpathlineto{\pgfqpoint{2.596369in}{1.995733in}}%
\pgfpathlineto{\pgfqpoint{2.611286in}{2.007436in}}%
\pgfpathlineto{\pgfqpoint{2.617916in}{2.011662in}}%
\pgfpathlineto{\pgfqpoint{2.648993in}{2.030598in}}%
\pgfpathlineto{\pgfqpoint{2.675927in}{2.043755in}}%
\pgfpathlineto{\pgfqpoint{2.702446in}{2.058908in}}%
\pgfpathlineto{\pgfqpoint{2.711148in}{2.062425in}}%
\pgfpathlineto{\pgfqpoint{2.723164in}{2.067310in}}%
\pgfpathlineto{\pgfqpoint{2.741396in}{2.076849in}}%
\pgfpathlineto{\pgfqpoint{2.760457in}{2.084622in}}%
\pgfpathlineto{\pgfqpoint{2.761285in}{2.084491in}}%
\pgfpathlineto{\pgfqpoint{2.773302in}{2.090519in}}%
\pgfpathlineto{\pgfqpoint{2.782418in}{2.095002in}}%
\pgfpathlineto{\pgfqpoint{2.784490in}{2.095239in}}%
\pgfpathlineto{\pgfqpoint{2.798578in}{2.100482in}}%
\pgfpathlineto{\pgfqpoint{2.803550in}{2.102340in}}%
\pgfpathlineto{\pgfqpoint{2.810595in}{2.104931in}}%
\pgfpathlineto{\pgfqpoint{2.825097in}{2.109598in}}%
\pgfpathlineto{\pgfqpoint{2.836699in}{2.114510in}}%
\pgfpathlineto{\pgfqpoint{2.838771in}{2.115337in}}%
\pgfpathlineto{\pgfqpoint{2.839600in}{2.115017in}}%
\pgfpathlineto{\pgfqpoint{2.847058in}{2.117182in}}%
\pgfpathlineto{\pgfqpoint{2.859489in}{2.121001in}}%
\pgfpathlineto{\pgfqpoint{2.860318in}{2.121344in}}%
\pgfpathlineto{\pgfqpoint{2.861147in}{2.121050in}}%
\pgfpathlineto{\pgfqpoint{2.871091in}{2.126269in}}%
\pgfpathlineto{\pgfqpoint{2.875235in}{2.128569in}}%
\pgfpathlineto{\pgfqpoint{2.876064in}{2.128340in}}%
\pgfpathlineto{\pgfqpoint{2.889738in}{2.135740in}}%
\pgfpathlineto{\pgfqpoint{2.890566in}{2.135395in}}%
\pgfpathlineto{\pgfqpoint{2.903412in}{2.141539in}}%
\pgfpathlineto{\pgfqpoint{2.904240in}{2.141098in}}%
\pgfpathlineto{\pgfqpoint{2.910456in}{2.143378in}}%
\pgfpathlineto{\pgfqpoint{2.917914in}{2.144794in}}%
\pgfpathlineto{\pgfqpoint{2.926202in}{2.146672in}}%
\pgfpathlineto{\pgfqpoint{2.933660in}{2.149326in}}%
\pgfpathlineto{\pgfqpoint{2.936561in}{2.150502in}}%
\pgfpathlineto{\pgfqpoint{2.937389in}{2.149927in}}%
\pgfpathlineto{\pgfqpoint{2.948163in}{2.154684in}}%
\pgfpathlineto{\pgfqpoint{2.950649in}{2.155978in}}%
\pgfpathlineto{\pgfqpoint{2.951478in}{2.155443in}}%
\pgfpathlineto{\pgfqpoint{2.958107in}{2.159533in}}%
\pgfpathlineto{\pgfqpoint{2.961422in}{2.161877in}}%
\pgfpathlineto{\pgfqpoint{2.961837in}{2.161121in}}%
\pgfpathlineto{\pgfqpoint{2.962251in}{2.161412in}}%
\pgfpathlineto{\pgfqpoint{2.970124in}{2.167268in}}%
\pgfpathlineto{\pgfqpoint{2.970538in}{2.166461in}}%
\pgfpathlineto{\pgfqpoint{2.970953in}{2.166809in}}%
\pgfpathlineto{\pgfqpoint{2.979240in}{2.173517in}}%
\pgfpathlineto{\pgfqpoint{2.979654in}{2.172621in}}%
\pgfpathlineto{\pgfqpoint{2.980069in}{2.172942in}}%
\pgfpathlineto{\pgfqpoint{2.998301in}{2.187229in}}%
\pgfpathlineto{\pgfqpoint{2.998715in}{2.186227in}}%
\pgfpathlineto{\pgfqpoint{2.999129in}{2.186550in}}%
\pgfpathlineto{\pgfqpoint{3.008245in}{2.193741in}}%
\pgfpathlineto{\pgfqpoint{3.008660in}{2.192666in}}%
\pgfpathlineto{\pgfqpoint{3.009074in}{2.193009in}}%
\pgfpathlineto{\pgfqpoint{3.018190in}{2.200273in}}%
\pgfpathlineto{\pgfqpoint{3.018604in}{2.199059in}}%
\pgfpathlineto{\pgfqpoint{3.019018in}{2.199389in}}%
\pgfpathlineto{\pgfqpoint{3.026849in}{2.205807in}}%
\pgfpathlineto{\pgfqpoint{3.026849in}{2.205807in}}%
\pgfusepath{stroke}%
\end{pgfscope}%
\begin{pgfscope}%
\pgfpathrectangle{\pgfqpoint{0.650000in}{0.420000in}}{\pgfqpoint{2.388110in}{1.994488in}}%
\pgfusepath{clip}%
\pgfsetrectcap%
\pgfsetroundjoin%
\pgfsetlinewidth{1.003750pt}%
\definecolor{currentstroke}{rgb}{0.501961,0.450980,0.674510}%
\pgfsetstrokecolor{currentstroke}%
\pgfsetdash{}{0pt}%
\pgfpathmoveto{\pgfqpoint{0.650000in}{0.512189in}}%
\pgfpathlineto{\pgfqpoint{0.662864in}{0.517315in}}%
\pgfpathlineto{\pgfqpoint{0.667861in}{0.520396in}}%
\pgfpathlineto{\pgfqpoint{0.672440in}{0.525697in}}%
\pgfpathlineto{\pgfqpoint{0.677020in}{0.532600in}}%
\pgfpathlineto{\pgfqpoint{0.686180in}{0.548582in}}%
\pgfpathlineto{\pgfqpoint{0.692842in}{0.561303in}}%
\pgfpathlineto{\pgfqpoint{0.721571in}{0.610830in}}%
\pgfpathlineto{\pgfqpoint{0.750299in}{0.651239in}}%
\pgfpathlineto{\pgfqpoint{0.796515in}{0.706225in}}%
\pgfpathlineto{\pgfqpoint{0.828159in}{0.740183in}}%
\pgfpathlineto{\pgfqpoint{0.928085in}{0.842173in}}%
\pgfpathlineto{\pgfqpoint{0.963891in}{0.877690in}}%
\pgfpathlineto{\pgfqpoint{0.978048in}{0.891944in}}%
\pgfpathlineto{\pgfqpoint{0.991787in}{0.905663in}}%
\pgfpathlineto{\pgfqpoint{1.069646in}{0.984467in}}%
\pgfpathlineto{\pgfqpoint{1.082137in}{0.994842in}}%
\pgfpathlineto{\pgfqpoint{1.093379in}{1.002986in}}%
\pgfpathlineto{\pgfqpoint{1.103371in}{1.009353in}}%
\pgfpathlineto{\pgfqpoint{1.113780in}{1.014782in}}%
\pgfpathlineto{\pgfqpoint{1.122524in}{1.018042in}}%
\pgfpathlineto{\pgfqpoint{1.135431in}{1.021399in}}%
\pgfpathlineto{\pgfqpoint{1.149587in}{1.025320in}}%
\pgfpathlineto{\pgfqpoint{1.157914in}{1.029138in}}%
\pgfpathlineto{\pgfqpoint{1.174568in}{1.038935in}}%
\pgfpathlineto{\pgfqpoint{1.196219in}{1.053533in}}%
\pgfpathlineto{\pgfqpoint{1.240769in}{1.088331in}}%
\pgfpathlineto{\pgfqpoint{1.259922in}{1.104161in}}%
\pgfpathlineto{\pgfqpoint{1.261587in}{1.105399in}}%
\pgfpathlineto{\pgfqpoint{1.279075in}{1.119296in}}%
\pgfpathlineto{\pgfqpoint{1.304056in}{1.139958in}}%
\pgfpathlineto{\pgfqpoint{1.337781in}{1.167640in}}%
\pgfpathlineto{\pgfqpoint{1.370257in}{1.193598in}}%
\pgfpathlineto{\pgfqpoint{1.407730in}{1.224108in}}%
\pgfpathlineto{\pgfqpoint{1.409395in}{1.225270in}}%
\pgfpathlineto{\pgfqpoint{1.423967in}{1.236268in}}%
\pgfpathlineto{\pgfqpoint{1.491001in}{1.284308in}}%
\pgfpathlineto{\pgfqpoint{1.554704in}{1.332426in}}%
\pgfpathlineto{\pgfqpoint{1.599671in}{1.361416in}}%
\pgfpathlineto{\pgfqpoint{1.600920in}{1.361950in}}%
\pgfpathlineto{\pgfqpoint{1.628816in}{1.379694in}}%
\pgfpathlineto{\pgfqpoint{1.637976in}{1.386050in}}%
\pgfpathlineto{\pgfqpoint{1.670452in}{1.407824in}}%
\pgfpathlineto{\pgfqpoint{1.718333in}{1.441178in}}%
\pgfpathlineto{\pgfqpoint{1.727077in}{1.446836in}}%
\pgfpathlineto{\pgfqpoint{1.735404in}{1.452083in}}%
\pgfpathlineto{\pgfqpoint{1.749560in}{1.459729in}}%
\pgfpathlineto{\pgfqpoint{1.767464in}{1.470610in}}%
\pgfpathlineto{\pgfqpoint{1.779955in}{1.477971in}}%
\pgfpathlineto{\pgfqpoint{1.797858in}{1.488614in}}%
\pgfpathlineto{\pgfqpoint{1.811598in}{1.497400in}}%
\pgfpathlineto{\pgfqpoint{1.812847in}{1.497936in}}%
\pgfpathlineto{\pgfqpoint{1.826587in}{1.507096in}}%
\pgfpathlineto{\pgfqpoint{1.848654in}{1.522976in}}%
\pgfpathlineto{\pgfqpoint{1.870304in}{1.538338in}}%
\pgfpathlineto{\pgfqpoint{1.871554in}{1.538971in}}%
\pgfpathlineto{\pgfqpoint{1.894453in}{1.555727in}}%
\pgfpathlineto{\pgfqpoint{1.901948in}{1.560777in}}%
\pgfpathlineto{\pgfqpoint{1.965651in}{1.601179in}}%
\pgfpathlineto{\pgfqpoint{1.970231in}{1.602990in}}%
\pgfpathlineto{\pgfqpoint{1.975227in}{1.604069in}}%
\pgfpathlineto{\pgfqpoint{1.978974in}{1.607349in}}%
\pgfpathlineto{\pgfqpoint{1.982721in}{1.610292in}}%
\pgfpathlineto{\pgfqpoint{1.991048in}{1.615091in}}%
\pgfpathlineto{\pgfqpoint{1.992298in}{1.615529in}}%
\pgfpathlineto{\pgfqpoint{2.008952in}{1.625419in}}%
\pgfpathlineto{\pgfqpoint{2.032268in}{1.641706in}}%
\pgfpathlineto{\pgfqpoint{2.033517in}{1.642287in}}%
\pgfpathlineto{\pgfqpoint{2.061413in}{1.660852in}}%
\pgfpathlineto{\pgfqpoint{2.079317in}{1.671800in}}%
\pgfpathlineto{\pgfqpoint{2.094722in}{1.680925in}}%
\pgfpathlineto{\pgfqpoint{2.126782in}{1.701489in}}%
\pgfpathlineto{\pgfqpoint{2.134692in}{1.706454in}}%
\pgfpathlineto{\pgfqpoint{2.151763in}{1.715312in}}%
\pgfpathlineto{\pgfqpoint{2.154261in}{1.716399in}}%
\pgfpathlineto{\pgfqpoint{2.156759in}{1.719033in}}%
\pgfpathlineto{\pgfqpoint{2.160923in}{1.723494in}}%
\pgfpathlineto{\pgfqpoint{2.162172in}{1.724169in}}%
\pgfpathlineto{\pgfqpoint{2.191317in}{1.744185in}}%
\pgfpathlineto{\pgfqpoint{2.202558in}{1.750289in}}%
\pgfpathlineto{\pgfqpoint{2.229622in}{1.764394in}}%
\pgfpathlineto{\pgfqpoint{2.248358in}{1.774697in}}%
\pgfpathlineto{\pgfqpoint{2.265429in}{1.786389in}}%
\pgfpathlineto{\pgfqpoint{2.277919in}{1.795130in}}%
\pgfpathlineto{\pgfqpoint{2.289578in}{1.802942in}}%
\pgfpathlineto{\pgfqpoint{2.303318in}{1.811085in}}%
\pgfpathlineto{\pgfqpoint{2.304150in}{1.811251in}}%
\pgfpathlineto{\pgfqpoint{2.324968in}{1.821778in}}%
\pgfpathlineto{\pgfqpoint{2.341206in}{1.829531in}}%
\pgfpathlineto{\pgfqpoint{2.387839in}{1.855704in}}%
\pgfpathlineto{\pgfqpoint{2.398248in}{1.862143in}}%
\pgfpathlineto{\pgfqpoint{2.412820in}{1.871478in}}%
\pgfpathlineto{\pgfqpoint{2.422813in}{1.880330in}}%
\pgfpathlineto{\pgfqpoint{2.438218in}{1.894410in}}%
\pgfpathlineto{\pgfqpoint{2.447794in}{1.901583in}}%
\pgfpathlineto{\pgfqpoint{2.469029in}{1.918480in}}%
\pgfpathlineto{\pgfqpoint{2.528152in}{1.970948in}}%
\pgfpathlineto{\pgfqpoint{2.541059in}{1.984827in}}%
\pgfpathlineto{\pgfqpoint{2.555215in}{2.002656in}}%
\pgfpathlineto{\pgfqpoint{2.582277in}{2.038173in}}%
\pgfpathlineto{\pgfqpoint{2.596432in}{2.054711in}}%
\pgfpathlineto{\pgfqpoint{2.611003in}{2.070342in}}%
\pgfpathlineto{\pgfqpoint{2.621411in}{2.079504in}}%
\pgfpathlineto{\pgfqpoint{2.635983in}{2.092679in}}%
\pgfpathlineto{\pgfqpoint{2.646391in}{2.102288in}}%
\pgfpathlineto{\pgfqpoint{2.659297in}{2.113089in}}%
\pgfpathlineto{\pgfqpoint{2.694268in}{2.140150in}}%
\pgfpathlineto{\pgfqpoint{2.710088in}{2.152542in}}%
\pgfpathlineto{\pgfqpoint{2.720913in}{2.160920in}}%
\pgfpathlineto{\pgfqpoint{2.725492in}{2.164333in}}%
\pgfpathlineto{\pgfqpoint{2.733403in}{2.170072in}}%
\pgfpathlineto{\pgfqpoint{2.740064in}{2.173138in}}%
\pgfpathlineto{\pgfqpoint{2.761296in}{2.182268in}}%
\pgfpathlineto{\pgfqpoint{2.774619in}{2.189103in}}%
\pgfpathlineto{\pgfqpoint{2.793770in}{2.198286in}}%
\pgfpathlineto{\pgfqpoint{2.808757in}{2.206001in}}%
\pgfpathlineto{\pgfqpoint{2.815835in}{2.207663in}}%
\pgfpathlineto{\pgfqpoint{2.823329in}{2.209617in}}%
\pgfpathlineto{\pgfqpoint{2.829157in}{2.212481in}}%
\pgfpathlineto{\pgfqpoint{2.854137in}{2.226447in}}%
\pgfpathlineto{\pgfqpoint{2.862880in}{2.231539in}}%
\pgfpathlineto{\pgfqpoint{2.869957in}{2.235571in}}%
\pgfpathlineto{\pgfqpoint{2.902014in}{2.249074in}}%
\pgfpathlineto{\pgfqpoint{2.910341in}{2.252362in}}%
\pgfpathlineto{\pgfqpoint{2.931573in}{2.263422in}}%
\pgfpathlineto{\pgfqpoint{2.964463in}{2.280495in}}%
\pgfpathlineto{\pgfqpoint{2.976536in}{2.288025in}}%
\pgfpathlineto{\pgfqpoint{2.993606in}{2.298015in}}%
\pgfpathlineto{\pgfqpoint{3.006095in}{2.302683in}}%
\pgfpathlineto{\pgfqpoint{3.011924in}{2.306587in}}%
\pgfpathlineto{\pgfqpoint{3.025246in}{2.316398in}}%
\pgfpathlineto{\pgfqpoint{3.032733in}{2.321424in}}%
\pgfpathlineto{\pgfqpoint{3.038110in}{2.323830in}}%
\pgfpathlineto{\pgfqpoint{3.038110in}{2.323830in}}%
\pgfusepath{stroke}%
\end{pgfscope}%
\begin{pgfscope}%
\pgfsetrectcap%
\pgfsetmiterjoin%
\pgfsetlinewidth{0.803000pt}%
\definecolor{currentstroke}{rgb}{0.000000,0.000000,0.000000}%
\pgfsetstrokecolor{currentstroke}%
\pgfsetdash{}{0pt}%
\pgfpathmoveto{\pgfqpoint{0.650000in}{0.420000in}}%
\pgfpathlineto{\pgfqpoint{0.650000in}{2.414488in}}%
\pgfusepath{stroke}%
\end{pgfscope}%
\begin{pgfscope}%
\pgfsetrectcap%
\pgfsetmiterjoin%
\pgfsetlinewidth{0.803000pt}%
\definecolor{currentstroke}{rgb}{0.000000,0.000000,0.000000}%
\pgfsetstrokecolor{currentstroke}%
\pgfsetdash{}{0pt}%
\pgfpathmoveto{\pgfqpoint{3.038110in}{0.420000in}}%
\pgfpathlineto{\pgfqpoint{3.038110in}{2.414488in}}%
\pgfusepath{stroke}%
\end{pgfscope}%
\begin{pgfscope}%
\pgfsetrectcap%
\pgfsetmiterjoin%
\pgfsetlinewidth{0.803000pt}%
\definecolor{currentstroke}{rgb}{0.000000,0.000000,0.000000}%
\pgfsetstrokecolor{currentstroke}%
\pgfsetdash{}{0pt}%
\pgfpathmoveto{\pgfqpoint{0.650000in}{0.420000in}}%
\pgfpathlineto{\pgfqpoint{3.038110in}{0.420000in}}%
\pgfusepath{stroke}%
\end{pgfscope}%
\begin{pgfscope}%
\pgfsetrectcap%
\pgfsetmiterjoin%
\pgfsetlinewidth{0.803000pt}%
\definecolor{currentstroke}{rgb}{0.000000,0.000000,0.000000}%
\pgfsetstrokecolor{currentstroke}%
\pgfsetdash{}{0pt}%
\pgfpathmoveto{\pgfqpoint{0.650000in}{2.414488in}}%
\pgfpathlineto{\pgfqpoint{3.038110in}{2.414488in}}%
\pgfusepath{stroke}%
\end{pgfscope}%
\end{pgfpicture}%
\makeatother%
\endgroup%

    \label{fig:R500deg18R100deg90_delta99}
\end{subfigure}
\\[-15pt]
\caption[First and second order statistics of R100deg90 and R500deg18 optimised simulations at $\mathrm{Re}_{\theta} = 2000$ with reference data taken from \citet{schlatter_assessment_2010} at $\mathrm{Re}_{\theta} = 2000$, and streamwise boundary layer thickness development normalised with reference boundary layer thickness $\delta_0$]{First and second order statistics R100deg90 and R500deg18 optimised simulations at $\mathrm{Re}_{\theta} = 2000$ with reference data taken from \citet{schlatter_assessment_2010} at $\mathrm{Re}_{\theta} = 2000$: (a) streamwise velocity profile in viscous scales; (b) Reynolds stresses $\langle u_i' u_j' \rangle$; (c) budget of \acrshort{tke} = $\frac{1}{2} \langle u_i' u_i' \rangle$, line styles denote budget terms according to \autoref{tab:budgetLinestyles}; (d) streamwise boundary layer thickness normalised with reference boundary layer thickness $\delta_0$ of \lswatch[solid]{delta99thetacolone} base /\lswatch[solid]{delta99thetacoltwo} optimised R100deg90, and \lswatch[solid]{delta99thetacolthree} base / \lswatch[solid]{delta99thetacolfour} optimised R500deg18.}
\label{fig:R500deg18R100deg90_Re2000stats}
\end{figure}


In \autoref{fig:R500deg18R100deg90_Re2000stats}a-c The statistics extracted at $\mathrm{Re}_{\theta}=2000$ are depicted in \autoref{fig:R500deg18R100deg90_Re2000stats}a-c.
As can be seen, the velocity profiles of the two cases are similar. 
Up to the log-law region, the profiles match the reference profile and the linear ad the log law perfectly.
However, in the wake region some differences become apparent.
The R100deg90 profile starts to surpass the R500deg18 profile at approximately $y^+ \approx 100$.
The gap between the profiles grows constantly, reaching its maximum at the transition to the outer layer.
In the outer layer, a significant decrease in the velocity of the R100deg90 profile is observed.
The  R500deg18 profile fits the reference data up to the beginning of the outer layer.
In the outer layer the velocity remains constant at a slightly higher value than that of the reference profile.



The velocity decrease in the outer layer of case R100deg90 can be induced by the \acrshort{bc} at the upper wall.
This phenomenon occurs when the imposed velocity at the upper wall is smaller than the free-stream velocity in the domain, which evokes a relative velocity reduction in the wall-normal direction.
In the velocity profile, however, this effect  is only visible if the influence of the \acrshort{bc} extends significantly into the domain up to the region that is already plotted as the outer region in the velocity profile.
A possible explanation for the observed behaviour is the limited height of the domain, resulting in only a small distance between the edge of the \acrshort{tbl} and the upper boundary. 
Consequently, the velocity gradient between the free stream and the upper wall becomes sufficiently large to influence a substantial portion of the free stream.
As stated in \autoref{tab:curved_cases} the domain of R100deg90 is less high than the one for R500deg18.


The Reynolds stresses are visualised in  \autoref{fig:R500deg18R100deg90_fluct}.
Overall, both simulations are in good agreement with the reference data for all stress components. 
The only remarkable divergence of the reference data occurs in the peak of the $\langle u'u' \rangle^+$-component
where the reference data shows a higher peak value than either  simulation performed.
 The reduction is particularly pronounced for case  R100deg90. 
This effect is indicated by the geometry of this case.
As already shown by \citet{so_experiment_1973} and \cite{schwarz_convex_1996}, a convex curved geometry reduces the peak of the $\langle u'u' \rangle^+$-component of the Reynolds stresses. 
Since convex curvature has a stabilising effect on the \acrshort{tbl}, the turbulent fluctuations are suppressed  and consequently, the Reynolds stresses \citet{yang_assessment_2016}.
This effect strengthens with increasing intensity of the curvature; thus,  the reduction is more pronounced for  case  R100deg90.
In \autoref{fig:R500deg18R100deg90_bud} the \acrshort{tke} budgets are shown.
As observed for case R100deg30,  the overall distributions of the individual budget terms agree well with the reference data. 
Nevertheless, the influence of convex curvature discussed above is also reflected in the \acrshort{tke} budget.


A reduction in the peak of the production term is observed, particularly for case R100deg90. This behaviour was also reported  by \citet{gillis_turbulent_1983} and  is attributed to the strong convex curvature in this configuration.
The streamwise development of the boundary layer thickness over the domain is recorded in \autoref{fig:R500deg18R100deg90_delta99}.
For both cases, the optimised and the base configurations are plotted to assess the effect of the optimisation on the boundary layer thickness.
The overall development is similar for all configurations and follows the same trend as observed for the development in the case R100deg30, visualised in \autoref{fig:R100deg30Re1200Thickness}.
In both  R100deg90 and R500deg18, the onset and the end of the curvature noticeable affects the growth of the \acrshort{tbl}.
A detailed discussion of these general characteristics has been provided in  \autoref{subsec:R100deg30}.

Compared with R100deg90,  case R500deg18 exhibits  a larger \acrshort{tbl} growth rate. This difference is consistent with the stronger stabilising effect of the more intense convex curvature in R100deg90, which suppresses boundary-layer growth \citep{gillis_turbulent_1983,so_experiment_1973}.

The optimisation affects the \acrshort{tbl} growth differently in case R100deg90 and R500deg18.
For case R500deg18, $\delta_{99}$ is less affected by the curvature in the optimised configuration.
A similar behaviour was already detected for case R100deg30.
A converse effect has been reported for a concave curved \acrshortpl{tbl} by \citet{appelbaum_systematic_2025}.
In contrast, the optimisation further suppresses the growth of $\delta_{99}$ in case R100deg90 when compared to the base configuration.
This indicates that the qualitative effect of the optimisation depends on the curvature intensity of the configuration.

\begin{figure}[H]
\centering
\begin{subfigure}[t]{0.495\textwidth}
    \centering
    \subcaption{}
    %% Creator: Matplotlib, PGF backend
%%
%% To include the figure in your LaTeX document, write
%%   \input{<filename>.pgf}
%%
%% Make sure the required packages are loaded in your preamble
%%   \usepackage{pgf}
%%
%% Also ensure that all the required font packages are loaded; for instance,
%% the lmodern package is sometimes necessary when using math font.
%%   \usepackage{lmodern}
%%
%% Figures using additional raster images can only be included by \input if
%% they are in the same directory as the main LaTeX file. For loading figures
%% from other directories you can use the `import` package
%%   \usepackage{import}
%%
%% and then include the figures with
%%   \import{<path to file>}{<filename>.pgf}
%%
%% Matplotlib used the following preamble
%%   \def\mathdefault#1{#1}
%%   \everymath=\expandafter{\the\everymath\displaystyle}
%%   \IfFileExists{scrextend.sty}{
%%     \usepackage[fontsize=11.000000pt]{scrextend}
%%   }{
%%     \renewcommand{\normalsize}{\fontsize{11.000000}{13.200000}\selectfont}
%%     \normalsize
%%   }
%%   
%%   \makeatletter\@ifpackageloaded{underscore}{}{\usepackage[strings]{underscore}}\makeatother
%%
\begingroup%
\makeatletter%
\begin{pgfpicture}%
\pgfpathrectangle{\pgfpointorigin}{\pgfqpoint{3.118110in}{2.494488in}}%
\pgfusepath{use as bounding box, clip}%
\begin{pgfscope}%
\pgfsetbuttcap%
\pgfsetmiterjoin%
\definecolor{currentfill}{rgb}{1.000000,1.000000,1.000000}%
\pgfsetfillcolor{currentfill}%
\pgfsetlinewidth{0.000000pt}%
\definecolor{currentstroke}{rgb}{1.000000,1.000000,1.000000}%
\pgfsetstrokecolor{currentstroke}%
\pgfsetdash{}{0pt}%
\pgfpathmoveto{\pgfqpoint{0.000000in}{0.000000in}}%
\pgfpathlineto{\pgfqpoint{3.118110in}{0.000000in}}%
\pgfpathlineto{\pgfqpoint{3.118110in}{2.494488in}}%
\pgfpathlineto{\pgfqpoint{0.000000in}{2.494488in}}%
\pgfpathlineto{\pgfqpoint{0.000000in}{0.000000in}}%
\pgfpathclose%
\pgfusepath{fill}%
\end{pgfscope}%
\begin{pgfscope}%
\pgfsetbuttcap%
\pgfsetmiterjoin%
\definecolor{currentfill}{rgb}{1.000000,1.000000,1.000000}%
\pgfsetfillcolor{currentfill}%
\pgfsetlinewidth{0.000000pt}%
\definecolor{currentstroke}{rgb}{0.000000,0.000000,0.000000}%
\pgfsetstrokecolor{currentstroke}%
\pgfsetstrokeopacity{0.000000}%
\pgfsetdash{}{0pt}%
\pgfpathmoveto{\pgfqpoint{0.650000in}{0.420000in}}%
\pgfpathlineto{\pgfqpoint{3.038110in}{0.420000in}}%
\pgfpathlineto{\pgfqpoint{3.038110in}{2.414488in}}%
\pgfpathlineto{\pgfqpoint{0.650000in}{2.414488in}}%
\pgfpathlineto{\pgfqpoint{0.650000in}{0.420000in}}%
\pgfpathclose%
\pgfusepath{fill}%
\end{pgfscope}%
\begin{pgfscope}%
\pgfsetbuttcap%
\pgfsetroundjoin%
\definecolor{currentfill}{rgb}{0.000000,0.000000,0.000000}%
\pgfsetfillcolor{currentfill}%
\pgfsetlinewidth{0.501875pt}%
\definecolor{currentstroke}{rgb}{0.000000,0.000000,0.000000}%
\pgfsetstrokecolor{currentstroke}%
\pgfsetdash{}{0pt}%
\pgfsys@defobject{currentmarker}{\pgfqpoint{0.000000in}{0.000000in}}{\pgfqpoint{0.000000in}{0.041667in}}{%
\pgfpathmoveto{\pgfqpoint{0.000000in}{0.000000in}}%
\pgfpathlineto{\pgfqpoint{0.000000in}{0.041667in}}%
\pgfusepath{stroke,fill}%
}%
\begin{pgfscope}%
\pgfsys@transformshift{0.650000in}{0.420000in}%
\pgfsys@useobject{currentmarker}{}%
\end{pgfscope}%
\end{pgfscope}%
\begin{pgfscope}%
\pgfsetbuttcap%
\pgfsetroundjoin%
\definecolor{currentfill}{rgb}{0.000000,0.000000,0.000000}%
\pgfsetfillcolor{currentfill}%
\pgfsetlinewidth{0.501875pt}%
\definecolor{currentstroke}{rgb}{0.000000,0.000000,0.000000}%
\pgfsetstrokecolor{currentstroke}%
\pgfsetdash{}{0pt}%
\pgfsys@defobject{currentmarker}{\pgfqpoint{0.000000in}{-0.041667in}}{\pgfqpoint{0.000000in}{0.000000in}}{%
\pgfpathmoveto{\pgfqpoint{0.000000in}{0.000000in}}%
\pgfpathlineto{\pgfqpoint{0.000000in}{-0.041667in}}%
\pgfusepath{stroke,fill}%
}%
\begin{pgfscope}%
\pgfsys@transformshift{0.650000in}{2.414488in}%
\pgfsys@useobject{currentmarker}{}%
\end{pgfscope}%
\end{pgfscope}%
\begin{pgfscope}%
\definecolor{textcolor}{rgb}{0.000000,0.000000,0.000000}%
\pgfsetstrokecolor{textcolor}%
\pgfsetfillcolor{textcolor}%
\pgftext[x=0.650000in,y=0.371389in,,top]{\color{textcolor}{\rmfamily\fontsize{11.000000}{13.200000}\selectfont\catcode`\^=\active\def^{\ifmmode\sp\else\^{}\fi}\catcode`\%=\active\def%{\%}0}}%
\end{pgfscope}%
\begin{pgfscope}%
\pgfsetbuttcap%
\pgfsetroundjoin%
\definecolor{currentfill}{rgb}{0.000000,0.000000,0.000000}%
\pgfsetfillcolor{currentfill}%
\pgfsetlinewidth{0.501875pt}%
\definecolor{currentstroke}{rgb}{0.000000,0.000000,0.000000}%
\pgfsetstrokecolor{currentstroke}%
\pgfsetdash{}{0pt}%
\pgfsys@defobject{currentmarker}{\pgfqpoint{0.000000in}{0.000000in}}{\pgfqpoint{0.000000in}{0.041667in}}{%
\pgfpathmoveto{\pgfqpoint{0.000000in}{0.000000in}}%
\pgfpathlineto{\pgfqpoint{0.000000in}{0.041667in}}%
\pgfusepath{stroke,fill}%
}%
\begin{pgfscope}%
\pgfsys@transformshift{1.493052in}{0.420000in}%
\pgfsys@useobject{currentmarker}{}%
\end{pgfscope}%
\end{pgfscope}%
\begin{pgfscope}%
\pgfsetbuttcap%
\pgfsetroundjoin%
\definecolor{currentfill}{rgb}{0.000000,0.000000,0.000000}%
\pgfsetfillcolor{currentfill}%
\pgfsetlinewidth{0.501875pt}%
\definecolor{currentstroke}{rgb}{0.000000,0.000000,0.000000}%
\pgfsetstrokecolor{currentstroke}%
\pgfsetdash{}{0pt}%
\pgfsys@defobject{currentmarker}{\pgfqpoint{0.000000in}{-0.041667in}}{\pgfqpoint{0.000000in}{0.000000in}}{%
\pgfpathmoveto{\pgfqpoint{0.000000in}{0.000000in}}%
\pgfpathlineto{\pgfqpoint{0.000000in}{-0.041667in}}%
\pgfusepath{stroke,fill}%
}%
\begin{pgfscope}%
\pgfsys@transformshift{1.493052in}{2.414488in}%
\pgfsys@useobject{currentmarker}{}%
\end{pgfscope}%
\end{pgfscope}%
\begin{pgfscope}%
\definecolor{textcolor}{rgb}{0.000000,0.000000,0.000000}%
\pgfsetstrokecolor{textcolor}%
\pgfsetfillcolor{textcolor}%
\pgftext[x=1.493052in,y=0.371389in,,top]{\color{textcolor}{\rmfamily\fontsize{11.000000}{13.200000}\selectfont\catcode`\^=\active\def^{\ifmmode\sp\else\^{}\fi}\catcode`\%=\active\def%{\%}100}}%
\end{pgfscope}%
\begin{pgfscope}%
\pgfsetbuttcap%
\pgfsetroundjoin%
\definecolor{currentfill}{rgb}{0.000000,0.000000,0.000000}%
\pgfsetfillcolor{currentfill}%
\pgfsetlinewidth{0.501875pt}%
\definecolor{currentstroke}{rgb}{0.000000,0.000000,0.000000}%
\pgfsetstrokecolor{currentstroke}%
\pgfsetdash{}{0pt}%
\pgfsys@defobject{currentmarker}{\pgfqpoint{0.000000in}{0.000000in}}{\pgfqpoint{0.000000in}{0.041667in}}{%
\pgfpathmoveto{\pgfqpoint{0.000000in}{0.000000in}}%
\pgfpathlineto{\pgfqpoint{0.000000in}{0.041667in}}%
\pgfusepath{stroke,fill}%
}%
\begin{pgfscope}%
\pgfsys@transformshift{2.336103in}{0.420000in}%
\pgfsys@useobject{currentmarker}{}%
\end{pgfscope}%
\end{pgfscope}%
\begin{pgfscope}%
\pgfsetbuttcap%
\pgfsetroundjoin%
\definecolor{currentfill}{rgb}{0.000000,0.000000,0.000000}%
\pgfsetfillcolor{currentfill}%
\pgfsetlinewidth{0.501875pt}%
\definecolor{currentstroke}{rgb}{0.000000,0.000000,0.000000}%
\pgfsetstrokecolor{currentstroke}%
\pgfsetdash{}{0pt}%
\pgfsys@defobject{currentmarker}{\pgfqpoint{0.000000in}{-0.041667in}}{\pgfqpoint{0.000000in}{0.000000in}}{%
\pgfpathmoveto{\pgfqpoint{0.000000in}{0.000000in}}%
\pgfpathlineto{\pgfqpoint{0.000000in}{-0.041667in}}%
\pgfusepath{stroke,fill}%
}%
\begin{pgfscope}%
\pgfsys@transformshift{2.336103in}{2.414488in}%
\pgfsys@useobject{currentmarker}{}%
\end{pgfscope}%
\end{pgfscope}%
\begin{pgfscope}%
\definecolor{textcolor}{rgb}{0.000000,0.000000,0.000000}%
\pgfsetstrokecolor{textcolor}%
\pgfsetfillcolor{textcolor}%
\pgftext[x=2.336103in,y=0.371389in,,top]{\color{textcolor}{\rmfamily\fontsize{11.000000}{13.200000}\selectfont\catcode`\^=\active\def^{\ifmmode\sp\else\^{}\fi}\catcode`\%=\active\def%{\%}200}}%
\end{pgfscope}%
\begin{pgfscope}%
\pgfsetbuttcap%
\pgfsetroundjoin%
\definecolor{currentfill}{rgb}{0.000000,0.000000,0.000000}%
\pgfsetfillcolor{currentfill}%
\pgfsetlinewidth{0.401500pt}%
\definecolor{currentstroke}{rgb}{0.000000,0.000000,0.000000}%
\pgfsetstrokecolor{currentstroke}%
\pgfsetdash{}{0pt}%
\pgfsys@defobject{currentmarker}{\pgfqpoint{0.000000in}{0.000000in}}{\pgfqpoint{0.000000in}{0.020833in}}{%
\pgfpathmoveto{\pgfqpoint{0.000000in}{0.000000in}}%
\pgfpathlineto{\pgfqpoint{0.000000in}{0.020833in}}%
\pgfusepath{stroke,fill}%
}%
\begin{pgfscope}%
\pgfsys@transformshift{0.818610in}{0.420000in}%
\pgfsys@useobject{currentmarker}{}%
\end{pgfscope}%
\end{pgfscope}%
\begin{pgfscope}%
\pgfsetbuttcap%
\pgfsetroundjoin%
\definecolor{currentfill}{rgb}{0.000000,0.000000,0.000000}%
\pgfsetfillcolor{currentfill}%
\pgfsetlinewidth{0.401500pt}%
\definecolor{currentstroke}{rgb}{0.000000,0.000000,0.000000}%
\pgfsetstrokecolor{currentstroke}%
\pgfsetdash{}{0pt}%
\pgfsys@defobject{currentmarker}{\pgfqpoint{0.000000in}{-0.020833in}}{\pgfqpoint{0.000000in}{0.000000in}}{%
\pgfpathmoveto{\pgfqpoint{0.000000in}{0.000000in}}%
\pgfpathlineto{\pgfqpoint{0.000000in}{-0.020833in}}%
\pgfusepath{stroke,fill}%
}%
\begin{pgfscope}%
\pgfsys@transformshift{0.818610in}{2.414488in}%
\pgfsys@useobject{currentmarker}{}%
\end{pgfscope}%
\end{pgfscope}%
\begin{pgfscope}%
\pgfsetbuttcap%
\pgfsetroundjoin%
\definecolor{currentfill}{rgb}{0.000000,0.000000,0.000000}%
\pgfsetfillcolor{currentfill}%
\pgfsetlinewidth{0.401500pt}%
\definecolor{currentstroke}{rgb}{0.000000,0.000000,0.000000}%
\pgfsetstrokecolor{currentstroke}%
\pgfsetdash{}{0pt}%
\pgfsys@defobject{currentmarker}{\pgfqpoint{0.000000in}{0.000000in}}{\pgfqpoint{0.000000in}{0.020833in}}{%
\pgfpathmoveto{\pgfqpoint{0.000000in}{0.000000in}}%
\pgfpathlineto{\pgfqpoint{0.000000in}{0.020833in}}%
\pgfusepath{stroke,fill}%
}%
\begin{pgfscope}%
\pgfsys@transformshift{0.987221in}{0.420000in}%
\pgfsys@useobject{currentmarker}{}%
\end{pgfscope}%
\end{pgfscope}%
\begin{pgfscope}%
\pgfsetbuttcap%
\pgfsetroundjoin%
\definecolor{currentfill}{rgb}{0.000000,0.000000,0.000000}%
\pgfsetfillcolor{currentfill}%
\pgfsetlinewidth{0.401500pt}%
\definecolor{currentstroke}{rgb}{0.000000,0.000000,0.000000}%
\pgfsetstrokecolor{currentstroke}%
\pgfsetdash{}{0pt}%
\pgfsys@defobject{currentmarker}{\pgfqpoint{0.000000in}{-0.020833in}}{\pgfqpoint{0.000000in}{0.000000in}}{%
\pgfpathmoveto{\pgfqpoint{0.000000in}{0.000000in}}%
\pgfpathlineto{\pgfqpoint{0.000000in}{-0.020833in}}%
\pgfusepath{stroke,fill}%
}%
\begin{pgfscope}%
\pgfsys@transformshift{0.987221in}{2.414488in}%
\pgfsys@useobject{currentmarker}{}%
\end{pgfscope}%
\end{pgfscope}%
\begin{pgfscope}%
\pgfsetbuttcap%
\pgfsetroundjoin%
\definecolor{currentfill}{rgb}{0.000000,0.000000,0.000000}%
\pgfsetfillcolor{currentfill}%
\pgfsetlinewidth{0.401500pt}%
\definecolor{currentstroke}{rgb}{0.000000,0.000000,0.000000}%
\pgfsetstrokecolor{currentstroke}%
\pgfsetdash{}{0pt}%
\pgfsys@defobject{currentmarker}{\pgfqpoint{0.000000in}{0.000000in}}{\pgfqpoint{0.000000in}{0.020833in}}{%
\pgfpathmoveto{\pgfqpoint{0.000000in}{0.000000in}}%
\pgfpathlineto{\pgfqpoint{0.000000in}{0.020833in}}%
\pgfusepath{stroke,fill}%
}%
\begin{pgfscope}%
\pgfsys@transformshift{1.155831in}{0.420000in}%
\pgfsys@useobject{currentmarker}{}%
\end{pgfscope}%
\end{pgfscope}%
\begin{pgfscope}%
\pgfsetbuttcap%
\pgfsetroundjoin%
\definecolor{currentfill}{rgb}{0.000000,0.000000,0.000000}%
\pgfsetfillcolor{currentfill}%
\pgfsetlinewidth{0.401500pt}%
\definecolor{currentstroke}{rgb}{0.000000,0.000000,0.000000}%
\pgfsetstrokecolor{currentstroke}%
\pgfsetdash{}{0pt}%
\pgfsys@defobject{currentmarker}{\pgfqpoint{0.000000in}{-0.020833in}}{\pgfqpoint{0.000000in}{0.000000in}}{%
\pgfpathmoveto{\pgfqpoint{0.000000in}{0.000000in}}%
\pgfpathlineto{\pgfqpoint{0.000000in}{-0.020833in}}%
\pgfusepath{stroke,fill}%
}%
\begin{pgfscope}%
\pgfsys@transformshift{1.155831in}{2.414488in}%
\pgfsys@useobject{currentmarker}{}%
\end{pgfscope}%
\end{pgfscope}%
\begin{pgfscope}%
\pgfsetbuttcap%
\pgfsetroundjoin%
\definecolor{currentfill}{rgb}{0.000000,0.000000,0.000000}%
\pgfsetfillcolor{currentfill}%
\pgfsetlinewidth{0.401500pt}%
\definecolor{currentstroke}{rgb}{0.000000,0.000000,0.000000}%
\pgfsetstrokecolor{currentstroke}%
\pgfsetdash{}{0pt}%
\pgfsys@defobject{currentmarker}{\pgfqpoint{0.000000in}{0.000000in}}{\pgfqpoint{0.000000in}{0.020833in}}{%
\pgfpathmoveto{\pgfqpoint{0.000000in}{0.000000in}}%
\pgfpathlineto{\pgfqpoint{0.000000in}{0.020833in}}%
\pgfusepath{stroke,fill}%
}%
\begin{pgfscope}%
\pgfsys@transformshift{1.324441in}{0.420000in}%
\pgfsys@useobject{currentmarker}{}%
\end{pgfscope}%
\end{pgfscope}%
\begin{pgfscope}%
\pgfsetbuttcap%
\pgfsetroundjoin%
\definecolor{currentfill}{rgb}{0.000000,0.000000,0.000000}%
\pgfsetfillcolor{currentfill}%
\pgfsetlinewidth{0.401500pt}%
\definecolor{currentstroke}{rgb}{0.000000,0.000000,0.000000}%
\pgfsetstrokecolor{currentstroke}%
\pgfsetdash{}{0pt}%
\pgfsys@defobject{currentmarker}{\pgfqpoint{0.000000in}{-0.020833in}}{\pgfqpoint{0.000000in}{0.000000in}}{%
\pgfpathmoveto{\pgfqpoint{0.000000in}{0.000000in}}%
\pgfpathlineto{\pgfqpoint{0.000000in}{-0.020833in}}%
\pgfusepath{stroke,fill}%
}%
\begin{pgfscope}%
\pgfsys@transformshift{1.324441in}{2.414488in}%
\pgfsys@useobject{currentmarker}{}%
\end{pgfscope}%
\end{pgfscope}%
\begin{pgfscope}%
\pgfsetbuttcap%
\pgfsetroundjoin%
\definecolor{currentfill}{rgb}{0.000000,0.000000,0.000000}%
\pgfsetfillcolor{currentfill}%
\pgfsetlinewidth{0.401500pt}%
\definecolor{currentstroke}{rgb}{0.000000,0.000000,0.000000}%
\pgfsetstrokecolor{currentstroke}%
\pgfsetdash{}{0pt}%
\pgfsys@defobject{currentmarker}{\pgfqpoint{0.000000in}{0.000000in}}{\pgfqpoint{0.000000in}{0.020833in}}{%
\pgfpathmoveto{\pgfqpoint{0.000000in}{0.000000in}}%
\pgfpathlineto{\pgfqpoint{0.000000in}{0.020833in}}%
\pgfusepath{stroke,fill}%
}%
\begin{pgfscope}%
\pgfsys@transformshift{1.661662in}{0.420000in}%
\pgfsys@useobject{currentmarker}{}%
\end{pgfscope}%
\end{pgfscope}%
\begin{pgfscope}%
\pgfsetbuttcap%
\pgfsetroundjoin%
\definecolor{currentfill}{rgb}{0.000000,0.000000,0.000000}%
\pgfsetfillcolor{currentfill}%
\pgfsetlinewidth{0.401500pt}%
\definecolor{currentstroke}{rgb}{0.000000,0.000000,0.000000}%
\pgfsetstrokecolor{currentstroke}%
\pgfsetdash{}{0pt}%
\pgfsys@defobject{currentmarker}{\pgfqpoint{0.000000in}{-0.020833in}}{\pgfqpoint{0.000000in}{0.000000in}}{%
\pgfpathmoveto{\pgfqpoint{0.000000in}{0.000000in}}%
\pgfpathlineto{\pgfqpoint{0.000000in}{-0.020833in}}%
\pgfusepath{stroke,fill}%
}%
\begin{pgfscope}%
\pgfsys@transformshift{1.661662in}{2.414488in}%
\pgfsys@useobject{currentmarker}{}%
\end{pgfscope}%
\end{pgfscope}%
\begin{pgfscope}%
\pgfsetbuttcap%
\pgfsetroundjoin%
\definecolor{currentfill}{rgb}{0.000000,0.000000,0.000000}%
\pgfsetfillcolor{currentfill}%
\pgfsetlinewidth{0.401500pt}%
\definecolor{currentstroke}{rgb}{0.000000,0.000000,0.000000}%
\pgfsetstrokecolor{currentstroke}%
\pgfsetdash{}{0pt}%
\pgfsys@defobject{currentmarker}{\pgfqpoint{0.000000in}{0.000000in}}{\pgfqpoint{0.000000in}{0.020833in}}{%
\pgfpathmoveto{\pgfqpoint{0.000000in}{0.000000in}}%
\pgfpathlineto{\pgfqpoint{0.000000in}{0.020833in}}%
\pgfusepath{stroke,fill}%
}%
\begin{pgfscope}%
\pgfsys@transformshift{1.830272in}{0.420000in}%
\pgfsys@useobject{currentmarker}{}%
\end{pgfscope}%
\end{pgfscope}%
\begin{pgfscope}%
\pgfsetbuttcap%
\pgfsetroundjoin%
\definecolor{currentfill}{rgb}{0.000000,0.000000,0.000000}%
\pgfsetfillcolor{currentfill}%
\pgfsetlinewidth{0.401500pt}%
\definecolor{currentstroke}{rgb}{0.000000,0.000000,0.000000}%
\pgfsetstrokecolor{currentstroke}%
\pgfsetdash{}{0pt}%
\pgfsys@defobject{currentmarker}{\pgfqpoint{0.000000in}{-0.020833in}}{\pgfqpoint{0.000000in}{0.000000in}}{%
\pgfpathmoveto{\pgfqpoint{0.000000in}{0.000000in}}%
\pgfpathlineto{\pgfqpoint{0.000000in}{-0.020833in}}%
\pgfusepath{stroke,fill}%
}%
\begin{pgfscope}%
\pgfsys@transformshift{1.830272in}{2.414488in}%
\pgfsys@useobject{currentmarker}{}%
\end{pgfscope}%
\end{pgfscope}%
\begin{pgfscope}%
\pgfsetbuttcap%
\pgfsetroundjoin%
\definecolor{currentfill}{rgb}{0.000000,0.000000,0.000000}%
\pgfsetfillcolor{currentfill}%
\pgfsetlinewidth{0.401500pt}%
\definecolor{currentstroke}{rgb}{0.000000,0.000000,0.000000}%
\pgfsetstrokecolor{currentstroke}%
\pgfsetdash{}{0pt}%
\pgfsys@defobject{currentmarker}{\pgfqpoint{0.000000in}{0.000000in}}{\pgfqpoint{0.000000in}{0.020833in}}{%
\pgfpathmoveto{\pgfqpoint{0.000000in}{0.000000in}}%
\pgfpathlineto{\pgfqpoint{0.000000in}{0.020833in}}%
\pgfusepath{stroke,fill}%
}%
\begin{pgfscope}%
\pgfsys@transformshift{1.998883in}{0.420000in}%
\pgfsys@useobject{currentmarker}{}%
\end{pgfscope}%
\end{pgfscope}%
\begin{pgfscope}%
\pgfsetbuttcap%
\pgfsetroundjoin%
\definecolor{currentfill}{rgb}{0.000000,0.000000,0.000000}%
\pgfsetfillcolor{currentfill}%
\pgfsetlinewidth{0.401500pt}%
\definecolor{currentstroke}{rgb}{0.000000,0.000000,0.000000}%
\pgfsetstrokecolor{currentstroke}%
\pgfsetdash{}{0pt}%
\pgfsys@defobject{currentmarker}{\pgfqpoint{0.000000in}{-0.020833in}}{\pgfqpoint{0.000000in}{0.000000in}}{%
\pgfpathmoveto{\pgfqpoint{0.000000in}{0.000000in}}%
\pgfpathlineto{\pgfqpoint{0.000000in}{-0.020833in}}%
\pgfusepath{stroke,fill}%
}%
\begin{pgfscope}%
\pgfsys@transformshift{1.998883in}{2.414488in}%
\pgfsys@useobject{currentmarker}{}%
\end{pgfscope}%
\end{pgfscope}%
\begin{pgfscope}%
\pgfsetbuttcap%
\pgfsetroundjoin%
\definecolor{currentfill}{rgb}{0.000000,0.000000,0.000000}%
\pgfsetfillcolor{currentfill}%
\pgfsetlinewidth{0.401500pt}%
\definecolor{currentstroke}{rgb}{0.000000,0.000000,0.000000}%
\pgfsetstrokecolor{currentstroke}%
\pgfsetdash{}{0pt}%
\pgfsys@defobject{currentmarker}{\pgfqpoint{0.000000in}{0.000000in}}{\pgfqpoint{0.000000in}{0.020833in}}{%
\pgfpathmoveto{\pgfqpoint{0.000000in}{0.000000in}}%
\pgfpathlineto{\pgfqpoint{0.000000in}{0.020833in}}%
\pgfusepath{stroke,fill}%
}%
\begin{pgfscope}%
\pgfsys@transformshift{2.167493in}{0.420000in}%
\pgfsys@useobject{currentmarker}{}%
\end{pgfscope}%
\end{pgfscope}%
\begin{pgfscope}%
\pgfsetbuttcap%
\pgfsetroundjoin%
\definecolor{currentfill}{rgb}{0.000000,0.000000,0.000000}%
\pgfsetfillcolor{currentfill}%
\pgfsetlinewidth{0.401500pt}%
\definecolor{currentstroke}{rgb}{0.000000,0.000000,0.000000}%
\pgfsetstrokecolor{currentstroke}%
\pgfsetdash{}{0pt}%
\pgfsys@defobject{currentmarker}{\pgfqpoint{0.000000in}{-0.020833in}}{\pgfqpoint{0.000000in}{0.000000in}}{%
\pgfpathmoveto{\pgfqpoint{0.000000in}{0.000000in}}%
\pgfpathlineto{\pgfqpoint{0.000000in}{-0.020833in}}%
\pgfusepath{stroke,fill}%
}%
\begin{pgfscope}%
\pgfsys@transformshift{2.167493in}{2.414488in}%
\pgfsys@useobject{currentmarker}{}%
\end{pgfscope}%
\end{pgfscope}%
\begin{pgfscope}%
\pgfsetbuttcap%
\pgfsetroundjoin%
\definecolor{currentfill}{rgb}{0.000000,0.000000,0.000000}%
\pgfsetfillcolor{currentfill}%
\pgfsetlinewidth{0.401500pt}%
\definecolor{currentstroke}{rgb}{0.000000,0.000000,0.000000}%
\pgfsetstrokecolor{currentstroke}%
\pgfsetdash{}{0pt}%
\pgfsys@defobject{currentmarker}{\pgfqpoint{0.000000in}{0.000000in}}{\pgfqpoint{0.000000in}{0.020833in}}{%
\pgfpathmoveto{\pgfqpoint{0.000000in}{0.000000in}}%
\pgfpathlineto{\pgfqpoint{0.000000in}{0.020833in}}%
\pgfusepath{stroke,fill}%
}%
\begin{pgfscope}%
\pgfsys@transformshift{2.504714in}{0.420000in}%
\pgfsys@useobject{currentmarker}{}%
\end{pgfscope}%
\end{pgfscope}%
\begin{pgfscope}%
\pgfsetbuttcap%
\pgfsetroundjoin%
\definecolor{currentfill}{rgb}{0.000000,0.000000,0.000000}%
\pgfsetfillcolor{currentfill}%
\pgfsetlinewidth{0.401500pt}%
\definecolor{currentstroke}{rgb}{0.000000,0.000000,0.000000}%
\pgfsetstrokecolor{currentstroke}%
\pgfsetdash{}{0pt}%
\pgfsys@defobject{currentmarker}{\pgfqpoint{0.000000in}{-0.020833in}}{\pgfqpoint{0.000000in}{0.000000in}}{%
\pgfpathmoveto{\pgfqpoint{0.000000in}{0.000000in}}%
\pgfpathlineto{\pgfqpoint{0.000000in}{-0.020833in}}%
\pgfusepath{stroke,fill}%
}%
\begin{pgfscope}%
\pgfsys@transformshift{2.504714in}{2.414488in}%
\pgfsys@useobject{currentmarker}{}%
\end{pgfscope}%
\end{pgfscope}%
\begin{pgfscope}%
\pgfsetbuttcap%
\pgfsetroundjoin%
\definecolor{currentfill}{rgb}{0.000000,0.000000,0.000000}%
\pgfsetfillcolor{currentfill}%
\pgfsetlinewidth{0.401500pt}%
\definecolor{currentstroke}{rgb}{0.000000,0.000000,0.000000}%
\pgfsetstrokecolor{currentstroke}%
\pgfsetdash{}{0pt}%
\pgfsys@defobject{currentmarker}{\pgfqpoint{0.000000in}{0.000000in}}{\pgfqpoint{0.000000in}{0.020833in}}{%
\pgfpathmoveto{\pgfqpoint{0.000000in}{0.000000in}}%
\pgfpathlineto{\pgfqpoint{0.000000in}{0.020833in}}%
\pgfusepath{stroke,fill}%
}%
\begin{pgfscope}%
\pgfsys@transformshift{2.673324in}{0.420000in}%
\pgfsys@useobject{currentmarker}{}%
\end{pgfscope}%
\end{pgfscope}%
\begin{pgfscope}%
\pgfsetbuttcap%
\pgfsetroundjoin%
\definecolor{currentfill}{rgb}{0.000000,0.000000,0.000000}%
\pgfsetfillcolor{currentfill}%
\pgfsetlinewidth{0.401500pt}%
\definecolor{currentstroke}{rgb}{0.000000,0.000000,0.000000}%
\pgfsetstrokecolor{currentstroke}%
\pgfsetdash{}{0pt}%
\pgfsys@defobject{currentmarker}{\pgfqpoint{0.000000in}{-0.020833in}}{\pgfqpoint{0.000000in}{0.000000in}}{%
\pgfpathmoveto{\pgfqpoint{0.000000in}{0.000000in}}%
\pgfpathlineto{\pgfqpoint{0.000000in}{-0.020833in}}%
\pgfusepath{stroke,fill}%
}%
\begin{pgfscope}%
\pgfsys@transformshift{2.673324in}{2.414488in}%
\pgfsys@useobject{currentmarker}{}%
\end{pgfscope}%
\end{pgfscope}%
\begin{pgfscope}%
\pgfsetbuttcap%
\pgfsetroundjoin%
\definecolor{currentfill}{rgb}{0.000000,0.000000,0.000000}%
\pgfsetfillcolor{currentfill}%
\pgfsetlinewidth{0.401500pt}%
\definecolor{currentstroke}{rgb}{0.000000,0.000000,0.000000}%
\pgfsetstrokecolor{currentstroke}%
\pgfsetdash{}{0pt}%
\pgfsys@defobject{currentmarker}{\pgfqpoint{0.000000in}{0.000000in}}{\pgfqpoint{0.000000in}{0.020833in}}{%
\pgfpathmoveto{\pgfqpoint{0.000000in}{0.000000in}}%
\pgfpathlineto{\pgfqpoint{0.000000in}{0.020833in}}%
\pgfusepath{stroke,fill}%
}%
\begin{pgfscope}%
\pgfsys@transformshift{2.841934in}{0.420000in}%
\pgfsys@useobject{currentmarker}{}%
\end{pgfscope}%
\end{pgfscope}%
\begin{pgfscope}%
\pgfsetbuttcap%
\pgfsetroundjoin%
\definecolor{currentfill}{rgb}{0.000000,0.000000,0.000000}%
\pgfsetfillcolor{currentfill}%
\pgfsetlinewidth{0.401500pt}%
\definecolor{currentstroke}{rgb}{0.000000,0.000000,0.000000}%
\pgfsetstrokecolor{currentstroke}%
\pgfsetdash{}{0pt}%
\pgfsys@defobject{currentmarker}{\pgfqpoint{0.000000in}{-0.020833in}}{\pgfqpoint{0.000000in}{0.000000in}}{%
\pgfpathmoveto{\pgfqpoint{0.000000in}{0.000000in}}%
\pgfpathlineto{\pgfqpoint{0.000000in}{-0.020833in}}%
\pgfusepath{stroke,fill}%
}%
\begin{pgfscope}%
\pgfsys@transformshift{2.841934in}{2.414488in}%
\pgfsys@useobject{currentmarker}{}%
\end{pgfscope}%
\end{pgfscope}%
\begin{pgfscope}%
\pgfsetbuttcap%
\pgfsetroundjoin%
\definecolor{currentfill}{rgb}{0.000000,0.000000,0.000000}%
\pgfsetfillcolor{currentfill}%
\pgfsetlinewidth{0.401500pt}%
\definecolor{currentstroke}{rgb}{0.000000,0.000000,0.000000}%
\pgfsetstrokecolor{currentstroke}%
\pgfsetdash{}{0pt}%
\pgfsys@defobject{currentmarker}{\pgfqpoint{0.000000in}{0.000000in}}{\pgfqpoint{0.000000in}{0.020833in}}{%
\pgfpathmoveto{\pgfqpoint{0.000000in}{0.000000in}}%
\pgfpathlineto{\pgfqpoint{0.000000in}{0.020833in}}%
\pgfusepath{stroke,fill}%
}%
\begin{pgfscope}%
\pgfsys@transformshift{3.010545in}{0.420000in}%
\pgfsys@useobject{currentmarker}{}%
\end{pgfscope}%
\end{pgfscope}%
\begin{pgfscope}%
\pgfsetbuttcap%
\pgfsetroundjoin%
\definecolor{currentfill}{rgb}{0.000000,0.000000,0.000000}%
\pgfsetfillcolor{currentfill}%
\pgfsetlinewidth{0.401500pt}%
\definecolor{currentstroke}{rgb}{0.000000,0.000000,0.000000}%
\pgfsetstrokecolor{currentstroke}%
\pgfsetdash{}{0pt}%
\pgfsys@defobject{currentmarker}{\pgfqpoint{0.000000in}{-0.020833in}}{\pgfqpoint{0.000000in}{0.000000in}}{%
\pgfpathmoveto{\pgfqpoint{0.000000in}{0.000000in}}%
\pgfpathlineto{\pgfqpoint{0.000000in}{-0.020833in}}%
\pgfusepath{stroke,fill}%
}%
\begin{pgfscope}%
\pgfsys@transformshift{3.010545in}{2.414488in}%
\pgfsys@useobject{currentmarker}{}%
\end{pgfscope}%
\end{pgfscope}%
\begin{pgfscope}%
\definecolor{textcolor}{rgb}{0.000000,0.000000,0.000000}%
\pgfsetstrokecolor{textcolor}%
\pgfsetfillcolor{textcolor}%
\pgftext[x=1.844055in,y=0.208426in,,top]{\color{textcolor}{\rmfamily\fontsize{12.000000}{14.400000}\selectfont\catcode`\^=\active\def^{\ifmmode\sp\else\^{}\fi}\catcode`\%=\active\def%{\%}$s/\delta_{0}$}}%
\end{pgfscope}%
\begin{pgfscope}%
\pgfsetbuttcap%
\pgfsetroundjoin%
\definecolor{currentfill}{rgb}{0.000000,0.000000,0.000000}%
\pgfsetfillcolor{currentfill}%
\pgfsetlinewidth{0.501875pt}%
\definecolor{currentstroke}{rgb}{0.000000,0.000000,0.000000}%
\pgfsetstrokecolor{currentstroke}%
\pgfsetdash{}{0pt}%
\pgfsys@defobject{currentmarker}{\pgfqpoint{0.000000in}{0.000000in}}{\pgfqpoint{0.041667in}{0.000000in}}{%
\pgfpathmoveto{\pgfqpoint{0.000000in}{0.000000in}}%
\pgfpathlineto{\pgfqpoint{0.041667in}{0.000000in}}%
\pgfusepath{stroke,fill}%
}%
\begin{pgfscope}%
\pgfsys@transformshift{0.650000in}{0.420000in}%
\pgfsys@useobject{currentmarker}{}%
\end{pgfscope}%
\end{pgfscope}%
\begin{pgfscope}%
\pgfsetbuttcap%
\pgfsetroundjoin%
\definecolor{currentfill}{rgb}{0.000000,0.000000,0.000000}%
\pgfsetfillcolor{currentfill}%
\pgfsetlinewidth{0.501875pt}%
\definecolor{currentstroke}{rgb}{0.000000,0.000000,0.000000}%
\pgfsetstrokecolor{currentstroke}%
\pgfsetdash{}{0pt}%
\pgfsys@defobject{currentmarker}{\pgfqpoint{-0.041667in}{0.000000in}}{\pgfqpoint{-0.000000in}{0.000000in}}{%
\pgfpathmoveto{\pgfqpoint{-0.000000in}{0.000000in}}%
\pgfpathlineto{\pgfqpoint{-0.041667in}{0.000000in}}%
\pgfusepath{stroke,fill}%
}%
\begin{pgfscope}%
\pgfsys@transformshift{3.038110in}{0.420000in}%
\pgfsys@useobject{currentmarker}{}%
\end{pgfscope}%
\end{pgfscope}%
\begin{pgfscope}%
\definecolor{textcolor}{rgb}{0.000000,0.000000,0.000000}%
\pgfsetstrokecolor{textcolor}%
\pgfsetfillcolor{textcolor}%
\pgftext[x=0.254976in, y=0.367193in, left, base]{\color{textcolor}{\rmfamily\fontsize{11.000000}{13.200000}\selectfont\catcode`\^=\active\def^{\ifmmode\sp\else\^{}\fi}\catcode`\%=\active\def%{\%}0.000}}%
\end{pgfscope}%
\begin{pgfscope}%
\pgfsetbuttcap%
\pgfsetroundjoin%
\definecolor{currentfill}{rgb}{0.000000,0.000000,0.000000}%
\pgfsetfillcolor{currentfill}%
\pgfsetlinewidth{0.501875pt}%
\definecolor{currentstroke}{rgb}{0.000000,0.000000,0.000000}%
\pgfsetstrokecolor{currentstroke}%
\pgfsetdash{}{0pt}%
\pgfsys@defobject{currentmarker}{\pgfqpoint{0.000000in}{0.000000in}}{\pgfqpoint{0.041667in}{0.000000in}}{%
\pgfpathmoveto{\pgfqpoint{0.000000in}{0.000000in}}%
\pgfpathlineto{\pgfqpoint{0.041667in}{0.000000in}}%
\pgfusepath{stroke,fill}%
}%
\begin{pgfscope}%
\pgfsys@transformshift{0.650000in}{0.752415in}%
\pgfsys@useobject{currentmarker}{}%
\end{pgfscope}%
\end{pgfscope}%
\begin{pgfscope}%
\pgfsetbuttcap%
\pgfsetroundjoin%
\definecolor{currentfill}{rgb}{0.000000,0.000000,0.000000}%
\pgfsetfillcolor{currentfill}%
\pgfsetlinewidth{0.501875pt}%
\definecolor{currentstroke}{rgb}{0.000000,0.000000,0.000000}%
\pgfsetstrokecolor{currentstroke}%
\pgfsetdash{}{0pt}%
\pgfsys@defobject{currentmarker}{\pgfqpoint{-0.041667in}{0.000000in}}{\pgfqpoint{-0.000000in}{0.000000in}}{%
\pgfpathmoveto{\pgfqpoint{-0.000000in}{0.000000in}}%
\pgfpathlineto{\pgfqpoint{-0.041667in}{0.000000in}}%
\pgfusepath{stroke,fill}%
}%
\begin{pgfscope}%
\pgfsys@transformshift{3.038110in}{0.752415in}%
\pgfsys@useobject{currentmarker}{}%
\end{pgfscope}%
\end{pgfscope}%
\begin{pgfscope}%
\definecolor{textcolor}{rgb}{0.000000,0.000000,0.000000}%
\pgfsetstrokecolor{textcolor}%
\pgfsetfillcolor{textcolor}%
\pgftext[x=0.254976in, y=0.699608in, left, base]{\color{textcolor}{\rmfamily\fontsize{11.000000}{13.200000}\selectfont\catcode`\^=\active\def^{\ifmmode\sp\else\^{}\fi}\catcode`\%=\active\def%{\%}0.001}}%
\end{pgfscope}%
\begin{pgfscope}%
\pgfsetbuttcap%
\pgfsetroundjoin%
\definecolor{currentfill}{rgb}{0.000000,0.000000,0.000000}%
\pgfsetfillcolor{currentfill}%
\pgfsetlinewidth{0.501875pt}%
\definecolor{currentstroke}{rgb}{0.000000,0.000000,0.000000}%
\pgfsetstrokecolor{currentstroke}%
\pgfsetdash{}{0pt}%
\pgfsys@defobject{currentmarker}{\pgfqpoint{0.000000in}{0.000000in}}{\pgfqpoint{0.041667in}{0.000000in}}{%
\pgfpathmoveto{\pgfqpoint{0.000000in}{0.000000in}}%
\pgfpathlineto{\pgfqpoint{0.041667in}{0.000000in}}%
\pgfusepath{stroke,fill}%
}%
\begin{pgfscope}%
\pgfsys@transformshift{0.650000in}{1.084829in}%
\pgfsys@useobject{currentmarker}{}%
\end{pgfscope}%
\end{pgfscope}%
\begin{pgfscope}%
\pgfsetbuttcap%
\pgfsetroundjoin%
\definecolor{currentfill}{rgb}{0.000000,0.000000,0.000000}%
\pgfsetfillcolor{currentfill}%
\pgfsetlinewidth{0.501875pt}%
\definecolor{currentstroke}{rgb}{0.000000,0.000000,0.000000}%
\pgfsetstrokecolor{currentstroke}%
\pgfsetdash{}{0pt}%
\pgfsys@defobject{currentmarker}{\pgfqpoint{-0.041667in}{0.000000in}}{\pgfqpoint{-0.000000in}{0.000000in}}{%
\pgfpathmoveto{\pgfqpoint{-0.000000in}{0.000000in}}%
\pgfpathlineto{\pgfqpoint{-0.041667in}{0.000000in}}%
\pgfusepath{stroke,fill}%
}%
\begin{pgfscope}%
\pgfsys@transformshift{3.038110in}{1.084829in}%
\pgfsys@useobject{currentmarker}{}%
\end{pgfscope}%
\end{pgfscope}%
\begin{pgfscope}%
\definecolor{textcolor}{rgb}{0.000000,0.000000,0.000000}%
\pgfsetstrokecolor{textcolor}%
\pgfsetfillcolor{textcolor}%
\pgftext[x=0.254976in, y=1.032023in, left, base]{\color{textcolor}{\rmfamily\fontsize{11.000000}{13.200000}\selectfont\catcode`\^=\active\def^{\ifmmode\sp\else\^{}\fi}\catcode`\%=\active\def%{\%}0.002}}%
\end{pgfscope}%
\begin{pgfscope}%
\pgfsetbuttcap%
\pgfsetroundjoin%
\definecolor{currentfill}{rgb}{0.000000,0.000000,0.000000}%
\pgfsetfillcolor{currentfill}%
\pgfsetlinewidth{0.501875pt}%
\definecolor{currentstroke}{rgb}{0.000000,0.000000,0.000000}%
\pgfsetstrokecolor{currentstroke}%
\pgfsetdash{}{0pt}%
\pgfsys@defobject{currentmarker}{\pgfqpoint{0.000000in}{0.000000in}}{\pgfqpoint{0.041667in}{0.000000in}}{%
\pgfpathmoveto{\pgfqpoint{0.000000in}{0.000000in}}%
\pgfpathlineto{\pgfqpoint{0.041667in}{0.000000in}}%
\pgfusepath{stroke,fill}%
}%
\begin{pgfscope}%
\pgfsys@transformshift{0.650000in}{1.417244in}%
\pgfsys@useobject{currentmarker}{}%
\end{pgfscope}%
\end{pgfscope}%
\begin{pgfscope}%
\pgfsetbuttcap%
\pgfsetroundjoin%
\definecolor{currentfill}{rgb}{0.000000,0.000000,0.000000}%
\pgfsetfillcolor{currentfill}%
\pgfsetlinewidth{0.501875pt}%
\definecolor{currentstroke}{rgb}{0.000000,0.000000,0.000000}%
\pgfsetstrokecolor{currentstroke}%
\pgfsetdash{}{0pt}%
\pgfsys@defobject{currentmarker}{\pgfqpoint{-0.041667in}{0.000000in}}{\pgfqpoint{-0.000000in}{0.000000in}}{%
\pgfpathmoveto{\pgfqpoint{-0.000000in}{0.000000in}}%
\pgfpathlineto{\pgfqpoint{-0.041667in}{0.000000in}}%
\pgfusepath{stroke,fill}%
}%
\begin{pgfscope}%
\pgfsys@transformshift{3.038110in}{1.417244in}%
\pgfsys@useobject{currentmarker}{}%
\end{pgfscope}%
\end{pgfscope}%
\begin{pgfscope}%
\definecolor{textcolor}{rgb}{0.000000,0.000000,0.000000}%
\pgfsetstrokecolor{textcolor}%
\pgfsetfillcolor{textcolor}%
\pgftext[x=0.254976in, y=1.364437in, left, base]{\color{textcolor}{\rmfamily\fontsize{11.000000}{13.200000}\selectfont\catcode`\^=\active\def^{\ifmmode\sp\else\^{}\fi}\catcode`\%=\active\def%{\%}0.003}}%
\end{pgfscope}%
\begin{pgfscope}%
\pgfsetbuttcap%
\pgfsetroundjoin%
\definecolor{currentfill}{rgb}{0.000000,0.000000,0.000000}%
\pgfsetfillcolor{currentfill}%
\pgfsetlinewidth{0.501875pt}%
\definecolor{currentstroke}{rgb}{0.000000,0.000000,0.000000}%
\pgfsetstrokecolor{currentstroke}%
\pgfsetdash{}{0pt}%
\pgfsys@defobject{currentmarker}{\pgfqpoint{0.000000in}{0.000000in}}{\pgfqpoint{0.041667in}{0.000000in}}{%
\pgfpathmoveto{\pgfqpoint{0.000000in}{0.000000in}}%
\pgfpathlineto{\pgfqpoint{0.041667in}{0.000000in}}%
\pgfusepath{stroke,fill}%
}%
\begin{pgfscope}%
\pgfsys@transformshift{0.650000in}{1.749659in}%
\pgfsys@useobject{currentmarker}{}%
\end{pgfscope}%
\end{pgfscope}%
\begin{pgfscope}%
\pgfsetbuttcap%
\pgfsetroundjoin%
\definecolor{currentfill}{rgb}{0.000000,0.000000,0.000000}%
\pgfsetfillcolor{currentfill}%
\pgfsetlinewidth{0.501875pt}%
\definecolor{currentstroke}{rgb}{0.000000,0.000000,0.000000}%
\pgfsetstrokecolor{currentstroke}%
\pgfsetdash{}{0pt}%
\pgfsys@defobject{currentmarker}{\pgfqpoint{-0.041667in}{0.000000in}}{\pgfqpoint{-0.000000in}{0.000000in}}{%
\pgfpathmoveto{\pgfqpoint{-0.000000in}{0.000000in}}%
\pgfpathlineto{\pgfqpoint{-0.041667in}{0.000000in}}%
\pgfusepath{stroke,fill}%
}%
\begin{pgfscope}%
\pgfsys@transformshift{3.038110in}{1.749659in}%
\pgfsys@useobject{currentmarker}{}%
\end{pgfscope}%
\end{pgfscope}%
\begin{pgfscope}%
\definecolor{textcolor}{rgb}{0.000000,0.000000,0.000000}%
\pgfsetstrokecolor{textcolor}%
\pgfsetfillcolor{textcolor}%
\pgftext[x=0.254976in, y=1.696852in, left, base]{\color{textcolor}{\rmfamily\fontsize{11.000000}{13.200000}\selectfont\catcode`\^=\active\def^{\ifmmode\sp\else\^{}\fi}\catcode`\%=\active\def%{\%}0.004}}%
\end{pgfscope}%
\begin{pgfscope}%
\pgfsetbuttcap%
\pgfsetroundjoin%
\definecolor{currentfill}{rgb}{0.000000,0.000000,0.000000}%
\pgfsetfillcolor{currentfill}%
\pgfsetlinewidth{0.501875pt}%
\definecolor{currentstroke}{rgb}{0.000000,0.000000,0.000000}%
\pgfsetstrokecolor{currentstroke}%
\pgfsetdash{}{0pt}%
\pgfsys@defobject{currentmarker}{\pgfqpoint{0.000000in}{0.000000in}}{\pgfqpoint{0.041667in}{0.000000in}}{%
\pgfpathmoveto{\pgfqpoint{0.000000in}{0.000000in}}%
\pgfpathlineto{\pgfqpoint{0.041667in}{0.000000in}}%
\pgfusepath{stroke,fill}%
}%
\begin{pgfscope}%
\pgfsys@transformshift{0.650000in}{2.082073in}%
\pgfsys@useobject{currentmarker}{}%
\end{pgfscope}%
\end{pgfscope}%
\begin{pgfscope}%
\pgfsetbuttcap%
\pgfsetroundjoin%
\definecolor{currentfill}{rgb}{0.000000,0.000000,0.000000}%
\pgfsetfillcolor{currentfill}%
\pgfsetlinewidth{0.501875pt}%
\definecolor{currentstroke}{rgb}{0.000000,0.000000,0.000000}%
\pgfsetstrokecolor{currentstroke}%
\pgfsetdash{}{0pt}%
\pgfsys@defobject{currentmarker}{\pgfqpoint{-0.041667in}{0.000000in}}{\pgfqpoint{-0.000000in}{0.000000in}}{%
\pgfpathmoveto{\pgfqpoint{-0.000000in}{0.000000in}}%
\pgfpathlineto{\pgfqpoint{-0.041667in}{0.000000in}}%
\pgfusepath{stroke,fill}%
}%
\begin{pgfscope}%
\pgfsys@transformshift{3.038110in}{2.082073in}%
\pgfsys@useobject{currentmarker}{}%
\end{pgfscope}%
\end{pgfscope}%
\begin{pgfscope}%
\definecolor{textcolor}{rgb}{0.000000,0.000000,0.000000}%
\pgfsetstrokecolor{textcolor}%
\pgfsetfillcolor{textcolor}%
\pgftext[x=0.254976in, y=2.029267in, left, base]{\color{textcolor}{\rmfamily\fontsize{11.000000}{13.200000}\selectfont\catcode`\^=\active\def^{\ifmmode\sp\else\^{}\fi}\catcode`\%=\active\def%{\%}0.005}}%
\end{pgfscope}%
\begin{pgfscope}%
\pgfsetbuttcap%
\pgfsetroundjoin%
\definecolor{currentfill}{rgb}{0.000000,0.000000,0.000000}%
\pgfsetfillcolor{currentfill}%
\pgfsetlinewidth{0.501875pt}%
\definecolor{currentstroke}{rgb}{0.000000,0.000000,0.000000}%
\pgfsetstrokecolor{currentstroke}%
\pgfsetdash{}{0pt}%
\pgfsys@defobject{currentmarker}{\pgfqpoint{0.000000in}{0.000000in}}{\pgfqpoint{0.041667in}{0.000000in}}{%
\pgfpathmoveto{\pgfqpoint{0.000000in}{0.000000in}}%
\pgfpathlineto{\pgfqpoint{0.041667in}{0.000000in}}%
\pgfusepath{stroke,fill}%
}%
\begin{pgfscope}%
\pgfsys@transformshift{0.650000in}{2.414488in}%
\pgfsys@useobject{currentmarker}{}%
\end{pgfscope}%
\end{pgfscope}%
\begin{pgfscope}%
\pgfsetbuttcap%
\pgfsetroundjoin%
\definecolor{currentfill}{rgb}{0.000000,0.000000,0.000000}%
\pgfsetfillcolor{currentfill}%
\pgfsetlinewidth{0.501875pt}%
\definecolor{currentstroke}{rgb}{0.000000,0.000000,0.000000}%
\pgfsetstrokecolor{currentstroke}%
\pgfsetdash{}{0pt}%
\pgfsys@defobject{currentmarker}{\pgfqpoint{-0.041667in}{0.000000in}}{\pgfqpoint{-0.000000in}{0.000000in}}{%
\pgfpathmoveto{\pgfqpoint{-0.000000in}{0.000000in}}%
\pgfpathlineto{\pgfqpoint{-0.041667in}{0.000000in}}%
\pgfusepath{stroke,fill}%
}%
\begin{pgfscope}%
\pgfsys@transformshift{3.038110in}{2.414488in}%
\pgfsys@useobject{currentmarker}{}%
\end{pgfscope}%
\end{pgfscope}%
\begin{pgfscope}%
\definecolor{textcolor}{rgb}{0.000000,0.000000,0.000000}%
\pgfsetstrokecolor{textcolor}%
\pgfsetfillcolor{textcolor}%
\pgftext[x=0.254976in, y=2.361681in, left, base]{\color{textcolor}{\rmfamily\fontsize{11.000000}{13.200000}\selectfont\catcode`\^=\active\def^{\ifmmode\sp\else\^{}\fi}\catcode`\%=\active\def%{\%}0.006}}%
\end{pgfscope}%
\begin{pgfscope}%
\pgfsetbuttcap%
\pgfsetroundjoin%
\definecolor{currentfill}{rgb}{0.000000,0.000000,0.000000}%
\pgfsetfillcolor{currentfill}%
\pgfsetlinewidth{0.401500pt}%
\definecolor{currentstroke}{rgb}{0.000000,0.000000,0.000000}%
\pgfsetstrokecolor{currentstroke}%
\pgfsetdash{}{0pt}%
\pgfsys@defobject{currentmarker}{\pgfqpoint{0.000000in}{0.000000in}}{\pgfqpoint{0.020833in}{0.000000in}}{%
\pgfpathmoveto{\pgfqpoint{0.000000in}{0.000000in}}%
\pgfpathlineto{\pgfqpoint{0.020833in}{0.000000in}}%
\pgfusepath{stroke,fill}%
}%
\begin{pgfscope}%
\pgfsys@transformshift{0.650000in}{0.486483in}%
\pgfsys@useobject{currentmarker}{}%
\end{pgfscope}%
\end{pgfscope}%
\begin{pgfscope}%
\pgfsetbuttcap%
\pgfsetroundjoin%
\definecolor{currentfill}{rgb}{0.000000,0.000000,0.000000}%
\pgfsetfillcolor{currentfill}%
\pgfsetlinewidth{0.401500pt}%
\definecolor{currentstroke}{rgb}{0.000000,0.000000,0.000000}%
\pgfsetstrokecolor{currentstroke}%
\pgfsetdash{}{0pt}%
\pgfsys@defobject{currentmarker}{\pgfqpoint{-0.020833in}{0.000000in}}{\pgfqpoint{-0.000000in}{0.000000in}}{%
\pgfpathmoveto{\pgfqpoint{-0.000000in}{0.000000in}}%
\pgfpathlineto{\pgfqpoint{-0.020833in}{0.000000in}}%
\pgfusepath{stroke,fill}%
}%
\begin{pgfscope}%
\pgfsys@transformshift{3.038110in}{0.486483in}%
\pgfsys@useobject{currentmarker}{}%
\end{pgfscope}%
\end{pgfscope}%
\begin{pgfscope}%
\pgfsetbuttcap%
\pgfsetroundjoin%
\definecolor{currentfill}{rgb}{0.000000,0.000000,0.000000}%
\pgfsetfillcolor{currentfill}%
\pgfsetlinewidth{0.401500pt}%
\definecolor{currentstroke}{rgb}{0.000000,0.000000,0.000000}%
\pgfsetstrokecolor{currentstroke}%
\pgfsetdash{}{0pt}%
\pgfsys@defobject{currentmarker}{\pgfqpoint{0.000000in}{0.000000in}}{\pgfqpoint{0.020833in}{0.000000in}}{%
\pgfpathmoveto{\pgfqpoint{0.000000in}{0.000000in}}%
\pgfpathlineto{\pgfqpoint{0.020833in}{0.000000in}}%
\pgfusepath{stroke,fill}%
}%
\begin{pgfscope}%
\pgfsys@transformshift{0.650000in}{0.552966in}%
\pgfsys@useobject{currentmarker}{}%
\end{pgfscope}%
\end{pgfscope}%
\begin{pgfscope}%
\pgfsetbuttcap%
\pgfsetroundjoin%
\definecolor{currentfill}{rgb}{0.000000,0.000000,0.000000}%
\pgfsetfillcolor{currentfill}%
\pgfsetlinewidth{0.401500pt}%
\definecolor{currentstroke}{rgb}{0.000000,0.000000,0.000000}%
\pgfsetstrokecolor{currentstroke}%
\pgfsetdash{}{0pt}%
\pgfsys@defobject{currentmarker}{\pgfqpoint{-0.020833in}{0.000000in}}{\pgfqpoint{-0.000000in}{0.000000in}}{%
\pgfpathmoveto{\pgfqpoint{-0.000000in}{0.000000in}}%
\pgfpathlineto{\pgfqpoint{-0.020833in}{0.000000in}}%
\pgfusepath{stroke,fill}%
}%
\begin{pgfscope}%
\pgfsys@transformshift{3.038110in}{0.552966in}%
\pgfsys@useobject{currentmarker}{}%
\end{pgfscope}%
\end{pgfscope}%
\begin{pgfscope}%
\pgfsetbuttcap%
\pgfsetroundjoin%
\definecolor{currentfill}{rgb}{0.000000,0.000000,0.000000}%
\pgfsetfillcolor{currentfill}%
\pgfsetlinewidth{0.401500pt}%
\definecolor{currentstroke}{rgb}{0.000000,0.000000,0.000000}%
\pgfsetstrokecolor{currentstroke}%
\pgfsetdash{}{0pt}%
\pgfsys@defobject{currentmarker}{\pgfqpoint{0.000000in}{0.000000in}}{\pgfqpoint{0.020833in}{0.000000in}}{%
\pgfpathmoveto{\pgfqpoint{0.000000in}{0.000000in}}%
\pgfpathlineto{\pgfqpoint{0.020833in}{0.000000in}}%
\pgfusepath{stroke,fill}%
}%
\begin{pgfscope}%
\pgfsys@transformshift{0.650000in}{0.619449in}%
\pgfsys@useobject{currentmarker}{}%
\end{pgfscope}%
\end{pgfscope}%
\begin{pgfscope}%
\pgfsetbuttcap%
\pgfsetroundjoin%
\definecolor{currentfill}{rgb}{0.000000,0.000000,0.000000}%
\pgfsetfillcolor{currentfill}%
\pgfsetlinewidth{0.401500pt}%
\definecolor{currentstroke}{rgb}{0.000000,0.000000,0.000000}%
\pgfsetstrokecolor{currentstroke}%
\pgfsetdash{}{0pt}%
\pgfsys@defobject{currentmarker}{\pgfqpoint{-0.020833in}{0.000000in}}{\pgfqpoint{-0.000000in}{0.000000in}}{%
\pgfpathmoveto{\pgfqpoint{-0.000000in}{0.000000in}}%
\pgfpathlineto{\pgfqpoint{-0.020833in}{0.000000in}}%
\pgfusepath{stroke,fill}%
}%
\begin{pgfscope}%
\pgfsys@transformshift{3.038110in}{0.619449in}%
\pgfsys@useobject{currentmarker}{}%
\end{pgfscope}%
\end{pgfscope}%
\begin{pgfscope}%
\pgfsetbuttcap%
\pgfsetroundjoin%
\definecolor{currentfill}{rgb}{0.000000,0.000000,0.000000}%
\pgfsetfillcolor{currentfill}%
\pgfsetlinewidth{0.401500pt}%
\definecolor{currentstroke}{rgb}{0.000000,0.000000,0.000000}%
\pgfsetstrokecolor{currentstroke}%
\pgfsetdash{}{0pt}%
\pgfsys@defobject{currentmarker}{\pgfqpoint{0.000000in}{0.000000in}}{\pgfqpoint{0.020833in}{0.000000in}}{%
\pgfpathmoveto{\pgfqpoint{0.000000in}{0.000000in}}%
\pgfpathlineto{\pgfqpoint{0.020833in}{0.000000in}}%
\pgfusepath{stroke,fill}%
}%
\begin{pgfscope}%
\pgfsys@transformshift{0.650000in}{0.685932in}%
\pgfsys@useobject{currentmarker}{}%
\end{pgfscope}%
\end{pgfscope}%
\begin{pgfscope}%
\pgfsetbuttcap%
\pgfsetroundjoin%
\definecolor{currentfill}{rgb}{0.000000,0.000000,0.000000}%
\pgfsetfillcolor{currentfill}%
\pgfsetlinewidth{0.401500pt}%
\definecolor{currentstroke}{rgb}{0.000000,0.000000,0.000000}%
\pgfsetstrokecolor{currentstroke}%
\pgfsetdash{}{0pt}%
\pgfsys@defobject{currentmarker}{\pgfqpoint{-0.020833in}{0.000000in}}{\pgfqpoint{-0.000000in}{0.000000in}}{%
\pgfpathmoveto{\pgfqpoint{-0.000000in}{0.000000in}}%
\pgfpathlineto{\pgfqpoint{-0.020833in}{0.000000in}}%
\pgfusepath{stroke,fill}%
}%
\begin{pgfscope}%
\pgfsys@transformshift{3.038110in}{0.685932in}%
\pgfsys@useobject{currentmarker}{}%
\end{pgfscope}%
\end{pgfscope}%
\begin{pgfscope}%
\pgfsetbuttcap%
\pgfsetroundjoin%
\definecolor{currentfill}{rgb}{0.000000,0.000000,0.000000}%
\pgfsetfillcolor{currentfill}%
\pgfsetlinewidth{0.401500pt}%
\definecolor{currentstroke}{rgb}{0.000000,0.000000,0.000000}%
\pgfsetstrokecolor{currentstroke}%
\pgfsetdash{}{0pt}%
\pgfsys@defobject{currentmarker}{\pgfqpoint{0.000000in}{0.000000in}}{\pgfqpoint{0.020833in}{0.000000in}}{%
\pgfpathmoveto{\pgfqpoint{0.000000in}{0.000000in}}%
\pgfpathlineto{\pgfqpoint{0.020833in}{0.000000in}}%
\pgfusepath{stroke,fill}%
}%
\begin{pgfscope}%
\pgfsys@transformshift{0.650000in}{0.818898in}%
\pgfsys@useobject{currentmarker}{}%
\end{pgfscope}%
\end{pgfscope}%
\begin{pgfscope}%
\pgfsetbuttcap%
\pgfsetroundjoin%
\definecolor{currentfill}{rgb}{0.000000,0.000000,0.000000}%
\pgfsetfillcolor{currentfill}%
\pgfsetlinewidth{0.401500pt}%
\definecolor{currentstroke}{rgb}{0.000000,0.000000,0.000000}%
\pgfsetstrokecolor{currentstroke}%
\pgfsetdash{}{0pt}%
\pgfsys@defobject{currentmarker}{\pgfqpoint{-0.020833in}{0.000000in}}{\pgfqpoint{-0.000000in}{0.000000in}}{%
\pgfpathmoveto{\pgfqpoint{-0.000000in}{0.000000in}}%
\pgfpathlineto{\pgfqpoint{-0.020833in}{0.000000in}}%
\pgfusepath{stroke,fill}%
}%
\begin{pgfscope}%
\pgfsys@transformshift{3.038110in}{0.818898in}%
\pgfsys@useobject{currentmarker}{}%
\end{pgfscope}%
\end{pgfscope}%
\begin{pgfscope}%
\pgfsetbuttcap%
\pgfsetroundjoin%
\definecolor{currentfill}{rgb}{0.000000,0.000000,0.000000}%
\pgfsetfillcolor{currentfill}%
\pgfsetlinewidth{0.401500pt}%
\definecolor{currentstroke}{rgb}{0.000000,0.000000,0.000000}%
\pgfsetstrokecolor{currentstroke}%
\pgfsetdash{}{0pt}%
\pgfsys@defobject{currentmarker}{\pgfqpoint{0.000000in}{0.000000in}}{\pgfqpoint{0.020833in}{0.000000in}}{%
\pgfpathmoveto{\pgfqpoint{0.000000in}{0.000000in}}%
\pgfpathlineto{\pgfqpoint{0.020833in}{0.000000in}}%
\pgfusepath{stroke,fill}%
}%
\begin{pgfscope}%
\pgfsys@transformshift{0.650000in}{0.885381in}%
\pgfsys@useobject{currentmarker}{}%
\end{pgfscope}%
\end{pgfscope}%
\begin{pgfscope}%
\pgfsetbuttcap%
\pgfsetroundjoin%
\definecolor{currentfill}{rgb}{0.000000,0.000000,0.000000}%
\pgfsetfillcolor{currentfill}%
\pgfsetlinewidth{0.401500pt}%
\definecolor{currentstroke}{rgb}{0.000000,0.000000,0.000000}%
\pgfsetstrokecolor{currentstroke}%
\pgfsetdash{}{0pt}%
\pgfsys@defobject{currentmarker}{\pgfqpoint{-0.020833in}{0.000000in}}{\pgfqpoint{-0.000000in}{0.000000in}}{%
\pgfpathmoveto{\pgfqpoint{-0.000000in}{0.000000in}}%
\pgfpathlineto{\pgfqpoint{-0.020833in}{0.000000in}}%
\pgfusepath{stroke,fill}%
}%
\begin{pgfscope}%
\pgfsys@transformshift{3.038110in}{0.885381in}%
\pgfsys@useobject{currentmarker}{}%
\end{pgfscope}%
\end{pgfscope}%
\begin{pgfscope}%
\pgfsetbuttcap%
\pgfsetroundjoin%
\definecolor{currentfill}{rgb}{0.000000,0.000000,0.000000}%
\pgfsetfillcolor{currentfill}%
\pgfsetlinewidth{0.401500pt}%
\definecolor{currentstroke}{rgb}{0.000000,0.000000,0.000000}%
\pgfsetstrokecolor{currentstroke}%
\pgfsetdash{}{0pt}%
\pgfsys@defobject{currentmarker}{\pgfqpoint{0.000000in}{0.000000in}}{\pgfqpoint{0.020833in}{0.000000in}}{%
\pgfpathmoveto{\pgfqpoint{0.000000in}{0.000000in}}%
\pgfpathlineto{\pgfqpoint{0.020833in}{0.000000in}}%
\pgfusepath{stroke,fill}%
}%
\begin{pgfscope}%
\pgfsys@transformshift{0.650000in}{0.951863in}%
\pgfsys@useobject{currentmarker}{}%
\end{pgfscope}%
\end{pgfscope}%
\begin{pgfscope}%
\pgfsetbuttcap%
\pgfsetroundjoin%
\definecolor{currentfill}{rgb}{0.000000,0.000000,0.000000}%
\pgfsetfillcolor{currentfill}%
\pgfsetlinewidth{0.401500pt}%
\definecolor{currentstroke}{rgb}{0.000000,0.000000,0.000000}%
\pgfsetstrokecolor{currentstroke}%
\pgfsetdash{}{0pt}%
\pgfsys@defobject{currentmarker}{\pgfqpoint{-0.020833in}{0.000000in}}{\pgfqpoint{-0.000000in}{0.000000in}}{%
\pgfpathmoveto{\pgfqpoint{-0.000000in}{0.000000in}}%
\pgfpathlineto{\pgfqpoint{-0.020833in}{0.000000in}}%
\pgfusepath{stroke,fill}%
}%
\begin{pgfscope}%
\pgfsys@transformshift{3.038110in}{0.951863in}%
\pgfsys@useobject{currentmarker}{}%
\end{pgfscope}%
\end{pgfscope}%
\begin{pgfscope}%
\pgfsetbuttcap%
\pgfsetroundjoin%
\definecolor{currentfill}{rgb}{0.000000,0.000000,0.000000}%
\pgfsetfillcolor{currentfill}%
\pgfsetlinewidth{0.401500pt}%
\definecolor{currentstroke}{rgb}{0.000000,0.000000,0.000000}%
\pgfsetstrokecolor{currentstroke}%
\pgfsetdash{}{0pt}%
\pgfsys@defobject{currentmarker}{\pgfqpoint{0.000000in}{0.000000in}}{\pgfqpoint{0.020833in}{0.000000in}}{%
\pgfpathmoveto{\pgfqpoint{0.000000in}{0.000000in}}%
\pgfpathlineto{\pgfqpoint{0.020833in}{0.000000in}}%
\pgfusepath{stroke,fill}%
}%
\begin{pgfscope}%
\pgfsys@transformshift{0.650000in}{1.018346in}%
\pgfsys@useobject{currentmarker}{}%
\end{pgfscope}%
\end{pgfscope}%
\begin{pgfscope}%
\pgfsetbuttcap%
\pgfsetroundjoin%
\definecolor{currentfill}{rgb}{0.000000,0.000000,0.000000}%
\pgfsetfillcolor{currentfill}%
\pgfsetlinewidth{0.401500pt}%
\definecolor{currentstroke}{rgb}{0.000000,0.000000,0.000000}%
\pgfsetstrokecolor{currentstroke}%
\pgfsetdash{}{0pt}%
\pgfsys@defobject{currentmarker}{\pgfqpoint{-0.020833in}{0.000000in}}{\pgfqpoint{-0.000000in}{0.000000in}}{%
\pgfpathmoveto{\pgfqpoint{-0.000000in}{0.000000in}}%
\pgfpathlineto{\pgfqpoint{-0.020833in}{0.000000in}}%
\pgfusepath{stroke,fill}%
}%
\begin{pgfscope}%
\pgfsys@transformshift{3.038110in}{1.018346in}%
\pgfsys@useobject{currentmarker}{}%
\end{pgfscope}%
\end{pgfscope}%
\begin{pgfscope}%
\pgfsetbuttcap%
\pgfsetroundjoin%
\definecolor{currentfill}{rgb}{0.000000,0.000000,0.000000}%
\pgfsetfillcolor{currentfill}%
\pgfsetlinewidth{0.401500pt}%
\definecolor{currentstroke}{rgb}{0.000000,0.000000,0.000000}%
\pgfsetstrokecolor{currentstroke}%
\pgfsetdash{}{0pt}%
\pgfsys@defobject{currentmarker}{\pgfqpoint{0.000000in}{0.000000in}}{\pgfqpoint{0.020833in}{0.000000in}}{%
\pgfpathmoveto{\pgfqpoint{0.000000in}{0.000000in}}%
\pgfpathlineto{\pgfqpoint{0.020833in}{0.000000in}}%
\pgfusepath{stroke,fill}%
}%
\begin{pgfscope}%
\pgfsys@transformshift{0.650000in}{1.151312in}%
\pgfsys@useobject{currentmarker}{}%
\end{pgfscope}%
\end{pgfscope}%
\begin{pgfscope}%
\pgfsetbuttcap%
\pgfsetroundjoin%
\definecolor{currentfill}{rgb}{0.000000,0.000000,0.000000}%
\pgfsetfillcolor{currentfill}%
\pgfsetlinewidth{0.401500pt}%
\definecolor{currentstroke}{rgb}{0.000000,0.000000,0.000000}%
\pgfsetstrokecolor{currentstroke}%
\pgfsetdash{}{0pt}%
\pgfsys@defobject{currentmarker}{\pgfqpoint{-0.020833in}{0.000000in}}{\pgfqpoint{-0.000000in}{0.000000in}}{%
\pgfpathmoveto{\pgfqpoint{-0.000000in}{0.000000in}}%
\pgfpathlineto{\pgfqpoint{-0.020833in}{0.000000in}}%
\pgfusepath{stroke,fill}%
}%
\begin{pgfscope}%
\pgfsys@transformshift{3.038110in}{1.151312in}%
\pgfsys@useobject{currentmarker}{}%
\end{pgfscope}%
\end{pgfscope}%
\begin{pgfscope}%
\pgfsetbuttcap%
\pgfsetroundjoin%
\definecolor{currentfill}{rgb}{0.000000,0.000000,0.000000}%
\pgfsetfillcolor{currentfill}%
\pgfsetlinewidth{0.401500pt}%
\definecolor{currentstroke}{rgb}{0.000000,0.000000,0.000000}%
\pgfsetstrokecolor{currentstroke}%
\pgfsetdash{}{0pt}%
\pgfsys@defobject{currentmarker}{\pgfqpoint{0.000000in}{0.000000in}}{\pgfqpoint{0.020833in}{0.000000in}}{%
\pgfpathmoveto{\pgfqpoint{0.000000in}{0.000000in}}%
\pgfpathlineto{\pgfqpoint{0.020833in}{0.000000in}}%
\pgfusepath{stroke,fill}%
}%
\begin{pgfscope}%
\pgfsys@transformshift{0.650000in}{1.217795in}%
\pgfsys@useobject{currentmarker}{}%
\end{pgfscope}%
\end{pgfscope}%
\begin{pgfscope}%
\pgfsetbuttcap%
\pgfsetroundjoin%
\definecolor{currentfill}{rgb}{0.000000,0.000000,0.000000}%
\pgfsetfillcolor{currentfill}%
\pgfsetlinewidth{0.401500pt}%
\definecolor{currentstroke}{rgb}{0.000000,0.000000,0.000000}%
\pgfsetstrokecolor{currentstroke}%
\pgfsetdash{}{0pt}%
\pgfsys@defobject{currentmarker}{\pgfqpoint{-0.020833in}{0.000000in}}{\pgfqpoint{-0.000000in}{0.000000in}}{%
\pgfpathmoveto{\pgfqpoint{-0.000000in}{0.000000in}}%
\pgfpathlineto{\pgfqpoint{-0.020833in}{0.000000in}}%
\pgfusepath{stroke,fill}%
}%
\begin{pgfscope}%
\pgfsys@transformshift{3.038110in}{1.217795in}%
\pgfsys@useobject{currentmarker}{}%
\end{pgfscope}%
\end{pgfscope}%
\begin{pgfscope}%
\pgfsetbuttcap%
\pgfsetroundjoin%
\definecolor{currentfill}{rgb}{0.000000,0.000000,0.000000}%
\pgfsetfillcolor{currentfill}%
\pgfsetlinewidth{0.401500pt}%
\definecolor{currentstroke}{rgb}{0.000000,0.000000,0.000000}%
\pgfsetstrokecolor{currentstroke}%
\pgfsetdash{}{0pt}%
\pgfsys@defobject{currentmarker}{\pgfqpoint{0.000000in}{0.000000in}}{\pgfqpoint{0.020833in}{0.000000in}}{%
\pgfpathmoveto{\pgfqpoint{0.000000in}{0.000000in}}%
\pgfpathlineto{\pgfqpoint{0.020833in}{0.000000in}}%
\pgfusepath{stroke,fill}%
}%
\begin{pgfscope}%
\pgfsys@transformshift{0.650000in}{1.284278in}%
\pgfsys@useobject{currentmarker}{}%
\end{pgfscope}%
\end{pgfscope}%
\begin{pgfscope}%
\pgfsetbuttcap%
\pgfsetroundjoin%
\definecolor{currentfill}{rgb}{0.000000,0.000000,0.000000}%
\pgfsetfillcolor{currentfill}%
\pgfsetlinewidth{0.401500pt}%
\definecolor{currentstroke}{rgb}{0.000000,0.000000,0.000000}%
\pgfsetstrokecolor{currentstroke}%
\pgfsetdash{}{0pt}%
\pgfsys@defobject{currentmarker}{\pgfqpoint{-0.020833in}{0.000000in}}{\pgfqpoint{-0.000000in}{0.000000in}}{%
\pgfpathmoveto{\pgfqpoint{-0.000000in}{0.000000in}}%
\pgfpathlineto{\pgfqpoint{-0.020833in}{0.000000in}}%
\pgfusepath{stroke,fill}%
}%
\begin{pgfscope}%
\pgfsys@transformshift{3.038110in}{1.284278in}%
\pgfsys@useobject{currentmarker}{}%
\end{pgfscope}%
\end{pgfscope}%
\begin{pgfscope}%
\pgfsetbuttcap%
\pgfsetroundjoin%
\definecolor{currentfill}{rgb}{0.000000,0.000000,0.000000}%
\pgfsetfillcolor{currentfill}%
\pgfsetlinewidth{0.401500pt}%
\definecolor{currentstroke}{rgb}{0.000000,0.000000,0.000000}%
\pgfsetstrokecolor{currentstroke}%
\pgfsetdash{}{0pt}%
\pgfsys@defobject{currentmarker}{\pgfqpoint{0.000000in}{0.000000in}}{\pgfqpoint{0.020833in}{0.000000in}}{%
\pgfpathmoveto{\pgfqpoint{0.000000in}{0.000000in}}%
\pgfpathlineto{\pgfqpoint{0.020833in}{0.000000in}}%
\pgfusepath{stroke,fill}%
}%
\begin{pgfscope}%
\pgfsys@transformshift{0.650000in}{1.350761in}%
\pgfsys@useobject{currentmarker}{}%
\end{pgfscope}%
\end{pgfscope}%
\begin{pgfscope}%
\pgfsetbuttcap%
\pgfsetroundjoin%
\definecolor{currentfill}{rgb}{0.000000,0.000000,0.000000}%
\pgfsetfillcolor{currentfill}%
\pgfsetlinewidth{0.401500pt}%
\definecolor{currentstroke}{rgb}{0.000000,0.000000,0.000000}%
\pgfsetstrokecolor{currentstroke}%
\pgfsetdash{}{0pt}%
\pgfsys@defobject{currentmarker}{\pgfqpoint{-0.020833in}{0.000000in}}{\pgfqpoint{-0.000000in}{0.000000in}}{%
\pgfpathmoveto{\pgfqpoint{-0.000000in}{0.000000in}}%
\pgfpathlineto{\pgfqpoint{-0.020833in}{0.000000in}}%
\pgfusepath{stroke,fill}%
}%
\begin{pgfscope}%
\pgfsys@transformshift{3.038110in}{1.350761in}%
\pgfsys@useobject{currentmarker}{}%
\end{pgfscope}%
\end{pgfscope}%
\begin{pgfscope}%
\pgfsetbuttcap%
\pgfsetroundjoin%
\definecolor{currentfill}{rgb}{0.000000,0.000000,0.000000}%
\pgfsetfillcolor{currentfill}%
\pgfsetlinewidth{0.401500pt}%
\definecolor{currentstroke}{rgb}{0.000000,0.000000,0.000000}%
\pgfsetstrokecolor{currentstroke}%
\pgfsetdash{}{0pt}%
\pgfsys@defobject{currentmarker}{\pgfqpoint{0.000000in}{0.000000in}}{\pgfqpoint{0.020833in}{0.000000in}}{%
\pgfpathmoveto{\pgfqpoint{0.000000in}{0.000000in}}%
\pgfpathlineto{\pgfqpoint{0.020833in}{0.000000in}}%
\pgfusepath{stroke,fill}%
}%
\begin{pgfscope}%
\pgfsys@transformshift{0.650000in}{1.483727in}%
\pgfsys@useobject{currentmarker}{}%
\end{pgfscope}%
\end{pgfscope}%
\begin{pgfscope}%
\pgfsetbuttcap%
\pgfsetroundjoin%
\definecolor{currentfill}{rgb}{0.000000,0.000000,0.000000}%
\pgfsetfillcolor{currentfill}%
\pgfsetlinewidth{0.401500pt}%
\definecolor{currentstroke}{rgb}{0.000000,0.000000,0.000000}%
\pgfsetstrokecolor{currentstroke}%
\pgfsetdash{}{0pt}%
\pgfsys@defobject{currentmarker}{\pgfqpoint{-0.020833in}{0.000000in}}{\pgfqpoint{-0.000000in}{0.000000in}}{%
\pgfpathmoveto{\pgfqpoint{-0.000000in}{0.000000in}}%
\pgfpathlineto{\pgfqpoint{-0.020833in}{0.000000in}}%
\pgfusepath{stroke,fill}%
}%
\begin{pgfscope}%
\pgfsys@transformshift{3.038110in}{1.483727in}%
\pgfsys@useobject{currentmarker}{}%
\end{pgfscope}%
\end{pgfscope}%
\begin{pgfscope}%
\pgfsetbuttcap%
\pgfsetroundjoin%
\definecolor{currentfill}{rgb}{0.000000,0.000000,0.000000}%
\pgfsetfillcolor{currentfill}%
\pgfsetlinewidth{0.401500pt}%
\definecolor{currentstroke}{rgb}{0.000000,0.000000,0.000000}%
\pgfsetstrokecolor{currentstroke}%
\pgfsetdash{}{0pt}%
\pgfsys@defobject{currentmarker}{\pgfqpoint{0.000000in}{0.000000in}}{\pgfqpoint{0.020833in}{0.000000in}}{%
\pgfpathmoveto{\pgfqpoint{0.000000in}{0.000000in}}%
\pgfpathlineto{\pgfqpoint{0.020833in}{0.000000in}}%
\pgfusepath{stroke,fill}%
}%
\begin{pgfscope}%
\pgfsys@transformshift{0.650000in}{1.550210in}%
\pgfsys@useobject{currentmarker}{}%
\end{pgfscope}%
\end{pgfscope}%
\begin{pgfscope}%
\pgfsetbuttcap%
\pgfsetroundjoin%
\definecolor{currentfill}{rgb}{0.000000,0.000000,0.000000}%
\pgfsetfillcolor{currentfill}%
\pgfsetlinewidth{0.401500pt}%
\definecolor{currentstroke}{rgb}{0.000000,0.000000,0.000000}%
\pgfsetstrokecolor{currentstroke}%
\pgfsetdash{}{0pt}%
\pgfsys@defobject{currentmarker}{\pgfqpoint{-0.020833in}{0.000000in}}{\pgfqpoint{-0.000000in}{0.000000in}}{%
\pgfpathmoveto{\pgfqpoint{-0.000000in}{0.000000in}}%
\pgfpathlineto{\pgfqpoint{-0.020833in}{0.000000in}}%
\pgfusepath{stroke,fill}%
}%
\begin{pgfscope}%
\pgfsys@transformshift{3.038110in}{1.550210in}%
\pgfsys@useobject{currentmarker}{}%
\end{pgfscope}%
\end{pgfscope}%
\begin{pgfscope}%
\pgfsetbuttcap%
\pgfsetroundjoin%
\definecolor{currentfill}{rgb}{0.000000,0.000000,0.000000}%
\pgfsetfillcolor{currentfill}%
\pgfsetlinewidth{0.401500pt}%
\definecolor{currentstroke}{rgb}{0.000000,0.000000,0.000000}%
\pgfsetstrokecolor{currentstroke}%
\pgfsetdash{}{0pt}%
\pgfsys@defobject{currentmarker}{\pgfqpoint{0.000000in}{0.000000in}}{\pgfqpoint{0.020833in}{0.000000in}}{%
\pgfpathmoveto{\pgfqpoint{0.000000in}{0.000000in}}%
\pgfpathlineto{\pgfqpoint{0.020833in}{0.000000in}}%
\pgfusepath{stroke,fill}%
}%
\begin{pgfscope}%
\pgfsys@transformshift{0.650000in}{1.616693in}%
\pgfsys@useobject{currentmarker}{}%
\end{pgfscope}%
\end{pgfscope}%
\begin{pgfscope}%
\pgfsetbuttcap%
\pgfsetroundjoin%
\definecolor{currentfill}{rgb}{0.000000,0.000000,0.000000}%
\pgfsetfillcolor{currentfill}%
\pgfsetlinewidth{0.401500pt}%
\definecolor{currentstroke}{rgb}{0.000000,0.000000,0.000000}%
\pgfsetstrokecolor{currentstroke}%
\pgfsetdash{}{0pt}%
\pgfsys@defobject{currentmarker}{\pgfqpoint{-0.020833in}{0.000000in}}{\pgfqpoint{-0.000000in}{0.000000in}}{%
\pgfpathmoveto{\pgfqpoint{-0.000000in}{0.000000in}}%
\pgfpathlineto{\pgfqpoint{-0.020833in}{0.000000in}}%
\pgfusepath{stroke,fill}%
}%
\begin{pgfscope}%
\pgfsys@transformshift{3.038110in}{1.616693in}%
\pgfsys@useobject{currentmarker}{}%
\end{pgfscope}%
\end{pgfscope}%
\begin{pgfscope}%
\pgfsetbuttcap%
\pgfsetroundjoin%
\definecolor{currentfill}{rgb}{0.000000,0.000000,0.000000}%
\pgfsetfillcolor{currentfill}%
\pgfsetlinewidth{0.401500pt}%
\definecolor{currentstroke}{rgb}{0.000000,0.000000,0.000000}%
\pgfsetstrokecolor{currentstroke}%
\pgfsetdash{}{0pt}%
\pgfsys@defobject{currentmarker}{\pgfqpoint{0.000000in}{0.000000in}}{\pgfqpoint{0.020833in}{0.000000in}}{%
\pgfpathmoveto{\pgfqpoint{0.000000in}{0.000000in}}%
\pgfpathlineto{\pgfqpoint{0.020833in}{0.000000in}}%
\pgfusepath{stroke,fill}%
}%
\begin{pgfscope}%
\pgfsys@transformshift{0.650000in}{1.683176in}%
\pgfsys@useobject{currentmarker}{}%
\end{pgfscope}%
\end{pgfscope}%
\begin{pgfscope}%
\pgfsetbuttcap%
\pgfsetroundjoin%
\definecolor{currentfill}{rgb}{0.000000,0.000000,0.000000}%
\pgfsetfillcolor{currentfill}%
\pgfsetlinewidth{0.401500pt}%
\definecolor{currentstroke}{rgb}{0.000000,0.000000,0.000000}%
\pgfsetstrokecolor{currentstroke}%
\pgfsetdash{}{0pt}%
\pgfsys@defobject{currentmarker}{\pgfqpoint{-0.020833in}{0.000000in}}{\pgfqpoint{-0.000000in}{0.000000in}}{%
\pgfpathmoveto{\pgfqpoint{-0.000000in}{0.000000in}}%
\pgfpathlineto{\pgfqpoint{-0.020833in}{0.000000in}}%
\pgfusepath{stroke,fill}%
}%
\begin{pgfscope}%
\pgfsys@transformshift{3.038110in}{1.683176in}%
\pgfsys@useobject{currentmarker}{}%
\end{pgfscope}%
\end{pgfscope}%
\begin{pgfscope}%
\pgfsetbuttcap%
\pgfsetroundjoin%
\definecolor{currentfill}{rgb}{0.000000,0.000000,0.000000}%
\pgfsetfillcolor{currentfill}%
\pgfsetlinewidth{0.401500pt}%
\definecolor{currentstroke}{rgb}{0.000000,0.000000,0.000000}%
\pgfsetstrokecolor{currentstroke}%
\pgfsetdash{}{0pt}%
\pgfsys@defobject{currentmarker}{\pgfqpoint{0.000000in}{0.000000in}}{\pgfqpoint{0.020833in}{0.000000in}}{%
\pgfpathmoveto{\pgfqpoint{0.000000in}{0.000000in}}%
\pgfpathlineto{\pgfqpoint{0.020833in}{0.000000in}}%
\pgfusepath{stroke,fill}%
}%
\begin{pgfscope}%
\pgfsys@transformshift{0.650000in}{1.816142in}%
\pgfsys@useobject{currentmarker}{}%
\end{pgfscope}%
\end{pgfscope}%
\begin{pgfscope}%
\pgfsetbuttcap%
\pgfsetroundjoin%
\definecolor{currentfill}{rgb}{0.000000,0.000000,0.000000}%
\pgfsetfillcolor{currentfill}%
\pgfsetlinewidth{0.401500pt}%
\definecolor{currentstroke}{rgb}{0.000000,0.000000,0.000000}%
\pgfsetstrokecolor{currentstroke}%
\pgfsetdash{}{0pt}%
\pgfsys@defobject{currentmarker}{\pgfqpoint{-0.020833in}{0.000000in}}{\pgfqpoint{-0.000000in}{0.000000in}}{%
\pgfpathmoveto{\pgfqpoint{-0.000000in}{0.000000in}}%
\pgfpathlineto{\pgfqpoint{-0.020833in}{0.000000in}}%
\pgfusepath{stroke,fill}%
}%
\begin{pgfscope}%
\pgfsys@transformshift{3.038110in}{1.816142in}%
\pgfsys@useobject{currentmarker}{}%
\end{pgfscope}%
\end{pgfscope}%
\begin{pgfscope}%
\pgfsetbuttcap%
\pgfsetroundjoin%
\definecolor{currentfill}{rgb}{0.000000,0.000000,0.000000}%
\pgfsetfillcolor{currentfill}%
\pgfsetlinewidth{0.401500pt}%
\definecolor{currentstroke}{rgb}{0.000000,0.000000,0.000000}%
\pgfsetstrokecolor{currentstroke}%
\pgfsetdash{}{0pt}%
\pgfsys@defobject{currentmarker}{\pgfqpoint{0.000000in}{0.000000in}}{\pgfqpoint{0.020833in}{0.000000in}}{%
\pgfpathmoveto{\pgfqpoint{0.000000in}{0.000000in}}%
\pgfpathlineto{\pgfqpoint{0.020833in}{0.000000in}}%
\pgfusepath{stroke,fill}%
}%
\begin{pgfscope}%
\pgfsys@transformshift{0.650000in}{1.882625in}%
\pgfsys@useobject{currentmarker}{}%
\end{pgfscope}%
\end{pgfscope}%
\begin{pgfscope}%
\pgfsetbuttcap%
\pgfsetroundjoin%
\definecolor{currentfill}{rgb}{0.000000,0.000000,0.000000}%
\pgfsetfillcolor{currentfill}%
\pgfsetlinewidth{0.401500pt}%
\definecolor{currentstroke}{rgb}{0.000000,0.000000,0.000000}%
\pgfsetstrokecolor{currentstroke}%
\pgfsetdash{}{0pt}%
\pgfsys@defobject{currentmarker}{\pgfqpoint{-0.020833in}{0.000000in}}{\pgfqpoint{-0.000000in}{0.000000in}}{%
\pgfpathmoveto{\pgfqpoint{-0.000000in}{0.000000in}}%
\pgfpathlineto{\pgfqpoint{-0.020833in}{0.000000in}}%
\pgfusepath{stroke,fill}%
}%
\begin{pgfscope}%
\pgfsys@transformshift{3.038110in}{1.882625in}%
\pgfsys@useobject{currentmarker}{}%
\end{pgfscope}%
\end{pgfscope}%
\begin{pgfscope}%
\pgfsetbuttcap%
\pgfsetroundjoin%
\definecolor{currentfill}{rgb}{0.000000,0.000000,0.000000}%
\pgfsetfillcolor{currentfill}%
\pgfsetlinewidth{0.401500pt}%
\definecolor{currentstroke}{rgb}{0.000000,0.000000,0.000000}%
\pgfsetstrokecolor{currentstroke}%
\pgfsetdash{}{0pt}%
\pgfsys@defobject{currentmarker}{\pgfqpoint{0.000000in}{0.000000in}}{\pgfqpoint{0.020833in}{0.000000in}}{%
\pgfpathmoveto{\pgfqpoint{0.000000in}{0.000000in}}%
\pgfpathlineto{\pgfqpoint{0.020833in}{0.000000in}}%
\pgfusepath{stroke,fill}%
}%
\begin{pgfscope}%
\pgfsys@transformshift{0.650000in}{1.949108in}%
\pgfsys@useobject{currentmarker}{}%
\end{pgfscope}%
\end{pgfscope}%
\begin{pgfscope}%
\pgfsetbuttcap%
\pgfsetroundjoin%
\definecolor{currentfill}{rgb}{0.000000,0.000000,0.000000}%
\pgfsetfillcolor{currentfill}%
\pgfsetlinewidth{0.401500pt}%
\definecolor{currentstroke}{rgb}{0.000000,0.000000,0.000000}%
\pgfsetstrokecolor{currentstroke}%
\pgfsetdash{}{0pt}%
\pgfsys@defobject{currentmarker}{\pgfqpoint{-0.020833in}{0.000000in}}{\pgfqpoint{-0.000000in}{0.000000in}}{%
\pgfpathmoveto{\pgfqpoint{-0.000000in}{0.000000in}}%
\pgfpathlineto{\pgfqpoint{-0.020833in}{0.000000in}}%
\pgfusepath{stroke,fill}%
}%
\begin{pgfscope}%
\pgfsys@transformshift{3.038110in}{1.949108in}%
\pgfsys@useobject{currentmarker}{}%
\end{pgfscope}%
\end{pgfscope}%
\begin{pgfscope}%
\pgfsetbuttcap%
\pgfsetroundjoin%
\definecolor{currentfill}{rgb}{0.000000,0.000000,0.000000}%
\pgfsetfillcolor{currentfill}%
\pgfsetlinewidth{0.401500pt}%
\definecolor{currentstroke}{rgb}{0.000000,0.000000,0.000000}%
\pgfsetstrokecolor{currentstroke}%
\pgfsetdash{}{0pt}%
\pgfsys@defobject{currentmarker}{\pgfqpoint{0.000000in}{0.000000in}}{\pgfqpoint{0.020833in}{0.000000in}}{%
\pgfpathmoveto{\pgfqpoint{0.000000in}{0.000000in}}%
\pgfpathlineto{\pgfqpoint{0.020833in}{0.000000in}}%
\pgfusepath{stroke,fill}%
}%
\begin{pgfscope}%
\pgfsys@transformshift{0.650000in}{2.015590in}%
\pgfsys@useobject{currentmarker}{}%
\end{pgfscope}%
\end{pgfscope}%
\begin{pgfscope}%
\pgfsetbuttcap%
\pgfsetroundjoin%
\definecolor{currentfill}{rgb}{0.000000,0.000000,0.000000}%
\pgfsetfillcolor{currentfill}%
\pgfsetlinewidth{0.401500pt}%
\definecolor{currentstroke}{rgb}{0.000000,0.000000,0.000000}%
\pgfsetstrokecolor{currentstroke}%
\pgfsetdash{}{0pt}%
\pgfsys@defobject{currentmarker}{\pgfqpoint{-0.020833in}{0.000000in}}{\pgfqpoint{-0.000000in}{0.000000in}}{%
\pgfpathmoveto{\pgfqpoint{-0.000000in}{0.000000in}}%
\pgfpathlineto{\pgfqpoint{-0.020833in}{0.000000in}}%
\pgfusepath{stroke,fill}%
}%
\begin{pgfscope}%
\pgfsys@transformshift{3.038110in}{2.015590in}%
\pgfsys@useobject{currentmarker}{}%
\end{pgfscope}%
\end{pgfscope}%
\begin{pgfscope}%
\pgfsetbuttcap%
\pgfsetroundjoin%
\definecolor{currentfill}{rgb}{0.000000,0.000000,0.000000}%
\pgfsetfillcolor{currentfill}%
\pgfsetlinewidth{0.401500pt}%
\definecolor{currentstroke}{rgb}{0.000000,0.000000,0.000000}%
\pgfsetstrokecolor{currentstroke}%
\pgfsetdash{}{0pt}%
\pgfsys@defobject{currentmarker}{\pgfqpoint{0.000000in}{0.000000in}}{\pgfqpoint{0.020833in}{0.000000in}}{%
\pgfpathmoveto{\pgfqpoint{0.000000in}{0.000000in}}%
\pgfpathlineto{\pgfqpoint{0.020833in}{0.000000in}}%
\pgfusepath{stroke,fill}%
}%
\begin{pgfscope}%
\pgfsys@transformshift{0.650000in}{2.148556in}%
\pgfsys@useobject{currentmarker}{}%
\end{pgfscope}%
\end{pgfscope}%
\begin{pgfscope}%
\pgfsetbuttcap%
\pgfsetroundjoin%
\definecolor{currentfill}{rgb}{0.000000,0.000000,0.000000}%
\pgfsetfillcolor{currentfill}%
\pgfsetlinewidth{0.401500pt}%
\definecolor{currentstroke}{rgb}{0.000000,0.000000,0.000000}%
\pgfsetstrokecolor{currentstroke}%
\pgfsetdash{}{0pt}%
\pgfsys@defobject{currentmarker}{\pgfqpoint{-0.020833in}{0.000000in}}{\pgfqpoint{-0.000000in}{0.000000in}}{%
\pgfpathmoveto{\pgfqpoint{-0.000000in}{0.000000in}}%
\pgfpathlineto{\pgfqpoint{-0.020833in}{0.000000in}}%
\pgfusepath{stroke,fill}%
}%
\begin{pgfscope}%
\pgfsys@transformshift{3.038110in}{2.148556in}%
\pgfsys@useobject{currentmarker}{}%
\end{pgfscope}%
\end{pgfscope}%
\begin{pgfscope}%
\pgfsetbuttcap%
\pgfsetroundjoin%
\definecolor{currentfill}{rgb}{0.000000,0.000000,0.000000}%
\pgfsetfillcolor{currentfill}%
\pgfsetlinewidth{0.401500pt}%
\definecolor{currentstroke}{rgb}{0.000000,0.000000,0.000000}%
\pgfsetstrokecolor{currentstroke}%
\pgfsetdash{}{0pt}%
\pgfsys@defobject{currentmarker}{\pgfqpoint{0.000000in}{0.000000in}}{\pgfqpoint{0.020833in}{0.000000in}}{%
\pgfpathmoveto{\pgfqpoint{0.000000in}{0.000000in}}%
\pgfpathlineto{\pgfqpoint{0.020833in}{0.000000in}}%
\pgfusepath{stroke,fill}%
}%
\begin{pgfscope}%
\pgfsys@transformshift{0.650000in}{2.215039in}%
\pgfsys@useobject{currentmarker}{}%
\end{pgfscope}%
\end{pgfscope}%
\begin{pgfscope}%
\pgfsetbuttcap%
\pgfsetroundjoin%
\definecolor{currentfill}{rgb}{0.000000,0.000000,0.000000}%
\pgfsetfillcolor{currentfill}%
\pgfsetlinewidth{0.401500pt}%
\definecolor{currentstroke}{rgb}{0.000000,0.000000,0.000000}%
\pgfsetstrokecolor{currentstroke}%
\pgfsetdash{}{0pt}%
\pgfsys@defobject{currentmarker}{\pgfqpoint{-0.020833in}{0.000000in}}{\pgfqpoint{-0.000000in}{0.000000in}}{%
\pgfpathmoveto{\pgfqpoint{-0.000000in}{0.000000in}}%
\pgfpathlineto{\pgfqpoint{-0.020833in}{0.000000in}}%
\pgfusepath{stroke,fill}%
}%
\begin{pgfscope}%
\pgfsys@transformshift{3.038110in}{2.215039in}%
\pgfsys@useobject{currentmarker}{}%
\end{pgfscope}%
\end{pgfscope}%
\begin{pgfscope}%
\pgfsetbuttcap%
\pgfsetroundjoin%
\definecolor{currentfill}{rgb}{0.000000,0.000000,0.000000}%
\pgfsetfillcolor{currentfill}%
\pgfsetlinewidth{0.401500pt}%
\definecolor{currentstroke}{rgb}{0.000000,0.000000,0.000000}%
\pgfsetstrokecolor{currentstroke}%
\pgfsetdash{}{0pt}%
\pgfsys@defobject{currentmarker}{\pgfqpoint{0.000000in}{0.000000in}}{\pgfqpoint{0.020833in}{0.000000in}}{%
\pgfpathmoveto{\pgfqpoint{0.000000in}{0.000000in}}%
\pgfpathlineto{\pgfqpoint{0.020833in}{0.000000in}}%
\pgfusepath{stroke,fill}%
}%
\begin{pgfscope}%
\pgfsys@transformshift{0.650000in}{2.281522in}%
\pgfsys@useobject{currentmarker}{}%
\end{pgfscope}%
\end{pgfscope}%
\begin{pgfscope}%
\pgfsetbuttcap%
\pgfsetroundjoin%
\definecolor{currentfill}{rgb}{0.000000,0.000000,0.000000}%
\pgfsetfillcolor{currentfill}%
\pgfsetlinewidth{0.401500pt}%
\definecolor{currentstroke}{rgb}{0.000000,0.000000,0.000000}%
\pgfsetstrokecolor{currentstroke}%
\pgfsetdash{}{0pt}%
\pgfsys@defobject{currentmarker}{\pgfqpoint{-0.020833in}{0.000000in}}{\pgfqpoint{-0.000000in}{0.000000in}}{%
\pgfpathmoveto{\pgfqpoint{-0.000000in}{0.000000in}}%
\pgfpathlineto{\pgfqpoint{-0.020833in}{0.000000in}}%
\pgfusepath{stroke,fill}%
}%
\begin{pgfscope}%
\pgfsys@transformshift{3.038110in}{2.281522in}%
\pgfsys@useobject{currentmarker}{}%
\end{pgfscope}%
\end{pgfscope}%
\begin{pgfscope}%
\pgfsetbuttcap%
\pgfsetroundjoin%
\definecolor{currentfill}{rgb}{0.000000,0.000000,0.000000}%
\pgfsetfillcolor{currentfill}%
\pgfsetlinewidth{0.401500pt}%
\definecolor{currentstroke}{rgb}{0.000000,0.000000,0.000000}%
\pgfsetstrokecolor{currentstroke}%
\pgfsetdash{}{0pt}%
\pgfsys@defobject{currentmarker}{\pgfqpoint{0.000000in}{0.000000in}}{\pgfqpoint{0.020833in}{0.000000in}}{%
\pgfpathmoveto{\pgfqpoint{0.000000in}{0.000000in}}%
\pgfpathlineto{\pgfqpoint{0.020833in}{0.000000in}}%
\pgfusepath{stroke,fill}%
}%
\begin{pgfscope}%
\pgfsys@transformshift{0.650000in}{2.348005in}%
\pgfsys@useobject{currentmarker}{}%
\end{pgfscope}%
\end{pgfscope}%
\begin{pgfscope}%
\pgfsetbuttcap%
\pgfsetroundjoin%
\definecolor{currentfill}{rgb}{0.000000,0.000000,0.000000}%
\pgfsetfillcolor{currentfill}%
\pgfsetlinewidth{0.401500pt}%
\definecolor{currentstroke}{rgb}{0.000000,0.000000,0.000000}%
\pgfsetstrokecolor{currentstroke}%
\pgfsetdash{}{0pt}%
\pgfsys@defobject{currentmarker}{\pgfqpoint{-0.020833in}{0.000000in}}{\pgfqpoint{-0.000000in}{0.000000in}}{%
\pgfpathmoveto{\pgfqpoint{-0.000000in}{0.000000in}}%
\pgfpathlineto{\pgfqpoint{-0.020833in}{0.000000in}}%
\pgfusepath{stroke,fill}%
}%
\begin{pgfscope}%
\pgfsys@transformshift{3.038110in}{2.348005in}%
\pgfsys@useobject{currentmarker}{}%
\end{pgfscope}%
\end{pgfscope}%
\begin{pgfscope}%
\definecolor{textcolor}{rgb}{0.000000,0.000000,0.000000}%
\pgfsetstrokecolor{textcolor}%
\pgfsetfillcolor{textcolor}%
\pgftext[x=0.227199in,y=1.417244in,,bottom,rotate=90.000000]{\color{textcolor}{\rmfamily\fontsize{12.000000}{14.400000}\selectfont\catcode`\^=\active\def^{\ifmmode\sp\else\^{}\fi}\catcode`\%=\active\def%{\%}$C_f$}}%
\end{pgfscope}%
\begin{pgfscope}%
\pgfpathrectangle{\pgfqpoint{0.650000in}{0.420000in}}{\pgfqpoint{2.388110in}{1.994488in}}%
\pgfusepath{clip}%
\pgfsetrectcap%
\pgfsetroundjoin%
\pgfsetlinewidth{1.003750pt}%
\definecolor{currentstroke}{rgb}{0.796078,0.094118,0.113725}%
\pgfsetstrokecolor{currentstroke}%
\pgfsetdash{}{0pt}%
\pgfpathmoveto{\pgfqpoint{0.650000in}{1.603969in}}%
\pgfpathlineto{\pgfqpoint{0.651792in}{1.604664in}}%
\pgfpathlineto{\pgfqpoint{0.653584in}{1.606751in}}%
\pgfpathlineto{\pgfqpoint{0.655735in}{1.611072in}}%
\pgfpathlineto{\pgfqpoint{0.658243in}{1.618642in}}%
\pgfpathlineto{\pgfqpoint{0.661469in}{1.632432in}}%
\pgfpathlineto{\pgfqpoint{0.665053in}{1.652924in}}%
\pgfpathlineto{\pgfqpoint{0.670071in}{1.688864in}}%
\pgfpathlineto{\pgfqpoint{0.677956in}{1.745147in}}%
\pgfpathlineto{\pgfqpoint{0.681182in}{1.769765in}}%
\pgfpathlineto{\pgfqpoint{0.684766in}{1.807150in}}%
\pgfpathlineto{\pgfqpoint{0.707346in}{2.062278in}}%
\pgfpathlineto{\pgfqpoint{0.710572in}{2.085961in}}%
\pgfpathlineto{\pgfqpoint{0.714873in}{2.107133in}}%
\pgfpathlineto{\pgfqpoint{0.727417in}{2.166266in}}%
\pgfpathlineto{\pgfqpoint{0.730284in}{2.173114in}}%
\pgfpathlineto{\pgfqpoint{0.732793in}{2.176428in}}%
\pgfpathlineto{\pgfqpoint{0.734944in}{2.177431in}}%
\pgfpathlineto{\pgfqpoint{0.736736in}{2.177097in}}%
\pgfpathlineto{\pgfqpoint{0.738886in}{2.175436in}}%
\pgfpathlineto{\pgfqpoint{0.741395in}{2.171972in}}%
\pgfpathlineto{\pgfqpoint{0.744621in}{2.165528in}}%
\pgfpathlineto{\pgfqpoint{0.749280in}{2.153500in}}%
\pgfpathlineto{\pgfqpoint{0.756807in}{2.130461in}}%
\pgfpathlineto{\pgfqpoint{0.771143in}{2.086298in}}%
\pgfpathlineto{\pgfqpoint{0.779745in}{2.063518in}}%
\pgfpathlineto{\pgfqpoint{0.788347in}{2.043944in}}%
\pgfpathlineto{\pgfqpoint{0.797666in}{2.025463in}}%
\pgfpathlineto{\pgfqpoint{0.809852in}{2.004096in}}%
\pgfpathlineto{\pgfqpoint{0.824905in}{1.980253in}}%
\pgfpathlineto{\pgfqpoint{0.837091in}{1.962637in}}%
\pgfpathlineto{\pgfqpoint{0.846051in}{1.951575in}}%
\pgfpathlineto{\pgfqpoint{0.858954in}{1.937896in}}%
\pgfpathlineto{\pgfqpoint{0.882251in}{1.914870in}}%
\pgfpathlineto{\pgfqpoint{0.894437in}{1.904370in}}%
\pgfpathlineto{\pgfqpoint{0.905906in}{1.895944in}}%
\pgfpathlineto{\pgfqpoint{0.924543in}{1.884184in}}%
\pgfpathlineto{\pgfqpoint{0.941030in}{1.873960in}}%
\pgfpathlineto{\pgfqpoint{0.980097in}{1.848595in}}%
\pgfpathlineto{\pgfqpoint{1.037443in}{1.818088in}}%
\pgfpathlineto{\pgfqpoint{1.047479in}{1.814421in}}%
\pgfpathlineto{\pgfqpoint{1.055722in}{1.812509in}}%
\pgfpathlineto{\pgfqpoint{1.073643in}{1.809849in}}%
\pgfpathlineto{\pgfqpoint{1.079736in}{1.810079in}}%
\pgfpathlineto{\pgfqpoint{1.084395in}{1.811264in}}%
\pgfpathlineto{\pgfqpoint{1.088696in}{1.813515in}}%
\pgfpathlineto{\pgfqpoint{1.094431in}{1.818015in}}%
\pgfpathlineto{\pgfqpoint{1.100882in}{1.822665in}}%
\pgfpathlineto{\pgfqpoint{1.106258in}{1.825144in}}%
\pgfpathlineto{\pgfqpoint{1.112351in}{1.826671in}}%
\pgfpathlineto{\pgfqpoint{1.120595in}{1.827496in}}%
\pgfpathlineto{\pgfqpoint{1.128838in}{1.827280in}}%
\pgfpathlineto{\pgfqpoint{1.150343in}{1.824837in}}%
\pgfpathlineto{\pgfqpoint{1.162887in}{1.822832in}}%
\pgfpathlineto{\pgfqpoint{1.168980in}{1.820867in}}%
\pgfpathlineto{\pgfqpoint{1.173998in}{1.818001in}}%
\pgfpathlineto{\pgfqpoint{1.178299in}{1.814274in}}%
\pgfpathlineto{\pgfqpoint{1.184392in}{1.807171in}}%
\pgfpathlineto{\pgfqpoint{1.191561in}{1.799337in}}%
\pgfpathlineto{\pgfqpoint{1.198012in}{1.793845in}}%
\pgfpathlineto{\pgfqpoint{1.204463in}{1.789802in}}%
\pgfpathlineto{\pgfqpoint{1.216650in}{1.784016in}}%
\pgfpathlineto{\pgfqpoint{1.242097in}{1.771575in}}%
\pgfpathlineto{\pgfqpoint{1.255717in}{1.763556in}}%
\pgfpathlineto{\pgfqpoint{1.268978in}{1.754185in}}%
\pgfpathlineto{\pgfqpoint{1.337077in}{1.705236in}}%
\pgfpathlineto{\pgfqpoint{1.345679in}{1.698675in}}%
\pgfpathlineto{\pgfqpoint{1.362524in}{1.688143in}}%
\pgfpathlineto{\pgfqpoint{1.373994in}{1.682832in}}%
\pgfpathlineto{\pgfqpoint{1.390481in}{1.676527in}}%
\pgfpathlineto{\pgfqpoint{1.402308in}{1.671973in}}%
\pgfpathlineto{\pgfqpoint{1.419871in}{1.664317in}}%
\pgfpathlineto{\pgfqpoint{1.426322in}{1.662707in}}%
\pgfpathlineto{\pgfqpoint{1.433849in}{1.660679in}}%
\pgfpathlineto{\pgfqpoint{1.448544in}{1.656431in}}%
\pgfpathlineto{\pgfqpoint{1.454279in}{1.654885in}}%
\pgfpathlineto{\pgfqpoint{1.491554in}{1.640564in}}%
\pgfpathlineto{\pgfqpoint{1.499439in}{1.637947in}}%
\pgfpathlineto{\pgfqpoint{1.507324in}{1.635195in}}%
\pgfpathlineto{\pgfqpoint{1.524169in}{1.628346in}}%
\pgfpathlineto{\pgfqpoint{1.529187in}{1.626524in}}%
\pgfpathlineto{\pgfqpoint{1.547108in}{1.623484in}}%
\pgfpathlineto{\pgfqpoint{1.555351in}{1.622128in}}%
\pgfpathlineto{\pgfqpoint{1.570405in}{1.616732in}}%
\pgfpathlineto{\pgfqpoint{1.583307in}{1.611932in}}%
\pgfpathlineto{\pgfqpoint{1.607680in}{1.606435in}}%
\pgfpathlineto{\pgfqpoint{1.616282in}{1.603817in}}%
\pgfpathlineto{\pgfqpoint{1.624525in}{1.601197in}}%
\pgfpathlineto{\pgfqpoint{1.628826in}{1.600280in}}%
\pgfpathlineto{\pgfqpoint{1.637428in}{1.599261in}}%
\pgfpathlineto{\pgfqpoint{1.653198in}{1.595228in}}%
\pgfpathlineto{\pgfqpoint{1.660367in}{1.593566in}}%
\pgfpathlineto{\pgfqpoint{1.666818in}{1.592306in}}%
\pgfpathlineto{\pgfqpoint{1.672911in}{1.591170in}}%
\pgfpathlineto{\pgfqpoint{1.679363in}{1.588358in}}%
\pgfpathlineto{\pgfqpoint{1.684739in}{1.585102in}}%
\pgfpathlineto{\pgfqpoint{1.692266in}{1.580239in}}%
\pgfpathlineto{\pgfqpoint{1.696925in}{1.578571in}}%
\pgfpathlineto{\pgfqpoint{1.707319in}{1.576812in}}%
\pgfpathlineto{\pgfqpoint{1.713412in}{1.575324in}}%
\pgfpathlineto{\pgfqpoint{1.717713in}{1.573187in}}%
\pgfpathlineto{\pgfqpoint{1.730616in}{1.565309in}}%
\pgfpathlineto{\pgfqpoint{1.735634in}{1.563937in}}%
\pgfpathlineto{\pgfqpoint{1.754988in}{1.560533in}}%
\pgfpathlineto{\pgfqpoint{1.771475in}{1.554583in}}%
\pgfpathlineto{\pgfqpoint{1.777568in}{1.554805in}}%
\pgfpathlineto{\pgfqpoint{1.786171in}{1.554943in}}%
\pgfpathlineto{\pgfqpoint{1.794414in}{1.552555in}}%
\pgfpathlineto{\pgfqpoint{1.801224in}{1.549177in}}%
\pgfpathlineto{\pgfqpoint{1.805525in}{1.547758in}}%
\pgfpathlineto{\pgfqpoint{1.819861in}{1.544195in}}%
\pgfpathlineto{\pgfqpoint{1.860720in}{1.532830in}}%
\pgfpathlineto{\pgfqpoint{1.879358in}{1.530026in}}%
\pgfpathlineto{\pgfqpoint{1.890110in}{1.528485in}}%
\pgfpathlineto{\pgfqpoint{1.903730in}{1.526094in}}%
\pgfpathlineto{\pgfqpoint{1.910181in}{1.526078in}}%
\pgfpathlineto{\pgfqpoint{1.916991in}{1.524119in}}%
\pgfpathlineto{\pgfqpoint{1.921650in}{1.522326in}}%
\pgfpathlineto{\pgfqpoint{1.933120in}{1.516931in}}%
\pgfpathlineto{\pgfqpoint{1.937779in}{1.516145in}}%
\pgfpathlineto{\pgfqpoint{1.954266in}{1.514366in}}%
\pgfpathlineto{\pgfqpoint{1.961792in}{1.512337in}}%
\pgfpathlineto{\pgfqpoint{1.973620in}{1.511075in}}%
\pgfpathlineto{\pgfqpoint{1.980429in}{1.511865in}}%
\pgfpathlineto{\pgfqpoint{1.991540in}{1.514564in}}%
\pgfpathlineto{\pgfqpoint{1.996916in}{1.514002in}}%
\pgfpathlineto{\pgfqpoint{2.005877in}{1.512753in}}%
\pgfpathlineto{\pgfqpoint{2.011611in}{1.512834in}}%
\pgfpathlineto{\pgfqpoint{2.018779in}{1.514244in}}%
\pgfpathlineto{\pgfqpoint{2.029173in}{1.517262in}}%
\pgfpathlineto{\pgfqpoint{2.032758in}{1.517268in}}%
\pgfpathlineto{\pgfqpoint{2.036700in}{1.516090in}}%
\pgfpathlineto{\pgfqpoint{2.041718in}{1.514652in}}%
\pgfpathlineto{\pgfqpoint{2.045660in}{1.514859in}}%
\pgfpathlineto{\pgfqpoint{2.049961in}{1.514798in}}%
\pgfpathlineto{\pgfqpoint{2.075050in}{1.508731in}}%
\pgfpathlineto{\pgfqpoint{2.080426in}{1.507140in}}%
\pgfpathlineto{\pgfqpoint{2.088311in}{1.506505in}}%
\pgfpathlineto{\pgfqpoint{2.095838in}{1.505005in}}%
\pgfpathlineto{\pgfqpoint{2.104798in}{1.503380in}}%
\pgfpathlineto{\pgfqpoint{2.113758in}{1.502565in}}%
\pgfpathlineto{\pgfqpoint{2.120568in}{1.500810in}}%
\pgfpathlineto{\pgfqpoint{2.126661in}{1.499564in}}%
\pgfpathlineto{\pgfqpoint{2.140997in}{1.499293in}}%
\pgfpathlineto{\pgfqpoint{2.146373in}{1.499739in}}%
\pgfpathlineto{\pgfqpoint{2.151749in}{1.498694in}}%
\pgfpathlineto{\pgfqpoint{2.167161in}{1.495083in}}%
\pgfpathlineto{\pgfqpoint{2.172179in}{1.493680in}}%
\pgfpathlineto{\pgfqpoint{2.176121in}{1.492718in}}%
\pgfpathlineto{\pgfqpoint{2.179347in}{1.493253in}}%
\pgfpathlineto{\pgfqpoint{2.184364in}{1.494072in}}%
\pgfpathlineto{\pgfqpoint{2.189023in}{1.493570in}}%
\pgfpathlineto{\pgfqpoint{2.195116in}{1.491784in}}%
\pgfpathlineto{\pgfqpoint{2.206584in}{1.486823in}}%
\pgfpathlineto{\pgfqpoint{2.210885in}{1.485414in}}%
\pgfpathlineto{\pgfqpoint{2.214111in}{1.485560in}}%
\pgfpathlineto{\pgfqpoint{2.221278in}{1.486542in}}%
\pgfpathlineto{\pgfqpoint{2.227371in}{1.485378in}}%
\pgfpathlineto{\pgfqpoint{2.238481in}{1.482423in}}%
\pgfpathlineto{\pgfqpoint{2.243498in}{1.481234in}}%
\pgfpathlineto{\pgfqpoint{2.256400in}{1.477983in}}%
\pgfpathlineto{\pgfqpoint{2.277903in}{1.473446in}}%
\pgfpathlineto{\pgfqpoint{2.286146in}{1.473237in}}%
\pgfpathlineto{\pgfqpoint{2.295464in}{1.473002in}}%
\pgfpathlineto{\pgfqpoint{2.309083in}{1.475105in}}%
\pgfpathlineto{\pgfqpoint{2.316609in}{1.474455in}}%
\pgfpathlineto{\pgfqpoint{2.338471in}{1.470942in}}%
\pgfpathlineto{\pgfqpoint{2.343130in}{1.471746in}}%
\pgfpathlineto{\pgfqpoint{2.351014in}{1.473267in}}%
\pgfpathlineto{\pgfqpoint{2.354957in}{1.472915in}}%
\pgfpathlineto{\pgfqpoint{2.369292in}{1.470160in}}%
\pgfpathlineto{\pgfqpoint{2.386853in}{1.469649in}}%
\pgfpathlineto{\pgfqpoint{2.391870in}{1.467971in}}%
\pgfpathlineto{\pgfqpoint{2.401189in}{1.464219in}}%
\pgfpathlineto{\pgfqpoint{2.408715in}{1.463632in}}%
\pgfpathlineto{\pgfqpoint{2.423409in}{1.464090in}}%
\pgfpathlineto{\pgfqpoint{2.428426in}{1.464885in}}%
\pgfpathlineto{\pgfqpoint{2.432010in}{1.464129in}}%
\pgfpathlineto{\pgfqpoint{2.437027in}{1.461776in}}%
\pgfpathlineto{\pgfqpoint{2.440611in}{1.460647in}}%
\pgfpathlineto{\pgfqpoint{2.443837in}{1.460896in}}%
\pgfpathlineto{\pgfqpoint{2.456022in}{1.463065in}}%
\pgfpathlineto{\pgfqpoint{2.463190in}{1.460684in}}%
\pgfpathlineto{\pgfqpoint{2.472508in}{1.457908in}}%
\pgfpathlineto{\pgfqpoint{2.482902in}{1.455132in}}%
\pgfpathlineto{\pgfqpoint{2.493295in}{1.451523in}}%
\pgfpathlineto{\pgfqpoint{2.499387in}{1.450154in}}%
\pgfpathlineto{\pgfqpoint{2.504047in}{1.448771in}}%
\pgfpathlineto{\pgfqpoint{2.508348in}{1.446276in}}%
\pgfpathlineto{\pgfqpoint{2.513365in}{1.441974in}}%
\pgfpathlineto{\pgfqpoint{2.546338in}{1.410987in}}%
\pgfpathlineto{\pgfqpoint{2.552072in}{1.408179in}}%
\pgfpathlineto{\pgfqpoint{2.556731in}{1.407132in}}%
\pgfpathlineto{\pgfqpoint{2.566408in}{1.405042in}}%
\pgfpathlineto{\pgfqpoint{2.575728in}{1.403417in}}%
\pgfpathlineto{\pgfqpoint{2.582896in}{1.402878in}}%
\pgfpathlineto{\pgfqpoint{2.587914in}{1.400857in}}%
\pgfpathlineto{\pgfqpoint{2.592933in}{1.399377in}}%
\pgfpathlineto{\pgfqpoint{2.596517in}{1.399524in}}%
\pgfpathlineto{\pgfqpoint{2.600818in}{1.400849in}}%
\pgfpathlineto{\pgfqpoint{2.606911in}{1.403967in}}%
\pgfpathlineto{\pgfqpoint{2.611571in}{1.407641in}}%
\pgfpathlineto{\pgfqpoint{2.619457in}{1.414330in}}%
\pgfpathlineto{\pgfqpoint{2.622683in}{1.415247in}}%
\pgfpathlineto{\pgfqpoint{2.634869in}{1.417066in}}%
\pgfpathlineto{\pgfqpoint{2.646697in}{1.421964in}}%
\pgfpathlineto{\pgfqpoint{2.656734in}{1.424728in}}%
\pgfpathlineto{\pgfqpoint{2.667128in}{1.427112in}}%
\pgfpathlineto{\pgfqpoint{2.671071in}{1.427852in}}%
\pgfpathlineto{\pgfqpoint{2.674297in}{1.429763in}}%
\pgfpathlineto{\pgfqpoint{2.682183in}{1.435036in}}%
\pgfpathlineto{\pgfqpoint{2.687201in}{1.437108in}}%
\pgfpathlineto{\pgfqpoint{2.696878in}{1.440098in}}%
\pgfpathlineto{\pgfqpoint{2.701896in}{1.441725in}}%
\pgfpathlineto{\pgfqpoint{2.707273in}{1.443380in}}%
\pgfpathlineto{\pgfqpoint{2.710140in}{1.445711in}}%
\pgfpathlineto{\pgfqpoint{2.720176in}{1.455648in}}%
\pgfpathlineto{\pgfqpoint{2.739532in}{1.470301in}}%
\pgfpathlineto{\pgfqpoint{2.745267in}{1.473506in}}%
\pgfpathlineto{\pgfqpoint{2.749568in}{1.477663in}}%
\pgfpathlineto{\pgfqpoint{2.760321in}{1.488949in}}%
\pgfpathlineto{\pgfqpoint{2.765339in}{1.490743in}}%
\pgfpathlineto{\pgfqpoint{2.771432in}{1.493320in}}%
\pgfpathlineto{\pgfqpoint{2.775017in}{1.496268in}}%
\pgfpathlineto{\pgfqpoint{2.784336in}{1.505174in}}%
\pgfpathlineto{\pgfqpoint{2.787562in}{1.506042in}}%
\pgfpathlineto{\pgfqpoint{2.793297in}{1.507365in}}%
\pgfpathlineto{\pgfqpoint{2.803333in}{1.510933in}}%
\pgfpathlineto{\pgfqpoint{2.809785in}{1.511381in}}%
\pgfpathlineto{\pgfqpoint{2.809785in}{1.511381in}}%
\pgfusepath{stroke}%
\end{pgfscope}%
\begin{pgfscope}%
\pgfpathrectangle{\pgfqpoint{0.650000in}{0.420000in}}{\pgfqpoint{2.388110in}{1.994488in}}%
\pgfusepath{clip}%
\pgfsetrectcap%
\pgfsetroundjoin%
\pgfsetlinewidth{1.003750pt}%
\definecolor{currentstroke}{rgb}{0.454902,0.678431,0.819608}%
\pgfsetstrokecolor{currentstroke}%
\pgfsetdash{}{0pt}%
\pgfpathmoveto{\pgfqpoint{0.650000in}{1.593310in}}%
\pgfpathlineto{\pgfqpoint{0.651792in}{1.593983in}}%
\pgfpathlineto{\pgfqpoint{0.653584in}{1.596010in}}%
\pgfpathlineto{\pgfqpoint{0.655735in}{1.600245in}}%
\pgfpathlineto{\pgfqpoint{0.658243in}{1.607718in}}%
\pgfpathlineto{\pgfqpoint{0.661469in}{1.621373in}}%
\pgfpathlineto{\pgfqpoint{0.665412in}{1.643875in}}%
\pgfpathlineto{\pgfqpoint{0.670788in}{1.682405in}}%
\pgfpathlineto{\pgfqpoint{0.677598in}{1.729811in}}%
\pgfpathlineto{\pgfqpoint{0.681182in}{1.756283in}}%
\pgfpathlineto{\pgfqpoint{0.684766in}{1.793279in}}%
\pgfpathlineto{\pgfqpoint{0.708063in}{2.054060in}}%
\pgfpathlineto{\pgfqpoint{0.710930in}{2.073683in}}%
\pgfpathlineto{\pgfqpoint{0.722041in}{2.131042in}}%
\pgfpathlineto{\pgfqpoint{0.725983in}{2.147952in}}%
\pgfpathlineto{\pgfqpoint{0.729209in}{2.157429in}}%
\pgfpathlineto{\pgfqpoint{0.731718in}{2.161922in}}%
\pgfpathlineto{\pgfqpoint{0.733868in}{2.163797in}}%
\pgfpathlineto{\pgfqpoint{0.735660in}{2.164070in}}%
\pgfpathlineto{\pgfqpoint{0.737452in}{2.163301in}}%
\pgfpathlineto{\pgfqpoint{0.739603in}{2.161198in}}%
\pgfpathlineto{\pgfqpoint{0.742470in}{2.156779in}}%
\pgfpathlineto{\pgfqpoint{0.746413in}{2.148532in}}%
\pgfpathlineto{\pgfqpoint{0.751789in}{2.134604in}}%
\pgfpathlineto{\pgfqpoint{0.763975in}{2.098616in}}%
\pgfpathlineto{\pgfqpoint{0.776878in}{2.062678in}}%
\pgfpathlineto{\pgfqpoint{0.785121in}{2.042740in}}%
\pgfpathlineto{\pgfqpoint{0.794440in}{2.023341in}}%
\pgfpathlineto{\pgfqpoint{0.806268in}{2.001322in}}%
\pgfpathlineto{\pgfqpoint{0.815228in}{1.986766in}}%
\pgfpathlineto{\pgfqpoint{0.825980in}{1.971301in}}%
\pgfpathlineto{\pgfqpoint{0.853578in}{1.935895in}}%
\pgfpathlineto{\pgfqpoint{0.861821in}{1.927587in}}%
\pgfpathlineto{\pgfqpoint{0.905906in}{1.886975in}}%
\pgfpathlineto{\pgfqpoint{0.925977in}{1.870567in}}%
\pgfpathlineto{\pgfqpoint{0.934579in}{1.865228in}}%
\pgfpathlineto{\pgfqpoint{0.958951in}{1.851747in}}%
\pgfpathlineto{\pgfqpoint{0.967553in}{1.848495in}}%
\pgfpathlineto{\pgfqpoint{0.973287in}{1.847476in}}%
\pgfpathlineto{\pgfqpoint{0.983323in}{1.847011in}}%
\pgfpathlineto{\pgfqpoint{0.988699in}{1.847984in}}%
\pgfpathlineto{\pgfqpoint{0.995509in}{1.850405in}}%
\pgfpathlineto{\pgfqpoint{1.002319in}{1.853980in}}%
\pgfpathlineto{\pgfqpoint{1.008053in}{1.858379in}}%
\pgfpathlineto{\pgfqpoint{1.023823in}{1.872856in}}%
\pgfpathlineto{\pgfqpoint{1.033501in}{1.880807in}}%
\pgfpathlineto{\pgfqpoint{1.047120in}{1.890327in}}%
\pgfpathlineto{\pgfqpoint{1.053572in}{1.893522in}}%
\pgfpathlineto{\pgfqpoint{1.058948in}{1.894948in}}%
\pgfpathlineto{\pgfqpoint{1.063966in}{1.895157in}}%
\pgfpathlineto{\pgfqpoint{1.070417in}{1.894204in}}%
\pgfpathlineto{\pgfqpoint{1.079736in}{1.891454in}}%
\pgfpathlineto{\pgfqpoint{1.088696in}{1.889319in}}%
\pgfpathlineto{\pgfqpoint{1.095864in}{1.887674in}}%
\pgfpathlineto{\pgfqpoint{1.099807in}{1.885380in}}%
\pgfpathlineto{\pgfqpoint{1.105183in}{1.880726in}}%
\pgfpathlineto{\pgfqpoint{1.111634in}{1.873649in}}%
\pgfpathlineto{\pgfqpoint{1.132422in}{1.849278in}}%
\pgfpathlineto{\pgfqpoint{1.141024in}{1.840807in}}%
\pgfpathlineto{\pgfqpoint{1.147834in}{1.835586in}}%
\pgfpathlineto{\pgfqpoint{1.155719in}{1.831070in}}%
\pgfpathlineto{\pgfqpoint{1.162887in}{1.828109in}}%
\pgfpathlineto{\pgfqpoint{1.169697in}{1.825162in}}%
\pgfpathlineto{\pgfqpoint{1.175073in}{1.821512in}}%
\pgfpathlineto{\pgfqpoint{1.180091in}{1.816719in}}%
\pgfpathlineto{\pgfqpoint{1.188335in}{1.808578in}}%
\pgfpathlineto{\pgfqpoint{1.195862in}{1.803350in}}%
\pgfpathlineto{\pgfqpoint{1.208406in}{1.795892in}}%
\pgfpathlineto{\pgfqpoint{1.215216in}{1.791413in}}%
\pgfpathlineto{\pgfqpoint{1.225610in}{1.782809in}}%
\pgfpathlineto{\pgfqpoint{1.243531in}{1.766673in}}%
\pgfpathlineto{\pgfqpoint{1.255000in}{1.755090in}}%
\pgfpathlineto{\pgfqpoint{1.294784in}{1.713169in}}%
\pgfpathlineto{\pgfqpoint{1.305178in}{1.704917in}}%
\pgfpathlineto{\pgfqpoint{1.310196in}{1.702470in}}%
\pgfpathlineto{\pgfqpoint{1.324174in}{1.697696in}}%
\pgfpathlineto{\pgfqpoint{1.334568in}{1.695099in}}%
\pgfpathlineto{\pgfqpoint{1.342453in}{1.694027in}}%
\pgfpathlineto{\pgfqpoint{1.350697in}{1.694197in}}%
\pgfpathlineto{\pgfqpoint{1.368259in}{1.695835in}}%
\pgfpathlineto{\pgfqpoint{1.378295in}{1.696820in}}%
\pgfpathlineto{\pgfqpoint{1.389764in}{1.696555in}}%
\pgfpathlineto{\pgfqpoint{1.396932in}{1.695393in}}%
\pgfpathlineto{\pgfqpoint{1.402308in}{1.693253in}}%
\pgfpathlineto{\pgfqpoint{1.442809in}{1.672753in}}%
\pgfpathlineto{\pgfqpoint{1.482593in}{1.650031in}}%
\pgfpathlineto{\pgfqpoint{1.489762in}{1.648173in}}%
\pgfpathlineto{\pgfqpoint{1.497288in}{1.647200in}}%
\pgfpathlineto{\pgfqpoint{1.519510in}{1.646279in}}%
\pgfpathlineto{\pgfqpoint{1.530262in}{1.645650in}}%
\pgfpathlineto{\pgfqpoint{1.539581in}{1.646511in}}%
\pgfpathlineto{\pgfqpoint{1.549617in}{1.646434in}}%
\pgfpathlineto{\pgfqpoint{1.556068in}{1.646397in}}%
\pgfpathlineto{\pgfqpoint{1.575781in}{1.643770in}}%
\pgfpathlineto{\pgfqpoint{1.585099in}{1.640940in}}%
\pgfpathlineto{\pgfqpoint{1.593343in}{1.638771in}}%
\pgfpathlineto{\pgfqpoint{1.608038in}{1.634003in}}%
\pgfpathlineto{\pgfqpoint{1.615923in}{1.631354in}}%
\pgfpathlineto{\pgfqpoint{1.630618in}{1.627961in}}%
\pgfpathlineto{\pgfqpoint{1.636353in}{1.627490in}}%
\pgfpathlineto{\pgfqpoint{1.644955in}{1.626390in}}%
\pgfpathlineto{\pgfqpoint{1.649973in}{1.626879in}}%
\pgfpathlineto{\pgfqpoint{1.676495in}{1.632274in}}%
\pgfpathlineto{\pgfqpoint{1.682588in}{1.631147in}}%
\pgfpathlineto{\pgfqpoint{1.698000in}{1.626677in}}%
\pgfpathlineto{\pgfqpoint{1.713412in}{1.625406in}}%
\pgfpathlineto{\pgfqpoint{1.717713in}{1.623660in}}%
\pgfpathlineto{\pgfqpoint{1.731691in}{1.616302in}}%
\pgfpathlineto{\pgfqpoint{1.738143in}{1.614493in}}%
\pgfpathlineto{\pgfqpoint{1.743877in}{1.612278in}}%
\pgfpathlineto{\pgfqpoint{1.751045in}{1.609853in}}%
\pgfpathlineto{\pgfqpoint{1.756063in}{1.607769in}}%
\pgfpathlineto{\pgfqpoint{1.763590in}{1.604480in}}%
\pgfpathlineto{\pgfqpoint{1.770400in}{1.602491in}}%
\pgfpathlineto{\pgfqpoint{1.775059in}{1.602551in}}%
\pgfpathlineto{\pgfqpoint{1.787963in}{1.602970in}}%
\pgfpathlineto{\pgfqpoint{1.805883in}{1.601457in}}%
\pgfpathlineto{\pgfqpoint{1.819861in}{1.603022in}}%
\pgfpathlineto{\pgfqpoint{1.862513in}{1.605354in}}%
\pgfpathlineto{\pgfqpoint{1.872906in}{1.606842in}}%
\pgfpathlineto{\pgfqpoint{1.882583in}{1.606298in}}%
\pgfpathlineto{\pgfqpoint{1.889035in}{1.604190in}}%
\pgfpathlineto{\pgfqpoint{1.904088in}{1.598526in}}%
\pgfpathlineto{\pgfqpoint{1.909823in}{1.597467in}}%
\pgfpathlineto{\pgfqpoint{1.914841in}{1.595008in}}%
\pgfpathlineto{\pgfqpoint{1.921650in}{1.590897in}}%
\pgfpathlineto{\pgfqpoint{1.929894in}{1.583397in}}%
\pgfpathlineto{\pgfqpoint{1.935628in}{1.578636in}}%
\pgfpathlineto{\pgfqpoint{1.941363in}{1.575761in}}%
\pgfpathlineto{\pgfqpoint{1.952474in}{1.570278in}}%
\pgfpathlineto{\pgfqpoint{1.964660in}{1.561723in}}%
\pgfpathlineto{\pgfqpoint{1.971828in}{1.558214in}}%
\pgfpathlineto{\pgfqpoint{1.977921in}{1.556741in}}%
\pgfpathlineto{\pgfqpoint{1.981863in}{1.556653in}}%
\pgfpathlineto{\pgfqpoint{1.990823in}{1.557201in}}%
\pgfpathlineto{\pgfqpoint{2.001217in}{1.557155in}}%
\pgfpathlineto{\pgfqpoint{2.007310in}{1.558782in}}%
\pgfpathlineto{\pgfqpoint{2.012686in}{1.561192in}}%
\pgfpathlineto{\pgfqpoint{2.018779in}{1.565025in}}%
\pgfpathlineto{\pgfqpoint{2.028456in}{1.572085in}}%
\pgfpathlineto{\pgfqpoint{2.032399in}{1.573487in}}%
\pgfpathlineto{\pgfqpoint{2.035983in}{1.573465in}}%
\pgfpathlineto{\pgfqpoint{2.041359in}{1.573184in}}%
\pgfpathlineto{\pgfqpoint{2.045302in}{1.574376in}}%
\pgfpathlineto{\pgfqpoint{2.049961in}{1.575574in}}%
\pgfpathlineto{\pgfqpoint{2.054979in}{1.575220in}}%
\pgfpathlineto{\pgfqpoint{2.071824in}{1.573554in}}%
\pgfpathlineto{\pgfqpoint{2.084368in}{1.569269in}}%
\pgfpathlineto{\pgfqpoint{2.089028in}{1.567570in}}%
\pgfpathlineto{\pgfqpoint{2.093329in}{1.564706in}}%
\pgfpathlineto{\pgfqpoint{2.106948in}{1.554491in}}%
\pgfpathlineto{\pgfqpoint{2.115909in}{1.548665in}}%
\pgfpathlineto{\pgfqpoint{2.128453in}{1.540070in}}%
\pgfpathlineto{\pgfqpoint{2.158918in}{1.527425in}}%
\pgfpathlineto{\pgfqpoint{2.165011in}{1.526350in}}%
\pgfpathlineto{\pgfqpoint{2.172179in}{1.525282in}}%
\pgfpathlineto{\pgfqpoint{2.176121in}{1.524612in}}%
\pgfpathlineto{\pgfqpoint{2.179705in}{1.525359in}}%
\pgfpathlineto{\pgfqpoint{2.185439in}{1.526619in}}%
\pgfpathlineto{\pgfqpoint{2.190098in}{1.526402in}}%
\pgfpathlineto{\pgfqpoint{2.194757in}{1.525059in}}%
\pgfpathlineto{\pgfqpoint{2.209810in}{1.519830in}}%
\pgfpathlineto{\pgfqpoint{2.213035in}{1.520187in}}%
\pgfpathlineto{\pgfqpoint{2.221995in}{1.522441in}}%
\pgfpathlineto{\pgfqpoint{2.230596in}{1.521400in}}%
\pgfpathlineto{\pgfqpoint{2.234897in}{1.521165in}}%
\pgfpathlineto{\pgfqpoint{2.242065in}{1.521577in}}%
\pgfpathlineto{\pgfqpoint{2.256400in}{1.519962in}}%
\pgfpathlineto{\pgfqpoint{2.262851in}{1.519957in}}%
\pgfpathlineto{\pgfqpoint{2.266793in}{1.519495in}}%
\pgfpathlineto{\pgfqpoint{2.273603in}{1.518263in}}%
\pgfpathlineto{\pgfqpoint{2.290447in}{1.518420in}}%
\pgfpathlineto{\pgfqpoint{2.295464in}{1.518125in}}%
\pgfpathlineto{\pgfqpoint{2.312309in}{1.519374in}}%
\pgfpathlineto{\pgfqpoint{2.316968in}{1.518322in}}%
\pgfpathlineto{\pgfqpoint{2.323060in}{1.515514in}}%
\pgfpathlineto{\pgfqpoint{2.328078in}{1.513684in}}%
\pgfpathlineto{\pgfqpoint{2.339904in}{1.511505in}}%
\pgfpathlineto{\pgfqpoint{2.351373in}{1.510422in}}%
\pgfpathlineto{\pgfqpoint{2.355673in}{1.508202in}}%
\pgfpathlineto{\pgfqpoint{2.365350in}{1.502585in}}%
\pgfpathlineto{\pgfqpoint{2.370367in}{1.501061in}}%
\pgfpathlineto{\pgfqpoint{2.376101in}{1.500611in}}%
\pgfpathlineto{\pgfqpoint{2.382194in}{1.500247in}}%
\pgfpathlineto{\pgfqpoint{2.386495in}{1.498791in}}%
\pgfpathlineto{\pgfqpoint{2.390795in}{1.496205in}}%
\pgfpathlineto{\pgfqpoint{2.403339in}{1.487029in}}%
\pgfpathlineto{\pgfqpoint{2.431293in}{1.474324in}}%
\pgfpathlineto{\pgfqpoint{2.435952in}{1.469400in}}%
\pgfpathlineto{\pgfqpoint{2.440970in}{1.464305in}}%
\pgfpathlineto{\pgfqpoint{2.445629in}{1.461435in}}%
\pgfpathlineto{\pgfqpoint{2.451721in}{1.457437in}}%
\pgfpathlineto{\pgfqpoint{2.454947in}{1.454063in}}%
\pgfpathlineto{\pgfqpoint{2.459965in}{1.446763in}}%
\pgfpathlineto{\pgfqpoint{2.482543in}{1.411514in}}%
\pgfpathlineto{\pgfqpoint{2.488277in}{1.400783in}}%
\pgfpathlineto{\pgfqpoint{2.495087in}{1.385249in}}%
\pgfpathlineto{\pgfqpoint{2.504047in}{1.362565in}}%
\pgfpathlineto{\pgfqpoint{2.510856in}{1.341538in}}%
\pgfpathlineto{\pgfqpoint{2.521250in}{1.304798in}}%
\pgfpathlineto{\pgfqpoint{2.550997in}{1.199076in}}%
\pgfpathlineto{\pgfqpoint{2.564616in}{1.158972in}}%
\pgfpathlineto{\pgfqpoint{2.575011in}{1.131238in}}%
\pgfpathlineto{\pgfqpoint{2.581821in}{1.116686in}}%
\pgfpathlineto{\pgfqpoint{2.589707in}{1.102187in}}%
\pgfpathlineto{\pgfqpoint{2.594366in}{1.095852in}}%
\pgfpathlineto{\pgfqpoint{2.599026in}{1.091425in}}%
\pgfpathlineto{\pgfqpoint{2.604761in}{1.087538in}}%
\pgfpathlineto{\pgfqpoint{2.608704in}{1.086045in}}%
\pgfpathlineto{\pgfqpoint{2.612288in}{1.085825in}}%
\pgfpathlineto{\pgfqpoint{2.626267in}{1.086479in}}%
\pgfpathlineto{\pgfqpoint{2.632360in}{1.086818in}}%
\pgfpathlineto{\pgfqpoint{2.637378in}{1.088172in}}%
\pgfpathlineto{\pgfqpoint{2.671788in}{1.101233in}}%
\pgfpathlineto{\pgfqpoint{2.688634in}{1.110461in}}%
\pgfpathlineto{\pgfqpoint{2.695803in}{1.113703in}}%
\pgfpathlineto{\pgfqpoint{2.701896in}{1.116339in}}%
\pgfpathlineto{\pgfqpoint{2.709423in}{1.119120in}}%
\pgfpathlineto{\pgfqpoint{2.722327in}{1.126617in}}%
\pgfpathlineto{\pgfqpoint{2.739173in}{1.133384in}}%
\pgfpathlineto{\pgfqpoint{2.745625in}{1.135481in}}%
\pgfpathlineto{\pgfqpoint{2.751002in}{1.138616in}}%
\pgfpathlineto{\pgfqpoint{2.759604in}{1.143737in}}%
\pgfpathlineto{\pgfqpoint{2.776451in}{1.149793in}}%
\pgfpathlineto{\pgfqpoint{2.782902in}{1.153065in}}%
\pgfpathlineto{\pgfqpoint{2.787203in}{1.153737in}}%
\pgfpathlineto{\pgfqpoint{2.796523in}{1.154672in}}%
\pgfpathlineto{\pgfqpoint{2.801899in}{1.155042in}}%
\pgfpathlineto{\pgfqpoint{2.809785in}{1.154863in}}%
\pgfpathlineto{\pgfqpoint{2.809785in}{1.154863in}}%
\pgfusepath{stroke}%
\end{pgfscope}%
\begin{pgfscope}%
\pgfpathrectangle{\pgfqpoint{0.650000in}{0.420000in}}{\pgfqpoint{2.388110in}{1.994488in}}%
\pgfusepath{clip}%
\pgfsetrectcap%
\pgfsetroundjoin%
\pgfsetlinewidth{1.003750pt}%
\definecolor{currentstroke}{rgb}{0.878431,0.509804,0.078431}%
\pgfsetstrokecolor{currentstroke}%
\pgfsetdash{}{0pt}%
\pgfpathmoveto{\pgfqpoint{0.650000in}{1.632281in}}%
\pgfpathlineto{\pgfqpoint{0.650925in}{1.633631in}}%
\pgfpathlineto{\pgfqpoint{0.653332in}{1.638620in}}%
\pgfpathlineto{\pgfqpoint{0.656176in}{1.647296in}}%
\pgfpathlineto{\pgfqpoint{0.659488in}{1.661856in}}%
\pgfpathlineto{\pgfqpoint{0.663218in}{1.684029in}}%
\pgfpathlineto{\pgfqpoint{0.668191in}{1.721115in}}%
\pgfpathlineto{\pgfqpoint{0.675235in}{1.772688in}}%
\pgfpathlineto{\pgfqpoint{0.678551in}{1.797634in}}%
\pgfpathlineto{\pgfqpoint{0.681866in}{1.832089in}}%
\pgfpathlineto{\pgfqpoint{0.705486in}{2.100872in}}%
\pgfpathlineto{\pgfqpoint{0.708387in}{2.120413in}}%
\pgfpathlineto{\pgfqpoint{0.720819in}{2.183073in}}%
\pgfpathlineto{\pgfqpoint{0.724134in}{2.194808in}}%
\pgfpathlineto{\pgfqpoint{0.727035in}{2.201118in}}%
\pgfpathlineto{\pgfqpoint{0.729107in}{2.203396in}}%
\pgfpathlineto{\pgfqpoint{0.730764in}{2.203957in}}%
\pgfpathlineto{\pgfqpoint{0.732422in}{2.203485in}}%
\pgfpathlineto{\pgfqpoint{0.734494in}{2.201588in}}%
\pgfpathlineto{\pgfqpoint{0.736980in}{2.197621in}}%
\pgfpathlineto{\pgfqpoint{0.740295in}{2.190008in}}%
\pgfpathlineto{\pgfqpoint{0.745268in}{2.175240in}}%
\pgfpathlineto{\pgfqpoint{0.753970in}{2.145076in}}%
\pgfpathlineto{\pgfqpoint{0.765574in}{2.105938in}}%
\pgfpathlineto{\pgfqpoint{0.774276in}{2.080599in}}%
\pgfpathlineto{\pgfqpoint{0.782149in}{2.060932in}}%
\pgfpathlineto{\pgfqpoint{0.790437in}{2.043046in}}%
\pgfpathlineto{\pgfqpoint{0.800797in}{2.023661in}}%
\pgfpathlineto{\pgfqpoint{0.815715in}{1.998266in}}%
\pgfpathlineto{\pgfqpoint{0.826075in}{1.982533in}}%
\pgfpathlineto{\pgfqpoint{0.850525in}{1.948598in}}%
\pgfpathlineto{\pgfqpoint{0.861299in}{1.935512in}}%
\pgfpathlineto{\pgfqpoint{0.874559in}{1.921137in}}%
\pgfpathlineto{\pgfqpoint{0.900252in}{1.896002in}}%
\pgfpathlineto{\pgfqpoint{0.916413in}{1.882100in}}%
\pgfpathlineto{\pgfqpoint{0.928016in}{1.873455in}}%
\pgfpathlineto{\pgfqpoint{0.981059in}{1.836347in}}%
\pgfpathlineto{\pgfqpoint{1.000536in}{1.824173in}}%
\pgfpathlineto{\pgfqpoint{1.027471in}{1.808537in}}%
\pgfpathlineto{\pgfqpoint{1.034930in}{1.805461in}}%
\pgfpathlineto{\pgfqpoint{1.067668in}{1.794372in}}%
\pgfpathlineto{\pgfqpoint{1.074712in}{1.793398in}}%
\pgfpathlineto{\pgfqpoint{1.083000in}{1.793277in}}%
\pgfpathlineto{\pgfqpoint{1.102477in}{1.794392in}}%
\pgfpathlineto{\pgfqpoint{1.113251in}{1.793727in}}%
\pgfpathlineto{\pgfqpoint{1.120710in}{1.792079in}}%
\pgfpathlineto{\pgfqpoint{1.135214in}{1.788443in}}%
\pgfpathlineto{\pgfqpoint{1.161321in}{1.781700in}}%
\pgfpathlineto{\pgfqpoint{1.169608in}{1.777859in}}%
\pgfpathlineto{\pgfqpoint{1.178310in}{1.772542in}}%
\pgfpathlineto{\pgfqpoint{1.186599in}{1.768036in}}%
\pgfpathlineto{\pgfqpoint{1.201517in}{1.761459in}}%
\pgfpathlineto{\pgfqpoint{1.216435in}{1.755036in}}%
\pgfpathlineto{\pgfqpoint{1.229696in}{1.748768in}}%
\pgfpathlineto{\pgfqpoint{1.254974in}{1.738018in}}%
\pgfpathlineto{\pgfqpoint{1.275694in}{1.730196in}}%
\pgfpathlineto{\pgfqpoint{1.297242in}{1.721752in}}%
\pgfpathlineto{\pgfqpoint{1.375563in}{1.691836in}}%
\pgfpathlineto{\pgfqpoint{1.382608in}{1.690795in}}%
\pgfpathlineto{\pgfqpoint{1.423219in}{1.686128in}}%
\pgfpathlineto{\pgfqpoint{1.431507in}{1.683604in}}%
\pgfpathlineto{\pgfqpoint{1.446425in}{1.679361in}}%
\pgfpathlineto{\pgfqpoint{1.457199in}{1.676761in}}%
\pgfpathlineto{\pgfqpoint{1.466730in}{1.674458in}}%
\pgfpathlineto{\pgfqpoint{1.476676in}{1.670303in}}%
\pgfpathlineto{\pgfqpoint{1.488279in}{1.665366in}}%
\pgfpathlineto{\pgfqpoint{1.513557in}{1.655156in}}%
\pgfpathlineto{\pgfqpoint{1.543394in}{1.646592in}}%
\pgfpathlineto{\pgfqpoint{1.591049in}{1.635478in}}%
\pgfpathlineto{\pgfqpoint{1.610111in}{1.633595in}}%
\pgfpathlineto{\pgfqpoint{1.624201in}{1.632934in}}%
\pgfpathlineto{\pgfqpoint{1.638705in}{1.630413in}}%
\pgfpathlineto{\pgfqpoint{1.652794in}{1.627475in}}%
\pgfpathlineto{\pgfqpoint{1.666469in}{1.626259in}}%
\pgfpathlineto{\pgfqpoint{1.685531in}{1.625529in}}%
\pgfpathlineto{\pgfqpoint{1.696720in}{1.626252in}}%
\pgfpathlineto{\pgfqpoint{1.714953in}{1.626430in}}%
\pgfpathlineto{\pgfqpoint{1.726142in}{1.627160in}}%
\pgfpathlineto{\pgfqpoint{1.741060in}{1.628772in}}%
\pgfpathlineto{\pgfqpoint{1.754321in}{1.628796in}}%
\pgfpathlineto{\pgfqpoint{1.764267in}{1.627705in}}%
\pgfpathlineto{\pgfqpoint{1.770897in}{1.625938in}}%
\pgfpathlineto{\pgfqpoint{1.785815in}{1.621796in}}%
\pgfpathlineto{\pgfqpoint{1.805706in}{1.617956in}}%
\pgfpathlineto{\pgfqpoint{1.863307in}{1.606011in}}%
\pgfpathlineto{\pgfqpoint{1.874496in}{1.605138in}}%
\pgfpathlineto{\pgfqpoint{1.882784in}{1.603134in}}%
\pgfpathlineto{\pgfqpoint{1.891072in}{1.601459in}}%
\pgfpathlineto{\pgfqpoint{1.900188in}{1.599353in}}%
\pgfpathlineto{\pgfqpoint{1.915521in}{1.595628in}}%
\pgfpathlineto{\pgfqpoint{1.929610in}{1.593975in}}%
\pgfpathlineto{\pgfqpoint{1.938727in}{1.594052in}}%
\pgfpathlineto{\pgfqpoint{1.947844in}{1.593936in}}%
\pgfpathlineto{\pgfqpoint{1.954060in}{1.592956in}}%
\pgfpathlineto{\pgfqpoint{1.980581in}{1.586851in}}%
\pgfpathlineto{\pgfqpoint{1.993842in}{1.585079in}}%
\pgfpathlineto{\pgfqpoint{2.007931in}{1.583408in}}%
\pgfpathlineto{\pgfqpoint{2.023678in}{1.579166in}}%
\pgfpathlineto{\pgfqpoint{2.032795in}{1.576788in}}%
\pgfpathlineto{\pgfqpoint{2.053515in}{1.573496in}}%
\pgfpathlineto{\pgfqpoint{2.066775in}{1.573387in}}%
\pgfpathlineto{\pgfqpoint{2.075063in}{1.573250in}}%
\pgfpathlineto{\pgfqpoint{2.081694in}{1.572016in}}%
\pgfpathlineto{\pgfqpoint{2.092882in}{1.569776in}}%
\pgfpathlineto{\pgfqpoint{2.095369in}{1.570667in}}%
\pgfpathlineto{\pgfqpoint{2.099513in}{1.573786in}}%
\pgfpathlineto{\pgfqpoint{2.104485in}{1.577208in}}%
\pgfpathlineto{\pgfqpoint{2.107801in}{1.578245in}}%
\pgfpathlineto{\pgfqpoint{2.110701in}{1.578014in}}%
\pgfpathlineto{\pgfqpoint{2.113602in}{1.576624in}}%
\pgfpathlineto{\pgfqpoint{2.118989in}{1.573547in}}%
\pgfpathlineto{\pgfqpoint{2.122719in}{1.573374in}}%
\pgfpathlineto{\pgfqpoint{2.126034in}{1.572628in}}%
\pgfpathlineto{\pgfqpoint{2.131007in}{1.571178in}}%
\pgfpathlineto{\pgfqpoint{2.136394in}{1.571363in}}%
\pgfpathlineto{\pgfqpoint{2.142610in}{1.571725in}}%
\pgfpathlineto{\pgfqpoint{2.146754in}{1.570837in}}%
\pgfpathlineto{\pgfqpoint{2.150069in}{1.568965in}}%
\pgfpathlineto{\pgfqpoint{2.158771in}{1.563403in}}%
\pgfpathlineto{\pgfqpoint{2.166645in}{1.561619in}}%
\pgfpathlineto{\pgfqpoint{2.170374in}{1.561946in}}%
\pgfpathlineto{\pgfqpoint{2.174104in}{1.562063in}}%
\pgfpathlineto{\pgfqpoint{2.177004in}{1.560692in}}%
\pgfpathlineto{\pgfqpoint{2.181977in}{1.558033in}}%
\pgfpathlineto{\pgfqpoint{2.187779in}{1.555788in}}%
\pgfpathlineto{\pgfqpoint{2.191508in}{1.552695in}}%
\pgfpathlineto{\pgfqpoint{2.197310in}{1.547647in}}%
\pgfpathlineto{\pgfqpoint{2.199796in}{1.546988in}}%
\pgfpathlineto{\pgfqpoint{2.202282in}{1.547507in}}%
\pgfpathlineto{\pgfqpoint{2.208084in}{1.549526in}}%
\pgfpathlineto{\pgfqpoint{2.210985in}{1.548646in}}%
\pgfpathlineto{\pgfqpoint{2.213886in}{1.548114in}}%
\pgfpathlineto{\pgfqpoint{2.222588in}{1.548728in}}%
\pgfpathlineto{\pgfqpoint{2.225489in}{1.547397in}}%
\pgfpathlineto{\pgfqpoint{2.230047in}{1.544992in}}%
\pgfpathlineto{\pgfqpoint{2.234191in}{1.545085in}}%
\pgfpathlineto{\pgfqpoint{2.237920in}{1.544320in}}%
\pgfpathlineto{\pgfqpoint{2.241236in}{1.544005in}}%
\pgfpathlineto{\pgfqpoint{2.244136in}{1.544932in}}%
\pgfpathlineto{\pgfqpoint{2.254496in}{1.550407in}}%
\pgfpathlineto{\pgfqpoint{2.260298in}{1.554868in}}%
\pgfpathlineto{\pgfqpoint{2.268171in}{1.559753in}}%
\pgfpathlineto{\pgfqpoint{2.269829in}{1.559199in}}%
\pgfpathlineto{\pgfqpoint{2.271901in}{1.557293in}}%
\pgfpathlineto{\pgfqpoint{2.276873in}{1.552046in}}%
\pgfpathlineto{\pgfqpoint{2.280189in}{1.550679in}}%
\pgfpathlineto{\pgfqpoint{2.286405in}{1.549213in}}%
\pgfpathlineto{\pgfqpoint{2.298837in}{1.549991in}}%
\pgfpathlineto{\pgfqpoint{2.302981in}{1.551299in}}%
\pgfpathlineto{\pgfqpoint{2.306296in}{1.550610in}}%
\pgfpathlineto{\pgfqpoint{2.312926in}{1.549576in}}%
\pgfpathlineto{\pgfqpoint{2.315827in}{1.549811in}}%
\pgfpathlineto{\pgfqpoint{2.323286in}{1.551562in}}%
\pgfpathlineto{\pgfqpoint{2.334061in}{1.551083in}}%
\pgfpathlineto{\pgfqpoint{2.339447in}{1.550378in}}%
\pgfpathlineto{\pgfqpoint{2.342763in}{1.548556in}}%
\pgfpathlineto{\pgfqpoint{2.349393in}{1.544219in}}%
\pgfpathlineto{\pgfqpoint{2.353123in}{1.543053in}}%
\pgfpathlineto{\pgfqpoint{2.356024in}{1.543252in}}%
\pgfpathlineto{\pgfqpoint{2.358510in}{1.544572in}}%
\pgfpathlineto{\pgfqpoint{2.366383in}{1.550044in}}%
\pgfpathlineto{\pgfqpoint{2.372185in}{1.549776in}}%
\pgfpathlineto{\pgfqpoint{2.375500in}{1.548589in}}%
\pgfpathlineto{\pgfqpoint{2.378401in}{1.547668in}}%
\pgfpathlineto{\pgfqpoint{2.380473in}{1.548136in}}%
\pgfpathlineto{\pgfqpoint{2.384203in}{1.549412in}}%
\pgfpathlineto{\pgfqpoint{2.385860in}{1.548708in}}%
\pgfpathlineto{\pgfqpoint{2.387932in}{1.546369in}}%
\pgfpathlineto{\pgfqpoint{2.404508in}{1.522444in}}%
\pgfpathlineto{\pgfqpoint{2.406994in}{1.521135in}}%
\pgfpathlineto{\pgfqpoint{2.409895in}{1.521227in}}%
\pgfpathlineto{\pgfqpoint{2.415282in}{1.522751in}}%
\pgfpathlineto{\pgfqpoint{2.419841in}{1.525263in}}%
\pgfpathlineto{\pgfqpoint{2.422741in}{1.526282in}}%
\pgfpathlineto{\pgfqpoint{2.425642in}{1.525877in}}%
\pgfpathlineto{\pgfqpoint{2.431858in}{1.523796in}}%
\pgfpathlineto{\pgfqpoint{2.436002in}{1.520557in}}%
\pgfpathlineto{\pgfqpoint{2.441389in}{1.514730in}}%
\pgfpathlineto{\pgfqpoint{2.445119in}{1.511511in}}%
\pgfpathlineto{\pgfqpoint{2.447605in}{1.510847in}}%
\pgfpathlineto{\pgfqpoint{2.452578in}{1.510194in}}%
\pgfpathlineto{\pgfqpoint{2.458794in}{1.508522in}}%
\pgfpathlineto{\pgfqpoint{2.461280in}{1.509196in}}%
\pgfpathlineto{\pgfqpoint{2.466253in}{1.510924in}}%
\pgfpathlineto{\pgfqpoint{2.469568in}{1.510374in}}%
\pgfpathlineto{\pgfqpoint{2.474955in}{1.509214in}}%
\pgfpathlineto{\pgfqpoint{2.477027in}{1.510125in}}%
\pgfpathlineto{\pgfqpoint{2.479513in}{1.512728in}}%
\pgfpathlineto{\pgfqpoint{2.485729in}{1.522028in}}%
\pgfpathlineto{\pgfqpoint{2.490288in}{1.527534in}}%
\pgfpathlineto{\pgfqpoint{2.492774in}{1.529159in}}%
\pgfpathlineto{\pgfqpoint{2.494846in}{1.529213in}}%
\pgfpathlineto{\pgfqpoint{2.496918in}{1.528119in}}%
\pgfpathlineto{\pgfqpoint{2.500233in}{1.524831in}}%
\pgfpathlineto{\pgfqpoint{2.503134in}{1.522510in}}%
\pgfpathlineto{\pgfqpoint{2.504792in}{1.522386in}}%
\pgfpathlineto{\pgfqpoint{2.511836in}{1.523698in}}%
\pgfpathlineto{\pgfqpoint{2.514737in}{1.523685in}}%
\pgfpathlineto{\pgfqpoint{2.516809in}{1.525028in}}%
\pgfpathlineto{\pgfqpoint{2.518881in}{1.527986in}}%
\pgfpathlineto{\pgfqpoint{2.524683in}{1.537623in}}%
\pgfpathlineto{\pgfqpoint{2.529655in}{1.541705in}}%
\pgfpathlineto{\pgfqpoint{2.535043in}{1.545996in}}%
\pgfpathlineto{\pgfqpoint{2.537943in}{1.546474in}}%
\pgfpathlineto{\pgfqpoint{2.540844in}{1.547212in}}%
\pgfpathlineto{\pgfqpoint{2.544988in}{1.548617in}}%
\pgfpathlineto{\pgfqpoint{2.549132in}{1.548034in}}%
\pgfpathlineto{\pgfqpoint{2.552862in}{1.548279in}}%
\pgfpathlineto{\pgfqpoint{2.556176in}{1.548172in}}%
\pgfpathlineto{\pgfqpoint{2.559491in}{1.546564in}}%
\pgfpathlineto{\pgfqpoint{2.562806in}{1.545288in}}%
\pgfpathlineto{\pgfqpoint{2.566121in}{1.545348in}}%
\pgfpathlineto{\pgfqpoint{2.575651in}{1.547221in}}%
\pgfpathlineto{\pgfqpoint{2.577309in}{1.548674in}}%
\pgfpathlineto{\pgfqpoint{2.579795in}{1.552760in}}%
\pgfpathlineto{\pgfqpoint{2.584767in}{1.561427in}}%
\pgfpathlineto{\pgfqpoint{2.590568in}{1.567547in}}%
\pgfpathlineto{\pgfqpoint{2.592226in}{1.568256in}}%
\pgfpathlineto{\pgfqpoint{2.593469in}{1.567845in}}%
\pgfpathlineto{\pgfqpoint{2.595126in}{1.566032in}}%
\pgfpathlineto{\pgfqpoint{2.600928in}{1.558031in}}%
\pgfpathlineto{\pgfqpoint{2.606314in}{1.554430in}}%
\pgfpathlineto{\pgfqpoint{2.608800in}{1.550022in}}%
\pgfpathlineto{\pgfqpoint{2.615016in}{1.537811in}}%
\pgfpathlineto{\pgfqpoint{2.617088in}{1.535922in}}%
\pgfpathlineto{\pgfqpoint{2.618745in}{1.535665in}}%
\pgfpathlineto{\pgfqpoint{2.620403in}{1.536619in}}%
\pgfpathlineto{\pgfqpoint{2.623303in}{1.540152in}}%
\pgfpathlineto{\pgfqpoint{2.625789in}{1.542450in}}%
\pgfpathlineto{\pgfqpoint{2.627032in}{1.542452in}}%
\pgfpathlineto{\pgfqpoint{2.628690in}{1.541089in}}%
\pgfpathlineto{\pgfqpoint{2.633662in}{1.535108in}}%
\pgfpathlineto{\pgfqpoint{2.634905in}{1.535240in}}%
\pgfpathlineto{\pgfqpoint{2.636563in}{1.536802in}}%
\pgfpathlineto{\pgfqpoint{2.639463in}{1.542047in}}%
\pgfpathlineto{\pgfqpoint{2.643607in}{1.549287in}}%
\pgfpathlineto{\pgfqpoint{2.645679in}{1.550893in}}%
\pgfpathlineto{\pgfqpoint{2.647336in}{1.550816in}}%
\pgfpathlineto{\pgfqpoint{2.649408in}{1.549204in}}%
\pgfpathlineto{\pgfqpoint{2.666811in}{1.528377in}}%
\pgfpathlineto{\pgfqpoint{2.670955in}{1.521372in}}%
\pgfpathlineto{\pgfqpoint{2.673027in}{1.520177in}}%
\pgfpathlineto{\pgfqpoint{2.675098in}{1.520569in}}%
\pgfpathlineto{\pgfqpoint{2.677584in}{1.522329in}}%
\pgfpathlineto{\pgfqpoint{2.679656in}{1.525005in}}%
\pgfpathlineto{\pgfqpoint{2.682142in}{1.530361in}}%
\pgfpathlineto{\pgfqpoint{2.689601in}{1.549036in}}%
\pgfpathlineto{\pgfqpoint{2.691258in}{1.549525in}}%
\pgfpathlineto{\pgfqpoint{2.692916in}{1.548625in}}%
\pgfpathlineto{\pgfqpoint{2.697059in}{1.545765in}}%
\pgfpathlineto{\pgfqpoint{2.698717in}{1.546290in}}%
\pgfpathlineto{\pgfqpoint{2.700374in}{1.548213in}}%
\pgfpathlineto{\pgfqpoint{2.703275in}{1.554019in}}%
\pgfpathlineto{\pgfqpoint{2.706175in}{1.559350in}}%
\pgfpathlineto{\pgfqpoint{2.707833in}{1.560410in}}%
\pgfpathlineto{\pgfqpoint{2.709076in}{1.560006in}}%
\pgfpathlineto{\pgfqpoint{2.710733in}{1.558099in}}%
\pgfpathlineto{\pgfqpoint{2.716120in}{1.550821in}}%
\pgfpathlineto{\pgfqpoint{2.718192in}{1.550372in}}%
\pgfpathlineto{\pgfqpoint{2.722336in}{1.550105in}}%
\pgfpathlineto{\pgfqpoint{2.723993in}{1.548606in}}%
\pgfpathlineto{\pgfqpoint{2.726065in}{1.545096in}}%
\pgfpathlineto{\pgfqpoint{2.731451in}{1.534849in}}%
\pgfpathlineto{\pgfqpoint{2.732695in}{1.534403in}}%
\pgfpathlineto{\pgfqpoint{2.733938in}{1.535247in}}%
\pgfpathlineto{\pgfqpoint{2.735595in}{1.538300in}}%
\pgfpathlineto{\pgfqpoint{2.737667in}{1.544851in}}%
\pgfpathlineto{\pgfqpoint{2.741396in}{1.561656in}}%
\pgfpathlineto{\pgfqpoint{2.747612in}{1.589155in}}%
\pgfpathlineto{\pgfqpoint{2.753413in}{1.608449in}}%
\pgfpathlineto{\pgfqpoint{2.757971in}{1.620943in}}%
\pgfpathlineto{\pgfqpoint{2.760457in}{1.624763in}}%
\pgfpathlineto{\pgfqpoint{2.762529in}{1.625973in}}%
\pgfpathlineto{\pgfqpoint{2.764600in}{1.625860in}}%
\pgfpathlineto{\pgfqpoint{2.768330in}{1.624223in}}%
\pgfpathlineto{\pgfqpoint{2.776617in}{1.620102in}}%
\pgfpathlineto{\pgfqpoint{2.778689in}{1.620482in}}%
\pgfpathlineto{\pgfqpoint{2.781589in}{1.622422in}}%
\pgfpathlineto{\pgfqpoint{2.784904in}{1.624224in}}%
\pgfpathlineto{\pgfqpoint{2.786976in}{1.624008in}}%
\pgfpathlineto{\pgfqpoint{2.789048in}{1.622612in}}%
\pgfpathlineto{\pgfqpoint{2.794849in}{1.617811in}}%
\pgfpathlineto{\pgfqpoint{2.796506in}{1.617839in}}%
\pgfpathlineto{\pgfqpoint{2.798578in}{1.618951in}}%
\pgfpathlineto{\pgfqpoint{2.802722in}{1.622895in}}%
\pgfpathlineto{\pgfqpoint{2.813081in}{1.634254in}}%
\pgfpathlineto{\pgfqpoint{2.822611in}{1.648326in}}%
\pgfpathlineto{\pgfqpoint{2.824683in}{1.649143in}}%
\pgfpathlineto{\pgfqpoint{2.826340in}{1.648743in}}%
\pgfpathlineto{\pgfqpoint{2.828412in}{1.646563in}}%
\pgfpathlineto{\pgfqpoint{2.830898in}{1.642000in}}%
\pgfpathlineto{\pgfqpoint{2.837114in}{1.629761in}}%
\pgfpathlineto{\pgfqpoint{2.840014in}{1.626906in}}%
\pgfpathlineto{\pgfqpoint{2.843744in}{1.625011in}}%
\pgfpathlineto{\pgfqpoint{2.857832in}{1.618797in}}%
\pgfpathlineto{\pgfqpoint{2.860318in}{1.619131in}}%
\pgfpathlineto{\pgfqpoint{2.862804in}{1.620607in}}%
\pgfpathlineto{\pgfqpoint{2.865290in}{1.623356in}}%
\pgfpathlineto{\pgfqpoint{2.868605in}{1.628805in}}%
\pgfpathlineto{\pgfqpoint{2.878964in}{1.647556in}}%
\pgfpathlineto{\pgfqpoint{2.882694in}{1.650723in}}%
\pgfpathlineto{\pgfqpoint{2.888909in}{1.655357in}}%
\pgfpathlineto{\pgfqpoint{2.891395in}{1.659014in}}%
\pgfpathlineto{\pgfqpoint{2.899682in}{1.673159in}}%
\pgfpathlineto{\pgfqpoint{2.901340in}{1.673242in}}%
\pgfpathlineto{\pgfqpoint{2.902997in}{1.671972in}}%
\pgfpathlineto{\pgfqpoint{2.905483in}{1.668410in}}%
\pgfpathlineto{\pgfqpoint{2.912113in}{1.658032in}}%
\pgfpathlineto{\pgfqpoint{2.914599in}{1.656254in}}%
\pgfpathlineto{\pgfqpoint{2.916671in}{1.656246in}}%
\pgfpathlineto{\pgfqpoint{2.918743in}{1.657645in}}%
\pgfpathlineto{\pgfqpoint{2.921229in}{1.660939in}}%
\pgfpathlineto{\pgfqpoint{2.930345in}{1.674832in}}%
\pgfpathlineto{\pgfqpoint{2.933246in}{1.677162in}}%
\pgfpathlineto{\pgfqpoint{2.935317in}{1.677631in}}%
\pgfpathlineto{\pgfqpoint{2.937804in}{1.676995in}}%
\pgfpathlineto{\pgfqpoint{2.945262in}{1.673139in}}%
\pgfpathlineto{\pgfqpoint{2.947748in}{1.670579in}}%
\pgfpathlineto{\pgfqpoint{2.950234in}{1.666430in}}%
\pgfpathlineto{\pgfqpoint{2.953549in}{1.658389in}}%
\pgfpathlineto{\pgfqpoint{2.958936in}{1.644581in}}%
\pgfpathlineto{\pgfqpoint{2.961422in}{1.641182in}}%
\pgfpathlineto{\pgfqpoint{2.963080in}{1.640438in}}%
\pgfpathlineto{\pgfqpoint{2.964737in}{1.640856in}}%
\pgfpathlineto{\pgfqpoint{2.967223in}{1.642931in}}%
\pgfpathlineto{\pgfqpoint{2.970538in}{1.645424in}}%
\pgfpathlineto{\pgfqpoint{2.972196in}{1.645467in}}%
\pgfpathlineto{\pgfqpoint{2.973853in}{1.644444in}}%
\pgfpathlineto{\pgfqpoint{2.975925in}{1.641780in}}%
\pgfpathlineto{\pgfqpoint{2.978826in}{1.635782in}}%
\pgfpathlineto{\pgfqpoint{2.982969in}{1.623873in}}%
\pgfpathlineto{\pgfqpoint{2.986698in}{1.614477in}}%
\pgfpathlineto{\pgfqpoint{2.988770in}{1.611867in}}%
\pgfpathlineto{\pgfqpoint{2.990428in}{1.611447in}}%
\pgfpathlineto{\pgfqpoint{2.992500in}{1.612514in}}%
\pgfpathlineto{\pgfqpoint{2.996229in}{1.616370in}}%
\pgfpathlineto{\pgfqpoint{3.001615in}{1.621863in}}%
\pgfpathlineto{\pgfqpoint{3.003687in}{1.622748in}}%
\pgfpathlineto{\pgfqpoint{3.005345in}{1.622516in}}%
\pgfpathlineto{\pgfqpoint{3.007002in}{1.621287in}}%
\pgfpathlineto{\pgfqpoint{3.009074in}{1.618197in}}%
\pgfpathlineto{\pgfqpoint{3.011560in}{1.612012in}}%
\pgfpathlineto{\pgfqpoint{3.016533in}{1.595413in}}%
\pgfpathlineto{\pgfqpoint{3.020674in}{1.583167in}}%
\pgfpathlineto{\pgfqpoint{3.023518in}{1.577878in}}%
\pgfpathlineto{\pgfqpoint{3.025924in}{1.575608in}}%
\pgfpathlineto{\pgfqpoint{3.026849in}{1.575195in}}%
\pgfpathlineto{\pgfqpoint{3.026849in}{1.575195in}}%
\pgfusepath{stroke}%
\end{pgfscope}%
\begin{pgfscope}%
\pgfpathrectangle{\pgfqpoint{0.650000in}{0.420000in}}{\pgfqpoint{2.388110in}{1.994488in}}%
\pgfusepath{clip}%
\pgfsetrectcap%
\pgfsetroundjoin%
\pgfsetlinewidth{1.003750pt}%
\definecolor{currentstroke}{rgb}{0.501961,0.450980,0.674510}%
\pgfsetstrokecolor{currentstroke}%
\pgfsetdash{}{0pt}%
\pgfpathmoveto{\pgfqpoint{0.650000in}{1.639743in}}%
\pgfpathlineto{\pgfqpoint{0.650930in}{1.641009in}}%
\pgfpathlineto{\pgfqpoint{0.653348in}{1.645711in}}%
\pgfpathlineto{\pgfqpoint{0.656205in}{1.653962in}}%
\pgfpathlineto{\pgfqpoint{0.659117in}{1.665922in}}%
\pgfpathlineto{\pgfqpoint{0.662864in}{1.686686in}}%
\pgfpathlineto{\pgfqpoint{0.667444in}{1.719033in}}%
\pgfpathlineto{\pgfqpoint{0.675771in}{1.778832in}}%
\pgfpathlineto{\pgfqpoint{0.679102in}{1.804462in}}%
\pgfpathlineto{\pgfqpoint{0.682849in}{1.844516in}}%
\pgfpathlineto{\pgfqpoint{0.703667in}{2.083386in}}%
\pgfpathlineto{\pgfqpoint{0.706998in}{2.111898in}}%
\pgfpathlineto{\pgfqpoint{0.709913in}{2.128010in}}%
\pgfpathlineto{\pgfqpoint{0.716574in}{2.161018in}}%
\pgfpathlineto{\pgfqpoint{0.721571in}{2.184613in}}%
\pgfpathlineto{\pgfqpoint{0.724902in}{2.195114in}}%
\pgfpathlineto{\pgfqpoint{0.727400in}{2.199735in}}%
\pgfpathlineto{\pgfqpoint{0.729482in}{2.201512in}}%
\pgfpathlineto{\pgfqpoint{0.731147in}{2.201661in}}%
\pgfpathlineto{\pgfqpoint{0.732812in}{2.200776in}}%
\pgfpathlineto{\pgfqpoint{0.734894in}{2.198355in}}%
\pgfpathlineto{\pgfqpoint{0.737809in}{2.192879in}}%
\pgfpathlineto{\pgfqpoint{0.741556in}{2.183202in}}%
\pgfpathlineto{\pgfqpoint{0.747385in}{2.164741in}}%
\pgfpathlineto{\pgfqpoint{0.769868in}{2.090249in}}%
\pgfpathlineto{\pgfqpoint{0.777779in}{2.069017in}}%
\pgfpathlineto{\pgfqpoint{0.786523in}{2.048671in}}%
\pgfpathlineto{\pgfqpoint{0.797348in}{2.026550in}}%
\pgfpathlineto{\pgfqpoint{0.808590in}{2.006086in}}%
\pgfpathlineto{\pgfqpoint{0.819415in}{1.988455in}}%
\pgfpathlineto{\pgfqpoint{0.830657in}{1.972269in}}%
\pgfpathlineto{\pgfqpoint{0.853556in}{1.941982in}}%
\pgfpathlineto{\pgfqpoint{0.865631in}{1.927437in}}%
\pgfpathlineto{\pgfqpoint{0.875623in}{1.916942in}}%
\pgfpathlineto{\pgfqpoint{0.888530in}{1.905177in}}%
\pgfpathlineto{\pgfqpoint{0.902270in}{1.893984in}}%
\pgfpathlineto{\pgfqpoint{0.935163in}{1.868114in}}%
\pgfpathlineto{\pgfqpoint{0.963059in}{1.849278in}}%
\pgfpathlineto{\pgfqpoint{1.003862in}{1.822090in}}%
\pgfpathlineto{\pgfqpoint{1.079223in}{1.782904in}}%
\pgfpathlineto{\pgfqpoint{1.105037in}{1.772787in}}%
\pgfpathlineto{\pgfqpoint{1.140427in}{1.755307in}}%
\pgfpathlineto{\pgfqpoint{1.172486in}{1.737259in}}%
\pgfpathlineto{\pgfqpoint{1.182063in}{1.730257in}}%
\pgfpathlineto{\pgfqpoint{1.192888in}{1.722892in}}%
\pgfpathlineto{\pgfqpoint{1.205379in}{1.715729in}}%
\pgfpathlineto{\pgfqpoint{1.223699in}{1.707042in}}%
\pgfpathlineto{\pgfqpoint{1.240353in}{1.698912in}}%
\pgfpathlineto{\pgfqpoint{1.251179in}{1.692366in}}%
\pgfpathlineto{\pgfqpoint{1.264502in}{1.684553in}}%
\pgfpathlineto{\pgfqpoint{1.282822in}{1.675674in}}%
\pgfpathlineto{\pgfqpoint{1.323625in}{1.658279in}}%
\pgfpathlineto{\pgfqpoint{1.334450in}{1.655427in}}%
\pgfpathlineto{\pgfqpoint{1.346941in}{1.653348in}}%
\pgfpathlineto{\pgfqpoint{1.367343in}{1.649651in}}%
\pgfpathlineto{\pgfqpoint{1.417722in}{1.637293in}}%
\pgfpathlineto{\pgfqpoint{1.435209in}{1.631909in}}%
\pgfpathlineto{\pgfqpoint{1.456860in}{1.624809in}}%
\pgfpathlineto{\pgfqpoint{1.471433in}{1.620568in}}%
\pgfpathlineto{\pgfqpoint{1.514734in}{1.607963in}}%
\pgfpathlineto{\pgfqpoint{1.528474in}{1.604080in}}%
\pgfpathlineto{\pgfqpoint{1.538883in}{1.601194in}}%
\pgfpathlineto{\pgfqpoint{1.549708in}{1.599610in}}%
\pgfpathlineto{\pgfqpoint{1.561782in}{1.597264in}}%
\pgfpathlineto{\pgfqpoint{1.583850in}{1.593479in}}%
\pgfpathlineto{\pgfqpoint{1.601336in}{1.591559in}}%
\pgfpathlineto{\pgfqpoint{1.615076in}{1.590450in}}%
\pgfpathlineto{\pgfqpoint{1.628400in}{1.588820in}}%
\pgfpathlineto{\pgfqpoint{1.638809in}{1.586130in}}%
\pgfpathlineto{\pgfqpoint{1.663790in}{1.578032in}}%
\pgfpathlineto{\pgfqpoint{1.708757in}{1.566030in}}%
\pgfpathlineto{\pgfqpoint{1.748728in}{1.556626in}}%
\pgfpathlineto{\pgfqpoint{1.767464in}{1.551432in}}%
\pgfpathlineto{\pgfqpoint{1.774542in}{1.549165in}}%
\pgfpathlineto{\pgfqpoint{1.812430in}{1.535570in}}%
\pgfpathlineto{\pgfqpoint{1.824089in}{1.533321in}}%
\pgfpathlineto{\pgfqpoint{1.850736in}{1.528312in}}%
\pgfpathlineto{\pgfqpoint{1.865308in}{1.525767in}}%
\pgfpathlineto{\pgfqpoint{1.884461in}{1.524156in}}%
\pgfpathlineto{\pgfqpoint{1.913190in}{1.516942in}}%
\pgfpathlineto{\pgfqpoint{1.929844in}{1.511933in}}%
\pgfpathlineto{\pgfqpoint{1.967316in}{1.498458in}}%
\pgfpathlineto{\pgfqpoint{1.973978in}{1.496494in}}%
\pgfpathlineto{\pgfqpoint{1.988134in}{1.493156in}}%
\pgfpathlineto{\pgfqpoint{2.001874in}{1.490481in}}%
\pgfpathlineto{\pgfqpoint{2.011866in}{1.490294in}}%
\pgfpathlineto{\pgfqpoint{2.039762in}{1.488623in}}%
\pgfpathlineto{\pgfqpoint{2.063079in}{1.486690in}}%
\pgfpathlineto{\pgfqpoint{2.078484in}{1.486453in}}%
\pgfpathlineto{\pgfqpoint{2.095138in}{1.486910in}}%
\pgfpathlineto{\pgfqpoint{2.099718in}{1.487034in}}%
\pgfpathlineto{\pgfqpoint{2.102633in}{1.488272in}}%
\pgfpathlineto{\pgfqpoint{2.114291in}{1.494823in}}%
\pgfpathlineto{\pgfqpoint{2.117622in}{1.494760in}}%
\pgfpathlineto{\pgfqpoint{2.120953in}{1.493413in}}%
\pgfpathlineto{\pgfqpoint{2.125949in}{1.491209in}}%
\pgfpathlineto{\pgfqpoint{2.130112in}{1.491337in}}%
\pgfpathlineto{\pgfqpoint{2.133443in}{1.490797in}}%
\pgfpathlineto{\pgfqpoint{2.138023in}{1.489974in}}%
\pgfpathlineto{\pgfqpoint{2.143436in}{1.490603in}}%
\pgfpathlineto{\pgfqpoint{2.150514in}{1.491713in}}%
\pgfpathlineto{\pgfqpoint{2.154678in}{1.491256in}}%
\pgfpathlineto{\pgfqpoint{2.158841in}{1.489384in}}%
\pgfpathlineto{\pgfqpoint{2.165503in}{1.486251in}}%
\pgfpathlineto{\pgfqpoint{2.176328in}{1.484299in}}%
\pgfpathlineto{\pgfqpoint{2.181324in}{1.483887in}}%
\pgfpathlineto{\pgfqpoint{2.184239in}{1.482203in}}%
\pgfpathlineto{\pgfqpoint{2.189652in}{1.478870in}}%
\pgfpathlineto{\pgfqpoint{2.195897in}{1.475938in}}%
\pgfpathlineto{\pgfqpoint{2.200061in}{1.472240in}}%
\pgfpathlineto{\pgfqpoint{2.205057in}{1.468019in}}%
\pgfpathlineto{\pgfqpoint{2.207555in}{1.467289in}}%
\pgfpathlineto{\pgfqpoint{2.210469in}{1.467724in}}%
\pgfpathlineto{\pgfqpoint{2.215049in}{1.468786in}}%
\pgfpathlineto{\pgfqpoint{2.217964in}{1.467957in}}%
\pgfpathlineto{\pgfqpoint{2.221295in}{1.467229in}}%
\pgfpathlineto{\pgfqpoint{2.230038in}{1.467588in}}%
\pgfpathlineto{\pgfqpoint{2.232953in}{1.466219in}}%
\pgfpathlineto{\pgfqpoint{2.237533in}{1.463628in}}%
\pgfpathlineto{\pgfqpoint{2.249607in}{1.461764in}}%
\pgfpathlineto{\pgfqpoint{2.253771in}{1.463116in}}%
\pgfpathlineto{\pgfqpoint{2.260432in}{1.465815in}}%
\pgfpathlineto{\pgfqpoint{2.272507in}{1.473629in}}%
\pgfpathlineto{\pgfqpoint{2.275421in}{1.474799in}}%
\pgfpathlineto{\pgfqpoint{2.277503in}{1.474289in}}%
\pgfpathlineto{\pgfqpoint{2.280002in}{1.472096in}}%
\pgfpathlineto{\pgfqpoint{2.284165in}{1.468034in}}%
\pgfpathlineto{\pgfqpoint{2.287079in}{1.466958in}}%
\pgfpathlineto{\pgfqpoint{2.292908in}{1.466210in}}%
\pgfpathlineto{\pgfqpoint{2.297488in}{1.467140in}}%
\pgfpathlineto{\pgfqpoint{2.307065in}{1.469156in}}%
\pgfpathlineto{\pgfqpoint{2.310812in}{1.470500in}}%
\pgfpathlineto{\pgfqpoint{2.324135in}{1.471575in}}%
\pgfpathlineto{\pgfqpoint{2.331630in}{1.474899in}}%
\pgfpathlineto{\pgfqpoint{2.341206in}{1.476686in}}%
\pgfpathlineto{\pgfqpoint{2.347451in}{1.477033in}}%
\pgfpathlineto{\pgfqpoint{2.350366in}{1.475875in}}%
\pgfpathlineto{\pgfqpoint{2.358277in}{1.472022in}}%
\pgfpathlineto{\pgfqpoint{2.362441in}{1.471374in}}%
\pgfpathlineto{\pgfqpoint{2.365355in}{1.472114in}}%
\pgfpathlineto{\pgfqpoint{2.368269in}{1.474179in}}%
\pgfpathlineto{\pgfqpoint{2.372433in}{1.477422in}}%
\pgfpathlineto{\pgfqpoint{2.374931in}{1.477816in}}%
\pgfpathlineto{\pgfqpoint{2.380761in}{1.476934in}}%
\pgfpathlineto{\pgfqpoint{2.388255in}{1.474661in}}%
\pgfpathlineto{\pgfqpoint{2.392835in}{1.475598in}}%
\pgfpathlineto{\pgfqpoint{2.394501in}{1.474626in}}%
\pgfpathlineto{\pgfqpoint{2.396998in}{1.471508in}}%
\pgfpathlineto{\pgfqpoint{2.410738in}{1.451581in}}%
\pgfpathlineto{\pgfqpoint{2.413653in}{1.448359in}}%
\pgfpathlineto{\pgfqpoint{2.416151in}{1.447081in}}%
\pgfpathlineto{\pgfqpoint{2.419898in}{1.446764in}}%
\pgfpathlineto{\pgfqpoint{2.425727in}{1.447262in}}%
\pgfpathlineto{\pgfqpoint{2.430724in}{1.448081in}}%
\pgfpathlineto{\pgfqpoint{2.433638in}{1.447025in}}%
\pgfpathlineto{\pgfqpoint{2.441133in}{1.442717in}}%
\pgfpathlineto{\pgfqpoint{2.446129in}{1.437848in}}%
\pgfpathlineto{\pgfqpoint{2.453623in}{1.430019in}}%
\pgfpathlineto{\pgfqpoint{2.456538in}{1.428920in}}%
\pgfpathlineto{\pgfqpoint{2.461118in}{1.427516in}}%
\pgfpathlineto{\pgfqpoint{2.467363in}{1.424950in}}%
\pgfpathlineto{\pgfqpoint{2.470278in}{1.425231in}}%
\pgfpathlineto{\pgfqpoint{2.474441in}{1.425862in}}%
\pgfpathlineto{\pgfqpoint{2.477773in}{1.424832in}}%
\pgfpathlineto{\pgfqpoint{2.484018in}{1.422540in}}%
\pgfpathlineto{\pgfqpoint{2.486100in}{1.423191in}}%
\pgfpathlineto{\pgfqpoint{2.488598in}{1.425197in}}%
\pgfpathlineto{\pgfqpoint{2.498590in}{1.434813in}}%
\pgfpathlineto{\pgfqpoint{2.501088in}{1.435651in}}%
\pgfpathlineto{\pgfqpoint{2.503170in}{1.435206in}}%
\pgfpathlineto{\pgfqpoint{2.505668in}{1.433240in}}%
\pgfpathlineto{\pgfqpoint{2.510248in}{1.427269in}}%
\pgfpathlineto{\pgfqpoint{2.512747in}{1.425056in}}%
\pgfpathlineto{\pgfqpoint{2.515244in}{1.424463in}}%
\pgfpathlineto{\pgfqpoint{2.518576in}{1.423689in}}%
\pgfpathlineto{\pgfqpoint{2.524821in}{1.420948in}}%
\pgfpathlineto{\pgfqpoint{2.526903in}{1.421834in}}%
\pgfpathlineto{\pgfqpoint{2.529817in}{1.424850in}}%
\pgfpathlineto{\pgfqpoint{2.532732in}{1.427381in}}%
\pgfpathlineto{\pgfqpoint{2.535646in}{1.428207in}}%
\pgfpathlineto{\pgfqpoint{2.541059in}{1.429514in}}%
\pgfpathlineto{\pgfqpoint{2.543973in}{1.429898in}}%
\pgfpathlineto{\pgfqpoint{2.556880in}{1.427586in}}%
\pgfpathlineto{\pgfqpoint{2.561044in}{1.426662in}}%
\pgfpathlineto{\pgfqpoint{2.564791in}{1.425895in}}%
\pgfpathlineto{\pgfqpoint{2.567705in}{1.424073in}}%
\pgfpathlineto{\pgfqpoint{2.572285in}{1.421222in}}%
\pgfpathlineto{\pgfqpoint{2.576864in}{1.420097in}}%
\pgfpathlineto{\pgfqpoint{2.586023in}{1.419298in}}%
\pgfpathlineto{\pgfqpoint{2.588105in}{1.421246in}}%
\pgfpathlineto{\pgfqpoint{2.595183in}{1.429474in}}%
\pgfpathlineto{\pgfqpoint{2.601011in}{1.432836in}}%
\pgfpathlineto{\pgfqpoint{2.602677in}{1.432293in}}%
\pgfpathlineto{\pgfqpoint{2.604758in}{1.429981in}}%
\pgfpathlineto{\pgfqpoint{2.609754in}{1.423579in}}%
\pgfpathlineto{\pgfqpoint{2.616831in}{1.418288in}}%
\pgfpathlineto{\pgfqpoint{2.621411in}{1.410055in}}%
\pgfpathlineto{\pgfqpoint{2.624742in}{1.405460in}}%
\pgfpathlineto{\pgfqpoint{2.626823in}{1.404137in}}%
\pgfpathlineto{\pgfqpoint{2.628489in}{1.404176in}}%
\pgfpathlineto{\pgfqpoint{2.630570in}{1.405573in}}%
\pgfpathlineto{\pgfqpoint{2.635566in}{1.409893in}}%
\pgfpathlineto{\pgfqpoint{2.637232in}{1.409446in}}%
\pgfpathlineto{\pgfqpoint{2.639313in}{1.407329in}}%
\pgfpathlineto{\pgfqpoint{2.642644in}{1.403995in}}%
\pgfpathlineto{\pgfqpoint{2.644309in}{1.403983in}}%
\pgfpathlineto{\pgfqpoint{2.645974in}{1.405300in}}%
\pgfpathlineto{\pgfqpoint{2.648472in}{1.408993in}}%
\pgfpathlineto{\pgfqpoint{2.653052in}{1.415876in}}%
\pgfpathlineto{\pgfqpoint{2.655133in}{1.417306in}}%
\pgfpathlineto{\pgfqpoint{2.656799in}{1.417316in}}%
\pgfpathlineto{\pgfqpoint{2.658880in}{1.416015in}}%
\pgfpathlineto{\pgfqpoint{2.670538in}{1.405734in}}%
\pgfpathlineto{\pgfqpoint{2.674701in}{1.402480in}}%
\pgfpathlineto{\pgfqpoint{2.681362in}{1.395249in}}%
\pgfpathlineto{\pgfqpoint{2.683027in}{1.395329in}}%
\pgfpathlineto{\pgfqpoint{2.685525in}{1.396901in}}%
\pgfpathlineto{\pgfqpoint{2.688439in}{1.400264in}}%
\pgfpathlineto{\pgfqpoint{2.690937in}{1.404964in}}%
\pgfpathlineto{\pgfqpoint{2.694684in}{1.415082in}}%
\pgfpathlineto{\pgfqpoint{2.698015in}{1.423137in}}%
\pgfpathlineto{\pgfqpoint{2.700097in}{1.425378in}}%
\pgfpathlineto{\pgfqpoint{2.701762in}{1.425537in}}%
\pgfpathlineto{\pgfqpoint{2.707590in}{1.424045in}}%
\pgfpathlineto{\pgfqpoint{2.709256in}{1.425459in}}%
\pgfpathlineto{\pgfqpoint{2.711754in}{1.429603in}}%
\pgfpathlineto{\pgfqpoint{2.716333in}{1.437991in}}%
\pgfpathlineto{\pgfqpoint{2.717998in}{1.438948in}}%
\pgfpathlineto{\pgfqpoint{2.719664in}{1.438402in}}%
\pgfpathlineto{\pgfqpoint{2.722162in}{1.435727in}}%
\pgfpathlineto{\pgfqpoint{2.725076in}{1.432837in}}%
\pgfpathlineto{\pgfqpoint{2.727158in}{1.432176in}}%
\pgfpathlineto{\pgfqpoint{2.732986in}{1.431585in}}%
\pgfpathlineto{\pgfqpoint{2.735068in}{1.429250in}}%
\pgfpathlineto{\pgfqpoint{2.738815in}{1.422517in}}%
\pgfpathlineto{\pgfqpoint{2.741313in}{1.419268in}}%
\pgfpathlineto{\pgfqpoint{2.742562in}{1.418880in}}%
\pgfpathlineto{\pgfqpoint{2.743811in}{1.419579in}}%
\pgfpathlineto{\pgfqpoint{2.745476in}{1.422153in}}%
\pgfpathlineto{\pgfqpoint{2.747558in}{1.427711in}}%
\pgfpathlineto{\pgfqpoint{2.750888in}{1.440422in}}%
\pgfpathlineto{\pgfqpoint{2.757549in}{1.465885in}}%
\pgfpathlineto{\pgfqpoint{2.762545in}{1.479859in}}%
\pgfpathlineto{\pgfqpoint{2.767541in}{1.491525in}}%
\pgfpathlineto{\pgfqpoint{2.770039in}{1.494874in}}%
\pgfpathlineto{\pgfqpoint{2.772121in}{1.495880in}}%
\pgfpathlineto{\pgfqpoint{2.774202in}{1.495587in}}%
\pgfpathlineto{\pgfqpoint{2.778366in}{1.493490in}}%
\pgfpathlineto{\pgfqpoint{2.785443in}{1.489919in}}%
\pgfpathlineto{\pgfqpoint{2.787941in}{1.489863in}}%
\pgfpathlineto{\pgfqpoint{2.790439in}{1.491084in}}%
\pgfpathlineto{\pgfqpoint{2.795019in}{1.493737in}}%
\pgfpathlineto{\pgfqpoint{2.797100in}{1.493725in}}%
\pgfpathlineto{\pgfqpoint{2.799598in}{1.492381in}}%
\pgfpathlineto{\pgfqpoint{2.804594in}{1.489092in}}%
\pgfpathlineto{\pgfqpoint{2.806676in}{1.489149in}}%
\pgfpathlineto{\pgfqpoint{2.809174in}{1.490564in}}%
\pgfpathlineto{\pgfqpoint{2.814170in}{1.495205in}}%
\pgfpathlineto{\pgfqpoint{2.822496in}{1.503878in}}%
\pgfpathlineto{\pgfqpoint{2.827492in}{1.511666in}}%
\pgfpathlineto{\pgfqpoint{2.831655in}{1.517441in}}%
\pgfpathlineto{\pgfqpoint{2.834153in}{1.519336in}}%
\pgfpathlineto{\pgfqpoint{2.836235in}{1.519628in}}%
\pgfpathlineto{\pgfqpoint{2.837900in}{1.518833in}}%
\pgfpathlineto{\pgfqpoint{2.839982in}{1.516453in}}%
\pgfpathlineto{\pgfqpoint{2.848725in}{1.504213in}}%
\pgfpathlineto{\pgfqpoint{2.851222in}{1.503360in}}%
\pgfpathlineto{\pgfqpoint{2.862463in}{1.502277in}}%
\pgfpathlineto{\pgfqpoint{2.869125in}{1.502417in}}%
\pgfpathlineto{\pgfqpoint{2.872039in}{1.503955in}}%
\pgfpathlineto{\pgfqpoint{2.874537in}{1.506454in}}%
\pgfpathlineto{\pgfqpoint{2.877867in}{1.511362in}}%
\pgfpathlineto{\pgfqpoint{2.882863in}{1.521270in}}%
\pgfpathlineto{\pgfqpoint{2.887443in}{1.529524in}}%
\pgfpathlineto{\pgfqpoint{2.891190in}{1.534043in}}%
\pgfpathlineto{\pgfqpoint{2.896602in}{1.538402in}}%
\pgfpathlineto{\pgfqpoint{2.899516in}{1.541427in}}%
\pgfpathlineto{\pgfqpoint{2.902431in}{1.546228in}}%
\pgfpathlineto{\pgfqpoint{2.909924in}{1.560013in}}%
\pgfpathlineto{\pgfqpoint{2.911590in}{1.560839in}}%
\pgfpathlineto{\pgfqpoint{2.913255in}{1.560398in}}%
\pgfpathlineto{\pgfqpoint{2.915337in}{1.558518in}}%
\pgfpathlineto{\pgfqpoint{2.922830in}{1.550350in}}%
\pgfpathlineto{\pgfqpoint{2.924912in}{1.549642in}}%
\pgfpathlineto{\pgfqpoint{2.926994in}{1.550163in}}%
\pgfpathlineto{\pgfqpoint{2.929075in}{1.552104in}}%
\pgfpathlineto{\pgfqpoint{2.931990in}{1.556820in}}%
\pgfpathlineto{\pgfqpoint{2.941565in}{1.574320in}}%
\pgfpathlineto{\pgfqpoint{2.944479in}{1.577400in}}%
\pgfpathlineto{\pgfqpoint{2.946977in}{1.578467in}}%
\pgfpathlineto{\pgfqpoint{2.950308in}{1.578366in}}%
\pgfpathlineto{\pgfqpoint{2.956969in}{1.577145in}}%
\pgfpathlineto{\pgfqpoint{2.959467in}{1.575171in}}%
\pgfpathlineto{\pgfqpoint{2.961965in}{1.571651in}}%
\pgfpathlineto{\pgfqpoint{2.965712in}{1.563656in}}%
\pgfpathlineto{\pgfqpoint{2.969875in}{1.555150in}}%
\pgfpathlineto{\pgfqpoint{2.972373in}{1.552584in}}%
\pgfpathlineto{\pgfqpoint{2.974038in}{1.552242in}}%
\pgfpathlineto{\pgfqpoint{2.975704in}{1.552944in}}%
\pgfpathlineto{\pgfqpoint{2.978618in}{1.555768in}}%
\pgfpathlineto{\pgfqpoint{2.981949in}{1.558546in}}%
\pgfpathlineto{\pgfqpoint{2.984030in}{1.558752in}}%
\pgfpathlineto{\pgfqpoint{2.986112in}{1.557558in}}%
\pgfpathlineto{\pgfqpoint{2.988610in}{1.554402in}}%
\pgfpathlineto{\pgfqpoint{2.991524in}{1.548576in}}%
\pgfpathlineto{\pgfqpoint{2.998185in}{1.534467in}}%
\pgfpathlineto{\pgfqpoint{2.999851in}{1.533206in}}%
\pgfpathlineto{\pgfqpoint{3.001516in}{1.533272in}}%
\pgfpathlineto{\pgfqpoint{3.003597in}{1.534840in}}%
\pgfpathlineto{\pgfqpoint{3.008593in}{1.540965in}}%
\pgfpathlineto{\pgfqpoint{3.012757in}{1.545367in}}%
\pgfpathlineto{\pgfqpoint{3.014838in}{1.546376in}}%
\pgfpathlineto{\pgfqpoint{3.016503in}{1.546230in}}%
\pgfpathlineto{\pgfqpoint{3.018169in}{1.545020in}}%
\pgfpathlineto{\pgfqpoint{3.020250in}{1.541915in}}%
\pgfpathlineto{\pgfqpoint{3.022748in}{1.535958in}}%
\pgfpathlineto{\pgfqpoint{3.028161in}{1.518773in}}%
\pgfpathlineto{\pgfqpoint{3.032320in}{1.507658in}}%
\pgfpathlineto{\pgfqpoint{3.035153in}{1.503136in}}%
\pgfpathlineto{\pgfqpoint{3.037652in}{1.501341in}}%
\pgfpathlineto{\pgfqpoint{3.038110in}{1.501193in}}%
\pgfpathlineto{\pgfqpoint{3.038110in}{1.501193in}}%
\pgfusepath{stroke}%
\end{pgfscope}%
\begin{pgfscope}%
\pgfsetrectcap%
\pgfsetmiterjoin%
\pgfsetlinewidth{0.803000pt}%
\definecolor{currentstroke}{rgb}{0.000000,0.000000,0.000000}%
\pgfsetstrokecolor{currentstroke}%
\pgfsetdash{}{0pt}%
\pgfpathmoveto{\pgfqpoint{0.650000in}{0.420000in}}%
\pgfpathlineto{\pgfqpoint{0.650000in}{2.414488in}}%
\pgfusepath{stroke}%
\end{pgfscope}%
\begin{pgfscope}%
\pgfsetrectcap%
\pgfsetmiterjoin%
\pgfsetlinewidth{0.803000pt}%
\definecolor{currentstroke}{rgb}{0.000000,0.000000,0.000000}%
\pgfsetstrokecolor{currentstroke}%
\pgfsetdash{}{0pt}%
\pgfpathmoveto{\pgfqpoint{3.038110in}{0.420000in}}%
\pgfpathlineto{\pgfqpoint{3.038110in}{2.414488in}}%
\pgfusepath{stroke}%
\end{pgfscope}%
\begin{pgfscope}%
\pgfsetrectcap%
\pgfsetmiterjoin%
\pgfsetlinewidth{0.803000pt}%
\definecolor{currentstroke}{rgb}{0.000000,0.000000,0.000000}%
\pgfsetstrokecolor{currentstroke}%
\pgfsetdash{}{0pt}%
\pgfpathmoveto{\pgfqpoint{0.650000in}{0.420000in}}%
\pgfpathlineto{\pgfqpoint{3.038110in}{0.420000in}}%
\pgfusepath{stroke}%
\end{pgfscope}%
\begin{pgfscope}%
\pgfsetrectcap%
\pgfsetmiterjoin%
\pgfsetlinewidth{0.803000pt}%
\definecolor{currentstroke}{rgb}{0.000000,0.000000,0.000000}%
\pgfsetstrokecolor{currentstroke}%
\pgfsetdash{}{0pt}%
\pgfpathmoveto{\pgfqpoint{0.650000in}{2.414488in}}%
\pgfpathlineto{\pgfqpoint{3.038110in}{2.414488in}}%
\pgfusepath{stroke}%
\end{pgfscope}%
\end{pgfpicture}%
\makeatother%
\endgroup%

    \label{fig:R500deg18R100deg90_cf}
\end{subfigure}
\hfill
\begin{subfigure}[t]{0.495\textwidth}
    \subcaption{}
    \centering
    %% Creator: Matplotlib, PGF backend
%%
%% To include the figure in your LaTeX document, write
%%   \input{<filename>.pgf}
%%
%% Make sure the required packages are loaded in your preamble
%%   \usepackage{pgf}
%%
%% Also ensure that all the required font packages are loaded; for instance,
%% the lmodern package is sometimes necessary when using math font.
%%   \usepackage{lmodern}
%%
%% Figures using additional raster images can only be included by \input if
%% they are in the same directory as the main LaTeX file. For loading figures
%% from other directories you can use the `import` package
%%   \usepackage{import}
%%
%% and then include the figures with
%%   \import{<path to file>}{<filename>.pgf}
%%
%% Matplotlib used the following preamble
%%   \def\mathdefault#1{#1}
%%   \everymath=\expandafter{\the\everymath\displaystyle}
%%   \IfFileExists{scrextend.sty}{
%%     \usepackage[fontsize=11.000000pt]{scrextend}
%%   }{
%%     \renewcommand{\normalsize}{\fontsize{11.000000}{13.200000}\selectfont}
%%     \normalsize
%%   }
%%   
%%   \makeatletter\@ifpackageloaded{underscore}{}{\usepackage[strings]{underscore}}\makeatother
%%
\begingroup%
\makeatletter%
\begin{pgfpicture}%
\pgfpathrectangle{\pgfpointorigin}{\pgfqpoint{3.118110in}{2.494488in}}%
\pgfusepath{use as bounding box, clip}%
\begin{pgfscope}%
\pgfsetbuttcap%
\pgfsetmiterjoin%
\definecolor{currentfill}{rgb}{1.000000,1.000000,1.000000}%
\pgfsetfillcolor{currentfill}%
\pgfsetlinewidth{0.000000pt}%
\definecolor{currentstroke}{rgb}{1.000000,1.000000,1.000000}%
\pgfsetstrokecolor{currentstroke}%
\pgfsetdash{}{0pt}%
\pgfpathmoveto{\pgfqpoint{0.000000in}{0.000000in}}%
\pgfpathlineto{\pgfqpoint{3.118110in}{0.000000in}}%
\pgfpathlineto{\pgfqpoint{3.118110in}{2.494488in}}%
\pgfpathlineto{\pgfqpoint{0.000000in}{2.494488in}}%
\pgfpathlineto{\pgfqpoint{0.000000in}{0.000000in}}%
\pgfpathclose%
\pgfusepath{fill}%
\end{pgfscope}%
\begin{pgfscope}%
\pgfsetbuttcap%
\pgfsetmiterjoin%
\definecolor{currentfill}{rgb}{1.000000,1.000000,1.000000}%
\pgfsetfillcolor{currentfill}%
\pgfsetlinewidth{0.000000pt}%
\definecolor{currentstroke}{rgb}{0.000000,0.000000,0.000000}%
\pgfsetstrokecolor{currentstroke}%
\pgfsetstrokeopacity{0.000000}%
\pgfsetdash{}{0pt}%
\pgfpathmoveto{\pgfqpoint{0.650000in}{0.420000in}}%
\pgfpathlineto{\pgfqpoint{3.038110in}{0.420000in}}%
\pgfpathlineto{\pgfqpoint{3.038110in}{2.414488in}}%
\pgfpathlineto{\pgfqpoint{0.650000in}{2.414488in}}%
\pgfpathlineto{\pgfqpoint{0.650000in}{0.420000in}}%
\pgfpathclose%
\pgfusepath{fill}%
\end{pgfscope}%
\begin{pgfscope}%
\pgfsetbuttcap%
\pgfsetroundjoin%
\definecolor{currentfill}{rgb}{0.000000,0.000000,0.000000}%
\pgfsetfillcolor{currentfill}%
\pgfsetlinewidth{0.501875pt}%
\definecolor{currentstroke}{rgb}{0.000000,0.000000,0.000000}%
\pgfsetstrokecolor{currentstroke}%
\pgfsetdash{}{0pt}%
\pgfsys@defobject{currentmarker}{\pgfqpoint{0.000000in}{0.000000in}}{\pgfqpoint{0.000000in}{0.041667in}}{%
\pgfpathmoveto{\pgfqpoint{0.000000in}{0.000000in}}%
\pgfpathlineto{\pgfqpoint{0.000000in}{0.041667in}}%
\pgfusepath{stroke,fill}%
}%
\begin{pgfscope}%
\pgfsys@transformshift{0.650000in}{0.420000in}%
\pgfsys@useobject{currentmarker}{}%
\end{pgfscope}%
\end{pgfscope}%
\begin{pgfscope}%
\pgfsetbuttcap%
\pgfsetroundjoin%
\definecolor{currentfill}{rgb}{0.000000,0.000000,0.000000}%
\pgfsetfillcolor{currentfill}%
\pgfsetlinewidth{0.501875pt}%
\definecolor{currentstroke}{rgb}{0.000000,0.000000,0.000000}%
\pgfsetstrokecolor{currentstroke}%
\pgfsetdash{}{0pt}%
\pgfsys@defobject{currentmarker}{\pgfqpoint{0.000000in}{-0.041667in}}{\pgfqpoint{0.000000in}{0.000000in}}{%
\pgfpathmoveto{\pgfqpoint{0.000000in}{0.000000in}}%
\pgfpathlineto{\pgfqpoint{0.000000in}{-0.041667in}}%
\pgfusepath{stroke,fill}%
}%
\begin{pgfscope}%
\pgfsys@transformshift{0.650000in}{2.414488in}%
\pgfsys@useobject{currentmarker}{}%
\end{pgfscope}%
\end{pgfscope}%
\begin{pgfscope}%
\definecolor{textcolor}{rgb}{0.000000,0.000000,0.000000}%
\pgfsetstrokecolor{textcolor}%
\pgfsetfillcolor{textcolor}%
\pgftext[x=0.650000in,y=0.371389in,,top]{\color{textcolor}{\rmfamily\fontsize{11.000000}{13.200000}\selectfont\catcode`\^=\active\def^{\ifmmode\sp\else\^{}\fi}\catcode`\%=\active\def%{\%}0}}%
\end{pgfscope}%
\begin{pgfscope}%
\pgfsetbuttcap%
\pgfsetroundjoin%
\definecolor{currentfill}{rgb}{0.000000,0.000000,0.000000}%
\pgfsetfillcolor{currentfill}%
\pgfsetlinewidth{0.501875pt}%
\definecolor{currentstroke}{rgb}{0.000000,0.000000,0.000000}%
\pgfsetstrokecolor{currentstroke}%
\pgfsetdash{}{0pt}%
\pgfsys@defobject{currentmarker}{\pgfqpoint{0.000000in}{0.000000in}}{\pgfqpoint{0.000000in}{0.041667in}}{%
\pgfpathmoveto{\pgfqpoint{0.000000in}{0.000000in}}%
\pgfpathlineto{\pgfqpoint{0.000000in}{0.041667in}}%
\pgfusepath{stroke,fill}%
}%
\begin{pgfscope}%
\pgfsys@transformshift{1.493052in}{0.420000in}%
\pgfsys@useobject{currentmarker}{}%
\end{pgfscope}%
\end{pgfscope}%
\begin{pgfscope}%
\pgfsetbuttcap%
\pgfsetroundjoin%
\definecolor{currentfill}{rgb}{0.000000,0.000000,0.000000}%
\pgfsetfillcolor{currentfill}%
\pgfsetlinewidth{0.501875pt}%
\definecolor{currentstroke}{rgb}{0.000000,0.000000,0.000000}%
\pgfsetstrokecolor{currentstroke}%
\pgfsetdash{}{0pt}%
\pgfsys@defobject{currentmarker}{\pgfqpoint{0.000000in}{-0.041667in}}{\pgfqpoint{0.000000in}{0.000000in}}{%
\pgfpathmoveto{\pgfqpoint{0.000000in}{0.000000in}}%
\pgfpathlineto{\pgfqpoint{0.000000in}{-0.041667in}}%
\pgfusepath{stroke,fill}%
}%
\begin{pgfscope}%
\pgfsys@transformshift{1.493052in}{2.414488in}%
\pgfsys@useobject{currentmarker}{}%
\end{pgfscope}%
\end{pgfscope}%
\begin{pgfscope}%
\definecolor{textcolor}{rgb}{0.000000,0.000000,0.000000}%
\pgfsetstrokecolor{textcolor}%
\pgfsetfillcolor{textcolor}%
\pgftext[x=1.493052in,y=0.371389in,,top]{\color{textcolor}{\rmfamily\fontsize{11.000000}{13.200000}\selectfont\catcode`\^=\active\def^{\ifmmode\sp\else\^{}\fi}\catcode`\%=\active\def%{\%}100}}%
\end{pgfscope}%
\begin{pgfscope}%
\pgfsetbuttcap%
\pgfsetroundjoin%
\definecolor{currentfill}{rgb}{0.000000,0.000000,0.000000}%
\pgfsetfillcolor{currentfill}%
\pgfsetlinewidth{0.501875pt}%
\definecolor{currentstroke}{rgb}{0.000000,0.000000,0.000000}%
\pgfsetstrokecolor{currentstroke}%
\pgfsetdash{}{0pt}%
\pgfsys@defobject{currentmarker}{\pgfqpoint{0.000000in}{0.000000in}}{\pgfqpoint{0.000000in}{0.041667in}}{%
\pgfpathmoveto{\pgfqpoint{0.000000in}{0.000000in}}%
\pgfpathlineto{\pgfqpoint{0.000000in}{0.041667in}}%
\pgfusepath{stroke,fill}%
}%
\begin{pgfscope}%
\pgfsys@transformshift{2.336103in}{0.420000in}%
\pgfsys@useobject{currentmarker}{}%
\end{pgfscope}%
\end{pgfscope}%
\begin{pgfscope}%
\pgfsetbuttcap%
\pgfsetroundjoin%
\definecolor{currentfill}{rgb}{0.000000,0.000000,0.000000}%
\pgfsetfillcolor{currentfill}%
\pgfsetlinewidth{0.501875pt}%
\definecolor{currentstroke}{rgb}{0.000000,0.000000,0.000000}%
\pgfsetstrokecolor{currentstroke}%
\pgfsetdash{}{0pt}%
\pgfsys@defobject{currentmarker}{\pgfqpoint{0.000000in}{-0.041667in}}{\pgfqpoint{0.000000in}{0.000000in}}{%
\pgfpathmoveto{\pgfqpoint{0.000000in}{0.000000in}}%
\pgfpathlineto{\pgfqpoint{0.000000in}{-0.041667in}}%
\pgfusepath{stroke,fill}%
}%
\begin{pgfscope}%
\pgfsys@transformshift{2.336103in}{2.414488in}%
\pgfsys@useobject{currentmarker}{}%
\end{pgfscope}%
\end{pgfscope}%
\begin{pgfscope}%
\definecolor{textcolor}{rgb}{0.000000,0.000000,0.000000}%
\pgfsetstrokecolor{textcolor}%
\pgfsetfillcolor{textcolor}%
\pgftext[x=2.336103in,y=0.371389in,,top]{\color{textcolor}{\rmfamily\fontsize{11.000000}{13.200000}\selectfont\catcode`\^=\active\def^{\ifmmode\sp\else\^{}\fi}\catcode`\%=\active\def%{\%}200}}%
\end{pgfscope}%
\begin{pgfscope}%
\pgfsetbuttcap%
\pgfsetroundjoin%
\definecolor{currentfill}{rgb}{0.000000,0.000000,0.000000}%
\pgfsetfillcolor{currentfill}%
\pgfsetlinewidth{0.401500pt}%
\definecolor{currentstroke}{rgb}{0.000000,0.000000,0.000000}%
\pgfsetstrokecolor{currentstroke}%
\pgfsetdash{}{0pt}%
\pgfsys@defobject{currentmarker}{\pgfqpoint{0.000000in}{0.000000in}}{\pgfqpoint{0.000000in}{0.020833in}}{%
\pgfpathmoveto{\pgfqpoint{0.000000in}{0.000000in}}%
\pgfpathlineto{\pgfqpoint{0.000000in}{0.020833in}}%
\pgfusepath{stroke,fill}%
}%
\begin{pgfscope}%
\pgfsys@transformshift{0.818610in}{0.420000in}%
\pgfsys@useobject{currentmarker}{}%
\end{pgfscope}%
\end{pgfscope}%
\begin{pgfscope}%
\pgfsetbuttcap%
\pgfsetroundjoin%
\definecolor{currentfill}{rgb}{0.000000,0.000000,0.000000}%
\pgfsetfillcolor{currentfill}%
\pgfsetlinewidth{0.401500pt}%
\definecolor{currentstroke}{rgb}{0.000000,0.000000,0.000000}%
\pgfsetstrokecolor{currentstroke}%
\pgfsetdash{}{0pt}%
\pgfsys@defobject{currentmarker}{\pgfqpoint{0.000000in}{-0.020833in}}{\pgfqpoint{0.000000in}{0.000000in}}{%
\pgfpathmoveto{\pgfqpoint{0.000000in}{0.000000in}}%
\pgfpathlineto{\pgfqpoint{0.000000in}{-0.020833in}}%
\pgfusepath{stroke,fill}%
}%
\begin{pgfscope}%
\pgfsys@transformshift{0.818610in}{2.414488in}%
\pgfsys@useobject{currentmarker}{}%
\end{pgfscope}%
\end{pgfscope}%
\begin{pgfscope}%
\pgfsetbuttcap%
\pgfsetroundjoin%
\definecolor{currentfill}{rgb}{0.000000,0.000000,0.000000}%
\pgfsetfillcolor{currentfill}%
\pgfsetlinewidth{0.401500pt}%
\definecolor{currentstroke}{rgb}{0.000000,0.000000,0.000000}%
\pgfsetstrokecolor{currentstroke}%
\pgfsetdash{}{0pt}%
\pgfsys@defobject{currentmarker}{\pgfqpoint{0.000000in}{0.000000in}}{\pgfqpoint{0.000000in}{0.020833in}}{%
\pgfpathmoveto{\pgfqpoint{0.000000in}{0.000000in}}%
\pgfpathlineto{\pgfqpoint{0.000000in}{0.020833in}}%
\pgfusepath{stroke,fill}%
}%
\begin{pgfscope}%
\pgfsys@transformshift{0.987221in}{0.420000in}%
\pgfsys@useobject{currentmarker}{}%
\end{pgfscope}%
\end{pgfscope}%
\begin{pgfscope}%
\pgfsetbuttcap%
\pgfsetroundjoin%
\definecolor{currentfill}{rgb}{0.000000,0.000000,0.000000}%
\pgfsetfillcolor{currentfill}%
\pgfsetlinewidth{0.401500pt}%
\definecolor{currentstroke}{rgb}{0.000000,0.000000,0.000000}%
\pgfsetstrokecolor{currentstroke}%
\pgfsetdash{}{0pt}%
\pgfsys@defobject{currentmarker}{\pgfqpoint{0.000000in}{-0.020833in}}{\pgfqpoint{0.000000in}{0.000000in}}{%
\pgfpathmoveto{\pgfqpoint{0.000000in}{0.000000in}}%
\pgfpathlineto{\pgfqpoint{0.000000in}{-0.020833in}}%
\pgfusepath{stroke,fill}%
}%
\begin{pgfscope}%
\pgfsys@transformshift{0.987221in}{2.414488in}%
\pgfsys@useobject{currentmarker}{}%
\end{pgfscope}%
\end{pgfscope}%
\begin{pgfscope}%
\pgfsetbuttcap%
\pgfsetroundjoin%
\definecolor{currentfill}{rgb}{0.000000,0.000000,0.000000}%
\pgfsetfillcolor{currentfill}%
\pgfsetlinewidth{0.401500pt}%
\definecolor{currentstroke}{rgb}{0.000000,0.000000,0.000000}%
\pgfsetstrokecolor{currentstroke}%
\pgfsetdash{}{0pt}%
\pgfsys@defobject{currentmarker}{\pgfqpoint{0.000000in}{0.000000in}}{\pgfqpoint{0.000000in}{0.020833in}}{%
\pgfpathmoveto{\pgfqpoint{0.000000in}{0.000000in}}%
\pgfpathlineto{\pgfqpoint{0.000000in}{0.020833in}}%
\pgfusepath{stroke,fill}%
}%
\begin{pgfscope}%
\pgfsys@transformshift{1.155831in}{0.420000in}%
\pgfsys@useobject{currentmarker}{}%
\end{pgfscope}%
\end{pgfscope}%
\begin{pgfscope}%
\pgfsetbuttcap%
\pgfsetroundjoin%
\definecolor{currentfill}{rgb}{0.000000,0.000000,0.000000}%
\pgfsetfillcolor{currentfill}%
\pgfsetlinewidth{0.401500pt}%
\definecolor{currentstroke}{rgb}{0.000000,0.000000,0.000000}%
\pgfsetstrokecolor{currentstroke}%
\pgfsetdash{}{0pt}%
\pgfsys@defobject{currentmarker}{\pgfqpoint{0.000000in}{-0.020833in}}{\pgfqpoint{0.000000in}{0.000000in}}{%
\pgfpathmoveto{\pgfqpoint{0.000000in}{0.000000in}}%
\pgfpathlineto{\pgfqpoint{0.000000in}{-0.020833in}}%
\pgfusepath{stroke,fill}%
}%
\begin{pgfscope}%
\pgfsys@transformshift{1.155831in}{2.414488in}%
\pgfsys@useobject{currentmarker}{}%
\end{pgfscope}%
\end{pgfscope}%
\begin{pgfscope}%
\pgfsetbuttcap%
\pgfsetroundjoin%
\definecolor{currentfill}{rgb}{0.000000,0.000000,0.000000}%
\pgfsetfillcolor{currentfill}%
\pgfsetlinewidth{0.401500pt}%
\definecolor{currentstroke}{rgb}{0.000000,0.000000,0.000000}%
\pgfsetstrokecolor{currentstroke}%
\pgfsetdash{}{0pt}%
\pgfsys@defobject{currentmarker}{\pgfqpoint{0.000000in}{0.000000in}}{\pgfqpoint{0.000000in}{0.020833in}}{%
\pgfpathmoveto{\pgfqpoint{0.000000in}{0.000000in}}%
\pgfpathlineto{\pgfqpoint{0.000000in}{0.020833in}}%
\pgfusepath{stroke,fill}%
}%
\begin{pgfscope}%
\pgfsys@transformshift{1.324441in}{0.420000in}%
\pgfsys@useobject{currentmarker}{}%
\end{pgfscope}%
\end{pgfscope}%
\begin{pgfscope}%
\pgfsetbuttcap%
\pgfsetroundjoin%
\definecolor{currentfill}{rgb}{0.000000,0.000000,0.000000}%
\pgfsetfillcolor{currentfill}%
\pgfsetlinewidth{0.401500pt}%
\definecolor{currentstroke}{rgb}{0.000000,0.000000,0.000000}%
\pgfsetstrokecolor{currentstroke}%
\pgfsetdash{}{0pt}%
\pgfsys@defobject{currentmarker}{\pgfqpoint{0.000000in}{-0.020833in}}{\pgfqpoint{0.000000in}{0.000000in}}{%
\pgfpathmoveto{\pgfqpoint{0.000000in}{0.000000in}}%
\pgfpathlineto{\pgfqpoint{0.000000in}{-0.020833in}}%
\pgfusepath{stroke,fill}%
}%
\begin{pgfscope}%
\pgfsys@transformshift{1.324441in}{2.414488in}%
\pgfsys@useobject{currentmarker}{}%
\end{pgfscope}%
\end{pgfscope}%
\begin{pgfscope}%
\pgfsetbuttcap%
\pgfsetroundjoin%
\definecolor{currentfill}{rgb}{0.000000,0.000000,0.000000}%
\pgfsetfillcolor{currentfill}%
\pgfsetlinewidth{0.401500pt}%
\definecolor{currentstroke}{rgb}{0.000000,0.000000,0.000000}%
\pgfsetstrokecolor{currentstroke}%
\pgfsetdash{}{0pt}%
\pgfsys@defobject{currentmarker}{\pgfqpoint{0.000000in}{0.000000in}}{\pgfqpoint{0.000000in}{0.020833in}}{%
\pgfpathmoveto{\pgfqpoint{0.000000in}{0.000000in}}%
\pgfpathlineto{\pgfqpoint{0.000000in}{0.020833in}}%
\pgfusepath{stroke,fill}%
}%
\begin{pgfscope}%
\pgfsys@transformshift{1.661662in}{0.420000in}%
\pgfsys@useobject{currentmarker}{}%
\end{pgfscope}%
\end{pgfscope}%
\begin{pgfscope}%
\pgfsetbuttcap%
\pgfsetroundjoin%
\definecolor{currentfill}{rgb}{0.000000,0.000000,0.000000}%
\pgfsetfillcolor{currentfill}%
\pgfsetlinewidth{0.401500pt}%
\definecolor{currentstroke}{rgb}{0.000000,0.000000,0.000000}%
\pgfsetstrokecolor{currentstroke}%
\pgfsetdash{}{0pt}%
\pgfsys@defobject{currentmarker}{\pgfqpoint{0.000000in}{-0.020833in}}{\pgfqpoint{0.000000in}{0.000000in}}{%
\pgfpathmoveto{\pgfqpoint{0.000000in}{0.000000in}}%
\pgfpathlineto{\pgfqpoint{0.000000in}{-0.020833in}}%
\pgfusepath{stroke,fill}%
}%
\begin{pgfscope}%
\pgfsys@transformshift{1.661662in}{2.414488in}%
\pgfsys@useobject{currentmarker}{}%
\end{pgfscope}%
\end{pgfscope}%
\begin{pgfscope}%
\pgfsetbuttcap%
\pgfsetroundjoin%
\definecolor{currentfill}{rgb}{0.000000,0.000000,0.000000}%
\pgfsetfillcolor{currentfill}%
\pgfsetlinewidth{0.401500pt}%
\definecolor{currentstroke}{rgb}{0.000000,0.000000,0.000000}%
\pgfsetstrokecolor{currentstroke}%
\pgfsetdash{}{0pt}%
\pgfsys@defobject{currentmarker}{\pgfqpoint{0.000000in}{0.000000in}}{\pgfqpoint{0.000000in}{0.020833in}}{%
\pgfpathmoveto{\pgfqpoint{0.000000in}{0.000000in}}%
\pgfpathlineto{\pgfqpoint{0.000000in}{0.020833in}}%
\pgfusepath{stroke,fill}%
}%
\begin{pgfscope}%
\pgfsys@transformshift{1.830272in}{0.420000in}%
\pgfsys@useobject{currentmarker}{}%
\end{pgfscope}%
\end{pgfscope}%
\begin{pgfscope}%
\pgfsetbuttcap%
\pgfsetroundjoin%
\definecolor{currentfill}{rgb}{0.000000,0.000000,0.000000}%
\pgfsetfillcolor{currentfill}%
\pgfsetlinewidth{0.401500pt}%
\definecolor{currentstroke}{rgb}{0.000000,0.000000,0.000000}%
\pgfsetstrokecolor{currentstroke}%
\pgfsetdash{}{0pt}%
\pgfsys@defobject{currentmarker}{\pgfqpoint{0.000000in}{-0.020833in}}{\pgfqpoint{0.000000in}{0.000000in}}{%
\pgfpathmoveto{\pgfqpoint{0.000000in}{0.000000in}}%
\pgfpathlineto{\pgfqpoint{0.000000in}{-0.020833in}}%
\pgfusepath{stroke,fill}%
}%
\begin{pgfscope}%
\pgfsys@transformshift{1.830272in}{2.414488in}%
\pgfsys@useobject{currentmarker}{}%
\end{pgfscope}%
\end{pgfscope}%
\begin{pgfscope}%
\pgfsetbuttcap%
\pgfsetroundjoin%
\definecolor{currentfill}{rgb}{0.000000,0.000000,0.000000}%
\pgfsetfillcolor{currentfill}%
\pgfsetlinewidth{0.401500pt}%
\definecolor{currentstroke}{rgb}{0.000000,0.000000,0.000000}%
\pgfsetstrokecolor{currentstroke}%
\pgfsetdash{}{0pt}%
\pgfsys@defobject{currentmarker}{\pgfqpoint{0.000000in}{0.000000in}}{\pgfqpoint{0.000000in}{0.020833in}}{%
\pgfpathmoveto{\pgfqpoint{0.000000in}{0.000000in}}%
\pgfpathlineto{\pgfqpoint{0.000000in}{0.020833in}}%
\pgfusepath{stroke,fill}%
}%
\begin{pgfscope}%
\pgfsys@transformshift{1.998883in}{0.420000in}%
\pgfsys@useobject{currentmarker}{}%
\end{pgfscope}%
\end{pgfscope}%
\begin{pgfscope}%
\pgfsetbuttcap%
\pgfsetroundjoin%
\definecolor{currentfill}{rgb}{0.000000,0.000000,0.000000}%
\pgfsetfillcolor{currentfill}%
\pgfsetlinewidth{0.401500pt}%
\definecolor{currentstroke}{rgb}{0.000000,0.000000,0.000000}%
\pgfsetstrokecolor{currentstroke}%
\pgfsetdash{}{0pt}%
\pgfsys@defobject{currentmarker}{\pgfqpoint{0.000000in}{-0.020833in}}{\pgfqpoint{0.000000in}{0.000000in}}{%
\pgfpathmoveto{\pgfqpoint{0.000000in}{0.000000in}}%
\pgfpathlineto{\pgfqpoint{0.000000in}{-0.020833in}}%
\pgfusepath{stroke,fill}%
}%
\begin{pgfscope}%
\pgfsys@transformshift{1.998883in}{2.414488in}%
\pgfsys@useobject{currentmarker}{}%
\end{pgfscope}%
\end{pgfscope}%
\begin{pgfscope}%
\pgfsetbuttcap%
\pgfsetroundjoin%
\definecolor{currentfill}{rgb}{0.000000,0.000000,0.000000}%
\pgfsetfillcolor{currentfill}%
\pgfsetlinewidth{0.401500pt}%
\definecolor{currentstroke}{rgb}{0.000000,0.000000,0.000000}%
\pgfsetstrokecolor{currentstroke}%
\pgfsetdash{}{0pt}%
\pgfsys@defobject{currentmarker}{\pgfqpoint{0.000000in}{0.000000in}}{\pgfqpoint{0.000000in}{0.020833in}}{%
\pgfpathmoveto{\pgfqpoint{0.000000in}{0.000000in}}%
\pgfpathlineto{\pgfqpoint{0.000000in}{0.020833in}}%
\pgfusepath{stroke,fill}%
}%
\begin{pgfscope}%
\pgfsys@transformshift{2.167493in}{0.420000in}%
\pgfsys@useobject{currentmarker}{}%
\end{pgfscope}%
\end{pgfscope}%
\begin{pgfscope}%
\pgfsetbuttcap%
\pgfsetroundjoin%
\definecolor{currentfill}{rgb}{0.000000,0.000000,0.000000}%
\pgfsetfillcolor{currentfill}%
\pgfsetlinewidth{0.401500pt}%
\definecolor{currentstroke}{rgb}{0.000000,0.000000,0.000000}%
\pgfsetstrokecolor{currentstroke}%
\pgfsetdash{}{0pt}%
\pgfsys@defobject{currentmarker}{\pgfqpoint{0.000000in}{-0.020833in}}{\pgfqpoint{0.000000in}{0.000000in}}{%
\pgfpathmoveto{\pgfqpoint{0.000000in}{0.000000in}}%
\pgfpathlineto{\pgfqpoint{0.000000in}{-0.020833in}}%
\pgfusepath{stroke,fill}%
}%
\begin{pgfscope}%
\pgfsys@transformshift{2.167493in}{2.414488in}%
\pgfsys@useobject{currentmarker}{}%
\end{pgfscope}%
\end{pgfscope}%
\begin{pgfscope}%
\pgfsetbuttcap%
\pgfsetroundjoin%
\definecolor{currentfill}{rgb}{0.000000,0.000000,0.000000}%
\pgfsetfillcolor{currentfill}%
\pgfsetlinewidth{0.401500pt}%
\definecolor{currentstroke}{rgb}{0.000000,0.000000,0.000000}%
\pgfsetstrokecolor{currentstroke}%
\pgfsetdash{}{0pt}%
\pgfsys@defobject{currentmarker}{\pgfqpoint{0.000000in}{0.000000in}}{\pgfqpoint{0.000000in}{0.020833in}}{%
\pgfpathmoveto{\pgfqpoint{0.000000in}{0.000000in}}%
\pgfpathlineto{\pgfqpoint{0.000000in}{0.020833in}}%
\pgfusepath{stroke,fill}%
}%
\begin{pgfscope}%
\pgfsys@transformshift{2.504714in}{0.420000in}%
\pgfsys@useobject{currentmarker}{}%
\end{pgfscope}%
\end{pgfscope}%
\begin{pgfscope}%
\pgfsetbuttcap%
\pgfsetroundjoin%
\definecolor{currentfill}{rgb}{0.000000,0.000000,0.000000}%
\pgfsetfillcolor{currentfill}%
\pgfsetlinewidth{0.401500pt}%
\definecolor{currentstroke}{rgb}{0.000000,0.000000,0.000000}%
\pgfsetstrokecolor{currentstroke}%
\pgfsetdash{}{0pt}%
\pgfsys@defobject{currentmarker}{\pgfqpoint{0.000000in}{-0.020833in}}{\pgfqpoint{0.000000in}{0.000000in}}{%
\pgfpathmoveto{\pgfqpoint{0.000000in}{0.000000in}}%
\pgfpathlineto{\pgfqpoint{0.000000in}{-0.020833in}}%
\pgfusepath{stroke,fill}%
}%
\begin{pgfscope}%
\pgfsys@transformshift{2.504714in}{2.414488in}%
\pgfsys@useobject{currentmarker}{}%
\end{pgfscope}%
\end{pgfscope}%
\begin{pgfscope}%
\pgfsetbuttcap%
\pgfsetroundjoin%
\definecolor{currentfill}{rgb}{0.000000,0.000000,0.000000}%
\pgfsetfillcolor{currentfill}%
\pgfsetlinewidth{0.401500pt}%
\definecolor{currentstroke}{rgb}{0.000000,0.000000,0.000000}%
\pgfsetstrokecolor{currentstroke}%
\pgfsetdash{}{0pt}%
\pgfsys@defobject{currentmarker}{\pgfqpoint{0.000000in}{0.000000in}}{\pgfqpoint{0.000000in}{0.020833in}}{%
\pgfpathmoveto{\pgfqpoint{0.000000in}{0.000000in}}%
\pgfpathlineto{\pgfqpoint{0.000000in}{0.020833in}}%
\pgfusepath{stroke,fill}%
}%
\begin{pgfscope}%
\pgfsys@transformshift{2.673324in}{0.420000in}%
\pgfsys@useobject{currentmarker}{}%
\end{pgfscope}%
\end{pgfscope}%
\begin{pgfscope}%
\pgfsetbuttcap%
\pgfsetroundjoin%
\definecolor{currentfill}{rgb}{0.000000,0.000000,0.000000}%
\pgfsetfillcolor{currentfill}%
\pgfsetlinewidth{0.401500pt}%
\definecolor{currentstroke}{rgb}{0.000000,0.000000,0.000000}%
\pgfsetstrokecolor{currentstroke}%
\pgfsetdash{}{0pt}%
\pgfsys@defobject{currentmarker}{\pgfqpoint{0.000000in}{-0.020833in}}{\pgfqpoint{0.000000in}{0.000000in}}{%
\pgfpathmoveto{\pgfqpoint{0.000000in}{0.000000in}}%
\pgfpathlineto{\pgfqpoint{0.000000in}{-0.020833in}}%
\pgfusepath{stroke,fill}%
}%
\begin{pgfscope}%
\pgfsys@transformshift{2.673324in}{2.414488in}%
\pgfsys@useobject{currentmarker}{}%
\end{pgfscope}%
\end{pgfscope}%
\begin{pgfscope}%
\pgfsetbuttcap%
\pgfsetroundjoin%
\definecolor{currentfill}{rgb}{0.000000,0.000000,0.000000}%
\pgfsetfillcolor{currentfill}%
\pgfsetlinewidth{0.401500pt}%
\definecolor{currentstroke}{rgb}{0.000000,0.000000,0.000000}%
\pgfsetstrokecolor{currentstroke}%
\pgfsetdash{}{0pt}%
\pgfsys@defobject{currentmarker}{\pgfqpoint{0.000000in}{0.000000in}}{\pgfqpoint{0.000000in}{0.020833in}}{%
\pgfpathmoveto{\pgfqpoint{0.000000in}{0.000000in}}%
\pgfpathlineto{\pgfqpoint{0.000000in}{0.020833in}}%
\pgfusepath{stroke,fill}%
}%
\begin{pgfscope}%
\pgfsys@transformshift{2.841934in}{0.420000in}%
\pgfsys@useobject{currentmarker}{}%
\end{pgfscope}%
\end{pgfscope}%
\begin{pgfscope}%
\pgfsetbuttcap%
\pgfsetroundjoin%
\definecolor{currentfill}{rgb}{0.000000,0.000000,0.000000}%
\pgfsetfillcolor{currentfill}%
\pgfsetlinewidth{0.401500pt}%
\definecolor{currentstroke}{rgb}{0.000000,0.000000,0.000000}%
\pgfsetstrokecolor{currentstroke}%
\pgfsetdash{}{0pt}%
\pgfsys@defobject{currentmarker}{\pgfqpoint{0.000000in}{-0.020833in}}{\pgfqpoint{0.000000in}{0.000000in}}{%
\pgfpathmoveto{\pgfqpoint{0.000000in}{0.000000in}}%
\pgfpathlineto{\pgfqpoint{0.000000in}{-0.020833in}}%
\pgfusepath{stroke,fill}%
}%
\begin{pgfscope}%
\pgfsys@transformshift{2.841934in}{2.414488in}%
\pgfsys@useobject{currentmarker}{}%
\end{pgfscope}%
\end{pgfscope}%
\begin{pgfscope}%
\pgfsetbuttcap%
\pgfsetroundjoin%
\definecolor{currentfill}{rgb}{0.000000,0.000000,0.000000}%
\pgfsetfillcolor{currentfill}%
\pgfsetlinewidth{0.401500pt}%
\definecolor{currentstroke}{rgb}{0.000000,0.000000,0.000000}%
\pgfsetstrokecolor{currentstroke}%
\pgfsetdash{}{0pt}%
\pgfsys@defobject{currentmarker}{\pgfqpoint{0.000000in}{0.000000in}}{\pgfqpoint{0.000000in}{0.020833in}}{%
\pgfpathmoveto{\pgfqpoint{0.000000in}{0.000000in}}%
\pgfpathlineto{\pgfqpoint{0.000000in}{0.020833in}}%
\pgfusepath{stroke,fill}%
}%
\begin{pgfscope}%
\pgfsys@transformshift{3.010545in}{0.420000in}%
\pgfsys@useobject{currentmarker}{}%
\end{pgfscope}%
\end{pgfscope}%
\begin{pgfscope}%
\pgfsetbuttcap%
\pgfsetroundjoin%
\definecolor{currentfill}{rgb}{0.000000,0.000000,0.000000}%
\pgfsetfillcolor{currentfill}%
\pgfsetlinewidth{0.401500pt}%
\definecolor{currentstroke}{rgb}{0.000000,0.000000,0.000000}%
\pgfsetstrokecolor{currentstroke}%
\pgfsetdash{}{0pt}%
\pgfsys@defobject{currentmarker}{\pgfqpoint{0.000000in}{-0.020833in}}{\pgfqpoint{0.000000in}{0.000000in}}{%
\pgfpathmoveto{\pgfqpoint{0.000000in}{0.000000in}}%
\pgfpathlineto{\pgfqpoint{0.000000in}{-0.020833in}}%
\pgfusepath{stroke,fill}%
}%
\begin{pgfscope}%
\pgfsys@transformshift{3.010545in}{2.414488in}%
\pgfsys@useobject{currentmarker}{}%
\end{pgfscope}%
\end{pgfscope}%
\begin{pgfscope}%
\definecolor{textcolor}{rgb}{0.000000,0.000000,0.000000}%
\pgfsetstrokecolor{textcolor}%
\pgfsetfillcolor{textcolor}%
\pgftext[x=1.844055in,y=0.208426in,,top]{\color{textcolor}{\rmfamily\fontsize{12.000000}{14.400000}\selectfont\catcode`\^=\active\def^{\ifmmode\sp\else\^{}\fi}\catcode`\%=\active\def%{\%}$s/\delta_{0}$}}%
\end{pgfscope}%
\begin{pgfscope}%
\pgfsetbuttcap%
\pgfsetroundjoin%
\definecolor{currentfill}{rgb}{0.000000,0.000000,0.000000}%
\pgfsetfillcolor{currentfill}%
\pgfsetlinewidth{0.501875pt}%
\definecolor{currentstroke}{rgb}{0.000000,0.000000,0.000000}%
\pgfsetstrokecolor{currentstroke}%
\pgfsetdash{}{0pt}%
\pgfsys@defobject{currentmarker}{\pgfqpoint{0.000000in}{0.000000in}}{\pgfqpoint{0.041667in}{0.000000in}}{%
\pgfpathmoveto{\pgfqpoint{0.000000in}{0.000000in}}%
\pgfpathlineto{\pgfqpoint{0.041667in}{0.000000in}}%
\pgfusepath{stroke,fill}%
}%
\begin{pgfscope}%
\pgfsys@transformshift{0.650000in}{0.420000in}%
\pgfsys@useobject{currentmarker}{}%
\end{pgfscope}%
\end{pgfscope}%
\begin{pgfscope}%
\pgfsetbuttcap%
\pgfsetroundjoin%
\definecolor{currentfill}{rgb}{0.000000,0.000000,0.000000}%
\pgfsetfillcolor{currentfill}%
\pgfsetlinewidth{0.501875pt}%
\definecolor{currentstroke}{rgb}{0.000000,0.000000,0.000000}%
\pgfsetstrokecolor{currentstroke}%
\pgfsetdash{}{0pt}%
\pgfsys@defobject{currentmarker}{\pgfqpoint{-0.041667in}{0.000000in}}{\pgfqpoint{-0.000000in}{0.000000in}}{%
\pgfpathmoveto{\pgfqpoint{-0.000000in}{0.000000in}}%
\pgfpathlineto{\pgfqpoint{-0.041667in}{0.000000in}}%
\pgfusepath{stroke,fill}%
}%
\begin{pgfscope}%
\pgfsys@transformshift{3.038110in}{0.420000in}%
\pgfsys@useobject{currentmarker}{}%
\end{pgfscope}%
\end{pgfscope}%
\begin{pgfscope}%
\definecolor{textcolor}{rgb}{0.000000,0.000000,0.000000}%
\pgfsetstrokecolor{textcolor}%
\pgfsetfillcolor{textcolor}%
\pgftext[x=0.407060in, y=0.367193in, left, base]{\color{textcolor}{\rmfamily\fontsize{11.000000}{13.200000}\selectfont\catcode`\^=\active\def^{\ifmmode\sp\else\^{}\fi}\catcode`\%=\active\def%{\%}\ensuremath{-}2}}%
\end{pgfscope}%
\begin{pgfscope}%
\pgfsetbuttcap%
\pgfsetroundjoin%
\definecolor{currentfill}{rgb}{0.000000,0.000000,0.000000}%
\pgfsetfillcolor{currentfill}%
\pgfsetlinewidth{0.501875pt}%
\definecolor{currentstroke}{rgb}{0.000000,0.000000,0.000000}%
\pgfsetstrokecolor{currentstroke}%
\pgfsetdash{}{0pt}%
\pgfsys@defobject{currentmarker}{\pgfqpoint{0.000000in}{0.000000in}}{\pgfqpoint{0.041667in}{0.000000in}}{%
\pgfpathmoveto{\pgfqpoint{0.000000in}{0.000000in}}%
\pgfpathlineto{\pgfqpoint{0.041667in}{0.000000in}}%
\pgfusepath{stroke,fill}%
}%
\begin{pgfscope}%
\pgfsys@transformshift{0.650000in}{0.918622in}%
\pgfsys@useobject{currentmarker}{}%
\end{pgfscope}%
\end{pgfscope}%
\begin{pgfscope}%
\pgfsetbuttcap%
\pgfsetroundjoin%
\definecolor{currentfill}{rgb}{0.000000,0.000000,0.000000}%
\pgfsetfillcolor{currentfill}%
\pgfsetlinewidth{0.501875pt}%
\definecolor{currentstroke}{rgb}{0.000000,0.000000,0.000000}%
\pgfsetstrokecolor{currentstroke}%
\pgfsetdash{}{0pt}%
\pgfsys@defobject{currentmarker}{\pgfqpoint{-0.041667in}{0.000000in}}{\pgfqpoint{-0.000000in}{0.000000in}}{%
\pgfpathmoveto{\pgfqpoint{-0.000000in}{0.000000in}}%
\pgfpathlineto{\pgfqpoint{-0.041667in}{0.000000in}}%
\pgfusepath{stroke,fill}%
}%
\begin{pgfscope}%
\pgfsys@transformshift{3.038110in}{0.918622in}%
\pgfsys@useobject{currentmarker}{}%
\end{pgfscope}%
\end{pgfscope}%
\begin{pgfscope}%
\definecolor{textcolor}{rgb}{0.000000,0.000000,0.000000}%
\pgfsetstrokecolor{textcolor}%
\pgfsetfillcolor{textcolor}%
\pgftext[x=0.407060in, y=0.865815in, left, base]{\color{textcolor}{\rmfamily\fontsize{11.000000}{13.200000}\selectfont\catcode`\^=\active\def^{\ifmmode\sp\else\^{}\fi}\catcode`\%=\active\def%{\%}\ensuremath{-}1}}%
\end{pgfscope}%
\begin{pgfscope}%
\pgfsetbuttcap%
\pgfsetroundjoin%
\definecolor{currentfill}{rgb}{0.000000,0.000000,0.000000}%
\pgfsetfillcolor{currentfill}%
\pgfsetlinewidth{0.501875pt}%
\definecolor{currentstroke}{rgb}{0.000000,0.000000,0.000000}%
\pgfsetstrokecolor{currentstroke}%
\pgfsetdash{}{0pt}%
\pgfsys@defobject{currentmarker}{\pgfqpoint{0.000000in}{0.000000in}}{\pgfqpoint{0.041667in}{0.000000in}}{%
\pgfpathmoveto{\pgfqpoint{0.000000in}{0.000000in}}%
\pgfpathlineto{\pgfqpoint{0.041667in}{0.000000in}}%
\pgfusepath{stroke,fill}%
}%
\begin{pgfscope}%
\pgfsys@transformshift{0.650000in}{1.417244in}%
\pgfsys@useobject{currentmarker}{}%
\end{pgfscope}%
\end{pgfscope}%
\begin{pgfscope}%
\pgfsetbuttcap%
\pgfsetroundjoin%
\definecolor{currentfill}{rgb}{0.000000,0.000000,0.000000}%
\pgfsetfillcolor{currentfill}%
\pgfsetlinewidth{0.501875pt}%
\definecolor{currentstroke}{rgb}{0.000000,0.000000,0.000000}%
\pgfsetstrokecolor{currentstroke}%
\pgfsetdash{}{0pt}%
\pgfsys@defobject{currentmarker}{\pgfqpoint{-0.041667in}{0.000000in}}{\pgfqpoint{-0.000000in}{0.000000in}}{%
\pgfpathmoveto{\pgfqpoint{-0.000000in}{0.000000in}}%
\pgfpathlineto{\pgfqpoint{-0.041667in}{0.000000in}}%
\pgfusepath{stroke,fill}%
}%
\begin{pgfscope}%
\pgfsys@transformshift{3.038110in}{1.417244in}%
\pgfsys@useobject{currentmarker}{}%
\end{pgfscope}%
\end{pgfscope}%
\begin{pgfscope}%
\definecolor{textcolor}{rgb}{0.000000,0.000000,0.000000}%
\pgfsetstrokecolor{textcolor}%
\pgfsetfillcolor{textcolor}%
\pgftext[x=0.525347in, y=1.364437in, left, base]{\color{textcolor}{\rmfamily\fontsize{11.000000}{13.200000}\selectfont\catcode`\^=\active\def^{\ifmmode\sp\else\^{}\fi}\catcode`\%=\active\def%{\%}0}}%
\end{pgfscope}%
\begin{pgfscope}%
\pgfsetbuttcap%
\pgfsetroundjoin%
\definecolor{currentfill}{rgb}{0.000000,0.000000,0.000000}%
\pgfsetfillcolor{currentfill}%
\pgfsetlinewidth{0.501875pt}%
\definecolor{currentstroke}{rgb}{0.000000,0.000000,0.000000}%
\pgfsetstrokecolor{currentstroke}%
\pgfsetdash{}{0pt}%
\pgfsys@defobject{currentmarker}{\pgfqpoint{0.000000in}{0.000000in}}{\pgfqpoint{0.041667in}{0.000000in}}{%
\pgfpathmoveto{\pgfqpoint{0.000000in}{0.000000in}}%
\pgfpathlineto{\pgfqpoint{0.041667in}{0.000000in}}%
\pgfusepath{stroke,fill}%
}%
\begin{pgfscope}%
\pgfsys@transformshift{0.650000in}{1.915866in}%
\pgfsys@useobject{currentmarker}{}%
\end{pgfscope}%
\end{pgfscope}%
\begin{pgfscope}%
\pgfsetbuttcap%
\pgfsetroundjoin%
\definecolor{currentfill}{rgb}{0.000000,0.000000,0.000000}%
\pgfsetfillcolor{currentfill}%
\pgfsetlinewidth{0.501875pt}%
\definecolor{currentstroke}{rgb}{0.000000,0.000000,0.000000}%
\pgfsetstrokecolor{currentstroke}%
\pgfsetdash{}{0pt}%
\pgfsys@defobject{currentmarker}{\pgfqpoint{-0.041667in}{0.000000in}}{\pgfqpoint{-0.000000in}{0.000000in}}{%
\pgfpathmoveto{\pgfqpoint{-0.000000in}{0.000000in}}%
\pgfpathlineto{\pgfqpoint{-0.041667in}{0.000000in}}%
\pgfusepath{stroke,fill}%
}%
\begin{pgfscope}%
\pgfsys@transformshift{3.038110in}{1.915866in}%
\pgfsys@useobject{currentmarker}{}%
\end{pgfscope}%
\end{pgfscope}%
\begin{pgfscope}%
\definecolor{textcolor}{rgb}{0.000000,0.000000,0.000000}%
\pgfsetstrokecolor{textcolor}%
\pgfsetfillcolor{textcolor}%
\pgftext[x=0.525347in, y=1.863059in, left, base]{\color{textcolor}{\rmfamily\fontsize{11.000000}{13.200000}\selectfont\catcode`\^=\active\def^{\ifmmode\sp\else\^{}\fi}\catcode`\%=\active\def%{\%}1}}%
\end{pgfscope}%
\begin{pgfscope}%
\pgfsetbuttcap%
\pgfsetroundjoin%
\definecolor{currentfill}{rgb}{0.000000,0.000000,0.000000}%
\pgfsetfillcolor{currentfill}%
\pgfsetlinewidth{0.501875pt}%
\definecolor{currentstroke}{rgb}{0.000000,0.000000,0.000000}%
\pgfsetstrokecolor{currentstroke}%
\pgfsetdash{}{0pt}%
\pgfsys@defobject{currentmarker}{\pgfqpoint{0.000000in}{0.000000in}}{\pgfqpoint{0.041667in}{0.000000in}}{%
\pgfpathmoveto{\pgfqpoint{0.000000in}{0.000000in}}%
\pgfpathlineto{\pgfqpoint{0.041667in}{0.000000in}}%
\pgfusepath{stroke,fill}%
}%
\begin{pgfscope}%
\pgfsys@transformshift{0.650000in}{2.414488in}%
\pgfsys@useobject{currentmarker}{}%
\end{pgfscope}%
\end{pgfscope}%
\begin{pgfscope}%
\pgfsetbuttcap%
\pgfsetroundjoin%
\definecolor{currentfill}{rgb}{0.000000,0.000000,0.000000}%
\pgfsetfillcolor{currentfill}%
\pgfsetlinewidth{0.501875pt}%
\definecolor{currentstroke}{rgb}{0.000000,0.000000,0.000000}%
\pgfsetstrokecolor{currentstroke}%
\pgfsetdash{}{0pt}%
\pgfsys@defobject{currentmarker}{\pgfqpoint{-0.041667in}{0.000000in}}{\pgfqpoint{-0.000000in}{0.000000in}}{%
\pgfpathmoveto{\pgfqpoint{-0.000000in}{0.000000in}}%
\pgfpathlineto{\pgfqpoint{-0.041667in}{0.000000in}}%
\pgfusepath{stroke,fill}%
}%
\begin{pgfscope}%
\pgfsys@transformshift{3.038110in}{2.414488in}%
\pgfsys@useobject{currentmarker}{}%
\end{pgfscope}%
\end{pgfscope}%
\begin{pgfscope}%
\definecolor{textcolor}{rgb}{0.000000,0.000000,0.000000}%
\pgfsetstrokecolor{textcolor}%
\pgfsetfillcolor{textcolor}%
\pgftext[x=0.525347in, y=2.361681in, left, base]{\color{textcolor}{\rmfamily\fontsize{11.000000}{13.200000}\selectfont\catcode`\^=\active\def^{\ifmmode\sp\else\^{}\fi}\catcode`\%=\active\def%{\%}2}}%
\end{pgfscope}%
\begin{pgfscope}%
\pgfsetbuttcap%
\pgfsetroundjoin%
\definecolor{currentfill}{rgb}{0.000000,0.000000,0.000000}%
\pgfsetfillcolor{currentfill}%
\pgfsetlinewidth{0.401500pt}%
\definecolor{currentstroke}{rgb}{0.000000,0.000000,0.000000}%
\pgfsetstrokecolor{currentstroke}%
\pgfsetdash{}{0pt}%
\pgfsys@defobject{currentmarker}{\pgfqpoint{0.000000in}{0.000000in}}{\pgfqpoint{0.020833in}{0.000000in}}{%
\pgfpathmoveto{\pgfqpoint{0.000000in}{0.000000in}}%
\pgfpathlineto{\pgfqpoint{0.020833in}{0.000000in}}%
\pgfusepath{stroke,fill}%
}%
\begin{pgfscope}%
\pgfsys@transformshift{0.650000in}{0.519724in}%
\pgfsys@useobject{currentmarker}{}%
\end{pgfscope}%
\end{pgfscope}%
\begin{pgfscope}%
\pgfsetbuttcap%
\pgfsetroundjoin%
\definecolor{currentfill}{rgb}{0.000000,0.000000,0.000000}%
\pgfsetfillcolor{currentfill}%
\pgfsetlinewidth{0.401500pt}%
\definecolor{currentstroke}{rgb}{0.000000,0.000000,0.000000}%
\pgfsetstrokecolor{currentstroke}%
\pgfsetdash{}{0pt}%
\pgfsys@defobject{currentmarker}{\pgfqpoint{-0.020833in}{0.000000in}}{\pgfqpoint{-0.000000in}{0.000000in}}{%
\pgfpathmoveto{\pgfqpoint{-0.000000in}{0.000000in}}%
\pgfpathlineto{\pgfqpoint{-0.020833in}{0.000000in}}%
\pgfusepath{stroke,fill}%
}%
\begin{pgfscope}%
\pgfsys@transformshift{3.038110in}{0.519724in}%
\pgfsys@useobject{currentmarker}{}%
\end{pgfscope}%
\end{pgfscope}%
\begin{pgfscope}%
\pgfsetbuttcap%
\pgfsetroundjoin%
\definecolor{currentfill}{rgb}{0.000000,0.000000,0.000000}%
\pgfsetfillcolor{currentfill}%
\pgfsetlinewidth{0.401500pt}%
\definecolor{currentstroke}{rgb}{0.000000,0.000000,0.000000}%
\pgfsetstrokecolor{currentstroke}%
\pgfsetdash{}{0pt}%
\pgfsys@defobject{currentmarker}{\pgfqpoint{0.000000in}{0.000000in}}{\pgfqpoint{0.020833in}{0.000000in}}{%
\pgfpathmoveto{\pgfqpoint{0.000000in}{0.000000in}}%
\pgfpathlineto{\pgfqpoint{0.020833in}{0.000000in}}%
\pgfusepath{stroke,fill}%
}%
\begin{pgfscope}%
\pgfsys@transformshift{0.650000in}{0.619449in}%
\pgfsys@useobject{currentmarker}{}%
\end{pgfscope}%
\end{pgfscope}%
\begin{pgfscope}%
\pgfsetbuttcap%
\pgfsetroundjoin%
\definecolor{currentfill}{rgb}{0.000000,0.000000,0.000000}%
\pgfsetfillcolor{currentfill}%
\pgfsetlinewidth{0.401500pt}%
\definecolor{currentstroke}{rgb}{0.000000,0.000000,0.000000}%
\pgfsetstrokecolor{currentstroke}%
\pgfsetdash{}{0pt}%
\pgfsys@defobject{currentmarker}{\pgfqpoint{-0.020833in}{0.000000in}}{\pgfqpoint{-0.000000in}{0.000000in}}{%
\pgfpathmoveto{\pgfqpoint{-0.000000in}{0.000000in}}%
\pgfpathlineto{\pgfqpoint{-0.020833in}{0.000000in}}%
\pgfusepath{stroke,fill}%
}%
\begin{pgfscope}%
\pgfsys@transformshift{3.038110in}{0.619449in}%
\pgfsys@useobject{currentmarker}{}%
\end{pgfscope}%
\end{pgfscope}%
\begin{pgfscope}%
\pgfsetbuttcap%
\pgfsetroundjoin%
\definecolor{currentfill}{rgb}{0.000000,0.000000,0.000000}%
\pgfsetfillcolor{currentfill}%
\pgfsetlinewidth{0.401500pt}%
\definecolor{currentstroke}{rgb}{0.000000,0.000000,0.000000}%
\pgfsetstrokecolor{currentstroke}%
\pgfsetdash{}{0pt}%
\pgfsys@defobject{currentmarker}{\pgfqpoint{0.000000in}{0.000000in}}{\pgfqpoint{0.020833in}{0.000000in}}{%
\pgfpathmoveto{\pgfqpoint{0.000000in}{0.000000in}}%
\pgfpathlineto{\pgfqpoint{0.020833in}{0.000000in}}%
\pgfusepath{stroke,fill}%
}%
\begin{pgfscope}%
\pgfsys@transformshift{0.650000in}{0.719173in}%
\pgfsys@useobject{currentmarker}{}%
\end{pgfscope}%
\end{pgfscope}%
\begin{pgfscope}%
\pgfsetbuttcap%
\pgfsetroundjoin%
\definecolor{currentfill}{rgb}{0.000000,0.000000,0.000000}%
\pgfsetfillcolor{currentfill}%
\pgfsetlinewidth{0.401500pt}%
\definecolor{currentstroke}{rgb}{0.000000,0.000000,0.000000}%
\pgfsetstrokecolor{currentstroke}%
\pgfsetdash{}{0pt}%
\pgfsys@defobject{currentmarker}{\pgfqpoint{-0.020833in}{0.000000in}}{\pgfqpoint{-0.000000in}{0.000000in}}{%
\pgfpathmoveto{\pgfqpoint{-0.000000in}{0.000000in}}%
\pgfpathlineto{\pgfqpoint{-0.020833in}{0.000000in}}%
\pgfusepath{stroke,fill}%
}%
\begin{pgfscope}%
\pgfsys@transformshift{3.038110in}{0.719173in}%
\pgfsys@useobject{currentmarker}{}%
\end{pgfscope}%
\end{pgfscope}%
\begin{pgfscope}%
\pgfsetbuttcap%
\pgfsetroundjoin%
\definecolor{currentfill}{rgb}{0.000000,0.000000,0.000000}%
\pgfsetfillcolor{currentfill}%
\pgfsetlinewidth{0.401500pt}%
\definecolor{currentstroke}{rgb}{0.000000,0.000000,0.000000}%
\pgfsetstrokecolor{currentstroke}%
\pgfsetdash{}{0pt}%
\pgfsys@defobject{currentmarker}{\pgfqpoint{0.000000in}{0.000000in}}{\pgfqpoint{0.020833in}{0.000000in}}{%
\pgfpathmoveto{\pgfqpoint{0.000000in}{0.000000in}}%
\pgfpathlineto{\pgfqpoint{0.020833in}{0.000000in}}%
\pgfusepath{stroke,fill}%
}%
\begin{pgfscope}%
\pgfsys@transformshift{0.650000in}{0.818898in}%
\pgfsys@useobject{currentmarker}{}%
\end{pgfscope}%
\end{pgfscope}%
\begin{pgfscope}%
\pgfsetbuttcap%
\pgfsetroundjoin%
\definecolor{currentfill}{rgb}{0.000000,0.000000,0.000000}%
\pgfsetfillcolor{currentfill}%
\pgfsetlinewidth{0.401500pt}%
\definecolor{currentstroke}{rgb}{0.000000,0.000000,0.000000}%
\pgfsetstrokecolor{currentstroke}%
\pgfsetdash{}{0pt}%
\pgfsys@defobject{currentmarker}{\pgfqpoint{-0.020833in}{0.000000in}}{\pgfqpoint{-0.000000in}{0.000000in}}{%
\pgfpathmoveto{\pgfqpoint{-0.000000in}{0.000000in}}%
\pgfpathlineto{\pgfqpoint{-0.020833in}{0.000000in}}%
\pgfusepath{stroke,fill}%
}%
\begin{pgfscope}%
\pgfsys@transformshift{3.038110in}{0.818898in}%
\pgfsys@useobject{currentmarker}{}%
\end{pgfscope}%
\end{pgfscope}%
\begin{pgfscope}%
\pgfsetbuttcap%
\pgfsetroundjoin%
\definecolor{currentfill}{rgb}{0.000000,0.000000,0.000000}%
\pgfsetfillcolor{currentfill}%
\pgfsetlinewidth{0.401500pt}%
\definecolor{currentstroke}{rgb}{0.000000,0.000000,0.000000}%
\pgfsetstrokecolor{currentstroke}%
\pgfsetdash{}{0pt}%
\pgfsys@defobject{currentmarker}{\pgfqpoint{0.000000in}{0.000000in}}{\pgfqpoint{0.020833in}{0.000000in}}{%
\pgfpathmoveto{\pgfqpoint{0.000000in}{0.000000in}}%
\pgfpathlineto{\pgfqpoint{0.020833in}{0.000000in}}%
\pgfusepath{stroke,fill}%
}%
\begin{pgfscope}%
\pgfsys@transformshift{0.650000in}{1.018346in}%
\pgfsys@useobject{currentmarker}{}%
\end{pgfscope}%
\end{pgfscope}%
\begin{pgfscope}%
\pgfsetbuttcap%
\pgfsetroundjoin%
\definecolor{currentfill}{rgb}{0.000000,0.000000,0.000000}%
\pgfsetfillcolor{currentfill}%
\pgfsetlinewidth{0.401500pt}%
\definecolor{currentstroke}{rgb}{0.000000,0.000000,0.000000}%
\pgfsetstrokecolor{currentstroke}%
\pgfsetdash{}{0pt}%
\pgfsys@defobject{currentmarker}{\pgfqpoint{-0.020833in}{0.000000in}}{\pgfqpoint{-0.000000in}{0.000000in}}{%
\pgfpathmoveto{\pgfqpoint{-0.000000in}{0.000000in}}%
\pgfpathlineto{\pgfqpoint{-0.020833in}{0.000000in}}%
\pgfusepath{stroke,fill}%
}%
\begin{pgfscope}%
\pgfsys@transformshift{3.038110in}{1.018346in}%
\pgfsys@useobject{currentmarker}{}%
\end{pgfscope}%
\end{pgfscope}%
\begin{pgfscope}%
\pgfsetbuttcap%
\pgfsetroundjoin%
\definecolor{currentfill}{rgb}{0.000000,0.000000,0.000000}%
\pgfsetfillcolor{currentfill}%
\pgfsetlinewidth{0.401500pt}%
\definecolor{currentstroke}{rgb}{0.000000,0.000000,0.000000}%
\pgfsetstrokecolor{currentstroke}%
\pgfsetdash{}{0pt}%
\pgfsys@defobject{currentmarker}{\pgfqpoint{0.000000in}{0.000000in}}{\pgfqpoint{0.020833in}{0.000000in}}{%
\pgfpathmoveto{\pgfqpoint{0.000000in}{0.000000in}}%
\pgfpathlineto{\pgfqpoint{0.020833in}{0.000000in}}%
\pgfusepath{stroke,fill}%
}%
\begin{pgfscope}%
\pgfsys@transformshift{0.650000in}{1.118071in}%
\pgfsys@useobject{currentmarker}{}%
\end{pgfscope}%
\end{pgfscope}%
\begin{pgfscope}%
\pgfsetbuttcap%
\pgfsetroundjoin%
\definecolor{currentfill}{rgb}{0.000000,0.000000,0.000000}%
\pgfsetfillcolor{currentfill}%
\pgfsetlinewidth{0.401500pt}%
\definecolor{currentstroke}{rgb}{0.000000,0.000000,0.000000}%
\pgfsetstrokecolor{currentstroke}%
\pgfsetdash{}{0pt}%
\pgfsys@defobject{currentmarker}{\pgfqpoint{-0.020833in}{0.000000in}}{\pgfqpoint{-0.000000in}{0.000000in}}{%
\pgfpathmoveto{\pgfqpoint{-0.000000in}{0.000000in}}%
\pgfpathlineto{\pgfqpoint{-0.020833in}{0.000000in}}%
\pgfusepath{stroke,fill}%
}%
\begin{pgfscope}%
\pgfsys@transformshift{3.038110in}{1.118071in}%
\pgfsys@useobject{currentmarker}{}%
\end{pgfscope}%
\end{pgfscope}%
\begin{pgfscope}%
\pgfsetbuttcap%
\pgfsetroundjoin%
\definecolor{currentfill}{rgb}{0.000000,0.000000,0.000000}%
\pgfsetfillcolor{currentfill}%
\pgfsetlinewidth{0.401500pt}%
\definecolor{currentstroke}{rgb}{0.000000,0.000000,0.000000}%
\pgfsetstrokecolor{currentstroke}%
\pgfsetdash{}{0pt}%
\pgfsys@defobject{currentmarker}{\pgfqpoint{0.000000in}{0.000000in}}{\pgfqpoint{0.020833in}{0.000000in}}{%
\pgfpathmoveto{\pgfqpoint{0.000000in}{0.000000in}}%
\pgfpathlineto{\pgfqpoint{0.020833in}{0.000000in}}%
\pgfusepath{stroke,fill}%
}%
\begin{pgfscope}%
\pgfsys@transformshift{0.650000in}{1.217795in}%
\pgfsys@useobject{currentmarker}{}%
\end{pgfscope}%
\end{pgfscope}%
\begin{pgfscope}%
\pgfsetbuttcap%
\pgfsetroundjoin%
\definecolor{currentfill}{rgb}{0.000000,0.000000,0.000000}%
\pgfsetfillcolor{currentfill}%
\pgfsetlinewidth{0.401500pt}%
\definecolor{currentstroke}{rgb}{0.000000,0.000000,0.000000}%
\pgfsetstrokecolor{currentstroke}%
\pgfsetdash{}{0pt}%
\pgfsys@defobject{currentmarker}{\pgfqpoint{-0.020833in}{0.000000in}}{\pgfqpoint{-0.000000in}{0.000000in}}{%
\pgfpathmoveto{\pgfqpoint{-0.000000in}{0.000000in}}%
\pgfpathlineto{\pgfqpoint{-0.020833in}{0.000000in}}%
\pgfusepath{stroke,fill}%
}%
\begin{pgfscope}%
\pgfsys@transformshift{3.038110in}{1.217795in}%
\pgfsys@useobject{currentmarker}{}%
\end{pgfscope}%
\end{pgfscope}%
\begin{pgfscope}%
\pgfsetbuttcap%
\pgfsetroundjoin%
\definecolor{currentfill}{rgb}{0.000000,0.000000,0.000000}%
\pgfsetfillcolor{currentfill}%
\pgfsetlinewidth{0.401500pt}%
\definecolor{currentstroke}{rgb}{0.000000,0.000000,0.000000}%
\pgfsetstrokecolor{currentstroke}%
\pgfsetdash{}{0pt}%
\pgfsys@defobject{currentmarker}{\pgfqpoint{0.000000in}{0.000000in}}{\pgfqpoint{0.020833in}{0.000000in}}{%
\pgfpathmoveto{\pgfqpoint{0.000000in}{0.000000in}}%
\pgfpathlineto{\pgfqpoint{0.020833in}{0.000000in}}%
\pgfusepath{stroke,fill}%
}%
\begin{pgfscope}%
\pgfsys@transformshift{0.650000in}{1.317520in}%
\pgfsys@useobject{currentmarker}{}%
\end{pgfscope}%
\end{pgfscope}%
\begin{pgfscope}%
\pgfsetbuttcap%
\pgfsetroundjoin%
\definecolor{currentfill}{rgb}{0.000000,0.000000,0.000000}%
\pgfsetfillcolor{currentfill}%
\pgfsetlinewidth{0.401500pt}%
\definecolor{currentstroke}{rgb}{0.000000,0.000000,0.000000}%
\pgfsetstrokecolor{currentstroke}%
\pgfsetdash{}{0pt}%
\pgfsys@defobject{currentmarker}{\pgfqpoint{-0.020833in}{0.000000in}}{\pgfqpoint{-0.000000in}{0.000000in}}{%
\pgfpathmoveto{\pgfqpoint{-0.000000in}{0.000000in}}%
\pgfpathlineto{\pgfqpoint{-0.020833in}{0.000000in}}%
\pgfusepath{stroke,fill}%
}%
\begin{pgfscope}%
\pgfsys@transformshift{3.038110in}{1.317520in}%
\pgfsys@useobject{currentmarker}{}%
\end{pgfscope}%
\end{pgfscope}%
\begin{pgfscope}%
\pgfsetbuttcap%
\pgfsetroundjoin%
\definecolor{currentfill}{rgb}{0.000000,0.000000,0.000000}%
\pgfsetfillcolor{currentfill}%
\pgfsetlinewidth{0.401500pt}%
\definecolor{currentstroke}{rgb}{0.000000,0.000000,0.000000}%
\pgfsetstrokecolor{currentstroke}%
\pgfsetdash{}{0pt}%
\pgfsys@defobject{currentmarker}{\pgfqpoint{0.000000in}{0.000000in}}{\pgfqpoint{0.020833in}{0.000000in}}{%
\pgfpathmoveto{\pgfqpoint{0.000000in}{0.000000in}}%
\pgfpathlineto{\pgfqpoint{0.020833in}{0.000000in}}%
\pgfusepath{stroke,fill}%
}%
\begin{pgfscope}%
\pgfsys@transformshift{0.650000in}{1.516968in}%
\pgfsys@useobject{currentmarker}{}%
\end{pgfscope}%
\end{pgfscope}%
\begin{pgfscope}%
\pgfsetbuttcap%
\pgfsetroundjoin%
\definecolor{currentfill}{rgb}{0.000000,0.000000,0.000000}%
\pgfsetfillcolor{currentfill}%
\pgfsetlinewidth{0.401500pt}%
\definecolor{currentstroke}{rgb}{0.000000,0.000000,0.000000}%
\pgfsetstrokecolor{currentstroke}%
\pgfsetdash{}{0pt}%
\pgfsys@defobject{currentmarker}{\pgfqpoint{-0.020833in}{0.000000in}}{\pgfqpoint{-0.000000in}{0.000000in}}{%
\pgfpathmoveto{\pgfqpoint{-0.000000in}{0.000000in}}%
\pgfpathlineto{\pgfqpoint{-0.020833in}{0.000000in}}%
\pgfusepath{stroke,fill}%
}%
\begin{pgfscope}%
\pgfsys@transformshift{3.038110in}{1.516968in}%
\pgfsys@useobject{currentmarker}{}%
\end{pgfscope}%
\end{pgfscope}%
\begin{pgfscope}%
\pgfsetbuttcap%
\pgfsetroundjoin%
\definecolor{currentfill}{rgb}{0.000000,0.000000,0.000000}%
\pgfsetfillcolor{currentfill}%
\pgfsetlinewidth{0.401500pt}%
\definecolor{currentstroke}{rgb}{0.000000,0.000000,0.000000}%
\pgfsetstrokecolor{currentstroke}%
\pgfsetdash{}{0pt}%
\pgfsys@defobject{currentmarker}{\pgfqpoint{0.000000in}{0.000000in}}{\pgfqpoint{0.020833in}{0.000000in}}{%
\pgfpathmoveto{\pgfqpoint{0.000000in}{0.000000in}}%
\pgfpathlineto{\pgfqpoint{0.020833in}{0.000000in}}%
\pgfusepath{stroke,fill}%
}%
\begin{pgfscope}%
\pgfsys@transformshift{0.650000in}{1.616693in}%
\pgfsys@useobject{currentmarker}{}%
\end{pgfscope}%
\end{pgfscope}%
\begin{pgfscope}%
\pgfsetbuttcap%
\pgfsetroundjoin%
\definecolor{currentfill}{rgb}{0.000000,0.000000,0.000000}%
\pgfsetfillcolor{currentfill}%
\pgfsetlinewidth{0.401500pt}%
\definecolor{currentstroke}{rgb}{0.000000,0.000000,0.000000}%
\pgfsetstrokecolor{currentstroke}%
\pgfsetdash{}{0pt}%
\pgfsys@defobject{currentmarker}{\pgfqpoint{-0.020833in}{0.000000in}}{\pgfqpoint{-0.000000in}{0.000000in}}{%
\pgfpathmoveto{\pgfqpoint{-0.000000in}{0.000000in}}%
\pgfpathlineto{\pgfqpoint{-0.020833in}{0.000000in}}%
\pgfusepath{stroke,fill}%
}%
\begin{pgfscope}%
\pgfsys@transformshift{3.038110in}{1.616693in}%
\pgfsys@useobject{currentmarker}{}%
\end{pgfscope}%
\end{pgfscope}%
\begin{pgfscope}%
\pgfsetbuttcap%
\pgfsetroundjoin%
\definecolor{currentfill}{rgb}{0.000000,0.000000,0.000000}%
\pgfsetfillcolor{currentfill}%
\pgfsetlinewidth{0.401500pt}%
\definecolor{currentstroke}{rgb}{0.000000,0.000000,0.000000}%
\pgfsetstrokecolor{currentstroke}%
\pgfsetdash{}{0pt}%
\pgfsys@defobject{currentmarker}{\pgfqpoint{0.000000in}{0.000000in}}{\pgfqpoint{0.020833in}{0.000000in}}{%
\pgfpathmoveto{\pgfqpoint{0.000000in}{0.000000in}}%
\pgfpathlineto{\pgfqpoint{0.020833in}{0.000000in}}%
\pgfusepath{stroke,fill}%
}%
\begin{pgfscope}%
\pgfsys@transformshift{0.650000in}{1.716417in}%
\pgfsys@useobject{currentmarker}{}%
\end{pgfscope}%
\end{pgfscope}%
\begin{pgfscope}%
\pgfsetbuttcap%
\pgfsetroundjoin%
\definecolor{currentfill}{rgb}{0.000000,0.000000,0.000000}%
\pgfsetfillcolor{currentfill}%
\pgfsetlinewidth{0.401500pt}%
\definecolor{currentstroke}{rgb}{0.000000,0.000000,0.000000}%
\pgfsetstrokecolor{currentstroke}%
\pgfsetdash{}{0pt}%
\pgfsys@defobject{currentmarker}{\pgfqpoint{-0.020833in}{0.000000in}}{\pgfqpoint{-0.000000in}{0.000000in}}{%
\pgfpathmoveto{\pgfqpoint{-0.000000in}{0.000000in}}%
\pgfpathlineto{\pgfqpoint{-0.020833in}{0.000000in}}%
\pgfusepath{stroke,fill}%
}%
\begin{pgfscope}%
\pgfsys@transformshift{3.038110in}{1.716417in}%
\pgfsys@useobject{currentmarker}{}%
\end{pgfscope}%
\end{pgfscope}%
\begin{pgfscope}%
\pgfsetbuttcap%
\pgfsetroundjoin%
\definecolor{currentfill}{rgb}{0.000000,0.000000,0.000000}%
\pgfsetfillcolor{currentfill}%
\pgfsetlinewidth{0.401500pt}%
\definecolor{currentstroke}{rgb}{0.000000,0.000000,0.000000}%
\pgfsetstrokecolor{currentstroke}%
\pgfsetdash{}{0pt}%
\pgfsys@defobject{currentmarker}{\pgfqpoint{0.000000in}{0.000000in}}{\pgfqpoint{0.020833in}{0.000000in}}{%
\pgfpathmoveto{\pgfqpoint{0.000000in}{0.000000in}}%
\pgfpathlineto{\pgfqpoint{0.020833in}{0.000000in}}%
\pgfusepath{stroke,fill}%
}%
\begin{pgfscope}%
\pgfsys@transformshift{0.650000in}{1.816142in}%
\pgfsys@useobject{currentmarker}{}%
\end{pgfscope}%
\end{pgfscope}%
\begin{pgfscope}%
\pgfsetbuttcap%
\pgfsetroundjoin%
\definecolor{currentfill}{rgb}{0.000000,0.000000,0.000000}%
\pgfsetfillcolor{currentfill}%
\pgfsetlinewidth{0.401500pt}%
\definecolor{currentstroke}{rgb}{0.000000,0.000000,0.000000}%
\pgfsetstrokecolor{currentstroke}%
\pgfsetdash{}{0pt}%
\pgfsys@defobject{currentmarker}{\pgfqpoint{-0.020833in}{0.000000in}}{\pgfqpoint{-0.000000in}{0.000000in}}{%
\pgfpathmoveto{\pgfqpoint{-0.000000in}{0.000000in}}%
\pgfpathlineto{\pgfqpoint{-0.020833in}{0.000000in}}%
\pgfusepath{stroke,fill}%
}%
\begin{pgfscope}%
\pgfsys@transformshift{3.038110in}{1.816142in}%
\pgfsys@useobject{currentmarker}{}%
\end{pgfscope}%
\end{pgfscope}%
\begin{pgfscope}%
\pgfsetbuttcap%
\pgfsetroundjoin%
\definecolor{currentfill}{rgb}{0.000000,0.000000,0.000000}%
\pgfsetfillcolor{currentfill}%
\pgfsetlinewidth{0.401500pt}%
\definecolor{currentstroke}{rgb}{0.000000,0.000000,0.000000}%
\pgfsetstrokecolor{currentstroke}%
\pgfsetdash{}{0pt}%
\pgfsys@defobject{currentmarker}{\pgfqpoint{0.000000in}{0.000000in}}{\pgfqpoint{0.020833in}{0.000000in}}{%
\pgfpathmoveto{\pgfqpoint{0.000000in}{0.000000in}}%
\pgfpathlineto{\pgfqpoint{0.020833in}{0.000000in}}%
\pgfusepath{stroke,fill}%
}%
\begin{pgfscope}%
\pgfsys@transformshift{0.650000in}{2.015590in}%
\pgfsys@useobject{currentmarker}{}%
\end{pgfscope}%
\end{pgfscope}%
\begin{pgfscope}%
\pgfsetbuttcap%
\pgfsetroundjoin%
\definecolor{currentfill}{rgb}{0.000000,0.000000,0.000000}%
\pgfsetfillcolor{currentfill}%
\pgfsetlinewidth{0.401500pt}%
\definecolor{currentstroke}{rgb}{0.000000,0.000000,0.000000}%
\pgfsetstrokecolor{currentstroke}%
\pgfsetdash{}{0pt}%
\pgfsys@defobject{currentmarker}{\pgfqpoint{-0.020833in}{0.000000in}}{\pgfqpoint{-0.000000in}{0.000000in}}{%
\pgfpathmoveto{\pgfqpoint{-0.000000in}{0.000000in}}%
\pgfpathlineto{\pgfqpoint{-0.020833in}{0.000000in}}%
\pgfusepath{stroke,fill}%
}%
\begin{pgfscope}%
\pgfsys@transformshift{3.038110in}{2.015590in}%
\pgfsys@useobject{currentmarker}{}%
\end{pgfscope}%
\end{pgfscope}%
\begin{pgfscope}%
\pgfsetbuttcap%
\pgfsetroundjoin%
\definecolor{currentfill}{rgb}{0.000000,0.000000,0.000000}%
\pgfsetfillcolor{currentfill}%
\pgfsetlinewidth{0.401500pt}%
\definecolor{currentstroke}{rgb}{0.000000,0.000000,0.000000}%
\pgfsetstrokecolor{currentstroke}%
\pgfsetdash{}{0pt}%
\pgfsys@defobject{currentmarker}{\pgfqpoint{0.000000in}{0.000000in}}{\pgfqpoint{0.020833in}{0.000000in}}{%
\pgfpathmoveto{\pgfqpoint{0.000000in}{0.000000in}}%
\pgfpathlineto{\pgfqpoint{0.020833in}{0.000000in}}%
\pgfusepath{stroke,fill}%
}%
\begin{pgfscope}%
\pgfsys@transformshift{0.650000in}{2.115315in}%
\pgfsys@useobject{currentmarker}{}%
\end{pgfscope}%
\end{pgfscope}%
\begin{pgfscope}%
\pgfsetbuttcap%
\pgfsetroundjoin%
\definecolor{currentfill}{rgb}{0.000000,0.000000,0.000000}%
\pgfsetfillcolor{currentfill}%
\pgfsetlinewidth{0.401500pt}%
\definecolor{currentstroke}{rgb}{0.000000,0.000000,0.000000}%
\pgfsetstrokecolor{currentstroke}%
\pgfsetdash{}{0pt}%
\pgfsys@defobject{currentmarker}{\pgfqpoint{-0.020833in}{0.000000in}}{\pgfqpoint{-0.000000in}{0.000000in}}{%
\pgfpathmoveto{\pgfqpoint{-0.000000in}{0.000000in}}%
\pgfpathlineto{\pgfqpoint{-0.020833in}{0.000000in}}%
\pgfusepath{stroke,fill}%
}%
\begin{pgfscope}%
\pgfsys@transformshift{3.038110in}{2.115315in}%
\pgfsys@useobject{currentmarker}{}%
\end{pgfscope}%
\end{pgfscope}%
\begin{pgfscope}%
\pgfsetbuttcap%
\pgfsetroundjoin%
\definecolor{currentfill}{rgb}{0.000000,0.000000,0.000000}%
\pgfsetfillcolor{currentfill}%
\pgfsetlinewidth{0.401500pt}%
\definecolor{currentstroke}{rgb}{0.000000,0.000000,0.000000}%
\pgfsetstrokecolor{currentstroke}%
\pgfsetdash{}{0pt}%
\pgfsys@defobject{currentmarker}{\pgfqpoint{0.000000in}{0.000000in}}{\pgfqpoint{0.020833in}{0.000000in}}{%
\pgfpathmoveto{\pgfqpoint{0.000000in}{0.000000in}}%
\pgfpathlineto{\pgfqpoint{0.020833in}{0.000000in}}%
\pgfusepath{stroke,fill}%
}%
\begin{pgfscope}%
\pgfsys@transformshift{0.650000in}{2.215039in}%
\pgfsys@useobject{currentmarker}{}%
\end{pgfscope}%
\end{pgfscope}%
\begin{pgfscope}%
\pgfsetbuttcap%
\pgfsetroundjoin%
\definecolor{currentfill}{rgb}{0.000000,0.000000,0.000000}%
\pgfsetfillcolor{currentfill}%
\pgfsetlinewidth{0.401500pt}%
\definecolor{currentstroke}{rgb}{0.000000,0.000000,0.000000}%
\pgfsetstrokecolor{currentstroke}%
\pgfsetdash{}{0pt}%
\pgfsys@defobject{currentmarker}{\pgfqpoint{-0.020833in}{0.000000in}}{\pgfqpoint{-0.000000in}{0.000000in}}{%
\pgfpathmoveto{\pgfqpoint{-0.000000in}{0.000000in}}%
\pgfpathlineto{\pgfqpoint{-0.020833in}{0.000000in}}%
\pgfusepath{stroke,fill}%
}%
\begin{pgfscope}%
\pgfsys@transformshift{3.038110in}{2.215039in}%
\pgfsys@useobject{currentmarker}{}%
\end{pgfscope}%
\end{pgfscope}%
\begin{pgfscope}%
\pgfsetbuttcap%
\pgfsetroundjoin%
\definecolor{currentfill}{rgb}{0.000000,0.000000,0.000000}%
\pgfsetfillcolor{currentfill}%
\pgfsetlinewidth{0.401500pt}%
\definecolor{currentstroke}{rgb}{0.000000,0.000000,0.000000}%
\pgfsetstrokecolor{currentstroke}%
\pgfsetdash{}{0pt}%
\pgfsys@defobject{currentmarker}{\pgfqpoint{0.000000in}{0.000000in}}{\pgfqpoint{0.020833in}{0.000000in}}{%
\pgfpathmoveto{\pgfqpoint{0.000000in}{0.000000in}}%
\pgfpathlineto{\pgfqpoint{0.020833in}{0.000000in}}%
\pgfusepath{stroke,fill}%
}%
\begin{pgfscope}%
\pgfsys@transformshift{0.650000in}{2.314764in}%
\pgfsys@useobject{currentmarker}{}%
\end{pgfscope}%
\end{pgfscope}%
\begin{pgfscope}%
\pgfsetbuttcap%
\pgfsetroundjoin%
\definecolor{currentfill}{rgb}{0.000000,0.000000,0.000000}%
\pgfsetfillcolor{currentfill}%
\pgfsetlinewidth{0.401500pt}%
\definecolor{currentstroke}{rgb}{0.000000,0.000000,0.000000}%
\pgfsetstrokecolor{currentstroke}%
\pgfsetdash{}{0pt}%
\pgfsys@defobject{currentmarker}{\pgfqpoint{-0.020833in}{0.000000in}}{\pgfqpoint{-0.000000in}{0.000000in}}{%
\pgfpathmoveto{\pgfqpoint{-0.000000in}{0.000000in}}%
\pgfpathlineto{\pgfqpoint{-0.020833in}{0.000000in}}%
\pgfusepath{stroke,fill}%
}%
\begin{pgfscope}%
\pgfsys@transformshift{3.038110in}{2.314764in}%
\pgfsys@useobject{currentmarker}{}%
\end{pgfscope}%
\end{pgfscope}%
\begin{pgfscope}%
\definecolor{textcolor}{rgb}{0.000000,0.000000,0.000000}%
\pgfsetstrokecolor{textcolor}%
\pgfsetfillcolor{textcolor}%
\pgftext[x=0.379282in,y=1.417244in,,bottom,rotate=90.000000]{\color{textcolor}{\rmfamily\fontsize{12.000000}{14.400000}\selectfont\catcode`\^=\active\def^{\ifmmode\sp\else\^{}\fi}\catcode`\%=\active\def%{\%}$\beta$}}%
\end{pgfscope}%
\begin{pgfscope}%
\pgfpathrectangle{\pgfqpoint{0.650000in}{0.420000in}}{\pgfqpoint{2.388110in}{1.994488in}}%
\pgfusepath{clip}%
\pgfsetrectcap%
\pgfsetroundjoin%
\pgfsetlinewidth{1.003750pt}%
\definecolor{currentstroke}{rgb}{0.796078,0.094118,0.113725}%
\pgfsetstrokecolor{currentstroke}%
\pgfsetdash{}{0pt}%
\pgfpathmoveto{\pgfqpoint{0.650000in}{1.395520in}}%
\pgfpathlineto{\pgfqpoint{0.652867in}{1.396134in}}%
\pgfpathlineto{\pgfqpoint{0.655735in}{1.397871in}}%
\pgfpathlineto{\pgfqpoint{0.659319in}{1.401498in}}%
\pgfpathlineto{\pgfqpoint{0.663978in}{1.408039in}}%
\pgfpathlineto{\pgfqpoint{0.674730in}{1.423712in}}%
\pgfpathlineto{\pgfqpoint{0.678315in}{1.426849in}}%
\pgfpathlineto{\pgfqpoint{0.681182in}{1.428042in}}%
\pgfpathlineto{\pgfqpoint{0.684049in}{1.428006in}}%
\pgfpathlineto{\pgfqpoint{0.686916in}{1.426839in}}%
\pgfpathlineto{\pgfqpoint{0.690859in}{1.423841in}}%
\pgfpathlineto{\pgfqpoint{0.708780in}{1.408139in}}%
\pgfpathlineto{\pgfqpoint{0.714156in}{1.405790in}}%
\pgfpathlineto{\pgfqpoint{0.719890in}{1.404476in}}%
\pgfpathlineto{\pgfqpoint{0.727059in}{1.404042in}}%
\pgfpathlineto{\pgfqpoint{0.737811in}{1.404674in}}%
\pgfpathlineto{\pgfqpoint{0.758240in}{1.405839in}}%
\pgfpathlineto{\pgfqpoint{0.782971in}{1.406037in}}%
\pgfpathlineto{\pgfqpoint{0.861463in}{1.403669in}}%
\pgfpathlineto{\pgfqpoint{0.890494in}{1.402797in}}%
\pgfpathlineto{\pgfqpoint{0.917375in}{1.401746in}}%
\pgfpathlineto{\pgfqpoint{0.938521in}{1.401677in}}%
\pgfpathlineto{\pgfqpoint{0.999810in}{1.398565in}}%
\pgfpathlineto{\pgfqpoint{1.012354in}{1.399575in}}%
\pgfpathlineto{\pgfqpoint{1.020239in}{1.399455in}}%
\pgfpathlineto{\pgfqpoint{1.030633in}{1.397996in}}%
\pgfpathlineto{\pgfqpoint{1.046403in}{1.395907in}}%
\pgfpathlineto{\pgfqpoint{1.062890in}{1.394100in}}%
\pgfpathlineto{\pgfqpoint{1.068983in}{1.392029in}}%
\pgfpathlineto{\pgfqpoint{1.075793in}{1.388438in}}%
\pgfpathlineto{\pgfqpoint{1.083678in}{1.382992in}}%
\pgfpathlineto{\pgfqpoint{1.089054in}{1.378061in}}%
\pgfpathlineto{\pgfqpoint{1.093714in}{1.372199in}}%
\pgfpathlineto{\pgfqpoint{1.098015in}{1.364937in}}%
\pgfpathlineto{\pgfqpoint{1.102674in}{1.354693in}}%
\pgfpathlineto{\pgfqpoint{1.107333in}{1.341499in}}%
\pgfpathlineto{\pgfqpoint{1.111993in}{1.324726in}}%
\pgfpathlineto{\pgfqpoint{1.116652in}{1.303685in}}%
\pgfpathlineto{\pgfqpoint{1.123820in}{1.264909in}}%
\pgfpathlineto{\pgfqpoint{1.128838in}{1.239991in}}%
\pgfpathlineto{\pgfqpoint{1.131705in}{1.230760in}}%
\pgfpathlineto{\pgfqpoint{1.133497in}{1.227813in}}%
\pgfpathlineto{\pgfqpoint{1.134931in}{1.227207in}}%
\pgfpathlineto{\pgfqpoint{1.136006in}{1.227783in}}%
\pgfpathlineto{\pgfqpoint{1.137440in}{1.229875in}}%
\pgfpathlineto{\pgfqpoint{1.139591in}{1.235600in}}%
\pgfpathlineto{\pgfqpoint{1.142458in}{1.247186in}}%
\pgfpathlineto{\pgfqpoint{1.147117in}{1.271979in}}%
\pgfpathlineto{\pgfqpoint{1.156436in}{1.322340in}}%
\pgfpathlineto{\pgfqpoint{1.161812in}{1.344547in}}%
\pgfpathlineto{\pgfqpoint{1.166472in}{1.358488in}}%
\pgfpathlineto{\pgfqpoint{1.171131in}{1.368968in}}%
\pgfpathlineto{\pgfqpoint{1.175790in}{1.376860in}}%
\pgfpathlineto{\pgfqpoint{1.181525in}{1.384234in}}%
\pgfpathlineto{\pgfqpoint{1.186543in}{1.388620in}}%
\pgfpathlineto{\pgfqpoint{1.191919in}{1.391866in}}%
\pgfpathlineto{\pgfqpoint{1.199087in}{1.394708in}}%
\pgfpathlineto{\pgfqpoint{1.206972in}{1.396269in}}%
\pgfpathlineto{\pgfqpoint{1.222026in}{1.397668in}}%
\pgfpathlineto{\pgfqpoint{1.245323in}{1.399222in}}%
\pgfpathlineto{\pgfqpoint{1.297651in}{1.397953in}}%
\pgfpathlineto{\pgfqpoint{1.320231in}{1.396646in}}%
\pgfpathlineto{\pgfqpoint{1.337435in}{1.396985in}}%
\pgfpathlineto{\pgfqpoint{1.348546in}{1.397058in}}%
\pgfpathlineto{\pgfqpoint{1.368259in}{1.397107in}}%
\pgfpathlineto{\pgfqpoint{1.396215in}{1.396271in}}%
\pgfpathlineto{\pgfqpoint{1.409118in}{1.396499in}}%
\pgfpathlineto{\pgfqpoint{1.410910in}{1.396657in}}%
\pgfpathlineto{\pgfqpoint{1.420587in}{1.396193in}}%
\pgfpathlineto{\pgfqpoint{1.434924in}{1.395280in}}%
\pgfpathlineto{\pgfqpoint{1.449261in}{1.394554in}}%
\pgfpathlineto{\pgfqpoint{1.477575in}{1.391812in}}%
\pgfpathlineto{\pgfqpoint{1.483310in}{1.392692in}}%
\pgfpathlineto{\pgfqpoint{1.498722in}{1.395133in}}%
\pgfpathlineto{\pgfqpoint{1.506965in}{1.394555in}}%
\pgfpathlineto{\pgfqpoint{1.523811in}{1.392802in}}%
\pgfpathlineto{\pgfqpoint{1.537431in}{1.392307in}}%
\pgfpathlineto{\pgfqpoint{1.552842in}{1.392858in}}%
\pgfpathlineto{\pgfqpoint{1.571121in}{1.393329in}}%
\pgfpathlineto{\pgfqpoint{1.581515in}{1.392481in}}%
\pgfpathlineto{\pgfqpoint{1.598002in}{1.390227in}}%
\pgfpathlineto{\pgfqpoint{1.625242in}{1.390561in}}%
\pgfpathlineto{\pgfqpoint{1.627034in}{1.390766in}}%
\pgfpathlineto{\pgfqpoint{1.638145in}{1.390660in}}%
\pgfpathlineto{\pgfqpoint{1.653915in}{1.389692in}}%
\pgfpathlineto{\pgfqpoint{1.672194in}{1.390868in}}%
\pgfpathlineto{\pgfqpoint{1.687248in}{1.388710in}}%
\pgfpathlineto{\pgfqpoint{1.706602in}{1.389043in}}%
\pgfpathlineto{\pgfqpoint{1.708394in}{1.389145in}}%
\pgfpathlineto{\pgfqpoint{1.732408in}{1.389085in}}%
\pgfpathlineto{\pgfqpoint{1.739935in}{1.389031in}}%
\pgfpathlineto{\pgfqpoint{1.758931in}{1.387521in}}%
\pgfpathlineto{\pgfqpoint{1.772909in}{1.388563in}}%
\pgfpathlineto{\pgfqpoint{1.782228in}{1.387430in}}%
\pgfpathlineto{\pgfqpoint{1.801224in}{1.385221in}}%
\pgfpathlineto{\pgfqpoint{1.808392in}{1.385542in}}%
\pgfpathlineto{\pgfqpoint{1.822728in}{1.386732in}}%
\pgfpathlineto{\pgfqpoint{1.838499in}{1.386247in}}%
\pgfpathlineto{\pgfqpoint{1.840291in}{1.386566in}}%
\pgfpathlineto{\pgfqpoint{1.853552in}{1.387472in}}%
\pgfpathlineto{\pgfqpoint{1.867172in}{1.387189in}}%
\pgfpathlineto{\pgfqpoint{1.886526in}{1.385083in}}%
\pgfpathlineto{\pgfqpoint{1.907314in}{1.386575in}}%
\pgfpathlineto{\pgfqpoint{1.916633in}{1.384990in}}%
\pgfpathlineto{\pgfqpoint{1.924159in}{1.384181in}}%
\pgfpathlineto{\pgfqpoint{1.933836in}{1.384177in}}%
\pgfpathlineto{\pgfqpoint{1.947456in}{1.385179in}}%
\pgfpathlineto{\pgfqpoint{1.963584in}{1.387013in}}%
\pgfpathlineto{\pgfqpoint{1.972903in}{1.385578in}}%
\pgfpathlineto{\pgfqpoint{1.978996in}{1.385337in}}%
\pgfpathlineto{\pgfqpoint{1.986881in}{1.386346in}}%
\pgfpathlineto{\pgfqpoint{1.995841in}{1.387167in}}%
\pgfpathlineto{\pgfqpoint{2.003368in}{1.386461in}}%
\pgfpathlineto{\pgfqpoint{2.020213in}{1.384496in}}%
\pgfpathlineto{\pgfqpoint{2.030965in}{1.383175in}}%
\pgfpathlineto{\pgfqpoint{2.041001in}{1.381848in}}%
\pgfpathlineto{\pgfqpoint{2.047094in}{1.382372in}}%
\pgfpathlineto{\pgfqpoint{2.067523in}{1.385182in}}%
\pgfpathlineto{\pgfqpoint{2.075408in}{1.384990in}}%
\pgfpathlineto{\pgfqpoint{2.082935in}{1.383584in}}%
\pgfpathlineto{\pgfqpoint{2.093329in}{1.381303in}}%
\pgfpathlineto{\pgfqpoint{2.100138in}{1.381188in}}%
\pgfpathlineto{\pgfqpoint{2.120568in}{1.382536in}}%
\pgfpathlineto{\pgfqpoint{2.131679in}{1.383163in}}%
\pgfpathlineto{\pgfqpoint{2.145298in}{1.383776in}}%
\pgfpathlineto{\pgfqpoint{2.154258in}{1.383914in}}%
\pgfpathlineto{\pgfqpoint{2.161785in}{1.382685in}}%
\pgfpathlineto{\pgfqpoint{2.172537in}{1.380526in}}%
\pgfpathlineto{\pgfqpoint{2.199058in}{1.378406in}}%
\pgfpathlineto{\pgfqpoint{2.205867in}{1.380409in}}%
\pgfpathlineto{\pgfqpoint{2.214827in}{1.382848in}}%
\pgfpathlineto{\pgfqpoint{2.221278in}{1.383139in}}%
\pgfpathlineto{\pgfqpoint{2.237764in}{1.382151in}}%
\pgfpathlineto{\pgfqpoint{2.244932in}{1.380436in}}%
\pgfpathlineto{\pgfqpoint{2.251383in}{1.379230in}}%
\pgfpathlineto{\pgfqpoint{2.257475in}{1.379148in}}%
\pgfpathlineto{\pgfqpoint{2.270736in}{1.380464in}}%
\pgfpathlineto{\pgfqpoint{2.283996in}{1.382299in}}%
\pgfpathlineto{\pgfqpoint{2.289372in}{1.381194in}}%
\pgfpathlineto{\pgfqpoint{2.297615in}{1.379359in}}%
\pgfpathlineto{\pgfqpoint{2.302632in}{1.379551in}}%
\pgfpathlineto{\pgfqpoint{2.315534in}{1.380717in}}%
\pgfpathlineto{\pgfqpoint{2.323060in}{1.379443in}}%
\pgfpathlineto{\pgfqpoint{2.332378in}{1.377846in}}%
\pgfpathlineto{\pgfqpoint{2.334170in}{1.377879in}}%
\pgfpathlineto{\pgfqpoint{2.342055in}{1.378105in}}%
\pgfpathlineto{\pgfqpoint{2.351014in}{1.378109in}}%
\pgfpathlineto{\pgfqpoint{2.361049in}{1.377970in}}%
\pgfpathlineto{\pgfqpoint{2.362841in}{1.378385in}}%
\pgfpathlineto{\pgfqpoint{2.372159in}{1.380623in}}%
\pgfpathlineto{\pgfqpoint{2.378969in}{1.381598in}}%
\pgfpathlineto{\pgfqpoint{2.385061in}{1.381364in}}%
\pgfpathlineto{\pgfqpoint{2.402980in}{1.379698in}}%
\pgfpathlineto{\pgfqpoint{2.412299in}{1.379283in}}%
\pgfpathlineto{\pgfqpoint{2.437027in}{1.376449in}}%
\pgfpathlineto{\pgfqpoint{2.444912in}{1.377224in}}%
\pgfpathlineto{\pgfqpoint{2.464265in}{1.379839in}}%
\pgfpathlineto{\pgfqpoint{2.471792in}{1.380851in}}%
\pgfpathlineto{\pgfqpoint{2.478242in}{1.383011in}}%
\pgfpathlineto{\pgfqpoint{2.496520in}{1.391036in}}%
\pgfpathlineto{\pgfqpoint{2.501896in}{1.394809in}}%
\pgfpathlineto{\pgfqpoint{2.507272in}{1.399889in}}%
\pgfpathlineto{\pgfqpoint{2.511932in}{1.405699in}}%
\pgfpathlineto{\pgfqpoint{2.515874in}{1.412303in}}%
\pgfpathlineto{\pgfqpoint{2.519816in}{1.421152in}}%
\pgfpathlineto{\pgfqpoint{2.523759in}{1.432926in}}%
\pgfpathlineto{\pgfqpoint{2.528059in}{1.449456in}}%
\pgfpathlineto{\pgfqpoint{2.533077in}{1.473632in}}%
\pgfpathlineto{\pgfqpoint{2.538453in}{1.505483in}}%
\pgfpathlineto{\pgfqpoint{2.543112in}{1.540444in}}%
\pgfpathlineto{\pgfqpoint{2.547413in}{1.581406in}}%
\pgfpathlineto{\pgfqpoint{2.552072in}{1.636869in}}%
\pgfpathlineto{\pgfqpoint{2.558165in}{1.725921in}}%
\pgfpathlineto{\pgfqpoint{2.564258in}{1.810781in}}%
\pgfpathlineto{\pgfqpoint{2.567125in}{1.838417in}}%
\pgfpathlineto{\pgfqpoint{2.569276in}{1.851481in}}%
\pgfpathlineto{\pgfqpoint{2.570710in}{1.855256in}}%
\pgfpathlineto{\pgfqpoint{2.571427in}{1.855578in}}%
\pgfpathlineto{\pgfqpoint{2.572143in}{1.854843in}}%
\pgfpathlineto{\pgfqpoint{2.573219in}{1.851787in}}%
\pgfpathlineto{\pgfqpoint{2.574652in}{1.843348in}}%
\pgfpathlineto{\pgfqpoint{2.576803in}{1.825363in}}%
\pgfpathlineto{\pgfqpoint{2.580029in}{1.786751in}}%
\pgfpathlineto{\pgfqpoint{2.586481in}{1.691134in}}%
\pgfpathlineto{\pgfqpoint{2.592933in}{1.603789in}}%
\pgfpathlineto{\pgfqpoint{2.597951in}{1.550486in}}%
\pgfpathlineto{\pgfqpoint{2.602610in}{1.512502in}}%
\pgfpathlineto{\pgfqpoint{2.607628in}{1.480840in}}%
\pgfpathlineto{\pgfqpoint{2.613005in}{1.454384in}}%
\pgfpathlineto{\pgfqpoint{2.618740in}{1.431650in}}%
\pgfpathlineto{\pgfqpoint{2.623758in}{1.415923in}}%
\pgfpathlineto{\pgfqpoint{2.628059in}{1.405996in}}%
\pgfpathlineto{\pgfqpoint{2.631643in}{1.400298in}}%
\pgfpathlineto{\pgfqpoint{2.635228in}{1.396634in}}%
\pgfpathlineto{\pgfqpoint{2.639170in}{1.394223in}}%
\pgfpathlineto{\pgfqpoint{2.650999in}{1.388111in}}%
\pgfpathlineto{\pgfqpoint{2.662827in}{1.379747in}}%
\pgfpathlineto{\pgfqpoint{2.666053in}{1.379203in}}%
\pgfpathlineto{\pgfqpoint{2.669637in}{1.379731in}}%
\pgfpathlineto{\pgfqpoint{2.680749in}{1.382370in}}%
\pgfpathlineto{\pgfqpoint{2.685408in}{1.381638in}}%
\pgfpathlineto{\pgfqpoint{2.693652in}{1.380168in}}%
\pgfpathlineto{\pgfqpoint{2.697595in}{1.380660in}}%
\pgfpathlineto{\pgfqpoint{2.703688in}{1.382457in}}%
\pgfpathlineto{\pgfqpoint{2.710499in}{1.384264in}}%
\pgfpathlineto{\pgfqpoint{2.714800in}{1.384283in}}%
\pgfpathlineto{\pgfqpoint{2.720176in}{1.383079in}}%
\pgfpathlineto{\pgfqpoint{2.729496in}{1.380443in}}%
\pgfpathlineto{\pgfqpoint{2.734514in}{1.380420in}}%
\pgfpathlineto{\pgfqpoint{2.742041in}{1.381729in}}%
\pgfpathlineto{\pgfqpoint{2.748851in}{1.382161in}}%
\pgfpathlineto{\pgfqpoint{2.762830in}{1.382095in}}%
\pgfpathlineto{\pgfqpoint{2.766056in}{1.383655in}}%
\pgfpathlineto{\pgfqpoint{2.771432in}{1.387137in}}%
\pgfpathlineto{\pgfqpoint{2.787203in}{1.398290in}}%
\pgfpathlineto{\pgfqpoint{2.789354in}{1.397935in}}%
\pgfpathlineto{\pgfqpoint{2.791505in}{1.396226in}}%
\pgfpathlineto{\pgfqpoint{2.793655in}{1.392906in}}%
\pgfpathlineto{\pgfqpoint{2.796164in}{1.386923in}}%
\pgfpathlineto{\pgfqpoint{2.800107in}{1.373982in}}%
\pgfpathlineto{\pgfqpoint{2.805125in}{1.357815in}}%
\pgfpathlineto{\pgfqpoint{2.807634in}{1.352598in}}%
\pgfpathlineto{\pgfqpoint{2.809426in}{1.350989in}}%
\pgfpathlineto{\pgfqpoint{2.809785in}{1.350903in}}%
\pgfpathlineto{\pgfqpoint{2.809785in}{1.350903in}}%
\pgfusepath{stroke}%
\end{pgfscope}%
\begin{pgfscope}%
\pgfpathrectangle{\pgfqpoint{0.650000in}{0.420000in}}{\pgfqpoint{2.388110in}{1.994488in}}%
\pgfusepath{clip}%
\pgfsetrectcap%
\pgfsetroundjoin%
\pgfsetlinewidth{1.003750pt}%
\definecolor{currentstroke}{rgb}{0.454902,0.678431,0.819608}%
\pgfsetstrokecolor{currentstroke}%
\pgfsetdash{}{0pt}%
\pgfpathmoveto{\pgfqpoint{0.650000in}{1.395523in}}%
\pgfpathlineto{\pgfqpoint{0.652867in}{1.396142in}}%
\pgfpathlineto{\pgfqpoint{0.655735in}{1.397894in}}%
\pgfpathlineto{\pgfqpoint{0.659319in}{1.401552in}}%
\pgfpathlineto{\pgfqpoint{0.663978in}{1.408147in}}%
\pgfpathlineto{\pgfqpoint{0.674372in}{1.423411in}}%
\pgfpathlineto{\pgfqpoint{0.677956in}{1.426620in}}%
\pgfpathlineto{\pgfqpoint{0.680823in}{1.427855in}}%
\pgfpathlineto{\pgfqpoint{0.683691in}{1.427839in}}%
\pgfpathlineto{\pgfqpoint{0.686558in}{1.426673in}}%
\pgfpathlineto{\pgfqpoint{0.690501in}{1.423666in}}%
\pgfpathlineto{\pgfqpoint{0.706987in}{1.409254in}}%
\pgfpathlineto{\pgfqpoint{0.712722in}{1.406541in}}%
\pgfpathlineto{\pgfqpoint{0.718815in}{1.404891in}}%
\pgfpathlineto{\pgfqpoint{0.725625in}{1.404223in}}%
\pgfpathlineto{\pgfqpoint{0.734227in}{1.404590in}}%
\pgfpathlineto{\pgfqpoint{0.749638in}{1.405357in}}%
\pgfpathlineto{\pgfqpoint{0.771860in}{1.405591in}}%
\pgfpathlineto{\pgfqpoint{0.794082in}{1.406137in}}%
\pgfpathlineto{\pgfqpoint{0.807701in}{1.405316in}}%
\pgfpathlineto{\pgfqpoint{0.825980in}{1.404378in}}%
\pgfpathlineto{\pgfqpoint{0.850711in}{1.403537in}}%
\pgfpathlineto{\pgfqpoint{0.873649in}{1.404392in}}%
\pgfpathlineto{\pgfqpoint{0.887627in}{1.403618in}}%
\pgfpathlineto{\pgfqpoint{0.906623in}{1.401530in}}%
\pgfpathlineto{\pgfqpoint{0.916658in}{1.400663in}}%
\pgfpathlineto{\pgfqpoint{0.939955in}{1.399668in}}%
\pgfpathlineto{\pgfqpoint{0.951066in}{1.397840in}}%
\pgfpathlineto{\pgfqpoint{0.962893in}{1.394491in}}%
\pgfpathlineto{\pgfqpoint{0.971495in}{1.391196in}}%
\pgfpathlineto{\pgfqpoint{0.977947in}{1.387489in}}%
\pgfpathlineto{\pgfqpoint{0.984040in}{1.382729in}}%
\pgfpathlineto{\pgfqpoint{0.989774in}{1.376793in}}%
\pgfpathlineto{\pgfqpoint{0.995509in}{1.369160in}}%
\pgfpathlineto{\pgfqpoint{1.001960in}{1.358583in}}%
\pgfpathlineto{\pgfqpoint{1.009129in}{1.344443in}}%
\pgfpathlineto{\pgfqpoint{1.019164in}{1.321686in}}%
\pgfpathlineto{\pgfqpoint{1.037802in}{1.278633in}}%
\pgfpathlineto{\pgfqpoint{1.042461in}{1.270866in}}%
\pgfpathlineto{\pgfqpoint{1.046403in}{1.266344in}}%
\pgfpathlineto{\pgfqpoint{1.049629in}{1.264306in}}%
\pgfpathlineto{\pgfqpoint{1.052138in}{1.263907in}}%
\pgfpathlineto{\pgfqpoint{1.054647in}{1.264668in}}%
\pgfpathlineto{\pgfqpoint{1.057156in}{1.266700in}}%
\pgfpathlineto{\pgfqpoint{1.060023in}{1.270693in}}%
\pgfpathlineto{\pgfqpoint{1.063249in}{1.277410in}}%
\pgfpathlineto{\pgfqpoint{1.066833in}{1.287618in}}%
\pgfpathlineto{\pgfqpoint{1.071134in}{1.303363in}}%
\pgfpathlineto{\pgfqpoint{1.076868in}{1.329023in}}%
\pgfpathlineto{\pgfqpoint{1.087262in}{1.382051in}}%
\pgfpathlineto{\pgfqpoint{1.095864in}{1.423689in}}%
\pgfpathlineto{\pgfqpoint{1.100882in}{1.443245in}}%
\pgfpathlineto{\pgfqpoint{1.104825in}{1.454678in}}%
\pgfpathlineto{\pgfqpoint{1.108050in}{1.460822in}}%
\pgfpathlineto{\pgfqpoint{1.110559in}{1.463301in}}%
\pgfpathlineto{\pgfqpoint{1.112351in}{1.463721in}}%
\pgfpathlineto{\pgfqpoint{1.114143in}{1.462947in}}%
\pgfpathlineto{\pgfqpoint{1.115935in}{1.460941in}}%
\pgfpathlineto{\pgfqpoint{1.118086in}{1.456897in}}%
\pgfpathlineto{\pgfqpoint{1.120953in}{1.448865in}}%
\pgfpathlineto{\pgfqpoint{1.124896in}{1.433865in}}%
\pgfpathlineto{\pgfqpoint{1.134214in}{1.396357in}}%
\pgfpathlineto{\pgfqpoint{1.137082in}{1.389315in}}%
\pgfpathlineto{\pgfqpoint{1.139591in}{1.385874in}}%
\pgfpathlineto{\pgfqpoint{1.141741in}{1.384819in}}%
\pgfpathlineto{\pgfqpoint{1.143891in}{1.385182in}}%
\pgfpathlineto{\pgfqpoint{1.146759in}{1.387114in}}%
\pgfpathlineto{\pgfqpoint{1.154285in}{1.392853in}}%
\pgfpathlineto{\pgfqpoint{1.156794in}{1.393283in}}%
\pgfpathlineto{\pgfqpoint{1.159303in}{1.392548in}}%
\pgfpathlineto{\pgfqpoint{1.162171in}{1.390471in}}%
\pgfpathlineto{\pgfqpoint{1.165755in}{1.386313in}}%
\pgfpathlineto{\pgfqpoint{1.172206in}{1.376497in}}%
\pgfpathlineto{\pgfqpoint{1.184751in}{1.355571in}}%
\pgfpathlineto{\pgfqpoint{1.189410in}{1.349709in}}%
\pgfpathlineto{\pgfqpoint{1.190127in}{1.349518in}}%
\pgfpathlineto{\pgfqpoint{1.195145in}{1.345034in}}%
\pgfpathlineto{\pgfqpoint{1.199446in}{1.342662in}}%
\pgfpathlineto{\pgfqpoint{1.203030in}{1.341927in}}%
\pgfpathlineto{\pgfqpoint{1.206256in}{1.342379in}}%
\pgfpathlineto{\pgfqpoint{1.209481in}{1.343954in}}%
\pgfpathlineto{\pgfqpoint{1.211632in}{1.346061in}}%
\pgfpathlineto{\pgfqpoint{1.214141in}{1.348970in}}%
\pgfpathlineto{\pgfqpoint{1.217366in}{1.353320in}}%
\pgfpathlineto{\pgfqpoint{1.221309in}{1.359430in}}%
\pgfpathlineto{\pgfqpoint{1.223101in}{1.362604in}}%
\pgfpathlineto{\pgfqpoint{1.230269in}{1.375346in}}%
\pgfpathlineto{\pgfqpoint{1.240664in}{1.394678in}}%
\pgfpathlineto{\pgfqpoint{1.249982in}{1.414629in}}%
\pgfpathlineto{\pgfqpoint{1.261452in}{1.438912in}}%
\pgfpathlineto{\pgfqpoint{1.268261in}{1.450635in}}%
\pgfpathlineto{\pgfqpoint{1.273279in}{1.457117in}}%
\pgfpathlineto{\pgfqpoint{1.278297in}{1.462153in}}%
\pgfpathlineto{\pgfqpoint{1.282956in}{1.465375in}}%
\pgfpathlineto{\pgfqpoint{1.287257in}{1.467051in}}%
\pgfpathlineto{\pgfqpoint{1.289766in}{1.467044in}}%
\pgfpathlineto{\pgfqpoint{1.293350in}{1.466581in}}%
\pgfpathlineto{\pgfqpoint{1.296935in}{1.464963in}}%
\pgfpathlineto{\pgfqpoint{1.300877in}{1.461904in}}%
\pgfpathlineto{\pgfqpoint{1.304820in}{1.457438in}}%
\pgfpathlineto{\pgfqpoint{1.311988in}{1.448281in}}%
\pgfpathlineto{\pgfqpoint{1.323099in}{1.432258in}}%
\pgfpathlineto{\pgfqpoint{1.329909in}{1.420844in}}%
\pgfpathlineto{\pgfqpoint{1.337077in}{1.406440in}}%
\pgfpathlineto{\pgfqpoint{1.363241in}{1.351162in}}%
\pgfpathlineto{\pgfqpoint{1.368259in}{1.343143in}}%
\pgfpathlineto{\pgfqpoint{1.374352in}{1.334517in}}%
\pgfpathlineto{\pgfqpoint{1.379370in}{1.329107in}}%
\pgfpathlineto{\pgfqpoint{1.383671in}{1.325955in}}%
\pgfpathlineto{\pgfqpoint{1.386897in}{1.324635in}}%
\pgfpathlineto{\pgfqpoint{1.387613in}{1.324975in}}%
\pgfpathlineto{\pgfqpoint{1.390839in}{1.324869in}}%
\pgfpathlineto{\pgfqpoint{1.394065in}{1.325855in}}%
\pgfpathlineto{\pgfqpoint{1.397649in}{1.328238in}}%
\pgfpathlineto{\pgfqpoint{1.400516in}{1.331426in}}%
\pgfpathlineto{\pgfqpoint{1.404817in}{1.336697in}}%
\pgfpathlineto{\pgfqpoint{1.406609in}{1.339519in}}%
\pgfpathlineto{\pgfqpoint{1.414495in}{1.351325in}}%
\pgfpathlineto{\pgfqpoint{1.416645in}{1.355003in}}%
\pgfpathlineto{\pgfqpoint{1.437791in}{1.387660in}}%
\pgfpathlineto{\pgfqpoint{1.443168in}{1.393631in}}%
\pgfpathlineto{\pgfqpoint{1.448544in}{1.398009in}}%
\pgfpathlineto{\pgfqpoint{1.454637in}{1.401555in}}%
\pgfpathlineto{\pgfqpoint{1.460013in}{1.403407in}}%
\pgfpathlineto{\pgfqpoint{1.464673in}{1.403834in}}%
\pgfpathlineto{\pgfqpoint{1.468973in}{1.403040in}}%
\pgfpathlineto{\pgfqpoint{1.473633in}{1.400939in}}%
\pgfpathlineto{\pgfqpoint{1.480801in}{1.396212in}}%
\pgfpathlineto{\pgfqpoint{1.491912in}{1.388712in}}%
\pgfpathlineto{\pgfqpoint{1.500514in}{1.384455in}}%
\pgfpathlineto{\pgfqpoint{1.524169in}{1.373102in}}%
\pgfpathlineto{\pgfqpoint{1.531337in}{1.369625in}}%
\pgfpathlineto{\pgfqpoint{1.541015in}{1.364775in}}%
\pgfpathlineto{\pgfqpoint{1.546391in}{1.363186in}}%
\pgfpathlineto{\pgfqpoint{1.552484in}{1.362553in}}%
\pgfpathlineto{\pgfqpoint{1.559294in}{1.362937in}}%
\pgfpathlineto{\pgfqpoint{1.564312in}{1.364230in}}%
\pgfpathlineto{\pgfqpoint{1.573989in}{1.367489in}}%
\pgfpathlineto{\pgfqpoint{1.581874in}{1.369543in}}%
\pgfpathlineto{\pgfqpoint{1.592985in}{1.372189in}}%
\pgfpathlineto{\pgfqpoint{1.599436in}{1.374964in}}%
\pgfpathlineto{\pgfqpoint{1.614490in}{1.382006in}}%
\pgfpathlineto{\pgfqpoint{1.622733in}{1.384230in}}%
\pgfpathlineto{\pgfqpoint{1.628826in}{1.384734in}}%
\pgfpathlineto{\pgfqpoint{1.635278in}{1.383992in}}%
\pgfpathlineto{\pgfqpoint{1.641012in}{1.381990in}}%
\pgfpathlineto{\pgfqpoint{1.647105in}{1.378614in}}%
\pgfpathlineto{\pgfqpoint{1.654274in}{1.373209in}}%
\pgfpathlineto{\pgfqpoint{1.664309in}{1.365568in}}%
\pgfpathlineto{\pgfqpoint{1.669685in}{1.362732in}}%
\pgfpathlineto{\pgfqpoint{1.676137in}{1.360684in}}%
\pgfpathlineto{\pgfqpoint{1.683305in}{1.358241in}}%
\pgfpathlineto{\pgfqpoint{1.687965in}{1.355436in}}%
\pgfpathlineto{\pgfqpoint{1.693341in}{1.350831in}}%
\pgfpathlineto{\pgfqpoint{1.704452in}{1.340956in}}%
\pgfpathlineto{\pgfqpoint{1.707319in}{1.339764in}}%
\pgfpathlineto{\pgfqpoint{1.710903in}{1.338903in}}%
\pgfpathlineto{\pgfqpoint{1.714487in}{1.339189in}}%
\pgfpathlineto{\pgfqpoint{1.718788in}{1.340780in}}%
\pgfpathlineto{\pgfqpoint{1.725240in}{1.344576in}}%
\pgfpathlineto{\pgfqpoint{1.745311in}{1.358133in}}%
\pgfpathlineto{\pgfqpoint{1.751045in}{1.363600in}}%
\pgfpathlineto{\pgfqpoint{1.757497in}{1.371355in}}%
\pgfpathlineto{\pgfqpoint{1.772192in}{1.389708in}}%
\pgfpathlineto{\pgfqpoint{1.778644in}{1.395377in}}%
\pgfpathlineto{\pgfqpoint{1.787246in}{1.401289in}}%
\pgfpathlineto{\pgfqpoint{1.794772in}{1.405314in}}%
\pgfpathlineto{\pgfqpoint{1.799073in}{1.406456in}}%
\pgfpathlineto{\pgfqpoint{1.802658in}{1.406276in}}%
\pgfpathlineto{\pgfqpoint{1.806600in}{1.404877in}}%
\pgfpathlineto{\pgfqpoint{1.811976in}{1.401606in}}%
\pgfpathlineto{\pgfqpoint{1.833839in}{1.386122in}}%
\pgfpathlineto{\pgfqpoint{1.840649in}{1.378906in}}%
\pgfpathlineto{\pgfqpoint{1.852477in}{1.366096in}}%
\pgfpathlineto{\pgfqpoint{1.857495in}{1.362401in}}%
\pgfpathlineto{\pgfqpoint{1.862513in}{1.360168in}}%
\pgfpathlineto{\pgfqpoint{1.868964in}{1.358697in}}%
\pgfpathlineto{\pgfqpoint{1.877924in}{1.357631in}}%
\pgfpathlineto{\pgfqpoint{1.879358in}{1.357759in}}%
\pgfpathlineto{\pgfqpoint{1.884734in}{1.358171in}}%
\pgfpathlineto{\pgfqpoint{1.888676in}{1.359571in}}%
\pgfpathlineto{\pgfqpoint{1.892261in}{1.361980in}}%
\pgfpathlineto{\pgfqpoint{1.896203in}{1.365981in}}%
\pgfpathlineto{\pgfqpoint{1.900862in}{1.372417in}}%
\pgfpathlineto{\pgfqpoint{1.907672in}{1.383618in}}%
\pgfpathlineto{\pgfqpoint{1.914482in}{1.396805in}}%
\pgfpathlineto{\pgfqpoint{1.920934in}{1.411760in}}%
\pgfpathlineto{\pgfqpoint{1.937420in}{1.451608in}}%
\pgfpathlineto{\pgfqpoint{1.942438in}{1.459861in}}%
\pgfpathlineto{\pgfqpoint{1.946739in}{1.464767in}}%
\pgfpathlineto{\pgfqpoint{1.951040in}{1.467956in}}%
\pgfpathlineto{\pgfqpoint{1.954983in}{1.469497in}}%
\pgfpathlineto{\pgfqpoint{1.958208in}{1.469681in}}%
\pgfpathlineto{\pgfqpoint{1.961434in}{1.468712in}}%
\pgfpathlineto{\pgfqpoint{1.964301in}{1.466603in}}%
\pgfpathlineto{\pgfqpoint{1.967527in}{1.463126in}}%
\pgfpathlineto{\pgfqpoint{1.971111in}{1.457670in}}%
\pgfpathlineto{\pgfqpoint{1.975053in}{1.449707in}}%
\pgfpathlineto{\pgfqpoint{1.979354in}{1.438536in}}%
\pgfpathlineto{\pgfqpoint{1.984372in}{1.422162in}}%
\pgfpathlineto{\pgfqpoint{1.991540in}{1.394239in}}%
\pgfpathlineto{\pgfqpoint{2.006952in}{1.333469in}}%
\pgfpathlineto{\pgfqpoint{2.018421in}{1.292986in}}%
\pgfpathlineto{\pgfqpoint{2.026665in}{1.267846in}}%
\pgfpathlineto{\pgfqpoint{2.033474in}{1.250296in}}%
\pgfpathlineto{\pgfqpoint{2.038492in}{1.240157in}}%
\pgfpathlineto{\pgfqpoint{2.042076in}{1.235097in}}%
\pgfpathlineto{\pgfqpoint{2.044944in}{1.232690in}}%
\pgfpathlineto{\pgfqpoint{2.047452in}{1.231900in}}%
\pgfpathlineto{\pgfqpoint{2.049961in}{1.232357in}}%
\pgfpathlineto{\pgfqpoint{2.052470in}{1.234010in}}%
\pgfpathlineto{\pgfqpoint{2.055338in}{1.237251in}}%
\pgfpathlineto{\pgfqpoint{2.058922in}{1.243070in}}%
\pgfpathlineto{\pgfqpoint{2.063581in}{1.252901in}}%
\pgfpathlineto{\pgfqpoint{2.070749in}{1.270935in}}%
\pgfpathlineto{\pgfqpoint{2.080068in}{1.297252in}}%
\pgfpathlineto{\pgfqpoint{2.088669in}{1.324935in}}%
\pgfpathlineto{\pgfqpoint{2.125227in}{1.446158in}}%
\pgfpathlineto{\pgfqpoint{2.130245in}{1.457587in}}%
\pgfpathlineto{\pgfqpoint{2.134904in}{1.465655in}}%
\pgfpathlineto{\pgfqpoint{2.139564in}{1.471583in}}%
\pgfpathlineto{\pgfqpoint{2.143506in}{1.474828in}}%
\pgfpathlineto{\pgfqpoint{2.152466in}{1.476450in}}%
\pgfpathlineto{\pgfqpoint{2.160351in}{1.476593in}}%
\pgfpathlineto{\pgfqpoint{2.167878in}{1.477072in}}%
\pgfpathlineto{\pgfqpoint{2.171103in}{1.476167in}}%
\pgfpathlineto{\pgfqpoint{2.174329in}{1.474017in}}%
\pgfpathlineto{\pgfqpoint{2.178272in}{1.469788in}}%
\pgfpathlineto{\pgfqpoint{2.194399in}{1.450274in}}%
\pgfpathlineto{\pgfqpoint{2.200492in}{1.443877in}}%
\pgfpathlineto{\pgfqpoint{2.204434in}{1.437999in}}%
\pgfpathlineto{\pgfqpoint{2.208735in}{1.429597in}}%
\pgfpathlineto{\pgfqpoint{2.214469in}{1.415687in}}%
\pgfpathlineto{\pgfqpoint{2.228804in}{1.379819in}}%
\pgfpathlineto{\pgfqpoint{2.233822in}{1.370699in}}%
\pgfpathlineto{\pgfqpoint{2.239556in}{1.362432in}}%
\pgfpathlineto{\pgfqpoint{2.250666in}{1.348855in}}%
\pgfpathlineto{\pgfqpoint{2.255325in}{1.344446in}}%
\pgfpathlineto{\pgfqpoint{2.258909in}{1.342534in}}%
\pgfpathlineto{\pgfqpoint{2.261776in}{1.342173in}}%
\pgfpathlineto{\pgfqpoint{2.265002in}{1.342917in}}%
\pgfpathlineto{\pgfqpoint{2.269661in}{1.345249in}}%
\pgfpathlineto{\pgfqpoint{2.271452in}{1.346339in}}%
\pgfpathlineto{\pgfqpoint{2.275753in}{1.347907in}}%
\pgfpathlineto{\pgfqpoint{2.279695in}{1.348133in}}%
\pgfpathlineto{\pgfqpoint{2.285430in}{1.347063in}}%
\pgfpathlineto{\pgfqpoint{2.292597in}{1.345972in}}%
\pgfpathlineto{\pgfqpoint{2.310158in}{1.345469in}}%
\pgfpathlineto{\pgfqpoint{2.313742in}{1.346662in}}%
\pgfpathlineto{\pgfqpoint{2.316968in}{1.348941in}}%
\pgfpathlineto{\pgfqpoint{2.320551in}{1.352993in}}%
\pgfpathlineto{\pgfqpoint{2.324852in}{1.359840in}}%
\pgfpathlineto{\pgfqpoint{2.332378in}{1.374627in}}%
\pgfpathlineto{\pgfqpoint{2.337754in}{1.383898in}}%
\pgfpathlineto{\pgfqpoint{2.341696in}{1.388643in}}%
\pgfpathlineto{\pgfqpoint{2.345280in}{1.391363in}}%
\pgfpathlineto{\pgfqpoint{2.349939in}{1.393438in}}%
\pgfpathlineto{\pgfqpoint{2.356749in}{1.396526in}}%
\pgfpathlineto{\pgfqpoint{2.363200in}{1.401078in}}%
\pgfpathlineto{\pgfqpoint{2.367859in}{1.403796in}}%
\pgfpathlineto{\pgfqpoint{2.371084in}{1.404422in}}%
\pgfpathlineto{\pgfqpoint{2.374310in}{1.403802in}}%
\pgfpathlineto{\pgfqpoint{2.378610in}{1.401604in}}%
\pgfpathlineto{\pgfqpoint{2.384344in}{1.398687in}}%
\pgfpathlineto{\pgfqpoint{2.387211in}{1.398375in}}%
\pgfpathlineto{\pgfqpoint{2.389720in}{1.399100in}}%
\pgfpathlineto{\pgfqpoint{2.392587in}{1.401158in}}%
\pgfpathlineto{\pgfqpoint{2.396171in}{1.405378in}}%
\pgfpathlineto{\pgfqpoint{2.401905in}{1.414330in}}%
\pgfpathlineto{\pgfqpoint{2.409790in}{1.426210in}}%
\pgfpathlineto{\pgfqpoint{2.416599in}{1.434227in}}%
\pgfpathlineto{\pgfqpoint{2.426276in}{1.445405in}}%
\pgfpathlineto{\pgfqpoint{2.429860in}{1.451448in}}%
\pgfpathlineto{\pgfqpoint{2.433443in}{1.459857in}}%
\pgfpathlineto{\pgfqpoint{2.437027in}{1.471392in}}%
\pgfpathlineto{\pgfqpoint{2.441328in}{1.489632in}}%
\pgfpathlineto{\pgfqpoint{2.446704in}{1.518328in}}%
\pgfpathlineto{\pgfqpoint{2.454947in}{1.569463in}}%
\pgfpathlineto{\pgfqpoint{2.466416in}{1.647717in}}%
\pgfpathlineto{\pgfqpoint{2.473942in}{1.706405in}}%
\pgfpathlineto{\pgfqpoint{2.480034in}{1.763186in}}%
\pgfpathlineto{\pgfqpoint{2.486127in}{1.830897in}}%
\pgfpathlineto{\pgfqpoint{2.492936in}{1.920246in}}%
\pgfpathlineto{\pgfqpoint{2.501538in}{2.049905in}}%
\pgfpathlineto{\pgfqpoint{2.524756in}{2.419488in}}%
\pgfpathmoveto{\pgfqpoint{2.603447in}{2.419488in}}%
\pgfpathlineto{\pgfqpoint{2.610854in}{2.298885in}}%
\pgfpathlineto{\pgfqpoint{2.619098in}{2.185300in}}%
\pgfpathlineto{\pgfqpoint{2.645622in}{1.845091in}}%
\pgfpathlineto{\pgfqpoint{2.654225in}{1.754885in}}%
\pgfpathlineto{\pgfqpoint{2.659960in}{1.705721in}}%
\pgfpathlineto{\pgfqpoint{2.664261in}{1.677665in}}%
\pgfpathlineto{\pgfqpoint{2.668204in}{1.658508in}}%
\pgfpathlineto{\pgfqpoint{2.673222in}{1.639933in}}%
\pgfpathlineto{\pgfqpoint{2.683258in}{1.604238in}}%
\pgfpathlineto{\pgfqpoint{2.690427in}{1.573290in}}%
\pgfpathlineto{\pgfqpoint{2.700104in}{1.531967in}}%
\pgfpathlineto{\pgfqpoint{2.705122in}{1.515494in}}%
\pgfpathlineto{\pgfqpoint{2.713008in}{1.494294in}}%
\pgfpathlineto{\pgfqpoint{2.720893in}{1.471340in}}%
\pgfpathlineto{\pgfqpoint{2.739173in}{1.417292in}}%
\pgfpathlineto{\pgfqpoint{2.744550in}{1.403591in}}%
\pgfpathlineto{\pgfqpoint{2.747776in}{1.398078in}}%
\pgfpathlineto{\pgfqpoint{2.750285in}{1.395705in}}%
\pgfpathlineto{\pgfqpoint{2.752436in}{1.395046in}}%
\pgfpathlineto{\pgfqpoint{2.754586in}{1.395562in}}%
\pgfpathlineto{\pgfqpoint{2.757095in}{1.397427in}}%
\pgfpathlineto{\pgfqpoint{2.760321in}{1.401317in}}%
\pgfpathlineto{\pgfqpoint{2.766414in}{1.410837in}}%
\pgfpathlineto{\pgfqpoint{2.773225in}{1.420794in}}%
\pgfpathlineto{\pgfqpoint{2.783261in}{1.432891in}}%
\pgfpathlineto{\pgfqpoint{2.786128in}{1.435135in}}%
\pgfpathlineto{\pgfqpoint{2.787920in}{1.435380in}}%
\pgfpathlineto{\pgfqpoint{2.789354in}{1.434587in}}%
\pgfpathlineto{\pgfqpoint{2.791146in}{1.432012in}}%
\pgfpathlineto{\pgfqpoint{2.792939in}{1.427374in}}%
\pgfpathlineto{\pgfqpoint{2.795448in}{1.417155in}}%
\pgfpathlineto{\pgfqpoint{2.798673in}{1.398361in}}%
\pgfpathlineto{\pgfqpoint{2.806200in}{1.351580in}}%
\pgfpathlineto{\pgfqpoint{2.808351in}{1.344883in}}%
\pgfpathlineto{\pgfqpoint{2.809785in}{1.343323in}}%
\pgfpathlineto{\pgfqpoint{2.809785in}{1.343323in}}%
\pgfusepath{stroke}%
\end{pgfscope}%
\begin{pgfscope}%
\pgfpathrectangle{\pgfqpoint{0.650000in}{0.420000in}}{\pgfqpoint{2.388110in}{1.994488in}}%
\pgfusepath{clip}%
\pgfsetrectcap%
\pgfsetroundjoin%
\pgfsetlinewidth{1.003750pt}%
\definecolor{currentstroke}{rgb}{0.878431,0.509804,0.078431}%
\pgfsetstrokecolor{currentstroke}%
\pgfsetdash{}{0pt}%
\pgfpathmoveto{\pgfqpoint{0.650000in}{1.391281in}}%
\pgfpathlineto{\pgfqpoint{0.653332in}{1.394162in}}%
\pgfpathlineto{\pgfqpoint{0.657003in}{1.398963in}}%
\pgfpathlineto{\pgfqpoint{0.661975in}{1.407672in}}%
\pgfpathlineto{\pgfqpoint{0.670677in}{1.423125in}}%
\pgfpathlineto{\pgfqpoint{0.674407in}{1.427596in}}%
\pgfpathlineto{\pgfqpoint{0.677307in}{1.429534in}}%
\pgfpathlineto{\pgfqpoint{0.679794in}{1.430011in}}%
\pgfpathlineto{\pgfqpoint{0.682280in}{1.429431in}}%
\pgfpathlineto{\pgfqpoint{0.685181in}{1.427612in}}%
\pgfpathlineto{\pgfqpoint{0.689325in}{1.423561in}}%
\pgfpathlineto{\pgfqpoint{0.700928in}{1.411401in}}%
\pgfpathlineto{\pgfqpoint{0.705901in}{1.408060in}}%
\pgfpathlineto{\pgfqpoint{0.710873in}{1.406023in}}%
\pgfpathlineto{\pgfqpoint{0.716261in}{1.405019in}}%
\pgfpathlineto{\pgfqpoint{0.722891in}{1.404959in}}%
\pgfpathlineto{\pgfqpoint{0.734908in}{1.406201in}}%
\pgfpathlineto{\pgfqpoint{0.757286in}{1.408608in}}%
\pgfpathlineto{\pgfqpoint{0.771789in}{1.408440in}}%
\pgfpathlineto{\pgfqpoint{0.797482in}{1.408245in}}%
\pgfpathlineto{\pgfqpoint{0.818616in}{1.407845in}}%
\pgfpathlineto{\pgfqpoint{0.833534in}{1.407767in}}%
\pgfpathlineto{\pgfqpoint{0.849696in}{1.407306in}}%
\pgfpathlineto{\pgfqpoint{0.894450in}{1.406695in}}%
\pgfpathlineto{\pgfqpoint{0.904396in}{1.406787in}}%
\pgfpathlineto{\pgfqpoint{0.942520in}{1.404929in}}%
\pgfpathlineto{\pgfqpoint{0.956195in}{1.404865in}}%
\pgfpathlineto{\pgfqpoint{0.972357in}{1.404374in}}%
\pgfpathlineto{\pgfqpoint{0.993491in}{1.404181in}}%
\pgfpathlineto{\pgfqpoint{1.010481in}{1.403086in}}%
\pgfpathlineto{\pgfqpoint{1.023327in}{1.402648in}}%
\pgfpathlineto{\pgfqpoint{1.035759in}{1.402216in}}%
\pgfpathlineto{\pgfqpoint{1.072640in}{1.397150in}}%
\pgfpathlineto{\pgfqpoint{1.083000in}{1.393716in}}%
\pgfpathlineto{\pgfqpoint{1.096261in}{1.388410in}}%
\pgfpathlineto{\pgfqpoint{1.101648in}{1.384954in}}%
\pgfpathlineto{\pgfqpoint{1.106206in}{1.380720in}}%
\pgfpathlineto{\pgfqpoint{1.111179in}{1.374443in}}%
\pgfpathlineto{\pgfqpoint{1.116981in}{1.365406in}}%
\pgfpathlineto{\pgfqpoint{1.123611in}{1.355671in}}%
\pgfpathlineto{\pgfqpoint{1.126926in}{1.352618in}}%
\pgfpathlineto{\pgfqpoint{1.129413in}{1.351609in}}%
\pgfpathlineto{\pgfqpoint{1.131899in}{1.351802in}}%
\pgfpathlineto{\pgfqpoint{1.134385in}{1.353146in}}%
\pgfpathlineto{\pgfqpoint{1.138944in}{1.358022in}}%
\pgfpathlineto{\pgfqpoint{1.156763in}{1.380599in}}%
\pgfpathlineto{\pgfqpoint{1.163393in}{1.385787in}}%
\pgfpathlineto{\pgfqpoint{1.170437in}{1.389878in}}%
\pgfpathlineto{\pgfqpoint{1.178310in}{1.393010in}}%
\pgfpathlineto{\pgfqpoint{1.187428in}{1.395334in}}%
\pgfpathlineto{\pgfqpoint{1.196544in}{1.396384in}}%
\pgfpathlineto{\pgfqpoint{1.212706in}{1.397965in}}%
\pgfpathlineto{\pgfqpoint{1.221408in}{1.398475in}}%
\pgfpathlineto{\pgfqpoint{1.240885in}{1.398568in}}%
\pgfpathlineto{\pgfqpoint{1.251659in}{1.399031in}}%
\pgfpathlineto{\pgfqpoint{1.266163in}{1.398908in}}%
\pgfpathlineto{\pgfqpoint{1.282324in}{1.401222in}}%
\pgfpathlineto{\pgfqpoint{1.288540in}{1.400669in}}%
\pgfpathlineto{\pgfqpoint{1.308846in}{1.397673in}}%
\pgfpathlineto{\pgfqpoint{1.319620in}{1.396963in}}%
\pgfpathlineto{\pgfqpoint{1.338682in}{1.396365in}}%
\pgfpathlineto{\pgfqpoint{1.352357in}{1.397010in}}%
\pgfpathlineto{\pgfqpoint{1.375149in}{1.398363in}}%
\pgfpathlineto{\pgfqpoint{1.407057in}{1.396773in}}%
\pgfpathlineto{\pgfqpoint{1.424462in}{1.398029in}}%
\pgfpathlineto{\pgfqpoint{1.432335in}{1.396792in}}%
\pgfpathlineto{\pgfqpoint{1.443939in}{1.394836in}}%
\pgfpathlineto{\pgfqpoint{1.450569in}{1.395054in}}%
\pgfpathlineto{\pgfqpoint{1.460514in}{1.395331in}}%
\pgfpathlineto{\pgfqpoint{1.474604in}{1.395123in}}%
\pgfpathlineto{\pgfqpoint{1.505269in}{1.395227in}}%
\pgfpathlineto{\pgfqpoint{1.518944in}{1.393650in}}%
\pgfpathlineto{\pgfqpoint{1.528061in}{1.394416in}}%
\pgfpathlineto{\pgfqpoint{1.542151in}{1.395304in}}%
\pgfpathlineto{\pgfqpoint{1.550439in}{1.395421in}}%
\pgfpathlineto{\pgfqpoint{1.568257in}{1.395795in}}%
\pgfpathlineto{\pgfqpoint{1.580275in}{1.395766in}}%
\pgfpathlineto{\pgfqpoint{1.588148in}{1.394642in}}%
\pgfpathlineto{\pgfqpoint{1.598923in}{1.391677in}}%
\pgfpathlineto{\pgfqpoint{1.606796in}{1.390024in}}%
\pgfpathlineto{\pgfqpoint{1.612598in}{1.389941in}}%
\pgfpathlineto{\pgfqpoint{1.620057in}{1.391061in}}%
\pgfpathlineto{\pgfqpoint{1.633317in}{1.392984in}}%
\pgfpathlineto{\pgfqpoint{1.663154in}{1.395665in}}%
\pgfpathlineto{\pgfqpoint{1.671856in}{1.394798in}}%
\pgfpathlineto{\pgfqpoint{1.699621in}{1.390878in}}%
\pgfpathlineto{\pgfqpoint{1.706251in}{1.390719in}}%
\pgfpathlineto{\pgfqpoint{1.714125in}{1.391850in}}%
\pgfpathlineto{\pgfqpoint{1.721584in}{1.392544in}}%
\pgfpathlineto{\pgfqpoint{1.728628in}{1.391957in}}%
\pgfpathlineto{\pgfqpoint{1.739403in}{1.390852in}}%
\pgfpathlineto{\pgfqpoint{1.746033in}{1.391474in}}%
\pgfpathlineto{\pgfqpoint{1.761780in}{1.393384in}}%
\pgfpathlineto{\pgfqpoint{1.777113in}{1.393584in}}%
\pgfpathlineto{\pgfqpoint{1.784572in}{1.392495in}}%
\pgfpathlineto{\pgfqpoint{1.800733in}{1.389393in}}%
\pgfpathlineto{\pgfqpoint{1.805292in}{1.390198in}}%
\pgfpathlineto{\pgfqpoint{1.812751in}{1.392882in}}%
\pgfpathlineto{\pgfqpoint{1.818552in}{1.394467in}}%
\pgfpathlineto{\pgfqpoint{1.823111in}{1.394579in}}%
\pgfpathlineto{\pgfqpoint{1.828498in}{1.393471in}}%
\pgfpathlineto{\pgfqpoint{1.843416in}{1.389391in}}%
\pgfpathlineto{\pgfqpoint{1.849218in}{1.389633in}}%
\pgfpathlineto{\pgfqpoint{1.857920in}{1.389946in}}%
\pgfpathlineto{\pgfqpoint{1.862893in}{1.388944in}}%
\pgfpathlineto{\pgfqpoint{1.868280in}{1.386660in}}%
\pgfpathlineto{\pgfqpoint{1.882369in}{1.379555in}}%
\pgfpathlineto{\pgfqpoint{1.885270in}{1.379869in}}%
\pgfpathlineto{\pgfqpoint{1.888585in}{1.381428in}}%
\pgfpathlineto{\pgfqpoint{1.893144in}{1.385026in}}%
\pgfpathlineto{\pgfqpoint{1.900188in}{1.390447in}}%
\pgfpathlineto{\pgfqpoint{1.904332in}{1.392099in}}%
\pgfpathlineto{\pgfqpoint{1.908062in}{1.392397in}}%
\pgfpathlineto{\pgfqpoint{1.913035in}{1.391559in}}%
\pgfpathlineto{\pgfqpoint{1.926295in}{1.388661in}}%
\pgfpathlineto{\pgfqpoint{1.932926in}{1.388715in}}%
\pgfpathlineto{\pgfqpoint{1.940799in}{1.389957in}}%
\pgfpathlineto{\pgfqpoint{1.956132in}{1.392559in}}%
\pgfpathlineto{\pgfqpoint{1.960690in}{1.391984in}}%
\pgfpathlineto{\pgfqpoint{1.965249in}{1.390246in}}%
\pgfpathlineto{\pgfqpoint{1.970636in}{1.386936in}}%
\pgfpathlineto{\pgfqpoint{1.978509in}{1.381964in}}%
\pgfpathlineto{\pgfqpoint{1.982653in}{1.380652in}}%
\pgfpathlineto{\pgfqpoint{1.986797in}{1.380435in}}%
\pgfpathlineto{\pgfqpoint{1.992598in}{1.381367in}}%
\pgfpathlineto{\pgfqpoint{2.007931in}{1.385273in}}%
\pgfpathlineto{\pgfqpoint{2.018291in}{1.388133in}}%
\pgfpathlineto{\pgfqpoint{2.023678in}{1.388271in}}%
\pgfpathlineto{\pgfqpoint{2.065532in}{1.385910in}}%
\pgfpathlineto{\pgfqpoint{2.071748in}{1.387242in}}%
\pgfpathlineto{\pgfqpoint{2.080865in}{1.389266in}}%
\pgfpathlineto{\pgfqpoint{2.085838in}{1.389064in}}%
\pgfpathlineto{\pgfqpoint{2.092468in}{1.387506in}}%
\pgfpathlineto{\pgfqpoint{2.103657in}{1.384847in}}%
\pgfpathlineto{\pgfqpoint{2.118161in}{1.383945in}}%
\pgfpathlineto{\pgfqpoint{2.123133in}{1.384969in}}%
\pgfpathlineto{\pgfqpoint{2.128935in}{1.387441in}}%
\pgfpathlineto{\pgfqpoint{2.134322in}{1.389465in}}%
\pgfpathlineto{\pgfqpoint{2.137637in}{1.389607in}}%
\pgfpathlineto{\pgfqpoint{2.140952in}{1.388579in}}%
\pgfpathlineto{\pgfqpoint{2.145925in}{1.385578in}}%
\pgfpathlineto{\pgfqpoint{2.152141in}{1.382059in}}%
\pgfpathlineto{\pgfqpoint{2.155456in}{1.381447in}}%
\pgfpathlineto{\pgfqpoint{2.159600in}{1.381880in}}%
\pgfpathlineto{\pgfqpoint{2.171203in}{1.383653in}}%
\pgfpathlineto{\pgfqpoint{2.183635in}{1.384759in}}%
\pgfpathlineto{\pgfqpoint{2.190265in}{1.386712in}}%
\pgfpathlineto{\pgfqpoint{2.198967in}{1.389412in}}%
\pgfpathlineto{\pgfqpoint{2.203111in}{1.389446in}}%
\pgfpathlineto{\pgfqpoint{2.208084in}{1.388271in}}%
\pgfpathlineto{\pgfqpoint{2.214714in}{1.386724in}}%
\pgfpathlineto{\pgfqpoint{2.219273in}{1.386923in}}%
\pgfpathlineto{\pgfqpoint{2.226317in}{1.387413in}}%
\pgfpathlineto{\pgfqpoint{2.230047in}{1.386334in}}%
\pgfpathlineto{\pgfqpoint{2.243722in}{1.380788in}}%
\pgfpathlineto{\pgfqpoint{2.248280in}{1.381091in}}%
\pgfpathlineto{\pgfqpoint{2.263199in}{1.382921in}}%
\pgfpathlineto{\pgfqpoint{2.279360in}{1.382574in}}%
\pgfpathlineto{\pgfqpoint{2.294278in}{1.385140in}}%
\pgfpathlineto{\pgfqpoint{2.298837in}{1.383892in}}%
\pgfpathlineto{\pgfqpoint{2.310854in}{1.379557in}}%
\pgfpathlineto{\pgfqpoint{2.314584in}{1.379823in}}%
\pgfpathlineto{\pgfqpoint{2.319971in}{1.381511in}}%
\pgfpathlineto{\pgfqpoint{2.327844in}{1.383889in}}%
\pgfpathlineto{\pgfqpoint{2.332817in}{1.384193in}}%
\pgfpathlineto{\pgfqpoint{2.337376in}{1.383284in}}%
\pgfpathlineto{\pgfqpoint{2.343177in}{1.380712in}}%
\pgfpathlineto{\pgfqpoint{2.348979in}{1.378475in}}%
\pgfpathlineto{\pgfqpoint{2.352708in}{1.378278in}}%
\pgfpathlineto{\pgfqpoint{2.356438in}{1.379213in}}%
\pgfpathlineto{\pgfqpoint{2.362654in}{1.382086in}}%
\pgfpathlineto{\pgfqpoint{2.376743in}{1.388730in}}%
\pgfpathlineto{\pgfqpoint{2.380887in}{1.389292in}}%
\pgfpathlineto{\pgfqpoint{2.385032in}{1.388734in}}%
\pgfpathlineto{\pgfqpoint{2.390004in}{1.386910in}}%
\pgfpathlineto{\pgfqpoint{2.403679in}{1.381089in}}%
\pgfpathlineto{\pgfqpoint{2.407823in}{1.380976in}}%
\pgfpathlineto{\pgfqpoint{2.419426in}{1.381624in}}%
\pgfpathlineto{\pgfqpoint{2.423570in}{1.379934in}}%
\pgfpathlineto{\pgfqpoint{2.431029in}{1.377114in}}%
\pgfpathlineto{\pgfqpoint{2.433930in}{1.377548in}}%
\pgfpathlineto{\pgfqpoint{2.436831in}{1.379117in}}%
\pgfpathlineto{\pgfqpoint{2.441389in}{1.383145in}}%
\pgfpathlineto{\pgfqpoint{2.447191in}{1.388054in}}%
\pgfpathlineto{\pgfqpoint{2.450506in}{1.389468in}}%
\pgfpathlineto{\pgfqpoint{2.453821in}{1.389623in}}%
\pgfpathlineto{\pgfqpoint{2.457965in}{1.388449in}}%
\pgfpathlineto{\pgfqpoint{2.470397in}{1.384012in}}%
\pgfpathlineto{\pgfqpoint{2.474127in}{1.384160in}}%
\pgfpathlineto{\pgfqpoint{2.477442in}{1.385382in}}%
\pgfpathlineto{\pgfqpoint{2.481586in}{1.388242in}}%
\pgfpathlineto{\pgfqpoint{2.495261in}{1.398698in}}%
\pgfpathlineto{\pgfqpoint{2.500648in}{1.403074in}}%
\pgfpathlineto{\pgfqpoint{2.505620in}{1.408789in}}%
\pgfpathlineto{\pgfqpoint{2.512665in}{1.418874in}}%
\pgfpathlineto{\pgfqpoint{2.520539in}{1.431917in}}%
\pgfpathlineto{\pgfqpoint{2.526340in}{1.443675in}}%
\pgfpathlineto{\pgfqpoint{2.530899in}{1.455378in}}%
\pgfpathlineto{\pgfqpoint{2.535871in}{1.471437in}}%
\pgfpathlineto{\pgfqpoint{2.546231in}{1.506574in}}%
\pgfpathlineto{\pgfqpoint{2.549132in}{1.512114in}}%
\pgfpathlineto{\pgfqpoint{2.551204in}{1.513885in}}%
\pgfpathlineto{\pgfqpoint{2.552862in}{1.513945in}}%
\pgfpathlineto{\pgfqpoint{2.554519in}{1.512762in}}%
\pgfpathlineto{\pgfqpoint{2.556591in}{1.509613in}}%
\pgfpathlineto{\pgfqpoint{2.559491in}{1.502419in}}%
\pgfpathlineto{\pgfqpoint{2.563221in}{1.489706in}}%
\pgfpathlineto{\pgfqpoint{2.581452in}{1.422265in}}%
\pgfpathlineto{\pgfqpoint{2.586011in}{1.410979in}}%
\pgfpathlineto{\pgfqpoint{2.589740in}{1.404557in}}%
\pgfpathlineto{\pgfqpoint{2.593055in}{1.400865in}}%
\pgfpathlineto{\pgfqpoint{2.596784in}{1.398346in}}%
\pgfpathlineto{\pgfqpoint{2.602585in}{1.396006in}}%
\pgfpathlineto{\pgfqpoint{2.609215in}{1.393094in}}%
\pgfpathlineto{\pgfqpoint{2.613773in}{1.389819in}}%
\pgfpathlineto{\pgfqpoint{2.618331in}{1.385127in}}%
\pgfpathlineto{\pgfqpoint{2.630347in}{1.371758in}}%
\pgfpathlineto{\pgfqpoint{2.633662in}{1.370234in}}%
\pgfpathlineto{\pgfqpoint{2.636977in}{1.370077in}}%
\pgfpathlineto{\pgfqpoint{2.641121in}{1.371161in}}%
\pgfpathlineto{\pgfqpoint{2.649408in}{1.373403in}}%
\pgfpathlineto{\pgfqpoint{2.657281in}{1.374079in}}%
\pgfpathlineto{\pgfqpoint{2.663082in}{1.373450in}}%
\pgfpathlineto{\pgfqpoint{2.668054in}{1.373068in}}%
\pgfpathlineto{\pgfqpoint{2.670955in}{1.373884in}}%
\pgfpathlineto{\pgfqpoint{2.675098in}{1.376445in}}%
\pgfpathlineto{\pgfqpoint{2.679242in}{1.378637in}}%
\pgfpathlineto{\pgfqpoint{2.681728in}{1.378733in}}%
\pgfpathlineto{\pgfqpoint{2.684214in}{1.377727in}}%
\pgfpathlineto{\pgfqpoint{2.693744in}{1.372372in}}%
\pgfpathlineto{\pgfqpoint{2.696645in}{1.372832in}}%
\pgfpathlineto{\pgfqpoint{2.700789in}{1.374923in}}%
\pgfpathlineto{\pgfqpoint{2.714048in}{1.382295in}}%
\pgfpathlineto{\pgfqpoint{2.717363in}{1.382780in}}%
\pgfpathlineto{\pgfqpoint{2.719849in}{1.382170in}}%
\pgfpathlineto{\pgfqpoint{2.722750in}{1.380274in}}%
\pgfpathlineto{\pgfqpoint{2.727722in}{1.375268in}}%
\pgfpathlineto{\pgfqpoint{2.731451in}{1.372141in}}%
\pgfpathlineto{\pgfqpoint{2.733938in}{1.371209in}}%
\pgfpathlineto{\pgfqpoint{2.736424in}{1.371387in}}%
\pgfpathlineto{\pgfqpoint{2.739739in}{1.372983in}}%
\pgfpathlineto{\pgfqpoint{2.748026in}{1.377679in}}%
\pgfpathlineto{\pgfqpoint{2.760042in}{1.381885in}}%
\pgfpathlineto{\pgfqpoint{2.764186in}{1.385786in}}%
\pgfpathlineto{\pgfqpoint{2.772888in}{1.394581in}}%
\pgfpathlineto{\pgfqpoint{2.775374in}{1.395510in}}%
\pgfpathlineto{\pgfqpoint{2.777860in}{1.395203in}}%
\pgfpathlineto{\pgfqpoint{2.780346in}{1.393651in}}%
\pgfpathlineto{\pgfqpoint{2.783661in}{1.390033in}}%
\pgfpathlineto{\pgfqpoint{2.794020in}{1.377312in}}%
\pgfpathlineto{\pgfqpoint{2.797335in}{1.375527in}}%
\pgfpathlineto{\pgfqpoint{2.800650in}{1.375131in}}%
\pgfpathlineto{\pgfqpoint{2.805208in}{1.376004in}}%
\pgfpathlineto{\pgfqpoint{2.811009in}{1.376898in}}%
\pgfpathlineto{\pgfqpoint{2.815153in}{1.376253in}}%
\pgfpathlineto{\pgfqpoint{2.823026in}{1.374495in}}%
\pgfpathlineto{\pgfqpoint{2.826340in}{1.375040in}}%
\pgfpathlineto{\pgfqpoint{2.831727in}{1.377402in}}%
\pgfpathlineto{\pgfqpoint{2.835871in}{1.378725in}}%
\pgfpathlineto{\pgfqpoint{2.839186in}{1.378628in}}%
\pgfpathlineto{\pgfqpoint{2.844158in}{1.377007in}}%
\pgfpathlineto{\pgfqpoint{2.848301in}{1.376119in}}%
\pgfpathlineto{\pgfqpoint{2.851616in}{1.376658in}}%
\pgfpathlineto{\pgfqpoint{2.855760in}{1.378682in}}%
\pgfpathlineto{\pgfqpoint{2.861561in}{1.381368in}}%
\pgfpathlineto{\pgfqpoint{2.865290in}{1.381706in}}%
\pgfpathlineto{\pgfqpoint{2.870677in}{1.380822in}}%
\pgfpathlineto{\pgfqpoint{2.880622in}{1.379318in}}%
\pgfpathlineto{\pgfqpoint{2.886423in}{1.379265in}}%
\pgfpathlineto{\pgfqpoint{2.889738in}{1.380244in}}%
\pgfpathlineto{\pgfqpoint{2.894710in}{1.383444in}}%
\pgfpathlineto{\pgfqpoint{2.900097in}{1.386952in}}%
\pgfpathlineto{\pgfqpoint{2.902997in}{1.387444in}}%
\pgfpathlineto{\pgfqpoint{2.905483in}{1.386710in}}%
\pgfpathlineto{\pgfqpoint{2.908384in}{1.384528in}}%
\pgfpathlineto{\pgfqpoint{2.914185in}{1.378119in}}%
\pgfpathlineto{\pgfqpoint{2.918329in}{1.374360in}}%
\pgfpathlineto{\pgfqpoint{2.922058in}{1.372539in}}%
\pgfpathlineto{\pgfqpoint{2.931588in}{1.369161in}}%
\pgfpathlineto{\pgfqpoint{2.937804in}{1.364120in}}%
\pgfpathlineto{\pgfqpoint{2.940704in}{1.362622in}}%
\pgfpathlineto{\pgfqpoint{2.942776in}{1.362624in}}%
\pgfpathlineto{\pgfqpoint{2.944848in}{1.363740in}}%
\pgfpathlineto{\pgfqpoint{2.947334in}{1.366527in}}%
\pgfpathlineto{\pgfqpoint{2.951063in}{1.372836in}}%
\pgfpathlineto{\pgfqpoint{2.957279in}{1.383356in}}%
\pgfpathlineto{\pgfqpoint{2.960179in}{1.386333in}}%
\pgfpathlineto{\pgfqpoint{2.963080in}{1.387692in}}%
\pgfpathlineto{\pgfqpoint{2.965566in}{1.387661in}}%
\pgfpathlineto{\pgfqpoint{2.968881in}{1.386455in}}%
\pgfpathlineto{\pgfqpoint{2.975511in}{1.383670in}}%
\pgfpathlineto{\pgfqpoint{2.978411in}{1.383731in}}%
\pgfpathlineto{\pgfqpoint{2.981312in}{1.384952in}}%
\pgfpathlineto{\pgfqpoint{2.985870in}{1.388276in}}%
\pgfpathlineto{\pgfqpoint{2.991256in}{1.391951in}}%
\pgfpathlineto{\pgfqpoint{2.995400in}{1.393322in}}%
\pgfpathlineto{\pgfqpoint{2.999129in}{1.393430in}}%
\pgfpathlineto{\pgfqpoint{3.002030in}{1.392499in}}%
\pgfpathlineto{\pgfqpoint{3.004516in}{1.390546in}}%
\pgfpathlineto{\pgfqpoint{3.007002in}{1.386910in}}%
\pgfpathlineto{\pgfqpoint{3.009488in}{1.380948in}}%
\pgfpathlineto{\pgfqpoint{3.011975in}{1.371955in}}%
\pgfpathlineto{\pgfqpoint{3.014875in}{1.357254in}}%
\pgfpathlineto{\pgfqpoint{3.019018in}{1.329753in}}%
\pgfpathlineto{\pgfqpoint{3.025333in}{1.287778in}}%
\pgfpathlineto{\pgfqpoint{3.026849in}{1.280157in}}%
\pgfpathlineto{\pgfqpoint{3.026849in}{1.280157in}}%
\pgfusepath{stroke}%
\end{pgfscope}%
\begin{pgfscope}%
\pgfpathrectangle{\pgfqpoint{0.650000in}{0.420000in}}{\pgfqpoint{2.388110in}{1.994488in}}%
\pgfusepath{clip}%
\pgfsetrectcap%
\pgfsetroundjoin%
\pgfsetlinewidth{1.003750pt}%
\definecolor{currentstroke}{rgb}{0.501961,0.450980,0.674510}%
\pgfsetstrokecolor{currentstroke}%
\pgfsetdash{}{0pt}%
\pgfpathmoveto{\pgfqpoint{0.650000in}{1.391096in}}%
\pgfpathlineto{\pgfqpoint{0.653348in}{1.393939in}}%
\pgfpathlineto{\pgfqpoint{0.657036in}{1.398696in}}%
\pgfpathlineto{\pgfqpoint{0.662031in}{1.407393in}}%
\pgfpathlineto{\pgfqpoint{0.671191in}{1.423668in}}%
\pgfpathlineto{\pgfqpoint{0.674939in}{1.428015in}}%
\pgfpathlineto{\pgfqpoint{0.677853in}{1.429766in}}%
\pgfpathlineto{\pgfqpoint{0.680351in}{1.430020in}}%
\pgfpathlineto{\pgfqpoint{0.682849in}{1.429174in}}%
\pgfpathlineto{\pgfqpoint{0.685764in}{1.427015in}}%
\pgfpathlineto{\pgfqpoint{0.690344in}{1.421990in}}%
\pgfpathlineto{\pgfqpoint{0.699504in}{1.411625in}}%
\pgfpathlineto{\pgfqpoint{0.704084in}{1.408148in}}%
\pgfpathlineto{\pgfqpoint{0.708664in}{1.406028in}}%
\pgfpathlineto{\pgfqpoint{0.713660in}{1.404928in}}%
\pgfpathlineto{\pgfqpoint{0.720322in}{1.404693in}}%
\pgfpathlineto{\pgfqpoint{0.730731in}{1.405580in}}%
\pgfpathlineto{\pgfqpoint{0.754047in}{1.407895in}}%
\pgfpathlineto{\pgfqpoint{0.777779in}{1.407123in}}%
\pgfpathlineto{\pgfqpoint{0.791935in}{1.405963in}}%
\pgfpathlineto{\pgfqpoint{0.815251in}{1.403551in}}%
\pgfpathlineto{\pgfqpoint{0.841898in}{1.402966in}}%
\pgfpathlineto{\pgfqpoint{0.871876in}{1.404288in}}%
\pgfpathlineto{\pgfqpoint{0.891861in}{1.404164in}}%
\pgfpathlineto{\pgfqpoint{0.914345in}{1.404454in}}%
\pgfpathlineto{\pgfqpoint{0.927252in}{1.403082in}}%
\pgfpathlineto{\pgfqpoint{0.955148in}{1.399261in}}%
\pgfpathlineto{\pgfqpoint{0.983877in}{1.398548in}}%
\pgfpathlineto{\pgfqpoint{1.000115in}{1.400520in}}%
\pgfpathlineto{\pgfqpoint{1.025096in}{1.404310in}}%
\pgfpathlineto{\pgfqpoint{1.040918in}{1.407420in}}%
\pgfpathlineto{\pgfqpoint{1.050494in}{1.407917in}}%
\pgfpathlineto{\pgfqpoint{1.065483in}{1.407431in}}%
\pgfpathlineto{\pgfqpoint{1.079639in}{1.405991in}}%
\pgfpathlineto{\pgfqpoint{1.088799in}{1.403997in}}%
\pgfpathlineto{\pgfqpoint{1.095877in}{1.401296in}}%
\pgfpathlineto{\pgfqpoint{1.102539in}{1.397507in}}%
\pgfpathlineto{\pgfqpoint{1.108368in}{1.392961in}}%
\pgfpathlineto{\pgfqpoint{1.114197in}{1.386985in}}%
\pgfpathlineto{\pgfqpoint{1.128353in}{1.371371in}}%
\pgfpathlineto{\pgfqpoint{1.131267in}{1.370404in}}%
\pgfpathlineto{\pgfqpoint{1.133766in}{1.370753in}}%
\pgfpathlineto{\pgfqpoint{1.136264in}{1.372179in}}%
\pgfpathlineto{\pgfqpoint{1.139595in}{1.375533in}}%
\pgfpathlineto{\pgfqpoint{1.144591in}{1.382416in}}%
\pgfpathlineto{\pgfqpoint{1.156665in}{1.399503in}}%
\pgfpathlineto{\pgfqpoint{1.163327in}{1.406713in}}%
\pgfpathlineto{\pgfqpoint{1.169156in}{1.411399in}}%
\pgfpathlineto{\pgfqpoint{1.175817in}{1.415243in}}%
\pgfpathlineto{\pgfqpoint{1.187892in}{1.420671in}}%
\pgfpathlineto{\pgfqpoint{1.194970in}{1.422818in}}%
\pgfpathlineto{\pgfqpoint{1.202881in}{1.423945in}}%
\pgfpathlineto{\pgfqpoint{1.227862in}{1.426940in}}%
\pgfpathlineto{\pgfqpoint{1.236606in}{1.427720in}}%
\pgfpathlineto{\pgfqpoint{1.243268in}{1.427131in}}%
\pgfpathlineto{\pgfqpoint{1.252011in}{1.425118in}}%
\pgfpathlineto{\pgfqpoint{1.276576in}{1.418928in}}%
\pgfpathlineto{\pgfqpoint{1.282822in}{1.418739in}}%
\pgfpathlineto{\pgfqpoint{1.288651in}{1.419660in}}%
\pgfpathlineto{\pgfqpoint{1.296978in}{1.422308in}}%
\pgfpathlineto{\pgfqpoint{1.306138in}{1.425008in}}%
\pgfpathlineto{\pgfqpoint{1.312383in}{1.425631in}}%
\pgfpathlineto{\pgfqpoint{1.319461in}{1.425144in}}%
\pgfpathlineto{\pgfqpoint{1.349439in}{1.420830in}}%
\pgfpathlineto{\pgfqpoint{1.371506in}{1.415281in}}%
\pgfpathlineto{\pgfqpoint{1.394406in}{1.408391in}}%
\pgfpathlineto{\pgfqpoint{1.403566in}{1.405438in}}%
\pgfpathlineto{\pgfqpoint{1.420637in}{1.399199in}}%
\pgfpathlineto{\pgfqpoint{1.427298in}{1.398274in}}%
\pgfpathlineto{\pgfqpoint{1.447284in}{1.397585in}}%
\pgfpathlineto{\pgfqpoint{1.452696in}{1.398197in}}%
\pgfpathlineto{\pgfqpoint{1.458109in}{1.399907in}}%
\pgfpathlineto{\pgfqpoint{1.476013in}{1.406668in}}%
\pgfpathlineto{\pgfqpoint{1.483507in}{1.407815in}}%
\pgfpathlineto{\pgfqpoint{1.493916in}{1.408407in}}%
\pgfpathlineto{\pgfqpoint{1.504325in}{1.407740in}}%
\pgfpathlineto{\pgfqpoint{1.512652in}{1.407631in}}%
\pgfpathlineto{\pgfqpoint{1.530139in}{1.408441in}}%
\pgfpathlineto{\pgfqpoint{1.538466in}{1.406645in}}%
\pgfpathlineto{\pgfqpoint{1.547626in}{1.405020in}}%
\pgfpathlineto{\pgfqpoint{1.568028in}{1.403488in}}%
\pgfpathlineto{\pgfqpoint{1.578437in}{1.403646in}}%
\pgfpathlineto{\pgfqpoint{1.607582in}{1.405767in}}%
\pgfpathlineto{\pgfqpoint{1.630065in}{1.409655in}}%
\pgfpathlineto{\pgfqpoint{1.663374in}{1.412803in}}%
\pgfpathlineto{\pgfqpoint{1.673783in}{1.413815in}}%
\pgfpathlineto{\pgfqpoint{1.679196in}{1.415552in}}%
\pgfpathlineto{\pgfqpoint{1.687939in}{1.419914in}}%
\pgfpathlineto{\pgfqpoint{1.693768in}{1.422152in}}%
\pgfpathlineto{\pgfqpoint{1.700014in}{1.423257in}}%
\pgfpathlineto{\pgfqpoint{1.710006in}{1.424912in}}%
\pgfpathlineto{\pgfqpoint{1.720831in}{1.428357in}}%
\pgfpathlineto{\pgfqpoint{1.729991in}{1.430760in}}%
\pgfpathlineto{\pgfqpoint{1.735820in}{1.431197in}}%
\pgfpathlineto{\pgfqpoint{1.743315in}{1.430452in}}%
\pgfpathlineto{\pgfqpoint{1.750393in}{1.430131in}}%
\pgfpathlineto{\pgfqpoint{1.761219in}{1.431235in}}%
\pgfpathlineto{\pgfqpoint{1.784535in}{1.433679in}}%
\pgfpathlineto{\pgfqpoint{1.792029in}{1.436328in}}%
\pgfpathlineto{\pgfqpoint{1.801189in}{1.439638in}}%
\pgfpathlineto{\pgfqpoint{1.806185in}{1.440209in}}%
\pgfpathlineto{\pgfqpoint{1.824921in}{1.440755in}}%
\pgfpathlineto{\pgfqpoint{1.835747in}{1.442534in}}%
\pgfpathlineto{\pgfqpoint{1.851568in}{1.444374in}}%
\pgfpathlineto{\pgfqpoint{1.858230in}{1.444783in}}%
\pgfpathlineto{\pgfqpoint{1.869472in}{1.443198in}}%
\pgfpathlineto{\pgfqpoint{1.876550in}{1.442832in}}%
\pgfpathlineto{\pgfqpoint{1.883628in}{1.443725in}}%
\pgfpathlineto{\pgfqpoint{1.910691in}{1.448664in}}%
\pgfpathlineto{\pgfqpoint{1.918602in}{1.449952in}}%
\pgfpathlineto{\pgfqpoint{1.939004in}{1.454803in}}%
\pgfpathlineto{\pgfqpoint{1.949829in}{1.454960in}}%
\pgfpathlineto{\pgfqpoint{1.959822in}{1.453283in}}%
\pgfpathlineto{\pgfqpoint{1.985219in}{1.446146in}}%
\pgfpathlineto{\pgfqpoint{2.003539in}{1.438133in}}%
\pgfpathlineto{\pgfqpoint{2.013948in}{1.434807in}}%
\pgfpathlineto{\pgfqpoint{2.019777in}{1.431451in}}%
\pgfpathlineto{\pgfqpoint{2.045592in}{1.414530in}}%
\pgfpathlineto{\pgfqpoint{2.053086in}{1.411186in}}%
\pgfpathlineto{\pgfqpoint{2.058915in}{1.409840in}}%
\pgfpathlineto{\pgfqpoint{2.075153in}{1.407948in}}%
\pgfpathlineto{\pgfqpoint{2.080566in}{1.408432in}}%
\pgfpathlineto{\pgfqpoint{2.086395in}{1.410151in}}%
\pgfpathlineto{\pgfqpoint{2.123034in}{1.423024in}}%
\pgfpathlineto{\pgfqpoint{2.127614in}{1.422727in}}%
\pgfpathlineto{\pgfqpoint{2.133027in}{1.421147in}}%
\pgfpathlineto{\pgfqpoint{2.143436in}{1.418019in}}%
\pgfpathlineto{\pgfqpoint{2.152179in}{1.415464in}}%
\pgfpathlineto{\pgfqpoint{2.172581in}{1.406966in}}%
\pgfpathlineto{\pgfqpoint{2.177577in}{1.402910in}}%
\pgfpathlineto{\pgfqpoint{2.183406in}{1.396604in}}%
\pgfpathlineto{\pgfqpoint{2.199644in}{1.377823in}}%
\pgfpathlineto{\pgfqpoint{2.203808in}{1.375568in}}%
\pgfpathlineto{\pgfqpoint{2.210053in}{1.373789in}}%
\pgfpathlineto{\pgfqpoint{2.214633in}{1.372246in}}%
\pgfpathlineto{\pgfqpoint{2.220462in}{1.369052in}}%
\pgfpathlineto{\pgfqpoint{2.239198in}{1.357839in}}%
\pgfpathlineto{\pgfqpoint{2.244611in}{1.355468in}}%
\pgfpathlineto{\pgfqpoint{2.249607in}{1.354506in}}%
\pgfpathlineto{\pgfqpoint{2.261682in}{1.352935in}}%
\pgfpathlineto{\pgfqpoint{2.268759in}{1.350090in}}%
\pgfpathlineto{\pgfqpoint{2.277919in}{1.346610in}}%
\pgfpathlineto{\pgfqpoint{2.285830in}{1.344872in}}%
\pgfpathlineto{\pgfqpoint{2.291659in}{1.344614in}}%
\pgfpathlineto{\pgfqpoint{2.296239in}{1.345499in}}%
\pgfpathlineto{\pgfqpoint{2.303318in}{1.348138in}}%
\pgfpathlineto{\pgfqpoint{2.307481in}{1.349524in}}%
\pgfpathlineto{\pgfqpoint{2.312478in}{1.350136in}}%
\pgfpathlineto{\pgfqpoint{2.325801in}{1.350756in}}%
\pgfpathlineto{\pgfqpoint{2.331213in}{1.352780in}}%
\pgfpathlineto{\pgfqpoint{2.338292in}{1.356779in}}%
\pgfpathlineto{\pgfqpoint{2.346202in}{1.360948in}}%
\pgfpathlineto{\pgfqpoint{2.359526in}{1.366977in}}%
\pgfpathlineto{\pgfqpoint{2.385757in}{1.386000in}}%
\pgfpathlineto{\pgfqpoint{2.395333in}{1.394041in}}%
\pgfpathlineto{\pgfqpoint{2.400330in}{1.396667in}}%
\pgfpathlineto{\pgfqpoint{2.404909in}{1.397840in}}%
\pgfpathlineto{\pgfqpoint{2.417400in}{1.400132in}}%
\pgfpathlineto{\pgfqpoint{2.430307in}{1.405212in}}%
\pgfpathlineto{\pgfqpoint{2.433638in}{1.404468in}}%
\pgfpathlineto{\pgfqpoint{2.441966in}{1.401868in}}%
\pgfpathlineto{\pgfqpoint{2.444880in}{1.402573in}}%
\pgfpathlineto{\pgfqpoint{2.447794in}{1.404556in}}%
\pgfpathlineto{\pgfqpoint{2.451542in}{1.408586in}}%
\pgfpathlineto{\pgfqpoint{2.461118in}{1.419583in}}%
\pgfpathlineto{\pgfqpoint{2.476523in}{1.434037in}}%
\pgfpathlineto{\pgfqpoint{2.482352in}{1.442303in}}%
\pgfpathlineto{\pgfqpoint{2.489430in}{1.454457in}}%
\pgfpathlineto{\pgfqpoint{2.507333in}{1.488572in}}%
\pgfpathlineto{\pgfqpoint{2.513163in}{1.502558in}}%
\pgfpathlineto{\pgfqpoint{2.525654in}{1.536762in}}%
\pgfpathlineto{\pgfqpoint{2.536479in}{1.567927in}}%
\pgfpathlineto{\pgfqpoint{2.541059in}{1.585004in}}%
\pgfpathlineto{\pgfqpoint{2.546888in}{1.611843in}}%
\pgfpathlineto{\pgfqpoint{2.553549in}{1.642146in}}%
\pgfpathlineto{\pgfqpoint{2.556880in}{1.652345in}}%
\pgfpathlineto{\pgfqpoint{2.558962in}{1.655912in}}%
\pgfpathlineto{\pgfqpoint{2.560628in}{1.657013in}}%
\pgfpathlineto{\pgfqpoint{2.561877in}{1.656739in}}%
\pgfpathlineto{\pgfqpoint{2.563542in}{1.655083in}}%
\pgfpathlineto{\pgfqpoint{2.565624in}{1.650964in}}%
\pgfpathlineto{\pgfqpoint{2.568538in}{1.641948in}}%
\pgfpathlineto{\pgfqpoint{2.572701in}{1.624568in}}%
\pgfpathlineto{\pgfqpoint{2.585607in}{1.568159in}}%
\pgfpathlineto{\pgfqpoint{2.592268in}{1.545646in}}%
\pgfpathlineto{\pgfqpoint{2.599762in}{1.524443in}}%
\pgfpathlineto{\pgfqpoint{2.607672in}{1.505266in}}%
\pgfpathlineto{\pgfqpoint{2.613917in}{1.492755in}}%
\pgfpathlineto{\pgfqpoint{2.620162in}{1.482913in}}%
\pgfpathlineto{\pgfqpoint{2.627240in}{1.471648in}}%
\pgfpathlineto{\pgfqpoint{2.631819in}{1.461911in}}%
\pgfpathlineto{\pgfqpoint{2.646391in}{1.428265in}}%
\pgfpathlineto{\pgfqpoint{2.651386in}{1.421277in}}%
\pgfpathlineto{\pgfqpoint{2.662211in}{1.406662in}}%
\pgfpathlineto{\pgfqpoint{2.673452in}{1.389882in}}%
\pgfpathlineto{\pgfqpoint{2.682195in}{1.376203in}}%
\pgfpathlineto{\pgfqpoint{2.690521in}{1.364773in}}%
\pgfpathlineto{\pgfqpoint{2.694684in}{1.356175in}}%
\pgfpathlineto{\pgfqpoint{2.700929in}{1.343134in}}%
\pgfpathlineto{\pgfqpoint{2.703844in}{1.339513in}}%
\pgfpathlineto{\pgfqpoint{2.706341in}{1.338113in}}%
\pgfpathlineto{\pgfqpoint{2.708839in}{1.338000in}}%
\pgfpathlineto{\pgfqpoint{2.712586in}{1.339130in}}%
\pgfpathlineto{\pgfqpoint{2.717998in}{1.340640in}}%
\pgfpathlineto{\pgfqpoint{2.721329in}{1.340306in}}%
\pgfpathlineto{\pgfqpoint{2.724243in}{1.338785in}}%
\pgfpathlineto{\pgfqpoint{2.727158in}{1.335932in}}%
\pgfpathlineto{\pgfqpoint{2.730905in}{1.330395in}}%
\pgfpathlineto{\pgfqpoint{2.740480in}{1.314983in}}%
\pgfpathlineto{\pgfqpoint{2.742978in}{1.313228in}}%
\pgfpathlineto{\pgfqpoint{2.745476in}{1.312790in}}%
\pgfpathlineto{\pgfqpoint{2.748807in}{1.313575in}}%
\pgfpathlineto{\pgfqpoint{2.753802in}{1.314761in}}%
\pgfpathlineto{\pgfqpoint{2.756717in}{1.314248in}}%
\pgfpathlineto{\pgfqpoint{2.761713in}{1.311835in}}%
\pgfpathlineto{\pgfqpoint{2.765043in}{1.310816in}}%
\pgfpathlineto{\pgfqpoint{2.767541in}{1.311180in}}%
\pgfpathlineto{\pgfqpoint{2.770039in}{1.312740in}}%
\pgfpathlineto{\pgfqpoint{2.772953in}{1.315928in}}%
\pgfpathlineto{\pgfqpoint{2.784610in}{1.330401in}}%
\pgfpathlineto{\pgfqpoint{2.787108in}{1.331121in}}%
\pgfpathlineto{\pgfqpoint{2.789606in}{1.330529in}}%
\pgfpathlineto{\pgfqpoint{2.792104in}{1.328710in}}%
\pgfpathlineto{\pgfqpoint{2.795435in}{1.324775in}}%
\pgfpathlineto{\pgfqpoint{2.807092in}{1.309445in}}%
\pgfpathlineto{\pgfqpoint{2.810006in}{1.307999in}}%
\pgfpathlineto{\pgfqpoint{2.812921in}{1.307936in}}%
\pgfpathlineto{\pgfqpoint{2.816251in}{1.309131in}}%
\pgfpathlineto{\pgfqpoint{2.822912in}{1.311814in}}%
\pgfpathlineto{\pgfqpoint{2.825827in}{1.311587in}}%
\pgfpathlineto{\pgfqpoint{2.828741in}{1.310199in}}%
\pgfpathlineto{\pgfqpoint{2.833737in}{1.306253in}}%
\pgfpathlineto{\pgfqpoint{2.839149in}{1.302488in}}%
\pgfpathlineto{\pgfqpoint{2.851222in}{1.295637in}}%
\pgfpathlineto{\pgfqpoint{2.858716in}{1.289997in}}%
\pgfpathlineto{\pgfqpoint{2.862047in}{1.289002in}}%
\pgfpathlineto{\pgfqpoint{2.865794in}{1.289106in}}%
\pgfpathlineto{\pgfqpoint{2.872871in}{1.289583in}}%
\pgfpathlineto{\pgfqpoint{2.877867in}{1.288492in}}%
\pgfpathlineto{\pgfqpoint{2.886194in}{1.286669in}}%
\pgfpathlineto{\pgfqpoint{2.892855in}{1.285061in}}%
\pgfpathlineto{\pgfqpoint{2.899516in}{1.283109in}}%
\pgfpathlineto{\pgfqpoint{2.902014in}{1.283471in}}%
\pgfpathlineto{\pgfqpoint{2.904512in}{1.284912in}}%
\pgfpathlineto{\pgfqpoint{2.907426in}{1.287939in}}%
\pgfpathlineto{\pgfqpoint{2.919083in}{1.301985in}}%
\pgfpathlineto{\pgfqpoint{2.922830in}{1.303709in}}%
\pgfpathlineto{\pgfqpoint{2.929908in}{1.306458in}}%
\pgfpathlineto{\pgfqpoint{2.934488in}{1.309939in}}%
\pgfpathlineto{\pgfqpoint{2.939900in}{1.313867in}}%
\pgfpathlineto{\pgfqpoint{2.942814in}{1.314657in}}%
\pgfpathlineto{\pgfqpoint{2.946145in}{1.314279in}}%
\pgfpathlineto{\pgfqpoint{2.951140in}{1.313500in}}%
\pgfpathlineto{\pgfqpoint{2.953638in}{1.314361in}}%
\pgfpathlineto{\pgfqpoint{2.956136in}{1.316577in}}%
\pgfpathlineto{\pgfqpoint{2.959051in}{1.320777in}}%
\pgfpathlineto{\pgfqpoint{2.968626in}{1.336211in}}%
\pgfpathlineto{\pgfqpoint{2.971124in}{1.337968in}}%
\pgfpathlineto{\pgfqpoint{2.973206in}{1.338327in}}%
\pgfpathlineto{\pgfqpoint{2.975287in}{1.337715in}}%
\pgfpathlineto{\pgfqpoint{2.977785in}{1.335843in}}%
\pgfpathlineto{\pgfqpoint{2.981532in}{1.331332in}}%
\pgfpathlineto{\pgfqpoint{2.989026in}{1.321676in}}%
\pgfpathlineto{\pgfqpoint{2.991940in}{1.319744in}}%
\pgfpathlineto{\pgfqpoint{2.994438in}{1.319357in}}%
\pgfpathlineto{\pgfqpoint{2.996936in}{1.320096in}}%
\pgfpathlineto{\pgfqpoint{3.000267in}{1.322485in}}%
\pgfpathlineto{\pgfqpoint{3.006512in}{1.328933in}}%
\pgfpathlineto{\pgfqpoint{3.010259in}{1.331944in}}%
\pgfpathlineto{\pgfqpoint{3.012340in}{1.332461in}}%
\pgfpathlineto{\pgfqpoint{3.014006in}{1.331918in}}%
\pgfpathlineto{\pgfqpoint{3.015671in}{1.330281in}}%
\pgfpathlineto{\pgfqpoint{3.017752in}{1.326306in}}%
\pgfpathlineto{\pgfqpoint{3.019834in}{1.319808in}}%
\pgfpathlineto{\pgfqpoint{3.022332in}{1.308238in}}%
\pgfpathlineto{\pgfqpoint{3.025246in}{1.289342in}}%
\pgfpathlineto{\pgfqpoint{3.029410in}{1.253961in}}%
\pgfpathlineto{\pgfqpoint{3.036253in}{1.195318in}}%
\pgfpathlineto{\pgfqpoint{3.038110in}{1.183641in}}%
\pgfpathlineto{\pgfqpoint{3.038110in}{1.183641in}}%
\pgfusepath{stroke}%
\end{pgfscope}%
\begin{pgfscope}%
\pgfsetrectcap%
\pgfsetmiterjoin%
\pgfsetlinewidth{0.803000pt}%
\definecolor{currentstroke}{rgb}{0.000000,0.000000,0.000000}%
\pgfsetstrokecolor{currentstroke}%
\pgfsetdash{}{0pt}%
\pgfpathmoveto{\pgfqpoint{0.650000in}{0.420000in}}%
\pgfpathlineto{\pgfqpoint{0.650000in}{2.414488in}}%
\pgfusepath{stroke}%
\end{pgfscope}%
\begin{pgfscope}%
\pgfsetrectcap%
\pgfsetmiterjoin%
\pgfsetlinewidth{0.803000pt}%
\definecolor{currentstroke}{rgb}{0.000000,0.000000,0.000000}%
\pgfsetstrokecolor{currentstroke}%
\pgfsetdash{}{0pt}%
\pgfpathmoveto{\pgfqpoint{3.038110in}{0.420000in}}%
\pgfpathlineto{\pgfqpoint{3.038110in}{2.414488in}}%
\pgfusepath{stroke}%
\end{pgfscope}%
\begin{pgfscope}%
\pgfsetrectcap%
\pgfsetmiterjoin%
\pgfsetlinewidth{0.803000pt}%
\definecolor{currentstroke}{rgb}{0.000000,0.000000,0.000000}%
\pgfsetstrokecolor{currentstroke}%
\pgfsetdash{}{0pt}%
\pgfpathmoveto{\pgfqpoint{0.650000in}{0.420000in}}%
\pgfpathlineto{\pgfqpoint{3.038110in}{0.420000in}}%
\pgfusepath{stroke}%
\end{pgfscope}%
\begin{pgfscope}%
\pgfsetrectcap%
\pgfsetmiterjoin%
\pgfsetlinewidth{0.803000pt}%
\definecolor{currentstroke}{rgb}{0.000000,0.000000,0.000000}%
\pgfsetstrokecolor{currentstroke}%
\pgfsetdash{}{0pt}%
\pgfpathmoveto{\pgfqpoint{0.650000in}{2.414488in}}%
\pgfpathlineto{\pgfqpoint{3.038110in}{2.414488in}}%
\pgfusepath{stroke}%
\end{pgfscope}%
\end{pgfpicture}%
\makeatother%
\endgroup%

    \label{fig:R500deg18R100deg90_beta}
\end{subfigure}
\\[-20pt]
\begin{subfigure}[t]{0.495\textwidth}
    \centering
    \subcaption{}
    %% Creator: Matplotlib, PGF backend
%%
%% To include the figure in your LaTeX document, write
%%   \input{<filename>.pgf}
%%
%% Make sure the required packages are loaded in your preamble
%%   \usepackage{pgf}
%%
%% Also ensure that all the required font packages are loaded; for instance,
%% the lmodern package is sometimes necessary when using math font.
%%   \usepackage{lmodern}
%%
%% Figures using additional raster images can only be included by \input if
%% they are in the same directory as the main LaTeX file. For loading figures
%% from other directories you can use the `import` package
%%   \usepackage{import}
%%
%% and then include the figures with
%%   \import{<path to file>}{<filename>.pgf}
%%
%% Matplotlib used the following preamble
%%   \def\mathdefault#1{#1}
%%   \everymath=\expandafter{\the\everymath\displaystyle}
%%   \IfFileExists{scrextend.sty}{
%%     \usepackage[fontsize=11.000000pt]{scrextend}
%%   }{
%%     \renewcommand{\normalsize}{\fontsize{11.000000}{13.200000}\selectfont}
%%     \normalsize
%%   }
%%   
%%   \makeatletter\@ifpackageloaded{underscore}{}{\usepackage[strings]{underscore}}\makeatother
%%
\begingroup%
\makeatletter%
\begin{pgfpicture}%
\pgfpathrectangle{\pgfpointorigin}{\pgfqpoint{3.118110in}{2.494488in}}%
\pgfusepath{use as bounding box, clip}%
\begin{pgfscope}%
\pgfsetbuttcap%
\pgfsetmiterjoin%
\definecolor{currentfill}{rgb}{1.000000,1.000000,1.000000}%
\pgfsetfillcolor{currentfill}%
\pgfsetlinewidth{0.000000pt}%
\definecolor{currentstroke}{rgb}{1.000000,1.000000,1.000000}%
\pgfsetstrokecolor{currentstroke}%
\pgfsetdash{}{0pt}%
\pgfpathmoveto{\pgfqpoint{0.000000in}{0.000000in}}%
\pgfpathlineto{\pgfqpoint{3.118110in}{0.000000in}}%
\pgfpathlineto{\pgfqpoint{3.118110in}{2.494488in}}%
\pgfpathlineto{\pgfqpoint{0.000000in}{2.494488in}}%
\pgfpathlineto{\pgfqpoint{0.000000in}{0.000000in}}%
\pgfpathclose%
\pgfusepath{fill}%
\end{pgfscope}%
\begin{pgfscope}%
\pgfsetbuttcap%
\pgfsetmiterjoin%
\definecolor{currentfill}{rgb}{1.000000,1.000000,1.000000}%
\pgfsetfillcolor{currentfill}%
\pgfsetlinewidth{0.000000pt}%
\definecolor{currentstroke}{rgb}{0.000000,0.000000,0.000000}%
\pgfsetstrokecolor{currentstroke}%
\pgfsetstrokeopacity{0.000000}%
\pgfsetdash{}{0pt}%
\pgfpathmoveto{\pgfqpoint{0.650000in}{0.420000in}}%
\pgfpathlineto{\pgfqpoint{3.038110in}{0.420000in}}%
\pgfpathlineto{\pgfqpoint{3.038110in}{2.414488in}}%
\pgfpathlineto{\pgfqpoint{0.650000in}{2.414488in}}%
\pgfpathlineto{\pgfqpoint{0.650000in}{0.420000in}}%
\pgfpathclose%
\pgfusepath{fill}%
\end{pgfscope}%
\begin{pgfscope}%
\pgfsetbuttcap%
\pgfsetroundjoin%
\definecolor{currentfill}{rgb}{0.000000,0.000000,0.000000}%
\pgfsetfillcolor{currentfill}%
\pgfsetlinewidth{0.501875pt}%
\definecolor{currentstroke}{rgb}{0.000000,0.000000,0.000000}%
\pgfsetstrokecolor{currentstroke}%
\pgfsetdash{}{0pt}%
\pgfsys@defobject{currentmarker}{\pgfqpoint{0.000000in}{0.000000in}}{\pgfqpoint{0.000000in}{0.041667in}}{%
\pgfpathmoveto{\pgfqpoint{0.000000in}{0.000000in}}%
\pgfpathlineto{\pgfqpoint{0.000000in}{0.041667in}}%
\pgfusepath{stroke,fill}%
}%
\begin{pgfscope}%
\pgfsys@transformshift{0.650000in}{0.420000in}%
\pgfsys@useobject{currentmarker}{}%
\end{pgfscope}%
\end{pgfscope}%
\begin{pgfscope}%
\pgfsetbuttcap%
\pgfsetroundjoin%
\definecolor{currentfill}{rgb}{0.000000,0.000000,0.000000}%
\pgfsetfillcolor{currentfill}%
\pgfsetlinewidth{0.501875pt}%
\definecolor{currentstroke}{rgb}{0.000000,0.000000,0.000000}%
\pgfsetstrokecolor{currentstroke}%
\pgfsetdash{}{0pt}%
\pgfsys@defobject{currentmarker}{\pgfqpoint{0.000000in}{-0.041667in}}{\pgfqpoint{0.000000in}{0.000000in}}{%
\pgfpathmoveto{\pgfqpoint{0.000000in}{0.000000in}}%
\pgfpathlineto{\pgfqpoint{0.000000in}{-0.041667in}}%
\pgfusepath{stroke,fill}%
}%
\begin{pgfscope}%
\pgfsys@transformshift{0.650000in}{2.414488in}%
\pgfsys@useobject{currentmarker}{}%
\end{pgfscope}%
\end{pgfscope}%
\begin{pgfscope}%
\definecolor{textcolor}{rgb}{0.000000,0.000000,0.000000}%
\pgfsetstrokecolor{textcolor}%
\pgfsetfillcolor{textcolor}%
\pgftext[x=0.650000in,y=0.371389in,,top]{\color{textcolor}{\rmfamily\fontsize{11.000000}{13.200000}\selectfont\catcode`\^=\active\def^{\ifmmode\sp\else\^{}\fi}\catcode`\%=\active\def%{\%}0}}%
\end{pgfscope}%
\begin{pgfscope}%
\pgfsetbuttcap%
\pgfsetroundjoin%
\definecolor{currentfill}{rgb}{0.000000,0.000000,0.000000}%
\pgfsetfillcolor{currentfill}%
\pgfsetlinewidth{0.501875pt}%
\definecolor{currentstroke}{rgb}{0.000000,0.000000,0.000000}%
\pgfsetstrokecolor{currentstroke}%
\pgfsetdash{}{0pt}%
\pgfsys@defobject{currentmarker}{\pgfqpoint{0.000000in}{0.000000in}}{\pgfqpoint{0.000000in}{0.041667in}}{%
\pgfpathmoveto{\pgfqpoint{0.000000in}{0.000000in}}%
\pgfpathlineto{\pgfqpoint{0.000000in}{0.041667in}}%
\pgfusepath{stroke,fill}%
}%
\begin{pgfscope}%
\pgfsys@transformshift{1.493052in}{0.420000in}%
\pgfsys@useobject{currentmarker}{}%
\end{pgfscope}%
\end{pgfscope}%
\begin{pgfscope}%
\pgfsetbuttcap%
\pgfsetroundjoin%
\definecolor{currentfill}{rgb}{0.000000,0.000000,0.000000}%
\pgfsetfillcolor{currentfill}%
\pgfsetlinewidth{0.501875pt}%
\definecolor{currentstroke}{rgb}{0.000000,0.000000,0.000000}%
\pgfsetstrokecolor{currentstroke}%
\pgfsetdash{}{0pt}%
\pgfsys@defobject{currentmarker}{\pgfqpoint{0.000000in}{-0.041667in}}{\pgfqpoint{0.000000in}{0.000000in}}{%
\pgfpathmoveto{\pgfqpoint{0.000000in}{0.000000in}}%
\pgfpathlineto{\pgfqpoint{0.000000in}{-0.041667in}}%
\pgfusepath{stroke,fill}%
}%
\begin{pgfscope}%
\pgfsys@transformshift{1.493052in}{2.414488in}%
\pgfsys@useobject{currentmarker}{}%
\end{pgfscope}%
\end{pgfscope}%
\begin{pgfscope}%
\definecolor{textcolor}{rgb}{0.000000,0.000000,0.000000}%
\pgfsetstrokecolor{textcolor}%
\pgfsetfillcolor{textcolor}%
\pgftext[x=1.493052in,y=0.371389in,,top]{\color{textcolor}{\rmfamily\fontsize{11.000000}{13.200000}\selectfont\catcode`\^=\active\def^{\ifmmode\sp\else\^{}\fi}\catcode`\%=\active\def%{\%}100}}%
\end{pgfscope}%
\begin{pgfscope}%
\pgfsetbuttcap%
\pgfsetroundjoin%
\definecolor{currentfill}{rgb}{0.000000,0.000000,0.000000}%
\pgfsetfillcolor{currentfill}%
\pgfsetlinewidth{0.501875pt}%
\definecolor{currentstroke}{rgb}{0.000000,0.000000,0.000000}%
\pgfsetstrokecolor{currentstroke}%
\pgfsetdash{}{0pt}%
\pgfsys@defobject{currentmarker}{\pgfqpoint{0.000000in}{0.000000in}}{\pgfqpoint{0.000000in}{0.041667in}}{%
\pgfpathmoveto{\pgfqpoint{0.000000in}{0.000000in}}%
\pgfpathlineto{\pgfqpoint{0.000000in}{0.041667in}}%
\pgfusepath{stroke,fill}%
}%
\begin{pgfscope}%
\pgfsys@transformshift{2.336103in}{0.420000in}%
\pgfsys@useobject{currentmarker}{}%
\end{pgfscope}%
\end{pgfscope}%
\begin{pgfscope}%
\pgfsetbuttcap%
\pgfsetroundjoin%
\definecolor{currentfill}{rgb}{0.000000,0.000000,0.000000}%
\pgfsetfillcolor{currentfill}%
\pgfsetlinewidth{0.501875pt}%
\definecolor{currentstroke}{rgb}{0.000000,0.000000,0.000000}%
\pgfsetstrokecolor{currentstroke}%
\pgfsetdash{}{0pt}%
\pgfsys@defobject{currentmarker}{\pgfqpoint{0.000000in}{-0.041667in}}{\pgfqpoint{0.000000in}{0.000000in}}{%
\pgfpathmoveto{\pgfqpoint{0.000000in}{0.000000in}}%
\pgfpathlineto{\pgfqpoint{0.000000in}{-0.041667in}}%
\pgfusepath{stroke,fill}%
}%
\begin{pgfscope}%
\pgfsys@transformshift{2.336103in}{2.414488in}%
\pgfsys@useobject{currentmarker}{}%
\end{pgfscope}%
\end{pgfscope}%
\begin{pgfscope}%
\definecolor{textcolor}{rgb}{0.000000,0.000000,0.000000}%
\pgfsetstrokecolor{textcolor}%
\pgfsetfillcolor{textcolor}%
\pgftext[x=2.336103in,y=0.371389in,,top]{\color{textcolor}{\rmfamily\fontsize{11.000000}{13.200000}\selectfont\catcode`\^=\active\def^{\ifmmode\sp\else\^{}\fi}\catcode`\%=\active\def%{\%}200}}%
\end{pgfscope}%
\begin{pgfscope}%
\pgfsetbuttcap%
\pgfsetroundjoin%
\definecolor{currentfill}{rgb}{0.000000,0.000000,0.000000}%
\pgfsetfillcolor{currentfill}%
\pgfsetlinewidth{0.401500pt}%
\definecolor{currentstroke}{rgb}{0.000000,0.000000,0.000000}%
\pgfsetstrokecolor{currentstroke}%
\pgfsetdash{}{0pt}%
\pgfsys@defobject{currentmarker}{\pgfqpoint{0.000000in}{0.000000in}}{\pgfqpoint{0.000000in}{0.020833in}}{%
\pgfpathmoveto{\pgfqpoint{0.000000in}{0.000000in}}%
\pgfpathlineto{\pgfqpoint{0.000000in}{0.020833in}}%
\pgfusepath{stroke,fill}%
}%
\begin{pgfscope}%
\pgfsys@transformshift{0.818610in}{0.420000in}%
\pgfsys@useobject{currentmarker}{}%
\end{pgfscope}%
\end{pgfscope}%
\begin{pgfscope}%
\pgfsetbuttcap%
\pgfsetroundjoin%
\definecolor{currentfill}{rgb}{0.000000,0.000000,0.000000}%
\pgfsetfillcolor{currentfill}%
\pgfsetlinewidth{0.401500pt}%
\definecolor{currentstroke}{rgb}{0.000000,0.000000,0.000000}%
\pgfsetstrokecolor{currentstroke}%
\pgfsetdash{}{0pt}%
\pgfsys@defobject{currentmarker}{\pgfqpoint{0.000000in}{-0.020833in}}{\pgfqpoint{0.000000in}{0.000000in}}{%
\pgfpathmoveto{\pgfqpoint{0.000000in}{0.000000in}}%
\pgfpathlineto{\pgfqpoint{0.000000in}{-0.020833in}}%
\pgfusepath{stroke,fill}%
}%
\begin{pgfscope}%
\pgfsys@transformshift{0.818610in}{2.414488in}%
\pgfsys@useobject{currentmarker}{}%
\end{pgfscope}%
\end{pgfscope}%
\begin{pgfscope}%
\pgfsetbuttcap%
\pgfsetroundjoin%
\definecolor{currentfill}{rgb}{0.000000,0.000000,0.000000}%
\pgfsetfillcolor{currentfill}%
\pgfsetlinewidth{0.401500pt}%
\definecolor{currentstroke}{rgb}{0.000000,0.000000,0.000000}%
\pgfsetstrokecolor{currentstroke}%
\pgfsetdash{}{0pt}%
\pgfsys@defobject{currentmarker}{\pgfqpoint{0.000000in}{0.000000in}}{\pgfqpoint{0.000000in}{0.020833in}}{%
\pgfpathmoveto{\pgfqpoint{0.000000in}{0.000000in}}%
\pgfpathlineto{\pgfqpoint{0.000000in}{0.020833in}}%
\pgfusepath{stroke,fill}%
}%
\begin{pgfscope}%
\pgfsys@transformshift{0.987221in}{0.420000in}%
\pgfsys@useobject{currentmarker}{}%
\end{pgfscope}%
\end{pgfscope}%
\begin{pgfscope}%
\pgfsetbuttcap%
\pgfsetroundjoin%
\definecolor{currentfill}{rgb}{0.000000,0.000000,0.000000}%
\pgfsetfillcolor{currentfill}%
\pgfsetlinewidth{0.401500pt}%
\definecolor{currentstroke}{rgb}{0.000000,0.000000,0.000000}%
\pgfsetstrokecolor{currentstroke}%
\pgfsetdash{}{0pt}%
\pgfsys@defobject{currentmarker}{\pgfqpoint{0.000000in}{-0.020833in}}{\pgfqpoint{0.000000in}{0.000000in}}{%
\pgfpathmoveto{\pgfqpoint{0.000000in}{0.000000in}}%
\pgfpathlineto{\pgfqpoint{0.000000in}{-0.020833in}}%
\pgfusepath{stroke,fill}%
}%
\begin{pgfscope}%
\pgfsys@transformshift{0.987221in}{2.414488in}%
\pgfsys@useobject{currentmarker}{}%
\end{pgfscope}%
\end{pgfscope}%
\begin{pgfscope}%
\pgfsetbuttcap%
\pgfsetroundjoin%
\definecolor{currentfill}{rgb}{0.000000,0.000000,0.000000}%
\pgfsetfillcolor{currentfill}%
\pgfsetlinewidth{0.401500pt}%
\definecolor{currentstroke}{rgb}{0.000000,0.000000,0.000000}%
\pgfsetstrokecolor{currentstroke}%
\pgfsetdash{}{0pt}%
\pgfsys@defobject{currentmarker}{\pgfqpoint{0.000000in}{0.000000in}}{\pgfqpoint{0.000000in}{0.020833in}}{%
\pgfpathmoveto{\pgfqpoint{0.000000in}{0.000000in}}%
\pgfpathlineto{\pgfqpoint{0.000000in}{0.020833in}}%
\pgfusepath{stroke,fill}%
}%
\begin{pgfscope}%
\pgfsys@transformshift{1.155831in}{0.420000in}%
\pgfsys@useobject{currentmarker}{}%
\end{pgfscope}%
\end{pgfscope}%
\begin{pgfscope}%
\pgfsetbuttcap%
\pgfsetroundjoin%
\definecolor{currentfill}{rgb}{0.000000,0.000000,0.000000}%
\pgfsetfillcolor{currentfill}%
\pgfsetlinewidth{0.401500pt}%
\definecolor{currentstroke}{rgb}{0.000000,0.000000,0.000000}%
\pgfsetstrokecolor{currentstroke}%
\pgfsetdash{}{0pt}%
\pgfsys@defobject{currentmarker}{\pgfqpoint{0.000000in}{-0.020833in}}{\pgfqpoint{0.000000in}{0.000000in}}{%
\pgfpathmoveto{\pgfqpoint{0.000000in}{0.000000in}}%
\pgfpathlineto{\pgfqpoint{0.000000in}{-0.020833in}}%
\pgfusepath{stroke,fill}%
}%
\begin{pgfscope}%
\pgfsys@transformshift{1.155831in}{2.414488in}%
\pgfsys@useobject{currentmarker}{}%
\end{pgfscope}%
\end{pgfscope}%
\begin{pgfscope}%
\pgfsetbuttcap%
\pgfsetroundjoin%
\definecolor{currentfill}{rgb}{0.000000,0.000000,0.000000}%
\pgfsetfillcolor{currentfill}%
\pgfsetlinewidth{0.401500pt}%
\definecolor{currentstroke}{rgb}{0.000000,0.000000,0.000000}%
\pgfsetstrokecolor{currentstroke}%
\pgfsetdash{}{0pt}%
\pgfsys@defobject{currentmarker}{\pgfqpoint{0.000000in}{0.000000in}}{\pgfqpoint{0.000000in}{0.020833in}}{%
\pgfpathmoveto{\pgfqpoint{0.000000in}{0.000000in}}%
\pgfpathlineto{\pgfqpoint{0.000000in}{0.020833in}}%
\pgfusepath{stroke,fill}%
}%
\begin{pgfscope}%
\pgfsys@transformshift{1.324441in}{0.420000in}%
\pgfsys@useobject{currentmarker}{}%
\end{pgfscope}%
\end{pgfscope}%
\begin{pgfscope}%
\pgfsetbuttcap%
\pgfsetroundjoin%
\definecolor{currentfill}{rgb}{0.000000,0.000000,0.000000}%
\pgfsetfillcolor{currentfill}%
\pgfsetlinewidth{0.401500pt}%
\definecolor{currentstroke}{rgb}{0.000000,0.000000,0.000000}%
\pgfsetstrokecolor{currentstroke}%
\pgfsetdash{}{0pt}%
\pgfsys@defobject{currentmarker}{\pgfqpoint{0.000000in}{-0.020833in}}{\pgfqpoint{0.000000in}{0.000000in}}{%
\pgfpathmoveto{\pgfqpoint{0.000000in}{0.000000in}}%
\pgfpathlineto{\pgfqpoint{0.000000in}{-0.020833in}}%
\pgfusepath{stroke,fill}%
}%
\begin{pgfscope}%
\pgfsys@transformshift{1.324441in}{2.414488in}%
\pgfsys@useobject{currentmarker}{}%
\end{pgfscope}%
\end{pgfscope}%
\begin{pgfscope}%
\pgfsetbuttcap%
\pgfsetroundjoin%
\definecolor{currentfill}{rgb}{0.000000,0.000000,0.000000}%
\pgfsetfillcolor{currentfill}%
\pgfsetlinewidth{0.401500pt}%
\definecolor{currentstroke}{rgb}{0.000000,0.000000,0.000000}%
\pgfsetstrokecolor{currentstroke}%
\pgfsetdash{}{0pt}%
\pgfsys@defobject{currentmarker}{\pgfqpoint{0.000000in}{0.000000in}}{\pgfqpoint{0.000000in}{0.020833in}}{%
\pgfpathmoveto{\pgfqpoint{0.000000in}{0.000000in}}%
\pgfpathlineto{\pgfqpoint{0.000000in}{0.020833in}}%
\pgfusepath{stroke,fill}%
}%
\begin{pgfscope}%
\pgfsys@transformshift{1.661662in}{0.420000in}%
\pgfsys@useobject{currentmarker}{}%
\end{pgfscope}%
\end{pgfscope}%
\begin{pgfscope}%
\pgfsetbuttcap%
\pgfsetroundjoin%
\definecolor{currentfill}{rgb}{0.000000,0.000000,0.000000}%
\pgfsetfillcolor{currentfill}%
\pgfsetlinewidth{0.401500pt}%
\definecolor{currentstroke}{rgb}{0.000000,0.000000,0.000000}%
\pgfsetstrokecolor{currentstroke}%
\pgfsetdash{}{0pt}%
\pgfsys@defobject{currentmarker}{\pgfqpoint{0.000000in}{-0.020833in}}{\pgfqpoint{0.000000in}{0.000000in}}{%
\pgfpathmoveto{\pgfqpoint{0.000000in}{0.000000in}}%
\pgfpathlineto{\pgfqpoint{0.000000in}{-0.020833in}}%
\pgfusepath{stroke,fill}%
}%
\begin{pgfscope}%
\pgfsys@transformshift{1.661662in}{2.414488in}%
\pgfsys@useobject{currentmarker}{}%
\end{pgfscope}%
\end{pgfscope}%
\begin{pgfscope}%
\pgfsetbuttcap%
\pgfsetroundjoin%
\definecolor{currentfill}{rgb}{0.000000,0.000000,0.000000}%
\pgfsetfillcolor{currentfill}%
\pgfsetlinewidth{0.401500pt}%
\definecolor{currentstroke}{rgb}{0.000000,0.000000,0.000000}%
\pgfsetstrokecolor{currentstroke}%
\pgfsetdash{}{0pt}%
\pgfsys@defobject{currentmarker}{\pgfqpoint{0.000000in}{0.000000in}}{\pgfqpoint{0.000000in}{0.020833in}}{%
\pgfpathmoveto{\pgfqpoint{0.000000in}{0.000000in}}%
\pgfpathlineto{\pgfqpoint{0.000000in}{0.020833in}}%
\pgfusepath{stroke,fill}%
}%
\begin{pgfscope}%
\pgfsys@transformshift{1.830272in}{0.420000in}%
\pgfsys@useobject{currentmarker}{}%
\end{pgfscope}%
\end{pgfscope}%
\begin{pgfscope}%
\pgfsetbuttcap%
\pgfsetroundjoin%
\definecolor{currentfill}{rgb}{0.000000,0.000000,0.000000}%
\pgfsetfillcolor{currentfill}%
\pgfsetlinewidth{0.401500pt}%
\definecolor{currentstroke}{rgb}{0.000000,0.000000,0.000000}%
\pgfsetstrokecolor{currentstroke}%
\pgfsetdash{}{0pt}%
\pgfsys@defobject{currentmarker}{\pgfqpoint{0.000000in}{-0.020833in}}{\pgfqpoint{0.000000in}{0.000000in}}{%
\pgfpathmoveto{\pgfqpoint{0.000000in}{0.000000in}}%
\pgfpathlineto{\pgfqpoint{0.000000in}{-0.020833in}}%
\pgfusepath{stroke,fill}%
}%
\begin{pgfscope}%
\pgfsys@transformshift{1.830272in}{2.414488in}%
\pgfsys@useobject{currentmarker}{}%
\end{pgfscope}%
\end{pgfscope}%
\begin{pgfscope}%
\pgfsetbuttcap%
\pgfsetroundjoin%
\definecolor{currentfill}{rgb}{0.000000,0.000000,0.000000}%
\pgfsetfillcolor{currentfill}%
\pgfsetlinewidth{0.401500pt}%
\definecolor{currentstroke}{rgb}{0.000000,0.000000,0.000000}%
\pgfsetstrokecolor{currentstroke}%
\pgfsetdash{}{0pt}%
\pgfsys@defobject{currentmarker}{\pgfqpoint{0.000000in}{0.000000in}}{\pgfqpoint{0.000000in}{0.020833in}}{%
\pgfpathmoveto{\pgfqpoint{0.000000in}{0.000000in}}%
\pgfpathlineto{\pgfqpoint{0.000000in}{0.020833in}}%
\pgfusepath{stroke,fill}%
}%
\begin{pgfscope}%
\pgfsys@transformshift{1.998883in}{0.420000in}%
\pgfsys@useobject{currentmarker}{}%
\end{pgfscope}%
\end{pgfscope}%
\begin{pgfscope}%
\pgfsetbuttcap%
\pgfsetroundjoin%
\definecolor{currentfill}{rgb}{0.000000,0.000000,0.000000}%
\pgfsetfillcolor{currentfill}%
\pgfsetlinewidth{0.401500pt}%
\definecolor{currentstroke}{rgb}{0.000000,0.000000,0.000000}%
\pgfsetstrokecolor{currentstroke}%
\pgfsetdash{}{0pt}%
\pgfsys@defobject{currentmarker}{\pgfqpoint{0.000000in}{-0.020833in}}{\pgfqpoint{0.000000in}{0.000000in}}{%
\pgfpathmoveto{\pgfqpoint{0.000000in}{0.000000in}}%
\pgfpathlineto{\pgfqpoint{0.000000in}{-0.020833in}}%
\pgfusepath{stroke,fill}%
}%
\begin{pgfscope}%
\pgfsys@transformshift{1.998883in}{2.414488in}%
\pgfsys@useobject{currentmarker}{}%
\end{pgfscope}%
\end{pgfscope}%
\begin{pgfscope}%
\pgfsetbuttcap%
\pgfsetroundjoin%
\definecolor{currentfill}{rgb}{0.000000,0.000000,0.000000}%
\pgfsetfillcolor{currentfill}%
\pgfsetlinewidth{0.401500pt}%
\definecolor{currentstroke}{rgb}{0.000000,0.000000,0.000000}%
\pgfsetstrokecolor{currentstroke}%
\pgfsetdash{}{0pt}%
\pgfsys@defobject{currentmarker}{\pgfqpoint{0.000000in}{0.000000in}}{\pgfqpoint{0.000000in}{0.020833in}}{%
\pgfpathmoveto{\pgfqpoint{0.000000in}{0.000000in}}%
\pgfpathlineto{\pgfqpoint{0.000000in}{0.020833in}}%
\pgfusepath{stroke,fill}%
}%
\begin{pgfscope}%
\pgfsys@transformshift{2.167493in}{0.420000in}%
\pgfsys@useobject{currentmarker}{}%
\end{pgfscope}%
\end{pgfscope}%
\begin{pgfscope}%
\pgfsetbuttcap%
\pgfsetroundjoin%
\definecolor{currentfill}{rgb}{0.000000,0.000000,0.000000}%
\pgfsetfillcolor{currentfill}%
\pgfsetlinewidth{0.401500pt}%
\definecolor{currentstroke}{rgb}{0.000000,0.000000,0.000000}%
\pgfsetstrokecolor{currentstroke}%
\pgfsetdash{}{0pt}%
\pgfsys@defobject{currentmarker}{\pgfqpoint{0.000000in}{-0.020833in}}{\pgfqpoint{0.000000in}{0.000000in}}{%
\pgfpathmoveto{\pgfqpoint{0.000000in}{0.000000in}}%
\pgfpathlineto{\pgfqpoint{0.000000in}{-0.020833in}}%
\pgfusepath{stroke,fill}%
}%
\begin{pgfscope}%
\pgfsys@transformshift{2.167493in}{2.414488in}%
\pgfsys@useobject{currentmarker}{}%
\end{pgfscope}%
\end{pgfscope}%
\begin{pgfscope}%
\pgfsetbuttcap%
\pgfsetroundjoin%
\definecolor{currentfill}{rgb}{0.000000,0.000000,0.000000}%
\pgfsetfillcolor{currentfill}%
\pgfsetlinewidth{0.401500pt}%
\definecolor{currentstroke}{rgb}{0.000000,0.000000,0.000000}%
\pgfsetstrokecolor{currentstroke}%
\pgfsetdash{}{0pt}%
\pgfsys@defobject{currentmarker}{\pgfqpoint{0.000000in}{0.000000in}}{\pgfqpoint{0.000000in}{0.020833in}}{%
\pgfpathmoveto{\pgfqpoint{0.000000in}{0.000000in}}%
\pgfpathlineto{\pgfqpoint{0.000000in}{0.020833in}}%
\pgfusepath{stroke,fill}%
}%
\begin{pgfscope}%
\pgfsys@transformshift{2.504714in}{0.420000in}%
\pgfsys@useobject{currentmarker}{}%
\end{pgfscope}%
\end{pgfscope}%
\begin{pgfscope}%
\pgfsetbuttcap%
\pgfsetroundjoin%
\definecolor{currentfill}{rgb}{0.000000,0.000000,0.000000}%
\pgfsetfillcolor{currentfill}%
\pgfsetlinewidth{0.401500pt}%
\definecolor{currentstroke}{rgb}{0.000000,0.000000,0.000000}%
\pgfsetstrokecolor{currentstroke}%
\pgfsetdash{}{0pt}%
\pgfsys@defobject{currentmarker}{\pgfqpoint{0.000000in}{-0.020833in}}{\pgfqpoint{0.000000in}{0.000000in}}{%
\pgfpathmoveto{\pgfqpoint{0.000000in}{0.000000in}}%
\pgfpathlineto{\pgfqpoint{0.000000in}{-0.020833in}}%
\pgfusepath{stroke,fill}%
}%
\begin{pgfscope}%
\pgfsys@transformshift{2.504714in}{2.414488in}%
\pgfsys@useobject{currentmarker}{}%
\end{pgfscope}%
\end{pgfscope}%
\begin{pgfscope}%
\pgfsetbuttcap%
\pgfsetroundjoin%
\definecolor{currentfill}{rgb}{0.000000,0.000000,0.000000}%
\pgfsetfillcolor{currentfill}%
\pgfsetlinewidth{0.401500pt}%
\definecolor{currentstroke}{rgb}{0.000000,0.000000,0.000000}%
\pgfsetstrokecolor{currentstroke}%
\pgfsetdash{}{0pt}%
\pgfsys@defobject{currentmarker}{\pgfqpoint{0.000000in}{0.000000in}}{\pgfqpoint{0.000000in}{0.020833in}}{%
\pgfpathmoveto{\pgfqpoint{0.000000in}{0.000000in}}%
\pgfpathlineto{\pgfqpoint{0.000000in}{0.020833in}}%
\pgfusepath{stroke,fill}%
}%
\begin{pgfscope}%
\pgfsys@transformshift{2.673324in}{0.420000in}%
\pgfsys@useobject{currentmarker}{}%
\end{pgfscope}%
\end{pgfscope}%
\begin{pgfscope}%
\pgfsetbuttcap%
\pgfsetroundjoin%
\definecolor{currentfill}{rgb}{0.000000,0.000000,0.000000}%
\pgfsetfillcolor{currentfill}%
\pgfsetlinewidth{0.401500pt}%
\definecolor{currentstroke}{rgb}{0.000000,0.000000,0.000000}%
\pgfsetstrokecolor{currentstroke}%
\pgfsetdash{}{0pt}%
\pgfsys@defobject{currentmarker}{\pgfqpoint{0.000000in}{-0.020833in}}{\pgfqpoint{0.000000in}{0.000000in}}{%
\pgfpathmoveto{\pgfqpoint{0.000000in}{0.000000in}}%
\pgfpathlineto{\pgfqpoint{0.000000in}{-0.020833in}}%
\pgfusepath{stroke,fill}%
}%
\begin{pgfscope}%
\pgfsys@transformshift{2.673324in}{2.414488in}%
\pgfsys@useobject{currentmarker}{}%
\end{pgfscope}%
\end{pgfscope}%
\begin{pgfscope}%
\pgfsetbuttcap%
\pgfsetroundjoin%
\definecolor{currentfill}{rgb}{0.000000,0.000000,0.000000}%
\pgfsetfillcolor{currentfill}%
\pgfsetlinewidth{0.401500pt}%
\definecolor{currentstroke}{rgb}{0.000000,0.000000,0.000000}%
\pgfsetstrokecolor{currentstroke}%
\pgfsetdash{}{0pt}%
\pgfsys@defobject{currentmarker}{\pgfqpoint{0.000000in}{0.000000in}}{\pgfqpoint{0.000000in}{0.020833in}}{%
\pgfpathmoveto{\pgfqpoint{0.000000in}{0.000000in}}%
\pgfpathlineto{\pgfqpoint{0.000000in}{0.020833in}}%
\pgfusepath{stroke,fill}%
}%
\begin{pgfscope}%
\pgfsys@transformshift{2.841934in}{0.420000in}%
\pgfsys@useobject{currentmarker}{}%
\end{pgfscope}%
\end{pgfscope}%
\begin{pgfscope}%
\pgfsetbuttcap%
\pgfsetroundjoin%
\definecolor{currentfill}{rgb}{0.000000,0.000000,0.000000}%
\pgfsetfillcolor{currentfill}%
\pgfsetlinewidth{0.401500pt}%
\definecolor{currentstroke}{rgb}{0.000000,0.000000,0.000000}%
\pgfsetstrokecolor{currentstroke}%
\pgfsetdash{}{0pt}%
\pgfsys@defobject{currentmarker}{\pgfqpoint{0.000000in}{-0.020833in}}{\pgfqpoint{0.000000in}{0.000000in}}{%
\pgfpathmoveto{\pgfqpoint{0.000000in}{0.000000in}}%
\pgfpathlineto{\pgfqpoint{0.000000in}{-0.020833in}}%
\pgfusepath{stroke,fill}%
}%
\begin{pgfscope}%
\pgfsys@transformshift{2.841934in}{2.414488in}%
\pgfsys@useobject{currentmarker}{}%
\end{pgfscope}%
\end{pgfscope}%
\begin{pgfscope}%
\pgfsetbuttcap%
\pgfsetroundjoin%
\definecolor{currentfill}{rgb}{0.000000,0.000000,0.000000}%
\pgfsetfillcolor{currentfill}%
\pgfsetlinewidth{0.401500pt}%
\definecolor{currentstroke}{rgb}{0.000000,0.000000,0.000000}%
\pgfsetstrokecolor{currentstroke}%
\pgfsetdash{}{0pt}%
\pgfsys@defobject{currentmarker}{\pgfqpoint{0.000000in}{0.000000in}}{\pgfqpoint{0.000000in}{0.020833in}}{%
\pgfpathmoveto{\pgfqpoint{0.000000in}{0.000000in}}%
\pgfpathlineto{\pgfqpoint{0.000000in}{0.020833in}}%
\pgfusepath{stroke,fill}%
}%
\begin{pgfscope}%
\pgfsys@transformshift{3.010545in}{0.420000in}%
\pgfsys@useobject{currentmarker}{}%
\end{pgfscope}%
\end{pgfscope}%
\begin{pgfscope}%
\pgfsetbuttcap%
\pgfsetroundjoin%
\definecolor{currentfill}{rgb}{0.000000,0.000000,0.000000}%
\pgfsetfillcolor{currentfill}%
\pgfsetlinewidth{0.401500pt}%
\definecolor{currentstroke}{rgb}{0.000000,0.000000,0.000000}%
\pgfsetstrokecolor{currentstroke}%
\pgfsetdash{}{0pt}%
\pgfsys@defobject{currentmarker}{\pgfqpoint{0.000000in}{-0.020833in}}{\pgfqpoint{0.000000in}{0.000000in}}{%
\pgfpathmoveto{\pgfqpoint{0.000000in}{0.000000in}}%
\pgfpathlineto{\pgfqpoint{0.000000in}{-0.020833in}}%
\pgfusepath{stroke,fill}%
}%
\begin{pgfscope}%
\pgfsys@transformshift{3.010545in}{2.414488in}%
\pgfsys@useobject{currentmarker}{}%
\end{pgfscope}%
\end{pgfscope}%
\begin{pgfscope}%
\definecolor{textcolor}{rgb}{0.000000,0.000000,0.000000}%
\pgfsetstrokecolor{textcolor}%
\pgfsetfillcolor{textcolor}%
\pgftext[x=1.844055in,y=0.208426in,,top]{\color{textcolor}{\rmfamily\fontsize{12.000000}{14.400000}\selectfont\catcode`\^=\active\def^{\ifmmode\sp\else\^{}\fi}\catcode`\%=\active\def%{\%}$s/\delta_{0}$}}%
\end{pgfscope}%
\begin{pgfscope}%
\pgfsetbuttcap%
\pgfsetroundjoin%
\definecolor{currentfill}{rgb}{0.000000,0.000000,0.000000}%
\pgfsetfillcolor{currentfill}%
\pgfsetlinewidth{0.501875pt}%
\definecolor{currentstroke}{rgb}{0.000000,0.000000,0.000000}%
\pgfsetstrokecolor{currentstroke}%
\pgfsetdash{}{0pt}%
\pgfsys@defobject{currentmarker}{\pgfqpoint{0.000000in}{0.000000in}}{\pgfqpoint{0.041667in}{0.000000in}}{%
\pgfpathmoveto{\pgfqpoint{0.000000in}{0.000000in}}%
\pgfpathlineto{\pgfqpoint{0.041667in}{0.000000in}}%
\pgfusepath{stroke,fill}%
}%
\begin{pgfscope}%
\pgfsys@transformshift{0.650000in}{0.420000in}%
\pgfsys@useobject{currentmarker}{}%
\end{pgfscope}%
\end{pgfscope}%
\begin{pgfscope}%
\pgfsetbuttcap%
\pgfsetroundjoin%
\definecolor{currentfill}{rgb}{0.000000,0.000000,0.000000}%
\pgfsetfillcolor{currentfill}%
\pgfsetlinewidth{0.501875pt}%
\definecolor{currentstroke}{rgb}{0.000000,0.000000,0.000000}%
\pgfsetstrokecolor{currentstroke}%
\pgfsetdash{}{0pt}%
\pgfsys@defobject{currentmarker}{\pgfqpoint{-0.041667in}{0.000000in}}{\pgfqpoint{-0.000000in}{0.000000in}}{%
\pgfpathmoveto{\pgfqpoint{-0.000000in}{0.000000in}}%
\pgfpathlineto{\pgfqpoint{-0.041667in}{0.000000in}}%
\pgfusepath{stroke,fill}%
}%
\begin{pgfscope}%
\pgfsys@transformshift{3.038110in}{0.420000in}%
\pgfsys@useobject{currentmarker}{}%
\end{pgfscope}%
\end{pgfscope}%
\begin{pgfscope}%
\definecolor{textcolor}{rgb}{0.000000,0.000000,0.000000}%
\pgfsetstrokecolor{textcolor}%
\pgfsetfillcolor{textcolor}%
\pgftext[x=0.212731in, y=0.367193in, left, base]{\color{textcolor}{\rmfamily\fontsize{11.000000}{13.200000}\selectfont\catcode`\^=\active\def^{\ifmmode\sp\else\^{}\fi}\catcode`\%=\active\def%{\%}\ensuremath{-}0.05}}%
\end{pgfscope}%
\begin{pgfscope}%
\pgfsetbuttcap%
\pgfsetroundjoin%
\definecolor{currentfill}{rgb}{0.000000,0.000000,0.000000}%
\pgfsetfillcolor{currentfill}%
\pgfsetlinewidth{0.501875pt}%
\definecolor{currentstroke}{rgb}{0.000000,0.000000,0.000000}%
\pgfsetstrokecolor{currentstroke}%
\pgfsetdash{}{0pt}%
\pgfsys@defobject{currentmarker}{\pgfqpoint{0.000000in}{0.000000in}}{\pgfqpoint{0.041667in}{0.000000in}}{%
\pgfpathmoveto{\pgfqpoint{0.000000in}{0.000000in}}%
\pgfpathlineto{\pgfqpoint{0.041667in}{0.000000in}}%
\pgfusepath{stroke,fill}%
}%
\begin{pgfscope}%
\pgfsys@transformshift{0.650000in}{1.084829in}%
\pgfsys@useobject{currentmarker}{}%
\end{pgfscope}%
\end{pgfscope}%
\begin{pgfscope}%
\pgfsetbuttcap%
\pgfsetroundjoin%
\definecolor{currentfill}{rgb}{0.000000,0.000000,0.000000}%
\pgfsetfillcolor{currentfill}%
\pgfsetlinewidth{0.501875pt}%
\definecolor{currentstroke}{rgb}{0.000000,0.000000,0.000000}%
\pgfsetstrokecolor{currentstroke}%
\pgfsetdash{}{0pt}%
\pgfsys@defobject{currentmarker}{\pgfqpoint{-0.041667in}{0.000000in}}{\pgfqpoint{-0.000000in}{0.000000in}}{%
\pgfpathmoveto{\pgfqpoint{-0.000000in}{0.000000in}}%
\pgfpathlineto{\pgfqpoint{-0.041667in}{0.000000in}}%
\pgfusepath{stroke,fill}%
}%
\begin{pgfscope}%
\pgfsys@transformshift{3.038110in}{1.084829in}%
\pgfsys@useobject{currentmarker}{}%
\end{pgfscope}%
\end{pgfscope}%
\begin{pgfscope}%
\definecolor{textcolor}{rgb}{0.000000,0.000000,0.000000}%
\pgfsetstrokecolor{textcolor}%
\pgfsetfillcolor{textcolor}%
\pgftext[x=0.331018in, y=1.032023in, left, base]{\color{textcolor}{\rmfamily\fontsize{11.000000}{13.200000}\selectfont\catcode`\^=\active\def^{\ifmmode\sp\else\^{}\fi}\catcode`\%=\active\def%{\%}0.00}}%
\end{pgfscope}%
\begin{pgfscope}%
\pgfsetbuttcap%
\pgfsetroundjoin%
\definecolor{currentfill}{rgb}{0.000000,0.000000,0.000000}%
\pgfsetfillcolor{currentfill}%
\pgfsetlinewidth{0.501875pt}%
\definecolor{currentstroke}{rgb}{0.000000,0.000000,0.000000}%
\pgfsetstrokecolor{currentstroke}%
\pgfsetdash{}{0pt}%
\pgfsys@defobject{currentmarker}{\pgfqpoint{0.000000in}{0.000000in}}{\pgfqpoint{0.041667in}{0.000000in}}{%
\pgfpathmoveto{\pgfqpoint{0.000000in}{0.000000in}}%
\pgfpathlineto{\pgfqpoint{0.041667in}{0.000000in}}%
\pgfusepath{stroke,fill}%
}%
\begin{pgfscope}%
\pgfsys@transformshift{0.650000in}{1.749659in}%
\pgfsys@useobject{currentmarker}{}%
\end{pgfscope}%
\end{pgfscope}%
\begin{pgfscope}%
\pgfsetbuttcap%
\pgfsetroundjoin%
\definecolor{currentfill}{rgb}{0.000000,0.000000,0.000000}%
\pgfsetfillcolor{currentfill}%
\pgfsetlinewidth{0.501875pt}%
\definecolor{currentstroke}{rgb}{0.000000,0.000000,0.000000}%
\pgfsetstrokecolor{currentstroke}%
\pgfsetdash{}{0pt}%
\pgfsys@defobject{currentmarker}{\pgfqpoint{-0.041667in}{0.000000in}}{\pgfqpoint{-0.000000in}{0.000000in}}{%
\pgfpathmoveto{\pgfqpoint{-0.000000in}{0.000000in}}%
\pgfpathlineto{\pgfqpoint{-0.041667in}{0.000000in}}%
\pgfusepath{stroke,fill}%
}%
\begin{pgfscope}%
\pgfsys@transformshift{3.038110in}{1.749659in}%
\pgfsys@useobject{currentmarker}{}%
\end{pgfscope}%
\end{pgfscope}%
\begin{pgfscope}%
\definecolor{textcolor}{rgb}{0.000000,0.000000,0.000000}%
\pgfsetstrokecolor{textcolor}%
\pgfsetfillcolor{textcolor}%
\pgftext[x=0.331018in, y=1.696852in, left, base]{\color{textcolor}{\rmfamily\fontsize{11.000000}{13.200000}\selectfont\catcode`\^=\active\def^{\ifmmode\sp\else\^{}\fi}\catcode`\%=\active\def%{\%}0.05}}%
\end{pgfscope}%
\begin{pgfscope}%
\pgfsetbuttcap%
\pgfsetroundjoin%
\definecolor{currentfill}{rgb}{0.000000,0.000000,0.000000}%
\pgfsetfillcolor{currentfill}%
\pgfsetlinewidth{0.501875pt}%
\definecolor{currentstroke}{rgb}{0.000000,0.000000,0.000000}%
\pgfsetstrokecolor{currentstroke}%
\pgfsetdash{}{0pt}%
\pgfsys@defobject{currentmarker}{\pgfqpoint{0.000000in}{0.000000in}}{\pgfqpoint{0.041667in}{0.000000in}}{%
\pgfpathmoveto{\pgfqpoint{0.000000in}{0.000000in}}%
\pgfpathlineto{\pgfqpoint{0.041667in}{0.000000in}}%
\pgfusepath{stroke,fill}%
}%
\begin{pgfscope}%
\pgfsys@transformshift{0.650000in}{2.414488in}%
\pgfsys@useobject{currentmarker}{}%
\end{pgfscope}%
\end{pgfscope}%
\begin{pgfscope}%
\pgfsetbuttcap%
\pgfsetroundjoin%
\definecolor{currentfill}{rgb}{0.000000,0.000000,0.000000}%
\pgfsetfillcolor{currentfill}%
\pgfsetlinewidth{0.501875pt}%
\definecolor{currentstroke}{rgb}{0.000000,0.000000,0.000000}%
\pgfsetstrokecolor{currentstroke}%
\pgfsetdash{}{0pt}%
\pgfsys@defobject{currentmarker}{\pgfqpoint{-0.041667in}{0.000000in}}{\pgfqpoint{-0.000000in}{0.000000in}}{%
\pgfpathmoveto{\pgfqpoint{-0.000000in}{0.000000in}}%
\pgfpathlineto{\pgfqpoint{-0.041667in}{0.000000in}}%
\pgfusepath{stroke,fill}%
}%
\begin{pgfscope}%
\pgfsys@transformshift{3.038110in}{2.414488in}%
\pgfsys@useobject{currentmarker}{}%
\end{pgfscope}%
\end{pgfscope}%
\begin{pgfscope}%
\definecolor{textcolor}{rgb}{0.000000,0.000000,0.000000}%
\pgfsetstrokecolor{textcolor}%
\pgfsetfillcolor{textcolor}%
\pgftext[x=0.331018in, y=2.361681in, left, base]{\color{textcolor}{\rmfamily\fontsize{11.000000}{13.200000}\selectfont\catcode`\^=\active\def^{\ifmmode\sp\else\^{}\fi}\catcode`\%=\active\def%{\%}0.10}}%
\end{pgfscope}%
\begin{pgfscope}%
\pgfsetbuttcap%
\pgfsetroundjoin%
\definecolor{currentfill}{rgb}{0.000000,0.000000,0.000000}%
\pgfsetfillcolor{currentfill}%
\pgfsetlinewidth{0.401500pt}%
\definecolor{currentstroke}{rgb}{0.000000,0.000000,0.000000}%
\pgfsetstrokecolor{currentstroke}%
\pgfsetdash{}{0pt}%
\pgfsys@defobject{currentmarker}{\pgfqpoint{0.000000in}{0.000000in}}{\pgfqpoint{0.020833in}{0.000000in}}{%
\pgfpathmoveto{\pgfqpoint{0.000000in}{0.000000in}}%
\pgfpathlineto{\pgfqpoint{0.020833in}{0.000000in}}%
\pgfusepath{stroke,fill}%
}%
\begin{pgfscope}%
\pgfsys@transformshift{0.650000in}{0.552966in}%
\pgfsys@useobject{currentmarker}{}%
\end{pgfscope}%
\end{pgfscope}%
\begin{pgfscope}%
\pgfsetbuttcap%
\pgfsetroundjoin%
\definecolor{currentfill}{rgb}{0.000000,0.000000,0.000000}%
\pgfsetfillcolor{currentfill}%
\pgfsetlinewidth{0.401500pt}%
\definecolor{currentstroke}{rgb}{0.000000,0.000000,0.000000}%
\pgfsetstrokecolor{currentstroke}%
\pgfsetdash{}{0pt}%
\pgfsys@defobject{currentmarker}{\pgfqpoint{-0.020833in}{0.000000in}}{\pgfqpoint{-0.000000in}{0.000000in}}{%
\pgfpathmoveto{\pgfqpoint{-0.000000in}{0.000000in}}%
\pgfpathlineto{\pgfqpoint{-0.020833in}{0.000000in}}%
\pgfusepath{stroke,fill}%
}%
\begin{pgfscope}%
\pgfsys@transformshift{3.038110in}{0.552966in}%
\pgfsys@useobject{currentmarker}{}%
\end{pgfscope}%
\end{pgfscope}%
\begin{pgfscope}%
\pgfsetbuttcap%
\pgfsetroundjoin%
\definecolor{currentfill}{rgb}{0.000000,0.000000,0.000000}%
\pgfsetfillcolor{currentfill}%
\pgfsetlinewidth{0.401500pt}%
\definecolor{currentstroke}{rgb}{0.000000,0.000000,0.000000}%
\pgfsetstrokecolor{currentstroke}%
\pgfsetdash{}{0pt}%
\pgfsys@defobject{currentmarker}{\pgfqpoint{0.000000in}{0.000000in}}{\pgfqpoint{0.020833in}{0.000000in}}{%
\pgfpathmoveto{\pgfqpoint{0.000000in}{0.000000in}}%
\pgfpathlineto{\pgfqpoint{0.020833in}{0.000000in}}%
\pgfusepath{stroke,fill}%
}%
\begin{pgfscope}%
\pgfsys@transformshift{0.650000in}{0.685932in}%
\pgfsys@useobject{currentmarker}{}%
\end{pgfscope}%
\end{pgfscope}%
\begin{pgfscope}%
\pgfsetbuttcap%
\pgfsetroundjoin%
\definecolor{currentfill}{rgb}{0.000000,0.000000,0.000000}%
\pgfsetfillcolor{currentfill}%
\pgfsetlinewidth{0.401500pt}%
\definecolor{currentstroke}{rgb}{0.000000,0.000000,0.000000}%
\pgfsetstrokecolor{currentstroke}%
\pgfsetdash{}{0pt}%
\pgfsys@defobject{currentmarker}{\pgfqpoint{-0.020833in}{0.000000in}}{\pgfqpoint{-0.000000in}{0.000000in}}{%
\pgfpathmoveto{\pgfqpoint{-0.000000in}{0.000000in}}%
\pgfpathlineto{\pgfqpoint{-0.020833in}{0.000000in}}%
\pgfusepath{stroke,fill}%
}%
\begin{pgfscope}%
\pgfsys@transformshift{3.038110in}{0.685932in}%
\pgfsys@useobject{currentmarker}{}%
\end{pgfscope}%
\end{pgfscope}%
\begin{pgfscope}%
\pgfsetbuttcap%
\pgfsetroundjoin%
\definecolor{currentfill}{rgb}{0.000000,0.000000,0.000000}%
\pgfsetfillcolor{currentfill}%
\pgfsetlinewidth{0.401500pt}%
\definecolor{currentstroke}{rgb}{0.000000,0.000000,0.000000}%
\pgfsetstrokecolor{currentstroke}%
\pgfsetdash{}{0pt}%
\pgfsys@defobject{currentmarker}{\pgfqpoint{0.000000in}{0.000000in}}{\pgfqpoint{0.020833in}{0.000000in}}{%
\pgfpathmoveto{\pgfqpoint{0.000000in}{0.000000in}}%
\pgfpathlineto{\pgfqpoint{0.020833in}{0.000000in}}%
\pgfusepath{stroke,fill}%
}%
\begin{pgfscope}%
\pgfsys@transformshift{0.650000in}{0.818898in}%
\pgfsys@useobject{currentmarker}{}%
\end{pgfscope}%
\end{pgfscope}%
\begin{pgfscope}%
\pgfsetbuttcap%
\pgfsetroundjoin%
\definecolor{currentfill}{rgb}{0.000000,0.000000,0.000000}%
\pgfsetfillcolor{currentfill}%
\pgfsetlinewidth{0.401500pt}%
\definecolor{currentstroke}{rgb}{0.000000,0.000000,0.000000}%
\pgfsetstrokecolor{currentstroke}%
\pgfsetdash{}{0pt}%
\pgfsys@defobject{currentmarker}{\pgfqpoint{-0.020833in}{0.000000in}}{\pgfqpoint{-0.000000in}{0.000000in}}{%
\pgfpathmoveto{\pgfqpoint{-0.000000in}{0.000000in}}%
\pgfpathlineto{\pgfqpoint{-0.020833in}{0.000000in}}%
\pgfusepath{stroke,fill}%
}%
\begin{pgfscope}%
\pgfsys@transformshift{3.038110in}{0.818898in}%
\pgfsys@useobject{currentmarker}{}%
\end{pgfscope}%
\end{pgfscope}%
\begin{pgfscope}%
\pgfsetbuttcap%
\pgfsetroundjoin%
\definecolor{currentfill}{rgb}{0.000000,0.000000,0.000000}%
\pgfsetfillcolor{currentfill}%
\pgfsetlinewidth{0.401500pt}%
\definecolor{currentstroke}{rgb}{0.000000,0.000000,0.000000}%
\pgfsetstrokecolor{currentstroke}%
\pgfsetdash{}{0pt}%
\pgfsys@defobject{currentmarker}{\pgfqpoint{0.000000in}{0.000000in}}{\pgfqpoint{0.020833in}{0.000000in}}{%
\pgfpathmoveto{\pgfqpoint{0.000000in}{0.000000in}}%
\pgfpathlineto{\pgfqpoint{0.020833in}{0.000000in}}%
\pgfusepath{stroke,fill}%
}%
\begin{pgfscope}%
\pgfsys@transformshift{0.650000in}{0.951863in}%
\pgfsys@useobject{currentmarker}{}%
\end{pgfscope}%
\end{pgfscope}%
\begin{pgfscope}%
\pgfsetbuttcap%
\pgfsetroundjoin%
\definecolor{currentfill}{rgb}{0.000000,0.000000,0.000000}%
\pgfsetfillcolor{currentfill}%
\pgfsetlinewidth{0.401500pt}%
\definecolor{currentstroke}{rgb}{0.000000,0.000000,0.000000}%
\pgfsetstrokecolor{currentstroke}%
\pgfsetdash{}{0pt}%
\pgfsys@defobject{currentmarker}{\pgfqpoint{-0.020833in}{0.000000in}}{\pgfqpoint{-0.000000in}{0.000000in}}{%
\pgfpathmoveto{\pgfqpoint{-0.000000in}{0.000000in}}%
\pgfpathlineto{\pgfqpoint{-0.020833in}{0.000000in}}%
\pgfusepath{stroke,fill}%
}%
\begin{pgfscope}%
\pgfsys@transformshift{3.038110in}{0.951863in}%
\pgfsys@useobject{currentmarker}{}%
\end{pgfscope}%
\end{pgfscope}%
\begin{pgfscope}%
\pgfsetbuttcap%
\pgfsetroundjoin%
\definecolor{currentfill}{rgb}{0.000000,0.000000,0.000000}%
\pgfsetfillcolor{currentfill}%
\pgfsetlinewidth{0.401500pt}%
\definecolor{currentstroke}{rgb}{0.000000,0.000000,0.000000}%
\pgfsetstrokecolor{currentstroke}%
\pgfsetdash{}{0pt}%
\pgfsys@defobject{currentmarker}{\pgfqpoint{0.000000in}{0.000000in}}{\pgfqpoint{0.020833in}{0.000000in}}{%
\pgfpathmoveto{\pgfqpoint{0.000000in}{0.000000in}}%
\pgfpathlineto{\pgfqpoint{0.020833in}{0.000000in}}%
\pgfusepath{stroke,fill}%
}%
\begin{pgfscope}%
\pgfsys@transformshift{0.650000in}{1.217795in}%
\pgfsys@useobject{currentmarker}{}%
\end{pgfscope}%
\end{pgfscope}%
\begin{pgfscope}%
\pgfsetbuttcap%
\pgfsetroundjoin%
\definecolor{currentfill}{rgb}{0.000000,0.000000,0.000000}%
\pgfsetfillcolor{currentfill}%
\pgfsetlinewidth{0.401500pt}%
\definecolor{currentstroke}{rgb}{0.000000,0.000000,0.000000}%
\pgfsetstrokecolor{currentstroke}%
\pgfsetdash{}{0pt}%
\pgfsys@defobject{currentmarker}{\pgfqpoint{-0.020833in}{0.000000in}}{\pgfqpoint{-0.000000in}{0.000000in}}{%
\pgfpathmoveto{\pgfqpoint{-0.000000in}{0.000000in}}%
\pgfpathlineto{\pgfqpoint{-0.020833in}{0.000000in}}%
\pgfusepath{stroke,fill}%
}%
\begin{pgfscope}%
\pgfsys@transformshift{3.038110in}{1.217795in}%
\pgfsys@useobject{currentmarker}{}%
\end{pgfscope}%
\end{pgfscope}%
\begin{pgfscope}%
\pgfsetbuttcap%
\pgfsetroundjoin%
\definecolor{currentfill}{rgb}{0.000000,0.000000,0.000000}%
\pgfsetfillcolor{currentfill}%
\pgfsetlinewidth{0.401500pt}%
\definecolor{currentstroke}{rgb}{0.000000,0.000000,0.000000}%
\pgfsetstrokecolor{currentstroke}%
\pgfsetdash{}{0pt}%
\pgfsys@defobject{currentmarker}{\pgfqpoint{0.000000in}{0.000000in}}{\pgfqpoint{0.020833in}{0.000000in}}{%
\pgfpathmoveto{\pgfqpoint{0.000000in}{0.000000in}}%
\pgfpathlineto{\pgfqpoint{0.020833in}{0.000000in}}%
\pgfusepath{stroke,fill}%
}%
\begin{pgfscope}%
\pgfsys@transformshift{0.650000in}{1.350761in}%
\pgfsys@useobject{currentmarker}{}%
\end{pgfscope}%
\end{pgfscope}%
\begin{pgfscope}%
\pgfsetbuttcap%
\pgfsetroundjoin%
\definecolor{currentfill}{rgb}{0.000000,0.000000,0.000000}%
\pgfsetfillcolor{currentfill}%
\pgfsetlinewidth{0.401500pt}%
\definecolor{currentstroke}{rgb}{0.000000,0.000000,0.000000}%
\pgfsetstrokecolor{currentstroke}%
\pgfsetdash{}{0pt}%
\pgfsys@defobject{currentmarker}{\pgfqpoint{-0.020833in}{0.000000in}}{\pgfqpoint{-0.000000in}{0.000000in}}{%
\pgfpathmoveto{\pgfqpoint{-0.000000in}{0.000000in}}%
\pgfpathlineto{\pgfqpoint{-0.020833in}{0.000000in}}%
\pgfusepath{stroke,fill}%
}%
\begin{pgfscope}%
\pgfsys@transformshift{3.038110in}{1.350761in}%
\pgfsys@useobject{currentmarker}{}%
\end{pgfscope}%
\end{pgfscope}%
\begin{pgfscope}%
\pgfsetbuttcap%
\pgfsetroundjoin%
\definecolor{currentfill}{rgb}{0.000000,0.000000,0.000000}%
\pgfsetfillcolor{currentfill}%
\pgfsetlinewidth{0.401500pt}%
\definecolor{currentstroke}{rgb}{0.000000,0.000000,0.000000}%
\pgfsetstrokecolor{currentstroke}%
\pgfsetdash{}{0pt}%
\pgfsys@defobject{currentmarker}{\pgfqpoint{0.000000in}{0.000000in}}{\pgfqpoint{0.020833in}{0.000000in}}{%
\pgfpathmoveto{\pgfqpoint{0.000000in}{0.000000in}}%
\pgfpathlineto{\pgfqpoint{0.020833in}{0.000000in}}%
\pgfusepath{stroke,fill}%
}%
\begin{pgfscope}%
\pgfsys@transformshift{0.650000in}{1.483727in}%
\pgfsys@useobject{currentmarker}{}%
\end{pgfscope}%
\end{pgfscope}%
\begin{pgfscope}%
\pgfsetbuttcap%
\pgfsetroundjoin%
\definecolor{currentfill}{rgb}{0.000000,0.000000,0.000000}%
\pgfsetfillcolor{currentfill}%
\pgfsetlinewidth{0.401500pt}%
\definecolor{currentstroke}{rgb}{0.000000,0.000000,0.000000}%
\pgfsetstrokecolor{currentstroke}%
\pgfsetdash{}{0pt}%
\pgfsys@defobject{currentmarker}{\pgfqpoint{-0.020833in}{0.000000in}}{\pgfqpoint{-0.000000in}{0.000000in}}{%
\pgfpathmoveto{\pgfqpoint{-0.000000in}{0.000000in}}%
\pgfpathlineto{\pgfqpoint{-0.020833in}{0.000000in}}%
\pgfusepath{stroke,fill}%
}%
\begin{pgfscope}%
\pgfsys@transformshift{3.038110in}{1.483727in}%
\pgfsys@useobject{currentmarker}{}%
\end{pgfscope}%
\end{pgfscope}%
\begin{pgfscope}%
\pgfsetbuttcap%
\pgfsetroundjoin%
\definecolor{currentfill}{rgb}{0.000000,0.000000,0.000000}%
\pgfsetfillcolor{currentfill}%
\pgfsetlinewidth{0.401500pt}%
\definecolor{currentstroke}{rgb}{0.000000,0.000000,0.000000}%
\pgfsetstrokecolor{currentstroke}%
\pgfsetdash{}{0pt}%
\pgfsys@defobject{currentmarker}{\pgfqpoint{0.000000in}{0.000000in}}{\pgfqpoint{0.020833in}{0.000000in}}{%
\pgfpathmoveto{\pgfqpoint{0.000000in}{0.000000in}}%
\pgfpathlineto{\pgfqpoint{0.020833in}{0.000000in}}%
\pgfusepath{stroke,fill}%
}%
\begin{pgfscope}%
\pgfsys@transformshift{0.650000in}{1.616693in}%
\pgfsys@useobject{currentmarker}{}%
\end{pgfscope}%
\end{pgfscope}%
\begin{pgfscope}%
\pgfsetbuttcap%
\pgfsetroundjoin%
\definecolor{currentfill}{rgb}{0.000000,0.000000,0.000000}%
\pgfsetfillcolor{currentfill}%
\pgfsetlinewidth{0.401500pt}%
\definecolor{currentstroke}{rgb}{0.000000,0.000000,0.000000}%
\pgfsetstrokecolor{currentstroke}%
\pgfsetdash{}{0pt}%
\pgfsys@defobject{currentmarker}{\pgfqpoint{-0.020833in}{0.000000in}}{\pgfqpoint{-0.000000in}{0.000000in}}{%
\pgfpathmoveto{\pgfqpoint{-0.000000in}{0.000000in}}%
\pgfpathlineto{\pgfqpoint{-0.020833in}{0.000000in}}%
\pgfusepath{stroke,fill}%
}%
\begin{pgfscope}%
\pgfsys@transformshift{3.038110in}{1.616693in}%
\pgfsys@useobject{currentmarker}{}%
\end{pgfscope}%
\end{pgfscope}%
\begin{pgfscope}%
\pgfsetbuttcap%
\pgfsetroundjoin%
\definecolor{currentfill}{rgb}{0.000000,0.000000,0.000000}%
\pgfsetfillcolor{currentfill}%
\pgfsetlinewidth{0.401500pt}%
\definecolor{currentstroke}{rgb}{0.000000,0.000000,0.000000}%
\pgfsetstrokecolor{currentstroke}%
\pgfsetdash{}{0pt}%
\pgfsys@defobject{currentmarker}{\pgfqpoint{0.000000in}{0.000000in}}{\pgfqpoint{0.020833in}{0.000000in}}{%
\pgfpathmoveto{\pgfqpoint{0.000000in}{0.000000in}}%
\pgfpathlineto{\pgfqpoint{0.020833in}{0.000000in}}%
\pgfusepath{stroke,fill}%
}%
\begin{pgfscope}%
\pgfsys@transformshift{0.650000in}{1.882625in}%
\pgfsys@useobject{currentmarker}{}%
\end{pgfscope}%
\end{pgfscope}%
\begin{pgfscope}%
\pgfsetbuttcap%
\pgfsetroundjoin%
\definecolor{currentfill}{rgb}{0.000000,0.000000,0.000000}%
\pgfsetfillcolor{currentfill}%
\pgfsetlinewidth{0.401500pt}%
\definecolor{currentstroke}{rgb}{0.000000,0.000000,0.000000}%
\pgfsetstrokecolor{currentstroke}%
\pgfsetdash{}{0pt}%
\pgfsys@defobject{currentmarker}{\pgfqpoint{-0.020833in}{0.000000in}}{\pgfqpoint{-0.000000in}{0.000000in}}{%
\pgfpathmoveto{\pgfqpoint{-0.000000in}{0.000000in}}%
\pgfpathlineto{\pgfqpoint{-0.020833in}{0.000000in}}%
\pgfusepath{stroke,fill}%
}%
\begin{pgfscope}%
\pgfsys@transformshift{3.038110in}{1.882625in}%
\pgfsys@useobject{currentmarker}{}%
\end{pgfscope}%
\end{pgfscope}%
\begin{pgfscope}%
\pgfsetbuttcap%
\pgfsetroundjoin%
\definecolor{currentfill}{rgb}{0.000000,0.000000,0.000000}%
\pgfsetfillcolor{currentfill}%
\pgfsetlinewidth{0.401500pt}%
\definecolor{currentstroke}{rgb}{0.000000,0.000000,0.000000}%
\pgfsetstrokecolor{currentstroke}%
\pgfsetdash{}{0pt}%
\pgfsys@defobject{currentmarker}{\pgfqpoint{0.000000in}{0.000000in}}{\pgfqpoint{0.020833in}{0.000000in}}{%
\pgfpathmoveto{\pgfqpoint{0.000000in}{0.000000in}}%
\pgfpathlineto{\pgfqpoint{0.020833in}{0.000000in}}%
\pgfusepath{stroke,fill}%
}%
\begin{pgfscope}%
\pgfsys@transformshift{0.650000in}{2.015590in}%
\pgfsys@useobject{currentmarker}{}%
\end{pgfscope}%
\end{pgfscope}%
\begin{pgfscope}%
\pgfsetbuttcap%
\pgfsetroundjoin%
\definecolor{currentfill}{rgb}{0.000000,0.000000,0.000000}%
\pgfsetfillcolor{currentfill}%
\pgfsetlinewidth{0.401500pt}%
\definecolor{currentstroke}{rgb}{0.000000,0.000000,0.000000}%
\pgfsetstrokecolor{currentstroke}%
\pgfsetdash{}{0pt}%
\pgfsys@defobject{currentmarker}{\pgfqpoint{-0.020833in}{0.000000in}}{\pgfqpoint{-0.000000in}{0.000000in}}{%
\pgfpathmoveto{\pgfqpoint{-0.000000in}{0.000000in}}%
\pgfpathlineto{\pgfqpoint{-0.020833in}{0.000000in}}%
\pgfusepath{stroke,fill}%
}%
\begin{pgfscope}%
\pgfsys@transformshift{3.038110in}{2.015590in}%
\pgfsys@useobject{currentmarker}{}%
\end{pgfscope}%
\end{pgfscope}%
\begin{pgfscope}%
\pgfsetbuttcap%
\pgfsetroundjoin%
\definecolor{currentfill}{rgb}{0.000000,0.000000,0.000000}%
\pgfsetfillcolor{currentfill}%
\pgfsetlinewidth{0.401500pt}%
\definecolor{currentstroke}{rgb}{0.000000,0.000000,0.000000}%
\pgfsetstrokecolor{currentstroke}%
\pgfsetdash{}{0pt}%
\pgfsys@defobject{currentmarker}{\pgfqpoint{0.000000in}{0.000000in}}{\pgfqpoint{0.020833in}{0.000000in}}{%
\pgfpathmoveto{\pgfqpoint{0.000000in}{0.000000in}}%
\pgfpathlineto{\pgfqpoint{0.020833in}{0.000000in}}%
\pgfusepath{stroke,fill}%
}%
\begin{pgfscope}%
\pgfsys@transformshift{0.650000in}{2.148556in}%
\pgfsys@useobject{currentmarker}{}%
\end{pgfscope}%
\end{pgfscope}%
\begin{pgfscope}%
\pgfsetbuttcap%
\pgfsetroundjoin%
\definecolor{currentfill}{rgb}{0.000000,0.000000,0.000000}%
\pgfsetfillcolor{currentfill}%
\pgfsetlinewidth{0.401500pt}%
\definecolor{currentstroke}{rgb}{0.000000,0.000000,0.000000}%
\pgfsetstrokecolor{currentstroke}%
\pgfsetdash{}{0pt}%
\pgfsys@defobject{currentmarker}{\pgfqpoint{-0.020833in}{0.000000in}}{\pgfqpoint{-0.000000in}{0.000000in}}{%
\pgfpathmoveto{\pgfqpoint{-0.000000in}{0.000000in}}%
\pgfpathlineto{\pgfqpoint{-0.020833in}{0.000000in}}%
\pgfusepath{stroke,fill}%
}%
\begin{pgfscope}%
\pgfsys@transformshift{3.038110in}{2.148556in}%
\pgfsys@useobject{currentmarker}{}%
\end{pgfscope}%
\end{pgfscope}%
\begin{pgfscope}%
\pgfsetbuttcap%
\pgfsetroundjoin%
\definecolor{currentfill}{rgb}{0.000000,0.000000,0.000000}%
\pgfsetfillcolor{currentfill}%
\pgfsetlinewidth{0.401500pt}%
\definecolor{currentstroke}{rgb}{0.000000,0.000000,0.000000}%
\pgfsetstrokecolor{currentstroke}%
\pgfsetdash{}{0pt}%
\pgfsys@defobject{currentmarker}{\pgfqpoint{0.000000in}{0.000000in}}{\pgfqpoint{0.020833in}{0.000000in}}{%
\pgfpathmoveto{\pgfqpoint{0.000000in}{0.000000in}}%
\pgfpathlineto{\pgfqpoint{0.020833in}{0.000000in}}%
\pgfusepath{stroke,fill}%
}%
\begin{pgfscope}%
\pgfsys@transformshift{0.650000in}{2.281522in}%
\pgfsys@useobject{currentmarker}{}%
\end{pgfscope}%
\end{pgfscope}%
\begin{pgfscope}%
\pgfsetbuttcap%
\pgfsetroundjoin%
\definecolor{currentfill}{rgb}{0.000000,0.000000,0.000000}%
\pgfsetfillcolor{currentfill}%
\pgfsetlinewidth{0.401500pt}%
\definecolor{currentstroke}{rgb}{0.000000,0.000000,0.000000}%
\pgfsetstrokecolor{currentstroke}%
\pgfsetdash{}{0pt}%
\pgfsys@defobject{currentmarker}{\pgfqpoint{-0.020833in}{0.000000in}}{\pgfqpoint{-0.000000in}{0.000000in}}{%
\pgfpathmoveto{\pgfqpoint{-0.000000in}{0.000000in}}%
\pgfpathlineto{\pgfqpoint{-0.020833in}{0.000000in}}%
\pgfusepath{stroke,fill}%
}%
\begin{pgfscope}%
\pgfsys@transformshift{3.038110in}{2.281522in}%
\pgfsys@useobject{currentmarker}{}%
\end{pgfscope}%
\end{pgfscope}%
\begin{pgfscope}%
\definecolor{textcolor}{rgb}{0.000000,0.000000,0.000000}%
\pgfsetstrokecolor{textcolor}%
\pgfsetfillcolor{textcolor}%
\pgftext[x=0.184953in,y=1.417244in,,bottom,rotate=90.000000]{\color{textcolor}{\rmfamily\fontsize{12.000000}{14.400000}\selectfont\catcode`\^=\active\def^{\ifmmode\sp\else\^{}\fi}\catcode`\%=\active\def%{\%}$p_w-\langle p_{c}\rangle$}}%
\end{pgfscope}%
\begin{pgfscope}%
\pgfpathrectangle{\pgfqpoint{0.650000in}{0.420000in}}{\pgfqpoint{2.388110in}{1.994488in}}%
\pgfusepath{clip}%
\pgfsetrectcap%
\pgfsetroundjoin%
\pgfsetlinewidth{1.003750pt}%
\definecolor{currentstroke}{rgb}{0.796078,0.094118,0.113725}%
\pgfsetstrokecolor{currentstroke}%
\pgfsetdash{}{0pt}%
\pgfpathmoveto{\pgfqpoint{0.650000in}{2.229023in}}%
\pgfpathlineto{\pgfqpoint{0.652150in}{2.223975in}}%
\pgfpathlineto{\pgfqpoint{0.655376in}{2.212991in}}%
\pgfpathlineto{\pgfqpoint{0.661111in}{2.193097in}}%
\pgfpathlineto{\pgfqpoint{0.662903in}{2.190360in}}%
\pgfpathlineto{\pgfqpoint{0.664336in}{2.190153in}}%
\pgfpathlineto{\pgfqpoint{0.668637in}{2.191504in}}%
\pgfpathlineto{\pgfqpoint{0.673655in}{2.189854in}}%
\pgfpathlineto{\pgfqpoint{0.675089in}{2.191178in}}%
\pgfpathlineto{\pgfqpoint{0.677239in}{2.195025in}}%
\pgfpathlineto{\pgfqpoint{0.681540in}{2.206115in}}%
\pgfpathlineto{\pgfqpoint{0.686200in}{2.216916in}}%
\pgfpathlineto{\pgfqpoint{0.689425in}{2.221670in}}%
\pgfpathlineto{\pgfqpoint{0.692293in}{2.224137in}}%
\pgfpathlineto{\pgfqpoint{0.695160in}{2.225158in}}%
\pgfpathlineto{\pgfqpoint{0.697669in}{2.224898in}}%
\pgfpathlineto{\pgfqpoint{0.700536in}{2.223481in}}%
\pgfpathlineto{\pgfqpoint{0.704479in}{2.220198in}}%
\pgfpathlineto{\pgfqpoint{0.709138in}{2.214829in}}%
\pgfpathlineto{\pgfqpoint{0.717381in}{2.203014in}}%
\pgfpathlineto{\pgfqpoint{0.736019in}{2.176958in}}%
\pgfpathlineto{\pgfqpoint{0.744621in}{2.167021in}}%
\pgfpathlineto{\pgfqpoint{0.760032in}{2.151016in}}%
\pgfpathlineto{\pgfqpoint{0.770785in}{2.140848in}}%
\pgfpathlineto{\pgfqpoint{0.803400in}{2.112308in}}%
\pgfpathlineto{\pgfqpoint{0.817378in}{2.100493in}}%
\pgfpathlineto{\pgfqpoint{0.835657in}{2.084967in}}%
\pgfpathlineto{\pgfqpoint{0.876875in}{2.051433in}}%
\pgfpathlineto{\pgfqpoint{0.908415in}{2.026782in}}%
\pgfpathlineto{\pgfqpoint{0.948199in}{1.996484in}}%
\pgfpathlineto{\pgfqpoint{0.961818in}{1.985924in}}%
\pgfpathlineto{\pgfqpoint{1.015222in}{1.945449in}}%
\pgfpathlineto{\pgfqpoint{1.028841in}{1.935970in}}%
\pgfpathlineto{\pgfqpoint{1.043178in}{1.924737in}}%
\pgfpathlineto{\pgfqpoint{1.072926in}{1.900473in}}%
\pgfpathlineto{\pgfqpoint{1.081528in}{1.891301in}}%
\pgfpathlineto{\pgfqpoint{1.090846in}{1.879702in}}%
\pgfpathlineto{\pgfqpoint{1.098015in}{1.869440in}}%
\pgfpathlineto{\pgfqpoint{1.102674in}{1.860807in}}%
\pgfpathlineto{\pgfqpoint{1.108409in}{1.847603in}}%
\pgfpathlineto{\pgfqpoint{1.113785in}{1.832396in}}%
\pgfpathlineto{\pgfqpoint{1.118803in}{1.814540in}}%
\pgfpathlineto{\pgfqpoint{1.123462in}{1.793973in}}%
\pgfpathlineto{\pgfqpoint{1.127763in}{1.769935in}}%
\pgfpathlineto{\pgfqpoint{1.131705in}{1.741637in}}%
\pgfpathlineto{\pgfqpoint{1.136365in}{1.698843in}}%
\pgfpathlineto{\pgfqpoint{1.142816in}{1.640682in}}%
\pgfpathlineto{\pgfqpoint{1.147476in}{1.609414in}}%
\pgfpathlineto{\pgfqpoint{1.152135in}{1.585251in}}%
\pgfpathlineto{\pgfqpoint{1.157153in}{1.564886in}}%
\pgfpathlineto{\pgfqpoint{1.162171in}{1.548907in}}%
\pgfpathlineto{\pgfqpoint{1.167547in}{1.535261in}}%
\pgfpathlineto{\pgfqpoint{1.173281in}{1.523794in}}%
\pgfpathlineto{\pgfqpoint{1.179733in}{1.513391in}}%
\pgfpathlineto{\pgfqpoint{1.186184in}{1.504976in}}%
\pgfpathlineto{\pgfqpoint{1.194069in}{1.496703in}}%
\pgfpathlineto{\pgfqpoint{1.207689in}{1.483954in}}%
\pgfpathlineto{\pgfqpoint{1.222026in}{1.471737in}}%
\pgfpathlineto{\pgfqpoint{1.229553in}{1.465910in}}%
\pgfpathlineto{\pgfqpoint{1.252491in}{1.449472in}}%
\pgfpathlineto{\pgfqpoint{1.276863in}{1.432025in}}%
\pgfpathlineto{\pgfqpoint{1.287257in}{1.424734in}}%
\pgfpathlineto{\pgfqpoint{1.294784in}{1.419462in}}%
\pgfpathlineto{\pgfqpoint{1.303744in}{1.413411in}}%
\pgfpathlineto{\pgfqpoint{1.333134in}{1.392539in}}%
\pgfpathlineto{\pgfqpoint{1.339586in}{1.387845in}}%
\pgfpathlineto{\pgfqpoint{1.349263in}{1.381085in}}%
\pgfpathlineto{\pgfqpoint{1.414853in}{1.338654in}}%
\pgfpathlineto{\pgfqpoint{1.429189in}{1.329650in}}%
\pgfpathlineto{\pgfqpoint{1.444243in}{1.319992in}}%
\pgfpathlineto{\pgfqpoint{1.454279in}{1.313442in}}%
\pgfpathlineto{\pgfqpoint{1.489762in}{1.288785in}}%
\pgfpathlineto{\pgfqpoint{1.510191in}{1.276931in}}%
\pgfpathlineto{\pgfqpoint{1.519510in}{1.271039in}}%
\pgfpathlineto{\pgfqpoint{1.534205in}{1.261548in}}%
\pgfpathlineto{\pgfqpoint{1.550692in}{1.251012in}}%
\pgfpathlineto{\pgfqpoint{1.566462in}{1.241747in}}%
\pgfpathlineto{\pgfqpoint{1.596927in}{1.222960in}}%
\pgfpathlineto{\pgfqpoint{1.606246in}{1.217082in}}%
\pgfpathlineto{\pgfqpoint{1.616282in}{1.210729in}}%
\pgfpathlineto{\pgfqpoint{1.633127in}{1.200236in}}%
\pgfpathlineto{\pgfqpoint{1.671477in}{1.177503in}}%
\pgfpathlineto{\pgfqpoint{1.676854in}{1.174511in}}%
\pgfpathlineto{\pgfqpoint{1.692624in}{1.164132in}}%
\pgfpathlineto{\pgfqpoint{1.742085in}{1.135081in}}%
\pgfpathlineto{\pgfqpoint{1.749612in}{1.130425in}}%
\pgfpathlineto{\pgfqpoint{1.770041in}{1.118084in}}%
\pgfpathlineto{\pgfqpoint{1.780794in}{1.112032in}}%
\pgfpathlineto{\pgfqpoint{1.814127in}{1.091332in}}%
\pgfpathlineto{\pgfqpoint{1.836707in}{1.078184in}}%
\pgfpathlineto{\pgfqpoint{1.847459in}{1.072092in}}%
\pgfpathlineto{\pgfqpoint{1.853911in}{1.068654in}}%
\pgfpathlineto{\pgfqpoint{1.862154in}{1.063867in}}%
\pgfpathlineto{\pgfqpoint{1.872548in}{1.058365in}}%
\pgfpathlineto{\pgfqpoint{1.912332in}{1.035938in}}%
\pgfpathlineto{\pgfqpoint{1.919141in}{1.031840in}}%
\pgfpathlineto{\pgfqpoint{1.930969in}{1.024861in}}%
\pgfpathlineto{\pgfqpoint{1.959283in}{1.008912in}}%
\pgfpathlineto{\pgfqpoint{1.967527in}{1.005104in}}%
\pgfpathlineto{\pgfqpoint{1.980788in}{0.997535in}}%
\pgfpathlineto{\pgfqpoint{1.989390in}{0.992978in}}%
\pgfpathlineto{\pgfqpoint{1.999067in}{0.988159in}}%
\pgfpathlineto{\pgfqpoint{2.019138in}{0.977656in}}%
\pgfpathlineto{\pgfqpoint{2.029890in}{0.971919in}}%
\pgfpathlineto{\pgfqpoint{2.037417in}{0.967572in}}%
\pgfpathlineto{\pgfqpoint{2.055338in}{0.957212in}}%
\pgfpathlineto{\pgfqpoint{2.096196in}{0.935237in}}%
\pgfpathlineto{\pgfqpoint{2.102289in}{0.931869in}}%
\pgfpathlineto{\pgfqpoint{2.111608in}{0.926746in}}%
\pgfpathlineto{\pgfqpoint{2.124511in}{0.919566in}}%
\pgfpathlineto{\pgfqpoint{2.168953in}{0.896616in}}%
\pgfpathlineto{\pgfqpoint{2.182214in}{0.889134in}}%
\pgfpathlineto{\pgfqpoint{2.189740in}{0.885084in}}%
\pgfpathlineto{\pgfqpoint{2.202642in}{0.876951in}}%
\pgfpathlineto{\pgfqpoint{2.216261in}{0.869913in}}%
\pgfpathlineto{\pgfqpoint{2.252100in}{0.851529in}}%
\pgfpathlineto{\pgfqpoint{2.257834in}{0.848184in}}%
\pgfpathlineto{\pgfqpoint{2.267152in}{0.843480in}}%
\pgfpathlineto{\pgfqpoint{2.278979in}{0.837023in}}%
\pgfpathlineto{\pgfqpoint{2.297973in}{0.827540in}}%
\pgfpathlineto{\pgfqpoint{2.306933in}{0.822627in}}%
\pgfpathlineto{\pgfqpoint{2.329511in}{0.811393in}}%
\pgfpathlineto{\pgfqpoint{2.336321in}{0.807373in}}%
\pgfpathlineto{\pgfqpoint{2.344922in}{0.802753in}}%
\pgfpathlineto{\pgfqpoint{2.363916in}{0.792381in}}%
\pgfpathlineto{\pgfqpoint{2.369292in}{0.789619in}}%
\pgfpathlineto{\pgfqpoint{2.376460in}{0.786349in}}%
\pgfpathlineto{\pgfqpoint{2.383986in}{0.782757in}}%
\pgfpathlineto{\pgfqpoint{2.394379in}{0.777836in}}%
\pgfpathlineto{\pgfqpoint{2.406923in}{0.771397in}}%
\pgfpathlineto{\pgfqpoint{2.413374in}{0.768273in}}%
\pgfpathlineto{\pgfqpoint{2.430935in}{0.758929in}}%
\pgfpathlineto{\pgfqpoint{2.442403in}{0.752714in}}%
\pgfpathlineto{\pgfqpoint{2.453872in}{0.746800in}}%
\pgfpathlineto{\pgfqpoint{2.487561in}{0.731482in}}%
\pgfpathlineto{\pgfqpoint{2.503688in}{0.725434in}}%
\pgfpathlineto{\pgfqpoint{2.512648in}{0.723657in}}%
\pgfpathlineto{\pgfqpoint{2.523400in}{0.722563in}}%
\pgfpathlineto{\pgfqpoint{2.528418in}{0.723435in}}%
\pgfpathlineto{\pgfqpoint{2.532360in}{0.725185in}}%
\pgfpathlineto{\pgfqpoint{2.536661in}{0.728374in}}%
\pgfpathlineto{\pgfqpoint{2.541678in}{0.733664in}}%
\pgfpathlineto{\pgfqpoint{2.546338in}{0.740123in}}%
\pgfpathlineto{\pgfqpoint{2.550997in}{0.748419in}}%
\pgfpathlineto{\pgfqpoint{2.554939in}{0.757659in}}%
\pgfpathlineto{\pgfqpoint{2.558881in}{0.769805in}}%
\pgfpathlineto{\pgfqpoint{2.562824in}{0.785713in}}%
\pgfpathlineto{\pgfqpoint{2.566767in}{0.806367in}}%
\pgfpathlineto{\pgfqpoint{2.577161in}{0.864005in}}%
\pgfpathlineto{\pgfqpoint{2.581462in}{0.880146in}}%
\pgfpathlineto{\pgfqpoint{2.585764in}{0.892055in}}%
\pgfpathlineto{\pgfqpoint{2.590423in}{0.902038in}}%
\pgfpathlineto{\pgfqpoint{2.594725in}{0.908635in}}%
\pgfpathlineto{\pgfqpoint{2.599384in}{0.913819in}}%
\pgfpathlineto{\pgfqpoint{2.604044in}{0.917512in}}%
\pgfpathlineto{\pgfqpoint{2.610496in}{0.920829in}}%
\pgfpathlineto{\pgfqpoint{2.615514in}{0.922368in}}%
\pgfpathlineto{\pgfqpoint{2.621249in}{0.923066in}}%
\pgfpathlineto{\pgfqpoint{2.626625in}{0.922486in}}%
\pgfpathlineto{\pgfqpoint{2.658168in}{0.915274in}}%
\pgfpathlineto{\pgfqpoint{2.664978in}{0.912465in}}%
\pgfpathlineto{\pgfqpoint{2.672146in}{0.909726in}}%
\pgfpathlineto{\pgfqpoint{2.680390in}{0.907542in}}%
\pgfpathlineto{\pgfqpoint{2.686125in}{0.905811in}}%
\pgfpathlineto{\pgfqpoint{2.707631in}{0.898477in}}%
\pgfpathlineto{\pgfqpoint{2.748134in}{0.885449in}}%
\pgfpathlineto{\pgfqpoint{2.754586in}{0.883802in}}%
\pgfpathlineto{\pgfqpoint{2.759604in}{0.882142in}}%
\pgfpathlineto{\pgfqpoint{2.768207in}{0.878734in}}%
\pgfpathlineto{\pgfqpoint{2.789354in}{0.873988in}}%
\pgfpathlineto{\pgfqpoint{2.796164in}{0.874169in}}%
\pgfpathlineto{\pgfqpoint{2.800824in}{0.873935in}}%
\pgfpathlineto{\pgfqpoint{2.803692in}{0.872687in}}%
\pgfpathlineto{\pgfqpoint{2.806559in}{0.870241in}}%
\pgfpathlineto{\pgfqpoint{2.809785in}{0.866619in}}%
\pgfpathlineto{\pgfqpoint{2.809785in}{0.866619in}}%
\pgfusepath{stroke}%
\end{pgfscope}%
\begin{pgfscope}%
\pgfpathrectangle{\pgfqpoint{0.650000in}{0.420000in}}{\pgfqpoint{2.388110in}{1.994488in}}%
\pgfusepath{clip}%
\pgfsetrectcap%
\pgfsetroundjoin%
\pgfsetlinewidth{1.003750pt}%
\definecolor{currentstroke}{rgb}{0.454902,0.678431,0.819608}%
\pgfsetstrokecolor{currentstroke}%
\pgfsetdash{}{0pt}%
\pgfpathmoveto{\pgfqpoint{0.962760in}{2.419488in}}%
\pgfpathlineto{\pgfqpoint{0.970029in}{2.394063in}}%
\pgfpathlineto{\pgfqpoint{0.975530in}{2.370892in}}%
\pgfpathlineto{\pgfqpoint{0.980707in}{2.343972in}}%
\pgfpathlineto{\pgfqpoint{0.987503in}{2.302046in}}%
\pgfpathlineto{\pgfqpoint{0.994298in}{2.253720in}}%
\pgfpathlineto{\pgfqpoint{1.001741in}{2.192243in}}%
\pgfpathlineto{\pgfqpoint{1.010478in}{2.110214in}}%
\pgfpathlineto{\pgfqpoint{1.028598in}{1.936904in}}%
\pgfpathlineto{\pgfqpoint{1.034747in}{1.888884in}}%
\pgfpathlineto{\pgfqpoint{1.039600in}{1.858700in}}%
\pgfpathlineto{\pgfqpoint{1.043807in}{1.838711in}}%
\pgfpathlineto{\pgfqpoint{1.047367in}{1.826538in}}%
\pgfpathlineto{\pgfqpoint{1.050602in}{1.819133in}}%
\pgfpathlineto{\pgfqpoint{1.053191in}{1.815598in}}%
\pgfpathlineto{\pgfqpoint{1.055456in}{1.814234in}}%
\pgfpathlineto{\pgfqpoint{1.057398in}{1.814369in}}%
\pgfpathlineto{\pgfqpoint{1.059339in}{1.815708in}}%
\pgfpathlineto{\pgfqpoint{1.061604in}{1.818677in}}%
\pgfpathlineto{\pgfqpoint{1.064840in}{1.824920in}}%
\pgfpathlineto{\pgfqpoint{1.077137in}{1.851083in}}%
\pgfpathlineto{\pgfqpoint{1.079402in}{1.853343in}}%
\pgfpathlineto{\pgfqpoint{1.081020in}{1.853722in}}%
\pgfpathlineto{\pgfqpoint{1.082314in}{1.853062in}}%
\pgfpathlineto{\pgfqpoint{1.083932in}{1.850812in}}%
\pgfpathlineto{\pgfqpoint{1.085874in}{1.845742in}}%
\pgfpathlineto{\pgfqpoint{1.088462in}{1.835141in}}%
\pgfpathlineto{\pgfqpoint{1.094934in}{1.806577in}}%
\pgfpathlineto{\pgfqpoint{1.098493in}{1.796737in}}%
\pgfpathlineto{\pgfqpoint{1.102700in}{1.788507in}}%
\pgfpathlineto{\pgfqpoint{1.113055in}{1.769798in}}%
\pgfpathlineto{\pgfqpoint{1.119851in}{1.754131in}}%
\pgfpathlineto{\pgfqpoint{1.125352in}{1.738410in}}%
\pgfpathlineto{\pgfqpoint{1.131823in}{1.716339in}}%
\pgfpathlineto{\pgfqpoint{1.158682in}{1.619156in}}%
\pgfpathlineto{\pgfqpoint{1.165477in}{1.600195in}}%
\pgfpathlineto{\pgfqpoint{1.170331in}{1.589726in}}%
\pgfpathlineto{\pgfqpoint{1.175508in}{1.581194in}}%
\pgfpathlineto{\pgfqpoint{1.180039in}{1.575516in}}%
\pgfpathlineto{\pgfqpoint{1.183275in}{1.573161in}}%
\pgfpathlineto{\pgfqpoint{1.187158in}{1.571811in}}%
\pgfpathlineto{\pgfqpoint{1.190717in}{1.571638in}}%
\pgfpathlineto{\pgfqpoint{1.193953in}{1.572675in}}%
\pgfpathlineto{\pgfqpoint{1.197836in}{1.575245in}}%
\pgfpathlineto{\pgfqpoint{1.201719in}{1.579045in}}%
\pgfpathlineto{\pgfqpoint{1.207544in}{1.586689in}}%
\pgfpathlineto{\pgfqpoint{1.213045in}{1.595345in}}%
\pgfpathlineto{\pgfqpoint{1.218546in}{1.606348in}}%
\pgfpathlineto{\pgfqpoint{1.227606in}{1.627423in}}%
\pgfpathlineto{\pgfqpoint{1.247346in}{1.674082in}}%
\pgfpathlineto{\pgfqpoint{1.253817in}{1.687491in}}%
\pgfpathlineto{\pgfqpoint{1.260289in}{1.698007in}}%
\pgfpathlineto{\pgfqpoint{1.266114in}{1.705575in}}%
\pgfpathlineto{\pgfqpoint{1.270321in}{1.709394in}}%
\pgfpathlineto{\pgfqpoint{1.273880in}{1.711229in}}%
\pgfpathlineto{\pgfqpoint{1.277440in}{1.711881in}}%
\pgfpathlineto{\pgfqpoint{1.280675in}{1.711304in}}%
\pgfpathlineto{\pgfqpoint{1.285206in}{1.709265in}}%
\pgfpathlineto{\pgfqpoint{1.288442in}{1.706705in}}%
\pgfpathlineto{\pgfqpoint{1.291678in}{1.702714in}}%
\pgfpathlineto{\pgfqpoint{1.296208in}{1.695072in}}%
\pgfpathlineto{\pgfqpoint{1.303003in}{1.681323in}}%
\pgfpathlineto{\pgfqpoint{1.308181in}{1.668344in}}%
\pgfpathlineto{\pgfqpoint{1.317241in}{1.641659in}}%
\pgfpathlineto{\pgfqpoint{1.341834in}{1.568658in}}%
\pgfpathlineto{\pgfqpoint{1.350571in}{1.548219in}}%
\pgfpathlineto{\pgfqpoint{1.356072in}{1.537698in}}%
\pgfpathlineto{\pgfqpoint{1.361573in}{1.529844in}}%
\pgfpathlineto{\pgfqpoint{1.373222in}{1.516078in}}%
\pgfpathlineto{\pgfqpoint{1.379047in}{1.510679in}}%
\pgfpathlineto{\pgfqpoint{1.385195in}{1.506949in}}%
\pgfpathlineto{\pgfqpoint{1.397492in}{1.499611in}}%
\pgfpathlineto{\pgfqpoint{1.402993in}{1.494805in}}%
\pgfpathlineto{\pgfqpoint{1.411730in}{1.485758in}}%
\pgfpathlineto{\pgfqpoint{1.416260in}{1.479041in}}%
\pgfpathlineto{\pgfqpoint{1.424673in}{1.463764in}}%
\pgfpathlineto{\pgfqpoint{1.430821in}{1.450365in}}%
\pgfpathlineto{\pgfqpoint{1.438264in}{1.431265in}}%
\pgfpathlineto{\pgfqpoint{1.447325in}{1.404746in}}%
\pgfpathlineto{\pgfqpoint{1.486155in}{1.287352in}}%
\pgfpathlineto{\pgfqpoint{1.499746in}{1.253937in}}%
\pgfpathlineto{\pgfqpoint{1.509130in}{1.234895in}}%
\pgfpathlineto{\pgfqpoint{1.524016in}{1.207291in}}%
\pgfpathlineto{\pgfqpoint{1.537283in}{1.183719in}}%
\pgfpathlineto{\pgfqpoint{1.545049in}{1.168208in}}%
\pgfpathlineto{\pgfqpoint{1.556051in}{1.143465in}}%
\pgfpathlineto{\pgfqpoint{1.592293in}{1.055632in}}%
\pgfpathlineto{\pgfqpoint{1.601677in}{1.035444in}}%
\pgfpathlineto{\pgfqpoint{1.608473in}{1.022070in}}%
\pgfpathlineto{\pgfqpoint{1.618181in}{1.005941in}}%
\pgfpathlineto{\pgfqpoint{1.624652in}{0.996838in}}%
\pgfpathlineto{\pgfqpoint{1.633066in}{0.987678in}}%
\pgfpathlineto{\pgfqpoint{1.657011in}{0.964386in}}%
\pgfpathlineto{\pgfqpoint{1.668661in}{0.956745in}}%
\pgfpathlineto{\pgfqpoint{1.682899in}{0.948862in}}%
\pgfpathlineto{\pgfqpoint{1.702638in}{0.941688in}}%
\pgfpathlineto{\pgfqpoint{1.711698in}{0.938774in}}%
\pgfpathlineto{\pgfqpoint{1.715258in}{0.936657in}}%
\pgfpathlineto{\pgfqpoint{1.722053in}{0.930809in}}%
\pgfpathlineto{\pgfqpoint{1.735644in}{0.916821in}}%
\pgfpathlineto{\pgfqpoint{1.752147in}{0.897160in}}%
\pgfpathlineto{\pgfqpoint{1.765414in}{0.877618in}}%
\pgfpathlineto{\pgfqpoint{1.775769in}{0.862514in}}%
\pgfpathlineto{\pgfqpoint{1.786771in}{0.851678in}}%
\pgfpathlineto{\pgfqpoint{1.791625in}{0.848762in}}%
\pgfpathlineto{\pgfqpoint{1.797126in}{0.846718in}}%
\pgfpathlineto{\pgfqpoint{1.801009in}{0.846535in}}%
\pgfpathlineto{\pgfqpoint{1.804892in}{0.847511in}}%
\pgfpathlineto{\pgfqpoint{1.809422in}{0.849908in}}%
\pgfpathlineto{\pgfqpoint{1.813953in}{0.853526in}}%
\pgfpathlineto{\pgfqpoint{1.819777in}{0.859814in}}%
\pgfpathlineto{\pgfqpoint{1.824954in}{0.867222in}}%
\pgfpathlineto{\pgfqpoint{1.834015in}{0.882286in}}%
\pgfpathlineto{\pgfqpoint{1.850194in}{0.908732in}}%
\pgfpathlineto{\pgfqpoint{1.856989in}{0.918282in}}%
\pgfpathlineto{\pgfqpoint{1.863461in}{0.925996in}}%
\pgfpathlineto{\pgfqpoint{1.867668in}{0.928933in}}%
\pgfpathlineto{\pgfqpoint{1.873816in}{0.931858in}}%
\pgfpathlineto{\pgfqpoint{1.877375in}{0.932327in}}%
\pgfpathlineto{\pgfqpoint{1.883847in}{0.931636in}}%
\pgfpathlineto{\pgfqpoint{1.887083in}{0.930417in}}%
\pgfpathlineto{\pgfqpoint{1.891613in}{0.927187in}}%
\pgfpathlineto{\pgfqpoint{1.897114in}{0.922330in}}%
\pgfpathlineto{\pgfqpoint{1.903910in}{0.914120in}}%
\pgfpathlineto{\pgfqpoint{1.910705in}{0.904820in}}%
\pgfpathlineto{\pgfqpoint{1.918471in}{0.892309in}}%
\pgfpathlineto{\pgfqpoint{1.923972in}{0.881439in}}%
\pgfpathlineto{\pgfqpoint{1.930444in}{0.866232in}}%
\pgfpathlineto{\pgfqpoint{1.948241in}{0.823101in}}%
\pgfpathlineto{\pgfqpoint{1.955360in}{0.810553in}}%
\pgfpathlineto{\pgfqpoint{1.959243in}{0.805631in}}%
\pgfpathlineto{\pgfqpoint{1.963449in}{0.802260in}}%
\pgfpathlineto{\pgfqpoint{1.967656in}{0.800254in}}%
\pgfpathlineto{\pgfqpoint{1.971539in}{0.799555in}}%
\pgfpathlineto{\pgfqpoint{1.975099in}{0.799947in}}%
\pgfpathlineto{\pgfqpoint{1.978335in}{0.801569in}}%
\pgfpathlineto{\pgfqpoint{1.981894in}{0.804764in}}%
\pgfpathlineto{\pgfqpoint{1.986101in}{0.810291in}}%
\pgfpathlineto{\pgfqpoint{1.994514in}{0.823366in}}%
\pgfpathlineto{\pgfqpoint{2.011017in}{0.851060in}}%
\pgfpathlineto{\pgfqpoint{2.015870in}{0.857199in}}%
\pgfpathlineto{\pgfqpoint{2.020077in}{0.860800in}}%
\pgfpathlineto{\pgfqpoint{2.025254in}{0.863754in}}%
\pgfpathlineto{\pgfqpoint{2.029137in}{0.864906in}}%
\pgfpathlineto{\pgfqpoint{2.032373in}{0.864664in}}%
\pgfpathlineto{\pgfqpoint{2.035932in}{0.863122in}}%
\pgfpathlineto{\pgfqpoint{2.040139in}{0.859944in}}%
\pgfpathlineto{\pgfqpoint{2.048228in}{0.852263in}}%
\pgfpathlineto{\pgfqpoint{2.055993in}{0.844691in}}%
\pgfpathlineto{\pgfqpoint{2.070554in}{0.831598in}}%
\pgfpathlineto{\pgfqpoint{2.098704in}{0.802732in}}%
\pgfpathlineto{\pgfqpoint{2.104852in}{0.794133in}}%
\pgfpathlineto{\pgfqpoint{2.110676in}{0.783754in}}%
\pgfpathlineto{\pgfqpoint{2.117148in}{0.770688in}}%
\pgfpathlineto{\pgfqpoint{2.127502in}{0.746147in}}%
\pgfpathlineto{\pgfqpoint{2.154358in}{0.678049in}}%
\pgfpathlineto{\pgfqpoint{2.161476in}{0.663493in}}%
\pgfpathlineto{\pgfqpoint{2.166330in}{0.655553in}}%
\pgfpathlineto{\pgfqpoint{2.174742in}{0.644686in}}%
\pgfpathlineto{\pgfqpoint{2.181537in}{0.637664in}}%
\pgfpathlineto{\pgfqpoint{2.189626in}{0.630898in}}%
\pgfpathlineto{\pgfqpoint{2.195774in}{0.625483in}}%
\pgfpathlineto{\pgfqpoint{2.201275in}{0.620377in}}%
\pgfpathlineto{\pgfqpoint{2.206775in}{0.615820in}}%
\pgfpathlineto{\pgfqpoint{2.215512in}{0.605890in}}%
\pgfpathlineto{\pgfqpoint{2.228778in}{0.588395in}}%
\pgfpathlineto{\pgfqpoint{2.240103in}{0.573562in}}%
\pgfpathlineto{\pgfqpoint{2.248839in}{0.563933in}}%
\pgfpathlineto{\pgfqpoint{2.254987in}{0.558438in}}%
\pgfpathlineto{\pgfqpoint{2.259193in}{0.556355in}}%
\pgfpathlineto{\pgfqpoint{2.265341in}{0.554410in}}%
\pgfpathlineto{\pgfqpoint{2.268900in}{0.554439in}}%
\pgfpathlineto{\pgfqpoint{2.272459in}{0.555643in}}%
\pgfpathlineto{\pgfqpoint{2.276019in}{0.557927in}}%
\pgfpathlineto{\pgfqpoint{2.279254in}{0.561284in}}%
\pgfpathlineto{\pgfqpoint{2.283785in}{0.567810in}}%
\pgfpathlineto{\pgfqpoint{2.287667in}{0.575014in}}%
\pgfpathlineto{\pgfqpoint{2.291227in}{0.584172in}}%
\pgfpathlineto{\pgfqpoint{2.297051in}{0.603031in}}%
\pgfpathlineto{\pgfqpoint{2.302228in}{0.623119in}}%
\pgfpathlineto{\pgfqpoint{2.307082in}{0.646602in}}%
\pgfpathlineto{\pgfqpoint{2.313229in}{0.681990in}}%
\pgfpathlineto{\pgfqpoint{2.318730in}{0.720740in}}%
\pgfpathlineto{\pgfqpoint{2.325525in}{0.776302in}}%
\pgfpathlineto{\pgfqpoint{2.332644in}{0.844584in}}%
\pgfpathlineto{\pgfqpoint{2.341704in}{0.943305in}}%
\pgfpathlineto{\pgfqpoint{2.355618in}{1.109168in}}%
\pgfpathlineto{\pgfqpoint{2.368237in}{1.274895in}}%
\pgfpathlineto{\pgfqpoint{2.379886in}{1.441791in}}%
\pgfpathlineto{\pgfqpoint{2.389594in}{1.581381in}}%
\pgfpathlineto{\pgfqpoint{2.395743in}{1.651604in}}%
\pgfpathlineto{\pgfqpoint{2.402862in}{1.719697in}}%
\pgfpathlineto{\pgfqpoint{2.410953in}{1.785971in}}%
\pgfpathlineto{\pgfqpoint{2.418396in}{1.838291in}}%
\pgfpathlineto{\pgfqpoint{2.426810in}{1.889077in}}%
\pgfpathlineto{\pgfqpoint{2.435547in}{1.935994in}}%
\pgfpathlineto{\pgfqpoint{2.443961in}{1.974708in}}%
\pgfpathlineto{\pgfqpoint{2.451404in}{2.004283in}}%
\pgfpathlineto{\pgfqpoint{2.458523in}{2.028064in}}%
\pgfpathlineto{\pgfqpoint{2.467261in}{2.053478in}}%
\pgfpathlineto{\pgfqpoint{2.475351in}{2.073441in}}%
\pgfpathlineto{\pgfqpoint{2.489266in}{2.104112in}}%
\pgfpathlineto{\pgfqpoint{2.494120in}{2.112038in}}%
\pgfpathlineto{\pgfqpoint{2.499945in}{2.119234in}}%
\pgfpathlineto{\pgfqpoint{2.505770in}{2.124685in}}%
\pgfpathlineto{\pgfqpoint{2.511595in}{2.128725in}}%
\pgfpathlineto{\pgfqpoint{2.517743in}{2.131603in}}%
\pgfpathlineto{\pgfqpoint{2.525186in}{2.133592in}}%
\pgfpathlineto{\pgfqpoint{2.528746in}{2.133644in}}%
\pgfpathlineto{\pgfqpoint{2.534895in}{2.133280in}}%
\pgfpathlineto{\pgfqpoint{2.558842in}{2.137035in}}%
\pgfpathlineto{\pgfqpoint{2.579876in}{2.144304in}}%
\pgfpathlineto{\pgfqpoint{2.583759in}{2.147144in}}%
\pgfpathlineto{\pgfqpoint{2.592821in}{2.154442in}}%
\pgfpathlineto{\pgfqpoint{2.595086in}{2.154436in}}%
\pgfpathlineto{\pgfqpoint{2.597351in}{2.153149in}}%
\pgfpathlineto{\pgfqpoint{2.599940in}{2.150782in}}%
\pgfpathlineto{\pgfqpoint{2.599940in}{2.150782in}}%
\pgfusepath{stroke}%
\end{pgfscope}%
\begin{pgfscope}%
\pgfpathrectangle{\pgfqpoint{0.650000in}{0.420000in}}{\pgfqpoint{2.388110in}{1.994488in}}%
\pgfusepath{clip}%
\pgfsetrectcap%
\pgfsetroundjoin%
\pgfsetlinewidth{1.003750pt}%
\definecolor{currentstroke}{rgb}{0.878431,0.509804,0.078431}%
\pgfsetstrokecolor{currentstroke}%
\pgfsetdash{}{0pt}%
\pgfpathmoveto{\pgfqpoint{0.650000in}{1.783355in}}%
\pgfpathlineto{\pgfqpoint{0.658660in}{1.750821in}}%
\pgfpathlineto{\pgfqpoint{0.660317in}{1.747719in}}%
\pgfpathlineto{\pgfqpoint{0.661560in}{1.747054in}}%
\pgfpathlineto{\pgfqpoint{0.663218in}{1.747697in}}%
\pgfpathlineto{\pgfqpoint{0.665290in}{1.748511in}}%
\pgfpathlineto{\pgfqpoint{0.666947in}{1.747994in}}%
\pgfpathlineto{\pgfqpoint{0.669848in}{1.746755in}}%
\pgfpathlineto{\pgfqpoint{0.671091in}{1.747257in}}%
\pgfpathlineto{\pgfqpoint{0.672749in}{1.749386in}}%
\pgfpathlineto{\pgfqpoint{0.675235in}{1.755263in}}%
\pgfpathlineto{\pgfqpoint{0.685181in}{1.782365in}}%
\pgfpathlineto{\pgfqpoint{0.688082in}{1.786047in}}%
\pgfpathlineto{\pgfqpoint{0.690982in}{1.787771in}}%
\pgfpathlineto{\pgfqpoint{0.693469in}{1.788105in}}%
\pgfpathlineto{\pgfqpoint{0.695955in}{1.787382in}}%
\pgfpathlineto{\pgfqpoint{0.698856in}{1.785268in}}%
\pgfpathlineto{\pgfqpoint{0.703829in}{1.779891in}}%
\pgfpathlineto{\pgfqpoint{0.710459in}{1.771154in}}%
\pgfpathlineto{\pgfqpoint{0.726620in}{1.748639in}}%
\pgfpathlineto{\pgfqpoint{0.740295in}{1.733376in}}%
\pgfpathlineto{\pgfqpoint{0.745683in}{1.728639in}}%
\pgfpathlineto{\pgfqpoint{0.770132in}{1.710650in}}%
\pgfpathlineto{\pgfqpoint{0.811157in}{1.683079in}}%
\pgfpathlineto{\pgfqpoint{0.823174in}{1.675710in}}%
\pgfpathlineto{\pgfqpoint{0.831877in}{1.669970in}}%
\pgfpathlineto{\pgfqpoint{0.845966in}{1.662007in}}%
\pgfpathlineto{\pgfqpoint{0.860470in}{1.653249in}}%
\pgfpathlineto{\pgfqpoint{0.902738in}{1.628689in}}%
\pgfpathlineto{\pgfqpoint{0.911855in}{1.623990in}}%
\pgfpathlineto{\pgfqpoint{0.946664in}{1.603500in}}%
\pgfpathlineto{\pgfqpoint{0.963240in}{1.594687in}}%
\pgfpathlineto{\pgfqpoint{0.977330in}{1.586825in}}%
\pgfpathlineto{\pgfqpoint{0.997220in}{1.576270in}}%
\pgfpathlineto{\pgfqpoint{1.048191in}{1.547766in}}%
\pgfpathlineto{\pgfqpoint{1.083415in}{1.524283in}}%
\pgfpathlineto{\pgfqpoint{1.102477in}{1.507061in}}%
\pgfpathlineto{\pgfqpoint{1.111594in}{1.496680in}}%
\pgfpathlineto{\pgfqpoint{1.116566in}{1.489546in}}%
\pgfpathlineto{\pgfqpoint{1.121953in}{1.479611in}}%
\pgfpathlineto{\pgfqpoint{1.127341in}{1.467431in}}%
\pgfpathlineto{\pgfqpoint{1.139772in}{1.437386in}}%
\pgfpathlineto{\pgfqpoint{1.145160in}{1.428338in}}%
\pgfpathlineto{\pgfqpoint{1.153033in}{1.417394in}}%
\pgfpathlineto{\pgfqpoint{1.162149in}{1.407018in}}%
\pgfpathlineto{\pgfqpoint{1.168365in}{1.401177in}}%
\pgfpathlineto{\pgfqpoint{1.184941in}{1.388274in}}%
\pgfpathlineto{\pgfqpoint{1.215606in}{1.369050in}}%
\pgfpathlineto{\pgfqpoint{1.228038in}{1.362485in}}%
\pgfpathlineto{\pgfqpoint{1.246272in}{1.352073in}}%
\pgfpathlineto{\pgfqpoint{1.264505in}{1.342529in}}%
\pgfpathlineto{\pgfqpoint{1.274451in}{1.337477in}}%
\pgfpathlineto{\pgfqpoint{1.285225in}{1.333130in}}%
\pgfpathlineto{\pgfqpoint{1.295170in}{1.328766in}}%
\pgfpathlineto{\pgfqpoint{1.358987in}{1.295241in}}%
\pgfpathlineto{\pgfqpoint{1.363960in}{1.292946in}}%
\pgfpathlineto{\pgfqpoint{1.372248in}{1.289499in}}%
\pgfpathlineto{\pgfqpoint{1.406229in}{1.273232in}}%
\pgfpathlineto{\pgfqpoint{1.414931in}{1.269177in}}%
\pgfpathlineto{\pgfqpoint{1.436894in}{1.259794in}}%
\pgfpathlineto{\pgfqpoint{1.445182in}{1.255441in}}%
\pgfpathlineto{\pgfqpoint{1.455127in}{1.250279in}}%
\pgfpathlineto{\pgfqpoint{1.463415in}{1.246927in}}%
\pgfpathlineto{\pgfqpoint{1.476261in}{1.240243in}}%
\pgfpathlineto{\pgfqpoint{1.486207in}{1.235817in}}%
\pgfpathlineto{\pgfqpoint{1.494909in}{1.231247in}}%
\pgfpathlineto{\pgfqpoint{1.511899in}{1.224012in}}%
\pgfpathlineto{\pgfqpoint{1.527232in}{1.216159in}}%
\pgfpathlineto{\pgfqpoint{1.572816in}{1.196499in}}%
\pgfpathlineto{\pgfqpoint{1.598094in}{1.185590in}}%
\pgfpathlineto{\pgfqpoint{1.632903in}{1.168184in}}%
\pgfpathlineto{\pgfqpoint{1.668955in}{1.153522in}}%
\pgfpathlineto{\pgfqpoint{1.680559in}{1.148862in}}%
\pgfpathlineto{\pgfqpoint{1.725728in}{1.128659in}}%
\pgfpathlineto{\pgfqpoint{1.736502in}{1.123908in}}%
\pgfpathlineto{\pgfqpoint{1.751835in}{1.116948in}}%
\pgfpathlineto{\pgfqpoint{1.794932in}{1.099829in}}%
\pgfpathlineto{\pgfqpoint{1.813994in}{1.090817in}}%
\pgfpathlineto{\pgfqpoint{1.837615in}{1.082606in}}%
\pgfpathlineto{\pgfqpoint{1.844659in}{1.078926in}}%
\pgfpathlineto{\pgfqpoint{1.851290in}{1.076060in}}%
\pgfpathlineto{\pgfqpoint{1.867866in}{1.069170in}}%
\pgfpathlineto{\pgfqpoint{1.881540in}{1.062014in}}%
\pgfpathlineto{\pgfqpoint{1.891072in}{1.055442in}}%
\pgfpathlineto{\pgfqpoint{1.896045in}{1.053245in}}%
\pgfpathlineto{\pgfqpoint{1.919251in}{1.044973in}}%
\pgfpathlineto{\pgfqpoint{1.938313in}{1.036617in}}%
\pgfpathlineto{\pgfqpoint{1.958204in}{1.029223in}}%
\pgfpathlineto{\pgfqpoint{1.968564in}{1.025693in}}%
\pgfpathlineto{\pgfqpoint{1.976437in}{1.022172in}}%
\pgfpathlineto{\pgfqpoint{1.995085in}{1.011804in}}%
\pgfpathlineto{\pgfqpoint{2.005859in}{1.006642in}}%
\pgfpathlineto{\pgfqpoint{2.013318in}{1.003178in}}%
\pgfpathlineto{\pgfqpoint{2.032381in}{0.995724in}}%
\pgfpathlineto{\pgfqpoint{2.084594in}{0.974401in}}%
\pgfpathlineto{\pgfqpoint{2.092882in}{0.971166in}}%
\pgfpathlineto{\pgfqpoint{2.135980in}{0.952783in}}%
\pgfpathlineto{\pgfqpoint{2.144682in}{0.950465in}}%
\pgfpathlineto{\pgfqpoint{2.152141in}{0.946865in}}%
\pgfpathlineto{\pgfqpoint{2.161672in}{0.941618in}}%
\pgfpathlineto{\pgfqpoint{2.167888in}{0.939777in}}%
\pgfpathlineto{\pgfqpoint{2.174518in}{0.936459in}}%
\pgfpathlineto{\pgfqpoint{2.183220in}{0.932887in}}%
\pgfpathlineto{\pgfqpoint{2.197310in}{0.927040in}}%
\pgfpathlineto{\pgfqpoint{2.215958in}{0.920978in}}%
\pgfpathlineto{\pgfqpoint{2.225074in}{0.916583in}}%
\pgfpathlineto{\pgfqpoint{2.231290in}{0.915899in}}%
\pgfpathlineto{\pgfqpoint{2.235434in}{0.913676in}}%
\pgfpathlineto{\pgfqpoint{2.242064in}{0.910009in}}%
\pgfpathlineto{\pgfqpoint{2.253253in}{0.905037in}}%
\pgfpathlineto{\pgfqpoint{2.259055in}{0.902723in}}%
\pgfpathlineto{\pgfqpoint{2.266514in}{0.899928in}}%
\pgfpathlineto{\pgfqpoint{2.280189in}{0.893651in}}%
\pgfpathlineto{\pgfqpoint{2.283918in}{0.891652in}}%
\pgfpathlineto{\pgfqpoint{2.304224in}{0.884742in}}%
\pgfpathlineto{\pgfqpoint{2.309611in}{0.882256in}}%
\pgfpathlineto{\pgfqpoint{2.318313in}{0.877486in}}%
\pgfpathlineto{\pgfqpoint{2.324944in}{0.875378in}}%
\pgfpathlineto{\pgfqpoint{2.329502in}{0.873096in}}%
\pgfpathlineto{\pgfqpoint{2.336132in}{0.870941in}}%
\pgfpathlineto{\pgfqpoint{2.342348in}{0.868844in}}%
\pgfpathlineto{\pgfqpoint{2.346907in}{0.866853in}}%
\pgfpathlineto{\pgfqpoint{2.358924in}{0.860201in}}%
\pgfpathlineto{\pgfqpoint{2.362654in}{0.859429in}}%
\pgfpathlineto{\pgfqpoint{2.366383in}{0.858272in}}%
\pgfpathlineto{\pgfqpoint{2.375086in}{0.854218in}}%
\pgfpathlineto{\pgfqpoint{2.382130in}{0.852921in}}%
\pgfpathlineto{\pgfqpoint{2.394148in}{0.848827in}}%
\pgfpathlineto{\pgfqpoint{2.399121in}{0.846873in}}%
\pgfpathlineto{\pgfqpoint{2.412381in}{0.840794in}}%
\pgfpathlineto{\pgfqpoint{2.424399in}{0.837128in}}%
\pgfpathlineto{\pgfqpoint{2.427300in}{0.836030in}}%
\pgfpathlineto{\pgfqpoint{2.431029in}{0.833055in}}%
\pgfpathlineto{\pgfqpoint{2.434759in}{0.830690in}}%
\pgfpathlineto{\pgfqpoint{2.448849in}{0.825354in}}%
\pgfpathlineto{\pgfqpoint{2.455479in}{0.825012in}}%
\pgfpathlineto{\pgfqpoint{2.463352in}{0.821493in}}%
\pgfpathlineto{\pgfqpoint{2.467082in}{0.820528in}}%
\pgfpathlineto{\pgfqpoint{2.486558in}{0.812832in}}%
\pgfpathlineto{\pgfqpoint{2.490702in}{0.812325in}}%
\pgfpathlineto{\pgfqpoint{2.494846in}{0.811882in}}%
\pgfpathlineto{\pgfqpoint{2.503963in}{0.809340in}}%
\pgfpathlineto{\pgfqpoint{2.513080in}{0.809134in}}%
\pgfpathlineto{\pgfqpoint{2.517224in}{0.809390in}}%
\pgfpathlineto{\pgfqpoint{2.521782in}{0.809279in}}%
\pgfpathlineto{\pgfqpoint{2.526755in}{0.811017in}}%
\pgfpathlineto{\pgfqpoint{2.534628in}{0.813290in}}%
\pgfpathlineto{\pgfqpoint{2.538358in}{0.814857in}}%
\pgfpathlineto{\pgfqpoint{2.543745in}{0.819372in}}%
\pgfpathlineto{\pgfqpoint{2.547060in}{0.822627in}}%
\pgfpathlineto{\pgfqpoint{2.555762in}{0.833447in}}%
\pgfpathlineto{\pgfqpoint{2.559077in}{0.836717in}}%
\pgfpathlineto{\pgfqpoint{2.564878in}{0.843071in}}%
\pgfpathlineto{\pgfqpoint{2.567364in}{0.843930in}}%
\pgfpathlineto{\pgfqpoint{2.572751in}{0.845345in}}%
\pgfpathlineto{\pgfqpoint{2.577723in}{0.847040in}}%
\pgfpathlineto{\pgfqpoint{2.584767in}{0.847379in}}%
\pgfpathlineto{\pgfqpoint{2.587668in}{0.846540in}}%
\pgfpathlineto{\pgfqpoint{2.592226in}{0.844825in}}%
\pgfpathlineto{\pgfqpoint{2.600513in}{0.843838in}}%
\pgfpathlineto{\pgfqpoint{2.606728in}{0.842451in}}%
\pgfpathlineto{\pgfqpoint{2.626618in}{0.836958in}}%
\pgfpathlineto{\pgfqpoint{2.631176in}{0.834522in}}%
\pgfpathlineto{\pgfqpoint{2.635734in}{0.832495in}}%
\pgfpathlineto{\pgfqpoint{2.642778in}{0.828365in}}%
\pgfpathlineto{\pgfqpoint{2.651894in}{0.825671in}}%
\pgfpathlineto{\pgfqpoint{2.661010in}{0.820758in}}%
\pgfpathlineto{\pgfqpoint{2.668054in}{0.819938in}}%
\pgfpathlineto{\pgfqpoint{2.670955in}{0.817231in}}%
\pgfpathlineto{\pgfqpoint{2.675098in}{0.813292in}}%
\pgfpathlineto{\pgfqpoint{2.677170in}{0.812650in}}%
\pgfpathlineto{\pgfqpoint{2.679656in}{0.813051in}}%
\pgfpathlineto{\pgfqpoint{2.683800in}{0.813995in}}%
\pgfpathlineto{\pgfqpoint{2.686286in}{0.813228in}}%
\pgfpathlineto{\pgfqpoint{2.689601in}{0.810652in}}%
\pgfpathlineto{\pgfqpoint{2.694573in}{0.806685in}}%
\pgfpathlineto{\pgfqpoint{2.697888in}{0.805463in}}%
\pgfpathlineto{\pgfqpoint{2.706175in}{0.803444in}}%
\pgfpathlineto{\pgfqpoint{2.714463in}{0.799452in}}%
\pgfpathlineto{\pgfqpoint{2.723993in}{0.797514in}}%
\pgfpathlineto{\pgfqpoint{2.731451in}{0.794863in}}%
\pgfpathlineto{\pgfqpoint{2.735181in}{0.791872in}}%
\pgfpathlineto{\pgfqpoint{2.738910in}{0.789294in}}%
\pgfpathlineto{\pgfqpoint{2.742225in}{0.788478in}}%
\pgfpathlineto{\pgfqpoint{2.748440in}{0.787290in}}%
\pgfpathlineto{\pgfqpoint{2.754656in}{0.784306in}}%
\pgfpathlineto{\pgfqpoint{2.765015in}{0.780426in}}%
\pgfpathlineto{\pgfqpoint{2.771645in}{0.778527in}}%
\pgfpathlineto{\pgfqpoint{2.787805in}{0.776216in}}%
\pgfpathlineto{\pgfqpoint{2.793191in}{0.773702in}}%
\pgfpathlineto{\pgfqpoint{2.808523in}{0.767024in}}%
\pgfpathlineto{\pgfqpoint{2.821782in}{0.762543in}}%
\pgfpathlineto{\pgfqpoint{2.832556in}{0.757206in}}%
\pgfpathlineto{\pgfqpoint{2.847473in}{0.752622in}}%
\pgfpathlineto{\pgfqpoint{2.850788in}{0.750695in}}%
\pgfpathlineto{\pgfqpoint{2.854517in}{0.748554in}}%
\pgfpathlineto{\pgfqpoint{2.862390in}{0.746973in}}%
\pgfpathlineto{\pgfqpoint{2.865290in}{0.746233in}}%
\pgfpathlineto{\pgfqpoint{2.877721in}{0.741281in}}%
\pgfpathlineto{\pgfqpoint{2.881036in}{0.740190in}}%
\pgfpathlineto{\pgfqpoint{2.884351in}{0.739557in}}%
\pgfpathlineto{\pgfqpoint{2.886837in}{0.738584in}}%
\pgfpathlineto{\pgfqpoint{2.891395in}{0.736004in}}%
\pgfpathlineto{\pgfqpoint{2.894710in}{0.734676in}}%
\pgfpathlineto{\pgfqpoint{2.898025in}{0.733214in}}%
\pgfpathlineto{\pgfqpoint{2.903826in}{0.732616in}}%
\pgfpathlineto{\pgfqpoint{2.906312in}{0.732282in}}%
\pgfpathlineto{\pgfqpoint{2.909213in}{0.732214in}}%
\pgfpathlineto{\pgfqpoint{2.912113in}{0.730501in}}%
\pgfpathlineto{\pgfqpoint{2.914599in}{0.729723in}}%
\pgfpathlineto{\pgfqpoint{2.917086in}{0.729144in}}%
\pgfpathlineto{\pgfqpoint{2.919572in}{0.727092in}}%
\pgfpathlineto{\pgfqpoint{2.922058in}{0.725497in}}%
\pgfpathlineto{\pgfqpoint{2.926616in}{0.723935in}}%
\pgfpathlineto{\pgfqpoint{2.929517in}{0.722245in}}%
\pgfpathlineto{\pgfqpoint{2.935732in}{0.720734in}}%
\pgfpathlineto{\pgfqpoint{2.938218in}{0.719456in}}%
\pgfpathlineto{\pgfqpoint{2.941533in}{0.718447in}}%
\pgfpathlineto{\pgfqpoint{2.943605in}{0.715992in}}%
\pgfpathlineto{\pgfqpoint{2.946505in}{0.712553in}}%
\pgfpathlineto{\pgfqpoint{2.948577in}{0.711946in}}%
\pgfpathlineto{\pgfqpoint{2.951063in}{0.711311in}}%
\pgfpathlineto{\pgfqpoint{2.954793in}{0.709668in}}%
\pgfpathlineto{\pgfqpoint{2.959765in}{0.709815in}}%
\pgfpathlineto{\pgfqpoint{2.963909in}{0.707443in}}%
\pgfpathlineto{\pgfqpoint{2.968466in}{0.707737in}}%
\pgfpathlineto{\pgfqpoint{2.972610in}{0.705635in}}%
\pgfpathlineto{\pgfqpoint{2.976339in}{0.705877in}}%
\pgfpathlineto{\pgfqpoint{2.977997in}{0.704362in}}%
\pgfpathlineto{\pgfqpoint{2.980897in}{0.701381in}}%
\pgfpathlineto{\pgfqpoint{2.982555in}{0.701210in}}%
\pgfpathlineto{\pgfqpoint{2.985870in}{0.701539in}}%
\pgfpathlineto{\pgfqpoint{2.990842in}{0.699788in}}%
\pgfpathlineto{\pgfqpoint{2.994571in}{0.700936in}}%
\pgfpathlineto{\pgfqpoint{2.996229in}{0.699627in}}%
\pgfpathlineto{\pgfqpoint{2.999129in}{0.697068in}}%
\pgfpathlineto{\pgfqpoint{3.000787in}{0.697231in}}%
\pgfpathlineto{\pgfqpoint{3.004102in}{0.698192in}}%
\pgfpathlineto{\pgfqpoint{3.006173in}{0.696891in}}%
\pgfpathlineto{\pgfqpoint{3.008245in}{0.695746in}}%
\pgfpathlineto{\pgfqpoint{3.009903in}{0.696073in}}%
\pgfpathlineto{\pgfqpoint{3.012803in}{0.697054in}}%
\pgfpathlineto{\pgfqpoint{3.014461in}{0.696002in}}%
\pgfpathlineto{\pgfqpoint{3.018190in}{0.692650in}}%
\pgfpathlineto{\pgfqpoint{3.019847in}{0.692977in}}%
\pgfpathlineto{\pgfqpoint{3.022316in}{0.693785in}}%
\pgfpathlineto{\pgfqpoint{3.023518in}{0.693088in}}%
\pgfpathlineto{\pgfqpoint{3.025001in}{0.690759in}}%
\pgfpathlineto{\pgfqpoint{3.026849in}{0.685870in}}%
\pgfpathlineto{\pgfqpoint{3.026849in}{0.685870in}}%
\pgfusepath{stroke}%
\end{pgfscope}%
\begin{pgfscope}%
\pgfpathrectangle{\pgfqpoint{0.650000in}{0.420000in}}{\pgfqpoint{2.388110in}{1.994488in}}%
\pgfusepath{clip}%
\pgfsetrectcap%
\pgfsetroundjoin%
\pgfsetlinewidth{1.003750pt}%
\definecolor{currentstroke}{rgb}{0.501961,0.450980,0.674510}%
\pgfsetstrokecolor{currentstroke}%
\pgfsetdash{}{0pt}%
\pgfpathmoveto{\pgfqpoint{0.650000in}{1.445952in}}%
\pgfpathlineto{\pgfqpoint{0.659117in}{1.412269in}}%
\pgfpathlineto{\pgfqpoint{0.660782in}{1.409726in}}%
\pgfpathlineto{\pgfqpoint{0.662031in}{1.409409in}}%
\pgfpathlineto{\pgfqpoint{0.666195in}{1.410257in}}%
\pgfpathlineto{\pgfqpoint{0.670359in}{1.408413in}}%
\pgfpathlineto{\pgfqpoint{0.671608in}{1.409243in}}%
\pgfpathlineto{\pgfqpoint{0.673273in}{1.411869in}}%
\pgfpathlineto{\pgfqpoint{0.675771in}{1.418344in}}%
\pgfpathlineto{\pgfqpoint{0.684515in}{1.443156in}}%
\pgfpathlineto{\pgfqpoint{0.687429in}{1.447872in}}%
\pgfpathlineto{\pgfqpoint{0.689928in}{1.449985in}}%
\pgfpathlineto{\pgfqpoint{0.692009in}{1.450465in}}%
\pgfpathlineto{\pgfqpoint{0.694507in}{1.449762in}}%
\pgfpathlineto{\pgfqpoint{0.697422in}{1.447648in}}%
\pgfpathlineto{\pgfqpoint{0.701586in}{1.443069in}}%
\pgfpathlineto{\pgfqpoint{0.707831in}{1.434187in}}%
\pgfpathlineto{\pgfqpoint{0.722820in}{1.412886in}}%
\pgfpathlineto{\pgfqpoint{0.730731in}{1.403168in}}%
\pgfpathlineto{\pgfqpoint{0.742389in}{1.390749in}}%
\pgfpathlineto{\pgfqpoint{0.749467in}{1.385095in}}%
\pgfpathlineto{\pgfqpoint{0.760292in}{1.376561in}}%
\pgfpathlineto{\pgfqpoint{0.779445in}{1.361154in}}%
\pgfpathlineto{\pgfqpoint{0.801095in}{1.343560in}}%
\pgfpathlineto{\pgfqpoint{0.831906in}{1.314888in}}%
\pgfpathlineto{\pgfqpoint{0.845646in}{1.302324in}}%
\pgfpathlineto{\pgfqpoint{0.856055in}{1.293944in}}%
\pgfpathlineto{\pgfqpoint{0.907267in}{1.258278in}}%
\pgfpathlineto{\pgfqpoint{0.921423in}{1.249390in}}%
\pgfpathlineto{\pgfqpoint{0.940575in}{1.236329in}}%
\pgfpathlineto{\pgfqpoint{0.990538in}{1.196241in}}%
\pgfpathlineto{\pgfqpoint{1.003862in}{1.187220in}}%
\pgfpathlineto{\pgfqpoint{1.021765in}{1.176634in}}%
\pgfpathlineto{\pgfqpoint{1.032591in}{1.171607in}}%
\pgfpathlineto{\pgfqpoint{1.049245in}{1.166023in}}%
\pgfpathlineto{\pgfqpoint{1.078390in}{1.156450in}}%
\pgfpathlineto{\pgfqpoint{1.100041in}{1.147088in}}%
\pgfpathlineto{\pgfqpoint{1.108784in}{1.141673in}}%
\pgfpathlineto{\pgfqpoint{1.115446in}{1.135956in}}%
\pgfpathlineto{\pgfqpoint{1.122524in}{1.128637in}}%
\pgfpathlineto{\pgfqpoint{1.127104in}{1.122130in}}%
\pgfpathlineto{\pgfqpoint{1.132517in}{1.112106in}}%
\pgfpathlineto{\pgfqpoint{1.138346in}{1.101855in}}%
\pgfpathlineto{\pgfqpoint{1.143342in}{1.095452in}}%
\pgfpathlineto{\pgfqpoint{1.148755in}{1.090380in}}%
\pgfpathlineto{\pgfqpoint{1.155416in}{1.085693in}}%
\pgfpathlineto{\pgfqpoint{1.160829in}{1.083164in}}%
\pgfpathlineto{\pgfqpoint{1.166241in}{1.081897in}}%
\pgfpathlineto{\pgfqpoint{1.177066in}{1.080943in}}%
\pgfpathlineto{\pgfqpoint{1.186643in}{1.080992in}}%
\pgfpathlineto{\pgfqpoint{1.193305in}{1.082181in}}%
\pgfpathlineto{\pgfqpoint{1.203714in}{1.084265in}}%
\pgfpathlineto{\pgfqpoint{1.211624in}{1.085838in}}%
\pgfpathlineto{\pgfqpoint{1.221201in}{1.087964in}}%
\pgfpathlineto{\pgfqpoint{1.229112in}{1.089984in}}%
\pgfpathlineto{\pgfqpoint{1.256591in}{1.097544in}}%
\pgfpathlineto{\pgfqpoint{1.263669in}{1.098541in}}%
\pgfpathlineto{\pgfqpoint{1.281573in}{1.099425in}}%
\pgfpathlineto{\pgfqpoint{1.291149in}{1.099558in}}%
\pgfpathlineto{\pgfqpoint{1.300725in}{1.100644in}}%
\pgfpathlineto{\pgfqpoint{1.310718in}{1.102903in}}%
\pgfpathlineto{\pgfqpoint{1.322376in}{1.105333in}}%
\pgfpathlineto{\pgfqpoint{1.338614in}{1.107662in}}%
\pgfpathlineto{\pgfqpoint{1.359848in}{1.109542in}}%
\pgfpathlineto{\pgfqpoint{1.372755in}{1.109671in}}%
\pgfpathlineto{\pgfqpoint{1.380666in}{1.108875in}}%
\pgfpathlineto{\pgfqpoint{1.392741in}{1.106937in}}%
\pgfpathlineto{\pgfqpoint{1.412726in}{1.102051in}}%
\pgfpathlineto{\pgfqpoint{1.431046in}{1.094611in}}%
\pgfpathlineto{\pgfqpoint{1.439789in}{1.091521in}}%
\pgfpathlineto{\pgfqpoint{1.449782in}{1.087238in}}%
\pgfpathlineto{\pgfqpoint{1.465187in}{1.081363in}}%
\pgfpathlineto{\pgfqpoint{1.495998in}{1.075566in}}%
\pgfpathlineto{\pgfqpoint{1.509738in}{1.073124in}}%
\pgfpathlineto{\pgfqpoint{1.521396in}{1.070786in}}%
\pgfpathlineto{\pgfqpoint{1.538466in}{1.068387in}}%
\pgfpathlineto{\pgfqpoint{1.549292in}{1.065644in}}%
\pgfpathlineto{\pgfqpoint{1.598006in}{1.053958in}}%
\pgfpathlineto{\pgfqpoint{1.609247in}{1.051535in}}%
\pgfpathlineto{\pgfqpoint{1.620073in}{1.049689in}}%
\pgfpathlineto{\pgfqpoint{1.647552in}{1.046302in}}%
\pgfpathlineto{\pgfqpoint{1.687523in}{1.043653in}}%
\pgfpathlineto{\pgfqpoint{1.714170in}{1.046436in}}%
\pgfpathlineto{\pgfqpoint{1.732073in}{1.049777in}}%
\pgfpathlineto{\pgfqpoint{1.751642in}{1.053831in}}%
\pgfpathlineto{\pgfqpoint{1.763717in}{1.056602in}}%
\pgfpathlineto{\pgfqpoint{1.802438in}{1.066204in}}%
\pgfpathlineto{\pgfqpoint{1.813680in}{1.070171in}}%
\pgfpathlineto{\pgfqpoint{1.828669in}{1.074739in}}%
\pgfpathlineto{\pgfqpoint{1.915271in}{1.107601in}}%
\pgfpathlineto{\pgfqpoint{1.931509in}{1.114813in}}%
\pgfpathlineto{\pgfqpoint{1.964402in}{1.130828in}}%
\pgfpathlineto{\pgfqpoint{1.995212in}{1.142756in}}%
\pgfpathlineto{\pgfqpoint{2.017695in}{1.148475in}}%
\pgfpathlineto{\pgfqpoint{2.028104in}{1.150019in}}%
\pgfpathlineto{\pgfqpoint{2.036848in}{1.150333in}}%
\pgfpathlineto{\pgfqpoint{2.055584in}{1.149484in}}%
\pgfpathlineto{\pgfqpoint{2.067242in}{1.148548in}}%
\pgfpathlineto{\pgfqpoint{2.077235in}{1.147344in}}%
\pgfpathlineto{\pgfqpoint{2.088477in}{1.146043in}}%
\pgfpathlineto{\pgfqpoint{2.115956in}{1.145769in}}%
\pgfpathlineto{\pgfqpoint{2.125949in}{1.146596in}}%
\pgfpathlineto{\pgfqpoint{2.135525in}{1.147436in}}%
\pgfpathlineto{\pgfqpoint{2.165087in}{1.146317in}}%
\pgfpathlineto{\pgfqpoint{2.169250in}{1.146197in}}%
\pgfpathlineto{\pgfqpoint{2.174247in}{1.146544in}}%
\pgfpathlineto{\pgfqpoint{2.192982in}{1.141898in}}%
\pgfpathlineto{\pgfqpoint{2.200893in}{1.138188in}}%
\pgfpathlineto{\pgfqpoint{2.208388in}{1.134393in}}%
\pgfpathlineto{\pgfqpoint{2.215049in}{1.131525in}}%
\pgfpathlineto{\pgfqpoint{2.221295in}{1.128680in}}%
\pgfpathlineto{\pgfqpoint{2.225458in}{1.126390in}}%
\pgfpathlineto{\pgfqpoint{2.232953in}{1.121471in}}%
\pgfpathlineto{\pgfqpoint{2.238782in}{1.118545in}}%
\pgfpathlineto{\pgfqpoint{2.245443in}{1.113422in}}%
\pgfpathlineto{\pgfqpoint{2.251272in}{1.109431in}}%
\pgfpathlineto{\pgfqpoint{2.267511in}{1.098853in}}%
\pgfpathlineto{\pgfqpoint{2.275838in}{1.092482in}}%
\pgfpathlineto{\pgfqpoint{2.282499in}{1.087551in}}%
\pgfpathlineto{\pgfqpoint{2.286663in}{1.084570in}}%
\pgfpathlineto{\pgfqpoint{2.293741in}{1.078673in}}%
\pgfpathlineto{\pgfqpoint{2.300403in}{1.073513in}}%
\pgfpathlineto{\pgfqpoint{2.306232in}{1.069281in}}%
\pgfpathlineto{\pgfqpoint{2.320388in}{1.059639in}}%
\pgfpathlineto{\pgfqpoint{2.326634in}{1.055066in}}%
\pgfpathlineto{\pgfqpoint{2.332879in}{1.051092in}}%
\pgfpathlineto{\pgfqpoint{2.337875in}{1.047399in}}%
\pgfpathlineto{\pgfqpoint{2.344953in}{1.043418in}}%
\pgfpathlineto{\pgfqpoint{2.350783in}{1.040518in}}%
\pgfpathlineto{\pgfqpoint{2.354946in}{1.038170in}}%
\pgfpathlineto{\pgfqpoint{2.360775in}{1.034667in}}%
\pgfpathlineto{\pgfqpoint{2.370352in}{1.030516in}}%
\pgfpathlineto{\pgfqpoint{2.374515in}{1.029382in}}%
\pgfpathlineto{\pgfqpoint{2.385341in}{1.024627in}}%
\pgfpathlineto{\pgfqpoint{2.392002in}{1.022828in}}%
\pgfpathlineto{\pgfqpoint{2.397831in}{1.021424in}}%
\pgfpathlineto{\pgfqpoint{2.414902in}{1.018663in}}%
\pgfpathlineto{\pgfqpoint{2.428642in}{1.016042in}}%
\pgfpathlineto{\pgfqpoint{2.436969in}{1.016358in}}%
\pgfpathlineto{\pgfqpoint{2.445296in}{1.013984in}}%
\pgfpathlineto{\pgfqpoint{2.452791in}{1.012625in}}%
\pgfpathlineto{\pgfqpoint{2.456538in}{1.012135in}}%
\pgfpathlineto{\pgfqpoint{2.459452in}{1.013087in}}%
\pgfpathlineto{\pgfqpoint{2.462783in}{1.013969in}}%
\pgfpathlineto{\pgfqpoint{2.466114in}{1.013437in}}%
\pgfpathlineto{\pgfqpoint{2.469862in}{1.012989in}}%
\pgfpathlineto{\pgfqpoint{2.481103in}{1.015054in}}%
\pgfpathlineto{\pgfqpoint{2.493594in}{1.018553in}}%
\pgfpathlineto{\pgfqpoint{2.499007in}{1.021144in}}%
\pgfpathlineto{\pgfqpoint{2.505252in}{1.024816in}}%
\pgfpathlineto{\pgfqpoint{2.510248in}{1.027446in}}%
\pgfpathlineto{\pgfqpoint{2.517743in}{1.033349in}}%
\pgfpathlineto{\pgfqpoint{2.522739in}{1.038374in}}%
\pgfpathlineto{\pgfqpoint{2.534397in}{1.051688in}}%
\pgfpathlineto{\pgfqpoint{2.543141in}{1.062626in}}%
\pgfpathlineto{\pgfqpoint{2.547304in}{1.068849in}}%
\pgfpathlineto{\pgfqpoint{2.556464in}{1.086724in}}%
\pgfpathlineto{\pgfqpoint{2.567705in}{1.112038in}}%
\pgfpathlineto{\pgfqpoint{2.575615in}{1.127052in}}%
\pgfpathlineto{\pgfqpoint{2.587272in}{1.142300in}}%
\pgfpathlineto{\pgfqpoint{2.597264in}{1.152594in}}%
\pgfpathlineto{\pgfqpoint{2.603509in}{1.157843in}}%
\pgfpathlineto{\pgfqpoint{2.608089in}{1.161403in}}%
\pgfpathlineto{\pgfqpoint{2.619746in}{1.167804in}}%
\pgfpathlineto{\pgfqpoint{2.634317in}{1.175309in}}%
\pgfpathlineto{\pgfqpoint{2.638480in}{1.175923in}}%
\pgfpathlineto{\pgfqpoint{2.646391in}{1.176633in}}%
\pgfpathlineto{\pgfqpoint{2.651803in}{1.176219in}}%
\pgfpathlineto{\pgfqpoint{2.662211in}{1.177211in}}%
\pgfpathlineto{\pgfqpoint{2.671370in}{1.174899in}}%
\pgfpathlineto{\pgfqpoint{2.677199in}{1.175613in}}%
\pgfpathlineto{\pgfqpoint{2.679697in}{1.174295in}}%
\pgfpathlineto{\pgfqpoint{2.686358in}{1.169600in}}%
\pgfpathlineto{\pgfqpoint{2.689272in}{1.169519in}}%
\pgfpathlineto{\pgfqpoint{2.693852in}{1.169711in}}%
\pgfpathlineto{\pgfqpoint{2.696350in}{1.168618in}}%
\pgfpathlineto{\pgfqpoint{2.699680in}{1.165810in}}%
\pgfpathlineto{\pgfqpoint{2.705925in}{1.160320in}}%
\pgfpathlineto{\pgfqpoint{2.710921in}{1.157765in}}%
\pgfpathlineto{\pgfqpoint{2.716749in}{1.154637in}}%
\pgfpathlineto{\pgfqpoint{2.725076in}{1.149582in}}%
\pgfpathlineto{\pgfqpoint{2.731321in}{1.146260in}}%
\pgfpathlineto{\pgfqpoint{2.739231in}{1.140628in}}%
\pgfpathlineto{\pgfqpoint{2.743394in}{1.136493in}}%
\pgfpathlineto{\pgfqpoint{2.748807in}{1.130797in}}%
\pgfpathlineto{\pgfqpoint{2.752970in}{1.128443in}}%
\pgfpathlineto{\pgfqpoint{2.758382in}{1.125308in}}%
\pgfpathlineto{\pgfqpoint{2.762962in}{1.120942in}}%
\pgfpathlineto{\pgfqpoint{2.772121in}{1.112481in}}%
\pgfpathlineto{\pgfqpoint{2.779198in}{1.107021in}}%
\pgfpathlineto{\pgfqpoint{2.786692in}{1.102575in}}%
\pgfpathlineto{\pgfqpoint{2.795851in}{1.096849in}}%
\pgfpathlineto{\pgfqpoint{2.803345in}{1.090927in}}%
\pgfpathlineto{\pgfqpoint{2.826243in}{1.071339in}}%
\pgfpathlineto{\pgfqpoint{2.833737in}{1.065467in}}%
\pgfpathlineto{\pgfqpoint{2.844145in}{1.055023in}}%
\pgfpathlineto{\pgfqpoint{2.854137in}{1.046574in}}%
\pgfpathlineto{\pgfqpoint{2.871206in}{1.028481in}}%
\pgfpathlineto{\pgfqpoint{2.877035in}{1.022968in}}%
\pgfpathlineto{\pgfqpoint{2.884945in}{1.013976in}}%
\pgfpathlineto{\pgfqpoint{2.889108in}{1.009936in}}%
\pgfpathlineto{\pgfqpoint{2.902014in}{0.996508in}}%
\pgfpathlineto{\pgfqpoint{2.906177in}{0.991315in}}%
\pgfpathlineto{\pgfqpoint{2.909508in}{0.987504in}}%
\pgfpathlineto{\pgfqpoint{2.921998in}{0.976915in}}%
\pgfpathlineto{\pgfqpoint{2.925328in}{0.973981in}}%
\pgfpathlineto{\pgfqpoint{2.928243in}{0.971652in}}%
\pgfpathlineto{\pgfqpoint{2.934488in}{0.965165in}}%
\pgfpathlineto{\pgfqpoint{2.937402in}{0.962584in}}%
\pgfpathlineto{\pgfqpoint{2.940732in}{0.959671in}}%
\pgfpathlineto{\pgfqpoint{2.945728in}{0.957182in}}%
\pgfpathlineto{\pgfqpoint{2.950308in}{0.953490in}}%
\pgfpathlineto{\pgfqpoint{2.952806in}{0.951317in}}%
\pgfpathlineto{\pgfqpoint{2.958634in}{0.944310in}}%
\pgfpathlineto{\pgfqpoint{2.962381in}{0.942237in}}%
\pgfpathlineto{\pgfqpoint{2.966128in}{0.939505in}}%
\pgfpathlineto{\pgfqpoint{2.970292in}{0.938292in}}%
\pgfpathlineto{\pgfqpoint{2.975704in}{0.934184in}}%
\pgfpathlineto{\pgfqpoint{2.978618in}{0.933424in}}%
\pgfpathlineto{\pgfqpoint{2.981116in}{0.931065in}}%
\pgfpathlineto{\pgfqpoint{2.983614in}{0.929307in}}%
\pgfpathlineto{\pgfqpoint{2.987361in}{0.928038in}}%
\pgfpathlineto{\pgfqpoint{2.989442in}{0.925283in}}%
\pgfpathlineto{\pgfqpoint{2.992357in}{0.921499in}}%
\pgfpathlineto{\pgfqpoint{2.994438in}{0.920632in}}%
\pgfpathlineto{\pgfqpoint{2.996520in}{0.919660in}}%
\pgfpathlineto{\pgfqpoint{2.999018in}{0.916781in}}%
\pgfpathlineto{\pgfqpoint{3.001516in}{0.914474in}}%
\pgfpathlineto{\pgfqpoint{3.006512in}{0.912234in}}%
\pgfpathlineto{\pgfqpoint{3.011091in}{0.907276in}}%
\pgfpathlineto{\pgfqpoint{3.016087in}{0.906188in}}%
\pgfpathlineto{\pgfqpoint{3.020250in}{0.902705in}}%
\pgfpathlineto{\pgfqpoint{3.024830in}{0.901489in}}%
\pgfpathlineto{\pgfqpoint{3.027744in}{0.897216in}}%
\pgfpathlineto{\pgfqpoint{3.029826in}{0.895021in}}%
\pgfpathlineto{\pgfqpoint{3.034763in}{0.891864in}}%
\pgfpathlineto{\pgfqpoint{3.036586in}{0.887714in}}%
\pgfpathlineto{\pgfqpoint{3.038110in}{0.882838in}}%
\pgfpathlineto{\pgfqpoint{3.038110in}{0.882838in}}%
\pgfusepath{stroke}%
\end{pgfscope}%
\begin{pgfscope}%
\pgfpathrectangle{\pgfqpoint{0.650000in}{0.420000in}}{\pgfqpoint{2.388110in}{1.994488in}}%
\pgfusepath{clip}%
\pgfsetrectcap%
\pgfsetroundjoin%
\pgfsetlinewidth{0.803000pt}%
\definecolor{currentstroke}{rgb}{0.500000,0.500000,0.500000}%
\pgfsetstrokecolor{currentstroke}%
\pgfsetdash{}{0pt}%
\pgfpathmoveto{\pgfqpoint{0.650000in}{1.084829in}}%
\pgfpathlineto{\pgfqpoint{3.038110in}{1.084829in}}%
\pgfusepath{stroke}%
\end{pgfscope}%
\begin{pgfscope}%
\pgfsetrectcap%
\pgfsetmiterjoin%
\pgfsetlinewidth{0.803000pt}%
\definecolor{currentstroke}{rgb}{0.000000,0.000000,0.000000}%
\pgfsetstrokecolor{currentstroke}%
\pgfsetdash{}{0pt}%
\pgfpathmoveto{\pgfqpoint{0.650000in}{0.420000in}}%
\pgfpathlineto{\pgfqpoint{0.650000in}{2.414488in}}%
\pgfusepath{stroke}%
\end{pgfscope}%
\begin{pgfscope}%
\pgfsetrectcap%
\pgfsetmiterjoin%
\pgfsetlinewidth{0.803000pt}%
\definecolor{currentstroke}{rgb}{0.000000,0.000000,0.000000}%
\pgfsetstrokecolor{currentstroke}%
\pgfsetdash{}{0pt}%
\pgfpathmoveto{\pgfqpoint{3.038110in}{0.420000in}}%
\pgfpathlineto{\pgfqpoint{3.038110in}{2.414488in}}%
\pgfusepath{stroke}%
\end{pgfscope}%
\begin{pgfscope}%
\pgfsetrectcap%
\pgfsetmiterjoin%
\pgfsetlinewidth{0.803000pt}%
\definecolor{currentstroke}{rgb}{0.000000,0.000000,0.000000}%
\pgfsetstrokecolor{currentstroke}%
\pgfsetdash{}{0pt}%
\pgfpathmoveto{\pgfqpoint{0.650000in}{0.420000in}}%
\pgfpathlineto{\pgfqpoint{3.038110in}{0.420000in}}%
\pgfusepath{stroke}%
\end{pgfscope}%
\begin{pgfscope}%
\pgfsetrectcap%
\pgfsetmiterjoin%
\pgfsetlinewidth{0.803000pt}%
\definecolor{currentstroke}{rgb}{0.000000,0.000000,0.000000}%
\pgfsetstrokecolor{currentstroke}%
\pgfsetdash{}{0pt}%
\pgfpathmoveto{\pgfqpoint{0.650000in}{2.414488in}}%
\pgfpathlineto{\pgfqpoint{3.038110in}{2.414488in}}%
\pgfusepath{stroke}%
\end{pgfscope}%
\begin{pgfscope}%
\pgfsetrectcap%
\pgfsetroundjoin%
\pgfsetlinewidth{1.003750pt}%
\definecolor{currentstroke}{rgb}{0.796078,0.094118,0.113725}%
\pgfsetstrokecolor{currentstroke}%
\pgfsetdash{}{0pt}%
\pgfpathmoveto{\pgfqpoint{1.529452in}{2.275599in}}%
\pgfpathlineto{\pgfqpoint{1.640563in}{2.275599in}}%
\pgfpathlineto{\pgfqpoint{1.751675in}{2.275599in}}%
\pgfusepath{stroke}%
\end{pgfscope}%
\begin{pgfscope}%
\definecolor{textcolor}{rgb}{0.000000,0.000000,0.000000}%
\pgfsetstrokecolor{textcolor}%
\pgfsetfillcolor{textcolor}%
\pgftext[x=1.840563in,y=2.236710in,left,base]{\color{textcolor}{\rmfamily\fontsize{8.000000}{9.600000}\selectfont\catcode`\^=\active\def^{\ifmmode\sp\else\^{}\fi}\catcode`\%=\active\def%{\%}R100deg90 base}}%
\end{pgfscope}%
\begin{pgfscope}%
\pgfsetrectcap%
\pgfsetroundjoin%
\pgfsetlinewidth{1.003750pt}%
\definecolor{currentstroke}{rgb}{0.454902,0.678431,0.819608}%
\pgfsetstrokecolor{currentstroke}%
\pgfsetdash{}{0pt}%
\pgfpathmoveto{\pgfqpoint{1.529452in}{2.120661in}}%
\pgfpathlineto{\pgfqpoint{1.640563in}{2.120661in}}%
\pgfpathlineto{\pgfqpoint{1.751675in}{2.120661in}}%
\pgfusepath{stroke}%
\end{pgfscope}%
\begin{pgfscope}%
\definecolor{textcolor}{rgb}{0.000000,0.000000,0.000000}%
\pgfsetstrokecolor{textcolor}%
\pgfsetfillcolor{textcolor}%
\pgftext[x=1.840563in,y=2.081772in,left,base]{\color{textcolor}{\rmfamily\fontsize{8.000000}{9.600000}\selectfont\catcode`\^=\active\def^{\ifmmode\sp\else\^{}\fi}\catcode`\%=\active\def%{\%}R100deg90 optimised}}%
\end{pgfscope}%
\begin{pgfscope}%
\pgfsetrectcap%
\pgfsetroundjoin%
\pgfsetlinewidth{1.003750pt}%
\definecolor{currentstroke}{rgb}{0.878431,0.509804,0.078431}%
\pgfsetstrokecolor{currentstroke}%
\pgfsetdash{}{0pt}%
\pgfpathmoveto{\pgfqpoint{1.529452in}{1.965723in}}%
\pgfpathlineto{\pgfqpoint{1.640563in}{1.965723in}}%
\pgfpathlineto{\pgfqpoint{1.751675in}{1.965723in}}%
\pgfusepath{stroke}%
\end{pgfscope}%
\begin{pgfscope}%
\definecolor{textcolor}{rgb}{0.000000,0.000000,0.000000}%
\pgfsetstrokecolor{textcolor}%
\pgfsetfillcolor{textcolor}%
\pgftext[x=1.840563in,y=1.926834in,left,base]{\color{textcolor}{\rmfamily\fontsize{8.000000}{9.600000}\selectfont\catcode`\^=\active\def^{\ifmmode\sp\else\^{}\fi}\catcode`\%=\active\def%{\%}R500deg18 base}}%
\end{pgfscope}%
\begin{pgfscope}%
\pgfsetrectcap%
\pgfsetroundjoin%
\pgfsetlinewidth{1.003750pt}%
\definecolor{currentstroke}{rgb}{0.501961,0.450980,0.674510}%
\pgfsetstrokecolor{currentstroke}%
\pgfsetdash{}{0pt}%
\pgfpathmoveto{\pgfqpoint{1.529452in}{1.810784in}}%
\pgfpathlineto{\pgfqpoint{1.640563in}{1.810784in}}%
\pgfpathlineto{\pgfqpoint{1.751675in}{1.810784in}}%
\pgfusepath{stroke}%
\end{pgfscope}%
\begin{pgfscope}%
\definecolor{textcolor}{rgb}{0.000000,0.000000,0.000000}%
\pgfsetstrokecolor{textcolor}%
\pgfsetfillcolor{textcolor}%
\pgftext[x=1.840563in,y=1.771895in,left,base]{\color{textcolor}{\rmfamily\fontsize{8.000000}{9.600000}\selectfont\catcode`\^=\active\def^{\ifmmode\sp\else\^{}\fi}\catcode`\%=\active\def%{\%}R500deg18 optimised}}%
\end{pgfscope}%
\end{pgfpicture}%
\makeatother%
\endgroup%

    \label{fig:R500deg18R100deg90_pres}
\end{subfigure}
\hfill
\begin{subfigure}[t]{0.495\textwidth}
    \centering
    \subcaption{}
    %% Creator: Matplotlib, PGF backend
%%
%% To include the figure in your LaTeX document, write
%%   \input{<filename>.pgf}
%%
%% Make sure the required packages are loaded in your preamble
%%   \usepackage{pgf}
%%
%% Also ensure that all the required font packages are loaded; for instance,
%% the lmodern package is sometimes necessary when using math font.
%%   \usepackage{lmodern}
%%
%% Figures using additional raster images can only be included by \input if
%% they are in the same directory as the main LaTeX file. For loading figures
%% from other directories you can use the `import` package
%%   \usepackage{import}
%%
%% and then include the figures with
%%   \import{<path to file>}{<filename>.pgf}
%%
%% Matplotlib used the following preamble
%%   \def\mathdefault#1{#1}
%%   \everymath=\expandafter{\the\everymath\displaystyle}
%%   \IfFileExists{scrextend.sty}{
%%     \usepackage[fontsize=11.000000pt]{scrextend}
%%   }{
%%     \renewcommand{\normalsize}{\fontsize{11.000000}{13.200000}\selectfont}
%%     \normalsize
%%   }
%%   
%%   \makeatletter\@ifpackageloaded{underscore}{}{\usepackage[strings]{underscore}}\makeatother
%%
\begingroup%
\makeatletter%
\begin{pgfpicture}%
\pgfpathrectangle{\pgfpointorigin}{\pgfqpoint{3.118110in}{2.494488in}}%
\pgfusepath{use as bounding box, clip}%
\begin{pgfscope}%
\pgfsetbuttcap%
\pgfsetmiterjoin%
\definecolor{currentfill}{rgb}{1.000000,1.000000,1.000000}%
\pgfsetfillcolor{currentfill}%
\pgfsetlinewidth{0.000000pt}%
\definecolor{currentstroke}{rgb}{1.000000,1.000000,1.000000}%
\pgfsetstrokecolor{currentstroke}%
\pgfsetdash{}{0pt}%
\pgfpathmoveto{\pgfqpoint{0.000000in}{0.000000in}}%
\pgfpathlineto{\pgfqpoint{3.118110in}{0.000000in}}%
\pgfpathlineto{\pgfqpoint{3.118110in}{2.494488in}}%
\pgfpathlineto{\pgfqpoint{0.000000in}{2.494488in}}%
\pgfpathlineto{\pgfqpoint{0.000000in}{0.000000in}}%
\pgfpathclose%
\pgfusepath{fill}%
\end{pgfscope}%
\begin{pgfscope}%
\pgfsetbuttcap%
\pgfsetmiterjoin%
\definecolor{currentfill}{rgb}{1.000000,1.000000,1.000000}%
\pgfsetfillcolor{currentfill}%
\pgfsetlinewidth{0.000000pt}%
\definecolor{currentstroke}{rgb}{0.000000,0.000000,0.000000}%
\pgfsetstrokecolor{currentstroke}%
\pgfsetstrokeopacity{0.000000}%
\pgfsetdash{}{0pt}%
\pgfpathmoveto{\pgfqpoint{0.650000in}{0.420000in}}%
\pgfpathlineto{\pgfqpoint{3.038110in}{0.420000in}}%
\pgfpathlineto{\pgfqpoint{3.038110in}{2.414488in}}%
\pgfpathlineto{\pgfqpoint{0.650000in}{2.414488in}}%
\pgfpathlineto{\pgfqpoint{0.650000in}{0.420000in}}%
\pgfpathclose%
\pgfusepath{fill}%
\end{pgfscope}%
\begin{pgfscope}%
\pgfsetbuttcap%
\pgfsetroundjoin%
\definecolor{currentfill}{rgb}{0.000000,0.000000,0.000000}%
\pgfsetfillcolor{currentfill}%
\pgfsetlinewidth{0.501875pt}%
\definecolor{currentstroke}{rgb}{0.000000,0.000000,0.000000}%
\pgfsetstrokecolor{currentstroke}%
\pgfsetdash{}{0pt}%
\pgfsys@defobject{currentmarker}{\pgfqpoint{0.000000in}{0.000000in}}{\pgfqpoint{0.000000in}{0.041667in}}{%
\pgfpathmoveto{\pgfqpoint{0.000000in}{0.000000in}}%
\pgfpathlineto{\pgfqpoint{0.000000in}{0.041667in}}%
\pgfusepath{stroke,fill}%
}%
\begin{pgfscope}%
\pgfsys@transformshift{0.650000in}{0.420000in}%
\pgfsys@useobject{currentmarker}{}%
\end{pgfscope}%
\end{pgfscope}%
\begin{pgfscope}%
\pgfsetbuttcap%
\pgfsetroundjoin%
\definecolor{currentfill}{rgb}{0.000000,0.000000,0.000000}%
\pgfsetfillcolor{currentfill}%
\pgfsetlinewidth{0.501875pt}%
\definecolor{currentstroke}{rgb}{0.000000,0.000000,0.000000}%
\pgfsetstrokecolor{currentstroke}%
\pgfsetdash{}{0pt}%
\pgfsys@defobject{currentmarker}{\pgfqpoint{0.000000in}{-0.041667in}}{\pgfqpoint{0.000000in}{0.000000in}}{%
\pgfpathmoveto{\pgfqpoint{0.000000in}{0.000000in}}%
\pgfpathlineto{\pgfqpoint{0.000000in}{-0.041667in}}%
\pgfusepath{stroke,fill}%
}%
\begin{pgfscope}%
\pgfsys@transformshift{0.650000in}{2.414488in}%
\pgfsys@useobject{currentmarker}{}%
\end{pgfscope}%
\end{pgfscope}%
\begin{pgfscope}%
\definecolor{textcolor}{rgb}{0.000000,0.000000,0.000000}%
\pgfsetstrokecolor{textcolor}%
\pgfsetfillcolor{textcolor}%
\pgftext[x=0.650000in,y=0.371389in,,top]{\color{textcolor}{\rmfamily\fontsize{11.000000}{13.200000}\selectfont\catcode`\^=\active\def^{\ifmmode\sp\else\^{}\fi}\catcode`\%=\active\def%{\%}0}}%
\end{pgfscope}%
\begin{pgfscope}%
\pgfsetbuttcap%
\pgfsetroundjoin%
\definecolor{currentfill}{rgb}{0.000000,0.000000,0.000000}%
\pgfsetfillcolor{currentfill}%
\pgfsetlinewidth{0.501875pt}%
\definecolor{currentstroke}{rgb}{0.000000,0.000000,0.000000}%
\pgfsetstrokecolor{currentstroke}%
\pgfsetdash{}{0pt}%
\pgfsys@defobject{currentmarker}{\pgfqpoint{0.000000in}{0.000000in}}{\pgfqpoint{0.000000in}{0.041667in}}{%
\pgfpathmoveto{\pgfqpoint{0.000000in}{0.000000in}}%
\pgfpathlineto{\pgfqpoint{0.000000in}{0.041667in}}%
\pgfusepath{stroke,fill}%
}%
\begin{pgfscope}%
\pgfsys@transformshift{1.493052in}{0.420000in}%
\pgfsys@useobject{currentmarker}{}%
\end{pgfscope}%
\end{pgfscope}%
\begin{pgfscope}%
\pgfsetbuttcap%
\pgfsetroundjoin%
\definecolor{currentfill}{rgb}{0.000000,0.000000,0.000000}%
\pgfsetfillcolor{currentfill}%
\pgfsetlinewidth{0.501875pt}%
\definecolor{currentstroke}{rgb}{0.000000,0.000000,0.000000}%
\pgfsetstrokecolor{currentstroke}%
\pgfsetdash{}{0pt}%
\pgfsys@defobject{currentmarker}{\pgfqpoint{0.000000in}{-0.041667in}}{\pgfqpoint{0.000000in}{0.000000in}}{%
\pgfpathmoveto{\pgfqpoint{0.000000in}{0.000000in}}%
\pgfpathlineto{\pgfqpoint{0.000000in}{-0.041667in}}%
\pgfusepath{stroke,fill}%
}%
\begin{pgfscope}%
\pgfsys@transformshift{1.493052in}{2.414488in}%
\pgfsys@useobject{currentmarker}{}%
\end{pgfscope}%
\end{pgfscope}%
\begin{pgfscope}%
\definecolor{textcolor}{rgb}{0.000000,0.000000,0.000000}%
\pgfsetstrokecolor{textcolor}%
\pgfsetfillcolor{textcolor}%
\pgftext[x=1.493052in,y=0.371389in,,top]{\color{textcolor}{\rmfamily\fontsize{11.000000}{13.200000}\selectfont\catcode`\^=\active\def^{\ifmmode\sp\else\^{}\fi}\catcode`\%=\active\def%{\%}100}}%
\end{pgfscope}%
\begin{pgfscope}%
\pgfsetbuttcap%
\pgfsetroundjoin%
\definecolor{currentfill}{rgb}{0.000000,0.000000,0.000000}%
\pgfsetfillcolor{currentfill}%
\pgfsetlinewidth{0.501875pt}%
\definecolor{currentstroke}{rgb}{0.000000,0.000000,0.000000}%
\pgfsetstrokecolor{currentstroke}%
\pgfsetdash{}{0pt}%
\pgfsys@defobject{currentmarker}{\pgfqpoint{0.000000in}{0.000000in}}{\pgfqpoint{0.000000in}{0.041667in}}{%
\pgfpathmoveto{\pgfqpoint{0.000000in}{0.000000in}}%
\pgfpathlineto{\pgfqpoint{0.000000in}{0.041667in}}%
\pgfusepath{stroke,fill}%
}%
\begin{pgfscope}%
\pgfsys@transformshift{2.336103in}{0.420000in}%
\pgfsys@useobject{currentmarker}{}%
\end{pgfscope}%
\end{pgfscope}%
\begin{pgfscope}%
\pgfsetbuttcap%
\pgfsetroundjoin%
\definecolor{currentfill}{rgb}{0.000000,0.000000,0.000000}%
\pgfsetfillcolor{currentfill}%
\pgfsetlinewidth{0.501875pt}%
\definecolor{currentstroke}{rgb}{0.000000,0.000000,0.000000}%
\pgfsetstrokecolor{currentstroke}%
\pgfsetdash{}{0pt}%
\pgfsys@defobject{currentmarker}{\pgfqpoint{0.000000in}{-0.041667in}}{\pgfqpoint{0.000000in}{0.000000in}}{%
\pgfpathmoveto{\pgfqpoint{0.000000in}{0.000000in}}%
\pgfpathlineto{\pgfqpoint{0.000000in}{-0.041667in}}%
\pgfusepath{stroke,fill}%
}%
\begin{pgfscope}%
\pgfsys@transformshift{2.336103in}{2.414488in}%
\pgfsys@useobject{currentmarker}{}%
\end{pgfscope}%
\end{pgfscope}%
\begin{pgfscope}%
\definecolor{textcolor}{rgb}{0.000000,0.000000,0.000000}%
\pgfsetstrokecolor{textcolor}%
\pgfsetfillcolor{textcolor}%
\pgftext[x=2.336103in,y=0.371389in,,top]{\color{textcolor}{\rmfamily\fontsize{11.000000}{13.200000}\selectfont\catcode`\^=\active\def^{\ifmmode\sp\else\^{}\fi}\catcode`\%=\active\def%{\%}200}}%
\end{pgfscope}%
\begin{pgfscope}%
\pgfsetbuttcap%
\pgfsetroundjoin%
\definecolor{currentfill}{rgb}{0.000000,0.000000,0.000000}%
\pgfsetfillcolor{currentfill}%
\pgfsetlinewidth{0.401500pt}%
\definecolor{currentstroke}{rgb}{0.000000,0.000000,0.000000}%
\pgfsetstrokecolor{currentstroke}%
\pgfsetdash{}{0pt}%
\pgfsys@defobject{currentmarker}{\pgfqpoint{0.000000in}{0.000000in}}{\pgfqpoint{0.000000in}{0.020833in}}{%
\pgfpathmoveto{\pgfqpoint{0.000000in}{0.000000in}}%
\pgfpathlineto{\pgfqpoint{0.000000in}{0.020833in}}%
\pgfusepath{stroke,fill}%
}%
\begin{pgfscope}%
\pgfsys@transformshift{0.818610in}{0.420000in}%
\pgfsys@useobject{currentmarker}{}%
\end{pgfscope}%
\end{pgfscope}%
\begin{pgfscope}%
\pgfsetbuttcap%
\pgfsetroundjoin%
\definecolor{currentfill}{rgb}{0.000000,0.000000,0.000000}%
\pgfsetfillcolor{currentfill}%
\pgfsetlinewidth{0.401500pt}%
\definecolor{currentstroke}{rgb}{0.000000,0.000000,0.000000}%
\pgfsetstrokecolor{currentstroke}%
\pgfsetdash{}{0pt}%
\pgfsys@defobject{currentmarker}{\pgfqpoint{0.000000in}{-0.020833in}}{\pgfqpoint{0.000000in}{0.000000in}}{%
\pgfpathmoveto{\pgfqpoint{0.000000in}{0.000000in}}%
\pgfpathlineto{\pgfqpoint{0.000000in}{-0.020833in}}%
\pgfusepath{stroke,fill}%
}%
\begin{pgfscope}%
\pgfsys@transformshift{0.818610in}{2.414488in}%
\pgfsys@useobject{currentmarker}{}%
\end{pgfscope}%
\end{pgfscope}%
\begin{pgfscope}%
\pgfsetbuttcap%
\pgfsetroundjoin%
\definecolor{currentfill}{rgb}{0.000000,0.000000,0.000000}%
\pgfsetfillcolor{currentfill}%
\pgfsetlinewidth{0.401500pt}%
\definecolor{currentstroke}{rgb}{0.000000,0.000000,0.000000}%
\pgfsetstrokecolor{currentstroke}%
\pgfsetdash{}{0pt}%
\pgfsys@defobject{currentmarker}{\pgfqpoint{0.000000in}{0.000000in}}{\pgfqpoint{0.000000in}{0.020833in}}{%
\pgfpathmoveto{\pgfqpoint{0.000000in}{0.000000in}}%
\pgfpathlineto{\pgfqpoint{0.000000in}{0.020833in}}%
\pgfusepath{stroke,fill}%
}%
\begin{pgfscope}%
\pgfsys@transformshift{0.987221in}{0.420000in}%
\pgfsys@useobject{currentmarker}{}%
\end{pgfscope}%
\end{pgfscope}%
\begin{pgfscope}%
\pgfsetbuttcap%
\pgfsetroundjoin%
\definecolor{currentfill}{rgb}{0.000000,0.000000,0.000000}%
\pgfsetfillcolor{currentfill}%
\pgfsetlinewidth{0.401500pt}%
\definecolor{currentstroke}{rgb}{0.000000,0.000000,0.000000}%
\pgfsetstrokecolor{currentstroke}%
\pgfsetdash{}{0pt}%
\pgfsys@defobject{currentmarker}{\pgfqpoint{0.000000in}{-0.020833in}}{\pgfqpoint{0.000000in}{0.000000in}}{%
\pgfpathmoveto{\pgfqpoint{0.000000in}{0.000000in}}%
\pgfpathlineto{\pgfqpoint{0.000000in}{-0.020833in}}%
\pgfusepath{stroke,fill}%
}%
\begin{pgfscope}%
\pgfsys@transformshift{0.987221in}{2.414488in}%
\pgfsys@useobject{currentmarker}{}%
\end{pgfscope}%
\end{pgfscope}%
\begin{pgfscope}%
\pgfsetbuttcap%
\pgfsetroundjoin%
\definecolor{currentfill}{rgb}{0.000000,0.000000,0.000000}%
\pgfsetfillcolor{currentfill}%
\pgfsetlinewidth{0.401500pt}%
\definecolor{currentstroke}{rgb}{0.000000,0.000000,0.000000}%
\pgfsetstrokecolor{currentstroke}%
\pgfsetdash{}{0pt}%
\pgfsys@defobject{currentmarker}{\pgfqpoint{0.000000in}{0.000000in}}{\pgfqpoint{0.000000in}{0.020833in}}{%
\pgfpathmoveto{\pgfqpoint{0.000000in}{0.000000in}}%
\pgfpathlineto{\pgfqpoint{0.000000in}{0.020833in}}%
\pgfusepath{stroke,fill}%
}%
\begin{pgfscope}%
\pgfsys@transformshift{1.155831in}{0.420000in}%
\pgfsys@useobject{currentmarker}{}%
\end{pgfscope}%
\end{pgfscope}%
\begin{pgfscope}%
\pgfsetbuttcap%
\pgfsetroundjoin%
\definecolor{currentfill}{rgb}{0.000000,0.000000,0.000000}%
\pgfsetfillcolor{currentfill}%
\pgfsetlinewidth{0.401500pt}%
\definecolor{currentstroke}{rgb}{0.000000,0.000000,0.000000}%
\pgfsetstrokecolor{currentstroke}%
\pgfsetdash{}{0pt}%
\pgfsys@defobject{currentmarker}{\pgfqpoint{0.000000in}{-0.020833in}}{\pgfqpoint{0.000000in}{0.000000in}}{%
\pgfpathmoveto{\pgfqpoint{0.000000in}{0.000000in}}%
\pgfpathlineto{\pgfqpoint{0.000000in}{-0.020833in}}%
\pgfusepath{stroke,fill}%
}%
\begin{pgfscope}%
\pgfsys@transformshift{1.155831in}{2.414488in}%
\pgfsys@useobject{currentmarker}{}%
\end{pgfscope}%
\end{pgfscope}%
\begin{pgfscope}%
\pgfsetbuttcap%
\pgfsetroundjoin%
\definecolor{currentfill}{rgb}{0.000000,0.000000,0.000000}%
\pgfsetfillcolor{currentfill}%
\pgfsetlinewidth{0.401500pt}%
\definecolor{currentstroke}{rgb}{0.000000,0.000000,0.000000}%
\pgfsetstrokecolor{currentstroke}%
\pgfsetdash{}{0pt}%
\pgfsys@defobject{currentmarker}{\pgfqpoint{0.000000in}{0.000000in}}{\pgfqpoint{0.000000in}{0.020833in}}{%
\pgfpathmoveto{\pgfqpoint{0.000000in}{0.000000in}}%
\pgfpathlineto{\pgfqpoint{0.000000in}{0.020833in}}%
\pgfusepath{stroke,fill}%
}%
\begin{pgfscope}%
\pgfsys@transformshift{1.324441in}{0.420000in}%
\pgfsys@useobject{currentmarker}{}%
\end{pgfscope}%
\end{pgfscope}%
\begin{pgfscope}%
\pgfsetbuttcap%
\pgfsetroundjoin%
\definecolor{currentfill}{rgb}{0.000000,0.000000,0.000000}%
\pgfsetfillcolor{currentfill}%
\pgfsetlinewidth{0.401500pt}%
\definecolor{currentstroke}{rgb}{0.000000,0.000000,0.000000}%
\pgfsetstrokecolor{currentstroke}%
\pgfsetdash{}{0pt}%
\pgfsys@defobject{currentmarker}{\pgfqpoint{0.000000in}{-0.020833in}}{\pgfqpoint{0.000000in}{0.000000in}}{%
\pgfpathmoveto{\pgfqpoint{0.000000in}{0.000000in}}%
\pgfpathlineto{\pgfqpoint{0.000000in}{-0.020833in}}%
\pgfusepath{stroke,fill}%
}%
\begin{pgfscope}%
\pgfsys@transformshift{1.324441in}{2.414488in}%
\pgfsys@useobject{currentmarker}{}%
\end{pgfscope}%
\end{pgfscope}%
\begin{pgfscope}%
\pgfsetbuttcap%
\pgfsetroundjoin%
\definecolor{currentfill}{rgb}{0.000000,0.000000,0.000000}%
\pgfsetfillcolor{currentfill}%
\pgfsetlinewidth{0.401500pt}%
\definecolor{currentstroke}{rgb}{0.000000,0.000000,0.000000}%
\pgfsetstrokecolor{currentstroke}%
\pgfsetdash{}{0pt}%
\pgfsys@defobject{currentmarker}{\pgfqpoint{0.000000in}{0.000000in}}{\pgfqpoint{0.000000in}{0.020833in}}{%
\pgfpathmoveto{\pgfqpoint{0.000000in}{0.000000in}}%
\pgfpathlineto{\pgfqpoint{0.000000in}{0.020833in}}%
\pgfusepath{stroke,fill}%
}%
\begin{pgfscope}%
\pgfsys@transformshift{1.661662in}{0.420000in}%
\pgfsys@useobject{currentmarker}{}%
\end{pgfscope}%
\end{pgfscope}%
\begin{pgfscope}%
\pgfsetbuttcap%
\pgfsetroundjoin%
\definecolor{currentfill}{rgb}{0.000000,0.000000,0.000000}%
\pgfsetfillcolor{currentfill}%
\pgfsetlinewidth{0.401500pt}%
\definecolor{currentstroke}{rgb}{0.000000,0.000000,0.000000}%
\pgfsetstrokecolor{currentstroke}%
\pgfsetdash{}{0pt}%
\pgfsys@defobject{currentmarker}{\pgfqpoint{0.000000in}{-0.020833in}}{\pgfqpoint{0.000000in}{0.000000in}}{%
\pgfpathmoveto{\pgfqpoint{0.000000in}{0.000000in}}%
\pgfpathlineto{\pgfqpoint{0.000000in}{-0.020833in}}%
\pgfusepath{stroke,fill}%
}%
\begin{pgfscope}%
\pgfsys@transformshift{1.661662in}{2.414488in}%
\pgfsys@useobject{currentmarker}{}%
\end{pgfscope}%
\end{pgfscope}%
\begin{pgfscope}%
\pgfsetbuttcap%
\pgfsetroundjoin%
\definecolor{currentfill}{rgb}{0.000000,0.000000,0.000000}%
\pgfsetfillcolor{currentfill}%
\pgfsetlinewidth{0.401500pt}%
\definecolor{currentstroke}{rgb}{0.000000,0.000000,0.000000}%
\pgfsetstrokecolor{currentstroke}%
\pgfsetdash{}{0pt}%
\pgfsys@defobject{currentmarker}{\pgfqpoint{0.000000in}{0.000000in}}{\pgfqpoint{0.000000in}{0.020833in}}{%
\pgfpathmoveto{\pgfqpoint{0.000000in}{0.000000in}}%
\pgfpathlineto{\pgfqpoint{0.000000in}{0.020833in}}%
\pgfusepath{stroke,fill}%
}%
\begin{pgfscope}%
\pgfsys@transformshift{1.830272in}{0.420000in}%
\pgfsys@useobject{currentmarker}{}%
\end{pgfscope}%
\end{pgfscope}%
\begin{pgfscope}%
\pgfsetbuttcap%
\pgfsetroundjoin%
\definecolor{currentfill}{rgb}{0.000000,0.000000,0.000000}%
\pgfsetfillcolor{currentfill}%
\pgfsetlinewidth{0.401500pt}%
\definecolor{currentstroke}{rgb}{0.000000,0.000000,0.000000}%
\pgfsetstrokecolor{currentstroke}%
\pgfsetdash{}{0pt}%
\pgfsys@defobject{currentmarker}{\pgfqpoint{0.000000in}{-0.020833in}}{\pgfqpoint{0.000000in}{0.000000in}}{%
\pgfpathmoveto{\pgfqpoint{0.000000in}{0.000000in}}%
\pgfpathlineto{\pgfqpoint{0.000000in}{-0.020833in}}%
\pgfusepath{stroke,fill}%
}%
\begin{pgfscope}%
\pgfsys@transformshift{1.830272in}{2.414488in}%
\pgfsys@useobject{currentmarker}{}%
\end{pgfscope}%
\end{pgfscope}%
\begin{pgfscope}%
\pgfsetbuttcap%
\pgfsetroundjoin%
\definecolor{currentfill}{rgb}{0.000000,0.000000,0.000000}%
\pgfsetfillcolor{currentfill}%
\pgfsetlinewidth{0.401500pt}%
\definecolor{currentstroke}{rgb}{0.000000,0.000000,0.000000}%
\pgfsetstrokecolor{currentstroke}%
\pgfsetdash{}{0pt}%
\pgfsys@defobject{currentmarker}{\pgfqpoint{0.000000in}{0.000000in}}{\pgfqpoint{0.000000in}{0.020833in}}{%
\pgfpathmoveto{\pgfqpoint{0.000000in}{0.000000in}}%
\pgfpathlineto{\pgfqpoint{0.000000in}{0.020833in}}%
\pgfusepath{stroke,fill}%
}%
\begin{pgfscope}%
\pgfsys@transformshift{1.998883in}{0.420000in}%
\pgfsys@useobject{currentmarker}{}%
\end{pgfscope}%
\end{pgfscope}%
\begin{pgfscope}%
\pgfsetbuttcap%
\pgfsetroundjoin%
\definecolor{currentfill}{rgb}{0.000000,0.000000,0.000000}%
\pgfsetfillcolor{currentfill}%
\pgfsetlinewidth{0.401500pt}%
\definecolor{currentstroke}{rgb}{0.000000,0.000000,0.000000}%
\pgfsetstrokecolor{currentstroke}%
\pgfsetdash{}{0pt}%
\pgfsys@defobject{currentmarker}{\pgfqpoint{0.000000in}{-0.020833in}}{\pgfqpoint{0.000000in}{0.000000in}}{%
\pgfpathmoveto{\pgfqpoint{0.000000in}{0.000000in}}%
\pgfpathlineto{\pgfqpoint{0.000000in}{-0.020833in}}%
\pgfusepath{stroke,fill}%
}%
\begin{pgfscope}%
\pgfsys@transformshift{1.998883in}{2.414488in}%
\pgfsys@useobject{currentmarker}{}%
\end{pgfscope}%
\end{pgfscope}%
\begin{pgfscope}%
\pgfsetbuttcap%
\pgfsetroundjoin%
\definecolor{currentfill}{rgb}{0.000000,0.000000,0.000000}%
\pgfsetfillcolor{currentfill}%
\pgfsetlinewidth{0.401500pt}%
\definecolor{currentstroke}{rgb}{0.000000,0.000000,0.000000}%
\pgfsetstrokecolor{currentstroke}%
\pgfsetdash{}{0pt}%
\pgfsys@defobject{currentmarker}{\pgfqpoint{0.000000in}{0.000000in}}{\pgfqpoint{0.000000in}{0.020833in}}{%
\pgfpathmoveto{\pgfqpoint{0.000000in}{0.000000in}}%
\pgfpathlineto{\pgfqpoint{0.000000in}{0.020833in}}%
\pgfusepath{stroke,fill}%
}%
\begin{pgfscope}%
\pgfsys@transformshift{2.167493in}{0.420000in}%
\pgfsys@useobject{currentmarker}{}%
\end{pgfscope}%
\end{pgfscope}%
\begin{pgfscope}%
\pgfsetbuttcap%
\pgfsetroundjoin%
\definecolor{currentfill}{rgb}{0.000000,0.000000,0.000000}%
\pgfsetfillcolor{currentfill}%
\pgfsetlinewidth{0.401500pt}%
\definecolor{currentstroke}{rgb}{0.000000,0.000000,0.000000}%
\pgfsetstrokecolor{currentstroke}%
\pgfsetdash{}{0pt}%
\pgfsys@defobject{currentmarker}{\pgfqpoint{0.000000in}{-0.020833in}}{\pgfqpoint{0.000000in}{0.000000in}}{%
\pgfpathmoveto{\pgfqpoint{0.000000in}{0.000000in}}%
\pgfpathlineto{\pgfqpoint{0.000000in}{-0.020833in}}%
\pgfusepath{stroke,fill}%
}%
\begin{pgfscope}%
\pgfsys@transformshift{2.167493in}{2.414488in}%
\pgfsys@useobject{currentmarker}{}%
\end{pgfscope}%
\end{pgfscope}%
\begin{pgfscope}%
\pgfsetbuttcap%
\pgfsetroundjoin%
\definecolor{currentfill}{rgb}{0.000000,0.000000,0.000000}%
\pgfsetfillcolor{currentfill}%
\pgfsetlinewidth{0.401500pt}%
\definecolor{currentstroke}{rgb}{0.000000,0.000000,0.000000}%
\pgfsetstrokecolor{currentstroke}%
\pgfsetdash{}{0pt}%
\pgfsys@defobject{currentmarker}{\pgfqpoint{0.000000in}{0.000000in}}{\pgfqpoint{0.000000in}{0.020833in}}{%
\pgfpathmoveto{\pgfqpoint{0.000000in}{0.000000in}}%
\pgfpathlineto{\pgfqpoint{0.000000in}{0.020833in}}%
\pgfusepath{stroke,fill}%
}%
\begin{pgfscope}%
\pgfsys@transformshift{2.504714in}{0.420000in}%
\pgfsys@useobject{currentmarker}{}%
\end{pgfscope}%
\end{pgfscope}%
\begin{pgfscope}%
\pgfsetbuttcap%
\pgfsetroundjoin%
\definecolor{currentfill}{rgb}{0.000000,0.000000,0.000000}%
\pgfsetfillcolor{currentfill}%
\pgfsetlinewidth{0.401500pt}%
\definecolor{currentstroke}{rgb}{0.000000,0.000000,0.000000}%
\pgfsetstrokecolor{currentstroke}%
\pgfsetdash{}{0pt}%
\pgfsys@defobject{currentmarker}{\pgfqpoint{0.000000in}{-0.020833in}}{\pgfqpoint{0.000000in}{0.000000in}}{%
\pgfpathmoveto{\pgfqpoint{0.000000in}{0.000000in}}%
\pgfpathlineto{\pgfqpoint{0.000000in}{-0.020833in}}%
\pgfusepath{stroke,fill}%
}%
\begin{pgfscope}%
\pgfsys@transformshift{2.504714in}{2.414488in}%
\pgfsys@useobject{currentmarker}{}%
\end{pgfscope}%
\end{pgfscope}%
\begin{pgfscope}%
\pgfsetbuttcap%
\pgfsetroundjoin%
\definecolor{currentfill}{rgb}{0.000000,0.000000,0.000000}%
\pgfsetfillcolor{currentfill}%
\pgfsetlinewidth{0.401500pt}%
\definecolor{currentstroke}{rgb}{0.000000,0.000000,0.000000}%
\pgfsetstrokecolor{currentstroke}%
\pgfsetdash{}{0pt}%
\pgfsys@defobject{currentmarker}{\pgfqpoint{0.000000in}{0.000000in}}{\pgfqpoint{0.000000in}{0.020833in}}{%
\pgfpathmoveto{\pgfqpoint{0.000000in}{0.000000in}}%
\pgfpathlineto{\pgfqpoint{0.000000in}{0.020833in}}%
\pgfusepath{stroke,fill}%
}%
\begin{pgfscope}%
\pgfsys@transformshift{2.673324in}{0.420000in}%
\pgfsys@useobject{currentmarker}{}%
\end{pgfscope}%
\end{pgfscope}%
\begin{pgfscope}%
\pgfsetbuttcap%
\pgfsetroundjoin%
\definecolor{currentfill}{rgb}{0.000000,0.000000,0.000000}%
\pgfsetfillcolor{currentfill}%
\pgfsetlinewidth{0.401500pt}%
\definecolor{currentstroke}{rgb}{0.000000,0.000000,0.000000}%
\pgfsetstrokecolor{currentstroke}%
\pgfsetdash{}{0pt}%
\pgfsys@defobject{currentmarker}{\pgfqpoint{0.000000in}{-0.020833in}}{\pgfqpoint{0.000000in}{0.000000in}}{%
\pgfpathmoveto{\pgfqpoint{0.000000in}{0.000000in}}%
\pgfpathlineto{\pgfqpoint{0.000000in}{-0.020833in}}%
\pgfusepath{stroke,fill}%
}%
\begin{pgfscope}%
\pgfsys@transformshift{2.673324in}{2.414488in}%
\pgfsys@useobject{currentmarker}{}%
\end{pgfscope}%
\end{pgfscope}%
\begin{pgfscope}%
\pgfsetbuttcap%
\pgfsetroundjoin%
\definecolor{currentfill}{rgb}{0.000000,0.000000,0.000000}%
\pgfsetfillcolor{currentfill}%
\pgfsetlinewidth{0.401500pt}%
\definecolor{currentstroke}{rgb}{0.000000,0.000000,0.000000}%
\pgfsetstrokecolor{currentstroke}%
\pgfsetdash{}{0pt}%
\pgfsys@defobject{currentmarker}{\pgfqpoint{0.000000in}{0.000000in}}{\pgfqpoint{0.000000in}{0.020833in}}{%
\pgfpathmoveto{\pgfqpoint{0.000000in}{0.000000in}}%
\pgfpathlineto{\pgfqpoint{0.000000in}{0.020833in}}%
\pgfusepath{stroke,fill}%
}%
\begin{pgfscope}%
\pgfsys@transformshift{2.841934in}{0.420000in}%
\pgfsys@useobject{currentmarker}{}%
\end{pgfscope}%
\end{pgfscope}%
\begin{pgfscope}%
\pgfsetbuttcap%
\pgfsetroundjoin%
\definecolor{currentfill}{rgb}{0.000000,0.000000,0.000000}%
\pgfsetfillcolor{currentfill}%
\pgfsetlinewidth{0.401500pt}%
\definecolor{currentstroke}{rgb}{0.000000,0.000000,0.000000}%
\pgfsetstrokecolor{currentstroke}%
\pgfsetdash{}{0pt}%
\pgfsys@defobject{currentmarker}{\pgfqpoint{0.000000in}{-0.020833in}}{\pgfqpoint{0.000000in}{0.000000in}}{%
\pgfpathmoveto{\pgfqpoint{0.000000in}{0.000000in}}%
\pgfpathlineto{\pgfqpoint{0.000000in}{-0.020833in}}%
\pgfusepath{stroke,fill}%
}%
\begin{pgfscope}%
\pgfsys@transformshift{2.841934in}{2.414488in}%
\pgfsys@useobject{currentmarker}{}%
\end{pgfscope}%
\end{pgfscope}%
\begin{pgfscope}%
\pgfsetbuttcap%
\pgfsetroundjoin%
\definecolor{currentfill}{rgb}{0.000000,0.000000,0.000000}%
\pgfsetfillcolor{currentfill}%
\pgfsetlinewidth{0.401500pt}%
\definecolor{currentstroke}{rgb}{0.000000,0.000000,0.000000}%
\pgfsetstrokecolor{currentstroke}%
\pgfsetdash{}{0pt}%
\pgfsys@defobject{currentmarker}{\pgfqpoint{0.000000in}{0.000000in}}{\pgfqpoint{0.000000in}{0.020833in}}{%
\pgfpathmoveto{\pgfqpoint{0.000000in}{0.000000in}}%
\pgfpathlineto{\pgfqpoint{0.000000in}{0.020833in}}%
\pgfusepath{stroke,fill}%
}%
\begin{pgfscope}%
\pgfsys@transformshift{3.010545in}{0.420000in}%
\pgfsys@useobject{currentmarker}{}%
\end{pgfscope}%
\end{pgfscope}%
\begin{pgfscope}%
\pgfsetbuttcap%
\pgfsetroundjoin%
\definecolor{currentfill}{rgb}{0.000000,0.000000,0.000000}%
\pgfsetfillcolor{currentfill}%
\pgfsetlinewidth{0.401500pt}%
\definecolor{currentstroke}{rgb}{0.000000,0.000000,0.000000}%
\pgfsetstrokecolor{currentstroke}%
\pgfsetdash{}{0pt}%
\pgfsys@defobject{currentmarker}{\pgfqpoint{0.000000in}{-0.020833in}}{\pgfqpoint{0.000000in}{0.000000in}}{%
\pgfpathmoveto{\pgfqpoint{0.000000in}{0.000000in}}%
\pgfpathlineto{\pgfqpoint{0.000000in}{-0.020833in}}%
\pgfusepath{stroke,fill}%
}%
\begin{pgfscope}%
\pgfsys@transformshift{3.010545in}{2.414488in}%
\pgfsys@useobject{currentmarker}{}%
\end{pgfscope}%
\end{pgfscope}%
\begin{pgfscope}%
\definecolor{textcolor}{rgb}{0.000000,0.000000,0.000000}%
\pgfsetstrokecolor{textcolor}%
\pgfsetfillcolor{textcolor}%
\pgftext[x=1.844055in,y=0.208426in,,top]{\color{textcolor}{\rmfamily\fontsize{12.000000}{14.400000}\selectfont\catcode`\^=\active\def^{\ifmmode\sp\else\^{}\fi}\catcode`\%=\active\def%{\%}$s/\delta_{0}$}}%
\end{pgfscope}%
\begin{pgfscope}%
\pgfsetbuttcap%
\pgfsetroundjoin%
\definecolor{currentfill}{rgb}{0.000000,0.000000,0.000000}%
\pgfsetfillcolor{currentfill}%
\pgfsetlinewidth{0.501875pt}%
\definecolor{currentstroke}{rgb}{0.000000,0.000000,0.000000}%
\pgfsetstrokecolor{currentstroke}%
\pgfsetdash{}{0pt}%
\pgfsys@defobject{currentmarker}{\pgfqpoint{0.000000in}{0.000000in}}{\pgfqpoint{0.041667in}{0.000000in}}{%
\pgfpathmoveto{\pgfqpoint{0.000000in}{0.000000in}}%
\pgfpathlineto{\pgfqpoint{0.041667in}{0.000000in}}%
\pgfusepath{stroke,fill}%
}%
\begin{pgfscope}%
\pgfsys@transformshift{0.650000in}{0.420000in}%
\pgfsys@useobject{currentmarker}{}%
\end{pgfscope}%
\end{pgfscope}%
\begin{pgfscope}%
\pgfsetbuttcap%
\pgfsetroundjoin%
\definecolor{currentfill}{rgb}{0.000000,0.000000,0.000000}%
\pgfsetfillcolor{currentfill}%
\pgfsetlinewidth{0.501875pt}%
\definecolor{currentstroke}{rgb}{0.000000,0.000000,0.000000}%
\pgfsetstrokecolor{currentstroke}%
\pgfsetdash{}{0pt}%
\pgfsys@defobject{currentmarker}{\pgfqpoint{-0.041667in}{0.000000in}}{\pgfqpoint{-0.000000in}{0.000000in}}{%
\pgfpathmoveto{\pgfqpoint{-0.000000in}{0.000000in}}%
\pgfpathlineto{\pgfqpoint{-0.041667in}{0.000000in}}%
\pgfusepath{stroke,fill}%
}%
\begin{pgfscope}%
\pgfsys@transformshift{3.038110in}{0.420000in}%
\pgfsys@useobject{currentmarker}{}%
\end{pgfscope}%
\end{pgfscope}%
\begin{pgfscope}%
\definecolor{textcolor}{rgb}{0.000000,0.000000,0.000000}%
\pgfsetstrokecolor{textcolor}%
\pgfsetfillcolor{textcolor}%
\pgftext[x=0.331018in, y=0.367193in, left, base]{\color{textcolor}{\rmfamily\fontsize{11.000000}{13.200000}\selectfont\catcode`\^=\active\def^{\ifmmode\sp\else\^{}\fi}\catcode`\%=\active\def%{\%}0.90}}%
\end{pgfscope}%
\begin{pgfscope}%
\pgfsetbuttcap%
\pgfsetroundjoin%
\definecolor{currentfill}{rgb}{0.000000,0.000000,0.000000}%
\pgfsetfillcolor{currentfill}%
\pgfsetlinewidth{0.501875pt}%
\definecolor{currentstroke}{rgb}{0.000000,0.000000,0.000000}%
\pgfsetstrokecolor{currentstroke}%
\pgfsetdash{}{0pt}%
\pgfsys@defobject{currentmarker}{\pgfqpoint{0.000000in}{0.000000in}}{\pgfqpoint{0.041667in}{0.000000in}}{%
\pgfpathmoveto{\pgfqpoint{0.000000in}{0.000000in}}%
\pgfpathlineto{\pgfqpoint{0.041667in}{0.000000in}}%
\pgfusepath{stroke,fill}%
}%
\begin{pgfscope}%
\pgfsys@transformshift{0.650000in}{0.918622in}%
\pgfsys@useobject{currentmarker}{}%
\end{pgfscope}%
\end{pgfscope}%
\begin{pgfscope}%
\pgfsetbuttcap%
\pgfsetroundjoin%
\definecolor{currentfill}{rgb}{0.000000,0.000000,0.000000}%
\pgfsetfillcolor{currentfill}%
\pgfsetlinewidth{0.501875pt}%
\definecolor{currentstroke}{rgb}{0.000000,0.000000,0.000000}%
\pgfsetstrokecolor{currentstroke}%
\pgfsetdash{}{0pt}%
\pgfsys@defobject{currentmarker}{\pgfqpoint{-0.041667in}{0.000000in}}{\pgfqpoint{-0.000000in}{0.000000in}}{%
\pgfpathmoveto{\pgfqpoint{-0.000000in}{0.000000in}}%
\pgfpathlineto{\pgfqpoint{-0.041667in}{0.000000in}}%
\pgfusepath{stroke,fill}%
}%
\begin{pgfscope}%
\pgfsys@transformshift{3.038110in}{0.918622in}%
\pgfsys@useobject{currentmarker}{}%
\end{pgfscope}%
\end{pgfscope}%
\begin{pgfscope}%
\definecolor{textcolor}{rgb}{0.000000,0.000000,0.000000}%
\pgfsetstrokecolor{textcolor}%
\pgfsetfillcolor{textcolor}%
\pgftext[x=0.331018in, y=0.865815in, left, base]{\color{textcolor}{\rmfamily\fontsize{11.000000}{13.200000}\selectfont\catcode`\^=\active\def^{\ifmmode\sp\else\^{}\fi}\catcode`\%=\active\def%{\%}0.95}}%
\end{pgfscope}%
\begin{pgfscope}%
\pgfsetbuttcap%
\pgfsetroundjoin%
\definecolor{currentfill}{rgb}{0.000000,0.000000,0.000000}%
\pgfsetfillcolor{currentfill}%
\pgfsetlinewidth{0.501875pt}%
\definecolor{currentstroke}{rgb}{0.000000,0.000000,0.000000}%
\pgfsetstrokecolor{currentstroke}%
\pgfsetdash{}{0pt}%
\pgfsys@defobject{currentmarker}{\pgfqpoint{0.000000in}{0.000000in}}{\pgfqpoint{0.041667in}{0.000000in}}{%
\pgfpathmoveto{\pgfqpoint{0.000000in}{0.000000in}}%
\pgfpathlineto{\pgfqpoint{0.041667in}{0.000000in}}%
\pgfusepath{stroke,fill}%
}%
\begin{pgfscope}%
\pgfsys@transformshift{0.650000in}{1.417244in}%
\pgfsys@useobject{currentmarker}{}%
\end{pgfscope}%
\end{pgfscope}%
\begin{pgfscope}%
\pgfsetbuttcap%
\pgfsetroundjoin%
\definecolor{currentfill}{rgb}{0.000000,0.000000,0.000000}%
\pgfsetfillcolor{currentfill}%
\pgfsetlinewidth{0.501875pt}%
\definecolor{currentstroke}{rgb}{0.000000,0.000000,0.000000}%
\pgfsetstrokecolor{currentstroke}%
\pgfsetdash{}{0pt}%
\pgfsys@defobject{currentmarker}{\pgfqpoint{-0.041667in}{0.000000in}}{\pgfqpoint{-0.000000in}{0.000000in}}{%
\pgfpathmoveto{\pgfqpoint{-0.000000in}{0.000000in}}%
\pgfpathlineto{\pgfqpoint{-0.041667in}{0.000000in}}%
\pgfusepath{stroke,fill}%
}%
\begin{pgfscope}%
\pgfsys@transformshift{3.038110in}{1.417244in}%
\pgfsys@useobject{currentmarker}{}%
\end{pgfscope}%
\end{pgfscope}%
\begin{pgfscope}%
\definecolor{textcolor}{rgb}{0.000000,0.000000,0.000000}%
\pgfsetstrokecolor{textcolor}%
\pgfsetfillcolor{textcolor}%
\pgftext[x=0.331018in, y=1.364437in, left, base]{\color{textcolor}{\rmfamily\fontsize{11.000000}{13.200000}\selectfont\catcode`\^=\active\def^{\ifmmode\sp\else\^{}\fi}\catcode`\%=\active\def%{\%}1.00}}%
\end{pgfscope}%
\begin{pgfscope}%
\pgfsetbuttcap%
\pgfsetroundjoin%
\definecolor{currentfill}{rgb}{0.000000,0.000000,0.000000}%
\pgfsetfillcolor{currentfill}%
\pgfsetlinewidth{0.501875pt}%
\definecolor{currentstroke}{rgb}{0.000000,0.000000,0.000000}%
\pgfsetstrokecolor{currentstroke}%
\pgfsetdash{}{0pt}%
\pgfsys@defobject{currentmarker}{\pgfqpoint{0.000000in}{0.000000in}}{\pgfqpoint{0.041667in}{0.000000in}}{%
\pgfpathmoveto{\pgfqpoint{0.000000in}{0.000000in}}%
\pgfpathlineto{\pgfqpoint{0.041667in}{0.000000in}}%
\pgfusepath{stroke,fill}%
}%
\begin{pgfscope}%
\pgfsys@transformshift{0.650000in}{1.915866in}%
\pgfsys@useobject{currentmarker}{}%
\end{pgfscope}%
\end{pgfscope}%
\begin{pgfscope}%
\pgfsetbuttcap%
\pgfsetroundjoin%
\definecolor{currentfill}{rgb}{0.000000,0.000000,0.000000}%
\pgfsetfillcolor{currentfill}%
\pgfsetlinewidth{0.501875pt}%
\definecolor{currentstroke}{rgb}{0.000000,0.000000,0.000000}%
\pgfsetstrokecolor{currentstroke}%
\pgfsetdash{}{0pt}%
\pgfsys@defobject{currentmarker}{\pgfqpoint{-0.041667in}{0.000000in}}{\pgfqpoint{-0.000000in}{0.000000in}}{%
\pgfpathmoveto{\pgfqpoint{-0.000000in}{0.000000in}}%
\pgfpathlineto{\pgfqpoint{-0.041667in}{0.000000in}}%
\pgfusepath{stroke,fill}%
}%
\begin{pgfscope}%
\pgfsys@transformshift{3.038110in}{1.915866in}%
\pgfsys@useobject{currentmarker}{}%
\end{pgfscope}%
\end{pgfscope}%
\begin{pgfscope}%
\definecolor{textcolor}{rgb}{0.000000,0.000000,0.000000}%
\pgfsetstrokecolor{textcolor}%
\pgfsetfillcolor{textcolor}%
\pgftext[x=0.331018in, y=1.863059in, left, base]{\color{textcolor}{\rmfamily\fontsize{11.000000}{13.200000}\selectfont\catcode`\^=\active\def^{\ifmmode\sp\else\^{}\fi}\catcode`\%=\active\def%{\%}1.05}}%
\end{pgfscope}%
\begin{pgfscope}%
\pgfsetbuttcap%
\pgfsetroundjoin%
\definecolor{currentfill}{rgb}{0.000000,0.000000,0.000000}%
\pgfsetfillcolor{currentfill}%
\pgfsetlinewidth{0.501875pt}%
\definecolor{currentstroke}{rgb}{0.000000,0.000000,0.000000}%
\pgfsetstrokecolor{currentstroke}%
\pgfsetdash{}{0pt}%
\pgfsys@defobject{currentmarker}{\pgfqpoint{0.000000in}{0.000000in}}{\pgfqpoint{0.041667in}{0.000000in}}{%
\pgfpathmoveto{\pgfqpoint{0.000000in}{0.000000in}}%
\pgfpathlineto{\pgfqpoint{0.041667in}{0.000000in}}%
\pgfusepath{stroke,fill}%
}%
\begin{pgfscope}%
\pgfsys@transformshift{0.650000in}{2.414488in}%
\pgfsys@useobject{currentmarker}{}%
\end{pgfscope}%
\end{pgfscope}%
\begin{pgfscope}%
\pgfsetbuttcap%
\pgfsetroundjoin%
\definecolor{currentfill}{rgb}{0.000000,0.000000,0.000000}%
\pgfsetfillcolor{currentfill}%
\pgfsetlinewidth{0.501875pt}%
\definecolor{currentstroke}{rgb}{0.000000,0.000000,0.000000}%
\pgfsetstrokecolor{currentstroke}%
\pgfsetdash{}{0pt}%
\pgfsys@defobject{currentmarker}{\pgfqpoint{-0.041667in}{0.000000in}}{\pgfqpoint{-0.000000in}{0.000000in}}{%
\pgfpathmoveto{\pgfqpoint{-0.000000in}{0.000000in}}%
\pgfpathlineto{\pgfqpoint{-0.041667in}{0.000000in}}%
\pgfusepath{stroke,fill}%
}%
\begin{pgfscope}%
\pgfsys@transformshift{3.038110in}{2.414488in}%
\pgfsys@useobject{currentmarker}{}%
\end{pgfscope}%
\end{pgfscope}%
\begin{pgfscope}%
\definecolor{textcolor}{rgb}{0.000000,0.000000,0.000000}%
\pgfsetstrokecolor{textcolor}%
\pgfsetfillcolor{textcolor}%
\pgftext[x=0.331018in, y=2.361681in, left, base]{\color{textcolor}{\rmfamily\fontsize{11.000000}{13.200000}\selectfont\catcode`\^=\active\def^{\ifmmode\sp\else\^{}\fi}\catcode`\%=\active\def%{\%}1.10}}%
\end{pgfscope}%
\begin{pgfscope}%
\pgfsetbuttcap%
\pgfsetroundjoin%
\definecolor{currentfill}{rgb}{0.000000,0.000000,0.000000}%
\pgfsetfillcolor{currentfill}%
\pgfsetlinewidth{0.401500pt}%
\definecolor{currentstroke}{rgb}{0.000000,0.000000,0.000000}%
\pgfsetstrokecolor{currentstroke}%
\pgfsetdash{}{0pt}%
\pgfsys@defobject{currentmarker}{\pgfqpoint{0.000000in}{0.000000in}}{\pgfqpoint{0.020833in}{0.000000in}}{%
\pgfpathmoveto{\pgfqpoint{0.000000in}{0.000000in}}%
\pgfpathlineto{\pgfqpoint{0.020833in}{0.000000in}}%
\pgfusepath{stroke,fill}%
}%
\begin{pgfscope}%
\pgfsys@transformshift{0.650000in}{0.519724in}%
\pgfsys@useobject{currentmarker}{}%
\end{pgfscope}%
\end{pgfscope}%
\begin{pgfscope}%
\pgfsetbuttcap%
\pgfsetroundjoin%
\definecolor{currentfill}{rgb}{0.000000,0.000000,0.000000}%
\pgfsetfillcolor{currentfill}%
\pgfsetlinewidth{0.401500pt}%
\definecolor{currentstroke}{rgb}{0.000000,0.000000,0.000000}%
\pgfsetstrokecolor{currentstroke}%
\pgfsetdash{}{0pt}%
\pgfsys@defobject{currentmarker}{\pgfqpoint{-0.020833in}{0.000000in}}{\pgfqpoint{-0.000000in}{0.000000in}}{%
\pgfpathmoveto{\pgfqpoint{-0.000000in}{0.000000in}}%
\pgfpathlineto{\pgfqpoint{-0.020833in}{0.000000in}}%
\pgfusepath{stroke,fill}%
}%
\begin{pgfscope}%
\pgfsys@transformshift{3.038110in}{0.519724in}%
\pgfsys@useobject{currentmarker}{}%
\end{pgfscope}%
\end{pgfscope}%
\begin{pgfscope}%
\pgfsetbuttcap%
\pgfsetroundjoin%
\definecolor{currentfill}{rgb}{0.000000,0.000000,0.000000}%
\pgfsetfillcolor{currentfill}%
\pgfsetlinewidth{0.401500pt}%
\definecolor{currentstroke}{rgb}{0.000000,0.000000,0.000000}%
\pgfsetstrokecolor{currentstroke}%
\pgfsetdash{}{0pt}%
\pgfsys@defobject{currentmarker}{\pgfqpoint{0.000000in}{0.000000in}}{\pgfqpoint{0.020833in}{0.000000in}}{%
\pgfpathmoveto{\pgfqpoint{0.000000in}{0.000000in}}%
\pgfpathlineto{\pgfqpoint{0.020833in}{0.000000in}}%
\pgfusepath{stroke,fill}%
}%
\begin{pgfscope}%
\pgfsys@transformshift{0.650000in}{0.619449in}%
\pgfsys@useobject{currentmarker}{}%
\end{pgfscope}%
\end{pgfscope}%
\begin{pgfscope}%
\pgfsetbuttcap%
\pgfsetroundjoin%
\definecolor{currentfill}{rgb}{0.000000,0.000000,0.000000}%
\pgfsetfillcolor{currentfill}%
\pgfsetlinewidth{0.401500pt}%
\definecolor{currentstroke}{rgb}{0.000000,0.000000,0.000000}%
\pgfsetstrokecolor{currentstroke}%
\pgfsetdash{}{0pt}%
\pgfsys@defobject{currentmarker}{\pgfqpoint{-0.020833in}{0.000000in}}{\pgfqpoint{-0.000000in}{0.000000in}}{%
\pgfpathmoveto{\pgfqpoint{-0.000000in}{0.000000in}}%
\pgfpathlineto{\pgfqpoint{-0.020833in}{0.000000in}}%
\pgfusepath{stroke,fill}%
}%
\begin{pgfscope}%
\pgfsys@transformshift{3.038110in}{0.619449in}%
\pgfsys@useobject{currentmarker}{}%
\end{pgfscope}%
\end{pgfscope}%
\begin{pgfscope}%
\pgfsetbuttcap%
\pgfsetroundjoin%
\definecolor{currentfill}{rgb}{0.000000,0.000000,0.000000}%
\pgfsetfillcolor{currentfill}%
\pgfsetlinewidth{0.401500pt}%
\definecolor{currentstroke}{rgb}{0.000000,0.000000,0.000000}%
\pgfsetstrokecolor{currentstroke}%
\pgfsetdash{}{0pt}%
\pgfsys@defobject{currentmarker}{\pgfqpoint{0.000000in}{0.000000in}}{\pgfqpoint{0.020833in}{0.000000in}}{%
\pgfpathmoveto{\pgfqpoint{0.000000in}{0.000000in}}%
\pgfpathlineto{\pgfqpoint{0.020833in}{0.000000in}}%
\pgfusepath{stroke,fill}%
}%
\begin{pgfscope}%
\pgfsys@transformshift{0.650000in}{0.719173in}%
\pgfsys@useobject{currentmarker}{}%
\end{pgfscope}%
\end{pgfscope}%
\begin{pgfscope}%
\pgfsetbuttcap%
\pgfsetroundjoin%
\definecolor{currentfill}{rgb}{0.000000,0.000000,0.000000}%
\pgfsetfillcolor{currentfill}%
\pgfsetlinewidth{0.401500pt}%
\definecolor{currentstroke}{rgb}{0.000000,0.000000,0.000000}%
\pgfsetstrokecolor{currentstroke}%
\pgfsetdash{}{0pt}%
\pgfsys@defobject{currentmarker}{\pgfqpoint{-0.020833in}{0.000000in}}{\pgfqpoint{-0.000000in}{0.000000in}}{%
\pgfpathmoveto{\pgfqpoint{-0.000000in}{0.000000in}}%
\pgfpathlineto{\pgfqpoint{-0.020833in}{0.000000in}}%
\pgfusepath{stroke,fill}%
}%
\begin{pgfscope}%
\pgfsys@transformshift{3.038110in}{0.719173in}%
\pgfsys@useobject{currentmarker}{}%
\end{pgfscope}%
\end{pgfscope}%
\begin{pgfscope}%
\pgfsetbuttcap%
\pgfsetroundjoin%
\definecolor{currentfill}{rgb}{0.000000,0.000000,0.000000}%
\pgfsetfillcolor{currentfill}%
\pgfsetlinewidth{0.401500pt}%
\definecolor{currentstroke}{rgb}{0.000000,0.000000,0.000000}%
\pgfsetstrokecolor{currentstroke}%
\pgfsetdash{}{0pt}%
\pgfsys@defobject{currentmarker}{\pgfqpoint{0.000000in}{0.000000in}}{\pgfqpoint{0.020833in}{0.000000in}}{%
\pgfpathmoveto{\pgfqpoint{0.000000in}{0.000000in}}%
\pgfpathlineto{\pgfqpoint{0.020833in}{0.000000in}}%
\pgfusepath{stroke,fill}%
}%
\begin{pgfscope}%
\pgfsys@transformshift{0.650000in}{0.818898in}%
\pgfsys@useobject{currentmarker}{}%
\end{pgfscope}%
\end{pgfscope}%
\begin{pgfscope}%
\pgfsetbuttcap%
\pgfsetroundjoin%
\definecolor{currentfill}{rgb}{0.000000,0.000000,0.000000}%
\pgfsetfillcolor{currentfill}%
\pgfsetlinewidth{0.401500pt}%
\definecolor{currentstroke}{rgb}{0.000000,0.000000,0.000000}%
\pgfsetstrokecolor{currentstroke}%
\pgfsetdash{}{0pt}%
\pgfsys@defobject{currentmarker}{\pgfqpoint{-0.020833in}{0.000000in}}{\pgfqpoint{-0.000000in}{0.000000in}}{%
\pgfpathmoveto{\pgfqpoint{-0.000000in}{0.000000in}}%
\pgfpathlineto{\pgfqpoint{-0.020833in}{0.000000in}}%
\pgfusepath{stroke,fill}%
}%
\begin{pgfscope}%
\pgfsys@transformshift{3.038110in}{0.818898in}%
\pgfsys@useobject{currentmarker}{}%
\end{pgfscope}%
\end{pgfscope}%
\begin{pgfscope}%
\pgfsetbuttcap%
\pgfsetroundjoin%
\definecolor{currentfill}{rgb}{0.000000,0.000000,0.000000}%
\pgfsetfillcolor{currentfill}%
\pgfsetlinewidth{0.401500pt}%
\definecolor{currentstroke}{rgb}{0.000000,0.000000,0.000000}%
\pgfsetstrokecolor{currentstroke}%
\pgfsetdash{}{0pt}%
\pgfsys@defobject{currentmarker}{\pgfqpoint{0.000000in}{0.000000in}}{\pgfqpoint{0.020833in}{0.000000in}}{%
\pgfpathmoveto{\pgfqpoint{0.000000in}{0.000000in}}%
\pgfpathlineto{\pgfqpoint{0.020833in}{0.000000in}}%
\pgfusepath{stroke,fill}%
}%
\begin{pgfscope}%
\pgfsys@transformshift{0.650000in}{1.018346in}%
\pgfsys@useobject{currentmarker}{}%
\end{pgfscope}%
\end{pgfscope}%
\begin{pgfscope}%
\pgfsetbuttcap%
\pgfsetroundjoin%
\definecolor{currentfill}{rgb}{0.000000,0.000000,0.000000}%
\pgfsetfillcolor{currentfill}%
\pgfsetlinewidth{0.401500pt}%
\definecolor{currentstroke}{rgb}{0.000000,0.000000,0.000000}%
\pgfsetstrokecolor{currentstroke}%
\pgfsetdash{}{0pt}%
\pgfsys@defobject{currentmarker}{\pgfqpoint{-0.020833in}{0.000000in}}{\pgfqpoint{-0.000000in}{0.000000in}}{%
\pgfpathmoveto{\pgfqpoint{-0.000000in}{0.000000in}}%
\pgfpathlineto{\pgfqpoint{-0.020833in}{0.000000in}}%
\pgfusepath{stroke,fill}%
}%
\begin{pgfscope}%
\pgfsys@transformshift{3.038110in}{1.018346in}%
\pgfsys@useobject{currentmarker}{}%
\end{pgfscope}%
\end{pgfscope}%
\begin{pgfscope}%
\pgfsetbuttcap%
\pgfsetroundjoin%
\definecolor{currentfill}{rgb}{0.000000,0.000000,0.000000}%
\pgfsetfillcolor{currentfill}%
\pgfsetlinewidth{0.401500pt}%
\definecolor{currentstroke}{rgb}{0.000000,0.000000,0.000000}%
\pgfsetstrokecolor{currentstroke}%
\pgfsetdash{}{0pt}%
\pgfsys@defobject{currentmarker}{\pgfqpoint{0.000000in}{0.000000in}}{\pgfqpoint{0.020833in}{0.000000in}}{%
\pgfpathmoveto{\pgfqpoint{0.000000in}{0.000000in}}%
\pgfpathlineto{\pgfqpoint{0.020833in}{0.000000in}}%
\pgfusepath{stroke,fill}%
}%
\begin{pgfscope}%
\pgfsys@transformshift{0.650000in}{1.118071in}%
\pgfsys@useobject{currentmarker}{}%
\end{pgfscope}%
\end{pgfscope}%
\begin{pgfscope}%
\pgfsetbuttcap%
\pgfsetroundjoin%
\definecolor{currentfill}{rgb}{0.000000,0.000000,0.000000}%
\pgfsetfillcolor{currentfill}%
\pgfsetlinewidth{0.401500pt}%
\definecolor{currentstroke}{rgb}{0.000000,0.000000,0.000000}%
\pgfsetstrokecolor{currentstroke}%
\pgfsetdash{}{0pt}%
\pgfsys@defobject{currentmarker}{\pgfqpoint{-0.020833in}{0.000000in}}{\pgfqpoint{-0.000000in}{0.000000in}}{%
\pgfpathmoveto{\pgfqpoint{-0.000000in}{0.000000in}}%
\pgfpathlineto{\pgfqpoint{-0.020833in}{0.000000in}}%
\pgfusepath{stroke,fill}%
}%
\begin{pgfscope}%
\pgfsys@transformshift{3.038110in}{1.118071in}%
\pgfsys@useobject{currentmarker}{}%
\end{pgfscope}%
\end{pgfscope}%
\begin{pgfscope}%
\pgfsetbuttcap%
\pgfsetroundjoin%
\definecolor{currentfill}{rgb}{0.000000,0.000000,0.000000}%
\pgfsetfillcolor{currentfill}%
\pgfsetlinewidth{0.401500pt}%
\definecolor{currentstroke}{rgb}{0.000000,0.000000,0.000000}%
\pgfsetstrokecolor{currentstroke}%
\pgfsetdash{}{0pt}%
\pgfsys@defobject{currentmarker}{\pgfqpoint{0.000000in}{0.000000in}}{\pgfqpoint{0.020833in}{0.000000in}}{%
\pgfpathmoveto{\pgfqpoint{0.000000in}{0.000000in}}%
\pgfpathlineto{\pgfqpoint{0.020833in}{0.000000in}}%
\pgfusepath{stroke,fill}%
}%
\begin{pgfscope}%
\pgfsys@transformshift{0.650000in}{1.217795in}%
\pgfsys@useobject{currentmarker}{}%
\end{pgfscope}%
\end{pgfscope}%
\begin{pgfscope}%
\pgfsetbuttcap%
\pgfsetroundjoin%
\definecolor{currentfill}{rgb}{0.000000,0.000000,0.000000}%
\pgfsetfillcolor{currentfill}%
\pgfsetlinewidth{0.401500pt}%
\definecolor{currentstroke}{rgb}{0.000000,0.000000,0.000000}%
\pgfsetstrokecolor{currentstroke}%
\pgfsetdash{}{0pt}%
\pgfsys@defobject{currentmarker}{\pgfqpoint{-0.020833in}{0.000000in}}{\pgfqpoint{-0.000000in}{0.000000in}}{%
\pgfpathmoveto{\pgfqpoint{-0.000000in}{0.000000in}}%
\pgfpathlineto{\pgfqpoint{-0.020833in}{0.000000in}}%
\pgfusepath{stroke,fill}%
}%
\begin{pgfscope}%
\pgfsys@transformshift{3.038110in}{1.217795in}%
\pgfsys@useobject{currentmarker}{}%
\end{pgfscope}%
\end{pgfscope}%
\begin{pgfscope}%
\pgfsetbuttcap%
\pgfsetroundjoin%
\definecolor{currentfill}{rgb}{0.000000,0.000000,0.000000}%
\pgfsetfillcolor{currentfill}%
\pgfsetlinewidth{0.401500pt}%
\definecolor{currentstroke}{rgb}{0.000000,0.000000,0.000000}%
\pgfsetstrokecolor{currentstroke}%
\pgfsetdash{}{0pt}%
\pgfsys@defobject{currentmarker}{\pgfqpoint{0.000000in}{0.000000in}}{\pgfqpoint{0.020833in}{0.000000in}}{%
\pgfpathmoveto{\pgfqpoint{0.000000in}{0.000000in}}%
\pgfpathlineto{\pgfqpoint{0.020833in}{0.000000in}}%
\pgfusepath{stroke,fill}%
}%
\begin{pgfscope}%
\pgfsys@transformshift{0.650000in}{1.317520in}%
\pgfsys@useobject{currentmarker}{}%
\end{pgfscope}%
\end{pgfscope}%
\begin{pgfscope}%
\pgfsetbuttcap%
\pgfsetroundjoin%
\definecolor{currentfill}{rgb}{0.000000,0.000000,0.000000}%
\pgfsetfillcolor{currentfill}%
\pgfsetlinewidth{0.401500pt}%
\definecolor{currentstroke}{rgb}{0.000000,0.000000,0.000000}%
\pgfsetstrokecolor{currentstroke}%
\pgfsetdash{}{0pt}%
\pgfsys@defobject{currentmarker}{\pgfqpoint{-0.020833in}{0.000000in}}{\pgfqpoint{-0.000000in}{0.000000in}}{%
\pgfpathmoveto{\pgfqpoint{-0.000000in}{0.000000in}}%
\pgfpathlineto{\pgfqpoint{-0.020833in}{0.000000in}}%
\pgfusepath{stroke,fill}%
}%
\begin{pgfscope}%
\pgfsys@transformshift{3.038110in}{1.317520in}%
\pgfsys@useobject{currentmarker}{}%
\end{pgfscope}%
\end{pgfscope}%
\begin{pgfscope}%
\pgfsetbuttcap%
\pgfsetroundjoin%
\definecolor{currentfill}{rgb}{0.000000,0.000000,0.000000}%
\pgfsetfillcolor{currentfill}%
\pgfsetlinewidth{0.401500pt}%
\definecolor{currentstroke}{rgb}{0.000000,0.000000,0.000000}%
\pgfsetstrokecolor{currentstroke}%
\pgfsetdash{}{0pt}%
\pgfsys@defobject{currentmarker}{\pgfqpoint{0.000000in}{0.000000in}}{\pgfqpoint{0.020833in}{0.000000in}}{%
\pgfpathmoveto{\pgfqpoint{0.000000in}{0.000000in}}%
\pgfpathlineto{\pgfqpoint{0.020833in}{0.000000in}}%
\pgfusepath{stroke,fill}%
}%
\begin{pgfscope}%
\pgfsys@transformshift{0.650000in}{1.516968in}%
\pgfsys@useobject{currentmarker}{}%
\end{pgfscope}%
\end{pgfscope}%
\begin{pgfscope}%
\pgfsetbuttcap%
\pgfsetroundjoin%
\definecolor{currentfill}{rgb}{0.000000,0.000000,0.000000}%
\pgfsetfillcolor{currentfill}%
\pgfsetlinewidth{0.401500pt}%
\definecolor{currentstroke}{rgb}{0.000000,0.000000,0.000000}%
\pgfsetstrokecolor{currentstroke}%
\pgfsetdash{}{0pt}%
\pgfsys@defobject{currentmarker}{\pgfqpoint{-0.020833in}{0.000000in}}{\pgfqpoint{-0.000000in}{0.000000in}}{%
\pgfpathmoveto{\pgfqpoint{-0.000000in}{0.000000in}}%
\pgfpathlineto{\pgfqpoint{-0.020833in}{0.000000in}}%
\pgfusepath{stroke,fill}%
}%
\begin{pgfscope}%
\pgfsys@transformshift{3.038110in}{1.516968in}%
\pgfsys@useobject{currentmarker}{}%
\end{pgfscope}%
\end{pgfscope}%
\begin{pgfscope}%
\pgfsetbuttcap%
\pgfsetroundjoin%
\definecolor{currentfill}{rgb}{0.000000,0.000000,0.000000}%
\pgfsetfillcolor{currentfill}%
\pgfsetlinewidth{0.401500pt}%
\definecolor{currentstroke}{rgb}{0.000000,0.000000,0.000000}%
\pgfsetstrokecolor{currentstroke}%
\pgfsetdash{}{0pt}%
\pgfsys@defobject{currentmarker}{\pgfqpoint{0.000000in}{0.000000in}}{\pgfqpoint{0.020833in}{0.000000in}}{%
\pgfpathmoveto{\pgfqpoint{0.000000in}{0.000000in}}%
\pgfpathlineto{\pgfqpoint{0.020833in}{0.000000in}}%
\pgfusepath{stroke,fill}%
}%
\begin{pgfscope}%
\pgfsys@transformshift{0.650000in}{1.616693in}%
\pgfsys@useobject{currentmarker}{}%
\end{pgfscope}%
\end{pgfscope}%
\begin{pgfscope}%
\pgfsetbuttcap%
\pgfsetroundjoin%
\definecolor{currentfill}{rgb}{0.000000,0.000000,0.000000}%
\pgfsetfillcolor{currentfill}%
\pgfsetlinewidth{0.401500pt}%
\definecolor{currentstroke}{rgb}{0.000000,0.000000,0.000000}%
\pgfsetstrokecolor{currentstroke}%
\pgfsetdash{}{0pt}%
\pgfsys@defobject{currentmarker}{\pgfqpoint{-0.020833in}{0.000000in}}{\pgfqpoint{-0.000000in}{0.000000in}}{%
\pgfpathmoveto{\pgfqpoint{-0.000000in}{0.000000in}}%
\pgfpathlineto{\pgfqpoint{-0.020833in}{0.000000in}}%
\pgfusepath{stroke,fill}%
}%
\begin{pgfscope}%
\pgfsys@transformshift{3.038110in}{1.616693in}%
\pgfsys@useobject{currentmarker}{}%
\end{pgfscope}%
\end{pgfscope}%
\begin{pgfscope}%
\pgfsetbuttcap%
\pgfsetroundjoin%
\definecolor{currentfill}{rgb}{0.000000,0.000000,0.000000}%
\pgfsetfillcolor{currentfill}%
\pgfsetlinewidth{0.401500pt}%
\definecolor{currentstroke}{rgb}{0.000000,0.000000,0.000000}%
\pgfsetstrokecolor{currentstroke}%
\pgfsetdash{}{0pt}%
\pgfsys@defobject{currentmarker}{\pgfqpoint{0.000000in}{0.000000in}}{\pgfqpoint{0.020833in}{0.000000in}}{%
\pgfpathmoveto{\pgfqpoint{0.000000in}{0.000000in}}%
\pgfpathlineto{\pgfqpoint{0.020833in}{0.000000in}}%
\pgfusepath{stroke,fill}%
}%
\begin{pgfscope}%
\pgfsys@transformshift{0.650000in}{1.716417in}%
\pgfsys@useobject{currentmarker}{}%
\end{pgfscope}%
\end{pgfscope}%
\begin{pgfscope}%
\pgfsetbuttcap%
\pgfsetroundjoin%
\definecolor{currentfill}{rgb}{0.000000,0.000000,0.000000}%
\pgfsetfillcolor{currentfill}%
\pgfsetlinewidth{0.401500pt}%
\definecolor{currentstroke}{rgb}{0.000000,0.000000,0.000000}%
\pgfsetstrokecolor{currentstroke}%
\pgfsetdash{}{0pt}%
\pgfsys@defobject{currentmarker}{\pgfqpoint{-0.020833in}{0.000000in}}{\pgfqpoint{-0.000000in}{0.000000in}}{%
\pgfpathmoveto{\pgfqpoint{-0.000000in}{0.000000in}}%
\pgfpathlineto{\pgfqpoint{-0.020833in}{0.000000in}}%
\pgfusepath{stroke,fill}%
}%
\begin{pgfscope}%
\pgfsys@transformshift{3.038110in}{1.716417in}%
\pgfsys@useobject{currentmarker}{}%
\end{pgfscope}%
\end{pgfscope}%
\begin{pgfscope}%
\pgfsetbuttcap%
\pgfsetroundjoin%
\definecolor{currentfill}{rgb}{0.000000,0.000000,0.000000}%
\pgfsetfillcolor{currentfill}%
\pgfsetlinewidth{0.401500pt}%
\definecolor{currentstroke}{rgb}{0.000000,0.000000,0.000000}%
\pgfsetstrokecolor{currentstroke}%
\pgfsetdash{}{0pt}%
\pgfsys@defobject{currentmarker}{\pgfqpoint{0.000000in}{0.000000in}}{\pgfqpoint{0.020833in}{0.000000in}}{%
\pgfpathmoveto{\pgfqpoint{0.000000in}{0.000000in}}%
\pgfpathlineto{\pgfqpoint{0.020833in}{0.000000in}}%
\pgfusepath{stroke,fill}%
}%
\begin{pgfscope}%
\pgfsys@transformshift{0.650000in}{1.816142in}%
\pgfsys@useobject{currentmarker}{}%
\end{pgfscope}%
\end{pgfscope}%
\begin{pgfscope}%
\pgfsetbuttcap%
\pgfsetroundjoin%
\definecolor{currentfill}{rgb}{0.000000,0.000000,0.000000}%
\pgfsetfillcolor{currentfill}%
\pgfsetlinewidth{0.401500pt}%
\definecolor{currentstroke}{rgb}{0.000000,0.000000,0.000000}%
\pgfsetstrokecolor{currentstroke}%
\pgfsetdash{}{0pt}%
\pgfsys@defobject{currentmarker}{\pgfqpoint{-0.020833in}{0.000000in}}{\pgfqpoint{-0.000000in}{0.000000in}}{%
\pgfpathmoveto{\pgfqpoint{-0.000000in}{0.000000in}}%
\pgfpathlineto{\pgfqpoint{-0.020833in}{0.000000in}}%
\pgfusepath{stroke,fill}%
}%
\begin{pgfscope}%
\pgfsys@transformshift{3.038110in}{1.816142in}%
\pgfsys@useobject{currentmarker}{}%
\end{pgfscope}%
\end{pgfscope}%
\begin{pgfscope}%
\pgfsetbuttcap%
\pgfsetroundjoin%
\definecolor{currentfill}{rgb}{0.000000,0.000000,0.000000}%
\pgfsetfillcolor{currentfill}%
\pgfsetlinewidth{0.401500pt}%
\definecolor{currentstroke}{rgb}{0.000000,0.000000,0.000000}%
\pgfsetstrokecolor{currentstroke}%
\pgfsetdash{}{0pt}%
\pgfsys@defobject{currentmarker}{\pgfqpoint{0.000000in}{0.000000in}}{\pgfqpoint{0.020833in}{0.000000in}}{%
\pgfpathmoveto{\pgfqpoint{0.000000in}{0.000000in}}%
\pgfpathlineto{\pgfqpoint{0.020833in}{0.000000in}}%
\pgfusepath{stroke,fill}%
}%
\begin{pgfscope}%
\pgfsys@transformshift{0.650000in}{2.015590in}%
\pgfsys@useobject{currentmarker}{}%
\end{pgfscope}%
\end{pgfscope}%
\begin{pgfscope}%
\pgfsetbuttcap%
\pgfsetroundjoin%
\definecolor{currentfill}{rgb}{0.000000,0.000000,0.000000}%
\pgfsetfillcolor{currentfill}%
\pgfsetlinewidth{0.401500pt}%
\definecolor{currentstroke}{rgb}{0.000000,0.000000,0.000000}%
\pgfsetstrokecolor{currentstroke}%
\pgfsetdash{}{0pt}%
\pgfsys@defobject{currentmarker}{\pgfqpoint{-0.020833in}{0.000000in}}{\pgfqpoint{-0.000000in}{0.000000in}}{%
\pgfpathmoveto{\pgfqpoint{-0.000000in}{0.000000in}}%
\pgfpathlineto{\pgfqpoint{-0.020833in}{0.000000in}}%
\pgfusepath{stroke,fill}%
}%
\begin{pgfscope}%
\pgfsys@transformshift{3.038110in}{2.015590in}%
\pgfsys@useobject{currentmarker}{}%
\end{pgfscope}%
\end{pgfscope}%
\begin{pgfscope}%
\pgfsetbuttcap%
\pgfsetroundjoin%
\definecolor{currentfill}{rgb}{0.000000,0.000000,0.000000}%
\pgfsetfillcolor{currentfill}%
\pgfsetlinewidth{0.401500pt}%
\definecolor{currentstroke}{rgb}{0.000000,0.000000,0.000000}%
\pgfsetstrokecolor{currentstroke}%
\pgfsetdash{}{0pt}%
\pgfsys@defobject{currentmarker}{\pgfqpoint{0.000000in}{0.000000in}}{\pgfqpoint{0.020833in}{0.000000in}}{%
\pgfpathmoveto{\pgfqpoint{0.000000in}{0.000000in}}%
\pgfpathlineto{\pgfqpoint{0.020833in}{0.000000in}}%
\pgfusepath{stroke,fill}%
}%
\begin{pgfscope}%
\pgfsys@transformshift{0.650000in}{2.115315in}%
\pgfsys@useobject{currentmarker}{}%
\end{pgfscope}%
\end{pgfscope}%
\begin{pgfscope}%
\pgfsetbuttcap%
\pgfsetroundjoin%
\definecolor{currentfill}{rgb}{0.000000,0.000000,0.000000}%
\pgfsetfillcolor{currentfill}%
\pgfsetlinewidth{0.401500pt}%
\definecolor{currentstroke}{rgb}{0.000000,0.000000,0.000000}%
\pgfsetstrokecolor{currentstroke}%
\pgfsetdash{}{0pt}%
\pgfsys@defobject{currentmarker}{\pgfqpoint{-0.020833in}{0.000000in}}{\pgfqpoint{-0.000000in}{0.000000in}}{%
\pgfpathmoveto{\pgfqpoint{-0.000000in}{0.000000in}}%
\pgfpathlineto{\pgfqpoint{-0.020833in}{0.000000in}}%
\pgfusepath{stroke,fill}%
}%
\begin{pgfscope}%
\pgfsys@transformshift{3.038110in}{2.115315in}%
\pgfsys@useobject{currentmarker}{}%
\end{pgfscope}%
\end{pgfscope}%
\begin{pgfscope}%
\pgfsetbuttcap%
\pgfsetroundjoin%
\definecolor{currentfill}{rgb}{0.000000,0.000000,0.000000}%
\pgfsetfillcolor{currentfill}%
\pgfsetlinewidth{0.401500pt}%
\definecolor{currentstroke}{rgb}{0.000000,0.000000,0.000000}%
\pgfsetstrokecolor{currentstroke}%
\pgfsetdash{}{0pt}%
\pgfsys@defobject{currentmarker}{\pgfqpoint{0.000000in}{0.000000in}}{\pgfqpoint{0.020833in}{0.000000in}}{%
\pgfpathmoveto{\pgfqpoint{0.000000in}{0.000000in}}%
\pgfpathlineto{\pgfqpoint{0.020833in}{0.000000in}}%
\pgfusepath{stroke,fill}%
}%
\begin{pgfscope}%
\pgfsys@transformshift{0.650000in}{2.215039in}%
\pgfsys@useobject{currentmarker}{}%
\end{pgfscope}%
\end{pgfscope}%
\begin{pgfscope}%
\pgfsetbuttcap%
\pgfsetroundjoin%
\definecolor{currentfill}{rgb}{0.000000,0.000000,0.000000}%
\pgfsetfillcolor{currentfill}%
\pgfsetlinewidth{0.401500pt}%
\definecolor{currentstroke}{rgb}{0.000000,0.000000,0.000000}%
\pgfsetstrokecolor{currentstroke}%
\pgfsetdash{}{0pt}%
\pgfsys@defobject{currentmarker}{\pgfqpoint{-0.020833in}{0.000000in}}{\pgfqpoint{-0.000000in}{0.000000in}}{%
\pgfpathmoveto{\pgfqpoint{-0.000000in}{0.000000in}}%
\pgfpathlineto{\pgfqpoint{-0.020833in}{0.000000in}}%
\pgfusepath{stroke,fill}%
}%
\begin{pgfscope}%
\pgfsys@transformshift{3.038110in}{2.215039in}%
\pgfsys@useobject{currentmarker}{}%
\end{pgfscope}%
\end{pgfscope}%
\begin{pgfscope}%
\pgfsetbuttcap%
\pgfsetroundjoin%
\definecolor{currentfill}{rgb}{0.000000,0.000000,0.000000}%
\pgfsetfillcolor{currentfill}%
\pgfsetlinewidth{0.401500pt}%
\definecolor{currentstroke}{rgb}{0.000000,0.000000,0.000000}%
\pgfsetstrokecolor{currentstroke}%
\pgfsetdash{}{0pt}%
\pgfsys@defobject{currentmarker}{\pgfqpoint{0.000000in}{0.000000in}}{\pgfqpoint{0.020833in}{0.000000in}}{%
\pgfpathmoveto{\pgfqpoint{0.000000in}{0.000000in}}%
\pgfpathlineto{\pgfqpoint{0.020833in}{0.000000in}}%
\pgfusepath{stroke,fill}%
}%
\begin{pgfscope}%
\pgfsys@transformshift{0.650000in}{2.314764in}%
\pgfsys@useobject{currentmarker}{}%
\end{pgfscope}%
\end{pgfscope}%
\begin{pgfscope}%
\pgfsetbuttcap%
\pgfsetroundjoin%
\definecolor{currentfill}{rgb}{0.000000,0.000000,0.000000}%
\pgfsetfillcolor{currentfill}%
\pgfsetlinewidth{0.401500pt}%
\definecolor{currentstroke}{rgb}{0.000000,0.000000,0.000000}%
\pgfsetstrokecolor{currentstroke}%
\pgfsetdash{}{0pt}%
\pgfsys@defobject{currentmarker}{\pgfqpoint{-0.020833in}{0.000000in}}{\pgfqpoint{-0.000000in}{0.000000in}}{%
\pgfpathmoveto{\pgfqpoint{-0.000000in}{0.000000in}}%
\pgfpathlineto{\pgfqpoint{-0.020833in}{0.000000in}}%
\pgfusepath{stroke,fill}%
}%
\begin{pgfscope}%
\pgfsys@transformshift{3.038110in}{2.314764in}%
\pgfsys@useobject{currentmarker}{}%
\end{pgfscope}%
\end{pgfscope}%
\begin{pgfscope}%
\definecolor{textcolor}{rgb}{0.000000,0.000000,0.000000}%
\pgfsetstrokecolor{textcolor}%
\pgfsetfillcolor{textcolor}%
\pgftext[x=0.303240in,y=1.417244in,,bottom,rotate=90.000000]{\color{textcolor}{\rmfamily\fontsize{12.000000}{14.400000}\selectfont\catcode`\^=\active\def^{\ifmmode\sp\else\^{}\fi}\catcode`\%=\active\def%{\%}$U/U_{0}$}}%
\end{pgfscope}%
\begin{pgfscope}%
\pgfpathrectangle{\pgfqpoint{0.650000in}{0.420000in}}{\pgfqpoint{2.388110in}{1.994488in}}%
\pgfusepath{clip}%
\pgfsetrectcap%
\pgfsetroundjoin%
\pgfsetlinewidth{1.003750pt}%
\definecolor{currentstroke}{rgb}{0.796078,0.094118,0.113725}%
\pgfsetstrokecolor{currentstroke}%
\pgfsetdash{}{0pt}%
\pgfpathmoveto{\pgfqpoint{0.650000in}{1.230245in}}%
\pgfpathlineto{\pgfqpoint{0.657168in}{1.231239in}}%
\pgfpathlineto{\pgfqpoint{0.666129in}{1.232398in}}%
\pgfpathlineto{\pgfqpoint{0.674014in}{1.232303in}}%
\pgfpathlineto{\pgfqpoint{0.685124in}{1.231000in}}%
\pgfpathlineto{\pgfqpoint{0.691217in}{1.230844in}}%
\pgfpathlineto{\pgfqpoint{0.697669in}{1.231728in}}%
\pgfpathlineto{\pgfqpoint{0.704120in}{1.233658in}}%
\pgfpathlineto{\pgfqpoint{0.712005in}{1.237151in}}%
\pgfpathlineto{\pgfqpoint{0.722399in}{1.243025in}}%
\pgfpathlineto{\pgfqpoint{0.741395in}{1.255344in}}%
\pgfpathlineto{\pgfqpoint{0.788347in}{1.286173in}}%
\pgfpathlineto{\pgfqpoint{0.793365in}{1.289339in}}%
\pgfpathlineto{\pgfqpoint{0.834224in}{1.314747in}}%
\pgfpathlineto{\pgfqpoint{0.869348in}{1.335595in}}%
\pgfpathlineto{\pgfqpoint{0.924543in}{1.367115in}}%
\pgfpathlineto{\pgfqpoint{0.985832in}{1.400595in}}%
\pgfpathlineto{\pgfqpoint{1.029200in}{1.423252in}}%
\pgfpathlineto{\pgfqpoint{1.058589in}{1.437596in}}%
\pgfpathlineto{\pgfqpoint{1.063607in}{1.439796in}}%
\pgfpathlineto{\pgfqpoint{1.077944in}{1.445510in}}%
\pgfpathlineto{\pgfqpoint{1.088338in}{1.448660in}}%
\pgfpathlineto{\pgfqpoint{1.096581in}{1.450193in}}%
\pgfpathlineto{\pgfqpoint{1.101957in}{1.450234in}}%
\pgfpathlineto{\pgfqpoint{1.107692in}{1.449472in}}%
\pgfpathlineto{\pgfqpoint{1.112710in}{1.447715in}}%
\pgfpathlineto{\pgfqpoint{1.117369in}{1.444882in}}%
\pgfpathlineto{\pgfqpoint{1.119878in}{1.442274in}}%
\pgfpathlineto{\pgfqpoint{1.124179in}{1.437345in}}%
\pgfpathlineto{\pgfqpoint{1.128480in}{1.430616in}}%
\pgfpathlineto{\pgfqpoint{1.133139in}{1.421001in}}%
\pgfpathlineto{\pgfqpoint{1.143891in}{1.397172in}}%
\pgfpathlineto{\pgfqpoint{1.148551in}{1.390254in}}%
\pgfpathlineto{\pgfqpoint{1.152852in}{1.385751in}}%
\pgfpathlineto{\pgfqpoint{1.157511in}{1.382507in}}%
\pgfpathlineto{\pgfqpoint{1.160378in}{1.381191in}}%
\pgfpathlineto{\pgfqpoint{1.161095in}{1.379070in}}%
\pgfpathlineto{\pgfqpoint{1.166113in}{1.377851in}}%
\pgfpathlineto{\pgfqpoint{1.171848in}{1.377646in}}%
\pgfpathlineto{\pgfqpoint{1.177941in}{1.378431in}}%
\pgfpathlineto{\pgfqpoint{1.178299in}{1.376472in}}%
\pgfpathlineto{\pgfqpoint{1.178658in}{1.376544in}}%
\pgfpathlineto{\pgfqpoint{1.187260in}{1.378884in}}%
\pgfpathlineto{\pgfqpoint{1.194428in}{1.381491in}}%
\pgfpathlineto{\pgfqpoint{1.194786in}{1.379536in}}%
\pgfpathlineto{\pgfqpoint{1.195145in}{1.379678in}}%
\pgfpathlineto{\pgfqpoint{1.209123in}{1.385768in}}%
\pgfpathlineto{\pgfqpoint{1.212707in}{1.387460in}}%
\pgfpathlineto{\pgfqpoint{1.213065in}{1.385514in}}%
\pgfpathlineto{\pgfqpoint{1.213424in}{1.385684in}}%
\pgfpathlineto{\pgfqpoint{1.230269in}{1.394020in}}%
\pgfpathlineto{\pgfqpoint{1.230628in}{1.392072in}}%
\pgfpathlineto{\pgfqpoint{1.230986in}{1.392254in}}%
\pgfpathlineto{\pgfqpoint{1.248549in}{1.401268in}}%
\pgfpathlineto{\pgfqpoint{1.248907in}{1.399320in}}%
\pgfpathlineto{\pgfqpoint{1.249266in}{1.399506in}}%
\pgfpathlineto{\pgfqpoint{1.266111in}{1.408244in}}%
\pgfpathlineto{\pgfqpoint{1.266469in}{1.406293in}}%
\pgfpathlineto{\pgfqpoint{1.266828in}{1.406478in}}%
\pgfpathlineto{\pgfqpoint{1.285465in}{1.416055in}}%
\pgfpathlineto{\pgfqpoint{1.285824in}{1.414103in}}%
\pgfpathlineto{\pgfqpoint{1.286182in}{1.414286in}}%
\pgfpathlineto{\pgfqpoint{1.304103in}{1.423590in}}%
\pgfpathlineto{\pgfqpoint{1.304461in}{1.421642in}}%
\pgfpathlineto{\pgfqpoint{1.304820in}{1.421830in}}%
\pgfpathlineto{\pgfqpoint{1.323816in}{1.431741in}}%
\pgfpathlineto{\pgfqpoint{1.324174in}{1.429784in}}%
\pgfpathlineto{\pgfqpoint{1.324532in}{1.429967in}}%
\pgfpathlineto{\pgfqpoint{1.344245in}{1.439924in}}%
\pgfpathlineto{\pgfqpoint{1.344604in}{1.437962in}}%
\pgfpathlineto{\pgfqpoint{1.344962in}{1.438142in}}%
\pgfpathlineto{\pgfqpoint{1.364675in}{1.447932in}}%
\pgfpathlineto{\pgfqpoint{1.365033in}{1.445968in}}%
\pgfpathlineto{\pgfqpoint{1.365392in}{1.446144in}}%
\pgfpathlineto{\pgfqpoint{1.384746in}{1.455713in}}%
\pgfpathlineto{\pgfqpoint{1.385104in}{1.453743in}}%
\pgfpathlineto{\pgfqpoint{1.385463in}{1.453920in}}%
\pgfpathlineto{\pgfqpoint{1.409118in}{1.465137in}}%
\pgfpathlineto{\pgfqpoint{1.409477in}{1.463153in}}%
\pgfpathlineto{\pgfqpoint{1.409835in}{1.463315in}}%
\pgfpathlineto{\pgfqpoint{1.433132in}{1.474025in}}%
\pgfpathlineto{\pgfqpoint{1.433490in}{1.472051in}}%
\pgfpathlineto{\pgfqpoint{1.433849in}{1.472219in}}%
\pgfpathlineto{\pgfqpoint{1.454995in}{1.482322in}}%
\pgfpathlineto{\pgfqpoint{1.455354in}{1.480345in}}%
\pgfpathlineto{\pgfqpoint{1.455712in}{1.480516in}}%
\pgfpathlineto{\pgfqpoint{1.479009in}{1.491576in}}%
\pgfpathlineto{\pgfqpoint{1.479368in}{1.489589in}}%
\pgfpathlineto{\pgfqpoint{1.479726in}{1.489755in}}%
\pgfpathlineto{\pgfqpoint{1.505173in}{1.501181in}}%
\pgfpathlineto{\pgfqpoint{1.505532in}{1.499186in}}%
\pgfpathlineto{\pgfqpoint{1.505890in}{1.499343in}}%
\pgfpathlineto{\pgfqpoint{1.530621in}{1.510382in}}%
\pgfpathlineto{\pgfqpoint{1.530979in}{1.508389in}}%
\pgfpathlineto{\pgfqpoint{1.531337in}{1.508553in}}%
\pgfpathlineto{\pgfqpoint{1.555351in}{1.519523in}}%
\pgfpathlineto{\pgfqpoint{1.555710in}{1.517524in}}%
\pgfpathlineto{\pgfqpoint{1.556068in}{1.517688in}}%
\pgfpathlineto{\pgfqpoint{1.579365in}{1.528300in}}%
\pgfpathlineto{\pgfqpoint{1.579723in}{1.526300in}}%
\pgfpathlineto{\pgfqpoint{1.580082in}{1.526465in}}%
\pgfpathlineto{\pgfqpoint{1.601587in}{1.536218in}}%
\pgfpathlineto{\pgfqpoint{1.601945in}{1.534212in}}%
\pgfpathlineto{\pgfqpoint{1.602303in}{1.534374in}}%
\pgfpathlineto{\pgfqpoint{1.625242in}{1.544611in}}%
\pgfpathlineto{\pgfqpoint{1.625600in}{1.542595in}}%
\pgfpathlineto{\pgfqpoint{1.625959in}{1.542748in}}%
\pgfpathlineto{\pgfqpoint{1.653198in}{1.554371in}}%
\pgfpathlineto{\pgfqpoint{1.653557in}{1.552351in}}%
\pgfpathlineto{\pgfqpoint{1.653915in}{1.552505in}}%
\pgfpathlineto{\pgfqpoint{1.679721in}{1.563805in}}%
\pgfpathlineto{\pgfqpoint{1.680079in}{1.561787in}}%
\pgfpathlineto{\pgfqpoint{1.680438in}{1.561946in}}%
\pgfpathlineto{\pgfqpoint{1.706602in}{1.573494in}}%
\pgfpathlineto{\pgfqpoint{1.706961in}{1.571470in}}%
\pgfpathlineto{\pgfqpoint{1.707319in}{1.571627in}}%
\pgfpathlineto{\pgfqpoint{1.730616in}{1.581849in}}%
\pgfpathlineto{\pgfqpoint{1.730974in}{1.579816in}}%
\pgfpathlineto{\pgfqpoint{1.731333in}{1.579970in}}%
\pgfpathlineto{\pgfqpoint{1.755705in}{1.590479in}}%
\pgfpathlineto{\pgfqpoint{1.756063in}{1.588435in}}%
\pgfpathlineto{\pgfqpoint{1.756422in}{1.588588in}}%
\pgfpathlineto{\pgfqpoint{1.783662in}{1.600071in}}%
\pgfpathlineto{\pgfqpoint{1.784020in}{1.598019in}}%
\pgfpathlineto{\pgfqpoint{1.784378in}{1.598170in}}%
\pgfpathlineto{\pgfqpoint{1.814127in}{1.610479in}}%
\pgfpathlineto{\pgfqpoint{1.814485in}{1.608419in}}%
\pgfpathlineto{\pgfqpoint{1.814843in}{1.608572in}}%
\pgfpathlineto{\pgfqpoint{1.838499in}{1.618476in}}%
\pgfpathlineto{\pgfqpoint{1.838857in}{1.616398in}}%
\pgfpathlineto{\pgfqpoint{1.839215in}{1.616539in}}%
\pgfpathlineto{\pgfqpoint{1.866455in}{1.627392in}}%
\pgfpathlineto{\pgfqpoint{1.866813in}{1.625298in}}%
\pgfpathlineto{\pgfqpoint{1.867172in}{1.625433in}}%
\pgfpathlineto{\pgfqpoint{1.895486in}{1.636453in}}%
\pgfpathlineto{\pgfqpoint{1.895845in}{1.634359in}}%
\pgfpathlineto{\pgfqpoint{1.896203in}{1.634503in}}%
\pgfpathlineto{\pgfqpoint{1.923442in}{1.645214in}}%
\pgfpathlineto{\pgfqpoint{1.923801in}{1.643098in}}%
\pgfpathlineto{\pgfqpoint{1.924159in}{1.643235in}}%
\pgfpathlineto{\pgfqpoint{1.952832in}{1.653796in}}%
\pgfpathlineto{\pgfqpoint{1.953191in}{1.651657in}}%
\pgfpathlineto{\pgfqpoint{1.953549in}{1.651781in}}%
\pgfpathlineto{\pgfqpoint{1.982222in}{1.661366in}}%
\pgfpathlineto{\pgfqpoint{1.989031in}{1.663671in}}%
\pgfpathlineto{\pgfqpoint{1.989390in}{1.661518in}}%
\pgfpathlineto{\pgfqpoint{1.989748in}{1.661642in}}%
\pgfpathlineto{\pgfqpoint{2.019496in}{1.672463in}}%
\pgfpathlineto{\pgfqpoint{2.019855in}{1.670315in}}%
\pgfpathlineto{\pgfqpoint{2.020213in}{1.670446in}}%
\pgfpathlineto{\pgfqpoint{2.051036in}{1.681886in}}%
\pgfpathlineto{\pgfqpoint{2.051395in}{1.679711in}}%
\pgfpathlineto{\pgfqpoint{2.051753in}{1.679842in}}%
\pgfpathlineto{\pgfqpoint{2.080784in}{1.690542in}}%
\pgfpathlineto{\pgfqpoint{2.081143in}{1.688352in}}%
\pgfpathlineto{\pgfqpoint{2.081501in}{1.688489in}}%
\pgfpathlineto{\pgfqpoint{2.108023in}{1.698424in}}%
\pgfpathlineto{\pgfqpoint{2.108382in}{1.696205in}}%
\pgfpathlineto{\pgfqpoint{2.108740in}{1.696337in}}%
\pgfpathlineto{\pgfqpoint{2.133112in}{1.705315in}}%
\pgfpathlineto{\pgfqpoint{2.133471in}{1.703097in}}%
\pgfpathlineto{\pgfqpoint{2.133829in}{1.703236in}}%
\pgfpathlineto{\pgfqpoint{2.157126in}{1.712093in}}%
\pgfpathlineto{\pgfqpoint{2.157484in}{1.709841in}}%
\pgfpathlineto{\pgfqpoint{2.157842in}{1.709973in}}%
\pgfpathlineto{\pgfqpoint{2.186515in}{1.720856in}}%
\pgfpathlineto{\pgfqpoint{2.186873in}{1.718583in}}%
\pgfpathlineto{\pgfqpoint{2.187231in}{1.718720in}}%
\pgfpathlineto{\pgfqpoint{2.212677in}{1.728417in}}%
\pgfpathlineto{\pgfqpoint{2.213035in}{1.726102in}}%
\pgfpathlineto{\pgfqpoint{2.213394in}{1.726230in}}%
\pgfpathlineto{\pgfqpoint{2.237047in}{1.734261in}}%
\pgfpathlineto{\pgfqpoint{2.246365in}{1.737465in}}%
\pgfpathlineto{\pgfqpoint{2.246724in}{1.735124in}}%
\pgfpathlineto{\pgfqpoint{2.247082in}{1.735247in}}%
\pgfpathlineto{\pgfqpoint{2.277187in}{1.746033in}}%
\pgfpathlineto{\pgfqpoint{2.277545in}{1.743657in}}%
\pgfpathlineto{\pgfqpoint{2.277903in}{1.743789in}}%
\pgfpathlineto{\pgfqpoint{2.304066in}{1.753200in}}%
\pgfpathlineto{\pgfqpoint{2.304424in}{1.750788in}}%
\pgfpathlineto{\pgfqpoint{2.304782in}{1.750920in}}%
\pgfpathlineto{\pgfqpoint{2.332378in}{1.761233in}}%
\pgfpathlineto{\pgfqpoint{2.332737in}{1.758777in}}%
\pgfpathlineto{\pgfqpoint{2.333095in}{1.758909in}}%
\pgfpathlineto{\pgfqpoint{2.355673in}{1.766699in}}%
\pgfpathlineto{\pgfqpoint{2.361049in}{1.768470in}}%
\pgfpathlineto{\pgfqpoint{2.361408in}{1.765955in}}%
\pgfpathlineto{\pgfqpoint{2.361766in}{1.766082in}}%
\pgfpathlineto{\pgfqpoint{2.391512in}{1.775661in}}%
\pgfpathlineto{\pgfqpoint{2.395454in}{1.776949in}}%
\pgfpathlineto{\pgfqpoint{2.395813in}{1.774354in}}%
\pgfpathlineto{\pgfqpoint{2.396171in}{1.774466in}}%
\pgfpathlineto{\pgfqpoint{2.430218in}{1.785540in}}%
\pgfpathlineto{\pgfqpoint{2.430576in}{1.782893in}}%
\pgfpathlineto{\pgfqpoint{2.430935in}{1.783013in}}%
\pgfpathlineto{\pgfqpoint{2.462115in}{1.793546in}}%
\pgfpathlineto{\pgfqpoint{2.462473in}{1.790821in}}%
\pgfpathlineto{\pgfqpoint{2.462832in}{1.790944in}}%
\pgfpathlineto{\pgfqpoint{2.483618in}{1.798489in}}%
\pgfpathlineto{\pgfqpoint{2.488994in}{1.800724in}}%
\pgfpathlineto{\pgfqpoint{2.489352in}{1.797972in}}%
\pgfpathlineto{\pgfqpoint{2.489711in}{1.798126in}}%
\pgfpathlineto{\pgfqpoint{2.503688in}{1.804658in}}%
\pgfpathlineto{\pgfqpoint{2.513007in}{1.810014in}}%
\pgfpathlineto{\pgfqpoint{2.513365in}{1.807249in}}%
\pgfpathlineto{\pgfqpoint{2.513724in}{1.807479in}}%
\pgfpathlineto{\pgfqpoint{2.520533in}{1.812434in}}%
\pgfpathlineto{\pgfqpoint{2.527343in}{1.818857in}}%
\pgfpathlineto{\pgfqpoint{2.533077in}{1.825736in}}%
\pgfpathlineto{\pgfqpoint{2.534152in}{1.827213in}}%
\pgfpathlineto{\pgfqpoint{2.534511in}{1.824705in}}%
\pgfpathlineto{\pgfqpoint{2.534869in}{1.825219in}}%
\pgfpathlineto{\pgfqpoint{2.540245in}{1.834011in}}%
\pgfpathlineto{\pgfqpoint{2.545262in}{1.844501in}}%
\pgfpathlineto{\pgfqpoint{2.548847in}{1.853746in}}%
\pgfpathlineto{\pgfqpoint{2.549205in}{1.851773in}}%
\pgfpathlineto{\pgfqpoint{2.549563in}{1.852821in}}%
\pgfpathlineto{\pgfqpoint{2.554222in}{1.868414in}}%
\pgfpathlineto{\pgfqpoint{2.558523in}{1.886787in}}%
\pgfpathlineto{\pgfqpoint{2.558881in}{1.888536in}}%
\pgfpathlineto{\pgfqpoint{2.559240in}{1.887461in}}%
\pgfpathlineto{\pgfqpoint{2.563541in}{1.912328in}}%
\pgfpathlineto{\pgfqpoint{2.577520in}{2.000546in}}%
\pgfpathlineto{\pgfqpoint{2.582179in}{2.020845in}}%
\pgfpathlineto{\pgfqpoint{2.584330in}{2.026857in}}%
\pgfpathlineto{\pgfqpoint{2.588990in}{2.040840in}}%
\pgfpathlineto{\pgfqpoint{2.593291in}{2.051093in}}%
\pgfpathlineto{\pgfqpoint{2.593649in}{2.050110in}}%
\pgfpathlineto{\pgfqpoint{2.598667in}{2.059474in}}%
\pgfpathlineto{\pgfqpoint{2.604044in}{2.067225in}}%
\pgfpathlineto{\pgfqpoint{2.607270in}{2.071065in}}%
\pgfpathlineto{\pgfqpoint{2.607628in}{2.069551in}}%
\pgfpathlineto{\pgfqpoint{2.607987in}{2.069941in}}%
\pgfpathlineto{\pgfqpoint{2.614080in}{2.075699in}}%
\pgfpathlineto{\pgfqpoint{2.620890in}{2.080715in}}%
\pgfpathlineto{\pgfqpoint{2.625192in}{2.083168in}}%
\pgfpathlineto{\pgfqpoint{2.625550in}{2.081214in}}%
\pgfpathlineto{\pgfqpoint{2.625908in}{2.081402in}}%
\pgfpathlineto{\pgfqpoint{2.635586in}{2.085868in}}%
\pgfpathlineto{\pgfqpoint{2.646697in}{2.090065in}}%
\pgfpathlineto{\pgfqpoint{2.647056in}{2.087704in}}%
\pgfpathlineto{\pgfqpoint{2.647414in}{2.087820in}}%
\pgfpathlineto{\pgfqpoint{2.658168in}{2.090765in}}%
\pgfpathlineto{\pgfqpoint{2.672146in}{2.094083in}}%
\pgfpathlineto{\pgfqpoint{2.672505in}{2.091240in}}%
\pgfpathlineto{\pgfqpoint{2.672863in}{2.091323in}}%
\pgfpathlineto{\pgfqpoint{2.694728in}{2.096132in}}%
\pgfpathlineto{\pgfqpoint{2.695086in}{2.092726in}}%
\pgfpathlineto{\pgfqpoint{2.695444in}{2.092800in}}%
\pgfpathlineto{\pgfqpoint{2.717309in}{2.096952in}}%
\pgfpathlineto{\pgfqpoint{2.717667in}{2.092814in}}%
\pgfpathlineto{\pgfqpoint{2.718026in}{2.092863in}}%
\pgfpathlineto{\pgfqpoint{2.741682in}{2.096229in}}%
\pgfpathlineto{\pgfqpoint{2.742041in}{2.091101in}}%
\pgfpathlineto{\pgfqpoint{2.742399in}{2.091130in}}%
\pgfpathlineto{\pgfqpoint{2.751002in}{2.092430in}}%
\pgfpathlineto{\pgfqpoint{2.764981in}{2.094664in}}%
\pgfpathlineto{\pgfqpoint{2.765339in}{2.088321in}}%
\pgfpathlineto{\pgfqpoint{2.765697in}{2.088372in}}%
\pgfpathlineto{\pgfqpoint{2.776809in}{2.090523in}}%
\pgfpathlineto{\pgfqpoint{2.783261in}{2.092043in}}%
\pgfpathlineto{\pgfqpoint{2.783619in}{2.084212in}}%
\pgfpathlineto{\pgfqpoint{2.783978in}{2.084301in}}%
\pgfpathlineto{\pgfqpoint{2.796164in}{2.087844in}}%
\pgfpathlineto{\pgfqpoint{2.801541in}{2.089952in}}%
\pgfpathlineto{\pgfqpoint{2.801899in}{2.080185in}}%
\pgfpathlineto{\pgfqpoint{2.802258in}{2.080334in}}%
\pgfpathlineto{\pgfqpoint{2.809785in}{2.083415in}}%
\pgfpathlineto{\pgfqpoint{2.809785in}{2.083415in}}%
\pgfusepath{stroke}%
\end{pgfscope}%
\begin{pgfscope}%
\pgfpathrectangle{\pgfqpoint{0.650000in}{0.420000in}}{\pgfqpoint{2.388110in}{1.994488in}}%
\pgfusepath{clip}%
\pgfsetrectcap%
\pgfsetroundjoin%
\pgfsetlinewidth{1.003750pt}%
\definecolor{currentstroke}{rgb}{0.454902,0.678431,0.819608}%
\pgfsetstrokecolor{currentstroke}%
\pgfsetdash{}{0pt}%
\pgfpathmoveto{\pgfqpoint{0.650000in}{1.197600in}}%
\pgfpathlineto{\pgfqpoint{0.657168in}{1.198586in}}%
\pgfpathlineto{\pgfqpoint{0.665770in}{1.199696in}}%
\pgfpathlineto{\pgfqpoint{0.673297in}{1.199645in}}%
\pgfpathlineto{\pgfqpoint{0.682974in}{1.198532in}}%
\pgfpathlineto{\pgfqpoint{0.690859in}{1.198177in}}%
\pgfpathlineto{\pgfqpoint{0.696952in}{1.198959in}}%
\pgfpathlineto{\pgfqpoint{0.703762in}{1.200895in}}%
\pgfpathlineto{\pgfqpoint{0.711647in}{1.204314in}}%
\pgfpathlineto{\pgfqpoint{0.722041in}{1.210123in}}%
\pgfpathlineto{\pgfqpoint{0.741753in}{1.222846in}}%
\pgfpathlineto{\pgfqpoint{0.796590in}{1.258471in}}%
\pgfpathlineto{\pgfqpoint{0.835657in}{1.282401in}}%
\pgfpathlineto{\pgfqpoint{0.868631in}{1.301654in}}%
\pgfpathlineto{\pgfqpoint{0.923468in}{1.332326in}}%
\pgfpathlineto{\pgfqpoint{0.952141in}{1.347697in}}%
\pgfpathlineto{\pgfqpoint{0.970420in}{1.356443in}}%
\pgfpathlineto{\pgfqpoint{0.981172in}{1.360653in}}%
\pgfpathlineto{\pgfqpoint{0.990133in}{1.363156in}}%
\pgfpathlineto{\pgfqpoint{0.995867in}{1.363978in}}%
\pgfpathlineto{\pgfqpoint{0.996943in}{1.363810in}}%
\pgfpathlineto{\pgfqpoint{1.003394in}{1.363530in}}%
\pgfpathlineto{\pgfqpoint{1.009845in}{1.362056in}}%
\pgfpathlineto{\pgfqpoint{1.012713in}{1.360589in}}%
\pgfpathlineto{\pgfqpoint{1.019523in}{1.358782in}}%
\pgfpathlineto{\pgfqpoint{1.023465in}{1.358675in}}%
\pgfpathlineto{\pgfqpoint{1.025616in}{1.359480in}}%
\pgfpathlineto{\pgfqpoint{1.027408in}{1.361188in}}%
\pgfpathlineto{\pgfqpoint{1.029200in}{1.364308in}}%
\pgfpathlineto{\pgfqpoint{1.032067in}{1.371902in}}%
\pgfpathlineto{\pgfqpoint{1.033142in}{1.376461in}}%
\pgfpathlineto{\pgfqpoint{1.033501in}{1.376050in}}%
\pgfpathlineto{\pgfqpoint{1.065399in}{1.340363in}}%
\pgfpathlineto{\pgfqpoint{1.089054in}{1.315454in}}%
\pgfpathlineto{\pgfqpoint{1.110201in}{1.294757in}}%
\pgfpathlineto{\pgfqpoint{1.128838in}{1.278033in}}%
\pgfpathlineto{\pgfqpoint{1.146042in}{1.264074in}}%
\pgfpathlineto{\pgfqpoint{1.156436in}{1.256790in}}%
\pgfpathlineto{\pgfqpoint{1.157870in}{1.257233in}}%
\pgfpathlineto{\pgfqpoint{1.159662in}{1.259012in}}%
\pgfpathlineto{\pgfqpoint{1.161812in}{1.262764in}}%
\pgfpathlineto{\pgfqpoint{1.164679in}{1.270137in}}%
\pgfpathlineto{\pgfqpoint{1.171131in}{1.291455in}}%
\pgfpathlineto{\pgfqpoint{1.177582in}{1.322406in}}%
\pgfpathlineto{\pgfqpoint{1.204822in}{1.459207in}}%
\pgfpathlineto{\pgfqpoint{1.209840in}{1.479877in}}%
\pgfpathlineto{\pgfqpoint{1.210198in}{1.479498in}}%
\pgfpathlineto{\pgfqpoint{1.213065in}{1.488381in}}%
\pgfpathlineto{\pgfqpoint{1.215216in}{1.495766in}}%
\pgfpathlineto{\pgfqpoint{1.215574in}{1.495272in}}%
\pgfpathlineto{\pgfqpoint{1.218083in}{1.503159in}}%
\pgfpathlineto{\pgfqpoint{1.218442in}{1.502588in}}%
\pgfpathlineto{\pgfqpoint{1.221309in}{1.510623in}}%
\pgfpathlineto{\pgfqpoint{1.221667in}{1.509955in}}%
\pgfpathlineto{\pgfqpoint{1.224893in}{1.517733in}}%
\pgfpathlineto{\pgfqpoint{1.225252in}{1.516951in}}%
\pgfpathlineto{\pgfqpoint{1.229194in}{1.524641in}}%
\pgfpathlineto{\pgfqpoint{1.229553in}{1.523707in}}%
\pgfpathlineto{\pgfqpoint{1.229911in}{1.524300in}}%
\pgfpathlineto{\pgfqpoint{1.234212in}{1.530242in}}%
\pgfpathlineto{\pgfqpoint{1.238155in}{1.533778in}}%
\pgfpathlineto{\pgfqpoint{1.241380in}{1.535340in}}%
\pgfpathlineto{\pgfqpoint{1.244606in}{1.535722in}}%
\pgfpathlineto{\pgfqpoint{1.245323in}{1.535650in}}%
\pgfpathlineto{\pgfqpoint{1.246040in}{1.537077in}}%
\pgfpathlineto{\pgfqpoint{1.249266in}{1.535827in}}%
\pgfpathlineto{\pgfqpoint{1.251416in}{1.534384in}}%
\pgfpathlineto{\pgfqpoint{1.251774in}{1.535692in}}%
\pgfpathlineto{\pgfqpoint{1.252133in}{1.535394in}}%
\pgfpathlineto{\pgfqpoint{1.256075in}{1.531289in}}%
\pgfpathlineto{\pgfqpoint{1.256792in}{1.530384in}}%
\pgfpathlineto{\pgfqpoint{1.257151in}{1.531554in}}%
\pgfpathlineto{\pgfqpoint{1.257509in}{1.531074in}}%
\pgfpathlineto{\pgfqpoint{1.265394in}{1.519463in}}%
\pgfpathlineto{\pgfqpoint{1.271487in}{1.506005in}}%
\pgfpathlineto{\pgfqpoint{1.272204in}{1.504264in}}%
\pgfpathlineto{\pgfqpoint{1.273279in}{1.499773in}}%
\pgfpathlineto{\pgfqpoint{1.281164in}{1.478383in}}%
\pgfpathlineto{\pgfqpoint{1.287616in}{1.459306in}}%
\pgfpathlineto{\pgfqpoint{1.288691in}{1.454009in}}%
\pgfpathlineto{\pgfqpoint{1.309838in}{1.391283in}}%
\pgfpathlineto{\pgfqpoint{1.316647in}{1.376162in}}%
\pgfpathlineto{\pgfqpoint{1.322024in}{1.366543in}}%
\pgfpathlineto{\pgfqpoint{1.322740in}{1.362904in}}%
\pgfpathlineto{\pgfqpoint{1.327400in}{1.356863in}}%
\pgfpathlineto{\pgfqpoint{1.331342in}{1.353509in}}%
\pgfpathlineto{\pgfqpoint{1.334568in}{1.352098in}}%
\pgfpathlineto{\pgfqpoint{1.337435in}{1.351920in}}%
\pgfpathlineto{\pgfqpoint{1.340303in}{1.352811in}}%
\pgfpathlineto{\pgfqpoint{1.343529in}{1.355135in}}%
\pgfpathlineto{\pgfqpoint{1.344962in}{1.356627in}}%
\pgfpathlineto{\pgfqpoint{1.345321in}{1.354405in}}%
\pgfpathlineto{\pgfqpoint{1.345679in}{1.354841in}}%
\pgfpathlineto{\pgfqpoint{1.349263in}{1.360184in}}%
\pgfpathlineto{\pgfqpoint{1.353206in}{1.368125in}}%
\pgfpathlineto{\pgfqpoint{1.357865in}{1.380230in}}%
\pgfpathlineto{\pgfqpoint{1.363241in}{1.397602in}}%
\pgfpathlineto{\pgfqpoint{1.363600in}{1.398879in}}%
\pgfpathlineto{\pgfqpoint{1.363958in}{1.397725in}}%
\pgfpathlineto{\pgfqpoint{1.370768in}{1.424743in}}%
\pgfpathlineto{\pgfqpoint{1.382596in}{1.477592in}}%
\pgfpathlineto{\pgfqpoint{1.387613in}{1.497418in}}%
\pgfpathlineto{\pgfqpoint{1.394782in}{1.525098in}}%
\pgfpathlineto{\pgfqpoint{1.400158in}{1.542206in}}%
\pgfpathlineto{\pgfqpoint{1.400516in}{1.541534in}}%
\pgfpathlineto{\pgfqpoint{1.404817in}{1.552522in}}%
\pgfpathlineto{\pgfqpoint{1.405176in}{1.551691in}}%
\pgfpathlineto{\pgfqpoint{1.409835in}{1.560817in}}%
\pgfpathlineto{\pgfqpoint{1.414136in}{1.566987in}}%
\pgfpathlineto{\pgfqpoint{1.414495in}{1.567409in}}%
\pgfpathlineto{\pgfqpoint{1.414853in}{1.566193in}}%
\pgfpathlineto{\pgfqpoint{1.415211in}{1.566587in}}%
\pgfpathlineto{\pgfqpoint{1.418079in}{1.569264in}}%
\pgfpathlineto{\pgfqpoint{1.418796in}{1.571443in}}%
\pgfpathlineto{\pgfqpoint{1.422738in}{1.573591in}}%
\pgfpathlineto{\pgfqpoint{1.426681in}{1.574448in}}%
\pgfpathlineto{\pgfqpoint{1.430981in}{1.574136in}}%
\pgfpathlineto{\pgfqpoint{1.433132in}{1.573573in}}%
\pgfpathlineto{\pgfqpoint{1.433490in}{1.575207in}}%
\pgfpathlineto{\pgfqpoint{1.433849in}{1.575086in}}%
\pgfpathlineto{\pgfqpoint{1.439225in}{1.572624in}}%
\pgfpathlineto{\pgfqpoint{1.446393in}{1.567971in}}%
\pgfpathlineto{\pgfqpoint{1.450336in}{1.565056in}}%
\pgfpathlineto{\pgfqpoint{1.451053in}{1.562593in}}%
\pgfpathlineto{\pgfqpoint{1.465748in}{1.552012in}}%
\pgfpathlineto{\pgfqpoint{1.472558in}{1.548489in}}%
\pgfpathlineto{\pgfqpoint{1.474708in}{1.547680in}}%
\pgfpathlineto{\pgfqpoint{1.475425in}{1.545321in}}%
\pgfpathlineto{\pgfqpoint{1.481160in}{1.544093in}}%
\pgfpathlineto{\pgfqpoint{1.486536in}{1.544154in}}%
\pgfpathlineto{\pgfqpoint{1.491554in}{1.545332in}}%
\pgfpathlineto{\pgfqpoint{1.496930in}{1.547798in}}%
\pgfpathlineto{\pgfqpoint{1.497288in}{1.548004in}}%
\pgfpathlineto{\pgfqpoint{1.497647in}{1.546031in}}%
\pgfpathlineto{\pgfqpoint{1.498005in}{1.546245in}}%
\pgfpathlineto{\pgfqpoint{1.503740in}{1.550442in}}%
\pgfpathlineto{\pgfqpoint{1.510191in}{1.556692in}}%
\pgfpathlineto{\pgfqpoint{1.514134in}{1.561197in}}%
\pgfpathlineto{\pgfqpoint{1.514492in}{1.559487in}}%
\pgfpathlineto{\pgfqpoint{1.514850in}{1.559924in}}%
\pgfpathlineto{\pgfqpoint{1.523094in}{1.570876in}}%
\pgfpathlineto{\pgfqpoint{1.534922in}{1.588766in}}%
\pgfpathlineto{\pgfqpoint{1.538147in}{1.593890in}}%
\pgfpathlineto{\pgfqpoint{1.538506in}{1.592460in}}%
\pgfpathlineto{\pgfqpoint{1.538864in}{1.593034in}}%
\pgfpathlineto{\pgfqpoint{1.559294in}{1.624690in}}%
\pgfpathlineto{\pgfqpoint{1.562878in}{1.629654in}}%
\pgfpathlineto{\pgfqpoint{1.563236in}{1.628267in}}%
\pgfpathlineto{\pgfqpoint{1.563595in}{1.628749in}}%
\pgfpathlineto{\pgfqpoint{1.572197in}{1.639288in}}%
\pgfpathlineto{\pgfqpoint{1.580082in}{1.647185in}}%
\pgfpathlineto{\pgfqpoint{1.585458in}{1.651544in}}%
\pgfpathlineto{\pgfqpoint{1.585816in}{1.649984in}}%
\pgfpathlineto{\pgfqpoint{1.586175in}{1.650239in}}%
\pgfpathlineto{\pgfqpoint{1.593343in}{1.654562in}}%
\pgfpathlineto{\pgfqpoint{1.600153in}{1.657330in}}%
\pgfpathlineto{\pgfqpoint{1.608038in}{1.659229in}}%
\pgfpathlineto{\pgfqpoint{1.611264in}{1.659714in}}%
\pgfpathlineto{\pgfqpoint{1.611981in}{1.657920in}}%
\pgfpathlineto{\pgfqpoint{1.623808in}{1.658685in}}%
\pgfpathlineto{\pgfqpoint{1.632410in}{1.659159in}}%
\pgfpathlineto{\pgfqpoint{1.633127in}{1.657250in}}%
\pgfpathlineto{\pgfqpoint{1.642804in}{1.658609in}}%
\pgfpathlineto{\pgfqpoint{1.650689in}{1.660840in}}%
\pgfpathlineto{\pgfqpoint{1.657499in}{1.663942in}}%
\pgfpathlineto{\pgfqpoint{1.662517in}{1.667093in}}%
\pgfpathlineto{\pgfqpoint{1.662876in}{1.665307in}}%
\pgfpathlineto{\pgfqpoint{1.663234in}{1.665566in}}%
\pgfpathlineto{\pgfqpoint{1.669685in}{1.670984in}}%
\pgfpathlineto{\pgfqpoint{1.676495in}{1.678241in}}%
\pgfpathlineto{\pgfqpoint{1.684380in}{1.688382in}}%
\pgfpathlineto{\pgfqpoint{1.693341in}{1.701886in}}%
\pgfpathlineto{\pgfqpoint{1.704452in}{1.720960in}}%
\pgfpathlineto{\pgfqpoint{1.705527in}{1.722903in}}%
\pgfpathlineto{\pgfqpoint{1.705885in}{1.721691in}}%
\pgfpathlineto{\pgfqpoint{1.706244in}{1.722345in}}%
\pgfpathlineto{\pgfqpoint{1.730974in}{1.766417in}}%
\pgfpathlineto{\pgfqpoint{1.732766in}{1.769242in}}%
\pgfpathlineto{\pgfqpoint{1.733125in}{1.768121in}}%
\pgfpathlineto{\pgfqpoint{1.733483in}{1.768676in}}%
\pgfpathlineto{\pgfqpoint{1.741368in}{1.779830in}}%
\pgfpathlineto{\pgfqpoint{1.748537in}{1.788081in}}%
\pgfpathlineto{\pgfqpoint{1.754630in}{1.793403in}}%
\pgfpathlineto{\pgfqpoint{1.760006in}{1.796717in}}%
\pgfpathlineto{\pgfqpoint{1.765382in}{1.798701in}}%
\pgfpathlineto{\pgfqpoint{1.766457in}{1.798951in}}%
\pgfpathlineto{\pgfqpoint{1.767174in}{1.797483in}}%
\pgfpathlineto{\pgfqpoint{1.772550in}{1.797833in}}%
\pgfpathlineto{\pgfqpoint{1.778285in}{1.796996in}}%
\pgfpathlineto{\pgfqpoint{1.784737in}{1.794840in}}%
\pgfpathlineto{\pgfqpoint{1.791905in}{1.791335in}}%
\pgfpathlineto{\pgfqpoint{1.797639in}{1.787976in}}%
\pgfpathlineto{\pgfqpoint{1.798356in}{1.785778in}}%
\pgfpathlineto{\pgfqpoint{1.814127in}{1.776762in}}%
\pgfpathlineto{\pgfqpoint{1.820936in}{1.774176in}}%
\pgfpathlineto{\pgfqpoint{1.826671in}{1.773180in}}%
\pgfpathlineto{\pgfqpoint{1.827030in}{1.773155in}}%
\pgfpathlineto{\pgfqpoint{1.827746in}{1.771196in}}%
\pgfpathlineto{\pgfqpoint{1.833122in}{1.771560in}}%
\pgfpathlineto{\pgfqpoint{1.838499in}{1.773126in}}%
\pgfpathlineto{\pgfqpoint{1.843516in}{1.775798in}}%
\pgfpathlineto{\pgfqpoint{1.845667in}{1.777323in}}%
\pgfpathlineto{\pgfqpoint{1.846025in}{1.775647in}}%
\pgfpathlineto{\pgfqpoint{1.846384in}{1.775929in}}%
\pgfpathlineto{\pgfqpoint{1.852118in}{1.781226in}}%
\pgfpathlineto{\pgfqpoint{1.857853in}{1.788129in}}%
\pgfpathlineto{\pgfqpoint{1.865021in}{1.798681in}}%
\pgfpathlineto{\pgfqpoint{1.873265in}{1.812862in}}%
\pgfpathlineto{\pgfqpoint{1.877924in}{1.821715in}}%
\pgfpathlineto{\pgfqpoint{1.878283in}{1.820688in}}%
\pgfpathlineto{\pgfqpoint{1.905163in}{1.872224in}}%
\pgfpathlineto{\pgfqpoint{1.911257in}{1.881742in}}%
\pgfpathlineto{\pgfqpoint{1.911615in}{1.880920in}}%
\pgfpathlineto{\pgfqpoint{1.911973in}{1.881415in}}%
\pgfpathlineto{\pgfqpoint{1.917349in}{1.887795in}}%
\pgfpathlineto{\pgfqpoint{1.919500in}{1.889735in}}%
\pgfpathlineto{\pgfqpoint{1.919858in}{1.888798in}}%
\pgfpathlineto{\pgfqpoint{1.920217in}{1.889075in}}%
\pgfpathlineto{\pgfqpoint{1.924159in}{1.891289in}}%
\pgfpathlineto{\pgfqpoint{1.927385in}{1.891899in}}%
\pgfpathlineto{\pgfqpoint{1.930611in}{1.891309in}}%
\pgfpathlineto{\pgfqpoint{1.932044in}{1.890628in}}%
\pgfpathlineto{\pgfqpoint{1.932761in}{1.889029in}}%
\pgfpathlineto{\pgfqpoint{1.935987in}{1.886145in}}%
\pgfpathlineto{\pgfqpoint{1.941005in}{1.879737in}}%
\pgfpathlineto{\pgfqpoint{1.944947in}{1.871144in}}%
\pgfpathlineto{\pgfqpoint{1.949965in}{1.857223in}}%
\pgfpathlineto{\pgfqpoint{1.956416in}{1.835563in}}%
\pgfpathlineto{\pgfqpoint{1.962868in}{1.811388in}}%
\pgfpathlineto{\pgfqpoint{1.963943in}{1.805449in}}%
\pgfpathlineto{\pgfqpoint{1.974695in}{1.766075in}}%
\pgfpathlineto{\pgfqpoint{1.980429in}{1.749088in}}%
\pgfpathlineto{\pgfqpoint{1.984730in}{1.739435in}}%
\pgfpathlineto{\pgfqpoint{1.988315in}{1.733751in}}%
\pgfpathlineto{\pgfqpoint{1.989390in}{1.732485in}}%
\pgfpathlineto{\pgfqpoint{1.990107in}{1.729271in}}%
\pgfpathlineto{\pgfqpoint{1.992974in}{1.727272in}}%
\pgfpathlineto{\pgfqpoint{1.995483in}{1.726806in}}%
\pgfpathlineto{\pgfqpoint{1.997992in}{1.727590in}}%
\pgfpathlineto{\pgfqpoint{2.000501in}{1.729653in}}%
\pgfpathlineto{\pgfqpoint{2.003368in}{1.733537in}}%
\pgfpathlineto{\pgfqpoint{2.006952in}{1.740496in}}%
\pgfpathlineto{\pgfqpoint{2.011253in}{1.751695in}}%
\pgfpathlineto{\pgfqpoint{2.016271in}{1.768224in}}%
\pgfpathlineto{\pgfqpoint{2.022722in}{1.793649in}}%
\pgfpathlineto{\pgfqpoint{2.032041in}{1.835517in}}%
\pgfpathlineto{\pgfqpoint{2.049961in}{1.916584in}}%
\pgfpathlineto{\pgfqpoint{2.057846in}{1.946793in}}%
\pgfpathlineto{\pgfqpoint{2.064656in}{1.968724in}}%
\pgfpathlineto{\pgfqpoint{2.070749in}{1.984858in}}%
\pgfpathlineto{\pgfqpoint{2.076484in}{1.996966in}}%
\pgfpathlineto{\pgfqpoint{2.081860in}{2.005709in}}%
\pgfpathlineto{\pgfqpoint{2.084010in}{2.008548in}}%
\pgfpathlineto{\pgfqpoint{2.084368in}{2.007931in}}%
\pgfpathlineto{\pgfqpoint{2.084727in}{2.008357in}}%
\pgfpathlineto{\pgfqpoint{2.089028in}{2.012610in}}%
\pgfpathlineto{\pgfqpoint{2.093329in}{2.015405in}}%
\pgfpathlineto{\pgfqpoint{2.097271in}{2.016747in}}%
\pgfpathlineto{\pgfqpoint{2.101214in}{2.016969in}}%
\pgfpathlineto{\pgfqpoint{2.105515in}{2.016089in}}%
\pgfpathlineto{\pgfqpoint{2.109815in}{2.014097in}}%
\pgfpathlineto{\pgfqpoint{2.114475in}{2.010740in}}%
\pgfpathlineto{\pgfqpoint{2.119851in}{2.005511in}}%
\pgfpathlineto{\pgfqpoint{2.120926in}{2.004291in}}%
\pgfpathlineto{\pgfqpoint{2.122002in}{2.001826in}}%
\pgfpathlineto{\pgfqpoint{2.128811in}{1.992606in}}%
\pgfpathlineto{\pgfqpoint{2.137772in}{1.978405in}}%
\pgfpathlineto{\pgfqpoint{2.143506in}{1.968400in}}%
\pgfpathlineto{\pgfqpoint{2.144581in}{1.965097in}}%
\pgfpathlineto{\pgfqpoint{2.160710in}{1.935683in}}%
\pgfpathlineto{\pgfqpoint{2.161785in}{1.932197in}}%
\pgfpathlineto{\pgfqpoint{2.172179in}{1.914704in}}%
\pgfpathlineto{\pgfqpoint{2.184723in}{1.896002in}}%
\pgfpathlineto{\pgfqpoint{2.190815in}{1.890063in}}%
\pgfpathlineto{\pgfqpoint{2.196191in}{1.886417in}}%
\pgfpathlineto{\pgfqpoint{2.201208in}{1.884373in}}%
\pgfpathlineto{\pgfqpoint{2.203717in}{1.883860in}}%
\pgfpathlineto{\pgfqpoint{2.204434in}{1.881976in}}%
\pgfpathlineto{\pgfqpoint{2.209452in}{1.881923in}}%
\pgfpathlineto{\pgfqpoint{2.214111in}{1.882906in}}%
\pgfpathlineto{\pgfqpoint{2.219486in}{1.885267in}}%
\pgfpathlineto{\pgfqpoint{2.225579in}{1.889196in}}%
\pgfpathlineto{\pgfqpoint{2.230955in}{1.893625in}}%
\pgfpathlineto{\pgfqpoint{2.231313in}{1.892286in}}%
\pgfpathlineto{\pgfqpoint{2.231672in}{1.892609in}}%
\pgfpathlineto{\pgfqpoint{2.239914in}{1.900851in}}%
\pgfpathlineto{\pgfqpoint{2.254250in}{1.917081in}}%
\pgfpathlineto{\pgfqpoint{2.269661in}{1.935018in}}%
\pgfpathlineto{\pgfqpoint{2.270019in}{1.934010in}}%
\pgfpathlineto{\pgfqpoint{2.270377in}{1.934424in}}%
\pgfpathlineto{\pgfqpoint{2.290089in}{1.956403in}}%
\pgfpathlineto{\pgfqpoint{2.300123in}{1.967017in}}%
\pgfpathlineto{\pgfqpoint{2.300482in}{1.966135in}}%
\pgfpathlineto{\pgfqpoint{2.300840in}{1.966496in}}%
\pgfpathlineto{\pgfqpoint{2.311233in}{1.976198in}}%
\pgfpathlineto{\pgfqpoint{2.319476in}{1.982392in}}%
\pgfpathlineto{\pgfqpoint{2.327002in}{1.986738in}}%
\pgfpathlineto{\pgfqpoint{2.330945in}{1.988473in}}%
\pgfpathlineto{\pgfqpoint{2.331303in}{1.987466in}}%
\pgfpathlineto{\pgfqpoint{2.331661in}{1.987601in}}%
\pgfpathlineto{\pgfqpoint{2.338471in}{1.989495in}}%
\pgfpathlineto{\pgfqpoint{2.345639in}{1.990283in}}%
\pgfpathlineto{\pgfqpoint{2.352806in}{1.989965in}}%
\pgfpathlineto{\pgfqpoint{2.356032in}{1.989489in}}%
\pgfpathlineto{\pgfqpoint{2.356749in}{1.988205in}}%
\pgfpathlineto{\pgfqpoint{2.365708in}{1.985846in}}%
\pgfpathlineto{\pgfqpoint{2.377177in}{1.981508in}}%
\pgfpathlineto{\pgfqpoint{2.377893in}{1.979987in}}%
\pgfpathlineto{\pgfqpoint{2.387570in}{1.976680in}}%
\pgfpathlineto{\pgfqpoint{2.393304in}{1.975938in}}%
\pgfpathlineto{\pgfqpoint{2.399038in}{1.976294in}}%
\pgfpathlineto{\pgfqpoint{2.399755in}{1.975157in}}%
\pgfpathlineto{\pgfqpoint{2.405489in}{1.976877in}}%
\pgfpathlineto{\pgfqpoint{2.410507in}{1.979553in}}%
\pgfpathlineto{\pgfqpoint{2.416241in}{1.983907in}}%
\pgfpathlineto{\pgfqpoint{2.423050in}{1.990411in}}%
\pgfpathlineto{\pgfqpoint{2.425200in}{1.992748in}}%
\pgfpathlineto{\pgfqpoint{2.425559in}{1.992084in}}%
\pgfpathlineto{\pgfqpoint{2.425917in}{1.992494in}}%
\pgfpathlineto{\pgfqpoint{2.438819in}{2.008088in}}%
\pgfpathlineto{\pgfqpoint{2.444912in}{2.014719in}}%
\pgfpathlineto{\pgfqpoint{2.452796in}{2.023135in}}%
\pgfpathlineto{\pgfqpoint{2.456381in}{2.026319in}}%
\pgfpathlineto{\pgfqpoint{2.457097in}{2.026342in}}%
\pgfpathlineto{\pgfqpoint{2.462473in}{2.029770in}}%
\pgfpathlineto{\pgfqpoint{2.467132in}{2.031441in}}%
\pgfpathlineto{\pgfqpoint{2.471075in}{2.031747in}}%
\pgfpathlineto{\pgfqpoint{2.472150in}{2.031309in}}%
\pgfpathlineto{\pgfqpoint{2.476092in}{2.030106in}}%
\pgfpathlineto{\pgfqpoint{2.480034in}{2.027657in}}%
\pgfpathlineto{\pgfqpoint{2.483260in}{2.024429in}}%
\pgfpathlineto{\pgfqpoint{2.487561in}{2.018829in}}%
\pgfpathlineto{\pgfqpoint{2.491503in}{2.011917in}}%
\pgfpathlineto{\pgfqpoint{2.496162in}{2.001584in}}%
\pgfpathlineto{\pgfqpoint{2.501180in}{1.987603in}}%
\pgfpathlineto{\pgfqpoint{2.506556in}{1.969361in}}%
\pgfpathlineto{\pgfqpoint{2.512648in}{1.944452in}}%
\pgfpathlineto{\pgfqpoint{2.522325in}{1.896036in}}%
\pgfpathlineto{\pgfqpoint{2.527701in}{1.863241in}}%
\pgfpathlineto{\pgfqpoint{2.533794in}{1.819746in}}%
\pgfpathlineto{\pgfqpoint{2.540603in}{1.763924in}}%
\pgfpathlineto{\pgfqpoint{2.550997in}{1.671168in}}%
\pgfpathlineto{\pgfqpoint{2.566408in}{1.539294in}}%
\pgfpathlineto{\pgfqpoint{2.571427in}{1.507086in}}%
\pgfpathlineto{\pgfqpoint{2.578595in}{1.474098in}}%
\pgfpathlineto{\pgfqpoint{2.584689in}{1.456388in}}%
\pgfpathlineto{\pgfqpoint{2.586481in}{1.452952in}}%
\pgfpathlineto{\pgfqpoint{2.587556in}{1.449552in}}%
\pgfpathlineto{\pgfqpoint{2.591499in}{1.444265in}}%
\pgfpathlineto{\pgfqpoint{2.591857in}{1.443878in}}%
\pgfpathlineto{\pgfqpoint{2.592574in}{1.441743in}}%
\pgfpathlineto{\pgfqpoint{2.596875in}{1.438505in}}%
\pgfpathlineto{\pgfqpoint{2.597951in}{1.437906in}}%
\pgfpathlineto{\pgfqpoint{2.598667in}{1.436332in}}%
\pgfpathlineto{\pgfqpoint{2.602969in}{1.434919in}}%
\pgfpathlineto{\pgfqpoint{2.604761in}{1.434625in}}%
\pgfpathlineto{\pgfqpoint{2.605478in}{1.433511in}}%
\pgfpathlineto{\pgfqpoint{2.611929in}{1.433604in}}%
\pgfpathlineto{\pgfqpoint{2.612646in}{1.433680in}}%
\pgfpathlineto{\pgfqpoint{2.613363in}{1.432964in}}%
\pgfpathlineto{\pgfqpoint{2.629851in}{1.436254in}}%
\pgfpathlineto{\pgfqpoint{2.642755in}{1.441022in}}%
\pgfpathlineto{\pgfqpoint{2.666770in}{1.451631in}}%
\pgfpathlineto{\pgfqpoint{2.677523in}{1.455580in}}%
\pgfpathlineto{\pgfqpoint{2.696161in}{1.460264in}}%
\pgfpathlineto{\pgfqpoint{2.706914in}{1.462924in}}%
\pgfpathlineto{\pgfqpoint{2.719101in}{1.467390in}}%
\pgfpathlineto{\pgfqpoint{2.726987in}{1.471418in}}%
\pgfpathlineto{\pgfqpoint{2.735589in}{1.476581in}}%
\pgfpathlineto{\pgfqpoint{2.778243in}{1.503769in}}%
\pgfpathlineto{\pgfqpoint{2.798315in}{1.515149in}}%
\pgfpathlineto{\pgfqpoint{2.804408in}{1.520613in}}%
\pgfpathlineto{\pgfqpoint{2.808351in}{1.523443in}}%
\pgfpathlineto{\pgfqpoint{2.809785in}{1.524048in}}%
\pgfpathlineto{\pgfqpoint{2.809785in}{1.524048in}}%
\pgfusepath{stroke}%
\end{pgfscope}%
\begin{pgfscope}%
\pgfpathrectangle{\pgfqpoint{0.650000in}{0.420000in}}{\pgfqpoint{2.388110in}{1.994488in}}%
\pgfusepath{clip}%
\pgfsetrectcap%
\pgfsetroundjoin%
\pgfsetlinewidth{1.003750pt}%
\definecolor{currentstroke}{rgb}{0.878431,0.509804,0.078431}%
\pgfsetstrokecolor{currentstroke}%
\pgfsetdash{}{0pt}%
\pgfpathmoveto{\pgfqpoint{0.650000in}{1.282889in}}%
\pgfpathlineto{\pgfqpoint{0.665704in}{1.284318in}}%
\pgfpathlineto{\pgfqpoint{0.674407in}{1.283889in}}%
\pgfpathlineto{\pgfqpoint{0.688496in}{1.282956in}}%
\pgfpathlineto{\pgfqpoint{0.695541in}{1.283795in}}%
\pgfpathlineto{\pgfqpoint{0.703000in}{1.285708in}}%
\pgfpathlineto{\pgfqpoint{0.712531in}{1.289226in}}%
\pgfpathlineto{\pgfqpoint{0.727449in}{1.296028in}}%
\pgfpathlineto{\pgfqpoint{0.804112in}{1.332546in}}%
\pgfpathlineto{\pgfqpoint{0.814058in}{1.337038in}}%
\pgfpathlineto{\pgfqpoint{0.851353in}{1.353431in}}%
\pgfpathlineto{\pgfqpoint{0.862542in}{1.358203in}}%
\pgfpathlineto{\pgfqpoint{0.919314in}{1.381707in}}%
\pgfpathlineto{\pgfqpoint{0.978987in}{1.405435in}}%
\pgfpathlineto{\pgfqpoint{1.045290in}{1.430173in}}%
\pgfpathlineto{\pgfqpoint{1.073469in}{1.439641in}}%
\pgfpathlineto{\pgfqpoint{1.091288in}{1.444761in}}%
\pgfpathlineto{\pgfqpoint{1.096261in}{1.445902in}}%
\pgfpathlineto{\pgfqpoint{1.108278in}{1.448050in}}%
\pgfpathlineto{\pgfqpoint{1.117809in}{1.448773in}}%
\pgfpathlineto{\pgfqpoint{1.125683in}{1.448428in}}%
\pgfpathlineto{\pgfqpoint{1.143502in}{1.447361in}}%
\pgfpathlineto{\pgfqpoint{1.158006in}{1.448810in}}%
\pgfpathlineto{\pgfqpoint{1.179554in}{1.453736in}}%
\pgfpathlineto{\pgfqpoint{1.230110in}{1.469180in}}%
\pgfpathlineto{\pgfqpoint{1.244200in}{1.473839in}}%
\pgfpathlineto{\pgfqpoint{1.245029in}{1.473798in}}%
\pgfpathlineto{\pgfqpoint{1.266163in}{1.480841in}}%
\pgfpathlineto{\pgfqpoint{1.266992in}{1.480806in}}%
\pgfpathlineto{\pgfqpoint{1.298900in}{1.491538in}}%
\pgfpathlineto{\pgfqpoint{1.299729in}{1.491508in}}%
\pgfpathlineto{\pgfqpoint{1.318791in}{1.498065in}}%
\pgfpathlineto{\pgfqpoint{1.319620in}{1.498029in}}%
\pgfpathlineto{\pgfqpoint{1.339511in}{1.504766in}}%
\pgfpathlineto{\pgfqpoint{1.340340in}{1.504724in}}%
\pgfpathlineto{\pgfqpoint{1.360645in}{1.511500in}}%
\pgfpathlineto{\pgfqpoint{1.361474in}{1.511455in}}%
\pgfpathlineto{\pgfqpoint{1.385923in}{1.519502in}}%
\pgfpathlineto{\pgfqpoint{1.386752in}{1.519451in}}%
\pgfpathlineto{\pgfqpoint{1.409958in}{1.526965in}}%
\pgfpathlineto{\pgfqpoint{1.410787in}{1.526908in}}%
\pgfpathlineto{\pgfqpoint{1.434407in}{1.534374in}}%
\pgfpathlineto{\pgfqpoint{1.435236in}{1.534314in}}%
\pgfpathlineto{\pgfqpoint{1.459271in}{1.541949in}}%
\pgfpathlineto{\pgfqpoint{1.460100in}{1.541888in}}%
\pgfpathlineto{\pgfqpoint{1.480405in}{1.548255in}}%
\pgfpathlineto{\pgfqpoint{1.481234in}{1.548195in}}%
\pgfpathlineto{\pgfqpoint{1.504026in}{1.555305in}}%
\pgfpathlineto{\pgfqpoint{1.504855in}{1.555234in}}%
\pgfpathlineto{\pgfqpoint{1.530547in}{1.562906in}}%
\pgfpathlineto{\pgfqpoint{1.531376in}{1.562825in}}%
\pgfpathlineto{\pgfqpoint{1.554168in}{1.569476in}}%
\pgfpathlineto{\pgfqpoint{1.554997in}{1.569400in}}%
\pgfpathlineto{\pgfqpoint{1.589806in}{1.579942in}}%
\pgfpathlineto{\pgfqpoint{1.590635in}{1.579879in}}%
\pgfpathlineto{\pgfqpoint{1.610111in}{1.586020in}}%
\pgfpathlineto{\pgfqpoint{1.610940in}{1.585954in}}%
\pgfpathlineto{\pgfqpoint{1.631660in}{1.592236in}}%
\pgfpathlineto{\pgfqpoint{1.644506in}{1.595636in}}%
\pgfpathlineto{\pgfqpoint{1.645335in}{1.595523in}}%
\pgfpathlineto{\pgfqpoint{1.668541in}{1.601516in}}%
\pgfpathlineto{\pgfqpoint{1.789130in}{1.633970in}}%
\pgfpathlineto{\pgfqpoint{1.789959in}{1.633859in}}%
\pgfpathlineto{\pgfqpoint{1.824354in}{1.643119in}}%
\pgfpathlineto{\pgfqpoint{1.825183in}{1.643025in}}%
\pgfpathlineto{\pgfqpoint{1.925052in}{1.672416in}}%
\pgfpathlineto{\pgfqpoint{1.925881in}{1.672293in}}%
\pgfpathlineto{\pgfqpoint{1.955303in}{1.680228in}}%
\pgfpathlineto{\pgfqpoint{1.956132in}{1.680105in}}%
\pgfpathlineto{\pgfqpoint{1.963176in}{1.681660in}}%
\pgfpathlineto{\pgfqpoint{1.966906in}{1.680956in}}%
\pgfpathlineto{\pgfqpoint{1.968978in}{1.681200in}}%
\pgfpathlineto{\pgfqpoint{1.972707in}{1.683673in}}%
\pgfpathlineto{\pgfqpoint{1.976023in}{1.685589in}}%
\pgfpathlineto{\pgfqpoint{1.981410in}{1.687409in}}%
\pgfpathlineto{\pgfqpoint{1.982239in}{1.687318in}}%
\pgfpathlineto{\pgfqpoint{2.005859in}{1.694603in}}%
\pgfpathlineto{\pgfqpoint{2.006688in}{1.694495in}}%
\pgfpathlineto{\pgfqpoint{2.022435in}{1.698862in}}%
\pgfpathlineto{\pgfqpoint{2.044812in}{1.704790in}}%
\pgfpathlineto{\pgfqpoint{2.045641in}{1.704647in}}%
\pgfpathlineto{\pgfqpoint{2.067604in}{1.710435in}}%
\pgfpathlineto{\pgfqpoint{2.068433in}{1.710300in}}%
\pgfpathlineto{\pgfqpoint{2.089982in}{1.716047in}}%
\pgfpathlineto{\pgfqpoint{2.090810in}{1.715898in}}%
\pgfpathlineto{\pgfqpoint{2.109873in}{1.720611in}}%
\pgfpathlineto{\pgfqpoint{2.128520in}{1.724791in}}%
\pgfpathlineto{\pgfqpoint{2.139709in}{1.726913in}}%
\pgfpathlineto{\pgfqpoint{2.142610in}{1.726787in}}%
\pgfpathlineto{\pgfqpoint{2.147583in}{1.725751in}}%
\pgfpathlineto{\pgfqpoint{2.149655in}{1.727156in}}%
\pgfpathlineto{\pgfqpoint{2.153384in}{1.729859in}}%
\pgfpathlineto{\pgfqpoint{2.156285in}{1.730818in}}%
\pgfpathlineto{\pgfqpoint{2.157114in}{1.730636in}}%
\pgfpathlineto{\pgfqpoint{2.203525in}{1.741406in}}%
\pgfpathlineto{\pgfqpoint{2.204354in}{1.741231in}}%
\pgfpathlineto{\pgfqpoint{2.244965in}{1.752259in}}%
\pgfpathlineto{\pgfqpoint{2.260298in}{1.756765in}}%
\pgfpathlineto{\pgfqpoint{2.261126in}{1.756617in}}%
\pgfpathlineto{\pgfqpoint{2.286819in}{1.763961in}}%
\pgfpathlineto{\pgfqpoint{2.305467in}{1.768810in}}%
\pgfpathlineto{\pgfqpoint{2.306296in}{1.768626in}}%
\pgfpathlineto{\pgfqpoint{2.333646in}{1.775749in}}%
\pgfpathlineto{\pgfqpoint{2.334475in}{1.775543in}}%
\pgfpathlineto{\pgfqpoint{2.353537in}{1.780323in}}%
\pgfpathlineto{\pgfqpoint{2.354366in}{1.780093in}}%
\pgfpathlineto{\pgfqpoint{2.382545in}{1.786221in}}%
\pgfpathlineto{\pgfqpoint{2.448020in}{1.801979in}}%
\pgfpathlineto{\pgfqpoint{2.463767in}{1.806645in}}%
\pgfpathlineto{\pgfqpoint{2.464595in}{1.806467in}}%
\pgfpathlineto{\pgfqpoint{2.499404in}{1.818962in}}%
\pgfpathlineto{\pgfqpoint{2.515152in}{1.826561in}}%
\pgfpathlineto{\pgfqpoint{2.528412in}{1.835014in}}%
\pgfpathlineto{\pgfqpoint{2.537115in}{1.843021in}}%
\pgfpathlineto{\pgfqpoint{2.545402in}{1.853142in}}%
\pgfpathlineto{\pgfqpoint{2.557834in}{1.870102in}}%
\pgfpathlineto{\pgfqpoint{2.564464in}{1.877504in}}%
\pgfpathlineto{\pgfqpoint{2.565707in}{1.878478in}}%
\pgfpathlineto{\pgfqpoint{2.572336in}{1.884426in}}%
\pgfpathlineto{\pgfqpoint{2.573580in}{1.885174in}}%
\pgfpathlineto{\pgfqpoint{2.581038in}{1.890582in}}%
\pgfpathlineto{\pgfqpoint{2.582281in}{1.891160in}}%
\pgfpathlineto{\pgfqpoint{2.591397in}{1.896409in}}%
\pgfpathlineto{\pgfqpoint{2.601342in}{1.901404in}}%
\pgfpathlineto{\pgfqpoint{2.602585in}{1.901743in}}%
\pgfpathlineto{\pgfqpoint{2.614187in}{1.906801in}}%
\pgfpathlineto{\pgfqpoint{2.615430in}{1.907047in}}%
\pgfpathlineto{\pgfqpoint{2.629933in}{1.912217in}}%
\pgfpathlineto{\pgfqpoint{2.646093in}{1.917266in}}%
\pgfpathlineto{\pgfqpoint{2.661839in}{1.921509in}}%
\pgfpathlineto{\pgfqpoint{2.678413in}{1.925302in}}%
\pgfpathlineto{\pgfqpoint{2.697474in}{1.929396in}}%
\pgfpathlineto{\pgfqpoint{2.698303in}{1.929252in}}%
\pgfpathlineto{\pgfqpoint{2.730208in}{1.936008in}}%
\pgfpathlineto{\pgfqpoint{2.731037in}{1.935826in}}%
\pgfpathlineto{\pgfqpoint{2.756728in}{1.940938in}}%
\pgfpathlineto{\pgfqpoint{2.771230in}{1.943249in}}%
\pgfpathlineto{\pgfqpoint{2.819711in}{1.952444in}}%
\pgfpathlineto{\pgfqpoint{2.838771in}{1.956213in}}%
\pgfpathlineto{\pgfqpoint{2.839600in}{1.955884in}}%
\pgfpathlineto{\pgfqpoint{2.860318in}{1.960349in}}%
\pgfpathlineto{\pgfqpoint{2.861147in}{1.959981in}}%
\pgfpathlineto{\pgfqpoint{2.875235in}{1.963022in}}%
\pgfpathlineto{\pgfqpoint{2.876064in}{1.962611in}}%
\pgfpathlineto{\pgfqpoint{2.889738in}{1.965530in}}%
\pgfpathlineto{\pgfqpoint{2.890566in}{1.965079in}}%
\pgfpathlineto{\pgfqpoint{2.903412in}{1.967679in}}%
\pgfpathlineto{\pgfqpoint{2.904240in}{1.967189in}}%
\pgfpathlineto{\pgfqpoint{2.950649in}{1.977123in}}%
\pgfpathlineto{\pgfqpoint{2.951478in}{1.976580in}}%
\pgfpathlineto{\pgfqpoint{2.970124in}{1.981235in}}%
\pgfpathlineto{\pgfqpoint{2.970953in}{1.980640in}}%
\pgfpathlineto{\pgfqpoint{2.979240in}{1.983363in}}%
\pgfpathlineto{\pgfqpoint{2.980069in}{1.982737in}}%
\pgfpathlineto{\pgfqpoint{2.988356in}{1.985737in}}%
\pgfpathlineto{\pgfqpoint{2.989185in}{1.985073in}}%
\pgfpathlineto{\pgfqpoint{2.998301in}{1.988585in}}%
\pgfpathlineto{\pgfqpoint{2.999129in}{1.987893in}}%
\pgfpathlineto{\pgfqpoint{3.008245in}{1.991929in}}%
\pgfpathlineto{\pgfqpoint{3.009074in}{1.991203in}}%
\pgfpathlineto{\pgfqpoint{3.018190in}{1.995543in}}%
\pgfpathlineto{\pgfqpoint{3.018604in}{1.994563in}}%
\pgfpathlineto{\pgfqpoint{3.019018in}{1.994772in}}%
\pgfpathlineto{\pgfqpoint{3.026849in}{1.999008in}}%
\pgfpathlineto{\pgfqpoint{3.026849in}{1.999008in}}%
\pgfusepath{stroke}%
\end{pgfscope}%
\begin{pgfscope}%
\pgfpathrectangle{\pgfqpoint{0.650000in}{0.420000in}}{\pgfqpoint{2.388110in}{1.994488in}}%
\pgfusepath{clip}%
\pgfsetrectcap%
\pgfsetroundjoin%
\pgfsetlinewidth{1.003750pt}%
\definecolor{currentstroke}{rgb}{0.501961,0.450980,0.674510}%
\pgfsetstrokecolor{currentstroke}%
\pgfsetdash{}{0pt}%
\pgfpathmoveto{\pgfqpoint{0.650000in}{1.282890in}}%
\pgfpathlineto{\pgfqpoint{0.669110in}{1.284296in}}%
\pgfpathlineto{\pgfqpoint{0.688262in}{1.282828in}}%
\pgfpathlineto{\pgfqpoint{0.696589in}{1.283922in}}%
\pgfpathlineto{\pgfqpoint{0.704500in}{1.286142in}}%
\pgfpathlineto{\pgfqpoint{0.714909in}{1.290213in}}%
\pgfpathlineto{\pgfqpoint{0.734062in}{1.299115in}}%
\pgfpathlineto{\pgfqpoint{0.795683in}{1.328357in}}%
\pgfpathlineto{\pgfqpoint{0.836069in}{1.346406in}}%
\pgfpathlineto{\pgfqpoint{0.883118in}{1.366344in}}%
\pgfpathlineto{\pgfqpoint{0.935579in}{1.387775in}}%
\pgfpathlineto{\pgfqpoint{0.949319in}{1.393270in}}%
\pgfpathlineto{\pgfqpoint{1.021765in}{1.421858in}}%
\pgfpathlineto{\pgfqpoint{1.045081in}{1.430009in}}%
\pgfpathlineto{\pgfqpoint{1.056739in}{1.432885in}}%
\pgfpathlineto{\pgfqpoint{1.066732in}{1.434238in}}%
\pgfpathlineto{\pgfqpoint{1.075475in}{1.434300in}}%
\pgfpathlineto{\pgfqpoint{1.083803in}{1.433240in}}%
\pgfpathlineto{\pgfqpoint{1.092130in}{1.431038in}}%
\pgfpathlineto{\pgfqpoint{1.100873in}{1.427561in}}%
\pgfpathlineto{\pgfqpoint{1.110866in}{1.422370in}}%
\pgfpathlineto{\pgfqpoint{1.123773in}{1.414371in}}%
\pgfpathlineto{\pgfqpoint{1.138762in}{1.403958in}}%
\pgfpathlineto{\pgfqpoint{1.165409in}{1.386284in}}%
\pgfpathlineto{\pgfqpoint{1.183312in}{1.376602in}}%
\pgfpathlineto{\pgfqpoint{1.198717in}{1.370491in}}%
\pgfpathlineto{\pgfqpoint{1.222033in}{1.363688in}}%
\pgfpathlineto{\pgfqpoint{1.248264in}{1.359361in}}%
\pgfpathlineto{\pgfqpoint{1.284904in}{1.356585in}}%
\pgfpathlineto{\pgfqpoint{1.297394in}{1.356313in}}%
\pgfpathlineto{\pgfqpoint{1.298643in}{1.355992in}}%
\pgfpathlineto{\pgfqpoint{1.336116in}{1.357156in}}%
\pgfpathlineto{\pgfqpoint{1.374004in}{1.361767in}}%
\pgfpathlineto{\pgfqpoint{1.423967in}{1.372722in}}%
\pgfpathlineto{\pgfqpoint{1.438956in}{1.376713in}}%
\pgfpathlineto{\pgfqpoint{1.439789in}{1.376625in}}%
\pgfpathlineto{\pgfqpoint{1.462273in}{1.382884in}}%
\pgfpathlineto{\pgfqpoint{1.463105in}{1.382801in}}%
\pgfpathlineto{\pgfqpoint{1.486422in}{1.389363in}}%
\pgfpathlineto{\pgfqpoint{1.487254in}{1.389293in}}%
\pgfpathlineto{\pgfqpoint{1.519730in}{1.398435in}}%
\pgfpathlineto{\pgfqpoint{1.520563in}{1.398364in}}%
\pgfpathlineto{\pgfqpoint{1.541381in}{1.404075in}}%
\pgfpathlineto{\pgfqpoint{1.542214in}{1.403994in}}%
\pgfpathlineto{\pgfqpoint{1.563448in}{1.409465in}}%
\pgfpathlineto{\pgfqpoint{1.564281in}{1.409365in}}%
\pgfpathlineto{\pgfqpoint{1.585931in}{1.414412in}}%
\pgfpathlineto{\pgfqpoint{1.599671in}{1.417180in}}%
\pgfpathlineto{\pgfqpoint{1.600920in}{1.417130in}}%
\pgfpathlineto{\pgfqpoint{1.624236in}{1.421495in}}%
\pgfpathlineto{\pgfqpoint{1.636311in}{1.423496in}}%
\pgfpathlineto{\pgfqpoint{1.637560in}{1.423422in}}%
\pgfpathlineto{\pgfqpoint{1.660460in}{1.427308in}}%
\pgfpathlineto{\pgfqpoint{1.661709in}{1.427214in}}%
\pgfpathlineto{\pgfqpoint{1.681694in}{1.430170in}}%
\pgfpathlineto{\pgfqpoint{1.682943in}{1.430043in}}%
\pgfpathlineto{\pgfqpoint{1.702095in}{1.432234in}}%
\pgfpathlineto{\pgfqpoint{1.703345in}{1.432059in}}%
\pgfpathlineto{\pgfqpoint{1.723746in}{1.433744in}}%
\pgfpathlineto{\pgfqpoint{1.741233in}{1.434788in}}%
\pgfpathlineto{\pgfqpoint{1.742482in}{1.434557in}}%
\pgfpathlineto{\pgfqpoint{1.764549in}{1.435621in}}%
\pgfpathlineto{\pgfqpoint{1.785367in}{1.436397in}}%
\pgfpathlineto{\pgfqpoint{1.786616in}{1.436140in}}%
\pgfpathlineto{\pgfqpoint{1.811598in}{1.437069in}}%
\pgfpathlineto{\pgfqpoint{1.812847in}{1.436816in}}%
\pgfpathlineto{\pgfqpoint{1.846572in}{1.438138in}}%
\pgfpathlineto{\pgfqpoint{1.847821in}{1.437902in}}%
\pgfpathlineto{\pgfqpoint{1.870304in}{1.439166in}}%
\pgfpathlineto{\pgfqpoint{1.871554in}{1.438942in}}%
\pgfpathlineto{\pgfqpoint{1.891122in}{1.440019in}}%
\pgfpathlineto{\pgfqpoint{1.892372in}{1.439784in}}%
\pgfpathlineto{\pgfqpoint{1.913606in}{1.440806in}}%
\pgfpathlineto{\pgfqpoint{1.914855in}{1.440558in}}%
\pgfpathlineto{\pgfqpoint{1.938587in}{1.441506in}}%
\pgfpathlineto{\pgfqpoint{1.939836in}{1.441248in}}%
\pgfpathlineto{\pgfqpoint{1.962320in}{1.442060in}}%
\pgfpathlineto{\pgfqpoint{1.963569in}{1.441807in}}%
\pgfpathlineto{\pgfqpoint{1.969398in}{1.441637in}}%
\pgfpathlineto{\pgfqpoint{1.973145in}{1.440006in}}%
\pgfpathlineto{\pgfqpoint{1.975227in}{1.439733in}}%
\pgfpathlineto{\pgfqpoint{1.978558in}{1.441088in}}%
\pgfpathlineto{\pgfqpoint{1.981472in}{1.442251in}}%
\pgfpathlineto{\pgfqpoint{1.986885in}{1.442748in}}%
\pgfpathlineto{\pgfqpoint{2.005621in}{1.443480in}}%
\pgfpathlineto{\pgfqpoint{2.006870in}{1.443234in}}%
\pgfpathlineto{\pgfqpoint{2.080149in}{1.448325in}}%
\pgfpathlineto{\pgfqpoint{2.098053in}{1.450452in}}%
\pgfpathlineto{\pgfqpoint{2.099302in}{1.450310in}}%
\pgfpathlineto{\pgfqpoint{2.149265in}{1.457382in}}%
\pgfpathlineto{\pgfqpoint{2.150514in}{1.456650in}}%
\pgfpathlineto{\pgfqpoint{2.154261in}{1.455645in}}%
\pgfpathlineto{\pgfqpoint{2.156343in}{1.456796in}}%
\pgfpathlineto{\pgfqpoint{2.160923in}{1.459986in}}%
\pgfpathlineto{\pgfqpoint{2.161756in}{1.459956in}}%
\pgfpathlineto{\pgfqpoint{2.169667in}{1.461760in}}%
\pgfpathlineto{\pgfqpoint{2.184239in}{1.464589in}}%
\pgfpathlineto{\pgfqpoint{2.185071in}{1.464447in}}%
\pgfpathlineto{\pgfqpoint{2.210053in}{1.469502in}}%
\pgfpathlineto{\pgfqpoint{2.210886in}{1.469354in}}%
\pgfpathlineto{\pgfqpoint{2.240031in}{1.475037in}}%
\pgfpathlineto{\pgfqpoint{2.241280in}{1.474977in}}%
\pgfpathlineto{\pgfqpoint{2.291659in}{1.485747in}}%
\pgfpathlineto{\pgfqpoint{2.303318in}{1.488590in}}%
\pgfpathlineto{\pgfqpoint{2.304150in}{1.488501in}}%
\pgfpathlineto{\pgfqpoint{2.353697in}{1.501638in}}%
\pgfpathlineto{\pgfqpoint{2.354530in}{1.501567in}}%
\pgfpathlineto{\pgfqpoint{2.419482in}{1.519718in}}%
\pgfpathlineto{\pgfqpoint{2.420731in}{1.519798in}}%
\pgfpathlineto{\pgfqpoint{2.437802in}{1.523934in}}%
\pgfpathlineto{\pgfqpoint{2.449044in}{1.526013in}}%
\pgfpathlineto{\pgfqpoint{2.460285in}{1.527419in}}%
\pgfpathlineto{\pgfqpoint{2.469029in}{1.527874in}}%
\pgfpathlineto{\pgfqpoint{2.479022in}{1.527704in}}%
\pgfpathlineto{\pgfqpoint{2.487765in}{1.526925in}}%
\pgfpathlineto{\pgfqpoint{2.496508in}{1.525592in}}%
\pgfpathlineto{\pgfqpoint{2.506501in}{1.523370in}}%
\pgfpathlineto{\pgfqpoint{2.517743in}{1.520080in}}%
\pgfpathlineto{\pgfqpoint{2.529401in}{1.515777in}}%
\pgfpathlineto{\pgfqpoint{2.541475in}{1.510576in}}%
\pgfpathlineto{\pgfqpoint{2.561044in}{1.501579in}}%
\pgfpathlineto{\pgfqpoint{2.563126in}{1.500547in}}%
\pgfpathlineto{\pgfqpoint{2.586023in}{1.491757in}}%
\pgfpathlineto{\pgfqpoint{2.599346in}{1.488189in}}%
\pgfpathlineto{\pgfqpoint{2.613085in}{1.485959in}}%
\pgfpathlineto{\pgfqpoint{2.624325in}{1.485697in}}%
\pgfpathlineto{\pgfqpoint{2.637232in}{1.486527in}}%
\pgfpathlineto{\pgfqpoint{2.650970in}{1.489004in}}%
\pgfpathlineto{\pgfqpoint{2.664293in}{1.492678in}}%
\pgfpathlineto{\pgfqpoint{2.681362in}{1.498785in}}%
\pgfpathlineto{\pgfqpoint{2.701762in}{1.507300in}}%
\pgfpathlineto{\pgfqpoint{2.752137in}{1.528848in}}%
\pgfpathlineto{\pgfqpoint{2.777949in}{1.538652in}}%
\pgfpathlineto{\pgfqpoint{2.848725in}{1.564554in}}%
\pgfpathlineto{\pgfqpoint{2.876202in}{1.574637in}}%
\pgfpathlineto{\pgfqpoint{2.910341in}{1.586742in}}%
\pgfpathlineto{\pgfqpoint{2.968210in}{1.607177in}}%
\pgfpathlineto{\pgfqpoint{2.984446in}{1.614773in}}%
\pgfpathlineto{\pgfqpoint{2.992357in}{1.618491in}}%
\pgfpathlineto{\pgfqpoint{2.999434in}{1.620377in}}%
\pgfpathlineto{\pgfqpoint{3.010259in}{1.623570in}}%
\pgfpathlineto{\pgfqpoint{3.026912in}{1.630242in}}%
\pgfpathlineto{\pgfqpoint{3.038110in}{1.635740in}}%
\pgfpathlineto{\pgfqpoint{3.038110in}{1.635740in}}%
\pgfusepath{stroke}%
\end{pgfscope}%
\begin{pgfscope}%
\pgfpathrectangle{\pgfqpoint{0.650000in}{0.420000in}}{\pgfqpoint{2.388110in}{1.994488in}}%
\pgfusepath{clip}%
\pgfsetrectcap%
\pgfsetroundjoin%
\pgfsetlinewidth{0.803000pt}%
\definecolor{currentstroke}{rgb}{0.500000,0.500000,0.500000}%
\pgfsetstrokecolor{currentstroke}%
\pgfsetdash{}{0pt}%
\pgfpathmoveto{\pgfqpoint{0.650000in}{1.417244in}}%
\pgfpathlineto{\pgfqpoint{3.038110in}{1.417244in}}%
\pgfusepath{stroke}%
\end{pgfscope}%
\begin{pgfscope}%
\pgfsetrectcap%
\pgfsetmiterjoin%
\pgfsetlinewidth{0.803000pt}%
\definecolor{currentstroke}{rgb}{0.000000,0.000000,0.000000}%
\pgfsetstrokecolor{currentstroke}%
\pgfsetdash{}{0pt}%
\pgfpathmoveto{\pgfqpoint{0.650000in}{0.420000in}}%
\pgfpathlineto{\pgfqpoint{0.650000in}{2.414488in}}%
\pgfusepath{stroke}%
\end{pgfscope}%
\begin{pgfscope}%
\pgfsetrectcap%
\pgfsetmiterjoin%
\pgfsetlinewidth{0.803000pt}%
\definecolor{currentstroke}{rgb}{0.000000,0.000000,0.000000}%
\pgfsetstrokecolor{currentstroke}%
\pgfsetdash{}{0pt}%
\pgfpathmoveto{\pgfqpoint{3.038110in}{0.420000in}}%
\pgfpathlineto{\pgfqpoint{3.038110in}{2.414488in}}%
\pgfusepath{stroke}%
\end{pgfscope}%
\begin{pgfscope}%
\pgfsetrectcap%
\pgfsetmiterjoin%
\pgfsetlinewidth{0.803000pt}%
\definecolor{currentstroke}{rgb}{0.000000,0.000000,0.000000}%
\pgfsetstrokecolor{currentstroke}%
\pgfsetdash{}{0pt}%
\pgfpathmoveto{\pgfqpoint{0.650000in}{0.420000in}}%
\pgfpathlineto{\pgfqpoint{3.038110in}{0.420000in}}%
\pgfusepath{stroke}%
\end{pgfscope}%
\begin{pgfscope}%
\pgfsetrectcap%
\pgfsetmiterjoin%
\pgfsetlinewidth{0.803000pt}%
\definecolor{currentstroke}{rgb}{0.000000,0.000000,0.000000}%
\pgfsetstrokecolor{currentstroke}%
\pgfsetdash{}{0pt}%
\pgfpathmoveto{\pgfqpoint{0.650000in}{2.414488in}}%
\pgfpathlineto{\pgfqpoint{3.038110in}{2.414488in}}%
\pgfusepath{stroke}%
\end{pgfscope}%
\end{pgfpicture}%
\makeatother%
\endgroup%

    \label{fig:R500deg18R100deg90_Ue}
\end{subfigure}
\caption{Budget of $\langle u_i' u_j' \rangle$ (component: uu) at $Re_\theta \approx 700$.
Linestyle denotes budget term:
Production \lswatch[solid]{black},
Convection \lswatch[dash pattern=on 5pt off 1pt]{black},
Dissipation \lswatch[dash pattern=on 3pt off 1pt on 1pt off 1pt]{black},
Vel-Press.\ \lswatch[dash pattern=on 1pt off 1pt]{black},
Turb.\ Trans.\ \lswatch[dash pattern=on 5pt off 5pt]{black},
Visc.\ Diff.\ \lswatch[dash pattern=on 3pt off 5pt on 1pt off 5pt]{black},
Residual \lswatch[dash pattern=on 1pt off 5pt]{black}.}
\label{fig:R500deg18R100deg90}
\end{figure}



The friction coefficient is depicted in \autoref{fig:R500deg18R100deg90_cf}.
For both cases, differences between the base and optimised configurations are observed.
All configurations exhibit a peak near the inflow, results from the tripping of the boundary layer \citep{schlatter_turbulent_2012}.
Then, the friction coefficient decreases monotonically in all simulations.

For the case R500deg18, a slight deviation from this decreasing trend occurs at approximately $s/\delta_0 \approx 75$, corresponding to the onset of the curved section.
In the optimised configuration, this deviation is absent, and the decreasing trend continues smoothly, resulting in a lower friction coefficient than in the base configuration. 
For a concave curvature, \citet{appelbaum_systematic_2025} reported an opposite effect, indicating that the present observations are consistent with their findings.

In contrast, no comparable difference between the base and optimised configurations is observed for case R100deg90. 
Both configurations exhibit a similar deviation at the beginning of the curved segment and subsequently follow nearly identical trends until the end of the curvature.


Finally, the quantities directly related to the optimisation are analysed.
The distribution of the Clauser parameter is shown in \autoref{fig:R500deg18R100deg90_beta}.
We first focus on the case R500deg18.
Similar to case R100deg30 discussed in \autoref{subsec:R100deg30}, both configurations  exhibit distinct peaks at the entrance and exit of the curved  segment. However, these peaks are less pronounced than those observed for R100deg30. 
Within the curved segment, the base configuration shows a smooth trend for $\beta$, whereas the optimised configuration exhibits small local fluctuations.
Compared to case R100deg30, these fluctuations are noticeably weaker.

For case R100deg90, the peaks at the transitions into and out of the curved section are substantially larger than those observed for R500deg18.  
Within the curved segment, the base configuration again follows the same smooth trend as in the other cases.
The optimised configuration reduces the global trend of $\beta$ towards zero. 
On a local level, however, pronounced local fluctuations remain. 
Their origin is discussed in conjunction with the pressure and free-stream velocity distributions below.

The wall-pressure distributions are visualised in \autoref{fig:R500deg18R100deg90_pres}.
For case R500deg18, the qualitative behaviour closely resembles that of R100deg30. 
As discussed previously, the base configuration exhibits an approximately linear pressure decrease over the entire domain length, interrupted only by local deviations at the beginning and end of the curved section. 
Compared to the case R100deg30, the overall slope is reduced.
In the optimised configuration, the pressure is shifted towards an approximately constant level close to zero throughout the curved section, while only minor local fluctuations remain. 

In the case R100deg90, the base configuration follows the same qualitative behaviour as the other cases. 
The slope of the pressure  decreases more rapidly over the curved section, which is consistent with the stronger wall curvature \citep{so_experiment_1973}.
The corresponding free-stream velocity, depicted in \autoref{fig:R500deg18R100deg90_Ue} shows analogous behaviour.
Again, both base configurations show smooth positive trends.
In case R500deg18, the optimised configuration results in an almost constant free-stream velocity of close to unity.
In contrast, the optimised configuration of R100deg90 still exhibits an increasing free-stream velocity, with an even steeper overall trend than the corresponding base configuration.
Additionally, pronounced local fluctuations are present, consistent with the behaviour observed in the distributions of $\beta$ and wall pressure.




The analysis of the pressure-related quantities for the cases R500deg18 and R100deg90 reveals substantial differences in the performance of the optimisation algorithm. 
For case R500deg18, the optimisation achieves results comparable to those obtained for R100deg30.
On a global level, the pressure field is successfully modified to establish an approximately constant pressure distribution throughout the curved section. 
Consequently, the free-stream velocity is balanced on a nearly constant level.
Although local fluctuations are still present in all analysed quantities, they are less pronounced than in case R100deg30.
This behaviour is consistent with the weaker curvature, which induces a smaller pressure gradient and therefore requires less extensive modifications of the pressure field.
In addition, the lower curvature may increase the effective influence of individual design points, leading to a smoother distribution of the Clauser parameter.

In contrast, the optimisation performs considerably less effectively for case R100deg90.
While the base configuration exhibits trends similar to those observed in the other cases, the optimised solution is dominated by strong local fluctuations, preventing the establishment of the desired pressure distribution.
One possible explanation is that the effective influence of each \benennung{control point} decreases with increasing curvature.
If this hypothesis is correct, geometries with stronger curvature require a larger number of design points to achieve the same level of pressure control. 
In addition, the stronger curvature provokes a larger pressure gradient demanding more substantial modifications of the pressure field.
This effect is particularly pronounced for R100deg90.
Although R100deg90 and R100deg30 share the same radius of curvature, the longer curved section of R100deg90 results in a larger cumulative pressure-gradient effect.


Overall, the results indicate that the proposed optimisation strategy is effective for configurations with mild convex curvature, where it successfully drives the pressure distribution towards \acrshort{zpg} conditions using a limited number of design points.
For geometries with stronger or more extended curvature, however, further development of the optimisation approach is required to achieve comparable performance.

